\documentclass[11pt,openany]{book}

\usepackage{fontspec}
\usepackage{amsmath,amssymb,amsfonts}
\usepackage{amsthm}
\usepackage[margin=1in,headheight=17pt,headsep=14pt]{geometry}
\usepackage{booktabs,longtable,array}
\usepackage{graphicx}
\usepackage{fancyhdr}
\usepackage{titlesec}
\usepackage{microtype}
\usepackage[dvipsnames]{xcolor}
\usepackage[colorlinks=true,linkcolor=NavyBlue,citecolor=NavyBlue,urlcolor=Mahogany]{hyperref}
\makeatletter
\@ifundefined{IfFontExistsTF}{\newcommand{\IfFontExistsTF}[3]{#3}}{}
\makeatother
\IfFontExistsTF{TeX Gyre Pagella}{\setmainfont{TeX Gyre Pagella}}{%
  \IfFontExistsTF{texgyrepagella-regular.otf}{\setmainfont{texgyrepagella-regular.otf}}{}}
\IfFontExistsTF{TeX Gyre Heros}{\setsansfont{TeX Gyre Heros}}{}
\IfFontExistsTF{Latin Modern Mono}{\setmonofont{Latin Modern Mono}[Scale=0.92]}{%
  \IfFontExistsTF{lmmono10-regular.otf}{\setmonofont{lmmono10-regular.otf}[Scale=0.92]}{}}
% --- amsthm setup -----------------------------------------------------------
% \theoremstyle{plain}-family style for theorem-like environments and
% \theoremstyle{definition}-family style for definition-like environments.
% The bodies are set upright (a `plain` italic body would italicise the
% corpus's extensive \text{...} proof justifications); head formatting and
% numbering follow the plain / definition styles respectively.
\newtheoremstyle{canonplain}%
  {7pt plus 2pt minus 1pt}{7pt plus 2pt minus 1pt}%
  {\normalfont}{0pt}{\bfseries}{.}%
  {0.6em plus 0.2em}%
  {\thmname{#1}\thmnumber{ #2}\thmnote{\normalfont\ (#3)}}
\newtheoremstyle{canondef}%
  {7pt plus 2pt minus 1pt}{7pt plus 2pt minus 1pt}%
  {\normalfont}{0pt}{\bfseries}{.}%
  {0.6em plus 0.2em}%
  {\thmname{#1}\thmnumber{ #2}\thmnote{\normalfont\ (#3)}}
\theoremstyle{canonplain}
\newtheorem{theorem}{Theorem}[chapter]
\newtheorem{corollary}{Corollary}[chapter]
\newtheorem{axiom}{Axiom}[chapter]
\newtheorem{convention}{Convention}[chapter]
\theoremstyle{canondef}
\newtheorem{definition}{Definition}[chapter]
\newenvironment{movedblock}
  {\par\medskip\begingroup\small\itshape\leftskip2.2em\rightskip2.2em}
  {\par\endgroup\medskip}
% run-in heading for theorem-class blocks that contain tables or verbatim
% material (longtable/verbatim cannot live inside amsthm environments); the
% printed form matches the environment heads and the counter is shared.
\newcommand{\pseudohead}[3]{\par\medskip\noindent\refstepcounter{#1}%
  \textbf{\csname the#1\endcsname\hspace{0.6em}(#2\hspace{0.35em}---\hspace{0.35em}#3).}%
  \par\nobreak\smallskip}
% ----------------------------------------------------------------------------
\pagestyle{fancy}
\fancyhf{}
\fancyhead[LE]{\small\thepage}
\fancyhead[RO]{\small\thepage}
\fancyhead[RE]{\nouppercase{\small\leftmark}}
\fancyhead[LO]{\nouppercase{\small\rightmark}}
\renewcommand{\headrulewidth}{0.4pt}
\setlength{\parskip}{2pt plus 1pt}
\renewcommand{\arraystretch}{1.08}
\fancypagestyle{plain}{\fancyhf{}\fancyfoot[C]{\thepage}}
\titleformat{\chapter}[hang]{\normalfont\Large\bfseries}{\thechapter}{1em}{}
\titleformat{\section}{\normalfont\large\bfseries}{\thesection}{1em}{}
\titleformat{\subsection}{\normalfont\normalsize\bfseries}{\thesubsection}{1em}{}
\titlespacing{\subsection}{0pt}{8pt plus 2pt minus 1pt}{3pt}
\setcounter{tocdepth}{1}
\setcounter{secnumdepth}{3}
\allowdisplaybreaks
\emergencystretch=3em
\hbadness=10000
\vbadness=10000
\hfuzz=2pt
\begin{document}


\begin{titlepage}
\centering
\vspace*{2.0cm}
{\Huge $\langle qu(b(i(t|t))ick)\rangle$\par}
\vspace{0.4cm}
{\LARGE The Four-Era Canon\par}
\vspace{0.3cm}
{\large A Machine-Verified, Discrete-Ledger Framework for Foundational Physics\par}
\vspace{0.5cm}
{\large Volume I --- The Four-Era Canon\par}
\vspace{0.8cm}
{\large Qubit\_era $\longrightarrow$ Bosonic\_era $\longrightarrow$ Fermion\_era $\longrightarrow$ Baryomesonic\_era\par}
\vspace{1.3cm}
\begin{tabular}{rl}
\toprule
axioms & A1--A9 \\
conventions & C1--C13 \\
definitions & D0.1--D4.3, Def-EM.20--22, Def-H.1--43, Def-BM.1--20 \\
theorem and data-cell tags & 481 \\
registry tags (total) & 623 \\
dependency edges (Req graph) & 8014 \\
free dimensionful scales & $1$ ($\mathfrak h$, P1) \\
verification suite & 80 checkers, exact arithmetic, no floating point \\
\bottomrule
\end{tabular}
\vfill
{\large Edition of 2026-10-09\par}
\vspace{0.2cm}
{\normalsize Machine-verified: \texttt{verification/run\_all.py} --- suite state 80/80 green.\par}
\vspace{0.3cm}
{\small The mathematics is transplanted unmodified from the canon files preserved in \texttt{working\_markdowns/}; this volume is the sanitized presentation layer.\par}
\end{titlepage}

\frontmatter

\chapter*{\texorpdfstring{Preface — Guide to this Volume}{Preface - Guide to this Volume}}
\addcontentsline{toc}{chapter}{Preface --- Guide to this Volume}

This volume is the presentation layer of the $\langle qu(b(i(t|t))ick)\rangle$
corpus (read \textit{qubitick}): a formal corpus
deriving, from nine axioms and one free dimensionful scale $\mathfrak h$, the
qubit, bosonic, fermionic and baryomesonic layers, with every theorem carrying
its requirement set (\textbf{Req}), its symbolic proof (\textbf{Pf}) and its
physical reading (\textbf{Phys}), and every identity machine-verified in exact
arithmetic. The corpus's own description of its subject stands: qubits
$\longrightarrow$ bosonic gauge/geometry $\longrightarrow$ fermionic bound
chains $\longrightarrow$ baryomesonic colour-sealed composites; pure symbolic
derivation; no floating point anywhere; no external physical constant enters
any derivation.

\medskip
\noindent\textbf{The volume.}\quad
Chapters 1--4 are the four eras (Part 1 · Qubit, Part 2 · Bosonic,
Part 3 · Fermion, Part 4 · Baryomesonic); Appendix A is the Census annex (the
25 computed data cells, cited row-wise by the parts); the Index of Corpus Tags
maps every canonical tag to its location in this volume or in the Addendum.

\noindent\textbf{Companion documents.}\quad
\textit{Addendum --- Structural Constraints, Refutations, and No-Go Theorems}
(\texttt{docs/Addendum\_Refutations.pdf}) is the volume's Appendix B: the
adjudication record, the refuted hypotheses, the no-go theorems and the trap
register, moved out of the main text; where such a block stood inline, a stub
cross-reference remains. \textit{Calibration Policy}
(\texttt{docs/Calibration\_Policy.pdf}) lists every calibration slot, its
locus, its consumers, and its status under P5. The repository
\texttt{README.md} describes the package; \texttt{working\_markdowns/}
preserves the raw working corpus byte-identically;
\texttt{verification/} contains the checker suite.

\noindent\textbf{Canonical tags and numbering.}\quad
Every theorem environment carries both its presentation number (Theorem 1.1,
Definition 3.12, \dots) and its canonical corpus tag (CT0.1, FT.54, BM.137,
\dots) in its head. \textbf{All cross-citations in Req lines, proofs and scope
clauses use the canonical tags}, which are unique across the corpus and are
resolved by the Index of Corpus Tags; tags moved to the Addendum keep their
identity there (Appendix B.$x$.$y$ entries are labelled with the tag).

\noindent\textbf{Citations and records.}\quad
This volume presents the corpus as a finished mathematical model: every
citation resolves to a canonical tag, a Census cell, a [math] class, or a
machine checker. The adjudications, refutations, no-go theorems and trap
register --- the record of how the model was won --- live in the Addendum;
the archival working record (including the historical ledger
\texttt{HistLog.md}) is preserved in \texttt{working\_markdowns/}. Paths of
the form \texttt{checks/verify\_*.py} in Req lines refer to
\texttt{verification/checks/verify\_*.py} in this repository.

\noindent\textbf{Verification.}\quad
The dependency graph over the registry tags is acyclic and is re-derived from
the canon at every suite run (\texttt{verify\_corpus\_graph.py}); the era
structure, document format and every algebraic identity are checked in exact
arithmetic (integer, \texttt{Fraction}, algebraic number fields, prime-log
rings, truncated power series) with no floating point in any assertion.

\tableofcontents
\mainmatter
% Part1_Qubit_era.md — sanitized presentation layer (generated; see source/CONVERSION_LOG.md)
\chapter{\texorpdfstring{$\langle qu(b(i(t|t))ick)\rangle$ — \textbf{Part 1 · Qubit\_era} — Axioms · Conventions · Qubit-Era Definitions · Pair Canon · Ledger Dynamics}{Qubitick - **Part 1 - Qubit _era** - Axioms - Conventions - Qubit-Era Definitions - Pair Canon - Ledger Dynamics}}

\textbf{0.1 Corpus map.} The canon is the era chain

%dsp-level:0
\[
\textbf{Part 1 · Qubit\_era}\ \longrightarrow\ \textbf{Part 2 · Bosonic\_era}\ \longrightarrow\ \textbf{Part 3 · Fermion\_era}\ \longrightarrow\ \textbf{Part 4 · Baryomesonic\_era},
\]

each part citing only its priors, itself (earlier tags), and the Census annex (the computed data cells, cited row-wise by tag). Part 1 (this chapter): §A axioms A1–A9 · §B conventions C1–C13 · §C the definition layer D0.1–D4.3 (the corpus-wide vocabulary root: Tiers 0–1 qubit-era, Tiers 2–4 naming the structures the bosonic era derives) · §D master symbol table · §E calibration policy P1–P5 · §F foundations (CT0.1–CT0.14, CT1.0–CT1.11b, L0–L5). Part 2 (section letters preserve the classical canon's): §C CT1.12–CT2.11 (geometry), §D EW.1–EW.36a (electroweak), §E EM.1–EM.26 (electromagnetic), §F CT3.1–CT3.22a (colour), §G CT4.1–CT4.13 (saturation), §Cal-2 (registry note). Part 3: Def-H.1–Def-H.43, FT.1–FT.114, FT.116–FT.131, FT.133–FT.139, FT.141–FT.145 (structural tags), FT.153–FT.160. Part 4: Def-BM.1–Def-BM.20, BM.1–BM.152. Census: the 25 data-cell tags (Category B of the separation charter).

\textbf{0.2 Citation closure.} $\forall$ claims in Parts 1–4 cite only: A1–A9, C1–C13, D0.1–D4.3, P1–P5 (Parts 1–2), previously proved tags, and Census cells (row-wise). $\neg\exists$ imported external physical law, constant, or measured value. Physics vocabulary $:=$ naming of derived structures only. External comparison data enter only quarantined audit rows (Census FT.114, FT.120; P1, P5).

\textbf{0.2a Math \& hypothesis policy.} Standard mathematics (linear algebra, representation theory, tensor calculus, classification theorems) is cited as $[\text{math}]$-class tags in \textbf{Req} and is not part of the physical axioms. Conditional theorems carry their hypotheses explicitly in \textbf{Thm} and \textbf{Req}. The complete list of extra-axiomatic inputs is empty: CT2.6 (ii), CT2.6 (iii) — the two continuum-limit conditions — are derived theorems of Part 4 (Tier BM.24, BM.151): (ii) is witnessed by the ledger's own unbounded ladder, (iii) by the density-matrix spectral bound on the D2.1 coefficient fields; no others exist. The declarations at use stand as scope-inheritance markers over the same loci. CT2.6 (i) is a definitional precondition: the valid-universe condition of D2.11 in application (“there is something”), $\neg$ an extra-axiomatic hypothesis. CT2.8 carries no hypotheses: its reflection and factorization clauses are derived (D2.7, D1.9, D1.3, D2.8). The edge phase is registration data: its alphabet is derived (CT0.9) and its per-edge value is determined at crystallization (Def-EM.20, EM.20). \textbf{Integration note.} The integer-edge hypothesis of CT1.29a is \textit{not} independent: it is realized on materialized closed chains of the fully antialigned class (CT1.19b(3)). The ambient-richness condition in the CT3.22a Rmk is a witness scope for that Rmk's non-locality construction only; no theorem clause of Parts 1–2 rests on it. The materialization domain of CT1.19b is CT2.5’s own pole condition (no new input); its domain-generalization step from EW.17’s wrapping chains is flagged in its Pf. The per-edge independence hypothesis of CT4.5a is an ensemble scope for that theorem alone — its dynamics realization (a crystallization-rate datum) is NOT declared in Parts 1–2. The crystallization–tilt coupling of L0 (a seal-time state update: one Gibbs tilt step per A4 event) and the finite-incidence scope of L4 are theorem-local scopes of the same species, declared at use. The interface definitions of Part 3 (Def-H.35–Def-H.43) carry their own non-derived clauses explicitly.

\textbf{0.3 Free data.} $\exists!\,$ free dimensionful scale: $\mathfrak h$ (D2.9, P1). Every other quantity: dimensionless-derived, or a power of $\mathfrak h$. Free internal configuration parameters (P4-class, $\neg$constants of the theory): the P4 core (§E) plus the era registries Cal-2 (Part 2), Cal-F (Part 3), Cal-BM (Part 4).

\textbf{0.4 Theorem format (all parts).} Per tag: \textbf{Thm} — statement, purely symbolic · \textbf{Req} — \{axioms, conventions, definitions, prior theorems\} used · \textbf{Pf} — derivation, purely symbolic, justifications in $[\cdot]$, end $\blacksquare$ · \textbf{Phys} — short English summary of the theorem and its physics. Corollaries: \textbf{Cor}. Definitions: \textbf{Def.} blocks; scope conditions and constitutive clauses are carried inside the \textbf{Def.}/\textbf{Thm.} text. Forward pointers (scope conditions cited at use) are marked “forward pointer” at the citing site.

\textbf{0.5 Symbol reconciliations.} $P_M\equiv\mathcal P_M$ (C7); $F_Z\equiv F_{W^3}$ (C8); $\tau(v)\in\Xi_n$ (type set) vs $v\in V_n$ (vertex set), D0.8.

\textbf{0.6 Machine verification.} \texttt{checks/\allowbreak{}}: dependency-free exact-arithmetic verification suite (\texttt{python3\ checks/\allowbreak{}run\_\allowbreak{}all.py}). \texttt{verify\_\allowbreak{}era\_\allowbreak{}structure.py}: era-chain law, citation closure, census separation, tag ownership, declaration hygiene, \textbf{Thm·Req·Pf·Phys} completeness, index regeneration. \texttt{verify\_\allowbreak{}B.py}–\texttt{verify\_\allowbreak{}G.py} and the \texttt{verify\_\allowbreak{}*} content suite: algebraic/geometry/arithmetical identities, exact, no floating point. See \texttt{checks/\allowbreak{}README.md}.

\phantomsection
\section[§A — Axioms]{\texorpdfstring{§A — Axioms}{SA - Axioms}}

\begin{axiom}[{\texorpdfstring{A1}{A1}}]\label{tag:A1}
%dsp-level:0
\[
\mathcal Q=\{q_n\}_{n\ge1};
\]

%dsp-level:0
\[
\qquad \mathcal H_q\cong\mathbb C^2;
\]

%dsp-level:0
\[
\qquad \forall(p,q)\in\mathcal Q\times\mathcal Q,\ p\ne q:\quad
\rho_{pq}\in\mathfrak H(\mathcal H_p\otimes\mathcal H_q),\ \ \rho_{pq}\succeq0,\ \ \mathfrak T(\rho_{pq})=1 .
\]

\end{axiom}

\begin{axiom}[{\texorpdfstring{A2}{A2}}]\label{tag:A2}
%dsp-level:0
\[
\forall(p,q):\qquad C(p,q)=1\,\mathfrak q .
\]

\end{axiom}

\begin{axiom}[{\texorpdfstring{A3}{A3}}]\label{tag:A3}
%dsp-level:0
\[
\mathcal H_p\otimes\mathcal H_q=\vee^2\oplus\wedge^2;
\]

%dsp-level:0
\[
S(\sigma_\mu\otimes\sigma_\nu)=\sigma_\nu\otimes\sigma_\mu,\quad S^2=\mathbb1;\qquad
T(\sigma_\mu\otimes\sigma_\nu)=\eta^\mu\eta^\nu\,(\sigma_\mu\otimes\sigma_\nu),\quad
T^2=\mathbb1;
\]

%dsp-level:0
\[
\eta\in\{\pm1\}^4;\qquad
\exists!\,\mathfrak o\in\{0,1,2,3\}:\ \ \eta^{\mathfrak o}=-\eta^{\mu}\ \ \forall\mu\ne\mathfrak o .
\]

\end{axiom}

\begin{axiom}[{\texorpdfstring{A4}{A4}}]\label{tag:A4}
%dsp-level:0
\[
(p,q)\in\mathcal R\iff\bigl|\Pi_{\mathcal V}\rho_{pq}\bigr|\le\varepsilon_q\ \ \vee\ \ \bigl|\Pi_{\mathcal V}\rho_{pq}\bigr|\ge1-\varepsilon_q;\qquad
n_{\mathcal R}(p)=\bigl|\{q:(p,q)\in\mathcal R\}\bigr|
\]

%dsp-level:0
\[
\alpha_{pq}:=\begin{cases}0&\bigl|\Pi_{\mathcal V}\rho_{pq}\bigr|\le\varepsilon_q\\[1mm] \pi&\bigl|\Pi_{\mathcal V}\rho_{pq}\bigr|\ge1-\varepsilon_q\end{cases};\qquad
\alpha_{pq}\in\{0,\pi\}\quad\forall(p,q)\in\mathcal R;
\]

%dsp-level:0
\[
\exists\,\rho_{pq}:\ \bigl|\Pi_{\mathcal V}\rho_{pq}\bigr|\le\varepsilon_q;\qquad
\exists\,\rho_{p'q'}:\ \bigl|\Pi_{\mathcal V}\rho_{p'q'}\bigr|\ge1-\varepsilon_q .
\]

\end{axiom}

\begin{axiom}[{\texorpdfstring{A5}{A5}}]\label{tag:A5}
%dsp-level:0
\[
\rho_\circ=\tfrac14 I;\qquad
\rho\mapsto U\rho U^\dagger,\quad U\sim\mu_\star\bigl(\mathrm{SU}(2)\bigr),\quad U\notin\mathfrak L .
\]

\end{axiom}

\begin{axiom}[{\texorpdfstring{A6}{A6}}]\label{tag:A6}
%dsp-level:0
\[
\mu_\star:=\mu_{\mathrm{Haar}}\bigl(\mathrm{SU}(2)\bigr);\qquad
U\sim\mu_\star;\qquad
U\notin\mathfrak L;\qquad
\mathcal T(\rho):=U\rho U^\dagger;\qquad
\mathcal T^{-1}(\rho)=U^\dagger\rho U
\]

%dsp-level:0
\[
\mathfrak T\bigl(\mathcal T(\rho)\bigr)=\mathfrak T(\rho);\qquad
\mathcal T^{2}\ne\mathbb1;\qquad
\tau_n\to\tau_{n+1}:\ \ \rho\mapsto\mathcal T(\rho);\qquad
\neg\exists\,\tau^{-1}
\]

%dsp-level:0
\[
\bigl|\mathfrak L(\tau_{n+1})\bigr|\ \ge\ \bigl|\mathfrak L(\tau_n)\bigr|;\qquad
\bigl|\mathfrak L(\tau_{n+1})\bigr|\ =\ \bigl|\mathfrak L(\tau_n)\bigr|\,+\,n_{\mathrm{real\_out}}(\tau_n) .
\]

\end{axiom}

\begin{axiom}[{\texorpdfstring{A7}{A7}}]\label{tag:A7}
%dsp-level:0
\[
\forall n\ge0:\quad\Gamma=(V,E),\qquad E\subseteq V\times V
\]

%dsp-level:0
\[
\tau(v):=\bigl(\partial^{\leftarrow}(v),\partial^{\rightarrow}(v)\bigr)\in\Xi_n;\qquad
e\in E_{\mathcal R}\mapsto\phi_e;\qquad
e\in E_{\mathcal R}\mapsto\alpha_e\in\{0,\pi\}
\]

%dsp-level:0
\[
\mathfrak L=\bigcup_{n\ge0}\mathfrak L(\tau_n);
\]

%dsp-level:0
\[
\qquad \mathfrak L(\tau_n)\subseteq\mathfrak L(\tau_{n+1});
\]

%dsp-level:0
\[
\qquad \bigl|\mathfrak L(\tau_0)\bigr|=0;
\]

%dsp-level:0
\[
\qquad \bigl|\mathfrak L(\tau_n)\bigr|=\sum_{k<n}n_{\mathrm{real\_out}}(\tau_k) .
\]

\end{axiom}

\begin{axiom}[{\texorpdfstring{A8}{A8}}]\label{tag:A8}
%dsp-level:0
\[
\mathbf F=\mathbf F_{(+,\mathcal R)}\oplus\mathbf F_{(-,\mathcal R)}\oplus\mathbf F_{(+,\mathcal V)}\oplus\mathbf F_{(-,\mathcal V)}\oplus\mathbf F_{\mathrm{st}};
\]

%dsp-level:0
\[
\mathbf F_{(+,\mathcal R)}=F_G\oplus F_H\oplus F_{D^+};\qquad
\mathbf F_{(-,\mathcal R)}=F_{W^+}\oplus F_{W^-}\oplus F_Z;\qquad
\mathbf F_{(+,\mathcal V)}=F_{J^+};
\]

%dsp-level:0
\[
\mathbf F_{(-,\mathcal V)}=F_{J^-};\qquad
\mathbf F_{\mathrm{st}}=F_{\mathrm{EM}}\oplus F_{\mathrm{strong}}\oplus F_{D^-};
\]

%dsp-level:0
\[
F_{D^+}=\Lambda\,\mathbb I;\qquad
F_{D^-}:=\bigoplus_{n\ge4}F_n;\qquad
\mathbf F_M=P_M\,\mathbf F\,P_M;
\]

%dsp-level:0
\[
w(C)=\mathcal J(C)\!\!\prod_{e\in E(C)}\!\!G(e)\ \cdot\ (-1)^{\aleph(\sigma,C)};\qquad
Z(\mathcal R)=\sum_{C\in\mathcal R}\ \sum_\sigma w(C) .
\]

\end{axiom}

\begin{axiom}[{\texorpdfstring{A9}{A9}}]\label{tag:A9}
%dsp-level:0
\[
\mathcal U_\odot\in\mathfrak S\bigl(Z[\mathbf F_\oplus,M_\oplus]\bigr);\qquad
\mathcal U_\odot\notin\mathfrak S\implies\neg\bigl(\mathrm{A1}\wedge\cdots\wedge\mathrm{A9}\bigr) .
\]

$\blacksquare$

\end{axiom}

\phantomsection
\section[§B — Conventions]{\texorpdfstring{§B — Conventions}{SB - Conventions}}

\begin{convention}[{\texorpdfstring{C1}{C1}}]\label{tag:C1}
$\mu,\nu,\rho,\ldots\in\{0,1,2,3\}$: Pauli basis $\sigma_0=I_2,\sigma_1,\sigma_2,\sigma_3$ / emergent spacetime coordinates where a metric exists; $i,j,k,\ldots\in\{1,2,3\}$: spatial/Pauli-spatial; $a,b,c,\ldots\in\{1,\dots,8\}$: colour adjoint. Repeated indices: summed.

\end{convention}

\begin{convention}[{\texorpdfstring{C2}{C2}}]\label{tag:C2}
$T_A^\dagger=T_A$; $\ [T_A,T_B]=i\,f_{AB}{}^{C}\,T_C$; $\ f_{AB}{}^{C}\in\mathbb R$.

\end{convention}

\begin{convention}[{\texorpdfstring{C3}{C3}}]\label{tag:C3}
$\mathcal E_{st}:=\{\,X\in\mathfrak H(\mathbb C^4):\ SX=sX\ \wedge\ TX=tX\,\}$, $s,t\in\{+,-\}$. Symbol $\wedge$: antisymmetrization only ($\wedge^2$, $\varepsilon$-contractions); $\neg$ a sector-label conjunct.

\end{convention}

\begin{convention}[{\texorpdfstring{C4}{C4}}]\label{tag:C4}
$\varkappa$: sector-to-metric embedding scale (D2.1); $\kappa$: gravitational coupling (D2.9, CT2.9); $\varkappa\ne\kappa$; $\neg$ interchangeable.

\end{convention}

\begin{convention}[{\texorpdfstring{C5}{C5}}]\label{tag:C5}
$\Sigma:=\sum_i\sigma_i\otimes\sigma_i$ (operator); $\Sigma_{\mathrm{tr}}:=\delta^{ij}q_{ij}\in\mathbb R$ (scalar, D2.1); $\neg$ interchangeable.

\end{convention}

\begin{convention}[{\texorpdfstring{C6}{C6}}]\label{tag:C6}
$\mathsf E_{\mathrm{tr}}$: transit entropy; $\mathsf E_{\mathrm{bd}}$: bound entropy (D1.5); $\neg$ squared-field notation for either.

\end{convention}

\begin{convention}[{\texorpdfstring{C7}{C7}}]\label{tag:C7}
$P_M\equiv\mathcal P_M$ (A8). Rung-resolved: $\mathcal P_M^{(n)}$ (D2.3); $\ \mathbf F_M^{(n)}=\mathcal P_M^{(n)}\,\mathbf F\,\mathcal P_M^{(n)}$.

\end{convention}

\begin{convention}[{\texorpdfstring{C8}{C8}}]\label{tag:C8}
$F_Z\equiv F_{W^3}$: neutral-weak aspect of A8; $\ F_W:=F_{W^+}\oplus F_{W^-}\oplus F_Z$; pre-mixing seed generator: $W_3$. Bare $Z$ in a spectrum: physical neutral eigenstate (D3.10, EW.24); $Z(\cdot)$ with argument: partition function (A8) / masked resolvent (D2.4).

\end{convention}

\begin{convention}[{\texorpdfstring{C9}{C9}}]\label{tag:C9}
Entropy in nats: $S=\ln\Omega$; \_bits\_ only where explicit. One real edge $=1$ bit $=\ln2$ nats (CT0.9, CT1.29).

\end{convention}

\begin{convention}[{\texorpdfstring{C10}{C10}}]\label{tag:C10}
Register accounting (reality principle): registrations convert quantum data to classical alternatives at the ledger. Registering a classical alternative set $X$ for single-shot linear readout ($\mathfrak L$ must distinguish which value was carried, D0.7) assigns the alternatives mutually orthogonal register directions: register cost $=|X|$. Shannon content $=\log_2|X|$ bits is a distinct quantity: source-availability register count $\ne$ channel bit throughput. One channel quantum supplies $C_{\mathrm{cap}}=\dim_{\mathbb R}\mathfrak H(\mathcal H_p\otimes\mathcal H_q)=16$ register directions (A2, D0.9). Ledger-level registration data are real: edge transport acts on the reality line by $\{\pm1\}$ (D1.1, CT0.9).

\end{convention}

\begin{convention}[{\texorpdfstring{C11}{C11}}]\label{tag:C11}
A5/A6 reading: the conjugation clause is node-level — $\rho\mapsto U\rho U^\dagger$ [A6: $\mathcal T$] with $\rho\in\mathfrak H(\mathcal H_q)\cong\mathfrak H(\mathbb C^2)$ and $U\sim\mu_\star(\mathrm{SU}(2))$ ($2\times2$: well-typed); $\rho_\circ=\tfrac14I$ is the vacuum pair state ($4\times4$). Node tick lifecycle: wait $\to$ data in $\to$ Haar draw/tick [A6: $\tau_n\to\tau_{n+1}$] (virtual states as basis) $\to$ data out/hop $\to$ wait. Draws are not retained: $U\notin\mathfrak L$ [A5, A6].

\end{convention}

\begin{convention}[{\texorpdfstring{C12}{C12}}]\label{tag:C12}
In A8's $Z(\mathcal R)=\sum_{C\in\mathcal R}\sum_\sigma w(C)$ the summation domain reads $\mathfrak C(\mathcal R)$ (D0.10): the configuration set on the real-edge set. $\mathcal R$ itself (A4, D0.1) is a set of pairs and is never a summation domain.

\end{convention}

\begin{convention}[{\texorpdfstring{C13}{C13}}]\label{tag:C13}
Signature convention. A3 fixes $T$ only through the products $\eta^\mu\eta^\nu$, and fixes $\eta$ only up to its odd slot $\mathfrak o$ and its overall sign. The corpus takes $\mathfrak o:=0$ (the identity slot $\sigma_0=I_2$, C1) and $\eta^{0}:=-1$, i.e. $\eta:=(-,+,+,+)$ (mostly-plus; D2.1: $\bar g_{00}:=-1$). The overall sign does not change $T$ [CT1.0(2)]; the three choices $\mathfrak o\in\{1,2,3\}$ are conjugate under $\mathrm{SU}(2)$ relabelling of the spatial Pauli axes [CT1.0(4)]; $\mathfrak o=0$ versus $\mathfrak o\ge1$ is the choice of the $\mathrm{SU}(2)$-covariant representative [CT1.0(5)], on which CT1.2's generator table rests [CT1.0 Rmk]. Uses of the representative carry the tag C13. $\neg$ an axiom.

\end{convention}

\phantomsection
\section[§C — The definition layer]{\texorpdfstring{§C — The definition layer}{SC - The definition layer}}

Definitions $:=$ vocabulary stipulations. $\forall$ content-bearing claims: theorems of Parts 1–4. Tiers 0–1 are qubit-era vocabulary; Tiers 2–4 name the structures the bosonic era derives — the definition layer is the corpus-wide vocabulary root (the classical canon's A §Part 3).

\phantomsection
\subsection[Tier 0 — bookkeeping]{\texorpdfstring{Tier 0 — bookkeeping}{Tier 0 - bookkeeping}}

\begin{definition}[{\texorpdfstring{D0.1}{D0.1}}]\label{tag:D0-1}
$\mathcal R:=\{(p,q):\ \mathrm{A4}\ \text{holds}\}$; $\ \mathcal V:=(\mathcal Q\times\mathcal Q)\setminus\mathcal R$; $\ n_{\mathcal R}:=\sum_p n_{\mathcal R}(p)$; global vacuum $\iff n_{\mathcal R}=0$.

\end{definition}

\begin{definition}[{\texorpdfstring{D0.2}{D0.2}}]\label{tag:D0-2}
$\bigl|\Pi X\bigr|:=\mathfrak T(\Pi X\Pi)=\mathfrak T(\Pi X)$; $\ \ \rho\succeq0\wedge\mathfrak T\rho=1\implies\bigl|\Pi\rho\bigr|\in[0,1]$.

\end{definition}

\begin{definition}[{\texorpdfstring{D0.3}{D0.3}}]\label{tag:D0-3}
$S_{\mathrm{raw}}(\rho):=-\mathfrak T(\rho\log\rho)=\mathfrak T(\rho K)$; $\ K:=-\log\rho$ (faithful $\rho$). $\ \gamma:=1-4\varepsilon_q$ [A4, D0.4, CT0.8]; $\ \gamma\in(0,1]$ for $\varepsilon_q\in[0,\tfrac14)$. Gibbs tilt (functional calculus): $\tilde{\mathcal T}_\gamma(X):=Xe^{-\gamma X}/\mathfrak T(Xe^{-\gamma X})$ where $\mathfrak T(Xe^{-\gamma X})\ne0$; on states: $\tilde{\mathcal T}_\gamma(\rho)=\rho\,e^{-\gamma\rho}/\mathfrak T(\rho\,e^{-\gamma\rho})$; $\ \gamma=0:$ the identity on trace-$1$ states; for faithful $\rho$: $\tilde{\mathcal T}_\gamma(K)=K\rho^{\gamma}/\mathfrak T(K\rho^{\gamma})$ $\ (\gamma=1:\ =\rho(-\log\rho)/S_{\mathrm{raw}}(\rho))$. Variance functional (reference state explicit): $\mathbb V_\rho(X):=\mathfrak T(\rho X^2)-\bigl(\mathfrak T(\rho X)\bigr)^2$. Lyapunov functional of the tilt: $V(\tilde K):=\mathfrak T(\tilde K^2)$, $\ \tilde K:=K-\mathfrak T(K\rho)\,I$ — monotonicity: L1 (§F, forward pointer). $\tilde{\mathcal T}_\gamma\ne\mathcal T$ [A6]; $\neg$ interchangeable.

\end{definition}

\begin{definition}[{\texorpdfstring{D0.4}{D0.4}}]\label{tag:D0-4}
$\varepsilon_q\in[0,1]$: threshold of A4; admissible range: CT0.8; sharp choice: $\varepsilon_q=0$.

\end{definition}

\begin{definition}[{\texorpdfstring{D0.5}{D0.5}}]\label{tag:D0-5}
$\Phi$: phase alphabet of a real edge (CT0.9: $\Phi=\{0,\pi\}$); $\ \sigma_e:=e^{i\phi_e}=(-1)^{\phi_e/\pi}\in\{+1,-1\}$.

\end{definition}

\begin{definition}[{\texorpdfstring{D0.6}{D0.6}}]\label{tag:D0-6}
hop $:=$ one traversal of one $e\in E_{\mathcal R}$; $\ \delta h(\gamma):=\#\{e\in\gamma\}$; $\ \mathfrak h:=$ length assigned to one hop.

\end{definition}

\begin{definition}[{\texorpdfstring{D0.7}{D0.7}}]\label{tag:D0-7}
$\mathfrak L:=$ ordered record of retained events [A6, A7: stratification $\mathfrak L=\bigcup_{n\ge0}\mathfrak L(\tau_n)$, $|\mathfrak L(\tau_0)|=0$, monotone under ticks]; $\ \delta\tau:=$ ledger ticks (retained updates along real edges); $\ U\notin\mathfrak L$ [A5, A6].

\end{definition}

\begin{definition}[{\texorpdfstring{D0.8}{D0.8}}]\label{tag:D0-8}
$n(v):=\partial^{\leftarrow}(v)+\partial^{\rightarrow}(v)$; $\ V_n:=\{v\in V:\ n(v)=n\}$ (vertex set); $\ \Xi_n:=\tau(V_n)=\{(a,b)\in\mathbb N^2:\ a+b=n\}$ (type set); “at rung $n$” $:=$ on $V_n$.

\end{definition}

\begin{definition}[{\texorpdfstring{D0.9}{D0.9}}]\label{tag:D0-9}
$D_{\mathrm{tick}}(n):=\#\{$independent real registers a vertex of $V_n$ pushes across one pair channel per tick$\}$ — counted by the register accounting of \textbf{C10} (one orthogonal direction per admissible classical alternative: source availability, single-shot linear readout; $\neg$ Shannon $\log_2$ counting). $\ C_{\mathrm{cap}}:=\dim_{\mathbb R}\mathfrak H(\mathcal H_p\otimes\mathcal H_q)$ (register directions supplied by $1\mathfrak q$, A2); $\ O(n):=\max\bigl(0,\ D_{\mathrm{tick}}(n)-C_{\mathrm{cap}}\bigr)$.

\end{definition}

\begin{definition}[{\texorpdfstring{D0.10}{D0.10}}]\label{tag:D0-10}
$\mathfrak C(\mathcal R):=$ configuration set on $E_{\mathcal R}$: finite oriented edge sets $E(C)\subseteq E_{\mathcal R}$ with sign assignments $\sigma$, seal data (D1.2) and vertex types (A7). $w(C)$ (A8) is defined on $\mathfrak C(\mathcal R)$; valid universes (D2.11) form a subset; $Z(\mathcal R):=\sum_{C\in\mathfrak C(\mathcal R)}\sum_\sigma w(C)$ (C12).

\end{definition}

\phantomsection
\subsection[Tier 1 — graph, chains, ledger dynamics]{\texorpdfstring{Tier 1 — graph, chains, ledger dynamics}{Tier 1 - graph, chains, ledger dynamics}}

\begin{definition}[{\texorpdfstring{D1.1}{D1.1}}]\label{tag:D1-1}
$e=(p,q)\in\mathcal R_{\rightarrow}(p)\ \longmapsto\ (\mathcal G_n,\phi_e)$; $\ \mathcal G_n:=$ rung class of $q$; $\ \phi_e:$ D0.5; $\ \neg$registration $\implies e\in\mathcal V$. Reality principle of registration (C10): $\phi_e:=$ phase of the registration transport's action on the reality line $\wedge^2$ (1-dimensional, since $\dim\mathcal H_q=2$ [A1]); the transport preserves the $(S,T)$-grading [A3, A4] and the real Pauli-frame structure — registrations are classical data. On a minimal real edge ($\widehat{ds}^{\,2}=1$ [D1.6]) the geometric factor of the transfer weight $G(e)$ [D3.5] is the edge's alignment holonomy $e^{i\alpha_e}$ — the binary registration read of the A4 branch index at the seal ($\alpha_e=\alpha_{pq}$ for $e=(p,q)$ [A4]; per-edge determination read in \textbf{E}: Def-EM.20, EM.20, forward pointer).

\end{definition}

\begin{definition}[{\texorpdfstring{D1.2}{D1.2}}]\label{tag:D1-2}
seal $:=$ transition of one pair across the A4 criterion (entry/exit of either branch): $\ J^+:\ n_{\mathcal R}:0\to1$ (crystallization); $\ J^-:\ n_{\mathcal R}:1\to0$ (dissolution).

\end{definition}

\begin{definition}[{\texorpdfstring{D1.3}{D1.3}}]\label{tag:D1-3}
Janus-balanced $\iff\ \#J^+=\#J^-$.

\end{definition}

\begin{definition}[{\texorpdfstring{D1.4}{D1.4}}]\label{tag:D1-4}
chain $c=(e_1,\dots,e_k)$: $e_i\in E_{\mathcal R}$, $\ t(e_i)=s(e_{i+1})$; $\ s(c):=s(e_1)$; $\ t(c):=t(e_k)$; $\ \partial c:=t(c)-s(c)$ (formal difference of vertices); closed $\iff\ \partial c=0\ \wedge\ \#J^+=\#J^-$.

\end{definition}

\begin{definition}[{\texorpdfstring{D1.5}{D1.5}}]\label{tag:D1-5}
bound $:=$ closed self-contained chain (content: fermionic); transit $:=$ propagating open chain (content: bosonic); $\ \Delta S=\mathsf E_{\mathrm{bd}}+\mathsf E_{\mathrm{tr}}$; $\ \mathsf E_{\mathrm{bd}}:=\Delta S_{\mathrm{bound}}$; $\ \mathsf E_{\mathrm{tr}}:=\Delta S_{\mathrm{transit}}$.

\end{definition}

\begin{definition}[{\texorpdfstring{D1.6}{D1.6}}]\label{tag:D1-6}
$ds^2(e):=g_{\mu\nu}(e)\,\Delta x^\mu_e\,\Delta x^\nu_e$ [$g_{\mu\nu}:$ D2.1]; $\ \widehat{ds}^{\,2}(e):=ds^2(e)/\mathfrak h^2$.

\end{definition}

\begin{definition}[{\texorpdfstring{D1.7}{D1.7}}]\label{tag:D1-7}
$\Theta(\gamma):=\sum_{e\in\gamma}\widehat{ds}^{\,2}(e)$; $\ \Phi(\gamma):=\sum_{e\in\gamma}\alpha_e\,\widehat{ds}^{\,2}(e)$ [$\alpha_e\in\Phi:$ per-edge phase, D3.5; uniform-phase case: $\Phi=\alpha\,\Theta$]; $\ \operatorname{Hol}_{\mathrm{spiral}}(\gamma):=e^{i\Phi(\gamma)}$; $\ \gamma\in\mathsf{Trivial}\iff\forall e\in\gamma:\ \widehat{ds}^{\,2}(e)=0$; wrapping $\iff$ spatial hop projection has winding $\ne0$ around a closed ledger cycle.

\end{definition}

\begin{definition}[{\texorpdfstring{D1.8}{D1.8}}]\label{tag:D1-8}
$\tau_n:=\sum\delta\tau$ accumulated along a causal chain of real edges; emergent coordinate $x^0:=\tau_n$.

\end{definition}

\begin{definition}[{\texorpdfstring{D1.9}{D1.9}}]\label{tag:D1-9}
tick locality: $\ \operatorname{dom}\bigl(\mathrm{tick}(v)\bigr)\subseteq\{v\}\cup\{e\in E:\ e\ni v\}$, pre- and post-tick data only; $\neg\exists$ read at edge-distance $>1$.

\end{definition}

\begin{definition}[{\texorpdfstring{D1.10}{D1.10}}]\label{tag:D1-10}
node force: $\ \mathbf F\!\restriction_v$ resolves local information content into channels \{bound, transit, holonomy, charge, flux\}; \{bound, transit\} $:=$ entropy channels [D1.5]; \{holonomy, charge, flux\} $:=$ registration data [D1.1] $\cup$ transition amplitude [D3.3].

\end{definition}

\begin{definition}[{\texorpdfstring{D1.11}{D1.11}}]\label{tag:D1-11}
accessible information content at $v$ $\leftrightarrow n(v)$; content beyond the capacity bound (CT0.11): $\neg$accessible [D4.1].

\end{definition}

\begin{definition}[{\texorpdfstring{D1.12}{D1.12}}]\label{tag:D1-12}
$\mathcal A(n):=\{$aspects of $\mathbf F$ acting nontrivially at rung $n\}$; value derived (CT0.13), $\neg$ stipulated.

\end{definition}

\phantomsection
\subsection[Tier 2 — masks, geometry, units]{\texorpdfstring{Tier 2 — masks, geometry, units}{Tier 2 - masks, geometry, units}}

\begin{definition}[{\texorpdfstring{D2.1}{D2.1}}]\label{tag:D2-1}
sector coefficients [basis: CT1.2]: $\ X_{(+)}=a_{00}E_0+q_{ij}E^{ij}$, $q_{ij}=q_{ji}$; $\ X_{(-)}=w_iW_i$;

%dsp-level:0
\[
\Sigma_{\mathrm{tr}}:=\delta^{ij}q_{ij};
\]

%dsp-level:0
\[
\quad \hat h_{ij}:=q_{ij}-\tfrac13\delta_{ij}\Sigma_{\mathrm{tr}};
\]

%dsp-level:0
\[
\quad h:=\tfrac13\Sigma_{\mathrm{tr}};
\]

%dsp-level:0
\[
\quad N_i:=\beta w_i;
\]

%dsp-level:0
\[
\quad \Omega:=e^{\alpha_\Omega a_{00}};
\]

%dsp-level:0
\[
\quad \alpha_\Omega,\beta,\varkappa\in\mathbb R;
\]

%dsp-level:0
\[
\bar g_{00}:=-1;\quad
\bar g_{0i}:=\varkappa N_i;\quad
\bar g_{ij}:=\delta_{ij}+\varkappa\bigl(\hat h_{ij}+\tfrac13\delta_{ij}\Sigma_{\mathrm{tr}}\bigr);\quad
g_{\mu\nu}:=\Omega^2\,\bar g_{\mu\nu}.
\]

\end{definition}

\begin{definition}[{\texorpdfstring{D2.2}{D2.2}}]\label{tag:D2-2}
$\Upsilon:=\ln\Omega$; $\ \Upsilon_\mu:=\bar\nabla_\mu\Upsilon$.

\end{definition}

\begin{definition}[{\texorpdfstring{D2.3}{D2.3}}]\label{tag:D2-3}
$\mathcal K:$ aspect space; $\ \mathcal K^{(A)}:$ block of aspect $A$; $\ \mathcal K_n:=\bigoplus_{A\in\mathcal A(n)}\mathcal K^{(A)}$; $\ \mathcal P_M^{(n)}:\mathcal K\to\mathcal K_n:$ orthogonal projection.

\end{definition}

\begin{definition}[{\texorpdfstring{D2.4}{D2.4}}]\label{tag:D2-4}
$Z(\mathcal R,n):=\bigl\langle\mathrm{out}\bigl|(1-\mathbf F_M^{(n)})^{-1}\bigr|\mathrm{in}\bigr\rangle$; poles of $Z(\mathcal R,n):=$ rung-$n$ spectrum; residues $:=$ amplitudes.

\end{definition}

\begin{definition}[{\texorpdfstring{D2.5}{D2.5}}]\label{tag:D2-5}
massless $\iff$ resolvent pole at $k^2=0$; massive $\iff$ pole at $k^2=m^2>0$; $m^2:=$ pole displacement.

\end{definition}

\begin{definition}[{\texorpdfstring{D2.6}{D2.6}}]\label{tag:D2-6}
$\omega:$ orientation-reversal involution: $s(c)\leftrightarrow t(c)$; chirality $\mapsto-$chirality; charge $\mapsto-$charge.

\end{definition}

\begin{definition}[{\texorpdfstring{D2.7}{D2.7}}]\label{tag:D2-7}
$\Theta_{\mathfrak L}:$ ledger reflection: reverses ledger order; $J^+\leftrightarrow J^-$; $G(e)\mapsto G(e)^\ast$.

\end{definition}

\begin{definition}[{\texorpdfstring{D2.8}{D2.8}}]\label{tag:D2-8}
convergence sheet: $1\mapsto1+i\varepsilon$, $\ \varepsilon\to0^+$, for $\forall$ resolvents and configuration sums.

\end{definition}

\begin{definition}[{\texorpdfstring{D2.9}{D2.9}}]\label{tag:D2-9}
units $:=$ powers of $\mathfrak q$; $\ c=\hbar=k_B=1$; $\ [\mathfrak h]=\mathfrak q^{-1/2}$; $\ \mathfrak h:=$ single free scale; $\ \kappa,\ G_N:$ names for combinations of $\mathfrak h$ (CT2.9), $\neg$ independent inputs.

\end{definition}

\begin{definition}[{\texorpdfstring{D2.10}{D2.10}}]\label{tag:D2-10}
$\mathfrak M=\begin{pmatrix}B&Q\\ R&F\end{pmatrix}$: even blocks $B,F$; odd blocks $Q,R$; Hermitian $\iff B^\dagger=B\wedge F^\dagger=F\wedge R=Q^\dagger$; $\ \operatorname{Str}\mathfrak M:=\mathfrak T B-\mathfrak T F$; $\ [X,Y\}:=XY-(-1)^{|X||Y|}YX$ ($X,Y$ homogeneous).

\end{definition}

\begin{definition}[{\texorpdfstring{D2.11}{D2.11}}]\label{tag:D2-11}
valid universe $\iff$ real-edge network closed $\ \wedge\ \#J^+=\#J^-$ [D1.3] $\ \wedge\ \in\operatorname{Im}($active mask$)$ $\ \wedge\ \neg\exists$ dangling ledger boundary beyond the sealed Janus boundary.

\end{definition}

\begin{definition}[{\texorpdfstring{D2.12}{D2.12}}]\label{tag:D2-12}
Aspect realization (slot$\to$aspect bridge): $\iota_n:\mathfrak u(n)\to\mathcal K$ — linear realization of the rung-$n$ slot algebra in aspect space, with $\iota_n\bigl(Z(\mathfrak u(2))\bigr)$ spanning $\mathcal K^{(\mathrm{EM})}$ ($n=2$) and $\iota_n\bigl(\mathfrak{su}(3)\bigr)$ spanning $\mathcal K^{(\mathrm{strong})}$ ($n=3$), per the derivation of CT0.13. Slot algebras and pair algebras remain distinct objects (F header); D3.2 “already realized at rung $m$” $:=$ the abstract isomorphism class lies in $\iota_m(\mathfrak u(m))$ or in the pair blocks of CT0.12. \textbf{Realization domain:} exactly $n\le3$ — for $n\ge4$ no aspect realization exists [CT4.4a]; the silence of this definition at $n\ge4$ is a non-existence consequence, not an underdetermination.

\end{definition}

\begin{definition}[{\texorpdfstring{D2.13}{D2.13}}]\label{tag:D2-13}
Subnetwork and microstate space \textit{(scope clauses)}. $\Sigma\subseteq U$ [D2.11] $:=$ a finite connected vertex set together with its interior edges; $\partial\Sigma:=$ edges with exactly one endpoint in $\Sigma$. $\mathcal M(\Sigma):=$ valid microscopic configurations on $\Sigma$: restrictions $C|_\Sigma$ of valid ambient configurations $C$ [D2.11], carrying edge signs $\sigma_e$ [D0.5, CT0.9], geometric data $\widehat{ds}^{\,2}(e)$ [D1.6], seal data [D1.2] and vertex types [A7, D0.8] — subject to the rung masks [CT0.14, CT2.2] and Janus balance [D1.3], $\neg$ to colour sealing (sealing is ambient-global [CT3.18]; the completeness filter acts downstream: CT3.22a Rmk). $|\Sigma|:=\#\{\text{real edges touching}\ \Sigma\}$ (interior $+$ boundary).

\end{definition}

\begin{definition}[{\texorpdfstring{D2.14}{D2.14}}]\label{tag:D2-14}
Colour boundary and the contraction-completeness projector. $\mathrm{Colour}(U):=\{u\in\mathcal Q:\ u\ \text{carries a colour attachment}\ [\mathrm{D3.8}]\}$; for valid $C$, $\pi_C(v)\subseteq\mathrm{Colour}(U):=$ the other members of the elementary invariant contraction group of $v$ under $C$ (one for a meson pair via $\delta^i{}_j$; two for a (anti)baryon triple via $\varepsilon_{ijk}/\varepsilon^{ijk}$ [CT3.17]); $\pi_C$ is total on $\mathrm{Colour}(U)$ for valid $C$ [CT3.18 $+$ D2.11 applied to the colour ledger]. For $\Sigma$ [D2.13]: $\mathrm{Colour}(\Sigma):=\mathrm{Colour}(U)\cap\Sigma$;

%dsp-level:0
\[
\partial_{\mathrm{col}}(\Sigma;C):=\{v\in\mathrm{Colour}(\Sigma):\ \pi_C(v)\not\subseteq\Sigma\};\qquad
\mathcal C(\Sigma;C):=\begin{cases}1&\partial_{\mathrm{col}}(\Sigma;C)=\varnothing\\ 0&\text{otherwise}\end{cases};\qquad
\Pi_{\mathcal C}(\Sigma;C):=\mathcal C(\Sigma;C)\cdot\mathbb I^{(\mathrm{strong})}_\Sigma
\]

(Hermitian scalar idempotent on the colour slot space of $\Sigma$ [D3.1]; matches the mask pattern $\mathcal P_M^{(n)}$ [CT0.14]).

\end{definition}

\phantomsection
\subsection[Tier 3 — rung-2 and rung-3 structures]{\texorpdfstring{Tier 3 — rung-2 and rung-3 structures}{Tier 3 - rung-2 and rung-3 structures}}

\begin{definition}[{\texorpdfstring{D3.1}{D3.1}}]\label{tag:D3-1}
$\mathcal H^{(n)}_{\mathrm{slot}}:=\operatorname{span}_{\mathbb C}\{|1\rangle,\dots,|n\rangle\}\cong\mathbb C^{n}$; $n=3$: slot index $i\in\{1,2,3\}$ $:=$ colour.

\end{definition}

\begin{definition}[{\texorpdfstring{D3.2}{D3.2}}]\label{tag:D3-2}
subset principle: strictly new at rung $n$ $:=$ part of the rung-$n$ structure $\neg$realized at any $m<n$.

\end{definition}

\begin{definition}[{\texorpdfstring{D3.3}{D3.3}}]\label{tag:D3-3}
at $v$ of type $(1,1)$, edge signs $\sigma_{\leftarrow}(v),\sigma_{\rightarrow}(v)$:

%dsp-level:0
\[
f_v:=\frac{\sigma_{\rightarrow}(v)-\sigma_{\leftarrow}(v)}{2}\in\{-1,0,+1\};\qquad
\aleph_v:=|f_v|;\qquad
\aleph(\sigma,C):=\sum_{v\in C}\aleph_v .
\]

\end{definition}

\begin{definition}[{\texorpdfstring{D3.4}{D3.4}}]\label{tag:D3-4}
$q_{\mathrm{EM}}(c):=\dfrac{\sigma_{s(c)}-\sigma_{t(c)}}{2}$ (oriented chain $c$). Closed chains: the basepoint core read $q_{\mathrm{core}}$ is defined, ternary, and charge-conjugation odd [EM.9a].

\end{definition}

\begin{definition}[{\texorpdfstring{D3.5}{D3.5}}]\label{tag:D3-5}
$G^{(n)}(e):=\sigma_e\,\exp\!\bigl(i\,\alpha_e\,\widehat{ds}^{\,2}(e)\bigr)$ [D0.5, D1.6]; $\ \alpha_e\in\Phi:$ per-edge dynamic phase $=\alpha_{pq}$ for $e=(p,q)$ — the A4 branch index (read in Part 2 §E: Def-EM.20, EM.20, forward pointer); dimensionless; the same alphabet at every rung. $\ \alpha_{\mathrm{EM}}^{(n)}:$ EM-block spectral read at rung $n$ (EM.17, forward pointer).

\end{definition}

\begin{definition}[{\texorpdfstring{D3.6}{D3.6}}]\label{tag:D3-6}
$\mathcal J(C):=\prod_{J^+\in C}(+1)\cdot\prod_{J^-\in C}(+1)^\ast$ $\ \implies\ \mathcal J(\Theta_{\mathfrak L}C)=\mathcal J(C)^\ast$.

\end{definition}

\begin{definition}[{\texorpdfstring{D3.7}{D3.7}}]\label{tag:D3-7}
photon $:=$ trivial chain [D1.7] $\subseteq V_2$; $\ $gluon $:=$ trivial chain $\cap\,V_3\ne\varnothing$.

\end{definition}

\begin{definition}[{\texorpdfstring{D3.8}{D3.8}}]\label{tag:D3-8}
colour triplet $q^i$: one fundamental slot index; colour antitriplet $r_i$: one dual slot index; quark $:=$ bound (fermionic) loop with a triplet attachment.

\end{definition}

\begin{definition}[{\texorpdfstring{D3.9}{D3.9}}]\label{tag:D3-9}
baryon $B:=\varepsilon_{ijk}\,q^iq^jq^k$; $\ $meson $M:=\delta^i{}_j\,q^jr_i$. [Disambiguation: $M$ here $\ne$ total mask $M_\oplus$ of A9.]

\end{definition}

\begin{definition}[{\texorpdfstring{D3.10}{D3.10}}]\label{tag:D3-10}
$W^\pm:=\tfrac1{\sqrt2}\bigl(W_1\mp iW_2\bigr)$; $\ Z:=\cos\theta_W\,W_3-\sin\theta_W\,Y$; $\ A:=\sin\theta_W\,W_3+\cos\theta_W\,Y$ [$\theta_W:$ EW.24].

\end{definition}

\phantomsection
\subsection[Tier 4 — saturation]{\texorpdfstring{Tier 4 — saturation}{Tier 4 - saturation}}

\begin{definition}[{\texorpdfstring{D4.1}{D4.1}}]\label{tag:D4-1}
$F_{D^-}=\bigoplus_{n\ge4}F_n$ [A8]: single possibility space receiving $\forall$ configurations with rung $n\ge4$.

\end{definition}

\begin{definition}[{\texorpdfstring{D4.2}{D4.2}}]\label{tag:D4-2}
$D^-$ absorbs information at rate $\propto O(n)$ [D0.9]; absorbed $\implies$ locked $:=$ retained in $\mathfrak L$, $\neg$accessible [D1.11].

\end{definition}

\begin{definition}[{\texorpdfstring{D4.3}{D4.3}}]\label{tag:D4-3}
$F_{D^+}=\Lambda\,\mathbb I$ [A8]: uniform scalar aspect; $\neg$configuration-dependent.

\end{definition}

\phantomsection
\section[§D — Master symbol table]{\texorpdfstring{§D — Master symbol table}{SD - Master symbol table}}

%tbl-level:2
{\footnotesize\setlength{\tabcolsep}{4pt}
\begin{longtable}{>{\raggedright\arraybackslash}p{\dimexpr0.2460\textwidth-2\tabcolsep\relax}|>{\raggedright\arraybackslash}p{\dimexpr0.6032\textwidth-2\tabcolsep\relax}|>{\raggedright\arraybackslash}p{\dimexpr0.1508\textwidth-2\tabcolsep\relax}}
\toprule
Symbol & Meaning & Declared \\
\addlinespace[2pt]
\midrule
\endfirsthead
\toprule
Symbol & Meaning & Declared \\
\addlinespace[2pt]
\midrule
\endhead
\bottomrule
\endlastfoot
$\mathcal Q$, $q_n$ & qubit set; $n$-th qubit & A1 \\
$\mathcal H_q$ & carrier space of $q$, $\cong\mathbb C^2$ & A1 \\
$\rho_{pq}$ & state of ordered pair $(p,\allowbreak q)$ & A1 \\
$\mathfrak H(X)$ & real space of Hermitian operators on $X$ & A1 \\
$\mathfrak T(A)$ & trace of $A$ & A1 \\
$\succeq0$ & positive semidefinite & A1 \\
$C(p,\allowbreak q)$, $\mathfrak q$ & channel capacity of a pair; channel quantum & A2 \\
$\vee^2$, $\wedge^2$ & symmetric, antisymmetric square & A3 \\
$S$, $T$ & swap involution; sign-grading involution ($T=\allowbreak \mathsf D_\zeta\otimes\mathsf D_\zeta$; metric-sign reading: C13) & A3, CT1.0 \\
$\sigma_\mu$ & Pauli basis, $\mu\in\allowbreak \{0,1,2,3\}$, $\sigma_0=\allowbreak I_2$ & A3, C1 \\
$\mathbb1$, $I$ & identity operator; identity matrix & A3, A5 \\
$\eta$ & sign vector of the $T$-grading, $\eta\in\allowbreak \{\pm1\}^4$, one odd entry; representative $(-,\allowbreak +\allowbreak ,\allowbreak +\allowbreak ,\allowbreak +\allowbreak )$ & A3, C13 \\
$\mathfrak o$, $\eta^{(\mathfrak o)}$ & odd-sign slot of $\eta$; sign vector with a single $-1$ at slot $\mathfrak o$ & A3, CT1.0 \\
$\zeta_i$, $\mathsf D_\zeta$, $\mathrm{sf}$ & $\eta^0\eta^i$; Bloch sign map $\sigma_\mu\mapsto\allowbreak \zeta_\mu\sigma_\mu$; spin flip $X\mapsto\allowbreak \sigma_2X^{\mathrm T}\sigma_2$ & CT1.0 \\
$\mathcal R$, $\mathcal V$ & real (crystallized) set; virtual set & A4, D0.1 \\
$\Pi_{\mathcal V}$, $\Pi_{\mathcal R}$ & projections onto $\wedge^2$, $\vee^2$ & A4, CT0.7 \\
$\varepsilon_q$ & crystallization threshold & A4, D0.4 \\
$n_{\mathcal R}(p)$, $n_{\mathcal R}$ & real-edge count at $p$; global count & A4, D0.1 \\
$\rho_\circ$ & vacuum pair state $I/4$ & A5 \\
$U$, $\mu_\star$, $\dagger$ & Haar draw on $\mathrm{SU}(2)$; Haar measure; conjugate transpose & A5, A6 \\
$\mathfrak L$, $\mathfrak L(\tau_n)$, $n_{\mathrm{real\_out}}$ & ledger; ledger stratum at tick $\tau_n$; real-out count per tick & A5, A6, A7, D0.7 \\
$\mathcal T$, $\mathcal T^{-1}$ & tick map (Haar conjugation); its inverse & A6 \\
$\tilde{\mathcal T}_\gamma$, $\gamma$, $K$, $S_{\mathrm{raw}}$, $\mathbb V_\rho$ & $\varepsilon_q$-dependent Gibbs tilt; tilt strength $1-4\varepsilon_q$; surprisal; raw entropy; variance & D0.3 \\
$\tilde K$, $V(\tilde K)$ & centered surprisal; Lyapunov functional & D0.3, L1 \\
$\Gamma=\allowbreak (V,\allowbreak E)$ & directed graph & A7 \\
$\tau(v)$, $\partial^{\leftarrow}$, $\partial^{\rightarrow}$ & vertex type; in-degree; out-degree & A7 \\
$n(v)$, $V_n$, $\Xi_n$ & rung of $v$; rung-$n$ vertex set; type set & D0.8 \\
$E_{\mathcal R}$, $\phi_e$, $\sigma_e$, $\Phi$ & real edges; edge phase; edge sign; phase alphabet & A7, D0.5 \\
$\delta h$, $\mathfrak h$ & hop count; hop length & D0.6 \\
$\delta\tau$, $\tau_n$ & ledger tick count; ledger time & D0.7, D1.8 \\
$D_{\mathrm{tick}}(n)$, $C_{\mathrm{cap}}$, $O(n)$ & tick load; channel capacity; oversubscription & D0.9 \\
$\mathbf F$, $F_A$ & force operator; aspect $A$ & A8 \\
$F_{W^+}$, $F_{W^-}$, $F_Z$, $F_W$ & charged weak; neutral weak ($F_Z\equiv F_{W^3}$); weak block & A8, C8 \\
$F_{J^+}$, $F_{J^-}$ & crystallization; dissolution aspects & A8, D1.2 \\
$F_{\mathrm{EM}}$, $F_{\mathrm{strong}}$ & electromagnetic; colour aspects & A8 \\
$\Lambda$, $\mathbb I$ & stretch scalar; identity on force space & A8 \\
$\mathcal U_\odot$, $\mathfrak S$, $\mathbf F_\oplus$, $M_\oplus$ & observed universe; solution set; total force; total mask & A9 \\
$\mathfrak C(\mathcal R)$ & configuration set on $E_{\mathcal R}$; summation domain of $Z$ & D0.10, C12 \\
$\tilde{\mathcal T}_\gamma(X)$ & operator Gibbs tilt (functional calculus) & D0.3 \\
$\mathcal E_{st}$ & joint $(S,\allowbreak T)$ eigenspace & C3 \\
$f_{AB}{}^{C}$ & structure constants (Hermitian convention) & C2 \\
$\mathsf E_{\mathrm{bd}}$, $\mathsf E_{\mathrm{tr}}$ & bound, transit entropy & C6, D1.5 \\
$E_0$, $E^{ij}$, $W_i$, $J_i$, $K_i$ & generator basis of $\mathfrak H(\mathbb C^4)$ & CT1.2 \\
$\mathcal A(n)$ & active aspect set at rung $n$ & D1.12, CT0.13 \\
$\mathcal G_n$ & target rung class of an edge & D1.1 \\
$J^+$, $J^-$ & crystallization seal; dissolution seal & D1.2 \\
$c$, $\gamma$, $s(c)$, $t(c)$, $\partial c$ & chain; source; target; boundary & D1.4 \\
$ds^2(e)$, $\widehat{ds}^{\,2}(e)$ & discrete line element; hop-normalized & D1.6 \\
$\Theta(\gamma)$, $\Phi(\gamma)$, $\operatorname{Hol}_{\mathrm{spiral}}$, $\mathsf{Trivial}$ & spiral sum; phase; holonomy; trivial class & D1.7 \\
$x^0$ & emergent time coordinate & D1.8 \\
$\ne\allowbreak g$, $\wedge$, $\vee$, $\implies$, $\iff$ & negation; conjunction; disjunction; implication; equivalence & — \\
$\oplus$, $\bigoplus$, $\otimes$ & direct sum; indexed direct sum; tensor product & A1, A3, A8 \\
$F_G$, $F_H$, $F_{D^+}$, $F_{D^-}$ & spin-2; scalar (Higgs); uniform dark; saturated dark & A8, D4.1, D4.3 \\
$P_M$, $\mathcal P_M^{(n)}$, $\mathbf F_M^{(n)}$ & mask projection; rung mask; masked force & A8, C7, D2.3 \\
$w(C)$, $\mathcal J(C)$, $G(e)$ & configuration weight; Janus factor; edge transfer weight & A8, D3.5, D3.6 \\
$f_v$, $\aleph_v$, $\aleph(\sigma,\allowbreak C)$ & vertex transition amplitude; indicator; count & D3.3 \\
$Z(\mathcal R)$, $Z(\mathcal R,\allowbreak n)$ & partition function; masked resolvent & A8, D2.4 \\
$\iota_n$ & aspect realization $\mathfrak u(n)\to\mathcal K$ (slot$\to$aspect bridge) & D2.12 \\
$\mathcal P^{(\mathrm{EM})}$, $Y_{\mathrm{asp}}$, $Y_{\mathrm{slot}}$, $Y_{\mathrm{car}}$ & EM block projection; the three realizations of $Y$ & EM.2, EW.8, D2.12 \\
$\varkappa$, $\kappa$, $G_N$ & embedding scale; gravitational coupling; Newton constant & C4, D2.1, D2.9, CT2.9 \\
$\Sigma$, $\Sigma_{\mathrm{tr}}$ & operator $\sum_i\sigma_i\otimes\sigma_i$; scalar $\delta^{ij}q_{ij}$ & C5, D2.1 \\
$a_{00}$, $q_{ij}$, $w_i$ & sector coefficients & D2.1 \\
$\hat h_{ij}$, $h$, $N_i$, $\Omega$, $\Upsilon$ & traceless shear; trace mode; shift; conformal factor; potential & D2.1, D2.2 \\
$\bar g_{\mu\nu}$, $g_{\mu\nu}$ & auxiliary metric; emergent metric & D2.1 \\
$\alpha_\Omega$, $\beta$ & metric embedding constants & D2.1 \\
$\mathcal K$, $\mathcal K^{(A)}$, $\mathcal K_n$ & aspect space; aspect block; rung-$n$ block & D2.3 \\
$m^2$ & mass squared (pole displacement) & D2.5 \\
$\omega$, $\Theta_{\mathfrak L}$ & orientation reversal; ledger reflection & D2.6, D2.7 \\
$\mathfrak M$, $\operatorname{Str}$, $[\cdot,\allowbreak \cdot\}$ & supermatrix; supertrace; graded bracket & D2.10 \\
$\mathcal H^{(n)}_{\mathrm{slot}}$ & slot space $\cong\mathbb C^n$ & D3.1 \\
$q_{\mathrm{EM}}(c)$ & chain end-charge & D3.4 \\
$\alpha_e$, $\alpha_{pq}$ & per-edge dynamic phase $\in\allowbreak \Phi$; branch phase of a crystallizing pair & A4, A7, D3.5 \\
$\alpha^{(n)}_{\mathrm{EM}}$ & EM-block spectral read at rung $n$ & EM.17 \\
$\mathcal P(\Sigma)$, $M(\mathcal P)$ & alignment pattern; mirror image of the pattern & Def-EM.21, Def-EM.22 \\
$q^i$, $r_i$ & colour triplet; antitriplet & D3.8 \\
$B$, $M$ & baryon; meson & D3.9 \\
$W^\pm$, $Z$, $A$, $Y$, $\theta_W$ & weak eigenstates; photon direction; hypercharge (realizations: EW.8); mixing angle & D3.10, EW.8, EW.24 \\
$T_a$, $f_{abc}$ & colour generators; colour structure constants & CT3.9, CT3.21 \\
$\varepsilon_{ijk}$, $\delta^i{}_j$ & colour invariant tensors & CT3.17 \\
$g$, $g'$, $e$, $h_0$ & weak, hypercharge, EM couplings; Higgs background & EW.33, P4 \\
$F_n$ & rung-$n$ saturation block, $n\ge\allowbreak 4$ & A8, D4.1 \\
$\Sigma$, $\mathcal M(\Sigma)$, $C|_\Sigma$ & subnetwork, its microstate space, configuration restriction & D2.13 \\
$\partial_{\mathrm{col}}$, $\mathcal C(\Sigma;C)$, $\Pi_{\mathcal C}$ & colour boundary, contraction completeness, its projector & D2.14 \\
$\kappa_e$, $\Theta(c)$, $\bar\kappa$ & per-edge line-element quantum, loop total, mean quantum & CT1.29a \\
\end{longtable}
}

\phantomsection
\section[§E — Calibration policy]{\texorpdfstring{§E — Calibration policy}{SE - Calibration policy}}

\phantomsection\label{tag:P1}
\textbf{P1.} $\forall$ quantities of the corpus: derived from A1–A9 $\wedge$ C1–C13 $\wedge$ D0.1–D4.3, except the hop length $\mathfrak h$; $[\mathfrak h]=\mathfrak q^{-1/2}$; $\mathfrak h$ fixes the absolute scale.

\phantomsection\label{tag:P2}
\textbf{P2.} Derived, zero freedom:

%dsp-level:0
\[
(\dim\mathcal E_{++},\dim\mathcal E_{+-},\dim\mathcal E_{-+},\dim\mathcal E_{--})=(7,3,3,3);\quad
C_{\mathrm{cap}}=16;\quad
D_{\mathrm{tick}}(n)=2^{n+1};\quad
n_{\max}=3;
\]

%dsp-level:0
\[
N_{\text{colour}}=3;\quad
N_{\text{gluon}}=8;\quad
N_{\gamma\text{-pol}}=2;\quad
q_{\mathrm{EM}}\in\{-1,0,+1\};\quad
\alpha_e\in\{0,\pi\}\ \text{per edge};
\]

%dsp-level:0
\[
m_W=m_Z\cos\theta_W;\quad
e=g\sin\theta_W;\quad
S=n_{\mathcal R}\ln2;\quad
w_{D^+}=-1 .
\]

\textbf{Audit note.} Every entry of this list is re-verified from first principles in exact arithmetic: \texttt{checks/\allowbreak{}verify\_\allowbreak{}merge\_\allowbreak{}tierB.py} §F (totals and cross-references; the identity-level checks live in \texttt{verify\_\allowbreak{}B.py}–\texttt{verify\_\allowbreak{}G.py}).

\phantomsection\label{tag:P3}
\textbf{P3.} Named combinations of $\mathfrak h$ (external comparison only; $\neg$ inputs):

%dsp-level:0
\[
G_N=\frac{\mathfrak h^2}{4\ln2};\qquad
\kappa^2=8\pi G_N=\frac{2\pi\,\mathfrak h^2}{\ln2}\qquad[\mathrm{CT2.9}].
\]

\phantomsection\label{tag:P4}
\textbf{P4.} Free internal configuration parameters ($\neg$ constants of the theory). Independent core: $g,\ g';\ h_0;\ c_W;\ \Lambda;\ \varkappa,\ \beta,\ \alpha_\Omega$. Functionally dependent members (properties of a solution, $\neg$ further freedom): $\theta_W(g,g')$ [EW.24]; $e(g,g')$ [EW.35]; $\mu_\star^2(c_W,\varkappa,\beta,h_0)$ [EW.33]; $m_W,\ m_Z$ [EW.20, EW.23]; $m_H$ [EW.27]. Predictions $:=$ dimensionless ratios among derived quantities at fixed core data.

\phantomsection\label{tag:P5}
\textbf{P5.} $\neg\exists$ quantity identified with a measured external constant. Comparison procedure: fix $\mathfrak h$ once; read P2 $\wedge$ P3; result $:=$ test of A9, $\neg$ input to A1–A8. \textbf{Era-extension note.} The P4 core lists the bosonic-era free internal parameters (Part 2 loci). The fermion-era parameters (the Janus rate ratio $r=x_+/x_-$; the defect residues $\hat\mu^2$, $\hat\mu'^2$; the Ansatz-B formal temperature $\alpha$; the channel scales $\mu$, $\mu_\nu$; the formal fugacities) are registered in Part 3 §Cal-F under this policy; the baryomesonic era introduces none (Part 4 §Cal-BM).

\phantomsection
\section[§F — Foundations: Bookkeeping, Pair Algebra, Ledger Dynamics]{\texorpdfstring{§F — Foundations: Bookkeeping, Pair Algebra, Ledger Dynamics}{SF - Foundations: Bookkeeping, Pair Algebra, Ledger Dynamics}}

\textit{Dep.: §A–§E of this part. Notation: as §A–§C. Machine: \texttt{checks/\allowbreak{}verify\_\allowbreak{}B.py}, \texttt{checks/\allowbreak{}verify\_\allowbreak{}L\_\allowbreak{}lyapunov.py}, \texttt{checks/\allowbreak{}verify\_\allowbreak{}L0\_\allowbreak{}ledger\_\allowbreak{}tilt.py}, \texttt{checks/\allowbreak{}verify\_\allowbreak{}merge\_\allowbreak{}tierB.py}, \texttt{checks/\allowbreak{}verify\_\allowbreak{}promotion\_\allowbreak{}B\_\allowbreak{}rung.py}.}

\phantomsection
\section[Tier 0 — Bookkeeping, Reality, and Capacity]{\texorpdfstring{Tier 0 — Bookkeeping, Reality, and Capacity}{Tier 0 - Bookkeeping, Reality, and Capacity}}

\begin{theorem}[{\texorpdfstring{CT0.1 — Real/virtual dichotomy}{CT0.1 - Real/virtual dichotomy}}]\label{tag:CT0-1}
\textbf{Thm.} $\ \forall(p,q)\in\mathcal Q\times\mathcal Q:\quad\bigl[(p,q)\in\mathcal R\bigr]\ \underline\vee\ \bigl[(p,q)\in\mathcal V\bigr].$

\textbf{Req.} A1, A3, A4; D0.1, D0.2, D0.4.

\textbf{Pf.}

%dsp-level:0
\[
\bigl|\Pi_{\mathcal V}\rho_{pq}\bigr|=\mathfrak T\bigl(\Pi_{\mathcal V}\rho_{pq}\bigr)\in[0,1]\subset\mathbb R
\]

%dsp-level:0
\[
\quad [\mathrm{D0.2},\ \mathrm{A1}:\ \rho_{pq}\succeq0\wedge\mathfrak T\rho_{pq}=1;\ \mathrm{A3}:\ \Pi_{\mathcal V}=\Pi_{\wedge^2}]
\]

%dsp-level:0
\[
x\in\mathbb R\implies\bigl[(x\le\varepsilon_q\ \vee\ x\ge1-\varepsilon_q)\ \underline\vee\ (\varepsilon_q<x<1-\varepsilon_q)\bigr]\quad[\text{total order of }\mathbb R];
\]

%dsp-level:0
\[
\qquad \varepsilon_q\in[0,1]\subset\mathbb R\ [\mathrm{D0.4}]
\]

%dsp-level:0
\[
\bigl[(p,q)\in\mathcal R\bigr]\iff\bigl[\bigl|\Pi_{\mathcal V}\rho_{pq}\bigr|\le\varepsilon_q\ \vee\ \bigl|\Pi_{\mathcal V}\rho_{pq}\bigr|\ge1-\varepsilon_q\bigr]\ [\mathrm{A4}];
\]

%dsp-level:0
\[
\qquad \mathcal V=(\mathcal Q\times\mathcal Q)\setminus\mathcal R\ [\mathrm{D0.1}]
\]

%dsp-level:0
\[
\Longrightarrow\quad\bigl[(p,q)\in\mathcal R\bigr]\ \underline\vee\ \bigl[(p,q)\in\mathcal V\bigr].\ \blacksquare
\]

\textbf{Phys.} Every qubit pair has exactly one status — crystallized (real) or uncrystallized (virtual) — decided by A4's two-branch threshold test on its antisymmetric weight. This dichotomy is the bookkeeping ground of the corpus: real pairs become edges of the emergent graph, virtual pairs remain vacuum.

\end{theorem}

\begin{theorem}[{\texorpdfstring{CT0.2 — Global vacuum}{CT0.2 - Global vacuum}}]\label{tag:CT0-2}
\textbf{Thm.} $\ n_{\mathcal R}=0\ \implies\ \forall(p,q):\ \bigl[(p,q)\in\mathcal V\ \wedge\ \rho_{pq}=\rho_\circ=\tfrac14 I\bigr]\ \wedge\ S_{\mathrm{raw}}(\rho_\circ)=\ln4=\max_{\rho}S_{\mathrm{raw}}(\rho),$ with $\arg\max=\{\tfrac14 I\}$.

\textbf{Req.} A1, A4, A5; D0.1, D0.3; CT0.1; Gibbs lemma: $\ S_{\mathrm{raw}}(\rho)\le\ln d$, $\ S_{\mathrm{raw}}(\rho)=\ln d\iff\rho=I/d$ [linear algebra].

\textbf{Pf.}

%dsp-level:0
\[
n_{\mathcal R}=0\implies\neg\,\exists(p,q)\in\mathcal R\ [\mathrm{D0.1}]\implies\forall(p,q):\ (p,q)\in\mathcal V\ [\mathrm{CT0.1}]
\]

%dsp-level:0
\[
\rho_{pq}=\rho_\circ=\tfrac14 I\quad[\mathrm{A5}]
\]

%dsp-level:0
\[
S_{\mathrm{raw}}(\rho_\circ)=-\mathfrak T\bigl(\tfrac14 I\ln\tfrac14 I\bigr)=-4\cdot\tfrac14\ln\tfrac14=\ln4\quad[\mathrm{D0.3}]
\]

%dsp-level:0
\[
\mathfrak H(\mathcal H_p\otimes\mathcal H_q)\cong\mathfrak H(\mathbb C^4),\ d=4:\quad S_{\mathrm{raw}}(\rho)\le\ln4,\quad S_{\mathrm{raw}}(\rho)=\ln4\iff\rho=\tfrac14 I\quad[\mathrm{A1},\ \text{Gibbs}].\ \blacksquare
\]

\textbf{Phys.} The empty configuration (no real edges) is all-virtual, with every pair at the maximally mixed state $I/4$ and entropy maximal at $\ln 4$: the vacuum of this theory is a state of maximal raw information. All structure — edges, particles, geometry — is a lowering of this entropy by crystallization.

\end{theorem}

\begin{theorem}[{\texorpdfstring{CT0.3 — Seal increments}{CT0.3 - Seal increments}}]\label{tag:CT0-3}
\textbf{Thm.} $\ J^{+}:\ (\Delta n_{\mathcal R},\ \delta h,\ \delta\tau)=(+1,+1,+1);\qquad J^{-}:\ (\Delta n_{\mathcal R},\ \delta h,\ \delta\tau)=(-1,-1,-1).$

\textbf{Req.} A4, A7; D0.1, D0.6, D0.7, D1.2, D1.9.

\textbf{Pf.}

%dsp-level:0
\[
J^\pm\ [\mathrm{D1.2}]:\ \ \exists!\,(p,q):\ \Delta\bigl[(p,q)\in\mathcal R\bigr]=\pm1\ [\mathrm{A4}]\implies\Delta n_{\mathcal R}=\pm1\ [\mathrm{D0.1}]
\]

%dsp-level:0
\[
\Delta\bigl|E_{\mathcal R}\bigr|=\Delta n_{\mathcal R}\ [\mathrm{D0.1}]\implies\delta h=\pm1\ [\mathrm{D0.6}:\ \delta h=\#\{e\in E_{\mathcal R}\ \text{traversed}\}]
\]

%dsp-level:0
\[
\text{seal}=\text{retained update along a real edge}\ [\mathrm{D0.7}]\implies\delta\tau=\pm1\ \text{(same sign)}
\]

%dsp-level:0
\[
\mathrm{D1.9}:\ \operatorname{dom}(\mathrm{tick})\subseteq\{v\}\cup\{e\ni v\}\implies\Delta(\text{all further counters})=0.\ \blacksquare
\]

\textbf{Phys.} Crystallization ($J^+$) and dissolution ($J^-$) are the atomic ledger events: each adds or removes exactly one real edge, one traversal unit, and one tick of ledger time. Every extensive quantity of the corpus moves in integer steps of these events.

\end{theorem}

\begin{theorem}[{\texorpdfstring{CT0.4 — Real outgoing edge registration}{CT0.4 - Real outgoing edge registration}}]\label{tag:CT0-4}
\textbf{Thm.} $\ e=(p,q):\quad e\in\mathcal R_{\rightarrow}(p)\iff e\longmapsto(\mathcal G_n,\phi_e).$

\textbf{Req.} A7; D0.5, D0.8, D1.1; CT0.1.

\textbf{Pf.}

%dsp-level:0
\[
(\Rightarrow):\ \ e\in E_{\mathcal R}\implies\exists\,\phi_e\ [\mathrm{A7}];\qquad q\in V_n\ \text{for}\ n=n(q)\ [\mathrm{D0.8}]\implies\exists\,\mathcal G_n;
\]

%dsp-level:0
\[
\mathrm{D1.1}:\ \ e\longmapsto(\mathcal G_n,\phi_e).\qquad
(\Leftarrow):\ \ \mathrm{D1.1}:\ \neg\bigl[e\longmapsto(\mathcal G_n,\phi_e)\bigr]\implies e\in\mathcal V;
\]

%dsp-level:0
\[
e\longmapsto(\mathcal G_n,\phi_e)\implies e\notin\mathcal V\implies e\in\mathcal R\quad[\mathrm{CT0.1}].\ \blacksquare
\]

\textbf{Phys.} “Real” and “registered in the ledger” are the same property of an edge: real edges are exactly those carrying ledger data (target rung class and phase); unregistered edges are virtual. Existence, in this corpus, is recordability.

\end{theorem}

\begin{theorem}[{\texorpdfstring{CT0.5 — Tick load at rung $n$}{CT0.5 - Tick load at rung}}]\label{tag:CT0-5}
\textbf{Thm.} $\ D_{\mathrm{tick}}(n)=2^{\,n+1}.$

\textbf{Req.} A7; C10; D0.7, D0.8, D0.9, D1.9; CT0.9.

\textbf{Pf.}

%dsp-level:0
\[
v\in V_n\implies\bigl|\{e\in E_{\mathcal R}:\ e\ni v\}\bigr|=n\quad[\mathrm{D0.8}]
\]

%dsp-level:0
\[
\phi_e\in\Phi,\ |\Phi|=2\ [\mathrm{CT0.9}]\implies(\phi_e)_{e\ni v}\in\Phi^{\,n},\quad\bigl|\Phi^{\,n}\bigr|=2^n
\]

%dsp-level:0
\[
\mathrm{D1.9}:\ \ \mathrm{tick\ datum}=(\mathrm{pre},\mathrm{post})\in\Phi^{\,n}\times\Phi^{\,n},\quad \mathrm{pre}\perp\mathrm{post}
\]

%dsp-level:0
\[
\mathrm{C10}/\mathrm{D0.9}\ (\text{register accounting}):\ \ \text{registration of}\ X\ \text{at single-shot linear readout}
\]

%dsp-level:0
\[
= |X|\ \text{orthogonal register directions}
\]

%dsp-level:0
\[
\text{Shannon content}=\log_2|X|=n\ \text{bits}\ \ne\ \text{register cost}\ 2^n\quad(\text{source availability}\ \ne\ \text{throughput})
\]

%dsp-level:0
\[
\Longrightarrow\quad D_{\mathrm{tick}}(n)=2^n+2^n=2^{\,n+1}.\ \blacksquare
\]

\textbf{Phys.} A vertex with $n$ real edges must make $2^{n+1}$ classical alternatives linearly distinguishable per tick — $2^n$ pre-tick and $2^n$ post-tick phase configurations (C10 register accounting: one orthogonal direction per admissible value, not $\log_2$ bit encoding). This exponential growth of source availability is what caps accessible complexity at rung 3 (CT0.11).

\end{theorem}

\begin{theorem}[{\texorpdfstring{CT0.6 — Channel capacity and oversubscription}{CT0.6 - Channel capacity and oversubscription}}]\label{tag:CT0-6}
\textbf{Thm.}

%dsp-level:0
\[
C_{\mathrm{cap}}=\dim_{\mathbb R}\mathfrak H(\mathbb C^4)=16;\qquad
O(n)=\max\bigl(0,\ 2^{\,n+1}-16\bigr);
\]

%dsp-level:0
\[
O(0)=O(1)=O(2)=O(3)=0;\qquad O(4)=16;\qquad O(n)=2^{\,n+1}-16\quad(n\ge4).
\]

\textbf{Req.} A1, A2; D0.9; CT0.5; $\dim_{\mathbb R}\mathfrak H(\mathbb C^d)=d^2$ [linear algebra].

\textbf{Pf.}

%dsp-level:0
\[
\mathcal H_p\otimes\mathcal H_q\cong\mathbb C^2\otimes\mathbb C^2\cong\mathbb C^4\ [\mathrm{A1}];\qquad
\dim_{\mathbb R}\mathfrak H(\mathbb C^d)=d^2\ [\text{linear algebra}]
\]

%dsp-level:0
\[
C_{\mathrm{cap}}=\dim_{\mathbb R}\mathfrak H(\mathbb C^4)=4^2=16\ [\mathrm{D0.9}];\qquad C(p,q)=1\,\mathfrak q\ \ \forall(p,q)\ [\mathrm{A2}]
\]

%dsp-level:0
\[
O(n)=\max\bigl(0,\ D_{\mathrm{tick}}(n)-C_{\mathrm{cap}}\bigr)=\max\bigl(0,\ 2^{\,n+1}-16\bigr)\quad[\mathrm{D0.9},\ \mathrm{CT0.5}]
\]

%dsp-level:0
\[
2^{\,n+1}\le16\iff n\le3\implies O(n)=0\ \ (n\le3);
\]

%dsp-level:0
\[
\qquad 2^{5}-16=16\implies O(4)=16;
\]

%dsp-level:0
\[
\qquad n\ge4\implies2^{\,n+1}\ge32>16\implies O(n)=2^{\,n+1}-16.\ \blacksquare
\]

\textbf{Phys.} One qubit–qubit channel carries exactly 16 real registers per tick — the dimension of the two-qubit Hermitian operator space. Rungs $0$–$3$ fit the budget, rung $3$ exactly saturates it, and rung $4$ doubles the demand to $32$, exceeding it by exactly one channel quantum. All finiteness of the accessible world starts here.

\end{theorem}

\begin{corollary}[{\texorpdfstring{CT0.6a — transit-hop register bound}{CT0.6a - transit-hop register bound}}]\label{tag:CT0-6a}
Let $\gamma=(e_1,\dots,e_{\delta h})$ be a transit chain [D1.5, D1.4] of distinct directed real edges with hop count $\delta h(\gamma)$ [D0.6] whose entire edge-register content is registered through one pair channel over $\delta\tau(\gamma)$ ledger ticks [A2, D0.7]. Then:

%dsp-level:0
\[
\mathsf E(\gamma)=\delta h(\gamma)\,\ln2\ \ \text{nats}\qquad(\delta h\ \ \text{independent binary registers}\ [\mathrm{CT0.9},\ \mathrm{C9}]);
\]

%dsp-level:0
\[
\text{single-channel supply over}\ \delta\tau\ \text{ticks}:\ \log_2C_{\mathrm{cap}}\cdot\delta\tau=4\,\delta\tau\ \ \text{bits}
\]

%dsp-level:0
\[
\qquad [\mathrm{CT0.6}:\ C_{\mathrm{cap}}=16;\ \mathrm{C10}:\ 16\ \text{directions}\iff2^{4}\ \text{alternatives}]
\]

%dsp-level:0
\[
\implies\quad\delta h(\gamma)\ \le\ 4\,\delta\tau(\gamma);\qquad
\mathsf E(\gamma)\ \le\ \min\bigl(\delta h(\gamma),\ 4\,\delta\tau(\gamma)\bigr)\,\ln2\ \ \text{nats}.
\]

A per-tick payload beyond $4$ hop-bits presents $2^{5}=32>16=C_{\mathrm{cap}}$ alternatives at the crossing — oversubscribed [CT0.6]; this is the payload counterpart of CT0.11's vertex-rung admissibility. The bound is sharp: for $\delta h\le4\,\delta\tau$ a schedule of $\delta\tau$ consecutive groups of $\le4$ hops presents $\le2^{4}=16=C_{\mathrm{cap}}$ alternatives per tick — at capacity, no oversubscription [CT0.6]. (Ledger accounting of registered content vs channel supply over a tick window [D0.7]; $\neg$a tick-to-tick transition law.)

\textbf{Req (Cor).} A2; C9, C10; D0.6, D0.7, D1.4, D1.5; CT0.6, CT0.9; exact binary capacity arithmetic [math; machine: \texttt{checks/\allowbreak{}verify\_\allowbreak{}promotion\_\allowbreak{}B\_\allowbreak{}rung.py} §B].

\textbf{Pf (Cor).}

%dsp-level:0
\[
\mathrm{CT0.9}:\ \sigma_e\in\{\pm1\}:\ 1\ \text{real edge}=1\ \text{bit}=\ln2\ \text{nats};
\]

%dsp-level:0
\[
\qquad \text{distinct edges}=\text{independent registers}\implies\mathsf E(\gamma)=\delta h\ln2\ [\mathrm{C9}]
\]

%dsp-level:0
\[
\mathrm{CT0.6}:\ C_{\mathrm{cap}}=16\ \text{register directions per channel crossing};
\]

%dsp-level:0
\[
\qquad \mathrm{C10}:\ |X|\ \text{alternatives}
\]

%dsp-level:0
\[
\mapsto |X|\ \text{directions},
\]

%dsp-level:0
\[
\ \text{Shannon}\ \log_2|X|:\ 16\ \text{directions}
\]

%dsp-level:0
\[
\iff |X|
\]

%dsp-level:0
\[
\le 16\iff
\]

%dsp-level:0
\[
\le 4\ \text{bits per tick}
\]

%dsp-level:0
\[
\delta\tau\ \text{ticks}:\ \text{supply}\le4\,\delta\tau\ \text{bits};
\]

%dsp-level:0
\[
\qquad \text{registered content}\le\text{supply}\implies\delta h\le4\,\delta\tau;
\]

%dsp-level:0
\[
\qquad \mathsf E\le\min(\delta h,\,4\,\delta\tau)\ln2
\]

%dsp-level:0
\[
>4\ \text{bits in one tick}:\ 2^{5}=32>16=C_{\mathrm{cap}}:\ \text{oversubscribed}\ [\mathrm{CT0.6}];
\]

%dsp-level:0
\[
\qquad \text{groups of}\ \le4:\ 2^{4}=16=C_{\mathrm{cap}}:\ \text{at capacity}\ [\mathrm{CT0.6}]\ \blacksquare
\]

\end{corollary}

\begin{theorem}[{\texorpdfstring{CT0.7 — Reality is the swap grading}{CT0.7 - Reality is the swap grading}}]\label{tag:CT0-7}
\textbf{Thm.} Let $\mathbb S\in\mathcal U(\mathcal H_p\otimes\mathcal H_q):\ |a\rangle\otimes|b\rangle\mapsto|b\rangle\otimes|a\rangle$. Then

%dsp-level:0
\[
(1)\quad \mathbb S=\tfrac12\bigl(E_0+\Sigma\bigr);
\]

%dsp-level:0
\[
\qquad \Pi_{\mathcal R}:=\Pi_{\vee^2}=\tfrac14\bigl(3E_0+\Sigma\bigr);
\]

%dsp-level:0
\[
\qquad \Pi_{\mathcal V}:=\Pi_{\wedge^2}=\tfrac14\bigl(E_0-\Sigma\bigr);
\]

%dsp-level:0
\[
\qquad \mathfrak T\Pi_{\mathcal R}=3,\ \ \mathfrak T\Pi_{\mathcal V}=1 .
\]

%dsp-level:0
\[
(2)\quad \forall X:\ S(X)=\mathbb S X\mathbb S .
\]

%dsp-level:0
\[
(3)\quad SX=X\implies[\mathbb S,X]=0\implies X(\vee^2)\subseteq\vee^2\ \wedge\ X(\wedge^2)\subseteq\wedge^2;
\]

%dsp-level:0
\[
SX=-X\implies\{\mathbb S,X\}=0\implies X(\vee^2)\subseteq\wedge^2\ \wedge\ X(\wedge^2)\subseteq\vee^2;
\]

%dsp-level:0
\[
\underbrace{\dim_{\mathbb R}\mathfrak H(\vee^2)+\dim_{\mathbb R}\mathfrak H(\wedge^2)}_{9+1=10}\ \oplus\ \underbrace{2\dim_{\mathbb C}\operatorname{Hom}(\wedge^2,\vee^2)}_{2\cdot3=6}\ =\ 16 .
\]

%dsp-level:0
\[
(4)\quad \mathbf F_{(+,\mathcal R)}=\mathcal E_{++};\quad
\mathbf F_{(-,\mathcal R)}=\mathcal E_{+-};\quad
\mathbf F_{(+,\mathcal V)}=\mathcal E_{-+};\quad
\mathbf F_{(-,\mathcal V)}=\mathcal E_{--}
\]

(first label of $\mathbf F_{(\pm,\cdot)}$: $T$-parity; second label: $S$-parity, $\mathcal R\leftrightarrow+$, $\mathcal V\leftrightarrow-$); $\ F_{J^+}=\mathcal E_{-+}=\operatorname{span}\{J_i\}$, $\ F_{J^-}=\mathcal E_{--}=\operatorname{span}\{K_i\}$: the $S$-odd blocks.

\textbf{Req.} A3, A8; C3, C8; D1.2, D3.10; CT0.6, CT1.1, CT1.2.

\textbf{Pf.}

%dsp-level:0
\[
(1)\ \ \mathbb S=\tfrac12\Bigl(I\otimes I+\textstyle\sum_i\sigma_i\otimes\sigma_i\Bigr)=\tfrac12\bigl(E_0+\Sigma\bigr)\quad[\mathrm{CT1.2}];
\]

basis check: $\ a=b:\ \tfrac12(1+1)|aa\rangle=|aa\rangle$; $\ a\ne b:\ \sum_i\sigma_i\otimes\sigma_i|ab\rangle=2|ba\rangle-|ab\rangle\implies\tfrac12(|ab\rangle+2|ba\rangle-|ab\rangle)=|ba\rangle$.

%dsp-level:0
\[
\mathbb S^2=\mathbb1\ \wedge\ \mathbb S^\dagger=\mathbb S\implies\operatorname{spec}(\mathbb S)=\{\pm1\},\ \ \text{eigenspaces}\ \vee^2,\ \wedge^2\quad[\mathrm{A3}]
\]

%dsp-level:0
\[
\Pi_{\mathcal R}=\mathbb1-\Pi_{\mathcal V}=\tfrac12(\mathbb1+\mathbb S)=\tfrac14\bigl(3E_0+\Sigma\bigr);
\]

%dsp-level:0
\[
\qquad \Pi_{\mathcal V}=\tfrac12(\mathbb1-\mathbb S)=\tfrac14\bigl(E_0-\Sigma\bigr)\quad[\mathrm{A3},\mathrm{A4}]
\]

%dsp-level:0
\[
\mathfrak T E_0=4,\ \mathfrak T\Sigma=0\implies\mathfrak T\Pi_{\mathcal R}=3=\dim\vee^2,\quad\mathfrak T\Pi_{\mathcal V}=1=\dim\wedge^2\quad[\mathrm{CT1.1}]
\]

%dsp-level:0
\[
(2)\ \ \mathbb S(\sigma_\mu\otimes\sigma_\nu)\mathbb S=\sigma_\nu\otimes\sigma_\mu=S(\sigma_\mu\otimes\sigma_\nu)\ [\mathrm{A3}];\quad\text{linearity}\implies S(X)=\mathbb S X\mathbb S\ \ \forall X
\]

%dsp-level:0
\[
(3)\ \ SX=X\Rightarrow\mathbb S X\mathbb S=X\Rightarrow\mathbb S X=X\mathbb S\ (\times\mathbb S\ \text{right});\qquad SX=-X\Rightarrow\mathbb S X=-X\mathbb S
\]

%dsp-level:0
\[
X=X^\dagger\ \wedge\ [\mathbb S,X]=0\implies X\ \text{preserves each}\ \mathbb S\text{-eigenspace};\qquad
\{\mathbb S,X\}=0\implies X:\ \vee^2\leftrightarrow\wedge^2
\]

%dsp-level:0
\[
\dim:\ \ \mathfrak H(\vee^2)\oplus\mathfrak H(\wedge^2)\cong\mathfrak H(\mathbb C^3)\oplus\mathfrak H(\mathbb C^1):\ 9+1=10;
\]

%dsp-level:0
\[
\qquad \operatorname{Hom}(\wedge^2,\vee^2)\cong\mathbb C^{3}:\ 2\cdot3=6;
\]

%dsp-level:0
\[
\qquad 10+6=16\ [\mathrm{CT0.6}]
\]

%dsp-level:0
\[
(4)\ \ S\text{-odd generators}=\{J_i\}\cup\{K_i\}\ (6)\ [\mathrm{CT1.2}];\quad S\text{-even}=\{E_0,\Sigma,\hat E^{ij},W_i\}\ (10)
\]

%dsp-level:0
\[
\mathrm{D1.2}:\ \text{seal moves A4-weight }\wedge^2\!\leftrightarrow\!\vee^2\implies\text{seal generators}
\]

%dsp-level:0
\[
= S\text{-odd}\ [(3)]
\]

%dsp-level:0
\[
\implies F_{J^\pm}
\]

%dsp-level:0
\[
= S\text{-odd blocks}\ [\mathrm{A8}]
\]

%dsp-level:0
\[
T\text{-parity}:\ J_i:+,\ K_i:-\ [\mathrm{CT1.2}]\implies F_{J^+}=\mathcal E_{-+},\ F_{J^-}=\mathcal E_{--}
\]

%dsp-level:0
\[
\mathcal E_{++}\supseteq\{E_0,\Sigma,\hat E^{ij}\}\mapsto\{F_{D^+},F_H,F_G\};\qquad
\mathcal E_{+-}=\operatorname{span}\{W_i\}\mapsto\{F_{W^+},F_{W^-},F_Z\}\quad[\mathrm{A8},\ \mathrm{C8},\ \mathrm{D3.10}].\ \blacksquare
\]

\end{theorem}

\begin{corollary}[{\texorpdfstring{CT0.7a — spin content}{CT0.7a - spin content}}]\label{tag:CT0-7a}
$\ \sum_iW_i^2=\tfrac32E_0+\tfrac12\Sigma$; $\ \sum_iW_i^2\big|_{\vee^2}=2=1(1{+}1)$; $\ \sum_iW_i^2\big|_{\wedge^2}=0=0(0{+}1)$; i.e. $\vee^2$: spin $1$, $\wedge^2$: spin $0$.

\textbf{Req (Cor).} CT0.7, CT1.2.

\textbf{Pf (Cor).}

%dsp-level:0
\[
\textstyle\sum_iW_i^2=\tfrac14\sum_i(\sigma_i\otimes\sigma_0+\sigma_0\otimes\sigma_i)^2=\tfrac14\bigl(3E_0+3E_0+2\Sigma\bigr)=\tfrac32E_0+\tfrac12\Sigma
\]

%dsp-level:0
\[
\Sigma=2\mathbb S-E_0\ [\mathrm{CT0.7}(1)];
\]

%dsp-level:0
\[
\qquad \vee^2:\ \mathbb S=+1\implies\Sigma=+1\implies\tfrac32+\tfrac12=2;
\]

%dsp-level:0
\[
\qquad \wedge^2:\ \mathbb S=-1\implies\Sigma=-3\implies\tfrac32-\tfrac32=0.\ \blacksquare
\]

\textbf{Phys.} The real/virtual split of A4 \textit{is} the exchange (swap) grading of the pair: real $=$ symmetric $=$ spin-1 triplet, virtual $=$ antisymmetric $=$ spin-0 singlet. The A8 force labels $(\pm,\mathcal R/\mathcal V)$ are precisely the two parities $(T,S)$; the Janus seals are exactly the operators that exchange the two grades, i.e. that create and destroy reality. “Existence” in this corpus is Bose exchange symmetry of a qubit pair.

\end{corollary}

\begin{theorem}[{\texorpdfstring{CT0.8 — Admissible range of the crystallization threshold}{CT0.8 - Admissible range of the crystallization threshold}}]\label{tag:CT0-8}
\textbf{Thm.} $\ \bigl[\mathrm{A4}\wedge\mathrm{A5}\wedge\mathrm{D0.1}\ \text{consistent}\bigr]\implies 0\le\varepsilon_q<\tfrac14;$ $\ \ $in particular $\varepsilon_q<1-\varepsilon_q$ (A4's two branches are disjoint); $\qquad\varepsilon_q=0\implies\bigl[(p,q)\in\mathcal R\iff\operatorname{supp}\rho_{pq}\subseteq\vee^2\ \ \vee\ \ \rho_{pq}=\Pi_{\mathcal V}\bigr]$ (branch reads $\alpha_{pq}=0\ /\ \pi$ [A4]); and $\forall$ theorems of the corpus: A4 enters only via CT0.1 (valid $\forall\varepsilon_q\in[0,1]$) and CT0.2 (valid $\iff\varepsilon_q\in[0,\tfrac14)$).

\textbf{Req.} A1, A4, A5; D0.1, D0.2, D0.4; CT0.2, CT0.7.

\textbf{Pf.}

%dsp-level:0
\[
\bigl|\Pi_{\mathcal V}\rho_\circ\bigr|=\mathfrak T\bigl(\Pi_{\mathcal V}\cdot\tfrac14 I\bigr)=\tfrac14\,\mathfrak T\Pi_{\mathcal V}=\tfrac14\quad[\mathrm{A5},\ \mathrm{CT0.7}(1):\ \mathfrak T\Pi_{\mathcal V}=1]
\]

%dsp-level:0
\[
\mathrm{D0.1}:\ \text{vacuum}\iff n_{\mathcal R}=0\implies\rho_\circ\notin\mathcal R\ [\mathrm{CT0.2}]\implies\tfrac14>\varepsilon_q\quad[\mathrm{A4}]
\]

%dsp-level:0
\[
\varepsilon_q\ge0\ [\mathrm{D0.4}];\qquad\bigl|\Pi_{\mathcal V}\rho\bigr|\ge0\ [\mathrm{D0.2}]\implies\varepsilon_q=0\ \text{admissible}
\]

%dsp-level:0
\[
\text{branch 2 at the vacuum}:\ \tfrac14
\]

%dsp-level:0
\[
\ge 1-\varepsilon_q\iff\varepsilon_q\ge\tfrac34:\ \text{excluded by}\ \varepsilon_q<\tfrac14\implies\rho_\circ\ \text{virtual in both branches}
\]

%dsp-level:0
\[
\text{branch disjointness}:\ 0\le\varepsilon_q<\tfrac14\implies\varepsilon_q<\tfrac12<1-\varepsilon_q\implies[0,\varepsilon_q]\cap[1-\varepsilon_q,1]=\varnothing
\]

%dsp-level:0
\[
\varepsilon_q=0:\ \ \mathfrak T\bigl(\Pi_{\mathcal V}\rho_{pq}\bigr)=0\ \wedge\ \Pi_{\mathcal V}\succeq0\ \wedge\ \rho_{pq}\succeq0\iff\Pi_{\mathcal V}\rho_{pq}\Pi_{\mathcal V}=0\iff\operatorname{supp}\rho_{pq}\subseteq\vee^2\quad[\mathrm{A1}]
\]

%dsp-level:0
\[
\varepsilon_q
\]

%dsp-level:0
\[
= 0,
\]

%dsp-level:0
\[
\ \text{branch 2}:\ \ \mathfrak T\bigl(\Pi_{\mathcal V}\rho_{pq}\bigr)
\]

%dsp-level:0
\[
= 1\ \wedge\ \mathfrak T\rho_{pq}
\]

%dsp-level:0
\[
= 1\ \wedge\ \rho_{pq}\succeq0\ \wedge\ \operatorname{rk}\Pi_{\mathcal V}
\]

%dsp-level:0
\[
= 1\ \ [\mathrm{CT0.7}(1)]\iff\rho_{pq}\qquad = \Pi_{\mathcal V}
\]

%dsp-level:0
\[
\text{A4-dependence of all later tags}=\{\mathrm{CT0.1},\ \mathrm{CT0.2}\}\implies\text{range}\ [0,\tfrac14)\ \text{suffices; sharp point}\ \varepsilon_q=0.\ \blacksquare
\]

\textbf{Phys.} The crystallization threshold is not a free constant: demanding that the vacuum be all-virtual confines it to $0\le\varepsilon_q<\tfrac14$ — which simultaneously keeps A4's two branches disjoint ($\varepsilon_q<1-\varepsilon_q$). The sharp choice $\varepsilon_q=0$ says a pair is real exactly when its antisymmetric weight is exactly zero (aligned branch, $\alpha_{pq}=0$) or the pair is the pure singlet (antialigned branch, $\alpha_{pq}=\pi$). No result depends on where inside the range the tolerance sits — the theory is tolerance-independent.

\end{theorem}

\begin{theorem}[{\texorpdfstring{CT0.9 — The edge phase alphabet is binary}{CT0.9 - The edge phase alphabet is binary}}]\label{tag:CT0-9}
\textbf{Thm.} $\ \phi_e\in\Phi=\{0,\pi\}\ (\mathrm{mod}\ 2\pi);\qquad\sigma_e=e^{i\phi_e}=(-1)^{\phi_e/\pi}\in\{+1,-1\};\qquad 1\ \text{real edge}=1\ \text{bit}=\ln2\ \text{nats}.$

\textbf{Req.} A1, A3, A4, A7; C9, C10; D0.5, D1.1; $\dim\wedge^2=\binom{d}{2}$ [linear algebra].

\textbf{Pf.}

%dsp-level:0
\[
e\in E_{\mathcal R}\mapsto\phi_e\ [\mathrm{A7}];
\]

%dsp-level:0
\[
\qquad \mathrm{D1.1}:\ \phi_e:=\text{phase of the registration transport's action on the reality line}\ \wedge^2
\]

%dsp-level:0
\[
\mathrm{A1}:\ \dim\mathcal H_q=d=2\implies\dim\wedge^2=\tbinom{2}{2}=1\quad(\text{the reality line is 1-dimensional}\iff\text{qubits are}\ \mathbb C^2)
\]

%dsp-level:0
\[
\mathrm{A4}:\ \text{reality test}
\]

%dsp-level:0
\[
= \text{extremality of the}\ \wedge^2\ \text{weight (either branch)}\implies\text{registration transport preserves the}\ (S,T)\text{-grading}\ [\mathrm{A3},\ \mathrm{D1.1}]
\]

%dsp-level:0
\[
\mathrm{D1.1}/\mathrm{C10}\ \text{reality principle}:\ \text{transport preserves the real frame structure}\implies\text{action on}\ \wedge^2_{\mathbb R}\ (1\text{-dim real line})
\]

%dsp-level:0
\[
= \text{multiplication by}\ c\in\mathbb R
\]

%dsp-level:0
\[
\text{unitarity}\implies|c|=1\implies c\in\{+1,-1\}\implies e^{i\phi_e}\in\{\pm1\}\implies\phi_e\in\{0,\pi\}\ (\mathrm{mod}\ 2\pi);
\]

%dsp-level:0
\[
\qquad \sigma_e=(-1)^{\phi_e/\pi}\ [\mathrm{D0.5}]
\]

%dsp-level:0
\[
|\Phi|=2\implies\log_2|\Phi|=1\ \text{bit}=\ln2\ \text{nats}\quad[\mathrm{C9}].\ \blacksquare
\]

\textbf{Rmk (repair).} The prior proof argued from the centralizer of $\{S,T\}$ alone — insufficient, since the complex centralizer of the swap is $U(3)\times U(1)$ and contains continuous phases. The binary alphabet follows once the reality principle of D1.1/C10 (registrations are real, grading-preserving frame data — quantum data classicalizes at registration) is explicit. The continuous counterexamples are complex-transport artifacts excluded by that principle.

\textbf{Phys.} The only internal datum a real edge can carry is a single sign — and this is \textit{qubit-specific}: the reality line $\wedge^2$ is one-dimensional exactly because the carriers are $\mathbb C^2$ ($\dim\wedge^2=\binom d2=1\iff d=2$); for higher-dimensional carriers continuous grading-preserving phases would survive. Every edge of the emergent graph stores exactly one bit; all later charge and phase structure is built from these signs.

\end{theorem}

\begin{theorem}[{\texorpdfstring{CT0.10 — The two sign structures are independent}{CT0.10 - The two sign structures are independent}}]\label{tag:CT0-10}
\textbf{Thm.} With $\ P(C):=\prod_{e\in E(C)}\sigma_e\ $ and $\ \pi(C):=(-1)^{\aleph(\sigma,C)}$ :

%dsp-level:0
\[
\neg\bigl[P(C)\ \text{determines}\ \pi(C)\bigr]\ \wedge\ \neg\bigl[\pi(C)\ \text{determines}\ P(C)\bigr];\qquad
L\ge4:\ \ (P,\pi)\twoheadrightarrow\{\pm1\}^2;
\]

%dsp-level:0
\[
\Longrightarrow\quad w(C)=\mathcal J(C)\prod_{e\in E(C)}G(e)\cdot(-1)^{\aleph(\sigma,C)}\ \ [\mathrm{A8}]\ \text{double-counts no edge-sign data}.
\]

\textbf{Req.} A8; D0.5, D3.3.

\textbf{Pf.}

%dsp-level:0
\[
\aleph(\sigma,C)=\#\{i<L:\ \sigma_i\ne\sigma_{i+1}\}\quad[\mathrm{D3.3}]\ :\ \text{function of consecutive differences only};
\]

%dsp-level:0
\[
\qquad P=\prod_{i}\sigma_i\ :\ \text{function of the signs}
\]

%dsp-level:0
\[
\text{witness}_1\ [\pi\ \text{equal},\ P\ \text{opposite}]:\quad(+,-,+):\ \aleph=2,\ P=-1;\qquad(-,+,-):\ \aleph=2,\ P=+1
\]

%dsp-level:0
\[
\text{witness}_2\ [P\ \text{equal},\ \pi\ \text{opposite}]:\quad(+,+,+,+):\ P=+1,\ \aleph=0;\qquad(+,+,-,-):\ P=+1,\ \aleph=1
\]

%dsp-level:0
\[
\Longrightarrow\quad P\not\Rightarrow\pi\ \wedge\ \pi\not\Rightarrow P;
\]

%dsp-level:0
\[
\qquad L\ge4:\ \{(P,\pi)\}=\{(+,+),(+,-),(-,+),(-,-)\}\ \text{attained}\implies(P,\pi)\ \text{surjective}
\]

%dsp-level:0
\[
\Longrightarrow\quad\text{the two factors of}\ w(C)\ \text{carry independent information}.\ \blacksquare
\]

\textbf{Phys.} The product of edge signs and the transition parity measure different things — net sign content versus the number of domain walls between signs — and neither fixes the other. The A8 weight is therefore free of double counting: bulk sign and wall-count contribute independently to every configuration amplitude.

\end{theorem}

\begin{theorem}[{\texorpdfstring{CT0.11 — Maximal accessible rung}{CT0.11 - Maximal accessible rung}}]\label{tag:CT0-11}
\textbf{Thm.} $\ \text{channel-admissible}(v)\iff D_{\mathrm{tick}}\bigl(n(v)\bigr)\le C_{\mathrm{cap}}\iff n(v)\le3;$ accessible vertex sets $=\{V_0,V_1,V_2,V_3\}$ (types $\{\Xi_0,\dots,\Xi_3\}$); $D_{\mathrm{tick}}(3)=C_{\mathrm{cap}}$ (saturation); $\forall v:\ n(v)\ge4\implies O\bigl(n(v)\bigr)>0\implies F_{D^-}$ [D4.1].

\textbf{Req.} D0.8, D0.9, D4.1; CT0.5, CT0.6.

\textbf{Pf.}

%dsp-level:0
\[
\text{admissible}\iff D_{\mathrm{tick}}(n)\le C_{\mathrm{cap}}\quad[\mathrm{D0.9}]
\]

%dsp-level:0
\[
2^{\,n+1}\le16\iff n+1\le4\iff n\le3\quad[\mathrm{CT0.5},\ \mathrm{CT0.6}]
\]

%dsp-level:0
\[
n=3:\ \ 2^{4}=16=C_{\mathrm{cap}}\ (\text{tight});\qquad
n\ge4:\ \ O(n)>0\ [\mathrm{CT0.6}]\implies\text{configuration}\in F_{D^-}\ [\mathrm{D4.1}].\ \blacksquare
\]

\end{theorem}

\begin{corollary}[{\texorpdfstring{CT0.11a — three slots}{CT0.11a - three slots}}]\label{tag:CT0-11a}
$\ \max\{\mathcal H^{(n)}_{\mathrm{slot}}:\ \text{admissible}\}=\mathcal H^{(3)}_{\mathrm{slot}}\cong\mathbb C^3;\qquad 3=\lfloor\log_2 C_{\mathrm{cap}}\rfloor-1=\lfloor\log_216\rfloor-1;$ i.e. $N_{\text{colour}}=3$ is derived, not chosen.

\textbf{Req (Cor).} D3.1; CT0.6, CT0.11.

\textbf{Pf (Cor).} $\mathcal H^{(n)}_{\mathrm{slot}}\cong\mathbb C^n$ [D3.1]; $n\le3$ [CT0.11] $\Rightarrow\max=\mathbb C^3$; $\lfloor\log_2 16\rfloor-1=3$. $\blacksquare$

\end{corollary}

\begin{corollary}[{\texorpdfstring{CT0.11b — directed rung of bidirectional skeletons}{CT0.11b - directed rung of bidirectional skeletons}}]\label{tag:CT0-11b}
A channel $\{p,q\}$, $p\ne q$, is \textit{bidirectional} when both orientations are real: $(p,q)\in\mathcal R\ \wedge\ (q,p)\in\mathcal R$ [A4]. Let $S$ be a finite real-edge subgraph whose every used channel is bidirectional, and $d_S(v):=\#\{\text{channels of}\ S\ \text{incident on}\ v\}$. Then:

%dsp-level:0
\[
n(v)\ \ge\ 2\,d_S(v)\qquad[\text{each bidirectional channel adds}\ 1\ \text{to}\ \partial^{\leftarrow}\ \text{and}\ 1\ \text{to}\ \partial^{\rightarrow}\ \text{at each endpoint}]
\]

%dsp-level:0
\[
S\ \text{exhausts the real edges at}\ v\implies n(v)=2\,d_S(v);
\]

%dsp-level:0
\[
\qquad d_S(v)\ge2\implies n(v)\ge4\implies D_{\mathrm{tick}}(n)=2^{n+1}\ge32>16\implies O\bigl(n(v)\bigr)>0\implies F_{D^-}\ [\mathrm{D4.1}]
\]

%dsp-level:0
\[
\text{unoversubscribed}\iff n(v)\le3\ [\mathrm{CT0.6}]:
\]

%dsp-level:0
\[
\qquad \text{bidirectional}\ S:\ d_S(v)\le1\ (n\in\{0,2\});
\]

%dsp-level:0
\[
\qquad \text{mixed}:\ n(v)=2\,d^{\mathrm{bi}}_S(v)+d^{\mathrm{uni}}_S(v)\le3 .
\]

\textbf{Req (Cor).} A4; D0.8, D4.1; CT0.5, CT0.6, CT0.11; directed-degree enumeration [math; machine: \texttt{checks/\allowbreak{}verify\_\allowbreak{}promotion\_\allowbreak{}B\_\allowbreak{}rung.py} §A].

\textbf{Pf (Cor).}

%dsp-level:0
\[
\mathrm{A4}:\ \mathcal R\text{-membership per ordered pair};
\]

%dsp-level:0
\[
\quad\text{bidirectional channel}
\]

%dsp-level:0
\[
= \{e_{pq},e_{qp}\}\subseteq\mathcal R:\ \partial^{\leftarrow}\!+1,
\]

%dsp-level:0
\[
\ \partial^{\rightarrow}\!+1\ \text{at each endpoint}
\]

%dsp-level:0
\[
\implies n(v)
\]

%dsp-level:0
\[
\ge 2\,
\]

%dsp-level:0
\[
d_S(v)\ [\mathrm{D0.8}]
\]

%dsp-level:0
\[
\text{equality}\iff S\ \text{exhausts}\ v\text{'s incident real edges};
\]

%dsp-level:0
\[
\qquad d_S\ge2:\ n\ge4:\ \mathrm{CT0.5}:\ D_{\mathrm{tick}}(n)=2^{n+1}\ge32;\quad\mathrm{CT0.6}:\ O(n)=2^{n+1}-16>0
\]

%dsp-level:0
\[
\mathrm{CT0.11}:\ n\ge4\implies F_{D^-}\ [\mathrm{D4.1}];
\]

%dsp-level:0
\[
\qquad O(n)=0\iff n\le3\ [\mathrm{CT0.6}];
\]

%dsp-level:0
\[
\qquad \text{mixed count}:\ \text{additivity of}\ \partial^{\leftarrow},\partial^{\rightarrow}\ [\mathrm{D0.8}]\ \blacksquare
\]

\textbf{Phys.} The 16-register channel budget admits vertices of degree at most three: rungs $0$–$3$ are accessible, rung $3$ exactly exhausts the channel, and everything of degree $\ge4$ is oversubscribed and falls to the dark sector. The number three — hence the threefold slot (“colour”) structure of the theory — is forced by capacity, not assumed.

\textbf{Note.} At rungs $n\ge4$ no aspect realization $\iota_n$ exists: the unique derived read of rung-$\ge4$ slot content is $D^-$-absorption at the ledger level [CT4.4a (Part 2) — forward pointer, non-load-bearing].

\end{corollary}

\begin{theorem}[{\texorpdfstring{CT0.12 — Aspect–capacity match}{CT0.12 - Aspect-capacity match}}]\label{tag:CT0-12}
\textbf{Thm.}

%dsp-level:0
\[
\begin{array}{c|c|c}
\text{aspect} & \text{block} & \dim_{\mathbb R}\\ \hline
F_{J^+} & \mathcal E_{-+}=\operatorname{span}\{J_i\} & 3\\
F_{J^-} & \mathcal E_{--}=\operatorname{span}\{K_i\} & 3\\
F_G & \operatorname{span}\{\hat E^{ij}\}\ (\text{traceless}) & 5\\
F_W & \mathcal E_{+-}=\operatorname{span}\{W_i\} & 3\\
F_H & \operatorname{span}\{\Sigma\} & 1\\
F_{D^+} & \operatorname{span}\{E_0\} & 1\\ \hline
\text{total} & \mathfrak H(\mathbb C^4) & 16=C_{\mathrm{cap}}
\end{array}
\]

\textbf{Req.} A8; C8; CT0.6, CT0.7(4), CT1.1, CT1.2.

\textbf{Pf.}

%dsp-level:0
\[
\text{block assignments}\quad[\mathrm{CT1.2},\ \mathrm{CT0.7}(4)];\qquad
3+3+5+3+1+1=16=\dim_{\mathbb R}\mathfrak H(\mathbb C^4)\quad[\mathrm{CT1.1}]
\]

%dsp-level:0
\[
16=C_{\mathrm{cap}}\quad[\mathrm{CT0.6}].\ \blacksquare
\]

\textbf{Rmk.} $\mathbf F_{\mathrm{st}}=F_{\mathrm{EM}}\oplus F_{\mathrm{strong}}\oplus F_{D^-}$ acts on vertex/slot registers [D3.1], not on the pair channel; the CT0.6 bound constrains the pair channel only.

\textbf{Phys.} The six pair-sector force aspects of A8 tile the 16 dimensions of the two-qubit operator space exactly: no part of the channel is forceless and no force overruns it. The stack aspects (electromagnetic, colour, dark-saturation) live on separate vertex registers and are not bounded by this count.

\end{theorem}

\begin{theorem}[{\texorpdfstring{CT0.13 — Aspect activation and the active sets}{CT0.13 - Aspect activation and the active sets}}]\label{tag:CT0-13}
\textbf{Thm.}

%dsp-level:0
\[
\mathcal A(0)=\{J\};\qquad
\mathcal A(1)=\{J,\,G,\,W,\,H,\,D^+\};\qquad
\mathcal A(2)=\mathcal A(1)\cup\{\mathrm{EM}\};
\]

%dsp-level:0
\[
\mathcal A(3)=\mathcal A(2)\cup\{\mathrm{strong}\};\qquad
\mathcal A(n\ge4)=\mathcal A(3)\cup\{D^-\};
\]

with $J:=\{J^+,J^-\}$ and $W:=F_W=F_{W^+}\oplus F_{W^-}\oplus F_Z$ [C8]; $\ \neg\exists$ independent aspect named $Z$ at any rung: neutral seed $W_3\in F_W$; physical $Z$: D3.10, EW.24.

\textbf{Req.} A7, A8; C8; D0.8, D1.12, D2.12, D3.1, D3.2, D3.3, D4.1, D4.3; CT0.1, CT0.7, CT0.11, CT1.2, CT1.4, CT1.10; $\bigl[L^\dagger L=I_n\iff L\in U(n);\ \ \mathfrak u(n)=\mathfrak{su}(n)\oplus\mathfrak u(1)_{\mathrm{centre}}$ — linear algebra; cf. CT3.4, CT3.5$\bigr]$.

\textbf{Pf.}

%dsp-level:0
\[
n=0:\ \ v\in V_0\implies\neg\,\exists e\in E_{\mathcal R}:\ e\ni v\ [\mathrm{D0.8}]\implies\forall\ \text{incident pairs}:\ \in\mathcal V\ [\mathrm{CT0.1}]
\]

%dsp-level:0
\[
\mathrm{CT0.7}:\ \text{operators on the virtual block}=S\text{-odd}=\{F_{J^+},F_{J^-}\}\ [\mathrm{A8}]\implies\mathcal A(0)=\{J\}
\]

%dsp-level:0
\[
n=1:\ \ \bigl|\{e\in E_{\mathcal R}:e\ni v\}\bigr|=1\implies\vee^2\ \text{populated};\qquad
\mathrm{CT0.7}:\ S\text{-even act within}\ \vee^2
\]

%dsp-level:0
\[
\mathrm{CT1.2},\ \mathrm{CT1.10}:\ S\text{-even}=\{E_0,\ \Sigma,\ \hat E^{ij}_{(5)},\ W_i\}\xrightarrow[\mathrm{A8}]{}\{F_{D^+},\ F_H,\ F_G,\ F_W\}
\]

%dsp-level:0
\[
\implies\mathcal A(1)=\{J,G,W,H,D^+\};\qquad
\dim F_W=3\ \text{single block}\implies W_3\ \neg\text{independent}\implies\neg\exists\ \text{aspect}\ Z
\]

%dsp-level:0
\[
\text{slot algebras}:\ \ \mathcal H^{(n)}_{\mathrm{slot}}\cong\mathbb C^n\ [\mathrm{D3.1}];\quad
\text{probability-preserving}\iff L^\dagger L=I_n\iff L\in U(n)
\]

%dsp-level:0
\[
X=\Bigl(X-\tfrac{\mathfrak T X}{n}I_n\Bigr)+\tfrac{\mathfrak T X}{n}I_n\ \text{(unique trace split)}\implies\mathfrak u(n)=\mathfrak{su}(n)\oplus\mathfrak u(1)_{\mathrm{centre}}
\]

%dsp-level:0
\[
\quad [\text{linear algebra; }n=3:\ \mathrm{CT3.4},\mathrm{CT3.5}]
\]

%dsp-level:0
\[
n = 1\ \text{slot}:\ \ \mathfrak u(1)\ \text{on one slot}
\]

%dsp-level:0
\[
= \text{global phase};
\]

%dsp-level:0
\[
\ \ \neg\exists\ \text{second slot}\implies\neg\text{configuration-dependent}
\]

%dsp-level:0
\[
\implies
\]

%dsp-level:0
\[
= F_{D^+}
\]

%dsp-level:0
\[
= \Lambda\mathbb I\ [\mathrm{D4.3}]\ \text{(counted)};
\]

%dsp-level:0
\[
\quad\neg\text{new}\ [\mathrm{D3.2}]
\]

%dsp-level:0
\[
n=2:\ \ \mathfrak u(2)=\mathfrak{su}(2)\oplus\mathfrak u(1):
\]

%dsp-level:0
\[
\quad \iota_2\bigl(\mathfrak{su}(2)\bigr)\cong\mathfrak{su}(2)\cong\operatorname{span}\{W_i\}\ [\mathrm{CT1.4}]:\ \text{same abstract class realized at rung}\ 1
\]

%dsp-level:0
\[
\implies\neg\text{strictly new}\ [\mathrm{D2.12}\ (\iota_n\ \text{bridge}),\ \mathrm{D3.2}]
\]

%dsp-level:0
\[
\text{centre measurable}:\ \ V_2\ni\text{type}(1,1)\ [\mathrm{A7}]\implies\text{relative in/out phase defined}=f_v\not\equiv0\ [\mathrm{D3.3}]\implies\text{strictly new}=F_{\mathrm{EM}}\ [\mathrm{A8}]
\]

%dsp-level:0
\[
\implies\mathcal A(2)=\mathcal A(1)\cup\{\mathrm{EM}\};
\]

%dsp-level:0
\[
\qquad f_v\ \text{requires}\ \partial^{\leftarrow}\!\ge1\ \wedge\ \partial^{\rightarrow}\!\ge1\implies n\ge2\ [\mathrm{D3.3},\mathrm{D0.8}]\implies\mathrm{EM}\notin\mathcal A(0)\cup\mathcal A(1)
\]

%dsp-level:0
\[
n=3:\ \ \mathfrak u(3)=\mathfrak{su}(3)\oplus\mathfrak u(1):\quad
\mathfrak u(1)\ \text{duplicates}\ F_{\mathrm{EM}}\ (n=2)\implies\neg\text{new};
\]

%dsp-level:0
\[
\qquad \mathfrak{su}(3)\ \text{strictly new}=F_{\mathrm{strong}}\ [\mathrm{A8}]
\]

%dsp-level:0
\[
\implies\mathcal A(3)=\mathcal A(2)\cup\{\mathrm{strong}\}
\]

%dsp-level:0
\[
n\ge4:\ \ \text{oversubscribed}\ [\mathrm{CT0.11}]\implies F_{D^-}\ \text{active there}\wedge\text{only there}\ [\mathrm{D4.1}];
\]

%dsp-level:0
\[
\qquad \text{rung-3 structure persists}\implies\mathcal A(n\ge4)=\mathcal A(3)\cup\{D^-\}.\ \blacksquare
\]

\textbf{Phys.} Forces switch on one at a time as vertices acquire edges. Rung $0$ has only the Janus seals that create and destroy reality; rung $1$ lights up gravity ($G$), weak ($W$), Higgs ($H$) and uniform dark ($D^+$); at rung $2$ the local phase along a pass-through becomes measurable — electromagnetism; at rung $3$ the traceless three-slot algebra appears — colour; beyond, configurations lock into dark saturation ($D^-$). The gauge sequence of the corpus is channel bookkeeping, and no rung ever carries an independent $Z$ aspect — only the seed $W_3$ inside $F_W$.

\end{theorem}

\begin{theorem}[{\texorpdfstring{CT0.14 — Mask flag and block dimensions}{CT0.14 - Mask flag and block dimensions}}]\label{tag:CT0-14}
\textbf{Thm.}

%dsp-level:0
\[
\mathcal K_0\subset\mathcal K_1\subset\mathcal K_2\subset\mathcal K_3\subset\mathcal K_{n\ge4}\quad(\text{strict});\qquad
\dim\mathcal K_0=6,\ \ \dim\mathcal K_1=16,\ \ \dim\mathcal K_2=17,\ \ \dim\mathcal K_3=25;
\]

%dsp-level:0
\[
\bigl(\mathcal P_M^{(n)}\bigr)^2=\mathcal P_M^{(n)}=\bigl(\mathcal P_M^{(n)}\bigr)^\dagger;\qquad
\mathfrak T\,\mathcal P_M^{(n)}=\dim\mathcal K_n;\qquad
\dim\mathcal K^{(\mathrm{EM})}=1,\ \ \dim\mathcal K^{(\mathrm{strong})}=8 .
\]

\textbf{Req.} A8; D2.3; CT0.12, CT0.13; CT3.8 (\textbf{F} — forward pointer: the block-dimension row re-verifies the gluon count).

\textbf{Pf.}

%dsp-level:0
\[
\mathrm{CT0.13}:\ \ \mathcal A(0)\subset\mathcal A(1)\subset\mathcal A(2)\subset\mathcal A(3)\subset\mathcal A(n\ge4)\ \ (\text{strict})
\]

%dsp-level:0
\[
\mathrm{D2.3}:\ \ \mathcal K_n=\bigoplus_{A\in\mathcal A(n)}\mathcal K^{(A)}\implies\mathcal K_0\subset\mathcal K_1\subset\mathcal K_2\subset\mathcal K_3\subset\mathcal K_{n\ge4}\ \ (\text{strict})
\]

%dsp-level:0
\[
\dim\mathcal K_0=3+3=6\ [\mathrm{CT0.12}];\qquad
\dim\mathcal K_1=6+(5+3+1+1)=16\ [\mathrm{CT0.12}]
\]

%dsp-level:0
\[
\dim\mathcal K^{(\mathrm{EM})}=\dim Z\bigl(\mathfrak u(2)\bigr)=1\ [\mathrm{CT0.13}\ \text{step}\ n{=}2:\ \text{centre spanned by}\ I_2]\implies\dim\mathcal K_2=17
\]

%dsp-level:0
\[
\dim\mathcal K^{(\mathrm{strong})}=\dim\mathfrak{su}(3)=8\ [\mathrm{CT3.8}]\implies\dim\mathcal K_3=25
\]

%dsp-level:0
\[
\mathrm{D2.3}:\ \mathcal P_M^{(n)}\ \text{orthogonal projection}\implies P^2=P=P^\dagger,\quad\mathfrak T P=\dim\operatorname{Im}P=\dim\mathcal K_n.\ \blacksquare
\]

\textbf{Rmk.} $\dim\mathcal K_1=16=C_{\mathrm{cap}}$: rung $1$ already exhausts the pair channel; growth at rungs $2,3$ is carried entirely by stack registers [CT0.12 Rmk].

\textbf{Phys.} The rung masks form a strict flag $6\subset16\subset17\subset25$: each new rung adds exactly the block of its new force — one dimension for electromagnetism at rung $2$, eight for colour at rung $3$. “What the universe can do at rung $n$” is literally a nested subspace of force space, and the masks are the orthogonal projections onto it.

\end{theorem}

\phantomsection
\section[Tier 1A — The Pair Operator Algebra]{\texorpdfstring{Tier 1A — The Pair Operator Algebra}{Tier 1A - The Pair Operator Algebra}}

\begin{theorem}[{\texorpdfstring{CT1.0 — Signature classification of the $T$-grading}{CT1.0 - Signature classification of the -grading}}]\label{tag:CT1-0}
\textbf{Thm.} Let $\eta\in\{\pm1\}^4$ have exactly one odd entry [A3]. Put $\eta^{(\mathfrak o)}:=(+,\dots,-_{\mathfrak o},\dots,+)$, $\zeta_i:=\eta^0\eta^i$, $\zeta_0:=1$, $\mathsf D_\zeta:\sigma_\mu\mapsto\zeta_\mu\sigma_\mu$, $\mathrm{sf}(X):=\sigma_2X^{\mathrm T}\sigma_2$. Then:

$(1)$ $\eta=\pm\eta^{(\mathfrak o)}$ for exactly one $\mathfrak o\in\{0,1,2,3\}$ and exactly one sign.

$(2)$ $T[\eta]=T[-\eta]$, and $T=\mathsf D_\zeta\otimes\mathsf D_\zeta$.

$(3)$ $[S,T]=0$ and $(\dim\mathcal E_{++},\dim\mathcal E_{+-},\dim\mathcal E_{-+},\dim\mathcal E_{--})=(7,3,3,3)$ for every $\mathfrak o$.

$(4)$ $\mathfrak o=0$: $T=\mathrm{sf}\otimes\mathrm{sf}$ (inversion of both Bloch vectors). $\mathfrak o=k\ge1$: $\mathsf D_\zeta$ reflects the Bloch axis $k$. The three maps $\mathfrak o\in\{1,2,3\}$ are conjugate under $\mathrm{Ad}_{U\otimes U}$, $U\in\mathrm{SU}(2)$; the map $\mathfrak o=0$ is conjugate to none of them.

$(5)$ $T\,\mathrm{Ad}_{U\otimes U}=\mathrm{Ad}_{U\otimes U}\,T$ for all $U\in\mathrm{SU}(2)$ $\iff\mathfrak o=0$.

\textbf{Req.} A3; C1; $[\text{math}]$: $\mathrm{Ad}_U$ fixes $\sigma_0$ and acts on $\operatorname{span}\{\sigma_i\}$ by $R(U)\in SO(3)$, every $R\in SO(3)$ arising; $\sigma_2U^{\mathrm T}\sigma_2=U^\dagger$ for $U\in\mathrm{SU}(2)$.

\textbf{Pf.}

%dsp-level:0
\[
\text{(1)}:\ \text{odd entry}=-1\implies\eta=\eta^{(\mathfrak o)};\quad\text{odd entry}=+1\implies\eta=-\eta^{(\mathfrak o)};\quad\mathfrak o\ \text{is the unique odd slot}\quad[\mathrm{A3}]
\]

%dsp-level:0
\[
\text{(2)}:\ \ \eta^\mu\eta^\nu=(-\eta^\mu)(-\eta^\nu)\implies T[\eta]=T[-\eta];
\]

%dsp-level:0
\[
\qquad \zeta_\mu\zeta_\nu=(\eta^0\eta^\mu)(\eta^0\eta^\nu)=\eta^\mu\eta^\nu\ \ [(\eta^0)^2=1;\ \zeta_0=\eta^0\eta^0]
\]

%dsp-level:0
\[
\implies(\mathsf D_\zeta\otimes\mathsf D_\zeta)(\sigma_\mu\otimes\sigma_\nu)=\zeta_\mu\zeta_\nu\,\sigma_\mu\otimes\sigma_\nu=T(\sigma_\mu\otimes\sigma_\nu)\ \Rightarrow\ T=\mathsf D_\zeta\otimes\mathsf D_\zeta
\]

%dsp-level:0
\[
\text{(3)}:\ \ ST(\sigma_\mu\otimes\sigma_\nu)=\eta^\mu\eta^\nu\,\sigma_\nu\otimes\sigma_\mu=TS(\sigma_\mu\otimes\sigma_\nu)\implies[S,T]=0;
\]

%dsp-level:0
\[
\qquad \sigma_\mu\otimes\sigma_\mu:\ (S,T)=(+,+);\quad
\sigma_\mu\otimes\sigma_\nu\pm\sigma_\nu\otimes\sigma_\mu\ (\mu<\nu):\ (S,T)=(\pm,\eta^\mu\eta^\nu)
\]

%dsp-level:0
\[
n_-:=\#\{\mu<\nu:\eta^\mu\eta^\nu=-1\}=\#\bigl\{\{\mathfrak o,\nu\}:\nu\ne\mathfrak o\bigr\}=3;\qquad n_+=6-n_-=3
\]

%dsp-level:0
\[
\dim\mathcal E_{++}=4+n_+=7,\quad\dim\mathcal E_{+-}=n_-=3,\quad\dim\mathcal E_{-+}=n_+=3,\quad\dim\mathcal E_{--}=n_-=3
\]

%dsp-level:0
\[
\text{(4)}:\ \ \mathfrak o=0:\ \zeta_i=\eta^0\eta^i=-1\ \forall i;
\]

%dsp-level:0
\[
\qquad \mathrm{sf}(\sigma_0)=\sigma_0,\ \ \mathrm{sf}(\sigma_i)=\sigma_2\sigma_i^{\mathrm T}\sigma_2=-\sigma_i\ (i=1,2,3)\implies\mathsf D_\zeta=\mathrm{sf}
\]

%dsp-level:0
\[
\mathfrak o=k\ge1:\ \zeta_k=\eta^0\eta^k=-1,\ \ \zeta_i=+1\ (i\ne k)\implies\mathsf D_\zeta=\text{reflection of Bloch axis }k
\]

%dsp-level:0
\[
\mathrm{Ad}_U\,\mathsf D_\zeta\,\mathrm{Ad}_U^{-1}=\mathsf D_{\zeta'}\ \text{with}\ \operatorname{diag}(\zeta')=R(U)\operatorname{diag}(\zeta)R(U)^{-1}\quad[\text{Req}]
\]

%dsp-level:0
\[
\text{axis reflections}:\ R\ \text{permuting the axes}\implies T^{(1)}\sim T^{(2)}\sim T^{(3)};
\]

%dsp-level:0
\[
\qquad \operatorname{spec}(-\mathbb1_3)=\{-1,-1,-1\}\ne\{-1,+1,+1\}\implies T^{(0)}\not\sim T^{(k)}
\]

%dsp-level:0
\[
\text{(5)}\ (\Leftarrow):\ \ \mathrm{sf}(UXU^\dagger)=\sigma_2\bar UX^{\mathrm T}U^{\mathrm T}\sigma_2=U\,\sigma_2X^{\mathrm T}\sigma_2\,U^\dagger=U\,\mathrm{sf}(X)\,U^\dagger
\]

%dsp-level:0
\[
\quad [\sigma_2U^{\mathrm T}\sigma_2=U^\dagger\Rightarrow\sigma_2\bar U\sigma_2=U]
\]

%dsp-level:0
\[
\implies\mathrm{sf}\otimes\mathrm{sf}\ \text{commutes with}\ \mathrm{Ad}_{U\otimes U};
\]

%dsp-level:0
\[
\qquad \text{(5)}\ (\Rightarrow):\ \ U_c:=\tfrac12\bigl(\mathbb1-i(\sigma_1+\sigma_2+\sigma_3)\bigr):\ \ \mathrm{Ad}_{U_c}:\ \sigma_1\mapsto\sigma_2\mapsto\sigma_3\mapsto\sigma_1
\]

%dsp-level:0
\[
T\,\mathrm{Ad}_{U_c\otimes U_c}(\sigma_0\otimes\sigma_i)=\eta^0\eta^{i+1}\,\sigma_0\otimes\sigma_{i+1};
\]

%dsp-level:0
\[
\qquad \mathrm{Ad}_{U_c\otimes U_c}\,T(\sigma_0\otimes\sigma_i)=\eta^0\eta^{i}\,\sigma_0\otimes\sigma_{i+1}\quad(i\bmod3)
\]

%dsp-level:0
\[
\implies\eta^1=\eta^2=\eta^3\ \ [\eta^0=\pm1];\qquad\text{exactly one odd entry}\implies\mathfrak o=0.\ \blacksquare
\]

\textbf{Phys.} A3 now says only that the sign pattern is $1+3$. The overall sign is the mostly-plus / mostly-minus convention; which spatial slot is odd is a relabelling. What remains is one selection: the odd slot is the identity slot, which is the unique choice making $T$ rotation-covariant, namely the spin flip $\mathrm{sf}\otimes\mathrm{sf}$ (time reversal on the pair). Any other slot gives a reflection that singles out a spatial direction.

\textbf{Rmk (what $\mathfrak o\ge1$ would do).} The sector dimensions $(7,3,3,3)$ do not change, but CT1.2's table does. For $\mathfrak o=1$: $W_2,W_3\in\mathcal E_{++}$ (not $\mathcal E_{+-}$); $J_2,J_3\in\mathcal E_{--}$ and $K_2,K_3\in\mathcal E_{-+}$ (swapped); $E^{12},E^{13}\in\mathcal E_{+-}$ (not $\mathcal E_{++}$). Only $W_1,J_1,K_1$ sit where CT1.2 puts them. Symmetric statements hold for $\mathfrak o=2,3$. (Numerical check over all four $\mathfrak o$ on the 16-element basis; machine: \texttt{checks/\allowbreak{}verify\_\allowbreak{}CT10\_\allowbreak{}signature.py}.) So the selection $\mathfrak o=0$ in C13 is load-bearing for CT1.2 and everything built on the $W_i$ triplet, not cosmetic.

\end{theorem}

\begin{theorem}[{\texorpdfstring{CT1.1 — Joint $(S,T)$ decomposition of $\mathfrak H(\mathbb C^4)$}{CT1.1 - Joint decomposition of}}]\label{tag:CT1-1}
\textbf{Thm.} $\ [S,T]=0;$ $\qquad\mathfrak H(\mathbb C^4)=\mathcal E_{++}\oplus\mathcal E_{+-}\oplus\mathcal E_{-+}\oplus\mathcal E_{--};$ $\qquad(\dim\mathcal E_{++},\dim\mathcal E_{+-},\dim\mathcal E_{-+},\dim\mathcal E_{--})=(7,3,3,3);\qquad7+3+3+3=16.$

\textbf{Req.} A3, C13; C3; CT1.0; Pauli basis facts: $\ \mathfrak T(\sigma_\mu\sigma_\alpha)=2\delta_{\mu\alpha}$, $\sigma_\mu^\dagger=\sigma_\mu$ [linear algebra].

\textbf{Pf.}

%dsp-level:0
\[
\bigl\{\sigma_\mu\otimes\sigma_\nu\bigr\}_{\mu,\nu=0}^{3}\subset\mathfrak H(\mathbb C^4);\qquad
\mathfrak T\bigl((\sigma_\mu\otimes\sigma_\nu)(\sigma_\alpha\otimes\sigma_\beta)\bigr)=4\,\delta_{\mu\alpha}\delta_{\nu\beta}
\]

%dsp-level:0
\[
\Longrightarrow\ \text{orthogonal}\implies\text{basis};\qquad\dim_{\mathbb R}\mathfrak H(\mathbb C^4)=16
\]

%dsp-level:0
\[
\mathrm{A3}:\ \ S:\ \sigma_\mu\otimes\sigma_\nu\mapsto\sigma_\nu\otimes\sigma_\mu\ (\text{permutation});
\]

%dsp-level:0
\[
\qquad T:\ \sigma_\mu\otimes\sigma_\nu\mapsto\eta^\mu\eta^\nu\,\sigma_\mu\otimes\sigma_\nu,\ \ \eta^\mu\eta^\nu\in\{\pm1\}
\]

%dsp-level:0
\[
\implies S^2=T^2=\mathbb1,\ \ \operatorname{spec}\subseteq\{\pm1\};
\]

%dsp-level:0
\[
\qquad ST(\sigma_\mu\otimes\sigma_\nu)=\eta^\mu\eta^\nu\,\sigma_\nu\otimes\sigma_\mu=TS(\sigma_\mu\otimes\sigma_\nu)\implies[S,T]=0
\]

%dsp-level:0
\[
[S,T]=0\ \wedge\ S^2=T^2=\mathbb1\implies\text{simultaneous eigenspaces}\ \mathcal E_{st}\ \text{span}\ \mathfrak H(\mathbb C^4)\quad[\mathrm{C3}]
\]

%dsp-level:0
\[
\eta:=(-,+,+,+)\ [\mathrm{C13}]:\quad T=+1\iff(\mu,\nu)\in\{0,0\}\cup\{i,j\};\qquad T=-1\iff(\mu,\nu)\in\{0,i\}\cup\{i,0\}
\]

%dsp-level:0
\[
\mathcal E_{++}:\ \{\sigma_\mu\otimes\sigma_\mu\}_{\mu=0}^{3}\ (4)\ \cup\ \{\sigma_i\otimes\sigma_j+\sigma_j\otimes\sigma_i\}_{i<j}\ (3)=7
\]

%dsp-level:0
\[
\mathcal E_{+-}:\ \{\sigma_0\otimes\sigma_i+\sigma_i\otimes\sigma_0\}_{i=1}^{3}=3;\qquad
\mathcal E_{-+}:\ \{\sigma_i\otimes\sigma_j-\sigma_j\otimes\sigma_i\}_{i<j}=3
\]

%dsp-level:0
\[
\mathcal E_{--}:\ \{\sigma_0\otimes\sigma_i-\sigma_i\otimes\sigma_0\}_{i=1}^{3}=3;\qquad
7+3+3+3=16\implies\mathfrak H(\mathbb C^4)=\bigoplus_{s,t}\mathcal E_{st}.\ \blacksquare
\]

\textbf{Phys.} The 16-dimensional space of two-qubit observables splits under exchange symmetry ($S$) and the sign grading $T$ ($\eta^\mu\eta^\nu$; metric-sign reading by C13) into blocks $7+3+3+3$. This grading is the skeleton of the whole force structure: the $7$-block supplies gravity, Higgs and uniform dark; the $3$-blocks supply weak, Janus-seal, and the remaining virtual sector.

\end{theorem}

\begin{theorem}[{\texorpdfstring{CT1.2 — Generator basis and sector assignment}{CT1.2 - Generator basis and sector assignment}}]\label{tag:CT1-2}
\textbf{Thm.} Define

%dsp-level:0
\[
E_0:=\sigma_0\otimes\sigma_0;\qquad
\Sigma:=\textstyle\sum_{i=1}^{3}\sigma_i\otimes\sigma_i;\qquad
W_i:=\tfrac12\bigl(\sigma_0\otimes\sigma_i+\sigma_i\otimes\sigma_0\bigr);
\]

%dsp-level:0
\[
J_i:=\tfrac12\bigl(\sigma_j\otimes\sigma_k-\sigma_k\otimes\sigma_j\bigr)\ [(i,j,k)\ \text{cyclic}];
\]

%dsp-level:0
\[
\qquad K_i:=\tfrac12\bigl(\sigma_0\otimes\sigma_i-\sigma_i\otimes\sigma_0\bigr);
\]

%dsp-level:0
\[
\qquad E^{ij}:=\tfrac12\bigl(\sigma_i\otimes\sigma_j+\sigma_j\otimes\sigma_i\bigr).
\]

Then

%dsp-level:0
\[
\begin{array}{c|c|c|c|c}
\text{generator} & \# & S & T & \text{sector}\\ \hline
E_0 & 1 & + & + & \mathcal E_{++}\\
\Sigma & 1 & + & + & \mathcal E_{++}\\
\hat E^{ij}\ (\text{traceless}) & 5 & + & + & \mathcal E_{++}\\
W_i & 3 & + & - & \mathcal E_{+-}\\
J_i & 3 & - & + & \mathcal E_{-+}\\
K_i & 3 & - & - & \mathcal E_{--}
\end{array}
\qquad 1+1+5+3+3+3=16\implies\text{basis of}\ \mathfrak H(\mathbb C^4);
\]

%dsp-level:0
\[
W_i=\tfrac12\bigl(\sigma_i^{(p)}+\sigma_i^{(q)}\bigr)=\text{total-spin generator}.
\]

\textbf{Req.} A3, C1, C3, C13; CT1.0, CT1.1.

\textbf{Pf.}

%dsp-level:0
\[
S E_0=E_0;\ \ \eta^0\eta^0=+1\implies E_0\in\mathcal E_{++}\quad[\mathrm{A3}]
\]

%dsp-level:0
\[
S E^{ij}=E^{ij}\ (\text{by construction});\ \ \eta^i\eta^j=+1\ [\mathrm{C13}]\implies E^{ij}\in\mathcal E_{++};\qquad\#\{E^{ij}\}_{i\le j}=6
\]

%dsp-level:0
\[
\delta_{ij}E^{ij}=\textstyle\sum_iE^{ii}=\Sigma\ (1\ \text{component})\ \oplus\ \hat E^{ij}\ (5);\qquad S,T\ \text{diagonal on index pairs}\implies\text{both}\subseteq\mathcal E_{++}
\]

%dsp-level:0
\[
W_i:\ S=+;\ \ \eta^0\eta^i=-1\ [\mathrm{C13}]\implies T=-\implies W_i\in\mathcal E_{+-}
\]

%dsp-level:0
\[
J_i:\ S=-;\ \ \eta^j\eta^k=+1\ [\mathrm{C13}]\implies T=+\implies J_i\in\mathcal E_{-+}
\]

%dsp-level:0
\[
K_i:\ S=-;\ \ \eta^0\eta^i=-1\ [\mathrm{C13}]\implies T=-\implies K_i\in\mathcal E_{--}
\]

%dsp-level:0
\[
\text{counts}\ (1,1,5,3,3,3)\ \text{match}\ (7{=}1{+}1{+}5,\ 3,\ 3,\ 3)\ [\mathrm{CT1.1}];
\]

%dsp-level:0
\[
\qquad \text{distinct orthogonal combinations}\implies\text{independent}\implies\text{basis}
\]

%dsp-level:0
\[
\sigma_0\otimes\sigma_i\equiv\sigma_i^{(p)},\ \ \sigma_i\otimes\sigma_0\equiv\sigma_i^{(q)}\implies W_i=\tfrac12\bigl(\sigma_i^{(p)}+\sigma_i^{(q)}\bigr).\ \blacksquare
\]

\textbf{Phys.} This fixes the named basis used throughout the corpus: identity $E_0$, exchange scalar $\Sigma$, total spin $W_i$ (weak sector), the antisymmetric spatial and temporal generators $J_i,K_i$ (Janus/seal sectors), and the symmetric spatial tensors $E^{ij}$ (gravity). Each generator’s $(S,T)$ parity decides which force aspect of A8 it belongs to.

\end{theorem}

\begin{theorem}[{\texorpdfstring{CT1.3 — Centrality of the identity}{CT1.3 - Centrality of the identity}}]\label{tag:CT1-3}
\textbf{Thm.} $\ \forall X\in\mathfrak H(\mathbb C^4):\quad[E_0,X]=0.$

\textbf{Req.} C1; CT1.2.

\textbf{Pf.}

%dsp-level:0
\[
E_0=\sigma_0\otimes\sigma_0=I_2\otimes I_2=I_4\quad[\mathrm{C1},\ \mathrm{CT1.2}];\qquad
[I_4,X]=I_4X-XI_4=0.\ \blacksquare
\]

\textbf{Phys.} The identity direction commutes with everything, so the aspect it carries — the uniform dark term $F_{D^+}=\Lambda\mathbb I$ — acts identically on every configuration: a cosmological-constant-type force.

\end{theorem}

\begin{theorem}[{\texorpdfstring{CT1.4 — The $W_i$ close $\mathfrak{su}(2)$}{CT1.4 - The close}}]\label{tag:CT1-4}
\textbf{Thm.} $\ [W_i,W_j]=i\,\epsilon_{ijk}W_k;\qquad\operatorname{span}_{\mathbb R}\{W_i\}\cong\mathfrak{su}(2)\ [\mathrm{C2}].$

\textbf{Req.} C2; CT1.2; $\ [\sigma_i,\sigma_j]=2i\,\epsilon_{ijk}\sigma_k$ [Pauli algebra].

\textbf{Pf.}

%dsp-level:0
\[
\text{definitions cyclically symmetric}\implies\text{suffices}:\ [W_1,W_2]=iW_3
\]

%dsp-level:0
\[
[W_1,W_2]=\tfrac14\bigl[\sigma_0\otimes\sigma_1+\sigma_1\otimes\sigma_0,\ \ \sigma_0\otimes\sigma_2+\sigma_2\otimes\sigma_0\bigr]\quad[\mathrm{CT1.2}]
\]

%dsp-level:0
\[
[\sigma_0\otimes\sigma_1,\ \sigma_0\otimes\sigma_2]=\sigma_0^2\otimes[\sigma_1,\sigma_2]=2i\,(\sigma_0\otimes\sigma_3)
\]

%dsp-level:0
\[
[\sigma_0\otimes\sigma_1,\ \sigma_2\otimes\sigma_0]=\sigma_2\otimes\sigma_1-\sigma_2\otimes\sigma_1=0;\qquad
[\sigma_1\otimes\sigma_0,\ \sigma_0\otimes\sigma_2]=0\quad[\text{distinct factors}]
\]

%dsp-level:0
\[
[\sigma_1\otimes\sigma_0,\ \sigma_2\otimes\sigma_0]=[\sigma_1,\sigma_2]\otimes\sigma_0^2=2i\,(\sigma_3\otimes\sigma_0)
\]

%dsp-level:0
\[
[W_1,W_2]=\tfrac14\bigl(2i\,\sigma_0\otimes\sigma_3+2i\,\sigma_3\otimes\sigma_0\bigr)=\tfrac i2\bigl(\sigma_0\otimes\sigma_3+\sigma_3\otimes\sigma_0\bigr)=iW_3.\ \blacksquare
\]

\textbf{Phys.} The three symmetric mixed generators close a perfect $\mathfrak{su}(2)$: they are the total spin of the qubit pair. This $\mathfrak{su}(2)$ becomes the weak-isospin gauge algebra of the corpus — the weak force is the pair’s rotation algebra, not an added postulate.

\end{theorem}

\begin{theorem}[{\texorpdfstring{CT1.5 — Commutators of $\Sigma$ with $J_i$, $K_i$, $W_i$}{CT1.5 - Commutators of with , ,}}]\label{tag:CT1-5}
\textbf{Thm.} $\ [\Sigma,J_i]=4i\,K_i;\qquad[\Sigma,K_i]=-4i\,J_i;\qquad[\Sigma,W_i]=0.$

\textbf{Req.} CT0.7(1),(3) $\ [\Sigma=2\mathbb S-E_0;\ \ S\text{-even}\Rightarrow[\mathbb S,\cdot]=0]$, CT1.2, CT1.3; Pauli algebra.

\textbf{Pf.}

%dsp-level:0
\[
\text{cyclicity}\implies\text{suffices}:\ i=1
\]

%dsp-level:0
\[
[\sigma_1\otimes\sigma_1,\ J_1]=0:
\]

%dsp-level:0
\[
\quad [\sigma_1\otimes\sigma_1,\ \sigma_2\otimes\sigma_3]=(i\sigma_3)\otimes(-i\sigma_2)-(-i\sigma_3)\otimes(i\sigma_2)=0;\quad\text{likewise}\ \sigma_3\otimes\sigma_2
\]

%dsp-level:0
\[
[\sigma_2\otimes\sigma_2,\ \sigma_2\otimes\sigma_3]=\sigma_0\otimes[\sigma_2,\sigma_3]=2i\,(\sigma_0\otimes\sigma_1);\qquad
[\sigma_2\otimes\sigma_2,\ \sigma_3\otimes\sigma_2]=2i\,(\sigma_1\otimes\sigma_0)
\]

%dsp-level:0
\[
\implies[\sigma_2\otimes\sigma_2,\ J_1]=\tfrac12\bigl(2i\,\sigma_0\otimes\sigma_1-2i\,\sigma_1\otimes\sigma_0\bigr)=2i\,K_1
\]

%dsp-level:0
\[
[\sigma_3\otimes\sigma_3,\ \sigma_2\otimes\sigma_3]=-2i\,(\sigma_1\otimes\sigma_0);
\]

%dsp-level:0
\[
\qquad [\sigma_3\otimes\sigma_3,\ \sigma_3\otimes\sigma_2]=-2i\,(\sigma_0\otimes\sigma_1)\implies[\sigma_3\otimes\sigma_3,\ J_1]=2i\,K_1
\]

%dsp-level:0
\[
[\Sigma,J_1]=0+2iK_1+2iK_1=4i\,K_1
\]

%dsp-level:0
\[
[\sigma_1\otimes\sigma_1,\ K_1]=0\ (\text{as above});
\]

%dsp-level:0
\[
\qquad [\sigma_2\otimes\sigma_2,\ \sigma_0\otimes\sigma_1]=-2i\,(\sigma_2\otimes\sigma_3);
\]

%dsp-level:0
\[
\qquad [\sigma_2\otimes\sigma_2,\ \sigma_1\otimes\sigma_0]=-2i\,(\sigma_3\otimes\sigma_2)
\]

%dsp-level:0
\[
\implies[\sigma_2\otimes\sigma_2,\ K_1]=\tfrac12\bigl(-2i\,\sigma_2\otimes\sigma_3+2i\,\sigma_3\otimes\sigma_2\bigr)=-2i\,J_1
\]

%dsp-level:0
\[
[\sigma_3\otimes\sigma_3,\ \sigma_0\otimes\sigma_1]=2i\,(\sigma_3\otimes\sigma_2);
\]

%dsp-level:0
\[
\qquad [\sigma_3\otimes\sigma_3,\ \sigma_1\otimes\sigma_0]=2i\,(\sigma_2\otimes\sigma_3)\implies[\sigma_3\otimes\sigma_3,\ K_1]=-2i\,J_1
\]

%dsp-level:0
\[
[\Sigma,K_1]=-4i\,J_1
\]

%dsp-level:0
\[
\Sigma=2\mathbb S-E_0\ [\mathrm{CT0.7}(1)];\qquad
SW_i=W_i\ (S\text{-even})\implies[\mathbb S,W_i]=0\ [\mathrm{CT0.7}(3)];\qquad
[E_0,W_i]=0\ [\mathrm{CT1.3}]
\]

%dsp-level:0
\[
\implies[\Sigma,W_i]=2[\mathbb S,W_i]-[E_0,W_i]=0.\ \blacksquare
\]

\textbf{Phys.} The exchange scalar $\Sigma$ rotates the two seal sectors into one another ($J\leftrightarrow K$, a $90^\circ$ rotation inside the virtual sector) while leaving total spin untouched. The volume/trace mode therefore couples to Janus (edge creation–annihilation) dynamics but not to spin.

\end{theorem}

\begin{theorem}[{\texorpdfstring{CT1.6 — Commutators $[J_i,J_j]$ and $[K_i,K_j]$}{CT1.6 - Commutators and}}]\label{tag:CT1-6}
\textbf{Thm.} $\ [J_i,J_j]=i\,\epsilon_{ijk}W_k;\qquad[K_i,K_j]=i\,\epsilon_{ijk}W_k.$

\textbf{Req.} CT1.2; Pauli algebra.

\textbf{Pf.}

%dsp-level:0
\[
\text{suffices}:\ [J_1,J_2]=iW_3=[K_1,K_2]
\]

%dsp-level:0
\[
[J_1,J_2]=\tfrac14\bigl[\sigma_2\otimes\sigma_3-\sigma_3\otimes\sigma_2,\ \ \sigma_3\otimes\sigma_1-\sigma_1\otimes\sigma_3\bigr]
\]

%dsp-level:0
\[
[\sigma_2\otimes\sigma_3,\ \sigma_3\otimes\sigma_1]=(i\sigma_1)\otimes(i\sigma_2)-(-i\sigma_1)\otimes(-i\sigma_2)=0
\]

%dsp-level:0
\[
[\sigma_2\otimes\sigma_3,\ -\sigma_1\otimes\sigma_3]=-[\sigma_2,\sigma_1]\otimes\sigma_3^2=2i\,(\sigma_3\otimes\sigma_0)
\]

%dsp-level:0
\[
[-\sigma_3\otimes\sigma_2,\ \sigma_3\otimes\sigma_1]=-\sigma_0\otimes[\sigma_2,\sigma_1]=2i\,(\sigma_0\otimes\sigma_3)
\]

%dsp-level:0
\[
[-\sigma_3\otimes\sigma_2,\ -\sigma_1\otimes\sigma_3]=(i\sigma_2)\otimes(i\sigma_1)-(-i\sigma_2)\otimes(-i\sigma_1)=0
\]

%dsp-level:0
\[
\implies[J_1,J_2]=\tfrac14\bigl(2i\,\sigma_3\otimes\sigma_0+2i\,\sigma_0\otimes\sigma_3\bigr)=iW_3
\]

%dsp-level:0
\[
[K_1,K_2]:\ \ \text{cross terms}=0\ [\text{distinct factors}];
\]

%dsp-level:0
\[
\qquad [K_1,K_2]=\tfrac14\Bigl([\sigma_0\otimes\sigma_1,\sigma_0\otimes\sigma_2]+[\sigma_1\otimes\sigma_0,\sigma_2\otimes\sigma_0]\Bigr)
\]

%dsp-level:0
\[
=\tfrac14\bigl(2i\,\sigma_0\otimes\sigma_3+2i\,\sigma_3\otimes\sigma_0\bigr)=iW_3.\ \blacksquare
\]

\textbf{Phys.} Neither seal sector closes on itself: brackets of two $J$’s or two $K$’s land in the total-spin (weak) sector. Reality-flipping operations source the rotational/weak structure of the pair — the first algebraic link between the Janus dynamics and gauge structure.

\end{theorem}

\begin{theorem}[{\texorpdfstring{CT1.7 — Mixed commutators $[J_i,K_j]$}{CT1.7 - Mixed commutators}}]\label{tag:CT1-7}
\textbf{Thm.}

%dsp-level:0
\[
[J_i,K_j]=i\bigl(\delta_{ij}\Sigma-E^{ij}\bigr)=\tfrac i2\Bigl(2\delta_{ij}\Sigma-\sigma_i\otimes\sigma_j-\sigma_j\otimes\sigma_i\Bigr);
\]

%dsp-level:0
\[
[J_i,K_i]=i\!\!\sum_{k\ne i}\!\!\sigma_k\otimes\sigma_k;\qquad
[J_i,K_j]=-iE^{ij}=-\tfrac i2\bigl(\sigma_i\otimes\sigma_j+\sigma_j\otimes\sigma_i\bigr)\ \ (i\ne j);
\]

%dsp-level:0
\[
[J_i,K_j]\in(\mathcal E_{++})_{\mathrm{sym\ spatial}};\qquad [J_i,K_j]\notin\operatorname{span}\{W_k\}.
\]

\textbf{Req.} CT1.1, CT1.2; $\epsilon$-identities: $\ \epsilon_{iab}\epsilon_{bjm}=\delta_{ij}\delta_{am}-\delta_{im}\delta_{aj}$, $\ \epsilon_{iab}\epsilon_{ajm}=\delta_{bj}\delta_{im}-\delta_{bm}\delta_{ij}$.

\textbf{Pf.}

%dsp-level:0
\[
J_i=\tfrac12\,\epsilon_{iab}\,\sigma_a\otimes\sigma_b\ (a,b\in\{1,2,3\})\quad[\mathrm{CT1.2}];\qquad
K_j=\tfrac12\bigl(\sigma_0\otimes\sigma_j-\sigma_j\otimes\sigma_0\bigr)
\]

%dsp-level:0
\[
[\sigma_a\otimes\sigma_b,\ \sigma_0\otimes\sigma_j]=\sigma_a\otimes[\sigma_b,\sigma_j]=2i\,\epsilon_{bjm}\,\sigma_a\otimes\sigma_m
\]

%dsp-level:0
\[
[\sigma_a\otimes\sigma_b,\ \sigma_j\otimes\sigma_0]=[\sigma_a,\sigma_j]\otimes\sigma_b=2i\,\epsilon_{ajm}\,\sigma_m\otimes\sigma_b
\]

%dsp-level:0
\[
\implies[J_i,K_j]=\tfrac i2\,\epsilon_{iab}\bigl(\epsilon_{bjm}\,\sigma_a\otimes\sigma_m-\epsilon_{ajm}\,\sigma_m\otimes\sigma_b\bigr)
\]

%dsp-level:0
\[
\epsilon_{iab}\epsilon_{bjm}\,\sigma_a\otimes\sigma_m=\delta_{ij}\Sigma-\sigma_j\otimes\sigma_i;\qquad
\epsilon_{iab}\epsilon_{ajm}\,\sigma_m\otimes\sigma_b=\sigma_i\otimes\sigma_j-\delta_{ij}\Sigma
\]

%dsp-level:0
\[
[J_i,K_j]
\]

%dsp-level:0
\[
= \tfrac i2\Bigl(\delta_{ij}\Sigma-\sigma_j\otimes\sigma_i-\sigma_i\otimes\sigma_j+\delta_{ij}\Sigma\Bigr)
\]

%dsp-level:0
\[
= \tfrac i2\bigl(2\delta_{ij}\Sigma-\sigma_i\otimes\sigma_j-\sigma_j\otimes\sigma_i\bigr)
\]

%dsp-level:0
\[
= i\bigl(\delta_{ij}\Sigma-E^{ij}\bigr)
\]

%dsp-level:0
\[
i=j:\ \ i\bigl(\Sigma-\sigma_i\otimes\sigma_i\bigr)=i\!\!\sum_{k\ne i}\!\!\sigma_k\otimes\sigma_k;\qquad
i\ne j:\ \ \delta_{ij}=0\implies-iE^{ij}
\]

%dsp-level:0
\[
\Sigma,\ E^{ij}\in\mathcal E_{++};\ \ W_k\in\mathcal E_{+-};\ \ \text{direct sum}\ [\mathrm{CT1.1},\ \mathrm{CT1.2}]\implies[J_i,K_j]\notin\operatorname{span}\{W_k\}.\ \blacksquare
\]

\textbf{Phys.} Bracketing a spatial seal with a temporal seal produces symmetric spatial — metric-type — data, never weak data. Mixed virtual processes feed the gravitational sector: the algebraic seed of “quantum fluctuations sourcing geometry.”

\end{theorem}

\begin{theorem}[{\texorpdfstring{CT1.8 — Action of $W_i$ on $J_j$ and $K_j$}{CT1.8 - Action of on and}}]\label{tag:CT1-8}
\textbf{Thm.} $\ [W_i,J_j]=i\,\epsilon_{ijk}J_k;\qquad[W_i,K_j]=i\,\epsilon_{ijk}K_k;$ i.e. $\{J_i\}$, $\{K_i\}$: $SO(3)$ vectors under $\operatorname{span}\{W_i\}$.

\textbf{Req.} CT1.2; Pauli algebra.

\textbf{Pf.}

%dsp-level:0
\[
\text{suffices}:\ [W_3,J_1]=iJ_2,\quad[W_3,K_1]=iK_2
\]

%dsp-level:0
\[
[W_3,J_1]=\tfrac14\bigl[\sigma_0\otimes\sigma_3+\sigma_3\otimes\sigma_0,\ \ \sigma_2\otimes\sigma_3-\sigma_3\otimes\sigma_2\bigr]
\]

%dsp-level:0
\[
[\sigma_0\otimes\sigma_3,\ \sigma_2\otimes\sigma_3]=\sigma_2\otimes[\sigma_3,\sigma_3]=0;
\]

%dsp-level:0
\[
\qquad [\sigma_0\otimes\sigma_3,\ -\sigma_3\otimes\sigma_2]=-\sigma_3\otimes[\sigma_3,\sigma_2]=2i\,(\sigma_3\otimes\sigma_1)
\]

%dsp-level:0
\[
[\sigma_3\otimes\sigma_0,\ \sigma_2\otimes\sigma_3]=[\sigma_3,\sigma_2]\otimes\sigma_3=-2i\,(\sigma_1\otimes\sigma_3);\qquad
[\sigma_3\otimes\sigma_0,\ -\sigma_3\otimes\sigma_2]=0
\]

%dsp-level:0
\[
\implies[W_3,J_1]=\tfrac i2\bigl(\sigma_3\otimes\sigma_1-\sigma_1\otimes\sigma_3\bigr)=iJ_2
\]

%dsp-level:0
\[
[W_3,K_1]=\tfrac14\bigl[\sigma_0\otimes\sigma_3+\sigma_3\otimes\sigma_0,\ \ \sigma_0\otimes\sigma_1-\sigma_1\otimes\sigma_0\bigr]
\]

%dsp-level:0
\[
[\sigma_0\otimes\sigma_3,\ \sigma_0\otimes\sigma_1]=2i\,(\sigma_0\otimes\sigma_2);\qquad
[\sigma_0\otimes\sigma_3,\ -\sigma_1\otimes\sigma_0]=0=[\sigma_3\otimes\sigma_0,\ \sigma_0\otimes\sigma_1]
\]

%dsp-level:0
\[
[\sigma_3\otimes\sigma_0,\ -\sigma_1\otimes\sigma_0]=-[\sigma_3,\sigma_1]\otimes\sigma_0=-2i\,(\sigma_2\otimes\sigma_0)
\]

%dsp-level:0
\[
\implies[W_3,K_1]=\tfrac14\bigl(2i\,\sigma_0\otimes\sigma_2-2i\,\sigma_2\otimes\sigma_0\bigr)=iK_2;\qquad
\text{cyclicity}\implies\text{general form}.\ \blacksquare
\]

\textbf{Phys.} Total spin rotates both virtual sectors as vectors: the seal generators carry spin $1$ under the weak $\mathfrak{su}(2)$. The reality/virtuality bookkeeping is rotation-covariant — angular momentum is built into the ledger before any dynamics is introduced.

\end{theorem}

\begin{theorem}[{\texorpdfstring{CT1.9 — Sector closure table and the global algebra}{CT1.9 - Sector closure table and the global algebra}}]\label{tag:CT1-9}
\textbf{Thm.}

%dsp-level:0
\[
\begin{array}{c|c}
\text{bracket} & \text{closes in}\\ \hline
[W_i,W_j] & \operatorname{span}\{W_k\}=\mathcal E_{+-}\\
{}[W_i,J_j] & \operatorname{span}\{J_k\}=\mathcal E_{-+}\\
{}[W_i,K_j] & \operatorname{span}\{K_k\}=\mathcal E_{--}\\
{}[J_i,J_j] & \operatorname{span}\{W_k\}\\
{}[K_i,K_j] & \operatorname{span}\{W_k\}\\
{}[J_i,K_j] & (\mathcal E_{++})_{\mathrm{sym\ spatial}}\\
{}[\Sigma,J_i] & \operatorname{span}\{K_i\}\\
{}[\Sigma,K_i] & \operatorname{span}\{J_i\}\\
{}[\Sigma,W_i] & 0\\
{}[E_0,X] & 0\ \ \forall X
\end{array}
\]

%dsp-level:0
\[
\bigl(\mathfrak H(\mathbb C^4),\ (X,Y)\mapsto-i[X,Y]\bigr)\cong\mathfrak u(4)\ [\mathrm{C2}];\qquad
Z\bigl(\mathfrak u(4)\bigr)=\operatorname{span}_{\mathbb R}\{E_0\};
\]

%dsp-level:0
\[
\mathfrak{su}(4)=\operatorname{span}_{\mathbb R}\bigl\{\Sigma,\ \hat E^{ij},\ W_i,\ J_i,\ K_i\bigr\};\qquad
\dim_{\mathbb R}\mathfrak{su}(4)=1+5+3+3+3=15 .
\]

\textbf{Req.} C2; CT1.1–CT1.8; $\ \mathfrak T(\sigma_\mu\otimes\sigma_\nu)=4\,\delta_{\mu0}\delta_{\nu0}$.

\textbf{Pf.}

%dsp-level:0
\[
\text{rows }1\text{–}5:\ [\mathrm{CT1.4},\ \mathrm{CT1.8},\ \mathrm{CT1.6}];\qquad
\text{row }6:\ [\mathrm{CT1.7}];\qquad
\text{rows }7\text{–}9:\ [\mathrm{CT1.5}];\qquad
\text{row }10:\ [\mathrm{CT1.3}]
\]

%dsp-level:0
\[
\text{sector labels}:\ [\mathrm{CT1.2}]
\]

%dsp-level:0
\[
X,Y\in\mathfrak H(\mathbb C^4)\implies(-i[X,Y])^\dagger=-i[X,Y];
\]

%dsp-level:0
\[
\qquad \text{bilinear}\ \wedge\ \text{antisymmetric}\ \wedge\ \text{Jacobi}\ [\text{associative commutator}]
\]

%dsp-level:0
\[
\dim_{\mathbb R}\mathfrak H(\mathbb C^4)=16\ [\mathrm{CT1.1}]\implies\bigl(\mathfrak H(\mathbb C^4),-i[\cdot,\cdot]\bigr)=\mathfrak u(4)\quad[\mathrm{C2}]
\]

%dsp-level:0
\[
X\in Z\iff[X,Y]=0\ \ \forall Y\in\mathfrak H(\mathbb C^4)\implies X\ \text{scalar}\implies Z=\operatorname{span}_{\mathbb R}\{E_0\}\quad[\mathrm{CT1.3}]
\]

%dsp-level:0
\[
\mathfrak T(\sigma_\mu\otimes\sigma_\nu)=4\,\delta_{\mu0}\delta_{\nu0}\implies E_0\ \text{alone has trace}\implies\text{traceless complement}=\mathfrak{su}(4),\ \dim=16-1=15
\]

%dsp-level:0
\[
\mathfrak{su}(4)=\operatorname{span}_{\mathbb R}\{\Sigma,\hat E^{ij},W_i,J_i,K_i\}:\quad1+5+3+3+3=15.\ \blacksquare
\]

\textbf{Phys.} The full commutator structure closes: the $16$ pair observables form $\mathfrak u(4)=\mathfrak u(1)\oplus\mathfrak{su}(4)$, with the central $\mathfrak u(1)$ ($E_0$) as the uniform dark aspect and the $15$ traceless directions carrying gravity ($5$), weak ($3$), the two seal sectors ($3{+}3$) and the exchange scalar ($1$). Every force algebra later in the corpus is a subalgebra of this one $16$-dimensional algebra.

\end{theorem}

\begin{theorem}[{\texorpdfstring{CT1.10 — $SO(3)$ decomposition of $\mathcal E_{++}$}{CT1.10 - decomposition of}}]\label{tag:CT1-10}
\textbf{Thm.} Under $\operatorname{span}\{W_i\}$:

%dsp-level:0
\[
\mathcal E_{++}\ \longrightarrow\ \mathbf 1\oplus\mathbf 5\oplus\mathbf 1\qquad(7=1+5+1);
\]

%dsp-level:0
\[
E_0:\ \mathbf 1;\qquad\Sigma:\ \mathbf 1;\qquad
\hat E^{ij}:=E^{ij}-\tfrac13\delta^{ij}\Sigma:\ \mathbf 5\ (\text{spin }2),\quad\delta_{ij}\hat E^{ij}=0 .
\]

\textbf{Req.} CT1.1, CT1.3, CT1.5, CT1.8; symmetric traceless rank-$2$ tensor $=\ell=2$ representation [$SO(3)$ representation theory].

\textbf{Pf.}

%dsp-level:0
\[
\dim\mathcal E_{++}=7\quad[\mathrm{CT1.1}]
\]

%dsp-level:0
\[
[E_0,W_i]=0\ [\mathrm{CT1.3}];\ \ E_0\ \text{index-free}\implies E_0:\ \mathbf 1
\]

%dsp-level:0
\[
E^{ij}=\underbrace{\tfrac13\delta^{ij}\Sigma}_{\text{trace}}+\underbrace{\hat E^{ij}}_{\text{traceless}}\ \ (\text{unique split});\qquad
\delta_{ij}\hat E^{ij}=\Sigma-\tfrac13\cdot3\cdot\Sigma=0
\]

%dsp-level:0
\[
[\Sigma,W_i]=0\ [\mathrm{CT1.5}]\implies\Sigma:\ \mathbf 1
\]

%dsp-level:0
\[
[W_a,\sigma_i\otimes\sigma_j]=i\,\epsilon_{aik}\,\sigma_k\otimes\sigma_j+i\,\epsilon_{ajk}\,\sigma_i\otimes\sigma_k\quad[\mathrm{CT1.2},\ \text{Pauli algebra}]
\]

%dsp-level:0
\[
\implies[W_a,E^{ij}]=i\bigl(\epsilon_{aik}E^{kj}+\epsilon_{ajk}E^{ik}\bigr);
\]

%dsp-level:0
\[
\qquad [W_a,\hat E^{ij}]=[W_a,E^{ij}]-\tfrac13\delta^{ij}\underbrace{[W_a,\Sigma]}_{=0\ [\mathrm{CT1.5}]}=i\bigl(\epsilon_{aik}\hat E^{kj}+\epsilon_{ajk}\hat E^{ik}\bigr)
\]

%dsp-level:0
\[
\implies\hat E^{ij}:\ \text{traceless symmetric rank-}2\ \text{tensor}=\ell=2:\ \mathbf 5\quad[SO(3)\ \text{representation theory}]
\]

%dsp-level:0
\[
1+1+5=7=\dim\mathcal E_{++}\implies\mathcal E_{++}=\mathbf 1\oplus\mathbf 5\oplus\mathbf 1.\ \blacksquare
\]

\textbf{Phys.} Seen by rotations, the real sector contains exactly two scalars and one five-component spin-$2$ object. The spin-$2$ quintet is the graviton field of the corpus; with the three $W_i$ as shift, the $(7,3)$ real blocks supply the ten components of a Lorentzian metric (CT1.12): conformal ($1$) $+$ traceless metric ($5$) $+$ volume ($1$) $+$ shift ($3$).

\end{theorem}

\begin{theorem}[{\texorpdfstring{CT1.11 — The two scalars are distinct}{CT1.11 - The two scalars are distinct}}]\label{tag:CT1-11}
\textbf{Thm.} $\ E_0\in Z\bigl(\mathfrak H(\mathbb C^4)\bigr);\qquad\Sigma\notin Z\bigl(\mathfrak H(\mathbb C^4)\bigr);\qquad\forall c\in\mathbb R\setminus\{0\}:\ E_0\ne c\,\Sigma;$ hence the A8 assignment $E_0\mapsto F_{D^+}$, $\Sigma\mapsto F_H$ [CT0.12] is well-defined and non-degenerate.

\textbf{Req.} A8; CT0.12, CT1.3, CT1.5.

\textbf{Pf.}

%dsp-level:0
\[
[E_0,X]=0\ \ \forall X\implies E_0\in Z\quad[\mathrm{CT1.3}]
\]

%dsp-level:0
\[
[\Sigma,J_1]=4i\,K_1\ne0\ \ (K_1\ne0)\implies\Sigma\notin Z\quad[\mathrm{CT1.5}]
\]

%dsp-level:0
\[
c\ne0:\ \ X\in Z\iff cX\in Z\implies\bigl[E_0\in Z\ \wedge\ \Sigma\notin Z\bigr]\implies E_0\ne c\,\Sigma\ \ \forall c\ne0
\]

%dsp-level:0
\[
\implies E_0,\ \Sigma:\ \text{independent directions of}\ \mathcal E_{++};\qquad
E_0\mapsto F_{D^+}=\Lambda\,\mathbb I,\ \ \Sigma\mapsto F_H\quad[\mathrm{A8},\ \mathrm{CT0.12}].\ \blacksquare
\]

\textbf{Phys.} The two rotation scalars are provably different objects: $E_0$ is central — it cannot distinguish configurations — and supplies the uniform dark aspect $\Lambda\mathbb I$; $\Sigma$ is non-central — it couples to the seals and to geometry — and supplies the scalar (Higgs-type) aspect. The corpus’s two scalar directions are established as distinct before either receives a physical name.

\end{theorem}

\begin{theorem}[{\texorpdfstring{CT1.11a — The involution table of the four sectors}{CT1.11a - The involution table of the four sectors}}]\label{tag:CT1-11a}
\textbf{Thm.} $(0)$ \textbf{Dictionary.} A8's four sectors are CT1.1's blocks under the identification

%dsp-level:0
\[
\mathbf F_{(+,\mathcal R)}=F_G\oplus F_H\oplus F_{D^+}\ (5{+}1{+}1)\ \leftrightarrow\ \mathcal E_{++};
\]

%dsp-level:0
\[
\qquad \mathbf F_{(-,\mathcal R)}=F_{W^+}\oplus F_{W^-}\oplus F_Z\ (3)\ \leftrightarrow\ \mathcal E_{+-};
\]

%dsp-level:0
\[
\mathbf F_{(+,\mathcal V)}=F_{J^+}\ (3)\ \leftrightarrow\ \mathcal E_{-+};\qquad
\mathbf F_{(-,\mathcal V)}=F_{J^-}\ (3)\ \leftrightarrow\ \mathcal E_{--}
\]

— first label $=$ the $T$-grading ($\eta^\mu\eta^\nu$ [A3, C13]), $\mathcal R/\mathcal V=$ the $S$-grading (reality is the swap grading [CT0.7]); machine-verified on CT1.2's generator table (17 generators: $E_0,\Sigma,\hat E^{ij}{\times}5,W_i,J_i,K_i$). $(1)$ \textbf{Status of the two gradings.} $S=\operatorname{Ad}(\mathbb S)$ with $\mathbb S=\tfrac12(E_0+\Sigma)$ [CT0.7(1)] — an \textit{inner algebra automorphism} ($\mathbb S^2=\mathbb S$? no: $\mathbb S$ unitary involution, machine). $T$ is a \textit{pure grading}: the diagonal $\pm1$-map $\sigma_\mu\otimes\sigma_\nu\mapsto\eta^\mu\eta^\nu\,\sigma_\mu\otimes\sigma_\nu$ [A3] — $\neg$ multiplicative (witness: $T(\sigma_0\otimes\sigma_1)\,T(\sigma_0\otimes\sigma_2)=(-1)(-1)\ [\mathrm{C13}]\ \sigma_0\otimes(\sigma_1\sigma_2)\ne T(\sigma_0\otimes\sigma_1\sigma_2)$-type failure, machine), hence $\neg$ inner. $[S,T]=0$ on all 16 basis elements [CT1.1]; the grading group on $\mathfrak H(\mathbb C^4)$ is $\mathbb Z_2\times\mathbb Z_2$ with simultaneous eigenspaces $\mathcal E_{st}$, $\dim=(7,3,3,3)$. $(2)$ \textbf{Table of the declared reversals.} With $\omega$ [D2.6] and $\Theta_{\mathfrak L}$ [D2.7]:

%dsp-level:0
\[
\begin{array}{c|cccc|c}
& \mathbf F_{(+,\mathcal R)} & \mathbf F_{(-,\mathcal R)} & \mathbf F_{(+,\mathcal V)} & \mathbf F_{(-,\mathcal V)} & \text{edge/chain data}\\ \hline
S=\operatorname{Ad}(\mathbb S) & + & + & - & - & \text{reality grading [CT0.7]}\\
T\ (\text{grading}) & + & - & + & - & \eta^\mu\eta^\nu\ [\mathrm{A3},\mathrm{C13}]\\
\omega & \text{fixed} & W^+\!\leftrightarrow\!W^-;\ Z\ \text{fixed} & \neg\text{declared} & \neg\text{declared} & s(c)\leftrightarrow t(c);\ \text{charge}\mapsto-\text{charge [D2.6, EM.9];}\ \omega^2=1;\ R\mapsto R^{\mathrm T};\ \omega\notin\operatorname{Aut}(\Gamma),\ |V|\ge3\\
\Theta_{\mathfrak L} & \text{fixed} & \text{fixed} & \leftrightarrow & \leftrightarrow & \text{ledger order}\mapsto\text{reverse};\ J^+\leftrightarrow J^-;\ G(e)\mapsto G(e)^*\ (\alpha_e\widehat{ds}^{\,2}(e)\mapsto-\alpha_e\widehat{ds}^{\,2}(e),\ \sigma_e\ \text{fixed [D3.5]});\ \mathcal J\mapsto\mathcal J^*\ \text{[D3.6]};\ w\mapsto w^*\ \text{[CT2.8(1)]};\ \Theta^2=1\\
\omega\Theta_{\mathfrak L} & \text{fixed} & W^+\!\leftrightarrow\!W^- & \leftrightarrow & \leftrightarrow & \text{the fourth class}
\end{array}
\]

$\omega$ acts on the edge set, $\Theta_{\mathfrak L}$ on phase/ledger order — disjoint data $\Rightarrow$ they commute, and $\langle\omega,\Theta_{\mathfrak L}\rangle\cong\mathbb Z_2\times\mathbb Z_2$ with classes $(1,\omega,\Theta_{\mathfrak L},\omega\Theta_{\mathfrak L})$; $\omega$ is charge-odd [D2.6; covariance of the closed-core charge read: EM.9a], $\Theta_{\mathfrak L}$ phase-odd and charge-even [D2.7, D3.5, D3.6]. $(3)$ \textbf{Edge-class answer (the sector question at edge level).} Real edges split into $\sigma_e=\pm1$ [CT0.9: $\Phi=\{0,\pi\}$]; virtual edges carry $\neg$registration [A5, A7, D0.5: $\phi_e$ defined on $E_{\mathcal R}$ only] — a proposed $(\pm,\mathcal V)$ split has no read: the virtual class is single and ungraded (weight $1$, no phase). The two sign structures ($\sigma$ vs $(-1)^{\aleph}$) are independent [CT0.10]. $(4)$ \textbf{[READING] rows (recorded, $\neg$adopted).} External CPT vocabulary maps as: $C\leftrightarrow\omega$ (the charge-odd reversal [D2.6]; sign-level realization $\sigma_e\mapsto-\sigma_e$: EM.24(1) — \textbf{E}, forward pointer); $T\leftrightarrow\Theta_{\mathfrak L}$ (ledger-time reversal [D2.7] — $\neg$ anti-unitary: no anti-unitary operator is declared; conjugation acts on weights/phases [CT2.8(1), D3.6]); $P\leftrightarrow$ the alignment-pattern reversal $\mathcal P\mapsto E_{\mathcal R}\setminus\mathcal P$ (equivalently $\alpha_e\mapsto\pi-\alpha_e$) with the mirror classification of Def-EM.21/Def-EM.22, EM.23 — \textbf{E}, forward pointers: a pattern-level reversal on configuration data (spacetime parity as a Lorentzian field-algebra read has count $0$ in Parts 1–2; the orientation read $\chi=\operatorname{sign}(k)$ remains [CANDIDATE], not adopted); "CPT" $\leftrightarrow$ the derived composition $C\circ P\circ T$ on $\mathfrak C(\mathcal R)$ with unconditional $|w|^2$-invariance [EM.24(4)–(5) — \textbf{E}, forward pointer]; at the involution-table level its reversal part is the $\omega\Theta_{\mathfrak L}$ class of $(2)$. This is the exact amount of reversal structure the axioms declare — $\neg$ a Lorentzian CPT theorem (no Lorentzian field algebra, no spinor bundle in the inventory).

\textbf{Req.} A3, A5, A7, A8; C3; D0.5, D2.6, D2.7, D3.5, D3.6; CT0.7, CT0.7a, CT0.9, CT0.10, CT0.12, CT1.1, CT1.2, CT1.6, CT1.9, CT1.11 (Part 1); CT2.8(1) (\textbf{C} — forward pointer, re-verification); EM.9 (\textbf{E} — forward pointer, re-verification); exact $\mathbb Q(i)$ matrix arithmetic, exhaustive pattern enumeration, non-multiplicativity witness [math; machine: \texttt{checks/\allowbreak{}verify\_\allowbreak{}merge\_\allowbreak{}tierA.py} §E].

\textbf{Pf.}

%dsp-level:0
\[
(0):\ \text{CT1.2's table verbatim: }\ E_0,
\]

%dsp-level:0
\[
\Sigma,
\]

%dsp-level:0
\[
\hat E^{ij}\in\mathcal E_{++};
\]

%dsp-level:0
\[
\ W_i\in\mathcal E_{+-};
\]

%dsp-level:0
\[
\ J_i\in\mathcal E_{-+};
\]

%dsp-level:0
\[
\ K_i\in\mathcal E_{--}\ \wedge\ \text{CT0.12's aspect rows}\ [\text{machine}]
\]

%dsp-level:0
\[
(1):\ \mathbb S=\tfrac12(E_0+\Sigma)\ [\mathrm{CT0.7}(1)];\quad \operatorname{Ad}(\mathbb S)(X)=\mathbb S X\mathbb S=S(X)\ [\mathrm{CT0.7}(2)];
\]

%dsp-level:0
\[
\qquad T^2=\mathbb 1,\ S^2=\mathbb 1,\ [S,T]=0\ [\mathrm{A3},\ \mathrm{CT1.1}]
\]

%dsp-level:0
\[
T\ \neg\text{multiplicative}:\ T(\sigma_0\otimes\sigma_1)\,
\]

%dsp-level:0
\[
T(\sigma_0\otimes\sigma_2)
\]

%dsp-level:0
\[
= (-\sigma_0\otimes\sigma_1)(-\sigma_0\otimes\sigma_2)
\]

%dsp-level:0
\[
= \sigma_0\otimes i\sigma_3
\]

%dsp-level:0
\[
\ne -\,
\]

%dsp-level:0
\[
i\,
\]

%dsp-level:0
\[
\sigma_0\otimes\sigma_3
\]

%dsp-level:0
\[
= T(\sigma_0\otimes i\sigma_3)\ [\text{machine witness}]
\]

%dsp-level:0
\[
(2):\ D2.6/D2.7\ \text{verbatim on data rows};
\]

%dsp-level:0
\[
\quad J^\pm\ \text{blocks}:\ \Theta_{\mathfrak L}:\ J^+\leftrightarrow J^-\ [\mathrm{D2.7}]\ \wedge\ \mathbf F_{(+,\mathcal V)}
\]

%dsp-level:0
\[
= F_{J^+},
\]

%dsp-level:0
\[
\ \mathbf F_{(-,\mathcal V)}
\]

%dsp-level:0
\[
= F_{J^-}\ [\mathrm{A8}]\implies\text{sector exchange}
\]

%dsp-level:0
\[
\omega:\ \text{charge}
\]

%dsp-level:0
\[
\mapsto -\text{charge}\ [\mathrm{D2.6}]
\]

%dsp-level:0
\[
\implies W^+\leftrightarrow W^-\ \text{inside}\ \mathbf F_{(-,\mathcal R)}\ [\text{charge eigen-decomposition of the weak block: }F_{W^+}\oplus F_{W^-}\oplus F_Z,\ \mathrm{A8}];
\]

%dsp-level:0
\[
\quad \omega\ \text{on}\ J^\pm:\ \neg\text{declared}
\]

%dsp-level:0
\[
\omega^2=1:\ (R^{\mathrm T})^{\mathrm T}=R;
\]

%dsp-level:0
\[
\qquad \omega\notin\operatorname{Aut}(\Gamma):\ p(a)
\]

%dsp-level:0
\[
= b\ \wedge\ p(b)
\]

%dsp-level:0
\[
= a\ \forall(a,b)
\]

%dsp-level:0
\[
\in E:\ \bot\ \text{for}\ |V|
\]

%dsp-level:0
\[
\ge 3\ [\text{machine: exhaustive on the tetrahedron}]
\]

%dsp-level:0
\[
\text{disjoint supports}\ (\text{edge set vs phase/order})\implies[\omega,\Theta_{\mathfrak L}]=0;\qquad \omega^2=\Theta_{\mathfrak L}^2=1\implies\mathbb Z_2\times\mathbb Z_2.\ \blacksquare
\]

\textbf{Phys.} The corpus declares exactly two reversals and they commute: one flips orientation and charge on the edge set, the other flips ledger order and conjugates phases — four classes in total. The four force sectors of A8 are the joint grading atlas of the pair register; the Janus pair is the only sector pair a declared reversal exchanges, and spacetime parity has no force-sector house here — its pattern-level read lives on the alignment data (EM.23, EM.24 — \textbf{E}). The table is the complete reversal inventory at sector level: nothing further is derivable without a new declaration.

\end{theorem}

\begin{theorem}[{\texorpdfstring{CT1.11b — The native Clifford frames of the pair register and the weak-covariance pincer}{CT1.11b - The native Clifford frames of the pair register and the weak-covariance pincer}}]\label{tag:CT1-11b}
\textbf{Thm.} On $\mathfrak H(\mathbb C^4)$ with CT1.2's named basis and CT0.7's $\mathbb S$, $\Pi_{\mathcal R}=\tfrac14(3E_0+\Sigma)$, $\Pi_{\mathcal V}=\tfrac14(E_0-\Sigma)$. Call the $15$ non-identity $\sigma_\mu\otimes\sigma_\nu$ [A3, CT1.1] the \textit{Pauli products} — each Hermitian, squaring to $+I_4$, traceless. A \textit{frame} is a $4$-subset of pairwise-anticommuting Pauli products; for a frame $\mathcal F=(e_1,e_2,e_3,e_4)$ let $\gamma_{\mathcal F}:=e_1e_2e_3e_4$ (its \textit{pseudoscalar}; frame- and order-dependent: every transposition flips its sign). $(1)$ \textbf{Native frames.} There are exactly $30$ frames (unordered; $720$ as ordered tuples) among the $\binom{15}{4}=1365$ candidates. For every frame: $\gamma_{\mathcal F}$ is Hermitian, $\gamma_{\mathcal F}^2=I_4$, $\mathfrak T\gamma_{\mathcal F}=0$, $\{\gamma_{\mathcal F},e_i\}=0$; and the $16$ real products $\{e_{i_1}\cdots e_{i_r}\}_{r=0..4}$ are $\mathbb R$-linearly independent — the frame generates a real subalgebra isomorphic to $Cl(4,0)$ acting on $\mathbb C^4\cong\mathbb R^8$, \textbf{with no new axiom}: A3's Pauli facts and CT1.1's basis suffice. The signature is Euclidean ($e_i^2=+I_4$); the representative $\eta=(-,+,+,+)$ [C13] of A3's sign vector labels the $T$-grading only [CT1.1] and enters no Clifford relation. No canonical pseudoscalar exists: a distinguished "$\gamma$-element" would be a choice, $\neg$a derivation. $(2)$ \textbf{The pseudoscalar catalog (exhaustive).} All $15$ Pauli products occur as $\pm\gamma_{\mathcal F}$, \textbf{each from exactly $2$ frames}, in three species:

%dsp-level:0
\[
\begin{array}{c|c|c|c}
\text{species} & \gamma_{\mathcal F}\ (\text{up to sign}) & [\gamma_{\mathcal F},\mathbb S] & \text{commuting}\ W_i\\ \hline
\text{mixed}\ (6) & \sigma_0\otimes\sigma_i\ /\ \sigma_i\otimes\sigma_0 & \pm2iJ_i & W_i\ \text{only}\\
\text{diagonal}\ (3) & \sigma_i\otimes\sigma_i\in\mathcal E_{++} & 0 & W_i\ \text{only (anticomm.}\ W_j,W_k)\\
\text{off-diag. spatial}\ (6) & \sigma_i\otimes\sigma_j\ (i\ne j) & \mp2iK_k,\ k=6{-}i{-}j & \text{none (anticomm.}\ W_k)
\end{array}
\]

No pseudoscalar commutes with more than one $W_i$; none anticommutes with more than two. Witness frame $\mathcal F_0=(\sigma_1\otimes\sigma_2,\ \sigma_2\otimes\sigma_2,\ \sigma_3\otimes\sigma_2,\ \sigma_0\otimes\sigma_1)$:

%dsp-level:0
\[
\gamma_{\mathcal F_0}=\sigma_0\otimes\sigma_3=\operatorname{diag}(1,-1,1,-1);
\]

%dsp-level:0
\[
\qquad P^{(\mathcal F_0)}_\pm:=\tfrac12(I_4\mp\gamma_{\mathcal F_0})=\operatorname{diag}(0,1,0,1)\ /\ \operatorname{diag}(1,0,1,0)
\]

— both idempotent, orthogonal, summing to $I_4$, trace $2$, rank $2$ (exact forms, machine). $(3)$ \textbf{The weak-covariance pincer.} No grading involution of the pair register can be weak-covariant, from either side: \textbf{(a) commuting side.} $\{M:[M,W_i]=0\ \forall i\}=\operatorname{span}_{\mathbb C}\{\Pi_{\mathcal R},\Pi_{\mathcal V}\}$ exactly (complex dimension $2$ — Schur on the $3\oplus1$ decomposition [CT1.10, CT0.7]; machine nullspace solve). Its involutions are $a\Pi_{\mathcal R}+b\Pi_{\mathcal V}$ with $a^2=b^2=1$, traces $3a+b\in\{\pm4,\pm2\}$; a pseudoscalar needs trace $0$ (cyclic flip: $\mathfrak T(e_1e_2e_3e_4)=-\mathfrak T(e_2e_3e_4e_1)$) — $3a+b=0\ \wedge\ a^2=b^2=1$ has \textbf{no solution}. \textbf{(b) anticommuting side.} $\{M:\{M,W_i\}=0\ \forall i\}=\mathbb C\cdot\Pi_{\mathcal V}$ exactly (complex dimension $1$; machine nullspace solve): rank $1$, trace $1$ — singular, never an involution. Structural reason: an invertible $M$ anticommuting with every $W_i$ would make $\operatorname{Ad}(M)$ a Lie automorphism of $\operatorname{span}\{W_i\}$ equal to $-\mathrm{id}$ — but $-\mathrm{id}$ reverses the bracket ($[W_1,W_2]=iW_3\mapsto-iW_3\ne iW_3$) [CT1.4]: contradiction. \textbf{(c) exhaustive.} Every one of the $15$ attained pseudoscalars fails (a) and (b) explicitly, entrywise: for each, $[\gamma,W_i]\ne0$ for some $i$ and $\{\gamma,W_i\}\ne0$ for some $i$. \textbf{(d) carrier side.} On any $2$-dim carrier with weak action $W_i\mapsto\tfrac12\sigma_i$ (the weak seed locus: EW.32): $\{m:[m,\sigma_i]=0\ \forall i\}=\mathbb C\cdot I_2$ and $\{m:\{m,\sigma_i\}=0\ \forall i\}=\{0\}$ — the weak seed representation is irreducible: no nontrivial projector commutes with the weak action, and no grade anticommutes with it. $(4)$ \textbf{The native selectivity (what the weak force does grade).} The weak force grades the pair arena by \textbf{$S$-parity}: $W_i\Pi_{\mathcal V}=0$ exactly $\forall i$; $\sum_iW_i^2=2\Pi_{\mathcal R}$ [Cor CT0.7a]; $\bigcap_i\ker W_i=\operatorname{im}\Pi_{\mathcal V}$ exactly (the $1$-dim singlet line; machine) — a $3$-active/$1$-inactive split. This is a force-arena selectivity; a matter grade would have to live on a carrier on which the gauge action is trivial — the opposite pattern. $(5)$ \textbf{Non-alignment.} $[\gamma_{\mathcal F},\mathbb S]=0\iff\gamma_{\mathcal F}=\pm\,\sigma_i\otimes\sigma_i$: the $S$-definite pseudoscalars are exactly the three diagonal ones — they live in $\mathcal E_{++}$, the gravity/Higgs block [CT1.1, CT1.2], and are still $\neg$weak-covariant (each anticommutes with the two non-own $W$: e.g.\ $[\sigma_3\otimes\sigma_3,\ W_1]=2iE^{23}$ exactly — the weak rotation turns an $S$-definite pseudoscalar into gravitational shear). Under the swap, off-diagonal pseudoscalar grades rotate into the Janus seal sector: $[\sigma_0\otimes\sigma_3,\ \mathbb S]=2iJ_3$. The frame splits and the $S$-split provably do not align: $\operatorname{rank}(P^{(\mathcal F_0)}_\pm\Pi_{\mathcal R})=2$, $\operatorname{rank}(P^{(\mathcal F_0)}_\pm\Pi_{\mathcal V})=1$ — cross-ranks $(2,1,2,1)$ against the $S$-split's $(3,1)$.

\textbf{Req.} A3; C3; CT0.7, CT0.7a, CT1.1, CT1.2, CT1.4, CT1.10, CT1.11a (Part 1); EW.32 (\textbf{D}; seed-carrier locus — forward pointer, corroboration only: clause (3d) is self-contained by Schur [math]); exhaustive frame enumeration, exact nullspace solves over $\mathbb Q(i)$, Schur's lemma [math; machine: \texttt{checks/\allowbreak{}verify\_\allowbreak{}merge\_\allowbreak{}tierC.py}].

\textbf{Pf.}

%dsp-level:0
\[
(1):\ \gamma_{\mathcal F}^2=(e_1e_2e_3e_4)(e_1e_2e_3e_4)=(-1)^{\binom42}e_1^2e_2^2e_3^2e_4^2=(-1)^6I_4=I_4;
\]

%dsp-level:0
\[
\qquad \mathfrak T\gamma_{\mathcal F}=-\mathfrak T(e_2e_3e_4e_1)=-\mathfrak T\gamma_{\mathcal F}\implies\mathfrak T\gamma_{\mathcal F}=0\ [\text{cyclic}]
\]

%dsp-level:0
\[
\{\gamma_{\mathcal F},e_i\}=0:\ e_i\ \text{anticommutes with the three other factors},\ \text{commutes with itself};
\]

%dsp-level:0
\[
\qquad 30\ \text{frames}:\ \text{exhaustive enumeration of}\ 1365\ \text{candidates}\ [\text{machine}]
\]

%dsp-level:0
\[
\mathbb R\text{-independence of the}\ 16\ \text{products}:\ \text{exact elimination on}\ \mathbb Q\text{-flattened components, all}\ 30\ \text{frames}\ [\text{machine}]
\]

%dsp-level:0
\[
(2):\ \text{canonical-sign census over the}\ 30\ \text{frames}\ [\text{machine}];
\]

%dsp-level:0
\[
\quad \text{species}=\text{index patterns};
\]

%dsp-level:0
\[
\quad [\sigma_0\otimes\sigma_3,\mathbb S]=2iJ_3,\ [\sigma_3\otimes\sigma_3,W_1]=2iE^{23}:\ \text{direct products}\ [\text{machine}]
\]

%dsp-level:0
\[
\mathcal F_0:\ \text{direct multiplication}:\ \gamma_{\mathcal F_0}=\sigma_0\otimes\sigma_3;
\]

%dsp-level:0
\[
\quad P_\pm^2=P_\pm,\ P_+P_-=0,\ P_++P_-=I_4,\ \mathfrak T P_\pm=2\ [\text{machine}]
\]

%dsp-level:0
\[
(3a):\ \text{nullspace of the}\ 3\times16\ \text{complex linear system}\ [X,W_i]
\]

%dsp-level:0
\[
= 0\ \text{in the Pauli basis}\
\]

%dsp-level:0
\[
= \ \operatorname{span}\{\Pi_{\mathcal R},\Pi_{\mathcal V}\}\ [\text{machine, exact}];
\]

%dsp-level:0
\[
\quad \text{Schur on}\ 3\oplus1\ [\mathrm{CT1.10};\ \text{math}]
\]

%dsp-level:0
\[
\text{involutions in the commutant}:\ (a\Pi_{\mathcal R}+b\Pi_{\mathcal V})^2=I_4\iff a^2=b^2=1;
\]

%dsp-level:0
\[
\quad \mathfrak T=3a+b\in\{4,2,-2,-4\}:\ \neg\exists\ \text{trace-0 involution}
\]

%dsp-level:0
\[
(3b):\ \text{nullspace of}\ \{X,W_i\}=0\ =\ \mathbb C\cdot\Pi_{\mathcal V}\ [\text{machine}];\qquad W_i\Pi_{\mathcal V}=0\ [(4)]\implies\operatorname{rank}\Pi_{\mathcal V}=1:\ \text{singular}
\]

%dsp-level:0
\[
\text{invertible case}:\ \operatorname{Ad}(M)W_i
\]

%dsp-level:0
\[
= -W_i\implies\operatorname{Ad}(M)[W_1,W_2]
\]

%dsp-level:0
\[
= [-W_1,-W_2]
\]

%dsp-level:0
\[
= [W_1,W_2]
\]

%dsp-level:0
\[
= iW_3,
\]

%dsp-level:0
\[
\ \text{but}\ \operatorname{Ad}(M)(iW_3)
\]

%dsp-level:0
\[
= -iW_3:\ \bot\ [\mathrm{CT1.4}]
\]

%dsp-level:0
\[
(3c):\ \text{entrywise for all}\ 15\ [\text{machine}];\qquad (3d):\ 2\times2\ \text{nullspace solves}\ [\text{machine};\ \text{Schur}]
\]

%dsp-level:0
\[
(4):\ W_i\ \text{preserves}\ \vee^2,\wedge^2\ [\mathrm{CT0.7}(3)];
\]

%dsp-level:0
\[
\quad W_i|_{\wedge^2}=0\ (\text{singlet});
\]

%dsp-level:0
\[
\quad \sum_iW_i^2=\tfrac32E_0+\tfrac12\Sigma=2\Pi_{\mathcal R}\ [\text{Cor CT0.7a}];
\]

%dsp-level:0
\[
\quad \bigcap\ker W_i:\ \text{machine}
\]

%dsp-level:0
\[
(5):\ [\sigma_\mu\otimes\sigma_\nu,\mathbb S]=0\iff\sigma_\mu\otimes\sigma_\nu=\sigma_\nu\otimes\sigma_\mu\iff\mu=\nu\ [\text{machine over all}\ 15];
\]

%dsp-level:0
\[
\quad \text{ranks by exact elimination}.\ \blacksquare
\]

\textbf{Phys.} The pair register carries a full Euclidean Clifford structure natively — thirty frames, fifteen pseudoscalars, all of it derived from A3's Pauli facts with no declaration anywhere — and the weak force selects exactly one split inside it: the $3$-active/$1$-inactive $S$-parity grading. But the pincer is exact: from the commuting side the only weak-covariant candidates are the two $S$-projections, whose involutions all have nonzero trace; from the anticommuting side the candidate set collapses to a singular singlet line. No pseudoscalar grade of this register can ever behave as a weak-covariant matter grade — the obstruction is not a shortage of matrices but the absence of a spacetime-side carrier on which the gauge action is trivial, and that absence is a declaration the corpus has not made.

\textbf{Rmk (deficit relocation record).} The completeness audit's first deficit is hereby \textit{relocated inside Parts 1–2}: native frames exist (clause (1)); their grades cannot be weak-covariant (clause (3)); a Dirac-type operator would additionally require a per-edge frame assignment $e\mapsto\mathcal F_e$ coupled to the EW.32 transport — an enumerated would-be declaration of the audit layer, recorded here as \textbf{not adopted, no tag issued}. The reversal structure declared by the corpus remains exactly $\langle\omega,\Theta_{\mathfrak L}\rangle\cong\mathbb Z_2\times\mathbb Z_2$ [CT1.11a].

\end{theorem}

\phantomsection
\section[Tier 2 — Lyapunov Stability of the Pair State]{\texorpdfstring{Tier 2 — Lyapunov Stability of the Pair State}{Tier 2 - Lyapunov Stability of the Pair State}}

\begin{theorem}[{\texorpdfstring{L0 — The ledger–tilt coupling}{L0 - The ledger-tilt coupling}}]\label{tag:L0}
\textbf{Thm.} \textit{Dynamical scope (declared at use — the crystallization–tilt coupling; CT4.5a precedent: A4–A7 declare the criterion, the tick and the ledger, $\neg$ a state update at the seal):} every A4 crystallization event applies exactly one Gibbs tilt step [D0.3] to the pair state of that event. Under A5 ($\rho_\circ=\tfrac14I$), A6 ($\mathcal T(\rho)=U\rho U^\dagger$, $U\sim\mu_\star$; $|\mathfrak L(\tau_{n+1})|=|\mathfrak L(\tau_n)|+n_{\mathrm{real\_out}}(\tau_n)$), A7 ($\mathfrak L=\bigcup_{n\ge0}\mathfrak L(\tau_n)$, $|\mathfrak L(\tau_n)|=\sum_{k<n}n_{\mathrm{real\_out}}(\tau_k)$), D0.3/D0.4 ($\gamma:=1-4\varepsilon_q$), CT0.8: $(1)$ \textbf{Tick preserves $V$.} $\ V\bigl(\tilde K(\mathcal T(\rho))\bigr)=V\bigl(\tilde K(\rho)\bigr)$ for every draw [A6; L1(2)]. $(2)$ \textbf{Events apply the tilt.} At tick $\tau_n$ the A4 criterion selects $n_{\mathrm{real\_out}}(\tau_n)$ edges into the ledger [A6]; under the declared scope each selection maps its pair state

%dsp-level:0
\[
\rho\ \longmapsto\ \tilde{\mathcal T}_\gamma(\rho)=\frac{\rho\,e^{-\gamma\rho}}{\mathfrak T(\rho\,e^{-\gamma\rho})},\qquad
\gamma=1-4\varepsilon_q\in(0,1]\quad[\mathrm{D0.3};\ \mathrm{CT0.8}];
\]

$Z_\gamma>0$ on faithful $\rho$ [Cor L0a] $\implies\tilde{\mathcal T}_\gamma(\rho)$ faithful, $\succeq0$, $\mathfrak T=1$ [D0.3]; $\ \tilde{\mathcal T}_\gamma\ne\mathcal T$ [D0.3]. $(3)$ \textbf{Tilt contracts $V$.} $\ V'\le V$, equality iff $\rho=\tfrac14I$ [A5; L1(5)]. $(4)$ \textbf{Ledger counts tilt steps.} $\#\{\text{tilt steps at }\tau_n\}=n_{\mathrm{real\_out}}(\tau_n)$; cumulative count $|\mathfrak L(\tau_n)|=\sum_{k<n}n_{\mathrm{real\_out}}(\tau_k)$ [A6, A7 — exact integer recursion]. On one pair channel with $k$ events interleaved with arbitrary draws, the $V$-sequence is that of $\tilde{\mathcal T}_\gamma^{k}(\rho_0)$ — the interleaving is spectrally inert by $(1)$; per fixed $k$: $\ \mu_k=\bigl(1-\tfrac\gamma4\bigr)^{k}\mu_0+O(\mu_0^2)$ [L2, iterated]. $(5)$ \textbf{Contraction is the ledger dynamics.} The deficits of L1–L3 are realized once per ledger increment: the deficit fraction per tilt step is $\tfrac{\gamma(8-\gamma)}{16}$ of $V$ to exact leading order [L2], accumulating as $V\mapsto\bigl(1-\tfrac\gamma4\bigr)^{2k}V+O(\delta\rho_0^3)$ after $k$ recorded events; the inertial read $m_i^2(q)=\Delta S(q)$ [L3(4)] is the deficit read per increment.

\textbf{Req.} A4; A5; A6; A7; D0.2; D0.3; D0.4; D1.2; CT0.8; L1; L2 (forward pointers — the derivation order of \texttt{index/\allowbreak{}dependency\_\allowbreak{}tree.md} places this tag after L2); the declared crystallization–tilt dynamical scope (theorem-local; see Status); exact integer recursion, exact Haar covariance over $\mathbb Q(i)$, exact rational $\log/\exp$-enclosures [math; machine: \texttt{checks/\allowbreak{}verify\_\allowbreak{}L0\_\allowbreak{}ledger\_\allowbreak{}tilt.py} §A–§D].

\textbf{Pf.}

%dsp-level:0
\[
(1):\ K(\mathcal T\rho)=UKU^\dagger\ [\text{functional-calculus covariance}];
\]

%dsp-level:0
\[
\qquad \mathfrak T\bigl(K(\mathcal T\rho)\,\mathcal T\rho\bigr)=\mathfrak T(K\rho)\ [\text{cyclicity}]\implies\tilde K(\mathcal T\rho)=U\tilde KU^\dagger\implies V'=V\ [\mathrm{L1}(2)]
\]

%dsp-level:0
\[
(2):\ \varepsilon_q\in[0,\tfrac14)\ [\mathrm{CT0.8}]\implies\gamma=1-4\varepsilon_q\in(0,1]\ [\text{exact endpoints}:\ \varepsilon_q=0\mapsto\gamma=1];
\]

%dsp-level:0
\[
\qquad Z_\gamma\in(0,\,e^{-\gamma/4}]\ [\text{Cor L0a}]\implies\lambda_i'=\lambda_ie^{-\gamma\lambda_i}/Z_\gamma>0\ [\exp>0\ \text{on exact enclosures}]
\]

%dsp-level:0
\[
\implies\tilde{\mathcal T}_\gamma(\rho)\ \text{faithful},\ \succeq0,\ \mathfrak T\bigl(\tilde{\mathcal T}_\gamma(\rho)\bigr)=1\ [\mathrm{D0.3}\ \text{normalization}]
\]

%dsp-level:0
\[
(3):\ \mathrm{L1}(5)\ [\text{Jensen on}\ \phi_\gamma;\ \text{zero locus}\ \rho=\tfrac14I]
\]

%dsp-level:0
\[
(4):\ |\mathfrak L(\tau_0)|
\]

%dsp-level:0
\[
= 0\ [\mathrm{A7}];
\]

%dsp-level:0
\[
\quad\text{recursion}\ [\mathrm{A6}]
\]

%dsp-level:0
\[
\implies |\mathfrak L(\tau_n)|
\]

%dsp-level:0
\[
= \textstyle\sum_{k<n}n_{\mathrm{real\_out}}(\tau_k)\ [\mathrm{A7}]\ \text{by induction on}\ n\ [\text{exact integers; machine}]
\]

%dsp-level:0
\[
\operatorname{spec}\bigl(\mathcal T(\rho)\bigr)
\]

%dsp-level:0
\[
= \operatorname{spec}(\rho):\ \mathfrak T\bigl((U\rho U^\dagger)^m\bigr)
\]

%dsp-level:0
\[
= \mathfrak T(\rho^m),
\]

%dsp-level:0
\[
\ m
\]

%dsp-level:0
\[
= 1..4\ [\text{cyclicity; machine, exact over}\ \mathbb Q(i)]
\]

%dsp-level:0
\[
\implies V\circ\mathcal T\qquad = V\ [(1)]
\]

%dsp-level:0
\[
\implies V\bigl(\tilde{\mathcal T}\,\mathcal T\,\tilde{\mathcal T}\cdots(\rho_0)\bigr)=V\bigl(\tilde{\mathcal T}^{\,k}(\rho_0)\bigr)\ [\text{composition}];
\]

%dsp-level:0
\[
\qquad \mu_k
\]

%dsp-level:0
\[
= \bigl(1-\tfrac\gamma4\bigr)^k\mu_0+O(\mu_0^2):\ \text{induction on L2's}\ \mu'
\]

%dsp-level:0
\[
= \bigl(1-\tfrac\gamma4\bigr)\mu+O(\mu^2)\ [\text{coefficient-wise composition; machine},\ k\le3]
\]

%dsp-level:0
\[
(5):\ \mathrm{L2}\ \text{per step}:\ V'-V=-\tfrac{\gamma(8-\gamma)}{16}V+O(\delta\rho^3);\qquad
\text{accumulation at leading order from}\ (4).\ \blacksquare
\]

\end{theorem}

\begin{corollary}[{\texorpdfstring{L0a — tilt modular update}{L0a - tilt modular update}}]\label{tag:L0a}
Faithful $\rho$, $K:=-\log\rho$ [D0.3], $K':=K\bigl(\tilde{\mathcal T}_\gamma(\rho)\bigr)$, $Z_\gamma:=\mathfrak T(\rho\,e^{-\gamma\rho})$:

%dsp-level:0
\[
K'=K+\gamma\rho+\log Z_\gamma\cdot I;\qquad
Z_\gamma\in(0,\,e^{-\gamma/4}];\qquad
\mathfrak T(K'-K)=\gamma+4\log Z_\gamma\le0,
\]

with equality iff $\rho=\tfrac14I$; $\ \gamma=1:\ Z_1\le e^{-1/4}$, $\ \mathfrak T(K'-K)=1+4\log Z_1\le0$.

\textbf{Req (Cor).} A1; A5; D0.3; CT0.8; Jensen's inequality for $p\,e^{-\gamma p}$, exact rational $\log/\exp$-enclosures [math; machine: \texttt{checks/\allowbreak{}verify\_\allowbreak{}L0\_\allowbreak{}ledger\_\allowbreak{}tilt.py} §B–§C].

\textbf{Pf (Cor).}

%dsp-level:0
\[
\rho,\ \tilde{\mathcal T}_\gamma(\rho)\ \text{simultaneously diagonalizable}\ [\text{functional calculus on a commuting family}];
\]

%dsp-level:0
\[
\qquad \lambda_i'=\lambda_ie^{-\gamma\lambda_i}/Z_\gamma
\]

%dsp-level:0
\[
\implies -\log\lambda_i'
\]

%dsp-level:0
\[
= -\log\lambda_i+\gamma\lambda_i+\log Z_\gamma\ [\text{exact spectral identity; machine at two enclosure tolerances}]
\]

%dsp-level:0
\[
\implies K'
\]

%dsp-level:0
\[
= K+\gamma\rho+\log Z_\gamma\cdot I
\]

%dsp-level:0
\[
g(p):=p\,e^{-\gamma p}:\ \ g''(p)=e^{-\gamma p}\bigl(\gamma^2p-2\gamma\bigr)\le0\ \ \text{on}\ \ [0,1],\ \gamma\in(0,1]
\]

%dsp-level:0
\[
\quad [\gamma p\le1<2;\ \text{exact rational sign}]
\]

%dsp-level:0
\[
\text{Jensen}:\ Z_\gamma
\]

%dsp-level:0
\[
= \textstyle\sum_ig(\lambda_i)
\]

%dsp-level:0
\[
\le 4\,
\]

%dsp-level:0
\[
g\bigl(\tfrac14\bigr)
\]

%dsp-level:0
\[
= e^{-\gamma/4},
\]

%dsp-level:0
\[
\ \ \text{equality iff}\ \lambda_i
\]

%dsp-level:0
\[
= \tfrac14\ \forall i\iff\rho
\]

%dsp-level:0
\[
= \tfrac14I\ [\text{strict concavity}];
\]

%dsp-level:0
\[
\qquad Z_\gamma>0\ [\lambda_i>0]
\]

%dsp-level:0
\[
\mathfrak T(K'-K)=\gamma\,\mathfrak T\rho+4\log Z_\gamma=\gamma+4\log Z_\gamma\le\gamma+4\bigl(-\tfrac\gamma4\bigr)=0;
\]

%dsp-level:0
\[
\qquad \gamma=1:\ Z_1\le e^{-1/4},\ \ \mathfrak T(K'-K)=1+4\log Z_1.\ \blacksquare
\]

\textbf{Phys.} The Haar tick rotates the pair state without touching the Lyapunov function; the ledger grows exactly at the A4 events; under the declared crystallization dynamics every event is one Gibbs tilt step — the only contraction a pair state undergoes. The monotonicity of L1–L3 is thereby the per-increment dynamics of the ledger: ticks rotate, events contract, and $|\mathfrak L(\tau_n)|$ is the tilt-step count. Each tilt step updates the modular Hamiltonian by $K\mapsto K+\gamma\rho+\log Z_\gamma\cdot I$, whose trace deficit $\gamma+4\log Z_\gamma\le0$ vanishes only at the A5 vacuum. No particle identification.

\end{corollary}

\begin{theorem}[{\texorpdfstring{L1 — The Lyapunov function of the pair state}{L1 - The Lyapunov function of the pair state}}]\label{tag:L1}
\textbf{Thm.} Faithful pair states $\rho$ [A1: $\rho\in\mathfrak H(\mathbb C^4)$, $\rho\succeq0$, $\mathfrak T\rho=1$, $\operatorname{spec}\rho\subset(0,1]$]; with

%dsp-level:0
\[
\delta\rho:=\rho-\tfrac14I\ \ [\mathrm{A5}:\ \rho_\circ=\tfrac14I];\qquad
\mu:=\operatorname{spec}(\delta\rho)\subset\bigl[-\tfrac14,\tfrac34\bigr];\qquad
\textstyle\sum_i\mu_i=0\ \ [\mathfrak T\rho=1];
\]

%dsp-level:0
\[
\gamma:=1-4\varepsilon_q\in(0,1]\ \ [\mathrm{D0.3};\ \mathrm{A4},\ \mathrm{D0.4},\ \mathrm{CT0.8}];
\]

%dsp-level:0
\[
K:=-\log\rho;\qquad
\tilde K:=K-\mathfrak T(K\rho)\,I\ \ [\mathrm{D0.3}];\qquad
V(\tilde K):=\mathfrak T\bigl(\tilde K^2\bigr)\ \ge0;
\]

$(1)$ \textbf{Zero locus.} $V=0\iff\tilde K=0\iff\rho=\tfrac14I$ — the A5 vacuum. $(2)$ \textbf{Tick invariance.} $V\bigl(\tilde K(\mathcal T(\rho))\bigr)=V\bigl(\tilde K(\rho)\bigr)$ for every draw [A6: $\mathcal T(\rho)=U\rho U^\dagger$]. $(3)$ \textbf{Decay identity ($\gamma$-dependent).} Under the Gibbs tilt $\rho':=\tilde{\mathcal T}_\gamma(\rho)$ [D0.3], with $V':=V\bigl(\tilde K(\rho')\bigr)$:

%dsp-level:0
\[
V'-V=-\gamma\sum_i\phi_\gamma(\mu_i)+O(\delta\rho^4);\qquad
\phi_\gamma(x):=2x\log(1+4x)-\gamma x^2;
\]

%dsp-level:0
\[
\text{leading}:\ \ V'-V=-\gamma(8-\gamma)\textstyle\sum_i\mu_i^2+O(\delta\rho^3);
\]

%dsp-level:0
\[
\qquad \gamma=1:\ \ V'-V=-\textstyle\sum_i\phi(\mu_i)+O(\delta\rho^4),\ \ \phi:=\phi_1 .
\]

$(4)$ \textbf{$\phi_\gamma$-calculus.} $\phi_\gamma(0)=\phi_\gamma'(0)=0$; $\ \phi_\gamma=\phi_1+(1-\gamma)x^2$;

%dsp-level:0
\[
\phi_\gamma''(x)=\frac{14+16x-32x^2}{(1+4x)^2}+2(1-\gamma)\quad[\text{exact rational identity}];
\]

%dsp-level:0
\[
\phi_\gamma\ \text{convex on}\ W_\gamma\supseteq W:=\Bigl[\tfrac{1-2\sqrt2}{4},\tfrac{1+2\sqrt2}{4}\Bigr]\supseteq\bigl[-\tfrac14,\tfrac34\bigr]\quad\forall\gamma\in(0,1];
\]

%dsp-level:0
\[
\qquad \phi_\gamma>0\ \ \text{on}\ \ \bigl[-\tfrac14,\tfrac34\bigr]\setminus\{0\}.
\]

$(5)$ \textbf{Monotonicity.} Jensen: $\ \sum_i\phi_\gamma(\mu_i)\ \ge\ 4\,\phi_\gamma\bigl(\tfrac14\sum_i\mu_i\bigr)=0\ $; hence, on the faithful simplex, for every $\gamma\in(0,1]$:

%dsp-level:0
\[
V'\le V,\qquad V'=V\iff\tilde K=0\iff\rho=\tfrac14I .
\]

\textbf{Req.} A1; A4; A5; A6; D0.3, D0.4; CT0.8 (Part 1); exact formal-power-series arithmetic over $\mathbb Q$ with rational $\gamma$, exact $\mathbb Q(\sqrt2)$ order, Jensen's inequality, exact rational $\log$-enclosures, continuity–density extension [math; machine: \texttt{checks/\allowbreak{}verify\_\allowbreak{}L\_\allowbreak{}lyapunov.py}].

\textbf{Pf.}

%dsp-level:0
\[
(1):\ \tilde K=0\iff K=\mathfrak T(K\rho)I\iff\rho=e^{-c}I\iff\rho=\tfrac14I\ [\mathfrak T\rho=1];
\]

%dsp-level:0
\[
\qquad V=\mathfrak T(\tilde K^2)\ge0\ [\tilde K\ \text{Hermitian}],\quad V=0\iff\tilde K=0
\]

%dsp-level:0
\[
(2):\ K(\mathcal T\rho)=-\log\bigl(U\rho U^\dagger\bigr)=UKU^\dagger\ [\text{functional-calculus covariance}];
\]

%dsp-level:0
\[
\qquad \mathfrak T\bigl(K(\mathcal T\rho)\,\mathcal T\rho\bigr)=\mathfrak T\bigl(UK\rho U^\dagger\bigr)=\mathfrak T(K\rho)\ [\text{cyclicity}]
\]

%dsp-level:0
\[
\implies\tilde K(\mathcal T\rho)=U\tilde KU^\dagger\implies V'=V\ [\text{cyclicity}]
\]

%dsp-level:0
\[
(3):\ \rho\ \text{faithful}\implies\rho,
\]

%dsp-level:0
\[
\ \tilde{\mathcal T}_\gamma(\rho)\ \text{simultaneously diagonalizable}:\ \rho
\]

%dsp-level:0
\[
= \operatorname{diag}(\lambda),
\]

%dsp-level:0
\[
\ \lambda_i\qquad = \tfrac14+\mu_i;
\]

%dsp-level:0
\[
\qquad \tilde K=\operatorname{diag}\bigl(-\log\lambda_i-S\bigr),\ \ S:=\mathfrak T(K\rho)=-\textstyle\sum_j\lambda_j\log\lambda_j
\]

%dsp-level:0
\[
\lambda_i'=\lambda_ie^{-\gamma\lambda_i}/Z_\gamma,\ \ Z_\gamma=\textstyle\sum_j\lambda_je^{-\gamma\lambda_j};
\]

%dsp-level:0
\[
\qquad \log\lambda_i'=\log\lambda_i-\gamma\lambda_i-\log Z_\gamma;
\]

%dsp-level:0
\[
\qquad e^{-\gamma/4}\ \text{cancels in the ratio}\implies\text{exact}\ \mathbb Q[[t]]\ \text{data}
\]

%dsp-level:0
\[
\text{formal series along}\ \rho(t)
\]

%dsp-level:0
\[
= \tfrac14I+t\,
\]

%dsp-level:0
\[
\delta_0\ (\delta_0\ \text{Hermitian traceless, exact rational data}),
\]

%dsp-level:0
\[
\ \gamma\in\{\tfrac14,\tfrac12,\tfrac34,1\}:
\]

%dsp-level:0
\[
V'-V+\gamma\textstyle\sum_i\phi_\gamma(\mu_i(t))\ \text{has vanishing}\ t^0..t^3\ \text{coefficients}\ [\text{machine, exact}];
\]

%dsp-level:0
\[
\qquad t^4\ \text{coefficient}\ne0\ \text{generically}\implies O(\delta\rho^4)\ \text{sharp}
\]

%dsp-level:0
\[
\phi_\gamma(x)
\]

%dsp-level:0
\[
= 2x\bigl(4x-8x^2+\tfrac{64}3x^3-\dots\bigr)-\gamma x^2
\]

%dsp-level:0
\[
= (8-\gamma)x^2-16x^3+O(x^4)\implies\text{leading}:\ V'-V
\]

%dsp-level:0
\[
= -\gamma(8-\gamma)\textstyle\sum_i\mu_i^2+O(\delta\rho^3)
\]

%dsp-level:0
\[
(4):\ \phi_\gamma'(x)=2\log(1+4x)+\tfrac{8x}{1+4x}-2\gamma x;
\]

%dsp-level:0
\[
\qquad \phi_\gamma''(x)=\tfrac8{1+4x}+\tfrac8{(1+4x)^2}-2\gamma=\tfrac{14+16x-32x^2}{(1+4x)^2}+2(1-\gamma)\ [\text{common denominator}]
\]

%dsp-level:0
\[
\gamma\in(0,1]:\ 2(1-\gamma)\ge0\implies\phi_\gamma''\ge\phi_1''\ \text{pointwise}\implies W_\gamma\supseteq W;
\]

%dsp-level:0
\[
\qquad 16x^2-8x-7=0\iff x=\tfrac{1\pm2\sqrt2}4\ [\mathbb Q(\sqrt2),\ \text{exact}];
\]

%dsp-level:0
\[
\qquad W\supseteq[-\tfrac14,\tfrac34]\ [\sqrt2\ge1]
\]

%dsp-level:0
\[
\phi_\gamma=\phi_1+(1-\gamma)x^2;
\]

%dsp-level:0
\[
\qquad \phi_1\ \text{strictly convex on}\ W,\ \phi_1(0)=\phi_1'(0)=0\implies\phi_1>0\ \text{off}\ 0\implies\phi_\gamma>0\ \text{off}\ 0\ \text{on}\ [-\tfrac14,\tfrac34]
\]

%dsp-level:0
\[
(5):\ \text{Jensen on convex}\ \phi_\gamma:\ \tfrac14\sum_i\phi_\gamma(\mu_i)\ge\phi_\gamma\bigl(\tfrac14\sum_i\mu_i\bigr)=\phi_\gamma(0)=0;
\]

%dsp-level:0
\[
\qquad =\iff\mu=0\ \ [\textstyle\sum\mu=0;\ \text{strict convexity}]
\]

%dsp-level:0
\[
\gamma>0\implies V'-V=-\gamma\textstyle\sum_i\phi_\gamma(\mu_i)+O(\delta\rho^4)\le0\ \text{near the vacuum};
\]

%dsp-level:0
\[
\qquad \text{globally: exact rational}\ \log\text{-enclosure sweep},\ \gamma\in\{\tfrac14,\tfrac12,\tfrac34,1\}\ [\text{machine}]
\]

%dsp-level:0
\[
\text{extension to every}\ \gamma
\]

%dsp-level:0
\[
\in (0,1]:\ V'-V\ \text{continuous in}\ (\operatorname{spec}\rho,\gamma)\ \text{on the faithful simplex}\times(0,1]\ [\log,\exp,\ \text{rational maps}];
\]

%dsp-level:0
\[
\ \
\]

%dsp-level:0
\[
\le 0\ \text{on a dense rational set}\implies
\]

%dsp-level:0
\[
\le 0\ \text{everywhere}\ [\text{math}]
\]

%dsp-level:0
\[
V'=V\iff\tilde K=0:\ (\Leftarrow)\ (1);
\]

%dsp-level:0
\[
\qquad (\Rightarrow)\ \text{near vacuum by}\ (3),(4)\ \text{strictness}\ (\gamma>0);\ \ \text{sweep points strict off-vacuum}\ [\text{machine}].\ \blacksquare
\]

\textbf{Phys.} The pair state carries a natural nonnegative stability functional: the trace-square of its centered surprisal. Ledger ticks rotate it without changing it — the Haar conjugation of A6 is spectral, not contractive. The Gibbs tilt — a separate operation from the tick, determined by $\varepsilon_q$ — strictly lowers it, by an exact per-eigenvalue deficit $\gamma\phi_\gamma$, driving every faithful state toward the A5 vacuum. The tilt strength $\gamma=1-4\varepsilon_q$ is the normalized distance from the crystallization threshold to the vacuum: at $\varepsilon_q=0$ (sharp crystallization) the tilt is at maximum strength; at the vacuum threshold $\varepsilon_q=\tfrac14$ it vanishes. The tilt is grounded in the ledger–tick coupling [L0]: it is applied once per A4 crystallization event, and the ledger counts the steps.

\end{theorem}

\begin{theorem}[{\texorpdfstring{L2 — The temporal decay law}{L2 - The temporal decay law}}]\label{tag:L2}
\textbf{Thm.} Under L1's setting, per tilt step [D0.3; L1(3)], to exact leading order:

%dsp-level:0
\[
V'-V=-\Bigl[1-\bigl(1-\tfrac\gamma4\bigr)^{2}\Bigr]V+O(\delta\rho^3)=-\frac{\gamma(8-\gamma)}{16}\,V+O(\delta\rho^3);\qquad
V'=\bigl(1-\tfrac\gamma4\bigr)^{2}V+O(\delta\rho^3);
\]

%dsp-level:0
\[
\gamma=1:\ \ V'-V=-\tfrac7{16}\,V+O(\delta\rho^3),\ \ V'=\tfrac9{16}V+O(\delta\rho^3);\qquad
\gamma\to0:\ \ V'=V\ \ (\text{identity-tilt limit});
\]

%dsp-level:0
\[
\mu'=\bigl(1-\tfrac\gamma4\bigr)\mu+O(\mu^2)\quad[\text{eigenvalue contraction}];
\]

%dsp-level:0
\[
\qquad \gamma=1:\ \ \mu'=\tfrac34\mu;
\]

%dsp-level:0
\[
\qquad 1-\bigl(1-\tfrac\gamma4\bigr)^2=\frac{\gamma(8-\gamma)}{16}\ \ [\text{exact rationals}].
\]

\textbf{Req.} A1; A4; A5; D0.3, D0.4; CT0.8 (Part 1); L1; exact formal-series arithmetic [math; machine: \texttt{checks/\allowbreak{}verify\_\allowbreak{}L\_\allowbreak{}lyapunov.py} §B].

\textbf{Pf.}

%dsp-level:0
\[
K=-\log\bigl(\tfrac14I+\delta\rho\bigr)=\ln4\,I-4\,\delta\rho+O(\delta\rho^2);\qquad \mathfrak T(K\rho)=\ln4+O(\delta\rho^2)\ \ [\mathfrak T\delta\rho=0]
\]

%dsp-level:0
\[
\implies\tilde K
\]

%dsp-level:0
\[
= K-\mathfrak T(K\rho)I
\]

%dsp-level:0
\[
= -4\,
\]

%dsp-level:0
\[
\delta\rho+O(\delta\rho^2)
\]

%dsp-level:0
\[
\implies V
\]

%dsp-level:0
\[
= 16\,
\]

%dsp-level:0
\[
\mathfrak T(\delta\rho^2)+O(\delta\rho^3)
\]

%dsp-level:0
\[
= 16\textstyle\sum_i\mu_i^2+O(\delta\rho^3)\ \ [\gamma\text{-free}]
\]

%dsp-level:0
\[
\mathrm{L1}(3):\ V'-V=-\gamma\textstyle\sum_i\phi_\gamma(\mu_i)+O(\delta\rho^4);\qquad \gamma\phi_\gamma(x)=\gamma(8-\gamma)x^2+O(x^3)
\]

%dsp-level:0
\[
\implies V'-V
\]

%dsp-level:0
\[
= -\gamma(8-\gamma)\textstyle\sum_i\mu_i^2+O(\delta\rho^3)
\]

%dsp-level:0
\[
= -\frac{\gamma(8-\gamma)}{16}\,
\]

%dsp-level:0
\[
V+O(\delta\rho^3)
\]

%dsp-level:0
\[
= -\Bigl[1-\bigl(1-\tfrac\gamma4\bigr)^2\Bigr]V+O(\delta\rho^3)
\]

%dsp-level:0
\[
\lambda'=\lambda e^{-\gamma\lambda}/Z_\gamma:\quad\text{series}:\ \ \mu'=\bigl(1-\tfrac\gamma4\bigr)\mu+O(\mu^2)\ [\text{exact coefficient; machine}];
\]

%dsp-level:0
\[
\qquad V'=16\bigl(1-\tfrac\gamma4\bigr)^2\textstyle\sum_i\mu_i^2+O(\delta\rho^3)
\]

%dsp-level:0
\[
\gamma=1:\ \ 1-\tfrac14=\tfrac34,\ \ \bigl(\tfrac34\bigr)^2=\tfrac9{16},\ \ 1-\tfrac9{16}=\tfrac7{16},\ \ \gamma(8-\gamma)=7;\qquad
\gamma\to0:\ \ \bigl(1-\tfrac\gamma4\bigr)^2\to1.\ \blacksquare
\]

\textbf{Phys.} Near the vacuum, every tilt step contracts the Lyapunov function by a $\gamma$-dependent exact rational factor: $\bigl(1-\tfrac\gamma4\bigr)^2$ of $V$ survives, $\tfrac{\gamma(8-\gamma)}{16}$ of it is lost. At $\gamma=1$ ($\varepsilon_q=0$, sharp crystallization) the contraction is $\tfrac9{16}$ and the deficit exactly $\tfrac7{16}$ of $V$; at the vacuum threshold ($\gamma\to0$) the contraction is $1$ and the deficit $0$. The $\tfrac7{16}$ is not a stipulation: it is the decay rate at sharp crystallization, derived from $\varepsilon_q=0$ [D0.4]. The tilt is grounded in the ledger–tick coupling [L0]: one step per recorded crystallization event, so the decay law is the per-increment dynamics of the ledger.

\end{theorem}

\begin{theorem}[{\texorpdfstring{L3 — Entropy deficit and the inertial mass reading}{L3 - Entropy deficit and the inertial mass reading}}]\label{tag:L3}
\textbf{Thm.} Under L1, L2: $(1)$ \textbf{Entropy deficit.} With $\delta\rho=\tfrac14\sum_ar_aG_a$ on the Pauli-tensor basis $\{G_a\}=\{\sigma_\mu\otimes\sigma_\nu\}_{(\mu,\nu)\ne(0,0)}$ [C1; $\mathfrak T(G_aG_b)=4\delta_{ab}$: CT1.1], $|\delta\rho|^2:=\mathfrak T(\delta\rho^2)$:

%dsp-level:0
\[
\Delta S_{\mathrm{pair}}:=2\,|\delta\rho|^2=\tfrac12\,|r|^2;
\]

%dsp-level:0
\[
\qquad S\bigl(\rho\,\big\|\,\tfrac14I\bigr)=\Delta S_{\mathrm{pair}}+O(\delta\rho^3)
\]

%dsp-level:0
\[
\qquad \bigl[S(\rho\|\sigma):=\mathfrak T\bigl(\rho(\log\rho-\log\sigma)\bigr)\bigr]
\]

$(2)$ \textbf{Lyapunov identification.} $\ \tilde K=-4\,\delta\rho+O(\delta\rho^2);\qquad V(\tilde K)=16\,|\delta\rho|^2+O(\delta\rho^3)=8\,\Delta S_{\mathrm{pair}}+O(\delta\rho^3).$ $(3)$ \textbf{Deficit decay ($\gamma$-dependent).} $\ V'-V=-\dfrac{\gamma(8-\gamma)}{2}\,\Delta S_{\mathrm{pair}}+O(\delta\rho^3);\qquad\gamma=1:\ \ V'-V=-\tfrac72\,\Delta S_{\mathrm{pair}}+O(\delta\rho^3).$ $(4)$ \textbf{Inertial read.} $\ m_i^2(q):=\Delta S(q)\quad[\Delta S(q):\ \text{the patch deficit, L4};\ \text{dimensionless, nats};\ \neg\text{a particle identification}].$

\textbf{Req.} A1; A5; C1; D0.3; CT1.1 (Part 1); L1, L2; exact formal-series arithmetic, relative-entropy eigenvalue identity [math; machine: \texttt{checks/\allowbreak{}verify\_\allowbreak{}L\_\allowbreak{}lyapunov.py} §B].

\textbf{Pf.}

%dsp-level:0
\[
(1):\ \delta\rho\ \text{Hermitian}\wedge\ \mathfrak T\delta\rho
\]

%dsp-level:0
\[
= 0\implies\delta\rho
\]

%dsp-level:0
\[
= \tfrac14\sum_ar_aG_a,
\]

%dsp-level:0
\[
\ \ r_a
\]

%dsp-level:0
\[
= \mathfrak T(\delta\rho\,G_a)\in\mathbb R\ \ [\mathrm{C1},\ \mathrm{CT1.1}:\ \mathfrak T(G_aG_b)=4\delta_{ab}]
\]

%dsp-level:0
\[
|\delta\rho|^2=\mathfrak T(\delta\rho^2)=\tfrac1{16}\sum_{ab}r_ar_b\,\mathfrak T(G_aG_b)=\tfrac14|r|^2\implies2|\delta\rho|^2=\tfrac12|r|^2
\]

%dsp-level:0
\[
S\bigl(\rho\|\tfrac14I\bigr)
\]

%dsp-level:0
\[
= \mathfrak T(\rho\log\rho)+\ln4
\]

%dsp-level:0
\[
= \textstyle\sum_i\bigl(\tfrac14+\mu_i\bigr)\log(1+4\mu_i)\ [\text{exact eigenvalue identity}]
\]

%dsp-level:0
\[
= 2\textstyle\sum_i\mu_i^2+O(\mu^3)
\]

%dsp-level:0
\[
= \Delta S_{\mathrm{pair}}+O(\delta\rho^3)
\]

%dsp-level:0
\[
(2):\ \mathrm{L2}\ \text{Pf}:\ \tilde K=-4\,\delta\rho+O(\delta\rho^2);\qquad V=16\,\mathfrak T(\delta\rho^2)+O(\delta\rho^3)=8\,\Delta S_{\mathrm{pair}}+O(\delta\rho^3)\ \ [(1)]
\]

%dsp-level:0
\[
(3):\ \mathrm{L2}:\ V'-V
\]

%dsp-level:0
\[
= -\tfrac{\gamma(8-\gamma)}{16}V+O(\delta\rho^3)
\]

%dsp-level:0
\[
= -\tfrac{\gamma(8-\gamma)}{16}\cdot8\,
\]

%dsp-level:0
\[
\Delta S_{\mathrm{pair}}+O(\delta\rho^3)
\]

%dsp-level:0
\[
= -\tfrac{\gamma(8-\gamma)}{2}\Delta S_{\mathrm{pair}}+O(\delta\rho^3)
\]

%dsp-level:0
\[
\gamma=1:\ \ \tfrac{\gamma(8-\gamma)}{2}=\tfrac72
\]

%dsp-level:0
\[
(4):\ \text{definition};\qquad\text{a read},\ \neg\text{an identification}.\ \blacksquare
\]

\textbf{Phys.} The pair state's departure from the vacuum is one object under three exact leading reads: a squared Bloch radius, an entropy deficit against the maximally mixed pair, and the Lyapunov functional itself — locked by exact rational factors that depend on the tilt strength $\gamma=1-4\varepsilon_q$. The inertial read names the deficit $m_i^2(q)=\Delta S(q)$: mass-as-missing-information per channel, dimensionless in nats — a reading, not a particle claim; no pole-mass identification is asserted. The tilt is grounded in the ledger–tick coupling [L0]; the patch-level deficit $\Delta S(q)$ with its integer degree factor is derived in L4.

\end{theorem}

\begin{theorem}[{\texorpdfstring{L4 — The patch entropy deficit}{L4 - The patch entropy deficit}}]\label{tag:L4}
\textbf{Thm.} Vertex $q$ of finite incidence (theorem-local scope); incident pair channels $\{(q,p)\}$ [A1]; $n(q)=\partial^{\leftarrow}(q)+\partial^{\rightarrow}(q)$ [D0.8]; $n_{\mathcal R}(q)$ [A4]; $n_{\mathcal V}(q):=n(q)-n_{\mathcal R}(q)$ [D0.1]. Patch deficit:

%dsp-level:0
\[
\Delta S(q)\ :=\ \sum_{(q,p)\in\mathcal V}S\bigl(\rho_{qp}\,\big\|\,\tfrac14I\bigr)
\]

%dsp-level:0
\[
\qquad [\text{sum over the virtual incident channels of}\ q;\ \ S(\cdot\|\cdot)\ \text{as in L3}(1)]
\]

$(1)$ \textbf{Channel sum.} $\ \Delta S(q)=\sum_{(q,p)\in\mathcal V}\Delta S_{\mathrm{pair}}(pq)+O(\delta\rho^3)=2\sum_{(q,p)\in\mathcal V}|\delta\rho_{pq}|^2+O(\delta\rho^3)$ [L3(1) per channel]. $(2)$ \textbf{Mean form.} For $n_{\mathcal V}(q)\ge1$, with $\langle\Delta S_{\mathrm{pair}}\rangle_q:=\frac1{n_{\mathcal V}(q)}\sum_{(q,p)\in\mathcal V}\Delta S_{\mathrm{pair}}(pq)$:

%dsp-level:0
\[
\Delta S(q)\
\]

%dsp-level:0
\[
= \ n_{\mathcal V}(q)\,
\]

%dsp-level:0
\[
\langle\Delta S_{\mathrm{pair}}\rangle_q+O(\delta\rho^3)\
\]

%dsp-level:0
\[
= \ 2\,
\]

%dsp-level:0
\[
n_{\mathcal V}(q)\,
\]

%dsp-level:0
\[
\langle|\delta\rho|^2\rangle_q+O(\delta\rho^3)\
\]

%dsp-level:0
\[
= \ \tfrac12\,
\]

%dsp-level:0
\[
n_{\mathcal V}(q)\,
\]

%dsp-level:0
\[
\langle|r|^2\rangle_q+O(\delta\rho^3);
\]

$n_{\mathcal V}(q)=0\implies\Delta S(q)=0$ [empty sum]. $(3)$ \textbf{Patch Lyapunov read.} $\ V_{\mathrm{patch}}(q):=\sum_{(q,p)\in\mathcal V}V(\tilde K_{pq})=8\,\Delta S(q)+O(\delta\rho^3)$; under per-channel tilt steps [D0.3]: $\ V_{\mathrm{patch}}'-V_{\mathrm{patch}}=-\frac{\gamma(8-\gamma)}2\,\Delta S(q)+O(\delta\rho^3)$ [L3(2), (3) summed]. $(4)$ \textbf{Inertial read, grounded.} $\ m_i^2(q)=\Delta S(q)$ [L3(4)] $=2\,n_{\mathcal V}(q)\,\langle|\delta\rho|^2\rangle_q+O(\delta\rho^3)$ — an integer multiplicity $n_{\mathcal V}(q)\in\mathbb Z_{\ge0}$ times the mean channel deficit; the read vanishes to leading order iff $n_{\mathcal V}(q)=0$ or every virtual-channel deviation vanishes.

\textbf{Req.} A1; A4; A5; C1; D0.1, D0.3, D0.8; L3; the finite-incidence scope (theorem-local); exact formal-series arithmetic over $\mathbb Q$, exact integer bookkeeping [math; machine: \texttt{checks/\allowbreak{}verify\_\allowbreak{}L0\_\allowbreak{}ledger\_\allowbreak{}tilt.py} §E].

\textbf{Pf.}

%dsp-level:0
\[
(1):\ S\bigl(\rho\,\|\,\tfrac14I\bigr)=\Delta S_{\mathrm{pair}}+O(\delta\rho^3)\ \text{per channel}\ [\mathrm{L3}(1),\ \text{sharp cubic remainder}];
\]

%dsp-level:0
\[
\qquad \text{finite sum over}\ n_{\mathcal V}(q)\ \text{channels preserves the order}\ [\text{linearity}]
\]

%dsp-level:0
\[
(2):\ n_{\mathcal V}\cdot\tfrac1{n_{\mathcal V}}\textstyle\sum=\sum\ [n_{\mathcal V}\ge1];
\]

%dsp-level:0
\[
\qquad |\delta\rho|^2=\tfrac14|r|^2\implies2\cdot\tfrac14=\tfrac12\ [\mathrm{L3}(1);\ \text{exact rationals}];
\]

%dsp-level:0
\[
\qquad n_{\mathcal V}=0:\ \text{empty sum}=0
\]

%dsp-level:0
\[
(3):\ \textstyle\sum_eV(\tilde K_e)=\sum_e\bigl(8\,\Delta S_{\mathrm{pair}}(e)+O(\delta\rho_e^3)\bigr)=8\,\Delta S(q)+O(\delta\rho^3)\ [\mathrm{L3}(2)];
\]

%dsp-level:0
\[
\textstyle\sum_e\bigl(V'(\tilde K_e)-V(\tilde K_e)\bigr)
\]

%dsp-level:0
\[
= -\tfrac{\gamma(8-\gamma)}2\sum_e\Delta S_{\mathrm{pair}}(e)+O(\delta\rho^3)
\]

%dsp-level:0
\[
= -\tfrac{\gamma(8-\gamma)}2\Delta S(q)+O(\delta\rho^3)\ [\mathrm{L3}(3);\ \text{channel-wise tilt, D0.3}]
\]

%dsp-level:0
\[
(4):\ (2)\ \text{into the L3}(4)\ \text{definition};
\]

%dsp-level:0
\[
\qquad n(q),\ n_{\mathcal R}(q)\ \text{integer counts}\ [\mathrm{D0.8},\ \mathrm{A4}]\implies n_{\mathcal V}(q)\in\mathbb Z_{\ge0}.\ \blacksquare
\]

\textbf{Phys.} The patch-level entropy deficit is the pair deficit counted with its multiplicity: every virtual channel incident to $q$ contributes its own missing information against the A5 vacuum, and a fully crystallized patch contributes zero. This grounds the symbol $\Delta S(q)$ of the inertial read $m_i^2(q)=\Delta S(q)$ [L3(4)]: mass-as-missing-information is an integer degree-deficit factor $n_{\mathcal V}(q)=n(q)-n_{\mathcal R}(q)$ times the mean channel deficit at leading order. A reading, not an identification.

\end{theorem}

\begin{theorem}[{\texorpdfstring{L5 — The second-order sector flow of the tilt}{L5 - The second-order sector flow of the tilt}}]\label{tag:L5}
\textbf{Thm.} Under L1's setting, with $c:=1-\tfrac\gamma4$, $\ d:=\tfrac{\gamma(8-\gamma)}8$ [exact rationals in $\gamma$; L2]: $(1)$ \textbf{Quadratic flow.} With $\delta\rho':=\tilde{\mathcal T}_\gamma(\rho)-\tfrac14I$ [D0.3]:

%dsp-level:0
\[
\delta\rho'\ =\ c\,\delta\rho\ -\ d\,\Bigl(\delta\rho^2-\tfrac14|\delta\rho|^2E_0\Bigr)\ +\ O(\delta\rho^3);\qquad
\mathfrak T\Bigl(\delta\rho^2-\tfrac14|\delta\rho|^2E_0\Bigr)=0;\qquad
\delta\rho^2\succeq0 .
\]

$(2)$ \textbf{Single-sector landing.} $\delta\rho\in\mathcal E_{st}$ (a single $(S,T)$ sector [C3; CT1.2]) $\implies S(\delta\rho^2)=\delta\rho^2\ \wedge\ T(\delta\rho^2)=\delta\rho^2:\ \delta\rho^2\in\mathcal E_{++}$; the traceless part lies in $\operatorname{span}\{\Sigma,\hat E^{ij}\}$ [CT1.2, CT1.10]; the square channel of $\mathcal E_{++}$ closes in $\mathcal E_{++}$ — the quadratic transport inside the real block is one-directional. $(3)$ \textbf{Partition function.} $\ \log Z_\gamma=-\tfrac\gamma4-d\,|\delta\rho|^2+O(\delta\rho^3)=-\tfrac\gamma4-\tfrac{\gamma(8-\gamma)}{16}\Delta S_{\mathrm{pair}}+O(\delta\rho^3)$ [Cor L0a; L3(1)]; $\ \gamma=1:\ \log Z_1=-\tfrac14-\tfrac78|\delta\rho|^2=-\tfrac14-\tfrac7{16}\Delta S_{\mathrm{pair}}+O(\delta\rho^3)$. $(4)$ \textbf{Modular trace deficit.} $\ \mathfrak T(K'-K)=\gamma+4\log Z_\gamma$ [Cor L0a] $=-4d\,|\delta\rho|^2+O(\delta\rho^3)=-\tfrac{\gamma(8-\gamma)}4\Delta S_{\mathrm{pair}}+O(\delta\rho^3)\le0$ near the vacuum; $\ \gamma=1:\ -\tfrac74\Delta S_{\mathrm{pair}}+O(\delta\rho^3)$. $(5)$ \textbf{Sharp-crystallization table ($\gamma=1$).} $\ c=\tfrac34;\quad d=\tfrac78;\quad1-c^2=\tfrac7{16}\ [\mathrm{L2}];\quad\tfrac d2=\tfrac7{16};\quad2d=\tfrac72\quad[\text{exact rationals}]$.

\textbf{Req.} A1; A5; C1, C3; D0.3, D0.4; CT0.8; CT1.1, CT1.2, CT1.10; L0 (Cor L0a); L1, L2, L3; exact formal-series arithmetic over $\mathbb Q$, exhaustive square-channel census over $\mathbb Q(i)$ (all $256$ basis pairs) [math; machine: \texttt{checks/\allowbreak{}verify\_\allowbreak{}L0\_\allowbreak{}ledger\_\allowbreak{}tilt.py} §F].

\textbf{Pf.}

%dsp-level:0
\[
(1):\ \rho\,e^{-\gamma\rho}=e^{-\gamma/4}\bigl(\tfrac14E_0+\delta\rho\bigr)e^{-\gamma\,\delta\rho}\ [\text{functional calculus on a commuting family}]
\]

%dsp-level:0
\[
=e^{-\gamma/4}\Bigl[\tfrac14E_0+c\,\delta\rho-\gamma\bigl(1-\tfrac\gamma8\bigr)\delta\rho^2+O(\delta\rho^3)\Bigr],
\]

%dsp-level:0
\[
\qquad \gamma\bigl(1-\tfrac\gamma8\bigr)=\tfrac{\gamma(8-\gamma)}8=d\ [\text{exact rationals}]
\]

%dsp-level:0
\[
Z_\gamma=e^{-\gamma/4}\bigl[1-d\,|\delta\rho|^2+O(\delta\rho^3)\bigr]\ [\mathfrak T\delta\rho=0,\ \mathfrak T\delta\rho^2=|\delta\rho|^2];
\]

%dsp-level:0
\[
\qquad \rho'=\bigl(\tfrac14E_0+c\,\delta\rho-d\,\delta\rho^2\bigr)\bigl(1+d\,|\delta\rho|^2\bigr)+O(\delta\rho^3)
\]

%dsp-level:0
\[
\implies\delta\rho'
\]

%dsp-level:0
\[
= c\,
\]

%dsp-level:0
\[
\delta\rho-d\bigl(\delta\rho^2-\tfrac14|\delta\rho|^2E_0\bigr)+O(\delta\rho^3)\ [\text{coefficient-wise machine:}\ \gamma\in\{\tfrac14,\tfrac12,\tfrac34,1\}\times5\ \text{exact rational families}];
\]

%dsp-level:0
\[
\mathfrak T\bigl(\delta\rho^2-\tfrac14|\delta\rho|^2E_0\bigr)=|\delta\rho|^2-|\delta\rho|^2=0;\qquad
\delta\rho^2\succeq0\ [\text{Hermitian square}]
\]

%dsp-level:0
\[
(2):\ X^2
\]

%dsp-level:0
\[
= \textstyle\sum_ax_a^2G_a^2+\sum_{a<b}x_ax_b\{G_a,G_b\}:\ \text{the commutator channel cancels in squares}\ [\text{ordered double sum}]
\]

%dsp-level:0
\[
\text{square channel}:\ \{G_a,G_b\}
\]

%dsp-level:0
\[
\ne 0
\]

%dsp-level:0
\[
\implies T(G_aG_b)
\]

%dsp-level:0
\[
= \vartheta_a\vartheta_b\,
\]

%dsp-level:0
\[
G_aG_b\ [\text{exhaustive: all}\ 136\ \text{nonvanishing pairs of the}\ 256;\ \text{machine}];
\]

%dsp-level:0
\[
\text{single sector}\implies\vartheta_a\vartheta_b=t^2=+1;\qquad
S=\operatorname{Ad}(\mathbb S)\ \text{inner}\implies S(G_aG_b)=S(G_a)S(G_b)=s^2G_aG_b\ [\mathrm{CT0.7}(2)]
\]

%dsp-level:0
\[
\implies T(X^2)=X^2\ \wedge\ S(X^2)=X^2\implies X^2\in\mathcal E_{++}\ [\mathrm{CT1.1}];
\]

%dsp-level:0
\[
\qquad \mathcal E_{++}=\operatorname{span}\{E_0,\Sigma,\hat E^{ij}\}\ [\mathrm{CT1.2},\ \mathrm{CT1.10}]\implies\text{traceless part}\subseteq\operatorname{span}\{\Sigma,\hat E^{ij}\}
\]

%dsp-level:0
\[
\text{one-directionality}: (+,+)\cdot(+,+)\ \text{square channel}
\]

%dsp-level:0
\[
\mapsto (+,+):\ \text{the same census restricted to}\ \mathcal E_{++}\ [\text{machine}];
\]

%dsp-level:0
\[
\qquad \text{the}\ \neg\text{multiplicativity witness of CT1.11a}(1)\ \text{lies on the commutator channel}\ (\{G_a,G_b\}=0)\ \text{and is reproduced}
\]

%dsp-level:0
\[
(3):\ \log Z_\gamma
\]

%dsp-level:0
\[
= -\tfrac\gamma4+\log\bigl(1-d|\delta\rho|^2+O(\delta\rho^3)\bigr)
\]

%dsp-level:0
\[
= -\tfrac\gamma4-d|\delta\rho|^2+O(\delta\rho^3)\ [\log(1+u)=u+O(u^2)]
\]

%dsp-level:0
\[
\Delta S_{\mathrm{pair}}=2|\delta\rho|^2\ [\mathrm{L3}(1)];\qquad\tfrac d2=\tfrac{\gamma(8-\gamma)}{16};\qquad\gamma=1:\ d=\tfrac78,\ \tfrac d2=\tfrac7{16}
\]

%dsp-level:0
\[
(4):\ \mathfrak T(K'-K)
\]

%dsp-level:0
\[
= \gamma+4\log Z_\gamma\ [\text{Cor L0a}]
\]

%dsp-level:0
\[
= \gamma+4\bigl(-\tfrac\gamma4-d|\delta\rho|^2\bigr)+O(\delta\rho^3)
\]

%dsp-level:0
\[
= -4d|\delta\rho|^2+O(\delta\rho^3)
\]

%dsp-level:0
\[
-4d=-\tfrac{\gamma(8-\gamma)}2;
\]

%dsp-level:0
\[
\qquad -4d|\delta\rho|^2=-2d\,\Delta S_{\mathrm{pair}}=-\tfrac{\gamma(8-\gamma)}4\Delta S_{\mathrm{pair}};
\]

%dsp-level:0
\[
\qquad \gamma=1:\ -\tfrac74\Delta S_{\mathrm{pair}};
\]

%dsp-level:0
\[
\qquad \le0\ \text{near the vacuum}\ [\Delta S_{\mathrm{pair}}\ge0,\ \text{zero only at}\ \tfrac14I:\ \mathrm{L1}(1),\ \mathrm{L3}(1)]
\]

%dsp-level:0
\[
(5):\ \text{exact Fraction table}\ [\text{machine}].\ \blacksquare
\]

\textbf{Phys.} To second order the tilt does two things: it shrinks the deviation uniformly by $c=1-\tfrac\gamma4$, and it transports the traceless square of the deviation with exact weight $d=\tfrac{\gamma(8-\gamma)}8$ into the $S$-even, $T$-even real block — a quadratic sector flow that lands wholly in $\mathcal E_{++}$ for any single-sector input and stays there under further squaring. At sharp crystallization the transport weight is exactly $\tfrac78$, the partition function reads $-\tfrac14-\tfrac7{16}\Delta S_{\mathrm{pair}}$, and the modular trace deficit reads $-\tfrac74\Delta S_{\mathrm{pair}}$ — one exact rational family with the Lyapunov deficit $\tfrac7{16}$ [L2]. A reading, not an identification.

\end{theorem}

\phantomsection
\section[Closure exports of Part 1 §F (consumed by Part 2)]{\texorpdfstring{Closure exports of Part 1 §F (consumed by Part 2)}{Closure exports of Part 1 SF (consumed by Part 2)}}

%dsp-level:0
\[
\begin{array}{lll}
\mathrm{CT0.1}\ \text{dichotomy} & \mathrm{CT0.5}\ D_{\mathrm{tick}}(n)=2^{n+1} & \mathrm{CT1.1}\ (7,3,3,3)\\
\mathrm{CT0.2}\ \rho_\circ,\ \ln4 & \mathrm{CT0.6(a)}\ C_{\mathrm{cap}}=16,\ O(n);\ \delta h\le4\,\delta\tau & \mathrm{CT1.2}\ \text{basis}\ \{E_0,\Sigma,W_i,J_i,K_i,\hat E^{ij}\}\\
\mathrm{CT0.3}\ J^\pm\ \text{increments} & \mathrm{CT0.7}\ \Pi_{\mathcal R},\Pi_{\mathcal V};\ \mathbf F_{(\pm,\cdot)}=\mathcal E_{st} & \mathrm{CT1.3}\text{–}\mathrm{CT1.9}\ \text{closure table};\ \mathfrak u(4)=\mathfrak u(1)\oplus\mathfrak{su}(4)\\
\mathrm{CT0.4}\ \text{real}=\text{registered} & \mathrm{CT0.8}\ 0\le\varepsilon_q<\tfrac14 & \mathrm{CT1.10}\ \mathcal E_{++}=\mathbf1\oplus\mathbf5\oplus\mathbf1\\
\mathrm{CT0.9}\ \Phi=\{0,\pi\};\ 1\ \text{edge}=1\ \text{bit} & \mathrm{CT0.11(a)}\ n\le3;\ \mathbb C^3\ \text{slots};\ \mathrm{CT0.11b}\ n=2d\ \text{(bidirectional)} & \mathrm{CT1.11}\ E_0\nparallel\Sigma\\
\mathrm{CT0.10}\ (P,\pi)\ \text{independent} & \mathrm{CT0.12}\ 3{+}3{+}5{+}3{+}1{+}1=16 & \mathrm{CT0.13}\ \mathcal A(n)\\
\mathrm{CT0.14}\ \dim\mathcal K_n=(6,16,17,25) & &
\end{array}
\]

$\blacksquare$

\phantomsection
\section[§G — Master index of results]{\texorpdfstring{§G — Master index of results}{SG - Master index of results}}

%tbl-level:3
{\scriptsize\setlength{\tabcolsep}{4pt}
\begin{longtable}{>{\raggedright\arraybackslash}p{\dimexpr0.1834\textwidth-2\tabcolsep\relax}|>{\raggedright\arraybackslash}p{\dimexpr0.7540\textwidth-2\tabcolsep\relax}|>{\raggedright\arraybackslash}p{\dimexpr0.0626\textwidth-2\tabcolsep\relax}}
\toprule
Tags & Content & Locus \\
\addlinespace[2pt]
\midrule
\endfirsthead
\toprule
Tags & Content & Locus \\
\addlinespace[2pt]
\midrule
\endhead
\bottomrule
\endlastfoot
A1–A9 · C1–C13 · D0.1–D4.3 · P1–P5 & axioms · conventions · the definition layer · calibration & Part 1 §A–§E \\
CT0.1–CT0.14 (+CT0.6a, CT0.7a, CT0.11a/b) & dichotomy, vacuum, seals, tick load $2^{n+1}$, $C_{\mathrm{cap}}=\allowbreak 16$, swap grading, $\varepsilon_q$ range, $\Phi=\allowbreak \{0,\pi\}$, $n\le\allowbreak 3$, $\mathcal A(n)$, mask blocks & Part 1 §F \\
CT1.0–CT1.11 (+CT1.11a/b) & signature classification, $(S,\allowbreak T)$ sector decomposition, generator basis, $\mathfrak{su}(2)$ closures, commutator tables, global algebra, involution table, native Clifford frames and the weak-covariance pincer & Part 1 §F \\
L0–L5 (+Cor L0a) & ledger–tilt coupling, tilt modular update, Lyapunov function, decay law $\gamma(8-\gamma)/16$, entropy deficit and the inertial read, patch deficit, second-order sector flow & Part 1 §F \\
CT1.12–CT1.29 (+CT1.12a, CT1.19a/b, CT1.29a) · CT2.1–CT2.11 & metric emergence, conformal rescaling, linearized spin-2, continuum limit, effective action, $T_{\mu\nu}$, field equations, conservation, $\Lambda=\allowbreak D^+$, entropy/calibration, one free scale & Part 2 §C \\
EW.1–EW.36 (+EW.5a, EW.36a) & electroweak: masks, $Q=\allowbreak T_3+\allowbreak Y/2$, abelian $\mathfrak u(1)$, scalar trace mode, rung-3 wrapping, $m_W$, neutral mass matrix, $\theta_W$, $m_W=\allowbreak m_Z\cos\theta_W$, $e=\allowbreak g\sin\theta_W$, spectrum & Part 2 §D \\
EM.1–EM.26 (+EM.9a/b, Def-EM.20–22) & electromagnetic: rung-2 $\mathfrak u(1)$, photon masslessness/polarizations/entropy, end-charge, quantization, conservation, $G(e)$, dynamic per-edge phase, alignment patterns and parity, CPT bridge, spectrum & Part 2 §E \\
CT3.1–CT3.22 (+CT3.10a, CT3.17a, CT3.22a) & colour: $N=\allowbreak 3$ forced, $\mathfrak u(3)=\allowbreak \mathfrak{su}(3)\oplus\mathfrak u(1)$, rung-3 algebra, $f_{abc}$, gluons $=\allowbreak 8$, masslessness, baryon/meson singlets, sealing grammar, confinement & Part 2 §F \\
CT4.1–CT4.13 (+CT4.1a, CT4.4a, CT4.5a, CT4.7a) & saturation: $O(n)\to D^-$, absorbing lock, $w_{D^-}$, dark contrast, entropy triangle, bounds, global closure, falsifiers, audit & Part 2 §G \\
Def-H.1–Def-H.43 · FT.1–FT.145 (structural) · FT.153–FT.155 & the fermion layer: measurement primitives, the fermion definition, floors, spin/stability, flicker, spectra, Clifford/branches, path integrals, family channels, carrier/parity, rung one; the interface definitions and their derived interface theorems & Part 3 \\
Def-BM.1–2 · BM.1–BM.6 & colour-sealed composites: baryons/mesons, junction grammar, observability, the forced family group, invariance laws, census consumption, construction specs & Part 4 \\
the 25 Category-B cells & the computed data cells and their consumption block & Census \\
\end{longtable}
}

$\blacksquare$


% Part2_Bosonic_era.md — sanitized presentation layer (generated; see source/CONVERSION_LOG.md)
\chapter{\texorpdfstring{$\langle qu(b(i(t|t))ick)\rangle$ — \textbf{Part 2 · Bosonic\_era} — Geometry · Electroweak · Electromagnetic · Colour · Saturation}{Qubitick - **Part 2 - Bosonic _era** - Geometry - Electroweak - Electromagnetic - Colour - Saturation}}

\textbf{Dep.} Part 1 (axioms A1–A9, conventions C1–C13, the definition layer D0.1–D4.3, calibration P1–P5, pair canon CT0.*/CT1.1–CT1.11b, ledger dynamics L0–L5) — cited by tag at use. Format: Part 1 §0.4. \textbf{Namespace.} This part owns: §C geometry CT1.12–CT1.29a, CT2.1–CT2.11; §D electroweak EW.1–EW.36a; §E electromagnetic EM.1–EM.26, Def-EM.20–Def-EM.22; §F colour CT3.1–CT3.22a; §G saturation CT4.1–CT4.13; §Cal-2 (the bosonic registry note). The definition layer (D2.1–D4.3) and the master symbol table are Part 1 §C–§D; section letters §C–§G preserve the classical canon's file letters. \textbf{DAG law.} Citations resolve to Part 1, to earlier tags of this part, or to marked forward pointers (restatements at use; corroboration-only rows). No citation of Parts 3–4. \textbf{Internal section dependencies.} §C depends on Part 1 §F (and on Part 1 §C's Tier 2–4 definitions). §D depends on §C. §E depends on §C, §D (EM.18 alone, restated at use). §F depends on Part 1 §F and §C; its two borrowed electroweak/EM results (colour masslessness EW.26, coefficient stability EM.18) are restated with full arguments at use, so §F may be read before §D, §E. §G depends on Part 1 §F, §C, §D (EW.27), §F.

\phantomsection
\section[§Cal-2 — Calibration note (bosonic era)]{\texorpdfstring{§Cal-2 — Calibration note (bosonic era)}{SCal-2 - Calibration note (bosonic era)}}

The P4 core ($g,\ g';\ h_0;\ c_W;\ \Lambda;\ \varkappa,\ \beta,\ \alpha_\Omega$) and P3 ($G_N$, $\kappa^2$ — named combinations of $\mathfrak h$) are Part 1 §E. Logged spectral datum of this part: $\alpha_{\mathrm{EM}}^{(2)}$ — the EM-block scale read [EM.17]: operatorially defined; no theorem of the corpus fixes its numerical value under the dynamic-phase regime [EM.25(8)]; it is a spectral datum of $\mathbf F$ at the solution level, $\neg$a fitted constant and $\neg$a P4 core member; no result consumes its value — every edge-level phase read uses the registration data $\alpha_e\in\{0,\pi\}$ [Def-EM.20, EM.20]. Any future use that needs a value requires a logged declaration first.

\phantomsection
\section[§C — Geometry: Emergent Metric, Masks, Field Equations,…]{\texorpdfstring{§C — Geometry: Emergent Metric, Masks, Field Equations, Calibration}{SC - Geometry: Emergent Metric, Masks, Field Equations, Calibration}}

\textit{Dep.: Part 1 (§A–§F). Notation: $\hat h_{ij}$ — spatial traceless shear (D2.1); $\hat h^{(4)}_{\mu\nu}$ — 4-dim $\eta$-traceless perturbation (CT1.18) — distinct symbols. Machine: \texttt{checks/\allowbreak{}verify\_\allowbreak{}C.py} (exact rational jet calculus; linearization via truncated series); CT1.29a: \texttt{checks/\allowbreak{}verify\_\allowbreak{}merge\_\allowbreak{}tierA.py}; CT1.19b/CT1.29: \texttt{checks/\allowbreak{}verify\_\allowbreak{}promotion\_\allowbreak{}}}.*

\phantomsection
\section[§C.I — Emergent Metric Geometry]{\texorpdfstring{§C.I — Emergent Metric Geometry}{SC.I - Emergent Metric Geometry}}

\begin{theorem}[{\texorpdfstring{CT1.12 — Metric emergence from sector components}{CT1.12 - Metric emergence from sector components}}]\label{tag:CT1-12}
\textbf{Thm.} With D2.1 data $\bigl(a_{00},q_{ij},w_i\bigr)\in\mathcal E_{++}\oplus\mathcal E_{+-}$:

%dsp-level:0
\[
\delta^{ij}\hat h_{ij}=0;
\]

%dsp-level:0
\[
\qquad \Omega=e^{\alpha_\Omega a_{00}}>0\ \ \forall a_{00}\in\mathbb R;
\]

%dsp-level:0
\[
\qquad \bar g_{ij}\succ0\iff\forall a:\ 1+\varkappa\lambda_a(q)>0\quad(\Leftarrow\ \varkappa>0\wedge\varkappa\|q\|_{\mathrm{op}}<1);
\]

%dsp-level:0
\[
\bar g_{ij}\succ0\implies\operatorname{sig}(\bar g_{\mu\nu})=\operatorname{sig}(g_{\mu\nu})=(-,+,+,+);
\]

%dsp-level:0
\[
\begin{array}{c|c|c}
\text{sector object} & \# & \text{metric role}\\ \hline
a_{00} & 1 & \text{conformal factor}\ \Omega\\
\Sigma_{\mathrm{tr}} & 1 & \text{spatial volume mode}\\
\hat h_{ij} & 5 & \text{spin-2 shear}\\
w_i & 3 & \text{shift}\ N_i=\beta w_i\\ \hline
\text{total} & 10 & \#\{g_{\mu\nu}=g_{\nu\mu}\}=\tfrac{4\cdot5}{2}
\end{array}
\]

\textbf{Req.} A3; D2.1; CT1.1, CT1.10; Schur identity $\det\begin{pmatrix}A&B\\ C&D\end{pmatrix}=\det D\cdot\det(A-BD^{-1}C)$, Sylvester criterion [linear algebra].

\textbf{Pf.}

%dsp-level:0
\[
\dim\mathcal E_{++}=7=1\oplus5\oplus1\ \text{with coefficients}\ (a_{00},\hat h_{ij},\Sigma_{\mathrm{tr}})\ [\mathrm{CT1.1},\mathrm{CT1.10}];
\]

%dsp-level:0
\[
\qquad \dim\mathcal E_{+-}=3\ \text{with coefficients}\ w_i\ [\mathrm{CT1.1}]
\]

%dsp-level:0
\[
7+3=10=\tfrac{4\cdot5}{2}\implies\text{dictionary is a bijection}
\]

%dsp-level:0
\[
\delta^{ij}\hat h_{ij}=\delta^{ij}q_{ij}-\tfrac13\delta^{ij}\delta_{ij}\Sigma_{\mathrm{tr}}=\Sigma_{\mathrm{tr}}-\tfrac13\cdot3\cdot\Sigma_{\mathrm{tr}}=0\ [\mathrm{D2.1}];
\]

%dsp-level:0
\[
\qquad q_{ij}=\hat h_{ij}+\tfrac13\delta_{ij}\Sigma_{\mathrm{tr}}\ \text{unique (trace functional)}
\]

%dsp-level:0
\[
a_{00}\in\mathbb R\implies\Omega=e^{\alpha_\Omega a_{00}}>0\implies\Omega^2>0
\]

%dsp-level:0
\[
\bar g=\begin{pmatrix}-1&\varkappa N_j\\ \varkappa N_i&\bar g_{ij}\end{pmatrix}:\quad
\det\bar g=\det(\bar g_{ij})\cdot\bigl(-1-\varkappa^2N_i(\bar g^{-1})^{ij}N_j\bigr)\quad[\text{Schur}]
\]

%dsp-level:0
\[
\bar g_{ij}\succ0\implies\det(\bar g_{ij})>0\ \wedge\ (\bar g^{-1})^{ij}\succ0\implies N(\bar g^{-1})N\ge0\implies(\cdot)\le-1<0\implies\det\bar g<0
\]

%dsp-level:0
\[
\bar g_{ij}\succ0:\ 3\ \text{positive directions};\quad\det\bar g<0:\ \text{4th eigenvalue}<0\implies\operatorname{sig}(\bar g)=(-,+,+,+)
\]

%dsp-level:0
\[
\bar g_{ij}=\delta_{ij}+\varkappa q_{ij}:\ \operatorname{spec}=\{1+\varkappa\lambda_a(q)\}\implies\succ0\iff\forall a:\ 1+\varkappa\lambda_a>0
\]

%dsp-level:0
\[
g_{\mu\nu}=\Omega^2\bar g_{\mu\nu},\ \Omega^2>0:\ \operatorname{spec}(g)=\Omega^2\operatorname{spec}(\bar g)\implies\operatorname{sig}(g)=(-,+,+,+).\ \blacksquare
\]

\textbf{Rmk (scale hygiene).} $\varkappa$: embedding scale (D2.1, C4); $\ \varkappa\ne\kappa$ (CT2.9); $\ [\varkappa]=[q]^{-1}$; $\neg\exists$ assumed relation $\varkappa\sim\kappa$.

\textbf{Phys.} The ten real sector coefficients of a pair — conformal amplitude, spatial trace, five traceless shears, three shifts — are exactly the ten components of a Lorentzian metric in conformal gauge. Gravity is not added to the theory: the metric is a repackaging of the $S$-even operator content, with signature $(-,+,+,+)$ guaranteed in the perturbative regime $\varkappa\|q\|_{\mathrm{op}}<1$.

\end{theorem}

\begin{theorem}[{\texorpdfstring{CT1.12a — The pair-channel audit of the metric dictionary: tiling, block split, and the three discrete read classes}{CT1.12a - The pair-channel audit of the metric dictionary: tiling, block split, and the three discrete read classes}}]\label{tag:CT1-12a}
\textbf{Thm.} $(1)$ \textbf{Tiling.} The $16$ pair-channel registers [CT0.6] tile by aspects:

%dsp-level:0
\[
(G_5+W_3+H_1)+(J^+_3+J^-_3)+D^+_1=16\qquad[\mathrm{CT0.12},\ \mathrm{CT1.1},\ \mathrm{CT1.2}]
\]

matching the $(S,T)$ block decomposition $7+3+3+3$ [CT1.1] under the dictionary $(G,H,D^+)\leftrightarrow\mathcal E_{++}$, $W\leftrightarrow\mathcal E_{+-}$, $J^+\leftrightarrow\mathcal E_{-+}$, $J^-\leftrightarrow\mathcal E_{--}$ [CT0.12; CT1.11a(0)]. $(2)$ \textbf{Block split of the dictionary.} CT1.12's dictionary rows decompose as

%dsp-level:0
\[
g_{\mu\nu}\ \leftrightarrow\ \hat h_{ij}\ (5,\ G:\ \text{spin-2 shear})\ \oplus\ \Sigma_{\mathrm{tr}}\ (1,\ H:\ \text{volume mode})\ \oplus\ a_{00}\ (1,\ D^+:\ \Omega\ \text{conformal})\ \oplus\ w_i\ (3,\ W:\ \text{shift}),
\]

%dsp-level:0
\[
\qquad 1+1+5+3=10=\tfrac{4\cdot5}{2}
\]

every carrier $S$-even [$\mathcal E_{++},\mathcal E_{+-}$ are $S$-even: CT1.1, CT1.2, C3; CT1.12's data live in $\mathcal E_{++}\oplus\mathcal E_{+-}$]; the $S$-odd blocks ($J^\pm$) carry \textbf{zero} metric coordinates. Consequence: $g_{\mu\nu}$ is orientation-even (used in CT1.19b(2)). $(3)$ \textbf{Read classes.} The corpus's discrete reads fall into exactly three classes, one per discreteness route: \textbf{(R1)} binary registration — $\sigma_e$ [CT0.9], $q_{\mathrm{EM}}\in\{-1,0,+1\}$ [EM.9, D3.4], wall parity $(-1)^{\aleph}$ [$\aleph_v$ from D3.3's $f_v$]; \textbf{(R2)} integer winding on materialized closed chains — $\kappa,\Theta,k$ [CT1.19b]; \textbf{(R3)} binary singlet admissibility — $\mathcal C\in\{0,1\}$ [D2.14; confinement semantics CT3.22a]. Every further dictionary/aspect datum is either continuous configuration data [P4: $g,g',h_0,c_W,\Lambda,\varkappa,\beta,\alpha_\Omega$] or structural non-observable (per-aspect observability verdicts: audit-layer record — not re-derived here); neither carries an independent quantum.

\textbf{Req.} A8; D2.1, D2.14, D3.3; P4; CT0.6, CT0.9, CT0.12, CT1.11a (Part 1); CT1.1, CT1.2, CT1.12; EM.9 (\textbf{E}; forward pointer — C DAG note); exact arithmetic [math; machine: \texttt{checks/\allowbreak{}verify\_\allowbreak{}merge\_\allowbreak{}tierB.py} §C].

\textbf{Pf.}

%dsp-level:0
\[
(1):\ \mathrm{CT0.12}'\text{s rows}:\ 3{+}3{+}5{+}3{+}1{+}1=16=C_{\mathrm{cap}}\ [\mathrm{CT0.6}];
\]

%dsp-level:0
\[
\qquad \text{regrouped by CT1.2's parity table}:\ (5{+}1{+}1){+}3{+}3{+}3=(7,3,3,3)\ [\mathrm{CT1.1}]
\]

%dsp-level:0
\[
(2):\ \mathrm{CT1.12}'\text{s row list}\ (a_{00},\Sigma_{\mathrm{tr}},\hat h_{ij},w_i)\ \text{with counts}\ 1{+}1{+}5{+}3=10=\#\{g_{\mu\nu}=g_{\nu\mu}\};
\]

%dsp-level:0
\[
\quad \text{carriers}:\ a_{00}\sim E_0,\ \Sigma_{\mathrm{tr}}\sim\Sigma,\ \hat h\subset\operatorname{span}\{\hat E^{ij}\},\ w\sim W_i\ \text{all}\ S\text{-even}\ [\mathrm{CT1.2};\ \text{machine}]
\]

%dsp-level:0
\[
J_i\in\mathcal E_{-+},\ K_i\in\mathcal E_{--}\ [\mathrm{CT1.2}]:\ S\text{-odd}\implies\neg\text{rows of an}\ S\text{-even dictionary}
\]

%dsp-level:0
\[
(3):\ \text{enumeration of the declared discrete reads}:\ \mathrm{CT0.9}\ (\sigma\ \text{binary}),
\]

%dsp-level:0
\[
\ \mathrm{EM.9}/\mathrm{D3.4}\ (\text{ternary from two binaries}),
\]

%dsp-level:0
\[
\ \mathrm{D3.3}\ (\aleph_v=|f_v|\ \text{from signs}),
\]

%dsp-level:0
\[
\ \mathrm{CT1.19b}\ (\kappa,\Theta,k\ \text{integers on-shell}),
\]

%dsp-level:0
\[
\ \mathrm{D2.14}\ (\mathcal C\in\{0,1\});
\]

%dsp-level:0
\[
\quad \text{continuous residue}=\mathrm{P4}'\text{s list verbatim}.\ \blacksquare
\]

\textbf{Phys.} The metric speaks only the swap-even dialect: all ten coordinates of spacetime geometry live in the gravity$+$Higgs$+$dark and weak blocks, and the Janus blocks contribute nothing to $g_{\mu\nu}$. Auditing the whole pair channel leaves exactly three routes to discreteness — registration bits, winding integers, sealing booleans — and everything else is continuous configuration data or structure. This is the census of what can be a "quantum" in this corpus.

\end{theorem}

\begin{theorem}[{\texorpdfstring{CT1.13 — Christoffel symbols under conformal rescaling}{CT1.13 - Christoffel symbols under conformal rescaling}}]\label{tag:CT1-13}
\textbf{Thm.} $\ g_{\mu\nu}=\Omega^2\bar g_{\mu\nu},\ \Upsilon=\ln\Omega\implies$

%dsp-level:0
\[
g^{\mu\nu}=\Omega^{-2}\bar g^{\mu\nu};
\]

%dsp-level:0
\[
\qquad \Gamma^\lambda{}_{\mu\nu}=\bar\Gamma^\lambda{}_{\mu\nu}+\delta^\lambda{}_\mu\bar\nabla_\nu\Upsilon+\delta^\lambda{}_\nu\bar\nabla_\mu\Upsilon-\bar g_{\mu\nu}\bar g^{\lambda\rho}\bar\nabla_\rho\Upsilon .
\]

\textbf{Req.} D2.1, D2.2; Levi-Civita formula $\Gamma^\lambda{}_{\mu\nu}=\tfrac12g^{\lambda\rho}(\partial_\mu g_{\rho\nu}+\partial_\nu g_{\rho\mu}-\partial_\rho g_{\mu\nu})$ [math].

\textbf{Pf.}

%dsp-level:0
\[
g^{\mu\alpha}g_{\alpha\nu}=\Omega^{-2}\Omega^2\delta^\mu{}_\nu\implies g^{\mu\nu}=\Omega^{-2}\bar g^{\mu\nu}
\]

%dsp-level:0
\[
\partial_\rho g_{\mu\nu}=2\Omega(\partial_\rho\Omega)\bar g_{\mu\nu}+\Omega^2\partial_\rho\bar g_{\mu\nu}=\Omega^2\bigl(2\Upsilon_\rho\bar g_{\mu\nu}+\partial_\rho\bar g_{\mu\nu}\bigr),
\]

%dsp-level:0
\[
\qquad \Upsilon_\rho=\partial_\rho\Omega/\Omega\ [\mathrm{D2.2}]
\]

%dsp-level:0
\[
\Gamma^\lambda{}_{\mu\nu}
\]

%dsp-level:0
\[
= \tfrac12\Omega^{-2}\bar g^{\lambda\rho}\Omega^2\bigl[(\partial_\mu\bar g_{\rho\nu}+2\bar g_{\rho\nu}\Upsilon_\mu)+(\partial_\nu\bar g_{\rho\mu}+2\bar g_{\rho\mu}\Upsilon_\nu)-(\partial_\rho\bar g_{\mu\nu}+2\bar g_{\mu\nu}\Upsilon_\rho)\bigr]
\]

%dsp-level:0
\[
= \bar\Gamma^\lambda{}_{\mu\nu}+\tfrac12\bar g^{\lambda\rho}\bigl(2\bar g_{\rho\nu}\Upsilon_\mu+2\bar g_{\rho\mu}\Upsilon_\nu-2\bar g_{\mu\nu}\Upsilon_\rho\bigr)
\]

%dsp-level:0
\[
= \bar\Gamma^\lambda{}_{\mu\nu}+\delta^\lambda{}_\nu\Upsilon_\mu+\delta^\lambda{}_\mu\Upsilon_\nu-\bar g_{\mu\nu}\bar g^{\lambda\rho}\Upsilon_\rho
\]

%dsp-level:0
\[
\Upsilon\ \text{scalar}\implies\Upsilon_\mu=\partial_\mu\Upsilon=\bar\nabla_\mu\Upsilon;\qquad\text{RHS symmetric in}\ \mu\leftrightarrow\nu.\ \blacksquare
\]

\textbf{Phys.} A conformal change of the metric shifts the connection by first derivatives of the conformal potential only. This single formula generates the whole curvature-transfer chain CT1.14–CT1.17 that relates emergent geometry to its auxiliary background.

\end{theorem}

\begin{theorem}[{\texorpdfstring{CT1.14 — Riemann tensor under conformal rescaling}{CT1.14 - Riemann tensor under conformal rescaling}}]\label{tag:CT1-14}
\textbf{Thm.} With $\ C^\rho{}_{\mu\nu}:=\Gamma^\rho{}_{\mu\nu}-\bar\Gamma^\rho{}_{\mu\nu}=\delta^\rho{}_\mu\Upsilon_\nu+\delta^\rho{}_\nu\Upsilon_\mu-\bar g_{\mu\nu}\Upsilon^\rho$ (tensor; $\Upsilon^\rho:=\bar g^{\rho\sigma}\Upsilon_\sigma$):

%dsp-level:0
\[
R^\rho{}_{\sigma\mu\nu}=\bar R^\rho{}_{\sigma\mu\nu}+2\bar\nabla_{[\mu}C^\rho{}_{\nu]\sigma}+2C^\rho{}_{[\mu|\lambda|}C^\lambda{}_{\nu]\sigma}.
\]

\textbf{Req.} CT1.13; Riemann definition $R^\rho{}_{\sigma\mu\nu}=\partial_\mu\Gamma^\rho{}_{\nu\sigma}-\partial_\nu\Gamma^\rho{}_{\mu\sigma}+\Gamma^\rho{}_{\mu\lambda}\Gamma^\lambda{}_{\nu\sigma}-\Gamma^\rho{}_{\nu\lambda}\Gamma^\lambda{}_{\mu\sigma}$; $\bar\nabla$ of a $(1,2)$-tensor [math].

\textbf{Pf.}

%dsp-level:0
\[
C=\Gamma-\bar\Gamma:\ \text{difference of connections}\implies\text{tensor};\qquad\text{explicit form}=\mathrm{CT1.13}
\]

%dsp-level:0
\[
\Gamma=\bar\Gamma+C\ \text{substituted into the definition}:\quad
O(C^0):\ \bar R^\rho{}_{\sigma\mu\nu}
\]

%dsp-level:0
\[
O(C^1):\ \ \partial_\mu C^\rho{}_{\nu\sigma}-\partial_\nu C^\rho{}_{\mu\sigma}+\bar\Gamma^\rho{}_{\mu\lambda}C^\lambda{}_{\nu\sigma}+C^\rho{}_{\mu\lambda}\bar\Gamma^\lambda{}_{\nu\sigma}-\bar\Gamma^\rho{}_{\nu\lambda}C^\lambda{}_{\mu\sigma}-C^\rho{}_{\nu\lambda}\bar\Gamma^\lambda{}_{\mu\sigma}
\]

%dsp-level:0
\[
-\bar\Gamma^\lambda{}_{\mu\nu}C^\rho{}_{\lambda\sigma}\ \text{cancels between the antisymmetrized halves}\ (\bar\Gamma^\lambda{}_{\mu\nu}=\bar\Gamma^\lambda{}_{\nu\mu})
\]

%dsp-level:0
\[
\implies O(C^1)
\]

%dsp-level:0
\[
= 2\bar\nabla_{[\mu}C^\rho{}_{\nu]\sigma}
\]

%dsp-level:0
\[
O(C^2):\ \ C^\rho{}_{\mu\lambda}C^\lambda{}_{\nu\sigma}-C^\rho{}_{\nu\lambda}C^\lambda{}_{\mu\sigma}=2C^\rho{}_{[\mu|\lambda|}C^\lambda{}_{\nu]\sigma}
\]

%dsp-level:0
\[
\quad (\lambda\ \text{excluded from antisymmetrization}).\ \blacksquare
\]

\textbf{Phys.} The curvature of the emergent metric differs from the background by a term linear in covariant derivatives of the conformal potential and a term quadratic in it — an exact, closed bookkeeping of how conformal (volume) data manufactures curvature.

\end{theorem}

\begin{theorem}[{\texorpdfstring{CT1.15 — Ricci tensor under conformal rescaling}{CT1.15 - Ricci tensor under conformal rescaling}}]\label{tag:CT1-15}
\textbf{Thm.} In $D$ dimensions:

%dsp-level:0
\[
R_{\mu\nu}
\]

%dsp-level:0
\[
= \bar R_{\mu\nu}-(D-2)\bar\nabla_\mu\bar\nabla_\nu\Upsilon-\bar g_{\mu\nu}\bar\Box\Upsilon+(D-2)\bigl(\Upsilon_\mu\Upsilon_\nu-\bar g_{\mu\nu}|\bar\nabla\Upsilon|^2\bigr).
\]

\textbf{Req.} CT1.14; $\bar\nabla_\mu\Upsilon_\nu=\bar\nabla_\nu\Upsilon_\mu$, $\bar\nabla\bar g=0$ [math].

\textbf{Pf.}

%dsp-level:0
\[
\mathrm{CT1.14}\ \text{contracted}\ \rho=\mu:
\]

%dsp-level:0
\[
\quad R_{\sigma\nu}=\bar R_{\sigma\nu}+\bar\nabla_\mu C^\mu{}_{\nu\sigma}-\bar\nabla_\nu C^\mu{}_{\mu\sigma}+C^\mu{}_{\mu\lambda}C^\lambda{}_{\nu\sigma}-C^\mu{}_{\nu\lambda}C^\lambda{}_{\mu\sigma}
\]

%dsp-level:0
\[
C^\mu{}_{\mu\sigma}=\Upsilon_\sigma+D\,\Upsilon_\sigma-\Upsilon_\sigma=D\,\Upsilon_\sigma
\]

%dsp-level:0
\[
\bar\nabla_\mu C^\mu{}_{\nu\sigma}=\bar\nabla_\nu\Upsilon_\sigma+\bar\nabla_\sigma\Upsilon_\nu-\bar g_{\nu\sigma}\bar\Box\Upsilon;\qquad
\bar\nabla_\nu C^\mu{}_{\mu\sigma}=D\,\bar\nabla_\nu\Upsilon_\sigma
\]

%dsp-level:0
\[
\bar\nabla_\nu\Upsilon_\sigma=\bar\nabla_\sigma\Upsilon_\nu\implies\text{derivative part}=-(D-2)\bar\nabla_\nu\Upsilon_\sigma-\bar g_{\nu\sigma}\bar\Box\Upsilon
\]

%dsp-level:0
\[
C^\mu{}_{\mu\lambda}C^\lambda{}_{\nu\sigma}
\]

%dsp-level:0
\[
= D\,
\]

%dsp-level:0
\[
\Upsilon_\lambda\bigl(\delta^\lambda{}_\nu\Upsilon_\sigma+\delta^\lambda{}_\sigma\Upsilon_\nu-\bar g_{\nu\sigma}\Upsilon^\lambda\bigr)
\]

%dsp-level:0
\[
= D\bigl(2\Upsilon_\nu\Upsilon_\sigma-\bar g_{\nu\sigma}|\bar\nabla\Upsilon|^2\bigr)
\]

%dsp-level:0
\[
C^\mu{}_{\nu\lambda}C^\lambda{}_{\mu\sigma}=(D+2)\,\Upsilon_\nu\Upsilon_\sigma-2\,\bar g_{\nu\sigma}|\bar\nabla\Upsilon|^2
\]

%dsp-level:0
\[
\quad [\text{9 products; coefficients}\ 1{+}1{+}D{+}1{-}1{-}1{+}1{=}D{+}2;\ {-}1{-}1{=}{-}2]
\]

%dsp-level:0
\[
\text{quadratic difference}=(D-2)\bigl(\Upsilon_\nu\Upsilon_\sigma-\bar g_{\nu\sigma}|\bar\nabla\Upsilon|^2\bigr);\qquad
+\bar R_{\sigma\nu},\ \text{rename}\implies\text{Thm}.\ \blacksquare
\]

\textbf{Phys.} The exact response of Ricci curvature to the conformal factor in any dimension. The $(D-2)$ coefficients are what makes $D=4$ (the arena fixed by CT1.12/CT2.6) the non-degenerate case; the identity feeds directly the Einstein-tensor law CT1.17.

\end{theorem}

\begin{theorem}[{\texorpdfstring{CT1.16 — Ricci scalar under conformal rescaling}{CT1.16 - Ricci scalar under conformal rescaling}}]\label{tag:CT1-16}
\textbf{Thm.}

%dsp-level:0
\[
R=\Omega^{-2}\Bigl[\bar R-2(D-1)\bar\Box\Upsilon-(D-1)(D-2)|\bar\nabla\Upsilon|^2\Bigr];\qquad
D=4:\quad R=\Omega^{-2}\Bigl[\bar R-6\bar\Box\Upsilon-6|\bar\nabla\Upsilon|^2\Bigr].
\]

\textbf{Req.} CT1.13 ($g^{\mu\nu}=\Omega^{-2}\bar g^{\mu\nu}$), CT1.15; $\bar g^{\mu\nu}\bar g_{\mu\nu}=D$ [math].

\textbf{Pf.}

%dsp-level:0
\[
R=g^{\mu\nu}R_{\mu\nu}=\Omega^{-2}\bar g^{\mu\nu}R_{\mu\nu}\quad[\mathrm{CT1.13}]
\]

%dsp-level:0
\[
\bar g^{\mu\nu}R_{\mu\nu}=\bar R-(D-2)\bar\Box\Upsilon-D\,\bar\Box\Upsilon+(D-2)|\bar\nabla\Upsilon|^2-D(D-2)|\bar\nabla\Upsilon|^2\quad[\mathrm{CT1.15}]
\]

%dsp-level:0
\[
-(D-2)-D=-2(D-1);\qquad(D-2)-D(D-2)=-(D-1)(D-2)
\]

%dsp-level:0
\[
\times\,\Omega^{-2};\qquad D=4:\ 2(D-1)=6=(D-1)(D-2).\ \blacksquare
\]

\textbf{Phys.} The scalar curvature transforms with the standard conformal weights; in $D=4$ the bracket $\bar R-6\bar\Box\Upsilon-6|\bar\nabla\Upsilon|^2$ is what multiplies $\Omega^{-2}$ — the input that makes the Einstein tensor law CT1.17 collapse to four terms.

\end{theorem}

\begin{theorem}[{\texorpdfstring{CT1.17 — Einstein tensor under conformal rescaling in $D=4$}{CT1.17 - Einstein tensor under conformal rescaling in}}]\label{tag:CT1-17}
\textbf{Thm.}

%dsp-level:0
\[
G_{\mu\nu}=\bar G_{\mu\nu}-2\bar\nabla_\mu\bar\nabla_\nu\Upsilon+2\bar g_{\mu\nu}\bar\Box\Upsilon+2\Upsilon_\mu\Upsilon_\nu+\bar g_{\mu\nu}|\bar\nabla\Upsilon|^2 .
\]

\textbf{Req.} CT1.15, CT1.16 ($D=4$); $G_{\mu\nu}:=R_{\mu\nu}-\tfrac12g_{\mu\nu}R$ [math].

\textbf{Pf.}

%dsp-level:0
\[
R_{\mu\nu}=\bar R_{\mu\nu}-2\bar\nabla_\mu\bar\nabla_\nu\Upsilon-\bar g_{\mu\nu}\bar\Box\Upsilon+2\Upsilon_\mu\Upsilon_\nu-2\bar g_{\mu\nu}|\bar\nabla\Upsilon|^2
\]

%dsp-level:0
\[
\quad [\mathrm{CT1.15},\ D=4]
\]

%dsp-level:0
\[
\tfrac12g_{\mu\nu}R
\]

%dsp-level:0
\[
= \tfrac12\Omega^2\bar g_{\mu\nu}\cdot\Omega^{-2}\bigl[\bar R-6\bar\Box\Upsilon-6|\bar\nabla\Upsilon|^2\bigr]
\]

%dsp-level:0
\[
= \tfrac12\bar g_{\mu\nu}\bigl[\bar R-6\bar\Box\Upsilon-6|\bar\nabla\Upsilon|^2\bigr]
\]

%dsp-level:0
\[
\quad [\mathrm{CT1.16}]
\]

%dsp-level:0
\[
G_{\mu\nu}
\]

%dsp-level:0
\[
= \underbrace{\bigl(\bar R_{\mu\nu}-\tfrac12\bar g_{\mu\nu}\bar R\bigr)}_{\bar G_{\mu\nu}}+(-1+3)\,
\]

%dsp-level:0
\[
\bar g_{\mu\nu}\bar\Box\Upsilon+(-2+3)\,
\]

%dsp-level:0
\[
\bar g_{\mu\nu}|\bar\nabla\Upsilon|^2-2\bar\nabla_\mu\bar\nabla_\nu\Upsilon+2\Upsilon_\mu\Upsilon_\nu.\ \blacksquare
\]

\textbf{Phys.} The exact transformation law of the Einstein tensor in the emergent conformal gauge: everything beyond the background Einstein tensor is built from the single field $\Upsilon=\ln\Omega$. The conformal (volume) degree of freedom of the pair algebra gravitates through a closed-form correction — no approximation involved.

\end{theorem}

\begin{theorem}[{\texorpdfstring{CT1.18 — Linearized spin-2 operator}{CT1.18 - Linearized spin-2 operator}}]\label{tag:CT1-18}
\textbf{Thm.} $\ \bar g_{\mu\nu}=\eta_{\mu\nu}+\varkappa\hat h^{(4)}_{\mu\nu}$, $\ \eta^{\mu\nu}\hat h^{(4)}_{\mu\nu}=0$ (as a field), $\ \Upsilon=N_i=\Sigma_{\mathrm{tr}}=0$. To $O(\varkappa)$:

%dsp-level:0
\[
\bar G_{\mu\nu}
\]

%dsp-level:0
\[
= \frac{\varkappa}{2}\Bigl[-\Box\hat h^{(4)}_{\mu\nu}+\partial_\mu\partial^\lambda\hat h^{(4)}_{\nu\lambda}+\partial_\nu\partial^\lambda\hat h^{(4)}_{\mu\lambda}-\eta_{\mu\nu}\partial^\lambda\partial^\sigma\hat h^{(4)}_{\lambda\sigma}\Bigr]+O(\varkappa^2),
\]

%dsp-level:0
\[
\qquad \Box=\eta^{\alpha\beta}\partial_\alpha\partial_\beta .
\]

With the full linearized Einstein tensor (general $h$)

%dsp-level:0
\[
\bar G^{(1)}[h]_{\mu\nu}:
\]

%dsp-level:0
\[
= \tfrac12\Bigl[-\Box h_{\mu\nu}-\partial_\mu\partial_\nu h^\lambda{}_\lambda+\partial_\mu\partial^\lambda h_{\nu\lambda}+\partial_\nu\partial^\lambda h_{\mu\lambda}+\eta_{\mu\nu}\bigl(\Box h^\lambda{}_\lambda-\partial^\lambda\partial^\sigma h_{\lambda\sigma}\bigr)\Bigr]:
\]

%dsp-level:0
\[
h^\lambda{}_\lambda=0\implies\bar G^{(1)}[h]=\text{bracket above}/2;\qquad
\bar G^{(1)}\bigl[h+\partial_\mu\xi_\nu+\partial_\nu\xi_\mu\bigr]=\bar G^{(1)}[h]\ \ \forall\xi;
\]

%dsp-level:0
\[
\partial^\lambda\hat h^{(4)}_{\lambda\nu}=0\ \wedge\ \hat h^{(4)\lambda}{}_\lambda=0\implies\bar G^{(1)}[\hat h^{(4)}]=-\tfrac12\Box\hat h^{(4)}_{\mu\nu}
\]

%dsp-level:0
\[
\quad (\text{massless wave operator}).
\]

\textbf{Req.} CT1.10 (spin-2 mode), CT1.12; Levi-Civita/Ricci expansion [math]; $\partial_\mu\partial_\nu=\partial_\nu\partial_\mu$.

\textbf{Pf.}

%dsp-level:0
\[
\bar g^{\lambda\rho}
\]

%dsp-level:0
\[
= \eta^{\lambda\rho}-\varkappa\hat h^{(4)\lambda\rho}+O(\varkappa^2)
\]

%dsp-level:0
\[
\implies \bar\Gamma^\lambda{}_{\mu\nu}
\]

%dsp-level:0
\[
= \tfrac\varkappa2\bigl(\partial_\mu\hat h^{(4)\lambda}{}_\nu+\partial_\nu\hat h^{(4)\lambda}{}_\mu-\partial^\lambda\hat h^{(4)}_{\mu\nu}\bigr)+O(\varkappa^2)
\]

%dsp-level:0
\[
\quad (\eta\text{-raising})
\]

%dsp-level:0
\[
\bar R_{\mu\nu}
\]

%dsp-level:0
\[
= \partial_\lambda\bar\Gamma^\lambda{}_{\mu\nu}-\partial_\nu\bar\Gamma^\lambda{}_{\lambda\mu}+O(\varkappa^2)
\]

%dsp-level:0
\[
= \tfrac\varkappa2\bigl(\partial_\mu\partial^\lambda\hat h^{(4)}_{\lambda\nu}+\partial_\nu\partial^\lambda\hat h^{(4)}_{\lambda\mu}-\Box\hat h^{(4)}_{\mu\nu}-\partial_\mu\partial_\nu\hat h^{(4)\lambda}{}_\lambda\bigr)
\]

%dsp-level:0
\[
\hat h^{(4)\lambda}{}_\lambda\equiv0\ \text{(field)}\implies\partial_\mu\partial_\nu\hat h^{(4)\lambda}{}_\lambda=0
\]

%dsp-level:0
\[
\bar R=\eta^{\mu\nu}\bar R_{\mu\nu}=\varkappa\bigl(\partial^\mu\partial^\nu\hat h^{(4)}_{\mu\nu}-\Box\hat h^{(4)\lambda}{}_\lambda\bigr)=\varkappa\,\partial^\lambda\partial^\sigma\hat h^{(4)}_{\lambda\sigma}
\]

%dsp-level:0
\[
\bar G_{\mu\nu}=\bar R_{\mu\nu}-\tfrac12\eta_{\mu\nu}\bar R+O(\varkappa^2)\implies\text{Thm formula}
\]

%dsp-level:0
\[
\bar G^{(1)}[h+\partial\xi+\partial\xi]-\bar G^{(1)}[h]:\ \xi\text{-terms}=\text{combinations of}\ \partial_{(\alpha}\partial_\beta\partial_{\gamma)}\xi_\delta=0\implies\text{invariance}
\]

%dsp-level:0
\[
\partial^\lambda\hat h^{(4)}_{\lambda\nu}
\]

%dsp-level:0
\[
= 0\implies\partial_\mu\partial^\lambda\hat h^{(4)}_{\nu\lambda}
\]

%dsp-level:0
\[
= \partial_\nu\partial^\lambda\hat h^{(4)}_{\mu\lambda}
\]

%dsp-level:0
\[
= \partial^\lambda\partial^\sigma\hat h^{(4)}_{\lambda\sigma}
\]

%dsp-level:0
\[
= 0\implies\bar G^{(1)}[\hat h^{(4)}]
\]

%dsp-level:0
\[
= -\tfrac12\Box\hat h^{(4)}_{\mu\nu}.\ \blacksquare
\]

Machine verification (incl. $\bar G^{(0)}=0$, $\bar G^{(1)}=L[\hat h^{(4)}]$ from the full nonlinear pipeline, gauge invariance): \texttt{checks/\allowbreak{}verify\_\allowbreak{}C.py}.

\textbf{Rmk (correction).} The traceless-projected transformation $\hat h^{(4)}\mapsto\hat h^{(4)}+\partial_\mu\xi_\nu+\partial_\nu\xi_\mu-\tfrac12\eta_{\mu\nu}\partial\!\cdot\!\xi$ preserves $\eta^{\mu\nu}\hat h^{(4)}_{\mu\nu}=0$ but shifts $\bar G^{(1)}$ by $\tfrac12\bigl(\partial_\mu\partial_\nu-\eta_{\mu\nu}\Box\bigr)\partial\!\cdot\!\xi\ne0$ (machine-verified): gauge invariance holds for the standard linearized diffeomorphism $\partial_\mu\xi_\nu+\partial_\nu\xi_\mu$, which does not preserve tracelessness and need not, since $\bar G^{(1)}$ is defined for general $h$. The prior codex claimed invariance of the projected form — false as stated.

\textbf{Phys.} With only the five-component traceless shear excited, the emergent Einstein tensor reduces at leading order to the massless spin-2 wave operator in transverse gauge. The graviton of the corpus is the $\mathcal E_{++}$ quintet of CT1.10 propagating on the auxiliary background; its gauge freedom is linearized diffeomorphism invariance.

\end{theorem}

\phantomsection
\section[§C.II — Masks, Manifold, and Closure]{\texorpdfstring{§C.II — Masks, Manifold, and Closure}{SC.II - Masks, Manifold, and Closure}}

\begin{theorem}[{\texorpdfstring{CT2.1 — Unit and dimension conventions}{CT2.1 - Unit and dimension conventions}}]\label{tag:CT2-1}
\textbf{Thm.} With $c=\hbar=k_B=1$, $[\mathfrak h]=\mathfrak q^{-1/2}$:

%dsp-level:0
\[
[\text{length}]=\mathfrak q^{-1/2};
\]

%dsp-level:0
\[
\quad [\text{area}]=\mathfrak q^{-1};
\]

%dsp-level:0
\[
\quad [\text{energy}]=\mathfrak q^{+1/2};
\]

%dsp-level:0
\[
\quad [T_{\mu\nu}]=\mathfrak q^{+2};
\]

%dsp-level:0
\[
\quad [G_{\mu\nu}]=[\Lambda]=\mathfrak q^{+1};
\]

%dsp-level:0
\[
\quad [\kappa^2]=[G_N]=\mathfrak q^{-1};
\]

%dsp-level:0
\[
\quad [\kappa]=\mathfrak q^{-1/2};
\]

%dsp-level:0
\[
\quad [S_{\mathrm{geom}}]=\mathfrak q^{0},
\]

and every equation of Parts III–IV balances.

\textbf{Req.} D0.6, D2.9; dimensional bookkeeping [math].

\textbf{Rmk (forward pointers, non-load-bearing).} The normalization of $\kappa$ is cross-checked against CT1.23 and $G_N\propto\kappa^2$ against CT2.9 — consistency pointers only.

\textbf{Pf.}

%dsp-level:0
\[
[\mathfrak h]=\mathfrak q^{-1/2}\ [\mathrm{D0.6},\mathrm{D2.9}]\implies[\text{length}]=\mathfrak q^{-1/2},\ [\text{area}]=\mathfrak q^{-1};
\]

%dsp-level:0
\[
\qquad \hbar=c=1\implies[\text{energy}]=[\text{length}]^{-1}=\mathfrak q^{+1/2}
\]

%dsp-level:0
\[
[\text{curvature}]=[\text{area}]^{-1}=\mathfrak q^{+1}=[G_{\mu\nu}];\qquad\Lambda g_{\mu\nu}\sim G_{\mu\nu}\implies[\Lambda]=\mathfrak q^{+1}
\]

%dsp-level:0
\[
[T_{\mu\nu}]=\frac{[\text{energy}]}{[\text{volume}]}=\frac{\mathfrak q^{1/2}}{\mathfrak q^{-3/2}}=\mathfrak q^{+2}
\]

%dsp-level:0
\[
[\kappa^2T_{\mu\nu}]=[G_{\mu\nu}]\ [\mathrm{CT1.23}]\implies[\kappa^2]=\mathfrak q^{1-2}=\mathfrak q^{-1},\ [\kappa]=\mathfrak q^{-1/2};
\]

%dsp-level:0
\[
\qquad G_N\propto\kappa^2\ [\mathrm{CT2.9}]\implies[G_N]=\mathfrak q^{-1}
\]

%dsp-level:0
\[
[d^4x]=\mathfrak q^{-2},\ [\sqrt{-g}]=\mathfrak q^0,\ [R]=\mathfrak q^{+1},\ [\kappa^{-2}]=\mathfrak q^{+1}\implies[S_{\mathrm{geom}}]=\mathfrak q^0\ (\hbar=1).\ \blacksquare
\]

\textbf{Phys.} One information unit $\mathfrak q$ grades the entire continuum layer: lengths $\mathfrak q^{-1/2}$, curvature and $\Lambda$ at $\mathfrak q^{+1}$, stress at $\mathfrak q^{+2}$, Newton’s constant an area $\mathfrak q^{-1}$. The dimensionless action confirms no dimensionful quantity beyond $\mathfrak h$ has been smuggled in.

\end{theorem}

\begin{theorem}[{\texorpdfstring{CT2.2 — The rung mask is a projection}{CT2.2 - The rung mask is a projection}}]\label{tag:CT2-2}
\textbf{Thm.}

%dsp-level:0
\[
\bigl(\mathcal P_M^{(n)}\bigr)^2=\mathcal P_M^{(n)}=\bigl(\mathcal P_M^{(n)}\bigr)^\dagger;
\]

%dsp-level:0
\[
\qquad \operatorname{Im}\mathcal P_M^{(n)}=\mathcal K_n;
\]

%dsp-level:0
\[
\qquad \mathfrak T\,\mathcal P_M^{(n)}=\dim\mathcal K_n\in\{6,16,17,25,\dots\}.
\]

\textbf{Req.} D2.3; CT0.13, CT0.14 (Part 1); orthogonal-projection facts [linear algebra].

\textbf{Pf.}

%dsp-level:0
\[
\mathrm{D2.3}:\ \mathcal P_M^{(n)}=\text{orthogonal projection}\ \mathcal K\to\mathcal K_n=\bigoplus_{A\in\mathcal A(n)}\mathcal K^{(A)}
\]

%dsp-level:0
\[
\text{orthogonal projection}\implies P^2=P=P^\dagger,\quad\mathfrak T P=\dim\operatorname{Im}P\quad[\text{linear algebra}]
\]

%dsp-level:0
\[
\operatorname{Im}=\mathcal K_n\ [\text{construction}];\qquad\dim\mathcal K_n\ [\mathrm{CT0.14}];\qquad\mathcal A(n)\ \text{derived}\ [\mathrm{CT0.13}].\ \blacksquare
\]

\textbf{Rmk.} Theorem, not stipulation: CT0.13 derives which blocks are active; D2.3 converts the list into a projector.

\textbf{Phys.} Rung selection is a genuine orthogonal projection in force space — idempotent, self-adjoint, with trace equal to the block dimension $6/16/17/25$. Filtering by rung twice is the same as once: the universe’s capability space grows by strict projections, never by deformations.

\end{theorem}

\begin{theorem}[{\texorpdfstring{CT2.3 — Masked force operator}{CT2.3 - Masked force operator}}]\label{tag:CT2-3}
\textbf{Thm.} $\ \mathbf F_M^{(n)}:=\mathcal P_M^{(n)}\mathbf F\,\mathcal P_M^{(n)}$ satisfies

%dsp-level:0
\[
\mathbf F_M^{(n)}\big|_{\mathcal K_n^\perp}=0;\qquad
\mathbf F_M^{(n)}(\mathcal K)\subseteq\mathcal K_n;\qquad
\mathbf F^\dagger=\mathbf F\implies\bigl(\mathbf F_M^{(n)}\bigr)^\dagger=\mathbf F_M^{(n)}.
\]

\textbf{Req.} A8, C7; D2.3; CT2.2.

\textbf{Pf.}

%dsp-level:0
\[
x\in\mathcal K_n^\perp\implies\mathcal P_M^{(n)}x=0\implies\mathbf F_M^{(n)}x=0
\]

%dsp-level:0
\[
\forall x:\ \ \mathbf F_M^{(n)}x=\mathcal P_M^{(n)}\bigl(\mathbf F\,\mathcal P_M^{(n)}x\bigr)\in\operatorname{Im}\mathcal P_M^{(n)}=\mathcal K_n\quad[\mathrm{CT2.2}]
\]

%dsp-level:0
\[
\bigl(\mathcal P\mathbf F\mathcal P\bigr)^\dagger=\mathcal P^\dagger\mathbf F^\dagger\mathcal P^\dagger=\mathcal P\mathbf F\mathcal P\quad[\mathrm{CT2.2};\ \mathbf F^\dagger=\mathbf F].\ \blacksquare
\]

\textbf{Phys.} The masked force is exactly the restriction of the full force operator to the rung’s active subspace: it annihilates everything the rung cannot access, stays inside what it can, and inherits Hermiticity — hence a real spectrum — from $\mathbf F$. A8’s $\mathbf F_M=P_M\mathbf F P_M$ is the rung-unresolved form (C7).

\end{theorem}

\begin{theorem}[{\texorpdfstring{CT2.4 — Force restriction by rung}{CT2.4 - Force restriction by rung}}]\label{tag:CT2-4}
\textbf{Thm.}

%dsp-level:0
\[
\mathbf F_M^{(n)}=\bigoplus_{A\in\mathcal A(n)}F_A\big|_{\mathcal K^{(A)}},
\]

%dsp-level:0
\[
\qquad \mathcal A(n)\ \text{as in CT0.13};
\]

%dsp-level:0
\[
\qquad \neg\exists\ \text{independent aspect}\ Z\ \text{at any rung}\ (W_3\ \text{seed}\in F_W)\ [\mathrm{C8}].
\]

\textbf{Req.} A8, C8; D2.3; CT0.13 (Part 1), CT2.2, CT2.3.

\textbf{Pf.}

%dsp-level:0
\[
\mathrm{A8}:\ \mathbf F=\bigoplus_AF_A,\quad\operatorname{supp}F_A\subseteq\mathcal K^{(A)};\qquad
\mathrm{D2.3}:\ \mathcal P_M^{(n)}\to\bigoplus_{A\in\mathcal A(n)}\mathcal K^{(A)}
\]

%dsp-level:0
\[
\text{block-diagonal conjugated by block-subset projection}\implies\text{retains exactly those blocks}
\]

%dsp-level:0
\[
\mathcal A(n)\ \text{list}=[\mathrm{CT0.13}];\qquad F_W=F_{W^+}\oplus F_{W^-}\oplus F_Z,\ F_Z\equiv F_{W^3}\ \text{single block}\ [\mathrm{C8}].\ \blacksquare
\]

\textbf{Phys.} Masking never mixes forces: what survives at rung $n$ is exactly the direct sum of the active aspects. This licenses speaking of gravity from rung 1, electromagnetism from rung 2 and colour from rung 3 as literal block content of one operator.

\end{theorem}

\begin{theorem}[{\texorpdfstring{CT2.5 — Masked resolvent as spectrum generator}{CT2.5 - Masked resolvent as spectrum generator}}]\label{tag:CT2-5}
\textbf{Thm.} $\ \mathbf F_M^{(n)}=\sum_a\lambda_aP_a$ (spectral decomposition; $P_aP_b=\delta_{ab}P_a$, $P_a^\dagger=P_a$). Then

%dsp-level:0
\[
Z(\mathcal R,n)=\Bigl\langle\mathrm{out}\Bigl|(1-\mathbf F_M^{(n)})^{-1}\Bigr|\mathrm{in}\Bigr\rangle=\sum_a\frac{\langle\mathrm{out}|P_a|\mathrm{in}\rangle}{1-\lambda_a};
\]

%dsp-level:0
\[
\text{poles}\iff1\in\operatorname{spec}\mathbf F_M^{(n)};
\]

%dsp-level:0
\[
\qquad \operatorname*{Res}_{\lambda_a\to1}=-\langle\mathrm{out}|P_a|\mathrm{in}\rangle=\text{transition amplitude};
\]

%dsp-level:0
\[
\qquad \mathrm{D2.8}:\ 1\mapsto1+i\varepsilon\ \text{selects the branch}.
\]

\textbf{Req.} D2.4, D2.5, D2.8; CT0.14 (\textbf{B}: $\dim\mathcal K_n<\infty$), CT2.3 (Hermitian); spectral theorem [math].

\textbf{Pf.}

%dsp-level:0
\[
(\mathbf F_M^{(n)})^\dagger=\mathbf F_M^{(n)}\ [\mathrm{CT2.3}]\ \wedge\ \dim\mathcal K_n<\infty\ [\mathrm{CT0.14}]\implies\mathbf F_M^{(n)}=\sum_a\lambda_aP_a\quad[\text{spectral theorem}]
\]

%dsp-level:0
\[
(1-\mathbf F_M^{(n)})^{-1}=\sum_a\frac{P_a}{1-\lambda_a},\quad\text{defined}\iff\forall a:\ \lambda_a\ne1;\quad\text{singular}\iff\exists a:\ \lambda_a=1
\]

%dsp-level:0
\[
Z(\mathcal R,n)=\sum_a\frac{\langle\mathrm{out}|P_a|\mathrm{in}\rangle}{1-\lambda_a}:\ \ \text{simple pole at}\ \lambda_a=1,\ \operatorname{Res}=-\langle\mathrm{out}|P_a|\mathrm{in}\rangle\ [\mathrm{D2.5}:\ \text{amplitude}]
\]

%dsp-level:0
\[
\mathrm{D2.8}:\ 1\mapsto1+i\varepsilon\implies\text{poles off the real axis, fixed direction};\quad\varepsilon\to0^+\ \text{after contour integration}.\ \blacksquare
\]

\textbf{Phys.} The rung-$n$ particle spectrum is literally the condition “the masked force has eigenvalue 1”: poles of the resolvent are masses (D2.5), residues are amplitudes. The $i\varepsilon$ sheet makes the sums convergent and selects the branch — the entire spectral content of the corpus sits in this one matrix element.

\end{theorem}

\begin{theorem}[{\texorpdfstring{CT2.6 — Continuum manifold closure}{CT2.6 - Continuum manifold closure}}]\label{tag:CT2-6}
\textbf{Thm.} Conditions: (i) valid universe [D2.11] — a definitional precondition (D2.11’s own definition in application, “there is something”), $\neg$ an extra-axiomatic hypothesis [A §0.2a]; (ii) coarse-graining cells contain unboundedly many real edges as $\mathfrak h\to0$ — derived (Part 4, Tier BM.24: the ledger's own unbounded ladder witnesses it [BM.151(1), forward pointer]); (iii) sector coefficient fields [D2.1] of bounded variation across neighbouring cells — derived (Part 4, Tier BM.24: the spectral bound $|r_a|\le\sqrt3$ with convex averaging [BM.151(2), forward pointer]). Then

%dsp-level:0
\[
M:=\Gamma_{\mathcal R}/\!\!\sim\ \ \text{is a smooth 4-manifold};\qquad
\chi([v])=(x^0,x^1,x^2,x^3)=(\tau_n,h_1,h_2,h_3)\ \ \text{charts};
\]

%dsp-level:0
\[
\text{transition maps}=\text{smooth diffeomorphisms};
\]

%dsp-level:0
\[
\qquad g_{\mu\nu}\ [\mathrm{CT1.12}]=\text{smooth Lorentzian metric on}\ M;
\]

%dsp-level:0
\[
\qquad \text{nondegeneracy}\iff\varkappa\|q\|_{\mathrm{op}}<1 .
\]

\textbf{Req.} A7; D1.3, D1.8, D1.9, D2.11; CT0.3 (Part 1), CT1.12; conditions (ii)–(iii) [derived — cf. Part 4 Tier BM.24, BM.151, forward pointer]; (i) $=$ D2.11 in application.

\textbf{Pf.}

%dsp-level:0
\[
(\mathrm{i}):\ \#J^+
\]

%dsp-level:0
\[
= \#J^-\ [\mathrm{D2.11},\mathrm{D1.3}]\implies\textstyle\sum\Delta n_{\mathcal R}
\]

%dsp-level:0
\[
= 0\ [\mathrm{CT0.3}]\implies\neg\exists\ \text{dangling ledger boundary beyond the sealed Janus boundary}
\]

%dsp-level:0
\[
\sim\ :=\ \text{equal coarse-grained ledger entry under mask-preserving relabeling};
\]

%dsp-level:0
\[
\qquad (\mathrm{ii})\implies\text{cells contain unboundedly many edges}\implies M\ \text{locally a continuum as}\ \mathfrak h\to0
\]

%dsp-level:0
\[
\text{charts}:\ x^0:=\tau_n\ [\mathrm{D1.8}];
\]

%dsp-level:0
\[
\quad x^i:=\text{hop counters projected on the 3 shift directions}\ (W_i\ \text{sector})\ [\mathrm{CT1.12}];\quad\chi:U\subseteq M\to\mathbb R^4;
\]

%dsp-level:0
\[
\quad 4=1+3\ [\mathrm{CT1.12}]
\]

%dsp-level:0
\[
\text{smoothness}:\ \text{coefficients}
\]

%dsp-level:0
\[
= \text{local ledger averages};
\]

%dsp-level:0
\[
\ \ \mathrm{D1.9}:\ \text{updates local, finite};
\]

%dsp-level:0
\[
\ \ (\mathrm{iii}):\ \text{jumps uniformly bounded}\implies\text{smooth interpolation, deviation}\to0\ \text{with}\ \mathfrak h
\]

%dsp-level:0
\[
\text{transitions}
\]

%dsp-level:0
\[
= \text{resynchronizations between valid observers}:\ \text{mask-}\wedge\text{Janus-preserving}\implies\text{finite compositions of translations/permutations/sign reversals}
\]

%dsp-level:0
\[
\implies\text{smooth, constant Jacobian},\ |\det J|=1;\quad\circ\ \text{interpolation}\implies\text{diffeomorphisms}
\]

%dsp-level:0
\[
\text{nondegeneracy}:\ \mathrm{CT1.12}\ (\varkappa\|q\|_{\mathrm{op}}<1)\iff J\ \text{invertible};
\]

%dsp-level:0
\[
\qquad \text{atlas}\implies M\ \text{smooth 4-manifold};\ g_{\mu\nu}\ \text{smooth Lorentzian}.\ \blacksquare
\]

\textbf{Phys.} Under three explicit coarse-graining hypotheses — and nothing else — the real-edge network of the ledger becomes a smooth four-dimensional spacetime carrying the emergent Lorentzian metric. This is the bridge theorem of the corpus: it converts combinatorics (edges, ticks, counters) into the differential geometry that Parts III–IV vary.

\textbf{Rmk (derivation locus: Part 4 Tier BM.24).} Conditions (ii)–(iii) carry no extra-axiomatic status [A §0.2a]: BM.151 derives both from the ledger's own structure: (ii) from the configuration-space ladder (A7; FT.32; FT.45–FT.46) with the ladder row $\#\{E_{\mathcal R}\ \text{per cell}\}\ge2sE\ln2/\mathfrak h\to\infty$, and (iii) from the pair-state spectral bound (A1; L1; L3(1)) with the convex averaging of the coefficient fields. This theorem's statement, proof and conclusions are unchanged; the point-of-use declarations across the corpus stand as scope-inheritance markers. The extra-axiomatic list of A §0.2a is empty.

\end{theorem}

\begin{theorem}[{\texorpdfstring{CT2.7 — Exact geometric action closure}{CT2.7 - Exact geometric action closure}}]\label{tag:CT2-7}
\textbf{Thm.} Under the hypotheses of CT2.6:

%dsp-level:0
\[
S_{\mathrm{geom}}[g]=\frac{1}{2\kappa^2}\int d^4x\,\sqrt{-g}\,\bigl(R-2\Lambda\bigr);\qquad
\neg\exists\ \text{additional local higher-curvature terms}.
\]

\textbf{Req.} A8; D1.9, D4.3; CT2.6; Lovelock classification $D=4$; Gauss–Bonnet $=$ total derivative in $D=4$ [math].

\textbf{Rmk (forward pointer, non-load-bearing).} The $\kappa$-normalization cross-check against CT1.23 is a consistency pointer, not a derivation input.

\textbf{Pf.}

%dsp-level:0
\[
\mathrm{D1.9}:\ \text{tick stencil}=1\ \text{time step}+1\ \text{edge}\implies\text{continuum E--L equations}\le2\ \text{nd order}
\]

%dsp-level:0
\[
\text{arena}=\text{smooth 4-dim Lorentzian manifold}\ [\mathrm{CT2.6}]
\]

%dsp-level:0
\[
\text{weights depend on}\ (\Gamma,\text{registrations})\ \text{only},
\]

%dsp-level:0
\[
\ \neg\text{vertex names}\implies\text{relabeling-invariant}\xrightarrow[\mathrm{CT2.6}]{}\text{diffeomorphism-invariant local functional of}\ g
\]

%dsp-level:0
\[
\text{Lovelock}\ (D=4):\ \bigl\{\sqrt{-g}R,\ \sqrt{-g},\ \sqrt{-g}\,\mathcal G\bigr\},\quad\mathcal G=R^2-4R_{\mu\nu}R^{\mu\nu}+R_{\mu\nu\rho\sigma}R^{\mu\nu\rho\sigma}\quad[\text{math}]
\]

%dsp-level:0
\[
\sqrt{-g}\,\mathcal G=\partial(\cdots)\ \text{in}\ D=4\implies\delta S\equiv0\implies\text{dropped}\quad[\text{math}]
\]

%dsp-level:0
\[
F_{D^+}=\Lambda\mathbb I\ [\mathrm{A8}],\ \neg\text{configuration-dependent}\ [\mathrm{D4.3}]\implies\text{zero-derivative density}\ \sqrt{-g}\,\Lambda\ \text{only}
\]

%dsp-level:0
\[
\text{collect at}\ \tfrac{1}{2\kappa^2}\ \text{normalization}\ [\mathrm{CT1.23}]\implies\text{Thm}.\ \blacksquare
\]

\textbf{Phys.} The geometric dynamics is closed exactly: tick-locality plus diffeomorphism invariance plus Lovelock leave only Einstein–Hilbert with a cosmological term, and that term is the uniform dark aspect itself. No higher-curvature corrections exist — the higher-derivative question is answered, not deferred.

\end{theorem}

\begin{theorem}[{\texorpdfstring{CT2.8 — Reflection positivity and Hilbert-space closure}{CT2.8 - Reflection positivity and Hilbert-space closure}}]\label{tag:CT2-8}
\textbf{Thm.} With $\ \langle A|B\rangle:=Z\bigl[\Theta_{\mathfrak L}A\cup B\bigr]\ $ on $\mathfrak C(\mathcal R)$ [D0.10]:

%dsp-level:0
\[
(1)\ \ w(\Theta_{\mathfrak L}C)=w(C)^\ast;\qquad
(2)\ \ w(C)=w(C_{\mathrm{left}})\cdot w(C_{\mathrm{right}})\ \ \text{across any interior ledger-time cut};
\]

%dsp-level:0
\[
(3)\ \ \langle A|A\rangle=\sum_C|w(C)|^2\ \ge\ 0;
\]

%dsp-level:0
\[
\qquad (4)\ \ \langle A|A\rangle=0\iff A\ \text{mask-excluded}\ \ (\text{null space}=\text{mask-excluded configurations});
\]

%dsp-level:0
\[
(5)\ \ \text{quotient}\wedge\text{completion}=\text{Hilbert space};\qquad
(6)\ \ Z\big|_{\text{reflected pair}}=\|\cdot\|^2,\ \ \text{convergent on the}\ \mathrm{D2.8}\ \text{sheet}.
\]

\textbf{Req.} A8; D0.7, D0.10; D1.2, D1.3, D1.9; D2.6, D2.7, D2.8; D3.3, D3.5, D3.6; CT2.2; Hilbert-space completion [math; machine: \texttt{checks/\allowbreak{}verify\_\allowbreak{}F33\_\allowbreak{}dynamic\_\allowbreak{}phase.py} §D].

\textbf{Pf.}

%dsp-level:0
\[
(1):\ \mathcal J(\Theta_{\mathfrak L}C)=\mathcal J(C)^\ast\ [\mathrm{D3.6}];
\]

%dsp-level:0
\[
\qquad G(\Theta_{\mathfrak L}e)=G(e)^\ast\ [\mathrm{D2.7},\ \mathrm{D3.5}:\ \alpha_e\widehat{ds}^{\,2}(e)\mapsto-\alpha_e\widehat{ds}^{\,2}(e),\ \sigma_e\ \text{fixed}]
\]

%dsp-level:0
\[
\omega:\ \text{type}(1,1)\mapsto\text{signs exchanged}
\]

%dsp-level:0
\[
\implies f_v
\]

%dsp-level:0
\[
\mapsto -f_v\implies\aleph_v
\]

%dsp-level:0
\[
= |f_v|\ \text{invariant}
\]

%dsp-level:0
\[
\implies (-1)^\aleph\in\mathbb R\ \text{invariant}\ [\mathrm{D2.6},\mathrm{D3.3}]
\]

%dsp-level:0
\[
\implies w(\Theta_{\mathfrak L}C)=w(C)^\ast\quad[\mathrm{A8}]
\]

%dsp-level:0
\[
(2):\ \mathrm{A8}:\ w(C)=\mathcal J(C)\prod_{e\in E(C)}G(e)\cdot(-1)^{\aleph(\sigma,C)};
\]

%dsp-level:0
\[
\qquad \text{interior cut at a ledger time}:\ E(C)
\]

%dsp-level:0
\[
= E_{\mathrm L}\sqcup E_{\mathrm R}\ [\mathrm{D0.7}:\ \text{each edge is registered between consecutive ticks, wholly on one side}]
\]

%dsp-level:0
\[
\prod_{e\in E(C)}G(e)=\prod_{E_{\mathrm L}}G\cdot\prod_{E_{\mathrm R}}G;
\]

%dsp-level:0
\[
\qquad (-1)^{\aleph(\sigma,C)}=\prod_v(-1)^{\aleph_v}:\ \aleph_v\ v\text{-local}\ [\mathrm{D3.3};\ \mathrm{D1.9}:\ \operatorname{dom}(\mathrm{tick}(v))\subseteq\{v\}\cup\{e\ni v\}]
\]

%dsp-level:0
\[
\text{interface}:\ \text{the seal data at the cut are boundary data of the halves}\ [\mathrm{D1.2};\ \mathrm{D1.3}:\ \#J^+=\#J^-\ \text{pair at boundaries}];
\]

%dsp-level:0
\[
\qquad \text{side convention: each vertex factor belongs to the side holding its tick}
\]

%dsp-level:0
\[
\implies w(C)=w(C_{\mathrm{left}})\cdot w(C_{\mathrm{right}})\quad[\text{product regrouping}]
\]

%dsp-level:0
\[
(3):\ \ A=B:\ \ \langle A|A\rangle=\sum_Cw(\Theta_{\mathfrak L}C)w(C)=\sum_C|w(C)|^2\ge0\quad[(1),(2)]
\]

%dsp-level:0
\[
(4):\ \ =0\implies\forall C:\ w(C)=0;
\]

%dsp-level:0
\[
\qquad \mathrm{A8}\ \text{weight}=\text{product of nonzero factors except where excluded}\ [|G(e)|=1\ \mathrm{(D3.5)};\ \mathcal J\ne0];
\]

%dsp-level:0
\[
\qquad \text{excluded}\iff\notin\operatorname{Im}\mathcal P_M^{(n)}\ [\mathrm{CT2.2}]
\]

%dsp-level:0
\[
(5):\ \text{PSD form}\implies\text{positive definite on quotient}\implies\text{completion}=\text{Hilbert space}
\]

%dsp-level:0
\[
\quad [\text{math: completion theorem; no network-specific input}]
\]

%dsp-level:0
\[
(6):\ \ \mathrm{D2.8}\ \text{sheet}\implies\sum|w|^2\ \text{converges};\ \neg\exists\ \text{indefinite oscillatory spoilage}. \blacksquare
\]

\textbf{Phys.} The ledger’s own reflection symmetry makes the configuration sum positive: the reflected pairing is a sum of $|w|^2$, its null space is exactly what the mask forbids, and the quotient completes to a genuine Hilbert space. The split along an interior ledger-time cut is a consequence of A8’s per-edge product structure and tick locality — the seals sit at boundaries [D1.2, D1.3], and no factor straddles an interior cut. This is the positivity property that licenses reading ledger weights as quantum theory.

\end{theorem}

\begin{theorem}[{\texorpdfstring{CT2.10 — Graded matrices and the supertrace}{CT2.10 - Graded matrices and the supertrace}}]\label{tag:CT2-10}
\textbf{Thm.} For homogeneous $X,Y$ [D2.10]:

%dsp-level:0
\[
\operatorname{Str}(XY)=(-1)^{|X||Y|}\operatorname{Str}(YX);\qquad
\operatorname{Str}[X,Y\}=0 .
\]

\textbf{Req.} D2.10; cyclicity of $\mathfrak T$ [linear algebra].

\textbf{Pf.}

%dsp-level:0
\[
|X|=|Y|=0:\ \ X=\operatorname{diag}(X_B,X_F),\ Y=\operatorname{diag}(Y_B,Y_F):
\]

%dsp-level:0
\[
\quad \operatorname{Str}(XY)=\mathfrak T(X_BY_B)-\mathfrak T(X_FY_F)=\mathfrak T(Y_BX_B)-\mathfrak T(Y_FX_F)=\operatorname{Str}(YX);\ (-1)^0=+1
\]

%dsp-level:0
\[
|X|\ne|Y|:\ \ XY,\ YX\ \text{odd}\implies\text{diagonal blocks}=0\implies\operatorname{Str}(XY)=\operatorname{Str}(YX)=0
\]

%dsp-level:0
\[
|X|=|Y|=1:\ \ X=\begin{pmatrix}0&X_1\\ X_2&0\end{pmatrix},\ Y=\begin{pmatrix}0&Y_1\\ Y_2&0\end{pmatrix}:\quad
XY=\operatorname{diag}(X_1Y_2,\ X_2Y_1),\ YX=\operatorname{diag}(Y_1X_2,\ Y_2X_1)
\]

%dsp-level:0
\[
\operatorname{Str}(XY)=\mathfrak T(X_1Y_2)-\mathfrak T(X_2Y_1)=\mathfrak T(Y_2X_1)-\mathfrak T(Y_1X_2)=-\operatorname{Str}(YX);\ (-1)^{1}=-1
\]

%dsp-level:0
\[
\implies\operatorname{Str}[X,Y\}=\operatorname{Str}(XY)-(-1)^{|X||Y|}\operatorname{Str}(YX)=0\ \ \forall\ \text{parities}.\ \blacksquare
\]

\textbf{Phys.} The supertrace is graded-cyclic: it vanishes on graded commutators and therefore descends to an invariant form on the graded (bosonic/fermionic) algebra — the tool that lets bound (fermionic) and transit (bosonic) content coexist in one trace formula.

\end{theorem}

\phantomsection
\section[§C.III — Continuum Limit, Effective Action, and Field…]{\texorpdfstring{§C.III — Continuum Limit, Effective Action, and Field Equations}{SC.III - Continuum Limit, Effective Action, and Field Equations}}

\begin{theorem}[{\texorpdfstring{CT1.19 — Continuum limit at leading order}{CT1.19 - Continuum limit at leading order}}]\label{tag:CT1-19}
\textbf{Thm.} $L=\mathfrak h\,n_{\mathcal R}$ fixed; $\ \mathfrak h\to0\ \wedge\ n_{\mathcal R}\to\infty\implies$

%dsp-level:0
\[
\sum_{e\in\gamma}ds^2(e)\ \longrightarrow\ \int_\gamma g_{\mu\nu}\,dx^\mu dx^\nu=L\int_0^1g_{\mu\nu}\bigl(x(\lambda)\bigr)\dot x^\mu\dot x^\nu\,d\lambda;
\]

%dsp-level:0
\[
\qquad \text{abs.\ error}=O(\mathfrak hL),\quad\text{rel.\ error}=O(\mathfrak h).
\]

\end{theorem}

\begin{corollary}[{\texorpdfstring{CT1.19a}{CT1.19a}}]\label{tag:CT1-19a}
$\ \widehat{ds}^{\,2}(e)=ds^2(e)/\mathfrak h^2:$ dimensionless, $O(1)$ generically; $\ \widehat{ds}^{\,2}(e)=0\iff\Delta x_e\ \text{null}$; trivial chain [D1.7] $=$ all-null-edge chain.

\textbf{Req.} D0.6, D1.6, D1.7; CT2.6; midpoint Taylor, Riemann sums [math].

\textbf{Pf.}

%dsp-level:0
\[
\mathrm{CT2.6}\implies g_{\mu\nu}\ \text{smooth along}\ \gamma;\qquad
\mathrm{D0.6}:\ |\Delta x^\mu_e|=O(\mathfrak h)
\]

%dsp-level:0
\[
\lambda\in[0,1],\ x^\mu(\lambda)\ \text{smooth}:\quad\Delta x^\mu_e=\mathfrak h\,\dot x^\mu(\lambda_e)+O(\mathfrak h^2)\quad[\text{midpoint Taylor}]
\]

%dsp-level:0
\[
\sum_{e}ds^2(e)=\sum_e g_{\mu\nu}(x_e)\,\mathfrak h^2\dot x^\mu_e\dot x^\nu_e+O(\mathfrak h^3n_{\mathcal R})\quad[\mathrm{D1.6}]
\]

%dsp-level:0
\[
= \mathfrak h\sum_e\mathfrak h\,
\]

%dsp-level:0
\[
g_{\mu\nu}(x_e)\dot x^\mu_e\dot x^\nu_e+O(\mathfrak h^2L)\xrightarrow[\text{Riemann, mesh}\ \mathfrak h/L]{}L\int_0^1g_{\mu\nu}\dot x^\mu\dot x^\nu d\lambda
\]

%dsp-level:0
\[
\text{leading}=O(\mathfrak hL)\implies\text{rel.}=O(\mathfrak h),\ \text{abs.}=O(\mathfrak hL)\to0.\ \blacksquare
\]

\textbf{Pf (Cor).} $[\widehat{ds}^{\,2}]=[\mathfrak h^0]$ [D1.6]; $\ \widehat{ds}^{\,2}(e)=0\iff ds^2(e)=0\iff\Delta x_e$ null; $\ \mathsf{Trivial}\iff\forall e:\ \widehat{ds}^{\,2}(e)=0$ [D1.7]. $\blacksquare$

\textbf{Phys.} At fixed macroscopic length, shrinking the hop length makes the sum of discrete line elements converge to the continuum path integral of the metric, with relative error $O(\mathfrak h)$. Spacetime intervals are literally Riemann sums over edge data — and null edges define the trivial chains that later become photons and gluons.

\textbf{Rmk (on-shell interface).} On-shell integer quantization is no obstruction to this limit: for materialized closed chains of the fully antialigned class (realized — CT1.19b(3)) each $\kappa_e\in\mathbb Z_{\ge0}$ samples $g_{\mu\nu}\dot x^\mu\dot x^\nu(\lambda_e)$ with $O(\mathfrak h)$ slack, and $\sum_e\kappa_e\,\mathfrak h^2=\sum_e ds^2(e)$ identically [D1.6] — so CT1.19 applies verbatim to on-shell chains (abs.\ error $O(\mathfrak hL)$). Machine: \texttt{checks/\allowbreak{}verify\_\allowbreak{}merge\_\allowbreak{}tierA.py} §G.

\end{corollary}

\begin{theorem}[{\texorpdfstring{CT1.19b — On-shell line-element quantization: even antialigned winding, orientation evenness, the geometry quantum, internal holonomy}{CT1.19b - On-shell line-element quantization: even antialigned winding, orientation evenness, the geometry quantum, internal holonomy}}]\label{tag:CT1-19b}
\textbf{Thm.} A closed chain $c$ [D1.4, D1.5] \textbf{materializes} $:=$ it contributes a spectral pole of the masked resolvent at the rung of its active mask [D2.4, D2.5, CT2.5, D1.12, CT0.13] — the corpus's own pole-as-materialization condition as applied in EW.18; instance: wrapping chains [EW.17]. Under EM.20's phase alphabet ($\alpha_e\in\{0,\pi\}$ per real edge, at every accessible rung): $(1)$ \textbf{Even antialigned winding.} For every materialized closed chain the spiral holonomy is single-valued:

%dsp-level:0
\[
\operatorname{Hol}_{\mathrm{spiral}}(c)=e^{i\Phi(c)}=+1,
\]

%dsp-level:0
\[
\qquad \Phi(c)=\sum_{e\in c}\alpha_e\,\widehat{ds}^{\,2}(e)=\pi\,\Theta_{\mathcal P}(c)=2\pi k,\quad k\in\mathbb Z\ \ (\text{winding}\ [\mathrm{D1.7},\ \mathrm{EW.18}])
\]

%dsp-level:0
\[
\Theta_{\mathcal P}(c):=\!\!\sum_{e\in\mathcal P(c)}\!\!\widehat{ds}^{\,2}(e)
\]

%dsp-level:0
\[
\qquad [\mathcal P(c):\ \text{the antialigned edges of}\ c,\ \mathrm{Def\text{-}EM.21};\quad \alpha_e\in\{0,\pi\}:\ \mathrm{EM.20}(1)]
\]

%dsp-level:0
\[
\implies\ \Theta_{\mathcal P}(c)=2k\in2\mathbb Z\qquad(n=2,3,\ge4);\qquad \mathcal P(c)=c\ (\text{fully antialigned class}):\ \ \Theta(c)=2k .
\]

On rung-3 wrapping chains this reduces verbatim to EW.18 (consistency). $(2)$ \textbf{Orientation evenness.} For a bidirectional channel — both orientations $e_{pq}=(p,q)$, $e_{qp}=(q,p)$ in $\mathcal R$:

%dsp-level:0
\[
\widehat{ds}^{\,2}(e_{pq})=\widehat{ds}^{\,2}(e_{qp})
\]

$(3)$ \textbf{The geometry quantum.} For a materialized 2-cycle $c=(e_1,e_2)$, $e_1=(p,q)$, $e_2=(q,p)$ — the minimal closed chain (A1 excludes self-loops [D1.4]) — writing $\kappa:=\widehat{ds}^{\,2}(e_1)=\widehat{ds}^{\,2}(e_2)$ [clause (2)] and $m:=|\mathcal P(c)|\in\{0,1,2\}$:

%dsp-level:0
\[
\Theta_{\mathcal P}(c)=m\kappa\in2\mathbb Z\implies\kappa\in\tfrac2m\mathbb Z\ \ (m\ge1);
\]

%dsp-level:0
\[
\qquad m=2:\ \ \kappa\in\mathbb Z;\quad ds^2(e)=\kappa\,\mathfrak h^2\in\mathfrak h^2\mathbb Z\ [\mathrm{D1.6}];\quad \mathfrak h^2=4G_N\ln2\ [\mathrm{CT2.9}]\implies ds^2(e)=4\kappa\,G_N\ln2
\]

— the interval quantum is one Bekenstein bit-area [CT1.29: $S=\ln2$ at $A=\mathfrak h^2$] per unit $\kappa$ on the fully antialigned class. Attainment and sign classes: $\kappa=1$ attained on the minimal real edge [D1.1] with $m=2$ (floor tight; both alignment values realizable: EM.20(5)); $\kappa=0\iff$ null edge [CT1.19a] (degenerate floor; all-null transit chains are the trivial chains — photons/gluons [D1.7, D3.7]); $\kappa<0$ admissible $=$ timelike [CT1.12: $\operatorname{sig}(g)=(-,+,+,+)$]. $m=0$: the winding condition is vacuous, $\kappa$ unconstrained by clause (1). $(4)$ \textbf{Internal holonomy.} If a subnetwork's interior edges $\Sigma\setminus\partial\Sigma$ [D2.13] decompose edge-disjointly into materialized closed chains $\{c_i\}$, then:

%dsp-level:0
\[
\Theta_{\mathcal P,\mathrm{int}}(\Sigma):=\!\!\sum_{e\ \mathrm{int},\ \alpha_e=\pi}\!\!\widehat{ds}^{\,2}(e)=\sum_i\Theta_{\mathcal P}(c_i)\in2\mathbb Z,
\]

%dsp-level:0
\[
\qquad \mathrm{Hol}_{\mathrm{int}}:=e^{i\Phi_{\mathrm{int}}}=e^{i\pi\Theta_{\mathcal P,\mathrm{int}}}=+1
\]

and the surviving on-shell integer label is the winding $k(c)=\Theta_{\mathcal P}(c)/2$. On \textit{partial} internal sets (proper subsets of a materialized chain's edges) $\Theta_{\mathcal P,\mathrm{int}}$ is a partial sum of integers with free parity: $\mathrm{Hol}_{\mathrm{int}}=\pm1$ well-defined, nontrivial (witness: a single $\kappa=1$ antialigned edge gives $(-1)^1=-1$). Off-shell, $\widehat{ds}^{\,2}(e)$ remains generic continuous [CT1.19a] — quantization is an on-shell (materialization) property, $\neg$a postulate on every edge.

\textbf{Req.} A1; D1.1, D1.4, D1.5, D1.6, D1.7, D1.12; D2.4, D2.5, D2.13; CT0.7(2), CT0.13, CT1.1, CT1.2, C3 (Part 1); CT1.12, CT1.19a, CT1.29, CT2.5, CT2.6 (charts; hypotheses — A §0.2a); D3.7; Def-EM.21, EM.20 (\textbf{E}; forward pointers — C DAG note); EW.17, EW.18 (\textbf{D}; forward pointers — wrapping instance and winding freedom); single-valued eigenmode condition, quadratic-form evenness [math; machine: \texttt{checks/\allowbreak{}verify\_\allowbreak{}merge\_\allowbreak{}tierB.py} §B].

\textbf{Pf.}

%dsp-level:0
\[
(1):\ \mathrm{D1.7}:\ \Theta,\ \Phi=\textstyle\sum_e\alpha_e\widehat{ds}^{\,2}(e),\ \operatorname{Hol}_{\mathrm{spiral}}=e^{i\Phi}\ \text{defined for every chain};
\]

%dsp-level:0
\[
\qquad \mathrm{EW.17}:\ \text{transport around a closed cycle multiplies the eigenmode by}\ \operatorname{Hol}_{\mathrm{spiral}}
\]

%dsp-level:0
\[
\mathrm{CT2.5}\ (\text{as applied in EW.18}):\ \text{pole}
\]

%dsp-level:0
\[
= \text{eigenvalue condition on a single-valued eigenmode}
\]

%dsp-level:0
\[
\implies e^{i\Phi(c)}\qquad = 1;
\]

%dsp-level:0
\[
\qquad \text{winding freedom}\ [\mathrm{EW.18}]:\ \Phi(c)=2\pi k,\ k\in\mathbb Z
\]

%dsp-level:0
\[
\text{domain step (flagged)}:\ \mathrm{EW.17}'\text{s rung-3 restriction is an existence claim about wrapping chains};
\]

%dsp-level:0
\[
\ \mathrm{CT2.5}'\text{s single-valuedness step is chain-generic}\implies\text{the Pf extends to every materialized closed chain with no added premise}
\]

%dsp-level:0
\[
\mathrm{EM.20}(1):\ \alpha_e\in\{0,\pi\}\ \text{at every accessible rung (the case split is rung-blind; CT0.9 alphabet)}\implies\Phi(c)
\]

%dsp-level:0
\[
= \pi\!\!\sum_{e\in\mathcal P(c)}\!\!\widehat{ds}^{\,2}(e)
\]

%dsp-level:0
\[
= \pi\Theta_{\mathcal P}(c)\implies\Theta_{\mathcal P}(c)\qquad = 2k
\]

%dsp-level:0
\[
(2):\ \mathrm{CT2.6}\ \text{charts}:\ \Delta x_{e_{qp}}=-\Delta x_{e_{pq}};
\]

%dsp-level:0
\[
\qquad g_{\mu\nu}\ \text{built solely from}\ S\text{-even data}\ [\mathrm{CT1.12};\ \mathrm{CT1.12a}(2);\ \mathcal E_{++},\mathcal E_{+-}\ S\text{-even: CT1.1, CT1.2, C3};\ S(X)=\mathbb S X\mathbb S:\ \mathrm{CT0.7}(2)]
\]

%dsp-level:0
\[
\mathrm{D1.6}:\ ds^2(e)
\]

%dsp-level:0
\[
= g_{\mu\nu}(e)\,
\]

%dsp-level:0
\[
\Delta x^\mu_e\Delta x^\nu_e\ \text{quadratic in}\ \Delta x\ (\text{every monomial of degree}\ 2,\ \text{incl. the shift cross-term}\ 2g_{0i}\Delta x^0\Delta x^i)\implies\Delta x
\]

%dsp-level:0
\[
\mapsto -\Delta x\ \text{fixes}\ ds^2
\]

%dsp-level:0
\[
(3):\ \Theta_{\mathcal P}(c)=m\kappa,\ m=|\mathcal P(c)|\ [(2):\ \text{both edges share}\ \kappa];
\]

%dsp-level:0
\[
\quad (1):\ m\kappa\in2\mathbb Z;
\]

%dsp-level:0
\[
\quad m=2:\ \kappa\in\mathbb Z;
\]

%dsp-level:0
\[
\quad \mathrm{D1.6}:\ ds^2=\kappa\mathfrak h^2;
\]

%dsp-level:0
\[
\quad \mathrm{CT2.9}:\ \mathfrak h^2=4G_N\ln2;
\]

%dsp-level:0
\[
\quad \mathrm{CT1.29}:\ A=\mathfrak h^2\Rightarrow S=\ln2
\]

%dsp-level:0
\[
\mathrm{D1.1}:\ \widehat{ds}^{\,2}=1\ \text{on the minimal real edge};
\]

%dsp-level:0
\[
\qquad \mathrm{EM.20}(5):\ \alpha_e=0\ \wedge\ \alpha_e=\pi\ \text{each realizable}\implies m=2\ \text{attained}\implies\kappa=1\ \text{attained};
\]

%dsp-level:0
\[
\quad \mathrm{CT1.19a}:\ \widehat{ds}^{\,2}=0\iff\text{null};
\]

%dsp-level:0
\[
\quad \mathrm{D1.7}/\mathrm{D3.7}:\ \text{all-null transit}=\text{trivial chains}
\]

%dsp-level:0
\[
(4):\ \text{edge-disjointness}:\ \Theta_{\mathcal P,\mathrm{int}}=\sum_i\Theta_{\mathcal P}(c_i)\in2\mathbb Z\ [(1)];
\]

%dsp-level:0
\[
\quad e^{i\pi\cdot2\mathbb Z}=+1;
\]

%dsp-level:0
\[
\quad \text{partial sets}:\ \text{single antialigned}\ \kappa=1\ \text{edge}\implies\Theta_{\mathcal P,\mathrm{int}}=1,\ (-1)^1=-1.\ \blacksquare
\]

\textbf{Rmk (locus of the quantization).} The winding constraint reads on the antialigned content: edges with $\alpha_e=0$ drop out of $\Phi(c)$ [D1.7], so $\Theta(c)$ is constrained exactly on the fully antialigned class $\mathcal P(c)=c$; on a general pattern the constrained sum is $\Theta_{\mathcal P}(c)$ [Def-EM.21]. The alphabet $\{0,\pi\}$ — derived, with no effectivity input [EM.20(1); A §0.2a] — is what makes the constraint integral (machine: §B parity sweep).

\textbf{Rmk (sign-factor separation).} Single-valuedness is imposed on the spiral/geometric transport [EW.18]; the registration signs $\sigma_e$ and the wall parity $(-1)^{\aleph}$ are ledger data of the weight $w(C)$ [A8, CT0.9, D3.3], $\neg$of the geometric eigenmode transport.

\textbf{Rmk (area language, scoped).} $\Theta_{\mathcal P}(c)\,\mathfrak h^2$ is the $\kappa$-weighted accumulated line element of the antialigned content — dimensionally an area, and equal to CT1.29's pierced-edge area $A=n_{\mathcal R}\mathfrak h^2$ exactly on the minimal-edge fully antialigned class $\kappa\equiv1$ (then $\Theta_{\mathcal P}=\delta h$; the exactness locus of the rate: CT1.29a(2)). It is $\neg$the metric-enclosed area of $c$: the floor 2-cycle traces a degenerate out-and-back locus — zero enclosed area — yet carries $\Theta_{\mathcal P}=2\kappa\ne0$ at $m=2$, because $ds^2$ is quadratic, hence orientation-blind [clause (2)].

\textbf{Phys.} The $ds^2$ question is settled on-shell by computation, with no new postulate: the antialigned content of a materialized closed chain winds in even quanta, both orientations of a channel agree on its length, the minimal materialized 2-cycle pins the geometry quantum at one Bekenstein bit-area per hop on the fully antialigned class, and internal holonomy reduces to a single integer winding per chain. Off-shell, nothing is quantized — integrality is a property of materialization, not an axiom about edges.

\textbf{Rmk (scope).} Round-trip-closed bidirectional channels realize $\kappa\in\mathbb Z$ per edge; bare cycles admit the rational ladder $\kappa\in\tfrac2L\mathbb Z$, outside hypothesis (ii). The effective-quantum label reads $\bar\kappa$ here [READING]; the transport vocabulary has no corpus referent (nearest declared objects: round-trip bidirectional channels).

\end{theorem}

\begin{theorem}[{\texorpdfstring{CT1.20 — Effective action from the partition function}{CT1.20 - Effective action from the partition function}}]\label{tag:CT1-20}
\textbf{Thm.} $\ Z[g]=\sum_{C:\,g(C)=g}|w(C)|\,e^{i\Phi(C)}$ [$\Phi:$ D1.7+D3.5+parity phases]. Semiclassically ($\delta\Phi\ne0$: pairwise cancellation; $\delta\Phi=0$: dominance):

%dsp-level:0
\[
S_{\mathrm{eff}}[g]:=-i\log Z[g]\ \ (\text{branch}\ [\mathrm{D2.8}]);
\]

%dsp-level:0
\[
\qquad \frac{\delta S_{\mathrm{eff}}}{\delta g^{\mu\nu}}=0\iff\frac{\delta Z}{\delta g^{\mu\nu}}=0\ \ (Z\ne0);
\]

%dsp-level:0
\[
\qquad \text{stationary points dominate}.
\]

\textbf{Req.} A8; D1.7, D2.8, D3.5; CT1.12, CT2.6 ($\Phi$ differentiable in coarse-grained variables); stationary phase [math].

\textbf{Pf.}

%dsp-level:0
\[
\mathrm{CT1.12}:\ C\mapsto g_{\mu\nu}(C)\implies Z\ \text{function of}\ g\ \text{after grouping}\ [\mathrm{A8}]
\]

%dsp-level:0
\[
w(C)=|w(C)|e^{i\Phi(C)};\qquad\Phi=\text{accumulated phase}\ [\mathrm{D1.7},\ \mathrm{D3.5},\ (-1)^\aleph]
\]

%dsp-level:0
\[
\delta\Phi\ne0\ \text{neighbourhoods cancel pairwise};\quad\text{survivors}=\{\delta\Phi=0\}
\]

%dsp-level:0
\[
\quad [\text{stationary phase};\ \Phi\ \text{differentiable}\ [\mathrm{CT2.6}]]
\]

%dsp-level:0
\[
S_{\mathrm{eff}}:=-i\log Z\ \ \text{single-valued on the D2.8 sheet}
\]

%dsp-level:0
\[
\delta Z=0\ \wedge\ Z\ne0\iff\delta\log Z=\delta Z/Z=0\iff\delta S_{\mathrm{eff}}=0.\ \blacksquare
\]

\textbf{Phys.} Grouping ledger configurations by their emergent metric turns the partition sum into a path-integral-like object dominated by stationary phase; its logarithm is an effective action for geometry. Classical spacetime = the stationary set of the ledger sum; quantum fluctuations = the rest.

\end{theorem}

\begin{theorem}[{\texorpdfstring{CT1.21 — The geometric part of the effective action}{CT1.21 - The geometric part of the effective action}}]\label{tag:CT1-21}
\textbf{Thm.} $\ S_{\mathrm{geom}}[g]=\dfrac{1}{2\kappa^2}\displaystyle\int d^4x\,\sqrt{-g}\,(R-2\Lambda).$

\textbf{Req.} D1.9, D4.3; CT1.20, CT2.7.

\textbf{Pf.}

%dsp-level:0
\[
\mathrm{CT1.20}\implies S_{\mathrm{eff}}\ \text{local, diffeomorphism-invariant},\ \le2\ \text{nd order}\ [\mathrm{D1.9}+\text{relabeling}]
\]

%dsp-level:0
\[
\mathrm{CT2.7}:\ \text{unique such functional in}\ D
\]

%dsp-level:0
\[
= 4\ (\text{topological density dropped};\ F_{D^+}=\Lambda\mathbb I\ \text{included}\ [\mathrm{D4.3}])\implies\text{Thm}.\ \blacksquare
\]

\textbf{Phys.} The geometry part of the effective action is exactly Einstein–Hilbert with a cosmological term — derived from the structure of the ledger sum, with $\Lambda$ supplied by the uniform dark aspect of A8.

\end{theorem}

\begin{theorem}[{\texorpdfstring{CT1.22 — Stress-energy tensor from the force operator}{CT1.22 - Stress-energy tensor from the force operator}}]\label{tag:CT1-22}
\textbf{Thm.} With $D^+$ removed from $Z_{\mathrm{force}}$ (carried by $\Lambda$ in $S_{\mathrm{geom}}$):

%dsp-level:0
\[
S_{\mathrm{force}}[g]:=-i\log Z_{\mathrm{force}}[g];
\]

%dsp-level:0
\[
\qquad Z_{\mathrm{force}}[g]=\bigl\langle\mathrm{out}\bigl|(1-\mathbf F_M^{(n)})^{-1}\bigr|\mathrm{in}\bigr\rangle;
\]

%dsp-level:0
\[
\qquad T_{\mu\nu}:=-\frac{2}{\sqrt{-g}}\frac{\delta S_{\mathrm{force}}}{\delta g^{\mu\nu}};
\]

%dsp-level:0
\[
T_{\mu\nu}=T^{\mathrm{bd}}_{\mu\nu}+T^{\mathrm{tr}}_{\mu\nu};
\]

%dsp-level:0
\[
\qquad T^{\mathrm{bd}}_{\mu\nu}=(\rho+p)u_\mu u_\nu+p\,g_{\mu\nu};
\]

%dsp-level:0
\[
\qquad T^{\mathrm{tr}}_{\mu\nu}=\sum_a\Bigl(F^{(a)}_{\mu\alpha}F^{(a)\ \alpha}_{\ \ \nu}-\tfrac14g_{\mu\nu}F^{(a)}_{\alpha\beta}F^{(a)\alpha\beta}\Bigr),
\]

the sum over aspects whose leading transit representation is a connection.

\textbf{Req.} A8; D1.1, D1.5, D1.10, D4.3; CT1.19, CT1.21, CT2.4; $\delta\sqrt{-g}=-\tfrac12\sqrt{-g}\,g_{\mu\nu}\delta g^{\mu\nu}$, $\delta(F_{\alpha\beta}F^{\alpha\beta})=2F_{\mu\rho}F_\nu{}^\rho\delta g^{\mu\nu}$ [math].

\textbf{Pf.}

%dsp-level:0
\[
\mathrm{A8}\ \text{blocks};\ \mathrm{CT2.4}:\ \text{mask retains exactly}\ \mathcal A(n);
\]

%dsp-level:0
\[
\qquad \Delta S=\mathsf E_{\mathrm{bd}}+\mathsf E_{\mathrm{tr}}\ [\mathrm{D1.5}]\implies\forall\ \text{contribution}\in\text{exactly one channel}\implies T=T^{\mathrm{bd}}+T^{\mathrm{tr}}
\]

%dsp-level:0
\[
\text{bd}:\ \text{closed chain}\implies\text{rest frame selects timelike}\ u^\mu,\ \neg\text{spatial direction};\ \ \text{content}=(\rho,p)\ [\mathrm{D1.10}]
\]

%dsp-level:0
\[
\text{most general symmetric 2-tensor from}\ (g,u,\rho,p),
\]

%dsp-level:0
\[
\ u\text{-rotation-invariant}
\]

%dsp-level:0
\[
= (\rho+p)u_\mu u_\nu+p\,\qquad g_{\mu\nu}
\]

%dsp-level:0
\[
\quad [\text{math}]
\]

%dsp-level:0
\[
\text{tr}:\ \mathrm{D1.1}:\ \text{edges carry phases}\implies\text{gauge-invariant local datum}
\]

%dsp-level:0
\[
= \text{holonomy};
\]

%dsp-level:0
\[
\quad\text{holonomy/area}\to F^{(a)}_{\mu\nu}\ [\mathrm{CT1.19}]
\]

%dsp-level:0
\[
\text{unique local, gauge-invariant},\ \le2\ \text{nd order, quadratic in}\ F:
\]

%dsp-level:0
\[
\quad S_a=-\tfrac14\int d^4x\sqrt{-g}\,F^{(a)}_{\mu\nu}F^{(a)\mu\nu}
\]

%dsp-level:0
\[
\delta S_a
\]

%dsp-level:0
\[
= -\tfrac14\int\sqrt{-g}\bigl[-\tfrac12g_{\mu\nu}F^2+2F_{\mu\rho}F_\nu{}^\rho\bigr]\delta g^{\mu\nu}
\]

%dsp-level:0
\[
\implies T^{(a)}_{\mu\nu}
\]

%dsp-level:0
\[
= F^{(a)}_{\mu\rho}F^{(a)\ \rho}_{\ \ \nu}-\tfrac14g_{\mu\nu}F^{(a)}_{\alpha\beta}F^{(a)\alpha\beta}
\]

%dsp-level:0
\[
\text{no double counting}:\ D^+\ \text{excluded from}\ Z_{\mathrm{force}}\ \text{by construction};
\]

%dsp-level:0
\[
\ \ \mathrm{D4.3}:\ \neg\text{config.-dependent}\implies\text{removal}
\]

%dsp-level:0
\[
= \text{constant factor}\implies\delta\ \text{unaffected}.\ \blacksquare
\]

\textbf{Rmk.} Abelian case written out in EM.15 (Part 2), which depends on this result, not the reverse.

\textbf{Phys.} The right-hand side of the field equations is built from the same force operator that built the geometry: bound (closed-chain, fermionic) content yields perfect-fluid stress; transit (connection-carrying, bosonic) content yields Maxwell-type stress. $T_{\mu\nu}$ is defined by metric variation of the masked-force resolvent, with the uniform dark part counted once, as $\Lambda$.

\end{theorem}

\begin{theorem}[{\texorpdfstring{CT1.23 — Field equations}{CT1.23 - Field equations}}]\label{tag:CT1-23}
\textbf{Thm.}

%dsp-level:0
\[
S_{\mathrm{eff}}=S_{\mathrm{geom}}+S_{\mathrm{force}};\qquad
\delta S_{\mathrm{eff}}=0\iff G_{\mu\nu}+\Lambda g_{\mu\nu}=\kappa^2T_{\mu\nu};
\]

%dsp-level:0
\[
\text{dimensions}\ [\mathrm{CT2.1}]:\quad\mathfrak q^{+1}=\mathfrak q^{-1}\cdot\mathfrak q^{+2}=\mathfrak q^{+1};
\]

%dsp-level:0
\[
\qquad \kappa:=\text{coefficient fixed by}\ S_{\mathrm{geom}}\ \text{normalization}.
\]

\textbf{Req.} CT1.21, CT1.22, CT2.1; $\delta(\sqrt{-g}R)=\sqrt{-g}\,G_{\mu\nu}\delta g^{\mu\nu}+\partial(\cdots)$, $\delta\sqrt{-g}=-\tfrac12\sqrt{-g}\,g_{\mu\nu}\delta g^{\mu\nu}$ [math].

\textbf{Pf.}

%dsp-level:0
\[
\delta S_{\mathrm{geom}}=\frac{1}{2\kappa^2}\int d^4x\sqrt{-g}\,\bigl(G_{\mu\nu}+\Lambda g_{\mu\nu}\bigr)\delta g^{\mu\nu}\quad[\mathrm{CT1.21};\ \delta\text{-identities}]
\]

%dsp-level:0
\[
\delta S_{\mathrm{force}}=-\tfrac12\int d^4x\sqrt{-g}\,T_{\mu\nu}\,\delta g^{\mu\nu}\quad[\mathrm{CT1.22}\ \text{definition}]
\]

%dsp-level:0
\[
\forall\,\delta g^{\mu\nu}:\ \ \delta S_{\mathrm{eff}}=0\iff\frac{1}{2\kappa^2}\bigl(G_{\mu\nu}+\Lambda g_{\mu\nu}\bigr)-\tfrac12T_{\mu\nu}=0\iff G_{\mu\nu}+\Lambda g_{\mu\nu}=\kappa^2T_{\mu\nu}
\]

%dsp-level:0
\[
\implies\kappa=\text{the}\ S_{\mathrm{geom}}\ \text{coefficient (no other normalization admits source}=T_{\mu\nu});
\]

%dsp-level:0
\[
\qquad \text{dimensions}\ [\mathrm{CT2.1}]:\ \text{both sides}\ \mathfrak q^{+1}.\ \blacksquare
\]

\textbf{Phys.} Varying the derived effective action yields Einstein-type field equations with a cosmological term and the ledger-sourced stress tensor — and simultaneously fixes $\kappa$ as that coefficient. In this corpus the field equations of gravity are a theorem with both sides constructed from A1–A9.

\textbf{Rmk (dependency audit).} The field-equation chain (CT1.20 $\to$ CT1.21 $\to$ CT1.22 $\to$ CT1.23, conservation CT1.24/CT1.26) rests solely on the declared inputs listed in A §0.2a (CT2.6 (i)–(iii) via CT1.19/CT2.6; Lovelock [math]) — no clause depends on persistence data, on any tick-rate/field coupling (count $0$ in Parts 1–2), or on the materialization-selection question (outside Parts 1–2; downstream of and disjoint from this derivation); the spatial components are no shakier than the temporal ones. On-shell interface: CT1.19 Rmk. The "chosen vs forced" upgrade claim has no corpus referent and is blocked by proven no-gos.

\end{theorem}

\begin{theorem}[{\texorpdfstring{CT1.24 — Conservation from the field equations}{CT1.24 - Conservation from the field equations}}]\label{tag:CT1-24}
\textbf{Thm.} $\ \Lambda=\mathrm{const}\implies\nabla^\mu T_{\mu\nu}=0.$

\textbf{Req.} CT1.23, CT2.9 ($\kappa\ne0$); contracted Bianchi $\nabla^\mu G_{\mu\nu}=0$; metric compatibility $\nabla^\mu g_{\mu\nu}=0$ [math].

\textbf{Pf.}

%dsp-level:0
\[
\nabla^\mu G_{\mu\nu}=0\ [\text{Bianchi}];\qquad\nabla^\mu(\Lambda g_{\mu\nu})=\Lambda\nabla^\mu g_{\mu\nu}=0\ [\text{compatibility};\ \Lambda=\mathrm{const}]
\]

%dsp-level:0
\[
\nabla^\mu(\mathrm{CT1.23}):\ \ 0+0=\kappa^2\nabla^\mu T_{\mu\nu};\qquad\kappa^2=\tfrac{2\pi\mathfrak h^2}{\ln2}>0\ [\mathrm{CT2.9}]\implies\nabla^\mu T_{\mu\nu}=0.\ \blacksquare
\]

\textbf{Phys.} The Bianchi identity forces covariant conservation of the ledger stress tensor: energy–momentum conservation is a geometric consequence of the field equations, never an extra assumption.

\end{theorem}

\begin{theorem}[{\texorpdfstring{CT1.25 — The cosmological constant is the $D^+$ aspect}{CT1.25 - The cosmological constant is the aspect}}]\label{tag:CT1-25}
\textbf{Thm.} $\ F_{D^+}=\Lambda\mathbb I\implies T^{(D^+)}_{\mu\nu}=-\dfrac{\Lambda}{\kappa^2}g_{\mu\nu};\qquad\rho_{D^+}=\dfrac{\Lambda}{\kappa^2};\quad p_{D^+}=-\rho_{D^+};\quad w_{D^+}:=\dfrac{p}{\rho}=-1.$

\textbf{Req.} A8; D4.3; CT1.3 (Part 1), CT1.21, CT1.22; $\delta\sqrt{-g}$ identity [math].

\textbf{Pf.}

%dsp-level:0
\[
E_0\in Z\ [\mathrm{CT1.3}]\implies D^+\ \text{commutes with every aspect}\implies\neg\exists\ \text{preferred direction}:\ \text{uniform}\wedge\text{isotropic}
\]

%dsp-level:0
\[
S_{\mathrm{geom}}\ni-\frac{1}{\kappa^2}\int d^4x\sqrt{-g}\,\Lambda\quad[\mathrm{CT1.21}]
\]

%dsp-level:0
\[
T^{(D^+)}_{\mu\nu}=-\frac{2}{\sqrt{-g}}\frac{\delta}{\delta g^{\mu\nu}}\Bigl(-\frac{\Lambda}{\kappa^2}\int\sqrt{-g}\Bigr)=-\frac{\Lambda}{\kappa^2}g_{\mu\nu}
\]

%dsp-level:0
\[
\quad [\mathrm{CT1.22}\ \text{definition};\ \delta\sqrt{-g}]
\]

%dsp-level:0
\[
\propto g_{\mu\nu}\implies\neg\exists\ \text{rest frame};\qquad
-\rho\,g_{\mu\nu}\overset!=(\rho+p)u_\mu u_\nu+p\,g_{\mu\nu}\implies\rho+p=0\implies p=-\rho
\]

%dsp-level:0
\[
\rho_{D^+}=\frac{\Lambda}{\kappa^2};\qquad w_{D^+}=\frac{p}{\rho}=-1\ \ (\text{derived},\ \text{zero freedom});\qquad\Lambda=\text{configuration data}\ [\mathrm{P4}].\ \blacksquare
\]

\textbf{Phys.} The uniform aspect behaves exactly as a cosmological fluid with equation of state $w=-1$: negative pressure equal in magnitude to its energy density — a derived value with no freedom. Only its magnitude $\Lambda$ remains configuration data (P4): the corpus predicts the dark-energy equation of state, not its size.

\end{theorem}

\begin{theorem}[{\texorpdfstring{CT1.26 — Conservation from mask invariance alone}{CT1.26 - Conservation from mask invariance alone}}]\label{tag:CT1-26}
\textbf{Thm.} $\bigl(\mathcal P_M^{(n)}\bigr)^2=\mathcal P_M^{(n)}$ $\ \wedge\ $ weight invariance under mask-preserving relabeling $\implies\nabla_\mu T^{\mu\nu}=0$, independently of CT1.23.

\textbf{Req.} CT1.22 ($T$ definition), CT2.2; $\delta_\xi g_{\mu\nu}=\nabla_\mu\xi_\nu+\nabla_\nu\xi_\mu$; integration by parts; fundamental lemma ($\forall\xi$ compact: $\int f\xi=0\Rightarrow f=0$) [math].

\textbf{Pf.}

%dsp-level:0
\[
\mathrm{CT2.2}:\ \text{sum runs over}\ \operatorname{Im}\mathcal P_M^{(n)}\implies\text{invariant under image-preserving transformations}
\]

%dsp-level:0
\[
\xi^\mu=\text{infinitesimal relabeling};\ \ \text{mask}\wedge\text{weights invariant}\ [\text{hyp}]\implies\delta_\xi Z=0
\]

%dsp-level:0
\[
\delta_\xi g_{\mu\nu}=\nabla_\mu\xi_\nu+\nabla_\nu\xi_\mu\quad[\text{Lie derivative}]
\]

%dsp-level:0
\[
0=\delta_\xi\log Z=\int d^4x\sqrt{-g}\,\tfrac12T^{\mu\nu}\bigl(\nabla_\mu\xi_\nu+\nabla_\nu\xi_\mu\bigr)=\int d^4x\sqrt{-g}\,T^{\mu\nu}\nabla_\mu\xi_\nu
\]

%dsp-level:0
\[
\quad [T^{\mu\nu}=T^{\nu\mu}]
\]

%dsp-level:0
\[
\xi\ \text{compactly supported}:\ \ \text{IBP}\implies0=-\int d^4x\sqrt{-g}\,\bigl(\nabla_\mu T^{\mu\nu}\bigr)\xi_\nu\implies\nabla_\mu T^{\mu\nu}=0
\]

%dsp-level:0
\[
\quad [\text{fundamental lemma}].\ \blacksquare
\]

\textbf{Phys.} Conservation has a second, independent origin: invariance of the ledger sum under relabelings — discrete diffeomorphism invariance. CT1.24 and CT1.26 are two separate routes to $\nabla T=0$; their agreement is a consistency check on the whole continuum layer.

\end{theorem}

\begin{theorem}[{\texorpdfstring{CT1.27 — Quantum geometry state}{CT1.27 - Quantum geometry state}}]\label{tag:CT1-27}
\textbf{Thm.} $A(C):=\sum_\sigma w(C)$; $\ \{|g(C)\rangle\}:\ \langle g(C)|g(C')\rangle=\delta_{CC'}$. Then

%dsp-level:0
\[
|\Psi\rangle:=\sum_CA(C)\,|g(C)\rangle\in\mathcal H\ [\mathrm{CT2.8}];
\]

%dsp-level:0
\[
\qquad Z_B=\langle B|\Psi\rangle\ \ \forall\ \text{boundary}\ |B\rangle;
\]

%dsp-level:0
\[
\qquad \langle B|=\Theta\text{-dual}\implies Z=\langle\Psi|\Psi\rangle=\sum_C|A(C)|^2\ge0;
\]

%dsp-level:0
\[
\text{classical metric}:=g^\ast_{\mu\nu}=g_{\mu\nu}(C^\ast),\ \ \delta\Phi(C^\ast)=0\ [\mathrm{CT1.20}];
\]

%dsp-level:0
\[
\qquad \text{fluctuations}:=\delta g_{\mu\nu}(C)=g_{\mu\nu}(C)-g^\ast_{\mu\nu}\ \ (\text{subleading}\ [\mathrm{CT1.19}]).
\]

\textbf{Req.} A8; CT1.12, CT1.19, CT1.20, CT2.8.

\textbf{Pf.}

%dsp-level:0
\[
Z(\mathcal R)=\sum_C\sum_\sigma w(C)=\sum_CA(C)\quad[\mathrm{A8}]\implies\text{coefficients}=\text{terms of}\ Z
\]

%dsp-level:0
\[
\mathrm{CT1.12}:\ C\mapsto g_{\mu\nu}(C);\quad\langle g(C)|g(C')\rangle=\delta_{CC'}\ [\text{declaration}]\implies|\Psi\rangle\in\mathcal H\ [\mathrm{CT2.8}]
\]

%dsp-level:0
\[
Z_B=\langle B|\Psi\rangle\quad[\text{linearity}];\qquad
\langle B|=\Theta\text{-dual}:\ \ Z=\sum_C|A(C)|^2=\langle\Psi|\Psi\rangle\ge0\quad[\mathrm{CT2.8}(3)]
\]

%dsp-level:0
\[
\mathrm{CT1.20}:\ \text{dominant}\ C^\ast:\ \delta\Phi(C^\ast)=0;
\]

%dsp-level:0
\[
\quad g^\ast:=g(C^\ast);\quad\text{non-stationary subleading in the}\ \mathfrak h\text{-expansion}\ [\mathrm{CT1.19}].\ \blacksquare
\]

\textbf{Phys.} The corpus carries genuine quantum states of geometry: superpositions of emergent metrics weighted by ledger amplitudes, with positive-definite inner product. Classical spacetime is the stationary-phase component; the rest of the ket is quantum geometry — superposition applied to metrics themselves, built from A8.

\end{theorem}

\begin{theorem}[{\texorpdfstring{CT1.28 — Singularity resolution by the saturation lock}{CT1.28 - Singularity resolution by the saturation lock}}]\label{tag:CT1-28}
\textbf{Thm.} $\kappa_{\mathrm{ind}}(v):=$ accessible curvature indicator at $v$ (built from accessible registrations):

%dsp-level:0
\[
\kappa_{\mathrm{ind}}(v)\le f\bigl(C_{\mathrm{cap}},\mathfrak h\bigr)=f(16,\mathfrak h)<\infty\quad\forall\,\mathfrak h>0;
\qquad\Longrightarrow\quad\neg\exists\ \text{accessible curvature singularity}.
\]

(Bound $\mathfrak h$-dependent: $\to\infty$ as $\mathfrak h\to0$; finite at fixed $\mathfrak h>0$.)

\textbf{Req.} D1.11, D4.1, D4.2; CT0.5, CT0.6, CT0.11 (Part 1), CT1.19.

\textbf{Pf.}

%dsp-level:0
\[
\mathrm{CT1.19}:\ \text{curvature}
\]

%dsp-level:0
\[
= \text{holonomies per unit hop area}
\]

%dsp-level:0
\[
= \text{function of registrations}\implies\text{singularity}\Rightarrow\text{accessible information density unbounded}
\]

%dsp-level:0
\[
\mathrm{D1.11}:\ \text{accessible content at}\ v\leftrightarrow n(v);\qquad
\mathrm{CT0.5}/\mathrm{CT0.6}:\ D_{\mathrm{tick}}(n)=2^{n+1},\ C_{\mathrm{cap}}=16
\]

%dsp-level:0
\[
\mathrm{CT0.11}:\ \text{admissible}\iff n\le3;\qquad n\ge4\implies O(n)=2^{n+1}-16>0
\]

%dsp-level:0
\[
\mathrm{D4.1}/\mathrm{D4.2}:\ n\ge4\implies\text{config}\in F_{D^-},\ \text{excess absorbed}\wedge\text{locked}\ (\in\mathfrak L,\ \neg\text{accessible})
\]

%dsp-level:0
\[
\implies\text{accessible}:\ n\le3\ \wedge\ \text{registers}\le16\implies\kappa_{\mathrm{ind}}\le f(16,\mathfrak h)<\infty\ \ (\mathfrak h>0).\ \blacksquare
\]

\textbf{Phys.} Divergent curvature would require unbounded accessible information at a vertex, but the channel caps access at rung 3 / 16 registers and routes all excess into the locked dark sector. Singularities are not cured by new forces — they are administratively excluded: oversubscription is absorbed before it can become geometry.

\end{theorem}

\phantomsection
\section[§C.IV — Entropy, Calibration, and the Single Free Scale]{\texorpdfstring{§C.IV — Entropy, Calibration, and the Single Free Scale}{SC.IV - Entropy, Calibration, and the Single Free Scale}}

\begin{theorem}[{\texorpdfstring{CT1.29 — Boundary entropy from channel capacity}{CT1.29 - Boundary entropy from channel capacity}}]\label{tag:CT1-29}
\textbf{Thm.} $\partial\mathcal B:$ closed real-edge surface separating a $D^-$-locked interior. Then

%dsp-level:0
\[
\Omega_{\partial\mathcal B}=2^{\,n_{\mathcal R}(\partial\mathcal B)};\qquad
S_{\partial\mathcal B}=\ln\Omega_{\partial\mathcal B}=n_{\mathcal R}(\partial\mathcal B)\ln2\ \ \text{nats}=n_{\mathcal R}(\partial\mathcal B)\ \ \text{bits};
\]

%dsp-level:0
\[
A:=n_{\mathcal R}(\partial\mathcal B)\,\mathfrak h^2\implies S_{\partial\mathcal B}=\frac{A}{\mathfrak h^2}\ln2 .
\]

\textbf{Req.} A2; C9; CT0.9 (Part 1), CT1.19.

\textbf{Pf.}

%dsp-level:0
\[
\mathrm{A2}:\ \text{each pair carries}\ 1\mathfrak q\implies\text{each boundary real edge}=1\ \text{independent channel}
\]

%dsp-level:0
\[
\mathrm{CT0.9}:\ \phi_e\in\{0,\pi\}\implies2\ \text{microstates per edge},\ \neg\text{more}
\]

%dsp-level:0
\[
\text{edges independent}\implies\Omega_{\partial\mathcal B}=2^{\,n_{\mathcal R}(\partial\mathcal B)};
\]

%dsp-level:0
\[
\qquad S_{\partial\mathcal B}=\ln\Omega_{\partial\mathcal B}=n_{\mathcal R}(\partial\mathcal B)\ln2\ \ \text{nats}=n_{\mathcal R}(\partial\mathcal B)\ \text{bits}\ [\mathrm{C9}]
\]

%dsp-level:0
\[
\mathrm{CT1.19}:\ \text{edge subtends hop length}\ \mathfrak h\implies\text{area per pierced boundary element}
\]

%dsp-level:0
\[
= \mathfrak h^2
\]

%dsp-level:0
\[
\implies A
\]

%dsp-level:0
\[
= n_{\mathcal R}(\partial\mathcal B)\,\qquad \mathfrak h^2
\]

%dsp-level:0
\[
\implies n_{\mathcal R}(\partial\mathcal B)=A/\mathfrak h^2\implies S_{\partial\mathcal B}=\frac{A}{\mathfrak h^2}\ln2\ \ (\neg\exists\ \text{constant beyond}\ \mathfrak h).\ \blacksquare
\]

\textbf{Phys.} Boundary entropy is the boundary’s channel capacity counted in bits: one bit per hop-area $\mathfrak h^2$. The area law is a theorem of capacity bookkeeping — its form $S\propto A$ is derived, and its coefficient sets the calibration of gravity (CT2.9).

\end{theorem}

\begin{theorem}[{\texorpdfstring{CT1.29a — The mean edge quantum: structure, sign, and the exactness locus of the entropy–area rate}{CT1.29a - The mean edge quantum: structure, sign, and the exactness locus of the entropy-area rate}}]\label{tag:CT1-29a}
\textbf{Thm.} Hypotheses: (i) $c$ a closed real-edge chain (bound content [D1.5]); (ii) per-edge integer quanta $\kappa_e:=\widehat{ds}^{\,2}(e)\in\mathbb Z_{\ge1}$ [D1.6, CT1.19a; all-nonnull class], with $\Theta(c):=\sum_{e\in c}\kappa_e$, $\delta h(c):=n_{\mathcal R}(c)$, $\bar\kappa(c):=\Theta/\delta h$. The hypothesis class is nonempty on materialized closed chains of the fully antialigned class [CT1.19b(3)]. Then, as exact identities over $\mathbb Q$:

%dsp-level:0
\[
(1)\quad S_{\mathrm{reg}}(c)=\delta h\ln2\ [\mathrm{CT0.9},\ \mathrm{C9}];
\]

%dsp-level:0
\[
\qquad A_{\mathrm{loop}}(c):=\Theta\,\mathfrak h^2\ [\mathrm{D1.6}];
\]

%dsp-level:0
\[
\qquad \frac{S_{\mathrm{reg}}}{A_{\mathrm{loop}}}=\frac{\ln2}{\mathfrak h^2}\cdot\frac1{\bar\kappa}=\frac{\delta h\ln2}{\Theta\,\mathfrak h^2};
\]

%dsp-level:0
\[
(2)\quad \text{CT1.29's rate }\ S=\frac{A}{\mathfrak h^2}\ln2\ \text{holds exactly on }\ c\iff\bar\kappa=1\iff\kappa_e\equiv1\ \text{pointwise};
\]

%dsp-level:0
\[
(3)\quad \bar\kappa\ge1\ \text{with equality}\iff\kappa\equiv1;
\]

%dsp-level:0
\[
\qquad \operatorname{sign}:\ \bar\kappa>0\ \text{always}\ (\Theta\ge\delta h\ge1);
\]

%dsp-level:0
\[
\qquad \bar\kappa\in\tfrac1{\delta h}\mathbb Z_{\ge\delta h};
\]

%dsp-level:0
\[
(4)\quad \text{uniform-}\kappa\ \text{chains}:\ \bar\kappa=\kappa\in\mathbb Z_{\ge1}\ (\text{exact rate only at }\kappa=1);
\]

%dsp-level:0
\[
\qquad \text{deviation}=\frac1{\bar\kappa}\in\mathbb Q:\ \neg\text{free parameter}\ [\mathrm{P1},\ \mathrm{P5}].
\]

\textbf{Req.} D1.5, D1.6, D3.7; C9; CT0.9 (Part 1); CT1.19a, CT1.29 (Part 2); hypotheses (i)–(ii) above (declared at use; realization on materialized closed chains of the fully antialigned class — CT1.19b(3), A §0.2a note); exact rational arithmetic [math; machine: \texttt{checks/\allowbreak{}verify\_\allowbreak{}merge\_\allowbreak{}tierA.py} §F].

\textbf{Pf.}

%dsp-level:0
\[
(1):\ 1\ \text{real edge}=\ln2\ \text{nats}\ [\mathrm{CT0.9}]\implies S_{\mathrm{reg}}=\delta h\ln2;
\]

%dsp-level:0
\[
\qquad A_{\mathrm{loop}}=\oint_c ds^2=\sum_e\widehat{ds}^{\,2}(e)\,\mathfrak h^2=\Theta\,\mathfrak h^2\ [\mathrm{D1.6}]
\]

%dsp-level:0
\[
\frac{S_{\mathrm{reg}}}{A_{\mathrm{loop}}}=\frac{\delta h\ln2}{\Theta\,\mathfrak h^2}=\frac{\ln2}{\mathfrak h^2}\cdot\frac{\delta h}{\Theta}=\frac{\ln2}{\mathfrak h^2}\cdot\frac1{\bar\kappa}.\ \blacksquare_{(1)}
\]

%dsp-level:0
\[
(2):\ \text{CT1.29}:\ S_{\partial\mathcal B}=\frac{A}{\mathfrak h^2}\ln2\iff S_{\mathrm{reg}}/A_{\mathrm{loop}}=\ln2/\mathfrak h^2\iff\bar\kappa=1\ (1);
\]

%dsp-level:0
\[
\qquad \kappa_e\ge1\ \forall e\implies\bar\kappa=1\iff\kappa_e=1\ \forall e
\]

%dsp-level:0
\[
(3):\ \Theta=\sum\kappa_e\ge\sum1=\delta h\ge1;\qquad \bar\kappa=\Theta/\delta h\in\tfrac1{\delta h}\mathbb Z;\qquad(4):\ \kappa_e\equiv\kappa\implies\Theta=\kappa\,\delta h.\ \blacksquare
\]

\textbf{Phys.} The Bekenstein-type rate of CT1.29 is exact precisely on minimal chains; every heavier integer winding dilutes it by the rational factor $1/\bar\kappa$ — computable from ledger data, never fitted. The sign question is settled with the structure: the mean edge quantum is a positive rational of the form (loop total)/(hop count), integer exactly on uniform chains.

\end{theorem}

\begin{theorem}[{\texorpdfstring{CT2.9 — Calibration closure}{CT2.9 - Calibration closure}}]\label{tag:CT2-9}
\textbf{Thm.} $G_N$ defined by $S_{\partial\mathcal B}\overset!=\dfrac{A}{4G_N}$; $\ \kappa$ by the CT1.23 normalization. Then

%dsp-level:0
\[
\mathfrak h^2=4\,G_N\ln2\iff G_N=\frac{\mathfrak h^2}{4\ln2};\qquad
\kappa^2=8\pi G_N=\frac{2\pi\,\mathfrak h^2}{\ln2};\qquad
\frac{A}{4G_N}=\frac{A\ln2}{\mathfrak h^2}=S_{\partial\mathcal B};
\]

%dsp-level:0
\[
[G_N]=[\kappa^2]=\mathfrak q^{-1}\ [\mathrm{CT2.1}];\qquad G_N,\kappa=\text{names for combinations of}\ \mathfrak h\ (\neg\text{independent}).
\]

\textbf{Req.} CT1.23, CT1.29, CT2.1.

\textbf{Pf.}

%dsp-level:0
\[
\mathrm{CT1.29}:\ \ S_{\partial\mathcal B}=\frac{A\ln2}{\mathfrak h^2};
\]

%dsp-level:0
\[
\qquad \frac{A\ln2}{\mathfrak h^2}=\frac{A}{4G_N}\iff\frac{\ln2}{\mathfrak h^2}=\frac{1}{4G_N}\iff\mathfrak h^2=4G_N\ln2\iff G_N=\frac{\mathfrak h^2}{4\ln2}
\]

%dsp-level:0
\[
\text{check}:\ \ \frac{A}{4\cdot\dfrac{\mathfrak h^2}{4\ln2}}=\frac{A\ln2}{\mathfrak h^2}=S_{\partial\mathcal B}\quad(\ln2\ \text{appears once},\ \neg\text{twice})
\]

%dsp-level:0
\[
\mathrm{CT1.23}:\ S_{\mathrm{geom}}\ \text{coefficient}
\]

%dsp-level:0
\[
= \tfrac{1}{2\kappa^2},
\]

%dsp-level:0
\[
\ \text{source}
\]

%dsp-level:0
\[
= T_{\mu\nu}\implies\kappa^2:
\]

%dsp-level:0
\[
= 8\pi G_N
\]

%dsp-level:0
\[
= 8\pi\cdot\frac{\mathfrak h^2}{4\ln2}
\]

%dsp-level:0
\[
= \frac{2\pi\,\mathfrak h^2}{\ln2}
\]

%dsp-level:0
\[
[\mathfrak h^2]=\mathfrak q^{-1}\implies[G_N]=[\kappa^2]=\mathfrak q^{-1}\ [\mathrm{CT2.1}];\qquad\mathfrak h\ \text{given}\implies G_N,\kappa\ \text{determined}.\ \blacksquare
\]

\textbf{Phys.} Demanding that the boundary-entropy law take the quarter-area form locks the hop length, Newton’s constant and the gravitational coupling together: $\mathfrak h^2=4G_N\ln2$. The $\ln2$ is the one bit per edge of CT0.9 — the theory’s information quantum appearing directly in the calibration of gravity.

\end{theorem}

\begin{theorem}[{\texorpdfstring{CT2.11 — One free scale}{CT2.11 - One free scale}}]\label{tag:CT2-11}
\textbf{Thm.}

%dsp-level:0
\[
\forall Q\ \text{dimensionful}:\ \ Q=c_Q\,\mathfrak h^{k_Q},\ c_Q\ \text{fixed};
\]

%dsp-level:0
\[
\qquad \forall q\ \text{dimensionless-derived}:\ \ q=\text{fixed number},\ \mathfrak h\text{-free};
\]

%dsp-level:0
\[
\mathfrak h\mapsto\lambda\mathfrak h:\ \ Q\mapsto\lambda^{k_Q}Q,\quad\text{P2 quantities invariant}\implies\text{rescaling}=\text{symmetry}\implies\exists!\ \text{calibration d.o.f.}
\]

\textbf{Req.} D2.9; P1, P2, P4; CT0.5, CT0.6, CT0.11(a), CT1.1 (Part 1); CT1.29, CT2.1, CT2.9; remaining P2 entries (\textbf{D}, \textbf{E}, \textbf{F}, \textbf{G}).

\textbf{Pf.}

%dsp-level:0
\[
\text{sole dimensionful input}=\mathfrak h\quad[\mathrm{D2.9},\ \mathrm{P1}]
\]

%dsp-level:0
\[
\mathrm{CT2.9}:\ G_N,
\]

%dsp-level:0
\[
\kappa^2
\]

%dsp-level:0
\[
= \text{fixed multiples of}\ \mathfrak h^2;
\]

%dsp-level:0
\[
\quad\mathrm{CT2.1}:\ \text{lengths/areas/energies/densities}
\]

%dsp-level:0
\[
= \text{fixed powers};
\]

%dsp-level:0
\[
\quad\mathrm{CT1.29}:\ S
\]

%dsp-level:0
\[
= \text{pure count}
\]

%dsp-level:0
\[
\mathrm{P2}:\ (7,3,3,3)\ [\mathrm{CT1.1}];
\]

%dsp-level:0
\[
\ 16\ [\mathrm{CT0.6}];
\]

%dsp-level:0
\[
\ 2^{n+1}\ [\mathrm{CT0.5}];
\]

%dsp-level:0
\[
\ n
\]

%dsp-level:0
\[
\le 3,
\]

%dsp-level:0
\[
\ \mathbb C^3\ [\mathrm{CT0.11(a)}];
\]

%dsp-level:0
\[
\ \text{rest}:\ \mathbf{D},
\]

%dsp-level:0
\[
\mathbf{E},
\]

%dsp-level:0
\[
\mathbf{F},
\]

%dsp-level:0
\[
\mathbf{G}\quad\text{--- all}\ \mathfrak h\text{-free}
\]

%dsp-level:0
\[
\mathfrak h\mapsto\lambda\mathfrak h:\ \ Q\mapsto\lambda^{k_Q}Q;
\]

%dsp-level:0
\[
\quad\mathrm{P2}\ \text{invariant}\implies\text{rescaling}
\]

%dsp-level:0
\[
= \text{symmetry of the corpus}\implies\mathfrak h
\]

%dsp-level:0
\[
= \text{genuine single d.o.f.}\ (\neg\text{internal redundancy})
\]

%dsp-level:0
\[
\mathrm{P4}:\ (g,g',h_0,\Lambda,\varkappa,\beta,\alpha_\Omega)=\text{solution data},\ \neg\text{theory constants};
\]

%dsp-level:0
\[
\qquad \text{predictions}=\text{dimensionless ratios at fixed P4 data}.\ \blacksquare
\]

\textbf{Phys.} The entire corpus has exactly one dial: the hop length. Turning it rescales every dimensionful quantity uniformly and leaves every pure-number prediction untouched. Comparing to external measurement means fixing $\mathfrak h$ once — an act that tests A9, never an input to A1–A8.

\end{theorem}

\phantomsection
\section[Closure exports of Part 2 §C (consumed by Part 2 §D–§G,…]{\texorpdfstring{Closure exports of Part 2 §C (consumed by Part 2 §D–§G, Parts 3–4)}{Closure exports of Part 2 SC (consumed by Part 2 SD-SG, Parts 3-4)}}

%dsp-level:0
\[
\begin{array}{lll}
\mathrm{CT1.12}\ g_{\mu\nu}=\Omega^2\bar g_{\mu\nu},\ \operatorname{sig}=(-,+,+,+) &
\mathrm{CT1.13\text{–}17}\ \text{conformal laws} &
\mathrm{CT1.18}\ \bar G^{(1)}=-\tfrac12\Box\hat h^{(4)}\ (\text{TT})\\
\mathrm{CT2.1}\ [\cdot]=\mathfrak q\text{-powers} &
\mathrm{CT2.2\text{–}2.4}\ \mathcal P_M^{(n)}\ \text{projection};\ \mathbf F_M^{(n)}=\bigoplus_{\mathcal A(n)}F_A &
\mathrm{CT2.5}\ Z(\mathcal R,n)\ \text{poles}=\text{spectrum}\\
\mathrm{CT2.6}\ M\ \text{smooth 4-manifold} &
\mathrm{CT2.7}/\mathrm{CT1.21}\ S_{\mathrm{geom}}=\tfrac{1}{2\kappa^2}\int\sqrt{-g}(R-2\Lambda) &
\mathrm{CT2.8}\ \text{Hilbert space};\ \langle A|A\rangle\ge0\\
\mathrm{CT1.19}\ \sum ds^2\to\int g\,dx\,dx &
\mathrm{CT1.20}\ S_{\mathrm{eff}}=-i\log Z &
\mathrm{CT1.22}\ T_{\mu\nu}=T^{\mathrm{bd}}+T^{\mathrm{tr}}\\
\mathrm{CT1.23}\ G_{\mu\nu}+\Lambda g_{\mu\nu}=\kappa^2T_{\mu\nu} &
\mathrm{CT1.24}/\mathrm{CT1.26}\ \nabla^\mu T_{\mu\nu}=0\ (\times2\ \text{routes}) &
\mathrm{CT1.25}\ w_{D^+}=-1\\
\mathrm{CT1.27}\ |\Psi\rangle=\sum_CA(C)|g(C)\rangle &
\mathrm{CT1.28}\ \kappa_{\mathrm{ind}}\le f(16,\mathfrak h) &
\mathrm{CT1.29}\ S_{\partial\mathcal B}=\tfrac{A}{\mathfrak h^2}\ln2\\
\mathrm{CT2.9}\ \mathfrak h^2=4G_N\ln2;\ \kappa^2=\tfrac{2\pi\mathfrak h^2}{\ln2} &
\mathrm{CT2.10}\ \operatorname{Str}[X,Y\}=0 &
\mathrm{CT2.11}\ \exists!\ \text{free scale}\ \mathfrak h
\end{array}
\]

$\blacksquare$

\phantomsection
\section[§D — Electroweak Sector]{\texorpdfstring{§D — Electroweak Sector}{SD - Electroweak Sector}}

\textit{Dep.: Part 1; Part 2 §C. EW.1–EW.4, EW.6, EW.19 are restatement stubs of prior tags (cited nowhere downstream); forward pointer: EW.36$\to$EM.9 (restated at use); EW.8$\to$EW.32 is intra-section. Notation: carrier images $W_i\mapsto\tfrac12\sigma_i$, $Y_{\mathrm{car}}=\tfrac12\sigma_0$ (EW.32); $Y$ — electromagnetic generator with three realizations ($Y_{\mathrm{asp}},Y_{\mathrm{slot}},Y_{\mathrm{car}}$). Machine: \texttt{checks/\allowbreak{}verify\_\allowbreak{}D.py}; EW.5a/EW.36a: \texttt{checks/\allowbreak{}verify\_\allowbreak{}merge\_\allowbreak{}tierA.py}.}

\phantomsection
\section[§D.I — Operator Preliminaries]{\texorpdfstring{§D.I — Operator Preliminaries}{SD.I - Operator Preliminaries}}

\begin{theorem}[{\texorpdfstring{EW.1 — Joint $(S,T)$ sector decomposition}{EW.1 - Joint sector decomposition}}]\label{tag:EW-1}
\textbf{Thm.} $\equiv$ CT1.1: $\ \mathfrak H(\mathbb C^4)=\mathcal E_{++}\oplus\mathcal E_{+-}\oplus\mathcal E_{-+}\oplus\mathcal E_{--}$, $\ \dim=(7,3,3,3)$.

\textbf{Req.} CT1.1 (Part 1).

\textbf{Pf.} Identity [CT1.1]. $\blacksquare$

\textbf{Phys.} The arena of the electroweak sector: the weak triplet will live in the $T$-odd symmetric block $\mathcal E_{+-}$ of the pair operator space.

\end{theorem}

\begin{theorem}[{\texorpdfstring{EW.2 — Generator basis and sector assignment}{EW.2 - Generator basis and sector assignment}}]\label{tag:EW-2}
\textbf{Thm.} $\equiv$ CT1.2: $\ \{E_0,\Sigma,\hat E^{ij}\}\subset\mathcal E_{++}\ (1,1,5);\ \ W_i\in\mathcal E_{+-}\ (3);\ \ J_i\in\mathcal E_{-+}\ (3);\ \ K_i\in\mathcal E_{--}\ (3);\ \ 1{+}1{+}5{+}3{+}3{+}3=16$.

\textbf{Req.} CT1.2 (Part 1).

\textbf{Pf.} Identity [CT1.2]. $\blacksquare$

\textbf{Phys.} Fixes the named generators ($W_i$ total spin, $J_i,K_i$ Janus seals, $\Sigma$ exchange scalar) used throughout this document.

\end{theorem}

\begin{theorem}[{\texorpdfstring{EW.3 — The $W_i$ close $\mathfrak{su}(2)$}{EW.3 - The close}}]\label{tag:EW-3}
\textbf{Thm.} $\equiv$ CT1.4: $\ [W_i,W_j]=i\,\epsilon_{ijk}W_k$; $\ \operatorname{span}_{\mathbb R}\{W_i\}\cong\mathfrak{su}(2)$ [C2].

\textbf{Req.} CT1.4 (Part 1).

\textbf{Pf.} Identity [CT1.4]. $\blacksquare$

\textbf{Phys.} The weak-isospin algebra is the total-spin algebra of the qubit pair — present from rung 1 (CT0.13), not installed at rung 2.

\end{theorem}

\begin{theorem}[{\texorpdfstring{EW.4 — Central, trace, and mixed commutators}{EW.4 - Central, trace, and mixed commutators}}]\label{tag:EW-4}
\textbf{Thm.} $\equiv$ CT1.3, CT1.5, CT1.6, CT1.7:

%dsp-level:0
\[
(1)\ [E_0,X]=0\ \ \forall X;\qquad
(2)\ [\Sigma,W_i]=0;\qquad
(3)\ [\Sigma,J_i]=4iK_i,\ \ [\Sigma,K_i]=-4iJ_i;
\]

%dsp-level:0
\[
(4)\ [J_i,J_j]=[K_i,K_j]=i\,\epsilon_{ijk}W_k;\qquad
(5)\ [J_i,K_j]=i\bigl(\delta_{ij}\Sigma-E^{ij}\bigr).
\]

\textbf{Req.} CT1.3, CT1.5, CT1.6, CT1.7 (Part 1).

\textbf{Pf.} Identities [CT1.3, CT1.5, CT1.6, CT1.7]. $\blacksquare$

\textbf{Phys.} The commutator backbone of the sector: the scalar $\Sigma$ commutes with the weak triplet but rotates the two Janus sectors into each other; seals close onto the weak triplet; mixed seals close into the metric block.

\end{theorem}

\begin{theorem}[{\texorpdfstring{EW.5 — Second-order Higgs-to-weak pathways}{EW.5 - Second-order Higgs-to-weak pathways}}]\label{tag:EW-5}
\textbf{Thm.}

%dsp-level:0
\[
\bigl[K_j,[\Sigma,J_i]\bigr]=4\,\epsilon_{ijk}W_k;\qquad
\bigl[J_j,[\Sigma,K_i]\bigr]=-4\,\epsilon_{ijk}W_k;
\]

%dsp-level:0
\[
[J_j,K_i]\in\mathcal E_{++}\ \wedge\ \mathcal E_{++}\cap\operatorname{span}\{W_k\}=\{0\}\implies\text{the}\ [J,K]\ \text{pathway}\ \neg\text{reaches the weak sector}.
\]

\textbf{Req.} EW.4 (= CT1.5, CT1.6, CT1.7); CT1.1 (direct sum), CT1.2.

\textbf{Pf.}

%dsp-level:0
\[
[\Sigma,J_i]=4iK_i\ \implies\ \bigl[K_j,[\Sigma,J_i]\bigr]=4i\,[K_j,K_i]=4i\cdot i\,\epsilon_{jik}W_k=-4\,\epsilon_{jik}W_k=4\,\epsilon_{ijk}W_k
\]

%dsp-level:0
\[
\quad [\mathrm{EW.4}(3),(4)]
\]

%dsp-level:0
\[
[\Sigma,K_i]=-4iJ_i\ \implies\ \bigl[J_j,[\Sigma,K_i]\bigr]=-4i\,[J_j,J_i]=-4i\cdot i\,\epsilon_{jik}W_k=4\,\epsilon_{jik}W_k=-4\,\epsilon_{ijk}W_k
\]

%dsp-level:0
\[
[J_j,K_i]=i\bigl(\delta_{ji}\Sigma-E^{ji}\bigr)\in\mathcal E_{++}\ [\mathrm{EW.4}(5),\ \mathrm{EW.2}];
\]

%dsp-level:0
\[
\qquad W_k\in\mathcal E_{+-};\quad\text{direct sum}\ [\mathrm{EW.1}]\implies\mathcal E_{++}\cap\operatorname{span}\{W_k\}=\{0\}.\ \blacksquare
\]

\textbf{Phys.} Although the scalar commutes with the weak triplet, it reaches it at second adjoint order — and only through the Janus seals. The scalar’s unique algebraic handle on the weak sector runs through crystallization/dissolution: the algebraic shadow of the wrapping mechanism that actually generates mass in Part IV.

\end{theorem}

\begin{theorem}[{\texorpdfstring{EW.5a — Convention exactness and the rigidity of the named basis}{EW.5a - Convention exactness and the rigidity of the named basis}}]\label{tag:EW-5a}
\textbf{Thm.} In the Pauli basis of $\mathfrak H(\mathbb C^4)$ [A3, CT1.1], with CT1.2's definitions verbatim ($E_0=\sigma_0\otimes\sigma_0$, $\Sigma=\sum_i\sigma_i\otimes\sigma_i$, $W_i=\tfrac12(\sigma_0\otimes\sigma_i+\sigma_i\otimes\sigma_0)$, $J_i=\tfrac12(\sigma_j\otimes\sigma_k-\sigma_k\otimes\sigma_j)$ cyclic, $K_i=\tfrac12(\sigma_0\otimes\sigma_i-\sigma_i\otimes\sigma_0)$, $E^{ij}=\tfrac12(\sigma_i\otimes\sigma_j+\sigma_j\otimes\sigma_i)$), as \textbf{exact matrix equalities over $\mathbb Q(i)$} (no tolerance, no floats): $(1)$ \textbf{The frozen identities hold exactly as stated.} CT1.3 (16 brackets), CT1.4 (6 ordered pairs), CT1.5 (9), CT1.6 with $+i\epsilon_{ijk}W_k$ (12; no sign residue), CT1.7 (9, each with all three $W$-components exactly zero), CT1.8 (12); \textbf{EW.5 both clauses at all nine $(i,j)$} (diagonal $i=j$: both sides exactly $0$); CT0.7a; CT0.7(1) (incl.\ $\mathfrak T\Pi_{\mathcal R}=3$, $\mathfrak T\Pi_{\mathcal V}=1$, ranks $3/1$); the associative relation $\Sigma^2=3E_0-2\Sigma$; $[S,T]=0$ on all 16 basis elements with $\eta=(-,+,+,+)$ [A3]; the full $(S,T)$ parity table of CT1.2 (17 generators). $(2)$ \textbf{The convention group is exactly $\mathbb Z_2$.} The group of diagonal sign redefinitions $(J,K,W,\Sigma,E)\mapsto(\lambda_JJ,\lambda_KK,\lambda_WW,\lambda_\Sigma\Sigma,\lambda_EE)$, $\lambda\in\{\pm1\}^5$, with $E_0=I_4$ fixed (the unit), preserving every identity of clause (1) — including $\Sigma^2=3E_0-2\Sigma$ and $\mathbb S=\tfrac12(E_0+\Sigma)$ — is exactly

%dsp-level:0
\[
\{\mathrm{id},\ \operatorname{Ad}(\mathbb S)\}\cong\mathbb Z_2,\qquad \operatorname{Ad}(\mathbb S):\ (J_i,K_i)\mapsto(-J_i,-K_i),\quad W_i,\Sigma,E^{ij},E_0\ \text{fixed}
\]

(exactly $2$ of the $32$ sign maps survive; exhaustive enumeration). $(3)$ \textbf{EW.5 is invariant under the nontrivial element} (each clause maps to itself: $J,K$ flip together). The naming exchange $J\leftrightarrow K$ maps clause 1 of EW.5 to clause 2 with the wrong sign and is \textit{excluded} by the frozen identity set; Hermiticity excludes continuous phases ($e^{i\varphi}K_i$ Hermitian $\Rightarrow\varphi\in\{0,\pi\}$ — the sign enumeration is exhaustive).

\textbf{Req.} A3; C1, C2, C3; CT0.7, CT0.7a, CT1.1–CT1.8 (Part 1); EW.4, EW.5; exact $\mathbb Q(i)$ matrix arithmetic, exhaustive $2^5$ sign enumeration [math; machine: \texttt{checks/\allowbreak{}verify\_\allowbreak{}merge\_\allowbreak{}tierA.py} §B].

\textbf{Pf.}

%dsp-level:0
\[
(1):\ \text{each identity evaluated entrywise over}\ \mathbb Q(i)\ \text{from CT1.2's definitions; zero residual}\ [\text{machine}]
\]

%dsp-level:0
\[
(2):\ (\lambda_\Sigma\Sigma)^2=3E_0-2\lambda_\Sigma\Sigma\ \wedge\ \lambda_\Sigma^2=1\ \wedge\ \Sigma^2=3E_0-2\Sigma\ \wedge\ \Sigma\ne0\implies\lambda_\Sigma=1
\]

%dsp-level:0
\[
\mathrm{CT1.5}:\ \lambda_\Sigma\lambda_J=\lambda_K\ \wedge\ \lambda_\Sigma\lambda_K=\lambda_J\implies\lambda_J=\lambda_K;\qquad \mathrm{CT1.7}:\ \lambda_J\lambda_K=\lambda_\Sigma=1\ \wedge\ \lambda_J\lambda_K=\lambda_E\implies\lambda_E=1
\]

%dsp-level:0
\[
\mathrm{CT1.6}:\ \lambda_J^2=\lambda_K^2=\lambda_W;
\]

%dsp-level:0
\[
\qquad \mathrm{CT1.4}:\ \lambda_W^2=\lambda_W\implies\lambda_W=1;
\]

%dsp-level:0
\[
\qquad \implies(\lambda_J,\lambda_K,\lambda_W,\lambda_\Sigma,\lambda_E)\in\{(1,1,1,1,1),(-1,-1,1,1,1)\}
\]

%dsp-level:0
\[
(-1,-1,1,1,1)=\operatorname{Ad}(\mathbb S):\ \mathbb S J_i\mathbb S=-J_i,\ \mathbb S K_i\mathbb S=-K_i\ [\mathrm{CT0.7}(3):\ J,K\ S\text{-odd}];
\]

%dsp-level:0
\[
\quad \mathbb S W_i\mathbb S=W_i,\ \mathbb S\Sigma\mathbb S=\Sigma\ [S\text{-even}]
\]

%dsp-level:0
\[
(3):\ [K_j,[\Sigma,J_i]]\mapsto[-K_j,[\Sigma,-J_i]]=[K_j,[\Sigma,J_i]]\ \checkmark;
\]

%dsp-level:0
\[
\qquad J\leftrightarrow K:\ [J_j,[\Sigma,K_i]]=+4\epsilon_{ijk}W_k\ne-4\epsilon_{ijk}W_k\implies\text{excluded}
\]

%dsp-level:0
\[
\text{Hermiticity}:\ (e^{i\varphi}K_i)^\dagger=e^{-i\varphi}K_i=e^{i\varphi}K_i\iff e^{2i\varphi}=1\iff\varphi\in\{0,\pi\}.\ \blacksquare
\]

\textbf{Phys.} The electroweak backbone was accused of a sign slip; the audit finds the opposite — a lock. Once the unit, the associative square of the exchange scalar, and the swap definition are held fixed, the only freedom left in the whole named basis is the simultaneous flip of $J$ and $K$, under which every frozen identity — EW.5 included — is exactly invariant. There is no phase convention left to fix; the basis is rigid up to the one declared $\mathbb Z_2$.

\end{theorem}

\phantomsection
\section[§D.II — Rung Masks and the Electroweak Algebra]{\texorpdfstring{§D.II — Rung Masks and the Electroweak Algebra}{SD.II - Rung Masks and the Electroweak Algebra}}

\begin{theorem}[{\texorpdfstring{EW.6 — Rung masks and active aspect sets}{EW.6 - Rung masks and active aspect sets}}]\label{tag:EW-6}
\textbf{Thm.} $\equiv$ CT2.2, CT2.3, CT0.13:

%dsp-level:0
\[
\bigl(\mathcal P_M^{(n)}\bigr)^2=\mathcal P_M^{(n)}=\bigl(\mathcal P_M^{(n)}\bigr)^\dagger;\qquad
\mathbf F_M^{(n)}=\mathcal P_M^{(n)}\mathbf F\,\mathcal P_M^{(n)};
\]

%dsp-level:0
\[
\mathcal A(0)=\{J\};
\]

%dsp-level:0
\[
\quad \mathcal A(1)=\{J,G,W,H,D^+\};
\]

%dsp-level:0
\[
\quad \mathcal A(2)=\mathcal A(1)\cup\{\mathrm{EM}\};
\]

%dsp-level:0
\[
\quad \mathcal A(3)=\mathcal A(2)\cup\{\mathrm{strong}\};
\]

%dsp-level:0
\[
\quad \mathcal A(n\ge4)=\mathcal A(3)\cup\{D^-\}.
\]

$W=\operatorname{span}\{W_1,W_2,W_3\}$: neutral seed $W_3$, $\neg$ independent aspect [C8]; scalar seed $=$ trace mode $h$ of D2.1.

\textbf{Req.} CT0.13, CT2.2, CT2.3 (\textbf{B}, \textbf{C}); C8; D2.1, D3.3.

\textbf{Pf.} Identities [CT2.2, CT2.3, CT0.13 — derived there from CT0.7, D3.3, D3.1, D3.2]. $\blacksquare$

\textbf{Phys.} Which forces exist at which rung is bookkeeping, already fixed in B: the electroweak content (weak triplet + Higgs scalar + dark + EM centre) is available from rungs 1–2; colour joins only at rung 3.

\end{theorem}

\begin{theorem}[{\texorpdfstring{EW.7 — The rung-2 electroweak mask}{EW.7 - The rung-2 electroweak mask}}]\label{tag:EW-7}
\textbf{Thm.}

%dsp-level:0
\[
\mathcal A(2)=\{J,\ G,\ W_1,\ W_2,\ W_3,\ H,\ D^+,\ \mathrm{EM}\};\qquad
\mathbf F_M^{(2)}=\mathcal P_M^{(2)}\mathbf F\,\mathcal P_M^{(2)}=\bigoplus_{A\in\mathcal A(2)}F_A\big|_{\mathcal K^{(A)}} .
\]

\textbf{Req.} EW.6; CT2.4 (Part 2).

\textbf{Pf.}

%dsp-level:0
\[
\mathrm{EW.6}:\ \mathcal A(2)=\mathcal A(1)\cup\{\mathrm{EM}\},\ \ \mathcal A(1)=\{J,G,W,H,D^+\},\ W=\operatorname{span}\{W_1,W_2,W_3\}
\]

%dsp-level:0
\[
\mathrm{CT2.4}:\ \text{mask retains exactly the blocks of}\ \mathcal A(2),\ \text{annihilates the rest}.\ \blacksquare
\]

\textbf{Phys.} The rung-2 world carries gravity, weak, Higgs, uniform dark and electromagnetism — and no colour: the electroweak sector is exactly the mask content at rung 2.

\end{theorem}

\begin{theorem}[{\texorpdfstring{EW.8 — The electromagnetic generator}{EW.8 - The electromagnetic generator}}]\label{tag:EW-8}
\textbf{Thm.} The electromagnetic generator $Y$ is one abstract generator with three realizations:

%dsp-level:0
\[
Y_{\mathrm{slot}}:=Z\bigl(\iota_2(\mathfrak u(2))\bigr)\ (1\text{-dim})\ [\mathrm{D2.12},\mathrm{CT0.13}];
\]

%dsp-level:0
\[
\qquad Y_{\mathrm{asp}}:=\mathcal P^{(\mathrm{EM})}:\ \mathcal K\to\mathcal K^{(\mathrm{EM})},\ \ \operatorname{rk}=1,\ \mathfrak T(Y_{\mathrm{asp}})=1\ [\mathrm{D2.3},\mathrm{CT0.14}];
\]

%dsp-level:0
\[
Y_{\mathrm{car}}:=\tfrac12\sigma_0\ \ \text{on}\ \mathcal H_{\mathsf v}\cong\mathbb C^2\ [\mathrm{EW.32}];
\]

%dsp-level:0
\[
[Y_{\mathrm{asp}},\mathbf F_M^{(n)}]=0\ (\text{block-diagonal});\qquad
[Y_{\mathrm{slot}},\mathfrak{su}(2)]=0;\qquad
[Y_{\mathrm{car}},W_i|_{\mathrm{car}}]=0 .
\]

\textbf{Rmk (separation).} $\ \mathcal P_M^{(2)}E_0\mathcal P_M^{(2)}=\mathcal P_M^{(2)}$, $\mathfrak T=17$ [CT2.2, CT0.14]: a mask fact about the \textit{pair} centre $E_0\leftrightarrow F_{D^+}$ [CT0.12] — $\neg$ a definition of $Y$ (the prior codex defined $Y$ this way; repaired).

\textbf{Req.} A8; D2.3, D2.12; CT0.12, CT0.13, CT0.14, CT1.2 (Part 1); CT2.2 (Part 2).

\textbf{Pf.}

%dsp-level:0
\[
\mathrm{CT0.13}:\ F_{\mathrm{EM}}=\iota_2\bigl(Z(\mathfrak u(2))\bigr),\ 1\text{-dim}\ [\mathrm{D2.12}];
\]

%dsp-level:0
\[
\qquad \mathrm{CT0.14}:\ \dim\mathcal K^{(\mathrm{EM})}=1\implies Y_{\mathrm{asp}}=\mathcal P^{(\mathrm{EM})}:\ \operatorname{rk}1,\ \mathfrak T=1\ [\mathrm{D2.3}]
\]

%dsp-level:0
\[
\mathrm{D2.3}:\ \mathcal K=\bigoplus_A\mathcal K^{(A)}\ \text{orthogonal}\implies[Y_{\mathrm{asp}},\ \text{block-diagonal operators}]=0\implies[Y_{\mathrm{asp}},\mathbf F_M^{(n)}]=0
\]

%dsp-level:0
\[
Z(\mathfrak u(2))\ \text{commutes with the}\ \mathfrak{su}(2)\ \text{factor}\implies[Y_{\mathrm{slot}},\mathfrak{su}(2)]=0;
\]

%dsp-level:0
\[
\qquad [\tfrac12\sigma_0,\tfrac12\sigma_i]=0\ [\text{machine}:\ \text{verify\_D}]
\]

%dsp-level:0
\[
E_0=I_4\ [\mathrm{CT1.2}]\implies\mathcal P_M^{(2)}E_0\mathcal P_M^{(2)}=\mathcal P_M^{(2)}\ [\mathrm{CT2.2}];
\]

%dsp-level:0
\[
\quad E_0\leftrightarrow F_{D^+}\ [\mathrm{CT0.12}]:\ \text{pair-centre fact, kept separate from}\ Y.\ \blacksquare
\]

\textbf{Rmk (forward pointer, non-load-bearing).} The carrier realization/transport locus is EW.32 (EW.32 cites EW.8 backward).

\textbf{Phys.} Hypercharge is one generator with three coordinated faces: the centre of the rung-2 slot algebra, the rank-one projection onto the EM block, and $\tfrac12\sigma_0$ on the transported seed. Keeping the three apart — and keeping the \textit{pair} centre $E_0$ (which belongs to $D^+$) apart from all of them — is what makes the EM.17 trace formula and the EW.32 charge statements each well-typed.

\end{theorem}

\begin{theorem}[{\texorpdfstring{EW.9 — The abelian algebra}{EW.9 - The abelian algebra}}]\label{tag:EW-9}
\textbf{Thm.} $\ \bigl(\operatorname{span}_{\mathbb R}\{Y\},[\cdot,\cdot]\bigr)\cong\mathfrak u(1):\quad[aY,bY]=ab[Y,Y]=0;\quad\dim=1.$

\textbf{Req.} EW.8; Lie-algebra axioms [math].

\textbf{Pf.}

%dsp-level:0
\[
\dim\operatorname{span}\{Y\}=1\implies[aY,bY]=ab[Y,Y]=0;\qquad\text{Jacobi}\equiv0;\qquad
1\text{-dim abelian}=\mathfrak u(1).\ \blacksquare
\]

\textbf{Phys.} One generator, vanishing bracket: the abelian factor of the electroweak algebra.

\end{theorem}

\begin{theorem}[{\texorpdfstring{EW.10 — The rung-2 electroweak algebra}{EW.10 - The rung-2 electroweak algebra}}]\label{tag:EW-10}
\textbf{Thm.}

%dsp-level:0
\[
\mathfrak g_{\mathrm{EW}}:=\operatorname{span}_{\mathbb R}\{W_1,W_2,W_3,Y\}=\mathfrak{su}(2)\oplus\mathfrak u(1);
\]

%dsp-level:0
\[
\qquad [W_i,W_j]=i\,\epsilon_{ijk}W_k;\quad[W_i,Y]=0;\quad[Y,Y]=0;\quad\dim\mathfrak g_{\mathrm{EW}}=4 .
\]

\textbf{Req.} EW.3, EW.8, EW.9; direct-sum criterion (two commuting ideals, trivial intersection) [math].

\textbf{Pf.}

%dsp-level:0
\[
\operatorname{span}\{W_i\}\cong\mathfrak{su}(2)\ [\mathrm{EW.3}];\qquad\operatorname{span}\{Y\}\cong\mathfrak u(1)\ [\mathrm{EW.9}]
\]

%dsp-level:0
\[
[Y,W_i]=0\ [\mathrm{EW.8}]\implies\text{each factor an ideal of the sum}
\]

%dsp-level:0
\[
Y\in Z(\mathfrak g_{\mathrm{EW}});\quad\text{no nonzero}\ \textstyle\sum c_iW_i\ \text{central}\ [\mathrm{EW.3}]\implies\operatorname{span}\{W\}\cap\operatorname{span}\{Y\}=\{0\}
\]

%dsp-level:0
\[
\implies\mathfrak g_{\mathrm{EW}}=\mathfrak{su}(2)\oplus\mathfrak u(1);\qquad\dim=3+1=4.\ \blacksquare
\]

\textbf{Phys.} The electroweak gauge algebra of the corpus — $\mathfrak{su}(2)\oplus\mathfrak u(1)$, dimension four — is derived as the algebra generated by the rung-2 mask content. Nothing resembling a gauge principle was assumed.

\end{theorem}

\begin{theorem}[{\texorpdfstring{EW.11 — The mechanism is projection, not path integration}{EW.11 - The mechanism is projection, not path integration}}]\label{tag:EW-11}
\textbf{Thm.} $\mathfrak g_{\mathrm{EW}}$ arises from $\bigl(\mathcal P_M^{(2)}\mathbf F\mathcal P_M^{(2)},\ [\cdot,\cdot]\bigr)$ (finite-dimensional restriction); $\ Z(\mathcal R)=\sum_{C}\sum_\sigma w(C)$ (countable sum, A8); continuum $\to$ only after coarse-graining (CT2.6) $\implies\mathfrak g_{\mathrm{EW}}$ fixed before that limit; $\neg\exists$ functional integral over gauge fields in the derivation.

\textbf{Req.} A8; CT2.3, CT2.6 (Part 2); EW.10; CT1.4 (Part 1).

\textbf{Pf.}

%dsp-level:0
\[
\mathrm{CT2.3}:\ \mathbf F_M^{(2)}=\mathcal P_M^{(2)}\mathbf F\mathcal P_M^{(2)}=\text{operator restriction}\ (\neg\text{integral})
\]

%dsp-level:0
\[
\mathrm{EW.10}\Leftarrow[\mathrm{CT1.4},\ \mathrm{EW.8}]=\text{finite-dimensional matrix identities}
\]

%dsp-level:0
\[
\mathrm{A8}:\ Z(\mathcal R)=\textstyle\sum_{C\in\mathcal R}\sum_\sigma w(C)=\text{sum over a countable set};
\]

%dsp-level:0
\[
\qquad \mathrm{CT2.6}:\ \text{continuum appears only at coarse graining}\implies\text{algebra fixed pre-limit}.\ \blacksquare
\]

\textbf{Phys.} The gauge structure is produced by masking a discrete, finite algebra — the continuum functional description is a derived approximation, never the defining mechanism. “Gauge theory” in this corpus names a spectral property of a masked operator sum.

\end{theorem}

\phantomsection
\section[§D.III — The Scalar Trace Mode]{\texorpdfstring{§D.III — The Scalar Trace Mode}{SD.III - The Scalar Trace Mode}}

\begin{theorem}[{\texorpdfstring{EW.12 — Metric expansion and the scalar trace mode}{EW.12 - Metric expansion and the scalar trace mode}}]\label{tag:EW-12}
\textbf{Thm.} With D2.1 data:

%dsp-level:0
\[
\delta^{ij}\hat h_{ij}=0;\qquad q_{ij}=\hat h_{ij}+\delta_{ij}h;\qquad h=\tfrac13\Sigma_{\mathrm{tr}};\qquad
\delta_{ij}h\,E^{ij}=h\sum_{i=1}^3E^{ii}=h\,\Sigma .
\]

\textbf{Req.} D2.1; C5; CT1.12 (Part 2), EW.2 (= CT1.2).

\textbf{Pf.}

%dsp-level:0
\[
\delta^{ij}\hat h_{ij}=\Sigma_{\mathrm{tr}}-\tfrac13\cdot3\cdot\Sigma_{\mathrm{tr}}=0\quad[\mathrm{D2.1},\ \mathrm{CT1.12}]
\]

%dsp-level:0
\[
h=\tfrac13\Sigma_{\mathrm{tr}}\ \text{substituted into}\ \hat h_{ij}=q_{ij}-\tfrac13\delta_{ij}\Sigma_{\mathrm{tr}}\implies q_{ij}=\hat h_{ij}+\delta_{ij}h
\]

%dsp-level:0
\[
E^{ii}=\tfrac12\bigl(\sigma_i\otimes\sigma_i+\sigma_i\otimes\sigma_i\bigr)=\sigma_i\otimes\sigma_i\implies\sum_iE^{ii}=\Sigma\quad[\mathrm{EW.2}].\ \blacksquare
\]

\textbf{Phys.} The trace part of the real sector is the operator $h\Sigma$: a scalar field $h$ multiplying the exchange scalar. C5 keeps the bookkeeping honest — $\Sigma$ is an operator, $\Sigma_{\mathrm{tr}}$ a number.

\end{theorem}

\begin{theorem}[{\texorpdfstring{EW.13 — Identification of the scalar mode}{EW.13 - Identification of the scalar mode}}]\label{tag:EW-13}
\textbf{Thm.} The electroweak scalar mode is

%dsp-level:0
\[
H:=h\,\Sigma\qquad(\text{field}\ h;\ \ \text{operator representative}\ h\Sigma),
\]

the noncentral rotation scalar; the central scalar $E_0$ is unavailable (allocated to $F_{D^+}$).

\textbf{Req.} CT1.10, CT1.11, CT0.12 (Part 1), CT1.25 (Part 2); EW.1 (= CT1.1), EW.4.

\textbf{Pf.}

%dsp-level:0
\[
\dim\mathcal E_{++}=7\ [\mathrm{EW.1}];\qquad\mathcal E_{++}\to\mathbf1\oplus\mathbf5\oplus\mathbf1\ [\mathrm{CT1.10}]:\ \ \text{scalars}=\{E_0,\ \Sigma\}
\]

%dsp-level:0
\[
E_0\in Z;\quad[\Sigma,J_1]=4iK_1\ne0\implies\Sigma\notin Z;\quad E_0\nparallel\Sigma\quad[\mathrm{EW.4}(1),(3);\ \mathrm{CT1.11}]
\]

%dsp-level:0
\[
E_0\ \text{already allocated}:\ F_{D^+}=\Lambda\mathbb I\ [\mathrm{CT0.12}]\to\text{cosmological term}\ [\mathrm{CT1.25}];\quad\text{reallocating}=\text{double counting}
\]

%dsp-level:0
\[
\implies\Sigma=\text{unique remaining rotation-scalar direction}=F_H\ [\mathrm{CT0.12}];\qquad H=h\Sigma.\ \blacksquare
\]

\textbf{Phys.} The Higgs-type scalar of the corpus is the exchange scalar weighted by the trace mode: forced, not chosen — once the central scalar has become the uniform dark aspect, exactly one rotation-scalar direction remains, and A8 already names it $F_H$.

\end{theorem}

\begin{theorem}[{\texorpdfstring{EW.14 — The scalar background commutes with the electroweak generators}{EW.14 - The scalar background commutes with the electroweak generators}}]\label{tag:EW-14}
\textbf{Thm.} $\ \forall h_0\in\mathbb R:\quad[h_0\Sigma,\ W_i]=0;\qquad[h_0\Sigma,\ Y_{\mathrm{asp}}]=[h_0\Sigma,\ Y_{\mathrm{car}}]=0.$

\textbf{Req.} EW.4(2) (= CT1.5), EW.8.

\textbf{Pf.}

%dsp-level:0
\[
h_0\in\mathbb R\implies[h_0\Sigma,X]=h_0[\Sigma,X];\qquad[\Sigma,W_i]=0\ [\mathrm{EW.4}(2)]
\]

%dsp-level:0
\[
[Y_{\mathrm{asp}},F_H\text{-block}]=0\ [\mathrm{EW.8}:\ \text{block-diagonal}];\qquad
[Y_{\mathrm{car}},h_0\Sigma|_{\mathrm{car}}]=0\ [\tfrac12\sigma_0\ \text{central}].\ \blacksquare
\]

\textbf{Phys.} A scalar background is invisible to the electroweak algebra at the commutator level — the first half of the statement that this theory has no adjoint Higgs mechanism.

\end{theorem}

\begin{theorem}[{\texorpdfstring{EW.15 — No adjoint electroweak breaking}{EW.15 - No adjoint electroweak breaking}}]\label{tag:EW-15}
\textbf{Thm.}

%dsp-level:0
\[
\operatorname{Stab}_{\mathfrak g_{\mathrm{EW}}}(h_0\Sigma):=\{X\in\mathfrak g_{\mathrm{EW}}:\ [h_0\Sigma,X]=0\}=\mathfrak g_{\mathrm{EW}}\qquad(\text{full stabilizer});
\]

%dsp-level:0
\[
\implies\neg\exists\ \text{proper subalgebra selected by an adjoint stabilizer mechanism}.
\]

\textbf{Consequence.} Any electroweak reduction is a statement about the spectrum of the masked resolvent (Part IV: EW.33 transport), $\neg$ about Lie-algebra stabilizers.

\textbf{Req.} EW.10, EW.14.

\textbf{Pf.}

%dsp-level:0
\[
\mathrm{EW.14}:\ [h_0\Sigma,X]=0\ \ \forall X\in\{W_1,W_2,W_3,Y\}\implies\operatorname{Stab}=\mathfrak g_{\mathrm{EW}}\implies\neg\exists\ \text{proper reduction by commutators}.\ \blacksquare
\]

\textbf{Phys.} The scalar background breaks nothing adjointly: the stabilizer is the whole algebra. Mass generation must therefore be — and in Part IV is — a transport effect on a carried seed, not a commutator effect on generators.

\end{theorem}

\begin{theorem}[{\texorpdfstring{EW.16 — Scalar coupling to the antisymmetric sectors}{EW.16 - Scalar coupling to the antisymmetric sectors}}]\label{tag:EW-16}
\textbf{Thm.}

%dsp-level:0
\[
[\Sigma,J_i]=4iK_i;\qquad[\Sigma,K_i]=-4iJ_i;\qquad
\bigl[K_j,[\Sigma,J_i]\bigr]=4\,\epsilon_{ijk}W_k;\qquad
\bigl[J_j,[\Sigma,K_i]\bigr]=-4\,\epsilon_{ijk}W_k .
\]

\textbf{Req.} EW.4(3), EW.5.

\textbf{Pf.} First-order: [EW.4(3)]; second-order: [EW.5]. $\blacksquare$

\textbf{Rmk.} $J_i,K_i\ =$ threshold-moving generators [CT0.7]: the scalar’s only algebraic handle on the weak sector runs through crystallization/dissolution — the algebraic shadow of the wrapping mechanism of Part IV.

\textbf{Phys.} The scalar talks to the Janus seals at first order and, through them, to the weak triplet at second order. Every algebraic route from the Higgs mode to the weak force passes through the reality-changing generators.

\end{theorem}

\phantomsection
\section[§D.IV — Rung-3 Wrapping, the Seed, and Mass Generation]{\texorpdfstring{§D.IV — Rung-3 Wrapping, the Seed, and Mass Generation}{SD.IV - Rung-3 Wrapping, the Seed, and Mass Generation}}

\begin{theorem}[{\texorpdfstring{EW.17 — Rung-3 closed wrapping chains}{EW.17 - Rung-3 closed wrapping chains}}]\label{tag:EW-17}
\textbf{Thm.} $\gamma=$ chain [D1.4]; closed $\iff\#J^+=\#J^-\ \wedge\ \partial\gamma=0$; wrapping $=$ nonzero spatial winding around a closed ledger cycle [D1.7]:

%dsp-level:0
\[
\Theta(\gamma)=\sum_{e\in\gamma}\widehat{ds}^{\,2}(e);
\]

%dsp-level:0
\[
\qquad \Phi(\gamma)=\sum_{e\in\gamma}\alpha_e\,\widehat{ds}^{\,2}(e)\ \ [\mathrm{D1.7};\ \text{uniform-phase case:}\ \Phi=\alpha\,\Theta];
\]

%dsp-level:0
\[
\qquad \operatorname{Hol}_{\mathrm{spiral}}(\gamma)=e^{i\Phi(\gamma)} .
\]

%dsp-level:0
\[
\max\{n:\ \exists\ \text{closed wrapping chain at rung}\ n\}=3 .
\]

\textbf{Req.} D1.4, D1.7; CT0.11 (Part 1).

\textbf{Pf.}

%dsp-level:0
\[
\mathrm{CT0.11}:\ \text{accessible rungs}=\{0,1,2,3\}
\]

%dsp-level:0
\[
\text{closed wrapping}\Rightarrow\exists v:\ \text{chain continues}\wedge\text{branches back onto itself}\Rightarrow\bigl|\{e\in E_{\mathcal R}:e\ni v\}\bigr|\ge3\Rightarrow n(v)\ge3
\]

%dsp-level:0
\[
\implies n=3.\ \blacksquare
\]

\textbf{Phys.} Rung 3 is the first — and by the capacity bound the last — rung at which closed wrapping chains exist. Mass generation is therefore a rung-3 phenomenon, which is also why it coexists with colour (activated at the same rung).

\end{theorem}

\begin{theorem}[{\texorpdfstring{EW.18 — Spiral pole condition}{EW.18 - Spiral pole condition}}]\label{tag:EW-18}
\textbf{Thm.} Closed wrapping chain materializes as a rung-3 pole $\implies$

%dsp-level:0
\[
\Phi(\gamma)\in2\pi\mathbb Z\qquad\iff\qquad\sum_{e\in\gamma}\alpha_e\,\widehat{ds}^{\,2}(e)=2\pi k,\ \ k\in\mathbb Z\qquad[\Phi:\ \mathrm{D1.7}].
\]

\textbf{Req.} EW.17; D2.4, CT2.5 (Part 2).

\textbf{Pf.}

%dsp-level:0
\[
\mathrm{EW.17}:\ \text{transport around the cycle}=\times\operatorname{Hol}_{\mathrm{spiral}}(\gamma)=e^{i\Phi(\gamma)}
\]

%dsp-level:0
\[
\mathrm{CT2.5}:\ \text{pole}=\text{eigenvalue condition on a single-valued eigenmode}\implies e^{i\Phi(\gamma)}=1
\]

%dsp-level:0
\[
\text{winding}\ k:\ \Phi\mapsto\Phi+2\pi k\implies\text{condition}=\Phi(\gamma)\in2\pi\mathbb Z.\ \blacksquare
\]

\textbf{Phys.} Materialization of a wrapping chain is quantized: only integer-winding spiral phases support resolvent poles — a single-valuedness (Bohr–Sommerfeld-type) condition derived from the ledger. With the per-edge phase $\alpha_e$ — A4's branch index, alphabet $\{0,\pi\}$ (read in E: Def-EM.20, EM.20) — the constraint reads on the antialigned content: $\Theta_{\mathcal P}(\gamma)=2k$ (CT1.19b).

\end{theorem}

\begin{theorem}[{\texorpdfstring{EW.19 — Mass as pole displacement}{EW.19 - Mass as pole displacement}}]\label{tag:EW-19}
\textbf{Thm.} $\equiv$ D2.5 (rung-3 context): massless $\iff$ pole at $k^2=0$; massive $\iff$ pole at $k^2=m^2>0$; $m^2:=$ pole displacement of $Z(\mathcal R,3)$ [D2.4].

\textbf{Req.} D2.4, D2.5; CT2.5 (Part 2).

\textbf{Pf.} Restatement [D2.5]. $\blacksquare$

\textbf{Phys.} In this corpus “mass” has one meaning only: how far a resolvent pole sits from the emergent light cone.

\end{theorem}

\begin{theorem}[{\texorpdfstring{EW.32 — The wrapping seed: carrier and normalization}{EW.32 - The wrapping seed: carrier and normalization}}]\label{tag:EW-32}
\textbf{Thm.} Carrier: $\ \mathcal H_{\mathsf v}\cong\mathbb C^2$ (single qubit). Generator images:

%dsp-level:0
\[
W_i\longmapsto\tfrac12\sigma_i;\qquad Y\longmapsto\tfrac12\sigma_0 ;
\]

%dsp-level:0
\[
\mathfrak T\bigl(\mathcal G_A\mathcal G_B\bigr)=\tfrac12\,\delta_{AB},\quad\mathcal G_A\in\{W_1,W_2,W_3,Y\}\qquad(\text{unique uniform trace normalization});
\]

%dsp-level:0
\[
T_a\,\mathsf v=0\qquad(a=1,\dots,8)\qquad(\text{colour neutrality}).
\]

\textbf{Req.} A1, A2; D1.4, D2.6, D2.12, P4; CT0.12 (Rmk), CT0.13, CT0.14, CT1.2 (Part 1); EW.8.

\textbf{Pf.}

%dsp-level:0
\[
\mathrm{A2}:\ \text{channel capacity}
\]

%dsp-level:0
\[
= 1\mathfrak q;
\]

%dsp-level:0
\[
\quad\mathrm{A1}:\ \text{carrier of one qubit}
\]

%dsp-level:0
\[
= \mathbb C^2;
\]

%dsp-level:0
\[
\quad\mathrm{D1.4}:\ \text{chain}
\]

%dsp-level:0
\[
= \text{sequence of channels}\implies\text{transported content}
\]

%dsp-level:0
\[
= \text{one qubit}:\ \mathcal H_{\mathsf v}\cong\mathbb C^2
\]

%dsp-level:0
\[
\mathrm{CT1.2}:\ W_i=\tfrac12\bigl(\sigma_i^{(p)}+\sigma_i^{(q)}\bigr);\quad\text{restriction to one transported factor}\implies W_i\mapsto\tfrac12\sigma_i
\]

%dsp-level:0
\[
\mathrm{CT0.13}/\mathrm{D2.12}:\ Y_{\mathrm{slot}}
\]

%dsp-level:0
\[
= Z\bigl(\iota_2(\mathfrak u(2))\bigr)\ [\mathrm{EW.8}]
\]

%dsp-level:0
\[
\implies Y_{\mathrm{car}}
\]

%dsp-level:0
\[
= y_0\sigma_0\ \ (y_0\in\mathbb R:\ \text{centre acts as scalar on}\ \mathbb C^2)
\]

%dsp-level:0
\[
\mathfrak T(W_iW_j)=\tfrac14\mathfrak T(\sigma_i\sigma_j)=\tfrac12\delta_{ij};
\]

%dsp-level:0
\[
\qquad \mathfrak T(W_iY)=\tfrac12y_0\,\mathfrak T(\sigma_i)=0;
\]

%dsp-level:0
\[
\qquad \mathfrak T(Y^2)=2y_0^2\overset!=\tfrac12\implies y_0=\tfrac12
\]

%dsp-level:0
\[
\text{positive root selected; negative root}=\omega\text{-image}\ [\mathrm{D2.6}]=\text{conjugate seed}
\]

%dsp-level:0
\[
\text{uniqueness}:\ y_0\ne\tfrac12\implies\neg\text{uniform normalization}\implies\text{relative ideal weight free}\ (\neg\mathrm{P4});
\]

%dsp-level:0
\[
\quad\text{remaining freedom}
\]

%dsp-level:0
\[
= (g,g')\ [\mathrm{P4}]
\]

%dsp-level:0
\[
\mathrm{CT0.12\ Rmk}:\ F_H
\]

%dsp-level:0
\[
= \text{pair block},
\]

%dsp-level:0
\[
\ F_{\mathrm{strong}}
\]

%dsp-level:0
\[
= \text{stack block on}\ \mathcal H^{(3)}_{\mathrm{slot}};
\]

%dsp-level:0
\[
\quad\mathrm{CT0.14}:\ \text{orthogonal summands of}\ \mathcal K
\]

%dsp-level:0
\[
\implies\text{seed from the scalar background carries}\ \neg\text{slot index}\implies T_a\mathsf v=0\ \ \forall a.\ \blacksquare
\]

\textbf{Phys.} The object transported around a wrapping chain is a single qubit carrying the fundamental of $\mathfrak u(2)$ in the unique uniformly trace-normalized representation — $W_i\to\sigma_i/2$, $Y\to\sigma_0/2$ — and no colour. Representation content is derived from the channel quantum A2; colour-neutrality is what keeps the gluons massless (EW.26).

\end{theorem}

\begin{theorem}[{\texorpdfstring{EW.33 — The wrapping mass form}{EW.33 - The wrapping mass form}}]\label{tag:EW-33}
\textbf{Thm.} With $\ \mathcal G=(gW_1,\ gW_2,\ gW_3,\ g'Y)$, $g,g'>0$:

%dsp-level:0
\[
M^2_{AB}=\mu_\star^2\;\operatorname{Re}\bigl\langle\mathcal G_A\mathsf v,\ \mathcal G_B\mathsf v\bigr\rangle,\qquad
\mu_\star^2=c_W\,\varkappa^2\beta^2h_0^2>0\quad(c_W>0:\ \text{rung-3 residue});
\]

%dsp-level:0
\[
M^2\ \text{real symmetric PSD (Gram)};\qquad
\ker M^2=\Bigl\{c\in\mathbb R^4:\ \Bigl(\textstyle\sum_Ac_A\mathcal G_A\Bigr)\mathsf v=0\Bigr\};\qquad
h_0=0\implies M^2\equiv0 .
\]

\textbf{Req.} EW.17, EW.32; D1.1; CT1.12, CT2.5, CT2.8 (Part 2); Gram-matrix facts [linear algebra].

\textbf{Pf.}

%dsp-level:0
\[
\mathrm{EW.17}:\ \text{transport}=\text{path-ordered product of edge transports}\ [\mathrm{D1.1}];
\]

%dsp-level:0
\[
\quad \text{leading order}:\ \ \mathcal U=\exp\Bigl(i\oint\textstyle\sum_Ac_A\mathcal G_A\,d\ell\Bigr)=1+i\,\ell_\gamma\textstyle\sum_Ac_A\mathcal G_A+O(c^2)
\]

%dsp-level:0
\[
\mathrm{CT2.5}:\ \text{pole displacement}=\text{quadratic single-valuedness defect}:
\]

%dsp-level:0
\[
\quad \bigl\|(\mathcal U-1)\mathsf v\bigr\|^2=\ell_\gamma^2\Bigl\|\textstyle\sum_Ac_A\mathcal G_A\mathsf v\Bigr\|^2+O(c^3)
\]

%dsp-level:0
\[
\Bigl\|\textstyle\sum_Ac_A\mathcal G_A\mathsf v\Bigr\|^2=\sum_{A,B}c_Ac_B\,\operatorname{Re}\langle\mathcal G_A\mathsf v,\mathcal G_B\mathsf v\rangle
\]

%dsp-level:0
\[
\qquad \bigl[c_Ac_B\ \text{sym};\ \operatorname{Im}\langle\cdot,\cdot\rangle\ \text{antisym}\implies\text{Im cancels}\bigr]
\]

%dsp-level:0
\[
\mathrm{CT1.12}:\ \bar g_{0i}
\]

%dsp-level:0
\[
= \varkappa N_i
\]

%dsp-level:0
\[
= \varkappa\beta w_i;
\]

%dsp-level:0
\[
\quad\bar g_{ij}\supset\varkappa h_0\delta_{ij}\implies\text{cycle amplitude}\propto\varkappa^2\beta^2h_0^2;
\]

%dsp-level:0
\[
\quad\ell_\gamma^2\ \text{absorbed into}\ c_W>0\implies\mu_\star^2
\]

%dsp-level:0
\[
= c_W\varkappa^2\beta^2h_0^2
\]

%dsp-level:0
\[
M^2=\mu_\star^2\cdot\text{Gram}\bigl(\{\mathcal G_A\mathsf v\},\operatorname{Re}\langle\cdot,\cdot\rangle\bigr)\implies\text{symmetric}\ \wedge\ \text{PSD}
\]

%dsp-level:0
\[
\quad [\text{linear algebra};\ \text{consistent with}\ \mathrm{CT2.8}:\ \text{residues}\ge0]
\]

%dsp-level:0
\[
c\in\ker M^2\iff c^{\mathsf T}M^2c=0\iff\Bigl\|\textstyle\sum_Ac_A\mathcal G_A\mathsf v\Bigr\|^2=0\iff\textstyle\sum_Ac_A\mathcal G_A\in\operatorname{Ann}(\mathsf v)
\]

%dsp-level:0
\[
h_0=0\implies\mu_\star^2=0\implies M^2\equiv0:\ \forall\ \text{electroweak modes massless}.\ \blacksquare
\]

\textbf{Phys.} Mass is a transport defect: the failure of a wrapped seed to return to itself, squared. It is automatically a positive-semidefinite Gram matrix, its kernel is the seed’s stabilizer, and it vanishes without the scalar background — the background is necessary for mass even though (EW.15) it breaks nothing by commutators. The overall scale $\mu_\star^2$ is the sector’s single free positive number (P4 configuration data); every ratio below is free of it.

\end{theorem}

\begin{theorem}[{\texorpdfstring{EW.20 — Weak mass generation}{EW.20 - Weak mass generation}}]\label{tag:EW-20}
\textbf{Thm.} $\langle h\rangle=h_0\ne0$; weak frame with $\mathsf v=$ lower weight state of $W_3$ (an $\mathrm{SU}(2)$ basis choice, always available):

%dsp-level:0
\[
W_3\mathsf v=-\tfrac12\mathsf v;\qquad Y\mathsf v=+\tfrac12\mathsf v;
\]

%dsp-level:0
\[
M^2=\frac{\mu_\star^2}{4}
\begin{pmatrix}
g^2&0&0&0\\
0&g^2&0&0\\
0&0&g^2&-gg'\\
0&0&-gg'&g'^2
\end{pmatrix}
\quad\text{in the ordered basis}\ (W_1,W_2,W_3,Y);
\]

%dsp-level:0
\[
m_W^2:=M^2_{11}=M^2_{22}=\frac{\mu_\star^2g^2}{4}=\frac{c_W\,\varkappa^2\beta^2h_0^2\,g^2}{4}>0 .
\]

\textbf{Req.} EW.3, EW.32, EW.33; Bloch-frame fact [linear algebra]; P4.

\textbf{Pf.}

%dsp-level:0
\[
\mathsf v\in\mathbb C^2,\ \|\mathsf v\|=1:\ \ \exists!\,\hat n\in S^2:\ (\hat n\cdot\vec\sigma)\mathsf v=-\mathsf v\ \ (\text{negative Bloch direction});
\]

%dsp-level:0
\[
\quad \exists\,U\in\mathrm{SU}(2):\ \hat n\mapsto e_3
\]

%dsp-level:0
\[
U\ \text{acting by conjugation}
\]

%dsp-level:0
\[
= \text{automorphism of}\ \mathfrak{su}(2)\ [\mathrm{EW.3}]\implies\text{in the rotated frame}\ \sigma_3\mathsf v
\]

%dsp-level:0
\[
= -\mathsf v
\]

%dsp-level:0
\[
\implies W_3\mathsf v
\]

%dsp-level:0
\[
= -\tfrac12\mathsf v\ [\mathrm{EW.32}]
\]

%dsp-level:0
\[
Y=\tfrac12\sigma_0\ [\mathrm{EW.32}]\implies Y\mathsf v=+\tfrac12\mathsf v\ \ \forall\mathsf v
\]

%dsp-level:0
\[
\qquad (\text{hypercharge}\ +\tfrac12:\ \text{derived from trace normalization},\ \neg\text{assigned})
\]

%dsp-level:0
\[
\text{neutral entries}:\ \ \mathcal G_3\mathsf v=-\tfrac g2\mathsf v,\ \ \mathcal G_Y\mathsf v=+\tfrac{g'}2\mathsf v\implies
M^2_{33}=\mu_\star^2\tfrac{g^2}4,
\]

%dsp-level:0
\[
\quad M^2_{YY}=\mu_\star^2\tfrac{g'^2}4,
\]

%dsp-level:0
\[
\quad M^2_{3Y}=\mu_\star^2\bigl(-\tfrac g2\bigr)\bigl(\tfrac{g'}2\bigr)=-\mu_\star^2\tfrac{gg'}4
\]

%dsp-level:0
\[
\text{charged entries}:\ \ \mathsf v
\]

%dsp-level:0
\[
= \binom01:\ \ \sigma_1\mathsf v
\]

%dsp-level:0
\[
= \binom10,
\]

%dsp-level:0
\[
\ \ \sigma_2\mathsf v
\]

%dsp-level:0
\[
= \binom{-i}0
\]

%dsp-level:0
\[
\implies \|W_1\mathsf v\|^2
\]

%dsp-level:0
\[
= \|W_2\mathsf v\|^2
\]

%dsp-level:0
\[
= \tfrac14;
\]

%dsp-level:0
\[
\quad\operatorname{Re}\langle W_1\mathsf v,
\]

%dsp-level:0
\[
W_2\mathsf v\rangle
\]

%dsp-level:0
\[
= \operatorname{Re}\bigl(\tfrac14\cdot(-i)\bigr)\qquad = 0
\]

%dsp-level:0
\[
\implies M^2_{11}=M^2_{22}=\tfrac{\mu_\star^2g^2}4,\quad M^2_{12}=0
\]

%dsp-level:0
\[
\text{cross blocks}:\ \ W_{1,2}\mathsf v\propto\binom10\perp\mathsf v;\quad W_3\mathsf v,\ Y\mathsf v\propto\mathsf v\implies M^2_{1,2;\ 3,Y}=0\implies M^2\ \text{block diagonal}
\]

%dsp-level:0
\[
\mu_\star^2>0\ (h_0\ne0,\ c_W>0)\ [\mathrm{EW.33}]\ \wedge\ g>0\implies m_W^2>0;\qquad
\text{undetermined content}=\mu_\star^2\ \text{only}\ [\mathrm{P4}].\ \blacksquare
\]

\textbf{Phys.} In the frame where the seed points down the third axis, the Gram mass form is block diagonal: two degenerate charged massive directions and a coupled neutral $2\times2$ block. The $W$ mass is derived; the only number left free is the overall wrapping scale $\mu_\star$.

\end{theorem}

\begin{theorem}[{\texorpdfstring{EW.21 — Charged weak bosons}{EW.21 - Charged weak bosons}}]\label{tag:EW-21}
\textbf{Thm.} With $W^\pm=\tfrac1{\sqrt2}(W_1\mp iW_2)$ [D3.10]:

%dsp-level:0
\[
m_{W^+}^2=m_{W^-}^2=m_W^2=\frac{\mu_\star^2g^2}{4}.
\]

\textbf{Req.} D3.10; EW.20; eigenvalue invariance of $c\,\mathbb1$ [linear algebra].

\textbf{Pf.}

%dsp-level:0
\[
\mathrm{EW.20}:\ M^2\big|_{\text{charged}}=\tfrac{\mu_\star^2g^2}{4}\,\mathbb1_2
\]

%dsp-level:0
\[
\operatorname{spec}\bigl(U(c\,\mathbb1_2)U^{-1}\bigr)=\{c,c\}\ \ \forall\ \text{invertible}\ U\ (\text{incl. complexification to}\ W^\pm)
\]

%dsp-level:0
\[
\bigl[\text{equiv.}:\ SO(2)\ \text{invariance generated by}\ W_3;\ W^\pm:\ \text{eigenvectors with phases}\ e^{\mp i\theta}\bigr]
\]

%dsp-level:0
\[
\implies m_{W^\pm}^2=m_W^2=\tfrac{\mu_\star^2g^2}{4}.\ \blacksquare
\]

\textbf{Phys.} The charged weak modes are exactly degenerate: the charged Gram block is a multiple of the identity, so the complex combinations $W^\pm$ inherit the same mass. Charge conjugates share mass by structure, not by assumption.

\end{theorem}

\phantomsection
\section[§D.V — The Neutral Sector]{\texorpdfstring{§D.V — The Neutral Sector}{SD.V - The Neutral Sector}}

\begin{theorem}[{\texorpdfstring{EW.22 — The unbroken generator}{EW.22 - The unbroken generator}}]\label{tag:EW-22}
\textbf{Thm.}

%dsp-level:0
\[
\ker M^2=\operatorname{span}\{(0,0,g',g)\}\quad\text{in the basis}\ (W_1,W_2,W_3,Y);\qquad
Q:=W_3+Y;\qquad Q\,\mathsf v=0;
\]

%dsp-level:0
\[
A:=\frac{g'W_3+gY}{\sqrt{g^2+g'^2}}\qquad(\text{normalized kernel direction}).
\]

\textbf{Req.} EW.20, EW.32, EW.33.

\textbf{Pf.}

%dsp-level:0
\[
\mathrm{EW.33}:\ \ker M^2=\{c:\ \textstyle\sum_Ac_A\mathcal G_A\mathsf v=0\}
\]

%dsp-level:0
\[
\mathrm{EW.20}:\ W_{1,2}\mathsf v\perp\mathsf v\ \wedge\ W_{1,2}\mathsf v\ne0;
\]

%dsp-level:0
\[
\quad W_3\mathsf v,\ Y\mathsf v\propto\mathsf v\implies c_1\ne0\ \vee\ c_2\ne0\implies\textstyle\sum c_A\mathcal G_A\mathsf v\ne0\implies\ker\subseteq\{c_1=c_2=0\}
\]

%dsp-level:0
\[
\bigl(c_3gW_3+c_Yg'Y\bigr)\mathsf v=\tfrac12\bigl(-c_3g+c_Yg'\bigr)\mathsf v=0\iff c_3g=c_Yg'\iff(c_3,c_Y)\propto(g',g)
\]

%dsp-level:0
\[
\implies\dim\ker M^2=1;\qquad
\text{operator form}:\ g'\cdot gW_3+g\cdot g'Y=gg'\,(W_3+Y)\implies Q=W_3+Y\ \ (gg'\ne0)
\]

%dsp-level:0
\[
Q\mathsf v=\bigl(-\tfrac12+\tfrac12\bigr)\mathsf v=0;\qquad A=(g'W_3+gY)/s,\ \ s=\sqrt{g^2+g'^2}.\ \blacksquare
\]

\textbf{Phys.} The massless direction is derived, not posited: it is the unique combination annihilating the seed. Electric charge $Q=W_3+Y$ appears here for the first time — as the stabilizer of the transported qubit.

\end{theorem}

\begin{theorem}[{\texorpdfstring{EW.23 — The neutral mass matrix}{EW.23 - The neutral mass matrix}}]\label{tag:EW-23}
\textbf{Thm.} In $(W_3,Y)$:

%dsp-level:0
\[
M^2_{\mathrm{neutral}}=\frac{\mu_\star^2}{4}
\begin{pmatrix}g^2&-gg'\\ -gg'&g'^2\end{pmatrix}
=\frac{\mu_\star^2}{4}\,uu^{\mathsf T},\quad u=(g,-g');
\]

%dsp-level:0
\[
\det=0;\qquad\operatorname{tr}=\frac{\mu_\star^2(g^2+g'^2)}{4}=:m_Z^2;\qquad
\operatorname{spec}=\{0,\ m_Z^2\};\qquad
\text{null}=(g',g);\quad\text{massive}=(g,-g').
\]

\textbf{Req.} EW.20, EW.22; rank-one facts [linear algebra].

\textbf{Pf.}

%dsp-level:0
\[
\text{entries}=[\mathrm{EW.20}];\qquad
\det=\tfrac{\mu_\star^4}{16}\bigl(g^2g'^2-g^2g'^2\bigr)=0;\qquad
\operatorname{tr}=\tfrac{\mu_\star^2}{4}\bigl(g^2+g'^2\bigr)>0
\]

%dsp-level:0
\[
\det=0\ \wedge\ \operatorname{tr}>0\implies\operatorname{rank}=1,\quad\operatorname{spec}=\bigl\{0,\ \tfrac{\mu_\star^2}{4}(g^2+g'^2)\bigr\}
\]

%dsp-level:0
\[
\text{null vector}=(g',g)\ [\mathrm{EW.22}];\qquad\text{massive eigenvector}\perp\text{null}=(g,-g')
\]

%dsp-level:0
\[
m_Z^2:=\tfrac{\mu_\star^2(g^2+g'^2)}{4};\qquad
M^2_{\mathrm{neutral}}=\tfrac{\mu_\star^2}{4}uu^{\mathsf T}\implies\text{PSD}\ [\mathrm{EW.33}].\ \blacksquare
\]

\textbf{Phys.} The neutral block is a rank-one Gram matrix — one massive mode ($Z$) and one massless mode — because a single seed direction is being probed by two neutral generators. $m_Z^2=\mu_\star^2(g^2+g'^2)/4$ is derived with zero freedom beyond $\mu_\star$.

\end{theorem}

\begin{theorem}[{\texorpdfstring{EW.24 — The mixing angle}{EW.24 - The mixing angle}}]\label{tag:EW-24}
\textbf{Thm.}

%dsp-level:0
\[
\cos\theta_W:=\frac{g}{s};\qquad\sin\theta_W:=\frac{g'}{s};\qquad\tan\theta_W=\frac{g'}{g};\qquad s:=\sqrt{g^2+g'^2};
\]

%dsp-level:0
\[
A=\sin\theta_W\,W_3+\cos\theta_W\,Y;\qquad
Z=\cos\theta_W\,W_3-\sin\theta_W\,Y\qquad[\mathrm{D3.10}];
\]

%dsp-level:0
\[
M^2_{\mathrm{neutral}}\,A=0;\qquad M^2_{\mathrm{neutral}}\,Z=m_Z^2\,Z .
\]

\textbf{Req.} EW.23; D3.10.

\textbf{Pf.}

%dsp-level:0
\[
A\leftrightarrow(\sin\theta_W,\cos\theta_W)\propto(g',g);\qquad Z\leftrightarrow(\cos\theta_W,-\sin\theta_W)\propto(g,-g');\qquad\text{orthonormal}
\]

%dsp-level:0
\[
M^2_{\mathrm{neutral}}A=\frac{m_Z^2}{s^2}\begin{pmatrix}g^2&-gg'\\ -gg'&g'^2\end{pmatrix}\frac1s\begin{pmatrix}g'\\ g\end{pmatrix}
=\frac{m_Z^2}{s^3}\begin{pmatrix}g^2g'-g^2g'\\ -gg'^2+g'g^2\end{pmatrix}=0
\]

%dsp-level:0
\[
M^2_{\mathrm{neutral}}Z=\frac{m_Z^2}{s^2}\begin{pmatrix}g^2&-gg'\\ -gg'&g'^2\end{pmatrix}\frac1s\begin{pmatrix}g\\ -g'\end{pmatrix}
=\frac{m_Z^2}{s^3}\begin{pmatrix}g(g^2+g'^2)\\ -g'(g^2+g'^2)\end{pmatrix}=m_Z^2\,(\cos\theta_W,-\sin\theta_W)=m_Z^2Z.\ \blacksquare
\]

\textbf{Phys.} The weak mixing angle is the geometry of the kernel: the angle between the neutral gauge plane and the seed stabilizer. $\tan\theta_W=g'/g$ is a derived relation — $\theta_W$ is not a new parameter, since $g,g'$ are the configuration data (P4).

\end{theorem}

\begin{theorem}[{\texorpdfstring{EW.25 — The photon direction is massless}{EW.25 - The photon direction is massless}}]\label{tag:EW-25}
\textbf{Thm.} $\ m_A^2=0$ (pole of $A$ at $k^2=0$).

\textbf{Req.} EW.19, EW.24.

\textbf{Pf.}

%dsp-level:0
\[
\mathrm{EW.24}:\ M^2A=0\implies\text{pole displacement}=0\implies\text{pole at}\ k^2=0\quad[\mathrm{EW.19}\equiv\mathrm{D2.5}].\ \blacksquare
\]

\textbf{Phys.} The photon direction is exactly massless — a protected zero of the Gram matrix (the seed’s stabilizer), not a tuned smallness.

\end{theorem}

\begin{theorem}[{\texorpdfstring{EW.26 — Colour modes remain massless at rung 3}{EW.26 - Colour modes remain massless at rung 3}}]\label{tag:EW-26}
\textbf{Thm.} Extending $M^2$ over $\mathcal G\cup\{T_a\}_{a=1}^{8}$:

%dsp-level:0
\[
M^2_{aB}=\mu_\star^2\operatorname{Re}\langle T_a\mathsf v,\ \mathcal G_B\mathsf v\rangle=0;\qquad
M^2_{ab}=0;\qquad
\{T_a\}_{a=1}^8\subseteq\ker M^2;\qquad
m_g^2=0 .
\]

\textbf{Req.} EW.6, EW.19, EW.32, EW.33; CT0.12 (Rmk), CT0.14 (Part 1).

\textbf{Pf.}

%dsp-level:0
\[
\mathrm{EW.6}:\ \mathcal A(3)=\mathcal A(2)\cup\{\mathrm{strong}\}\implies T_a\ \text{active at rung 3, in the connection}
\]

%dsp-level:0
\[
\mathrm{EW.32}:\ T_a\mathsf v=0\ \ \forall a\implies M^2_{aB}=\mu_\star^2\operatorname{Re}\langle0,\ \mathcal G_B\mathsf v\rangle=0;\quad M^2_{ab}=0
\]

%dsp-level:0
\[
\implies\text{colour block}\equiv0\implies8\ \text{colour directions}\subseteq\ker M^2\implies\text{poles at}\ k^2=0\quad[\mathrm{EW.19}]
\]

%dsp-level:0
\[
\text{structural reason}:\ F_H=\text{pair block}\perp F_{\mathrm{strong}}=\text{stack block}\quad[\mathrm{CT0.12\ Rmk},\ \mathrm{CT0.14}].\ \blacksquare
\]

\textbf{Phys.} Gluons stay massless because the wrapping seed is colour-blind: the scalar lives in the pair register and colour in the orthogonal slot register. A structural orthogonality — not a stipulation — protects the colour gauge modes. (Cited in \textbf{F}.)

\end{theorem}

\begin{theorem}[{\texorpdfstring{EW.27 — The scalar excitation}{EW.27 - The scalar excitation}}]\label{tag:EW-27}
\textbf{Thm.} $S_{\mathrm{eff}}[h]:=$ CT1.20 action restricted to the trace mode; $h_0:$ nondegenerate stable stationary point. Then

%dsp-level:0
\[
m_H^2:=\left.\frac{\delta^2S_{\mathrm{eff}}}{\delta h^2}\right|_{h=h_0}>0;\qquad
\text{scalar pole at}\ k^2=m_H^2\ \ (\text{simple});\qquad
m_H^2=\text{configuration data}\ [\mathrm{P4}].
\]

\textbf{Req.} CT1.20, CT2.5 (Part 2); stability criterion [math]; P4.

\textbf{Pf.}

%dsp-level:0
\[
\text{stable stationary}:\quad\left.\frac{\delta S_{\mathrm{eff}}}{\delta h}\right|_{h_0}=0\ \wedge\ \left.\frac{\delta^2S_{\mathrm{eff}}}{\delta h^2}\right|_{h_0}>0
\]

%dsp-level:0
\[
\text{quadratic expansion about}\ h_0:\ \ \text{coefficient}=m_H^2\implies\text{resolvent pole displaced by}\ m_H^2\quad[\mathrm{CT2.5}]
\]

%dsp-level:0
\[
\text{nondegenerate}\implies\text{pole simple};\qquad m_H^2\ \text{value}=\text{shape of}\ S_{\mathrm{eff}}\implies\text{P4 data},\ \neg\text{in the derived ratios}.\ \blacksquare
\]

\textbf{Phys.} The Higgs-like particle is the trace-mode fluctuation about a stable background; its mass is configuration data and never enters any of the derived dimensionless ratios. (Cited in \textbf{G}.)

\end{theorem}

\phantomsection
\section[§D.VI — Derived Electroweak Relations]{\texorpdfstring{§D.VI — Derived Electroweak Relations}{SD.VI - Derived Electroweak Relations}}

\begin{theorem}[{\texorpdfstring{EW.28 — The gauge spectrum}{EW.28 - The gauge spectrum}}]\label{tag:EW-28}
\textbf{Thm.}

%dsp-level:0
\[
\begin{array}{c|c|c}
\text{mode} & \text{combination} & m^2\\ \hline
W^+ & \tfrac1{\sqrt2}(W_1-iW_2) & m_W^2=\mu_\star^2g^2/4>0\\
W^- & \tfrac1{\sqrt2}(W_1+iW_2) & m_W^2\\
Z & \cos\theta_W\,W_3-\sin\theta_W\,Y & m_Z^2=\mu_\star^2(g^2+g'^2)/4>0\\
A & \sin\theta_W\,W_3+\cos\theta_W\,Y & 0
\end{array}
\qquad\tan\theta_W=\frac{g'}{g};
\]

exhaustive: $\dim\mathfrak g_{\mathrm{EW}}=4$ and the four directions are independent.

\textbf{Req.} EW.10, EW.20, EW.21, EW.23, EW.24, EW.25; D3.10.

\textbf{Pf.} Charged: [EW.20, EW.21]; neutral: [EW.23, EW.24]; $m_A^2=0$: [EW.25]; exhaustiveness: $\dim\mathfrak g_{\mathrm{EW}}=4$ [EW.10] $\wedge$ independence $\{W^\pm,Z,A\}$. $\blacksquare$

\textbf{Phys.} Exactly four gauge directions materialize: two degenerate charged massive, one neutral massive, one neutral massless — the electroweak spectrum, all eigen-data of one Gram matrix.

\end{theorem}

\begin{theorem}[{\texorpdfstring{EW.34 — The mass ratio}{EW.34 - The mass ratio}}]\label{tag:EW-34}
\textbf{Thm.}

%dsp-level:0
\[
\frac{m_W^2}{m_Z^2}=\frac{g^2}{g^2+g'^2}=\cos^2\theta_W;\qquad
m_W=m_Z\cos\theta_W;\qquad
\rho:=\frac{m_W^2}{m_Z^2\cos^2\theta_W}=1;
\]

%dsp-level:0
\[
\rho=1\iff j=\tfrac12\iff\dim\mathcal H_{\mathsf v}=2\iff\mathrm{A2}\ (1\mathfrak q\ \text{per channel});
\]

%dsp-level:0
\[
\qquad \rho\ \text{independent of}\ \mu_\star^2,h_0,\varkappa,\beta,c_W,\mathfrak h .
\]

\textbf{Req.} A2; EW.20, EW.23, EW.24, EW.32; spin-$j$ ladder algebra [math].

\textbf{Pf.}

%dsp-level:0
\[
m_W^2=\tfrac{\mu_\star^2g^2}4,
\]

%dsp-level:0
\[
\quad m_Z^2
\]

%dsp-level:0
\[
= \tfrac{\mu_\star^2(g^2+g'^2)}4\ \ \text{from one Gram matrix}\ [\mathrm{EW.20},\mathrm{EW.23}]
\]

%dsp-level:0
\[
\implies \frac{m_W^2}{m_Z^2}
\]

%dsp-level:0
\[
= \frac{g^2}{g^2+g'^2}
\]

%dsp-level:0
\[
= \cos^2\theta_W\ [\mathrm{EW.24}]
\]

%dsp-level:0
\[
\implies\rho\equiv1\ \ \text{(identically in}\ g,g');\qquad\mu_\star^2,h_0,\varkappa,\beta,c_W,\mathfrak h\ \text{cancel}
\]

%dsp-level:0
\[
\text{driver}:\ \ \frac{m_W^2}{m_Z^2}\Big/\cos^2\theta_W\ \text{traces to}\ \frac{\|W_1\mathsf v\|^2}{\|W_3\mathsf v\|^2};
\]

%dsp-level:0
\[
\qquad \text{spin-}j\ \text{lowest weight}:\ \|W_1\mathsf v\|^2=\tfrac14\bigl[j(j{+}1)-(-j)(-j{+}1)\bigr]=\tfrac j2,\quad\|W_3\mathsf v\|^2=j^2
\]

%dsp-level:0
\[
\frac{\|W_1\mathsf v\|^2}{\|W_3\mathsf v\|^2}=\frac{1}{2j}=1\iff j=\tfrac12\iff\mathcal H_{\mathsf v}=\mathbb C^2\ [\mathrm{EW.32}]\iff1\mathfrak q\ \text{per channel}\ [\mathrm{A2}].\ \blacksquare
\]

\textbf{Phys.} The tree-level $\rho$ parameter equals one because the seed carrier is a single qubit: for spin $j$ the charged-to-neutral Gram ratio is $1/(2j)$, forced to $1$ exactly at $j=\tfrac12$, which A2 demands. $m_W=m_Z\cos\theta_W$ is a dimensionless prediction with zero freedom — independent of every configuration parameter and of the calibration scale. (Cited in \textbf{G}.)

\end{theorem}

\begin{theorem}[{\texorpdfstring{EW.35 — The residual coupling}{EW.35 - The residual coupling}}]\label{tag:EW-35}
\textbf{Thm.} Neutral connection $\ \mathcal A_{\mathrm{neutral}}=g\,w^3W_3+g'\,y\,Y$; in the mass eigenbasis ($w^3=\sin\theta_W\,\mathsf A+\cos\theta_W\,\mathsf Z$; $y=\cos\theta_W\,\mathsf A-\sin\theta_W\,\mathsf Z$):

%dsp-level:0
\[
\mathsf A\text{-coefficient}=e\,Q;\qquad
e:=\frac{gg'}{\sqrt{g^2+g'^2}}=g\sin\theta_W=g'\cos\theta_W;\qquad
\frac1{e^2}=\frac1{g^2}+\frac1{g'^2}.
\]

\textbf{Req.} EW.22, EW.24.

\textbf{Pf.}

%dsp-level:0
\[
\text{rotation to mass basis}\ [\mathrm{EW.24}]:
\]

%dsp-level:0
\[
\quad \mathcal A_{\mathrm{neutral}}=\bigl(g\sin\theta_W\,W_3+g'\cos\theta_W\,Y\bigr)\mathsf A+\bigl(g\cos\theta_W\,W_3-g'\sin\theta_W\,Y\bigr)\mathsf Z
\]

%dsp-level:0
\[
g\sin\theta_W\,W_3+g'\cos\theta_W\,Y=\frac{gg'}{s}\,W_3+\frac{g'g}{s}\,Y=\frac{gg'}{s}\,(W_3+Y)=e\,Q
\]

%dsp-level:0
\[
\quad [\mathrm{EW.22},\ \mathrm{EW.24}:\ \sin\theta_W=\tfrac{g'}s,\ \cos\theta_W=\tfrac gs]
\]

%dsp-level:0
\[
e=\frac{gg'}{s}=g\sin\theta_W=g'\cos\theta_W;\qquad
\frac1{e^2}=\frac{g^2+g'^2}{g^2g'^2}=\frac1{g^2}+\frac1{g'^2}.\ \blacksquare
\]

\textbf{Phys.} The coupling of the massless direction is the harmonic combination of $g$ and $g'$: $e=g\sin\theta_W$ is a derived identity, not a definition. (Cited in \textbf{E}, \textbf{G}; the broken `$\ne$’ glyphs of the prior codex are the letter $e$.)

\end{theorem}

\begin{theorem}[{\texorpdfstring{EW.36 — The charge spectrum and its consistency with the chain end-charge}{EW.36 - The charge spectrum and its consistency with the chain end-charge}}]\label{tag:EW-36}
\textbf{Thm.}

%dsp-level:0
\[
Q\big|_{\mathcal H_{\mathsf v}}=\tfrac12\sigma_3+\tfrac12\sigma_0=\operatorname{diag}(1,0);\qquad
\operatorname{spec}Q=\{0,+1\};\qquad
\omega\ [\mathrm{D2.6}]:\ \operatorname{spec}\mapsto\{0,-1\};
\]

%dsp-level:0
\[
\text{charge alphabet}=q\in\{-1,0,+1\}=q_{\mathrm{EM}}(c)\ \text{alphabet of D3.4}\ [\mathrm{EM.9}].
\]

\end{theorem}

\begin{corollary}[{\texorpdfstring{EW.36a}{EW.36a}}]\label{tag:EW-36a}
On the carrier: $Q=W_3+Y$; seed hypercharge $=+\tfrac12$ [EW.20].

\textbf{Req.} EW.20, EW.22, EW.32; D2.6, D3.4; CT0.9 (Part 1); EM.9 (\textbf{E}, forward pointer).

\textbf{Pf.}

%dsp-level:0
\[
\mathrm{EW.32}:\ W_3=\tfrac12\sigma_3,\ Y=\tfrac12\sigma_0\implies Q=\tfrac12(\sigma_3+\sigma_0)=\operatorname{diag}(1,0)
\]

%dsp-level:0
\[
\operatorname{spec}\operatorname{diag}(1,0)=\{1,0\};\qquad\text{eigenvalue}\ 1:\ \text{upper weight state};\quad0:\ \mathsf v\ [\mathrm{EW.22}:\ Q\mathsf v=0]
\]

%dsp-level:0
\[
\mathrm{D2.6}:\ \omega\ \text{reverses charge}\implies\text{reversed carrier}:\ \operatorname{spec}=\{0,-1\}\implies\text{alphabet}=\{-1,0,+1\}
\]

%dsp-level:0
\[
\text{independently}:\ \ \mathrm{D3.4}:\ q_{\mathrm{EM}}(c)
\]

%dsp-level:0
\[
= \tfrac12\bigl(\sigma_{s(c)}-\sigma_{t(c)}\bigr),
\]

%dsp-level:0
\[
\ \ \sigma\in\{+1,-1\}\ [\mathrm{CT0.9}]
\]

%dsp-level:0
\[
\implies q_{\mathrm{EM}}\in\{-1,0,+1\}\ [\mathrm{EM.9}]
\]

%dsp-level:0
\[
\text{inputs disjoint}\ (\text{seed carrier}+\text{trace normalization}\ \ \text{vs}\ \ \text{binary edge phase})\implies\text{agreement}
\]

%dsp-level:0
\[
= \text{consistency check}.\ \blacksquare
\]

\textbf{Phys.} Two independent constructions of charge — the operator route (stabilizer of the seed) and the combinatorial route (difference of edge signs at a chain’s ends) — yield the same ternary alphabet $\{-1,0,+1\}$. Charge quantization is triality of a qubit seed, mirrored by the binary edge register.

\end{corollary}

\begin{theorem}[{\texorpdfstring{EW.29 — The full electroweak spectrum}{EW.29 - The full electroweak spectrum}}]\label{tag:EW-29}
\textbf{Thm.} Rung-3 materialized poles:

%dsp-level:0
\[
\{W^+,\ W^-,\ Z,\ A,\ H\};\qquad
m_{W^\pm}^2=m_W^2>0;\quad m_Z^2=\frac{m_W^2}{\cos^2\theta_W}>0;\quad m_A^2=0;\quad m_H^2>0\ (h_0\ne0\ \text{stable});
\]

exhaustive: gauge directions span $\mathfrak g_{\mathrm{EW}}$ ($\dim=4$) $\oplus$ scalar direction $F_H$ ($\dim=1$).

\textbf{Req.} CT0.12 (Part 1); EW.10, EW.27, EW.28, EW.34.

\textbf{Pf.} Gauge poles + masses: [EW.28, EW.34]; scalar pole: [EW.27]; exhaustiveness: $\dim\mathfrak g_{\mathrm{EW}}=4$ [EW.10] $+\ \dim F_H=1$ [CT0.12]. $\blacksquare$

\textbf{Phys.} The electroweak particle content materialized at rung 3 is exactly: two charged massive, one neutral massive, one neutral massless, one scalar — and nothing else.

\end{theorem}

\begin{theorem}[{\texorpdfstring{EW.30 — Spectral reduction}{EW.30 - Spectral reduction}}]\label{tag:EW-30}
\textbf{Thm.}

%dsp-level:0
\[
\text{rung 2}:\ \mathfrak g_{\mathrm{EW}}=\mathfrak{su}(2)\oplus\mathfrak u(1)\ [\mathrm{EW.10}];\qquad
\text{rung 3}:\ \operatorname{spec}=\{m_W^2,m_W^2,m_Z^2,0\}\ [\mathrm{EW.28}];
\]

%dsp-level:0
\[
\text{“}\mathrm{SU}(2)\times\mathrm{U}(1)\to\mathrm{U}(1)_Q\text{”}\ :=\ \text{naming of the spectral reduction};
\]

%dsp-level:0
\[
\qquad \neg\text{Lie-algebra breaking by commutators}\ [\mathrm{EW.15}];
\]

%dsp-level:0
\[
[Q,Q]=0\implies\text{residual structure}=\mathfrak u(1)_Q\ \ (\text{genuine abelian algebra}).
\]

\textbf{Req.} EW.10, EW.15, EW.22, EW.28, EW.33.

\textbf{Pf.}

%dsp-level:0
\[
\mathrm{EW.15}:\ \operatorname{Stab}_{\mathrm{adj}}=\mathfrak g_{\mathrm{EW}}\implies\text{algebra intact}
\]

%dsp-level:0
\[
\mathrm{EW.33}:\ \text{reduction occurs in}\ \ker M^2=\text{stabilizer of the transported seed};
\]

%dsp-level:0
\[
\quad \mathrm{EW.22}:\ \dim\ker=1,\ \text{spanned by}\ Q
\]

%dsp-level:0
\[
\mathrm{EW.28}:\ \text{spectrum}=3\ \text{massive}+1\ \text{massless pole};
\]

%dsp-level:0
\[
\qquad [Q,Q]=0\implies\mathfrak u(1)_Q\ \text{abelian}\implies\text{name accurate as spectral description}.\ \blacksquare
\]

\textbf{Phys.} The correct statement of “electroweak symmetry breaking” in this corpus: the algebra never breaks — what reduces is the spectrum. Three poles leave the light cone, one stays, and the residual structure is the abelian stabilizer of the seed. The conventional name is retained as a label only.

\end{theorem}

\phantomsection
\section[§D.VII — Completion]{\texorpdfstring{§D.VII — Completion}{SD.VII - Completion}}

\begin{theorem}[{\texorpdfstring{EW.31 — Self-containment of the electroweak sector}{EW.31 - Self-containment of the electroweak sector}}]\label{tag:EW-31}
\textbf{Thm.}

%dsp-level:0
\[
\forall\ \text{statements of}\ \mathbf D:\quad\vdash\ \mathrm{A1\text{–}A9}\ \cup\ \mathrm{C1\text{–}C9}\ \cup\ \mathrm{D0.1\text{–}D4.3}\ \cup\ \{\mathbf B,\mathbf C\ \text{theorems}\};
\]

%dsp-level:0
\[
\neg\exists\ \text{external electroweak law, parameter, or measured value used};
\]

%dsp-level:0
\[
\qquad \text{undetermined}=\{g,\ g',\ h_0,\ c_W,\ \varkappa,\ \beta\}\subset\mathrm{P4};
\]

%dsp-level:0
\[
\forall\ \text{quoted ratios}\ (m_W/m_Z,\ e/g,\ \rho,\ \text{charge alphabet}):\ \ \text{independent of P4 data and of}\ \mathfrak h .
\]

\textbf{Req.} EW.1–EW.30; CT-tags as cited; P4.

\textbf{Pf.} Dependency trace:

%dsp-level:0
\[
\mathrm{EW.2}\equiv\mathrm{CT1.2}\Leftarrow\mathrm{A3};\quad
\mathrm{EW.3}\equiv\mathrm{CT1.4};\quad
\mathrm{EW.6/7}\Leftarrow\mathrm{CT0.13}\Leftarrow\{\mathrm{CT0.7},\mathrm{D3.1},\mathrm{D3.2},\mathrm{D3.3}\}
\]

%dsp-level:0
\[
\mathrm{EW.8\text{–}10}:\ Y=\text{masked centre}\ [\mathrm{CT2.2},\mathrm{CT0.13}];\ \text{commutation}\Leftarrow\text{centrality}
\]

%dsp-level:0
\[
\mathrm{EW.13}:\ \text{scalar}=\text{unique unallocated rotation scalar}\ [\mathrm{CT1.10},\mathrm{CT1.11},\mathrm{CT0.12}]
\]

%dsp-level:0
\[
\mathrm{EW.14/15}:\ \text{background commutes}\implies\neg\text{adjoint breaking};\quad
\mathrm{EW.33}:\ \text{mass}=\text{transport Gram form}
\]

%dsp-level:0
\[
\mathrm{EW.32}:\ \text{carrier}=\mathbb C^2\Leftarrow\mathrm{A2};\ \text{normalization}\Leftarrow\mathfrak T(\mathcal G_A\mathcal G_B)=\tfrac12\delta_{AB}\ (\text{unique})
\]

%dsp-level:0
\[
\mathrm{EW.33}:\ \text{Gram}\implies\text{PSD};\ \ker=\text{stabilizer};\quad
\mathrm{EW.22\text{–}25}:\ \text{kernel}\implies\theta_W,\ m_Z,\ m_A=0
\]

%dsp-level:0
\[
\mathrm{EW.26}:\ F_H\perp F_{\mathrm{strong}}\ [\mathrm{CT0.14}]\implies m_g=0;
\]

%dsp-level:0
\[
\quad \mathrm{EW.34/35}:\ \text{ratios free of P4}\wedge\mathfrak h;
\]

%dsp-level:0
\[
\quad \mathrm{EW.36}:\ \text{alphabets agree (disjoint inputs)}.\ \blacksquare
\]

\textbf{Phys.} Closure audit of the sector: every electroweak statement traces to A1–A9 through B and C; the only undetermined quantities are the P4 configuration data $(g,g',h_0,c_W,\varkappa,\beta)$, and every quoted ratio is independent of all of them.

\end{theorem}

\phantomsection
\section[Closure exports of Part 2 §D (consumed by Part 2 §E–§G,…]{\texorpdfstring{Closure exports of Part 2 §D (consumed by Part 2 §E–§G, Parts 3–4)}{Closure exports of Part 2 SD (consumed by Part 2 SE-SG, Parts 3-4)}}

%dsp-level:0
\[
\begin{array}{lll}
\mathrm{EW.10}\ \mathfrak g_{\mathrm{EW}}=\mathfrak{su}(2)\oplus\mathfrak u(1),\ \dim=4 &
\mathrm{EW.18}\ \Phi(\gamma)\in2\pi\mathbb Z &
\mathrm{EW.20}\ M^2\ \text{block form};\ m_W^2=\mu_\star^2g^2/4\\
\mathrm{EW.22}\ Q=W_3+Y;\ Q\mathsf v=0 &
\mathrm{EW.24}\ \tan\theta_W=g'/g;\ A,\ Z\ \text{eigenbasis} &
\mathrm{EW.25}\ m_A^2=0\\
\mathrm{EW.26}\ m_g^2=0\ (\text{cited in F}) &
\mathrm{EW.27}\ m_H^2>0\ (\text{cited in G}) &
\mathrm{EW.28}\ \text{spectrum}\ \{W^\pm,Z,A\}\\
\mathrm{EW.32}\ \mathcal H_{\mathsf v}=\mathbb C^2;\ W_i\mapsto\tfrac12\sigma_i;\ Y\mapsto\tfrac12\sigma_0;\ T_a\mathsf v=0 &
\mathrm{EW.33}\ M^2=\mu_\star^2\,\mathrm{Gram};\ \ker=\text{stabilizer} &
\mathrm{EW.34}\ \rho=1;\ m_W=m_Z\cos\theta_W\\
\mathrm{EW.35}\ e=gg'/\sqrt{g^2+g'^2};\ 1/e^2=1/g^2+1/g'^2 &
\mathrm{EW.36(a)}\ \operatorname{spec}Q=\{0,1\};\ q\in\{-1,0,+1\} &
\mathrm{EW.29/30}\ \text{full spectrum};\ \text{reduction}=\text{spectral}
\end{array}
\]

$\blacksquare$

\phantomsection
\section[§E — Electromagnetic Sector]{\texorpdfstring{§E — Electromagnetic Sector}{SE - Electromagnetic Sector}}

\textit{Dep.: Part 1; Part 2 §C, §D. Forward cross-section pointers, corroboration only ($\neg$load-bearing): EM.18, EM.26 $\to$ CT3.7 (Part 2 §F; the main arguments are CT0.14-based). Notation: $\mathsf E_{\mathrm{bd}},\mathsf E_{\mathrm{tr}}$ [C6]; $\mathcal P_M^{(n)},\mathbf F_M^{(n)}$ [C7]; $Z$ [C8]; $\alpha_e$ — per-edge dynamic phase [D3.5, Def-EM.20, EM.20]. Machine: \texttt{checks/\allowbreak{}verify\_\allowbreak{}E.py}; Def-EM.20/EM.20/EM.23/EM.24: \texttt{checks/\allowbreak{}verify\_\allowbreak{}F33\_\allowbreak{}dynamic\_\allowbreak{}phase.py}.}

\phantomsection
\section[§E.I — Rung-2 Existence and Algebra]{\texorpdfstring{§E.I — Rung-2 Existence and Algebra}{SE.I - Rung-2 Existence and Algebra}}

\begin{theorem}[{\texorpdfstring{EM.1 — The electromagnetic aspect activates at rung 2}{EM.1 - The electromagnetic aspect activates at rung 2}}]\label{tag:EM-1}
\textbf{Thm.}

%dsp-level:0
\[
\mathrm{EM}\in\mathcal A(2);\qquad\mathrm{EM}\notin\mathcal A(1);\qquad\mathrm{EM}\notin\mathcal A(0);
\]

%dsp-level:0
\[
\neg\exists\ \text{independent aspect named}\ Z\ \text{at rungs}\ \le2:\quad\text{neutral seed}=W_3\in F_W\ [\mathrm{C8}].
\]

\textbf{Req.} A7, C8; D0.8, D3.3; CT0.13 (Part 1).

\textbf{Pf.}

%dsp-level:0
\[
\mathrm{CT0.13}:\ \mathcal A(0)=\{J\};\quad\mathcal A(1)=\{J,G,W,H,D^+\};\quad\mathcal A(2)=\mathcal A(1)\cup\{\mathrm{EM}\}\implies\mathrm{EM}\in\mathcal A(2)
\]

%dsp-level:0
\[
\mathrm{D3.3}:\ f_v\ \text{defined}\iff\partial^{\leftarrow}(v)\ge1\ \wedge\ \partial^{\rightarrow}(v)\ge1\implies n(v)\ge2\quad[\mathrm{D0.8}]
\]

%dsp-level:0
\[
n \le 1:\ \neg\exists v:\ (\partial^{\leftarrow}\ge1\wedge\partial^{\rightarrow}\ge1)
\]

%dsp-level:0
\[
\implies f_v\ \text{undefined}\ \forall v\implies\mathrm{EM}\ \text{has no support}\implies\mathrm{EM}\notin\mathcal A(0)\cup\mathcal A(1)
\]

%dsp-level:0
\[
\mathrm{CT0.13}:\ W_3\in F_W\ (3\text{-dim block});\quad F_{\mathrm{EM}}=\text{distinct stack block}\implies\neg\exists\ \text{aspect}\ Z\ \text{at rungs}\le2.\ \blacksquare
\]

\textbf{Phys.} Electromagnetism switches on exactly at rung 2, because measuring a relative phase needs both an incoming and an outgoing real edge at one vertex. Its existence — and its absence at rungs 0–1 — is a theorem of degree bookkeeping.

\end{theorem}

\begin{theorem}[{\texorpdfstring{EM.2 — The electromagnetic generator}{EM.2 - The electromagnetic generator}}]\label{tag:EM-2}
\textbf{Thm.} [EW.8 realizations]

%dsp-level:0
\[
Y_{\mathrm{asp}}=\mathcal P^{(\mathrm{EM})}:\ \operatorname{rk}=1;\qquad
\mathfrak T\bigl(Y_{\mathrm{asp}}\bigr)=\dim\mathcal K^{(\mathrm{EM})}=1;\qquad
[Y_{\mathrm{asp}},\ \text{any block-diagonal operator}]=0;
\]

%dsp-level:0
\[
Y_{\mathrm{car}}=\tfrac12\sigma_0;\qquad[Y_{\mathrm{car}},W_i|_{\mathrm{car}}]=0\ [\mathrm{EW.32}];
\]

%dsp-level:0
\[
\text{mask fact}\ (\neg Y):\quad\mathcal P_M^{(2)}E_0\mathcal P_M^{(2)}=\mathcal P_M^{(2)},\quad\mathfrak T\bigl(\mathcal P_M^{(2)}\bigr)=\dim\mathcal K_2=17
\]

%dsp-level:0
\[
\quad (E_0\leftrightarrow F_{D^+}\ [\mathrm{CT0.12}]).
\]

\textbf{Req.} CT0.12, CT0.13, CT0.14, CT1.2 (Part 1); CT2.2, D2.3 (Part 2); D2.12; EW.8, EW.32 (Part 2).

\textbf{Pf.}

%dsp-level:0
\[
\dim\mathcal K^{(\mathrm{EM})}=1\ [\mathrm{CT0.14}]\implies Y_{\mathrm{asp}}=\mathcal P^{(\mathrm{EM})}:\ \operatorname{rk}1,\ \ \mathfrak T=\dim\operatorname{Im}=1\ [\mathrm{D2.3},\ \mathrm{CT2.2}]
\]

%dsp-level:0
\[
\mathrm{D2.3}:\ \mathcal K=\bigoplus_A\mathcal K^{(A)}\ \text{orthogonal}\implies[\mathcal P^{(\mathrm{EM})},\ \text{block-diagonal}]=0
\]

%dsp-level:0
\[
Y_{\mathrm{car}}=\tfrac12\sigma_0:\ [\mathrm{EW.32}:\ \text{uniform trace normalization}\ \mathfrak T(\mathcal G_A\mathcal G_B)=\tfrac12\delta_{AB}];
\]

%dsp-level:0
\[
\qquad [\tfrac12\sigma_0,\tfrac12\sigma_i]=0
\]

%dsp-level:0
\[
E_0=I_4\ [\mathrm{CT1.2}]\implies\mathcal P_M^{(2)}E_0\mathcal P_M^{(2)}=\mathcal P_M^{(2)}\ [\mathrm{CT2.2}];\quad\mathfrak T=17\ [\mathrm{CT0.14}];
\]

%dsp-level:0
\[
\quad E_0\leftrightarrow F_{D^+}\ [\mathrm{CT0.12}].\ \blacksquare
\]

\textbf{Rmk (repair).} The prior codex defined $Y:=\mathcal P_M^{(2)}E_0\mathcal P_M^{(2)}=\mathcal P_M^{(2)}$ (rank $17$) and then used it both as the abelian generator and inside the EM.17 trace formula. The rank-$17$ object is the rung-2 \textit{mask}; the abelian generator is the rank-one block projection $Y_{\mathrm{asp}}$ (aspect space) and $\tfrac12\sigma_0$ (carrier). The three realizations $Y_{\mathrm{slot}},Y_{\mathrm{asp}},Y_{\mathrm{car}}$ are coordinated by EW.8/D2.12.

\textbf{Phys.} One generator, three well-typed faces: the centre of the rung-2 slot algebra, the rank-one projection onto the EM block (trace $1$), and $\tfrac12\sigma_0$ on the seed. The trace formula (EM.17) uses the block face; charge statements (EW.36) use the carrier face; the number $17$ belongs to the mask, not to $Y$.

\end{theorem}

\begin{theorem}[{\texorpdfstring{EM.3 — The electromagnetic Lie algebra}{EM.3 - The electromagnetic Lie algebra}}]\label{tag:EM-3}
\textbf{Thm.} $\equiv$ EW.9: $\ \bigl(\operatorname{span}_{\mathbb R}\{Y\},[\cdot,\cdot]\bigr)\cong\mathfrak u(1);\quad[aY,bY]=0;\quad\dim=1.$

\textbf{Req.} EW.9 (Part 2).

\textbf{Pf.} Identity [EW.9]. $\blacksquare$

\textbf{Phys.} The electromagnetic block is the one-dimensional abelian factor: bracket identically zero.

\end{theorem}

\begin{theorem}[{\texorpdfstring{EM.4 — The rung-2 algebra containing the abelian factor}{EM.4 - The rung-2 algebra containing the abelian factor}}]\label{tag:EM-4}
\textbf{Thm.} $\equiv$ EW.10: $\ \mathfrak g_{\mathrm{EW}}=\operatorname{span}_{\mathbb R}\{W_1,W_2,W_3,Y\}=\mathfrak{su}(2)\oplus\mathfrak u(1);\quad\dim=4.$

\textbf{Req.} EW.10 (Part 2); EM.2, EM.3.

\textbf{Pf.} Identity [EW.10: CT1.4 $\mathfrak{su}(2)$; EM.3 $\mathfrak u(1)$; EM.2 commutation; commuting ideals, trivial intersection]. $\blacksquare$

\textbf{Phys.} The abelian factor sits inside the rung-2 electroweak algebra as its commuting centre-plus-identity summand.

\end{theorem}

\phantomsection
\section[§E.II — The Photon]{\texorpdfstring{§E.II — The Photon}{SE.II - The Photon}}

\begin{theorem}[{\texorpdfstring{EM.5 — Photon and gluon as trivial chains}{EM.5 - Photon and gluon as trivial chains}}]\label{tag:EM-5}
\textbf{Thm.} [D1.7, D3.7]

%dsp-level:0
\[
\gamma_{\gamma}\in\mathsf{Trivial}\ \wedge\ \gamma_{\gamma}\subseteq V_2\qquad(\text{photon});\qquad
\gamma_g\in\mathsf{Trivial}\ \wedge\ \gamma_g\cap V_3\ne\varnothing\qquad(\text{gluon});
\]

%dsp-level:0
\[
\forall\gamma\in\mathsf{Trivial}:\ \ \Theta(\gamma)=0\ \wedge\ \operatorname{Hol}_{\mathrm{spiral}}(\gamma)=1;\qquad
V_2\cap V_3=\varnothing;
\]

%dsp-level:0
\[
\text{distinguishing datum}=\text{possibility space only}\ (V_2\ \text{vs}\ V_3\ \text{availability});\qquad\neg\text{bound loops}.
\]

\textbf{Req.} D0.8, D1.7, D3.7; CT0.13 (Part 1).

\textbf{Pf.}

%dsp-level:0
\[
\mathrm{D1.7}:\ \gamma\in\mathsf{Trivial}\iff\forall e\in\gamma:\ \widehat{ds}^{\,2}(e)=0\implies\Theta(\gamma)=\sum_e0=0\implies\Phi(\gamma)=\sum_e\alpha_e\cdot0=0\implies\operatorname{Hol}=e^{i\Phi}=1
\]

%dsp-level:0
\[
\mathrm{D0.8}:\ V_n=\{v:n(v)=n\}\implies V_2\cap V_3=\varnothing
\]

%dsp-level:0
\[
\mathrm{CT0.13}:\ \mathrm{strong}\in\mathcal A(3)\setminus\mathcal A(2)\implies\gamma\subseteq V_2:\ \neg\text{colour index};\quad\gamma\cap V_3\ne\varnothing:\ \text{colour possible}
\]

%dsp-level:0
\[
\text{same class}\ \mathsf{Trivial};\quad\neg\exists\ \text{further distinguishing datum}.\ \blacksquare
\]

\textbf{Phys.} Photon and gluon are one kind of object — a null (trivial) chain of the graph — distinguished only by which vertex rungs it may visit. “Being a photon” is a possibility-space property of rung-2 chains, not a different substance. (Cited in \textbf{F}.)

\end{theorem}

\begin{theorem}[{\texorpdfstring{EM.6 — Photon masslessness}{EM.6 - Photon masslessness}}]\label{tag:EM-6}
\textbf{Thm.} $\ m_\gamma^2=0$, with two independent supports:

%dsp-level:0
\[
(\mathrm{i})\quad A=\sin\theta_W W_3+\cos\theta_W Y\ \text{spans}\ \ker M^2\ [\mathrm{EW.22},\mathrm{EW.24}];
\]

%dsp-level:0
\[
(\mathrm{ii})\quad\gamma_{\gamma}\in\mathsf{Trivial}
\]

%dsp-level:0
\[
\subseteq V_2:\ \ \Theta(\gamma)
\]

%dsp-level:0
\[
= 0\ \wedge\ \text{winding}
\]

%dsp-level:0
\[
= 0\implies\neg\exists\ \text{wrapping holonomy}\implies\neg\exists\ \text{wrapping contribution}\ [\mathrm{EM.5},\mathrm{EW.18}].
\]

\textbf{Req.} D2.5; EM.5; EW.18, EW.22, EW.24, EW.25, EW.33 (Part 2).

\textbf{Pf.}

%dsp-level:0
\[
\mathrm{D2.5}:\ m^2=\text{pole displacement};
\]

%dsp-level:0
\[
\qquad \mathrm{EW.33}:\ \text{displacement source}=\mu_\star^2\Bigl\|\textstyle\sum_Ac_A\mathcal G_A\mathsf v\Bigr\|^2\ (\text{rung-3 wrapping form})
\]

%dsp-level:0
\[
\mathrm{EW.22}:\ \ker M^2=\operatorname{span}\{(g',g)\}_{(W_3,Y)}=A\text{-direction}\ [\mathrm{EW.24}]
\]

%dsp-level:0
\[
(\mathrm{ii}):\ \ \text{trivial}\implies\Theta=0\ [\mathrm{EM.5}];\quad\text{winding}=0\implies e^{i\Phi}=1\ \text{vacuous}\implies\neg\text{materialization from wrapping}\ [\mathrm{EW.18}]
\]

%dsp-level:0
\[
\text{neutral direction with}\ \neg\text{wrapping}=\ker\text{-direction}\implies\gamma=A\implies\text{displacement}=0\implies m_\gamma^2=0\quad[=\mathrm{EW.25}].\ \blacksquare
\]

\textbf{Phys.} The photon is massless twice over: algebraically it spans the kernel of the Gram mass form (the seed’s stabilizer), and geometrically a null chain confined to rung-2 vertices cannot wind, so the wrapping defect never sees it. (Cited in \textbf{F}, \textbf{G}.)

\end{theorem}

\begin{theorem}[{\texorpdfstring{EM.7 — Photon entropy budget}{EM.7 - Photon entropy budget}}]\label{tag:EM-7}
\textbf{Thm.} Photon chain of hop count $\delta h$:

%dsp-level:0
\[
\mathsf E_{\mathrm{tr}}=\delta h\,\ln2\ \ \text{nats}=\delta h\ \ \text{bits};\qquad
\mathsf E_{\mathrm{bd}}=0 .
\]

\textbf{Req.} C9; D1.5; CT0.9 (Part 1); EM.5.

\textbf{Pf.}

%dsp-level:0
\[
\mathrm{CT0.9}+\mathrm{C9}:\ 1\ \text{real edge}=1\ \text{bit}=\ln2\ \text{nats};\qquad
\text{edges}=\text{independent registers}\implies\delta h\ \text{hops}:\ \delta h\ln2
\]

%dsp-level:0
\[
\mathrm{D1.5}:\ \text{open propagating chain}=\text{transit}\ [\mathrm{EM.5}]\implies\mathsf E_{\mathrm{tr}}=\delta h\ln2
\]

%dsp-level:0
\[
\mathsf E_{\mathrm{bd}}=\Delta S_{\mathrm{bound}}\ [\mathrm{D1.5}];\quad\text{photon}\ \neg\text{closed}\implies\mathsf E_{\mathrm{bd}}=0.\ \blacksquare
\]

\textbf{Phys.} A photon carries exactly one bit of transit entropy per hop and zero bound entropy: bosonic transport content, measured by the ledger’s own edge count.

\end{theorem}

\begin{theorem}[{\texorpdfstring{EM.8 — Photon polarization count}{EM.8 - Photon polarization count}}]\label{tag:EM-8}
\textbf{Thm.} $D=4$; $k_\mu\ne0$ null $\bigl(k^\mu k_\mu=0\bigr)$; $\epsilon_\mu\sim\epsilon_\mu+k_\mu$ (gauge, EM.14); transversality $k^\mu\epsilon_\mu=0$ (source-free equation, CT1.22). Then

%dsp-level:0
\[
N_{\mathrm{pol}}=\dim\bigl(k^\perp/\operatorname{span}\{k\}\bigr)=(4-1)-1=2,\qquad
k^\perp:=\{v:\ k_\mu v^\mu=0\}.
\]

\textbf{Req.} CT1.22, CT2.6 (Part 2); EM.6, EM.14; rank–nullity [linear algebra].

\textbf{Pf.}

%dsp-level:0
\[
\mathrm{CT1.22}:\ F_{\mu\nu}=2\nabla_{[\mu}A_{\nu]};
\]

%dsp-level:0
\[
\qquad \mathrm{EM.14}:\ A_\mu\sim A_\mu+\nabla_\mu\chi\implies\epsilon_\mu\propto k_\mu:\ 2k_{[\mu}k_{\nu]}=0\implies F=0:\ \text{pure gauge}
\]

%dsp-level:0
\[
\text{source-free equation}\ [\mathrm{CT1.22}]\implies k^\mu\epsilon_\mu=0\implies\epsilon\in k^\perp
\]

%dsp-level:0
\[
k_\mu\ne0\implies\operatorname{rank}(k_\cdot)=1\implies\dim k^\perp=4-1=3\quad[\text{rank--nullity}]
\]

%dsp-level:0
\[
k^\mu k_\mu=0\implies k\in k^\perp;\quad k\ne0\implies\dim\bigl(k^\perp/\operatorname{span}\{k\}\bigr)=3-1=2.\ \blacksquare
\]

\textbf{Phys.} A massless spin-1 mode in four dimensions has exactly two physical polarizations: transversality gives three, gauge redundancy removes one. The count is derived from dimension arithmetic plus nullness — the same $N_{\gamma\text{-pol}}=2$ of policy P2. (Cited in \textbf{G}.)

\end{theorem}

\phantomsection
\section[§E.III — Electromagnetic Charge]{\texorpdfstring{§E.III — Electromagnetic Charge}{SE.III - Electromagnetic Charge}}

\begin{theorem}[{\texorpdfstring{EM.9 — Chain end-charge}{EM.9 - Chain end-charge}}]\label{tag:EM-9}
\textbf{Thm.} [D3.4]

%dsp-level:0
\[
q_{\mathrm{EM}}(c)=\frac{\sigma_{s(c)}-\sigma_{t(c)}}{2}\in\{-1,0,+1\};\qquad
\text{all three values attained};\qquad
\begin{array}{c|cccc}
(\sigma_s,\sigma_t) & (+,+) & (+,-) & (-,+) & (-,-)\\ \hline
q_{\mathrm{EM}} & 0 & +1 & -1 & 0
\end{array}
\]

\textbf{Req.} D3.4; CT0.9 (Part 1).

\textbf{Pf.}

%dsp-level:0
\[
\mathrm{CT0.9}:\ \sigma_e\in\{+1,-1\}\implies\sigma_{s(c)},
\]

%dsp-level:0
\[
\sigma_{t(c)}\in\{+1,-1\}
\]

%dsp-level:0
\[
\implies q_{\mathrm{EM}}\in\bigl\{\tfrac{1-1}2,\tfrac{1-(-1)}2,\tfrac{-1-1}2,\tfrac{-1-(-1)}2\bigr\}\qquad = \{0,+1,-1\}
\]

%dsp-level:0
\[
\text{end phases}=\text{independent registers}\implies\forall(\sigma_s,\sigma_t)\in\{\pm1\}^2\ \text{realizable}.\ \blacksquare
\]

\textbf{Phys.} The charge of an open chain is the half-difference of its two end-signs: a ternary alphabet arising from binary edges. (Cited in \textbf{D} as EW.36’s independent counterpart, and in \textbf{G}.)

\end{theorem}

\begin{corollary}[{\texorpdfstring{EM.9b — chain-wall telescoping}{EM.9b - chain-wall telescoping}}]\label{tag:EM-9b}
Let $c=(e_1,\dots,e_k)$, $k\ge2$, be an oriented chain of distinct real edges [D1.4] with interior vertices $v_1,\dots,v_{k-1}$, $v_i:=t(e_i)=s(e_{i+1})$, each of type $(1,1)$ [A7, D0.8] — the chain's two edges are $v_i$'s only incident real edges. Then the D3.3 vertex reads telescope:

%dsp-level:0
\[
\sum_{i=1}^{k-1}f_{v_i}\ =\ \tfrac12\bigl(\sigma_{e_k}-\sigma_{e_1}\bigr)\ =\ -\,q_{\mathrm{EM}}(c)
\]

%dsp-level:0
\[
\qquad \bigl[\sigma_{s(c)}=\sigma_{e_1},\ \ \sigma_{t(c)}=\sigma_{e_k}:\ \mathrm{D3.4}/\mathrm{EM.9}\ \text{end registers}\bigr]
\]

— the total interior vertex activity of a chain is carried by its two end registers alone. In particular:

%dsp-level:0
\[
\sum_{i=1}^{k-1}f_{v_i}=0\iff\sigma_{e_1}=\sigma_{e_k}\iff q_{\mathrm{EM}}(c)=0;
\]

%dsp-level:0
\[
\qquad \sum_{i=1}^{k-1}\aleph_{v_i}=\#\{i:\ \sigma_{e_{i+1}}\ne\sigma_{e_i}\}\le k-1
\]

%dsp-level:0
\[
\qquad [\mathrm{D3.3}:\ \aleph_v=|f_v|;\ \mathrm{CT0.9}].
\]

\textbf{Req (Cor).} A7; D0.8, D1.4, D3.3, D3.4; CT0.9 (Part 1); EM.9; finite telescoping [math; machine: \texttt{checks/\allowbreak{}verify\_\allowbreak{}promotion\_\allowbreak{}E\_\allowbreak{}telescoping.py}].

\textbf{Pf (Cor).}

%dsp-level:0
\[
\text{type}(1,1)\ [\mathrm{A7},\mathrm{D0.8}]:\ \sigma_{\leftarrow}(v_i)=\sigma_{e_i},\ \ \sigma_{\rightarrow}(v_i)=\sigma_{e_{i+1}};
\]

%dsp-level:0
\[
\qquad \mathrm{D3.3}:\ f_{v_i}=\tfrac12\bigl(\sigma_{\rightarrow}(v_i)-\sigma_{\leftarrow}(v_i)\bigr)=\tfrac12\bigl(\sigma_{e_{i+1}}-\sigma_{e_i}\bigr)
\]

%dsp-level:0
\[
\sum_{i=1}^{k-1}f_{v_i}=\tfrac12\sum_{i=1}^{k-1}\bigl(\sigma_{e_{i+1}}-\sigma_{e_i}\bigr)=\tfrac12\bigl(\sigma_{e_k}-\sigma_{e_1}\bigr)\quad[\text{telescoping}];
\]

%dsp-level:0
\[
\qquad \mathrm{D3.4}/\mathrm{EM.9}:\ q_{\mathrm{EM}}(c)=\tfrac12\bigl(\sigma_{e_1}-\sigma_{e_k}\bigr)=-\sum_{i}f_{v_i}
\]

%dsp-level:0
\[
\mathrm{CT0.9}:\ \sigma\in\{\pm1\}:\ f_{v_i}\ne0\iff\sigma_{e_{i+1}}\ne\sigma_{e_i};\qquad\aleph_{v_i}=|f_{v_i}|\in\{0,1\}\ [\mathrm{D3.3}]\ \blacksquare
\]

\end{corollary}

\begin{theorem}[{\texorpdfstring{EM.9a — The closed-core charge read and the end-register count}{EM.9a - The closed-core charge read and the end-register count}}]\label{tag:EM-9a}
\textbf{Thm.} $(1)$ \textbf{Closed-core read.} For a closed oriented chain $c=(e_1,\dots,e_k)$ based at $v=s(c)=t(c)$ [D1.4: ordered tuple] with first edge $e_1$ (out of $v$) and last edge $e_k$ (into $v$):

%dsp-level:0
\[
q_{\mathrm{core}}(c;v):=\frac{\sigma_{e_1}-\sigma_{e_k}}{2}\ \in\ \{-1,0,+1\}
\]

is well-defined [D3.4 applied to closed $c$; CT0.9: per-edge independent registers; EM.9: alphabet, all values attained]. $(2)$ \textbf{Charge-conjugation law.} Under rebase / orientation reversal the roles of $e_1$ and $e_k$ swap:

%dsp-level:0
\[
q_{\mathrm{core}}\ \mapsto\ -\,q_{\mathrm{core}}
\]

— consistent with, and an instance of, D2.6's $\omega$ (charge $\mapsto$ $-$charge): the corpus's own charge-conjugation law, verbatim on the core read. At the 2-cycle floor: $q_{\mathrm{core}}=\tfrac12(\sigma_{e_1}-\sigma_{e_2})$ at base $p$, negated at base $q$. $(3)$ \textbf{End-register count.} The EM.9/D3.4 read consumes exactly the two end registers of the chain: $N_{\mathrm{charge}}=2$ is a theorem, $\neg$an enumeration convention; the ternary alphabet is the difference of two binary registers [CT0.9], with all three values attained (witnesses $(\sigma_1,\sigma_k)=(+,-)\mapsto+1$, $(+,+)\mapsto0$, $(-,+)$ or $(-,-)\mapsto$ the remaining values; surjectivity machine-checked). $(4)$ \textbf{Not neutral by closure.} The closed core is $\neg$forced charge-neutral: its ternary read is basepoint data. The observer-facing charge of a transit chain remains the EM.9 end-read.

\textbf{Req.} A1; D1.4, D2.6, D3.4; CT0.9 (Part 1); EM.9; witness enumeration [math; machine: \texttt{checks/\allowbreak{}verify\_\allowbreak{}merge\_\allowbreak{}tierB.py} §D].

\textbf{Pf.}

%dsp-level:0
\[
(1):\ \mathrm{D3.4}\ \text{is stated for oriented chains generally};
\]

%dsp-level:0
\[
\quad \text{closed}\ c:\ s(c)
\]

%dsp-level:0
\[
= t(c)
\]

%dsp-level:0
\[
= v\ \text{with distinguished first/last edges}\ [\mathrm{D1.4}]\implies\sigma_{s(c)}:
\]

%dsp-level:0
\[
= \sigma_{e_1},
\]

%dsp-level:0
\[
\ \sigma_{t(c)}:
\]

%dsp-level:0
\[
= \sigma_{e_k}\ \text{well-defined}
\]

%dsp-level:0
\[
\mathrm{CT0.9}:\ \sigma_{e_1},\sigma_{e_k}\in\{\pm1\}\ \text{independent};
\]

%dsp-level:0
\[
\qquad q_{\mathrm{core}}\in\{\tfrac12(\pm1\mp1)\}=\{-1,0,+1\};
\]

%dsp-level:0
\[
\qquad \mathrm{EM.9}:\ \text{all attained}
\]

%dsp-level:0
\[
(2):\ \text{rebase/orientation reversal}:\ e_1\leftrightarrow e_k\ \text{roles swap}
\]

%dsp-level:0
\[
\implies q_{\mathrm{core}}\mapsto\tfrac12(\sigma_{e_k}-\sigma_{e_1})
\]

%dsp-level:0
\[
= -q_{\mathrm{core}};
\]

%dsp-level:0
\[
\qquad \mathrm{D2.6}:\ \omega\ \text{maps charge}\mapsto-\text{charge}:\ \text{consistent}
\]

%dsp-level:0
\[
(3):\ \mathrm{D3.4}'\text{s formula contains exactly}\ \sigma_{s(c)},\sigma_{t(c)}:\ \text{two registers};\ \text{interior signs}\ \neg\text{enter};
\]

%dsp-level:0
\[
\quad \#\{\tfrac12(\sigma-\sigma'):\ \sigma,\sigma'\in\{\pm1\}\}=3
\]

%dsp-level:0
\[
(4):\ (1)'\text{s witnesses include}\ q_{\mathrm{core}}=\pm1\ne0.\ \blacksquare
\]

\textbf{Phys.} A closed chain carries a ternary charge-type read at its basepoint even though it has no ends: the difference of the two registers where it leaves and re-enters its base. Orientation reversal negates it exactly — charge conjugation acting on the only charge-like label closed content can carry. The count is rigid: two registers make a charge, as a theorem.

\end{theorem}

\begin{theorem}[{\texorpdfstring{EM.10 — Charge quantization}{EM.10 - Charge quantization}}]\label{tag:EM-10}
\textbf{Thm.}

%dsp-level:0
\[
S^2
\]

%dsp-level:0
\[
= T^2
\]

%dsp-level:0
\[
= \mathbb1\ [\mathrm{A3}]\implies\phi_e\in\{0,\pi\}\implies\sigma_e\in\{+1,-1\}\implies\sigma_{s}-\sigma_{t}\in\{-2,0,+2\}
\]

%dsp-level:0
\[
\implies q_{\mathrm{EM}}\in\mathbb Z\cap\{-1,0,+1\};
\]

%dsp-level:0
\[
q\notin\mathbb Z\implies\phi_e\notin\Phi\implies\neg\mathrm{CT0.9}\qquad(\neg\exists\ \text{finer values}).
\]

\textbf{Req.} A3; D0.5, D3.4; CT0.9 (Part 1).

\textbf{Pf.}

%dsp-level:0
\[
\mathrm{CT0.9}\ [\Leftarrow\mathrm{D0.5},\ \mathrm{A3}:\ S^2=T^2=\mathbb1\implies\text{residual holonomy}\ \mathbb Z_2]:\ \ \phi_e\in\{0,\pi\}\implies\sigma_e=e^{i\phi_e}\in\{\pm1\}
\]

%dsp-level:0
\[
\sigma_s,\sigma_t\in\{\pm1\}\implies\sigma_s-\sigma_t\in\{-2,0,+2\}\subset2\mathbb Z\implies q_{\mathrm{EM}}=\tfrac12(\sigma_s-\sigma_t)\in\{-1,0,+1\}\subset\mathbb Z
\]

%dsp-level:0
\[
q_{\mathrm{EM}}\notin\mathbb Z\implies\sigma_s-\sigma_t\ \text{odd}\implies\sigma\notin\{\pm1\}\implies\phi_e\notin\{0,\pi\}\implies\neg\mathrm{CT0.9}\implies\bot.\ \blacksquare
\]

\textbf{Phys.} Charge quantization is not imposed — it is the arithmetic shadow of the binary edge phase, which is itself forced by the involution axiom $S^2=T^2=\mathbb1$. Integers all the way down. (Cited in \textbf{G}.)

\end{theorem}

\begin{theorem}[{\texorpdfstring{EM.11 — Global charge conservation}{EM.11 - Global charge conservation}}]\label{tag:EM-11}
\textbf{Thm.} $U$ valid [D2.11] $\implies\ \displaystyle Q_{\mathrm{EM}}:=\sum_{c\in U}q_{\mathrm{EM}}(c)=0.$

\textbf{Req.} D1.2, D1.3, D1.4, D2.11, D3.4; EM.9.

\textbf{Pf.}

%dsp-level:0
\[
\mathrm{D3.4}:\ q_{\mathrm{EM}}(c)=\tfrac12\sigma_{s(c)}-\tfrac12\sigma_{t(c)}\implies
Q_{\mathrm{EM}}=\tfrac12\!\!\sum_{\text{source-ends}}\!\!\sigma\ -\ \tfrac12\!\!\sum_{\text{target-ends}}\!\!\sigma
\]

%dsp-level:0
\[
\qquad (\text{regroup by ends})
\]

%dsp-level:0
\[
\mathrm{D2.11}:\ \neg\exists\ \text{dangling ledger boundary}\implies\forall\ \text{ends sealed}
\]

%dsp-level:0
\[
\mathrm{D1.2}:\ \text{seal joins two edge-ends};
\]

%dsp-level:0
\[
\quad\mathrm{D1.4}:\ \text{chains oriented}\implies\text{seal joins target-end}\to\text{source-end}\ \ (\text{same-orientation join}\implies\neg\text{continuable orientation},\ \bot\,\mathrm{D1.4})
\]

%dsp-level:0
\[
\mathrm{D1.3}:\ \#J^+=\#J^-\implies\text{pairing covers every end exactly once}
\]

%dsp-level:0
\[
\text{seal identifies joined ends}\implies\text{equal}\ \sigma\implies\text{per-seal contribution}=\tfrac12\sigma-\tfrac12\sigma=0
\]

%dsp-level:0
\[
Q_{\mathrm{EM}}=\sum_{\text{seals}}0=0.\ \blacksquare
\]

\textbf{Phys.} Total charge vanishes in every valid universe: each chain end is sealed against an oppositely oriented end carrying the same sign. Charge conservation is Janus topology — seals pair sources with sinks — not a Noether afterthought.

\end{theorem}

\phantomsection
\section[§E.IV — Field Strength and Dynamics]{\texorpdfstring{§E.IV — Field Strength and Dynamics}{SE.IV - Field Strength and Dynamics}}

\begin{theorem}[{\texorpdfstring{EM.12 — Discrete field strength from phase transitions}{EM.12 - Discrete field strength from phase transitions}}]\label{tag:EM-12}
\textbf{Thm.} [D3.3]

%dsp-level:0
\[
f_v=\tfrac12\bigl(\sigma_{\rightarrow}(v)-\sigma_{\leftarrow}(v)\bigr)\in\{-1,0,+1\};\qquad\aleph_v=|f_v|\in\{0,1\};
\]

%dsp-level:0
\[
\aleph(\sigma,C)=0\iff\forall v:\ \sigma_{\leftarrow}(v)=\sigma_{\rightarrow}(v)\iff\sigma\ \text{constant on connected components}\ \ (\text{flat});
\]

%dsp-level:0
\[
\aleph_v=1\iff\text{oriented transition at}\ v;
\]

%dsp-level:0
\[
\qquad \text{coarse-graining}\ [\mathrm{CT1.19}]:\ f_v\longrightarrow F_{\mu\nu}\ \ (\text{circulation per plaquette})\ [\mathrm{CT1.22}].
\]

\textbf{Req.} D3.3; CT0.9 (Part 1); CT1.19, CT1.22 (Part 2).

\textbf{Pf.}

%dsp-level:0
\[
\sigma_{\leftarrow},\sigma_{\rightarrow}\in\{\pm1\}\ [\mathrm{CT0.9}]\implies f_v\in\{-1,0,+1\},\ \ |f_v|\in\{0,1\}=\aleph_v\quad[\mathrm{D3.3}]
\]

%dsp-level:0
\[
\aleph=0\iff\forall v:\ \aleph_v=0\iff\forall v:\ \sigma_{\leftarrow}=\sigma_{\rightarrow}\iff\sigma\ \text{constant along components}
\]

%dsp-level:0
\[
\quad (\text{zero discrete curvature})
\]

%dsp-level:0
\[
\aleph_v=1\iff\sigma_{\leftarrow}\ne\sigma_{\rightarrow}\iff f_v=\pm1:\ \text{oriented transition}
\]

%dsp-level:0
\[
\mathrm{CT1.19}:\ \text{holonomy per unit area}\to F_{\mu\nu}\ [\mathrm{CT1.22}];\quad\text{vertex sign flip}=\text{nonzero circulation}.\ \blacksquare
\]

\textbf{Phys.} The electromagnetic field on the ledger is the sign mismatch at vertices: $f_v$ is the oriented discrete field strength and $\aleph_v$ its magnitude. Zero excitation means signs flow unflipped; each flip is one excitation quantum.

\end{theorem}

\begin{theorem}[{\texorpdfstring{EM.13 — The transition-parity factor in the weight}{EM.13 - The transition-parity factor in the weight}}]\label{tag:EM-13}
\textbf{Thm.} $w(C)=\mathcal J(C)\prod_{e}G(e)\cdot(-1)^{\aleph(\sigma,C)}$ [A8]:

%dsp-level:0
\[
(-1)^{\aleph}=\begin{cases}+1&\aleph\ \text{even}\\ -1&\aleph\ \text{odd}\end{cases};\qquad
\prod_eG(e)\supseteq\prod_e\sigma_e\ [\mathrm{D3.5}];\qquad
\Bigl(\prod_e\sigma_e\Bigr)\ \text{independent of}\ (-1)^{\aleph}\ [\mathrm{CT0.10}];
\]

%dsp-level:0
\[
Z(\mathcal R)=\sum_C\sum_\sigma w(C):\ \ \Delta\aleph\ \text{odd}\implies\text{sign flip}\implies\text{destructive interference}\ \ (\text{discrete EM dynamics});
\]

%dsp-level:0
\[
\quad \text{coarse-graining}\implies\text{relative signs of}\ F_{\mu\nu}.
\]

\textbf{Req.} A8; D3.5; CT0.9, CT0.10 (Part 1).

\textbf{Pf.}

%dsp-level:0
\[
(-1)^{\aleph}\in\{+1,-1\}\ \text{by parity}\ [\text{definition}]
\]

%dsp-level:0
\[
\mathrm{D3.5}:\ G(e)=\sigma_e\exp\bigl(i\alpha_e\widehat{ds}^{\,2}(e)\bigr)\implies\prod_eG(e)\supseteq\prod_e\sigma_e
\]

%dsp-level:0
\[
\mathrm{CT0.10}:\ \Bigl(\prod_e\sigma_e,\ (-1)^{\aleph}\Bigr)\ \text{independent},\ \text{joint image}=\{\pm1\}^2\implies\text{two distinct sign data};\ \neg\text{double counting}
\]

%dsp-level:0
\[
Z:\ C\mapsto C'\ \text{with}\ \Delta\aleph\ \text{odd}\implies w(C')=-w(C)\ \text{(ceteris paribus)}\implies\text{cancellation}.\ \blacksquare
\]

\textbf{Phys.} The A8 weight assigns a strict minus sign to every odd number of domain walls in the sign field: this destructive interference between one-transition-apart configurations \textit{is} the lattice dynamics of electromagnetism, and it survives coarse-graining as the Maxwell sign structure.

\end{theorem}

\begin{theorem}[{\texorpdfstring{EM.14 — Continuum gauge invariance}{EM.14 - Continuum gauge invariance}}]\label{tag:EM-14}
\textbf{Thm.}

%dsp-level:0
\[
A_\mu\sim A_\mu+\nabla_\mu\chi;\qquad
F_{\mu\nu}=2\nabla_{[\mu}A_{\nu]}\ \mapsto\ F_{\mu\nu}+2\nabla_{[\mu}\nabla_{\nu]}\chi=F_{\mu\nu};
\]

%dsp-level:0
\[
\implies\text{gradient direction}=\text{redundancy};\quad N_{\mathrm{pol}}\ \text{count unaffected}\ [\mathrm{EM.8}].
\]

\textbf{Req.} CT1.22, CT2.6 (Part 2); EM.12; $\nabla_{[\mu}\nabla_{\nu]}\chi=0$ for scalars [math].

\textbf{Pf.}

%dsp-level:0
\[
\mathrm{EM.12}:\ \text{physical datum}=f_v\propto\sigma_{\rightarrow}-\sigma_{\leftarrow}:\ \text{differences only};\ \neg\text{dependence on common offset}
\]

%dsp-level:0
\[
\text{smooth offset field}\xrightarrow[\mathrm{CT1.19}]{}\chi:\quad A_\mu\mapsto A_\mu+\nabla_\mu\chi
\]

%dsp-level:0
\[
F_{\mu\nu}\mapsto F_{\mu\nu}+2\nabla_{[\mu}\nabla_{\nu]}\chi=F_{\mu\nu}\quad[\nabla_\mu\nabla_\nu\chi=\nabla_\nu\nabla_\mu\chi\ \text{on scalars}].\ \blacksquare
\]

\textbf{Phys.} Gauge invariance is the continuum image of a ledger triviality: vertices record sign \textit{changes}, never absolute signs. The redundancy is derived from the discreteness, not postulated.

\end{theorem}

\begin{theorem}[{\texorpdfstring{EM.15 — Electromagnetic stress-energy}{EM.15 - Electromagnetic stress-energy}}]\label{tag:EM-15}
\textbf{Thm.}

%dsp-level:0
\[
S_{\mathrm{EM}}[A,g]=-\tfrac14\int d^4x\sqrt{-g}\,F_{\mu\nu}F^{\mu\nu};
\]

%dsp-level:0
\[
\qquad T^{(\mathrm{EM})}_{\mu\nu}:=-\frac{2}{\sqrt{-g}}\frac{\delta S_{\mathrm{EM}}}{\delta g^{\mu\nu}}=F_{\mu\rho}F_\nu{}^{\rho}-\tfrac14g_{\mu\nu}F_{\alpha\beta}F^{\alpha\beta};
\]

%dsp-level:0
\[
g^{\mu\nu}T^{(\mathrm{EM})}_{\mu\nu}=0\qquad(D=4).
\]

\textbf{Req.} CT1.22 (Part 2); EM.6; $\delta\sqrt{-g}=-\tfrac12\sqrt{-g}\,g_{\mu\nu}\delta g^{\mu\nu}$, $\delta\bigl(g^{\mu\alpha}g^{\nu\beta}F_{\mu\nu}F_{\alpha\beta}\bigr)=2F_{\mu\rho}F_\nu{}^\rho\,\delta g^{\mu\nu}$ [math; machine: \texttt{checks/\allowbreak{}verify\_\allowbreak{}E.py}].

\textbf{Pf.}

%dsp-level:0
\[
\mathrm{CT1.22}:\ S_{\mathrm{EM}}=\text{unique local, diff-}\wedge\text{gauge-invariant},\ \le2\text{nd-order functional quadratic in}\ F
\]

%dsp-level:0
\[
\delta S_{\mathrm{EM}}=-\tfrac14\int\sqrt{-g}\Bigl[\underbrace{-\tfrac12g_{\mu\nu}F^2}_{\delta\sqrt{-g}}+\underbrace{2F_{\mu\rho}F_\nu{}^{\rho}}_{\delta(F^2)}\Bigr]\delta g^{\mu\nu}
\]

%dsp-level:0
\[
\qquad [F_{\mu\nu}\ \text{fixed lower indices}]
\]

%dsp-level:0
\[
T^{(\mathrm{EM})}_{\mu\nu}
\]

%dsp-level:0
\[
= -\tfrac{2}{\sqrt{-g}}\cdot(-\tfrac14)\sqrt{-g}\bigl[2F_{\mu\rho}F_\nu{}^{\rho}-\tfrac12g_{\mu\nu}F^2\bigr]
\]

%dsp-level:0
\[
= F_{\mu\rho}F_\nu{}^{\rho}-\tfrac14g_{\mu\nu}F_{\alpha\beta}F^{\alpha\beta}
\]

%dsp-level:0
\[
\text{trace}:\ \ g^{\mu\nu}T_{\mu\nu}=F_{\mu\rho}F^{\mu\rho}-\tfrac14\cdot4\cdot F^2=0
\]

%dsp-level:0
\[
\qquad [\text{consistency}:\ \text{massless transit aspect},\ \mathrm{EM.6}].\ \blacksquare
\]

\textbf{Phys.} The Maxwell stress tensor is derived as the metric variation of the unique admissible quadratic action — the $a=\mathrm{EM}$ instance of CT1.22 — and its vanishing trace is the local imprint of photon masslessness.

\end{theorem}

\phantomsection
\section[§E.V — The Dynamic Phase]{\texorpdfstring{§E.V — The Dynamic Phase}{SE.V - The Dynamic Phase}}

\begin{theorem}[{\texorpdfstring{EM.16 — Edge transfer weight and the per-edge read-back}{EM.16 - Edge transfer weight and the per-edge read-back}}]\label{tag:EM-16}
\textbf{Thm.} [D3.5] For every real edge $e$:

%dsp-level:0
\[
G(e)=\sigma_e\,\exp\!\bigl(i\,\alpha_e\,\widehat{ds}^{\,2}(e)\bigr);\qquad
\alpha_e=\frac1i\,\frac{\partial}{\partial\,\widehat{ds}^{\,2}(e)}\log\!\bigl(\sigma_e^{-1}G(e)\bigr);\qquad
[\alpha_e]=\mathfrak q^{0}.
\]

\textbf{Req.} D0.5, D1.6, D2.8, D2.9, D3.5; CT2.1 (Part 2).

\textbf{Pf.}

%dsp-level:0
\[
\sigma_e^{-1}=\sigma_e\ (\sigma_e^2=1)\implies\sigma_e^{-1}G(e)=\exp\bigl(i\alpha_e\widehat{ds}^{\,2}(e)\bigr)
\]

%dsp-level:0
\[
\log\ \text{on the}\ \mathrm{D2.8}\ \text{branch}=i\alpha_e\widehat{ds}^{\,2}(e);
\]

%dsp-level:0
\[
\qquad \tfrac1i\,\partial_{\widehat{ds}^{\,2}}\implies\alpha_e\quad(\text{constant along the edge: the read-back is per-edge})
\]

%dsp-level:0
\[
[\widehat{ds}^{\,2}]=\frac{[ds^2]}{[\mathfrak h^2]}=\frac{\mathfrak q^{-1}}{\mathfrak q^{-1}}=\mathfrak q^{0}\ [D1.6,D2.9];\qquad[\text{phase}]=\mathfrak q^0\implies[\alpha_e]=\mathfrak q^{0}. \blacksquare
\]

\textbf{Phys.} The transfer weight of an edge factors cleanly into a binary registration sign and a dimensionless geometric phase carrying the edge's own $\alpha_e$; hop-normalization makes every per-edge phase a pure number from the outset — not a later repair.

\end{theorem}

\begin{theorem}[{\texorpdfstring{EM.17 — Spectral read of the electromagnetic block}{EM.17 - Spectral read of the electromagnetic block}}]\label{tag:EM-17}
\textbf{Thm.} $\ \forall n\ge2:$ the masked force acts as a scalar on the one-dimensional EM block [CT0.14, D2.3]:

%dsp-level:0
\[
\mathbf F_M^{(n)}\big|_{\mathcal K^{(\mathrm{EM})}}=\alpha_{\mathrm{EM}}^{(n)}\cdot\mathrm{id};
\]

%dsp-level:0
\[
\qquad \alpha_{\mathrm{EM}}^{(n)}=\frac{\mathfrak T\bigl(Y_{\mathrm{asp}}\,\mathbf F_M^{(n)}\,Y_{\mathrm{asp}}\bigr)}{\mathfrak T\bigl(Y_{\mathrm{asp}}\bigr)}=\mathfrak T\bigl(Y_{\mathrm{asp}}\,\mathbf F_M^{(n)}\,Y_{\mathrm{asp}}\bigr);
\]

%dsp-level:0
\[
\qquad \mathfrak T\bigl(Y_{\mathrm{asp}}\bigr)=1.
\]

$\alpha_{\mathrm{EM}}^{(n)}$ is the force-space block-scale read at rung $n$; the edge phases $\alpha_e$ are per-edge registration data [D3.5, EM.20].

\textbf{Req.} D2.3; CT0.14 (Part 1); CT2.4 (Part 2); EM.2; EW.8 (Part 2); machine: \texttt{checks/\allowbreak{}verify\_\allowbreak{}E.py}.

\textbf{Pf.}

%dsp-level:0
\[
\dim\mathcal K^{(\mathrm{EM})}=1\ [CT0.14]\implies\mathbf F_M^{(n)}\big|_{\mathcal K^{(\mathrm{EM})}}=\alpha_{\mathrm{EM}}^{(n)}\cdot\mathrm{id}
\]

%dsp-level:0
\[
\quad [\text{every operator on a 1-dim space is a scalar};\ \alpha_{\mathrm{EM}}^{(n)}:=\text{that eigenvalue}]
\]

%dsp-level:0
\[
\mathrm{D2.3}\ \text{orthogonality}:\quad A\ne\mathrm{EM}:\ \ Y_{\mathrm{asp}}F_AY_{\mathrm{asp}}=\mathcal P^{(\mathrm{EM})}F_A\mathcal P^{(\mathrm{EM})}=0\implies Y_{\mathrm{asp}}\mathbf F_M^{(n)}Y_{\mathrm{asp}}=\alpha_{\mathrm{EM}}^{(n)}Y_{\mathrm{asp}}
\]

%dsp-level:0
\[
\mathfrak T:\quad\mathfrak T\bigl(Y_{\mathrm{asp}}\mathbf F_M^{(n)}Y_{\mathrm{asp}}\bigr)=\alpha_{\mathrm{EM}}^{(n)}\,\mathfrak T\bigl(Y_{\mathrm{asp}}\bigr)=\alpha_{\mathrm{EM}}^{(n)}\quad[\mathfrak T(Y_{\mathrm{asp}})=1,\ \mathrm{EM.2}].\ \blacksquare
\]

\textbf{Rmk (repair).} The prior formula $\mathfrak T(Y\mathbf F_M^{(n)}Y)/\mathfrak T(Y)$ with $Y=\mathcal P_M^{(2)}$ ($\mathfrak T=17$) is false as written: the rank-$2$ mask retains $J,G,W,H,D^+,\mathrm{EM}$, whose sum is not a scalar on $\mathcal K_2$. With the rank-one $Y_{\mathrm{asp}}$ the identity is exact — a 1-dimensional block has one eigenvalue and nothing else can contribute.

\textbf{Phys.} The electromagnetic block scale is read off operatorially: the single eigenvalue of the masked force on the one-dimensional EM block. A spectral formula replaces any fitting procedure. It is force-space data, distinct from the per-edge registration phases of EM.20. (Cited in \textbf{F}.)

\end{theorem}

\begin{theorem}[{\texorpdfstring{EM.18 — The colour sector does not shift the block read}{EM.18 - The colour sector does not shift the block read}}]\label{tag:EM-18}
\textbf{Thm.} $ \alpha_{\mathrm{EM}}^{(3)}=\alpha_{\mathrm{EM}}^{(2)}.$

\textbf{Req.} CT0.13, CT0.14 (Part 1); CT2.4 (Part 2); EM.17; [CT3.7 (Part 2) — forward, corroboration only].

\textbf{Pf.}

%dsp-level:0
\[
\mathrm{CT0.13}:\ \mathcal A(3)=\mathcal A(2)\cup\{\mathrm{strong}\}\implies\mathbf F_M^{(3)}=\mathbf F_M^{(2)}\oplus F_{\mathrm{strong}}\quad[CT2.4]
\]

%dsp-level:0
\[
\mathrm{CT0.14}:\ \mathcal K^{(\mathrm{strong})}\perp\mathcal K^{(\mathrm{EM})}\implies
Y_{\mathrm{asp}}\,F_{\mathrm{strong}}\,Y_{\mathrm{asp}}=\mathcal P^{(\mathrm{EM})}F_{\mathrm{strong}}\mathcal P^{(\mathrm{EM})}=0
\]

%dsp-level:0
\[
\bigl[\text{corroboration}: \mathfrak T(T_a)=0\ [CT3.7]\implies\mathfrak T(F_{\mathrm{strong}})=\textstyle\sum_af_a\mathfrak T(T_a)=0\ \text{--- independent argument}\bigr]
\]

%dsp-level:0
\[
\implies Y_{\mathrm{asp}}\mathbf F_M^{(3)}Y_{\mathrm{asp}}=Y_{\mathrm{asp}}\mathbf F_M^{(2)}Y_{\mathrm{asp}}\implies\alpha_{\mathrm{EM}}^{(3)}=\alpha_{\mathrm{EM}}^{(2)}\quad[EM.17].\ \blacksquare
\]

\textbf{Phys.} Switching colour on at rung 3 adds traceless structure orthogonal to the electromagnetic block: the block read does not run. Two disjoint arguments — block orthogonality and colour tracelessness — agree. (Cited in \textbf{F}.)

\end{theorem}

\begin{theorem}[{\texorpdfstring{EM.19 — The saturation sector does not shift the block read}{EM.19 - The saturation sector does not shift the block read}}]\label{tag:EM-19}
\textbf{Thm.} On the accessible rung-2 subspace: $ \alpha_{\mathrm{EM}}^{(n)}=\alpha_{\mathrm{EM}}^{(2)}\quad\forall n\ge4.$

\textbf{Req.} CT0.6, CT0.11, CT0.13, CT0.14 (Part 1); D4.1; EM.17, EM.18.

\textbf{Pf.}

%dsp-level:0
\[
\mathrm{CT0.13}:\ \mathcal A(n\ge4)=\mathcal A(3)\cup\{D^-\}
\]

%dsp-level:0
\[
\mathrm{D4.1}:\ F_{D^-}=\bigoplus_{m\ge4}F_m\ \text{acts only on oversubscribed rungs}\ m\ge4\ [CT0.6]
\]

%dsp-level:0
\[
\mathrm{CT0.11}:\ \mathcal K_2\ \neg\text{oversubscribed};\quad\mathrm{CT0.14}:\ \mathcal K^{(\mathrm{EM})}\perp\mathcal K^{(D^-)}\implies Y_{\mathrm{asp}}F_{D^-}Y_{\mathrm{asp}}=0
\]

%dsp-level:0
\[
\implies Y_{\mathrm{asp}}\mathbf F_M^{(n)}Y_{\mathrm{asp}}=Y_{\mathrm{asp}}\mathbf F_M^{(3)}Y_{\mathrm{asp}}\overset{\mathrm{EM.18}}{=}Y_{\mathrm{asp}}\mathbf F_M^{(2)}Y_{\mathrm{asp}}\implies\alpha_{\mathrm{EM}}^{(n)}=\alpha_{\mathrm{EM}}^{(2)}
\]

%dsp-level:0
\[
\quad [EM.17].\ \blacksquare
\]

\textbf{Phys.} The saturation lock absorbs oversubscribed information without renormalizing the accessible electromagnetic block read: $\alpha_{\mathrm{EM}}$ is rung-independent across all accessible rungs.

\end{theorem}

\begin{definition}[{\texorpdfstring{Def-EM.20 — The per-edge phase at crystallization}{Def-EM.20 - The per-edge phase at crystallization}}]\label{tag:Def-EM-20}
\textbf{Def.} For each real edge $e=(s,t)$ [A4, D0.1], the phase $\alpha_e\in\Phi$ [D0.5] is the branch index of the A4 crystallization criterion for the pair state $\rho_{st}$ at the seal [A4, D1.2]:

%dsp-level:0
\[
\alpha_e\ :=\ \alpha_{st}\ :=\ \begin{cases} 0 & \text{if } \bigl|\Pi_{\mathcal V}\rho_{st}\bigr|\le\varepsilon_q \quad \text{(aligned branch)} \\ \pi & \text{if } \bigl|\Pi_{\mathcal V}\rho_{st}\bigr|\ge 1-\varepsilon_q \quad \text{(antialigned branch)} \end{cases}
\]

— A4 declares the branch phase $\alpha_{pq}$ and both branches' non-emptiness; A7 declares the edge map $e\mapsto\alpha_e\in\{0,\pi\}$; this definition is the identification $\alpha_e=\alpha_{st}$. The branch windows are disjoint exactly when $\varepsilon_q<1-\varepsilon_q$, guaranteed by the admissible range $\varepsilon_q\in[0,\tfrac14)$ [CT0.8]. The branch read is a registration alternative [D1.1: the reality principle — registration transport acts on the reality line by $\{+1,-1\}$]: aligned $\leftrightarrow+1$, antialigned $\leftrightarrow-1$, so $e^{i\alpha_e}\in\{+1,-1\}$ is the alignment holonomy. Virtual pairs ($\varepsilon_q<\bigl|\Pi_{\mathcal V}\rho_{st}\bigr|<1-\varepsilon_q$) carry no $\alpha_e$ [A4, CT0.1; clause (3) of EM.20]

\end{definition}

\begin{theorem}[{\texorpdfstring{EM.20 — The dynamic edge phase}{EM.20 - The dynamic edge phase}}]\label{tag:EM-20}
\textbf{Thm.} With $\alpha_e$ as in Def-EM.20, under A1 ($\mathbb C^2$ qubits), A4 (two-branch crystallization; branch phase $\alpha_{pq}$; both branches populated), A5 (vacuum $\rho_\circ=I/4$; Haar draws), A6 (tick dynamics), A7 (edge phase map), CT0.9 ($\Phi=\{0,\pi\}$), D0.4 ($\varepsilon_q$), D1.6 ($\widehat{ds}^{\,2}$), D3.5 (edge weight): $(1)$ \textbf{Alphabet preservation.} $\alpha_e\in\{0,\pi\}=\Phi$ for every real edge $e$; the phase alphabet is unchanged. [A4, A7; CT0.9] $(2)$ \textbf{Phase determination.} $\alpha_e=\alpha_{st}$ is the branch index of the A4 criterion for the pair state at the moment of crystallization [A4, D1.2]: the phase is a measurement outcome — registration data read at the seal [D1.1] — $\neg$ a declared constant. $(3)$ \textbf{Virtual edges.} For $e\in\mathcal V$ [CT0.1: $\varepsilon_q<\bigl|\Pi_{\mathcal V}\rho_{st}\bigr|<1-\varepsilon_q$] the pair carries the full $\mathbb C^2\otimes\mathbb C^2$ statespace [A1]; $\alpha_e$ is undefined. The binary restriction scopes to real (crystallized) edges. $(4)$ \textbf{Edge weight.} $G(e)=\sigma_e\cdot e^{i\alpha_e\,\widehat{ds}^{\,2}(e)}$ with $\alpha_e$ per-edge. [D3.5] $(5)$ \textbf{Both values attained; distribution governed.} Both A4 branches are populated [A4 populatedness clause]; witnesses: a symmetric-dominated pair ($\bigl|\Pi_{\mathcal V}\rho\bigr|=0\Rightarrow\alpha_e=0$) and the pure singlet ($\bigl|\Pi_{\mathcal V}\rho\bigr|=1\Rightarrow\alpha_e=\pi$). Across $E_{\mathcal R}$ the distribution of $\alpha_e$ is governed by the Haar randomization [A5, A6] and the crystallization criterion [A4].

\textbf{Req.} A1; A4; A5; A6; A7; CT0.9 (Part 1); D0.4, D1.1, D1.2, D1.6, D3.5; Def-EM.20; exhaustive exact pair-state census [math; machine: \texttt{checks/\allowbreak{}verify\_\allowbreak{}F33\_\allowbreak{}dynamic\_\allowbreak{}phase.py} §A].

\textbf{Pf.}

%dsp-level:0
\[
(1):\ \mathrm{A4}:\ \alpha_{pq}\in\{0,\pi\}\ \forall(p,q)\in\mathcal R;\qquad \mathrm{A7}:\ e\in E_{\mathcal R}\mapsto\alpha_e\in\{0,\pi\};\qquad \mathrm{CT0.9}:\ \Phi=\{0,\pi\}\implies\alpha_e\in\Phi
\]

%dsp-level:0
\[
(2):\ \mathrm{A4}/\mathrm{D1.2}:\ e\in\mathcal R\iff\text{branch membership at a seal};\qquad \mathrm{Def\text{-}EM.20}:\ \alpha_e=\alpha_{st}=\text{branch index}
\]

%dsp-level:0
\[
\mathrm{D1.1}:\ \text{registration transport acts on the reality line by}\ \{\pm1\}
\]

%dsp-level:0
\[
\implies e^{i\alpha_e}
\]

%dsp-level:0
\[
= \text{alignment holonomy}\in\{+1,-1\}:\ \text{a function of the registered pair state},
\]

%dsp-level:0
\[
\ \neg\text{a constant}
\]

%dsp-level:0
\[
(3):\ \mathrm{CT0.1}/\mathrm{A4}:\ e\in\mathcal V\iff\varepsilon_q<\bigl|\Pi_{\mathcal V}\rho_{st}\bigr|<1-\varepsilon_q\iff\neg\bigl(\text{registration}\bigr);
\]

%dsp-level:0
\[
\qquad \mathrm{A1}:\ \rho_{st}\in\mathfrak H(\mathbb C^2\otimes\mathbb C^2)\ \text{retained without branch read}\implies\alpha_e\ \text{undefined}
\]

%dsp-level:0
\[
(4):\ \mathrm{D3.5}:\ G(e)=\sigma_e\exp\bigl(i\alpha_e\widehat{ds}^{\,2}(e)\bigr)\ \text{with per-edge}\ \alpha_e=\alpha_{st}\ [\mathrm{A4}];
\]

%dsp-level:0
\[
\qquad \mathrm{D1.6}:\ \widehat{ds}^{\,2}(e)\ \text{well-defined per edge}
\]

%dsp-level:0
\[
(5):\ \mathrm{A4}:\ \exists\,\rho:\ \bigl|\Pi_{\mathcal V}\rho\bigr|\le\varepsilon_q\ \wedge\ \exists\,\rho':\ \bigl|\Pi_{\mathcal V}\rho'\bigr|\ge1-\varepsilon_q;
\]

%dsp-level:0
\[
\qquad \text{witnesses}:\ \rho=\Pi_{\mathcal R}/3\ (w=0,\ \alpha=0);\quad \rho=\Pi_{\mathcal V}\ (w=1,\ \alpha=\pi)\ [\mathrm{CT0.7};\ \text{machine}]
\]

%dsp-level:0
\[
\mathrm{A5}/\mathrm{A6}:\ \text{node states Haar-drawn per tick}\ (U\notin\mathfrak L:\ \text{draws are not retained; registered reads are});
\]

%dsp-level:0
\[
\qquad \mathrm{A4}:\ \text{which pairs crystallize, and on which branch}.\ \blacksquare
\]

\textbf{Phys.} Every real edge carries its own geometric phase: the branch index of the pair's crystallization criterion — symmetric-dominated pairs register $0$, singlet-dominated pairs register $\pi$ — valued in exactly the CT0.9 alphabet, with the two branch windows disjoint on the admissible threshold range ($\varepsilon_q<1-\varepsilon_q$; CT0.8). Which branch an edge takes is measurement outcome data, governed by the Haar randomization and the criterion. All downstream phase structure — the alignment patterns and parity classes (EM.23), the CPT bridge (EM.24) — reads this per-edge datum.

\end{theorem}

\begin{theorem}[{\texorpdfstring{EM.21 — Dimensional status of the per-edge phase}{EM.21 - Dimensional status of the per-edge phase}}]\label{tag:EM-21}
\textbf{Thm.} $ \ [\alpha_e]=\mathfrak q^{0};\qquad \mathfrak h\mapsto\lambda\mathfrak h:\ \ \{0,\pi\}\mapsto\{0,\pi\};$ i.e. the phase alphabet is a derived pure-number set in the sense of CT2.11, and each $\alpha_e$ is an $\mathfrak h$-free registration value [EM.20(2)].

\textbf{Req.} D1.6; CT2.11 (Part 2); EM.16, EM.20.

\textbf{Pf.}

%dsp-level:0
\[
[\widehat{ds}^{\,2}]=\mathfrak q^0\ [EM.16];\qquad[\alpha_e\widehat{ds}^{\,2}]=[\text{phase}]=\mathfrak q^0\implies[\alpha_e]=\mathfrak q^0
\]

%dsp-level:0
\[
\mathfrak h\mapsto\lambda\mathfrak h:\ \ ds^2\mapsto\lambda^2ds^2,\ \mathfrak h^2\mapsto\lambda^2\mathfrak h^2\implies\widehat{ds}^{\,2}\ \text{invariant};
\]

%dsp-level:0
\[
\qquad \text{alphabet fixed by}\ \mathfrak h\text{-free data}\ [CT0.9,\ EM.20(1)]\implies\alpha_e\in\{0,\pi\}\ \text{in every calibration}
\]

%dsp-level:0
\[
\mathrm{CT2.11}\ \text{criterion}:\ \text{dimensionless}\ \wedge\ \text{rescaling-invariant}\implies\text{derived pure-number alphabet}. \blacksquare
\]

\textbf{Phys.} The phase alphabet $\{0,\pi\}$ is calibration-independent: it survives the rescaling of the single free scale, marking it as a genuine prediction rather than a unit convention. Individual phase values are $\mathfrak h$-free too — binary registration outcomes.

\end{theorem}

\phantomsection
\section[§E.VI — Parity and CPT]{\texorpdfstring{§E.VI — Parity and CPT}{SE.VI - Parity and CPT}}

\begin{definition}[{\texorpdfstring{Def-EM.21 — Alignment pattern}{Def-EM.21 - Alignment pattern}}]\label{tag:Def-EM-21}
\textbf{Def.} For a subnetwork $\Sigma$ [D2.13] with real edge set $E_{\mathcal R}(\Sigma):=\{e\in E_{\mathcal R}:\ e\ \text{touches}\ \Sigma\}$ ($|\Sigma|=\#E_{\mathcal R}(\Sigma)$, D2.13), the alignment pattern is

%dsp-level:0
\[
\mathcal P(\Sigma)\ :=\ \bigl\{\,e\in E_{\mathcal R}(\Sigma):\ \alpha_e=\pi\,\bigr\}
\]

— the set of antialigned edges [Def-EM.20]. $\mathcal P(\Sigma)\subseteq E_{\mathcal R}(\Sigma)$; on a chain $c$ [D1.4]: $\mathcal P(c):=\mathcal P(\Sigma)\cap c$; on a configuration $C$ [D0.10]: $\mathcal P(C):=\{e\in E(C):\ \alpha_e=\pi\}$.

\end{definition}

\begin{definition}[{\texorpdfstring{Def-EM.22 — Mirror operation}{Def-EM.22 - Mirror operation}}]\label{tag:Def-EM-22}
\textbf{Def.} A mirror $M$ on $\Sigma$ is an automorphism of the underlying graph of $\Sigma$ that reverses orientation [D2.6 $\omega$-type: $s\leftrightarrow t$]. $M$ acts on the alignment pattern:

%dsp-level:0
\[
M(\mathcal P(\Sigma))\ :=\ \bigl\{M(e):\ e\in\mathcal P(\Sigma)\bigr\}
\]

where $\Sigma$ admits such an automorphism (witnesses: path and cycle reflections — machine, \texttt{checks/\allowbreak{}verify\_\allowbreak{}F33\_\allowbreak{}dynamic\_\allowbreak{}phase.py} §B).

\end{definition}

\begin{theorem}[{\texorpdfstring{EM.23 — Parity classification}{EM.23 - Parity classification}}]\label{tag:EM-23}
\textbf{Thm.} Under Def-EM.21 and Def-EM.22, for a subnetwork $\Sigma$ with mirror $M$: $(1)$ \textbf{Parity-preserving.} $\Sigma$ is parity-even $\iff M(\mathcal P(\Sigma))=\mathcal P(\Sigma)$: the alignment pattern is invariant under the mirror. $(2)$ \textbf{Parity-reversing.} $\Sigma$ is parity-odd $\iff M(\mathcal P(\Sigma))=E_{\mathcal R}(\Sigma)\setminus\mathcal P(\Sigma)$: the alignment pattern inverts under the mirror. $(3)$ \textbf{Distinguishability.} Parity-even and parity-odd subnetworks are distinguishable $\iff\mathcal P(\Sigma)\ne E_{\mathcal R}(\Sigma)\setminus\mathcal P(\Sigma)$ — the alignment pattern is not self-complementary — which holds exactly when $E_{\mathcal R}(\Sigma)\ne\varnothing$: at least one real edge ("there is something" — CT2.6(i) as the definitional precondition it is, D2.11). $(4)$ \textbf{Degrees of freedom.} The number of distinct alignment patterns on $\Sigma$ is $2^{|E_{\mathcal R}(\Sigma)|}$: each real edge carries an independent binary phase [CT0.9, EM.20(5)]. Each pattern is a distinct parity configuration. The uniform assignments ($\mathcal P=\varnothing$ and $\mathcal P=E_{\mathcal R}$) are two points of this space; any single fixed pattern occupies one point.

\textbf{Req.} Def-EM.20, Def-EM.21, Def-EM.22; A4; D2.11, D2.13; CT0.9 (Part 1); EM.20; independent-register enumeration, orbit counting [math; machine: \texttt{checks/\allowbreak{}verify\_\allowbreak{}F33\_\allowbreak{}dynamic\_\allowbreak{}phase.py} §B].

\textbf{Pf.}

%dsp-level:0
\[
(1),(2):\ \mathrm{Def\text{-}EM.22}:\ M\ \text{permutes}\ E_{\mathcal R}(\Sigma)\implies M(\mathcal P)\subseteq E_{\mathcal R}(\Sigma);
\]

%dsp-level:0
\[
\qquad \text{even}:=\text{fixed point of the induced action};\quad\text{odd}:=\text{complement image}
\]

%dsp-level:0
\[
(3):\ e\in E_{\mathcal R}(\Sigma)\ne\varnothing:\ \ e\in\mathcal P\implies e\notin E_{\mathcal R}\setminus\mathcal P\implies\mathcal P\ne E_{\mathcal R}\setminus\mathcal P;
\]

%dsp-level:0
\[
\qquad E_{\mathcal R}=\varnothing:\ \mathcal P=\varnothing=E_{\mathcal R}\setminus\mathcal P
\]

%dsp-level:0
\[
\text{even}\wedge\text{odd jointly}\iff M(\mathcal P)=\mathcal P=E_{\mathcal R}\setminus\mathcal P\iff E_{\mathcal R}=\varnothing\implies\text{distinguishable}\iff E_{\mathcal R}(\Sigma)\ne\varnothing
\]

%dsp-level:0
\[
(4):\ \mathrm{CT0.9}/\mathrm{EM.20}(5):\ \text{each}\ e
\]

%dsp-level:0
\[
\in E_{\mathcal R}(\Sigma)\ \text{carries an independent}\ \alpha_e\in\{0,\pi\}\implies\#\{\mathcal P(\Sigma)\}
\]

%dsp-level:0
\[
= 2^{|E_{\mathcal R}(\Sigma)|}\ [\text{power-set count}]
\]

%dsp-level:0
\[
\text{uniform assignments}:\ \alpha\equiv0\iff\mathcal P=\varnothing;\quad\alpha\equiv\pi\iff\mathcal P=E_{\mathcal R}:\ \text{two of the}\ 2^{|E_{\mathcal R}|}.\ \blacksquare
\]

\textbf{Phys.} Parity is read off the ledger, not added to it: mirror-invariance of the antialigned set classifies every subnetwork even, odd, or unclassified, and the classification has teeth exactly when there is at least one real edge. The dynamic phase gives the pattern space its full $2^{|E_{\mathcal R}|}$ cardinality — uniform configurations are two points of it.

\end{theorem}

\begin{theorem}[{\texorpdfstring{EM.24 — The CPT bridge}{EM.24 - The CPT bridge}}]\label{tag:EM-24}
\textbf{Thm.} On the configuration space $\mathfrak C(\mathcal R)$ [D0.10]: $(1)$ \textbf{C (charge conjugation).} $q_{\mathrm{EM}}=\tfrac12(\sigma_s-\sigma_t)\in\{-1,0,+1\}$ [D3.4, EM.9]; charge conjugation $C:\sigma_e\mapsto-\sigma_e\ \forall e\in E(C)$ gives $q\mapsto-q$, $f_v\mapsto-f_v$, $\aleph_v$ invariant [D3.3]; the sign-level realization of D2.6's $\omega$ on end and core reads [EM.9a(2)]. Derived. $(2)$ \textbf{P (parity).} The alignment pattern $\mathcal P(\Sigma)$ [Def-EM.21]; parity reversal $P:\mathcal P\mapsto E_{\mathcal R}\setminus\mathcal P$ — equivalently $\alpha_e\mapsto\pi-\alpha_e$ per edge — the inversion locus of EM.23(2). Derived from the dynamic phase [Def-EM.20, EM.20(1)]. $(3)$ \textbf{T (time reversal).} The Janus reflection $\Theta_{\mathfrak L}$ [D2.7]: $w(\Theta_{\mathfrak L}C)=w(C)^\ast$ [CT2.8(1)]. Derived. $(4)$ \textbf{CPT composition.} $\mathrm{CPT}:=C\circ P\circ T$: well-defined on $\mathfrak C(\mathcal R)$; each component an involution ($C^2=P^2=T^2=\mathrm{id}$); each component derived from the corpus — no external input. $(5)$ \textbf{CPT invariance.} The exact phase identity and its consequences:

%dsp-level:0
\[
w(\mathrm{CPT}\,C)=(-1)^{|E(C)|}\,e^{i\pi\Theta(C)}\,w(C),\qquad\Theta(C):=\!\!\sum_{e\in E(C)}\!\!\widehat{ds}^{\,2}(e)\in\mathbb R;
\]

%dsp-level:0
\[
\bigl|w(\mathrm{CPT}\,C)\bigr|^2=\bigl|w(C)\bigr|^2\qquad\forall C\ \ (\text{unconditional});
\]

%dsp-level:0
\[
\text{on-shell}\ [\widehat{ds}^{\,2}(e)\in\mathbb Z\ \forall e]:\ \ w(C)\in\mathbb R\ \ \wedge\ \ \bigl(|E(C)|+\Theta(C)\equiv0\ (\mathrm{mod}\ 2)\bigr)\implies w(\mathrm{CPT}\,C)=w(C)^\ast .
\]

\textbf{Req.} A8; D0.5, D0.10; D1.6, D2.6, D2.7; D3.3, D3.4, D3.5, D3.6; CT0.9 (Part 1); CT2.8 (Part 2); Def-EM.20, Def-EM.21; EM.9, EM.9a, EM.20, EM.23; exhaustive small-model enumeration over $\mathbb Q(i)$ [math; machine: \texttt{checks/\allowbreak{}verify\_\allowbreak{}F33\_\allowbreak{}dynamic\_\allowbreak{}phase.py} §C].

\textbf{Pf.}

%dsp-level:0
\[
(1):\ \mathrm{CT0.9}:\ \sigma_e\in\{\pm1\}\implies-\sigma_e\in\{\pm1\}:\ C\ \text{closed on sign assignments};
\]

%dsp-level:0
\[
\qquad \mathrm{D3.4}:\ q\mapsto\tfrac12(-\sigma_s+\sigma_t)=-q;
\]

%dsp-level:0
\[
\qquad \mathrm{D3.3}:\ f_v\mapsto-f_v,\ \aleph_v=|f_v|\ \text{invariant}
\]

%dsp-level:0
\[
\mathrm{D2.6}/\mathrm{EM.9a}(2):\ \omega\ \text{maps charge}\mapsto-\text{charge on the end and core reads}:\ C=\text{sign-level realization}
\]

%dsp-level:0
\[
(2):\ \mathrm{EM.20}(1):\ \alpha_e\in\{0,\pi\}\implies\pi-\alpha_e\in\{0,\pi\}:\ P\ \text{closed on patterns};
\]

%dsp-level:0
\[
\qquad \mathrm{EM.23}(2):\ \text{the odd locus is}\ M(\mathcal P)=E_{\mathcal R}\setminus\mathcal P
\]

%dsp-level:0
\[
(3):\ \mathrm{D2.7}:\ \Theta_{\mathfrak L}\ \text{reverses the ledger order},\ J^+\leftrightarrow J^-,\ G(e)\mapsto G(e)^\ast;
\]

%dsp-level:0
\[
\qquad \mathrm{CT2.8}(1):\ w(\Theta_{\mathfrak L}C)=w(C)^\ast;
\]

%dsp-level:0
\[
\qquad \Theta_{\mathfrak L}^2=\mathrm{id}\ [\text{double reversal}]
\]

%dsp-level:0
\[
(4):\ (1),(2),(3)\ \text{map}\ \mathfrak C(\mathcal R)\to\mathfrak C(\mathcal R)\ [\mathrm{D0.10}];
\]

%dsp-level:0
\[
\qquad C^2:\ \sigma\mapsto-\sigma\mapsto\sigma;\quad P^2:\ \alpha\mapsto\pi-\alpha\mapsto\alpha;\quad T^2:\ \mathrm{id}\ [(3)]
\]

%dsp-level:0
\[
(5):\ w(C)=\mathcal J(C)\prod_{e}\sigma_e\,e^{i\alpha_e\widehat{ds}^{\,2}(e)}\cdot(-1)^{\aleph(\sigma,C)}=:\mathcal J\,S\,\Gamma\,A\quad[\mathrm{A8},\ \mathrm{D3.5},\ \mathrm{D3.6}]
\]

%dsp-level:0
\[
T:\ \Gamma\mapsto\bar\Gamma\ (\text{per-edge conjugation}),
\]

%dsp-level:0
\[
\ \ \mathcal J,
\]

%dsp-level:0
\[
S,
\]

%dsp-level:0
\[
A\in\mathbb R\ [\mathrm{D2.7},\ \mathrm{D3.6};\ \aleph\in\mathbb Z]
\]

%dsp-level:0
\[
\implies w(\Theta_{\mathfrak L}C)\qquad = w(C)^\ast
\]

%dsp-level:0
\[
P\circ T:\ \alpha_e
\]

%dsp-level:0
\[
\mapsto -\alpha_e\mapsto\pi+\alpha_e\ (\mathrm{mod}\ 2\pi):\ \ e^{-i\alpha_e\widehat{ds}^{\,2}}
\]

%dsp-level:0
\[
\mapsto e^{i(\pi+\alpha_e)\widehat{ds}^{\,2}}
\]

%dsp-level:0
\[
= e^{i\pi\widehat{ds}^{\,2}}\,
\]

%dsp-level:0
\[
e^{i\alpha_e\widehat{ds}^{\,2}}\implies\bar\Gamma
\]

%dsp-level:0
\[
\mapsto e^{i\pi\Theta}\,
\]

%dsp-level:0
\[
\Gamma
\]

%dsp-level:0
\[
\implies w(PT\,C)
\]

%dsp-level:0
\[
= e^{i\pi\Theta}\,\qquad w(C)
\]

%dsp-level:0
\[
C:\ S\mapsto(-1)^{|E|}S;\quad A\ \text{invariant}\ [\aleph_v=|-f_v|];\quad\mathcal J\ \text{invariant}\implies w(\mathrm{CPT}\,C)=(-1)^{|E|}\,e^{i\pi\Theta}\,w(C)
\]

%dsp-level:0
\[
\Theta\in\mathbb R\ [\mathrm{D1.6}:\ ds^2\ \text{real}]\implies\bigl|(-1)^{|E|}e^{i\pi\Theta}\bigr|=1\implies|w(\mathrm{CPT}\,C)|^2=|w(C)|^2
\]

%dsp-level:0
\[
\text{on-shell}:\ \widehat{ds}^{\,2}\in\mathbb Z\implies e^{i\alpha_e\widehat{ds}^{\,2}}=(\pm1)^{\widehat{ds}^{\,2}}\in\{\pm1\}\implies\Gamma\in\{\pm1\}\implies w\in\mathbb R:\ w^\ast=w;
\]

%dsp-level:0
\[
\qquad |E|+\Theta\equiv0\ (2)\implies(-1)^{|E|}e^{i\pi\Theta}=1
\]

%dsp-level:0
\[
\implies w(\mathrm{CPT}\,C)=w(C)=w(C)^\ast .\ \blacksquare
\]

\textbf{Phys.} Charge conjugation, parity, and time reversal each exist as derived operations on the ledger's configuration space: sign flip on every edge, complement of the antialigned set, Janus reflection of the record. Their composition leaves every configuration probability $|w|^2$ invariant — the corpus's own CPT theorem — with the conjugate-weight identity exact on the on-shell even locus.

\textbf{Rmk (clause-(5) scope).} $|w|^2$-invariance is unconditional; $w(\mathrm{CPT}\,C)=w(C)^\ast$ holds on the stated on-shell even locus.

\end{theorem}

\phantomsection
\section[§E.VII — Spectrum and Completion]{\texorpdfstring{§E.VII — Spectrum and Completion}{SE.VII - Spectrum and Completion}}

\begin{theorem}[{\texorpdfstring{EM.22 — The electromagnetic spectrum}{EM.22 - The electromagnetic spectrum}}]\label{tag:EM-22}
\textbf{Thm.}

%dsp-level:0
\[
\begin{array}{c|c|c|c}
\text{object} & \text{type} & \text{mass/charge} & \text{data}\\ \hline
\text{photon} & \text{trivial chain}\subseteq V_2 & m^2=0;\ q_{\mathrm{EM}}=0 & N_{\mathrm{pol}}=2\\
\text{field excitation} & \text{transition vertex} & \aleph_v=1 & f_v=\pm1\\
\text{charged chain end} & \text{chain terminus} & q_{\mathrm{EM}}\in\{+1,-1\} & \text{binary sign}
\end{array}
\]

%dsp-level:0
\[
\neg\exists\ \text{massive abelian gauge boson};\qquad\text{list exhaustive}:\ \dim(\text{abelian block})=1\ [EM.3].
\]

\textbf{Req.} D3.4, D3.7; EM.3, EM.5, EM.6, EM.8, EM.9, EM.10, EM.12; EW.22, EW.28 (Part 2).

\textbf{Pf.}

%dsp-level:0
\[
\text{photon}: \mathrm{EM.5}\ (\text{trivial}\ V_2);\ \ \mathrm{EM.6}\ (m^2=0);\ \ \mathrm{EM.8}\ (2\ \text{pol})
\]

%dsp-level:0
\[
q_\gamma=0:\ \ \text{trivial}\implies\neg\text{transitions}\ [EM.12]\implies\sigma_{\leftarrow}=\sigma_{\rightarrow}\ \forall v\implies\sigma_{s(c)}=\sigma_{t(c)}\implies q=0\ [\mathrm{D3.4}]
\]

%dsp-level:0
\[
\text{excitations}: \mathrm{EM.12}\ (\aleph_v=1,\ f_v=\pm1);\qquad\text{ends}: \mathrm{EM.9},\ \mathrm{EM.10}
\]

%dsp-level:0
\[
\neg\exists\ \text{massive abelian}: \ \mathrm{EW.22}:\ \ker M^2
\]

%dsp-level:0
\[
= \text{1-dim abelian direction}\ A;
\]

%dsp-level:0
\[
\quad\text{massive neutral}
\]

%dsp-level:0
\[
= Z\in\text{weak part}\ [EW.28]
\]

%dsp-level:0
\[
\text{exhaustive}: 1\ \text{generator}\ [EM.3]\implies\{\text{transit mode}=\gamma,\ \text{curvature data}=\aleph_v,\ \text{sources}=\text{ends}\}. \blacksquare
\]

\textbf{Phys.} The complete electromagnetic bestiary: one massless neutral two-polarization transit mode, sign-field vertex excitations, and quantized chain-end charges. The photon’s own charge vanishes identically — a null chain carries no sign flips.

\end{theorem}

\begin{theorem}[{\texorpdfstring{EM.26 — Cross-checks}{EM.26 - Cross-checks}}]\label{tag:EM-26}
\textbf{Thm.} Four agreements between derivations with pairwise-disjoint inputs:

%dsp-level:0
\[
(1)\ \{-1,0,+1\}:\ \ \mathrm{EM.9}\ [\Leftarrow\mathrm{CT0.9}]\ =\ \mathrm{EW.36}\ [\Leftarrow\mathrm{EW.32}];
\]

%dsp-level:0
\[
(2)\ m_\gamma^2=0:\ \ \mathrm{EM.6(ii)}\ [\Leftarrow\mathrm{D1.7},\mathrm{EW.18}]\ =\ \mathrm{EW.22}\ [\Leftarrow\mathrm{EW.33}];
\]

%dsp-level:0
\[
(3)\ \alpha_{\mathrm{EM}}^{(3)}=\alpha_{\mathrm{EM}}^{(2)}:\ \ \mathrm{EM.18}\ \text{step 2}\ [\Leftarrow\mathrm{CT0.14}]\ =\ \mathrm{EM.18}\ \text{step 3}\ [\Leftarrow\mathrm{CT3.7}];
\]

%dsp-level:0
\[
(4)\ g^{\mu\nu}T_{\mu\nu}=0:\ \ \mathrm{EM.15}\ \text{(explicit variation)}\ =\ \mathrm{EM.6}\ \text{(massless transit}\implies\text{traceless stress)}.
\]

\textbf{Req.} CT0.9, CT0.14 (Part 1); CT3.7 (\textbf{F}, corroboration); EM.6, EM.9, EM.15, EM.18; EW.18, EW.22, EW.32, EW.33 (Part 2).

\textbf{Pf.}

%dsp-level:0
\[
\text{input pairs}: (1)\ \mathrm{CT0.9}\ \text{vs}\ \mathrm{EW.32};
\]

%dsp-level:0
\[
\quad (2)\ \{\mathrm{D1.7},\mathrm{EW.18}\}\ \text{vs}\ \mathrm{EW.33};
\]

%dsp-level:0
\[
\quad (3)\ \mathrm{CT0.14}\ \text{vs}\ \mathrm{CT3.7};
\]

%dsp-level:0
\[
\quad (4)\ \text{variation vs pole location}
\]

%dsp-level:0
\[
\text{shared intermediate results}=\varnothing\ \text{in each pair}\implies\text{agreements}=\text{genuine consistency tests};
\]

%dsp-level:0
\[
\qquad \forall4:\ \text{hold}\ [\text{cited tags}]. \blacksquare
\]

\textbf{Phys.} The sector cross-validates at four joints against the combinatorics of \textbf{B}, the resolvent apparatus of \textbf{C}, and the seed calculus of \textbf{D} — each check by two routes sharing no intermediate result.

\end{theorem}

\begin{theorem}[{\texorpdfstring{EM.25 — What the corpus does and does not fix about the abelian sector}{EM.25 - What the corpus does and does not fix about the abelian sector}}]\label{tag:EM-25}
\textbf{Thm.} Fixed with zero freedom:

%dsp-level:0
\[
(1)\ q\in\{-1,0,+1\}\ [\mathrm{EM.9},\mathrm{EM.10},\mathrm{EW.36}];\quad
(2)\ m_\gamma^2=0\ [\mathrm{EM.6}];\quad
(3)\ N_{\mathrm{pol}}=2\ [\mathrm{EM.8}];
\]

%dsp-level:0
\[
(4)\ \alpha_e\in\{0,\pi\}\ \text{per real edge}\ [\mathrm{EM.20}];
\]

%dsp-level:0
\[
\quad (5)\ \alpha_{\mathrm{EM}}^{(n)}=\alpha_{\mathrm{EM}}^{(2)}\ \forall n\ \ (\text{block read runs at no accessible rung})\ [\mathrm{EM.18},\mathrm{EM.19}];
\]

%dsp-level:0
\[
\quad (6)\ e=g\sin\theta_W=g'\cos\theta_W\ [\mathrm{EW.35}].
\]

$\neg$Fixed:

%dsp-level:0
\[
(7)\ \text{the numerical value of}\ e:\quad
\neg\exists\ \text{theorem asserting a value of}\ e;
\]

%dsp-level:0
\[
\qquad e
\]

%dsp-level:0
\[
= e(g,g'),
\]

%dsp-level:0
\[
\ (g,g')\ \text{free}\ [\mathrm{P4}]\implies\text{family of valid configurations with distinct}\ e,
\]

%dsp-level:0
\[
\ \text{all satisfying}\ (1)\text{--}(6);
\]

%dsp-level:0
\[
\text{determining}\ e\ \text{requires one further datum}\ (\theta_W\ \text{or}\ g)
\]

%dsp-level:0
\[
= \text{calibration step}\ [\mathrm{P5}];
\]

%dsp-level:0
\[
\ \text{then}\ (6)\implies\text{rest}
\]

%dsp-level:0
\[
= \text{predictions};
\]

%dsp-level:0
\[
(8)\ \text{the per-edge values}\ \alpha_e:\ \text{registration outcomes governed by}\ \mathrm{A4}/\mathrm{A5}/\mathrm{A6}\ [\mathrm{EM.20}(2),(5)],\ \neg\text{parameters};
\]

%dsp-level:0
\[
\qquad \alpha_{\mathrm{EM}}^{(2)}\ \text{as a number}:\ \text{operator read}\ [\mathrm{EM.17}],\ \neg\exists\ \text{theorem fixing its value}.
\]

Separation and closure:

%dsp-level:0
\[
(9)\ \neg\bigl(\alpha\text{-data determines}\ e\bigr);
\]

%dsp-level:0
\[
\qquad \neg\bigl(e\ \text{determines}\ \alpha\text{-data}\bigr);
\]

%dsp-level:0
\[
\qquad \neg\exists\ \text{corpus equation relating them}\ [g,g'\ \text{appear nowhere in EM.16--EM.21}];
\]

%dsp-level:0
\[
(10)\ \forall\ \text{statements of}\ \mathbf E:\quad\vdash\ \mathrm{A1\text{–}A9}\cup\mathrm{C1\text{–}C12}\cup\mathrm{D0.1\text{–}D4.3}\cup\{\mathbf B,\mathbf C,\mathbf D\ \text{theorems}\};
\]

%dsp-level:0
\[
\qquad \neg\exists\ \text{external EM law, parameter, or measured value}.
\]

\textbf{Req.} P4, P5; EM.6, EM.8, EM.9, EM.10, EM.16, EM.17, EM.18, EM.19, EM.20, EM.21, EM.23, EM.24; EW.24, EW.35, EW.36 (Part 2); citation closure [machine: \texttt{checks/\allowbreak{}verify\_\allowbreak{}corpus\_\allowbreak{}graph.py}].

\textbf{Pf.}

%dsp-level:0
\[
(1)\text{--}(6):\ \text{cited proofs use only A}\cup\text{derived results};\ \neg\exists\ \text{parameter introduced};
\]

%dsp-level:0
\[
\qquad (4):\ \mathrm{EM.20}(1);
\]

%dsp-level:0
\[
\qquad (5):\ \mathrm{EM.18},\ \mathrm{EM.19}
\]

%dsp-level:0
\[
(7):\ \mathrm{P4}:\ (g,g')
\]

%dsp-level:0
\[
\in (\mathbb R_{>0})^2\ \text{unconstrained}\implies\forall(g,g'):\ \text{configuration with}\ e
\]

%dsp-level:0
\[
= gg'/\sqrt{g^2+g'^2}\ \text{satisfying}\ (1)\text{--}(6);
\]

%dsp-level:0
\[
\qquad \tan\theta_W=g'/g\ [\mathrm{EW.24}]
\]

%dsp-level:0
\[
\implies\neg\exists\ \text{function of the axioms alone}\mapsto\text{unique}\ e;\qquad\mathrm{P5}:\ \text{one datum}\ +(6)\implies\text{predictions}
\]

%dsp-level:0
\[
(8):\ \mathrm{EM.20}(2),(5):\ \alpha_e=\text{A4 branch index at the seal}\ [\mathrm{Def\text{-}EM.20}];
\]

%dsp-level:0
\[
\qquad \mathrm{EM.17}:\ \alpha_{\mathrm{EM}}^{(n)}=\text{eigenvalue of}\ \mathbf F_M^{(n)}\big|_{\mathcal K^{(\mathrm{EM})}}:\ \text{operator datum},\ \neg\text{a fixed pure number}
\]

%dsp-level:0
\[
(9):\ \mathrm{EM.16}\text{--}\mathrm{EM.21}\ \text{derivations contain}\ \neg\{g,g'\};
\]

%dsp-level:0
\[
\qquad \mathrm{EW.35}:\ e=gg'/\sqrt{g^2+g'^2};
\]

%dsp-level:0
\[
\qquad (g,g')\mapsto(g_2,g_2'):\ e\mapsto e_2\ne e\ \wedge\ \alpha\text{-data unchanged}\implies\text{independent variation}
\]

%dsp-level:0
\[
(10):\ \text{dependency trace}:\ \text{activation}\Leftarrow\mathrm{CT0.13}+\mathrm{D3.3}\ [\mathrm{EM.1}\text{--}\mathrm{EM.4}];
\]

%dsp-level:0
\[
\quad\gamma:\ \text{trivial}\ V_2\ \text{chain},
\]

%dsp-level:0
\[
\ m
\]

%dsp-level:0
\[
= 0,
\]

%dsp-level:0
\[
\ 2\ \text{pol},
\]

%dsp-level:0
\[
\ \mathsf E_{\mathrm{tr}}\ \text{only}\ [\mathrm{EM.5}\text{--}\mathrm{EM.8}]
\]

%dsp-level:0
\[
\text{charge}:\ \text{end-difference},
\]

%dsp-level:0
\[
\ \mathbb Z\text{-quantized by the binary phase},
\]

%dsp-level:0
\[
\ Q_{\mathrm{EM}}
\]

%dsp-level:0
\[
= 0\ \text{by seal pairing}\ [\mathrm{EM.9}\text{--}\mathrm{EM.11}];
\]

%dsp-level:0
\[
\qquad f_v;\ \text{sign structures independent};\ \text{gauge}=\text{offset redundancy}\ [\mathrm{EM.12}\text{--}\mathrm{EM.14}]
\]

%dsp-level:0
\[
T_{\mu\nu}=F\text{-bilinear},\ \text{trace}\ 0\ [\mathrm{EM.15}];
\]

%dsp-level:0
\[
\qquad \alpha_e:\ \text{per-edge registration read},\ \text{alphabet}\ \{0,\pi\},\ \mathfrak q^0\ [\mathrm{EM.16}\text{--}\mathrm{EM.21}]
\]

%dsp-level:0
\[
\text{parity}:\ \text{pattern data},\ 2^{|E_{\mathcal R}|}\ \text{classes}\ [\mathrm{EM.23}];
\]

%dsp-level:0
\[
\qquad \mathrm{CPT}:\ \text{derived components},\ |w|^2\text{-invariance}\ [\mathrm{EM.24}];
\]

%dsp-level:0
\[
\qquad \text{spectrum exhaustive}\ [\mathrm{EM.22}];\quad4\ \text{cross-checks}\ [\mathrm{EM.26}].\ \blacksquare
\]

\textbf{Phys.} The corpus states exactly which abelian quantities are derived — alphabet, masslessness, polarizations, the per-edge phase alphabet, block-read stability, and the identity $e=g\sin\theta_W$ — and exactly which data stay open: the numerical coupling (configuration data), the per-edge phase values (registration outcomes), and the block-scale number (operator read). It makes no claim to derive a fine-structure constant, and no equation ties phase data to the coupling: an honesty clause that is itself a theorem about the axiom system. Closure audit: every electromagnetic statement traces to A1–A9.

\end{theorem}

\phantomsection
\section[Closure exports of Part 2 §E (consumed by Part 2 §F–§G,…]{\texorpdfstring{Closure exports of Part 2 §E (consumed by Part 2 §F–§G, Parts 3–4)}{Closure exports of Part 2 SE (consumed by Part 2 SF-SG, Parts 3-4)}}

%dsp-level:0
\[
\begin{array}{lll}
\mathrm{EM.2}\ Y_{\mathrm{asp}}=\mathcal P^{(\mathrm{EM})},\ \mathfrak T=1;\ Y_{\mathrm{car}}=\tfrac12\sigma_0 &
\mathrm{EM.5}\ \gamma_\gamma,\ \gamma_g\in\mathsf{Trivial}\ (V_2/V_3) &
\mathrm{EM.6}\ m_\gamma^2=0\ (\times2\ \text{routes})\\
\mathrm{EM.8}\ N_{\mathrm{pol}}=2 &
\mathrm{EM.9/10}\ q_{\mathrm{EM}}\in\{-1,0,+1\}\subset\mathbb Z;\ \mathrm{EM.9b}:\ \textstyle\sum_if_{v_i}=\tfrac12(\sigma_{e_k}-\sigma_{e_1}) &
\mathrm{EM.11}\ Q_{\mathrm{EM}}=0\\
\mathrm{EM.12/13}\ f_v;\ (-1)^\aleph\ \text{dynamics} &
\mathrm{EM.14}\ A_\mu\sim A_\mu+\nabla_\mu\chi &
\mathrm{EM.15}\ T^{(\mathrm{EM})}_{\mu\nu};\ \text{traceless}\\
\mathrm{EM.16}\ G(e)=\sigma_e e^{i\alpha_e\widehat{ds}^{\,2}} &
\mathrm{EM.17}\ \alpha_{\mathrm{EM}}^{(n)}=\mathfrak T(Y_{\mathrm{asp}}\mathbf F_MY_{\mathrm{asp}}) &
\mathrm{EM.18/19}\ \alpha_{\mathrm{EM}}^{(n)}=\alpha_{\mathrm{EM}}^{(2)}\ \forall n\\
\mathrm{Def\text{-}EM.20}/\mathrm{EM.20}\ \alpha_e=\alpha_{pq}\in\{0,\pi\}:\ \text{A4 branch, per edge} &
\mathrm{EM.21}\ [\alpha_e]=\mathfrak q^0 &
\mathrm{Def\text{-}EM.21/22}\ \mathcal P(\Sigma);\ M(\mathcal P)\\
\mathrm{EM.23}\ \text{parity classes};\ 2^{|E_{\mathcal R}|}\ \text{patterns} &
\mathrm{EM.24}\ \mathrm{CPT}=C\circ P\circ T;\ |w(\mathrm{CPT}\,C)|^2=|w(C)|^2 &
\mathrm{EM.22}\ \text{spectrum}\ \{\gamma,\ \aleph_v,\ \text{ends}\}\\
\mathrm{EM.25}\ e=\text{P4 data};\ \text{self-containment} &
\mathrm{EM.26}\ 4\ \text{cross-checks} &
\end{array}
\]

$\blacksquare$

\phantomsection
\section[§F — Colour Sector]{\texorpdfstring{§F — Colour Sector}{SF - Colour Sector}}

\textit{Dep.: Part 1; Part 2 §C. This section may be read before §D, §E: its two borrowed results — colour masslessness (EW.26) and coefficient stability (EM.18) — are restated with full arguments at use. Notation: $\mathcal H^{(n)}_{\mathrm{slot}}$ [D3.1] $\ne\mathfrak H(\mathbb C^4)$ [A1] — slot space and pair operator space are never identified; $T_a$, $f_{abc}$: colour generators and structure constants in the Hermitian convention [C2]. Machine: \texttt{checks/\allowbreak{}verify\_\allowbreak{}F.py} — exact arithmetic in $\mathbb Q(\sqrt3,i)$; CT3.22a/CT3.17a: \texttt{checks/\allowbreak{}verify\_\allowbreak{}merge\_\allowbreak{}tierA.py}; CT3.10a: \texttt{checks/\allowbreak{}verify\_\allowbreak{}promotion\_\allowbreak{}F\_\allowbreak{}colour.py}.}

\phantomsection
\section[§F.I — Slot Structure]{\texorpdfstring{§F.I — Slot Structure}{SF.I - Slot Structure}}

\begin{theorem}[{\texorpdfstring{CT3.20 — The number of colours is forced}{CT3.20 - The number of colours is forced}}]\label{tag:CT3-20}
\textbf{Thm.}

%dsp-level:0
\[
\text{admissible}(v)\iff n(v)\le3;\qquad
\max_{\text{admissible}}\mathcal H^{(n)}_{\mathrm{slot}}=\mathcal H^{(3)}_{\mathrm{slot}}\cong\mathbb C^3;
\]

%dsp-level:0
\[
N_{\text{colour}}=3=\log_2C_{\mathrm{cap}}-1=\log_216-1\qquad(\text{derived from A2}+\mathrm{CT0.9};\ \neg\text{chosen}).
\]

\textbf{Req.} A2; D3.1; CT0.5, CT0.6, CT0.9, CT0.11 (Part 1).

\textbf{Pf.}

%dsp-level:0
\[
\mathrm{D3.1}:\ \text{rung-}n\ \text{vertex carries}\ n\ \text{slots}
\]

%dsp-level:0
\[
\mathrm{CT0.5}/\mathrm{CT0.6}:\ D_{\mathrm{tick}}(n)=2^{n+1},\ C_{\mathrm{cap}}=16;\qquad
\mathrm{CT0.11}:\ \text{admissible}\iff2^{n+1}\le16\iff n\le3
\]

%dsp-level:0
\[
\implies\mathcal H^{(n)}_{\mathrm{slot}}\in\{\mathbb C^0,\mathbb C^1,\mathbb C^2,\mathbb C^3\};\qquad\max=\mathbb C^3
\]

%dsp-level:0
\[
n\le\log_2C_{\mathrm{cap}}-1=3:\quad\text{inputs}=\mathrm{A2}\ (1\mathfrak q\ \text{per channel})\ +\ \mathrm{CT0.9}\ (\Phi=\{0,\pi\});\quad\neg\exists\ \text{further input}.\ \blacksquare
\]

\textbf{Phys.} Three is not a choice: a rung-$n$ vertex loads $2^{n+1}$ registers against a capacity of $16$, so three slots are the most any admissible vertex can carry. The number of colours is fixed by the channel quantum and the binary phase alphabet, and by nothing else.

\end{theorem}

\begin{theorem}[{\texorpdfstring{CT3.1 — Classical trivalent symmetry}{CT3.1 - Classical trivalent symmetry}}]\label{tag:CT3-1}
\textbf{Thm.} Trivalent vertex with distinguishable slots $\{s_1,s_2,s_3\}$:

%dsp-level:0
\[
\text{relabeling symmetry}=\operatorname{Sym}\{s_1,s_2,s_3\}=S_3;\qquad
|S_3|=3!=6;\qquad
S_3=\{e,(12),(13),(23),(123),(132)\}.
\]

\textbf{Req.} D3.1; bijection/group facts [math].

\textbf{Pf.}

%dsp-level:0
\[
\mathrm{D3.1}:\ \text{rung-3 vertex}:\ \text{slots}=\{s_1,s_2,s_3\}
\]

%dsp-level:0
\[
\text{relabeling}=\text{invertible},\ \text{slot-preserving map}=\text{bijection}\ \{s_1,s_2,s_3\}\to\{s_1,s_2,s_3\}
\]

%dsp-level:0
\[
\text{bijections under composition}=S_3;\qquad|S_3|=3!=6.\ \blacksquare
\]

\textbf{Phys.} The classical symmetry of a trivalent vertex is the full permutation group of its three slots — six relabelings, no more.

\end{theorem}

\begin{theorem}[{\texorpdfstring{CT3.2 — Quantum slot space}{CT3.2 - Quantum slot space}}]\label{tag:CT3-2}
\textbf{Thm.}

%dsp-level:0
\[
|\psi\rangle=\alpha_1|1\rangle+\alpha_2|2\rangle+\alpha_3|3\rangle,\ \ \alpha_i\in\mathbb C;\qquad
\langle\psi|\phi\rangle=\textstyle\sum_i\alpha_i^\ast\beta_i;\qquad
\textstyle\sum_i|\alpha_i|^2=1;
\]

%dsp-level:0
\[
\mathcal H^{(3)}_{\mathrm{slot}}=\operatorname{span}_{\mathbb C}\{|1\rangle,|2\rangle,|3\rangle\}\cong\mathbb C^3 .
\]

\textbf{Req.} A1; D3.1.

\textbf{Pf.}

%dsp-level:0
\[
\mathrm{A1}:\ \text{carrier spaces}=\text{complex Hilbert spaces}\implies\text{slot excitations admit complex superpositions}
\]

%dsp-level:0
\[
\text{span of 3 independent basis states over}\ \mathbb C=\mathbb C^3\quad[\mathrm{D3.1}].\ \blacksquare
\]

\textbf{Phys.} Superposition promotes the discrete slot label to a complex three-dimensional carrier: colour becomes a quantum number with amplitudes.

\end{theorem}

\begin{theorem}[{\texorpdfstring{CT3.3 — Classical permutations embed unitarily}{CT3.3 - Classical permutations embed unitarily}}]\label{tag:CT3-3}
\textbf{Thm.} $P_\pi|i\rangle:=|\pi(i)\rangle$, $\ (P_\pi)^i{}_j=\delta^i_{\pi(j)}$. Then

%dsp-level:0
\[
P_\pi^\dagger P_\pi=I_3\implies P_\pi\in U(3);
\]

%dsp-level:0
\[
\qquad P_\pi P_\rho=P_{\pi\circ\rho}\ \ (\text{homomorphism});
\]

%dsp-level:0
\[
\qquad P_\pi=I_3\iff\pi=e\ \ (\text{faithful});
\]

%dsp-level:0
\[
\qquad S_3\hookrightarrow U(3).
\]

\textbf{Req.} CT3.1, CT3.2; unitarity criterion [math].

\textbf{Pf.}

%dsp-level:0
\[
P_\pi:\ \text{ON basis}\mapsto\text{ON basis}\implies\langle P_\pi\psi,P_\pi\phi\rangle=\langle\psi,\phi\rangle\implies P_\pi^\dagger P_\pi=I_3\implies P_\pi\in U(3)
\]

%dsp-level:0
\[
P_\pi P_\rho|i\rangle=P_\pi|\rho(i)\rangle=|\pi(\rho(i))\rangle=P_{\pi\circ\rho}|i\rangle\implies\text{homomorphism}
\]

%dsp-level:0
\[
P_\pi=I_3\implies\forall i:\ \pi(i)=i\implies\pi=e\implies\text{injective}.\ \blacksquare
\]

\textbf{Phys.} The classical relabeling symmetry sits inside the quantum one as the six permutation matrices — a faithful discrete subgroup of $U(3)$.

\end{theorem}

\begin{theorem}[{\texorpdfstring{CT3.4 — Probability preservation gives $U(3)$}{CT3.4 - Probability preservation gives}}]\label{tag:CT3-4}
\textbf{Thm.}

%dsp-level:0
\[
\bigl\{L\ \text{linear}:\ \langle L\psi,L\phi\rangle=\langle\psi,\phi\rangle\ \ \forall\psi,\phi\bigr\}
=\bigl\{L:\ L^\dagger L=I_3\bigr\}=U(3);\qquad
S_3\subseteq U(3)\ [\mathrm{CT3.3}].
\]

\textbf{Req.} CT3.2, CT3.3; polarization/inner-product matrix form [math].

\textbf{Pf.}

%dsp-level:0
\[
\text{probabilities}=\text{moduli of inner products}\implies\text{preservation}\iff\langle L\psi,L\phi\rangle=\langle\psi,\phi\rangle\ \forall\psi,\phi
\]

%dsp-level:0
\[
\quad [\text{polarization}]
\]

%dsp-level:0
\[
\psi^\dagger L^\dagger L\phi=\psi^\dagger\phi\ \ \forall\psi,\phi\iff L^\dagger L=I_3\quad[\psi,\phi\ \text{over a basis}]
\]

%dsp-level:0
\[
\{L:\ L^\dagger L=I_3\}=:U(3)\ \ \text{(group)};\qquad S_3\hookrightarrow U(3)\ [\mathrm{CT3.3}].\ \blacksquare
\]

\textbf{Phys.} The full quantum symmetry of the slot carrier is exactly $U(3)$: the linear maps preserving every probability. The classical $S_3$ is its discrete relic.

\end{theorem}

\begin{theorem}[{\texorpdfstring{CT3.5 — The algebra $\mathfrak u(3)$ and its decomposition}{CT3.5 - The algebra and its decomposition}}]\label{tag:CT3-5}
\textbf{Thm.} [C2]

%dsp-level:0
\[
\mathfrak u(3)=\{X\in M_3(\mathbb C):\ X^\dagger=X\};\qquad
\mathfrak u(3)=\mathfrak{su}(3)\oplus\mathfrak u(1);
\]

%dsp-level:0
\[
\mathfrak{su}(3)=\{X\in\mathfrak u(3):\ \mathfrak T(X)=0\};\qquad
\mathfrak u(1)=\operatorname{span}_{\mathbb R}\{I_3\};\qquad
[\mathfrak{su}(3),\mathfrak u(1)]=0\ \ (\text{commuting ideals});
\]

%dsp-level:0
\[
\forall X\ \exists!\ (X_0,c):\ X=X_0+cI_3,\ \ \mathfrak T(X_0)=0,\ \ c=\tfrac{\mathfrak T(X)}{3}.
\]

\textbf{Req.} C2; CT3.4; trace cyclicity [linear algebra].

\textbf{Pf.}

%dsp-level:0
\[
X^\dagger=X\wedge Y^\dagger=Y\implies[X,Y]^\dagger=Y^\dagger X^\dagger-X^\dagger Y^\dagger=-[X,Y]\implies-i[X,Y]\ \text{Hermitian}:\ \text{closure}\ [\mathrm{C2}]
\]

%dsp-level:0
\[
\mathfrak T[X,Y]=\mathfrak T(XY)-\mathfrak T(YX)=0\implies\text{traceless Hermitian}=\text{subalgebra}
\]

%dsp-level:0
\[
X_0:=X-\tfrac{\mathfrak T(X)}{3}I_3:\quad\mathfrak T(X_0)=\mathfrak T(X)-\tfrac{\mathfrak T(X)}{3}\cdot3=0\implies X_0\in\mathfrak{su}(3)
\]

%dsp-level:0
\[
\text{uniqueness}:\ \ X=S+cI_3,\ \mathfrak T(S)=0\implies\mathfrak T(X)=3c\implies c=\tfrac{\mathfrak T(X)}3
\]

%dsp-level:0
\[
I_3\in Z\implies[\mathfrak u(1),\mathfrak u(3)]=0\implies\text{both summands ideals};\qquad
cI_3\ \text{traceless}\iff c=0\implies\mathfrak{su}(3)\cap\mathfrak u(1)=\{0\}.\ \blacksquare
\]

\textbf{Phys.} The slot symmetry algebra splits canonically into traceless plus centre — the template from which colour and electromagnetism are cut by the subset principle.

\end{theorem}

\phantomsection
\section[§F.II — The Rung-3 Gauge Algebra]{\texorpdfstring{§F.II — The Rung-3 Gauge Algebra}{SF.II - The Rung-3 Gauge Algebra}}

\begin{theorem}[{\texorpdfstring{CT3.6 — Rung-3 active set}{CT3.6 - Rung-3 active set}}]\label{tag:CT3-6}
\textbf{Thm.} $\ \mathcal A(3)=\mathcal A(2)\cup\{\mathrm{strong}\};\qquad\mathbf F_M^{(3)}=\mathcal P_M^{(3)}\mathbf F\,\mathcal P_M^{(3)}=\mathbf F\big|_{\mathcal K_3}.$

\textbf{Req.} CT0.13 (Part 1); CT2.2, CT2.3, CT2.4 (Part 2).

\textbf{Pf.} Active set: [CT0.13 step $n{=}3$]; projector: [CT2.2]; restriction: [CT2.3, CT2.4]. $\blacksquare$

\textbf{Phys.} At rung 3 the mask admits one new block — colour — on top of the rung-2 content.

\end{theorem}

\begin{theorem}[{\texorpdfstring{CT3.7 — The strictly new rung-3 algebra is $\mathfrak{su}(3)$}{CT3.7 - The strictly new rung-3 algebra is}}]\label{tag:CT3-7}
\textbf{Thm.} Strictly new gauge structure at rung 3 [D3.2]:

%dsp-level:0
\[
\mathfrak u(3)\big/\mathfrak u(1)\cong\mathfrak{su}(3);
\]

%dsp-level:0
\[
\qquad \pi:\mathfrak u(3)\to\mathfrak{su}(3),\ \ \pi(X)=X-\tfrac{\mathfrak T(X)}{3}I_3:\ \ \text{Lie homomorphism},\ \ \ker\pi=\operatorname{span}_{\mathbb R}\{I_3\},\ \ \text{surjective};
\]

%dsp-level:0
\[
\mathfrak T(T_a)=0\quad\forall a\ \ (\neg\exists\ \text{central component}).
\]

\textbf{Req.} D3.2, D2.12; CT0.13 (Part 1); CT3.4, CT3.5; EM.2, EW.8 (\textbf{D}/\textbf{E}, restated: centre spent at rung 2).

\textbf{Pf.}

%dsp-level:0
\[
\mathrm{CT3.4}:\ \text{rung-3 slot symmetry}=U(3)\implies\text{algebra}=\mathfrak u(3);\qquad
\mathrm{CT3.5}:\ \mathfrak u(3)=\mathfrak{su}(3)\oplus\mathfrak u(1)_{\langle I_3\rangle}
\]

%dsp-level:0
\[
\iota_2\bigl(Z(\mathfrak u(2))\bigr)=F_{\mathrm{EM}}\ \text{realized at rung}\ 2\ [\mathrm{D2.12},\ \mathrm{CT0.13},\ \mathrm{EM.2}\ \text{restated}:\ Y_{\mathrm{asp}}=\mathcal P^{(\mathrm{EM})}]
\]

%dsp-level:0
\[
\mathrm{D3.2}\ [\mathrm{D2.12}:\ \text{“realized”}=\text{abstract class in}\ \iota_m(\mathfrak u(m)),\ m<n]:\quad\iota_3\bigl(Z(\mathfrak u(3))\bigr)\cong\mathfrak u(1)\cong\iota_2\bigl(Z(\mathfrak u(2))\bigr)\implies\neg\text{strictly new}
\]

%dsp-level:0
\[
\implies\text{strictly new}=\mathfrak u(3)/\mathfrak u(1)
\]

%dsp-level:0
\[
\pi\ \text{well-defined Lie homomorphism}\ (I_3\ \text{central});
\]

%dsp-level:0
\[
\quad \ker\pi=\{X:\ X=\tfrac{\mathfrak T X}3I_3\}=\operatorname{span}_{\mathbb R}\{I_3\};\quad\text{surjective}\ [\mathrm{CT3.5}\ \text{split}]
\]

%dsp-level:0
\[
\mathfrak{su}(3)=\{\mathfrak T=0\}\implies\mathfrak T(T_a)=0\ \ \forall a.\ \blacksquare
\]

\textbf{Phys.} The new force at rung 3 is exactly the traceless slot algebra $\mathfrak su(3)$: eight generators, no central component — the ninth direction of $\mathfrak u(3)$ was already spent on electromagnetism at rung 2 by the subset principle.

\end{theorem}

\begin{theorem}[{\texorpdfstring{CT3.8 — Dimension of the colour algebra}{CT3.8 - Dimension of the colour algebra}}]\label{tag:CT3-8}
\textbf{Thm.} $\ \dim_{\mathbb R}\mathfrak u(3)=3+3\cdot2=9;\qquad\mathfrak T(X)=0:\ 1\ \text{nontrivial real constraint}\ (\mathfrak T(I_3)=3\ne0);\qquad\dim_{\mathbb R}\mathfrak{su}(3)=9-1=8=3^2-1.$

\textbf{Req.} CT3.5; Hermitian-matrix parameter count [linear algebra].

\textbf{Pf.}

%dsp-level:0
\[
X^\dagger=X:\ 3\ \text{real diagonal}\ +\ 3\ \text{complex off-diagonal}\ (2\ \text{real each})\implies\dim_{\mathbb R}\mathfrak u(3)=9
\]

%dsp-level:0
\[
\mathfrak T:\ \mathfrak u(3)\to\mathbb R\ \text{linear},\ \mathfrak T(I_3)=3\ne0\implies\text{surjective}\implies\dim\ker=9-1=8.\ \blacksquare
\]

\textbf{Phys.} The colour algebra is eight-dimensional — the gluon count of CT3.13 and the colour-block dimension of CT0.14 ($\dim\mathcal K_3=25=17+8$).

\end{theorem}

\begin{theorem}[{\texorpdfstring{CT3.9 — Explicit generator basis}{CT3.9 - Explicit generator basis}}]\label{tag:CT3-9}
\textbf{Thm.} With the Gell-Mann matrices

%dsp-level:0
\[
\lambda_1=\begin{pmatrix}0&1&0\\1&0&0\\0&0&0\end{pmatrix},\
\lambda_2=\begin{pmatrix}0&-i&0\\i&0&0\\0&0&0\end{pmatrix},\
\lambda_3=\begin{pmatrix}1&0&0\\0&-1&0\\0&0&0\end{pmatrix},\
\lambda_4=\begin{pmatrix}0&0&1\\0&0&0\\1&0&0\end{pmatrix},
\]

%dsp-level:0
\[
\lambda_5=\begin{pmatrix}0&0&-i\\0&0&0\\i&0&0\end{pmatrix},\
\lambda_6=\begin{pmatrix}0&0&0\\0&0&1\\0&1&0\end{pmatrix},\
\lambda_7=\begin{pmatrix}0&0&0\\0&0&-i\\0&i&0\end{pmatrix},\
\lambda_8=\frac1{\sqrt3}\begin{pmatrix}1&0&0\\0&1&0\\0&0&-2\end{pmatrix},
\]

and $T_a:=\tfrac12\lambda_a$:

%dsp-level:0
\[
T_a^\dagger=T_a;
\]

%dsp-level:0
\[
\qquad \mathfrak T(T_a)=0;
\]

%dsp-level:0
\[
\qquad \{T_a\}_{a=1}^{8}:\ \mathbb R\text{-basis of}\ \mathfrak{su}(3);
\]

%dsp-level:0
\[
\qquad [T_a,T_b]=i\,f_{abc}\,T_c,\ \ f_{abc}\in\mathbb R;
\]

%dsp-level:0
\[
\qquad \mathfrak T(T_aT_b)=\tfrac12\,\delta_{ab}.
\]

\textbf{Req.} C2; CT3.5, CT3.8; direct computation [machine: \texttt{checks/\allowbreak{}verify\_\allowbreak{}F.py}, exact in $\mathbb Q(\sqrt3,i)$].

\textbf{Pf.}

%dsp-level:0
\[
\lambda_a^\dagger=\lambda_a\ \ \forall a\ [\text{inspection}];\qquad
\mathfrak T\lambda_a=0:\ a\le7\ \text{visually};\quad\mathfrak T\lambda_8=\tfrac1{\sqrt3}(1+1-2)=0
\]

%dsp-level:0
\[
\text{supports}:\ \lambda_{1,2}\leftrightarrow(1,2)\text{-entry}\ \{\mathrm{Re},\mathrm{Im}\};
\]

%dsp-level:0
\[
\ \ \lambda_{4,5}\leftrightarrow(1,3);
\]

%dsp-level:0
\[
\ \ \lambda_{6,7}\leftrightarrow(2,3);
\]

%dsp-level:0
\[
\ \ \lambda_{3},
\]

%dsp-level:0
\[
\lambda_{8}\leftrightarrow\text{traceless diagonal}\ (\neg\text{proportional})
\]

%dsp-level:0
\[
\implies8\ \text{independent};\qquad\dim\mathfrak{su}(3)=8\ [\mathrm{CT3.8}]\implies\text{basis}
\]

%dsp-level:0
\[
\mathrm{CT3.5}:\ -i[T_a,T_b]\ \text{Hermitian}\wedge\text{traceless}\implies\in\mathfrak{su}(3)=\operatorname{span}_{\mathbb R}\{T_c\}\implies[T_a,T_b]=if_{abc}T_c,\ \ f_{abc}\in\mathbb R
\]

%dsp-level:0
\[
\mathfrak T(\lambda_a\lambda_b)
\]

%dsp-level:0
\[
= 2\delta_{ab}\ [\text{direct computation, machine-verified}]\implies\mathfrak T(T_aT_b)
\]

%dsp-level:0
\[
= \tfrac14\cdot2\delta_{ab}
\]

%dsp-level:0
\[
= \tfrac12\delta_{ab}
\]

%dsp-level:0
\[
\quad [=\text{uniform normalization of EW.32}].\ \blacksquare
\]

\textbf{Phys.} The explicit normalized colour octet: same uniform trace normalization $\mathfrak T(\mathcal G_A\mathcal G_B)=\tfrac12\delta_{AB}$ as the electroweak generators — one gauge convention across the corpus.

\end{theorem}

\begin{theorem}[{\texorpdfstring{CT3.21 — Structure constants}{CT3.21 - Structure constants}}]\label{tag:CT3-21}
\textbf{Thm.} $\ f_{abc}:=-2i\,\mathfrak T\bigl([T_a,T_b]T_c\bigr)$ is totally antisymmetric, with independent nonzero values

%dsp-level:0
\[
f_{123}=1;\qquad
f_{147}=f_{246}=f_{257}=f_{345}=\tfrac12;\qquad
f_{156}=f_{367}=-\tfrac12;\qquad
f_{458}=f_{678}=\frac{\sqrt3}{2};
\]

$\forall$ other independent components $=0$.

\textbf{Req.} CT3.9; trace cyclicity [linear algebra]; machine: \texttt{checks/\allowbreak{}verify\_\allowbreak{}F.py} (all 512 components + antisymmetry + all 512 Jacobi triples, exact).

\textbf{Pf.}

%dsp-level:0
\[
[T_a,T_b]
\]

%dsp-level:0
\[
= if_{abd}T_d;
\]

%dsp-level:0
\[
\quad\mathfrak T(T_dT_c)
\]

%dsp-level:0
\[
= \tfrac12\delta_{dc}\ [\mathrm{CT3.9}]\implies\mathfrak T\bigl([T_a,T_b]T_c\bigr)
\]

%dsp-level:0
\[
= \tfrac i2f_{abc}
\]

%dsp-level:0
\[
\implies f_{abc}
\]

%dsp-level:0
\[
= -2i\,
\]

%dsp-level:0
\[
\mathfrak T\bigl([T_a,T_b]T_c\bigr)
\]

%dsp-level:0
\[
a\leftrightarrow b:\ \text{antisymmetry of}\ [\cdot,\cdot];
\]

%dsp-level:0
\[
\qquad \mathfrak T\bigl([T_a,T_b]T_c\bigr)=\mathfrak T(T_aT_bT_c)-\mathfrak T(T_bT_aT_c)=\mathfrak T\bigl(T_a[T_b,T_c]\bigr)\implies\text{antisymmetry in}\ b\leftrightarrow c
\]

%dsp-level:0
\[
\implies f\ \text{totally antisymmetric}
\]

%dsp-level:0
\[
\text{sample }1:\ \ [\lambda_1,\lambda_2]=2i\lambda_3\implies[T_1,T_2]=\tfrac14\cdot2i\lambda_3=iT_3\implies f_{123}=1
\]

%dsp-level:0
\[
\text{sample }2:\ \ [\lambda_4,\lambda_5]
\]

%dsp-level:0
\[
= 2i\,
\]

%dsp-level:0
\[
\operatorname{diag}(1,0,-1)
\]

%dsp-level:0
\[
= i\bigl(\lambda_3+\sqrt3\,\lambda_8\bigr)
\]

%dsp-level:0
\[
\implies [T_4,T_5]
\]

%dsp-level:0
\[
= \tfrac i2\bigl(T_3+\sqrt3\,T_8\bigr)
\]

%dsp-level:0
\[
\implies f_{453}
\]

%dsp-level:0
\[
= \tfrac12,
\]

%dsp-level:0
\[
\ f_{458}\qquad = \tfrac{\sqrt3}2
\]

%dsp-level:0
\[
\forall\ (a,b,c):\ \text{complete exact list machine-verified}\ [\text{verify\_F}];
\]

%dsp-level:0
\[
\qquad \text{Jacobi}:\ \sum_{\mathrm{cyclic}}\bigl[[T_a,T_b],T_c\bigr]=0\ \ \forall(a,b,c)\ [\text{machine, exact}].\ \blacksquare
\]

\textbf{Phys.} The colour algebra is fully determined — basis, brackets, normalization, structure constants down to the $\sqrt3/2$ fingerprints of $\lambda_8$. Zero continuous freedom: the strongest rigidity result of the corpus.

\end{theorem}

\begin{theorem}[{\texorpdfstring{CT3.10 — The colour aspect has no central component}{CT3.10 - The colour aspect has no central component}}]\label{tag:CT3-10}
\textbf{Thm.}

%dsp-level:0
\[
F_{\mathrm{strong}}=\textstyle\sum_af_aT_a\implies\mathfrak T(F_{\mathrm{strong}})=0;
\]

%dsp-level:0
\[
\qquad Y_{\mathrm{asp}}F_{\mathrm{strong}}Y_{\mathrm{asp}}=0;
\]

%dsp-level:0
\[
\qquad \alpha_{\mathrm{EM}}^{(3)}=\alpha_{\mathrm{EM}}^{(2)}\quad[=\mathrm{EM.18},\ \text{restated}].
\]

\textbf{Req.} CT0.14 (Part 1); CT3.7; EM.2, EM.17 (\textbf{E}, restated at use).

\textbf{Pf.}

%dsp-level:0
\[
\mathrm{CT3.7}:\ F_{\mathrm{strong}}=\textstyle\sum_{a=1}^{8}f_aT_a,\ f_a\in\mathbb R;\qquad
\mathfrak T(T_a)=0\ \forall a\implies\mathfrak T(F_{\mathrm{strong}})=\textstyle\sum_af_a\cdot0=0
\]

%dsp-level:0
\[
\mathrm{CT0.14}:\ \mathcal K^{(\mathrm{strong})}\perp\mathcal K^{(\mathrm{EM})};
\]

%dsp-level:0
\[
\qquad \mathrm{EM.2}\ \text{(restated)}:\ Y_{\mathrm{asp}}=\mathcal P^{(\mathrm{EM})}\implies Y_{\mathrm{asp}}F_{\mathrm{strong}}Y_{\mathrm{asp}}=\mathcal P^{(\mathrm{EM})}F_{\mathrm{strong}}\mathcal P^{(\mathrm{EM})}=0
\]

%dsp-level:0
\[
\mathrm{EM.17}\ \text{(restated)}:\ \alpha_{\mathrm{EM}}^{(n)}=\mathfrak T\bigl(Y_{\mathrm{asp}}\mathbf F_M^{(n)}Y_{\mathrm{asp}}\bigr)/\mathfrak T\bigl(Y_{\mathrm{asp}}\bigr);
\]

%dsp-level:0
\[
\quad \text{colour contributes}\ 0\implies\alpha_{\mathrm{EM}}^{(3)}=\alpha_{\mathrm{EM}}^{(2)}.\ \blacksquare
\]

\textbf{Phys.} Colour adds only traceless structure orthogonal to the rung-2 subspace: it neither carries a central charge nor disturbs the electromagnetic block read $\alpha_{\mathrm{EM}}$.

\end{theorem}

\begin{corollary}[{\texorpdfstring{CT3.10a — no per-edge abelian strong phase}{CT3.10a - no per-edge abelian strong phase}}]\label{tag:CT3-10a}
$(1)$ \textbf{Coefficient-read obstruction.} EM.17's read $\alpha^{(n)}=\mathfrak T\bigl(Y_{\mathrm{asp}}\mathbf F_M^{(n)}Y_{\mathrm{asp}}\bigr)/\mathfrak T\bigl(Y_{\mathrm{asp}}\bigr)$ yields a single number exactly because the EM block is one-dimensional [CT0.14: $\dim\mathcal K^{(\mathrm{EM})}=1$; EM.17]. The strong block is eight-dimensional [CT0.14, CT3.8: $\dim\mathcal K_3=25=17+8$], and

%dsp-level:0
\[
F_{\mathrm{strong}}=\textstyle\sum_af_aT_a=c\,I_3\quad\iff\quad c=0\ \wedge\ f_a=0\ \ \forall a
\]

[CT3.7: $\mathfrak T(T_a)=0\implies3c=\mathfrak T(F_{\mathrm{strong}})=0$; CT3.9: the $T_a$ are linearly independent] — the read returns no scalar on the strong block: $\neg\exists$ a per-edge abelian strong phase coefficient $\alpha_{\mathrm{strong}}$ defined by derived structure. $(2)$ \textbf{Non-canonicity.} $\mathfrak{su}(3)$ contains $\mathfrak u(1)$ subalgebras [$\operatorname{span}_{\mathbb R}\{T_a\}$, CT3.9], but no generator line is fixed by the adjoint action — all maximal tori are conjugate [math] — so no $\mathfrak u(1)$ choice is distinguished by derived structure. Explicit witnesses (integer-entry $SU(3)$ elements of the CT3.9 basis [CT3.4; $\lambda_a=2T_a$; unitary, $\det=1$ — machine]):

%dsp-level:0
\[
U_2=i\lambda_2+I_3-\lambda_2^{2},\qquad U_5=i\lambda_5+I_3-\lambda_5^{2}\qquad[\lambda_2^{2}=\operatorname{diag}(1,1,0),\ \ \lambda_5^{2}=\operatorname{diag}(1,0,1)]
\]

%dsp-level:0
\[
\operatorname{Ad}_{U_5}(T_1)=U_5\,T_1\,U_5^{\dagger}=-\,T_6\ \notin\ \operatorname{span}_{\mathbb R}\{T_1\};
\]

%dsp-level:0
\[
\qquad \forall a\in\{1,\dots,8\}\ \exists U\in\{U_2,U_5\}:\ \operatorname{Ad}_U(T_a)\notin\operatorname{span}_{\mathbb R}\{T_a\}
\]

%dsp-level:0
\[
\qquad [\text{machine census}].
\]

$(3)$ \textbf{Consequence (pointer).} With $\alpha_{\mathrm{EM}}^{(3)}=\alpha_{\mathrm{EM}}^{(2)}$ [CT3.10; EM.18 restated], the per-edge scalar factor $e^{i\alpha_e\widehat{ds}^{\,2}(e)}$ of the transfer weight is entirely electromagnetic [D3.5: $\alpha_e\in\Phi$, the same alphabet at every rung], and CT3.22a(0)'s decomposition — EM factor $\times$ $\{0,1\}$ singlet-sealing admissibility [D2.14; CT3.17, CT3.18 — forward in-document pointers, non-load-bearing] — exhausts the per-edge strong datum: no phase factor remains.

\textbf{Req (Cor).} D2.14, D3.5; CT0.14 (Part 1); CT3.4, CT3.7, CT3.8, CT3.9, CT3.10; EM.17, EM.18 (\textbf{E}, restated at use); torus conjugacy, linear independence [math; machine: \texttt{checks/\allowbreak{}verify\_\allowbreak{}promotion\_\allowbreak{}F\_\allowbreak{}colour.py}].

\textbf{Pf (Cor).}

%dsp-level:0
\[
(1):\ \mathfrak T(cI_3)=3c;\quad\mathfrak T(F_{\mathrm{strong}})=\textstyle\sum_af_a\,\mathfrak T(T_a)=0\ [\mathrm{CT3.7}]\implies c=0;
\]

%dsp-level:0
\[
\qquad \textstyle\sum_af_aT_a=0\implies f=0\ [\mathrm{CT3.9}:\ \text{basis}]
\]

%dsp-level:0
\[
\dim\mathcal K^{(\mathrm{EM})}
\]

%dsp-level:0
\[
= 1\ \ \text{vs}\ \ \dim\mathcal K^{(\mathrm{strong})}
\]

%dsp-level:0
\[
= 8\ [\mathrm{CT0.14},\ \mathrm{CT3.8}]:\ \mathrm{EM.17}\text{'s single-coefficient read (single eigenvalue on a 1-dim block)}\ \neg\text{replicable on the strong block}
\]

%dsp-level:0
\[
(2):\ \lambda_2^{2}=\operatorname{diag}(1,1,0),\ \lambda_5^{2}=\operatorname{diag}(1,0,1)\implies U_2,U_5\ \text{real orthogonal},\ \det=1\ [\text{machine}];
\]

%dsp-level:0
\[
\qquad U_5\,T_1\,U_5^{\dagger}=-T_6;\quad U_5\,T_2\,U_5^{\dagger}=T_7\ \ \text{etc.}\ [\text{machine census}]
\]

%dsp-level:0
\[
\text{maximal tori conjugate}\ [\text{math}]\implies\neg\exists\ \text{distinguished}\ \mathfrak u(1)\subseteq\mathfrak{su}(3);
\]

%dsp-level:0
\[
\qquad (3):\ \mathrm{CT3.10}:\ \alpha^{(3)}=\alpha^{(2)};\quad\mathrm{D3.5}:\ G^{(n)}(e)=\sigma_e\,e^{i\alpha^{(n)}_{\mathrm{EM}}\widehat{ds}^{\,2}(e)}\ \blacksquare
\]

\end{corollary}

\phantomsection
\section[§F.III — Colour Objects and Spectrum]{\texorpdfstring{§F.III — Colour Objects and Spectrum}{SF.III - Colour Objects and Spectrum}}

\begin{theorem}[{\texorpdfstring{CT3.11 — Colour objects}{CT3.11 - Colour objects}}]\label{tag:CT3-11}
\textbf{Thm.} Within the rung-3 structure:

%dsp-level:0
\[
(1)\ \ \text{colour}:=\text{slot index}\ i\in\{1,2,3\}\ [\mathrm{D3.1}],\ \ \text{fundamental of}\ \mathfrak{su}(3);
\]

%dsp-level:0
\[
(2)\ \ \text{triplet}:\ q^i\mapsto U^i{}_jq^j\ (U\in SU(3));\qquad
(3)\ \ \text{antitriplet}:\ r_i\mapsto r_j(U^{-1})^j{}_i;
\]

%dsp-level:0
\[
(4)\ \ \text{gluon}:=\gamma\in\mathsf{Trivial}\wedge\gamma\cap V_3\ne\varnothing;\qquad
(5)\ \ \text{photon}:=\gamma\in\mathsf{Trivial}\wedge\gamma\subseteq V_2\qquad[\mathrm{D3.7}];
\]

%dsp-level:0
\[
(6)\ \ \text{quark}:=\text{bound loop}\ [\mathrm{D1.5}:\ \text{fermionic}]\ \text{with triplet attachment}\ [\mathrm{D3.8}];
\]

%dsp-level:0
\[
(7)\ \ B:=\varepsilon_{ijk}q^iq^jq^k\ (\text{baryon});\qquad M:=\delta^i{}_jq^jr_i\ (\text{meson})\qquad[\mathrm{D3.9}].
\]

\textbf{Req.} D1.5, D3.1, D3.7, D3.8, D3.9; CT3.2, CT3.4, CT3.12, CT3.15, CT3.17.

\textbf{Pf.}

%dsp-level:0
\[
(1)\text{--}(3):\ \ \mathrm{CT3.2}/\mathrm{CT3.4}:\ \text{carrier}\ \mathbb C^3,\ \ SU(3)\ \text{defining action};\quad\text{dual law}=\mathrm{CT3.17}\ \text{step 1}
\]

%dsp-level:0
\[
(4),(5):\ \ \mathrm{D3.7};\quad\text{distinguished only by vertex type}=\mathrm{CT3.12}
\]

%dsp-level:0
\[
(6):\ \ \mathrm{D3.8};\quad\text{bound}\implies\text{fermionic}\ [\mathrm{D1.5}]
\]

%dsp-level:0
\[
(7):\ \ \mathrm{D3.9};\quad\text{both contractions invariant}:\ \mathrm{CT3.15},\ \mathrm{CT3.17}.\ \blacksquare
\]

\textbf{Phys.} The colour vocabulary, fixed: every term — colour, triplet, antitriplet, gluon, photon, quark, baryon, meson — names a structure already derived from the axioms.

\end{theorem}

\begin{theorem}[{\texorpdfstring{CT3.12 — Photon and gluon are distinguished by possibility space}{CT3.12 - Photon and gluon are distinguished by possibility space}}]\label{tag:CT3-12}
\textbf{Thm.}

%dsp-level:0
\[
\gamma_\gamma,\ \gamma_g\in\mathsf{Trivial}:\ \ \widehat{ds}^{\,2}(e)=0\ \ \forall e,\ \ \Theta=0,\ \ \operatorname{Hol}_{\mathrm{spiral}}=1;
\]

%dsp-level:0
\[
\text{distinction}:\ \ \gamma_\gamma\subseteq V_2\quad\text{vs}\quad\gamma_g\cap V_3\ne\varnothing;
\]

%dsp-level:0
\[
\mathrm{strong}\in\mathcal A(3)\setminus\mathcal A(2)\implies\gamma_\gamma:\ \neg\text{colour index};\quad\gamma_g:\ \text{slot index possible}\ [\mathrm{D3.1}];
\]

%dsp-level:0
\[
\text{class, entropy channel (transit, D1.5), mass}\ (0)\ \text{coincide}.
\]

\textbf{Req.} D0.8, D1.5, D1.7, D3.1, D3.7; CT0.13 (Part 1); EM.5, EM.6 (Part 2); CT3.14.

\textbf{Pf.}

%dsp-level:0
\[
\mathrm{D3.7}:\ \text{both}\in\mathsf{Trivial};\qquad\mathrm{D1.7}:\ \widehat{ds}^{\,2}\equiv0\implies\Theta=0,\ \operatorname{Hol}=e^{i\alpha\cdot0}=1
\]

%dsp-level:0
\[
\mathrm{D0.8}:\ V_2\cap V_3=\varnothing;\qquad\mathrm{CT0.13}:\ \mathrm{strong}\ \text{active only from rung}\ 3
\]

%dsp-level:0
\[
\text{no further distinguishing property}:\ \text{class}\ \mathsf{Trivial};
\]

%dsp-level:0
\[
\ \ \text{channel}
\]

%dsp-level:0
\[
= \text{transit}\ [\mathrm{D1.5}];
\]

%dsp-level:0
\[
\ \ m^2
\]

%dsp-level:0
\[
= 0\ [\mathrm{CT3.14},\ \mathrm{EM.6}].\ \blacksquare
\]

\textbf{Phys.} Photon and gluon are the same kind of object — a null chain — differing only in where each is allowed to travel. The gluon’s colour is a consequence of touching trivalent vertices; the photon’s neutrality, of never doing so. (= EM.5 seen from the colour side.)

\end{theorem}

\begin{theorem}[{\texorpdfstring{CT3.13 — Gluon count}{CT3.13 - Gluon count}}]\label{tag:CT3-13}
\textbf{Thm.} $\ N_{\mathrm{gluon}}=\dim\mathfrak{su}(3)=8;\qquad\neg\exists\ \text{ninth mode}\ (\text{centre of}\ \mathfrak u(3)\ \text{allocated to}\ F_{\mathrm{EM}}).$

\textbf{Req.} CT3.7, CT3.8; adjoint-representation dimension [math].

\textbf{Pf.}

%dsp-level:0
\[
\mathrm{CT3.7}:\ \text{colour gauge algebra}=\mathfrak{su}(3);\qquad
\text{gauge mode}=\text{adjoint direction};\quad\dim\mathrm{adj}=\dim\mathfrak{su}(3)=8\ [\mathrm{CT3.8}]
\]

%dsp-level:0
\[
\text{9th direction of}\ \mathfrak u(3)=\text{centre}\to F_{\mathrm{EM}}\ [\mathrm{CT3.7}\ \text{step 3}];\quad\text{reallocation}=\text{double counting}.\ \blacksquare
\]

\textbf{Phys.} Eight gluons — one per $\mathfrak{su}(3)$ generator. The ninth possible mode is electromagnetism, already counted at rung 2: the subset principle at work.

\end{theorem}

\begin{theorem}[{\texorpdfstring{CT3.14 — Gluons remain massless}{CT3.14 - Gluons remain massless}}]\label{tag:CT3-14}
\textbf{Thm.} $\ m_g^2=0$, by two independent routes:

%dsp-level:0
\[
(\mathrm{i})\quad T_a\mathsf v=0\ [\mathrm{EW.32}]\implies M^2_{aB}=M^2_{ab}=0\implies\{T_a\}\subseteq\ker M^2\implies\text{zero pole displacement}\quad[=\mathrm{EW.26},\ \text{restated}];
\]

%dsp-level:0
\[
(\mathrm{ii})\quad\gamma_g\in\mathsf{Trivial}\implies\Theta(\gamma_g)=0\implies\neg\text{wrapping holonomy}\implies\neg\text{wrapping contribution}\ [\mathrm{EW.18}].
\]

\textbf{Req.} D2.5, D3.7; CT0.14 (Part 1); EW.18, EW.26, EW.32, EW.33 (\textbf{D}, restated at use).

\textbf{Pf.}

%dsp-level:0
\[
(\mathrm{i})\ \ \mathrm{D2.5}:\ m^2=\text{pole displacement};
\]

%dsp-level:0
\[
\qquad \mathrm{EW.33}\ \text{(restated)}:\ \text{rung-3 displacement source}=M^2_{AB}=\mu_\star^2\operatorname{Re}\langle\mathcal G_A\mathsf v,\mathcal G_B\mathsf v\rangle
\]

%dsp-level:0
\[
\mathrm{EW.32}\ \text{(restated)}:\ F_H=\text{pair block}\perp F_{\mathrm{strong}}=\text{stack block}\ [\mathrm{CT0.14}]\implies\text{seed has}\ \neg\text{slot index}\implies T_a\mathsf v=0
\]

%dsp-level:0
\[
\implies M^2_{aB}=\mu_\star^2\operatorname{Re}\langle0,\mathcal G_B\mathsf v\rangle=0,\quad M^2_{ab}=0\implies\{T_a\}_{a=1}^{8}\subseteq\ker M^2\implies m_g^2=0
\]

%dsp-level:0
\[
(\mathrm{ii})\ \ \mathrm{D3.7}:\ \gamma_g\in\mathsf{Trivial}\implies\Theta=0;
\]

%dsp-level:0
\[
\qquad \mathrm{EW.18}\ \text{(restated)}:\ \text{materialization requires}\ \Phi\in2\pi\mathbb Z\ \text{from winding}\implies\neg\text{wrapping}\implies\neg\text{contribution}.\ \blacksquare
\]

\textbf{Phys.} Gluons are massless structurally: the seed that generates mass is colour-blind (route i), and null chains cannot wind (route ii). Two independent arguments, agreeing — masslessness of the strong gauge modes is as robust as the photon’s.

\end{theorem}

\begin{theorem}[{\texorpdfstring{CT3.15 — Baryon invariance}{CT3.15 - Baryon invariance}}]\label{tag:CT3-15}
\textbf{Thm.} $\ B=\varepsilon_{ijk}q^iq^jq^k;\qquad U\in SU(3):\quad B\mapsto\varepsilon_{ijk}U^i{}_aU^j{}_bU^k{}_c\,q^aq^bq^c=\det(U)\,\varepsilon_{abc}\,q^aq^bq^c=B.$

\textbf{Req.} D3.9; CT3.2; Leibniz identity $\varepsilon_{ijk}U^i{}_aU^j{}_bU^k{}_c=\det(U)\varepsilon_{abc}$ [math]; machine: \texttt{checks/\allowbreak{}verify\_\allowbreak{}F.py}.

\textbf{Pf.}

%dsp-level:0
\[
\mathrm{CT3.2}:\ q^i\mapsto U^i{}_aq^a\implies B\mapsto\varepsilon_{ijk}U^i{}_aU^j{}_bU^k{}_c\,q^aq^bq^c
\]

%dsp-level:0
\[
\varepsilon_{ijk}U^i{}_aU^j{}_bU^k{}_c=\det(U)\,\varepsilon_{abc}\quad[\text{Leibniz expansion of}\ \det;\ \text{machine-verified on random exact}\ U]
\]

%dsp-level:0
\[
U\in SU(3)\implies\det U=1\implies B\mapsto\varepsilon_{abc}q^aq^bq^c=B.\ \blacksquare
\]

\textbf{Phys.} The totally antisymmetric three-index contraction is a colour singlet: the baryon — three triplets sealed by $\varepsilon$ — is the minimal neutral object of the strong sector.

\end{theorem}

\begin{theorem}[{\texorpdfstring{CT3.16 — Infinitesimal baryon invariance}{CT3.16 - Infinitesimal baryon invariance}}]\label{tag:CT3-16}
\textbf{Thm.} $\ U=I_3+i\theta^aT_a+O(\theta^2),\ \ \mathfrak T(T_a)=0\implies\det U=1+O(\theta^2)\implies B\mapsto B+O(\theta^2)$ (first-order variation vanishes).

\textbf{Req.} CT3.7 ($\mathfrak T T_a=0$), CT3.15; $\det=\exp\mathfrak T\log$ [math].

\textbf{Pf.}

%dsp-level:0
\[
\log U=i\theta^aT_a+O(\theta^2);\qquad
\det U=\exp\mathfrak T(\log U)=\exp\bigl(i\theta^a\underbrace{\mathfrak T(T_a)}_{=0}+O(\theta^2)\bigr)=1+O(\theta^2)
\]

%dsp-level:0
\[
\mathrm{CT3.15}:\ B\mapsto\det(U)\,B=B+O(\theta^2).\ \blacksquare
\]

\textbf{Phys.} The singlet property holds at the algebra level: tracelessness of the colour generators is exactly the condition that makes baryon sealing infinitesimally exact.

\end{theorem}

\begin{theorem}[{\texorpdfstring{CT3.17 — Invariant tensors}{CT3.17 - Invariant tensors}}]\label{tag:CT3-17}
\textbf{Thm.} The invariant tensors for sealing colour indices are generated by

%dsp-level:0
\[
\delta^i{}_j,\qquad\varepsilon_{ijk},\qquad\varepsilon^{ijk};
\]

%dsp-level:0
\[
q^ir_i\mapsto U^i{}_aq^a\,r_b(U^{-1})^b{}_i=q^ar_a;
\]

%dsp-level:0
\[
\qquad \varepsilon_{ijk}:\ 3\ \text{fundamentals}\ [\mathrm{CT3.15}];
\]

%dsp-level:0
\[
\qquad \varepsilon^{ijk}:\ 3\ \text{duals}\ \bigl(\det(U^{-1})^{\mathsf T}=1\bigr);
\]

%dsp-level:0
\[
\forall\ \text{invariants of}\ (\mathbf3,\bar{\mathbf3}):\ \ \text{linear combinations of products of}\ \{\delta,\varepsilon,\bar\varepsilon\}.
\]

\textbf{Req.} CT3.2, CT3.15; determinant identities [math]; machine: \texttt{checks/\allowbreak{}verify\_\allowbreak{}F.py}.

\textbf{Pf.}

%dsp-level:0
\[
\delta:\ \ U^i{}_a(U^{-1})^b{}_i=\delta^b{}_a\implies q^ir_i\ \text{invariant}\quad[\text{machine}:\ U^{-1}U=I]
\]

%dsp-level:0
\[
\varepsilon_{ijk}:\ \mathrm{CT3.15};\qquad
\varepsilon^{ijk}:\ \text{same identity at}\ (U^{-1})^{\mathsf T},\ \ \det(U^{-1})^{\mathsf T}=1
\]

%dsp-level:0
\[
\text{generation}:\ \delta\ \text{seals matched pairs};
\]

%dsp-level:0
\[
\ \ \varepsilon/\bar\varepsilon\ \text{seal residues mod}\ 3\ \text{of unmatched indices}\implies\forall\ \text{invariant}
\]

%dsp-level:0
\[
= \text{combination of products}.\ \blacksquare
\]

\textbf{Phys.} Only three tensors can seal colour: pair-annihilation $\delta$, three-quark $\varepsilon$, three-antiquark $\bar\varepsilon$. The entire hadronic vocabulary — mesons, baryons, antibaryons — is built from these.

\end{theorem}

\begin{corollary}[{\texorpdfstring{CT3.17a — sealing partition rule; grammar counts}{CT3.17a - sealing partition rule; grammar counts}}]\label{tag:CT3-17a}
$(1)$ \textbf{Partition rule.} If a colour-attachment set is sealed — $\mathcal C(\Sigma;C)=1$ [D2.14] — then $\mathrm{Colour}(\Sigma)$ partitions into elementary contraction groups of size $2$ ($\delta$-type) and size $3$ ($\varepsilon/\bar\varepsilon$-type) [CT3.17: the invariant tensors are $\delta^{i}{}_{j}$ ($2$ members), $\varepsilon_{ijk}$, $\varepsilon^{ijk}$ ($3$ members); D2.14: $\pi_C(v)=$ the other members of $v$'s elementary group]. A lone attachment is never sealed:

%dsp-level:0
\[
|\mathrm{Colour}(\Sigma)|=1\implies\mathcal C(\Sigma;C)=0
\]

[a contraction partner is a distinct node — A1: $p\ne q$ — hence $\pi_C(v)\not\subseteq\{v\}$; equivalently, CT3.17's list contains no invariant of a single slot]. $(2)$ \textbf{Grammar count.} The number $g(N)$ of elementary groupings of a labelled $N$-element attachment set into groups of size $2$ and $3$ (labelled set partitions with parts in $\{2,3\}$) is

%dsp-level:0
\[
g(N)=\!\!\sum_{\substack{a,b\ge0\  2a+3b=N}}\!\!\frac{N!}{a!\,b!\,2^{a}\,6^{b}};
\]

%dsp-level:0
\[
\qquad g(1)=0,\quad g(2)=1,\quad g(3)=1,\quad g(4)=3,\quad g(5)=10,\quad g(6)=25\ (15+10),\quad g(7)=105 .
\]

$(3)$ \textbf{Type selection.} Attachment-type assignment (triplet/antitriplet per attachment [D3.8]) selects the type-compatible groupings — $\delta^{i}{}_{j}$ consumes one attachment of each type, $\varepsilon_{ijk}$ ($\varepsilon^{ijk}$) three of the same type [CT3.17] — configuration data refining the count of $(2)$.

\textbf{Req (Cor).} A1; D2.14, D3.8; CT3.17; labelled set-partition enumeration [math; machine: \texttt{checks/\allowbreak{}verify\_\allowbreak{}promotion\_\allowbreak{}F\_\allowbreak{}colour.py}].

\textbf{Pf (Cor).}

%dsp-level:0
\[
(1):\ \mathrm{CT3.17}:\ \text{every invariant}
\]

%dsp-level:0
\[
= \text{combination of products of}\ \{\delta,\varepsilon,\bar\varepsilon\}:\ \text{an elementary group consumes}\ 2\ \text{or}\ 3\ \text{slots};
\]

%dsp-level:0
\[
\qquad \mathrm{D2.14}:\ \mathcal C=1\iff\forall v:\ \pi_C(v)\subseteq\Sigma
\]

%dsp-level:0
\[
\implies\mathrm{Colour}(\Sigma)=\bigsqcup\ \text{groups of size}\ 2,3;
\]

%dsp-level:0
\[
\qquad N=1:\ \mathrm{A1}:\ p\ne q\implies\pi_C(v)\ \text{requires a distinct member}\implies\pi_C(v)\not\subseteq\{v\}\implies\mathcal C=0\ [\mathrm{D2.14}]
\]

%dsp-level:0
\[
(2):\ \text{fix}\ (a,b),\ 2a+3b=N:\ \ \#\{\text{partitions into}\ a\ \text{pairs}+b\ \text{triples}\}=\frac{N!}{a!\,b!\,(2!)^{a}\,(3!)^{b}}
\]

%dsp-level:0
\[
\quad [\text{math}:\ \text{labelled elements; block order and within-block order unlabelled}]
\]

%dsp-level:0
\[
(2!)^{a}=2^{a},\ (3!)^{b}=6^{b};
\]

%dsp-level:0
\[
\qquad g(4)=\tfrac{4!}{2!\,2^{2}}=3;
\]

%dsp-level:0
\[
\quad g(5)=\tfrac{5!}{1!\,1!\,2\cdot6}=10;
\]

%dsp-level:0
\[
\quad g(6)=\tfrac{6!}{3!\,2^{3}}+\tfrac{6!}{2!\,6^{2}}=15+10=25;
\]

%dsp-level:0
\[
\quad g(7)=\tfrac{7!}{2!\,1!\,2^{2}\cdot6}=105\ [\text{machine}]
\]

%dsp-level:0
\[
(3):\ \mathrm{D3.8}:\ \text{attachment types};
\]

%dsp-level:0
\[
\quad\delta^{i}{}_{j}:\ (\mathbf3,\bar{\mathbf3});
\]

%dsp-level:0
\[
\quad\varepsilon_{ijk}:\ (\mathbf3,\mathbf3,\mathbf3);
\]

%dsp-level:0
\[
\quad\varepsilon^{ijk}:\ (\bar{\mathbf3},\bar{\mathbf3},\bar{\mathbf3})\ [\mathrm{CT3.17}]\ \blacksquare
\]

\end{corollary}

\begin{theorem}[{\texorpdfstring{CT3.18 — Colour-singlet closure}{CT3.18 - Colour-singlet closure}}]\label{tag:CT3-18}
\textbf{Thm.}

%dsp-level:0
\[
\text{free colour index}=\text{dangling boundary of the colour ledger}\ \neg\text{admitted by}\ \mathcal P_M^{(3)};
\]

%dsp-level:0
\[
\mathcal P_M^{(3)}\big|_{\text{free index}}=0\implies\neg\exists\ \text{pole}\ [\mathrm{CT2.5}];
\]

%dsp-level:0
\[
\qquad \text{surviving configurations}:\ \text{fully contracted by}\ \{\delta,\varepsilon,\bar\varepsilon\}\implies\text{colour singlets};
\]

%dsp-level:0
\[
\text{minimal singlets}:\ M=\delta^i{}_jq^jr_i,\quad B=\varepsilon_{ijk}q^iq^jq^k\ [\mathrm{D3.9}].
\]

\textbf{Req.} D2.3, D2.11, D3.1, D3.9; CT2.5 (Part 2); CT3.17.

\textbf{Pf.}

%dsp-level:0
\[
\mathrm{D2.11}:\ \text{valid universe}:\ \text{closed}\ \wedge\ \text{balanced seals}\ \wedge\ \neg\exists\ \text{dangling ledger boundary}
\]

%dsp-level:0
\[
\mathrm{D3.1}:\ \text{colour indices}
\]

%dsp-level:0
\[
= \text{ledger labels at trivalent vertices}\implies\text{uncontracted index}
\]

%dsp-level:0
\[
= \text{dangling colour-ledger boundary}
\]

%dsp-level:0
\[
\mathrm{D2.3}:\ \mathcal P_M^{(3)}\ \text{projects onto valid configurations}\implies\text{free-index configurations annihilated};
\]

%dsp-level:0
\[
\qquad \mathrm{CT2.5}:\ \neg\text{pole}
\]

%dsp-level:0
\[
\text{survivors fully contracted};\quad\mathrm{CT3.17}:\ \text{contractions}\in\{\delta,\varepsilon,\bar\varepsilon\}\implies\text{invariant}=\text{colour singlet}.\ \blacksquare
\]

\textbf{Phys.} The mask admits only colour-sealed configurations: every observable bound rung-3 object is a colour singlet, minimally a meson or a baryon. Validity of the ledger \textit{is} the singlet condition.

\end{theorem}

\begin{theorem}[{\texorpdfstring{CT3.22 — Sealing is confinement}{CT3.22 - Sealing is confinement}}]\label{tag:CT3-22}
\textbf{Thm.}

%dsp-level:0
\[
\neg\exists\ \text{isolated triplet pole of}\ Z(\mathcal R,3);\qquad
\text{isolated triplet}:\ \text{physical norm}=0\ [\mathrm{CT2.8}];
\]

%dsp-level:0
\[
\text{colour charge}\ \neg\text{observable as end-charge}\ [\mathrm{D3.4}];
\]

%dsp-level:0
\[
\qquad \text{contrast}:\ q_{\mathrm{EM}}\ \text{observable on open chains}\ (\text{per-edge registration},\ \neg\text{contraction}).
\]

\textbf{Req.} D3.4; CT0.9 (Part 1); CT2.3, CT2.5, CT2.8 (Part 2); CT3.18.

\textbf{Pf.}

%dsp-level:0
\[
\mathrm{CT3.18}:\ \text{free colour index}\implies\notin\operatorname{Im}\mathcal P_M^{(3)}
\]

%dsp-level:0
\[
\mathrm{CT2.3}:\ \mathbf F_M^{(3)}\big|_{\mathcal K_3^\perp}=0;
\]

%dsp-level:0
\[
\qquad \mathrm{CT2.5}:\ \text{poles}=\text{eigenvalue conditions on}\ \mathbf F_M^{(3)}\implies\text{excluded configurations support}\ \neg\text{poles}
\]

%dsp-level:0
\[
\mathrm{CT2.8}\ \text{step 3}:\ \text{mask-excluded}=\text{null vectors of}\ \langle\cdot|\cdot\rangle\implies\text{isolated triplet}:\ \text{norm}=0
\]

%dsp-level:0
\[
\text{contrast}:\ \mathrm{D3.4}:\ q_{\mathrm{EM}}=\tfrac12(\sigma_s-\sigma_t)\ \text{from edge signs};
\]

%dsp-level:0
\[
\quad \mathrm{CT0.9}:\ \sigma_e\ \text{registered per real edge},\ \neg\text{contraction required}\implies q_{\mathrm{EM}}\ \text{observable on open chains}.\ \blacksquare
\]

\textbf{Phys.} Confinement is a theorem: isolated colour has no pole and no norm — it is filtered out by the mask, not suppressed by a force law. The asymmetry with visible electric charge is exact: a sign is registered per edge; an index must await contraction.

\end{theorem}

\begin{theorem}[{\texorpdfstring{CT3.22a — Subnetwork confinement: the domain restriction, all-or-nothing, and the converse}{CT3.22a - Subnetwork confinement: the domain restriction, all-or-nothing, and the converse}}]\label{tag:CT3-22a}
\textbf{Thm.} Let $\Sigma$ be a subnetwork [D2.13] with $\mathcal C(\Sigma;C)$, $\Pi_{\mathcal C}(\Sigma;C)$ [D2.14], and $F_{\mathrm{strong}}|_\Sigma:=$ the strong force operator restricted to the colour slots of $\Sigma$ under the rung mask [CT2.3, CT2.4; D3.1]. $(0)$ \textbf{Global sum-restriction.}

%dsp-level:0
\[
Z^{(3)}(\mathcal R)=\sum_{\substack{C\in\mathfrak C(\mathcal R):\ \text{colour-sealed}}}\ \sum_\sigma\mathcal J(C)\!\!\prod_{e\in E(C)}\!\!G^{(3)}(e)\ (-1)^{\aleph(\sigma,C)}
\]

%dsp-level:0
\[
\qquad [\mathrm{D0.10},\ \mathrm{C12}]
\]

— confinement restricts the summation domain, $\neg$any individual edge factor: the per-edge geometric/registration data are rung-independent once the block read is rung-stable ($\alpha_{\mathrm{EM}}^{(3)}=\alpha_{\mathrm{EM}}^{(2)}$ — the EM.18/EM.19 mechanism) and the per-edge phase alphabet is rung-blind ($\alpha_e\in\Phi=\{0,\pi\}$ at every rung [D0.5; CT0.9; EM.20(1)]). $(1)$ \textbf{All-or-nothing.} With $F^\Sigma_{\mathrm{strong}}(C):=\Pi_{\mathcal C}(\Sigma;C)\,F_{\mathrm{strong}}|_\Sigma\,\Pi_{\mathcal C}(\Sigma;C)$:

%dsp-level:0
\[
\mathcal C(\Sigma;C)=0\implies F^\Sigma_{\mathrm{strong}}(C)=0
\]

— regardless of how many colour-carrying vertices $\Sigma$ contains, and even if $|\partial_{\mathrm{col}}(\Sigma;C)|=1$. $(2)$ \textbf{Converse with the exact exception locus.} $\Pi_{\mathcal C}=\mathcal C\cdot\mathbb I$ with $\mathcal C\in\{0,1\}$ [D2.14: scalar idempotent] $\implies F^\Sigma_{\mathrm{strong}}=\mathcal C\cdot F_{\mathrm{strong}}|_\Sigma$; hence

%dsp-level:0
\[
\mathcal C
\]

%dsp-level:0
\[
= 1
\]

%dsp-level:0
\[
\implies F^\Sigma_{\mathrm{strong}}
\]

%dsp-level:0
\[
= F_{\mathrm{strong}}|_\Sigma
\]

%dsp-level:0
\[
\ne 0\quad\text{for nonzero aspect data}\ (f_a\ne0\ \text{on}\ \ge1\ \text{colour slot}\ [\mathrm{CT3.9},\ \mathrm{CT3.10}])
\]

with exact exception locus $\{f=0\}\ \cup\ \{\mathrm{Colour}(\Sigma)=\varnothing\}$. $(3)$ \textbf{Singlet diagonal (consistency).} $T_a|\text{singlet}\rangle=0$ [CT3.15, CT3.17 invariance] $\implies\langle\text{singlet}|F_{\mathrm{strong}}|\text{singlet}\rangle=0$ — consistent with CT3.22 (colour is $\neg$observable as end-charge [D3.4]): the nonzero content of the strong block is the admissibility structure and the gluon mode count [CT3.13], $\neg$singlet expectation values.

\textbf{Req.} D2.13, D2.14; D0.10, D3.1, D3.4, D3.5, C12; CT2.3, CT2.4, CT2.5, CT2.8 (Part 2); CT3.2, CT3.9, CT3.10, CT3.13, CT3.15, CT3.17, CT3.18, CT3.22; EM.18, EM.19, EM.20(1) for clause (0) note (Part 2); idempotent algebra, the tensor null-factor argument [math; machine: \texttt{checks/\allowbreak{}verify\_\allowbreak{}merge\_\allowbreak{}tierA.py} §A].

\textbf{Pf.}

%dsp-level:0
\[
(0):\ \mathrm{CT3.18}:\ \mathcal P_M^{(3)}\ \text{admits only colour-sealed}\ C\ (\delta^i{}_j,\ \varepsilon_{ijk},\ \varepsilon^{ijk}\ \text{contractions}\ [\mathrm{CT3.17}]);
\]

%dsp-level:0
\[
\qquad \mathrm{CT3.22}:\ \neg\exists\ \text{isolated triplet pole}
\]

%dsp-level:0
\[
\text{free-index}\ C\notin\operatorname{Im}\mathcal P_M^{(3)};
\]

%dsp-level:0
\[
\quad \mathrm{CT2.5}:\ \text{poles}=\text{eigenvalue conditions on}\ \mathbf F_M^{(3)};
\]

%dsp-level:0
\[
\quad \mathrm{CT2.3}:\ \mathbf F_M^{(3)}\big|_{\mathcal K_3^\perp}=0\implies\text{excluded}\ C:\ \neg\text{pole}\implies\text{domain restriction}
\]

%dsp-level:0
\[
\text{edge factors}:\ G^{(n)}(e)=\sigma_e\exp(i\alpha_e\widehat{ds}^{\,2}(e))\ [\mathrm{D3.5}];
\]

%dsp-level:0
\[
\quad \alpha_e\in\{0,\pi\}\ \text{at every rung}\ [\mathrm{EM.20}(1);\ \mathrm{CT0.9}];
\]

%dsp-level:0
\[
\quad \alpha_{\mathrm{EM}}^{(3)}=\alpha_{\mathrm{EM}}^{(2)}\ [\mathrm{EM.18}/\mathrm{EM.19}]\implies G^{(3)}=G^{(2)}\implies\neg\text{weight modification}
\]

%dsp-level:0
\[
(1):\ \mathcal C=0:\ \exists v\in\mathrm{Colour}(\Sigma):\ \pi_C(v)\not\subseteq\Sigma\ [\mathrm{D2.14}]\implies v\ \text{carries an index uncontracted relative to}\ \Sigma
\]

%dsp-level:0
\[
\mathrm{CT3.18}\ \text{with}\ \Sigma\ \text{as ambient}:\ \text{uncontracted index}
\]

%dsp-level:0
\[
= \text{dangling colour-ledger boundary}\implies\text{mask-excluded};
\]

%dsp-level:0
\[
\qquad \mathrm{CT3.22}+\mathrm{CT2.5}:\ \neg\text{pole of the}\ \Sigma\text{-restricted resolvent}
\]

%dsp-level:0
\[
\mathrm{CT2.8}:\ \text{physical norm}\ 0\implies\text{the}\ \Sigma\text{-restricted colour state of}\ v\ \text{is a null vector};
\]

%dsp-level:0
\[
\qquad \text{colour sector of}\ \Sigma\ =\bigotimes_{\mathrm{Colour}(\Sigma)}(\text{slot spaces})\ [\mathrm{D3.1},\ \mathrm{CT3.2}]
\]

%dsp-level:0
\[
\text{one null factor}\implies\text{whole product null}\implies F_{\mathrm{strong}}|_\Sigma\ \text{has vanishing matrix elements};
\]

%dsp-level:0
\[
\qquad \Pi_{\mathcal C}=0\implies F^\Sigma_{\mathrm{strong}}=0\cdot F_{\mathrm{strong}}|_\Sigma\cdot0=0
\]

%dsp-level:0
\[
(2):\ \mathcal C\in\{0,1\}\implies\mathcal C^2=\mathcal C\implies\Pi_{\mathcal C}\,F\,\Pi_{\mathcal C}=\mathcal C^2\,F|_\Sigma=\mathcal C\,F|_\Sigma;
\]

%dsp-level:0
\[
\qquad \mathcal C=1:\ T_a\ne0\ \text{on}\ \mathbb C^3\ [\mathrm{CT3.9}:\ \text{explicit Gell-Mann basis}]\ \wedge\ f\ne0\ [\mathrm{CT3.10}]\implies F|_\Sigma\ne0
\]

%dsp-level:0
\[
\text{exception locus}:\ f=0\implies F|_\Sigma=0\ \text{by linearity};\quad \mathrm{Colour}(\Sigma)=\varnothing\implies F|_\Sigma\ \text{has no argument}
\]

%dsp-level:0
\[
(3):\ \varepsilon_{ijk}U^i{}_aU^j{}_bU^k{}_c=\varepsilon_{abc}\det U=\varepsilon_{abc}\quad\forall U\in SU(3)\ [\mathrm{CT3.15}];
\]

%dsp-level:0
\[
\qquad \delta^i{}_j\ \text{likewise invariant}\ [\mathrm{CT3.17}]
\]

%dsp-level:0
\[
\implies\text{singlets fixed under}\ SU(3)\implies\mathfrak{su}(3)\ \text{annihilates them}\implies\langle\text{singlet}|F_{\mathrm{strong}}|\text{singlet}\rangle=0.\ \blacksquare
\]

\textbf{Phys.} Confinement becomes a two-sided exact law at subnetwork level: the colour block of a region lives precisely when the region self-seals against the ambient ledger — one dangling attachment anywhere zeroes the entire colour sector, no matter how many safely paired vertices sit inside. Unlike the additive per-vertex read of EM [D3.3], the strong read is $\neg$gradable: a sealed object never carries a colour number, only the grammar of how it composes with its environment.

\textbf{Rmk (non-locality and filter placement; witness scope).} $\mathcal C(\Sigma;C)$ is $\neg$a function of $\Sigma$-local data: valid ambient pairs $C,C'$ with identical restrictions $C|_\Sigma=C'|_\Sigma$ yet $\mathcal C=1$ vs $0$ exist once the ambient supplies further triplet carriers to seal (explicit finite witness: two valid ambient groupings, identical $\Sigma$-occupation, opposite seals) — an ambient-richness scope for that witness, $\neg$derivable from A1 alone; no theorem clause above rests on it (A §0.2a note). Hence the completeness filter belongs \textit{downstream} of any $\mathcal M(\Sigma)$ weighting: pre-weighting by $\mathcal C$ would delete the physical EM/$G$/$H$ reads of externally-sealed restrictions, which are colour-independent by construction [D3.3, CT1.12]. Universe-level validity remains CT3.18/CT2.8(3)'s question — distinct from $\Sigma$-relativity.

\end{theorem}

\phantomsection
\section[§F.IV — Completion]{\texorpdfstring{§F.IV — Completion}{SF.IV - Completion}}

\begin{theorem}[{\texorpdfstring{CT3.19 — Self-containment of the colour sector}{CT3.19 - Self-containment of the colour sector}}]\label{tag:CT3-19}
\textbf{Thm.}

%dsp-level:0
\[
\forall\ \text{statements of}\ \mathbf F:\quad\vdash\ \mathrm{A1\text{–}A9}\cup\mathrm{C1\text{–}C9}\cup\mathrm{D0.1\text{–}D4.3}\cup\{\mathbf B,\mathbf C\ \text{theorems}\}\cup\{\mathrm{EW.26},\mathrm{EM.18}\ \text{restated}\};
\]

%dsp-level:0
\[
\neg\exists\ \text{external strong-interaction law, parameter, or measured value};\qquad
\text{open quantities}=\varnothing:
\]

%dsp-level:0
\[
N_{\text{colour}}=3,\quad N_{\mathrm{gluon}}=8,\quad f_{abc},\quad\mathfrak T(T_aT_b)=\tfrac12\delta_{ab},\quad m_g^2=0\qquad\text{all fixed};
\]

%dsp-level:0
\[
\text{colour sector contributes}\ \neg\mathrm{P4}\ \text{data}.
\]

\textbf{Req.} CT3.1–CT3.22; P4.

\textbf{Pf.} Dependency trace:

%dsp-level:0
\[
3\ \text{slots}\Leftarrow2^{n+1}\le16\ [\mathrm{CT3.20}\Leftarrow\mathrm{CT0.5},\mathrm{CT0.6},\mathrm{CT0.11}]
\]

%dsp-level:0
\[
S_3\hookrightarrow U(3)\ [\mathrm{CT3.1}\text{--}\mathrm{CT3.4}];
\]

%dsp-level:0
\[
\quad \mathfrak u(3)=\mathfrak{su}(3)\oplus\mathfrak u(1)\ [\mathrm{CT3.5}];
\]

%dsp-level:0
\[
\quad \text{centre spent at rung 2}\implies\text{new}=\mathfrak{su}(3)\ [\mathrm{CT3.7},\mathrm{D3.2}]
\]

%dsp-level:0
\[
\dim=8\ [\mathrm{CT3.8}];\quad\text{normalized basis}\ [\mathrm{CT3.9}];\quad f_{abc}\ \text{determined}\ [\mathrm{CT3.21}];\quad\text{normalization}=\mathrm{EW.32}\text{'s}
\]

%dsp-level:0
\[
\mathfrak T=0\implies\neg\text{abelian carry/shift}\ [\mathrm{CT3.10}]
\]

%dsp-level:0
\[
\text{gluons}:\ 8\ \text{massless trivial chains}\ [\mathrm{CT3.13},\mathrm{CT3.14}];
\]

%dsp-level:0
\[
\quad\gamma/g\ \text{differ by possibility space only}\ [\mathrm{CT3.12}]
\]

%dsp-level:0
\[
\{\delta,\varepsilon,\bar\varepsilon\}\ [\mathrm{CT3.17}];
\]

%dsp-level:0
\[
\quad\text{mask admits only singlets}\ [\mathrm{CT3.18}];
\]

%dsp-level:0
\[
\quad\text{isolated colour}:\ \neg\text{pole},
\]

%dsp-level:0
\[
\ \text{norm}\ 0\ [\mathrm{CT3.22}].\ \blacksquare
\]

\textbf{Phys.} Closure audit: the colour sector is the most rigid in the corpus — number of colours, number of gluons, all structure constants, normalization and masslessness are fixed by the axioms, with not a single configuration parameter of its own.

\end{theorem}

\phantomsection
\section[Closure exports of Part 2 §F (consumed by Part 2 §E, §G,…]{\texorpdfstring{Closure exports of Part 2 §F (consumed by Part 2 §E, §G, Parts 3–4)}{Closure exports of Part 2 SF (consumed by Part 2 SE, SG, Parts 3-4)}}

%dsp-level:0
\[
\begin{array}{lll}
\mathrm{CT3.20}\ N_{\text{colour}}=3\ (\text{forced}) &
\mathrm{CT3.7}\ \text{new algebra}=\mathfrak{su}(3),\ \mathfrak T(T_a)=0\ [\text{cited in E}] &
\mathrm{CT3.8}\ \dim=8\ [\text{cited in B: CT0.14}]\\
\mathrm{CT3.9}\ \{T_a\}:\ \mathfrak T(T_aT_b)=\tfrac12\delta_{ab} &
\mathrm{CT3.21}\ f_{abc}\ \text{exact list} &
\mathrm{CT3.10(a)}\ \mathfrak T(F_{\mathrm{strong}})=0;\ \alpha^{(3)}=\alpha^{(2)};\ \neg\exists\,\alpha_{\mathrm{strong}}\\
\mathrm{CT3.13}\ N_{\mathrm{gluon}}=8 &
\mathrm{CT3.14}\ m_g^2=0\ (\times2\ \text{routes}) &
\mathrm{CT3.15}\text{--}\mathrm{CT3.17}\ B,\ M\ \text{singlets};\ \{\delta,\varepsilon,\bar\varepsilon\};\ \mathrm{CT3.17a}:\ g(N)\ \text{grammar}\\
\mathrm{CT3.18}\ \text{mask admits only singlets} &
\mathrm{CT3.22}\ \text{confinement}:\ \neg\text{triplet pole},\ \text{norm}\ 0 &
\mathrm{CT3.19}\ \text{sector rigid},\ \neg\mathrm{P4}\ \text{data}
\end{array}
\]

$\blacksquare$

\phantomsection
\section[§G — Saturation Lock, Dark Aspects, Bounds]{\texorpdfstring{§G — Saturation Lock, Dark Aspects, Bounds}{SG - Saturation Lock, Dark Aspects, Bounds}}

\textit{Dep.: Part 1; Part 2 §C–§F (EW.27 at use). Order: §G.I the saturation lock (CT4.1, +CT4.1a, CT4.2, CT4.3, CT4.4, +CT4.4a) · §G.II the two dark aspects (CT4.5, +CT4.5a, CT4.6, CT4.7, +CT4.7a, CT4.10) · §G.III bounds and boundaries (CT4.9, CT4.8, CT4.11) · §G.IV falsifiability and the closure audit (CT4.12, CT4.13). Notation as Part 2 §C–§F. Machine: \texttt{checks/\allowbreak{}verify\_\allowbreak{}G.py}; CT4.4a/CT4.7a: \texttt{checks/\allowbreak{}verify\_\allowbreak{}merge\_\allowbreak{}tierA.py}; CT4.5a/CT4.1a: \texttt{checks/\allowbreak{}verify\_\allowbreak{}promotion\_\allowbreak{}G\_\allowbreak{}saturation.py}.}

\phantomsection
\section[§G.I — The Saturation Lock]{\texorpdfstring{§G.I — The Saturation Lock}{SG.I - The Saturation Lock}}

\begin{theorem}[{\texorpdfstring{CT4.1 — Oversubscription falls to $D^-$}{CT4.1 - Oversubscription falls to}}]\label{tag:CT4-1}
\textbf{Thm.}

%dsp-level:0
\[
\forall n\ge4:\quad D_{\mathrm{tick}}(n)=2^{n+1}\ge32>16=C_{\mathrm{cap}}\implies O(n)=2^{n+1}-16>0;
\]

%dsp-level:0
\[
O(4)=32-16=16=C_{\mathrm{cap}}\quad(\text{exactly one further channel quantum per tick});
\]

%dsp-level:0
\[
F_{D^-}=\bigoplus_{n\ge4}F_n\ [\mathrm{D4.1}];\qquad
\text{absorption rate}\propto O(n)\sim2^{n+1}\ [\mathrm{D4.2}];\qquad
\text{last admissible rung}=3 .
\]

%dsp-level:0
\[
\begin{array}{c|c|c|c}
n & D_{\mathrm{tick}}(n)=2^{n+1} & O(n) & \text{status}\\ \hline
0 & 2 & 0 & \text{virtual only},\ \mathcal A(0)=\{J\}\\
1 & 4 & 0 & \text{pair aspects active}\\
2 & 8 & 0 & \text{abelian aspect activates}\\
3 & 16 & 0 & \text{exact saturation};\ \text{colour activates}\\
4 & 32 & 16 & \text{locked in}\ D^-\\
5 & 64 & 48 & \text{locked in}\ D^-\\
n\ge4 & 2^{n+1} & 2^{n+1}-16 & \text{locked in}\ D^-
\end{array}
\]

\textbf{Req.} A2; D4.1, D4.2; CT0.5, CT0.6, CT0.11 (Part 1).

\textbf{Pf.}

%dsp-level:0
\[
D_{\mathrm{tick}}(n)=2^{n+1}\ [\mathrm{CT0.5}];\qquad C_{\mathrm{cap}}=\dim_{\mathbb R}\mathfrak H(\mathbb C^4)=16\ [\mathrm{CT0.6}\Leftarrow\mathrm{A2}]
\]

%dsp-level:0
\[
n=4:\ \ O(4)=2^{5}-16=16=C_{\mathrm{cap}};\qquad
n\ge4:\ \ 2^{n+1}\ge32>16\implies O(n)=2^{n+1}-16>0
\]

%dsp-level:0
\[
\mathrm{D4.1}:\ n\ge4\implies\text{config}\in F_{D^-}=\bigoplus_{n\ge4}F_n;\qquad
\mathrm{D4.2}:\ \text{rate}\propto O(n)\implies\neg\exists\ \text{rung}\ n\ge4\ \text{escaping}
\]

%dsp-level:0
\[
\mathrm{CT0.11}:\ \text{admissible}\iff n\le3.\ \blacksquare
\]

\textbf{Phys.} Rung 4 demands $32$ registers against a budget of $16$ — exactly one extra channel quantum — and every higher rung demands more. All such configurations collapse into a single sink $D^-$ whose absorption grows as $2^{n+1}$: rung 3 is the last rung the universe can access, and above it there is only the lock.

\end{theorem}

\begin{corollary}[{\texorpdfstring{CT4.1a — universal saturation of bidirectional real-edge subgraphs}{CT4.1a - universal saturation of bidirectional real-edge subgraphs}}]\label{tag:CT4-1a}
Let $S$ be a bidirectional real-edge subgraph — both orientations of every used channel in $\mathcal R$ [A4] — exhausting the incident real edges at its vertices, with channel degree $d_S(v)$. Then $n(v)=2\,d_S(v)$ [CT0.11b (Part 1), D0.8], and: $(1)\ \ d_S(v)\ge2\implies n(v)\ge4\implies O\bigl(n(v)\bigr)>0$: every such vertex falls to $F_{D^-}$ [CT4.1, D4.1]; its content is absorbed, retained in $\mathfrak L$ and counted, $\neg$accessible [CT4.2, CT4.3, D4.2, D1.11] — no bidirectional vertex of channel degree $\ge2$ is unoversubscribed. $(2)\ \ S\ 2$-\textit{regular} (bidirectional cycles, any length $\ge3$): uniform minimal saturation $n(v)=4\ \forall v$; $O(4)=16=C_{\mathrm{cap}}$ — exactly one channel quantum over budget per vertex, independent of cycle length [CT0.6, CT4.1]. $(3)\ \ S\ d$-\textit{regular}: graded oversubscription $O(2d)=2^{2d+1}-16$ [CT0.6]:

%dsp-level:0
\[
O(4)=16,\qquad O(6)=112,\qquad O(8)=496,\qquad O(10)=2032,\qquad O(12)=8176 .
\]

$(4)$ The unoversubscribed real-edge family under directed rung counting is exactly $n(v)\le3$ [CT0.11b, CT0.6]: on bidirectional $S$ only single-channel vertices ($n=2$); mixed skeletons admit unidirectional attachments up to $n=3$; every vertex with $n(v)\ge4$ is saturated [CT0.11, CT4.1].

\textbf{Req (Cor).} A4; D0.8, D1.11, D4.1, D4.2; CT0.6, CT0.11, CT0.11b (Part 1); CT4.1, CT4.2, CT4.3; exact ladder arithmetic [math; machine: \texttt{checks/\allowbreak{}verify\_\allowbreak{}promotion\_\allowbreak{}G\_\allowbreak{}saturation.py}].

\textbf{Pf (Cor).}

%dsp-level:0
\[
n(v)=2\,d_S(v)\ [\mathrm{CT0.11b},\ \mathrm{D0.8}];\qquad
(1):\ d_S\ge2:\ n\ge4:\ O(n)=2^{n+1}-16>0\ [\mathrm{CT0.6}]\implies F_{D^-}\ [\mathrm{CT4.1},\ \mathrm{D4.1}]
\]

%dsp-level:0
\[
\mathrm{CT4.2}:\ \text{absorbing};\qquad\mathrm{CT4.3}/\mathrm{D4.2}:\ \text{retained}\subseteq\mathfrak L,\ \text{counted},\ \neg\text{accessible}\ [\mathrm{D1.11}]
\]

%dsp-level:0
\[
(2):\ d_S\equiv2:\ n\equiv4;\quad O(4)=2^{5}-16=16=C_{\mathrm{cap}}\ [\mathrm{CT0.6},\ \mathrm{CT4.1}];
\]

%dsp-level:0
\[
\qquad (3):\ O(2d)=2^{2d+1}-16:\ \ 2^{7}{-}16=112,\ \ 2^{9}{-}16=496,\ \ 2^{11}{-}16=2032,\ \ 2^{13}{-}16=8176\ [\text{exact}]
\]

%dsp-level:0
\[
(4):\ \mathrm{CT0.11b}\ \text{verbatim};\qquad n\ge4\implies O>0\implies\text{saturated}\ [\mathrm{CT0.11},\ \mathrm{CT4.1}]\ \blacksquare
\]

\end{corollary}

\begin{theorem}[{\texorpdfstring{CT4.2 — The lock is absorbing}{CT4.2 - The lock is absorbing}}]\label{tag:CT4-2}
\textbf{Thm.}

%dsp-level:0
\[
\forall n\le3:\quad\mathcal P_M^{(n)}\,F_{D^-}\,\mathcal P_M^{(n)}=0;
\]

%dsp-level:0
\[
\text{locked content}:\ \neg\exists\ \text{pole of}\ Z(\mathcal R,n)\ \ \forall n\le3
\]

%dsp-level:0
\[
\qquad (\text{one-way transfer w.r.t. the accessible sector}).
\]

\textbf{Req.} D4.1; CT0.13, CT0.14 (Part 1); CT2.5 (Part 2).

\textbf{Pf.}

%dsp-level:0
\[
\mathrm{D4.1}:\ \operatorname{supp}F_{D^-}\subseteq\mathcal K^{(D^-)};\qquad
\mathrm{CT0.13}:\ D^-\in\mathcal A(n)\iff n\ge4
\]

%dsp-level:0
\[
\mathrm{CT0.14}:\ \text{aspect blocks}=\text{orthogonal summands of}\ \mathcal K;
\]

%dsp-level:0
\[
\quad \mathcal K^{(D^-)}\not\subseteq\mathcal K_n\ \ (n\le3)\implies\mathcal P_M^{(n)}\big|_{\mathcal K^{(D^-)}}=0\implies\mathcal P_M^{(n)}F_{D^-}\mathcal P_M^{(n)}=0
\]

%dsp-level:0
\[
\mathrm{CT2.5}:\ \text{pole}\implies\text{nonzero support in}\ \mathcal K_n;\qquad
\text{locked content has none}\implies\neg\exists\ \text{accessible pole at any}\ n\le3.\ \blacksquare
\]

\textbf{Phys.} The lock is a one-way membrane: no accessible mask can bring $D^-$ content back, and locked content supports no accessible spectral pole at any rung. What falls in stays in — with respect to everything an accessible observer can measure.

\end{theorem}

\begin{theorem}[{\texorpdfstring{CT4.3 — Locked, not destroyed}{CT4.3 - Locked, not destroyed}}]\label{tag:CT4-3}
\textbf{Thm.}

%dsp-level:0
\[
\text{absorption conserves ledger content}:\quad\text{absorbed}\subseteq\mathfrak L;\qquad\neg\text{accessible}\ [\mathrm{D1.11}];\qquad\neg\exists\ \text{deleted entry};
\]

%dsp-level:0
\[
\text{oversubscription}\ \neg\text{seal}:\ \Delta n_{\mathcal R}=0;
\]

%dsp-level:0
\[
\text{locked content}:\ \ \text{physical norm}
\]

%dsp-level:0
\[
= 0\ \text{at accessible rungs}\ [\mathrm{CT2.8}]\ \ \wedge\ \ w(C)
\]

%dsp-level:0
\[
\ne 0\ \text{in the unmasked A8 sum}\ (\text{record survives}).
\]

\textbf{Req.} A5, A8; D0.7, D1.11, D4.2; CT0.3 (Part 1); CT2.8 (Part 2); CT4.1, CT4.2.

\textbf{Pf.}

%dsp-level:0
\[
\mathrm{D0.7}:\ \mathfrak L=\text{ordered record of retained events};\qquad
\mathrm{A5}:\ \text{sole exclusion}=U\ (\text{Haar draw}),\ \ U\notin\mathfrak L
\]

%dsp-level:0
\[
\mathrm{D4.2}:\ \text{absorbed}=\text{locked}:=\text{retained in}\ \mathfrak L\ \wedge\ \neg\text{accessible}\ [\mathrm{D1.11}]
\]

%dsp-level:0
\[
\mathrm{CT0.3}:\ \text{counters change only via seals};
\]

%dsp-level:0
\[
\qquad \mathrm{CT4.1}:\ \text{oversubscription}
\]

%dsp-level:0
\[
= \text{load/capacity mismatch at fixed rung}\implies\Delta n_{\mathcal R}
\]

%dsp-level:0
\[
= 0\implies\neg\text{seal}\implies\neg\exists\ \text{removed entry}
\]

%dsp-level:0
\[
\mathrm{CT2.8}:\ \text{null vectors}=\text{mask-excluded}\implies\text{locked}:\ \text{norm}\ 0\ \text{at accessible rungs}\ [\mathrm{CT4.2}]
\]

%dsp-level:0
\[
\mathrm{A8}:\ w(C)\ \text{retains nonzero value in the unmasked sum}\implies\text{record survives}.\ \blacksquare
\]

\textbf{Phys.} The lock is a loss of access, not of information: the ledger keeps every absorbed register with nonzero weight; only its recoverability — its norm in the accessible physical inner product — is forfeited. (Used again in CT4.7, CT4.8.)

\end{theorem}

\begin{theorem}[{\texorpdfstring{CT4.4 — The locked sector carries no gauge structure}{CT4.4 - The locked sector carries no gauge structure}}]\label{tag:CT4-4}
\textbf{Thm.}

%dsp-level:0
\[
D^-:\ \ \neg\exists\ \text{nonabelian gauge algebra};
\]

%dsp-level:0
\[
\qquad \neg\exists\ \text{abelian gauge algebra};
\]

%dsp-level:0
\[
\qquad \text{contributes}\ \neg\text{generator to any accessible algebra}.
\]

\textbf{Req.} D3.2; CT0.11, CT0.12, CT0.13 (Part 1); CT3.20 (Part 2); CT4.2.

\textbf{Pf.}

%dsp-level:0
\[
\text{gauge structure arises only from}:\ \ (a)\ \text{pair block};\quad (b)\ \text{slot algebras}\ \mathfrak u(n)
\]

%dsp-level:0
\[
(a):\ \ \mathrm{CT0.12}:\ 16\ \text{pair directions fully allocated to}\ \{J^\pm,G,W,H,D^+\}\implies\neg\exists\ \text{direction left for}\ D^-
\]

%dsp-level:0
\[
(b):\ \ \mathrm{CT3.20}/\mathrm{CT0.11}:\ \text{slot algebra realized only for}\ n\le3;
\]

%dsp-level:0
\[
\qquad \mathrm{CT0.13}+\mathrm{D3.2}:\ \text{strictly new content at rungs}\ 2,3\ \text{allocated to}\ \{\mathrm{EM},\mathrm{strong}\}
\]

%dsp-level:0
\[
n\ge4:\ \neg\text{channel-admissible}\implies\text{slot algebra}\ \neg\text{realized}
\]

%dsp-level:0
\[
\implies F_{D^-}:\ \text{block disjoint from every gauge block}=\text{scalar sink},\ \neg\text{internal directions};
\]

%dsp-level:0
\[
\qquad \mathrm{CT4.2}:\ \neg\text{probing by any accessible mask}.\ \blacksquare
\]

\textbf{Phys.} The dark sink is structureless: no charges, no generators, no self-interaction visible from the accessible side. That structurelessness is precisely what allows it to behave as a single simple fluid (CT4.5).

\end{theorem}

\medskip\noindent\textit{\texorpdfstring{CT4.4a — The realization domain of $\iota_n$: non-existence at $n\ge4$, persistence refuted}{CT4.4a - The realization domain of : non-existence at , persistence refuted}.}\ Structural-constraints record: see Appendix B, Section B.3.1 of the Addendum.\par\medskip

\phantomsection
\section[§G.II — The Two Dark Aspects]{\texorpdfstring{§G.II — The Two Dark Aspects}{SG.II - The Two Dark Aspects}}

\begin{theorem}[{\texorpdfstring{CT4.5 — Equation of state of the locked sector}{CT4.5 - Equation of state of the locked sector}}]\label{tag:CT4-5}
\textbf{Thm.}

%dsp-level:0
\[
T^{(D^-)}_{\mu\nu}=\rho_{D^-}\,u_\mu u_\nu;\qquad
p_{D^-}=0;\qquad
w_{D^-}=\frac{p}{\rho}=0;\qquad
\rho_{D^-}\ge0\ \ \bigl(=0\iff\neg\exists v:\ n(v)\ge4\bigr).
\]

\textbf{Req.} D1.5; CT1.22, CT1.23 (Part 2); CT4.1, CT4.2, CT4.3.

\textbf{Pf.}

%dsp-level:0
\[
\mathrm{CT1.22}:\ \ T=T^{\mathrm{bd}}+T^{\mathrm{tr}}\quad[\text{channels of}\ \mathrm{D1.5}]
\]

%dsp-level:0
\[
T^{\mathrm{tr}}\ \text{requires propagating chain with connection}\ [\mathrm{CT1.22}];
\]

%dsp-level:0
\[
\qquad \mathrm{CT4.2}:\ \text{locked}:\ \neg\text{accessible pole}\implies\neg\text{propagating mode}\implies\mathsf E_{\mathrm{tr}}=0\implies T^{\mathrm{tr}}=0
\]

%dsp-level:0
\[
\implies T^{(D^-)}=T^{\mathrm{bd}}=(\rho+p)u_\mu u_\nu+p\,g_{\mu\nu}\quad[\mathrm{CT1.22};\ u^\mu=\text{rest frame of the locked cell}]
\]

%dsp-level:0
\[
p=\text{spatial momentum flux in rest frame};\quad\neg\text{propagation}\implies\text{flux}\equiv0\implies p_{D^-}=0\implies T^{(D^-)}_{\mu\nu}=\rho_{D^-}u_\mu u_\nu;\qquad w=0/\rho=0
\]

%dsp-level:0
\[
\rho_{D^-}=\text{absorbed registers per volume}=\textstyle\sum O(n)\ge0\ [\mathrm{CT4.1}];
\]

%dsp-level:0
\[
\qquad \mathrm{CT4.3}:\ \neg\text{removal}\implies\rho_{D^-}\ge0;\quad=0\iff\neg\exists v:\ n(v)\ge4.\ \blacksquare
\]

\textbf{Phys.} Locked content gravitates as pressureless dust: no propagation, no momentum flux, $w=0$ — derived, not fitted. It couples to geometry through the field equations CT1.23 and through nothing else. Contrast with $D^+$: $w=-1$ (CT1.25, CT4.6).

\end{theorem}

\begin{theorem}[{\texorpdfstring{CT4.5a — The flicker-dust ensemble: exact density polynomial and rate-invariant $w=0$}{CT4.5a - The flicker-dust ensemble: exact density polynomial and rate-invariant}}]\label{tag:CT4-5a}
\textbf{Thm.} Hypotheses (ensemble scope, stated): $\Gamma$ a finite skeleton [A7] with potential degrees $n(v)$; every potential edge crystallizes independently with common probability $p\in[0,1]$ (a per-edge Bernoulli ensemble; the \textit{dynamics} realizing such an ensemble is not declared in Parts 1–2 — Merge note). Then: $(1)$ \textbf{Exact dust density.} $k_v\sim\mathrm{Bin}(n(v),p)$, and with $f_{\mathrm{adm}}(n,p):=\sum_{k\le3}\binom{n}{k}p^k(1-p)^{n-k}$ [CT0.11: admissible $\iff k\le3$], $O(n):=\max(0,\,2^{n+1}-16)$ [CT4.1, CT0.6]:

%dsp-level:0
\[
\langle\rho_{D^-}\rangle_\Gamma(p)=\sum_v O(n_v)\,\bigl(1-f_{\mathrm{adm}}(n_v,p)\bigr)
\]

— an exact polynomial in $p$ with \textbf{integer coefficients} [CT4.5's source criterion $\rho_{D^-}>0\iff\exists v:\ k_v\ge4$, read in expectation]. $(2)$ \textbf{Anchors.} $\langle\rho\rangle_\Gamma(1)=\sum_vO(n_v)$ — Platonic values $(448,\ 896,\ 2976,\ 2240,\ 24384)$ for (tetra, cube, octa, dodeca, icosa) [$(V,\deg)=(4,3),(8,3),(6,4),(20,3),(12,5)$; math]; at $p=\tfrac12$: $\bigl(154,\ 308,\ \tfrac{15159}{8},\ 770,\ 20193\bigr)$; the polygon family $P_n$ ($n\ge3$, $\deg=2$): the exact monomial $16\,n\,p^4$; the digon identically $0$ ($k=2$ never locks [CT0.6: $O(2)=0$]); $\langle\rho\rangle_\Gamma(0)=0$; strict increase on $(0,1)$ for every lockable $\Gamma$ ($\exists v:\ n(v)\ge4$). $(3)$ \textbf{The flicker modulates $\rho$, never $w$.} CT4.5 applies per configuration: every locked vertex contributes pressureless dust ($w=0$, source side, coupled to geometry through CT1.23 and nothing else). A mixture of $w=0$ components is $w=0$ on the locked part at every $p$ [CT1.22's linearity in the source]; $D^+$'s $w=-1$ [CT1.25, D4.3] is configuration-independent and untouched — CT4.6's contrast table holds ensemble-wise term by term. $(4)$ \textbf{No fluid.} The locked sector supports no further equation-of-state data: a structureless absorbing sink — one possibility space [D4.1], no gauge content [CT4.4], absorption one-way [CT4.2], content retained-and-inaccessible [CT4.3, D1.11]. Statements assigning $D^-$ a pressure response, polarization, or re-emission have no defined ledger referent (the referent would contradict CT4.2/D1.11); "the vacuum is a $D^-$ fluid" is inadmissible as stated. What is derived: locked content gravitates as dust, with the density computed here.

\textbf{Req.} Hypotheses (stated above); A7; D1.11, D4.1, D4.2, D4.3; CT0.6, CT0.11 (Part 1); CT1.22, CT1.23, CT1.25 (Part 2); CT4.1, CT4.2, CT4.3, CT4.4, CT4.5, CT4.6; binomial law under per-edge independence, exact polynomial arithmetic [math; machine: \texttt{checks/\allowbreak{}verify\_\allowbreak{}merge\_\allowbreak{}tierB.py} §E].

\textbf{Pf.}

%dsp-level:0
\[
(1):\ \text{locked}\iff k_v\ge4\ [\mathrm{CT4.5}\ \text{iff};\ \mathrm{CT0.11}:\ \text{admissible}\iff k\le3];
\]

%dsp-level:0
\[
\quad \mu_p(k_v\ge4)=1-f_{\mathrm{adm}}(n_v,p)\ [\mathrm{Bin}(n_v,p),\ \text{hypothesis}];
\]

%dsp-level:0
\[
\quad \bigl\langle\textstyle\sum_v O(n_v)\mathbf 1_{k_v\ge4}\bigr\rangle=\sum_v O(n_v)(1-f_{\mathrm{adm}})\ [\text{linearity};\ \mathrm{CT4.5}]
\]

%dsp-level:0
\[
\text{integer coefficients}:\ O(n)\in\mathbb Z\ [\mathrm{CT4.1}]\ \wedge\ \binom{n}{k}\in\mathbb Z
\]

%dsp-level:0
\[
\implies (1-f_{\mathrm{adm}})\ \text{expands with}\ \mathbb Z\text{-coefficients}
\]

%dsp-level:0
\[
(2):\ p{=}1:\ 1-f_{\mathrm{adm}}(n,1)=\mathbf 1_{n\ge4};\quad O(2d)=(16,112,496,2032)\ \text{at}\ 2d=(4,6,8,10)\ [\mathrm{CT4.1}]
\]

%dsp-level:0
\[
\text{Platonic sums}:\ (4\cdot112,\ 8\cdot112,\ 6\cdot496,\ 20\cdot112,\ 12\cdot2032)
\]

%dsp-level:0
\[
= (448,896,2976,2240,24384)\ [\text{machine}]
\]

%dsp-level:0
\[
p=\tfrac12:\ \text{exact binomial evaluation}\ [\text{machine}];
\]

%dsp-level:0
\[
\qquad P_n:\ n(v)=4\ \forall v\implies1-f_{\mathrm{adm}}(4,p)=p^4\implies\textstyle\sum_v16\,p^4=16np^4;
\]

%dsp-level:0
\[
\qquad \text{digon}:\ O(2)=0
\]

%dsp-level:0
\[
\text{monotonicity}:\ \tfrac{d}{dp}P(\mathrm{Bin}(n,p)\ge4)>0\ \text{on}\ (0,1)\ \text{for}\ n\ge4\ [\text{math};\ \text{machine on the rational grid}]
\]

%dsp-level:0
\[
(3):\ w=0\ \text{per configuration}\ [\mathrm{CT4.5}];
\]

%dsp-level:0
\[
\quad \langle T_{\mu\nu}\rangle=\textstyle\sum_c\mu(c)\,T_{\mu\nu}(c)\ [\mathrm{CT1.22}]:\ p_{\mathrm{eff}}=\sum_c\mu_c\cdot0=0\ \text{on the locked part};
\]

%dsp-level:0
\[
\quad \mathrm{CT1.25}/\mathrm{D4.3}:\ D^+\ \text{constant}
\]

%dsp-level:0
\[
(4):\ \text{the cited tags exhaust the}\ D^-\ \text{inventory};
\]

%dsp-level:0
\[
\quad \mathrm{CT4.2}:\ \text{one-way};
\]

%dsp-level:0
\[
\quad \mathrm{D1.11}:\ \text{inaccessible}\implies\neg\exists\ \text{response read}.\ \blacksquare
\]

\textbf{Phys.} Shake a crystal at any rate and its dark content has an exact polynomial density — integer coefficients, closed form, every Platonic and polygon anchor reproduced — and its equation of state never moves: dust at every rate. What the flicker changes is \textit{how much} content is locked, never \textit{what locked content does}. And what locked content does is fixed: it gravitates, it absorbs, it does not respond.

\end{theorem}

\begin{theorem}[{\texorpdfstring{CT4.6 — The two dark aspects contrasted}{CT4.6 - The two dark aspects contrasted}}]\label{tag:CT4-6}
\textbf{Thm.}

%dsp-level:0
\[
\begin{array}{c|c|c}
& D^+ & D^-\\ \hline
\text{definition} & F_{D^+}=\Lambda\,\mathbb I\ [\mathrm{D4.3}] & F_{D^-}=\bigoplus_{n\ge4}F_n\ [\mathrm{D4.1}]\\
\text{block} & \text{central pair direction}\ E_0\ [\mathrm{CT0.12}] & \text{stack block}\perp\text{pair space}\ [\mathrm{CT0.12\ Rmk}]\\
\text{activation} & n=1\ [\mathrm{CT0.13}] & n\ge4\ [\mathrm{CT0.13}]\\
\text{config.\ dependence} & \text{none}\ [\mathrm{D4.3}] & \text{full}\ [\mathrm{CT4.1}]\\
\text{rest frame} & \text{none; isotropic}\ [\mathrm{CT1.25}] & \text{one per locked cell}\ [\mathrm{CT4.5}]\\
\text{equation of state} & w=-1\ [\mathrm{CT1.25}] & w=0\ [\mathrm{CT4.5}]\\
\text{gauge content} & \text{none; central}\ [\mathrm{CT1.3}] & \text{none; disjoint}\ [\mathrm{CT4.4}]\\
\text{entropy} & \text{none}\ [\mathrm{CT4.7}] & \text{maximal, locked}\ [\mathrm{CT4.7}]
\end{array}
\]

%dsp-level:0
\[
\forall\ \text{rows differ}\implies D^+\ne D^-\ (\neg\text{two descriptions of one thing});
\]

%dsp-level:0
\[
\mathrm{CT1.23}:\ D^+\ \text{geometric side}\ (\Lambda g_{\mu\nu});\quad D^-\ \text{source side}\ (\rho_{D^-}u_\mu u_\nu);\qquad\neg\text{double counting}\ [\mathrm{CT1.22}].
\]

\textbf{Req.} D4.1, D4.3; CT0.12, CT0.13, CT1.3 (Part 1); CT1.22, CT1.23, CT1.25 (Part 2); CT4.1, CT4.4, CT4.5, CT4.7.

\textbf{Pf.} Each cell $=$ cited tag; rows differ pointwise; distinct entry into CT1.23 with no double count [CT1.22 step 5]. $\blacksquare$

\textbf{Phys.} The corpus’s two dark aspects differ on every structural axis: one is uniform cosmological tension on the geometric side ($w=-1$), the other is oversubscribed ledger sediment on the source side ($w=0$). “Dark” names an observational status, not a shared nature.

\end{theorem}

\begin{theorem}[{\texorpdfstring{CT4.7 — The entropy triangle}{CT4.7 - The entropy triangle}}]\label{tag:CT4-7}
\textbf{Thm.} In the plane $\bigl(S_{\mathrm{raw}},\ S_{\mathfrak L}\bigr)$ with $S_{\mathrm{raw}}(\rho)=-\mathfrak T(\rho\log\rho)$ [D0.3] and $S_{\mathfrak L}=n_{\mathcal R}\ln2$ [CT1.29], there are exactly three scalar extremes, mutually exclusive:

%dsp-level:0
\[
(1)\ \ \text{Janus vacuum}:\ n_{\mathcal R}=0:\quad\bigl(S_{\mathrm{raw}},S_{\mathfrak L}\bigr)=\bigl(\ln4\ \text{per pair (max)},\ 0\bigr);
\]

%dsp-level:0
\[
(2)\ \ D^+:\ \text{one configuration}:\quad(0,\ 0)\qquad(\ln1=0;\ \neg\text{real edges});
\]

%dsp-level:0
\[
(3)\ \ D^-:\ \text{locked}:\quad\bigl(\min_{\mathrm{acc}},\ \max\ \wedge\ \neg\text{accessible}\bigr).
\]

\textbf{Req.} A4; D0.3, D4.3; CT0.1, CT0.2, CT0.8 (Part 1); CT1.29 (Part 2); CT4.1, CT4.3, CT4.9.

\textbf{Pf.}

%dsp-level:0
\[
\text{two independent measures}:\ \ S_{\mathrm{raw}}=\text{pair mixedness}\ [\mathrm{D0.3}];\quad S_{\mathfrak L}=\text{registered real edges}\ [\mathrm{CT1.29}]
\]

%dsp-level:0
\[
(1):\ \ \mathrm{CT0.2}:\ n_{\mathcal R}=0\implies\rho=\rho_\circ=I/4,\ \ S_{\mathrm{raw}}=\ln4=\max\ (d=4);\qquad S_{\mathfrak L}=0\cdot\ln2=0
\]

%dsp-level:0
\[
(2):\ \ \mathrm{D4.3}:\ \text{single direction, single value}\ \Lambda\implies\ln1=0;\qquad\neg\text{own real edges}\implies S_{\mathfrak L}\text{-contribution}=0
\]

%dsp-level:0
\[
(3):\ \ \mathrm{CT4.1}:\ \text{locked vertex carries}\ 2^{n+1}>16\ \text{registers};
\]

%dsp-level:0
\[
\quad\mathrm{CT4.3}:\ \neg\text{deletion}
\]

%dsp-level:0
\[
\implies S_{\mathfrak L}^{\text{locked}}>\text{accessible bound}\ [\mathrm{CT4.9}],
\]

%dsp-level:0
\[
\ \ \to\infty\ \text{in}\ n
\]

%dsp-level:0
\[
\mathrm{CT0.1}+\mathrm{A4}:\ \text{interior pairs crystallized};\quad
\mathrm{CT0.8}\ (\varepsilon_q=0):\ \operatorname{supp}\rho\subseteq\vee^2\implies S_{\mathrm{raw}}\le\ln3<\ln4:\ \min_{\mathrm{acc}}
\]

%dsp-level:0
\[
\text{exclusivity}:\ \ (n_{\mathcal R}=0)\ \wedge\ (\exists v:\ n\ge4)\ \text{contradictory};\qquad D^+:\ \neg\text{configuration dependence at all}.\ \blacksquare
\]

\textbf{Phys.} The scalar landscape is a triangle in the (pair-entropy, ledger-entropy) plane: the vacuum at (max, 0), $D^+$ at (0, 0), $D^-$ at (min, max-and-hidden). Three distinct ways of being featureless — and by CT4.10, no fourth.

\end{theorem}

\begin{theorem}[{\texorpdfstring{CT4.7a — The vacuum census, closed, and the baseline equality}{CT4.7a - The vacuum census, closed, and the baseline equality}}]\label{tag:CT4-7a}
\textbf{Thm.} $(1)$ \textbf{The complete census.} The three scalar extremes of CT4.7, with CT4.10's exhaustiveness, are the complete inventory of vacuum-type scalar reads, mutually exclusive:

%dsp-level:0
\[
\begin{array}{c|cc|l}
& S_{\mathrm{raw}} & S_{\mathfrak L} & \text{structure}\\ \hline
\text{Janus vacuum}\ (n_{\mathcal R}=0) & \ln4\ \text{per pair (max, argmax unique}\ \tfrac14I) & 0 & [\mathrm{CT0.2},\ \mathrm{A5},\ \mathrm{D0.3};\ \mathrm{CT4.7}(1)]\\
D^+ & 0 & 0 & \text{pure, uniform, configuration-blind};\ w=-1\ [\mathrm{D4.3},\ \mathrm{CT1.25};\ \mathrm{CT4.7}(2)]\\
D^-\ \text{locked} & \min_{\mathrm{acc}} & \max\ \wedge\ \neg\text{accessible} & [\mathrm{CT4.2},\ \mathrm{CT4.3},\ \mathrm{D1.11};\ \mathrm{CT4.7}(3)]
\end{array}
\]

$\neg\exists$ a fourth extreme [CT4.10]; $F_H$ is a dynamical field, $\neg$an extreme [EW.27]. $(2)$ \textbf{Baseline equality.} For every finite skeleton $S$ [A7], let $\mathrm{VV}(S):=$ the all-virtual state on its potential-edge set ($n_{\mathcal R}=0$, every pair at $\rho_\circ=\tfrac14I$ [A5, CT0.2]) and $V_{\mathrm{raw}}:=$ the global vacuum [CT0.2]. Then every constructible vacuum read coincides:

%dsp-level:0
\[
S_{\mathfrak L}(\mathrm{VV}(S))=n_{\mathcal R}\ln2=0=S_{\mathfrak L}(V_{\mathrm{raw}})\ [\mathrm{CT1.29}];
\]

%dsp-level:0
\[
\qquad \Delta S_{\mathrm{raw}}=0\ \text{on both sides}\ [\mathrm{CT0.2}:\ \text{every pair at the unique argmax}]
\]

%dsp-level:0
\[
\rho_{D^-}=0\ \text{on both sides}\ [\mathrm{CT4.5}\ \text{iff}:\ \neg\exists v:\ k_v\ge4];
\]

%dsp-level:0
\[
\qquad S_{\partial\mathcal B}\ \text{has no argument on either side}\ (\neg\exists\ \text{closed real-edge surface}\ [\mathrm{CT1.29}\ \text{domain}])
\]

hence the vacuum-polarization deficit vanishes identically: $\Delta S_{\text{vac-pol}}(S)\equiv0\ \forall S$. The reason is registration, $\neg$arithmetic: virtual edges carry no registration — $\phi_e/\sigma_e$ are defined on $E_{\mathcal R}$ only [A7, D0.5], and the sole ledger exclusion is the Haar draw [A5: $U\notin\mathfrak L$]. $(3)$ \textbf{Honest scope.} The \textit{absolute} raw entropy $|\mathrm{pairs}|\cdot\ln4$ is shape-dependent; every constructible \textit{deficit} and read-difference is shape-independent and zero. No particle or cosmological identification [P5].

\textbf{Req.} A4, A5, A7; D0.3, D0.5, D1.11, D4.3; CT0.1, CT0.2, CT0.8 (Part 1); CT1.25, CT1.29 (Part 2); EW.27 (Part 2); CT4.1, CT4.2, CT4.3, CT4.5, CT4.7, CT4.9, CT4.10; exact arithmetic [math; machine: \texttt{checks/\allowbreak{}verify\_\allowbreak{}merge\_\allowbreak{}tierA.py} §D].

\textbf{Pf.}

%dsp-level:0
\[
(1):\ \mathrm{CT4.7}\ \text{rows verbatim};
\]

%dsp-level:0
\[
\qquad \mathrm{CT4.10}:\ \{\text{rotation-scalar directions}\}=\{E_0\to F_{D^+},\ \Sigma\to F_H,\ \text{vacuum configuration},\ D^-\ \text{sink}\};
\]

%dsp-level:0
\[
\quad F_H\ \neg\text{extreme}\ [\mathrm{EW.27}]
\]

%dsp-level:0
\[
(2):\ n_{\mathcal R}=0\implies S_{\mathfrak L}=n_{\mathcal R}\ln2=0\ [\mathrm{CT1.29}];
\]

%dsp-level:0
\[
\qquad \text{every pair at}\ \rho_\circ=\tfrac14I\ [\mathrm{A5}]:\ S_{\mathrm{raw}}=\ln4\ \text{per pair}=\max,\ \text{argmax unique}\ [\mathrm{CT0.2},\ \mathrm{D0.3}]\implies\Delta S_{\mathrm{raw}}=0
\]

%dsp-level:0
\[
k_v=0<4\ \forall v\implies\rho_{D^-}=0\ [\mathrm{CT4.5}\ \text{iff}];\qquad E_{\mathcal R}=\varnothing\implies\neg\exists\ \partial\mathcal B\implies S_{\partial\mathcal B}\ \text{undefined}
\]

%dsp-level:0
\[
U\notin\mathfrak L\ [\mathrm{A5}]:\ \text{Haar draws leave no ledger record}\implies\text{all-virtual states are ledger-indistinguishable from}\ V_{\mathrm{raw}}\ \text{on every}\ n_{\mathcal R}\text{-based read}.\ \blacksquare
\]

\textbf{Phys.} The vacuum ledger is closed: exactly three scalar ways to be empty, and every constructible read agrees across all of them — an all-virtual skeleton reads identically to the bare vacuum, shape by shape, because virtual edges were never registered at all. Raw entropy knows the shape; no constructible read consults it.

\end{theorem}

\begin{theorem}[{\texorpdfstring{CT4.10 — The triangle is exhaustive}{CT4.10 - The triangle is exhaustive}}]\label{tag:CT4-10}
\textbf{Thm.}

%dsp-level:0
\[
\neg\exists\ \text{fourth scalar extreme}:
\]

%dsp-level:0
\[
\qquad \{\text{rotation-scalar directions of the corpus}\}
\]

%dsp-level:0
\[
= \{E_0\to F_{D^+},\ \ \Sigma\to F_H,\ \ \text{vacuum configuration},\ \ D^-\ \text{sink}\};
\]

%dsp-level:0
\[
F_H\ \neg\text{extreme}\ (\text{dynamical field}).
\]

\textbf{Req.} CT0.2, CT0.12, CT1.2, CT1.10, CT1.11 (Part 1); EW.27 (Part 2); CT3.7 (Part 2); CT4.4.

\textbf{Pf.}

%dsp-level:0
\[
\mathrm{CT1.10}:\ \ \mathcal E_{++}\to\mathbf1\oplus\mathbf5\oplus\mathbf1\implies\text{pair-space rotation scalars}=\{E_0,\ \Sigma\}\ \text{exactly}
\]

%dsp-level:0
\[
\mathrm{CT1.11}:\ E_0\nparallel\Sigma;\qquad\mathrm{CT0.12}:\ E_0\mapsto F_{D^+},\ \ \Sigma\mapsto F_H
\]

%dsp-level:0
\[
\mathrm{EW.27}:\ F_H:\ \text{stable stationary background}+\text{massive excitation}\implies\text{dynamical field},\ \neg\text{limiting configuration}
\]

%dsp-level:0
\[
\mathrm{CT1.2}:\ J,\ W,\ G\ \text{carry spatial indices}\implies\neg\text{scalars}
\]

%dsp-level:0
\[
\mathrm{CT4.4}:\ D^-\ \text{supplies}\ \neg\text{generator};\qquad
\mathrm{CT3.7}:\ \mathfrak T(T_a)=0\implies\text{colour}\ \neg\text{scalar direction};\qquad
\text{abelian}=\text{centre, counted}
\]

%dsp-level:0
\[
\text{vacuum corner}=n_{\mathcal R}=0\ \text{configuration}\ (\neg\text{generator}),\ \ \text{unique}\ [\mathrm{CT0.2}]
\]

%dsp-level:0
\[
\implies\{\text{vacuum},\ D^+,\ D^-\}\ \text{complete}.\ \blacksquare
\]

\textbf{Phys.} No fourth scalar extreme exists: the operator content of the corpus offers exactly two rotation scalars (one spent on $D^+$, one dynamical as $F_H$), the vacuum is a unique configuration, and the sink is unique by construction. The entropy triangle closes.

\end{theorem}

\phantomsection
\section[§G.III — Bounds and Boundaries]{\texorpdfstring{§G.III — Bounds and Boundaries}{SG.III - Bounds and Boundaries}}

\begin{theorem}[{\texorpdfstring{CT4.9 — Accessible entropy bounds}{CT4.9 - Accessible entropy bounds}}]\label{tag:CT4-9}
\textbf{Thm.}

%dsp-level:0
\[
(1)\ \ \text{per vertex}:\quad\text{accessible}(v)\implies S(v)\le3\ln2\ \ \text{nats}=3\ \ \text{bits}\qquad(\text{attained at}\ n=3);
\]

%dsp-level:0
\[
(2)\ \ \text{per region}\ \mathcal B,\ \partial\mathcal B\ \text{closed real-edge surface of area}\ A:
\]

%dsp-level:0
\[
\quad S_{\mathrm{acc}}(\mathcal B)\le\frac{A}{\mathfrak h^2}\ \ \text{bits}=\frac{A}{\mathfrak h^2}\ln2\ \ \text{nats}.
\]

\textbf{Req.} A2; D0.8, D1.9, D4.2; CT0.9, CT0.11 (Part 1); CT1.29 (Part 2); CT4.1.

\textbf{Pf.}

%dsp-level:0
\[
(1):\ \ \mathrm{CT0.11}:\ n\le3\implies\bigl|\{e\in E_{\mathcal R}:e\ni v\}\bigr|\le3\ [\mathrm{D0.8}];
\]

%dsp-level:0
\[
\qquad \mathrm{CT0.9}:\ 1\ \text{bit/edge}\implies2^3\ \text{configs}\implies S\le3\ln2;\quad n=3:\ \text{attained}
\]

%dsp-level:0
\[
(2):\ \ \mathrm{CT1.29}:\ n_{\mathcal R}(\partial\mathcal B)=A/\mathfrak h^2;\quad S_{\partial\mathcal B}=n_{\mathcal R}(\partial\mathcal B)\ln2
\]

%dsp-level:0
\[
\mathrm{A2}+\mathrm{D1.9}:\ \text{ticks local};\ \text{registers cross only boundary edges}\implies S_{\mathrm{acc}}(\mathcal B)\le\text{boundary register count}
\]

%dsp-level:0
\[
\mathrm{CT4.1}+\mathrm{D4.2}:\ \text{interior excess}\implies\text{oversubscribed}\implies\text{locked}\implies\neg\text{contributing}.\ \blacksquare
\]

\textbf{Phys.} Accessible information is bounded twice over: three bits at any vertex, one bit per hop-area of any bounding surface — a holographic-type bound whose coefficient is exactly $\ln2$, derived from channel capacity alone.

\end{theorem}

\begin{theorem}[{\texorpdfstring{CT4.8 — Locked regions and their boundaries}{CT4.8 - Locked regions and their boundaries}}]\label{tag:CT4-8}
\textbf{Thm.} Interior locked in $D^-$, boundary $\partial\mathcal B$ of area $A$:

%dsp-level:0
\[
S_{\partial\mathcal B}=\frac{A}{\mathfrak h^2}\ln2\quad\Bigl(=\frac{A}{4G_N}\ [\mathrm{CT2.9}\ \text{naming}]\Bigr);
\]

%dsp-level:0
\[
\text{exterior measurements}=\text{residues of the accessible resolvent}\ [\mathrm{CT2.5}]:
\]

%dsp-level:0
\[
\quad \text{equal boundary register counts}\implies\text{indistinguishable interiors};
\]

%dsp-level:0
\[
\neg\exists\ \text{accessible interior singularity}\ [\mathrm{CT1.28}].
\]

\textbf{Req.} CT1.28, CT1.29, CT2.5, CT2.9 (Part 2); CT4.1, CT4.2, CT4.3, CT4.9.

\textbf{Pf.}

%dsp-level:0
\[
\mathrm{CT4.2}:\ \text{locked}:\ \neg\text{accessible pole};
\]

%dsp-level:0
\[
\qquad \mathrm{CT2.5}:\ \text{exterior measurement}=\text{residue of accessible resolvent}\implies\neg\exists\ \text{direct interior signal}
\]

%dsp-level:0
\[
\mathrm{CT4.9}\ \text{step}(2):\ \text{accessible registers of}\ \mathcal B=\text{boundary edges only};\qquad
\mathrm{CT1.29}:\ S_{\partial\mathcal B}=\frac{A}{\mathfrak h^2}\ln2
\]

%dsp-level:0
\[
\text{equal boundary content}\implies\text{equal accessible data}\implies\text{interiors indistinguishable}
\]

%dsp-level:0
\[
\mathrm{CT1.28}:\ \text{accessible rung}\le3\implies\neg\text{accessible curvature singularity};
\]

%dsp-level:0
\[
\qquad \mathrm{CT4.1}:\ \text{would-be divergence transferred to}\ D^-;
\]

%dsp-level:0
\[
\qquad \mathrm{CT2.9}:\ \tfrac{A}{\mathfrak h^2}\ln2=\tfrac{A}{4G_N}.\ \blacksquare
\]

\textbf{Phys.} A locked region is characterized, for the exterior, by its boundary area alone: exact boundary entropy, indistinguishable interiors, no interior singularity. The corpus’s black-hole-type statement — derived from masks and capacity, with the quarter-area form supplied by the CT2.9 calibration.

\end{theorem}

\begin{theorem}[{\texorpdfstring{CT4.11 — Global closure}{CT4.11 - Global closure}}]\label{tag:CT4-11}
\textbf{Thm.} With $\ \mathbf F_\oplus=\bigoplus_{\text{all aspects}}F_A$ and $M_\oplus=\bigoplus_n\mathcal P_M^{(n)}$:

%dsp-level:0
\[
Z\bigl[\mathbf F_\oplus,M_\oplus\bigr]\ \ \text{well defined};\qquad
\mathfrak S\bigl(Z[\mathbf F_\oplus,M_\oplus]\bigr)\ne\varnothing\qquad(\text{the object of A9 exists}).
\]

\textbf{Req.} A2, A8, A9; C12; D0.10, D2.8, D4.1; CT0.2, CT0.12, CT0.13, CT0.14 (Part 1); CT2.2, CT2.5, CT2.8 (Part 2); CT4.2.

\textbf{Pf.}

%dsp-level:0
\[
\text{domain}:\ \ Z\ \text{sums over}\ \mathfrak C(\mathcal R)\ [\mathrm{D0.10},\ \mathrm{C12}]
\]

%dsp-level:0
\[
\text{aspects}:\ \ \mathrm{CT0.12}:\ \text{pair}\ 16
\]

%dsp-level:0
\[
= \{J^\pm,G,W,H,D^+\};
\]

%dsp-level:0
\[
\quad\mathrm{CT0.13}:\ \text{stack}
\]

%dsp-level:0
\[
= \{\mathrm{EM},\mathrm{strong},D^-\}\implies\text{exactly the A8 list}\ (\neg\text{undefined},\ \neg\text{twice})
\]

%dsp-level:0
\[
\text{masks}:\ \ \mathrm{CT2.2}:\ \mathcal P_M^{(n)}\ \text{Hermitian idempotents};\quad\mathrm{CT0.14}:\ \text{strict flag},\ \dim=(6,16,17,25)<\infty
\]

%dsp-level:0
\[
\text{convergence}:\ \ \mathrm{D2.8}:\ \text{convergent sheet};
\]

%dsp-level:0
\[
\quad \mathrm{CT2.5}:\ \text{resolvent finite off poles}\ (\dim\mathcal K_n<\infty);
\]

%dsp-level:0
\[
\quad \mathrm{CT2.8}:\ \text{reflected sum}=\sum|w|^2\ge0,\ \text{finite for square-summable halves}
\]

%dsp-level:0
\[
\text{residual divergence}:\ \ \bigoplus_{n\ge4}F_n\ \text{collected into}\ D^-\ [\mathrm{D4.1}];
\]

%dsp-level:0
\[
\quad \mathrm{CT4.2}:\ D^-\perp\forall\ \text{accessible mask}\implies\text{accessible part of}\ Z=\text{finite blocks only}
\]

%dsp-level:0
\[
\text{nonemptiness}:\ \ \mathrm{CT0.2}:\ \text{global vacuum}\ (n_{\mathcal R}=0,\ \forall\ \text{pairs virtual},\ \mathcal A(0)=\{J\})\ \text{satisfies every axiom}\implies\in\mathfrak S\implies\mathfrak S\ne\varnothing.\ \blacksquare
\]

\textbf{Phys.} The total masked partition function named by A9 is well defined — every aspect accounted for exactly once, every mask a finite projection, every sum on the convergent sheet, the infinite tower sealed inside $D^-$ — and its solution set contains at least the empty universe. The falsifiability axiom has mathematical content.

\end{theorem}

\phantomsection
\section[§G.IV — Falsifiability and the Closure Audit]{\texorpdfstring{§G.IV — Falsifiability and the Closure Audit}{SG.IV - Falsifiability and the Closure Audit}}

\begin{theorem}[{\texorpdfstring{CT4.12 — Falsifiers}{CT4.12 - Falsifiers}}]\label{tag:CT4-12}
\textbf{Thm.} Observation of any row $\implies\neg(\mathrm{A1}\wedge\cdots\wedge\mathrm{A9})$ [A9 contrapositive; each row negates a derived parameter-free statement]:

%dsp-level:0
\[
\begin{array}{c|c}
\text{falsifier} & \text{contradicted tag}\\ \hline
\text{4th accessible slot index}\ (\text{colour algebra}\supsetneq\mathfrak{su}(3)) & \mathrm{CT3.20},\ \mathrm{CT0.11}\\
N_{\mathrm{gluon}}\ne8 & \mathrm{CT3.13}\\
\text{isolated unsealed colour charge} & \mathrm{CT3.18},\ \mathrm{CT3.22}\\
q\notin\{-1,0,+1\} & \mathrm{EM.9},\ \mathrm{EM.10},\ \mathrm{EW.36}\\
m_\gamma\ne0\ \ \vee\ \ N_{\mathrm{pol}}\ne2 & \mathrm{EM.6},\ \mathrm{EM.8}\\
m_W\ne m_Z\cos\theta_W\ \ (\rho\ne1) & \mathrm{EW.34}\\
e\ne g\sin\theta_W & \mathrm{EW.35}\\
\text{uniform vacuum aspect with}\ w\ne-1 & \mathrm{CT1.25}\\
\text{locked non-propagating component with}\ w\ne0 & \mathrm{CT4.5}\\
S_{\partial\mathcal B}\not\propto A\ \ \vee\ \ \text{coefficient}\ne\ln2\ \text{per hop area} & \mathrm{CT1.29}\\
\text{local higher-curvature corrections in}\ D=4 & \mathrm{CT2.7}\\
\nabla^\mu T_{\mu\nu}\ne0\ (\Lambda=\mathrm{const}) & \mathrm{CT1.24},\ \mathrm{CT1.26}\\
\text{accessible curvature singularity} & \mathrm{CT1.28}\\
S_{\mathrm{acc}}>1\ \text{bit per hop area of bounding surface} & \mathrm{CT4.9}
\end{array}
\]

\textbf{Rmk ($\neg$falsifiers).} Measured values of $g,\ g',\ \theta_W,\ e,\ h_0,\ \Lambda,\ m_W,\ m_Z,\ m_H,\ \mathfrak h$: configuration/calibration data [P1, P4] — $\neg$falsifying alone; only relations among them falsify. Falsifier rows resting on EM.20 inherit its derived phase alphabet $\alpha_e\in\{0,\pi\}$ [EM.20(1)]; no extra-axiomatic input enters there [A §0.2a].

\textbf{Req.} A9; P1, P4; all cited tags.

\textbf{Pf.}

%dsp-level:0
\[
\forall\ \text{rows}:\ \text{cited tag}=\text{theorem proved without free parameters}
\]

%dsp-level:0
\[
\text{inspection of cited proofs}:\ \ \neg\exists\ \mathfrak h\text{-dependence surviving}\ (\text{entropy rows fix the coefficient},\ \neg\text{the scale});
\]

%dsp-level:0
\[
\quad\neg\exists\ \mathrm{P4}\ \text{data used}
\]

%dsp-level:0
\[
\text{parameter-free derived statement}\ \vdash\ \mathrm{A1}\wedge\cdots\wedge\mathrm{A9}\implies\text{its negation}\implies\neg(\mathrm{A1}\wedge\cdots\wedge\mathrm{A9})
\]

%dsp-level:0
\[
\quad [\mathrm{A9}\ \text{contrapositive}].\ \blacksquare
\]

\textbf{Phys.} The corpus is falsifiable in the strict sense: fourteen crisp observation types, any single one of which demolishes the axiom set. Magnitudes never falsify — only relations and structures do. That asymmetry is deliberate and stated.

\end{theorem}

\begin{theorem}[{\texorpdfstring{CT4.13 — Closure audit}{CT4.13 - Closure audit}}]\label{tag:CT4-13}
\textbf{Thm.} Every quantity of the corpus lies in exactly one of three categories; the categorization is complete:

%dsp-level:0
\[
(\mathbf{I})\ \ \text{derived parameter-free data}\ (\neg\text{freedom},\ \neg\mathfrak h\text{-dependence};\ \text{numbers and constraints}):
\]

%dsp-level:0
\[
(7,3,3,3)\ [\mathrm{CT1.1}];
\]

%dsp-level:0
\[
\quad C_{\mathrm{cap}}=16\ [\mathrm{CT0.6}];
\]

%dsp-level:0
\[
\quad D_{\mathrm{tick}}(n)=2^{n+1}\ [\mathrm{CT0.5}];
\]

%dsp-level:0
\[
\quad n\le3\ [\mathrm{CT0.11}];
\]

%dsp-level:0
\[
\quad N_{\text{colour}}=3\ [\mathrm{CT3.20}];
\]

%dsp-level:0
\[
\quad N_{\mathrm{gluon}}=8\ [\mathrm{CT3.13}];
\]

%dsp-level:0
\[
N_{\mathrm{pol}}=2\ [\mathrm{EM.8}];\quad q\in\{-1,0,+1\}\ [\mathrm{EM.9},\mathrm{EW.36}];\quad\alpha_e\in\{0,\pi\}\ \text{per edge}\ [\mathrm{EM.20}];\quad\rho=1\ [\mathrm{EW.34}];
\]

%dsp-level:0
\[
w_{D^+}=-1\ [\mathrm{CT1.25}];
\]

%dsp-level:0
\[
\quad w_{D^-}=0\ [\mathrm{CT4.5}];
\]

%dsp-level:0
\[
\quad S=n_{\mathcal R}\ln2\ [\mathrm{CT1.29}];
\]

%dsp-level:0
\[
\quad f_{abc}\ [\mathrm{CT3.21}];\quad\mathfrak T(\mathcal G_A\mathcal G_B)=\tfrac12\delta_{AB}\ [\mathrm{EW.32},\mathrm{CT3.9}];
\]

%dsp-level:0
\[
\text{derived constraint}\ (\text{interval; sharp point}\ 0;\ \neg\exists\ \text{theorem depending on the value within it}):\quad\varepsilon_q
\]

%dsp-level:0
\[
\in [0,\tfrac14)\ [\mathrm{CT0.8}];
\]

%dsp-level:0
\[
(\mathbf{II})\ \ \text{named combinations of}\ \mathfrak h\ (\text{external comparison only}):
\]

%dsp-level:0
\[
\qquad G_N=\frac{\mathfrak h^2}{4\ln2};\quad\kappa^2=8\pi G_N=\frac{2\pi\,\mathfrak h^2}{\ln2}
\]

%dsp-level:0
\[
\qquad [\mathrm{CT2.9};\ \neg\exists\ \text{further dimensionful freedom}:\ \mathrm{CT2.11}];
\]

%dsp-level:0
\[
(\mathbf{IIIa})\ \ \text{configuration data --- independent core}\ (\mathrm{P4}):
\]

%dsp-level:0
\[
\qquad g,\ g'\ [\mathrm{EW.33}];\quad h_0\ [\mathrm{EW.20}];\quad c_W\ [\mathrm{EW.33}];\quad\Lambda\ [\mathrm{CT1.25}];\quad\varkappa,\ \beta,\ \alpha_\Omega\ [\mathrm{D2.1}];
\]

%dsp-level:0
\[
(\mathbf{IIIb})\ \ \text{configuration data --- functionally dependent}\ (\neg\text{further freedom}):
\]

%dsp-level:0
\[
\qquad \theta_W=\arctan(g'/g)\ [\mathrm{EW.24}];\quad e=\frac{gg'}{\sqrt{g^2+g'^2}}\ [\mathrm{EW.35}];\quad\mu_\star^2=c_W\varkappa^2\beta^2h_0^2\ [\mathrm{EW.33}];
\]

%dsp-level:0
\[
m_W^2=\mu_\star^2g^2/4\ [\mathrm{EW.20}];\quad m_Z^2=\mu_\star^2(g^2+g'^2)/4\ [\mathrm{EW.23}];\quad m_H^2=\delta^2S_{\mathrm{eff}}/\delta h^2\big|_{h_0}\ [\mathrm{EW.27}].
\]

%dsp-level:0
\[
\text{completeness}:\ \ \forall\ \text{symbol of the A-Part-4 master table}\in\mathbf I\cup\mathbf{II}\cup\mathbf{IIIa}\cup\mathbf{IIIb}\cup\{\text{structural objects}\};
\]

%dsp-level:0
\[
\qquad \mathbf I,\ \mathbf{II},\ \mathbf{IIIa}\ \text{pairwise disjoint};\quad\mathbf{IIIb}=f(\mathbf{IIIa})\ (\neg\text{independent data}).
\]

\textbf{Req.} P1–P5; every cited derivation tag (listed inline).

\textbf{Pf.}

%dsp-level:0
\[
\mathbf I:\ \ \forall\ \text{entries proved at the inline-cited tag from A1--A9 + earlier derived only}
\]

%dsp-level:0
\[
\mathbf{II}:\ \ \mathrm{CT2.9};\qquad\mathrm{CT2.11}:\ \mathfrak h\mapsto\lambda\mathfrak h=\text{symmetry}\implies\neg\exists\ \text{further dimensionful freedom}
\]

%dsp-level:0
\[
\mathbf{IIIa}/\mathbf{IIIb}:\ \ \mathrm{P4}\ (\text{core})\ +\ \text{functional dependencies at the inline-cited tags}\ (\mathbf{IIIb})
\]

%dsp-level:0
\[
\text{completeness}:\ \ \text{scan of the master symbol table}:\ \forall\ \text{symbol}
\]

%dsp-level:0
\[
= \text{structural}\ \vee\ \mathbf I\ \vee\ \mathbf{II}\ \vee\ \mathbf{III};
\]

%dsp-level:0
\[
\quad\neg\exists\ \text{symbol in two categories}
\]

%dsp-level:0
\[
\implies\text{corpus}
\]

%dsp-level:0
\[
= 1\ \text{calibration d.o.f.}\ (\mathfrak h)\ +\ \text{finite independent core}\ (\mathbf{IIIa})\ +\ \text{closed parameter-free prediction list}\ (\mathbf I).\ \blacksquare
\]

\textbf{Phys.} The final audit: every quantity the corpus derives is either a fixed number, a name for the single scale $\mathfrak h$, or data of a particular solution — nothing floats in between. The boundary between prediction and parameter is drawn exactly, and it is the same boundary CT4.12 turns into falsifiers.

\end{theorem}

\phantomsection
\section[Closure exports of Part 2 §G (consumed by Parts 3–4)]{\texorpdfstring{Closure exports of Part 2 §G (consumed by Parts 3–4)}{Closure exports of Part 2 SG (consumed by Parts 3-4)}}

%dsp-level:0
\[
\begin{array}{lll}
\mathrm{CT4.1(a)}\ O(n)=2^{n+1}-16;\ O(4)=16;\ n_{\mathrm{bidir}}=2d\ge4 &
\mathrm{CT4.2}\ \mathcal P_M^{(n\le3)}F_{D^-}\mathcal P_M^{(n)}=0 &
\mathrm{CT4.3}\ \text{locked}\subseteq\mathfrak L;\ w(C)\ne0\\
\mathrm{CT4.4}\ D^-:\ \neg\text{gauge structure} &
\mathrm{CT4.5}\ T^{(D^-)}_{\mu\nu}=\rho u_\mu u_\nu;\ w_{D^-}=0 &
\mathrm{CT4.6}\ D^+\ne D^-\ (8\ \text{axes})\\
\mathrm{CT4.7}/\mathrm{CT4.10}\ \text{entropy triangle}\ \{vac, D^+, D^-\}\ \text{complete} &
\mathrm{CT4.9}\ S_v\le3\ln2;\ S_{\mathcal B}\le\tfrac{A}{\mathfrak h^2}\ln2 &
\mathrm{CT4.8}\ \text{locked region}\equiv\text{boundary area}\\
\mathrm{CT4.11}\ Z[\mathbf F_\oplus,M_\oplus]\ \text{well defined};\ \mathfrak S\ne\varnothing &
\mathrm{CT4.12}\ 14\ \text{falsifiers} &
\mathrm{CT4.13}\ \text{categories}\ \mathbf I/\mathbf{II}/\mathbf{III}\ \text{exhaustive}
\end{array}
\]

\textbf{Corpus status.} $\mathbf A\to\mathbf B\to\mathbf C\to\mathbf D\to\mathbf E\to\mathbf F\to\mathbf G$ complete; A9 test surface $=$ CT4.12; audit $=$ CT4.13. Verification suite: \texttt{checks/\allowbreak{}run\_\allowbreak{}all.py} (exact arithmetic, no floating point).

$\blacksquare$


% Part3_Fermion_era.md — sanitized presentation layer (generated; see source/CONVERSION_LOG.md)
\chapter{\texorpdfstring{$\langle qu(b(i(t|t))ick)\rangle$ — \textbf{Part 3 · Fermion\_era} — Bound Chains, Termini, Floors, Carriers, Channels}{Qubitick - **Part 3 - Fermion _era** - Bound Chains, Termini, Floors, Carriers, Channels}}

\textbf{Dep.} Part 1 (axioms, conventions, qubit-era definitions, calibration, pair canon, ledger dynamics L0–L5); Part 2 (definitions D2.1–D4.3, geometry, electroweak, electromagnetic, colour, saturation) — cited by tag at use. Annex cited row-wise: \textbf{the Census annex} (Appendix A of this volume; the computed data cells: the Category-B tags FT.28, FT.32, FT.35, FT.42–FT.45, FT.49–FT.51, FT.63–FT.65, FT.69, FT.76, FT.78, FT.81, FT.83, FT.84, FT.96, FT.98, FT.99, FT.105, FT.114, FT.120). Format: Part 1 §0.4. \textbf{Namespace.} This part owns: Def-H.1–Def-H.43; FT.1–FT.114 (structural tags), FT.116–FT.131, FT.133–FT.139, FT.141–FT.145, FT.153–FT.160. Tiers F.1–F.42, F.44 in series order. \textbf{DAG law.} Citations resolve to Parts 1–2, Census cells, or earlier tags of this part; the marked forward pointers are scope conditions cited at use (the FT.14$\to$FT.22/FT.24 pattern). No citation of Part 4. \textbf{0.1 Conventions and units.} Bits regime: $D_{\mathrm{tick}}(n)=n{+}1$ bits [CT0.5]; one pair channel supplies $C_{\mathrm{cap}}=16$ register directions per tick [CT0.6, D0.9] $=4$ bits of Shannon content [C10: register count $\ne$ bit throughput]; all entropy/flux bounds in nats via $\ln2$ [C9]. $\widehat{ds}^{\,2}$ is dimensionless [D1.6, CT1.19a]; $\mathfrak h$ is the single free scale [D2.9, CT2.11]. Geodesic and continuum-limit forms inherit CT2.6's hypotheses (i)–(iii) at the point of use. \textbf{0.2 Resolution rules.} (i) This canon carries no ledger citations: every citation resolves to a tag of Parts 1–2, a Census cell, or a [math]/[machine] justification. (ii) Scope tokens: “Comm-min”, “Eff” mark the fixed-phase read of the tiers that carry them; both premises are dissolved at corpus level; a tag carrying them retains its fixed-phase scope, and the dynamic-phase translation is: quantization reads apply to the antialigned content ($\Theta\to\Theta_{\mathcal P}$ [CT1.19b]), with the uniform-antialigned specialization $\mathcal P=E_{\mathcal R}$ reproducing the literal forms [Def-EM.20, EM.20]. (iii) “M1/M2” are the materialization scope clauses of Def-H.20 (spectral hypotheses, constitutive of the Platonic cards). (iv) The interface surface of this part is Def-H.35–Def-H.43 and the constitutive/scope clauses cited therein. \textbf{0.3 Machine verification.} \texttt{checks/\allowbreak{}run\_\allowbreak{}all.py}; the per-tier content checkers \texttt{checks/\allowbreak{}verify\_\allowbreak{}H\_\allowbreak{}*.py} (exact arithmetic, no floating point); the era-structure checker \texttt{checks/\allowbreak{}verify\_\allowbreak{}era\_\allowbreak{}structure.py}.

\phantomsection
\section[§Cal-F — Free internal parameters of the fermion era…]{\texorpdfstring{§Cal-F — Free internal parameters of the fermion era (P4-class)}{SCal-F - Free internal parameters of the fermion era (P4-class)}}

Policy: Part 1 §E (P4: free internal configuration parameters, $\neg$constants of the theory; P5: no external identification; P1: $\mathfrak h$ the single free dimensionful scale). Every entry: dimensionless internal configuration data unless a dimension is stated; no value is fitted to external data; any numerical use beyond the loci listed requires a logged declaration first.

%tbl-level:3
{\scriptsize\setlength{\tabcolsep}{4pt}
\begin{longtable}{>{\raggedright\arraybackslash}p{\dimexpr0.2312\textwidth-2\tabcolsep\relax}|>{\raggedright\arraybackslash}p{\dimexpr0.1559\textwidth-2\tabcolsep\relax}|>{\raggedright\arraybackslash}p{\dimexpr0.1344\textwidth-2\tabcolsep\relax}|>{\raggedright\arraybackslash}p{\dimexpr0.4785\textwidth-2\tabcolsep\relax}}
\toprule
parameter & locus & dimension & consumers \\
\addlinespace[2pt]
\midrule
\endfirsthead
\toprule
parameter & locus & dimension & consumers \\
\addlinespace[2pt]
\midrule
\endhead
\bottomrule
\endlastfoot
$r=\allowbreak x_+\allowbreak /x_-$ (Janus rate ratio; equivalently $x_\pm$) & Def-H.26; FT.68(5); FT.70–FT.71 & $\mathfrak q^0$ & the flicker ensemble; the equilibrium form $p=\allowbreak x_+\allowbreak /(x_+\allowbreak +\allowbreak x_-)$ is derived [FT.108(4)], the ratio is the datum \\
$\hat\mu^2$ & Def-H.37(1) & $\mathfrak q^2$ [EW.33 packaging; D2.9] & the leading-order pole read [FT.33] \\
$\hat\mu'^2$ & Def-H.37(2) & $\mathfrak q$ [FT.54(4)] & the exact pole read [FT.54(4); Census FT.76] \\
$\alpha$ (Ansatz-B formal temperature) & Def-H.28 & $\mathfrak q^0$ & class masses $m_B=\allowbreak (K-\alpha\ln\Omega)\ln2/\mathfrak h$; no value declared; Ansatz-B results stand as exact $(K,\allowbreak a,\allowbreak b)$ classes \\
$\mu$ (family-channel overall scale) & FT.113; FT.119 & — & channel spectra; every ratio is $\mu$-free \\
$\mu_\nu$ (rung-one channel scale) & FT.143 & — & the rung-one spectrum; ratios exact, scale undeclared \\
$t,\allowbreak u,\allowbreak x_\pm,\allowbreak y_d,\allowbreak z_d,\allowbreak \tau_{\mathcal R},\allowbreak \tau_{\mathcal V}$ (formal fugacities) & Def-H.32(3) & $\mathfrak q^0$ & formal indeterminates over $\mathbb Q$; specializations marked at use [FT.108(4); Cor L0a] \\
\end{longtable}
}

\phantomsection
\section[Tier F.1 — Information Flux and Measurement Primitivity]{\texorpdfstring{Tier F.1 — Information Flux and Measurement Primitivity}{Tier F.1 - Information Flux and Measurement Primitivity}}

\begin{definition}[{\texorpdfstring{Def-H.1 — Observer–subnetwork channel}{Def-H.1 - Observer-subnetwork channel}}]\label{tag:Def-H-1}
\textbf{Def.} $o\in\mathcal Q$ a single node; $\Sigma\subset\mathcal Q$, $\Sigma\cap\{o\}=\varnothing$. The channel is

%dsp-level:0
\[
\gamma(o,\Sigma):=\bigl\{e\in\mathcal R:\ e\ \text{has one endpoint}\ o,\ \text{the other in}\ \Sigma\bigr\}\ \cup\ \bigcup_{P}E(P),
\]

the union over every valid real-edge path $P$ from $o$ to $\Sigma$ lying entirely in $\mathcal Q\setminus\Sigma$ except its terminal edge. Always:

%dsp-level:0
\[
\gamma(o,\Sigma)\cap\bigl(\Sigma\setminus\partial\Sigma\bigr)=\varnothing\qquad(\text{no internal edge of}\ \Sigma\ \text{is ever a member}).
\]

Hop and tick counts over an interval: $ \delta h(\gamma):=\sum_{e\in\gamma}\delta h(e)$ [D0.6]; $ \delta\tau(\gamma)$ identically from D0.7. Bookkeeping extension of A2 by summation over the boundary edge set — $\neg$new primitive (uses A2, D0.6, D0.7 only).

\textbf{Rmk (no internal-reaching reading).} The alternative reading of the second clause — once $o$ touches $\Sigma$, all of $\Sigma$'s internal edges qualify — is \textbf{excluded on citation-discipline grounds}: CT1.29 defines $\partial\mathcal B$ as a \textit{separating} real-edge surface, and CT4.9/CT4.8 prove from A2 that externally accessible content is capped by boundary registers alone. Def-H.1 extends A2 without new primitives, so it cannot overturn closed theorems. Reductio: the internal reading collapses $I_{\mathrm{hid}}\to0$ for any well-connected $\Sigma$ with one external contact, contradicting CT4.8's generically-nonzero area law.

\vspace{2pt}\noindent\rule{\textwidth}{0.4pt}\vspace{2pt}

\end{definition}

\begin{definition}[{\texorpdfstring{Def-H.2 — Physical (metric) channel separation}{Def-H.2 - Physical (metric) channel separation}}]\label{tag:Def-H-2}
\textbf{Def.} In a region where $g_{\mu\nu}=\Omega^2\bar g_{\mu\nu}$ [CT1.12] is defined:

%dsp-level:0
\[
\Delta\ell(\gamma):=\int_\gamma\sqrt{\bar g_{ij}\,dx^i\,dx^j}\;\Omega;\qquad
\Delta t(\gamma):=\int_\gamma\Omega\,dx^0 .
\]

\textbf{Rmk (role of $D^+$).} $\Omega=e^{\alpha_\Omega a_{00}}$ [CT1.12] and $F_{D^+}=\Lambda\mathbb I$ acts along the same central $E_0$ direction [CT1.25]: $D^+$'s contribution to channel separation is carried inside $\Omega$; $\neg$additive term. Built entirely from CT1.12's proven components; $\neg\exists$ new free parameter beyond $(\alpha_\Omega,\beta,\varkappa)$ [P4].

\vspace{2pt}\noindent\rule{\textwidth}{0.4pt}\vspace{2pt}

\end{definition}

\begin{definition}[{\texorpdfstring{Def-H.3 — Extractable information}{Def-H.3 - Extractable information}}]\label{tag:Def-H-3}
\textbf{Def.}

%dsp-level:0
\[
I_{o\leftarrow\Sigma}(\gamma):=\sup\bigl\{\mathsf E_{\mathrm{tr}}(\gamma'):\ \gamma'\subseteq\gamma\ \text{valid transit configuration over the interval}\bigr\},
\]

the supremum of transit entropy carried [EM.7: $\mathsf E_{\mathrm{tr}}=\delta h\ln2$, extended additively over $\gamma$]. The identification of extractable information with carried transit entropy is constitutive of this definition; coherence with the computed cards: channel payload stays FT.1-bounded [FT.36(2)], tuple labels are superselection data [Def-H.15(4)].

\vspace{2pt}\noindent\rule{\textwidth}{0.4pt}\vspace{2pt}

\end{definition}

\begin{theorem}[{\texorpdfstring{FT.1 — Information flux bound}{FT.1 - Information flux bound}}]\label{tag:FT-1}
\textbf{Thm.} \textit{(observer-channel form; the chain-level bound itself — jaws, entropy read, threshold, sharpness — is CT0.6a (Part 1)):} $(1)$ \textbf{Observer-channel extension.} With $\mathsf B(\gamma):=$ CT0.6a's bound evaluated on $\gamma$'s hop/tick counts [D0.6, D0.7]:

%dsp-level:0
\[
I_{o\leftarrow\Sigma}(\gamma)\ \le\ \mathsf B(\gamma)
\]

%dsp-level:0
\[
\qquad [\mathrm{Def\text{-}H.3}:\ \text{the sup over valid transit subconfigurations inherits the memberwise bound}]
\]

$(2)$ \textbf{Geodesic gauge form.} Geodesic $\gamma$, slowly varying $\Omega,\bar g_{ij}$: $ \delta h=\Delta\ell/(\Omega\,\mathfrak h)$, $ \delta\tau=\Delta t/\Omega$ [Def-H.2], hence

%dsp-level:0
\[
I_{o\leftarrow\Sigma}(\gamma)\ \le\ \frac{1}{\Omega}\,\min\Bigl(\frac{\Delta\ell(\gamma)}{\mathfrak h},\ 4\,\Delta t(\gamma)\Bigr)\ln2 .
\]

\textbf{Req.} Def-H.1–H.3; CT0.6a (Part 1); D0.6, D0.7; CT2.6 step 5 (\textbf{C}, piecewise smoothness).

\textbf{Pf.}

%dsp-level:0
\[
(1):\ \mathrm{Def\text{-}H.3}:\ I_{o\leftarrow\Sigma}(\gamma)
\]

%dsp-level:0
\[
= \sup\bigl\{\mathsf E_{\mathrm{tr}}(\gamma'):\ \gamma'\subseteq\gamma\ \text{valid transit configuration over the interval}\bigr\};
\]

%dsp-level:0
\[
\qquad \text{each}\ \gamma'\ \text{is an admissible transit path configuration}:\ \mathsf E_{\mathrm{tr}}(\gamma')\le\mathsf B(\gamma')\ [\mathrm{CT0.6a}]
\]

%dsp-level:0
\[
\text{the bound is monotone in}\ (\delta h,\delta\tau):\ \gamma'\subseteq\gamma\implies\mathsf B(\gamma')\le\mathsf B(\gamma)\implies\sup\le\mathsf B(\gamma)
\]

%dsp-level:0
\[
(2):\ \text{substitution}\ \delta h=\Delta\ell/(\Omega\mathfrak h),\ \delta\tau=\Delta t/\Omega\ [\mathrm{Def\text{-}H.2};\ \mathrm{CT2.6}\ \text{step}\ 5]\implies\text{physical form}.\ \blacksquare
\]

\textbf{Phys.} Knowledge through a channel is squeezed by two independent jaws — total edge entropy (spatial) and per-tick capacity (temporal). The jaws and their sharpness are pure ledger arithmetic at the foundation layer (CT0.6a); what this theorem adds is the observer-channel sup read (Def-H.1–H.3 layer) and $D^+$'s $1/\Omega$ rescaling of the physical form.

\vspace{2pt}\noindent\rule{\textwidth}{0.4pt}\vspace{2pt}

\end{theorem}

\begin{theorem}[{\texorpdfstring{FT.2 — Exhaustiveness of the real-edge measurement primitive}{FT.2 - Exhaustiveness of the real-edge measurement primitive}}]\label{tag:FT-2}
\textbf{Thm.} Within A1–A9, the only axiom-level mechanism increasing $o$'s information about $\Sigma$ is the addition of valid real edges to $\gamma(o,\Sigma)$ [Def-H.1]. Hence FT.1's bound is exhaustive over every axiomatic route to knowledge, and any two coarse-graining procedures on the same $\gamma$ yield the same bound.

\textbf{Req.} A1–A9; C11, C12; D0.10; Def-H.1.

\textbf{Pf.}

%dsp-level:0
\[
\mathrm{A2}:\ C(p,q)=1\mathfrak q\ \text{is the only channel-defining pair relation};
\]

%dsp-level:0
\[
\qquad \mathrm{A4}:\ \mathcal R\text{-membership is the sole knowability criterion}
\]

%dsp-level:0
\[
\mathrm{A5}\ [\mathrm{C11}:\ \text{node-level}\ \rho\mapsto U\rho U^\dagger,\ U\sim\mu_\star(\mathrm{SU}(2))]:\ \text{undriven single-node evolution};
\]

%dsp-level:0
\[
\ U\notin\mathfrak L:\ \text{structurally unrecorded}\implies\neg\text{readable}
\]

%dsp-level:0
\[
\mathrm{A6}\ [\mathrm{D0.3}:\ \mathcal T\ \text{on operators};\ \text{variance clause reserved}]:\ \text{internal map on an existing}\ \rho;
\]

%dsp-level:0
\[
\ \neg\exists\ \text{new inter-node relation}
\]

%dsp-level:0
\[
\mathrm{A7}:\ \phi:E_{\mathcal R}\to\Phi\ \text{defined only on existing real edges}:\ \text{structure carried by an edge},
\]

%dsp-level:0
\[
\ \neg\text{a means of establishing one}
\]

%dsp-level:0
\[
\mathrm{A8}\ [\mathrm{D0.10},\ \mathrm{C12}:\ Z\ \text{sums over}\ \mathfrak C(\mathcal R)]:\ \text{classifies, weights, sums existing configurations};
\]

%dsp-level:0
\[
\ \neg\text{manufactures edges}
\]

%dsp-level:0
\[
\mathrm{A9}:\ \text{global solution-membership condition}:\ \text{meta-level},\ \neg\text{per-observer channel}
\]

%dsp-level:0
\[
\implies\neg\exists\ \text{route to}\ o\text{'s knowledge of}\ \Sigma\ \text{bypassing real-edge membership}\ [\mathrm{A2},\mathrm{A4}].\ \blacksquare
\]

\textbf{Phys.} Measurement in this corpus has exactly one primitive: the real edge. Every claim about observation, information and hiding reduces to channel edge-counts — which is what makes the fermion program computational.

\vspace{2pt}\noindent\rule{\textwidth}{0.4pt}\vspace{2pt}

\end{theorem}

\begin{theorem}[{\texorpdfstring{FT.3 — Unresolvability of $D^-$-sourced subnetwork structure}{FT.3 - Unresolvability of -sourced subnetwork structure}}]\label{tag:FT-3}
\textbf{Thm.} $\Sigma,\Sigma'$ with $(\delta h,\delta\tau)$-identical channels to $o$, $\Sigma$ containing a $D^-$-configured node ($n\ge4$, CT4.1) and $\Sigma'$ none:

%dsp-level:0
\[
I_{o\leftarrow\Sigma}(\gamma)=I_{o\leftarrow\Sigma'}(\gamma);\qquad
I_{o\leftarrow\Sigma''}(\gamma)=I\bigl(\delta h(\gamma),\delta\tau(\gamma)\bigr)\ \ \forall\Sigma''\in[\gamma],
\]

%dsp-level:0
\[
[\gamma]:=\bigl\{\Sigma'':\ \gamma(o,\Sigma'')\ \text{has hop/tick counts}\ (\delta h(\gamma),\delta\tau(\gamma))\bigr\}
\]

%dsp-level:0
\[
\qquad (\text{exact and total over}\ [\gamma]).
\]

\textbf{Req.} Def-H.1, Def-H.3; EM.7 (Part 2); FT.1.

\textbf{Pf.}

%dsp-level:0
\[
\mathrm{EM.7}:\ D^-\text{-sourced edges}
\]

%dsp-level:0
\[
= \text{ordinary real edges}:\ \text{identical received data}\ (\neg\exists\ D^-\text{-emission rule})
\]

%dsp-level:0
\[
\mathrm{FT.1}:\ I\ \text{bounded by a function of}\ (\delta h,\delta\tau)\ \text{only}:\ \neg\text{reference to internal node count/topology/rung structure}
\]

%dsp-level:0
\[
\text{hypothesis}:\ (\delta h,\delta\tau)\ \text{identical}\implies\text{bounds identical};
\]

%dsp-level:0
\[
\qquad \mathrm{Def\text{-}H.3}:\ I=\sup\ \text{over configurations constrained only by these bounds}\implies I\ \text{identical on}\ [\gamma].\ \blacksquare
\]

\textbf{Cor FT.3.1 (computational form).} For any candidate fermion terminus $f$:

%dsp-level:0
\[
I_o(f)\ \le\ \min\bigl(\delta h(\gamma(o,f)),\ 4\,\delta\tau(\gamma(o,f))\bigr)\ln2,
\]

and this bound is the only quantity distinguishing $f$ from any terminus sharing its channel signature.

\textbf{Phys.} Locked ($D^-$) structure is invisible from outside by construction: two subnetworks differing arbitrarily in locked content but sharing a channel signature are observationally identical. Darkness is indistinguishability, exactly.

\vspace{2pt}\noindent\rule{\textwidth}{0.4pt}\vspace{2pt}

\end{theorem}

\phantomsection
\section[Tier F.2 — The Fermion Definition]{\texorpdfstring{Tier F.2 — The Fermion Definition}{Tier F.2 - The Fermion Definition}}

\begin{definition}[{\texorpdfstring{Def-H.4 — Fermion}{Def-H.4 - Fermion}}]\label{tag:Def-H-4}
\textbf{Def.} Let $c$ be a Janus-closed oriented chain [D1.4]: $\partial c=0\wedge\#J^+=\#J^-$ — a \textbf{bound} chain, hence fermionic content by D1.5 — realized at rung $n(v)\ge2$. A \textbf{fermion} is the terminus (basepoint) $v\in\{s(c),t(c)\}$ of $c$: a bound-chain endpoint, at either end, for any admissible $\delta h(c)$, characterized externally by whatever transit (bosonic-content, D1.5) channels connect $v$ to an observer [Def-H.1, Def-H.3] and by which aspects of A8 act nontrivially there.

\textbf{Rmk (floor: $\delta h(c)\ge2$, $n(v)\ge2$).} A $\delta h=1$/rung-1 instance is \textbf{excluded by Parts 1–2}:

%dsp-level:0
\[
\mathrm{A1}:\ p\ne q\implies\neg\exists\ \text{self-loops};\qquad
\mathrm{D1.4}:\ \text{closed}\iff\partial c=t(c)-s(c)=0\implies s(c)=t(c)=v
\]

%dsp-level:0
\[
\implies\delta h(c)=1\ \text{closed chain}=\text{self-loop}=\bot\implies\delta h(c)\ge2\ (\text{minimal closed chain}=\text{2-cycle}\ p\to q\to p)
\]

%dsp-level:0
\[
\text{closed chain based at}\ v\implies v\ \text{has}\ \ge1\ \text{in- and}\ \ge1\ \text{out-edge}\implies n(v)\ge2\ [\mathrm{D0.8}]
\]

%dsp-level:0
\[
\mathrm{D1.5}:\ \text{open propagating chains}
\]

%dsp-level:0
\[
= \text{transit}
\]

%dsp-level:0
\[
= \text{bosonic content}\implies\text{rung-1 nodes on open chains are}\ \neg\text{fermion termini}.
\]

There is \textit{one} object — the bound-chain terminus — varying only in $\delta h(c)\ge2$; FT.14's hypothesis $\delta h(c)\ge2$ is canonical, not ad hoc. No new primitive: $s(c),t(c)$ [D1.4], Janus-closed [D1.3], bosonic aspect [A8]; this definition fixes the referent of \textit{fermion} in the corpus; consumers: FT.4, FT.28, FT.144(1). Terminology per D1.5: the closed chain \textit{is} the fermionic content; what is bosonic/transit is the external channel connecting the terminus to an observer (Def-H.1/Def-H.3), not the chain itself. The floor is computed on in Tier F.8 (FT.24).

\vspace{2pt}\noindent\rule{\textwidth}{0.4pt}\vspace{2pt}

\end{definition}

\begin{theorem}[{\texorpdfstring{FT.4 — External state characterization of a fermion terminus}{FT.4 - External state characterization of a fermion terminus}}]\label{tag:FT-4}
\textbf{Thm.} $v$ a fermion [Def-H.4], $\gamma(o,v)$ per Def-H.1:

%dsp-level:0
\[
I_o(v)\le\min\bigl(\delta h(\gamma(o,v)),\ 4\,\delta\tau(\gamma(o,v))\bigr)\ln2,
\]

and this bound is the entire externally available description: $v,v'$ with the same $(\delta h,\delta\tau)$ are observationally identical to $o$.

\textbf{Req.} Def-H.1, Def-H.4; FT.1, FT.3.

\textbf{Pf.}

%dsp-level:0
\[
\mathrm{Def\text{-}H.4}:\ v=\text{single vertex}=\text{degenerate case}\ \Sigma=\{v\}\implies\mathrm{FT.1}\ \text{applies}
\]

%dsp-level:0
\[
\mathrm{FT.3}/\mathrm{Cor\ FT.3.1}:\ I\ \text{constant on}\ [\gamma];
\]

%dsp-level:0
\[
\quad\mathrm{Def\text{-}H.4}\ \text{constrains}\ \neg\text{internal structure}
\]

%dsp-level:0
\[
\implies v,
\]

%dsp-level:0
\[
v'\ \text{with equal signature indistinguishable}.\ \blacksquare
\]

\textbf{Cor FT.4.1 (state-counting ceiling).} Externally distinguishable states through a fixed channel signature:

%dsp-level:0
\[
N_{\mathrm{states}}\le2^{\,\min(\delta h,\ 4\,\delta\tau)}\qquad\bigl(=\min(\delta h,4\,\delta\tau)\ \text{bits of externally resolvable state}\bigr).
\]

$\neg$yet a derivation of spin, generation number, or any specific quantum number — a ceiling on the labels any channel could support.

\textbf{Phys.} A fermion seen from outside is its channel signature and nothing more: the corpus's measurement theory turns "particle" into "equivalence class of subnetworks under $(\delta h,\delta\tau)$".

\vspace{2pt}\noindent\rule{\textwidth}{0.4pt}\vspace{2pt}

\end{theorem}

\phantomsection
\section[Tier F.3 — Information Flux as Rate]{\texorpdfstring{Tier F.3 — Information Flux as Rate}{Tier F.3 - Information Flux as Rate}}

\begin{definition}[{\texorpdfstring{Def-H.5 — Information flux}{Def-H.5 - Information flux}}]\label{tag:Def-H-5}
\textbf{Def.}

%dsp-level:0
\[
\mathcal F(\gamma):=\frac{I_{o\leftarrow\Sigma}(\gamma)}{\delta\tau(\gamma)}\quad[\text{nats/tick};\ \mathrm{C9}];
\]

%dsp-level:0
\[
\qquad \mathcal F_{\mathrm{phys}}(\gamma):=\frac{I_{o\leftarrow\Sigma}(\gamma)}{\Delta t(\gamma)}=\frac{\mathcal F(\gamma)}{\Omega}
\]

%dsp-level:0
\[
\quad [\Delta t=\Omega\,\delta\tau,\ \Omega\ \text{constant along}\ \gamma,\ \mathrm{Def\text{-}H.2}].
\]

\textbf{Rmk ($\Omega$-dependence).} $\mathcal F$ is $\Omega$-invariant (pure graph rate); $\mathcal F_{\mathrm{phys}}\propto1/\Omega$; $\Omega$ is set by $D^+$ [CT1.12, CT1.25] — the precise mechanism coupling information flux to geometry.

\vspace{2pt}\noindent\rule{\textwidth}{0.4pt}\vspace{2pt}

\end{definition}

\begin{definition}[{\texorpdfstring{Def-H.6 — Information flux density}{Def-H.6 - Information flux density}}]\label{tag:Def-H-6}
\textbf{Def.} $N_\perp$ parallel real edges (transverse width); $ A_\perp=N_\perp\,\mathfrak h^2$ [CT1.29]; $ \mathcal F_\rho(\gamma):=\dfrac{\mathcal F(\gamma)}{A_\perp(\gamma)}$ (local, per unit area).

\vspace{2pt}\noindent\rule{\textwidth}{0.4pt}\vspace{2pt}

\end{definition}

\begin{theorem}[{\texorpdfstring{FT.5 — Flux bound (rate form of FT.1)}{FT.5 - Flux bound (rate form of FT.1)}}]\label{tag:FT-5}
\textbf{Thm.}

%dsp-level:0
\[
\mathcal F(\gamma)\le\min\!\Bigl(\frac{\delta h(\gamma)}{\delta\tau(\gamma)}\ln2,\ \ C_{\mathrm{cap}}\ln2\Bigr)=\min\!\Bigl(\frac{\delta h}{\delta\tau},\ 4\Bigr)\ln2;
\]

%dsp-level:0
\[
\qquad \mathcal F_{\mathrm{phys}}(\gamma)\le\frac{1}{\Omega}\min\!\Bigl(\frac{\delta h}{\delta\tau},\ 4\Bigr)\ln2 .
\]

\textbf{Req.} FT.1; Def-H.2, Def-H.5.

\textbf{Pf.} Divide FT.1 by $\delta\tau$; both constraints apply to every member of the supremum set [FT.1]; convert via $\Delta t=\Omega\,\delta\tau$. $\blacksquare$

\textbf{Phys.} The hop-rate caps flux combinatorially; the 4-bit channel caps it per tick; geometry ($\Omega$, i.e. $D^+$) rescales the physical rate without touching the combinatorial one.

\vspace{2pt}\noindent\rule{\textwidth}{0.4pt}\vspace{2pt}

\end{theorem}

\begin{theorem}[{\texorpdfstring{FT.6 — Conformal invariance of combinatorial flux}{FT.6 - Conformal invariance of combinatorial flux}}]\label{tag:FT-6}
\textbf{Thm.} $g_{\mu\nu}\mapsto\Omega'^2g_{\mu\nu}:\quad\mathcal F(\gamma)\ \text{invariant};\qquad\mathcal F_{\mathrm{phys}}\mapsto\mathcal F_{\mathrm{phys}}\cdot\Omega/\Omega'.$

\textbf{Req.} Def-H.5; CT0.9 (Part 1); D0.7.

\textbf{Pf.} $\mathcal F=I/\delta\tau$; $I$ bounded by edge counts [CT0.9], $\delta\tau$ a tick count [D0.7]: both $\Omega$-free [combinatorial]; $\mathcal F_{\mathrm{phys}}=\mathcal F/\Omega$ [Def-H.5]. $\blacksquare$

\textbf{Phys.} The ledger rate is conformally blind; only the physical rate feels $D^+$ — information bookkeeping is pre-geometric.

\vspace{2pt}\noindent\rule{\textwidth}{0.4pt}\vspace{2pt}

\end{theorem}

\begin{theorem}[{\texorpdfstring{FT.7 — Universal flux density bound}{FT.7 - Universal flux density bound}}]\label{tag:FT-7}
\textbf{Thm.}

%dsp-level:0
\[
\mathcal F_\rho(\gamma)\le\frac{\ln2}{\mathfrak h^2}\quad\text{per tick};\qquad
\mathcal F_{\rho,\mathrm{phys}}(\gamma)\le\frac{1}{4\,\Omega\,G_N}\quad\text{per unit proper time}.
\]

\textbf{Req.} CT0.9 (Part 1); EM.7 (Part 2); CT1.29, CT2.9 (Part 2); Def-H.2, Def-H.6.

\textbf{Pf.}

%dsp-level:0
\[
\text{each of}\ N_\perp\ \text{parallel edges}\le\ln2\ \text{per tick}\ [\mathrm{CT0.9},\mathrm{EM.7}]\implies\mathcal F\le N_\perp\ln2
\]

%dsp-level:0
\[
A_\perp=N_\perp\mathfrak h^2\ [\mathrm{CT1.29}]\implies\mathcal F_\rho\le\frac{N_\perp\ln2}{N_\perp\mathfrak h^2}=\frac{\ln2}{\mathfrak h^2}\qquad(N_\perp\ \text{cancels}:\ \text{universal})
\]

%dsp-level:0
\[
\mathfrak h^2=4G_N\ln2\ [\mathrm{CT2.9}];\quad/\Omega\ [\mathrm{Def\text{-}H.5}]:\quad\mathcal F_{\rho,\mathrm{phys}}\le\frac{\ln2}{\Omega\cdot4G_N\ln2}=\frac{1}{4\,\Omega\,G_N}.\ \blacksquare
\]

\textbf{Cor FT.7.1 (Bekenstein rate bound).} The universal flux-density bound is the Bekenstein entropy bound expressed as a rate: maximal information flux density $1/(4G_N)$ per unit proper time, up to the conformal factor; $D^+$ modulates the ceiling locally.

\textbf{Phys.} One bit per hop-area per tick — universally, independent of channel width — and in physical units the quarter-$G_N$ law appears as a \textit{rate}. No dependence on the detailed form of $G(e)$ at any rung: safe from the rung-3 edge-structure question.

\vspace{2pt}\noindent\rule{\textwidth}{0.4pt}\vspace{2pt}

\end{theorem}

\begin{theorem}[{\texorpdfstring{FT.8 — Flux-mediated coarse-graining connection}{FT.8 - Flux-mediated coarse-graining connection}}]\label{tag:FT-8}
\textbf{Thm.} Coarse-graining $\Sigma\to v$ is geometry-dependent: surviving information is bounded by the boundary flux, which depends on the geometry along $\gamma$ (via $\Omega$, $\mathfrak h$).

\textbf{Req.} FT.5, FT.7; Def-H.5, Def-H.6.

\textbf{Pf.} Coarse-graining $:=$ replace $\Sigma$'s internals by a terminus, retaining only what flows through $\gamma(o,v)$; retained rate $\le\mathcal F_{\mathrm{phys}}(\gamma)$ [FT.5], density $\le\ln2/(\Omega\mathfrak h^2)$ [FT.7] — both geometry-dependent $\implies$ same $\Sigma$, different ambient geometry $\Rightarrow$ different surviving content: coarse-graining is geometry-mediated, $\neg$ purely combinatorial. $\blacksquare$

\textbf{Phys.} This is the bridge the fermion construction needs: what a particle "keeps" of its subnetwork under coarse-graining is set by the geometry of the channel through which it is seen.

\vspace{2pt}\noindent\rule{\textwidth}{0.4pt}\vspace{2pt}

\vspace{2pt}\noindent\rule{\textwidth}{0.4pt}\vspace{2pt}

\end{theorem}

\phantomsection
\section[Tier F.4 — Rung-Dependent Edge Structure]{\texorpdfstring{Tier F.4 — Rung-Dependent Edge Structure}{Tier F.4 - Rung-Dependent Edge Structure}}

\begin{definition}[{\texorpdfstring{Def-H.7 — Rung-dependent edge data}{Def-H.7 - Rung-dependent edge data}}]\label{tag:Def-H-7}
\textbf{Def.} $G^{(n)}(e)$ rung-dependent; EM part [D3.5, EM.16] for $n\ge2$:

%dsp-level:0
\[
G^{(n)}_{\mathrm{EM}}(e)=\sigma_e\,\exp\!\bigl(i\,\alpha_{\mathrm{EM}}^{(n)}\,\widehat{ds}^{\,2}(e)\bigr),\qquad
\alpha_{\mathrm{EM}}^{(n)}=\pi\quad[\mathrm{FT.9};\ \mathrm{EM.20}\ \text{under Comm-min}+\text{Eff}].
\]

\textbf{On-shell/off-shell scope [FT.22–FT.25; CT1.19a].} On edges of materialized minimal bound chains — and, pending the completion argument, on on-shell edges generally: $\widehat{ds}^{\,2}(e)\in\mathbb Z$ [FT.24], hence the geometric factor $e^{i\pi\widehat{ds}^{\,2}(e)}=(-1)^{\widehat{ds}^{\,2}(e)}\in\{\pm1\}$ and $G(e)\in\{+1,-1\}$; on every materialized closed chain, the loop product is $+1$ [FT.22: $\Theta(c)\in2\mathbb Z$]. On generic off-shell configurations in $Z$'s summation domain: $\widehat{ds}^{\,2}(e)$ is a generic $O(1)$ real [CT1.19a] and $G(e)$ retains the full phase $e^{i\pi\widehat{ds}^{\,2}(e)}$ — materialization selects the integer sublattice, exactly as CT2.5 selects discrete spectra from continuous operator families.

%dsp-level:0
\[
\begin{array}{c|c|c}
\text{rung} & \text{active stack aspects} & \text{edge data}\\ \hline
2 & \mathrm{EM} & \sigma_e\ (\text{sign})\ +\ (-1)^{\widehat{ds}^{\,2}(e)}\ (\text{geometric parity, on-shell})\ [\mathrm{FT.9},\ \mathrm{FT.24}]\\
3 & \mathrm{EM}+\mathrm{strong} & \sigma_e\ +\ \widehat{ds}^{\,2}(e)\in\mathbb Z\ (\text{on-shell})\ +\ \tau(e):\ \text{colour transport label}\ [\mathrm{CT3.11};\ \mathrm{Def\text{-}H.17}]\ —\ \neg\exists\ \text{independent strong phase}\ [\mathrm{FT.29}(2)]\\
\ge4 & \mathrm{EM}+\mathrm{strong}+D^- & +\ D^-\ \text{saturation content}\ [\mathrm{D4.1},\ \mathrm{CT4.1}]\ (\text{locked})
\end{array}
\]

On the minimal real edge alone ($\widehat{ds}^{\,2}(e)=1$ [D1.1, Comm-min]), $G(e)=-\sigma_e$ needs only EM.20/FT.9 — no Tier F.8 input. Rung-3 strong factor: $G^{(3)}(e)=G^{(2)}(e)\otimes\tau(e)$ [Def-H.17, FT.29] — the slot algebra's $\mathfrak u(1)$ centre \textit{is} the EM phase already carried by $G^{(2)}$ [CT3.7], so the strictly new strong content is the traceless $\mathfrak{su}(3)$ label, observable only in singlet sealings [FT.29(3)].

\vspace{2pt}\noindent\rule{\textwidth}{0.4pt}\vspace{2pt}

\end{definition}

\begin{theorem}[{\texorpdfstring{FT.9 — EM part of edge weight is rung-stable}{FT.9 - EM part of edge weight is rung-stable}}]\label{tag:FT-9}
\textbf{Thm.} $ \alpha_{\mathrm{EM}}^{(n)}=\alpha_{\mathrm{EM}}^{(2)}=\pi\quad(n=2,3,\ge4;\ \text{under EM.20's scope tokens Comm-min}+\text{Eff}).$

\textbf{Req.} EM.17, EM.18, EM.19, EM.20 (Part 2); CT0.14 (Part 1); CT3.7, CT3.10 (Part 2); D4.1.

\textbf{Pf.}

%dsp-level:0
\[
\mathrm{EM.20}\ [\text{Comm-min}+\text{Eff}]:\ \alpha_{\mathrm{EM}}^{(2)}=\pi
\]

%dsp-level:0
\[
\mathrm{EM.18}:\ \alpha^{(3)}=\alpha^{(2)}:
\]

%dsp-level:0
\[
\quad Y_{\mathrm{asp}}F_{\mathrm{strong}}Y_{\mathrm{asp}}=\mathcal P^{(\mathrm{EM})}F_{\mathrm{strong}}\mathcal P^{(\mathrm{EM})}=0\ [\mathrm{CT0.14}:\ \mathcal K^{(\mathrm{strong})}\perp\mathcal K^{(\mathrm{EM})}]
\]

%dsp-level:0
\[
\text{independently}:\ \mathfrak T(T_a)=0\ [\mathrm{CT3.7}]\implies\mathfrak T(F_{\mathrm{strong}})=0\ [\mathrm{CT3.10}]
\]

%dsp-level:0
\[
\mathrm{EM.19}:\ n\ge4:\ \ Y_{\mathrm{asp}}F_{D^-}Y_{\mathrm{asp}}=0\ [\mathrm{D4.1}:\ F_{D^-}\ \text{acts only on}\ n\ge4;\ \ \mathrm{CT0.11}:\ \mathcal K_2\ \neg\text{oversubscribed};\ \mathrm{CT0.14}:\ \perp]
\]

%dsp-level:0
\[
\mathrm{EM.17}:\ \alpha^{(n)}=\mathfrak T\bigl(Y_{\mathrm{asp}}\mathbf F_M^{(n)}Y_{\mathrm{asp}}\bigr)/\mathfrak T\bigl(Y_{\mathrm{asp}}\bigr);
\]

%dsp-level:0
\[
\quad Y_{\mathrm{asp}}\mathbf F_M^{(n)}Y_{\mathrm{asp}}=Y_{\mathrm{asp}}\mathbf F_M^{(2)}Y_{\mathrm{asp}}\implies\alpha^{(n)}=\alpha^{(2)}=\pi.\ \blacksquare
\]

\textbf{Phys.} The binary geometric phase of every real edge is the same $\pi$ at every accessible rung: colour and saturation add structure without touching the electromagnetic bit.

\vspace{2pt}\noindent\rule{\textwidth}{0.4pt}\vspace{2pt}

\end{theorem}

\begin{theorem}[{\texorpdfstring{FT.10 — Confinement is a global sum-restriction}{FT.10 - Confinement is a global sum-restriction}}]\label{tag:FT-10}
\textbf{Thm.}

%dsp-level:0
\[
Z^{(3)}(\mathcal R)=\sum_{\substack{C\in\mathfrak C(\mathcal R):\ \text{colour-sealed}}}\ \sum_\sigma\mathcal J(C)\!\!\prod_{e\in E(C)}\!\!G^{(3)}(e)\ (-1)^{\aleph(\sigma,C)}
\]

%dsp-level:0
\[
\qquad [\mathrm{D0.10},\ \mathrm{C12}];
\]

confinement restricts the summation domain; $\neg$modification of any individual edge weight.

\textbf{Req.} CT3.17, CT3.18, CT3.22 (Part 2); CT2.3, CT2.5 (Part 2); FT.9; D0.10.

\textbf{Pf.}

%dsp-level:0
\[
\mathrm{CT3.18}:\ \mathcal P_M^{(3)}\ \text{admits only colour-sealed}\ C\ (\delta^i{}_j,\varepsilon_{ijk},\varepsilon^{ijk}\ \text{contractions},\ \mathrm{CT3.17});
\]

%dsp-level:0
\[
\qquad \mathrm{CT3.22}:\ \neg\exists\ \text{isolated triplet pole}
\]

%dsp-level:0
\[
\text{free-index}\ C\notin\operatorname{Im}\mathcal P_M^{(3)};
\]

%dsp-level:0
\[
\quad\mathrm{CT2.5}:\ \text{poles}
\]

%dsp-level:0
\[
= \text{eigenvalue conditions on}\ \mathbf F_M^{(3)};
\]

%dsp-level:0
\[
\quad\mathrm{CT2.3}:\ \mathbf F_M^{(3)}\big|_{\mathcal K_3^\perp}
\]

%dsp-level:0
\[
= 0\implies\text{excluded}\ C:\ \neg\text{pole}
\]

%dsp-level:0
\[
G^{(3)}(e)
\]

%dsp-level:0
\[
= G^{(2)}(e)\otimes|i(e)\rangle\text{-label}:\ \text{per-edge weight unmodified}\ [\mathrm{FT.9}]\implies\text{confinement}
\]

%dsp-level:0
\[
= \text{domain restriction}.\ \blacksquare
\]

\textbf{Phys.} Confinement acts globally — on which configurations exist in the sum — not locally on edge amplitudes: the partition function itself only contains colour singlets.

\vspace{2pt}\noindent\rule{\textwidth}{0.4pt}\vspace{2pt}

\end{theorem}

\phantomsection
\section[Tier F.5 — Spin–Charge Decomposition via Complete Force…]{\texorpdfstring{Tier F.5 — Spin–Charge Decomposition via Complete Force Operator}{Tier F.5 - Spin-Charge Decomposition via Complete Force Operator}}

\begin{theorem}[{\texorpdfstring{FT.11 — Complete force operator charge audit}{FT.11 - Complete force operator charge audit}}]\label{tag:FT-11}
\textbf{Thm.} Per-aspect charge determination for $\mathbf F$ [A8]:

%dsp-level:0
\[
\begin{array}{c|c|c}
\text{aspect} & \text{independent charge?} & \text{basis}\\ \hline
F_{\mathrm{EM}} & \text{yes} & \mathrm{EM.9}:\ q_{\mathrm{EM}}=(\sigma_s-\sigma_t)/2\in\{-1,0,+1\},\ \text{additive end-charge}\\
F_{\mathrm{strong}} & \text{no (confined)} & \mathrm{CT3.22};\ \mathfrak T(T_a)=0\ [\mathrm{CT3.7}];\ \mathfrak T(F_{\mathrm{strong}})=0\ [\mathrm{CT3.10}]\\
F_G & \text{yes} & \mathrm{D2.5}:\ m^2=\text{pole displacement}\\
F_H & \text{yes (via mass)} & \mathrm{EW.13}:\ H=h\Sigma;\ \mathrm{EW.20},\mathrm{EW.33}\\
F_{D^+} & \text{no} & \mathrm{D4.3}:\ \Lambda\mathbb I,\ \neg\text{config.-dependent}\\
F_{D^-} & \text{no} & \mathrm{CT4.4}:\ \neg\text{gauge structure}\\
F_{J^\pm} & \text{no} & \mathrm{CT0.7}:\ \text{seals},\ \neg\text{charges}\\
F_W & \text{no} & \mathrm{FT.12}:\ \text{no weak self-charge}
\end{array}
\]

\textbf{Req.} A8, C8; CT0.7, CT0.12 (Part 1); D2.5, D4.3; CT3.7, CT3.10, CT3.22 (Part 2); CT4.4 (Part 2); EW.13, EW.20, EW.33 (Part 2); EM.9 (Part 2); FT.12.

\textbf{Pf.}

%dsp-level:0
\[
\mathrm{A8}\ [\mathrm{C8}]:\ \mathbf F
\]

%dsp-level:0
\[
= \mathbf F_{(+,\mathcal R)}\oplus\mathbf F_{(-,\mathcal R)}\oplus\mathbf F_{(+,\mathcal V)}\oplus\mathbf F_{(-,\mathcal V)}\oplus\mathbf F_{\mathrm{st}};
\]

%dsp-level:0
\[
\quad\mathbf F_{(-,\mathcal R)}
\]

%dsp-level:0
\[
= F_{W^+}\oplus F_{W^-}\oplus F_Z
\]

%dsp-level:0
\[
\mathrm{CT0.12}:\ \text{pair aspects exhaust}\ 16
\]

%dsp-level:0
\[
= 2^{C_{\mathrm{cap}}}\ \text{states}\ (3{+}3{+}5{+}3{+}1{+}1);
\]

%dsp-level:0
\[
\quad\text{stack aspects act on vertex/slot registers}\ [\mathrm{CT0.12\ Rmk}]
\]

%dsp-level:0
\[
\text{rows}:\ \text{as cited (each row one frozen-corpus tag)};\quad F_W:\ \mathrm{FT.12}.\ \blacksquare
\]

\textbf{Phys.} A complete audit of which forces can carry an independently readable charge: only EM (boundary-readable), gravity and Higgs (as mass). Everything else is structurally chargeless — the input to the spin–charge split.

\vspace{2pt}\noindent\rule{\textwidth}{0.4pt}\vspace{2pt}

\end{theorem}

\begin{theorem}[{\texorpdfstring{FT.12 — Weak has no self-charge}{FT.12 - Weak has no self-charge}}]\label{tag:FT-12}
\textbf{Thm.} The weak sector produces $\neg$independent abelian charge; its charge content reduces entirely to $Q=W_3+Y$ (the EM generator).

\textbf{Req.} EW.3, EW.22, EW.36, EW.36a (Part 2); D3.4.

\textbf{Pf.}

%dsp-level:0
\[
\mathrm{EW.22}:\ \ker M^2=\operatorname{span}\{Q\}_{(W_3,Y)},\ \ Q:=W_3+Y\ \ (1\text{-dim})
\]

%dsp-level:0
\[
\mathrm{EW.36}:\ \operatorname{spec}Q|_{\mathcal H_{\mathsf v}}
\]

%dsp-level:0
\[
= \{0,+1\};
\]

%dsp-level:0
\[
\quad\omega\ [\mathrm{D2.6}]:\ \{0,-1\};
\]

%dsp-level:0
\[
\quad\text{alphabet}
\]

%dsp-level:0
\[
= \{-1,0,+1\}
\]

%dsp-level:0
\[
= \mathrm{D3.4}\ \text{end-charge alphabet}
\]

%dsp-level:0
\[
\mathrm{EW.36a}:\ Q=W_3+Y;\quad W_3\ \text{alone}\ \neg\text{charge};\quad Q=\text{EM generator}
\]

%dsp-level:0
\[
\mathfrak{su}(2)\ \text{nonabelian},
\]

%dsp-level:0
\[
\ Z(\mathfrak{su}(2))
\]

%dsp-level:0
\[
= \{0\}\ [\mathrm{EW.3}:\ [W_i,W_j]=i\epsilon_{ijk}W_k]\implies\neg\exists\ \text{central}\ \mathfrak u(1)\subseteq\mathfrak{su}(2)\implies\neg\text{independent abelian weak charge}.\ \blacksquare
\]

\textbf{Phys.} Weak carries no charge of its own — its only abelian direction \textit{is} electromagnetism. Chargelessness of the weak sector is a Lie-algebra fact (centre-free $\mathfrak{su}(2)$), not a convention.

\vspace{2pt}\noindent\rule{\textwidth}{0.4pt}\vspace{2pt}

\end{theorem}

\begin{theorem}[{\texorpdfstring{FT.13 — Weak $=$ spin (algebraically)}{FT.13 - Weak spin (algebraically)}}]\label{tag:FT-13}
\textbf{Thm.} $W_i=\tfrac12\bigl(\sigma_i^{(p)}+\sigma_i^{(q)}\bigr)=$ total-spin generators of the pair; Casimir eigenvalues $2=j(j{+}1)|_{j=1}$ on $\vee^2$, $0$ on $\wedge^2$; $W_i\in\mathcal E_{+-}$ (weak sector). Hence weak $\equiv$ spin, algebraically.

\textbf{Req.} CT0.7a, CT1.1, CT1.2 (Part 1).

\textbf{Pf.} Identities [CT1.2 (definition), CT0.7a (Casimir), CT1.1 (sector)]. $\blacksquare$

\textbf{Phys.} With FT.12 (no self-charge), the weak sector's entire coarse-grained contribution is spin — a positive identification, not a residual.

\vspace{2pt}\noindent\rule{\textwidth}{0.4pt}\vspace{2pt}

\end{theorem}

\begin{definition}[{\texorpdfstring{Def-H.8 — Spin (emergent weak contribution)}{Def-H.8 - Spin (emergent weak contribution)}}]\label{tag:Def-H-8}
\textbf{Def.} For a fermion terminus $v$ [Def-H.4, $\delta h(c)\ge2$]: spin $:=$ the coarse-grained emergent-spacetime property arising from the weak sector's contribution to $c$. Justification: weak has no self-charge [FT.12] $\Rightarrow$ its contribution cannot be charge; $W_i$ are total-spin generators [FT.13] $\Rightarrow$ the contribution \textit{is} spin. \textbf{Rmk (coarse-graining hierarchy).} Pre-dimensionally: spin $=$ interior edge phase carried by the weak sector. Coarse-grained to the terminus: interior hidden [FT.3] $\Rightarrow$ cumulative interior phase manifests as an intrinsic localized quantum number. Further coarse-grained along a geodesic: angular momentum.

\vspace{2pt}\noindent\rule{\textwidth}{0.4pt}\vspace{2pt}

\end{definition}

\begin{theorem}[{\texorpdfstring{FT.14 — Spin–charge decomposition}{FT.14 - Spin-charge decomposition}}]\label{tag:FT-14}
\textbf{Thm.} $c$ oriented chain, $\delta h(c)\ge2$ [Def-H.4], proper separation $\Delta\ell(c)$ [CT1.12]. Decompose $N_{\mathrm{visible}}=\delta h(c)$ bits under the complete force audit [FT.11]:

%dsp-level:0
\[
(1)\ \ N_{\mathrm{spin}}=N_{\mathrm{weak}};
\]

%dsp-level:0
\[
\qquad (2)\ \ N_{\mathrm{charge}}=N_{\mathrm{visible}}-N_{\mathrm{weak}};
\]

%dsp-level:0
\[
\qquad (3)\ \ f_{\mathrm{charge}}=\frac{N_{\mathrm{charge}}}{N_{\mathrm{visible}}},\quad f_{\mathrm{spin}}=\frac{N_{\mathrm{weak}}}{N_{\mathrm{visible}}};
\]

%dsp-level:0
\[
\mathfrak h\to0,\ \Delta\ell\ \text{fixed}\ [\mathrm{CT1.19}]:\quad f_{\mathrm{charge}}\to0,\quad f_{\mathrm{spin}}\to1 .
\]

\textbf{Req.} A1, A2, A4; C8, C9, C10; CT0.1, CT0.6, CT0.7, CT0.9 (Part 1); CT1.12, CT1.19, CT1.25 (Part 2); D3.5; EW.13, EW.17 (Part 2); EM.9, EM.16 (Part 2); FT.9; FT.22–FT.24 (on-shell line-element quantization — the parity-read step below is scoped to materialized edges); FT.11, FT.12, FT.13; Def-H.8; rung-2 decomposition (Def-H.8 attribution clause; the counting half is derived — EM.9/D3.4 end registers). Table evaluated at $\Omega\equiv1$ (stated simplification, not a general claim).

\textbf{Pf.}

%dsp-level:0
\[
\text{qubit}\to\text{bit}:\ \ \mathrm{A1}:\ \mathcal H_q\cong\mathbb C^2;
\]

%dsp-level:0
\[
\ \ \mathrm{A2}:\ C(p,q)
\]

%dsp-level:0
\[
= 1\mathfrak q;
\]

%dsp-level:0
\[
\ \ \mathrm{CT0.6}:\ C_{\mathrm{cap}}
\]

%dsp-level:0
\[
= 4\ \text{bits}\ (16\ \text{states});
\]

%dsp-level:0
\[
\ \ \mathrm{A4}/\mathrm{CT0.1}:\ \text{dichotomy};
\]

%dsp-level:0
\[
\ \ \mathrm{CT0.7}:\ \text{reality}
\]

%dsp-level:0
\[
= \text{swap grading}
\]

%dsp-level:0
\[
\mathrm{CT0.9}:\ \Phi
\]

%dsp-level:0
\[
= \{0,\pi\},
\]

%dsp-level:0
\[
\ \sigma_e\in\{\pm1\}:\ 1\ \text{bit per real edge}\ (\mathbb C^2\to\mathbb Z_2\ \text{reduction}:\ \text{primordial node-level entropy cost});
\]

%dsp-level:0
\[
\quad\mathrm{C9}:\ 1\ \text{bit}
\]

%dsp-level:0
\[
= \ln2\ \text{nats}
\]

%dsp-level:0
\[
N_{\mathrm{visible}}=\delta h(c)\ \text{bits};
\]

%dsp-level:0
\[
\qquad \mathrm{Def\text{-}H.2}\,/\,\mathrm{CT1.19}\ (\text{geodesic}):\ \delta h(c)=\Delta\ell(c)/(\Omega\,\mathfrak h)\ \ [\Omega\equiv1\ \text{in the table below}]
\]

%dsp-level:0
\[
\mathrm{D3.5}/\mathrm{EM.16},\ \mathrm{FT.9}\,(\alpha_{\mathrm{EM}}=\pi),\ \mathrm{FT.24}\,(\widehat{ds}^{\,2}\in\mathbb Z\ \text{on-shell}):
\]

%dsp-level:0
\[
\quad G^{(n)}(e)
\]

%dsp-level:0
\[
= \sigma_e\,
\]

%dsp-level:0
\[
(-1)^{\widehat{ds}^{\,2}(e)}\in\{\pm1\}:\ \text{one transit bit per materialized edge, factored into two register reads}
\]

%dsp-level:0
\[
\text{sign read}\ \sigma_e\to\text{charge (additive boundary read, EM.9)};
\]

%dsp-level:0
\[
\qquad \text{parity read}\ (-1)^{\widehat{ds}^{\,2}}\to\text{holonomy (cumulative read; per-edge}\ \alpha=\pi\ [\mathrm{FT.9}]\ \text{with on-shell integrality}\ [\mathrm{FT.24}/\mathrm{FT.25}])
\]

%dsp-level:0
\[
\text{bookkeeping}:\ \text{the per-edge visible datum is the product}\ G^{(n)}(e)\in\{\pm1\}\ —\ \text{one bit};
\]

%dsp-level:0
\[
\quad N_{\mathrm{visible}}=\delta h(c)\ \text{counts exactly these}\ [\mathrm{CT0.9}],\ \neg\text{the two factors separately}
\]

%dsp-level:0
\[
\sigma_e\ \text{and}\ (-1)^{\widehat{ds}^{\,2}(e)}\ \text{are its two factor registers}\ [\mathrm{CT0.9}\ \text{vs}\ \mathrm{FT.24}],
\]

%dsp-level:0
\[
\ \text{separately readable only where not cancelled by the other}:\ \text{signs at the ends}\ [\mathrm{EM.9}],
\]

%dsp-level:0
\[
\ \text{parity cumulatively}\ [\mathrm{D1.7}];
\]

%dsp-level:0
\[
\quad\text{the full geometric ledger}\ \kappa_e\in\mathbb Z\ [\mathrm{FT.24}]\ \text{exceeds the binary weight}
\]

%dsp-level:0
\[
\mathrm{FT.11}:\ \text{audit};\qquad\mathrm{FT.12}:\ \text{weak}\ \neg\text{charge};\qquad\mathrm{FT.13}:\ W_i=\text{spin}\implies N_{\mathrm{spin}}=N_{\mathrm{weak}}\ (\text{positive identification})
\]

%dsp-level:0
\[
N_{\mathrm{charge}}=N_{\mathrm{visible}}-N_{\mathrm{weak}}\quad(\text{decomposition})
\]

%dsp-level:0
\[
\text{metric dependence (general, }\Omega\ \text{restored)}:\ \ \Delta\ell
\]

%dsp-level:0
\[
= \Omega\sqrt{\bar g_{ij}\Delta x^i\Delta x^j},
\]

%dsp-level:0
\[
\ \ \Omega
\]

%dsp-level:0
\[
= e^{\alpha_\Omega a_{00}},
\]

%dsp-level:0
\[
\ \ a_{00}\leftrightarrow D^+\ [\mathrm{CT1.12},\mathrm{CT1.25}]
\]

%dsp-level:0
\[
\implies D^+\ \text{enters via}\ \Omega;
\]

%dsp-level:0
\[
\qquad \delta h(c)
\]

%dsp-level:0
\[
= \Delta\ell(c)/(\Omega\mathfrak h)\ \text{throughout — the table's}\ \Omega
\]

%dsp-level:0
\[
\equiv 1\ \text{is a simplification, not the general relation}
\]

%dsp-level:0
\[
\mathfrak h\to0,
\]

%dsp-level:0
\[
\ \Delta\ell\ \text{fixed}:\ \ \text{charge}
\]

%dsp-level:0
\[
= \text{boundary effect}\ (O(1)\ \text{edge data});
\]

%dsp-level:0
\[
\ \ \text{spin}
\]

%dsp-level:0
\[
= \text{bulk}\ (\delta h-O(1)\to\infty)
\]

%dsp-level:0
\[
\implies f_{\mathrm{charge}}\to0,
\]

%dsp-level:0
\[
\ f_{\mathrm{spin}}\to1
\]

%dsp-level:0
\[
\textbf{rung-2 decomposition}\ (\text{Def-H.8 attribution};\ \text{counting}:\ \text{FT.41(a)},\ \text{forward}):\ \ N_{\mathrm{weak}}
\]

%dsp-level:0
\[
= \delta h(c)-2\ (\text{interior edges}),
\]

%dsp-level:0
\[
\ \ N_{\mathrm{charge}}
\]

%dsp-level:0
\[
= 2\ (\text{end bits};\ \mathrm{EM.9}\ \text{reads exactly the two end signs})
\]

%dsp-level:0
\[
\begin{array}{c|c|c|c|c}
\delta h(c) & \Delta\ell\ (\Omega\equiv1) & f_{\mathrm{charge}} & f_{\mathrm{spin}} & \text{content}\\ \hline
2 & 2\mathfrak h & 1 & 0 & \text{pure charge}\\
3 & 3\mathfrak h & 2/3 & 1/3 & \text{minimal spin}\\
4 & 4\mathfrak h & 1/2 & 1/2 & \text{equal}\\
\to\infty & \to\infty & \to0 & \to1 & \text{pure spin}
\end{array}\qquad(\text{general}\ \Omega:\ \text{multiply the}\ \Delta\ell\ \text{column by}\ \Omega)\ \blacksquare
\]

\textbf{Phys.} The visible bit budget of a bound chain splits into charge (its two ends — a boundary effect) and spin (its interior — a bulk effect): long chains are dominated by spin, and the continuum limit is a spin/holonomy limit. The split is a positive identification via the complete force audit, not a subtraction residue.

\vspace{2pt}\noindent\rule{\textwidth}{0.4pt}\vspace{2pt}

\end{theorem}

\begin{definition}[{\texorpdfstring{Def-H.9 — Subnetwork coarse-graining projection}{Def-H.9 - Subnetwork coarse-graining projection}}]\label{tag:Def-H-9}
\textbf{Def.} $\Sigma$ bounded subnetwork, designated terminus $v$ [Def-H.4]. $\mathcal M(\Sigma):=$ valid microscopic configurations (signs $\sigma_e$ [CT0.9], geometric data $\widehat{ds}^{\,2}(e)$ [D1.6] — integer on-shell [FT.24–FT.25], generic $O(1)$ real off-shell [CT1.19a] — subject to rung masks and Janus balance D1.3). Projection $\Pi_v:\mathcal M(\Sigma)\to\mathcal Q_v=(q,\mathbf S,m_{\mathrm{hid}})$:

%dsp-level:0
\[
(1)\ \ q:=\tfrac12(\sigma_s-\sigma_t)\ [\mathrm{EM.9}];
\]

%dsp-level:0
\[
(2)\ \ E_{\mathrm{tot}}(C):=|\Sigma|\,\ln2,
\]

%dsp-level:0
\[
\quad |\Sigma|:
\]

%dsp-level:0
\[
= \#\{\text{real edges touching}\ \Sigma,\ \text{internal}+\text{boundary}\}\ \ (\text{ledger-retained bits, locked content included}\ [\mathrm{CT4.3}]);
\]

%dsp-level:0
\[
E_{\mathrm{vis}}(\gamma):=\min\bigl(\delta h(\gamma),\ 4\,\delta\tau(\gamma)\bigr)\ln2\ [\mathrm{FT.1}];\qquad
I_{\mathrm{hid}}:=E_{\mathrm{tot}}-E_{\mathrm{vis}};\qquad
m_{\mathrm{hid}}:=I_{\mathrm{hid}}/\mathfrak h;
\]

%dsp-level:0
\[
(3)\ \ \Theta_{\mathrm{int}}(C):=\!\!\sum_{e\in\Sigma\setminus\partial\Sigma}\!\!\widehat{ds}^{\,2}(e);\qquad
\mathrm{Hol}_{\mathrm{int}}:=e^{i\pi\Theta_{\mathrm{int}}}=(-1)^{\Theta_{\mathrm{int}}}\in\{\pm1\}
\]

%dsp-level:0
\[
[\text{on-shell}\ \Theta_{\mathrm{int}}\in\mathbb Z\ \text{— derived from materialization, not postulated}:\ \mathrm{FT.25};\ \alpha=\pi:\ \mathrm{FT.9}];
\]

%dsp-level:0
\[
\mathbf S:\ \text{the}\ \mathfrak{su}(2)\ \text{Casimir label}\ j(j{+}1)\ \text{associated with}\ \mathrm{Hol}_{\mathrm{int}}\ \text{under the weak action}\ [\mathrm{FT.13}]
\]

%dsp-level:0
\[
\quad (\text{the extraction rule is application-layer content}).
\]

\textbf{Dimensional check.} $[I_{\mathrm{hid}}]=\mathfrak q^0$ (information); $[\mathfrak h]=\mathfrak q^{-1/2}$ [D2.9] $\Rightarrow[m_{\mathrm{hid}}]=\mathfrak q^{+1/2}=$ energy [CT2.1] $\checkmark$ (machine: \texttt{verify\_\allowbreak{}H.py}). \textbf{Notation.} $m_{\mathrm{hid}}$ — the informational mass read — is kept symbolically distinct from the pole mass $m$ of D2.5/EW.19 ($m^2:=$ pole displacement): the two are different constructions, and their identification or relation is nowhere assumed in this document.

\vspace{2pt}\noindent\rule{\textwidth}{0.4pt}\vspace{2pt}

\end{definition}

\begin{theorem}[{\texorpdfstring{FT.15 — Coarse-graining invariance}{FT.15 - Coarse-graining invariance}}]\label{tag:FT-15}
\textbf{Thm.} $C_1,C_2\in\mathcal M(\Sigma)$ sharing boundary signature $(\delta h,\delta\tau)$, boundary charge $q$, and $\Theta_{\mathrm{int}}$:

%dsp-level:0
\[
\Pi_v(C_1)=\Pi_v(C_2)\qquad(\text{invariance under internal rearrangement}).
\]

\textbf{Req.} Def-H.9; EM.9 (Part 2); FT.1; FT.13. The $\mathbf S$ clause additionally requires FT.25's on-shell scope for $\Theta_{\mathrm{int}}$ ($\mathrm{Hol}_{\mathrm{int}}$ well-definedness); the $q,m_{\mathrm{hid}}$ clauses do not.

\textbf{Pf.}

%dsp-level:0
\[
q=\tfrac12(\sigma_s-\sigma_t):\ \text{boundary signs only}\ [\mathrm{EM.9}]\implies q(C_1)=q(C_2)
\]

%dsp-level:0
\[
m_{\mathrm{hid}}
\]

%dsp-level:0
\[
= (|\Sigma|\ln2-E_{\mathrm{vis}})/\mathfrak h:\ \ |\Sigma|\ \text{fixed volume datum};
\]

%dsp-level:0
\[
\ E_{\mathrm{vis}}
\]

%dsp-level:0
\[
= f(\delta h,\delta\tau)\ [\mathrm{FT.1}]
\]

%dsp-level:0
\[
\implies m_{\mathrm{hid}}(C_1)
\]

%dsp-level:0
\[
= m_{\mathrm{hid}}(C_2)
\]

%dsp-level:0
\[
\mathbf S\ \text{from}\ \mathrm{Hol}_{\mathrm{int}}
\]

%dsp-level:0
\[
= (-1)^{\Theta_{\mathrm{int}}}\ [\mathrm{FT.25}]:\ \Theta_{\mathrm{int}}(C_1)
\]

%dsp-level:0
\[
= \Theta_{\mathrm{int}}(C_2)\implies\mathbf S(C_1)
\]

%dsp-level:0
\[
= \mathbf S(C_2)\ \ (\text{same weak action},\ \mathrm{FT.13})
\]

%dsp-level:0
\[
\implies\Pi_v\ \text{depends only on}\ (\delta h,\delta\tau,\sigma_s,\sigma_t,\Theta_{\mathrm{int}}):\ \neg\text{internal topology},\ \neg\text{phase distribution}.\ \blacksquare
\]

\textbf{Cor FT.15.1 (self-averaging).} $(q,\mathbf S,m_{\mathrm{hid}})$ are the unique self-averaging macroscopic observables of $\Sigma$: as temperature depends only on total kinetic energy, invariant under permutation of molecular velocities, the fermionic quantum numbers depend only on boundary channel capacity and total internal holonomy — invariant under permutation of individual internal edge phases.

\textbf{Phys.} A coarse-grained fermion forgets everything about its internal wiring except three numbers: its boundary charge, its hidden-information mass, and (on-shell, FT.25) a spin label built from total internal holonomy. Two microscopically different internal configurations that agree on boundary data and total holonomy are the same particle from outside — coarse-graining here plays exactly the role bulk thermodynamic variables play for a gas of molecules.

\vspace{2pt}\noindent\rule{\textwidth}{0.4pt}\vspace{2pt}

\end{theorem}

\begin{theorem}[{\texorpdfstring{FT.16 — The single-vertex hidden-information gap is generic; strong's signature is contraction-incompleteness}{FT.16 - The single-vertex hidden-information gap is generic; strong's signature is contraction-incompleteness}}]\label{tag:FT-16}
\textbf{Thm.} $v$ observed by a single neighbour through one incident edge ($\delta\tau\ge1$):

%dsp-level:0
\[
I_{\mathrm{hid}}(\{v\})=(n(v)-1)\ln2\qquad(\text{aspect-blind});
\]

$\neg$strong-specific (nonzero already at $n=2$, where strong $\notin\mathcal A(2)$ [CT0.13]); strong's distinguishing signature is structural: contraction-completeness [CT3.17/18/22], $\neg$ a bit count.

\textbf{Req.} CT0.6, CT0.9, CT0.13 (Part 1); CT1.29 (Part 2); CT3.17, CT3.18, CT3.22 (Part 2); D3.3, D3.9; Def-H.1, Def-H.9; FT.1, FT.10.

\textbf{Pf.}

%dsp-level:0
\[
\Sigma=\{v\},\ n(v)=n:\ \ E_{\mathrm{tot}}=n\ln2\ [\mathrm{Def\text{-}H.9};\ 1\ \text{bit per incident edge},\ \mathrm{CT0.9}]
\]

%dsp-level:0
\[
o\ \text{single neighbour}:\ \gamma=\{1\ \text{edge}\},\ \delta h=1:\ \ E_{\mathrm{vis}}=\min(1\cdot\ln2,\ \delta\tau\cdot4\ln2)=\ln2\ (\delta\tau\ge1)\ [\mathrm{FT.1}]
\]

%dsp-level:0
\[
\implies I_{\mathrm{hid}}=(n-1)\ln2;\qquad n=1,2,3,4:\ 0,\ \ln2,\ 2\ln2,\ 3\ln2
\]

%dsp-level:0
\[
n=2:\ I_{\mathrm{hid}}>0\ \wedge\ \mathrm{strong}\notin\mathcal A(2)\ [\mathrm{CT0.13}]\implies I_{\mathrm{hid}}\ \text{aspect-blind}\ (\text{Def-H.1 locality})
\]

%dsp-level:0
\[
\text{full access}:\ \delta h
\]

%dsp-level:0
\[
= n:\ \ E_{\mathrm{vis}}
\]

%dsp-level:0
\[
= \min(n,4\,\delta\tau)\ln2\overset{n\le3,\ \delta\tau\ge1}{=}n\ln2
\]

%dsp-level:0
\[
= E_{\mathrm{tot}}
\]

%dsp-level:0
\[
\implies I_{\mathrm{hid}}\to0\ \ \forall\ \text{accessible rung}\ [\mathrm{CT0.6}:\ O(n)=0]
\]

%dsp-level:0
\[
\text{asymmetry}\ [\mathrm{CT3.22},\ \mathrm{CT3.18}]:\ \ F_{\mathrm{EM}}\ \text{at}\ v=f_v\ (\text{own two signs};\ \text{complete locally},\ \mathrm{D3.3});
\]

%dsp-level:0
\[
\text{colour at}\ v:\ \text{physical only contracted}\ (\delta^i{}_j,\varepsilon_{ijk},\varepsilon^{ijk};\ \text{baryon}:3\ \text{triplets},\ \text{meson}:\ 3\oplus\bar3\ [\mathrm{D3.9}]);
\]

%dsp-level:0
\[
\ \ \text{isolated triplet}:\ \text{norm}\ 0,
\]

%dsp-level:0
\[
\ \neg\text{pole}\ [\mathrm{CT3.22}]
\]

%dsp-level:0
\[
\implies I_{\mathrm{hid}}\ \neg\text{carvable per-aspect by subtraction (aspect-blind; vanishes under full access)};
\]

%dsp-level:0
\[
\qquad \text{strong's signature}=\text{contraction-completeness}\ (\text{constructed: Def-H.10},\ \mathrm{FT.18}\text{--}\mathrm{FT.19}).\ \blacksquare
\]

\textbf{Phys.} Hidden information at a vertex is generic bookkeeping, not a force fingerprint: it vanishes under full local access for every accessible rung. What genuinely distinguishes the strong sector is that its observables are incomplete until contracted against partners possibly outside the object — the seed of confinement as an information-theoretic property.

\vspace{2pt}\noindent\rule{\textwidth}{0.4pt}\vspace{2pt}

\end{theorem}

\begin{definition}[{\texorpdfstring{Def-H.35 — Channel-visible content (the boundary-transport read)}{Def-H.35 - Channel-visible content (the boundary-transport read)}}]\label{tag:Def-H-35}
\textbf{Def.} $\Sigma$ a finite vertex set; $o\notin\Sigma$ a single external vertex whose only contact is its own incident edges into $\Sigma$, so $\gamma(o,\Sigma)=E_{\mathrm{bd}}$ [Def-H.1], over an interval $(\delta h,\delta\tau)$ [D0.6, D0.7]. The \textbf{channel-visible content} of $\Sigma$ to $o$ is the boundary registration sum

%dsp-level:0
\[
E_{\mathrm{vis}}(\gamma)\ :=\ \sum_{e\in E_{\mathrm{bd}}}1\ \text{bit}\ =\ |E_{\mathrm{bd}}|\ln2\ \ \text{nats}
\]

%dsp-level:0
\[
\qquad [\mathrm{CT0.9}:\ 1\ \text{real edge}=1\ \text{bit};\ \mathrm{D1.1}:\ \text{registrations are classical edge data};\ \mathrm{C9}],
\]

attained whenever a hop schedule within the channel capacity exists [$\delta h\le4\,\delta\tau$ jaw, CT0.6a; attainment on boundary sets: FT.153]. $I_{\mathrm{hid}}:=E_{\mathrm{tot}}-E_{\mathrm{vis}}$ [Def-H.9(2)] is evaluated with $E_{\mathrm{vis}}$ of this definition. \textbf{Locus-separation clause (constitutive).} Two distinct reads exist in the corpus and are never interchanged: $(1)$ \textbf{the edge-transport read} (this definition): content carried across a boundary edge by a hop. A hop reads only the edge's own registration [D1.9: tick/transport locality — $\operatorname{dom}\subseteq\{v\}\cup\{e\ni v\}$, pre/post-tick data only, $\neg\exists$ read at edge-distance $>1$]; the registration lives in the pair channel of the edge's endpoints [A2: $C(p,q)=1\mathfrak q$ per pair; CT0.4: real $=$ registered]. $(2)$ \textbf{the vertex-access read} [CT4.9(1)]: content accessible \textit{at} a vertex $v$ — capped, $S(v)\le3\ln2$ for accessible $v$ ($n\le3$ [CT0.11]); at locked $v$ ($n\ge4$) the excess is absorbed and retained in the ledger [CT4.1, CT4.3, D4.2: retained, $\neg$accessible; D1.11]. The cap of $(2)$ bounds reads at vertices; it does not bound transports across edges: a locked boundary vertex retains its excess under $(2)$ while each incident real boundary edge carries its own pair-channel registration under $(1)$. The two reads answer different questions — “what can $o$ receive through $\gamma$” versus “what is accessible at $v$” — and every theorem states which it uses.

\end{definition}

\begin{theorem}[{\texorpdfstring{FT.153 — Boundary-schedule attainment}{FT.153 - Boundary-schedule attainment}}]\label{tag:FT-153}
\textbf{Thm.} Single-contact setting [Def-H.35]: $\gamma(o,\Sigma)=E_{\mathrm{bd}}$, $\delta h(\gamma)=|E_{\mathrm{bd}}|$. If $|E_{\mathrm{bd}}|\le4\,\delta\tau(\gamma)$ [CT0.6a capacity jaw], then a valid transit configuration $\gamma'\subseteq\gamma$ exists whose hop schedule carries every boundary registration within $\delta\tau$ ticks at $\le4$ hops/tick; hence

%dsp-level:0
\[
E_{\mathrm{vis}}(\gamma)=|E_{\mathrm{bd}}|\ln2=\min\bigl(\delta h(\gamma),\,4\,\delta\tau(\gamma)\bigr)\ln2=\mathsf B(\gamma)
\]

%dsp-level:0
\[
\qquad [\mathrm{FT.1}'\text{s bound met with equality};\ \mathrm{Def\text{-}H.3}'\text{s sup attained}],
\]

regardless of the lock status of the boundary vertices. On channels whose edges are bound-locked in every configuration the transit domain is empty and the received content is $0$ (attainment vacuous; the bound trivially valid) [FT.36(2) clause].

\textbf{Req.} Def-H.1, Def-H.3, Def-H.35; A2; D0.10, D1.1, D1.5, D1.9; CT0.4, CT0.6a, CT0.9 (Part 1); CT4.3 (Part 2); EM.5, EM.6, EM.7 (Part 2); FT.1; the schedule construction generalized by FT.36(2) (forward pointer, Tier F.13).

\textbf{Pf.}

%dsp-level:0
\[
E_{\mathrm{bd}}=\{e_1,\dots,e_m\},\ m\le4\,\delta\tau:\ \text{partition into}\ \lceil m/4\rceil\le\delta\tau\ \text{per-tick groups of}\ \le4
\]

%dsp-level:0
\[
\text{per group: run transit (photon) content on the group's edges}\ [\mathrm{EM.5},\mathrm{EM.6}:\ \text{trivial chains}\subseteq V_2\ \text{valid transit content};\ \mathrm{D1.5}:\ \text{open propagating}=\text{transit};\ \mathrm{D0.10}:\ \mathfrak C(\mathcal R)\ \text{admits transit assignments}]
\]

%dsp-level:0
\[
\text{each hop carries its edge's registration bit}\ [\mathrm{CT0.9};\ \mathrm{D1.1}];
\]

%dsp-level:0
\[
\qquad \text{hops read edge data only}\ [\mathrm{D1.9}]:\ \neg\text{vertex read invoked}\implies\text{lock status of endpoints irrelevant to the transport}\ [\mathrm{Def\text{-}H.35}\ \text{locus separation};\ \mathrm{CT4.3}:\ \text{retained excess stays ledger-internal}]
\]

%dsp-level:0
\[
\mathsf E_{\mathrm{tr}}(\gamma')=m\ln2\ [\mathrm{EM.7}\ \text{additive}];
\]

%dsp-level:0
\[
\qquad \mathrm{FT.1}:\ \sup\le\min(m,4\,\delta\tau)\ln2=m\ln2;
\]

%dsp-level:0
\[
\qquad \text{construction attains}\ \implies E_{\mathrm{vis}}=m\ln2.\ \blacksquare
\]

\textbf{Phys.} The channel bound is not merely valid on boundary sets but reached: the observer receives the full boundary registration sum, scheduled at channel speed — four bits per tick — and vertex locks eat none of it, because a hop reads its edge and never its vertex. This is the attainment half of the internal-edge lemma's hypotheses; FT.36(2) generalizes the same construction to arbitrary channels.

\end{theorem}

\begin{theorem}[{\texorpdfstring{FT.17 — Internal-edge lemma}{FT.17 - Internal-edge lemma}}]\label{tag:FT-17}
\textbf{Thm.} $\Sigma$ finite vertex set; $o\notin\Sigma$ a single external vertex whose only contact is its own incident edges into $\Sigma$ (so $\gamma(o,\Sigma)=$ the $E_{\mathrm{bd}}$ boundary edges; no internal edges). Hypotheses: $E_{\mathrm{bd}}\le4\,\delta\tau(\gamma)$ [channel-capacity jaw, CT0.6a: the channel does not clip $o$'s view]; channel-visible content per Def-H.35's edge-transport read, attained by FT.153. Then

%dsp-level:0
\[
I_{\mathrm{hid}}(\Sigma)=E_{\mathrm{int}}\,\ln2\qquad(\text{exactly; aspect-blind}).
\]

\textbf{Req.} CT0.6, CT0.9 (Part 1); CT1.29 (Part 2); CT4.2, CT4.3, CT4.9 (Part 2); D4.2; Def-H.1, Def-H.9; FT.1.

\textbf{Pf.}

%dsp-level:0
\[
\mathrm{Def\text{-}H.9}\ (|\Sigma|=\text{internal}+\text{boundary}:\ \text{internal-only reading gives}\ I_{\mathrm{hid}}<0\ \text{on proper sub-selections}):
\]

%dsp-level:0
\[
\quad E_{\mathrm{tot}}=(E_{\mathrm{int}}+E_{\mathrm{bd}})\ln2\ [\mathrm{CT0.9}/\mathrm{CT1.29}\ \text{counting};\ \text{locked content retained},\ \mathrm{CT4.3}]
\]

%dsp-level:0
\[
\delta h(\gamma)=E_{\mathrm{bd}};\qquad E_{\mathrm{vis}}=\min\bigl(E_{\mathrm{bd}},\ 4\,\delta\tau\bigr)\ln2\overset{\text{hyp}}{=}E_{\mathrm{bd}}\ln2\ (\delta\tau\ge1)
\]

%dsp-level:0
\[
\text{hyp 2 (Def-H.35)}:\ E_{\mathrm{vis}}=\text{actual accessible content}\ (\neg\text{merely an upper bound})
\]

%dsp-level:0
\[
I_{\mathrm{hid}}=E_{\mathrm{tot}}-E_{\mathrm{vis}}=(E_{\mathrm{int}}+E_{\mathrm{bd}})\ln2-E_{\mathrm{bd}}\ln2=E_{\mathrm{int}}\ln2.\ \blacksquare
\]

\textbf{Cross-checks} (independently computed; machine: \texttt{verify\_\allowbreak{}H.py}): 1. Oriented 3-cycle, external observer: $E_{\mathrm{int}}=3$, $E_{\mathrm{bd}}=1\le4\cdot1\ \checkmark\Rightarrow I_{\mathrm{hid}}=3$ bits. 2. Bidirectional 3-vertex face $\{A,B,C\}$ of the bidirectional $K_4$ ("tetrahedron", $n(v)=6$), observed by $D$: $E_{\mathrm{int}}=6$, $E_{\mathrm{bd}}=6$: hypothesis requires $\delta\tau\ge2$ ($6\le4\cdot2=8$). Vertices oversubscribed ($n=6\ge4$): their excess is locked but \textit{retained} [CT4.3, D4.2] and counted in $E_{\mathrm{tot}}$; Def-H.35's locus separation concerns channel-visible content only. $I_{\mathrm{hid}}=6$ bits. 3. Asymmetric single-point contact (whole $K_4$, $\Sigma=\{A,B,C,D\}$, external $o$ joined only to $A$ by one bidirectional pair): $E_{\mathrm{int}}=\binom42\times2=12$, $E_{\mathrm{bd}}=2\le4\cdot1\ \checkmark$; $A$'s own accessible content 2 bits $<3$ [CT4.9] $\checkmark\Rightarrow I_{\mathrm{hid}}=12$ bits, $m_{\mathrm{hid}}=12\ln2/\mathfrak h$ [Def-H.9]. \textit{(note [FT.62(4)]: at the corrected directed rung $n(A)=8$ the contact vertex is locked; the per-vertex-cap check is carried by Def-H.35's locus separation and CT4.3 retention — $I_{\mathrm{hid}}=12$ bits unchanged.)}

\textbf{Rmk (where the $I_{\mathrm{hid}}$ comes from).} Not $D^-$ (binding constraint is the channel bound, reached before per-vertex caps bind) and not colour (the algebra never references $\mathcal A(n)$): $I_{\mathrm{hid}}$ — and hence $m_{\mathrm{hid}}=I_{\mathrm{hid}}/\mathfrak h$ — is exactly $\Sigma$'s own internal connectivity, invisible to any single exterior observer by Def-H.1 alone. A purely topological fact, prior to any question of which forces act where. Cross-check 3 shows independence of contact distribution.

\textbf{Phys.} Mass, in this reading, is nothing but the information a single external vantage point cannot recover about a subnetwork's inside — a topological fact about connectivity, not a substance. The three cross-checks show the same total hidden content surviving no matter how the internal edges are arranged or where the single external contact point sits, as long as the channel-capacity hypothesis holds: hiddenness here is a property of the boundary/interior split, not of any particular internal configuration.

\vspace{2pt}\noindent\rule{\textwidth}{0.4pt}\vspace{2pt}

\vspace{2pt}\noindent\rule{\textwidth}{0.4pt}\vspace{2pt}

\end{theorem}

\phantomsection
\section[Tier F.5b — The Mass-Read Registry and the Tilt Composition]{\texorpdfstring{Tier F.5b — The Mass-Read Registry and the Tilt Composition}{Tier F.5b - The Mass-Read Registry and the Tilt Composition}}

\begin{definition}[{\texorpdfstring{Def-H.39 — The mass-read registry; bridge slots}{Def-H.39 - The mass-read registry; bridge slots}}]\label{tag:Def-H-39}
\textbf{Def.} The corpus's mass reads form a registry of species; a \textbf{bridge} between two species is a map carrying a dimensional factor (a power of $\mathfrak h$) and a numeric coefficient. The registry records the loci; no bridge is constructed, and none may be silently assumed. \textbf{Nomenclature rule.} Every occurrence of \textit{mass} in this corpus names one read of this registry at the species declared at its locus; no unqualified \textit{mass} crosses species; a computed variant value is a \textit{price} — an exact named functional of the declaring data — and $\neg$a mass read.

%dsp-level:0
\[
(1)\ \text{informational}:\quad m_{\mathrm{hid}}=I_{\mathrm{hid}}/\mathfrak h\qquad[\mathrm{Def\text{-}H.9}(2);\ \mathrm{FT.16},\mathrm{FT.17};\ \ [m_{\mathrm{hid}}]=\mathfrak q^{+1/2}];
\]

%dsp-level:0
\[
(2)\ \text{pole/defect}:\quad m^2_{\mathrm{pole}}
\]

%dsp-level:0
\[
\qquad [\mathrm{Def\text{-}H.37}\ \text{— forward pointer, Tier F.17};\ \mathrm{D2.5}\ \text{vocabulary};\ \ [m^2_{\mathrm{pole}}]=\mathfrak q];
\]

%dsp-level:0
\[
(3)\ \text{inertial/tilt — the Lyapunov-contraction read}:\quad m_i^2(q)=\Delta S(q)
\]

%dsp-level:0
\[
\qquad [\mathrm{L3}(4),\mathrm{L4};\ \text{Part 1: the deficit accumulated under the L1 contraction law};\ \text{dimensionless nats}];
\]

%dsp-level:0
\[
(4)\ \text{sector-expectation}:\quad m^2_B(q,S)
\]

%dsp-level:0
\[
\qquad [\mathrm{Def\text{-}H.32}(6)\ \text{— forward pointer, Tier F.36};\ \text{dimensionless nats}];
\]

%dsp-level:0
\[
(5)\ \text{1-loop}:\quad m_i^2=V_{\mathrm{eff}}(\varepsilon_q^{(i)})/8
\]

%dsp-level:0
\[
\qquad [\mathrm{Def\text{-}H.31}(6)\ \text{— forward pointer, Tier F.35};\ \mathrm{L3}(2):\ V=8\,\Delta S;\ \text{dimensionless nats}].
\]

\textbf{Bridge slots.} $\beta_{31}$ (inertial$\leftrightarrow$informational): unconstructed — a construction must supply the power of $\mathfrak h$ and the per-edge numeric coefficient reconciling L4's deficits (exact rationals, state-dependent: the branch-state values $(\tfrac34,\tfrac1{12})$ [FT.107(2) — forward pointer]) with FT.17's $\ln2$-per-hidden-bit count; the read-identification locus (“the same object read differently”, at pair level versus subnetwork level) is this slot. $\beta_{21}$, $\beta_{23}$ (pole$\leftrightarrow$informational / inertial): the relating-equation locus. \textbf{Kind discipline:} a dimensionless-nats read and a dimensionful read are comparable only through a constructed $\beta$ [FT.89(2) — forward pointer]; any quantitative use of a bridge requires its construction first. \textbf{Ceilings (derived; cited at use).} Species (2): $m^2_{\mathrm{pole}}\in[0,4\hat\mu'^2]$ [FT.76(3) — forward pointer]; the cover-ladder read $K=2|\mathcal C|\le3V/2$ [FT.79(2) — forward pointer]; the form-factor species $M\le c_wV$ [FT.97(5) — forward pointer]; species (5): bare range linear in cell size, composed-tilt range exponential in $E_{\mathrm{int}}$ under Def-H.40 [FT.104(2) — forward pointer]; species (4): $\le\tfrac32B(B-1)$, quadratic in budget [FT.107(4) — forward pointer].

\end{definition}

\begin{definition}[{\texorpdfstring{Def-H.40 — Chain-level tilt composition}{Def-H.40 - Chain-level tilt composition}}]\label{tag:Def-H-40}
\textbf{Def.} Let $c$ be a materialized bound chain / floor with $E_{\mathrm{int}}$ internal edges [Def-H.9]. Ledger facts (derived): every internal edge entered $\mathcal R$ through exactly one $J^+$ seal [D1.2: crystallization $n_{\mathcal R}:0\to1$; D1.4: chain edges $\in E_{\mathcal R}$]; every A4 crystallization event applies exactly one Gibbs tilt step to the pair state of that event [L0(2), its declared theorem-local dynamical scope]; the ledger counts tilt steps [L0(4)]. \textbf{Composition clause (constitutive scope).} The tilt steps of $c$'s sealing history compose on one chain-level aggregate state object — the per-event iteration of L0(4)'s $k$-event law — giving the chain's Lyapunov composition

%dsp-level:0
\[
V_{E_{\mathrm{int}}}\ :=\ \bigl(1-\tfrac\gamma4\bigr)^{2E_{\mathrm{int}}}V_0\qquad[\mathrm{L2}\ \text{per-step law iterated};\ \gamma=1-4\varepsilon_q\in(0,1],\ \mathrm{D0.3};\ \mathrm{CT0.8}],
\]

sharp point $\gamma=1$ (D0.4): $\ V_{E_{\mathrm{int}}}=(\tfrac9{16})^{E_{\mathrm{int}}}V_0$ [exact rationals, L2]. \textbf{Read.} The fermion-mass read of the contraction rate is species (3) of Def-H.39 [L3(4): $m_i^2=\Delta S$]; its composed range is the exponential species of Def-H.39's ceiling row [FT.104(2)]. \textbf{Scope clauses.} (i) The aggregate state object is this definition's sole non-derived locus: the inventory's states are pair states [A1]; L0(4) proves the $k$-event iteration on one pair channel; the chain-level aggregate extends that iteration to the chain's sealing history. (ii) No calibration claim: no datum fixes $(\tfrac9{16})^{E_{\mathrm{int}}}$ to any observed hierarchy [P1, P5]. (iii) No particle-family identification [FT.77, FT.79(4)–(6); Census FT.84(2)]. (iv) The read is dimensionless nats; comparison with dimensionful species requires Def-H.39's $\beta$ slots.

\end{definition}

\begin{theorem}[{\texorpdfstring{FT.154 — The tilt-composition law}{FT.154 - The tilt-composition law}}]\label{tag:FT-154}
\textbf{Thm.} Under Def-H.40, for a materialized bound chain with $E_{\mathrm{int}}$ internal edges and tilt strength $\gamma\in(0,1]$:

%dsp-level:0
\[
(1)\quad\textbf{event count:}\ \ \text{the chain's tilt-step count is exactly}\ E_{\mathrm{int}};
\]

%dsp-level:0
\[
(2)\quad\textbf{per-step contraction:}\ \ V'=\bigl(1-\tfrac\gamma4\bigr)^2V+O(\delta\rho^3)\qquad[\mathrm{L2}];
\]

%dsp-level:0
\[
(3)\quad\textbf{composition:}\ \ V_{E_{\mathrm{int}}}=\bigl(1-\tfrac\gamma4\bigr)^{2E_{\mathrm{int}}}V_0\ \text{with per-step remainders accumulated};
\]

%dsp-level:0
\[
\qquad \gamma=1:\ \ V_{E_{\mathrm{int}}}=\bigl(\tfrac9{16}\bigr)^{E_{\mathrm{int}}}V_0;
\]

%dsp-level:0
\[
\text{eigenvalue composition:}\ \ \mu_k=\bigl(1-\tfrac\gamma4\bigr)^k\mu_0+O(\mu_0^2)\ \ \text{after}\ k\ \text{recorded events}\qquad[\mathrm{L0}(4)];
\]

%dsp-level:0
\[
(4)\quad\textbf{range consequences:}\ \ \text{the composed read's spans and thresholds are the exact arithmetic of}\ \bigl(\tfrac{16}9\bigr)^{E_{\mathrm{int}}}\ \text{tabulated at}\ \mathrm{FT.104}(2)\ (\text{forward pointer, Tier F.35}).
\]

\textbf{Req.} Def-H.40; Def-H.9 ($E_{\mathrm{int}}$); D0.3, D0.4, D1.2, D1.4; CT0.8; L0(1), L0(2), L0(4); L2; L3(4) (Part 1); exact rational arithmetic.

\textbf{Pf.}

%dsp-level:0
\[
(1):\ e\ \text{internal}
\]

%dsp-level:0
\[
\implies e
\]

%dsp-level:0
\[
\in E_{\mathcal R}\ [\mathrm{D1.4}]
\]

%dsp-level:0
\[
\implies e\ \text{crossed the A4 criterion once, at its}\ J^+\ \text{seal}\ [\mathrm{D1.2}:\ n_{\mathcal R}:0\to1;\ \mathrm{CT0.3}:\ \text{seal increments}]
\]

%dsp-level:0
\[
\mathrm{L0}(2):\ \text{one tilt step per event};
\]

%dsp-level:0
\[
\qquad \mathrm{L0}(4):\ |\mathfrak L(\tau_n)|
\]

%dsp-level:0
\[
= \textstyle\sum_{k<n}n_{\mathrm{real\_out}}(\tau_k)\ \text{is the tilt-step count}\ [\text{exact integer recursion},\ \mathrm{A6},\mathrm{A7}];
\]

%dsp-level:0
\[
\qquad \text{Def-H.40's aggregate scope composes the chain's sealing history}\implies\#\text{steps}=E_{\mathrm{int}}
\]

%dsp-level:0
\[
(2):\ \mathrm{L2}\ \text{verbatim}:\ V'-V=-\tfrac{\gamma(8-\gamma)}{16}V+O(\delta\rho^3)=-[1-(1-\tfrac\gamma4)^2]V+O(\delta\rho^3)
\]

%dsp-level:0
\[
(3):\ \text{iteration of}\ (2)\ E_{\mathrm{int}}\ \text{times under Def-H.40's composition clause}\ [\mathrm{L0}(4):\ \tilde{\mathcal T}_\gamma^{\,k}\ \text{composition};\ \text{interleaved Haar draws spectrally inert},\ \mathrm{L0}(1)];
\]

%dsp-level:0
\[
\qquad \gamma=1:\ (1-\tfrac14)^2=\tfrac9{16}\ [\mathrm{L2};\ \text{exact rationals}]
\]

%dsp-level:0
\[
\mu_k:\ \mathrm{L0}(4)\ \text{verbatim (induction on L2's}\ \mu'=(1-\tfrac\gamma4)\mu+O(\mu^2))
\]

%dsp-level:0
\[
(4):\ \text{arithmetic table of}\ \mathrm{FT.104}(2)\ \text{re-cited};
\]

%dsp-level:0
\[
\ \text{its thresholds}\ E_{\mathrm{int}}
\]

%dsp-level:0
\[
\ge 16/27/35/42\ \text{and the proportionality obstruction are computed there under this law}.\ \blacksquare
\]

\textbf{Phys.} The ledger gives a chain exactly one contraction event per internal edge, and each event contracts by an exact rational factor — so the geometric law $(\tfrac9{16})^{E_{\mathrm{int}}}$ at sharp crystallization is seal arithmetic, not a fitted model: the count comes from the ledger recursion, the factor from the tilt's exact leading law, and the definition supplies only the aggregate on which the count composes. The mass interpretation of the rate is the registry's inertial species; the exponential range it buys is priced, with its obstruction, at FT.104(2).

\end{theorem}

\phantomsection
\section[Tier F.6 — Subnetwork-Level Force Structure: Confinement…]{\texorpdfstring{Tier F.6 — Subnetwork-Level Force Structure: Confinement and the Active Aspect Set}{Tier F.6 - Subnetwork-Level Force Structure: Confinement and the Active Aspect Set}}

\begin{theorem}[{\texorpdfstring{FT.18 — Contraction-completeness is not $\Sigma$-local}{FT.18 - Contraction-completeness is not -local}}]\label{tag:FT-18}
\textbf{Thm.} $\Sigma\supseteq\{v_1,v_2,v_3\}$ triplet-carrying [D3.8], no antitriplet in $\Sigma$. $\exists\,C,C'\in\mathcal M(U)$ [D2.11] with $C|_\Sigma=C'|_\Sigma$ (identical restriction) such that $\mathcal C(\Sigma;C)=1$ and $\mathcal C(\Sigma;C')=0$ [Def-H.10]. Contraction-completeness is $\neg$ a function of $\mathcal M(\Sigma)$ alone.

\textbf{Req.} A1; D2.11, D3.8, D3.9; CT3.17, CT3.18 (Part 2); \textbf{hypothesis (stated, not derived from A1 alone):} $U\setminus\Sigma$ contains further triplet carriers $w_1,w_2$ available to seal with $v_1$ — A1's countable infinitude of $\mathcal Q$ secures infinitely many vertices but not that any are unsealed triplet carriers positioned to form $C'$; this is an additional genericity assumption.

\textbf{Pf.}

%dsp-level:0
\[
U\setminus\Sigma\ \text{contains further triplet carriers}\ w_1,w_2\ [\text{hypothesis, above}]
\]

%dsp-level:0
\[
C:\ \{v_1,v_2,v_3\}\ \text{sealed as one baryon}\ \varepsilon_{ijk}q_1^iq_2^jq_3^k\ \text{inside}\ \Sigma\ [\mathrm{D3.9}]:\ \text{valid}\ [\mathrm{CT3.18}];
\]

%dsp-level:0
\[
\ \ \partial_{\mathrm{col}}(\Sigma;C)\qquad = \varnothing
\]

%dsp-level:0
\[
C':\ v_1\ \text{sealed with}\ w_1,
\]

%dsp-level:0
\[
w_2\notin\Sigma;
\]

%dsp-level:0
\[
\ \ v_2,
\]

%dsp-level:0
\[
v_3\ \text{sealed elsewhere in}\ U\setminus\Sigma:\ \text{equally valid};
\]

%dsp-level:0
\[
\ \ \pi_{C'}(v_i)\not\subseteq\Sigma\ \forall i\implies\partial_{\mathrm{col}}(\Sigma;C')\qquad = \{v_1,v_2,v_3\}
\]

%dsp-level:0
\[
C|_\Sigma
\]

%dsp-level:0
\[
= C'|_\Sigma\ (\text{same vertices, edges, attachments, local entries})\implies\mathcal C\ \text{undetermined by}\ \Sigma\text{-local data}.\ \blacksquare
\]

\textbf{Rmk.} The colour-sector analogue of D2.11's dangling-boundary condition being a \textit{global} fact: whether an end is sealed depends on whether its partner exists somewhere — no $\Sigma$-local inspection certifies it.

\textbf{Phys.} Whether a colour-carrying vertex reads as "confined within $\Sigma$" or "open toward the rest of $U$" is not a fact about $\Sigma$'s own microstate at all — it is a fact about the ambient universe's contraction pattern, which $\Sigma$-local data cannot see. Confinement bookkeeping is therefore irreducibly non-local in exactly the sense D2.11 already makes boundary-sealing non-local.

\vspace{2pt}\noindent\rule{\textwidth}{0.4pt}\vspace{2pt}

\end{theorem}

\begin{definition}[{\texorpdfstring{Def-H.10 — Colour boundary and the contraction-completeness projector}{Def-H.10 - Colour boundary and the contraction-completeness projector}}]\label{tag:Def-H-10}
\textbf{Def.} $C\in\mathcal M(U)$ valid. For $v\in\mathrm{Colour}(U):=\{u\in\mathcal Q:\ u\ \text{carries a colour attachment}\ [\mathrm{D3.8}]\}$, let $\pi_C(v)\subseteq\mathrm{Colour}(U)$ be the other members of $v$'s elementary invariant contraction group under $C$ (one for a meson pair via $\delta^i{}_j$; two for a (anti)baryon triple via $\varepsilon_{ijk}/\varepsilon^{ijk}$). $\pi_C$ total on $\mathrm{Colour}(U)$ for valid $C$ [CT3.18 + D2.11 applied to the colour ledger]. For $\Sigma$: $\mathrm{Colour}(\Sigma):=\mathrm{Colour}(U)\cap\Sigma$;

%dsp-level:0
\[
\partial_{\mathrm{col}}(\Sigma;C):=\{v\in\mathrm{Colour}(\Sigma):\ \pi_C(v)\not\subseteq\Sigma\};\qquad
\mathcal C(\Sigma;C):=\begin{cases}1&\partial_{\mathrm{col}}(\Sigma;C)=\varnothing\\ 0&\text{otherwise}\end{cases};
\]

%dsp-level:0
\[
\Pi_{\mathcal C}(\Sigma;C):=\mathcal C(\Sigma;C)\cdot\mathbb I_\Sigma^{(\mathrm{strong})}
\]

%dsp-level:0
\[
\qquad (\text{Hermitian idempotent};\ \text{matches}\ \mathcal P_M^{(n)}\ [\mathrm{CT0.14}]\ \text{and}\ \Pi_v\ [\mathrm{Def\text{-}H.9}]).
\]

\textbf{Rmk (why $\mathcal C$ is defined on $\pi_C$).} A predicate counting only $\mathcal M(\Sigma)$-internal data is insufficient (FT.18): $\mathcal C$ is defined on the contraction partners $\pi_C$ — manifestly a function of the ambient configuration $C$, not of $\mathcal M(\Sigma)$ alone.

\vspace{2pt}\noindent\rule{\textwidth}{0.4pt}\vspace{2pt}

\end{definition}

\begin{theorem}[{\texorpdfstring{FT.19 — Confinement is an all-or-nothing subnetwork property}{FT.19 - Confinement is an all-or-nothing subnetwork property}}]\label{tag:FT-19}
\textbf{Thm.}

%dsp-level:0
\[
F_{\mathrm{strong}}^\Sigma(C):=\Pi_{\mathcal C}(\Sigma;C)\,F_{\mathrm{strong}}\big|_\Sigma\,\Pi_{\mathcal C}(\Sigma;C);\qquad
\mathcal C(\Sigma;C)=0\implies F_{\mathrm{strong}}^\Sigma(C)=0
\]

— regardless of how many colour-carrying vertices $\Sigma$ contains, and even if $|\partial_{\mathrm{col}}(\Sigma;C)|=1$.

\textbf{Req.} Def-H.9, Def-H.10; CT2.5, CT2.8 (Part 2); CT3.2, CT3.18, CT3.22 (Part 2); D3.1.

\textbf{Pf.}

%dsp-level:0
\[
\mathcal C
\]

%dsp-level:0
\[
= 0:\ \exists v\in\mathrm{Colour}(\Sigma):\ \pi_C(v)\not\subseteq\Sigma\ [\mathrm{Def\text{-}H.10}]
\]

%dsp-level:0
\[
\implies v\ \text{has}\ \neg\text{partner inside}\ \Sigma:\ \text{uncontracted index relative to}\ \Sigma
\]

%dsp-level:0
\[
\mathrm{CT3.18}:\ \text{uncontracted index}
\]

%dsp-level:0
\[
= \text{dangling colour-ledger boundary}\implies\text{mask-excluded}\ (\text{same argument with}\ \Sigma\ \text{as ambient})
\]

%dsp-level:0
\[
\mathrm{CT3.22}:\ \text{excluded}\implies\neg\text{pole}\ [\mathrm{CT2.5}]\ \wedge\ \text{physical norm}\ 0\ [\mathrm{CT2.8}]:\ v\text{'s}\ \Sigma\text{-restricted colour state}
\]

%dsp-level:0
\[
= \text{null vector}
\]

%dsp-level:0
\[
\Sigma\text{'s colour sector}
\]

%dsp-level:0
\[
= \bigotimes_{\mathrm{Colour}(\Sigma)}\ (\text{slot spaces})\ [\mathrm{D3.1},\mathrm{CT3.2}]:\ \text{one null factor}\implies\text{whole product null}
\]

%dsp-level:0
\[
\implies F_{\mathrm{strong}}\big|_\Sigma\text{'s matrix elements vanish}
\]

%dsp-level:0
\[
\implies\Pi_{\mathcal C}=0\implies F_{\mathrm{strong}}^\Sigma=0\cdot F_{\mathrm{strong}}\big|_\Sigma\cdot0=0.\ \blacksquare
\]

\textbf{Rmk.} Sharpens FT.16 step 4 to subnetworks: unlike EM's $f_v$ (additive, per-vertex local [D3.3, CT3.22]), strong's observability on $\Sigma$ is $\neg$gradable — one unsealed attachment anywhere zeroes $\Sigma$'s \textit{entire} colour sector. No conflict with FT.15: the tuple $(q,\mathbf S,m_{\mathrm{hid}})$ never carried a strong component [FT.11]; $\mathcal C$ tracks admissibility, not a value.

\textbf{Rmk (definition interface).} Def-H.9's $\mathcal M(\Sigma)$ is subject to rung masks $+$ Janus balance only, $\neg$ colour-sealing explicitly: an $\mathcal M(\Sigma)$-microstate with an incomplete attachment is not excluded by Def-H.9 yet has zero physical norm when $\Sigma$ is examined alone. Whether to weight $\mathcal M(\Sigma)$ by $\mathcal C(\Sigma;C)$ or keep $\Pi_{\mathcal C}$ as a downstream filter is \textbf{decided (FT.39(3))}: the downstream filter is the corpus-consistent placement — pre-weighting would delete physical EM/$G$/$H$ reads of valid ambient restrictions that seal outside $\Sigma$ [FT.18]; Def-H.9 stands unweighted.

\textbf{Phys.} Strong interactions are the one place in this decomposition where a subnetwork cannot certify its own physical activity from what it contains: one dangling colour attachment anywhere zeroes the whole sector, regardless of how many other colour-carrying vertices sit safely paired inside $\Sigma$. This is stronger than the usual notion of a boundary effect — it is not that strong force is \textit{weaker} at an unsealed boundary, it is that the entire subnetwork's colour amplitude vanishes identically.

\vspace{2pt}\noindent\rule{\textwidth}{0.4pt}\vspace{2pt}

\end{theorem}

\begin{definition}[{\texorpdfstring{Def-H.11 — Active aspect set of a subnetwork}{Def-H.11 - Active aspect set of a subnetwork}}]\label{tag:Def-H-11}
\textbf{Def.}

%dsp-level:0
\[
\mathcal A(\Sigma;C):
\]

%dsp-level:0
\[
= \bigl\{A\in\mathbf F:\ A\ \text{contributes nonzero physically observable content to}\ \Sigma\text{'s coarse-grained description under}\ C\bigr\},
\]

read against two channels: nonzero value in $\Sigma$'s coarse-grained tuple $(q,\mathbf S,m_{\mathrm{hid}})$ [Def-H.9] (\textit{charge-type}), or nonzero surviving support under $\Pi_{\mathcal C}(\Sigma;C)$ [FT.19] (\textit{selection-type}: strong). Subnetwork analogue of D1.12's $\mathcal A(n)$: $\mathcal A(n)$ asks which aspects act nontrivially on one vertex's slot algebra; $\mathcal A(\Sigma)$ asks which survive coarse-graining and masking as observables on $\Sigma$ as a whole.

\vspace{2pt}\noindent\rule{\textwidth}{0.4pt}\vspace{2pt}

\end{definition}

\begin{definition}[{\texorpdfstring{Def-H.36 — The scalar-row fold (the H read at $\Sigma$-level)}{Def-H.36 - The scalar-row fold (the H read at -level)}}]\label{tag:Def-H-36}
\textbf{Def.} The coarse-grained tuple $(q,\mathbf S,m_{\mathrm{hid}})$ [Def-H.9] carries no independent H slot. The $\Sigma$-level observable content of the scalar-trace-mode aspect — $F_H=\operatorname{span}\{\Sigma\}$, $\dim F_H=1$ [A8, CT0.12] — is read through the mass component:

%dsp-level:0
\[
\mathrm H\in\mathcal A(\Sigma;C)\iff\text{the G-row criterion holds}\ \bigl(I_{\mathrm{hid}}(\Sigma;o)\ne0\bigr)\qquad[\mathrm{FT.20}\ \text{table}].
\]

\textbf{Derived grounds (cited; the fold is the read convention this definition constitutes).} $H=h\,\Sigma$ [EW.13]; the scalar background enters the pole masses only via $M^2\propto h_0^2$ [EW.20, EW.33] — H's tuple-level contribution is via the mass read; $h_0$ remains a P4-class datum [P4]. \textbf{Distinctness clause.} The fold is a read convention for the coarse-grained tuple, $\neg$ a block identification: $\mathcal K^{(H)}$ remains a distinct block inside $\mathcal K_\Sigma$ [Def-H.15 — forward pointer, Tier F.9; Census FT.28(4)].

\end{definition}

\begin{theorem}[{\texorpdfstring{FT.20 — Reduction of $\mathcal A(\Sigma)$ to proven per-aspect criteria}{FT.20 - Reduction of to proven per-aspect criteria}}]\label{tag:FT-20}
\textbf{Thm.}

%dsp-level:0
\[
\begin{array}{c|c|c|c}
\text{aspect} & \in\mathcal A(\Sigma;C)\iff & \text{locality} & \text{basis}\\ \hline
\mathrm{EM} & \exists v\in\Sigma:\ f_v\ne0 & \Sigma\text{-local, additive} & \mathrm{D3.3},\ \mathrm{EM.9},\ \mathrm{CT3.22}\ \text{step 4}\\
G\ (\text{mass}) & I_{\mathrm{hid}}(\Sigma;o)\ne0 & \text{observer-relative}\ (\text{inherits Def-H.9}) & \mathrm{Def\text{-}H.9},\ \mathrm{FT.16},\ \mathrm{FT.17}\\
H & \text{folded into}\ G\ [\mathrm{Def\text{-}H.36}] & \text{as}\ G & \mathrm{EW.13},\ \mathrm{EW.20},\ \mathrm{EW.33}\\
W\ (\text{spin}) & \Theta_{\mathrm{int}}(\Sigma)\ne0\implies\Sigma\setminus\partial\Sigma\ne\varnothing\ \ (\text{one-directional}) & \Sigma\text{-local, intrinsic} & \mathrm{Def\text{-}H.9}(3),\ \mathrm{FT.13},\ \mathrm{FT.25}\\
D^+ & \text{never (independently)} & \text{uniform} & \mathrm{D4.3}\\
D^- & \text{never (independently)} & — & \mathrm{CT4.4}\\
J^\pm & \text{never} & \text{seals},\ \neg\text{charges} & \mathrm{CT0.7}\\
\mathrm{strong} & \mathrm{Colour}(\Sigma)\ne\varnothing\ \wedge\ \mathcal C(\Sigma;C)=1 & \text{genuinely global} & \mathrm{Def\text{-}H.10},\ \mathrm{FT.19}
\end{array}
\]

\textbf{Req.} FT.11, FT.13, FT.16–FT.19; Def-H.9–H.11; FT.25; all cited tags.

\textbf{Pf.} All rows but strong $=$ FT.11's audit extended verbatim to general $\Sigma$ (their bases reference per-vertex/uniform facts only, never size or topology); $G=$ Def-H.9's mass definition (already general); $H=$ fold into the mass read [Def-H.36; grounds: EW.13 $H=h\Sigma$; EW.20/EW.33: $M^2\propto h_0^2$ — H's tuple-level contribution is via mass]; $W=$ Def-H.9(3): $\Theta_{\mathrm{int}}\ne0\implies\Sigma$ has $\ge1$ internal edge carrying nonzero $\widehat{ds}^{\,2}$ (converse fails: a sum of nonzero terms can cancel, so $\ge1$ internal edge does not imply $\Theta_{\mathrm{int}}\ne0$); strong $=$ FT.19. $\blacksquare$

\textbf{Rmk (the point of the exercise).} Exactly one aspect — strong — requires genuinely non-$\Sigma$-local data for activity: confinement [CT3.22] is the only aspect-defining mechanism in the corpus requiring contraction against data possibly outside the object described [D3.9]. $J$ and $D^-$ never appear in $\mathcal A(\Sigma)$ despite local presence in $\mathcal A(n)$: they are structural/bookkeeping aspects (seals, locks), not observables — a distinction $\mathcal A(n)$'s representation-theoretic definition never needed, but observability-based $\mathcal A(\Sigma)$ must draw.

\textbf{Phys.} Collapsing $\mathcal A(\Sigma)$ to per-aspect activity criteria makes explicit what kind of fact each force \textit{is}: EM and mass/spin are things a subnetwork simply has, readable off its own boundary and interior; strong confinement is instead a fact about how $\Sigma$ sits inside the rest of the universe. This table is the subnetwork-level generalization of $\mathcal A(n)$ that the coarse-graining program needs next — the single-vertex active-aspect question, answered for arbitrary bounded regions.

\vspace{2pt}\noindent\rule{\textwidth}{0.4pt}\vspace{2pt}

\end{theorem}

\phantomsection
\section[Tier F.7 — The Coarse-Grained Slot Algebra: $\mathcal A(\Sigma)$ Properly Derived]{\texorpdfstring{Tier F.7 — The Coarse-Grained Slot Algebra: $\mathcal A(\Sigma)$ Properly Derived}{Tier F.7 - The Coarse-Grained Slot Algebra: Properly Derived}}

\begin{definition}[{\texorpdfstring{Def-H.12 — Coarse-grained slot algebra}{Def-H.12 - Coarse-grained slot algebra}}]\label{tag:Def-H-12}
\textbf{Def.} $\Sigma$ a bounded subnetwork, $\partial\Sigma:=\{e\in\mathcal R:\ e$ has exactly one endpoint in $\Sigma\}$ [D2.11/CT1.29 separating-surface reading, as already used in Def-H.9(2)/Def-H.10]. Let $N(\Sigma):=|\partial\Sigma|$. The \textbf{coarse-grained slot algebra} of $\Sigma$ is

%dsp-level:0
\[
\mathcal H_{\mathrm{slot}}(\Sigma):=\operatorname{span}_{\mathbb C}\{|e\rangle:\ e\in\partial\Sigma\}\cong\mathbb C^{N(\Sigma)},
\]

one basis ket per boundary edge — the direct generalization of D3.1's "one ket per incident real edge," taken over $\Sigma$'s external interface in place of a single vertex's incident set.

\textbf{Degenerate-case check ($\Sigma=\{v\}$).} $\Sigma=\{v\}$ has no internal edges [A1: $p\ne q$, no self-loops, so no edge has both endpoints in $\{v\}$] $\implies\partial\Sigma=$ the full incident edge set of $v$ $\implies N(\Sigma)=n(v)$ [D0.8] $\implies\mathcal H_{\mathrm{slot}}(\Sigma)=\mathcal H^{(n(v))}_{\mathrm{slot}}$ exactly [D3.1]. The construction specializes to the single-vertex slot algebra on the nose, not merely by analogy.

\textbf{Rmk (relation to $N_{\mathrm{visible}}$).} At a single tick ($\delta\tau=1$), $\delta h(e)=1$ per real edge [CT0.9], so $N(\Sigma)=\delta h(\gamma(o,\Sigma))=N_{\mathrm{visible}}$ [FT.1, Def-H.9] for any single external observer whose channel is exactly $\partial\Sigma$: the slot-algebra dimension and the already-established externally-visible bit count coincide. No new counting primitive introduced.

\vspace{2pt}\noindent\rule{\textwidth}{0.4pt}\vspace{2pt}

\end{definition}

\begin{definition}[{\texorpdfstring{Def-H.13 — Compositional bridge map $\iota_\Sigma$}{Def-H.13 - Compositional bridge map}}]\label{tag:Def-H-13}
\textbf{Def.} For each $e\in\partial\Sigma$, let $w_e$ be $e$'s endpoint outside $\Sigma$, with its own rung $n(w_e)$, existing bridge map $\iota_{n(w_e)}$ [D2.12], and $e$'s own edge-local pair-channel decomposition [CT0.12, rung-independent, fixed 16-dim block]. Define $\iota_\Sigma:\mathfrak u(N(\Sigma))\to\mathcal K$ by requiring, on the slot direction $|e\rangle$, that $\iota_\Sigma$ \textbf{factor through} the pair $\bigl(\text{CT0.12's block for }e,\ \iota_{n(w_e)}\bigr)$:

%dsp-level:0
\[
\iota_\Sigma\big|_{|e\rangle}:=\bigl(\text{CT0.12 pair-block map for }e\bigr)\,\otimes\,\iota_{n(w_e)}\qquad(e\in\partial\Sigma),
\]

extended linearly across $\partial\Sigma$. \textbf{No independent embedding is postulated}: every component of $\iota_\Sigma$ is inherited, edge-by-edge, from machinery that already exists at the single-edge (CT0.12) and single-vertex (D2.12) level.

\textbf{"Strictly new at $\Sigma$" (generalizes D3.2).} A piece of $\mathfrak u(N(\Sigma))$ is strictly new at $\Sigma$ iff it does \textbf{not} factor through $\iota_\Sigma$ as defined above — i.e. iff it is a genuinely collective effect of $\Sigma$'s boundary edges taken together, not a sum/tensor of independent single-edge and single-vertex facts already accounted for at lower structural levels. This is the precise sense in which the coarse-graining machine treats $\Sigma$ "as a single quantum object": an aspect belongs to $\mathcal A(\Sigma)$ properly only if its action on $\mathcal H_{\mathrm{slot}}(\Sigma)$ resists the boundary-edgewise factorization $\iota_\Sigma$ provides for free.

\textbf{Degenerate-case check.} At $\Sigma=\{v\}$, $\partial\Sigma$ is $v$'s own incident edge set and each $w_e$ ranges over $v$'s neighbours; $\iota_\Sigma$ reduces to exactly the rung-$n(v)$ construction CT0.13 already performs (its Pf. builds $\mathfrak u(n)$ directly from the slot algebra and reads off $S$-even/odd blocks — the same data $\iota_\Sigma$ assembles edge-by-edge here). $\iota_{\{v\}}=\iota_{n(v)}$.

\textbf{Rmk (strictly new at $\Sigma$ — first counts).} Def-H.13's filter evaluated at small $N(\Sigma)$: $N=1$: $\mathfrak u(1)$ factors per-edge entirely — nothing strictly new; single-contact floors carry no collective boundary mode.

\vspace{2pt}\noindent\rule{\textwidth}{0.4pt}\vspace{2pt}

\end{definition}

\begin{theorem}[{\texorpdfstring{FT.21 — $\mathcal A(\Sigma)$'s charge-type rows, properly derived}{FT.21 - 's charge-type rows, properly derived}}]\label{tag:FT-21}
\textbf{Thm.} $(i)$ \textbf{EM — boundary-exact.} FT.20's EM-row \textit{observable} is boundary-expressible: the tuple read $q=\tfrac12(\sigma_s-\sigma_t)$ [Def-H.9(1), EM.9] reads end signs on $\partial\Sigma$'s terminal edges, and interior vertex activity telescopes into end reads [EM.9b (Part 2): the identity, with its type-$(1,1)$ domain] — so the row's criterion ($\exists v\in\Sigma:\ f_v\ne0$) can exceed what is separately observable at the boundary, while its observable content factors through $\iota_\Sigma$'s per-edge CT0.12 factor [Def-H.13]. $(ii)$ \textbf{$G/H$ — boundary-expressible modulo one declared interior label.} $I_{\mathrm{hid}}(\Sigma;o)=(E_{\mathrm{int}}+E_{\mathrm{bd}})\ln2-E_{\mathrm{vis}}$, where $E_{\mathrm{bd}}=N(\Sigma)=\dim\mathcal H_{\mathrm{slot}}(\Sigma)$ [Def-H.12 Rmk] and $E_{\mathrm{vis}}=\min(\delta h(\gamma),4\,\delta\tau)\ln2$ is a boundary/channel sum [FT.1]; $E_{\mathrm{int}}=|\Sigma\setminus\partial\Sigma|$ is $\Sigma$'s own interior volume datum — a single integer carried as a c-number label alongside $\mathcal H_{\mathrm{slot}}(\Sigma)$, of the same species as Def-H.15(4)'s superselection labels $(\Theta_{\mathrm{int}},k)$, and $\neg$factored through $\iota_\Sigma$ edge-wise. $(iii)$ Consequently both rows are instances of CT0.13's active-aspect machine on $\mathcal H_{\mathrm{slot}}(\Sigma)$ — EM exactly, $G/H$ modulo the declared interior label — not case-by-case arithmetic extensions.

\textbf{Req.} Def-H.9(1)–(2), Def-H.12, Def-H.13, Def-H.15(4); FT.1, FT.20 (EM and $G$ rows); D2.12, D3.1, D3.2, D3.3; CT0.12, CT0.13 (Part 1); EM.9, EM.9b (\textbf{E}; telescoping identity; H-side record: \texttt{verify\_\allowbreak{}H\_\allowbreak{}ASigma.py}).

\textbf{Pf.}

%dsp-level:0
\[
\mathrm{EM}:\ q=\tfrac12(\sigma_s-\sigma_t)\ [\mathrm{Def\text{-}H.9}(1),\ \mathrm{EM.9}]:\ \text{end signs on}\ \partial\Sigma\text{'s terminal edges}\implies\text{boundary read}
\]

%dsp-level:0
\[
\text{interior telescoping}:\ \mathrm{EM.9b}\ (\text{E})\ [\text{H-side record}:\ \mathrm{verify\_H\_ASigma.py}]
\]

%dsp-level:0
\[
G:\ I_{\mathrm{hid}}=(E_{\mathrm{int}}+E_{\mathrm{bd}})\ln2-E_{\mathrm{vis}}\ [\mathrm{Def\text{-}H.9}(2),\ \mathrm{FT.17}\ \text{form}];
\]

%dsp-level:0
\[
\quad E_{\mathrm{bd}}=N(\Sigma)=\dim\mathcal H_{\mathrm{slot}}(\Sigma)\ [\mathrm{Def\text{-}H.12}\ \mathrm{Rmk}];
\]

%dsp-level:0
\[
\quad E_{\mathrm{vis}}=\text{boundary/channel sum}\ [\mathrm{FT.1}];
\]

%dsp-level:0
\[
\quad E_{\mathrm{int}}=|\Sigma\setminus\partial\Sigma|=\Sigma\text{'s fixed volume datum}
\]

%dsp-level:0
\[
\implies\text{EM row factors through}\ \iota_\Sigma\ \text{exactly};
\]

%dsp-level:0
\[
\quad G\ \text{row factors through}\ \iota_\Sigma\ \text{extended by the c-number label}\ E_{\mathrm{int}}\ [\mathrm{Def\text{-}H.15}(4)\ \text{pattern}];
\]

%dsp-level:0
\[
\quad\text{both are instances of CT0.13's machine on}\ \mathcal H_{\mathrm{slot}}(\Sigma),
\]

%dsp-level:0
\[
\ \neg\text{separate arguments}.\ \blacksquare
\]

\textbf{Rmk (why the $W$/spin row resists boundary treatment — and where its quantization comes from instead).} FT.20's row $W$ depends on $\Theta_{\mathrm{int}}(\Sigma)=\sum_{e\in\Sigma\setminus\partial\Sigma}\widehat{ds}^{\,2}(e)$ [Def-H.9(3)] — a sum over \textbf{internal} edges, which are by construction outside $\mathcal H_{\mathrm{slot}}(\Sigma)$'s span [Def-H.12]. The spin/holonomy datum is therefore not expressible in $\iota_\Sigma$'s boundary-edgewise vocabulary: it is a global fact about $\Sigma$, structurally of the same kind as colour contraction-completeness (FT.18/FT.19: non-locality in ambient data). Contrast FT.21(ii): $E_{\mathrm{int}}$ is a \textit{fixed} volume label, while $\Theta_{\mathrm{int}}$ is configuration-dependent interior data — the $W$ row resists even the label-extended factorization. Tier F.8 supplies the quantization this structure calls for: internal-edge line elements are quantized neither by boundary factorization nor by a single-edge reality-line argument (which reaches $\sigma_e$ [CT0.9], not the metric coordinates), but by the corpus's own materialization condition on closed chains (FT.22, via EW.18/CT2.5) evaluated on bound structures — at the fermion floor yielding per-edge integrality (FT.23–FT.24) and a well-defined $\mathrm{Hol}_{\mathrm{int}}$ (FT.25). The internal-edge slot object itself: Def-H.14.

\textbf{Phys.} $\Sigma$ becomes a single quantum object exactly to the extent that its boundary edges can be treated as one composite slot register — and the EM read turns out to be expressible in that register exactly, the mass read modulo $\Sigma$'s own interior volume label $E_{\mathrm{int}}$ (clause (ii)), inherited edge-by-edge from machinery already proven at the single-edge and single-vertex level, with nothing new postulated. Spin is different in kind: it is a fact about $\Sigma$'s \textit{inside}, invisible to the boundary register by construction — and its quantization accordingly comes from the global materialization condition on closed chains (Tier F.8), not from boundary-register arithmetic.

\vspace{2pt}\noindent\rule{\textwidth}{0.4pt}\vspace{2pt}

\end{theorem}

\phantomsection
\section[Tier F.8 — Line-Element Quantization from Materialization]{\texorpdfstring{Tier F.8 — Line-Element Quantization from Materialization}{Tier F.8 - Line-Element Quantization from Materialization}}

\begin{definition}[{\texorpdfstring{Def-H.14 — General-edge sector space (internal-edge slot object)}{Def-H.14 - General-edge sector space (internal-edge slot object)}}]\label{tag:Def-H-14}
\textbf{Def.} For any real edge $e$ (boundary or internal to any $\Sigma$), let $\mathcal E(e):=\mathcal E_{++}(e)\oplus\mathcal E_{+-}(e)\subset\mathfrak H(\mathbb C^4)_e$ be $e$'s own $S$-even pair sector [CT1.1, CT1.2] — the space with coordinates $(a_{00}(e),q_{ij}(e),w_i(e))$ [D2.1] that already fix $g_{\mu\nu}(e)$ [CT1.12] and hence $\widehat{ds}^{\,2}(e)$ [D1.6]. For $\Sigma$ a bounded subnetwork:

%dsp-level:0
\[
\mathcal H_{\mathrm{slot}}^{\mathrm{int}}(\Sigma):=\bigoplus_{e\in\Sigma\setminus\partial\Sigma}\mathcal E(e),\qquad\dim_{\mathbb R}\mathcal H_{\mathrm{slot}}^{\mathrm{int}}(\Sigma)=10\,|\Sigma\setminus\partial\Sigma|.
\]

\textbf{No bridge map needed.} Unlike $\iota_\Sigma$ [Def-H.13], which had to reach across each boundary edge to the \textit{external} endpoint's vertex-level bridge map $\iota_{n(w_e)}$ (because a boundary edge's far side is genuinely outside $\Sigma$), $\mathcal E(e)$ is already fully edge-local [CT0.12: rung-independent, fixed 16-dim block] — it references neither endpoint's rung nor vertex-level slot algebra. An internal edge's two endpoints are both already inside $\Sigma$, which is exactly why no external anchor is available or needed: $\mathcal E(e)$ is defined identically whether $e$ is internal or a boundary edge, and $\mathcal H_{\mathrm{slot}}^{\mathrm{int}}(\Sigma)$ is simply this data restricted to $\Sigma\setminus\partial\Sigma$.

\textbf{Degenerate-case check ($\Sigma=\{v\}$).} No internal edges [A1, as in Def-H.12] $\implies\mathcal H_{\mathrm{slot}}^{\mathrm{int}}(\{v\})=0$, consistent with Def-H.12's finding at the same degenerate case.

\textbf{Rmk (relation to $\Theta_{\mathrm{int}}$).} $\Theta_{\mathrm{int}}(\Sigma)=\sum_{e\in\Sigma\setminus\partial\Sigma}\widehat{ds}^{\,2}(e)$ [Def-H.9(3)] is, by CT1.12/D1.6, a fixed function of the coordinates of $\mathcal H_{\mathrm{slot}}^{\mathrm{int}}(\Sigma)$ together with the embedding data $\{\Delta x_e^\mu\}$. The internal-edge slot object is therefore exactly this space: compositional, edge-local, and implicit in Def-H.9(3).

\textbf{Rmk (ambient continuity).} $\mathcal E(e)$ is a continuous configuration arena: $a_{00}\in\mathbb R$ unrestricted [CT1.12], $\widehat{ds}^{\,2}(e)$ a generic $O(1)$ real [CT1.19a]. CT0.9's reality-line mechanism — the corpus's edge-level discreteness result for $\phi_e$ — acts on the 1-dimensional single-qubit line $\wedge^2\mathcal H_q$ [$\dim\mathcal H_q=2$, A1], not on these coordinates [CT0.9 Req/Pf]; a reality-line analogue cannot reach $\widehat{ds}^{\,2}$. It does not need to: the corpus's pole mechanism (CT2.5) already produces discreteness from continuous families — EW.18 quantizes wrapping loops by exactly this route. FT.22–FT.25 compute that mechanism on bound chains.

\vspace{2pt}\noindent\rule{\textwidth}{0.4pt}\vspace{2pt}

\end{definition}

\begin{theorem}[{\texorpdfstring{FT.22 — Loop quantization of materialized bound chains}{FT.22 - Loop quantization of materialized bound chains}}]\label{tag:FT-22}
\textbf{Thm.} Let $c$ be a closed chain [D1.4] that \textbf{materializes} — contributes a spectral pole of the masked resolvent $Z(\mathcal R,n)$ [D2.4, D2.5, CT2.5] at the rung $n$ of its active mask [D1.12, CT0.13]; instances: materialized wrapping chains [EW.17, EW.18], bound chains of fermion realizations [Def-H.4, D1.5]. Then the spiral holonomy around $c$ is single-valued:

%dsp-level:0
\[
\operatorname{Hol}_{\mathrm{spiral}}(c)=e^{i\Phi(c)}=+1,\qquad\Phi(c)=\alpha\,\Theta(c)=2\pi k,\quad k\in\mathbb Z\ \ (\text{winding});
\]

and under EM.20's scope tokens (Comm-min $+$ Eff), with rung stability [FT.9: $\alpha_{\mathrm{EM}}=\pi$ at every accessible rung]:

%dsp-level:0
\[
\Theta(c)=\sum_{e\in c}\widehat{ds}^{\,2}(e)=2k\in2\mathbb Z .
\]

On rung-3 wrapping chains this reduces verbatim to EW.18 (consistency check).

\textbf{Req.} D1.4, D1.5, D1.7; D2.4, D2.5; CT2.5; EW.17, EW.18 (mechanism; domain generalized from wrapping chains to materialized closed chains — the single H-level step, flagged in Pf); EM.20 (Comm-min, Eff; A §0.2a); FT.9; Def-H.4; D1.12, CT0.13.

\textbf{Pf.}

%dsp-level:0
\[
\mathrm{D1.7}:\ \Theta(\gamma)=\textstyle\sum_{e\in\gamma}\widehat{ds}^{\,2}(e),\ \Phi=\alpha\Theta,\ \operatorname{Hol}_{\mathrm{spiral}}=e^{i\Phi}\ \text{are defined for every chain};
\]

%dsp-level:0
\[
\qquad \mathrm{EW.17}:\ \text{transport around a closed cycle multiplies the eigenmode by}\ \operatorname{Hol}_{\mathrm{spiral}}
\]

%dsp-level:0
\[
\mathrm{CT2.5}\ (\text{as applied in EW.18}):\ \text{pole}
\]

%dsp-level:0
\[
= \text{eigenvalue condition on a single-valued eigenmode}
\]

%dsp-level:0
\[
\implies e^{i\Phi(c)}\qquad = 1;
\]

%dsp-level:0
\[
\qquad \text{winding freedom}\ [\mathrm{EW.18}]:\ \Phi(c)=2\pi k,\ k\in\mathbb Z
\]

%dsp-level:0
\[
\text{domain step (H-level, flagged)}:\ \mathrm{EW.17}\text{'s rung-3 restriction is an existence claim about wrapping chains};
\]

%dsp-level:0
\[
\ \mathrm{CT2.5}\text{'s single-valuedness step is chain-generic}\implies\text{the domain extends to every materialized closed chain with no added premise}
\]

%dsp-level:0
\[
\mathrm{EM.20}\ (\text{Comm-min}+\text{Eff}):\ \alpha=\pi;
\]

%dsp-level:0
\[
\qquad \mathrm{FT.9}:\ \alpha^{(n)}=\pi\ \forall\ \text{accessible rungs}\implies\Phi(c)=\pi\,\Theta(c)=2\pi k\implies\Theta(c)=2k\in2\mathbb Z.\ \blacksquare
\]

\textbf{Rmk (oddness of $\alpha$).} $\Theta(c)$ is constrained exactly when $\alpha/\pi$ is odd: Comm-min alone gives $\alpha\in\pi\mathbb Z$ [EM.20], and even branches ($\alpha\in2\pi\mathbb Z$) leave $\Theta(c)$ unconstrained. Eff — the corpus's sole effectivity hypothesis — is what selects $\alpha=\pi$ and makes the quantization bite.

\textbf{Rmk (sign-factor separation).} Following EW.18, single-valuedness is imposed on the spiral/geometric transport; the registration signs $\sigma_e$ and the wall parity $(-1)^{\aleph}$ are ledger data of the weight $w(C)$ [A8, CT0.9, D3.3], not of the geometric eigenmode transport.

\textbf{Rmk (area language, scoped).} $\Theta(c)\,\mathfrak h^2=\sum_{e\in c}\kappa_e\,\mathfrak h^2$ is the $\kappa$-\textit{weighted accumulated line element} along $c$ — dimensionally an area, and equal to CT1.29's pierced-edge area $A=n_{\mathcal R}\,\mathfrak h^2$ exactly on the minimal-edge class $\kappa_e\equiv1$ [FT.31(4): $\Theta=\delta h$]. It is $\neg$the metric-enclosed area of $c$: the floor loop (2-cycle) traces a degenerate out-and-back locus — zero enclosed area — yet carries $\Theta=2\kappa\ne0$, because $ds^2$ is quadratic, hence orientation-blind [FT.23–FT.24]. Nor is $\sqrt{\Theta}\,\mathfrak h$ a proper length of an inhomogeneous path: proper length reads $\sum_e\sqrt{\kappa_e}\,\mathfrak h$ (sum of square roots $\ne$ square root of the sum). Program-level "cell"/"enclosed area" phrasing must be read as this weighted accumulation.

\textbf{Phys.} Every bound structure that appears in the spectrum obeys a Bohr–Sommerfeld-type selection rule: its loop line-element sum is an even integer. In continuum language [CT1.19, CT2.9]: $\oint_c ds^2=\Theta(c)\,\mathfrak h^2=2k\,\mathfrak h^2=8k\,G_N\ln2$ — the accumulated line element of an on-shell loop is an even multiple of the Bekenstein bit-area $\mathfrak h^2$ [CT1.29, CT2.9]; on the minimal-edge class this equals CT1.29's pierced-edge area [FT.31(4)], in general it is a $\kappa$-weighted accumulation, $\neg$a metric-enclosed area (Rmk below). Smooth geometry unfolds as the large-$k$ limit; integrality persists at every scale as the rule selecting which loops materialize.

\vspace{2pt}\noindent\rule{\textwidth}{0.4pt}\vspace{2pt}

\end{theorem}

\begin{theorem}[{\texorpdfstring{FT.23 — Orientation-evenness of the line element}{FT.23 - Orientation-evenness of the line element}}]\label{tag:FT-23}
\textbf{Thm.} For a bidirectional channel — both orientations $e_{pq}=(p,q)$, $e_{qp}=(q,p)$ in $\mathcal R$:

%dsp-level:0
\[
\widehat{ds}^{\,2}(e_{pq})=\widehat{ds}^{\,2}(e_{qp}).
\]

\textbf{Req.} D1.6; CT2.6 (charts; $g_{\mu\nu}$ a single smooth field on $M$; hypotheses (i)–(iii)); CT1.12 ($g_{\mu\nu}$ built entirely from $S$-even sector data); CT1.1, CT1.2, C3 ($\mathcal E_{++},\mathcal E_{+-}$ are $S$-even); CT0.7(2) ($S(X)=\mathbb S X\mathbb S$: the swap action); endpoint-difference displacement [math].

\textbf{Pf.}

%dsp-level:0
\[
\mathrm{CT2.6}\ \text{charts}:\ \Delta x_e:=x_{t(e)}-x_{s(e)}\implies\Delta x_{e_{qp}}=-\Delta x_{e_{pq}}\ \ (\text{all four components})
\]

%dsp-level:0
\[
\qquad [\text{consistent with CT1.19's}\ \Delta x^\mu_e=\mathfrak h\,\dot x^\mu(\lambda_e)+O(\mathfrak h^2)]
\]

%dsp-level:0
\[
g_{\mu\nu}(e)\ \text{built solely from}\ (a_{00},q_{ij},w_i)\in\mathcal E_{++}\oplus\mathcal E_{+-}\ [\mathrm{CT1.12},\ \mathrm{D2.1}];
\]

%dsp-level:0
\[
\qquad \mathcal E_{++},\ \mathcal E_{+-}\ \text{are}\ S\text{-even}\ [\mathrm{CT1.1},\ \mathrm{CT1.2},\ \mathrm{C3}]
\]

%dsp-level:0
\[
S(X)=\mathbb S X\mathbb S\ [\mathrm{CT0.7}(2)];
\]

%dsp-level:0
\[
\qquad SX
\]

%dsp-level:0
\[
= X\implies\mathbb S X\mathbb S
\]

%dsp-level:0
\[
= X:\ \text{metric data invariant under}\ p\leftrightarrow q
\]

%dsp-level:0
\[
\implies g_{\mu\nu}(e_{pq})
\]

%dsp-level:0
\[
= g_{\mu\nu}(e_{qp})\ \ (\text{same channel locus}\ [\mathrm{CT2.6}])
\]

%dsp-level:0
\[
\mathrm{D1.6}:\ ds^2(e)
\]

%dsp-level:0
\[
= g_{\mu\nu}(e)\,
\]

%dsp-level:0
\[
\Delta x^\mu_e\Delta x^\nu_e\ \text{is quadratic in}\ \Delta x:\ \text{every monomial has}\ \Delta x\text{-degree}\ 2\implies\Delta x
\]

%dsp-level:0
\[
\mapsto -\Delta x\ \text{fixes}\ ds^2
\]

%dsp-level:0
\[
(\text{incl. the shift cross-term}\ 2g_{0i}\Delta x^0\Delta x^i:\ \text{both factors flip});
\]

%dsp-level:0
\[
\qquad \implies ds^2(e_{pq})=ds^2(e_{qp})\implies\widehat{ds}^{\,2}(e_{pq})=\widehat{ds}^{\,2}(e_{qp})\quad[\mathrm{D1.6}:\ /\mathfrak h^2].\ \blacksquare
\]

\textbf{Phys.} Both orientations of one channel carry a single geometric datum: the line element is a property of the channel; orientation lives in the registration signs and transport order, not in the metric. This is what lets the minimal closed chain — the 2-cycle — speak about one number $\widehat{ds}^{\,2}$.

\vspace{2pt}\noindent\rule{\textwidth}{0.4pt}\vspace{2pt}

\end{theorem}

\begin{theorem}[{\texorpdfstring{FT.24 — The minimal fermion pins the geometry quantum}{FT.24 - The minimal fermion pins the geometry quantum}}]\label{tag:FT-24}
\textbf{Thm.} Let $c=(e_1,e_2)$, $e_1=(p,q)$, $e_2=(q,p)$, be a materialized 2-cycle — the minimal closed chain [A1 excludes self-loops; D1.4; Def-H.4 floor: $\delta h(c)\ge2$]. Write $\widehat{ds}^{\,2}(e_1)=\widehat{ds}^{\,2}(e_2)=:\kappa$ [FT.23]. Then:

%dsp-level:0
\[
(1)\quad\Theta(c)=2\kappa\in2\mathbb Z\ \ [\mathrm{FT.22}]\implies\kappa\in\mathbb Z
\]

%dsp-level:0
\[
\qquad (\text{per-edge integrality at the materialized fermion floor});
\]

%dsp-level:0
\[
(2)\quad ds^2(e)=\kappa\,\mathfrak h^2\in\mathfrak h^2\,\mathbb Z\qquad(\text{line element quantized in hop quanta}\ [\mathrm{D1.6}]);
\]

%dsp-level:0
\[
(3)\quad\mathfrak h^2=4\,G_N\ln2\ \ [\mathrm{CT2.9}]\implies ds^2(e)=4\kappa\,G_N\ln2
\]

%dsp-level:0
\[
\qquad (\text{interval quantum}=\text{one Bekenstein bit-area}\ [\mathrm{CT1.29}]\ \text{per unit}\ \kappa);
\]

%dsp-level:0
\[
(4)\quad\text{minimal real edge}\ [\mathrm{D1.1},\ \text{Comm-min}]:\ \kappa=1\ \text{attained}\ (\text{floor tight});
\]

%dsp-level:0
\[
\qquad \kappa
\]

%dsp-level:0
\[
= 0\iff\text{null edge}\ [\mathrm{CT1.19a}]\ (\text{degenerate floor};\ \text{all-null transit chains are the trivial chains}=\text{photons/gluons}\ [\mathrm{D1.7},\ \mathrm{D3.7}]);
\]

%dsp-level:0
\[
\qquad \kappa<0:\ \text{timelike}\ [\mathrm{CT1.12}:\ \operatorname{sig}(g)=(-,+,+,+)].
\]

\textbf{Req.} FT.22, FT.23; A1; D1.1, D1.4, D1.6, D1.7; Def-H.4; CT1.12, CT1.19a, CT1.29, CT2.9 (Part 2); D3.7; EM.20's scope tokens (inherited via FT.22).

\textbf{Pf.}

%dsp-level:0
\[
(1):\ \Theta(c)=\widehat{ds}^{\,2}(e_1)+\widehat{ds}^{\,2}(e_2)=2\kappa\ [\mathrm{FT.23}];\qquad\mathrm{FT.22}:\ \Theta(c)\in2\mathbb Z\implies\kappa\in\mathbb Z
\]

%dsp-level:0
\[
(2):\ \mathrm{D1.6}:\ ds^2(e)=\widehat{ds}^{\,2}(e)\,\mathfrak h^2=\kappa\,\mathfrak h^2;
\]

%dsp-level:0
\[
\qquad (3):\ \mathrm{CT2.9}\ \text{substitution};\quad\mathrm{CT1.29}:\ S=\frac{A\ln2}{\mathfrak h^2}\ \text{at}\ A=\mathfrak h^2:\ S=\ln2=1\ \text{bit per unit}\ \kappa
\]

%dsp-level:0
\[
(4):\ \mathrm{D1.1}:\ \text{minimal real edge has}\ \widehat{ds}^{\,2}=1\implies\kappa=1;
\]

%dsp-level:0
\[
\qquad \mathrm{CT1.19a}:\ \widehat{ds}^{\,2}=0\iff\Delta x_e\ \text{null};\quad\mathrm{D1.7}/\mathrm{D3.7}:\ \text{all-null chains}=\text{photons/gluons}
\]

%dsp-level:0
\[
\kappa\in\mathbb Z\ \wedge\ \operatorname{sig}(g)=(-,+,+,+)\ [\mathrm{CT1.12}]\implies\kappa<0\ \text{admissible (timelike edges)}.\ \blacksquare
\]

\textbf{Rmk (two integer reads at the floor).} The materialized minimal bound structure carries two derived integer reads: the geometric winding $\kappa=\Theta(c)/2$ [FT.22, clause (1)] and the informational mass $m_{\mathrm{hid}}=I_{\mathrm{hid}}/\mathfrak h$ [Def-H.9(2), FT.17: $I_{\mathrm{hid}}=E_{\mathrm{int}}\ln2$, an internal-edge count].

\textbf{Phys.} Computing the smallest fermion the corpus allows — the materialized 2-cycle — settles the $ds^2$ question on-shell: geometry comes in integer multiples of the hop area, and the hop area is itself the calibrated information quantum $4G_N\ln2$. The spacetime quantum is not a new postulate; it is the fermion floor read through the corpus's own pole condition. Off-shell configurations in $Z$ keep their generic continuous $\widehat{ds}^{\,2}$ [CT1.19a]: integrality is a property of the spectrum, not of the summation domain — exactly as in CT2.5.

\vspace{2pt}\noindent\rule{\textwidth}{0.4pt}\vspace{2pt}

\end{theorem}

\begin{theorem}[{\texorpdfstring{FT.25 — Internal holonomy is well-defined on-shell}{FT.25 - Internal holonomy is well-defined on-shell}}]\label{tag:FT-25}
\textbf{Thm.} $\Sigma$ a bounded subnetwork [Def-H.9]. Then:

%dsp-level:0
\[
(1)\quad\Sigma\setminus\partial\Sigma\ \text{decomposes edge-disjointly into materialized closed chains}\ \{c_i\}:
\]

%dsp-level:0
\[
\qquad \Theta_{\mathrm{int}}(\Sigma)=\sum_i\Theta(c_i)\in2\mathbb Z,
\]

%dsp-level:0
\[
\qquad \mathrm{Hol}_{\mathrm{int}}=(-1)^{\Theta_{\mathrm{int}}}=+1\ \ [\mathrm{FT.22}];
\]

%dsp-level:0
\[
(2)\quad\text{every internal edge carries on-shell integer line element}\ [\mathrm{FT.24}]:
\]

%dsp-level:0
\[
\qquad \Theta_{\mathrm{int}}(\Sigma)\in\mathbb Z,
\]

%dsp-level:0
\[
\qquad \mathrm{Hol}_{\mathrm{int}}=(-1)^{\Theta_{\mathrm{int}}}\in\{\pm1\}\ \ \text{well-defined};
\]

%dsp-level:0
\[
(3)\quad\text{the surviving on-shell integer label is the winding}\ k(c)
\]

%dsp-level:0
\[
= \Theta(c)/2\in\mathbb Z\ \ [\mathrm{FT.22}];
\]

%dsp-level:0
\[
\ \ \text{at the floor}\ k
\]

%dsp-level:0
\[
= \kappa\ \ [\mathrm{FT.24}(1)];
\]

%dsp-level:0
\[
\ \ \text{full materialized loops have}\ \mathrm{Hol}_{\mathrm{spiral}}
\]

%dsp-level:0
\[
= +1,
\]

%dsp-level:0
\[
\ \text{so the binary}\ \mathrm{Hol}_{\mathrm{int}}\ \text{is nontrivial only on partial-chain internal sets (proper subsets of materialized chain edges), where}\ \Theta_{\mathrm{int}}\ \text{is a partial sum of integers with free parity}.
\]

\textbf{Req.} FT.22, FT.24; Def-H.9, Def-H.14; D1.7; FT.9 ($\alpha=\pi$); CT1.19a (off-shell contrast).

\textbf{Pf.}

%dsp-level:0
\[
(1):\ \text{edge-disjoint decomposition}:\ \Theta_{\mathrm{int}}
\]

%dsp-level:0
\[
= \!\!\sum_{e\in\Sigma\setminus\partial\Sigma}\!\!\widehat{ds}^{\,2}(e)
\]

%dsp-level:0
\[
= \sum_i\sum_{e\in c_i}\widehat{ds}^{\,2}(e)
\]

%dsp-level:0
\[
= \sum_i\Theta(c_i)
\]

%dsp-level:0
\[
\in 2\mathbb Z\ \ [\mathrm{FT.22}];
\]

%dsp-level:0
\[
\qquad (-1)^{2\mathbb Z}=+1
\]

%dsp-level:0
\[
(2):\ \text{finite sum of integers}\ [\mathrm{FT.24}];
\]

%dsp-level:0
\[
\qquad \mathrm{Hol}_{\mathrm{int}}=e^{i\pi\Theta_{\mathrm{int}}}=(-1)^{\Theta_{\mathrm{int}}}\ \text{well-defined}\ [\alpha=\pi:\ \mathrm{FT.9}]
\]

%dsp-level:0
\[
(3):\ k:=\Theta/2\in\mathbb Z\ [\mathrm{FT.22}];\quad\text{floor}:\ \Theta=2\kappa\implies k=\kappa\ [\mathrm{FT.24}(1)];
\]

%dsp-level:0
\[
\qquad \text{full loop}:\ \mathrm{Hol}_{\mathrm{spiral}}=e^{i\pi\cdot2k}=+1
\]

%dsp-level:0
\[
\text{partial set}\ P\subsetneq E(c):\ \Theta_{\mathrm{int}}
\]

%dsp-level:0
\[
= \sum_{e\in P}\widehat{ds}^{\,2}(e)\in\mathbb Z\ \text{by}\ (2),
\]

%dsp-level:0
\[
\ \text{parity free}\ (\text{witness}: c\ \text{with per-edge}\ (1,1,1,3),\ \Theta=6\in2\mathbb Z;\ P=\text{first three}:\ \Theta_{\mathrm{int}}=3,\ \mathrm{Hol}_{\mathrm{int}}=-1;\ \texttt{verify\_H\_ds2.py}).\ \blacksquare
\]

\textbf{Rmk (no layer-local integrality postulate).} No layer-local postulate of general-edge line-element integrality — $\widehat{ds}^{\,2}(e)\in\mathbb Z$ for \textit{every} real edge — is used. Its on-shell content is derived (FT.22–FT.25, from the corpus's own materialization mechanism); its off-shell generalization is neither derivable nor needed — every use site in this document (Def-H.7, Def-H.9(3), FT.14's parity read, FT.15's $\mathbf S$ clause, FT.20 row $W$) concerns bound/fermion content and is scoped on-shell at its site. The axioms remain self-contained: no step above adds a premise beyond the corpus's declared ones (Comm-min, Eff — A §0.2a).

\textbf{Rmk (spin-extraction consequence).} For $\Sigma$ swept entirely by fermion chains, clause (1) forces $\mathrm{Hol}_{\mathrm{int}}=+1$ on-shell: the \textit{parity} of a full materialized loop cannot carry the spin label. The live candidates are the integer data — the winding $k(c)$, partial-sum $\Theta_{\mathrm{int}}$ [clause (3)], weak-sector insertion multiplicities [FT.13], the transported seed's $j=\tfrac12$ carrier [EW.32].

\textbf{Phys.} The internal geometric ledger of a coarse-grained fermion is integer-valued wherever the corpus's own spectral condition reaches: loops materialize in even hop-area quanta, the minimal loop is exactly one quantum per channel, and Def-H.9's binary holonomy label is a well-defined derived object rather than a postulate-conditional one. The remaining question is not whether on-shell geometry is quantized — it is — but which integer label carries spin, and how far the on-shell sublattice extends across the graph.

\vspace{2pt}\noindent\rule{\textwidth}{0.4pt}\vspace{2pt}

\end{theorem}

\phantomsection
\section[Tier F.9 — The Activated Aspect Operator: $\mathcal A(\Sigma)$ as a Computable Mask]{\texorpdfstring{Tier F.9 — The Activated Aspect Operator: $\mathcal A(\Sigma)$ as a Computable Mask}{Tier F.9 - The Activated Aspect Operator: as a Computable Mask}}

\begin{definition}[{\texorpdfstring{Def-H.15 — Coarse-grained aspect space and masked force operator}{Def-H.15 - Coarse-grained aspect space and masked force operator}}]\label{tag:Def-H-15}
\textbf{Def.} $\Sigma$ a bounded subnetwork; $C\in\mathcal M(U)$ the ambient configuration [D2.11].

%dsp-level:0
\[
(1)\quad\mathcal K_\Sigma(C):=\bigoplus_{A\in\mathcal A(\Sigma;C)}\mathcal K^{(A)}
\]

%dsp-level:0
\[
\qquad [\text{the aspect blocks of D2.3; membership by FT.20's criteria table}];
\]

%dsp-level:0
\[
(2)\quad\mathcal P_M^{(\Sigma)}(C):\mathcal K\to\mathcal K_\Sigma(C)\qquad(\text{orthogonal projection}\ [\mathrm{D2.3}/\mathrm{CT2.2}\ \text{pattern}]);
\]

%dsp-level:0
\[
(3)\quad\mathbf F_M^{(\Sigma)}(C):=\mathcal P_M^{(\Sigma)}\,\mathbf F\,\mathcal P_M^{(\Sigma)}\qquad[\mathrm{A8}/\mathrm{C7},\ \mathrm{CT2.3}\ \text{pattern}],
\]

with the strong summand admissibility-filtered:

%dsp-level:0
\[
F_{\mathrm{strong}}\ \longmapsto\ \Pi_{\mathcal C}(\Sigma;C)\,F_{\mathrm{strong}}\big|_\Sigma\,\Pi_{\mathcal C}(\Sigma;C)\qquad[\mathrm{Def\text{-}H.10},\ \mathrm{FT.19}];
\]

%dsp-level:0
\[
(4)\quad\text{realization}:\ \iota_\Sigma\ [\mathrm{Def\text{-}H.13}]\ \text{maps}\ \mathfrak u(N(\Sigma))\ \text{into}\ \mathcal K_\Sigma\text{'s boundary-coupled part}\ [\mathrm{FT.21}];
\]

%dsp-level:0
\[
\ \text{on-shell internal labels}\ (\Theta_{\mathrm{int}},k)\ [\mathrm{FT.25},\ \mathrm{FT.22}]\ \text{superselect the}\ W/G\ \text{reads of Def-H.9's tuple; off-shell, blocks carry arena data}\ [\mathrm{CT1.19a},\ \mathrm{Def\text{-}H.18}].
\]

\textbf{Rmk (configuration-relativity).} Unlike the rung masks $\mathcal P_M^{(n)}$ [D2.3], which are state-independent, $\mathcal K_\Sigma(C)$ depends on the ambient configuration $C$ — FT.20's criteria read $f_v$, $I_{\mathrm{hid}}$, $\Theta_{\mathrm{int}}$, $\mathcal C(\Sigma;C)$, which are state data.

\vspace{2pt}\noindent\rule{\textwidth}{0.4pt}\vspace{2pt}

\end{definition}

\begin{definition}[{\texorpdfstring{Def-H.16 — Coarse-grained resolvent and the $\Sigma$-spectrum}{Def-H.16 - Coarse-grained resolvent and the -spectrum}}]\label{tag:Def-H-16}
\textbf{Def.}

%dsp-level:0
\[
Z(\mathcal R,\Sigma):=\Bigl\langle\mathrm{out}\Bigl|\bigl(1-\mathbf F_M^{(\Sigma)}\bigr)^{-1}\Bigr|\mathrm{in}\Bigr\rangle\qquad[\mathrm{D2.4}\ \text{pattern};\ \mathrm{D2.8}\ \text{sheet}],
\]

with $|\mathrm{in}\rangle,|\mathrm{out}\rangle$ read on boundary-channel states via $\iota_\Sigma$ [Def-H.13] — a constitutive clause of this definition (the $\iota_\Sigma$ boundary-state identification). Poles of $Z(\mathcal R,\Sigma)$ $:=$ the $\Sigma$-spectrum [D2.5 pattern]: pole $\iff1\in\operatorname{spec}\mathbf F_M^{(\Sigma)}$ [CT2.5 pattern]; residues $=$ transition amplitudes; $m^2:=$ pole displacement [D2.5]. A bound chain $c$ interior to $\Sigma$ \textbf{materializes} (FT.22's sense) iff it contributes to a pole of $Z(\mathcal R,\Sigma)$: FT.22's single-valuedness condition \textit{is} this resolvent's eigenmode condition [CT2.5 as applied in EW.18]. \textbf{Rmk (two pole readings, scoped).} The corpus reads "pole" two ways. \textbf{(i) On/off eigenvalue condition} — pole $\iff1\in\operatorname{spec}\mathbf F_M$ [CT2.5]: no momentum label required; this is the reading Def-H.16/FT.26/FT.28(7) use, and FT.22's single-valuedness is exactly this reading applied to the chain's spiral transport. \textbf{(ii) Displacement value} — D2.5's "$m^2:=$ pole displacement" vocabulary presupposes a $k^2$-parametrized resolvent family; it is a corpus-wide definitional layer (D2.4's own $Z(\mathcal R,n)$ likewise carries no explicit $k$), whose values are extracted in \textbf{D} by a transport-defect computation, not a literal momentum-space two-point function: EW.33's $\mathcal U=\exp\bigl(i\oint\sum_Ac_A\mathcal G_A\,d\ell\bigr)=1+i\,\ell_\gamma\sum_Ac_A\mathcal G_A+O(c^2)$ — the \textit{weak/gauge} transport, an object distinct from the spiral factor $e^{i\Phi}$ constrained to $+1$ on-shell by EW.18/FT.22; the two conditions coexist on the same materialized wrapping chains [EW.18 $\wedge$ EW.20: $m_W^2>0$]. At the $\Sigma$-level, reading (ii) is \textbf{not yet constructed}: the premise-free candidate is an EW.33-pattern defect with one interface transfer choice; a literal momentum-space resolvent would need a new declared choice.

\vspace{2pt}\noindent\rule{\textwidth}{0.4pt}\vspace{2pt}

\end{definition}

\begin{theorem}[{\texorpdfstring{FT.26 — Structural transfer: the $\Sigma$-mask inherits CT2.3–CT2.5}{FT.26 - Structural transfer: the -mask inherits CT2.3-CT2.5}}]\label{tag:FT-26}
\textbf{Thm.} $\mathbf F_M^{(\Sigma)}(C)$ is: $(1)$ Hermitian; $(2)$ block-diagonal, $\mathbf F_M^{(\Sigma)}=\bigoplus_{A\in\mathcal A(\Sigma;C)}F_A\big|_{\mathcal K^{(A)}}$, the strong summand $\equiv0$ whenever $\mathcal C(\Sigma;C)=0$; $(3)$ finite-dimensional with

%dsp-level:0
\[
\dim\mathcal K_\Sigma\ \le\ 1+5+1+3+8\ =\ 18\ <\ \infty .
\]

Hence the spectral theorem applies and $Z(\mathcal R,\Sigma)$ has simple poles exactly at $1\in\operatorname{spec}\mathbf F_M^{(\Sigma)}$, residues $-\langle\mathrm{out}|P_a|\mathrm{in}\rangle$ — CT2.5's statement transferred verbatim with $n\mapsto\Sigma$.

\textbf{Req.} Def-H.15, Def-H.16; A8; CT0.12, CT0.14 (block dimensions); CT2.2, CT2.3, CT2.4, CT2.5 (Part 2); FT.19, FT.20; spectral theorem [math].

\textbf{Pf.}

%dsp-level:0
\[
(1):\ \bigl(\mathcal P\mathbf F\mathcal P\bigr)^\dagger=\mathcal P^\dagger\mathbf F^\dagger\mathcal P^\dagger=\mathcal P\mathbf F\mathcal P
\]

%dsp-level:0
\[
\quad [\mathrm{CT2.3}\ \text{Pf pattern}:\ \mathcal P_M^{(\Sigma)}\ \text{orthogonal}\ [\mathrm{D2.3}\ \text{pattern}],\ \mathbf F^\dagger=\mathbf F]
\]

%dsp-level:0
\[
(2):\ \mathrm{A8}:\ \mathbf F=\bigoplus_AF_A,\ \operatorname{supp}F_A\subseteq\mathcal K^{(A)};
\]

%dsp-level:0
\[
\qquad \text{block-subset projection retains exactly the members}\ [\mathrm{CT2.4}\ \text{Pf pattern}]
\]

%dsp-level:0
\[
\text{strong summand}:\ \mathcal C(\Sigma;C)=0\implies\Pi_{\mathcal C}=0\implies F_{\mathrm{strong}}^\Sigma=0\ \ [\mathrm{FT.19}\ \text{Thm};\ \mathrm{Def\text{-}H.15}(3)]
\]

%dsp-level:0
\[
(3):\ \mathcal A(\Sigma;C)\subseteq\{\mathrm{EM},G,H,W,\mathrm{strong}\}\ [\mathrm{FT.20}:\ J^\pm,D^\pm\ \text{never}]\implies\dim\mathcal K_\Sigma
\]

%dsp-level:0
\[
\le 1+5+1+3+8
\]

%dsp-level:0
\[
= 18\ [\mathrm{CT0.12},\ \mathrm{CT0.14}:\ \dim\mathcal K^{(\mathrm{EM})}=1,\ \dim\mathcal K^{(\mathrm{strong})}=8]
\]

%dsp-level:0
\[
\text{finite-dim Hermitian}\implies\text{spectral decomposition}
\]

%dsp-level:0
\[
\implies (1-\mathbf F_M^{(\Sigma)})^{-1}
\]

%dsp-level:0
\[
= \sum_a\frac{P_a}{1-\lambda_a}:\ \text{poles}\iff\lambda_a
\]

%dsp-level:0
\[
= 1;
\]

%dsp-level:0
\[
\ \ \mathrm{D2.8}\ \text{sheet}\ [\mathrm{CT2.5}\ \text{Pf, verbatim}].\ \blacksquare
\]

\textbf{Rmk (the observability deficit, in dimensions).} The rung-3 mask carries $\dim\mathcal K_3=25$ [CT0.14]; the maximal observable $\Sigma$-mask carries $18$. The deficit $25-18=7=3+3+1$ is exactly $J^+\oplus J^-\oplus D^+$ [CT0.12] — the structural aspects FT.20 reclassifies as non-observable (seals, uniform centre). At the $\Sigma$-level, action and observability are different masks, and the difference is a measurable dimension.

\textbf{Phys.} A coarse-grained subnetwork now has the same spectral machine a single vertex has: a Hermitian masked force on a finite aspect space, a resolvent whose poles are the $\Sigma$-particles, and confinement built into the block structure — one unsealed colour attachment anywhere and the entire strong summand vanishes [FT.19]. This is the activation of $\mathcal A(\Sigma)$: membership is a computable nonzero block, not merely a criteria table.

\vspace{2pt}\noindent\rule{\textwidth}{0.4pt}\vspace{2pt}

\end{theorem}

\begin{theorem}[{\texorpdfstring{FT.27 — Single-vertex degeneration; action versus observable}{FT.27 - Single-vertex degeneration; action versus observable}}]\label{tag:FT-27}
\textbf{Thm.} At $\Sigma=\{v\}$: $(a)$ the spaces degenerate on the nose — $\mathcal H_{\mathrm{slot}}(\{v\})=\mathcal H^{(n(v))}_{\mathrm{slot}}$ [Def-H.12], $\mathcal H^{\mathrm{int}}_{\mathrm{slot}}(\{v\})=0$ [Def-H.14], $\iota_{\{v\}}=\iota_{n(v)}$ [Def-H.13]; $(b)$ FT.20's criteria evaluated at $\{v\}$ agree with CT0.13's active set on the EM row (for type-$(1,1)$ states), and refine it on four rows — each refinement a proven fact, not a deviation:

%dsp-level:0
\[
\mathrm{EM}\in\mathcal A(\{v\};C)\iff f_v\ne0\qquad(\text{CT0.13's type-}(1,1)\ \text{condition}\ [\mathrm{D3.3}]);
\]

%dsp-level:0
\[
G\in\mathcal A(\{v\};C)\iff I_{\mathrm{hid}}(\{v\};o)\ne0\iff n(v)\ge2
\]

%dsp-level:0
\[
\qquad [\mathrm{FT.16}:\ I_{\mathrm{hid}}=(n{-}1)\ln2]\ —\ \text{refines}\ G\in\mathcal A(1):\ \text{at}\ n=1\ \text{the block acts but no observable exists};
\]

%dsp-level:0
\[
W\in\mathcal A(\{v\};C)\iff\text{never}
\]

%dsp-level:0
\[
\qquad [\Theta_{\mathrm{int}}(\{v\})=0:\ \text{empty sum},\ \mathrm{Def\text{-}H.14}]\ —\ \text{weak acts from rung}\ 1\ [\mathrm{CT0.13}]\ \text{but its observable is born only at the bound floor}\ (\mathrm{FT.28});
\]

%dsp-level:0
\[
\text{strong}\in\mathcal A(\{v\};C)\iff\text{never}
\]

%dsp-level:0
\[
\qquad [\text{case split}:\ \mathrm{Colour}(\{v\})=\varnothing\ (v\ \text{colourless})\implies\text{FT.20's first conjunct fails}\ \wedge\ F_{\mathrm{strong}}\big|_{\{v\}}=0\ (\neg\exists\ \text{colour content});\quad\mathrm{Colour}(\{v\})=\{v\}\implies\pi_C(v)\subseteq\{v\}\ \text{impossible}\ [\mathrm{A1}:\ \text{partners distinct}]\implies\mathcal C(\{v\};C)=0\implies\text{block}\equiv0\ [\mathrm{FT.19}]]\ \text{—}\ \text{either way the strong block is zero at a single vertex};
\]

%dsp-level:0
\[
\ \text{the smallest colour-observable}\ \Sigma\ \text{is the sealed meson-type pair}\ [\mathrm{D3.9},\ \mathrm{CT3.17};\ \text{realized at}\ \mathrm{FT.28}\ \mathrm{Rmk}];
\]

%dsp-level:0
\[
J^\pm,D^\pm\notin\mathcal A(\{v\};C)\ \ \forall n\qquad[\mathrm{FT.20}\ \text{reclassification}].
\]

\textbf{Req.} Def-H.10, Def-H.12, Def-H.13, Def-H.14, Def-H.15; A1; CT0.13 (Part 1); D0.8, D3.1, D3.3, D3.8, D3.9; CT3.11, CT3.17 (Part 2); FT.16, FT.19, FT.20.

\textbf{Pf.}

%dsp-level:0
\[
(a):\ \text{degenerate-case checks of Def-H.12/Def-H.13/Def-H.14, verbatim}
\]

%dsp-level:0
\[
(b):\ \text{FT.20's table at}\ \Sigma=\{v\}:\ \ \mathrm{EM}\ \text{row}=f_v\ne0\ [\mathrm{D3.3}];
\]

%dsp-level:0
\[
\qquad G\ \text{row}:\ I_{\mathrm{hid}}(\{v\};o)=(n{-}1)\ln2\ne0\iff n\ge2\ [\mathrm{FT.16}]
\]

%dsp-level:0
\[
W\ \text{row}:\ \Sigma\setminus\partial\Sigma=\varnothing\ [\mathrm{A1}:\ \neg\exists\ \text{self-loops}]\implies\Theta_{\mathrm{int}}=0\ (\text{empty sum}):\ \text{criterion cannot fire}
\]

%dsp-level:0
\[
\text{strong row (case split)}:
\]

%dsp-level:0
\[
\qquad \mathrm{Colour}(\{v\})
\]

%dsp-level:0
\[
= \varnothing\ (v\ \text{colourless},\ \text{e.g.}\ n(v)\ne3\ [\mathrm{D3.1}]):\ \text{FT.20's first conjunct fails}\ \wedge\ F_{\mathrm{strong}}\big|_{\{v\}}
\]

%dsp-level:0
\[
= 0\ (\neg\exists\ \text{colour content to act on}\ [\mathrm{D3.8}])\implies\text{block}\qquad \equiv 0
\]

%dsp-level:0
\[
\mathrm{Colour}(\{v\})
\]

%dsp-level:0
\[
= \{v\}\ (v\ \text{colour-carrying}):\ \pi_C(v)\ne\varnothing\ \text{requires a distinct partner}\ [\mathrm{A1}:\ p\ne q]\implies\pi_C(v)\not\subseteq\{v\}\implies\partial_{\mathrm{col}}(\{v\};C)
\]

%dsp-level:0
\[
= \{v\}\implies\mathcal C(\{v\};C)
\]

%dsp-level:0
\[
= 0\ [\mathrm{Def\text{-}H.10}]\implies\text{block}
\]

%dsp-level:0
\[
\equiv 0\ [\mathrm{FT.19}]
\]

%dsp-level:0
\[
\text{deltas vs.}\ \mathrm{CT0.13}\ \text{are exactly}:\ G@n{=}1\ (\text{acts},\ \neg\text{observable}\ [\mathrm{FT.16}]);
\]

%dsp-level:0
\[
\ \ W@\neg\exists\text{-interior}\ (\text{acts},\ \neg\text{observable});
\]

%dsp-level:0
\[
\ \ \text{strong@single vertex}\ (\text{acts at}\ n{=}3\ [\mathrm{CT0.13}],\ \neg\text{observable}:\ \text{colourless}\ —\ \neg\exists\text{content};\ \text{colour-carrying}\ —\ \text{partner distinctness}\ [\mathrm{A1}]);
\]

%dsp-level:0
\[
\ \ J^\pm,
\]

%dsp-level:0
\[
D^\pm\ (\text{structural reclassification}\ [\mathrm{FT.20}\ \mathrm{Rmk};\ \mathrm{CT0.7}/\mathrm{FT.11}\ \text{seals-}\neg\text{charges};\ \mathrm{FT.3}/\mathrm{CT4.4}\ D^-\ \text{darkness}]) — \text{each a proven fact}.\ \blacksquare
\]

\textbf{Rmk (the birth of spin).} The weak observable requires an interior edge; the first interior exists at the bound-structure floor (2-cycle, Def-H.4). A single vertex has weak \textit{action} but no weak \textit{observable}: spin — already identified algebraically as the weak contribution [FT.13, Def-H.8] — is an observable of coarse-graining itself.

\textbf{Phys.} At a single vertex the activated operator reduces to the rung machine CT0.13–CT2.5 on everything observable, and the four deviations are theorems, not defects: gravity becomes observable exactly when something can be hidden ($n\ge2$), spin exactly when something can wind (bound floor), colour exactly when a sealing partner is inside the region read — never at one vertex, whether colourless (no content to act on) or colour-carrying (no inside partner), and seals and locks are never observable at all. Action and observability are two different masks — H now carries both.

\vspace{2pt}\noindent\rule{\textwidth}{0.4pt}\vspace{2pt}

\end{theorem}

\phantomsection
\section[Tier F.10 — The Strong Edge Weight]{\texorpdfstring{Tier F.10 — The Strong Edge Weight}{Tier F.10 - The Strong Edge Weight}}

\begin{definition}[{\texorpdfstring{Def-H.17 — Strong edge transport}{Def-H.17 - Strong edge transport}}]\label{tag:Def-H-17}
\textbf{Def.} For $e\in E_{\mathcal R}$ incident on colour-carrying vertices (rung-3 attachments [D3.8]):

%dsp-level:0
\[
G^{(3)}(e)\ :==\ G^{(2)}(e)\ \otimes\ \tau(e),\qquad
G^{(2)}(e)=\sigma_e\,e^{i\pi\,\widehat{ds}^{\,2}(e)}\quad[\mathrm{D3.5},\ \mathrm{FT.9}],
\]

$\tau(e):=$ the slot-index transport of $e$ on $\mathcal H^{(3)}_{\mathrm{slot}}\cong\mathbb C^3$ [D3.1, CT3.2], acting in the defining representation [CT3.11(2)–(3): triplet $\mapsto U^i{}_jq^j$, antitriplet $\mapsto r_j(U^{-1})^j{}_i$, $U\in SU(3)$], algebra image in $\iota_3(\mathfrak u(3))$ [D2.12]. Strictly-new strong content: $\pi(\tau)\in\mathfrak{su}(3)$, $\pi(X)=X-\tfrac{\mathfrak T(X)}{3}I_3$ the trace projection [CT3.7, D3.2] — the $\mathfrak u(1)$ centre of $\mathfrak u(3)$ is EM's, spent at rung 2 [CT3.7, CT3.13; CT0.13 step $n{=}3$], and already carried by $G^{(2)}$'s phase factor.

\vspace{2pt}\noindent\rule{\textwidth}{0.4pt}\vspace{2pt}

\end{definition}

\begin{theorem}[{\texorpdfstring{FT.29 — Strong weight: factorization, no $\alpha_{\mathrm{strong}}$, observable sealing}{FT.29 - Strong weight: factorization, no , observable sealing}}]\label{tag:FT-29}
\textbf{Thm.}

%dsp-level:0
\[
(1)\quad\text{factorization}:\ G^{(3)}(e)=G^{(2)}(e)\otimes\tau(e)\ [\mathrm{Def\text{-}H.17}],
\]

%dsp-level:0
\[
\qquad \alpha_{\mathrm{EM}}^{(3)}=\alpha_{\mathrm{EM}}^{(2)}=\pi\ [\mathrm{EM.18},\ \mathrm{CT3.10}];
\]

%dsp-level:0
\[
(2)\quad\text{non-existence}:\ \neg\exists\ \text{per-edge abelian strong phase coefficient}\ (\text{an}\ ``\alpha_{\mathrm{strong}}\text{''})\ [\textbf{CT3.10a (F)}:\ \text{trace route};\ \text{eight-dimensional obstruction to a scalar coefficient read};\ \text{non-canonicity of}\ \mathfrak u(1)\subseteq\mathfrak{su}(3)\ \text{choices (explicit}\ SU(3)\ \text{witnesses)};\ \text{centre identification via CT3.7/CT3.13}];
\]

%dsp-level:0
\[
(3)\quad\text{observable content}:\ \text{under any pole read}\ [\mathrm{CT2.5}],
\]

%dsp-level:0
\[
\ \tau(e)\text{-factors survive only inside invariant contractions}\ \{\delta^i{}_j,\varepsilon_{ijk},\varepsilon^{ijk}\}\ [\mathrm{CT3.17}];
\]

%dsp-level:0
\[
\ \text{free indices are mask-annihilated}\ [\mathrm{CT3.18}];
\]

%dsp-level:0
\[
\ \text{isolated triplets have}\ \neg\text{pole and norm}\ 0\ [\mathrm{CT3.22},\ \mathrm{CT2.8}];
\]

%dsp-level:0
\[
\ \text{at}\ \Sigma\text{-level the strong summand is gated by}\ \Pi_{\mathcal C}(\Sigma;C)\in\{0,1\}\ [\mathrm{FT.19},\ \mathrm{FT.26}(2)].\ \text{Hence the effective strong weight of any materialized configuration}
\]

%dsp-level:0
\[
= (\text{its}\ G^{(2)}\ \text{factors})\times(\text{singlet-sealing factor}):\ \neg\exists\ \text{per-edge strong phase read}.
\]

\textbf{Req.} Def-H.17; D2.12, D3.1, D3.2, D3.5, D3.8, D3.9; CT0.13, CT0.14 (Part 1); CT2.5, CT2.8 (Part 2); CT3.2, CT3.4, CT3.7, CT3.8, CT3.9, CT3.10, CT3.10a, CT3.11, CT3.13, CT3.17, CT3.18, CT3.22 (Part 2); EM.17, EM.18, EM.20 (Part 2); FT.9, FT.10, FT.19, FT.26.

\textbf{Pf.}

%dsp-level:0
\[
(1):\ \mathrm{Def\text{-}H.17};
\]

%dsp-level:0
\[
\qquad \text{coefficient stability}:\ \mathrm{EM.18}\ [\mathrm{CT0.14}:\ \mathcal K^{(\mathrm{strong})}\perp\mathcal K^{(\mathrm{EM})}]\ \wedge\ \mathrm{CT3.10}\ [\mathfrak T(F_{\mathrm{strong}})=0:\ \text{independent argument}]\implies\alpha^{(3)}
\]

%dsp-level:0
\[
= \alpha^{(2)}
\]

%dsp-level:0
\[
= \pi\ [\mathrm{FT.9}/\mathrm{EM.20}]
\]

%dsp-level:0
\[
\text{consistency}:\ \mathrm{FT.10}\ \text{Pf already computes with this factorization}\ (G^{(3)}=G^{(2)}\otimes\text{label}:\ \text{per-edge weight unmodified})
\]

%dsp-level:0
\[
(2):\ \mathrm{CT3.10a}\ (\text{F})
\]

%dsp-level:0
\[
(3):\ \mathrm{CT3.17}:\ \forall\ \text{invariants}=\text{combinations of}\ \{\delta,\varepsilon,\bar\varepsilon\};
\]

%dsp-level:0
\[
\qquad \mathrm{CT3.18}:\ \mathcal P_M^{(3)}\big|_{\text{free index}}=0\implies\neg\text{pole}\ [\mathrm{CT2.5}]
\]

%dsp-level:0
\[
\mathrm{CT3.22}:\ \text{isolated triplet norm}\ 0\ [\mathrm{CT2.8}];
\]

%dsp-level:0
\[
\qquad \mathrm{FT.19}/\mathrm{FT.26}(2):\ \Pi_{\mathcal C}\ \text{gates the entire}\ \Sigma\text{-level summand};
\]

%dsp-level:0
\[
\qquad \text{sealed}:\ \tau\text{-content enters as singlet contraction amplitudes only}\ [\mathrm{CT3.17},\ \mathrm{D3.9}].\ \blacksquare
\]

\textbf{Rmk (stipulation-dependence).} The centre \textit{identification} route — the unique canonical centre is EM's — cites CT3.7's reading of D2.12's realization clause ("the abstract isomorphism class lies in $\iota_m(\mathfrak u(m))$"): the one stipulative layer in the chain [CT3.13: $\neg\exists$ ninth gluon mode]. The non-existence core (trace route; dimension obstruction; non-canonicity witnesses) is independent of that reading, and the theorem's force does not rest on the identification.

\textbf{Rmk (three routes to discreteness).} The corpus's quantization mechanisms are now three, on three different objects: \textbf{(R1) reality line} — CT0.9: $\sigma_e\in\{\pm1\}$, single-qubit $\wedge^2\mathcal H_q$; \textbf{(R2) materialization} — EW.18/FT.22–FT.25: $\kappa,k,\Theta_{\mathrm{int}}\in\mathbb Z$, pole single-valuedness on loops; \textbf{(R3) selection} — CT3.18/FT.19: admissibility $\in\{0,1\}$, mask annihilation of free colour. The strong edge weight is governed by R3 alone: it registers no per-edge number. (Cross-ref: FT.30's verdict column.)

\textbf{Phys.} The strong edge weight is the EM weight tensored with a colour-transport label — and the label carries no Abelian phase of its own (non-existence frozen: CT3.10a). Colour becomes observable only sealed into singlets: the per-edge strong "quantum" is the sealing itself. Confinement is not a force law added to $G(e)$; it is the observable content of $G(e)$'s strong factor.

\vspace{2pt}\noindent\rule{\textwidth}{0.4pt}\vspace{2pt}

\end{theorem}

\phantomsection
\section[Tier F.11 — Sector Quantum Audit and the Continuum Ledger]{\texorpdfstring{Tier F.11 — Sector Quantum Audit and the Continuum Ledger}{Tier F.11 - Sector Quantum Audit and the Continuum Ledger}}

\begin{definition}[{\texorpdfstring{Def-H.18 — Arena and spectrum (the two-locus vocabulary)}{Def-H.18 - Arena and spectrum (the two-locus vocabulary)}}]\label{tag:Def-H-18}
\textbf{Def.} The \textbf{arena} of a quantity $:=$ the configuration domain over which $Z(\mathcal R)$ sums [A8, C12: $\mathfrak C(\mathcal R)$] — off-shell; continuous wherever CT1.12/CT1.19a make it so. The \textbf{spectrum} of a quantity $:=$ its values on materialized configurations — those contributing to poles of the masked resolvent [D2.4, CT2.5; $\Sigma$-level: Def-H.16] — on-shell. \textbf{Two-locus principle} (vocabulary for an already-proven distinction): discreteness claims in this corpus are spectrum claims; continuity claims are arena claims. One and the same object may be continuous in the arena and quantized in the spectrum — CT2.5 (continuous operator families, discrete pole set) and FT.22–FT.25 (continuous $\mathcal E(e)$ arena [Def-H.14 Rmk], integer winding spectrum) are the two proven instances. Vocabulary only — no new premise.

\vspace{2pt}\noindent\rule{\textwidth}{0.4pt}\vspace{2pt}

\end{definition}

\begin{theorem}[{\texorpdfstring{FT.30 — Pair-channel sector quantum audit: $\bigl((5{+}3{+}1)+(3{+}3)+1\bigr)$}{FT.30 - Pair-channel sector quantum audit:}}]\label{tag:FT-30}
\textbf{Thm.} The $16$ pair-channel registers [CT0.6] tile as $(G_5+W_3+H_1)+(J^+_3+J^-_3)+D^+_1$ [CT0.12, CT1.1, CT1.2], and CT1.12's metric dictionary decomposes block-wise:

%dsp-level:0
\[
g_{\mu\nu}\ \leftrightarrow\ \underbrace{\hat h_{ij}\ (5)}_{G:\ \text{spin-2 shear}}\ \oplus\ \underbrace{\Sigma_{\mathrm{tr}}\ (1)}_{H:\ \text{volume mode}}\ \oplus\ \underbrace{a_{00}\ (1)}_{D^+:\ \Omega\ \text{conformal}}\ \oplus\ \underbrace{w_i\ (3)}_{W:\ \text{shift}};
\]

%dsp-level:0
\[
\qquad J^\pm\ (S\text{-odd}):\ 0\ \text{metric coordinates}\ [\mathrm{Def\text{-}H.14}:\ \mathcal E(e)\ \text{is}\ S\text{-even}].
\]

Independent-quanta audit:

%dsp-level:0
\[
\begin{array}{c|c|c|c}
\text{block} & \text{arena data} & \text{spectrum data} & \text{independent quantum?}\\ \hline
G\ (5) & \text{continuous shear}\ [\mathrm{CT1.12}] & \text{winding}\ \kappa,k\in\mathbb Z\ [\mathrm{FT.22}/\mathrm{FT.24}],\ \text{unit}\ \mathfrak h^2=4G_N\ln2\ [\mathrm{CT2.9}] & \textbf{yes — geometric}\ (\text{joint read, }\neg\text{per-block})\\
W\ (3) & \text{shift}\ w_i\ [\mathrm{CT1.12}];\ \text{Casimir fixed}:\ \vee^2\,j{=}1,\ \wedge^2\,j{=}0\ [\mathrm{CT0.7a}];\ \text{seed}\ j{=}\tfrac12\ [\mathrm{EW.32}] & \text{integer labels}\ k,\ \Theta_{\mathrm{int}}\ [\mathrm{FT.25}(3),\ \mathrm{FT.28}(2)] & \textbf{candidate — spin}\ (\text{extraction open, Gaps 11});\ \neg\text{self-charge}\ [\mathrm{FT.12}]\\
H\ (1) & h_0\ \text{seed background},\ \text{continuous P4 data}\ [\mathrm{EW.13},\ \mathrm{P4}];\ M^2\propto h_0^2\ [\mathrm{EW.33}] & \text{via mass read only}\ [\mathrm{FT.20}\ \text{row}\ H,\ \mathrm{Def-H.36}] & \text{no — continuous configuration datum}\\
J^+\ (3) & \text{seals}\ n_{\mathcal R}:0\to1\ [\mathrm{D1.2}];\ S\text{-odd}=\text{reality grading}\ [\mathrm{CT0.7}] & \neg\text{observable}\ [\mathrm{FT.20}] & \text{no — structural};\ \#J^+=\#J^-\ [\mathrm{D1.3}/\mathrm{D2.11}];\ \neg\text{charges}\ [\mathrm{FT.11}]\\
J^-\ (3) & \text{dissolution seals}\ [\mathrm{D1.2}] & \neg\text{observable}\ [\mathrm{FT.20}] & \text{no — structural, as}\ J^+\\
D^+\ (1) & \Lambda\ \text{uniform centre}\ [\mathrm{A8},\ \mathrm{D4.3}];\ w=-1\ [\mathrm{CT1.25}];\ \Omega\ \text{channel entry}\ [\mathrm{Def\text{-}H.2}] & \neg\text{observable independently}\ [\mathrm{FT.20}] & \text{no — P4 datum}\\ \hline
\multicolumn{4}{c}{\text{stack aspects (outside the}\ 16\ [\mathrm{CT0.12\ Rmk}]):}\\
\mathrm{EM} & \sigma_e\ \text{binary}\ [\mathrm{CT0.9}];\ q\in\{-1,0,+1\}\ [\mathrm{EM.9}/\mathrm{D3.4}] & \text{parity}\ (-1)^\kappa\ \text{on-shell}\ [\mathrm{FT.24}] & \textbf{yes — charge}\ (\text{R1 reality line})\\
\text{strong} & \tau(e)\ \text{defining-rep label}\ [\mathrm{Def\text{-}H.17}];\ \neg\exists\ \text{per-edge phase}\ [\mathrm{FT.29}(2)] & \text{singlet sealing}\ \{0,1\}\ [\mathrm{FT.19}/\mathrm{FT.29}(3)] & \textbf{yes — admissibility}\ (\text{R3 selection})\\
D^- & \text{locked content}\ [\mathrm{D4.2}] & \neg\text{observable}\ [\mathrm{FT.3},\ \mathrm{CT4.4}] & \text{no}
\end{array}
\]

\textbf{Verdict.} Independent quanta exist in exactly three read classes, one per discreteness route [FT.29 Rmk]: \textbf{(R1)} binary registration — reality line [CT0.9]: $\sigma_e$, $q$, $(-1)^{\aleph}$; \textbf{(R2)} integer winding — materialization [EW.18, FT.22–FT.25]: $\kappa$, $k$, $\Theta_{\mathrm{int}}$, unit $\mathfrak h^2=4G_N\ln2$ [CT2.9, CT1.29]; \textbf{(R3)} singlet admissibility — mask selection [CT3.18, FT.19]: $\{0,1\}$. No pair-channel block carries a block-local quantum: the geometric quantum is a \textit{joint} $S$-even read — $\widehat{ds}^{\,2}(e)=g_{\mu\nu}(e)\Delta x^\mu_e\Delta x^\nu_e/\mathfrak h^2$ [D1.6] mixes all ten metric coordinates across $G,H,D^+,W$ — and the $\bigl((5{+}3{+}1)+(3{+}3)+1\bigr)$ blocks are the arena decomposition of that read.

\textbf{Rmk (the emergence chain, assembled).} How spacetime "truly emerges", in six cited steps: (1) per-edge $S$-even data $(a_{00},q_{ij},w_i)\leftrightarrow$ metric dictionary $g_{\mu\nu}$ [CT1.12, D2.1 — block-wise per this Thm]; (2) coarse-grained ledger averages $\to$ smooth $(M,g)$ under hypotheses (i)–(iii) [CT2.6]; (3) line-element sums $\to$ path integrals, rel. error $O(\mathfrak h)$ [CT1.19]; (4) on-shell selection: $\oint_c ds^2\in2\mathfrak h^2\mathbb Z$ [FT.22], per-edge integers at the floor [FT.24], classical limit $=$ large-$k$ [FT.31(3)]; (5) stationary set of the ledger sum $=$ classical geometry [CT1.20], Einstein–Hilbert closure [CT1.21, CT2.7], calibration $\mathfrak h^2=4G_N\ln2$ [CT2.9]; (6) particles $=$ poles of the $\Sigma$-resolvent [Def-H.16, FT.26]. Two loci throughout [Def-H.18]: the arena continuous, the spectrum integer — no step adds a premise beyond the corpus's declared set (A §0.2a).

\textbf{Req.} CT0.6, CT0.7, CT0.7a, CT0.12, CT0.14 (Part 1); CT1.1, CT1.2, CT1.12, CT1.19, CT1.20, CT1.21, CT1.25, CT1.29, CT2.5, CT2.6, CT2.7, CT2.9 (Part 2); D1.2, D1.3, D1.6, D2.11, D4.2, D4.3; EW.13, EW.18, EW.32, EW.33 (Part 2); EM.9 (Part 2); CT3.18, CT3.22 (Part 2); CT4.4 (Part 2); FT.3, FT.9, FT.11, FT.12, FT.20, FT.22, FT.24, FT.25, FT.29; Def-H.2, Def-H.14, Def-H.17, Def-H.18; P4.

\textbf{Pf.}

%dsp-level:0
\[
\text{tiling}:\ \mathrm{CT0.12}\ \text{table}:\ 3{+}3{+}5{+}3{+}1{+}1=16=C_{\mathrm{cap}}\ [\mathrm{CT0.6}];
\]

%dsp-level:0
\[
\qquad \text{sector parities}:\ \mathrm{CT1.1}/\mathrm{CT1.2}:\ \mathcal E_{++}
\]

%dsp-level:0
\[
= \{E_0,\Sigma,\hat E^{ij}\}\ (1{+}1{+}5),
\]

%dsp-level:0
\[
\ \mathcal E_{+-}
\]

%dsp-level:0
\[
= \{W_i\}\ (3),
\]

%dsp-level:0
\[
\ \mathcal E_{-+}
\]

%dsp-level:0
\[
= \{J_i\}\ (3),
\]

%dsp-level:0
\[
\ \mathcal E_{--}\qquad = \{K_i\}\ (3)
\]

%dsp-level:0
\[
\text{dictionary blocks}:\ a_{00}\sim E_0
\]

%dsp-level:0
\[
= D^+\ \text{block};
\]

%dsp-level:0
\[
\ \ \Sigma_{\mathrm{tr}}\sim\Sigma
\]

%dsp-level:0
\[
= H\ \text{block};
\]

%dsp-level:0
\[
\ \ \hat h_{ij}\sim\hat E^{ij}
\]

%dsp-level:0
\[
= G\ \text{block};
\]

%dsp-level:0
\[
\ \ w_i\sim W_i\qquad = W\ \text{block}
\]

%dsp-level:0
\[
\quad [\mathrm{CT1.12}\ \text{table}\ +\ \mathrm{CT0.12}\ \text{block table};\ \mathrm{D2.1}\ \text{coefficients}]
\]

%dsp-level:0
\[
J^\pm:\ S\text{-odd}\ [\mathrm{CT1.1}]\implies\notin\mathcal E(e)=\mathcal E_{++}\oplus\mathcal E_{+-}\ [\mathrm{Def\text{-}H.14}]\implies0\ \text{metric coordinates}
\]

%dsp-level:0
\[
\text{table rows}:\ \text{as cited (each cell one frozen-corpus or proven-H tag)}
\]

%dsp-level:0
\[
\neg\text{block-locality}:\ \widehat{ds}^{\,2}(e)
\]

%dsp-level:0
\[
= g_{\mu\nu}\Delta x^\mu\Delta x^\nu/\mathfrak h^2\ [\mathrm{D1.6}]\ \text{with}\ g_{\mu\nu}\ \text{spanning all four blocks}\ [\mathrm{CT1.12}\ \text{dictionary}]\implies\kappa\ \text{is a joint read}.\ \blacksquare
\]

\textbf{Phys.} Read the 16-matrix as the ledger's periodic table: nine $S$-even force directions build the metric arena (shear $+$ volume $+$ conformal $+$ shift), six $S$-odd directions make and unmake the reality of edges, one centre rescales uniformly — and not one of those blocks holds a quantum by itself. The quanta live in the reads: a sign at registration, a winding at materialization, a sealing at selection. That is the precise sense in which spacetime "truly emerges": an $S$-even arena on which materialization writes the integers, while the $S$-odd sectors run the registration that makes edges — and hence geometry — exist at all.

\vspace{2pt}\noindent\rule{\textwidth}{0.4pt}\vspace{2pt}

\end{theorem}

\begin{theorem}[{\texorpdfstring{FT.31 — The continuum limit survives on-shell quantization}{FT.31 - The continuum limit survives on-shell quantization}}]\label{tag:FT-31}
\textbf{Thm.} Let $\{c_{\mathfrak h}\}$ be a family of materialized chains [FT.22] with $L=\mathfrak h\,n_{\mathcal R}$ fixed, $\mathfrak h\to0$ [CT1.19's regime]. Then:

%dsp-level:0
\[
(1)\quad\text{per edge}:\ \ \kappa_e
\]

%dsp-level:0
\[
= g_{\mu\nu}\dot x^\mu\dot x^\nu(\lambda_e)+O(\mathfrak h)\ \wedge\ \kappa_e\in\mathbb Z\ \text{on-shell}\ [\mathrm{FT.24}\ \text{scope}]:\ \text{integers sampling a smooth function with}\ O(\mathfrak h)\ \text{slack};
\]

%dsp-level:0
\[
(2)\quad\text{sums}:\ \ \sum_e\kappa_e\mathfrak h^2
\]

%dsp-level:0
\[
= \sum_e ds^2(e)\ \text{identically}\ [\mathrm{D1.6}]\implies\mathrm{CT1.19}\ \text{applies verbatim}:\ \text{abs. error}\ O(\mathfrak hL),
\]

%dsp-level:0
\[
\ \text{rel.}\ O(\mathfrak h):\ \text{integrality adds no error term — selection acts on family membership, not on sums};
\]

%dsp-level:0
\[
(3)\quad\text{selection in the limit}:\ \ \oint_c ds^2\in2\mathfrak h^2\mathbb Z\ [\mathrm{FT.22}\ \mathrm{Phys}];
\]

%dsp-level:0
\[
\qquad k
\]

%dsp-level:0
\[
= \frac{1}{2\mathfrak h^2}\oint_c ds^2\to\infty\ \text{at fixed classical loop}:\ \text{the lattice is classically invisible, quantally exact};
\]

%dsp-level:0
\[
(4)\quad\text{floor relating equation (Gaps 2, minimal-edge class)}:\ \ \kappa_e
\]

%dsp-level:0
\[
\equiv 1\ [\mathrm{D1.1},\ \text{Comm-min}]\implies\Theta(c)
\]

%dsp-level:0
\[
= \delta h(c)\ \text{and}\ \ \frac{S_{\mathrm{reg}}(c)}{\oint_c ds^2}
\]

%dsp-level:0
\[
= \frac{\delta h\ln2}{\delta h\,\mathfrak h^2}
\]

%dsp-level:0
\[
= \frac{\ln2}{\mathfrak h^2}
\]

%dsp-level:0
\[
= \frac{S_{\min}}{A_{\min}}
\]

%dsp-level:0
\[
\qquad [\mathrm{CT0.9},\ \mathrm{C9},\ \mathrm{CT1.29}].
\]

\textbf{Req.} FT.22, FT.24; CT1.19, CT1.19a (Part 2); CT2.6 hypotheses (i)–(iii) (inherited); D1.1, D1.6, D1.7; CT0.9 (Part 1); C9; CT1.29; Comm-min (via D1.1/EM.20).

\textbf{Pf.}

%dsp-level:0
\[
(1):\ \mathrm{CT1.19}\ \text{Pf}:\ \Delta x^\mu_e=\mathfrak h\,\dot x^\mu(\lambda_e)+O(\mathfrak h^2)\implies\widehat{ds}^{\,2}(e)=g_{\mu\nu}\dot x^\mu\dot x^\nu(\lambda_e)+O(\mathfrak h);
\]

%dsp-level:0
\[
\qquad \text{on-shell}\ \widehat{ds}^{\,2}(e)=\kappa_e\in\mathbb Z\ [\mathrm{FT.24}]
\]

%dsp-level:0
\[
(2):\ \kappa_e\mathfrak h^2
\]

%dsp-level:0
\[
= ds^2(e)\ \text{by definition}\ [\mathrm{D1.6}:\ \widehat{ds}^{\,2}=ds^2/\mathfrak h^2]\implies\text{the sums are the same sums}:\ \mathrm{CT1.19}\ \text{unmodified}
\]

%dsp-level:0
\[
(3):\ \Theta(c)\mathfrak h^2
\]

%dsp-level:0
\[
= \sum_e ds^2(e)\to\oint_c ds^2\ [\mathrm{CT1.19}]\ \wedge\ \Theta(c)
\]

%dsp-level:0
\[
= 2k\ [\mathrm{FT.22}]\implies\oint_c ds^2
\]

%dsp-level:0
\[
\in 2\mathfrak h^2\mathbb Z\ \text{at leading order};
\]

%dsp-level:0
\[
\qquad \oint\ \text{fixed},\ \mathfrak h\to0:\ k\to\infty
\]

%dsp-level:0
\[
(4):\ \kappa_e\equiv1\implies\Theta=\sum_e1=\delta h(c)\ [\mathrm{D1.7}];
\]

%dsp-level:0
\[
\qquad S_{\mathrm{reg}}(c)=\delta h\cdot\ln2\ [\mathrm{CT0.9}/\mathrm{C9}:\ 1\ \text{real edge}=1\ \text{bit}];
\]

%dsp-level:0
\[
\qquad \oint_c ds^2=\Theta\mathfrak h^2=\delta h\,\mathfrak h^2
\]

%dsp-level:0
\[
\frac{S_{\mathrm{reg}}}{\oint ds^2}=\frac{\delta h\ln2}{\delta h\,\mathfrak h^2}=\frac{\ln2}{\mathfrak h^2};
\]

%dsp-level:0
\[
\qquad \mathrm{CT1.29}:\ S_{\min}=\ln2\ (1\ \text{edge}),\ A_{\min}=\mathfrak h^2\ (1\ \text{hop area})\implies\frac{S_{\min}}{A_{\min}}=\frac{\ln2}{\mathfrak h^2}:\ \text{identical}.\ \blacksquare
\]

\textbf{Rmk (decoupling of the continuum hypotheses).} The continuum-limit layer — CT2.6's hypotheses (i)–(iii), CT1.19's convergence — is now fully decoupled from line-element quantization: the hypotheses concern smoothness of $g$ and cell refinement (\textbf{arena}); quantization concerns materialization (\textbf{spectrum}) [Def-H.18]. Neither supports nor threatens the other, and no result of this document makes the continuum limit depend on a discreteness postulate.

\textbf{Phys.} Quantized summands do not break the Riemann machine: the integers $\kappa_e$ track the smooth sample $g_{\mu\nu}\dot x^\mu\dot x^\nu$ to $O(\mathfrak h)$ per edge, so CT1.19's convergence passes through untouched and the on-shell condition becomes, in the continuum, the Bohr–Sommerfeld rule $\oint ds^2\in2\mathfrak h^2\mathbb Z$. At the minimal-edge floor the information/area relating equation $\ln2/\mathfrak h^2$ is an identity — one bit per hop area [CT1.29] is the same statement as one winding quantum per minimal edge. Classical spacetime is the large-winding limit of an integer ledger.

\vspace{2pt}\noindent\rule{\textwidth}{0.4pt}\vspace{2pt}

\end{theorem}

\phantomsection
\section[Tier F.12 — The Cycle Floor Family, the Defect Bridge, and…]{\texorpdfstring{Tier F.12 — The Cycle Floor Family, the Defect Bridge, and Composite Floors}{Tier F.12 - The Cycle Floor Family, the Defect Bridge, and Composite Floors}}

\begin{definition}[{\texorpdfstring{Def-H.19 — Stationary, stable, and unstable jumps (vocabulary)}{Def-H.19 - Stationary, stable, and unstable jumps (vocabulary)}}]\label{tag:Def-H-19}
\textbf{Def.} A \textbf{jump} $:=$ a transit path between two termini in the emergent chart [D0.6, CT2.6] — program vocabulary, no new structure. A \textbf{stationary jump}: $\delta\!\int_\gamma ds^2=0$ among fixed-endpoint variations — an affinely parametrized geodesic of $g_{\mu\nu}$ [CT1.20's stationary-phase read; Euler–Lagrange for $\int g_{\mu\nu}\dot x^\mu\dot x^\nu d\lambda$ — math]. A \textbf{stable jump}: a stationary jump that is a local minimum of $\int ds^2$ on its interval (no conjugate points [math: Morse/Synge]); \textbf{unstable}: stationary but saddle or maximal. Lorentzian read [CT1.12: $\operatorname{sig}=(-,+,+,+)$]: on timelike paths, minimizing $\int ds^2$ is \textit{maximizing} proper time — the classical stable trajectory; on spacelike paths, shortest length. Program vocabulary (constitutive; Morse-theoretic [math]); the "relatively stable" idiom of program notes means: measured against the other stationary paths with the same endpoints — an internal comparison, no external scale. \textbf{No coupling to materialization is asserted.} CT1.20's stationary phase weights minima and saddles alike (every $\delta\Phi=0$ configuration dominates non-stationary ones); whether the pole condition [CT2.5] or the defect criterion [FT.33] \textit{selects} the stable subclass is a theorem-shaped open question, not a definition. \textbf{(refinement [Def-H.24(5), FT.66]):} for degenerate stationary families (isometry orbits) the vocabulary is Morse–Bott — "stable" $:=\operatorname{ind}_B=0$ on a clean critical manifold, "unstable" $:=\operatorname{ind}_B\ge1$; and the coupling question is answered: no index criterion operates on closed chains within the weight inventory [FT.66(2)–(3)], and absolute index-0 selection is \textit{inadmissible} for bound chains on Synge-class slices [FT.66(4)]; the open-channel residue is tracked under Open Gaps 13.

\vspace{2pt}\noindent\rule{\textwidth}{0.4pt}\vspace{2pt}

\end{definition}

\begin{theorem}[{\texorpdfstring{FT.33 — The defect bridge: pole-side mass candidate in well-posed ratio form}{FT.33 - The defect bridge: pole-side mass candidate in well-posed ratio form}}]\label{tag:FT-33}
\textbf{Thm.} Under Def-H.37(1) (interface transfer: the floor fermion's pole displacement takes EW.33's leading-order defect form for the seed $\mathsf v$ [EW.32] weak-transported around its materialized bound chain), for a materialized \textbf{spacelike} single-contact cycle floor $\Sigma_L$ with $\kappa_e>0$:

%dsp-level:0
\[
(1)\quad m_{\mathrm{pole}}^2(\Sigma_L)=\hat\mu^2\,\ell_c^2\;G(c),\qquad\ell_c=\mathfrak h\,\Lambda_c,\quad\Lambda_c:=\sum_{e\in c}\sqrt{\kappa_e}\ \ (\text{dimensionless});
\]

%dsp-level:0
\[
G(c)=\sum_{A,B}c_Ac_B\operatorname{Re}\langle\mathcal G_A\mathsf v,\mathcal G_B\mathsf v\rangle\ \ [\mathrm{EW.33};\ \mathrm{EW.20}\ \text{Gram data}];
\]

%dsp-level:0
\[
\qquad \hat\mu^2>0:\ \text{floor-independent residue packaging}\ \varkappa^2\beta^2h_0^2\ \ [\mathrm{EW.33}\ \text{with}\ \ell^2\ \text{factored out}];
\]

%dsp-level:0
\[
(2)\quad\frac{m_{\mathrm{pole}}^2}{m_{\mathrm{hid}}^2}=\frac{\hat\mu^2\,\mathfrak h^4\,G(c)}{(\ln2)^2}\cdot\Bigl(\frac{\Lambda_c}{L}\Bigr)^{2}\ \in\ \mathfrak q^0
\]

%dsp-level:0
\[
\qquad [\text{well-posed};\ m_{\mathrm{hid}}^2=L^2(\ln2)^2/\mathfrak h^2\ \ \mathrm{FT.32}(1)];
\]

%dsp-level:0
\[
(3)\quad\textbf{length-universality}:\ \text{at fixed}\ \kappa\text{-profile shape and Gram data the ratio is independent of}\ L;
\]

%dsp-level:0
\[
\quad\text{uniform}\ \kappa:\ \text{ratio}
\]

%dsp-level:0
\[
= \frac{\hat\mu^2\mathfrak h^4G(c)}{(\ln2)^2}\,
\]

%dsp-level:0
\[
\kappa\ —\ \text{linear in the edge quantum, floor-free};
\]

%dsp-level:0
\[
(4)\quad\text{Cauchy–Schwarz}:\ \Lambda_c^2
\]

%dsp-level:0
\[
\le L\,
\]

%dsp-level:0
\[
\Theta(c),
\]

%dsp-level:0
\[
\ \text{equality}\iff\kappa_e\ \text{uniform}\ [\text{math}]\implies\text{at fixed}\ (\Theta,L):\ \text{uniform floors maximize the bridge ratio};
\]

%dsp-level:0
\[
(5)
\]

%dsp-level:0
\[
\quad h_0
\]

%dsp-level:0
\[
= 0\implies\hat\mu^2
\]

%dsp-level:0
\[
= 0
\]

%dsp-level:0
\[
\implies m_{\mathrm{pole}}^2
\]

%dsp-level:0
\[
= 0\ \ [\mathrm{EW.33}]\ \wedge\ m_{\mathrm{hid}}
\]

%dsp-level:0
\[
\ne 0\ \ [\mathrm{FT.32}(1)]:\ \text{the two reads decouple without the scalar background};
\]

%dsp-level:0
\[
(6)\quad\textbf{kernel coherence}:\ c\in\ker M^2
\]

%dsp-level:0
\[
= \operatorname{span}\{Q\}\ [\mathrm{EW.22}]
\]

%dsp-level:0
\[
\implies G(c)
\]

%dsp-level:0
\[
= 0
\]

%dsp-level:0
\[
\implies m_{\mathrm{pole}}^2
\]

%dsp-level:0
\[
= 0\ [\mathrm{EW.33}:\ \text{Gram PSD}]\ —\ \text{charge-aligned (photon-direction) transport is defect-massless at every floor, matching}\ m_\gamma^2
\]

%dsp-level:0
\[
= 0\ [\mathrm{EM.6}].
\]

No identification $m_{\mathrm{pole}}=m_{\mathrm{hid}}$ is asserted; the ratio is the bridge datum.

\textbf{Req.} Def-H.37(1) (forward pointer — Tier F.17); EW.20, EW.22, EW.32, EW.33 (Part 2); EM.6 (Part 2); FT.23, FT.24, FT.32; Def-H.2, Def-H.9(2); D1.6; CT2.9 (Part 2); D2.9 (units); Cauchy–Schwarz [math].

\textbf{Pf.}

%dsp-level:0
\[
(1):\ \mathrm{EW.33}:\ \mathcal U=\exp\bigl(i\oint\textstyle\sum_Ac_A\mathcal G_A\,d\ell\bigr)=1+i\,\ell_c\sum_Ac_A\mathcal G_A+O(c^2);
\]

%dsp-level:0
\[
\qquad \mathrm{CT2.5}\ \text{read}\ [\mathrm{EW.33}]:\ m^2=\|(\mathcal U-1)\mathsf v\|^2=\ell_c^2\,G(c)+O(c^3)
\]

%dsp-level:0
\[
\ell_c
\]

%dsp-level:0
\[
= \textstyle\sum_e\sqrt{\kappa_e}\,
\]

%dsp-level:0
\[
\mathfrak h\ [\mathrm{Def\text{-}H.2}\ \text{piecewise-straight length};\ \mathrm{FT.24}(2):\ ds^2(e)=\kappa_e\mathfrak h^2,\ \text{spacelike}\ \kappa_e>0];
\]

%dsp-level:0
\[
\qquad \hat\mu^2:\ \mathrm{EW.33}\ \text{packaging},\ \ell^2\ \text{factored out}\ [\mathrm{Def\text{-}H.37(1)}]
\]

%dsp-level:0
\[
(2):\ \frac{m_{\mathrm{pole}}^2}{m_{\mathrm{hid}}^2}=\frac{\hat\mu^2\mathfrak h^2\Lambda_c^2\,G}{L^2(\ln2)^2/\mathfrak h^2}=\frac{\hat\mu^2\mathfrak h^4G}{(\ln2)^2}\Bigl(\frac{\Lambda_c}{L}\Bigr)^2;
\]

%dsp-level:0
\[
\qquad [\hat\mu^2]=\frac{[m^2]}{[\ell^2]}=\mathfrak q\cdot\mathfrak q=\mathfrak q^2\ [\mathrm{EW.33}\ \text{packaging};\ \mathrm{D2.9}]\implies[\hat\mu^2\mathfrak h^4]=\mathfrak q^0\ \checkmark
\]

%dsp-level:0
\[
(3):\ \text{uniform}\ \kappa:\ \Lambda_c=L\sqrt\kappa\implies(\Lambda_c/L)^2=\kappa\ [\text{exact squaring},\ \neg\text{numerical roots}];\quad L\ \text{cancels}
\]

%dsp-level:0
\[
(4):\ \bigl(\sum_e\sqrt{\kappa_e}\bigr)^2
\]

%dsp-level:0
\[
\le L\sum_e\kappa_e
\]

%dsp-level:0
\[
= L\,
\]

%dsp-level:0
\[
\Theta\ [\text{Cauchy–Schwarz};\ \text{equality}\iff\sqrt{\kappa_e}\ \text{proportional to}\ 1\iff\kappa\ \text{uniform}]
\]

%dsp-level:0
\[
(5):\ \mathrm{EW.33}:\ \mu_\star^2=c_W\varkappa^2\beta^2h_0^2,\ \ h_0=0\implies M^2\equiv0;\qquad\mathrm{FT.32}(1):\ m_{\mathrm{hid}}=L\ln2/\mathfrak h\ \neg\text{references}\ h_0
\]

%dsp-level:0
\[
(6):\ G(c)=c^{\mathsf T}\bigl(M^2/\mu_\star^2\bigr)c\ [\mathrm{EW.20}/\mathrm{EW.33}];
\]

%dsp-level:0
\[
\quad M^2\ \text{PSD}\ [\mathrm{EW.33}:\ \text{Gram}]\implies G(c)=0\iff c\in\ker M^2=\operatorname{span}\{Q\}\ [\mathrm{EW.22}];
\]

%dsp-level:0
\[
\qquad Q\text{-aligned}
\]

%dsp-level:0
\[
= \text{photon direction}\ [\mathrm{D3.10}:\ A=\sin\theta_WW_3+\cos\theta_WY;\ \mathrm{EW.24}]
\]

%dsp-level:0
\[
\implies m_{\mathrm{pole}}^2
\]

%dsp-level:0
\[
= 0\ [\mathrm{EM.6}\ \text{consistency}].\ \blacksquare
\]

\textbf{Rmk (wrapping-chain consistency).} On rung-3 wrapping chains, Def-H.37(2)s formula \textit{is} EW.33 verbatim — the transfer adds nothing there; FT.33's content is the ratio against $m_{\mathrm{hid}}$ and its floor-universality. The Gram quadratic form $G(c)$ remains configuration data (the chain's weak-transport profile $c_A$); floor-independence in clause (3) is at fixed $G$ and fixed profile shape. \textbf{Genericity upgrade:} on-shell rigidity [FT.38(2), forward] makes uniform-$\kappa$ — hence clauses (3)–(4) — continuum-generic within the integrality classes, not a specialization.

\textbf{Phys.} If the fermion's pole mass is the seed's failure to return to itself around the bound chain, then the informational and pole reads of mass are bridged by a single dimensionless ratio: a fixed P4/Gram factor times the squared mean root of the edge quanta. Cycle length drops out completely — the ladder of hidden bits and the ladder of geometric length grow at exactly the same rate — and the ratio is maximized by uniform floors, vanishes with the scalar background, and vanishes identically on charge-aligned transport, reproducing the massless photon at every floor. The integers enter the bridge as $\overline{\sqrt\kappa}^{\,2}$: the one place where the discrete ledger meets irrational lengths.

\vspace{2pt}\noindent\rule{\textwidth}{0.4pt}\vspace{2pt}

\end{theorem}

\begin{theorem}[{\texorpdfstring{FT.34 — Round-trip quantization between floors; integrality propagation; sealing partition rule}{FT.34 - Round-trip quantization between floors; integrality propagation; sealing partition rule}}]\label{tag:FT-34}
\textbf{Thm.} Let $\Sigma^{(1)},\Sigma^{(2)}$ carry materialized 2-cycles $c_1,c_2$ [$\Theta(c_i)=2\kappa_i$, $\kappa_i\in\mathbb Z$: FT.24], and let $P,P'$ be connector chains with $\Gamma:=c_1\cup P\cup c_2\cup P'$ a valid closed chain [D1.4: concatenation, Janus balance], edge-disjoint union. If $\Gamma$ materializes:

%dsp-level:0
\[
(1)\quad\Theta(P)+\Theta(P')=\Theta(\Gamma)-2\kappa_1-2\kappa_2\in2\mathbb Z\qquad[\mathrm{FT.22},\ \mathrm{D1.7}\ \text{additivity}];
\]

%dsp-level:0
\[
(2)\quad P'=\bar P\ (\text{same bidirectional channels reversed}\ [\mathrm{FT.23}\ \text{per edge}]):\quad\Theta(P)\in\mathbb Z;
\]

%dsp-level:0
\[
\qquad d_{12}:=\Theta(P)=\textbf{geometric distance (integer)};
\]

%dsp-level:0
\[
(3)\quad\textbf{additive winding law}:\ k(\Gamma)=\frac{\Theta(\Gamma)}{2}=\kappa_1+\kappa_2+d_{12}=k(c_1)+k(c_2)+d_{12}
\]

%dsp-level:0
\[
\qquad [\mathrm{FT.25}(3):\ k=\kappa\ \text{at 2-cycle floors}];
\]

%dsp-level:0
\[
(4)\quad\textbf{round-trip integrality (propagation lemma)}:\ \text{let}\ e,
\]

%dsp-level:0
\[
\bar e\ \text{be both real},
\]

%dsp-level:0
\[
\ \Gamma\ \text{materialized closed with edge multiset}\ \{e,\bar e\}\sqcup E_0,
\]

%dsp-level:0
\[
\ E_0\ \text{decomposing edge-disjointly into materialized closed chains}:
\]

%dsp-level:0
\[
\quad 2\kappa_e=\Theta(\Gamma)-\Theta(E_0)\in2\mathbb Z\implies\kappa_e\in\mathbb Z;
\]

%dsp-level:0
\[
(5)\quad\textbf{sealing partition rule}\ [\textbf{CT3.17a(1) (F)}]:\ \text{the group-size structure of sealed colour-attachment sets, and the non-sealing of a lone attachment}.
\]

\textbf{Req.} FT.22, FT.23, FT.24, FT.25; D1.4, D1.7; Def-H.10; CT3.17, CT3.17a (Part 2); FT.19, FT.27; elementary number theory [math].

\textbf{Pf.}

%dsp-level:0
\[
(1):\ \mathrm{D1.7}:\ \Theta(\Gamma)=\Theta(c_1)+\Theta(P)+\Theta(c_2)+\Theta(P')\ (\text{edge-disjoint});
\]

%dsp-level:0
\[
\qquad \mathrm{FT.22}:\ \Theta(\Gamma)\in2\mathbb Z;\quad\mathrm{FT.24}(1):\ \Theta(c_i)=2\kappa_i\in2\mathbb Z
\]

%dsp-level:0
\[
\implies\Theta(P)+\Theta(P')=\Theta(\Gamma)-2(\kappa_1+\kappa_2)\in2\mathbb Z-2\mathbb Z=2\mathbb Z
\]

%dsp-level:0
\[
(2):\ P'=\bar P:\ \Theta(P')=\Theta(P)\ [\mathrm{FT.23}\ \text{applied per edge of}\ P]\implies2\Theta(P)\in2\mathbb Z\implies\Theta(P)\in\mathbb Z;
\]

%dsp-level:0
\[
\qquad d_{12}:=\Theta(P)
\]

%dsp-level:0
\[
(3):\ k(\Gamma)=\Theta(\Gamma)/2=(2\kappa_1+2\kappa_2+2d_{12})/2=\kappa_1+\kappa_2+d_{12}\in\mathbb Z
\]

%dsp-level:0
\[
(4):\ \Theta(\Gamma)=\kappa_e+\kappa_{\bar e}+\Theta(E_0)=2\kappa_e+\Theta(E_0)\ [\mathrm{FT.23}:\ \kappa_{\bar e}=\kappa_e;\ \mathrm{D1.7}];
\]

%dsp-level:0
\[
\qquad \Theta(\Gamma)\in2\mathbb Z\ [\mathrm{FT.22}];\quad\Theta(E_0)=\sum_i\Theta(c_i)\in2\mathbb Z\ [\mathrm{FT.22}\ \text{per chain}]
\]

%dsp-level:0
\[
\implies2\kappa_e\in2\mathbb Z\implies\kappa_e\in\mathbb Z
\]

%dsp-level:0
\[
(5):\ \mathrm{CT3.17a}\ (\text{F}).\ \blacksquare
\]

\textbf{Rmk (scope of the distance).} Clause (2)'s integer distance requires the \textit{bidirectional} round-trip closure: a single-orientation connector gives only the symmetrized constraint $\Theta(P)+\Theta(P')\in2\mathbb Z$ — the one-way open path remains unconstrained (no winding without closure; the handoff-level claim "$\int_\gamma ds^2=2k\mathfrak h^2$ for geodesics between jumps" holds only for the closed round-trip, where it reads $\oint_\Gamma ds^2=2k(\Gamma)\mathfrak h^2$). Refinement: the quantization \textit{law} is refinement-exact [FT.22 at every level]; the integer \textit{value} scales as $\mathfrak h\to0$ — the classical limit is large $k$ [FT.31(3)].

\textbf{Phys.} Two floor fermions joined by a bidirectional channel measure their separation in integers: the round-trip's line-element sum is even-quantized, the one-way accumulation is an integer, and the composite loop's winding number is the \textit{sum} of the two floor windings and the distance between them — an additive composition law for the geometric ledger. The same clause propagates integrality: any bidirectional channel that closes through materialized floors into a materialized loop inherits $\kappa\in\mathbb Z$ — the on-shell sublattice grows outward from the 2-cycle floors. (The colour sealing partition rule: CT3.17a(1) (Part 2).)

\vspace{2pt}\noindent\rule{\textwidth}{0.4pt}\vspace{2pt}

\end{theorem}

\phantomsection
\section[Tier F.13 — Gap-Closure Pass: Flux Ledger, Relating…]{\texorpdfstring{Tier F.13 — Gap-Closure Pass: Flux Ledger, Relating Equation, Rigidity, Confinement Converse, Core Charge, Aspect Ledger}{Tier F.13 - Gap-Closure Pass: Flux Ledger, Relating Equation, Rigidity, Confinement Converse, Core Charge, Aspect Ledger}}

\begin{theorem}[{\texorpdfstring{FT.36 — Flux ledger completion}{FT.36 - Flux ledger completion}}]\label{tag:FT-36}
\textbf{Thm.} $(1)$ \textbf{Saturation locus.} Jaw/threshold arithmetic — where the temporal jaw binds, and the threshold $=$ the channel constant: CT0.6a (Part 1). In FT.1's geodesic slowly-varying gauge ($\delta h=\Delta\ell/(\Omega\mathfrak h)$, $\delta\tau=\Delta t/\Omega$, $\Omega>0$ [CT1.12]) the locus reads

%dsp-level:0
\[
\frac{\Delta\ell(\gamma)}{\Delta t(\gamma)}=4\,\mathfrak h\qquad(\text{channel speed}=4\ \text{hop lengths per tick}),
\]

with $\Omega$ cancelling: \textbf{the saturation scale is a conformal invariant} [FT.6 pattern]; only the saturated value is $\Omega$-modulated ($\mathcal F=4\ln2$, $\mathcal F_{\mathrm{phys}}=4\ln2/\Omega$ [FT.5, Def-H.5]). $(2)$ \textbf{Attainment (FT.36(2) discharged).} Whenever Def-H.3's transit domain is nonempty, its sup is attained by an explicit photon-content schedule: partition $\gamma$'s edges into per-tick groups of $\le4$; run transit (photon) content [EM.5, EM.6: trivial chains $\subseteq V_2$; D1.5: open propagating $=$ transit] on $\mathsf B(\gamma)/\ln2$ scheduled hops [FT.1's contracted notation] — per-tick group capacity: CT0.6a's sharpness clause (Part 1) — $\mathsf E_{\mathrm{tr}}=\delta h'\ln2$ [EM.7] $\Rightarrow\sup=\mathsf B(\gamma)$ [FT.1's bound met with equality]. On channels whose edges are bound-locked in every configuration the transit domain is empty and $I=0$ (attainment vacuous; bound trivially valid). $(3)$ Hence FT.1/FT.3/FT.4/FT.5's bounds are sup-attained equalities on exhibited configurations wherever transit content is admissible.

\textbf{Req.} FT.1, FT.5, FT.6; Def-H.2, Def-H.3, Def-H.5; C9; CT0.6a (Part 1); CT1.12 (Part 2); D0.10, D1.5; EM.5, EM.6, EM.7 (Part 2).

\textbf{Pf.}

%dsp-level:0
\[
(1):\ \text{jaw/threshold arithmetic}:\ \mathrm{CT0.6a}\ (\text{B});
\]

%dsp-level:0
\[
\qquad \frac{\delta h}{\delta\tau}=\frac{\Delta\ell/(\Omega\mathfrak h)}{\Delta t/\Omega}=\frac{\Delta\ell}{\mathfrak h\,\Delta t}\ [\Omega>0\ \text{cancels};\ \mathrm{CT1.12}:\ \Omega=e^{\alpha_\Omega a_{00}}>0]
\]

%dsp-level:0
\[
\text{locus}:\ \delta h/\delta\tau=4\iff\Delta\ell/\Delta t=4\mathfrak h
\]

%dsp-level:0
\[
\quad (\text{dimension check in hop exponents}:\ [\Delta\ell]=\mathfrak h^1,\ [\Delta t]=\mathfrak h^0\ [\mathrm{D1.8}:\ x^0=\tau_n\ \text{tick count}]\ \checkmark)
\]

%dsp-level:0
\[
(2):\ \gamma'\subseteq\gamma\ \text{with}\ \delta h'
\]

%dsp-level:0
\[
= \min(\delta h,4\,\delta\tau)\ \text{hops carrying photon content is a valid transit configuration}\ [\mathrm{D0.10}:\ \mathfrak C(\mathcal R)\ \text{admits transit assignments};\ \mathrm{D1.5};\ \mathrm{EM.5}/\mathrm{EM.6}]:\ \mathsf E_{\mathrm{tr}}(\gamma')
\]

%dsp-level:0
\[
= \delta h'\ln2\ [\mathrm{EM.7}]
\]

%dsp-level:0
\[
\text{schedule exactness (per-tick groups at channel capacity)}:\ \mathrm{CT0.6a}\ (\text{B});
\]

%dsp-level:0
\[
\qquad \sup\ge\min(\delta h,4\,\delta\tau)\ln2\ \wedge\ \mathrm{FT.1}:\ \sup\le\ \text{same}\implies\text{equality, attained}
\]

%dsp-level:0
\[
(3):\ \text{from}\ (2)\ \text{at each bound's locus}.\ \blacksquare
\]

\textbf{Phys.} The flux ledger closes: the crossover where knowledge switches from space-limited to time-limited is the conformally invariant channel speed — four hop lengths per tick, the same four that is the channel's bit capacity — and the bound is \textit{reached}, not merely valid, by photon schedules running the ledger at capacity. What FT.36(2) assumed, FT.36 exhibits.

\vspace{2pt}\noindent\rule{\textwidth}{0.4pt}\vspace{2pt}

\end{theorem}

\begin{theorem}[{\texorpdfstring{FT.37 — The relating equation, complete}{FT.37 - The relating equation, complete}}]\label{tag:FT-37}
\textbf{Thm.} For a materialized closed chain $c$ [FT.22], with $S_{\mathrm{reg}}(c)=\delta h(c)\,\ln2$ [CT0.9, C9] and $A_{\mathrm{loop}}(c):=\oint_c ds^2=\Theta(c)\,\mathfrak h^2=2k\,\mathfrak h^2$ [D1.6, FT.22], and mean edge quantum $\bar\kappa(c):=\Theta(c)/\delta h(c)$:

%dsp-level:0
\[
\frac{S_{\mathrm{reg}}(c)}{A_{\mathrm{loop}}(c)}=\frac{\ln2}{\mathfrak h^2}\cdot\frac{1}{\bar\kappa(c)}=\frac{\delta h(c)\,\ln2}{2k\,\mathfrak h^2};
\]

%dsp-level:0
\[
\qquad \frac{S_{\mathrm{reg}}}{A_{\mathrm{loop}}}=\frac{S_{\min}}{A_{\min}}=\frac{\ln2}{\mathfrak h^2}\ [\mathrm{CT1.29}]\iff\Theta(c)=\delta h(c)\ (\bar\kappa=1).
\]

On all-nonnull per-edge-integer chains [FT.24/FT.34(4) scope]: $\bar\kappa\ge1$, with equality $\iff\kappa_e\equiv1$; null edges or class-(iii) rational profiles [FT.38(3)] admit $\bar\kappa<1$ (witness $(1,1,0)$: $\bar\kappa=\tfrac23$).

\textbf{Req.} FT.22, FT.24, FT.31, FT.34; CT0.9 (Part 1); C9; CT1.29 (Part 2); D1.6, D1.7.

\textbf{Pf.}

%dsp-level:0
\[
\frac{S_{\mathrm{reg}}}{A_{\mathrm{loop}}}=\frac{\delta h\ln2}{\Theta\,\mathfrak h^2}=\frac{\ln2}{\mathfrak h^2}\cdot\frac{\delta h}{\Theta};
\]

%dsp-level:0
\[
\qquad \Theta=2k\ [\mathrm{FT.22}]\implies\text{second form};
\]

%dsp-level:0
\[
\qquad \bar\kappa:=\Theta/\delta h\implies\text{first form}
\]

%dsp-level:0
\[
\mathrm{CT1.29}\ \text{at the single-edge floor}:\ S_{\min}=\ln2,\ A_{\min}=\mathfrak h^2\implies\text{rate}\ \ln2/\mathfrak h^2;\qquad\text{recovered}\iff\delta h/\Theta=1
\]

%dsp-level:0
\[
\text{integer non-null profiles}:\ \kappa_e\ge1\implies\Theta=\sum\kappa_e\ge\delta h\implies\bar\kappa\ge1,\ =\iff\kappa_e\equiv1;
\]

%dsp-level:0
\[
\qquad (1,1,0):\ \Theta=2,\ \delta h=3:\ \bar\kappa=\tfrac23<1.\ \blacksquare
\]

\textbf{Phys.} The loop-area/qubit-power identification is complete: CT1.29's one-bit-per-hop-area rate holds for every materialized loop, deformed by exactly the loop's own mean edge quantum. Minimal-edge loops are the undeformed case; the winding $k$ reads the deformation in the second form. Gap 2's "completeness unresolved" is resolved: the correction factor \textit{is} derived data.

\vspace{2pt}\noindent\rule{\textwidth}{0.4pt}\vspace{2pt}

\end{theorem}

\begin{theorem}[{\texorpdfstring{FT.38 — Non-geodesic tightening and on-shell rigidity}{FT.38 - Non-geodesic tightening and on-shell rigidity}}]\label{tag:FT-38}
\textbf{Thm.} $(1)$ \textbf{Tightening.} For any spacelike path $\gamma$ with on-shell integer edge data $\kappa_e\in\mathbb Z_{\ge1}$ [FT.24/FT.34(4) scope]:

%dsp-level:0
\[
I_{o\leftarrow\Sigma}(\gamma)\ \le\ \min\bigl(\delta h(\gamma),\ 4\,\delta\tau(\gamma)\bigr)\ln2\ \le\ \min\Bigl(\frac{\Delta\ell(\gamma)}{\mathfrak h},\ 4\,\delta\tau(\gamma)\Bigr)\ln2,
\]

%dsp-level:0
\[
\qquad \Delta\ell(\gamma)=\mathfrak h\sum_{e\in\gamma}\sqrt{\kappa_e}\ \ [\mathrm{Def\text{-}H.2},\ \mathrm{D1.6},\ \mathrm{FT.24}(2)],
\]

with no geodesic and no slow-variation hypothesis; spatial equality $\iff\kappa_e\equiv1$. FT.1's geodesic physical form is the special case; the bound is tightest on stable jumps [Def-H.19: local $\Delta\ell$-minimizers; Morse — math]. $(2)$ \textbf{On-shell rigidity.} In CT1.19's family regime ($L=\mathfrak h n_{\mathcal R}$ fixed, $\mathfrak h\to0$, $g$ smooth along $\gamma$ [CT2.6]): $\kappa_e=f(\lambda_e)+O(\mathfrak h)$ with $f:=g_{\mu\nu}\dot x^\mu\dot x^\nu$ continuous [CT1.19's Pf, midpoint Taylor; per-edge remainder uniform]. For per-edge-integer classes (i)/(ii) : $\exists\,\mathfrak h_0$ such that $\forall\,\mathfrak h<\mathfrak h_0$, $\kappa_e$ is \textbf{constant along each connected materialized chain} — signature included (spacelike $\kappa\ge1$, null $0$, timelike $\kappa\le-1$ are $\ge1$ apart). Hence uniform-$\kappa$ is continuum-generic on-shell: FT.33(3)'s length-universality and FT.33(4)'s Cauchy–Schwarz equality are the generic on-shell case, not specializations; signature changes are confined to $O(\mathfrak h)$ boundary layers. $(3)$ \textbf{Rational ladder (class-(iii) arena).} Without per-edge integrality, FT.22 alone permits uniform $\kappa\in(2/L)\mathbb Z$ on an $L$-cycle ($\Theta=L\kappa\in2\mathbb Z$; witness $L=3$, $\kappa=\tfrac23$, $\Theta=2$): per-edge integrality is special to $L=2$ [FT.24: $2\mathbb Z/2=\mathbb Z$] and round-trip-closed channels [FT.34(4)]; clause (2)'s integer-step argument does not reach spacings $2/L<1$.

\textbf{Req.} FT.1, FT.22, FT.24, FT.33, FT.34; Def-H.2, Def-H.19; D1.6; CT0.9 (Part 1); CT1.12, CT1.19 (Part 2); CT2.6 hypotheses (inherited); integer/rational arithmetic, Morse minimality [math].

\textbf{Pf.}

%dsp-level:0
\[
(1):\ \mathrm{FT.1}\ \text{first form: hypothesis-free};
\]

%dsp-level:0
\[
\qquad \kappa_e\ge1\implies\sqrt{\kappa_e}\ge1\implies\Delta\ell/\mathfrak h=\sum_e\sqrt{\kappa_e}\ \ge\ \delta h\implies\text{chain};
\]

%dsp-level:0
\[
\qquad \text{equality}\iff\kappa_e\equiv1
\]

%dsp-level:0
\[
\Delta\ell(\gamma)
\]

%dsp-level:0
\[
= \mathfrak h\textstyle\sum\sqrt{\kappa_e}:\ \text{per-edge}\ g\text{-length}\ \sqrt{ds^2(e)}
\]

%dsp-level:0
\[
= \sqrt{\kappa_e}\,
\]

%dsp-level:0
\[
\mathfrak h\ [\mathrm{FT.24}(2),\ \text{spacelike}];
\]

%dsp-level:0
\[
\ \text{piecewise additivity}\ [\mathrm{Def\text{-}H.2}];
\]

%dsp-level:0
\[
\qquad \text{stable jumps minimize}\ \Delta\ell\ \text{locally}\ [\mathrm{Def\text{-}H.19};\ \text{math}]
\]

%dsp-level:0
\[
(2):\ |\kappa_{e+1}-\kappa_e|\le|f(\lambda_{e+1})-f(\lambda_e)|+2C\mathfrak h\le\bigl(\tfrac{\sup|f'|}{L}+2C\bigr)\,\mathfrak h\ \to0;
\]

%dsp-level:0
\[
\qquad \mathfrak h<\bigl(\tfrac{\sup|f'|}{L}+2C\bigr)^{-1}\implies|\kappa_{e+1}-\kappa_e|<1\wedge\kappa\in\mathbb Z\implies\kappa_{e+1}=\kappa_e
\]

%dsp-level:0
\[
\text{signature}:\  \kappa\ge1\ \vee\ \kappa=0\ \vee\ \kappa\le-1\ \text{pairwise distance}\ \ge1\implies\text{constancy includes signature}
\]

%dsp-level:0
\[
(3):\ \Theta=L\kappa\in2\mathbb Z\ [\mathrm{FT.22}]\implies\kappa\in(2/L)\mathbb Z;
\]

%dsp-level:0
\[
\qquad L=3:\ \kappa=\tfrac23:\ \Theta=2\in2\mathbb Z\ \wedge\ \kappa\notin\mathbb Z;
\]

%dsp-level:0
\[
\qquad L=2:\ \kappa\in2\mathbb Z/2=\mathbb Z\ [\mathrm{FT.24}\ \text{recovered}].\ \blacksquare
\]

\textbf{Rmk (scope discipline).} Clause (2) upgrades the uniform-$\kappa$ regime to generic on-shell status \textit{within the integrality classes}; class (iii) is exactly where rigidity can fail — a second structural signature distinguishing the classes (the first being FT.34(4)'s round-trip criterion). FT.32(6)'s uniform/all-nonnull regimes are scoped accordingly.

\textbf{Phys.} The information bound needs no geodesic: any spacelike on-shell path obeys it with its own proper length, tightest along stable jumps. And the continuum limit does not erode the integer ledger — a slowly varying geometry cannot move an integer, so materialized chains of the integrality classes become uniformly quantized along their length in the classical regime: integrality plus smoothness equals rigidity. Where integrality is not yet derived, the ledger can refine to spacing $2/L$ — and that is precisely the open frontier.

\vspace{2pt}\noindent\rule{\textwidth}{0.4pt}\vspace{2pt}

\end{theorem}

\begin{theorem}[{\texorpdfstring{FT.39 — Confinement converse, singlet diagonal, and filter placement}{FT.39 - Confinement converse, singlet diagonal, and filter placement}}]\label{tag:FT-39}
\textbf{Thm.} $(1)$ \textbf{Converse (Gap 8).} $\Pi_{\mathcal C}(\Sigma;C)=\mathcal C(\Sigma;C)\cdot\mathbb I$ with $\mathcal C\in\{0,1\}$ [Def-H.10: scalar idempotent] $\implies F_{\mathrm{strong}}^\Sigma(C)=\mathcal C\cdot F_{\mathrm{strong}}\big|_\Sigma$ [FT.19's definition]. Hence $\mathcal C=1\implies F_{\mathrm{strong}}^\Sigma=F_{\mathrm{strong}}\big|_\Sigma\ne0$ for nonzero aspect data ($f_a\ne0$ [CT3.9, CT3.10]) on $\ge1$ colour slot — generic nonvanishing with exact exception locus $\{f=0\}\ \cup\ \{\mathrm{Colour}(\Sigma)=\varnothing\}$. At support level [Def-H.11's selection-type read]: $\Pi_{\mathcal C}=\mathbb I\ne0$ whenever $\mathcal C=1$ — unconditional. $(2)$ \textbf{Singlet diagonal (consistency).} On sealed (singlet) states the Lie action annihilates: $T_a|\text{singlet}\rangle=0$ [CT3.15, CT3.17 invariance] $\implies\langle\text{singlet}|F_{\mathrm{strong}}|\text{singlet}\rangle=0$ — consistent with FT.11 (no independent colour charge) and CT3.22 (colour is not an end-charge): the strong block's nonzero content is the admissibility structure and the gluon mode count [CT3.13], not singlet expectation values. $(3)$ \textbf{Filter placement (Gap 9, decision derived).} Def-H.9's $\mathcal M(\Sigma)$ must \textbf{not} be pre-weighted by $\mathcal C(\Sigma;C)$: its members are restrictions $C|_\Sigma$ of valid ambient configurations [FT.18's setup; D2.11], whose colour sealing is ambient-global [CT3.18] — $\mathcal C(\Sigma;C)=0$ records that $\Sigma$'s partners lie outside [FT.18's non-locality], not that $C$ is invalid. Pre-weighting would delete the EM/$G$/$H$ reads of externally-sealed configurations, contradicting their colour-independence [FT.11, FT.21(i)–(ii)]; the downstream $\Pi_{\mathcal C}$ filter [Def-H.15(3)] zeroes exactly the strong block [FT.19] and keeps the rest — the corpus-consistent placement. CT2.8(3)'s null-quotient handles \textit{universe-level} invalidity [CT3.18], a different question from $\Sigma$-relativity.

\textbf{Req.} Def-H.9, Def-H.10, Def-H.11, Def-H.15; FT.11, FT.18, FT.19, FT.21; D2.11; CT2.8 (Part 2); CT3.9, CT3.10, CT3.13, CT3.15, CT3.17, CT3.18, CT3.22 (Part 2).

\textbf{Pf.}

%dsp-level:0
\[
(1):\ \mathcal C\in\{0,1\}\implies\mathcal C^2=\mathcal C\implies\Pi_{\mathcal C}F\Pi_{\mathcal C}=\mathcal C\,F\big|_\Sigma;
\]

%dsp-level:0
\[
\qquad \mathcal C=1:\ F^\Sigma=F\big|_\Sigma;\quad T_a\ne0\ \text{on}\ \mathbb C^3\ [\mathrm{CT3.9}:\ \text{explicit basis}]\ \wedge\ \text{slot present}\implies F\big|_\Sigma\ne0\ \text{for}\ f\ne0
\]

%dsp-level:0
\[
(2):\ \varepsilon_{ijk}U^i{}_aU^j{}_bU^k{}_c
\]

%dsp-level:0
\[
= \varepsilon_{abc}\det U
\]

%dsp-level:0
\[
= \varepsilon_{abc}\ [\mathrm{CT3.15}];
\]

%dsp-level:0
\[
\quad\delta^i{}_j\ \text{likewise}\ [\mathrm{CT3.17}]\implies\text{singlets fixed under}\ SU(3)\implies\mathfrak{su}(3)\ \text{annihilates them}
\]

%dsp-level:0
\[
(3):\ \text{validity is ambient}\ [\mathrm{D2.11};\ \mathrm{CT3.18}:\ \mathcal P_M^{(3)}\ \text{excludes free indices at universe level}];
\]

%dsp-level:0
\[
\qquad \mathrm{FT.18}:\ \exists\,C,C'\in\mathcal M(U)\ \text{valid with}\ C|_\Sigma=C'|_\Sigma,\ \mathcal C(\Sigma;C)=1,\ \mathcal C(\Sigma;C')=0
\]

%dsp-level:0
\[
C'\ \text{valid yet}\ \Sigma\text{-unsealed}:\ \text{pre-weighting deletes}\ C'|_\Sigma\ \text{entirely} — \text{but its EM}/G/H\ \text{reads are physical and colour-independent}\ [\mathrm{FT.21}]\implies\text{pre-weighting loses physical data};
\]

%dsp-level:0
\[
\qquad \text{downstream}\ \Pi_{\mathcal C}:\ \text{strong block}\to0\ [\mathrm{FT.19}],\ \text{rest kept}.\ \blacksquare
\]

\textbf{Phys.} Confinement is now a two-sided exact statement at subnetwork level: the colour block lives precisely when the region self-seals, dies entirely when any attachment reaches outside, and even when alive its singlet-diagonal read vanishes — colour is never a number a sealed object carries, only the grammar of how it composes with its environment. And the filter belongs downstream: sealing is a fact about $\Sigma$'s embedding, so it must gate $\Sigma$'s strong block, never erase $\Sigma$'s other reads.

\vspace{2pt}\noindent\rule{\textwidth}{0.4pt}\vspace{2pt}

\end{theorem}

\begin{theorem}[{\texorpdfstring{FT.40 — Closed-core charge read}{FT.40 - Closed-core charge read}}]\label{tag:FT-40}
\textbf{Thm.} For a closed oriented chain $c$ based at $v=s(c)=t(c)$ with first edge $e_1$ (out of $v$) and last edge $e_k$ (into $v$) [D1.4's ordered tuple]:

%dsp-level:0
\[
q_{\mathrm{core}}(c;v):=\frac{\sigma_{e_1}-\sigma_{e_k}}{2}\ \in\ \{-1,0,+1\}
\]

%dsp-level:0
\[
\qquad [\mathrm{D3.4}\ \text{applied to closed}\ c;\ \mathrm{CT0.9}:\ \text{per-edge independent registers};\ \mathrm{EM.9}:\ \text{alphabet, all values attained}]
\]

is well-defined; under basepoint/orientation reversal [D2.6's $\omega$: $s\leftrightarrow t$, charge $\mapsto$ $-$charge]: $q_{\mathrm{core}}\mapsto-q_{\mathrm{core}}$ — the corpus's own charge-conjugation law, verbatim on the core read. At the 2-cycle floor: $q_{\mathrm{core}}=(\sigma_{e_1}-\sigma_{e_2})/2$ at base $p$, negated at base $q$. The closed core is therefore \textbf{not} charge-neutral by closure; its ternary read is basepoint data, and the observer-facing charge remains the EM.9 end-read of the full transit chain through the channel [Def-H.4].

\textbf{Req.} D1.4, D2.6, D3.4; CT0.9 (Part 1); EM.9 (Part 2); FT.28, Def-H.4.

\textbf{Pf.}

%dsp-level:0
\[
\mathrm{D3.4}\ \text{is stated for oriented chains}\ c\ \text{generally};
\]

%dsp-level:0
\[
\ \text{closed}\ c\ \text{has}\ s(c)
\]

%dsp-level:0
\[
= t(c)
\]

%dsp-level:0
\[
= v\ \text{with distinguished first/last edges}\ [\mathrm{D1.4}:\ c=(e_1,\dots,e_k)\ \text{ordered}]\implies\sigma_{s(c)}:
\]

%dsp-level:0
\[
= \sigma_{e_1},
\]

%dsp-level:0
\[
\ \sigma_{t(c)}:
\]

%dsp-level:0
\[
= \sigma_{e_k}\ \text{well-defined}
\]

%dsp-level:0
\[
\mathrm{CT0.9}:\ \sigma_{e_1},\sigma_{e_k}\in\{\pm1\}\ \text{independent registers};\qquad\mathrm{EM.9}:\ \text{ternary alphabet},\ \text{all attained}
\]

%dsp-level:0
\[
\text{rebase/orientation reversal}:\ e_1\leftrightarrow e_k\ \text{roles swap}\implies q\mapsto-q
\]

%dsp-level:0
\[
\qquad [\mathrm{D2.6}:\ \omega\ \text{maps charge}\mapsto\text{-charge}:\ \text{consistent}].\ \blacksquare
\]

\textbf{Phys.} A fermion core carries a ternary charge datum — the basepoint sign difference — that flips under reversal exactly as the corpus's orientation involution demands. What an observer measures through a channel remains the end-read of the full transit chain; the core datum is the bound ledger's contribution to that read, not an independent free charge [CT3.22's asymmetry untouched: this is registration-sign data, not colour].

\vspace{2pt}\noindent\rule{\textwidth}{0.4pt}\vspace{2pt}

\end{theorem}

\begin{theorem}[{\texorpdfstring{FT.41 — The per-aspect ledger}{FT.41 - The per-aspect ledger}}]\label{tag:FT-41}
\textbf{Thm.} $(a)$ \textbf{Charge count derived.} EM.9/D3.4's read consumes exactly the two end registers of a chain; $N_{\mathrm{charge}}=2$ in FT.14's decomposition is a theorem and $N_{\mathrm{weak}}=\delta h-2$ its arithmetic complement. the residual interface content is the \textit{attribution} clause of Def-H.8 (interior weight bits read through the weak/spin channel [Def-H.8]) — reduced — Def-H.8's attribution clause. $(b)$ \textbf{Colour ledger counted.} The sealing grammar of $\mathrm{Colour}(\Sigma)$ — admissible groupings [FT.34(5), Def-H.10], closed form, small-$N$ counts $g$, and attachment-type refinement — is \textbf{CT3.17a(2)–(3) (F)}. The ledger read defined here: $N_{\mathrm{strong}}(\Sigma):=\log_2 g\bigl(|\mathrm{Colour}(\Sigma)|\bigr)$ — the sealing-pattern information of the grammar count (a registration-type ledger, per the trichotomy below). $(c)$ \textbf{Gravity/Higgs ledger reframed.} $G$'s on-shell ledger is the integer $\kappa$-profile modulo the single loop constraint [FT.22]: $L-1$ free integers on a materialized $L$-cycle — unbounded, hence $\neg$bit-type [C10: bit counts attach to finite registered alternative sets]; $H$'s ledger is the continuous P4 datum $h_0$ [EW.13, P4]. Per-aspect "bit counts" are well-posed exactly for the registration-type ledgers $(a),(b)$ — the trichotomy answers the original question. $(d)$ \textbf{Strictly new at $\Sigma$, first counts [Def-H.13's filter].} $N(\Sigma)=1$: $\mathfrak u(1)$ factors per-edge entirely — nothing strictly new; single-contact floors carry no collective boundary mode. $N(\Sigma)=2$: $\dim\mathfrak u(2)=4$, factored diagonal $=2$, strictly new $=N^2-N=2$ real dimensions — the off-diagonal boundary mixings $|e_1\rangle\langle e_2|+\mathrm{h.c.}$, $i(|e_1\rangle\langle e_2|-\mathrm{h.c.})$: a collective inter-contact exchange mode, first available at two-contact floors [FT.28's variant Rmk].

\textbf{Req.} EM.9 (Part 2); D3.4, D3.8; CT0.9, CT0.13 (Part 1); C10; CT3.17, CT3.17a (Part 2); FT.14, FT.22, FT.28, FT.34; Def-H.8, Def-H.10, Def-H.13; EW.13 (Part 2); P4.

\textbf{Pf.}

%dsp-level:0
\[
(a):\ \mathrm{EM.9}'\text{s table reads}\ (\sigma_s,\sigma_t)\ \text{only};\ \mathrm{D3.4}'\text{s formula references no interior register}\implies\text{count}=2;
\]

%dsp-level:0
\[
\qquad N_{\mathrm{visible}}=\delta h\ [\mathrm{FT.14}]\implies N_{\mathrm{weak}}=\delta h-2\ \text{arithmetic}
\]

%dsp-level:0
\[
(b):\ \text{grammar}+\text{counts}+\text{type refinement}:\ \mathrm{CT3.17a}\ (\text{F});
\]

%dsp-level:0
\[
\qquad N_{\mathrm{strong}}\ \text{read}:\ \text{definition on the cited counts}\ [\mathrm{C10}/\mathrm{C9}]
\]

%dsp-level:0
\[
(c):\ \kappa\text{-profile}:\ L\ \text{integers},
\]

%dsp-level:0
\[
\ 1\ \text{constraint}\ \Theta
\]

%dsp-level:0
\[
\in 2\mathbb Z\ [\mathrm{FT.22}]
\]

%dsp-level:0
\[
\implies L{-}1\ \text{free integers},
\]

%dsp-level:0
\[
\ \text{unbounded}\implies\neg\text{finite register count}\ [\mathrm{C10}];
\]

%dsp-level:0
\[
\qquad h_0\in\mathbb R\ [\mathrm{P4}]
\]

%dsp-level:0
\[
(d):\ \iota_\Sigma\ \text{edge-wise}\ [\mathrm{Def\text{-}H.13}]:\ \text{factored subalgebra}=\bigoplus_e\mathfrak u(1)_e:\ \dim=N;
\]

%dsp-level:0
\[
\qquad \dim\mathfrak u(N)=N^2\implies\text{strictly new}=N^2-N:\ N{=}1:\ 0;\quad N{=}2:\ 2\ \text{off-diagonal generators}.\ \blacksquare
\]

\textbf{Phys.} "How many bits per force?" now has a typed answer: charge costs two end bits (theorem); colour costs the log of the sealing grammar [counts: CT3.17a(2) (Part 2)]; gravity keeps an unbounded integer ledger; Higgs keeps a continuous dial. And collective boundary structure — genuinely new physics at $\Sigma$ — begins at exactly two contacts, in two dimensions of exchange modes: the first place where a subnetwork can do something its edges cannot.

\vspace{2pt}\noindent\rule{\textwidth}{0.4pt}\vspace{2pt}

\vspace{2pt}\noindent\rule{\textwidth}{0.4pt}\vspace{2pt}

\end{theorem}

\phantomsection
\section[Tier F.14 — Platonic Floor Cards: The Five Solids as…]{\texorpdfstring{Tier F.14 — Platonic Floor Cards: The Five Solids as Fermion Floors}{Tier F.14 - Platonic Floor Cards: The Five Solids as Fermion Floors}}

\begin{definition}[{\texorpdfstring{Def-H.20 — Platonic floors (setup vocabulary)}{Def-H.20 - Platonic floors (setup vocabulary)}}]\label{tag:Def-H-20}
\textbf{Def.} Let $S$ be a Platonic solid with 1-skeleton $\Gamma_S$: $V$ vertices, $E$ undirected edges, $F$ faces of uniform length $L$ ($FL=2E$), regular vertex degree $d$ ($Vd=2E$); cardinals [math]:

%dsp-level:0
\[
\text{tetrahedron}\ (4,6,4,3);
\]

%dsp-level:0
\[
\quad\text{cube}\ (8,12,6,3);
\]

%dsp-level:0
\[
\quad\text{octahedron}\ (6,12,8,4);
\]

%dsp-level:0
\[
\quad\text{dodecahedron}\ (20,30,12,3);
\]

%dsp-level:0
\[
\quad\text{icosahedron}\ (12,30,20,5).
\]

The \textbf{Platonic floor} $\Sigma_S$ $:=$ the bounded subnetwork on $\Gamma_S$'s vertex set whose real-edge content realizes every undirected edge as a \textbf{bidirectional channel} (both orientations in $\mathcal R$) — the full bidirectional 1-skeleton. Variants: \textbf{dark} ($\partial\Sigma_S=\varnothing$: no observer contact [Def-H.1]) and \textbf{single-contact observed} (one external $o$ joined by exactly one real edge to one $v_1$). Card hypotheses, declared at each use: \textbf{(M1)} every face cycle materializes [FT.22's scope]; \textbf{(M2)} every bidirectional pair materializes as a 2-cycle [FT.24's scope] — both are spectral questions the corpus does not decide (FT.32's Status pattern). Sub-floor vocabulary: $\Sigma_{\min}$ is the $E{=}1$ cell [FT.28], $\Sigma_{12}$ the $E{=}3$ cell [FT.35], the cycle ladder $\Sigma_L$ the unidirectional family [FT.32].

\vspace{2pt}\noindent\rule{\textwidth}{0.4pt}\vspace{2pt}

\end{definition}

\phantomsection
\section[Tier F.15 — The Regular-Polygon Floors: Full Ladder,…]{\texorpdfstring{Tier F.15 — The Regular-Polygon Floors: Full Ladder, Sectors, and the Non-Existence Ceiling}{Tier F.15 - The Regular-Polygon Floors: Full Ladder, Sectors, and the Non-Existence Ceiling}}

\begin{theorem}[{\texorpdfstring{FT.46 — Distinctness is unbounded: no upper limit on distinct regular polygons}{FT.46 - Distinctness is unbounded: no upper limit on distinct regular polygons}}]\label{tag:FT-46}
\textbf{Thm.} $(1)$ \textbf{Combinatorial distinctness.} The regular $n$-gons are pairwise non-isomorphic for all $n\ge3$ ($n$ is the complete invariant of the 1-skeleton); $\Sigma_{P_n}$ is a valid bounded subnetwork configuration class for every $n$ [A1: $\mathcal Q$ countably infinite supplies vertices; D1.2: seals supply edges; no mask or capacity clause restricts $n(v)=4$ bidirectional vertices — oversubscription is absorbed and retained [CT4.1, CT4.3; rung correction, FT.62], CT0.6: $O(4)=16$ uniform]. $(2)$ \textbf{Ledger injectivity.} On the uniform-$\kappa$ minima, the card invariants separate every member:

%dsp-level:0
\[
m_{\mathrm{hid}}(n)=\frac{2n\ln2}{\mathfrak h}\ \text{strictly increasing};
\]

%dsp-level:0
\[
\qquad E_{\mathrm{int}}(n)=2n\ \text{injective};
\]

%dsp-level:0
\[
\qquad \beta_1(n)=n+1\ \text{injective};
\]

%dsp-level:0
\[
\qquad \text{null-carrying dimension}\ n-2\ \text{injective};
\]

%dsp-level:0
\[
\Theta_{\min}(n)=\begin{cases}2 & \text{spacelike-or-null / signature-mixed}\\ n+(n\ \mathrm{mod}\ 2) & \text{strictly-spacelike}\\ n\ \text{or}\ 2n & \text{uniform (even/odd)}\end{cases}\quad\text{regime-separating};\qquad\text{lattice index}=2\ \forall n\ [\mathrm{FT.45}(3)]\ (\text{uniform, non-separating}).
\]

\textbf{No upper limit exists on the number of distinct regular-polygon floors.} $(3)$ \textbf{Embeddability (declared, scoped).} For every $n$ there exist configurations in which $\Sigma_{P_n}$ embeds: A1's countable infinitude, A2's channel universality ($C(p,q)=1\,\mathfrak q$ for all pairs), and D1.2's seal dynamics impose no cardinality bound; the A4 reality test is threshold data, not a prohibition. This is a configuration-class existence statement, not a claim about any particular solution of A9. $(4)$ \textbf{Non-degeneracy of the ladder.} $m_{\mathrm{hid}}$ and $\beta_1$ increase by fixed steps ($2\ln2/\mathfrak h$ and $1$); $\Theta_{\min}$ is constant ($2$) in the null-carrying/mixed regimes, but those floors are separated by the null-carrying dimension $n-2$; the parity lattice index is constant ($2$) — the family's parity ledger is uniform, and distinctness rides on rank, mass, and null structure, not sector count. No two members are ledger-degenerate.

\textbf{Req.} A1, A2; D0.8, D1.2; CT0.5, CT0.6, CT0.11 (Part 1); FT.42, FT.45; Def-H.20; graph-isomorphism fact for cycles [math].

\textbf{Pf.}

%dsp-level:0
\[
(1):\ C_n\not\cong C_m\ (n\ne m):\ \text{vertex/edge counts differ}\ [\text{math}];
\]

%dsp-level:0
\[
\qquad \text{validity}:\ \Sigma_{P_n}\ \text{is a finite bounded vertex set with real-edge content}\ [\mathrm{D0.10}\ \text{configuration class}];
\]

%dsp-level:0
\[
\ \ \neg\exists\ \text{clause in}\ \mathrm{A1}\text{--}\mathrm{A9}\ \text{bounding}\ n
\]

%dsp-level:0
\[
\mathrm{CT0.6}:\ O(2)=0:\ \text{no oversubscription at any polygon vertex, dark}\ [\mathrm{FT.45}(2)]
\]

%dsp-level:0
\[
(2):\ m_{\mathrm{hid}}(n)=2n\ln2/\mathfrak h\ [\mathrm{FT.45}(1)]:\ n\mapsto2n\ \text{injective};
\]

%dsp-level:0
\[
\qquad \beta_1(n)=n+1\ [\mathrm{FT.45}(3)];
\]

%dsp-level:0
\[
\qquad \text{null-carrying dimension}\ n-2\ \text{injective};
\]

%dsp-level:0
\[
\qquad \text{regime table}\ [\mathrm{FT.45}(4)]
\]

%dsp-level:0
\[
(3):\ \text{declared embedding clause}\ [\text{see Req}]:\ \text{A1 countable infinitude},
\]

%dsp-level:0
\[
\ \text{A2 pair universality},
\]

%dsp-level:0
\[
\ \text{D1.2 seals};
\]

%dsp-level:0
\[
\qquad \text{no cardinality obstruction in}\ \mathrm{A1}\text{--}\mathrm{A9}
\]

%dsp-level:0
\[
(4):\ \text{direct from the regime table and}\ (2).\ \blacksquare
\]

\textbf{Rmk (what "distinct" does not yet mean).} Clause (2) separates the floors as \textit{ledger objects}. No generation claim is made here.

\textbf{Phys.} The regular-polygon family is unlimited: every $n$ gives a valid, uniformly minimally-locked, colourless fermion floor with its own mass rung, its own winding rank $n{+}1$, and its own null-carrying dimension. If the polygon ladder is ever read as a family structure, nothing in the axioms caps its length — the cap, if any, must come from the spectrum side, not the ledger side.

\vspace{2pt}\noindent\rule{\textwidth}{0.4pt}\vspace{2pt}

\end{theorem}

\begin{theorem}[{\texorpdfstring{FT.47 — No geometric ceiling: the Gauss–Bonnet bound binds area, never $n$}{FT.47 - No geometric ceiling: the Gauss-Bonnet bound binds area, never}}]\label{tag:FT-47}
\textbf{Thm.} $(1)$ \textbf{Gauss–Bonnet area bound.} A geodesic $n$-gon enclosing a topological disk in a smooth Riemannian surface of constant curvature $\bar K$, interior angles $\alpha_i$:

%dsp-level:0
\[
\sum_{i=1}^n\alpha_i=(n-2)\pi+\bar K\,A\qquad[\text{Gauss–Bonnet};\ \text{geodesic edges}:\ \kappa_g\equiv0;\ \text{math}],
\]

and convexity ($\alpha_i<\pi\ \forall i$) forces $\bar K A<2\pi$ — a bound on the \textbf{enclosed area per polygon, with no $n$-dependence}. $(2)$ \textbf{Model-geometry existence — every $n\ge3$, no ceiling.} Flat ($\bar K=0$): the Euclidean regular $n$-gon, $\alpha=\frac{(n-2)\pi}{n}\in(0,\pi)$ for every $n\ge3$ [math]. Spherical ($\bar K>0$): for every $A\in(0,2\pi/\bar K)$ a convex regular geodesic $n$-gon with $\alpha=\frac{(n-2)\pi+\bar KA}{n}\in\bigl(\frac{(n-2)\pi}{n},\pi\bigr)$ exists for every $n\ge3$ [math: polar construction on $S^2$; small-area members are convex]. Hyperbolic ($\bar K<0$): every $\alpha\in\bigl(0,\frac{(n-2)\pi}{n}\bigr)$ is realized by a unique regular $n$-gon of area $A=\frac{n\alpha-(n-2)\pi}{\bar K}>0$, for every $n\ge3$ [math]. $(3)$ \textbf{Local existence in the general emergent geometry.} CT2.6's spatial slices are smooth Riemannian in the perturbative regime [CT1.12: $\bar g_{ij}\succ0$]; every point has a convex normal neighborhood [math: standard existence], inside which any $n\ge3$ points on a small geodesic circle join by unique minimizing geodesics — an \textbf{equilateral geodesic $n$-gon for every $n$}. Exact regularity (equal angles too) needs a local symmetry hypothesis; equilateral local existence does not. $(4)$ \textbf{Transfer.} Nothing in A1–A9 or Parts 1–3's definitions/theorems bounds polygon $n$: CT2.7's dynamics (Einstein–Hilbert, $\Lambda$ a P4 datum) fixes $\bar K$ per solution but no polygon ceiling — for $\bar K\le0$ clause (2) is global, for $\bar K>0$ local (small-area members), in every case for all $n\ge3$. The corpus's only numeric caps bind different objects: per-vertex accessibility ($n(v)\le3$ unoversubscribed [CT0.11, CT4.1] — polygon floors sit at $n(v)\in\{2,3\}$ [FT.45(2)], never touching it) and the spectrum (M1$_n$/M2$_n$ materialization — solution data, not a bound on $n$). $(5)$ \textbf{Verdict.} The question "do distinct regular polygons have an upper limit" is answered \textbf{no} at all three levels — combinatorial [FT.46], model-geometric (clause 2), local emergent-geometric (clause 3) — with the two bounds that do exist precisely located: the Gauss–Bonnet area bound (per polygon, curvature-signed) and the materialization filter (spectral, per solution).

\textbf{Req.} FT.45, FT.46; Def-H.19 (geodesic vocabulary); CT1.12, CT2.6, CT2.7 (Part 2); CT0.11 (Part 1); CT4.1 (Part 2); Gauss–Bonnet, model-geometry polygon constructions, convex-normal-neighborhood existence [math].

\textbf{Pf.}

%dsp-level:0
\[
(1):\ \text{Gauss–Bonnet on the enclosed disk}:\ \int_{\partial}\kappa_g\,ds+\iint\bar K\,dA+\sum_i(\pi-\alpha_i)=2\pi\ [\text{math}];
\]

%dsp-level:0
\[
\qquad \kappa_g\equiv0\ [\text{geodesic edges}]\implies\sum_i\alpha_i=(n-2)\pi+\bar KA
\]

%dsp-level:0
\[
\alpha_i<\pi\ \forall i\implies\sum_i\alpha_i<n\pi\implies(n-2)\pi+\bar KA<n\pi\implies\bar KA<2\pi\qquad(n\text{-free})
\]

%dsp-level:0
\[
(2):\ \bar K=0:\ \alpha=\tfrac{(n-2)\pi}{n}\in(0,\pi)\ \forall n\ge3;
\]

%dsp-level:0
\[
\qquad \bar K>0:\ A\in(0,2\pi/\bar K)\implies\alpha\in\bigl(\tfrac{(n-2)\pi}{n},\pi\bigr):\ \text{convex valid}\ \forall n\ \ [\text{constructions cited}]
\]

%dsp-level:0
\[
\bar K<0:\ A=\tfrac{n\alpha-(n-2)\pi}{\bar K}>0\iff\alpha<\tfrac{(n-2)\pi}{n};
\]

%dsp-level:0
\[
\qquad \bigl(0,\tfrac{(n-2)\pi}{n}\bigr)\ne\varnothing\ \forall n\ge3\ \text{and lies inside}\ (0,\pi)
\]

%dsp-level:0
\[
(3):\ \text{convex normal neighborhoods exist}\ [\text{math}];
\]

%dsp-level:0
\[
\quad\text{within one}:\ \exists!\ \text{minimizing geodesic per point pair}\implies\text{equilateral}\ n\text{-gons}\ \forall n
\]

%dsp-level:0
\[
\ge 3\ (\text{vertices on a small geodesic circle})
\]

%dsp-level:0
\[
(4)\text{--}(5):\ \mathrm{CT2.6}:\ (M,g)\ \text{smooth};
\]

%dsp-level:0
\[
\ \mathrm{CT1.12}:\ \bar g_{ij}\succ0\ \text{perturbatively};
\]

%dsp-level:0
\[
\quad\neg\exists\ n\text{-bounding clause in Parts 1–3};
\]

%dsp-level:0
\[
\quad\mathrm{CT2.7}/\mathrm{P4}:\ \bar K\ \text{solution data};
\]

%dsp-level:0
\[
\qquad \mathrm{FT.45}(2):\ n(v)\in\{2,3\}\ \forall\ \text{polygon floors}.\ \blacksquare
\]

\textbf{Rmk (what a real ceiling would require).} An upper limit on the polygon family could only arise from the \textit{spectral} side: a proof that M1$_n$/M2$_n$ fails for large $n$ in every valid universe (e.g. if $Z(\mathcal R,\Sigma_{P_n})$'s pole structure [Def-H.16] degenerated with $n$). The program-note intuition "higher levels make the impossible possible" receives its precise counterpart here: the levels (larger $n$) are always \textit{valid}; what remains undetermined per level is \textit{materialization}.

\textbf{Phys.} Ask the ledger how large a polygon floor can be and the answer is unbounded — every $n$ is a valid, distinct, lock-free configuration class. Ask the geometry and the answer is again unbounded: Gauss–Bonnet caps the \textit{area} a geodesic polygon can enclose ($\bar KA<2\pi$), never its number of sides; flat and hyperbolic slices realize every regular $n$-gon globally, positively curved slices realize small ones locally. The only gate left is spectral — which polygons materialize — and that gate is the corpus's own open question, not a ceiling.

\vspace{2pt}\noindent\rule{\textwidth}{0.4pt}\vspace{2pt}

\end{theorem}

\begin{theorem}[{\texorpdfstring{FT.48 — Stability structure of floor-to-floor channels (Gap 18, partial)}{FT.48 - Stability structure of floor-to-floor channels (Gap 18, partial)}}]\label{tag:FT-48}
\textbf{Thm.} Let $\gamma$ be a channel path between two floor termini, on-shell (edges of integrality classes (i)/(ii) ) in the continuum regime of CT1.19/CT2.6. $(1)$ \textbf{Stationarity is geodesy.} $\delta\!\int_\gamma ds^2=0$ among fixed-endpoint variations $\iff$ $\gamma$ is an affinely parametrized geodesic of $g_{\mu\nu}$ [CT1.20's stationary-phase locus; FT.31(2): the on-shell sums \textit{are} the Riemann sums, so the ledger's stationary condition and the metric's Euler–Lagrange equation coincide; math]. $(2)$ \textbf{Two-sided distance law.} For a bidirectional round-trip connector $P$ of $m$ channels [FT.34(2)]: the ledger side is the integer distance $d_{12}=\Theta(P)=\sum_{e\in P}\kappa_e\in\mathbb Z$; the metric side is the proper length $\ell(P)=\mathfrak h\sum_{e\in P}\sqrt{\kappa_e}$ [FT.33(1) pattern, spacelike]. They are bound by

%dsp-level:0
\[
\ell(P)^2\ \le\ m\cdot d_{12}\cdot\mathfrak h^2\qquad[\text{Cauchy–Schwarz};\ \text{math}],\qquad\text{equality}\iff\kappa_e\ \text{uniform on}\ P;
\]

uniform connectors obey $\ell^2=m\,d_{12}\,\mathfrak h^2$ exactly. Both sides are derived: the integer from materialization [FT.34(2)], the length from the metric dictionary [D1.6, FT.24(2)] — two reads of one channel, dimensionally consistent ($[\ell^2]=[\mathfrak h^2]=\mathfrak q^{-1}$, $d_{12}$ dimensionless). $(3)$ \textbf{Stable-channel trichotomy.} Within a fixed endpoint class in the emergent geometry [Def-H.19 vocabulary]: index-0 stationary channels (no conjugate points) are exactly the local $\ell$-minimizers [math: Morse/Synge]; on-shell with $\kappa_e>0$, local $\ell$-minimality at fixed channel count $m$ coincides with local $d_{12}$-minimality along uniform families [clause (2): $\ell^2=m\,d_{12}\,\mathfrak h^2$ strictly monotone in $d_{12}$ at fixed $m$]. Saddle/maximal stationary channels exist wherever conjugate points do; CT1.20's stationary phase weights them \textit{alike} with the minima. $(4)$ \textbf{Derived vs. application-layer.} Clauses (1)–(3) close the variational side of the stability program: stable channels are geodesics, their ledger distance is integer, and the two distance reads are Cauchy–Schwarz-bound with equality at uniformity. The selection side — whether the pole condition [CT2.5/Def-H.16] or the defect criterion [FT.33] \textit{selects} index-0 channels — remains open : it requires fluctuation-sign/Maslov-type structure on the $Z$ side, which Parts 1–3 does not develop.

\textbf{Req.} FT.22, FT.24, FT.31, FT.33, FT.34, FT.38; Def-H.19; D1.6; CT1.19, CT1.20, CT2.5, CT2.6 (Part 2); Cauchy–Schwarz, Morse/Synge basics [math].

\textbf{Pf.}

%dsp-level:0
\[
(1):\ \mathrm{FT.31}(2):\ \sum_e\kappa_e\mathfrak h^2
\]

%dsp-level:0
\[
= \sum_e ds^2(e)\ \text{identically}\implies\text{stationary locus of the ledger phase}\ \Phi
\]

%dsp-level:0
\[
= \pi\Theta\ [\mathrm{FT.9}/\mathrm{EM.20}]\
\]

%dsp-level:0
\[
= \ \text{stationary locus of}\ \textstyle\int g_{\mu\nu}\dot x^\mu\dot x^\nu d\lambda\ [\mathrm{CT1.19}]
\]

%dsp-level:0
\[
\delta\!\int g_{\mu\nu}\dot x^\mu\dot x^\nu d\lambda=0\iff\text{geodesic, affinely parametrized}\ [\text{Euler–Lagrange};\ \text{math}]
\]

%dsp-level:0
\[
(2):\ d_{12}=\Theta(P)=\sum_e\kappa_e\in\mathbb Z\ [\mathrm{FT.34}(2)];
\]

%dsp-level:0
\[
\qquad \ell(P)
\]

%dsp-level:0
\[
= \mathfrak h\sum_e\sqrt{\kappa_e}\ [\mathrm{D1.6}/\mathrm{FT.24}(2):\ ds^2(e)=\kappa_e\mathfrak h^2,\ \text{spacelike};\ \text{piecewise additivity}\ \mathrm{Def\text{-}H.2}]
\]

%dsp-level:0
\[
\bigl(\textstyle\sum_e\sqrt{\kappa_e}\bigr)^2\le m\sum_e\kappa_e\ [\text{Cauchy–Schwarz};\ \text{equality}\iff\sqrt{\kappa_e}\ \text{constant}]\implies\ell^2\le m\,d_{12}\,\mathfrak h^2;
\]

%dsp-level:0
\[
\qquad [\ell^2]=\mathfrak q^{-1}=[m\,d_{12}\,\mathfrak h^2]\ \checkmark\ [\mathrm{D2.9}]
\]

%dsp-level:0
\[
(3):\ \text{index-0}\iff\text{local}\ \ell\text{-minimizer}\ [\text{Morse/Synge};\ \mathrm{Def\text{-}H.19}];
\]

%dsp-level:0
\[
\qquad \text{uniform families at fixed}\ m:\ \ell^2=m\,d_{12}\,\mathfrak h^2\ \text{strictly monotone in}\ d_{12}\implies\text{minimizers coincide}
\]

%dsp-level:0
\[
(4):\ \mathrm{CT1.20}:\ \delta\Phi
\]

%dsp-level:0
\[
= 0\ \text{configurations dominate},
\]

%dsp-level:0
\[
\ \delta\Phi
\]

%dsp-level:0
\[
\ne 0\ \text{cancel pairwise}:\ \text{minima and saddles weighted alike};
\]

%dsp-level:0
\[
\qquad \neg\exists\ \text{fluctuation-sign machinery in Parts 1–3}\implies\text{selection open}.\ \blacksquare
\]

\textit{The two distance claims of the derivation record (the winding-integer distance and the mass-as-distance identification) are resolved in the Addendum: see Appendix B, Section B.2.1.}

\textbf{Phys.} A channel between two fermion floors carries two distances at once: an integer the ledger can only write in whole quanta, and a length the metric reads in generally irrational square roots — locked together by Cauchy–Schwarz, equal only when every hop carries the same quantum. Stable channels are the geodesics; the integer distance is the ledger's shadow of the metric one. The single thing the corpus cannot yet say is whether nature prefers the stable ones: that is the selection problem, and it is named.

\vspace{2pt}\noindent\rule{\textwidth}{0.4pt}\vspace{2pt}

\end{theorem}

\phantomsection
\section[Tier F.16 — Composition, Enumeration, and the…]{\texorpdfstring{Tier F.16 — Composition, Enumeration, and the Spin-Connection Obstruction}{Tier F.16 - Composition, Enumeration, and the Spin-Connection Obstruction}}

\begin{theorem}[{\texorpdfstring{FT.52 — Spin extraction: the constraint system, uniqueness of the stretched rule, and the localized obstruction}{FT.52 - Spin extraction: the constraint system, uniqueness of the stretched rule, and the localized obstruction}}]\label{tag:FT-52}
\textbf{Thm.} Let $R$ be an extraction rule in the sense of Def-H.9(3) — assigning admissible $\mathfrak{su}(2)$ labels $j$ to on-shell internal ledger data, under the weak action [FT.13]. Then:

%dsp-level:0
\[
(1)
\]

%dsp-level:0
\[
\quad R\ \text{cannot factor through}\ \mathrm{Hol}_{\mathrm{int}}\ \text{alone on chain-swept interiors}:\ (-1)^{\Theta_{\mathrm{int}}}
\]

%dsp-level:0
\[
= +1\ \text{on-shell}\ [\mathrm{FT.25}(1)];
\]

%dsp-level:0
\[
(2)
\]

%dsp-level:0
\[
\quad R\ \text{must respect composition}:\ k(\Gamma)
\]

%dsp-level:0
\[
= k_1+k_2+d_{12}\ [\mathrm{FT.34}(3),\ \mathrm{FT.50}(2)]\ \text{and}\ k_f\text{-sums}\ \sum_fk_f
\]

%dsp-level:0
\[
= K_S\ [\mathrm{FT.43}(1)]\implies\text{the stretched values add}:\ \text{any rank-1 additive}\ R\ \text{satisfies}\ R(k_1{+}k_2)
\]

%dsp-level:0
\[
= R(k_1){+}R(k_2)\ \text{on aligned compositions}\ [\text{math}:\ \text{Cauchy equation on}\ \mathbb Z];
\]

%dsp-level:0
\[
(3)\quad\textbf{uniqueness of the stretched rule}:\ \text{the unique rank-1 additive rule anchored at the seed}\ R(1)
\]

%dsp-level:0
\[
= \tfrac12\ [\mathrm{EW.32}]\ \text{is}\ R(k)
\]

%dsp-level:0
\[
= k/2\ \text{on the nonnegative winding cone},
\]

%dsp-level:0
\[
\ \text{with admissible labels}\ j
\]

%dsp-level:0
\[
= |k|/2\ (\text{admissibility correction}:\ \mathfrak{su}(2)\ \text{labels require}\ j\ge0,\ \text{and the Casimir distinguishes}\ \pm k/2:\ (-\tfrac32)(-\tfrac12)=\tfrac34\ne\tfrac{15}{4}=(\tfrac32)(\tfrac52)\ —\ \text{the absolute value is forced, machine-witnessed};\ \text{stretched additivity holds where all contributions are}\ \ge0\ —\ \text{spacelike-dominant compositions};\ k\ \text{is a signature datum, not orientation}\ [\mathrm{FT.23}]);
\]

%dsp-level:0
\[
\qquad k\ \text{odd}\Rightarrow j\in\tfrac12+\mathbb Z_{\ge0},\quad k\ \text{even}\Rightarrow j\in\mathbb Z_{\ge0};
\]

%dsp-level:0
\[
(4)\quad\textbf{the obstruction}:\ \text{promoting}\ R(k)
\]

%dsp-level:0
\[
= k/2\ \text{from candidate to theorem requires one missing lemma — a corpus-level spin connection}: \text{an identification of the weak transport holonomy}\ \mathcal U_c
\]

%dsp-level:0
\[
= \mathrm P\exp\bigl(i\oint_c\textstyle\sum_Ac_A\mathcal G_A\,d\ell\bigr)\ [\mathrm{EW.33}] \text{with the frame rotation of the emergent metric along}\ c\ [\mathrm{CT1.13}\text{--}\mathrm{CT1.17}'\text{s Levi-Civita data lifted to the}\ SU(2)\ \text{carrier}\ [\mathrm{EW.32}]].
\]

%dsp-level:0
\[
\text{Parts 1–3 contains every ingredient but the lift}:\ g_{\mu\nu}\leftrightarrow(a_{00},q_{ij},w_i)\ \text{dictionary with}\ w_i\ \text{in the}\ W\ \text{block}\ [\mathrm{CT1.12},\ \mathrm{FT.30}];
\]

%dsp-level:0
\[
\ \ \Gamma^\lambda{}_{\mu\nu},
\]

%dsp-level:0
\[
\ R^\rho{}_{\sigma\mu\nu}\ [\mathrm{CT1.13},\ \mathrm{CT1.14}];
\]

%dsp-level:0
\[
\ \ W_i
\]

%dsp-level:0
\[
= \tfrac12(\sigma_i^{(p)}+\sigma_i^{(q)})\ \text{total-spin generators}\ [\mathrm{FT.13},\ \mathrm{CT1.2}];
\]

%dsp-level:0
\[
\ \ \text{seed transport}\ [\mathrm{EW.32}/\mathrm{EW.33}]\ —\ \text{but}\ \neg\exists\ \text{tag constructing the spinor connection or tying}\ \oint\text{(weak)}\ \text{to}\ \oint\text{(frame rotation)}. \text{Without it,}\ k\ \text{(geometric winding)}\ \text{and}\ j\ \text{(weak Casimir)}\ \text{are independently derived integers/half-integers with no derived bridge};
\]

%dsp-level:0
\[
(5)\quad\textbf{verdict}:\ j
\]

%dsp-level:0
\[
= k/2\ \text{remains the leading candidate}\ —\ \text{it passes every available consistency test}:\ (1)\ \text{parity triviality (it reads}\ k,\ \neg\text{parity});
\]

%dsp-level:0
\[
\ (2)\ \text{additive composition}\ [\mathrm{FT.34}(3)];
\]

%dsp-level:0
\[
\ (3)\ \text{seed anchor}\ R(1)
\]

%dsp-level:0
\[
= \tfrac12\ [\mathrm{EW.32}];
\]

%dsp-level:0
\[
\ (4)\ \text{half-integrality from odd winding}\ [\text{clause}(3)].\ \text{Promotion blocked exactly at Open Gaps 19 (spin connection)}\ —\ \text{not at any arithmetic step}.
\]

\textbf{Req.} FT.13, FT.25, FT.33, FT.34, FT.43, FT.50; Def-H.8, Def-H.9(3); CT1.2 (Part 1); CT1.12, CT1.13, CT1.14 (Part 2); EW.32, EW.33 (Part 2); Cauchy functional equation on $\mathbb Z$, Clebsch–Gordan basics [math].

\textbf{Pf.}

%dsp-level:0
\[
(1):\ \mathrm{FT.25}(1):\ \Theta_{\mathrm{int}}
\]

%dsp-level:0
\[
\in 2\mathbb Z\ \text{on chain-swept interiors}
\]

%dsp-level:0
\[
\implies (-1)^{\Theta_{\mathrm{int}}}
\]

%dsp-level:0
\[
\equiv +1:\ \text{constant}\implies\neg\text{label-carrying}
\]

%dsp-level:0
\[
(2):\ \text{aligned composition}:\ k(\Gamma)
\]

%dsp-level:0
\[
= k_1+k_2+d_{12}\ [\mathrm{FT.34}(3)];
\]

%dsp-level:0
\[
\quad\text{additive}\ R:\ R(k(\Gamma))
\]

%dsp-level:0
\[
= R(k_1)+R(k_2)+R(d_{12})\ \text{required for the stretched (maximal-weight) component of the CG decomposition}\ [\text{math}:\ j_{\max}=j_1+j_2]
\]

%dsp-level:0
\[
(3):\ R:\ \mathbb Z\to\tfrac12\mathbb Z\ \text{additive}\implies R(k)=k\,R(1)\ [\text{Cauchy on}\ \mathbb Z];
\]

%dsp-level:0
\[
\qquad R(1)=\tfrac12\ [\mathrm{EW.32}\ \text{anchor}]\implies R(k)=k/2;
\]

%dsp-level:0
\[
\qquad \text{parity claims immediate}
\]

%dsp-level:0
\[
(4):\ \text{ingredient inventory as stated}:\ \mathrm{CT1.12}\ (\text{dictionary}),
\]

%dsp-level:0
\[
\ \mathrm{CT1.13}/\mathrm{CT1.14}\ (\text{Levi-Civita data}),
\]

%dsp-level:0
\[
\ \mathrm{FT.13}/\mathrm{CT1.2}\ (W_i\ \text{as spin generators}),
\]

%dsp-level:0
\[
\ \mathrm{EW.32}/\mathrm{EW.33}\ (\text{seed transport});
\]

%dsp-level:0
\[
\qquad \text{the lift}\ \Gamma\mapsto\omega^{ab}{}_\mu\mapsto SU(2)\ \text{connection on}\ \mathcal H_{\mathsf v}:\ \neg\exists\ \text{in Parts 1–3}\ [\text{tag inventory complete}]
\]

%dsp-level:0
\[
(5):\ \text{conjunction of}\ (1)\text{--}(4).\ \blacksquare
\]

\textbf{Rmk (why this is progress, not failure).} Before this tier, "$j=k/2$" was one candidate among four. After it: every rank-1 additive rule \textit{is} $k/2$ (uniqueness), all available consistency tests are passed, and the single missing link is named, located, and minimal — a spin-connection lemma, not a new postulate and not a search space. Gap 11 is reduced to Gap 19; Gap 19 is a construction question about lifting existing corpus data (CT1.13's $\Gamma$, EW.32's carrier), which is exactly the kind of step A §0.2a admits as [math]-class work if executed without new physical input.

\textbf{Phys.} If the weak transport of the seed around a bound chain is the frame rotation of the geometry the chain lives in, then spin is half the winding: odd windings give half-integer spin, even windings integer spin, aligned compositions add — and the whole fermion program's spin question reduces to one unbuilt bridge between two objects the corpus already has. The bridge is named; it is not yet built, and nothing here claims it is.

\vspace{2pt}\noindent\rule{\textwidth}{0.4pt}\vspace{2pt}

\end{theorem}

\phantomsection
\section[Tier F.17 — The Stationarity Loop: Constancy, the Spin…]{\texorpdfstring{Tier F.17 — The Stationarity Loop: Constancy, the Spin Connection, and the No-Go}{Tier F.17 - The Stationarity Loop: Constancy, the Spin Connection, and the No-Go}}

\begin{theorem}[{\texorpdfstring{FT.53 — Geodesic constancy: the stationary on-shell ledger is constant-quantum}{FT.53 - Geodesic constancy: the stationary on-shell ledger is constant-quantum}}]\label{tag:FT-53}
\textbf{Thm.} $(1)$ \textbf{Conservation.} Along any affinely parametrized geodesic of $g_{\mu\nu}$ [CT2.6]: $\frac{d}{d\lambda}\bigl(g_{\mu\nu}\dot x^\mu\dot x^\nu\bigr)=2\,g_{\mu\nu}\dot x^\mu(\nabla_{\dot x}\dot x)^\nu=0$ [math] — the sampled quantity $\kappa(\lambda)=g_{\mu\nu}\dot x^\mu\dot x^\nu$ is constant along the path. $(2)$ \textbf{Constancy of the stationary on-shell ledger.} In CT1.19's regime, materialized chains contributing to the semiclassical spectrum sit at $\delta\Phi=0$ [CT1.20: $\delta\Phi\ne0$ cancels pairwise; $\delta\Phi=0$ dominates] $=$ geodesy [FT.48(1)]; CT1.19's sampling $\kappa_e=g_{\mu\nu}\dot x^\mu\dot x^\nu(\lambda_e)+O(\mathfrak h)$ then gives $\kappa_e=\kappa_0+O(\mathfrak h)$ with a single constant $\kappa_0$; on integrality-class edges [FT.24/FT.34(4): $\kappa_e\in\mathbb Z$] and small $\mathfrak h$: $\kappa_e=\kappa_0\in\mathbb Z$ \textbf{exactly} — FT.38(2)'s rigidity strengthened from eventual-local to exact-global on the stationary locus, without the integer-step $\varepsilon$-argument. $(3)$ \textbf{Equality saturation.} On stationary on-shell channels: FT.33(4)'s Cauchy–Schwarz inequality and FT.48(2)'s two-sided law hold with \textbf{equality}: $\Lambda_c=L\sqrt{\kappa_0}$, $\ell^2=m\,d_{12}\,\mathfrak h^2$, $\overline{\sqrt\kappa}^{\,2}=\kappa_0$ — the FT.33 bridge ratio becomes a single-integer function: $m_{\mathrm{pole}}^2/m_{\mathrm{hid}}^2=\hat\mu^2\mathfrak h^4G(c)\,\kappa_0/(\ln2)^2$. $(4)$ \textbf{Semiclassical exclusion of class (iii).} FT.38(3)'s rational non-uniform ladders ($\kappa_e\in(2/L)\mathbb Z$, varying) require $\kappa(\lambda)$ non-constant along the chain $\Rightarrow$ [contrapositive of (1)] non-geodesic support $\Rightarrow$ off-stationary $\Rightarrow$ pairwise cancellation in $Z$ [CT1.20] $\Rightarrow$ no semiclassical spectral contribution.

\textbf{Req.} FT.22, FT.24, FT.31, FT.33, FT.34, FT.38, FT.48; CT1.12, CT1.19, CT1.20, CT2.6 (Part 2); D1.6; geodesic conservation, contrapositive logic [math].

\textbf{Pf.}

%dsp-level:0
\[
(1):\ \tfrac{d}{d\lambda}g_{\mu\nu}\dot x^\mu\dot x^\nu
\]

%dsp-level:0
\[
= (\partial_\rho g_{\mu\nu})\dot x^\rho\dot x^\mu\dot x^\nu+2g_{\mu\nu}\dot x^\mu\ddot x^\nu
\]

%dsp-level:0
\[
= 2g_{\mu\nu}\dot x^\mu(\ddot x^\nu+\Gamma^\nu{}_{\rho\sigma}\dot x^\rho\dot x^\sigma)
\]

%dsp-level:0
\[
= 2g_{\mu\nu}\dot x^\mu(\nabla_{\dot x}\dot x)^\nu
\]

%dsp-level:0
\[
= 0\ [\text{math};\ \Gamma\ \text{from}\ \mathrm{CT1.13}]
\]

%dsp-level:0
\[
(2):\ \mathrm{CT1.20}:\ \text{spectral dominance at}\ \delta\Phi
\]

%dsp-level:0
\[
= 0;
\]

%dsp-level:0
\[
\quad\mathrm{FT.48}(1):\ \delta\Phi
\]

%dsp-level:0
\[
= 0\iff\text{geodesy (on-shell)};
\]

%dsp-level:0
\[
\quad\mathrm{CT1.19}\ \text{Pf}:\ \kappa_e
\]

%dsp-level:0
\[
= g(\dot x,\dot x)(\lambda_e)+O(\mathfrak h)\xrightarrow{(1)}\kappa_0+O(\mathfrak h);
\]

%dsp-level:0
\[
\quad\kappa_e\in\mathbb Z\wedge|\kappa_e-\kappa_0|<\tfrac12\ (\mathfrak h\ \text{small})\implies\kappa_e
\]

%dsp-level:0
\[
= \kappa_0\ \forall e
\]

%dsp-level:0
\[
(3):\ \kappa_e\equiv\kappa_0:\ \Lambda_c
\]

%dsp-level:0
\[
= \textstyle\sum_e\sqrt{\kappa_0}
\]

%dsp-level:0
\[
= L\sqrt{\kappa_0};
\]

%dsp-level:0
\[
\quad\bigl(\textstyle\sum\sqrt{\kappa_e}\bigr)^2
\]

%dsp-level:0
\[
= L^2\kappa_0
\]

%dsp-level:0
\[
= L\sum_e\kappa_e:\ \text{Cauchy–Schwarz equality}\ [\mathrm{FT.33}(4)];
\]

%dsp-level:0
\[
\quad\text{connector of}\ m\ \text{channels}:\ d_{12}
\]

%dsp-level:0
\[
= m\kappa_0,
\]

%dsp-level:0
\[
\ \ell
\]

%dsp-level:0
\[
= \mathfrak h\,
\]

%dsp-level:0
\[
m\sqrt{\kappa_0}\implies\ell^2
\]

%dsp-level:0
\[
= \mathfrak h^2m^2\kappa_0
\]

%dsp-level:0
\[
= m\,
\]

%dsp-level:0
\[
d_{12}\,
\]

%dsp-level:0
\[
\mathfrak h^2\ [\mathrm{FT.48}(2)\ \text{equality}];
\]

%dsp-level:0
\[
\quad\overline{\sqrt\kappa}^{\,2}\qquad = \kappa_0
\]

%dsp-level:0
\[
(4):\ \kappa\ \text{non-constant along}\ c\overset{(1)}{\implies}\neg\text{geodesic}\overset{\mathrm{FT.48}(1)}{\implies}\neg\text{stationary}\overset{\mathrm{CT1.20}}{\implies}\text{pairwise cancellation};
\]

%dsp-level:0
\[
\qquad \text{residue}:\ \text{finite-}\mathfrak h\ \text{non-stationary poles}\ [\text{Open Gaps 13}/18]
\]

%dsp-level:0
\[
\text{bridge ratio at}\ (3):\ \frac{m_{\mathrm{pole}}^2}{m_{\mathrm{hid}}^2}
\]

%dsp-level:0
\[
= \frac{\hat\mu^2\mathfrak h^4G(c)}{(\ln2)^2}\,
\]

%dsp-level:0
\[
\overline{\sqrt\kappa}^{\,2}\ [\mathrm{FT.33}(2)]\xrightarrow{\kappa\equiv\kappa_0}\frac{\hat\mu^2\mathfrak h^4G(c)\,\kappa_0}{(\ln2)^2}.\ \blacksquare
\]

\textbf{Rmk (what became exact).} Three inequalities/qualifiers of the earlier passes were artifacts of allowing non-stationary profiles: FT.33's "$\overline{\sqrt\kappa}^{\,2}$ (CS-maximal at uniform)", FT.48(2)'s "$\le$", and FT.38(2)'s "eventually locally constant". On the stationary on-shell locus all three are equalities/identities: uniformity is not generic — it is \textit{forced} by geodesy plus conservation.

\textbf{Phys.} The classical limit does more than make windings large: it makes them rigid. A materialized chain that survives to the continuum is a geodesic, and a geodesic samples one number — its constant quantum $\kappa_0$ — so the ledger of the stationary sector is a single integer per chain, the distance laws become equalities, and the exotic rational ladders live only in the cancelling off-stationary foam.

\vspace{2pt}\noindent\rule{\textwidth}{0.4pt}\vspace{2pt}

\end{theorem}

\begin{definition}[{\texorpdfstring{Def-H.21 — The spin connection on the emergent geometry (construction)}{Def-H.21 - The spin connection on the emergent geometry (construction)}}]\label{tag:Def-H-21}
\textbf{Def.} On a spatial slice of CT2.6's $(M,g)$ [CT1.12: $\bar g_{ij}\succ0$ perturbatively], within a convex normal neighborhood [FT.47(3)] or, globally, on any orientable slice [math: orientable 3-manifolds are spin]:

%dsp-level:0
\[
(1)\quad\text{orthonormal frames}\ e_a=e_a{}^\mu\partial_\mu,\quad g_{\mu\nu}e_a{}^\mu e_b{}^\nu=\delta_{ab}\ (\text{spatial});
\]

%dsp-level:0
\[
\qquad (2)\quad\text{Levi-Civita connection}\ \Gamma^\lambda{}_{\mu\nu}\ [\mathrm{CT1.13}]\ \to\ \text{spin connection}\ \omega_\mu{}^{ab}
\]

%dsp-level:0
\[
= e^a{}_\lambda\nabla_\mu e^{b\lambda}\ [\text{math}];
\]

%dsp-level:0
\[
(3)\quad SU(2)\to SO(3)\ \text{double cover}\ [\text{math}]\ \to\ \mathfrak{su}(2)\text{-connection}\ \tilde\omega;
\]

%dsp-level:0
\[
\qquad (4)\quad\text{carrier}:\ \mathcal H_{\mathsf v}\cong\mathbb C^2\ [\mathrm{EW.32}],
\]

%dsp-level:0
\[
\ \text{with}\ W_i\mapsto\tfrac12\sigma_i\ [\mathrm{EW.32}/\mathrm{CT1.2}:\ \text{the corpus's own basis match}];
\]

%dsp-level:0
\[
(5)\quad\text{spinor holonomy around a closed loop}\ c:\ \mathcal U^{\mathrm{spin}}_c=\mathrm P\exp\bigl(i\!\oint_c\tilde\omega\bigr)\in SU(2);
\]

%dsp-level:0
\[
\qquad (6)\ \text{frame rotation angle}\ \Theta_R(c):\ \text{the}\ SO(3)\ \text{angle of the projected holonomy};
\]

%dsp-level:0
\[
\ \text{for}\ c
\]

%dsp-level:0
\[
= \partial S,
\]

%dsp-level:0
\[
\ S\ \text{disk-like in a constant-curvature slice}:\ \Theta_R
\]

%dsp-level:0
\[
= \iint_S K\,
\]

%dsp-level:0
\[
dA\ [\text{Gauss–Bonnet};\ \text{math}].
\]

\textbf{No physical identification is made here}: $\mathcal U^{\mathrm{spin}}$ is \textit{not} asserted to be EW.33's weak transport $\mathcal U$ (couplings $c_A$, P4 data) — that identification is the declared content of Def-H.37(2), and the obstruction to deriving it is Open Gaps 19. The construction is [math]-class over derived corpus objects [CT1.13, CT2.6, EW.32]; it adds no premise. A5/C11's node-level $SU(2)$ action is the corpus precedent for $SU(2)$ carriers.

\vspace{2pt}\noindent\rule{\textwidth}{0.4pt}\vspace{2pt}

\end{definition}

\begin{definition}[{\texorpdfstring{Def-H.37 — The defect-transfer interface}{Def-H.37 - The defect-transfer interface}}]\label{tag:Def-H-37}
\textbf{Def.} The pole-side mass read of a materialized bound chain. Two clauses: $(1)$ \textbf{Leading-order transfer.} For $c$ a materialized \textbf{spacelike} bound chain with weak-transport profile $c_A$ over the seed generator set [EW.33's $c_A$; EW.32's seed $\mathsf v\in\mathcal H_{\mathsf v}$], the chain's pole displacement takes EW.33's defect form for the seed weak-transported around $c$:

%dsp-level:0
\[
m^2_{\mathrm{pole}}(c)=\hat\mu^2\,\ell_c^2\,G(c),
\]

%dsp-level:0
\[
\qquad G(c)=\sum_{A,B}c_Ac_B\operatorname{Re}\langle\mathcal G_A\mathsf v,\mathcal G_B\mathsf v\rangle\ \ [\mathrm{EW.33}\ \text{Gram data};\ \mathrm{EW.20}],
\]

with $\hat\mu^2>0$ the floor-independent residue packaging $\varkappa^2\beta^2h_0^2$ (EW.33's $c_W$ with $\ell^2$ factored out), $[\hat\mu^2]=\mathfrak q^2$ [EW.33 packaging; D2.9; Cal-F], and uniform coupling profile along the chain (EW.33's own leading-order regime). $(2)$ \textbf{Holonomy refinement.} The transfer of $(1)$ is Def-H.21's spin-connection holonomy: $\mathcal U_c=\mathcal U^{\mathrm{spin}}_c$ with frame angle $\Theta_R(c)$ [Def-H.21(5)–(6)]; then exactly

%dsp-level:0
\[
m^2_{\mathrm{pole}}(c)=\hat\mu'^2\cdot4\sin^2\!\Bigl(\frac{\Theta_R(c)}{4}\Bigr),\qquad[\hat\mu'^2]=\mathfrak q\ \ [\mathrm{FT.54}(4)\ \text{packaging pattern};\ \text{Cal-F}],
\]

and clause $(2)$ implies clause $(1)$ at leading order under the matching $\hat\mu'^2\Theta_R^2/4=\hat\mu^2\ell_c^2G(c)$ — recorded, not imposed [FT.54(2), FT.33]. \textbf{Interface status.} The non-derived content of this definition is exactly: the extension of EW.33's defect read from wrapping-chain gauge modes [EW.33's locus] to materialized bound-chain termini, the holonomy identification of $(2)$, and the two P4-class residues ($\hat\mu^2$, $\hat\mu'^2$; Cal-F). All further content is forced: the exact state- and axis-independent shape $4\sin^2(\cdot/4)$ [FT.54(1)]; EW.33 as its Taylor term [FT.54(2)]; the exclusion of rational-coupling combinatorial replacements [FT.73(2)–(3) — forward pointer, Tier F.24]; the range ceiling $[0,4\hat\mu'^2]$ [FT.76(3) — forward pointer, Tier F.26]. Minimality audit: FT.155. \textbf{Locus clause.} No identification $m_{\mathrm{pole}}=m_{\mathrm{hid}}$ is asserted anywhere; the read registry and its bridge slots: Def-H.39. On rung-3 wrapping chains clause (1) is EW.33 verbatim — the transfer adds nothing there [FT.33 Rmk].

\end{definition}

\begin{theorem}[{\texorpdfstring{FT.54 — Exact holonomy-defect formula and the closed-form pole side (under Def-H.37(2))}{FT.54 - Exact holonomy-defect formula and the closed-form pole side (under Def-H.37(2))}}]\label{tag:FT-54}
\textbf{Thm.} $(1)$ \textbf{Exact SU(2) defect (state- and axis-independent).} For $U=\cos\phi\,\mathbb I+i\sin\phi\,\hat n\cdot\vec\sigma\in SU(2)$ and any unit $v\in\mathbb C^2$:

%dsp-level:0
\[
(U-1)^\dagger(U-1)=\bigl(2-2\cos\phi\bigr)\mathbb I=4\sin^2\!\bigl(\tfrac\phi2\bigr)\,\mathbb I
\]

%dsp-level:0
\[
\qquad \implies
\]

%dsp-level:0
\[
\qquad \|(U-1)v\|^2=2-2\cos\phi\ [\text{math; exact Pauli algebra}].
\]

$(2)$ \textbf{EW.33 recovered as Taylor term.} Any traceless Hermitian $X=\vec x\cdot\vec\sigma$ satisfies $X^2=|\vec x|^2\mathbb I$; with $U=e^{iX}$, $\phi=|\vec x|$: defect $=2-2\cos|\vec x|=|\vec x|^2+O(|\vec x|^4)=\|Xv\|^2+O(X^4)$ — EW.33's leading-order defect $\ell_\gamma^2\|\sum_Ac_A\mathcal G_A\mathsf v\|^2$ is the expansion of the same exact quantity [consistency derived, not assumed]. $(3)$ \textbf{Spin-connection holonomy.} $\mathcal U^{\mathrm{spin}}_c=U(\phi=\Theta_R(c)/2)$ [Def-H.21(5)–(6); double cover: frame angle $\Theta_R$, spinor angle $\Theta_R/2$; sign convention irrelevant — $\cos$ even] $\Rightarrow$ defect$_{\mathrm{spin}}(c)=2-2\cos\bigl(\tfrac{\Theta_R(c)}{2}\bigr)=4\sin^2\bigl(\tfrac{\Theta_R(c)}{4}\bigr)$. $(4)$ \textbf{Closed-form pole side (Def-H.37(2)).} Under Def-H.37(2) (clause (1)'s defect transport is the spin-connection holonomy):

%dsp-level:0
\[
m_{\mathrm{pole}}^2(c)=\hat\mu'^2\cdot4\sin^2\!\Bigl(\frac{\Theta_R(c)}{4}\Bigr),\qquad[\hat\mu'^2]=[m^2]=\mathfrak q\ [\mathrm{EW.33}\ \text{packaging pattern}];
\]

for geodesic polygonal loops bounding disk-like $S$ in constant-curvature slices: $\Theta_R=\bar KA=\sum_i\alpha_i-(n-2)\pi$ [Gauss–Bonnet; FT.47(1)] — mass from the angle defect. Bridge ratio (Gap 7's well-posed object):

%dsp-level:0
\[
\frac{m_{\mathrm{pole}}^2}{m_{\mathrm{hid}}^2}=\frac{4\,\hat\mu'^2\,\mathfrak h^2\,\sin^2(\Theta_R/4)}{E_{\mathrm{int}}^2\,(\ln2)^2}\in\mathfrak q^0
\]

%dsp-level:0
\[
\qquad [\text{dimension check}:\ [\hat\mu'^2\mathfrak h^2]=\mathfrak q\cdot\mathfrak q^{-1}=\mathfrak q^0].
\]

$(5)$ \textbf{Exact witness.} Unit-sphere octant triangle ($n=3$, $\alpha_i=\pi/2$, $A=\pi/2$): $\Theta_R=\pi/2$; defect $=2-2\cos(\pi/4)=2-\sqrt2\in\mathbb Q(\sqrt2)$ — an exact irrational ledger value, in the corpus's own radical-arithmetic class [checks/README pattern].

\textbf{Req.} Def-H.37(2) (clause 4 only); Def-H.21; EW.32, EW.33 (Part 2); FT.47(1); FT.33, FT.42(1), FT.45(1) ($m_{\mathrm{hid}}$ forms); D2.9, CT2.9 (\textbf{C}, units); Pauli algebra, half-angle identities, Gauss–Bonnet [math].

\textbf{Pf.}

%dsp-level:0
\[
(1):\ U^\dagger U=\mathbb I\implies(U-1)^\dagger(U-1)=2\mathbb I-U-U^\dagger;
\]

%dsp-level:0
\[
\qquad U+U^\dagger=2\cos\phi\,\mathbb I\ (\text{since}\ i\sin\phi\,\hat n\cdot\vec\sigma\ \text{is antihermitian})\implies(2-2\cos\phi)\mathbb I
\]

%dsp-level:0
\[
\|(U-1)v\|^2=\langle v|(U-1)^\dagger(U-1)|v\rangle=(2-2\cos\phi)\|v\|^2=2-2\cos\phi;
\]

%dsp-level:0
\[
\qquad 2-2\cos\phi=4\sin^2(\phi/2)\ [\text{half-angle};\ \text{math}]
\]

%dsp-level:0
\[
(2):\ X=\vec x\cdot\vec\sigma:\ X^2=(\vec x\cdot\vec x)\mathbb I\ [\text{Pauli};\ \text{math}];
\]

%dsp-level:0
\[
\quad e^{iX}=\cos|x|\,\mathbb I+i\sin|x|\,\hat x\cdot\vec\sigma\implies\text{defect}=2-2\cos|x|=|x|^2+O(|x|^4)=\|Xv\|^2+O(X^4)\ [\|v\|=1]
\]

%dsp-level:0
\[
\mathrm{EW.33}:\ \|(U-1)\mathsf v\|^2
\]

%dsp-level:0
\[
= \ell_\gamma^2\|\textstyle\sum c_A\mathcal G_A\mathsf v\|^2+O(c^3)
\]

%dsp-level:0
\[
= \text{the same expansion with}\ X
\]

%dsp-level:0
\[
= \ell_\gamma\textstyle\sum c_A\mathcal G_A
\]

%dsp-level:0
\[
(3):\ \text{Def-H.21}(5)\text{--}(6):\ \phi=\Theta_R/2;
\]

%dsp-level:0
\[
\qquad (4):\ \text{Def-H.37(2)}\ \text{substitution};\quad\Theta_R=\bar KA=\sum\alpha_i-(n-2)\pi\ [\mathrm{FT.47}(1),\ \text{Gauss–Bonnet}];
\]

%dsp-level:0
\[
\qquad \text{ratio}:\ m_{\mathrm{hid}}^2=E_{\mathrm{int}}^2(\ln2)^2/\mathfrak h^2\ [\mathrm{FT.42}(1)\ \text{template}]
\]

%dsp-level:0
\[
(5):\ \sum\alpha=3\cdot\tfrac\pi2,\ (n-2)\pi=\tfrac\pi2\implies\Theta_R=\tfrac\pi2;
\]

%dsp-level:0
\[
\qquad 2-2\cos\tfrac\pi4=2-\sqrt2;
\]

%dsp-level:0
\[
\qquad (2-\sqrt2)(2+\sqrt2)=2\ \checkmark\ [\mathbb Q(\sqrt2)\ \text{exact}].\ \blacksquare
\]

\textbf{Rmk (Def-H.37(1) vs Def-H.37(2)).} Def-H.37(1) (leading order, weak-coupling profile $c_A$) and Def-H.37(2) (exact, spin-connection) agree to leading order when $\hat\mu'^2\Theta_R^2/4=\hat\mu^2\ell_c^2G(c)$ — a matching condition on the P4-class packaging (Cal-F), recorded not imposed.

\textbf{Phys.} If the fermion's pole mass is the seed's failure to return to itself under frame transport, then it is a pure trigonometric function of the curvature flux the loop encloses — $4\sin^2(\Theta_R/4)$, exact, axis-free, state-free — and on geodesic polygons the curvature flux is the angle defect, the oldest invariant in geometry. The octant's mass ledger entry is $2-\sqrt2$: the first exact irrational number in the program's cards.

\vspace{2pt}\noindent\rule{\textwidth}{0.4pt}\vspace{2pt}

\end{theorem}

\begin{theorem}[{\texorpdfstring{FT.55 — The no-go and the surviving spin realization}{FT.55 - The no-go and the surviving spin realization}}]\label{tag:FT-55}
\textbf{Thm.} $(0)$ \textbf{Admissibility (correction).} Throughout this tier and FT.52, the admissible spin label is $j=|k|/2$ with carrier $\mathrm{Sym}^{|k|}$: $\mathfrak{su}(2)$ labels require $j\ge0$ and the Casimir $j(j{+}1)$ distinguishes $\pm k/2$; stretched additivity is confined to the nonnegative cone (spacelike-dominant compositions) [FT.52(3) patched form]. $(1)$ \textbf{Calibration no-go.} The identification $\Theta_R(c)=2\pi k(c)$ (frame-holonomy angle $=$ winding $\times2\pi$) is \textbf{ruled out} as a corpus relation: $k=\oint ds^2/(2\mathfrak h^2)$ carries $\mathfrak h$-weight $-2$ [FT.31(3); D2.9], while $\Theta_R=\iint_S K\,dA$ is $\mathfrak h$-weight $0$ for fixed classical regions [Def-H.21(6)]; a proportionality between quantities of different calibration weight is not $\mathfrak h$-free and hence not a derived pure-number relation [CT2.11's criterion]. Consequently \textbf{the spin label cannot be read as the spin-connection holonomy angle of the same loop}. $(2)$ \textbf{The surviving realization.} FT.52's unique rank-1 additive rule $R(k)=k/2$ has exactly one natural representation-theoretic realization: $j=|k|/2\iff$ carrier $\mathrm{Sym}^{|k|}(\mathcal H_{\mathsf v})$ — the $(|k|{+}1)$-dimensional irreducible $\mathfrak{su}(2)$ representation [math: classification; $k\ge0$ on the cone of clause (0)]. It satisfies every derived constraint: anchor $\mathrm{Sym}^1=\mathcal H_{\mathsf v}$, $j=\tfrac12$ [EW.32]; composition $\mathrm{Sym}^{k_1}\otimes\mathrm{Sym}^{k_2}\supseteq\mathrm{Sym}^{k_1+k_2}$ as maximal weight [math: Clebsch–Gordan] $=$ FT.34(3)'s additive winding law $k(\Gamma)=k_1+k_2+d_{12}$ on the stretched channel; parity triviality respected (the rule reads $k$, not $\mathrm{Hol}_{\mathrm{int}}$) [FT.52(1)]; half-integrality from odd $k$ [FT.52(3)]. $(3)$ \textbf{Scope consistency.} Elementary floors carry small $k$ [FT.28: $k=\kappa=1$ floor; FT.45/FT.42 minima: $k=O(1)$] $\Rightarrow$ small $j$; the large-$k$ regime is the classical-geometry limit [FT.31(3)], not particle territory — no macroscopic-spin contradiction. $(4)$ \textbf{Residual, relocated.} What remains missing for promotion is neither the connection (Def-H.21 constructs it) nor the algebra (FT.52 settles it): it is a \textbf{ledger mechanism assigning $\mathrm{Sym}^k(\mathcal H_{\mathsf v})$ to a materialized chain of winding $k$} — a $k$-fold symmetric seed-tensoring tied to the winding count. EW.32 transports \textit{one} qubit around the chain; nothing in Parts 1–3 tensors it $k$ times.

\textbf{Req.} FT.22, FT.28, FT.31, FT.34, FT.45, FT.52; Def-H.21; D2.9; CT2.11 (Part 2); EW.32 (Part 2); $\mathfrak{su}(2)$ representation classification, Clebsch–Gordan maximal weight [math].

\textbf{Pf.}

%dsp-level:0
\[
(1):\ k=\Theta/2,\ \Theta\mathfrak h^2=\oint ds^2\ [\mathrm{FT.22}/\mathrm{FT.31}(3)]\implies[k]=[\mathfrak h^{-2}]=\mathfrak q^{+1}\Rightarrow k\ \text{has}\ \mathfrak h\text{-weight}-2;
\]

%dsp-level:0
\[
\qquad [\Theta_R]=\mathfrak h^0\ [\text{Def-H.21}(6):\ K\,dA\ \text{classical}]
\]

%dsp-level:0
\[
\Theta_R
\]

%dsp-level:0
\[
= 2\pi k\ \text{would require a constant of weight}\ +2\ \text{in}\ \mathfrak h:\ \neg\text{pure number}\implies\neg\text{derived relation}\ [\mathrm{CT2.11}:\ \text{rescaling}\ \mathfrak h\mapsto\lambda\mathfrak h\ \text{is a symmetry};\ \text{predictions are}\ \mathfrak h\text{-free ratios}]
\]

%dsp-level:0
\[
(2):\ \dim\mathrm{Sym}^k(\mathbb C^2)=k+1;\quad\text{highest weight}=k\cdot\tfrac12=\tfrac k2\ [\text{math}];
\]

%dsp-level:0
\[
\qquad \mathrm{Sym}^{k_1}\otimes\mathrm{Sym}^{k_2}=\bigoplus_{j=|k_1-k_2|/2}^{(k_1+k_2)/2}\mathcal D_j\supset\mathrm{Sym}^{k_1+k_2}\ [\text{CG};\ \text{math}]
\]

%dsp-level:0
\[
\text{anchor}:\ k=1:\ \mathrm{Sym}^1=\mathcal H_{\mathsf v}\cong\mathbb C^2,\ j=\tfrac12\ [\mathrm{EW.32}];
\]

%dsp-level:0
\[
\qquad \text{constraints}:\ \mathrm{FT.52}(1)\text{--}(3)\ \text{each verified above}
\]

%dsp-level:0
\[
(3):\ \text{floor minima}:\ k\in\{1,2,3\}\ [\mathrm{FT.28}(2),\ \mathrm{FT.45}(4),\ \mathrm{FT.42}(5)\ \text{uniform minima}]
\]

%dsp-level:0
\[
\implies j\in\{\tfrac12,1,\tfrac32\};
\]

%dsp-level:0
\[
\qquad k\to\infty:\ \text{classical regime}\ [\mathrm{FT.31}(3)]
\]

%dsp-level:0
\[
(4):\ \text{inventory}:\ \text{connection}\ [\text{Def-H.21}]\ \checkmark;
\]

%dsp-level:0
\[
\ \ \text{algebra}\ [\mathrm{FT.52}]\ \checkmark;
\]

%dsp-level:0
\[
\ \ \text{tensoring mechanism}:\ \neg\exists\ \text{tag}\ [\text{Parts 1–3 inventory}].\ \blacksquare
\]

\textbf{Rmk (why the no-go is progress).} Before this tier, "read spin from holonomy" was an open direction among several. The calibration argument closes it with a weight mismatch — the same kind of dimensional discipline that bars the malformed comparison. What survives is sharper and smaller: one construction question, with its target representation already unique.

\textbf{Phys.} The geometry refuses to hand over spin directly: the frame's rotation angle and the ledger's winding number live at different calibration weights, and no constant can bridge them. What the ledger can do is count copies — a winding-$k$ chain as $k$ symmetric copies of the seed qubit, spin $k/2$, combining exactly as windings compose. The rule is unique, the space is unique, the anchor is the corpus's own seed; only the copying machine is missing.

\vspace{2pt}\noindent\rule{\textwidth}{0.4pt}\vspace{2pt}

\end{theorem}

\phantomsection
\section[Tier F.18 — Native Rungs, Boundary Exchange Shape, and the…]{\texorpdfstring{Tier F.18 — Native Rungs, Boundary Exchange Shape, and the Admissibility Correction}{Tier F.18 - Native Rungs, Boundary Exchange Shape, and the Admissibility Correction}}

\begin{theorem}[{\texorpdfstring{FT.56 — Native ladder rungs, the spin–orbit decomposition, and the missing map}{FT.56 - Native ladder rungs, the spin-orbit decomposition, and the missing map}}]\label{tag:FT-56}
\textbf{Thm.} $(1)$ \textbf{The bottom rungs are native.} The candidate spin carriers for $j=0,\tfrac12,1$ are existing corpus objects with the correct dimensions $2j{+}1$ and Casimir labels:

%dsp-level:0
\[
j=0:\ \wedge^2\mathcal H_q\cong\mathrm{Sym}^0\ (\text{trivial},\ 1\text{-dim})\ —\ \text{the reality line}\ [\mathrm{A3},\ \mathrm{CT0.7}(1):\ \Pi_{\mathcal V}=\tfrac14(E_0-\Sigma),\ \mathfrak T\Pi_{\mathcal V}=1;\ \text{Casimir}\ 0\ [\mathrm{CT0.7a}];\ \sigma_e\ \text{lives here}\ [\mathrm{CT0.9}]];
\]

%dsp-level:0
\[
j=\tfrac12:\ \mathcal H_{\mathsf v}\cong\mathbb C^2,\ W_i\mapsto\tfrac12\sigma_i\ —\ \text{the transported seed}\ [\mathrm{EW.32}];
\]

%dsp-level:0
\[
j = 1:\ \vee^2\mathcal H_q
\]

%dsp-level:0
\[
= \mathrm{Sym}^2(\mathcal H_q)\ (3\text{-dim})\ —\ \text{the pair channel's reality sector}\ [\mathrm{A3}:\ \mathcal H_p\otimes\mathcal H_q=\vee^2\oplus\wedge^2;\ \mathrm{CT0.7}(1):\ \Pi_{\mathcal R}=\tfrac14(3E_0+\Sigma),\ \mathfrak T\Pi_{\mathcal R}=3;\ \mathrm{CT0.7a}:\ \textstyle\sum_iW_i^2\big|_{\vee^2}=2=1(1{+}1)];
\]

%dsp-level:0
\[
\text{and the reality grading is the placement rule}:\ \text{real edges carry}\ \vee^2\text{-weight}\ [\mathrm{A4}:\ |\Pi_{\mathcal V}\rho_{pq}|\le\varepsilon_q]\ —\ \text{the symmetric square is the corpus's own}\ \textbf{real}\ \text{sector}.
\]

$(2)$ \textbf{Spin–orbit decomposition of the additive law.} $k(\Gamma)=k_1+k_2+d_{12}$ [FT.34(3)] reads, under the stretched rule on the nonnegative cone [FT.52(3), corrected form]: $j(\Gamma)=j(c_1)+j(c_2)+j_{\mathrm{rel}}$, $j_{\mathrm{rel}}=d_{12}/2$ — the connector's distance quanta contribute as a \textit{relative} (orbital-type) term, the floors' own windings as intrinsic terms. On the cone this is exactly maximal-weight Clebsch–Gordan composition [math]; $\Theta$'s orientation-blindness [FT.23] makes $k$ (hence $j_{\mathrm{rel}}$) a \textbf{signature datum} — negative only for timelike-dominant connectors, where the cone condition fails and the read is out of scope. $(3)$ \textbf{The missing map, stated exactly.} EW.32 transports \textit{one} qubit around a materialized chain, regardless of its winding; no tag in Parts 1–3 (inventory) couples a $\Theta$-type ledger count to carrier multiplicity. What any candidate map winding $\to$ carrier must satisfy: (i) $k{=}1\mapsto\mathcal H_{\mathsf v}$ [EW.32 anchor]; (ii) $k{=}2$ must land in the channel's own reality sector $\vee^2=\mathrm{Sym}^2$ [clause (1) — the rung exists natively; the map must select it]; (iii) nonnegative-cone additivity [clause (2)]; (iv) admissibility $j=|k|/2$ [admissibility correction]. Gap 19 is exactly the construction of a map satisfying (i)–(iv) — or a no-go on such maps, which would reopen FT.52's constraint list to higher-rank rules.

\textbf{Req.} A3, A4; CT0.7, CT0.7a, CT0.9, CT1.2 (Part 1); EW.32 (Part 2); FT.13, FT.23, FT.34, FT.52, FT.55; Def-H.8; $\wedge^2\cong\mathrm{Sym}^0$, $\vee^2\cong\mathrm{Sym}^2$ of $\mathbb C^2$, maximal-weight CG [math].

\textbf{Pf.}

%dsp-level:0
\[
(1):\ \dim\wedge^2=\tbinom{2}{2}=1,\ \dim\vee^2=3\ [\mathrm{A3};\ \mathrm{CT0.7}(1):\ \mathfrak T\Pi_{\mathcal V}=1,\ \mathfrak T\Pi_{\mathcal R}=3];
\]

%dsp-level:0
\[
\qquad \textstyle\sum_iW_i^2=\tfrac32E_0+\tfrac12\Sigma\ [\mathrm{CT0.7a}]
\]

%dsp-level:0
\[
\Sigma=2\mathbb S-E_0\ [\mathrm{CT0.7}(1):\ \mathbb S=\tfrac12(E_0+\Sigma)]\implies\textstyle\sum_iW_i^2\big|_{\vee^2}=\tfrac32+\tfrac12(2\cdot1-1)=2;
\]

%dsp-level:0
\[
\qquad \big|_{\wedge^2}=\tfrac32+\tfrac12(2(-1)-1)=0
\]

%dsp-level:0
\[
\qquad [\text{machine}:4\times4\ \text{exact}]
\]

%dsp-level:0
\[
2=1(1{+}1)\implies j=1\ \text{on}\ \vee^2;\quad0=0(0{+}1)\implies j=0\ \text{on}\ \wedge^2\ [\mathrm{CT0.7a}\ \text{verbatim}];
\]

%dsp-level:0
\[
\qquad \mathcal H_{\mathsf v}:\ j=\tfrac12\ [\mathrm{EW.32}];
\]

%dsp-level:0
\[
\qquad \dim=2j{+}1:\ 1,2,3\ \checkmark
\]

%dsp-level:0
\[
(2):\ \mathrm{FT.34}(3)\ \text{additivity}+\ \text{stretched rule on the cone}\ [\mathrm{FT.52}(3)\ \text{form}]:\ j(\Gamma)
\]

%dsp-level:0
\[
= \tfrac{k_1+k_2+d_{12}}{2}
\]

%dsp-level:0
\[
= j_1+j_2+j_{\mathrm{rel}};
\]

%dsp-level:0
\[
\qquad \mathrm{Sym}^{k_1}\otimes\mathrm{Sym}^{k_2}\supseteq\mathrm{Sym}^{k_1+k_2}\ [\text{maximal weight};\ \text{math}]
\]

%dsp-level:0
\[
\text{sign of}\ k:\ \Theta
\]

%dsp-level:0
\[
= \sum\kappa_e\ \text{orientation-blind}\ [\mathrm{FT.23}]
\]

%dsp-level:0
\[
\implies k<0\iff\text{timelike-dominant}\ [\mathrm{CT1.12}\ \text{signature}]\ —\ \text{cone fails};
\]

%dsp-level:0
\[
\ \text{read out of scope}
\]

%dsp-level:0
\[
(3):\ \mathrm{EW.32}:\ \text{carrier}=\mathbb C^2\ \text{per chain}\ (\neg\text{per winding unit});
\]

%dsp-level:0
\[
\qquad \text{tag inventory}:\ \neg\exists\ \text{multiplicity coupling};
\]

%dsp-level:0
\[
\qquad (i)\text{--}(iv):\ \text{as stated}.\ \blacksquare
\]

\textbf{Rmk (why native rungs matter).} Before this tier, "$j=k/2$" was a rule looking for carriers. The carriers for its first three values are now located \textit{inside the frozen structure}: the virtual line, the seed, and the real sector of every pair channel — with the reality grading (A4/CT0.7) already deciding which of $\mathrm{Sym}^0/\mathrm{Sym}^2$ is "real". The missing map does not have to invent representation spaces; it has to select among existing ones.

\textbf{Phys.} The bottom of the spin ladder is not hypothetical: spin 0 is the ledger's virtual line, spin $\tfrac12$ is the seed it transports, spin 1 is the symmetric square its own reality test selects — three objects Parts 1–2 already proved, wearing the exact dimensions and Casimiers the winding rule asks for. Composite floors add their windings like angular momenta, with the channel between them contributing the relative term. What is still missing is the one mechanism that would say \textit{why} a winding-$k$ loop carries $k$ symmetric copies — and that mechanism is now the whole of the question.

\vspace{2pt}\noindent\rule{\textwidth}{0.4pt}\vspace{2pt}

\end{theorem}

\begin{theorem}[{\texorpdfstring{FT.57 — The boundary exchange algebra: root-space shape}{FT.57 - The boundary exchange algebra: root-space shape}}]\label{tag:FT-57}
\textbf{Thm.} $(1)$ \textbf{Decomposition.} For $N:=N(\Sigma)\ge1$: $\mathfrak u(N)=\mathfrak u(1)_{\mathrm{center}}\oplus\mathfrak{su}(N)$ [math; CT3.5's pattern at $N{=}3$]; relative to $\iota_\Sigma$'s factored diagonal $\bigoplus_e\mathfrak u(1)_e$ [Def-H.13: edge-wise assembly]:

%dsp-level:0
\[
\text{factored}=\text{diagonal}\ (N\ \text{dims},\ \text{incl. the centre});
\]

%dsp-level:0
\[
\qquad \text{strictly new}=\mathfrak{su}(N)\setminus\text{Cartan}=\text{root spaces}:\ (N^2{-}1)-(N{-}1)=N^2-N\ \ [\text{matches FT.41(d) exactly}].
\]

$(2)$ \textbf{Shape at $N=2$.} The strictly-new sector $\{|e_1\rangle\langle e_2|+\mathrm{h.c.},\ i(|e_1\rangle\langle e_2|-\mathrm{h.c.})\}$ [FT.41(d)] is, with the factored traceless diagonal, an $\mathfrak{su}(2)$ triplet: $E_\pm=\tfrac12(\sigma_1\pm i\sigma_2)$, $[E_+,E_-]=\sigma_3$, $[\sigma_3,E_\pm]=\pm2E_\pm$ [math; machine-exact over $\mathbb Q(i)$] — the raising/lowering (charged-current) shape of D3.10's $W^\pm=\tfrac1{\sqrt2}(W_1\mp iW_2)$ [C8: $F_W=F_{W^+}\oplus F_{W^-}\oplus F_Z$], as algebra shape, not as identification. $(3)$ \textbf{Hosting question, localized.} Which aspect block $\mathcal K^{(A)}$ the exchange modes land in under $\iota_\Sigma$ is \textbf{not determined by the current tags}: Def-H.13 defines $\iota_\Sigma$ by a factorization \textit{requirement}, not by explicit matrices; the image of the root spaces in $\mathcal K=\bigoplus_A\mathcal K^{(A)}$ needs an explicit edge-block realization of Def-H.13's assembly [CT0.12 blocks $\otimes$ $\iota_{n(w_e)}$] — one named construction away. Candidates remain: a second EM-type centre read is \textit{excluded} by (1)–(2) (the exchange sector is nonabelian root-space, not central); a $W$-type hosting matches the charged-current shape; a genuinely new collective aspect remains possible.

\textbf{Req.} Def-H.13; FT.41(d); CT0.12; CT3.5 (\textbf{F}, pattern); C8; D3.10; $\mathfrak u(N)$ decomposition, $\mathfrak{su}(2)$ commutators [math].

\textbf{Pf.}

%dsp-level:0
\[
(1):\ \mathfrak u(N)=\mathfrak{su}(N)\oplus\mathfrak u(1)_{\mathrm{center}}\ [\text{math};\ \mathrm{CT3.5}\ \text{at}\ N{=}3];
\]

%dsp-level:0
\[
\qquad \iota_\Sigma|_{|e\rangle}=(\text{CT0.12 block})\otimes\iota_{n(w_e)}\ [\mathrm{Def\text{-}H.13}]\implies\text{diagonal}\ \bigoplus_e\mathfrak u(1)_e\ \text{factored}
\]

%dsp-level:0
\[
\dim\mathfrak{su}(N)=N^2{-}1,\ \dim\text{Cartan}=N{-}1\implies\dim\text{root spaces}=N^2-N\ [\text{math}]=\text{FT.41(d)'s strictly-new count}\ \checkmark
\]

%dsp-level:0
\[
(2):\ N
\]

%dsp-level:0
\[
= 2:\ E_\pm
\]

%dsp-level:0
\[
= \tfrac12(\sigma_1\pm i\sigma_2):\ [E_+,E_-]
\]

%dsp-level:0
\[
= \tfrac14[\sigma_1+i\sigma_2,\ \sigma_1-i\sigma_2]
\]

%dsp-level:0
\[
= \tfrac14(-2i[\sigma_1,\sigma_2])
\]

%dsp-level:0
\[
= \tfrac14(-2i\cdot2i\sigma_3)
\]

%dsp-level:0
\[
= \sigma_3\ [\mathrm{CT1.2}\text{-pattern commutators};\ \text{machine over}\ \mathbb Q(i)]
\]

%dsp-level:0
\[
[\sigma_3,E_\pm]=\pm2E_\pm\ [\text{machine}];
\]

%dsp-level:0
\[
\qquad \text{shape match to}\ W^\pm\ [\mathrm{D3.10}:\ \text{same raising/lowering combinations}]
\]

%dsp-level:0
\[
(3):\ \mathrm{Def\text{-}H.13}\ \text{is a factorization requirement};
\]

%dsp-level:0
\[
\ \neg\exists\ \text{explicit matrix tag for}\ \iota_\Sigma\ [\text{inventory}]\implies\text{image decomposition in}\ \mathcal K\ \text{undetermined};
\]

%dsp-level:0
\[
\qquad \text{centre read excluded}:\ \text{exchange sector nonabelian by}\ (2).\ \blacksquare
\]

\textbf{Phys.} The first genuinely collective structure of the coarse-graining program — the thing a subnetwork can do that no single edge can — has the algebraic shape of a charged current: raising and lowering operators flipping slot excitation between boundary contacts, completing an $\mathfrak{su}(2)$ with the factored diagonal. Two contacts can exchange; one cannot. Where this exchange lives among the forces is the next named construction, not a guess: the nonabelian shape already rules out the abelian candidate.

\vspace{2pt}\noindent\rule{\textwidth}{0.4pt}\vspace{2pt}

\end{theorem}

\begin{definition}[{\texorpdfstring{Def-H.22 — The explicit boundary realization (corolla and separated patterns)}{Def-H.22 - The explicit boundary realization (corolla and separated patterns)}}]\label{tag:Def-H-22}
\textbf{Def.} For $\Sigma$ with $N:=N(\Sigma)$ boundary edges $\partial\Sigma=\{e_1,\dots,e_N\}$, $\mathcal H_{\mathrm{slot}}(\Sigma)=\operatorname{span}\{|e_i\rangle\}$ [Def-H.12], $\iota_\Sigma$ is made explicit by contact pattern: \textbf{(a) Corolla} (all boundary edges share one external vertex $w$; slots edge-labelled [A7: $\tau(v)$ records incident edges; D0.8] $\Rightarrow$ canonical inclusion $\mathcal H_{\mathrm{slot}}(\Sigma)\hookrightarrow\mathcal H^{(n(w))}_{\mathrm{slot}}$, equality iff $n(w)=N$):

%dsp-level:0
\[
\iota_\Sigma^{\mathrm{cor}}:=\iota_{n(w)}\big|_{\mathfrak u(N)}\qquad[\mathrm{D2.12}\ \text{applied to the subalgebra on the included slots}].
\]

\textbf{(scope [FT.59])}: the formula is defined exactly for $n(w)\le3$ — at $n(w)\ge4$ no realization $\iota_{n(w)}$ exists [CT4.4 Pf(b): slot algebras realized only for $n\le3$]; the corolla read then degenerates to the registration/sealing ledger ($\sigma_e$, chain-local $f_v$, $q_{\mathrm{core}}$, $\mathcal C$) plus $D^-$-absorption [D4.1, D4.2], with the collective $\mathfrak u(N)$ wholly unhosted [FT.59(4)]. \textbf{(b) Separated} (boundary edges with distinct external vertices $w_1\ne w_2\cdots$): only the factored diagonal is defined — each edge's $\mathfrak u(1)_e$ maps through its own CT0.12 block and $\iota_{n(w_e)}$'s centre [D2.12: centre $\mapsto\mathcal K^{(\mathrm{EM})}$]; since $\dim\mathcal K^{(\mathrm{EM})}=1$ [CT0.14], all edge centres share one image line:

%dsp-level:0
\[
\iota_\Sigma^{\mathrm{sep}}\Bigl(\sum_ec_eD_e\Bigr)=\Bigl(\sum_ec_e\Bigr)Y_{\mathrm{asp}}
\]

%dsp-level:0
\[
\qquad [\mathrm{EM.2}:\ \mathfrak T(Y_{\mathrm{asp}})=1];
\]

%dsp-level:0
\[
\qquad \text{off-diagonal}\ |e_i\rangle\langle e_j|\ (i\ne j):\ \neg\exists\ \text{factored image (no cross-vertex bridge in the inventory)}.
\]

No new choice is declared: (a) is the canonical edge-labelled slot inclusion composed with D2.12; (b) is forced by the 1-dimensional EM block. Every computed floor instantiates a pattern: single-contact floors [FT.28/FT.32/FT.42]: $N=1$; the face floor [FT.49]: corolla, $N=6$, $n(D)=6$ (equality) — in the void regime $n(w)\ge4$ [FT.59(4a)]: no realization, exchange sector kernel/absorbed; the two-contact $\Sigma_{\min}$ variant [FT.28 Rmk]: separated, $N=2$.

\vspace{2pt}\noindent\rule{\textwidth}{0.4pt}\vspace{2pt}

\end{definition}

\begin{theorem}[{\texorpdfstring{FT.58 — Hosting classification and the no-hosting theorem}{FT.58 - Hosting classification and the no-hosting theorem}}]\label{tag:FT-58}
\textbf{Thm.} $(1)$ \textbf{Corolla hosting follows the CT0.13 ladder at the contact vertex.} Under Def-H.22(a): the centre of $\mathfrak u(N)$ maps to $\mathcal K^{(\mathrm{EM})}$ [D2.12]; and:

%dsp-level:0
\[
n(w)=2:\quad\iota_\Sigma(\mathfrak{su}(2))\subseteq\operatorname{span}\{W_i\}=\mathcal K^{(W)}\qquad[\mathrm{CT1.4}:\ W_i\ \text{close}\ \mathfrak{su}(2);\ \mathrm{D2.12}];
\]

%dsp-level:0
\[
n(w)=3:\quad\iota_\Sigma(\mathfrak{su}(3))=\mathcal K^{(\mathrm{strong})}
\]

%dsp-level:0
\[
\qquad [\mathrm{D2.12}:\ \iota_3(\mathfrak{su}(3))\ \text{spans}\ \mathcal K^{(\mathrm{strong})};\ \mathrm{CT3.7}/\mathrm{CT3.8}:\ 8=8];
\]

%dsp-level:0
\[
\text{the exchange generators host}\ \textbf{literally}:\quad\frac{\lambda_1+i\lambda_2}{2}
\]

%dsp-level:0
\[
= E_{12}\ (3\times3\ \text{matrix unit}):\ \ |e_1\rangle\langle e_2|\ \mapsto\ \text{the}\ \lambda_1,
\]

%dsp-level:0
\[
\lambda_2\ \text{directions};
\]

%dsp-level:0
\[
\quad D_1-D_2\mapsto\lambda_3\text{-type}\ [\mathrm{CT3.9};\ \text{machine}];
\]

%dsp-level:0
\[
n(w)
\]

%dsp-level:0
\[
\ge 4:\ \textbf{resolved (FT.59)}\ —\ \neg\exists\ \iota_n\ \text{realization at all for}\ n
\]

%dsp-level:0
\[
\ge 4\ [\mathrm{CT4.4}\ \text{Pf}(b):\ \text{slot algebras realized only for}\ n\le3];
\]

%dsp-level:0
\[
\ \text{the suggested strong-hosting of}\ \mathfrak{su}(3)\ \text{subalgebras is impossible}\ [\mathrm{FT.59}(3)];
\]

%dsp-level:0
\[
\ \text{such corollas are wholly kernel}/D^-\text{-absorbed}\ [\mathrm{FT.59}(1),(4)].
\]

$(2)$ \textbf{Separated contacts: the collective part is unhosted.} Under Def-H.22(b): $\operatorname{im}\iota_\Sigma^{\mathrm{sep}}=\mathcal K^{(\mathrm{EM})}$ (the total centre read, 1-dim); $\ker\iota_\Sigma^{\mathrm{sep}}\supseteq$ traceless diagonal ($D_1-D_2$: the two edge centres coincide in the 1-dim block); the off-diagonal exchange generators ($N^2-N$ dims [FT.41(d), FT.57(1)]) have \textbf{no factored image and no derived host}: no tag bridges distinct vertex slot algebras [inventory]. Ledger at $N=2$: $4=1\ (\text{image})+1\ (\text{kernel})+2\ (\text{unhosted})$. $(3)$ \textbf{No-hosting theorem ($N\ge4$).} The observable aspect capacity is $\mathcal K_{\mathrm{pair}}^{\mathrm{obs}}\oplus\mathcal K_{\mathrm{stack}}^{\mathrm{obs}}$ with $\dim\mathcal K_{\mathrm{pair}}^{\mathrm{obs}}=5{+}1{+}3=9$ [$G,H,W$; $J^\pm,D^\pm$ never observable — FT.20] and $\dim\mathcal K_{\mathrm{stack}}^{\mathrm{obs}}=1{+}8=9$ [EM, strong; $D^-$ never — FT.20], the two commuting (distinct registers [CT0.12 Rmk]). For $N\ge4$: $\mathfrak{su}(N)$ is simple [math] with $\dim=N^2-1\ge15>9$; any Lie homomorphism $\varphi:\mathfrak{su}(N)\to\mathcal K_{\mathrm{pair}}^{\mathrm{obs}}\oplus\mathcal K_{\mathrm{stack}}^{\mathrm{obs}}$ has projections $\pi_{\mathrm{pair}}\circ\varphi$, $\pi_{\mathrm{stack}}\circ\varphi$ that are homomorphic images of a simple algebra (the factors commute) — hence trivial or injective; injective is impossible by dimension ($15>9$); therefore $\varphi\equiv0$ on $\mathfrak{su}(N)$: \textbf{the entire collective nonabelian boundary algebra maps to zero under every Lie-type-respecting hosting} [D2.12/D3.2's realization semantics] — not partially hosted: wholly kernel. The centre $\mathfrak u(1)$ still hosts in EM. $(4)$ \textbf{Capacity ceiling, exact.} Nontrivial Lie hosting of $\mathfrak{su}(N)$ exists within the current inventory only for $N\le3$ [$\dim\mathfrak{su}(3)=8\le9$; $N{=}2$: $W$-hosting, $N{=}3$: strong-hosting, both realized by (1)]. For $N\ge4$ the coarse-grained read either kernels the collective structure (FT.26's mask keeps only hosted blocks) or the aspect inventory must grow by declaration — no derivation in Parts 1–3 supplies it. The kernel \textit{is} the coarse-graining: $\dim\mathfrak u(N)=N^2$ vs $\dim\mathcal K_\Sigma\le18$ [FT.26(3)] forces $\ker\ne0$ already at $N=5$ by dimension alone.

\textbf{Req.} Def-H.13, Def-H.22; Def-H.15, Def-H.16; FT.20, FT.26, FT.41(d), FT.57; D2.12, D3.1, D3.2, D3.9; A7, D0.8; CT0.12 (Rmk), CT0.13, CT0.14 (Part 1); CT1.4 (Part 1); CT3.7, CT3.8, CT3.9 (Part 2); EM.2 (Part 2); simplicity of $\mathfrak{su}(N)$, projection argument for commuting factors [math].

\textbf{Pf.}

%dsp-level:0
\[
(1):\ \text{corolla inclusion canonical}\ [\mathrm{A7}/\mathrm{D0.8}:\ \text{slots edge-labelled}]\implies\iota_\Sigma^{\mathrm{cor}}
\]

%dsp-level:0
\[
= \iota_{n(w)}|_{\mathfrak u(N)}\ [\mathrm{Def\text{-}H.22}(a)];
\]

%dsp-level:0
\[
\qquad n(w)=2:\ \iota_2(\mathfrak{su}(2))=\operatorname{span}\{W_i\}\ [\mathrm{CT1.4},\ \mathrm{D2.12}];
\]

%dsp-level:0
\[
\quad n(w)=3:\ \iota_3(\mathfrak{su}(3))\ \text{spans}\ \mathcal K^{(\mathrm{strong})}\ [\mathrm{D2.12}],\ \dim\ 8=8\ [\mathrm{CT3.8}]
\]

%dsp-level:0
\[
\text{literal exchange hosting}:\ \lambda_1
\]

%dsp-level:0
\[
= E_{12}+E_{21},
\]

%dsp-level:0
\[
\ \lambda_2
\]

%dsp-level:0
\[
= -iE_{12}+iE_{21}\ [\mathrm{CT3.9}\ \text{explicit basis}]\implies\tfrac12(\lambda_1+i\lambda_2)
\]

%dsp-level:0
\[
= \tfrac12(E_{12}+E_{21}+E_{12}-E_{21})
\]

%dsp-level:0
\[
= E_{12}\ [\text{machine}]
\]

%dsp-level:0
\[
(2):\ \dim\mathcal K^{(\mathrm{EM})}
\]

%dsp-level:0
\[
= 1\ [\mathrm{CT0.14}]\implies\iota_\Sigma^{\mathrm{sep}}(D_1),
\]

%dsp-level:0
\[
\ \iota_\Sigma^{\mathrm{sep}}(D_2)\ \text{collinear}\ [\mathrm{EM.2}\ \text{normalization}]
\]

%dsp-level:0
\[
\implies D_1-D_2\in\ker;
\]

%dsp-level:0
\[
\qquad \text{off-diagonals}:\ \mathrm{Def\text{-}H.13}'\text{s factorization is per-edge}\ [\text{CT0.12 block}\otimes\iota_{n(w_e)}]
\]

%dsp-level:0
\[
\implies |e_1\rangle\langle e_2|\ \text{has no per-edge form};
\]

%dsp-level:0
\[
\ \neg\exists\ \text{cross-vertex bridge tag}
\]

%dsp-level:0
\[
(3):\ \mathcal K^{\mathrm{obs}}
\]

%dsp-level:0
\[
= \mathcal K_{\mathrm{pair}}^{\mathrm{obs}}\oplus\mathcal K_{\mathrm{stack}}^{\mathrm{obs}},
\]

%dsp-level:0
\[
\ [\mathcal K_{\mathrm{pair}}^{\mathrm{obs}},\mathcal K_{\mathrm{stack}}^{\mathrm{obs}}]
\]

%dsp-level:0
\[
= 0\ [\mathrm{CT0.12}\ \mathrm{Rmk}:\ \text{distinct registers}];
\]

%dsp-level:0
\[
\quad\dim
\]

%dsp-level:0
\[
= 9,
\]

%dsp-level:0
\[
\ 9\ [\mathrm{CT0.12}\ \text{block table};\ \mathrm{FT.20}:\ J^\pm,D^\pm,D^-\ \text{never}]
\]

%dsp-level:0
\[
\varphi:\ \mathfrak{su}(N)\to\mathcal K^{\mathrm{obs}}\ \text{Lie hom.}:\ \pi_{\mathrm{pair}}\circ\varphi\ \text{hom. image of simple}\implies\ker
\]

%dsp-level:0
\[
= \{0\}\ \vee\ \mathfrak{su}(N);
\]

%dsp-level:0
\[
\quad\ker
\]

%dsp-level:0
\[
= \{0\}\implies\dim\mathfrak{su}(N)
\]

%dsp-level:0
\[
\le 9
\]

%dsp-level:0
\[
\implies N^2-1
\]

%dsp-level:0
\[
\le 9
\]

%dsp-level:0
\[
\implies N\qquad \le 3;
\]

%dsp-level:0
\[
\qquad N\ge4:\ \text{both projections trivial}\implies\varphi\equiv0
\]

%dsp-level:0
\[
(4):\ N\le3:\ \text{realized by}\ (1);
\]

%dsp-level:0
\[
\qquad N\ge4:\ (3);
\]

%dsp-level:0
\[
\qquad \dim\ \text{forcing}:\ N^2>18\iff N\ge5:\ \ker\ne0\ \text{for any linear map}\ [\mathrm{FT.26}(3)].\ \blacksquare
\]

\textbf{Rmk (what would change the verdict).} Clause (3) is an inventory theorem, not a permanent barrier: a declared exchange block $\mathcal K^{(\mathrm{exch})}$ (new aspect candidate, interface-definition species) or a derived cross-vertex bridge would reopen hosting for separated contacts and $N\ge4$ corollas. Neither is constructed or declared here; the theorem states exactly what the \textit{current} inventory forces. Note the physical reading of the corolla cases: at $n(w)=2$ the exchange flip is the environment vertex's own weak isospin flip [CT1.4] — the "collective mode of $\Sigma$" is an environment mode in disguise; only separated contacts (and large corollas) carry structure with no home in the existing forces.

\textbf{Phys.} The coarse-graining machine has a hard collective ceiling, and it is the colour algebra: two boundary contacts exchange through the weak block, three through the colour block — four and beyond, the entire nonabelian exchange structure falls into the kernel of every available read. The kernel is not a defect; it is what coarse-graining means at the operator level. Whether nature's larger boundaries declare a new block or simply forget their exchange structure is the sharpest extension question the program has produced — and it is now a theorem that the existing sixteen registers cannot answer it.

\vspace{2pt}\noindent\rule{\textwidth}{0.4pt}\vspace{2pt}

\end{theorem}

\phantomsection
\section[Tier F.19 — The Three Tails: Realization Non-Existence, the…]{\texorpdfstring{Tier F.19 — The Three Tails: Realization Non-Existence, the Multiplicity Map, and the Selection Obstruction}{Tier F.19 - The Three Tails: Realization Non-Existence, the Multiplicity Map, and the Selection Obstruction}}

\medskip\noindent\textit{\texorpdfstring{FT.59 — The $\iota_{n\ge4}$ specification: non-existence, persistence refuted, hosting ceiling}{FT.59 - The specification: non-existence, persistence refuted, hosting ceiling}.}\ Structural-constraints record: see Appendix B, Section B.3.2 of the Addendum.\par\medskip

\begin{definition}[{\texorpdfstring{Def-H.38 — The multiplicity–winding coupling}{Def-H.38 - The multiplicity-winding coupling}}]\label{tag:Def-H-38}
\textbf{Def.} A materialized closed bound chain $c$ [D1.5: bound content fermionic; FT.22: materialized] with winding $k(c)=\Theta(c)/2\ge0$ [FT.22, FT.25(3); the nonnegative (spacelike-dominant) cone, FT.52(3)] carries $k(c)$ symmetric copies of the transported seed:

%dsp-level:0
\[
\textbf{multiplicity rule (constitutive clause):}\qquad \operatorname{copies}(c)\ :=\ k(c),
\]

realized on the carrier $\Psi(c)=\mathrm{Sym}^{k(c)}(\mathcal H_{\mathsf v})$ with the diagonal transport $\mathcal U_c^{\otimes k}\big|_{\mathrm{Sym}^k}$ [Def-H.23 — immediate below; EW.32: $\mathcal H_{\mathsf v}\cong\mathbb C^2$ the transported seed]. \textbf{Derived content (not part of the interface).} Uniqueness of $\Psi$ given the rule [FT.52(3), FT.55(2), FT.60(2)]; the four conditions (i)–(iv), including the corpus-pinned placement at $k{=}2$ — real $\iff$ support in $\vee^2$ at $\varepsilon_q=0$, so two real-registered seeds grade symmetric [FT.60(1); A4, CT0.7, CT0.8]; the Casimir ladder $\tfrac k2(\tfrac k2{+}1)$ generalizing CT0.7a [FT.60(3)]; the weight ladder and the defect multiplets [FT.60(4)–(5), under Def-H.37(2)]; composition as maximal-weight Clebsch–Gordan on the cone [FT.60(6)]; the spin label $j=|k|/2$ [FT.55(2), FT.60(1)(iv)]. \textbf{Minimality (derived).} The coupling is irreducible to prior corpus content: EW.32 transports one seed per chain regardless of winding; the uniform-minimal triangle ($k=3$) carries no three materialized winding-$1$ subloops; admissible decompositions are basis-relative [FT.60(7a)–(b) — forward pointers in-tier]. One integer coupling; no construction below it. \textbf{Scope.} Bound chains on the nonnegative cone only; transit/trivial chains excluded — photon helicity [EM.8] is a transit-side read, no identification. Carriers for $k\ge3$ ($\dim\ge4$) exceed the registered pair-channel inventory [CT0.12: blocks $3,3,5,3,1,1$] — no aspect claim at $k\ge3$; the inventory-extension question is application-layer content [FT.60(8)].

\end{definition}

\begin{definition}[{\texorpdfstring{Def-H.23 — The winding-multiplicity carrier map (under Def-H.38)}{Def-H.23 - The winding-multiplicity carrier map (under Def-H.38)}}]\label{tag:Def-H-23}
\textbf{Def.} For $c$ a \textbf{materialized closed bound chain} [D1.5: bound content fermionic; FT.22: materialized] with winding $k(c)=\Theta(c)/2\ge0$ [FT.22, FT.25(3)] — the nonnegative (spacelike-dominant) cone [FT.52(3) form]:

%dsp-level:0
\[
\Psi(c)\ :=\ \mathrm{Sym}^{k(c)}\bigl(\mathcal H_{\mathsf v}\bigr)\ \subseteq\ \mathcal H_{\mathsf v}^{\otimes k(c)}
\]

%dsp-level:0
\[
\qquad [\mathrm{Def-H.38}:\ \text{multiplicity}=\text{winding};\ \ \mathcal H_{\mathsf v}\cong\mathbb C^2\ \text{the transported seed},\ \mathrm{EW.32}],
\]

where $\mathrm{Sym}^{k}$ is the common $(+1)$-eigenspace of all pairwise swaps on $\mathcal H_{\mathsf v}^{\otimes k}$ [math: the symmetric subspace; the swaps act as CT0.7's $\mathbb S$ on each tensor pair, adjacent transpositions generating $S_k$], and transport is the \textbf{diagonal action} $\mathcal U_c^{\otimes k}\big|_{\mathrm{Sym}^k}$ of EW.32's seed transport $\mathcal U_c$ around $c$ [EW.32, EW.33; Def-H.21(5) for the spin-connection form]. Boundary cases: $k=0$: $\mathrm{Sym}^0=\mathbb C$, the trivial carrier — representation-identified with $\wedge^2\cong\mathrm{Sym}^0$ [FT.56(1)], $\neg$subspace-identified; scope note: all-null trivial chains are transit objects (photons/gluons [D1.7, D3.7]) whose helicity content is a transit-side read [EM.8] — outside $\Psi$, no identification made. $k<0$ (timelike-dominant): outside the cone, $\Psi$ undefined [FT.52(3); FT.56(2)]. \textbf{Interface status (Def-H.38).} The corpus proves one seed transported per chain regardless of winding [EW.32]; the coupling \textit{multiplicity $=k(c)$} is Def-H.38's constitutive clause — FT.60(7) proves it irreducible to decomposition, so the clause is load-bearing, not packaging.

\vspace{2pt}\noindent\rule{\textwidth}{0.4pt}\vspace{2pt}

\end{definition}

\begin{theorem}[{\texorpdfstring{FT.60 — The multiplicity map: conditions, ladders, basis-relativity, and the honesty witnesses}{FT.60 - The multiplicity map: conditions, ladders, basis-relativity, and the honesty witnesses}}]\label{tag:FT-60}
\textbf{Thm.} Under Def-H.38, with Def-H.21's spin connection and EW.32's normalization: $(1)$ \textbf{The four conditions hold.} $\Psi$ satisfies FT.56(3)(i)–(iv): (i) $k{=}1:\ \mathrm{Sym}^1=\mathcal H_{\mathsf v}$, $j=\tfrac12$ [EW.32 anchor]; (ii) $k{=}2:\ \mathrm{Sym}^2=\vee^2\mathcal H_{\mathsf v}$ — the channel's own reality sector, with projector $\Pi_{\mathcal R}=\tfrac14(3E_0+\Sigma)$ [CT0.7(1); A3] — and the \textit{placement} is corpus-pinned: real $\iff$ support in $\vee^2$ at $\varepsilon_q=0$ [A4, CT0.8], so once two seeds share a materialized (real) register their symmetric grading is forced [CT0.7 Phys: reality $=$ Bose exchange symmetry]; only the multiplicity is interface content [Def-H.38]; (iii) cone-additivity: $\mathrm{Sym}^{k_1}\otimes\mathrm{Sym}^{k_2}\supseteq\mathrm{Sym}^{k_1+k_2}$ as maximal weight [math: CG] $=$ the additive winding law on the stretched channel [FT.34(3), FT.50(2), FT.56(2)]; (iv) admissibility: carrier dimension $k+1=2j+1$ with $j=|k|/2$ on the cone [FT.52(3) form; FT.55(0)]. $(2)$ \textbf{Uniqueness.} Any carrier assignment satisfying (i)–(iv) equals $\Psi$: the $(k{+}1)$-dim irreducible $\mathfrak{su}(2)$-module is unique up to isomorphism [math: classification; FT.55(2)] and the anchor fixes the seed [EW.32]. Gap 19's map exists and is unique under Def-H.38. $(3)$ \textbf{Casimir ladder (CT0.7a generalized).} With $W_i\mapsto$ total-spin generators on $\mathcal H_{\mathsf v}$ [EW.32/CT1.2: the corpus's own basis match] and the induced diagonal action on $\mathrm{Sym}^k$:

%dsp-level:0
\[
\textstyle\sum_i\bigl(W_i^{(k)}\bigr)^2=\frac{k}{2}\Bigl(\frac{k}{2}+1\Bigr)\,\mathbb I_{k+1}
\]

%dsp-level:0
\[
\qquad [\text{machine: exact rational matrices,}\ k\le6;\ \text{physics normalization}\ W=\tfrac12\cdot\text{generators}],
\]

recovering the native rungs verbatim: $k{=}0:\ 0$ on $\wedge^2$; $k{=}1:\ \tfrac34$ on the seed; $k{=}2:\ 2$ on $\vee^2$ — CT0.7a's eigenvalues are the $k\le2$ instances of one ladder [FT.56(1) extended from three rungs to all]. $(4)$ \textbf{Weight ladder.} For $\mathcal U=\exp(i\phi\,\hat n\cdot\vec\sigma)$ [FT.54's convention]: $\mathrm{Sym}^k(\mathcal U)$ has eigenvalues $e^{im\phi}$, $m\in\{k,k-2,\dots,-k\}$ — i.e. $\mathfrak{su}(2)$-weights $m/2\in\{j,j-1,\dots,-j\}$ at $j=k/2$, character $\sum_m e^{im\phi}$; on spin-connection transport $\phi=\Theta_R/2$ [Def-H.21(6)] [machine: diagonal transport over $\mathbb Q(i)$ at exact circle points; Laurent identity $(t^m-1)(t^{-m}-1)=2-t^m-t^{-m}$]. $(5)$ \textbf{Defect multiplet ladder (under Def-H.37(2)).} The CT2.5 defect read $\|(\mathcal U^{(k)}-1)\psi\|^2$ on the weight-$m$ vector [$m=k-2r$ in the $\phi$-labelling]:

%dsp-level:0
\[
m^2_{\mathrm{pole}}(m)=\hat\mu'^2\bigl(2-2\cos(m\,\Theta_R/2)\bigr)=4\,\hat\mu'^2\sin^2\bigl(m\,\Theta_R/4\bigr),\qquad m\in\{k,k-2,\dots,-k\};
\]

$k{=}1$ reproduces FT.54(4) exactly and its \textit{state-independence is special to the seed} ($m=\pm1$ agree by evenness); $k\ge2$ \textbf{splits into a multiplet}; every even $k$ carries a weight $m{=}0$ that is defect-massless at every curvature flux — the kernel coherence of FT.33(6)/FT.54 generalized to a standing zero of the ladder [machine: exact over $\mathbb Q(i)$; Chebyshev identity $2-2\cos(m\phi)=2-2T_m(\cos\phi)$ verified as integer-polynomial recursion]. $(6)$ \textbf{Composition at carrier level.} For $\Gamma=c_1\cdot P\cdot c_2\cdot\bar P$ materialized with $d_{12}=\Theta(P)\ge0$: $\Psi(\Gamma)=\mathrm{Sym}^{k_1+k_2+d_{12}}\subseteq\mathrm{Sym}^{k_1}\otimes\mathrm{Sym}^{k_2}\otimes\mathrm{Sym}^{d_{12}}$ as maximal weight [math: CG] — the spin–orbit decomposition $j(\Gamma)=j_1+j_2+j_{\mathrm{rel}}$, $j_{\mathrm{rel}}=d_{12}/2$ [FT.56(2)], now an inclusion of carriers, on the cone. $(7)$ \textbf{Two honesty witnesses.} (a) \textit{Basis-relativity}: the carrier assignment is per-chain, and a floor admits several admissible chain decompositions with the same total winding but different tensor carriers — witness (machine): the tetrahedron at uniform $\kappa_0=2$: the digon basis (M2) gives six chains of $k=2$: $\bigotimes^6\mathrm{Sym}^2$, dimension $3^6=729$; the face basis (M1) gives four faces of $k_f=3$: $\bigotimes^4\mathrm{Sym}^3$, dimension $4^4=256$; both have $\sum k=12=K_S$ [FT.43(1)'s sum rule is the invariance mechanism]. \textbf{Decomposition-invariant content}: $K_S$ and the stretched total $j_{\max}=K_S/2$ (and per-chain multisets up to the homology choice); the full floor-level tensor product is basis-relative — recorded, no canonical choice imposed. Only edge-disjoint decompositions [FT.25(1)] (the digon basis) support an unrestricted $\otimes$; face bases overlap on channels and their $\otimes$ is formal. (b) \textit{Irreducibility of the multiplicity datum}: $\Psi$ is not "count materialized subloops" — the uniform-minimal triangle has $k=3$ [FT.45(4): $\kappa\equiv2$, $\Theta=6$] yet carries no materialized winding-$1$ \textit{face} subloop ($\kappa\equiv1$ gives $\Theta=3$, odd, off-shell [FT.22; machine]), and the profile identity $(2,2,2)=(1,1,0)+(0,1,1)+(1,0,1)$ decomposes into null-carrying digon-supported profiles — different chains, not three winding-$1$ loops. EW.32 transports one seed per chain regardless of $k$; no tag ties tensor multiplicity to winding [FT.56(3)]. Def-H.38's single datum is therefore genuinely new — a transfer of exactly Def-H.37(2)s species. $(8)$ \textbf{Inventory consequence, flagged.} $\mathrm{Sym}^k$ for $k\ge3$ ($\dim\ge4$) lies in multi-seed tensor spaces with no counterpart among the frozen $16$ pair registers [CT0.12 block dimensions $3,3,5,3,1,1$: no $4$-dim irreducible] nor the stack blocks: the carriers for $k\ge3$ exist as [math] objects on $\mathcal H_{\mathsf v}^{\otimes k}$ but exceed the registered aspect inventory — an inventory-extension consequence of Def-H.38, recorded here and joined to Gap 16's residue-(b) declaration question ($\mathcal K^{(\mathrm{exch})}$-species). No aspect claim is made for $k\ge3$ carriers.

\textbf{Req.} Def-H.38; Def-H.23, Def-H.21; EW.32, EW.33 (Part 2); A3, A4; CT0.7, CT0.7a, CT0.8, CT0.12 (Part 1); CT1.2 (Part 1); CT2.5 (Part 2); D1.5, D1.7, D3.7; EM.8 (Part 2); FT.22, FT.25, FT.33, FT.34, FT.43, FT.45, FT.49, FT.50, FT.52, FT.54, FT.55, FT.56; $\mathfrak{su}(2)$ classification, Clebsch–Gordan, symmetric-subspace characterization, Chebyshev recursion [math].

\textbf{Pf.}

%dsp-level:0
\[
(1):\ \mathrm{Sym}^1=\mathcal H_{\mathsf v}\ [\text{math}];\ j=\tfrac12\ [\mathrm{EW.32}];
\]

%dsp-level:0
\[
\qquad \mathrm{Sym}^2(\mathbb C^2)=\vee^2\ [\mathrm{A3}];\ \Pi_{\mathcal R}=\tfrac14(3E_0+\Sigma),\ \mathfrak T=3\ [\mathrm{CT0.7}(1);\ \text{machine regression}]
\]

%dsp-level:0
\[
\varepsilon_q
\]

%dsp-level:0
\[
= 0:\ (p,q)\in\mathcal R\iff\operatorname{supp}\rho_{pq}\subseteq\vee^2\ [\mathrm{CT0.8}]:\ \text{two real-registered seeds grade symmetric}\ [\mathrm{CT0.7\ Phys};\ \text{math}:\ \text{common}\ {+}1\ \text{eigenspace of transpositions}=\mathrm{Sym}^k]
\]

%dsp-level:0
\[
\mathrm{CG}:\ \mathrm{Sym}^{k_1}\otimes\mathrm{Sym}^{k_2}=\bigoplus_{j=|k_1-k_2|/2}^{(k_1+k_2)/2}\mathcal D_j\supseteq\mathrm{Sym}^{k_1+k_2}\ [\text{math}];
\]

%dsp-level:0
\[
\qquad k(\Gamma)=k_1+k_2+d_{12}\ [\mathrm{FT.34}(3),\mathrm{FT.50}(2)]
\]

%dsp-level:0
\[
(2):\ \dim=k+1\ \wedge\ \text{irreducible}\ [\text{math}]\implies\text{unique isomorphism type per}\ k;
\]

%dsp-level:0
\[
\quad R(1)=\tfrac12\ \text{anchor}\ [\mathrm{EW.32}];
\]

%dsp-level:0
\[
\quad j=|k|/2\ [\mathrm{FT.52}(3)]
\]

%dsp-level:0
\[
(3):\ \text{monomial basis}\ x^{k-r}y^r:\ H
\]

%dsp-level:0
\[
= x\partial_x-y\partial_y,
\]

%dsp-level:0
\[
\ E
\]

%dsp-level:0
\[
= x\partial_y,
\]

%dsp-level:0
\[
\ F
\]

%dsp-level:0
\[
= y\partial_x:\ [H,E]
\]

%dsp-level:0
\[
= 2E,
\]

%dsp-level:0
\[
\ [H,F]
\]

%dsp-level:0
\[
= -2F,
\]

%dsp-level:0
\[
\ [E,F]
\]

%dsp-level:0
\[
= H\ [\text{machine},\ k\le6]
\]

%dsp-level:0
\[
\textstyle\sum_i(W_i^{(k)})^2=\tfrac12(EF+FE+\tfrac12H^2)=\tfrac k2(\tfrac k2+1)\,\mathbb I\ [\text{machine exact over}\ \mathbb Q];
\]

%dsp-level:0
\[
\qquad k=0,1,2:\ 0,\tfrac34,2\ [\mathrm{CT0.7a}\ \text{eigenvalues recovered}]
\]

%dsp-level:0
\[
(4):\ \mathcal U
\]

%dsp-level:0
\[
= \exp(i\phi\hat n\cdot\vec\sigma):\ \text{eigenvalues}\ e^{\pm i\phi}\ \text{on}\ \mathcal H_{\mathsf v}
\]

%dsp-level:0
\[
\implies e^{i(k-2r)\phi}\ \text{on}\ \mathrm{Sym}^k\ [\text{math};\ \text{machine over}\ \mathbb Q(i),\ \text{diagonal}\ \mathcal U,\ \text{exact circle points}]
\]

%dsp-level:0
\[
(5):\ \mathrm{CT2.5}\ \text{defect read}\ [\mathrm{EW.33}\ \text{pattern}]:\ \|(U^{(k)}-1)\psi_m\|^2
\]

%dsp-level:0
\[
= |e^{im\phi}-1|^2
\]

%dsp-level:0
\[
= 2-2\cos(m\phi);
\]

%dsp-level:0
\[
\quad\phi
\]

%dsp-level:0
\[
= \Theta_R/2\ [\mathrm{Def\text{-}H.21}(6)]
\]

%dsp-level:0
\[
k=1:\ m=\pm1:\ 2-2\cos(\Theta_R/2)=4\sin^2(\Theta_R/4)\ [\mathrm{FT.54}(4)\ \text{recovered};\ \text{both weights equal}];
\]

%dsp-level:0
\[
\qquad m=0\ (k\ \text{even}):\ 2-2\cos0=0\ \text{exactly}
\]

%dsp-level:0
\[
2-2\cos(m\phi)
\]

%dsp-level:0
\[
= 2-2T_m(\cos\phi):\ T_0
\]

%dsp-level:0
\[
= 1,
\]

%dsp-level:0
\[
\ T_1
\]

%dsp-level:0
\[
= c,
\]

%dsp-level:0
\[
\ T_{m+1}
\]

%dsp-level:0
\[
= 2cT_m-T_{m-1}\ [\text{machine integer-polynomial recursion over}\ \mathbb Q]
\]

%dsp-level:0
\[
(6):\ (1)\text{'s CG at three factors on the cone};\qquad d_{12}\ge0\ [\mathrm{FT.34}(2)\ \text{spacelike connector}]
\]

%dsp-level:0
\[
(7a):\ \kappa_0
\]

%dsp-level:0
\[
= 2\ \text{tetrahedron}\ [\mathrm{FT.42}(5)]:\ \text{digons}\ \Theta
\]

%dsp-level:0
\[
= 2\kappa_0
\]

%dsp-level:0
\[
= 4,
\]

%dsp-level:0
\[
\ k
\]

%dsp-level:0
\[
= \kappa_0
\]

%dsp-level:0
\[
= 2\ [\mathrm{FT.24}(1),\mathrm{FT.25}(3)],
\]

%dsp-level:0
\[
\ 6\ \text{chains};
\]

%dsp-level:0
\[
\ \text{faces}\ \Theta_f
\]

%dsp-level:0
\[
= 2L
\]

%dsp-level:0
\[
= 6,
\]

%dsp-level:0
\[
\ k_f
\]

%dsp-level:0
\[
= 3,
\]

%dsp-level:0
\[
\ 4\ \text{faces};
\]

%dsp-level:0
\[
\ \sum
\]

%dsp-level:0
\[
= 12
\]

%dsp-level:0
\[
= K_S\ [\mathrm{FT.43}(1)]
\]

%dsp-level:0
\[
\dim:\ 3^6=729\ne256=4^4\ [\text{machine}];\qquad j_{\max}=K_S/2=6\ \text{both bases}\ [\text{stretched};\ \text{math}]
\]

%dsp-level:0
\[
(7b):\ \kappa\equiv1\ \text{on}\ 3\text{-cycle}:\ \Theta=3\notin2\mathbb Z\ [\mathrm{FT.22}];
\]

%dsp-level:0
\[
\qquad (2,2,2)=(1,1,0)+(0,1,1)+(1,0,1)\ [\text{machine}];\ (1,1,0)\text{-class null-carrying}\ [\mathrm{FT.45}(4)]
\]

%dsp-level:0
\[
(8):\ \dim\mathrm{Sym}^3
\]

%dsp-level:0
\[
= 4;
\]

%dsp-level:0
\[
\quad\mathrm{CT0.12}\ \text{blocks}\ \{3,3,5,3,1,1\}:\ \neg\exists\ 4\text{-dim irreducible register block};
\]

%dsp-level:0
\[
\quad\text{no aspect claim made}.\ \blacksquare
\]

\textbf{Rmk (Def-H.38 and the no-go, consistent).} FT.55(1) ruled out $\Theta_R=2\pi k$ — spin as the frame-holonomy \textit{angle} of the same loop. Def-H.38 makes no such identification: the winding enters as \textit{copy count} of the transported seed, and the transport phases are EW.32-side data; under Def-H.37(2) the defect ladder (5) reads in $\Theta_R$ at fixed weight $m$ with no $k$–$\Theta_R$ proportionality anywhere — the calibration weights stay separate ($[k]=\mathfrak h$-weight $-2$; $[\Theta_R]=0$; $[\hat\mu'^2]=\mathfrak q$). The no-go constrains the \textit{form} of the read; Def-H.38 supplies the \textit{multiplicity}; they meet only in (5), consistently.

\textbf{Rmk (relation to Def-H.37(1)).} Def-H.37 fixes the defect \textit{transfer} (pole side $=$ transport defect of the seed); Def-H.38 fixes the defect \textit{carrier} (which space the seed's copies occupy).

\textbf{Phys.} A winding-$k$ fermion loop carries $k$ symmetric copies of the qubit the corpus already transports — spin $k/2$, Casimir $\tfrac k2(\tfrac k2+1)$, a multiplet of $k+1$ weights, and, if the pole mass is the transport defect, a mass multiplet $4\hat\mu'^2\sin^2(m\Theta_R/4)$ whose middle state is exactly massless whenever $k$ is even. The bottom three rungs are not added structure: they are the virtual line, the seed, and the real sector of the pair channel — objects Parts 1–2 already proved, wearing the labels CT0.7a already computed. What the declaration contributes is one number tied to another: copies $=$ winding. That the tie cannot be unpacked from anything already proven is itself proven — the triangle's minimal winding of three has no three subloops to count.

\vspace{2pt}\noindent\rule{\textwidth}{0.4pt}\vspace{2pt}

\end{theorem}

\begin{theorem}[{\texorpdfstring{FT.61 — The selection side: phase inventory, index-blindness, and the two exits (Gap 18)}{FT.61 - The selection side: phase inventory, index-blindness, and the two exits (Gap 18)}}]\label{tag:FT-61}
\textbf{Thm.} $(1)$ \textbf{Complete on-shell phase inventory.} For a configuration $C$ with on-shell integer edge data [FT.24/FT.34(4) scope], A8's weight assembles as

%dsp-level:0
\[
w(C)=\mathcal J(C)\prod_{e\in E(C)}G(e)\cdot(-1)^{\aleph(\sigma,C)},
\]

%dsp-level:0
\[
\qquad \mathcal J(C)=1\ [\mathrm{D3.6}:\ (+1)\text{-factors only}],
\]

%dsp-level:0
\[
\qquad G(e)=\sigma_e\,e^{i\pi\kappa_e}\ [\mathrm{D3.5},\ \mathrm{FT.9}:\ \alpha=\pi;\ \mathrm{EM.20}],
\]

so the full phase is $\Phi(C)=\pi\bigl(\Theta(C)+\aleph(\sigma,C)+n_-(\sigma)\bigr)\ (\mathrm{mod}\ 2\pi)$ with $n_-:=\#\{e:\sigma_e=-1\}$: \textbf{the path-sum phase has exactly three $\mathbb Z_2$ registers — geometric ($\Theta$), wall ($\aleph$ [D3.3]), sign ($n_-$)} [CT0.9, CT0.10: the sign structures are independent; machine: exhaustive on $\Sigma_{\min}$ configurations]. $(2)$ \textbf{Geometric triviality on-shell.} On materialized closed chains $\Theta\in2\mathbb Z$ [FT.22] $\implies(-1)^{\Theta}=+1$: the surviving weight phase of any materialized-loop configuration is registration-only, $(-1)^{\aleph+n_-}$ — \textit{no variational continuous component remains in the on-shell phase}. In particular the stationary channels of CT1.20's semiclassical locus carry constant $\kappa_0$ [FT.53(2)] and their complete geometric ledger record is the single integer $\kappa_0$ (plus $\delta h$ and $\Theta=2k$). $(3)$ \textbf{Index-blindness witness.} There exist CT2.6-class geometries and pairs of closed geodesics $(\gamma_A,\gamma_B)$ with \textit{identical} ledger records and \textit{different} Morse indices: witness $M=S^2(R)\times S^1(2R)$ with the product metric [math: smooth complete Riemannian 3-manifold — a spatial-slice class in Def-H.21's sense; whether emergent dynamics realizes this slice is P4/solution data, recorded not claimed]: $\gamma_A=$ a great circle traversed twice — length $4\pi R$, Morse index $3$ [math: Morse index theorem; conjugate points at arc-distances $\pi R,2\pi R,3\pi R$ in the open interval; iterates of $S^2$ great circles: index $2m-1$]; $\gamma_B=\{\mathrm{pt}\}\times S^1(2R)$ — length $2\pi\cdot2R=4\pi R$ (equal), Morse index $0$ [flat factor; product formula; $\pi_1$-nontrivial so index $0$ is admissible — Synge's obstruction applies only to simply-connected slices]. At ledger scale $\mathfrak h=2\pi R/N$: both discretize to $\kappa\equiv1$ minimal-edge chains with $\delta h=\Theta=2N\in2\mathbb Z$ [FT.31(4), FT.22: materialization-eligible], $\sigma\equiv+1\implies\aleph=0,\ n_-=0$; by (1) both weights are $+1$, and $|w|^2=1$ for both [CT2.8's norms equally blind]. \textbf{No function of the weight data separates index $0$ from index $\ge1$ stationary channels.} \textit{([FT.66(2)]: both witness members are degenerate — they sit in clean isometry-orbit critical manifolds ($\mathcal C_A=\mathbb{RP}^2\times S^1\times S^1$, nullity $4$; $\mathcal C_B=S^2\times S^1$, nullity $3$); the Bott indices remain distinct ($3$ vs $0$), the monodromy types differ ($I_2$ semisimple vs unipotent parabolic, Jordan-separated at identical trace/det), and the records and weights above are unchanged: this clause holds with "Morse index" read throughout as "Bott index" [Def-H.24].)} $(4)$ \textbf{Two exits (theorem-shaped localization).} A selection criterion implementing "stable jumps materialize" [Def-H.19 vocabulary] must read the Morse index [math: index is the fluctuation-sign datum]. By (1)–(3) no weight-level or ledger-level read can supply it, and the mask/pole side reads stationary channels only through hosted blocks and superselection labels $(\Theta_{\mathrm{int}},k,q_{\mathrm{core}},m_{\mathrm{hid}},\mathcal C)$ [FT.26, Def-H.15(4), FT.58's inventory] — every entry index-blind by (2). Hence Gap 18's selection side is closed \textit{against the current inventory}, with exactly two exits: \textbf{(a)} a new variational registration — a per-channel conjugate-point/Jacobi count appended to the ledger (a interface-definition species extending D0.7's retained-event vocabulary; nothing of the kind exists [inventory]) — **(reclassification [FT.66(3)]: the datum is reconstructable in principle from the ambient configuration's metric content [CT1.20's $C\mapsto g_{\mu\nu}(C)$; CT1.12/CT1.14/CT2.6] and eventually constant in $\mathfrak h$ away from bifurcation loci; the gap is a structural \textit{coupling} — no A8/D2.11/D2.3/CT2.5 clause reads it — not a registration gap; exit (a) becomes a declared weight extension $w\mapsto w\,(-1)^{\operatorname{ind}_B}$, interface-definition species, not adopted)\textbf{; or }(b)** a spectral-side argument — positivity structure on $Z(\mathcal R,\Sigma)$ [CT2.8/CT2.5/Def-H.16] constraining which stationary families contribute poles without per-channel signs (not constructed; CT2.8 supplies norms, not pole-selection). The earlier "bridge construction" residue decomposes accordingly: the \textit{leading} bridge is derived — materialization $\iff$ pole $\iff$ single-valued eigenmode at constant-$\kappa$ stationary channels [Def-H.16 closing line; FT.22; FT.53(2); CT2.5] — and the \textit{next order} (fluctuation prefactors and signs) requires exactly (a) or (b). $(5)$ \textbf{Consistency with CT1.20.} The witness pair sits inside CT1.20's own statement: stationary points dominate and are weighted alike — (3) upgrades "alike" from an absence of distinction to an identity of weights, at every ledger scale. Gap 13's finite-$\mathfrak h$ residue [FT.53(4)] inherits the localization: the excluded class-(iii) profiles were killed by \textit{non-stationarity}; index-$\ge1$ stationary saddles are killed by nothing available.

\textbf{Req.} A8; D0.7, D1.7, D2.7, D3.3, D3.5, D3.6; CT0.9, CT0.10 (Part 1); CT1.20, CT2.5, CT2.6, CT2.8 (Part 2); EM.20 (\textbf{E}; via FT.9's $\alpha=\pi$); FT.9, FT.22, FT.24, FT.26, FT.31, FT.34, FT.48, FT.53, FT.58; Def-H.15, Def-H.16, Def-H.19, Def-H.21; P4; Morse index theorem, conjugate-point counts for sphere iterates, product formula, Synge [math].

\textbf{Pf.}

%dsp-level:0
\[
(1):\ \mathcal J(C)=\prod_{J^+}(+1)\prod_{J^-}(+1)^*=1\ [\mathrm{D3.6}];
\]

%dsp-level:0
\[
\qquad G(e)=\sigma_e\exp(i\alpha_{\mathrm{EM}}\widehat{ds}^{\,2}(e))=\sigma_e(-1)^{\kappa_e}\ [\mathrm{D3.5};\ \mathrm{FT.9}/\mathrm{EM.20}:\ \alpha=\pi;\ \kappa_e\in\mathbb Z\ \text{on-shell}]
\]

%dsp-level:0
\[
\prod_eG(e)=(-1)^{n_-}\cdot(-1)^{\Theta}\ [\mathrm{D1.7}:\ \Theta=\textstyle\sum_e\widehat{ds}^{\,2}(e)];
\]

%dsp-level:0
\[
\qquad w\text{-phase}=(-1)^{\Theta+\aleph+n_-};
\]

%dsp-level:0
\[
\qquad \text{independence of registers}\ [\mathrm{CT0.10}];\ \text{machine}:\ 24/24\ \Sigma_{\min}\ \text{configurations}
\]

%dsp-level:0
\[
(2):\ \Theta\in2\mathbb Z\ [\mathrm{FT.22}]\implies(-1)^\Theta=1;
\]

%dsp-level:0
\[
\qquad \text{stationary record}:\ \kappa_e\equiv\kappa_0\ [\mathrm{FT.53}(2)];\quad\Theta=2k\ [\mathrm{FT.22}];\quad\delta h\ [\mathrm{D0.6}]
\]

%dsp-level:0
\[
(3):\ \ell(\gamma_A)=2\cdot2\pi R=4\pi R=\ell(\gamma_B);
\]

%dsp-level:0
\[
\qquad \text{index}(\gamma_A)=\#\{j\pi R:\ 1\le j\le2m-1,\ m=2\}=3\ [\text{Morse index thm.; machine: integer count}]
\]

%dsp-level:0
\[
\text{index}(\gamma_B)=0\ [\text{flat factor},\ \text{product formula};\ \gamma_B\ \text{generates}\ \pi_1\implies\neg\text{contractible}\implies\neg\text{Synge-forced}];
\]

%dsp-level:0
\[
\qquad \mathfrak h=2\pi R/N:\ \delta h=\ell/\mathfrak h=2N\ \text{both};\ \kappa\equiv1\implies\Theta=\delta h\ [\mathrm{FT.31}(4)]
\]

%dsp-level:0
\[
\text{ledger records}:\ (\kappa_0,\delta h,\Theta,\sigma,\aleph)=(1,2N,2N,+1,0)\ \text{both};
\]

%dsp-level:0
\[
\qquad w=(+1)\ \text{both by}\ (1);\quad|w|^2=1\ \text{both}\ [\mathrm{CT2.8}(2)\ \text{norm side}];\qquad \qquad 3\ne0
\]

%dsp-level:0
\[
(4):\ \text{index}=\text{sign of fluctuation determinant}\ [\text{math}];
\]

%dsp-level:0
\[
\quad (1)\text{--}(3):\ \neg\exists\ \text{weight function distinguishing};
\]

%dsp-level:0
\[
\quad\text{mask reads}:\ \mathrm{FT.26}(2)\ \text{block-diagonal on}\ \mathcal A(\Sigma;C),
\]

%dsp-level:0
\[
\ \text{labels}\ \mathrm{Def\text{-}H.15}(4):\ \text{all from}\ (2)\text{'s record}
\]

%dsp-level:0
\[
\text{exits}:\ (a)\ \mathrm{D0.7}\ \text{ledger}=\text{retained events};\ \neg\exists\ \text{variational datum}\ [\text{inventory}];
\]

%dsp-level:0
\[
\quad (b)\ \mathrm{CT2.8}:\ \langle A|A\rangle=\sum|w|^2\ge0:\ \text{norms};\ \neg\exists\ \text{per-saddle phase/pole-selection tag}
\]

%dsp-level:0
\[
(5):\ \mathrm{CT1.20}:\ \delta\Phi=0\ \text{dominates, weighted alike};
\]

%dsp-level:0
\[
\quad (3)\ \text{exhibits the "alike" as identical weights};
\]

%dsp-level:0
\[
\qquad \mathrm{FT.53}(4):\ \text{class-(iii) excluded by non-stationarity},\ \neg\text{by index}.\ \blacksquare
\]

\textbf{Rmk (scope of the witness).} The witness geometry is offered within the smooth-geometry class the ledger program works in [CT2.6's hypotheses are regularity clauses on $(M,g)$; Def-H.21 operates on arbitrary orientable slices]. If emergent dynamics were later shown to realize only simply-connected or curvature-restricted slices, the witness would need rebuilding in that class — the \textit{structural} half of the theorem (clauses (1)–(2): the on-shell phase has no variational component at all, and the stationary record is one integer) is witness-independent and would survive any such restriction.

\textbf{Phys.} On-shell, the ledger's phase bookkeeping is finished: geometry contributes nothing (its factor is $+1$ on every materialized loop), and what remains — wall crossings and registration signs — is classical data fixed before the variation ever happens. Two loops can then be indistinguishable to the last bit of weight while one is the stablest path in the geometry and the other a thrice-unstable saddle: on a sphere-times-circle they are the same length, the same hop count, the same everything the ledger writes down. No reading of the sum can prefer the minimum; only new variational bookkeeping or a positivity argument from the spectrum side can — and both are named, neither is built.

\vspace{2pt}\noindent\rule{\textwidth}{0.4pt}\vspace{2pt}

\vspace{2pt}\noindent\rule{\textwidth}{0.4pt}\vspace{2pt}

\end{theorem}

\phantomsection
\section[Tier F.20 — The Bidirectional Rung Audit and the Platonic…]{\texorpdfstring{Tier F.20 — The Bidirectional Rung Audit and the Platonic Physical Cards}{Tier F.20 - The Bidirectional Rung Audit and the Platonic Physical Cards}}

\begin{theorem}[{\texorpdfstring{FT.62 — The bidirectional rung audit (correction theorem)}{FT.62 - The bidirectional rung audit (correction theorem)}}]\label{tag:FT-62}
\textbf{Thm.} $(1)$ \textbf{Convention re-derivation; floor-card rungs.} The directed-rung law — $n(v)=\partial^{\leftarrow}(v)+\partial^{\rightarrow}(v)$ counts \textit{directed real edges} [D0.8, A7]; the tick load is built from $\Phi^n\times\Phi^n$ sign configurations on the $n$ incident real edges [CT0.5 Pf; CT0.9: one bit per \textit{directed} real edge]; every bidirectional channel contributes both orientations — is \textbf{CT0.11b (B)}. Floor-card rows [FT.28(5)'s explicit enumeration $n(p)=3=2{+}1$; Def-H.20's bidirectionality]:

%dsp-level:0
\[
\text{single-directed-edge contact}:\ n(v_1)=2d+1;\qquad\text{bidirectional contact}:\ n(v_1)=2d+2;
\]

%dsp-level:0
\[
\text{polygon}\ \Sigma_{P_n}:\ \text{dark}\ n(v)=4\ \forall v\ (\text{type}\ (2,2)),\ \text{contact}\ n(v_1)=5;
\]

%dsp-level:0
\[
\qquad \text{Platonic}\ \Sigma_S:\ \text{dark}\ n(v)=2d\in\{6,6,8,6,10\}.
\]

$(2)$ \textbf{Corrected lock table} [$O(n)=\max(0,2^{n+1}-16)$: CT0.6]:

%dsp-level:0
\[
\begin{array}{c|c|c|c|c|c}
S & d & O(2d)\ \text{dark/vertex} & O(2d{+}1)\ \text{contact} & \sum_vO\ \text{dark} & \sum_vO\ \text{obs (1 edge)}\\ \hline
\text{tetrahedron} & 3 & 112 & 240 & 448 & 576\\
\text{cube} & 3 & 112 & 240 & 896 & 1024\\
\text{octahedron} & 4 & 496 & 1008 & 2976 & 3488\\
\text{dodecahedron} & 3 & 112 & 240 & 2240 & 2368\\
\text{icosahedron} & 5 & 2032 & 4080 & 24384 & 26432
\end{array}
\]

bidirectional-contact variants: contact $O(2d{+}2)\in\{496,496,2032,496,8176\}$, sums $\{832,1280,4512,2624,30528\}$ [machine]. $(3)$ \textbf{Saturation universality; the trichotomy retracted and replaced.} Universality — bidirectional real-edge subgraphs of channel degree $\ge2$ saturate: $D^-$-configured [CT4.1], content locked, retained, and counted [CT4.3, D4.2], inaccessible [D1.11] — the uniform minimal-saturation $2$-regular family (exactly one channel quantum over budget, length-independent), the degree-graded ladder, and the exact lock-free characterization (the \textit{unidirectional} family on bidirectional skeletons): \textbf{CT4.1a (G)} [with CT0.11b (Part 1)]. The retraction record — FT.42(2)'s "accessibility trichotomy by degree" and FT.45(2)'s "no lock at any $n$ — the unlimited-accessibility family" are \textbf{retracted}, replaced by the frozen ladder above, with $\Sigma_{\min}$ ($n\in\{2,3\}$) and $\Sigma_L$ ($n\le3$) as its floor instances: $O=0$ [FT.28(5), FT.32(2) — unchanged by this audit]. Duality splitting survives and sharpens [FT.44(2) corrected]: cube $112$ vs octahedron $496$; dodecahedron $112$ vs icosahedron $2032$ — saturation grade splits both dual pairs. $(4)$ \textbf{Invariance theorem.} $E_{\mathrm{int}}$, $I_{\mathrm{hid}}$, $m_{\mathrm{hid}}$, $\Theta_{\mathrm{int}}$, $K_S$, the $\kappa$-lattices and their indices, the winding sum rule, $g(V)$, $\Pi$, $S_{\mathrm{hid}}$, and every defect/holonomy entry are functions of edge-, channel-, and face-counts and of $\kappa$-profiles — none references $n(v)$: FT.42(1),(3),(4),(5) (lock columns excepted), FT.43, FT.44(1),(4), FT.45(1),(3),(4), FT.49–FT.51 stand unchanged [machine: value-by-value re-derivation]. The observed-variant information counts stand on retention: locked vertices retain and count [CT4.3; FT.17 cross-check 2 precedent]; the CT4.9(1) "visible $<3$" arithmetic anchors in FT.42(1)'s Pf (and FT.17 cross-check 3's) are \textbf{re-scoped}: at corrected contact rungs $n\ge7$ the contact vertex is locked, so its channel-visible content is governed by Def-H.35's locus-separation clause (per-vertex caps do not further reduce channel-visible content) and CT4.9(2)'s boundary form — not by CT4.9(1)'s accessible-vertex branch. Results unchanged; anchors corrected. $(5)$ \textbf{Type ledger and the EM row (scope note).} Bidirectional-skeleton vertices with $d\ge2$ channels have type $(d,d)$ [D0.8/A7]; D3.3's literal $f_v$ domain (type-$(1,1)$ vertices) is empty on dark polygon and Platonic floors (the digon $\Sigma_{\min}$ keeps type-$(1,1)$ vertices — $\Sigma_{\min}$'s and $\Sigma_L$'s EM rows are untouched by this audit): the EM row of FT.42(4)/FT.20 evaluates on observed variants through chain-local pass-through reads along transit chains (chain-local telescoping: EM.9b (Part 2); EM.9 end-reads); dark floors carry the registration registers ($q_{\mathrm{core}}$ per closed core [FT.40]) with no channel read at all [FT.3]. $(6)$ \textbf{Materialization sharpening.} M1/M2 at Platonic/polygon vertices are now explicitly spectral hypotheses at \textit{locked} vertices ($n=2d\ge4$): strictly stronger than the text suggested. No contradiction with EW.17 (whose rung-$3$ claim concerns wrapping chains at channel-admissible vertices — a different existence locus): face/digon materialization, if it occurs, occurs through the generic CT2.5/Def-H.16 mechanism.

\textbf{Req.} D0.8, D0.9; A7; CT0.5, CT0.6, CT0.9, CT0.10, CT0.11, CT0.11b (Part 1); CT4.1, CT4.2, CT4.3, CT4.9, CT4.1a (Part 2); EM.9b (Part 2); D1.11, D4.1, D4.2; FT.3, FT.17, FT.20, FT.21, FT.28, FT.32, FT.40, FT.42–FT.45, FT.49–FT.51; Def-H.20; EW.17 (Part 2); CT2.5 (Part 2); Def-H.16.

\textbf{Pf.}

%dsp-level:0
\[
(1):\ \text{general law}:\ \mathrm{CT0.11b}\ (\text{B});
\]

%dsp-level:0
\[
\qquad \text{cells re-enumerated}\ [\text{machine}:\ \Sigma_{\min},\ \Sigma_L,\ \Sigma_{12},\ \Sigma_\triangle,\ \text{bowtie},\ \text{patch unchanged}]
\]

%dsp-level:0
\[
(2):\ O(n)
\]

%dsp-level:0
\[
= 2^{n+1}-16\ (n\ge4)\ [\mathrm{CT0.6}];
\]

%dsp-level:0
\[
\quad\text{table arithmetic}\ [\text{machine}:\ O(6)=112,O(7)=240,O(8)=496,O(9)=1008,O(10)=2032,O(11)=4080,O(12)=8176]
\]

%dsp-level:0
\[
(3):\ \text{general ladder/universality}:\ \mathrm{CT4.1a}\ (\text{G});
\]

%dsp-level:0
\[
\text{unidirectional cells}:\ \mathrm{FT.28}(5):\ n\in\{2,3\};\ \mathrm{FT.32}(2):\ n\le3:\ O=0\ [\mathrm{CT4.1a}(4)];
\]

%dsp-level:0
\[
\qquad \text{dual pairs}:\ O(6)=112\ne496=O(8);\quad112\ne2032=O(10)
\]

%dsp-level:0
\[
(4):\ \text{each quantity's defining formula inspected}:\ \mathrm{FT.42}(1)\ (E_{\mathrm{int}}=2E),
\]

%dsp-level:0
\[
\ (3)\ (\Theta_{\mathrm{int}}=2K_S),
\]

%dsp-level:0
\[
\ (5)\ (\text{face regimes}:\ \kappa\text{-data}),
\]

%dsp-level:0
\[
\ \mathrm{FT.43}\ (\text{incidence}),
\]

%dsp-level:0
\[
\ \mathrm{FT.41}(b)\ (g(V)=V\text{-grammar}),
\]

%dsp-level:0
\[
\ \mathrm{FT.49}\text{--}\mathrm{FT.51}:\ \neg\exists\ n(v)\ \text{reference};
\]

%dsp-level:0
\[
\qquad \text{retention}:\ \mathrm{CT4.3}/\mathrm{D4.2}
\]

%dsp-level:0
\[
\text{anchor re-scope}:\ \mathrm{CT4.9}(1)\ \text{reads}\ \text{accessible}(v);
\]

%dsp-level:0
\[
\quad n\ge7:\ \neg\text{accessible}\ [\mathrm{CT0.11}]\implies\text{hypothesis check passes to Def-H.35 (interface clause)} + \mathrm{CT4.9}(2)\ [\text{boundary form}]
\]

%dsp-level:0
\[
(5):\ \text{type}(d,d):\ \partial^{\leftarrow}=\partial^{\rightarrow}=d\ [\mathrm{D0.8}];\quad\mathrm{D3.3}\ \text{domain};
\]

%dsp-level:0
\[
\qquad \mathrm{FT.21}(i)\ \text{chain-local telescoping}\ [\text{machine in}\ \mathrm{verify\_H\_ASigma}\ \text{record}];
\]

%dsp-level:0
\[
\qquad (6):\ \mathrm{Def\text{-}H.20}\ \text{M1/M2 declared};\ \mathrm{EW.17}:\ \text{rung-3 wrapping locus}.\ \blacksquare
\]

\textbf{Rmk (integrity ledger).} The fork entered with FT.42(2) ("dark $n(v)=d$") and FT.45(2) ("$n(v)=2$"), was inherited by the then-current summary claims ("$d{=}3$ solids run lock-free in the dark"; "the unlimited-accessibility family"), and survived because the later cells (face floor, bowtie, patch) were computed with the \textit{correct} directed convention — so no theorem built on FT.49/FT.50 is affected, and the two conventions never met in one computation until this pass. All affected clauses are patched in place with (rung correction) markers; the retracted epithets are recorded here and at their sites. This is the document's fourth recorded self-catch (after the mod-4 anomaly, the earlier $\sigma_2$/tuple-sum pair, and the display-balance latent defect) — and the first that changes a \textit{physical} qualitative claim: saturation is universal on bidirectional skeletons.

\textbf{Phys.} Count the ledger's own bits and the picture inverts: a bidirectional channel is two real edges, each with its own registration sign, and a vertex hosting even two of them already exceeds the channel budget. There are no clean-running bidirectional floors — polygons saturate at exactly one quantum over budget per vertex, solids by degree — and the old trichotomy collapses into a single law with a graded shadow. What survives untouched is everything the integers were made of: masses, windings, lattices, sealings, defects. The lock column was the only column that miscounted, and the correction makes the dark sector larger, not smaller.

\vspace{2pt}\noindent\rule{\textwidth}{0.4pt}\vspace{2pt}

\end{theorem}

\phantomsection
\section[Tier F.22 — The Bott Correction: Degenerate Witnesses…]{\texorpdfstring{Tier F.22 — The Bott Correction: Degenerate Witnesses Re-Evaluated, Reconstructability, and the Selection Dilemma}{Tier F.22 - The Bott Correction: Degenerate Witnesses Re-Evaluated, Reconstructability, and the Selection Dilemma}}

\begin{definition}[{\texorpdfstring{Def-H.24 — Bott index, clean critical manifolds, and monodromy type for ledger channels (construction)}{Def-H.24 - Bott index, clean critical manifolds, and monodromy type for ledger channels (construction)}}]\label{tag:Def-H-24}
\textbf{Def.} Let $\gamma$ be a stationary channel [Def-H.19: $\delta\!\int ds^2=0$; FT.48(1)/FT.53(2): geodesic, constant $\kappa_0$] in a CT2.6 slice $(M,g)$ of the emergent geometry, closed (bound side) or fixed-endpoint (transit side). All items below are [math] constructions on derived data — $g$ from sector coordinates [CT1.12, D2.1], $\Gamma$ [CT1.13], $R^\rho{}_{\sigma\mu\nu}$ [CT1.14], sampling [CT1.19], coarse-graining [CT2.6 hypotheses (i)–(iii)] — no new physical input, in Def-H.21's construction class: $(1)$ \textbf{Critical manifold.} $\mathcal C\ni\gamma$: the connected component of the stationary locus of the energy/ledger functional $\Theta\leftrightarrow\int ds^2$ [FT.31(2): the functionals coincide on-shell] in the coarse-grained loop space (closed case: free loops; transit case: fixed-endpoint paths). Isometry orbits of $(M,g)$ through $\gamma$ lie in $\mathcal C$; phase-shift and translation families likewise. $(2)$ \textbf{Normal Hessian and Bott index.} $N_\gamma:=T_\gamma\Lambda/T_\gamma\mathcal C$ the normal space; $Q_\gamma$ the second variation (index form) restricted to periodic normal sections:

%dsp-level:0
\[
Q_\gamma(V,V)=\int_0^L\Bigl(|\nabla_{\dot\gamma}V|^2-\bigl\langle R(V,\dot\gamma)\dot\gamma,\ V\bigr\rangle\Bigr)\,dt\qquad[R\ \text{from CT1.14 data};\ \text{math}],
\]

%dsp-level:0
\[
\operatorname{ind}_B(\gamma):=\dim(\text{maximal negative-definite subspace of}\ Q_\gamma),\qquad\operatorname{nul}(\gamma):=\dim\ker Q_\gamma;
\]

$\gamma$ is \textbf{clean} iff $\ker Q_\gamma$ equals the orbit/tangential nullity (all kernel directions integrate into $\mathcal C$) — equivalently $\mathcal C$ is a nondegenerate critical submanifold [math: Morse–Bott]. $\mathcal C$-constant $\operatorname{ind}_B$; the naive Morse data of FT.61/Def-H.19 are the $\operatorname{nul}=0$ special case. $(3)$ \textbf{Equivariant refinement (recorded, not needed scalar-wise).} Where $(M,g)$ carries an isometry group $\mathcal G$ with $\mathcal C=\mathcal G\cdot\gamma$ (the witness case), the normal Hessian is $\mathcal G_\gamma$-equivariant and the nullity carries the isotropy representation [math: equivariant Morse–Bott theory]; the selection-relevant scalar of stationary-phase weighting over $\mathcal C$ is the sign $(-1)^{\operatorname{ind}_B(\gamma)}$ with amplitude $\int_{\mathcal C}|\det{}'Q|^{-1/2}$ [math: Bott's explicit formula] — the sign datum Gap 18 asked about. $(4)$ \textbf{Monodromy type.} The linearized Poincaré return map $P_\gamma$ on the symplectic normal space of the geodesic flow ($2(n{-}1)$-dim; $2$-dim for surface slices): semisimple/elliptic ($|\operatorname{tr}|<2$ or diagonalizable), parabolic-unipotent ($P\ne I$, $(P-I)^2=0$), hyperbolic ($|\operatorname{tr}|>2$), or $P=I$ (fully degenerate). Iterates: $\operatorname{ind}_B(\gamma^m)$ is determined by $P_\gamma$'s conjugacy class through Bott's iteration theory [math]; for $P=I$ on a surface the ladder follows from first principles by normal-mode counting — $\operatorname{ind}_B(\gamma^m)=2m-1$ (FT.66(2), machine), matching the classical spherical iterate indices [math: Bott's computation for great circles; cited as consistency, not re-derived here]. $(5)$ \textbf{Vocabulary fold.} Def-H.19's "stable" $:=$ $\operatorname{ind}_B=0$ (clean), "unstable" $:=$ $\operatorname{ind}_B\ge1$; the interface content of Def-H.19 shrinks accordingly (the Morse-theoretic vocabulary is now [math]-constructed; only the \textit{non-coupling} clause remains declared — itself narrowed in part by FT.66(4)'s derived anti-coupling for absolute rules on closed chains).

\vspace{2pt}\noindent\rule{\textwidth}{0.4pt}\vspace{2pt}

\end{definition}

\begin{theorem}[{\texorpdfstring{FT.66 — The Bott correction: witness re-evaluation, reconstructability, and the inadmissibility of absolute selection}{FT.66 - The Bott correction: witness re-evaluation, reconstructability, and the inadmissibility of absolute selection}}]\label{tag:FT-66}
\textbf{Thm.} $(1)$ \textbf{Framework.} Under Def-H.24, stationary ledger channels carry the Bott data $(\operatorname{ind}_B,\operatorname{nul},\mathcal C,[P_\gamma])$; for degenerate (isometry-orbit) families the stationary-phase weight is Bott's: sign $(-1)^{\operatorname{ind}_B}$, amplitude over $\mathcal C$ [Def-H.24(3); math]. The naive conjugate-point count of FT.61(3) is the fixed-endpoint special case; for closed channels the periodic count and the orbit quotient must replace it — executed in (2). $(2)$ \textbf{The witness, re-evaluated (degeneracy critique absorbed; blindness Bott-stable).} On $M=S^2(R)\times S^1(2R)$ [FT.61(3)'s geometry]:

%dsp-level:0
\[
\gamma_A\ (m{=}2\ \text{great circle}):\quad\mathcal C_A=\mathbb{RP}^2\times S^1_{\mathrm{phase}}\times S^1_{\mathrm{trans}}\ (\dim 4),
\]

%dsp-level:0
\[
\qquad \operatorname{ind}_B=3\ (=2m{-}1),
\]

%dsp-level:0
\[
\qquad \operatorname{nul}=4\ (\text{clean}),
\]

%dsp-level:0
\[
\qquad P_{\gamma_A}=I_2;
\]

%dsp-level:0
\[
\gamma_B\ (\text{fiber circle}):\quad\mathcal C_B=S^2\times S^1_{\mathrm{phase}}\ (\dim 3),\qquad\operatorname{ind}_B=0\ (\forall\ \text{iterates}),\qquad\operatorname{nul}=3\ (\text{clean}),\qquad P_{\gamma_B}=\begin{pmatrix}1&L\\0&1\end{pmatrix}\ (\text{unipotent, parabolic});
\]

derivations (machine, exact): $\gamma_A$'s normal modes on the rank-1 normal bundle: $Q$-eigenvalues $\lambda_n\propto n^2-m^2$ on period-$L=2\pi Rm$ Fourier modes — negatives $\{|n|<m\}$: $2m-1$ of them; zeros $\{|n|=m\}$: the $2$-dim $\mathbb{RP}^2$ family motion; plus phase and translation tangents: $\operatorname{nul}=4=\dim\mathcal C_A$ — independent cross-check: Bott's iteration theorem at $P=I$ [math]. $\gamma_B$'s normal bundle is the trivial $TS^2$-pair with vanishing mixed curvature [product; math]: $Q=\int|J'|^2\succeq0$, kernel $=$ the $2$ constant basepoint fields: $\operatorname{ind}_B=0$, $\operatorname{nul}=3=\dim\mathcal C_B$; monodromy from $J=a+bt$: unipotent, iterates $\bigl(\begin{smallmatrix}1&mL\\0&1\end{smallmatrix}\bigr)$, Jordan type separating it from $P_{\gamma_A}=I_2$ at identical trace/det $(2,1)$ [machine]. \textbf{Verdicts}: (a) the Bott indices still differ, $3$ vs $0$, and the semiclassical signs would differ, $(-1)^3=-1$ vs $(-1)^0=+1$; (b) the ledger records and weights remain identical at every scale $\mathfrak h=2\pi R/N$: $(\kappa_0,\delta h,\Theta,n_-,\aleph)=(1,2N,2N,0,0)$, $w=+1$ both [FT.61(1),(3) re-run]; (c) hence \textbf{FT.61(3)–(4) hold with "Morse index" read as "Bott index"} — index-blindness of the weight inventory is Bott-stable, and the monodromy type (a conjugacy-class datum strictly finer than the index) is equally unregistered: no factor of $w(C)$ [A8; FT.61(1)'s complete inventory] can carry $(-1)^{\operatorname{ind}_B}$ or $[P_\gamma]$. (d) \textit{Homotopy-class note}: $\gamma_A$ is contractible, $\gamma_B$ is not — the class is combinatorially present in the configuration [D0.10: $E(C)$ ordered], but (i) no validity/weight/mask clause reads it [A8, D2.11, D2.3/CT2.4: inventory], and (ii) it is not the index: within a fixed free homotopy class, stationary members of different $\operatorname{ind}_B$ exist [math: length-tuning by intermediate value + neck-deformation focusing], so no class read implements index selection. $(3)$ \textbf{Reconstructability, and the true shape of exit (a).} The ambient configuration determines the metric: $C\mapsto g_{\mu\nu}(C)$ [CT1.20's own grouping map; CT1.12: $g$ from per-edge sector coordinates $(a_{00},q_{ij},w_i)$ [D2.1, D1.6, Def-H.14's arena]], sampled at $O(\mathfrak h)$ [CT1.19], coarse-grained to $(M,g)$ under CT2.6's hypotheses; hence $R^\rho{}_{\sigma\mu\nu}$ [CT1.14] along and near the channel, the normal Jacobi operator, and $(\operatorname{ind}_B,\operatorname{nul},[P_\gamma])$ are \textbf{in-principle functionals of the ambient configuration} in the CT2.6 regime — the datum is reconstructable, \textit{not} missing. Integer-valuedness plus $O(\mathfrak h)$ convergence gives eventual $\mathfrak h$-stability away from clean-bifurcation loci [FT.38(2)'s rigidity pattern; machine: minimal mode gap $|n^2-m^2|\ge1$ in scaled units — perturbations below half the gap cannot flip signs]; at the loci themselves the index jumps — those are precisely the finite-$\mathfrak h$ residue /18. What is missing is a \textbf{structural coupling}: no clause of A8's weight [FT.61(1)], D2.11's validity, D2.3's masks, or CT2.5's pole condition reads the datum; and the resolvent side is \textit{sign-free by hermiticity} — poles are conditions $\lambda_a=1$ on a real spectrum with residues $-\langle\mathrm{out}|P_a|\mathrm{in}\rangle$ [FT.26(1), CT2.5]: no per-saddle sign can live there without extending the structure. FT.61(4)'s exit (a) is accordingly reclassified: from "new variational registration (interface-definition species extending D0.7)" to "\textbf{a declared weight extension} $w\mapsto w\cdot(-1)^{\operatorname{ind}_B}$ (interface-definition species; well-definedness in the continuum regime supplied by the eventual-constancy lemma above) — constructible, not adopted here"; exit (b) (spectral-side positivity) unchanged, not constructed. $(4)$ \textbf{The selection dilemma (absolute index-0 rules are inadmissible for closed chains).} Let a floor's ambient slice realization be a topological sphere carrying $K\ge0$ with $K>0$ somewhere — the class instantiated by every Platonic cone-sphere realization: faces flat, conical defects $\delta_v>0$ at vertices, total $V\delta_v=4\pi>0$ [FT.63(4); smoothed approximation; math]. On any such slice:

%dsp-level:0
\[
\text{every non-constant closed geodesic has}\ \operatorname{ind}_B\ge1
\]

%dsp-level:0
\[
\qquad [\text{math: Synge/shortening — every closed geodesic on}\ S^2\ \text{bounds a disk and is exceeded by the disk-sweep};\ \text{machine: the latitude-push second variation}\ L''(0)=-2\pi/R<0\ \text{exactly},\ \text{and the round-sphere ladder}\ 2m{-}1\ge1\ \forall m].
\]

Hence an absolute selection rule — "index-0 stationary channels materialize; $\operatorname{ind}_B\ge1$ channels are suppressed" — admits \textbf{no} closed stationary channel on such a slice. Dilemma against M1/M2 (declared face/digon materialization [Def-H.20]): \textbf{horn 1} — if the declared materializations contribute semiclassically [FT.53(2)'s regime: stationary $=$ geodesic support], the absolute rule suppresses them all, contradicting M1 — the rule is inadmissible; \textbf{horn 2} — if they do not (Platonic face loops are corner loops of the cone metric, $\neg$geodesic: non-stationary in the smooth read), then closed-loop materialization lives entirely in the finite-$\mathfrak h$ non-stationary residue [FT.53(4)], where no stationary-index criterion applies at all — stability selection is \textit{vacuous} for bound chains. On either horn, \textbf{"stable jumps materialize" is refuted as the mechanism for closed (bound) chains} in this curvature class. $(5)$ \textbf{Verdict.} The four sub-questions  now carry derived answers: variational side — closed [FT.48(1)–(3)]; weight/ledger blindness — closed [FT.61(1)–(3)], Bott-stable [(2) here]; reconstructability — settled [(3) here]: the datum exists in principle, the coupling does not, and the resolvent side is sign-free; selection for \textbf{closed chains} — answered \textbf{negatively} [(4)]: no index criterion can operate, and the absolute rule is inadmissible against M1. This residue is the chain's honest end: closing it further requires a construction or declaration outside the frozen inventory — the same species boundary as Def-H.37(1)/Def-H.37(2)/Def-H.38. $(6)$ \textbf{Folds.} Def-H.19's vocabulary reads as Morse–Bott [Def-H.24(5)] — its [math] content constructed; the non-coupling clause is narrowed in part by clause (4)'s derived anti-coupling for absolute closed-chain rules; FT.48(3)'s trichotomy is scoped to clean fixed-endpoint problems (conjugate endpoints degenerate: read via Def-H.24); FT.61(3)–(4) hold per (2)–(3).

\textbf{Req.} Def-H.19, Def-H.21, Def-H.24; A8; D0.10, D1.6, D2.1, D2.3, D2.11; CT1.12, CT1.13, CT1.14, CT1.19, CT1.20, CT2.3, CT2.4, CT2.5, CT2.6, CT2.8 (Part 2); FT.22, FT.26, FT.31, FT.33, FT.38, FT.48, FT.53, FT.61, FT.62, FT.63; Def-H.14, Def-H.16, Def-H.20; M1/M2 (interface clause); Def-H.19; Morse–Bott theory, Bott's iteration theorem for spherical geodesics, Synge's theorem / disk-sweep shortening, Jordan-form conjugacy invariants, Gauss–Bonnet for cone metrics, intermediate-value length tuning [math].

\textbf{Pf.}

%dsp-level:0
\[
(1):\ \mathrm{Def\text{-}H.24}(2)\text{--}(3);
\]

%dsp-level:0
\[
\qquad \text{stationary-phase over a clean critical manifold}:\ \text{sign}
\]

%dsp-level:0
\[
= (-1)^{\operatorname{ind}_B},
\]

%dsp-level:0
\[
\ \text{amplitude}
\]

%dsp-level:0
\[
= \textstyle\int_{\mathcal C}|\det{}'Q|^{-1/2}\ [\text{math: Bott}]
\]

%dsp-level:0
\[
(2)\ \gamma_A:\ Q\ \text{on periodic normal sections}:\ V
\]

%dsp-level:0
\[
= e^{2\pi i nt/L}e_2(t):\ Q(V)
\]

%dsp-level:0
\[
= (\tfrac{n^2}{m^2}-1)\tfrac1{R^2}\|V\|^2\ [\text{Jacobi}\ f''+R^{-2}f=0\ \text{on the sphere};\ \text{math}]\implies\operatorname{ind}_B
\]

%dsp-level:0
\[
= \#\{|n|<m\}
\]

%dsp-level:0
\[
= 2m-1,
\]

%dsp-level:0
\[
\ \ker\qquad = \{|n|=m\}:\dim2
\]

%dsp-level:0
\[
\operatorname{nul}(\gamma_A)=2+1+1=4=\dim(\mathbb{RP}^2\times S^1\times S^1)\ [\text{orbit tangents};\ \text{machine}];
\]

%dsp-level:0
\[
\qquad \text{consistency}:\ \text{periodic count}
\]

%dsp-level:0
\[
= \text{fixed-endpoint conjugate count}
\]

%dsp-level:0
\[
= \#\{j\pi R:\ 1\le j\le2m-1\}
\]

%dsp-level:0
\[
= 2m-1\ [\text{machine}];
\]

%dsp-level:0
\[
\ \text{ladder}\ 1,
\]

%dsp-level:0
\[
3,
\]

%dsp-level:0
\[
5,
\]

%dsp-level:0
\[
\dots\ \text{matches the classical spherical iterates}\ [\text{math}]
\]

%dsp-level:0
\[
\gamma_B:\ R(\cdot,\dot\gamma_B)\dot\gamma_B|_{TS^2}=0\ [\text{product curvature};\ \text{math}]\implies J=a+bt;\ \text{periodic}\implies b=0;
\]

%dsp-level:0
\[
\quad Q=\int|J'|^2\succeq0\implies\operatorname{ind}_B=0;\quad\operatorname{nul}=2+1=3=\dim(S^2\times S^1)
\]

%dsp-level:0
\[
P_{\gamma_A}=I_2:\ J(2\pi R)=J(0)\cos2\pi+RJ'(0)\sin2\pi=J(0)\ [\cos2\pi=1,\sin2\pi=0\ \text{exact}];\qquad P_{\gamma_B}^m=\bigl(\begin{smallmatrix}1&mL\\0&1\end{smallmatrix}\bigr)\ [\text{machine: integer powers}]
\]

%dsp-level:0
\[
\text{Jordan separation}:\ P_{\gamma_A}-I
\]

%dsp-level:0
\[
= 0\ \wedge\ P_{\gamma_B}-I
\]

%dsp-level:0
\[
\ne 0\ \wedge\ (P_{\gamma_B}-I)^2
\]

%dsp-level:0
\[
= 0:\ \neg\text{conjugate despite}\ \operatorname{tr}
\]

%dsp-level:0
\[
= 2,
\]

%dsp-level:0
\[
\det
\]

%dsp-level:0
\[
= 1\ \text{both}\ [\text{machine}]
\]

%dsp-level:0
\[
\text{blindness}:\ \text{records}\ (1,2N,2N,0,0)\ \text{equal},\ w=(+1)\ \text{equal}\ [\mathrm{FT.61}(1),(3);\ \text{machine re-run}];
\]

%dsp-level:0
\[
\qquad (-1)^3\ne(-1)^0;\quad\neg\exists\ \text{factor of}\ w(C)\ \text{carrying it}\ [\mathrm{FT.61}(1)\ \text{inventory complete}]
\]

%dsp-level:0
\[
(2d):\ \pi_1(S^2\times S^1)=\mathbb Z:\ \gamma_A\mapsto0,\ \gamma_B\mapsto1\ [\text{math}];
\]

%dsp-level:0
\[
\qquad \text{class unread}:\ \mathrm{A8}\ w(C)\ \text{factors},\ \mathrm{D2.11}\ \text{clauses},\ \mathrm{D2.3}\ \text{blocks}:\ \neg\exists\ \text{homotopy entry}\ [\text{inventory}]
\]

%dsp-level:0
\[
\text{class}\ne\text{index}:\ \text{fixed class},
\]

%dsp-level:0
\[
\ \text{neck-bump family}:\ L_{\text{over}}(t)\ \text{continuous},
\]

%dsp-level:0
\[
\ \text{crosses}\ L_{\text{flat}}\ [\mathrm{IVT}];
\]

%dsp-level:0
\[
\ \text{focusing gives}\ \operatorname{ind}_B
\]

%dsp-level:0
\[
\ge 1\ \text{on the over-route}\ [\text{math: comparison}]
\]

%dsp-level:0
\[
(3):\ C
\]

%dsp-level:0
\[
\mapsto g_{\mu\nu}(C)\ [\mathrm{CT1.20}\ \mathrm{Pf}\ \text{step 1};\ \mathrm{CT1.12}/\mathrm{D2.1}:\ \text{sector coordinates};\ \mathrm{Def\text{-}H.14}:\ \text{edge arena}];
\]

%dsp-level:0
\[
\quad\Gamma\ [\mathrm{CT1.13}],
\]

%dsp-level:0
\[
\ R\ [\mathrm{CT1.14}]
\]

%dsp-level:0
\[
\implies Q_\gamma,
\]

%dsp-level:0
\[
\ P_\gamma\ \text{determined in the CT2.6 regime}
\]

%dsp-level:0
\[
\text{stability}:\ \lambda\text{-gaps}\ge1\ \text{in scaled units}\ [\text{machine}:\ \min|n^2-m^2|=1];
\]

%dsp-level:0
\[
\quad |\Delta\lambda|=O(\mathfrak h)\ [\mathrm{CT1.19}]\implies\exists\mathfrak h_0:\ \text{signs fixed}\ \forall\mathfrak h<\mathfrak h_0\ [\mathrm{FT.38}(2)\ \text{pattern}]
\]

%dsp-level:0
\[
\text{no reader}:\ \mathrm{A8}\ [\mathrm{FT.61}(1)];
\]

%dsp-level:0
\[
\quad\mathrm{D2.11}:\ \text{closedness/Janus/mask};
\]

%dsp-level:0
\[
\quad\mathrm{CT2.5}/\mathrm{FT.26}(1):\ \text{Hermitian}\implies\lambda_a\in\mathbb R,
\]

%dsp-level:0
\[
\ \text{pole}
\]

%dsp-level:0
\[
= \{\lambda_a=1\}:\ \text{sign-free}
\]

%dsp-level:0
\[
(4):\ \text{cone sphere}:\ K
\]

%dsp-level:0
\[
= 0\ \text{on faces},
\]

%dsp-level:0
\[
\ \delta_v>0\ \text{at vertices},
\]

%dsp-level:0
\[
\ \textstyle\sum\delta_v
\]

%dsp-level:0
\[
= 4\pi\ [\mathrm{FT.63}(4);\ \text{machine}];
\]

%dsp-level:0
\[
\quad\text{smoothing}:\ K
\]

%dsp-level:0
\[
\ge 0,
\]

%dsp-level:0
\[
\ K>0\ \text{somewhere}\ [\text{math}]
\]

%dsp-level:0
\[
\text{Synge}:\ \text{closed geodesic}\ \gamma\ \text{bounds a disk}\ D:\ \text{disk-sweep shortening}\implies\operatorname{ind}_B\ge1\ [\text{math}];
\]

%dsp-level:0
\[
\qquad \text{machine instances}:\ L''(0)=-2\pi/R<0;\quad2m-1\ge1\ \forall m\ge1
\]

%dsp-level:0
\[
\text{horn 1}:\ \text{absolute rule}\implies\neg\exists\ \text{admissible closed stationary channel}\implies\text{M1's semiclassical reading empty}:\ \text{rule inadmissible};
\]

%dsp-level:0
\[
\qquad \text{horn 2}:\ \text{face loops}
\]

%dsp-level:0
\[
= \text{corner loops}\ [\text{interior angles}<\pi\ \text{at cone vertices}]:\ \neg\text{geodesic}\implies\text{finite-}\mathfrak h\ \text{residue}\ [\mathrm{FT.53}(4)]
\]

%dsp-level:0
\[
(4a):\ \pi_1=0\implies\text{single free class}:\ \text{class-min}=\text{global min}=\text{constant loops (degenerate)}\ [\text{math}];
\]

%dsp-level:0
\[
\qquad (4b):\ \text{open minimizers exist}\ [\mathrm{FT.48}(3):\ \text{index-0}\iff\text{local}\ \ell\text{-minimizer, clean case}]
\]

%dsp-level:0
\[
(5):\ \text{conjunction of FT.48, FT.61, (2)--(4)};\qquad(6):\ \text{registry/patch moves as listed}.\ \blacksquare
\]

\textbf{Rmk (what would reverse the verdicts).} Clause (2)'s blindness is an inventory theorem: a derived cross-vertex or cross-channel bridge registering curvature along channels would reopen it — none exists [inventory; FT.59's absorption picture at locked rungs points the other way]. Clause (4) is conditional on (i) M1/M2 (interface clause) and (ii) the Synge-class slice realization of the floors — the class their own defect cards instantiate [FT.63(4)]; a flat or negatively curved ambient realization of a Platonic floor would evade the dilemma's horn structure (index-0 closed geodesics exist there [flat tori; math]) — but then the floor's \textit{own} cone-sphere ledger ($\sum\delta_v=4\pi$) would not be the ambient curvature read, and the defect card's $\Theta_R$ values would belong to a different geometric object than the channels' Jacobi data: recorded as a coherence condition any such alternative must address, not dismissed.

\textbf{Rmk (the adversarial exchange, recorded).} The critique — degenerate witnesses invalidate naive-index statements — is correct as to framework and void as to conclusion: the witness pair is clean Morse–Bott, the indices survive the correction unchanged ($3$ vs $0$), and every downstream claim of FT.61 holds with the refined index. The correction nevertheless produced three results the naive framing could not reach: the reconstructability lemma (exit (a) reclassified), the Jordan-type refinement (monodromy strictly finer than index, equally unregistered), and the Synge dilemma (absolute selection inadmissible, not merely unavailable) — the critique is folded in as FT.66, per the program's rigor-both-ways practice.

\textbf{Phys.} Two loops can wind through the geometry side by side, one thrice-unstable and one the very floor of its class, and the ledger writes down the same four symbols for both — not because the ledger is coarse, but because stability is a fact about the space around the loop, and the loop's own record carries only its length quanta and its signs. Worse for the hopeful selector: on any sphere-like slice with the Platonic floors' own positive curvature budget, \textit{no} closed geodesic is stable at all — every one of them can be shortened somewhere — so a rule that materialized only stable closed loops would materialize nothing, contradicting the very floors the program computes. Stability, then, cannot be the gate for bound chains; it remains a live candidate only for open jumps between termini, where short arcs genuinely minimize — and there the gate, if any, is either a declared sign carried by the weight or a positivity fact on the spectrum side. Both are named; neither is built; the boundary is exactly the species boundary of this document's other interface transfers.

\vspace{2pt}\noindent\rule{\textwidth}{0.4pt}\vspace{2pt}

\end{theorem}

\phantomsection
\section[Tier F.23 — The Virtual Bypass: Path Integration over…]{\texorpdfstring{Tier F.23 — The Virtual Bypass: Path Integration over Crystallization Patterns}{Tier F.23 - The Virtual Bypass: Path Integration over Crystallization Patterns}}

\begin{definition}[{\texorpdfstring{Def-H.25 — Skeleton configuration sums and the bypass vocabulary}{Def-H.25 - Skeleton configuration sums and the bypass vocabulary}}]\label{tag:Def-H-25}
\textbf{Def.} A \textbf{skeleton} $\Gamma$ $:=$ a finite bidirectional potential-edge set — the 1-skeletons of Def-H.20 (Platonic floors) and FT.45 (polygon floors), or any finite subgraph thereof. A \textbf{crystallization pattern} is a subset $R\subseteq E_{\mathrm{dir}}(\Gamma)$ [A4 per directed pair; D0.1: $\mathcal R$ membership is per-pair]; at $v$: $k_v:=\#\{e\in R: e\ni v\}=\partial^{\leftarrow}_R(v)+\partial^{\rightarrow}_R(v)$ [D0.8], potential degree $n(v)=| \{e\in E_{\mathrm{dir}}(\Gamma): e\ni v\} |$. Classification [CT0.11, CT4.1]: $v$ \textbf{channel-admissible} $\iff k_v\le3$; $v$ \textbf{$D^-$-locked} $\iff k_v\ge4$. A pattern is a \textbf{bypass configuration} iff every vertex is admissible while some has $n(v)\ge4$; the \textbf{bypass deficit} at $v$ is $n(v)-k_v$ (the virtual edges at $v$). The \textbf{skeleton configuration sum} (A8's weight restricted to skeleton-supported configurations [A8, D0.10], at CT1.20's grouping layer):

%dsp-level:0
\[
Z_\Gamma\ :=\ \sum_{R\subseteq E_{\mathrm{dir}}(\Gamma)}\ \sum_{\sigma\in\{\pm1\}^R} w(R,\sigma),
\]

%dsp-level:0
\[
\qquad w(R,\sigma)=\mathcal J\cdot\prod_{e\in R}G(e)\cdot(-1)^{\aleph(\sigma,R)},
\]

%dsp-level:0
\[
\qquad G(e)=\sigma_e\,e^{i\pi\kappa_e}\ [\mathrm{D3.5},\ \mathrm{FT.9}],
\]

split $Z_\Gamma=Z_{\mathrm{byp}}(\Gamma)+Z_{\mathrm{lock}}(\Gamma)$ by lock status (any $k_v\ge4$). \textbf{Scope declarations, marked at every use}: \textbf{(S1)} minimal-edge class $\kappa_e\equiv1$ on real edges [D1.1, Comm-min]; for general on-shell integer $\kappa$: every theorem below is unchanged with values scaled by $(-1)^{\sum_e\kappa_e}$ — the cover classification is $\kappa$-independent. \textbf{(S2)} $\aleph$ in D3.3's literal type-$(1,1)$ domain; $\aleph_v=0$ at all other types; the chain-local extension [FT.21(i), Open Gaps 16(d)] would change the survivor class — flagged load-bearing (Rmk of FT.67). \textbf{(S3)} seal/Janus bookkeeping suppressed: the sum is over static patterns; $\mathcal J=1$ identically [D3.6]; universe-level balance $\#J^+=\#J^-$ [D1.3, D2.11] is ambient data, not a skeleton constraint. \textbf{(S4)} virtual edges contribute weight $1$ and no phase — no registration exists for them [A4, D0.5/A7: $\phi_e$ is defined on $e\in E_{\mathcal R}$]. Finite skeleton $\Rightarrow$ finite sum: convergence is trivial; growth-rate statements are made per family (FT.68(7)).

\vspace{2pt}\noindent\rule{\textwidth}{0.4pt}\vspace{2pt}

\end{definition}

\begin{theorem}[{\texorpdfstring{FT.67 — The exact-cover theorem and the single-vertex bypass}{FT.67 - The exact-cover theorem and the single-vertex bypass}}]\label{tag:FT-67}
\textbf{Thm.} Under Def-H.25 with (S1)–(S4): $(1)$ \textbf{The $\aleph$ identity.} At a type-$(1,1)$ vertex $v$ with in-edge $e_\leftarrow$, out-edge $e_\rightarrow$: $(-1)^{\aleph_v}=(-1)^{|f_v|}=\sigma_{e_\leftarrow}\sigma_{e_\rightarrow}$ [D3.3: $f_v=\tfrac12(\sigma_\rightarrow-\sigma_\leftarrow)\in\{-1,0,1\}$; $\sigma$'s binary [CT0.9]; machine]. $(2)$ \textbf{Exact-cover theorem.} For $e\in R$ let $\operatorname{cov}(e):=\#\{(1,1)\text{-type endpoints of }e\text{ in }R\}\in\{0,1,2\}$. Then the weight is the monomial

%dsp-level:0
\[
w(R,\sigma)=(-1)^{\sum_e\kappa_e}\prod_{e\in R:\ \operatorname{cov}(e)\ \mathrm{even}}\sigma_e\ ,
\]

%dsp-level:0
\[
\qquad [\text{display correction: the}\ \sigma_e\text{-exponent is}\ 1+\operatorname{cov}(e),\ \text{so the surviving factors are the even-coverage edges — as this clause's own Pf records; classification and values unchanged}]
\]

and the sign sum evaluates exactly:

%dsp-level:0
\[
\sum_{\sigma\in\{\pm1\}^R}w(R,\sigma)=\begin{cases}(-1)^{\sum_e\kappa_e}\,2^{|R|}=4^{\,|\mathcal C|}\,(-1)^{\sum\kappa} & \text{if}\ \operatorname{cov}(e)=1\ \forall e\in R\quad(\textbf{pass-through cover};\ \mathcal C:=\{(1,1)\text{-vertices}\}),\\[2pt] 0 & \text{otherwise,}\end{cases}
\]

with $|R|=2|\mathcal C|$ forced on covers [each center contributes its in- and out-edge; every edge once]. At (S1): value $=4^{|\mathcal C|}>0$. [Machine: exhaustive on the tetrahedron (all $4096$ patterns $=$ all $3^{12}$ (pattern,sign) pairs), on $P_3,P_4,P_5$ exhaustively, on the cube exhaustively for $K\le4$ ($12831$ configs) plus a deterministic stride sample.] $(3)$ \textbf{Structural corollaries (what cancels, what survives).} (a) Every directed cycle cancels: along a cycle of $(1,1)$-vertices each edge has $\operatorname{cov}=2$ — \textbf{no bound cell survives at the amplitude layer}: digons (the $\Sigma_{\min}$ pattern [FT.28]), directed cycles (the $\Sigma_L$ pattern [FT.32]), face boundaries, and the fully crystallized Def-H.20 floors (all vertices type $(d,d)$, $d\ge2$: no centers at all) all have $\sum_\sigma w=0$; (b) single real edges and length-1 threads cancel ($\operatorname{cov}=0$); (c) threads of length $\ge3$ cancel (interior edges double-covered); (d) \textbf{survivors $=$ V-thread networks}: each center $v\in\mathcal C$ pairs one in-edge $(s\to v)$ with one out-edge $(v\to t)$; non-centers are sources $(0,1)$, sinks $(1,0)$, or hubs of any type $\ne(1,1)$ — \textit{including locked hubs} ($k\ge4$) whose every edge is paired by a distinct center; every survivor has even $|R|=2|\mathcal C|$ and value $4^{|\mathcal C|}$. $(4)$ \textbf{Single-vertex enumeration ($n\in\{4,5,6\}$, balanced types).} $2^n$ patterns, $\binom nk$ at size $k$; admissibility fractions

%dsp-level:0
\[
f_{\mathrm{adm}}(4)=\tfrac{15}{16},\quad f_{\mathrm{adm}}(5)=\tfrac{13}{16},\quad f_{\mathrm{adm}}(6)=\tfrac{21}{32}\qquad\Bigl[\ =\sum_{k\le3}\binom nk/2^n\Bigr];
\]

sign-summed amplitudes: $Z_{\mathrm{byp}}^{\pm}(v,n)=1+4\,n_{\mathrm{in}}n_{\mathrm{out}}=(17,25,37)$ — vacuum $+$ the $\binom{n_{\mathrm{in}}}{1}\binom{n_{\mathrm{out}}}{1}$ pass-through patterns at $4$ each; \textbf{only $k\in\{0,2\}$ survive}; $Z_{\mathrm{lock}}^{\pm}(v,n)=0$ (locked $\Rightarrow$ type $\ne(1,1)$ $\Rightarrow$ uncovered). At $\sigma\equiv+1$: $Z_{\mathrm{byp}}=-\binom{n-1}{3}\in\{-1,-4,-10\}$, $Z_{\mathrm{lock}}=+\binom{n-1}{3}$, total $(1-1)^n=0$ [alternating partial-binomial identity; binomial theorem; machine]. Aspect activation per pattern [CT0.13]: $\mathcal A(0)=\{J\}$, $\mathcal A(1)=\{J,G,W,H,D^+\}$, $\mathcal A(2)=\mathcal A(1)\cup\{\mathrm{EM}\}$, $\mathcal A(3)\cup\{\text{strong}\}$; mask dimensions $\dim\mathcal K_k=(6,16,17,25)$ [CT0.14]; observable strong never fires at a single vertex [FT.27]; EM is state-dependent: $f_v\ne0$ on exactly half of the pass-through sign assignments [D3.3; machine]. $(5)$ \textbf{Virtual bypass theorem (single-vertex form).} Every vertex of potential degree $n\ge4$ with $n_{\mathrm{in}},n_{\mathrm{out}}\ge1$ has strictly positive admissible-sector amplitude, $Z_{\mathrm{byp}}^{\pm}=1+4n_{\mathrm{in}}n_{\mathrm{out}}>0$, independent of the $\kappa$-class up to the global sign $(-1)^{\sum\kappa}$; its locked sector vanishes identically. \textbf{[scope: the locked-sector clause is a literal-domain (S2) theorem — in the chain-local domains E1/E2 of Def-H.27 an isolated locked vertex has strictly positive amplitude, $Z^{\pm}_{\mathrm{lock}}=(32,96,672)$ for $n=(4,5,6)$, with the locked digon-hub family at $(!a)4^a$ [FT.75(1)]. The bypass clause is domain-invariant.]} \textbf{$D^-$ locking is avoided with nonzero amplitude by partial crystallization} — the bypass mechanism is real at the amplitude layer, and its surviving content is exactly the pass-through (transit) sector [D1.5: open, propagating $=$ bosonic content]. $(6)$ \textbf{Consistency anchors and the norm layer.} (a) The $k{=}2$ survivor is the $\Sigma_{\min}$ pass-through structure verbatim [FT.28(5): $q$ of type $(1,1)$]; (b) FT.61(1)'s phase inventory $e^{i\Phi}=(-1)^{\Theta+\aleph+n_-}$ is reproduced termwise inside the monomial of (2) [machine]; (c) \textbf{no conflict with positivity}: $|w(R,\sigma)|=1$ for every pair, so CT2.8's norms $\langle A|A\rangle=\sum|w|^2>0$ never cancel — the cancellations of (2)–(3) are amplitude-layer ($Z$) facts; the spectral layer [Def-H.16: poles of the masked resolvent] and the materialization declarations M1/M2 are untouched (the layer bridge is item 13's residue — see FT.68(4)).

\textbf{Req.} A4, A7, A8; D0.1, D0.5, D0.8, D0.10, D1.1, D1.2, D1.3, D1.5, D2.11, D3.3, D3.5, D3.6; CT0.9, CT0.10, CT0.11, CT0.13, CT0.14 (Part 1); CT1.20, CT2.8 (Part 2); CT4.1, CT4.4 (Part 2); FT.9, FT.22, FT.24, FT.27, FT.28, FT.32, FT.61, FT.62; Def-H.16, Def-H.18, Def-H.20, Def-H.25; alternating partial-binomial sums [math].

\textbf{Pf.}

%dsp-level:0
\[
(1):\ \sigma\in\{\pm1\}\ [\mathrm{CT0.9}]:\ f_v
\]

%dsp-level:0
\[
= \tfrac12(\sigma_\rightarrow-\sigma_\leftarrow):\ |f_v|
\]

%dsp-level:0
\[
= 1\iff\sigma_\rightarrow\ne\sigma_\leftarrow\iff\sigma_\leftarrow\sigma_\rightarrow
\]

%dsp-level:0
\[
= -1
\]

%dsp-level:0
\[
= (-1)^{|f_v|};
\]

%dsp-level:0
\[
\quad\sigma_\leftarrow
\]

%dsp-level:0
\[
= \sigma_\rightarrow:\ |f_v|
\]

%dsp-level:0
\[
= 0,
\]

%dsp-level:0
\[
\ \sigma_\leftarrow\sigma_\rightarrow\qquad = +1\ \checkmark
\]

%dsp-level:0
\[
(2):\ w
\]

%dsp-level:0
\[
= (-1)^{\sum\kappa}\prod_{e\in R}\sigma_e\cdot\!\!\prod_{v\in\mathcal C}\!\!(\sigma_{e_\leftarrow(v)}\sigma_{e_\rightarrow(v)})\ [(1);\ \mathrm{D3.6}:\ \mathcal J=1]:\ \sigma_e\text{-exponent}
\]

%dsp-level:0
\[
= 1+\operatorname{cov}(e)
\]

%dsp-level:0
\[
\operatorname{cov}(e)=0\ \text{or}\ 2\ \text{for some}\ e:\ \text{flip}\ \sigma_e:\ \text{summand}\mapsto-\text{summand},\ \text{involution on}\ \{\pm1\}^R\implies\text{sum}=0;
\]

%dsp-level:0
\[
\qquad \operatorname{cov}\equiv1:\ \text{every}\ \sigma_e\ \text{appears squared}\implies w=(-1)^{\sum\kappa}\ \sigma\text{-constant}
\]

%dsp-level:0
\[
\implies\sum_\sigma w=(-1)^{\sum\kappa}2^{|R|};
\]

%dsp-level:0
\[
\qquad |R|
\]

%dsp-level:0
\[
= 2|\mathcal C|:\ \text{centers pair disjoint edges, covers all}\ [\text{each center reads exactly its 2 edges; each edge read once}];
\]

%dsp-level:0
\[
\quad (-1)^{2|\mathcal C|}2^{2|\mathcal C|}=4^{|\mathcal C|}\ \text{at}\ \kappa\text{-even sums}
\]

%dsp-level:0
\[
\kappa\text{-invariance}:\ \sigma_e\ \text{enters linearly}\ \forall\kappa_e\in\mathbb Z:\ \text{classification unaffected};
\]

%dsp-level:0
\[
\ \text{value scaled by}\ (-1)^{\sum\kappa}\ [\mathrm{S1}\ \text{scope}]
\]

%dsp-level:0
\[
(3a):\ \text{cycle}:\ \text{every vertex type}\ (1,1)\ \text{or higher along a closed walk};
\]

%dsp-level:0
\[
\ \text{cycle edges read by both endpoints when}\ (1,1)\ \text{at both}\implies\operatorname{cov}
\]

%dsp-level:0
\[
= 2;
\]

%dsp-level:0
\[
\ \text{any non-}(1,1)\ \text{vertex on the cycle leaves its edges with}\ \operatorname{cov}
\]

%dsp-level:0
\[
\ne 1\ \text{or forces type violation}\ [\text{machine-exhaustive}];
\]

%dsp-level:0
\[
\ \text{Def-H.20 floors}:\ \text{types}\ (d,d),
\]

%dsp-level:0
\[
\ d
\]

%dsp-level:0
\[
\ge 2:\ \mathcal C
\]

%dsp-level:0
\[
= \varnothing,
\]

%dsp-level:0
\[
\ R\ne\varnothing\qquad \implies 0
\]

%dsp-level:0
\[
(3b,c):\ \text{single edge}:\ \text{endpoints}\ (0,1),(1,0):\ \operatorname{cov}=0;
\]

%dsp-level:0
\[
\qquad \text{length-}\ge3\ \text{thread}:\ \text{interior vertices}\ (1,1)\ \text{read the shared interior edges twice}\implies\operatorname{cov}=2
\]

%dsp-level:0
\[
(3d):\ \text{cover}\iff\text{the pairing structure described};
\]

%dsp-level:0
\[
\ \text{hubs}:\ \text{non-centers with all edges center-covered},
\]

%dsp-level:0
\[
\ \text{any type}
\]

%dsp-level:0
\[
\ne (1,1),
\]

%dsp-level:0
\[
\ \text{locked allowed}\ [\text{machine witnesses: tetrahedron}\ K_{1,3}\ \text{stars},\ \text{cube locked coverings}]
\]

%dsp-level:0
\[
(4):\ \text{enumeration}\ [\text{machine exhaustive}];
\]

%dsp-level:0
\[
\qquad Z^{\pm}_{\mathrm{byp}}
\]

%dsp-level:0
\[
= 1\ (\text{vacuum})+\binom{n_{\mathrm{in}}}1\binom{n_{\mathrm{out}}}1\cdot4\ [\text{pass-throughs};\ \text{single-vertex model: no accidental}(1,1)\ \text{non-centers}];
\]

%dsp-level:0
\[
\quad k\ge4:\ \text{type}\ne(1,1):\ 0\ [(2)]
\]

%dsp-level:0
\[
\sigma\equiv+1:\ w=(-1)^k\ [\aleph=0]:\ Z_{\mathrm{byp}}=\textstyle\sum_{k\le3}\binom nk(-1)^k=(-1)^3\binom{n-1}3\ [\text{math}];
\]

%dsp-level:0
\[
\quad Z_{\mathrm{lock}}=\textstyle\sum_{k\ge4}\binom nk(-1)^k=+\binom{n-1}3;\quad\text{total}=(1-1)^n=0
\]

%dsp-level:0
\[
\text{aspects}:\ \mathrm{CT0.13}\ \text{rows verbatim at}\ k;\ \dim\mathcal K_k\ [\mathrm{CT0.14}];\ \text{strong}\ \neg\text{observable at}\ \{v\}\ [\mathrm{FT.27}];
\]

%dsp-level:0
\[
\quad f_v\ne0\iff\sigma_\leftarrow\ne\sigma_\rightarrow:\ 2/4\ \text{of pass-through signs}\ [\mathrm{D3.3}]
\]

%dsp-level:0
\[
(5):\ 1+4n_{\mathrm{in}}n_{\mathrm{out}}>0\ \forall n_{\mathrm{in}},n_{\mathrm{out}}\ge1;
\]

%dsp-level:0
\[
\qquad (6):\ \text{anchors as cited};\ |w|=1:\ \text{CT2.8 norms cancellation-free}.\ \blacksquare
\]

\textbf{Rmk (domain sensitivity — load-bearing).} The exact-cover classification is a theorem \textit{within scope (S2)}: it uses $\aleph_v$ at type-$(1,1)$ vertices only. A chain-local extension of $f_v$ to higher types [FT.21(i)'s telescoping reads] would let higher-type vertices pin edges, enlarging the survivor class (and un-cancelling some bound content). The computation therefore sharpens 16(d) from a bookkeeping question to an amplitude-level one: the two $\aleph$-domains give \textit{provably different} path integrals. No extension is adopted or guessed here. \textbf{[both canonical chain-local extensions are now defined and computed exactly — Def-H.27 (E1 chain-dominant, E2 literal-forced), FT.74 (extended exact-cover theorems, brute-force verified), FT.75 (consequences: floor revival, lock revival, the even-$L$ parity theorem, the triangle census, norm-layer invariance). They remain computed variants, not adopted; (S2) stays the frozen domain of record.]}

\textbf{Rmk (locked hubs — locking is not amplitude death).} Clause (3d) exhibits locked vertices \textit{inside} surviving configurations: a hub threaded by pass-through centers carries nonzero amplitude while remaining $D^-$-configured — locked content is retained and counted [CT4.3], inaccessible [D1.11], and now: amplitude-visible only through the threads that cross it. The single-vertex vanishing $Z^{\pm}_{\mathrm{lock}}(v,n)=0$ is a statement about isolated vertices; joint skeletons admit locked survivors (tetrahedron: $448$; cube: $1\,150\,848$ [FT.68]).

\textbf{Phys.} Sum the ledger's own signs over every way a skeleton can partially crystallize and a stark rule emerges: a configuration contributes if and only if its pass-through vertices tile its edges exactly once — every edge read by exactly one local phase-comparison, no edge read twice or never. Loops cannot be so tiled (their edges are read at both ends); neither can loose single edges. What survives is a gas of open two-step threads $s\to v\to t$, each contributing exactly $+4$, with locked hubs admissible only as threading junctions. Crystallization flicker therefore does not dress the solids — it dissolves them at the amplitude layer and leaves transit threads: existence, at this layer, is a perfect matching of edges onto pass-through comparisons.

\vspace{2pt}\noindent\rule{\textwidth}{0.4pt}\vspace{2pt}

\end{theorem}

\begin{theorem}[{\texorpdfstring{FT.68 — Joint skeletons: tetrahedron, the polygon ladder, the cube; the layer dichotomy; flicker}{FT.68 - Joint skeletons: tetrahedron, the polygon ladder, the cube; the layer dichotomy; flicker}}]\label{tag:FT-68}
\textbf{Thm.} Under Def-H.25/FT.67 (same scopes): $(1)$ \textbf{Tetrahedron (exhaustive).} $V=4$, potential $n(v)=6$ [FT.62(1)], $12$ directed edges. Consistency (each directed edge real at both endpoints or virtual at both — reality is one global fact per pair [A4, D0.1]) reduces the $4$-tuple space $2^{24}$ to $2^{12}=4096$ joint patterns. Bypass configurations: $1230$; \textbf{bypass probability} $\tfrac{615}{2048}$; $K\le6$ forced ($\sum_vk_v=2K\le4\cdot3$); counts by $K$: $(1,12,66,220,435,408,88)$ for $K=0..6$ (all-patterns: the symmetric row $\binom{12}K$ minus locked); orbits under $S_4$: $70$ total, by $K$: $(1,1,5,13,23,21,6)$. Amplitudes:

%dsp-level:0
\[
Z^{\pm}_{\mathrm{byp}}=1057=1+24\cdot4+60\cdot16,\qquad Z^{\pm}_{\mathrm{lock}}=448=12\cdot16+4\cdot64,\qquad Z^{\pm}_{\mathrm{all}}=1505;
\]

covering polynomials (zeta-derived values, FT.68(7)): $\mathcal Z_{\mathrm{byp}}(y)=1+24y+60y^2$, $\mathcal Z_{\mathrm{lock}}(y)=12y^2+4y^3$ — the locked sector includes the four fully bidirectional $K_{1,3}$ stars (locked $(3,3)$ hub threaded by three centers, $4^3=64$ each). \textbf{Effective mass} $m_{\mathrm{eff}}=\langle K\rangle\ln2/\mathfrak h$ (dark read: $E_{\mathrm{tot}}=K$ directed real edges [FT.17 dark pattern; Def-H.9(2)]), in three readings: (i) uniform over bypass configurations $\langle K\rangle=\tfrac{852}{205}$; (ii) amplitude-weighted $\bigl(\sum K|A_C|^2/\sum|A_C|^2\bigr)=\tfrac{62208}{15745}$; (iii) Cal-F equilibrium at $r{=}1$: equals (i). \textbf{Fractional contributions}: the equilibrium measure of a pattern with $K$ real edges is $p^K(1-p)^{12-K}$, $p=r/(1+r)$ — the per-$K$ coefficients above are the exact fractional-contribution table. Aspect-activation frequency over bypass configs by $\max_vk_v$: $(1,24,245,960)$ for $\max k=(0,1,2,3)$. Bound-content census (configurations containing $\ge1$): digons $501$, directed $3$-cycles $252$, directed $4$-cycles $90$ — with $\Theta\in\{2,3,4\}$ at (S1): digons and $4$-cycles are on-shell materializable candidates [$\Theta\in2\mathbb Z$: FT.22; $4$-cycles: $k=2$, $j=1$ under Def-H.38; digons: the $\Sigma_{\min}$ card, $j=\tfrac12$], $3$-cycles off-shell arena content [$\Theta=3$ odd; Def-H.18] — and \textbf{the bound cells themselves all cancel in amplitude} [FT.67(3a)]: standalone bound content is configuration-present, $Z$-absent. \textbf{(correction [FT.69(4)]: the census configurations are not all $Z$-absent — $48$ of the $501$ digon-containing bypass configurations are covers: dissolved digons inside thread networks [witness in FT.69(4)]; the $3$-cycle and $4$-cycle rows cancel completely — odd cycles never enter covers, and no tetrahedron cover contains a $4$-cycle [FT.69(3)].)} The fully crystallized floor cancels [machine]. $(2)$ \textbf{Polygon ladder $P_n$ (potential $n(v)=4$; transfer matrix).} Bypass $\iff$ every vertex drops $\ge1$ of its $4$ directed edges. With channel states $s_i=(f_i,r_i)\in\{0,1\}^2$ and the vertex constraint $\operatorname{pop}(s_i)+\operatorname{pop}(s_{i+1})\le3$:

%dsp-level:0
\[
B_{P_n}(t)=\operatorname{tr}\bigl(T(t)^n\bigr),\qquad T(t)[s][s']=t^{|s'|}\,[\,|s|+|s'|\le3\,]\qquad[\text{math: transfer matrix}],
\]

counts $n=3..10$: $(54,207,783,2970,11259,42687,161838,613575)$; fractions $(\tfrac{27}{32},\tfrac{207}{256},\tfrac{783}{1024},\tfrac{1485}{2048},\dots)$ decaying; the polynomials $B_{P_n}(t)$ are tabulated by the machine [zeta-derived values]. \textbf{Dynamical zeta}: $\operatorname{charpoly}T(1)=x^4-3x^3-3x^2$, $\det(I-zT(1))=1-3z-3z^2$ — the bypass counts satisfy the exact recursion $\operatorname{tr}_n=3\operatorname{tr}_{n-1}+3\operatorname{tr}_{n-2}$ [machine: $n=3..12$], growth rate $\lambda_{\max}=\tfrac{3+\sqrt{21}}2<4$: \textbf{the bypass fraction decays exponentially in $n$} — bypass is a finite-cell phenomenon in the thermodynamic limit (the amplitude sector itself keeps growing, at the covering growth rate). Sign-summed amplitudes (bypass, locked): $P_3$: $(25,48)$; $P_4$: $(353,64)$; $P_5$: $(1161,80)$; $P_6$: $(3121,1632)$; $P_7$: $(11929,10864)$; $P_8$: $(53825,43136)$; $P_9$: $(229417,159888)$; $P_{10}$: $(844881,775840)$ — the locked sector is nonzero for every $n\ge3$ (single-hub coverings exist on every polygon). $(3)$ \textbf{Cube (orbit-reduced; Directive bound respected).} Potential $n(v)=6$, $24$ directed edges, $\operatorname{Aut}$ of order $48$. Bypass configurations: $1\,580\,092$, by $K$: $(1,24,276,2024,10506,40296,115796,247704,383913,408464,270096,91392,9600)$ for $K=0..12$ (cap-admissible DP; $K\le12$ forced). \textbf{Orbit structure} (cap-aware Burnside DP over conjugacy classes): by $K$: $(1,1,11,49,250,863,2509,5211,8163,8571,5768,1922,245)$, total $33\,564>10^4$ — per the bounded-enumeration rule, orbits are reported by $K$ only, with individual weights computed for coverings alone. Covering polynomials: $\mathcal Z_{\mathrm{byp}}(y)=1+48y+648y^2+2672y^3+3624y^4$ ($6993$ configs), $\mathcal Z_{\mathrm{lock}}(y)=24y^2+728y^3+3192y^4+264y^5+4y^6$ ($4212$ configs); amplitudes $Z^{\pm}_{\mathrm{byp}}=1\,109\,313$, $Z^{\pm}_{\mathrm{lock}}=1\,150\,848$, $Z^{\pm}_{\mathrm{all}}=2\,260\,161$; the fully crystallized cube cancels. The exact-cover theorem is exhaustively verified for all bypass configs with $K\le4$ and stride-sampled beyond. Dodecahedron/icosahedron: \textbf{not run wholesale}; their covering enumeration (sparse by FT.67(3d)) is the natural next cell, deferred and recorded. $(4)$ \textbf{Layer dichotomy (the flicker hypothesis, resolved).} (a) \textit{Bypass is real}: $Z^{\pm}_{\mathrm{byp}}>0$ on every computed skeleton — single vertex $1+n_{\mathrm{in}}n_{\mathrm{out}}\cdot4$, tetrahedron $1057$, polygons growing, cube $1\,109\,313$: the admissible sector of an $n\ge4$ skeleton has strictly positive amplitude at the A8 configuration-sum layer. (b) \textit{What it produces}: the survivors are exactly open V-thread networks — $1$-dimensional transit-type content [D1.5: open propagating $=$ bosonic]; the effective dimensionality of the path-integrated bypass configuration is \textbf{thread-like: neither polygon-like $2$D (faces cancel) nor solid-like $3$D (solids cancel)}; its ledger reads are $m_{\mathrm{eff}}$ of (1), aspect activation concentrated at $\max k_v\in\{2,3\}$ ($\mathcal A\supseteq\{\mathrm{EM}\}$-capable and strong-\textit{activated} patterns [CT0.13], with observable strong still requiring $\Sigma$-level sealing [FT.27, FT.20]). (c) \textit{What it does not produce}: bound floors — $\Sigma_{\min}$, $\Sigma_L$, $\Sigma_{P_n}$, $\Sigma_S$ — cancel exactly at this layer [(FT.67(3a))]; they live at the spectral layer [Def-H.16 poles; M1/M2 declarations], where CT2.8's norms $\sum|w|^2$ are strictly positive and no cancellation occurs. (d) \textbf{Verdict on "the path integral fuzzes the flickering 2D shape into a higher-tier solid"}: flickering thread gas — \textit{yes}, exact and strictly positive; fuzzing into solids — \textit{no} at the amplitude layer \textbf{[scope: a literal-domain (S2) verdict — in the chain-local domains every computed Platonic floor and every even polygon floor carries strictly positive amplitude (tetra $2^{17}$, cube $2304\cdot2^{24}$, octa $72384\cdot2^{24}$; $4\cdot2^{2L}$ for even $L$, $0$ for odd), so the verdict is domain-relative [FT.75(2)]; the thread-gas clause and the norm-layer clause (c) are domain-invariant [FT.75(6)]]}; the solid floors enter through the spectral layer, and the bridge between the two layers is precisely the open residue  [FT.66(3)–(5): structural coupling absent; exit (b) not constructed]. No identification between layers is made or forced. $(5)$ \textbf{Flicker rate (Cal-F, declared).} The corpus supplies no $J^\pm$ rates — D1.2 is transition vocabulary, A5's draws are unrecorded, and no rate tag exists [inventory]. Declare $r:=\Gamma_{J^+}/\Gamma_{J^-}$ a free internal dynamics datum (P4-class, $\neg$constant of the theory) with per-edge independent detailed-balance equilibrium $p=r/(1+r)$ [Cal-F]. Then the vertex-admissibility fractions are the exact polynomials

%dsp-level:0
\[
f_{\mathrm{adm}}(n,p)=\sum_{k\le3}\binom nk p^k(1-p)^{n-k}:\quad1-p^4\ (n{=}4),\quad1-5p^4+4p^5\ (n{=}5),\quad1-15p^4+24p^5-10p^6\ (n{=}6),
\]

with samples $f_{\mathrm{adm}}(6;r{=}1)=\tfrac{21}{32}$, $f_{\mathrm{adm}}(6;r{=}3)=\tfrac{347}{2048}$ [machine]; skeleton bypass probabilities are $B_\Gamma(p)=\sum_Kc_K\,p^K(1-p)^{|E_{\mathrm{dir}}|-K}$ with $c_K$ the by-$K$ counts of (1)–(3) — the exact fractional-contribution tables. Time-fraction reading: at equilibrium each vertex spends the fraction $f_{\mathrm{adm}}(n(v),p)$ of time channel-admissible; the fully crystallized Def-H.20 floor occupies measure $p^{2E}\to$ small as $p<1$ — the static floor is a measure-zero-corner of its own skeleton's flicker ensemble, while its spectral status is $p$-independent (declared, M1/M2). $(6)$ \textbf{Bott connection.} Amplitude survivors contain no closed chains [FT.67(3a)] $\Rightarrow$ no closed stationary channels at the amplitude layer $\Rightarrow$ Def-H.24's Bott data is \textit{undefined on survivors} (no loops to index); the thread networks are transit-side objects, where fixed-endpoint index-0 minimizers exist [FT.48(3)] — the amplitude layer's survivor class and the variational layer's stable class are both transit-side, a derived consistency with FT.66(4b) ("stability is transit-side vocabulary"). Bound content inside bypass configurations (the census of (1)) carries its Bott datum at the spectral layer, where FT.66(4) applies unchanged: on Synge-class slices every closed geodesic has $\operatorname{ind}_B\ge1$ — \textbf{the bypass does not evade the selection dilemma; it relocates the surviving content to the one sector where index-0 stability exists}. $(7)$ \textbf{Convergence and zeta ledger.} Every $Z$ above is a finite exact integer (finite skeletons): the path integral converges trivially and has well-defined values — no regularization was needed and none is claimed. Derived zeta-type values recorded: polygon dynamical zeta $\det(I-zT(1))=1-3z-3z^2$ with recursion $\operatorname{tr}_n=3\operatorname{tr}_{n-1}+3\operatorname{tr}_{n-2}$ and $\lambda_{\max}=\tfrac{3+\sqrt{21}}2$; the covering polynomials $\mathcal Z_\Gamma(y)$ (tetrahedron/cube, per sector); the bypass polynomials $B_{P_n}(t)$, $n\le10$ — all exact integer data, machine-verified.

\textbf{Req.} Def-H.25, Def-H.24; FT.17, FT.20, FT.22, FT.27, FT.28, FT.32, FT.42, FT.45, FT.48, FT.61, FT.62, FT.63, FT.66, FT.67; A4, A5, A8; D0.1, D0.10, D1.2, D1.5, D2.11; CT0.11, CT0.13 (Part 1); CT1.20, CT2.8 (Part 2); CT4.1, CT4.3 (Part 2); Def-H.16, Def-H.18, Def-H.20; M1/M2 and Def-H.38/Cal-F (interface rows; census and flicker clauses only); transfer matrices, Burnside with constraint DP, Newton recursions, binomial identities [math].

\textbf{Pf.}

%dsp-level:0
\[
(1):\ \text{consistency}:\ \text{reality per pair}\ [\mathrm{A4}/\mathrm{D0.1}]\implies\text{joint patterns}=2^{12};
\]

%dsp-level:0
\[
\qquad \text{counts/orbits/census}:\ \text{machine exhaustive}\ (4096\ \text{patterns};\ 3^{12}\ \text{sign pairs in the theorem checks})
\]

%dsp-level:0
\[
Z^{\pm}\ \text{values}:\ \mathrm{FT.67}(2)\ \text{value formula summed over covering classes}\ [\text{machine}:\ \text{two independent enumerators agree}];
\]

%dsp-level:0
\[
\qquad K_{1,3}\ \text{star hubs}:\ \text{witnesses exhibited}\ [\text{machine}]
\]

%dsp-level:0
\[
m_{\mathrm{eff}}:\ \text{readings (i)--(iii) computed from the census}\ [\text{machine exact rationals}];
\]

%dsp-level:0
\[
\qquad \text{fractional contributions}:\ \text{Bernoulli measure coefficients}\ [\mathrm{Cal-F}]
\]

%dsp-level:0
\[
(2):\ T(t)\ \text{construction}:\ \text{vertex constraint}\iff\operatorname{pop}(s_i)+\operatorname{pop}(s_{i+1})\le3;
\]

%dsp-level:0
\[
\qquad B_{P_n}(t)=\operatorname{tr}T(t)^n\ [\text{math}];\quad\text{cross-check}:\ \text{cap-DFS brute force},\ n\le6\ [\text{machine}]
\]

%dsp-level:0
\[
\operatorname{charpoly}T(1)=x^2(x^2-3x-3)\ [\text{machine: exact interpolation}];
\]

%dsp-level:0
\[
\qquad \det(I-zT)=z^4\,\mathrm{ch}(1/z)=1-3z-3z^2;
\]

%dsp-level:0
\[
\qquad \text{Newton}:\ \operatorname{tr}_n=3\operatorname{tr}_{n-1}+3\operatorname{tr}_{n-2}\ [\text{machine}\ n\le12]
\]

%dsp-level:0
\[
Z^{\pm}\ \text{per sector}:\ \text{covering enumeration (escape-pruned backtracking)},
\]

%dsp-level:0
\[
\ n
\]

%dsp-level:0
\[
\le 10\ [\text{machine};\ \text{theorem verified exhaustively}\ n\le5]
\]

%dsp-level:0
\[
(3):\ \text{counts}:\ \text{cap-DP over}\ 24\ \text{edges}\ [\text{machine}];
\]

%dsp-level:0
\[
\qquad \text{orbits}:\ \text{Burnside, cap-aware DP per conjugacy class, memoized}\ [\text{machine}];
\]

%dsp-level:0
\[
\qquad 33564>10^4:\ \text{STOP rule}:\ \text{counts by}\ K\ \text{only}\ (K\le12\ \text{automatic}),\ \text{weights for coverings only}
\]

%dsp-level:0
\[
(4):\ (a),(b):\ \mathrm{FT.67}(2),(3),(5)+\text{values of}\ (1)\text{--}(3);
\]

%dsp-level:0
\[
\qquad (c):\ \mathrm{FT.67}(3a)+\mathrm{CT2.8}(2);
\]

%dsp-level:0
\[
\qquad (d):\ \text{conjunction};\ \text{bridge}:\ \mathrm{FT.66}(3)\text{--}(5),\ \text{item}\ 13
\]

%dsp-level:0
\[
(5):\ \neg\exists\ \text{rate tag}\ [\text{inventory}];
\]

%dsp-level:0
\[
\qquad f_{\mathrm{adm}}(n,p)\ \text{expansion}:\ \text{machine-verified coefficient identities}\ \{1,-,\dots\};
\]

%dsp-level:0
\[
\qquad \text{samples exact}
\]

%dsp-level:0
\[
(6):\ \mathrm{FT.67}(3a)\implies\neg\exists\ \text{closed chains in survivors};
\]

%dsp-level:0
\[
\qquad \mathrm{FT.48}(3):\ \text{open minimizers};
\]

%dsp-level:0
\[
\qquad \mathrm{FT.66}(4):\ \text{spectral-layer bound unchanged}
\]

%dsp-level:0
\[
(7):\ \text{finiteness}:\ |E_{\mathrm{dir}}|<\infty\implies2^{|E|}\cdot2^{|E|}\ \text{finite sum};\qquad\text{zeta values as tabulated}\ [\text{machine}].\ \blacksquare
\]

\textbf{Rmk (relation to Def-H.20/FT.62 — no retraction).} The Def-H.20 floors are the fully crystallized points of their skeletons' pattern lattices; FT.62's rung and lock ledgers describe those points and are untouched. This tier adds the amplitude-layer anatomy of the \textit{neighbourhood} those points sit in: the floor itself cancels in the sign sum, its bypass neighbourhood does not, and the two layers (amplitude/spectral) meet only through the open bridge . Nothing here identifies bypass survivors with named particles  or derives the J$\pm$ rates (Cal-F remains declared).

\textbf{Phys.} Let the edges of a locked skeleton flicker — crystallize, dissolve, crystallize — and sum the ledger's own weights over every pattern: the answer is a strict rule with a strict survivor. Only open two-step threads pass, each worth exactly four; every loop, every face, every fully grown solid sums to nothing. The tetrahedron's flicker ensemble weighs $1230$ admissible patterns at probability $615/2048$, and its amplitude budget — $1057$ of it — is spent entirely on thread networks threading sources, sinks, and the occasional locked hub. Polygons obey a two-term recursion with a rational zeta, $1-3z-3z^2$, and their admissible fraction dies exponentially with size: bypass is the physics of small cells. The cube's $1.58$ million admissible patterns collapse to $33{,}564$ symmetry orbits. And the old insight gets its two-layer answer: the flicker is real, and it is a gas of threads — the solids are not what the amplitude sums to; they are what the spectrum declares, on the other side of the one bridge this document has not built.

\vspace{2pt}\noindent\rule{\textwidth}{0.4pt}\vspace{2pt}

\end{theorem}

\phantomsection
\section[Tier F.24 — The Flicker Phase Structure, the Cycle Grammar…]{\texorpdfstring{Tier F.24 — The Flicker Phase Structure, the Cycle Grammar of Covers, the Dust Equation of State, Face Decomposition, and the Bridge Audit}{Tier F.24 - The Flicker Phase Structure, the Cycle Grammar of Covers, the Dust Equation of State, Face Decomposition, and the Bridge Audit}}

\begin{definition}[{\texorpdfstring{Def-H.26 — Flicker-ensemble phase reads}{Def-H.26 - Flicker-ensemble phase reads}}]\label{tag:Def-H-26}
\textbf{Def.} For a skeleton $\Gamma$ [Def-H.25] with $2E$ directed edges and equilibrium parameter $p\in[0,1]$ [$p=x_+/(x_++x_-)$: the unique stationary measure of the per-pair Janus generator — FT.108(4), forward pointer Tier F.36; $r:=x_+/x_-$ a P4-class rate datum (Cal-F); the per-edge independent Bernoulli scope is constitutive of the ensemble (CT4.5a-species ensemble scope)], the \textbf{flicker ensemble} is the Bernoulli product measure $\mu_p$ on crystallization patterns $R\subseteq E_{\mathrm{dir}}(\Gamma)$ (each directed edge real independently with probability $p$). Pattern-layer reads (all exact polynomials/rationals in $p$; $\neg$spectral reads):

%dsp-level:0
\[
B_\Gamma(p):=\mu_p(\text{bypass})=\sum_K c_K\,p^K(1-p)^{2E-K}\quad[c_K:\ \text{cap census, CT0.11}];
\]

%dsp-level:0
\[
\qquad F_\Gamma(p):=p^{2E}\quad[\text{floor measure}];
\]

%dsp-level:0
\[
L_\Gamma(p):=1-B_\Gamma(p)\quad[\text{lock fraction}];\qquad \langle K\rangle_{\mathrm{byp}}(p):=\frac{\sum_K K\,c_K\,p^K(1-p)^{2E-K}}{B_\Gamma(p)};
\]

%dsp-level:0
\[
p^*_\Gamma:=\text{the unique root of }B_\Gamma(p)=F_\Gamma(p)\text{ in }(0,1)\ \ [\text{lockable }\Gamma;\ \text{FT.70(1)}];
\]

%dsp-level:0
\[
\qquad \langle\rho_{D^-}\rangle_\Gamma(p):=\sum_v O(n_v)\,\mu_p(k_v\ge4)\ \ [\mathrm{CT4.1},\mathrm{CT4.5}].
\]

\textbf{Rmk (layer scope).} Every read above is a $\mu_p$-average of pattern-layer ledger entries; the spectral layer (Def-H.16 poles; M1/M2 materialization) is $p$-independent by declaration, and the two layers meet only through the inter-family bridge (application-layer) [FT.68(4)]. Static census data ($c_K$ rows, coverings, crossovers as polynomial facts) are declaration-free; only the equilibrium \textit{interpretation} uses Cal-F (as in FT.68(5)). $\langle\rho_{D^-}\rangle$ inherits CT4.5's per-configuration $w=0$ for its locked part; unlocked configurations carry channel stresses [CT1.22] outside this definition's scope.

\vspace{2pt}\noindent\rule{\textwidth}{0.4pt}\vspace{2pt}

\end{definition}

\begin{theorem}[{\texorpdfstring{FT.70 — The flicker phase diagram: unique crossover, computed boundaries, the polygon limit}{FT.70 - The flicker phase diagram: unique crossover, computed boundaries, the polygon limit}}]\label{tag:FT-70}
\textbf{Thm.} Under Def-H.26 (equilibrium readings under the Cal-F rate datum; the polynomial data declaration-free): $(1)$ \textbf{Existence and uniqueness of the phase boundary.} The admissible family is downward-closed (subsets of admissible patterns are admissible: $k_v$ decreases), nonempty ($\varnothing$), and proper (the floor is locked iff $\Gamma$ is lockable, i.e. $\max_v n(v)\ge4$). Hence $B_\Gamma$ is strictly decreasing on $(0,1)$ [math: monotone coupling; strictness witness: any $3$ edges at a vertex of $n(v)\ge4$ give an admissible $R$ with $R\cup\{e\}$ locked], while $F_\Gamma(p)=p^{2E}$ is strictly increasing; $B(0)-F(0)=1>0$, $B(1)-F(1)=c_{2E}-1=-1<0$. \textbf{Every lockable skeleton has a unique crossover $p^*_\Gamma\in(0,1)$}; non-lockable skeletons (all $n(v)\le3$: the digon, unidirectional cells) have $B\equiv1$: no interior boundary, thread gas at every $p<1$. $(2)$ \textbf{Computed boundaries} (exact dyadic isolating intervals of width $2^{-40}$; machine; decimal display $\approx$):

%dsp-level:0
\[
p^*:\quad\text{icosa}\ 0.652034\ <\ \text{octa}\ 0.679217\ <\ \text{tetra}\ 0.723525\ <\ \text{dodeca}\ 0.724519\ <\ \text{cube}\ 0.725372\ <\ P_3\ 0.807901\ <\ P_5\ <\ P_7\ <\ P_9\ <\ \sqrt{\tfrac23}\ <\ P_{10}\ <\ P_8\ <\ P_6\ <\ P_4\ 0.819671 .
\]

Regions: $p<p^*_\Gamma$ thread-gas-dominant ($B>F$), $p>p^*_\Gamma$ floor-dominant ($F>B$); the locked fraction $L_\Gamma(p)=1-B_\Gamma(p)$ rises $0\to1$ across the diagram. The two degree-$\ge4$ solids cross earliest — the crossover order reproduces the lock-proneness order of clause (4). $(3)$ \textbf{Polygon limit, closed form.} $\operatorname{charpoly}T(t)=x^2\bigl(x^2-(1+2t)x-t^2(1+2t)\bigr)$ [machine: exact at $t\in\{1,2,\tfrac13,\tfrac52\}$; generalizes FT.68(2)'s $t{=}1$ value]. Asymptotically $B_{P_n}(p)\sim(1-p)^{2n}\lambda_{\max}(t)^n$, $t=p/(1-p)$: the crossing equation tends to $\lambda_{\max}(s)=s^2$, i.e. $s^4=2s^2(1+2s)$, i.e. $s^2-4s-2=0$ with $s=p/(1-p)$: $s=2+\sqrt6$ and

%dsp-level:0
\[
p^*_\infty=\frac{s}{1+s}=\frac{2+\sqrt6}{3+\sqrt6}=\frac{\sqrt6}{3}=\sqrt{\tfrac23}
\]

%dsp-level:0
\[
\qquad \Bigl[\ \bigl(\sqrt\tfrac23\bigr)^2=\tfrac23;\ (2{+}\sqrt6)^2=10{+}4\sqrt6=4(2{+}\sqrt6){+}2\ \checkmark\ \text{in}\ \mathbb Z[\sqrt6]\Bigr],
\]

approached from below by odd $n$ and from above by even $n$ (interlaced; machine through $n=10$). $(4)$ \textbf{Lock-proneness order (the corrected "forbidden zone").} Exact identity [machine, symbolic $n=3..10$]:

%dsp-level:0
\[
f_{\mathrm{adm}}(n{+}1,p)=f_{\mathrm{adm}}(n,p)-\binom{n}{3}p^4(1-p)^{n-3}\ <\ f_{\mathrm{adm}}(n,p)\qquad\forall p\in(0,1),
\]

so admissibility is strictly decreasing in potential degree: $f_{\mathrm{adm}}(2d,\tfrac12)=(1,\tfrac{15}{16},\tfrac{21}{32},\tfrac{93}{256},\tfrac{11}{64})$ for $2d=(2,4,6,8,10)$: digon $>$ polygons $>$ \{tetra, cube, dodeca\} $>$ octa $>$ icosa. \textbf{The proposed octahedron "forbidden zone" is refuted}: $f_{\mathrm{adm}}(8,p)\ge(1-p)^8>0$ for every $p<1$ and $318\,523$ admissible configurations exist at the cap level [(1)]. What is true: at fixed $p<1$ the octahedron/icosahedron are the most lock-prone Platonic skeletons, and $p=1$ locks \textit{every} skeleton — the static floor is the $p\to1$ limit of its own ensemble [FT.68(5)], with amplitude zero there [FT.67(3a)]. $(5)$ \textbf{Effective edge content.} $\langle K\rangle_{\mathrm{byp}}(p)$ is strictly increasing on the exact grid $p\in\{\tfrac1{20},\dots,\tfrac{19}{20}\}$ (machine: tetra, cube, icosa, $P_5$), with limits $0$ (vacuum-dominated, $p\to0$) and $K_{\max}=\lfloor3V/2\rfloor$ (top bypass row, $p\to1$); at $p=\tfrac12$: $\bigl(1,\tfrac{852}{205},\tfrac{3311952}{395023},\tfrac{2124600}{318523},\tfrac{4007491093935690}{191727103387307},\tfrac{215801166207}{15519497071}\bigr)$ for digon/tetra/cube/octa/dodeca/icosa (the tetra entry is FT.68(1)'s $m_{\mathrm{eff}}$ reading (i) verbatim — $p{=}\tfrac12$ is the uniform measure) and $\bigl(\tfrac83,\tfrac{248}{69},\tfrac{130}{29},\tfrac{296}{55},\tfrac{2618}{417},\tfrac{11344}{1581},\tfrac{896}{111},\tfrac{13588}{1515}\bigr)$ for $P_3..P_{10}$. The dark mass read is $m_{\mathrm{eff}}(p)=\langle K\rangle_{\mathrm{byp}}(p)\ln2/\mathfrak h$ [FT.68(1) pattern]: light at high flicker, heavy toward the static limit, exact at every $p$.

\textbf{Req.} Def-H.25, Def-H.26; FT.42, FT.45, FT.62, FT.63, FT.67, FT.68; A4, D0.1; CT0.11 (Part 1); Cal-F (equilibrium readings); monotone coupling for product measures, transfer-matrix asymptotics, dyadic bisection [math].

\textbf{Pf.}

%dsp-level:0
\[
(1):\ R'\subseteq R:\ k_v(R')\le k_v(R):\ \text{downward-closed};
\]

%dsp-level:0
\[
\qquad \text{coupling}\ p_1<p_2:\ R_1\subseteq R_2:\ \mathbf 1_{\mathrm{adm}}(R_1)\ge\mathbf 1_{\mathrm{adm}}(R_2):\ B(p_1)\ge B(p_2)
\]

%dsp-level:0
\[
\text{strictness witness}:\ |R|=3\ \text{at}\ v,\ n(v)\ge4:\ R\ \text{admissible},\ R\cup\{e_v\}:\ k_v=4:\ \text{locked};
\]

%dsp-level:0
\[
\qquad B(0)=c_0=1,\ B(1)=c_{2E}=0\ [\text{floor locked}]
\]

%dsp-level:0
\[
(2):\ \text{exact integer bisection on}\ G(p)
\]

%dsp-level:0
\[
= B(p)-p^{2E}:\ G(m/2^{40})\ \text{as an integer via}\ 2^{40\cdot 2E}\ \text{clearing};
\]

%dsp-level:0
\[
\ \text{single sign flip on a 51-point grid}\ [\text{machine}]
\]

%dsp-level:0
\[
(3):\ T(t)\ \text{rows}\ (00,01,10)\ \text{equal on}\ \{00,01,10,11\}\text{-pattern}\implies\operatorname{charpoly}
\]

%dsp-level:0
\[
= x^2\chi_2,
\]

%dsp-level:0
\[
\ \chi_2
\]

%dsp-level:0
\[
= x^2-(1{+}2t)x-t^2(1{+}2t)\ [\text{machine at 4 exact}\ t];
\]

%dsp-level:0
\[
\quad\lambda_{\max}(s)
\]

%dsp-level:0
\[
= s^2
\]

%dsp-level:0
\[
\implies s^2\qquad = 2(1{+}2s)
\]

%dsp-level:0
\[
s=2+\sqrt6>0;\qquad p^*_\infty=\tfrac s{1+s}=\tfrac{2+\sqrt6}{3+\sqrt6}\cdot\tfrac{3-\sqrt6}{3-\sqrt6}=\tfrac{\sqrt6}3;\qquad\bigl(\tfrac{\sqrt6}3\bigr)^2=\tfrac23
\]

%dsp-level:0
\[
(4):\ \text{binomial recursion}\ \mathrm{Bin}(n{+}1,p)
\]

%dsp-level:0
\[
= \mathrm{Bin}(n,p)+\mathrm{Ber}(p):\ f(n{+}1,p)
\]

%dsp-level:0
\[
= (1-p)f(n,p)+p\,
\]

%dsp-level:0
\[
P(\mathrm{Bin}(n,p)\le2)
\]

%dsp-level:0
\[
= f(n,p)-p\,
\]

%dsp-level:0
\[
P(\mathrm{Bin}(n,p)=3)
\]

%dsp-level:0
\[
p\,
\]

%dsp-level:0
\[
P(\mathrm{Bin}(n,p){=}3)
\]

%dsp-level:0
\[
= \binom n3p^4(1-p)^{n-3}>0\ \ \forall p
\]

%dsp-level:0
\[
\in (0,1),
\]

%dsp-level:0
\[
\ n
\]

%dsp-level:0
\[
\ge 3\ [\text{machine: symbolic coefficient identity}]
\]

%dsp-level:0
\[
(5):\ \text{exact rational evaluation from}\ c_K\ \text{rows};
\]

%dsp-level:0
\[
\ \text{grid monotonicity machine};
\]

%dsp-level:0
\[
\ p\to0:\ c_0\ \text{term dominates}:\ \langle K\rangle\to0;
\]

%dsp-level:0
\[
\quad p\to1:\ c_{K_{\max}}\ \text{term dominates}.\ \blacksquare
\]

\textbf{Phys.} Let the crystallization rate run from never to always, and every skeleton shows the same two-phase portrait with one sharp boundary: below it the ledger's amplitude lives in a gas of open threads, above it in the static floor — and the boundary sits at an exactly computable flicker rate, about $0.72$ for the degree-three solids, earliest ($0.65$) for the icosahedron, latest for small polygons, converging for the polygon family to the closed value $\sqrt{2/3}$. No solid is forbidden from flickering: the octahedron's rumored "must lock" is false at every rate below one — it is only the \textit{most} lock-prone per vertex, and at rate one everything locks, floors included, whose amplitude there is zero anyway. The effective dark mass of the gas climbs smoothly from nothing to the floor's own ledger as the flicker slows.

\vspace{2pt}\noindent\rule{\textwidth}{0.4pt}\vspace{2pt}

\end{theorem}

\begin{theorem}[{\texorpdfstring{FT.71 — The flicker-dust equation of state}{FT.71 - The flicker-dust equation of state}}]\label{tag:FT-71}
\textbf{Thm.} Under Def-H.26 (equilibrium readings under the Cal-F rate datum), for every skeleton $\Gamma$: $(1)$ \textbf{Exact dust density.} Per-edge independence gives $k_v\sim\mathrm{Bin}(n(v),p)$, hence

%dsp-level:0
\[
\langle\rho_{D^-}\rangle_\Gamma(p)=\sum_v O(n_v)\,\bigl(1-f_{\mathrm{adm}}(n_v,p)\bigr),\qquad O(n)=2^{n+1}-16\ \ [\mathrm{CT4.1};\ \mathrm{FT.62}(2)],
\]

an exact polynomial in $p$ with integer coefficients: Platonic values at $p=\tfrac12$: $\bigl(154,\ 308,\ \tfrac{15159}{8},\ 770,\ 20193\bigr)$ for (tetra, cube, octa, dodeca, icosa); polygon family the exact monomial $\langle\rho_{D^-}\rangle_{P_n}(p)=16\,n\,p^4$ (at $p{=}\tfrac12$: $n$); digon identically $0$ (never locks). $(2)$ \textbf{Anchors.} $\langle\rho_{D^-}\rangle_\Gamma(1)=\sum_vO(n_v)=$ FT.63(5)'s dark sums $(448, 896, 2976, 2240, 24384)$ exactly [machine]; $\langle\rho_{D^-}\rangle_\Gamma(0)=0$; strict increase on $(0,1)$ for lockable $\Gamma$ [clause (4)'s $f_{\mathrm{adm}}$ strict decrease, FT.70(4)]. $(3)$ \textbf{The flicker modulates $\rho$, never $w$.} CT4.5 applies per configuration: every locked vertex contributes pressureless dust ($w=0$, source side, coupled to geometry through CT1.23 and nothing else). The ensemble's locked content is therefore dust at every $p$; the flicker parameter moves only the \textit{density} — along the exact polynomial of (1). $D^+$'s $w=-1$ [CT1.25, D4.3] is configuration-independent and untouched: the CT4.6 contrast table holds ensemble-wise term by term. $(4)$ \textbf{No fluid.} The locked sector supports no further equation-of-state data: it is a structureless absorbing sink — one possibility space [D4.1], no gauge content [CT4.4], absorption one-way [CT4.2], content inaccessible [D1.11] and unresolvable [FT.3]. Statements assigning $D^-$ a pressure response, polarization, or re-emission have no defined ledger referent (the referent would contradict CT4.2/D1.11); "the vacuum is a $D^-$ fluid" is inadmissible as stated — what the corpus derives is: locked content gravitates as dust with the density computed here.

\textbf{Req.} Def-H.26; FT.62, FT.63, FT.70; CT1.22, CT1.23, CT1.25 (Part 2); CT4.1, CT4.2, CT4.3, CT4.4, CT4.5, CT4.6 (Part 2); D1.11, D4.1, D4.2, D4.3; FT.3; binomial law under per-edge independence [math]; Cal-F (equilibrium readings).

\textbf{Pf.}

%dsp-level:0
\[
(1):\ \mu_p(k_v\ge4)=P(\mathrm{Bin}(n_v,p)\ge4)=1-f_{\mathrm{adm}}(n_v,p)\ [\text{per-edge independence, Def-H.26 scope}];
\]

%dsp-level:0
\[
\qquad \text{linearity of expectation over}\ v\ [\mathrm{CT4.5}:\ \rho_{D^-}=\textstyle\sum O(n)]
\]

%dsp-level:0
\[
O(2d)=2^{2d+1}-16:\ (16,112,496,2032)\ \text{at}\ 2d=(4,6,8,10)\ [\mathrm{FT.62}(2)\ \text{table}];
\]

%dsp-level:0
\[
\qquad P_n:\ \text{all}\ n(v)=4:\ \textstyle\sum_v16\,p^4=16np^4
\]

%dsp-level:0
\[
(2):\ p{=}1:\ 1-f_{\mathrm{adm}}=1:\ \textstyle\sum_vO(n_v)=(4\cdot112,\ 8\cdot112,\ 6\cdot496,\ 20\cdot112,\ 12\cdot2032)\ \checkmark;
\]

%dsp-level:0
\[
\qquad p{=}0:\ 0;\quad\text{monotone by FT.70(4)}
\]

%dsp-level:0
\[
(3):\ \text{CT4.5 per configuration};\ \text{mixture of}\ w{=}0\ \text{components is}\ w{=}0\ \text{on the locked part};
\]

%dsp-level:0
\[
\qquad (4):\ \text{cited tags exhaust the}\ D^-\ \text{inventory}\ [\text{Parts 1–2}];\ \text{no response read exists}.\ \blacksquare
\]

\textbf{Phys.} Shake a crystal at rate $r$ and its dark content has an exact expected density: zero when nothing crystallizes, the full retained-register sum $448\ldots24384$ per tick when everything does, and a smooth polynomial in between — $16np^4$ for a polygon ring, $\tfrac{15159}{8}$ for the octahedron at equal rates. Whatever the rate, that content gravitates the same way: pressureless dust on the source side of the field equations. The flicker changes how much sediment the ledger carries, never how the sediment weighs.

\vspace{2pt}\noindent\rule{\textwidth}{0.4pt}\vspace{2pt}

\end{theorem}

\begin{theorem}[{\texorpdfstring{FT.72 — Two-fold face decomposition; the Bernoulli marginals ("coarse-grained fuzz as data")}{FT.72 - Two-fold face decomposition; the Bernoulli marginals ("coarse-grained fuzz as data")}}]\label{tag:FT-72}
\textbf{Thm.} Let $\Sigma_S$ be a Platonic skeleton (sphere), $E,F$ its undirected edge/face counts, all arithmetic over $\mathrm{GF}(2)$. $(1)$ \textbf{Two-fold face-decomposition theorem.} The face-boundary map $\partial_2:\mathrm{GF}(2)^F\to\mathrm{GF}(2)^E$ has image exactly the cycle space $Z_1=\ker\partial_1$ (rank $\beta_1=F-1$ [FT.43(5), FT.64; machine: all five]) and kernel $\{0,\sum_f f\}$ (every edge lies in exactly two faces). Hence \textbf{every cycle $z\in Z_1$ has exactly two face-decompositions, $z=\partial_2 S=\partial_2(F\setminus S)$} — a set of faces and its complement. "Every subnetwork is a composition of triangles" holds exactly for the triangulated skeletons (tetrahedron $4$, octahedron $8$, icosahedron $20$ triangles [machine census]); the cube's faces are squares, and \textbf{the dodecahedron — girth $5$: zero triangles and zero $4$-cycles [machine] — admits no triangle decomposition within its skeleton at all}: its face grammar is pentagonal, and triangle compositions would require diagonals outside $\mathcal R$. $(2)$ \textbf{Hodge-type split (corrected form).} $Z_1=\mathrm{Cut}^\perp$ under the standard pairing; $\dim Z_1+\dim\mathrm{Cut}=\beta_1+(V{-}1)=E$ [Euler], but the sum is \textbf{not} direct in general: $j:=\dim(Z_1\cap\mathrm{Cut})=(2,3,3,5,5)$ for (tetra, cube, octa, dodeca, icosa) [machine]. A pattern $R$ is decomposable ($R\in Z_1+\mathrm{Cut}$: cycle content $+$ boundary content) iff its syndrome is orthogonal to $Z_1\cap\mathrm{Cut}$; the decomposition is unique up to $2^j$ gauge choices. $(3)$ \textbf{Exact Bernoulli marginals.} For $R\sim\mu_p$ [Def-H.26], with $q:=1-2p$:

%dsp-level:0
\[
P_p(R\in Z_1)=\frac{1}{2^{V-1}}\sum_{c\in\mathrm{Cut}}q^{|c|},
\]

%dsp-level:0
\[
\qquad P_p(R\ \text{decomposable})=\frac{1}{2^{j}}\sum_{u\in Z_1\cap\mathrm{Cut}}q^{|u|},
\]

%dsp-level:0
\[
\qquad P_p(R\in\mathrm{Cut})=\frac{1}{2^{\beta_1}}\sum_{z\in Z_1}q^{|z|}
\]

[MacWilliams for $Z_1=\mathrm{Cut}^\perp$ and $(Z_1+\mathrm{Cut})^\perp=Z_1\cap\mathrm{Cut}$; machine: direct enumeration cross-check at $\beta_1\le5$]. At $p=\tfrac12$: $P(R\in Z_1)=2^{-(V-1)}=(\tfrac18,\tfrac1{128},\tfrac1{32},\tfrac1{524288},\tfrac1{2048})$ and $P(\text{decomposable})=2^{-j}=(\tfrac14,\tfrac18,\tfrac18,\tfrac1{32},\tfrac1{32})$; samples $p=\tfrac14$: $(\tfrac{27}{128},\tfrac{19}{64})$, $p=\tfrac34$: $(\tfrac{11}{128},\tfrac{19}{64})$ for the tetrahedron [machine table for all five]. $(4)$ \textbf{Dual signatures at the Hodge level.} The intersection weight-enumerators of dual pairs coincide exactly: cube $\equiv$ octahedron $\{0{:}1,6{:}4,8{:}3\}$, dodeca $\equiv$ icosa $\{0{:}1,10{:}6,16{:}15,18{:}10\}$ [machine] — a new degeneracy of the FT.44 species — while the cut enumerators split them: the octahedron's cut space is even-weight (all degrees $4$: every star even; the subgraph-level echo of FT.63(3)'s black–white parity law) while the cube's carries odd stars (weight $3$). $(5)$ \textbf{"Fuzz as data", formalized.} The complete flicker datum of a skeleton is the measure $\mu_p$ on $\mathrm{GF}(2)^E$; every coarse-grained read computed in this document is a named marginal of it: $K\mapsto I_{\mathrm{hid}}$-type edge content [FT.68(1)], admissibility $\mapsto$ bypass/lock sectors [Def-H.25], cycle/cut class $\mapsto$ (3) above, cover support $\mapsto$ amplitude [FT.67(2)], winding lattices $\mapsto$ FT.64's index-$2^{F-1}$ sublattices (on-shell $\kappa$-side, not $\mu_p$-side — layer marked). No read is lost by coarse-graining that the measure retains; nothing in the measure is readable that the ledger's registers admit [C10 reality principle; FT.3 for locked content].

\textbf{Req.} FT.43, FT.44, FT.63, FT.64, FT.67; Def-H.25, Def-H.26; D0.1; CT1.29 (Part 2); GF(2) linear algebra, MacWilliams identity, Euler characteristic [math]; Cal-F (only for clause (5)'s measure).

\textbf{Pf.}

%dsp-level:0
\[
(1):\ \text{each edge in exactly 2 faces (sphere)}\implies\partial_2\bigl(\textstyle\sum_f f\bigr)=0;
\]

%dsp-level:0
\[
\qquad \operatorname{rk}\partial_2\le\beta_1\ [\text{faces are cycles}];\ \operatorname{rk}\partial_2\ge F-1\ \text{and}\ \dim\ker\partial_2=1\ [\text{machine: elimination, all five}]
\]

%dsp-level:0
\[
\implies\operatorname{im}\partial_2\subseteq Z_1,\ \dim\operatorname{im}=F-1=\beta_1=\dim Z_1\ [\mathrm{FT.43}(5):\ \beta_1=E-V+1=F-1]\implies\operatorname{im}\partial_2=Z_1;\qquad\#\partial_2^{-1}(z)=|\ker\partial_2|=2
\]

%dsp-level:0
\[
\text{girth census}:\ \text{dodeca triangles}=0,\ 4\text{-cycles}=0\ [\text{machine}];\qquad\text{triangle counts}\ (4,0,8,0,20)\ \checkmark
\]

%dsp-level:0
\[
(2):\ Z_1=\ker\partial_1=(\operatorname{im}\partial_1^T)^\perp=\mathrm{Cut}^\perp\ [\text{linear algebra}];\qquad j=\dim Z_1+\dim\mathrm{Cut}-\dim(Z_1+\mathrm{Cut})\ [\text{machine ranks}]
\]

%dsp-level:0
\[
(3):\ \text{MacWilliams}\ W_{C^\perp}(x,y)=\tfrac1{|C|}W_C(y-x,y+x);\quad x=p,y=1-p:\ y-x=q,\ y+x=1;
\]

%dsp-level:0
\[
\qquad (Z_1+\mathrm{Cut})^\perp=Z_1^\perp\cap\mathrm{Cut}^\perp=\mathrm{Cut}\cap Z_1
\]

%dsp-level:0
\[
p=\tfrac12:\ q=0:\ \text{only the empty word survives}:\ 2^{-(V-1)},\ 2^{-j};
\]

%dsp-level:0
\[
\qquad \text{direct enumeration cross-check: tetra/cube}\ [\text{machine}]
\]

%dsp-level:0
\[
(4):\ \text{enumerator tables machine};
\]

%dsp-level:0
\[
\qquad \text{octa stars}:\ \deg
\]

%dsp-level:0
\[
= 4
\]

%dsp-level:0
\[
\implies |\delta(v)|
\]

%dsp-level:0
\[
= 4\ \text{even};
\]

%dsp-level:0
\[
\ \text{even cuts closed under}\ \triangle\implies\text{even-weight cut space};
\]

%dsp-level:0
\[
\ \text{cube stars}\ |3|\ \text{odd}
\]

%dsp-level:0
\[
(5):\ \text{marginal list as cited};\ \text{layer marks at each site}.\ \blacksquare
\]

\textbf{Rmk (relation to the proposal transcript).} The proposal "any network subgraph is a composition of triangles with free real/virtual choice per edge" is \textit{corrected, not adopted}: the free binary choice per directed edge is exactly A4/D0.1/Def-H.25 (frozen and declared respectively); the composition theorem holds with each skeleton's own face type and a two-fold ambiguity per cycle; the triangle-only form fails on the cube (squares) and maximally on the dodecahedron (girth $5$). The proposal's ensemble sense — "the coarse-grained fuzz is the data about the system" — is clause (5): the measure is the datum, the ledger reads are its marginals.

\textbf{Phys.} Every loop you can draw on a solid is a sum of faces in exactly two ways — a set of faces, or everything else. That is the whole content of "decomposing into triangles" where triangles exist, and on the dodecahedron no triangles exist to decompose into. Let the edges flicker independently and the fuzz becomes an exact probability measure: the chance that a random pattern is a pure cycle is one in $2^{V-1}$ at equal rates — one in half a million on the dodecahedron — and every ledger read this document has ever computed is a shadow of that single measure.

\vspace{2pt}\noindent\rule{\textwidth}{0.4pt}\vspace{2pt}

\end{theorem}

\medskip\noindent\textit{\texorpdfstring{FT.73 — The bridge audit: dielectric non-existence, the field-degree no-go, and the surviving species}{FT.73 - The bridge audit: dielectric non-existence, the field-degree no-go, and the surviving species}.}\ Structural-constraints record: see Appendix B, Section B.2.2 of the Addendum.\par\medskip

\phantomsection
\section[Tier F.25 — The Chain-Local -Domains: Extended Exact-Cover…]{\texorpdfstring{Tier F.25 — The Chain-Local $\aleph$-Domains: Extended Exact-Cover Theorems and the Domain-Relativity of the Amplitude Layer}{Tier F.25 - The Chain-Local -Domains: Extended Exact-Cover Theorems and the Domain-Relativity of the Amplitude Layer}}

\begin{definition}[{\texorpdfstring{Def-H.27 — The $\aleph$-domains: literal, chain-dominant, literal-forced}{Def-H.27 - The -domains: literal, chain-dominant, literal-forced}}]\label{tag:Def-H-27}
\textbf{Def.} Let $R\subseteq E_{\mathrm{dir}}(\Gamma)$ be a crystallization pattern on a skeleton $\Gamma$ [Def-H.25] under scopes (S1), (S3), (S4). Three domains for the vertex-amplitude exponent $\aleph$ entering A8's weight $w=\mathcal J\prod_{e\in R}G(e)(-1)^{\aleph}$, $G(e)=\sigma_e e^{i\pi\kappa_e}$ [D3.5, FT.9]: \textbf{(a) Literal (S2) — frozen.} $\aleph_v=|f_v|$ [D3.3] at type-$(1,1)$ vertices of $R$, $f_v=\tfrac12(\sigma_{\rightarrow}-\sigma_{\leftarrow})$; $\aleph_v=0$ at every other type. No chain data. [This is the domain of FT.67–FT.73 and of the frozen corpus.] \textbf{(b) E1 — chain-dominant.} Configurations $(R,\sigma,D)$ with $D$ a \textbf{chain decomposition}: a partition of $R$ into directed walks (edge sequences with $t(e_i)=s(e_{i+1})$ [D1.4]). For each walk $c$: open ($s(c)\ne t(c)$) $\Rightarrow$ interior vertices $=$ all but the two endpoints; closed ($s(c)=t(c)$; D1.4's Janus clause is trivial at (S3)) $\Rightarrow$ \textbf{every} vertex of $c$ is interior, the start vertex included — arrival on the last edge, departure on the first (the \textbf{wrap read}). Then

%dsp-level:0
\[
\aleph(\sigma,R,D)=\sum_{c\in D}\ \sum_{v\in\mathrm{int}(c)}\bigl|f_v(c)\bigr|,
\]

%dsp-level:0
\[
\qquad f_v(c)=\tfrac12\bigl(\sigma_{\mathrm{out}(c,v)}-\sigma_{\mathrm{in}(c,v)}\bigr)\ \text{in }c\text{'s own edges at }v,
\]

and $Z^{E1}_\Gamma=\sum_R\sum_D\sum_\sigma w^{E1}$. \textbf{(c) E2 — literal-forced.} At type-$(1,1)$ vertices the literal D3.3 read is retained \textit{regardless of $D$} — the chain choice there is $\aleph$-gauge: such a vertex has a unique in- and a unique out-edge in $R$, so every passage through it gives the same $f_v$ and the datum is type-determined. At higher-type vertices the reads come from chain passages. Equivalently — and this is the form computed — the summed datum is the \textbf{partial-matching system} $M=(M_v)_{v\ \mathrm{higher}}$, each $M_v$ a partial matching between $\mathrm{in}_R(v)$ and $\mathrm{out}_R(v)$; the orbits of the induced partial bijection on $R$ are the chains, and an orbit that is a \textit{path returning to its start vertex} is a closed chain in the sense of (b) and gains the wrap read there:

%dsp-level:0
\[
\aleph(\sigma,R,M)=\sum_{v\ (1,1)}|f_v|+\sum_{v\ \mathrm{higher}}\ \sum_{(e,e')\in M_v}\bigl|f(e,e')\bigr|+\sum_{\mathrm{wraps}}|f|,
\]

%dsp-level:0
\[
\qquad Z^{E2}_\Gamma=\sum_R\sum_M\sum_\sigma w^{E2}.
\]

\textbf{Coverage convention (all three domains).} For $e\in R$ let $\operatorname{cov}_{\mathrm{dom}}(e):=\#\{\aleph\text{-reads of }e\text{ in the domain}\}$ (literal: the type-$(1,1)$ endpoints; E1: the walk-interior and wrap reads; E2: the forced literal reads, the $M$-passages, and the wraps). Since $(-1)^{\aleph_v}=\sigma_{\mathrm{in}}\sigma_{\mathrm{out}}$ [FT.67(1)] and $\prod_{e\in R}G(e)=(-1)^{\sum_e\kappa_e}\prod_{e\in R}\sigma_e$ at (S1),

%dsp-level:0
\[
w^{\mathrm{dom}}=(-1)^{\sum_e\kappa_e}\prod_{e\in R}\sigma_e^{\,1+\operatorname{cov}_{\mathrm{dom}}(e)}
=(-1)^{\sum_e\kappa_e}\!\!\prod_{e:\ \operatorname{cov}_{\mathrm{dom}}(e)\ \mathrm{even}}\!\!\sigma_e,
\]

%dsp-level:0
\[
\qquad \sum_\sigma w^{\mathrm{dom}}\ne0\iff\operatorname{cov}_{\mathrm{dom}}(e)\ \text{odd}\ \ \forall e\in R,
\]

the $\sigma_e$-exponent being $1+\operatorname{cov}_{\mathrm{dom}}(e)$ in every domain [FT.67(2)'s Pf, domain-blind]. **[typographical correction, applied in place at FT.67(2): its displayed monomial subscript read $\operatorname{cov}(e)$ \textit{odd}; the exponent identity $1+\operatorname{cov}(e)$ recorded in its own Pf forces \textit{even}. The classification, the values $4^{|\mathcal C|}(-1)^{\sum\kappa}$, the Pf, and every census built on them are unchanged — the slip is display-only.]**

\textbf{Rmk (status).} Domains (b), (c) are \textbf{computed variants, not adopted}: Parts 1–2 define $\aleph$ by (a) [D3.3 via A8]; domains (b), (c) await the domain choice. The extension is \textit{forced data} wherever the corpus itself uses chain-local pass-through reads — FT.21(i)'s telescoping identity, FT.62(5)'s observed-variant EM rows — so (b)/(c) are the two canonical completions of that usage to the weight; which (if either) governs A8's $\aleph$ is an author's declaration, not derivable here, and no interface definition is introduced (nothing below is used as a premise by any other tier; the literal-domain record of Tiers F.23–F.24 stands untouched). Both variants are evaluated under the same (S1)/(S3)/(S4) scopes; $\kappa$-invariance of the classifications and the $(-1)^{\sum\kappa}$ value scaling carry over verbatim [FT.67(2) Pf pattern].

\textbf{Rmk (finiteness).} For finite $R$ the sets of walk partitions $D$ and of partial-matching systems $M$ are finite, so $Z^{E1}_\Gamma$, $Z^{E2}_\Gamma$ are finite sums over a finite skeleton — convergence trivial, as in Def-H.25.

\vspace{2pt}\noindent\rule{\textwidth}{0.4pt}\vspace{2pt}

\end{definition}

\begin{theorem}[{\texorpdfstring{FT.74 — The extended exact-cover theorems}{FT.74 - The extended exact-cover theorems}}]\label{tag:FT-74}
\textbf{Thm.} Under Def-H.27, at (S1); for general on-shell integer $\kappa$ every value below scales by $(-1)^{\sum_e\kappa_e}$ and no classification changes. Write $\varepsilon(R):=(-1)^{|R|}2^{|R|}$: on the support of either classifier $|R|$ is even [each thread contributes $2$ edges], so $\varepsilon(R)=2^{|R|}>0$ and the stated values are the exact amplitudes. $(1)$ \textbf{E1 classification.}

%dsp-level:0
\[
\sum_D\ \sum_\sigma w^{E1}(R,\sigma,D)=N_{E1}(R)\cdot\varepsilon(R),
\]

where $N_{E1}(R)$ is the number of partitions of $R$ into \textbf{open directed $2$-threads} $s\to v\to t$ ($s\ne t$; the junction $v$ of any type — locked junctions included). Every other decomposition cancels: a walk of length $\ge3$ double-reads its interior edges; a length-$1$ walk leaves its edge unread ($\operatorname{cov}=0$); a closed walk double-reads every edge (all vertices interior, wrap included); a closed $2$-thread — a digon — is a closed chain [D1.4] and cancels by the same wrap read. $(2)$ \textbf{E2 classification.}

%dsp-level:0
\[
\sum_M\ \sum_\sigma w^{E2}(R,\sigma,M)=N_{E2}(R)\cdot\varepsilon(R),
\]

with $N_{E2}(R)=0$ unless \textbf{(a)} no edge of $R$ joins two type-$(1,1)$ vertices (such an edge is read literally twice), and \textbf{(b)} every literal thread at a type-$(1,1)$ vertex is open — a literal thread whose two ends coincide is a closed chain [D1.4], and its wrap read double-reads both edges. Otherwise $N_{E2}(R)$ is the number of open-$2$-thread partitions of the \textbf{higher-only remainder} $R_{\mathrm{hh}}:=\{e\in R:\ \text{neither endpoint is type-}(1,1)\}$, junctions necessarily higher-type; the forced literal threads cover exactly the edges touching type-$(1,1)$ vertices, once each. $(3)$ \textbf{Pointwise order, and incomparability of the three domains.} $N_{E2}(R)\le N_{E1}(R)$ for every $R$ on every skeleton (an E2 structure maps injectively to its thread partition: forced literal threads are open by (2b), edges with a $(1,1)$ endpoint are covered exactly once by (2a)-(2b), and the $R_{\mathrm{hh}}$ threads are E1 threads) — machine-asserted with zero violations on $9472$ patterns (the tetrahedron and bidirectional $P_4/P_5/P_6$, exhaustively), together with $|R|$ even on both supports. The literal domain embeds into neither extension and neither into it. Witnesses (triples $(\mathrm{lit},E2,E1)$; machine: classifier $=$ brute-force sum on each row):

%dsp-level:0
\[
\begin{array}{l|cccccc}
R & \text{digon} & \text{dir. }3\text{-cycle} & \text{dir. }4\text{-cycle}^{*} & \text{double-thread hub} & \text{digon-pair hub} & \text{bidir }P_4\text{ floor}\\
\hline
(\mathrm{lit},E2,E1) & (0,0,0) & (0,0,0) & (0,0,32) & (0,32,32) & (16,0,16) & (0,1024,1024)
\end{array}
\]

$^{*}$on the $4$-cycle skeleton. The \textbf{digon-pair hub} $R=\{(0{\to}1),(1{\to}0),(0{\to}2),(2{\to}0)\}$, types $\bigl((2,2),(1,1),(1,1)\bigr)$ — the FT.69(4)/FT.69(5) witness family — has literal value $16$, is killed by E2 (both literal threads are closed) and revived by E1 at the alternative threading $\{(2{\to}0),(0{\to}1)\},\{(1{\to}0),(0{\to}2)\}$: $N_{\mathrm{lit}}\not\le N_{E2}$. The \textbf{directed $4$-cycle} gives $N_{E2}=0<N_{E1}=2$ (strictness of (3)'s order). The \textbf{double-thread hub} $\{(0{\to}2),(1{\to}2),(2{\to}3),(2{\to}4)\}$ has literal value $0$ and extended value $32$: both extensions exceed the literal domain. \textbf{The digon — the $\Sigma_{\min}$ pattern [FT.28] — cancels in all three domains}: the minimal floor's absence from the amplitude layer is domain-invariant. $(4)$ \textbf{Equality locus.} On patterns with no type-$(1,1)$ vertex, E2 $\equiv$ E1 exactly (no literal reads exist; in particular on every fully crystallized floor, all of whose vertices are type $(d,d)$, $d\ge2$ [FT.62]). There the extended value is $N\cdot\varepsilon(R)$ with $N$ the open-thread partition count, and the literal value is $0$ whenever $R\ne\varnothing$ [FT.67(3a)]. $(5)$ \textbf{Machine verification.} The classifications are verified against brute-force sign sums taken straight from Def-H.27's semantics — E1 by explicit enumeration of \textit{all} walk partitions $D$ with interior/wrap reads, E2 by explicit enumeration of \textit{all} partial-matching systems $M$ with orbit-wrap closure — on \textbf{every} pattern of seven families: the digon ($2^2$), the bidirectional triangle ($2^6$), bidirectional $P_4$ ($2^8$), the single-vertex stars $n\in\{4,5,6\}$ ($2^4,2^5,2^6$), and all tetrahedron patterns with $K\le4$ ($794$): \textbf{$2460$ pattern-domain checks, zero mismatches}.

\textbf{Req.} Def-H.25, Def-H.27; FT.21, FT.22, FT.28, FT.62, FT.67, FT.68, FT.69; A4, A8; D0.1, D0.8, D1.4, D1.5, D3.3, D3.5; CT0.9 (Part 1); involution/sign-sum arguments, walk and partial-matching enumeration [math; machine].

\textbf{Pf.}

%dsp-level:0
\[
(1):\ \text{walk}\ (e_1,\dots,e_k),
\]

%dsp-level:0
\[
\ k
\]

%dsp-level:0
\[
\ge 3:\ e_j\ (1<j<k)\ \text{read at both its ends}:\ 1+\operatorname{cov}
\]

%dsp-level:0
\[
= 3:\ \sigma_{e_j}
\]

%dsp-level:0
\[
\mapsto -\sigma_{e_j}\ \text{involution}\qquad \implies 0
\]

%dsp-level:0
\[
k=1:\ e_1\ \text{unread}:\ \operatorname{cov}=0:\ 1+\operatorname{cov}=1\implies0;
\]

%dsp-level:0
\[
\qquad \text{closed walk}:\ \text{every edge read at both ends (wrap at the start)}:\ \operatorname{cov}=2\implies0
\]

%dsp-level:0
\[
k=2\ \text{open}:\ \text{each edge read exactly once}:\ 1+\operatorname{cov}=2:\ w=\varepsilon(R)\ \sigma\text{-constant}:\ \sum_\sigma w=\varepsilon(R);
\]

%dsp-level:0
\[
\qquad k=2\ \text{closed (digon)}:\ \operatorname{cov}=2\implies0
\]

%dsp-level:0
\[
\text{surviving}\ D\iff\text{partition into open}\ 2\text{-threads};\quad\text{distinct surviving}\ D\ \text{give identical weights}\implies\text{count}\ N_{E1}(R);
\]

%dsp-level:0
\[
\qquad |R|=2\#D\implies\varepsilon(R)=2^{|R|}
\]

%dsp-level:0
\[
(2):\ v\ \text{type}(1,1):\ \text{literal read forced}:\ \text{covers both edges of}\ v;
\]

%dsp-level:0
\[
\quad e=(x{\to}y),\ x,y\ \text{both}(1,1):\ \operatorname{cov}=2\implies0
\]

%dsp-level:0
\[
v\ (1,1),
\]

%dsp-level:0
\[
\ \text{literal thread closed}:\ \text{walk}\ s\to v\to s:\ s(c)
\]

%dsp-level:0
\[
= t(c):\ \text{closed}\ [\mathrm{D1.4}]:\ \text{wrap read at}\ s:\ \text{both edges}\ \operatorname{cov}
\]

%dsp-level:0
\[
= 2\qquad \implies 0
\]

%dsp-level:0
\[
e = (x{\to}y),
\]

%dsp-level:0
\[
\ x\ (1,1),
\]

%dsp-level:0
\[
\ y\ \text{higher}:\ e\ \text{read literally at}\ x:\ \text{any}\ M\text{-read of}\ e\ \text{at}\ y:\ \operatorname{cov}
\]

%dsp-level:0
\[
= 2
\]

%dsp-level:0
\[
\implies 0:\ \text{surviving}\ M\ \text{pair only within}\ R_{\mathrm{hh}}
\]

%dsp-level:0
\[
\text{orbits of}\ M|_{R_{\mathrm{hh}}}:\ \text{as in}\ (1):\ \text{only open}\ 2\text{-threads survive};
\]

%dsp-level:0
\[
\quad\text{forced literal threads cover}\ R\setminus R_{\mathrm{hh}}\ \text{once each}
\]

%dsp-level:0
\[
\implies N_{E2}(R)
\]

%dsp-level:0
\[
= \#\{\text{open thread partitions of}\ R_{\mathrm{hh}}\}
\]

%dsp-level:0
\[
(3):\ N_{E2}\le N_{E1}:\ M\mapsto\{\text{literal threads}\}\cup\{R_{\mathrm{hh}}\ \text{threads}\}:\ \text{well-defined and injective by}\ (2);
\]

%dsp-level:0
\[
\text{witness rows}:\ \text{direct computation}\ [\text{machine}:\ \text{classifier}=\text{brute force on each}];
\]

%dsp-level:0
\[
\qquad (4):\ \text{no}(1,1)\ \text{vertices}:\ R_{\mathrm{hh}}=R:\ \text{E2 sum}=\text{E1 sum}\ [\text{as in}\ (1)]
\]

%dsp-level:0
\[
(5):\ \text{exhaustive enumeration as scoped}\ [\text{machine}].\ \blacksquare
\]

\textbf{Phys.} Let the ledger compare phases along chains instead of at vertex types, and the exact-cover rule keeps its shape — edges tiled by two-step threads, each read exactly once — but the tiling grammar changes: \textit{who} may host a thread junction (any vertex, or only the over-subscribed ones) and \textit{which} loops count as closed become the defining choices. The three grammars provably disagree, in both directions, on four-edge patterns. One verdict is common to all three: the digon alone still sums to zero — the minimal fermion floor's exclusion from the amplitude layer is not an artifact of the literal domain.

\vspace{2pt}\noindent\rule{\textwidth}{0.4pt}\vspace{2pt}

\end{theorem}

\begin{theorem}[{\texorpdfstring{FT.75 — Domain-relativity of the amplitude layer: lock revival, floor revival, parity, and the invariant bridge gap}{FT.75 - Domain-relativity of the amplitude layer: lock revival, floor revival, parity, and the invariant bridge gap}}]\label{tag:FT-75}
\textbf{Thm.} Under Def-H.27/FT.74, at (S1): $(1)$ \textbf{Single vertex — the locked sector comes alive, in two closed forms.} For the balanced single-vertex star ($n_{\mathrm{in}}+n_{\mathrm{out}}=n\in\{4,5,6\}$, external endpoints simple), $Z^{\pm}_{\mathrm{byp}}$ is domain-invariant $=1+4n_{\mathrm{in}}n_{\mathrm{out}}=(17,25,37)$ [FT.67(5)], while

%dsp-level:0
\[
Z^{\pm}_{\mathrm{lock}}:\quad\text{literal }0;\qquad E2=E1=\sum_{a\ge2}\binom{n_{\mathrm{in}}}a\binom{n_{\mathrm{out}}}a\,a!\,4^{a}=(32,\ 96,\ 672),
\]

totalling $(49,121,709)$ across both sectors. \textbf{FT.67(5)'s "its locked sector vanishes identically" is a literal-domain theorem}: in the chain-local domains an isolated locked vertex has strictly positive amplitude — the type-$(a,a)$ patterns thread \textit{through the vertex itself}. Second family (the locked digon hub): a locked $(a,a)$ hub hosting $a$ digons with $a$ simple centres has

%dsp-level:0
\[
\text{literal}=4^{a},\qquad E2=0\ \ (\text{every literal thread closed}),\qquad E1=\bigl(!a\bigr)\,4^{a},
\]

$!a$ the derangement number (each thread must leave its own centre: a fixed-point-free permutation of the $a$ digons), machine-verified for $a=2,3,4,5$ with brute force at $a\le3$: $(16,128,2304,45056)$ against literal $(16,64,256,1024)$. The tetrahedron's $16$ locked digon-covers are exactly the $a=2$ ($12$; value $16$) and $a=3$ ($4$; value $128$) instances [FT.69(5)]. $(2)$ \textbf{Floor revival — the dichotomy is domain-relative.} Every fully crystallized floor has no type-$(1,1)$ vertex, so E2 $\equiv$ E1 there [FT.74(4)], and the extended amplitude is strictly positive whenever the thread-partition count is:

%dsp-level:0
\[
N(\text{tetra floor})=32\ \ (\text{value}\ 2^{17}),
\]

%dsp-level:0
\[
\quad N(\text{cube floor})=2304\ \ (2304\cdot2^{24}),
\]

%dsp-level:0
\[
\quad N(\text{octa floor})=72384\ \ (72384\cdot2^{24})
\]

— versus literal amplitude $0$ [FT.67(3a)]. Dodecahedron/icosahedron floor counts deferred (branch enumeration beyond the bounded rule; recorded, not guessed). \textbf{Bidirectional polygon floors obey an exact parity theorem}:

%dsp-level:0
\[
N\bigl(\Sigma_{P_L}\ \text{floor}\bigr)=\begin{cases}4 & L\ \text{even}\quad(\text{value}\ 4\cdot2^{2L}),\\[2pt] 0 & L\ \text{odd},\end{cases}
\]

machine-verified $L=3..10$ $(0,4,0,4,0,4,0,4)$. \textbf{FT.68(4)'s verdict "fuzzing into solids — no at the amplitude layer" is the literal-domain verdict; in the chain-local domains every computed even polygon floor and the three computed Platonic floors carry strictly positive amplitude} (odd polygon floors stay dead by parity; the dodecahedron/icosahedron floors are open counts, not known-zero). $(3)$ \textbf{Directed cycles split the domains.} The single-orientation directed cycle $C_L$: literal $0$; E2 $0$ for every $L\ge2$ (all vertices type-$(1,1)$: conditions (2a)-(2b) fail); E1 $=2\cdot2^{L}$ for even $L$ — the cycle's two perfect matchings — and $0$ for odd $L$ ($|R|$ odd). Machine: $L=3..10$, values $(0,32,0,128,0,512,0,2048)$. The on-shell winding ladder's even rungs [FT.22] are exactly the E1-survivors; the digon ($L=2$) is dead in all three [FT.74(3)]. $(4)$ \textbf{Bound content returns — the triangle census, and the digon-cover split.} On the tetrahedron (exhaustive, all $4096$ patterns): $2397$ patterns contain a directed $3$-cycle; of those, nonzero amplitude — literal $0$, E2 $833$, E1 $1025$; restricted to the bypass sector: $0/40/88$ (the FT.68(1)/FT.69(4) anchor "$0$ nonzero of $252$" is the bypass-only carrier count and stands). Exhibit: $R=\{(0{\to}1),(0{\to}3),(1{\to}2),(2{\to}0)\}$, threads $\{0{\to}1{\to}2\},\{2{\to}0{\to}3\}$, E1 value $16$ (classifier $=$ brute force). \textbf{The cycle-alternation lemma [FT.69(3)] is a literal-domain theorem and fails in both extensions} — chain-local threading dissolves odd cycles into open threads. The digon rows split by sector (machine, exhaustive): the $48$ \textit{bypass} digon-covers of FT.69(4) are killed by \textbf{both} extensions ($N_{E2}=N_{E1}=0$ on all $48$ — a dissolved digon's centre has no alternative threading), while the $16$ \textit{locked} digon-covers are killed by E2 and kept by E1 [clause (1)'s derangement family]. $(5)$ \textbf{Skeleton totals (exhaustive; the literal column reproduces the stored anchors exactly).}

%dsp-level:0
\[
\begin{array}{l|ccc}
& Z^{\pm}_{\mathrm{byp}} & Z^{\pm}_{\mathrm{lock}} & \#\text{nonzero (byp/lock)}\\
\hline
\text{tetra, lit} & 1057 & 448 & 85/16\\
\text{tetra, E2} & 5921 & 966656 & 173/934\\
\text{tetra, E1} & 6881 & 1001536 & 227/1166\\
\hline
P_3\ \text{bidir, lit/E2/E1} & 25/121/121 & 48/0/48 & 7/13/13\ \text{byp};\ 3/0/3\ \text{lock}\\
P_4\ \text{bidir, lit/E2/E1} & 353/225/289 & 64/1536/2112 & 29/21/23;\ 4/9/13\\
P_5\ \text{bidir, lit/E2/E1} & 1161/361/1801 & 80/5760/8400 & 51/31/61;\ 5/30/60\\
P_6\ \text{bidir, lit/E2/E1} & 3121/3697/7985 & 1632/50944/80224 & 97/79/135;\ 21/67/154
\end{array}
\]

The $P_3$ locked row exhibits the second incomparability at skeleton level: E2 kills every literal locked cover there (each is a digon-pair hub — a closed literal thread), E1 keeps all $48$. $(6)$ \textbf{Norm-layer invariance, with the exact configuration count.} In every domain every factor of $w$ is a sign — $|w(R,\sigma)|=|w(R,\sigma,D)|=|w(R,\sigma,M)|=1$ (machine: asserted on every $(R,D,\sigma)$ of the bidirectional triangle) — so CT2.8's norms $\sum|w|^2$ count configurations: strictly positive, cancellation-blind, and \textit{form-identical across domains}. At the pattern-sign layer the count is exactly

%dsp-level:0
\[
\#\{(R,\sigma)\}=\sum_{R\subseteq E_{\mathrm{dir}}}2^{|R|}=3^{\,|E_{\mathrm{dir}}(\Gamma)|}
\]

%dsp-level:0
\[
\qquad (\text{machine:}\ 3^6=729,\ 3^8=6561,\ 3^{12}=531\,441\ \text{for}\ P_3,P_4,\text{tetra}),
\]

each directed edge contributing (absent, $+$, $-$); the extended domains refine this set by an inner sum over $D$/$M$ and cannot change it. \textbf{Item 13's bridge question is therefore invariant under the $\aleph$-domain choice}: no domain rescues or worsens it — exit (b) must operate on resolvent structure in every domain alike. $(7)$ \textbf{Gap 16(d) verdict.} The $\aleph$-domain question of the EM row is now exact data: three domains, three provably distinct path integrals, a complete witness catalogue [(3)–(5)], the frozen answer unchanged (D3.3's literal domain governs A8's $\aleph$), and the two chain-local variants computed but not adopted. The residue — whether FT.62(5)'s observed-variant EM-row reads should govern the \textit{weight} as well — is an author's declaration whose consequences are now fully priced on both sides. $(8)$ \textbf{What is not domain-relative.} The cap census $c_K$ [FT.69(1)], the admissibility fractions $f_{\mathrm{adm}}(n,p)$, the phase-diagram reads $B_\Gamma(p)$, $L_\Gamma(p)$, $p^*_\Gamma$, $\langle K\rangle_{\mathrm{byp}}(p)$ [Def-H.26, FT.70] and the dust equation of state $\langle\rho_{D^-}\rangle_\Gamma(p)$ with $w=0$ [FT.71, CT4.5] are counting/type-level objects: none involves $\aleph$, so all are identical in the three domains — the phase boundaries of FT.70(2) and the dust densities of FT.71 stand as computed. Only the amplitude-side sums $Z^{\pm}$, the survivor grammar, and the censuses built from them are domain-relative. Covering \textit{enumeration} likewise splits: the literal covers of FT.69(2) are a proper subset of the E1-survivors on every computed skeleton [clause (5)].

\textbf{Req.} Def-H.25, Def-H.26, Def-H.27; FT.21, FT.22, FT.28, FT.34, FT.45, FT.62, FT.63, FT.67, FT.68, FT.69, FT.70, FT.71, FT.74; A8; D1.4, D1.5, D3.3; CT2.8 (Part 2); CT4.1, CT4.5 (Part 2); perfect matchings of cycles, derangement counts, parity arguments [math]; machine: exhaustive enumeration as scoped.

\textbf{Pf.}

%dsp-level:0
\[
(1):\ \text{type}(a,a)\ \text{pattern at the star}:\ a!\ \text{bijections in}\leftrightarrow\text{out},
\]

%dsp-level:0
\[
\ \text{all threads open (distinct externals)}:\ a!\,
\]

%dsp-level:0
\[
4^{a};
\]

%dsp-level:0
\[
\quad\text{pattern count}\ \binom{n_{\mathrm{in}}}a\binom{n_{\mathrm{out}}}a
\]

%dsp-level:0
\[
\text{types}\ (a,b),\ a\ne b:\ |R|=a+b\ \text{odd or unthreadable}:\ 0;
\]

%dsp-level:0
\[
\qquad \text{locked}\iff a+b\ge4\iff a\ge2\ \text{on the balanced support}
\]

%dsp-level:0
\[
\text{digon hub}:\ \text{literal}=4^{a}\ [\mathrm{FT.67}(2),\ |\mathcal C|=a];
\]

%dsp-level:0
\[
\quad E2:\ \text{centre}\ c_i\ \text{type}(1,1),\ \text{thread}\ h\to c_i\to h\ \text{closed}\ [\mathrm{D1.4}]:\ N_{E2}=0\ [(2b)]
\]

%dsp-level:0
\[
E1:\ \text{a thread through the hub is}\ (c_i{\to}h),
\]

%dsp-level:0
\[
(h{\to}c_j),
\]

%dsp-level:0
\[
\ i
\]

%dsp-level:0
\[
\ne j:\ \text{partitions}
\]

%dsp-level:0
\[
= \text{fixed-point-free permutations}:\ N_{E1}
\]

%dsp-level:0
\[
= \ !a:\ \text{value}\ !a\cdot2^{2a}\qquad = \ !a\cdot4^{a}
\]

%dsp-level:0
\[
\text{tetrahedron instances}:\ a=2\ (12\ \text{covers}),\ a=3\ (4)\ [\mathrm{FT.69}(5);\ \text{machine}]
\]

%dsp-level:0
\[
(2):\ \text{floors}:\ \text{types}(d,d),
\]

%dsp-level:0
\[
\ d
\]

%dsp-level:0
\[
\ge 2:\ \text{no}(1,1)\ \text{vertices}:\ E2
\]

%dsp-level:0
\[
\equiv E1\ [\mathrm{FT.74}(4)];
\]

%dsp-level:0
\[
\quad\text{counts}:\ \text{thread-partition enumeration}\ [\text{machine}]
\]

%dsp-level:0
\[
\text{parity}:\ \text{mixed thread at}\ i:\ \text{in-edge}\ (i{-}1{\to}i)\ [\text{fwd}]\ \text{with out-edge}\ (i{\to}i{-}1)\ [\text{bwd}]:\ \text{walk}\ i{-}1\to i\to i{-}1:\ \text{closed}:\ 0;
\]

%dsp-level:0
\[
\quad\text{likewise bwd-in with fwd-out}
\]

%dsp-level:0
\[
\implies\text{every thread pure-orientation}\ [\text{machine witness at}\ L{=}4:\ 4\ \text{partitions},\ \#\mathrm{FF}=\#\mathrm{BB}=L/2=2\ \text{in each}];
\]

%dsp-level:0
\[
\quad L=\#\text{fwd edges}=2\,\#\mathrm{FF}\implies L\ \text{even}
\]

%dsp-level:0
\[
\text{even}\ L:\ \text{tilings of a directed}\ L\text{-cycle by}\ 2\text{-paths}
\]

%dsp-level:0
\[
= \text{perfect matchings of}\ C_L
\]

%dsp-level:0
\[
= 2\ [\text{math}]:\ N
\]

%dsp-level:0
\[
= 2_{\mathrm{fwd}}\cdot2_{\mathrm{bwd}}\qquad = 4;
\]

%dsp-level:0
\[
\qquad \text{odd}\ L:\ |R|=2L\ \text{even but}\ L\ne2\#\mathrm{FF}:\ N=0
\]

%dsp-level:0
\[
(3):\ C_L:\ \text{every vertex type}(1,1):\ \text{literal/E2 forced reads}:\ \operatorname{cov}=2\ \text{per edge}:\ 0;
\]

%dsp-level:0
\[
\qquad E1:\ \text{tilings of the directed cycle}:\ \text{even}\ L:\ 2;\ \text{odd}\ L:\ L\ne2\#\mathrm{FF}:\ 0
\]

%dsp-level:0
\[
(4),
\]

%dsp-level:0
\[
(5):\ \text{exhaustive machine enumeration with the FT.74 classifiers (brute-force verified)};
\]

%dsp-level:0
\[
\qquad \text{exhibit threads checked edge-by-edge}
\]

%dsp-level:0
\[
(6):\ |\sigma_e|=1,\ \bigl|(-1)^{\aleph}\bigr|=1,\ \mathcal J=1\ [(S3)]\implies|w|=1;
\]

%dsp-level:0
\[
\qquad \#\{(R,\sigma)\}=\textstyle\sum_k\binom{|E_{\mathrm{dir}}|}k2^{k}=(1+2)^{|E_{\mathrm{dir}}|}\ [\text{binomial theorem}]
\]

%dsp-level:0
\[
(7):\ \text{conjunction of}\ (1)\text{–}(6);
\]

%dsp-level:0
\[
\qquad (8):\ \text{each read named in Def-H.26/FT.69(1)/FT.71 is a}\ \mu_p\text{-average of type-level data}:\ \aleph\ \text{absent}.\ \blacksquare
\]

\textbf{Phys.} Let the phase-comparisons ride chains instead of vertex types, and the amplitude layer changes character: locked vertices acquire weight of their own — a single over-subscribed vertex hums with amplitude $32$, $96$, $672$ as its degree climbs, and a hub holding $a$ digons answers with the derangement count times its literal voice — and the fully grown crystals, silent in the literal ledger, ring out: the tetrahedron floor at $2^{17}$, the cube at $2304\cdot2^{24}$, the octahedron at $72384\cdot2^{24}$, every even polygon ring at four times its full register. Odd rings stay dark by a pure parity argument, triangles come back dissolved, and the digon stays dark in every grammar.

\vspace{2pt}\noindent\rule{\textwidth}{0.4pt}\vspace{2pt}

\end{theorem}

\phantomsection
\section[Tier F.26 — The Platonic Pole Spectrum: Exact Multiplets,…]{\texorpdfstring{Tier F.26 — The Platonic Pole Spectrum: Exact Multiplets, the Bounded-Hierarchy Ceiling, and the SM-Hierarchy Audit}{Tier F.26 - The Platonic Pole Spectrum: Exact Multiplets, the Bounded-Hierarchy Ceiling, and the SM-Hierarchy Audit}}

\begin{theorem}[{\texorpdfstring{FT.155 — Minimality of the defect interface}{FT.155 - Minimality of the defect interface}}]\label{tag:FT-155}
\textbf{Thm.} The non-derived content of Def-H.37 is exactly one positive scalar ($\hat\mu'^2$; equivalently $\hat\mu^2$ at leading order, related by the recorded matching) and the holonomy identification $\mathcal U_c=\mathcal U^{\mathrm{spin}}_c$. Precisely:

%dsp-level:0
\[
(1)\quad\textbf{shape uniqueness:}\ \ \text{any state-independent defect read on}\ SU(2)\ \text{transport has the form}\ m^2_{\mathrm{pole}}
\]

%dsp-level:0
\[
= F(\Theta_R),
\]

%dsp-level:0
\[
\ \ F
\]

%dsp-level:0
\[
= 4\hat\mu'^2\sin^2(\cdot/4)\ \text{exact},
\]

%dsp-level:0
\[
\ \sim\hat\mu'^2(\cdot)^2/4\ \text{leading}
\]

%dsp-level:0
\[
\qquad\bigl[\mathrm{FT.54}(1):\ (U-1)^\dagger(U-1)=4\sin^2(\tfrac\phi2)\,\mathbb I\ \text{axis- and state-free};\ \mathrm{FT.54}(2):\ X^2=|\vec x|^2\mathbb I\ \text{fixes the leading term};\ \mathrm{EW.33}\ \text{is that Taylor term}\bigr];
\]

%dsp-level:0
\[
(2)\quad\textbf{combinatorial exclusion:}\ \ \text{no rational-coupling formula on integer-valued combinatorial reads reaches the computed pole factors — the five Platonic factors}\ \{2,\ 2{-}\sqrt2,\ 1,\ 2{-}2\cos\tfrac\pi{10},\ 2{-}\sqrt3\}\ \text{generate a degree-}16\ \text{extension of}\ \mathbb Q
\]

%dsp-level:0
\[
\qquad\bigl[\mathrm{FT.73}(2)\text{–}(3):\ \text{irreducibility machine-verified};\ \sqrt5\notin\mathbb Q(\sqrt2,\sqrt3);\ \mathbb Q(2\cos\tfrac\pi{10})\ \text{has unique quadratic subfield}\ \mathbb Q(\sqrt5);\ \text{single quadratic couplings excluded by the exact ratio}\ 4{-}2\sqrt2{+}2\sqrt3{-}\sqrt6\bigr];
\]

%dsp-level:0
\[
\text{the}\ D^-\text{-drag/dielectric candidates are independently refuted by response non-existence}\ [\mathrm{FT.73}(1);\ \mathrm{CT4.2},\ \mathrm{D1.11},\ \mathrm{CT4.4}];
\]

%dsp-level:0
\[
(3)\quad\textbf{range:}\ \ \text{on the Platonic skeletons the read takes finitely many values in}\ [0,4\hat\mu'^2]\ (\text{period}\ P_S=4\pi/\delta_v;\ P_S/2+1\ \text{distinct values};\ \text{monotone to the ceiling});
\]

%dsp-level:0
\[
\ \text{widest intra-solid ratio}\ \sqrt{4/D_S}
\]

%dsp-level:0
\[
\le 6.3925\ [\mathrm{FT.76}(3)];
\]

%dsp-level:0
\[
(4)\quad\textbf{rotor reparametrization:}\ \ \text{any}\ k^2/R^2\text{-type rotational mass formula is a reparametrization of Def-H.37(2)'s leading regime}\ (R^{-2}\equiv\hat\mu'^2\Theta_R^2/4),
\]

%dsp-level:0
\[
\ \neg\text{an independent mechanism};
\]

%dsp-level:0
\[
\ \text{its approximation error at the documented uniform-minimal windings is}\ 18.9\%\text{–}91.0\%\ [\mathrm{FT.87}(1),(4)].
\]

Hence the replacement search space of the interface contains no combinatorial, dielectric, or rotor species; the derivation of the transfer itself — why the seed's transport defect is the terminus's pole displacement — remains application-layer content.

\textbf{Req.} Def-H.37; Def-H.21; EW.32, EW.33 (Part 2); FT.54(1)–(2); FT.73(1)–(3); FT.76(3); FT.87(1), (4); CT4.2, CT4.4 (Part 2); D1.11; field-degree and enclosure arithmetic [math; machine: \texttt{checks/\allowbreak{}verify\_\allowbreak{}H\_\allowbreak{}stationarity.py}, \texttt{checks/\allowbreak{}verify\_\allowbreak{}H\_\allowbreak{}phasediagram.py}, \texttt{checks/\allowbreak{}verify\_\allowbreak{}H\_\allowbreak{}polespectrum.py}, \texttt{checks/\allowbreak{}verify\_\allowbreak{}H\_\allowbreak{}tumbling.py} pins].

\textbf{Pf.}

%dsp-level:0
\[
(1):\ \mathrm{FT.54}(1)\ \text{verbatim}:\ \|(U-1)v\|^2
\]

%dsp-level:0
\[
= 2-2\cos\phi\ \forall\ \text{unit}\ v,
\]

%dsp-level:0
\[
\ \text{axis-free}
\]

%dsp-level:0
\[
\implies F(\Theta_R)\ \text{shape with}\ \phi
\]

%dsp-level:0
\[
= \Theta_R/2\ [\mathrm{Def\text{-}H.21}(6)];
\]

%dsp-level:0
\[
\qquad (2):\ \mathrm{FT.54}(2)
\]

%dsp-level:0
\[
(2):\ \mathrm{FT.73}(2)\text{–}(3)\ \text{verbatim}:\ \text{degree-}16\ \text{extension}\ [y^4{-}5y^2{+}5,\ x^4{-}8x^3{+}19x^2{-}12x{+}1\ \text{irreducible};\ \text{machine}];
\]

%dsp-level:0
\[
\qquad \mathrm{FT.73}(1):\ \text{response non-existence}
\]

%dsp-level:0
\[
(3):\ \mathrm{FT.76}(3)\ \text{verbatim}:\ \text{finite value set, monotone cone, ceiling}\ 4\hat\mu'^2;
\]

%dsp-level:0
\[
\qquad (4):\ \mathrm{FT.87}(1):\ \text{Taylor identification};\ \mathrm{FT.87}(4):\ \text{error table}\ [\text{exact enclosures}]
\]

%dsp-level:0
\[
\text{search-space clause}: (1)\ \text{fixes the shape};
\]

%dsp-level:0
\[
\ (2)\ \text{excludes the combinatorial species};
\]

%dsp-level:0
\[
\ (4)\ \text{absorbs the rotor species};
\]

%dsp-level:0
\[
\ \text{the residue is the transfer's derivation}\.\ \blacksquare
\]

\textbf{Phys.} The defect interface is priced: one scalar and one identification, and every competing shape is excluded by theorem — each alternative mass formula the corpus has tested is either inside its Taylor regime, outside its range, or outside its field. What remains open is not the form of the transfer but its locus: why the seed's failure to return to itself around a bound chain is that chain's pole mass.

\end{theorem}

\medskip\noindent\textit{\texorpdfstring{FT.77 — The SM-hierarchy audit: a bounded-hierarchy negative result}{FT.77 - The SM-hierarchy audit: a bounded-hierarchy negative result}.}\ Structural-constraints record: see Appendix B, Section B.2.3 of the Addendum.\par\medskip

\phantomsection
\section[Tier F.27 — The Cover Orbit Decomposition: the Ansatz…]{\texorpdfstring{Tier F.27 — The Cover Orbit Decomposition: the Ansatz Split, the Combinatorial Ceiling, and the Bridge Census}{Tier F.27 - The Cover Orbit Decomposition: the Ansatz Split, the Combinatorial Ceiling, and the Bridge Census}}

\begin{definition}[{\texorpdfstring{Def-H.28 — Cover classes, and the two class-mass ansätze}{Def-H.28 - Cover classes, and the two class-mass anstze}}]\label{tag:Def-H-28}
\textbf{Def.} Let $\Gamma$ be a skeleton [Def-H.25] with directed edge set $E_{\mathrm{dir}}(\Gamma)$ and automorphism group $\operatorname{Aut}(\Gamma)$ acting on vertices, hence on $E_{\mathrm{dir}}(\Gamma)$ by $(a,b)\mapsto(p(a),p(b))$. A \textbf{cover class} (orbit class) is an $\operatorname{Aut}(\Gamma)$-orbit $\mathcal O$ of crystallization patterns $R\subseteq E_{\mathrm{dir}}(\Gamma)$; its \textbf{multiplicity} is $\Omega_{\mathcal O}=|\mathcal O|$, with $\Omega_{\mathcal O}\cdot|\operatorname{Stab}(R)|=|\operatorname{Aut}(\Gamma)|$ [orbit–stabilizer] and hence $\Omega_{\mathcal O}\mid|\operatorname{Aut}(\Gamma)|$. Invariants of a class (all constant on $\mathcal O$, machine-verified): $K=|R|$; admissibility ($k_v\le3\ \forall v$) and lock status; cover status and $|\mathcal C|$; amplitude $4^{|\mathcal C|}$ [FT.67(2)]; the canonical directed-cycle content of $R$ (its \textbf{signature}: which cycles the cover shadows); the aspect mask $\mathcal A(\max_v k_v)$ and observable mask dimension [CT0.13, CT0.14]. The two class-mass reads:

%dsp-level:0
\[
\textbf{A (corpus-native)}:\quad m_A(\mathcal O)=K\,\frac{\ln2}{\mathfrak h}
\]

%dsp-level:0
\[
\qquad [\text{FT.68}(1)/\text{FT.70}(5)\text{'s}\ \langle K\rangle\ln2/\mathfrak h\ \text{read, per class};\ \neg\text{new datum}],
\]

%dsp-level:0
\[
\textbf{B (external ansatz, NOT a corpus read)}:\quad m_B(\mathcal O)=\Bigl(K-\alpha\ln\Omega_{\mathcal O}\Bigr)\frac{\ln2}{\mathfrak h},
\]

%dsp-level:0
\[
\qquad \alpha\ \text{a formal parameter, no value declared}.
\]

\textbf{Rmk (B's exact form).} $\Omega\mid|\operatorname{Aut}|$ and $|\operatorname{Aut}|=2^3\cdot3$ (tetrahedron, $24$) or $2^4\cdot3$ (cube, $48$) force $\Omega=2^a3^b$, so $\ln\Omega=a\ln2+b\ln3$ \textbf{exactly} and every Ansatz-B mass lies in the $\mathbb Q$-span of $\{1,\ln2,\ln3\}$ with the single free coefficient $\alpha$: $(K,a,b)$ is a complete exact label. No floating-point logarithm is evaluated anywhere. The corpus supplies no free-energy/entropic mass read (CT4.7/CT4.10's entropy triangle and FT.71's dust equation of state are not mass formulas), so B is recorded as an \textit{ansatz}, computed as directed, never adopted, and no result below is used as a premise by any other tier. \textbf{Rmk ($p$-independence, verified).} Because $K$ is constant on $\mathcal O$, the directive's orbit-wise weighted average reduces identically:

%dsp-level:0
\[
\langle K\rangle_{\mathcal O}=\frac{\sum_{\sigma\in\mathcal O}K\,p^{K}(1-p)^{|E|-K}}{\sum_{\sigma\in\mathcal O}p^{K}(1-p)^{|E|-K}}=K\qquad\forall p\in(0,1),
\]

so Ansatz A is $p$-independent by proof, not by observation. What Cal-F's $r$ (hence $p=r/(1+r)$) does move is the \textit{measure} of each class, $\mu_p(\mathcal O)=\Omega_{\mathcal O}\,p^{K}(1-p)^{|E|-K}$, and with it every ensemble read — FT.70(5)'s $\langle K\rangle_{\mathrm{byp}}(p)$, FT.71's $\langle\rho_{D^-}\rangle_\Gamma(p)$, and the phase boundary $p^*_\Gamma$ [FT.70(2)]. Class masses and class weights are different objects; only the latter flicker.

\vspace{2pt}\noindent\rule{\textwidth}{0.4pt}\vspace{2pt}

\end{definition}

\begin{theorem}[{\texorpdfstring{FT.79 — The combinatorial ceiling and the cover-hierarchy audit (second clean negative)}{FT.79 - The combinatorial ceiling and the cover-hierarchy audit (second clean negative)}}]\label{tag:FT-79}
\textbf{Thm.} Under Def-H.28; the SM hierarchies of clause (4) are \textbf{external measured data}, marked, entering no derivation and no identification (A §0.2). $(1)$ \textbf{Ansatz A spectrum (exact).} Cover masses are $m_A=2|\mathcal C|\ln2/\mathfrak h$ [FT.67(2): $|R|=2|\mathcal C|$ on covers], so the nonzero values are the even ladder $K\in\{2,4,6\}$ (tetrahedron) and $\{2,4,6,8,10,12\}$ (cube); pooled, $6$ distinct masses with exactly $11$ pairwise ratios

%dsp-level:0
\[
\Bigl\{\tfrac65,\tfrac54,\tfrac43,\tfrac32,\tfrac53,2,\tfrac52,3,4,5,6\Bigr\}\subset\mathbb Q,
\]

lightest $2\ln2/\mathfrak h$, heaviest $12\ln2/\mathfrak h$, median $7$ (of the distinct values) / $8$ (orbit-weighted on the cube). Every ratio is a quotient of integers $\le6$. $(2)$ \textbf{The ceiling theorem.} For any cover on any skeleton $\Gamma$: $K=2|\mathcal C|$ [FT.67(2)]; admissibility gives $\sum_vk_v=2K\le3V$ [CT0.11], so

%dsp-level:0
\[
|\mathcal C|\le\tfrac{3V}{4}\ (\text{bypass}),\qquad |\mathcal C|\le V\ (\text{locked, centres are distinct vertices}),
\]

and the widest Ansatz-A mass ratio available on $\Gamma$ is $\le3V/4$ (bypass) or $\le V$ (locked). Spanning $1{:}R$ therefore \textbf{requires} $V\ge4R/3$ (bypass) or $V\ge R$ (locked):

%dsp-level:0
\[
\begin{array}{c|cccccc}
R & 20 & 207 & 600 & 900 & 3477 & 80000\\
\hline
V\ge\ (4R/3) & 27 & 276 & 800 & 1200 & 4636 & 106667
\end{array}
\]

against $V=4$ (tetrahedron, max ratio $3$) and $V=8$ (cube, max ratio $6$). The cover sector's dynamic range is bounded by \textbf{skeleton size}, exactly as the pole sector's was bounded by the defect angle [FT.76(3)] — and both bounds evaluate to $\approx6$ on the computed cells. $(3)$ \textbf{Ansatz B: no $\alpha$-free hierarchy.} The classes are the $(K,a,b)$ triples with $\ln\Omega=a\ln2+b\ln3$: $18$ on the tetrahedron's bypass sector — $19$ including its two locked cover classes, the extra one being $(K,a,b)=(6,2,0)$ — with $K\le6$, $a\le3$, $b\le1$, and $14$ on the cube's cover sector with $K\le12$, $a\le4$, $b\le1$, so $m_B/\!(\ln2/\mathfrak h)=K-\alpha(a\ln2+b\ln3)$ lies between $K-\alpha\ln|\operatorname{Aut}|$ and $K$. A ratio exceeding (2)'s bound can arise only where some class's mass passes through zero, i.e. at the tuned values $\alpha=K/(a\ln2+b\ln3)$ — a finite set, one per class, at each of which the ratio is \textit{infinite} rather than $207$. \textbf{The signature structure does not supply the $\ge3$ orders of magnitude required by Gap 7's residue}: the informational datum is capped by $3V/2$ and the entropic one by $\ln|\operatorname{Aut}|\le\ln48$; what the signature supplies is degeneracy, not range. $(4)$ \textbf{The audit (external data: charged leptons $1{:}207{:}3477$; up-type $1{:}600{:}80000$; down-type $1{:}20{:}900$).} Pool $=\{2,4,6,8,10,12\}$ (all cover masses of both skeletons, both sectors):

\begin{itemize}
\item pairwise, the closest cover ratio to each SM ratio is always the pool maximum $6$, giving discrepancies $\times2.7995$ (vs $m_\tau/m_\mu=16.797$), $\times3.333$ (vs $20$), $\times7.5$ (vs $45$), $\times22.22$ (vs $133.3$), $\times34.5$ (vs $207$), $\times100$ (vs $600$), $\times150$ (vs $900$), $\times579.5$ (vs $3477$), $\times13333.3$ (vs $80000$). The best, $\times2.7995$, is within an order of magnitude and is structurally void: it pairs our maximum with their minimum and extends to no triple.
\item aligned triples, exhaustive: the best is $(2,4,12)\mapsto1{:}2{:}6$ against every family, with discrepancy $\times150$ (down-type), $\times579.5$ (leptons), $\times13333.3$ (up-type). \textbf{No aligned triple is within a factor of $10$}, and the floor is exact: $900/6=150$.
\end{itemize}

$(5)$ \textbf{Dimension trap \#2, disclosed and rejected.} The amplitudes $4^{|\mathcal C|}\in\{1,4,16,64,256,1024,4096\}$ \textit{do} span $3.6$ orders, and produce the only near-coincidences in the whole audit: $256$ vs $207$ ($\times1.2367$), $4096$ vs $3477$ ($\times1.1780$), $16$ vs $16.797$ ($\times1.0498$), $1024$ vs $900$ ($\times1.1378$), $64$ vs $45$ ($\times1.4222$). All are rejected: $4^{|\mathcal C|}$ is the $Z$-weight of FT.67(2), not a mass; no corpus tag maps an amplitude to a mass, Def-H.37(2)/Def-H.38 declare no such read, and the square-root-corrected values ($16,64,4,32,8$) miss by $\times4.2$–$\times54$. Same species as the $m^2$-vs-$m$ trap [FT.77(3)] — recorded because it is exactly what a forced fit would seize on. $(6)$ \textbf{Verdict.} The cover hypothesis fails at the mass layer, for a structural reason dual to FT.77's and with no parameter fitted: the observable-sector masses are the even ladder $2|\mathcal C|\ln2/\mathfrak h$ capped by the cap census, so the reachable ratio range is $\le\max(3V/4,V)$ — six on the cube, three on the tetrahedron, against SM spans of $20$–$80000$. The messiness is genuine but sits in the degeneracies ($\Omega$ up to $48$; $264$ cube cover classes; $21$ signatures; chirality pairs), and degeneracy carries no hierarchy. \textbf{Neither the roof nor the fuzz yields the SM spectrum on the computed skeletons}; both sectors point to the same requirement — a much larger subnetwork ($V\ge1200$ for the down-type family alone) or a mass read of a different species.

\textbf{Req.} Def-H.25, Def-H.28; FT.44, FT.67, FT.68, FT.69, FT.70, FT.76, FT.77, FT.78; A §0.2 (policy); CT0.11 (Part 1); CT1.20 (Part 2); cap-census bound, exhaustive triple search, exact rational ratios [math; machine]. External SM hierarchies: comparison only, marked at (4).

\textbf{Pf.}

%dsp-level:0
\[
(1):\ |R|=2|\mathcal C|\ [\mathrm{FT.67}(2)];
\]

%dsp-level:0
\[
\qquad \text{ratios}=\{b/a:\ a<b\in\{2,4,6,8,10,12\}\}:\ \binom62=15\ \text{pairs},\ 11\ \text{distinct values}
\]

%dsp-level:0
\[
(2):\ 2K=\textstyle\sum_vk_v\le3V\ [\mathrm{CT0.11}:\ k_v\le3];
\]

%dsp-level:0
\[
\qquad |\mathcal C|\le V:\ \text{centres are distinct vertices, each contributing its two edges}
\]

%dsp-level:0
\[
\text{ratio}\le|\mathcal C|_{\max}/1;\qquad|\mathcal C|_{\max}\ge R\implies V\ge\tfrac{4R}3\ (\text{bypass}),\ V\ge R\ (\text{locked});\qquad\text{cube}:\ \tfrac{3\cdot8}4=6\ \checkmark
\]

%dsp-level:0
\[
(3):\ \Omega\mid|\operatorname{Aut}|=2^a3^b\implies\ln\Omega=a\ln2+b\ln3;
\]

%dsp-level:0
\[
\qquad m_B=0\iff\alpha=K/(a\ln2+b\ln3):\ \text{finitely many tuned values}
\]

%dsp-level:0
\[
\alpha\ \text{away from them}:\ |m_B|\ge\min_{\text{classes}}|K-\alpha L|>0:\ \text{ratios bounded by}\ \tfrac{K_{\max}+\alpha\ln48}{\min|K-\alpha L|}:\ \text{no}\ \alpha\text{-free large ratio}
\]

%dsp-level:0
\[
(4):\ \text{exhaustive over the pool}:\ 11\ \text{pairwise},\ \binom63=20\ \text{triples per family};
\]

%dsp-level:0
\[
\qquad \text{floor}:\ 900/6=150\ \text{exactly}
\]

%dsp-level:0
\[
(5):\ \text{direct evaluation of the stated pairs};\qquad\text{rejection}:\ 4^{|\mathcal C|}\ \text{is}\ Z\text{-weight}\ [\mathrm{FT.67}(2)],\ \neg\text{mass read}
\]

%dsp-level:0
\[
(6):\ \text{conjunction of}\ (2)\text{–}(5).\ \blacksquare
\]

\textbf{Phys.} Ask the covers for a mass spectrum and they answer with the even ladder two, four, six, eight, ten, twelve — six notes, all ratios of small integers, exactly the kind of clean answer the combinatorial hypothesis was meant to avoid. The reason is not that the covers are simple; there are two hundred and sixty-four of them on the cube alone, in chiral pairs, shadowing ten different cycle patterns. It is that a mass counted in crystallized edges cannot outrun the edge budget: three edges per vertex, four vertices, twelve in all. Hierarchy of the observed kind needs either a thousand-vertex subnetwork or a different question.

\vspace{2pt}\noindent\rule{\textwidth}{0.4pt}\vspace{2pt}

\end{theorem}

\begin{theorem}[{\texorpdfstring{FT.80 — The bridge census: covers are amplitude-real, norm-real, and spectrally un-indexed}{FT.80 - The bridge census: covers are amplitude-real, norm-real, and spectrally un-indexed}}]\label{tag:FT-80}
\textbf{Thm.} Under Def-H.16/Def-H.28, at (S1)–(S4): $(1)$ \textbf{What a cover class supplies, exactly.} For every class $\mathcal O$: the real-edge set and $K$; $|\mathcal C|$ open V-threads (each centre pairs one in-edge with one out-edge [FT.67(3d)]) $=$ transit content, bosonic-side by D1.5; the dissolved even closed chains with their exact $\delta h$ [D0.6: $\delta h=\#\{e\in\gamma\}$] $=2$ (tetrahedron) and $\in\{2,4\}$ (cube); the aspect mask $\mathcal A(\max_vk_v)$ and observable mask dimension $\in\{6,16,17,25\}$ [CT0.13, CT0.14]; the amplitude $4^{|\mathcal C|}\ne0$ [FT.67(2)]; and the CT2.8 norm weight $\Omega\cdot2^{K}$ (the class's $(R,\sigma)$ count, cancellation-blind). $(2)$ \textbf{What it does not supply.} Def-H.16's resolvent $Z(\mathcal R,\Sigma)=\langle\mathrm{out}|(1-\mathbf F_M^{(\Sigma)})^{-1}|\mathrm{in}\rangle$ with pole $\iff1\in\operatorname{spec}\mathbf F_M^{(\Sigma)}$ takes as arguments a \textbf{subnetwork $\Sigma$ with an observer channel $C$} (Def-H.1/Def-H.13, Def-H.16), the \textbf{colour data} that $\Pi_{\mathcal C}$ requires (FT.19/FT.27 — not determined by a crystallization pattern), and a \textbf{materialization} clause for whatever bound chains are to contribute poles (Def-H.16's own iff; M1/M2 species). None of the three is a function of $R$: no rule of Parts 1–3 maps a crystallization pattern to a $(\Sigma,C)$ pair with colour assignment. $(3)$ \textbf{Census (machine, every class).} Tetrahedron $7/7$ and cube $264/264$ cover classes have: transit content $|\mathcal C|\cdot\Omega$ threads ($180$ and $40\,200$ thread instances); bound-chain candidates exactly the even cycles ($84$ and $14\,832$ instances; every $\delta h\ge2$, satisfying Def-H.4's floor, at rung $n(v)=6\ge2$); amplitude nonzero; norm weight nonzero; and \textbf{zero spectral poles constructible within the declared inventory}. Covers are therefore neither amplitude ghosts nor poles: they are amplitude-real and norm-real, spectrally \textit{un-indexed}. $(4)$ \textbf{The exact resolvent condition a declaration would have to supply} (recorded, \textbf{not} adopted). Item 13's exit (b) is constructible only as a triple of declarations: (i) a rule $\mathcal O\mapsto(\Sigma(\mathcal O),C(\mathcal O))$; (ii) colour data making $\Pi_{\mathcal C}(\Sigma;C)$ computable, hence $F^{\Sigma}_{\mathrm{strong}}$ and the mask $\mathbf F_M^{(\Sigma)}$ [FT.19, FT.26]; (iii) a materialization clause for dissolved cycles (M2-species), after which the pole condition is exactly $1\in\operatorname{spec}\mathbf F_M^{(\Sigma(\mathcal O))}$ with FT.22's single-valuedness on each materialized chain. This is a interface-definition species extension of Def-H.16; it is enumerated here and \textbf{not} introduced. $(5)$ \textbf{Content bifurcation (the $\delta h/\delta\tau$ question).} The split is structural and exact, but it is a \textit{content-type} split, not a statistics theorem: every cover decomposes as $|\mathcal C|$ transit threads (D1.5: bosonic content) plus dissolved even closed chains (D1.5: fermionic content \textit{if materialized}; Def-H.4's termini are then candidate fermion loci, their $\delta h\ge2$ and rung $\ge2$ floors both satisfied). $\delta h$ is exact per chain ($2$ or $4$); $\delta\tau$ is a \textbf{schedule} datum (D0.7: retained ledger updates) that $R$ does not determine — Cal-F declares the flicker \textit{rates} $r=\Gamma_{J^+}/\Gamma_{J^-}$, not per-tick schedules — so the "finite $\delta h$ with $\delta\tau=0$ persistence" test is not evaluable from a pattern alone. The corpus contains no exchange operator and no permutation representation of particle interchange, so \textbf{the spin–statistics connection remains an external reading}; what is derived here is the bound/transit content split per class, and the FT.1/FT.38 flux bound $I\le\min(\delta h/\mathfrak h,4\delta\tau)\ln2$ applies to it only once a schedule is declared. $(6)$ \textbf{Consistency with the selection side.} Covers contain no \textit{materialized} closed chain, hence carry no Bott datum [Def-H.24, FT.66(4)]; their content sits wholly in the transit sector, which is the one sector where index-$0$ stability exists on such slices [FT.48(3), FT.66(4b)] — FT.68(4)(d)'s verdict recovered at class level, and the reason the cover sector is where any future selection mechanism would have to act.

\textbf{Req.} Def-H.1, Def-H.4, Def-H.13, Def-H.16, Def-H.24, Def-H.25, Def-H.28; FT.1, FT.19, FT.22, FT.26, FT.27, FT.38, FT.40, FT.48, FT.66, FT.67, FT.68, FT.78; D0.6, D0.7, D1.5, D2.6; A8; CT2.5, CT2.8 (Part 2); CT0.13, CT0.14 (Part 1); Def-H.16, Cal-F (referenced); census by exhaustive class enumeration [machine].

\textbf{Pf.}

%dsp-level:0
\[
(1):\ \text{each centre pairs one in- with one out-edge}\implies|\mathcal C|\ \text{threads}\ [\mathrm{FT.67}(3d)];
\]

%dsp-level:0
\[
\qquad \delta h(\gamma)=\#\{e\in\gamma\}\ [\mathrm{D0.6}]
\]

%dsp-level:0
\[
\text{norm weight}:\ \#\{(R,\sigma):\ R\in\mathcal O\}=\Omega\cdot2^{K}\ [\text{per-edge binary freedom}];\qquad|w|=1\ [\mathrm{FT.67}(6c)]
\]

%dsp-level:0
\[
(2):\ \text{arguments of}\ Z(\mathcal R,\Sigma):\ (\Sigma,C,\text{colour},\text{materialization});
\]

%dsp-level:0
\[
\quad R\mapsto\Sigma:\ \neg\exists\ \text{rule in Parts 1–3}\ [\text{inventory}]
\]

%dsp-level:0
\[
\Pi_{\mathcal C}\ \text{needs}\ \mathrm{Colour}(\Sigma)\ \text{and its sealing}\ [\mathrm{FT.19},\mathrm{FT.27}]:\ \text{colour is slot-algebra data},
\]

%dsp-level:0
\[
\ \neg\text{a crystallization datum}
\]

%dsp-level:0
\[
(3):\ \text{exhaustive per class}\ [\text{machine}];\qquad\text{amplitude}\ne0:\ \mathrm{FT.67}(2);\qquad\text{norm}>0:\ \mathrm{CT2.8},\ |w|=1
\]

%dsp-level:0
\[
(4):\ \text{Def-H.16's pole condition read with its own argument list};
\]

%dsp-level:0
\[
\qquad \text{FT.22 single-valuedness on materialized chains}
\]

%dsp-level:0
\[
(5):\ \mathrm{D1.5}\ \text{bound/transit}=\text{closed/open};\qquad\text{Def-H.4 floor}:\ \delta h\ge2\ \checkmark\ (2,4),\ n(v)\ge2\ \checkmark\ (6)
\]

%dsp-level:0
\[
\delta\tau:\ \text{schedule datum}\ [\mathrm{D0.7}];\quad\text{Cal-F supplies the rate datum},\ \neg\text{schedules};
\]

%dsp-level:0
\[
\qquad \neg\exists\ \text{exchange operator in the inventory}\implies\text{statistics external}
\]

%dsp-level:0
\[
(6):\ \neg\exists\ \text{materialized closed chain on a cover}\implies\neg\exists\ \operatorname{ind}_B\ \text{datum}\ [\mathrm{FT.66}(4)];
\]

%dsp-level:0
\[
\qquad \text{transit sector}:\ \mathrm{FT.48}(3).\ \blacksquare
\]

\textbf{Phys.} Ask whether the roof hangs from the fuzz, and the census answers precisely: the fuzz is real where the ledger can see it — every cover class carries nonzero amplitude and nonzero norm, thousands of open threads, and thousands of dissolved even loops with exact hop counts — and silent where a particle is defined. A pole needs a subnetwork, an observer channel, colour sealing and a materialization clause; a crystallization pattern supplies none of the four. So the covers are not ghosts and not particles: they are the amplitude-and-norm substrate on which a declaration would have to place a spectrum. The content split they \textit{do} carry is clean: everything threading is transit-side, every dissolved loop is a bound-side candidate with two or four hops, and the statistics names attached to those two sides remain vocabulary until an exchange structure is declared.

\vspace{2pt}\noindent\rule{\textwidth}{0.4pt}\vspace{2pt}

\end{theorem}

\phantomsection
\section[Tier F.28 — The Full-SM Completeness Audit: Particle…]{\texorpdfstring{Tier F.28 — The Full-SM Completeness Audit: Particle Content, the Gauge Blocks, the Rung Ledger, and the Deficit Table}{Tier F.28 - The Full-SM Completeness Audit: Particle Content, the Gauge Blocks, the Rung Ledger, and the Deficit Table}}

\begin{definition}[{\texorpdfstring{Def-H.29 — Completeness-audit vocabulary: read classes, census statuses, the external-data clause}{Def-H.29 - Completeness-audit vocabulary: read classes, census statuses, the external-data clause}}]\label{tag:Def-H-29}
\textbf{Def.} A \textbf{read class} is a species of quantity the corpus can attach to a configuration, chain or subnetwork: \textbf{(R1) binary registration} — $\sigma_e\in\{\pm1\}$ [CT0.9, A7/A8], the end-charge $q_{\mathrm{EM}}(c)=\tfrac12(\sigma_{s(c)}-\sigma_{t(c)})\in\{-1,0,+1\}$ [D3.4, EM.9], the vertex phase-comparison $f_v$ and $(-1)^{\aleph_v}$ [D3.3, FT.67(1)], the edge weight $G(e)=\sigma_e e^{i\pi\kappa_e}$ [D3.5]. \textbf{(R2) integer winding and ledger} — $\widehat{ds}^{\,2}$, $\Theta(\gamma)$ [D1.6, D1.7], the on-shell quantum $\kappa_e\in\mathbb Z$ [FT.22], the winding $k$ and $\Theta_{\mathrm{int}}$ [FT.25(3), FT.34, FT.43], $\delta h$ [D0.6], the spin ladder $j=|k|/2$ on $\mathrm{Sym}^k$ carriers [CT0.7a, EW.32, FT.52, FT.55, FT.60, Def-H.38]. \textbf{(R3) singlet admissibility} — the sealing indicator $\mathcal C(\Sigma;C)\in\{0,1\}$ and its idempotent $\Pi_{\mathcal C}$ [Def-H.9, Def-H.10, FT.19, CT3.18, CT3.22], with the partition rule ($|\mathrm{Colour}(\Sigma)|$ into $2$s and $3$s) [FT.34(5), CT3.17]. Observability is decided per aspect by FT.20's criteria table and FT.26's mask ($\dim\mathcal K_\Sigma\le18$), not by block membership [CT0.12/CT0.13/CT0.14]. \textbf{Census statuses.} An SM entry is \textbf{PRESENT} iff a cited tag supplies a read for it; \textbf{PARTIAL} iff a read exists for part of its defining data (e.g. the representation but not the chirality); \textbf{ABSENT} iff no read exists — and every ABSENT is either a \textit{machine-checked inventory fact} (a scan of the frozen Parts 1–2 text for the defining vocabulary) or a \textit{cited structural no-go}. The \textbf{completeness fraction} is $(\#\mathrm{PRESENT}+\tfrac12\#\mathrm{PARTIAL})/\#\mathrm{entries}$. \textbf{External-data clause.} Every SM datum used in this tier — field content, representation dimensions, hypercharge assignments, anomaly sums, parameter counts, the mass hierarchies of FT.77/FT.79 — is \textit{external measured or textbook data}, marked at each use, entering no derivation and no definition (A §0.2; P5: $\neg\exists$ quantity identified with a measured external constant; the comparison procedure is "fix $\mathfrak h$ once, read P2 $\wedge$ P3, the result is a test of A9"). No species is identified with a corpus object. \textbf{Rmk (citation correction).} The audit proposal attributed the read classes and the independent-quanta verdict to "FT.33/FT.34". In the current numbering those are the defect bridge and the round-trip quantization; the read classes are as listed above, and the independent-quanta audit is \textbf{FT.30} (with FT.31's continuum ledger), the pair-channel sector audit \textbf{FT.11/FT.12}, and the block decomposition \textbf{FT.30}'s $(G_5+W_3+H_1)+(J^+_3+J^-_3)+D^+_1$ tiling of the $16$ registers [CT0.12, CT1.1]. Recorded so that no future pass inherits the mis-citation.

\vspace{2pt}\noindent\rule{\textwidth}{0.4pt}\vspace{2pt}

\end{definition}

\begin{theorem}[{\texorpdfstring{FT.82 — Gauge-block reconstruction (audit pass B): the algebra closes, the coupling structure does not}{FT.82 - Gauge-block reconstruction (audit pass B): the algebra closes, the coupling structure does not}}]\label{tag:FT-82}
\textbf{Thm.} Under Def-H.29: $(1)$ \textbf{The register arithmetic} (all derived, zero freedom [P2]): $\dim\mathfrak H(\mathbb C^4)=16=(7,3,3,3)$ [CT1.1]; the aspect–capacity match $3{+}3{+}5{+}3{+}1{+}1=16=C_{\mathrm{cap}}$ [CT0.12]; the pair-channel tiling $(G_5{+}W_3{+}H_1)+(J^+_3{+}J^-_3)+D^+_1=16$ [FT.30]; the mask flag $\dim\mathcal K_n=(6,16,17,25)$ strict [CT0.14]; the observable bound $\dim\mathcal K_\Sigma\le1{+}5{+}1{+}3{+}8=18$ [FT.26(3)]. $(2)$ \textbf{$\mathfrak{su}(2)$ closes exactly} on the $W(3)$ block: with $W_i=\sigma_i/2$ over $\mathbb Q(i)$, $[W_i,W_j]=i\varepsilon_{ijk}W_k$ (all $9$ pairs, machine), $\mathfrak T(W_i)=0$, and EW.32's uniform normalization $\mathfrak T(\mathcal G_A\mathcal G_B)=\tfrac12\delta_{AB}$ on $\{W_1,W_2,W_3,Y\}$ with $Y=\sigma_0/2$ (all $16$ pairs, machine). The carrier is $\mathcal H_{\mathsf v}\cong\mathbb C^2$ — a genuine doublet read [EW.32]. $(3)$ \textbf{$\mathfrak{su}(3)$ closes exactly} on the colour algebra: in the matrix-unit (root-space) basis $\{E_{12},E_{21},E_{13},E_{31},E_{23},E_{32},H_1,H_2\}$ over $\mathbb Q$ — all $8$ traceless, all $64$ commutators in the span, the Jacobi identity on all $8^3$ triples, $[H_1,H_2]=0$, the $A_2$ root system $\{(2,-1),(1,1),(-1,2)\}\cup(-)$ in the $(H_1,H_2)$ basis, and a nondegenerate Killing form (machine, exact rationals). $\dim=8=3^2{-}1$ [CT3.8], $N_{\mathrm{gluon}}=8$ with \textbf{no ninth mode} — the $\mathfrak u(3)$ centre is allocated to $F_{\mathrm{EM}}$ [CT3.13] — and $N_{\text{colour}}=3=\log_2C_{\mathrm{cap}}-1$ is \textit{derived from A2 $+$ CT0.9, not chosen} [CT3.20]. FT.58's hosting identity $(\lambda_1+i\lambda_2)/2=E_{12}$ re-verified over $\mathbb Q(i)$. $(4)$ \textbf{The electroweak identities hold as exact polynomial identities} in $\mathbb Q[g,g',s]/(s^2-g^2-g'^2)$ (machine): EW.20's mass matrix in $(W_1,W_2,W_3,Y)$ has $W_{1,2}$ decoupled and a neutral block of \textbf{rank one} — $\det=0$ exactly, $\operatorname{tr}=g^2+g'^2=:m_Z^2$ (times $\mu_\star^2/4$ [EW.33]), null vector $(g',g)$, massive vector $(g,-g')$ [EW.23]; hence \textbf{the photon is exactly massless as the null direction} and $A=\sin\theta_WW_3+\cos\theta_WY$, $Z=\cos\theta_WW_3-\sin\theta_WY$ [D3.10, EW.24] with $\cos\theta_W=g/s$; $m_W=m_Z\cos\theta_W$ [P2] and $e=gg'/s=g\sin\theta_W=g'\cos\theta_W$, $1/e^2=1/g^2+1/g'^2$ [EW.35] — all verified with denominators cleared. The seed satisfies $W_3\mathsf v=-\tfrac12\mathsf v$, $Y\mathsf v=+\tfrac12\mathsf v$ [EW.20] and $T_a\mathsf v=0$ for $a=1..8$ [EW.32] — \textbf{the vev is a weak doublet component and a colour singlet}, exactly the SM pattern. $(5)$ \textbf{The $\mathrm U(1)$ direction} lives in EW.8 (the abstract generator $Y$ with three realizations: slot, aspect, $\dots$), D3.10/EW.24 (the mixing), and CT0.14's $\dim\mathcal K^{(\mathrm{EM})}=1$; EM.20 fixes the \textit{phase coefficient} $\alpha_{\mathrm{EM}}=\pi$ (a pure number under Comm-min $+$ Eff), which is \textbf{not} the fine-structure constant — the coupling is $e=g\sin\theta_W$ [EW.35], a P4-dependent quantity whose \textit{value} is not derived. (This is the standing trap recorded in FT.73(5): comparing $\pi$ to $1/137$ is a category error.) $(6)$ \textbf{Independent quanta: the SM's particle–boson picture is not reproduced.} FT.30's audit stands at block level: $G(5)$ is an independent quantum only as a \textit{joint geometric read} (winding $\kappa,k\in\mathbb Z$; unit $\mathfrak h^2=4G_N\ln2$ [CT2.9/P3]); $W(3)$ is a \textit{candidate} spin quantum with extraction closed under Def-H.38 [FT.55/FT.56/FT.60]; $H(1)$ is \textbf{not} an independent quantum (a continuous configuration datum, read only through the mass fold [Def-H.36, FT.20's $H$ row]); $J^\pm$ are structural and $\neg$observable [FT.20]; $D^+$ is uniform [D4.3]. So gauge bosons emerge as \textbf{aspect-block reads with per-aspect observability criteria} [FT.20], not as quanta of a gauged symmetry acting on chiral matter. $(7)$ \textbf{Two structural mismatches, derived not assumed.} (a) \textbf{No gauged chiral coupling}: the blocks act on the $16$-dimensional pair-channel register [CT0.6/CT0.12] and on slot algebras [D3.1, CT3.20]; nothing in Parts 1–2 couples a gauge block to a chiral fermion multiplet, because no chiral fermion multiplet exists [FT.81(5)]. The SM's defining feature — $\mathrm{SU}(2)$ acting on left-handed doublets only — has no read. (b) \textbf{The scalar block is $1$-dimensional}: $\dim F_H=1$ [CT0.12, derived: all $16$ directions allocated] against the SM Higgs doublet's $4$ real dof; the \textit{counting} of the Higgs mechanism is nonetheless right ($\operatorname{rk}M^2=3$: three massive vectors, one massless), because the breaking is carried by $\mu_\star^2\propto h_0^2$ [EW.33] rather than by four scalar fields. $(8)$ \textbf{The colour ceiling is the SM's own largest factor.} Slot algebras are realized only for $n\le3$ [CT3.20/CT4.4, P2's $n_{\max}=3$], so $\mathfrak{su}(3)$ is simultaneously the hosting ceiling and the realization ceiling [FT.58/FT.59]; at $n\ge4$, $D^-$ carries \textbf{no gauge structure at all} — no nonabelian and no abelian algebra, no generator contributed [CT4.4, derived from CT0.12's full allocation]. Hence the corpus cannot host a gauge factor larger than $\mathfrak{su}(3)$, which is exactly the SM's largest — a structural coincidence worth recording, and a hard bound against any GUT-scale reading.

\textbf{Req.} Def-H.29; A2, A8; C7, C8; D2.3, D3.1, D3.2, D3.10, D4.1, D4.3; CT0.5–CT0.14 (Part 1); CT1.1, CT1.2, CT1.10–CT1.12, CT2.9 (Part 2); CT3.2, CT3.7–CT3.9, CT3.13, CT3.17–CT3.22 (Part 2); CT4.1, CT4.4 (Part 2); EW.8, EW.13, EW.20, EW.22–EW.24, EW.27, EW.32, EW.33, EW.35 (Part 2); EM.9, EM.17–EM.20 (Part 2); FT.11, FT.12, FT.19, FT.20, FT.26, FT.27, FT.30, FT.31, FT.52–FT.60; P2, P3, P4; Pauli/Gell-Mann basis facts, root-space and Killing-form criteria, rank-one matrix facts [math; machine].

\textbf{Pf.}

%dsp-level:0
\[
(1):\ \text{each sum asserted exactly}\ [\text{machine}];\qquad 7+3+3+3=16=3+3+5+3+1+1=(5+3+1)+(3+3)+1
\]

%dsp-level:0
\[
(2):\ \sigma_i\sigma_j=\delta_{ij}+i\varepsilon_{ijk}\sigma_k\implies[\tfrac{\sigma_i}2,\tfrac{\sigma_j}2]=i\varepsilon_{ijk}\tfrac{\sigma_k}2;
\]

%dsp-level:0
\[
\qquad \mathfrak T(\tfrac{\sigma_A}2\tfrac{\sigma_B}2)=\tfrac14\cdot2\delta_{AB}=\tfrac12\delta_{AB}
\]

%dsp-level:0
\[
(3):\ [E_{ij},E_{kl}]=\delta_{jk}E_{il}-\delta_{li}E_{kj};
\]

%dsp-level:0
\[
\quad [H_1,H_2]=0;
\]

%dsp-level:0
\[
\quad [H_i,E_\alpha]=\alpha(H_i)E_\alpha;\quad\text{span closure by exact coordinate solve}
\]

%dsp-level:0
\[
\text{Killing}:\ K_{ab}=\textstyle\sum_{c}\operatorname{tr}(\operatorname{ad}_{E_c}\operatorname{ad}_{E_{\bar c}})\ \text{form};\ \det\ne0\implies\text{semisimple}\ [\text{machine, exact Fractions}]
\]

%dsp-level:0
\[
(\lambda_1+i\lambda_2)/2=(E_{12}{+}E_{21}+i(-iE_{12}{+}iE_{21}))/2=E_{12}\ \checkmark\ [\mathrm{FT.58}]
\]

%dsp-level:0
\[
(4):\ M^2_{\mathrm{neutral}}=\tfrac{\mu_\star^2}4(g,-g')^{\mathsf T}(g,-g'):\ \det=0,\ \operatorname{tr}=g^2+g'^2;
\]

%dsp-level:0
\[
\quad M^2(g',g)^{\mathsf T}=0;
\]

%dsp-level:0
\[
\quad M^2(g,-g')^{\mathsf T}=(g^2+g'^2)(g,-g')^{\mathsf T}
\]

%dsp-level:0
\[
m_W^2s^2=g^2(g^2+g'^2)=m_Z^2g^2\implies m_W=m_Z\cos\theta_W;\qquad e^2s^2=g^2g'^2\iff\tfrac1{e^2}=\tfrac1{g^2}+\tfrac1{g'^2}
\]

%dsp-level:0
\[
W_3\mathsf v=-\tfrac12\mathsf v,\ Y\mathsf v=+\tfrac12\mathsf v:\ \text{direct}\ 2\times2\ \text{multiplication over}\ \mathbb Q(i);\qquad T_a\mathsf v=0\ [\mathrm{EW.32}]
\]

%dsp-level:0
\[
(5)\text{–}(8):\ \text{each row cites its tag};\qquad(7a):\ \text{inventory scan}\ [\mathrm{FT.81}(5)];\qquad(7b):\ \mathrm{CT0.12}\ \text{allocation}
\]

%dsp-level:0
\[
(8):\ \mathrm{CT4.4}\ \text{Pf}:\ (a)\ 16\ \text{pair directions fully allocated}\implies\neg\exists\ \text{direction for}\ D^-;
\]

%dsp-level:0
\[
\ (b)\ \text{slot algebras only}\ n
\]

%dsp-level:0
\[
\le 3.\ \blacksquare
\]

\textbf{Phys.} Ask whether the sixteen registers rebuild the Standard Model's gauge group, and the arithmetic says yes in a way that is not a naming trick: an eight-dimensional colour algebra that closes with the right root system and no ninth gluon, a three-dimensional weak algebra whose doublet carries the vacuum exactly where the Standard Model puts it, and a neutral block of rank one whose null direction is the photon — with $\sin\theta_W$, $\cos\theta_W$ and $e$ falling out as polynomial identities rather than fits. Ask instead whether those blocks \textit{act on matter} the way a gauge group does, and the answer is no: there is no chiral multiplet for them to act on, the scalar block is one-dimensional instead of four, and the quanta are aspect reads with observability criteria, not gauge bosons coupled to fermions. The gauge skeleton is derived; the coupling to chiral matter is absent.

\vspace{2pt}\noindent\rule{\textwidth}{0.4pt}\vspace{2pt}

\end{theorem}

\phantomsection
\section[Tier F.29 — Gravitational Reads in Three States, the…]{\texorpdfstring{Tier F.29 — Gravitational Reads in Three States, the Flicker Ripple, and the Tumbling-Model Adjudication}{Tier F.29 - Gravitational Reads in Three States, the Flicker Ripple, and the Tumbling-Model Adjudication}}

\begin{definition}[{\texorpdfstring{Def-H.30 — The three shape states}{Def-H.30 - The three shape states}}]\label{tag:Def-H-30}
\textbf{Def.} For a skeleton $\Gamma$ [Def-H.25] with potential degrees $n(v)$: \textbf{(S0) VV / all-virtual.} $R=\varnothing$: every potential edge virtual, $n_{\mathcal R}=0$, every pair at $\rho_\circ=\tfrac14 I$ [A5, CT0.2]. The shape-derived vacuum. \textbf{(S1) flicker / bypass.} $R\ne\varnothing$ admissible: $k_v\le3$ at every vertex [CT0.11], i.e. the bypass sector of Def-H.25, with its ensemble reads at $p$ [Def-H.26, Cal-F] and its cover sub-sector [FT.67(2)]. \textbf{(S2) $D^-$ trapped.} Some vertex locked ($k_v\ge4$), up to the fully crystallized floor $R=E_{\mathrm{dir}}(\Gamma)$: content absorbed at rate $\propto O(n)$ [CT4.1/D4.2], retained in $\mathfrak L$ and \textbf{not accessible} [D1.11/CT4.3], with no accessible pole [CT4.2]. \textbf{Rmk (status).} The three states are the proposal's Punnett vocabulary (VV / RV,VR / RR) mapped onto corpus objects; the mapping is recorded, not adopted as a definition of anything physical. S1 and S2 are Def-H.25's own sectors; S0 is CT0.2's vacuum restricted to a skeleton's potential-edge set.

\vspace{2pt}\noindent\rule{\textwidth}{0.4pt}\vspace{2pt}

\end{definition}

\begin{theorem}[{\texorpdfstring{FT.85 — Vacuum census, baseline equality, and the constructible gravitational reads}{FT.85 - Vacuum census, baseline equality, and the constructible gravitational reads}}]\label{tag:FT-85}
\textbf{Thm.} Under Def-H.30: $(1)$ \textbf{Vacuum census.} The corpus declares exactly three scalar extremes and proves there is no fourth: the \textbf{Janus vacuum} $n_{\mathcal R}=0$ with $(S_{\mathrm{raw}},S_{\mathfrak L})=(\ln4\ \text{per pair (max, argmax unique}=\tfrac14I),0)$ [CT0.2, CT4.7(1)]; \textbf{$D^+$} with $(0,0)$ — a \textit{pure}, uniform, non-configuration-dependent scalar [CT4.7(2), D4.3], $w=-1$ [CT1.25]; \textbf{$D^-$ locked} with $(\min_{\mathrm{acc}},\max\wedge\neg\text{accessible})$ [CT4.7(3), D1.11, CT4.2, CT4.3]; and $\neg\exists$ a fourth extreme [CT4.10], $F_H$ being a dynamical field [EW.27]. $(2)$ \textbf{Baseline equality (proven, not measured).} For every shape $S$: $S_{\mathfrak L}(\mathrm{VV}(S))=n_{\mathcal R}\ln2=0=S_{\mathfrak L}(V_{\mathrm{raw}})$ [CT1.29]; $\Delta S_{\mathrm{raw}}=0$ on both sides [CT0.2: every pair at the argmax]; $m_{\mathrm{hid}}=S_{\mathfrak L}/\mathfrak h=0$ [FT.44(4)]; $\rho_{D^-}=0$ [CT4.5's iff]; and CT1.29's $S_{\partial\mathcal B}$ has no argument (no closed real-edge surface). Hence

%dsp-level:0
\[
m_g^2\bigl(\mathrm{VV}(S)\bigr)=m_g^2\bigl(V_{\mathrm{raw}}\bigr)\quad\forall S,
\]

on every constructible read, and $\Delta S_{\text{vacuum polarization}}(S)\equiv0$. The reason is registration, not arithmetic: virtual pairs carry \textbf{no registration} [A7/D0.5: $\phi_e$ is defined on $e\in E_{\mathcal R}$; A5: the sole ledger exclusion is the Haar draw; Def-H.25(S4): virtual edges weight $1$, no phase]. Honest remark: the \textit{absolute} raw entropy is $|\mathrm{pairs}|\cdot\ln4$ and is shape-dependent; the \textit{deficit} is shape-independent and zero. The raw vacuum's mass reads vanish from three independent derived reads — nothing is set to zero by hand. $(3)$ \textbf{$m_g^2:=\Delta S$ is not constructible.} (a) The inventory contains no inertial/gravitational mass distinction and no equivalence theorem (\texttt{inertial}, \texttt{Lyapunov}, \texttt{maximal\ mixing}, \texttt{equivalence\ principle}, \texttt{Theorem\ 11}: count $0$ in Parts 1–2); (b) $\Delta S_{\mathrm{raw}}$ of a \textit{registered} configuration is bounded but undetermined — A4/D0.4 force a real pair strictly below $\ln4$ and no tag declares its $\rho$, so $\Delta S_{\mathrm{raw}}\in(0,n_{\mathcal R}\ln4]$ with no value; (c) the dimensionally correct entropy-based mass$^2$ \textit{density} is $S/A=\ln2/\mathfrak h^2$ [CT1.29/CT2.9], a \textbf{state-independent constant} [FT.37: corrected by $1/\bar\kappa$], so state-dependence can enter only through $A$ or $\bar\kappa$. $(4)$ \textbf{The constructible reads, used for every table below.} (i) $m_{\mathrm{hid}}=S_{\mathfrak L}/\mathfrak h=E_{\mathrm{int}}\ln2/\mathfrak h$ — the universal taxonomy law [FT.44(4)] — with $m_{\mathrm{hid}}^2=E_{\mathrm{int}}^2(\ln2)^2/\mathfrak h^2$, dimension $\mathfrak q$ [D2.9/CT2.1]; (ii) the locked sector's stress-energy $T^{(D^-)}_{\mu\nu}=\rho_{D^-}u_\mu u_\nu$, $p=0$, $w=0$, $\rho_{D^-}=\sum_{v\ \mathrm{locked}}O(n_v)\ge0$ [CT4.1/CT4.5] — the corpus's one derived gravitational source; (iii) $S_{\partial\mathcal B}=A\ln2/\mathfrak h^2=A/(4G_N)$ with $A=n_{\mathcal R}(\partial\mathcal B)\mathfrak h^2$ and $G_N=\mathfrak h^2/(4\ln2)$, $\kappa^2=8\pi G_N$ \textbf{derived} [CT1.29/CT2.9/P3], bounded by CT4.9 ($\le3$ bits per accessible vertex; $S_{\mathrm{acc}}(\mathcal B)\le A/\mathfrak h^2$ bits); (iv) the pole side $m^2_{\mathrm{pole}}=4\hat\mu'^2\sin^2(\Theta_R/4)$ under Def-H.37(2) [FT.54/FT.76] — which is also what a two-patch field-source read reduces to on the cone metric [Def-H.21/FT.47], i.e. clause (iii) of the audit's B3 is already bounded [FT.76(3)]. $(5)$ \textbf{The source criterion, exact.} $\rho_{D^-}>0\iff\exists v:\ k_v\ge4$ [CT4.5's iff] — machine-exhaustive over the tetrahedron's $4096$ and $P_4$'s $256$ patterns. Consequences: \textbf{every S1 configuration has exactly zero dark-dust stress-energy}, and only S2 states are dust sources; those are retained-but-inaccessible [CT4.2/CT4.3/D1.11], amplitude-zero in the literal domain [FT.67(3a)] and revived only in the chain-local domains [FT.75(2)]. So the flicker gas does not gravitate through $D^-$ at all; its only mass read is informational. $(6)$ \textbf{The state table} ($14$ shapes; $\langle K\rangle$ at $p=\tfrac12$ reproduces FT.70(5)/FT.68(1) verbatim):

%dsp-level:0
\[
\begin{array}{c|ccc|ccc}
& \multicolumn3{c|}{m_{\mathrm{hid}}\,\mathfrak h/\ln2} & \multicolumn3{c}{\rho_{D^-}}\\
\text{shape} & \text{S0} & \text{S1}\ \langle K\rangle(\tfrac12) & \text{S2}\ E_{\mathrm{int}} & \text{S0} & \text{S1} & \text{S2}\\
\hline
\text{digon} & 0 & 1 & 2 & 0 & 0 & 0\ (n(v)=2\ \text{never locks})\\
P_3..P_{10} & 0 & \tfrac83,\tfrac{248}{69},\tfrac{130}{29},\tfrac{296}{55},\tfrac{2618}{417},\tfrac{11344}{1581},\tfrac{896}{111},\tfrac{13588}{1515} & 2n & 0 & 0 & 16n\\
\text{tetra} & 0 & \tfrac{852}{205} & 12 & 0 & 0 & 448\\
\text{cube} & 0 & \tfrac{3311952}{395023} & 24 & 0 & 0 & 896\\
\text{octa} & 0 & \tfrac{2124600}{318523} & 24 & 0 & 0 & 2976\\
\text{dodeca} & 0 & \tfrac{4007491093935690}{191727103387307} & 60 & 0 & 0 & 2240\\
\text{icosa} & 0 & \tfrac{215801166207}{15519497071} & 60 & 0 & 0 & 24384
\end{array}
\]

with $S_{\mathfrak L}=(\text{that row})\cdot\ln2$, $A/\mathfrak h^2=E_{\mathrm{int}}$ in S2, and per-sector entropies \textbf{not available} (clause (7)). $(7)$ \textbf{No per-sector $\Delta S_A$.} Four cited reasons: no declared state $\Psi_X$ for a crystallization pattern, so $\|F_A\Psi_X\|^2$ has no argument; FT.30's audit leaves only the $G$ block with a geometric read and marks it \textbf{joint} ($\neg$ per-block); $H$, $J^\pm$, $D^+$ are explicitly not independent quanta [FT.30]; CT0.14/FT.26 supply block \textit{dimensions} $(6,16,17,25)$ and the $\le18$ observable bound, which is all the sector map can carry. The only declared per-register decomposition is $S_{\mathfrak L}=K\ln2=\sum_{e\in R}\ln2$ [CT0.9: one real edge $=$ one bit] — machine-verified per configuration.

\textbf{Req.} Def-H.25, Def-H.26, Def-H.30; FT.16, FT.30, FT.37, FT.44, FT.47, FT.54, FT.63, FT.67, FT.70, FT.71, FT.75, FT.76; A4, A5, A7; D0.5, D1.11, D4.2, D4.3; CT0.2, CT0.9, CT0.11, CT0.12, CT0.14 (Part 1); CT1.22, CT1.25, CT1.29, CT2.1, CT2.8, CT2.9 (Part 2); CT4.1–CT4.5, CT4.7, CT4.9, CT4.10 (Part 2); P1, P3, P5; Gibbs lemma [math]; exhaustive pattern enumeration [machine].

\textbf{Pf.}

%dsp-level:0
\[
(1):\ \text{CT4.7's three extremes with CT4.10's exhaustiveness};
\]

%dsp-level:0
\[
\qquad (2):\ n_{\mathcal R}=0\implies S_{\mathfrak L}=0\ [\mathrm{CT1.29}];\ \rho_{pq}=\tfrac14I\ \forall(p,q)\ [\mathrm{A5},\mathrm{CT0.2}]
\]

%dsp-level:0
\[
\text{registration}:\ \phi_e\ \text{defined on}\ e\in E_{\mathcal R}\ [\mathrm{A7},\mathrm{D0.5}];\quad\text{virtual weight}\ 1,\ \text{no phase}\ [\mathrm{Def\text{-}H.25(S4)}]
\]

%dsp-level:0
\[
(3a):\ \text{inventory scan};
\]

%dsp-level:0
\[
\qquad (3b):\ \mathrm{A4}/\mathrm{D0.4}:\ \text{real}\iff|\Pi_{\mathcal V}\rho|>\varepsilon_q:\ \text{strict sub-maximality, no declared}\ \rho
\]

%dsp-level:0
\[
(3c):\ S/A=\ln2/\mathfrak h^2\ [\mathrm{CT1.29}+\mathrm{CT2.9}:\ A/(4G_N)=A\ln2/\mathfrak h^2];
\]

%dsp-level:0
\[
\qquad \text{FT.37}:\ S_{\mathrm{reg}}/A_{\mathrm{loop}}=(\ln2/\mathfrak h^2)/\bar\kappa
\]

%dsp-level:0
\[
(4):\ [\mathrm{FT.44}(4)]:\ m_{\mathrm{hid}}=E_{\mathrm{int}}\ln2/\mathfrak h\ \text{on every computed cell};\qquad[\mathrm{CT4.5}]:\ T^{(D^-)}=\rho u u,\ p=0
\]

%dsp-level:0
\[
(5):\ \mathrm{CT4.5}\ \text{iff}:\ \rho_{D^-}=0\iff\neg\exists v:\ n(v)\ge4;\qquad\text{S1 admissible}\implies k_v\le3\ \forall v\ [\mathrm{CT0.11}]\implies\rho_{D^-}=0
\]

%dsp-level:0
\[
(6):\ \langle K\rangle_{\mathrm{byp}}(\tfrac12)=\frac{\sum_K Kc_K2^{-2E}}{\sum_Kc_K2^{-2E}}\ [\text{cap census, exact}];
\]

%dsp-level:0
\[
\qquad \rho_{D^-}(\text{floor})=V\cdot O(2d)\ [\mathrm{FT.63}(5)]
\]

%dsp-level:0
\[
(7):\ \text{the four cited obstructions};\qquad S_{\mathfrak L}=\sum_{e\in R}\ln2\ [\mathrm{CT0.9}].\ \blacksquare
\]

\textbf{Phys.} Ask the corpus for the weight of empty space and it answers with a proof rather than a parameter: the vacuum is a state of maximal raw information with no ledger content at all, and a shape whose edges are all still virtual is that same vacuum — not a slightly different one. Mass enters only where edges crystallize, one bit each, and gravity's own coupling is fixed by the single scale. The trapped sector is the only place with a derived stress-energy: pressureless dust that curves geometry while remaining, by theorem, out of reach. The flickering gas between the two carries information and amplitude but sources no dust whatsoever.

\vspace{2pt}\noindent\rule{\textwidth}{0.4pt}\vspace{2pt}

\end{theorem}

\begin{theorem}[{\texorpdfstring{FT.86 — The flicker ripple: ultralocal raw dynamics, cap-induced alternation, and the quantum–classical gradient}{FT.86 - The flicker ripple: ultralocal raw dynamics, cap-induced alternation, and the quantum-classical gradient}}]\label{tag:FT-86}
\textbf{Thm.} Under Def-H.26/Def-H.30 at $p=\tfrac12$ (every value exact; general $p$ by the same census): $(1)$ \textbf{The raw flicker has no spatial ripple.} $\mu_p$ is a Bernoulli \textbf{product} measure on directed edges [Def-H.26], so every connected edge–edge correlation vanishes identically: $\langle x_ex_{e'}\rangle=p^2=\langle x_e\rangle\langle x_{e'}\rangle$ for $e\ne e'$ (machine-exhaustive on $P_3$, $P_4$). Any waviness in the flicker geometry is \textbf{constraint-induced}, not dynamical. $(2)$ \textbf{The admissible flicker carries an exact period-2 alternating ripple.} Conditioning on admissibility ($k_v\le3$, CT0.11) couples neighbours; the connected correlations, by \textit{true ring distance} (the two directions of one undirected edge are distance $0$ — a digon — not distance $1$), are exact rationals alternating in sign:

%dsp-level:0
\[
\begin{array}{c|cccccc}
& d=0 & d=1 & d=2 & d=3 & d=4 & d=5\\
\hline
P_4 & -\tfrac{133}{4761} & -\tfrac{64}{4761} & +\tfrac{28}{4761} & & & \\
P_8 & -\tfrac{71068}{2499561} & -\tfrac{33124}{2499561} & +\tfrac{6928}{2499561} & -\tfrac{1504}{2499561} & +\tfrac{604}{2499561} & \\
P_{10} & -\tfrac{1631509}{57380625} & -\tfrac{760384}{57380625} & +\tfrac{158716}{57380625} & -\tfrac{33184}{57380625} & +\tfrac{7216}{57380625} & -\tfrac{2884}{57380625}
\end{array}
\]

Short-distance values converge to $L$-independent limits ($d{=}0\to-0.0284331$, $d{=}1\to-0.0132516$): the cap makes neighbouring crystallizations anticorrelate and next-nearest correlate. $(3)$ \textbf{The decay rate is the corpus's own transfer matrix.} With $\operatorname{charpoly}T(1)=x^2-3x-3$ [FT.68(2)/FT.70(3)], $\lambda_{1,2}=(3\pm\sqrt{21})/2$ and

%dsp-level:0
\[
\frac{\lambda_2}{\lambda_1}=\frac{\sqrt{21}-5}{2}=-0.2087122\ldots\in\mathbb Q(\sqrt{21}),
\]

%dsp-level:0
\[
\qquad \xi=-1/\ln|\lambda_2/\lambda_1|=0.6382\ \text{directed edges},
\]

matching the measured ratios $C(2)/C(1)=-\tfrac{39679}{190096}=-0.2087314$ ($P_{10}$, \textbf{4 significant figures}) and $C(3)/C(2)=-\tfrac{8296}{39679}=-0.2090779$ (3 s.f.); $P_9$ gives $-\tfrac{2759}{13225}=-0.208620$. The \textbf{negative sign of $\lambda_2/\lambda_1$ is the alternation}. Deviations grow near the antipode ($d\approx L/2$) as finite-size wrap. $(4)$ \textbf{Exact moments from the cap census} (never a Bernoulli proxy): $\langle K\rangle$, $\langle K^2\rangle$, $\sigma_K^2$ for all $14$ shapes, with $\sigma_K^2/\sigma_{K,\mathrm{ind}}^2$ (the \textbf{cap factor}, $\sigma^2_{\mathrm{ind}}=2Ep(1-p)=E/2$) equal to $1$ for the digon (the cap never binds at $n(v)=2$), $\to0.7002$ for large polygons, and $0.155$–$0.373$ for the solids (icosa $0.1547$, octa $0.2228$). Anchors: digon $\langle K\rangle=1,\langle K^2\rangle=\tfrac32,\sigma^2=\tfrac12$; $P_3$: $\tfrac83,\tfrac{73}9,1$. $(5)$ \textbf{The ripple in the mass reads.} $\Delta S=S_{\mathfrak L}=K\ln2\Rightarrow\delta(\Delta S)=\sigma_K\ln2$ exactly; $\delta(m_{\mathrm{hid}})/m_{\mathrm{hid}}=\sigma_K/\langle K\rangle$ and $\delta(m_{\mathrm{hid}}^2)/m_{\mathrm{hid}}^2=2\sigma_K/\langle K\rangle$ (the factor $2$ exact). The raw-vacuum and VV baselines give \textbf{identical} $\delta(m_g^2)$ — by FT.85(2)'s proof, not by numerical agreement. $(6)$ \textbf{Quantum–classical gradient.} $\sigma_K/\langle K\rangle$ decreases monotonically from the digon's $1/\sqrt2=0.707107$ (exact) through $P_3$'s $3/8=0.375$ (exact) to the icosahedron's $0.10955$; the $1/\sqrt{\langle K\rangle}$ law holds with a \textbf{family-dependent coefficient} $\sigma_K/\sqrt{\langle K\rangle}$: $1/\sqrt2$ (digon), $0.6124$ ($P_3$), $0.6248$ converged ($P_4$–$P_{10}$), $0.41$–$0.52$ (solids) — the deviation from the binomial $1/\sqrt2$ being exactly the cap factor. \textbf{Near-pure-wave limit:} $\sigma_K/\langle K\rangle=1$ requires a \textit{single} directed edge at $p=\tfrac12$ — amplitude-zero [FT.67(3b): $\operatorname{cov}=0$], hence not a bound shape. No shape in the inventory is near-pure-wave. $(7)$ \textbf{No frequency is constructible; one period is derived.} the Cal-F rate datum is the ratio $r=\Gamma_{J^+}/\Gamma_{J^-}$ and per-edge detailed balance, and records that the corpus supplies no $J^\pm$ \textit{rates}, so $f_{\mathrm{ripple}}=\Gamma_{\mathrm{tot}}/2\pi$, $\tau_c=1/\Gamma_{\mathrm{tot}}$ are reported \textbf{in units of the undeclared total rate} and no absolute frequency is claimed. Rate-free and derived: one $J^\pm$ event changes $(n_{\mathcal R},\delta h,\delta\tau)$ by $(\pm1,\pm1,\pm1)$ [CT0.3], so \textbf{in ledger-tick units $\tau_c=1$ tick} and the Lorentzian half-width is $1$ rad/tick; propagation obeys $\delta\tau\ge\delta h/4$ (at most $4$ hops/tick) [FT.1/FT.38, CT0.6, EM.5–7]. Leading order (single-flip granularity) the cap correlation renormalizes the Lorentzian \textbf{amplitude} by the cap factor and leaves the \textbf{width} at $\Gamma_{\mathrm{tot}}$; the full constrained spectrum is multi-timescale and not constructible without rates. $(8)$ \textbf{Self-modulation is not constructible.} A tick-rate/field coupling $r_q(g_q)$ has no referent: \texttt{clock}, \texttt{dilation}, \texttt{tick\ rate}, \texttt{r\_\allowbreak{}q}, \texttt{g\_\allowbreak{}q}, \texttt{local\ rate}, \texttt{running\ coupling}, \texttt{beta\ function} all count $0$ in Parts 1–2. The corpus has the conformal factor $\Omega=e^{\alpha_\Omega a_{00}}$ and the shift $N_i=\beta w_i$ [CT1.12], and $\delta\tau$ is defined as retained ledger updates along real edges [D0.7], but nothing couples the two. Logged as a gap (a new interface-definition species), not adopted. $(9)$ \textbf{The trapped sector does not ripple.} CT4.5's own Pf: locked $\Rightarrow$ no accessible pole [CT4.2] $\Rightarrow$ no propagating mode $\Rightarrow\mathsf E_{\mathrm{tr}}=0\Rightarrow T^{\mathrm{tr}}=0$ — S2 states are \textbf{non-oscillatory by theorem} (pure dust, $p=0$, $w=0$).

\textbf{Req.} Def-H.26, Def-H.30; FT.1, FT.38, FT.67, FT.68, FT.70, FT.85; D0.7, D1.2; CT0.3, CT0.6, CT0.11 (Part 1); CT1.22, CT1.29 (Part 2); CT4.2, CT4.5 (Part 2); EM.5–EM.7 (Part 2); CT1.12; Cal-F; transfer-matrix spectra, exhaustive enumeration [math; machine].

\textbf{Pf.}

%dsp-level:0
\[
(1):\ \mu_p=\bigotimes_e\mathrm{Bernoulli}(p)\ [\mathrm{Def\text{-}H.26}]\implies\langle x_ex_{e'}\rangle=p^2\ \forall e\ne e'\ [\text{machine-exhaustive}]
\]

%dsp-level:0
\[
(2):\ \text{conditioning on}\ \{k_v\le3\}\ [\mathrm{CT0.11}]:\ \langle\cdot\rangle_{\mathrm{byp}}
\]

%dsp-level:0
\[
= \langle\cdot\,
\]

%dsp-level:0
\[
\mathbf1_{\mathrm{adm}}\rangle/B_\Gamma(p):\ \text{exact rationals by enumeration}
\]

%dsp-level:0
\[
(3):\ \det(xI-T(1))
\]

%dsp-level:0
\[
= x^2(x^2-3x-3)\ [\mathrm{FT.68}(2)];
\]

%dsp-level:0
\[
\quad\lambda_2/\lambda_1
\]

%dsp-level:0
\[
= (3-\sqrt{21})/(3+\sqrt{21})
\]

%dsp-level:0
\[
= (\sqrt{21}-5)/2\ \text{in}\ \mathbb Q(\sqrt{21})
\]

%dsp-level:0
\[
(4):\ \text{cap-DP census rows}\ [\mathrm{FT.68}(3),\mathrm{FT.69}(1)];\qquad\sigma^2=\langle K^2\rangle-\langle K\rangle^2\ \text{exact}
\]

%dsp-level:0
\[
(5):\ \Delta S=K\ln2\ [\mathrm{CT1.29}]\ \text{linear in}\ K;\qquad m_{\mathrm{hid}}^2\propto K^2\implies\text{relative ripple}\times2
\]

%dsp-level:0
\[
(6):\ \text{binomial reference}\ \sigma/\sqrt{\langle K\rangle}=\sqrt{1-p}=1/\sqrt2;\qquad\text{cap factor}=\sigma^2/\sigma^2_{\mathrm{ind}}
\]

%dsp-level:0
\[
(7):\ \mathrm{CT0.3}:\ J^\pm=(\pm1,\pm1,\pm1);\qquad\mathrm{FT.1}/\mathrm{FT.38}:\ I\le\min(\Delta\ell/\mathfrak h,4\delta\tau)\ln2
\]

%dsp-level:0
\[
(8):\ \text{inventory scan};\qquad(9):\ \mathrm{CT4.5}\ \text{Pf verbatim}.\ \blacksquare
\]

\textbf{Phys.} A flickering shape does not shimmer the way intuition suggests. Edge by edge its raw dynamics is completely uncorrelated — there is no wave in the dynamics at all. The waviness is created by the \textit{budget}: because no vertex may hold more than three crystallized edges, a real edge here forbids one there, and the forbiddenness propagates around the ring with an exact alternation, negative–positive–negative, dying away by a factor of one fifth per step. That single number is the same eigenvalue ratio the corpus derived for its polygon transfer matrix years earlier, now showing up as a physical ripple. Smaller shapes ripple harder, and the ripple fades as the shape grows — a quantum-to-classical gradient with no parameter in it. What no shape achieves is pure wave: that would take a single edge, and a single edge is nothing at all.

\vspace{2pt}\noindent\rule{\textwidth}{0.4pt}\vspace{2pt}

\end{theorem}

\medskip\noindent\textit{\texorpdfstring{FT.87 — The tumbling-model adjudication (external proposal, row by row)}{FT.87 - The tumbling-model adjudication (external proposal, row by row)}.}\ Structural-constraints record: see Appendix B, Section B.2.4 of the Addendum.\par\medskip

\begin{theorem}[{\texorpdfstring{FT.88 — The involution algebra of the two declared reversals}{FT.88 - The involution algebra of the two declared reversals}}]\label{tag:FT-88}
\textbf{Thm.} Let $\omega$ [D2.6] and $\Theta_{\mathfrak L}$ [D2.7] be the corpus's two declared involutions. $(1)$ \textbf{$\omega$ on patterns.} $\omega$ reverses orientation: $s(c)\leftrightarrow t(c)$, charge $\mapsto-$charge, chirality $\mapsto-$chirality [D2.6]. On a skeleton it acts as $R\mapsto R^{\mathrm T}=\{(b,a):(a,b)\in R\}$, with $\omega^2=1$ (machine: on all patterns sampled and on the tetrahedron's full orbit set), and $\omega\notin\operatorname{Aut}(\Gamma)$ — no vertex permutation induces it [FT.78(6)], where it was shown to exchange the two $\Omega=24$ dissolved-digon cover classes. $(2)$ \textbf{$\Theta_{\mathfrak L}$ on registrations.} $\Theta_{\mathfrak L}$ reverses ledger order with $J^+\leftrightarrow J^-$ and $G(e)\mapsto G(e)^\ast$ [D2.7], i.e. $\sigma_e$ fixed and $\kappa_e\mapsto-\kappa_e$ on the edge weight $G(e)=\sigma_e e^{i\pi\kappa_e}$ [D3.5]; $\mathcal J(C)\mapsto\mathcal J(C)^\ast$ [D3.6]. $(3)$ \textbf{The algebra.} $\omega$ acts on the \textit{edge set} and $\Theta_{\mathfrak L}$ on the \textit{phase/ledger order}: disjoint data, hence they commute, and

%dsp-level:0
\[
\langle\omega,\Theta_{\mathfrak L}\rangle\cong\mathbb Z_2\times\mathbb Z_2
\]

on skeleton registration data, with the four classes $(1,\ \omega,\ \Theta_{\mathfrak L},\ \omega\Theta_{\mathfrak L})$. $\omega$ is charge-odd [D2.6; FT.40's $\omega$-covariance of $q_{\mathrm{core}}$]; $\Theta_{\mathfrak L}$ is charge-even and phase-odd. $(4)$ \textbf{Scope, and what this is not.} The inventory contains no C, P or T \textit{operator}, no Lorentzian field algebra and no spinor bundle, so no CPT theorem is claimed or claimable; $\mathbb Z_2\times\mathbb Z_2$ is the exact group generated by the two declared reversals and nothing more. In particular $\omega$ is \textbf{not} a chirality projector: it is an orientation reversal whose fixed-point set is empty on chiral pairs [FT.78(6)], whereas a projector needs an eigenspace decomposition of a \textit{state} space that the corpus does not supply for patterns (application-layer content). The proposal's $\chi=\operatorname{sign}(k)$ is the $\omega$-odd coordinate of $(3)$: a legitimate C-odd ledger read, still not a projector.

\textbf{Req.} D2.6, D2.7, D3.5, D3.6; FT.40, FT.64, FT.78; CT0.9 (Part 1); group generation, exhaustive pattern enumeration [math; machine].

\textbf{Pf.}

%dsp-level:0
\[
(1):\ (R^{\mathrm T})^{\mathrm T}=R;
\]

%dsp-level:0
\[
\qquad \omega\in\operatorname{Aut}\iff\exists p:\ (p(a),p(b))=(b,a)\ \forall(a,b):\ p(a)=b\wedge p(b)=a\ \forall\ \text{pairs}:\ \bot\ \text{for}\ |V|\ge3
\]

%dsp-level:0
\[
\text{machine}: \#\{p\in\operatorname{Aut}:\ p\ \text{induces the reversal}\}=0\ \text{on the tetrahedron}
\]

%dsp-level:0
\[
(2):\ \mathrm{D2.7}\ \text{verbatim};\qquad G(e)^\ast=\sigma_e e^{-i\pi\kappa_e}\ [\sigma_e\in\{\pm1\}\ \text{real},\ \mathrm{CT0.9}]
\]

%dsp-level:0
\[
(3):\ \text{disjoint supports}\implies[\omega,\Theta_{\mathfrak L}]=0;\qquad\omega^2=\Theta_{\mathfrak L}^2=1\implies\mathbb Z_2\times\mathbb Z_2
\]

%dsp-level:0
\[
(4):\ \text{inventory scan (no C/P/T operator)};
\]

%dsp-level:0
\[
\qquad \text{projector needs}\ \Psi\ \text{eigenspaces}:\ \neg\exists\ \text{declared}\ \Psi_X\ [\mathrm{FT.85}(7)].\ \blacksquare
\]

\textbf{Phys.} The corpus has exactly two mirrors: one flips every arrow in the graph, the other runs the ledger backwards and conjugates every phase. They do not interfere — one acts on edges, the other on phases — so together they make the smallest possible non-cyclic group, four elements. That group already knows which cover shapes are left-handed and which are right-handed, and it already knows that charge changes sign in the first mirror. It is not the CPT of a field theory, and it is not a chirality projector; it is the exact amount of reversal structure the axioms declare, no more.

\vspace{2pt}\noindent\rule{\textwidth}{0.4pt}\vspace{2pt}

\end{theorem}

\phantomsection
\section[Tier F.30 — The Koide Adjudication, the Resolution Ladder,…]{\texorpdfstring{Tier F.30 — The Koide Adjudication, the Resolution Ladder, and the Minimal Bridge Declaration}{Tier F.30 - The Koide Adjudication, the Resolution Ladder, and the Minimal Bridge Declaration}}

\medskip\noindent\textit{\texorpdfstring{FT.89 — The Koide/phase-boundary adjudication: exact arithmetic, category trap}{FT.89 - The Koide/phase-boundary adjudication: exact arithmetic, category trap}.}\ Structural-constraints record: see Appendix B, Section B.1.13 of the Addendum.\par\medskip

\medskip\noindent\textit{\texorpdfstring{FT.90 — The resolution ladder: no-hair from the channel bound, and the chiral merger refuted}{FT.90 - The resolution ladder: no-hair from the channel bound, and the chiral merger refuted}.}\ Structural-constraints record: see Appendix B, Section B.2.5 of the Addendum.\par\medskip

\begin{theorem}[{\texorpdfstring{FT.91 — The minimal spectral-bridge form}{FT.91 - The minimal spectral-bridge form}}]\label{tag:FT-91}
\textbf{Thm.} Under Def-H.16/FT.80(4): $(1)$ \textbf{The three declarations restated.} Item 13's exit (b) requires exactly: \textbf{(i)} a rule $\mathcal O\mapsto(\Sigma(\mathcal O),C(\mathcal O))$ assigning a subnetwork and observer channel to a cover class; \textbf{(ii)} colour data making $\Pi_{\mathcal C}(\Sigma;C)$ computable [FT.19/FT.27]; \textbf{(iii)} a materialization clause for the class's dissolved even cycles (M2 species), after which the pole condition is the existing one, $1\in\operatorname{spec}\mathbf F_M^{(\Sigma)}$ [Def-H.16/CT2.5] with FT.22's single-valuedness on each materialized chain. $(2)$ \textbf{No poles without a declaration.} A cover pattern supplies $R$, $|\mathcal C|$ open threads, the dissolved even cycles with their $\delta h$, the aspect mask and the norm weight $\Omega2^K$ [FT.80(1)] — and none of $\Sigma$, $C$, colour, materialization. The amplitude $4^{|\mathcal C|}$ is a weight in $Z_\Gamma$ [Def-H.25/FT.67(2)], not an eigenvalue of any declared operator, and $|w|=1$ [FT.67(6c)/FT.75(6)] makes the norm layer a \textit{count}, not a spectrum. $(3)$ \textbf{Minimality refinement (new).} For \textbf{colour-free} covers, $\mathrm{Colour}(\Sigma)=\varnothing\implies\text{strong}\notin\mathcal A(\Sigma;C)$ [FT.20's criterion] and the strong summand of $\mathbf F_M^{(\Sigma)}$ is $\equiv0$ [FT.26(2)] — so declaration \textbf{(ii) is vacuous}. Declaration \textbf{(iii)} is needed only for the \textit{bound/fermionic} side, and covers carry no materialized closed chain [FT.80(6)]. Hence the \textbf{minimal bridge is (i) alone, restricted to colour-free covers}. $(4)$ \textbf{The candidate, stated exactly — NOT adopted.} \textbf{The minimal bridge map.} For a cover class $\mathcal O$ on a skeleton $\Gamma$ and an observer $o\notin\operatorname{supp}(\mathcal O)$: $\Sigma(\mathcal O):=$ the vertex support of any representative $R\in\mathcal O$; $C(\mathcal O):=$ the boundary channel induced by $\gamma(o,\Sigma(\mathcal O))$ [Def-H.1]; the spectral side of $\mathcal O$ is $Z(\mathcal R,\Sigma(\mathcal O))$ with poles at $1\in\operatorname{spec}\mathbf F_M^{(\Sigma(\mathcal O))}$ [Def-H.16], aspects by FT.20's criteria.

\begin{itemize}
\item \textbf{Declares:} one map (cover classes $\to$ subnetwork $+$ channel). No mass formula, no materialization, no colour data, no rates, no second scale.
\item \textbf{Buys:} a spectral side on the same objects that carry amplitude, so the inter-family bridge becomes computable; class-invariance is preserved because the vertex support is $\operatorname{Aut}(\Gamma)$-invariant.
\item \textbf{Costs / forbids:} (a) the resulting spectrum is \textbf{observer-relative}, inheriting Def-H.9/FT.3's observer-relativity of $m_{\mathrm{hid}}$ (consistent, but it must be stated); (b) \textbf{no fermionic poles} — the bound side still needs (iii); (c) it \textbf{lifts no ceiling}: $\dim\mathcal K_\Sigma\le18$ [FT.26(3)], and any Def-H.37(2)-based values remain bounded by FT.76(3); (d) traps barred: $4^{|\mathcal C|}$, $\Omega$, $Z^\pm$ are never masses or eigenvalues [FT.79(5)]; $\delta h$ is never a scale [D0.8]; aspects come from FT.20's criteria, never from block membership [FT.27].
\end{itemize}

$(5)$ \textbf{Compatibility.} With FT.90: \textbf{compatible} — the declaration adds no resolution, it only names $\Sigma$ and $C$, so FT.3/FT.4's bounds continue to apply to whatever is read; a bridge cannot exceed the channel budget. With FT.89: \textbf{vacuous} — the Koide comparison yields nothing admissible, so no compatibility condition arises. $(6)$ \textbf{Adoption status.} The map is constructed; no tag of Parts 1–4 imports it; its species is a \textbf{map}, contrasted with the transfer interfaces Def-H.37/Def-H.38 and the Cal-F dynamics datum. Its consequences are enumerable in advance: the pole set of each colour-free cover class, the observer-relativity of that set, and the exact statement of which classes remain spectrally empty.

\textbf{Req.} Def-H.1, Def-H.9, Def-H.16, Def-H.25, Def-H.30; FT.3, FT.19, FT.20, FT.22, FT.26, FT.27, FT.67, FT.75, FT.76, FT.79, FT.80, FT.89, FT.90; CT2.5, CT2.8 (Part 2); D0.8, D1.11; no external data.

\textbf{Pf.}

%dsp-level:0
\[
(1):\ \mathrm{FT.80}(4)\ \text{verbatim};
\]

%dsp-level:0
\[
\qquad (2):\ \text{argument list of}\ Z(\mathcal R,\Sigma)\ \text{vs the data a pattern supplies}\ [\mathrm{FT.80}(1)]
\]

%dsp-level:0
\[
4^{|\mathcal C|}\ \text{is}\ Z_\Gamma\text{'s weight}\ [\mathrm{FT.67}(2)];\qquad|w|=1\implies\sum|w|^2=\#\text{configurations}\ [\mathrm{FT.75}(6),\ \mathrm{CT2.8}]
\]

%dsp-level:0
\[
(3):\ \mathrm{Colour}(\Sigma)=\varnothing\implies\text{strong}\notin\mathcal A(\Sigma;C)\ [\mathrm{FT.20}]\ \wedge\ F^{\Sigma}_{\mathrm{strong}}\equiv0\ [\mathrm{FT.26}(2)]\implies(ii)\ \text{vacuous}
\]

%dsp-level:0
\[
\text{covers carry no materialized closed chain}\ [\mathrm{FT.80}(6)]
\]

%dsp-level:0
\[
\implies (iii)\ \text{not required for transit-side poles}
\]

%dsp-level:0
\[
(4):\ \operatorname{supp}(R)\ \text{is}\ \operatorname{Aut}\text{-invariant}\implies\Sigma(\mathcal O)\ \text{well-defined on the class};\qquad o\notin\Sigma\ [\mathrm{Def\text{-}H.1}]
\]

%dsp-level:0
\[
\text{observer-relativity}:\ \mathrm{Def\text{-}H.9}/\mathrm{FT.3};\qquad\text{bound}:\ \mathrm{FT.26}(3);\qquad(5),(6):\ \text{conjunction}.\ \blacksquare
\]

\textbf{Phys.} The bridge from "the fuzz has amplitude" to "the fuzz has a particle" needs one thing named: which subnetwork a cover class \textit{is}, seen from which outside point. Everything else — the pole condition, the aspect criteria, the mask bound — is already proved and waiting. Colour-free covers need nothing further, because a colourless subnetwork has no strong block to seal; covers with colour would need the sealing data too, and anything fermionic would need the materialization clause as well. The declaration is small, its consequences are enumerable before it is made, and it lifts no ceiling: it would give the flicker gas a spectrum, not a hierarchy.

\vspace{2pt}\noindent\rule{\textwidth}{0.4pt}\vspace{2pt}

\end{theorem}

\phantomsection
\section[Tier F.31 — The Native Clifford Census, the Weak-Chirality…]{\texorpdfstring{Tier F.31 — The Native Clifford Census, the Weak-Chirality Pincer, and the Reversal Census of the Covers}{Tier F.31 - The Native Clifford Census, the Weak-Chirality Pincer, and the Reversal Census of the Covers}}

\begin{theorem}[{\texorpdfstring{FT.92 — EW.5 is convention-exact: re-verification and the rigidity of the named basis}{FT.92 - EW.5 is convention-exact: re-verification and the rigidity of the named basis}}]\label{tag:FT-92}
\textbf{Thm.} In the Pauli basis of $\mathfrak H(\mathbb C^4)$ [A3, CT1.1], with CT1.2's definitions verbatim ($E_0=\sigma_0\otimes\sigma_0$, $\Sigma=\sum_i\sigma_i\otimes\sigma_i$, $W_i=\tfrac12(\sigma_0\otimes\sigma_i+\sigma_i\otimes\sigma_0)$, $J_i=\tfrac12(\sigma_j\otimes\sigma_k-\sigma_k\otimes\sigma_j)$ cyclic, $K_i=\tfrac12(\sigma_0\otimes\sigma_i-\sigma_i\otimes\sigma_0)$, $E^{ij}=\tfrac12(\sigma_i\otimes\sigma_j+\sigma_j\otimes\sigma_i)$), as \textbf{exact matrix equalities over $\mathbb Q(i)$} (no tolerance, no floats): $(1)$ \textbf{The frozen identities all hold exactly as stated.} CT1.3 ($[E_0,\cdot]=0$: 16 brackets); CT1.4 ($[W_i,W_j]=i\epsilon_{ijk}W_k$: all 6 ordered pairs); CT1.5 ($[\Sigma,J_i]=4iK_i$, $[\Sigma,K_i]=-4iJ_i$, $[\Sigma,W_i]=0$: 9); CT1.6 ($[J_i,J_j]=[K_i,K_j]=i\epsilon_{ijk}W_k$: 12); CT1.7 ($[J_i,K_j]=i(\delta_{ij}\Sigma-E^{ij})$: 9, each with all three $W$-components exactly zero); CT1.8 ($[W_i,J_j]=i\epsilon_{ijk}J_k$, $[W_i,K_j]=i\epsilon_{ijk}K_k$: 12); \textbf{EW.5 both clauses at all nine $(i,j)$}: $[K_j,[\Sigma,J_i]]=4\epsilon_{ijk}W_k$ and $[J_j,[\Sigma,K_i]]=-4\epsilon_{ijk}W_k$ (diagonal $i=j$: both sides exactly $0$); Cor CT0.7a ($\sum_iW_i^2=\tfrac32E_0+\tfrac12\Sigma=2\Pi_{\mathcal R}$; $\sum_iW_i^2|_{\wedge^2}=0$); CT0.7(1) ($\mathbb S=\tfrac12(E_0+\Sigma)$ as matrices, $\Pi_{\mathcal R/V}$ idempotent, orthogonal, sum $I_4$, traces $3/1$, ranks $3/1$); the associative relation $\Sigma^2=3E_0-2\Sigma$; $[S,T]=0$ on all 16 basis elements with $\eta=(-,+,+,+)$ [A3]; the full $(S,T)$ parity table of CT1.2 ($W\in\mathcal E_{+-}$, $J\in\mathcal E_{-+}$, $K\in\mathcal E_{--}$, $E_0,\Sigma,E^{ij}\in\mathcal E_{++}$: 17 generators). $(2)$ \textbf{The phase-convention premise is refuted.} CT1.6 holds with $+i\epsilon_{ijk}W_k$ \textbf{exactly} in the frozen convention — there is no "up to a sign" residue to fix; EW.5's own Pf is the two-line consequence $[K_j,4iK_i]=4i\cdot i\,\epsilon_{jik}W_k=4\epsilon_{ijk}W_k$, and D's header already lists EW.5 among the tags machine-verified by \texttt{verify\_\allowbreak{}D.py} in exact complex-rational arithmetic. \textbf{No edit was made or is needed} (Parts 1–2 are frozen; this clause is an layer-local re-verification record). $(3)$ \textbf{Convention rigidity (new).} The group of diagonal sign redefinitions $(J,K,W,\Sigma,E)\mapsto(\lambda_JJ,\lambda_KK,\lambda_WW,\lambda_\Sigma\Sigma,\lambda_EE)$, $\lambda\in\{\pm1\}$, preserving every identity of clause (1) \textit{including the associative relation $\Sigma^2=3E_0-2\Sigma$ and the definition $\mathbb S=\tfrac12(E_0+\Sigma)$} (with $E_0=I_4$ fixed, being the unit) is exactly

%dsp-level:0
\[
\{\mathrm{id},\ \operatorname{Ad}(\mathbb S)\}\cong\mathbb Z_2,\qquad \operatorname{Ad}(\mathbb S):\ (J_i,K_i)\mapsto(-J_i,-K_i),\quad W_i,\Sigma,E^{ij},E_0\ \text{fixed}
\]

(machine: exactly $2$ of the $32$ sign maps survive). Proof chain: $\Sigma^2$-relation $\Rightarrow\lambda_\Sigma=1$; CT1.5 $\Rightarrow\lambda_J=\lambda_K$; CT1.7's $\Sigma$-coefficient $\Rightarrow\lambda_J\lambda_K=\lambda_\Sigma=1$ and its $E$-coefficient $\Rightarrow\lambda_E=1$; CT1.6 $\Rightarrow\lambda_W=\lambda_J^2=1$. \textbf{EW.5 is invariant under the nontrivial element} (each clause maps to itself); the naming exchange $J\leftrightarrow K$ maps clause 1 to clause 2 and is \textit{excluded} by the same associative relation. Hermiticity excludes continuous phases ($e^{i\varphi}K_i$ Hermitian $\Rightarrow\varphi\in\{0,\pi\}$ — the sign enumeration is exhaustive). So the convention space of the named basis is one $\mathbb Z_2$, both members satisfy every frozen identity exactly, and the requested "fix" is vacuous — what the audit adds is the rigidity statement itself.

\textbf{Req.} A3; C1, C2, C3; CT0.7, CT0.7a, CT1.1–CT1.8 (Part 1); EW.4, EW.5 (Part 2); exact $\mathbb Q(i)$ matrix arithmetic, exhaustive sign enumeration [math; machine].

\textbf{Pf.}

%dsp-level:0
\[
(1):\ \text{each identity evaluated entrywise over}\ \mathbb Q(i)\ \text{from CT1.2's definitions; zero residual}\ [\text{machine: suite §A}]
\]

%dsp-level:0
\[
(2):\ \mathrm{CT1.6}\ \text{Thm verbatim};\qquad [K_1,K_2]=\tfrac14\bigl(2i\,\sigma_0\otimes\sigma_3+2i\,\sigma_3\otimes\sigma_0\bigr)=iW_3\ \checkmark\ (+i,\ \text{exact})
\]

%dsp-level:0
\[
(3):\ \Sigma^2=3E_0-2\Sigma\ \wedge\ \lambda_\Sigma\Sigma\ne0\implies\lambda_\Sigma^2\Sigma^2=3E_0-2\lambda_\Sigma\Sigma\implies\lambda_\Sigma=1;
\]

%dsp-level:0
\[
\qquad \mathrm{CT1.5}:\ \lambda_\Sigma\lambda_J=\lambda_K,\ \lambda_\Sigma\lambda_K=\lambda_J
\]

%dsp-level:0
\[
\mathrm{CT1.7}:\ \lambda_J\lambda_K=\lambda_\Sigma=1\ \wedge\ \lambda_J\lambda_K=\lambda_E;\qquad \mathrm{CT1.6}:\ \lambda_J^2=\lambda_K^2=\lambda_W;\qquad\mathrm{CT1.4}:\ \lambda_W^2=\lambda_W\implies\lambda_W=1
\]

%dsp-level:0
\[
\implies(\lambda_J,\lambda_K,\lambda_W,\lambda_\Sigma,\lambda_E)\in\{(1,1,1,1,1),(-1,-1,1,1,1)\};
\]

%dsp-level:0
\[
\qquad \operatorname{Ad}(\mathbb S)J_i\mathbb S=-J_i\ [\mathrm{CT0.7}(3):\ J\ S\text{-odd}]\ \blacksquare
\]

\textbf{Phys.} The electroweak backbone was accused of a sign slip; the audit finds the opposite of a slip — a lock. Once the identity generator, the exchange scalar's associative square, and the swap definition are held fixed, the only remaining freedom in the whole named basis is the simultaneous flip of the two Janus triplets, which is the swap conjugation itself, and every frozen identity — including both clauses of the second-order Higgs-to-weak pathway — reads identically in both conventions. There is no phase to choose.

\vspace{2pt}\noindent\rule{\textwidth}{0.4pt}\vspace{2pt}

\end{theorem}

\begin{theorem}[{\texorpdfstring{FT.93 — The native Clifford census and the weak-chirality pincer (Tasks 1–3)}{FT.93 - The native Clifford census and the weak-chirality pincer (Tasks 1-3)}}]\label{tag:FT-93}
\textbf{Thm.} On $\mathfrak H(\mathbb C^4)$ with CT1.2's basis and CT0.7's $\mathbb S$, $\Pi_{\mathcal R}=\tfrac14(3E_0+\Sigma)$, $\Pi_{\mathcal V}=\tfrac14(E_0-\Sigma)$: $(1)$ \textbf{Clifford frames are native (positive, new).} Among the $15$ non-identity Pauli products (each Hermitian, squaring to $+I_4$, traceless) there are exactly \textbf{$30$ pairwise-anticommuting $4$-frames} (unordered; $720$ ordered, $24$ orderings per frame). For every frame: $\gamma_{\mathcal F}$ is Hermitian, $\gamma_{\mathcal F}^2=I_4$, $\operatorname{Tr}\gamma_{\mathcal F}=0$, $\{\gamma_{\mathcal F},e_i\}=0$; the $16$ real products $\{e_{i_1}\cdots e_{i_r}\}$ are $\mathbb R$-linearly independent (rank $16$ on the example frame) — the frame generates a real subalgebra $\cong Cl(4,0)$ acting on $\mathbb C^4\cong\mathbb R^8$, \textbf{with no new axiom} (A3's Pauli facts + CT1.1's basis suffice). The override's items 1–6 are verified with two precisions: $\gamma_{\mathcal F}$ is \textit{frame-dependent} — every ordering sign flip flips $\gamma_{\mathcal F}$ ($12/12$ witness), and $15$ distinct values occur (below) — so "$\gamma_5=\omega$" names a \textit{choice}, not a canonical element; and A3's $\eta=(-,+,+,+)$ labels the $T$-involution only — it enters no Clifford relation (the signature is Euclidean, $e_i^2=+I_4$, as the override requires; no Minkowski structure is used anywhere). $(2)$ \textbf{The full pseudoscalar catalog (exhaustive).} All $15$ values are attained, \textbf{each by exactly $2$ frames}, in three species:

%dsp-level:0
\[
\begin{array}{c|c|c|c|c}
\text{species} & \gamma_{\mathcal F}\ (\text{up to sign}) & (S,T) & \text{commuting}\ W_i & [\gamma_{\mathcal F},\mathbb S]\\
\hline
\text{mixed}\ (6) & \sigma_0\otimes\sigma_i\ /\ \sigma_i\otimes\sigma_0 & (\neg\text{definite},-) & W_i\ \text{only} & \pm2iJ_i\\
\text{diagonal}\ (3) & \sigma_i\otimes\sigma_i & (+,+)\in\mathcal E_{++} & W_i\ \text{only}\ (\text{acomm.}\ W_j,W_k) & 0\\
\text{off-diag. spatial}\ (6) & \sigma_i\otimes\sigma_j\ (i\ne j) & (\neg\text{definite},+) & \text{none}\ (\text{acomm.}\ W_k,\ k=6{-}i{-}j) & \mp2iK_k\ (k=6{-}i{-}j)
\end{array}
\]

The example frame of the override, $\mathcal F_0=(\sigma_1\otimes\sigma_2,\ \sigma_2\otimes\sigma_2,\ \sigma_3\otimes\sigma_2,\ \sigma_0\otimes\sigma_1)$, gives exactly

%dsp-level:0
\[
\gamma_{\mathcal F_0}=\sigma_0\otimes\sigma_3=\operatorname{diag}(1,-1,1,-1);\qquad
P^{(\mathcal F_0)}_+=\operatorname{diag}(0,1,0,1);\qquad
P^{(\mathcal F_0)}_-=\operatorname{diag}(1,0,1,0),
\]

both idempotent, orthogonal, summing to $I_4$, trace $2$, rank $2$ — the requested exact matrix forms. \textbf{No $\gamma_{\mathcal F}$ commutes with more than one $W_i$; none anticommutes with more than two.} $(3)$ \textbf{The pincer (the sharpened deficit-1 obstruction).} No chirality grade on the pair register can be weak-covariant, from either side: \textbf{(a) commuting side.} $\{M:[M,W_i]=0\ \forall i\}=\operatorname{span}_{\mathbb C}\{\Pi_{\mathcal R},\Pi_{\mathcal V}\}$ (complex dimension $2$ — Schur on $3\oplus1$; machine-solved exactly). Its involutions are $\pm\Pi_{\mathcal R}\pm\Pi_{\mathcal V}$ with traces $3a+b\in\{\pm4,\pm2\}$; a pseudoscalar needs $\operatorname{Tr}=0$ (general: $\operatorname{Tr}(e_1e_2e_3e_4)=-\operatorname{Tr}(e_2e_3e_4e_1)$) — $3a+b=0\wedge a^2=b^2=1$ has \textbf{no solution}. \textbf{(b) anticommuting side.} $\{M:\{M,W_i\}=0\ \forall i\}=\mathbb C\cdot\Pi_{\mathcal V}$ \textbf{exactly} (complex dimension $1$; machine-solved): rank $1$, trace $1$ — singular, never an involution. Structural reason: invertible $M$ with $\{M,W_i\}=0$ makes $\operatorname{Ad}(M)$ a Lie automorphism equal to $-\mathrm{id}$ on $\mathfrak{su}(2)$, but $-\mathrm{id}$ reverses the bracket ($[W_1,W_2]=iW_3\mapsto-iW_3\ne iW_3$) — contradiction. \textbf{(c) exhaustive.} Every one of the $15$ attained $\gamma_{\mathcal F}$ fails (a) and (b) explicitly: $[\gamma_{\mathcal F},W_i]\ne0$ for some $i$ and $\{\gamma_{\mathcal F},W_i\}\ne0$ for some $i$ — all verified entrywise. \textbf{(d) seed side.} On the actual matter carrier $\mathcal H_{\mathsf v}\cong\mathbb C^2$ with $W_i\mapsto\tfrac12\sigma_i$ [EW.32]: $\{m:[m,\sigma_i]=0\}=\mathbb C\cdot I_2$ and $\{m:\{m,\sigma_i\}=0\}=\{0\}$ — the weak seed representation is \textbf{irreducible}: no nontrivial projector on the matter carrier commutes with the weak action, and no grade anticommutes with it either. $(4)$ \textbf{Task 2's claim refuted, per frame and in general.} For $\mathcal F_0$: $W_1,W_2$ \textit{flip} the chiral halves ($P_\mp W_{1,2}P_\pm\ne0$ entrywise) and $W_3$ (the only commuting one) preserves them with $P_-W_3P_-=\operatorname{diag}(1,0,0,0)$ — nonzero, non-scalar on the $2$-dimensional image: "annihilates or acts trivially on the right eigenspace" is false for every $i$. In general, no frame can make a chiral half weak-inactive: $W_iP=0\ \forall i$ forces $\operatorname{im}P\subseteq\bigcap_i\ker W_i=\wedge^2$ [clause (5)], which is $1$-dimensional, while $\operatorname{rank}P^{(\mathcal F)}_\pm=2$. Hence \textbf{$SU(2)_L$-type chiral coupling is NOT natively derivable}, and B6's target $8+7=15$ is not reachable: on $\mathbb C^4\otimes(\mathbb C^3\oplus\mathbb C^1)$ the native counts are $8+8=16$; the missing right state is the [EXT] SM species fact ($\nu_R$ absent — no species list, no neutrino read: application-layer data). $(5)$ \textbf{The native selectivity that does exist (positive, derived).} The weak force grades the pair arena by \textbf{$S$-parity}, not by chirality: $W_i\Pi_{\mathcal V}=0$ exactly $\forall i$ (the weak annihilates the singlet $\wedge^2$), $\sum_iW_i^2=2\Pi_{\mathcal R}$ [Cor CT0.7a] (spin $1$ on $\vee^2$, spin $0$ on $\wedge^2$) — a $3$-active/$1$-inactive split, and $\bigcap_i\ker W_i=\operatorname{im}\Pi_{\mathcal V}$ exactly (machine: the joint kernel is the $1$-dimensional singlet line). This is a \textit{force-arena} selectivity; matter chirality would have to grade a carrier on which the gauge action is trivial — the opposite pattern. $(6)$ \textbf{Task 3, exactly.} $[\gamma_{\mathcal F},\mathbb S]=0\iff\gamma_{\mathcal F}=\pm\sigma_i\otimes\sigma_i$ — the $S$-aligned chiralities are exactly the three \textit{diagonal} ones, and they live in $\mathcal E_{++}$, the gravity/Higgs block [CT1.1, CT1.2], each anticommuting with the two non-own $W$'s (never weak-covariant). The full commutator table is clause (2)'s last column; example: $[\sigma_0\otimes\sigma_3,\ \mathbb S]=2iJ_3$ — \textbf{chirality rotates into the Janus seal sector} under the swap. The $S$-aligned three are never weak-covariant either: $[\sigma_3\otimes\sigma_3,\ W_1]=2iE^{23}$ exactly (machine) — the weak rotation turns an $S$-aligned chirality into gravitational shear. The two splittings provably do not align: $\operatorname{rank}(P^{(\mathcal F_0)}_\pm\Pi_{\mathcal R})=2$, $\operatorname{rank}(P^{(\mathcal F_0)}_\pm\Pi_{\mathcal V})=1$ — each chiral half cuts the triplet $2{+}1$ and the singlet $1{+}1$... (the singlet line meets \textit{both} halves), so neither grading refines the other, and $P^{(\mathcal F_0)}_\pm\notin\operatorname{span}\{\Pi_{\mathcal R},\Pi_{\mathcal V}\}$ ($[P_\pm,\mathbb S]\ne0$). The override's hypothesis — "left-handed fermions are the S-even/S-odd states the Weak force can see" — is \textbf{false as stated}: the weak-visible split is $(\vee^2,\wedge^2)=(3,1)$ by $S$-parity, while every chiral split is $(2,2)$; the dimensions alone forbid the identification. $(7)$ \textbf{Diagnosis and relocation (the chirality deficit: the synthesis).} The override's correction is accepted where it is true: the "missing spinor bundle" framing was too strong — fiberwise, the Clifford algebra and its chiral projectors are \textbf{constructible from A1/A3 with no declaration} (clause (1)), and Def-H.14's $\mathcal E(e)\subset\mathfrak H(\mathbb C^4)_e$ already makes the register a per-edge arena bundle. But $\mathbb C^4$ is the \textbf{internal/force arena}: it carries the weak as $3\oplus1$ (clause (5)), while a Dirac-type $\gamma_5$ must commute with the gauge action on a \textit{spacetime-side} factor. The pincer (clause (3)) is the exact incompatibility: no element of $\operatorname{End}(\mathbb C^4)$ is simultaneously a pseudoscalar and weak-scalar, and the matter carrier $\mathbb C^2$ is weak-irreducible. \textbf{Deficit 1 is therefore relocated, not closed}: what is missing is not matrices but a spacetime spin factor (a carrier on which the gauge action is trivial and a Clifford module over the emergent geometry acts) — and a global Dirac-type operator would additionally require a per-edge frame assignment $e\mapsto\mathcal F_e$ (fiberwise frames exist by clause (1); the assignment is undeclared datum) coupled to the EW.32 transport. Recorded in item 21 as an enumerated would-be declaration, \textbf{not adopted, no declaration-entry number} (the pincer blocks the chiral use regardless of the assignment, and no computable consequence was identified). The reversal structure declared by the corpus remains exactly $\langle\omega,\Theta_{\mathfrak L}\rangle\cong\mathbb Z_2\times\mathbb Z_2$ [FT.88]; the $15$ register-level $\mathbb Z_2$-grades $\gamma_{\mathcal F}$ are new \textit{constructible} objects with no declared coupling to $\omega$ or $\Theta_{\mathfrak L}$ (different spaces: patterns vs fiber). $(8)$ \textbf{The inventory gloss.} The No Dirac/Clifford algebra is \textit{declared}, but $Cl(4,0)$ frames are \textit{derived} (clause (1)) — the question was never the algebra; it is the weak-covariant grade (clause (3), proven absent here) and the spacetime factor (clause (7), absent from the inventory). The chirality-projector non-constructibility result stands, on the proven pincer.

\textbf{Req.} A1, A3; C1, C2, C3; CT0.6, CT0.7, CT0.7a, CT0.12, CT0.14, CT1.1, CT1.2, CT1.4–CT1.8 (Part 1); EW.5, EW.32 (Part 2); Def-H.14; FT.78(6), FT.84, FT.87(8), FT.88, FT.89(2), FT.90(5); frame enumeration, exact kernel solves over $\mathbb Q(i)$, exhaustive catalog verification [math; machine]. External SM counts ($8/7/15$, $\nu_R$): marked, used in no derivation.

\textbf{Pf.}

%dsp-level:0
\[
(1):\ \#\{\text{4-subsets of the }15\}=1365;\qquad\text{pairwise }\{e_a,e_b\}=0\ \text{filter}\implies30\ [\text{machine}]
\]

%dsp-level:0
\[
\gamma_{\mathcal F}^2=(-1)^{\binom42}e_1^2e_2^2e_3^2e_4^2\cdot(\text{reordering sign})=I_4;
\]

%dsp-level:0
\[
\qquad \operatorname{Tr}\gamma=-\operatorname{Tr}\gamma=0;
\]

%dsp-level:0
\[
\qquad \text{rank}_{\mathbb R}\{e_{i_1}\!\cdots e_{i_r}\}=16
\]

%dsp-level:0
\[
(2):\ \text{Pauli-coefficient extraction}\ \operatorname{Tr}(\gamma P_{\mu\nu})/4:\ \text{single}\ \pm P_{\mu\nu}\ \text{per frame};
\]

%dsp-level:0
\[
\qquad \text{frame multiplicity }2\ \text{per value}\ [\text{machine, all }30]
\]

%dsp-level:0
\[
(3a):\ \text{exact nullspace of the }96\times32\ \text{real system}\ [M,W_i]=0:\ \dim_{\mathbb R}4=\dim_{\mathbb C}2;
\]

%dsp-level:0
\[
\qquad M=a\Pi_{\mathcal R}+b\Pi_{\mathcal V}:\ M^2=I\Rightarrow a^2=b^2=1;\ \operatorname{Tr}M=0\Rightarrow3a+b=0:\ \varnothing
\]

%dsp-level:0
\[
(3b):\ \text{exact nullspace of}\ \{M,W_i\}=0:\ \dim_{\mathbb C}1=\mathbb C\cdot\Pi_{\mathcal V}\ (\operatorname{rk}1,\operatorname{Tr}1);
\]

%dsp-level:0
\[
\qquad \operatorname{Ad}(M)=-\mathrm{id}\ \neg\in\operatorname{Aut}(\mathfrak{su}(2)):\ [W_1,W_2]=iW_3\ne0
\]

%dsp-level:0
\[
(3d):\ 2\times2\ \text{exact solves}:\ \dim_{\mathbb C}\{m:[m,\sigma_i]=0\}=1;\qquad\dim_{\mathbb C}\{m:\{m,\sigma_i\}=0\}=0
\]

%dsp-level:0
\[
(4):\ \text{block entries of}\ P_\mp W_{1,2}P_\pm\ \text{nonzero exact};
\]

%dsp-level:0
\[
\qquad P_-W_3P_-=\operatorname{diag}(1,0,0,0);
\]

%dsp-level:0
\[
\qquad \dim\bigcap_i\ker W_i=1<2=\operatorname{rk}P_\pm
\]

%dsp-level:0
\[
(5):\ W_i\Pi_{\mathcal V}=0\ [\text{machine}];\qquad\sum W_i^2=\tfrac32E_0+\tfrac12\Sigma\ [\mathrm{Cor\ CT0.7a}]=2\Pi_{\mathcal R}
\]

%dsp-level:0
\[
(6):\ \mathbb S P_{\mu\nu}\mathbb S=P_{\nu\mu}\implies[\gamma,\mathbb S]=0\iff\mu=\nu;
\]

%dsp-level:0
\[
\qquad [\sigma_0\otimes\sigma_3,\mathbb S]=\tfrac12[\sigma_0\otimes\sigma_3,\Sigma]=2iJ_3\ [\text{machine}]
\]

%dsp-level:0
\[
\operatorname{rank}(P_\pm\Pi_{\mathcal R})=2,\ \operatorname{rank}(P_\pm\Pi_{\mathcal V})=1\ \forall\pm\implies\neg\text{refinement either way};\qquad(3,1)\ne(2,2)\implies\neg\text{identification}
\]

%dsp-level:0
\[
(7),(8):\ \text{conjunction of}\ (1),(3),(5),(6)\ \text{with the cited inventory tags.}\ \blacksquare
\]

\textbf{Phys.} The Imaginator was half right, and the half that is right matters: the pair register really does carry a native Euclidean Clifford algebra — thirty frames, fifteen possible "$\gamma$$_5$'s", exact chiral projectors, all constructible from the axioms with nothing added. But every one of those fifteen chirality operators is blind to, or at war with, the weak force: the best of them commutes with a single Cartan generator while the other two rotate left into right, and the swap rotates all of them into the Janus seals. The weak force does have a native selectivity — it sees the symmetric triplet and annihilates the antisymmetric singlet, three against one — but that is a grading of the force arena by exchange parity, not a handedness of matter. A Dirac $\gamma$$_5$ needs a spacetime-side factor on which the gauge group acts trivially; the corpus has no such factor, and now there is a theorem saying the pair register cannot fake one. Chirality is still missing — but the search space has collapsed from "unknown primitive" to one named absent object.

\vspace{2pt}\noindent\rule{\textwidth}{0.4pt}\vspace{2pt}

\end{theorem}

\medskip\noindent\textit{\texorpdfstring{FT.94 — The chiral mass spectrum: not constructible as chiral, exact as a statistic, and the fourth hierarchy negative}{FT.94 - The chiral mass spectrum: not constructible as chiral, exact as a statistic, and the fourth hierarchy negative}.}\ Structural-constraints record: see Appendix B, Section B.2.6 of the Addendum.\par\medskip

\begin{theorem}[{\texorpdfstring{FT.95 — Flicker chirality: the tautology, the absent trajectory, and the zero-expectation theorem}{FT.95 - Flicker chirality: the tautology, the absent trajectory, and the zero-expectation theorem}}]\label{tag:FT-95}
\textbf{Thm.} Under Cal-F, Def-H.25, D2.6: $(1)$ \textbf{D4/D5 are tautological.} With $k(t)=(-1)^tk_0$ and $\mathrm{coupling}(t)=\mathrm{wb}\cdot\mathrm{hc}\cdot k(t)$: $\operatorname{sign}(k(t))\cdot k(t)=|k(t)|=|k_0|$ for every $t$ (machine: all $k\in\{\pm1,\dots,\pm5\}$), so $\langle\chi\rangle_T=\mathrm{wb}\cdot\mathrm{hc}\cdot|k_0|$ \textbf{for every $T\ge1$} — independent of the dynamics, the tick count and the phase boundary. "$\langle\chi\rangle_T\ne0$" is built in by the formula, not derived: the same species as FT.87(3)'s Regge-tautology verdict. D5 is void; D6 substitutes a measure parameter ($p^*_\infty=\sqrt{2/3}$ — itself FT.89(2d)'s cap-3 output, and a phase-boundary datum of the bypass measure [Def-H.26/FT.70], not a dynamics input) into the tautology and returns the tautology. $(2)$ \textbf{D1–D3 have no flicker referent.} The winding $k$ is a read of \textit{materialized} closed chains [EW.18, FT.22]; flicker states carry none — bound content survives only dissolved, in even cycles [FT.69's cycle-alternation lemma], and covers carry no materialized closed chain [FT.80(6)] — so no trajectory $k(t)$ exists in the flicker sector. Cal-F declares the \textit{rates} $r=\Gamma_{J^+}/\Gamma_{J^-}$ (per-edge independent detailed balance), not per-tick schedules; $\delta\tau$ is schedule data that $R$ does not determine [FT.80(5)]; and no oscillation frequency is constructible [FT.86(7): only $\tau_c=1$ ledger tick, CT0.3]. $\mathrm{weak\_bias}$ and $\mathrm{higgs\_coupling}$ count \textbf{0} in frozen Parts 1–2 (machine-asserted). The corpus's genuine period-2 phenomenon is FT.86(2)'s alternating \textit{correlation ripple in the crystallization indicators} (cap-induced, exact decay rate $(\sqrt{21}-5)/2$) — a different object: it is $C$-even, orientation-blind, and carries no chirality; D2's deterministic sign-alternating winding conflates the bypass and bound sectors. $(3)$ \textbf{The zero-expectation theorem (the positive replacement for D5).} Global reversal $R\mapsto R^{\mathrm T}$ preserves admissibility, $K$, cover status and the amplitude $A(R)=(-1)^{|R|}2^{|R|}$ — \textbf{exhaustively verified on the tetrahedron}: all $1230$ bypass configurations, all $101$ covers, and all $4096$ patterns; the norm-layer class data $(K,\Omega)$ are reversal-invariant by FT.80(1). Hence both layer weights of the declared flicker ensemble are reversal-invariant, and \textbf{every $C$-odd read has exactly zero ensemble expectation}: for $f(R^{\mathrm T})=-f(R)$,

%dsp-level:0
\[
\sum_R A(R)\,f(R)=0
\]

%dsp-level:0
\[
\qquad (\text{machine witness: the orientation-bias read}\ f(R)=\#\{(a,b)\in R:a<b\}-\#\{(a,b)\in R:a>b\},\ \text{exact}\ 0\ \text{over all}\ 4096).
\]

In particular $\langle q_{\mathrm{core}}\rangle=0$ ($\omega$-odd: FT.40/D2.6) and $\langle\operatorname{sign}(k)\rangle=0$ ($C$-odd: FT.87(8)) in the declared ensemble. \textbf{Parity violation is not constructible from the flicker}: producing $\langle C\text{-odd}\rangle\ne0$ requires declaring a $C$-odd state or a $C$-odd dynamics — application-layer territory; the exact declared reversal structure remains $\langle\omega,\Theta_{\mathfrak L}\rangle\cong\mathbb Z_2\times\mathbb Z_2$ [FT.88], whose elements are symmetries of the registration data, and by this clause of the equilibrium weights.

\textbf{Req.} Cal-F; Def-H.25, Def-H.26; FT.22, FT.40, FT.69, FT.80, FT.86, FT.87, FT.88, FT.89(2); D2.6; CT0.3; EW.18 (Part 2); exhaustive tetrahedron enumeration, exact integer sums [math; machine].

\textbf{Pf.}

%dsp-level:0
\[
(1):\ \operatorname{sign}(x)x=|x|\ \forall x\in\mathbb Z\setminus\{0\};\qquad\langle\chi\rangle_T=\tfrac1T\sum_{t=0}^{T-1}|k_0|\cdot\mathrm{wb}\cdot\mathrm{hc}\ [\text{machine},\ T=1..6]
\]

%dsp-level:0
\[
(2):\ \mathrm{FT.69}\ \text{(dissolved only)}+\mathrm{FT.80}(6)\implies\neg\exists\ \text{materialized}\ k(t);\qquad\text{grep counts}=0\ [\text{machine}]
\]

%dsp-level:0
\[
(3):\ A(R^{\mathrm T})=(-1)^{|R^{\mathrm T}|}2^{|R^{\mathrm T}|}=A(R);
\]

%dsp-level:0
\[
\qquad R\mapsto R^{\mathrm T}\ \text{a bijection of the ensemble}\implies\sum_RA(R)f(R)=\tfrac12\sum_R\bigl(A(R)f(R)+A(R^{\mathrm T})f(R^{\mathrm T})\bigr)=0\ \blacksquare
\]

\textbf{Phys.} The proposed flicker chirality is a formula that multiplies a sign by the thing the sign came from — the answer is the absolute value, every time, for every period, at every phase boundary. Strip that away and ask the real question — can the flickering gas prefer one orientation? — and the corpus answers with a clean no-theorem: reversing every arrow leaves the weights of both layers untouched, configuration by configuration, so anything odd under reversal averages to exactly zero. The gas is born parity-symmetric. A parity-violating world needs a chiral state or a chiral rule, and neither is declared; that absence is the chirality assignment's dynamics-side face, now with an exact vanishing theorem behind it.

\vspace{2pt}\noindent\rule{\textwidth}{0.4pt}\vspace{2pt}

\end{theorem}

\phantomsection
\section[Tier F.32 — The Multi-Node Form Factor: Formalization,…]{\texorpdfstring{Tier F.32 — The Multi-Node Form Factor: Formalization, Exact Reductions, and the Third Ceiling}{Tier F.32 - The Multi-Node Form Factor: Formalization, Exact Reductions, and the Third Ceiling}}

\begin{theorem}[{\texorpdfstring{FT.97 — The form factor formalized: node ontology, three underdeterminations, reduction identities, and the general ceiling}{FT.97 - The form factor formalized: node ontology, three underdeterminations, reduction identities, and the general ceiling}}]\label{tag:FT-97}
\textbf{Thm.} Under Def-H.25/FT.67, at (S1)–(S4): $(1)$ \textbf{Node ontology (premise correction).} For every flicker configuration $R$ on a skeleton $\Gamma$: $V(R)\subseteq V(\Gamma)$ — the nodes a configuration "passes through" are exactly the skeleton vertices its edges touch. The roles are the derived vertex types $(a,b)=(\mathrm{indeg}_R,\mathrm{outdeg}_R)$: \textbf{centre} $(1,1)$ (the threading junction; every cover edge is claimed by exactly one centre [FT.67(3d)]), \textbf{endpoint} $(a,0)/(0,b)$, \textbf{hub} $a+b\ge4$ (locked, $D^-$-configured [FT.3/CT4.5]), \textbf{unclaimed} $(0,0)$. "Intermediate nodes" do not exist as nodes; a form factor over subdivided skeletons has no corpus referent (no refinement mechanism is declared). $(2)$ \textbf{Three underdeterminations, exposed and fixed by convention (none is a free fit).} \textbf{(a) Per-node information assignment.} Registration bits live on \textit{edges} [CT1.29: $S_{\mathfrak L}=n_{\mathcal R}\ln2$; D0.5]; assigning them to nodes requires a convention — tail ($i_v=\mathrm{outdeg}_R(v)$), end-sum ($a{+}b$) or half-split ($\tfrac{a+b}{2}$) — and the first two differ by an \textbf{exact factor $2$} on covers ($\sum_v\deg_R(v)=2|R|$). No rule of Parts 1–3 selects one (this is FT.85(7)'s per-sector non-constructibility in miniature); the \textbf{tail convention is primary} here because $\sum_v i_v=|R|=n_{\mathcal R}$ matches CT1.29's count exactly, and the factor-$2$ freedom is reported, not hidden. \textbf{(b) The role weight $w$.} Five derived candidates, all functions of the finite type set: $w_U=1_{\mathrm{participating}}$, $w_K=k_v=a{+}b$ [CT0.11's capped quantity], $w_C=1_{(1,1)}$ (centre), $w_E=1_{\mathrm{endpoint}}$, $w_H=\max(0,a{+}b{-}2)$ (hub excess). Any other $w$ is a free function — a P4-class parameter whose selection against [EXT] targets would violate P5; \textbf{all five are reported and none is selected by any target}. \textbf{(c) Sector scope} (bypass / locked / all): values differ by $O(2)$ and all three are tabulated. $(3)$ \textbf{Reduction identities (new; exact, per configuration).} For every strict cover $R$:

%dsp-level:0
\[
S_{w_U}(R)=\sum_v \mathrm{outdeg}_R(v)=|R|=K(R);\qquad S_{w_C}(R)=\#\{(1,1)\ \text{vertices}\}=|\mathcal C(R)|=\tfrac{K}{2}
\]

(machine-asserted on all $17\,042$ covers of the three families). Hence the directive's form factor at its constructible core is

%dsp-level:0
\[
M_{\mathrm{form}}^{(w_U)}(\Gamma)=\frac{\sum_R|A(R)|^2K(R)}{\sum_R|A(R)|^2}\,\ln2=\langle K\rangle_{|A|^2}\ln2\qquad(|A|^2=16^{|\mathcal C|}\ \text{on covers}),
\]

\textbf{exactly FT.68(1)'s reading (ii)} — the amplitude-weighted mean winding, tabulated for the tetrahedron bypass sector since FT.68. Anchor reproduced exactly from independent enumeration: $M_{\mathrm{form}}^{(w_U)}(\mathrm{tetra},\mathrm{byp})=62208/15745$; and $M^{(w_C)}=M^{(w_U)}/2$ identically, per family. The "paradigm shift" is a \textit{reweighting of tabulated census data}, not a new read species. $(4)$ \textbf{The typed tuple remains non-constructible and ill-typed.} Per component [FT.94(1)'s citations, per-node scope added]: Topological — winding is per materialized \textit{chain} [FT.22], no per-node referent, and flicker edges are unmaterialized [FT.80(6)]; $G$ — the geometric quantum is a \textit{joint} $S$-even read over all ten metric coordinates [FT.30's verdict], its ledger unbounded-integer, $\neg$bit-type [FT.41(c)/C10]; $H$ — the continuous global P4 datum $h_0$ [EW.13], not per-node; $W$ — continuous per-edge shift data [CT1.12]; Strong — sealing bits are well-posed [FT.41(b)] but colour data is not determined by crystallization patterns [FT.80(2)]. A sum $\mathrm{Topological}+G+H+W+\mathrm{Strong}$ adds bits to unbounded integers to reals — \textbf{ill-typed} [C10]. The only constructible per-node information read is the registration bit count — a bit sum, which the directive's own D3 forbids: \textbf{D3 cannot be satisfied as written}; the constructible version is computed instead and flagged. $(5)$ \textbf{The Form-Factor Ceiling (the third size theorem).} $M_{\mathrm{form}}$ is a weighted \textit{mean}, so $M_{\mathrm{form}}\le\max_R S_w(R)$; and $S_w(R)\le V\cdot c_w$ with the finite-type maxima

%dsp-level:0
\[
c_w:\quad c_{w_U}=6,\quad c_{w_K}=36,\quad c_{w_C}=1,\quad c_{w_E}=6,\quad c_{w_H}=24\qquad(\text{universal over the type set}\ \{0<a+b\le6\}),
\]

so $M_{\mathrm{form}}(\Gamma)\le c_wV(\Gamma)$ — \textbf{linear in skeleton size for every role weight}. With $|\mathcal C|\le3V/4$ (bypass) / $V$ (locked) [CT0.11/FT.79(2)], the $16^{|\mathcal C|}$ enhancement \textbf{saturates at the cap}: exponential \textit{inside} a skeleton, linear \textit{across} skeletons. Spanning $1{:}R$ off the $P_3$ anchor requires $V\ge R\cdot M^{(w)}(P_3)/c_w$; at $R=3477$ [EXT]:

%dsp-level:0
\[
V\ \ge\ \begin{cases}2187 & (w_U)\\ 1062 & (w_K)\\ 6561 & (w_C)\\ 65 & (w_E)\\ 515 & (w_H)\end{cases}
\qquad(\text{exact: } \lceil 3477\cdot\tfrac{3264}{865\cdot6}\rceil,\ \lceil 3477\cdot\tfrac{9504}{865\cdot36}\rceil,\ \lceil 3477\cdot\tfrac{1632}{865}\rceil,\ \lceil 3477\cdot\tfrac{96}{865\cdot6}\rceil,\ \lceil 3477\cdot\tfrac{3072}{865\cdot24}\rceil)
\]

(the $w_E$ bound is loose — its $P_3$ anchor $96/865$ is tiny; its \textit{computed} family maximum is $7.681$, FT.98(4)). This is the third instance of the size-bound species [FT.76(3): pole ceiling; FT.79(2): cover ceiling]: \textbf{the two-ceiling synthesis [FT.84(2)] becomes a three-ceiling synthesis} — counting reads of every constructed species on every computed skeleton are $O(V)$. $(6)$ \textbf{Mass status.} $M_{\mathrm{form}}$ is a flicker-ensemble \textbf{statistic} (the species of FT.94(3)); promoting it to a particle mass requires a declaration linking ensemble averages to pole structure [Def-H.16] — none exists, none is adopted. $|A|^2$ enters only as a weight (no amplitude-to-mass promotion; the FT.79(5)/FT.87(9) boundary is respected). "Mass is a form factor" is recorded as a proposed identification, \textbf{not adopted}.

\textbf{Req.} Def-H.25; FT.3, FT.16, FT.22, FT.41, FT.42, FT.67, FT.68(1), FT.76(3), FT.79, FT.80, FT.84(2), FT.85(7), FT.94; D0.5; CT0.11 (Part 1); CT1.12, CT1.29 (Part 2); EW.13 (Part 2); C10; P4, P5; finite-type maximization, exact Fraction census averages [math; machine]. External hierarchies marked at (5); no identification.

\textbf{Pf.}

%dsp-level:0
\[
(1):\ \mathrm{Def\text{-}H.25}:\ R\subseteq\mathcal R_{\rightarrow}(\Gamma);\qquad\text{roles}=\mathrm{types\_of}(R)\ [\text{derived}]
\]

%dsp-level:0
\[
(3):\ \text{strict covers}: \text{every edge claimed once by a }(1,1)\text{ centre}
\]

%dsp-level:0
\[
\implies 2\#\mathrm{centres}
\]

%dsp-level:0
\[
= |R|\ [\mathrm{FT.67}(2)/(3d)]
\]

%dsp-level:0
\[
\sum_v\mathrm{outdeg}_v=|R|\ \text{(each edge has one tail)};
\]

%dsp-level:0
\[
\qquad M^{(w_U)}(\mathrm{tetra,byp})=\tfrac{62208}{15745}\ [\text{machine}=\mathrm{FT.68}(1)(ii)]
\]

%dsp-level:0
\[
(5):\ \text{mean}\le\max;\qquad S_w(R)=\sum_vw(v)i_v\le V\max_{(a,b)}w(a,b)\cdot b;\qquad\text{type set finite}\ (a+b\le6)
\]

%dsp-level:0
\[
|\mathcal C|\le3V/4\ [\mathrm{FT.79}(2)]\implies 16^{|\mathcal C|}\le16^{3V/4}:\ \text{weights saturate};
\]

%dsp-level:0
\[
\qquad \text{required-}V:\ \text{ceiling rearranged at the exact anchors}\ \blacksquare
\]

\textbf{Phys.} The proposal asked the gas to weigh its nodes instead of its edges. Taken literally, the nodes turn out to be the same old skeleton vertices wearing role costumes — centres, endpoints, hubs — and once the costumes are counted honestly, the new "form factor" collapses, identity by identity, onto a number the corpus has carried since FT.68: the amplitude-weighted mean winding of the covers. That is not a failure; it is a localization. Every way of building the proposal out of declared reads lands in the same bounded band, and the bound is a theorem: a weighted average cannot outrun its largest summand, its largest summand cannot outrun the per-vertex cap, and the per-vertex cap cannot outrun the skeleton's size. Three ceilings, one species. Hierarchy-hunting on four- and eight-vertex skeletons is over; the numbers say what would be needed, exactly.

\vspace{2pt}\noindent\rule{\textwidth}{0.4pt}\vspace{2pt}

\end{theorem}

\phantomsection
\section[Tier F.33 — The Dynamic Phase, Parity, and the CPT Bridge:…]{\texorpdfstring{Tier F.33 — The Dynamic Phase, Parity, and the CPT Bridge: the Mass Refinement}{Tier F.33 - The Dynamic Phase, Parity, and the CPT Bridge: the Mass Refinement}}

\begin{theorem}[{\texorpdfstring{FT.100 — The alignment pattern refines the informational mass}{FT.100 - The alignment pattern refines the informational mass}}]\label{tag:FT-100}
\textbf{Thm.} (candidate.) Under Def-H.9 (coarse-graining projection), the FT.16/FT.44(4) mass read ($m_{\mathrm{hid}}=E_{\mathrm{int}}\ln2/\mathfrak h$), and Def-EM.20 (dynamic phase): $(1)$ \textbf{Edge count.} The informational mass counts internal real edges: each real edge contributes $\ln2/\mathfrak h$; $m_{\mathrm{hid}}=E_{\mathrm{int}}\ln2/\mathfrak h$, $E_{\mathrm{int}}\in\mathbb Z_{\ge0}$ — on the floor families whose $E_{\mathrm{int}}$ read is computed [FT.32, FT.44(4)]. $(2)$ \textbf{Phase-blindness.} The dynamic phase $\alpha_e$ does not change $m_{\mathrm{hid}}$: the mass counts edges, $\neg$ phases. An edge with $\alpha_e=0$ and an edge with $\alpha_e=\pi$ each contribute $\ln2/\mathfrak h$; $m_{\mathrm{hid}}$ is constant across all $2^{E_{\mathrm{int}}}$ phase assignments. $(3)$ \textbf{Refinement.} The alignment pattern $\mathcal P(\Sigma)$ [Def-EM.21] is a refinement of the mass read: two subnetworks with the same $m_{\mathrm{hid}}$ (same $E_{\mathrm{int}}$) can carry different alignment patterns — $2^{E_{\mathrm{int}}}$ of them [EM.23(4)] — and different parity classes [EM.23(1)–(2)]. The phase adds structure without changing mass.

\textbf{Req.} Def-H.9; FT.1, FT.15, FT.16 (\textbf{Tiers F.1–F.11}); FT.32 (Tiers F.12–F.13); FT.44 (Tiers F.14–F.16); Def-EM.20, Def-EM.21; EM.20, EM.23 (Part 2); C9; CT0.9 (Part 1); exhaustive enumeration in exact integer arithmetic [math; machine: \texttt{checks/\allowbreak{}verify\_\allowbreak{}F33\_\allowbreak{}dynamic\_\allowbreak{}phase.py} §E].

\textbf{Pf.}

%dsp-level:0
\[
(1):\ \mathrm{Def\text{-}H.9}(2):\ E_{\mathrm{tot}}=|\Sigma|\ln2\ [\mathrm{CT0.9},\ \mathrm{C9}:\ 1\ \text{real edge}=1\ \text{bit}=\ln2\ \text{nats}];
\]

%dsp-level:0
\[
\qquad \mathrm{FT.32}/\mathrm{FT.44}(4):\ m_{\mathrm{hid}}=S_{\mathfrak L}/\mathfrak h=E_{\mathrm{int}}\ln2/\mathfrak h\ \text{on the computed floors}
\]

%dsp-level:0
\[
(2):\ m_{\mathrm{hid}}=(E_{\mathrm{tot}}-E_{\mathrm{vis}})/\mathfrak h;\quad E_{\mathrm{vis}}=\min(\delta h,\ 4\,\delta\tau)\ln2\ [\mathrm{FT.1}]:\ \text{neither term reads}\ \alpha_e;
\]

%dsp-level:0
\[
\qquad \mathrm{FT.15}:\ \Pi_v\ \text{depends on}\ (\delta h,\delta\tau,\sigma_s,\sigma_t,\Theta_{\mathrm{int}}),\ \neg\text{on the phase distribution}
\]

%dsp-level:0
\[
\mathrm{EM.20}(1):\ \alpha_e\in\{0,\pi\}\ \forall e
\]

%dsp-level:0
\[
\in E_{\mathcal R}:\ \text{every assignment leaves}\ E_{\mathrm{int}},
\]

%dsp-level:0
\[
\ \text{hence}\ m_{\mathrm{hid}},
\]

%dsp-level:0
\[
\ \text{fixed}\ [\text{machine: all}\ 2^{E_{\mathrm{int}}}\ \text{assignments, exact}]
\]

%dsp-level:0
\[
(3):\ \text{same}\ E_{\mathrm{int}}\implies\text{same}\ m_{\mathrm{hid}}\ [(1)];
\]

%dsp-level:0
\[
\qquad \mathrm{EM.23}(4):\ \#\{\mathcal P(\Sigma)\}=2^{|E_{\mathcal R}(\Sigma)|}:\ \text{fiber cardinality over each mass value};
\]

%dsp-level:0
\[
\qquad \mathrm{EM.23}(1)\text{--}(2):\ \text{the patterns split into parity classes under}\ M\ [\mathrm{Def\text{-}EM.22}].\ \blacksquare
\]

\textbf{Phys.} What a coarse-grained fermion forgets, the alignment pattern remembers: the mass counts the hidden edges — one bit each — blind to which way each edge's phase points; parity reads exactly that pointing data. Same mass, $2^{E_{\mathrm{int}}}$ patterns, split into mirror-even and mirror-odd classes: a degeneracy-refinement template of the same species as FT.44's dual-pair splitting. No physical identification claimed.

\vspace{2pt}\noindent\rule{\textwidth}{0.4pt}\vspace{2pt}

\end{theorem}

\phantomsection
\section[Tier F.34 — The Weak Branch Grading of the Two-Branch…]{\texorpdfstring{Tier F.34 — The Weak Branch Grading of the Two-Branch Crystallization}{Tier F.34 - The Weak Branch Grading of the Two-Branch Crystallization}}

\begin{theorem}[{\texorpdfstring{FT.101 — The weak branch grading: the two crystallization branches are weak-inequivalent}{FT.101 - The weak branch grading: the two crystallization branches are weak-inequivalent}}]\label{tag:FT-101}
\textbf{Thm.} Under A3 ($\mathcal H_p\otimes\mathcal H_q=\vee^2\oplus\wedge^2$), A4 (branches: $\alpha_{pq}=0$ on $|\Pi_{\mathcal V}\rho_{pq}|\le\varepsilon_q$; $\alpha_{pq}=\pi$ on $|\Pi_{\mathcal V}\rho_{pq}|\ge1-\varepsilon_q$), $\Pi_{\mathcal R}=\tfrac14(3E_0+\Sigma)$, $\Pi_{\mathcal V}=\tfrac14(E_0-\Sigma)$ [CT0.7], $W_i=\tfrac12(\sigma_0\otimes\sigma_i+\sigma_i\otimes\sigma_0)$ [CT1.2, FT.13], $\Phi=\{0,\pi\}$ [D0.5, CT0.9]: $(1)$ \textbf{Aligned branch: weak-active.} $\ [\mathbb S,W_i]=0$; $\ W_i\Pi_{\mathcal R}=W_i=\Pi_{\mathcal R}W_i\Pi_{\mathcal R}\ne0$ for every $i$ (entrywise witnesses); $\ \sum_iW_i^2=2\,\Pi_{\mathcal R}$; Casimir read on $\vee^2$: $2=1(1{+}1)$ [CT0.7a]. $(2)$ \textbf{Antialigned branch: weak-blind.} $\ W_i\Pi_{\mathcal V}=0=\Pi_{\mathcal V}W_i$ for every $i$; $\ \bigcap_i\ker W_i=\operatorname{im}\Pi_{\mathcal V}$ (the singlet line, dimension $1$); Casimir read on $\wedge^2$: $0$ [CT0.7a, FT.93(5)]. $(3)$ \textbf{The branch flip is label-level.} The unique nontrivial automorphism $\iota$ of $\Phi$: $\iota(\alpha)=\pi-\alpha$; it exchanges the A4 branch labels ($\alpha=0\leftrightarrow\alpha=\pi$). On configuration data $\iota$ is the alignment-pattern complement $\mathcal P\mapsto E_{\mathcal R}\setminus\mathcal P$ [Def-EM.21, EM.23, EM.24(2); CT1.11a(4)] — exhaustive machine census: the complement flips $\alpha_e$ on every edge and exchanges the uniform patterns $\varnothing\leftrightarrow E_{\mathcal R}$ [EM.23(4)]. \textbf{No Hilbert-space implementation exists:} $\neg\exists\,M\in\operatorname{End}(\mathbb C^4)$ invertible with $M\Pi_{\mathcal R}M^{-1}=\Pi_{\mathcal V}$ — conjugation preserves trace and rank, and $\mathfrak T\Pi_{\mathcal R}=3\ne1=\mathfrak T\Pi_{\mathcal V}$, $\operatorname{rk}\Pi_{\mathcal R}=3\ne1=\operatorname{rk}\Pi_{\mathcal V}$ [CT0.7(1); exact witness]. $(4)$ \textbf{The weak support is not flip-invariant.} The weak support is $\Pi_{\mathcal R}$: $W_i=\Pi_{\mathcal R}W_i\Pi_{\mathcal R}$ exactly [(1)]; $\iota$ maps the branch label $\mathcal R\mapsto\mathcal V$, on which the weak read is $0$ [(2)]: weak-activity is branch-selective and $\iota$-asymmetric — the flip maps the weak-active branch to the weak-blind branch. Sharp point $\varepsilon_q=0$: branch $\alpha=0\iff\operatorname{supp}\rho\subseteq\vee^2$, branch $\alpha=\pi\iff\rho=\Pi_{\mathcal V}$ [CT0.8]; witnesses: $W_3\cdot(\Pi_{\mathcal R}/3)\ne0$, $W_i\,\Pi_{\mathcal V}=0\ \forall i$ (exact). $(5)$ \textbf{Structural parity non-invariance.} The weak couples to one A4 branch and annihilates the other: a $\mathbb Z_2$-asymmetry of the force on the registration alphabet of the real edges; the branch exchange does not preserve the weak action. \textbf{Scope:} the force-arena $S$-parity grading [FT.93(5), FT.93(7)] — $\neg$ a matter chirality projector: the pincer [FT.93(3)] and the relocation [FT.93(7)] stand unmodified, and the chirality-assignment question is untouched by this theorem.

\textbf{Req.} A3; A4; C1; D0.5; CT0.7, CT0.7a, CT0.8, CT0.9, CT1.1, CT1.2, CT1.11a (Part 1); Def-EM.21, EM.23, EM.24 (Part 2); FT.13 (Tiers F.1–F.11); FT.93; exact matrix arithmetic over $\mathbb Q(i)$, exhaustive small-configuration enumeration, exact rank/trace invariants [math; machine: \texttt{checks/\allowbreak{}verify\_\allowbreak{}H\_\allowbreak{}branch\_\allowbreak{}parity.py}].

\textbf{Pf.}

%dsp-level:0
\[
(1):\ [\Sigma,W_i]=0\ [\mathrm{CT1.5}];
\]

%dsp-level:0
\[
\qquad \mathbb S=\tfrac12(E_0+\Sigma)\ [\mathrm{CT0.7}(1)]\implies[\mathbb S,W_i]=0\implies W_i\ \text{preserves}\ \vee^2,\ \wedge^2\ [\mathrm{CT0.7}(3)]
\]

%dsp-level:0
\[
W_i\Pi_{\mathcal V}
\]

%dsp-level:0
\[
= 0\ [(2)]
\]

%dsp-level:0
\[
\implies W_i
\]

%dsp-level:0
\[
= W_i(\Pi_{\mathcal R}+\Pi_{\mathcal V})
\]

%dsp-level:0
\[
= W_i\Pi_{\mathcal R}
\]

%dsp-level:0
\[
= \Pi_{\mathcal R}W_i\Pi_{\mathcal R}\ [\mathrm{CT0.7}(1):\ \Pi_{\mathcal R}+\Pi_{\mathcal V}=E_0];
\]

%dsp-level:0
\[
\qquad \text{nonzero witnesses entrywise}\ [\text{machine}]
\]

%dsp-level:0
\[
\textstyle\sum_iW_i^2=\tfrac32E_0+\tfrac12\Sigma=2\,\Pi_{\mathcal R}\ [\mathrm{CT0.7a}];\qquad\vee^2:\ \mathbb S=+1\implies\Sigma=+1\implies\tfrac32+\tfrac12=2
\]

%dsp-level:0
\[
(2):\ W_i\,\Pi_{\mathcal V}=0\ \forall i:\ \text{direct exact computation on the CT1.2 basis}\ [\text{machine}];
\]

%dsp-level:0
\[
\qquad \wedge^2=\mathbb C\,|\Psi^-\rangle,\ \ \Pi_{\mathcal V}=|\Psi^-\rangle\langle\Psi^-|\ [\mathrm{CT0.7}(1);\ \text{machine}]
\]

%dsp-level:0
\[
\dim\bigcap_i\ker W_i=4-\operatorname{rk}\binom{W_1}{W_2}{W_3}=4-3=1\ [\text{exact elimination; machine}];
\]

%dsp-level:0
\[
\qquad \wedge^2:\ \tfrac32+\tfrac12(-3)=0\ [\mathrm{CT0.7a}:\ \Sigma|_{\wedge^2}=-3]
\]

%dsp-level:0
\[
(3):\ |\Phi|=2\ [\mathrm{CT0.9}]\implies\operatorname{Aut}(\Phi)=\mathbb Z_2=\{\mathrm{id},\iota\},\ \iota(\alpha)=\pi-\alpha,\ \iota^2=\mathrm{id};
\]

%dsp-level:0
\[
\qquad \mathrm{A4}:\ \alpha=0\leftrightarrow\alpha=\pi\ \text{exchanges the branch labels}
\]

%dsp-level:0
\[
\mathcal P=\{e:\alpha_e=\pi\}\ [\mathrm{Def\text{-}EM.21}];
\]

%dsp-level:0
\[
\qquad \alpha_e\mapsto\pi-\alpha_e
\]

%dsp-level:0
\[
\iff e\ \text{changes membership in}\ \mathcal P\iff\mathcal P
\]

%dsp-level:0
\[
\mapsto E_{\mathcal R}\setminus\mathcal P\ [\mathrm{EM.24}(2);\ \text{exhaustive},\ |E_{\mathcal R}|\le5;\ \text{machine}]
\]

%dsp-level:0
\[
\neg\exists M:\ \mathfrak T(MXM^{-1})=\mathfrak TX\ \wedge\ \operatorname{rk}(MXM^{-1})=\operatorname{rk}X\ [\text{conjugation invariants}];
\]

%dsp-level:0
\[
\qquad 3\ne1\implies\Pi_{\mathcal R}\not\sim\Pi_{\mathcal V}\ [\mathrm{CT0.7}(1);\ \text{exact}]
\]

%dsp-level:0
\[
(4):\ \text{support identity}\ (1);
\]

%dsp-level:0
\[
\qquad \text{flip image}\ (2),(3);
\]

%dsp-level:0
\[
\qquad \text{sharp loci}\ [\mathrm{CT0.8}];
\]

%dsp-level:0
\[
\qquad \text{witness states}\ \Pi_{\mathcal R}/3,\ \Pi_{\mathcal V}:\ \text{exact}\ [\text{machine}]
\]

%dsp-level:0
\[
(5):\ (1)\text{--}(4)\ \text{conjoined};
\]

%dsp-level:0
\[
\qquad \text{scope clause cites FT.93}(3),(5),(7)\ \text{unmodified}\ [\text{no carrier grading is constructed here}].\ \blacksquare
\]

\textbf{Phys.} The two crystallization outcomes are inequivalent for the weak force: the aligned (triplet) branch is where the total-spin generators act — Casimir $2$ — and the antialigned (singlet) branch is annihilated by them — Casimir $0$. The flip of the binary branch label exchanges the two, so weak activity is not flip-invariant: a structural $\mathbb Z_2$ asymmetry carried by the registration alphabet of the real edges. This is the force-arena $S$-parity grading read against the ledger's own branch data — not a handedness of matter: the flip has no Hilbert-space implementation ($3$ vs $1$), and no chirality projector on a matter carrier is constructed or implied. Barred: naming either branch "left" or "right"; any SM species reading; any claim that the SM weak interaction is reproduced.

\vspace{2pt}\noindent\rule{\textwidth}{0.4pt}\vspace{2pt}

\end{theorem}

\phantomsection
\section[Tier F.35 — The 1-Loop Bridge: the Swap-Graded Effective…]{\texorpdfstring{Tier F.35 — The 1-Loop Bridge: the Swap-Graded Effective Potential on the Flicker Phase Diagram}{Tier F.35 - The 1-Loop Bridge: the Swap-Graded Effective Potential on the Flicker Phase Diagram}}

\begin{definition}[{\texorpdfstring{Def-H.31 — The bridge: dictionary, family, background, fluctuation operator, effective potential}{Def-H.31 - The bridge: dictionary, family, background, fluctuation operator, effective potential}}]\label{tag:Def-H-31}
\textbf{Def.} A construction class on the prior canon (Def-H.24 pattern: no new axiom); clauses (1)–(6) are constitutive interface definitions. Domain: $\varepsilon_q\in(0,\tfrac14)$ [CT0.8]; $\ \gamma:=1-4\varepsilon_q$, $\ q:=4\varepsilon_q$ [D0.3, D0.4].

$(1)$ \textbf{Dictionary (Def-H.31(1)).} $\ p:=\gamma=1-4\varepsilon_q;\qquad 1-p=q=4\varepsilon_q;\qquad s:=p/q;\qquad y:=16s^2=\dfrac{(1-4\varepsilon_q)^2}{\varepsilon_q^2}.$ Anchors (exact): $\varepsilon_q=\tfrac18\iff p=\tfrac12\iff y=16$ — the point at which FT.68(1)'s amplitude-weighted reading is the uniform measure and $\sum|A_C|^2=15745$ on the tetrahedron [CT2.8 norm layer: $|A_C|^2=16^{|\mathcal C|}$; FT.67(2)'s value formula]; $\varepsilon_q\to0\iff p\to1$ — the static-floor corner [FT.68(5); D0.4's sharp point $\gamma=1$]; $\varepsilon_q\to\tfrac14\iff p\to0$ — the identity-tilt, vacuum-dominated corner [L2 ($\gamma\to0$); FT.70(5) ($\langle K\rangle\to0$)].

$(2)$ \textbf{Family and weights (Def-H.31(2)).} $\ \mathbb F:=\{\text{tetra},\text{cube},\text{octa}\}$ — the skeletons whose two-sector covering polynomials are complete [FT.68(1), FT.68(3), FT.69(2)]:

%dsp-level:0
\[
\mathcal Z^{\mathrm{tot}}_{\mathrm{tetra}}(y)=1+24y+72y^2+4y^3;\qquad
\mathcal Z^{\mathrm{tot}}_{\mathrm{cube}}(y)=1+48y+672y^2+3400y^3+6816y^4+264y^5+4y^6;
\]

%dsp-level:0
\[
\mathcal Z^{\mathrm{tot}}_{\mathrm{octa}}(y)=1+72y+888y^2+2160y^3+840y^4
\qquad(\mathcal Z^{\mathrm{tot}}=\mathcal Z_{\mathrm{byp}}+\mathcal Z_{\mathrm{lock}},\ \text{rows pooled}).
\]

With $2E_\Gamma\in(12,24,24)$ the directed-edge counts and $K=2|\mathcal C|$ [FT.79(2)], the \textbf{cover-weighted Bernoulli weight} of $\Gamma$ is

%dsp-level:0
\[
W_\Gamma(p)\ :
\]

%dsp-level:0
\[
= \ \sum_k c_k(\Gamma)\,
\]

%dsp-level:0
\[
16^k\,
\]

%dsp-level:0
\[
p^{2k}\,
\]

%dsp-level:0
\[
q^{2E_\Gamma-2k}\
\]

%dsp-level:0
\[
= \ q^{2E_\Gamma}\,
\]

%dsp-level:0
\[
\mathcal Z^{\mathrm{tot}}_\Gamma\bigl(16p^2/q^2\bigr)\
\]

%dsp-level:0
\[
= \ \mu_p\text{-average of }|A_C|^2\text{ over covers }C\subseteq\Gamma,
\]

an integer polynomial in $p$ with $W_\Gamma(0)=1$ (the vacuum cover) and $W_\Gamma(1)=0$ (the fully crystallized floor is no cover — it cancels [FT.67(3a)]). The \textbf{family measure} is

%dsp-level:0
\[
P(\Gamma\mid\varepsilon_q)\ :=\ \frac{W_\Gamma(p)}{\sum_{\Gamma'\in\mathbb F}W_{\Gamma'}(p)}
\]

%dsp-level:0
\[
\qquad (\text{ensemble scope of CT4.5a/Def-H.26: per-edge Bernoulli }p).
\]

$(3)$ \textbf{Background state.} Cover-conditioned edge content and mean real fraction:

%dsp-level:0
\[
\langle K\rangle^{\mathrm{cov}}_\Gamma(p):=\frac{N_\Gamma(p)}{W_\Gamma(p)},\quad N_\Gamma:=\sum_k 2k\,c_k 16^k p^{2k}q^{2E-2k};
\]

%dsp-level:0
\[
\qquad \hat p_\Gamma:=\frac{\langle K\rangle^{\mathrm{cov}}_\Gamma}{2E_\Gamma}\in(0,\ K_{\mathrm{top}}/2E_\Gamma].
\]

The ensemble-averaged pair state (linear in the pattern measure; per-pattern state $=$ the branch state of its real/virtual edges, Cal-F detailed balance):

%dsp-level:0
\[
\bar\rho_\Gamma(\varepsilon_q)\ :=\ \hat p_\Gamma\cdot\frac{\Pi_{\mathcal R}}{3}\ +\ (1-\hat p_\Gamma)\,\Pi_{\mathcal V}
\]

%dsp-level:0
\[
\qquad [\mathrm{CT0.7}(1):\ \Pi_{\mathcal R}=\tfrac14(3E_0+\Sigma),\ \Pi_{\mathcal V}=\tfrac14(E_0-\Sigma)];
\]

$\operatorname{spec}(\bar\rho_\Gamma)=\bigl(\tfrac{\hat p_\Gamma}3,\tfrac{\hat p_\Gamma}3,\tfrac{\hat p_\Gamma}3,\ 1-\hat p_\Gamma\bigr)$ — faithful on the whole domain; $\mathfrak T\bar\rho_\Gamma=1$.

$(4)$ \textbf{Fluctuation operator (Def-H.31(4)).} $\Delta^{-1}_\Gamma(\varepsilon_q)$: the block-diagonal operator on the $16$-mode space $\mathfrak H(\mathbb C^4)$ [CT0.6] in the $\vee^2/\wedge^2$ decomposition of the swap grading [CT0.7(3)], with the A8 aspect labels of CT0.7(4):

%dsp-level:0
\[
\textbf{S-even block}\ (10=9+1=\dim\mathfrak H(\vee^2)+\dim\mathfrak H(\wedge^2);\ \text{aspects }F_{D^+},F_H,F_G,F_{W^\pm},F_Z):
\]

the Bogoliubov–Kubo–Mori Hessian of the relative entropy $S(\rho\|\bar\rho_\Gamma)$ at $\bar\rho_\Gamma$ [math: the standard quadratic fluctuation kernel of the state-space information geometry; eigenvalue $1/\lambda$ on diagonal block directions]:

%dsp-level:0
\[
M^2_{\vee^2}=3/\hat p_\Gamma\ (\times 9),\qquad M^2_{\wedge^2}=1/(1-\hat p_\Gamma)\ (\times 1);
\]

%dsp-level:0
\[
\textbf{S-odd block}\ (6=3+3;\ \text{the seal forces }F_{J^+}=\operatorname{span}\{J_i\},\ F_{J^-}=\operatorname{span}\{K_i\}\ [\mathrm{CT0.7}(4)]):
\]

%dsp-level:0
\[
M^2_{J^+}=p\cdot G_\Gamma(p)^2\ (\times 3),\qquad M^2_{J^-}=q\cdot G_\Gamma(p)^2\ (\times 3),\qquad G_\Gamma:=B_\Gamma-p^{2E_\Gamma}\ \ [\mathrm{Def\text{-}H.26},\ \mathrm{FT.70}(1)],
\]

with the detailed-balance occupations $(p,q)$ of the two seal channels [Cal-F; A8: $\mathbf F_{(+,\mathcal V)}=F_{J^+}$, $\mathbf F_{(-,\mathcal V)}=F_{J^-}$]. $B_\Gamma(p)=\sum_K c_Kp^Kq^{2E-K}$ with the frozen cap-census rows $c_K$ [FT.68(1), FT.68(3), FT.69(1)].

$(5)$ \textbf{Effective potential.} With the \textbf{swap supertrace} $\operatorname{Str}_{\mathbb S}X:=\mathfrak T(\mathbb S X)$ (defined here; $\ne$ the C-sector supertrace [CT2.10] — collision-avoidance declared once; $\operatorname{Str}_{\mathbb S}$ grades $\mathfrak H(\mathbb C^4)$ by CT0.7(4)'s $S$-parity, $10-6=4$):

%dsp-level:0
\[
V_{\text{1-loop},\Gamma}\ :
\]

%dsp-level:0
\[
= \ \tfrac12\operatorname{Str}_{\mathbb S}\log\Delta^{-1}_\Gamma\
\]

%dsp-level:0
\[
= \ \tfrac12\Bigl(\sum_{\text{S-even}}\log M^2_i-\sum_{\text{S-odd}}\log M^2_j\Bigr)
\
\]

%dsp-level:0
\[
= \ \tfrac92\log\tfrac3{\hat p_\Gamma}-\tfrac12\log(1-\hat p_\Gamma)-\tfrac32\log(pq)-6\log|G_\Gamma|;
\]

%dsp-level:0
\[
V_{\text{tree},\Gamma}\ :=\ -\log W_\Gamma(p);
\]

%dsp-level:0
\[
\qquad V_{\mathrm{eff}}(\varepsilon_q)\ :=\ \sum_{\Gamma\in\mathbb F}P(\Gamma\mid\varepsilon_q)\,\bigl[V_{\text{tree},\Gamma}(\varepsilon_q)+V_{\text{1-loop},\Gamma}(\varepsilon_q)\bigr].
\]

Every log argument is positive on $(0,\tfrac14)\setminus\{\varepsilon^*_\Gamma\}$, $\varepsilon^*_\Gamma:=(1-p^*_\Gamma)/4$; $V_{\mathrm{eff}}$ is smooth on the wall complement and $V_{\mathrm{eff}}\to+\infty$ at both endpoints and at every wall (FT.102(4)).

$(6)$ \textbf{Mass read (Def-H.31(6)).} For each vacuum $\varepsilon_q^{(i)}$ of $V_{\mathrm{eff}}$ (FT.103):

%dsp-level:0
\[
m_i^2\ :=\ \Delta S\bigl(\varepsilon_q^{(i)}\bigr)\ :=\ \frac{V_{\mathrm{eff}}\bigl(\varepsilon_q^{(i)}\bigr)}{8}
\]

%dsp-level:0
\[
\qquad [\mathrm{L3}(2):\ V=8\,\Delta S_{\mathrm{pair}}\ \text{exact leading};\ \mathrm{L3}(4):\ m_i^2=\Delta S\ \text{inertial read};\ \text{dimensionless, nats}].
\]

The phase-band $\leftrightarrow$ generation-label correspondence is a labelling clause of this definition, not a derived selection; no particle identification; no coupling to $m_{\mathrm{pole}}$/$m_{\mathrm{hid}}$ beyond Def-H.39's bridge slots (unconstructed).

\textbf{Rmk (what is new, what is frozen).} All numerical coefficients are prior-canon data: the covering rows $(1,24,60\,|\,12,4)$, $(1,48,648,2672,3624\,|\,24,728,3192,264,4)$, $(1,72,852,1152\,|\,36,1008,840)$ [FT.68, FT.69]; the cap-census rows $c_K$ [FT.68(1),(3), FT.69(1)]; $16=4^2$ [FT.67(2) value $\times$ CT2.8 norm]; $K=2|\mathcal C|$ [FT.79(2)]; the $10/6$ split [CT0.7(3)-(4)]; $\tfrac34$ and $\gamma(8-\gamma)$-family coefficients [L2, L5]; the factor $8$ [L3(2)]; the domain $[0,\tfrac14)$ [CT0.8]. New content: the assembly (Def-H.31(1)–e) and everything proved from it.

\vspace{2pt}\noindent\rule{\textwidth}{0.4pt}\vspace{2pt}

\end{definition}

\begin{theorem}[{\texorpdfstring{FT.102 — The fluctuation spectrum, the tree-level limit, and the seal-softening walls}{FT.102 - The fluctuation spectrum, the tree-level limit, and the seal-softening walls}}]\label{tag:FT-102}
\textbf{Thm.} Under Def-H.31: $(1)$ \textbf{Exact spectrum.} For every $\Gamma\in\mathbb F$ and $\varepsilon_q\in(0,\tfrac14)\setminus\{\varepsilon^*_\Gamma\}$:

%dsp-level:0
\[
\operatorname{Spec}\bigl(\Delta^{-1}_\Gamma(\varepsilon_q)\bigr)\
\]

%dsp-level:0
\[
= \ \underbrace{\Bigl\{\tfrac3{\hat p_\Gamma}\times9,\ \ \tfrac1{1-\hat p_\Gamma}\times1\Bigr\}}_{\text{S-even},\ 10}\ \sqcup\ \underbrace{\bigl\{p\,G_\Gamma(p)^2\times3,\ \ q\,G_\Gamma(p)^2\times3\bigr\}}_{\text{S-odd},\ 6}
\]

%dsp-level:0
\[
\qquad (16\ \text{values}),
\]

with $\hat p_\Gamma\in(0,K_{\mathrm{top}}/2E_\Gamma]$, $K_{\mathrm{top}}/2E_\Gamma=(\tfrac12,\tfrac12,\tfrac13)$ for (tetra, cube, octa) — the supremum attained exactly in the limit $p\to1$ (vanishing order $2E-K_{\mathrm{top}}+2$ of $N_\Gamma-K_{\mathrm{top}}W_\Gamma$ at $p=1$; machine). The S-even block is strictly positive on the whole domain; the S-odd block vanishes exactly at $p\in\{p^*_\Gamma\}\cup\{0,1\}$ and nowhere else. $(2)$ \textbf{The grading is matrix-exact.} On the CT1.2 basis over $\mathbb Q(i)$: $\operatorname{Ad}(\mathbb S)$ has eigenvalue $+1$ on $\{E_0,\Sigma,\hat E^{ij}\,(5),W_i\,(3)\}$ — $10$ modes — and $-1$ on $\{J_i\}\cup\{K_i\}$ — $6$ modes [CT0.7(4) membership reproduced from the matrices]; S-even $\iff[\mathbb S,X]=0\iff X$ preserves $\vee^2,\wedge^2$; S-odd $\iff\{\mathbb S,X\}=0\iff X$ swaps them [CT0.7(3)]; basis rank $16$ [CT0.6]; $\operatorname{Str}_{\mathbb S}(I_{16})=10-6=4$; on any degenerate spectrum $\lambda I$: $V_{\text{1-loop}}=\tfrac12(10-6)\log\lambda=2\log\lambda$ — the super-dimension is the interference coefficient the trace log carries. $(3)$ \textbf{Tree-level limit $=$ the FT.68/FT.69 covering polynomials.} Setting $V_{\text{1-loop}}\equiv0$: $e^{-V_{\text{tree},\Gamma}(p)}=W_\Gamma(p)=q^{2E}\mathcal Z^{\mathrm{tot}}_\Gamma(y)$ \textit{identically} — verified two ways: the polynomial expansion vs the sum over the \textit{enumerated} covers (two independent enumerators per cell, all covers: $101/11205/3961$). Static evaluation anchors (exact integers):

%dsp-level:0
\[
\begin{array}{c|ccc}
&\text{tetra}&\text{cube}&\text{octa}\\ \hline
\mathcal Z^{\mathrm{tot}}(4)\ (Z^\pm_{\mathrm{all}}) & 1505 & 2260161 & 367777\\
\mathcal Z_{\mathrm{byp}}(4)\ /\ \mathcal Z_{\mathrm{lock}}(4) & 1057\ /\ 448 & 1109313\ /\ 1150848 & 87649\ /\ 280128\\
W_\Gamma(\tfrac12)\cdot2^{2E}=\mathcal Z^{\mathrm{tot}}(16) & 35201 & 804725505 & 64126081\\
\mathcal Z_{\mathrm{byp}}(16)\ /\ \mathcal Z_{\mathrm{lock}}(16) & 15745\ /\ 19456 & 248613633\ /\ 556111872 & 4937857\ /\ 59188224
\end{array}
\]

the tetrahedron's bypass norm sum $15745$ and amplitude-weighted content $\langle K\rangle=62208/15745$ are FT.68(1)(ii) verbatim; the uniform readings $852/205$, $3311952/395023$, $2124600/318523$ at $p=\tfrac12$ are FT.68(1)(i)/FT.70(5) verbatim; the single-centre closed form $c_1=Vd(d-1)=(24,48,72)$ is FT.69(2) verbatim. $(4)$ \textbf{Seal-softening walls.} $V_{\mathrm{eff}}\to+\infty$ at each $\varepsilon^*_\Gamma=(1-p^*_\Gamma)/4$, $p^*_\Gamma$ the unique zero of $G_\Gamma$ [FT.70(1)–(2)] — dyadic isolation at width $2^{-40}$ reproduces the frozen boundaries exactly: $p^*=(0.725372,\ 0.723525,\ 0.679217)$ for (cube, tetra, octa), i.e. $\varepsilon^*=(0.068657016,\ 0.069118773,\ 0.080195792)$, order octa $>$ tetra $>$ cube in $\varepsilon^*$ [FT.70(2)'s order in $p^*$ inverted by the dictionary]. At a wall the six S-odd masses vanish quadratically — the phase boundary is exactly where the reality-exchanging (seal) fluctuations become marginal. $V_{\mathrm{eff}}\to+\infty$ also at both domain endpoints: the floor side ($\varepsilon_q\to0$) from $W_\Gamma\to0$ (the floor cancels [FT.67(3a)]) and $-\tfrac32\log q\to+\infty$; the threshold side ($\varepsilon_q\to\tfrac14$) from $3/\hat p_\Gamma\to\infty$ and $-\tfrac32\log p\to+\infty$. $(5)$ \textbf{Directive form.} $V_{\text{1-loop},\Gamma}=\tfrac12\bigl(\sum_{\text{S-even}}\log M_i^2-\sum_{\text{S-odd}}\log M_j^2\bigr)$ — the block assembly and the $16$-term eigenvalue assembly agree by enclosures; the minus sign on the $6$ odd terms is the swap-supertrace grading, and the odd block is precisely the seal force pair $F_{J^\pm}$ [CT0.7(4): "the $S$-odd blocks"; D1.2: the seals move A4-weight $\wedge^2\!\leftrightarrow\!\vee^2$].

\textbf{Req.} Def-H.31; A4, A8; D0.3, D0.4, D1.2; CT0.6, CT0.7, CT0.7a, CT0.8, CT0.11, CT1.1, CT1.2 (Part 1); EW.12 (Part 2); CT2.8 (Part 2); CT4.5a (Part 2); Def-H.25, Def-H.26; FT.67, FT.68, FT.69, FT.70; FT.79(2); FT.93(5); L2, L3, L5; Cal-F; exact $\mathbb Q(i)$ matrix arithmetic, two-enumerator cover re-derivation, integer-Horner dyadic bisection, rational log-enclosures [math; machine: \texttt{checks/\allowbreak{}verify\_\allowbreak{}H\_\allowbreak{}1loop\_\allowbreak{}bridge.py} §A–§D].

\textbf{Pf.}

%dsp-level:0
\[
(1):\ \bar\rho_\Gamma\ \text{faithful}\ (\hat p_\Gamma\in(0,\tfrac12])\implies\text{BKM diagonal coefficients }1/\lambda:\ 1/(\hat p/3)
\]

%dsp-level:0
\[
= 3/\hat p\ (\times9),
\]

%dsp-level:0
\[
\ 1/(1-\hat p)\ (\times1)\ [\text{math}]
\]

%dsp-level:0
\[
\hat p\text{-bounds}:\ N_\Gamma/W_\Gamma\le K_{\mathrm{top}}\ \text{termwise (each cover }K\le K_{\mathrm{top}});
\]

%dsp-level:0
\[
\qquad \hat p(p\to1)
\]

%dsp-level:0
\[
= K_{\mathrm{top}}/2E_\Gamma:\ \text{vanishing-order census of }N_\Gamma-K_{\mathrm{top}}W_\Gamma\ \text{at }p
\]

%dsp-level:0
\[
= 1\ [\text{machine: exact derivative ladder}]
\]

%dsp-level:0
\[
\text{odd zeros}:\ pG^2=0\iff p=0\vee G=0;\ qG^2=0\iff p=1\vee G=0;
\]

%dsp-level:0
\[
\qquad G_\Gamma\ \text{has exactly one root on the resolved grid}\ [\mathrm{FT.70}(1);\ \text{machine}]
\]

%dsp-level:0
\[
(2):\ \mathbb S=\tfrac12(E_0+\Sigma)\ \text{as matrices}\ [\mathrm{CT0.7}(1)];
\]

%dsp-level:0
\[
\qquad \operatorname{Ad}(\mathbb S)X\ \text{computed on all }16\ \text{basis elements over }\mathbb Q(i)\ [\text{machine}]
\]

%dsp-level:0
\[
\text{membership}=\mathrm{CT0.7}(4)\text{'s frozen lists};
\]

%dsp-level:0
\[
\qquad \operatorname{Str}_{\mathbb S}(I_{16})=\#\{+1\}-\#\{-1\}=10-6=4;
\]

%dsp-level:0
\[
\qquad \tfrac12\operatorname{Str}_{\mathbb S}\log(\lambda I)=\tfrac{10-6}2\log\lambda=2\log\lambda
\]

%dsp-level:0
\[
(3):\ W_\Gamma(p)
\]

%dsp-level:0
\[
= \sum_{\text{covers }C}16^{|C|}p^{2|C|}q^{2E-2|C|}:\ \text{polynomial expansion}
\]

%dsp-level:0
\[
= \text{enumerated sum at }p
\]

%dsp-level:0
\[
= \tfrac37\ [\text{machine, both enumerators}];
\]

%dsp-level:0
\[
\quad \text{anchors: exact integer evaluations}
\]

%dsp-level:0
\[
15745=\mathcal Z_{\mathrm{byp}}(16)=1+24\cdot16+60\cdot256;
\]

%dsp-level:0
\[
\qquad 62208=\sum_k2k\,c_k16^k\big|_{\mathrm{byp}}=2\cdot24\cdot16+4\cdot60\cdot256\ [\mathrm{FT.68}(1)(ii)]
\]

%dsp-level:0
\[
(4):\ G_\Gamma(0)=1>0>G_\Gamma(1)=-1\ [\mathrm{FT.70}(1)];
\]

%dsp-level:0
\[
\qquad \text{bisection at }2^{-40}:\ \text{intervals reproduce the frozen decimals}\ [\text{machine}]
\]

%dsp-level:0
\[
-6\log|G_\Gamma|\to+\infty\ \text{at }G_\Gamma=0;
\]

%dsp-level:0
\[
\qquad W_\Gamma(1)=0\ (\text{every term carries }q^{2E-2k},\ 2E-2k\ge2E-K_{\mathrm{top}}>0);
\]

%dsp-level:0
\[
\qquad \hat p\to0\implies3/\hat p\to\infty\ [\varepsilon_q\to\tfrac14]
\]

%dsp-level:0
\[
(5):\ \text{two assemblies of the same }16\text{-term sum, evaluated by independent code paths at }\varepsilon_q
\]

%dsp-level:0
\[
= \tfrac1{16}:\ \text{difference enclosure brackets }0\ [\text{machine}].\ \blacksquare
\]

\textbf{Phys.} The tree layer knows only polynomial censuses; the loop layer turns the census into geometry. Ten reality-preserving modes carry the information metric of the averaged crystallization state; six reality-exchanging seal modes carry the phase balance — and the walls of the effective potential sit exactly at the phase diagram's crossover boundaries, because a seal fluctuation has no restoring force where thread-gas and floor weigh equal. The grade difference $10-6=4$ — the swap super-dimension of the pair arena — is the coefficient the trace log carries into every degenerate limit. Structure class of a 1-loop effective potential (trace log of the fluctuation operator; Coleman–Weinberg species [physics vocabulary — analogy only; no SM field content, no identification]).

\vspace{2pt}\noindent\rule{\textwidth}{0.4pt}\vspace{2pt}

\end{theorem}

\begin{theorem}[{\texorpdfstring{FT.103 — The vacuum structure: four exact vacua, three wall-delimited phase bands, dominance windows, the susceptibility census}{FT.103 - The vacuum structure: four exact vacua, three wall-delimited phase bands, dominance windows, the susceptibility census}}]\label{tag:FT-103}
\textbf{Thm.} Under Def-H.31: $(1)$ \textbf{Vacuum census (complete at the resolved scale).} $dV_{\mathrm{eff}}/d\varepsilon_q=0$ has exactly four solutions in $(0,\tfrac14)$ on the wall-delimited bands — one per band, each a strict minimum ($V'$ passes $-\to+$; exact sign census by rational $+$ log-enclosure evaluation on dyadic grids; bisection to width $2^{-36}$):

%dsp-level:0
\[
\begin{array}{c|c|c|c|c}
&\text{band}&\varepsilon_q^{(i)}\ (\approx)&V_{\mathrm{eff}}\ (\approx)&m_i^2=V_{\mathrm{eff}}/8\ (\approx)\\ \hline
\mathrm V_1&(0,\varepsilon^*_{\mathrm{cube}})\ \text{floor band}&0.021337258636&24.271850813&3.033981352\\
\mathrm V_2&(\varepsilon^*_{\mathrm{cube}},\varepsilon^*_{\mathrm{tetra}})\ \text{bicritical sliver}&0.068783298916&59.645122359&7.455640295\\
\mathrm V_3&(\varepsilon^*_{\mathrm{tetra}},\varepsilon^*_{\mathrm{octa}})\ \text{mixed band}&0.080091490970&38.713071069&4.839133884\\
\mathrm V_4&(\varepsilon^*_{\mathrm{octa}},\tfrac14)\ \text{thread-gas band}&0.190792364908&10.655753803&1.331969225
\end{array}
\]

All four values are strictly positive — $\Delta S$-admissible in nats [L3]; all four positions lie in CT0.8's physical domain $(0,\tfrac14)$. $\mathrm V_3$ sits $1.04\times10^{-4}$ below the octahedron wall (wall-adjacent: its depth is set by the seal-softening balance); $\mathrm V_2$ is the sliver bottom between the near-degenerate cube/tetrahedron boundaries — its existence and depth are controlled by the frozen splitting $|\varepsilon^*_{\mathrm{tetra}}-\varepsilon^*_{\mathrm{cube}}|=4.6176\times10^{-4}$ [FT.70(2)'s own numbers]. $(2)$ \textbf{Mass order and ratios (enclosure-exact).} $m^2(\mathrm V_4)<m^2(\mathrm V_1)<m^2(\mathrm V_3)<m^2(\mathrm V_2)$ (disjoint enclosures). The phase triple $(\mathrm V_4,\mathrm V_1,\mathrm V_3)$ in mass order:

%dsp-level:0
\[
m^2:\ 1:2.277816:3.633067;\qquad m:\ 1:1.509244:1.906061;
\]

all four: $m^2:\ 1:2.277816:3.633067:5.597457$ (bare $m^2$-span $<6$ — clause (6)). $(3)$ \textbf{Dominance windows (exact).} The pairwise $W$-crossings isolate (width $2^{-30}$): tetra$|$cube at $p=0.662192$ and $p=0.035764$ ($\varepsilon_q=0.084410$ and $0.241059$); cube$|$octa at $p=0.091544$ ($\varepsilon_q=0.227062$). Windows: tetra-dominant $(0,0.084410)\cup(0.241059,\tfrac14)$; cube-dominant $(0.084410,0.227062)$; octa-dominant $(0.227062,0.241059)$. Dominant skeleton at the vacua: $\mathrm V_1,\mathrm V_2,\mathrm V_3$ tetrahedron ($P_{\mathrm{tetra}}\approx0.999/0.720/0.556$); $\mathrm V_4$ cube ($P_{\mathrm{cube}}\approx0.727$). The octa window and the tetra tail sliver carry \textbf{no vacuum}: $V'>0$ throughout both (machine census) — the octahedron enters the vacuum structure through its wall $\varepsilon^*_{\mathrm{octa}}$, not through a dominance minimum. $(4)$ \textbf{Susceptibility census.} Per skeleton, $\chi_\Gamma(y):=d\langle K\rangle^{\mathrm{cov}}_\Gamma/dy=2\operatorname{Var}_y(|\mathcal C|)/y$ [exact identity; math] is \textbf{strictly decreasing} on the exact grid $y=2^{i/4}$, $|i|\le40$ — no interior peaks: the van-Hove structure sits in $\operatorname{Var}_y(|\mathcal C|)$ (the log-fugacity susceptibility $d\langle K\rangle/d\ln y=2\operatorname{Var}$). The cover-condensation transition scales $y=c_k/c_{k+1}$ and their exact $\varepsilon$-images ($\varepsilon=1/(4+\sqrt y)$; each the unique root in $(0,\tfrac14)$ of the integer quadratic $(16d-n)\varepsilon^2-8d\varepsilon+d=0$, $y=n/d$; machine):

%dsp-level:0
\[
\begin{array}{c|c|c}
\Gamma&y=c_k/c_{k+1}&\varepsilon_q=1/(4+\sqrt y)\ (\approx)\\ \hline
\text{tetra}&\tfrac1{24},\ \tfrac13,\ 18&0.237862,\ 0.218467,\ 0.121320\\
\text{cube}&\tfrac1{48},\ \tfrac1{14},\ \tfrac{84}{425},\ \tfrac{425}{852},\ \tfrac{284}{11},\ 66&0.241293,\ 0.234342,\ 0.224993,\ 0.212482,\ 0.110118,\ 0.082481\\
\text{octa}&\tfrac1{72},\ \tfrac3{37},\ \tfrac{37}{90},\ \tfrac{18}7&0.242845,\ 0.233386,\ 0.215462,\ 0.178458
\end{array}
\]

The family trajectory $\chi_{\mathbb F}(\varepsilon_q)=d\langle K\rangle_{\mathrm{fam}}/dy$ (with $\langle K\rangle_{\mathrm{fam}}=\sum_\Gamma P_\Gamma\langle K\rangle^{\mathrm{cov}}_\Gamma$, rational in $p$) has exactly \textbf{two} interior extrema on the resolved $2^{-13}$ grid (degree-$96$ integer-Horner census): $\varepsilon_q\approx0.019724$ (adjacent to $\mathrm V_1$) and $\varepsilon_q\approx0.083699$ (adjacent to the tetra$|$cube dominance boundary $0.084410$) — a census, not an identity: the extrema do not coincide with the vacua. $(5)$ \textbf{Phase-band reading (Def-H.31(6)).} The three wall-delimited bands that carry vacua — floor band ($\mathrm V_1$; static-crystallization side), crossover band ($\mathrm V_2,\mathrm V_3$; sliver $+$ mixed), thread-gas band ($\mathrm V_4$) — are the three topological phases of the construction; the labelling of the phase triple as three generations is declared, not derived. The band count equals $|\mathbb F|$, the declared family size: the selection principle for \textit{three} is not supplied; the sliver vacuum $\mathrm V_2$ is a bicritical artifact of FT.70(2)'s near-degenerate boundaries — reported, not suppressed. $(6)$ \textbf{Bare-span bound on the computed cells.} The all-four $m^2$-span is $5.597457<6$ — the same order as FT.79(1)'s pool maximum $6$ and FT.76(3)'s widest ratio $6.3925$: a third independent appearance of the $\approx6$ scale on the computed cells (FT.104(1) quantifies the [EXT] consequence).

\textbf{Req.} Def-H.31; FT.102; CT0.8 (Part 1); Def-H.26; FT.67, FT.68, FT.69, FT.70; FT.76(3), FT.79(1)–(2); L3; exact rational $+$ log-enclosure bisection, integer-Horner dyadic sign census, quadratic-surd verification, variance identities [math; machine: \texttt{checks/\allowbreak{}verify\_\allowbreak{}H\_\allowbreak{}1loop\_\allowbreak{}bridge.py} §E–§F].

\textbf{Pf.}

%dsp-level:0
\[
(1):\ V'
\]

%dsp-level:0
\[
= \text{rational}(\varepsilon_q)+\sum_j\text{rational}_j(\varepsilon_q)\cdot\ln\bigl(\text{rational}_j(\varepsilon_q)\bigr):\ \text{each }\ln\text{ by rational atanh-enclosure (exact tails)}
\]

%dsp-level:0
\[
\text{band grids (dyadic, wall-graded near }\varepsilon^*_{\mathrm{octa}}\text{)}:\ \text{exactly one }-\to+\text{ sign flip per band}\ [\text{machine}];
\]

%dsp-level:0
\[
\qquad \text{bisection to }2^{-36};
\]

%dsp-level:0
\[
\qquad V(\varepsilon^{(i)})>0\ \text{by enclosure}
\]

%dsp-level:0
\[
V''\text{-flatness at the minimum}:\ |V(\text{mid})-V(\varepsilon^{(i)})|\le\tfrac12\sup|V''|\cdot(2^{-37})^2\ll2^{-44}:\ \text{enclosure values valid at width}
\]

%dsp-level:0
\[
\le 2^{-44}\ [\text{machine}]
\]

%dsp-level:0
\[
(2):\ \text{interval division of disjoint positive enclosures};
\]

%dsp-level:0
\[
\qquad m=\sqrt{m^2}:\ \text{dyadic }\sqrt{\ }\text{-enclosures (integer squaring)}
\]

%dsp-level:0
\[
(3):\ W_\Gamma-W_{\Gamma'}\ \text{integer polynomials}:\ \text{root census }2/1/2\ \text{on the resolved grid};
\]

%dsp-level:0
\[
\qquad \text{dominance at vacua}: \text{exact rational comparison of }W_\Gamma(\text{mid})
\]

%dsp-level:0
\[
V'>0\ \text{on both no-vacuum windows}:\ \text{enclosure census at }41\ \text{dyadic points each}\ [\text{machine}]
\]

%dsp-level:0
\[
(4):\ \langle k\rangle_y=y\mathcal Z'/\mathcal Z\implies y\,d\langle k\rangle/dy=\operatorname{Var}_y(k)\ [\text{math}];
\]

%dsp-level:0
\[
\qquad \chi_\Gamma=2\operatorname{Var}/y:\ \text{strict decrease on the exact grid}\ [\text{machine}]
\]

%dsp-level:0
\[
y = (1-4\varepsilon)^2/\varepsilon^2
\]

%dsp-level:0
\[
\iff y\varepsilon^2
\]

%dsp-level:0
\[
= (1-4\varepsilon)^2
\]

%dsp-level:0
\[
\iff (16-y)\varepsilon^2-8\varepsilon+1
\]

%dsp-level:0
\[
= 0
\]

%dsp-level:0
\[
\iff (16d-n)\varepsilon^2-8d\varepsilon+d
\]

%dsp-level:0
\[
= 0\ (y=n/d):\ \text{unique root in }(0,\tfrac14)\ [\text{machine}]
\]

%dsp-level:0
\[
\chi_{\mathbb F}\ \text{extrema}:\ \text{numerator of }d\chi_{\mathbb F}/dp\ \text{is an integer polynomial of degree }96:\ \text{sign census at }p
\]

%dsp-level:0
\[
= m/2^{13}\ \text{by integer Horner}:\ 2\ \text{flips}\ [\text{machine}].\ \blacksquare
\]

\textbf{Phys.} The effective potential has four exact vacua, bracketed by walls that \textit{are} the phase diagram: every crossover boundary FT.70 computed is now a divergent barrier at which the six seal modes go soft, and every vacuum is a phase — the crystallization-side well, the thread-gas well, and two crossover-band wells whose depths are set by how close the Platonic boundaries sit to one another. The lightest vacuum is the thread-gas phase (cube-dominant); the heaviest is the bicritical sliver between the near-degenerate cube/tetrahedron boundaries. The mass read is the entropy-deficit read in nats [L3(4)]: no pole identification is made or implied relating-equation residue applies unchanged (FT.104). No particle identification.

\vspace{2pt}\noindent\rule{\textwidth}{0.4pt}\vspace{2pt}

\end{theorem}

\begin{theorem}[{\texorpdfstring{FT.104 — The [EXT] audit, the range-species classification, and the fourth read of Gap 7}{FT.104 - The [EXT] audit, the range-species classification, and the fourth read of Gap 7}}]\label{tag:FT-104}
\textbf{Thm.} \textbf{External-data clause} (FT.73(5)/FT.79(4) pattern): the hierarchies compared below — charged leptons $1{:}207{:}3477$, down-type $1{:}20{:}900$, up-type $1{:}600{:}80000$ — are \textit{supplied measured data} [EXT]; they enter no derivation, no definition, no coefficient of the construction; this document identifies no particles ; no parameter is fitted (P1, P5; FT.87(11) precedent). $(1)$ \textbf{Bare 1-loop read: exact shortfalls.} The phase triple $(m_1,m_2,m_3)=(1,\,1.509244,\,1.906061)$ in $m$, $(1,\,2.277816,\,3.633067)$ in $m^2$ [FT.103(2)] against [EXT]:

%dsp-level:0
\[
\begin{array}{c|cccc}
&\multicolumn2c{m\text{-read shortfall}}&\multicolumn2c{m^2\text{-read shortfall}}\\
\text{[EXT] family}&\text{middle}&\text{outer}&\text{middle}&\text{outer}\\ \hline
\text{leptons }(207,\,3477)&\times137.2&\times1824.2&\times90.9&\times957.0\\
\text{down-type }(20,\,900)&\times13.3&\times472.2&\times8.8&\times247.7\\
\text{up-type }(600,\,80000)&\times397.6&\times41971.4&\times263.4&\times22020.0
\end{array}
\]

(enclosure-proven lower bounds: every middle shortfall $>10$, every outer $>100$.) The $m$-span $1.906$ is of the \textit{same order} as the tree proxy's flat span $2$ and the all-four $m^2$-span $5.597$ of the same order as FT.79(1)'s pool maximum $6$: \textbf{the loop correction re-shapes the vacuum geometry — bands, walls, dominance windows — without extending the computed cells' numeric range.} $(2)$ \textbf{Range-species classification.} (a) \textit{Bare read:} the vacuum-depth range on a family grows with the cells' size — floor-side barrier instance at $p_0=\tfrac{15}{16}$: $-\ln W_\Gamma(p_0)=(7.31,\ 16.01,\ 27.05)$ with exponents $2E-K_{\mathrm{top}}=(6,12,16)$ and the exact lower-bound law $-\ln W_\Gamma(p_0)\ge(2E-K_{\mathrm{top}})\ln\tfrac1{q_0}-\ln\mathcal Z^{\mathrm{tot}}_\Gamma(16)$ [machine] — \textbf{linear in size, unbounded in the family}: unlike the defect-angle read (hard ceiling $4\hat\mu'^2$ [FT.76(3)]) and of the same order class as the cover ladder — spanning $1{:}R$ needs $V\ge4R/3$ [FT.79(2): $V\ge276$ down-type, $V\ge4636$ leptons]. (b) \textit{Composed tilt read} (Def-H.40 species at general $\gamma$, composed with the vacua — declared composition, NOT adopted): per-edge contraction $c(\varepsilon_q)=1-\tfrac\gamma4=\tfrac34+\varepsilon_q$ [L2], composed over $E_{\mathrm{int}}$ internal edges [Def-H.40's composition clause]:

%dsp-level:0
\[
T_i(E):=c\bigl(\varepsilon_q^{(i)}\bigr)^{2E};\qquad \bigl(c(\mathrm V_1),c(\mathrm V_3),c(\mathrm V_4)\bigr)=(0.771337,\ 0.830091,\ 0.940792);
\]

%dsp-level:0
\[
\beta_{\mathrm{mid}}:=\tfrac{c(\mathrm V_3)}{c(\mathrm V_1)}=1.076172;\qquad \beta_{\mathrm{out}}:=\tfrac{c(\mathrm V_4)}{c(\mathrm V_1)}=1.219690;
\]

the range is \textbf{exponential in $E$} — the first read in Gap 7's comparison set that provably escapes both hard ceilings — with exact thresholds (each strict: $E{-}1$ falls below; enclosure-proven):

%dsp-level:0
\[
\beta_{\mathrm{out}}^E\ge20\iff E\ge16;\quad \ge207\iff E\ge27;\quad \ge900\iff E\ge35;\quad \ge3477\iff E\ge42.
\]

Against FT.79(2)'s $V\ge276$ (down-type, cover ladder) and $V\ge4636$ (leptons): a \textbf{two-order compression} of the size requirement, $E_{\mathrm{int}}\approx16$–$42$ — within the corpus's computed-cell scale class. $(3)$ \textbf{Proportionality obstruction (exact).} No single $E$ carries any [EXT] triple through (2b): the triple's shape is fixed by the base pair $(\beta_{\mathrm{mid}},\beta_{\mathrm{out}})$, and

%dsp-level:0
\[
\frac{\ln\beta_{\mathrm{mid}}}{\ln\beta_{\mathrm{out}}}=0.369645\qquad\text{vs}\qquad\frac{\ln207}{\ln3477}=0.654006\quad[\mathrm{EXT}]
\]

(disjoint enclosures; mismatch factor $1.7693$). At the outer threshold $E=42$: $\beta_{\mathrm{mid}}^{42}=21.83<207$ — middle shortfall $\times9.5$. The obstruction datum is the \textbf{base clustering}: the walls cluster in a band of width $0.0115$ [FT.70(2)], so the contraction bases span only $(0.771,0.941)$; an [EXT]-shaped triple would need the vacua's $\ln c$-spreads proportional to $(\ln207,\ln3477)$ — missed by the exact factor above. This is the 1-loop species' own structural obstruction, dual to FT.76(3)'s angle bound and FT.79(2)'s edge budget. $(4)$ \textbf{Adjudication summary.} The bridge is constructed and machine-exact: the tree limit reproduces the covering census [FT.102(3)]; the loop layer produces the corpus's first multi-vacuum spectrum — four exact vacua, three wall-delimited phase bands, walls at FT.70's boundaries, dominance windows with exact crossings [FT.103]; the mass read is $\Delta S$-admissible (all positive, nats) [L3]. It does not, on the computed cells: reach the [EXT] spans (clause (1)'s exact shortfalls); derive the number three (band count $=|\mathbb F|$, declared [FT.103(5)]); relate to $m_{\mathrm{pole}}$/$m_{\mathrm{hid}}$ (Def-H.39 unchanged; the category discipline of FT.89(2) applies to every cross-read). The exponential mechanism is real but sits in the composed species (clause (2b)), whose size thresholds are exact and whose proportionality datum (clause (3)) is the new obstruction. Item 7's comparison set gains the fourth read; the mass-hierarchy record (two hard ceilings $+$ linear bare $+$ exponential composed, all computed cells $\le6.39$ in span), 2 and 3 unchanged.

\textbf{Req.} FT.102, FT.103; FT.73(5); FT.76(3); FT.77; FT.79(1), (2), (4)–(6); FT.84(2); FT.87(9), (11); FT.89(2); L2, L3; Def-H.39, Def-H.40; P1, P5; [EXT] data clause; rational log-enclosures for every threshold and shortfall [math; machine: \texttt{checks/\allowbreak{}verify\_\allowbreak{}H\_\allowbreak{}1loop\_\allowbreak{}bridge.py} §G–§H].

\textbf{Pf.}

%dsp-level:0
\[
(1):\ \text{shortfall}
\]

%dsp-level:0
\[
= [\mathrm{EXT}]/\text{ratio}:\ \text{interval division of enclosure-proven ratios by exact integers}\ [\text{machine}]
\]

%dsp-level:0
\[
(2a):\ W_\Gamma(p)\le\mathcal Z^{\mathrm{tot}}_\Gamma(16)\cdot q^{2E-K_{\mathrm{top}}}\ \text{for }p
\]

%dsp-level:0
\[
\ge p_0\ (\text{termwise},\ p\le1):\ \text{log-enclosure comparison at }p_0
\]

%dsp-level:0
\[
= \tfrac{15}{16}\ [\text{machine}]
\]

%dsp-level:0
\[
(2b):\ \beta_{\mathrm{out}}^E
\]

%dsp-level:0
\[
\ge R
\]

%dsp-level:0
\[
\iff E\ln\beta_{\mathrm{out}}\ge\ln R:\ \text{both sides by rational log-enclosures; strictness at }E-1\ [\text{machine}];
\]

%dsp-level:0
\[
\qquad c=\tfrac34+\varepsilon_q\ [\mathrm{L2}:\ 1-\tfrac\gamma4,\ \gamma=1-4\varepsilon_q]
\]

%dsp-level:0
\[
(3):\ \ln\beta\ \text{enclosures at tolerance }2^{-50}:\ \text{quotient intervals disjoint};
\]

%dsp-level:0
\[
\qquad 42\ln\beta_{\mathrm{mid}}<\ln207\ [\text{machine}]
\]

%dsp-level:0
\[
(4):\ \text{conjunction of FT.102/FT.103 clauses and (1)–(3)};
\]

%dsp-level:0
\[
\qquad \text{no identification, no calibration invoked}\ [\mathrm{P1},\mathrm{P5}].\ \blacksquare
\]

\textbf{Phys.} The bridge carries the census across into geometry: walls where the phase diagram has boundaries, vacua where the ensemble has phases, masses as missing information. What it does not carry is the observed hierarchy: on the computed cells the span stays of order six in $m^2$ — the third independent appearance of that scale — and the exponential range that \textit{does} exist lives in the composed tilt size, reaching the lepton span at $E_{\mathrm{int}}=42$ while missing the middle value by the exact factor $9.5$. The ledger has learned to speak geometry; the hierarchy question now has a fourth, unbounded-range species with quantified size requirements and an exact proportionality obstruction any future family must beat. No particle identification; no calibration; the pole$\leftrightarrow$info relational question remains open, its residue sharpened.

\vspace{2pt}\noindent\rule{\textwidth}{0.4pt}\vspace{2pt}

\end{theorem}

\phantomsection
\section[Tier F.36 — The Topology-Blind Path Integral:…]{\texorpdfstring{Tier F.36 — The Topology-Blind Path Integral: Observable-Conditioned Sectors and the Janus Sum}{Tier F.36 - The Topology-Blind Path Integral: Observable-Conditioned Sectors and the Janus Sum}}

\begin{definition}[{\texorpdfstring{Def-H.32 — The topology-blind path integral}{Def-H.32 - The topology-blind path integral}}]\label{tag:Def-H-32}
\textbf{Def.} A construction class on the prior canon (Def-H.24 pattern: no new axiom); clauses (1)–(7) are constitutive interface definitions. Scopes (S1)–(S4) of Def-H.25 are inherited at every use ($\kappa_e\equiv1$ minimal-edge class on real edges unless a $\kappa$-class is named; literal $\aleph$-domain; seal/Janus bookkeeping static, $\mathcal J=1$; virtual edges weight $1$, no phase).

$(1)$ \textbf{The topology space and the budget filtration (Def-H.32(1)).} $\mathcal G:=$ the isomorphism classes of finite connected simple graphs (self-loops excluded [A1]; skeletons read as their bidirectional digraphs, $2E$ directed edges [A7, D0.1]; connectedness per D2.13's subnetwork scope). Regularization $=$ the declared \textbf{budget filtration}

%dsp-level:0
\[
\mathcal G_B:=\{\Gamma\in\mathcal G:\ V(\Gamma)\le B\},\qquad
\mathfrak C_B(\Gamma):=\{R\subseteq E_{\mathrm{dir}}(\Gamma):\ |R|\le B_K\}\quad[B_K\ \text{default}\ 2E(\Gamma)],
\]

the budget being the ledger's own content measure: $|R|=n_{\mathcal R}$, $S_{\mathfrak L}=n_{\mathcal R}\ln2$ [CT1.29]. Every filtered sum below is a finite exact sum [FT.68(7) pattern]; the $B\to\infty$ limit is not constructed (growth data: FT.105).

$(2)$ \textbf{Configurations and the two weight layers.} A configuration on $\Gamma$ is $C=(R,\sigma)\in\mathfrak C(\Gamma)$ [D0.10]: pattern $R$ with per-edge signs $\sigma_e\in\{\pm1\}$ [CT0.9, D3.4]; vertex types $(n_{\mathrm{in}},n_{\mathrm{out}})$ [A7, D0.8]; locked ($D^-$-configured) $\iff k_v\ge4$ [CT4.5 source criterion; CT0.11]. The A8 edge weight is $w(C)=\mathcal J(C)\prod_{e\in E(C)}G(e)\cdot(-1)^{\aleph(\sigma,C)}$ [A8; at (S1): $G(e)=-\sigma_e$; $\aleph=\sum_{\text{centres}}|f_v|$, $f_v=\tfrac{s_{\mathrm{out}}-s_{\mathrm{in}}}2$ [D3.3]]. Two layers [CT2.8]:

%dsp-level:0
\[
\textbf{amplitude layer}:\quad v_\Gamma(R):=\sum_\sigma w(R,\sigma)\qquad[\mathrm{FT.67}(2):\ v=2^{|R|}(-1)^{\sum\kappa}\ \text{on covers},\ 0\ \text{else}];
\]

%dsp-level:0
\[
\textbf{norm layer}:\quad \nu_\Gamma(R):=\sum_\sigma|w(R,\sigma)|^2=2^{|R|}
\]

%dsp-level:0
\[
\qquad [|w(R,\sigma)|=1\ \text{for every valid pair};\ \mathrm{FT.68}\ \text{Rmk}(6c)].
\]

$(3)$ \textbf{Fugacities (Def-H.32(3) — formal indeterminates over $\mathbb Q$; no rate datum, no declared value).} Per configuration:

%dsp-level:0
\[
\mathcal W^{\mathrm{amp}}_\Gamma(R)\ :=\ v_\Gamma(R)^2\;t^{|R|}\,u^{L(R)};\qquad
\mathcal W^{\mathrm{norm}}_\Gamma(R)\ :=\ \nu_\Gamma(R)\;t^{|R|}\,u^{L(R)},
\]

$L(R):=\#\{v:\ k_v\ge4\}$ the locked census; $t$ the crystallization fugacity, $u$ the $D^-$-retention fugacity [CT4.1–CT4.5 vocabulary: absorption/emission/dissolution/retention enter as the static fate fugacity; the emission/dissolution \textit{moves} live in clause (7)]. Specializations are marked at every use: $t=p/q$ with $p=x_+/(x_++x_-)$ [FT.108(4), derived]; tilt refinement $\mathcal W\mapsto\mathcal W\cdot\tau_{\mathcal R}^{|R|}\tau_{\mathcal V}^{2E-|R|}$ with spectral values $\tau_{\mathcal R}=e^{-\gamma/3}$, $\tau_{\mathcal V}=e^{-\gamma}$ [Cor L0a at the branch states; FT.108(7)].

$(4)$ \textbf{Observables (Def-H.32(4) — internal bookkeeping labels; no identification).} A configuration carries chains [D1.4: ordered tuples of distinct real edges]: \textbf{closed} directed simple cycles $c$, length $L\ge2$; \textbf{open} directed simple paths, length $\ge2$. Readings:

%dsp-level:0
\[
q(c,\sigma):=\frac{\sigma_{e_1}-\sigma_{e_k}}{2}\in\{-1,0,+1\}
\]

%dsp-level:0
\[
\quad [\mathrm{EM.9}\ \text{end registers};\ \mathrm{FT.40}\ \text{closed-core form at the tuple's own base};\ \mathrm{D3.4}],
\]

%dsp-level:0
\[
S(c):\ \text{the on-shell winding }k(c)=\Theta(c)/2\in\mathbb Z\quad[\mathrm{FT.22},\ \mathrm{FT.25}(3)],
\]

%dsp-level:0
\[
\qquad \text{minimal ladder}\ [\mathrm{FT.32}(6)]:\ \ L\ \text{even}\Rightarrow(\kappa\equiv1,\ k=L/2);\ \ L\ \text{odd}\Rightarrow(\kappa\equiv2,\ k=L);
\]

open chains carry $S=0$ (no closed core — bookkeeping label, $\neg$ a spin value; the $j=|k|/2$ map is Def-H.38-declared content and is \textit{not} imported). A \textbf{chain instance} is a pair $(c,\text{base})$ (closed: base $=$ each of the $L$ cyclic starts; open: the tuple itself). The $\delta$-constraints of the sector sum are realized by summing over instances with the readings fixed.

$(5)$ \textbf{The sector partition functions.} For labels $(q,S)$ and budget $(B,B_K)$:

%dsp-level:0
\[
Z^{\bullet}_{B}(q,S)\ :
\]

%dsp-level:0
\[
= \ \sum_{\Gamma\in\mathcal G_B}\ \sum_{R\in\mathfrak C_B(\Gamma)}\ \mathcal W^{\bullet}_\Gamma(R)\cdot\#\{(c,\text{base}):\ q(c,\sigma)\text{-sector}=q,\ S(c)=S\}\Big|_{\sigma\text{-resolved}},
\]

$\bullet\in\{\mathrm{amp},\mathrm{norm}\}$: the amplitude layer with $\sigma$ resolved inside the sector condition and $w$-weights; the norm layer with $|w|^2$-weights (closed form: $m_q\,2^{|R|-2}$ per instance, $m_q=(2,1,1)$ at $q=(0,\pm1)$ [FT.106(3)]). Both are exact polynomials in $\mathbb Z[t,u]$ per budget cell.

$(6)$ \textbf{The mass expectation (Def-H.32(6)).} For sectors with $Z^{\mathrm{norm}}_B(q,S)\ne0$:

%dsp-level:0
\[
m^2_B(q,S)\ :=\ \frac{1}{Z^{\mathrm{norm}}_B(q,S)}\sum_{\Gamma,R}\ \mathcal W^{\mathrm{norm}}_\Gamma(R)\cdot\#\text{inst}(q,S;\Gamma,R)\cdot\Delta S(\Gamma,R),
\]

%dsp-level:0
\[
\Delta S(\Gamma,R)\ :=\ 2\,n_{\mathcal V}(\Gamma,R)\,\bigl\langle|\delta\rho|^2\bigr\rangle_{\text{branch states}}\ =\ \tfrac32\,n_{\mathcal V}(\Gamma,R)\ =\ \tfrac32\bigl(2E(\Gamma)-|R|\bigr)
\]

%dsp-level:0
\[
\qquad [\mathrm{L4};\ \mathrm{L3}(1)\ \text{exact leading};\ \mathrm{n}_\mathrm{V}:=\#\text{virtual channels}],
\]

the branch-state values $|\delta\rho|^2=(\tfrac34,\tfrac1{12})$ (virtual $\Pi_{\mathcal V}$, real $\Pi_{\mathcal R}/3$) computed matrix-exact [FT.107(2)]; dimensionless nats [L3(4) species]; no coupling to $m_{\mathrm{pole}}/m_{\mathrm{hid}}$ (mass-form bridge residue; Def-H.39's bridge slots unconstructed); the exact relative-entropy variant ($\ln4$ per virtual channel) is a specialization remark, not the read.

$(7)$ \textbf{The Janus move system and the transition operator (Def-H.32(7)).} States $s=(\Gamma,R)$, $\Gamma$ any finite simple graph on $\le B$ labelled nodes (the path space admits disconnected and edgeless states; $\mathcal G$-membership is imposed at the sector sums). Moves, with formal weights:

%dsp-level:0
\[
J^+(e):\ R\mapsto R\cup\{e\},\ e\notin R\ [\text{weight }x_+];
\]

%dsp-level:0
\[
\qquad J^-(e):\ R\mapsto R\setminus\{e\}\ [\text{weight }x_-]
\]

%dsp-level:0
\[
\qquad [\mathrm{D1.2}:\ \text{one pair across A4 per event}];
\]

%dsp-level:0
\[
N^+(v,S):\ \text{add node }v\text{ with }d=|S|\ge1\ \text{virtual edges to existing nodes}\ [\text{weight }y_d];
\]

%dsp-level:0
\[
N^-(v):\ \text{remove node }v\text{ of degree }d\text{ together with its }r(v)\ \text{real edges}\ [\text{weight }z_d\,x_-^{r(v)};\ \mathrm{D1.2}:\ \text{reality is per pair}].
\]

$\hat U:=\sum_{\text{moves}}w\,|t\rangle\langle s|$ over the budget-truncated state space; $\langle s'|\hat U^n|s\rangle\in\mathbb Z[x_\pm,y_d,z_d]$; the topology sum of clause (5) and the transition sum $\sum_n\langle s'|\hat U^n|s\rangle$ are the static and dynamical faces of the same move content (FT.108).

\textbf{Rmk (interfaces and non-imports).} Frozen data only: A8/D3.3 weights; FT.67(2)'s classifier; CT0.9/CT0.11 registers and cap; CT1.29 ledger nats; EM.9/FT.40/FT.22/FT.25/FT.32 readings; CT4.1–CT4.5 lock vocabulary; Cor L0a tilt spectral values; L3/L4 deficit reads. Not imported: any declaration-entry entry (Def-H.38's $j=|k|/2$, Cal-F's rate datum, the non-adopted bridge-map's observer map, Def-H.39's bridge, Def-H.40's composition, Def-H.31's dictionary — the ensemble parameter $p$ and the dictionary $p=1-4\varepsilon_q$ are \textit{derived} here as the stationary measure of clause (7)'s edge chain and its specialization [FT.108(4)]). Not declared: aspect-resolved couplings inside $w$ (the gauge-block coupling structure does not close [FT.82]; the aspects of A8/CT0.7(4) label the move algebra only [FT.108(1)]); observer-channel relativity of the readings (Def-H.1 vocabulary; the sector sums are channel-free instance censuses — a declaration would be needed to make them observer-relative, and none is made). Topology labels (names of cells) are internal bookkeeping throughout.

\vspace{2pt}\noindent\rule{\textwidth}{0.4pt}\vspace{2pt}

\end{definition}

\begin{theorem}[{\texorpdfstring{FT.106 — Weight transport and the two-layer sector structure}{FT.106 - Weight transport and the two-layer sector structure}}]\label{tag:FT-106}
\textbf{Thm.} Under Def-H.32: $(1)$ \textbf{The exact-cover classifier is topology-blind.} FT.67(2)'s classification — $v_\Gamma(R)=2^{|R|}(-1)^{\sum\kappa}$ iff the type-$(1,1)$ centres exactly single-cover $R$, else $0$ — holds verbatim on every $\Gamma\in\mathcal G$: its involution proof consumes only local vertex types [D3.3] and sign flips, no global topology. Machine: brute-force sign sums equal the classifier on \textbf{all $5845$ patterns of all ten $V\le4$ cells} and on $5888$ further $V{=}5$ witness patterns ($E\le5$). $(2)$ \textbf{Closed cells die in amplitude on every topology; covers thread them.} Every pure directed cycle (any length, any cell) has $v=0$ [FT.67(3a) transport]; on every enumerated cover containing directed cycles the cycle-alternation lemma holds exactly (centres/non-centres strictly alternate; all cycle lengths even; no odd cycle inside any cover) [FT.69(3) transport]; census: the only closed content inside covers at $B=4$ is dissolved digons (and their locked-hub families [FT.69(5)]). $(3)$ \textbf{The two-layer sector structure (exact census at $B=4$; $\sigma$-resolved scope: all cells with $2E\le8$ exhaustive, $K_4$ at $K\le4$ [FT.74(5) precedent], diamond at $K\le5$).} Sector totals over $\mathcal G_4$, $Z^{\mathrm{norm}}$ complete / $Z^{\mathrm{amp}}$ in scope:

%dsp-level:0
\[
\begin{array}{c|cc|cc}
(q,S)&Z^{\mathrm{norm}}&&Z^{\mathrm{amp}}&\\ \hline
(\pm1,0)\ \text{transit}&197912&&1196&\\
(0,0)\ \text{transit}&395824&&5336&m_q=(2,1,1)\\
(\pm1,1)\ \text{digon winding}&787358&&802&\text{dissolved digons in covers}\\
(0,1)&1574716&&4796&\\
(\pm1,2)\ (C_4\text{-type winding})&655776&&0&\text{no cover contains a 4-cycle [FT.69(3)(v)]}\\
(0,2)&1311552&&0&\\
(\pm1,3)\ (\text{triangle winding},\ \kappa\equiv2)&1000512&&0&\text{odd cycles never enter covers [FT.69(3)(iv)]}\\
(0,3)&2001024&&0&
\end{array}
\]

Structural identities (all machine): \textbf{(a) norm-layer closed form} — per chain instance the $\sigma$-sum factorizes as $m_q2^{|R|-2}$, $m_q=(2,1,1)$ at $q=(0,\pm1)$ [EM.9's four end-pairs]; closed form $=$ census in scope, every cell and sector. \textbf{(b) amplitude $q$-sum identity} — for every winding $k$: $\sum_qZ^{\mathrm{amp}}_B(q,k)=\sum_{R\ \mathrm{cover}}v_\Gamma(R)\cdot\#\mathrm{inst}_k(R)$ exactly: the charge-resolved amplitude layer admits non-cover support ($18$ nonzero cancelling entries at $B=4$ — a sector-level refinement of FT.67's total), and the non-cover parts cancel exactly across $q$. \textbf{(c) amplitude death of $k\ge2$} — $Z^{\mathrm{amp}}(q,k)=0$ in total for $k=2$ and odd $k\ge3$; the winding-$1$ sector is amplitude-live exactly through dissolved-digon covers. \textbf{(d) $\omega$-symmetry} — $Z^{\bullet}(q,S)=Z^{\bullet}(-q,S)$: norm layer unconditional ($|w|$-invariance [EM.24(5) pattern]); amplitude layer on the cover support ($|R|=2|\mathcal C|$ even). \textbf{(e) charge alphabet} — $\{-1,0,+1\}$ attained in every chain-carrying cell [EM.9 transport]; every sector value is a nonnegative integer (norm) resp. integer (amp).

\textbf{Req.} Def-H.32; A7, A8; D0.8, D0.10, D3.3, D3.4; CT0.9, CT0.11 (Part 1); CT2.8 (Part 2); EM.9, EM.24(5) (Part 2); Def-H.25; FT.67, FT.68, FT.69(3)–(5), FT.74(3),(5); exhaustive enumeration, involution transport, factorized $\sigma$-sums [math; machine: \texttt{checks/\allowbreak{}verify\_\allowbreak{}H\_\allowbreak{}topology\_\allowbreak{}blind.py} §B–§C].

\textbf{Pf.}

%dsp-level:0
\[
(1):\ \mathrm{FT.67}(2)\text{'s involution flips }\sigma_e\text{ on unread/double-read edges per vertex — local data }[\mathrm{D3.3}];
\]

%dsp-level:0
\[
\qquad \text{brute force}=\text{classifier on }5845+5888\ \text{patterns}\ [\text{machine}]
\]

%dsp-level:0
\[
(2):\ \text{pure cycles}:\ \text{every vertex interior}\Rightarrow\text{cov}\equiv0\ \vee\ 2:\ v=0\ [\mathrm{FT.67}(3a)\ \text{proof transport}];
\]

%dsp-level:0
\[
\qquad \text{alternation census on every cover with cycles at }B=4\ [\text{machine}]
\]

%dsp-level:0
\[
(3a):\ q(c,\sigma)\ \text{reads two edge signs}:\ \#\{\sigma:\ q(c,\sigma)=q\}
\]

%dsp-level:0
\[
= m_q2^{L-2}\ \text{on the chain};
\]

%dsp-level:0
\[
\ \text{rest free}:\ 2^{|R|-L}:\ \text{product}\qquad = m_q2^{|R|-2};
\]

%dsp-level:0
\[
\qquad \text{census}=\text{closed form}\ [\text{machine}]
\]

%dsp-level:0
\[
(3b):\ \sum_q\sum_{\sigma:\,q(c,\sigma)=q}w(R,\sigma)=\sum_\sigma w(R,\sigma)=v_\Gamma(R)\ [\text{instances are }\sigma\text{-independent}];
\]

%dsp-level:0
\[
\qquad v=0\ \text{off covers}\ [(1)]
\]

%dsp-level:0
\[
(3c):\ \text{covers contain no 4-cycles and no odd cycles at }B
\]

%dsp-level:0
\[
= 4\ [\text{census}]
\]

%dsp-level:0
\[
\implies \#\mathrm{inst}_{k\ge2}
\]

%dsp-level:0
\[
= 0\ \text{on covers}\implies\sum_qZ^{\mathrm{amp}}\qquad = 0\ [(3b)]
\]

%dsp-level:0
\[
(3d):\ \sigma\mapsto-\sigma:\ q\mapsto-q;
\]

%dsp-level:0
\[
\quad |w(R,-\sigma)|=|w(R,\sigma)|;
\]

%dsp-level:0
\[
\quad w(R,-\sigma)=(-1)^{|R|}w(R,\sigma);
\]

%dsp-level:0
\[
\quad |R|=2|\mathcal C|\ \text{even on covers}\ [\mathrm{FT.79}(2)]
\]

%dsp-level:0
\[
(3e):\ m_q>0\ \forall q;\ \ \text{instance existence per cell}\ [\text{machine census}].\ \blacksquare
\]

\textbf{Phys.} Weighing every topology against every pattern does not blur the frozen anatomy — it confirms it cell by cell: the amplitude layer remains a thread gas with dissolved even cores, the norm layer remains strictly positive, and the odd cores stay amplitude-dead on all ten topologies at once. The new datum is finer than the total: resolved by charge, the amplitude layer carries cancelling non-cover parts whose sum is exactly the cover census — FT.67's cancellation, seen from inside. No particle identification.

\vspace{2pt}\noindent\rule{\textwidth}{0.4pt}\vspace{2pt}

\end{theorem}

\begin{theorem}[{\texorpdfstring{FT.107 — The observable sectors and the mass expectation}{FT.107 - The observable sectors and the mass expectation}}]\label{tag:FT-107}
\textbf{Thm.} Under Def-H.32: $(1)$ \textbf{On-shell ladder (transport of FT.32(6)).} A directed $L$-cycle is on-shell at $(\kappa\equiv1,\ k=L/2)$ for $L$ even and at $(\kappa\equiv2,\ k=L)$ for $L$ odd; instances at $B=4$: $k=1$ (digons, $\kappa\equiv1$), $k=2$ ($4$-cycles, $\kappa\equiv1$), $k=3$ (triangles, $\kappa\equiv2$); the floor card $L=2$, $k=1$ is FT.28's $\Sigma_{\min}$ datum reproduced by the census. $(2)$ \textbf{Branch-state deficits (matrix-exact).} At the canonical branch states [CT0.8 sharp readings]: $\rho_{\mathcal V}=\Pi_{\mathcal V}$: $|\delta\rho|^2=\tfrac34$; $\rho_{\mathcal R}=\Pi_{\mathcal R}/3$: $|\delta\rho|^2=\tfrac1{12}$; leading deficits $2|\delta\rho|^2=(\tfrac32,\ \tfrac16)$ [L3(1)]; hence $\Delta S(\Gamma,R)=\tfrac32\,n_{\mathcal V}(\Gamma,R)=\tfrac32(2E-|R|)$ — an exact rational per configuration. $(3)$ \textbf{The sector mass table ($B=4$, $(t,u)=(1,1)$, norm layer; exact rationals).} $m^2_B(q,S)$ is independent of $q$ (charge-degenerate at the leading deficit read — the topology-blind companion of FT.100's phase-blindness):

%dsp-level:0
\[
\begin{array}{c|cccc}
S\ (\text{winding})&0\ (\text{transit})&1&2&3\\ \hline
m^2_4(\cdot,S)&\dfrac{843957}{98956}&\dfrac{1921662}{393679}&\dfrac{91}{23}&\dfrac{13713}{3088}\\[2mm]
&\approx8.5286&\approx4.8813&\approx3.9565&\approx4.4407
\end{array}
\]

all strictly positive ($\Delta S$-admissible, nats [L3(4)]); the ordering places the transit sector heaviest — thread-network instances ride sparser patterns (larger $n_{\mathcal V}$), the deficit read weights missing information. $(4)$ \textbf{Quadratic budget law.} For every sector and budget: $0<m^2_B(q,S)\le\tfrac32\max_\Gamma n_{\mathcal V}\le\tfrac32\,B(B-1)$ [exact bound; $\max$ at $R=\varnothing$ on the densest cell]: \textbf{quadratic in the node budget} — at $B=4$: bound $18$, attained census maximum $\tfrac{843957}{98956}<9$. The winding spread at $B=4$ is $\tfrac{91}{23}\le m^2\le\tfrac{843957}{98956}$ (span $\approx2.16$ in $m^2$): the computed-cell span remains order-one while the \textit{bound} grows as $\Theta(B^2)$ — against the fixed-topology bare read's linear-in-size range [FT.104(2a)] and the two hard ceilings [FT.76(3), FT.79(2), FT.97(5)]. The fugacity axes add an unbounded structural direction (FT.108(5)); no axis value is declared. No [EXT] datum enters this part; no comparison is made.

\textbf{Req.} Def-H.32; CT0.7, CT0.8 (Part 1); L3, L4; FT.22, FT.24, FT.25(3), FT.28, FT.32(6); FT.100; FT.104(2); FT.76(3), FT.79(2), FT.97(5); exact rational matrix arithmetic, sector census algebra [math; machine: \texttt{checks/\allowbreak{}verify\_\allowbreak{}H\_\allowbreak{}topology\_\allowbreak{}blind.py} §C–§D].

\textbf{Pf.}

%dsp-level:0
\[
(1):\ \Theta=\sum_{e\in c}\widehat{ds}^{\,2}(e):\ \kappa\equiv1:\ \Theta=L;\ \ \kappa\equiv2:\ \Theta=2L;
\]

%dsp-level:0
\[
\qquad \text{on-shell}\iff\Theta\in2\mathbb Z\ [\mathrm{FT.22}];
\]

%dsp-level:0
\[
\qquad k_{\min}=\Theta_{\min}/2\ [\mathrm{FT.32}(6)]
\]

%dsp-level:0
\[
(2):\ \delta\rho=\Pi_{\mathcal V}-\tfrac14I:\ \operatorname{spec}=(-\tfrac14,-\tfrac14,-\tfrac14,\tfrac34):\ \mathfrak T(\delta\rho^2)=\tfrac3{16}+\tfrac9{16}=\tfrac34;
\]

%dsp-level:0
\[
\qquad \delta\rho
\]

%dsp-level:0
\[
= \tfrac{\Pi_{\mathcal R}}3-\tfrac14I:\ \operatorname{spec}
\]

%dsp-level:0
\[
= (\tfrac1{12},\tfrac1{12},\tfrac1{12},-\tfrac14):\ \mathfrak T(\delta\rho^2)
\]

%dsp-level:0
\[
= \tfrac3{144}+\tfrac1{16}
\]

%dsp-level:0
\[
= \tfrac1{12}\ [\text{machine}]
\]

%dsp-level:0
\[
(3):\ m^2
\]

%dsp-level:0
\[
= \frac{\sum_{\Gamma,R}\nu\,t^{|R|}u^{L}\#\mathrm{inst}\cdot\tfrac32(2E-|R|)}{\sum_{\Gamma,R}\nu\,t^{|R|}u^{L}\#\mathrm{inst}}\Bigg|_{t=u=1}:\ \text{exact rational division of census integers};
\]

%dsp-level:0
\[
\qquad q\text{-independence}:\ \Delta S(R)\ \text{is }\sigma\text{-free and }m_q\text{-symmetric}\ [(3a),(3d)\ \text{of FT.106}]
\]

%dsp-level:0
\[
(4):\ n_{\mathcal V}\le2E\le B(B-1)\ [\text{simple graphs}];
\]

%dsp-level:0
\[
\qquad \text{census max}<9<18\ [\text{machine}];
\]

%dsp-level:0
\[
\qquad \Theta(B^2)\ \text{vs}\ \Theta(B)\ [\mathrm{FT.104}(2a)]\ \text{and vs}\ 4\hat\mu'^2/3V\text{-ladder ceilings}.\ \blacksquare
\]

\textbf{Phys.} The mass of an observable-conditioned sector is the ensemble average of missing information over every topology that carries it — one exact rational per (charge, winding) label at each budget, charge-blind at this order, ordered by how sparse the carrying patterns are. The transit sector weighs heaviest at $B=4$: open threads live on the emptier configurations, and emptiness is what the deficit read weighs. The quadratic budget law is the mechanism statement: topology-blindness moves the range axis from skeleton size to the budget itself — how far that axis must be pushed to reach any external span is arithmetic on $\tfrac32B(B-1)$, not physics claimed here. No particle identification; no external datum used.

\vspace{2pt}\noindent\rule{\textwidth}{0.4pt}\vspace{2pt}

\end{theorem}

\begin{theorem}[{\texorpdfstring{FT.108 — The Janus sum: move algebra, budget conservation, derived detailed balance, and the Tier F.35 boundary}{FT.108 - The Janus sum: move algebra, budget conservation, derived detailed balance, and the Tier F.35 boundary}}]\label{tag:FT-108}
\textbf{Thm.} Under Def-H.32: $(1)$ \textbf{Move-label algebra.} The move classes carry CT0.7(4)'s aspect labels — $J^\pm\leftrightarrow F_{J^\pm}=\mathcal E_{-+}/\mathcal E_{--}$ (S-odd), node moves $\leftrightarrow$ the seal-generated topology change — and the frozen identities hold matrix-exact over $\mathbb Q(i)$ on the label algebra: $[\Sigma,J_i]=4iK_i$, $[\Sigma,K_i]=-4iJ_i$ [CT1.5]; $\{\mathbb S,J_i\}=0=\{\mathbb S,K_i\}$ [CT0.7(3)]. $(2)$ \textbf{Budget conservation (the coarse-grained invariant).} The grade map $\mathrm{gr}:(\Gamma,R)\mapsto\bigl(V,\,E_{\mathrm{pot}},\,|R|\bigr)$ is a homomorphism of the move system: every monomial $x_+^{a_+}x_-^{a_-}\prod y_d^{m_d}\prod z_d^{n_d}$ of $\langle s'|\hat U^n|s\rangle$ satisfies

%dsp-level:0
\[
a_+-a_-=|R'|-|R|;\qquad \textstyle\sum_d m_d-\sum_d n_d=V'-V;\qquad \sum_d d\,(m_d-n_d)=E'_{\mathrm{pot}}-E_{\mathrm{pot}}
\]

— exhaustive on every path monomial at $n\le3$ from every $V\le2$ source [machine]. Node dissolution carries its incident real edges' own $J^-$ factors [D1.2]: $z_d\,x_-^{r(v)}$ — without them the grade map fails (witnessed). $(3)$ \textbf{Path counts (two enumerators).} Operator iteration and depth-first path enumeration agree coefficient-wise on $\hat U^n$, $n\le4$; witness: $\langle\text{digon},R{=}E\,|\,\hat U^3\,|\,\bullet\rangle=2\,x_+^2y_1+\cdots$ — the two admissible move orderings [machine]. $(4)$ \textbf{Detailed balance derived; the dictionary is a specialization.} The per-pair Janus dynamics [D1.2] on (virtual, real) has master generator $Q=\bigl(\begin{smallmatrix}-x_+&x_+\\ x_-&-x_-\end{smallmatrix}\bigr)$ over $\mathbb Z[x_+,x_-]$; $\ker Q$ is one-dimensional with span $(x_-,x_+)$ [exact: $(x_-,x_+)Q=0$; $\det Q=0\ne Q$] — the unique stationary measure, hence

%dsp-level:0
\[
p\ =\ \frac{x_+}{x_++x_-}\qquad[\text{derived},\ \neg\text{imported}];
\]

the substitution $x_+=1-4\varepsilon_q$, $x_-=4\varepsilon_q$ gives $p=1-4\varepsilon_q$ identically [machine] — the fixed-topology ensemble parameter is a specialization of the derived stationary ratio, and the sealed universe-level balance $\#J^+=\#J^-$ [D1.3] is the balanced-sector restriction of $(2)$'s first identity. $(5)$ \textbf{Resolvent structure.} On a single pair channel, $\sum_n\langle\varnothing|\hat U^n|\varnothing\rangle w^n=\dfrac{1}{1-w^2x_+x_-}$ in $\mathbb Q[w]/(w^{12})$ [machine: $(1-w^2x_+x_-)\cdot\sum\equiv1$ coefficient-wise] — the transition sum is a geometric resolvent with pole at $w^2x_+x_-=1$: an unbounded structural range direction in the formal fugacities, no value declared (P4-class datum absent; any numerical use requires a logged declaration first). $(6)$ \textbf{The Tier F.35 boundary limit.} With node moves off ($y_d,z_d\mapsto0$) the state space decomposes into fixed-topology sectors, and on each sector $\Gamma_0$ the amplitude-layer cover generating function

%dsp-level:0
\[
\Phi_{\Gamma_0}(t)\ :=\ \sum_{R\ \mathrm{cover}}v_{\Gamma_0}(R)^2\,t^{|R|}\ =\ \sum_kc_k16^kt^{2k}
\]

satisfies the exact polynomial identity

%dsp-level:0
\[
q^{2E}\,\Phi_{\Gamma_0}\bigl(p/q\bigr)\ =\ W_{\Gamma_0}(p)
\]

%dsp-level:0
\[
\qquad [\text{coefficient-wise; machine on tetra/cube/octa against the frozen FT.68(1),(3)/FT.69(2) rows and on all nine live }V\le4\text{ cells}]
\]

— the fixed-topology tree layer of Def-H.31(2) [Tier F.35] is recovered identically; the lock refinement census at $(t,u)$: the tetrahedron's $16$ locked covers carry lock count $1$ each [FT.69(5) family]. Consequently the fixed-topology slice of the topology-blind integral carries Tier F.35's boundary data — its effective potential, four vacua, phase triple and FT.70 walls [FT.102–FT.104, under Def-H.31's interface definitions, cited as boundary content, not re-derived] — as the $y_d,z_d\to0$ restriction, per clause $(4)$'s specialization of $t$. $(7)$ \textbf{Tilt spectral specialization.} At the branch states, $Z_\gamma=\mathfrak T(\rho e^{-\gamma\rho})$ [Cor L0a] evaluates exactly by spectral calculus: real branch $e^{-\gamma/3}$, virtual branch $e^{-\gamma}$; at $\gamma=1$: disjoint rational enclosures $e^{-1/3}\in[0.716531311,\ldots]$, $e^{-1}\in[0.367879441,\ldots]$, strict order enclosure-proven [machine].

\textbf{Req.} Def-H.32; A4, A8; D0.3, D1.2, D1.3, D3.3; CT0.7, CT0.7a, CT1.2, CT1.5 (Part 1); L0 (Cor L0a); Def-H.25, Def-H.26, Def-H.31 (Tier F.35, boundary citation); FT.67, FT.68, FT.69(2),(5), FT.70, FT.74; FT.102, FT.103, FT.104(2a) (Tier F.35); exact multivariate polynomial arithmetic over $\mathbb Z$, master-equation kernel algebra, truncated power series, rational exp-enclosures [math; machine: \texttt{checks/\allowbreak{}verify\_\allowbreak{}H\_\allowbreak{}topology\_\allowbreak{}blind.py} §E–§G].

\textbf{Pf.}

%dsp-level:0
\[
(1):\ \text{matrices over }\mathbb Q(i)\ \text{from CT1.2's definitions verbatim};
\]

%dsp-level:0
\[
\qquad \text{all commutator/anticommutator entries compared exactly}\ [\text{machine}]
\]

%dsp-level:0
\[
(2):\ J^\pm:\ \Delta|R|=\pm1,\ \Delta V=\Delta E_{\mathrm{pot}}=0;
\]

%dsp-level:0
\[
\quad N^+(d):\ (\Delta V,\Delta E_{\mathrm{pot}},\Delta|R|)=(1,d,0);
\]

%dsp-level:0
\[
\quad N^-(v):\ (-1,-d,-r(v))\ \text{with weight }z_dx_-^{r(v)}:\ \text{grades match exponents monomial-wise}\ [\text{machine, exhaustive at }n\le3]
\]

%dsp-level:0
\[
(3):\ \text{dict-based operator iteration}
\]

%dsp-level:0
\[
= \text{recursive DFS aggregation}\ [\text{machine, coefficient-wise equality of multivariate polynomials}]
\]

%dsp-level:0
\[
(4):\ (x_-,x_+)Q=(x_-(-x_+)+x_+x_-,\ x_-x_+-x_+(x_-))=(0,0)\ [\text{exact in }\mathbb Z[x_+,x_-]];
\]

%dsp-level:0
\[
\qquad \det Q=x_+x_--x_-x_+=0,\ Q\ne0:\ \operatorname{rk}=1
\]

%dsp-level:0
\[
p=\tfrac{x_+}{x_++x_-};\qquad x_+=1-4\varepsilon_q,\ x_-=4\varepsilon_q:\ x_++x_-=1\implies p=1-4\varepsilon_q\ [\text{polynomial identity; machine}]
\]

%dsp-level:0
\[
(5):\ (\hat U^{2m})_{\varnothing\varnothing}=(x_+x_-)^m,\ (\hat U^{2m+1})_{\varnothing\varnothing}=0\ [\text{single channel}];
\]

%dsp-level:0
\[
\qquad (1-w^2x_+x_-)\sum_{m}(w^2x_+x_-)^m=1\ \text{in}\ \mathbb Q[w]/(w^{12})\ [\text{machine}]
\]

%dsp-level:0
\[
(6):\ v^2=4^{|R|}=16^{|\mathcal C|}\ [\mathrm{FT.67}(2),\ \mathrm{FT.79}(2)];
\]

%dsp-level:0
\[
\qquad q^{2E}\textstyle\sum_kc_k16^k(p/q)^{2k}=\sum_kc_k16^kp^{2k}q^{2E-2k}=W_{\Gamma_0}(p)\ [\mathrm{Def\text{-}H.31}(2)\ \text{form}];
\]

%dsp-level:0
\[
\qquad \text{coefficient-wise on frozen and live rows}\ [\text{machine}]
\]

%dsp-level:0
\[
(7):\ \operatorname{spec}(\Pi_{\mathcal R}/3)=(\tfrac13,\tfrac13,\tfrac13,0):\ \textstyle\sum\lambda e^{-\gamma\lambda}=e^{-\gamma/3};
\]

%dsp-level:0
\[
\qquad \operatorname{spec}(\Pi_{\mathcal V})=(0,0,0,1):\ e^{-\gamma};
\]

%dsp-level:0
\[
\qquad \gamma=1\ \text{enclosures}\ [\text{machine}].\ \blacksquare
\]

\textbf{Phys.} The Janus seals do what D1.2 always said they do — move one pair across the crystallization criterion — and once node events are admitted with the same per-pair discipline, topology itself becomes a walk in an exactly graded state space. The grade map is the conserved information budget; nothing else is conserved, and nothing else is needed: the stationary measure of a single pair channel comes out as the ratio of the two seal weights, so the ensemble parameter every earlier tier declared is here derived, and the fixed-topology world is the slice on which node moves are switched off. On that slice the boundary data reproduce identically — every covering polynomial, every wall, every vacuum of the fixed-topology construction is the restriction of the topology-blind integral, not a competing answer. The transition sum is a resolvent: its pole is the unbounded direction, its location a datum the corpus does not declare. No particle identification.

\end{theorem}

\phantomsection
\section[Tier F.37 — The Koide Channel: The -Covariant Seal Bilinear…]{\texorpdfstring{Tier F.37 — The Koide Channel: The $C_3$-Covariant Seal Bilinear and the Bound-Chain Mixing Angle}{Tier F.37 - The Koide Channel: The -Covariant Seal Bilinear and the Bound-Chain Mixing Angle}}

\begin{theorem}[{\texorpdfstring{FT.109 — The diagonal seal commutator}{FT.109 - The diagonal seal commutator}}]\label{tag:FT-109}
\textbf{Thm.} Under CT1.7, with the CT1.2 basis $\{E_0,\Sigma,\hat E^{ij},W_i,J_i,K_i\}$:

%dsp-level:0
\[
(1)\quad\sum_{i=1}^{3}[J_i,K_i]=i\sum_{i=1}^{3}\bigl(\Sigma-E^{ii}\bigr)=i\bigl(3\Sigma-\Sigma\bigr)=2i\Sigma ;
\]

%dsp-level:0
\[
(2)\quad\sum_{i=1}^{3}E^{ii}=\sum_{i=1}^{3}\sigma_i\otimes\sigma_i=\Sigma\qquad\bigl[E^{ii}=\tfrac12\bigl(\sigma_i\otimes\sigma_i+\sigma_i\otimes\sigma_i\bigr)=\sigma_i\otimes\sigma_i\bigr];
\]

%dsp-level:0
\[
(3)\quad\text{basis read: in the 16-element CT1.2 basis the sum }\textstyle\sum_i[J_i,K_i]\text{ has exactly one nonzero component — the }\Sigma\text{ component, with coefficient exactly }2i;
\]

%dsp-level:0
\[
(4)\quad 2i\Sigma\in\mathcal E_{++}\ \text{(the Higgs/gravity block)},\qquad 2i\Sigma\notin\operatorname{span}\{W_k\}.
\]

\textbf{Req.} CT1.2, CT1.7 (Part 1); exact $\mathbb Q(i)$ matrix arithmetic on all $16$ basis elements, basis-subset disjointness [math; machine: \texttt{checks/\allowbreak{}verify\_\allowbreak{}H\_\allowbreak{}koide\_\allowbreak{}channel.py} §A–§B].

\textbf{Pf.}

%dsp-level:0
\[
(1):\ [J_i,K_i]=i\bigl(\delta_{ii}\Sigma-E^{ii}\bigr)=i\bigl(\Sigma-E^{ii}\bigr)\ [\mathrm{CT1.7}];
\]

%dsp-level:0
\[
\qquad \sum_{i=1}^{3}\bigl(\Sigma-E^{ii}\bigr)=3\Sigma-\sum_iE^{ii}=3\Sigma-\Sigma=2\Sigma\ [(2)]
\]

%dsp-level:0
\[
(2):\ E^{ii}=\sigma_i\otimes\sigma_i\ \text{and}\ \Sigma=\textstyle\sum_i\sigma_i\otimes\sigma_i\ [\mathrm{CT1.2}\ \text{definitions, entry-wise over}\ \mathbb Q(i)]
\]

%dsp-level:0
\[
(3):\ \textstyle\sum_i[J_i,K_i]
\]

%dsp-level:0
\[
= 2i\Sigma\ \text{entry-wise exact over}\ \mathbb Q(i);
\]

%dsp-level:0
\[
\ \ \text{the 16 CT1.2 elements are a basis of}\ \mathfrak H(\mathbb C^4)\implies\text{component reads are unique}\ [\mathrm{CT1.2};\ \text{machine §A}]
\]

%dsp-level:0
\[
(4):\ S\Sigma S=\Sigma\ \wedge\ T\Sigma T=\Sigma\implies 2i\Sigma\in\mathcal E_{++}\ [\mathrm{CT1.2}\ \text{sector table}];\qquad SW_kS=-W_k\ [\mathrm{CT1.2}]
\]

%dsp-level:0
\[
\operatorname{span}\{E_0,\Sigma,\hat E^{ij}\}\cap\operatorname{span}\{W_k\}=\{0\}\ [\text{disjoint basis subsets};\ \mathrm{CT1.2}];
\]

%dsp-level:0
\[
\qquad 2i\Sigma\ne0\implies 2i\Sigma\notin\operatorname{span}\{W_k\}\ [\text{machine §A}].\ \blacksquare
\]

\textbf{Phys.} Bracketing each spatial seal with its own temporal seal and summing over the triplet cancels the entire symmetric-tensor content and leaves exactly twice the exchange scalar: the diagonal of the mixed seal commutator is a pure Higgs/gravity-block object, with exact integer coefficient $2$ and zero weak component. This integer $2$ is the sole numerator entering the mixing angle of FT.112 — no datum, no fit.

\vspace{2pt}\noindent\rule{\textwidth}{0.4pt}\vspace{2pt}

\end{theorem}

\begin{definition}[{\texorpdfstring{Def-H.41 — The bound-chain arena}{Def-H.41 - The bound-chain arena}}]\label{tag:Def-H-41}
\textbf{Def.} The \textbf{bound-chain arena} is the minimal counting space combining the two independently present structures of a bound chain $c$ [D1.4]:

%dsp-level:0
\[
\mathcal H_c\ :=\ \mathcal H_{\mathrm{seal}}\otimes\mathcal H_{\mathrm{poly}},
\]

%dsp-level:0
\[
\qquad \mathcal H_{\mathrm{seal}}:\ \text{the seal triplet at the chain's termini}\ [\mathrm{CT0.12},\ \mathrm{CT0.7}(4):\ F_{J^\pm}=\operatorname{span}\{J_i\}/\operatorname{span}\{K_i\},\ \dim=3],
\]

%dsp-level:0
\[
\mathcal H_{\mathrm{poly}}:\ \text{the polygonal transfer structure along its length}\ [\mathrm{FT.68}(2):\ \text{the transfer rule}\ T(t);\ \operatorname{tr}T(1)=3];
\]

%dsp-level:0
\[
\qquad \dim_{\mathrm{eff}}\mathcal H_c=3\times3=9\ \ [\mathrm{FT.110}(1)].
\]

\textbf{Species (constitutive).} A counting space: no inner product, no dynamics, no state, no Hamiltonian is declared on $\mathcal H_c$; its dimension enters the channel theorems exactly once — as the denominator of the mixing angle [FT.112: $\delta_\ell=\tfrac29$ with FT.109's diagonal seal coefficient $2$ the numerator]. \textbf{Bars (standing).} $\dim\mathcal H_c$ is $\neg$a state count; the ratios $r_i^2$ [FT.113] reproduce no measured hierarchy — the exact nonzero shortfalls stand [FT.114 — forward pointer; Census row]; the overall scale $\mu$ is undeclared (P4-class, Cal-F); external calibration barred [P1, P5]; particle identification barred.

\end{definition}

\begin{theorem}[{\texorpdfstring{FT.110 — The bound-chain arena dimension}{FT.110 - The bound-chain arena dimension}}]\label{tag:FT-110}
\textbf{Thm.} The bound-chain arena $\mathcal H_c$ is the tensor product of two independent $3$-dimensional structures:

%dsp-level:0
\[
(1)\quad\dim_{\mathrm{eff}}\mathcal H_c=\dim F_{J^\pm}\times\operatorname{tr}T(1)=3\times3=9;
\]

%dsp-level:0
\[
(2)\quad\operatorname{tr}T(1)=3
\]

%dsp-level:0
\[
\qquad \bigl[\mathrm{FT.68}(2):\ \operatorname{charpoly}T(1)=x^2\bigl(x^2-3x-3\bigr)=x^4-3x^3-3x^2;\ \ \operatorname{tr}T(1)=-[x^3]\operatorname{charpoly}T(1)=3\bigr].
\]

\textbf{Rmk (tensor-product justification).} The tensor product is constitutive of Def-H.41 (immediately above): the minimal combination consistent with both structures being independently present — the seal triplet at the chain's termini [CT0.12, CT0.7(4); chains as in Def-H.32(4)] and the polygonal transfer structure along its length [FT.68(2)]. Both factors are derived; the product is the interface clause. $\mathcal H_c$ is an arena — a counting space: no inner product, no dynamics, no state on it is declared, and no Parts 1–2 result uses the declaration.

\textbf{Req.} CT0.7(4), CT0.12 (Part 1); FT.68(2) (Tiers F.23–F.25); Def-H.32 (Tier F.36, chain vocabulary); Vieta coefficient read, exact integer characteristic polynomial of the $4\times4$ transfer rule [math; machine: \texttt{checks/\allowbreak{}verify\_\allowbreak{}H\_\allowbreak{}koide\_\allowbreak{}channel.py} §C].

\textbf{Pf.}

%dsp-level:0
\[
(2):\ T(1)[s][s']
\]

%dsp-level:0
\[
= [\,|s|+|s'|\le3\,],
\]

%dsp-level:0
\[
\ s,
\]

%dsp-level:0
\[
s'\in\{0,1\}^2\ [\mathrm{FT.68}(2)\ \text{transfer rule at}\ t=1]:\ \text{the single excluded pair is}\ (|s|,|s'|)\qquad = (2,2)
\]

%dsp-level:0
\[
\implies T(1)=\mathbb 1_{4\times4}-[\,\cdot\,]_{(2,2)}\ \text{on the diagonal}:\ \operatorname{tr}T(1)=4-1=3\ [\text{exact integers; machine §C}]
\]

%dsp-level:0
\[
\operatorname{charpoly}T(1)
\]

%dsp-level:0
\[
= x^4-3x^3-3x^2
\]

%dsp-level:0
\[
= x^2\bigl(x^2-3x-3\bigr)\ [\mathrm{FT.68}(2);\ \text{machine: exact re-derivation from the}\ 4\times4\ \text{rule}];
\]

%dsp-level:0
\[
\qquad \operatorname{tr}T(1)=-[x^3]=3\ [\text{Vieta}]
\]

%dsp-level:0
\[
(1):\ \dim F_{J^\pm}=3\ [\mathrm{CT0.12}\ \text{table}:\ F_{J^+}=\mathcal E_{-+}=\operatorname{span}\{J_i\},\ F_{J^-}=\mathcal E_{--}=\operatorname{span}\{K_i\}\ [\mathrm{CT0.7}(4)]];
\]

%dsp-level:0
\[
\qquad 3\times3=9\ [\text{exact}].\ \blacksquare
\]

\textbf{Phys.} A bound chain carries a three-dimensional seal choice at its boundary and a three-dimensional transfer weight along its length; the arena that holds both bookkeepings at once is nine-dimensional. The count $9$ is the product of two independently derived threes — the seal triplet dimension and the transfer trace — and enters the tier exactly once: as the denominator of the mixing angle (FT.112).

\vspace{2pt}\noindent\rule{\textwidth}{0.4pt}\vspace{2pt}

\end{theorem}

\begin{theorem}[{\texorpdfstring{FT.111 — The $C_3$ spectrum of the covariant seal bilinear}{FT.111 - The spectrum of the covariant seal bilinear}}]\label{tag:FT-111}
\textbf{Thm.} The mixed seal commutator $[J_i,K_j]$ is a bilinear in two seal triplets. Let $C_3$ be the cyclic permutation of the seal index, generated by the $SO(3)$ rotation by $2\pi/3$ about the $(1,1,1)$ axis — a subgroup of the $SO(3)$ under which $\{J_i\}$ and $\{K_i\}$ transform as vectors [CT1.8]; inner form: $U=u\otimes u$ with $u=\tfrac12\bigl(I_2-i(\sigma_1+\sigma_2+\sigma_3)\bigr)$, exact over $\mathbb Q(i)$: $U^3=I_4$, $UJ_iU^\dagger=J_{i+1}$, $UK_iU^\dagger=K_{i+1}$ (indices mod $3$). Then:

%dsp-level:0
\[
(1)\quad\text{bilinear}
\]

%dsp-level:0
\[
= \text{singlet}\oplus\text{doublet}:\ \ \text{the singlet is the}\ \Sigma\ \text{part (}C_3\text{-invariant); the doublet is the}\ \hat E^{ij}\ \text{part (}C_3\text{-cyclically permuted).}
\]

Concretely: the diagonal triple $[J_i,K_i]/i=\Sigma-E^{ii}=\tfrac23\Sigma-\hat E^{ii}$ splits into the fixed direction $\Sigma$ and the zero-sum plane of $\bigl(\hat E^{11},\hat E^{22},\hat E^{33}\bigr)$ — $\sum_i\hat E^{ii}=0$ — on which the cyclic shift $P$ acts with minimal polynomial $x^2+x+1$, irreducible over $\mathbb Q$: the doublet.

%dsp-level:0
\[
(2)\quad\text{Democratic diagonalization: any}\ C_3\text{-covariant Hermitian operator on the triplet —}\ [H,P]
\]

%dsp-level:0
\[
= 0,
\]

%dsp-level:0
\[
\ H^\dagger
\]

%dsp-level:0
\[
= H,
\]

%dsp-level:0
\[
\ P\ \text{the cyclic shift — has the form}
\]

%dsp-level:0
\[
H=a\,I_3+z\,P+\bar z\,P^{T}\qquad(a\in\mathbb R,\ z\in\mathbb C),
\]

%dsp-level:0
\[
\operatorname{spec}H:\quad\lambda_i=A+B\cos\bigl(\theta+2\pi i/3\bigr),\qquad i=0,1,2,\qquad A=a,\ \ B=2|z|,\ \ \theta=\arg z,
\]

equivalently the depressed cubic $\ (\lambda-A)^3-\tfrac{3B^2}{4}(\lambda-A)-\tfrac{B^3}{4}\cos3\theta=0$.

%dsp-level:0
\[
(3)\quad\text{Normalization:}\ \ A=1\ \ \Bigl[\sum_i\lambda_i=3A=3=\dim F_{J^\pm}\Bigr];
\]

%dsp-level:0
\[
\qquad B=\sqrt2\ \ \Bigl[\sum_i(\lambda_i-A)^2=B^2\sum_i\cos^2(\theta+2\pi i/3)=B^2\cdot\tfrac32=3\implies B^2=2\Bigr];
\]

%dsp-level:0
\[
\theta=\delta:\ \text{the mixing angle, determined in FT.112.}
\]

\textbf{Req.} CT1.2, CT1.8, CT0.12 (Part 1); FT.109; standard $C_3$ representation theory (commutant of the cyclic shift, Fourier diagonalization, triple-angle identity) [math]; exact $\mathbb Q(i)$ conjugation, exact symbolic identities in $\mathbb Q(\sqrt3)[c,s]/(c^2{+}s^2{-}1)$, exact $\mathbb Q$ linear algebra, Zariski-dense rational sampling of the charpoly identity [machine: \texttt{checks/\allowbreak{}verify\_\allowbreak{}H\_\allowbreak{}koide\_\allowbreak{}channel.py} §D].

\textbf{Pf.}

%dsp-level:0
\[
(1):\ U^3=I_4,\ UU^\dagger=I_4,\ \det u=1\ [\text{exact over}\ \mathbb Q(i);\ \text{machine}];
\]

%dsp-level:0
\[
\qquad UJ_iU^\dagger=J_{i+1},\ UK_iU^\dagger=K_{i+1}\ [\mathrm{CT1.8}\ \text{exponentiated};\ \text{machine}]
\]

%dsp-level:0
\[
U\Sigma U^\dagger=\Sigma;
\]

%dsp-level:0
\[
\qquad UE^{ii}U^\dagger=E^{i+1,i+1};
\]

%dsp-level:0
\[
\qquad U\hat E^{ii}U^\dagger=\hat E^{i+1,i+1};
\]

%dsp-level:0
\[
\qquad U\hat E^{12}U^\dagger=\hat E^{23},\ U\hat E^{23}U^\dagger=\hat E^{13},\ U\hat E^{13}U^\dagger=\hat E^{12}\ [\text{machine}]
\]

%dsp-level:0
\[
\bigl[J_i,K_j\bigr]\ \text{equivariance}:\ U\,[J_i,K_j]\,U^\dagger=[J_{i+1},K_{j+1}]\ \ \forall(i,j)\ [\text{machine, all 9 pairs}]
\]

%dsp-level:0
\[
\sum_i\hat E^{ii}=\sum_iE^{ii}-3\cdot\tfrac13\Sigma=\Sigma-\Sigma=0\ [\mathrm{FT.109}(2);\ \mathrm{CT1.2}\ \text{traceless split}]
\]

%dsp-level:0
\[
P^3=I_3,\ \operatorname{tr}P=0;
\]

%dsp-level:0
\[
\qquad P|_{\{\sum x_i=0\}}:\ \text{charpoly}=\text{minimal polynomial}=x^2+x+1\ (P|_{\text{plane}}\ne I);
\]

%dsp-level:0
\[
\qquad x^2+x+1\ \text{has no rational root}\ (\pm1\mapsto3\ne0)\implies\text{irreducible over}\ \mathbb Q:\ \text{doublet}\ [\text{machine}]
\]

%dsp-level:0
\[
(2):\ [X,P]
\]

%dsp-level:0
\[
= 0
\]

%dsp-level:0
\[
\iff X\ \text{circulant}:\ \dim_{\mathbb Q}\operatorname{comm}(P)
\]

%dsp-level:0
\[
= 3,
\]

%dsp-level:0
\[
\ \text{basis}\ \{I_3,P,P^2\}\ [\text{exact}\ \mathbb Q\ \text{elimination};\ \text{machine}]
\]

%dsp-level:0
\[
H=H^\dagger\in\operatorname{comm}(P)\iff H=aI_3+zP+\bar zP^T;
\]

%dsp-level:0
\[
\qquad \operatorname{charpoly}H
\]

%dsp-level:0
\[
= (x-a)^3-3|z|^2(x-a)-2\operatorname{Re}(z^3)\ [\text{machine at 24 exact rational points}\ (a,\operatorname{Re}z,\operatorname{Im}z):\ \text{Zariski-dense}\implies\text{identical}]
\]

%dsp-level:0
\[
A=a,\ B=2|z|,\ \cos3\theta=\operatorname{Re}(z^3)/|z|^3:\quad\tfrac{3B^2}4=3|z|^2,\quad\tfrac{B^3}4\cos3\theta=2\operatorname{Re}(z^3)\implies\text{depressed cubic form}\ \checkmark
\]

%dsp-level:0
\[
\prod_{i=0}^{2}\bigl(x-A-B\cos(\theta+\tfrac{2\pi i}3)\bigr)
\]

%dsp-level:0
\[
= (x-A)^3-\tfrac{3B^2}4(x-A)-\tfrac{B^3}4\cos3\theta\ [\text{exact in}\ \mathbb Q(\sqrt3)[A,B,x,c,s]/(c^2{+}s^2{-}1);\ \text{machine}];
\]

%dsp-level:0
\[
\qquad \cos3\theta=4c^3-3c\ [\text{same ring}]
\]

%dsp-level:0
\[
\text{covariance witness}:\ M(\theta):
\]

%dsp-level:0
\[
= A\,
\]

%dsp-level:0
\[
I_3+B\bigl(\cos\theta\,D_c-\sin\theta\,D_s\bigr),
\]

%dsp-level:0
\[
\ \ D_c
\]

%dsp-level:0
\[
= \operatorname{diag}\bigl(1,-\tfrac12,-\tfrac12\bigr),
\]

%dsp-level:0
\[
\ \ D_s
\]

%dsp-level:0
\[
= \operatorname{diag}\bigl(0,\tfrac{\sqrt3}2,-\tfrac{\sqrt3}2\bigr):
\]

%dsp-level:0
\[
\operatorname{spec}M(\theta)=\bigl\{A+B\cos(\theta+\tfrac{2\pi i}3)\bigr\}_{i=0,1,2};
\]

%dsp-level:0
\[
\qquad M\bigl(\theta+\tfrac{2\pi}3\bigr)=P^{T}M(\theta)\,P\ [\text{entry-wise exact in the same ring};\ \text{machine}]
\]

%dsp-level:0
\[
(3):\ \sum_i\cos(\theta+\tfrac{2\pi i}3)
\]

%dsp-level:0
\[
= 0,
\]

%dsp-level:0
\[
\quad\sum_i\cos^2(\theta+\tfrac{2\pi i}3)
\]

%dsp-level:0
\[
= \tfrac32\ \ \forall\theta\ [\text{exact identities in}\ \mathbb Q(\sqrt3)[c,s]/(c^2{+}s^2{-}1):\ \text{coefficients}\ (1,-\tfrac12,-\tfrac12),\ (0,\tfrac{\sqrt3}2,-\tfrac{\sqrt3}2);\ \text{machine}]
\]

%dsp-level:0
\[
\sum_i\lambda_i=3A=3\implies A=1\ [\mathrm{CT0.12}:\ \dim F_{J^\pm}=3];
\]

%dsp-level:0
\[
\qquad \sum_i(\lambda_i-A)^2=B^2\cdot\tfrac32=3\implies B^2=2,\ B=\sqrt2\ [(\sqrt2)^2=2\ \text{exact in}\ \mathbb Q(\sqrt2);\ \text{machine}].\ \blacksquare
\]

\textbf{Rmk.} This is representation theory of $C_3$, not the Koide ansatz. Any $C_3$-covariant triplet operator has this spectrum. The specific form $r_i=1+\sqrt2\cos(\delta+2\pi i/3)$ follows from the normalization conditions alone.

\textbf{Phys.} A $120^\circ$ rotation about the diagonal axis cycles the three seal directions, and the rotation is inner in the pair algebra: it conjugates each seal generator into the next. The mixed bilinear therefore splits into one rotation-invariant piece — the exchange scalar — and one irreducible rotating plane, the traceless diagonal tensors. Requiring the triplet's trace and its variance to each equal the seal dimension $3$ fixes the spectrum's offset and spread to $1$ and $\sqrt2$; the only remaining content is the phase $\delta$.

\vspace{2pt}\noindent\rule{\textwidth}{0.4pt}\vspace{2pt}

\end{theorem}

\begin{theorem}[{\texorpdfstring{FT.112 — The mixing angle}{FT.112 - The mixing angle}}]\label{tag:FT-112}
\textbf{Thm.} The mixing angle of the $C_3$-covariant seal bilinear is the exact rational

%dsp-level:0
\[
\delta=\frac{\sum_i[J_i,K_i]\big|_\Sigma}{\dim F_{J^\pm}\cdot\operatorname{tr}T(1)}=\frac{2}{3\cdot3}=\frac29,
\]

where the numerator $2$ is the scalar coefficient of the diagonal seal commutator [FT.109(3): $\sum_i[J_i,K_i]=2i\Sigma$] and the denominator $9$ is the bound-chain arena dimension [FT.110(1)].

\textbf{Rmk (cross-check, not a derivation of equality).} The same rational $2/3$ appears in three independent frozen results: $(p^*_\infty)^2=2/3$ [FT.70(3): the polygon phase boundary squared, exact in $\mathbb Z[\sqrt6]$]; $\delta_v/\pi=2/3$ [FT.63(4): the octahedron vertex defect, exact in $\pi$-units]; $f_{\mathrm{charge}}(\delta h=3)=2/3$ [FT.14: the $3$-chain charge fraction, exact rational]. The mixing angle is $\delta=2/9=(2/3)/3$: the polygon phase boundary divided by the number of generations (the seal dimension). Recorded as a cross-check — three independent appearances of the same rational — not as a derivation of their equality.

\textbf{Req.} FT.109, FT.110; FT.14 (Tiers F.1–F.11), FT.63(4) (Tiers F.20–F.22), FT.70(3) (Tiers F.23–F.25) for the cross-check row; exact rational arithmetic, exact $\mathbb Z[\sqrt6]$ arithmetic [math; machine: \texttt{checks/\allowbreak{}verify\_\allowbreak{}H\_\allowbreak{}koide\_\allowbreak{}channel.py} §E].

\textbf{Pf.}

%dsp-level:0
\[
\delta=\frac{2}{\dim F_{J^\pm}\cdot\operatorname{tr}T(1)}=\frac{2}{3\cdot3}=\frac29\ [\mathrm{FT.109}(3);\ \mathrm{FT.110};\ \text{exact Fractions; machine §E}]
\]

%dsp-level:0
\[
\text{cross-check rows}:\quad\bigl(\tfrac{\sqrt6}3\bigr)^2=\tfrac23\ [\text{exact in}\ \mathbb Z[\sqrt6];\ \mathrm{FT.70}(3)];
\]

%dsp-level:0
\[
\qquad \frac{\delta_v}\pi=2-\frac{d(L-2)}{L}\Big|_{d=4,\,L=3}=2-\tfrac43=\tfrac23\ [\mathrm{FT.63}(4);\ \text{exact integers}]
\]

%dsp-level:0
\[
f_{\mathrm{charge}}(\delta h=3)=\frac{N_{\mathrm{charge}}}{N_{\mathrm{visible}}}=\frac23\ [\mathrm{FT.14}\ \text{table row}\ \delta h=3];\qquad\frac{2/3}{3}=\frac29\ [\text{exact}].\ \blacksquare
\]

\textbf{Phys.} One exact integer from the seal algebra over one exact integer from the arena: the phase is $2/9$, a dimensionless rational with no adjustable content — the bilinear's coefficient divided by the arena's dimension. Its relation to the three independent $2/3$ results is a coincidence of rationals recorded for the ledger; nothing is calibrated to anything.

\vspace{2pt}\noindent\rule{\textwidth}{0.4pt}\vspace{2pt}

\end{theorem}

\begin{theorem}[{\texorpdfstring{FT.113 — The bilinear spectrum and the mass-squared ratios}{FT.113 - The bilinear spectrum and the mass-squared ratios}}]\label{tag:FT-113}
\textbf{Thm.} With $\delta=2/9$ [FT.112] and $A=1$, $B=\sqrt2$ [FT.111(3)], the three eigenvalues of the $C_3$-covariant bilinear are

%dsp-level:0
\[
r_i=1+\sqrt2\cos\Bigl(\frac29+\frac{2\pi i}{3}\Bigr),\qquad i=0,1,2,
\]

with exact radical-trigonometric closed forms ($c=\cos\tfrac29$, $s=\sin\tfrac29$):

%dsp-level:0
\[
r_0=1+\sqrt2\,c;\qquad r_1=1-\frac{\sqrt2\,c+\sqrt6\,s}{2};\qquad r_2=1-\frac{\sqrt2\,c-\sqrt6\,s}{2};
\]

%dsp-level:0
\[
\sum_{i}r_i=3;\qquad\sum_{i<j}r_ir_j=\frac32;\qquad 2(r-1)^3-3(r-1)-\sqrt2\cos\frac23=0\ \ \text{for each}\ r=r_i\ \ [3\delta=\tfrac23].
\]

The mass-squared ratios are $r_i^2$. Ordered by magnitude — $r_1^2<r_2^2<r_0^2$, enclosure-proven —

%dsp-level:0
\[
m_1^2:m_2^2:m_3^2=r_{\min}^2:r_{\mathrm{mid}}^2:r_{\max}^2=r_1^2:r_2^2:r_0^2=1:\frac{r_2^2}{r_1^2}:\frac{r_0^2}{r_1^2},
\]

%dsp-level:0
\[
\frac{r_2^2}{r_1^2}=206.770315972729\ldots,
\]

%dsp-level:0
\[
\qquad \frac{r_0^2}{r_1^2}=3477.472837104\ldots
\]

%dsp-level:0
\[
\qquad [\text{rational trigonometric enclosures, width}<2^{-50};\ \text{machine}].
\]

The overall scale $\mu$ in $m_i^2=\mu\,r_i^2$ is undeclared (P4-class free internal parameter): the ratios are this tag's only content. The $m^2$ vocabulary is the corpus's own inertial read [L3(4), L4: dimensionless entropy deficit in nats; $\neg$a particle identification].

\textbf{Req.} FT.111, FT.112; L3, L4 (\textbf{B}, the $m^2$ read vocabulary); P4 (\textbf{A}, the free-scale clause); exact rational trigonometric enclosures (Taylor remainder bounds), exact integer square-root enclosures (Newton with certificates), monotone rational interval arithmetic [math; machine: \texttt{checks/\allowbreak{}verify\_\allowbreak{}H\_\allowbreak{}koide\_\allowbreak{}channel.py} §F].

\textbf{Pf.}

%dsp-level:0
\[
r_i=A+B\cos(\delta+\tfrac{2\pi i}3)\ \text{at}\ A=1,\ B=\sqrt2,\ \delta=\tfrac29\ [\mathrm{FT.111}(3),\ \mathrm{FT.112}]
\]

%dsp-level:0
\[
\cos\bigl(\tfrac29+\tfrac{2\pi}3\bigr)=-\tfrac c2-\tfrac{\sqrt3}2s,
\]

%dsp-level:0
\[
\qquad \cos\bigl(\tfrac29+\tfrac{4\pi}3\bigr)=-\tfrac c2+\tfrac{\sqrt3}2s\ [\text{addition formulas}];
\]

%dsp-level:0
\[
\qquad \sqrt2\cdot\tfrac{\sqrt3}2=\tfrac{\sqrt6}2
\]

%dsp-level:0
\[
\sum_ir_i=3+\sqrt2\sum_i\cos(\delta+\tfrac{2\pi i}3)=3\ [\mathrm{FT.111}(3)\ \text{sum identity}];
\]

%dsp-level:0
\[
\qquad \text{Vieta pair sum and depressed cubic at}\ B^2=2,\ \cos3\delta=\cos\tfrac23\ [\mathrm{FT.111}(2)\ \text{charpoly identity}]
\]

%dsp-level:0
\[
c,s\ \text{enclosed by Taylor partial sums with exact rational remainder bounds (widths}<2^{-80}\text{)};
\]

%dsp-level:0
\[
\qquad \sqrt2,\sqrt6,\cos\tfrac23\ \text{enclosed exactly (integer certificates)}
\]

%dsp-level:0
\[
r_1^2<r_2^2<r_0^2\ [\text{disjoint enclosures}];
\]

%dsp-level:0
\[
\qquad \text{ratios by monotone interval division, width}<2^{-50}\ [\text{machine §F}];
\]

%dsp-level:0
\[
\qquad 0\in 2y_i^3-3y_i-\sqrt2\cos\tfrac23\ \text{enclosure for each}\ y_i=r_i-1\ [\text{machine}].\ \blacksquare
\]

\textbf{Phys.} The spectrum carries one large, one middle, and one nearly-vanished eigenvalue: a hierarchy produced by a rational phase inside a cosine, with closed forms in $\sqrt2$, $\sqrt6$, and $\cos\tfrac29$ alone. The proximity of the two ratios to the supplied external hierarchy is quantified in FT.114 and used nowhere else. Every quantity here is dimensionless; the scale $\mu$ converting ratios into a read is a free internal parameter (P4), and the $m^2$ name is the entropy-deficit read of L3(4)/L4 — never a particle.

\vspace{2pt}\noindent\rule{\textwidth}{0.4pt}\vspace{2pt}

\end{theorem}

\phantomsection
\section[Tier F.38 — The Colour Channel: The -Selected , Confinement…]{\texorpdfstring{Tier F.38 — The Colour Channel: The $\varepsilon$-Selected $C_3$, Confinement Blindness, and the Triplet-Attached Mixing Angle}{Tier F.38 - The Colour Channel: The -Selected , Confinement Blindness, and the Triplet-Attached Mixing Angle}}

\begin{theorem}[{\texorpdfstring{FT.116 — The relabelling lattice, the two-level obstruction, and the $\varepsilon$-selection of $C_3$}{FT.116 - The relabelling lattice, the two-level obstruction, and the -selection of}}]\label{tag:FT-116}
\textbf{Thm.} On the slot carrier $\mathcal H^{(3)}_{\mathrm{slot}}\cong\mathbb C^3$ [CT3.2] with the classical relabelling embedding $P_\pi|i\rangle=|\pi(i)\rangle$, $(P_\pi)^i{}_j=\delta^i_{\pi(j)}$, $S_3\hookrightarrow U(3)$ [CT3.3]:

%dsp-level:0
\[
(1)\quad S_3\ \text{has exactly six subgroups, of orders}\ (1,2,2,2,3,6):\ \{e\},\ \text{three}\ \mathbb Z_2=\langle\text{transposition}\rangle,\ A_3\cong C_3,\ S_3
\]

%dsp-level:0
\[
\qquad [\text{machine: closure of all}\ 2^6\ \text{generator subsets}];
\]

%dsp-level:0
\[
(2)\quad\textbf{commutant dimensions}\ (\text{Hermitian real form}):\qquad\{e\}\mapsto9,\qquad\mathbb Z_2\mapsto5,\qquad C_3\mapsto3,\qquad S_3\mapsto2;
\]

%dsp-level:0
\[
(3)\quad C_3\text{-commutant}=\{aI_3+zP+z^{*}P^{T}:\ a\in\mathbb R,\ z\in\mathbb C\}\ (P=P_{(123)}),
\]

%dsp-level:0
\[
\qquad S_3\text{-commutant}=\operatorname{span}\{I_3,J\}\ (J=\text{all-ones});
\]

%dsp-level:0
\[
(4)\quad\textbf{the two-level obstruction:}\ \ \operatorname{spec}(aI_3+bJ)
\]

%dsp-level:0
\[
= \{a+3b,\ a,\ a\}\ \text{exactly}\ \forall(a,b)\in\mathbb R^2\ (P_\pi\text{-covariance }\Rightarrow\text{ at most two distinct eigenvalues});
\]

%dsp-level:0
\[
(5)\quad\textbf{reflection rigidity:}\ \ THT=H\ \text{for a transposition}\ T\iff z\in\mathbb R
\]

%dsp-level:0
\[
\qquad [THT=aI_3+zP^{T}+z^{*}P:\ \text{conjugation by a transposition inverts the}\ 3\text{-cycle}];
\]

%dsp-level:0
\[
\operatorname{disc}=108\,(\operatorname{Im}z^3)^2=108\,(\operatorname{Im}z)^2\bigl(3(\operatorname{Re}z)^2-(\operatorname{Im}z)^2\bigr)\ \text{vanishes}\iff z^3\in\mathbb R\iff\arg z\in\tfrac\pi3\mathbb Z
\]

%dsp-level:0
\[
\qquad [\text{the degeneracy locus: three lines through the origin}];
\]

%dsp-level:0
\[
S_3\text{-covariance forces}\ z\in\mathbb R\ (\arg z\in\{0,\pi\})\subset\tfrac\pi3\mathbb Z\implies\text{every}\ S_3\text{-covariant circulant has a double root};
\]

%dsp-level:0
\[
\ \text{a three-line spectrum requires}\ \operatorname{Im}z^3
\]

%dsp-level:0
\[
\ne 0,
\]

%dsp-level:0
\[
\ \text{i.e.}\ S_3\text{-covariance must be broken to the cyclic subgroup};
\]

%dsp-level:0
\[
(6)\quad\textbf{the}\ \varepsilon\textbf{-selection:}\ \ \varepsilon_{ijk}P^i{}_aP^j{}_bP^k{}_c
\]

%dsp-level:0
\[
= \det(P_\pi)\,
\]

%dsp-level:0
\[
\varepsilon_{abc}
\]

%dsp-level:0
\[
= \operatorname{sgn}(\pi)\,
\]

%dsp-level:0
\[
\varepsilon_{abc}\ \forall\pi
\]

%dsp-level:0
\[
\in S_3\ [\mathrm{CT3.15}\ \text{on permutation matrices}];
\]

%dsp-level:0
\[
\qquad \det P_\pi=\operatorname{sgn}(\pi);
\]

%dsp-level:0
\[
\operatorname{Stab}_{S_3}(\varepsilon)=\{\pi:\det P_\pi=1\}=A_3=C_3=S_3\cap SU(3);
\]

%dsp-level:0
\[
\qquad \text{the}\ \varepsilon\text{-preserving subgroups are exactly}\ \{e\}\subset A_3,\ \text{so}\ A_3\ \text{is the unique maximal one};
\]

%dsp-level:0
\[
(7)\quad\textbf{seal-side companion census:}\ \ \text{on the seal triplet,}\ T\cong A_4\ (|T|=12,\ T\subset SO(3))\ \text{has commutant}\ \mathbb R\cdot I_3\ (\text{full degeneracy});
\]

%dsp-level:0
\[
\ \ V_4\ \text{has commutant the diagonal matrices}\ (3\ \text{parameters, no cyclic relation});
\]

%dsp-level:0
\[
\ \ C_3\ \text{(the}\ 2\pi/3\ \text{rotation about}\ (1,1,1)\ [\mathrm{CT1.8}])\ \text{has commutant the circulants}\ (3\ \text{parameters, fixed democratic eigenvectors}).
\]

\textbf{Req.} CT3.1, CT3.2, CT3.3, CT3.15 (Part 2); D3.9 (Parts 1–2); CT1.8 (Part 1); FT.111 (\textbf{Tier F.37}, the $C_3$ democratic form and its normalization); subgroup-lattice closure, exact nullspace dimension over $\mathbb Q$, Leibniz determinant identity, charpoly of the circulant and of $aI_3+bJ$, discriminant of the depressed cubic [math; machine: \texttt{checks/\allowbreak{}verify\_\allowbreak{}H\_\allowbreak{}colour\_\allowbreak{}channel.py} §A].

\textbf{Pf.}

%dsp-level:0
\[
(1):\ \text{closure of each of the}\ 2^6\ \text{generator subsets under}\ \pi\circ\rho\ \text{yields}\ 6\ \text{distinct subgroups; orders sorted}\ (1,2,2,2,3,6)\ [\text{machine}]
\]

%dsp-level:0
\[
(2),
\]

%dsp-level:0
\[
(3):\ \{X:[X,P_\pi]=0\ \forall\pi\in G\}\ \text{solved as an exact}\ 9\text{-unknown nullspace over}\ \mathbb Q\ \text{per subgroup};
\]

%dsp-level:0
\[
\ \dim
\]

%dsp-level:0
\[
= (9,5,3,2);
\]

%dsp-level:0
\[
\ \text{transpose-closed}\ \forall G\ [\text{machine}]
\]

%dsp-level:0
\[
C_3:\ \text{basis}\ \{I_3,P,P^2\}\ \text{(real form}\ \{I_3,\ P+P^T,\ P-P^T\}\ \leftrightarrow\ \{a,\ 2\operatorname{Re}z,\ -2\operatorname{Im}z\});
\]

%dsp-level:0
\[
\qquad S_3:\ \text{nullspace basis}\ \{I_3,J\}\ [\text{machine: exact rank}\ 2]
\]

%dsp-level:0
\[
(4):\ \operatorname{charpoly}(aI_3+bJ)
\]

%dsp-level:0
\[
= (x-a-3b)(x-a)^2
\]

%dsp-level:0
\[
= x^3-(3a+3b)x^2+(3a^2+6ab)x-(a^3+3a^2b)\ \text{at}\ 3\ \text{exact rational points}\ (a,b)\ [\text{machine}]
\]

%dsp-level:0
\[
(5):\ THT=aI_3+zP^{T}+z^{*}P\ (TPT=P^{T}:\ \text{a transposition inverts the}\ 3\text{-cycle});\qquad THT=H\iff z=z^{*}\iff\operatorname{Im}z=0
\]

%dsp-level:0
\[
\operatorname{charpoly}(aI_3+zP+z^{*}P^{T})=(x-a)^3-3|z|^2(x-a)-2\operatorname{Re}(z^3)\ [\text{machine at}\ 15\ \text{exact rational points}];
\]

%dsp-level:0
\[
\qquad \operatorname{disc}=-4(-3|z|^2)^3-27(-2\operatorname{Re}z^3)^2=108\bigl(|z|^6-(\operatorname{Re}z^3)^2\bigr)=108(\operatorname{Im}z^3)^2\ [\text{exact}]
\]

%dsp-level:0
\[
\operatorname{Im}z^3
\]

%dsp-level:0
\[
= 3(\operatorname{Re}z)^2\operatorname{Im}z-(\operatorname{Im}z)^3
\]

%dsp-level:0
\[
= \operatorname{Im}z\,
\]

%dsp-level:0
\[
\bigl(3(\operatorname{Re}z)^2-(\operatorname{Im}z)^2\bigr)\ [\text{polynomial identity at exact rational points, Zariski-dense};\ \text{machine}]
\]

%dsp-level:0
\[
\operatorname{disc}=0\iff(\operatorname{Im}z^3)^2=0\iff\operatorname{Im}z=0\ \vee\ (\operatorname{Im}z)^2=3(\operatorname{Re}z)^2\iff z^3\in\mathbb R\iff\arg z\in\tfrac\pi3\mathbb Z;
\]

%dsp-level:0
\[
\qquad z\in\mathbb R\implies\operatorname{disc}=0\implies\text{a double root}\ [\text{machine at every real-}z\ \text{point}]
\]

%dsp-level:0
\[
(6):\ \varepsilon_{ijk}P^i{}_aP^j{}_bP^k{}_c=\operatorname{sgn}(\pi)\varepsilon_{abc}\ \text{entrywise on all}\ 27\ \text{triples}\times6\ \text{relabellings};
\]

%dsp-level:0
\[
\qquad \det P_\pi=\operatorname{sgn}(\pi)\ [\text{exact}\ \mathbb Q\ \text{expansion};\ \text{machine}]
\]

%dsp-level:0
\[
\pi\in\operatorname{Stab}(\varepsilon)\iff\operatorname{sgn}(\pi)=1\iff\pi\in A_3;
\]

%dsp-level:0
\[
\qquad \text{lattice scan}: \{g\le S_3:g\subseteq A_3\}=\{\{e\},A_3\}\ [\text{machine}];
\]

%dsp-level:0
\[
\qquad A_3=S_3\cap SU(3)\ [\det=1]
\]

%dsp-level:0
\[
(7):\ T
\]

%dsp-level:0
\[
= \langle P_{(123)},
\]

%dsp-level:0
\[
\ \operatorname{diag}(1,-1,-1),
\]

%dsp-level:0
\[
\ \operatorname{diag}(-1,1,-1),
\]

%dsp-level:0
\[
\ \operatorname{diag}(-1,-1,1)\rangle:\ |T|
\]

%dsp-level:0
\[
= 12,
\]

%dsp-level:0
\[
\ \det
\]

%dsp-level:0
\[
= 1\ \forall T\ [\text{machine closure}];
\]

%dsp-level:0
\[
\qquad \dim\operatorname{comm}(T)=1,\ \dim\operatorname{comm}(V_4)=3,\ \dim\operatorname{comm}(C_3)=3\ [\text{exact nullspaces};\ \text{machine}].\ \blacksquare
\]

\textbf{Phys.} Ask the colour sector which of its own symmetries may carry a mass family, and the answer is forced twice over. The full relabelling group is too symmetric: an $S_3$-covariant operator on three slots has a doubly-degenerate spectrum, identically, for every choice of its two parameters, and any operator invariant under a single transposition has a double root — a three-line family cannot be $S_3$-covariant, and the phase $\arg z$ is precisely the datum that breaks it. Too little symmetry is equally useless: with no covariance at all nine parameters are free and nothing is predicted. The cyclic subgroup is the unique intermediate, and it is not selected by hand — the baryon seal $\varepsilon_{ijk}$, the tensor that makes three triplets into an observable [CT3.15, CT3.18], is preserved exactly by the even relabellings, so $A_3\cong C_3=S_3\cap SU(3)$ is the largest relabelling symmetry a sealed object can have. The same $C_3$ that FT.111 found inside the seal algebra is therefore also the maximal covariance group of the colour algebra, reached from the opposite side: there by the $2\pi/3$ rotation about $(1,1,1)$, here by the Levi-Civita tensor. The companion census closes the search on the seal side as well — the tetrahedral group is too big (its commutant is the scalars, so every $T$-covariant operator is fully degenerate) and the Klein group is the wrong shape (diagonal commutant, no cyclic relation among eigenvalues).

\vspace{2pt}\noindent\rule{\textwidth}{0.4pt}\vspace{2pt}

\end{theorem}

\medskip\noindent\textit{\texorpdfstring{FT.117 — Confinement blindness: the one-body no-go on colour-sealed singlets}{FT.117 - Confinement blindness: the one-body no-go on colour-sealed singlets}.}\ Structural-constraints record: see Appendix B, Section B.3.3 of the Addendum.\par\medskip

\begin{definition}[{\texorpdfstring{Def-H.42 — The colour-dressed arena}{Def-H.42 - The colour-dressed arena}}]\label{tag:Def-H-42}
\textbf{Def.} A triplet-attached bound chain [D3.8] carries, simultaneously and independently: the seal triplet at its termini [CT0.12, CT0.7(4)], the polygonal transfer structure along its length [FT.68(2)], and the colour slot at its attachment [CT3.2: $\mathbb C^3$]. The \textbf{colour-dressed arena} is their minimal combined counting space:

%dsp-level:0
\[
\mathcal H_c^{(q)}\ :=\ \mathcal H_c\otimes\mathbb C^3_{\mathrm{slot}}\qquad[\mathrm{Def\text{-}H.41}],\qquad
\dim_{\mathrm{eff}}\mathcal H_c^{(q)}=3\times3\times3=27\ \ [\mathrm{FT.118}(2)],
\]

entering exactly once — as the denominator of the triplet-attached mixing angle [FT.119: $\delta_q=\tfrac2{27}$ with FT.118(4)'s spectator numerator $2$]. Singlet-sealed chains keep Def-H.41's $9$: the invariant lines of a sealed object are $1$-dimensional [FT.117(5)], so its colour factor is $1$. \textbf{Species and bars.} As Def-H.41: a counting space — no inner product, no dynamics, no state, no Hamiltonian; dimension $\neg$a state count; the overall scale $\mu$ undeclared (P4-class, Cal-F); external calibration barred [P1, P5; FT.120 — forward pointer; Census row]; the adjoint-dressed numerator $8/3$ [FT.118(6)] and the alternative table rows [FT.119(7)] are logged alternatives, $\neg$adopted; particle identification barred. \textbf{Locus remark (structural).} Colour sits in an orthogonal stack block of the force arena [FT.93(4); CT0.12 Rmk: $\mathbb C^4\otimes\mathbb C^3$ is not a corpus construction], and EW.32's seed is colour-neutral by theorem — this definition counts D3.8's attachment choices in the chain's own bookkeeping, $\neg$a state-space composition of the force arena. The derived content compatible with both records is FT.118(4)'s spectator dressing $\sum_i[J_i\otimes I_3,\,K_i\otimes I_3]=2i\,\Sigma\otimes I_3$.

\end{definition}

\begin{theorem}[{\texorpdfstring{FT.118 — The colour-dressed arena and the spectator numerator}{FT.118 - The colour-dressed arena and the spectator numerator}}]\label{tag:FT-118}
\textbf{Thm.} With the transfer rule $T(1)[s][s']=[\,|s|+|s'|\le3\,]$, $s,s'\in\{0,1\}^2$ [FT.68(2)], re-derived at the point of use:

%dsp-level:0
\[
(1)\quad\operatorname{tr}T(1)=3
\]

%dsp-level:0
\[
\qquad [\text{the single excluded diagonal pair is}\ (|s|,|s'|)=(2,2)];
\]

%dsp-level:0
\[
\qquad \operatorname{charpoly}T(1)=x^4-3x^3-3x^2=x^2\bigl(x^2-3x-3\bigr)\ [\mathrm{FT.68}(2);\ \text{machine: Faddeev–LeVerrier over}\ \mathbb Z];
\]

%dsp-level:0
\[
(2)\quad\textbf{the colour-dressed arena:}\qquad\dim_{\mathrm{eff}}\mathcal H_c^{(q)}=\dim F_{J^\pm}\times\operatorname{tr}T(1)\times N_{\text{colour}}=3\times3\times3=27,
\]

%dsp-level:0
\[
\text{the third factor being}\ N_{\text{colour}}=3=\log_2C_{\mathrm{cap}}-1=\log_216-1\ [\mathrm{CT3.20};\ \mathrm{CT0.5},\ \mathrm{CT0.6},\ \mathrm{CT0.11}];
\]

%dsp-level:0
\[
\qquad \text{scope: triplet-attached chains}\ [\mathrm{D3.8}];
\]

%dsp-level:0
\[
(3)\quad\textbf{sealed attachments do not dress:}
\]

%dsp-level:0
\[
\qquad \text{colour factor of a sealed object}
\]

%dsp-level:0
\[
= 1\ [\mathrm{FT.117}(5):\ \text{the invariant lines are}\ 1\text{-dimensional}]\implies\dim_{\mathrm{eff}}\mathcal H_c^{\mathrm{sealed}}
\]

%dsp-level:0
\[
= 3\times3\times1
\]

%dsp-level:0
\[
= 9\ [\mathrm{FT.110}(1)\ \text{recovered}];
\]

%dsp-level:0
\[
(4)\quad\textbf{the spectator numerator:}
\]

%dsp-level:0
\[
\qquad \sum_{i=1}^{3}\bigl[J_i\otimes I_3,\ K_i\otimes I_3\bigr]=\Bigl(\sum_i[J_i,K_i]\Bigr)\otimes I_3=2i\,\Sigma\otimes I_3
\]

%dsp-level:0
\[
\qquad [12\times12\ \text{entrywise exact over}\ \mathbb Q(i);\ \mathrm{FT.109}\ \text{re-derived}];
\]

%dsp-level:0
\[
(5)\quad\text{the diagonal coefficient of the colour-dressed seal commutator is therefore exactly}\ 2\ \text{— unchanged from}\ \mathrm{FT.109}(3):\ \text{colour is a spectator of the seal algebra}\ [\mathrm{EW.32}:\ T_a\mathsf v=0;\ \mathrm{CT0.12}\ \mathrm{Rmk}:\ \text{colour sits in an orthogonal stack block}];
\]

%dsp-level:0
\[
(6)\quad\textbf{adjoint dressing (computed, NOT adopted):}
\]

%dsp-level:0
\[
\qquad \text{with}\ X^2
\]

%dsp-level:0
\[
= \textstyle\sum_aT_aT_a
\]

%dsp-level:0
\[
= C_FI_3
\]

%dsp-level:0
\[
= \tfrac43I_3\ [\text{machine: from the rational}\ E_{ij}\ \text{basis}]\ \text{and}\ [A\otimes X,\,B\otimes X]
\]

%dsp-level:0
\[
= [A,B]\otimes X^2,
\]

%dsp-level:0
\[
\sum_i\bigl[J_i\otimes X,\ K_i\otimes X\bigr]
\]

%dsp-level:0
\[
= 2i\,
\]

%dsp-level:0
\[
\Sigma\otimes X^2
\]

%dsp-level:0
\[
= \tfrac83\,
\]

%dsp-level:0
\[
i\,
\]

%dsp-level:0
\[
\Sigma\otimes I_3\implies\text{numerator}\ 2C_F
\]

%dsp-level:0
\[
= \tfrac83:\ \text{a different rational, and no corpus tag supplies the coupling};
\]

%dsp-level:0
\[
(7)\quad\textbf{the two candidate mixing angles:}
\]

%dsp-level:0
\[
\qquad \delta_q=\frac{2}{27}\ \ (\text{spectator, adopted in FT.119});
\]

%dsp-level:0
\[
\qquad \frac{2C_F}{27}=\frac{8/3}{27}=\frac8{81}\ \ (\text{adjoint-dressed, logged not adopted}).
\]

\textbf{Rmk (interface definition Def-H.42).} Clause (2)'s tensor product is the same species of interface definition as Def-H.41, with one new factor: a triplet-attached chain carries, simultaneously and independently, a seal triplet at its termini [CT0.12, CT0.7(4)], a polygonal transfer structure along its length [FT.68(2)], and a colour slot at its attachment [CT3.2, D3.8]; $\mathcal H_c^{(q)}$ is their minimal combined counting space. It is an arena — no inner product, no dynamics, no state, no Hamiltonian — and the count $27$ enters the tier exactly once, as the denominator of FT.119's mixing angle. \textbf{Tension logged:} FT.93(4) records colour as an orthogonal stack block (a $\oplus$-summand of $\mathcal K_\Sigma$, not a tensor factor [CT0.12 Rmk, CT0.14]) and the seed as colour-neutral by theorem [EW.32], so "$\mathbb C^4\otimes\mathbb C^3$" is not a corpus construction of the force arena; the declared factor counts D3.8's attachment choices in the chain's own bookkeeping, and clause (4)'s spectator dressing is the derived content compatible with both records. Clause (6) is recorded so that any future use of the adjoint-dressed numerator $8/3$ is a adoption rather than a silent one; the spectator form (4)–(5) is the derived content and needs no declaration beyond Def-H.42.

\textbf{Req.} CT3.2, CT3.20 (Part 2); CT0.5, CT0.6, CT0.11, CT0.12, CT0.14, CT0.7(4) (Part 1); D3.8 (Parts 1–2); EW.32 (Part 2); FT.68(2) (Tiers F.23–F.25); FT.93(4) (Tiers F.31–F.34, the tension note); FT.109, FT.110 (Tier F.37); FT.117(5) (this tier); Vieta coefficient read, exact integer characteristic polynomial, exact $\mathbb Q(i)$ Kronecker arithmetic, quadratic-Casimir normalization [math; machine: \texttt{checks/\allowbreak{}verify\_\allowbreak{}H\_\allowbreak{}colour\_\allowbreak{}channel.py} §C].

\textbf{Pf.}

%dsp-level:0
\[
(1):\ T(1)=\mathbb 1_{4\times4}-[\,\cdot\,]_{(2,2)}\ \text{on the diagonal}:\ \operatorname{tr}T(1)=4-1=3\ [\text{exact integers};\ \text{machine}]
\]

%dsp-level:0
\[
\operatorname{charpoly}T(1)
\]

%dsp-level:0
\[
= [1,-3,-3,0,0]\ \text{by Faddeev–LeVerrier over}\ \mathbb Z
\]

%dsp-level:0
\[
= x^4-3x^3-3x^2\ [\mathrm{FT.68}(2);\ \text{machine}]
\]

%dsp-level:0
\[
(2):\ \dim F_{J^\pm}=3\ [\mathrm{CT0.12}\ \text{table}:\ F_{J^+}=\mathcal E_{-+}=\operatorname{span}\{J_i\},\ F_{J^-}=\mathcal E_{--}=\operatorname{span}\{K_i\}];
\]

%dsp-level:0
\[
\qquad 3\cdot3\cdot3=27\ [\text{exact}]
\]

%dsp-level:0
\[
N_{\text{colour}}=3=\log_216-1\ [\mathrm{CT3.20}:\ \text{admissible}\iff2^{n+1}\le C_{\mathrm{cap}}=16\iff n\le3]
\]

%dsp-level:0
\[
(3):\ 3\cdot3\cdot1=9\ [\mathrm{FT.117}(5):\ \dim(3\otimes3\otimes3)^{\mathfrak{sl}(3)}=\dim(3\otimes\bar3)^{\mathfrak{sl}(3)}=1;\ \text{machine}];
\]

%dsp-level:0
\[
\qquad \mathrm{FT.110}(1)\ \text{value recovered}
\]

%dsp-level:0
\[
(4):\ [J_i\otimes I_3,\ K_i\otimes I_3]=[J_i,K_i]\otimes I_3\ \text{(Kronecker multiplicativity)};
\]

%dsp-level:0
\[
\qquad \sum_i[J_i,K_i]=2i\Sigma\ [\mathrm{FT.109}(3)\ \text{re-derived entrywise over}\ \mathbb Q(i);\ \text{machine}]
\]

%dsp-level:0
\[
\implies\sum_i[J_i\otimes I_3,K_i\otimes I_3]=(2i\Sigma)\otimes I_3\ [12\times12\ \text{equality, exact};\ \text{machine}]
\]

%dsp-level:0
\[
(5):\ \text{the}\ \Sigma\text{-coefficient of}\ (4)\ \text{is}\ 2\ [\text{unique nonzero basis component};\ \mathrm{FT.109}(3)];
\]

%dsp-level:0
\[
\qquad T_a\mathsf v=0\ [\mathrm{EW.32}]\ \text{and the stack-block orthogonality}\ [\mathrm{CT0.12}\ \mathrm{Rmk}]\ \text{are the structural reasons}
\]

%dsp-level:0
\[
(6):\ \tfrac12\bigl(\sum_{ij}E_{ij}E_{ji}-\tfrac13(\sum_iE_{ii})^2\bigr)
\]

%dsp-level:0
\[
= \tfrac12(3I_3-\tfrac13I_3)
\]

%dsp-level:0
\[
= \tfrac43I_3
\]

%dsp-level:0
\[
= C_FI_3\ [\text{exact}\ \mathbb Q;\ \text{machine}]
\]

%dsp-level:0
\[
[A\otimes X,\ B\otimes X]
\]

%dsp-level:0
\[
= (A\otimes X)(B\otimes X)-(B\otimes X)(A\otimes X)
\]

%dsp-level:0
\[
= AB\otimes X^2-BA\otimes X^2
\]

%dsp-level:0
\[
= [A,B]\otimes X^2\ [\text{machine at}\ X^2=\tfrac43I_3]
\]

%dsp-level:0
\[
\implies\sum_i[J_i\otimes X,K_i\otimes X]=2i\Sigma\otimes X^2=2i\cdot\tfrac43\,\Sigma\otimes I_3=\tfrac83i\,\Sigma\otimes I_3;
\]

%dsp-level:0
\[
\qquad 2\cdot\tfrac43=\tfrac83\ [\text{exact Fractions}]
\]

%dsp-level:0
\[
(7):\ \tfrac{2}{3\cdot3\cdot3}=\tfrac2{27};\qquad\tfrac{8/3}{27}=\tfrac8{81}\ [\text{exact Fractions};\ \text{machine}].\ \blacksquare
\]

\textbf{Phys.} A chain that grips a trivalent vertex carries one more independent bookkeeping than a chain that does not: which of the three slots it grips. That third ledger is the only new content of the rung-3 arena, and it multiplies — $3\times3$ becomes $3\times3\times3$. Sealing cancels the multiplication: because the baryon and meson invariant lines are one-dimensional, a sealed object's arena is back to nine, which is why FT.110's lepton-side count and this tier's differ at all. The numerator does not move. The seed that generates mass is colour-blind by theorem, so bracketing each spatial seal with its temporal partner and summing passes straight through the slot factor and returns the same integer $2$ over $\Sigma$ — the colour dressing is a spectator, and the alternative in which the strong generators themselves enter the bracket is a different rational ($8/3$), computed here and adopted nowhere, because no tag of the corpus couples the seal commutator to the slot algebra.

\vspace{2pt}\noindent\rule{\textwidth}{0.4pt}\vspace{2pt}

\end{theorem}

\begin{theorem}[{\texorpdfstring{FT.119 — The triplet-attached mixing angle $\delta_q=2/27$ and its spectrum}{FT.119 - The triplet-attached mixing angle and its spectrum}}]\label{tag:FT-119}
\textbf{Thm.} With the spectator numerator $2$ [FT.118(4)–(5)] over the colour-dressed arena $27$ [FT.118(2)], and the FT.111(3) normalization $A=1$, $B=\sqrt2$ (fixed by $\sum_i\lambda_i=3=\dim F_{J^\pm}$ and $\sum_i(\lambda_i-A)^2=3$, both seal-side and hence colour-blind):

%dsp-level:0
\[
(1)\quad\delta_q=\frac{2}{3\cdot3\cdot3}=\frac2{27}=0.074074074074\ldots\qquad[\text{exact rational};\ \text{machine}\ \S D];
\]

%dsp-level:0
\[
(2)\quad3\delta_q=\delta_\ell=\frac29\qquad[\mathrm{FT.112}:\ \delta_\ell=\tfrac2{9}\ \text{from the undressed arena}\ 3\times3];
\]

%dsp-level:0
\[
(3)\quad\textbf{the spectrum:}\qquad r^{(q)}_i=1+\sqrt2\cos\Bigl(\frac2{27}+\frac{2\pi i}3\Bigr),\qquad i=0,1,2,
\]

%dsp-level:0
\[
\text{with}\ c=\cos\tfrac2{27},\ s=\sin\tfrac2{27}:\qquad r^{(q)}_0=1+\sqrt2\,c=2.410335463511094\ldots,
\]

%dsp-level:0
\[
r^{(q)}_1=1-\frac{\sqrt2\,c+\sqrt6\,s}{2}=0.204193367808662\ldots,
\]

%dsp-level:0
\[
\qquad r^{(q)}_2=1-\frac{\sqrt2\,c-\sqrt6\,s}{2}=0.385471168680242\ldots;
\]

%dsp-level:0
\[
\text{all three strictly positive}\ [\text{disjoint enclosures};\ \text{machine}]\ —\ \text{no sign branch, unlike the}\ \delta>\pi/12\ \text{regime};
\]

%dsp-level:0
\[
(4)\quad\textbf{Vieta and the phase-free invariant:}\qquad\sum_ir^{(q)}_i=3,\qquad\sum_{i<j}r^{(q)}_ir^{(q)}_j=\frac32,\qquad\sum_i\bigl(r^{(q)}_i\bigr)^2=6,
\]

%dsp-level:0
\[
2\bigl(r^{(q)}-1\bigr)^3-3\bigl(r^{(q)}-1\bigr)-\sqrt2\cos\frac29=0\quad\text{for each}\ r^{(q)}=r^{(q)}_i
\]

%dsp-level:0
\[
\qquad [\cos3\delta_q=\cos\tfrac29=\cos\delta_\ell:\ \text{clause}\ (2)\ \text{inside the cubic}];
\]

%dsp-level:0
\[
Q\ :=\ \frac{\sum_im_i}{\bigl(\sum_i\sqrt{m_i}\bigr)^2}\ \overset{m_i=\mu (r^{(q)}_i)^2}{=}\ \frac{\sum_i(r^{(q)}_i)^2}{\bigl(\sum_ir^{(q)}_i\bigr)^2}=\frac6{9}=\frac23
\]

%dsp-level:0
\[
\qquad [\text{exact, and independent of}\ \delta\ \text{and of}\ \mu];
\]

%dsp-level:0
\[
(5)\quad\textbf{the mass ratios}\ (m_i=\mu\,(r^{(q)}_i)^2,\ \mu\ \text{undeclared},\ \mathrm{P4}):
\]

%dsp-level:0
\[
\qquad (r^{(q)}_1)^2<(r^{(q)}_2)^2<(r^{(q)}_0)^2\ [\text{enclosure-proven}],
\]

%dsp-level:0
\[
m_1:m_2:m_3
\]

%dsp-level:0
\[
= 1:\frac{(r^{(q)}_2)^2}{(r^{(q)}_1)^2}:\frac{(r^{(q)}_0)^2}{(r^{(q)}_1)^2}
\]

%dsp-level:0
\[
= 1:3.563695074946\ldots:139.3386880284\ldots
\]

%dsp-level:0
\[
\qquad [\text{widths}<2^{-50};\ \text{machine}];
\]

%dsp-level:0
\[
(6)\quad\textbf{anchored forms}\ (\text{scale-free, top-normalized}):
\]

%dsp-level:0
\[
\qquad \frac{m_2}{m_3}=\frac{(r^{(q)}_2)^2}{(r^{(q)}_0)^2}=0.025575776012\ldots,
\]

%dsp-level:0
\[
\qquad \frac{m_1}{m_3}=\frac{(r^{(q)}_1)^2}{(r^{(q)}_0)^2}=0.007176757684\ldots;
\]

%dsp-level:0
\[
(7)\quad\textbf{alternative-hypothesis table}\ (\text{computed, none adopted}):
\]

%dsp-level:0
\[
\qquad \delta=\tfrac2{81}\ (\text{per-baryon arena}\ 3\cdot27):\ 1:1.5130:84.316;
\]

%dsp-level:0
\[
\delta=\tfrac1{36}=\tfrac2{72}\ (\text{arena}\ 9\ \text{dressed by the adjoint multiplicity}\ 8\ [\mathrm{CT3.13}]):\ 1:1.5939:86.746;
\]

%dsp-level:0
\[
\qquad \delta=\tfrac8{81}\ (\mathrm{FT.118}(6)\ \text{numerator over}\ 27):\ 1:5.6438:187.997.
\]

\textbf{Rmk (cross-check, not a derivation of equality).} Clause (2) is the tier's structural coincidence, recorded in the FT.112 discipline: the colour channel's triple angle is the lepton channel's angle, so both channels share one depressed-cubic constant $\sqrt2\cos\tfrac29$ and differ only in which of its three roots they read. Nothing is calibrated to anything; the equality $3\cdot\tfrac2{27}=\tfrac29$ is arithmetic on two independently derived rationals.

\textbf{Req.} FT.118; FT.111(2)–(3), FT.112, FT.113 (\textbf{Tier F.37}: the democratic spectrum form, the normalization, the angle discipline, the $m=\mu r^2$ vocabulary); CT0.12 (Part 1); CT3.13 (\textbf{F}, the alternative table's adjoint multiplicity); L3, L4 (\textbf{B}, the inertial read vocabulary); P4 (\textbf{A}, the free-scale clause); exact rational trigonometric enclosures (Taylor remainder bounds), exact integer square-root enclosures with certificates, monotone rational interval arithmetic, Vieta reads [math; machine: \texttt{checks/\allowbreak{}verify\_\allowbreak{}H\_\allowbreak{}colour\_\allowbreak{}channel.py} §D].

\textbf{Pf.}

%dsp-level:0
\[
(1):\ \delta_q
\]

%dsp-level:0
\[
= \frac{2}{\dim F_{J^\pm}\cdot\operatorname{tr}T(1)\cdot N_{\text{colour}}}
\]

%dsp-level:0
\[
= \frac{2}{3\cdot3\cdot3}
\]

%dsp-level:0
\[
= \frac2{27}\ [\mathrm{FT.118}(1),(2),(4);\ \text{exact Fractions};\ \text{machine}]
\]

%dsp-level:0
\[
(2):\ 3\cdot\tfrac2{27}=\tfrac6{27}=\tfrac29=\delta_\ell\ [\mathrm{FT.112};\ \text{exact}]
\]

%dsp-level:0
\[
(3):\ r^{(q)}_i=A+B\cos(\delta_q+\tfrac{2\pi i}3)\ \text{at}\ A=1,\ B=\sqrt2\ [\mathrm{FT.111}(3)];
\]

%dsp-level:0
\[
\qquad \cos(\tfrac2{27}+\tfrac{2\pi}3)
\]

%dsp-level:0
\[
= -\tfrac c2-\tfrac{\sqrt3}2s,
\]

%dsp-level:0
\[
\quad\cos(\tfrac2{27}+\tfrac{4\pi}3)
\]

%dsp-level:0
\[
= -\tfrac c2+\tfrac{\sqrt3}2s\ [\text{addition formulas}];
\]

%dsp-level:0
\[
\ \sqrt2\cdot\tfrac{\sqrt3}2\qquad = \tfrac{\sqrt6}2
\]

%dsp-level:0
\[
c,s\ \text{enclosed by alternating Taylor partial sums with exact rational remainder bounds (widths}<2^{-90}\text{)};
\]

%dsp-level:0
\[
\qquad \sqrt2,\sqrt6\ \text{enclosed by integer Newton certificates};
\]

%dsp-level:0
\[
\qquad r^{(q)}_1,r^{(q)}_2,r^{(q)}_0>0\ [\text{positive lower endpoints};\ \text{machine}]
\]

%dsp-level:0
\[
(4):\ \sum_i\cos(\delta_q+\tfrac{2\pi i}3)
\]

%dsp-level:0
\[
= 0,
\]

%dsp-level:0
\[
\ \sum_i\cos^2(\delta_q+\tfrac{2\pi i}3)
\]

%dsp-level:0
\[
= \tfrac32\ [\mathrm{FT.111}(3)\ \text{identities}]\implies\sum_ir^{(q)}_i
\]

%dsp-level:0
\[
= 3,
\]

%dsp-level:0
\[
\ \sum_i(r^{(q)}_i)^2
\]

%dsp-level:0
\[
= 3+2\sqrt2\cdot0+2\cdot\tfrac32\qquad = 6
\]

%dsp-level:0
\[
\text{pair sum by Vieta}=\tfrac32;
\]

%dsp-level:0
\[
\qquad \text{depressed cubic}\ (\lambda-1)^3-\tfrac32(\lambda-1)-\tfrac{\sqrt2}2\cos3\delta_q
\]

%dsp-level:0
\[
= 0\ [\mathrm{FT.111}(2)\ \text{at}\ B^2=2];
\]

%dsp-level:0
\[
\quad\cos3\delta_q
\]

%dsp-level:0
\[
= \cos\tfrac29\ [(2)];
\]

%dsp-level:0
\[
\quad\times2:\ 2y^3-3y-\sqrt2\cos\tfrac29
\]

%dsp-level:0
\[
= 0,
\]

%dsp-level:0
\[
\ y\qquad = r^{(q)}-1
\]

%dsp-level:0
\[
0\in2y_i^3-3y_i-\sqrt2\cos\tfrac29\ \text{enclosure for each}\ i\ [\text{machine}];\qquad Q=\tfrac6{9}=\tfrac23\ [\text{exact Fractions}]
\]

%dsp-level:0
\[
(5):\ (r^{(q)}_1)^2<(r^{(q)}_2)^2<(r^{(q)}_0)^2\ [\text{disjoint enclosures}];
\]

%dsp-level:0
\[
\qquad \text{ratios by monotone interval division, widths}<2^{-50}\ [\text{machine}]
\]

%dsp-level:0
\[
(6):\ \text{same divisions against}\ (r^{(q)}_0)^2\ [\text{machine}];
\]

%dsp-level:0
\[
\qquad (7):\ \text{each row by the identical enclosure pipeline at the stated rational}\ \delta\ [\text{machine}].\ \blacksquare
\]

\textbf{Phys.} One integer over three integers: the phase of the colour channel is the seal coefficient divided by a seal factor, a transfer factor and a colour factor, and it needs no new input beyond the rung-3 arena count. Its consequence is compact — triple the angle and the lepton channel's phase appears, so the two hierarchies are the three roots of one cubic read at phases differing by exactly the colour count. The spectrum itself is milder than the lepton-side one: no eigenvalue nearly vanishes, the middle line sits a factor $3.56$ above the inner and the outer a factor $139$ above it, and the anchored forms trade the undeclared scale for two pure numbers. The invariant $Q=\tfrac23$ is not a prediction of this tier — it holds for every phase, being the normalization content $A=1$, $B=\sqrt2$ — so what $\delta_q$ actually fixes is the position on that surface: the two ratios of clause (5). Everything here is dimensionless; $\mu$ is a free internal parameter (P4) and the $m$ vocabulary is the entropy-deficit read of L3(4)/L4, never a particle.

\vspace{2pt}\noindent\rule{\textwidth}{0.4pt}\vspace{2pt}

\end{theorem}

\phantomsection
\section[Tier F.39 — Chirality: The Orientation Doublet of the Bound…]{\texorpdfstring{Tier F.39 — Chirality: The Orientation Doublet of the Bound Chain and Its Exchange-Form Grade}{Tier F.39 - Chirality: The Orientation Doublet of the Bound Chain and Its Exchange-Form Grade}}

\begin{definition}[{\texorpdfstring{Def-H.33 — The chain orientation space and its chirality operator}{Def-H.33 - The chain orientation space and its chirality operator}}]\label{tag:Def-H-33}
\textbf{Def.} For a bound chain $c$ [D1.4: ordered tuple of distinct real edges $c=(e_1,\dots,e_k)$; D1.5: bound $=$ closed self-contained, fermionic content]:

\begin{itemize}
\item \textbf{Orientations.} $c$ and $\omega(c)$ [D2.6: $s(c)\leftrightarrow t(c)$, reversal of the traversal; $\omega^2=\mathrm{id}$]. As ordered tuples, $c\ne\omega(c)$ for every admissible $k\ge2$ [D1.4: distinct edges].
\item \textbf{Orientation space.}
\end{itemize}

%dsp-level:0
\[
\mathcal H_\omega(c)\ :=\ \operatorname{span}_{\mathbb C}\bigl\{|c\rangle,\ |\omega(c)\rangle\bigr\}\ \cong\ \mathbb C^2,
\]

with the Hermitian inner product making $\{|c\rangle,|\omega(c)\rangle\}$ an orthonormal basis.

\begin{itemize}
\item \textbf{Orientation involution.} $\hat\omega:\mathcal H_\omega\to\mathcal H_\omega$, $\hat\omega|c\rangle=|\omega(c)\rangle$, $\hat\omega|\omega(c)\rangle=|c\rangle$: then $\hat\omega^2=\mathrm{id}$, $\hat\omega^\dagger=\hat\omega$ [D2.6: $\omega$ is an involution].
\item \textbf{Chirality operator (the exchange form).}
\end{itemize}

%dsp-level:0
\[
\boxed{\ \hat\chi\ :=\ i\bigl(|c\rangle\langle\omega(c)|-|\omega(c)\rangle\langle c|\bigr)\ }
\]

\begin{itemize}
\item \textbf{Chiral projectors.} $P_L:=\tfrac12(\mathrm{id}+\hat\chi)$, $P_R:=\tfrac12(\mathrm{id}-\hat\chi)$.
\item \textbf{Chiral eigenbasis.}
\end{itemize}

%dsp-level:0
\[
\boxed{\ |L\rangle=\tfrac1{\sqrt2}\bigl(|c\rangle-i\,|\omega(c)\rangle\bigr),\qquad |R\rangle=\tfrac1{\sqrt2}\bigl(|c\rangle+i\,|\omega(c)\rangle\bigr)\ }
\]

— the $+1$ and $-1$ eigenvectors of $\hat\chi$ respectively; the labels are pinned by the boxed $\hat\chi$ together with $P_L:=\tfrac12(\mathrm{id}+\hat\chi)$ [FT.122(3), (6)].

\textbf{Rmk (why $\mathcal H_\omega$ and not the pair register).} $\mathcal H_\omega(c)$ is indexed by the chain $c$, not by any pair $(p,q)$. It is \textit{not} a subspace of $\operatorname{End}(\mathbb C^4_{\mathrm{pair}})$: its basis states are \textit{chains} — configuration-level objects [D1.4, D0.10] — not operators on the pair space. The pincer [FT.93(3)] and the trace/rank obstruction [FT.101(3)] are statements about $\operatorname{End}(\mathbb C^4_{\mathrm{pair}})$ and do not reach $\mathcal H_\omega$ (FT.124 makes the locus separation exact).

\textbf{Rmk (the physical content).} A closed chain's worldline is topologically $S^1$ with two traversal senses. The orientation doublet is the one-dimensional analogue of the linear-polarization basis; the exchange form $\hat\chi=i\cdot$(antisymmetric exchange) is the Hermitian generator of the exchange, and its $\pm1$ rays are the analogue of the two circular polarizations — the two chiralities of a loop mode. Orientation reversal exchanges the two rays [FT.122(4): $\hat\omega|L\rangle=-i|R\rangle$, $\hat\omega|R\rangle=+i|L\rangle$], realizing D2.6's "chirality $\mapsto-$chirality" as an operator identity. \textbf{The spin-structure torsor is a distinct object:} fermion boundary conditions on $S^1$ form $H^1(S^1;\mathbb Z_2)\cong\mathbb Z_2$, which orientation reversal fixes pointwise [FT.122(7); machine] — the torsor $\mathbb Z_2$ and the chiral-ray $\mathbb Z_2$ are different kinds, and only the latter is load-bearing. The phase $i$ in $\hat\chi$ is what turns the antisymmetric exchange form into a Hermitian involution, and it is what makes $\hat\chi$ odd under $\hat\omega$.

\textbf{Rmk (canonicity and the naming residue).} Among Hermitian involutions on $\mathcal H_\omega$ that are $\hat\omega$-odd and diagonal-free in the orientation basis, $\hat\chi$ is unique up to overall sign [FT.122(6); machine: without the diagonal-free condition the family is the circle $M(a,y)=\bigl(\begin{smallmatrix}a&iy\\-iy&-a\end{smallmatrix}\bigr)$, $a^2+y^2=1$]. The sign itself is a naming choice: the map $(\hat\chi,P_L,|L\rangle)\mapsto(-\hat\chi,P_R,|R\rangle)$ leaves every identity of FT.122–FT.123 invariant (the adjoint action is linear: $\hat\omega(-\hat\chi)\hat\omega^{-1}=-(-\hat\chi)$); the boxed definition fixes it once and for all, and given it the labels $L/R$ are forced up to phase. The \textit{physical} naming — which ray is "left" — is fixed by no algebra here: it is naming-convention content reserved to the matter-assignment rule (application-layer; its non-derivability here is proven [FT.125(4)]).

\textbf{Rmk (achiral configurations).} For a physical state invariant under an orientation-reversing symmetry of its configuration, $\langle\hat\chi\rangle=0$ exactly [machine §A; the FT.95(3) pattern] — the grade is sharp per oriented chain, not per reversal-symmetric ensemble; a chiral read requires an oriented chain as boundary datum. This is the same observer-relativity as every chain-level read (Def-H.1 vocabulary); no new declaration.

\vspace{2pt}\noindent\rule{\textwidth}{0.4pt}\vspace{2pt}

\end{definition}

\begin{theorem}[{\texorpdfstring{FT.122 — The orientation doublet is corpus-constructible}{FT.122 - The orientation doublet is corpus-constructible}}]\label{tag:FT-122}
\textbf{Thm.} For every bound chain $c$ [D1.4, D1.5], the orientation space $\mathcal H_\omega(c)$, the involution $\hat\omega$, the chirality operator $\hat\chi$, and the projectors $P_L,P_R$ are constructible from Parts 1–2 data alone, with no new axiom and no declared parameter:

%dsp-level:0
\[
(1)\quad \hat\chi^\dagger=\hat\chi,\qquad \hat\chi^2=\mathrm{id},\qquad \operatorname{spec}\hat\chi=\{+1,-1\},\qquad\operatorname{tr}\hat\chi=0;
\]

%dsp-level:0
\[
(2)\quad P_L^2=P_L,\quad P_R^2=P_R,\quad P_LP_R=0,\quad P_L+P_R=\mathrm{id},\quad \operatorname{tr}P_L=\operatorname{tr}P_R=1;
\]

%dsp-level:0
\[
(3)\quad \hat\chi|L\rangle=+|L\rangle,\quad \hat\chi|R\rangle=-|R\rangle,\quad P_L|L\rangle=|L\rangle,\quad P_L|R\rangle=0,\quad P_R|R\rangle=|R\rangle,\quad P_R|L\rangle=0;
\]

%dsp-level:0
\[
\text{equivalently, without the}\ \sqrt2\ \text{normalizers:}
\]

%dsp-level:0
\[
\quad 2P_L
\]

%dsp-level:0
\[
= |L_u\rangle\langle L_u|,
\]

%dsp-level:0
\[
\ \ 2P_R
\]

%dsp-level:0
\[
= |R_u\rangle\langle R_u|\ \ \text{with}\ \ |L_u\rangle
\]

%dsp-level:0
\[
= |c\rangle-i|\omega(c)\rangle,
\]

%dsp-level:0
\[
\ |R_u\rangle
\]

%dsp-level:0
\[
= |c\rangle+i|\omega(c)\rangle\ \ [\text{exact over}\ \mathbb Q(i)];
\]

%dsp-level:0
\[
(4)\quad\textbf{the covariance identities and the winding read:}
\]

%dsp-level:0
\[
\qquad \hat\omega\hat\chi\hat\omega^{-1}=-\hat\chi,
\]

%dsp-level:0
\[
\qquad \{\hat\chi,\hat\omega\}=0,
\]

%dsp-level:0
\[
\qquad \hat\omega P_L\hat\omega^{-1}=P_R,
\]

%dsp-level:0
\[
\qquad \hat\omega|L\rangle=-i|R\rangle,\quad\hat\omega|R\rangle=+i|L\rangle;
\]

%dsp-level:0
\[
\text{for a materialized chain}\ c:\quad k(c)=\Theta(c)/2\in\mathbb Z\ [\mathrm{FT.22},\ \mathrm{FT.25}(3)];
\]

%dsp-level:0
\[
\qquad \omega:\ k\mapsto-k\ [\mathrm{D2.6};\ \mathrm{EM.24}(1):\ C\ \text{is}\ \omega\text{'s sign-level realization};\ \mathrm{FT.87}(8):\ \operatorname{sign}(k)\ \text{is}\ C\text{-odd}];
\]

%dsp-level:0
\[
(5)\quad\text{the construction is exact over}\ \mathbb Q(i)\ \text{on all matrix elements in the basis}\ \{|c\rangle,|\omega(c)\rangle\}:\quad\hat\chi=i\bigl(\begin{smallmatrix}0&1\\-1&0\end{smallmatrix}\bigr),\quad\hat\omega=\bigl(\begin{smallmatrix}0&1\\1&0\end{smallmatrix}\bigr);
\]

%dsp-level:0
\[
(6)\quad\textbf{canonicity:}
\]

%dsp-level:0
\[
\qquad \bigl\{M=M^\dagger\in\operatorname{End}(\mathcal H_\omega):\ \hat\omega M\hat\omega=-M,\ M^2=\mathrm{id},\ \langle c|M|c\rangle=\langle\omega(c)|M|\omega(c)\rangle=0\bigr\}=\{+\hat\chi,\ -\hat\chi\};
\]

%dsp-level:0
\[
\text{without the diagonal-free condition the}\ \hat\omega\text{-odd Hermitian involutions are the circle}\ M(a,y)=\bigl(\begin{smallmatrix}a&iy\\-iy&-a\end{smallmatrix}\bigr),\ a^2+y^2=1\ \ [\text{exact rational witnesses incl.}\ (a,y)=(\tfrac35,\tfrac45),(\tfrac5{13},\tfrac{12}{13}),(\tfrac8{17},\tfrac{15}{17});\ \text{machine}];
\]

%dsp-level:0
\[
(7)\quad\textbf{the kind flag (torsor distinction):}
\]

%dsp-level:0
\[
\qquad H^1(S^1;\mathbb Z_2)\cong\operatorname{Hom}(\mathbb Z,\mathbb Z_2)=\{h_{\mathrm R}:n\mapsto0,\ \ h_{\mathrm{NS}}:n\mapsto n\bmod2\},
\]

%dsp-level:0
\[
\text{and the orientation-reversal pullback}\ n
\]

%dsp-level:0
\[
\mapsto -n\ \text{fixes both pointwise}\ (h(-n)=h(n)\ \text{in}\ \mathbb Z_2):\ \textbf{the spin-structure torsor is reversal-invariant, NOT exchanged};
\]

%dsp-level:0
\[
\text{the}\ \hat\omega\text{-exchanged objects are the chiral rays of clause}\ (3)\text{: the torsor}\ \mathbb Z_2\ \text{(boundary conditions) and the ray}\ \mathbb Z_2\ \text{(the}\ \hat\chi\text{-grade) are distinct kinds}\ [\mathrm{FT.89}(2)\ \text{discipline}];
\]

%dsp-level:0
\[
\ \text{the R/NS analogy is recorded, never load-bearing}.
\]

\textbf{Req.} D1.4, D1.5, D2.6 (Parts 1–2); FT.22, FT.25 (Tiers F.1–F.11); EM.24 (Part 2); FT.87(8), FT.89(2) (Tiers F.29–F.30); Def-H.1 (\textbf{Tiers F.1–F.11}, the observer-relativity vocabulary of the closing Rmk); Hermitian-involution spectral facts, exact canonicity solve over $\mathbb Q$, $\operatorname{Hom}(\mathbb Z,\mathbb Z_2)$ arithmetic [math; machine: \texttt{checks/\allowbreak{}verify\_\allowbreak{}H\_\allowbreak{}chirality.py} §A].

\textbf{Pf.}

%dsp-level:0
\[
(1):\ \hat\chi^\dagger
\]

%dsp-level:0
\[
= \bigl[i\bigl(|c\rangle\langle\omega(c)|-|\omega(c)\rangle\langle c|\bigr)\bigr]^\dagger
\]

%dsp-level:0
\[
= (-i)\bigl(|\omega(c)\rangle\langle c|-|c\rangle\langle\omega(c)|\bigr)
\]

%dsp-level:0
\[
= i\bigl(|c\rangle\langle\omega(c)|-|\omega(c)\rangle\langle c|\bigr)
\]

%dsp-level:0
\[
= \hat\chi\ \ [(iA)^\dagger=-iA^\dagger;\ (|c\rangle\langle\omega(c)|)^\dagger=|\omega(c)\rangle\langle c|]
\]

%dsp-level:0
\[
E:=|c\rangle\langle\omega(c)|-|\omega(c)\rangle\langle c|:
\]

%dsp-level:0
\[
\quad E^2
\]

%dsp-level:0
\[
= \underbrace{|c\rangle\langle\omega(c)|c\rangle\langle\omega(c)|}_{0}-|c\rangle\underbrace{\langle\omega(c)|\omega(c)\rangle}_{1}\langle c|-\ |\omega(c)\rangle\underbrace{\langle c|c\rangle}_{1}\langle\omega(c)|+\underbrace{|\omega(c)\rangle\langle c|\omega(c)\rangle\langle c|}_{0}
\]

%dsp-level:0
\[
= -\bigl(|c\rangle\langle c|+|\omega(c)\rangle\langle\omega(c)|\bigr)\qquad = -\mathrm{id}
\]

%dsp-level:0
\[
\hat\chi^2=i^2E^2=(-1)(-\mathrm{id})=\mathrm{id};
\]

%dsp-level:0
\[
\qquad \operatorname{tr}\hat\chi=0\ [\text{off-diagonal}];
\]

%dsp-level:0
\[
\qquad \text{charpoly}(\hat\chi)=x^2-\operatorname{tr}(\hat\chi)x+\det\hat\chi=x^2-1\ [\det\hat\chi=0\cdot0-i\cdot(-i)=-1]:\ \operatorname{spec}=\{+1,-1\}
\]

%dsp-level:0
\[
(2):\ P_L^2=\tfrac14(\mathrm{id}+\hat\chi)^2=\tfrac14(\mathrm{id}+2\hat\chi+\hat\chi^2)=\tfrac14(2\,\mathrm{id}+2\hat\chi)=P_L;\quad\text{likewise}\ P_R;
\]

%dsp-level:0
\[
\qquad P_LP_R=\tfrac14(\mathrm{id}-\hat\chi^2)=0;
\]

%dsp-level:0
\[
\qquad P_L+P_R=\mathrm{id};
\]

%dsp-level:0
\[
\qquad \operatorname{tr}P_L=\tfrac12(2+0)=1
\]

%dsp-level:0
\[
(3):\ \hat\chi|c\rangle=i\bigl(|c\rangle\langle\omega(c)|c\rangle-|\omega(c)\rangle\langle c|c\rangle\bigr)=-i|\omega(c)\rangle;\qquad\hat\chi|\omega(c)\rangle=+i|c\rangle
\]

%dsp-level:0
\[
\hat\chi|L_u\rangle=\hat\chi|c\rangle-i\,\hat\chi|\omega(c)\rangle=-i|\omega(c)\rangle-i\cdot i|c\rangle=|c\rangle-i|\omega(c)\rangle=|L_u\rangle;
\]

%dsp-level:0
\[
\qquad \hat\chi|R_u\rangle=-i|\omega(c)\rangle+i\cdot i|c\rangle=-|R_u\rangle
\]

%dsp-level:0
\[
|L_u\rangle\langle L_u|
\]

%dsp-level:0
\[
= (|c\rangle-i|\omega(c)\rangle)(\langle c|+i\langle\omega(c)|)
\]

%dsp-level:0
\[
= \mathrm{id}+i\bigl(|c\rangle\langle\omega(c)|-|\omega(c)\rangle\langle c|\bigr)
\]

%dsp-level:0
\[
= \mathrm{id}+\hat\chi
\]

%dsp-level:0
\[
= 2P_L\ [\text{exact over}\ \mathbb Q(i);\ \langle L_u|L_u\rangle=2];
\]

%dsp-level:0
\[
\qquad \text{likewise}\ |R_u\rangle\langle R_u|=2P_R
\]

%dsp-level:0
\[
\implies P_L|L\rangle=|L\rangle,\ P_L|R\rangle=0,\ P_R|R\rangle=|R\rangle,\ P_R|L\rangle=0\ \checkmark
\]

%dsp-level:0
\[
(4):\ \hat\omega\hat\chi\hat\omega
\]

%dsp-level:0
\[
= i\bigl(\hat\omega|c\rangle\langle\omega(c)|\hat\omega-\hat\omega|\omega(c)\rangle\langle c|\hat\omega\bigr)
\]

%dsp-level:0
\[
= i\bigl(|\omega(c)\rangle\langle c|-|c\rangle\langle\omega(c)|\bigr)
\]

%dsp-level:0
\[
= -\hat\chi\ \ [\hat\omega^{-1}=\hat\omega];
\]

%dsp-level:0
\[
\qquad \{\hat\chi,\hat\omega\}=\hat\chi\hat\omega+\hat\omega\hat\chi=\hat\chi\hat\omega-\hat\chi\hat\omega=0
\]

%dsp-level:0
\[
\hat\omega P_L\hat\omega=\tfrac12(\mathrm{id}+\hat\omega\hat\chi\hat\omega)=\tfrac12(\mathrm{id}-\hat\chi)=P_R;
\]

%dsp-level:0
\[
\qquad \hat\omega|L_u\rangle=|\omega(c)\rangle-i|c\rangle=-i\bigl(|c\rangle+i|\omega(c)\rangle\bigr)=-i|R_u\rangle;
\]

%dsp-level:0
\[
\qquad \hat\omega|R_u\rangle=+i|L_u\rangle
\]

%dsp-level:0
\[
k=\Theta/2\in\mathbb Z\ [\mathrm{FT.22}:\ \Theta\in2\mathbb Z\ \text{on materialized chains};\ \mathrm{FT.25}(3)];
\]

%dsp-level:0
\[
\qquad k
\]

%dsp-level:0
\[
\mapsto -k\ \text{under}\ \omega:\ \mathrm{D2.6}\ \text{(reversal)}+\mathrm{EM.24}(1)\ (C:\sigma_e\mapsto-\sigma_e\ \text{is}\ \omega\text{'s sign-level realization})+\mathrm{FT.87}(8)\ (\operatorname{sign}(k)\ \text{is}\ C\text{-odd})
\]

%dsp-level:0
\[
(5):\ \text{entrywise in the ordered basis}\ (|c\rangle,|\omega(c)\rangle)\ [\text{machine}]
\]

%dsp-level:0
\[
(6):\ M=\bigl(\begin{smallmatrix}a&z\\\bar z&b\end{smallmatrix}\bigr),\ a,b\in\mathbb Q,\ z=x+iy:\quad\hat\omega M\hat\omega=\bigl(\begin{smallmatrix}b&\bar z\\z&a\end{smallmatrix}\bigr)=-M\iff b=-a\wedge x=0\iff M=M(a,y);\quad M(a,y)^2=(a^2+y^2)\,\mathrm{id}
\]

%dsp-level:0
\[
M^2=\mathrm{id}\iff a^2+y^2=1;
\]

%dsp-level:0
\[
\qquad \text{diagonal-free}
\]

%dsp-level:0
\[
\iff a
\]

%dsp-level:0
\[
= 0
\]

%dsp-level:0
\[
\iff y^2
\]

%dsp-level:0
\[
= 1
\]

%dsp-level:0
\[
\iff y
\]

%dsp-level:0
\[
= \pm1
\]

%dsp-level:0
\[
\iff M
\]

%dsp-level:0
\[
= \pm\hat\chi\ [\text{exact rational census of}\ a^2+y^2=1\ \text{at denominators}\ \le20;\ \text{machine}]
\]

%dsp-level:0
\[
(7):\ \operatorname{Hom}(\mathbb Z,\mathbb Z_2)=\{h_0:n\mapsto0,\ h_1:n\mapsto n\bmod2\}\ [\text{a homomorphism is fixed by its value at}\ 1];
\]

%dsp-level:0
\[
\qquad (r^*h)(n)=h(-n):\quad h_1(-n)=(-n)\bmod2=n\bmod2=h_1(n)\ [\mathbb Z_2];\quad h_0(-n)=0=h_0(n)
\]

%dsp-level:0
\[
\implies r^*=\mathrm{id}\ \text{on the torsor: fixed pointwise}\ [\text{machine on}\ |n|\le6,\ \text{general by parity}];
\]

%dsp-level:0
\[
\qquad \text{the exchanged objects are the}\ \hat\chi\text{-rays}\ [(4)];
\]

%dsp-level:0
\[
\ \text{two distinct}\ \mathbb Z_2\text{'s}\ [\mathrm{FT.89}(2):\ \text{kind distinction recorded}].\ \blacksquare
\]

\textbf{Phys.} A bound chain is a closed loop, and a closed loop has two traversal senses. The orientation doublet holds them; the exchange form grades it into two chiral rays — exactly one such grade exists up to sign once the basis is treated symmetrically, so the construction carries no arbitrary choice beyond the naming. Reversing the loop exchanges the rays, and D2.6's vocabulary clause becomes an identity of matrices. The fermion boundary conditions on the loop — the Ramond and Neveu–Schwarz sectors — form a second, reversal-invariant $\mathbb Z_2$ that is \textit{not} this grading; the tier keeps the two kinds apart rather than trading on their shared cardinality. All of it is built from two corpus objects that already exist — the chain and its reversal — with no axiom added and no parameter fitted.

\vspace{2pt}\noindent\rule{\textwidth}{0.4pt}\vspace{2pt}

\end{theorem}

\begin{theorem}[{\texorpdfstring{FT.123 — Weak covariance and $\omega$-covariance on the extended carrier}{FT.123 - Weak covariance and -covariance on the extended carrier}}]\label{tag:FT-123}
\textbf{Thm.} Under Def-H.34's constitutive carrier clause, the matter carrier of a bound chain extends from the seed $\mathcal H_{\mathsf v}\cong\mathbb C^2$ [EW.32] to $\mathcal H_{\mathrm{int}}\otimes\mathcal H_\omega(c)$, with the weak action on the first factor alone, $W_i\mapsto\tfrac12\sigma_i\otimes I_\omega$ [EW.32], and the chiral grade on the second alone:

%dsp-level:0
\[
(1)\quad\textbf{weak covariance:}\qquad[W_i\otimes I_\omega,\ I_{\mathrm{int}}\otimes P_{L/R}]=0\qquad\forall i\in\{1,2,3\};
\]

%dsp-level:0
\[
\textbf{locus sharpening:}\ \ \sum_i\bigl(\tfrac12\sigma_i\otimes I_2\bigr)^2
\]

%dsp-level:0
\[
= \tfrac34I_4\ (\text{type}\ 2\oplus2,\ \text{multiplicity}\ 2),
\]

%dsp-level:0
\[
\quad\text{while on the pair register}\ \sum_iW_i^2
\]

%dsp-level:0
\[
= 2\Pi_{\mathcal R}\ (\text{type}\ 3\oplus1)\ [\text{Cor CT0.7a}]:
\]

%dsp-level:0
\[
\text{the two weak actions are inequivalent}\ (\tfrac34I_4\ne2\Pi_{\mathcal R},\ \text{machine}),
\]

%dsp-level:0
\[
\ \text{so FT.93}(3)\text{'s pincer classifies the pair-register action and places no constraint here};
\]

%dsp-level:0
\[
\mathrm{FT.93}(3d)\text{'s seed irreducibility}\ (\{m:[m,\sigma_i]=0\}=\mathbb CI_2)\ \text{is discharged by multiplicity, not violated:}\ \{M:[M,\sigma_i\otimes I_2]=0\ \forall i\}
\]

%dsp-level:0
\[
= I_{\mathrm{int}}\otimes\operatorname{End}(\mathcal H_\omega),
\]

%dsp-level:0
\[
\ \dim_{\mathbb C}
\]

%dsp-level:0
\[
= 4\ [\text{machine}];
\]

%dsp-level:0
\[
(2)\quad\textbf{the}\ \hat\omega\textbf{-covariance:}
\]

%dsp-level:0
\[
\qquad (I\otimes\hat\omega)\,(I\otimes P_L)\,(I\otimes\hat\omega)^{-1}=I\otimes P_R,
\]

%dsp-level:0
\[
\qquad (I\otimes\hat\omega)\,(I\otimes\hat\chi)\,(I\otimes\hat\omega)^{-1}=-(I\otimes\hat\chi);
\]

%dsp-level:0
\[
(3)\quad\textbf{the weight-cardinality correspondence}\ [\text{conditional on Def-H.38}]:\ \ \text{a materialized chain of winding}\ k\ [\mathrm{FT.22},\ \mathrm{FT.25}(3)]\ \text{carries}\ \mathrm{Sym}^{k}(\mathcal H_{\mathsf v})\ \text{with weights}\ j,
\]

%dsp-level:0
\[
j{-}1,
\]

%dsp-level:0
\[
\dots,
\]

%dsp-level:0
\[
-j,
\]

%dsp-level:0
\[
\ j
\]

%dsp-level:0
\[
= |k|/2\ [\text{Def-H.38};\ \mathrm{FT.60}(3)\ \text{ladder}];
\]

%dsp-level:0
\[
\text{the floor fermion}\ (\text{digon},\ k=1\ [\mathrm{FT.28}])\ \text{has}\ j
\]

%dsp-level:0
\[
= \tfrac12:\ \text{exactly two weights}\ \{+\tfrac12,-\tfrac12\}\ [\text{machine}]\ —\ \text{an unordered-pair cardinality correspondence with the two chiral rays}\ \{+1,-1\};
\]

%dsp-level:0
\[
\text{the label-to-weight attachment is}\ \textbf{not derived}\;
\]

%dsp-level:0
\[
\qquad \text{``chirality}
\]

%dsp-level:0
\[
= \text{helicity on-shell'' is marked external vocabulary — the corpus derives no dispersion relation, and transit-side helicity reads}\ [\mathrm{EM.8}]\ \text{are outside the bound scope}\ [\text{Def-H.38 scope note}];
\]

%dsp-level:0
\[
(4)\quad\textbf{massive modes (structural):}\ \ \text{the}\ \hat\chi\text{-eigenspaces are defined for any mode; helicity is a frame-dependent refinement}\ [\text{math}];
\]

%dsp-level:0
\[
\ \text{on a massive mode chirality and helicity are distinct — the Dirac pattern};
\]

%dsp-level:0
\[
(5)\quad\dim\bigl(\mathcal H_{\mathrm{int}}\otimes\mathcal H_\omega\bigr)
\]

%dsp-level:0
\[
= 2\times2
\]

%dsp-level:0
\[
= 4:\ \text{the Dirac-type carrier dimension of one massive mode (two Weyl halves)};
\]

%dsp-level:0
\[
\qquad \text{a weak doublet of Dirac-type carriers is}\ 2\times4=8.
\]

\textbf{Rmk (interface definition Def-H.34).} The tensor product $\mathcal H_{\mathrm{int}}\otimes\mathcal H_\omega$ is the same species of interface definition as Def-H.41/Def-H.42: the minimal carrier combination consistent with both structures being independently present — the seed with its weak action [EW.32] and the orientation doublet with its exchange-form grade [Def-H.33]. It is exactly the "extended carrier" locus that FT.93(7) identified and item 21(i) enumerated without adopting. The orientation space itself, and every operator identity of FT.122, are declaration-free. Not declared and not constructed: the per-edge frame section $e\mapsto\mathcal F_e$ and its coupling to the EW.32 transport (item 21(i)'s residue), a global Dirac-type operator on the emergent geometry, and any chiral gauge coupling (application-layer content).

\textbf{Req.} Def-H.33; FT.122; Def-H.34 (carrier clause); A8 (Parts 1–2); EW.32, EW.33 (Part 2); D2.6 (Parts 1–2); CT0.7, Cor CT0.7a (Part 1); FT.22, FT.25 (Tiers F.1–F.11); FT.28 (Tiers F.1–F.11); FT.60 (Tiers F.20–F.22); Def-H.38; EM.8 (Part 2); Kronecker-product commutant algebra, Schur's lemma, exact nullspace solves over $\mathbb Q$ [math; machine: \texttt{checks/\allowbreak{}verify\_\allowbreak{}H\_\allowbreak{}chirality.py} §B].

\textbf{Pf.}

%dsp-level:0
\[
(1):\ [W_i\otimes I,\ I\otimes P]
\]

%dsp-level:0
\[
= (W_iI)\otimes(IP)-(IW_i)\otimes(PI)
\]

%dsp-level:0
\[
= 0\ \text{(each factor commutes with an identity)}\ [\text{machine, all}\ i,\ \text{both}\ P]
\]

%dsp-level:0
\[
\sum_i\bigl(\tfrac{\sigma_i}2\otimes I_2\bigr)^2=\bigl(\tfrac14\sum_i\sigma_i^2\bigr)\otimes I_2=\tfrac34I_2\otimes I_2=\tfrac34I_4;
\]

%dsp-level:0
\[
\qquad \sum_iW_i^2=\tfrac32E_0+\tfrac12\Sigma=2\Pi_{\mathcal R}\ [\text{Cor CT0.7a re-derived};\ \text{machine}];
\]

%dsp-level:0
\[
\qquad \tfrac34I_4\ne2\Pi_{\mathcal R}\ [\operatorname{tr}:\ 3\ne4]
\]

%dsp-level:0
\[
\{M:[M,\sigma_i\otimes I_2]=0\}:\ \text{exact}\ 32\text{-unknown real nullspace, dim}_{\mathbb Q}=8\ (=4\ \text{over}\ \mathbb C);
\]

%dsp-level:0
\[
\qquad \{I_2\otimes E_{ab}\}_{a,b}\ \text{lies in it and is}\ \mathbb Q\text{-independent (8 dims)}\implies\text{equality};
\]

%dsp-level:0
\[
\qquad I\otimes P_L\ \text{is a nontrivial member}
\]

%dsp-level:0
\[
\{m:[m,\sigma_i]=0\}=\mathbb CI_2,\quad\{m:\{m,\sigma_i\}=0\}=0\ [\mathrm{FT.93}(3d)\ \text{re-derived};\ \text{exact nullspaces}]
\]

%dsp-level:0
\[
(2):\ \text{clause (4) of FT.122 tensored with}\ I_{\mathrm{int}}\ [\text{machine}];
\]

%dsp-level:0
\[
\qquad (3):\ \tfrac{\sigma_3}2\ \text{on}\ \mathrm{Sym}^1(\mathbb C^2)
\]

%dsp-level:0
\[
= \mathbb C^2:\ \text{eigenvalues}\ \pm\tfrac12\ [\text{machine}];
\]

%dsp-level:0
\[
\ \text{the ladder for general}\ k\ \text{is}\ \mathrm{FT.60}(3)\text{'s (machine-verified there)}
\]

%dsp-level:0
\[
(4):\ \hat\chi\text{-eigenspaces are basis data of}\ \mathcal H_\omega\ [\mathrm{FT.122}(1),(3)];
\]

%dsp-level:0
\[
\ \text{helicity is the frame-dependent sign of the momentum-axis projection}\ [\text{math vocabulary}]
\]

%dsp-level:0
\[
(5):\ 2\cdot2=4;\qquad2\cdot4=8\ [\text{exact}].\ \blacksquare
\]

\textbf{Phys.} The weak force lives on the internal seed; the chiral grade lives on the chain's own orientation bookkeeping; disjoint factors commute. That is the whole mechanism — and it is precisely the ingredient whose absence FT.93(7) diagnosed: a carrier on which the gauge action is trivial \textit{for the grade in question}. The Casimir census keeps the accounting honest: the extended weak action ($\tfrac34I_4$, two doublets) is not the pair-register action ($2\Pi_{\mathcal R}$, a triplet and a singlet), so nothing proven about the latter binds the former, and the seed's irreducibility — the Schur wall behind the pincer — is lifted by multiplicity the moment the orientation factor is adjoined. What is \textit{not} obtained is the Standard Model's chiral coupling: here the weak force is chirality-blind, the chiral halves are weak-blind rather than weak-inactive (the extended weak kernel is $\{0\}$), and the rule that would make one ray gauge-active is application-layer content, $\neg$derived here.

\textbf{Rmk (forward pointer, non-load-bearing).} The chirality-flipping block of the complete static force matrix on the extended carrier — the off-diagonal complement of clause (1)'s chirality blindness — is computed in Tier F.40: at the derived locus $F_{LR}=0$ exactly (FT.128, the no-flip theorem: every derived aspect action is internal-only), and under the declared read promotion Def-H.43 the flipping block is $F_{LR}=q_{\mathrm{core}}I_2$ with flipping set $\{F_{\mathrm{EM}}\ \text{core read}\}$ and chiral-partner splitting $2|q|\in\{0,2\}$ (FT.129–FT.130). The pointer adds no load-bearing content to this tier.

\vspace{2pt}\noindent\rule{\textwidth}{0.4pt}\vspace{2pt}

\end{theorem}

\medskip\noindent\textit{\texorpdfstring{FT.124 — The prior obstructions do not reach the orientation space}{FT.124 - The prior obstructions do not reach the orientation space}.}\ Structural-constraints record: see Appendix B, Section B.3.4 of the Addendum.\par\medskip

\medskip\noindent\textit{\texorpdfstring{FT.125 — The deficit ledger: deficit 1 closed at the projector level; deficit 1$'$ opened}{FT.125 - The deficit ledger: deficit 1 closed at the projector level; deficit 1 opened}.}\ Structural-constraints record: see Appendix B, Section B.4.1 of the Addendum.\par\medskip

\begin{theorem}[{\texorpdfstring{FT.126 — Weak-covariant chiral projectors on the chain carrier}{FT.126 - Weak-covariant chiral projectors on the chain carrier}}]\label{tag:FT-126}
\textbf{Thm.} Combining FT.123's tensor-factor separation with the branch-grading locus of FT.101: the orientation-space chirality lifts to the chain-carried Dirac-type carrier $\mathcal H_{\mathrm{int}}\otimes\mathcal H_\omega(c)$:

%dsp-level:0
\[
(1)\quad P_{L/R}^{\mathrm{car}}\ :=\ I_{\mathrm{int}}\otimes P_{L/R}:
\]

%dsp-level:0
\[
\qquad (P^{\mathrm{car}}_{L/R})^2=P^{\mathrm{car}}_{L/R},
\]

%dsp-level:0
\[
\qquad P^{\mathrm{car}}_LP^{\mathrm{car}}_R=0,
\]

%dsp-level:0
\[
\qquad P^{\mathrm{car}}_L+P^{\mathrm{car}}_R=\mathrm{id},
\]

%dsp-level:0
\[
\qquad \operatorname{tr}P^{\mathrm{car}}_{L/R}=2;
\]

%dsp-level:0
\[
(2)\quad[W_i\otimes I_\omega,\ P^{\mathrm{car}}_{L/R}]=0\ \ \forall i\qquad[\text{the weak isospin acts chirality-blind}];
\]

%dsp-level:0
\[
(3)\quad\text{each Weyl half}\ \operatorname{im}P^{\mathrm{car}}_{L/R}\ \text{carries exactly one chiral ray of the orientation factor and the full internal seed action (weak-blind, not weak-inactive:}\ \bigcap_i\ker(W_i\otimes I_\omega)=\{0\}\ \text{[machine])};
\]

%dsp-level:0
\[
(4)\quad\textbf{independence of the two}\ \mathbb Z_2\textbf{'s:}\ \ \text{the A4 branch grading of FT.101 is label-level registration data}\ [\mathrm{FT.101}(3):\ \text{no Hilbert-space implementation on the pair register}];
\]

%dsp-level:0
\[
\ \text{the chiral grading is an implemented involution on the orientation factor. Any internal involution}\ G\otimes I_\omega\ \text{commutes with}\ I_{\mathrm{int}}\otimes\hat\omega\ \text{and their product is a third, distinct involution:}\ \mathbb Z_2\times\mathbb Z_2\ [\text{machine on}\ G\in\{\sigma_3,\sigma_1\};\ \mathrm{FT.88}\text{'s}\ \langle\omega,\Theta_{\mathfrak L}\rangle=\mathbb Z_2\times\mathbb Z_2\ \text{is the frozen precedent for the product structure}].
\]

\textbf{Req.} FT.122, FT.123, FT.124 (this tier); FT.88, FT.101 (\textbf{Tiers F.29–F.30}, \textbf{Tiers F.31–F.34}); EW.32 (Part 2); Def-H.34 (carrier clause); Kronecker-product identities [math; machine: \texttt{checks/\allowbreak{}verify\_\allowbreak{}H\_\allowbreak{}chirality.py} §B].

\textbf{Pf.}

%dsp-level:0
\[
(1):\ \mathrm{FT.122}(2)\ \text{tensored with}\ I_{\mathrm{int}}:\ (I\otimes P)^2=I\otimes P^2=I\otimes P;\quad\operatorname{tr}(I_2\otimes P_{L/R})=2\cdot1=2\ [\text{machine}]
\]

%dsp-level:0
\[
(2):\ \mathrm{FT.123}(1)\ \text{verbatim};
\]

%dsp-level:0
\[
\qquad (3):\ \operatorname{im}(I\otimes P_L)
\]

%dsp-level:0
\[
= \mathcal H_{\mathrm{int}}\otimes\mathbb C|L\rangle;
\]

%dsp-level:0
\[
\quad\bigcap_i\ker(\sigma_i\otimes I)
\]

%dsp-level:0
\[
= \bigl(\bigcap_i\ker\sigma_i\bigr)\otimes\mathcal H_\omega
\]

%dsp-level:0
\[
= \{0\}\ [\sigma_i\ \text{invertible};\ \text{exact}\ 24\times8\ \text{rank}\ 8;\ \text{machine}]
\]

%dsp-level:0
\[
(4):\ [G\otimes I,\ I\otimes\hat\omega]=0\ [\text{machine}];
\]

%dsp-level:0
\[
\qquad (G\otimes I)(I\otimes\hat\omega)=G\otimes\hat\omega\ne G\otimes I,\ I\otimes\hat\omega\ [\text{entrywise};\ \text{machine}];
\]

%dsp-level:0
\[
\qquad (G\otimes\hat\omega)^2=G^2\otimes\hat\omega^2=\mathrm{id}\ \blacksquare
\]

\textbf{Phys.} The weak force and the chiral grade do not interfere: the branch grading of FT.101 and the chiral grading of this tier are two independent $\mathbb Z_2$'s on disjoint data, and the lifted projectors give each Weyl half of the chain carrier its own orientation ray. Which branch, which ray, and which species go together is a \textit{third} rule — the application-layer assignment rule of FT.125(4).

\vspace{2pt}\noindent\rule{\textwidth}{0.4pt}\vspace{2pt}

\end{theorem}

\phantomsection
\section[Tier F.40 — The Static Force Matrix on the Extended…]{\texorpdfstring{Tier F.40 — The Static Force Matrix on the Extended Carrier: The Aspect Census, the No-Flip Theorem, and the Charge–Parity Splitting}{Tier F.40 - The Static Force Matrix on the Extended Carrier: The Aspect Census, the No-Flip Theorem, and the Charge-Parity Splitting}}

\begin{definition}[{\texorpdfstring{Def-H.34 — The extended carrier basis, the chiral transform, and the partner involution}{Def-H.34 - The extended carrier basis, the chiral transform, and the partner involution}}]\label{tag:Def-H-34}
\textbf{Def.} For a bound chain $c$ [D1.4: ordered tuple of distinct real edges; D1.5: bound $=$ closed self-contained, fermionic content]:

\begin{itemize}
\item \textbf{Extended carrier (constitutive carrier clause).} $\mathcal H_{\mathrm{ext}}:=\mathcal H_{\mathrm{int}}\otimes\mathcal H_\omega(c)\cong\mathbb C^2\otimes\mathbb C^2\cong\mathbb C^4$, with $\mathcal H_{\mathrm{int}}=\mathcal H_{\mathsf v}$ the seed carrier, orthonormal basis $\{|\mathsf v\rangle,|\mathsf v^\perp\rangle\}$ [EW.32; the complement ray is a basis convention — all statements below are basis-covariant], and $\mathcal H_\omega$, $\hat\omega$, $\hat\chi$, $P_{L/R}$, $|L\rangle$, $|R\rangle$ as in Def-H.33.
\item \textbf{Ordered orientation basis} (internal-major Kronecker order; composite index $=2\cdot\mathrm{int}+\omega$):
\end{itemize}

%dsp-level:0
\[
|1\rangle:=|\mathsf v\rangle\otimes|c\rangle,
\]

%dsp-level:0
\[
\qquad |2\rangle:=|\mathsf v\rangle\otimes|\omega(c)\rangle,
\]

%dsp-level:0
\[
\qquad |3\rangle:=|\mathsf v^\perp\rangle\otimes|c\rangle,
\]

%dsp-level:0
\[
\qquad |4\rangle:=|\mathsf v^\perp\rangle\otimes|\omega(c)\rangle.
\]

\begin{itemize}
\item \textbf{Chiral carrier basis.} $\{|\mathsf v\rangle\otimes|L\rangle,\ |\mathsf v\rangle\otimes|R\rangle,\ |\mathsf v^\perp\rangle\otimes|L\rangle,\ |\mathsf v^\perp\rangle\otimes|R\rangle\}$, with $|L\rangle=\tfrac1{\sqrt2}(|c\rangle-i|\omega(c)\rangle)$, $|R\rangle=\tfrac1{\sqrt2}(|c\rangle+i|\omega(c)\rangle)$ [Def-H.33].
\item \textbf{The $S$-transform.} $S:=\bigl(\begin{smallmatrix}1&1\\-i&i\end{smallmatrix}\bigr)$ (columns $=$ unnormalized chiral kets): $S^\dagger S=2I_2$; $U:=I_2\otimes S$; the conjugation $M\mapsto\tfrac12U^\dagger MU$ is exact over $\mathbb Q(i)$ and rotates the $\omega$-factor Pauli images cyclically:
\end{itemize}

%dsp-level:0
\[
\hat\omega\ (\sigma_1)\ \mapsto\ -\sigma_2,\qquad\hat\chi\ (-\sigma_2)\ \mapsto\ \sigma_3,\qquad\hat\kappa\ (\sigma_3)\ \mapsto\ \sigma_1.
\]

\begin{itemize}
\item \textbf{Partner involution.} $\hat\kappa:=|c\rangle\langle c|-|\omega(c)\rangle\langle\omega(c)|$ — the unique diagonal $\hat\omega$-odd Hermitian involution up to sign [FT.122(6)'s circle at $(a,y)=(\pm1,0)$; machine]. Quaternionic triple with Def-H.33's operators:
\end{itemize}

%dsp-level:0
\[
\hat\omega^2=\hat\chi^2=\hat\kappa^2=\mathrm{id};
\]

%dsp-level:0
\[
\qquad \{\hat\omega,\hat\chi\}=\{\hat\chi,\hat\kappa\}=\{\hat\kappa,\hat\omega\}=0;
\]

%dsp-level:0
\[
\qquad \hat\omega\hat\chi=-i\hat\kappa,\quad\hat\chi\hat\kappa=-i\hat\omega,\quad\hat\kappa\hat\omega=-i\hat\chi.
\]

\begin{itemize}
\item \textbf{Force-profile and read data.} $\mathbf c=(c_0,c_1,c_2,c_3)\in\mathbb Q^4$: the weak-transport profile of the chain over the seed generator set $\{W_1,W_2,W_3,Y\}$ [EW.33's $c_A$; EW.32's images] — configuration data, formal indeterminates in the Def-H.32(3) discipline, never fitted [P1, P4];
\end{itemize}

%dsp-level:0
\[
h(\mathbf c)=\tfrac12\sum_{a=0}^{3}c_a\,\sigma_a\qquad(\sigma_0:=I_2);\qquad|\mathbf c|^2:=c_1^2+c_2^2+c_3^2;
\]

$q:=q_{\mathrm{core}}(c)\in\{-1,0,+1\}$: the ternary core-charge read of the oriented chain [D3.4, EM.9, FT.40] — registration data, $\omega$-odd [FT.40].

\textbf{Rmk (the three bases of the triple).} Each basis of $\mathcal H_\omega$ diagonalizes one member of the triple: the orientation basis $\{|c\rangle,|\omega(c)\rangle\}$ diagonalizes $\hat\kappa$; the chiral basis $\{|L\rangle,|R\rangle\}$ diagonalizes $\hat\chi$; the parity basis $\{(|c\rangle+|\omega(c)\rangle)/\sqrt2,\ (|c\rangle-|\omega(c)\rangle)/\sqrt2\}$ diagonalizes $\hat\omega$ [machine]. Every ray of one member is an equal-weight superposition in the eigenbasis of the other two — in particular $|c\rangle=(|L\rangle+|R\rangle)/\sqrt2$: the orientation rays, which FT.129(5) makes stationary, carry no definite chirality. No new structure beyond Def-H.33: $\hat\kappa$ completes its operator triple.

\vspace{2pt}\noindent\rule{\textwidth}{0.4pt}\vspace{2pt}

\end{definition}

\begin{theorem}[{\texorpdfstring{FT.127 — The aspect census on the extended carrier}{FT.127 - The aspect census on the extended carrier}}]\label{tag:FT-127}
\textbf{Thm.} For the complete force operator of A8 — aspect inventory: the six pair-sector aspects [CT0.12], the rung masks and stack blocks [CT0.13, CT0.14], the two dark aspects [CT4.6, D4.1, D4.3], the Janus move classes [D1.2, FT.108(1)] — the derived-action census on $\mathcal H_{\mathrm{ext}}$ is:

%dsp-level:0
\[
(1)\quad F_W\ (\mathcal E_{+-}=\operatorname{span}\{W_i\},\ 3\text{-dim}\ [\mathrm{CT0.12}]):\ \textbf{acts}.
\]

%dsp-level:0
\[
\qquad W_i^{(\mathrm{ext})}=\tfrac12\sigma_i\otimes I_\omega
\]

%dsp-level:0
\[
\qquad [\mathrm{EW.32}\ \text{seed images};\ \text{Def-H.34}]\ \ [\text{exact}\ 4\times4\ \text{over}\ \mathbb Q(i)];
\]

%dsp-level:0
\[
(2)
\]

%dsp-level:0
\[
\quad F_{\mathrm{EM}}\ (\text{stack},\ 1\text{-dim}\ [\mathrm{CT0.14}]):\ \textbf{acts on the internal factor only},
\]

%dsp-level:0
\[
\ \text{via the neutral mixing}\ A
\]

%dsp-level:0
\[
= \sin\theta_W\,
\]

%dsp-level:0
\[
W_3+\cos\theta_W\,
\]

%dsp-level:0
\[
Y\ [\mathrm{D3.10},\ \mathrm{EW.24}]:
\]

%dsp-level:0
\[
A^{(\mathrm{ext})}=\tfrac12\bigl(\sin\theta_W\,\sigma_3+\cos\theta_W\,\sigma_0\bigr)\otimes I_\omega
\]

%dsp-level:0
\[
\qquad [\mathrm{EW.32}:\ Y\mapsto\tfrac12\sigma_0]\ —\ \text{inside clause}\ (1)\text{'s span};\ \text{chirality-blind}\ (\otimes I_\omega);
\]

%dsp-level:0
\[
\text{the ternary end-register read}\ q_{\mathrm{core}}\ [\mathrm{D3.4},\ \mathrm{EM.9},\ \mathrm{FT.40}]\ \text{is}\ \omega\text{-odd configuration data: an orientation-side action requires the read promotion of FT.129 (Def-H.43)};
\]

%dsp-level:0
\[
\text{a bridge between the seed neutral content and the ternary read is not derived (the ternary-read situation is $\neg$derived here)};
\]

%dsp-level:0
\[
(3)
\]

%dsp-level:0
\[
\quad F_G\ (\operatorname{span}\{\hat E^{ij}\}\ \text{traceless},\ 5\text{-dim}\ [\mathrm{CT0.12}]):\ \textbf{does not act on}\ \mathcal H_{\mathrm{ext}}\ —\ \text{pair-register block}\ [\mathrm{CT1.2}];
\]

%dsp-level:0
\[
\ \text{chain-level content}
\]

%dsp-level:0
\[
= \text{geometric defect data}\ (\Theta_R\text{-type scalars under the Def-H.37(2)}\ \text{transfer}\ [\mathrm{FT.33},\ \mathrm{FT.54}(4)])\ —\ \text{scalar class};
\]

%dsp-level:0
\[
(4)
\]

%dsp-level:0
\[
\quad F_H\ (\operatorname{span}\{\Sigma\},\ H=h\Sigma\ [\mathrm{EW.13}]):\ \textbf{does not act as a carrier matrix}\ —\ \Sigma\in\operatorname{End}(\mathbb C^4_{\mathrm{pair}})\ [\mathrm{CT1.2}];
\]

%dsp-level:0
\[
\ \text{Higgs content on chains}
\]

%dsp-level:0
\[
= \text{scalar background}\ (h_0\ \text{P4-class};\ \text{entering only through the EW.20/EW.33 defect-residue packaging})\ —\ \text{scalar class};
\]

%dsp-level:0
\[
(5)
\]

%dsp-level:0
\[
\quad F_{D^+}\ (\operatorname{span}\{E_0\}=\Lambda\,\mathbb I\ [\mathrm{D4.3},\ \mathrm{CT4.6}]):\ \textbf{scalar class}\ —\ \text{central identity direction, configuration-independent};
\]

%dsp-level:0
\[
\ \text{if admitted on a carrier it is}\ \lambda I_4,
\]

%dsp-level:0
\[
\ \text{re-absorbed by}\ h
\]

%dsp-level:0
\[
\mapsto h+\lambda I_2\ [\mathrm{FT.130}(7)];
\]

%dsp-level:0
\[
\ \text{geometric side}\ \Lambda g_{\mu\nu}\ [\mathrm{CT1.23}],
\]

%dsp-level:0
\[
\ w
\]

%dsp-level:0
\[
= -1\ [\mathrm{CT1.25}];
\]

%dsp-level:0
\[
(6)
\]

%dsp-level:0
\[
\quad F_{J^+},
\]

%dsp-level:0
\[
F_{J^-}\ (\operatorname{span}\{J_i\},\ \operatorname{span}\{K_i\}\ [\mathrm{CT0.12},\ \mathrm{CT0.7}(4)]):\ \textbf{do not act on}\ \mathcal H_{\mathrm{ext}}\ —\ \text{pair-register matrices}\ [\mathrm{CT1.2}]\ \text{and graph-level Janus move generators}\ [\mathrm{D1.2},\ \mathrm{FT.108}(1)];
\]

%dsp-level:0
\[
\ \text{no seed or orientation action derived};
\]

%dsp-level:0
\[
\ \text{their bilinear channel content is the seal-triplet reads of Tiers F.37–F.38}\ [\mathrm{FT.109},\ \mathrm{FT.116}];
\]

%dsp-level:0
\[
(7)
\]

%dsp-level:0
\[
\quad F_{\mathrm{strong}}\ (8\text{-dim stack}\ [\mathrm{CT3.8},\ \mathrm{CT3.13},\ \mathrm{CT0.14}]):\ \textbf{does not act on}\ \mathcal H_{\mathrm{ext}}\ —\ \text{acts on the colour-slot carrier}\ \mathbb C^3\ \text{at trivalent vertices}\ [\mathrm{CT3.2},\ \mathrm{D3.1}];
\]

%dsp-level:0
\[
\ \text{the seed is colour-neutral by theorem}\ [\mathrm{EW.32}:\ T_a\mathsf v=0];
\]

%dsp-level:0
\[
\ \text{colour sits in an orthogonal stack block}\ [\mathrm{CT0.12}\ \mathrm{Rmk};\ \mathrm{FT.93}(4)];
\]

%dsp-level:0
\[
\ \mathcal H_{\mathrm{ext}}\ \text{has no slot factor};
\]

%dsp-level:0
\[
\ \text{sealed ends are blind in any case}\ [\mathrm{FT.117}];
\]

%dsp-level:0
\[
(8)
\]

%dsp-level:0
\[
\quad F_{D^-}\ \bigl(\bigoplus_{n\ge4}F_n\ [\mathrm{D4.1},\ \mathrm{CT4.1},\ \mathrm{CT4.6}]\bigr):\ \textbf{does not act}\ —\ \text{ledger sink/retention verbs}\ [\mathrm{CT4.2}\text{–}\mathrm{CT4.5};\ \text{the retention census is}\ u\text{-fugacity bookkeeping}\ [\mathrm{FT.108}(3)]];
\]

%dsp-level:0
\[
\ \text{no carrier representation};
\]

%dsp-level:0
\[
(9)\quad\textbf{census conclusion:}\ \ \text{the derived static force matrix is exactly the EW seed content},
\]

%dsp-level:0
\[
\mathbf F_{\mathrm{ext}}(\mathbf c)=h(\mathbf c)\otimes I_\omega,\qquad h(\mathbf c)=\tfrac12\sum_{a=0}^{3}c_a\,\sigma_a,
\]

%dsp-level:0
\[
\text{and the four images}\ \{\tfrac12\sigma_a\otimes I_2\}_{a=0}^{3}\ \text{are}\ \mathbb C\text{-independent with span}\ \operatorname{End}(\mathbb C^2)\otimes I_\omega\ [\text{machine}]:\ \text{every derived aspect action is chirality-blind}.
\]

\textbf{Req.} CT0.7, CT0.12, CT0.13, CT0.14, CT1.2, CT1.3 (Part 1); CT1.23, CT1.25 (Part 2); A8, D1.2, D3.1, D3.4, D3.10, D4.1, D4.3, P1, P4 (Parts 1–2); EW.8, EW.13, EW.20, EW.24, EW.32, EW.33 (Part 2); EM.9 (Part 2); CT3.2, CT3.8, CT3.13 (Part 2); CT4.1, CT4.2, CT4.3, CT4.5, CT4.6 (Part 2); FT.20 (Tiers F.1–F.11); FT.33 (Tiers F.12–F.13); FT.54 (Tiers F.17–F.19); FT.93 (Tiers F.31–F.34); FT.108 (Tier F.36); FT.109 (Tier F.37); FT.116, FT.117 (Tier F.38); Def-H.33, FT.123 (Tier F.39); Def-H.37(2), Def-H.32, Def-H.34; exact $\mathbb Q(i)$ independence rank (real-form flattening with $i$-doubles) [math; machine: \texttt{checks/\allowbreak{}verify\_\allowbreak{}H\_\allowbreak{}force\_\allowbreak{}matrix.py} §B].

\textbf{Pf.}

%dsp-level:0
\[
(1):\ \mathrm{EW.32}:\ W_i\mapsto\tfrac12\sigma_i\ \text{on}\ \mathcal H_{\mathsf v};
\]

%dsp-level:0
\[
\qquad \text{Def-H.34: the extended action is}\ W_i\mapsto\tfrac12\sigma_i\otimes I_\omega;
\]

%dsp-level:0
\[
\qquad \text{the three}\ 4\times4\ \text{matrices written entrywise over}\ \mathbb Q(i)\ [\text{machine}]
\]

%dsp-level:0
\[
(2):\ A=\sin\theta_WW_3+\cos\theta_WY\ [\mathrm{EW.24};\ \mathrm{D3.10}];
\]

%dsp-level:0
\[
\qquad A^{(\mathrm{ext})}=\sin\theta_W\,W_3^{(\mathrm{ext})}+\cos\theta_W\,Y^{(\mathrm{ext})}\ [\text{linearity}];
\]

%dsp-level:0
\[
\qquad Y\mapsto\tfrac12\sigma_0\ [\mathrm{EW.32}];
\]

%dsp-level:0
\[
\qquad \text{unit normalization pin:}\ \mathfrak T\bigl((\tfrac12(\sin\theta_W\sigma_3+\cos\theta_W\sigma_0))^2\bigr)=\tfrac12\ [\sin^2+\cos^2=1;\ \mathfrak T\sigma_3=0]
\]

%dsp-level:0
\[
(9):\ \text{flattening}\ \{\tfrac12\sigma_a\otimes I_2\}\ \text{with}\ i\text{-doubles: real rank}\ 8=4\ \text{over}\ \mathbb C\ [\text{machine}];
\]

%dsp-level:0
\[
\qquad \operatorname{span}\{\sigma_0,\sigma_1,\sigma_2,\sigma_3\}=\operatorname{End}(\mathbb C^2)\ [\mathrm{CT1.1}\text{'s Pauli basis};\ \text{machine}]
\]

%dsp-level:0
\[
(3)\text{–}(8):\ \text{each row's locus is the cited frozen tag verbatim: pair-register membership}\ [\mathrm{CT1.2}\ \text{basis table}],
\]

%dsp-level:0
\[
\ \text{slot carriers}\ [\mathrm{CT3.2}],
\]

%dsp-level:0
\[
\ \text{ledger sinks}\ [\mathrm{D4.1},\ \mathrm{CT4.1}],
\]

%dsp-level:0
\[
\ \text{scalar/central directions}\ [\mathrm{D4.3},\ \mathrm{EW.13},\ \mathrm{EW.20}];
\]

%dsp-level:0
\[
\qquad \neg\exists\ \text{derived seed or orientation action for any of the eight}\ [\text{inventory scan: the only carrier bridge in Parts 1–2 is EW.32's}].\ \blacksquare
\]

\textbf{Phys.} Ask the nine aspects of the complete force where they act on the extended carrier, and eight answer elsewhere — on the pair register, in the colour slots, in the ledger, or as scalar background. The census is the theorem: the derived force content of a chain carrier is exactly the electroweak seed image, one simple tensor with the orientation identity, hence chirality-blind by construction. No aspect of A8 reaches the orientation factor; the only canonical candidate that does is the chain's own $\omega$-odd read data, taken up in FT.129.

\vspace{2pt}\noindent\rule{\textwidth}{0.4pt}\vspace{2pt}

\end{theorem}

\begin{theorem}[{\texorpdfstring{FT.128 — The no-flip theorem at the derived locus}{FT.128 - The no-flip theorem at the derived locus}}]\label{tag:FT-128}
\textbf{Thm.} In the chiral basis of Def-H.34:

%dsp-level:0
\[
(1)\quad\textbf{form invariance:}
\]

%dsp-level:0
\[
\qquad U^{-1}\bigl(h(\mathbf c)\otimes I_\omega\bigr)U=h(\mathbf c)\otimes I_2
\]

%dsp-level:0
\[
\qquad [\text{the conjugation touches only the}\ \omega\text{-factor, where the derived matrix is the identity}];
\]

%dsp-level:0
\[
(2)\quad\textbf{the block decomposition:}\qquad F_{LL}=F_{RR}=h(\mathbf c),\qquad F_{LR}=F_{RL}=0\qquad\text{exactly, for every profile}\ \mathbf c;
\]

%dsp-level:0
\[
(3)\quad\textbf{the weak flip contribution vanishes by computation:}
\]

%dsp-level:0
\[
\qquad P_L^{\mathrm{car}}\,W_i^{(\mathrm{ext})}\,P_R^{\mathrm{car}}=0\ \ \forall i
\]

%dsp-level:0
\[
\qquad [\text{entrywise on the explicit}\ 4\times4\ \text{matrices};\ \mathrm{FT.123}(1)\ \text{re-verified, not invoked}];
\]

%dsp-level:0
\[
(4)\quad\bigl[\mathbf F_{\mathrm{ext}}(\mathbf c),\ I_2\otimes\hat\chi\bigr]=0\ \ \forall\mathbf c:
\]

%dsp-level:0
\[
\qquad \text{chirality is an exact quantum number of the derived static force matrix};
\]

%dsp-level:0
\[
(5)\quad\textbf{the derived spectrum:}\qquad\operatorname{charpoly}\mathbf F_{\mathrm{ext}}(\mathbf c)=\Bigl(x^2-c_0x+\tfrac{c_0^2-|\mathbf c|^2}{4}\Bigr)^2,
\]

%dsp-level:0
\[
\operatorname{spec}
\]

%dsp-level:0
\[
= \Bigl\{\tfrac12\bigl(c_0+|\mathbf c|\bigr),\ \tfrac12\bigl(c_0-|\mathbf c|\bigr)\Bigr\},
\]

%dsp-level:0
\[
\ \text{each with multiplicity}\ 2\ \text{across the chiral rays}\ —\ \textbf{spectral parity doubling};
\]

%dsp-level:0
\[
\qquad |\mathbf c|=\sqrt{|\mathbf c|^2}\in\text{a quadratic extension of}\ \mathbb Q(\mathbf c);
\]

%dsp-level:0
\[
(6)\quad\textbf{the static shadow, derived locus:}
\]

%dsp-level:0
\[
\qquad \langle L|\mathbf F_{\mathrm{ext}}|R\rangle
\]

%dsp-level:0
\[
= 0\ —\ \text{the force aspects alone produce no parity-flipping seed on the extended carrier of a bound chain};
\]

%dsp-level:0
\[
\qquad \text{the flipping set at the derived locus}=\varnothing.
\]

\textbf{Req.} FT.127; Def-H.34; FT.122(1)–(3), FT.123(1) (Tier F.39); CT1.2 (Part 1); Zariski-dense rational sampling (entries linear in $\mathbf c$: vanishing on a dense set implies identically zero), Faddeev–LeVerrier over $\mathbb Q(i)$ [math; machine: \texttt{checks/\allowbreak{}verify\_\allowbreak{}H\_\allowbreak{}force\_\allowbreak{}matrix.py} §B].

\textbf{Pf.}

%dsp-level:0
\[
(1):\ U=I_2\otimes S:\quad U^{-1}(h\otimes I_2)U=h\otimes(S^{-1}I_2S)=h\otimes I_2\ [\text{machine on}\ 108\ \text{exact rational profiles}]
\]

%dsp-level:0
\[
(2):\ \text{block extraction in the chiral carrier basis}\ \{|\mathsf v L\rangle,|\mathsf v R\rangle,|\mathsf v^\perp L\rangle,|\mathsf v^\perp R\rangle\}:\ F_{LR}
\]

%dsp-level:0
\[
= P_L^{\mathrm{car}}\mathbf F_{\mathrm{ext}}P_R^{\mathrm{car}}
\]

%dsp-level:0
\[
= 0\ \text{entrywise at every sample profile};
\]

%dsp-level:0
\[
\qquad \text{entries of}\ F_{LR}\ \text{are linear in}\ \mathbf c:\ \text{vanishing on the Zariski-dense sample set}\implies F_{LR}\equiv0\ [\text{machine}]
\]

%dsp-level:0
\[
(3):\ P_L^{\mathrm{car}}=I_2\otimes P_L,\ P_R^{\mathrm{car}}=I_2\otimes P_R\ [\mathrm{FT.126}(1)]:
\]

%dsp-level:0
\[
\quad (I\otimes P_L)(\tfrac{\sigma_i}2\otimes I)(I\otimes P_R)=\tfrac{\sigma_i}2\otimes P_LP_R=0\ [\mathrm{FT.122}(2):\ P_LP_R=0];
\]

%dsp-level:0
\[
\qquad \text{machine: entrywise zero on all three}\ 4\times4\ \text{products}
\]

%dsp-level:0
\[
(4):\ [h\otimes I_2,\ I_2\otimes\hat\chi]
\]

%dsp-level:0
\[
= [h,I_2]\otimes[I_2,\hat\chi]\text{-type disjointness}:\ (hI_2)\otimes(I_2\hat\chi)-(I_2h)\otimes(\hat\chi I_2)
\]

%dsp-level:0
\[
= 0\ [\text{machine}]
\]

%dsp-level:0
\[
(5):\ \operatorname{charpoly}h(\mathbf c)=x^2-c_0x+\tfrac{c_0^2-|\mathbf c|^2}4\ [\det h=\tfrac{c_0^2-|\mathbf c|^2}4;\ \mathfrak T h=c_0];
\]

%dsp-level:0
\[
\qquad \operatorname{charpoly}(h\otimes I_2)
\]

%dsp-level:0
\[
= (\operatorname{charpoly}h)^2\ [\text{Kronecker spectrum};\ \text{machine: Faddeev–LeVerrier coefficient identity at all}\ 108\ \text{profiles}]
\]

%dsp-level:0
\[
(6):\ \langle L|\mathbf F_{\mathrm{ext}}|R\rangle=(I_2\text{-block of}\ F_{LR})=0\ [(2)].\ \blacksquare
\]

\textbf{Phys.} The derived force matrix is blind to the very distinction the carrier was built to carry: it preserves each chiral ray, doubles every level across the two rays, and commutes with the chirality grade. That is not a failure of the census — it is its content. Every force the corpus derives onto this carrier arrives through the seed, and the seed action is internal: the orientation factor sees only the identity. A parity flip, if the corpus is to have one, must come from the orientation factor's own read structure, and the no-flip theorem narrows the search to a single question: does any $\omega$-odd datum of the bound chain act there?

\vspace{2pt}\noindent\rule{\textwidth}{0.4pt}\vspace{2pt}

\end{theorem}

\begin{definition}[{\texorpdfstring{Def-H.43 — The read promotion on the orientation doublet}{Def-H.43 - The read promotion on the orientation doublet}}]\label{tag:Def-H-43}
\textbf{Def.} An $\omega$-odd A8-aspect read carried by a bound chain acts on the orientation doublet $\mathcal H_\omega$ [Def-H.33] as the multiplication operator of its eigenvalues on the orientation basis rays:

%dsp-level:0
\[
\text{read}\ q\ \longmapsto\ \hat q\ :=\ q\cdot\hat\kappa
\]

%dsp-level:0
\[
\qquad (\omega\text{-odd reads}),
\]

%dsp-level:0
\[
\qquad \hat\kappa\
\]

%dsp-level:0
\[
= \ |c\rangle\langle c|-|\omega(c)\rangle\langle\omega(c)|\ \ [\mathrm{Def\text{-}H.34}:\ \text{the unique diagonal}\ \hat\omega\text{-odd Hermitian involution up to sign},\ \mathrm{FT.122}(6)],
\]

EM core read: $\ \hat Q_{\mathrm{EM}}=q_{\mathrm{core}}(c)\,\hat\kappa$, with $q_{\mathrm{core}}\in\{-1,0,+1\}$ ternary registration data [D3.4, EM.9, FT.40] — $\neg$fitted. $\omega$-even reads shift (no flip); the parity split is derived [FT.129(3) — immediate below], $\neg$part of this definition. The promotion is the sole source of a chirality-flipping block [FT.129(4)–(5)]. \textbf{Derived under the promotion (cited).} Flipping block $F_{LR}=q_{\mathrm{core}}I_2$; flipping set $\{F_{\mathrm{EM}}\ \text{core read}\}$ — minimal and unique; chiral-partner splitting $2|q|\in\{0,2\}$ [FT.129(5), FT.130(5)]; at the derived locus (no promotion): the no-flip theorem holds unconditionally [FT.128]. \textbf{Bars (standing).} $2|q|$ is $\neg$a frequency — no carrier-time parameter exists; emergent time is a prerequisite [FT.131(3)]; $\hat\kappa$-/$\hat\chi$-eigenrays are $\neg$species and $\neg$SM chiralities — the assignment rule is application-layer content; $\omega$-even reads are not promoted to $\hat\kappa$-multiples [FT.129(3) is derived]; a transit-chain flipping census is out of scope [FT.131(1)]; external calibration barred [P1, P5]. \textbf{Tension remark (structural).} The chiral rays superpose $\omega$-opposed core charges $(q,-q)$ [FT.40] while channel-level coarse-graining is sign-blind and fuses only equal-$q$ states [FT.15's hypothesis; FT.90(5)'s audit] — the carrier-level doublet and channel fusion are distinct structures; their relation is bridge content [Def-H.39's slots], open, and no theorem of this part identifies them [FT.131(5)].

\end{definition}

\begin{theorem}[{\texorpdfstring{FT.129 — The canonical flip candidate and the read promotion}{FT.129 - The canonical flip candidate and the read promotion}}]\label{tag:FT-129}
\textbf{Thm.}

%dsp-level:0
\[
(1)\quad\textbf{canonicity:}
\]

%dsp-level:0
\[
\qquad \bigl\{M=M^\dagger:\ M\ \text{diagonal in the orientation basis},\ \hat\omega M\hat\omega=-M,\ M^2=\mathrm{id}\bigr\}=\{+\hat\kappa,\ -\hat\kappa\}
\]

%dsp-level:0
\[
\qquad [\mathrm{FT.122}(6)\text{'s circle at}\ (a,y)=(\pm1,0);\ \text{exact rational census};\ \text{machine}];
\]

%dsp-level:0
\[
(2)\quad\textbf{chiral images:}
\]

%dsp-level:0
\[
\qquad \hat\kappa|L\rangle=|R\rangle,\quad\hat\kappa|R\rangle=|L\rangle\ (\text{ray swap, no phase});
\]

%dsp-level:0
\[
\qquad \text{in the chiral basis:}\ \hat\kappa\mapsto\sigma_1\text{-type},
\]

%dsp-level:0
\[
\ \hat\chi\mapsto\sigma_3\text{-type},
\]

%dsp-level:0
\[
\ \hat\omega
\]

%dsp-level:0
\[
\mapsto -\sigma_2\text{-type}\ [\text{the}\ S\text{-rotation of Def-H.34};\ \text{machine}];
\]

%dsp-level:0
\[
(3)\quad\textbf{the parity split of the reads on the doublet:}
\]

%dsp-level:0
\[
\qquad \omega\text{-odd}:\ q_{\mathrm{core}}\ [\mathrm{FT.40}],\ k\ \text{and}\ \operatorname{sign}(k)\ [\mathrm{D2.6};\ \mathrm{EM.24}(1);\ \mathrm{FT.87}(8)];
\]

%dsp-level:0
\[
\qquad \omega\text{-even}:\ |k|,\ \Theta^2,\ \text{the sign-blind reads}\ [\mathrm{FT.90}(5)];
\]

%dsp-level:0
\[
\text{an even read, promoted, multiplies}\ I_2\ (\text{chirality-blind});
\]

%dsp-level:0
\[
\ \text{an odd read, promoted, multiplies}\ \hat\kappa\ (\text{chirality-flipping}):\ \text{the flipping structure of any scalar read class is decided by its}\ \omega\text{-parity};
\]

%dsp-level:0
\[
(4)\quad\textbf{Def-H.43 (read promotion, Def-H.43):}\ \ \text{an}\ \omega\text{-odd A8-aspect read carried by a bound chain acts on}\ \mathcal H_\omega\ \text{as the multiplication operator of its eigenvalues on the basis rays. For the EM core read:}
\]

%dsp-level:0
\[
\hat Q_{\mathrm{EM}}:=q_{\mathrm{core}}(c)\,\hat\kappa
\]

%dsp-level:0
\[
\qquad [\text{the unique minimal promotion consistent with FT.40's eigenvalues}\ (q,-q);\ q\in\{-1,0,+1\}\ \text{registration data},\ \neg\text{fitted}];
\]

%dsp-level:0
\[
\text{no such action is derived in Parts 1–2 (reads are labels}\ [\mathrm{Def\text{-}H.15}(4)]\text{), hence declared.}\ \textbf{Tension logged:}\ \text{the chiral rays superpose}\ \omega\text{-opposed core charges}\ (q,-q);
\]

%dsp-level:0
\[
\ \text{channel-level coarse-graining fuses only equal-}q\ \text{states}\ [\mathrm{FT.15}\ \text{hypothesis};\ \mathrm{FT.90}(5)\ \text{audit}]\ —\ \text{the relation between the carrier-level doublet and channel fusion is item 7/item 13 bridge content, open};
\]

%dsp-level:0
\[
(5)\quad\textbf{the promoted force matrix and its flipping block:}\qquad\text{under Def-H.43,}
\]

%dsp-level:0
\[
\mathbf F_{\mathrm{ext}}(\mathbf c,q)=h(\mathbf c)\otimes I_2+q\,I_2\otimes\hat\kappa,\qquad\bigl[h(\mathbf c)\otimes I_2,\ q\,I_2\otimes\hat\kappa\bigr]=0\ [\text{machine}];
\]

%dsp-level:0
\[
F_{LL}=F_{RR}=h(\mathbf c),\qquad F_{LR}=q\,I_2=F_{RL}:\qquad\text{nonzero}\iff q\ne0;
\]

%dsp-level:0
\[
\textbf{the flipping set}
\]

%dsp-level:0
\[
= \{F_{\mathrm{EM}}\ \text{core read}\}:\ \text{minimal and unique — removing the}\ \hat\kappa\ \text{term kills}\ F_{LR}\ \text{for every profile}\ [\mathrm{FT.128}(2)];
\]

%dsp-level:0
\[
\ \text{removing every internal term leaves}\ F_{LR}
\]

%dsp-level:0
\[
= qI_2\ [\text{machine}];
\]

%dsp-level:0
\[
\text{the off-diagonal matrix elements encoding the flipping structure are the core-charge read itself:}
\]

%dsp-level:0
\[
\qquad \langle L|\mathbf F_{\mathrm{ext}}|R\rangle=q\,I_2.
\]

\textbf{Req.} Def-H.34; FT.128; FT.122(6)–(7), FT.123, FT.126 (Tier F.39); Def-H.33 (Tier F.39); FT.15 (Tiers F.1–F.11); FT.40 (Tiers F.12–F.13); Def-H.15 (Tiers F.1–F.11); FT.87(8), FT.89(2), FT.90(5) (Tiers F.29–F.30); EM.9, EM.24 (Part 2); D2.6, D3.4 (Parts 1–2); Def-H.43; exact canonicity census, Kronecker block algebra [math; machine: \texttt{checks/\allowbreak{}verify\_\allowbreak{}H\_\allowbreak{}force\_\allowbreak{}matrix.py} §A, §C].

\textbf{Pf.}

%dsp-level:0
\[
(1):\ M=\operatorname{diag}(a,b):\ \hat\omega M\hat\omega=\operatorname{diag}(b,a)=-M\iff b=-a;
\]

%dsp-level:0
\[
\qquad M^2=\mathrm{id}\iff a^2=1\iff a=\pm1\iff M=\pm\hat\kappa\ [\text{exact rational census};\ \text{machine}]
\]

%dsp-level:0
\[
(2):\ \hat\kappa|L_u\rangle=\hat\kappa(|c\rangle-i|\omega(c)\rangle)=|c\rangle+i|\omega(c)\rangle=|R_u\rangle\ [\text{unnormalized, exact}];
\]

%dsp-level:0
\[
\qquad \text{likewise}\ \hat\kappa|R_u\rangle=|L_u\rangle;
\]

%dsp-level:0
\[
\qquad \tfrac12S^\dagger\sigma_3S=\sigma_1\ [\mathrm{Def\text{-}H.34}\ \text{rotation};\ \text{machine}]
\]

%dsp-level:0
\[
(3):\ q_{\mathrm{core}}(\omega(c))=-q_{\mathrm{core}}(c)\ [\mathrm{FT.40}];
\]

%dsp-level:0
\[
\qquad k\mapsto-k\ [\mathrm{D2.6};\ \mathrm{EM.24}(1);\ \mathrm{FT.87}(8)];
\]

%dsp-level:0
\[
\qquad |k|,\ \Theta^2\ \text{invariant};\ \text{sign-blindness of the listed reads}\ [\mathrm{FT.90}(5)]
\]

%dsp-level:0
\[
\text{promotion linearity: an}\ \omega\text{-even read}\ r\ (\text{values}\ (r,r))\ \mapsto\ rI_2;
\]

%dsp-level:0
\[
\ \text{an}\ \omega\text{-odd read}\ r\ (\text{values}\ (r,-r))\
\]

%dsp-level:0
\[
\mapsto r\hat\kappa\ [\text{diagonal multiplication operators on the orientation basis};\ \text{uniqueness of the diagonal form: clause}\ (1)]
\]

%dsp-level:0
\[
(4):\ \hat Q_{\mathrm{EM}}|c\rangle=q|c\rangle,\quad\hat Q_{\mathrm{EM}}|\omega(c)\rangle=-q|\omega(c)\rangle\ [\mathrm{FT.40}\ \text{eigenvalues}];
\]

%dsp-level:0
\[
\qquad \text{declaration status and tension: recorded at Def-H.43}
\]

%dsp-level:0
\[
(5):\ \bigl[h\otimes I_2,\ qI_2\otimes\hat\kappa\bigr]=q\bigl(hI_2\otimes I_2\hat\kappa-I_2h\otimes\hat\kappa I_2\bigr)=0\ [\text{machine}];
\]

%dsp-level:0
\[
\qquad \tfrac12U^\dagger(I_2\otimes\hat\kappa)U=I_2\otimes\sigma_1\ [\mathrm{Def\text{-}H.34}]
\]

%dsp-level:0
\[
\implies\ \text{chiral blocks of}\ h\otimes I_2+q\,I_2\otimes\sigma_1:\ F_{LL}=F_{RR}=h,
\]

%dsp-level:0
\[
\quad F_{LR}=F_{RL}=qI_2\ [\text{machine on}\ 48\ \text{profiles}\times q\in\{-1,0,1\}];
\]

%dsp-level:0
\[
\qquad \text{minimality rows: machine}\ \blacksquare
\]

\textbf{Phys.} The orientation doublet carries exactly two canonical involutions — the exchange form, which grades chirality, and its partner, which grades the traversal sense — and the partner is the flip operator: it swaps the chiral rays with no phase. The corpus's own reads split cleanly across this pair: everything $\omega$-even can only shift levels; everything $\omega$-odd can only flip chirality. Promoting the ternary core-charge read — one interface definition, its tension with equal-$q$ channel fusion logged rather than hidden — the complete static force matrix acquires its only off-diagonal content, and the flipping amplitude is the charge read itself: a charged core flips, a neutral core does not.

\vspace{2pt}\noindent\rule{\textwidth}{0.4pt}\vspace{2pt}

\end{theorem}

\begin{theorem}[{\texorpdfstring{FT.130 — Tensor structure, Schmidt rank, and the static spectrum}{FT.130 - Tensor structure, Schmidt rank, and the static spectrum}}]\label{tag:FT-130}
\textbf{Thm.}

%dsp-level:0
\[
(1)\quad\textbf{derived locus:}
\]

%dsp-level:0
\[
\qquad \mathbf F_{\mathrm{ext}}(\mathbf c)=h(\mathbf c)\otimes I_\omega\ \text{is a single simple tensor for every}\ \mathbf c:
\]

%dsp-level:0
\[
\qquad \text{operator-Schmidt rank}=1\ (\text{rank}\ 0\iff\mathbf c=0)\ [\text{exact Pauli-bilinear flattening};\ \text{machine}];
\]

%dsp-level:0
\[
(2)\quad\textbf{promoted locus:}
\]

%dsp-level:0
\[
\qquad \mathbf F_{\mathrm{ext}}(\mathbf c,q)
\]

%dsp-level:0
\[
= h\otimes I_2+q\,
\]

%dsp-level:0
\[
I_2\otimes\hat\kappa:\ \text{the two terms are Hilbert–Schmidt orthogonal}\ (\mathfrak T\hat\kappa=0);
\]

%dsp-level:0
\[
\qquad \text{Schmidt rank}=2\iff q\ne0\ \wedge\ (c_1,c_2,c_3)\ne0;
\]

%dsp-level:0
\[
\text{rank}\ 1\ \text{when}\ h\propto I_2\ (\text{the terms share the internal factor and recombine:}\ \mathbf F_{\mathrm{ext}}=I_2\otimes(\lambda I_2+q\hat\kappa))\ \text{or exactly one term survives};
\]

%dsp-level:0
\[
\qquad \text{rank}\ 0\iff\mathbf c=0=q;
\]

%dsp-level:0
\[
\textbf{no entangling term at any locus:}\ \ \mathbf F_{\mathrm{ext}}\ \text{is always a sum of}\
\]

%dsp-level:0
\[
\le 2\ \text{simple tensors with one-sided factors};
\]

%dsp-level:0
\[
(3)\quad\textbf{the flipping criterion (iff):}
\]

%dsp-level:0
\[
\qquad F_{LR}\ne0\iff\text{some Schmidt term has a chiral-off-diagonal}\ \omega\text{-factor}\iff\text{the}\ \hat\kappa\ \text{term is present with}\ q\ne0
\]

%dsp-level:0
\[
\qquad [I_2\ \text{is chiral-diagonal};\ \hat\kappa\ \text{is chiral-off-diagonal}\ [\mathrm{FT.129}(2)]];
\]

%dsp-level:0
\[
(4)\quad\textbf{derived spectrum:}
\]

%dsp-level:0
\[
\qquad \operatorname{spec}=\bigl\{\tfrac12(c_0\pm|\mathbf c|)\bigr\},\ \text{each}\times2;
\]

%dsp-level:0
\[
\qquad \text{field}\ \mathbb Q(\mathbf c)\bigl[\sqrt{|\mathbf c|^2}\bigr]:\ \text{quadratic extension};
\]

%dsp-level:0
\[
\ \text{charpoly exact over}\ \mathbb Q(\mathbf c)[x]\ [\mathrm{FT.128}(5)];
\]

%dsp-level:0
\[
(5)\quad\textbf{promoted spectrum:}
\]

%dsp-level:0
\[
\qquad \text{the terms commute}\ [\mathrm{FT.129}(5)]:
\]

%dsp-level:0
\[
\qquad \operatorname{spec}=\bigl\{\tfrac12\bigl(c_0+s|\mathbf c|\bigr)+tq:\ s,t\in\{\pm1\}\bigr\},
\]

%dsp-level:0
\[
\operatorname{charpoly}=\bigl((y-u)^2-q^2\bigr)\bigl((y+u)^2-q^2\bigr),\qquad y=x-\tfrac{c_0}2,\quad u^2=\tfrac{|\mathbf c|^2}4\qquad[\text{machine on all profiles}\times\text{ternary}\ q];
\]

%dsp-level:0
\[
\textbf{the chiral-partner splitting}=2|q|\in\{0,2\}
\]

%dsp-level:0
\[
\qquad [\text{the ternary read}\ [\mathrm{EM.9}]:\ \text{charged cores split the parity doublet by exactly two read units; neutral cores keep the doubling}];
\]

%dsp-level:0
\[
\text{simultaneous eigenbasis}
\]

%dsp-level:0
\[
= \text{internal Bloch rays}\otimes\text{orientation rays}\ \{|c\rangle,|\omega(c)\rangle\}\ (\hat\kappa\text{-eigenrays},\ \kappa\text{-eigenvalues}\ \pm1);
\]

%dsp-level:0
\[
\qquad |c\rangle=(|L\rangle+|R\rangle)/\sqrt2\ [\text{machine}]:
\]

%dsp-level:0
\[
\text{the stationary rays carry no definite chirality — chirality is the precessing quantity, and}\ 2|q|\ \text{is the static shadow of the parity-flipping frequency (the dynamical read is deferred, FT.131}(3));
\]

%dsp-level:0
\[
(6)\quad\textbf{collision census (exact):}
\]

%dsp-level:0
\[
\qquad q=0:\ \text{parity doubling};
\]

%dsp-level:0
\[
\qquad |\mathbf c|=0:\ \operatorname{spec}=\{\tfrac{c_0}2\pm q\},\ \text{each}\times2;
\]

%dsp-level:0
\[
\qquad \tfrac{|\mathbf c|}2=|q|:\ \text{the central eigenvalue}\ \tfrac{c_0}2\ \text{is double};
\]

%dsp-level:0
\[
\qquad \text{no further degeneracies at ternary}\ q;
\]

%dsp-level:0
\[
(7)\quad\textbf{scalar-shift clause:}
\]

%dsp-level:0
\[
\qquad \mathbf F_{\mathrm{ext}}\mapsto\mathbf F_{\mathrm{ext}}+\lambda I_4\iff h\mapsto h+\lambda I_2:
\]

%dsp-level:0
\[
\qquad F_{LR},
\]

%dsp-level:0
\[
\ \text{Schmidt ranks, commutants and the splitting}\ 2|q|\ \text{invariant};
\]

%dsp-level:0
\[
\ \text{the spectrum translates by}\ \lambda
\]

%dsp-level:0
\[
[\text{any scalar-class contribution}\ —\ F_{D^+}\ \text{if admitted},\ h_0\text{-type backgrounds},\ m^2\text{-type reads as identity multiples}\ —\ \text{is invisible to the flipping structure}].
\]

\textbf{Req.} FT.128, FT.129; Def-H.34; CT1.1, CT1.2 (\textbf{B}, the Pauli basis); EM.9 (Part 2); Hilbert–Schmidt and operator-Schmidt rank facts, exact flattening ranks over $\mathbb Q$, Faddeev–LeVerrier, Vieta reads [math; machine: \texttt{checks/\allowbreak{}verify\_\allowbreak{}H\_\allowbreak{}force\_\allowbreak{}matrix.py} §D].

\textbf{Pf.}

%dsp-level:0
\[
(1):\ \mathbf F_{\mathrm{ext}}=\sum_{ab}t_{ab}\,\sigma_a\otimes\sigma_b,\quad t_{ab}=\tfrac14\mathfrak T\bigl(\mathbf F_{\mathrm{ext}}(\sigma_a\otimes\sigma_b)\bigr):
\]

%dsp-level:0
\[
\qquad t_{a0}=c_a/2,\ \text{all other}\ t_{ab}=0\ [\text{Pauli orthogonality}\ \mathfrak T(\sigma_a\sigma_b)=2\delta_{ab}]
\]

%dsp-level:0
\[
\operatorname{rk}T=1\iff\mathbf c\ne0\ [\text{exact rational flattening ranks at all sample profiles};\ \text{machine}]
\]

%dsp-level:0
\[
(2):\ t_{a0}=c_a/2,\quad t_{03}=q\ [\hat\kappa=\sigma_3\ \text{in the orientation basis}]:
\]

%dsp-level:0
\[
\qquad \operatorname{rk}T=2\iff q\ne0\wedge(c_1,c_2,c_3)\ne0;
\]

%dsp-level:0
\[
\quad h
\]

%dsp-level:0
\[
= \lambda I_2\ (\text{i.e.}\ c_1=c_2=c_3=0):\ \mathbf F_{\mathrm{ext}}
\]

%dsp-level:0
\[
= I_2\otimes(\lambda I_2+q\sigma_3):\ \operatorname{rk}T
\]

%dsp-level:0
\[
= 1\ [\text{machine, degenerate corners included}]
\]

%dsp-level:0
\[
\mathfrak T\bigl((h\otimes I_2)^\dagger(qI_2\otimes\hat\kappa)\bigr)
\]

%dsp-level:0
\[
= q\,
\]

%dsp-level:0
\[
\mathfrak T(h)\,
\]

%dsp-level:0
\[
\mathfrak T(\hat\kappa)
\]

%dsp-level:0
\[
= 0\ [\mathfrak T\hat\kappa=0;\ \text{machine witness}]:\ \text{the two Schmidt terms are orthogonal}
\]

%dsp-level:0
\[
(3):\ F_{LR}=qI_2\ [\mathrm{FT.129}(5)];
\]

%dsp-level:0
\[
\qquad \text{chiral images:}\ I_2\mapsto I_2\ (\text{diagonal}),\ \hat\kappa\mapsto\sigma_1\ (\text{off-diagonal})\ [\mathrm{Def\text{-}H.34}]\ \implies\text{both directions exact}
\]

%dsp-level:0
\[
(4),
\]

%dsp-level:0
\[
(5):\ \text{Faddeev–LeVerrier coefficient identity against}\ \bigl(x^2-c_0x+\tfrac{c_0^2-|\mathbf c|^2}4\bigr)^2\ \text{and}\ \bigl((y-u)^2-q^2\bigr)\bigl((y+u)^2-q^2\bigr)\ \text{at all profiles}\times\text{ternary}\ q\ [\text{machine}]
\]

%dsp-level:0
\[
\text{splitting: fixed}\ s:\ \bigl(\tfrac12(c_0+s|\mathbf c|)+q\bigr)-\bigl(\tfrac12(c_0+s|\mathbf c|)-q\bigr)=2q;
\]

%dsp-level:0
\[
\qquad q\in\{-1,0,+1\}\implies2|q|\in\{0,2\}\ [\mathrm{EM.9};\ \text{exact}]
\]

%dsp-level:0
\[
(6):\ \text{collision loci from the closed form:}\ q=0;
\]

%dsp-level:0
\[
\quad |\mathbf c|=0;
\]

%dsp-level:0
\[
\quad u^2=q^2\ [\text{exhaustive over the closed form};\ \text{machine census at rational}\ u\in\{\tfrac12,1,\tfrac32\}\times\text{ternary}\ q]
\]

%dsp-level:0
\[
(7):\ \lambda I_4=(\lambda I_2)\otimes I_2:\ h+\lambda I_2\ \text{re-absorbs the shift}\ [\text{machine}];
\]

%dsp-level:0
\[
\qquad P_L^{\mathrm{car}}(\lambda I_4)P_R^{\mathrm{car}}=\lambda P_LP_R=0\ [\mathrm{FT.122}(2)].\ \blacksquare
\]

\textbf{Phys.} The static force matrix never entangles the internal and orientation books beyond a single product term, and the whole flipping phenomenon reduces to one ternary number: the core charge. Its level splitting $2|q|$ is the only frequency-shaped datum the static tier can produce — the eigenvalues of the promoted matrix encode, exactly, the splitting that a dynamical read would turn into the parity-flipping frequency. Scalar reads of every species shift the ladder and change nothing else: the flip structure is parity-graded, and parity is decided by $\omega$-oddness alone.

\vspace{2pt}\noindent\rule{\textwidth}{0.4pt}\vspace{2pt}

\end{theorem}

\begin{theorem}[{\texorpdfstring{FT.131 — Scope: bound chains, no isolated charge, the deferrals}{FT.131 - Scope: bound chains, no isolated charge, the deferrals}}]\label{tag:FT-131}
\textbf{Thm.}

%dsp-level:0
\[
(1)\quad\textbf{bound-chain scope:}
\]

%dsp-level:0
\[
\qquad \mathbf F_{\mathrm{ext}}\ \text{is computed for bound chains}\ [\mathrm{D1.5}];
\]

%dsp-level:0
\[
\ \text{the orientation doublet requires}\ k
\]

%dsp-level:0
\[
\ge 2\ \text{distinct edges as an ordered tuple}\ [\mathrm{D1.4}];
\]

%dsp-level:0
\[
\ \text{transit/trivial chains carry no bound orientation doublet}\ [\mathrm{Def\text{-}H.33}\ \text{scope};\ \text{Def-H.38 scope note}]:\ \text{their flipping census is a separate computation, not performed here};
\]

%dsp-level:0
\[
(2)\quad\textbf{no isolated charge:}
\]

%dsp-level:0
\[
\qquad \text{a charge read requires two end registers}\ [\mathrm{D3.4},\ \mathrm{EM.9}];
\]

%dsp-level:0
\[
\ \text{a bound chain carries}\ q_{\mathrm{core}}\ [\mathrm{FT.40}]\ \text{but no transit charge};
\]

%dsp-level:0
\[
\ \ F_{LR}
\]

%dsp-level:0
\[
= q_{\mathrm{core}}I_2\ \text{is the flipping structure}\ \textbf{of the bound chain},
\]

%dsp-level:0
\[
\ \text{not of an isolated carrier};
\]

%dsp-level:0
\[
\ \text{isolated colour is excluded in any case}\ [\mathrm{CT3.22}:\ \neg\text{pole},\ \text{norm}\ 0];
\]

%dsp-level:0
\[
(3)\quad\textbf{the deferral and the CPT adjudication:}
\]

%dsp-level:0
\[
\qquad \text{a frequency read requires a time parameter coupling}\ \mathbf F_{\mathrm{ext}}\ \text{to an evolution. The CPT-equivalent record was consulted:}\ \mathrm{EM.24}(5)\text{'s exact phase identity}
\]

%dsp-level:0
\[
w(\mathrm{CPT}\,C)=(-1)^{|E(C)|}e^{i\pi\Theta(C)}w(C),\qquad|w|^2\ \text{unconditionally invariant},
\]

%dsp-level:0
\[
\text{relates reversal-paired configuration amplitudes — the configuration-level flipping locus}\ (\text{the pairing exists, its phase is exact})\ \text{but supplies no carrier-time parameter};
\]

%dsp-level:0
\[
\text{the only derived clock is the ledger tick}\ [\mathrm{CT0.3}:\ (\Delta n_{\mathcal R},\delta h,\delta\tau)=\pm(1,1,1)\ \text{per Janus move}],
\]

%dsp-level:0
\[
\ \text{whose rate has no field coupling}\ [\mathrm{FT.86}(8)];
\]

%dsp-level:0
\[
\ \text{proper-time vocabulary exists}\ [\mathrm{Def\text{-}H.19}\ \text{stationary jumps};\ \text{Morse–Bott refinement}\ \mathrm{Def\text{-}H.24}]\ \text{with no materialization coupling}\ [\text{Def-H.19}];
\]

%dsp-level:0
\[
\text{hence: emergent time is a prerequisite; nothing in the CPT record discharges it; the frequency,}\ e^{-i\mathbf F t},
\]

%dsp-level:0
\[
\ \text{and the time-resolved chiral coupling are deferred — the splitting}\ 2|q|\ [\mathrm{FT.130}(5)]\ \text{is the complete static deliverable};
\]

%dsp-level:0
\[
(4)\quad\textbf{bars:}
\]

%dsp-level:0
\[
\qquad \text{no particle identification}\ (\text{the matrices are for an arbitrary bound chain with arbitrary profile and registration data});
\]

%dsp-level:0
\[
\ \text{no external data}\ —\ \text{the marker [EXT] appears in this clause alone; no external value enters any clause of this document};
\]

%dsp-level:0
\[
\ \text{no Yukawa/mass mechanism}\ (\text{mass-type scalars enter only through FT.130}(7)\text{'s identity-shift clause, never as}\ L\text{–}R\ \text{couplings});
\]

%dsp-level:0
\[
\ \text{no binding energy};
\]

%dsp-level:0
\[
\ \text{no fixed topology or privileged skeleton}\ (\mathbf c,q,k,\Theta\ \text{are chain and registration data — the matrices are skeleton-independent});
\]

%dsp-level:0
\[
\ \text{single-aspect intermediates are marked scope restrictions};
\]

%dsp-level:0
\[
(5)\quad\textbf{the standing superselection tension (scope restatement):}
\]

%dsp-level:0
\[
\qquad \text{the chiral rays superpose}\ \omega\text{-opposed core charges}\ (q,-q)\ [\mathrm{FT.40}];
\]

%dsp-level:0
\[
\ \text{channel reads are sign-blind and fuse only equal-}q\ \text{states}\ [\mathrm{FT.15};\ \mathrm{FT.90}(5)];
\]

%dsp-level:0
\[
\ \text{Def-H.43's promotion is carrier-level; its relation to channel-level fusion is open bridge content}\ \ —\ \text{no theorem here resolves it}.
\]

\textbf{Req.} D1.4, D1.5, D3.4 (Parts 1–2); EM.9, EM.24 (Part 2); CT0.3 (Part 1); CT3.22 (Part 2); FT.15 (Tiers F.1–F.11); FT.40 (Tiers F.12–F.13); Def-H.19 (Tiers F.12–F.13); Def-H.24 (Tiers F.20–F.22); FT.86, FT.90 (Tiers F.29–F.30); Def-H.19, Def-H.38, Def-H.43; exact bookkeeping; machine: \texttt{checks/\allowbreak{}verify\_\allowbreak{}H\_\allowbreak{}force\_\allowbreak{}matrix.py} §E (citation pins).

\textbf{Pf.}

%dsp-level:0
\[
(1):\ \mathrm{D1.5}:\ \text{bound}=\text{closed self-contained};
\]

%dsp-level:0
\[
\qquad \mathrm{D1.4}:\ c=(e_1,\dots,e_k),\ e_i\ \text{distinct};\ k\ge2\ [\mathrm{A1}:\ \neg\text{self-loops}];
\]

%dsp-level:0
\[
\qquad \mathrm{Def\text{-}H.33}\ \text{constructs}\ \mathcal H_\omega(c)\ \text{for bound}\ c
\]

%dsp-level:0
\[
(2):\ \mathrm{D3.4}:\ q_{\mathrm{EM}}=\tfrac12(\sigma_s-\sigma_t)\ \text{on end registers};
\]

%dsp-level:0
\[
\qquad \mathrm{FT.40}:\ q_{\mathrm{core}}(c;v)=\tfrac12(\sigma_{e_1}-\sigma_{e_k})\ \text{on the closed core};
\]

%dsp-level:0
\[
\qquad \mathrm{CT3.22}:\ \neg\exists\ \text{isolated triplet pole}
\]

%dsp-level:0
\[
(3):\ \mathrm{EM.24}(5)\ \text{verbatim (frozen)};
\]

%dsp-level:0
\[
\qquad \mathrm{CT0.3}\ \text{verbatim};
\]

%dsp-level:0
\[
\qquad \mathrm{FT.86}(8):\ \neg\exists\ \text{tag coupling}\ \delta\tau\ \text{to}\ \Omega/N_i;
\]

%dsp-level:0
\[
\qquad \text{Def-H.19}:\ \text{Def-H.19 vocabulary asserts no coupling to materialization}
\]

%dsp-level:0
\[
(4),(5):\ \text{citation conjunctions}\ [\mathrm{FT.15}\ \text{hypothesis verbatim};\ \mathrm{FT.90}(5)\ \text{audit verbatim}].\ \blacksquare
\]

\textbf{Phys.} The tier knows exactly where it stops. The matrix is static because the corpus has no time to evolve it with — and the consultation of the CPT record confirms that the configuration-level pairing, though exact, cannot supply one: reversal phases are amplitudes, not clocks. What survives as the complete static deliverable is sharp: bound chains only; charge only as the core read; flipping only from $\omega$-odd data; frequency only after time.

\textbf{Rmk (forward pointer, non-load-bearing).} The parity structures invoked by this part's census rows receive their systematic treatment in Tier F.41: the configuration-space reversal calculus with the exact commutations and both frozen CPT lifts (FT.133); the carrier bookkeeping — C, P, T fix the doublet rays and CPT acts as the scalar character $\lambda=(-1)^{|E|}e^{i\pi\Theta}$ (FT.134), so this part's splitting $2|q|$ is parity-even and charge-odd (FT.139(3)); the branch, cover, flicker, pole and channel censuses (FT.135–FT.139). The carrier-level parity-eigenvalue assignment to species is application-layer content.

\vspace{2pt}\noindent\rule{\textwidth}{0.4pt}\vspace{2pt}

\end{theorem}

\phantomsection
\section[Tier F.41 — Parity: The Configuration Calculus, the Carrier…]{\texorpdfstring{Tier F.41 — Parity: The Configuration Calculus, the Carrier Bookkeeping, and the Sector Census}{Tier F.41 - Parity: The Configuration Calculus, the Carrier Bookkeeping, and the Sector Census}}

\begin{theorem}[{\texorpdfstring{FT.133 — The configuration-space reversal calculus}{FT.133 - The configuration-space reversal calculus}}]\label{tag:FT-133}
\textbf{Thm.} On the configuration space $\mathfrak C(\mathcal R)$ [D0.10], with configuration data $(R,\ \sigma,\ \alpha,\ \kappa/\text{ledger})$ [A4, A8, D0.5, D3.5, Def-EM.20]:

%dsp-level:0
\[
(1)\quad\textbf{the four register involutions}\ [\text{each frozen, re-derived at the point of use}]:
\]

%dsp-level:0
\[
\omega\text{-pattern}:\ R\mapsto R^{\mathrm T}=\{(b,a):(a,b)\in R\}\ \text{with edge data carried}\ [\mathrm{D2.6};\ \mathrm{FT.88}(1)];
\]

%dsp-level:0
\[
\qquad C:\ \sigma_e\mapsto-\sigma_e\ [\mathrm{EM.24}(1)];
\]

%dsp-level:0
\[
P:\ \alpha_e\mapsto\pi-\alpha_e\ \ (\mathcal P\mapsto E_{\mathcal R}\setminus\mathcal P)\ [\mathrm{EM.24}(2),\ \mathrm{Def\text{-}EM.21}];
\]

%dsp-level:0
\[
\qquad T=\Theta_{\mathfrak L}:\ \text{ledger order reversed},\ J^+\leftrightarrow J^-,\ \kappa_e\mapsto-\kappa_e,\ \sigma\ \text{fixed}\ [\mathrm{D2.7};\ \mathrm{FT.88}(2)];
\]

%dsp-level:0
\[
\text{each is an involution on configuration data}\ [\text{machine: exhaustive on}\ 17^2+5^6\ \text{configurations}\ +\ 3000\ \text{deterministic off-shell samples}];
\]

%dsp-level:0
\[
(2)\quad\textbf{register-disjoint commutation:}\qquad[C,P]=[C,T]=[P,T]=[C,\omega]=[P,\omega]=[T,\omega]=0
\]

%dsp-level:0
\[
[\text{the four act on pairwise-disjoint data coordinates: edge-direction set};\ \sigma;\ \alpha;\ (\kappa,\text{ledger})]\implies\langle C,
\]

%dsp-level:0
\[
P,
\]

%dsp-level:0
\[
T\rangle\cong\mathbb Z_2^3,
\]

%dsp-level:0
\[
\qquad \langle\omega\text{-pattern},\,C,\,P,\,T\rangle\cong\mathbb Z_2^4;
\]

%dsp-level:0
\[
\mathrm{FT.88}(3)\text{'s}\ \langle\omega,
\]

%dsp-level:0
\[
\Theta_{\mathfrak L}\rangle\cong\mathbb Z_2\times\mathbb Z_2\ \text{is the frozen two-generator subalgebra}\ [\text{machine: exhaustive commutation census}];
\]

%dsp-level:0
\[
(3)\quad\textbf{component weight identities}\ (w=\mathcal J\cdot S\cdot\Gamma\cdot A,\ S=\prod\sigma_e,\ \Gamma=\prod e^{i\alpha_e\widehat{ds}^{\,2}(e)},\ A=(-1)^{\aleph}\ [\mathrm{A8},\ \mathrm{D3.5},\ \mathrm{D3.6}]):
\]

%dsp-level:0
\[
w(C\mathcal C)=(-1)^{|E|}\,w(\mathcal C)\ [\aleph\ \text{invariant}:\ \aleph_v=|f_v|,\ f_v\mapsto-f_v];
\]

%dsp-level:0
\[
\qquad w(T\mathcal C)=w(\mathcal C)^{*}\ [\mathrm{CT2.8}(1)];
\]

%dsp-level:0
\[
\qquad w(P\mathcal C)=e^{i\pi\Theta}\,w(\mathcal C)^{*};
\]

%dsp-level:0
\[
\Theta:
\]

%dsp-level:0
\[
= \!\!\sum_{e\in R}\!\!\widehat{ds}^{\,2}(e)\ [\mathrm{EM.24}(5)]:\quad\Theta\ \text{is}\ P\text{-even},
\]

%dsp-level:0
\[
\ T\text{-odd}\ (\Theta\mapsto-\Theta),
\]

%dsp-level:0
\[
\ \omega\text{-even}\ (\text{an edge-set function};\ \widehat{ds}^{\,2}\ \text{direction-symmetric});
\]

%dsp-level:0
\[
(4)\quad\textbf{the composition, in both frozen lifts:}\qquad\mathrm{CPT}:=C\circ P\circ T\ [\mathrm{EM.24}(4)];
\]

%dsp-level:0
\[
\text{conjugation lift}\ [\mathrm{EM.24}(5)\ \mathrm{Pf}:\ T\ \text{as}\ \Gamma\mapsto\bar\Gamma]:\qquad w(\mathrm{CPT}\,\mathcal C)=(-1)^{|E|}\,e^{+i\pi\Theta}\,w(\mathcal C);
\]

%dsp-level:0
\[
\kappa\text{-flip lift}\ [\mathrm{FT.88}(2)\ \text{register picture}:\ T:\kappa\mapsto-\kappa]:\qquad w(\mathrm{CPT}\,\mathcal C)=(-1)^{|E|}\,e^{-i\pi\Theta}\,w(\mathcal C);
\]

%dsp-level:0
\[
\text{the two lifts agree exactly on the integer-}\Theta\ \text{locus}\ (e^{2i\pi\Theta}=1),
\]

%dsp-level:0
\[
\ \text{which contains every materialized closed chain}\ [\mathrm{FT.22}:\ \Theta\in2\mathbb Z];
\]

%dsp-level:0
\[
\qquad \bigl|w(\mathrm{CPT}\,\mathcal C)\bigr|^2=\bigl|w(\mathcal C)\bigr|^2\ \text{unconditionally in both lifts}\ [\text{machine}];
\]

%dsp-level:0
\[
(5)\quad\textbf{on-shell reductions:}
\]

%dsp-level:0
\[
\qquad \Theta\in2\mathbb Z\implies e^{i\pi\Theta}=1:\ \text{the CPT character reduces to}\ (-1)^{|E|};
\]

%dsp-level:0
\[
\qquad |E|+\Theta\equiv0\ (\mathrm{mod}\ 2)\implies w(\mathrm{CPT}\,\mathcal C)=w(\mathcal C)=w(\mathcal C)^{*}\ [\mathrm{EM.24}(5)\ \text{third line, re-derived}];
\]

%dsp-level:0
\[
(6)\quad\textbf{the read-covariance table}\ [\text{exact; machine on closed digon and triangle chains}]:
\]

%dsp-level:0
\[
\begin{array}{c|cccc}
\text{read} & C & P & T & \omega\\ \hline
q_{\mathrm{core}}\ [\mathrm{FT.40}] & - & + & + & -\\
\Theta=\sum\widehat{ds}^{\,2} & + & + & - & +\\
|E|,\ K=|R| & + & + & + & +\\
\aleph\ [\mathrm{EM.24}(1)] & + & + & + & +\\
\mathcal P\ [\mathrm{Def\text{-}EM.21}] & + & \text{complement} & + & \text{transported}
\end{array}
\]

%dsp-level:0
\[
(7)\quad\textbf{the parity naming (derived record):}\ \ \text{the parity operation on configuration space is}\ \mathrm{EM.24}(2)\text{'s}\ P\ —\ \textbf{derived}\ \text{from the dynamic phase}\ [\mathrm{Def\text{-}EM.20},\ \mathrm{EM.20}(1)]\ \text{and named inside the derived CPT bridge};
\]

%dsp-level:0
\[
\ \mathrm{FT.101}(3)\text{'s branch flip}\ \iota(\alpha)
\]

%dsp-level:0
\[
= \pi-\alpha\ \textbf{is}\ P|_{\Phi}\ \text{on the A4 label alphabet}\ [\mathrm{FT.101}(3)\ \text{cites}\ \mathrm{EM.24}(2),\ \mathrm{Def\text{-}EM.21}]:\ \text{the identification is derived content and no declaration issues}.
\]

\textbf{Req.} A4, A8, D0.5, D0.10, D1.6, D2.6, D2.7, D3.3, D3.4, D3.5, D3.6 (Parts 1–2); CT2.8 (Part 2); Def-EM.20, Def-EM.21, EM.9, EM.9a, EM.20, EM.23, EM.24 (Part 2); FT.22 (Tiers F.1–F.11); FT.40 (Tiers F.12–F.13); FT.88 (Tiers F.29–F.30); FT.101 (Tiers F.31–F.34); exhaustive enumeration, exact $\mathbb Q(i)$ weight arithmetic [math; machine: \texttt{checks/\allowbreak{}verify\_\allowbreak{}H\_\allowbreak{}parity\_\allowbreak{}identification.py} §A].

\textbf{Pf.}

%dsp-level:0
\[
(1):\ \text{each map is an involution on its register}:\ (-\sigma)
\]

%dsp-level:0
\[
= \sigma;
\]

%dsp-level:0
\[
\ \pi-(\pi-\alpha)
\]

%dsp-level:0
\[
= \alpha;
\]

%dsp-level:0
\[
\ (-\kappa)\mapsto\kappa;
\]

%dsp-level:0
\[
\ (R^{\mathrm T})^{\mathrm T}
\]

%dsp-level:0
\[
= R;
\]

%dsp-level:0
\[
\ \text{data transport is a bijection}\ [\text{machine, exhaustive}]
\]

%dsp-level:0
\[
(2):\ \text{disjoint coordinates}\implies\text{maps commute pointwise}\ [\text{machine on all enumerated configurations}];
\]

%dsp-level:0
\[
\qquad \text{group generation}:\ \mathbb Z_2^3,\ \mathbb Z_2^4\ [\text{distinct nonidentity elements counted}]
\]

%dsp-level:0
\[
(3):\ C:\ S\mapsto(-1)^{|E|}S,\ \Gamma,\mathcal J,A\ \text{fixed}\ [\aleph_v=|f_v|,\ f_v\mapsto-f_v\ [\mathrm{D3.3},\ \mathrm{EM.24}(1)]];
\]

%dsp-level:0
\[
T:\ \Gamma\mapsto\bar\Gamma\ [\kappa\mapsto-\kappa:\ e^{i\alpha(-\kappa)\pi\cdots}=\overline{e^{i\alpha\kappa\cdots}};\ \mathcal J\mapsto\mathcal J^{*}\ [\mathrm{D2.7},\ \mathrm{D3.6}];\ S,A\in\{\pm1\}];
\]

%dsp-level:0
\[
\qquad P:\ e^{i(\pi-\alpha_e)\widehat{ds}^{\,2}(e)}=e^{i\pi\widehat{ds}^{\,2}(e)}\,e^{-i\alpha_e\widehat{ds}^{\,2}(e)}\implies\Gamma\mapsto e^{i\pi\Theta}\bar\Gamma
\]

%dsp-level:0
\[
(4):\ \text{composition of}\ (3)\ \text{in each lift};
\]

%dsp-level:0
\[
\qquad e^{+i\pi\Theta}=e^{-i\pi\Theta}\iff e^{2i\pi\Theta}=1\iff\Theta\in\mathbb Z;
\]

%dsp-level:0
\[
\qquad \Theta\in2\mathbb Z\ \text{on materialized closed chains}\ [\mathrm{FT.22}]
\]

%dsp-level:0
\[
|(-1)^{|E|}e^{\pm i\pi\Theta}|
\]

%dsp-level:0
\[
= 1
\]

%dsp-level:0
\[
\implies |w(\mathrm{CPT}\,\mathcal C)|^2
\]

%dsp-level:0
\[
= |w(\mathcal C)|^2\ [\text{machine, both lifts, all enumerated configurations}]
\]

%dsp-level:0
\[
(5):\ \Theta\in2\mathbb Z:\ e^{i\pi\Theta}=(-1)^{\Theta}=1;
\]

%dsp-level:0
\[
\qquad |E|+\Theta\equiv0\ (2):\ (-1)^{|E|}e^{i\pi\Theta}=(-1)^{|E|+\Theta}=1;
\]

%dsp-level:0
\[
\qquad \text{on-shell}\ w\in\mathbb R\ [\mathrm{EM.24}(5)\ \text{Pf re-run}]
\]

%dsp-level:0
\[
(6):\ q_{\mathrm{core}}=\tfrac12(\sigma_{e_1}-\sigma_{e_k})\ [\mathrm{FT.40}]:\ C\ \text{flips both}\ \sigma;\ P,T\ \text{fix}\ \sigma;\ \omega\ \text{swaps}\ e_1\leftrightarrow e_k;
\]

%dsp-level:0
\[
\qquad \Theta,|E|:\ \text{edge-set functions};\ \kappa\mapsto-\kappa\ \text{under}\ T\ [\mathrm{FT.88}(2)]\ \blacksquare
\]

\textbf{Phys.} Parity in this corpus is a register operation, not a postulate: it complements the alignment pattern, and it has done so since the dynamic phase was derived. The four reversals live on four disjoint shelves of the configuration data, so they commute by arithmetic — the elementary abelian group is a census, not an assumption. The CPT identity is the corpus's own CPT theorem; this tier pins its two frozen lifts and their exact agreement locus, which is precisely the locus where materialized chains live. The covariance table is the working ledger of the part: every later clause reads its parities off this table.

\vspace{2pt}\noindent\rule{\textwidth}{0.4pt}\vspace{2pt}

\end{theorem}

\begin{theorem}[{\texorpdfstring{FT.134 — The carrier bookkeeping: the scalar-CPT theorem}{FT.134 - The carrier bookkeeping: the scalar-CPT theorem}}]\label{tag:FT-134}
\textbf{Thm.} On the orientation doublet $\mathcal H_\omega$ [Def-H.33] and the extended carrier [Def-H.34]:

%dsp-level:0
\[
(1)\quad\textbf{the ray census:}\ \ \omega\text{-pattern is the only involution of the FT.133 family that moves the doublet basis}\ (\hat\omega|c\rangle=|\omega(c)\rangle\ [\mathrm{FT.122}(4)]);
\]

%dsp-level:0
\[
\ \ C,
\]

%dsp-level:0
\[
P,
\]

%dsp-level:0
\[
T\ \text{fix the chain tuple pointwise}\ [\mathrm{FT.133}(1)\ \text{registers}]:\ \text{their doublet action is the identity on rays, with a label map};
\]

%dsp-level:0
\[
\text{hence no operator identification}\ \hat\omega
\]

%dsp-level:0
\[
= C\ \text{or}\ \hat\kappa
\]

%dsp-level:0
\[
= P\ \text{exists: the two sides of each proposed equality act on different registers}\ [\text{machine witnesses}];
\]

%dsp-level:0
\[
(2)\quad\textbf{the read-level covariance (the content the identifications reach):}\ \ \text{under Def-H.43's promotion}\ \hat Q_{\mathrm{EM}}
\]

%dsp-level:0
\[
= q_{\mathrm{core}}\hat\kappa\ [\mathrm{FT.129}(4)]:\ \text{the family}\ \{\hat Q_q=q\hat\kappa\}\ \text{is}\ C\text{-odd}\ (\hat Q_{-q}=-\hat Q_q,\ \text{matching}\ q\mapsto-q),
\]

%dsp-level:0
\[
\ P\text{-even and}\ T\text{-even}\ (q_{\mathrm{core}}\ \text{is}\ P/T\text{-even}\ [\mathrm{FT.133}(6)]),
\]

%dsp-level:0
\[
\ \text{and}\ \omega\text{-covariant}\ (\hat\kappa\ \hat\omega\text{-odd}\ \wedge\ q\ \omega\text{-odd}:\ \text{the signs cancel});
\]

%dsp-level:0
\[
\text{consequently Tier F.40's flipping block}\ F_{LR}=q\,I_2\ \text{is}\ C\text{-odd},\ P\text{-even},\ T\text{-even}\ [\mathrm{FT.129}(5)];
\]

%dsp-level:0
\[
(3)\quad\textbf{the scalar-CPT theorem:}\ \ \text{the CPT character}\ \lambda
\]

%dsp-level:0
\[
= (-1)^{|E|}e^{i\pi\Theta}\ \text{takes the same value on both doublet rays}\ (|E|,\Theta\ \text{are edge-set functions}\ [\mathrm{FT.133}(6)]),
\]

%dsp-level:0
\[
\ |\lambda|\qquad = 1:
\]

%dsp-level:0
\[
\mathrm{CPT}\big|_{\mathcal H_\omega}=\lambda\,\mathrm{id},\qquad\lambda=(-1)^{|E|}e^{i\pi\Theta};
\]

%dsp-level:0
\[
\text{chirality is CPT-even: the}\ \hat\chi\text{-eigenspaces are CPT-stable};\ \hat\kappa\ \text{and}\ \hat\omega\ \text{likewise};
\]

%dsp-level:0
\[
\qquad \textbf{refutation witness:}\ \operatorname{tr}(\lambda\,\mathrm{id})=2\lambda\ne0=\operatorname{tr}(\pm\hat\kappa):\ \text{no}\ \hat\kappa\text{-multiple equals the CPT action};
\]

%dsp-level:0
\[
(4)\quad\textbf{the uniqueness pair}\ [\text{re-derived at the point of use}]:\ \ \pm\hat\chi\ \text{are the unique diagonal-free}\ \hat\omega\text{-odd Hermitian involutions}\ [\mathrm{FT.122}(6)];
\]

%dsp-level:0
\[
\ \pm\hat\kappa\ \text{the unique diagonal ones}\ [\mathrm{FT.129}(1)];
\]

%dsp-level:0
\[
\ \text{quaternionic triple}\ [\mathrm{Def\text{-}H.34}]:\ \hat\omega\hat\chi
\]

%dsp-level:0
\[
= -i\hat\kappa,
\]

%dsp-level:0
\[
\ \hat\chi\hat\kappa
\]

%dsp-level:0
\[
= -i\hat\omega,
\]

%dsp-level:0
\[
\ \hat\kappa\hat\omega\qquad = -i\hat\chi;
\]

%dsp-level:0
\[
(5)\quad\textbf{independence:}\ \ \text{any internal involution}\ G\otimes I_2\ \text{commutes with}\ I_2\otimes\hat\omega\ \text{and their product is a third distinct involution: the configuration-side and carrier-side}\ \mathbb Z_2\text{'s generate}\ \mathbb Z_2\times\mathbb Z_2\ [\mathrm{FT.126}(4)\ \text{pattern};\ \mathrm{FT.88}(3)\ \text{precedent};\ \text{machine}].
\]

\textbf{Req.} Def-H.33, FT.122, FT.126 (Tier F.39); Def-H.34, FT.128, FT.129 (Tier F.40); FT.133; FT.40 (Tiers F.12–F.13); Def-H.34, Def-H.43; exact $\mathbb Q(i)$ arithmetic, trace witnesses, rational canonicity census [math; machine: \texttt{checks/\allowbreak{}verify\_\allowbreak{}H\_\allowbreak{}parity\_\allowbreak{}identification.py} §B].

\textbf{Pf.}

%dsp-level:0
\[
(1):\ \text{register census of FT.133(1):}\ C,
\]

%dsp-level:0
\[
P,
\]

%dsp-level:0
\[
T\ \text{act on}\ (\sigma),
\]

%dsp-level:0
\[
(\alpha),
\]

%dsp-level:0
\[
(\kappa,\text{ledger})\ \text{and leave the edge tuple fixed}\ [\text{machine: key sets preserved on all enumerated configurations}];
\]

%dsp-level:0
\[
\qquad \omega\text{-pattern transposes it}\ [\text{machine}]
\]

%dsp-level:0
\[
\hat\omega\ne\pm\mathrm{id}\ [\text{entrywise}];
\]

%dsp-level:0
\[
\qquad \text{an operator equality between a ray-swapping and a ray-fixing map is impossible}
\]

%dsp-level:0
\[
(2):\ \hat Q_{-q}=-\hat Q_q\ [\text{linearity}];
\]

%dsp-level:0
\[
\qquad q_{\mathrm{core}}(P\mathcal C)=q_{\mathrm{core}}(T\mathcal C)=q_{\mathrm{core}}(\mathcal C),\quad q_{\mathrm{core}}(\omega\mathcal C)=-q_{\mathrm{core}}(\mathcal C)\ [\mathrm{FT.133}(6)];
\]

%dsp-level:0
\[
\qquad \hat\omega\hat\kappa\hat\omega=-\hat\kappa\ [\mathrm{Def\text{-}H.34}]
\]

%dsp-level:0
\[
\implies\ \text{the coupled operator}\ q\hat\kappa\ \text{is}\ \omega\text{-invariant (both signs cancel) and}\ C\text{-odd};
\]

%dsp-level:0
\[
\qquad F_{LR}=qI_2\ \text{inherits the characters}\ [\mathrm{FT.129}(5)]
\]

%dsp-level:0
\[
(3):\ \lambda(\text{ray}\ c)
\]

%dsp-level:0
\[
= (-1)^{|E(c)|}e^{i\pi\Theta(c)}
\]

%dsp-level:0
\[
= (-1)^{|E(\omega c)|}e^{i\pi\Theta(\omega c)}
\]

%dsp-level:0
\[
= \lambda(\text{ray}\ \omega(c))\ [\mathrm{FT.133}(6):\ \text{both reads edge-set functions};\ \text{machine on closed digon/triangle chains}]
\]

%dsp-level:0
\[
\mathrm{CPT}\ \text{fixes each ray and multiplies the weight by}\ \lambda\ [\mathrm{EM.24}(5);\ \mathrm{FT.133}(4)]\implies\mathrm{CPT}\big|_{\mathcal H_\omega}=\lambda\,\mathrm{id};
\]

%dsp-level:0
\[
\qquad |\lambda|^2=1\ [\text{machine}]
\]

%dsp-level:0
\[
[\lambda\,\mathrm{id},\ \hat\chi]=[\lambda\,\mathrm{id},\ \hat\kappa]=[\lambda\,\mathrm{id},\ \hat\omega]=0;\qquad\operatorname{tr}(\lambda\,\mathrm{id})=2\lambda\ne0=\operatorname{tr}(\hat\kappa)=\operatorname{tr}(-\hat\kappa)\ [\lambda\ne0]
\]

%dsp-level:0
\[
(4):\ \text{rational census: diagonal}\ \hat\omega\text{-odd involutions}\ \{a,-a\}:\ a^2
\]

%dsp-level:0
\[
= 1
\]

%dsp-level:0
\[
\implies a
\]

%dsp-level:0
\[
= \pm1;
\]

%dsp-level:0
\[
\ \text{diagonal-free}\ \{iy,-iy\}:\ y^2
\]

%dsp-level:0
\[
= 1
\]

%dsp-level:0
\[
\implies y
\]

%dsp-level:0
\[
= \pm1\ [\text{machine}];
\]

%dsp-level:0
\[
\qquad \text{triple products entrywise}\ [\text{machine}]
\]

%dsp-level:0
\[
(5):\ [G\otimes I,\ I\otimes\hat\omega]=0\ [\text{Kronecker disjointness}];
\]

%dsp-level:0
\[
\qquad (G\otimes\hat\omega)^2=G^2\otimes\hat\omega^2=\mathrm{id};
\]

%dsp-level:0
\[
\qquad G\otimes\hat\omega\ne G\otimes I\ne I\otimes\hat\omega\ [\text{entrywise witnesses};\ \text{machine}].\ \blacksquare
\]

\textbf{Phys.} The carrier keeps its own books. Charge conjugation, parity and time reversal act on the registration shelves; only the traversal reversal moves the rays of the doublet. The composition CPT — the operation the configuration calculus proves exact — reaches the doublet as a pure phase: the quietest possible action, leaving chirality, the orientation sign and the exchange grade all untouched. What the proposed operator identifications legitimately reach is the read level, and there the census is sharp: the promoted flip block is charge-odd and parity-even, the only parity-sensitive structure in the fermion tower.

\vspace{2pt}\noindent\rule{\textwidth}{0.4pt}\vspace{2pt}

\end{theorem}

\begin{theorem}[{\texorpdfstring{FT.135 — The weak branch grading under parity}{FT.135 - The weak branch grading under parity}}]\label{tag:FT-135}
\textbf{Thm.}

%dsp-level:0
\[
(1)\quad\iota(0)=\pi,\ \iota(\pi)=0\ \text{on the branch alphabet}\ \Phi=\{0,\pi\}\ [\mathrm{FT.101}(3);\ \mathrm{D0.5}];
\]

%dsp-level:0
\[
\qquad \iota
\]

%dsp-level:0
\[
= P|_{\Phi}\ [\mathrm{EM.24}(2);\ \mathrm{FT.101}(3)\ \text{citation chain}]:\ \text{parity exchanges the two crystallization branches at the label level};
\]

%dsp-level:0
\[
(2)\quad\textbf{aligned branch}\ (\alpha=0,\ \vee^2,\ \text{Casimir}\ 2):\ W_i
\]

%dsp-level:0
\[
= \Pi_{\mathcal R}W_i\Pi_{\mathcal R}
\]

%dsp-level:0
\[
\ne 0,
\]

%dsp-level:0
\[
\ \sum_iW_i^2
\]

%dsp-level:0
\[
= 2\Pi_{\mathcal R}\ [\mathrm{FT.101}(1),\ \text{Cor}\ \mathrm{CT0.7a}\ \text{re-derived}];
\]

%dsp-level:0
\[
\qquad \textbf{antialigned branch}\ (\alpha=\pi,\ \wedge^2,\ \text{Casimir}\ 0):\ W_i\Pi_{\mathcal V}
\]

%dsp-level:0
\[
= 0,
\]

%dsp-level:0
\[
\ \bigcap_i\ker W_i
\]

%dsp-level:0
\[
= \operatorname{im}\Pi_{\mathcal V}\ [\mathrm{FT.101}(2),\ \mathrm{FT.93}(5)\ \text{re-derived}];
\]

%dsp-level:0
\[
(3)
\]

%dsp-level:0
\[
\quad P\ \text{maps the weak-active branch label to the weak-blind one: the weak support is not flip-invariant}\ [\mathrm{FT.101}(4)];
\]

%dsp-level:0
\[
\ \text{structural parity non-invariance of the force arena}\ [\mathrm{FT.101}(5)];
\]

%dsp-level:0
\[
\ \text{scope: arena}\ S\text{-parity grading},
\]

%dsp-level:0
\[
\ \neg\ \text{matter chirality}\ [\mathrm{FT.93}(3),\ \mathrm{FT.93}(7)];
\]

%dsp-level:0
\[
(4)\quad\textbf{the exchange is label-level only:}
\]

%dsp-level:0
\[
\qquad \mathfrak T\Pi_{\mathcal R}=3\ne1=\mathfrak T\Pi_{\mathcal V};
\]

%dsp-level:0
\[
\qquad M=a\Pi_{\mathcal R}+b\Pi_{\mathcal V}:\ M^2=\mathrm{id}\implies a,b\in\{\pm1\}\implies\operatorname{tr}M=3a+b\in\{\pm4,\pm2\}\ne0:
\]

%dsp-level:0
\[
\text{no traceless involution implements the branch exchange on the pair register}\ [\mathrm{FT.101}(3),\ \mathrm{FT.93}(3a)\ \text{re-derived}];
\]

%dsp-level:0
\[
\qquad \mathbb S\Pi_{\mathcal R}=+\Pi_{\mathcal R},\ \mathbb S\Pi_{\mathcal V}=-\Pi_{\mathcal V}\ [\text{the}\ S\text{-eigenvalue labels}].
\]

\textbf{Req.} A3, A4 (Parts 1–2); CT0.7, Cor CT0.7a, CT0.8, CT1.2 (Part 1); Def-EM.21, EM.24 (Part 2); FT.93, FT.101 (Tiers F.31–F.34); exact $\mathbb Q(i)$ matrix arithmetic, trace/rank invariants, exact nullspace [math; machine: \texttt{checks/\allowbreak{}verify\_\allowbreak{}H\_\allowbreak{}parity\_\allowbreak{}identification.py} §C].

\textbf{Pf.}

%dsp-level:0
\[
(1):\ \iota(\alpha)=\pi-\alpha\ [\mathrm{FT.101}(3):\ \text{the unique nontrivial automorphism of}\ \Phi];
\]

%dsp-level:0
\[
\qquad P:\ \alpha_e\mapsto\pi-\alpha_e\ [\mathrm{EM.24}(2)]:\ \iota=P|_{\Phi}\ [\text{same map on the same alphabet}]
\]

%dsp-level:0
\[
(2):\ \Pi_{\mathcal R}
\]

%dsp-level:0
\[
= \tfrac12(E_0+\mathbb S)
\]

%dsp-level:0
\[
= \tfrac14(3E_0+\Sigma),
\]

%dsp-level:0
\[
\ \Pi_{\mathcal V}
\]

%dsp-level:0
\[
= \tfrac14(E_0-\Sigma)\ [\mathrm{CT0.7}(1)]:\ \text{idempotent, orthogonal, complete, traces}\ (3,1)\ [\text{machine}]
\]

%dsp-level:0
\[
W_i\Pi_{\mathcal V}=0\ \forall i;
\]

%dsp-level:0
\[
\quad W_i
\]

%dsp-level:0
\[
= \Pi_{\mathcal R}W_i\Pi_{\mathcal R};
\]

%dsp-level:0
\[
\quad\textstyle\sum_iW_i^2
\]

%dsp-level:0
\[
= \tfrac32E_0+\tfrac12\Sigma
\]

%dsp-level:0
\[
= 2\Pi_{\mathcal R};
\]

%dsp-level:0
\[
\quad\dim\ker(W\text{-stack})
\]

%dsp-level:0
\[
= 1\ (\text{over}\ \mathbb C)
\]

%dsp-level:0
\[
= \operatorname{im}\Pi_{\mathcal V}\ [\text{exact nullspace};\ \text{machine}]
\]

%dsp-level:0
\[
(3):\ \mathrm{FT.101}(4)\text{–}(5)\ \text{verbatim (frozen)};
\]

%dsp-level:0
\[
\qquad \text{the}\ \neg\text{-scope is FT.93's pincer and relocation}\ [\mathrm{FT.125}(1)\ \text{status stands: the carrier chirality of Tier F.39 is a distinct structure}]
\]

%dsp-level:0
\[
(4):\ \text{sign enumeration over}\ \{a,b\}\subset\{\pm1\}^2:\ \operatorname{tr}\in\{4,2,-2,-4\}\ [\text{machine}].\ \blacksquare
\]

\textbf{Phys.} Parity has a derived, exact and destructive action on the weak force: it swaps the branch the weak sees for the branch it annihilates. The non-invariance is structural — a $\mathbb{Z}$$_2$-asymmetry of the force on the registration alphabet — and it has no operator on the pair register, by trace and rank. What the force arena loses under parity, the carrier never had: the doublet's chirality grade commutes with the weak action entirely [FT.123(1)], and the two parity stories meet only in the assignment rule that does not yet exist.

\vspace{2pt}\noindent\rule{\textwidth}{0.4pt}\vspace{2pt}

\end{theorem}

\begin{theorem}[{\texorpdfstring{FT.136 — The rung-one S-grading: the 10/6 split}{FT.136 - The rung-one S-grading: the 10/6 split}}]\label{tag:FT-136}
\textbf{Thm.}

%dsp-level:0
\[
(1)\quad\textbf{the}\ \operatorname{Ad}(\mathbb S)\ \textbf{census on the CT1.2 generators}\ [\text{machine: all}\ 17]:
\]

%dsp-level:0
\[
\qquad S\text{-even}\ (+1):\ \{E_0,\ \Sigma,\ \hat E^{ij}\ (5\ \text{dims}),\ W_i\ (3)\}=10\ \text{dims}=\underbrace{G_5+W_3+H_1}_{\text{forces}}+\underbrace{D^+_1}_{\text{central}};
\]

%dsp-level:0
\[
S\text{-odd}\ (-1):\ \{J_i\ (3),\ K_i\ (3)\}=6\ \text{dims}=J^+_3+J^-_3
\]

%dsp-level:0
\[
\qquad [\mathrm{CT1.2}\ (S,T)\ \text{table};\ \mathrm{CT0.12};\ \mathrm{FT.30}\ \text{tiling}\ (5{+}3{+}1)+(3{+}3)+1];
\]

%dsp-level:0
\[
(2)\quad\mathcal A(1)
\]

%dsp-level:0
\[
= \{J,\,G,\,W,\,H,\,D^+\}\ [\mathrm{CT0.13}]:\ \text{the rung-one active set}\ \textbf{is}\ \text{the}\ 10/6\ \text{split};
\]

%dsp-level:0
\[
\ \text{the metric dictionary}\ [\mathrm{FT.30}]:\ \text{the}\ S\text{-even block carries every metric coordinate}\ (\hat h_{ij}\ 5,\ \Sigma_{\mathrm{tr}}\ 1,\ a_{00}\ 1,\ w_i\ 3);
\]

%dsp-level:0
\[
\ \text{the}\ S\text{-odd Janus seals carry none};
\]

%dsp-level:0
\[
(3)\quad\textbf{register separation:}\ \ P\ (\alpha\text{-complement})\ \text{acts on configuration data and has no action on the pair algebra}\ [\mathrm{FT.133}(2)];
\]

%dsp-level:0
\[
\ \text{the only derived bridge between parity and the}\ S\text{-grading is FT.101's label}\leftrightarrow\text{arena coupling}\ (\text{branch}\ \alpha\leftrightarrow\vee^2/\wedge^2\ \text{weak support}):\ \text{a naming of the}\ S\text{-blocks as}\ P\text{-sectors has no operator content — the two}\ \mathbb Z_2\text{'s act on different spaces};
\]

%dsp-level:0
\[
(4)\quad\textbf{the independence census:}\ \ \text{three pairwise-distinct}\ \mathbb Z_2\ \text{gradings stand in the record: the algebra}\ S\text{-grading}\ (\operatorname{Ad}\mathbb S),
\]

%dsp-level:0
\[
\ \text{the branch-label parity}\ (P/\iota\ [\mathrm{FT.101}]),
\]

%dsp-level:0
\[
\ \text{the carrier chirality}\ (\hat\chi\ [\mathrm{FT.122}])\ —\ \text{pairwise commuting where jointly defined, none reducible to another}\ [\mathrm{FT.126}(4);\ \text{machine}]:
\]

%dsp-level:0
\[
\text{the}\ 10/6\ \text{asymmetry is the force/dynamics split of rung one; it is not the source of the carrier chirality grade};
\]

%dsp-level:0
\[
(5)\quad\textbf{the weak location:}\ \ W\ \text{is}\ S\text{-even}\ [\text{clause}\ (1)];
\]

%dsp-level:0
\[
\ \text{its chiral selectivity is the arena}\ 3\oplus1\ \text{grading}\ [\mathrm{FT.93}(5):\ W_i\Pi_{\mathcal V}=0,\ \sum W_i^2=2\Pi_{\mathcal R}]\ —\ \text{a third structure, distinct from both the}\ S\text{-grading of the operator algebra and the}\ \hat\chi\text{-grading of the carrier}.
\]

\textbf{Req.} CT0.7, CT0.12, CT0.13, CT1.1, CT1.2 (Part 1); FT.30 (Tiers F.1–F.11); FT.93, FT.101 (Tiers F.31–F.34); FT.122, FT.126 (Tier F.39); FT.133 (this tier); exact $\operatorname{Ad}(\mathbb S)$ eigen-census over $\mathbb Q(i)$ [math; machine: \texttt{checks/\allowbreak{}verify\_\allowbreak{}H\_\allowbreak{}parity\_\allowbreak{}identification.py} §D].

\textbf{Pf.}

%dsp-level:0
\[
(1):\ \mathbb S=\tfrac12(E_0+\Sigma),\ \mathbb S^2=E_0\ [\mathrm{CT0.7}(1);\ \Sigma^2=3E_0-2\Sigma];
\]

%dsp-level:0
\[
\qquad \operatorname{Ad}(\mathbb S)X=\mathbb SX\mathbb S=\pm X\ \text{on each generator}\ [\text{machine: all}\ 17,\ \text{entrywise}]
\]

%dsp-level:0
\[
\text{membership:}\ \mathbb S E_0\mathbb S
\]

%dsp-level:0
\[
= E_0,
\]

%dsp-level:0
\[
\ \mathbb S\Sigma\mathbb S
\]

%dsp-level:0
\[
= \Sigma\ [\mathrm{CT1.3},\ \mathrm{CT1.5}];
\]

%dsp-level:0
\[
\quad\mathbb S\hat E^{ij}\mathbb S
\]

%dsp-level:0
\[
= \hat E^{ij}\ [\mathrm{CT1.2}\ \text{table}];
\]

%dsp-level:0
\[
\quad\mathbb SW_i\mathbb S
\]

%dsp-level:0
\[
= W_i;
\]

%dsp-level:0
\[
\quad\mathbb SJ_i\mathbb S
\]

%dsp-level:0
\[
= -J_i,
\]

%dsp-level:0
\[
\ \mathbb SK_i\mathbb S
\]

%dsp-level:0
\[
= -K_i\ [\mathrm{CT0.7}(3)]
\]

%dsp-level:0
\[
\text{dimension count}:\ 1+1+5+3=10;\quad3+3=6;\qquad10+6=16=C_{\mathrm{cap}}\ [\mathrm{CT0.12};\ \mathrm{CT0.6}]
\]

%dsp-level:0
\[
(2):\ \mathrm{CT0.13}\ \text{verbatim (frozen)};\qquad\mathrm{FT.30}\ \text{dictionary rows verbatim}
\]

%dsp-level:0
\[
(3):\ \mathrm{FT.133}(1)\text{–}(2):\ P\ \text{acts on}\ \alpha;
\]

%dsp-level:0
\[
\ \text{the pair algebra is spanned by}\ \{E_0,\Sigma,\hat E^{ij},W_i,J_i,K_i\}\ \text{with no}\ \alpha\ \text{argument};
\]

%dsp-level:0
\[
\qquad \mathrm{FT.101}(1)\text{–}(2):\ \text{the label-arena coupling}
\]

%dsp-level:0
\[
(4):\ [G\otimes I,\ I\otimes\hat\omega]=0,\ (G\otimes\hat\omega)^2=\mathrm{id},\ G\otimes\hat\omega\ne G\otimes I\ne I\otimes\hat\omega\ [\mathrm{FT.134}(5);\ \text{machine}];
\]

%dsp-level:0
\[
\qquad \operatorname{Ad}(\mathbb S)\ \text{acts on}\ \operatorname{End}(\mathbb C^4_{\mathrm{pair}}),\ \iota\ \text{on labels},\ \hat\chi\ \text{on}\ \mathcal H_\omega:\ \text{three distinct spaces}
\]

%dsp-level:0
\[
(5):\ W_i\in\mathcal E_{+-}\ [\mathrm{CT1.2}\ \text{table}];\qquad\mathrm{FT.93}(5)\ \text{data re-derived in}\ \S\mathrm{C}.\ \blacksquare
\]

\textbf{Phys.} At the first rung the algebra already divides: ten dimensions that speak to the metric — shear, volume, conformal, shift — against six that only move the ledger. That is the rung-one asymmetry, and it is exact. Parity names the branch labels that select the weak support inside the arena; the S-grading sorts the operators; the carrier chirality grades the chains. The three $\mathbb{Z}$$_2$'s commute and never collapse into one another — which is why no rung-one count generates a family, and why the family channels of Tiers F.37–F.38 run on their own integers.

\vspace{2pt}\noindent\rule{\textwidth}{0.4pt}\vspace{2pt}

\end{theorem}

\begin{theorem}[{\texorpdfstring{FT.137 — The cover and flicker census under the reversal family}{FT.137 - The cover and flicker census under the reversal family}}]\label{tag:FT-137}
\textbf{Thm.}

%dsp-level:0
\[
(1)\quad\textbf{pattern reversal}\ (\omega\text{-side}\ [\mathrm{FT.88}(1):\ \omega\ \text{acts on skeletons as}\ R\mapsto R^{\mathrm T}]):\ \ \text{preserves admissibility},
\]

%dsp-level:0
\[
\ K
\]

%dsp-level:0
\[
= |R|,
\]

%dsp-level:0
\[
\ \text{cover status and the amplitude}:\ v(R^{\mathrm T})
\]

%dsp-level:0
\[
= v(R)\ \text{on all}\ 4096\ \text{tetrahedron patterns}\ [\mathrm{FT.95}(3)\ \text{re-derived};\ \text{machine}];
\]

%dsp-level:0
\[
\ 101\ \text{covers reproduced};
\]

%dsp-level:0
\[
(2)\quad\textbf{the chiral cover pair:}\ \ \text{the}\ 48\ \text{dissolved-digon bypass covers}\ [\mathrm{FT.69}(4)]\ \text{split into two}\ \Omega
\]

%dsp-level:0
\[
= 24\ \operatorname{Aut}\text{-orbits exchanged by reversal};
\]

%dsp-level:0
\[
\ \text{no orbit is reversal-invariant}\ (R\mapsto R^{\mathrm T}\notin\operatorname{Aut}(\Gamma))\ [\mathrm{FT.78}(6)\ \text{re-derived};\ \text{machine}];
\]

%dsp-level:0
\[
(3)\quad\textbf{P (}\alpha\text{-complement) on the cover sector:}\ \ \text{the exact sign character}
\]

%dsp-level:0
\[
\qquad v_{\mathbf a}(R)=(-1)^{\sum_{e\in R}a_e}\,v_{\mathbf 0}(R)
\]

%dsp-level:0
\[
\qquad [\text{machine: exhaustive triangle census over all}\ (R,\mathbf a,\sigma)];
\]

%dsp-level:0
\[
\text{hence cover status}\ (v\ne0),\ |v|,\ K,\ 4^{|\mathcal C|},\ \text{and the norm layer}\ 2^{|R|}\ \text{are}\ P\text{-invariant};
\]

%dsp-level:0
\[
\qquad \text{weights:}\ w(P\mathcal C)
\]

%dsp-level:0
\[
= e^{i\pi\Theta}w(\mathcal C)^{*}\ [\mathrm{FT.133}(3)]:\ |w|^2\ \text{invariant};
\]

%dsp-level:0
\[
\ \text{materialized even-}\Theta\ \text{locus}:\ w(P\mathcal C)
\]

%dsp-level:0
\[
= w(\mathcal C)\ \text{exactly};
\]

%dsp-level:0
\[
(4)\quad\textbf{the flicker reads:}\ \ \text{the bypass/lock decomposition}\ [\mathrm{Def\text{-}H.25},\ \mathrm{FT.68}],
\]

%dsp-level:0
\[
\ B_\Gamma(p)\ \text{and the phase boundary}\ p^*_\Gamma\ [\mathrm{FT.70}(1)],
\]

%dsp-level:0
\[
\ f_{\mathrm{adm}}(n,p)\ [\mathrm{FT.70}(4)],
\]

%dsp-level:0
\[
\ \text{the dust density}\ \langle\rho_{D^-}\rangle_\Gamma(p)
\]

%dsp-level:0
\[
= \sum_vO(n_v)(1-f_{\mathrm{adm}}(n_v,p)),
\]

%dsp-level:0
\[
\ O(n)
\]

%dsp-level:0
\[
= 2^{n+1}-16\ [\mathrm{FT.71}(1),\ \mathrm{CT4.1}]
\]

%dsp-level:0
\[
\text{are functions of}\ (\Gamma,R,p,n_v)\ \text{alone}\ —\ \alpha\text{-blind and}\ \sigma\text{-blind by definitional census}:\ P\text{-},
\]

%dsp-level:0
\[
\ C\text{-},
\]

%dsp-level:0
\[
\ T\text{- and}\ \omega\text{-invariant};
\]

%dsp-level:0
\[
\ \text{the only}\ P\text{-sensitive quantity in the sector is clause}\ (3)\text{'s weight character};
\]

%dsp-level:0
\[
(5)\quad\textbf{the attribution census:}\ \ \text{the cover sector's nontrivial chirality}\ —\ \text{clause}\ (2)\text{'s exchanged pair}\ —\ \text{belongs to the pattern reversal}\ [\mathrm{FT.78}(6):\ \text{``the pattern-level form of charge conjugation''};\ \mathrm{FT.88}(1):\ \omega\ \text{acts as}\ R\mapsto R^{\mathrm T}];
\]

%dsp-level:0
\[
\ \ P\ \text{acts on the census only through the sign character of clause}\ (3):
\]

%dsp-level:0
\[
\text{reversal and parity are distinct involutions with distinct footprints}\ [\mathrm{FT.89}(2)\ \text{discipline}].
\]

\textbf{Req.} FT.67, FT.68, FT.69, FT.70, FT.71, Def-H.25, Def-H.26 (Tiers F.23–F.25); FT.78 (Tiers F.26–F.28); FT.88 (Tiers F.29–F.30); FT.95 (Tiers F.31–F.34); CT4.1 (Part 2); FT.133 (this tier); exhaustive pattern enumeration, exact brute-force sign sums, $\operatorname{Aut}$-orbit machinery [math; machine: \texttt{checks/\allowbreak{}verify\_\allowbreak{}H\_\allowbreak{}parity\_\allowbreak{}identification.py} §E].

\textbf{Pf.}

%dsp-level:0
\[
(1):\ v(R)
\]

%dsp-level:0
\[
= \sum_\sigma w(R,\sigma)\ \text{by brute-force sign sum}\ [\mathrm{FT.67}(2)\ \text{classifier transport}]:\ v(R^{\mathrm T})
\]

%dsp-level:0
\[
= v(R)\ \text{on all}\ 4096\ [\text{machine}];
\]

%dsp-level:0
\[
\qquad \text{types transpose}:\ (n_{\mathrm{in}},n_{\mathrm{out}})\mapsto(n_{\mathrm{out}},n_{\mathrm{in}}):\ \text{admissibility}\ k_v\le3\ \text{invariant}
\]

%dsp-level:0
\[
\ell\mapsto|\ell|\ \text{at type-}(1,1)\ \text{vertices under}\ \sigma\text{-transport}:\ \aleph\ \text{invariant};
\]

%dsp-level:0
\[
\qquad |R^{\mathrm T}|=|R|;
\]

%dsp-level:0
\[
\qquad \text{cover status}=(v\ne0)\ \text{invariant};
\]

%dsp-level:0
\[
\qquad 101\ \text{covers}\ [\text{machine census}]
\]

%dsp-level:0
\[
(2):\ |\operatorname{Aut}|=24\ [\text{machine}];
\]

%dsp-level:0
\[
\qquad \text{orbit census of the}\ 48:\ \text{two orbits, sizes}\ (24,24);
\]

%dsp-level:0
\[
\qquad \text{reversal maps each orbit bijectively onto the other};
\]

%dsp-level:0
\[
\qquad \text{no fixed orbit}\ [\text{machine}]
\]

%dsp-level:0
\[
(3):\ w(R,\sigma,\mathbf a)=\prod_{e\in R}\sigma_e(-1)^{a_e}\cdot(-1)^{\aleph(\sigma)}\ [\mathrm{A8};\ \kappa\equiv1\ \text{class}]:
\]

%dsp-level:0
\[
\qquad \sum_\sigma w(R,\sigma,\mathbf a)=(-1)^{\sum a_e}\sum_\sigma w(R,\sigma,\mathbf 0)\ [\text{factorization};\ \text{machine exhaustive}]
\]

%dsp-level:0
\[
\text{status}/|v|/K/4^{|\mathcal C|}/2^{|R|}:\ \text{functions of}\ (R,\ |v|)\ \text{only};
\]

%dsp-level:0
\[
\qquad w(P\mathcal C)=e^{i\pi\Theta}w(\mathcal C)^*\ [\mathrm{FT.133}(3)];
\]

%dsp-level:0
\[
\qquad \Theta\in2\mathbb Z\ \text{and real weights}:\ w(P\mathcal C)=w(\mathcal C)
\]

%dsp-level:0
\[
(4):\ \text{definitional census: each formula's argument list}\ [\mathrm{FT.68}(1)\text{–}(5),\ \mathrm{FT.70}(1),(4),\ \mathrm{FT.71}(1),\ \mathrm{CT4.1}]\ \text{contains no}\ \alpha\ \text{and no}\ \sigma;
\]

%dsp-level:0
\[
\qquad (5):\ \mathrm{FT.78}(6),\ \mathrm{FT.88}(1)\ \text{verbatim (frozen)}.\ \blacksquare
\]

\textbf{Phys.} The cover world is chirally paired — but by reversal, not by parity. The two $\Omega$=24 orbits are each other's mirror under the transposition of every arrow, and no automorphism of the skeleton can tell them apart from within: the cover sector carries a genuine $\mathbb{Z}$$_2$ chirality, and it is $\omega$-side. Parity's entire footprint on this sector is one exact sign per configuration — invisible to every status, count, boundary and density the flicker calculus reads.

\vspace{2pt}\noindent\rule{\textwidth}{0.4pt}\vspace{2pt}

\end{theorem}

\begin{theorem}[{\texorpdfstring{FT.138 — The pole spectrum census}{FT.138 - The pole spectrum census}}]\label{tag:FT-138}
\textbf{Thm.}

%dsp-level:0
\[
(1)\quad\textbf{the defect formula is register-blind:}
\]

%dsp-level:0
\[
\qquad D_S(m)=\frac{m^2_{\mathrm{pole}}}{\hat\mu'^2}=2-2\cos\bigl(m\,\delta_v/2\bigr)=4\sin^2\bigl(m\,\delta_v/4\bigr)
\]

%dsp-level:0
\[
\qquad [\mathrm{FT.76},\ \mathrm{FT.60}(5);\ \text{Def-H.37(2)}\ \text{transfer}]:
\]

%dsp-level:0
\[
\text{reads}\ (m,\delta_v)\ \text{alone}\ —\ \text{winding and vertex-figure defect}\ [\mathrm{FT.63}(4)]:\ \alpha\text{-blind},
\]

%dsp-level:0
\[
\ \sigma\text{-blind}
\]

%dsp-level:0
\[
\implies C\text{- and}\ P\text{-invariant};
\]

%dsp-level:0
\[
\qquad T:\ \kappa\mapsto-\kappa\ [\mathrm{FT.88}(2)]\ \text{maps}\ m\mapsto-m:\ D(-m)=D(m)\ \text{exactly}\ (\cos\ \text{even});
\]

%dsp-level:0
\[
(2)\quad\textbf{the five Platonic multiplets, exact in algebraic number fields}\ [\text{machine}]:
\]

%dsp-level:0
\[
\text{tetrahedron}\ (\delta_v=\pi,\ k=3):\ D\equiv2\ \text{on}\ m\in\{3,1,-1,-3\}\ —\ \text{fully degenerate}\ [\mathrm{FT.76}(5a)];
\]

%dsp-level:0
\[
\qquad \text{cube}\ (\pi/2,\ k=2):\ \{2,0,2\}\ —\ \text{the exactly massless even-winding weight}\ [\mathrm{FT.76}(5b)];
\]

%dsp-level:0
\[
\text{octahedron}\ (2\pi/3,\ k=3):\ \{4,1,1,4\}\in\mathbb Q;
\]

%dsp-level:0
\[
\qquad \text{dodecahedron}\ (\pi/5,\ k=5):\ D(m)
\]

%dsp-level:0
\[
= 2-2T_{|m|}(c),
\]

%dsp-level:0
\[
\ c
\]

%dsp-level:0
\[
= \cos\tfrac\pi{10},
\]

%dsp-level:0
\[
\ 16c^4-20c^2+5
\]

%dsp-level:0
\[
= 0\ [\mathbb Q[c]/(\cdot);\ T_5(c)=0\ \text{exact}];
\]

%dsp-level:0
\[
\text{icosahedron}\ (\pi/3,\ k=3):\ \{2,\ 2-\sqrt3,\ 2-\sqrt3,\ 2\}\in\mathbb Q(\sqrt3),\quad(2-\sqrt3)(2+\sqrt3)=1;
\]

%dsp-level:0
\[
\text{every multiplet is}\ m\leftrightarrow-m\ \text{paired}: D(m)=D(-m)\ \text{in each field};
\]

%dsp-level:0
\[
(3)\quad\textbf{pole condition and weights:}\ \ \text{on the materialized even-}\Theta\ \text{locus},
\]

%dsp-level:0
\[
\ P\ \text{acts trivially on weights}\ [\mathrm{FT.137}(3)]:\ \text{the pole census}\ (1\in\operatorname{spec}\mathbf F_M^{(\Sigma)}\ [\mathrm{CT2.5},\ \mathrm{Def\text{-}H.16}])\ \text{is}\ P\text{-invariant there};
\]

%dsp-level:0
\[
\ \text{off-shell the weights acquire the unit character}\ e^{i\pi\Theta}:\ |w|^2\text{-level invariance is unconditional}\ [\mathrm{EM.24}(5)];
\]

%dsp-level:0
\[
\ \text{the character multiplies weights, not the mask}\ [\mathrm{CT2.3}];
\]

%dsp-level:0
\[
(4)\quad\textbf{ceilings are parity-blind:}\ \ \mathrm{FT.76}(3)\text{'s}\ m^2_{\mathrm{pole}}
\]

%dsp-level:0
\[
\in [0,4\hat\mu'^2]\ \text{and FT.77's negative results are range statements about}\ D\ \text{values}\ —\ \text{invariant under the whole reversal family}.
\]

\textbf{Req.} FT.60, FT.63 (Tiers F.20–F.22); FT.76, FT.77 (Tiers F.26–F.28); FT.22 (Tiers F.1–F.11); CT2.3, CT2.5 (Part 2); Def-H.16 (Tiers F.1–F.11); Def-H.37(2); FT.133, FT.137 (this tier); exact Chebyshev arithmetic in $\mathbb Q[c]/(16c^4-20c^2+5)$, exact $\mathbb Q(\sqrt3)$ and $\mathbb Q$ multiplets [math; machine: \texttt{checks/\allowbreak{}verify\_\allowbreak{}H\_\allowbreak{}parity\_\allowbreak{}identification.py} §F].

\textbf{Pf.}

%dsp-level:0
\[
(1):\ \text{argument census of the frozen formula}\ [\mathrm{FT.76}(2)]:\ (m,\delta_v);
\]

%dsp-level:0
\[
\qquad \cos(-x)=\cos x\ [\text{even series}];
\]

%dsp-level:0
\[
\qquad m\mapsto-m\ \text{under}\ \kappa\mapsto-\kappa\ [\mathrm{FT.88}(2);\ k=\Theta/2]
\]

%dsp-level:0
\[
(2):\ \text{tetra}:\ \cos(m\pi/2)=0\ (m\ \text{odd}):\ D=2;
\]

%dsp-level:0
\[
\qquad \text{cube}:\ \cos(m\pi/4)\in\{0,1\}\ \text{at}\ m\in\{2,0,-2\}:\ D\in\{2,0\};
\]

%dsp-level:0
\[
\qquad \text{octa}:\ \cos(m\pi/3)\in\{-1,\tfrac12\}:\ D\in\{4,1\}
\]

%dsp-level:0
\[
\text{dodeca}:\ T_0
\]

%dsp-level:0
\[
= 1,
\]

%dsp-level:0
\[
\ T_1
\]

%dsp-level:0
\[
= c,
\]

%dsp-level:0
\[
\ T_{m+1}
\]

%dsp-level:0
\[
= 2cT_m-T_{m-1}\ \text{in}\ \mathbb Q[c]/(16c^4-20c^2+5):\ T_5
\]

%dsp-level:0
\[
= c(16c^4-20c^2+5)
\]

%dsp-level:0
\[
= 0\ [\text{machine}];
\]

%dsp-level:0
\[
\qquad \text{icosa}:\ \cos(\pi/6)=\tfrac{\sqrt3}2:\ D(\pm1)=2-\sqrt3,\ D(\pm3)=2;\ \text{norm}\ (2-\sqrt3)(2+\sqrt3)=4-3=1
\]

%dsp-level:0
\[
(3):\ \mathrm{FT.137}(3):\ \Theta\in2\mathbb Z\ \text{and real weights}\implies w(P\mathcal C)=w(\mathcal C);
\]

%dsp-level:0
\[
\qquad \text{off-shell}:\ |e^{i\pi\Theta}|=1;
\]

%dsp-level:0
\[
\qquad \mathcal P_M\ \text{reads}\ (R,\ \text{seals}),\ \neg\alpha\ [\mathrm{CT2.3},\ \mathrm{D2.11}]
\]

%dsp-level:0
\[
(4):\ 0\le D\le4\iff0\le4\sin^2\le4\ [\text{range of}\ \sin^2];\qquad\text{the ceiling is a function of}\ D\ \text{values alone}.\ \blacksquare
\]

\textbf{Phys.} The pole ladder is cut from geometry and combinatorics alone — windings and vertex angles — and no register reversal in the corpus touches either. Parity walks through the whole pole spectrum without a footprint; the only motion is $m\mapsto-m$ under the ledger reflection, and every mass in every Platonic multiplet is even in it. The bounded-hierarchy verdicts of the SM audit are parity statements before the name existed: they constrain ranges, and ranges are what reversals preserve.

\vspace{2pt}\noindent\rule{\textwidth}{0.4pt}\vspace{2pt}

\end{theorem}

\begin{theorem}[{\texorpdfstring{FT.139 — The channel census: Koide and colour channels under the reversal family}{FT.139 - The channel census: Koide and colour channels under the reversal family}}]\label{tag:FT-139}
\textbf{Thm.}

%dsp-level:0
\[
(1)\quad\textbf{the seal bilinear is}\ S\textbf{-even:}
\]

%dsp-level:0
\[
\qquad \operatorname{Ad}(\mathbb S)[J_i,K_j]=[J_i,K_j]\ \ \forall(i,j)\ [\text{machine: all}\ 9\ \text{pairs}];
\]

%dsp-level:0
\[
\qquad \sum_i[J_i,K_i]=2i\Sigma\ [\mathrm{FT.109}\ \text{re-derived}];
\]

%dsp-level:0
\[
\qquad 2i\Sigma\in\mathcal E_{++}\ [\mathrm{FT.109}(4)];
\]

%dsp-level:0
\[
(2)\quad\textbf{the channel constants are register-blind:}\ \ \text{the seal coefficient}\ 2\ [\mathrm{FT.109}(3)],
\]

%dsp-level:0
\[
\ \dim F_{J^\pm}
\]

%dsp-level:0
\[
= 3\ [\mathrm{CT0.12}],
\]

%dsp-level:0
\[
\ \operatorname{tr}T(1)
\]

%dsp-level:0
\[
= 3\ [\mathrm{FT.68}(2):\ \text{the cap-3 transfer rule reads}\ (R,\ \text{vertex degrees})],
\]

%dsp-level:0
\[
\ N_{\text{colour}}
\]

%dsp-level:0
\[
= 3\ [\mathrm{CT3.20}:\ \text{the cap count}]:
\]

%dsp-level:0
\[
\text{no configuration register}\ (\sigma,\alpha,\kappa,\text{ledger})\ \text{enters any constituent}\implies\delta_\ell
\]

%dsp-level:0
\[
= \tfrac29\ [\mathrm{FT.112}],
\]

%dsp-level:0
\[
\ \delta_q
\]

%dsp-level:0
\[
= \tfrac2{27}\ [\mathrm{FT.119}],
\]

%dsp-level:0
\[
\ (A,B)
\]

%dsp-level:0
\[
= (1,\sqrt2)\ [\mathrm{FT.111}(3)],
\]

%dsp-level:0
\[
\ \text{the spectra}\ r_i,
\]

%dsp-level:0
\[
\ Q
\]

%dsp-level:0
\[
= \tfrac23,
\]

%dsp-level:0
\[
\ \text{and all ratios are invariant under the full}\ \mathbb Z_2^4\ \text{family}\ [\mathrm{FT.133}(2)];
\]

%dsp-level:0
\[
(3)\quad\textbf{the contrast census:}\ \ \text{Tier F.40's flipping block is not invariant:}\ F_{LR}
\]

%dsp-level:0
\[
= q\,
\]

%dsp-level:0
\[
I_2\ \text{is}\ C\text{-odd},
\]

%dsp-level:0
\[
\ P\text{-even},
\]

%dsp-level:0
\[
\ T\text{-even}\ [\mathrm{FT.134}(2);\ \mathrm{FT.133}(6)]:
\]

%dsp-level:0
\[
\text{the channels are parity-even and register-blind; the flipping block is charge-odd — two exact, distinct parity characters in the same tower}.
\]

\textbf{Req.} FT.109, FT.111, FT.112 (Tier F.37); FT.116, FT.119 (Tier F.38); FT.129 (Tier F.39); FT.68 (Tiers F.23–F.25); CT0.12 (Part 1); CT3.20 (Part 2); FT.133, FT.134 (this tier); exact $\mathbb Q(i)$ conjugation census [math; machine: \texttt{checks/\allowbreak{}verify\_\allowbreak{}H\_\allowbreak{}parity\_\allowbreak{}identification.py} §G].

\textbf{Pf.}

%dsp-level:0
\[
(1):\ \operatorname{Ad}(\mathbb S)J_i
\]

%dsp-level:0
\[
= -J_i,
\]

%dsp-level:0
\[
\ \operatorname{Ad}(\mathbb S)K_i
\]

%dsp-level:0
\[
= -K_i\ [\mathrm{CT0.7}(3)]:\ \operatorname{Ad}(\mathbb S)[J_i,K_j]
\]

%dsp-level:0
\[
= [-J_i,-K_j]
\]

%dsp-level:0
\[
= [J_i,K_j]\ [\text{machine, all}\ 9];
\]

%dsp-level:0
\[
\qquad \sum_i[J_i,K_i]=2i\Sigma\ [\text{entrywise re-derivation};\ \text{machine}]
\]

%dsp-level:0
\[
\operatorname{Ad}(\mathbb S)\Sigma=\Sigma\ [\mathrm{CT1.5}:\ [\Sigma,\mathbb S]=0]
\]

%dsp-level:0
\[
(2):\ \text{constituent census: each integer's definition cites only cap counts, block dimensions and transfer rules}\ [\mathrm{FT.109}(3),\ \mathrm{CT0.12},\ \mathrm{FT.68}(2),\ \mathrm{CT3.20}];
\]

%dsp-level:0
\[
\qquad \text{invariance of constants under any data map}
\]

%dsp-level:0
\[
(3):\ \mathrm{FT.134}(2)\ \text{characters verbatim};\qquad q_{\mathrm{core}}(C\mathcal C)=-q_{\mathrm{core}}(\mathcal C)\ [\mathrm{FT.133}(6)].\ \blacksquare
\]

\textbf{Phys.} The numbers that carry the seal-side and colour-side hierarchies are cut from counts no reversal can reach — caps, dimensions, traces — so both channels sit outside the entire parity calculus, exactly invariant under all four involutions at once. Where parity does bite is the flip block: charge-odd, parity-even. A future dynamical read of the flipping frequency will inherit that character as a selection rule, and any application-layer assignment rule will have to reconcile it with the branch-selective weak force of FT.135.

\vspace{2pt}\noindent\rule{\textwidth}{0.4pt}\vspace{2pt}

\end{theorem}

\phantomsection
\section[Tier F.42 — The Rung-One Chirality, the Janus Arrow, and…]{\texorpdfstring{Tier F.42 — The Rung-One Chirality, the Janus Arrow, and the Family Channel of the Sixteen-Dimensional Arena}{Tier F.42 - The Rung-One Chirality, the Janus Arrow, and the Family Channel of the Sixteen-Dimensional Arena}}

\begin{theorem}[{\texorpdfstring{FT.141 — The rung-one chirality theorem: the $10/6$ grading is the chirality of the complete force complex}{FT.141 - The rung-one chirality theorem: the grading is the chirality of the complete force complex}}]\label{tag:FT-141}
\textbf{Thm.} With $\mathbb S=\tfrac12(E_0+\Sigma)$ and $S(X)=\mathbb SX\mathbb S=\operatorname{Ad}(\mathbb S)X$ [A3; CT0.7(1)–(2)] on the pair channel $\mathfrak H(\mathbb C^4)$:

%dsp-level:0
\[
(1)\quad\textbf{the grading census}\ [\mathrm{FT.136}(1)\ \text{re-derived at the point of use; machine: all}\ 17\ \mathrm{CT1.2}\ \text{generators}]:
\]

%dsp-level:0
\[
\qquad S\text{-even}\ (+1):\ \{E_0,\ \Sigma,\ \hat E^{ij}\ (5),\ W_i\ (3)\}=10\ \text{dims}=G_5+W_3+H_1+D^+_1;
\]

%dsp-level:0
\[
S\text{-odd}\ (-1):\ \{J_i\ (3),\ K_i\ (3)\}=6\ \text{dims}=J^+_3+J^-_3;
\]

%dsp-level:0
\[
\qquad 10+6=16=C_{\mathrm{cap}}\ [\mathrm{CT0.12};\ \mathrm{CT0.6}];
\]

%dsp-level:0
\[
\qquad \text{the asymmetry}\ 10\ne6\ \text{is exact};
\]

%dsp-level:0
\[
(2)\quad\textbf{the chirality of the complete force complex:}\ \ \text{the A8 block labels are readings of the}\ (S,T)\ \text{grades}\ [\mathrm{CT0.7}(4)]:\ \mathbf F_{(+,\mathcal R)}
\]

%dsp-level:0
\[
= F_G\oplus F_H\oplus F_{D^+}\ \text{and}\ \mathbf F_{(-,\mathcal R)}
\]

%dsp-level:0
\[
= F_W\ \text{are the}\ S\text{-even aspects},
\]

%dsp-level:0
\[
\ \mathbf F_{(+,\mathcal V)}
\]

%dsp-level:0
\[
= F_{J^+},
\]

%dsp-level:0
\[
\ \mathbf F_{(-,\mathcal V)}
\]

%dsp-level:0
\[
= F_{J^-}\ \text{the}\ S\text{-odd aspects};
\]

%dsp-level:0
\[
\text{the chiral structure of the complete force operator}\ \mathbf F\ [\mathrm{A8}:\ \text{always applied in full}]\ \text{at rung}\ 1\ \text{is the ordered pair}\ \bigl(\dim\mathcal K_1^{+},\ \dim\mathcal K_1^{-}\bigr)
\]

%dsp-level:0
\[
= (10,\ 6)\ [\mathrm{CT0.14}:\ \mathcal K_1=\text{the rung-one mask}]:
\]

%dsp-level:0
\[
\text{a property of the pair algebra}\ [\mathrm{A3}'\text{s swap};\ \mathrm{CT1.1},\ \mathrm{CT1.2}'\text{s basis}],
\]

%dsp-level:0
\[
\ \neg\ \text{of any single block:}\ \neg\exists\ \text{an aspect that is}\ \textbf{the}\ \text{chiral force};
\]

%dsp-level:0
\[
(3)\quad\textbf{deletion invariance:}\ \ \forall\,
\]

%dsp-level:0
\[
\mathcal Y\subseteq\{G,\ W,\ H,\ D^+,\ J^+,\ J^-\}\ \text{(the six rung-one blocks)}:\ \text{the}\ S\text{-eigenvalue of every surviving generator is unchanged}\ (\text{the eigen-relation}\ \mathbb SX\mathbb S=\pm X\ \text{is a two-operator relation: deletion touches neither}\ \mathbb S\ \text{nor}\ X);
\]

%dsp-level:0
\[
\text{the surviving split is}\ (10-d^+_{\mathcal Y},\ 6-d^-_{\mathcal Y}),\ d^{\pm}_{\mathcal Y}\ \text{the deleted dimensions in each grade};
\]

%dsp-level:0
\[
\qquad \text{the grading is nontrivial (both grades occupied)}\iff
\]

%dsp-level:0
\[
\ge 1\ S\text{-even block and}\
\]

%dsp-level:0
\[
\ge 1\ S\text{-odd block survive}:\ 45\ \text{of the}\ 2^6
\]

%dsp-level:0
\[
= 64\ \text{subsets}\ [\text{machine: exhaustive}];
\]

%dsp-level:0
\[
\text{the six single-block deletions leave}\ (5,6),
\]

%dsp-level:0
\[
\ (7,6),
\]

%dsp-level:0
\[
\ (9,6),
\]

%dsp-level:0
\[
\ (9,6),
\]

%dsp-level:0
\[
\ (10,3),
\]

%dsp-level:0
\[
\ (10,3):\ \text{each nontrivial, each asymmetric};
\]

%dsp-level:0
\[
\qquad \text{the full}\ 10/6\ \text{is recovered}\iff\mathcal Y=\varnothing;
\]

%dsp-level:0
\[
(4)\quad\textbf{the weak shadow:}\ \ W_i\in\mathcal E_{+-}:\ S\text{-even}\ [\mathrm{CT1.2}\ \text{table}];
\]

%dsp-level:0
\[
\qquad \operatorname{span}\{J_i,K_i\}\cap\operatorname{span}\{W_k\}
\]

%dsp-level:0
\[
= \{0\}\ [\text{disjoint basis subsets};\ \text{machine: rank}\ 9=6+3]:\ \text{the}\ S\text{-odd block carries no weak content};
\]

%dsp-level:0
\[
\text{the branch selectivity of the weak}\ [\mathrm{FT.93}(5),\ \mathrm{FT.101}]\ \text{sits at the state-space locus of the same swap}:
\]

%dsp-level:0
\[
\qquad \Pi_{\mathcal R}=\tfrac12(E_0+\mathbb S),\quad\Pi_{\mathcal V}=\tfrac12(E_0-\mathbb S)\quad[\mathrm{CT0.7}(1)],
\]

%dsp-level:0
\[
W_i\Pi_{\mathcal V}=0\ \ \forall i\ [\mathrm{FT.101}(2)\ \text{re-derived}];
\]

%dsp-level:0
\[
\qquad \textstyle\sum_iW_i^2=\tfrac32E_0+\tfrac12\Sigma\ [\mathrm{Cor}\ \mathrm{CT0.7a}]=E_0+\mathbb S=2\,\Pi_{\mathcal R}
\]

%dsp-level:0
\[
\qquad [\Sigma=2\mathbb S-E_0;\ \text{machine, entrywise}]:
\]

%dsp-level:0
\[
\text{the swap operator itself enters the weak Casimir identity — the branch asymmetry}\ (2\ \text{vs}\ 0)\ \text{is derived from the}\ S\text{-grading, and}\ \mathbb S\ \text{is A3 data, prior to every force assignment}\ [\mathrm{CT0.7}(4):\ \text{the A8 labels read the grades}];
\]

%dsp-level:0
\[
\text{locus discipline}\ [\mathrm{FT.136}(5)\ \text{stands}]:\ \text{the}\ 3\oplus1\ \text{branch split}\ (\operatorname{spec}\mathbb S\ \text{on}\ \mathbb C^4_{\mathrm{pair}})\ \text{and the}\ 10/6\ \text{operator split}\ (\operatorname{spec}\operatorname{Ad}(\mathbb S)\ \text{on}\ \operatorname{End}(\mathbb C^4_{\mathrm{pair}}))\ \text{are two loci of one swap, and neither is the carrier}\ \hat\chi\text{-grade}\ [\mathrm{FT.122}];
\]

%dsp-level:0
\[
(5)\quad\textbf{no}\ S\textbf{-neutral observable at rung one:}\ \ \operatorname{Ad}(\mathbb S)^2
\]

%dsp-level:0
\[
= \mathrm{id}\implies\operatorname{spec}\operatorname{Ad}(\mathbb S)\subseteq\{\pm1\}\ \text{and}\ \ker\operatorname{Ad}(\mathbb S)
\]

%dsp-level:0
\[
= \{0\}\ [\operatorname{Ad}(\mathbb S)\ \text{is its own inverse};\ \text{machine: rank}\ 16\ \text{on the flattened}\ 16\text{-basis};\ \operatorname{tr}\operatorname{Ad}(\mathbb S)=10-6=4];
\]

%dsp-level:0
\[
\text{every rung-one observable carries the unique grade decomposition}\ X=X_++X_-,
\]

%dsp-level:0
\[
\quad X_{\pm}=\tfrac12\bigl(\mathrm{id}\pm\operatorname{Ad}(\mathbb S)\bigr)X;
\]

%dsp-level:0
\[
\text{the metric reads}\ (\hat h_{ij}\ 5,\ \Sigma_{\mathrm{tr}}\ 1,\ a_{00}\ 1)\ \text{and the weak reads}\ (w_i\ 3)\ \text{are}\ S\text{-even};
\]

%dsp-level:0
\[
\ \text{the seal reads}\ (J^{\pm})\ \text{are}\ S\text{-odd with}\ 0\ \text{metric coordinates}\ [\mathrm{FT.30}\ \text{dictionary}]:
\]

%dsp-level:0
\[
\text{chirality at rung}\ 1\ \text{is}\ \neg\ \text{a selection rule applied by one force — it is the grading of the entire observable algebra}\ \mathcal K_1\ [\mathrm{CT0.14}\ \mathrm{Rmk}:\ \dim\mathcal K_1=16=C_{\mathrm{cap}}:\ \text{rung one exhausts the pair channel}];
\]

%dsp-level:0
\[
(6)\quad\textbf{reconciliation — the two independent}\ \mathbb Z_2\textbf{'s:}\ \ \text{the force-arena}\ S\text{-grading acts on}\ \operatorname{End}(\mathbb C^4_{\mathrm{pair}});
\]

%dsp-level:0
\[
\ \text{the carrier chirality}\ \hat\chi\ \text{acts on}\ \mathcal H_\omega\ [\mathrm{Def\text{-}H.33},\ \mathrm{FT.122};\ \text{Def-H.34}]:\ \text{pairwise distinct, commuting where jointly defined, neither reducible to the other}\ [\mathrm{FT.126}(4),\ \mathrm{FT.136}(4)\ \text{re-derived};\ \text{machine witnesses}];
\]

%dsp-level:0
\[
\text{the division of labour: the}\ S\text{-grading determines WHICH observables are chiral (the grade of every read); the carrier chirality determines WHICH chain orientations are chiral (the doublet's rays);}
\]

%dsp-level:0
\[
\text{the rule coupling the two}\ \mathbb Z_2\text{'s — and the branch-label}\ \mathbb Z_2\ [\mathrm{FT.101},\ \mathrm{FT.135}]\ \text{and the Janus arrow}\ [\mathrm{FT.144}(5)] — \text{to species is the deficit}\ 1'\ \text{assignment rule}\ [\mathrm{FT.125}(4)]:\ \neg\text{constructed here};
\]

%dsp-level:0
\[
\ \text{no new declaration issues}.
\]

\textbf{Req.} A3, A8 (Parts 1–2); CT0.6, CT0.7, Cor CT0.7a, CT0.12, CT0.13, CT0.14, CT1.1, CT1.2, CT1.3, CT1.5 (Part 1); FT.30 (Tiers F.1–F.11); FT.93, FT.101 (Tiers F.31–F.34); Def-H.33, FT.122, FT.125, FT.126 (Tier F.39); FT.135, FT.136 (Tier F.41); Def-H.34 (carrier clause); exact $\mathbb Q(i)$ matrix arithmetic, exhaustive $2^6$ subset census, rank/trace invariants, inclusion–exclusion [math; machine: \texttt{checks/\allowbreak{}verify\_\allowbreak{}H\_\allowbreak{}rung1\_\allowbreak{}neutrino.py} §A].

\textbf{Pf.}

%dsp-level:0
\[
(1):\ \mathbb S=\tfrac12(E_0+\Sigma),\ \mathbb S^2=E_0\ [\mathrm{CT0.7}(1);\ \Sigma^2=3E_0-2\Sigma];
\]

%dsp-level:0
\[
\qquad \mathbb SE_0\mathbb S=E_0,\ \mathbb S\Sigma\mathbb S=\Sigma\ [\mathrm{CT1.3},\ \mathrm{CT1.5}];
\]

%dsp-level:0
\[
\qquad \mathbb S\hat E^{ij}\mathbb S=\hat E^{ij},\ \mathbb SW_i\mathbb S=W_i;
\]

%dsp-level:0
\[
\qquad \mathbb SJ_i\mathbb S=-J_i,\ \mathbb SK_i\mathbb S=-K_i\ [\mathrm{CT0.7}(3);\ \mathrm{CT1.2}\ \text{table}]
\]

%dsp-level:0
\[
\text{entrywise on all}\ 17\ \text{generators over}\ \mathbb Q(i)\ [\text{machine}];
\]

%dsp-level:0
\[
\qquad \text{counts}:\ 1+1+5+3=10;\quad3+3=6;
\]

%dsp-level:0
\[
\qquad 10+6=16=C_{\mathrm{cap}}\ [\mathrm{CT0.12};\ \mathrm{CT0.6}]
\]

%dsp-level:0
\[
(2):\ \text{block}\leftrightarrow\text{aspect rows verbatim}\ [\mathrm{CT0.12}\ \text{table};\ \mathrm{CT0.7}(4)];
\]

%dsp-level:0
\[
\qquad \text{the}\ (S,T)\ \text{refinement}\ (7,3,3,3)\ [\mathrm{CT1.1}]:\ S\text{-even}=7+3=10,\ S\text{-odd}=3+3=6\ [\text{exact integers}];
\]

%dsp-level:0
\[
\qquad \text{A8}: \mathbf F\ \text{is the full direct sum — no sub-block restriction is an A8 force}
\]

%dsp-level:0
\[
(3):\ \text{the eigen-relation is binary in}\ (\mathbb S,X):\ \text{deleting}\ Y
\]

%dsp-level:0
\[
\ne X\ \text{leaves it untouched}\ [\text{machine: the grade table is identical on all}\ 64\ \text{subsets}];
\]

%dsp-level:0
\[
\qquad \text{both-grade census}:\ 2^6-2^4-2^2+1
\]

%dsp-level:0
\[
= 64-16-4+1
\]

%dsp-level:0
\[
= 45\ [\text{inclusion–exclusion over the}\ 4\ \text{even and}\ 2\ \text{odd blocks};\ \text{machine}]
\]

%dsp-level:0
\[
\text{single deletions}\ (d^+,d^-)\in\{(5,0),(3,0),(1,0),(1,0),(0,3),(0,3)\}
\]

%dsp-level:0
\[
\implies (5,6),
\]

%dsp-level:0
\[
(7,6),
\]

%dsp-level:0
\[
(9,6),
\]

%dsp-level:0
\[
(9,6),
\]

%dsp-level:0
\[
(10,3),
\]

%dsp-level:0
\[
(10,3)\ [\text{machine}];
\]

%dsp-level:0
\[
\qquad (10,6)\ \text{uniquely at}\ \mathcal Y=\varnothing\ [\text{machine}]
\]

%dsp-level:0
\[
(4):\ W_i\in\mathcal E_{+-}\ [\mathrm{CT1.2}];
\]

%dsp-level:0
\[
\qquad \operatorname{rank}\{J_1,J_2,J_3,K_1,K_2,K_3,W_1,W_2,W_3\}
\]

%dsp-level:0
\[
= 9
\]

%dsp-level:0
\[
= 6+3\ [\text{exact flattening rank};\ \text{machine}]\implies\text{the intersection is}\ \{0\}
\]

%dsp-level:0
\[
\Pi_{\mathcal R},
\]

%dsp-level:0
\[
\Pi_{\mathcal V}:\ \text{idempotent, orthogonal, complete},
\]

%dsp-level:0
\[
\ \mathfrak T\Pi_{\mathcal R}
\]

%dsp-level:0
\[
= 3,
\]

%dsp-level:0
\[
\ \mathfrak T\Pi_{\mathcal V}
\]

%dsp-level:0
\[
= 1\ [\mathrm{CT0.7}(1);\ \text{machine}];
\]

%dsp-level:0
\[
\qquad W_i\Pi_{\mathcal V}=0\ \forall i\ [\mathrm{FT.101}(2)\ \text{re-derived entrywise};\ \text{machine}]
\]

%dsp-level:0
\[
\textstyle\sum_iW_i^2=\tfrac14\sum_i(\sigma_i\otimes\sigma_0+\sigma_0\otimes\sigma_i)^2=\tfrac32E_0+\tfrac12\Sigma\ [\mathrm{Cor}\ \mathrm{CT0.7a}\ \text{re-derived}];
\]

%dsp-level:0
\[
\qquad \Sigma
\]

%dsp-level:0
\[
= 2\mathbb S-E_0\ [\mathrm{CT0.7}(1)]\implies\tfrac32E_0+\tfrac12(2\mathbb S-E_0)
\]

%dsp-level:0
\[
= E_0+\mathbb S
\]

%dsp-level:0
\[
= 2\cdot\tfrac12(E_0+\mathbb S)
\]

%dsp-level:0
\[
= 2\Pi_{\mathcal R}\ [\text{machine}]
\]

%dsp-level:0
\[
\text{locus rows}:\ \operatorname{spec}\mathbb S\ \text{on}\ \mathbb C^4_{\mathrm{pair}}
\]

%dsp-level:0
\[
= (3,1)\ [\mathfrak T\Pi_{\mathcal R}=3,\ \mathfrak T\Pi_{\mathcal V}=1];
\]

%dsp-level:0
\[
\quad\operatorname{spec}\operatorname{Ad}(\mathbb S)\ \text{on}\ \operatorname{End}(\mathbb C^4_{\mathrm{pair}})
\]

%dsp-level:0
\[
= (10,6)\ [(1)];
\]

%dsp-level:0
\[
\quad\hat\chi\ \text{on}\ \mathcal H_\omega\ [\mathrm{FT.136}(4)\ \text{re-run}];
\]

%dsp-level:0
\[
\qquad \text{three distinct spaces}
\]

%dsp-level:0
\[
(5):\ \operatorname{Ad}(\mathbb S)^2X=\mathbb S(\mathbb SX\mathbb S)\mathbb S=\mathbb S^2X\mathbb S^2=E_0XE_0=X\ [\mathrm{CT1.3}];
\]

%dsp-level:0
\[
\qquad \text{self-inverse}\implies\text{invertible}\implies\ker=\{0\};
\]

%dsp-level:0
\[
\qquad \text{rank}\ 16\ \text{on the flattened basis}\ [\text{machine}];
\]

%dsp-level:0
\[
\qquad \operatorname{tr}\operatorname{Ad}(\mathbb S)=\#\{+1\}-\#\{-1\}=10-6=4\ [\text{machine}]
\]

%dsp-level:0
\[
X_{\pm}=\tfrac12(X\pm\mathbb SX\mathbb S):\ \operatorname{Ad}(\mathbb S)X_{\pm}=\pm X_{\pm},\ X_++X_-=X\ [\mathbb S^2=E_0];
\]

%dsp-level:0
\[
\qquad \text{uniqueness: eigenspace decomposition of an involution}\ [\text{machine on the}\ 16\text{-basis}];
\]

%dsp-level:0
\[
\qquad \mathrm{FT.30}\ \text{dictionary rows verbatim}
\]

%dsp-level:0
\[
(6):\ [G\otimes I_2,\ I_2\otimes\hat\omega]=0,
\]

%dsp-level:0
\[
\quad (G\otimes\hat\omega)^2=\mathrm{id},
\]

%dsp-level:0
\[
\quad G\otimes\hat\omega\ne G\otimes I_2\ne I_2\otimes\hat\omega\ [\mathrm{FT.126}(4)/\mathrm{FT.134}(5)\ \text{pattern};\ \text{machine on}\ G\in\{\sigma_1,\sigma_3\}];
\]

%dsp-level:0
\[
\qquad \operatorname{Ad}(\mathbb S)\ \text{on}\ \operatorname{End}(\mathbb C^4_{\mathrm{pair}})\ \text{vs}\ \hat\chi\ \text{on}\ \mathcal H_\omega:\ \text{distinct carriers}\ [\mathrm{FT.136}(4)]
\]

%dsp-level:0
\[
\text{coupling rule}:\ \neg\exists\ \text{tag assigning grades to species}\ [\text{inventory scan};\ \mathrm{FT.125}(4)].\ \blacksquare
\]

\textbf{Phys.} At the first rung the entire force complex is chiral, and its chirality is a count: ten dimensions speak to the metric — shear, volume, conformal, shift — and six only move the ledger. The count belongs to no force: it is the swap grading of the pair algebra, and every aspect of the complete force sits inside it with its grade already decided. The weak force's one-sidedness is a shadow of the same swap — its Casimir identity contains the swap operator literally — so no isolation of the weak reproduces the asymmetry, and no deletion of any force block removes it: the grade of each generator is a two-operator fact, invariant under everything except deleting the generator itself. Every observable at rung one is graded; none is neutral. The carrier keeps a second, independent chirality of its own, and the rule that would marry the two gradings to species is application-layer content.

\vspace{2pt}\noindent\rule{\textwidth}{0.4pt}\vspace{2pt}

\end{theorem}

\begin{theorem}[{\texorpdfstring{FT.142 — The rung-one arena and the channel angle $\delta_\nu=\tfrac18$}{FT.142 - The rung-one arena and the channel angle}}]\label{tag:FT-142}
\textbf{Thm.} The rung-one channel — the deficit-6 ledger row, internal channel subscript $\nu$ [$\neg$a species name; Rmk below]:

%dsp-level:0
\[
(1)\quad\textbf{the channel class:}\ \ \text{a rung-one bound chain is a bound chain}\ [\mathrm{D1.5}]\ \text{whose terminus-active aspects lie in}\ \mathcal A(1)
\]

%dsp-level:0
\[
= \{J,\ G,\ W,\ H,\ D^+\}\ [\mathrm{CT0.13}]:
\]

%dsp-level:0
\[
\text{EM unactivated}\ (f_v=0\ \text{reads}\ [\mathrm{D3.3};\ \mathrm{FT.28}(3):\ \mathrm{EM}\in\mathcal A\iff f_p\ne0\ \vee\ f_q\ne0]);
\]

%dsp-level:0
\[
\qquad \text{colour block vanishing}\ (\text{no slot index below trivalence}\ [\mathrm{CT3.11}(1)];\ \text{single-contact floors:}\ \mathcal C=0\implies\text{block}\equiv0\ [\mathrm{FT.19},\ \mathrm{FT.28}(3)]);
\]

%dsp-level:0
\[
\text{the Def-H.4 floor stands:}\ \delta h(c)\ge2,\ n(v)\ge2\ [\text{self-loops}\ \bot\ \text{by A1}];
\]

%dsp-level:0
\[
\qquad \text{the minimal instance is the digon}\ \Sigma_{\min}\ [\mathrm{FT.28}];
\]

%dsp-level:0
\[
(2)\quad\textbf{the arena:}\ \ \text{the pair-channel arena on which the complete rung-one force acts is the rung-one mask}\ \mathcal K_1:
\]

%dsp-level:0
\[
\qquad \dim\mathcal K_1=16=C_{\mathrm{cap}}\quad[\mathrm{CT0.14},\ \mathrm{CT0.12};\ \text{machine: exact rank}];
\]

%dsp-level:0
\[
\text{the complete-force standing rule}\ [\mathrm{A8}]\ \text{fixes the arena as the full mask: no sub-block arena — isolating a force is not an A8 action};
\]

%dsp-level:0
\[
(3)\quad\textbf{the arena adjudication}\ (\text{computed candidates; one adopted}):
\]

%dsp-level:0
\[
\begin{array}{c|c|c}\text{candidate} & \text{source} & \text{verdict}\\ \hline 16=\dim\mathcal K_1 & \mathrm{CT0.14}\ \text{(derived mask)} & \textbf{adopted}\\ 10=S\text{-even block} & \mathrm{FT.141}(1) & \text{rejected: deleting the}\ S\text{-odd blocks violates A8's completeness}\ [\mathrm{CT0.13}:\ J\in\mathcal A(1)]\\ 9=\text{seal}\times\text{transfer} & \mathrm{FT.110}\ [\text{Def-H.41}] & \text{rejected: a interface definition for transfer-dressed chains}\ [\text{the transfer trace}\ \operatorname{tr}T(1)=3\ \text{is polygon-channel data},\ \mathrm{FT.68}(2)];\ \text{not imported in this part}\\ 27=\text{seal}\times\text{transfer}\times\text{colour} & \mathrm{FT.118}(2)\ [\text{Def-H.42}] & \text{rejected: colour slots are rung-three content}\ [\mathrm{CT0.13}];\ \text{clause}\ (1)\ \text{excludes them};\ \text{not imported}\\ 6=\dim\mathcal K_0 & \mathrm{CT0.14} & \text{rejected: seals alone; a terminus with}\ n(v)\ge2\ \text{carries metric and weak reads}\ [\mathrm{FT.28}(3):\ \{G,H,W\}\subseteq\mathcal A(\Sigma_{\min})]\end{array}
\]

%dsp-level:0
\[
\text{the adopted arena is the unique candidate that is derived content with no interface declaration:}\ \delta_\nu\ \text{is the first channel angle whose denominator requires no declaration-entry construction};
\]

%dsp-level:0
\[
(4)\quad\textbf{the seal coefficient, undressed:}\ \ [J_i,K_j]
\]

%dsp-level:0
\[
= i\bigl(\delta_{ij}\Sigma-E^{ij}\bigr)\ [\mathrm{CT1.7}\ \text{re-derived, all}\ 9\ \text{pairs}];
\]

%dsp-level:0
\[
\qquad \textstyle\sum_i[J_i,K_i]=i\bigl(3\Sigma-\textstyle\sum_iE^{ii}\bigr)=2i\Sigma\quad[\textstyle\sum_iE^{ii}=\Sigma];
\]

%dsp-level:0
\[
\text{in the}\ 16\text{-element CT1.2 basis the sum has exactly one nonzero component — the}\ \Sigma\text{-component, coefficient exactly}\ 2\ [\mathrm{FT.109}(3)\ \text{pattern};\ \text{machine: exact linear solve, rank}\ 16];
\]

%dsp-level:0
\[
\text{the rung-one channel class carries no dressing factor}\ (\neg\text{colour slots by clause}\ (1);\ \neg\text{transfer declaration imported by clause}\ (3)):\ \text{the numerator is}\ 2,
\]

%dsp-level:0
\[
\ \text{undressed}\ [\mathrm{FT.118}(5)\ \text{spectator pattern}];
\]

%dsp-level:0
\[
(5)\quad\textbf{the channel angle:}
\]

%dsp-level:0
\[
\qquad \delta_\nu\ :=\ \frac{\sum_i[J_i,K_i]\big|_\Sigma}{\dim\mathcal K_1}\ =\ \frac{2}{16}\ =\ \frac18
\]

%dsp-level:0
\[
\qquad [\text{exact rational};\ \mathrm{FT.112}\ \text{pattern}];
\]

%dsp-level:0
\[
\text{the democratic normalization}\ A
\]

%dsp-level:0
\[
= 1,
\]

%dsp-level:0
\[
\ B
\]

%dsp-level:0
\[
= \sqrt2\ \text{is seal-side data}\ [\mathrm{FT.111}(3):\ \textstyle\sum_i\lambda_i=3=\dim F_{J^\pm};\ \sum_i(\lambda_i-A)^2=3]\ \text{and hence channel-blind}\ [\mathrm{FT.119}\ \text{precedent}]:\ \text{the rung-one channel inherits it unchanged}.
\]

\textbf{Rmk (channel-label discipline).} The subscript $\nu$ names the deficit-6 ledger row — internal bookkeeping, joining $\ell$ [FT.112] and $q$ [FT.119] as a channel subscript; all three denote ledger rows, $\neg$particles. No species identification is made or implied [P1, P5]. The arena $\mathcal H_c$-vocabulary of FT.110/FT.118 (counting spaces without inner product, dynamics or states) applies verbatim: $\mathcal K_1$ enters this tier exactly once, as the denominator of clause (5).

\textbf{Req.} A1, A8, D0.8, D1.5, D3.3, P1, P4, P5 (Parts 1–2); CT0.6, CT0.12, CT0.13, CT0.14, CT1.2, CT1.7 (Part 1); CT3.11 (Part 2); Def-H.4 (Tier F.2); FT.19, FT.28 (Tiers F.1–F.11); FT.68(2) (Tiers F.23–F.25); FT.109, FT.110, FT.111, FT.112 (Tier F.37); FT.118 (Tier F.38); FT.141 (this tier); Def-H.41, Def-H.42 cited in the rejection rows and imported nowhere; exact $\mathbb Q(i)$ arithmetic, exact linear solve on the $16$-basis, exact flattening ranks [math; machine: \texttt{checks/\allowbreak{}verify\_\allowbreak{}H\_\allowbreak{}rung1\_\allowbreak{}neutrino.py} §B].

\textbf{Pf.}

%dsp-level:0
\[
(1):\ \mathrm{D1.5}:\ \text{bound}=\text{closed self-contained},\ \text{fermionic content};
\]

%dsp-level:0
\[
\qquad \mathcal A(1)\ \text{verbatim}\ [\mathrm{CT0.13}];
\]

%dsp-level:0
\[
\qquad \mathrm{FT.28}(3)\ \text{rows verbatim}\ [\mathrm{D3.3};\ \mathrm{CT3.11}(1);\ \mathrm{FT.19}]
\]

%dsp-level:0
\[
\mathrm{Def\text{-}H.4}\ \text{Rmk}:\ \delta h=1\ \text{closed chain}=\text{self-loop}=\bot\ [\mathrm{A1}];
\]

%dsp-level:0
\[
\qquad \text{closed chain based at}\ v\implies n(v)\ge2\ [\mathrm{D0.8}];
\]

%dsp-level:0
\[
\qquad \Sigma_{\min}=\text{the}\ 2\text{-cycle floor}\ [\mathrm{FT.28}]
\]

%dsp-level:0
\[
(2):\ \dim\mathcal K_1=6+(5+3+1+1)=16\ [\mathrm{CT0.14};\ \mathrm{CT0.12}\ \text{table}];
\]

%dsp-level:0
\[
\qquad \text{machine: flattening rank of the}\ 16\ \mathrm{CT1.2}\ \text{elements}=16\ \text{over}\ \mathbb Q\ (32\ \text{real components});
\]

%dsp-level:0
\[
\qquad 16=C_{\mathrm{cap}}\ [\mathrm{CT0.6}]
\]

%dsp-level:0
\[
\text{A8}: \mathbf F
\]

%dsp-level:0
\[
= F_{(+,\mathcal R)}\oplus F_{(-,\mathcal R)}\oplus F_{(+,\mathcal V)}\oplus F_{(-,\mathcal V)}\oplus F_{\mathrm{st}}\ \text{always applies in full}\implies\text{the rung-one read arena is the full mask}\ \mathcal K_1\ [\mathrm{CT0.14}:\ \mathcal P_M^{(1)}\ \text{projects onto it}]
\]

%dsp-level:0
\[
(3):\ \text{candidate rows}:\ 10=\operatorname{rank}\{E_0,\Sigma,\hat E^{ij},W_i\}\ [\text{machine}];
\]

%dsp-level:0
\[
\qquad 9=3\cdot3\ [\mathrm{FT.110}(1)\ \text{under Def-H.41's declared tensor product}];
\]

%dsp-level:0
\[
\qquad 27=3\cdot3\cdot3\ [\mathrm{FT.118}(2)\ \text{under Def-H.42}];
\]

%dsp-level:0
\[
\qquad 6=\operatorname{rank}\{J_i,K_i\}\ [\mathrm{CT0.14};\ \text{machine}]
\]

%dsp-level:0
\[
\text{rejection grounds}:\ \text{the}\ S\text{-even sub-arena omits}\ J\in\mathcal A(1)\ [\mathrm{CT0.13}]\ \text{against A8};
\]

%dsp-level:0
\[
\qquad \text{the}\ 9/27\ \text{arenas are declared constructions}\ [\mathrm{FT.110}\ \mathrm{Rmk};\ \mathrm{FT.118}\ \mathrm{Rmk}]\ \text{and this part's import discipline admits}\ \{ \text{Def-H.4},\ \text{Def-H.37(2)},\ \text{Def-H.38},\ \text{Def-H.34},\ \text{Def-H.43}\}\ \text{only};
\]

%dsp-level:0
\[
\qquad \mathcal K_0\ \text{omits the metric/weak reads of any}\ n\ge2\ \text{terminus}\ [\mathrm{FT.28}(3)]
\]

%dsp-level:0
\[
(4):\ \mathrm{CT1.7}\ \text{entrywise on all}\ 9\ \text{pairs over}\ \mathbb Q(i)\ [\text{machine}];
\]

%dsp-level:0
\[
\qquad \textstyle\sum_i(\Sigma-E^{ii})=3\Sigma-\Sigma=2\Sigma\ [\mathrm{FT.109}(2)\ \text{re-derived}];
\]

%dsp-level:0
\[
\qquad \text{component read: exact solve on the rank-}16\ \text{basis}\implies(0,\dots,0,\ 2i,\ 0,\dots,0)\ [\text{machine}]
\]

%dsp-level:0
\[
(5):\ \delta_\nu=\tfrac{2}{16}=\tfrac18\ [\text{exact Fractions}];
\]

%dsp-level:0
\[
\qquad A=1:\ \textstyle\sum_i\lambda_i=3A=3=\dim F_{J^\pm}\ [\mathrm{CT0.12}\ \text{table}];
\]

%dsp-level:0
\[
\qquad B^2=2:\ \textstyle\sum_i(\lambda_i-A)^2=B^2\textstyle\sum_i\cos^2(\theta+\tfrac{2\pi i}3)=B^2\cdot\tfrac32=3\ [\mathrm{FT.111}(3)\ \text{identities}].\ \blacksquare
\]

\textbf{Phys.} The rung-one channel is the class of bound chains that live on the pretime forces alone — no charge read firing, no colour slot gripped — and its arena needs no construction: the rung-one mask already is the whole pair channel, sixteen dimensions, derived. The seed integer of the family mechanism, the diagonal seal coefficient $2$, reappears undressed, because rung one dresses nothing. One integer over one derived dimension: the channel angle is $1/8$, the first of the three family angles whose denominator is a theorem rather than a interface definition.

\vspace{2pt}\noindent\rule{\textwidth}{0.4pt}\vspace{2pt}

\end{theorem}

\pseudohead{theorem}{\texorpdfstring{FT.143}{FT.143}}{\texorpdfstring{The rung-one family spectrum}{The rung-one family spectrum}}\label{tag:FT-143}
\textbf{Thm.} With $\delta_\nu=\tfrac18$ [FT.142(5)] and $A=1$, $B=\sqrt2$ [FT.111(3)]:

%dsp-level:0
\[
(1)\quad\textbf{the spectrum:}\qquad r_i^{(\nu)}=1+\sqrt2\cos\Bigl(\frac18+\frac{2\pi i}{3}\Bigr),\qquad i=0,1,2,
\]

%dsp-level:0
\[
\text{with exact radical-trigonometric closed forms}\ (c=\cos\tfrac18,\ s=\sin\tfrac18):
\]

%dsp-level:0
\[
\qquad r_0^{(\nu)}=1+\sqrt2\,c;
\]

%dsp-level:0
\[
\qquad r_1^{(\nu)}=1-\frac{\sqrt2\,c+\sqrt6\,s}{2};
\]

%dsp-level:0
\[
\qquad r_2^{(\nu)}=1-\frac{\sqrt2\,c-\sqrt6\,s}{2};
\]

%dsp-level:0
\[
r_0^{(\nu)}=2.403179397550\ldots,
\]

%dsp-level:0
\[
\qquad r_1^{(\nu)}=0.145715560918\ldots,
\]

%dsp-level:0
\[
\qquad r_2^{(\nu)}=0.451105041530\ldots
\]

%dsp-level:0
\[
\qquad [\text{all three strictly positive; disjoint enclosures; machine}];
\]

%dsp-level:0
\[
(2)\quad\textbf{Vieta and the channel-own cubic:}
\]

%dsp-level:0
\[
\qquad \textstyle\sum_ir_i^{(\nu)}=3,
\]

%dsp-level:0
\[
\qquad \sum_{i<j}r_i^{(\nu)}r_j^{(\nu)}=\frac32,
\]

%dsp-level:0
\[
\qquad \sum_i\bigl(r_i^{(\nu)}\bigr)^2=6
\]

%dsp-level:0
\[
\qquad [\text{exact ring identities};\ \text{machine}];
\]

%dsp-level:0
\[
2\bigl(r^{(\nu)}-1\bigr)^3-3\bigl(r^{(\nu)}-1\bigr)-\sqrt2\cos\frac38=0\quad\text{for each}\ r^{(\nu)}=r_i^{(\nu)}
\]

%dsp-level:0
\[
\qquad \bigl[\cos3\delta_\nu=\cos\tfrac38=4c^3-3c;\ \text{machine: symbolic + enclosure}\bigr];
\]

%dsp-level:0
\[
\cos\tfrac38\ne\cos\tfrac29\ [\text{disjoint enclosures}]:\ \text{the cubic constant}\ \sqrt2\cos\tfrac38\ \text{is channel-own — the rung-one channel shares no depressed cubic with the}\ \ell/q\ \text{pair}\ [\mathrm{FT.119}(4):\ \cos3\delta_q=\cos\delta_\ell=\cos\tfrac29];
\]

%dsp-level:0
\[
(3)\quad\textbf{the mass-squared ratios:}
\]

%dsp-level:0
\[
\qquad \bigl(r_1^{(\nu)}\bigr)^2<\bigl(r_2^{(\nu)}\bigr)^2<\bigl(r_0^{(\nu)}\bigr)^2\ [\text{enclosure-proven}];
\]

%dsp-level:0
\[
\qquad m_i=\mu_\nu\bigl(r_i^{(\nu)}\bigr)^2,\ \ \mu_\nu\ \text{undeclared}\ (\mathrm{P4});
\]

%dsp-level:0
\[
\qquad \text{the}\ m^2\ \text{vocabulary is the corpus's own inertial read}\ [\mathrm{L3}(4),\ \mathrm{L4}];\ \neg\text{a particle identification}:
\]

%dsp-level:0
\[
m_1^2:m_2^2:m_3^2
\]

%dsp-level:0
\[
= 1:\frac{\bigl(r_2^{(\nu)}\bigr)^2}{\bigl(r_1^{(\nu)}\bigr)^2}:\frac{\bigl(r_0^{(\nu)}\bigr)^2}{\bigl(r_1^{(\nu)}\bigr)^2}
\]

%dsp-level:0
\[
= 1:9.583926992379\ldots:271.994748749143\ldots
\]

%dsp-level:0
\[
\qquad [\text{widths}<2^{-50};\ \text{machine}];
\]

%dsp-level:0
\[
\textbf{anchored forms}\ (\text{scale-free, top-normalized}):
\]

%dsp-level:0
\[
\qquad \frac{m_2}{m_3}=\frac{\bigl(r_2^{(\nu)}\bigr)^2}{\bigl(r_0^{(\nu)}\bigr)^2}=0.035235705970\ldots,
\]

%dsp-level:0
\[
\qquad \frac{m_1}{m_3}=\frac{\bigl(r_1^{(\nu)}\bigr)^2}{\bigl(r_0^{(\nu)}\bigr)^2}=0.003676541567\ldots;
\]

%dsp-level:0
\[
(4)\quad\textbf{the Q-invariant:}
\]

%dsp-level:0
\[
\qquad Q\ :=\ \frac{\sum_im_i}{\bigl(\sum_i\sqrt{m_i}\bigr)^2}\ \overset{m_i=\mu_\nu(r_i^{(\nu)})^2}{=}\ \frac{\sum_i(r_i^{(\nu)})^2}{\bigl(\sum_ir_i^{(\nu)}\bigr)^2}\ =\ \frac69\ =\ \frac23
\]

%dsp-level:0
\[
\qquad [\text{exact, and independent of}\ \delta\ \text{and of}\ \mu\ [\mathrm{FT.119}(4)\ \text{pattern}]]:
\]

%dsp-level:0
\[
\text{the channel normalization}\ Q=\tfrac23\ \text{holds at rung one};
\]

%dsp-level:0
\[
(5)\quad\textbf{the angle ledger}\ (\text{exact rational arithmetic; no equality beyond those displayed}):
\]

%dsp-level:0
\[
16\,\delta_\nu=9\,\delta_\ell=27\,\delta_q=2
\]

%dsp-level:0
\[
\qquad [\text{the shared seal numerator is the only cross-channel identity};\ \mathrm{FT.112},\ \mathrm{FT.119}(1)];
\]

%dsp-level:0
\[
3\delta_\nu=\frac38\notin\Bigl\{\frac29,\ \frac2{27}\Bigr\}=\{\delta_\ell,\ \delta_q\};
\]

%dsp-level:0
\[
\qquad \delta_\nu\notin\Bigl\{3\delta_\ell,\ 3\delta_q,\ \frac{\delta_\ell}3,\ \frac{\delta_q}3\Bigr\}:\ \text{the triple-angle link}\ 3\delta_q
\]

%dsp-level:0
\[
= \delta_\ell\ [\mathrm{FT.119}(2)]\ \text{has no rung-one analogue};
\]

%dsp-level:0
\[
\delta_\nu-\delta_\ell=-\frac7{72};
\]

%dsp-level:0
\[
\qquad \delta_\nu-\delta_q=\frac{11}{216};
\]

%dsp-level:0
\[
\qquad 3\delta_\nu-\delta_\ell=\frac{11}{72}=3\,(\delta_\nu-\delta_q)\quad[\text{consistent with}\ \delta_\ell=3\delta_q];
\]

%dsp-level:0
\[
\delta_\nu:\delta_\ell:\delta_q=27:48:16\quad[\times216];\qquad\delta_q<\delta_\nu<\delta_\ell;
\]

%dsp-level:0
\[
(6)\quad\textbf{the three-channel table}\ (\text{frozen rows cited; rung-one row derived here}):
\]

%tbl-level:1
{\small\setlength{\tabcolsep}{4pt}
\begin{longtable}{>{\raggedright\arraybackslash}p{\dimexpr0.1582\textwidth-2\tabcolsep\relax}|>{\raggedright\arraybackslash}p{\dimexpr0.2101\textwidth-2\tabcolsep\relax}|>{\raggedright\arraybackslash}p{\dimexpr0.0437\textwidth-2\tabcolsep\relax}|>{\raggedright\arraybackslash}p{\dimexpr0.0684\textwidth-2\tabcolsep\relax}|>{\raggedright\arraybackslash}p{\dimexpr0.1043\textwidth-2\tabcolsep\relax}|>{\raggedright\arraybackslash}p{\dimexpr0.0594\textwidth-2\tabcolsep\relax}|>{\raggedright\arraybackslash}p{\dimexpr0.3559\textwidth-2\tabcolsep\relax}}
\toprule
channel (ledger row) & arena & $\dim$ & $\delta$ & cubic constant & $Q$ & $m_1^2:m_2^2:m_3^2$ \\
\addlinespace[2pt]
\midrule
\endfirsthead
\toprule
channel (ledger row) & arena & $\dim$ & $\delta$ & cubic constant & $Q$ & $m_1^2:m_2^2:m_3^2$ \\
\addlinespace[2pt]
\midrule
\endhead
\bottomrule
\endlastfoot
$\ell$ [FT.109–FT.113] & seal$\times$transfer [Def-H.41] & $9$ & $\tfrac29$ & $\sqrt2\cos\tfrac29$ & $\tfrac23$ & $1:206.770315972729\ldots:3477.472837104\ldots$ \\
$q$ [FT.116–FT.119] & seal$\times$transfer$\times$colour [Def-H.42] & $27$ & $\tfrac2{27}$ & $\sqrt2\cos\tfrac29$ & $\tfrac23$ & $1:3.563695074946\ldots:139.3386880284\ldots$ \\
$\nu$ (this tier) & $\mathcal K_1$ [derived] & $16$ & $\tfrac18$ & $\sqrt2\cos\tfrac38$ & $\tfrac23$ & $1:9.583926992379\ldots:271.994748749143\ldots$ \\
\end{longtable}
}

%dsp-level:0
\[
\text{the family count is}\ 3\ \text{in every channel}:\ 3
\]

%dsp-level:0
\[
= \dim F_{J^\pm}\ [\mathrm{CT0.12};\ \mathrm{FT.111}(3):\ \textstyle\sum_i\lambda_i=3]\ \text{— a datum for deficit}\ 3,
\]

%dsp-level:0
\[
\ \neg\ \text{a selection principle}\ [\text{item}\ 21(d)\ \text{stands}].
\]

\textbf{Req.} FT.142; FT.111(2)–(3), FT.112, FT.113 (\textbf{Tier F.37}: the democratic spectrum form, the normalization, the angle discipline, the $m=\mu r^2$ vocabulary and its enclosure pipeline); FT.119 (\textbf{Tier F.38}: the phase-free $Q$, the anchored forms, the triple-angle precedent); L3, L4 (\textbf{B}: the $m^2$ inertial-read vocabulary); P4 (\textbf{A}: the free-scale clause); exact symbolic identities in $\mathbb Q[\sqrt2,\sqrt3][c,s]/(c^2{+}s^2{-}1)$, rational trigonometric enclosures (Taylor remainder bounds), integer Newton certificates, monotone rational interval arithmetic, Vieta reads [math; machine: \texttt{checks/\allowbreak{}verify\_\allowbreak{}H\_\allowbreak{}rung1\_\allowbreak{}neutrino.py} §C].

\textbf{Pf.}

%dsp-level:0
\[
(1):\ r_i^{(\nu)}=A+B\cos(\delta_\nu+\tfrac{2\pi i}3)\ \text{at}\ A=1,\ B=\sqrt2\ [\mathrm{FT.111}(3)];
\]

%dsp-level:0
\[
\qquad \cos(\tfrac18+\tfrac{2\pi}3)=-\tfrac c2-\tfrac{\sqrt3}2s,\quad\cos(\tfrac18+\tfrac{4\pi}3)=-\tfrac c2+\tfrac{\sqrt3}2s\ [\text{addition formulas}];
\]

%dsp-level:0
\[
\qquad \sqrt2\cdot\tfrac{\sqrt3}2=\tfrac{\sqrt6}2
\]

%dsp-level:0
\[
c, s\ \text{enclosed by alternating Taylor partial sums with exact rational remainder bounds at}\ x
\]

%dsp-level:0
\[
= \tfrac18\ (\text{widths}<2^{-90});
\]

%dsp-level:0
\[
\qquad \sqrt2,\sqrt6\ \text{enclosed by integer Newton certificates};
\]

%dsp-level:0
\[
\qquad r_1^{(\nu)},r_2^{(\nu)}>0\ [\text{positive lower endpoints};\ \text{machine}]
\]

%dsp-level:0
\[
(2):\ \textstyle\sum_i\cos(\delta_\nu+\tfrac{2\pi i}3)
\]

%dsp-level:0
\[
= 0,
\]

%dsp-level:0
\[
\ \sum_i\cos^2(\delta_\nu+\tfrac{2\pi i}3)
\]

%dsp-level:0
\[
= \tfrac32\ [\mathrm{FT.111}(3)\ \text{identities, arbitrary phase}]\implies\sum_ir^{(\nu)}
\]

%dsp-level:0
\[
= 3,
\]

%dsp-level:0
\[
\ \sum_i(r^{(\nu)})^2
\]

%dsp-level:0
\[
= 3+2\sqrt2\cdot0+2\cdot\tfrac32\qquad = 6;
\]

%dsp-level:0
\[
\qquad \text{pair sum by Vieta}=\tfrac32
\]

%dsp-level:0
\[
\text{depressed cubic}\ (\lambda-1)^3-\tfrac32(\lambda-1)-\tfrac{\sqrt2}2\cos3\delta_\nu=0\ [\mathrm{FT.111}(2)\ \text{at}\ B^2=2];
\]

%dsp-level:0
\[
\qquad \cos3\delta_\nu=\cos\tfrac38=4c^3-3c\ [\text{triple angle};\ \text{machine: ring identity}];
\]

%dsp-level:0
\[
\qquad \times2:\ 2y^3-3y-\sqrt2\cos\tfrac38=0,\ y=r^{(\nu)}-1
\]

%dsp-level:0
\[
0\in2y_i^3-3y_i-\sqrt2\cos\tfrac38\ \text{enclosure for each}\ i\ [\text{machine}];
\]

%dsp-level:0
\[
\qquad \cos\tfrac38\ \text{vs}\ \cos\tfrac29:\ \text{disjoint enclosures}\ [\text{machine}];
\]

%dsp-level:0
\[
\qquad \text{ring-side: all identities exact in}\ \mathbb Q[\sqrt2,\sqrt3][c,s]/(c^2{+}s^2{-}1)\ [\text{machine}]
\]

%dsp-level:0
\[
(3):\ (r_1^{(\nu)})^2<(r_2^{(\nu)})^2<(r_0^{(\nu)})^2\ [\text{disjoint enclosures}];
\]

%dsp-level:0
\[
\qquad \text{ratios by monotone interval division, widths}<2^{-50}\ [\text{machine}];
\]

%dsp-level:0
\[
\qquad \text{anchored forms: the same divisions against}\ (r_0^{(\nu)})^2\ [\text{machine}]
\]

%dsp-level:0
\[
(4):\ Q=\tfrac{\sum_i(r^{(\nu)})^2}{(\sum_ir^{(\nu)})^2}=\tfrac69=\tfrac23\ [\text{exact Fractions from}\ (2)];
\]

%dsp-level:0
\[
\qquad \mu_\nu\ \text{cancels}\ [\text{homogeneity}]
\]

%dsp-level:0
\[
(5):\ 16\cdot\tfrac18=2,\quad9\cdot\tfrac29=2,\quad27\cdot\tfrac2{27}=2;
\]

%dsp-level:0
\[
\qquad 3\cdot\tfrac18=\tfrac38\ne\tfrac29\ne\tfrac2{27};
\]

%dsp-level:0
\[
\qquad \tfrac38-\tfrac29=\tfrac{27-16}{72}=\tfrac{11}{72};
\]

%dsp-level:0
\[
\qquad \tfrac18-\tfrac29=\tfrac{9-16}{72}=-\tfrac7{72};
\]

%dsp-level:0
\[
\qquad \tfrac18-\tfrac2{27}=\tfrac{27-16}{216}=\tfrac{11}{216}\ [\text{exact Fractions};\ \text{machine}]
\]

%dsp-level:0
\[
216\cdot(\tfrac18,\tfrac29,\tfrac2{27})=(27,48,16);
\]

%dsp-level:0
\[
\qquad \tfrac2{27}<\tfrac18<\tfrac29\ [\text{cross-multiplication};\ \text{machine}];
\]

%dsp-level:0
\[
\qquad (6):\ \text{frozen rows verbatim}\ [\mathrm{FT.113},\ \mathrm{FT.119}(3),(5);\ \text{machine: the cited decimals re-verified as enclosure floors at their stated widths}].\ \blacksquare
\]

\textbf{Phys.} The third channel runs the same democratic machine on the derived sixteen: one large line, one middle line, one small line, all exact in $\sqrt2$, $\sqrt6$ and $\cos\tfrac18$, with the universal $Q=\tfrac23$ holding again — the invariant knows the normalization, not the phase. Its cubic is its own: triple the angle and nothing from the other two channels appears, so the rung-one hierarchy is structurally distinct from the linked $\ell/q$ pair, sharing with them exactly one integer, the seal coefficient $2$. The angle sits between its two siblings and the hierarchy with them: milder than the $\ell$-channel's $3477$-fold span, steeper than the $q$-channel's $139$-fold span. Every number here is dimensionless; the scale $\mu_\nu$ is free internal data, and the $m^2$ name is the corpus's own inertial read — never a particle.

\vspace{2pt}\noindent\rule{\textwidth}{0.4pt}\vspace{2pt}

\begin{theorem}[{\texorpdfstring{FT.144 — The Janus-arrow theorem: no backward materialization, and single-sector occupancy at rung one}{FT.144 - The Janus-arrow theorem: no backward materialization, and single-sector occupancy at rung one}}]\label{tag:FT-144}
\textbf{Thm.}

%dsp-level:0
\[
(1)\quad\textbf{the valid-universe condition:}
\]

%dsp-level:0
\[
\qquad \text{valid universe}\iff\text{real-edge network closed}\ \wedge\ \#J^+
\]

%dsp-level:0
\[
= \#J^-\ \wedge\ \in\operatorname{Im}(\text{active mask})\ \wedge\ \neg\exists\ \text{dangling ledger boundary}\ [\mathrm{D2.11},\ \mathrm{D1.3}];
\]

%dsp-level:0
\[
\text{observable fermions are termini of Janus-closed bound chains}\ [\mathrm{Def\text{-}H.4}];
\]

%dsp-level:0
\[
\qquad \text{closed}\iff\partial c=0\ \wedge\ \#J^+=\#J^-\ [\mathrm{D1.4}];
\]

%dsp-level:0
\[
\qquad \text{bound content is fermionic}\ [\mathrm{D1.5}];
\]

%dsp-level:0
\[
(2)\quad\textbf{the forward arrow}\ (\text{per-pair transition domains}):
\]

%dsp-level:0
\[
\qquad J^+:\ n_{\mathcal R}:0\to1\ (\text{crystallization});
\]

%dsp-level:0
\[
\qquad J^-:\ n_{\mathcal R}:1\to0\ (\text{dissolution})
\]

%dsp-level:0
\[
\qquad [\mathrm{D1.2}];
\]

%dsp-level:0
\[
J^-\ \text{is undefined at}\ n_{\mathcal R}
\]

%dsp-level:0
\[
= 0\ \ \wedge\ \ J^+\ \text{is undefined at}\ n_{\mathcal R}
\]

%dsp-level:0
\[
= 1:\ \text{the domains force}\ \textbf{strict alternation}\ \text{— the only admissible single-pair histories are the prefixes of the alternating word}\ (J^+J^-J^+J^-\cdots);
\]

%dsp-level:0
\[
\text{every admissible history obeys the prefix law}\ \ \#J^+(\text{prefix})\ \ge\ \#J^-(\text{prefix}),
\]

%dsp-level:0
\[
\ \ \text{with equality}
\]

%dsp-level:0
\[
\iff n_{\mathcal R}\qquad = 0;
\]

%dsp-level:0
\[
\qquad \text{the completed single-pair history of length}\ 2m\ \text{is unique:}\ (J^+J^-)^m,\ \#J^+=\#J^-=m;
\]

%dsp-level:0
\[
\text{two-pair completed interleavings of total length}\ 2n:
\]

%dsp-level:0
\[
\qquad \sum_{m_0+m_1=n}\frac{(2n)!}{(2m_0)!\,(2m_1)!}\ =\ 2,\ 8,\ 32\quad(2n=2,4,6)
\]

%dsp-level:0
\[
\qquad [\text{machine: exhaustive enumeration}\ \le8\ (\text{single pair}),\ \le6\ (\text{two pairs})\ \text{agrees with the closed form}];
\]

%dsp-level:0
\[
\text{a materialized closed chain with any dissolution-first edge}\ (J^-\to J^+\ \text{as its leading local transition})\ \text{violates the prefix law at its first event: it is outside the transition domains, hence excluded from every valid universe}\ [\mathrm{D2.11}];
\]

%dsp-level:0
\[
\text{the sealed universe-level balance}\ \#J^+
\]

%dsp-level:0
\[
= \#J^-\ [\mathrm{D1.3}]\ \text{is the balanced-sector restriction of the move algebra's budget identity}\ a_+-a_-
\]

%dsp-level:0
\[
= |R'|-|R|\ [\mathrm{FT.108}(4)];
\]

%dsp-level:0
\[
(3)\quad\textbf{not a mass statement:}\ \ \text{the obstruction of}\ (2)\ \text{is the transition domain of D1.2 conjoined with the closure of D1.4/D2.11 — a combinatorial locus statement: no scale, no P4 parameter and no mass enters it}\ [\text{inventory: the premises contain}\ \neg\mathfrak h,\ \neg\mu,\ \neg h_0];
\]

%dsp-level:0
\[
\neg\exists\ \text{a value of any configuration datum that un-excludes the backward materialization};
\]

%dsp-level:0
\[
\qquad \text{contrast class: the deficit-12 ceilings are range bounds on reads}\ [\mathrm{FT.76}(3),\ \mathrm{FT.79}(2)]\ \text{— the present obstruction is a domain restriction, prior to every read};
\]

%dsp-level:0
\[
(4)\quad\textbf{the vacuum residence:}\ \ \text{the reversed reading exists only as formal bookkeeping. Two distinct reversals, adjudicated separately:}
\]

%dsp-level:0
\[
\text{(i)}\ \Theta_{\mathfrak L}\ (T):\ \text{ledger order reversed}\ \wedge\ J^+\leftrightarrow J^-\ [\mathrm{D2.7};\ \mathrm{FT.133}(1)]:\ \text{an involution on configuration data}\ [\text{machine: it bijects the completed admissible sequences onto themselves, counts preserved — the formal mirror}];
\]

%dsp-level:0
\[
\text{(ii) the pure order reversal}\ (\text{the ledger read backwards, labels unchanged — ``running the arrow in reverse''}):\ \text{every completed nonempty admissible history ends with}\ J^-\ (n_{\mathcal R}\ \text{returns to}\ 0\ \text{only via}\ 1\to0),
\]

%dsp-level:0
\[
\ \text{so its pure reversal leads with}\ J^-:\ \text{inadmissible at the first event}\ [\text{machine: exhaustive}];
\]

%dsp-level:0
\[
\text{and the tick sequence admits no inverse in either case}\ [\mathrm{A6}:\ \neg\exists\,\tau^{-1}]:\ \text{the ledger count}\ |\mathfrak L(\tau_n)|
\]

%dsp-level:0
\[
= \sum_{k<n}n_{\mathrm{real\_out}}(\tau_k)\ \text{is non-decreasing along every admissible history}\ (n_{\mathrm{real\_out}}\ge0)\ \text{and strictly increases at every crystallization: the reversed reading of a multi-crystallization history strictly decreases somewhere, and by (ii) the reversal of every completed word is inadmissible at its first event}\ [\text{machine: exhaustive}]\implies\text{no forward tick realization of any reversed reading exists};
\]

%dsp-level:0
\[
\text{the fluctuation content resides in the}\ \mathcal V\ \text{sector}\ [\mathrm{CT0.1}]:
\]

%dsp-level:0
\[
\quad n_{\mathcal R}=0\implies\forall(p,q):\ (p,q)\in\mathcal V\ \wedge\ \rho_{pq}=\rho_\circ=\tfrac14I,
\]

%dsp-level:0
\[
\quad S_{\mathrm{raw}}(\rho_\circ)=\ln4=\max,\quad\arg\max=\{\tfrac14I\}
\]

%dsp-level:0
\[
\quad [\mathrm{A5},\ \mathrm{CT0.2};\ \text{machine:}\ \mathfrak T\rho_\circ=1,\ \operatorname{charpoly}=(x-\tfrac14)^4,\ \text{dyadic entropy witnesses}<\ln4];
\]

%dsp-level:0
\[
\text{virtual pairs are unregistered}\ [\mathrm{D1.1}:\ \neg\text{registration}\implies e\in\mathcal V];
\]

%dsp-level:0
\[
\qquad \text{chains require real edges}\ [\mathrm{D1.4}:\ e_i\in E_{\mathcal R}]:\ |E_{\mathcal R}|
\]

%dsp-level:0
\[
= 0\implies\text{the closed-chain count is}\ 0\implies\neg\exists\ \text{fermion terminus}\ [\mathrm{Def\text{-}H.4};\ \text{machine}];
\]

%dsp-level:0
\[
\text{the fluctuation resolvent is even in balanced pairs:}\quad\sum_n\langle\varnothing|\hat U^n|\varnothing\rangle\,w^n=\frac1{1-w^2x_+x_-}
\]

%dsp-level:0
\[
\quad [\mathrm{FT.108}(5)\ \text{re-derived};\ \text{machine in}\ \mathbb Q[w,x_+,x_-]/(w^{12})]:
\]

%dsp-level:0
\[
\text{every nonzero return monomial carries equal}\ x_+\text{- and}\ x_-\text{-degree}\ [\text{machine degree census}]:\ \text{a}\ J^-\text{-leading monomial has zero support};
\]

%dsp-level:0
\[
\text{the zero-projection/no-pole pattern applies}\ [\mathrm{CT4.2}:\ \mathcal P_MF\mathcal P_M=0\implies\neg\exists\ \text{accessible pole};\ \mathrm{CT2.5}:\ \text{poles}\iff1\in\operatorname{spec}\mathbf F_M]:\ \text{unregistered content supports no accessible pole at any rung};
\]

%dsp-level:0
\[
(5)\quad\textbf{carrier reconciliation — single-sector occupancy:}\ \ P_L,
\]

%dsp-level:0
\[
P_R
\]

%dsp-level:0
\[
= \tfrac12(\mathrm{id}\pm\hat\chi)\ \text{are constructible, exact, rank-}1,
\]

%dsp-level:0
\[
\ \hat\omega\text{-exchanged}\ [\mathrm{Def\text{-}H.33},\ \mathrm{FT.122}(2)\text{–}(4);\ \text{machine}]:\ \text{mathematical objects on}\ \mathcal H_\omega\ \text{whether or not occupied};
\]

%dsp-level:0
\[
\textbf{unconditional}\ (\text{derived}):\ \text{the rung-one materialized locus is single-arrow by}\ (2)\ \text{— one orientation sector of the doublet basis}\ (|c\rangle\ \text{carries the materialized ledger-ordered tuple}\ [\mathrm{D1.4};\ \mathrm{FT.134}(1):\ C,P,T\ \text{fix the chain tuple pointwise}];\ |\omega(c)\rangle\ \text{is the reversal image}\ [\mathrm{D2.6};\ \mathrm{FT.88}(1):\ R\mapsto R^{\mathrm T}]);
\]

%dsp-level:0
\[
\text{its complement is unoccupied at rung one, and by}\ (4)\ \text{the reversed content is vacuum-sector content}\ (\text{all-virtual, unregistered, no accessible pole});
\]

%dsp-level:0
\[
\textbf{conditional}\ (\text{the coupling}):
\]

%dsp-level:0
\[
\quad |c\rangle=\tfrac1{\sqrt2}\bigl(|L\rangle+|R\rangle\bigr)
\]

%dsp-level:0
\[
\quad [\mathrm{Def\text{-}H.34}\ \mathrm{Rmk};\ \text{machine}]:\ \text{the orientation rays carry no definite chirality}\implies\text{transferring single-arrow occupancy to single-ray occupancy requires an arrow}\leftrightarrow\text{ray coupling — the deficit}\ 1'\ \text{assignment rule}\ [\mathrm{FT.125}(4)];
\]

%dsp-level:0
\[
\text{the coupling's forced shape is derived: any}\ \hat\omega\text{-covariant coupling maps the single-arrow locus into a single chiral ray}\ (\text{the locus is one sector;}\ \hat\omega\ \text{is the ray exchange}\ [\mathrm{FT.122}(4)];\ \text{machine: exhaustive on the}\ 2!\ \text{bijections}),
\]

%dsp-level:0
\[
\ \text{and the complementary ray has occupancy zero at rung one — the vacuum sector of the doublet};
\]

%dsp-level:0
\[
\text{which of}\ P_L/P_R\ \text{names the occupied ray is reserved naming content}\ [\mathrm{Def\text{-}H.33}\ \text{canonicity Rmk}]:\ \text{the occupancy pair}\ (1,0),
\]

%dsp-level:0
\[
\ \text{up to ray exchange, is this clause's entire content};
\]

%dsp-level:0
\[
(6)\quad\textbf{rung scope — the force-side half opens at rung two:}\ \ \text{at rung one the static force matrix on the extended carrier is}\ \mathbf F_{\mathrm{ext}}(\mathbf c)
\]

%dsp-level:0
\[
= h(\mathbf c)\otimes I_\omega\ \text{with}\ F_{LR}
\]

%dsp-level:0
\[
= F_{RL}
\]

%dsp-level:0
\[
= 0\ \text{exactly}\ [\mathrm{FT.127}(9),\ \mathrm{FT.128}(2):\ \text{the no-flip theorem};\ \text{machine}]:\ \text{the two chiral sectors are force-decoupled; the flipping set at the derived locus}
\]

%dsp-level:0
\[
= \varnothing\ [\mathrm{FT.128}(6)];
\]

%dsp-level:0
\[
\text{at rung two EM activates}\ [\mathrm{CT0.13}]\ \text{and the promoted core read opens the unique flip channel:}\quad\mathbf F_{\mathrm{ext}}(\mathbf c,q)
\]

%dsp-level:0
\[
= h(\mathbf c)\otimes I_2+q\,
\]

%dsp-level:0
\[
I_2\otimes\hat\kappa,
\]

%dsp-level:0
\[
\qquad F_{LR}=q\,I_2\ne0\iff q\ne0\quad[\mathrm{FT.129}(4)\text{–}(5),\ \text{Def-H.43};\ \mathrm{EM.9}:\ q\in\{-1,0,+1\}];
\]

%dsp-level:0
\[
\operatorname{charpoly}\mathbf F_{\mathrm{ext}}(\mathbf c,q)
\]

%dsp-level:0
\[
= \operatorname{charpoly}_h(x-q)\cdot\operatorname{charpoly}_h(x+q):\ \text{the chiral-partner splitting}\ 2|q|\in\{0,2\}\ [\mathrm{FT.130}(5);\ \text{machine: Faddeev–LeVerrier over}\ \mathbb Q(i)]:\ \text{both chiral sectors become force-coupled for charged cores};
\]

%dsp-level:0
\[
\text{at rung three colour adds no carrier flip: sealed ends are blind}\ [\mathrm{FT.117}],
\]

%dsp-level:0
\[
\ \text{colour acts on the slot carrier}\ [\mathrm{FT.127}(7)],
\]

%dsp-level:0
\[
\ \text{the seed is colour-neutral by theorem}\ [\mathrm{EW.32}]:\ \text{the flipping set remains}\ \{F_{\mathrm{EM}}\ \text{core read}\}\ \text{— minimal and unique}\ [\mathrm{FT.129}(5);\ \text{machine: removing the}\ \hat\kappa\ \text{term kills}\ F_{LR}\ \text{at every profile}];
\]

%dsp-level:0
\[
\text{the arrow law of}\ (2)\ \text{is rung-blind}\ [\mathrm{D1.2}'\text{s domains carry no rung argument}]:\ \text{what is rung-one specific is the coincidence of the two obstructions — no flip channel (force side)}\ \wedge\ \text{single-arrow locus (ledger side)};
\]

%dsp-level:0
\[
\ \text{the single-sector theorem is a rung-one theorem},
\]

%dsp-level:0
\[
\ \neg\ \text{a universal law}.
\]

\textbf{Req.} A5, A6, A8 (Parts 1–2); D1.1, D1.2, D1.3, D1.4, D1.5, D2.6, D2.7, D2.11 (Parts 1–2); CT0.1, CT0.2, CT0.13 (Part 1); CT2.5 (Part 2); EW.32 (Part 2); EM.9 (Part 2); CT4.2 (Part 2); Def-H.4 (Tier F.2); FT.40 (Tiers F.12–F.13); FT.76, FT.79 (Tiers F.26–F.28); FT.88 (Tiers F.29–F.30); FT.108 (Tier F.36); FT.117 (Tier F.38); Def-H.33, FT.122 (Tier F.39); Def-H.34, FT.127, FT.128, FT.129, FT.130 (Tier F.40); FT.133, FT.134 (Tier F.41); Def-H.4, Def-H.34, Def-H.43; exhaustive history enumeration, multinomial interleaving counts, exact polynomial resolvent identity in $\mathbb Q[w,x_+,x_-]/(w^{12})$, exact $\mathbb Q(i)$ block algebra, Faddeev–LeVerrier, exact rational $\log2$ enclosures [math; machine: \texttt{checks/\allowbreak{}verify\_\allowbreak{}H\_\allowbreak{}rung1\_\allowbreak{}neutrino.py} §D–§E].

\textbf{Pf.}

%dsp-level:0
\[
(1):\ \mathrm{D2.11},\ \mathrm{D1.3},\ \mathrm{D1.4},\ \mathrm{D1.5},\ \mathrm{Def\text{-}H.4}\ \text{verbatim (frozen)}
\]

%dsp-level:0
\[
(2):\ \mathrm{D1.2}:\ \text{the move domains are}\ (0\to1)\ \text{for}\ J^+,\ (1\to0)\ \text{for}\ J^-;
\]

%dsp-level:0
\[
\qquad \text{induction on prefixes: the state after any admissible prefix is}\ n_{\mathcal R}
\]

%dsp-level:0
\[
= \#J^+-\#J^-\in\{0,1\},
\]

%dsp-level:0
\[
\ \text{and the next move is forced}\implies\text{strict alternation and the prefix law}\ [\text{machine: exhaustive at lengths}\le8]
\]

%dsp-level:0
\[
\text{completed}\iff n_{\mathcal R}=0\iff\#J^+=\#J^-;
\]

%dsp-level:0
\[
\qquad \text{the completed single-pair word of length}\ 2m\ \text{is unique:}\ (J^+J^-)^m\ [\text{machine}];
\]

%dsp-level:0
\[
\qquad \text{two-pair interleavings}:\ \text{each pair's projection is an alternating word};
\]

%dsp-level:0
\[
\ \text{completed counts at total length}\ 2n:\ \textstyle\sum_{m_0+m_1=n}\binom{2n}{2m_0}
\]

%dsp-level:0
\[
= 2,
\]

%dsp-level:0
\[
8,
\]

%dsp-level:0
\[
32\ (2n=2,4,6)\ [\text{machine enumeration}=\text{multinomial closed form}];
\]

%dsp-level:0
\[
\ \text{global balance at completion}\ [\text{machine}]
\]

%dsp-level:0
\[
\text{dissolution-first chain}\implies\text{prefix law violated at the first event of that edge}\ [\text{machine witness}]\implies\neg\text{admissible}\implies\ \bot\ \text{in every valid universe}\ [\mathrm{D2.11}];
\]

%dsp-level:0
\[
\qquad \mathrm{FT.108}(4)\ \text{identity re-cited (frozen)}
\]

%dsp-level:0
\[
(3):\ \text{premise inventory of}\ (2):\ \{\mathrm{D1.2}\ \text{domains},\ \mathrm{D1.3}/\mathrm{D1.4}/\mathrm{D2.11}\ \text{closure}\}:\ \text{no dimensionful or P4 quantity occurs}\ [\text{scan}];
\]

%dsp-level:0
\[
\qquad \text{contrast rows}:\ \mathrm{FT.76}(3),\ \mathrm{FT.79}(2)\ \text{are read-range bounds}\ [\text{cited, not re-derived}]
\]

%dsp-level:0
\[
(4):\ \Theta_{\mathfrak L}
\]

%dsp-level:0
\[
= \text{reverse}+\text{label swap}\ [\mathrm{D2.7};\ \mathrm{FT.133}(1)]:\ \text{involution}\ [\text{machine}];
\]

%dsp-level:0
\[
\ \text{completed admissible words biject onto themselves}\ [\text{machine: counts preserved at every length}\le8]
\]

%dsp-level:0
\[
\text{pure order reversal}:\ \text{every completed nonempty admissible word ends in}\ J^-\ [\text{state}\ 0\ \text{is reached only via}\ 1\to0;\ \text{machine: exhaustive}]\implies\text{its reversal leads with}\ J^-:\ \text{inadmissible at the first event}\ [(2)\ \text{domains};\ \text{machine}]
\]

%dsp-level:0
\[
\mathrm{A6}:\ |\mathfrak L(\tau_{n+1})|=|\mathfrak L(\tau_n)|+n_{\mathrm{real\_out}}(\tau_n),\ n_{\mathrm{real\_out}}\ge0\implies\text{monotone non-decreasing};
\]

%dsp-level:0
\[
\qquad \text{the reversed reading of a multi-crystallization history strictly decreases somewhere, and every completed word's reversal is inadmissible from the start}\implies\neg\text{forward-realizable}\ [\text{machine: exhaustive};\ \neg\exists\,\tau^{-1}\ \text{axiomatic}]
\]

%dsp-level:0
\[
\rho_\circ=\tfrac14I_4:\ \mathfrak T=1,\ \operatorname{charpoly}=(x-\tfrac14)^4\ [\text{machine over}\ \mathbb Q];
\]

%dsp-level:0
\[
\qquad S_{\mathrm{raw}}(\rho_\circ)=-4\cdot\tfrac14\ln\tfrac14=\ln4=2\ln2\ [\text{exact rational}\ \log\ \text{enclosure, width}<2^{-80};\ \text{machine}];
\]

%dsp-level:0
\[
\qquad \text{dyadic witnesses}:\ S(\operatorname{diag}(\tfrac12,\tfrac14,\tfrac18,\tfrac18))
\]

%dsp-level:0
\[
= \tfrac74\ln2<2\ln2\ [\text{exact coefficient comparison};\ \text{machine}];
\]

%dsp-level:0
\[
\qquad \arg\max\ \text{uniqueness}\ [\text{Gibbs lemma};\ \mathrm{CT0.2}]
\]

%dsp-level:0
\[
|E_{\mathcal R}|
\]

%dsp-level:0
\[
= 0\implies\#\{\text{tuples}\ (e_1,\dots,e_k)\subseteq E_{\mathcal R}\}
\]

%dsp-level:0
\[
= 0\ \forall k
\]

%dsp-level:0
\[
\ge 1\ [\text{machine}]\implies\neg\exists\ \text{closed chain}\ [\mathrm{D1.4}]\implies\neg\exists\ \text{terminus}\ [\mathrm{Def\text{-}H.4}]
\]

%dsp-level:0
\[
\hat U=\bigl(\begin{smallmatrix}0&x_-\\ x_+&0\end{smallmatrix}\bigr)\ [\mathrm{FT.108}(4)\ \text{generator at single-pair scope}]:\quad\langle\varnothing|\hat U^n|\varnothing\rangle=[\,n\ \text{even}\,]\,(x_+x_-)^{n/2}\ [\text{machine to}\ n=12];\qquad(1-w^2x_+x_-)\sum_n\langle\varnothing|\hat U^n|\varnothing\rangle w^n\equiv1\ [\mathbb Q[w,x_+,x_-]/(w^{12});\ \text{machine coefficient-wise}]
\]

%dsp-level:0
\[
\deg_{x_+}=\deg_{x_-}\ \text{on every nonzero monomial}\ [\text{machine census}];
\]

%dsp-level:0
\[
\qquad \ker Q=\operatorname{span}\{(x_-,x_+)\}:\ (x_-,x_+)Q=0,\ \det Q=0\ne Q\implies\text{rank}\ 1\ [\mathrm{FT.108}(4)\ \text{re-derived};\ \text{machine}]
\]

%dsp-level:0
\[
(5):\ \mathrm{FT.122}(2)\text{–}(4)\ \text{re-derived entrywise over}\ \mathbb Q(i)\ [\text{machine}]:\ P_{L/R}^2
\]

%dsp-level:0
\[
= P_{L/R},
\]

%dsp-level:0
\[
\ P_LP_R
\]

%dsp-level:0
\[
= 0,
\]

%dsp-level:0
\[
\ P_L+P_R
\]

%dsp-level:0
\[
= \mathrm{id},
\]

%dsp-level:0
\[
\ \mathfrak T P_{L/R}\qquad = 1;
\]

%dsp-level:0
\[
\qquad \hat\omega P_L\hat\omega=P_R;
\]

%dsp-level:0
\[
\qquad |c\rangle=\tfrac12(|L_u\rangle+|R_u\rangle)\ \text{unnormalized}\ [\text{machine}]
\]

%dsp-level:0
\[
\text{arrow-sector census}:\ \text{the materialized locus}=\{|\,c\,\rangle\text{-sector}\}\ \text{by}\ (2);
\]

%dsp-level:0
\[
\qquad \hat\omega\text{-covariant bijections}\ \{\text{arrow sectors}\}\to\{\text{rays}\}:\ 2!\ \text{candidates, each maps the single-arrow locus to one ray and its complement to the other}\ [\text{machine: exhaustive}];
\]

%dsp-level:0
\[
\qquad \text{naming residue}:\ \mathrm{Def\text{-}H.33}\ \text{canonicity Rmk verbatim}
\]

%dsp-level:0
\[
(6):\ \mathrm{FT.127}(9),\ \mathrm{FT.128}(2),(6)\ \text{re-derived at}\ 36\ \text{exact rational profiles}:\ F_{LR}=0\ \text{entrywise};
\]

%dsp-level:0
\[
\qquad [\mathbf F_{\mathrm{ext}},\ I_2\otimes\hat\chi]=0\ [\text{machine}]
\]

%dsp-level:0
\[
\mathbf F_{\mathrm{ext}}(\mathbf c,q)\ \text{blocks in the chiral carrier basis}\ [\mathrm{Def\text{-}H.34}'\text{s}\ S\text{-transform};\ \text{machine}]:\ F_{LL}
\]

%dsp-level:0
\[
= F_{RR}
\]

%dsp-level:0
\[
= h,
\]

%dsp-level:0
\[
\ F_{LR}
\]

%dsp-level:0
\[
= F_{RL}
\]

%dsp-level:0
\[
= qI_2\ \text{at}\ q\in\{-1,0,1\}\times36\ \text{profiles};
\]

%dsp-level:0
\[
\qquad \operatorname{charpoly}\ \text{factorization by Faddeev–LeVerrier coefficient identity}\ [\text{machine}]
\]

%dsp-level:0
\[
\text{minimality rows}:\ \text{removing}\ qI_2\otimes\hat\kappa\ \text{kills}\ F_{LR}\ \text{at every profile};
\]

%dsp-level:0
\[
\ \text{keeping only it leaves}\ F_{LR}
\]

%dsp-level:0
\[
= qI_2\ [\mathrm{FT.129}(5);\ \text{machine}];
\]

%dsp-level:0
\[
\qquad \text{colour rows}:\ \mathrm{FT.117},\ \mathrm{FT.127}(7),\ \mathrm{EW.32}\ \text{verbatim (frozen)};
\]

%dsp-level:0
\[
\qquad \mathrm{D1.2}\ \text{carries no rung argument}\ [\text{verbatim scan}].\ \blacksquare
\]

\textbf{Phys.} Ask the ledger which direction a real edge may be born in, and the answer is a domain, not a preference: dissolution presupposes crystallization, so every materialized history is crystallization-first on every edge, and the prefix law is exhaustive at every length checked. A backward-arrow fermion is therefore not heavy, not suppressed, not decoupled — it is ungrammatical: no ledger history spells it, at any value of any parameter. What the reversal does spell is bookkeeping: the formal T-images and the balanced fluctuation pairs of the all-virtual state, which carry no registration, no chain, no terminus and no accessible pole. On the carrier the consequence is an occupancy statement — one sector materializes, its complement is vacuum — and because the orientation rays superpose both chiral rays equally, the last step from arrow to ray is exactly the missing assignment rule; the theorem fixes its shape and its output, $(1,0)$ up to exchange, and leaves the naming where the corpus has always left it. At rung two the force side opens: the promoted core read is the unique flip channel, and the coincidence that makes rung one single-sector does not recur.

\vspace{2pt}\noindent\rule{\textwidth}{0.4pt}\vspace{2pt}

\end{theorem}

\pseudohead{theorem}{\texorpdfstring{FT.145}{FT.145}}{\texorpdfstring{The rung-one completion}{The rung-one completion}}\label{tag:FT-145}
\textbf{Thm.}

%dsp-level:0
\[
(1)\quad\textbf{the content summary:}\ \ \text{the rung-one fermion sector consists of: one family channel}\ [\mathrm{FT.142}\text{–}\mathrm{FT.143}:\ \delta_\nu=\tfrac18,\ \text{triplet}\ r_i^{(\nu)},\ \text{ratios}\ 1:9.58\ldots:272.0\ldots,\ Q=\tfrac23];
\]

%dsp-level:0
\[
\ \text{its chirality structure}\ [\mathrm{FT.141}:\ \text{the complete-force}\ 10/6\ \text{grading},\ \text{no}\ S\text{-neutral observable},\ \text{the weak selectivity a state-locus shadow}];
\]

%dsp-level:0
\[
\ \text{its occupancy structure}\ [\mathrm{FT.144}:\ \text{single-arrow materialization},\ \text{single-ray under any}\ \hat\omega\text{-covariant coupling},\ \text{complement vacuum}];
\]

%dsp-level:0
\[
(2)\quad\textbf{the mass read:}
\]

%dsp-level:0
\[
\qquad m_{\mathrm{hid}}=E_{\mathrm{int}}\ln2/\mathfrak h
\]

%dsp-level:0
\[
\qquad [\mathrm{FT.16};\ \mathrm{FT.44}(4):\ \text{universal across every computed floor};\ \mathrm{Def\text{-}H.9}(2):\ m_{\mathrm{hid}}=I_{\mathrm{hid}}/\mathfrak h];
\]

%dsp-level:0
\[
\text{the minimal rung-one bound chain — the digon,}\ E_{\mathrm{int}}=2\ [\mathrm{FT.28}(1)]\ \text{— reads}\quad m_{\mathrm{hid}}=\frac{2\ln2}{\mathfrak h}
\]

%dsp-level:0
\[
\qquad [\text{machine: exact integer coefficient}];
\]

%dsp-level:0
\[
\text{the channel ratios of FT.143 are dimensionless with overall scale}\ \mu_\nu\ \text{undeclared}\ (\mathrm{P4});
\]

%dsp-level:0
\[
\qquad \text{the relating equation between the channel read and the}\ m_{\mathrm{hid}}/\text{pole/sector reads is undeclared}\ [\mathrm{FT.113}\ \text{Status pattern}];
\]

%dsp-level:0
\[
\text{amplitude}\ne\text{mass};\qquad\text{rung}\ne\text{scale}\qquad[\text{standing rules}];
\]

%dsp-level:0
\[
(3)\quad\textbf{the completeness fraction:}
\]

%dsp-level:0
\[
\qquad J^+_3+J^-_3+G_5+W_3+H_1+D^+_1=16=C_{\mathrm{cap}}=\dim\mathcal K_1
\]

%dsp-level:0
\[
\qquad [\mathrm{CT0.12},\ \mathrm{CT0.13},\ \mathrm{CT0.14};\ \text{machine}]:
\]

%dsp-level:0
\[
\mathcal A(1)\ \text{accounts for every pair-channel dimension — no part of the rung-one channel is forceless and no force overruns it}\ [\mathrm{CT0.12}\ \text{Phys pattern}];
\]

%dsp-level:0
\[
\text{the rung-one fermion sector carries:}\ \neg\mathrm{EM}\ (\text{activation at rung}\ 2\ [\mathrm{CT0.13}];\ \text{unactivated in the channel class}\ [\mathrm{FT.142}(1)]),
\]

%dsp-level:0
\[
\quad\neg\text{colour}\ (\text{rung}\ 3\ [\mathrm{CT0.13}];\ \text{structurally colourless}\ [\mathrm{FT.142}(1)]),
\]

%dsp-level:0
\[
\neg\text{Yukawa sector}\ (\text{deficit}\ 2\ \text{stands}),
\]

%dsp-level:0
\[
\qquad \neg\text{mixing matrix}\ (\text{deficit}\ 4\ \text{stands}),
\]

%dsp-level:0
\[
\qquad \neg\text{generation selection}\ (\text{deficit}\ 3\ \text{stands};\ \text{item}\ 21(d));
\]

%dsp-level:0
\[
(4)\quad\textbf{the bridge to rung two:}
\]

%dsp-level:0
\[
\qquad \mathcal K_1\subset\mathcal K_2\ (\text{strict}),
\]

%dsp-level:0
\[
\qquad \dim\mathcal K_2=17=\dim\mathcal K_1+1
\]

%dsp-level:0
\[
\qquad [\mathrm{CT0.14}:\ \text{the EM centre is the new summand};\ \text{machine}];
\]

%dsp-level:0
\[
\text{at rung two EM activates}\ [\mathrm{CT0.13}]\ \text{and the flip channel opens},
\]

%dsp-level:0
\[
\ F_{LR}
\]

%dsp-level:0
\[
= qI_2\ [\mathrm{FT.129}(5),\ \text{Def-H.43};\ \mathrm{FT.144}(6)]:\ \text{both chiral sectors become force-coupled};
\]

%dsp-level:0
\[
\text{the}\ 10/6\ \text{split persists on}\ \mathcal K_1:\ \text{the EM summand is a stack-register block}\ [\mathrm{CT0.12}\ \mathrm{Rmk};\ \mathrm{CT0.14}],
\]

%dsp-level:0
\[
\ \neg\ \text{a pair-channel direction}:\ \text{the}\ \operatorname{Ad}(\mathbb S)\ \text{census of the pair}\ 16\ \text{is unchanged};
\]

%dsp-level:0
\[
\text{the}\ \ell\text{-channel angle}\ \delta_\ell
\]

%dsp-level:0
\[
= \tfrac29\ [\mathrm{FT.112}\ \text{— Tier F.37, under Def-H.41's interface definition; cited as frozen channel datum, imported as no premise of this part}]\ \text{governs the rung-two family ratios};
\]

%dsp-level:0
\[
\qquad \text{the two angles are distinct rationals sharing the numerator}\ 2\ [\mathrm{FT.143}(5)];
\]

%dsp-level:0
\[
(5)\quad\textbf{the adjudication ledger:}
\]

%tbl-level:2
{\footnotesize\setlength{\tabcolsep}{4pt}
\begin{longtable}{>{\raggedright\arraybackslash}p{\dimexpr0.4878\textwidth-2\tabcolsep\relax}|>{\raggedright\arraybackslash}p{\dimexpr0.1356\textwidth-2\tabcolsep\relax}|>{\raggedright\arraybackslash}p{\dimexpr0.3766\textwidth-2\tabcolsep\relax}}
\toprule
Statement & Source & Status \\
\addlinespace[2pt]
\midrule
\endfirsthead
\toprule
Statement & Source & Status \\
\addlinespace[2pt]
\midrule
\endhead
\bottomrule
\endlastfoot
$10/6=\allowbreak \operatorname{Ad}(\mathbb S)$ census on $\mathcal K_1$ & FT.141(1) & Derived (FT.136(1) re-derived) \\
Chirality $=$ the complete-force grading & FT.141(2) & Derived \\
Deletion invariance; $45/64$ nontrivial subsets & FT.141(3) & Derived (exhaustive) \\
Weak branch selectivity an $S$-grading consequence & FT.141(4) & Derived ($\sum_iW_i^2=\allowbreak E_0+\allowbreak \mathbb S$) \\
No $S$-neutral observable at rung one & FT.141(5) & Derived \\
$S$-grading vs carrier $\hat\chi$: independent $\mathbb Z_2$'s & FT.141(6) & Derived (FT.126(4)/FT.136(4) re-run) \\
Arena $\mathcal K_1$, $\dim=\allowbreak 16$; candidates $10/9/27/6$ rejected & FT.142(2)–(3) & Derived (no declaration) \\
Seal coefficient $2$, undressed & FT.142(4) & Derived (FT.109(3) re-derived) \\
$\delta_\nu=\allowbreak 1/8$ & FT.142(5) & Derived \\
$r_i^{(\nu)}$; ratios $1:9.58\ldots:272.0\ldots$; $Q=\allowbreak \tfrac23$ & FT.143 & Derived (scale $\mu_\nu$: P4) \\
$3\delta_\nu=\allowbreak \tfrac38\notin\{\delta_\ell,\delta_q\}$; shared numerator $2$ & FT.143(2), (5) & Derived (exact rationals) \\
Forward arrow $J^+\allowbreak \to J^-$; strict alternation; prefix law & FT.144(2) & Derived (exhaustive) \\
Backward materialization excluded; $\ne\allowbreak g$a mass statement & FT.144(2)–(3) & Derived \\
Vacuum residence; balanced fluctuation pairs; no accessible pole & FT.144(4) & Derived \\
Single-ray occupancy under $\hat\omega$-covariant coupling & FT.144(5) & Derived (coupling-conditional; naming reserved) \\
Flip channel opens at rung two & FT.144(6) & Derived (under Def-H.43) \\
$m_{\mathrm{hid}}(\text{digon})=\allowbreak 2\ln2/\mathfrak h$ & FT.145(2) & Derived (FT.28(1) re-cited) \\
Completeness $3{+}3{+}5{+}3{+}1{+}1=\allowbreak 16=\allowbreak C_{\mathrm{cap}}$ & FT.145(3) & Derived \\
Overall scale $\mu_\nu$; relating equation & — & Free (P4); open \\
Species identification of the triplet & — & NOT derived (barred: P1, P5) \\
Arrow$\le\allowbreak ftrightarrow$ray assignment rule & — & Application-layer content \\
\end{longtable}
}

%dsp-level:0
\[
(6)\quad\textbf{the scope bars:}\ \ \text{no particle identification — the channel subscript}\ \nu\ \text{is ledger vocabulary}\ [\mathrm{FT.142}\ \mathrm{Rmk}];
\]

%dsp-level:0
\[
\qquad \text{no external data}\ ([\text{EXT}]\ \text{enters no clause of this document — this standing-rule record is its sole marker});
\]

%dsp-level:0
\[
\text{no new declaration issues: imports are confined to Def-H.4},
\]

%dsp-level:0
\[
\ \text{Def-H.34}\ (\text{extended carrier}),
\]

%dsp-level:0
\[
\ \text{Def-H.43}\ (\text{read promotion — rung-two bridge rows only});
\]

%dsp-level:0
\[
\qquad \text{Def-H.41/Def-H.42 are cited in the rejection rows and imported nowhere};
\]

%dsp-level:0
\[
\text{no fitted parameters}\ (\text{one scale}\ \mathfrak h,\ \text{P1});
\]

%dsp-level:0
\[
\qquad \text{all topology labels are internal bookkeeping};
\]

%dsp-level:0
\[
\qquad \text{no H-sector interpretation enters any theorem statement}.
\]

\textbf{Req.} FT.141–FT.144 (this tier); FT.16, Def-H.9 (Tiers F.1–F.11); FT.28 (Tiers F.1–F.11); FT.44 (Tiers F.14–F.16); CT0.12, CT0.13, CT0.14 (Part 1); FT.112 (\textbf{Tier F.37}, frozen-datum citation); FT.129 (Tier F.40); P1, P4, P5 (Parts 1–2); exact integer arithmetic, exact rational $\log2$ enclosures, exact ranks [math; machine: \texttt{checks/\allowbreak{}verify\_\allowbreak{}H\_\allowbreak{}rung1\_\allowbreak{}neutrino.py} §F].

\textbf{Pf.}

%dsp-level:0
\[
(1):\ \text{assembly of FT.141–FT.144's adjudicated rows}\ [\text{machine: all sections green}]
\]

%dsp-level:0
\[
(2):\ \mathrm{FT.44}(4):\ m_{\mathrm{hid}}=E_{\mathrm{int}}\ln2/\mathfrak h\ \text{across every computed floor};
\]

%dsp-level:0
\[
\qquad \mathrm{FT.28}(1):\ E_{\mathrm{int}}
\]

%dsp-level:0
\[
= 2\ (\text{both cycle edges internal})
\]

%dsp-level:0
\[
\implies m_{\mathrm{hid}}
\]

%dsp-level:0
\[
= 2\ln2/\mathfrak h\ [\text{exact integer coefficient};\ \text{machine}];
\]

%dsp-level:0
\[
\qquad \mu_\nu:\ \neg\exists\ \text{tag fixing it}\ [\text{P4}]
\]

%dsp-level:0
\[
\text{relating equation}:\ \neg\exists\ \text{tag relating the channel ratios to}\ m_{\mathrm{hid}},
\]

%dsp-level:0
\[
\ m^2_{\mathrm{pole}}\ \text{or}\ m^2_B(q,S)\ [\text{inventory scan};\ \mathrm{FT.113}\ \text{Status}]
\]

%dsp-level:0
\[
(3):\ 3+3+5+3+1+1=16\ [\mathrm{CT0.12}\ \text{table};\ \text{machine}];
\]

%dsp-level:0
\[
\qquad \dim\mathcal K_1=16=C_{\mathrm{cap}}\ [\mathrm{CT0.14};\ \mathrm{CT0.6}];
\]

%dsp-level:0
\[
\qquad \neg\text{-rows}:\ \mathrm{CT0.13}\ \text{activation rungs verbatim};
\]

%dsp-level:0
\[
\ \mathrm{FT.142}(1)\ \text{channel class};
\]

%dsp-level:0
\[
\ \text{deficits}\ 2/3/4:\ \neg\exists\ \text{Yukawa, flavour-index or mixing-matrix tag}\ [\text{inventory scan}]
\]

%dsp-level:0
\[
(4):\ \dim\mathcal K_2=\dim\mathcal K_1+\dim Z(\mathfrak u(2))=16+1=17\ [\mathrm{CT0.14};\ \text{machine: strict flag}\ 6\subset16\subset17\subset25];
\]

%dsp-level:0
\[
\qquad \text{EM summand}:\ \text{stack register}\ [\mathrm{CT0.12}\ \mathrm{Rmk}]\implies\operatorname{Ad}(\mathbb S)\ \text{on the pair}\ 16\ \text{untouched}\ [\mathrm{FT.141}(1)\ \text{census unchanged}]
\]

%dsp-level:0
\[
\text{flip row}:\ \mathrm{FT.144}(6)\ \text{verbatim};
\]

%dsp-level:0
\[
\qquad \delta_\ell\ \text{row}:\ \mathrm{FT.112}\ \text{value cited with its interface provenance}\ [\text{Def-H.41 — Tier F.37's;}\ \neg\text{imported here}]
\]

%dsp-level:0
\[
(5):\ \text{each row is the adjudication of its own tag}:\ \mathrm{FT.141}\ (\S\mathrm{A}),
\]

%dsp-level:0
\[
\ \mathrm{FT.142}\ (\S\mathrm{B}),
\]

%dsp-level:0
\[
\ \mathrm{FT.143}\ (\S\mathrm{C}),
\]

%dsp-level:0
\[
\ \mathrm{FT.144}\ (\S\mathrm{D}\text{–}\S\mathrm{E})\ [\text{machine: all sections green}]
\]

%dsp-level:0
\[
(6):\ \text{standing-rule record};
\]

%dsp-level:0
\[
\qquad \text{import scan}:\ \text{the Req lines of this part name}\ \{\text{Def-H.4},\ \text{Def-H.34},\ \text{Def-H.43}\}\ \text{as imports and}\ \{\text{Def-H.41},\ \text{Def-H.42}\}\ \text{as cited-not-imported}\ [\text{machine scan}].\ \blacksquare
\]

\textbf{Phys.} Rung one closes as a sector: sixteen dimensions, every one spoken for by the five pretime aspects; one family channel with its own angle, its own cubic and the universal two-thirds; one chirality, belonging to the whole force complex rather than to any of its members; and one occupancy law that leaves a single sector materialized and its partner empty — not heavy, empty. The digon fixes the informational read at two bits over the single scale; the family ratios remain dimensionless pending the relating equation the ledger still owes. At the next rung the mask grows by exactly one dimension, the flip channel opens with the core charge, and the sector that was vacuum at rung one becomes force-reachable: the completion is a floor, not a ceiling.

\vspace{2pt}\noindent\rule{\textwidth}{0.4pt}\vspace{2pt}

\phantomsection
\section[Tier F.44 — The 3+1 Chiral Asymmetry: The Manifold Locus of…]{\texorpdfstring{Tier F.44 — The 3+1 Chiral Asymmetry: The Manifold Locus of Chirality, the Spatial Provenance of the Channel Angle, and the Emergent-Time Identity}{Tier F.44 - The 3+1 Chiral Asymmetry: The Manifold Locus of Chirality, the Spatial Provenance of the Channel Angle, and the Emergent-Time Identity}}

\begin{theorem}[{\texorpdfstring{FT.156 — The 3+1 asymmetry theorem: one signature datum at two loci, and the manifold structure as the chiral structure}{FT.156 - The 3+1 asymmetry theorem: one signature datum at two loci, and the manifold structure as the chiral structure}}]\label{tag:FT-156}
\textbf{Thm.} On the pair channel $\mathfrak H(\mathbb C^4)$ and the emergent geometry [CT2.6; conditions (ii)–(iii) inherited at the point of use — A §0.2a]:

%dsp-level:0
\[
(1)\quad\textbf{one signature datum, two loci:}
\]

%dsp-level:0
\[
\qquad \mathrm{A3}\ \text{declares the register involution;}\ \mathrm{C13}\ \text{fixes the Lorentz-signature representative:}
\]

%dsp-level:0
\[
\quad T(\sigma_\mu\otimes\sigma_\nu)=\eta^\mu\eta^\nu\,(\sigma_\mu\otimes\sigma_\nu),\quad\eta=(-,+,+,+),
\]

%dsp-level:0
\[
\quad T^2=\mathbb1;
\]

%dsp-level:0
\[
\mathrm{CT1.12}\ \text{derives the signature of the emergent metric:}\quad\operatorname{sig}(g_{\mu\nu})=\operatorname{sig}(\bar g_{\mu\nu})=(-,+,+,+)
\]

%dsp-level:0
\[
\quad [\varkappa\|q\|_{\mathrm{op}}<1\ \text{regime}];
\]

%dsp-level:0
\[
\text{the same}\ (-,+,+,+)\ \text{sign pattern grades the register basis}\ [\mathrm{C3}:\ \mathcal E_{st}]\ \text{and the metric on}\ M\ [\mathrm{CT2.6}]:\ \text{one temporal sign and three spatial signs at each locus}\ [\text{machine: the}\ T\text{-table on the}\ 16\text{-element CT1.2 basis via exact Pauli-tensor decomposition over}\ \mathbb Q(i)];
\]

%dsp-level:0
\[
(2)\quad\textbf{the block census is the component census; the}\ 7/3\ T\textbf{-split is the}\ 3{+}1\ \textbf{split of the dictionary:}
\]

%dsp-level:0
\[
\qquad (7,3,3,3)\ [\mathrm{CT1.1}];
\]

%dsp-level:0
\[
\qquad S\text{-even}=10=7+3\ [\mathrm{CT1.2};\ \mathrm{FT.141}(1)];
\]

%dsp-level:0
\[
S\text{-even}= \text{the reality-preserving sector}\ ([\mathbb S,X]=0\ \text{preserves}\ \vee^2,\ \wedge^2\ [\mathrm{CT0.7}(3)]);
\]

%dsp-level:0
\[
\qquad S\text{-odd}=6=\{J_i\ (3),\ K_i\ (3)\}= \text{the reality-exchanging sector}\ (\text{the seals}\ [\mathrm{CT0.7}(3)\text{–}(4);\ \mathrm{D1.2}]);
\]

%dsp-level:0
\[
\text{within the}\ S\text{-even block:}
\]

%dsp-level:0
\[
\quad T\text{-even}\ 7=\{E_0,\ \Sigma,\ \hat E^{ij}\ (5)\}=\{D^+,\ H,\ G\}\ [\mathrm{CT0.12};\ \mathrm{CT1.10}:\ \mathbf 1\oplus\mathbf 5\oplus\mathbf 1];
\]

%dsp-level:0
\[
\qquad T\text{-odd}\ 3=\{W_i\}=F_W;
\]

%dsp-level:0
\[
\text{the dictionary}\ [\mathrm{D2.1},\ \mathrm{CT1.12}]:
\]

%dsp-level:0
\[
\qquad (a_{00},\ q_{ij})
\]

%dsp-level:0
\[
= (1{+}6)\ T\text{-even coefficients}\ \longmapsto\ \bigl(\Omega=e^{\alpha_\Omega a_{00}},\ \ \bar g_{ij}=\delta_{ij}+\varkappa q_{ij}\bigr)
\]

%dsp-level:0
\[
= \text{lapse/conformal}\ +\ \text{spatial metric};
\]

%dsp-level:0
\[
w_i\ (3)\ T\text{-odd coefficients}\ \longmapsto\ N_i=\beta w_i=\text{shift}=\text{the time–space cross components}\ \bar g_{0i};
\]

%dsp-level:0
\[
\qquad \text{components:}\quad g_{00}:1;\quad g_{ij}:6=\tfrac{3\cdot4}2;\quad g_{0i}:3;\quad\text{total}\ 10=\tfrac{4\cdot5}2\ [\mathrm{CT1.12}];
\]

%dsp-level:0
\[
\text{the temporal direction contributes exactly one component and one coefficient}\ (g_{00}=-\Omega^2\ \text{from the single}\ T\text{-even conformal coefficient}\ a_{00});
\]

%dsp-level:0
\[
\ \text{the spatial triad contributes the symmetric sextuple}\ q_{ij}\ \text{and the three shifts}\ w_i:
\]

%dsp-level:0
\[
\text{the}\ 7/3\ T\text{-split of the}\ S\text{-even block is the operator-side encoding of the}\ (1;6;3)\ \text{component split with signature}\ (-,+,+,+);
\]

%dsp-level:0
\[
\qquad \text{the}\ S\text{-odd block has}\ 0\ \text{metric coordinates}\ [\mathrm{FT.141}(5);\ \mathrm{FT.30}\ \text{dictionary}];
\]

%dsp-level:0
\[
(3)\quad\textbf{the reversal transport:}\ \ \text{the}\ T\text{-grading transports through the dictionary exactly to the time-reversal pullback:}\quad\text{with}\ J
\]

%dsp-level:0
\[
= \operatorname{diag}(-1,1,1,1)\ \text{and the}\ T\text{-action on sector data}\ (a_{00},q,w)
\]

%dsp-level:0
\[
\mapsto (a_{00},q,-w)
\]

%dsp-level:0
\[
[T(E_0)=E_0,\ T(E^{ij})=E^{ij},\ T(W_i)=-W_i;\ \text{machine, entrywise}]:
\]

%dsp-level:0
\[
\qquad \bar g_{\mu\nu}(a_{00},q,-w)=J^{\rho}{}_{\mu}\,\bar g_{\rho\sigma}(a_{00},q,w)\,J^{\sigma}{}_{\nu}
\]

%dsp-level:0
\[
\qquad [\text{machine: exact rational witnesses}];
\]

%dsp-level:0
\[
\text{lapse and spatial metric fixed, shift negated:}\quad\Omega\mapsto\Omega,\quad\bar g_{ij}\mapsto\bar g_{ij},\quad N_i\mapsto-N_i
\]

%dsp-level:0
\[
\qquad \Longleftrightarrow
\]

%dsp-level:0
\[
\qquad g_{00},\ g_{ij}\ \text{fixed},\quad g_{0i}\mapsto-g_{0i}\ [\Omega\ \text{depends on the}\ T\text{-even}\ a_{00}\ \text{alone}\ [\mathrm{D2.1}]];
\]

%dsp-level:0
\[
\text{the signature is reversal-invariant:}\quad\operatorname{sig}\bar g(w)=\operatorname{sig}\bar g(-w)=(-,+,+,+)
\]

%dsp-level:0
\[
\qquad [\text{machine: exact rational}\ \mathrm{LDL}^\top\ \text{inertia on a}\ 12\text{-point witness lattice}\ (q,w);\ \text{Sylvester criterion on}\ \bar g_{ij};\ \det\bar g<0\ \text{via the CT1.12 Schur form}];
\]

%dsp-level:0
\[
(4)\quad\textbf{space is symmetric and reversible:}\ \ \text{the spatial triad rotates innerly:}\quad\operatorname{span}\{W_i\}\cong\mathfrak{su}(2)\ [\mathrm{CT1.4}];
\]

%dsp-level:0
\[
\qquad J,K:\ \text{vectors},\quad\hat E^{ij}:\ \mathbf 5,\quad E_0,\Sigma:\ \mathbf 1\ \ \text{under the adjoint action}\ [\mathrm{CT1.8},\ \mathrm{CT1.10},\ \mathrm{CT1.5}];
\]

%dsp-level:0
\[
\text{every permutation of the three triad indices is realized by an inner unitary that is a rational polynomial in the weak generators alone:}
\]

%dsp-level:0
\[
\text{cyclic:}\quad U=u\otimes u,\ u=\tfrac12\bigl(I_2-i(\sigma_1+\sigma_2+\sigma_3)\bigr)\ [\mathrm{FT.111}]:
\]

%dsp-level:0
\[
\qquad U=\exp\!\bigl(-\tfrac{2\pi i}{3}\,\hat n\cdot W\bigr)=E_0-\tfrac12\mathcal W^2-\tfrac i2\mathcal W,\quad\mathcal W:=W_1{+}W_2{+}W_3,\quad\hat n=\tfrac{(1,1,1)}{\sqrt3},
\]

%dsp-level:0
\[
\text{via the polynomial identity}\ \mathcal W^3=3\mathcal W\ [\operatorname{spec}(\hat n\cdot W)=\{1,0,-1,0\};\ \text{machine}];
\]

%dsp-level:0
\[
\qquad U^3=E_0;
\]

%dsp-level:0
\[
\qquad U X_i U^\dagger=X_{i+1}\ \ \forall X\in\{W,J,K\}\ (\text{indices mod}\ 3)\ [\mathrm{FT.111}(1)\ \text{re-derived, plus the}\ W\text{-row};\ \text{machine}];
\]

%dsp-level:0
\[
\text{transpositions:}\quad U_{ij}=\exp\bigl(-i\pi\,\tfrac{W_i+W_j}{\sqrt2}\bigr)=E_0-(W_i+W_j)^2:
\]

%dsp-level:0
\[
\qquad U_{ij}^2=E_0;
\]

%dsp-level:0
\[
\qquad U_{ij}:\ X_i\leftrightarrow X_j,
\]

%dsp-level:0
\[
\ X_k
\]

%dsp-level:0
\[
\mapsto -X_k\ \ \forall X\in\{W,J,K\}\ [\text{machine: the full conjugation table on}\ \{W,J,K,\hat E\}];
\]

%dsp-level:0
\[
\text{the reversal involutions are exact: on the register}\ [S^2=T^2=\mathbb1\ [\mathrm{A3}]],
\]

%dsp-level:0
\[
\ \text{on configuration data}\ [\langle C,P,T,\omega\text{-pattern}\rangle\cong\mathbb Z_2^4,\ \text{register-disjoint}\ [\mathrm{FT.133}(1)\text{–}(2)]];
\]

%dsp-level:0
\[
(5)\quad\textbf{time is asymmetric and strictly forward:}\ \ \neg\exists\,\tau^{-1}\ [\mathrm{A6}:\ \tau_n\to\tau_{n+1}:\ \rho\mapsto\mathcal T(\rho)];
\]

%dsp-level:0
\[
\qquad |\mathfrak L(\tau_{n+1})|=|\mathfrak L(\tau_n)|+n_{\mathrm{real\_out}}(\tau_n)\ge|\mathfrak L(\tau_n)|\ [\mathrm{A6},\ \mathrm{A7}:\ \mathfrak L(\tau_n)\subseteq\mathfrak L(\tau_{n+1})];
\]

%dsp-level:0
\[
\text{the D1.2 seal domains force strict alternation and the prefix law; every completed nonempty history ends in}\ J^-,
\]

%dsp-level:0
\[
\ \text{and its pure order reversal is inadmissible at the first event}\ [\mathrm{FT.144}(2),\ (4)(\mathrm{ii});\ \text{machine: exhaustive at lengths}\le8];
\]

%dsp-level:0
\[
\Theta_{\mathfrak L}\ [\mathrm{D2.7}]\ \text{is the formal mirror — an involution on configuration data bijecting the completed admissible set onto itself}\ [\mathrm{FT.144}(4)(\mathrm{i});\ \text{machine}]\ —\ \neg\text{a dynamical inverse:}
\]

%dsp-level:0
\[
\text{the asymmetry is exact: the spatial-axis reversals are inner unitaries on the force arena}\ [(4)];
\]

%dsp-level:0
\[
\qquad \text{the temporal advance reversal has no implementation}\ [\mathrm{A6};\ \mathrm{FT.144}(3):\ \text{a domain restriction, prior to every read}];
\]

%dsp-level:0
\[
(6)\quad\textbf{the chiral structure is the manifold structure:}\ \ \text{the chirality of the complete force complex at rung one is the}\ 10/6\ S\text{-grading}\ [\mathrm{FT.141}(2)];
\]

%dsp-level:0
\[
\ \text{its}\ (S,T)\ \text{refinement}\ (7,3,3,3)\ [\mathrm{CT1.1}]\ \text{is the dictionary carrier of the metric}\ [\mathrm{CT1.12};\ (2)];
\]

%dsp-level:0
\[
\text{the weak block is the}\ T\text{-odd part of the}\ S\text{-even block}
\]

%dsp-level:0
\[
= \text{the metric's own time–space cross block}\ [\mathrm{CT0.12}:\ F_W=\mathcal E_{+-}=\operatorname{span}\{W_i\};\ \mathrm{D2.1}:\ X_{(-)}=w_iW_i;\ \mathrm{CT1.12}:\ w_i\mapsto N_i];
\]

%dsp-level:0
\[
\text{the chiral asymmetry}\ (10\ne6)\ \text{and the}\ 3{+}1\ \text{asymmetry are one graded structure read at three derived loci:}
\]

%dsp-level:0
\[
\text{the operator locus}\ (10,6)=(7{+}3,\ 3{+}3);
\]

%dsp-level:0
\[
\qquad \text{the component locus}\ (1;6;3)\ \text{with}\ \operatorname{sig}=(-,+,+,+);
\]

%dsp-level:0
\[
\qquad \text{the state locus}\ \operatorname{spec}\mathbb S=(3,1)\ [\mathrm{CT0.7}(1);\ \mathrm{FT.157}];
\]

%dsp-level:0
\[
\text{locus discipline}\ [\mathrm{FT.136}(5);\ \mathrm{FT.141}(4)]:\ \text{three distinct spaces}\ (\operatorname{End}(\mathbb C^4_{\mathrm{pair}}),\ \text{the metric on}\ M,\ \mathbb C^4_{\mathrm{pair}});
\]

%dsp-level:0
\[
\ \text{the counts are the derived content; no carrier identification is made}.
\]

\textbf{Rmk (label discipline).} The labels \textit{spatial}/\textit{temporal} of this tier are internal bookkeeping attached to derived T-parities and dictionary roles [CT1.2's table; CT1.12's table]; they name no particle and no external structure [P1, P5]. Naming either branch \textit{left} or \textit{right} remains barred [FT.101 Phys bars stand]. All topology and branch labels are internal bookkeeping; no H-sector interpretation enters any clause.

\textbf{Req.} A3, A6, A7, A8 (Part 1); C3 (Part 1); D1.2, D2.1, D2.7 (Part 1); CT0.7, CT0.12, CT1.1, CT1.2, CT1.4, CT1.5, CT1.8, CT1.10 (Part 1); CT1.12, CT2.6 (Part 2; CT2.6's conditions (ii)–(iii) at the point of use — A §0.2a); FT.30 (Tiers F.1–F.11); FT.101, FT.111 (Tiers F.34, F.37); FT.133 (Tier F.41); FT.136, FT.141, FT.144 (Tiers F.41–F.42); exact $\mathbb Q(i)$ matrix arithmetic (Pauli-tensor decomposition, conjugation tables, polynomial rotation identities), exact rational $\mathrm{LDL}^\top$ inertia, Sylvester/Schur witnesses, exhaustive ledger-word enumeration [math; machine: \texttt{checks/\allowbreak{}verify\_\allowbreak{}H\_\allowbreak{}3p1\_\allowbreak{}chiral\_\allowbreak{}asymmetry.py} §A].

\textbf{Pf.}

%dsp-level:0
\[
(1):\ \eta=(-,+,+,+)\ \text{verbatim}\ [\mathrm{A3},\ \mathrm{C13}];
\]

%dsp-level:0
\[
\qquad \operatorname{sig}(g)=(-,+,+,+)\ \text{verbatim}\ [\mathrm{CT1.12}:\ \text{Schur}\ +\ \text{Sylvester}\ +\ \text{conformal covariance}];
\]

%dsp-level:0
\[
\qquad T(X)
\]

%dsp-level:0
\[
= \sum_{\mu,\nu}c_{\mu\nu}\,
\]

%dsp-level:0
\[
\eta^\mu\eta^\nu\,
\]

%dsp-level:0
\[
\sigma_\mu\otimes\sigma_\nu,
\]

%dsp-level:0
\[
\ \ c_{\mu\nu}
\]

%dsp-level:0
\[
= \tfrac14\mathfrak T\bigl((\sigma_\mu\otimes\sigma_\nu)X\bigr)\ [\text{exact decomposition, unique on the orthogonal Pauli-tensor basis}\ [\mathrm{CT1.1}\ \text{Pf pattern}];\ \text{machine on the}\ 16\text{-basis}]
\]

%dsp-level:0
\[
(2):\ \mathrm{CT1.2}\ \text{table verbatim}: (S,T)\ \text{grades}\ (1{+}1{+}5,\ 3;\ 3;\ 3);
\]

%dsp-level:0
\[
\qquad 7=1{+}1{+}5\ [\mathrm{CT1.10}];
\]

%dsp-level:0
\[
\qquad \text{dictionary rows verbatim}\ [\mathrm{CT1.12}\ \text{table}]:\ a_{00}\mapsto\Omega,\ \Sigma_{\mathrm{tr}}\mapsto\text{volume},\ \hat h_{ij}\mapsto\text{shear},\ w_i\mapsto\text{shift}
\]

%dsp-level:0
\[
q_{ij}=q_{ji}:\ 6\ \text{components}\ [\mathrm{D2.1}];
\]

%dsp-level:0
\[
\qquad 1+6+3=10=\tfrac{4\cdot5}2\ [\mathrm{CT1.12}];
\]

%dsp-level:0
\[
\qquad g_{00}=\Omega^2\bar g_{00}=-\Omega^2\ [\mathrm{D2.1}:\ \bar g_{00}=-1];
\]

%dsp-level:0
\[
\qquad S\text{-odd}:\ 0\ \text{metric coordinates}\ [\mathrm{FT.141}(5);\ \mathrm{FT.30}];
\]

%dsp-level:0
\[
\qquad \text{reality-preserving/exchanging}:\ \mathrm{CT0.7}(3)\ \text{eigen-relations verbatim}
\]

%dsp-level:0
\[
(3):\ \bar g_{0i}=\varkappa\beta w_i\ [\mathrm{D2.1}:\ N_i=\beta w_i]\implies\bar g_{0i}(a_{00},q,-w)=-\bar g_{0i}(a_{00},q,w);
\]

%dsp-level:0
\[
\qquad \bar g_{00}=-1,\ \bar g_{ij}=\delta_{ij}+\varkappa q_{ij}\ \text{are}\ w\text{-free}\implies\text{fixed};
\]

%dsp-level:0
\[
\qquad (J^\top\bar g J)_{00}
\]

%dsp-level:0
\[
= \bar g_{00},
\]

%dsp-level:0
\[
\ (J^\top\bar g J)_{ij}
\]

%dsp-level:0
\[
= \bar g_{ij},
\]

%dsp-level:0
\[
\ (J^\top\bar g J)_{0i}
\]

%dsp-level:0
\[
= -\bar g_{0i}\ [\text{machine: exact equality on the witness lattice}]
\]

%dsp-level:0
\[
\text{inertia}: \mathrm{LDL}^\top\ \text{over}\ \mathbb Q\ \text{on}\ \bar g(q,w):\ D\ \text{signs}\ (-,+,+,+)\ \text{on all}\ 12\ \text{witnesses and at}\ w
\]

%dsp-level:0
\[
\mapsto -w\ [\text{machine}];
\]

%dsp-level:0
\[
\qquad \det\bar g
\]

%dsp-level:0
\[
= \det(\bar g_{ij})\cdot\bigl(-1-\varkappa^2N_i(\bar g^{-1})^{ij}N_j\bigr)<0\ [\mathrm{CT1.12}\ \text{Pf re-run};\ \text{machine: exact}\ 4\times4\ \text{determinants}]
\]

%dsp-level:0
\[
(4):\ [W_i,W_j]=i\epsilon_{ijk}W_k\ [\mathrm{CT1.4}\ \text{re-derived, all pairs};\ \text{machine}];
\]

%dsp-level:0
\[
\qquad [W_i,J_j]=i\epsilon_{ijk}J_k,\quad[W_i,K_j]=i\epsilon_{ijk}K_k\ [\mathrm{CT1.8}\ \text{re-derived, all}\ 9{+}9\ \text{pairs};\ \text{machine}];
\]

%dsp-level:0
\[
\qquad [W_a,\hat E^{ij}]=i(\epsilon_{aik}\hat E^{kj}+\epsilon_{ajk}\hat E^{ik}),\quad[W_i,\Sigma]=0\ [\mathrm{CT1.10},\ \mathrm{CT1.5}\ \text{re-derived};\ \text{machine}]
\]

%dsp-level:0
\[
\mathcal W^3
\]

%dsp-level:0
\[
= 3\mathcal W\ [\text{machine}]
\]

%dsp-level:0
\[
\implies X^3
\]

%dsp-level:0
\[
= X\ \text{at}\ X
\]

%dsp-level:0
\[
= \hat n\cdot W
\]

%dsp-level:0
\[
= \mathcal W/\sqrt3\implies\exp(-i\theta X)
\]

%dsp-level:0
\[
= E_0+(\cos\theta-1)X^2-i\sin\theta\,
\]

%dsp-level:0
\[
X\ [\text{polynomial identity on}\ \operatorname{spec}\{1,0,-1\}];
\]

%dsp-level:0
\[
\qquad \theta=\tfrac{2\pi}3:\ \cos\theta-1=-\tfrac32,\ \ \sin\theta\cdot\tfrac1{\sqrt3}=\tfrac12\implies U=E_0-\tfrac12\mathcal W^2-\tfrac i2\mathcal W
\]

%dsp-level:0
\[
U=u\otimes u\ \text{entrywise}\ [\mathrm{FT.111}\ \text{form};\ \text{machine}];
\]

%dsp-level:0
\[
\qquad U^3=E_0,\ U U^\dagger=E_0,\ \det u=1\ [\text{machine}];
\]

%dsp-level:0
\[
\qquad U X_i U^\dagger=X_{i+1},\ X\in\{W,J,K\}\ [\text{machine, all}\ 9];
\]

%dsp-level:0
\[
\qquad U\hat E^{ii}U^\dagger=\hat E^{i+1,i+1}\ [\text{machine}]
\]

%dsp-level:0
\[
U_{ij}=E_0-(W_i+W_j)^2\ [\theta=\pi:\ \cos\pi-1=-2,\ X=(W_i{+}W_j)/\sqrt2,\ X^3=X];
\]

%dsp-level:0
\[
\qquad U_{ij}^2=E_0,\ U_{ij}U_{ij}^\dagger=E_0\ [\text{machine}];
\]

%dsp-level:0
\[
\qquad \text{conjugation tables: all}\ 3\ \text{transpositions}\times\{W,J,K\}\ \text{and the}\ \hat E\text{-action}\ [\text{machine, exact}]
\]

%dsp-level:0
\[
(5):\ \mathrm{A6}/\mathrm{A7}\ \text{verbatim};
\]

%dsp-level:0
\[
\qquad \text{prefix law, uniqueness of}\ (J^+J^-)^m,
\]

%dsp-level:0
\[
\ \text{inadmissibility of pure reversals}\ [\mathrm{FT.144}(2),\ (4)\ \text{machine re-run: exhaustive at lengths}\le8];
\]

%dsp-level:0
\[
\qquad \Theta_{\mathfrak L}^2=\mathrm{id},\ \text{bijection on the completed set}\ [\mathrm{FT.144}(4)(\mathrm{i})\ \text{re-run};\ \text{machine}]
\]

%dsp-level:0
\[
(6):\ \text{conjunction of}\ (1)\text{–}(5)\ \text{with}\ \mathrm{FT.141}(2),\ \mathrm{CT0.12},\ \mathrm{CT0.7}(1);
\]

%dsp-level:0
\[
\qquad \text{locus rows verbatim}\ [\mathrm{FT.136}(5);\ \mathrm{FT.141}(4)].\ \blacksquare
\]

\textbf{Phys.} The corpus declares one sign pattern, $(-,+,+,+)$, at the register [A3's $T$; C13] and derives the same sign pattern for the emergent metric [CT1.12] — one datum at two loci. The ten reality-preserving dimensions split $7/3$ under that grading, and the split is exactly the metric's: the seven T-even coefficients build the lapse and the spatial metric, the three T-odd coefficients build the shift — the time–space cross terms. The grading transports through the dictionary as the time-reversal pullback: lapse and spatial metric stand, the shift flips, the signature stands. Space shows its symmetry as inner rotations — every triad permutation is a rational polynomial in the weak generators — and its reversibility as exact involutions; time shows its asymmetry as the absence of any inverse for the tick advance, the monotone ledger, and the inadmissibility of every reversed history. Chirality is the same structure: the $10/6$ split of the complete force complex, refined by the signature grading, with the weak block sitting precisely in the metric's time–space sector.

\vspace{2pt}\noindent\rule{\textwidth}{0.4pt}\vspace{2pt}

\end{theorem}

\begin{theorem}[{\texorpdfstring{FT.157 — The weak-selectivity theorem: the weak force is the shift block — the spatial force on the spatial branch}{FT.157 - The weak-selectivity theorem: the weak force is the shift block - the spatial force on the spatial branch}}]\label{tag:FT-157}
\textbf{Thm.} With $\mathbb S$, $\Pi_{\mathcal R}=\tfrac14(3E_0+\Sigma)$, $\Pi_{\mathcal V}=\tfrac14(E_0-\Sigma)$ [CT0.7(1)], $W_i=\tfrac12(\sigma_0\otimes\sigma_i+\sigma_i\otimes\sigma_0)$ [CT1.2], $\Phi=\{0,\pi\}$ [D0.5, CT0.9]:

%dsp-level:0
\[
(1)\quad\textbf{the weak block is the shift block:}
\]

%dsp-level:0
\[
\qquad F_W=\mathbf F_{(-,\mathcal R)}=\mathcal E_{+-}=\operatorname{span}\{W_i\}\quad[\mathrm{A8};\ \mathrm{CT0.7}(4);\ \mathrm{C8}:\ F_W=F_{W^+}\oplus F_{W^-}\oplus F_Z];
\]

%dsp-level:0
\[
\text{the same three generators carry the metric shift}\ [\mathrm{D2.1}:\ X_{(-)}=w_iW_i;\ \mathrm{CT1.12}:\ w_i\mapsto N_i]:
\]

%dsp-level:0
\[
\qquad g_{0i}=\varkappa\beta\,\Omega^2\,w_i\ [\mathrm{D2.1}]:
\]

%dsp-level:0
\[
\text{the weak aspect is the metric's time–space coupling block — one}\ 3\text{-dimensional block, two derived readings}\ [\text{machine: sector membership of every basis element}\ +\ \text{the dictionary row count}];
\]

%dsp-level:0
\[
(2)\quad\textbf{the weak acts on the spatial orientations and has no temporal one to act on:}
\]

%dsp-level:0
\[
\qquad W_i=\tfrac12\bigl(\sigma_i^{(p)}+\sigma_i^{(q)}\bigr)=\text{total spin}\ [\mathrm{CT1.2};\ \mathrm{FT.13}];
\]

%dsp-level:0
\[
\qquad [W_i,W_j]=i\epsilon_{ijk}W_k\ [\mathrm{CT1.4}];
\]

%dsp-level:0
\[
J,K:\ \text{vectors under}\ \operatorname{span}\{W\}\ [\mathrm{CT1.8}];\qquad\text{the triad index}\ i\ \text{is the rotation index}\ [\mathrm{FT.156}(4)];
\]

%dsp-level:0
\[
\textstyle\sum_iW_i^2=\tfrac32E_0+\tfrac12\Sigma=E_0+\mathbb S=2\,\Pi_{\mathcal R}\ [\mathrm{Cor}\ \mathrm{CT0.7a};\ \mathrm{FT.141}(4)]:
\]

%dsp-level:0
\[
\qquad \text{Casimir}\ 2=1(1{+}1)\ \text{on}\ \vee^2\ —\ \text{the spin-1 (vector) read};
\]

%dsp-level:0
\[
W_i\Pi_{\mathcal V}=0=\Pi_{\mathcal V}W_i\ \forall i,
\]

%dsp-level:0
\[
\qquad \bigcap_i\ker W_i=\operatorname{im}\Pi_{\mathcal V}\ [\mathrm{FT.101}(1)\text{–}(2)\ \text{re-derived};\ \text{machine}]:
\]

%dsp-level:0
\[
\qquad \text{Casimir}\ 0\ \text{on}\ \wedge^2\ —\ \text{the spin-0 (scalar) read, no orientation};
\]

%dsp-level:0
\[
\text{the derived dimension match}\ [\mathrm{CT0.7}(1):\ \mathfrak T\Pi_{\mathcal R}=3=\dim\vee^2,\ \mathfrak T\Pi_{\mathcal V}=1=\dim\wedge^2;\ \mathrm{CT1.12}:\ 3\ \text{positive},\ 1\ \text{negative eigenvalue}]:
\]

%dsp-level:0
\[
\textbf{spatial branch}\ :=\ \text{the aligned branch}\ (\alpha=0,\ \vee^2,\ \text{Casimir}\ 2)\ —\ \text{weak-active};
\]

%dsp-level:0
\[
\qquad \textbf{temporal branch}\ :=\ \text{the antialigned branch}\ (\alpha=\pi,\ \wedge^2,\ \text{Casimir}\ 0)\ —\ \text{weak-blind}
\]

%dsp-level:0
\[
\qquad [\mathrm{A4};\ \mathrm{FT.101}];
\]

%dsp-level:0
\[
\text{the labels record the derived matches}\ \dim\vee^2
\]

%dsp-level:0
\[
= 3
\]

%dsp-level:0
\[
= \#\{\text{spatial directions}\}\ \text{and}\ \dim\wedge^2
\]

%dsp-level:0
\[
= 1
\]

%dsp-level:0
\[
= \#\{\text{the temporal direction}\}\ [\mathrm{FT.156}(2)\text{–}(3)];
\]

%dsp-level:0
\[
\ \text{they are internal bookkeeping}\ [\text{Rmk below}];
\]

%dsp-level:0
\[
(3)\quad\textbf{the selectivity is the}\ 3{+}1\ \textbf{asymmetry at the state locus:}
\]

%dsp-level:0
\[
\qquad \text{the weak sees the spatial branch and annihilates the temporal branch}\ [\mathrm{FT.101}(1)\text{–}(2);\ \text{machine}];
\]

%dsp-level:0
\[
\text{its support is fixed by the}\ S\text{-grading}\ (\mathbb S\ \text{enters the Casimir identity literally}\ [\mathrm{FT.141}(4)])\ \text{and refined by the}\ T\text{-grading: the weak support is the}\ S\text{-even}\ T\text{-odd block}
\]

%dsp-level:0
\[
= \text{the sector coupling the temporal register direction to the spatial triad}\ [\mathrm{FT.156}(2)\text{–}(3)];
\]

%dsp-level:0
\[
\text{the branch flip}\ \iota(\alpha)
\]

%dsp-level:0
\[
= \pi-\alpha\ \text{has no Hilbert-space implementation}\ (\mathfrak T\Pi_{\mathcal R}=3\ne1=\mathfrak T\Pi_{\mathcal V};\ \operatorname{rk}3\ne1)\ [\mathrm{FT.101}(3)\ \text{re-derived}]:\ \text{the one-sidedness is exact on the pair register};
\]

%dsp-level:0
\[
(4)\quad\textbf{the carrier-side ledger, unchanged:}\ \ \text{on the extended carrier the weak acts seed}\otimes\text{trivial and commutes with the chiral projectors}\ [\mathrm{FT.126}(2):\ \text{blindness},\ \neg\text{coupling}];
\]

%dsp-level:0
\[
\text{the four gradings — the algebra}\ S\text{-grading, the branch-label parity, the carrier}\ \hat\chi\text{-grade, the Janus arrow — remain pairwise distinct}\ [\mathrm{FT.136};\ \mathrm{FT.141}(6);\ \mathrm{FT.144}(5)];
\]

%dsp-level:0
\[
\text{this tag adds the manifold-side read of the branch-label}\ \mathbb Z_2\ (\text{spatial}=\text{aligned}=\text{weak-active};\ \text{temporal}=\text{antialigned}=\text{weak-blind}):\ \text{a force-side input to the deficit}\ 1'\ \text{assignment rule}\ [\mathrm{FT.125}(4);\ \mathrm{FT.141}(6)];
\]

%dsp-level:0
\[
\text{the rule coupling the gradings to species is}\ \neg\text{constructed here; no species is assigned};
\]

%dsp-level:0
\[
(5)\quad\textbf{the one-sidedness, answered:}
\]

%dsp-level:0
\[
\qquad \text{space:}\ 3\text{-dimensional}\ [\mathrm{CT1.12}:\ \text{three positive eigenvalues};\ \mathrm{FT.156}(2):\ \text{the triad counts}],
\]

%dsp-level:0
\[
\ \text{triad-symmetric}\ [\mathrm{FT.156}(4):\ S_3\ \text{inner}],
\]

%dsp-level:0
\[
\ \text{reversible under exact involutions}\ [\mathrm{A3};\ \mathrm{FT.133}(1)];
\]

%dsp-level:0
\[
\text{time:}\ 1\text{-dimensional}\ [\mathrm{CT1.12}:\ \text{the single negative eigenvalue}]\ \text{and strictly forward}\ [\mathrm{A6}:\ \neg\exists\,\tau^{-1};\ \mathrm{FT.144}(2)\text{–}(3):\ \text{no backward materialization}];
\]

%dsp-level:0
\[
\text{the weak coupling is the triad-indexed vector coupling}\ [\mathrm{CT1.2}:\ W_i,\ i=1,2,3;\ \mathrm{CT1.8}]:\ \text{it couples to the reversible}\ 3\text{-dimensional sector, and the}\ 1\text{-dimensional strictly-forward sector carries no orientation object for it to act on}\ [\wedge^2:\ \text{spin}\ 0;\ W_i\Pi_{\mathcal V}=0]:
\]

%dsp-level:0
\[
\text{the chiral asymmetry of the weak is the}\ 3{+}1\ \text{asymmetry of the arena}.
\]

\textbf{Rmk (naming and identification bars).} The labels \textit{spatial branch}/\textit{temporal branch} attach to the A4 branch and swap-eigenspace loci only — internal bookkeeping for the derived dimension and Casimir matches. Naming either branch \textit{left} or \textit{right} is barred [FT.101 Phys bars stand]; the carrier-ray names $P_L/P_R$ remain algebraic up to the application-layer assignment rule [Def-H.33 canonicity Rmk]; no SM species reading; no claim that the SM weak interaction is reproduced [P1, P5].

\textbf{Req.} A3, A4, A6, A8 (Part 1); C8 (Part 1); D0.5, D2.1 (Part 1); CT0.7, Cor CT0.7a, CT0.9, CT0.12, CT1.2, CT1.4, CT1.8 (Part 1); CT1.12 (Part 2); FT.13 (Tiers F.1–F.11); FT.101 (Tier F.34); FT.125, FT.126 (Tier F.39); FT.133 (Tier F.41); FT.136, FT.141, FT.144 (Tiers F.41–F.42); Def-H.33, Def-H.34 cited within the frozen FT.126/FT.144 clauses; FT.156 (this tier); exact $\mathbb Q(i)$ arithmetic, rank/trace invariants, Casimir witnesses on exact rational states [math; machine: \texttt{checks/\allowbreak{}verify\_\allowbreak{}H\_\allowbreak{}3p1\_\allowbreak{}chiral\_\allowbreak{}asymmetry.py} §B].

\textbf{Pf.}

%dsp-level:0
\[
(1):\ \mathrm{A8}/\mathrm{CT0.7}(4)\ \text{block rows verbatim};
\]

%dsp-level:0
\[
\qquad \mathrm{D2.1}:\ X_{(-)}=w_iW_i;
\]

%dsp-level:0
\[
\qquad \mathrm{CT1.12}\ \text{table row}\ w_i\mapsto N_i=\beta w_i;
\]

%dsp-level:0
\[
\qquad g_{0i}=\Omega^2\bar g_{0i}=\varkappa\beta\Omega^2 w_i\ [\mathrm{D2.1}];
\]

%dsp-level:0
\[
\qquad \text{machine:}\ W_i\in\mathcal E_{+-}\ \text{on the}\ T\text{-table}\ [\mathrm{FT.156}(1)],\ \#\{T\text{-odd}\ S\text{-even basis elements}\}=3
\]

%dsp-level:0
\[
(2):\ \mathrm{CT1.2}/\mathrm{FT.13},\ \mathrm{CT1.4},\ \mathrm{CT1.8}\ \text{verbatim};
\]

%dsp-level:0
\[
\qquad \textstyle\sum_iW_i^2=\tfrac32E_0+\tfrac12\Sigma\ [\mathrm{Cor}\ \mathrm{CT0.7a}\ \text{re-derived};\ \text{machine}];
\]

%dsp-level:0
\[
\qquad \Sigma=2\mathbb S-E_0\ [\mathrm{CT0.7}(1)]\implies\tfrac32E_0+\tfrac12\Sigma=E_0+\mathbb S=2\Pi_{\mathcal R}\ [\text{machine}]
\]

%dsp-level:0
\[
\vee^2:\ \mathbb S=+1\implies\text{Casimir}\ 2;
\]

%dsp-level:0
\[
\qquad \wedge^2:\ \Sigma|_{\wedge^2}=-3\implies\tfrac32-\tfrac32=0\ [\mathrm{Cor}\ \mathrm{CT0.7a}];
\]

%dsp-level:0
\[
\qquad W_i\Pi_{\mathcal V}=0\ \forall i,\quad\dim\bigcap_i\ker W_i=4-\operatorname{rk}\binom{W_1}{W_2}{W_3}=1\ [\mathrm{FT.101}(2)\ \text{re-derived};\ \text{machine}];
\]

%dsp-level:0
\[
\qquad \mathfrak T\Pi_{\mathcal R}=3,\ \mathfrak T\Pi_{\mathcal V}=1\ [\mathrm{CT0.7}(1);\ \text{machine}]
\]

%dsp-level:0
\[
(3):\ \mathrm{FT.101}(1)\text{–}(3)\ \text{re-derived witnesses};
\]

%dsp-level:0
\[
\qquad \text{support identity}\ W_i=\Pi_{\mathcal R}W_i\Pi_{\mathcal R}\ [\text{machine}];
\]

%dsp-level:0
\[
\qquad \text{flip non-implementation: conjugation preserves}\ \mathfrak T\ \text{and}\ \operatorname{rk};\ 3\ne1\ [\text{machine}]
\]

%dsp-level:0
\[
(4):\ \mathrm{FT.126}(2),\ \mathrm{FT.136},\ \mathrm{FT.141}(6),\ \mathrm{FT.144}(5)\ \text{verbatim (frozen)};
\]

%dsp-level:0
\[
\qquad \neg\exists\ \text{assignment tag}\ [\text{inventory scan};\ \mathrm{FT.125}(4)]
\]

%dsp-level:0
\[
(5):\ \text{rows cite}\ \mathrm{CT1.12},\ \mathrm{FT.156}(2)\text{–}(5),\ \mathrm{A6},\ \mathrm{FT.144}(2)\text{–}(3),\ \mathrm{CT1.2},\ \mathrm{CT1.8};
\]

%dsp-level:0
\[
\qquad W_i\Pi_{\mathcal V}=0\ [(2);\ \text{machine}].\ \blacksquare
\]

\textbf{Phys.} The three generators the corpus calls weak are the same three coefficients the metric calls shift: the weak force is literally the time–space coupling block of the emergent geometry — a spatial force, triad-indexed, rotating everything it touches as vectors. On states, it lives entirely on the three-dimensional aligned branch — Casimir $2$, the spin-1 read — and annihilates the one-dimensional antialigned branch — Casimir $0$, no orientation to act on. Space is three-dimensional, fully symmetric under inner triad rotations, and reversible under exact involutions; time is one-dimensional and strictly forward, with no inverse for its advance. The weak couples to the spatial orientations because the spatial sector is where its vector action has a referent; it is blind to the temporal branch because the temporal sector has none. That is the chiral selectivity of the force — the $3{+}1$ asymmetry read on states — while the carrier-side gradings stay distinct and the species rule stays open.

\vspace{2pt}\noindent\rule{\textwidth}{0.4pt}\vspace{2pt}

\end{theorem}

\begin{theorem}[{\texorpdfstring{FT.158 — The spatial provenance of the channel angle: $C_3$ is the weak rotation of the triad, and the Koide-form spectrum is forced}{FT.158 - The spatial provenance of the channel angle: is the weak rotation of the triad, and the Koide-form spectrum is forced}}]\label{tag:FT-158}
\textbf{Thm.} With the seal bilinear $[J_i,K_j]$ [FT.109], the $C_3$-covariant spectrum [FT.111], the mixing angle $\delta=\tfrac29$ [FT.112], and the bilinear spectrum $r_i$ [FT.113]:

%dsp-level:0
\[
(1)\quad\textbf{the}\ C_3\ \textbf{is generated by the weak generators:}\ \ \mathrm{FT.111}\text{'s inner rotation}\ U
\]

%dsp-level:0
\[
= u\otimes u,
\]

%dsp-level:0
\[
\ u
\]

%dsp-level:0
\[
= \tfrac12\bigl(I_2-i(\sigma_1+\sigma_2+\sigma_3)\bigr),
\]

%dsp-level:0
\[
\ \text{is the}\ \tfrac{2\pi}3\ \text{rotation about the}\ (1,1,1)\ \text{axis generated by the total-spin/weak generators:}
\]

%dsp-level:0
\[
U=\exp\!\bigl(-\tfrac{2\pi i}3\,\hat n\cdot W\bigr)=E_0-\tfrac12\mathcal W^2-\tfrac i2\mathcal W,
\]

%dsp-level:0
\[
\qquad \mathcal W=W_1{+}W_2{+}W_3,
\]

%dsp-level:0
\[
\qquad \hat n=\tfrac{(1,1,1)}{\sqrt3}
\]

%dsp-level:0
\[
\qquad [\text{exact polynomial identity over}\ \mathbb Q(i),\ \text{via}\ \mathcal W^3=3\mathcal W;\ \text{machine}];
\]

%dsp-level:0
\[
U^3=E_0;
\]

%dsp-level:0
\[
\qquad U J_i U^\dagger=J_{i+1},\quad U K_i U^\dagger=K_{i+1},\quad U W_i U^\dagger=W_{i+1}\quad(\text{indices mod}\ 3)
\]

%dsp-level:0
\[
\qquad [\mathrm{FT.111}(1)\ \text{re-derived, plus the}\ W\text{-row};\ \text{machine}]:
\]

%dsp-level:0
\[
C_3
\]

%dsp-level:0
\[
\le SO(3)\ —\ \text{the cyclic permutation of the three spatial directions}\ [\mathrm{CT1.8}:\ \text{the triad}\ SO(3);\ \mathrm{CT1.4}:\ \operatorname{span}\{W\}\cong\mathfrak{su}(2)]\ —\ \text{is inner in the weak rotation algebra};
\]

%dsp-level:0
\[
(2)\quad\textbf{the numerator}\ 2=3-1\ \textbf{is the doublet dimension of the spatial triad:}
\]

%dsp-level:0
\[
\qquad \text{the diagonal bilinear splits singlet}\ \oplus\ \text{doublet}\ [\mathrm{FT.111}(1)]:\ \ \text{the fixed direction}\ \Sigma\ \oplus\ \text{the zero-sum plane}\ \textstyle\sum_i\hat E^{ii}
\]

%dsp-level:0
\[
= 0,
\]

%dsp-level:0
\[
\ \dim
\]

%dsp-level:0
\[
= 2\ [\text{machine: exact rank}];
\]

%dsp-level:0
\[
\text{the cyclic shift on the plane:}\quad\operatorname{charpoly}
\]

%dsp-level:0
\[
= \text{minimal polynomial}
\]

%dsp-level:0
\[
= x^2+x+1,
\]

%dsp-level:0
\[
\ \ \text{irreducible over}\ \mathbb Q\ (\pm1\mapsto3\ne0)\ [\mathrm{FT.111}(1)\ \text{re-derived};\ \text{machine}];
\]

%dsp-level:0
\[
\textstyle\sum_i[J_i,K_i]=i\sum_i\bigl(\Sigma-E^{ii}\bigr)=i\bigl(3\Sigma-\Sigma\bigr)=2i\Sigma\qquad[\mathrm{FT.109}(1)\text{–}(3)\ \text{re-derived};\ \text{machine}]:
\]

%dsp-level:0
\[
\text{the coefficient}\ 2
\]

%dsp-level:0
\[
= 3-1\ —\ \text{the triad count minus its trace}: \text{one integer at two derived reads}\ (\dim\ \text{doublet};\ \Sigma\text{-coefficient})\ [\text{exact integers}];
\]

%dsp-level:0
\[
(3)\quad\textbf{the denominator}\ 9=3\times3;\ \ \delta=\tfrac29\ \textbf{is a ratio of spatial counts:}
\]

%dsp-level:0
\[
\qquad \dim F_{J^\pm}
\]

%dsp-level:0
\[
= 3\ [\mathrm{CT0.12};\ \text{machine: exact rank}]\ —\ \text{the seal-triplet dimension}
\]

%dsp-level:0
\[
= \text{the triad dimension on which}\ SO(3)\ \text{acts}\ [\mathrm{CT1.8}];
\]

%dsp-level:0
\[
\operatorname{tr}T(1)
\]

%dsp-level:0
\[
= 3\ [\mathrm{FT.110}(2):\ \operatorname{charpoly}T(1)=x^4-3x^3-3x^2;\ \text{machine: exact re-derivation from the transfer rule}\ T(1)[s][s']=[\,|s|+|s'|\le3\,]];
\]

%dsp-level:0
\[
\delta=\frac{\sum_i[J_i,K_i]\big|_\Sigma}{\dim F_{J^\pm}\cdot\operatorname{tr}T(1)}=\frac{2}{3\cdot3}=\frac29
\]

%dsp-level:0
\[
\qquad [\mathrm{FT.112}\ \text{re-derived};\ \text{machine: exact Fractions}];
\]

%dsp-level:0
\[
\text{the arena reading carries Tier F.37's interface provenance}\ [\mathrm{Def\text{-}H.41}\ —\ \text{cited},\ \neg\text{imported here}]:
\]

%dsp-level:0
\[
\qquad \text{the angle's constituents are the doublet dimension, the seal dimension, the transfer trace — integer counts of the spatial channel, no datum, no fit};
\]

%dsp-level:0
\[
(4)\quad\textbf{the Koide-form spectrum is forced:}
\]

%dsp-level:0
\[
\qquad r_i
\]

%dsp-level:0
\[
= 1+\sqrt2\cos\bigl(\tfrac29+\tfrac{2\pi i}3\bigr)\ [\mathrm{FT.113}]:\ \ \text{the shape from}\ C_3\text{-covariance}\ [\mathrm{FT.111}(2):\ \lambda_i=A+B\cos(\theta+\tfrac{2\pi i}3)\ \text{for every}\ C_3\text{-covariant Hermitian triplet operator}];
\]

%dsp-level:0
\[
A = 1,
\]

%dsp-level:0
\[
\ B
\]

%dsp-level:0
\[
= \sqrt2\ \text{from the democratic normalization to the seal dimension}\ 3\ [\mathrm{FT.111}(3):\ \textstyle\sum_i\lambda_i=3=\dim F_{J^\pm},\ \sum_i(\lambda_i-A)^2=3];
\]

%dsp-level:0
\[
\qquad \theta=\delta\ \text{from}\ (2)/(3);
\]

%dsp-level:0
\[
\text{machine:}\quad\textstyle\sum_ir_i=3,\quad\sum_{i<j}r_ir_j=\tfrac32,\quad 2(r_i-1)^3-3(r_i-1)-\sqrt2\cos\tfrac23=0\ \ \forall i
\]

%dsp-level:0
\[
\qquad [\text{exact in}\ \mathbb Q(\sqrt2,\sqrt3)[c,s]/(c^2{+}s^2{-}1);\ \cos\tfrac23=4c^3-3c\ \text{at}\ c=\cos\tfrac29];
\]

%dsp-level:0
\[
\text{the Koide formula is a consequence of the}\ C_3\ \text{symmetry of the}\ 3\ \text{spatial directions};
\]

%dsp-level:0
\[
\qquad \text{the scale}\ \mu\ (\text{P4-class})\ \text{enters no ratio}\ [\mathrm{FT.113}];
\]

%dsp-level:0
\[
(5)\quad\textbf{the classification of the rational}\ \tfrac23\textbf{:}
\]

%dsp-level:0
\[
\qquad \tfrac23
\]

%dsp-level:0
\[
= \tfrac{3-1}3
\]

%dsp-level:0
\[
= \frac{\dim\ \text{doublet}}{\dim\ \text{triplet}}\ —\ \text{the decomposition ratio of the spatial seal triad under its inner rotation subgroup}\ [\mathrm{FT.111}(1):\ \mathbf 3=\mathbf 1\oplus\mathbf 2;\ \text{machine: exact Fractions}];
\]

%dsp-level:0
\[
\delta_\ell=\tfrac29=(\tfrac23)/3\ [\mathrm{FT.112}\ \mathrm{Rmk}]:
\]

%dsp-level:0
\[
\qquad \text{the three cross-check rationals}\ (p^*_\infty)^2
\]

%dsp-level:0
\[
= \tfrac23\ [\mathrm{FT.70}(3);\ \text{exact in}\ \mathbb Z[\sqrt6]],
\]

%dsp-level:0
\[
\quad\delta_v/\pi
\]

%dsp-level:0
\[
= 2-\tfrac43
\]

%dsp-level:0
\[
= \tfrac23\ [\mathrm{FT.63}(4)],
\]

%dsp-level:0
\[
\quad f_{\mathrm{charge}}(\delta h{=}3)=\tfrac23\ [\mathrm{FT.14}]
\]

%dsp-level:0
\[
\text{equal}\ (3-1)/3\ \text{exactly}\ [\text{machine}]:
\]

%dsp-level:0
\[
\qquad \mathrm{FT.89}(2\mathrm d)\text{'s classification of the rational}\ \tfrac23\ \text{``cap-3 combinatorial output''}\ \text{and}\ (5)\text{'s verdict on the Koide-form structure}\ \text{``coincidence;\ no structural claim''}\ \textbf{do not stand}:
\]

%dsp-level:0
\[
\text{the corpus-side Koide-form angle and spectrum have the derived spatial provenance}\ (1)\text{–}(4);
\]

%dsp-level:0
\[
\qquad \textbf{unchanged and standing:}\ \ \mathrm{FT.89}(1)\text{'s arithmetic}\ (\text{exact, frozen});
\]

%dsp-level:0
\[
\ \ (3)\text{'s refutation};
\]

%dsp-level:0
\[
\ \ (4)\text{'s negatives};
\]

%dsp-level:0
\[
\text{the inadmissibility grounds}\ (2\mathrm a)\text{–}(2\mathrm c)\ \text{and the identification bars: no family-mass read exists in the corpus},
\]

%dsp-level:0
\[
\ \neg\exists\ \text{a declared map relating derived rationals to measured mass ratios},
\]

%dsp-level:0
\[
\ \text{external data enters no derivation}\ [\mathrm{P1},\ \mathrm{P5}];
\]

%dsp-level:0
\[
\text{the cap-sensitivity arithmetic}\ [(2\mathrm d):\ c=2\to0.508004\ldots;\ c=4\to\neg\exists\ \text{interior boundary}]\ \text{stands as the exact polygon-ensemble record};
\]

%dsp-level:0
\[
\qquad \text{the scale}\ \mu\ \text{remains undeclared}\ (\text{P4}).
\]

\textbf{Rmk (channel labels).} The subscripts $\ell,q,\nu$ name ledger rows [FT.112, FT.119, FT.142 Rmks]; “Koide-form” names the corpus-internal functional shape $1+\sqrt2\cos(\delta+\tfrac{2\pi i}3)$; no particle identification is made or implied [P1, P5].

\textbf{Req.} CT0.12, CT1.2, CT1.4, CT1.8 (Part 1); FT.14 (Tiers F.1–F.11); FT.63(4) (Tiers F.20–F.22); FT.68(2), FT.70(3) (Tiers F.23–F.25); FT.89 (Tier F.30); FT.109, FT.110, FT.111, FT.112, FT.113, FT.119 (Tiers F.37–F.38); FT.142 (Tier F.42; the Rmk row); Def-H.41 cited as FT.110's arena provenance and imported nowhere; P4, P5 (Parts 1–2); FT.156 (this tier); exact $\mathbb Q(i)$ conjugation and polynomial rotation identities, exact $\mathbb Q$ ranks and characteristic polynomials, $\mathbb Z[\sqrt6]$ arithmetic, exact identities in $\mathbb Q(\sqrt2,\sqrt3)[c,s]/(c^2{+}s^2{-}1)$, Fractions [math; machine: \texttt{checks/\allowbreak{}verify\_\allowbreak{}H\_\allowbreak{}3p1\_\allowbreak{}chiral\_\allowbreak{}asymmetry.py} §C].

\textbf{Pf.}

%dsp-level:0
\[
(1):\ \mathcal W^3
\]

%dsp-level:0
\[
= 3\mathcal W\ [\text{machine}]
\]

%dsp-level:0
\[
\implies (\hat n\cdot W)^3
\]

%dsp-level:0
\[
= \hat n\cdot W\implies\exp(-i\theta\,\hat n\cdot W)
\]

%dsp-level:0
\[
= E_0+(\cos\theta-1)(\hat n\cdot W)^2-i\sin\theta\,
\]

%dsp-level:0
\[
(\hat n\cdot W)\ [\text{polynomial on}\ \operatorname{spec}\{1,0,-1\}];
\]

%dsp-level:0
\[
\qquad \theta=\tfrac{2\pi}3\implies E_0-\tfrac12\mathcal W^2-\tfrac i2\mathcal W
\]

%dsp-level:0
\[
=u\otimes u\ \text{entrywise over}\ \mathbb Q(i)\ [\text{machine}];
\]

%dsp-level:0
\[
\qquad U^3=E_0,\ UJ_iU^\dagger=J_{i+1},\ UK_iU^\dagger=K_{i+1}\ [\mathrm{FT.111}(1)\ \text{Pf re-run};\ \text{machine}];
\]

%dsp-level:0
\[
\qquad UW_iU^\dagger=W_{i+1}\ [\text{machine}];
\]

%dsp-level:0
\[
\qquad C_3\le SO(3)\ [\mathrm{CT1.8}\ \text{vector structure};\ \mathrm{FT.111}\ \text{Thm line}]
\]

%dsp-level:0
\[
(2):\ \hat E^{ii}=E^{ii}-\tfrac13\Sigma,\ \ \textstyle\sum_i\hat E^{ii}=0\ [\mathrm{FT.109}(2);\ \mathrm{CT1.2}\ \text{traceless split};\ \text{machine}];
\]

%dsp-level:0
\[
\qquad \operatorname{rk}\{\hat E^{11},\hat E^{22},\hat E^{33}\}=2\ [\text{exact flattening rank};\ \text{machine}];
\]

%dsp-level:0
\[
\qquad P^3=I_3,\ \operatorname{tr}P=0:\ x^2+x+1\ \text{irreducible over}\ \mathbb Q\ [\mathrm{FT.111}(1)\ \text{Pf re-run};\ \text{machine}]
\]

%dsp-level:0
\[
\textstyle\sum_i[J_i,K_i]
\]

%dsp-level:0
\[
= i\sum_i(\Sigma-E^{ii})
\]

%dsp-level:0
\[
= i(3\Sigma-\Sigma)
\]

%dsp-level:0
\[
= 2i\Sigma\ [\mathrm{CT1.7}\ \text{diagonal};\ \mathrm{FT.109}(1)\text{–}(2);\ \text{machine, entrywise}];
\]

%dsp-level:0
\[
\qquad 2=3-1\ [\text{exact integers}]
\]

%dsp-level:0
\[
(3):\ \dim\operatorname{span}\{J_i\}=3\ [\text{machine: exact rank}];
\]

%dsp-level:0
\[
\qquad T(1)
\]

%dsp-level:0
\[
= \mathbb 1_{4\times4}-[\,\cdot\,]_{(2,2)}\ \text{on the diagonal}:\ \operatorname{tr}T(1)
\]

%dsp-level:0
\[
= 4-1
\]

%dsp-level:0
\[
= 3;
\]

%dsp-level:0
\[
\quad\operatorname{charpoly}
\]

%dsp-level:0
\[
= x^4-3x^3-3x^2\ [\mathrm{FT.110}(2)\ \text{Pf re-run};\ \text{machine: exact Newton identities}];
\]

%dsp-level:0
\[
\qquad \delta=\tfrac{2}{3\cdot3}=\tfrac29\ [\text{Fractions}]
\]

%dsp-level:0
\[
(4):\ \text{closed forms}\ r_0
\]

%dsp-level:0
\[
= 1+\sqrt2 c,
\]

%dsp-level:0
\[
\ r_{1,2}
\]

%dsp-level:0
\[
= 1-\tfrac12(\sqrt2 c\pm\sqrt6 s)\ [\mathrm{FT.113}\ \text{re-derived by the addition formulas}];
\]

%dsp-level:0
\[
\qquad \sum_ir_i=3,\ \sum_{i<j}r_ir_j=\tfrac32\ [\text{exact in}\ \mathbb Q(\sqrt2,\sqrt3)[c,s]/(c^2{+}s^2{-}1);\ \text{machine}]
\]

%dsp-level:0
\[
2y^3-3y-\sqrt2(4c^3-3c)=0\ \ \forall y\in\{r_i-1\}\ [\text{same ring};\ \text{machine}];
\]

%dsp-level:0
\[
\qquad \cos\tfrac23=4c^3-3c\ \text{at}\ 3\delta=\tfrac23\ [\text{triple-angle identity in the ring};\ \mathrm{FT.113}\ \text{Pf pattern}]
\]

%dsp-level:0
\[
(5):\ \tfrac{3-1}3=\tfrac23\ [\text{Fractions}];
\]

%dsp-level:0
\[
\qquad \bigl(\tfrac{\sqrt6}3\bigr)^2=\tfrac69=\tfrac23\ [\mathbb Z[\sqrt6];\ \text{machine}];
\]

%dsp-level:0
\[
\qquad s=2+\sqrt6:\ s^2-4s-2=0,\ \tfrac{s}{1+s}=\tfrac{\sqrt6}3\ [\mathrm{FT.89}(1)\ \text{Pf re-run};\ \text{machine}];
\]

%dsp-level:0
\[
\qquad 2-\tfrac43=\tfrac23\ [\mathrm{FT.63}(4);\ \text{exact integers}];
\]

%dsp-level:0
\[
\qquad \tfrac{2/3}3=\tfrac29\ [\text{Fractions}]
\]

%dsp-level:0
\[
\text{classification row}: (1)\text{–}(4)\ \text{conjoined};
\]

%dsp-level:0
\[
\qquad \text{standing rows}: \mathrm{FT.89}(1)\text{–}(4)\ \text{cited unmodified};\ \mathrm{P1}/\mathrm{P5}\ \text{bars verbatim}.\ \blacksquare
\]

\textbf{Phys.} The $120^\circ$ rotation that shapes the triplet spectrum is a rotation of the three spatial directions performed by the weak generators themselves — an exact polynomial in $W_1+W_2+W_3$ with rational coefficients. The channel angle is the doublet dimension over the arena: $(3-1)$ over $3\times3$ — a ratio of spatial counts, with the same integer $2$ appearing independently as the seal-commutator coefficient. The famous cosine form is then forced three times over: covariance fixes the shape, the triad dimension fixes the normalization, the counts fix the angle — nothing is adjustable except an overall scale that no ratio sees. The classification of two thirds as a bookkeeping coincidence does not stand, at exactly that clause: the same rational is the singlet-complement fraction of the spatial triad under its own rotation group, and the Koide-form structure has a provenance theorem. Every bar against reading any of it as a measured mass ratio stands untouched.

\vspace{2pt}\noindent\rule{\textwidth}{0.4pt}\vspace{2pt}

\end{theorem}

\begin{theorem}[{\texorpdfstring{FT.159 — The emergent-time identity theorem: gravitational time and CPT time are the ledger order}{FT.159 - The emergent-time identity theorem: gravitational time and CPT time are the ledger order}}]\label{tag:FT-159}
\textbf{Thm.}

%dsp-level:0
\[
(1)\quad\textbf{gravitational time is the tick coordinate:}
\]

%dsp-level:0
\[
\qquad x^0:
\]

%dsp-level:0
\[
= \tau_n\ [\mathrm{D1.8}:\ \tau_n=\sum\delta\tau\ \text{accumulated along a causal chain of real edges; the emergent coordinate}];
\]

%dsp-level:0
\[
\chi([v])
\]

%dsp-level:0
\[
= (x^0,x^1,x^2,x^3)
\]

%dsp-level:0
\[
= (\tau_n,h_1,h_2,h_3)\ [\mathrm{CT2.6}:\ M=\Gamma_{\mathcal R}/\!\!\sim\ \text{a smooth}\ 4\text{-manifold};\ g_{\mu\nu}\ \text{a smooth Lorentzian metric}\ [\mathrm{CT1.12}];\ \text{conditions (ii)–(iii) at the point of use — A §0.2a}];
\]

%dsp-level:0
\[
g_{00}=\Omega^2\bar g_{00}=-\Omega^2,
\]

%dsp-level:0
\[
\qquad \Omega=e^{\alpha_\Omega a_{00}}\ [\mathrm{D2.1}:\ \bar g_{00}=-1];
\]

%dsp-level:0
\[
\qquad a_{00}=\text{the}\ E_0\text{-coefficient of}\ X_{(+)}\in\mathcal E_{++}:\ S\text{-even},\ T\text{-even}\ [\mathrm{CT1.2}\ \text{table};\ \text{machine}]:
\]

%dsp-level:0
\[
\text{the lapse/conformal read is built from the}\ S\text{-even}\ T\text{-even block}\ [\mathrm{CT1.12}\ \text{table}:\ a_{00}\mapsto\text{conformal factor}];
\]

%dsp-level:0
\[
\qquad \text{the metric's time coordinate is the ledger's tick structure}\ [\mathrm{A7}:\ \mathfrak L=\bigcup_{n\ge0}\mathfrak L(\tau_n)];
\]

%dsp-level:0
\[
(2)\quad\textbf{CPT time is the ledger reversal:}
\]

%dsp-level:0
\[
\qquad \Theta_{\mathfrak L}\ [\mathrm{D2.7}]:\ \text{ledger order reversed},\ J^+\leftrightarrow J^-,\ G(e)\mapsto G(e)^\ast;
\]

%dsp-level:0
\[
T=\Theta_{\mathfrak L}:\ \text{the}\ T\text{-register of the configuration calculus is the ledger reversal}\ [\mathrm{FT.133}(1):\ \text{order reversed},\ J^+\leftrightarrow J^-,\ \kappa_e\mapsto-\kappa_e,\ \sigma\ \text{fixed};\ \mathrm{FT.88}(2)];
\]

%dsp-level:0
\[
w(\Theta_{\mathfrak L}C)
\]

%dsp-level:0
\[
= w(C)^\ast\ [\mathrm{CT2.8}(1);\ \text{machine: exact witnesses over}\ \mathbb Q(i)\ —\ G(e)=\sigma_e e^{i\pi\kappa_e}\ [\mathrm{FT.88}(2);\ \mathrm{D3.5}],\ e^{i\pi\kappa}\in\{1,i,-1,-i\}\ \text{on half-integer}\ \kappa\ \text{samples},\ \kappa\mapsto-\kappa=\text{conjugation},\ \sigma\ \text{fixed},\ \mathcal J\mapsto\mathcal J^\ast\ [\mathrm{D3.6}];\ \text{exhaustive on the}\ 3\text{-edge configurations}];
\]

%dsp-level:0
\[
\langle A|B\rangle:=Z\bigl[\Theta_{\mathfrak L}A\cup B\bigr]=\sum_C|w(C)|^2\ge0,
\]

%dsp-level:0
\[
\qquad w(C)
\]

%dsp-level:0
\[
= w(C_{\mathrm{left}})\cdot w(C_{\mathrm{right}})\ \ \text{across any interior ledger-time cut}\ [\mathrm{CT2.8}(2)\text{–}(3);\ \text{machine: exact regrouping witnesses}]:
\]

%dsp-level:0
\[
\text{the Hilbert-space structure is glued across the same ledger time by the same reversal}\ [\mathrm{CT2.8}(5)\text{–}(6)];
\]

%dsp-level:0
\[
(3)\quad\textbf{the identity:}
\]

%dsp-level:0
\[
\qquad \text{gravitational time}\ (1)\ \text{and CPT time}\ (2)\ \text{are the same entity — the ledger order parameter — read once as the coordinate and once as the reversal:}
\]

%dsp-level:0
\[
x^0=\tau_n\ [\mathrm{D1.8};\ \mathrm{CT2.6}]\qquad\text{and}\qquad\text{its reversal}=\Theta_{\mathfrak L}\ [\mathrm{D2.7};\ \mathrm{FT.133}(1):\ T=\Theta_{\mathfrak L}];
\]

%dsp-level:0
\[
\text{the metric side transports: the register involution is declared with the metric signature}\ [\mathrm{A3},\ \mathrm{C13}:\ \eta=(-,+,+,+)]\ \text{and its}\ T\text{-parity transport through the dictionary is the time-reversal pullback}\ [\mathrm{FT.156}(3):\ (a_{00},q,w)\mapsto(a_{00},q,-w)\iff g_{00},g_{ij}\ \text{fixed},\ g_{0i}\mapsto-g_{0i};\ \text{machine}];
\]

%dsp-level:0
\[
\text{the record side conjugates:}\quad w(\Theta_{\mathfrak L}C)=w(C)^\ast\ [\mathrm{CT2.8}(1);\ (2)]:
\]

%dsp-level:0
\[
\qquad \text{the time reversal of the emergent metric and the}\ T\text{-register of the CPT calculus are the reversal of the same ledger order at two loci — the metric dictionary and the configuration record};
\]

%dsp-level:0
\[
\text{locus discipline}\ [\mathrm{FT.134}(1)\ \text{pattern}]:\ \text{the register grading}\ [\mathrm{A3}]\ \text{and the record reversal}\ [\mathrm{D2.7}]\ \text{act on different spaces — the identity is of the time entity}\ (\text{the order parameter}),
\]

%dsp-level:0
\[
\ \neg\text{an operator equality};
\]

%dsp-level:0
\[
\text{inventory: one order parameter, four tags}\ [\mathrm{D0.7}\ \text{ticks};\ \mathrm{D1.8}\ x^0;\ \mathrm{CT2.6}\ \text{charts};\ \mathrm{D2.7}\ \Theta_{\mathfrak L}]\ —\ \text{no second time entity exists in the corpus}\ [\text{scan}];
\]

%dsp-level:0
\[
(4)\quad\textbf{the blockage removed; the residue restated:}
\]

%dsp-level:0
\[
\qquad \text{the ``no emergent time'' prerequisite}\ [\text{item}\ 21(\mathrm b)\ \text{row};\ \mathrm{FT.131}(3):\ \text{emergent time a named prerequisite}]\ \textbf{is discharged by}\ (1)\text{–}(3):
\]

%dsp-level:0
\[
\text{the time parameter is derived content — the coordinate}\ x^0
\]

%dsp-level:0
\[
= \tau_n\ [\mathrm{D1.8};\ \mathrm{CT2.6}],
\]

%dsp-level:0
\[
\ \text{the lapse}\ g_{00}
\]

%dsp-level:0
\[
= -\Omega^2\ [\mathrm{D2.1};\ \mathrm{CT1.12}],
\]

%dsp-level:0
\[
\ \text{the reversal}\ \Theta_{\mathfrak L}
\]

%dsp-level:0
\[
= T\ [\mathrm{D2.7};\ \mathrm{FT.133}(1)],
\]

%dsp-level:0
\[
\ \text{the cut-glued inner product}\ [\mathrm{CT2.8}(2)\text{–}(3)];
\]

%dsp-level:0
\[
\text{the line-element channel contracts metric data with tick data:}
\]

%dsp-level:0
\[
\quad ds^2(e)
\]

%dsp-level:0
\[
= g_{\mu\nu}(e)\,
\]

%dsp-level:0
\[
\Delta x^\mu_e\,
\]

%dsp-level:0
\[
\Delta x^\nu_e,
\]

%dsp-level:0
\[
\quad\widehat{ds}^{\,2}
\]

%dsp-level:0
\[
= ds^2/\mathfrak h^2\ [\mathrm{D1.6}],
\]

%dsp-level:0
\[
\quad\Delta x^0
\]

%dsp-level:0
\[
= \text{a tick increment}\ [\mathrm{D1.8};\ \mathrm{CT2.8}(2)\ \text{Pf}:\ \text{each edge is registered between consecutive ticks}]:
\]

%dsp-level:0
\[
\Omega\ \text{and}\ \Delta\tau\ \text{meet in every}\ \widehat{ds}^{\,2}(e)\ \text{with a temporal increment}\ [\text{machine: the bilinear-form witness}\ ds^2=g_{00}(\Delta x^0)^2\ \text{on purely temporal increments}];
\]

%dsp-level:0
\[
\qquad \text{the derived clock:}\ (\Delta n_{\mathcal R},\delta h,\delta\tau)=\pm(1,1,1)\ \text{per Janus move}\ [\mathrm{CT0.3}];
\]

%dsp-level:0
\[
\text{restated precisely — the unconstructed residue is the dynamical rate modulation: a flicker-frequency read of}\ \delta\tau\ [\mathrm{D0.7}]\ \text{against}\ \Omega/N_i\ \text{along materialized worldlines (a redshifted or self-modulating rate)}\ [\mathrm{FT.86}(8)\text{'s inventory row stands as the construction record}];
\]

%dsp-level:0
\[
\text{it requires no declaration of time — time is derived}\ (1)\text{–}(3)\ —\ \text{and no declaration issues};
\]

%dsp-level:0
\[
\qquad \mathrm{FT.131}(3)\text{'s deferred frequency read has its time parameter to be posed in};
\]

%dsp-level:0
\[
(5)\quad\textbf{strict forwardness stands:}
\]

%dsp-level:0
\[
\qquad \neg\exists\,\tau^{-1}\ [\mathrm{A6}];
\]

%dsp-level:0
\[
\qquad \Theta_{\mathfrak L}\ \text{is the formal mirror on records},
\]

%dsp-level:0
\[
\ \neg\text{a dynamical inverse}\ [\mathrm{FT.144}(4):\ \text{the pure order reversal of every completed nonempty history is inadmissible at its first event};\ \text{machine}];
\]

%dsp-level:0
\[
\text{materialization is single-arrow}\ [\mathrm{FT.144}(2)\text{–}(3)];
\]

%dsp-level:0
\[
\qquad \text{the ledger count is monotone}\ [\mathrm{A7};\ \text{machine: exact samples}]:
\]

%dsp-level:0
\[
\text{the CPT time reversal is a register operation on configuration data; the gravitational arrow is the ledger's forward advance — the same entity on the record side and on the dynamical side}.
\]

\textbf{Req.} A3, A6, A7 (Part 1); D0.5, D0.7, D1.2, D1.6, D1.8, D2.1, D2.7, D3.3, D3.5, D3.6 (Part 1); CT0.3 (Part 1); CT1.2, CT1.12, CT2.6 (conditions (ii)–(iii) at the point of use — A §0.2a), CT2.8 (Part 2); FT.86(8) (Tier F.29); FT.88 (Tiers F.29–F.30); FT.131(3) (Tier F.40); FT.133, FT.134(1) (Tier F.41); FT.144 (Tier F.42); FT.156 (this tier); exact $\mathbb Q(i)$ weight arithmetic on the half-integer $\kappa$ ladder, exhaustive word enumeration, exact rational regrouping and transport identities [math; machine: \texttt{checks/\allowbreak{}verify\_\allowbreak{}H\_\allowbreak{}3p1\_\allowbreak{}chiral\_\allowbreak{}asymmetry.py} §D].

\textbf{Pf.}

%dsp-level:0
\[
(1):\ \mathrm{D1.8}\ \text{verbatim}:\ x^0:=\tau_n;
\]

%dsp-level:0
\[
\qquad \mathrm{CT2.6}\ \text{chart line verbatim}:\ \chi([v])=(\tau_n,h_1,h_2,h_3);
\]

%dsp-level:0
\[
\qquad g_{00}=\Omega^2\bar g_{00}=-\Omega^2\ [\mathrm{D2.1}];
\]

%dsp-level:0
\[
\qquad E_0\in\mathcal E_{++}:\ (S,T)=(+,+)\ [\mathrm{CT1.2}\ \text{table};\ \text{machine}];
\]

%dsp-level:0
\[
\qquad \mathrm{CT1.12}\ \text{row}\ a_{00}\mapsto\Omega
\]

%dsp-level:0
\[
(2):\ \mathrm{D2.7}\ \text{verbatim};
\]

%dsp-level:0
\[
\qquad \mathrm{FT.133}(1):\ T=\Theta_{\mathfrak L}\ \text{register line verbatim}\ [\mathrm{FT.88}(2):\ \kappa\mapsto-\kappa];
\]

%dsp-level:0
\[
\qquad w(\Theta_{\mathfrak L}C)
\]

%dsp-level:0
\[
= w(C)^\ast\ [\mathrm{CT2.8}(1)\ \text{frozen};\ \text{machine re-run:}\ e^{i\pi\kappa}\mapsto e^{-i\pi\kappa}=(e^{i\pi\kappa})^\ast\ \text{on}\ \kappa\in\tfrac12\mathbb Z:\ \{1,i,-1,-i\}\subset\mathbb Q(i);\ \sigma\ \text{fixed};\ (-1)^\aleph\ \text{invariant}\ [\mathrm{D3.3}:\ \aleph_v=|f_v|]]
\]

%dsp-level:0
\[
\langle A|A\rangle=\sum_C w(\Theta_{\mathfrak L}C)w(C)=\sum_C|w(C)|^2\ [\mathrm{CT2.8}(3)\ \text{Pf pattern}];
\]

%dsp-level:0
\[
\qquad \text{cut factorization}: w
\]

%dsp-level:0
\[
= \mathcal J\prod_eG(e)\cdot(-1)^\aleph,
\]

%dsp-level:0
\[
\ \text{factors tick-local}\ [\mathrm{D0.7};\ \mathrm{CT2.8}(2)\ \text{Pf regrouping};\ \text{machine: exact products on witnessed cuts}]
\]

%dsp-level:0
\[
(3):\ (1)+(2)\ \text{conjoined};
\]

%dsp-level:0
\[
\qquad \text{transport row}\ [\mathrm{FT.156}(3);\ \text{machine}];
\]

%dsp-level:0
\[
\qquad \text{inventory scan}:\ \{\mathrm{D0.7},\ \mathrm{D1.8},\ \mathrm{CT2.6},\ \mathrm{D2.7}\}\ \text{name one order parameter}\ [\text{no second time coordinate in any tag}]
\]

%dsp-level:0
\[
(4):\ \text{item}\ 21(\mathrm b)\ \text{prerequisite text cited as discharged by}\ (1)\text{–}(3);
\]

%dsp-level:0
\[
\qquad \mathrm{D1.6}\circ\mathrm{D1.8}:\ g_{\mu\nu}\Delta x^\mu\Delta x^\nu\big|_{\Delta x=(\Delta\tau,0,0,0)}
\]

%dsp-level:0
\[
= g_{00}(\Delta\tau)^2\ [\text{bilinear substitution};\ \text{machine witness over}\ \mathbb Q];
\]

%dsp-level:0
\[
\qquad g_{00}=-\Omega^2\ [\mathrm{D2.1}];
\]

%dsp-level:0
\[
\qquad \mathrm{FT.86}(8)\ \text{row cited unchanged};
\]

%dsp-level:0
\[
\qquad \neg\exists\ \text{new declaration}\ [\text{scan of this tier's Req}]
\]

%dsp-level:0
\[
(5):\ \mathrm{A6}\ \text{verbatim};
\]

%dsp-level:0
\[
\qquad \mathrm{FT.144}(4)\ \text{machine re-run (exhaustive, lengths}\le8);
\]

%dsp-level:0
\[
\qquad \mathrm{A7}\ \text{monotonicity on exact samples}.\ \blacksquare
\]

\textbf{Phys.} The corpus has one time. The manifold's time coordinate is the ledger's own tick index, declared as such and charted onto the emergent geometry; its time-time metric component is the conformal read of the reality-preserving, reversal-even block. The CPT calculus's $T$ is the ledger reversal — the order of the record run backwards, seals exchanged, phases conjugated — and the same reversal, transported through the metric dictionary, is exactly the time-reversal pullback: lapse and spatial metric standing still, shift flipping sign, signature untouched. Gravitational time and CPT time are therefore not two entities needing a bridge; they are the ledger order read as coordinate and read as reversal. The blockage that deferred every frequency read for want of a time parameter is removed — the parameter exists and is derived. What remains unconstructed is honest and small: a dynamical rate modulation of the flicker against the conformal factor — a redshift — which is now a construction problem inside derived time, not a missing declaration. And the arrow stands: the reversal is a mirror on records; the advance has no inverse.

\vspace{2pt}\noindent\rule{\textwidth}{0.4pt}\vspace{2pt}

\end{theorem}

\medskip\noindent\textit{\texorpdfstring{FT.160 — The 3+1 completion: adjudication and scope}{FT.160 - The 3+1 completion: adjudication and scope}.}\ Structural-constraints record: see Appendix B, Section B.1.14 of the Addendum.\par\medskip


% Part4_Baryomesonic_era.md — sanitized presentation layer (generated; see source/CONVERSION_LOG.md)
\chapter{\texorpdfstring{$\langle qu(b(i(t|t))ick)\rangle$ — \textbf{Part 4 · Baryomesonic\_era} — Colour-Sealed Composites: Baryons, Mesons, Junctions}{Qubitick - **Part 4 - Baryomesonic _era** - Colour-Sealed Composites: Baryons, Mesons, Junctions}}

\textbf{Dep.} Part 1 (axioms A1–A9; conventions C1–C13; qubit-era definitions D0.1–D1.12; calibration P1–P5; pair canon CT0.\textit{/CT1.1–CT1.11b; ledger dynamics L0–L5); Part 2 (definitions D2.1–D4.3; geometry CT1.12–CT2.11; electroweak EW.}; electromagnetic EM.\textit{; colour CT3.1–CT3.22a; saturation CT4.}); Part 3 (the fermion layer: Def-H.1–Def-H.43, FT.1–FT.155); \textbf{the Census annex} (Appendix A of this volume; the computed cells, cited row-wise). Format: Part 1 §0.4. \textbf{Namespace.} This part owns: Def-BM.1–Def-BM.20; BM.1–BM.152. Tiers 4.0, BM.1–BM.24 in series order. \textbf{DAG law.} Citations resolve to Parts 1–3, Census cells, or earlier tags of this part. Nothing in the canon cites this part. \textbf{0.1 Charter.} The baryomesonic era is the layer of colour-sealed composites: the algebraic sealings of D3.9 (baryon $\varepsilon_{ijk}q^iq^jq^k$; meson $\delta^i{}_j\,q^jr_i$) realized on bound-chain constituents (Def-H.4) at shared junctions. Its foundation stack, by citation: (i) the colour and sealing machinery [CT3.1–CT3.22a; D3.8, D3.9; D2.14; CT4.1–CT4.7]; (ii) the derived sealing grammar and junction structure [FT.29, FT.34(5), FT.19, FT.39]; (iii) the computed baryon-type cells [Census FT.35: the first colour-active, first $D^-$-carrying floor; Census FT.49: the baryon-type sealing phase $\Theta_{\mathrm{int}}\in4\mathbb Z$]; (iv) the channel machinery [FT.109–FT.119: the $\ell$/q-channels with the $\varepsilon$-selected $C_3$ and confinement blindness; FT.141–FT.145: the rung-one channel and the arrow law]; (v) the carrier and parity structure [FT.122–FT.134]; (vi) the sector formalisms [FT.102–FT.108]; (vii) the conservation and symmetry layer [Tier BM.4: the Janus balance and the telescoping endpoint read, the continuum promotion of the hop, the $C_3$ and the slot algebras, the stress-energy tensor of a sealed composite, the canonical-pair obstructions, the differential conservation law]; (viii) the calibration layer [Tier BM.5: the Born width, the coherence volume and energy density, the curvature of a sealed composite, the dimensionless lock identities, the neutron row of the identification layer]; (ix) the flicker layer [Tier BM.6: the topology-blind ensemble of a sealed composite and its derived weights, the path-integrated mass from the curvature law, the binding read, the static force matrix on the two charge classes, the transition with its forced carrier, the family naming rows]; (x) the CPT-sector and charge-restriction layers [Tier BM.7: the EM$^\pm$ sector labels and the EM scar, the mirror involution and its exact phase, the proton and antiproton rows, and the charge-constrained probability space with its three mass routes]; (xi) the additive charge read and the mass-route adjudication [Tier BM.8]; (xii) the mesonic flicker test [Tier BM.9: the contraction grammar, the colour projector, and the protected junction]; (xiii) the metric-energy and vacuum-entropy layer [Tier BM.10: the dispersion read of a sealed composite from the emergent metric with its kinetic energy and zero-point bound, the vacuum-balanced-entropy reframing of the rest mass, the force-structure ladder with its strength bar and the inter-family constraint, and the route-artifact adjudication closing the counting trap]; (xiv) the decay-vertex layer [Tier BM.11: the family transport-profile map discharging the Yukawa-labelled gap, the family mass form with its positivity, splitting law and sum rule, the rung-1 isospin row and the neutrality of the charge vertex, the forced neutral member and the Q-value sign theorem, the angular-momentum budget with its absorption adjudications, the member-wise curvature rows, and the per-tick time evolution of the assembled reaction]; (xv) the electromagnetic self-interaction layer [Tier BM.12: the chart-side EM source read of a sealed composite, the scar's field energy and the mass-splitting theorem, the infinite-range seal asymmetry with its UV/IR finiteness rows, the two-body potential with the binding category bar, and the beta-decay re-adjudication with the endpoint-to-average refutation]; (xvi) the full-scope mass-splitting layer [Tier BM.13: the Lyapunov foundation with its iterated contraction law and class-degenerate zero, the complete-force-operator sector census with its five zeros and one monomial, the metric background's zero contribution and the non-uniform clock with its Mandelstam–Tamm resolvability bound, the two binding zeros and the unconstructed atom, and the exact symbolic ratio with its named order residue]; (xvii) the constituent-transport-splitting layer [Tier BM.14: the continuation-symmetry census of the two vertex types with its three annihilation rows, the observation frames and the charge-blindness of the metric, the constituent transport splitting that fills the residue BM.55(4) names, the six-state quark assembly on the three Platonic mass floors with its branch order and its two inverted hierarchies, the coupling-blindness theorem and the named obstruction to $\alpha_B$, the extended six-term decomposition with its exact critical strength and sign theorem, and the decay re-adjudication with its exact closure margin]; (xviii) the organic-bridge layer [Tier BM.15: the operator lift of the EM correction with its charge-blind linear trace and its quadratic invariant, the derived $\mathbb Z_2$ Wilson sector with the two-torsion and calibration-weight obstructions to the real holonomy lift and its structured obstruction record, the constrained joint entropy as a fiber cardinality with its stationary geometric factor of one and its named missing enumeration, the bridge convergence with the non-derivation theorem reducing Gap 7 to one named scale ratio and two named attachments, the two-bit daughter-floor theorem with the formed-closed dichotomy and its exact margin ladder, and the five-route opening-term census closing the beta-decay energy audit]; (xix) the chiral-assignment layer [Tier BM.16: the left–right assignment rule coupling the four derived gradings and establishing the assignment rule's forced shape, the ray-resolved spin census with the opposite-ray vertex table, the weight–branch–core attachment theorem with the sign-coherent quark mass splitting, the nilpotent chiral vertex operator with its flip, rate and arrow rows, and the vertex read-specification with the exact trichotomy of the beta-decay opening branch]. (xx) the mass-ratio-census layer [Tier BM.17: the Platonic winding census with the ladder law on both priced windings, the Chebyshev tower closing at the dodecahedral top weight, and the intra-generational ratios as exact sine ratios with their minimal polynomials; the inter-generational ratios as the three roots of one cubic form with three channel constants, the cross-channel identity between the lepton channel's heaviest line and the quark channel's cubic constant, and the scale-free hierarchy-order theorem; the six-state ratio census with the splitting order, the doublet ratios, the two inverted hierarchies and the three pairwise non-proportional ladders; the bridge census with its six exact factors, the intra-currency closure and the localized single residue; the inter-family bridge in full with the four-currency census, the exact incommensurability of the two formed ratios, and the three-number localization; the neutron-proton mass ratio with its exact range, its irrational critical locus and its total trichotomy; the Q-value in two formed reads with both margin ladders and the read-relative verdict; and the census adjudication] (xxi) the electromagnetic self-energy correction layer [Tier BM.18: the category adjudication voiding the chart-side distribution factor and deriving the formed exterior value one half on two distribution-free routes; the running coupling as an exact geometric profile on the envelope ratio, the lapse and the curvature lock, vanishing at the member's own Born width and growing with the energy scale; the self-energy as an ensemble average over the flicker with one derived attenuation, the crystallization weight; the branch-blindness theorem refuting a defect-conditioned retention on four grounds; the corrected neutron–proton splitting with its parameter-free bar, its universal ratio bound and its ensemble-robust incommensurability; the Q-value on both formed reads with the uniform closure and the untouched specified branch; and the correction's adjudication] (xxii) the flicker-integral and rest-mass layer [Tier BM.19: the flicker integral replacing the crystallized snapshot, with the derived ensemble distribution, the materialization restriction and the two crystallization weights; the vertex superposition with its exact fibre weights and the two positive parts of the transport splitting; the rung-one $\nu$ line's seal-derived rest mass, its strict positivity, its single-ray occupancy and the four bars against a seesaw-type mechanism; the rest-mass Q-value on the specified and pure-pole reads with the exact critical locus; the relating equation in four exact forms with its flicker weighting and its integer bound; the materialized splitting with its two-sided parameter-free bar and its two tightened bounds; and the tier's adjudication] (xxiii) the six-objective-resolution layer [Tier BM.20: the energy-scale graph of the forward and inverse reads with the derived crossing strength, its exact enclosure and the composition-dependent crossing count; the dimensional bridge, with the $\mathfrak q$-exponent census of the inventory, the seven rows derived or adjudicated, the two bridges separated and the one calibration slot declared; the three-constituent assembly's two formed eigenvalue equations and the tensor-product referent census; the curvature composition's back-reaction fixed point with its three exact solution rows, its pole locus and the non-selection theorem; the general crystallization weight with its measure-free bounds, its threshold family and the refuted defect conditioning; and the kind structure of the arrow, with the spin-parity refutation at the floor and the exchange-structure referent census] (xxiv) the antimatter-and-resonance layer [Tier BM.21: the CP-conjugate chain derived from the three reversals with the T-block on the forward-only tick and the Lyapunov contraction, its exact phase, the mass-degeneracy theorem and the forcing theorem from the ledger's own total-charge-zero balance, with the Dirac parallel as an inventory statement; the CP-pair vacuum fluctuation with the Janus resolvent's own creation weight, its exact maximum one quarter and its two endpoint zeros; the two-chain state space discharging BM.6(2)'s exchange-space specification, the meson as the unique $\delta$-sealed $2$-juncture, the CP-pair saturation of the neutral class at exactly six, the two additive mass routes with their vanishing seal-binding and the named two-chain ensemble residue, and the amplitude/norm reconciliation on the locus distinction; the exchange operator as the flip composed with the CP transport, an unconditional involution with its eigenspace census and its statistics-blind orientation block; the spin–statistics connection $\varepsilon=(-1)^{|k|}=(-1)^{2j}$ from the fourth-root phase alphabet and the winding additivity, closing deficit 11 with the sector choice priced; the resonance spectrum with the transport characteristic polynomial $(z+1)^{2n}/z^{n}$ at every chain count, the degree-16 eigenvalue field, the multiset-orbit law with stabilizer multiplicities, the one exact three-chain collision and the minimal-polynomial degree census; the materialization threshold in its three exact readings with the never-empty locus and the no-forbidden-state theorem; the decay width with its three exact zeros, its bound and its three-tick lifetime floor; and the groupings, the ceilings and the octet adjudication] (xxv) the annihilation-channel and concrete-spin layer [Tier BM.22: the Janus register with its two nilpotent crossing moves, its master generator's decomposition, kernel, spectrum and derived detailed balance; the c-threshold adjudication with the exact subluminal deficit $1-v=m^2/(E(E+p))>0$ and the register boundary that replaces it; the zig-zag as the resolvent's evenness, every leg forward and the record involution fixing each completed word; the pair channel with its charge balance, exact threshold, weight and three CP legs, and the two-photon bar on four grounds; the annihilation channel with its exact direction ratio, detailed-balance indicator and entropy destination, closing residue (R9) at the channel level; the spin–winding relation $\Theta=4j$, the multiplet defect ladder in Chebyshev form with its standing zero and ceiling, and the dimensional bar against fusing the winding with the frame-holonomy angle; the refutation of a spin-gated tunnel as hypothesis (W7) on the four standing grounds, with the bound/transit domain row in its place; the winding-parity identification with the statistics and kind reads separated; the three-chain exchange structure with the $S_n$ character formula $d^{\,c(\sigma)}$, the vanishing carrier wedge against the one-dimensional colour wedge, and the multiplicative exchange law discharging (R7) at $n=3$; and the meson's unconditional boson sign against the baryon's exact odd-count fermion criterion] (xxvi) the closing layer [Tier BM.23: the Haar–Poisson duality with its twirl fixed point, complementary factorization, unique common fixed point and proper-time contraction law; the derived lapse $N(k,E)=\sqrt{1-G/(4\pi k)}$ with its formed domain and its envelope row, closing the lapse gap; the curvature composition selected as $\mathsf g=1/(1+G)$ by Lyapunov contractivity and minimal lapse-domain loss, completing the running coupling; the momentum-space $\Sigma$-resolvent with its form uniqueness, pole reads, on-shell measure and two-body phase space, closing item 17; the multi-chain charge census with the stretched states and the gauged read, closing Q10 and deficit 10; the $S_3$–colour-vertex equivariance with the exchange structure at every $n$ and the statistics per class, closing (R7$'$); the two-chain ensemble measure with the general $w_{\mathrm{mat}}^{(n)}$ law, the sector-mass pincer and the formed meson binding, closing (R8); the seed-tensoring mechanism as a finite-group twirl, closing item 19; the four convention declarations (A-BM.10–A-BM.13), closing item 16; the dodecahedron selection theorem with the five-currency ceiling census, closing deficit 12 as adjudicated; the Jensen residue as an exact rational function with its two-chain companion, closing (R6); the nilpotent supply vertex with the pair amplitude and the two symbolic cross-sections; the house cleaning and the calibration-register consolidation; and the era closure with the balanced ledger] (xxvii) the continuum-limit layer [Tier BM.24: the complete force operator acting over the graph, with the exact block census and the eight-step composed action; the continuum-limit function of content and span with its five closed forms and four formed-domain conditions; the metric–energy–continuum map with the dictionary bijection, the pure-trace linearized solve and the on-shell inverse read; the breaking threshold $E_{\mathrm{crit}}=2r/\mathfrak h$ with the envelope corollary; the locked-sector complement at the breaking locus with the absorption anchors, the one-way membrane and the retention theorem; the universality of the limit on one operator, one arena, one action and three fingerprints; and the derivation of $\mathrm{CT2.6}(\mathrm{ii})\text{--}(\mathrm{iii})$ as theorems — the unbounded ladder and the $\sqrt3$ spectral bound — leaving the extra-axiomatic list empty and closing the foundational era] Its enumerative substrate is the Census block: 25 cells, each carrying its consumption row (Census §Consumption). \textbf{0.2 Standing bars.} Species names are admitted in this part through the identification layer Def-BM.3 alone: a name is a label on a derived structural tuple, it carries no datum, and it occurs in no \textbf{Thm} clause of Parts 1–3, whose own label-discipline rows stand unmodified. No external data outside quarantined audit rows [Census FT.114, FT.120; FT.73(5) pattern]; no measured external constant is identified with any corpus quantity [P5]; topology labels are internal bookkeeping [Def-H.32(4)]; amplitude $\ne$ mass [FT.79(5)]; rung $\ne$ scale [D0.8].

\phantomsection
\section[§Cal-BM — Free internal parameters of the baryomesonic era…]{\texorpdfstring{§Cal-BM — Free internal parameters of the baryomesonic era (P4-class)}{SCal-BM - Free internal parameters of the baryomesonic era (P4-class)}}

None introduced by this part. Inherited and undeclared: the family-channel overall scale $\mu$ [Cal-F; FT.113, FT.119 — every ratio is $\mu$-free]. \textbf{Declared slot, value $\neg$issued (A-BM.9, Tier BM.20):} the bridge's dimensionless calibration ratio $\rho_{\mathrm B}=\hat\mu'^2\mathfrak h^2/(\ln2)^2\in\mathfrak q^0$ [BM.69(2); BM.98; BM.102(5)] — a packaging of Cal-F's $\hat\mu'^2$ and P1's $\mathfrak h$, hence $\neg$a new free parameter and $\neg$a new dimensionful scale; no numerical value attaches to it here or anywhere [P5; BM.69(4)'s four grounds stand: the slot is declared, $\neg$derived]. No second dimensionful scale exists [P1; the derived no-go: Part 1 §E]. Any numerical use of an inherited scale beyond its Part-3 loci requires a logged declaration first.

\textbf{Consolidated calibration register (Tier BM.23, BM.143(4) — a consolidation, $\neg$an addition; the parameter count is unchanged).}

%dsp-level:0
\[
(\mathrm{C\text{-}BM.1})\quad\rho_{\mathrm B}
\]

%dsp-level:0
\[
= \hat\mu'^2\mathfrak h^2/(\ln2)^2\in\mathfrak q^0:\ \ \text{the bridge ratio; slot declared}\ [\mathrm{A\text{-}BM.9}],
\]

%dsp-level:0
\[
\ \text{value}\ \neg\text{issued}\ [\mathrm{BM.143}(4)];
\]

%dsp-level:0
\[
(\mathrm{C\text{-}BM.2})
\]

%dsp-level:0
\[
\quad e_0:\ \ \text{the numerical value of the electric coupling}\ [\mathrm{EM.25}(7);\ \mathrm{BM.88}(8);\ \mathrm{BM.133}(6)];
\]

%dsp-level:0
\[
\ \ \text{value}\ \neg\text{issued};
\]

%dsp-level:0
\[
\ \neg\text{fixed by measurement}\ [\mathrm{P5}];
\]

%dsp-level:0
\[
(\mathrm{C\text{-}BM.3})\quad\mu_\nu,
\]

%dsp-level:0
\[
\ \mu,
\]

%dsp-level:0
\[
\ \rho_\ell,
\]

%dsp-level:0
\[
\ \rho_\nu:\ \ \text{the channel scales, inherited from}\ \mathrm{Cal\text{-}F},
\]

%dsp-level:0
\[
\ \neg\text{declared in this part}\ [\mathrm{FT.145}(3);\ \mathrm{BM.83}(5),(7)].
\]

Each row is a calibration datum of the P4/A-BM.9 species, $\neg$a structural gap; no row carries a value; $\mathfrak h$ remains the single dimensionful scale [P1].

\phantomsection
\section[Tier 4.0 — Composite Vocabulary]{\texorpdfstring{Tier 4.0 — Composite Vocabulary}{Tier 4.0 - Composite Vocabulary}}

\begin{definition}[{\texorpdfstring{Def-BM.1 — Colour-sealed chain composites}{Def-BM.1 - Colour-sealed chain composites}}]\label{tag:Def-BM-1}
\textbf{Def.} \textbf{Constituent} $:=$ a triplet-attached bound chain [D3.8: one fundamental slot index $q^i$ (or dual $r_i$) at the attachment; Def-H.4: a Janus-closed bound chain, terminus locus; $\delta h\ge2$]. \textbf{Junction} $:=$ a shared vertex hosting one elementary invariant contraction of the constituents' slot indices [CT3.17: $\delta^i{}_j$ for a pair, $\varepsilon_{ijk}/\varepsilon^{ijk}$ for a triple; trivalent slot structure CT3.2].

%dsp-level:0
\[
\textbf{baryon:}\quad B(c_1,c_2,c_3)\ :=\ \varepsilon_{ijk}\,c_1^i\,c_2^j\,c_3^k
\]

%dsp-level:0
\[
\qquad [\mathrm{D3.9}\ \text{algebraic form, chain-realized at one junction}];
\]

%dsp-level:0
\[
\textbf{meson:}\quad M(c,\bar c)\ :=\ \delta^i{}_j\,c^j\,\bar c_i
\]

%dsp-level:0
\[
\qquad [\mathrm{D3.9};\ \bar c:\ \text{the conjugate constituent}\ (\omega\text{-reversed chain with dual attachment},\ \mathrm{D2.6},\ \mathrm{D3.8})].
\]

\textbf{Validity clause.} For valid ambient $C$ [D2.11], $\pi_C$ is total on $\mathrm{Colour}(U)$ [CT3.18 $+$ D2.11 applied to the colour ledger; D2.14]: every attachment has its contraction partners inside the composite; a restriction $\Sigma$ with $\partial_{\mathrm{col}}(\Sigma;C)\ne\varnothing$ has $\mathcal C(\Sigma;C)=0$ and its entire colour sector vanishes [D2.14, FT.19: all-or-nothing] — the mask admits only singlets [CT3.18]. \textbf{Read clause.} Each constituent retains its own terminus tuple $(q,\mathbf S,m_{\mathrm{hid}})$ [Def-H.9]; the junction adds the sealing-phase class and the lock content [Census FT.49: $\Theta_{\mathrm{int}}\in4\mathbb Z$ at the baryon-type cell; Census FT.35: junction $D^-$ content]; the card assembly is Def-BM.2.

\end{definition}

\begin{definition}[{\texorpdfstring{Def-BM.2 — The composite card}{Def-BM.2 - The composite card}}]\label{tag:Def-BM-2}
\textbf{Def.} The \textbf{card} of a sealed composite refines the floor-card row template [Census FT.28 — the minimal fermion cell; every sealed-composite card refines its row grammar: registration data $(\sigma_e,\kappa)$, $\Theta_{\mathrm{int}}$, winding, holonomy, aspect rows, mask dimension, rung counts, ternary core read] with:

%dsp-level:0
\[
(1)\ \text{constituent rows: the}\ \Pi_v\ \text{projections of Def-H.9 per constituent};
\]

%dsp-level:0
\[
\qquad (2)\ \text{the junction row: sealing-phase class and lock content}\ [\mathrm{Def\text{-}BM.1}\ \text{read clause}];
\]

%dsp-level:0
\[
(3)\ \text{the sealing-grammar row:}\ N_{\mathrm{strong}}
\]

%dsp-level:0
\[
= \log_2 g\bigl(|\mathrm{Colour}(\Sigma)|\bigr)\ \text{with the grammar counts}\ g\ [\mathrm{CT3.17a}(2)\text{–}(3);\ \mathrm{FT.41}(b)];
\]

%dsp-level:0
\[
(4)\ \text{the sealing-profile row}\ [\mathrm{FT.34}(5)\ \text{partition rule}];
\]

%dsp-level:0
\[
\qquad (5)\ \text{the winding-ledger rows}\ [\text{additive composition}\ \mathrm{FT.34}(3);\ \text{Census FT.43},\ \text{FT.51}].
\]

Card \textbf{values} are Census-cell data at the computed cells; card \textbf{grammar} is this definition. No card row is a species identification [P1, P5].

\end{definition}

\phantomsection
\section[Tier BM.1 — Observability of Sealed Composites]{\texorpdfstring{Tier BM.1 — Observability of Sealed Composites}{Tier BM.1 - Observability of Sealed Composites}}

\begin{theorem}[{\texorpdfstring{BM.1 — Confinement inheritance and the trace-only read}{BM.1 - Confinement inheritance and the trace-only read}}]\label{tag:BM-1}
\textbf{Thm.}

%dsp-level:0
\[
(1)\quad\textbf{sealing is the observability switch:}\ \ \text{a sealed composite is colour-complete on its own support}\ (\partial_{\mathrm{col}}=\varnothing\implies\mathcal C=1)\ \text{and its colour sector survives masking generically}
\]

%dsp-level:0
\[
\qquad\bigl[\mathrm{FT.39}(1):\ \mathcal C=1\Rightarrow F^\Sigma_{\mathrm{strong}}=F_{\mathrm{strong}}|_\Sigma\ne0\ \text{generically, exact exception locus}\ \{f=0\}\cup\{\mathrm{Colour}=\varnothing\}\bigr];
\]

%dsp-level:0
\[
\text{any unsealed constituent read vanishes: one dangling attachment zeroes the whole colour sector}\ [\mathrm{FT.19}];
\]

%dsp-level:0
\[
\ \neg\exists\ \text{triplet pole, physical norm}\ 0\ [\mathrm{CT3.22}];
\]

%dsp-level:0
\[
(2)\quad\textbf{the one-body trace law:}\ \ \text{every Hermitian slot operator on a sealed singlet reads only its trace — baryon:}\ \mathfrak T(H)
\]

%dsp-level:0
\[
= 3a\ \text{for the covariant class}\ H
\]

%dsp-level:0
\[
= aI_3+zP+z^*P^T;
\]

%dsp-level:0
\[
\ \text{meson:}\ 0\ \text{(including the trace part)}\ [\mathrm{FT.117}(1)\text{–}(3)];
\]

%dsp-level:0
\[
\text{hence the}\ C_3\text{-doublet datum}\ (B,\theta)\ \text{is annihilated on every colour-sealed object, and hadron-level triplets cannot carry the family channel}\ [\mathrm{FT.117}(2),(7)];
\]

%dsp-level:0
\[
(3)\quad\textbf{the sealed colour factor:}\ \ \text{singlet-sealed chains keep the undressed arena}\ 9\ \text{— the invariant lines are}\ 1\text{-dimensional}\ [\mathrm{FT.117}(5):\ \dim(\mathbb C^3)^{\otimes3\,\mathfrak{sl}(3)}=\dim(\mathbb C^3\otimes\bar{\mathbb C}^3)^{\mathfrak{sl}(3)}=1;\ \mathrm{Def\text{-}H.42}];
\]

%dsp-level:0
\[
(4)\quad\textbf{the visibility locus:}\ \ \text{the same covariant data are nonzero on the adjoint channel}\ ([H,E_{12}]\ne0),
\]

%dsp-level:0
\[
\ \text{but the adjoint is not a sealed singlet and carries no pole}\ [\mathrm{FT.117}(4);\ \mathrm{CT3.18};\ \mathrm{CT3.22}]\ \text{— the constituent-level channel data}\ (\delta_q=\tfrac2{27}\ [\mathrm{FT.119}])\ \text{are ledger-internal at hadron level}.
\]

\textbf{Req.} D2.11, D2.14, D3.8, D3.9 (Part 2); CT3.17, CT3.18, CT3.22 (Part 2); Def-BM.1; FT.19, FT.39, FT.116(3), FT.117, FT.119, Def-H.9, Def-H.42 (Part 3).

\textbf{Pf.}

%dsp-level:0
\[
(1):\ \mathrm{Def\text{-}BM.1}\ \text{validity clause}:\ \pi_C\ \text{total}\implies\partial_{\mathrm{col}}(\Sigma;C)
\]

%dsp-level:0
\[
= \varnothing\ \text{on the composite's own support}\implies\mathcal C(\Sigma;C)
\]

%dsp-level:0
\[
= 1\ [\mathrm{D2.14}];
\]

%dsp-level:0
\[
\qquad \mathrm{FT.39}(1):\ \text{generic nonvanishing with the exact exception locus}
\]

%dsp-level:0
\[
\mathrm{FT.19}:\ \mathcal C=0\implies F^\Sigma_{\mathrm{strong}}=0\ (\text{one null factor kills the slot product});\qquad\mathrm{CT3.22}:\ \neg\text{triplet pole}\ \wedge\ \text{norm}\ 0
\]

%dsp-level:0
\[
(2)\text{–}(3):\ \mathrm{FT.117}(1)\text{–}(3),
\]

%dsp-level:0
\[
(5)\ \text{verbatim at the composite support}:\ \varepsilon\ \text{spans}\ \wedge^3\mathbb C^3\ (\text{the determinant representation});
\]

%dsp-level:0
\[
\ \text{the one-body sum acts by}\ \mathfrak T(H)\ [\mathrm{CT3.16};\ \mathrm{CT3.7}:\ \mathfrak T(T_a)=0];
\]

%dsp-level:0
\[
\ \text{meson}:\ (H\otimes I-I\otimes H^T)\,
\]

%dsp-level:0
\[
\mathrm{vec}(I)
\]

%dsp-level:0
\[
= 0\ [\text{machine rows re-cited}]
\]

%dsp-level:0
\[
(4):\ \mathrm{FT.117}(4)\ \text{verbatim}:\ \text{adjoint}\notin\operatorname{Im}\mathcal P_M^{(3)}\ [\mathrm{CT3.18}];
\]

%dsp-level:0
\[
\qquad \text{ledger-internality}: (2)\text{'s annihilation on sealed supports}.\ \blacksquare
\]

\textbf{Phys.} Sealing is the observability switch: the colour sector becomes physical exactly when the $\varepsilon/\delta$ grammar completes — and then reads a single number, the trace. Everything the constituent ledger carries (channel angles, doublet data, spectator dressings) stays behind the seal: the hadron sees sums, never slots. This is confinement as an information property, run forward from the subnetwork theorems to the composite level.

\end{theorem}

\begin{theorem}[{\texorpdfstring{BM.2 — Junction phase and winding grammar}{BM.2 - Junction phase and winding grammar}}]\label{tag:BM-2}
\textbf{Thm.}

%dsp-level:0
\[
(1)\quad\textbf{baryon-type sealing phase:}\ \ \text{at the computed baryon-type junction cell,}\ \Theta_{\mathrm{int}}\in4\mathbb Z
\]

%dsp-level:0
\[
\qquad [\text{Census FT.49};\ \text{per materialized closed chain}\ \Theta(c)\in2\mathbb Z,\ \mathrm{FT.22}];
\]

%dsp-level:0
\[
(2)\quad\textbf{additive winding under composition:}\ \ \Gamma
\]

%dsp-level:0
\[
= c_1\cdot P\cdot c_2\cdot\bar P\ \text{materialized with}\ d_{12}
\]

%dsp-level:0
\[
= \Theta(P)\qquad \ge 0:
\]

%dsp-level:0
\[
\quad k(\Gamma)=k_1+k_2+d_{12}
\]

%dsp-level:0
\[
\qquad [\mathrm{FT.34}(3),\ \text{Census FT.50}(2)];
\]

%dsp-level:0
\[
\text{carrier level:}\ \ \Psi(\Gamma)
\]

%dsp-level:0
\[
= \mathrm{Sym}^{k_1+k_2+d_{12}}\subseteq\mathrm{Sym}^{k_1}\otimes\mathrm{Sym}^{k_2}\otimes\mathrm{Sym}^{d_{12}}\ \text{as maximal weight}\ [\mathrm{FT.60}(6),\ \text{under Def-H.38}];
\]

%dsp-level:0
\[
(3)\quad\textbf{sealing partition:}\ \ \text{attachment partitions into contraction groups follow the}\ \mathrm{CT3.17a}\ \text{grammar with counts}\ g(N)
\]

%dsp-level:0
\[
= 1,
\]

%dsp-level:0
\[
1,
\]

%dsp-level:0
\[
3,
\]

%dsp-level:0
\[
10,
\]

%dsp-level:0
\[
25,
\]

%dsp-level:0
\[
105\ (N=2..7)\ [\mathrm{CT3.17a}(2);\ \mathrm{FT.41}(b)];
\]

%dsp-level:0
\[
\ \text{the gluing law of composed floor cells is the constraint-rank law}\ [\text{Census FT.50}];
\]

%dsp-level:0
\[
(4)\quad\textbf{junction lock content:}\ \ \text{a junction at}\ n
\]

%dsp-level:0
\[
\ge 4\ \text{carries}\ D^-\ \text{lock content — retained,}\ \neg\text{accessible}\ [\text{Census FT.35};\ \mathrm{CT4.1}\text{–}\mathrm{CT4.3};\ \mathrm{D4.2};\ \mathrm{D1.11}].
\]

\textbf{Req.} FT.22, FT.34, FT.60, Def-H.38 (Part 3); CT3.17a, CT4.1–CT4.3 (Part 2); D4.2, D1.11 (Parts 1–2); Def-BM.1; Census FT.35, FT.49, FT.50; exact integer arithmetic.

\textbf{Pf.}

%dsp-level:0
\[
(1):\ \text{Census FT.49's card row (computed from FT.22/FT.24's quantization on the tetrahedron-face sealing grammar)};
\]

%dsp-level:0
\[
\qquad \text{no new arithmetic}
\]

%dsp-level:0
\[
(2):\ \mathrm{FT.34}(3)\ \text{verbatim}:\ \Theta(\Gamma)
\]

%dsp-level:0
\[
= \Theta(c_1)+\Theta(P)+\Theta(c_2)+\Theta(\bar P)
\]

%dsp-level:0
\[
= 2(k_1+d_{12}+k_2)\ [\Theta(P)=\Theta(\bar P)=d_{12};\ \mathrm{FT.23}];
\]

%dsp-level:0
\[
\qquad \mathrm{FT.60}(6):\ \text{CG inclusion on the cone}
\]

%dsp-level:0
\[
(3):\ \mathrm{CT3.17a}(2)\ \text{counts re-cited};
\]

%dsp-level:0
\[
\ \text{Census FT.50's junction-saturation and constraint-rank rows}
\]

%dsp-level:0
\[
(4):\ \text{Census FT.35's junction row};\ \mathrm{CT4.1}\text{–}\mathrm{CT4.3}:\ \text{absorption/retention at}\ n\ge4.\ \blacksquare
\]

\textbf{Phys.} The hadron ledger is additive exactly where the fermion ledger is: windings compose, phase classes stack in $4\mathbb Z$ at baryon-type junctions, and the junction is where the first lock content sits — the baryon-type cell is the corpus's first $D^-$-carrying floor. What the junction does not add is new grammar: the sealing partitions are the frozen colour counts, and the gluing is the constraint-rank law the floor cells already proved.

\end{theorem}

\phantomsection
\section[Tier BM.2 — The Family Group and the Invariance Laws of the…]{\texorpdfstring{Tier BM.2 — The Family Group and the Invariance Laws of the Seal}{Tier BM.2 - The Family Group and the Invariance Laws of the Seal}}

\begin{theorem}[{\texorpdfstring{BM.3 — The family group of the seal}{BM.3 - The family group of the seal}}]\label{tag:BM-3}
\textbf{Thm.}

%dsp-level:0
\[
(1)\quad\textbf{the forced group:}\ \ \text{the stabilizer of the baryon seal}\ \varepsilon_{ijk}\ \text{inside the native slot-relabelling group}\ S_3\ \text{is}\ A_3\cong C_3
\]

%dsp-level:0
\[
= S_3\cap SU(3)\ \text{— the unique maximal}\ \varepsilon\text{-preserving subgroup};
\]

%dsp-level:0
\[
\ \text{full}\ S_3\ \text{covariance is two-level-only}\ (\text{every}\ S_3\text{-covariant slot operator has spectrum}\ \{a+3b,a,a\})\ [\mathrm{FT.116}];
\]

%dsp-level:0
\[
(2)\quad\textbf{channel inheritance:}\ \ \text{every triplet-attached constituent carries the democratic}\ C_3\ \text{mechanism with angle}\ \delta_q
\]

%dsp-level:0
\[
= \tfrac2{27}\ [\mathrm{FT.119};\ \mathrm{Def\text{-}H.42}];
\]

%dsp-level:0
\[
\ \text{the triple-angle link}\ 3\delta_q
\]

%dsp-level:0
\[
= \delta_\ell
\]

%dsp-level:0
\[
= \tfrac29\ [\mathrm{FT.119}(4);\ \mathrm{FT.112}];
\]

%dsp-level:0
\[
\ \text{the cross-channel numerator identity}\ 16\delta_\nu
\]

%dsp-level:0
\[
= 9\delta_\ell
\]

%dsp-level:0
\[
= 27\delta_q
\]

%dsp-level:0
\[
= 2\ [\mathrm{FT.143}];
\]

%dsp-level:0
\[
(3)\quad\textbf{exclusions:}\ \ \text{the light-triplet read}\ (A,B)
\]

%dsp-level:0
\[
= (1,\sqrt2)\ \text{is excluded —}\ Q\ \text{is confined to}\ (0.54,0.59)\ \text{across the allowed ratio box while the channel normalization forces}\ Q
\]

%dsp-level:0
\[
= \tfrac23\ \text{phase-independently}\ [\mathrm{FT.120}(6);\ \text{Census row, quarantined external data}];
\]

%dsp-level:0
\[
\ \text{the heavy-triplet proximity rows are audit-only, no identification}\ [\text{Census FT.120};\ \mathrm{P1},\ \mathrm{P5}].
\]

\textbf{Req.} CT3.17 (Part 2); FT.112, FT.116, FT.119, FT.143, Def-H.42 (Part 3); Census FT.120; exact $\mathbb Q(i)$/$\mathbb Q(\sqrt3)$ arithmetic (rows re-cited).

\textbf{Pf.} $(1)$: FT.116 verbatim — the relabelling lattice, the two-level obstruction, the $\varepsilon$-selection [machine rows re-cited]. $(2)$: FT.119(4)/FT.112(3)/FT.143 angle ledgers verbatim; per-constituent application by Def-BM.1's constituent structure (each constituent is a triplet-attached bound chain in FT.119's sense). $(3)$: Census FT.120's exclusion rows (quarantine discipline [FT.73(5) pattern]). $\blacksquare$

\textbf{Phys.} The seal picks its own symmetry: the only slot-relabelling subgroup preserving the baryon contraction is the cyclic three — the same three that carries the channel angles on every constituent. Hadron-level family structure is not imposed on the junction; it is the junction's stabilizer, with its phase fixed to two-thirds and its alternatives excluded by theorem.

\end{theorem}

\begin{theorem}[{\texorpdfstring{BM.4 — The invariance laws of the seal}{BM.4 - The invariance laws of the seal}}]\label{tag:BM-4}
\textbf{Thm.}

%dsp-level:0
\[
(1)\quad\textbf{colour invariance:}\ \ B,
\]

%dsp-level:0
\[
\ M\ \text{are colour singlets}\ [\mathrm{CT3.15}\text{–}\mathrm{CT3.17}:\ \{B,M\}\ \text{singlets};\ \{\delta,\varepsilon,\bar\varepsilon\}\ \text{the invariant grammar}];
\]

%dsp-level:0
\[
\ \text{the mask admits only singlets}\ [\mathrm{CT3.18}];
\]

%dsp-level:0
\[
\ \text{the colour sector is rigid and contributes}\ \neg\mathrm{P4}\ \text{data}\ [\mathrm{CT3.19}];
\]

%dsp-level:0
\[
(2)\quad\textbf{transport compatibility:}\ \ \text{the seed is colour-neutral by theorem}\ [\mathrm{EW.32}:\ T_a\mathsf v=0];
\]

%dsp-level:0
\[
\ \text{the weak action on the extended carrier is seed-factor-only}\ [W_i\mapsto\tfrac12\sigma_i\otimes I_\omega,\ \mathrm{FT.123}(1);\ \text{Def-H.34}]\ \text{and the diagonal}\ \mathrm{Sym}^k\ \text{transport is colour-blind}\ [\mathrm{FT.60};\ \mathrm{Def\text{-}H.38}]\ \implies\ \text{constituent channel data and multiplicity coupling compose with the seal without colour obstruction};
\]

%dsp-level:0
\[
(3)\quad\textbf{the parity side:}\ \ \text{the scalar-CPT action on each constituent's orientation doublet,}\ \mathrm{CPT}|_{\mathcal H_\omega}
\]

%dsp-level:0
\[
= (-1)^{|E|}e^{i\pi\Theta}\,
\]

%dsp-level:0
\[
\mathrm{id},
\]

%dsp-level:0
\[
\ \text{is chirality-even}\ [\mathrm{FT.134}(3)];
\]

%dsp-level:0
\[
\ \text{pole multiplets are}\ m\leftrightarrow-m\ \text{symmetric and the ceilings parity-blind}\ [\mathrm{FT.138}]\ \implies\ \text{both laws pass to sealed composites constituent-wise}\ [\mathrm{Def\text{-}BM.1}\ \text{product structure}].
\]

\textbf{Req.} CT3.15–CT3.19 (Part 2); EW.32 (Part 2); FT.60, FT.123, FT.134, FT.138, Def-H.34, Def-H.38 (Part 3); Def-BM.1.

\textbf{Pf.}

%dsp-level:0
\[
(1):\ \mathrm{CT3.15}\text{–}\mathrm{CT3.17}\ \text{singlet rows re-cited};
\]

%dsp-level:0
\[
\ \mathrm{CT3.18}:\ \operatorname{Im}\mathcal P_M^{(3)}\ \text{admits only singlets};
\]

%dsp-level:0
\[
\ \mathrm{CT3.19}:\ \text{open quantities}\qquad = \varnothing
\]

%dsp-level:0
\[
(2):\ \mathrm{EW.32}:\ T_a\mathsf v=0\ [\text{the seed's colour neutrality}];
\]

%dsp-level:0
\[
\qquad \mathrm{FT.123}(1):\ [W_i\otimes I_\omega,\ \cdot]\ \text{acts on the seed/orientation factors only};
\]

%dsp-level:0
\[
\qquad \mathrm{Def\text{-}BM.1}:\ \text{constituents are independent chains with own carriers}\implies\text{product composition}
\]

%dsp-level:0
\[
(3):\ \mathrm{FT.134}(3),
\]

%dsp-level:0
\[
\ \mathrm{FT.138}\ \text{verbatim per constituent};
\]

%dsp-level:0
\[
\ \text{composite-level passage}: \mathrm{Def\text{-}BM.1}'\text{s product structure}\ (\text{each factor carries its own doublet and}\ \Theta\text{-summand};\ (-1)^{|E|}e^{i\pi\Theta}\ \text{multiplicative over disjoint constituents}).\ \blacksquare
\]

\textbf{Phys.} Both algebras that act on a composite leave the seal intact: colour SU(3) has nothing to do inside a singlet, and weak transport never touches the colour factor — so a baryon's card assembles constituent by constituent, and the parity laws pass through the seal unbroken. The rigidity theorem closes the ledger: the colour side of the hadron grammar contributes no free datum at all.

\end{theorem}

\phantomsection
\section[Tier BM.3 — Consumption and Open Constructions]{\texorpdfstring{Tier BM.3 — Consumption and Open Constructions}{Tier BM.3 - Consumption and Open Constructions}}

\pseudohead{theorem}{\texorpdfstring{BM.5}{BM.5}}{\texorpdfstring{The census-consumption law}{The census-consumption law}}\label{tag:BM-5}
\textbf{Thm.}

%dsp-level:0
\[
(1)\quad\textbf{the substrate:}\ \ \text{the enumerative substrate of this era is the Census block — the}\ 25\ \text{Category-B cells with their configuration spaces, counts, and consumption rows}\ [\text{Census charter and §Consumption}];
\]

%dsp-level:0
\[
\ \text{every construction of this part consumes cells by tag citation; no cell is peripheral};
\]

%dsp-level:0
\[
(2)\quad\textbf{the family table:}
\]

%tbl-level:3
{\scriptsize\setlength{\tabcolsep}{4pt}
\begin{longtable}{>{\raggedright\arraybackslash}p{\dimexpr0.1200\textwidth-2\tabcolsep\relax}|>{\raggedright\arraybackslash}p{\dimexpr0.1127\textwidth-2\tabcolsep\relax}|>{\raggedright\arraybackslash}p{\dimexpr0.2302\textwidth-2\tabcolsep\relax}|>{\raggedright\arraybackslash}p{\dimexpr0.5371\textwidth-2\tabcolsep\relax}}
\toprule
census family & cells & configuration space supplied & consumption at the composite level \\
\addlinespace[2pt]
\midrule
\endfirsthead
\toprule
census family & cells & configuration space supplied & consumption at the composite level \\
\addlinespace[2pt]
\midrule
\endhead
\bottomrule
\endlastfoot
floor cards and ladders & FT.28, FT.32, FT.35, FT.42–FT.45, FT.49–FT.51, FT.63 & the computed fermion floors with full read columns & the cell grammar of sealed composites (Def-BM.2): baryon-type junctions are gluings of these cards; the winding ledgers and duality splits are the multiplet bookkeeping; the $\Sigma_{12}$/tetrahedron-face cells carry the first colour-active and $4\mathbb Z$-phased junctions \\
loop and cover classes & FT.64, FT.65, FT.69, FT.78, FT.96 & admissible-loop lattices, coarse-grained degeneracies, amplitude censuses, $\operatorname{Aut}$-orbits with chirality tags & the enumeration substrate of sealed-object path integrals: sector conditioning $(q,\allowbreak S)$ [Def-H.32(4)] runs over these pools; the chirality-tagged orbits feed the parity reads [BM.4(3)] \\
pole and mass multiplets & FT.76 & defect-angle pole reads under both convention variants, with the ceiling & the pole-side mass cells and the hard range ceiling every composite read inherits [Def-H.39(2)]; the convention variants are priced — the choice converts them into data [BM.6(4)] \\
form-factor tables & FT.98, FT.99 & polygon and Platonic form factors with growth and dominance rows & the size-bound range checks for extended composite skeletons [Def-H.39 ceilings] \\
audit ledgers & FT.81, FT.83, FT.84, FT.114, FT.120 & the read-class inventory, the rung-emergence ledger, the ceiling taxonomy, the quarantined comparison rows & the card grammar and the P1/P5-compliant comparison template; the activation map at $n=\allowbreak 3$/$n\ge\allowbreak 4$; the read-species taxonomy any new read must join \\
topology-blind cells & FT.105 & graph counts and canonical skeletons from first principles, budget-regularized & the regulated topology pool of the composite path integral — the space its fugacity formalism sums over [Def-H.32(1)] \\
\end{longtable}
}

%dsp-level:0
\[
(3)\quad\textbf{the conditioning law:}\ \ \text{sector-conditioned composite sums run on Def-H.32(4)'s labels}\ (q,S)\ \text{over}\ (2)\text{'s pools};
\]

%dsp-level:0
\[
\ \text{the amplitude/norm dichotomy and the amplitude}\leftrightarrow\text{spectral bridge residue are item-13 content}\ [\mathrm{FT.67}\text{–}\mathrm{FT.68};\ \mathrm{FT.80}].
\]

\textbf{Req.} The Census block (all 25 cells); Def-H.32, FT.67, FT.68, FT.80 (Part 3); Def-H.39 (Part 3); Def-BM.2

\textbf{Pf.} $(2)$: the consumption map — each family row references its cells by tag; every Census tag occurs in exactly one family row [exact set arithmetic: $11+5+1+2+5+1=25$]. $(1),(3)$: charter and citation scans. $\blacksquare$

\textbf{Phys.} The census is not an appendix: it is the table of configuration spaces the baryomesonic era enumerates — cards to glue, orbits to condition on, multiplets whose ceilings every composite read inherits. The consumption law states, family by family, which space feeds which construction, and the count closes: twenty-five cells, no residue.

\begin{theorem}[{\texorpdfstring{BM.6 — The construction interfaces (specifications)}{BM.6 - The construction interfaces (specifications)}}]\label{tag:BM-6}
\textbf{Def.} The constructions this era requires and does not yet carry, each stated as a specification (what a construction must supply); none is a declaration, and none is adopted:

%dsp-level:0
\[
(1)\quad\textbf{the multi-chain additive charge read}:\ \text{a read}\ Q\ \text{on sealed composites with}\ (i)\ Q|_{\text{single chain}}
\]

%dsp-level:0
\[
= q_{\mathrm{core}}\ [\mathrm{FT.40};\ \mathrm{EM.9}],
\]

%dsp-level:0
\[
\ (ii)\ \text{additivity under disjoint sealing},
\]

%dsp-level:0
\[
\ (iii)\ \text{observability only on sealed supports}\ [\mathrm{CT3.22}\ \text{discipline}].\ \text{Status: not constructed; the single-chain read is derived; two end registers per chain are counted}\ [\mathrm{FT.41}(a)];
\]

%dsp-level:0
\[
(2)\quad\textbf{the two-chain exchange space}:\ \text{a state space on ordered pairs of bound chains with an exchange involution; the single-chain chiral grade is}\ \neg\text{an exchange sign}\ [\mathrm{FT.125}(5);\ \mathrm{FT.80}(5):\ \neg\exists\ \text{exchange operator in the inventory};\ \mathrm{FT.89}(2)\ \text{kind discipline}];
\]

%dsp-level:0
\[
\ \text{spin–statistics content is external until constructed};
\]

%dsp-level:0
\[
(3)\quad\textbf{the}\ \aleph\text{-domain choice}:\ \text{literal}\ [\mathrm{D3.3}]\ \text{vs chain-local}\ \mathrm{E1}/\mathrm{E2}\ [\mathrm{Def\text{-}H.27}]\ \text{— both path integrals computed and pairwise inequivalent}\ [\mathrm{FT.67},\ \mathrm{FT.74},\ \mathrm{FT.75}];
\]

%dsp-level:0
\[
\ \text{sealed junctions are the natural site of the residual class-(iii) edges}\;
\]

%dsp-level:0
\[
\ \text{the choice requires a logged declaration before use};
\]

%dsp-level:0
\[
(4)\quad\textbf{the defect corner convention}:\ \text{variant A (face winding}\ k_f)\ \text{/ variant B (vertex-figure winding}\ k_{vf})\ \text{— both multiplet tables priced}\ [\text{Census FT.76}(1)\text{–}(2)];
\]

%dsp-level:0
\[
\ \text{a declaration converts them into data immediately};
\]

%dsp-level:0
\[
(5)\quad\textbf{the dynamics rates}\ [\text{Cal-F}\ r]:\ \text{a dynamics construction supplying}\ J^\pm\ \text{rates, or the free datum retained; the tick-rate/field coupling requires emergent time}\ [\mathrm{FT.131}(3)];
\]

%dsp-level:0
\[
(6)\quad\textbf{the momentum-space}\ \Sigma\text{-resolvent}:\ \text{plane-wave boundary-channel states on the coarse-grained channel}\ [\mathrm{Def\text{-}H.13}\ \text{pattern}];
\]

%dsp-level:0
\[
\ \text{constructed at BM.134 (Def-BM.20); consumed by scattering reads of sealed objects};
\]

%dsp-level:0
\[
(7)\quad\textbf{the embedding datum}:\ \text{a circumradius or moment-of-inertia read for rotor-type reads; the ratios are embedding-free and computed}\ [\mathrm{FT.87}(6)];
\]

%dsp-level:0
\[
\ \text{adoption would not lift the}\ 4\hat\mu'^2\ \text{ceiling}\ [\mathrm{FT.76}(3);\ \mathrm{FT.87}(1)];
\]

%dsp-level:0
\[
(8)\quad\textbf{the frame section}:\ e\mapsto\mathcal F_e\ \text{over each skeleton}\ [30\ \text{native frames catalogued},\ \mathrm{FT.93}]\ \text{plus its}\ \mathrm{EW.32}\text{-transport coupling and a global Dirac-type operator on the extended carrier}\ [\mathrm{Def\text{-}H.34}];
\]

%dsp-level:0
\[
\ \text{the weak-chirality pincer on the pair arena is untouched}\ [\mathrm{FT.93}(3)].
\]

\textbf{Req.} The cited tags

\textbf{Pf.} Specification content: each row states its requirement and cites its anchors; no truth claim is made; anchor existence machine-checked [\texttt{checks/\allowbreak{}verify\_\allowbreak{}era\_\allowbreak{}structure.py}]. $\blacksquare$

\textbf{Phys.} The era's debts are itemized as build-specs: what a charge read must satisfy, what an exchange space must carry, which two conventions must be chosen and what each choice costs. Nothing here is adopted; everything here is priced. The rung closes with its obligations named.

\end{theorem}

\phantomsection
\section[Tier BM.4 — Noether's Theorem from the Janus Seal Topology:…]{\texorpdfstring{Tier BM.4 — Noether's Theorem from the Janus Seal Topology: the Discrete Conservation Law, the Continuum Promotion, the Stress-Energy of a Sealed Composite, and the Canonical-Pair Obstructions}{Tier BM.4 - Noether's Theorem from the Janus Seal Topology: the Discrete Conservation Law, the Continuum Promotion, the Stress-Energy of a Sealed Composite, and the Canonical-Pair Obstructions}}

\begin{definition}[{\texorpdfstring{Def-BM.3 — The species-identification layer (declared)}{Def-BM.3 - The species-identification layer (declared)}}]\label{tag:Def-BM-3}
\textbf{Def.} A \textbf{species name} $:=$ a label attached to a derived structural tuple of a sealed composite. The layer is constitutive in exactly one respect — the name$\leftrightarrow$tuple assignment — and non-derived in that same respect: no theorem of Parts 1–4 forces any name onto any tuple, and no name carries a datum into a derivation.

%dsp-level:0
\[
(1)\quad\textbf{the tuple a name attaches to:}
\]

%dsp-level:0
\[
\qquad \bigl(\text{junction grammar}\ [\mathrm{Def\text{-}BM.1}],\ \ \text{sealing-phase class}\ [\mathrm{BM.2}(1)],\ \ \text{charge class}\ [\mathrm{EM.9};\ \text{clause}\ (3)],\ \ m_{\mathrm{hid}}=E_{\mathrm{int}}\ln2/\mathfrak h\ [\mathrm{FT.44}(4)],\ \ w=p/\rho\ [\mathrm{CT4.5};\ \mathrm{BM.10}(3)]\bigr);
\]

%dsp-level:0
\[
(2)\quad\textbf{the declared map (baryon-type, }\varepsilon\text{-sealed 3-junctures):}
\]

%dsp-level:0
\[
\qquad \textbf{neutron}\ :=\ \varepsilon_{ijk}c_1^ic_2^jc_3^k\ \text{in the neutral charge class}\ (q_{\mathrm{net}}=0);
\]

%dsp-level:0
\[
\textbf{proton}\ :=\ \varepsilon_{ijk}c_1^ic_2^jc_3^k\ \text{in the unit charge class}\ (q_{\mathrm{net}}=+1);
\]

%dsp-level:0
\[
\qquad \textbf{antiproton}\ :=\ \omega\text{-image of the proton}\ [\mathrm{D2.6}];
\]

%dsp-level:0
\[
(3)\quad\textbf{the declared map (meson-type, }\delta\text{-sealed 2-junctures):}
\]

%dsp-level:0
\[
\qquad \pi^0\ :=\ \delta^i{}_j\,c^j\bar c_i\ \text{in the neutral class};
\]

%dsp-level:0
\[
\qquad \pi^{\pm}\ :=\ \delta^i{}_j\,c^j\bar c_i\ \text{in the}\ q_{\mathrm{net}}=\pm1\ \text{classes};
\]

%dsp-level:0
\[
(4)\quad\textbf{the charge-class clause (conditional, declared):}
\]

%dsp-level:0
\[
\qquad q_{\mathrm{net}}\ \text{is the additive multi-chain read specified but}\ \neg\text{constructed}\ [\mathrm{BM.6}(1)]:
\]

%dsp-level:0
\[
\text{the class labels}\ 0,
\]

%dsp-level:0
\[
\ \pm1\ \text{of clauses}\ (2)\text{–}(3)\ \text{are the ternary alphabet of a single end read}\ [q_{\mathrm{EM}}\in\{-1,0,+1\}\ [\mathrm{D3.4},\ \mathrm{EM.9},\ \mathrm{P2}]]\ \text{carried to the composite by declaration};
\]

%dsp-level:0
\[
\text{the constituent alphabet remains ternary: the fractional read}\ (\pm\tfrac13,\ \pm\tfrac23)\ \text{is deficit}\ 1'\ \text{content}\ [\mathrm{FT.125}(4)]\ —\ \neg\text{supplied,}\ \neg\text{assumed here};
\]

%dsp-level:0
\[
(5)\quad\textbf{the naming-precedent clause:}
\]

%dsp-level:0
\[
\qquad \mathrm{D3.7}\ \text{names two transit structures (photon}\ :=\ \text{trivial chain}\ \subseteq V_2;\ \text{gluon}\ :=\ \text{trivial chain}\ \cap\,V_3\ne\varnothing)\ \text{and}\ \mathrm{D3.9}\ \text{names the two sealings (baryon, meson)}:
\]

%dsp-level:0
\[
\text{this definition extends that precedent from transit structures and sealings to charge classes of sealed composites};
\]

%dsp-level:0
\[
\qquad \text{a name is a definition-side label,}\ \neg\text{a measurement};
\]

%dsp-level:0
\[
(6)\quad\textbf{the data bar (standing):}
\]

%dsp-level:0
\[
\qquad \neg\exists\ \text{a measured external value in any clause of Parts 1–4}\ [\mathrm{P5}\ \text{stands}];
\]

%dsp-level:0
\[
\qquad \text{external comparison data remain quarantined to the audit rows}\ [\text{Census FT.114},\ \text{FT.120}];
\]

%dsp-level:0
\[
\text{the names of clauses}\ (2)\text{–}(3)\ \text{carry}\ \neg\text{mass value,}\ \neg\text{coupling value,}\ \neg\text{lifetime,}\ \neg\text{radius}:\ \text{every number attached to a named tuple is corpus-derived}\ (m_{\mathrm{hid}},\ \Theta_{\mathrm{int}},\ w,\ E_{\mathrm{int}});
\]

%dsp-level:0
\[
(7)\quad\textbf{scope:}
\]

%dsp-level:0
\[
\qquad \text{species names occur in Part 4 only}\ —\ \text{in this definition, in the}\ \textbf{Phys}\ \text{and}\ \textbf{Rmk}\ \text{prose of this part, and nowhere in a}\ \textbf{Thm}\ \text{clause of Parts 1–3}\ [\text{those tiers' own label-discipline rows stand unmodified}];
\]

%dsp-level:0
\[
\text{extension of the map beyond clauses}\ (2)\text{–}(3)\ \text{requires a logged declaration first}\ [\text{Part 4 §Cal-BM rule}];
\]

%dsp-level:0
\[
\qquad \text{the layer adds}\ \neg\text{dimensionful scale}\ [\mathrm{P1}]\ \text{and}\ \neg\text{P4 parameter}.
\]

\textbf{Rmk (identification discipline).} The layer answers \textit{which name sits on which derived tuple}; it does not answer \textit{which species sits in which gauge representation on which chiral ray} — that assignment rule is application-layer content; its non-derivability here is proven [FT.125(4)]. Naming the neutron is a labelling act on the $\varepsilon$-sealed neutral class; it derives no neutron property, imports no neutron datum, and closes no open row.

\textbf{Req.} D2.6, D3.4, D3.7, D3.9 (Parts 1–2); CT4.5 (Part 2); EM.9 (Part 2); Def-BM.1, BM.2(1), BM.6(1) (this part); FT.44(4), FT.125(4) (Part 3, Census); P1, P2, P5 (Parts 1–2); Census FT.114, FT.120 (quarantine rows)

\textbf{Pf.}

%dsp-level:0
\[
\text{Definition block: the constitutive clause is the name}\leftrightarrow\text{tuple assignment of}\ (2)\text{–}(3);
\]

%dsp-level:0
\[
\ \text{the non-derived clauses are}\ (4)\ \text{(the charge class is declared pending BM.6(1)) and the name}\leftrightarrow\text{tuple direction of}\ (1)\ [\text{no theorem forces it}];
\]

%dsp-level:0
\[
\ \text{the standing bars}\ (6)\text{–}(7)\ \text{are scope clauses cited at use}.\ \blacksquare
\]

\textbf{Phys.} The baryomesonic era names its objects. A name is pinned to a tuple the corpus already derives — the junction grammar, the sealing phase class, the charge class, the informational mass, the equation of state — and the pinning itself is the declared part. The neutron is the $\varepsilon$-sealed three-juncture in the neutral charge class; the proton is the same structure in the unit class; the pion-type triplet is the $\delta$-sealed pair in its three classes. Nothing measured enters: the mass read is $E_{\mathrm{int}}\ln2/\mathfrak h$, the phase class is $4\mathbb Z$, the equation of state is $w=0$, and the charge class is the ternary alphabet carried up by declaration while the additive read is still an open construction. The layer extends a precedent the corpus already set when it named the photon and the gluon as chain classes.

\end{definition}

\begin{theorem}[{\texorpdfstring{BM.7 — The discrete conservation law is Janus topology: no symmetry argument, no action, no limit}{BM.7 - The discrete conservation law is Janus topology: no symmetry argument, no action, no limit}}]\label{tag:BM-7}
\textbf{Thm.}

%dsp-level:0
\[
(1)\quad\textbf{the balance is definitional, and the seal ledger realizes it uniquely:}
\]

%dsp-level:0
\[
\qquad \text{valid universe}\iff\text{network closed}\ \wedge\ \#J^+
\]

%dsp-level:0
\[
= \#J^-\ \wedge\ \in\operatorname{Im}(\text{active mask})\ \wedge\ \neg\exists\,
\]

%dsp-level:0
\[
\text{dangling boundary}\ [\mathrm{D2.11},\ \mathrm{D1.3}];
\]

%dsp-level:0
\[
J^+:\ (\Delta n_{\mathcal R},\delta h,\delta\tau)=(+1,+1,+1),
\]

%dsp-level:0
\[
\qquad J^-:\ (-1,-1,-1)
\]

%dsp-level:0
\[
\qquad [\mathrm{CT0.3}]:
\]

%dsp-level:0
\[
\qquad \text{the seal excess}\ \#J^+-\#J^-=n_{\mathcal R}\ \text{is the ledger's own occupation count};
\]

%dsp-level:0
\[
\text{the D1.2 transition domains force strict alternation}\ (J^-\ \text{undefined at}\ n_{\mathcal R}=0;\ J^+\ \text{undefined at}\ n_{\mathcal R}=1):\ \text{the single-pair history of length}\ n\ \text{is unique}\ \forall n\ [\mathrm{FT.144}(2);\ \text{machine}];
\]

%dsp-level:0
\[
\text{the completed history of length}\ 2m\ \text{is unique and balanced:}\quad(J^+J^-)^m,
\]

%dsp-level:0
\[
\qquad \#J^+=\#J^-=m
\]

%dsp-level:0
\[
\qquad [\mathrm{FT.144}(2),\ \mathrm{FT.156}(5);\ \text{machine}];
\]

%dsp-level:0
\[
\text{the prefix law}\ \#J^+(\text{prefix})\ge\#J^-(\text{prefix})\ \text{holds on every admissible history, with equality}
\]

%dsp-level:0
\[
\iff n_{\mathcal R}
\]

%dsp-level:0
\[
= 0\ [\mathrm{FT.144}(2);\ \text{machine}];
\]

%dsp-level:0
\[
\qquad \text{the pure order reversal of every nonempty completed history leads with}\ J^-:\ \text{inadmissible at its first event}\ [\mathrm{FT.144}(4)(\mathrm{ii});\ \text{machine}];
\]

%dsp-level:0
\[
\textbf{the relaxation census}\ (\text{prefix law alone, alternation dropped}):\quad\#\{\text{completed words of length}\ 2m\}
\]

%dsp-level:0
\[
= \mathrm{Cat}_m
\]

%dsp-level:0
\[
= \frac1{m+1}\binom{2m}m
\]

%dsp-level:0
\[
= 1,
\]

%dsp-level:0
\[
1,
\]

%dsp-level:0
\[
2,
\]

%dsp-level:0
\[
5,
\]

%dsp-level:0
\[
14,\qquad 42
\]

%dsp-level:0
\[
\qquad [\text{machine, exhaustive}\ \le12];
\]

%dsp-level:0
\[
\mathrm{Cat}_m>1
\]

%dsp-level:0
\[
\iff m
\]

%dsp-level:0
\[
\ge 2:\ \text{the D1.2 domains are strictly stronger than the prefix law from}\ 2m
\]

%dsp-level:0
\[
= 4\ \text{on}\ —\ \text{the uniqueness of}\ (1)\ \text{is domain content,}\ \neg\text{balance content};
\]

%dsp-level:0
\[
\textbf{two pairs:}\quad\#\{\text{completed interleavings of total length}\ 2n\}
\]

%dsp-level:0
\[
= \sum_{m_0+m_1=n}\frac{(2n)!}{(2m_0)!\,(2m_1)!}
\]

%dsp-level:0
\[
= 2,
\]

%dsp-level:0
\[
\ 8,
\]

%dsp-level:0
\[
\ 32\quad(2n=2,4,6)
\]

%dsp-level:0
\[
\qquad [\mathrm{FT.144}(2);\ \text{machine: exhaustive}];
\]

%dsp-level:0
\[
\text{balance and the prefix law survive every interleaving}\ [\text{machine}];
\]

%dsp-level:0
\[
\qquad \text{the balanced sector is}\ a_+-a_-=|\mathcal R'|-|\mathcal R|=0\ \text{of the budget identity}\ [\mathrm{FT.108}(4)];
\]

%dsp-level:0
\[
(2)\quad\textbf{the endpoint read is a boundary term, exactly:}
\]

%dsp-level:0
\[
\qquad q_{\mathrm{EM}}(c)=\frac{\sigma_{s(c)}-\sigma_{t(c)}}2\in\{-1,0,+1\}
\]

%dsp-level:0
\[
\qquad [\mathrm{D3.4},\ \mathrm{EM.9};\ \text{all three values attained}];
\]

%dsp-level:0
\[
\textbf{the telescoping identity:}\quad\text{on a materialized chain}\ c=(e_1,\dots,e_k),\ \ t(e_i)=s(e_{i+1})\ [\mathrm{D1.4}]:
\]

%dsp-level:0
\[
\qquad q_{\mathrm{EM}}(c)=\sum_{i=1}^{k}j(e_i),
\]

%dsp-level:0
\[
\qquad j(e):=\frac{\sigma_{s(e)}-\sigma_{t(e)}}2
\]

%dsp-level:0
\[
\qquad [\text{machine: exhaustive}\ k\le7];
\]

%dsp-level:0
\[
\text{interior vertex signs cancel pairwise; the sum is the difference of the two end reads}\ —\ \text{the charge at the start equals the charge at the end};
\]

%dsp-level:0
\[
\textbf{the closed-chain read:}\quad\partial c=0\implies s(c)=t(c)\implies\sigma_{s(c)}=\sigma_{t(c)}\implies q_{\mathrm{EM}}(c)=0\ \text{exactly}
\]

%dsp-level:0
\[
\qquad [\mathrm{D1.4};\ \text{machine: every closed configuration}];
\]

%dsp-level:0
\[
\text{total charge vanishes in every valid universe: each end is sealed against an oppositely oriented end carrying the same sign}\ [\mathrm{EM.11}];
\]

%dsp-level:0
\[
(3)\quad\textbf{the ground is the chain complex, not an extremum:}
\]

%dsp-level:0
\[
\qquad \partial c:=t(c)-s(c)\ \text{(formal difference of vertices)},
\]

%dsp-level:0
\[
\qquad \partial^2=0\ \text{since}\ \partial(\text{vertex})=0
\]

%dsp-level:0
\[
\qquad [\mathrm{D1.4};\ \text{machine}];
\]

%dsp-level:0
\[
\text{closed}\iff\partial c=0\ \wedge\ \#J^+=\#J^-\ [\mathrm{D1.4}];
\]

%dsp-level:0
\[
\qquad \ker\partial=\{\text{closed chains}\},
\]

%dsp-level:0
\[
\qquad \text{every boundary has total multiplicity}\ 0\ [\text{machine}];
\]

%dsp-level:0
\[
\text{conservation is the kernel statement of an integer boundary map: no action functional, no variational derivative and no continuous parameter enters}\ (1)\text{–}(3);
\]

%dsp-level:0
\[
(4)\quad\textbf{the Noether route is unavailable for this read (obstruction):}
\]

%dsp-level:0
\[
\qquad \text{a Noether charge requires a one-parameter group action}\ \mathbb R\to\operatorname{Aut}(\text{the read's value space})\ [\text{math}];
\]

%dsp-level:0
\[
\text{the value space of}\ q_{\mathrm{EM}}\ \text{is the ternary set}\ \{-1,0,+1\},
\]

%dsp-level:0
\[
\ \text{so the action lands in}\ \operatorname{Perm}\{-1,0,+1\}
\]

%dsp-level:0
\[
= S_3\ [\mathrm{EM.9}];
\]

%dsp-level:0
\[
U(1)
\]

%dsp-level:0
\[
= \varinjlim_m\mathbb Z/m
\]

%dsp-level:0
\[
= \mathbb Q/\mathbb Z\ \text{is divisible}\ [\text{math}]:\ \text{its image in a finite group is divisible, hence trivial}\ —\ \text{no transposition of}\ S_3\ \text{has a square root, no}\ 3\text{-cycle has a cube root}\ [\text{machine: the full root census of}\ S_3];
\]

%dsp-level:0
\[
|\operatorname{hom}(\mathbb Z/m,S_3)|
\]

%dsp-level:0
\[
= \#\{g\in S_3:\operatorname{ord}(g)\mid m\}
\]

%dsp-level:0
\[
= 1,
\]

%dsp-level:0
\[
4,
\]

%dsp-level:0
\[
3,
\]

%dsp-level:0
\[
4,
\]

%dsp-level:0
\[
6,
\]

%dsp-level:0
\[
6\ (m=1,2,3,4,6,12)\ [\text{machine}]:\ \text{every finite stage admits nontrivial homs, and the only projectively compatible family is the trivial one};
\]

%dsp-level:0
\[
\therefore\ \neg\exists\,\text{a nontrivial continuous}\ U(1)\ \text{action on the charge alphabet}:\qquad q_{\mathrm{EM}}\ \text{is}\ \neg\text{a Noether charge};
\]

%dsp-level:0
\[
(5)\quad\textbf{priority:}
\]

%dsp-level:0
\[
\qquad (1)\text{–}(3)\ \text{are exact, discrete, integer-valued and hold at every scale, including}\ \mathfrak h\ne0;
\]

%dsp-level:0
\[
\qquad \text{the differential law}\ \partial_\mu J^\mu=0\ \text{is their continuum limit}\ [\mathrm{BM.11}]\ —\ \text{derived from them,}\ \neg\text{converse};
\]

%dsp-level:0
\[
\text{the converse fails:}\ \partial_\mu J^\mu
\]

%dsp-level:0
\[
= 0\ \text{is a density statement and does not force the integer balance at finite cell size}\ [\mathrm{BM.11}(5)].
\]

\textbf{Rmk (label discipline).} $J^{\pm}$, \textit{seal}, \textit{source}, \textit{sink} and \textit{current} name ledger structures [D1.2, D1.3, CT0.3]; the topology labels are internal bookkeeping [Def-H.32(4) pattern]. Where clause (2) is read on a named composite, the name is Def-BM.3's label on the derived tuple; the theorem's content is the tuple.

\textbf{Req.} A4, A6, A7 (Part 1); D1.2, D1.3, D1.4, D2.11, D3.3, D3.4 (Parts 1–2); CT0.3, CT0.9 (Part 1); EM.9, EM.11 (Part 2); P2 (Part 1); Def-BM.1, Def-BM.3 (this part); FT.108(4), FT.144(2), FT.144(4)(ii), FT.156(5) (Part 3); exhaustive ledger-word enumeration, the Catalan census, multinomial interleaving counts, hom-counts and the root census of $S_3$, divisibility of $\mathbb Q/\mathbb Z$ [math; machine: \texttt{checks/\allowbreak{}verify\_\allowbreak{}H\_\allowbreak{}noether.py} §A].

\textbf{Pf.}

%dsp-level:0
\[
(1):\ \mathrm{D2.11}/\mathrm{D1.3}\ \text{verbatim};
\]

%dsp-level:0
\[
\qquad \mathrm{CT0.3}\ \text{increments verbatim}\implies\#J^+-\#J^-=n_{\mathcal R}\ [\text{exact integers}];
\]

%dsp-level:0
\[
\qquad \text{alternation, uniqueness, the prefix law, the reversal inadmissibility}: \mathrm{FT.144}(2),
\]

%dsp-level:0
\[
(4)(\mathrm{ii})\ \text{re-derived by exhaustive enumeration at lengths}\
\]

%dsp-level:0
\[
\le 12\ [\text{machine}]
\]

%dsp-level:0
\[
\text{relaxation census}: \#\{\text{prefix-admissible completed words of length}\ 2m\}
\]

%dsp-level:0
\[
= \mathrm{Cat}_m\ [\text{machine enumeration} = \text{the closed form}\ \tfrac1{m+1}\binom{2m}m];
\]

%dsp-level:0
\[
\qquad \mathrm{Cat}_m=1\iff m\le1\ [\text{exact integers}];
\]

%dsp-level:0
\[
\qquad \text{two pairs}: \binom{2n}{2m_0}\ \text{labelled interleaving patterns summed over}\ m_0+m_1
\]

%dsp-level:0
\[
= n\ [\text{machine}]
\]

%dsp-level:0
\[
= \text{the closed form}
\]

%dsp-level:0
\[
(2):\ \mathrm{EM.9}\ \text{table verbatim}\ (\text{machine: the four sign pairs});
\]

%dsp-level:0
\[
\qquad \text{telescoping}: \sum_{i=1}^k\tfrac12(\sigma_{v_{i-1}}-\sigma_{v_i})
\]

%dsp-level:0
\[
= \tfrac12(\sigma_{v_0}-\sigma_{v_k})\ [\text{pairwise cancellation of the}\ k-1\ \text{interior terms};\ \text{machine: all}\ \sigma\in\{\pm1\}^{k+1},\ k\le7];
\]

%dsp-level:0
\[
\qquad \partial c=0\implies v_0=v_k\implies q_{\mathrm{EM}}=0\ [\text{machine}]
\]

%dsp-level:0
\[
(3):\ \partial c
\]

%dsp-level:0
\[
= t(c)-s(c)\ \text{is a signed pair of vertices};
\]

%dsp-level:0
\[
\ \partial(\text{vertex})
\]

%dsp-level:0
\[
= 0\implies\partial^2c
\]

%dsp-level:0
\[
= 0\ [\mathrm{D1.4};\ \text{machine: total multiplicity}\ 0\ \text{on every enumerated path}];
\]

%dsp-level:0
\[
\qquad \ker\partial=\{s(c)=t(c)\}\ [\text{machine: equality exactly on the closed paths}]
\]

%dsp-level:0
\[
(4):\ \text{a Noether charge is the generator of a one-parameter symmetry action}\ [\text{math}];
\]

%dsp-level:0
\[
\qquad \text{the value space is}\ \{-1,0,+1\}\ [\mathrm{EM.9}]\implies\text{the action factors through}\ S_3;
\]

%dsp-level:0
\[
\qquad U(1)\cong\mathbb Q/\mathbb Z\ \text{divisible}\ [\text{math}];
\]

%dsp-level:0
\[
\ \text{a divisible image in}\ S_3\ \text{consists of elements having}\ m\text{-th roots}\ \forall m:\ \{e\}\ \text{only}\ [\text{machine: the root census}];
\]

%dsp-level:0
\[
\qquad \text{hom-counts}: g^m=e\iff\operatorname{ord}(g)\mid m\ [\text{machine}]
\]

%dsp-level:0
\[
(5):\ \text{conjunction of}\ (1)\text{–}(4);
\]

%dsp-level:0
\[
\ \text{the limit step is BM.11's content, cited as a forward pointer and imported nowhere here}.\ \blacksquare
\]

\textbf{Phys.} Charge conservation in this corpus is a fact about integers, not about symmetries. A valid universe is one in which every crystallization has its dissolution — the definition says so, and the transition domains make the completed history unique: seal, unseal, seal, unseal. The endpoint charge read telescopes: sum the per-edge source along a materialized chain and every interior sign cancels, leaving the two ends; close the chain and the read is exactly zero. The ground is a boundary map with $\partial^2=0$ — the kernel statement of an integer chain complex. And the Noether route is not merely unnecessary here, it is closed: the charge alphabet has three elements, so any continuous charge-rotation action would have to land in a three-element permutation group, and a divisible group has no nontrivial finite image. There is no $U(1)$ to rotate the ternary read, hence no Noether current to build it from. What the continuum limit delivers later is the density form of a law that already holds, exactly, at every cell size.

\end{theorem}

\begin{theorem}[{\texorpdfstring{BM.8 — The continuum promotion: the hop becomes a translation, the $C_3$ becomes $SO(3)$, the tick stays an arrow, and the gauge algebras were never discrete}{BM.8 - The continuum promotion: the hop becomes a translation, the becomes , the tick stays an arrow, and the gauge algebras were never discrete}}]\label{tag:BM-8}
\textbf{Thm.}

%dsp-level:0
\[
(1)\quad\textbf{spatial translations from the discrete hop:}
\]

%dsp-level:0
\[
\qquad \text{on the emergent chart}\ \chi([v])
\]

%dsp-level:0
\[
= (x^0,x^1,x^2,x^3)\ [\mathrm{CT2.6};\ \mathrm{D1.8}:\ x^0:=\tau_n],
\]

%dsp-level:0
\[
\ \text{let}\ T_\ell f(x):
\]

%dsp-level:0
\[
= f(x+\ell)\ \text{be the hop of length}\ \ell\ \text{and}\ D_\ell:
\]

%dsp-level:0
\[
= \tfrac1\ell(T_\ell-1);
\]

%dsp-level:0
\[
\textbf{the exact commutator identity:}\qquad[\hat x,\,D_\ell]=-T_\ell\qquad\text{at every}\ \ell\ne0\qquad[\text{machine: exact on}\ x^k,\ k\le7];
\]

%dsp-level:0
\[
\textbf{the exact binomial expansion:}
\]

%dsp-level:0
\[
\qquad D_\ell\,x^k=\sum_{j=1}^{k}\binom{k}j x^{k-j}\ell^{\,j-1}=k\,x^{k-1}+\ell\cdot(\text{exact rational polynomial})
\]

%dsp-level:0
\[
\qquad [\text{machine}];
\]

%dsp-level:0
\[
\textbf{the Weyl transport:}
\]

%dsp-level:0
\[
\qquad T_\ell\,\hat x=(\hat x+\ell)\,T_\ell
\]

%dsp-level:0
\[
\qquad [\text{machine}];
\]

%dsp-level:0
\[
\qquad \ell\to0:\quad D_\ell\to\partial_x,\quad[\hat x,\partial_x]=-1,\quad\hat p:=-i\partial_x\implies[\hat x,\hat p]=i\ \ (\hbar=1\ [\mathrm{D2.9}]);
\]

%dsp-level:0
\[
\textbf{the dictionary carrier:}
\]

%dsp-level:0
\[
\qquad \text{the three}\ T\text{-odd coefficients}\ w_i\ [\mathrm{D2.1};\ \mathrm{CT1.2}]\ \longmapsto\ N_i=\beta w_i\ \longmapsto\ \bar g_{0i}=\varkappa N_i\ [\mathrm{CT1.12};\ \mathrm{D2.1}];
\]

%dsp-level:0
\[
N_i\ \text{is the Lagrange multiplier of the momentum constraint, whose charge generates spatial diffeomorphisms}\ [\text{math: ADM}]:
\]

%dsp-level:0
\[
\qquad \text{the conserved quantity is the momentum}\ P_i,
\]

%dsp-level:0
\[
\ \text{and}\ P_i\ \text{is}\ T\text{-odd, matching its carrier}\ [w_i\mapsto-w_i\ [\mathrm{FT.156}(3)]];
\]

%dsp-level:0
\[
(2)\quad\textbf{time translations: the tick is a semigroup, and the dynamics strictly contracts}\ —\ \textbf{an obstruction with a conditional charge:}
\]

%dsp-level:0
\[
\qquad \neg\exists\,\tau^{-1}\ [\mathrm{A6};\ \mathrm{FT.156}(5)];
\]

%dsp-level:0
\[
\qquad |\mathfrak L(\tau_{n+1})|\ge|\mathfrak L(\tau_n)|\ [\mathrm{A7};\ \text{machine}],\ \text{strict at every crystallization};
\]

%dsp-level:0
\[
V(\tilde K)\mapsto\bigl(1-\tfrac\gamma4\bigr)^2V(\tilde K),
\]

%dsp-level:0
\[
\qquad \text{deficit}=\tfrac{\gamma(8-\gamma)}{16},
\]

%dsp-level:0
\[
\qquad \gamma=1-4\varepsilon_q\in(0,1]\ [\mathrm{L2},\ \mathrm{D0.3},\ \mathrm{CT0.8};\ \text{machine}];
\]

%dsp-level:0
\[
\bigl(1-\tfrac\gamma4\bigr)^2+\tfrac{\gamma(8-\gamma)}{16}=1\ \text{exactly};
\]

%dsp-level:0
\[
\qquad \tfrac{\gamma(8-\gamma)}{16}=0\iff\gamma\in\{0,8\}\ [\text{machine}];
\]

%dsp-level:0
\[
\qquad \gamma=0\iff\varepsilon_q=\tfrac14\notin[0,\tfrac14)\ [\mathrm{CT0.8}];
\]

%dsp-level:0
\[
\therefore\ \text{on the whole admissible range the tick map strictly contracts}\ V:\quad\neg\exists\,
\]

%dsp-level:0
\[
\text{a nonzero tick-invariant of the pair-state dynamics}\implies\neg\exists\,
\]

%dsp-level:0
\[
\text{an exact time-translation invariance of the microscopic dynamics};
\]

%dsp-level:0
\[
\textbf{the conditional energy read}\ (\text{declared scope, theorem-local, L0's dynamical-scope pattern}):\ \text{under the stationarity hypothesis — the coarse-grained coefficient fields}\ (a_{00},q_{ij},w_i)\ \text{of}\ \mathrm{D2.1}\ \text{are tick-independent across the cells of}\ \mathrm{CT2.6}(\mathrm{ii})\text{–}(\mathrm{iii})\ —\ \text{the continuum limit carries a time-translation one-parameter group whose charge is the energy read};
\]

%dsp-level:0
\[
\text{on the sealed sector the hypothesis is structural rather than imposed: a bound chain's reads are tick-independent}\ [\mathrm{D1.5};\ \text{forward pointer:}\ \mathrm{BM.9}(3)];
\]

%dsp-level:0
\[
\qquad \text{at}\ \gamma\to0\ \text{the contraction}\to1\ \text{and the deficit}\to0\ [\text{exact rationals}];
\]

%dsp-level:0
\[
(3)\quad\textbf{rotations: the}\ C_3\ \textbf{is a rational point of a one-parameter subgroup already present in the register:}
\]

%dsp-level:0
\[
\qquad \mathcal W:=W_1{+}W_2{+}W_3,
\]

%dsp-level:0
\[
\qquad \mathcal W^3=3\mathcal W
\]

%dsp-level:0
\[
\qquad [\mathrm{FT.158}(1);\ \text{machine}];
\]

%dsp-level:0
\[
X:=\hat n\cdot W=\mathcal W/\sqrt3,\quad\hat n=\tfrac{(1,1,1)}{\sqrt3}:
\]

%dsp-level:0
\[
\qquad X^3=X,
\]

%dsp-level:0
\[
\qquad \operatorname{charpoly}(\mathcal W)=x^4-3x^2\implies\operatorname{spec}(X)=\{1,0,-1,0\}
\]

%dsp-level:0
\[
\qquad [\text{machine: Faddeev–LeVerrier, exact traces}];
\]

%dsp-level:0
\[
\textbf{the promotion identity:}
\]

%dsp-level:0
\[
\qquad u(\theta):=1+(\cos\theta-1)X^2-i\sin\theta\,X
\]

%dsp-level:0
\[
\qquad \text{satisfies}
\]

%dsp-level:0
\[
\qquad u(\theta)u(\phi)=u(\theta+\phi)\ \ \forall\theta,\phi
\]

%dsp-level:0
\[
\text{exactly in}\ \mathbb Q[i][c,s,c',s',X]/(c^2{+}s^2{-}1,\ c'^2{+}s'^2{-}1,\ X^3{-}X)\ [\text{machine}]:
\]

%dsp-level:0
\[
\qquad u(0)=1,\quad u(\theta)u(-\theta)=1,\quad u(\theta)^3=u(3\theta)\ [\text{triple-angle polynomials}];
\]

%dsp-level:0
\[
\text{the group law closes in}\ \operatorname{span}\{1,X,X^2\}:\ X^3=X\ \text{is the only relation used}\ [\text{machine: degree census}];
\]

%dsp-level:0
\[
\textbf{the discrete elements are its rational points:}
\]

%dsp-level:0
\[
\qquad \theta=\tfrac{2\pi}3:
\]

%dsp-level:0
\[
\quad U=E_0-\tfrac12\mathcal W^2-\tfrac i2\mathcal W=u\otimes u,
\]

%dsp-level:0
\[
\quad U^3=E_0,
\]

%dsp-level:0
\[
\quad UX_iU^\dagger=X_{i+1}\ (X\in\{W,J,K\})\ [\mathrm{FT.111},\ \mathrm{FT.158}(1);\ \text{machine}];
\]

%dsp-level:0
\[
\theta=\pi:\quad U_{ij}=E_0-(W_i+W_j)^2,\quad U_{ij}^2=E_0,\quad U_{ij}:\ X_i\leftrightarrow X_j,\ X_k\mapsto-X_k\ [\mathrm{FT.156}(4);\ \text{machine}];
\]

%dsp-level:0
\[
C_3\le SO(3)\ \text{is inner in the weak rotation algebra}\ [\mathrm{CT1.4}:\ \operatorname{span}\{W_i\}\cong\mathfrak{su}(2);\ \mathrm{CT1.8}:\ J,K\ \text{vectors}]:
\]

%dsp-level:0
\[
\qquad \text{the promotion is exact,}\ \neg\text{approximate};
\]

%dsp-level:0
\[
\textbf{the charge:}
\]

%dsp-level:0
\[
\qquad \textstyle\sum_iW_i^2=\tfrac32E_0+\tfrac12\Sigma=E_0+\mathbb S=2\,\Pi_{\mathcal R}\ [\mathrm{FT.157}(2);\ \text{machine}]:
\]

%dsp-level:0
\[
\qquad \text{Casimir}\ 2\ \text{on}\ \vee^2\ (j{=}1),\quad 0\ \text{on}\ \wedge^2\ (j{=}0)\ [\text{machine}];
\]

%dsp-level:0
\[
\text{the integer weight lattice}\ \{1,0,-1\}\ \text{is what makes the rational angles admissible:}\ u(\theta+2\pi)
\]

%dsp-level:0
\[
= u(\theta)\ \text{on integer weights}\ [\text{math}];
\]

%dsp-level:0
\[
(4)\quad\textbf{gauge: the slot algebras are natively continuous — no limit is needed for the algebra:}
\]

%dsp-level:0
\[
\qquad [W_i,W_j]=i\epsilon_{ijk}W_k\ \text{on all}\ 9\ \text{pairs}\ [\mathrm{CT1.4};\ \text{machine}];
\]

%dsp-level:0
\[
\mathfrak u(2)=\mathfrak{su}(2)\oplus\mathfrak u(1)_{\mathrm{centre}}:
\]

%dsp-level:0
\[
\quad Y=\tfrac12\sigma_0\ \text{commutes with all}\ W_i,
\]

%dsp-level:0
\[
\quad 3\ \text{traceless}\ +\ 1\ \text{central}=4\ [\mathrm{CT0.13},\ \mathrm{EM.2};\ \text{machine}];
\]

%dsp-level:0
\[
\pi:\mathfrak u(3)\to\mathfrak{su}(3),
\]

%dsp-level:0
\[
\ \ \pi(X)
\]

%dsp-level:0
\[
= X-\tfrac{\mathfrak T(X)}3I_3:\ \ \text{Lie homomorphism on all}\ 81\ \text{basis pairs},
\]

%dsp-level:0
\[
\ \ker\pi
\]

%dsp-level:0
\[
= \operatorname{span}\{I_3\}\ [\mathrm{CT3.7};\ \text{machine}];
\]

%dsp-level:0
\[
\text{the colour block: all}\ 8\ \text{generators traceless},
\]

%dsp-level:0
\[
\ \text{all}\ 64\ \text{commutators in the span},
\]

%dsp-level:0
\[
\ \text{Jacobi on all}\ 512\ \text{triples},
\]

%dsp-level:0
\[
\ [H_1,H_2]
\]

%dsp-level:0
\[
= 0,
\]

%dsp-level:0
\[
\ \det\mathrm K
\]

%dsp-level:0
\[
\ne 0\ [\mathrm{CT3.8},\ \mathrm{FT.82}(3);\ \text{machine: exact rationals}];
\]

%dsp-level:0
\[
\textbf{what the continuum limit supplies is the gauging,}\ \neg\text{the algebra: the parameters become fields on}\ M\ [\mathrm{CT2.6}]\ \text{and the covariant derivative is the chart-side object}\ [\text{math}];
\]

%dsp-level:0
\[
(5)\quad\textbf{the three charge statuses are distinct:}
\]

%dsp-level:0
\[
\qquad \textbf{colour:}\ \text{on a sealed composite every Hermitian slot operator reads only its trace}\ \mathfrak T(H)
\]

%dsp-level:0
\[
= 3a\ [\mathrm{BM.1}(2);\ \text{machine}]\ \text{and the invariant line is}\ 1\text{-dimensional}\ [\mathrm{BM.1}(3);\ \text{machine}:\ 27-\operatorname{rk}=1]:\ \text{the colour charge of a sealed singlet is}\ 0\ \text{exactly};
\]

%dsp-level:0
\[
\textbf{weak isospin:}\ \text{the derived Casimir of}\ (3)\ —\ 2\ \text{on the symmetric sector,}\ 0\ \text{on the antisymmetric one}\ [\mathrm{FT.157}(2)];
\]

%dsp-level:0
\[
\textbf{electric charge:}\ \neg\text{a Noether charge}\ [\mathrm{BM.7}(4)]:\ \text{the ternary read is topological};
\]

%dsp-level:0
\[
\qquad \text{the gauge Noether story of the corpus therefore has two continuous charges and one discrete invariant, not three of a kind}.
\]

\textbf{Rmk (promotion discipline).} \textit{Momentum}, \textit{energy} and \textit{angular momentum} name the charges of the promoted one-parameter groups; the names are physics vocabulary for derived structures [A §0.2]. Clause (2)'s stationarity hypothesis is a theorem-local dynamical scope of the species of L0's crystallization–tilt coupling and CT4.5a's ensemble scope, declared at the point of use; it is not an axiom, and it is not discharged by any tag of Parts 1–3.

\textbf{Req.} A6, A7 (Part 1); D0.3, D1.5, D1.8, D2.1 (Parts 1–2); CT0.8, CT0.13, CT1.2, CT1.4, CT1.8 (Part 1); CT1.12, CT2.6, EM.2 (Part 2; CT2.6's conditions (ii)–(iii) at the point of use — A §0.2a); CT3.7, CT3.8 (Part 2); L2 (Part 1); Def-BM.3, BM.1(2), BM.1(3) (this part); FT.82(3), FT.111, FT.144(2), FT.156(3), FT.156(4), FT.156(5), FT.157(2), FT.158(1) (Part 3); exact $\mathbb Q[x,\ell]$ hop calculus, exact identities in $\mathbb Q[i][c,s,c',s',X]/(c^2{+}s^2{-}1,X^3{-}X)$, Faddeev–LeVerrier characteristic polynomials, exact $\mathbb Q(i)$ conjugation tables, exact $\mathbb Q$ structure constants and Killing determinants, exact rational contraction arithmetic [math; machine: \texttt{checks/\allowbreak{}verify\_\allowbreak{}H\_\allowbreak{}noether.py} §B].

\textbf{Pf.}

%dsp-level:0
\[
(1):\ (T_\ell-1)f\ \text{is divisible by}\ \ell\ \text{in}\ \mathbb Q[x,\ell]\ \text{for every polynomial}\ f\ [\text{machine}];
\]

%dsp-level:0
\[
\qquad \hat x D_\ell f-D_\ell(\hat x f)
\]

%dsp-level:0
\[
= \tfrac1\ell\bigl[x(f(x{+}\ell)-f(x))-((x{+}\ell)f(x{+}\ell)-xf(x))\bigr]
\]

%dsp-level:0
\[
= -\tfrac{\ell f(x+\ell)}\ell
\]

%dsp-level:0
\[
= -T_\ell f\ [\text{exact}]
\]

%dsp-level:0
\[
D_\ell x^k
\]

%dsp-level:0
\[
= \tfrac{(x+\ell)^k-x^k}\ell
\]

%dsp-level:0
\[
= \sum_{j\ge1}\binom kj x^{k-j}\ell^{j-1}\ [\text{binomial theorem; machine coefficient census}];
\]

%dsp-level:0
\[
\qquad \ell\to0\ \text{leaves}\ kx^{k-1}\ \text{and}\ [\hat x,\partial_x]=-1;
\]

%dsp-level:0
\[
\qquad \hat p=-i\partial_x\implies[\hat x,\hat p]=i\ [\hbar=1,\ \mathrm{D2.9}]
\]

%dsp-level:0
\[
\text{dictionary rows verbatim}\ [\mathrm{D2.1},\ \mathrm{CT1.12}];
\]

%dsp-level:0
\[
\qquad N_i\ \text{as the momentum-constraint multiplier and}\ P_i\ \text{as its charge}\ [\text{math: ADM}];
\]

%dsp-level:0
\[
\qquad w_i\mapsto-w_i\ \text{under}\ T\ [\mathrm{FT.156}(3)\ \text{re-cited}]
\]

%dsp-level:0
\[
(2):\ \mathrm{A6}/\mathrm{A7}\ \text{verbatim}\ [\text{machine: monotone and strictly increasing samples}];
\]

%dsp-level:0
\[
\qquad \mathrm{L2}\ \text{contraction re-derived as an exact polynomial identity in}\ \gamma;
\]

%dsp-level:0
\[
\qquad \text{zero set}\ \{0,8\}\ [\text{machine over a rational sweep and by factorization}];
\]

%dsp-level:0
\[
\qquad \gamma\in(0,1]\ \text{on}\ \varepsilon_q\in[0,\tfrac14)\ [\mathrm{D0.3},\ \mathrm{CT0.8}]
\]

%dsp-level:0
\[
\text{strict contraction}\ \forall\gamma
\]

%dsp-level:0
\[
\in (0,1]
\]

%dsp-level:0
\[
\implies V(\mathcal T^n\rho)\to0\ \text{and no nonzero invariant}\ [\mathrm{L1}\ \text{monotone};\ \text{exact rationals}];
\]

%dsp-level:0
\[
\qquad \text{the conditional clause states its hypothesis explicitly and cites no tag for it}
\]

%dsp-level:0
\[
(3):\ \mathcal W^3=3\mathcal W\ [\text{machine}]\implies X^3=X;
\]

%dsp-level:0
\[
\qquad \text{charpoly}\ x^4-3x^2\ [\text{Faddeev–LeVerrier, exact};\ \operatorname{tr}\mathcal W=0,\ \operatorname{tr}\mathcal W^2=6,\ \operatorname{tr}\mathcal W^4=18];
\]

%dsp-level:0
\[
\qquad u(\theta)u(\phi)=u(\theta+\phi)\ [\text{machine in the quotient ring, using the addition formulas and}\ X^4=X^2]
\]

%dsp-level:0
\[
\theta
\]

%dsp-level:0
\[
= \tfrac{2\pi}3:\ (\cos\theta-1)
\]

%dsp-level:0
\[
= -\tfrac32,
\]

%dsp-level:0
\[
\ \sin\theta/\sqrt3
\]

%dsp-level:0
\[
= \tfrac12
\]

%dsp-level:0
\[
\implies u
\]

%dsp-level:0
\[
= E_0-\tfrac12\mathcal W^2-\tfrac i2\mathcal W
\]

%dsp-level:0
\[
= u\otimes u\ [\text{machine entrywise}];
\]

%dsp-level:0
\[
\qquad U^3=E_0,\ UX_iU^\dagger=X_{i+1}\ [\text{machine, all}\ 9];
\]

%dsp-level:0
\[
\qquad \theta=\pi:\ (\cos\pi-1)=-2,\ X=(W_i{+}W_j)/\sqrt2\implies U_{ij}=E_0-(W_i{+}W_j)^2\ [\text{machine}]
\]

%dsp-level:0
\[
\textstyle\sum_iW_i^2=\tfrac32E_0+\tfrac12\Sigma\ [\mathrm{FT.157}(2)\ \text{re-derived};\ \text{machine}];
\]

%dsp-level:0
\[
\qquad \text{Casimir reads on}\ \Pi_{\mathcal R},\ \Pi_{\mathcal V}\ [\text{machine}]
\]

%dsp-level:0
\[
(4):\ \text{all rows re-derived entrywise over}\ \mathbb Q(i)\ \text{and}\ \mathbb Q\ [\text{machine}];
\]

%dsp-level:0
\[
\qquad \text{the homomorphism and kernel rows of}\ \mathrm{CT3.7}\ \text{on the}\ \mathfrak u(3)\ \text{basis}\ \{T_a\}\cup\{I_3\}\ [\text{machine: all}\ 81\ \text{pairs}]
\]

%dsp-level:0
\[
(5):\ \mathrm{BM.1}(2)\text{–}(3)\ \text{re-derived on exact witnesses}\ [\text{machine}:\ \langle\varepsilon|\sum_iH^{(i)}|\varepsilon\rangle=\mathfrak T(H)\langle\varepsilon|\varepsilon\rangle;\ (\ H\otimes I-I\otimes H^{\mathrm T})\operatorname{vec}(I)=0;\ 27-\operatorname{rk}=1];
\]

%dsp-level:0
\[
\qquad \mathrm{BM.7}(4)\ \text{re-cited}.\ \blacksquare
\]

\textbf{Phys.} Three of the four continuum symmetries are already sitting in the register; the fourth is not, and the structure yields why. The hop is a translation with an exact error term: the commutator of position with the finite difference is minus the shift operator at every step size, and the binomial expansion hands over the derivative as its leading term, so the momentum generator and its canonical commutator are the $\ell\to0$ read of an identity that never needed a limit to be true. Rotation is stronger still: because the weak generator satisfies $X^3=X$, the exponential of $-i\theta X$ is a quadratic polynomial in $X$ with $\cos\theta$ and $\sin\theta$ as coefficients, and the group law holds as a ring identity — the $120^\circ$ seal rotation and the $180^\circ$ transpositions are just two rational points of a one-parameter group that was there all along, with integer weights making the periodicity exact. The gauge algebras were never discrete: $\mathfrak{su}(2)$, $\mathfrak{su}(3)$ and the rung-2 centre close as Lie algebras over the rationals, and what the continuum adds is the gauging, not the algebra. Time is the exception, and it is an honest one — the tick has no inverse and the ledger dynamics strictly contracts its Lyapunov function at every admissible threshold, so the microscopic theory has no exact time-translation symmetry and the energy charge exists only under a declared stationarity scope, which the sealed sector satisfies structurally. The charge ledger then splits three ways: colour reads zero on every sealed singlet, weak isospin reads the derived Casimir, and electric charge is not a Noether charge at all.

\end{theorem}

\begin{theorem}[{\texorpdfstring{BM.9 — The stress-energy tensor of a sealed composite: the two reads, their exact sign reconciliation, and the dust source}{BM.9 - The stress-energy tensor of a sealed composite: the two reads, their exact sign reconciliation, and the dust source}}]\label{tag:BM-9}
\textbf{Thm.}

%dsp-level:0
\[
(1)\quad\textbf{the corpus definition and the field equations:}
\]

%dsp-level:0
\[
\qquad T_{\mu\nu}:=-\frac{2}{\sqrt{-g}}\frac{\delta S_{\mathrm{force}}}{\delta g^{\mu\nu}}
\]

%dsp-level:0
\[
\qquad [\mathrm{CT1.22}];
\]

%dsp-level:0
\[
S_{\mathrm{eff}}=S_{\mathrm{geom}}+S_{\mathrm{force}},\qquad S_{\mathrm{geom}}[g]=\frac1{2\kappa^2}\int d^4x\,\sqrt{-g}\,(R-2\Lambda)\qquad[\mathrm{CT1.21}];
\]

%dsp-level:0
\[
\delta S_{\mathrm{eff}}=0\iff G_{\mu\nu}+\Lambda g_{\mu\nu}=\kappa^2T_{\mu\nu},\qquad\kappa^2=8\pi G_N=\frac{2\pi\,\mathfrak h^2}{\ln2}\qquad[\mathrm{CT1.23},\ \mathrm{CT2.9},\ \mathrm{P3}];
\]

%dsp-level:0
\[
(2)\quad\textbf{the canonical read and the exact sign reconciliation:}
\]

%dsp-level:0
\[
\qquad \text{for}\ \mathcal L=-\tfrac12g^{\alpha\beta}\partial_\alpha\phi\,\partial_\beta\phi-\tfrac12m^2\phi^2:
\]

%dsp-level:0
\[
\qquad \frac{\partial\mathcal L}{\partial(\partial^\mu\phi)}=-\partial_\mu\phi\ \ [\text{exact}];
\]

%dsp-level:0
\[
T^{\mathrm{can}}_{\mu\nu}:=\frac{\partial\mathcal L}{\partial(\partial^\mu\phi)}\,\partial_\nu\phi-g_{\mu\nu}\mathcal L=-\partial_\mu\phi\,\partial_\nu\phi-g_{\mu\nu}\mathcal L;
\]

%dsp-level:0
\[
\qquad T^{\mathrm{Hilb}}_{\mu\nu}=\partial_\mu\phi\,\partial_\nu\phi+g_{\mu\nu}\mathcal L;
\]

%dsp-level:0
\[
\boxed{\ T^{\mathrm{can}}_{\mu\nu}=-\,T^{\mathrm{Hilb}}_{\mu\nu}\ }
\]

%dsp-level:0
\[
\qquad \text{exactly, at signature}\ (-,+,+,+)\ [\mathrm{A3},\ \mathrm{C13},\ \mathrm{CT1.12}]\ \text{with this}\ \mathcal L\ \text{sign}
\]

%dsp-level:0
\[
\qquad [\text{machine: all}\ 16\ \text{components on}\ 3\ \text{rational metric witnesses}];
\]

%dsp-level:0
\[
\textbf{the convention is fixed by the energy sign:}\quad\text{at rest}\ T^{\mathrm{Hilb}}_{00}
\]

%dsp-level:0
\[
= \tfrac12\sum_i(\partial_i\phi)^2+\tfrac12m^2\phi^2\ \text{with every coefficient}\ +\tfrac12,
\]

%dsp-level:0
\[
\ \text{while}\ T^{\mathrm{can}}_{00}\ \text{carries}\ -\tfrac12\ [\text{machine: exact coefficient census}];
\]

%dsp-level:0
\[
\therefore\ \mathrm{CT1.22}\text{'s minus is the positive-energy read at}\ (-,+,+,+);
\]

%dsp-level:0
\[
\qquad \text{the variational composition is exact:}\quad\frac{\delta(\sqrt{-g}\,\mathcal L)}{\delta g^{\mu\nu}}
\]

%dsp-level:0
\[
= -\tfrac12\sqrt{-g}\,
\]

%dsp-level:0
\[
\bigl(\partial_\mu\phi\,\partial_\nu\phi+g_{\mu\nu}\mathcal L\bigr)\ [\text{machine}];
\]

%dsp-level:0
\[
(3)\quad\textbf{the Noether identity (translation current):}
\]

%dsp-level:0
\[
\qquad \partial^\mu T^{\mathrm{Hilb}}_{\mu\nu}=(\Box\phi-m^2\phi)\,\partial_\nu\phi
\]

%dsp-level:0
\[
\qquad \text{as an exact jet identity}\ [\text{machine: all}\ 4\ \text{components},\ 3\ \text{metric witnesses}];
\]

%dsp-level:0
\[
\Box\phi-m^2\phi=0\ \text{(the field equation of this}\ \mathcal L)\implies\partial^\mu T^{\mathrm{Hilb}}_{\mu\nu}=0:
\]

%dsp-level:0
\[
\qquad T_{\mu\nu}\ \text{is the Noether current of the translation group promoted in}\ \mathrm{BM.8}(1);
\]

%dsp-level:0
\[
\textbf{the trace identity:}\qquad g^{\mu\nu}T^{\mathrm{Hilb}}_{\mu\nu}=-(\partial\phi)^2-2m^2\phi^2\qquad[\text{machine: exact in}\ 4\ \text{dimensions}];
\]

%dsp-level:0
\[
(4)\quad\textbf{the sealed composite at rest — the neutron}\ [\mathrm{Def\text{-}BM.3}:\ \text{the}\ \varepsilon\text{-sealed 3-juncture in the neutral charge class}]:
\]

%dsp-level:0
\[
\qquad T^{\mathrm{bd}}_{\mu\nu}=(\rho+p)u_\mu u_\nu+p\,g_{\mu\nu}
\]

%dsp-level:0
\[
\qquad [\mathrm{CT1.22}\ \text{bound part}];
\]

%dsp-level:0
\[
u=(1,0,0,0),\quad g_{\mu\nu}=\mathrm{diag}(-1,1,1,1),\quad c=1\ [\mathrm{D2.9}]:
\]

%dsp-level:0
\[
\qquad T_{\mu\nu}=\mathrm{diag}(\rho,\,p,\,p,\,p)
\]

%dsp-level:0
\[
\qquad [\text{machine: exact component identity}];
\]

%dsp-level:0
\[
\textbf{fully sealed}\implies p=0,\quad w=\tfrac p\rho=0
\]

%dsp-level:0
\[
\qquad [\mathrm{CT4.5}:\ T^{(D^-)}_{\mu\nu}=\rho\,u_\mu u_\nu,\ p_{D^-}=0]:
\]

%dsp-level:0
\[
\qquad T_{\mu\nu}=\mathrm{diag}(\rho,0,0,0)\ \ [\text{machine}];
\]

%dsp-level:0
\[
\textbf{static:}\ \ \partial_0T_{\mu\nu}=0;
\]

%dsp-level:0
\[
\qquad \textbf{spherically symmetric:}\ \ \text{the isotropic spatial block}\ p\,
\]

%dsp-level:0
\[
\delta_{ij}\ \text{is invariant under all}\ 48\ \text{signed triad permutations}\ \subset O(3),
\]

%dsp-level:0
\[
\ \text{hence under the promoted}\ SO(3)\ [\mathrm{BM.8}(3);\ \text{machine}];
\]

%dsp-level:0
\[
\textbf{trace and energy conditions:}
\]

%dsp-level:0
\[
\qquad g^{\mu\nu}T_{\mu\nu}=-\rho+3p\overset{p=0}{=}-\rho;
\]

%dsp-level:0
\[
\qquad \text{WEC},\ \text{DEC},\ \text{SEC hold exactly at}\ \rho\ge0,\ p=0\ [\text{machine: rational witnesses}];
\]

%dsp-level:0
\[
\rho\ge0,\qquad\rho=0\iff\neg\exists v:\ n(v)\ge4\qquad[\mathrm{CT4.5};\ \text{machine on enumerated rung profiles}];
\]

%dsp-level:0
\[
(5)\quad\textbf{the source read and its dimensions:}
\]

%dsp-level:0
\[
\qquad \rho=\frac{m_{\mathrm{hid}}c^2}{V}=\frac{m_{\mathrm{hid}}}{V}\ \ [c=1,\ \mathrm{D2.9}],
\]

%dsp-level:0
\[
\qquad m_{\mathrm{hid}}=\frac{E_{\mathrm{int}}\ln2}{\mathfrak h}
\]

%dsp-level:0
\[
\qquad [\mathrm{FT.44}(4),\ \mathrm{FT.42}(1)];
\]

%dsp-level:0
\[
[\mathfrak h]=\mathfrak q^{-1/2},\quad[\text{length}]=\mathfrak q^{-1/2},\quad[\text{energy}]=\mathfrak q^{+1/2},\quad[T_{\mu\nu}]=\mathfrak q^{+2}\qquad[\mathrm{CT2.1},\ \mathrm{D2.9}];
\]

%dsp-level:0
\[
[m_{\mathrm{hid}}]=[E_{\mathrm{int}}]-[\mathfrak h]=0+\tfrac12=\mathfrak q^{+1/2}=[\text{energy}];
\]

%dsp-level:0
\[
\qquad [V]=\mathfrak q^{-3/2};
\]

%dsp-level:0
\[
\qquad [\rho]=\mathfrak q^{1/2-(-3/2)}=\mathfrak q^{+2}=[T_{\mu\nu}]
\]

%dsp-level:0
\[
\qquad [\text{machine: exact exponent arithmetic}];
\]

%dsp-level:0
\[
\textbf{the}\ \ln2\ \textbf{cancellation:}
\]

%dsp-level:0
\[
\qquad \kappa^2\rho=\frac{2\pi\,\mathfrak h^2}{\ln2}\cdot\frac{E_{\mathrm{int}}\ln2}{\mathfrak h\,V}=\frac{2\pi\,\mathfrak h\,E_{\mathrm{int}}}{V}
\]

%dsp-level:0
\[
\qquad [\text{machine: exact, denominators cleared}]:
\]

%dsp-level:0
\[
\text{the}\ \ln2\ \text{of the entropy calibration}\ [\mathrm{CT2.9}]\ \text{and the}\ \ln2\ \text{of the informational mass}\ [\mathrm{FT.44}(4)]\ \text{cancel in the sourced field equation}:G_{\mu\nu}+\Lambda g_{\mu\nu}
\]

%dsp-level:0
\[
= \tfrac{2\pi\mathfrak h E_{\mathrm{int}}}{V}\,
\]

%dsp-level:0
\[
u_\mu u_\nu\ \text{at}\ p\qquad = 0;
\]

%dsp-level:0
\[
(6)\quad\textbf{the width (forward pointer:}\ \mathrm{BM.10}(6)\textbf{):}
\]

%dsp-level:0
\[
\qquad V=(\Delta x)^3,
\]

%dsp-level:0
\[
\qquad \Delta x\ge\frac1{2\,\Delta p}\ \ (\hbar=1):
\]

%dsp-level:0
\[
\qquad \rho=\frac{m_{\mathrm{hid}}}{V}\le8\,m_{\mathrm{hid}}(\Delta p)^3
\]

%dsp-level:0
\[
\qquad [\text{machine: exact coefficients}];
\]

%dsp-level:0
\[
\rho\ \text{is finite for every}\ \Delta p>0:
\]

%dsp-level:0
\[
\qquad T_{\mu\nu}\ \text{is a distribution of nonzero width},\ \neg\text{a}\ \delta\text{-source};
\]

%dsp-level:0
\[
\qquad \text{the width is the scale at which the coarse-grained fields are}\ \mathrm{CT2.6}(\mathrm{ii})\text{-regular}.
\]

\textbf{Rmk (convention record).} Clause (2) prices both sign conventions instead of choosing silently: the canonical read of the classical formula and CT1.22's Hilbert read differ by an exact minus at signature $(-,+,+,+)$ with this $\mathcal L$ sign, and the corpus convention is the one with positive rest energy. Every downstream row — the dust form, the trace, the field equations, the energy conditions — is stated in the CT1.22 convention. No measured value enters any clause [P5; Def-BM.3(6)].

\textbf{Req.} A3 (Part 1); D1.5, D2.9 (Parts 1–2); CT1.21, CT1.22, CT1.23, CT2.1, CT2.9, CT4.5 (Part 2); Def-BM.1, Def-BM.3, BM.8(1), BM.8(3) (this part); FT.42(1) (Part 3); Census FT.44(4); P3, P5 (Parts 1–2); exact multivariate jet-ring polynomial identities, exact rational metric inversion by Cramer, exact coefficient censuses, Fractions [math; machine: \texttt{checks/\allowbreak{}verify\_\allowbreak{}H\_\allowbreak{}noether.py} §C].

\textbf{Pf.}

%dsp-level:0
\[
(1):\ \mathrm{CT1.21}/\mathrm{CT1.22}/\mathrm{CT1.23}\ \text{verbatim};
\]

%dsp-level:0
\[
\qquad \kappa^2=8\pi G_N=8\pi\cdot\tfrac{\mathfrak h^2}{4\ln2}=\tfrac{2\pi\mathfrak h^2}{\ln2}\ [\mathrm{CT2.9};\ \text{machine, denominators cleared}]
\]

%dsp-level:0
\[
(2):\ \partial\mathcal L/\partial(\partial^\mu\phi)
\]

%dsp-level:0
\[
= -g^{\mu\alpha}\partial_\alpha\phi\cdot\tfrac12\cdot2
\]

%dsp-level:0
\[
= -\partial^\mu\phi\ \text{and}\ \partial^\mu\phi\ \text{lowers to}\ \partial_\mu\phi\ \text{in the first slot}\ [\text{exact}];
\]

%dsp-level:0
\[
\qquad T^{\mathrm{can}}+T^{\mathrm{Hilb}}
\]

%dsp-level:0
\[
= 0\ \text{componentwise}\ [\text{machine:}\ 3\ \text{rational metrics, diagonal and fully off-diagonal}];
\]

%dsp-level:0
\[
\qquad \delta\sqrt{-g}
\]

%dsp-level:0
\[
= -\tfrac12\sqrt{-g}\,
\]

%dsp-level:0
\[
g_{\mu\nu}\delta g^{\mu\nu},
\]

%dsp-level:0
\[
\ \delta\mathcal L/\delta g^{\mu\nu}
\]

%dsp-level:0
\[
= -\tfrac12\partial_\mu\phi\partial_\nu\phi\ [\text{math}]\implies\text{the variational composition}\ [\text{machine}]
\]

%dsp-level:0
\[
(3):\ \partial_\mu(d_\mu\phi\,d_\nu\phi)=\Box\phi\,\partial_\nu\phi+\partial^\mu\phi\,\partial_\mu\partial_\nu\phi;
\]

%dsp-level:0
\[
\qquad g_{\mu\nu}\partial^\mu\mathcal L=-\partial_\nu\phi\,(\partial^\alpha\phi)-m^2\phi\,\partial_\nu\phi\ \text{at constant metric};
\]

%dsp-level:0
\[
\qquad \text{the two}\ \partial\partial\phi\ \text{terms cancel by symmetry of}\ \partial_\mu\partial_\nu\ \text{and relabeling}\ [\text{machine: the residual is exactly}\ (\Box\phi-m^2\phi)\partial_\nu\phi]
\]

%dsp-level:0
\[
(4):\ (\rho+p)u_\mu u_\nu+pg_{\mu\nu}\ \text{at}\ u
\]

%dsp-level:0
\[
= (1,0,0,0),
\]

%dsp-level:0
\[
\ g
\]

%dsp-level:0
\[
= \mathrm{diag}(-1,1,1,1):\ (0,0)\text{-entry}\ \rho,
\]

%dsp-level:0
\[
\ (0,i)\text{-entries}\ 0,
\]

%dsp-level:0
\[
\ (i,j)\text{-entry}\ p\delta_{ij}\ [\text{machine}];
\]

%dsp-level:0
\[
\qquad p=0,\ w=0\ [\mathrm{CT4.5}\ \text{re-cited}];
\]

%dsp-level:0
\[
\qquad R\,\mathrm{diag}(\rho,p,p,p)R^{\mathrm T}=\mathrm{diag}(\rho,p,p,p)\ \forall R\in\{\text{signed permutations}\}\ [\text{machine: all}\ 48];
\]

%dsp-level:0
\[
\qquad \text{trace and the three energy conditions at}\ p=0\ [\text{machine}]
\]

%dsp-level:0
\[
(5):\ c=1\ [\mathrm{D2.9}];
\]

%dsp-level:0
\[
\qquad m_{\mathrm{hid}}=E_{\mathrm{int}}\ln2/\mathfrak h\ [\mathrm{FT.44}(4)\ \text{re-cited}];
\]

%dsp-level:0
\[
\qquad \text{exponents}: \tfrac12-3(-\tfrac12)=2\ [\text{machine}];
\]

%dsp-level:0
\[
\qquad \kappa^2\rho:\ (2\pi\mathfrak h^2/\ln2)(E_{\mathrm{int}}\ln2/(\mathfrak h V))
\]

%dsp-level:0
\[
= 2\pi\mathfrak h E_{\mathrm{int}}/V\ [\text{machine: cross-multiplied identity}]
\]

%dsp-level:0
\[
(6):\ \Delta x\,\Delta p=\tfrac12\ \text{on rational witnesses}\ [\mathrm{BM.10}(6),\ \text{forward pointer}];
\]

%dsp-level:0
\[
\qquad (1/(2\Delta p))^3=1/(8(\Delta p)^3)\ [\text{machine}];
\]

%dsp-level:0
\[
\qquad \rho V=m_{\mathrm{hid}}\ \text{with}\ V>0\implies\rho<\infty\ [\text{exact}].\ \blacksquare
\]

\textbf{Phys.} A sealed composite gravitates as dust, and the corpus can now say so with both sign conventions on the table. The classical canonical formula and CT1.22's metric definition differ by an exact minus at the corpus's own signature; the minus is what makes the rest energy positive, so the convention is decided by the energy condition rather than by taste. The same tensor is the translation current: its divergence is the field equation times the field gradient, an identity in the jet variables that vanishes on-shell — Noether's theorem, run in reverse from the symmetry the hop promoted. For the neutron, the fully sealed three-juncture in the neutral class, the bound-part form collapses to $\mathrm{diag}(\rho,0,0,0)$: static, spherically symmetric, pressure-free, exactly the $D^-$ equation of state, and satisfying every energy condition at $\rho\ge0$. Its source is the informational mass over the width volume, and the dimensions balance to $\mathfrak q^{+2}$ without adjustment. The two $\ln2$ factors — one from the horizon-entropy calibration, one from the bit-count mass read — cancel in the sourced field equation, leaving $2\pi\mathfrak h E_{\mathrm{int}}/V$. And because the width is strictly positive, the source is a distribution, not a point: the smoothing scale is where the coarse-graining hypothesis of the continuum limit takes over.

\end{theorem}

\medskip\noindent\textit{\texorpdfstring{BM.10 — The canonical pair: the finite-dimension obstruction, the discrete Weyl pre-limit, and the three exact bounds}{BM.10 - The canonical pair: the finite-dimension obstruction, the discrete Weyl pre-limit, and the three exact bounds}.}\ Structural-constraints record: see Appendix B, Section B.3.5 of the Addendum.\par\medskip

\begin{theorem}[{\texorpdfstring{BM.11 — The continuum limit of the Janus topology: the differential conservation law and its two agreeing routes}{BM.11 - The continuum limit of the Janus topology: the differential conservation law and its two agreeing routes}}]\label{tag:BM-11}
\textbf{Thm.}

%dsp-level:0
\[
(1)\quad\textbf{the discrete divergence theorem carries no correction term:}
\]

%dsp-level:0
\[
\qquad \sum_{j=0}^{n-1}\ell\,(D_\ell f)(x_0+j\ell)=f(x_0+n\ell)-f(x_0)
\]

%dsp-level:0
\[
\qquad \text{EXACTLY on every rational grid}\ [\text{machine:}\ f=x^k,\ k\le5;\ n\le6;\ (x_0,\ell)\in\{(0,1),(\tfrac13,\tfrac25),(-\tfrac74,\tfrac12)\}];
\]

%dsp-level:0
\[
\text{the same total is the exact continuum integral of the derivative:}
\]

%dsp-level:0
\[
\qquad \sum_j\ell\,(D_\ell f)(x_j)=\int_{x_0}^{x_0+n\ell}\!f'(x)\,dx
\]

%dsp-level:0
\[
\qquad [\text{machine: exact equality at every}\ \ell,\ \neg\text{a limit statement}];
\]

%dsp-level:0
\[
\therefore\ \ell\to0\ \text{changes no term of the conservation identity}\ [\mathrm{BM.8}(1);\ \mathrm{CT2.6}];
\]

%dsp-level:0
\[
(2)\quad\textbf{the region law:}
\]

%dsp-level:0
\[
\qquad \text{for a materialized chain}\ c\ \text{and every sub-interval}\ [i,j]:
\]

%dsp-level:0
\[
\qquad \sum_{e\in[i,j]}j(e)=\frac{\sigma_{v_i}-\sigma_{v_j}}2=\text{the boundary read}
\]

%dsp-level:0
\[
\qquad [\text{machine: all}\ (\text{chain},\text{interval})\ \text{pairs},\ k\le5];
\]

%dsp-level:0
\[
\text{the full-chain case is}\ q_{\mathrm{EM}}(c)\ [\mathrm{BM.7}(2)];
\]

%dsp-level:0
\[
\qquad \text{the closed case reads}\ 0\ [\mathrm{D1.4};\ \mathrm{EM.11}:\ U\ \text{valid}\implies Q_{\mathrm{EM}}=\sum_{c\in U}q_{\mathrm{EM}}(c)=0];
\]

%dsp-level:0
\[
\textbf{the two-dimensional form:}
\]

%dsp-level:0
\[
\qquad \sum_{\text{cells}\in R}\bigl(\nabla_{\!\ell}\!\cdot j\bigr)=\text{the boundary flux of}\ R
\]

%dsp-level:0
\[
\qquad [\text{machine: every rectangle of a}\ 6\times5\ \text{rational lattice}];
\]

%dsp-level:0
\[
\textbf{the closed-network form:}
\]

%dsp-level:0
\[
\qquad \text{on a periodic lattice the total divergence is exactly}\ 0\ [\text{machine}]:
\]

%dsp-level:0
\[
\qquad \text{the discrete image of}\ \#J^+=\#J^-\ \text{with no boundary term}\ [\mathrm{D1.3},\ \mathrm{D2.11}];
\]

%dsp-level:0
\[
\text{the seal excess summed over every completed admissible history of length}\
\]

%dsp-level:0
\[
\le 12\ \text{is exactly}\ 0\ [\mathrm{BM.7}(1);\ \text{machine}];
\]

%dsp-level:0
\[
(3)\quad\textbf{the limit and the pointwise law:}
\]

%dsp-level:0
\[
\qquad \text{under}\ \mathrm{CT2.6}(\mathrm{ii})\text{–}(\mathrm{iii})\ [\text{extra-axiomatic, declared at the point of use — A §0.2a}]\ \text{the cells carry unboundedly many real edges and the coefficient fields are of bounded variation}:
\]

%dsp-level:0
\[
\sum_{\text{cells}\in U}(\cdots)\longrightarrow\int_U\!\sqrt{-g}\,\partial_\mu J^\mu\,d^4x,\qquad\text{boundary flux}\longrightarrow\oint_{\partial U}\!\sqrt{-g}\,J^\mu n_\mu\,d^3\Sigma;
\]

%dsp-level:0
\[
\text{the interval-sum moment matrix is nonsingular exactly for polynomial densities of degree}\
\]

%dsp-level:0
\[
\le 4\ [\text{machine: exact determinants}]\implies\text{the interval sums determine the density coefficients uniquely};
\]

%dsp-level:0
\[
\text{since}\ (2)\ \text{holds for every region, the fundamental lemma of the calculus of variations}\ [\text{math}]\ \text{gives}
\]

%dsp-level:0
\[
\qquad \boxed{\ \partial_\mu J^\mu=0\ }
\]

%dsp-level:0
\[
\qquad \text{pointwise};
\]

%dsp-level:0
\[
\textbf{the contrapositive is witnessed:}\quad\text{every nonzero polynomial density of degree}\
\]

%dsp-level:0
\[
\le 3\ \text{has a nonzero interval sum on a rational grid}\ [\text{machine}];
\]

%dsp-level:0
\[
\textbf{the current dictionary:}
\]

%dsp-level:0
\[
\qquad J^0=\text{the seal density}\ (\#J^+-\#J^-\ \text{per cell}=n_{\mathcal R}\ \text{density}\ [\mathrm{CT0.3}]);
\]

%dsp-level:0
\[
\qquad J^i=\text{the hop current on the}\ w_i\ \text{carrier}\ [\mathrm{D2.1};\ \mathrm{CT1.12}];
\]

%dsp-level:0
\[
(4)\quad\textbf{the stress-energy conservation, by two independent routes that agree:}
\]

%dsp-level:0
\[
\qquad \textbf{(i) geometry:}\ \nabla^\mu G_{\mu\nu}
\]

%dsp-level:0
\[
= 0\ \text{on the exact pp-wave witness}\ [\text{machine: nonlinear, all}\ 4\ \text{components, as polynomials in}\ u,v,x,y];
\]

%dsp-level:0
\[
g_{\mu\nu}:\ g_{uu}=H,\ g_{uv}=g_{vu}=-1,\ g_{xx}=g_{yy}=1\implies\det g=-1\ \text{exactly, independent of}\ H\ [\text{machine}];
\]

%dsp-level:0
\[
\qquad \Gamma^\mu_{\mu\lambda}=\partial_\lambda\ln\sqrt{|g|}=0\ [\text{machine}];
\]

%dsp-level:0
\[
\nabla_\lambda g_{\mu\nu}=0\ \text{on all}\ 64\ \text{components}\ [\text{machine}]\implies\nabla_\mu(\Lambda g^{\mu\nu})=0\ \text{on all}\ 4\ [\text{machine}]:
\]

%dsp-level:0
\[
\qquad \mathrm{CT1.23}\ \text{forces}\ \kappa^2\nabla_\mu T^{\mu\nu}=0;
\]

%dsp-level:0
\[
R_{\mu\nu}=-\tfrac12(\partial_x^2+\partial_y^2)H\,\delta^u_\mu\delta^u_\nu,\quad R=0,\quad G_{uu}=R_{uu}\ [\text{machine}];
\]

%dsp-level:0
\[
\qquad \Delta_\perp H=0\implies R_{\mu\nu}=0:\ \text{the vacuum pp-wave}\ [\text{machine}];
\]

%dsp-level:0
\[
\textbf{(ii) Noether:}\ \partial^\mu T_{\mu\nu}
\]

%dsp-level:0
\[
= (\Box\phi-m^2\phi)\partial_\nu\phi\ [\mathrm{BM.9}(3)]\ \text{on an exact on-shell witness}\ (\phi=t,\ m=0:\ \Box\phi=0):
\]

%dsp-level:0
\[
\qquad \partial^\mu T_{\mu\nu}=0\ [\text{machine}];
\]

%dsp-level:0
\[
\textbf{the dust read:}
\]

%dsp-level:0
\[
\qquad \partial_\mu T^{\mu\nu}
\]

%dsp-level:0
\[
= (\partial_0\rho,\,0,\,0,\,0)\ \text{for the static pressureless dust of}\ \mathrm{BM.9}(4):\ \text{the three spatial derivatives drop out at}\ p
\]

%dsp-level:0
\[
= 0\ [\text{machine}];
\]

%dsp-level:0
\[
\text{a static inhomogeneous dust is exactly conserved};
\]

%dsp-level:0
\[
\qquad \text{a time-dependent}\ \rho\ \text{leaves the exact residual}\ \partial_0\rho\ [\text{machine}];
\]

%dsp-level:0
\[
(5)\quad\textbf{the priority clause:}
\]

%dsp-level:0
\[
\qquad (1)\text{–}(3)\ \text{derive the differential law from the integer law}: \partial_\mu J^\mu
\]

%dsp-level:0
\[
= 0\ \text{is the density read of}\ \#J^+
\]

%dsp-level:0
\[
= \#J^-\ [\mathrm{BM.7}(5)];
\]

%dsp-level:0
\[
\text{the converse fails at finite cell size: the integer balance constrains the completed-history census}\ [\mathrm{BM.7}(1):\ \mathrm{Cat}_m,\ \text{the interleaving counts}]\ \text{and the density law does not};
\]

%dsp-level:0
\[
\text{no variational principle is used anywhere in}\ (1)\text{–}(3):\ \text{the input is}\ \mathrm{D1.2}\text{'s transition domains,}\ \mathrm{D1.4}\text{'s boundary map and}\ \mathrm{D2.11}\text{'s validity condition}.
\]

\textbf{Rmk (limit discipline).} Clause (3) is the tier's only use of the extra-axiomatic continuum-limit conditions, inherited at the point of use in the manner of A §0.2a; the discrete rows (1), (2) and (5) are unconditional and hold at every $\mathfrak h\ne0$. The pp-wave witness of clause (4) is a rational-coometric test object, not a claim about the corpus's own emergent metric: it verifies the identity $\nabla^\mu G_{\mu\nu}=0$, which is metric-independent [math].

\textbf{Req.} D1.2, D1.3, D1.4, D2.1, D2.11 (Parts 1–2); CT0.3 (Part 1); CT1.12, CT1.22, CT1.23, CT2.6 (Part 2; CT2.6's conditions (ii)–(iii) at the point of use — A §0.2a); CT4.5, EM.11 (Part 2); BM.7(1), BM.7(2), BM.7(5), BM.8(1), BM.9(3), BM.9(4) (this part); exact rational grid telescoping, exact moment-matrix determinants, exact multivariate pp-wave curvature and covariant-derivative algebra, the fundamental lemma of the calculus of variations [math; machine: \texttt{checks/\allowbreak{}verify\_\allowbreak{}H\_\allowbreak{}noether.py} §E].

\textbf{Pf.}

%dsp-level:0
\[
(1):\ \ell(D_\ell f)(x_j)
\]

%dsp-level:0
\[
= f(x_j+\ell)-f(x_j)\ [\mathrm{BM.8}(1)]\implies\text{the sum telescopes to}\ f(x_0+n\ell)-f(x_0)\ [\text{machine, exact rationals}];
\]

%dsp-level:0
\[
\qquad \int_a^b f'(x)\,dx=f(b)-f(a)\ \text{for}\ f=x^k\ [\text{the exact polynomial antiderivative}\ \tfrac{x^{k+1}}{k+1}];
\]

%dsp-level:0
\[
\qquad \text{the two totals are the same rational number at every}\ \ell\ [\text{machine}]
\]

%dsp-level:0
\[
(2):\ \text{pairwise cancellation of the interior signs}\ [\mathrm{BM.7}(2)\ \text{re-run per interval};\ \text{machine: exhaustive}];
\]

%dsp-level:0
\[
\qquad \text{the}\ 2\mathrm D\ \text{lattice}: \sum_{i,k}[(j^x_{i+1,k}-j^x_{i,k})+(j^y_{i,k+1}-j^y_{i,k})]\ \text{telescopes in each direction separately to the four boundary sums}\ [\text{machine: all rectangles}];
\]

%dsp-level:0
\[
\qquad \text{periodic identification cancels the boundary sums pairwise}\ [\text{machine}]
\]

%dsp-level:0
\[
(3):\ \text{cell}\to\text{chart convergence under}\ \mathrm{CT2.6}(\mathrm{ii})\text{–}(\mathrm{iii})\ [\text{math}];
\]

%dsp-level:0
\[
\qquad \text{moment matrix}: M_{b,m}
\]

%dsp-level:0
\[
= \sum_{j=0}^{b-1}\ell\,
\]

%dsp-level:0
\[
(j\ell)^m,
\]

%dsp-level:0
\[
\ b
\]

%dsp-level:0
\[
= 1..d{+}1:\ \det M
\]

%dsp-level:0
\[
\ne 0\ [\text{machine, exact Fractions}]\implies\text{injectivity of the coefficient map};
\]

%dsp-level:0
\[
\qquad \text{vanishing of all interval sums}\implies\text{all coefficients zero}\implies\partial_\mu J^\mu=0\ [\text{fundamental lemma};\ \text{math}]
\]

%dsp-level:0
\[
(4)\text{(i)}:\ \Gamma^\rho_{\mu\nu}
\]

%dsp-level:0
\[
= \tfrac12g^{\rho\sigma}(\partial_\mu g_{\nu\sigma}+\partial_\nu g_{\mu\sigma}-\partial_\sigma g_{\mu\nu}),
\]

%dsp-level:0
\[
\ R^\rho_{\sigma\mu\nu}
\]

%dsp-level:0
\[
= \partial_\mu\Gamma^\rho_{\nu\sigma}-\partial_\nu\Gamma^\rho_{\mu\sigma}+\Gamma^\rho_{\mu\lambda}\Gamma^\lambda_{\nu\sigma}-\Gamma^\rho_{\nu\lambda}\Gamma^\lambda_{\mu\sigma},
\]

%dsp-level:0
\[
\ R_{\sigma\nu}
\]

%dsp-level:0
\[
= R^\mu_{\sigma\mu\nu}\ [\text{math}];
\]

%dsp-level:0
\[
\qquad \text{all steps exact in}\ \mathbb Q[u,v,x,y]\ [\text{machine}];
\]

%dsp-level:0
\[
\qquad \nabla_\mu G^{\mu\nu}=\partial_\mu G^{\mu\nu}+\Gamma^\mu_{\mu\lambda}G^{\lambda\nu}+\Gamma^\nu_{\mu\lambda}G^{\mu\lambda}=0\ [\text{machine}]
\]

%dsp-level:0
\[
(4)\text{(ii)}:\ \mathrm{BM.9}(3)\text{'s jet identity evaluated at}\ \phi
\]

%dsp-level:0
\[
= t,
\]

%dsp-level:0
\[
\ m
\]

%dsp-level:0
\[
= 0\ [\Box\phi=g^{\mu\nu}\partial_\mu\partial_\nu t=0]:\ \text{eom}
\]

%dsp-level:0
\[
= 0\ \text{and}\ \partial^\mu T_{\mu\nu}
\]

%dsp-level:0
\[
= 0\ [\text{machine}];
\]

%dsp-level:0
\[
\qquad \text{dust}: T^{00}=\rho(x),\ T^{0i}=T^{ij}=0\implies\partial_\mu T^{\mu0}=\partial_0\rho,\ \partial_\mu T^{\mu i}=0\ [\text{machine}]
\]

%dsp-level:0
\[
(5):\ \text{conjunction of}\ (1)\text{–}(3)\ \text{with}\ \mathrm{BM.7}(1)\text{'s census rows}.\ \blacksquare
\]

\textbf{Phys.} The continuum conservation law is a limit of a telescoping sum, and the telescoping is exact — not asymptotic. Sum the per-edge source over any stretch of a materialized chain and the interior signs cancel pairwise, leaving only the two ends; do it on a lattice rectangle and the same cancellation leaves only the boundary flux; close the network and the total is zero, which is the integer balance wearing a density's clothes. Because the interval sums determine a polynomial density uniquely — the moment matrix has an exact nonzero determinant — vanishing on every region forces the density itself to vanish, and $\partial_\mu J^\mu=0$ follows pointwise. The stress-energy tensor is conserved twice over and the two derivations agree: geometrically, because the Einstein tensor is divergence-free (verified nonlinearly on an exact rational pp-wave, where $\det g=-1$ identically and metric compatibility kills the cosmological term), and dynamically, because the jet identity of the translation current vanishes on-shell. Static dust stays conserved however lumpy it is; only $\partial_0\rho$ can spoil it. And the arrow of explanation runs one way: the integers ground the density, never the reverse.

\end{theorem}

\medskip\noindent\textit{\texorpdfstring{BM.12 — The Noether status of the corpus: adjudication and scope}{BM.12 - The Noether status of the corpus: adjudication and scope}.}\ Structural-constraints record: see Appendix B, Section B.1.15 of the Addendum.\par\medskip

\phantomsection
\section[Tier BM.5 — The Calibration Lock: the Born Width, the…]{\texorpdfstring{Tier BM.5 — The Calibration Lock: the Born Width, the Coherence Volume, the Curvature of a Sealed Composite, and the Neutron Row}{Tier BM.5 - The Calibration Lock: the Born Width, the Coherence Volume, the Curvature of a Sealed Composite, and the Neutron Row}}

\begin{definition}[{\texorpdfstring{Def-BM.4 — The Born width and the coherence volume}{Def-BM.4 - The Born width and the coherence volume}}]\label{tag:Def-BM-4}
\textbf{Def.} Let $E_{\mathrm{int}}\in\mathbb Z_{\ge2}$ be the internal directed-edge count of a sealed composite [D0.8; FT.44(4)]: $E_{\mathrm{int}}=2E_{\mathrm{bidir}}+E_{\mathrm{uni}}$; the minimal floor is the digon, $E_{\mathrm{int}}=2$ [FT.28] and $m_{\mathrm{hid}}=E_{\mathrm{int}}\ln2/\mathfrak h$ its informational mass [FT.44(4), FT.42(1)].

%dsp-level:0
\[
(1)\quad\textbf{the Born width:}
\]

%dsp-level:0
\[
\qquad \Delta x_{\mathrm B}\ :=\ \frac1{m_{\mathrm{hid}}}\ =\ \frac{\mathfrak h}{E_{\mathrm{int}}\ln2}
\]

%dsp-level:0
\[
\qquad [\hbar=c=1,\ \mathrm{D2.9};\ \text{the Compton-like length of}\ \mathrm{FT.90}(6)\ \text{re-cited as an existing read}];
\]

%dsp-level:0
\[
\textbf{constitutive clause:}\ \ \Delta x_{\mathrm B}\cdot m_{\mathrm{hid}}=1\ \text{is the definition};
\]

%dsp-level:0
\[
\qquad \textbf{non-derived clause:}\ \ \text{the dimensionless normalization is fixed to}\ 1\ \text{(the Compton normalization),}\ \neg\text{derived};
\]

%dsp-level:0
\[
(2)\quad\textbf{the dimensional-uniqueness clause:}
\]

%dsp-level:0
\[
\qquad \exists!\ \text{free dimensionful scale}\ \mathfrak h\ [\mathrm{P1},\ \mathrm{CT2.11}]:\ \ \forall Q\ \text{dimensionful}:\ Q
\]

%dsp-level:0
\[
= c_Q\,
\]

%dsp-level:0
\[
\mathfrak h^{k_Q},
\]

%dsp-level:0
\[
\ c_Q\ \text{fixed}\ [\mathrm{CT2.11}];
\]

%dsp-level:0
\[
[\text{length}]=\mathfrak q^{-1/2}=[\mathfrak h],
\]

%dsp-level:0
\[
\qquad [m_{\mathrm{hid}}]=\mathfrak q^{+1/2}\ [\mathrm{CT2.1}]\implies\text{every length built from}\ m_{\mathrm{hid}}\ \text{alone is}\ c/m_{\mathrm{hid}}\ \text{with}\ c\ \text{dimensionless}:
\]

%dsp-level:0
\[
\text{the Born width is the unique such length at}\ c=1;
\]

%dsp-level:0
\[
\qquad \text{any other}\ c\ \text{is a P4-class internal choice and is priced, not adopted}\ [\text{clause}\ (4)];
\]

%dsp-level:0
\[
(3)\quad\textbf{the coherence volume:}\qquad V_{\mathrm{eff}}\ :=\ (\Delta x_{\mathrm B})^3\ =\ \frac{\mathfrak h^3}{(E_{\mathrm{int}}\ln2)^3}\qquad[\alpha=1:\ \text{the cube convention}];
\]

%dsp-level:0
\[
(4)\quad\textbf{the geometric-factor clause (priced, not adopted):}
\]

%dsp-level:0
\[
\qquad V=\alpha\,(\Delta x_{\mathrm B})^3\ \text{with}\ \alpha\ \text{dimensionless}:
\]

%dsp-level:0
\[
\qquad \rho(\alpha)=\rho(1)/\alpha,
\]

%dsp-level:0
\[
\qquad \alpha=\tfrac{4\pi}3:\ \rho\mapsto\tfrac3{4\pi}\rho\ \ \text{exactly};
\]

%dsp-level:0
\[
\text{every ratio in which}\ \alpha\ \text{appears in both entries is}\ \alpha\text{-free};
\]

%dsp-level:0
\[
\qquad \alpha\ \text{is}\ \neg\text{a dimensionful datum and}\ \neg\text{a P4 core member};
\]

%dsp-level:0
\[
(5)\quad\textbf{the locus clause:}
\]

%dsp-level:0
\[
\qquad \Delta x_{\mathrm B}\ \text{is a chart-side envelope on}\ M\ [\mathrm{CT2.6}]\ \text{in the sense of BM.10(6)'s width law},
\]

%dsp-level:0
\[
\ \neg\text{a materialized edge length}:
\]

%dsp-level:0
\[
\text{the minimal materialized line element is}\ \widehat{ds}^{\,2}=1\implies ds^2=\mathfrak h^2\ [\mathrm{D1.6};\ \mathrm{CT1.19a}]:
\]

%dsp-level:0
\[
\qquad \Delta x_{\mathrm B}<\mathfrak h\iff E_{\mathrm{int}}\ln2>1\iff E_{\mathrm{int}}\ge2\ [\text{forward pointer:}\ \mathrm{BM.13}(3);\ \text{machine}];
\]

%dsp-level:0
\[
\text{every admissible cell is sub-hop, so the width is a coarse-graining envelope and}\ \neg\text{a ledger distance}\ [\text{locus discipline},\ \mathrm{FT.136}(5)\ \text{pattern}];
\]

%dsp-level:0
\[
(6)\quad\textbf{the scope clause:}
\]

%dsp-level:0
\[
\qquad \mathrm{FT.90}(6)\ \text{refutes a Compton-wavelength relation to the generation threshold}\ (\delta h\ge\lceil\log_2\#\text{classes}\rceil\ \text{depends on the class count alone}):
\]

%dsp-level:0
\[
\text{that refutation is untouched here}\ —\ \text{this definition uses the length as a spread read, and derives no relation between it and any threshold}.
\]

\textbf{Rmk (naming).} \textit{Born width} names the spatial spread of the sealed composite's density read at the Compton normalization; \textit{coherence volume} names its cube. Both are internal bookkeeping on derived quantities [A §0.2]; neither is a measurement, and no external length datum enters [P5; Def-BM.3(6)].

\textbf{Req.} D0.6, D0.8, D2.9 (Parts 1–2); D1.6 (Part 1); CT1.19a, CT2.1, CT2.6, CT2.11 (Part 2); P1, P5 (Parts 1–2); Def-BM.3 (this part); FT.28, FT.42(1), FT.44(4) (Part 3, Census); FT.90(6) (Part 3); BM.10(6) (this part); the sub-hop row cited in clause (5) is a marked forward pointer at the citing site and is imported in no \textbf{Req} line; exact Laurent-monomial arithmetic, exact rational $\ln2$ enclosures [math; machine: \texttt{checks/\allowbreak{}verify\_\allowbreak{}H\_\allowbreak{}calibration\_\allowbreak{}lock.py} §A].

\textbf{Pf.}

%dsp-level:0
\[
\text{Definition block: the constitutive clause is}\ \Delta x_{\mathrm B}\cdot m_{\mathrm{hid}}
\]

%dsp-level:0
\[
= 1;
\]

%dsp-level:0
\[
\ \text{the non-derived clauses are the normalization}\ c
\]

%dsp-level:0
\[
= 1\ \text{of}\ (1)\text{–}(2),
\]

%dsp-level:0
\[
\ \text{the geometric factor}\ \alpha\ \text{of}\ (4),
\]

%dsp-level:0
\[
\ \text{and the locus reading of}\ (5);
\]

%dsp-level:0
\[
\ \text{clause}\ (6)\ \text{is a scope record}.\ \blacksquare
\]

\textbf{Phys.} The sealed composite has one length and it is the reciprocal of its own informational mass: $\mathfrak h$ over the edge count times $\ln2$. That is not a new postulate — the corpus already carried the Compton-like length as an existing read, and with one free scale in the theory there is no other length a mass can build. The cube of that width is the volume the density spreads over, with a dimensionless geometric factor left free and priced rather than chosen. One clause keeps the reading honest: the width is a chart-side envelope, not a hop. Every admissible cell turns out to be \textit{shorter} than a single materialized line element, which is exactly why the width belongs to the continuum chart and not to the ledger.

\end{definition}

\begin{theorem}[{\texorpdfstring{BM.13 — The Born width: the exact substitution, the quantized ladder, and the sub-hop locus}{BM.13 - The Born width: the exact substitution, the quantized ladder, and the sub-hop locus}}]\label{tag:BM-13}
\textbf{Thm.}

%dsp-level:0
\[
(1)\quad\textbf{the exact substitution:}
\]

%dsp-level:0
\[
\qquad \Delta x_{\mathrm B}=\frac1{m_{\mathrm{hid}}}=\frac{\mathfrak h}{E_{\mathrm{int}}\ln2}
\]

%dsp-level:0
\[
\qquad [\mathrm{Def\text{-}BM.4}(1);\ \mathrm{FT.44}(4);\ \text{machine}:\ \Delta x_{\mathrm B}\cdot m_{\mathrm{hid}}=1\ \text{as a Laurent identity}];
\]

%dsp-level:0
\[
\text{the monomial census is exact:}\quad\Delta x_{\mathrm B}=\mathfrak h^1E_{\mathrm{int}}^{-1}(\ln2)^{-1}\pi^0
\]

%dsp-level:0
\[
\qquad [\text{machine: one term, coefficient}\ 1];
\]

%dsp-level:0
\[
[\Delta x_{\mathrm B}]=[\mathfrak h]=\mathfrak q^{-1/2}=[\text{length}]
\]

%dsp-level:0
\[
\qquad [\mathrm{CT2.1};\ \text{machine: exact exponent arithmetic}];
\]

%dsp-level:0
\[
\qquad E_{\mathrm{int}},\ \ln2:\ \text{dimensionless};
\]

%dsp-level:0
\[
(2)\quad\textbf{quantization:}
\]

%dsp-level:0
\[
\qquad \Delta x_{\mathrm B}\in\Bigl\{\frac{\mathfrak h}{n\ln2}:\ n\in\mathbb Z_{\ge2}\Bigr\}:
\]

%dsp-level:0
\[
\qquad \text{a discrete ladder},\ \neg\text{a continuous range}\ [\text{machine: injective in}\ n];
\]

%dsp-level:0
\[
\textbf{the computed cells:}
\]

%dsp-level:0
\[
\quad E_{\mathrm{int}}
\]

%dsp-level:0
\[
= 2\ [\Sigma_{\min};\ \mathrm{FT.28}],
\]

%dsp-level:0
\[
\ 6\ [\Sigma_{12};\ \text{Census FT.35}(2)],
\]

%dsp-level:0
\[
\ 12,
\]

%dsp-level:0
\[
\ 24,
\]

%dsp-level:0
\[
\ 24,
\]

%dsp-level:0
\[
\ 60,
\]

%dsp-level:0
\[
\ 60\ [\text{the five solids};\ \text{Census FT.42}(5)];
\]

%dsp-level:0
\[
\text{the ladder is strictly decreasing with disjoint exact enclosures}\ [\text{machine: rational}\ \ln2\ \text{enclosures of width}\ 2^{-12}];
\]

%dsp-level:0
\[
\qquad \text{the cycle ladder}\ E_{\mathrm{int}}=L\ [\mathrm{FT.32}]\ \text{has no repeated value}\ [\text{machine}];
\]

%dsp-level:0
\[
(3)\quad\textbf{the sub-hop locus:}
\]

%dsp-level:0
\[
\qquad \Delta x_{\mathrm B}<\mathfrak h\iff\frac1{E_{\mathrm{int}}\ln2}<1\iff E_{\mathrm{int}}\ln2>1\iff E_{\mathrm{int}}\ge2
\]

%dsp-level:0
\[
\qquad [\text{machine: exact enclosure}\ \tfrac12<\ln2<1\ \text{drives every step}];
\]

%dsp-level:0
\[
\text{the minimal admissible count is}\ E_{\mathrm{int}}
\]

%dsp-level:0
\[
= 2\ [\mathrm{FT.28}:\ \text{the digon floor};\ \text{a materialized closed chain carries}\ \ge2\ \text{directed edges}]:
\]

%dsp-level:0
\[
\therefore\ \textbf{every admissible sealed composite is sub-hop:}\quad\Delta x_{\mathrm B}\le\frac{\mathfrak h}{2\ln2}<\mathfrak h
\]

%dsp-level:0
\[
\qquad [\text{machine: all computed cells}];
\]

%dsp-level:0
\[
\text{at the inadmissible}\ E_{\mathrm{int}}
\]

%dsp-level:0
\[
= 1\ \text{the order reverses}\ (\Delta x_{\mathrm B}>\mathfrak h):\ \text{the threshold is sharp and sits exactly at the floor}\ [\text{machine}];
\]

%dsp-level:0
\[
\textbf{locus discipline:}\ \ \Delta x_{\mathrm B}\ \text{is a}\ \mathrm{CT2.6}\ \text{chart envelope},
\]

%dsp-level:0
\[
\ \neg\text{a ledger distance}\ [\mathrm{Def\text{-}BM.4}(5);\ \mathrm{D1.6}:\ \text{the minimal materialized element is}\ ds^2=\mathfrak h^2];
\]

%dsp-level:0
\[
(4)\quad\textbf{the Robertson consistency:}
\]

%dsp-level:0
\[
\qquad \Delta x\,\Delta p\ge\tfrac12\ [\mathrm{BM.10}(3),(6);\ \hbar=1]\ \text{at}\ \Delta x=\Delta x_{\mathrm B}\iff\Delta p\ge\frac1{2\Delta x_{\mathrm B}}=\frac{m_{\mathrm{hid}}}2
\]

%dsp-level:0
\[
\qquad [\text{machine}];
\]

%dsp-level:0
\[
\Delta p=\tfrac12m_{\mathrm{hid}}\ \text{saturates Robertson exactly at the Born width}:
\]

%dsp-level:0
\[
\qquad \text{the width is the saturating spread for a momentum spread of half the informational mass};
\]

%dsp-level:0
\[
(5)\quad\textbf{the}\ \alpha\textbf{-variant priced:}
\]

%dsp-level:0
\[
\qquad V=\alpha(\Delta x_{\mathrm B})^3\implies\rho(\alpha)=\rho(1)/\alpha;
\]

%dsp-level:0
\[
\qquad \alpha=\tfrac{4\pi}3\implies\rho\mapsto\tfrac3{4\pi}\rho\ \ \text{exactly}\ [\text{machine: a single Laurent monomial}].
\]

\textbf{Rmk (comparison discipline).} No clause reads $\Delta x_{\mathrm B}$ against an external length. The comparisons of this tag are internal: against the hop $\mathfrak h$ [D0.6], against the minimal materialized element [D1.6], and against the Robertson bound [BM.10(6)]. External comparison remains quarantined [P5; Census FT.114, FT.120].

\textbf{Req.} D0.6, D0.8, D1.6, D2.9 (Parts 1–2); CT1.19a, CT2.1, CT2.6 (Part 2); P1, P5 (Parts 1–2); Def-BM.4 (this part); FT.28, FT.32, FT.42(1), FT.44(4), FT.90(6) (Part 3, Census); Census FT.35(2), FT.42(5) (row-wise); BM.10(3), BM.10(6) (this part); exact Laurent-monomial identities, exact rational $\ln2$ enclosures via $\ln2=\sum_{n\ge1}1/(n2^n)$ with tail $<2^{-N}$ [math; machine: \texttt{checks/\allowbreak{}verify\_\allowbreak{}H\_\allowbreak{}calibration\_\allowbreak{}lock.py} §A].

\textbf{Pf.}

%dsp-level:0
\[
(1):\ m_{\mathrm{hid}}
\]

%dsp-level:0
\[
= E_{\mathrm{int}}\ln2/\mathfrak h\ [\mathrm{FT.44}(4)\ \text{re-cited}]
\]

%dsp-level:0
\[
\implies 1/m_{\mathrm{hid}}
\]

%dsp-level:0
\[
= \mathfrak h/(E_{\mathrm{int}}\ln2)\ [\text{machine: the product is}\ 1\ \text{in}\ \mathbb Q[\mathfrak h,E,\ln2]^{\pm}];
\]

%dsp-level:0
\[
\qquad \text{exponents}: [\mathfrak h]-0-0=\mathfrak q^{-1/2}\ [\mathrm{CT2.1}]
\]

%dsp-level:0
\[
(2):\ \Delta x_{\mathrm B}(n)
\]

%dsp-level:0
\[
= \mathfrak h/(n\ln2):\ n_1
\]

%dsp-level:0
\[
\ne n_2\implies\Delta x_{\mathrm B}(n_1)\ne\Delta x_{\mathrm B}(n_2)\ [\text{machine: injectivity on}\ n\in2..59];
\]

%dsp-level:0
\[
\qquad \text{the five computed cells}: \text{disjoint exact enclosures}\ [\text{machine}:\ \ln2\in(\textstyle\sum_1^{12}\tfrac1{n2^n},\ \cdot+2^{-12})]
\]

%dsp-level:0
\[
(3):\ \tfrac12<\ln2<1\ [\text{machine: the alternating-free series with exact rational partial sums}]
\]

%dsp-level:0
\[
\implies E\ln2>1
\]

%dsp-level:0
\[
\iff E
\]

%dsp-level:0
\[
\ge 2\ [\text{machine: every}\ E\in1..59];
\]

%dsp-level:0
\[
\qquad E_{\mathrm{int}}\ge2\ [\mathrm{FT.28}]\implies\Delta x_{\mathrm B}\le\mathfrak h/(2\ln2)<\mathfrak h
\]

%dsp-level:0
\[
(4):\ \Delta x\,
\]

%dsp-level:0
\[
\Delta p\ge\tfrac12\ [\mathrm{BM.10}(3):\ \text{Robertson};\ \mathrm{BM.10}(6):\ \text{the width law}]\ \text{with}\ \Delta x
\]

%dsp-level:0
\[
= \Delta x_{\mathrm B}
\]

%dsp-level:0
\[
= 1/m_{\mathrm{hid}}\implies\Delta p
\]

%dsp-level:0
\[
\ge m_{\mathrm{hid}}/2\ [\text{machine: exact rationals}];
\]

%dsp-level:0
\[
\qquad \text{equality at}\ \Delta p=m_{\mathrm{hid}}/2\ [\mathrm{BM.10}(3)\ \text{saturation row}]
\]

%dsp-level:0
\[
(5):\ \rho=m_{\mathrm{hid}}/(\alpha\Delta x_{\mathrm B}^3)=\rho(1)/\alpha\ [\text{machine}];
\]

%dsp-level:0
\[
\qquad \alpha=4\pi/3:\ 3/(4\pi)\ \text{is a single Laurent monomial}\ [\text{machine}].\ \blacksquare
\]

\textbf{Phys.} The width of a sealed composite is a reciprocal integer count times the hop, divided by $\ln2$ — one Laurent monomial, no mixture of scales, nothing to tune. The ladder it walks is strictly discrete: two, six, twelve, twenty-four, sixty, each step shorter than the last with exact rational enclosures that never touch. And every rung of it sits \textit{below} a single hop, which is the sharpest locus statement in the tier: the Born width is not a distance the ledger can materialize, it is the envelope the continuum chart assigns to a bound whose internal edges all cancel pairwise. Consistency with the uncertainty bound is exact rather than approximate — half the informational mass is precisely the momentum spread that saturates Robertson at this width.

\end{theorem}

\begin{theorem}[{\texorpdfstring{BM.14 — The coherence volume and the energy density}{BM.14 - The coherence volume and the energy density}}]\label{tag:BM-14}
\textbf{Thm.}

%dsp-level:0
\[
(1)\quad\textbf{the exact volume:}
\]

%dsp-level:0
\[
\qquad V_{\mathrm{eff}}=(\Delta x_{\mathrm B})^3=\frac{\mathfrak h^3}{(E_{\mathrm{int}}\ln2)^3}
\]

%dsp-level:0
\[
\qquad [\mathrm{Def\text{-}BM.4}(3);\ \text{machine: one Laurent monomial}\ \mathfrak h^3E^{-3}(\ln2)^{-3}];
\]

%dsp-level:0
\[
[V_{\mathrm{eff}}]=\mathfrak q^{-3/2}=[\text{length}]^3\qquad[\mathrm{CT2.1};\ \text{machine:}\ 3\cdot(-\tfrac12)=-\tfrac32];
\]

%dsp-level:0
\[
(2)\quad\textbf{the exact energy density:}
\]

%dsp-level:0
\[
\qquad \rho=\frac{m_{\mathrm{hid}}}{V_{\mathrm{eff}}}=\frac{E_{\mathrm{int}}\ln2/\mathfrak h}{\mathfrak h^3/(E_{\mathrm{int}}\ln2)^3}=\frac{(E_{\mathrm{int}}\ln2)^4}{\mathfrak h^4}
\]

%dsp-level:0
\[
\qquad [\text{machine: Laurent identity, denominators cleared}];
\]

%dsp-level:0
\[
\rho\,\mathfrak h^4=(E_{\mathrm{int}}\ln2)^4\ \ \text{exactly};
\]

%dsp-level:0
\[
\qquad [\rho]=[\text{energy}]-[V]=\mathfrak q^{1/2+3/2}=\mathfrak q^{+2}=[T_{\mu\nu}]
\]

%dsp-level:0
\[
\qquad [\mathrm{CT2.1};\ \mathrm{BM.9}(5);\ \text{machine}];
\]

%dsp-level:0
\[
(3)\quad\textbf{the agreement with the general source read:}
\]

%dsp-level:0
\[
\qquad \kappa^2\rho=\frac{2\pi\,\mathfrak h\,E_{\mathrm{int}}}{V}\bigg|_{V=V_{\mathrm{eff}}}=\frac{2\pi\,E_{\mathrm{int}}^4(\ln2)^3}{\mathfrak h^2}
\]

%dsp-level:0
\[
\qquad [\mathrm{BM.9}(5);\ \text{machine: the two forms are the same monomial}];
\]

%dsp-level:0
\[
\kappa^2=8\pi G_N=\frac{2\pi\,\mathfrak h^2}{\ln2},\qquad G_N=\frac{\mathfrak h^2}{4\ln2}\qquad[\mathrm{CT1.23},\ \mathrm{CT2.9},\ \mathrm{P3};\ \text{machine: exact}];
\]

%dsp-level:0
\[
(4)\quad\textbf{the}\ \ln2\ \textbf{ledger:}
\]

%dsp-level:0
\[
\qquad m_{\mathrm{hid}}:\ (\ln2)^1\ [\mathrm{FT.44}(4)];
\]

%dsp-level:0
\[
\qquad V_{\mathrm{eff}}:\ (\ln2)^{-3};
\]

%dsp-level:0
\[
\qquad \rho:\ (\ln2)^4=(\ln2)^1\cdot(\ln2)^3;
\]

%dsp-level:0
\[
\qquad \kappa^2:\ (\ln2)^{-1};
\]

%dsp-level:0
\[
\qquad \kappa^2\rho:\ (\ln2)^3\cdot2\pi;
\]

%dsp-level:0
\[
\text{at a general}\ V\ [\mathrm{BM.9}(5)]\ \text{the}\ \ln2\ \text{of}\ m_{\mathrm{hid}}\ \text{and of}\ \kappa^2\ \text{cancel and no}\ \ln2\ \text{survives};
\]

%dsp-level:0
\[
\qquad \text{at}\ V=V_{\mathrm{eff}}\ \text{the volume contributes}\ (\ln2)^{-3}\ \text{and three powers survive}:
\]

%dsp-level:0
\[
\text{both statements are the same exact identity read at two volumes}\ [\text{machine: the monomial censuses agree}];
\]

%dsp-level:0
\[
(5)\quad\textbf{the}\ \alpha\textbf{-variant:}
\]

%dsp-level:0
\[
\qquad \rho(\alpha)=\frac{(E_{\mathrm{int}}\ln2)^4}{\alpha\,\mathfrak h^4},
\]

%dsp-level:0
\[
\qquad \kappa^2\rho(\alpha)=\frac{2\pi E_{\mathrm{int}}^4(\ln2)^3}{\alpha\,\mathfrak h^2};
\]

%dsp-level:0
\[
\qquad \alpha=\tfrac{4\pi}3\implies\kappa^2\rho\mapsto\frac{3\,E_{\mathrm{int}}^4(\ln2)^3}{2\,\mathfrak h^2}\ \ [\text{machine: the}\ \pi\ \text{cancels exactly}].
\]

\textbf{Rmk (volume convention).} Clause (5) prices the spherical convention instead of adopting it: at $\alpha=4\pi/3$ the $\pi$ of $\kappa^2$ cancels and the curvature scale becomes $\pi$-free. Both conventions are exact; the cube convention of Def-BM.4(3) is the one carried downstream, and every $\alpha$-free ratio is unaffected.

\textbf{Req.} D2.9 (Parts 1–2); CT1.23, CT2.1, CT2.9 (Part 2); P3 (Parts 1–2); Def-BM.4 (this part); FT.42(1), FT.44(4) (Part 3, Census); BM.9(5), BM.13(1) (this part); exact Laurent-monomial arithmetic with denominators cleared [math; machine: \texttt{checks/\allowbreak{}verify\_\allowbreak{}H\_\allowbreak{}calibration\_\allowbreak{}lock.py} §B].

\textbf{Pf.}

%dsp-level:0
\[
(1):\ \text{cubing}\ \mathrm{BM.13}(1)\text{'s monomial}\ [\text{machine}];\qquad3\cdot(-\tfrac12)=-\tfrac32\ [\mathrm{CT2.1}]
\]

%dsp-level:0
\[
(2):\ \rho=\frac{E\ln2}{\mathfrak h}\cdot\frac{(E\ln2)^3}{\mathfrak h^3}=\frac{(E\ln2)^4}{\mathfrak h^4}\ [\text{machine}];
\]

%dsp-level:0
\[
\qquad [\rho]=\mathfrak q^{1/2-(-3/2)}=\mathfrak q^{+2}=[T_{\mu\nu}]\ [\mathrm{CT2.1};\ \mathrm{BM.9}(5)\ \text{row}]
\]

%dsp-level:0
\[
(3):\ \kappa^2\rho=\tfrac{2\pi\mathfrak h^2}{\ln2}\cdot\tfrac{(E\ln2)^4}{\mathfrak h^4}=\tfrac{2\pi E^4(\ln2)^3}{\mathfrak h^2}\ [\text{machine}];
\]

%dsp-level:0
\[
\qquad \mathrm{BM.9}(5):\ \kappa^2\rho
\]

%dsp-level:0
\[
= 2\pi\mathfrak h E/V\ \text{at}\ V
\]

%dsp-level:0
\[
= \mathfrak h^3/(E\ln2)^3\implies\tfrac{2\pi\mathfrak h E(E\ln2)^3}{\mathfrak h^3}
\]

%dsp-level:0
\[
= \tfrac{2\pi E^4(\ln2)^3}{\mathfrak h^2}\ [\text{machine: identical monomials}]
\]

%dsp-level:0
\[
(4):\ \text{monomial censuses}\ [\text{machine}]:\ m_{\mathrm{hid}}\to(0,1,1,0),
\]

%dsp-level:0
\[
\ V_{\mathrm{eff}}\to(3,-3,-3,0),
\]

%dsp-level:0
\[
\ \rho\to(-4,4,4,0),
\]

%dsp-level:0
\[
\ \kappa^2\to(2,0,-1,1),
\]

%dsp-level:0
\[
\ \kappa^2\rho\to(-2,4,3,1)\ \text{in the basis}\ (\mathfrak h,E,\ln2,\pi)
\]

%dsp-level:0
\[
(5):\ \text{division by}\ \alpha\ [\text{machine}];\qquad\alpha=4\pi/3:\ (2\pi)/(4\pi/3)=3/2\ \text{exactly}.\ \blacksquare
\]

\textbf{Phys.} Cube the width and the density follows by one exact division: the fourth power of the information content over the fourth power of the hop. The dimensions land precisely on the stress-energy tensor's own, with no adjustment and no second scale. The same quantity arrives a second way — through the general source read of the stress-energy tier with the volume now supplied by the width — and the two routes give the identical Laurent monomial, which is the first sign that the calibration is closing rather than accumulating. The $\ln2$ bookkeeping is explicit about where its powers live: one from the mass, three from the volume, minus one from the coupling, three surviving in the curvature. And the geometric factor is priced rather than hidden: choose the sphere and the $\pi$ of the coupling cancels outright.

\end{theorem}

\begin{theorem}[{\texorpdfstring{BM.15 — The curvature of a sealed composite: the exact trace of the field equations, and the curvature scale}{BM.15 - The curvature of a sealed composite: the exact trace of the field equations, and the curvature scale}}]\label{tag:BM-15}
\textbf{Thm.} In the orthonormal rest frame $g_{\mu\nu}=\mathrm{diag}(-1,1,1,1)$ [CT1.12]: $\operatorname{sig}=(-,+,+,+)$; [A3, C13], with $T_{\mu\nu}=\mathrm{diag}(\rho,p,p,p)$ [BM.9(4)] and $G_{\mu\nu}+\Lambda g_{\mu\nu}=\kappa^2T_{\mu\nu}$ [CT1.23]:

%dsp-level:0
\[
(1)\quad\textbf{the dust trace:}
\]

%dsp-level:0
\[
\qquad T=g^{\mu\nu}T_{\mu\nu}=-\rho+3p
\]

%dsp-level:0
\[
\qquad [\text{machine: exact component algebra}];
\]

%dsp-level:0
\[
\qquad p=0\ [\mathrm{CT4.5}]\implies T=-\rho;
\]

%dsp-level:0
\[
(2)\quad\textbf{the exact trace of the field equations:}\qquad G^\mu{}_\mu=-R\ \ \text{in}\ 4\ \text{dimensions}\ [\text{math}]\implies-R+4\Lambda=\kappa^2T:
\]

%dsp-level:0
\[
\boxed{\ R=\kappa^2(\rho-3p)+4\Lambda\ }\qquad\text{and at}\ p=0:\qquad R=\kappa^2\rho+4\Lambda\qquad[\text{machine: exact in}\ \mathbb Q[\rho,p,\Lambda,\kappa^2]];
\]

%dsp-level:0
\[
\textbf{the inadmissible form:}\quad R=-\kappa^2\rho\ \text{would require}\ 2\kappa^2\rho-3\kappa^2p+4\Lambda=0\ \text{identically}:
\]

%dsp-level:0
\[
\qquad \textbf{refuted exactly}\ [\text{machine: the residual is a nonzero polynomial}];
\]

%dsp-level:0
\[
\text{the sign of the dust curvature is fixed by CT1.23's own trace}: \ \kappa^2>0,
\]

%dsp-level:0
\[
\ \rho
\]

%dsp-level:0
\[
\ge 0\ [\mathrm{CT4.5}]
\]

%dsp-level:0
\[
\implies R
\]

%dsp-level:0
\[
\ge 4\Lambda,
\]

%dsp-level:0
\[
\ \text{and}\ R
\]

%dsp-level:0
\[
= \kappa^2\rho>0\ \text{at}\ \Lambda\qquad = 0;
\]

%dsp-level:0
\[
(3)\quad\textbf{the component round trip:}
\]

%dsp-level:0
\[
\qquad R_{\mu\nu}=\kappa^2\bigl(T_{\mu\nu}-\tfrac12g_{\mu\nu}T\bigr)+\Lambda g_{\mu\nu}
\]

%dsp-level:0
\[
\qquad [\text{machine: all}\ 4\ \text{components exact}];
\]

%dsp-level:0
\[
p=0:\quad R_{00}=\tfrac12\kappa^2\rho-\Lambda,
\]

%dsp-level:0
\[
\qquad R_{ii}=\tfrac12\kappa^2\rho+\Lambda\ \ (i=1,2,3);
\]

%dsp-level:0
\[
\qquad g^{\mu\nu}R_{\mu\nu}=\kappa^2\rho+4\Lambda=R\ \ [\text{machine: the round trip closes}];
\]

%dsp-level:0
\[
G_{00}=\kappa^2\rho+\Lambda,
\]

%dsp-level:0
\[
\qquad G_{ii}=\kappa^2p-\Lambda\ \overset{p=0}{=}\ -\Lambda;
\]

%dsp-level:0
\[
\qquad G_{\mu\nu}=\kappa^2T_{\mu\nu}-\Lambda g_{\mu\nu}\ \ [\text{machine: all}\ 4\ \text{components}];
\]

%dsp-level:0
\[
(4)\quad\textbf{the curvature scale of the sealed composite:}
\]

%dsp-level:0
\[
\qquad \kappa^2\rho=\frac{2\pi\,E_{\mathrm{int}}^4(\ln2)^3}{\mathfrak h^2}
\]

%dsp-level:0
\[
\qquad [\mathrm{BM.14}(3);\ \text{machine}];
\]

%dsp-level:0
\[
[\kappa^2\rho]=\mathfrak q^{-1}\cdot\mathfrak q^{+2}=\mathfrak q^{+1}=[G_{\mu\nu}]=[\Lambda]=[R]
\]

%dsp-level:0
\[
\qquad [\mathrm{CT2.1};\ \text{machine: exact exponent arithmetic}];
\]

%dsp-level:0
\[
\text{the scale is a pure function of the integer count}\ E_{\mathrm{int}}\ \text{and the two universal factors}\ \mathfrak h,
\]

%dsp-level:0
\[
\ \ln2\ [\mathrm{P1};\ \mathrm{FT.44}(4)]:\ \neg\text{a continuous mass variable enters};
\]

%dsp-level:0
\[
(5)\quad\textbf{the curvature radius:}
\]

%dsp-level:0
\[
\qquad \ell_R:=\frac1{\sqrt{\kappa^2\rho}}:
\]

%dsp-level:0
\[
\qquad \Bigl(\frac{\Delta x_{\mathrm B}}{\ell_R}\Bigr)^{\!2}=(\Delta x_{\mathrm B})^2\,\kappa^2\rho=2\pi\,E_{\mathrm{int}}^2\ln2
\]

%dsp-level:0
\[
\qquad [\text{machine: one Laurent monomial}];
\]

%dsp-level:0
\[
2\pi\ln2>1\ [\text{machine: exact enclosures}\ \pi>\tfrac{223}{71},\ \ln2>\textstyle\sum_1^{12}\tfrac1{n2^n}]\implies\Bigl(\frac{\Delta x_{\mathrm B}}{\ell_R}\Bigr)^{\!2}>1\ \ \forall E_{\mathrm{int}}\ge1:
\]

%dsp-level:0
\[
\qquad \ell_R<\Delta x_{\mathrm B}\ \ \text{on the whole ladder};
\]

%dsp-level:0
\[
\text{the ratio is strictly increasing in}\ E_{\mathrm{int}}\ \text{with disjoint exact enclosures}\ [\text{machine}]:
\]

%dsp-level:0
\[
\qquad \text{the self-curvature of a sealed composite tightens as its information content grows};
\]

%dsp-level:0
\[
(6)\quad\textbf{the}\ \Lambda\ \textbf{clause:}
\]

%dsp-level:0
\[
\qquad \Lambda\ \text{is a P4 core member}\ [\mathrm{P4}]:
\]

%dsp-level:0
\[
\qquad R\ \text{carries exactly one additive free term},\ 4\Lambda,\ \text{with exact coefficient}\ 4\ [\text{machine}];
\]

%dsp-level:0
\[
\text{the locked part is}\ \kappa^2\rho=R-4\Lambda\ \ \text{at}\ p=0:
\]

%dsp-level:0
\[
\qquad \text{the curvature generated by the composite is calibration-locked; the ambient term is not}.
\]

\textbf{Rmk (frame and scope).} Clauses (1)–(3) are component algebra in one orthonormal frame with a diagonal metric; the identities are rational in $(\rho,p,\Lambda,\kappa^2)$ and hold in every frame [math: tensoriality]. Clause (5) compares lengths by squares, so no radical enters any assertion. No solution of the field equations is constructed here: the pp-wave witness of BM.11(4) carries the differential identities, and this tag carries the algebraic dust read.

\textbf{Req.} A3 (Part 1); CT1.12, CT1.23, CT2.1, CT2.9 (Part 2); CT4.5 (Part 2); P1, P4 (Parts 1–2); Def-BM.4 (this part); FT.44(4) (Census); BM.9(4), BM.11(4), BM.13(1), BM.14(3) (this part); exact multivariate polynomial algebra over $\mathbb Q$, Archimedes' rational bounds for $\pi$, exact rational $\ln2$ enclosures [math; machine: \texttt{checks/\allowbreak{}verify\_\allowbreak{}H\_\allowbreak{}calibration\_\allowbreak{}lock.py} §C].

\textbf{Pf.}

%dsp-level:0
\[
(1):\ g^{\mu\nu}T_{\mu\nu}=(-1)\rho+1\cdot p+1\cdot p+1\cdot p=-\rho+3p\ [\text{machine}];\qquad p=0\ [\mathrm{CT4.5}:\ w=p/\rho=0]
\]

%dsp-level:0
\[
(2):\ G_{\mu\nu}=R_{\mu\nu}-\tfrac12g_{\mu\nu}R\implies G^\mu{}_\mu=R-2R=-R\ [\text{math},\ 4\ \text{dimensions}];
\]

%dsp-level:0
\[
\qquad \mathrm{CT1.23}\ \text{traced}:\ -R+4\Lambda=\kappa^2(-\rho+3p)\implies R=\kappa^2(\rho-3p)+4\Lambda\ [\text{machine}]
\]

%dsp-level:0
\[
\text{residual of}\ R
\]

%dsp-level:0
\[
= -\kappa^2\rho:\ \bigl[\kappa^2(\rho-3p)+4\Lambda\bigr]+\kappa^2\rho
\]

%dsp-level:0
\[
= 2\kappa^2\rho-3\kappa^2p+4\Lambda
\]

%dsp-level:0
\[
\ne 0\ \text{as a polynomial}\ [\text{machine}]
\]

%dsp-level:0
\[
(3):\ R_{\mu\nu}
\]

%dsp-level:0
\[
= \kappa^2T_{\mu\nu}+\tfrac12g_{\mu\nu}R-\Lambda g_{\mu\nu}
\]

%dsp-level:0
\[
= \kappa^2(T_{\mu\nu}-\tfrac12g_{\mu\nu}T)+\Lambda g_{\mu\nu}\ [\text{substituting}\ (2);\ \text{machine}];
\]

%dsp-level:0
\[
\qquad G_{\mu\nu}=R_{\mu\nu}-\tfrac12g_{\mu\nu}R=\kappa^2T_{\mu\nu}-\Lambda g_{\mu\nu}\ [\text{machine: all}\ 4];
\]

%dsp-level:0
\[
\qquad \text{the dust values as displayed}\ [\text{machine}]
\]

%dsp-level:0
\[
(4):\ \mathrm{BM.14}(3)\ \text{re-cited};
\]

%dsp-level:0
\[
\qquad \text{exponents}: -1+2=+1\ [\mathrm{CT2.1}:\ [\kappa^2]=\mathfrak q^{-1},\ [T_{\mu\nu}]=\mathfrak q^{+2},\ [G_{\mu\nu}]=[\Lambda]=\mathfrak q^{+1}]
\]

%dsp-level:0
\[
(5):\ (\Delta x_{\mathrm B})^2\kappa^2\rho=\tfrac{\mathfrak h^2}{(E\ln2)^2}\cdot\tfrac{2\pi E^4(\ln2)^3}{\mathfrak h^2}=2\pi E^2\ln2\ [\text{machine}];
\]

%dsp-level:0
\[
\qquad 2\pi\ln2>2\cdot\tfrac{223}{71}\cdot\tfrac{405}{584}>1\ [\text{machine: exact rationals}]\implies\ell_R<\Delta x_{\mathrm B}
\]

%dsp-level:0
\[
(6):\ \text{the coefficient of}\ \Lambda\ \text{in}\ R\ \text{is}\ 4\ [\text{machine: monomial census}];
\]

%dsp-level:0
\[
\qquad \Lambda\in\mathrm{P4}\ \text{core}\ [\mathrm{P4}\ \text{re-cited}];
\]

%dsp-level:0
\[
\qquad R-4\Lambda=\kappa^2\rho\ \text{at}\ p=0\ [\text{machine}].\ \blacksquare
\]

\textbf{Phys.} Trace the field equations and the sign of a dust curvature is not a matter of convention: $R=\kappa^2(\rho-3p)+4\Lambda$, positive for a positive density at vanishing $\Lambda$, and the opposite sign is refuted as a polynomial identity rather than argued away. The component round trip closes — Ricci, scalar, Einstein, back to $\kappa^2T$ — in exact rational arithmetic on all four diagonal entries. Feeding in the width-derived density gives the curvature scale as $2\pi E_{\mathrm{int}}^4(\ln2)^3/\mathfrak h^2$: a pure function of an integer count, dimensionally exactly a curvature. Its radius is always shorter than the Born width, by a factor that grows with the count, so the more information a composite seals, the tighter its own curvature sits inside its spread. One term is not locked, and it is named: the cosmological constant enters the scalar additively with coefficient four, a free internal parameter, and the locked part is what remains after subtracting it.

\end{theorem}

\begin{theorem}[{\texorpdfstring{BM.16 — The calibration lock}{BM.16 - The calibration lock}}]\label{tag:BM-16}
\textbf{Thm.}

%dsp-level:0
\[
(1)\quad\textbf{the lock identity:}
\]

%dsp-level:0
\[
\qquad \boxed{\ (\Delta x_{\mathrm B})^2\cdot\kappa^2\rho\ =\ 2\pi\,E_{\mathrm{int}}^2\ln2\ }
\]

%dsp-level:0
\[
\qquad [\text{machine: one Laurent monomial, coefficient}\ 2];
\]

%dsp-level:0
\[
\mathfrak h\text{-free}\ [\text{machine: the}\ \mathfrak h\text{-exponent is}\ 0];
\]

%dsp-level:0
\[
\qquad \Lambda\text{-free}\ [\text{machine}];
\]

%dsp-level:0
\[
\qquad \text{P4-free}\ [\mathrm{P4}\ \text{core}:\ g,g',h_0,c_W,\Lambda,\varkappa,\beta,\alpha_\Omega\ —\ \neg\text{one occurs}];
\]

%dsp-level:0
\[
\textbf{the three dimensionless reads:}
\]

%dsp-level:0
\[
\qquad \frac{\Delta x_{\mathrm B}}{\mathfrak h}=\frac1{E_{\mathrm{int}}\ln2};
\]

%dsp-level:0
\[
\qquad \kappa^2\rho\,\mathfrak h^2=2\pi E_{\mathrm{int}}^4(\ln2)^3;
\]

%dsp-level:0
\[
\qquad \rho\,\mathfrak h^4=(E_{\mathrm{int}}\ln2)^4
\]

%dsp-level:0
\[
\qquad [\text{machine: all}\ \mathfrak h\text{-free}];
\]

%dsp-level:0
\[
(2)\quad\textbf{the rescaling covariance:}
\]

%dsp-level:0
\[
\qquad \mathfrak h\mapsto\lambda\mathfrak h:
\]

%dsp-level:0
\[
\qquad \Delta x_{\mathrm B}\mapsto\lambda\,\Delta x_{\mathrm B},\quad V_{\mathrm{eff}}\mapsto\lambda^3V_{\mathrm{eff}},\quad m_{\mathrm{hid}}\mapsto\lambda^{-1}m_{\mathrm{hid}},
\]

%dsp-level:0
\[
\rho\mapsto\lambda^{-4}\rho,\qquad\kappa^2\rho\mapsto\lambda^{-2}\kappa^2\rho,\qquad G_N\mapsto\lambda^2G_N,\qquad\kappa^2\mapsto\lambda^2\kappa^2\qquad[\text{machine: exact weight census}];
\]

%dsp-level:0
\[
\text{every dimensionless product of}\ (1)\ \text{is invariant}\ [\text{machine}]:
\]

%dsp-level:0
\[
\qquad \text{rescaling is a symmetry}\ \exists!\ \text{calibration d.o.f.}\ [\mathrm{CT2.11}]\ \text{and the lock rides above it};
\]

%dsp-level:0
\[
(3)\quad\textbf{no continuous mass variable:}
\]

%dsp-level:0
\[
\qquad \text{every read of this tier is a single Laurent monomial in}\ (\mathfrak h,E_{\mathrm{int}},\ln2,\pi)\ [\text{machine: monomial census}]:
\]

%dsp-level:0
\[
\neg\exists\ \text{a sum of scales,}\ \neg\exists\ \text{a tunable mixture};
\]

%dsp-level:0
\[
\qquad m_{\mathrm{hid}}\mathfrak h=E_{\mathrm{int}}\ln2\ \text{is injective in the integer}\ E_{\mathrm{int}}\ [\text{machine}];
\]

%dsp-level:0
\[
\textbf{quantization:}
\]

%dsp-level:0
\[
\qquad m_{\mathrm{hid}}\in\Bigl\{\frac{n\ln2}{\mathfrak h}:n\in\mathbb Z_{\ge2}\Bigr\},
\]

%dsp-level:0
\[
\qquad \Delta x_{\mathrm B}\in\Bigl\{\frac{\mathfrak h}{n\ln2}\Bigr\},
\]

%dsp-level:0
\[
\qquad \kappa^2\rho\in\Bigl\{\frac{2\pi n^4(\ln2)^3}{\mathfrak h^2}\Bigr\}:
\]

%dsp-level:0
\[
\qquad \text{three discrete ladders, strictly monotone}\ [\text{machine}];
\]

%dsp-level:0
\[
(4)\quad\textbf{the conjugate comparison scales}\ [\mathrm{P3}:\ G_N\ \text{is a named combination for external comparison},\ \neg\text{an input}]:
\]

%dsp-level:0
\[
\qquad r_s:=2G_Nm_{\mathrm{hid}}=\frac{\mathfrak h\,E_{\mathrm{int}}}2
\]

%dsp-level:0
\[
\qquad [\text{machine: exact}];
\]

%dsp-level:0
\[
\boxed{\ \Delta x_{\mathrm B}\cdot r_s=2G_N=\frac{\mathfrak h^2}{2\ln2}\ }
\]

%dsp-level:0
\[
\qquad \text{independent of}\ E_{\mathrm{int}}\ [\text{machine: the}\ E\text{-exponent is}\ 0]:
\]

%dsp-level:0
\[
\qquad 2\ln2\cdot(\Delta x_{\mathrm B}r_s)=\mathfrak h^2;
\]

%dsp-level:0
\[
\text{the two length reads are conjugate with a fixed product}: \ E_{\mathrm{int}}\uparrow\implies\Delta x_{\mathrm B}\downarrow,
\]

%dsp-level:0
\[
\ r_s\uparrow\ \text{along one hyperbola}\ [\text{machine}];
\]

%dsp-level:0
\[
(5)\quad\textbf{the crossing gap:}
\]

%dsp-level:0
\[
\qquad \ell_G:=\sqrt{G_N}=\frac{\mathfrak h}{2\sqrt{\ln2}}\ [\mathrm{CT2.9}]:
\]

%dsp-level:0
\[
\qquad \Bigl(\frac{\Delta x_{\mathrm B}}{\ell_G}\Bigr)^{\!2}=\frac{\Delta x_{\mathrm B}^2}{G_N}=\frac4{E_{\mathrm{int}}^2\ln2}
\]

%dsp-level:0
\[
\qquad [\text{machine: squares only, no radical in any assertion}];
\]

%dsp-level:0
\[
\Delta x_{\mathrm B}<\ell_G\iff E_{\mathrm{int}}^2\ln2>4\iff E_{\mathrm{int}}\ge3\qquad[\text{machine: exact enclosures, every}\ E\in1..60];
\]

%dsp-level:0
\[
\frac{r_s}{\Delta x_{\mathrm B}}=\frac{E_{\mathrm{int}}^2\ln2}2
\]

%dsp-level:0
\[
\qquad [\text{machine}]:
\]

%dsp-level:0
\[
\qquad r_s=\Delta x_{\mathrm B}\iff E_{\mathrm{int}}^2\ln2=2\iff E_{\mathrm{int}}=\sqrt{\tfrac2{\ln2}}\in(1,2)\ [\text{machine:}\ 1<\tfrac2{\ln2}<4];
\]

%dsp-level:0
\[
\textbf{the crossing lies strictly below the minimal floor}\ E_{\mathrm{int}}=2\ [\mathrm{FT.28}]:
\]

%dsp-level:0
\[
\qquad \neg\exists\,E_{\mathrm{int}}\in\mathbb Z\ \text{at the crossing};
\]

%dsp-level:0
\[
\qquad r_s>\Delta x_{\mathrm B}\ \ \forall E_{\mathrm{int}}\ge2\ [\text{machine}];
\]

%dsp-level:0
\[
(6)\quad\textbf{the not-locked ledger:}
\]

%dsp-level:0
\[
\qquad \Lambda\ [\mathrm{P4}:\ \text{one additive term}\ 4\Lambda\ \text{in}\ R,\ \mathrm{BM.15}(6)];
\]

%dsp-level:0
\[
\qquad \alpha\ [\mathrm{Def\text{-}BM.4}(4):\ \text{priced},\ \neg\text{adopted}];
\]

%dsp-level:0
\[
\qquad \mathfrak h\ [\mathrm{P1},\ \mathrm{CT2.11}:\ \text{the single free scale}\ —\ \text{absolute values ride on it}];
\]

%dsp-level:0
\[
\textbf{the comparison discipline:}
\]

%dsp-level:0
\[
\quad r_s,\ \ell_G\ \text{are built from}\ G_N,\ \text{a named combination of}\ \mathfrak h\ \text{for external comparison only}\ [\mathrm{P3}]:
\]

%dsp-level:0
\[
\qquad \neg\text{a collapse reading,}\ \neg\text{a horizon reading,}\ \neg\text{an identification is made}\ [\mathrm{P5};\ \mathrm{Def\text{-}BM.3}(6)];
\]

%dsp-level:0
\[
\text{the inequalities of}\ (5)\ \text{are order relations between derived lengths, and no external datum enters them}.
\]

\textbf{Rmk (lock scope).} The lock of clause (1) is a statement about dimensionless products: it says the width and the curvature cannot be varied independently of the integer count. It does not say the corpus predicts a length in metres — the single free scale $\mathfrak h$ fixes the unit [P1, P5], and every external comparison stays quarantined [Census FT.114, FT.120].

\textbf{Req.} D2.9 (Parts 1–2); CT2.1, CT2.9, CT2.11 (Part 2); P1, P3, P4, P5 (Parts 1–2); Def-BM.3(6), Def-BM.4 (this part); FT.28, FT.44(4) (Part 3, Census); Census FT.114, FT.120 (quarantine rows); BM.13(1), BM.14(2), BM.14(3), BM.15(4), BM.15(6) (this part); exact Laurent-monomial arithmetic, exact rational enclosures of $\ln2$ and Archimedes' bounds for $\pi$, square-comparison inequalities [math; machine: \texttt{checks/\allowbreak{}verify\_\allowbreak{}H\_\allowbreak{}calibration\_\allowbreak{}lock.py} §D].

\textbf{Pf.}

%dsp-level:0
\[
(1):\ (\Delta x_{\mathrm B})^2\kappa^2\rho=\tfrac{\mathfrak h^2}{(E\ln2)^2}\cdot\tfrac{2\pi E^4(\ln2)^3}{\mathfrak h^2}=2\pi E^2\ln2\ [\text{machine}];
\]

%dsp-level:0
\[
\qquad \text{the}\ \mathfrak h\text{-exponent census}: 2-2=0\ [\text{machine}];
\]

%dsp-level:0
\[
\qquad \text{the P4 core members occur in no monomial}\ [\text{machine scan of the tier's expressions}]
\]

%dsp-level:0
\[
(2):\ \text{each read is}\ c\,
\]

%dsp-level:0
\[
\mathfrak h^{k}\ \text{with}\ k\in\{1,3,-1,-4,-2,2,2\}\ [\text{machine: weight census}]\implies\mathfrak h\mapsto\lambda\mathfrak h\ \text{acts by}\ \lambda^k\ [\mathrm{CT2.11}\ \text{re-cited}];
\]

%dsp-level:0
\[
\qquad \text{the products of}\ (1)\ \text{have}\ k=0\implies\text{invariant}
\]

%dsp-level:0
\[
(3):\ \text{monomial census}\ [\text{machine}];
\]

%dsp-level:0
\[
\qquad n_1\ln2=n_2\ln2\iff n_1=n_2\ [\text{machine on}\ n\in1..39];
\]

%dsp-level:0
\[
\qquad \text{monotonicity of all three ladders}\ [\text{machine}]
\]

%dsp-level:0
\[
(4):\ r_s=2\cdot\tfrac{\mathfrak h^2}{4\ln2}\cdot\tfrac{E\ln2}{\mathfrak h}=\tfrac{\mathfrak h E}2\ [\text{machine}];
\]

%dsp-level:0
\[
\qquad \Delta x_{\mathrm B}r_s=\tfrac{\mathfrak h}{E\ln2}\cdot\tfrac{\mathfrak h E}2=\tfrac{\mathfrak h^2}{2\ln2}=2G_N\ [\text{machine}];
\]

%dsp-level:0
\[
\qquad \text{the}\ E\text{-exponent is}\ -1+1=0
\]

%dsp-level:0
\[
(5):\ \Delta x_{\mathrm B}^2/G_N=\tfrac{\mathfrak h^2}{(E\ln2)^2}\cdot\tfrac{4\ln2}{\mathfrak h^2}=\tfrac4{E^2\ln2}\ [\text{machine}];
\]

%dsp-level:0
\[
\qquad \tfrac4{E^2\ln2}<1
\]

%dsp-level:0
\[
\iff E^2\ln2>4:\ E
\]

%dsp-level:0
\[
= 2\to4\ln2<4\ (\ln2<1),
\]

%dsp-level:0
\[
\ E
\]

%dsp-level:0
\[
= 3\to9\ln2>4\ (\ln2>\tfrac49)\ [\text{machine: exact enclosures}];
\]

%dsp-level:0
\[
r_s/\Delta x_{\mathrm B}=\tfrac{\mathfrak h E}2\cdot\tfrac{E\ln2}{\mathfrak h}=\tfrac{E^2\ln2}2\ [\text{machine}];
\]

%dsp-level:0
\[
\qquad =1\iff E^2=\tfrac2{\ln2}\in(1,4)\ [\text{machine}]\implies E\in(1,2);
\]

%dsp-level:0
\[
\qquad E\ge2\implies\tfrac{E^2\ln2}2\ge2\ln2>1\ [\ln2>\tfrac12;\ \text{machine}]
\]

%dsp-level:0
\[
(6):\ \mathrm{BM.15}(6),\ \mathrm{Def\text{-}BM.4}(4),\ \mathrm{P1}/\mathrm{CT2.11},\ \mathrm{P3}/\mathrm{P5}\ \text{re-cited}.\ \blacksquare
\]

\textbf{Phys.} Two numbers in this tier carry neither the free scale nor a single internal parameter: the product of the squared Born width with the curvature scale, which is $2\pi$ times the squared edge count times $\ln2$, and the product of the Born width with the comparison radius, which is exactly twice Newton's constant and does not know the edge count at all. That is the lock — width and curvature are conjugate along a hyperbola whose constant is pure calibration, so no sealed composite can slide its size without sliding its field by the reciprocal amount. Rescaling the one free scale moves every dimensionful read by its exact weight and leaves every one of these products standing, which is what makes the lock a statement about the theory rather than about units. The ladder has a hole in it too: the point where the two lengths would coincide sits at an edge count between one and two, strictly below the smallest chain the ledger admits, so no admissible composite is ever at its own comparison radius. What is not locked is listed: the cosmological term, the geometric factor, and the hop itself.

\end{theorem}

\pseudohead{theorem}{\texorpdfstring{BM.17}{BM.17}}{\texorpdfstring{The neutron row of the identification layer}{The neutron row of the identification layer}}\label{tag:BM-17}
\textbf{Thm.}

%dsp-level:0
\[
(1)\quad\textbf{the structure:}
\]

%dsp-level:0
\[
\qquad \Sigma_{12}
\]

%dsp-level:0
\[
= \{p_1,q_1\}\cup\{p_2,q_2\}\ \text{with materialized 2-cycles joined by a bidirectional connector, single observer contact}\ [\text{Census FT.35}];
\]

%dsp-level:0
\[
E_{\mathrm{int}}=6\ (c_1{:}\,2,\ c_2{:}\,2,\ \text{connector}{:}\,2),
\]

%dsp-level:0
\[
\qquad I_{\mathrm{hid}}=6\ln2,
\]

%dsp-level:0
\[
\qquad m_{\mathrm{hid}}=\frac{6\ln2}{\mathfrak h}
\]

%dsp-level:0
\[
\qquad [\text{Census FT.35}(2);\ \text{cited row-wise}];
\]

%dsp-level:0
\[
\mathcal C(\Sigma_{12};C)=1\iff\text{all three attachments seal inside as one}\ \varepsilon\text{-type (baryon-like) triple}
\]

%dsp-level:0
\[
\qquad [\text{Census FT.35}(5);\ \mathrm{D3.9};\ \mathrm{CT3.17};\ \mathrm{FT.34}(5)];
\]

%dsp-level:0
\[
\Theta_{\mathrm{int}}=2k(\Gamma)\in2\mathbb Z,
\]

%dsp-level:0
\[
\qquad \Pi=72\,k(\Gamma)\,(\ln2)^2
\]

%dsp-level:0
\[
\qquad [\text{Census FT.35}(3)];
\]

%dsp-level:0
\[
\qquad \text{the junction vertices are}\ D^-\text{-configured}\ [\text{Census FT.35}(4);\ \mathrm{CT4.1}];
\]

%dsp-level:0
\[
(2)\quad\textbf{the identification:}
\]

%dsp-level:0
\[
\qquad \mathrm{Def\text{-}BM.3}(2)\ \text{labels the}\ \varepsilon\text{-sealed 3-juncture in the neutral charge class}\ q_{\mathrm{net}}
\]

%dsp-level:0
\[
= 0\ \text{as the}\ \textbf{neutron};
\]

%dsp-level:0
\[
\text{the theorem content of every clause below is the tuple}\ \bigl(\varepsilon\text{-junction},\ \Theta_{\mathrm{int}}\in2\mathbb Z,\ q_{\mathrm{net}}=0,\ E_{\mathrm{int}}=6,\ w=0\bigr);
\]

%dsp-level:0
\[
\qquad \text{the name is the layer's label}\ [\mathrm{Def\text{-}BM.3}(1),(7)];
\]

%dsp-level:0
\[
\text{the charge class is}\ \mathrm{Def\text{-}BM.3}(4)\text{'s declared label, pending BM.6(1)'s additive read}\;
\]

%dsp-level:0
\[
\qquad \text{deficit}\ 1'\ \text{is untouched}\ [\mathrm{FT.125}(4)];
\]

%dsp-level:0
\[
(3)\quad\textbf{the derived row at}\ E_{\mathrm{int}}^{(n)}=6\ \text{(all exact):}
\]

%dsp-level:0
\[
\Delta x_{\mathrm B}^{(n)}=\frac{\mathfrak h}{6\ln2}\ (<\mathfrak h\ \text{by the exact factor}\ \tfrac1{6\ln2});
\]

%dsp-level:0
\[
\qquad V_{\mathrm{eff}}^{(n)}=\frac{\mathfrak h^3}{216(\ln2)^3};
\]

%dsp-level:0
\[
\qquad m_{\mathrm{hid}}^{(n)}=\frac{6\ln2}{\mathfrak h};
\]

%dsp-level:0
\[
\rho^{(n)}=\frac{1296(\ln2)^4}{\mathfrak h^4};
\]

%dsp-level:0
\[
\qquad \kappa^2\rho^{(n)}=\frac{2592\,\pi(\ln2)^3}{\mathfrak h^2}
\]

%dsp-level:0
\[
\qquad [2\cdot6^4=2592;\ \text{machine}];
\]

%dsp-level:0
\[
\qquad \ell_R^{(n)}=\frac{\mathfrak h}{6^2\sqrt{2\pi(\ln2)^3}}<\Delta x_{\mathrm B}^{(n)};
\]

%dsp-level:0
\[
r_s^{(n)}=3\,\mathfrak h
\]

%dsp-level:0
\[
\qquad [\text{machine}];
\]

%dsp-level:0
\[
\qquad \Delta x_{\mathrm B}^{(n)}\cdot r_s^{(n)}=\frac{\mathfrak h^2}{2\ln2}=2G_N
\]

%dsp-level:0
\[
\qquad [\text{machine: the}\ E\text{-independent product}];
\]

%dsp-level:0
\[
(\Delta x_{\mathrm B}^{(n)})^2\,\kappa^2\rho^{(n)}=72\,\pi\ln2
\]

%dsp-level:0
\[
\qquad [\text{machine}]:
\]

%dsp-level:0
\[
\qquad \text{the integer}\ 72
\]

%dsp-level:0
\[
= 2\cdot6^2\ \text{is Census FT.35}(3)\text{'s}\ \Pi\text{-coefficient}\ —\ \text{recorded as a cross-check},
\]

%dsp-level:0
\[
\ \neg\text{a derivation of equality}\ [\mathrm{FT.112}\ \mathrm{Rmk}\ \text{pattern}];
\]

%dsp-level:0
\[
\Delta x_{\mathrm B}^{(n)}<\ell_G\iff36\ln2>4\iff\ln2>\tfrac19
\]

%dsp-level:0
\[
\qquad [\text{machine: exact enclosure}];
\]

%dsp-level:0
\[
\qquad w^{(n)}=0,\quad T^{(n)}_{\mu\nu}=\mathrm{diag}(\rho^{(n)},0,0,0)
\]

%dsp-level:0
\[
\qquad [\mathrm{CT4.5};\ \mathrm{BM.9}(4)];
\]

%dsp-level:0
\[
(4)\quad\textbf{the family table (the five computed cells):}
\]

%tbl-level:2
{\footnotesize\setlength{\tabcolsep}{4pt}
\begin{longtable}{>{\raggedright\arraybackslash}p{\dimexpr0.2770\textwidth-2\tabcolsep\relax}|>{\raggedright\arraybackslash}p{\dimexpr0.0529\textwidth-2\tabcolsep\relax}|>{\raggedright\arraybackslash}p{\dimexpr0.1010\textwidth-2\tabcolsep\relax}|>{\raggedright\arraybackslash}p{\dimexpr0.1411\textwidth-2\tabcolsep\relax}|>{\raggedright\arraybackslash}p{\dimexpr0.1330\textwidth-2\tabcolsep\relax}|>{\raggedright\arraybackslash}p{\dimexpr0.1491\textwidth-2\tabcolsep\relax}|>{\raggedright\arraybackslash}p{\dimexpr0.0690\textwidth-2\tabcolsep\relax}|>{\raggedright\arraybackslash}p{\dimexpr0.0770\textwidth-2\tabcolsep\relax}}
\toprule
cell & $E_{\mathrm{int}}$ & $\Delta x_{\mathrm B}/\mathfrak h$ & $V_{\mathrm{eff}}/\mathfrak h^3$ & $\rho\,\allowbreak \mathfrak h^4$ & $\kappa^2\rho\,\allowbreak \mathfrak h^2$ & $r_s/\mathfrak h$ & $\Delta x_{\mathrm B}$ vs $\ell_G$ \\
\addlinespace[2pt]
\midrule
\endfirsthead
\toprule
cell & $E_{\mathrm{int}}$ & $\Delta x_{\mathrm B}/\mathfrak h$ & $V_{\mathrm{eff}}/\mathfrak h^3$ & $\rho\,\allowbreak \mathfrak h^4$ & $\kappa^2\rho\,\allowbreak \mathfrak h^2$ & $r_s/\mathfrak h$ & $\Delta x_{\mathrm B}$ vs $\ell_G$ \\
\addlinespace[2pt]
\midrule
\endhead
\bottomrule
\endlastfoot
$\Sigma_{\min}$ (the digon) [FT.28] & 2 & $\tfrac1{2\ln2}$ & $\tfrac1{8(\ln2)^3}$ & $16(\ln2)^4$ & $32\pi(\ln2)^3$ & 1 & $\Delta x_{\mathrm B}>\allowbreak \ell_G$ \\
$\Sigma_{12}$ (the $\varepsilon$-type junction) [Census FT.35] & 6 & $\tfrac1{6\ln2}$ & $\tfrac1{216(\ln2)^3}$ & $1296(\ln2)^4$ & $2592\pi(\ln2)^3$ & 3 & $\Delta x_{\mathrm B}<\allowbreak \ell_G$ \\
tetrahedron [Census FT.42(5)] & 12 & $\tfrac1{12\ln2}$ & $\tfrac1{1728(\ln2)^3}$ & $20736(\ln2)^4$ & $41472\pi(\ln2)^3$ & 6 & $\Delta x_{\mathrm B}<\allowbreak \ell_G$ \\
cube / octahedron [Census FT.42(5)] & 24 & $\tfrac1{24\ln2}$ & $\tfrac1{13824(\ln2)^3}$ & $331776(\ln2)^4$ & $663552\pi(\ln2)^3$ & 12 & $\Delta x_{\mathrm B}<\allowbreak \ell_G$ \\
dodecahedron / icosahedron [Census FT.42(5)] & 60 & $\tfrac1{60\ln2}$ & $\tfrac1{216000(\ln2)^3}$ & $12960000(\ln2)^4$ & $25920000\pi(\ln2)^3$ & 30 & $\Delta x_{\mathrm B}<\allowbreak \ell_G$ \\
\end{longtable}
}

%dsp-level:0
\[
\text{the width column is strictly decreasing and the curvature column strictly increasing, with disjoint exact enclosures}\ [\text{machine}];
\]

%dsp-level:0
\[
\qquad r_s/\mathfrak h=E_{\mathrm{int}}/2\ \text{across the family}\ [\text{machine}];
\]

%dsp-level:0
\[
\Sigma_{\min}\ \text{is the unique computed cell with}\ \Delta x_{\mathrm B}>\ell_G\ [\mathrm{BM.16}(5);\ \text{machine}];
\]

%dsp-level:0
\[
\qquad \text{the duality-degenerate pairs share}\ E_{\mathrm{int}}\ \text{and hence every read of this tier}\ [\mathrm{FT.44}(1);\ \text{machine}];
\]

%dsp-level:0
\[
(5)\quad\textbf{not a point mass:}
\]

%dsp-level:0
\[
\qquad V_{\mathrm{eff}}^{(n)}>0,\quad\rho^{(n)}<\infty,\quad T^{(n)}_{\mu\nu}\ \text{is a finite-width distribution}\ [\mathrm{BM.9}(6);\ \mathrm{BM.10}(6)];
\]

%dsp-level:0
\[
\text{the width is sub-hop}\ [(3)]\implies\text{it is a chart envelope},
\]

%dsp-level:0
\[
\ \neg\text{a materialized distance}\ [\mathrm{Def\text{-}BM.4}(5);\ \mathrm{D1.6}];
\]

%dsp-level:0
\[
(6)\quad\textbf{the openness clause:}
\]

%dsp-level:0
\[
\qquad E_{\mathrm{int}}^{(n)}
\]

%dsp-level:0
\[
= 6\ \text{is Census-cell data at the computed}\ \varepsilon\text{-type cell}\ [\text{Census FT.35}(2);\ \text{row-wise consumption},\ \mathrm{BM.5}(2)\ \text{pattern}];
\]

%dsp-level:0
\[
\text{a general composite additivity law for}\ E_{\mathrm{int}}\ \text{beyond the computed cells is}\ \neg\text{constructed}\ [\mathrm{Def\text{-}BM.1}\ \text{read clause};\ \mathrm{BM.6}(1),(5)]:
\]

%dsp-level:0
\[
\text{every read of this tier is exact in}\ E_{\mathrm{int}}^{(n)}\ \text{and parametric wherever no cell supplies it}\ —\ \neg\text{a numeric value is invented};
\]

%dsp-level:0
\[
\neg\exists\ \text{a measured value in any clause}\ [\mathrm{P5};\ \mathrm{Def\text{-}BM.3}(6)].
\]

\textbf{Rmk (identification discipline).} Clause (3) is a row of derived exact expressions attached to a Census cell and labelled by Def-BM.3. It derives no neutron property that the tuple does not already carry, imports no neutron datum, and closes no row beyond its own: the additive charge read [BM.6(1)] and the assignment rule [FT.125(4)] stand as they are.

\textbf{Req.} D1.6, D2.9, D3.9 (Parts 1–2); CT3.17, CT4.1 (Part 2); P5 (Parts 1–2); Def-BM.1, Def-BM.3, Def-BM.4, BM.5(2), BM.6(1), BM.6(5) (this part); FT.28, FT.34(5), FT.44(1), FT.125(4) (Part 3); Census FT.35(2), FT.35(3), FT.35(4), FT.35(5), FT.42(5), FT.44(4) (row-wise); BM.9(4), BM.9(6), BM.10(6), BM.13, BM.14, BM.15, BM.16 (this part); exact rational arithmetic, exact $\ln2$ enclosures [math; machine: \texttt{checks/\allowbreak{}verify\_\allowbreak{}H\_\allowbreak{}calibration\_\allowbreak{}lock.py} §E].

\textbf{Pf.}

%dsp-level:0
\[
(1):\ \text{Census FT.35's rows re-cited verbatim}\ (2),
\]

%dsp-level:0
\[
(3),
\]

%dsp-level:0
\[
(4),
\]

%dsp-level:0
\[
(5)\ —\ \text{row-wise consumption},
\]

%dsp-level:0
\[
\ \neg\text{premise import}\ [\mathrm{BM.5}(2);\ \text{A §0.1}]
\]

%dsp-level:0
\[
(2):\ \mathrm{Def\text{-}BM.3}(2)\ \text{re-cited};
\]

%dsp-level:0
\[
\qquad \text{the tuple is the theorem content}\ [\mathrm{Def\text{-}BM.3}(1),(7)];
\]

%dsp-level:0
\[
\qquad \text{the conditional charge class}\ [\mathrm{Def\text{-}BM.3}(4);\ \mathrm{BM.6}(1)]
\]

%dsp-level:0
\[
(3):\ \text{substitution of}\ E_{\mathrm{int}}
\]

%dsp-level:0
\[
= 6\ \text{into BM.13(1), BM.14(1),(2),(3), BM.15(4),(5), BM.16(4)}\ [\text{machine: every coefficient exact}];
\]

%dsp-level:0
\[
\qquad 2\cdot6^4=2592,\ 6^4=1296,\ 6^3=216,\ 2\cdot6^2=72\ [\text{machine}];
\]

%dsp-level:0
\[
\qquad 36\ln2>4\iff\ln2>\tfrac19\ [\text{machine: enclosure}]
\]

%dsp-level:0
\[
(4):\ \text{the same substitutions at}\ E_{\mathrm{int}}\in\{2,6,12,24,60\}\ [\text{Census FT.42}(5)\ \text{card table};\ \mathrm{FT.28}];
\]

%dsp-level:0
\[
\qquad \text{monotonicity and disjointness}\ [\text{machine}];
\]

%dsp-level:0
\[
\qquad r_s/\mathfrak h=E/2\ [\text{machine}];
\]

%dsp-level:0
\[
\qquad \text{duality degeneracy}: E_{\mathrm{int}}=2E\ \text{shared by dual pairs}\ [\mathrm{FT.44}(1)\ \text{re-cited}]
\]

%dsp-level:0
\[
(5):\ \mathrm{BM.9}(6),\ \mathrm{BM.10}(6),\ \mathrm{Def\text{-}BM.4}(5)\ \text{re-cited}
\]

%dsp-level:0
\[
(6):\ \text{openness rows re-cited unmodified};
\]

%dsp-level:0
\[
\qquad \mathrm{P5}\ \text{scan of this tag's text}\ [\text{machine: zero external numeric tokens}].\ \blacksquare
\]

\textbf{Phys.} The corpus has one computed $\varepsilon$-type junction cell, and it carries six internal directed edges: two per cycle plus the two of the connector. Everything this tier derived then evaluates on it in exact rationals — a width of one hop over six $\ln2$, a volume of one hop cubed over $216(\ln2)^3$, a density of $1296(\ln2)^4$ per hop to the fourth, a curvature scale of $2592\pi(\ln2)^3$ per hop squared, and a comparison radius of exactly three hops. The lock reads $72\pi\ln2$ on this cell, and $72$ is the same integer the Census found in its $\Pi$ row — recorded as a coincidence of counts, not promoted to an identity. Def-BM.3's label attaches to that tuple, and the label changes nothing in it. The neutron of this corpus is not a point: it is a finite density over a strictly positive, sub-hop volume, with a curvature radius shorter than its own spread, sitting on a ladder whose rungs are integers and whose crossing point no integer reaches.

\medskip\noindent\textit{\texorpdfstring{BM.18 — The status of the calibration lock: adjudication and scope}{BM.18 - The status of the calibration lock: adjudication and scope}.}\ Structural-constraints record: see Appendix B, Section B.1.16 of the Addendum.\par\medskip

\phantomsection
\section[Tier BM.6 — The Topology-Blind Neutron and the…]{\texorpdfstring{Tier BM.6 — The Topology-Blind Neutron and the Topology-Blind + Static Proton: the Flicker Ensemble, the Path-Integrated Mass, the Binding Read, and the Forced Carrier}{Tier BM.6 - The Topology-Blind Neutron and the Topology-Blind + Static Proton: the Flicker Ensemble, the Path-Integrated Mass, the Binding Read, and the Forced Carrier}}

\begin{definition}[{\texorpdfstring{Def-BM.5 — The flicker ensemble of a sealed composite}{Def-BM.5 - The flicker ensemble of a sealed composite}}]\label{tag:Def-BM-5}
\textbf{Def.} A construction class on Def-H.32's topology-blind path integral (no new axiom; Def-H.32's scopes (S1)–(S4) are inherited at every use).

%dsp-level:0
\[
(1)\quad\textbf{the ensemble:}\qquad\mathcal E(\mathcal X)\ :=\ \bigl\{R\subseteq E_{\mathrm{dir}}(\Gamma):\ \text{the boundary data of}\ \mathcal X\ \text{is preserved on}\ R\bigr\}
\]

%dsp-level:0
\[
\text{the boundary data}\ :
\]

%dsp-level:0
\[
= \ \bigl(\text{charge class}\ [\mathrm{Def\text{-}BM.3}(2)\text{–}(3)],\ \ \Theta_{\mathrm{int}}\ \text{class}\ [\mathrm{BM.2}(1)],\ \ S\ \text{winding class}\ [\mathrm{Def\text{-}H.32}(4):\ \text{the}\ (q,S)\ \text{labels}]\bigr);
\]

%dsp-level:0
\[
\textbf{constitutive clause:}\ \ \text{the ensemble is the sector-conditioned configuration set of}\ \mathrm{Def\text{-}H.32}(1)\text{–}(2)\ \text{restricted to a composite's boundary data};
\]

%dsp-level:0
\[
\qquad \text{the sum below is finite and exact at every budget}\ B\ [\mathrm{Def\text{-}H.32}(1):\ \mathcal G_B=\{\Gamma:\ V(\Gamma)\le B\},\ \mathfrak C_B(\Gamma)=\{R:\ |R|\le B_K\}];
\]

%dsp-level:0
\[
(2)\quad\textbf{the path integral:}
\]

%dsp-level:0
\[
\qquad Z(\mathcal X)\ =\ \sum_{R\in\mathcal E(\mathcal X)}\ \sum_{\sigma}\ w(R,\sigma),
\]

%dsp-level:0
\[
\qquad w(C)=\mathcal J(C)\prod_{e\in E(C)}G(e)\cdot(-1)^{\aleph(\sigma,C)}
\]

%dsp-level:0
\[
\qquad [\mathrm{Def\text{-}H.32}(2);\ \mathrm{A8}];
\]

%dsp-level:0
\[
\text{the two layers (norm and amplitude) are kept separate}\ [\mathrm{CT2.8};\ \mathrm{Def\text{-}H.32}(2)]\ \text{and no layer is substituted for the other};
\]

%dsp-level:0
\[
(3)\quad\textbf{the weights are derived, not declared:}\qquad\text{the per-pair Janus master generator}\ Q=\bigl(\begin{smallmatrix}-x_+&x_+\\ x_-&-x_-\end{smallmatrix}\bigr)\ \text{over}\ \mathbb Z[x_+,x_-]\ \text{has a one-dimensional kernel spanned by}\ (x_-,x_+)\ [\mathrm{FT.108}(4)]:
\]

%dsp-level:0
\[
p\ =\ \frac{x_+}{x_++x_-}\qquad[\text{derived}];\qquad x_+=1-4\varepsilon_q,\quad x_-=4\varepsilon_q\implies p=1-4\varepsilon_q\ \ \text{identically}\qquad[\mathrm{A4},\ \mathrm{D0.4};\ \text{machine}];
\]

%dsp-level:0
\[
\textbf{no P-class datum is added:}\ \ \varepsilon_q\ \text{is the A4 crystallization threshold}\ [\mathrm{D0.4}:\ \text{admissible range}\ [0,\tfrac14)\ [\mathrm{CT0.8}];\ \text{sharp choice}\ \varepsilon_q=0];
\]

%dsp-level:0
\[
(4)\quad\textbf{the information content is distributed:}
\]

%dsp-level:0
\[
\qquad E_{\mathrm{int}}(R)
\]

%dsp-level:0
\[
= |R|
\]

%dsp-level:0
\[
= n_{\mathcal R}\ [\mathrm{Def\text{-}H.32}(1):\ \text{the budget is the ledger's own content measure};\ \mathrm{CT1.29}:\ S_{\mathfrak L}=n_{\mathcal R}\ln2];
\]

%dsp-level:0
\[
\text{the crystallized value}\ E_{\mathrm{int}}
\]

%dsp-level:0
\[
= 6\ [\text{Census FT.35}(2)]\ \text{is one point of the distribution},
\]

%dsp-level:0
\[
\ \neg\text{the ensemble};
\]

%dsp-level:0
\[
\qquad \text{the move system conserves the grade}\ a_+-a_-=|R'|-|R|\ [\mathrm{FT.108}(2)];
\]

%dsp-level:0
\[
(5)\quad\textbf{the no-flicker limit (scope):}
\]

%dsp-level:0
\[
\qquad \varepsilon_q
\]

%dsp-level:0
\[
= 0
\]

%dsp-level:0
\[
\implies x_-
\]

%dsp-level:0
\[
= 0\ \text{exactly}\implies\neg\exists\,
\]

%dsp-level:0
\[
\text{dissolution move}\implies\mathcal E(\mathcal X)\ \text{collapses to its fully crystallized member}\ [\text{machine}];
\]

%dsp-level:0
\[
\text{every}\ \varepsilon_q>0\ \text{gives}\ x_->0\ \text{and a strictly positive flicker weight};
\]

%dsp-level:0
\[
\qquad \text{the}\ B\to\infty\ \text{limit is}\ \neg\text{constructed}\ [\mathrm{Def\text{-}H.32}(1);\ \text{growth data: FT.105}].
\]

\textbf{Rmk (ensemble discipline).} The ensemble is a sum over crystallization patterns at fixed boundary data, not a sum over species: every member carries the composite's charge class, sealing-phase class and winding class. The topology labels are internal bookkeeping [Def-H.32(4) pattern], and no member of any ensemble is identified with an external object [P5; Def-BM.3(6)].

\textbf{Req.} A4, A8 (Part 1); D0.4, D2.9 (Parts 1–2); CT0.8, CT1.29 (Part 1); CT2.8 (Part 2); Def-H.32, FT.105, FT.108(2), FT.108(4), FT.108(5) (Part 3); Def-BM.3, BM.2(1) (this part); Census FT.35(2) (row-wise); exact polynomial kernel arithmetic over $\mathbb Q[x_+,x_-]$, exact rational specialization [math; machine: \texttt{checks/\allowbreak{}verify\_\allowbreak{}H\_\allowbreak{}topo\_\allowbreak{}blind\_\allowbreak{}neutron.py} §A].

\textbf{Pf.}

%dsp-level:0
\[
\text{Definition block: the constitutive clause is}\ (1)\text{'s sector-conditioned configuration set};
\]

%dsp-level:0
\[
\ \text{the inherited clauses are Def-H.32's}\ (1)\text{–}(2)\ \text{and scopes}\ (\mathrm{S1})\text{–}(\mathrm{S4}),
\]

%dsp-level:0
\[
\ \text{already declared};
\]

%dsp-level:0
\[
\ \text{clause}\ (3)\ \text{is derived content recorded at the definition site};
\]

%dsp-level:0
\[
\ (5)\ \text{is a scope clause cited at use}.\ \blacksquare
\]

\textbf{Phys.} A sealed composite is not one topology with a number attached; it is a finite exact sum over crystallization patterns that share its boundary data — the same topology-blind path integral the fermion era built, now conditioned on a composite's charge class, sealing phase and winding class. Nothing is postulated to weight the sum: the Janus move generator has a one-dimensional kernel, its span is the stationary measure, and substituting the crystallization threshold gives the flicker weight as $4\varepsilon_q$ with no new free parameter anywhere. Set that threshold to its sharp value and the ensemble collapses to the single crystallized configuration — the flicker is a consequence of the threshold being nonzero, not an extra ingredient.

\end{definition}

\begin{definition}[{\texorpdfstring{Def-BM.6 — The family naming rows (extension of the identification layer)}{Def-BM.6 - The family naming rows (extension of the identification layer)}}]\label{tag:Def-BM-6}
\textbf{Def.} An extension of Def-BM.3's identification layer to the three derived channel triplets, recorded by a logged declaration as Def-BM.3(7) requires. A name attaches to a \textit{line of a derived spectrum}, ordered by that spectrum; no name attaches to a mass value.

%dsp-level:0
\[
(1)\quad\textbf{the}\ \ell\text{-channel triplet}\ [\mathrm{FT.109}\text{–}\mathrm{FT.113}:\ r_i=1+\sqrt2\cos(\delta_\ell+\tfrac{2\pi i}3),\ \delta_\ell=\tfrac29]:
\]

%dsp-level:0
\[
\qquad \textbf{electron}\ :=\ \text{the}\ i=0\ \text{line},\quad\textbf{muon}\ :=\ \text{the}\ i=1\ \text{line},\quad\textbf{tau}\ :=\ \text{the}\ i=2\ \text{line};
\]

%dsp-level:0
\[
\text{the order rule is the spectrum's own:}\ r_0>r_1>r_2\ \text{at}\ \delta_\ell=\tfrac29\ [\mathrm{FT.113};\ \text{the span is}\ 3477\text{-fold}];
\]

%dsp-level:0
\[
(2)\quad\textbf{the}\ \nu\text{-channel triplet}\ [\mathrm{FT.142}\text{–}\mathrm{FT.144}:\ \delta_\nu=\tfrac18\ \text{on the derived}\ 16\text{-arena}]:
\]

%dsp-level:0
\[
\qquad \textbf{electron neutrino},\ \textbf{muon neutrino},\ \textbf{tau neutrino}\ :=\ \text{its three lines in the same spectral order};
\]

%dsp-level:0
\[
(3)\quad\textbf{the}\ q\text{-channel triplet}\ [\mathrm{FT.118},\ \mathrm{FT.119}:\ \delta_q=\tfrac2{27},\ \text{triplet-attached}]:
\]

%dsp-level:0
\[
\qquad \textbf{the quark triplet}\ :=\ \text{its three lines in the same order};
\]

%dsp-level:0
\[
\qquad \text{the cross-channel identity}\ 16\delta_\nu
\]

%dsp-level:0
\[
= 9\delta_\ell
\]

%dsp-level:0
\[
= 27\delta_q
\]

%dsp-level:0
\[
= 2\ [\mathrm{FT.143}]\ \text{is the derived relation between the three rows};
\]

%dsp-level:0
\[
(4)\quad\textbf{the derived content the names attach to:}
\]

%dsp-level:0
\[
\qquad \sum_ir_i=3=\dim F_{J^\pm}\ [\mathrm{FT.111}(3)];
\]

%dsp-level:0
\[
\qquad \sum_ir_i^2=6;
\]

%dsp-level:0
\[
\qquad Q:=\frac{\sum_ir_i^2}{(\sum_ir_i)^2}=\frac69=\frac23\ \ \textbf{exactly, for every phase}\ [\text{machine};\ \mathrm{FT.120}(6)];
\]

%dsp-level:0
\[
\text{the cosine form is forced by}\ C_3\text{-covariance and the normalization by the seal dimension}\ [\mathrm{FT.111}(2)\text{–}(3);\ \mathrm{FT.158}(4)];
\]

%dsp-level:0
\[
\qquad \text{the angle is a ratio of spatial counts}\ \delta_\ell=(3-1)/(3\cdot3)\ [\mathrm{FT.158}(2)\text{–}(3)];
\]

%dsp-level:0
\[
(5)\quad\textbf{the non-derived clause (the mass form):}
\]

%dsp-level:0
\[
\qquad \text{the bridge from a spectral line to a mass value}\ (\sqrt{m}\propto r\ \text{vs}\ m\propto r\ \text{vs}\ \neg\text{a bridge})\ \text{is item-13 content and is}\ \neg\text{adopted here}\;
\]

%dsp-level:0
\[
\text{the standing bar}\ \text{amplitude}\ne\text{mass}\ [\mathrm{FT.79}(5)]\ \text{is untouched};
\]

%dsp-level:0
\[
\qquad \text{the channel scale}\ \mu\ \text{remains an undeclared P4-class parameter}\ [\text{Cal-F}]:\ \neg\text{a mass value is derived for any named line};
\]

%dsp-level:0
\[
(6)\quad\textbf{the data bar:}
\]

%dsp-level:0
\[
\qquad \neg\exists\ \text{a measured external value in any clause}\ [\mathrm{P5}];
\]

%dsp-level:0
\[
\qquad \text{the ratios}\ (Q,\ \text{the span, the angle identities})\ \text{are corpus-derived and are}\ \neg\text{compared to measurement here};
\]

%dsp-level:0
\[
(7)\quad\textbf{scope:}
\]

%dsp-level:0
\[
\qquad \text{the names occur in Part 4 only}\ [\mathrm{Def\text{-}BM.3}(7)\ \text{extended}];
\]

%dsp-level:0
\[
\qquad \text{further extension requires a further logged declaration}.
\]

\textbf{Rmk (what the naming does and does not do).} Naming the $\ell$-channel lines \textit{electron, muon, tau} records that the corpus derives a three-line $C_3$-covariant spectrum whose Koide ratio is exactly $2/3$ at every phase — the ratio the family is famous for — and that the lines are strictly ordered with a 3477-fold span. It does not derive their masses, it does not adopt a mass form, and it does not compare any ratio to a measured value. The naming is Def-BM.3's labelling act applied to spectral lines instead of junctions.

\textbf{Req.} FT.79(5), FT.109, FT.110, FT.111, FT.112, FT.113, FT.118, FT.119, FT.120(6), FT.142, FT.143, FT.144, FT.158 (Part 3); Def-BM.3 (this part); P4, P5 (Parts 1–2); Cal-F; exact identities in $\mathbb Q[c,s,r]/(c^2{+}s^2{-}1,\ r^2{-}3)$ [math; machine: \texttt{checks/\allowbreak{}verify\_\allowbreak{}H\_\allowbreak{}topo\_\allowbreak{}blind\_\allowbreak{}neutron.py} §F].

\textbf{Pf.}

%dsp-level:0
\[
\text{Definition block: the constitutive clause is the name}\leftrightarrow\text{line attachment by spectral order};
\]

%dsp-level:0
\[
\ \text{the non-derived clause is}\ (5)\ \text{(the mass form is item-13 content,}\ \neg\text{adopted)};
\]

%dsp-level:0
\[
\ \text{clause}\ (4)\text{'s identities are derived}\ [\text{machine:}\ \textstyle\sum_i\cos(\theta+\tfrac{2\pi i}3)=0,\ \sum_i\cos^2=\tfrac32\ \text{in the quotient ring}]\ \text{and are recorded here as the content the names attach to}.\ \blacksquare
\]

\textbf{Phys.} Three channels, three triplets, one mechanism: a $C_3$-covariant seal bilinear whose shape covariance forces, whose normalization the seal dimension fixes, and whose angle is a ratio of spatial counts. The ratio the charged-lepton family is known for falls out as $6/9$ — exactly two thirds, at every phase, with nothing fitted — and the lines are strictly ordered across a 3477-fold span. The names now attach to those lines in spectral order. What the names do not carry is a mass: the bridge from a spectral line to a mass value is the corpus's open item, the amplitude-is-not-mass bar stands, and the overall scale is still an undeclared internal parameter.

\end{definition}

\begin{theorem}[{\texorpdfstring{BM.19 — The flicker ensemble is derived: the stationary measure, the threshold specialization, and the no-flicker limit}{BM.19 - The flicker ensemble is derived: the stationary measure, the threshold specialization, and the no-flicker limit}}]\label{tag:BM-19}
\textbf{Thm.}

%dsp-level:0
\[
(1)\quad\textbf{the measure is forced:}\qquad Q=\begin{pmatrix}-x_+&x_+\\ x_-&-x_-\end{pmatrix}\ \text{over}\ \mathbb Z[x_+,x_-]:\qquad(x_-,\,x_+)\,Q=(0,0)\ \ \text{exactly}\ [\text{machine}];
\]

%dsp-level:0
\[
\neg\exists\ \text{a nonzero constant vector in}\ \ker Q\ [\text{machine: exhaustive on}\ [-3,3]^2]:
\]

%dsp-level:0
\[
\qquad \dim\ker Q=1\ \text{over the fraction field}\implies p=\frac{x_+}{x_++x_-}\ \text{is the unique stationary measure}\ [\mathrm{FT.108}(4)];
\]

%dsp-level:0
\[
(2)\quad\textbf{the specialization is an identity:}
\]

%dsp-level:0
\[
\qquad x_+=1-4\varepsilon_q,\quad x_-=4\varepsilon_q\implies p=\frac{1-4\varepsilon_q}{1-4\varepsilon_q+4\varepsilon_q}=1-4\varepsilon_q
\]

%dsp-level:0
\[
\qquad [\text{machine: exact at six rational thresholds}];
\]

%dsp-level:0
\[
\text{the flicker (dissolution) weight is}\ x_-=4\varepsilon_q:
\]

%dsp-level:0
\[
\qquad \text{admissible on}\ \varepsilon_q\in[0,\tfrac14)\ [\mathrm{CT0.8}]\implies x_-\in[0,1),\ x_+\in(0,1]\ [\text{machine}];
\]

%dsp-level:0
\[
\textbf{no P-class datum enters:}\ \ \varepsilon_q\ \text{is the A4 threshold}\ [\mathrm{D0.4}];
\]

%dsp-level:0
\[
\qquad \neg\text{a new fugacity,}\ \neg\text{a new rate,}\ \neg\text{a P4 addition}\ [\text{Cal-BM stays empty}];
\]

%dsp-level:0
\[
(3)\quad\textbf{the no-flicker limit:}
\]

%dsp-level:0
\[
\qquad \varepsilon_q
\]

%dsp-level:0
\[
= 0\ [\mathrm{D0.4}\ \text{sharp choice}]
\]

%dsp-level:0
\[
\implies x_-
\]

%dsp-level:0
\[
= 0\ \text{exactly}\implies\neg\exists\,
\]

%dsp-level:0
\[
\text{a dissolution move}\implies\mathcal E(\mathcal X)
\]

%dsp-level:0
\[
= \{\text{the crystallized configuration}\}\ [\mathrm{Def\text{-}BM.5}(5);\ \text{machine}];
\]

%dsp-level:0
\[
\varepsilon_q>0\implies x_->0:
\]

%dsp-level:0
\[
\qquad \text{the flicker is strictly present at every nonzero threshold and strictly absent at the sharp one};
\]

%dsp-level:0
\[
(4)\quad\textbf{the ensemble is finite and exact at each budget:}
\]

%dsp-level:0
\[
\qquad |\mathfrak C_B(\Gamma)|<\infty\ [\mathrm{Def\text{-}H.32}(1)];
\]

%dsp-level:0
\[
\qquad \text{the}\ B=4\ \text{sector census is exact integers}\ [\mathrm{FT.106}(3):\ Z^{\mathrm{norm}}(q,S)\ \text{and}\ Z^{\mathrm{amp}}(q,S)\ \text{per sector}];
\]

%dsp-level:0
\[
Z^{\mathrm{norm}}(q,S)
\]

%dsp-level:0
\[
= Z^{\mathrm{norm}}(-q,S)\ \text{unconditionally and}\ Z^{\mathrm{amp}}(q,S)
\]

%dsp-level:0
\[
= Z^{\mathrm{amp}}(-q,S)\ \text{on the cover support}\ [\mathrm{FT.106}(3\mathrm d)];
\]

%dsp-level:0
\[
\qquad Z^{\mathrm{amp}}(q,k)=0\ \text{for}\ k=2\ \text{and odd}\ k\ge3\ [\mathrm{FT.106}(3\mathrm c)];
\]

%dsp-level:0
\[
(5)\quad\textbf{the information content is a distribution:}
\]

%dsp-level:0
\[
\qquad E_{\mathrm{int}}(R)=|R|\ [\mathrm{Def\text{-}BM.5}(4)]:
\]

%dsp-level:0
\[
\qquad \text{the ensemble spreads}\ E_{\mathrm{int}}\ \text{over}\ \{0,\dots,B_K\}\ \text{and the crystallized value is its top point};
\]

%dsp-level:0
\[
a_+-a_-=|R'|-|R|,
\]

%dsp-level:0
\[
\qquad \textstyle\sum_dm_d-\sum_dn_d=V'-V,
\]

%dsp-level:0
\[
\qquad \sum_d d\,(m_d-n_d)=E'_{\mathrm{pot}}-E_{\mathrm{pot}}
\]

%dsp-level:0
\[
\qquad [\mathrm{FT.108}(2):\ \text{the grade map is a homomorphism of the move system}]:
\]

%dsp-level:0
\[
\text{every flicker move changes}\ E_{\mathrm{int}}\ \text{by exactly the}\ J^\pm\ \text{excess it carries}\ —\ \text{the distribution is the move algebra's orbit, not an ansatz};
\]

%dsp-level:0
\[
(6)\quad\textbf{the composite is not a fixed topology:}
\]

%dsp-level:0
\[
\qquad \mathcal E(\mathcal X)\ \text{has more than one member}\iff x_->0\iff\varepsilon_q>0\ [(3)]:
\]

%dsp-level:0
\[
\qquad \text{at every admissible nonzero threshold the neutron and the proton are ensembles}.
\]

\textbf{Req.} A4 (Part 1); D0.4 (Part 1); CT0.8, CT1.29 (Part 1); CT2.8 (Part 2); Def-H.32, FT.105, FT.106(3), FT.108(2), FT.108(4), FT.108(5) (Part 3); Def-BM.5 (this part); P4 (Parts 1–2); §Cal-BM; exact polynomial arithmetic over $\mathbb Q[x_+,x_-]$, exhaustive integer kernel census, exact rational specialization [math; machine: \texttt{checks/\allowbreak{}verify\_\allowbreak{}H\_\allowbreak{}topo\_\allowbreak{}blind\_\allowbreak{}neutron.py} §A].

\textbf{Pf.}

%dsp-level:0
\[
(1):\ (x_-)(-x_+)+(x_+)(x_-)=0\ \text{and}\ (x_-)(x_+)+(x_+)(-x_-)=0\ \text{in}\ \mathbb Q[x_+,x_-]\ [\text{machine}];
\]

%dsp-level:0
\[
\qquad \text{constant vectors}: a(-x_+)+bx_-=0\ \wedge\ ax_++b(-x_-)=0\implies a=b=0\ [\text{machine: exhaustive}]
\]

%dsp-level:0
\[
(2):\ \text{substitution}\ [\text{machine at}\ \varepsilon_q\in\{\tfrac1{64},\tfrac1{16},\tfrac18,\tfrac3{16},\tfrac15,\tfrac{63}{256}\}];
\]

%dsp-level:0
\[
\qquad \varepsilon_q\in[0,\tfrac14)\ [\mathrm{CT0.8}]\implies4\varepsilon_q\in[0,1)\ [\text{machine}]
\]

%dsp-level:0
\[
(3):\ x_-=4\cdot0=0\ [\text{machine}];
\]

%dsp-level:0
\[
\qquad \text{a move system with no}\ J^-\ \text{generator admits no dissolution}\ [\mathrm{D1.2}:\ J^-:n_{\mathcal R}:1\to0];
\]

%dsp-level:0
\[
\qquad \text{the ensemble reduces to the crystallized pattern}\ [\mathrm{Def\text{-}BM.5}(5)]
\]

%dsp-level:0
\[
(4):\ \mathrm{Def\text{-}H.32}(1)\text{'s filtration and}\ \mathrm{FT.106}(3)\text{'s census rows re-cited}\ (\text{exact integers, frozen});
\]

%dsp-level:0
\[
\qquad (5):\ \mathrm{FT.108}(2)\text{'s three grade identities re-cited}\ [\text{machine rows frozen at Tier F.36}]
\]

%dsp-level:0
\[
(6):\ |\mathcal E(\mathcal X)|>1\iff\exists\,\text{a dissolution move}\iff x_->0\ [(2),(3)].\ \blacksquare
\]

\textbf{Phys.} Nothing in the flicker is chosen. The Janus move generator has a kernel, that kernel is a line, and its span is the stationary measure — so the weight of a crystallization against a dissolution is the ratio $x_+/(x_++x_-)$ before anyone names a rate. Substituting the crystallization threshold the axioms already carry turns that ratio into $1-4\varepsilon_q$, an identity rather than a fit, and the flicker weight into $4\varepsilon_q$. At the sharp threshold the dissolution weight is exactly zero and the ensemble collapses to one configuration; at any nonzero threshold the composite is genuinely an ensemble, and the move algebra's grade identities say precisely how each flicker step moves the information content. That is the whole mechanism of "the neutron is not a fixed topology", with no new parameter anywhere in it.

\end{theorem}

\begin{theorem}[{\texorpdfstring{BM.20 — The mass from the curvature law on the path-integrated route}{BM.20 - The mass from the curvature law on the path-integrated route}}]\label{tag:BM-20}
\textbf{Thm.}

%dsp-level:0
\[
(1)\quad\textbf{the mass is the curvature flux:}\qquad R=\kappa^2(\rho-3p)+4\Lambda\ [\mathrm{BM.15}(2)];\qquad p=0,\ \Lambda=0:\quad R=\kappa^2\rho,\quad\rho V=m:
\]

%dsp-level:0
\[
\boxed{\ R\,V=\kappa^2 m\ }\qquad\text{i.e.}\qquad m=\frac1{\kappa^2}\int_\Sigma R\,dV\qquad[\text{machine: exact monomial identity}];
\]

%dsp-level:0
\[
(2)\quad\textbf{the}\ \Lambda\ \textbf{clause:}
\]

%dsp-level:0
\[
\qquad (R-4\Lambda)V=\kappa^2m\ \ \text{identically}\ [\text{machine}]:
\]

%dsp-level:0
\[
\qquad \Lambda\ \text{is the one additive freedom}\ [\mathrm{P4};\ \mathrm{BM.15}(6)];
\]

%dsp-level:0
\[
(3)\quad\textbf{the ensemble form (the path-integrated route):}
\]

%dsp-level:0
\[
\qquad \langle m\rangle_{\mathcal X}=\frac1{\kappa^2}\Bigl\langle\int_\Sigma R\,dV\Bigr\rangle_{\mathcal X}=\sum_{R\in\mathcal E(\mathcal X)}w(R)\,m(R),
\]

%dsp-level:0
\[
\qquad \sum_Rw(R)=1
\]

%dsp-level:0
\[
\qquad [\text{machine: the normalized average equals the averaged flux}];
\]

%dsp-level:0
\[
\text{the mass of a sealed composite is the ensemble average of the curvature flux over its flicker}\ —\ \neg\text{the flux of the crystallized member};
\]

%dsp-level:0
\[
(4)\quad\textbf{the derived sector read:}
\]

%dsp-level:0
\[
\qquad m^2_B(q,S)=\frac{\sum_{\Gamma,R}\nu\,t^{|R|}u^L\#\mathrm{inst}\cdot\tfrac32(2E-|R|)}{\sum_{\Gamma,R}\nu\,t^{|R|}u^L\#\mathrm{inst}}\bigg|_{t=u=1}
\]

%dsp-level:0
\[
\qquad [\mathrm{FT.107}(3);\ \Delta S=\tfrac32n_{\mathcal V}\ [\mathrm{FT.107}(2)]]:
\]

%dsp-level:0
\[
m^2_4(\cdot,0)=\frac{843957}{98956},
\]

%dsp-level:0
\[
\quad m^2_4(\cdot,1)=\frac{1921662}{393679},
\]

%dsp-level:0
\[
\quad m^2_4(\cdot,2)=\frac{91}{23},
\]

%dsp-level:0
\[
\quad m^2_4(\cdot,3)=\frac{13713}{3088}
\]

%dsp-level:0
\[
\qquad [\text{exact rationals};\ \text{machine}];
\]

%dsp-level:0
\[
\text{charge-degenerate in}\ q\ [\mathrm{FT.106}(3\mathrm a),(3\mathrm d)];
\]

%dsp-level:0
\[
\qquad \text{all strictly positive};
\]

%dsp-level:0
\[
\qquad \text{the transit sector}\ (S=0)\ \text{is heaviest}\ [\text{machine}];
\]

%dsp-level:0
\[
(5)\quad\textbf{the bound:}
\]

%dsp-level:0
\[
\qquad 0<m^2_B(q,S)\le\tfrac32\max_\Gamma n_{\mathcal V}\le\tfrac32B(B-1)=18\ \ (B=4);
\]

%dsp-level:0
\[
\qquad \text{the census maximum is}<9
\]

%dsp-level:0
\[
\qquad [\mathrm{FT.107}(4);\ \text{machine}];
\]

%dsp-level:0
\[
(6)\quad\textbf{the curvature scale of the ensemble read:}
\]

%dsp-level:0
\[
\qquad \rho_B=\frac{m_B}{V_{\mathrm{eff}}},
\]

%dsp-level:0
\[
\qquad \kappa^2\rho_B=\frac{2\pi\,\mathfrak h\,m_B}{V_{\mathrm{eff}}}\ \Big/\ \frac{\mathfrak h}{1}=\frac{2\pi\,E^4_{\mathrm{eff}}(\ln2)^3}{\mathfrak h^2}\bigg|_{E_{\mathrm{eff}}:=m_B\mathfrak h/\ln2}
\]

%dsp-level:0
\[
\text{i.e. the sector read enters BM.14's and BM.15's formulas through the effective count}\ E_{\mathrm{eff}}:
\]

%dsp-level:0
\[
= m_B\mathfrak h/\ln2\ [\text{exact rational per sector}]:
\]

%dsp-level:0
\[
\qquad \text{the lock of BM.16(1) applies sector by sector}.
\]

\textbf{Rmk (route discipline).} Clause (1) is the curvature law read as a mass definition — the same law BM.15 derived from CT1.23 — applied configuration by configuration and then averaged, which is what "the path-integrated route" means here. It does not import a mass datum: $m(R)$ at each member is the informational read [FT.44(4)], and the ensemble average is an exact rational combination of integers times $\ln2/\mathfrak h$; the sector read of clause (4) is the deficit-average route [FT.107(3)], and the two coincide where both are defined.

\textbf{Req.} D2.9 (Parts 1–2); CT1.23, CT2.1, CT2.9 (Part 2); P4 (Parts 1–2); Def-H.32, FT.106(3), FT.107(2), FT.107(3), FT.107(4) (Part 3); FT.44(4) (Census); Def-BM.4, Def-BM.5, BM.14(2), BM.14(3), BM.15(2), BM.15(6), BM.16(1), BM.19 (this part); exact Laurent-monomial identities, exact rational ensemble averages [math; machine: \texttt{checks/\allowbreak{}verify\_\allowbreak{}H\_\allowbreak{}topo\_\allowbreak{}blind\_\allowbreak{}neutron.py} §B].

\textbf{Pf.}

%dsp-level:0
\[
(1):\ \mathrm{BM.15}(2)\ \text{at}\ p=0,\ \Lambda=0:\ R=\kappa^2\rho;
\]

%dsp-level:0
\[
\qquad \rho
\]

%dsp-level:0
\[
= m/V\ [\mathrm{BM.14}(2)\ \text{pattern}]
\]

%dsp-level:0
\[
\implies RV
\]

%dsp-level:0
\[
= \kappa^2m\ [\text{machine: substitution of the two defining relations into}\ RV-\kappa^2m\ \text{gives}\ 0\ \text{identically}]
\]

%dsp-level:0
\[
(2):\ R=\kappa^2\rho+4\Lambda\implies(R-4\Lambda)V=\kappa^2\rho V=\kappa^2m\ [\text{machine: the residual polynomial is}\ 0]
\]

%dsp-level:0
\[
(3):\ \text{linearity of the flux in}\ m\ \text{at fixed}\ V,\kappa^2\ [(1)]\implies\langle RV\rangle/\kappa^2=\langle m\rangle;
\]

%dsp-level:0
\[
\qquad \text{normalization}\ \sum w=1\ [\text{machine witness}]
\]

%dsp-level:0
\[
(4)\text{–}(5):\ \mathrm{FT.107}(2)\text{–}(4)\ \text{rows re-cited verbatim}\ (\text{exact rationals, frozen});
\]

%dsp-level:0
\[
\qquad \text{positivity, the ordering and the bound re-verified on the stated fractions}\ [\text{machine}]
\]

%dsp-level:0
\[
(6):\ E_{\mathrm{eff}}:=m_B\mathfrak h/\ln2\ \text{is the count whose informational mass is}\ m_B\ [\mathrm{FT.44}(4)\ \text{inverted}];
\]

%dsp-level:0
\[
\qquad \mathrm{BM.14}(3)/\mathrm{BM.15}(4)\ \text{re-cited at}\ E=E_{\mathrm{eff}}.\ \blacksquare
\]

\textbf{Phys.} Mass is curvature flux. Trace the field equations for a pressureless composite and the scalar curvature times the volume is $\kappa^2$ times the mass, with the cosmological term the only additive ambiguity — so the mass read the corpus derived from the dust stress tensor is the same object as the integral of its own curvature. Run that identity through the flicker and the mass of a composite becomes the ensemble average of the curvature flux, not the flux of its crystallized snapshot. The topology-blind census supplies the average independently, as an exact rational per sector: four fractions at budget four, all positive, all charge-degenerate, with the transit sector heaviest and the whole table under the quadratic budget bound. Both routes meet in one formula, and the calibration lock of the previous tier applies sector by sector through the effective count.

\end{theorem}

\begin{theorem}[{\texorpdfstring{BM.21 — The binding energy: additive in the count, non-additive in the ensemble}{BM.21 - The binding energy: additive in the count, non-additive in the ensemble}}]\label{tag:BM-21}
\textbf{Thm.}

%dsp-level:0
\[
(1)\quad\textbf{the informational count is exactly additive:}
\]

%dsp-level:0
\[
\qquad E_{\mathrm{int}}(\Sigma_{12})=6=(c_1{:}\,2)+(c_2{:}\,2)+(\text{connector}{:}\,2)
\]

%dsp-level:0
\[
\qquad [\text{Census FT.35}(2);\ \text{machine}];
\]

%dsp-level:0
\[
m_{\mathrm{hid}}
\]

%dsp-level:0
\[
= E_{\mathrm{int}}\ln2/\mathfrak h\ [\mathrm{FT.44}(4)]\ \text{is linear in the count}\implies\textbf{the informational binding read vanishes exactly:}
\]

%dsp-level:0
\[
\quad E^{\mathrm{info}}_{\mathrm{bind}}=\bigl(6-(2{+}2{+}2)\bigr)\tfrac{\ln2}{\mathfrak h}=0
\]

%dsp-level:0
\[
\qquad [\text{machine}];
\]

%dsp-level:0
\[
(2)\quad\textbf{the ensemble read is not additive:}
\]

%dsp-level:0
\[
\qquad \mathcal B(S_j;S_1,S_2,S_3)\ :=\ m^2_B(\cdot,S_j)-\sum_{i=1}^{3}m^2_B(\cdot,S_i)
\]

%dsp-level:0
\[
\qquad [\mathrm{BM.20}(4)]:
\]

%dsp-level:0
\[
\mathcal B(1;1,1,1)=-2\,m^2_4(\cdot,1)=-\frac{3843324}{393679};
\]

%dsp-level:0
\[
\qquad \mathcal B(2;1,1,1)=\frac{91}{23}-\frac{5764986}{393679};
\]

%dsp-level:0
\[
\qquad \mathcal B(3;1,1,1)=\frac{13713}{3088}-\frac{5764986}{393679};
\]

%dsp-level:0
\[
\mathcal B(0;1,1,1)=\frac{843957}{98956}-\frac{5764986}{393679};
\]

%dsp-level:0
\[
\qquad \textbf{every one is a nonzero exact rational}\ [\text{machine}];
\]

%dsp-level:0
\[
(3)\quad\textbf{the binding is the flicker's, not the count's:}
\]

%dsp-level:0
\[
\qquad \varepsilon_q
\]

%dsp-level:0
\[
= 0\implies\mathcal E\ \text{collapses}\ [\mathrm{BM.19}(3)]\implies\text{the ensemble average reduces to the crystallized member}\implies\mathcal B\to0\ \text{with}\ (1);
\]

%dsp-level:0
\[
\varepsilon_q>0\implies\mathcal B\ne0\ \text{exactly}\ [(2)]:
\]

%dsp-level:0
\[
\qquad \textbf{a non-flickering composite has no binding read, and a flickering one has a nonzero exact one};
\]

%dsp-level:0
\[
(4)\quad\textbf{the sign structure:}
\]

%dsp-level:0
\[
\qquad \text{all four combinations are bound:}\quad\mathcal B(S_j;1,1,1)<0\ \ \forall S_j\in\{0,1,2,3\}
\]

%dsp-level:0
\[
\qquad [\text{machine: exact rational sign census}];
\]

%dsp-level:0
\[
\text{the exact order of the sector masses is}\ m^2_4(\cdot,0)>m^2_4(\cdot,1)>m^2_4(\cdot,3)>m^2_4(\cdot,2)\ [\text{machine}]:
\]

%dsp-level:0
\[
\qquad \text{three separated constituents outweigh one junction in every sector combination};
\]

%dsp-level:0
\[
\text{the like-sector case is exactly}\ -2\,m^2_4(\cdot,1)\ [\text{machine}];
\]

%dsp-level:0
\[
\qquad |\mathcal B(S_j)|
\]

%dsp-level:0
\[
= 3m^2_4(\cdot,1)-m^2_4(\cdot,S_j)\ \text{is largest for the lightest junction sector}\ (S_j=2)\ \text{and smallest for the heaviest}\ (S_j=0)\ [\text{machine}];
\]

%dsp-level:0
\[
(5)\quad\textbf{the ceiling inheritance:}
\]

%dsp-level:0
\[
\qquad |\mathcal B|
\]

%dsp-level:0
\[
\le 4\cdot\tfrac32B(B-1)
\]

%dsp-level:0
\[
= 72\ \ \text{at}\ B
\]

%dsp-level:0
\[
= 4\ \text{from}\ \mathrm{FT.107}(4)\text{'s bound applied to all four terms}\ [\text{machine: exact bound arithmetic}];
\]

%dsp-level:0
\[
\text{the winding spread at}\ B=4\ \text{is}\ \tfrac{91}{23}\le m^2\le\tfrac{843957}{98956}\ [\mathrm{FT.107}(4)]:
\]

%dsp-level:0
\[
\qquad \text{the binding read stays order-one while the bound grows as}\ \Theta(B^2);
\]

%dsp-level:0
\[
(6)\quad\textbf{the read discipline:}
\]

%dsp-level:0
\[
\qquad \mathcal B\ \text{is a difference of the corpus's own}\ m^2_B\ \text{reads}\ [\mathrm{FT.107}(3)];
\]

%dsp-level:0
\[
\qquad \neg\text{a mass value in MeV,}\ \neg\text{an external datum},\ \neg\text{a comparison}\ [\mathrm{P5}];
\]

%dsp-level:0
\[
\text{the pole-mass side}\ (\text{defect angles, ceilings})\ \text{is a different read and remains item-13 bridge content}\ [\mathrm{FT.76}(3);\ \mathrm{Def\text{-}H.39}].
\]

\textbf{Req.} D2.9 (Parts 1–2); CT2.1 (Part 2); P5 (Parts 1–2); Def-H.32, Def-H.39, FT.76(3), FT.107(2), FT.107(3), FT.107(4) (Part 3); FT.44(4) (Census); Census FT.35(2) (row-wise); Def-BM.5, BM.19(3), BM.20(1), BM.20(4) (this part); exact rational arithmetic [math; machine: \texttt{checks/\allowbreak{}verify\_\allowbreak{}H\_\allowbreak{}topo\_\allowbreak{}blind\_\allowbreak{}neutron.py} §C].

\textbf{Pf.}

%dsp-level:0
\[
(1):\ \text{Census FT.35}(2)\text{'s decomposition row re-cited}\ (2{+}2{+}2=6);
\]

%dsp-level:0
\[
\qquad m_{\mathrm{hid}}\ \text{linear in}\ E_{\mathrm{int}}\ [\mathrm{FT.44}(4)]\implies\text{the difference is}\ 0\cdot\tfrac{\ln2}{\mathfrak h}\ [\text{machine}]
\]

%dsp-level:0
\[
(2):\ \text{exact rational subtraction of}\ \mathrm{FT.107}(3)\text{'s four frozen values}\ [\text{machine}];
\]

%dsp-level:0
\[
\qquad \text{nonvanishing: each numerator is nonzero after clearing denominators}\ [\text{machine}]
\]

%dsp-level:0
\[
(3):\ \mathrm{BM.19}(3)\ \text{re-cited};
\]

%dsp-level:0
\[
\qquad \text{a one-member ensemble has}\ \langle m\rangle=m\ \text{and the count-additivity of}\ (1)\ \text{applies}
\]

%dsp-level:0
\[
(4):\ \mathrm{FT.107}(3)\text{'s four values in exact descending order}\ \tfrac{843957}{98956}>\tfrac{1921662}{393679}>\tfrac{13713}{3088}>\tfrac{91}{23},
\]

%dsp-level:0
\[
\ \text{i.e.}\ S
\]

%dsp-level:0
\[
= 0>1>3>2\ [\text{machine: exact rational comparison}];
\]

%dsp-level:0
\[
\qquad \text{the sign census and the like-sector identity}\ [\text{machine}]
\]

%dsp-level:0
\[
(5):\ 4\cdot\tfrac32\cdot4\cdot3=72\ [\text{machine}];\qquad\mathrm{FT.107}(4)\text{'s spread row re-cited}
\]

%dsp-level:0
\[
(6):\ \mathrm{P5}\ \text{scan of this tag's text}\ [\text{machine: zero external numeric tokens}].\ \blacksquare
\]

\textbf{Phys.} Adding the edges yields six; adding the masses yields nothing. The informational count is exactly additive at the crystallized cell, so the binding read built from it vanishes identically — which is why a binding energy cannot be an integer summation. The flicker ensemble gives a different answer: averaged over the sector censuses, three separated constituents and one junction do not add up, and the mismatch is a nonzero exact rational in every combination of sectors. Turn the threshold to its sharp value and the ensemble collapses, the mismatch collapses with it, and the binding is exactly zero again — so the binding is a property of the flicker itself. Its magnitude stays order-one at budget four against a bound growing quadratically, its sign follows the transit sector's weight, and it remains a read of the corpus's own $m^2_B$ values with no external number anywhere near it.

\end{theorem}

\begin{theorem}[{\texorpdfstring{BM.22 — The static force matrix on the two charge classes}{BM.22 - The static force matrix on the two charge classes}}]\label{tag:BM-22}
\textbf{Thm.} With $h(\mathbf c)=\tfrac12(c_0I_2+\mathbf c\cdot\boldsymbol\sigma)$ and $\hat\kappa$ the ray swap [FT.129(2)], under Def-H.43's declared read promotion:

%dsp-level:0
\[
(1)\quad\textbf{the promoted force matrix:}
\]

%dsp-level:0
\[
\qquad \mathbf F_{\mathrm{ext}}(\mathbf c,q)=h(\mathbf c)\otimes I_2+q\,I_2\otimes\hat\kappa,
\]

%dsp-level:0
\[
\qquad \bigl[h\otimes I_2,\ q\,I_2\otimes\hat\kappa\bigr]=0
\]

%dsp-level:0
\[
\qquad [\mathrm{FT.129}(5);\ \text{machine: all profiles}\times\text{ternary}\ q];
\]

%dsp-level:0
\[
(2)\quad\textbf{the block read:}
\]

%dsp-level:0
\[
\qquad F_{LL}=F_{RR}=h(\mathbf c),
\]

%dsp-level:0
\[
\qquad F_{LR}=F_{RL}=q\,I_2
\]

%dsp-level:0
\[
\qquad [\text{machine: exact block extraction on}\ 4\ \text{profiles}\times3\ \text{charge values}];
\]

%dsp-level:0
\[
\text{the flipping block is the core read itself}:\ \langle L|\mathbf F_{\mathrm{ext}}|R\rangle=q\,I_2\ [\mathrm{FT.129}(5)];
\]

%dsp-level:0
\[
\qquad \text{the flipping set is}\ \{F_{\mathrm{EM}}\ \text{core read}\},\ \text{minimal and unique}\ [\mathrm{FT.129}(5)];
\]

%dsp-level:0
\[
(3)\quad\textbf{the spectrum:}
\]

%dsp-level:0
\[
\qquad \operatorname{spec}=\bigl\{\tfrac12(c_0+s|\mathbf c|)+tq:\ s,t\in\{\pm1\}\bigr\},
\]

%dsp-level:0
\[
\qquad \operatorname{charpoly}=\bigl((y-u)^2-q^2\bigr)\bigl((y+u)^2-q^2\bigr),\quad y=x-\tfrac{c_0}2,\quad u^2=\tfrac{|\mathbf c|^2}4
\]

%dsp-level:0
\[
[\mathrm{FT.130}(5);\ \text{machine: the monic quartic built from the four roots equals the Faddeev–LeVerrier charpoly, and equals the closed form, on every witness}];
\]

%dsp-level:0
\[
(4)\quad\textbf{the two classes:}\qquad\textbf{the chiral-partner splitting is}\ 2|q|\in\{0,2\}\qquad[\mathrm{FT.130}(5);\ \text{machine}]:
\]

%dsp-level:0
\[
\textbf{neutron class}\ (q_{\mathrm{net}}=0\ [\mathrm{Def\text{-}BM.3}(2)]):\quad F_{LR}=0,\quad\text{splitting}\ 0,\quad\text{parity doubling}\ [\mathrm{FT.130}(6)];
\]

%dsp-level:0
\[
\textbf{proton class}\ (q_{\mathrm{net}}=+1\ [\mathrm{Def\text{-}BM.3}(2)]):
\]

%dsp-level:0
\[
\quad F_{LR}=I_2,\quad\text{splitting}\ 2,\quad\text{the doublet is split by exactly two read units};
\]

%dsp-level:0
\[
\textbf{the EM difference between the two classes is exactly the flipping block}\ [(2),(4)]:\ \text{one carries it, the other is annihilated by it};
\]

%dsp-level:0
\[
(5)\quad\textbf{the derived locus and the promotion status:}
\]

%dsp-level:0
\[
\qquad \text{at the derived locus}\ \mathbf F_{\mathrm{ext}}(\mathbf c)=h\otimes I_\omega\ \text{and}\ F_{LR}=0\ \text{exactly}\ [\mathrm{FT.128}(2):\ \text{the no-flip theorem}]:
\]

%dsp-level:0
\[
\text{the flipping block exists only under Def-H.43's declared read promotion}\ [\text{the Def-H.43 declaration}];
\]

%dsp-level:0
\[
\qquad \textbf{status: conditional on that declaration, not on any new one};
\]

%dsp-level:0
\[
(6)\quad\textbf{the invariances:}
\]

%dsp-level:0
\[
\qquad \mathbf F_{\mathrm{ext}}\mapsto\mathbf F_{\mathrm{ext}}+\lambda I_4
\]

%dsp-level:0
\[
\implies F_{LR},
\]

%dsp-level:0
\[
\ \text{Schmidt ranks, commutants and}\ 2|q|\ \text{invariant; the spectrum translates by}\ \lambda
\]

%dsp-level:0
\[
\qquad [\mathrm{FT.130}(7)];
\]

%dsp-level:0
\[
\text{operator-Schmidt rank}=2\iff q\ne0\ \wedge\ \mathbf c\ne0;\quad1\ \text{when}\ h\propto I_2;\quad0\iff\mathbf c=0=q\ [\mathrm{FT.130}(2)];
\]

%dsp-level:0
\[
\qquad \neg\text{an entangling term at any locus}\ [\mathrm{FT.130}(2)];
\]

%dsp-level:0
\[
(7)\quad\textbf{the charge-label status:}
\]

%dsp-level:0
\[
\qquad q_{\mathrm{net}}\ \text{is}\ \mathrm{Def\text{-}BM.3}(4)\text{'s declared class label, pending the additive read}\ [\mathrm{BM.6}(1)]:
\]

%dsp-level:0
\[
\text{clauses}\ (1)\text{–}(4),(6)\ \text{are exact in}\ q\in\{-1,0,+1\}\ [\mathrm{EM.9}];
\]

%dsp-level:0
\[
\qquad \text{the attachment of}\ q=0\ \text{to the neutron class and}\ q=+1\ \text{to the proton class is the declaration}.
\]

\textbf{Req.} CT0.12, CT1.1, CT1.2 (Part 1); EM.9 (Part 2); Def-H.34, Def-H.43, FT.122(6), FT.123, FT.128(2), FT.129(2), FT.129(5), FT.130(2), FT.130(5), FT.130(6), FT.130(7) (Part 3); Def-BM.3 (this part); the Def-H.43 declaration; exact $\mathbb Q(i)$ block algebra, Faddeev–LeVerrier characteristic polynomials over $\mathbb Q$, exact quartic root reconstruction [math; machine: \texttt{checks/\allowbreak{}verify\_\allowbreak{}H\_\allowbreak{}topo\_\allowbreak{}blind\_\allowbreak{}neutron.py} §D].

\textbf{Pf.}

%dsp-level:0
\[
(1)\text{–}(2):\ \text{block extraction of}\ h\otimes I_2+qI_2\otimes\hat\kappa\ \text{with the chiral factor inner}: F_{ab}
\]

%dsp-level:0
\[
= h\,
\]

%dsp-level:0
\[
\delta\text{-block}+q\hat\kappa_{ab}
\]

%dsp-level:0
\[
\implies F_{LL}
\]

%dsp-level:0
\[
= F_{RR}
\]

%dsp-level:0
\[
= h\ (\hat\kappa_{00}=\hat\kappa_{11}=0),
\]

%dsp-level:0
\[
\ F_{LR}
\]

%dsp-level:0
\[
= F_{RL}
\]

%dsp-level:0
\[
= qI_2\ (\hat\kappa_{01}=\hat\kappa_{10}=1)\ [\text{machine}];
\]

%dsp-level:0
\[
\qquad \text{commutation}\ [\text{machine}]
\]

%dsp-level:0
\[
(3):\ \text{the four roots}\ \tfrac12(c_0+s|\mathbf c|)+tq\ \text{build a monic quartic over}\ \mathbb Q\ \text{at witnesses with}\ |\mathbf c|^2\ \text{a perfect square rational}\ (u\in\{\tfrac12,0,1,\tfrac37\});
\]

%dsp-level:0
\[
\ \text{Faddeev–LeVerrier on}\ \mathbf F_{\mathrm{ext}}\ \text{agrees, and the closed form}\ (y^2+u^2-q^2)^2-4u^2y^2\ \text{agrees with both}\ [\text{machine}]
\]

%dsp-level:0
\[
(4):\ 2|q|\ \text{at}\ q\in\{-1,0,+1\}\ [\text{machine}];\qquad\text{the class attachment}\ [\mathrm{Def\text{-}BM.3}(2)\ \text{re-cited}]
\]

%dsp-level:0
\[
(5)\text{–}(6):\ \mathrm{FT.128}(2),\ \mathrm{FT.130}(2),(7)\ \text{rows re-cited}\ (\text{frozen});\qquad(7):\ \mathrm{Def\text{-}BM.3}(4)\ \text{re-cited}.\ \blacksquare
\]

\textbf{Phys.} Put the static force matrix on the two charge classes and they come apart in exactly one place. Both classes share the same internal block $h$; the promoted read adds a chiral-flipping block equal to the core charge times the identity, so the neutral class has no flipping block at all and keeps its parity doubling, while the unit class splits its chiral partners by precisely two read units. That splitting — zero against two — is the whole electromagnetic difference between the neutron row and the proton row in this corpus, and it is a matrix identity, not a coupling constant. The promotion that allows a flipping block is a declared read, and the tier says so; everything downstream of it (the spectrum, the charpoly, the Schmidt ranks, the scalar-shift invariance) is exact. What is still declared rather than derived is which class label sits on which junction: the additive charge read remains an open construction.

\end{theorem}

\begin{theorem}[{\texorpdfstring{BM.23 — The transition: the ladder rule, the forced carrier, and the resolvent rate}{BM.23 - The transition: the ladder rule, the forced carrier, and the resolvent rate}}]\label{tag:BM-23}
\textbf{Thm.}

%dsp-level:0
\[
(1)\quad\textbf{the weak ladder:}
\]

%dsp-level:0
\[
\qquad E_\pm=\tfrac12(\sigma_1\pm i\sigma_2)=\text{the two matrix units in the}\ \sigma_3\ \text{eigenbasis};
\]

%dsp-level:0
\[
\qquad [E_+,E_-]=\sigma_3,
\]

%dsp-level:0
\[
\qquad [\sigma_3,E_\pm]=\pm2E_\pm
\]

%dsp-level:0
\[
\qquad [\mathrm{FT.41}(\mathrm d);\ \text{machine: exact over}\ \mathbb Q(i)];
\]

%dsp-level:0
\[
\text{the charged-current shape}\ W^\pm
\]

%dsp-level:0
\[
= \tfrac1{\sqrt2}(W_1\mp iW_2)\ [\mathrm{D3.10}]\ \text{is the same ladder on the pair register}\ [\mathrm{C8}];
\]

%dsp-level:0
\[
(2)\quad\textbf{the selection rule:}
\]

%dsp-level:0
\[
\qquad q_{\mathrm{core}}=\tfrac12(\sigma_{s}-\sigma_{t})\ [\mathrm{D3.4}]\ \text{and the ladder moves}\ \sigma_3\ \text{by}\ \pm2\ [(1)]\implies\boxed{\ \Delta q=\pm1\ }
\]

%dsp-level:0
\[
\qquad [\text{machine}];
\]

%dsp-level:0
\[
(3)\quad\textbf{the alphabet overflow:}
\]

%dsp-level:0
\[
\qquad \text{shifting the ternary read}\ \{-1,0,+1\}\ [\mathrm{EM.9}]\ \text{by}\ +1\ \text{gives}\ \{0,+1,+2\}\ \text{and by}\ -1\ \text{gives}\ \{-2,-1,0\}:
\]

%dsp-level:0
\[
\qquad \text{both leave the alphabet}\ [\text{machine}];
\]

%dsp-level:0
\[
\therefore\ \text{a charge-changing transition cannot stay inside the composite: it must emit a carrier};
\]

%dsp-level:0
\[
(4)\quad\textbf{the compensation is unique:}
\]

%dsp-level:0
\[
\qquad Q_{\mathrm{EM}}
\]

%dsp-level:0
\[
= \sum_{c\in U}q_{\mathrm{EM}}(c)
\]

%dsp-level:0
\[
= 0\ \text{in every valid universe}\ [\mathrm{EM.11};\ \mathrm{D2.11},\ \mathrm{D1.3}]\implies\text{a composite read changing by}\ \Delta q\ \text{requires a compensating read of exactly}\ -\Delta q:
\]

%dsp-level:0
\[
\#\{q\in\{-1,0,+1\}:\ q=-\Delta q\}=1\qquad[\text{machine}]:\qquad\textbf{the emitted carrier's charge is forced, not chosen};
\]

%dsp-level:0
\[
\text{the compensated total is exactly}\ 0:\ \text{EM.11's balance survives the transition}\ [\text{machine}];
\]

%dsp-level:0
\[
(5)\quad\textbf{the carrier selection theorem:}
\]

%dsp-level:0
\[
\qquad \text{the carrier must satisfy}\ (\mathrm i)\ q
\]

%dsp-level:0
\[
= -1\ \text{for}\ \Delta q
\]

%dsp-level:0
\[
= +1\ [(4)];
\]

%dsp-level:0
\[
\ (\mathrm{ii})\ \text{no colour attachment}\ [\text{else the composite's colour seal breaks}:\ \mathcal C=0\implies\text{the block vanishes}\ [\mathrm{FT.19};\ \mathrm{CT3.18},\ \mathrm{D2.14}]];
\]

%dsp-level:0
\[
\ (\mathrm{iii})\ \text{EM activated}\ (f_v\ne0)\ [\mathrm{D3.3};\ \text{else it cannot carry the read}];
\]

%dsp-level:0
\[
\text{in the channel inventory: the}\ q\text{-channel is triplet-attached}\ [\mathrm{FT.119}]\ \text{(fails ii)};
\]

%dsp-level:0
\[
\qquad \text{the}\ \nu\text{-channel is EM-unactivated}\ [\mathrm{FT.142}(1):\ f_v=0\ \text{reads}]\ \text{(fails iii)};
\]

%dsp-level:0
\[
\therefore\ \textbf{the}\ \ell\text{-channel is the unique carrier class}\ [\mathrm{FT.109}\text{–}\mathrm{FT.113}:\ \text{colour-blind, ternary read active}]:
\]

%dsp-level:0
\[
\qquad \text{the emitted line is Def-BM.6(1)'s electron row};
\]

%dsp-level:0
\[
(6)\quad\textbf{the neutral partner is admissible, not forced:}
\]

%dsp-level:0
\[
\qquad \neg\exists\ \text{a derived law requiring a neutral emission}\ [\text{no lepton-number read exists in the inventory}];
\]

%dsp-level:0
\[
\text{the}\ \nu\text{-channel}\ [\mathrm{FT.142}\text{–}\mathrm{FT.144}]\ \text{is EM-neutral and colour-blind, so it can carry the deficit read}\ \Delta S
\]

%dsp-level:0
\[
= \tfrac32n_{\mathcal V}\ [\mathrm{FT.107}(2)]\ \text{without disturbing the charge balance}:
\]

%dsp-level:0
\[
\textbf{status: admissible carrier for the deficit read; its emission is not a theorem of this corpus};
\]

%dsp-level:0
\[
(7)\quad\textbf{the rate:}
\]

%dsp-level:0
\[
\qquad \sum_n\langle\varnothing|\hat U^n|\varnothing\rangle w^n=\frac1{1-w^2x_+x_-}\ \ \text{in}\ \mathbb Q[w,x_+,x_-]/(w^{12})
\]

%dsp-level:0
\[
\qquad [\mathrm{FT.108}(5);\ \text{machine}];
\]

%dsp-level:0
\[
x_+x_-=4\varepsilon_q(1-4\varepsilon_q)\le\tfrac14,
\]

%dsp-level:0
\[
\qquad \text{equality at}\ \varepsilon_q=\tfrac18
\]

%dsp-level:0
\[
\qquad [\text{machine: exact maximum over the admissible range}]\implies\text{the pole}\ w^2x_+x_-=1\ \text{sits at}\ |w|\ge2;
\]

%dsp-level:0
\[
\text{the per-tick transition weight is an exact rational function of}\ \varepsilon_q\ \text{alone}\ [\mathrm{BM.19}(2)]:
\]

%dsp-level:0
\[
\qquad \Gamma\ \text{is}\ \neg\text{a declared number and}\ \neg\text{a new parameter};
\]

%dsp-level:0
\[
(8)\quad\textbf{the no-flicker theorem:}
\]

%dsp-level:0
\[
\qquad \varepsilon_q=0\implies x_+x_-=0\implies\text{the resolvent is}\ 1\implies\Gamma=0\ \text{exactly}
\]

%dsp-level:0
\[
\qquad [\mathrm{BM.19}(3);\ \text{machine}]:
\]

%dsp-level:0
\[
\textbf{a non-flickering composite never decays};
\]

%dsp-level:0
\[
\qquad \text{the flicker is the mechanism, and the threshold is its only control};
\]

%dsp-level:0
\[
(9)\quad\textbf{the arrow:}
\]

%dsp-level:0
\[
\qquad \text{the transition is forward-only}\ [\mathrm{FT.144}(2):\ \text{the domains force strict alternation;}\ \mathrm{A6}:\ \neg\exists\tau^{-1}]:
\]

%dsp-level:0
\[
\qquad \neg\exists\ \text{an inverse decay in the ledger};
\]

%dsp-level:0
\[
\text{the time coordinate is the ledger order}\ [x^0:=\tau_n,\ \mathrm{D1.8}]:
\]

%dsp-level:0
\[
\qquad \Gamma\ \text{is a weight per tick},\ \neg\text{a rate per external second}\ [\mathrm{P5}].
\]

\textbf{Rmk (provenance of each piece).} Clause (5) is a selection theorem: three conditions on the carrier, two channels eliminated by their own frozen rows, one class left. Clause (6) records honestly what is \textit{not} forced — the neutral partner — because no lepton-number or energy-sharing law exists in the inventory; naming it would be a declaration, and none is issued. Clause (7)'s $\Gamma$ is exact in $\varepsilon_q$ and carries no time unit beyond the tick [D1.8]; no lifetime is derived and none is compared [P5].

\textbf{Rmk (the forcing status of clause (6)).} The neutral emission is conditionally forced: BM.41(2)–(3) derives that every kinematically admissible transition vertex carries a third member, and that its class is unique among the constructed carrier classes — the $\nu$-channel line: a two-body vertex at the equal rest masses of BM.28(1) has no real on-shell solution for any positive member mass ($E_p-m_p=-m_\ell^2/(2m_n)<0$ identically), and the transit (EM-radiation) alternative is $\neg$constructed [Def-BM.8(8)]. Clause (6)'s inventory row is unchanged: $\neg\exists$ a lepton-number read — the forcing is kinematic plus uniqueness, $\neg$a number conservation; its hypothesis is the vertex's kinematic admissibility, which BM.41(4) shows no adopted read satisfies.

\textbf{Req.} A6 (Part 1); D1.2, D1.8, D3.3, D3.4, D3.10 (Parts 1–2); C8, CT0.8, CT3.18 (Part 1, Part 2); D2.11, D2.14 (Part 2); EM.9, EM.11 (Part 2); P5 (Parts 1–2); Def-H.32, FT.19, FT.41(d), FT.107(2), FT.108(3), FT.108(5), FT.109, FT.113, FT.119, FT.142(1), FT.144(2) (Part 3); Def-BM.5, Def-BM.6, BM.19(2), BM.19(3), BM.22 (this part); exact $\mathbb Q(i)$ ladder algebra, exhaustive ternary-alphabet censuses, exact rational resolvent identities and maxima [math; machine: \texttt{checks/\allowbreak{}verify\_\allowbreak{}H\_\allowbreak{}topo\_\allowbreak{}blind\_\allowbreak{}neutron.py} §D–§E].

\textbf{Pf.}

%dsp-level:0
\[
(1):\ E_\pm=\tfrac12(\sigma_1\pm i\sigma_2)\ \text{entrywise}\ [\text{machine}];
\]

%dsp-level:0
\[
\qquad [E_+,E_-]=\sigma_3,\ [\sigma_3,E_\pm]=\pm2E_\pm\ [\text{machine: all commutators exact over}\ \mathbb Q(i)]
\]

%dsp-level:0
\[
(2):\ \text{a}\ \pm2\ \text{move on}\ \sigma_s\ \text{changes}\ \tfrac12(\sigma_s-\sigma_t)\ \text{by}\ \pm1\ [\text{machine}];
\]

%dsp-level:0
\[
\qquad (3):\ \text{the shifted sets}\ \{0,1,2\},\ \{-2,-1,0\}\ \text{against the alphabet}\ \{-1,0,1\}\ [\text{machine}]
\]

%dsp-level:0
\[
(4):\ \mathrm{EM.11}\ \text{re-cited};\qquad\#\{q:q=-\Delta q\}=1\ \text{for}\ \Delta q=\pm1\ [\text{machine}];\qquad\Delta q+(-\Delta q)=0\ [\text{machine}]
\]

%dsp-level:0
\[
(5):\ (\mathrm{ii})\ \mathrm{FT.19}/\mathrm{CT3.18}/\mathrm{D2.14}\ \text{re-cited};\ (\mathrm{iii})\ \mathrm{D3.3}\ \text{re-cited};
\]

%dsp-level:0
\[
\qquad \text{eliminations}: \mathrm{FT.119}\ (\text{triplet-attached}),\ \mathrm{FT.142}(1)\ (\text{EM unactivated})\ \text{re-cited};
\]

%dsp-level:0
\[
\qquad \text{uniqueness}: \text{the channel inventory has three triplets}\ [\mathrm{FT.112},\ \mathrm{FT.119},\ \mathrm{FT.142}]
\]

%dsp-level:0
\[
(6):\ \text{inventory scan: no lepton-number read exists}\ [\mathrm{FT.81}(4)\ \text{pattern}];
\]

%dsp-level:0
\[
\qquad \Delta S=\tfrac32n_{\mathcal V}\ [\mathrm{FT.107}(2)\ \text{re-cited}]
\]

%dsp-level:0
\[
(7):\ \mathrm{FT.108}(5)\text{'s resolvent re-derived coefficient-wise at four rational thresholds}\ [\text{machine}];
\]

%dsp-level:0
\[
\qquad x_+x_-=4\varepsilon_q(1-4\varepsilon_q):\ \text{maximum}\ \tfrac14\ \text{at}\ \varepsilon_q=\tfrac18\ [\text{machine over}\ \varepsilon_q=k/256]
\]

%dsp-level:0
\[
(8):\ \varepsilon_q=0\implies x_+x_-=0\implies1/(1-0)=1\implies\text{no transition weight beyond the identity}\ [\text{machine}];
\]

%dsp-level:0
\[
\qquad (9):\ \mathrm{FT.144}(2),\ \mathrm{A6},\ \mathrm{D1.8}\ \text{re-cited}.\ \blacksquare
\]

\textbf{Phys.} The decay is assembled from three independent necessities and one honest gap. The weak ladder moves the $\sigma_3$ eigenvalue by two units, and the core charge is half a $\sigma$ difference, so any charge-changing transition moves the ternary read by exactly one — which pushes the endpoints straight out of the alphabet, so the change cannot stay inside the composite. Total charge vanishes in every valid universe, and the alphabet contains exactly one value that compensates a unit change, so the emitted carrier's charge is forced. Then the inventory does the rest: the quark-side channel is colour-attached and would break the seal, the neutrino-side channel is electromagnetically inert and could not carry the read, which leaves precisely one class — the $\ell$-channel line the naming layer calls the electron. The neutral partner is admissible and useful for carrying the deficit read, but no law in the corpus forces it, and the tier says so instead of naming a neutrino into the theorem. The rate is a geometric resolvent in the derived weights, its pole safely outside the unit disk, and it vanishes identically at the sharp crystallization threshold: a composite that cannot flicker cannot decay. The whole transition runs forward only, per tick, with no lifetime and no external unit anywhere in it.

\end{theorem}

\medskip\noindent\textit{\texorpdfstring{BM.24 — The status of the topology-blind composite: adjudication and scope}{BM.24 - The status of the topology-blind composite: adjudication and scope}.}\ Structural-constraints record: see Appendix B, Section B.1.17 of the Addendum.\par\medskip

\phantomsection
\section[Tier BM.7 — The Proton in the Positive CPT Sector and Its…]{\texorpdfstring{Tier BM.7 — The Proton in the Positive CPT Sector and Its Mirror: the Sector Labels, the Exact Phase, and the Charge-Blind Geometry}{Tier BM.7 - The Proton in the Positive CPT Sector and Its Mirror: the Sector Labels, the Exact Phase, and the Charge-Blind Geometry}}

\begin{definition}[{\texorpdfstring{Def-BM.7 — The EM$^\pm$ sector labels and the EM scar}{Def-BM.7 - The EM sector labels and the EM scar}}]\label{tag:Def-BM-7}
\textbf{Def.} Part-4 labels for two derived structures: the CPT-orientation classes of a charged sealed composite, and its endpoint charge read. Both labels are naming; both objects are derived.

%dsp-level:0
\[
(1)\quad\textbf{the sectors:}\qquad\mathrm{EM}^+\ :=\ \text{the CPT-orientation class with}\ q_{\mathrm{net}}=+1\ \text{and the unstarred weight}\ w(C);
\]

%dsp-level:0
\[
\mathrm{EM}^-\ :=\ \text{the class with}\ q_{\mathrm{net}}=-1\ \text{and the conjugated weight}\ w(C)^\ast;
\]

%dsp-level:0
\[
\textbf{constitutive clause:}\ \ \text{the two labels name the two orbits of the CPT involution on the charged classes}\ [\mathrm{EM.24}(4);\ \mathrm{BM.25}(4)];
\]

%dsp-level:0
\[
\qquad \textbf{non-derived clause:}\ \ \text{the signs}\ \pm\ \text{in the labels are a convention};
\]

%dsp-level:0
\[
(2)\quad\textbf{the scope clause (the mirror is formal, not dynamical):}
\]

%dsp-level:0
\[
\qquad \text{``backward time''}\ :=\ \Theta_{\mathfrak L}\ [\mathrm{D2.7}:\ \text{ledger order reversed},\ J^+\leftrightarrow J^-,\ G(e)\mapsto G(e)^\ast],
\]

%dsp-level:0
\[
\text{an involution on configuration data bijecting the completed admissible set onto itself}\ [\mathrm{FT.144}(4)(\mathrm i);\ \mathrm{FT.156}(5);\ \text{machine}]\ —\ \neg\text{a dynamical inverse}:
\]

%dsp-level:0
\[
\neg\exists\,\tau^{-1}\ [\mathrm{A6}];
\]

%dsp-level:0
\[
\qquad \text{the pure order reversal (labels unchanged) of every nonempty completed history is inadmissible at its first event}\ [\mathrm{FT.144}(4)(\mathrm{ii});\ \text{machine}];
\]

%dsp-level:0
\[
\therefore\ \mathrm{EM}^-\ \text{is a bookkeeping orbit of the mirror,}\ \neg\text{a universe running backwards};
\]

%dsp-level:0
\[
\qquad \text{no clause of this part asserts a backward evolution};
\]

%dsp-level:0
\[
(3)\quad\textbf{the EM scar:}
\]

%dsp-level:0
\[
\qquad \text{scar}(c)\ :=\ q_{\mathrm{core}}(c)=\tfrac12\bigl(\sigma_{s(c)}-\sigma_{t(c)}\bigr)\in\{-1,0,+1\}
\]

%dsp-level:0
\[
\qquad [\mathrm{D3.4},\ \mathrm{EM.9},\ \mathrm{EM.9a};\ \mathrm{FT.40}]:
\]

%dsp-level:0
\[
\text{the endpoint charge read of a materialized chain}\ —\ \text{registration data on the two ends}\ [\mathrm{D3.3},\ \mathrm{D1.1}],
\]

%dsp-level:0
\[
\ \neg\text{an internal edge}\ [\mathrm{D0.8}];
\]

%dsp-level:0
\[
\textbf{the zero-contribution clause:}\quad\text{the scar enters}\ E_{\mathrm{int}}\ \text{with coefficient}\ 0\ [\mathrm{BM.28}(3);\ \text{machine}:\ \text{Census FT.35}(2)\text{'s decomposition is}\ 2{+}2{+}2,\ \text{all edge counts}];
\]

%dsp-level:0
\[
(4)\quad\textbf{the neutral class:}
\]

%dsp-level:0
\[
\qquad q_{\mathrm{net}}=0\ \text{is the unique fixed point of the charge map}\ q\mapsto-q\ [\mathrm{EM.24}(1);\ \text{machine}]:
\]

%dsp-level:0
\[
\qquad \text{the neutral sector is charge-self-conjugate};
\]

%dsp-level:0
\[
\text{no configuration is fixed pointwise, since}\ \Theta_{\mathfrak L}\ \text{conjugates the weight}\ [\mathrm{CT2.8}(1)]:
\]

%dsp-level:0
\[
\qquad \text{self-conjugacy is a class property,}\ \neg\text{a pointwise one}\ [\mathrm{BM.25}(4)];
\]

%dsp-level:0
\[
(5)\quad\textbf{the data bar:}
\]

%dsp-level:0
\[
\qquad \text{``matter''/``antimatter'' are the labels of}\ (1)\ \text{and carry no external datum}\ [\mathrm{P5}];
\]

%dsp-level:0
\[
\qquad \text{no annihilation channel, no production channel and no cross-section is constructed in this corpus}\ [\mathrm{BM.29}(5)].
\]

\textbf{Rmk (vocabulary).} $\mathrm{EM}^\pm$ are Part-4 labels for CPT-orientation classes; they are not the register's $(S,T)$ sector symbols $\mathcal E_{\pm\pm}$ [CT0.12, CT1.1] and no clause identifies them. The pair-register sectors grade operators; the labels here grade composite configurations.

\textbf{Req.} A6, A8 (Part 1); D0.8, D1.1, D2.7, D3.3, D3.4 (Parts 1–2); CT2.8(1) (Part 2); EM.9, EM.9a, EM.24(1), EM.24(4) (Part 2); P5 (Parts 1–2); FT.40, FT.144(4), FT.156(5) (Part 3); Def-BM.3 (this part); clauses (3)–(4) carry marked forward pointers at their citing sites, imported in no \textbf{Req} line; Census FT.35(2) (row-wise); exhaustive enumeration of ledger records, exact $\mathbb Q(i)$ weight algebra [math; machine: \texttt{checks/\allowbreak{}verify\_\allowbreak{}H\_\allowbreak{}proton\_\allowbreak{}cpt.py} §A, §C].

\textbf{Pf.}

%dsp-level:0
\[
\text{Definition block: the constitutive clause is}\ (1)\text{'s orbit naming};
\]

%dsp-level:0
\[
\ \text{the non-derived clauses are the sign convention of}\ (1),
\]

%dsp-level:0
\[
\ \text{the scope reading of}\ (2)\ \text{and the vocabulary of}\ (5);
\]

%dsp-level:0
\[
\ (3)\text{'s zero-contribution clause and}\ (4)\text{'s fixed-point census are derived and machine-checked}.\ \blacksquare
\]

\textbf{Phys.} Two names are added and nothing else. The positive and negative CPT-orientation sectors label the two orbits that the corpus's own CPT involution makes of the charged classes, and the EM scar is the endpoint charge read the fermion era already derived. The one clause that does real work is the scope clause: the mirror is an involution on configuration data — reverse the ledger order and swap the seal labels — and it is not a way of running the universe backwards, because the tick has no inverse and every reversed history is inadmissible at its first event. So the negative sector is a bookkeeping orbit with equal amplitude, not a backwards-running world. The neutral class is the unique charge fixed point of the mirror, self-conjugate as a class while no single configuration is fixed, since the mirror always conjugates the weight.

\end{definition}

\begin{theorem}[{\texorpdfstring{BM.25 — The CPT involution on the sealed composite: orbits, the exact phase, and the census ratio}{BM.25 - The CPT involution on the sealed composite: orbits, the exact phase, and the census ratio}}]\label{tag:BM-25}
\textbf{Thm.}

%dsp-level:0
\[
(1)\quad\textbf{the three involutions are derived:}
\]

%dsp-level:0
\[
\qquad C:\ \sigma_e\mapsto-\sigma_e\ \forall e\implies q\mapsto-q,\quad f_v\mapsto-f_v,\quad\aleph_v\ \text{invariant}
\]

%dsp-level:0
\[
\qquad [\mathrm{EM.24}(1);\ \mathrm{D3.3}];
\]

%dsp-level:0
\[
P:\ \mathcal P\mapsto E_{\mathcal R}\setminus\mathcal P,\ \ \text{equivalently}\ \alpha_e\mapsto\pi-\alpha_e\qquad[\mathrm{EM.24}(2);\ \mathrm{Def\text{-}EM.21},\ \mathrm{EM.20}(1)];
\]

%dsp-level:0
\[
T:\ \Theta_{\mathfrak L},\qquad w(\Theta_{\mathfrak L}C)=w(C)^\ast\qquad[\mathrm{EM.24}(3);\ \mathrm{D2.7};\ \mathrm{CT2.8}(1)];
\]

%dsp-level:0
\[
\mathrm{CPT}:=C\circ P\circ T,\qquad C^2=P^2=T^2=\mathrm{id}\implies\mathrm{CPT}^2=\mathrm{id}\qquad[\mathrm{EM.24}(4):\ \text{each component derived, no external input}];
\]

%dsp-level:0
\[
(2)\quad\textbf{the mirror on the ledger is an exact involution with an unexpected fixed set:}
\]

%dsp-level:0
\[
\qquad \Theta_{\mathfrak L}^2=\mathrm{id}\ \text{on every admissible single-pair history}\ [\text{machine}];
\]

%dsp-level:0
\[
\Theta_{\mathfrak L}\bigl((J^+J^-)^m\bigr)=(J^+J^-)^m\ \ \forall m:
\]

%dsp-level:0
\[
\qquad \text{reversal composed with the}\ J^+\leftrightarrow J^-\ \text{relabelling fixes every completed word pointwise}\ [\text{machine: all}\ m\le6];
\]

%dsp-level:0
\[
\text{the pure order reversal (labels unchanged) is inadmissible at its first event}\ [\mathrm{FT.144}(4)(\mathrm{ii});\ \text{machine}]:
\]

%dsp-level:0
\[
\qquad \textbf{the formal mirror exists and the dynamical reversal does not};
\]

%dsp-level:0
\[
\Theta_{\mathfrak L}\ \text{bijects the two-pair completed interleavings onto themselves at}\ 2n
\]

%dsp-level:0
\[
= 2,
\]

%dsp-level:0
\[
4,
\]

%dsp-level:0
\[
6\ [\text{machine: set equality and cardinality}];
\]

%dsp-level:0
\[
\qquad \#J^+,\ \#J^-\ \text{are preserved}\implies\text{the Janus balance survives the mirror}\ [\mathrm{D1.3};\ \text{machine}];
\]

%dsp-level:0
\[
(3)\quad\textbf{the weight conjugation and the exact phase:}
\]

%dsp-level:0
\[
\qquad G(e)=\sigma_e\exp\!\bigl(i\alpha_e\widehat{ds}^{\,2}(e)\bigr),\quad\alpha_e\in\{0,\pi\}\ [\mathrm{D3.5},\ \mathrm{CT0.9}]\implies G(e)=\sigma_e\,i^{m_e},\quad m_e\in\mathbb Z:
\]

%dsp-level:0
\[
\text{the phase alphabet is the fourth roots of unity}\ —\ \text{exact in}\ \mathbb Q(i),\ \neg\text{a transcendental};
\]

%dsp-level:0
\[
\qquad \Theta_{\mathfrak L}:\ m_e\mapsto-m_e,\ \sigma_e\ \text{fixed},\ \mathcal J(C)\mapsto\mathcal J(C)^\ast\ [\mathrm{D3.6}];
\]

%dsp-level:0
\[
\boxed{\ w(\Theta_{\mathfrak L}C)=w(C)^\ast\ }
\]

%dsp-level:0
\[
\qquad [\mathrm{CT2.8}(1);\ \text{machine:}\ 24696\ \text{enumerated configurations}\ (\text{sign assignments}\times\text{half-integer}\ \kappa\ \text{ladder}\times\mathcal J\text{-phases})];
\]

%dsp-level:0
\[
w(\mathrm{CPT}\,C)=(-1)^{|E(C)|}\,e^{i\pi\Theta(C)}\,w(C),
\]

%dsp-level:0
\[
\qquad \lambda:=(-1)^{|E|}e^{i\pi\Theta}\in\{+1,-1,+i,-i\},
\]

%dsp-level:0
\[
\qquad |\lambda|^2=1
\]

%dsp-level:0
\[
\qquad [\mathrm{EM.24}(5),\ \mathrm{FT.134}(3);\ \text{machine}];
\]

%dsp-level:0
\[
\bigl|w(\mathrm{CPT}\,C)\bigr|^2=\bigl|w(C)\bigr|^2\ \ \forall C\qquad[\mathrm{EM.24}(5):\ \text{unconditional};\ \text{machine}];
\]

%dsp-level:0
\[
\textbf{the on-shell refinement:}\quad\widehat{ds}^{\,2}(e)\in\mathbb Z\ \forall e\implies\Theta\in\mathbb Z\implies\lambda=(-1)^{|E|+\Theta}\in\{\pm1\}:
\]

%dsp-level:0
\[
\qquad w(\mathrm{CPT}\,C)=w(C)^\ast\iff|E|+\Theta\equiv0\ (\mathrm{mod}\ 2)
\]

%dsp-level:0
\[
\qquad [\mathrm{EM.24}(5);\ \text{machine}];
\]

%dsp-level:0
\[
(4)\quad\textbf{the sector orbit census:}
\]

%dsp-level:0
\[
\qquad \mathrm{CPT}:\ (q,\ \text{orientation})\mapsto(-q,\ \text{flipped})\ \ \text{on}\ \{-1,0,+1\}\times\{w,w^\ast\}:
\]

%dsp-level:0
\[
\qquad \text{an involution with}\ \textbf{no fixed sector}\ [\text{machine}];
\]

%dsp-level:0
\[
\text{the orbit decomposition is exactly three pairs of size two};
\]

%dsp-level:0
\[
\qquad q\mapsto-q\ \text{has exactly one fixed point},\ q=0\ [\text{machine}]:
\]

%dsp-level:0
\[
\textbf{the neutral class is charge-self-conjugate and every charged class is exchanged}\ —\ \mathrm{EM}^+\leftrightarrow\mathrm{EM}^-\ [\mathrm{Def\text{-}BM.7}(1),(4)];
\]

%dsp-level:0
\[
(5)\quad\textbf{the census ratio (the sectors are equal in weight, the neutral is twice either):}
\]

%dsp-level:0
\[
\qquad Z^{\mathrm{norm}}(0,S)
\]

%dsp-level:0
\[
= 2\,
\]

%dsp-level:0
\[
Z^{\mathrm{norm}}(+1,S)\ \ \text{exactly for}\ S
\]

%dsp-level:0
\[
= 0,
\]

%dsp-level:0
\[
1,
\]

%dsp-level:0
\[
2,
\]

%dsp-level:0
\[
3\ [\text{machine}:\ 395824=2\cdot197912,\ 1574716=2\cdot787358,\ 1311552=2\cdot655776,\ 2001024=2\cdot1000512];
\]

%dsp-level:0
\[
Z^{\mathrm{norm}}(+1,S)=Z^{\mathrm{norm}}(-1,S)\ \ \text{exactly}\ [\mathrm{FT.106}(3\mathrm d);\ \text{machine}]:
\]

%dsp-level:0
\[
\qquad \textbf{the two charged sectors have identical weight};
\]

%dsp-level:0
\[
\text{the ratio is the factorization's own:}\ m_q2^{|R|-2}\ \text{with}\ m_q
\]

%dsp-level:0
\[
= (2,1,1)\ \text{at}\ q
\]

%dsp-level:0
\[
= (0,+1,-1)
\]

%dsp-level:0
\[
\implies m_0/m_{\pm1}
\]

%dsp-level:0
\[
= 2\ [\mathrm{FT.106}(3\mathrm a);\ \text{machine}];
\]

%dsp-level:0
\[
Z^{\mathrm{amp}}(+1,S)=Z^{\mathrm{amp}}(-1,S)\ \text{on the cover support},
\]

%dsp-level:0
\[
\qquad Z^{\mathrm{amp}}(q,S)=0\ \ \forall q,\ S\in\{2,3\}
\]

%dsp-level:0
\[
\qquad [\mathrm{FT.106}(3\mathrm c),(3\mathrm d);\ \text{machine}];
\]

%dsp-level:0
\[
(6)\quad\textbf{the baryon-type character:}
\]

%dsp-level:0
\[
\qquad \Theta_{\mathrm{int}}\in4\mathbb Z\ [\text{Census FT.49}]\implies e^{i\pi\Theta}=1;
\]

%dsp-level:0
\[
\qquad \Theta_{\mathrm{int}}\in2\mathbb Z\ [\text{Census FT.35}(3);\ \mathrm{BM.2}(1)]\implies e^{i\pi\Theta}=1\ \ \text{as well}\ [\text{machine}]:
\]

%dsp-level:0
\[
\text{on the}\ \varepsilon\text{-sealed junction class the CPT character is exactly}
\]

%dsp-level:0
\[
\qquad \lambda=(-1)^{|E|}
\]

%dsp-level:0
\[
\qquad [\mathrm{FT.134}(3);\ \text{machine}];
\]

%dsp-level:0
\[
(7)\quad\textbf{the norm pairing:}
\]

%dsp-level:0
\[
\qquad \langle A|A\rangle
\]

%dsp-level:0
\[
= \sum_C|w(C)|^2
\]

%dsp-level:0
\[
\ge 0\ [\mathrm{CT2.8}(3)]\ \text{is CPT-invariant term by term}:\ |w(C)|^2
\]

%dsp-level:0
\[
= |w(\Theta_{\mathfrak L}C)|^2\ [\text{machine on all enumerated pairs}]:
\]

%dsp-level:0
\[
\textbf{the two sectors contribute equally to the norm}\ —\ \text{the amplitude-level statement of}\ \mathrm{EM}^+\leftrightarrow\mathrm{EM}^-.
\]

\textbf{Rmk (locus discipline).} Clause (2)'s fixed set is a statement about \textit{words}: the mirror fixes each completed alternating word because reversing and relabelling commute to the same word. It is not a statement that any dynamical process is reversible — A6 and FT.144(4)(ii) bar that reading, and clause (2) records both rows side by side.

\textbf{Req.} A6, A8 (Part 1); D1.3, D2.7, D3.3, D3.4, D3.5, D3.6 (Parts 1–2); CT0.9 (Part 1); CT2.8(1), CT2.8(3), EM.9, EM.9a, EM.20(1), EM.24(1), EM.24(2), EM.24(3), EM.24(4), EM.24(5), Def-EM.21 (Part 2); Def-BM.7 (this part); FT.40, FT.106(3), FT.134(3), FT.144(4), FT.156(5) (Part 3); Census FT.35(3), FT.49 (row-wise); BM.2(1), BM.7(1) (this part); exact $\mathbb Q(i)$ weight algebra on the fourth-root alphabet, exhaustive ledger-record enumeration, exact integer census arithmetic [math; machine: \texttt{checks/\allowbreak{}verify\_\allowbreak{}H\_\allowbreak{}proton\_\allowbreak{}cpt.py} §A].

\textbf{Pf.}

%dsp-level:0
\[
(1):\ \mathrm{EM.24}(1)\text{–}(4)\ \text{rows re-cited}\ (\text{frozen, machine-verified at Tier F.33});
\]

%dsp-level:0
\[
\qquad \mathrm{CPT}^2=C^2P^2T^2=\mathrm{id}\ [\text{each an involution}]
\]

%dsp-level:0
\[
(2):\ \Theta_{\mathfrak L}(w)
\]

%dsp-level:0
\[
= \text{reverse}(w)\ \text{with labels swapped}: \text{on}\ (J^+J^-)^m\ \text{the reversal gives}\ (J^-J^+)^m\ \text{and the swap returns}\ (J^+J^-)^m\ [\text{machine}];
\]

%dsp-level:0
\[
\qquad \Theta_{\mathfrak L}^2=\mathrm{id}\ [\text{machine}];
\]

%dsp-level:0
\[
\qquad \text{pure reversal}: \text{leads with}\ J^-,\ \text{outside D1.2's domain}\ [\mathrm{FT.144}(4)(\mathrm{ii});\ \text{machine}]
\]

%dsp-level:0
\[
(3):\ \alpha_e\in\{0,\pi\}\ [\mathrm{CT0.9}]\ \wedge\ \widehat{ds}^{\,2}\in\tfrac12\mathbb Z\ [\mathrm{FT.22}\text{–}\mathrm{FT.24}\ \text{ladder}]\implies\alpha_e\widehat{ds}^{\,2}\in\tfrac\pi2\mathbb Z\implies\exp(i\alpha_e\widehat{ds}^{\,2})
\]

%dsp-level:0
\[
= i^{m_e}\ [\text{machine}];
\]

%dsp-level:0
\[
\qquad \Theta_{\mathfrak L}:\ G\mapsto G^\ast\ \text{per edge and}\ \mathcal J\mapsto\mathcal J^\ast\ [\mathrm{D2.7},\ \mathrm{D3.6}]\implies w\mapsto w^\ast\ [\text{machine}]
\]

%dsp-level:0
\[
\lambda=(-1)^{|E|}e^{i\pi\Theta}=i^{2|E|}i^{2\Theta}\in\{i^k\}\ [\text{machine}];
\]

%dsp-level:0
\[
\qquad |\lambda|^2=\lambda\lambda^\ast=1\ [\text{machine}];
\]

%dsp-level:0
\[
\qquad \text{on-shell}\ \Theta\in\mathbb Z\implies i^{2\Theta}=(-1)^\Theta\ [\text{machine}]
\]

%dsp-level:0
\[
(4):\ \text{the map on six sectors, its square, and its fixed-point census}\ [\text{machine}];
\]

%dsp-level:0
\[
\qquad q\mapsto-q\ \text{fixed points}: \{0\}\ [\text{machine}]
\]

%dsp-level:0
\[
(5):\ \mathrm{FT.106}(3)\text{'s census integers re-cited and the ratios verified exactly}\ [\text{machine}];
\]

%dsp-level:0
\[
\qquad m_q=(2,1,1)\ [\mathrm{FT.106}(3\mathrm a)]\implies m_0/m_{\pm1}=2\ \text{and}\ m_{+1}=m_{-1}\ [\text{machine}]
\]

%dsp-level:0
\[
(6):\ e^{i\pi\Theta}=i^{2\Theta}=1\ \text{at}\ \Theta\in2\mathbb Z\ [\text{machine}];\qquad\lambda=(-1)^{|E|}\cdot1\ [\mathrm{FT.134}(3)\ \text{re-cited}]
\]

%dsp-level:0
\[
(7):\ \mathrm{CT2.8}(3)\ \text{re-cited};\qquad|w|^2\ \text{equality term by term from}\ (3)\ [\text{machine}].\ \blacksquare
\]

\textbf{Phys.} The mirror does something sharper than exchange two sectors. On the ledger it is an involution that reverses order and swaps the seal labels — and on a completed history those two acts cancel, so every completed word is fixed pointwise, while the same history read backwards with its labels intact is inadmissible at its first event. The formal mirror exists; the dynamical reversal does not. Weights conjugate exactly, because the phase alphabet is the fourth roots of unity: with $\alpha_e\in\{0,\pi\}$ and a half-integer line element, every edge weight is a sign times a power of $i$, and negating the winding conjugates it. The CPT phase is therefore one of four exact values, its modulus is one, and the amplitude is invariant unconditionally — with the on-shell refinement that the phase is real and the weight literally conjugates exactly when the edge count plus the winding is even. The census closes the picture: the two charged sectors carry identical weight at every winding, the neutral sector exactly twice either, which is the factorization's own $m_q=(2,1,1)$ and not a coincidence. On the baryon-type junction, whose phase class lives in $4\mathbb Z$, the whole CPT character collapses to a sign counting edges.

\end{theorem}

\begin{theorem}[{\texorpdfstring{BM.26 — The proton row: the positive sector}{BM.26 - The proton row: the positive sector}}]\label{tag:BM-26}
\textbf{Thm.}

%dsp-level:0
\[
(1)\quad\textbf{the structure:}
\]

%dsp-level:0
\[
\qquad \text{the}\ \varepsilon\text{-sealed 3-juncture in the unit charge class}\ q_{\mathrm{net}}=+1
\]

%dsp-level:0
\[
\qquad [\mathrm{Def\text{-}BM.3}(2);\ \mathrm{D3.9}:\ B=\varepsilon_{ijk}c_1^ic_2^jc_3^k];
\]

%dsp-level:0
\[
\mathcal C(\Sigma;C)
\]

%dsp-level:0
\[
= 1\iff\text{all three attachments seal inside as one}\ \varepsilon\text{-type triple}\ [\text{Census FT.35}(5);\ \mathrm{CT3.17};\ \mathrm{FT.34}(5)]:
\]

%dsp-level:0
\[
\qquad \text{the colour seal is the observability switch}\ [\mathrm{BM.1}(1)];
\]

%dsp-level:0
\[
(2)\quad\textbf{the sector membership:}
\]

%dsp-level:0
\[
\qquad q_{\mathrm{net}}=+1\ \wedge\ w(C)\implies\ \text{the row lies in}\ \mathrm{EM}^+
\]

%dsp-level:0
\[
\qquad [\mathrm{Def\text{-}BM.7}(1);\ \text{the class label is}\ \mathrm{Def\text{-}BM.3}(4)\text{'s, pending BM.6}(1)];
\]

%dsp-level:0
\[
(3)\quad\textbf{the EM scar:}
\]

%dsp-level:0
\[
\qquad \text{scar}=q_{\mathrm{core}}=\tfrac12(\sigma_s-\sigma_t)\in\{-1,0,+1\}\ [\mathrm{D3.4},\ \mathrm{EM.9}];
\]

%dsp-level:0
\[
\qquad \text{the unit class carries}\ q_{\mathrm{core}}=+1\ \text{at its read end};
\]

%dsp-level:0
\[
\text{the scar is}\ \omega\text{-odd}\ [\mathrm{D2.6};\ \mathrm{FT.129}(3);\ \text{machine}]\ \text{and contributes}\ 0\ \text{to}\ E_{\mathrm{int}}\ [\mathrm{Def\text{-}BM.7}(3);\ \mathrm{BM.28}(3);\ \text{machine}];
\]

%dsp-level:0
\[
(4)\quad\textbf{the information content and the six reads:}
\]

%dsp-level:0
\[
\qquad E_{\mathrm{int}}^{(p)}=6\ (c_1{:}\,2,\ c_2{:}\,2,\ \text{connector}{:}\,2)
\]

%dsp-level:0
\[
\qquad [\text{Census FT.35}(2);\ \text{row-wise}];
\]

%dsp-level:0
\[
\Delta x_{\mathrm B}^{(p)}=\frac{\mathfrak h}{6\ln2},
\]

%dsp-level:0
\[
\qquad V_{\mathrm{eff}}^{(p)}=\frac{\mathfrak h^3}{216(\ln2)^3},
\]

%dsp-level:0
\[
\qquad m_{\mathrm{hid}}^{(p)}=\frac{6\ln2}{\mathfrak h},
\]

%dsp-level:0
\[
\qquad \rho^{(p)}=\frac{1296(\ln2)^4}{\mathfrak h^4};
\]

%dsp-level:0
\[
\kappa^2\rho^{(p)}=\frac{2592\,\pi(\ln2)^3}{\mathfrak h^2},
\]

%dsp-level:0
\[
\qquad r_s^{(p)}=3\,\mathfrak h,
\]

%dsp-level:0
\[
\qquad (\Delta x_{\mathrm B}^{(p)})^2\kappa^2\rho^{(p)}=72\,\pi\ln2
\]

%dsp-level:0
\[
\qquad [\mathrm{BM.13},\ \mathrm{BM.14},\ \mathrm{BM.15}(4),\ \mathrm{BM.16}(1),(4);\ \text{machine}];
\]

%dsp-level:0
\[
(5)\quad\textbf{the EM coupling:}
\]

%dsp-level:0
\[
\qquad F_{LR}=q\,I_2=I_2,
\]

%dsp-level:0
\[
\qquad \text{the chiral-partner splitting}\ 2|q|=2
\]

%dsp-level:0
\[
\qquad [\mathrm{BM.22}(2),(4);\ \mathrm{FT.129}(5),\ \mathrm{FT.130}(5);\ \text{machine}];
\]

%dsp-level:0
\[
\text{the promoted force matrix is}\ \mathbf F_{\mathrm{ext}}
\]

%dsp-level:0
\[
= h\otimes I_2+I_2\otimes\hat\kappa\ \text{with}\ \operatorname{spec}
\]

%dsp-level:0
\[
= \{\tfrac12(c_0+s|\mathbf c|)+t:\ s,t\in\{\pm1\}\}\ [\mathrm{BM.22}(3);\ \text{machine}];
\]

%dsp-level:0
\[
(6)\quad\textbf{the amplitude:}
\]

%dsp-level:0
\[
\qquad |w(\mathrm{CPT}\cdot C^{(p)})|^2=|w(C^{(p)})|^2\ [\mathrm{BM.25}(3)];
\]

%dsp-level:0
\[
\qquad \text{on-shell with}\ |E|+\Theta\equiv0\ (2):\ w(\mathrm{CPT}\cdot C^{(p)})=w(C^{(p)})^\ast;
\]

%dsp-level:0
\[
\text{the CPT image of this row is BM.27's}\ —\ \text{no independent datum is carried by either}.
\]

\textbf{Rmk (route scope of clause (4)).} The six reads of clause (4) are functions of $E_{\mathrm{int}}$ alone, so they are the crystallized-count route's rows: they are identical across the charge classes and clause (4) states exactly that. The ensemble-weighted mass read is route-relative — degenerate on the normalized route and split by the exact factor $\tfrac43$ between the neutral and the unit class on the total-deficit route [BM.28(4)]. Clause (4) is unchanged on its own route; the ensemble route is BM.28(4)'s and BM.30's content.

\textbf{Rmk (identification discipline).} The name attaches by Def-BM.3(2); the theorem content of every clause is the tuple $(\varepsilon\text{-junction},\ q_{\mathrm{net}}=+1,\ E_{\mathrm{int}}=6,\ \Theta_{\mathrm{int}}\in4\mathbb Z,\ 2|q|=2)$. No measured value enters [P5]; the charge class is the declared label pending the additive read [Def-BM.3(4); BM.6(1)].

\textbf{Req.} D3.4, D3.9 (Parts 1–2); CT3.17 (Part 2); EM.9 (Part 2); P5 (Parts 1–2); Def-BM.3, Def-BM.7, BM.1(1), BM.2(1), BM.13, BM.14, BM.15(4), BM.16(1), BM.16(4), BM.22(2), BM.22(3), BM.22(4) (this part); FT.34(5), FT.40, FT.129(3), FT.129(5), FT.130(5) (Part 3); Census FT.35(2), FT.35(3), FT.35(5), FT.49 (row-wise); BM.25(3) (this part); exact $\mathbb Q(i)$ block algebra, exact rational reads [math; machine: \texttt{checks/\allowbreak{}verify\_\allowbreak{}H\_\allowbreak{}proton\_\allowbreak{}cpt.py} §B–§C].

\textbf{Pf.}

%dsp-level:0
\[
(1)\text{–}(2):\ \mathrm{Def\text{-}BM.3}(2)\ \text{and}\ \mathrm{Def\text{-}BM.7}(1)\ \text{re-cited};\qquad\text{Census FT.35}(5)\text{'s sealing row re-cited}
\]

%dsp-level:0
\[
(3):\ \mathrm{D3.4}/\mathrm{EM.9}\ \text{re-cited};
\]

%dsp-level:0
\[
\qquad \omega\text{-oddness}: q_{\mathrm{core}}(\omega c)=\tfrac12(\sigma_t-\sigma_s)=-q_{\mathrm{core}}(c)\ [\mathrm{D2.6};\ \text{machine}];
\]

%dsp-level:0
\[
\qquad \text{the zero contribution}: \mathrm{BM.28}(3)
\]

%dsp-level:0
\[
(4):\ \text{substitution of}\ E_{\mathrm{int}}
\]

%dsp-level:0
\[
= 6\ \text{into BM.13(1), BM.14(1),(2),(3), BM.15(4), BM.16(1),(4)}\ [\text{machine: every coefficient exact}]
\]

%dsp-level:0
\[
(5):\ \mathrm{BM.22}(2)\text{–}(4)\ \text{re-derived at}\ q=+1\ [\text{machine}];\qquad(6):\ \mathrm{BM.25}(3)\ \text{re-cited}.\ \blacksquare
\]

\textbf{Phys.} The proton row is the neutron row with one entry changed. Same junction grammar, same six internal directed edges, same sealing phase class, same width, volume, mass, density, curvature scale, comparison radius and lock — because none of those reads contains a charge symbol. What changes is the flipping block: the promoted static force matrix carries $q I_2$ off-diagonal, so the unit class splits its chiral partners by exactly two read units where the neutral class splits them by zero. The scar is the endpoint read, it is odd under the orientation reversal, and it contributes nothing to the edge count — which is precisely why the geometry is blind to it. The amplitude of the row equals the amplitude of its mirror, unconditionally, and on-shell the mirror's weight is literally the complex conjugate.

\end{theorem}

\begin{theorem}[{\texorpdfstring{BM.27 — The antiproton row: the CPT image, and the resolution of the identity question}{BM.27 - The antiproton row: the CPT image, and the resolution of the identity question}}]\label{tag:BM-27}
\textbf{Thm.}

%dsp-level:0
\[
(1)\quad\textbf{the image:}
\]

%dsp-level:0
\[
\qquad C^{(\bar p)}:=\mathrm{CPT}\cdot C^{(p)}:
\]

%dsp-level:0
\[
\qquad q_{\mathrm{net}}\mapsto-q_{\mathrm{net}}=-1\ [\mathrm{EM.24}(1)],\quad w\mapsto w^\ast\ [\mathrm{CT2.8}(1)],\quad\mathcal P\mapsto E_{\mathcal R}\setminus\mathcal P\ [\mathrm{EM.24}(2)];
\]

%dsp-level:0
\[
\text{the image lies in}\ \mathrm{EM}^-\qquad[\mathrm{Def\text{-}BM.7}(1);\ \mathrm{BM.25}(4):\ \text{the two charged classes form one orbit}];
\]

%dsp-level:0
\[
(2)\quad\textbf{the structure is the same junction:}
\]

%dsp-level:0
\[
\qquad \varepsilon\text{-sealed 3-juncture in the}\ q_{\mathrm{net}}=-1\ \text{class}\ [\mathrm{Def\text{-}BM.3}(2):\ \text{the antiproton row}];
\]

%dsp-level:0
\[
E_{\mathrm{int}}^{(\bar p)}=E_{\mathrm{int}}^{(p)}=6:
\]

%dsp-level:0
\[
\qquad \Theta_{\mathfrak L}\ \text{preserves the record size and exchanges}\ \#J^+\ \text{with}\ \#J^-\ [\mathrm{BM.25}(2);\ \text{machine}]
\]

%dsp-level:0
\[
\implies |R|\ \text{and hence}\ E_{\mathrm{int}}
\]

%dsp-level:0
\[
= |R|\ \text{are invariant}\ [\mathrm{Def\text{-}BM.5}(4);\ \text{machine}];
\]

%dsp-level:0
\[
(3)\quad\textbf{every read is the image's, so every read is identical:}
\]

%dsp-level:0
\[
\qquad \Delta x_{\mathrm B}^{(\bar p)}=\Delta x_{\mathrm B}^{(p)},
\]

%dsp-level:0
\[
\quad V_{\mathrm{eff}}^{(\bar p)}=V_{\mathrm{eff}}^{(p)},
\]

%dsp-level:0
\[
\quad m_{\mathrm{hid}}^{(\bar p)}=m_{\mathrm{hid}}^{(p)},\quad\rho^{(\bar p)}=\rho^{(p)},\quad\kappa^2\rho^{(\bar p)}=\kappa^2\rho^{(p)},
\]

%dsp-level:0
\[
\quad r_s^{(\bar p)}=r_s^{(p)};
\]

%dsp-level:0
\[
\text{the reason is structural, not numerical: none of the six monomials carries a charge exponent}\ [\mathrm{BM.28}(1);\ \text{machine: the}\ q\text{-exponent is}\ 0\ \text{in all six}];
\]

%dsp-level:0
\[
(4)\quad\textbf{the amplitude and the character:}
\]

%dsp-level:0
\[
\qquad |w(C^{(\bar p)})|^2=|w(C^{(p)})|^2\ [\mathrm{EM.24}(5);\ \text{machine}];
\]

%dsp-level:0
\[
\qquad w(C^{(\bar p)})=\lambda\,w(C^{(p)}),\quad\lambda=(-1)^{|E|}e^{i\pi\Theta};
\]

%dsp-level:0
\[
\Theta_{\mathrm{int}}\in4\mathbb Z\ [\text{Census FT.49}]\implies\lambda=(-1)^{|E|}\in\{\pm1\}\qquad[\mathrm{BM.25}(6);\ \text{machine}];
\]

%dsp-level:0
\[
\text{the two rows have equal norm contribution to}\ \langle A|A\rangle
\]

%dsp-level:0
\[
= \sum_C|w(C)|^2\ [\mathrm{CT2.8}(3);\ \mathrm{BM.25}(7);\ \text{machine}];
\]

%dsp-level:0
\[
(5)\quad\textbf{the difference set is exactly two entries:}
\]

%dsp-level:0
\[
\qquad \text{scar}^{(\bar p)}=-\,\text{scar}^{(p)}\ [\omega\text{-odd};\ \text{machine}]\ \text{and the sector orientation}\ w\leftrightarrow w^\ast;
\]

%dsp-level:0
\[
F_{LR}^{(\bar p)}=-I_2=-F_{LR}^{(p)}\ \text{but}\ 2|q^{(\bar p)}|=2=2|q^{(p)}|:
\]

%dsp-level:0
\[
\qquad \textbf{the flipping block is odd and the splitting is even}\ [\mathrm{BM.22}(2),(4);\ \text{machine}];
\]

%dsp-level:0
\[
(6)\quad\textbf{the identity adjudication:}
\]

%dsp-level:0
\[
\qquad C^{(\bar p)}
\]

%dsp-level:0
\[
\ne C^{(p)}\ \text{as configurations}\ [\text{the scar and the orientation differ};\ \mathrm{BM.25}(4):\ \neg\exists\ \text{a fixed charged sector}];
\]

%dsp-level:0
\[
C^{(\bar p)}
\]

%dsp-level:0
\[
= \mathrm{CPT}\cdot C^{(p)}\ \text{with}\ \mathrm{CPT}^2
\]

%dsp-level:0
\[
= \mathrm{id}\ [\mathrm{EM.24}(4)]\implies\textbf{the antiproton row carries no independent datum}: \text{it is one orbit of size two, not two structures};
\]

%dsp-level:0
\[
\text{both statements hold because they are about different objects}\ —\ \text{configurations}\ (6)\text{'s first line, structure}\ (6)\text{'s second};
\]

%dsp-level:0
\[
\qquad \text{no clause asserts an identity of configurations};
\]

%dsp-level:0
\[
(7)\quad\textbf{what is not constructed:}
\]

%dsp-level:0
\[
\qquad \neg\exists\ \text{an annihilation channel},
\]

%dsp-level:0
\[
\ \neg\exists\ \text{a pair-production channel},
\]

%dsp-level:0
\[
\ \neg\exists\ \text{a cross-section between}\ \mathrm{EM}^+\ \text{and}\ \mathrm{EM}^-\ [\mathrm{Def\text{-}BM.7}(5);\ \mathrm{BM.29}(5)]:
\]

%dsp-level:0
\[
\text{the two sectors are amplitude-equal configuration orbits; their interaction is not a theorem of this corpus}.
\]

\textbf{Rmk (route scope of clause (3)).} The six identical reads of clause (3) carry no charge exponent [BM.28(1)], and the equality of the proton and antiproton rows holds on \textbf{every} mass route: the charge census is symmetric, $W(+1)=W(-1)=15$ exactly, and both ensemble layers are $\omega$-symmetric [BM.28(5); FT.106(3d)]. What the census does split is the neutral class against the charged pair, by the exact factor $\tfrac43$ on the total-deficit route [BM.28(4)(ii)] — a splitting that leaves this tag's CPT rows untouched.

\textbf{Rmk (mirror scope).} The image is configuration data under $\Theta_{\mathfrak L}$ composed with the sign and parity maps. It is not a time-reversed process: A6 gives the tick no inverse, and the pure order reversal of a completed history is inadmissible at its first event [BM.25(2)]. Reading clause (1) as a dynamical evolution would contradict FT.144(2), and no clause here does.

\textbf{Req.} A6 (Part 1); D2.6, D2.7 (Parts 1–2); CT2.8(1), CT2.8(3) (Part 2); EM.24(1), EM.24(2), EM.24(4), EM.24(5) (Part 2); P5 (Parts 1–2); Def-BM.3, Def-BM.5(4), Def-BM.7, BM.22(2), BM.22(4), BM.25(2), BM.25(4), BM.25(6), BM.25(7), BM.26 (this part); FT.129(3), FT.130(5), FT.144(2), FT.144(4) (Part 3); Census FT.49 (row-wise); clause (3) carries a marked forward pointer at its citing site, imported in no \textbf{Req} line; exact $\mathbb Q(i)$ arithmetic, exhaustive enumeration [math; machine: \texttt{checks/\allowbreak{}verify\_\allowbreak{}H\_\allowbreak{}proton\_\allowbreak{}cpt.py} §A–§C].

\textbf{Pf.}

%dsp-level:0
\[
(1):\ \mathrm{EM.24}(1),(2)\ \text{and}\ \mathrm{CT2.8}(1)\ \text{composed}\ [\mathrm{EM.24}(4)];\qquad\text{orbit membership}: \mathrm{BM.25}(4)\text{'s census}
\]

%dsp-level:0
\[
(2):\ \text{Census FT.35}(2)\text{'s row re-cited};
\]

%dsp-level:0
\[
\qquad |\Theta_{\mathfrak L}(w)|=|w|\ \text{and}\ \#J^+(\Theta w)=\#J^-(w)\ [\text{machine}]\implies E_{\mathrm{int}}=|R|\ \text{invariant}
\]

%dsp-level:0
\[
(3):\ \text{the six monomials of BM.13(1), BM.14(1),(2), BM.15(4), BM.16(4) have}\ q\text{-exponent}\ 0\ [\text{machine: monomial census}]\implies\text{substituting}\ q
\]

%dsp-level:0
\[
\mapsto -q\ \text{fixes each}
\]

%dsp-level:0
\[
(4):\ \mathrm{EM.24}(5)\text{'s identity re-cited and re-verified}\ [\text{machine}];
\]

%dsp-level:0
\[
\qquad \lambda=(-1)^{|E|}\ \text{at}\ \Theta\in4\mathbb Z\ [\mathrm{BM.25}(6)]
\]

%dsp-level:0
\[
(5):\ q_{\mathrm{core}}\mapsto-q_{\mathrm{core}}\ [\mathrm{D2.6};\ \text{machine}];\qquad F_{LR}=qI_2\mapsto-qI_2\ \text{while}\ 2|q|\mapsto2|q|\ [\text{machine}]
\]

%dsp-level:0
\[
(6):\ \neg\exists\ \text{a fixed charged sector}\ [\mathrm{BM.25}(4);\ \text{machine}]\implies\text{distinct configurations};
\]

%dsp-level:0
\[
\qquad \mathrm{CPT}^2=\mathrm{id}\ [\mathrm{EM.24}(4)]\implies\text{one orbit of size two}
\]

%dsp-level:0
\[
(7):\ \text{inventory scan: no annihilation, production or inter-sector transition operator exists}\ [\mathrm{BM.23}(6)\ \text{pattern}];
\]

%dsp-level:0
\[
\qquad \mathrm{Def\text{-}BM.7}(5)\ \text{re-cited}.\ \blacksquare
\]

\textbf{Phys.} The antiproton is what the mirror does to the proton, and the mirror is cheap: flip every edge sign, complement the alignment pattern, reverse the ledger while swapping the seal labels. Everything the corpus derives about the composite survives that map untouched, because not one of the six geometric reads contains a charge — the width, the volume, the mass, the density, the curvature scale and the comparison radius come out identical, and the amplitude comes out equal in modulus with a phase that is one of four exact roots of unity, collapsing to a bare sign on the baryon-type phase class. Only two entries change: the scar flips sign and the weight conjugates. The flipping block flips with the scar while the splitting, being $2|q|$, does not — so the proton and antiproton split their chiral partners by the same two units in opposite directions. Whether they are "the same object" depends on whether one is counting configurations or structures: as configurations they differ and no charged sector is fixed; as structure they are one orbit of size two, and the second carries no datum the first does not. What the corpus does not have is any process connecting the two sectors — no annihilation, no pair production, no cross-section — and that absence is recorded rather than filled.

\end{theorem}

\begin{theorem}[{\texorpdfstring{BM.28 — The charge-constrained probability space: the census, the three mass routes, and the splitting}{BM.28 - The charge-constrained probability space: the census, the three mass routes, and the splitting}}]\label{tag:BM-28}
\textbf{Thm.}

%dsp-level:0
\[
(1)\quad\textbf{the geometric monomials are charge-blind (the crystallized route):}
\]

%dsp-level:0
\[
\qquad \Delta x_{\mathrm B}=\mathfrak h^1E^{-1}(\ln2)^{-1},
\]

%dsp-level:0
\[
\quad V_{\mathrm{eff}}=\mathfrak h^3E^{-3}(\ln2)^{-3},\quad\rho=\mathfrak h^{-4}E^4(\ln2)^4,\quad\kappa^2\rho=2\pi\,\mathfrak h^{-2}E^4(\ln2)^3;
\]

%dsp-level:0
\[
r_s=\tfrac12\mathfrak h E,
\]

%dsp-level:0
\[
\qquad (\Delta x_{\mathrm B})^2\kappa^2\rho=2\pi E^2\ln2,
\]

%dsp-level:0
\[
\qquad \Delta x_{\mathrm B}r_s=\tfrac{\mathfrak h^2}{2\ln2}
\]

%dsp-level:0
\[
\qquad [\mathrm{BM.13}(1),\ \mathrm{BM.14}(1)\text{–}(3),\ \mathrm{BM.15}(4),\ \mathrm{BM.16}(1),(4)]:
\]

%dsp-level:0
\[
\textbf{the charge exponent is}\ 0\ \textbf{in all seven}\ [\text{machine: monomial census}]\implies\text{at fixed}\ E_{\mathrm{int}}\ \text{these seven reads are identical on every charge class};
\]

%dsp-level:0
\[
\textbf{scope:}\ \text{this clause governs reads that are functions of}\ E_{\mathrm{int}}\ \text{alone}\ —\ \text{the crystallized-count route};
\]

%dsp-level:0
\[
\qquad \text{it says nothing about ensemble-weighted reads}\ [(2)\text{–}(5)];
\]

%dsp-level:0
\[
(2)\quad\textbf{the charge constraint is a filter on the probability space, and the filter differs by class:}
\]

%dsp-level:0
\[
\qquad q_i=\tfrac12(\sigma_{s(c_i)}-\sigma_{t(c_i)})\in\{-1,0,+1\}\ [\mathrm{D3.4},\ \mathrm{EM.9}];
\]

%dsp-level:0
\[
\qquad Q=q_1+q_2+q_3\ [\mathrm{D3.9};\ \text{the additive read of}\ \mathrm{BM.6}(1)\text{'s spec}];
\]

%dsp-level:0
\[
\textbf{the end-pair census:}\quad m_q=(2,1,1)\ \text{at}\ q=(0,+1,-1)
\]

%dsp-level:0
\[
\qquad [\mathrm{EM.9}\text{'s four end-pairs};\ \mathrm{FT.106}(3\mathrm a);\ \text{machine}]:
\]

%dsp-level:0
\[
\qquad \text{the neutral read has two sign realizations, each charged read one};
\]

%dsp-level:0
\[
\textbf{the full census of the 3-constituent composite}\ (4^3=64\ \text{end-sign assignments}):
\]

%dsp-level:0
\[
\qquad \#\{Q=k\}=(1,\,6,\,15,\,\mathbf{20},\,15,\,6,\,1)\ \ \text{for}\ k=-3,\dots,+3
\]

%dsp-level:0
\[
\qquad [\text{machine: exhaustive}];
\]

%dsp-level:0
\[
\text{the generating function}\ (2+z+z^{-1})^3\ \text{reproduces it exactly and sums to}\ 64\ [\text{machine}];
\]

%dsp-level:0
\[
\textbf{the ordered charge-triple census:}\quad\#\{(q_1,q_2,q_3):\ \textstyle\sum q_i=k\}=(1,3,6,\mathbf7,6,3,1):
\]

%dsp-level:0
\[
\qquad \textbf{7 neutral triples against 6 charged}\ [\text{machine}];
\]

%dsp-level:0
\[
\textbf{the probability-space sizes:}
\]

%dsp-level:0
\[
\qquad W(0)=20,
\]

%dsp-level:0
\[
\qquad W(+1)=W(-1)=15,
\]

%dsp-level:0
\[
\qquad \boxed{\ \frac{W(0)}{W(\pm1)}=\frac{20}{15}=\frac43\ }
\]

%dsp-level:0
\[
\qquad [\text{machine: exact}];
\]

%dsp-level:0
\[
\text{one chain restricts}\ 2{:}1\ [m_q];\ \text{three chains under the additive constraint restrict}\ 4{:}3:
\]

%dsp-level:0
\[
\qquad \textbf{the constraint is what changes the ratio};
\]

%dsp-level:0
\[
(3)\quad\textbf{the scar is a boundary condition, not an internal edge:}
\]

%dsp-level:0
\[
\qquad E_{\mathrm{int}}=\#\{\text{internal directed edges}\}\ [\mathrm{D0.8}];
\]

%dsp-level:0
\[
\qquad \text{scar}=q_{\mathrm{core}}\ \text{is endpoint registration data}\ [\mathrm{D1.1},\ \mathrm{D3.3},\ \mathrm{D3.4}];
\]

%dsp-level:0
\[
\text{Census FT.35}(2)\text{'s decomposition}\ 6
\]

%dsp-level:0
\[
= (c_1{:}\,2)+(c_2{:}\,2)+(\text{connector}{:}\,2)\ \text{has}\ \textbf{no endpoint summand}\ [\text{machine}];
\]

%dsp-level:0
\[
\qquad E_{\mathrm{int}}\ \text{is unchanged across all four endpoint sign pairs}\ [\text{machine}];
\]

%dsp-level:0
\[
\therefore\ \text{the scar shifts no crystallized-count read}\ [(1)]\ —\ \textbf{and it does not follow that the scar shifts no read}: \text{the charge constraint acts on the ensemble, not on the count}\ [(2)];
\]

%dsp-level:0
\[
(4)\quad\textbf{the route theorem (the mass read is route-relative):}
\]

%dsp-level:0
\[
\qquad \textbf{(i) the normalized ensemble average}\ \ m^2_B(q,S)
\]

%dsp-level:0
\[
= \dfrac{\sum_R\nu\,t^{|R|}u^L\#\mathrm{inst}\cdot\Delta S(R)}{\sum_R\nu\,t^{|R|}u^L\#\mathrm{inst}}\ \bigg|_{t=u=1}\ [\mathrm{FT.107}(3)]:
\]

%dsp-level:0
\[
\text{a uniform sector factor cancels in a ratio}\ —\ \text{verified exactly for factors}\ 1,
\]

%dsp-level:0
\[
2,
\]

%dsp-level:0
\[
3,
\]

%dsp-level:0
\[
15,
\]

%dsp-level:0
\[
20,
\]

%dsp-level:0
\[
\tfrac73\ [\text{machine}]\implies\textbf{charge-degeneracy on this route is structural},
\]

%dsp-level:0
\[
\ \neg\text{an oversight}:
\]

%dsp-level:0
\[
\Delta S(\Gamma,R)
\]

%dsp-level:0
\[
= \tfrac32n_{\mathcal V}(\Gamma,R)
\]

%dsp-level:0
\[
= \tfrac32(2E-|R|)\ \text{is}\ \sigma\text{-free}\ [\mathrm{FT.107}(2)],
\]

%dsp-level:0
\[
\ \text{so the restriction rescales the measure without reshaping it};
\]

%dsp-level:0
\[
\textbf{(ii) the unnormalized total deficit}\ \ T(q,S)
\]

%dsp-level:0
\[
= \sum_R w(R)\,
\]

%dsp-level:0
\[
\Delta S(R)\ [\text{the}\ I_{\mathrm{hid}}\text{-type read};\ \mathrm{FT.42}(1)\ \text{pattern}]:
\]

%dsp-level:0
\[
\qquad T(0,S)=\tfrac43\,T(\pm1,S)\ \ \text{exactly}\ [\text{machine}]:
\]

%dsp-level:0
\[
\boxed{\ \Delta m\ne0\ \text{on this route},\qquad\frac{\langle m\rangle^{(n)}}{\langle m\rangle^{(p)}}=\frac43,\qquad\text{the neutral class is heavier}\ }
\]

%dsp-level:0
\[
\qquad [\text{machine: exact, ordering derived}];
\]

%dsp-level:0
\[
\textbf{(iii) the amplitude layer}\ [\mathrm{Def\text{-}H.32}(2);\ \mathrm{FT.106}(3)]:
\]

%dsp-level:0
\[
\qquad \frac{Z^{\mathrm{amp}}(0,0)}{Z^{\mathrm{amp}}(\pm1,0)}=\frac{5336}{1196}=\frac{58}{13},
\]

%dsp-level:0
\[
\qquad \frac{Z^{\mathrm{amp}}(0,1)}{Z^{\mathrm{amp}}(\pm1,1)}=\frac{4796}{802}=\frac{2398}{401}
\]

%dsp-level:0
\[
\qquad [\text{machine}]:
\]

%dsp-level:0
\[
\text{a third ratio, winding-dependent,}\ \neg\text{equal to}\ \tfrac43\ \text{or}\ 2;
\]

%dsp-level:0
\[
\qquad \text{its bridge to a mass read is item-13 content and is}\ \neg\text{constructed}\;
\]

%dsp-level:0
\[
(5)\quad\textbf{what the splitting does not touch:}
\]

%dsp-level:0
\[
\qquad W(+1)=W(-1)\ \text{exactly}\ [\text{machine}]\ \text{and}\ Z^{\bullet}(+1,S)=Z^{\bullet}(-1,S)\ [\mathrm{FT.106}(3\mathrm d)]:
\]

%dsp-level:0
\[
\textbf{the proton and the antiproton remain degenerate on all three routes}\ —\ \text{the charge restriction splits the neutral class from the charged pair and never splits the pair}\ [\mathrm{BM.27}(3),(4)];
\]

%dsp-level:0
\[
\text{CPT maps the}\ Q
\]

%dsp-level:0
\[
= +1\ \text{triple set bijectively onto the}\ Q
\]

%dsp-level:0
\[
= -1\ \text{set, and the neutral set onto itself with}\ (0,0,0)\ \text{the}\ \textbf{unique CPT-fixed charge triple}\ [\text{machine}];
\]

%dsp-level:0
\[
(6)\quad\textbf{the identity and difference sets:}
\]

%dsp-level:0
\[
\qquad \textbf{identical on every route:}\ E_{\mathrm{int}},
\]

%dsp-level:0
\[
\ \Delta x_{\mathrm B},
\]

%dsp-level:0
\[
\ V_{\mathrm{eff}},
\]

%dsp-level:0
\[
\ \rho,
\]

%dsp-level:0
\[
\ \kappa^2\rho,
\]

%dsp-level:0
\[
\ r_s,
\]

%dsp-level:0
\[
\ \text{both locks},
\]

%dsp-level:0
\[
\ \Theta_{\mathrm{int}}\ \text{class},
\]

%dsp-level:0
\[
\ |w|^2,
\]

%dsp-level:0
\[
\ 2|q|\ \text{restricted to}\ \{|q|=0\}\ \text{vs}\ \{|q|=1\}\ [(1)];
\]

%dsp-level:0
\[
\textbf{route-dependent:}\ \ \text{the ensemble mass read}\ —\ \text{degenerate on}\ (4)(\mathrm i),
\]

%dsp-level:0
\[
\ \text{split by}\ \tfrac43\ \text{on}\ (4)(\mathrm{ii}),
\]

%dsp-level:0
\[
\ \text{split by}\ \tfrac{58}{13}/\tfrac{2398}{401}\ \text{on}\ (4)(\mathrm{iii});
\]

%dsp-level:0
\[
\textbf{distinguishing on every route:}\ \ F_{LR}
\]

%dsp-level:0
\[
= qI_2,
\]

%dsp-level:0
\[
\ q_{\mathrm{core}},
\]

%dsp-level:0
\[
\ k,
\]

%dsp-level:0
\[
\ \operatorname{sign}(k),
\]

%dsp-level:0
\[
\ \text{the CPT orientation},
\]

%dsp-level:0
\[
\ \text{the sector label},
\]

%dsp-level:0
\[
\ \text{and now}\ W(Q)\ [\text{machine: the sets are disjoint}];
\]

%dsp-level:0
\[
(7)\quad\textbf{no fractional charge exists, structurally:}
\]

%dsp-level:0
\[
\qquad q\in\{-1,0,+1\}\subset\mathbb Z\ [\mathrm{EM.9},\ \mathrm{P2}]\ \text{and}\ Q=\sum_iq_i\implies Q\in\mathbb Z\ \text{with}\ Q\in\{-3,\dots,+3\}\ [\text{machine}]:
\]

%dsp-level:0
\[
\tfrac13,\ \tfrac23\notin\{-3,\dots,+3\}:
\]

%dsp-level:0
\[
\qquad \textbf{fractional charges are impossible under an integer-valued additive read},\ \neg\text{merely absent};
\]

%dsp-level:0
\[
\text{deficit}\ 1'\text{'s charge side}\ [\mathrm{FT.125}(4)]\ \text{is therefore not closable by this read}: \text{a fractional constituent alphabet requires a different read, and none is constructed};
\]

%dsp-level:0
\[
(8)\quad\textbf{the read-choice row:}
\]

%dsp-level:0
\[
\qquad \text{which of}\ (4)\text{'s three routes is ``the mass'' is a read specification},
\]

%dsp-level:0
\[
\ \neg\text{a theorem}\ [\text{Census FT.76's priced-variant pattern}];
\]

%dsp-level:0
\[
\text{all three are exact and machine-verified};
\]

%dsp-level:0
\[
\qquad \text{no route imports a measured value}\ [\mathrm{P5}];
\]

%dsp-level:0
\[
\qquad \neg\exists\ \text{a P4 parameter and}\ \neg\exists\ \text{a new dimensionful scale}\ [\mathrm{P1};\ \text{§Cal-BM stays empty}].
\]

\textbf{Rmk (route discipline).} Clause (1) and clause (4)(i) are not in tension with clause (4)(ii): they are statements about different functionals of the same ensemble. A normalized average is blind to any uniform rescaling of the measure; a total is not. The charge constraint of clause (2) is exactly such a rescaling — uniform within a sector, different between sectors — so it is invisible to the first and visible to the second, with the exact factor $4/3$. Any reading of this tag that quotes one route against another is reading a scope clause as a claim.

\textbf{Req.} D0.8, D1.1, D2.6, D3.3, D3.4, D3.9 (Parts 1–2); CT2.1, CT2.9 (Part 2); EM.9, EM.24(1) (Part 2); P1, P2, P5 (Parts 1–2); Def-BM.4, Def-BM.7, BM.13(1), BM.14(1), BM.14(2), BM.14(3), BM.15(4), BM.16(1), BM.16(4), BM.22(2), BM.22(4), BM.27(3), BM.27(4) (this part); Def-H.32, FT.42(1), FT.84(1), FT.90(5), FT.106(3), FT.107(2), FT.107(3), FT.125(4), FT.129(3) (Part 3); Census FT.35(2) (row-wise); the additive read of BM.6(1)'s specification (constructed in Tier BM.8, cited there as a forward pointer and imported in no \textbf{Req} line); exhaustive end-sign enumeration, generating-function convolution, exact rational ratio arithmetic [math; machine: \texttt{checks/\allowbreak{}verify\_\allowbreak{}H\_\allowbreak{}probability\_\allowbreak{}space\_\allowbreak{}restriction.py} §A–§C].

\textbf{Pf.}

%dsp-level:0
\[
(1):\ \text{the seven monomials re-cited from Tiers BM.5's rows and their exponent vectors read off}\ [\text{machine}];
\]

%dsp-level:0
\[
\qquad q\text{-exponent}\ 0\ \text{in each}\ [\text{machine}]
\]

%dsp-level:0
\[
(2):\ \mathrm{EM.9}\text{'s table}\ \to m_q=(2,1,1)\ [\text{machine}];
\]

%dsp-level:0
\[
\qquad \text{exhaustive enumeration of}\ 4^3
\]

%dsp-level:0
\[
= 64\ \text{end-sign assignments}\ \to\ \#\{Q=k\}
\]

%dsp-level:0
\[
= (1,6,15,20,15,6,1)\ [\text{machine}];
\]

%dsp-level:0
\[
\qquad (2+z+z^{-1})^3\ \text{convolved three times reproduces the same census and sums to}\ 64\ [\text{machine}]
\]

%dsp-level:0
\[
\text{triples}: \#\{\textstyle\sum q_i=k\}=(1,3,6,7,6,3,1),\ \text{summing to}\ 27=3^3\ [\text{machine}];
\]

%dsp-level:0
\[
\qquad W(k)=\sum_{\sum q_i=k}\prod_im_{q_i}\implies W(0)=8+12=20,\ W(\pm1)=12+3=15\ [\text{machine}];
\]

%dsp-level:0
\[
\qquad 20/15=4/3\ [\text{exact}]
\]

%dsp-level:0
\[
(3):\ \mathrm{D0.8}\ \text{counts internal directed edges};\ \mathrm{D3.4}\ \text{reads two endpoint signs};
\]

%dsp-level:0
\[
\qquad \text{Census FT.35}(2)\text{'s three summands are edge counts}\ [\text{machine}];
\]

%dsp-level:0
\[
\qquad \text{invariance across the four sign pairs}\ [\text{machine}]
\]

%dsp-level:0
\[
(4)(\mathrm i):\ \text{write}\ w_q(R)
\]

%dsp-level:0
\[
= c_qv(R)\ \text{with}\ v\ q\text{-free}: \frac{\sum c_qvm}{\sum c_qv}
\]

%dsp-level:0
\[
= \frac{\sum vm}{\sum v}\ \text{for every}\ c_q
\]

%dsp-level:0
\[
\ne 0\ [\text{machine at}\ c\in\{1,2,3,15,20,\tfrac73\}];
\]

%dsp-level:0
\[
\qquad \Delta S=\tfrac32(2E-|R|)\ \text{is}\ \sigma\text{-free}\ [\mathrm{FT.107}(2)]
\]

%dsp-level:0
\[
(4)(\mathrm{ii}):\ T(q,S)
\]

%dsp-level:0
\[
= W(q)\cdot\overline{\Delta S}(S)\ \text{with the same}\ \overline{\Delta S}\ [(\mathrm i)]
\]

%dsp-level:0
\[
\implies T(0)/T(\pm1)
\]

%dsp-level:0
\[
= W(0)/W(\pm1)
\]

%dsp-level:0
\[
= 4/3\ [\text{machine}];
\]

%dsp-level:0
\[
\qquad 4/3>1\implies\text{the neutral class is heavier}\ [\text{machine}]
\]

%dsp-level:0
\[
(4)(\mathrm{iii}):\ \mathrm{FT.106}(3)\text{'s amplitude integers re-cited};
\]

%dsp-level:0
\[
\qquad 5336/1196=58/13\ \text{and}\ 4796/802=2398/401\ [\text{machine: exact reduction}];
\]

%dsp-level:0
\[
\qquad 58/13\ne4/3\ne2\ [\text{machine}]
\]

%dsp-level:0
\[
(5):\ W(+1)=W(-1)=15\ [\text{machine}];
\]

%dsp-level:0
\[
\qquad \mathrm{FT.106}(3\mathrm d)\text{'s}\ \omega\text{-symmetry re-cited};
\]

%dsp-level:0
\[
\qquad \text{CPT}: q_i\mapsto-q_i\ \text{gives}\ Q\mapsto-Q\ \text{and a bijection of triple sets};\ \text{fixed triples}: \{(0,0,0)\}\ [\text{machine}]
\]

%dsp-level:0
\[
(6):\ \text{set arithmetic on the three lists}\ [\text{machine}];
\]

%dsp-level:0
\[
\qquad (7):\ Q\in\mathbb Z\ \text{as a sum of integers, range}\ [-3,3]\ [\text{machine}];
\]

%dsp-level:0
\[
\qquad \tfrac13\notin\mathbb Z\ [\text{exact}]
\]

%dsp-level:0
\[
(8):\ \text{Census FT.76's priced-variant pattern re-cited};
\]

%dsp-level:0
\[
\qquad \mathrm{P1}/\mathrm{P5}\ \text{scan of this tag's text}\ [\text{machine}].\ \blacksquare
\]

\textbf{Phys.} The charge constraint does exactly what a constraint on a sum does: it changes how many microscopic sign assignments land in each class. For one chain the neutral read has two realizations and each charged read one; for three chains bound by $q_1+q_2+q_3=Q$ the full census is $1,6,15,20,15,6,1$ across $Q=-3$ to $+3$, sixty-four assignments in all, and the neutral class ends up with twenty against the unit class's fifteen. That ratio, four thirds, is the whole content of the restriction. Whether it reaches the mass depends on what functional the mass is. A normalized ensemble average cannot see it — a uniform rescaling of the measure cancels between numerator and denominator, which is why the sector mass table is charge-degenerate as a matter of structure rather than luck. A total deficit can see it, and does: on that route the neutral class carries exactly four thirds the information content of either charged class, and is heavier. The amplitude layer carries a third ratio, fifty-eight thirteenths at zero winding and a different one at winding one, and its bridge to a mass read is still an open item. Two things survive every route: the proton and the antiproton stay degenerate, because the census is symmetric under charge conjugation and $(0,0,0)$ is the only self-conjugate triple; and no fractional charge can appear at all, since an integer sum of ternary reads cannot land on a third. What the corpus does not do is pick a route for the reader: all three are exact, all three are priced, and the choice is a read specification.

\end{theorem}

\medskip\noindent\textit{\texorpdfstring{BM.29 — The status of the CPT sectors: adjudication and scope}{BM.29 - The status of the CPT sectors: adjudication and scope}.}\ Structural-constraints record: see Appendix B, Section B.1.18 of the Addendum.\par\medskip

\phantomsection
\section[Tier BM.8 — The Additive Charge Read and the Mass-Route…]{\texorpdfstring{Tier BM.8 — The Additive Charge Read and the Mass-Route Adjudication}{Tier BM.8 - The Additive Charge Read and the Mass-Route Adjudication}}

\begin{definition}[{\texorpdfstring{Def-BM.8 — The additive charge read (the construction BM.6(1) specifies)}{Def-BM.8 - The additive charge read (the construction BM.6(1) specifies)}}]\label{tag:Def-BM-8}
\textbf{Def.} A read $Q$ on sealed composites, constructed to BM.6(1)'s three specification clauses; no new axiom, no new datum.

%dsp-level:0
\[
(1)\quad\textbf{the read:}
\]

%dsp-level:0
\[
\qquad Q(c_1,\dots,c_k)\ :=\ \sum_{i=1}^{k}q_{\mathrm{core}}(c_i),
\]

%dsp-level:0
\[
\qquad q_{\mathrm{core}}(c)=\tfrac12\bigl(\sigma_{s(c)}-\sigma_{t(c)}\bigr)\in\{-1,0,+1\}
\]

%dsp-level:0
\[
\qquad [\mathrm{D3.4},\ \mathrm{EM.9},\ \mathrm{EM.9a},\ \mathrm{FT.40}];
\]

%dsp-level:0
\[
(2)\quad\textbf{spec clause (i) — restriction:}\qquad Q|_{\text{single chain}}=q_{\mathrm{core}}\qquad[\text{immediate at}\ k=1];
\]

%dsp-level:0
\[
(3)\quad\textbf{spec clause (ii) — additivity:}
\]

%dsp-level:0
\[
\qquad Q(\mathcal X\sqcup\mathcal Y)=Q(\mathcal X)+Q(\mathcal Y)\ \ \text{for disjoint sealings}
\]

%dsp-level:0
\[
\qquad [\text{the sum is over disjoint constituent sets}];
\]

%dsp-level:0
\[
(4)\quad\textbf{spec clause (iii) — observability on sealed supports:}
\]

%dsp-level:0
\[
\qquad Q\ \text{is defined only where}\ \mathcal C(\Sigma;C)=1\ [\mathrm{D2.14}:\ \partial_{\mathrm{col}}(\Sigma;C)=\varnothing];
\]

%dsp-level:0
\[
\qquad \text{one dangling attachment zeroes the whole colour sector}\ [\mathrm{FT.19}]\ \text{and the read with it};
\]

%dsp-level:0
\[
(5)\quad\textbf{the derived range and parity:}
\]

%dsp-level:0
\[
\qquad k=3:\quad Q\in\{-3,-2,-1,0,+1,+2,+3\}\subset\mathbb Z
\]

%dsp-level:0
\[
\qquad [\text{machine}];
\]

%dsp-level:0
\[
\qquad k=2\ (\delta\text{-seal}):\ Q\in\{-2,\dots,+2\};
\]

%dsp-level:0
\[
\textbf{integer-valued at every}\ k:\qquad\tfrac13,\ \tfrac23\ \text{are unreachable}\ [\mathrm{BM.28}(7)];
\]

%dsp-level:0
\[
(6)\quad\textbf{the probability-space weights it induces:}
\]

%dsp-level:0
\[
\qquad W(Q):=\#\{(\sigma\text{-assignments}):\ \textstyle\sum_iq_i=Q\}=\bigl[z^{Q}\bigr](2+z+z^{-1})^{k}
\]

%dsp-level:0
\[
\qquad [\mathrm{BM.28}(2);\ \text{machine}]:
\]

%dsp-level:0
\[
k=3:\quad W=(1,6,15,20,15,6,1)\ \text{over}\ Q=-3,\dots,+3,\qquad\sum_QW(Q)=4^3=64\qquad[\text{machine}];
\]

%dsp-level:0
\[
(7)\quad\textbf{the CPT action:}
\]

%dsp-level:0
\[
\qquad \mathrm{CPT}:\ q_i\mapsto-q_i\ [\mathrm{EM.24}(1)]\implies Q\mapsto-Q;
\]

%dsp-level:0
\[
\qquad W(Q)=W(-Q)\ [\text{machine}];
\]

%dsp-level:0
\[
\qquad \text{the unique CPT-fixed triple at}\ k=3\ \text{is}\ (0,0,0)\ [\text{machine}];
\]

%dsp-level:0
\[
(8)\quad\textbf{status and scope:}
\]

%dsp-level:0
\[
\qquad \textbf{constructed}\ —\ \text{BM.6}(1)\text{'s specification is discharged for the sealed classes}\ k\le3;
\]

%dsp-level:0
\[
\qquad \text{the charge-class labels of}\ \mathrm{Def\text{-}BM.3}(4)\ \text{now rest on a derived read rather than on a declaration};
\]

%dsp-level:0
\[
\neg\text{constructed: the coupling of}\ Q\ \text{to a gauge field}\ [\text{deficit}\ 1',\ \mathrm{FT.125}(4);\ \mathrm{FT.82}(7)(\mathrm a)],
\]

%dsp-level:0
\[
\ \text{and any dynamics of}\ Q\ [\mathrm{BM.6}(5)].
\]

\textbf{Rmk (what the construction does not do).} $Q$ is a sum of the ternary end reads. It supplies no fractional constituent alphabet, no gauge coupling, and no mass: the mass routes are BM.28(4)'s and BM.30's, and the read-choice among them is not made here.

\textbf{Req.} D2.14, D3.4, D3.9 (Parts 1–2); EM.9, EM.9a, EM.24(1) (Part 2); P2 (Part 1); Def-BM.3(4), BM.6(1), BM.28(2), BM.28(7) (this part); FT.19, FT.40, FT.82(7), FT.125(4) (Part 3); generating-function convolution, exhaustive end-sign enumeration [math; machine: \texttt{checks/\allowbreak{}verify\_\allowbreak{}H\_\allowbreak{}probability\_\allowbreak{}space\_\allowbreak{}restriction.py} §A, §D].

\textbf{Pf.}

%dsp-level:0
\[
(1)\text{–}(4):\ \text{the three specification clauses of BM.6}(1)\ \text{discharged clause by clause};
\]

%dsp-level:0
\[
\qquad \text{observability}: \mathcal C(\Sigma;C)\ \text{is the}\ \Pi_{\mathcal C}\ \text{projector's support}\ [\mathrm{D2.14}];
\]

%dsp-level:0
\[
\ \mathrm{FT.19}\text{'s all-or-nothing row re-cited}
\]

%dsp-level:0
\[
(5):\ \text{a sum of}\ k\ \text{integers each in}\ \{-1,0,1\}\ \text{lies in}\ \{-k,\dots,k\}\ [\text{machine at}\ k=2,3];
\]

%dsp-level:0
\[
\qquad (6):\ (2+z+z^{-1})^3\ \text{convolved}\ [\text{machine}];
\]

%dsp-level:0
\[
\qquad \textstyle\sum_QW(Q)=(2+1+1)^3=64\ [\text{machine}]
\]

%dsp-level:0
\[
(7):\ q_i
\]

%dsp-level:0
\[
\mapsto -q_i\ \text{componentwise}\ [\mathrm{EM.24}(1)]
\]

%dsp-level:0
\[
\implies Q
\]

%dsp-level:0
\[
\mapsto -Q\ \text{and}\ W(Q)
\]

%dsp-level:0
\[
= W(-Q)\ \text{by the census's symmetry}\ [\text{machine}];
\]

%dsp-level:0
\[
\qquad \text{fixed triples}: -q_i=q_i\ \forall i\iff q_i=0\ \forall i\ [\text{machine}].\ \blacksquare
\]

\textbf{Phys.} The additive charge read the era listed as an open construction is now built, and it costs nothing: sum the ternary endpoint reads of the constituents, restrict to sealed supports, and the three clauses the specification demanded are satisfied by inspection. Two things come out of it that the specification did not ask for. The read is integer-valued at every constituent number, so a fractional charge is not missing from this corpus but unreachable by it. And the read carries its own statistics: the number of sign assignments realizing each net charge is the coefficient of a cubic generating function, giving twenty for the neutral class against fifteen for either unit class, which is the exact measure of how differently the charge constraint filters the two probability spaces.

\end{definition}

\begin{theorem}[{\texorpdfstring{BM.30 — The mass-route adjudication: three exact ratios, one read choice}{BM.30 - The mass-route adjudication: three exact ratios, one read choice}}]\label{tag:BM-30}
\textbf{Thm.}

%dsp-level:0
\[
(1)\quad\textbf{the three routes and their exact charge ratios:}\qquad\begin{array}{l|c|c}\text{route} & \text{functional} & \text{neutral}:\text{charged}\\\hline(\mathrm i)\ \text{normalized ensemble average}\ [\mathrm{FT.107}(3)] & \dfrac{\sum_Rw\,\Delta S}{\sum_Rw} & 1\ \ (\text{degenerate})\\[2mm](\mathrm{ii})\ \text{total deficit}\ [I_{\mathrm{hid}}\text{-type};\ \mathrm{FT.42}(1)] & \sum_Rw\,\Delta S & \tfrac43\\[1mm](\mathrm{iii})\ \text{amplitude layer}\ [\mathrm{Def\text{-}H.32}(2);\ \mathrm{FT.106}(3)] & \sum_Rv_\Gamma(R)\,\Delta S & \tfrac{58}{13}\ (S{=}0),\ \tfrac{2398}{401}\ (S{=}1)\end{array}
\]

%dsp-level:0
\[
\text{all three exact}\ [\text{machine}];\qquad\text{the three ratios are pairwise unequal}\ [\text{machine}];
\]

%dsp-level:0
\[
(2)\quad\textbf{the splitting, where it exists:}
\]

%dsp-level:0
\[
\qquad \text{on route}\ (\mathrm{ii}):
\]

%dsp-level:0
\[
\qquad \frac{\langle m\rangle^{(n)}}{\langle m\rangle^{(p)}}=\frac{W(0)}{W(+1)}=\frac{20}{15}=\frac43,
\]

%dsp-level:0
\[
\qquad \Delta m\ne0,
\]

%dsp-level:0
\[
\qquad \textbf{the neutral class is heavier}\ [\mathrm{BM.28}(4)(\mathrm{ii});\ \text{machine}];
\]

%dsp-level:0
\[
\text{the mechanism is the measure, not the deficit:}\ \Delta S
\]

%dsp-level:0
\[
= \tfrac32(2E-|R|)\ \text{is}\ \sigma\text{-free}\ [\mathrm{FT.107}(2)],
\]

%dsp-level:0
\[
\ \text{so the restriction rescales without reshaping};
\]

%dsp-level:0
\[
(3)\quad\textbf{the degeneracy, where it is structural:}
\]

%dsp-level:0
\[
\qquad \text{on route}\ (\mathrm i)\ \text{any uniform sector factor cancels in the ratio}\ [\text{machine: factors}\ 1,2,3,15,20,\tfrac73]:
\]

%dsp-level:0
\[
\qquad \textbf{charge-degeneracy is a theorem about normalized averages},\ \neg\text{an assumption about composites};
\]

%dsp-level:0
\[
(4)\quad\textbf{what survives every route:}
\]

%dsp-level:0
\[
\qquad W(+1)=W(-1)\ \text{and}\ Z^{\bullet}(+1,S)=Z^{\bullet}(-1,S)\ [\mathrm{BM.28}(5);\ \mathrm{FT.106}(3\mathrm d)]:
\]

%dsp-level:0
\[
\qquad \textbf{the proton and the antiproton are degenerate on all three routes};
\]

%dsp-level:0
\[
\text{the crystallized-count reads are charge-blind}\ [\mathrm{BM.28}(1)]:
\]

%dsp-level:0
\[
\qquad E_{\mathrm{int}},\ \Delta x_{\mathrm B},\ V_{\mathrm{eff}},\ \rho,\ \kappa^2\rho,\ r_s,\ \text{both locks are shared by all three classes};
\]

%dsp-level:0
\[
(5)\quad\textbf{the read-choice row:}
\]

%dsp-level:0
\[
\qquad \neg\exists\ \text{a theorem of Parts 1–4 selecting one route as ``the mass''}\ [\text{inventory: the}\ m_{\mathrm{hid}},\ \text{pole, sector-expectation and}\ m^2_B\ \text{reads are distinct}\ [\mathrm{FT.84}(1);\ \mathrm{FT.87}(8)]];
\]

%dsp-level:0
\[
\text{adoption requires a logged declaration}\ [\text{Census FT.76's priced-variant pattern};\ \mathrm{BM.6}(4)\ \text{precedent}];
\]

%dsp-level:0
\[
\qquad \textbf{no route is adopted here};
\]

%dsp-level:0
\[
(6)\quad\textbf{the ordering is derived and uncompared:}
\]

%dsp-level:0
\[
\qquad \langle m\rangle^{(n)}>\langle m\rangle^{(p)}\ \text{on route}\ (\mathrm{ii})\ [\text{machine:}\ \tfrac43>1];
\]

%dsp-level:0
\[
\qquad \neg\exists\ \text{a measured mass value in any clause and}\ \neg\exists\ \text{a comparison}\ [\mathrm{P5}];
\]

%dsp-level:0
\[
(7)\quad\textbf{the status of the adjudication:}
\]

%dsp-level:0
\[
\qquad \textbf{derived:}\ (1)\text{–}(4),\ (6);
\]

%dsp-level:0
\[
\qquad \textbf{open:}\ \text{the route choice}\ (5),\ \text{the amplitude}\leftrightarrow\text{mass bridge}\,\ \text{the gauge coupling of}\ Q\;
\]

%dsp-level:0
\[
\textbf{closed by this tier:}\ \mathrm{BM.6}(1)\text{'s construction spec}\ [\mathrm{Def\text{-}BM.8}(8)];
\]

%dsp-level:0
\[
\qquad \textbf{unclosable by this read:}\ \text{the fractional constituent alphabet}\ [\mathrm{BM.28}(7);\ \mathrm{Def\text{-}BM.8}(5)].
\]

\textbf{Rmk (no route is privileged).} Clause (2)'s splitting and clause (3)'s degeneracy are both exact and neither is a correction of the other: they are different functionals of one ensemble. A corpus that adopted route (ii) would have a neutron heavier than its proton by the factor four thirds; a corpus that adopted route (i) would have them equal; both statements are machine-verified, and the choice between them is the read specification that clause (5) declines to make.

\textbf{Req.} D2.9 (Parts 1–2); CT2.1 (Part 2); P2, P5 (Parts 1–2); Def-H.32, FT.42(1), FT.84(1), FT.87(8), FT.106(3), FT.107(2), FT.107(3), FT.125(4) (Part 3); Def-BM.8, BM.6(1), BM.6(4), BM.28(1), BM.28(2), BM.28(4), BM.28(5), BM.28(7) (this part); Census FT.76 (the priced-variant pattern); exact rational ratio arithmetic, exhaustive enumeration [math; machine: \texttt{checks/\allowbreak{}verify\_\allowbreak{}H\_\allowbreak{}probability\_\allowbreak{}space\_\allowbreak{}restriction.py} §B–§D].

\textbf{Pf.}

%dsp-level:0
\[
(1):\ \mathrm{BM.28}(4)\text{'s three rows re-cited with their machine values};
\]

%dsp-level:0
\[
\qquad \tfrac{58}{13}\ne\tfrac43\ \text{since}\ 58\cdot3=174\ne13\cdot4=52;
\]

%dsp-level:0
\[
\qquad \tfrac43\ne1\ [\text{machine}]
\]

%dsp-level:0
\[
(2):\ T(q,S)
\]

%dsp-level:0
\[
= W(q)\overline{\Delta S}(S)\ \text{with a common}\ \overline{\Delta S}\ [\mathrm{BM.28}(4)(\mathrm i),(\mathrm{ii})]\implies\text{the ratio is}\ W(0)/W(+1)
\]

%dsp-level:0
\[
= 4/3\ [\text{machine}]
\]

%dsp-level:0
\[
(3):\ \text{the cancellation lemma}\ [\mathrm{BM.28}(4)(\mathrm i);\ \text{machine}]
\]

%dsp-level:0
\[
(4):\ \mathrm{BM.28}(5)\ \text{and}\ \mathrm{BM.28}(1)\ \text{re-cited}
\]

%dsp-level:0
\[
(5):\ \text{inventory scan of the mass reads}\ [\mathrm{FT.84}(1)\text{'s read-class taxonomy};\ \mathrm{FT.87}(8)];
\]

%dsp-level:0
\[
\qquad \text{Census FT.76's variant-pricing pattern re-cited}
\]

%dsp-level:0
\[
(6):\ \tfrac43>1\ [\text{machine}];\qquad\mathrm{P5}\ \text{scan of this tag's text}\ [\text{machine: zero external numeric tokens}]
\]

%dsp-level:0
\[
(7):\ \text{conjunction of}\ (1)\text{–}(6)\ \text{with}\ \mathrm{Def\text{-}BM.8}(8).\ \blacksquare
\]

\textbf{Phys.} One ensemble, three functionals, three answers — and all three are exact. Normalize and the charge constraint disappears, because a uniform rescaling of the measure cancels between numerator and denominator; that is why the sector mass table is charge-degenerate, and it is a structural fact rather than a lucky one. Sum instead of averaging and the constraint is all that matters: twenty sign assignments realize the neutral class against fifteen for either charged class, so the neutral composite carries four thirds the information deficit and weighs more. Take the amplitude layer and a third ratio appears, fifty-eight thirteenths at zero winding and something else at winding one, with its bridge to a mass read still open. What no route disturbs is the mirror: the census is symmetric under charge conjugation, so the proton and antiproton stay degenerate everywhere, and the seven crystallized reads stay blind to charge entirely. The corpus prices the three routes and adopts none, because adopting one is a read specification and not a theorem — and it records the ordering it does derive, neutral heavier, without comparing it to anything.

\end{theorem}

\phantomsection
\section[Tier BM.9 — The Mesonic Flicker Test: the Contraction…]{\texorpdfstring{Tier BM.9 — The Mesonic Flicker Test: the Contraction Grammar, the Colour Projector, and the Protected Junction}{Tier BM.9 - The Mesonic Flicker Test: the Contraction Grammar, the Colour Projector, and the Protected Junction}}

\begin{definition}[{\texorpdfstring{Def-BM.9 — The extended flicker ensemble and the sealing filter}{Def-BM.9 - The extended flicker ensemble and the sealing filter}}]\label{tag:Def-BM-9}
\textbf{Def.} A construction class on Def-BM.5 (no new axiom; Def-H.32's scopes (S1)–(S4) and Def-BM.5's inherited clauses stand).

%dsp-level:0
\[
(1)\quad\textbf{the two restricted ensembles:}
\]

%dsp-level:0
\[
\qquad \mathcal E_\varepsilon(\mathcal X):=\{R\in\mathcal E(\mathcal X):\ \text{the sealing of}\ R\ \text{is}\ \varepsilon_{ijk}\text{-type}\},
\]

%dsp-level:0
\[
\qquad \mathcal E_\delta(\mathcal X):=\{R\in\mathcal E(\mathcal X):\ \delta^i{}_j\text{-type}\}
\]

%dsp-level:0
\[
\qquad [\mathrm{D3.9};\ \mathrm{CT3.17}];
\]

%dsp-level:0
\[
\mathcal E_{\mathrm{ext}}(\mathcal X):=\mathcal E_\varepsilon(\mathcal X)\cup\mathcal E_\delta(\mathcal X);
\]

%dsp-level:0
\[
\qquad \textbf{constitutive clause:}\ \text{the union is taken inside}\ \mathrm{Def\text{-}H.32}(1)\text{'s topology space, which is colour-blind}\ [\mathrm{D2.13}:\ \mathcal M(\Sigma)\ \text{is subject to the rung masks and Janus balance},\ \neg\text{to colour sealing}];
\]

%dsp-level:0
\[
(2)\quad\textbf{the filter clause (the load-bearing one):}\qquad\text{colour sealing is ambient-global}\ [\mathrm{CT3.18}]\ \text{and acts}\ \textbf{downstream}\ \text{as the projector}\ \Pi_{\mathcal C}(\Sigma;C)=\mathcal C(\Sigma;C)\cdot\mathbb I^{(\mathrm{strong})}_\Sigma,\qquad\mathcal C(\Sigma;C)=\begin{cases}1&\partial_{\mathrm{col}}(\Sigma;C)=\varnothing\\0&\text{otherwise}\end{cases}\qquad[\mathrm{D2.14}];
\]

%dsp-level:0
\[
\text{one dangling attachment zeroes the entire colour block}\ [\mathrm{FT.19}:\ \text{all-or-nothing}];
\]

%dsp-level:0
\[
\qquad \mathcal C\ \text{is an idempotent scalar},\ \mathcal C^2=\mathcal C\in\{0,1\}\ [\text{machine}];
\]

%dsp-level:0
\[
\textbf{non-derived clause:}\ \ \text{none}\ —\ (2)\ \text{is}\ \mathrm{D2.14}\text{'s derived projector, cited at the site where the union of}\ (1)\ \text{would otherwise be read as a sum over sealings};
\]

%dsp-level:0
\[
(3)\quad\textbf{the grouping grammar:}
\]

%dsp-level:0
\[
\qquad g(N):=\#\{\text{partitions of}\ N\ \text{attachments into contraction groups of size}\ 2\ (\delta)\ \text{and}\ 3\ (\varepsilon)\}
\]

%dsp-level:0
\[
\qquad [\mathrm{CT3.17a}(2);\ \mathrm{FT.41}(\mathrm b)]:
\]

%dsp-level:0
\[
g(N)=(1,\,1,\,3,\,10,\,25,\,105)\quad\text{for}\ N=2,\dots,7
\]

%dsp-level:0
\[
\qquad [\text{machine: exhaustive partition enumeration reproduces all six values}];
\]

%dsp-level:0
\[
g_\delta(N):=\#\{\text{groupings containing at least one}\ \delta\}=(1,\,\mathbf0,\,3,\,10,\,15,\,105),
\]

%dsp-level:0
\[
\qquad g_\varepsilon(N):=\#\{\text{groupings of}\ 3\text{-blocks only}\}=(0,\,1,\,0,\,0,\,10,\,0);
\]

%dsp-level:0
\[
g_\delta(N)+g_\varepsilon(N)=g(N)\ \ \forall N\ [\text{machine}];\qquad g_\varepsilon(N)\ne0\iff3\mid N\ [\text{machine on the frozen range}];
\]

%dsp-level:0
\[
(4)\quad\textbf{the boundary-data tuple:}
\]

%dsp-level:0
\[
\qquad \text{the labels that condition}\ \mathcal E(\mathcal X)\ \text{are}\ \bigl(Q,\ \Theta_{\mathrm{int}}\ \text{class},\ S,\ \text{attachment count}\ N\bigr)\ [\mathrm{Def\text{-}BM.5}(1);\ \mathrm{Def\text{-}BM.8}]:
\]

%dsp-level:0
\[
\textbf{the attachment count is a boundary datum},
\]

%dsp-level:0
\[
\ \neg\text{a summed-over structure}\ —\ \text{a}\ 3\text{-junction and a}\ 2\text{-junction are different supports, not two members of one ensemble};
\]

%dsp-level:0
\[
(5)\quad\textbf{the scope clause:}
\]

%dsp-level:0
\[
\qquad \text{no clause of this part sums over sealings of a fixed support as if they were alternatives}: \text{the grammar of}\ (3)\ \text{decides which sealings a support admits at all}.
\]

\textbf{Rmk (why the filter, not a cancellation).} The topology-blind path integral is colour-blind by construction — that is what "topology-blind" means here [D2.13's scope note; CT3.22a Rmk]. Colour enters only through $\Pi_{\mathcal C}$, a $0/1$ projector applied after the sum. So the question "do $\delta$-sealed states cancel in the baryon's integral?" has a category answer before it has an arithmetic one: they are not terms that could cancel, they are terms the projector never admits.

\textbf{Req.} D2.13, D2.14, D3.9 (Parts 1–2); CT3.17, CT3.17a(2), CT3.18, CT3.22a (Part 2); Def-H.32 (Part 3); FT.19, FT.41(b) (Part 3); Def-BM.5, Def-BM.8 (this part); exhaustive set-partition enumeration over blocks of size 2 and 3, exact integer arithmetic [math; machine: \texttt{checks/\allowbreak{}verify\_\allowbreak{}H\_\allowbreak{}mesonic\_\allowbreak{}flicker\_\allowbreak{}test.py} §A].

\textbf{Pf.}

%dsp-level:0
\[
\text{Definition block: the constitutive clause is}\ (1)\text{'s union inside the colour-blind topology space};
\]

%dsp-level:0
\[
\ (2)\ \text{is}\ \mathrm{D2.14}\text{'s projector re-cited at the site where the union is formed};
\]

%dsp-level:0
\[
\ (3)\text{'s counts are machine-reproduced from}\ \mathrm{CT3.17a}(2)\text{'s frozen list};
\]

%dsp-level:0
\[
\ (4)\text{–}(5)\ \text{are scope clauses}.\ \blacksquare
\]

\textbf{Phys.} Before any sum is taken, one structural fact settles the question. The topology-blind path integral does not filter by colour — it filters by topology, and colour enters afterwards as a zero-one projector on the support. A sealing is therefore not an alternative term competing inside one integral; it is a property of the support that either passes the projector or annihilates the whole colour block. And the grammar that decides it is small enough to enumerate: partition the attachments into groups of two and three, and count.

\end{definition}

\begin{theorem}[{\texorpdfstring{BM.31 — The boundary-data comparison and the grammar verdict}{BM.31 - The boundary-data comparison and the grammar verdict}}]\label{tag:BM-31}
\textbf{Thm.}

%dsp-level:0
\[
(1)\quad\textbf{the grammar at}\ N=3\ \textbf{is forced:}
\]

%dsp-level:0
\[
\qquad g(3)=1,
\]

%dsp-level:0
\[
\qquad \text{the unique grouping is the}\ \varepsilon\text{-triple},
\]

%dsp-level:0
\[
\qquad g_\delta(3)=\mathbf0
\]

%dsp-level:0
\[
\qquad [\mathrm{Def\text{-}BM.9}(3);\ \text{machine: exhaustive}];
\]

%dsp-level:0
\[
\text{a}\ \delta\ \text{sub-pairing on a}\ 3\text{-attachment support consumes two attachments and leaves exactly one dangling}\ (3-2=1):
\]

%dsp-level:0
\[
\qquad \partial_{\mathrm{col}}\ne\varnothing\implies\mathcal C=0\implies\text{the colour block vanishes}\ [\mathrm{D2.14};\ \mathrm{FT.19}];
\]

%dsp-level:0
\[
\textbf{there is no}\ \delta\text{-sealed member of a}\ 3\text{-junction's ensemble to cancel}: \mathcal E_\delta(\mathcal X_{N=3})=\varnothing;
\]

%dsp-level:0
\[
(2)\quad\textbf{the protection is unique to}\ N=3:
\]

%dsp-level:0
\[
\qquad g_\delta(N)=(1,0,3,10,15,105)\ \text{for}\ N=2,\dots,7\ [\text{machine}]:
\]

%dsp-level:0
\[
\qquad g_\delta(N)>0\ \ \forall N\ge4;
\]

%dsp-level:0
\[
g_\varepsilon(N)\ne0\iff3\mid N:\qquad g_\varepsilon(3)=1,\quad g_\varepsilon(6)=10,\quad g_\varepsilon(N)=0\ \text{at}\ N\in\{2,4,5,7\}\ [\text{machine}];
\]

%dsp-level:0
\[
\text{at}\ N
\]

%dsp-level:0
\[
= 6\ \text{the ensemble genuinely mixes:}\ \tfrac{15}{25}
\]

%dsp-level:0
\[
= \tfrac35\ \text{of the groupings contain a}\ \delta,
\]

%dsp-level:0
\[
\ \tfrac{10}{25}
\]

%dsp-level:0
\[
= \tfrac25\ \text{do not}\ [\text{machine}];
\]

%dsp-level:0
\[
\qquad \text{at}\ N=4\ \text{every grouping does}\ (g_\delta=g=3);
\]

%dsp-level:0
\[
(3)\quad\textbf{the charge classes overlap, so charge does not separate the sealings:}
\]

%dsp-level:0
\[
\qquad Q_{N=3}\in\{-3,\dots,+3\},\quad W_3=(1,6,15,20,15,6,1),\quad\sum W_3=64=4^3\ [\mathrm{Def\text{-}BM.8}(6);\ \text{machine}];
\]

%dsp-level:0
\[
Q_{N=2}\in\{-2,\dots,+2\},\quad W_2=(1,4,6,4,1),\quad\sum W_2=16=4^2\ [\text{machine}];
\]

%dsp-level:0
\[
\text{the overlap is}\ \{-2,-1,0,+1,+2\}:
\]

%dsp-level:0
\[
\qquad \textbf{at}\ Q
\]

%dsp-level:0
\[
= 0\ \text{and}\ Q
\]

%dsp-level:0
\[
= +1\ \text{the two sealings share a charge class}\ —\ \text{the neutron and proton classes are not separated from mesonic ones by charge};
\]

%dsp-level:0
\[
\text{only}\ |Q|=3\ \text{is baryon-exclusive}\ [\text{machine}];
\]

%dsp-level:0
\[
\qquad W_3(Q)
\]

%dsp-level:0
\[
\ne W_2(Q)\ \text{on the whole overlap}\ [\text{machine}]:\ \text{the probability spaces differ even where the classes coincide};
\]

%dsp-level:0
\[
(4)\quad\textbf{the phase class separates only partially:}
\]

%dsp-level:0
\[
\qquad \text{baryon-type}\ \Theta_{\mathrm{int}}=2\Theta_\triangle=4k_\triangle\in4\mathbb Z\ [\text{Census FT.49}(3);\ \mathrm{FT.22}];
\]

%dsp-level:0
\[
\qquad \text{per materialized closed chain}\ \Theta(c)\in2\mathbb Z\ [\mathrm{FT.22}];
\]

%dsp-level:0
\[
\qquad \Theta_{\mathrm{int}}=2k(\Gamma)\in2\mathbb Z\ [\mathrm{FT.34}(3);\ \text{Census FT.35}(3)]:
\]

%dsp-level:0
\[
4\mathbb Z\subsetneq2\mathbb Z\ [\text{machine}]:
\]

%dsp-level:0
\[
\qquad \Theta_{\mathrm{int}}
\]

%dsp-level:0
\[
= 2\ \text{is meson-admissible and baryon-excluded, while}\ \Theta_{\mathrm{int}}
\]

%dsp-level:0
\[
\in 4\mathbb Z\ \text{is admissible to both}\ —\ \textbf{the phase label is a partial separator};
\]

%dsp-level:0
\[
(5)\quad\textbf{the representation content does not separate them either:}
\]

%dsp-level:0
\[
\qquad \dim\bigl((\mathbb C^3)^{\otimes3}\bigr)^{\mathfrak{sl}(3)}=27-26=1,
\]

%dsp-level:0
\[
\qquad \dim\bigl(\mathbb C^3\otimes\bar{\mathbb C}^3\bigr)^{\mathfrak{sl}(3)}=9-8=1
\]

%dsp-level:0
\[
\qquad [\mathrm{BM.1}(3);\ \text{machine: exact rank over}\ \mathbb Q];
\]

%dsp-level:0
\[
\text{both contents carry exactly one singlet line}:
\]

%dsp-level:0
\[
\qquad \textbf{the separator is the attachment count}\ N,
\]

%dsp-level:0
\[
\ \text{carried in the boundary data}\ [\mathrm{Def\text{-}BM.9}(4)],
\]

%dsp-level:0
\[
\ \neg\text{charge},
\]

%dsp-level:0
\[
\ \neg\text{phase alone},
\]

%dsp-level:0
\[
\ \neg\text{representation content};
\]

%dsp-level:0
\[
(6)\quad\textbf{the verdict on the extended ensemble:}
\]

%dsp-level:0
\[
\qquad \mathcal E_{\mathrm{ext}}(\mathcal X_{N=3})=\mathcal E_\varepsilon(\mathcal X_{N=3})\sqcup\varnothing=\mathcal E_\varepsilon(\mathcal X_{N=3})
\]

%dsp-level:0
\[
\qquad [\text{machine: the grammar census}]:
\]

%dsp-level:0
\[
\textbf{the baryon's flicker ensemble is not extended by mesonic members}\ —\ \text{not because they cancel, but because a}\ 3\text{-junction admits no}\ \delta\text{-sealing at all}.
\]

\textbf{Rmk (what would extend it).} Clause (2) is the operative scope: at $N\ge4$ the union is nontrivial and the extended ensemble genuinely mixes sealings. Any composite whose support carries four or more colour attachments — the Census FT.50-type junction gluings [BM.2(3)] — is in that regime, and this tag's protection does not apply to it.

\textbf{Req.} D2.14, D3.9 (Parts 1–2); CT3.17, CT3.17a(2), CT3.18 (Part 2); EM.9 (Part 2); P2 (Part 1); Def-BM.8(6), Def-BM.9 (this part); FT.19, FT.22, FT.34(3), FT.41(b) (Part 3); BM.1(3), BM.2(1), BM.2(3) (this part); Census FT.35(3), FT.49(3) (row-wise); exhaustive set-partition enumeration, exact rational rank over $\mathbb Q$, generating-function convolution [math; machine: \texttt{checks/\allowbreak{}verify\_\allowbreak{}H\_\allowbreak{}mesonic\_\allowbreak{}flicker\_\allowbreak{}test.py} §A–§B].

\textbf{Pf.}

%dsp-level:0
\[
(1):\ \text{partitions of a}\ 3\text{-set into blocks of size}\ \in\{2,3\}:\ \text{the only one is}\ \{1,2,3\}\ [\text{machine}];
\]

%dsp-level:0
\[
\qquad \text{a}\ 2\text{-block leaves}\ 3-2=1\ \text{unpaired}\implies\partial_{\mathrm{col}}\ne\varnothing\ [\mathrm{D2.14}]\implies\mathcal C=0\implies\text{block}\equiv0\ [\mathrm{FT.19}]
\]

%dsp-level:0
\[
(2):\ \text{the censuses}\ g_\delta,\ g_\varepsilon\ \text{by exhaustive enumeration for}\ N=2..7\ [\text{machine}];
\]

%dsp-level:0
\[
\qquad g_\varepsilon(N)\ne0\implies N=3m\ [\text{a partition into}\ 3\text{-blocks requires}\ 3\mid N];
\]

%dsp-level:0
\[
\qquad \tfrac{15}{25}=\tfrac35,\ \tfrac{10}{25}=\tfrac25\ [\text{machine}]
\]

%dsp-level:0
\[
(3):\ (2+z+z^{-1})^3\ \text{and}\ (2+z+z^{-1})^2\ \text{convolved}\ [\mathrm{Def\text{-}BM.8}(6);\ \text{machine}];
\]

%dsp-level:0
\[
\qquad \text{the overlap and the exclusivity of}\ |Q|=3\ [\text{machine: set arithmetic}]
\]

%dsp-level:0
\[
(4):\ \text{Census FT.49}(3)\text{'s}\ 4k_\triangle\ \text{row re-cited};
\]

%dsp-level:0
\[
\qquad \mathrm{FT.22}\text{'s}\ \Theta(c)\in2\mathbb Z\ \text{and}\ \mathrm{FT.34}(3)\text{'s additivity re-cited};
\]

%dsp-level:0
\[
\qquad 4\mathbb Z\subsetneq2\mathbb Z\ \text{with}\ 2\in2\mathbb Z\setminus4\mathbb Z\ [\text{machine}]
\]

%dsp-level:0
\[
(5):\ \text{the common null space of the}\ 8\ \text{traceless generators' actions}: 27-26
\]

%dsp-level:0
\[
= 1\ \text{and}\ 9-8
\]

%dsp-level:0
\[
= 1\ [\text{machine: exact rank over}\ \mathbb Q];
\]

%dsp-level:0
\[
\qquad \mathrm{BM.1}(3)\ \text{re-derived}
\]

%dsp-level:0
\[
(6):\ (1)\text{'s}\ \mathcal E_\delta=\varnothing\ \text{substituted into}\ \mathrm{Def\text{-}BM.9}(1)\text{'s union}.\ \blacksquare
\]

\textbf{Phys.} Three of the four boundary labels fail to tell a baryon from a meson. Their charge classes overlap on everything except $|Q|=3$, so the neutron's neutral class and the proton's unit class are shared with mesonic configurations; their phase classes differ only in half the cases, since $4\mathbb Z$ sits strictly inside $2\mathbb Z$; and their colour representation content is identical in the only respect that matters, one singlet line each. What does separate them is the attachment count, and it separates them absolutely: three attachments admit exactly one contraction grouping, the epsilon triple, and a delta sub-pairing would leave one attachment dangling, which the all-or-nothing law turns into a vanished colour block. So the baryon's flicker ensemble has no mesonic members to cancel. That protection is not general — it is a property of the number three, and from four attachments upward the groupings genuinely mix, fifteen of twenty-five at six.

\end{theorem}

\begin{theorem}[{\texorpdfstring{BM.32 — The mesonic contribution: zero by filtration, and the amplitude layer is orthogonal}{BM.32 - The mesonic contribution: zero by filtration, and the amplitude layer is orthogonal}}]\label{tag:BM-32}
\textbf{Thm.}

%dsp-level:0
\[
(1)\quad\textbf{the contribution vanishes, and the mechanism is the projector:}
\]

%dsp-level:0
\[
\qquad Z^{(n)}_\delta:=\sum_{R\in\mathcal E_\delta(\mathcal X_{N=3})}\sum_\sigma w(R,\sigma)=\sum_{R\in\varnothing}(\cdots)=0
\]

%dsp-level:0
\[
\qquad [\mathrm{BM.31}(1),(6);\ \text{machine}];
\]

%dsp-level:0
\[
Z^{(n)}_{\mathrm{ext}}=Z^{(n)}_\varepsilon+Z^{(n)}_\delta=Z^{(n)}_\varepsilon\qquad[\text{machine}];
\]

%dsp-level:0
\[
\textbf{the zero is a filtration, not a cancellation:}\ \ \neg\exists\ \text{a pair of mesonic terms whose weights sum to zero}\ —\ \neg\exists\ \text{mesonic terms};
\]

%dsp-level:0
\[
(2)\quad\textbf{the projector form of the same statement:}
\]

%dsp-level:0
\[
\qquad \mathcal C(\Sigma;C)\in\{0,1\},\quad\mathcal C^2=\mathcal C\ [\mathrm{D2.14};\ \text{machine}]:
\]

%dsp-level:0
\[
\qquad \Pi_{\mathcal C}\ \text{is an idempotent scalar on the colour slot space};
\]

%dsp-level:0
\[
\text{a}\ \delta\text{-sealed member of a}\ 3\text{-support has}\ \mathcal C=0\implies\Pi_{\mathcal C}=0\implies F^\Sigma_{\mathrm{strong}}=0\ [\mathrm{FT.19}]:
\]

%dsp-level:0
\[
\qquad \text{its weight is multiplied by zero before any sum is formed};
\]

%dsp-level:0
\[
(3)\quad\textbf{both mass routes are untouched:}
\]

%dsp-level:0
\[
\qquad \text{adding zero-weight members leaves a normalized average invariant}\ \bigl(\tfrac{\sum w m}{\sum w}\bigr)\ \textbf{and}\ \text{a total invariant}\ \bigl(\sum wm\bigr)
\]

%dsp-level:0
\[
\qquad [\text{machine: exact on a rational witness ensemble}]:
\]

%dsp-level:0
\[
\langle m\rangle_{\mathrm{ext}}=\langle m\rangle_\varepsilon\ \ \text{on route}\ (\mathrm i),\qquad T_{\mathrm{ext}}=T_\varepsilon\ \ \text{on route}\ (\mathrm{ii})\qquad[\mathrm{BM.28}(4);\ \mathrm{BM.20}(3);\ \text{machine}];
\]

%dsp-level:0
\[
\therefore\ \mathrm{BM.20}\text{'s curvature-flux mass and BM.28's three ratios are}\ \textbf{unchanged}\ \text{by the extension};
\]

%dsp-level:0
\[
(4)\quad\textbf{the amplitude-layer mechanisms are orthogonal to this question:}
\]

%dsp-level:0
\[
\qquad \mathrm{FT.67}(3\mathrm a)\ \text{and}\ \mathrm{FT.106}(2)\ \text{kill}\ \textbf{pure directed cycles}\ (\text{any length, any cell});
\]

%dsp-level:0
\[
\qquad \mathrm{FT.71}\ \text{cancels every odd cycle and every undissolved even cycle};
\]

%dsp-level:0
\[
\qquad \mathrm{FT.69}(3)(\mathrm{iv}),(\mathrm v)\ \text{exclude odd cycles and}\ 4\text{-cycles from covers}:
\]

%dsp-level:0
\[
\text{these are statements about cycles in}\ \Gamma,\ \neg\text{about colour sealings of a support}\ [\mathrm{Def\text{-}BM.9}(2)\ \mathrm{Rmk}];
\]

%dsp-level:0
\[
\qquad \text{they neither produce nor require}\ Z_\delta=0;
\]

%dsp-level:0
\[
(5)\quad\textbf{where mesonic content does survive:}
\]

%dsp-level:0
\[
\qquad Z^{\mathrm{amp}}(q,1)=802\ (q=\pm1),\ 4796\ (q=0):
\]

%dsp-level:0
\[
\qquad \textbf{dissolved digons inside covers}\ [\mathrm{FT.106}(3);\ \text{machine}];
\]

%dsp-level:0
\[
Z^{\mathrm{amp}}(q,S)=0\ \ \forall q,\ S\in\{2,3\}\ [\mathrm{FT.106}(3\mathrm c);\ \text{machine}]:
\]

%dsp-level:0
\[
\qquad \text{a standalone digon}\ —\ \text{a pure directed cycle}\ —\ \text{dies at the amplitude layer};
\]

%dsp-level:0
\[
Z^{\mathrm{amp}}(q,S)=Z^{\mathrm{amp}}(-q,S)\ [\mathrm{FT.106}(3\mathrm d);\ \text{machine}];
\]

%dsp-level:0
\[
\qquad Z^{\mathrm{norm}}(0,S)
\]

%dsp-level:0
\[
= 2Z^{\mathrm{norm}}(\pm1,S)\ \text{at every}\ S\ [\text{machine}]:\ \text{the norm layer keeps its}\ 2{:}1\ \text{ratio where the amplitude layer dies};
\]

%dsp-level:0
\[
(6)\quad\textbf{the two readings reconciled:}
\]

%dsp-level:0
\[
\qquad \text{a meson as a}\ \textbf{sealed composite}\ \text{is excluded from a}\ 3\text{-junction's ensemble by the grammar}\ [(1)];
\]

%dsp-level:0
\[
\qquad \text{mesonic}\ \textbf{chain content}\ \text{enters the amplitude layer only as dissolved digons inside covers}\ [(5)]:
\]

%dsp-level:0
\[
\text{the two statements concern different objects}\ —\ \text{a support's sealing and a cover's cycle structure}\ —\ \text{and neither contradicts the other}.
\]

\textbf{Rmk (no protection claim beyond N = 3).} Clause (3)'s invariance is exact because the added weight is exactly zero. At $N\ge4$ [BM.31(2)] the added weight is nonzero, the average genuinely moves, and no clause of this tag applies.

\textbf{Req.} D2.14 (Parts 1–2); CT2.5, CT2.8, CT3.18 (Part 2); Def-H.32, FT.19, FT.67(2), FT.67(3), FT.69(3), FT.71, FT.106(2), FT.106(3) (Part 3); Def-BM.9, BM.20(3), BM.28(4), BM.31(1), BM.31(2), BM.31(6) (this part); exact rational ensemble algebra, exact integer census arithmetic [math; machine: \texttt{checks/\allowbreak{}verify\_\allowbreak{}H\_\allowbreak{}mesonic\_\allowbreak{}flicker\_\allowbreak{}test.py} §C–§D].

\textbf{Pf.}

%dsp-level:0
\[
(1):\ \mathcal E_\delta=\varnothing\ [\mathrm{BM.31}(1)]\implies\text{the sum is empty}\implies0\ [\text{machine}];\qquad Z_{\mathrm{ext}}=Z_\varepsilon+0\ [\text{machine}]
\]

%dsp-level:0
\[
(2):\ \mathrm{D2.14}\text{'s}\ \mathcal C\ \text{definition re-cited};
\]

%dsp-level:0
\[
\qquad \mathcal C^2=\mathcal C\ \text{at}\ \mathcal C\in\{0,1\}\ [\text{machine}];
\]

%dsp-level:0
\[
\qquad \mathrm{FT.19}\text{'s all-or-nothing row re-cited}
\]

%dsp-level:0
\[
(3):\ \text{adding terms with}\ w
\]

%dsp-level:0
\[
= 0\ \text{changes neither}\ \sum wm\ \text{nor}\ \sum w\implies\text{neither the ratio nor the total}\ [\text{machine: exact Fractions on a}\ 5\text{-member witness ensemble}]
\]

%dsp-level:0
\[
(4):\ \mathrm{FT.67}(3\mathrm a),\ \mathrm{FT.71},\ \mathrm{FT.69}(3)\ \text{rows re-cited};
\]

%dsp-level:0
\[
\qquad \text{their objects are cycles of}\ \Gamma\ [\mathrm{Def\text{-}H.32}(1)\text{'s topology space}],
\]

%dsp-level:0
\[
\ \text{while}\ \mathcal C\ \text{acts on colour attachments}\ [\mathrm{D2.14}]
\]

%dsp-level:0
\[
(5):\ \mathrm{FT.106}(3)\text{'s census integers re-cited and their relations verified exactly}\ [\text{machine}:\ 802,\ 4796,\ \text{the zeros at}\ S\in\{2,3\},\ \text{the}\ \omega\text{-symmetry, the}\ 2{:}1\ \text{norm ratio at all four windings}]
\]

%dsp-level:0
\[
(6):\ (1)\ \text{and}\ (5)\ \text{conjoined with}\ \mathrm{Def\text{-}BM.9}(2)\text{'s locus distinction}.\ \blacksquare
\]

\textbf{Phys.} The answer to the test is zero, and the reason is more interesting than the number. Nothing cancels: there is nothing to cancel. A three-attachment support admits exactly one contraction, so the mesonic subset of its ensemble is empty before any weight is written down, and the colour projector would have annihilated those terms individually anyway — one dangling attachment and the whole block goes. Because the added weight is exactly zero, both mass functionals stay put: the normalized average and the total deficit are each invariant, so the curvature-flux mass and the four-thirds ratio of the previous tier survive the extension untouched. The amplitude-layer cancellations the corpus already proved are about a different object entirely — directed cycles in the skeleton — and they cut the other way: a standalone digon dies there, while dissolved digons inside covers are precisely what survives, with exact census weights. Mesons as sealed composites and mesonic chain content in covers are two different things, and the tier keeps them apart.

\end{theorem}

\medskip\noindent\textit{\texorpdfstring{BM.33 — The status of the mesonic flicker test}{BM.33 - The status of the mesonic flicker test}.}\ Structural-constraints record: see Appendix B, Section B.1.19 of the Addendum.\par\medskip

\phantomsection
\section[Tier BM.10 — The Kinetic Energy from the Metric and the…]{\texorpdfstring{Tier BM.10 — The Kinetic Energy from the Metric and the Vacuum-Balanced Entropy: the Dispersion Read of a Sealed Composite, the Deficit Reframing of the Rest Mass, the Force-Structure Ladder, and the Route-Artifact Adjudication}{Tier BM.10 - The Kinetic Energy from the Metric and the Vacuum-Balanced Entropy: the Dispersion Read of a Sealed Composite, the Deficit Reframing of the Rest Mass, the Force-Structure Ladder, and the Route-Artifact Adjudication}}

\begin{definition}[{\texorpdfstring{Def-BM.10 — The metric-energy read of a sealed composite (the dispersion route)}{Def-BM.10 - The metric-energy read of a sealed composite (the dispersion route)}}]\label{tag:Def-BM-10}
\textbf{Def.} A read constructed on the chart side of a sealed composite from CT1.12's dictionary; no new axiom, no new datum, no new parameter. Local notation (internal bookkeeping): $\gamma_{ij}:=g_{ij}$ (the spatial block), $\gamma^{ij}:=$ the inverse matrix of the $3\times3$ block ($\gamma^{ik}g_{kj}=\delta^i{}_j$), $g^{\mu\nu}:=$ the inverse of the full $4\times4$ metric.

%dsp-level:0
\[
(1)\quad\textbf{locus and hypotheses:}
\]

%dsp-level:0
\[
\qquad \text{the read lives on the emergent chart}\ M\ [\mathrm{CT2.6}],
\]

%dsp-level:0
\[
\ \text{inheriting hypotheses (i)--(iii) at the point of use}\ [\text{A §0.2a}];
\]

%dsp-level:0
\[
\text{the}\ 16\text{-dimensional pair register carries no canonical pair}\ [\mathrm{BM.10}(1)];
\]

%dsp-level:0
\[
\qquad \text{the register-side precursors are the discrete Weyl pairs}\ [\mathrm{BM.10}(2)]\ \text{and the width law}\ [\mathrm{BM.10}(6)];
\]

%dsp-level:0
\[
(2)\quad\textbf{the metric data and the derived ADM split:}
\]

%dsp-level:0
\[
\qquad \text{for the coarse-grained sector coefficients}\ (a_{00},q_{ij},w_i)\ [\mathrm{D2.1}]:
\]

%dsp-level:0
\[
\Omega=e^{\alpha_\Omega a_{00}},\qquad g_{00}=-\Omega^2,\qquad g_{0i}=\Omega^2\varkappa\beta w_i,\qquad g_{ij}=\Omega^2(\delta_{ij}+\varkappa q_{ij})\qquad[\mathrm{CT1.12};\ \mathrm{FT.157}(1)];
\]

%dsp-level:0
\[
\textbf{lapse:}\quad N^2:=\Omega^2+\gamma^{ij}N_iN_j=-\frac{\det g}{\det\gamma}>0;
\]

%dsp-level:0
\[
\qquad \textbf{shift covector:}\quad N_i:=g_{0i}=\Omega^2\varkappa\beta w_i;
\]

%dsp-level:0
\[
\qquad \textbf{shift vector:}\quad N^i:=\gamma^{ij}N_j;
\]

%dsp-level:0
\[
\textbf{the line element (an exact identity):}
\]

%dsp-level:0
\[
\qquad ds^2=-N^2d\tau^2+g_{ij}\bigl(dx^i+N^id\tau\bigr)\bigl(dx^j+N^jd\tau\bigr)=g_{\mu\nu}dx^\mu dx^\nu
\]

%dsp-level:0
\[
\qquad [\text{machine: Schur, denominators cleared}];
\]

%dsp-level:0
\[
(3)\quad\textbf{the momenta:}
\]

%dsp-level:0
\[
\qquad \textbf{the unit-mass coordinate-time read:}\quad\pi_i:=g_{ij}\Bigl(\frac{dx^j}{dx^0}+N^j\Bigr)
\]

%dsp-level:0
\[
\qquad [x^0=\tau_n\ \text{the chart time};\ \mathrm{D1.8};\ \mathrm{FT.159}];
\]

%dsp-level:0
\[
\textbf{the physical spatial momentum:}\qquad p_i:=m\,u^0\,\pi_i,\qquad u^\mu:=\frac{dx^\mu}{d\lambda},\qquad g_{\mu\nu}u^\mu u^\nu=-1;
\]

%dsp-level:0
\[
\text{the shift contribution is}\ g_{ij}N^j=N_i=g_{0i}\ [\text{machine: exact polynomial identity}]:
\]

%dsp-level:0
\[
\qquad \text{the metric's time--space block enters the momentum exactly}\ [\mathrm{FT.157}(1)];
\]

%dsp-level:0
\[
(4)\quad\textbf{the rest input:}
\]

%dsp-level:0
\[
\qquad m:=m_{\mathrm{hid}}=\frac{E_{\mathrm{int}}\ln2}{\mathfrak h}
\]

%dsp-level:0
\[
\qquad [\mathrm{FT.44}(4)\ \text{(Census)}];
\]

%dsp-level:0
\[
\qquad \text{the ensemble route:}\quad m\ \mapsto\ \langle m\rangle\ [\mathrm{BM.20}(3)];
\]

%dsp-level:0
\[
\text{the vacuum-deficit reframing of the rest read is}\ \mathrm{BM.35}\ —\ \text{a marked forward pointer at the citing site, imported in no}\ \textbf{Req}\ \text{line};
\]

%dsp-level:0
\[
(5)\quad\textbf{the mass shell and the energy:}\qquad g^{\mu\nu}p_\mu p_\nu=-m^2\qquad[\text{signature}\ (-,+,+,+);\ \mathrm{CT1.12}];
\]

%dsp-level:0
\[
E:=-p_0\quad\text{is}\ \mathrm{BM.8}(2)\text{'s conditional energy read under its declared stationarity scope (theorem-local, inherited at the point of use)};
\]

%dsp-level:0
\[
\textbf{the Eulerian form:}
\]

%dsp-level:0
\[
\qquad E=N\sqrt{m^2+\gamma^{ij}p_ip_j}-N^ip_i;
\]

%dsp-level:0
\[
\qquad \textbf{the kinetic-energy read:}
\]

%dsp-level:0
\[
\qquad \mathrm{KE}:=E-N\,m\ \ge0\ \ \text{where}\ g_{ij}\succ0\ [\mathrm{CT1.12};\ \text{machine}];
\]

%dsp-level:0
\[
(6)\quad\textbf{status and scope:}
\]

%dsp-level:0
\[
\qquad \textbf{constructed}\ —\ \text{an abbreviation of derived objects:}\ (2)\text{'s split is linear algebra on}\ \mathrm{CT1.12}\text{'s dictionary};
\]

%dsp-level:0
\[
\qquad (3)\text{--}(5)\ \text{are pointwise identities}\ [\text{machine}];
\]

%dsp-level:0
\[
\neg\text{a new parameter}\ [\text{§Cal-BM stays empty}];
\]

%dsp-level:0
\[
\qquad \neg\text{a declaration};
\]

%dsp-level:0
\[
\qquad \text{the names (lapse, shift, kinetic energy) are internal bookkeeping}\ [\text{A §0.2}];
\]

%dsp-level:0
\[
\neg\text{constructed: the back-reaction of the composite's motion on its own sector coefficients}\ [\mathrm{BM.15}\text{'s read is the rest-frame one}]\ \text{and any inter-sector dynamics}\ [\mathrm{BM.29}(5)\ \text{stands}].
\]

\textbf{Rmk (the unit-mass normalization).} The leading kinetic read is $\mathrm{KE}=\tfrac12N\,\gamma^{ij}p_ip_j/m+O(p^4)$: the factor $1/m$ is dimensionally forced, since $[\gamma^{ij}p_ip_j]=\mathfrak q^{+1}=\text{energy}^2$ while $[\mathrm{KE}]=\mathfrak q^{+1/2}=\text{energy}$ [CT2.1; machine: exponent arithmetic]. The unit-mass unit-lapse normalization ($m=1$, $N=1$, $w_i=0$) gives $\mathrm{KE}=\tfrac12\gamma^{ij}p_ip_j$. And $\gamma^{ij}\ne g^{ij}$ where the shift is nonzero: $g^{ij}=\gamma^{ij}-N^iN^j/N^2$ [machine] — the contraction that inverts the spatial momentum is $\gamma^{ij}$; the two coincide exactly in the shift-free gauge.

\textbf{Req.} D1.8, D2.1, D2.9 (Parts 1–2); CT1.12, CT2.1, CT2.6 (Part 2); P1, P5 (Parts 1–2); FT.157(1), FT.159 (Part 3); FT.44(4) (Census); BM.8(2), BM.10(1), BM.10(2), BM.10(6), BM.20(3) (this part); A §0.2a (the continuum-limit hypotheses at the point of use); Schur decomposition, adjugate calculus [math; machine: \texttt{checks/\allowbreak{}verify\_\allowbreak{}H\_\allowbreak{}kinetic\_\allowbreak{}energy\_\allowbreak{}metric.py} §A].

\textbf{Pf.}

%dsp-level:0
\[
\text{Definition block: the constitutive clauses are the displayed identities of}\ (2)\text{--}(5)\ [\text{machine: exact polynomial identities over}\ \mathbb Q[\Omega^2,\varkappa,\beta,w_i,q_{ij},p_\mu,m]\ \text{with denominators cleared by}\ \det\gamma];
\]

%dsp-level:0
\[
\neg\exists\ \text{a non-derived clause: the ADM split is forced by the Schur decomposition of}\ g_{\mu\nu}\ [\text{math}];
\]

%dsp-level:0
\[
\qquad (1),(6)\ \text{are scope records}.\ \blacksquare
\]

\textbf{Phys.} The energy of a moving sealed composite is read off the metric the corpus already derived, not added to it. The dictionary's ten sector coefficients split into lapse, shift, and spatial metric by linear algebra alone — the Schur decomposition of the very matrix CT1.12 builds — and the same three shift coefficients that the weak block carries enter the spatial momentum as its drag term. The mass-shell constraint of the Lorentzian metric converts these data into one exact dispersion identity, whose rest input is the informational mass the calibration tiers locked. Two normalizations — unit mass and coordinate time — are fixed as bookkeeping, and the read is chart-side: the register itself carries no canonical pair, and the definition keeps the obstruction exactly where Tier BM.4 placed it.

\end{definition}

\begin{theorem}[{\texorpdfstring{BM.34 — The kinetic energy from the emergent metric: the ADM split, the mass shell, the dispersion relation, and the zero-point bound}{BM.34 - The kinetic energy from the emergent metric: the ADM split, the mass shell, the dispersion relation, and the zero-point bound}}]\label{tag:BM-34}
\textbf{Thm.}

%dsp-level:0
\[
(1)\quad\textbf{the split is derived:}
\]

%dsp-level:0
\[
\qquad \gamma_{ij}=\Omega^2(\delta_{ij}+\varkappa q_{ij}),\quad N_i=\Omega^2\varkappa\beta w_i,\quad N^i=\gamma^{ij}N_j,\quad N^2=\Omega^2+\gamma^{ij}N_iN_j;
\]

%dsp-level:0
\[
\boxed{\ N^2\det\gamma=-\det g\ }\qquad[\text{Schur; machine: exact polynomial identity over}\ \mathbb Q[\Omega^2,\varkappa,\beta,w_i,q_{ij}]];
\]

%dsp-level:0
\[
\textbf{the inverse metric in ADM form:}
\]

%dsp-level:0
\[
\qquad g^{00}=-\frac1{N^2},
\]

%dsp-level:0
\[
\qquad g^{0i}=\frac{N^i}{N^2},
\]

%dsp-level:0
\[
\qquad g^{ij}=\gamma^{ij}-\frac{N^iN^j}{N^2}
\]

%dsp-level:0
\[
\qquad [\text{machine: cleared adjugate identities}\ +\ 3\ \text{rational witnesses, exact Fractions}];
\]

%dsp-level:0
\[
(2)\quad\textbf{the line element and the vacuum:}
\]

%dsp-level:0
\[
\qquad ds^2=-N^2d\tau^2+g_{ij}\bigl(dx^i+N^id\tau\bigr)\bigl(dx^j+N^jd\tau\bigr)\ \ \text{identically}\ [\text{machine}];
\]

%dsp-level:0
\[
(a_{00},q_{ij},w_i)=0\implies(\Omega,N,N^i,g_{ij})=(1,1,0,\delta_{ij}):
\]

%dsp-level:0
\[
\qquad \textbf{the vacuum's emergent metric is the flat metric}\ [\mathrm{CT1.12};\ \text{machine: exact substitution}];
\]

%dsp-level:0
\[
(3)\quad\textbf{the momentum and its carrier:}
\]

%dsp-level:0
\[
\qquad p_i=m\,u^0\,g_{ij}\Bigl(\frac{dx^j}{dx^0}+N^j\Bigr)\ [\mathrm{Def\text{-}BM.10}(3)];
\]

%dsp-level:0
\[
\qquad g_{ij}N^j=N_i=g_{0i}=\Omega^2\varkappa\beta w_i\ [\text{machine}];
\]

%dsp-level:0
\[
\text{the drag term and the metric's time--space block are the same three coefficients}\ w_i\ —\ \text{the weak aspect, one block, two derived reads}\ [\mathrm{FT.157}(1);\ \mathrm{CT1.12a}(2)];
\]

%dsp-level:0
\[
\textbf{the T-parity of the carrier:}
\]

%dsp-level:0
\[
\quad w\mapsto-w\implies N_i\mapsto-N_i,
\]

%dsp-level:0
\[
\quad N^i\mapsto-N^i,
\]

%dsp-level:0
\[
\quad (N,\gamma_{ij})\ \text{fixed}\ [\mathrm{FT.156}(3);\ \mathrm{FT.159}(3);\ \text{machine}]:
\]

%dsp-level:0
\[
\qquad p_i\ \text{is T-odd, matching its carrier}\ [\mathrm{BM.8}(1)];
\]

%dsp-level:0
\[
(4)\quad\textbf{the mass shell}\iff\textbf{the dispersion relation:}\qquad g^{\mu\nu}p_\mu p_\nu=-m^2\iff\bigl(p_0-N^ip_i\bigr)^2=N^2\bigl(m^2+\gamma^{ij}p_ip_j\bigr)
\]

%dsp-level:0
\[
[\text{machine: the master polynomial identity, cleared by}\ (\det\gamma)^2];
\]

%dsp-level:0
\[
\qquad E:=-p_0,\ \text{the positive-energy branch}:
\]

%dsp-level:0
\[
\qquad \boxed{\ E=N\sqrt{m^2+p^2}-N^ip_i\ },
\]

%dsp-level:0
\[
\qquad p^2:=\gamma^{ij}p_ip_j;
\]

%dsp-level:0
\[
\textbf{the gauges:}\qquad w=0:\quad E^2=N^2(m^2+p^2);\qquad\Omega=1\wedge w=0:\quad\boxed{\ E^2=m^2+p^2\ }\qquad[c=1,\ \mathrm{D2.9}];
\]

%dsp-level:0
\[
\textbf{the dispersion is T-even:}\quad\text{the mass shell is invariant under the simultaneous}\ (N^i,p_i)
\]

%dsp-level:0
\[
\mapsto (-N^i,-p_i)\ [\text{machine: substitution in the master identity}];
\]

%dsp-level:0
\[
\qquad \text{the signature is invariant}\ [\mathrm{FT.159}(3)\ \text{re-cited}];
\]

%dsp-level:0
\[
(5)\quad\textbf{the kinetic energy and the non-relativistic read:}
\]

%dsp-level:0
\[
\qquad \mathrm{KE}:=E-Nm=N\bigl(\sqrt{m^2+p^2}-m\bigr)\ \ge0,
\]

%dsp-level:0
\[
\qquad \mathrm{KE}=0\iff p_i=0;
\]

%dsp-level:0
\[
\text{on the perturbative regime}\quad g_{ij}\succ0\iff\forall a:\ 1+\varkappa\lambda_a(q)>0
\]

%dsp-level:0
\[
\qquad [\mathrm{CT1.12};\ \text{machine: exact rational Sylvester/LDL}^\top\ \text{witnesses}];
\]

%dsp-level:0
\[
\sqrt{1+x}=1+\frac x2-\frac{x^2}8+\frac{x^3}{16}-\frac{5x^4}{128}+\frac{7x^5}{256}+O(x^6)
\]

%dsp-level:0
\[
\qquad [\text{machine: exact rational coefficients;}\ (\cdot)^2\equiv1+x\ \bmod x^6];
\]

%dsp-level:0
\[
\boxed{\ \mathrm{KE}=N\Bigl(\frac{p^2}{2m}-\frac{p^4}{8m^3}+\frac{p^6}{16m^5}\Bigr)+O(p^8)\ },\qquad\text{the leading read:}\quad\mathrm{KE}=\frac{N\,\gamma^{ij}p_ip_j}{2m};
\]

%dsp-level:0
\[
\text{the unit normalization}\ (m{=}1,N{=}1,w{=}0):\quad\mathrm{KE}=\tfrac12\gamma^{ij}p_ip_j\ [\mathrm{Def\text{-}BM.10}\ \mathrm{Rmk}];
\]

%dsp-level:0
\[
\qquad [\mathrm{KE}]=\mathfrak q^{+1/2},\quad[\gamma^{ij}p_ip_j]=\mathfrak q^{+1}\ [\mathrm{CT2.1};\ \text{machine: exponent arithmetic}];
\]

%dsp-level:0
\[
(6)\quad\textbf{the zero-point bound:}
\]

%dsp-level:0
\[
\qquad \Delta x\,\Delta p\ge\tfrac12\ [\mathrm{BM.10}(3),(6);\ \hbar=1];
\]

%dsp-level:0
\[
\qquad \Delta x=\Delta x_{\mathrm B}:\quad\Delta p\ge\frac1{2\Delta x_{\mathrm B}}=\frac{m_{\mathrm{hid}}}2\ [\mathrm{BM.13}(4)];
\]

%dsp-level:0
\[
\text{the isotropic cube convention}\ [\mathrm{Def\text{-}BM.4}(3)],\ \text{the rest read}\ \langle p_i\rangle=0:
\]

%dsp-level:0
\[
\qquad \langle p^2\rangle\ \ge\ 3\Bigl(\frac{m_{\mathrm{hid}}}2\Bigr)^2=\frac34\,m_{\mathrm{hid}}^2
\]

%dsp-level:0
\[
\qquad [\text{machine: exact Fractions}];
\]

$\text{in the vacuum gauge of}\ (4)\ \bigl(E^2=m^2+\langle p^2\rangle\bigr),\ \text{at}\ m=m_{\mathrm{hid}}:\qquad\boxed{\ E^2\ \ge\ \frac74\,m_{\mathrm{hid}}^2\ },\qquad\mathrm{KE}_{\mathrm{NR}}\ \ge\ \frac{\langle p^2\rangle}{2m}\ \ge\ \frac38\,m_{\mathrm{hid}},\qquad E\ \ge\ \frac{\sqrt7}2\,m_{\mathrm{hid}}\ \ \text{exact in}\ \mathbb Q(\sqrt7)$

%dsp-level:0
\[
[\text{machine: the quadratic ring, minimal polynomial}\ x^2-7,\ \text{comparison by squaring}];
\]

%dsp-level:0
\[
\qquad m_{\mathrm{hid}}=E_{\mathrm{int}}\ln2/\mathfrak h\ \text{substituted}:
\]

%dsp-level:0
\[
\quad E^2\ge\frac74\,\frac{(E_{\mathrm{int}}\ln2)^2}{\mathfrak h^2}\ [\text{machine: one Laurent monomial, coefficient}\ \tfrac74];
\]

%dsp-level:0
\[
(7)\quad\textbf{the rest slice and the dust closure:}
\]

%dsp-level:0
\[
\qquad p_i=0:\quad E=Nm;
\]

%dsp-level:0
\[
\qquad \text{the vacuum gauge:}\quad E=m=m_{\mathrm{hid}}=\frac{E_{\mathrm{int}}\ln2}{\mathfrak h}\ [\mathrm{FT.44}(4)];
\]

%dsp-level:0
\[
RV=\varkappa^2m\ [\mathrm{BM.20}(1)],
\]

%dsp-level:0
\[
\qquad \rho=\frac mV,\quad\kappa^2\rho=\frac{2\pi\,\mathfrak h\,E_{\mathrm{int}}}V\ [\mathrm{BM.9}(5);\ \mathrm{BM.15}(2)]:
\]

%dsp-level:0
\[
\qquad \textbf{the informational, the curvature-flux and the dispersion rest reads agree on-shell}\ [\text{machine: Laurent identities}];
\]

%dsp-level:0
\[
\text{KE is the}\ p
\]

%dsp-level:0
\[
\ne 0\ \text{completion the rest-frame dust read does not carry}\ [\mathrm{BM.9}(4):\ T_{\mu\nu}=\mathrm{diag}(\rho,0,0,0)\ \text{at}\ p=0]:
\]

%dsp-level:0
\[
\qquad \text{the kinetic term is the whole composite's transport on the chart}\ [\mathrm{BM.8}(1)].
\]

\textbf{Rmk (chart-side discipline).} The reads of (4)–(6) are continuum-chart objects under CT2.6's hypotheses at the point of use; no canonical pair exists in the finite register [BM.10(1)] and no clause imports one. The pointwise mass-shell identity needs no stationarity; the conserved-charge read $E=-p_0$ inherits BM.8(2)'s theorem-local stationarity hypothesis at the point of use. The bound of (6) applies the Robertson inequality chart-side, at the width law's locus — Def-BM.4(5)'s envelope discipline: the zero-point read is a coarse-graining envelope, never a ledger distance.

\textbf{Req.} A3 (Part 1); D1.8, D2.1, D2.9 (Parts 1–2); CT1.12, CT1.12a(2), CT2.1, CT2.6 (Part 2); P1, P5 (Parts 1–2); FT.44(4) (Census); FT.156(3), FT.157(1), FT.159 (Part 3); Def-BM.4, Def-BM.10, BM.8(1), BM.8(2), BM.9(4), BM.9(5), BM.10(1), BM.10(2), BM.10(3), BM.10(6), BM.13(4), BM.15(2), BM.20(1), BM.20(3) (this part); A §0.2a (the hypotheses at the point of use); Schur decomposition, adjugate calculus, the Sylvester criterion, exact rational $\mathrm{LDL}^\top$ inertia, the binomial series over $\mathbb Q$, the quadratic ring $\mathbb Q(\sqrt7)$, Laurent monomials [math; machine: \texttt{checks/\allowbreak{}verify\_\allowbreak{}H\_\allowbreak{}kinetic\_\allowbreak{}energy\_\allowbreak{}metric.py} §A].

\textbf{Pf.}

%dsp-level:0
\[
(1):\ \det g=\det\gamma\cdot\bigl(g_{00}-g_{0i}\gamma^{ij}g_{0j}\bigr)\ [\text{Schur}];
\]

%dsp-level:0
\[
\qquad g_{00}=-\Omega^2,\ g_{0i}=N_i\implies-\det g=\bigl(\Omega^2+\gamma^{ij}N_iN_j\bigr)\det\gamma=N^2\det\gamma\ [\text{machine}]
\]

%dsp-level:0
\[
\text{the inverse rows in cleared adjugate form}\ (\operatorname{adj}\gamma:=\det\gamma\cdot\gamma^{-1},\ D:=\det\gamma):
\]

%dsp-level:0
\[
\quad (\operatorname{adj}g)_{00}=D;
\]

%dsp-level:0
\[
\quad (\operatorname{adj}g)_{0i}=-\bigl((\operatorname{adj}\gamma)^{ij}N_j\bigr);
\]

%dsp-level:0
\[
\quad D\,
\]

%dsp-level:0
\[
(\operatorname{adj}g)_{ij}
\]

%dsp-level:0
\[
= (\det g)\,
\]

%dsp-level:0
\[
(\operatorname{adj}\gamma)^{ij}+\bigl((\operatorname{adj}\gamma)^{ik}N_k\bigr)\bigl((\operatorname{adj}\gamma)^{j\ell}N_\ell\bigr)\ [\text{machine: three exact polynomial identities}]
\]

%dsp-level:0
\[
\text{the three rational witnesses}\ (\Omega^2,\varkappa,\beta,w_i,q_{ij}\ \text{rational,}\ q\ \text{off-diagonal}):\ \text{every row as exact Fraction equality, including}\ g^{\mu\alpha}g_{\alpha\nu}
\]

%dsp-level:0
\[
= \delta^\mu_\nu\ \text{and}\ \det g\qquad = -N^2\det\gamma
\]

%dsp-level:0
\[
(2):\ \text{expansion in}\ (d\tau,dx^i):\quad g_{00}=-N^2+N_iN^i,\quad g_{0i}=\gamma_{ij}N^j=N_i,\quad g_{ij}=\gamma_{ij}\ [\text{machine}];
\]

%dsp-level:0
\[
\qquad \Omega=e^{\alpha_\Omega\cdot0}=1,\ q=0,\ w=0\implies\text{the flat metric}\ [\text{machine: exact substitution}]
\]

%dsp-level:0
\[
(3):\ \gamma_{ij}N^j=N_i\ [\text{machine: cleared by}\ \det\gamma];
\]

%dsp-level:0
\[
\qquad w
\]

%dsp-level:0
\[
\mapsto -w:\ N_i
\]

%dsp-level:0
\[
= \Omega^2\varkappa\beta w_i
\]

%dsp-level:0
\[
\mapsto -N_i\ \text{and}\ (N,\gamma_{ij})\ \text{carry}\ w\ \text{only through}\ N_iN_j\ \text{and}\ N_i\ [\text{machine:}\ N^2,\gamma_{ij}\ \text{fixed}]
\]

%dsp-level:0
\[
(4):\ g^{\mu\nu}p_\mu p_\nu
\]

%dsp-level:0
\[
= -\frac1{N^2}\bigl(p_0-N^ip_i\bigr)^2+\gamma^{ij}p_ip_j\ [(1)\text{'s inverse rows}]\implies\bigl(p_0-N^ip_i\bigr)^2
\]

%dsp-level:0
\[
= N^2\bigl(m^2+\gamma^{ij}p_ip_j\bigr)\ \text{on shell}\ [\text{machine: cleared by}\ (\det\gamma)^2]
\]

%dsp-level:0
\[
\text{the branch: at}\ p_i=0,\ w=0:\ p_0^2=N^2m^2\implies p_0=\mp Nm;
\]

%dsp-level:0
\[
\qquad E=-p_0=+Nm\ \text{on the positive branch}\ [\text{machine: rational witnesses}];
\]

%dsp-level:0
\[
\qquad \text{the gauges: substitution}\ [\text{machine}];
\]

%dsp-level:0
\[
\qquad \text{T-evenness:}\ (N^i,p_i)\mapsto(-N^i,-p_i)\ \text{fixes}\ (p_0-N^ip_i)^2\ \text{and}\ \gamma^{ij}p_ip_j\ [\text{machine}]
\]

%dsp-level:0
\[
(5):\ \mathrm{KE}=N\bigl(\sqrt{m^2+p^2}-m\bigr)=Nm\bigl(\sqrt{1+p^2/m^2}-1\bigr)\ [\text{exact}];
\]

%dsp-level:0
\[
\qquad \text{the series}\ a=\bigl(1,\tfrac12,-\tfrac18,\tfrac1{16},-\tfrac5{128},\tfrac7{256}\bigr):
\]

%dsp-level:0
\[
\quad (a*a)_n=\delta_{n0}+\delta_{n1}\ \forall n\le5\ [\text{machine: Fractions}]
\]

%dsp-level:0
\[
\text{positivity:}\ g_{ij}\succ0\ [\mathrm{CT1.12}]\implies\gamma^{ij}\succ0
\]

%dsp-level:0
\[
\implies p^2
\]

%dsp-level:0
\[
\ge 0\ \text{with equality iff}\ p_i
\]

%dsp-level:0
\[
= 0\ [\text{machine:}\ \mathrm{LDL}^\top\ \text{on exact rational witnesses}];
\]

%dsp-level:0
\[
\qquad [\mathrm{KE}]=[\gamma^{ij}p_ip_j]-[m]=\mathfrak q^{1-1/2}=\mathfrak q^{+1/2}\ [\mathrm{CT2.1};\ \text{machine}]
\]

%dsp-level:0
\[
(6):\ \Delta p\ge\tfrac1{2\Delta x_{\mathrm B}}=\tfrac{m_{\mathrm{hid}}}2\ [\mathrm{BM.13}(4)\ \text{re-cited}];
\]

%dsp-level:0
\[
\qquad \langle p^2\rangle=\sum_i\bigl((\Delta p_i)^2+\langle p_i\rangle^2\bigr)=\sum_i(\Delta p_i)^2\ge3\bigl(\tfrac{m_{\mathrm{hid}}}2\bigr)^2\ \text{at}\ \langle p_i\rangle=0\ [\text{machine}]
\]

%dsp-level:0
\[
E^2=m^2+\langle p^2\rangle\ \text{in the vacuum gauge of}\ (4):\quad E^2\ge\bigl(1+\tfrac34\bigr)m_{\mathrm{hid}}^2\ \text{at}\ m=m_{\mathrm{hid}}\ [\text{machine}];
\]

%dsp-level:0
\[
\qquad \mathrm{KE}_{\mathrm{NR}}=\tfrac{\langle p^2\rangle}{2m}\ge\tfrac{3/4}{2}\,m_{\mathrm{hid}}=\tfrac38m_{\mathrm{hid}}\ [\text{machine}];
\]

%dsp-level:0
\[
\qquad \bigl(\tfrac{\sqrt7}2\bigr)^2=\tfrac74,\ x^2-7\ \text{irreducible over}\ \mathbb Q\ [\text{machine: the quadratic ring}];
\]

%dsp-level:0
\[
\qquad E^2\ge\tfrac74m_{\mathrm{hid}}^2\wedge E>0\implies E\ge\tfrac{\sqrt7}2m_{\mathrm{hid}}\ [\text{squaring comparison}]
\]

%dsp-level:0
\[
(7):\ p=0\ \text{in}\ (4):\ E^2=N^2m^2\implies E=Nm\ [\text{the positive branch}];
\]

%dsp-level:0
\[
\qquad RV=\varkappa^2m\ \text{and}\ \rho V=m:\ \text{the same monomial}\ [\text{machine}];
\]

%dsp-level:0
\[
\qquad \kappa^2\rho
\]

%dsp-level:0
\[
= \tfrac{2\pi\mathfrak h^2}{\ln2}\cdot\tfrac{E_{\mathrm{int}}\ln2}{\mathfrak h V}
\]

%dsp-level:0
\[
= \tfrac{2\pi\mathfrak h E_{\mathrm{int}}}V\ [\mathrm{BM.9}(5)\ \text{re-verified; machine: denominators cleared}].\ \blacksquare
\]

\textbf{Phys.} Motion enters the corpus through the metric the corpus already built. The dictionary's ten coefficients split into lapse, shift, and spatial metric by the Schur complement alone, and the mass-shell constraint of the Lorentzian metric turns that split into one exact identity: the energy conjugate to ledger time, corrected by the shift's drag, is the lapse times the square root of the rest mass squared plus the spatial momentum squared. In the vacuum gauge that is precisely the dispersion relation, and its expansion gives the kinetic energy with rational coefficients no one fitted — one half, minus one eighth, plus one sixteenth — nonnegative, vanishing only at zero momentum, wherever the dictionary's perturbative positivity holds. Two exact consequences deserve names. The width law converts the uncertainty bound at the Born width into a zero-point bound: a sealed composite spread over its own width carries at least seven quarters of its rest mass squared in energy squared — three eighths of the rest mass in the non-relativistic kinetic read, $\sqrt7/2$ times it in the exact relativistic one — every constant an exact rational over the ledger count. And the rest slice closes a triangle the corpus drew twice before: at zero momentum, the dispersion energy, the informational mass, and the curvature flux are one and the same monomial.

\end{theorem}

\begin{theorem}[{\texorpdfstring{BM.35 — The vacuum-balanced entropy: the rest mass as the deficit against the all-virtual vacuum}{BM.35 - The vacuum-balanced entropy: the rest mass as the deficit against the all-virtual vacuum}}]\label{tag:BM-35}
\textbf{Thm.}

%dsp-level:0
\[
(1)\quad\textbf{the vacuum is the unique entropy maximum:}
\]

%dsp-level:0
\[
\qquad n_{\mathcal R}=0\implies\rho_{pq}=\rho_\circ=\tfrac14I\ \ \forall(p,q),
\]

%dsp-level:0
\[
\qquad S_{\mathrm{raw}}(\rho_\circ)=\ln4=\max_\rho S_{\mathrm{raw}}(\rho),
\]

%dsp-level:0
\[
\qquad \arg\max=\bigl\{\tfrac14I\bigr\}
\]

%dsp-level:0
\[
\qquad [\mathrm{CT0.2}];
\]

%dsp-level:0
\[
\textbf{the machine form:}
\]

%dsp-level:0
\[
\qquad \lambda=\bigl(\tfrac{k_1}8,\dots,\tfrac{k_4}8\bigr),\ \textstyle\sum_ik_i=8:
\]

%dsp-level:0
\[
\qquad S(\lambda)\le\ln4\iff\prod_ik_i^{k_i}\ge256\quad(0^0:=1)
\]

%dsp-level:0
\[
\qquad [\text{integer comparison;}\ \exp\ \text{strictly monotone}];
\]

%dsp-level:0
\[
\text{the exhaustive census of all}\ \tbinom{11}3
\]

%dsp-level:0
\[
= 165\ \text{compositions: the inequality holds at every one, with equality only at}\ (2,2,2,2)\ [\text{machine}]:
\]

%dsp-level:0
\[
\qquad \textbf{maximality and uniqueness without one transcendental evaluated};
\]

%dsp-level:0
\[
(2)\quad\textbf{the vacuum is flat:}
\]

%dsp-level:0
\[
\qquad (a_{00},q_{ij},w_i)=0\implies(N,N^i,g_{ij})=(1,0,\delta_{ij})\ [\mathrm{BM.34}(2)];
\]

%dsp-level:0
\[
\qquad \Theta_R=0\ \text{on geodesic polygons in flat slices}\ [\mathrm{FT.47}(1)];
\]

%dsp-level:0
\[
\textbf{the vacuum is entropy-maximal and defect-free}\ —\ \text{the two reference points against which every deficit read is read};
\]

%dsp-level:0
\[
(3)\quad\textbf{the ledger deficit and the balance identity:}
\]

%dsp-level:0
\[
\qquad \text{one real edge}=1\ \text{bit}=\ln2\ \text{nats}\ [\mathrm{C9};\ \mathrm{CT0.9};\ \mathrm{CT1.29}]:
\]

%dsp-level:0
\[
\qquad S_{\mathfrak L}=n_{\mathcal R}\ln2;
\]

%dsp-level:0
\[
\qquad \text{the vacuum:}\ n_{\mathcal R}=0\implies S_{\mathfrak L}(\mathrm{vac})=0\ [\mathrm{CT0.2};\ \mathrm{D0.1}];
\]

%dsp-level:0
\[
\Delta S_{\mathfrak L}:=S_{\mathfrak L}-S_{\mathfrak L}(\mathrm{vac})=E_{\mathrm{int}}\ln2=I_{\mathrm{hid}}\qquad[\mathrm{FT.42}(1);\ \mathrm{FT.44}(4);\ \text{machine: exact in}\ \mathbb Q\cdot\ln2];
\]

%dsp-level:0
\[
\textbf{the balance identity:}
\]

%dsp-level:0
\[
\qquad \boxed{\ \ln4=2\ln2\ }
\]

%dsp-level:0
\[
\qquad [\text{machine}]:
\]

%dsp-level:0
\[
\qquad \text{each bidirectionally crystallized channel}\ (2\ \text{directed edges}\ [\mathrm{D0.8}])\ \text{books ledger entropy}\ 2\ln2
\]

%dsp-level:0
\[
= \ln4
\]

%dsp-level:0
\[
= \text{the vacuum pair's own entropy};
\]

%dsp-level:0
\[
E_{\mathrm{int}}
\]

%dsp-level:0
\[
= 2E_{\mathrm{bidir}}\ [\text{the bidirectional composites: the digon}\ [\mathrm{FT.28}],\ \text{the polygon ladder}\ [\mathrm{FT.45}(1)],\ \text{the Platonic floors}\ [\mathrm{FT.42}(1)],\ \text{the neutron row}\ [\mathrm{BM.17}]]:
\]

%dsp-level:0
\[
\boxed{\ \Delta S_{\mathfrak L}=E_{\mathrm{bidir}}\ln4=E_{\mathrm{bidir}}\,S_{\mathrm{raw}}(\rho_\circ)\ }
\]

%dsp-level:0
\[
\qquad [\text{machine: exact in}\ \mathbb Q\cdot\ln2]:
\]

%dsp-level:0
\[
\qquad \textbf{crystallization re-books the vacuum's pair entropy into the ledger in full}\ —\ \text{the balance is exact};
\]

%dsp-level:0
\[
(4)\quad\textbf{the reframing identity:}
\]

%dsp-level:0
\[
\qquad m_{\mathrm{rest}}:=\frac{\Delta S_{\mathfrak L}}{\mathfrak h}=\frac{E_{\mathrm{int}}\ln2}{\mathfrak h}=m_{\mathrm{hid}}
\]

%dsp-level:0
\[
\qquad [\mathrm{FT.44}(4);\ \text{machine: Laurent identity}];
\]

%dsp-level:0
\[
[\Delta S_{\mathfrak L}]=\mathfrak q^0,\quad[\mathfrak h]=\mathfrak q^{-1/2}\implies[m_{\mathrm{rest}}]=\mathfrak q^{+1/2}=[\text{energy}]
\]

%dsp-level:0
\[
\qquad [\mathrm{CT2.1};\ \mathrm{D2.9};\ \text{machine: exact exponent arithmetic}];
\]

%dsp-level:0
\[
\textbf{the substance clause:}
\]

%dsp-level:0
\[
\qquad \text{the rest mass is the deficit read of ledger data against the vacuum}\ —\ \neg\text{an independent entity: the ontology adds nothing};
\]

%dsp-level:0
\[
\text{the rest energy:}\quad E_{\mathrm{rest}}=m_{\mathrm{rest}}c^2\ \ \text{with}\ \ c=1\ [\mathrm{D2.9}]:
\]

%dsp-level:0
\[
\qquad \text{the}\ c^2\ \text{is bookkeeping restoration only}\ [\text{machine: the}\ c\text{-exponent is}\ 0];
\]

%dsp-level:0
\[
(5)\quad\textbf{the geometric member of the deficit family:}
\]

%dsp-level:0
\[
\qquad m^2_{\mathrm{pole}}=4\hat\mu'^2\sin^2\Bigl(\frac{\Theta_R}4\Bigr),
\]

%dsp-level:0
\[
\qquad \Theta_R=\textstyle\sum_i\alpha_i-(n-2)\pi
\]

%dsp-level:0
\[
\qquad [\mathrm{Def\text{-}H.37}(2);\ \mathrm{FT.54}(4);\ \mathrm{FT.47}(1)]:
\]

%dsp-level:0
\[
\text{the pole mass is the angle-deficit read against}\ (2)\text{'s flat vacuum};
\]

%dsp-level:0
\[
\qquad \text{the range}\ [0,4\hat\mu'^2]\ [\mathrm{FT.76}(3)];
\]

%dsp-level:0
\[
\qquad \text{the exact witness}\ 2-\sqrt2\in\mathbb Q(\sqrt2)\ [\mathrm{FT.54}(5)];
\]

%dsp-level:0
\[
\neg\exists\ \text{a theorem equating the bit currency}\ (\ln2\ \text{per directed edge})\ \text{and the angle currency}\ (\Theta_R):
\]

%dsp-level:0
\[
\qquad \text{the relating equation is application-layer content};
\]

%dsp-level:0
\[
(6)\quad\textbf{the ensemble refinement and the third currency:}
\]

%dsp-level:0
\[
\qquad \langle m\rangle=\sum_Rw(R)\,m(R)\ [\mathrm{BM.20}(3)];
\]

%dsp-level:0
\[
\qquad \Delta S(\Gamma,R)=\tfrac32\bigl(2E-|R|\bigr)\ \text{nats}\ [\mathrm{FT.107}(2);\ \mathrm{L3}];
\]

%dsp-level:0
\[
\textbf{the three deficit functionals are pairwise distinct:}
\]

%dsp-level:0
\[
\qquad \text{the digon skeleton}\ (2E=2):\quad|R|=2:\ \ \Delta S_{\mathfrak L}=2\ln2\ \ \text{vs}\ \ \Delta S=0;
\]

%dsp-level:0
\[
\qquad R=\varnothing:\ \ \Delta S_{\mathfrak L}=0\ \ \text{vs}\ \ \Delta S=3
\]

%dsp-level:0
\[
\qquad [\text{machine}];
\]

%dsp-level:0
\[
2\ln2\ne3\ \text{by the exact enclosure}\ \tfrac12<\ln2<1\ [\text{machine:}\ \ln2=\textstyle\sum_{n\ge1}\tfrac1{n2^n},\ \text{tail}<2^{-20}];
\]

%dsp-level:0
\[
\qquad \mathrm{BM.20}\ \mathrm{Rmk}\text{'s agreement row stands untouched};
\]

%dsp-level:0
\[
(7)\quad\textbf{the total energy assembly:}
\]

%dsp-level:0
\[
\qquad E_{\mathrm{total}}^2=m_{\mathrm{rest}}^2+p^2\ [\mathrm{BM.34}(4)\ \text{vacuum gauge}];
\]

%dsp-level:0
\[
\qquad E_{\mathrm{total}}=N\sqrt{m_{\mathrm{rest}}^2+p^2}-N^ip_i\ [\mathrm{Def\text{-}BM.10}(5)];
\]

%dsp-level:0
\[
\textbf{NR:}
\]

%dsp-level:0
\[
\qquad \boxed{\ E_{\mathrm{total}}=m_{\mathrm{rest}}+\mathrm{KE}=\frac{\Delta S_{\mathfrak L}}{\mathfrak h}+\frac{N\,\gamma^{ij}p_ip_j}{2m}+O(p^4)\ }
\]

%dsp-level:0
\[
\qquad [\text{machine: exact series assembly}];
\]

%dsp-level:0
\[
(8)\quad\textbf{status and scope:}
\]

%dsp-level:0
\[
\qquad \text{the reframing adds no content: it is an identity of derived reads}\ ((3),(4))\ \text{plus a name}\ [\text{A §0.2 bookkeeping}];
\]

%dsp-level:0
\[
\neg\text{a declaration},\quad\neg\text{a P4 member},\quad\text{§Cal-BM stays empty};
\]

%dsp-level:0
\[
\qquad (3)\text{'s balance identity holds exactly on the bidirectionally crystallized composites};
\]

%dsp-level:0
\[
\text{a unidirectional edge books}\ \ln2\ \text{against the pair freedom}\ \ln4\ (\text{half-balance}):
\]

%dsp-level:0
\[
\qquad \text{the general directed form is}\ (4)\text{'s}\ m_{\mathrm{hid}}\ \text{identity}\ [\mathrm{FT.44}(4)].
\]

\textbf{Rmk (three currencies, one reference point).} The ledger deficit counts the crystallized bits — $\ln2$ per directed edge; the branch-state deficit counts the uncrystallized channels at $\tfrac32$ nats each [FT.107(2)]; the angle deficit counts the holonomy against flatness [FT.54(4)]. All three are reads of structure against the same vacuum — the all-virtual entropy maximum, which is also the flat metric — and no theorem of Parts 1–4 converts one currency into another: item 7 holds the relating-equation question, and this tier does not touch it. The balance identity is the sharpest single row the corpus can state about why the rest read deserves the entropy name: on a bidirectional crystallization the ledger books exactly the entropy the vacuum pair carried, so the mass counts re-booked pair freedom, not destroyed information — $\ln4$ per pair, $2\ln2$ per channel, $E_{\mathrm{int}}\ln2$ in total, one $\mathfrak h$ away from an energy.

\textbf{Req.} C9, D0.1, D0.8, D2.9 (Parts 1–2); CT0.2, CT0.9 (Part 1); CT1.29, CT2.1 (Part 2); P1, P5 (Parts 1–2); Def-H.37(2), FT.47(1), FT.54(4), FT.54(5), FT.107(2), L3 (Part 3); Census FT.28, FT.42(1), FT.44(4), FT.45(1), FT.76(3) (row-wise); Def-BM.4, Def-BM.5, Def-BM.10, BM.17, BM.20(3), BM.34(2), BM.34(4) (this part); exact $\mathbb Q\cdot\ln2$ arithmetic, the exhaustive integer-grid Gibbs census, exact rational $\ln2$ enclosures, Laurent monomials [math; machine: \texttt{checks/\allowbreak{}verify\_\allowbreak{}H\_\allowbreak{}kinetic\_\allowbreak{}energy\_\allowbreak{}metric.py} §B].

\textbf{Pf.}

%dsp-level:0
\[
(1):\ S(\lambda)=-\sum_i\tfrac{k_i}8\ln\tfrac{k_i}8=\ln8-\tfrac18\sum_ik_i\ln k_i;
\]

%dsp-level:0
\[
\qquad S\le\ln4=\ln8-\ln2\iff\textstyle\sum_ik_i\ln k_i\ge8\ln2=\ln256\iff\prod_ik_i^{k_i}\ge256\ [\exp\ \text{strictly monotone}];
\]

%dsp-level:0
\[
\qquad \text{the census over all}\ 165\ \text{compositions of}\ 8\ \text{into}\ 4\ \text{parts: integers only}\ [\text{machine}]
\]

%dsp-level:0
\[
(2):\ \mathrm{BM.34}(2)\ \text{re-cited}\ [\text{machine}];\qquad\Theta_R=\bar KA=0\ \text{at}\ \bar K=0\ [\mathrm{FT.47}(1)\ \text{re-cited}]
\]

%dsp-level:0
\[
(3):\ S_{\mathfrak L}=n_{\mathcal R}\ln2\ [\mathrm{C9};\ \mathrm{CT0.9};\ \mathrm{CT1.29}];
\]

%dsp-level:0
\[
\qquad n_{\mathcal R}=E_{\mathrm{int}}\ [\mathrm{D0.8};\ \mathrm{Def\text{-}BM.5}(4)\ \text{pattern}];
\]

%dsp-level:0
\[
\qquad 2\ln2=\ln4\ \text{and}\ E_{\mathrm{bidir}}\cdot2\ln2=E_{\mathrm{int}}\ln2\ \text{at}\ E_{\mathrm{int}}=2E_{\mathrm{bidir}}\ [\text{machine:}\ \mathbb Q\cdot\ln2]
\]

%dsp-level:0
\[
(4):\ \mathrm{FT.44}(4)\ \text{substituted}\ [\text{machine}];
\]

%dsp-level:0
\[
\qquad \text{the exponents:}\ 0-(-\tfrac12)=\tfrac12\ [\mathrm{CT2.1};\ \text{machine}];
\]

%dsp-level:0
\[
\qquad c=1\ [\mathrm{D2.9}]:\ \text{the restored}\ c^2\ \text{carries}\ \mathfrak q\text{-exponent}\ 0\ [\text{machine}]
\]

%dsp-level:0
\[
(5):\ \mathrm{Def\text{-}H.37}(2)/\mathrm{FT.54}(4)\ \text{re-cited with}\ \mathrm{FT.47}(1)\text{'s}\ \Theta_R;
\]

%dsp-level:0
\[
\qquad \text{the range and the witness re-cited}\ [\mathrm{FT.76}(3);\ \mathrm{FT.54}(5):\ (2-\sqrt2)(2+\sqrt2)=2\ \text{in}\ \mathbb Q(\sqrt2)]
\]

%dsp-level:0
\[
(6):\ \text{the witness evaluations:}\ 2\ln2\ \text{vs}\ \tfrac32(2-2)=0;
\]

%dsp-level:0
\[
\qquad 0\ \text{vs}\ \tfrac32(2-0)=3\ [\text{machine}];
\]

%dsp-level:0
\[
\qquad \text{the enclosure: the partial sum}\ \textstyle\sum_{n=1}^{20}\tfrac1{n2^n}\ \text{with tail}<2^{-20}\implies\tfrac12<\ln2<1\ [\text{machine: exact Fractions}]\implies2\ln2<2<3
\]

%dsp-level:0
\[
(7):\ \mathrm{BM.34}(4)\text{'s gauge row}\ +\ (5)\text{'s series composed}\ [\text{machine}];\qquad(8):\ \text{the scope record}.\ \blacksquare
\]

\textbf{Phys.} The rest mass of the corpus has always been an entropy in disguise, and the disguise comes off here. The vacuum is the unique state of maximal raw information, and it is flat: maximal entropy, zero defect — one reference point in two currencies. Structure is a lowering against that reference, and the ledger books the lowering in bits, one per crystallized directed edge. The balance identity is why the reframing earns its name: a bidirectionally crystallized channel books exactly the entropy its vacuum pair carried — two bits, $\ln4$ — so the ledger deficit is a complete re-booking, not a residue, and division by the single free scale turns it into the informational mass the floors have always carried. The reframing costs nothing: same monomial, same dimensions, no new entity, no declaration. What it clarifies is the family structure of the mass reads — the pole mass reads the angle deficit against the same flat vacuum, the ensemble reads the branch-state deficit against the same all-virtual reference — three currencies, one reference point, no exchange rate derived. The total energy of a moving composite is then the sum of two derived reads: the deficit against the vacuum, at rest; and the metric's kinetic term, in motion.

\end{theorem}

\begin{theorem}[{\texorpdfstring{BM.36 — The force-structure ladder: the mask dimensions, the transport channel, and the inter-family constraint}{BM.36 - The force-structure ladder: the mask dimensions, the transport channel, and the inter-family constraint}}]\label{tag:BM-36}
\textbf{Thm.}

%dsp-level:0
\[
(1)\quad\textbf{the flag and its increments:}
\]

%dsp-level:0
\[
\qquad \dim\mathcal K_n=(6,16,17,25)\ \ \text{strict}\ [\mathrm{CT0.14}];
\]

%dsp-level:0
\[
\qquad (G_5{+}W_3{+}H_1)+(J^+_3{+}J^-_3)+D^+_1=16=C_{\mathrm{cap}}\ [\mathrm{CT1.12a}(1);\ \mathrm{CT0.12}];
\]

%dsp-level:0
\[
\text{rung}\ 1\to2:\quad+1=\dim Z\bigl(\iota_2(\mathfrak u(2))\bigr)=\dim\mathcal K^{(\mathrm{EM})}\ [\mathrm{D2.12};\ \mathrm{CT0.14}];
\]

%dsp-level:0
\[
\qquad \text{rung}\ 2\to3:\quad+8=\dim\mathfrak{su}(3)=\dim\mathcal K^{(\mathrm{strong})}\ [\mathrm{D2.12};\ \mathrm{CT3.8}];
\]

%dsp-level:0
\[
\text{the observable bound:}\quad\dim\mathcal K_\Sigma\le1+5+1+3+8=18\ [\mathrm{FT.26}(3)]
\]

%dsp-level:0
\[
\qquad [\text{machine: exact integer censuses}];
\]

%dsp-level:0
\[
(2)\quad\textbf{the derived structure ratios:}
\]

%dsp-level:0
\[
\qquad \frac{\dim\mathcal K^{(\mathrm{EM})}}{\dim F_W}=\frac13\quad[F_W=F_{W^+}\oplus F_{W^-}\oplus F_Z,\ \dim=3;\ \mathrm{C8}];
\]

%dsp-level:0
\[
\qquad \frac{\dim\mathcal K^{(\mathrm{strong})}}{\dim\mathcal K^{(\mathrm{EM})}}=\frac81=8;
\]

%dsp-level:0
\[
\qquad \frac{\dim\mathcal K^{(\mathrm{EM})}}{\dim\mathcal K_3}=\frac1{25}
\]

%dsp-level:0
\[
\qquad [\text{machine: exact Fractions}];
\]

%dsp-level:0
\[
\textbf{these are the alphabet-capacity ratios of the rung-activation order}\ [\mathrm{CT0.13}]\ —\ \text{the only force-structure ratios the corpus derives};
\]

%dsp-level:0
\[
(3)\quad\textbf{the strength bar:}
\]

%dsp-level:0
\[
\qquad \neg\exists\ \text{a theorem fixing any numerical coupling strength}:
\]

%dsp-level:0
\[
\qquad (g,g')\ \text{are P4-class}\ [\mathrm{EW.33};\ \mathrm{P4}];
\]

%dsp-level:0
\[
\qquad \alpha^{(2)}_{\mathrm{EM}}\ \text{is an operator block read whose numerical value no theorem of the corpus fixes}\ [\mathrm{EM.17};\ \mathrm{EM.25}(8);\ \text{Part 2 §Cal-2's logged row}];
\]

%dsp-level:0
\[
\text{the identification}\ \alpha_X:=\text{a dimension ratio is}\ \neg\text{derived and}\ \neg\text{adopted}:
\]

%dsp-level:0
\[
\qquad \text{reading}\ (2)\text{'s rationals as coupling strengths is a category error of the}\ \pi\text{-against-}1/137\ \text{species}\ [\mathrm{FT.73}(5)\ \text{trap pattern}];
\]

%dsp-level:0
\[
[\text{machine: the text pins}\ —\ \text{the bar on record in this tag; zero numerical strength values in the tier}];
\]

%dsp-level:0
\[
(4)\quad\textbf{the transport channel (the mass-generation locus):}
\]

%dsp-level:0
\[
\qquad \operatorname{Stab}_{\mathfrak g_{\mathrm{EW}}}(h_0\Sigma)
\]

%dsp-level:0
\[
= \mathfrak g_{\mathrm{EW}}:\ \ \neg\text{adjoint breaking}\implies\text{mass generation is a transport effect on a carried seed,}\ \neg\text{a commutator effect on generators}\ [\mathrm{EW.15}];
\]

%dsp-level:0
\[
\text{the seed:}\quad\mathcal H_{\mathsf v}\cong\mathbb C^2,\quad W_i\mapsto\tfrac12\sigma_i,\quad Y\mapsto\tfrac12\sigma_0,\quad T_a\mathsf v=0\ [\mathrm{EW.32}];
\]

%dsp-level:0
\[
\qquad M^2_{AB}=\mu_\star^2\operatorname{Re}\langle\mathcal G_A\mathsf v,\mathcal G_B\mathsf v\rangle,\quad\ker M^2=\text{the stabilizer},\quad h_0=0\implies M^2\equiv0\ [\mathrm{EW.33}];
\]

%dsp-level:0
\[
\textbf{the exact recomputation in the weak frame}\ [\mathrm{EW.20}]:\qquad M^2=\frac{\mu_\star^2}4\begin{pmatrix}g^2&0&0&0\\0&g^2&0&0\\0&0&g^2&-gg'\\0&0&-gg'&g'^2\end{pmatrix}\qquad[\text{machine: Pauli spinors over}\ \mathbb Q(i)\otimes\mathbb Q[g,g']:\ \text{the real parts give the block matrix; the imaginary parts are antisymmetric and cancel in the symmetric form}\ c_Ac_B\ [\mathrm{EW.33}]];
\]

%dsp-level:0
\[
\text{the spectrum}\ \Bigl\{\frac{g^2}4,\ \frac{g^2}4,\ 0,\ \frac{g^2+g'^2}4\Bigr\}\mu_\star^2\ [\text{machine: the charpoly factorization exact in}\ \mathbb Q[g,g',x]];
\]

%dsp-level:0
\[
\qquad \text{the kernel line}\ c\propto(0,0,g',g)\ [\text{machine: exact annihilation}];
\]

%dsp-level:0
\[
\frac{M_Z^2}{M_W^2}=\frac{g^2+g'^2}{g^2}\ \ \text{exact in}\ \mathbb Q(g^2,g'^2)\ [\text{machine: cross-multiplied}];
\]

%dsp-level:0
\[
\qquad e=g\sin\theta_W\ [\mathrm{EW.35}]\ \text{stands untouched};
\]

%dsp-level:0
\[
(5)\quad\textbf{the inter-family constraint theorem:}
\]

%dsp-level:0
\[
\qquad \text{every fermion pole mass takes the form}\quad m^2_{\mathrm{pole}}(c)=\hat\mu^2\,\ell_c^2\,G(c),
\]

%dsp-level:0
\[
\qquad G(c)=\sum_{A,B}c_Ac_B\operatorname{Re}\langle\mathcal G_A\mathsf v,\mathcal G_B\mathsf v\rangle
\]

%dsp-level:0
\[
\qquad [\mathrm{Def\text{-}H.37}(1)]:
\]

%dsp-level:0
\[
\text{an inter-family mass ratio on this channel is a ratio of}\ G(c)\,
\]

%dsp-level:0
\[
\ell_c^2\ \text{values across transport profiles}\ [\text{structural consequence}];
\]

%dsp-level:0
\[
\qquad \textbf{the derived sandwich:}
\]

%dsp-level:0
\[
\qquad 0\ \le\ G(c)\ \le\ \mu_\star^2\,\frac{g^2+g'^2}4\,\|c\|^2
\]

%dsp-level:0
\[
\qquad [\text{machine: the difference form's PSD factorization}];
\]

%dsp-level:0
\[
\textbf{the exact massless line:}\qquad G(c)=0\iff\textstyle\sum_Ac_A\mathcal G_A\in\operatorname{Ann}(\mathsf v)\iff c\parallel(0,0,g',g)\qquad[\mathrm{EW.33};\ \text{machine}];
\]

%dsp-level:0
\[
\textbf{the range ceilings on the computed skeletons:}
\]

%dsp-level:0
\[
\qquad \text{the pole ladder lies in}\ [0,4\hat\mu'^2]\ [\mathrm{FT.76}(3)];
\]

%dsp-level:0
\[
\qquad \neg\exists\ \text{an assignment of defect reads reaching a multi-order family ratio}\ [\mathrm{FT.77},\ \text{the quarantined audit re-cited}];
\]

%dsp-level:0
\[
\text{a cover read spanning}\ 1{:}R\ \text{requires skeleton scale}\ V\ge\tfrac43R\ [\mathrm{FT.79}(6)\ \text{re-cited}]:
\]

%dsp-level:0
\[
\qquad \textbf{the dynamical range must come from the transport data}\ (\ell_c,c_A)\ —\ \neg\text{from the order-one structure of}\ (1)\text{--}(2);
\]

%dsp-level:0
\[
(6)\quad\textbf{the status row (the inter-family ratios):}
\]

%dsp-level:0
\[
\qquad \textbf{constrained,}\ \neg\textbf{derived:}
\]

%dsp-level:0
\[
\qquad \text{the channel}\ ((4)\text{'s transport}),
\]

%dsp-level:0
\[
\ \text{the form (the Gram)},
\]

%dsp-level:0
\[
\ \text{the kernel (the exact massless line)},
\]

%dsp-level:0
\[
\ \text{the sandwich}\ ((5)),
\]

%dsp-level:0
\[
\ \text{the ceilings}\ ((5))\ \text{are derived};
\]

%dsp-level:0
\[
\neg\exists\ \text{a map in Parts 1--4 assigning transport profiles}\ c_A\ \text{to family lines}\ [\text{inventory scan; machine: no clause of}\ \mathrm{Def\text{-}BM.6}\ \text{or of the channel tags outputs a}\ c_A]:
\]

%dsp-level:0
\[
\boxed{\ \textbf{the Yukawa-like coupling}\ —\ \text{the family-profile map and its scale}\ —\ \textbf{is application-layer content}\ —\ \neg\textbf{a structural absence of the base layer}\ };
\]

%dsp-level:0
\[
\text{its construction locus is}\ (4)\text{'s seed transport}\ [\mathrm{EW.32},\ \mathrm{EW.33},\ \mathrm{Def\text{-}H.37}(1)];
\]

%dsp-level:0
\[
\qquad \text{the Koide-form invariant}\ Q
\]

%dsp-level:0
\[
= \tfrac23\ [\mathrm{Def\text{-}BM.6}(4);\ \mathrm{FT.158}(4)]\ \text{stands as a structural invariant of the naming rows};
\]

%dsp-level:0
\[
\qquad \text{its mass-form bridge}\ (\sqrt m\propto r\ \text{vs}\ m\propto r)\ \text{is}\ \neg\text{adopted}\ [\mathrm{Def\text{-}BM.6}(5)].
\]

\textbf{Rmk (what the dimension ratios can and cannot carry).} The ladder $6\subset16\subset17\subset25$ and its ratios $(\tfrac13,\,8,\,\tfrac1{25})$ count activation capacity — how many generator directions the ledger can switch on at each rung. A strength read would require a coupling theorem relating those counts to resolvent weights; no such theorem exists in Parts 1–4, and the couplings the corpus does carry are free positive P4 numbers whose values no clause fixes. The ceilings of (5) sharpen the bar into a quantitative shape: order-one structure ratios cannot carry a multi-order hierarchy on the computed cells, so an inter-family mass generation must come from the transport data — profiles and lengths — whose derivation is application-layer construction, not from the flag.

\textbf{Req.} C8 (Part 1); D2.12 (Part 1); CT0.12, CT0.13, CT0.14 (Part 1); CT1.12a(1), CT3.8, EW.15, EW.20, EW.32, EW.33, EW.35, EM.17, EM.25(8) (Part 2); P4, P5 (Parts 1–2); Def-H.37(1), FT.26(3), FT.73(5), FT.77, FT.79(6), FT.158(4) (Part 3); Census FT.76(3) (row-wise); Def-BM.6 (this part); exact Pauli/Gram algebra over $\mathbb Q(i)\otimes\mathbb Q[g,g']$, charpoly factorization in $\mathbb Q[g,g',x]$, exact Fractions, text pins [math; machine: \texttt{checks/\allowbreak{}verify\_\allowbreak{}H\_\allowbreak{}kinetic\_\allowbreak{}energy\_\allowbreak{}metric.py} §C].

\textbf{Pf.}

%dsp-level:0
\[
(1):\ \mathrm{CT0.14}\text{'s flag re-cited};
\]

%dsp-level:0
\[
\qquad 3+3+5+3+1+1=16=C_{\mathrm{cap}},\quad16+1=17,\quad17+8=25,\quad6<16<17<25,\quad1+5+1+3+8=18\ [\text{machine}]
\]

%dsp-level:0
\[
(2):\ \dim F_W=3\ [\mathrm{C8}];\qquad\tfrac{1}{3},\ \tfrac{8}{1},\ \tfrac{1}{25}\ [\text{machine: exact Fractions}]
\]

%dsp-level:0
\[
(3):\ \mathrm{EW.33}\text{'s P4 row and Part 2 §Cal-2's logged}\ \alpha^{(2)}_{\mathrm{EM}}\ \text{row re-cited}\ [\text{text pin: ``no theorem of the corpus fixes its numerical value''}];
\]

%dsp-level:0
\[
\qquad \mathrm{FT.73}(5)\text{'s trap pattern re-cited}
\]

%dsp-level:0
\[
(4):\ \mathsf v=(0,1)^{\mathsf T}:\quad\sigma_1\mathsf v=(1,0)^{\mathsf T},\ \ \sigma_2\mathsf v=(-i,0)^{\mathsf T},\ \ \sigma_3\mathsf v=(0,-1)^{\mathsf T},\ \ \sigma_0\mathsf v=\mathsf v;
\]

%dsp-level:0
\[
\qquad \mathcal G=(gW_1,gW_2,gW_3,g'Y)\ [\mathrm{EW.33}]
\]

%dsp-level:0
\[
\operatorname{Re}\langle\mathcal G_A\mathsf v,\mathcal G_B\mathsf v\rangle=\tfrac14\,\bigl(\text{the displayed block matrix}\bigr)_{AB};
\]

%dsp-level:0
\[
\qquad \operatorname{Im}\langle\mathcal G_A\mathsf v,
\]

%dsp-level:0
\[
\mathcal G_B\mathsf v\rangle
\]

%dsp-level:0
\[
= -\operatorname{Im}\langle\mathcal G_B\mathsf v,
\]

%dsp-level:0
\[
\mathcal G_A\mathsf v\rangle\ \text{and cancels identically in}\ \textstyle\sum_{A,B}c_Ac_B(\cdot)\ [\mathrm{EW.33};\ \text{machine}];
\]

%dsp-level:0
\[
\qquad \det(xI-\mathrm{Gram})=\bigl(x-\tfrac{g^2}4\bigr)^2\,x\,\bigl(x-\tfrac{g^2+g'^2}4\bigr)\ [\text{machine: expanded and compared in}\ \mathbb Q[g,g',x]]
\]

%dsp-level:0
\[
\mathrm{Gram}\cdot(0,0,g',g)^{\mathsf T}=0\ [\text{machine}];
\]

%dsp-level:0
\[
\qquad \mu_\star^2=c_W\varkappa^2\beta^2h_0^2\big|_{h_0=0}=0\implies M^2\equiv0\ [\text{machine: substitution}];
\]

%dsp-level:0
\[
\qquad \tfrac{(g^2+g'^2)/4}{g^2/4}=\tfrac{g^2+g'^2}{g^2}\ [\text{machine: cross-multiplied}]
\]

%dsp-level:0
\[
(5):\ \mathrm{Def\text{-}H.37}(1)\text{'s form re-cited};
\]

%dsp-level:0
\[
\qquad \tfrac{g^2+g'^2}4I-\mathrm{Gram}:\ \ \text{the spectrum}\ \bigl\{\tfrac{g'^2}4,\tfrac{g'^2}4,\tfrac{g^2+g'^2}4,\tfrac{g^2}4\bigr\}\ [\text{machine: charpoly factorization}]\implies\text{the sandwich by the Rayleigh range};
\]

%dsp-level:0
\[
\qquad \mathrm{FT.76}(3),\ \mathrm{FT.77},\ \mathrm{FT.79}(6)\ \text{re-cited}\ (\neg\text{recomputed; the quarantined rows stay quarantined})
\]

%dsp-level:0
\[
(6):\ \text{the inventory scan of}\ \mathrm{Def\text{-}BM.6}\ \text{and the channel tags}\ \mathrm{FT.109}\text{--}\mathrm{FT.119},
\]

%dsp-level:0
\[
\ \mathrm{FT.142}\text{--}\mathrm{FT.145}:\ \neg\exists\ \text{a clause outputting a transport profile}\ [\text{machine: token scan}];\qquad \ \blacksquare
\]

\textbf{Phys.} The force side of the mass question splits cleanly into what the corpus derives and what it does not. Derived is the structure: the mask flag grows six to sixteen to seventeen to twenty-five by exact increments — one electromagnetic centre, eight gluons — and the only ratios the ladder supplies are capacity ratios: one third, eight, one twenty-fifth. Derived is also the channel: mass can only be a transport defect of a carried seed, because the scalar background breaks nothing adjointly, and the seed's Gram form comes with an exactly computed spectrum — two equal charged directions, one massive neutral direction, one exact kernel direction along which the transported family is massless whatever the couplings do. Not derived is any coupling value: no theorem turns a count into a strength, and the tier records the bar rather than let a rational drift across the gap. And not derived is the family map: nothing in Parts 1–4 assigns transport profiles to family lines, while the ceilings on the computed skeletons say that order-one structure cannot carry a multi-order hierarchy. The inter-family mass ratios are therefore constrained — channel, form, kernel, sandwich, ceiling — and open, with the construction locus named: the seed transport itself.

\end{theorem}

\begin{theorem}[{\texorpdfstring{BM.37 — The mass-route artifacts: the disjoint input alphabets and the closure of the counting trap}{BM.37 - The mass-route artifacts: the disjoint input alphabets and the closure of the counting trap}}]\label{tag:BM-37}
\textbf{Thm.}

%dsp-level:0
\[
(1)\quad\textbf{the three routes, three exact values:}
\]

%dsp-level:0
\[
\qquad (\mathrm i)\ \ \text{the normalized ensemble average}\ [\mathrm{Def\text{-}H.32}(6);\ \mathrm{FT.107}(3)]:\quad\text{neutral}{:}\text{charged}
\]

%dsp-level:0
\[
= 1\ (\text{degenerate})\ —\ \text{a uniform sector factor cancels at}\ \{1,2,3,15,20,\tfrac73\}\ [\mathrm{BM.30}(3);\ \text{machine}];
\]

%dsp-level:0
\[
(\mathrm{ii})\ \ \text{the total deficit}\ [\mathrm{BM.28}(4)(\mathrm{ii})]:
\]

%dsp-level:0
\[
\qquad \frac{W(0)}{W(+1)}=\frac{20}{15}=\frac43
\]

%dsp-level:0
\[
\qquad [\text{machine: the re-enumeration of}\ (2+z+z^{-1})^3\ \text{and the exhaustive}\ 27\text{-triple census}];
\]

%dsp-level:0
\[
(\mathrm{iii})\ \ \text{the amplitude layer}\ [\mathrm{Def\text{-}H.32}(2),(5);\ \mathrm{BM.28}(4)(\mathrm{iii})]:
\]

%dsp-level:0
\[
\qquad \frac{5336}{1196}=\frac{58}{13}\ (S{=}0),
\]

%dsp-level:0
\[
\qquad \frac{4796}{802}=\frac{2398}{401}\ (S{=}1)
\]

%dsp-level:0
\[
\qquad [\mathrm{FT.106}(3)\text{'s frozen integers; machine: exact reductions}];
\]

%dsp-level:0
\[
\text{pairwise unequal:}\qquad1\ \ne\ \frac43\ \ne\ \frac{58}{13}\ \ne\ \frac{2398}{401}\qquad[\text{machine: cross-multiplications in}\ \mathbb Z];
\]

%dsp-level:0
\[
(2)\quad\textbf{the disjointness theorem (the artifact classification):}
\]

%dsp-level:0
\[
\qquad \text{routes}\ (\mathrm i),
\]

%dsp-level:0
\[
(\mathrm{ii})\ \text{are measure functionals of}\ (\Gamma,R):\ \ \text{the weights}\ (\nu,\ t^{|R|},\ u^L,\ \#\mathrm{inst})\ \text{and the deficit}\ \Delta S\ [\mathrm{Def\text{-}H.32}(3),(6);\ \mathrm{FT.107}(2)];
\]

%dsp-level:0
\[
\text{route}\ (\mathrm{iii})\ \text{is the amplitude functional of}\ v_\Gamma(R)\ [\mathrm{Def\text{-}H.32}(2),(5)];
\]

%dsp-level:0
\[
\qquad \text{the dispersion read is the metric functional of}\ (g_{ij},N,N^i,p_i,m)\ [\mathrm{Def\text{-}BM.10}(2)\text{--}(5)];
\]

%dsp-level:0
\[
\textbf{the machine token census on the canonical defining strings:}
\]

%dsp-level:0
\[
\qquad \text{metric-alphabet occurrences inside the defining clauses of routes}\ (\mathrm i)\text{--}(\mathrm{iii}):\ \ \mathbf 0;
\]

%dsp-level:0
\[
\text{measure-alphabet occurrences inside}\ \mathrm{Def\text{-}BM.10}(2),
\]

%dsp-level:0
\[
(3),
\]

%dsp-level:0
\[
(5)\ \text{(the kinematic clauses; clause}\ (4)\text{'s rest input is the shared symbol}\ m):\ \ \mathbf 0;
\]

%dsp-level:0
\[
\boxed{\ \textbf{the input alphabets are disjoint:}\quad\neg\exists\ \text{a route output that is a function of metric data};\qquad\neg\exists\ \text{a dispersion energy that is a function of measure data}\ }
\]

%dsp-level:0
\[
(3)\quad\textbf{the trap and its closure:}
\]

%dsp-level:0
\[
\qquad \textbf{the counting trap}\ :
\]

%dsp-level:0
\[
= \ \text{reading a route ratio}\ \bigl(1,\ \tfrac43,\ \tfrac{58}{13},\ \tfrac{2398}{401}\bigr)\ \text{or the Koide-form invariant}\ Q
\]

%dsp-level:0
\[
= \tfrac23\ [\mathrm{Def\text{-}BM.6}(4);\ \mathrm{FT.158}(4)]\ \text{as a physical mass ratio}\ [\mathrm{FT.89}\text{'s category-trap species}];
\]

%dsp-level:0
\[
\textbf{barred by three independent rows:}
\]

%dsp-level:0
\[
\qquad (\mathrm a)\ \ \text{the disjointness}\ [(2)]:\ \ \text{a route output carries no momentum dependence and cannot transform as a dispersion energy};
\]

%dsp-level:0
\[
(\mathrm b)\ \ \text{amplitude}\ne\text{mass}\ [\mathrm{FT.79}(5)]:\qquad\text{route}\ (\mathrm{iii})\text{'s bridge is item-13 content,}\ \neg\text{constructed}\;
\]

%dsp-level:0
\[
(\mathrm c)\ \ \text{the route choice is a read specification,}\ \neg\text{a theorem}\ [\mathrm{BM.30}(5)]:
\]

%dsp-level:0
\[
\qquad \text{no route is adopted, and none is retracted};
\]

%dsp-level:0
\[
Q = \tfrac23
\]

%dsp-level:0
\[
= \tfrac{3-1}3\ \text{is the doublet/triplet decomposition ratio of the spatial seal triad}\ [\mathrm{FT.158}(5)]:
\]

%dsp-level:0
\[
\qquad \text{a structural invariant of the naming rows, whose mass-form bridge stays}\ \neg\text{adopted}\ [\mathrm{Def\text{-}BM.6}(5)];
\]

%dsp-level:0
\[
(4)\quad\textbf{the completion row:}\qquad\text{the physical mass-energy of a sealed composite is the dispersion read:}
\]

%dsp-level:0
\[
E=N\sqrt{m_{\mathrm{rest}}^2+p^2}-N^ip_i,\qquad m_{\mathrm{rest}}=\frac{\Delta S_{\mathfrak L}}{\mathfrak h}\qquad[\mathrm{BM.34}(4);\ \mathrm{BM.35}(4),(7)];
\]

%dsp-level:0
\[
\text{the routes are static}\ (p\text{-free})\ \text{slices: each evaluates an ensemble with no momentum argument}\ [(2)];
\]

%dsp-level:0
\[
\qquad \textbf{KE}\ —\ \text{the metric's motion term}\ —\ \textbf{is what no route can express};
\]

%dsp-level:0
\[
\text{at}\ p
\]

%dsp-level:0
\[
= 0\ \text{the dispersion rest term coincides on-shell with the curvature-flux and the informational reads}\ [\mathrm{BM.34}(7)]:
\]

%dsp-level:0
\[
\qquad \text{the routes price the ensemble measure's restrictions around the same rest object};
\]

%dsp-level:0
\[
\mathrm{BM.30}(2)\text{'s splitting and}\ \mathrm{BM.30}(3)\text{'s degeneracy stand derived and exact as}\ \textbf{measure rows}\ —\ \text{the closure classifies, it does not retract};
\]

%dsp-level:0
\[
(5)\quad\textbf{the standing rules:}
\]

%dsp-level:0
\[
\qquad \neg\text{a P4 addition}\ [\text{§Cal-BM stays empty}];
\]

%dsp-level:0
\[
\qquad \neg\text{a new dimensionful scale}\ [\mathrm{P1}];
\]

%dsp-level:0
\[
\qquad \neg\text{a measured value and}\ \neg\text{a decimal numeral in this tier's text}\ [\mathrm{P5};\ \text{machine scans}];
\]

%dsp-level:0
\[
\neg\text{a new declaration};
\]

%dsp-level:0
\[
\qquad \text{species names confined to the identification layers}\ [\mathrm{Def\text{-}BM.3}(7),\ \mathrm{Def\text{-}BM.6}(7)];
\]

%dsp-level:0
\[
\qquad \text{topology labels are internal bookkeeping}\ [\mathrm{Def\text{-}H.32}(4)];
\]

%dsp-level:0
\[
(6)\quad\textbf{the adjudication:}
\]

%dsp-level:0
\[
\qquad \boxed{\ \text{the three mass-route ratios are counting and amplitude artifacts}\ —\ \text{exact as measure prices, barred as physical masses}\ };
\]

%dsp-level:0
\[
\text{the counting trap is closed at the category level by the disjoint input alphabets}\ [(2)];
\]

%dsp-level:0
\[
\text{the physical mass-energy read is the metric dispersion route: the rest term is the vacuum-balanced entropy}\ [\mathrm{BM.35}],
\]

%dsp-level:0
\[
\ \text{the motion term is the shift-transported kinetic energy}\ [\mathrm{BM.34}];
\]

%dsp-level:0
\[
\text{the inter-family content is application-layer}\ [\mathrm{BM.36}(6)].
\]

\textbf{Rmk (what the closure does not do).} The closure retracts no route: every route value stands as a derived exact rational, and BM.30(5)'s read-choice row is untouched — adopting any route as "the mass" still requires a logged declaration. What the closure establishes is the classification: the three functionals price measures and amplitudes, and a measure price is not a dispersion energy however clean its rational looks. The Koide-form invariant belongs to the same classification — it is the decomposition ratio the spatial seal triad forces [FT.158(5)] — and converting it into a mass ratio needs exactly the mass-form bridge — application-layer content. The physical reading runs the other way: the rest mass is the vacuum-deficit read, the motion is the metric's kinetic term, and the route ratios stay what they always were — exact prices of how the charge constraint and the sign sum filter one ensemble.

\textbf{Req.} P1, P4, P5 (Parts 1–2); Def-H.32(2), Def-H.32(3), Def-H.32(4), Def-H.32(5), Def-H.32(6), FT.79(5), FT.89, FT.106(3), FT.107(2), FT.107(3), FT.158(4), FT.158(5) (Part 3); Def-BM.3(7), Def-BM.6(4), Def-BM.6(5), Def-BM.6(7), Def-BM.10, BM.28(2), BM.28(4), BM.30(1), BM.30(2), BM.30(3), BM.30(5), BM.34(4), BM.34(7), BM.35(4), BM.35(7), BM.36(6) (this part); exact rational reductions, the exhaustive triple census, generating-function convolution, token censuses [math; machine: \texttt{checks/\allowbreak{}verify\_\allowbreak{}H\_\allowbreak{}kinetic\_\allowbreak{}energy\_\allowbreak{}metric.py} §D].

\textbf{Pf.}

%dsp-level:0
\[
(1):\ (\mathrm i)\ \ \tfrac{\sum_Rc\,v(R)\Delta S(R)}{\sum_Rc\,v(R)}
\]

%dsp-level:0
\[
= \tfrac{\sum_Rv(R)\Delta S(R)}{\sum_Rv(R)}\ \ \forall c
\]

%dsp-level:0
\[
\ne 0\ [\text{machine at}\ c\in\{1,2,3,15,20,\tfrac73\}\ \text{on rational witness data}]
\]

%dsp-level:0
\[
(\mathrm{ii}):\ W(Q)=\bigl[z^Q\bigr](2+z+z^{-1})^3\ \text{re-enumerated}:
\]

%dsp-level:0
\[
\quad W=(1,6,15,20,15,6,1),
\]

%dsp-level:0
\[
\quad W(0)=8+12=20,
\]

%dsp-level:0
\[
\quad W(\pm1)=12+3=15,\quad\tfrac{20}{15}=\tfrac43\ [\text{machine}];
\]

%dsp-level:0
\[
\qquad (\mathrm{iii}):\ \tfrac{5336}{1196}=\tfrac{58}{13}\ (5336\cdot13=1196\cdot58),\ \ \tfrac{4796}{802}=\tfrac{2398}{401}\ [\text{machine: exact}]
\]

%dsp-level:0
\[
\text{the pairwise inequalities:}\ 3\cdot1\ne4\cdot1,
\]

%dsp-level:0
\[
\quad 4\cdot13\ne3\cdot58,
\]

%dsp-level:0
\[
\quad 58\cdot401\ne13\cdot2398,\quad\text{and the remaining pairs}\ [\text{machine: integer cross-multiplications}]
\]

%dsp-level:0
\[
(2):\ \text{the canonical defining strings extracted from the canon files}\ [\mathrm{Def\text{-}H.32}(2),(3),(5),(6);\ \mathrm{FT.107}(2),(3);\ \mathrm{BM.28}(4);\ \mathrm{Def\text{-}BM.10}(2)\text{--}(5)];
\]

%dsp-level:0
\[
\qquad \text{the metric alphabet}\ \{\Omega,\varkappa,\beta,w_i,q_{ij},\gamma,N,N^i,p_i,g_{0i},\text{lapse},\text{shift}\}:\ 0\ \text{occurrences in the route clauses}\ [\text{machine}]
\]

%dsp-level:0
\[
\text{the measure alphabet}\ \{W(Q),\#\mathrm{inst},\nu,v_\Gamma,\Delta S,t^{|R|},u^L,Z^{\bullet},\text{ensemble},\text{flicker}\}:\ 0\ \text{occurrences in}\ \mathrm{Def\text{-}BM.10}(2)\text{--}(5)\ \text{beyond the shared symbol}\ m\ [\text{machine}]
\]

%dsp-level:0
\[
(3):\ (\mathrm a)\ [(2)\ \text{conjoined with}\ \mathrm{Def\text{-}BM.10}(5)\text{'s}\ p\text{-dependence}];
\]

%dsp-level:0
\[
\qquad (\mathrm b)\ \mathrm{FT.79}(5)\ \text{re-cited};
\]

%dsp-level:0
\[
\qquad (\mathrm c)\ \mathrm{BM.30}(5)\ \text{re-cited};
\]

%dsp-level:0
\[
\qquad Q=\tfrac23\ \text{and}\ \tfrac{3-1}3\ [\mathrm{FT.158}(5);\ \mathrm{Def\text{-}BM.6}(4)\ \text{re-cited}]
\]

%dsp-level:0
\[
(4):\ \mathrm{BM.34}(4),\ \mathrm{BM.35}(4),(7)\ \text{conjoined};
\]

%dsp-level:0
\[
\qquad \mathrm{BM.34}(7)\text{'s on-shell agreement row re-cited};
\]

%dsp-level:0
\[
\qquad (5):\ \text{the standing-rule record}\ [\mathrm{BM.33}(5)\ \text{pattern}]\ +\ \text{the machine scans}
\]

%dsp-level:0
\[
(6):\ \text{conjunction of}\ (1)\text{--}(5).\ \blacksquare
\]

\textbf{Phys.} Three exact rationals — one, four thirds, fifty-eight thirteenths — came out of the charge-constrained ensemble, and the temptation is to read them as mass ratios. The corpus closes that trap the only way it closes anything: by classification. Each of the three routes is a functional of the ensemble's combinatorics alone — pattern counts, norm weights, sign-summed amplitudes — and a machine census of the defining clauses shows their input alphabet contains not one metric token, while the dispersion read's definition contains not one measure token beyond the shared rest symbol. Disjoint inputs, disjoint outputs: no route value can transform as a dispersion energy, because no route has ever seen a momentum. What the routes price is the measure — how the charge constraint rescales the probability space, how the sign sum weights the amplitudes — and those prices stand, exact and unretracted. The physical mass-energy is assembled elsewhere: the rest term from the vacuum-balanced entropy, the motion term from the metric's kinetic energy, the two joined by the mass-shell identity. The Koide-form invariant joins the same classification — a decomposition ratio the spatial triad forces, waiting on a mass-form bridge that item 13 holds open. The counting was never wrong; reading the counts as masses was.

\end{theorem}

\medskip\noindent\textit{\texorpdfstring{BM.38 — The status of the metric-energy tier: adjudication and scope}{BM.38 - The status of the metric-energy tier: adjudication and scope}.}\ Structural-constraints record: see Appendix B, Section B.1.20 of the Addendum.\par\medskip

\phantomsection
\section[Tier BM.11 — The Neutron Decay Vertex: the Family Transport…]{\texorpdfstring{Tier BM.11 — The Neutron Decay Vertex: the Family Transport Profiles, the Rung-1 Isospin Row, the Forced Neutral Member, the Angular-Momentum Budget, and the Full Kinematic Balance}{Tier BM.11 - The Neutron Decay Vertex: the Family Transport Profiles, the Rung-1 Isospin Row, the Forced Neutral Member, the Angular-Momentum Budget, and the Full Kinematic Balance}}

\begin{definition}[{\texorpdfstring{Def-BM.11 — The family transport-profile map (the constructed bridge)}{Def-BM.11 - The family transport-profile map (the constructed bridge)}}]\label{tag:Def-BM-11}
\textbf{Def.} A construction on EW.32/EW.33's seed transport and FT.158(1)'s inner $C_3$ rotation; no new axiom, no new dimensionful scale, no Yukawa-type object — none exists in the inventory [EW.15; FT.145(3)]. Local notation: $\mathrm{Gr}_{AB}:=\operatorname{Re}\langle\mathcal G_A\mathsf v,\mathcal G_B\mathsf v\rangle$ (EW.33's Gram data).

%dsp-level:0
\[
(1)\quad\textbf{the seed data:}
\]

%dsp-level:0
\[
\qquad \mathcal G=(gW_1,\ gW_2,\ gW_3,\ g'Y)\ [\mathrm{EW.33}];
\]

%dsp-level:0
\[
\qquad \mathsf v:\ \ W_3\mathsf v=-\tfrac12\mathsf v,\ \ Y\mathsf v=+\tfrac12\mathsf v\ [\mathrm{EW.20}]:
\]

%dsp-level:0
\[
\qquad \textbf{the hypercharge coordinate}\ y_0=\tfrac12\ \textbf{is derived}\ [\mathrm{EW.20}:\ \text{trace normalization},\ \neg\text{assigned}];
\]

%dsp-level:0
\[
(2)\quad\textbf{the family rotation:}
\]

%dsp-level:0
\[
\qquad U=\exp\!\bigl(-\tfrac{2\pi i}3\,\hat n\cdot W\bigr)=E_0-\tfrac12\mathcal W^2-\tfrac i2\mathcal W,\quad\mathcal W=W_1{+}W_2{+}W_3:
\]

%dsp-level:0
\[
\qquad UW_iU^\dagger=W_{i+1},\quad U^3=E_0\ [\mathrm{FT.158}(1);\ \text{machine: exact over}\ \mathbb Q(i),\ \text{via}\ \mathcal W^3=3\mathcal W];
\]

%dsp-level:0
\[
Y\ \text{is the centre of}\ \iota_2(\mathfrak u(2))\implies[Y,W_i]=0\implies UYU^\dagger=Y\ [\mathrm{EW.8};\ \mathrm{D2.12};\ \text{machine}]:
\]

%dsp-level:0
\[
\qquad \textbf{the}\ C_3\ \textbf{cycles the weak triplet coordinates of a transport profile and fixes the hypercharge coordinate};
\]

%dsp-level:0
\[
(3)\quad\textbf{the profile class:}
\]

%dsp-level:0
\[
\qquad \text{every}\ C_3\text{-covariant triplet profile takes the cosine form}\ \lambda_i
\]

%dsp-level:0
\[
= A+B\cos(\theta+\tfrac{2\pi i}3)\ [\mathrm{FT.111}(2)];
\]

%dsp-level:0
\[
\textbf{unit normalization (constitutive convention, the Def-BM.4(1) pattern):}\quad\text{the weak coordinates carry radius}\ 1;
\]

%dsp-level:0
\[
\textbf{the base profile:}\qquad c(0):=\bigl(\cos\delta,\ \cos(\delta+\tfrac{2\pi}3),\ \cos(\delta+\tfrac{4\pi}3),\ y_0\bigr)\in\mathbb R^4;
\]

%dsp-level:0
\[
\textbf{the orbit:}
\]

%dsp-level:0
\[
\qquad c(i):=\operatorname{Ad}(U^i)\,c(0)\ =\ \text{the cyclic shift of the weak triplet,}\ y_0\ \text{fixed}\ [(2);\ \text{machine}]:
\]

%dsp-level:0
\[
\qquad c(i)_3=\cos\bigl(\delta+\tfrac{2\pi(2-i)}3\bigr);
\]

%dsp-level:0
\[
\text{the}\ W_3\ \text{slot cycles through the triplet};
\]

%dsp-level:0
\[
\qquad \text{the Gram-value set}\ \bigl\{M^2\bigl(A_0-B_0\cos(\delta-\tfrac{2\pi j}3)\bigr)\bigr\}_j\ \text{is the orbit's content under}\ (4)\text{'s orbit-label convention};
\]

%dsp-level:0
\[
(4)\quad\textbf{the angle identification (the single declared clause):}
\]

%dsp-level:0
\[
\qquad \delta:=\delta_{\mathrm{channel}}\ —\ \text{per channel:}\quad\ell:\ \tfrac29\ [\mathrm{FT.112}];
\]

%dsp-level:0
\[
\qquad q:\ \tfrac{2}{27}\ [\mathrm{FT.119}(1)];
\]

%dsp-level:0
\[
\qquad \nu:\ \tfrac18\ [\mathrm{FT.142}(5)];
\]

%dsp-level:0
\[
\text{the orbit-label orientation convention}\ (\delta\mapsto-\delta\ \text{in the}\ W_3\ \text{slot})\ \text{is internal bookkeeping};
\]

%dsp-level:0
\[
\qquad \text{the lines attach to}\ \mathrm{Def\text{-}BM.6}\text{'s named rows by its own spectral order};
\]

%dsp-level:0
\[
(5)\quad\textbf{the mass evaluation:}
\]

%dsp-level:0
\[
\qquad m^2(\text{line}\ i)=\hat\mu^2\,\ell_c^2\,G\bigl(c(i)\bigr),
\]

%dsp-level:0
\[
\qquad G(c)=\sum_{A,B}c_Ac_B\,\mathrm{Gr}_{AB}
\]

%dsp-level:0
\[
\qquad [\mathrm{Def\text{-}H.37}(1);\ \mathrm{EW.33}]:
\]

%dsp-level:0
\[
\boxed{\ m^2_i=\hat\mu^2\ell_c^2\Bigl[\Bigl(\frac{3g^2}8+\frac{g'^2}{16}\Bigr)-\frac{gg'}4\,\cos\Bigl(\delta-\frac{2\pi i}3\Bigr)\Bigr]\ }=:M^2\bigl(A_0-B_0\cos\varphi_i\bigr),
\]

%dsp-level:0
\[
\qquad A_0=\frac{3g^2}8+\frac{g'^2}{16},\quad B_0=\frac{gg'}4
\]

%dsp-level:0
\[
\qquad [\text{machine: exact in}\ \mathbb Q[g,g']\otimes\mathbb Q[c,s,r_3]/(c^2{+}s^2{-}1,r_3^2{-}3)];
\]

%dsp-level:0
\[
\textbf{the affine form:}
\]

%dsp-level:0
\[
\qquad r_2\cdot m^2_i=M^2\bigl(A_0r_2+B_0-B_0\,r_i\bigr),
\]

%dsp-level:0
\[
\qquad r_i=1+r_2\cos\Bigl(\delta+\frac{2\pi i}3\Bigr),\quad r_2^2=2
\]

%dsp-level:0
\[
\qquad [\mathrm{FT.113}\ \text{vocabulary rows};\ \text{machine}];
\]

%dsp-level:0
\[
(6)\quad\textbf{interface status:}
\]

%dsp-level:0
\[
\qquad \textbf{derived:}\ \text{the orbit structure}\ [(2)],
\]

%dsp-level:0
\[
\ \text{the cosine shape}\ [\mathrm{FT.111}(2)],
\]

%dsp-level:0
\[
\ \text{the Gram evaluation and the affine form}\ [(5);\ \text{machine}],
\]

%dsp-level:0
\[
\ \text{positivity, the splitting law, the sum rule, the order inversion}\ [\mathrm{BM.39}];
\]

%dsp-level:0
\[
\textbf{non-derived:}\ \text{exactly one angle identification per channel}\ (4)\ \text{and the unit normalization}\ (3);
\]

%dsp-level:0
\[
\qquad \textbf{scales:}\ \hat\mu^2\in\mathrm{Cal\text{-}F}\ (\text{logged});\ \ \ell_c\ \text{is chain data}\ [\mathrm{D1.4}];\ \ (g,g')\ \text{are P4}\ [\mathrm{EW.33}]:
\]

%dsp-level:0
\[
\qquad \neg\text{a new parameter},\ \text{§Cal-BM stays empty};
\]

%dsp-level:0
\[
\text{declared}\ (\text{A-BM.4})\ \text{before use};
\]

%dsp-level:0
\[
(7)\quad\textbf{the locus clause:}
\]

%dsp-level:0
\[
\qquad \text{the construction lives in the pole/transport read species}\ [\mathrm{Def\text{-}H.37}(1)];
\]

%dsp-level:0
\[
\qquad \mathrm{FT.113}/\mathrm{FT.143}(3)\text{'s inertial vocabulary}\ m^2=\mu r^2\ \text{is a distinct read species on the same lines}:
\]

%dsp-level:0
\[
\text{the two species order the lines oppositely}\ [\mathrm{BM.39}(5)];
\]

%dsp-level:0
\[
\qquad \text{the relating equation is item-7 content, untouched}\;
\]

%dsp-level:0
\[
\textbf{BM.36}(6)\textbf{'s gap label is discharged:}\quad\text{the family-profile map exists as this construction};
\]

%dsp-level:0
\[
\qquad \text{the scale data is P4-class calibration}\ (\neg\text{declared}),\ \neg\text{the map}.
\]

\textbf{Rmk (what the construction is).} The map answers BM.36(6)'s named gap with derived material only: the rotation is FT.158(1)'s — the same $C_3$ that forces the channel spectra, inner in the weak algebra; the shape is FT.111(2)'s — every $C_3$-covariant triplet is a cosine triplet; the evaluation is Def-H.37(1)'s Gram transport; the hypercharge coordinate is EW.20's derived $\tfrac12$. The single declared content is that the profile's phase is the channel's own angle — one identification per channel, and the angles are derived rationals. The payoff is structural: the intra-family splitting is exactly $gg'$ — the weak–hypercharge mixed strength — acting on the transport phase, which is the derived form of "the relative force strength acting on the seed's transport profile." No commutator breaking enters anywhere [EW.15], and no Yukawa-type term exists to enter.

\textbf{Req.} D1.4, D2.12 (Parts 1–2); EW.8, EW.20, EW.32, EW.33 (Part 2); P4, P5 (Parts 1–2); Def-H.37(1), FT.111(2), FT.112, FT.113, FT.119(1), FT.142(5), FT.143, FT.145(3), FT.158(1) (Part 3); Def-BM.4(1), Def-BM.6, BM.36(4), BM.36(5), BM.36(6) (this part); Cal-F; the declaration addendum A-BM.4; exact quotient-ring trigonometric identities in $\mathbb Q[g,g',c,s,r_3,r_2]/(c^2{+}s^2{-}1,\ r_3^2{-}3,\ r_2^2{-}2)$, Gram algebra over $\mathbb Q[g,g']$ [math; machine: \texttt{checks/\allowbreak{}verify\_\allowbreak{}H\_\allowbreak{}neutron\_\allowbreak{}decay\_\allowbreak{}vertex.py} §A].

\textbf{Pf.}

%dsp-level:0
\[
\text{Definition block: the constitutive clauses are}\ (3)\text{'s unit normalization and}\ (4)\text{'s angle identification};
\]

%dsp-level:0
\[
\qquad \text{every displayed evaluation is a machine-verified exact identity}\ [\text{§A}];
\]

%dsp-level:0
\[
\qquad (6),(7)\ \text{are the status and locus records}.\ \blacksquare
\]

\textbf{Phys.} The map the corpus named as its open construction is built, and it is made of pieces the derivations already produced. The family rotation is the inner $C_3$ of the weak algebra — the same rotation that forces the cosine shape of every channel spectrum — so a transport profile carried around the orbit is automatically a cosine triplet with the channel's own angle. Feed the triplet through the seed's Gram form and the mass-squared of each line comes out affine in the channel's spectral invariant: a constant set by the couplings, minus the mixed weak–hypercharge strength times the line's cosine. The only thing declared is the identification of the profile's phase with the channel angle; everything else — the orbit, the shape, the evaluation, the positivity, the sum rule — is forced. What the construction does not do is convert currencies: the transport species and the inertial $m=\mu r^2$ vocabulary still wait on the relating equation, and the overall scales remain the undeclared internal numbers they have always been.

\end{definition}

\begin{theorem}[{\texorpdfstring{BM.39 — The family mass form: positivity, the splitting law, the sum rule, the order inversion, and the area ladder}{BM.39 - The family mass form: positivity, the splitting law, the sum rule, the order inversion, and the area ladder}}]\label{tag:BM-39}
\textbf{Thm.}

%dsp-level:0
\[
(1)\quad\textbf{the exact evaluation and positivity:}
\]

%dsp-level:0
\[
\qquad m^2_i=M^2\bigl(A_0-B_0\cos\varphi_i\bigr),
\]

%dsp-level:0
\[
\qquad M^2:=\hat\mu^2\ell_c^2,\quad A_0=\frac{3g^2}8+\frac{g'^2}{16},\quad B_0=\frac{gg'}4
\]

%dsp-level:0
\[
\qquad [\mathrm{Def\text{-}BM.11}(5);\ \text{machine}];
\]

%dsp-level:0
\[
A_0-B_0=\frac38\Bigl(g-\frac{g'}3\Bigr)^2+\frac{g'^2}{48},
\]

%dsp-level:0
\[
\qquad A_0+B_0=\frac38\Bigl(g+\frac{g'}3\Bigr)^2+\frac{g'^2}{48}
\]

%dsp-level:0
\[
\qquad [\text{machine: exact sum-of-squares decompositions}];
\]

%dsp-level:0
\[
\implies M^2(A_0-B_0)\ \le\ m^2_i\ \le\ M^2(A_0+B_0),
\]

%dsp-level:0
\[
\qquad m^2_i>0\ \ \forall i,\ \forall\delta,\ \forall(g,g')\ne(0,0):
\]

%dsp-level:0
\[
\qquad \textbf{every line of every channel is massive under the construction};
\]

%dsp-level:0
\[
\textbf{the kernel line is unreachable:}
\]

%dsp-level:0
\[
\quad c\parallel(0,0,g',g)\ [\mathrm{EW.22}:\ \ker M^2=\operatorname{span}\{Q\},\ Q=W_3+Y;\ \mathrm{FT.12}]\ \text{requires two vanishing cosines}:
\]

%dsp-level:0
\[
\cos\varphi
\]

%dsp-level:0
\[
= 0\wedge\cos(\varphi+\tfrac{2\pi}3)
\]

%dsp-level:0
\[
= 0\ \text{has no solution}\ [\text{machine:}\ c=0\implies s^2=1\implies\cos(\varphi+\tfrac{2\pi}3)=\mp\tfrac{r_3}2,\ r_3^2=3\ne0]:
\]

%dsp-level:0
\[
\qquad \mathrm{BM.36}(5)\text{'s exact massless line carries no family row};
\]

%dsp-level:0
\[
(2)\quad\textbf{the splitting law:}\qquad m^2_i-m^2_j=M^2B_0\bigl(\cos\varphi_j-\cos\varphi_i\bigr)\qquad[\text{machine: exact factorization}]:
\]

%dsp-level:0
\[
\textbf{the intra-family splitting is exactly the mixed weak--hypercharge strength}\ gg'\ \textbf{acting on the transport phase}\ [\mathrm{Def\text{-}BM.11}\ \mathrm{Rmk}];
\]

%dsp-level:0
\[
\text{degenerate}\iff gg'=0\ \vee\ \cos\varphi_i=\cos\varphi_j;
\]

%dsp-level:0
\[
\qquad \textbf{the democratic witness}\ \delta
\]

%dsp-level:0
\[
= 0:\ \ \{\cos\varphi_i\}
\]

%dsp-level:0
\[
= \{1,-\tfrac12,-\tfrac12\}:\ \ \text{two lines degenerate, the third split by}\ \tfrac32M^2B_0\ [\text{machine}];
\]

%dsp-level:0
\[
(3)\quad\textbf{the sum rule:}\qquad\sum_im^2_i=3M^2A_0\qquad[\textstyle\sum_i\cos\varphi_i=0:\ \mathrm{FT.111}(3)/\mathrm{FT.158}(4)\ \text{identity};\ \text{machine}]:
\]

%dsp-level:0
\[
\textbf{phase-independent and channel-independent}\ —\ \text{the seed is channel-blind}\ [\mathrm{EW.32}:\ \neg\text{colour},\ \neg\text{slot index}]:
\]

%dsp-level:0
\[
\qquad \text{all three channels share}\ (A_0,B_0);\ \text{only}\ \delta\in\{\tfrac29,\tfrac{2}{27},\tfrac18\}\ \text{differs};
\]

%dsp-level:0
\[
(4)\quad\textbf{the cross-channel row:}
\]

%dsp-level:0
\[
\qquad 16\,\delta_\nu=9\,\delta_\ell=27\,\delta_q=2,
\]

%dsp-level:0
\[
\qquad 3\delta_q=\delta_\ell
\]

%dsp-level:0
\[
\qquad [\mathrm{FT.143}(5);\ \mathrm{FT.119}(2);\ \text{machine: exact Fractions}]:
\]

%dsp-level:0
\[
\text{the three channels' phase data are exact rational multiples of one seed numerator}\ —\ \text{the inter-channel mass relations are symbolic rows in}\ (A_0,B_0,M^2),
\]

%dsp-level:0
\[
\ \neg\text{numbers}\ [\mathrm{P4};\ \mathrm{P5}];
\]

%dsp-level:0
\[
(5)\quad\textbf{the order inversion:}
\]

%dsp-level:0
\[
\qquad m^2_i\ \text{is affine strictly decreasing in}\ r_i\ [\mathrm{Def\text{-}BM.11}(5)\ \text{affine form};\ B_0>0\ \text{at}\ g,g'>0\ [\mathrm{EW.33}:\ g,g'>0];\ \text{machine: sign census}]:
\]

%dsp-level:0
\[
\textbf{the spectral-maximal line receives the mass-minimum};
\]

%dsp-level:0
\[
\qquad \text{the frozen chains re-cited:}\ \ r_1^2<r_2^2<r_0^2\ [\mathrm{FT.113}\ (\ell);\ \mathrm{FT.143}(1),(3)\ (\nu);\ \text{enclosure-proven, frozen}];
\]

%dsp-level:0
\[
\textbf{the read-species row:}\quad\mathrm{FT.113}\text{'s inertial vocabulary}\ m^2
\]

%dsp-level:0
\[
= \mu r^2\ \text{orders with the spectrum; the transport species orders against it}:
\]

%dsp-level:0
\[
\text{the order is read-relative; neither species is adopted as ``the mass''}\ [\mathrm{Def\text{-}BM.11}(7)];
\]

%dsp-level:0
\[
\qquad Q=\tfrac23\ \text{is the structure invariant of}\ \mathrm{Def\text{-}BM.6}(4),\ \neg\text{a mass claim}\ [\mathrm{FT.143}(4)];
\]

%dsp-level:0
\[
(6)\quad\textbf{the inter-family constraint stands:}
\]

%dsp-level:0
\[
\qquad \mathrm{BM.36}(5)\text{'s rows are untouched by the construction:}\quad0\le G(c)\le\mu_\star^2\tfrac{g^2+g'^2}4\|c\|^2,
\]

%dsp-level:0
\[
\qquad \text{the pole ladder}\subseteq[0,4\hat\mu'^2]\ [\mathrm{FT.76}(3)];
\]

%dsp-level:0
\[
\text{the multi-order hierarchies remain beyond the computed cells}\ [\mathrm{FT.77},\ \mathrm{FT.79}(6)\ \text{re-cited; quarantined rows stay quarantined}]:
\]

%dsp-level:0
\[
\qquad \text{the construction supplies the form,}\ \neg\text{a scale}\ [\mathrm{P4};\ \mathrm{P5}];
\]

%dsp-level:0
\[
(7)\quad\textbf{the area ladder (the holonomy species):}
\]

%dsp-level:0
\[
\qquad m^2_{\mathrm{pole}}=4\hat\mu'^2\sin^2\bigl(\tfrac{\Theta_R}4\bigr),
\]

%dsp-level:0
\[
\qquad \Theta_R=\bar KA=\textstyle\sum_i\alpha_i-(n-2)\pi
\]

%dsp-level:0
\[
\qquad [\mathrm{Def\text{-}H.37}(2);\ \mathrm{FT.54}(4);\ \mathrm{FT.47}(1);\ \text{machine: the exact SU(2) defect re-derived}];
\]

%dsp-level:0
\[
\textbf{the extent-dependent form:}\quad\text{a loop of spatial extent}\ r\ \text{enclosing area}\ A(r)\ \text{in a constant-curvature slice:}
\]

%dsp-level:0
\[
\qquad m^2(r)=4\hat\mu'^2\sin^2\Bigl(\frac{\bar K\,A(r)}4\Bigr)
\]

%dsp-level:0
\[
\qquad [\text{machine: formal identity}];
\]

%dsp-level:0
\[
\textbf{the ratio formula:}
\]

%dsp-level:0
\[
\qquad \frac{m^2_1}{m^2_2}=\frac{\sin^2(\Theta_{R,1}/4)}{\sin^2(\Theta_{R,2}/4)}
\]

%dsp-level:0
\[
\qquad [\text{machine: symbolic}];
\]

%dsp-level:0
\[
\qquad \text{monotone in}\ \Theta_R\in[0,2\pi],\ \ \text{range}\ [0,4\hat\mu'^2]\ [\mathrm{FT.76}(3)\ \text{re-cited}];
\]

%dsp-level:0
\[
\textbf{the winding multiplet:}\quad m^2(m)=4\hat\mu'^2\sin^2\bigl(\tfrac{m\,\Theta_R}4\bigr),\quad m\in\{k,k{-}2,\dots,-k\}
\]

%dsp-level:0
\[
\qquad [\mathrm{FT.60}(5)\ \text{re-cited; Def-H.38 ladder}];
\]

%dsp-level:0
\[
(8)\quad\textbf{the distinction row:}
\]

%dsp-level:0
\[
\qquad \Theta_R\ \text{is the angle defect}\ [\mathrm{FT.47}(1)];
\]

%dsp-level:0
\[
\qquad \text{the winding}\ k=\Theta/2\ \text{is the ledger read}\ [\mathrm{Def\text{-}H.32}(4);\ \mathrm{FT.22}];
\]

%dsp-level:0
\[
\text{the identification}\ \Theta_R=2\pi k\ \text{is ruled out on calibration weight}\ [\mathrm{FT.55}(1)\ \text{re-cited}]:
\]

%dsp-level:0
\[
\qquad \textbf{the two ladders are distinct reads, and the tier keeps them apart};
\]

%dsp-level:0
\[
(9)\quad\textbf{the matching row:}
\]

%dsp-level:0
\[
\qquad \text{the Gram species}\ ((1)\text{--}(5))\ \text{and the holonomy species}\ ((7))\ \text{are related by}\ \mathrm{Def\text{-}H.37}\text{'s leading-order matching}\ \hat\mu'^2\Theta_R^2/4
\]

%dsp-level:0
\[
= \hat\mu^2\ell_c^2G(c)\ [\mathrm{FT.54}(2);\ \mathrm{FT.33}]:
\]

%dsp-level:0
\[
\textbf{recorded, not imposed}\ [\mathrm{Def\text{-}H.37}\ \text{locus clause}];
\]

%dsp-level:0
\[
\qquad \text{the exact-matching derivation remains item-7 content}\ [\text{pointer}].
\]

\textbf{Rmk (two mechanisms, one form.)} Within a channel the families split by the Gram phase — the mixed strength times a cosine; along a winding or area ladder the levels split by the holonomy defect — a sine-squared of the enclosed curvature. The two are the leading-order and exact forms of one transfer [Def-H.37(1)–(2)], and the tier uses each where it is derived: the cosine triplet for the family structure, the defect ladder for the extent structure. Neither carries a scale: $\hat\mu,\hat\mu',\ell_c,(g,g')$ are the logged internal data, and no clause converts them into a number.

\textbf{Req.} P4, P5 (Parts 1–2); EW.22, EW.32, EW.33 (Part 2); Def-H.32(4), Def-H.37, Def-H.38, FT.12, FT.22, FT.33, FT.47(1), FT.54(2), FT.54(4), FT.55(1), FT.60(5), FT.77, FT.79(6), FT.111(3), FT.112, FT.113, FT.119(1), FT.119(2), FT.142(5), FT.143(1), FT.143(3), FT.143(4), FT.143(5), FT.158(4) (Part 3); Census FT.76(3) (row-wise); Def-BM.6(4), Def-BM.11, BM.36(5), BM.36(6) (this part); Cal-F; sum-of-squares positivity, exact quotient-ring identities, symbolic ratio algebra, exact Pauli defect identities [math; machine: \texttt{checks/\allowbreak{}verify\_\allowbreak{}H\_\allowbreak{}neutron\_\allowbreak{}decay\_\allowbreak{}vertex.py} §A].

\textbf{Pf.}

%dsp-level:0
\[
(1):\ \mathrm{Def\text{-}BM.11}(5)\text{'s evaluation re-cited};
\]

%dsp-level:0
\[
\qquad \tfrac38(g\mp\tfrac{g'}3)^2=\tfrac{3g^2}8\mp\tfrac{gg'}4+\tfrac{g'^2}{24}:\quad+\tfrac{g'^2}{48}\ \text{gives}\ A_0\mp B_0\ [\text{machine: expansion}];
\]

%dsp-level:0
\[
\qquad |\cos|\le1\implies\text{the sandwich};
\]

%dsp-level:0
\[
\qquad \text{positivity: a sum of squares with positive rational coefficients vanishes only at}\ (g,g')
\]

%dsp-level:0
\[
= (0,0)\ [\text{machine}]
\]

%dsp-level:0
\[
\text{the kernel row:}\ c=0\wedge s^2=1\implies\cos(\varphi+\tfrac{2\pi}3)=-\tfrac c2-\tfrac{r_3s}2=\mp\tfrac{r_3}2\ne0\ [r_3^2=3;\ \text{machine}]
\]

%dsp-level:0
\[
(2):\ \text{difference of}\ (1)\text{'s forms}\ [\text{machine}];
\]

%dsp-level:0
\[
\qquad \delta=0:\ c=1,s=0\implies\{1,-\tfrac12,-\tfrac12\}\ [\text{machine:}\ -\tfrac{c\pm r_3s}2\ \text{at}\ (c,s)=(1,0)]
\]

%dsp-level:0
\[
(3):\ \textstyle\sum_i\cos\varphi_i=0\ [\mathrm{FT.111}(3)\ \text{identity, generic phase};\ \text{machine}];
\]

%dsp-level:0
\[
\qquad (A_0,B_0)\ \text{carry no channel argument}\ [\mathrm{EW.32}\ \text{re-cited}]
\]

%dsp-level:0
\[
(4):\ 16\cdot\tfrac18=9\cdot\tfrac29=27\cdot\tfrac{2}{27}=2;\quad3\cdot\tfrac{2}{27}=\tfrac29\ [\text{machine: exact Fractions}]
\]

%dsp-level:0
\[
(5):\ \text{the affine coefficient}\ -M^2B_0/r_2<0\ \text{at}\ g,g'>0\ [\text{machine: sign census of the monomial}];
\]

%dsp-level:0
\[
\qquad \mathrm{FT.113}/\mathrm{FT.143}\text{'s chains re-cited (frozen, enclosure-proven)};
\]

%dsp-level:0
\[
\qquad Q=\tfrac69=\tfrac23\ [\mathrm{FT.143}(4)\ \text{re-cited}]
\]

%dsp-level:0
\[
(6):\ \mathrm{BM.36}(5),\ \mathrm{FT.76}(3),\ \mathrm{FT.77},\ \mathrm{FT.79}(6)\ \text{re-cited}\ (\neg\text{recomputed})
\]

%dsp-level:0
\[
(7):\ (U-1)^\dagger(U-1)
\]

%dsp-level:0
\[
= (2-2\cos\phi)\mathbb I
\]

%dsp-level:0
\[
= 4\sin^2(\tfrac\phi2)\mathbb I\ \text{at}\ U
\]

%dsp-level:0
\[
= \cos\phi\,
\]

%dsp-level:0
\[
\mathbb I+i\sin\phi\,
\]

%dsp-level:0
\[
\hat n\cdot\vec\sigma\ [\mathrm{FT.54}(1)\ \text{re-derived};\ \text{machine: exact Pauli algebra in}\ \mathbb Q[C,S]/(C^2{+}S^2{-}1)];
\]

%dsp-level:0
\[
\qquad \phi=\tfrac{\Theta_R}2\ [\mathrm{FT.54}(3);\ \mathrm{Def\text{-}H.37}(2)]
\]

%dsp-level:0
\[
(8):\ \mathrm{FT.55}(1)\text{'s calibration-weight obstruction re-cited};
\]

%dsp-level:0
\[
\qquad (9):\ \mathrm{Def\text{-}H.37}\text{'s matching row re-cited}.\ \blacksquare
\]

\textbf{Phys.} The family mass form comes out exact and tight. Positivity is a sum of two squares with rational coefficients, so no line of any channel can fall to zero through this map — the massless direction of the Gram form is the electromagnetic generator, and a cosine triplet can never align with it, because two cosines a third of a turn apart cannot vanish together. The splitting between lines is precisely the product of the two electroweak strengths times the cosine separation: the force structure acting on the transport phase, with no Yukawa object anywhere in it. The three channels share the whole form and differ only in their derived angles, so their orbit sums are one expression and their splittings are rational multiples of one another. The order runs against the spectrum: the brightest spectral line is the lightest mass under this read, and the tier records that reversal instead of hiding it, because the inertial vocabulary orders the other way and no theorem yet converts the two. Alongside the Gram mechanism stands the holonomy one: mass as the sine-squared of a quarter angle defect, growing with the area a loop encloses — the extent-dependent form the corpus can derive exactly, with the winding identification ruled out and the two ladders kept apart.

\end{theorem}

\begin{theorem}[{\texorpdfstring{BM.40 — The rung-1 isospin row: the Casimir identity, the branch selectivity, and the neutrality of the $\sigma$-vertex}{BM.40 - The rung-1 isospin row: the Casimir identity, the branch selectivity, and the neutrality of the -vertex}}]\label{tag:BM-40}
\textbf{Thm.}

%dsp-level:0
\[
(1)\quad\textbf{the Casimir identity:}
\]

%dsp-level:0
\[
\qquad W_i=\tfrac12\bigl(\sigma_i\otimes\sigma_0+\sigma_0\otimes\sigma_i\bigr)\ [\mathrm{CT1.2}];
\]

%dsp-level:0
\[
\qquad [W_i,W_j]=i\epsilon_{ijk}W_k\ [\mathrm{CT1.4};\ \text{machine}];
\]

%dsp-level:0
\[
\boxed{\ \textstyle\sum_iW_i^2=\tfrac32E_0+\tfrac12\Sigma=E_0+\mathbb S=2\,\Pi_{\mathcal R}\ }
\]

%dsp-level:0
\[
\qquad [\mathrm{Cor}\ \mathrm{CT0.7a};\ \mathrm{FT.141}(4);\ \mathrm{FT.157}(2);\ \text{machine: entrywise over}\ \mathbb Q(i),\ \text{all}\ 16\ \text{components}];
\]

%dsp-level:0
\[
(2)\quad\textbf{the branch selectivity:}
\]

%dsp-level:0
\[
\qquad \textstyle\sum_iW_i^2\big|_{\vee^2}=2=1(1{+}1),
\]

%dsp-level:0
\[
\qquad \textstyle\sum_iW_i^2\big|_{\wedge^2}=0=0(0{+}1)
\]

%dsp-level:0
\[
\qquad [\mathrm{Cor}\ \mathrm{CT0.7a};\ \text{machine: exact on the branch bases}];
\]

%dsp-level:0
\[
W_i\Pi_{\mathcal V}=0=\Pi_{\mathcal V}W_i\ \forall i,\qquad\bigcap_i\ker W_i=\operatorname{im}\Pi_{\mathcal V}\qquad[\mathrm{FT.101}(1)\text{--}(2);\ \mathrm{FT.157}(2);\ \text{machine}];
\]

%dsp-level:0
\[
\Pi_{\mathcal R}=\tfrac14(3E_0+\Sigma),\quad\Pi_{\mathcal V}=\tfrac14(E_0-\Sigma):
\]

%dsp-level:0
\[
\qquad \text{idempotent, orthogonal, complete},\quad\mathfrak T\Pi_{\mathcal R}=3,\quad\mathfrak T\Pi_{\mathcal V}=1
\]

%dsp-level:0
\[
\qquad [\mathrm{CT0.7}(1);\ \text{machine}];
\]

%dsp-level:0
\[
(3)\quad\textbf{the neutrality of the}\ \sigma\textbf{-vertex:}
\]

%dsp-level:0
\[
\qquad \text{the transition's charge move is an end-sign flip:}
\]

%dsp-level:0
\[
\quad q_{\mathrm{core}}=\tfrac12(\sigma_s-\sigma_t)\ \text{is registration data on the reality line}\ [\mathrm{D3.4};\ \mathrm{D1.1};\ \mathrm{CT0.9};\ \mathrm{BM.23}(2)];
\]

%dsp-level:0
\[
\text{the Casimir}\ 2\Pi_{\mathcal R}\ \text{is a function of the branch datum alone}\ [(1)];
\]

%dsp-level:0
\[
\qquad \text{the gradings are pairwise distinct}\ [\mathrm{FT.136};\ \mathrm{FT.141}(6);\ \mathrm{FT.144}(5)];
\]

%dsp-level:0
\[
\qquad \text{the register separation holds}\ [\mathrm{FT.136}(3)];
\]

%dsp-level:0
\[
\boxed{\ \textbf{the rung-1 Casimir deficit at the}\ \sigma\textbf{-vertex is identically zero}\ }
\]

%dsp-level:0
\[
\qquad [\text{machine: the token census}\ —\ \text{zero branch tokens in}\ \mathrm{D3.4}\text{'s defining row, zero registration tokens in}\ \mathrm{CT0.7}(1)/\mathrm{FT.101}\text{'s defining rows}];
\]

%dsp-level:0
\[
\text{the conservation law holds at the vertex}\ —\ \text{it neither forces nor forbids an emission}:
\]

%dsp-level:0
\[
(4)\quad\textbf{the carrier-side rows:}
\]

%dsp-level:0
\[
\qquad \text{every emitted chain is Janus-closed:}\quad\#J^+=\#J^-\ [\mathrm{D1.4};\ \mathrm{FT.144}(1)\text{--}(2);\ \text{the prefix law}]:
\]

%dsp-level:0
\[
\qquad \text{each carrier is internally Casimir-balanced};
\]

%dsp-level:0
\[
\text{the neutral carrier class is rung-one content:}\quad\mathcal A(1)\ \text{aspects},
\]

%dsp-level:0
\[
\ \mathcal K_1\ \text{arena},
\]

%dsp-level:0
\[
\ \dim
\]

%dsp-level:0
\[
= 16\ [\mathrm{FT.142}(1)\text{--}(2)];
\]

%dsp-level:0
\[
\qquad \text{EM-unactivated}\ (f_v=0),\ \text{colour-blind}\ [\mathrm{FT.142}(1)];
\]

%dsp-level:0
\[
\text{at rung one the chiral sectors are force-decoupled:}
\]

%dsp-level:0
\[
\quad F_{LR}=0\ [\mathrm{FT.144}(6);\ \mathrm{FT.128}(2);\ \text{the no-flip theorem}]:
\]

%dsp-level:0
\[
\qquad \textbf{the neutral carrier's orientation sector is frozen along its flight}\ [\mathrm{FT.144}(5):\ \text{single-sector occupancy}];
\]

%dsp-level:0
\[
(5)\quad\textbf{the status:}
\]

%dsp-level:0
\[
\qquad \text{the forcing of the neutral member comes}\ \neg\text{from the Casimir row}\ —\ (3)\ \text{adjudicates the deficit route neutral}\ —\ \text{but from}\ \mathrm{BM.41}\text{'s kinematic theorem};
\]

%dsp-level:0
\[
\mathrm{BM.23}(6)\text{'s clause stands as written and receives an on-record Rmk}\ [\text{the FT.89 Rmk pattern}];
\]

%dsp-level:0
\[
\qquad \neg\exists\ \text{a lepton-number read in the inventory}\ [\mathrm{BM.23}(6);\ \mathrm{FT.81}(4)\ \text{pattern}]\ —\ \text{the row is untouched}.
\]

\textbf{Rmk (why the proposed deficit does not arise).} The isospin bookkeeping the vertex could upset lives in the branch datum — aligned pairs carry Casimir two, antialigned pairs zero. The charge move that defines the transition lives in the registration sign — the reality-line datum the end reads consume. The corpus proved these gradings pairwise distinct and their registers separate, so no $\sigma$-flip can move a $\Pi_{\mathcal R}$: the deficit is zero as an operator identity, not as a small number. That is why the neutrino's necessity, where it is necessity, must come from somewhere else — and BM.41 shows it does: from the kinematics of equal rest masses and the uniqueness of the neutral class.

\textbf{Req.} A3 (Part 1); CT0.7(1), CT0.9, CT1.2, CT1.4, Cor CT0.7a, D1.1, D1.4, D3.4 (Parts 1–2); P5 (Parts 1–2); FT.12, FT.101(1), FT.101(2), FT.128(2), FT.136(3), FT.141(4), FT.141(6), FT.142(1), FT.142(2), FT.144(1), FT.144(2), FT.144(5), FT.144(6), FT.157(2) (Part 3); BM.7, BM.23(2), BM.23(6) (this part); the forcing theorem of this tier's vertex tag is a marked forward pointer at the citing site and is imported in no \textbf{Req} line; exact $\mathbb Q(i)$ $4\times4$ algebra, projector algebra, token censuses [math; machine: \texttt{checks/\allowbreak{}verify\_\allowbreak{}H\_\allowbreak{}neutron\_\allowbreak{}decay\_\allowbreak{}vertex.py} §B].

\textbf{Pf.}

%dsp-level:0
\[
(1):\ \textstyle\sum_iW_i^2
\]

%dsp-level:0
\[
= \tfrac14\sum_i(\sigma_i\otimes\sigma_0+\sigma_0\otimes\sigma_i)^2
\]

%dsp-level:0
\[
= \tfrac14(3E_0+3E_0+2\Sigma)
\]

%dsp-level:0
\[
= \tfrac32E_0+\tfrac12\Sigma\ [\mathrm{Cor}\ \mathrm{CT0.7a}\ \text{re-derived};\ \text{machine}];
\]

%dsp-level:0
\[
\qquad \Sigma=2\mathbb S-E_0\ [\mathrm{CT0.7}(1)]\implies=E_0+\mathbb S=2\Pi_{\mathcal R}\ [\text{machine}];
\]

%dsp-level:0
\[
\qquad [W_i,W_j]=i\epsilon_{ijk}W_k\ [\text{machine: all}\ 9\ \text{pairs}]
\]

%dsp-level:0
\[
(2):\ \vee^2\ \text{basis}\ \{(1,0,0,0),(0,1,1,0),(0,0,0,1)\}:\ 2\cdot v\ [\text{machine}];
\]

%dsp-level:0
\[
\qquad \wedge^2\ \text{basis}\ \{(0,1,-1,0)\}:\ 0\ [\text{machine}];
\]

%dsp-level:0
\[
\qquad W_i\Pi_{\mathcal V}=0=\Pi_{\mathcal V}W_i\ [\text{machine: all}\ 16\ \text{components each}];
\]

%dsp-level:0
\[
\qquad \text{the projector algebra and the traces}\ [\text{machine}]
\]

%dsp-level:0
\[
(3):\ \text{the defining strings extracted}: \mathrm{D3.4}\text{'s row vs}\ \mathrm{CT0.7}(1)/\mathrm{FT.101}\text{'s rows};
\]

%dsp-level:0
\[
\qquad \text{the branch alphabet}\ \{\Pi_{\mathcal R},\Pi_{\mathcal V},\alpha,\text{branch}\}:\ 0\ \text{occurrences in}\ \mathrm{D3.4};
\]

%dsp-level:0
\[
\qquad \text{the registration alphabet}\ \{\sigma_s,\sigma_t,q_{\mathrm{core}},\text{reality line}\}:\ 0\ \text{occurrences in the branch rows}\ [\text{machine}]
\]

%dsp-level:0
\[
(4):\ \mathrm{D1.4}/\mathrm{FT.144}(1)\text{--}(2)\ \text{re-cited (frozen)};
\]

%dsp-level:0
\[
\qquad F_{LR}=0\ [\mathrm{FT.128}(2)\ \text{re-cited}];
\]

%dsp-level:0
\[
\qquad (5):\ \text{the status record}.\ \blacksquare
\]

\textbf{Phys.} The weak isospin of the pair arena is an exact operator identity — the Casimir is twice the projector onto the real branch — and it splits the pair register into a spin-one triplet the weak force sees and a spin-zero singlet it annihilates. The question the vertex poses is whether a charge flip disturbs that bookkeeping, and the answer is a clean no: the flip moves a registration sign, the Casimir reads a branch label, and the corpus long ago proved those gradings live in separate registers. The deficit is zero identically, so the isospin law neither demands nor forbids the neutral emission. What the rung-one rows do supply is the carrier's pedigree: closed under Janus balance, electromagnetically inert, colour-blind, force-decoupled across its chiral sectors at rung one, single-arrow in occupancy — a neutral member that travels without turning. The necessity itself waits for the kinematics.

\end{theorem}

\begin{theorem}[{\texorpdfstring{BM.41 — The decay vertex: the forced members, the two-body impossibility, and the Q-value sign theorem}{BM.41 - The decay vertex: the forced members, the two-body impossibility, and the Q-value sign theorem}}]\label{tag:BM-41}
\textbf{Thm.}

%dsp-level:0
\[
(1)\quad\textbf{the vertex grammar:}\qquad\begin{array}{l|l}\text{the parent} & \text{the neutral-class}\ \varepsilon\text{-sealed 3-juncture},\ E_{\mathrm{int}}=6,\ Q=0\ [\mathrm{BM.17};\ \mathrm{Def\text{-}BM.3}(2);\ \text{Census FT.35}(2)]\\\text{the daughter} & \text{the unit-class}\ \varepsilon\text{-sealed 3-juncture},\ E_{\mathrm{int}}=6,\ Q=+1\ [\mathrm{BM.26};\ \mathrm{Def\text{-}BM.3}(3)]\\\text{the charged member} & \text{the}\ \ell\text{-channel line},\ q=-1:\ \text{forced and class-unique}\ [\mathrm{BM.23}(2)\text{--}(5):\ \Delta q=\pm1,\ \text{the unique compensating value},\ \text{the selection theorem}]\\\text{the neutral member} & \text{the}\ \nu\text{-channel line},\ q=0:\ \text{forced by}\ (2)\text{--}(3);\ \text{class-unique}\ [\mathrm{FT.142}(1);\ \mathrm{BM.23}(5)\text{'s eliminations}]\end{array}
\]

%dsp-level:0
\[
\textbf{the charge balance:}\quad0=(+1)+(-1)+0\ [\mathrm{EM.11};\ \mathrm{BM.23}(4);\ \text{machine: integers}];
\]

%dsp-level:0
\[
\qquad \textbf{the sealing phase:}\ \Theta_{\mathrm{int}}\in4\mathbb Z\ \text{on both baryon-type junctions}\ [\mathrm{BM.2}(1);\ \text{Census FT.49}(3)];
\]

%dsp-level:0
\[
\textbf{the arrow:}\ \text{forward-only}\ [\mathrm{BM.23}(9);\ \mathrm{FT.144}(2);\ \mathrm{A6}:\ \neg\exists\tau^{-1}];
\]

%dsp-level:0
\[
\qquad \textbf{the transition weight:}\ \varepsilon_q=0\implies\Gamma=0\ \text{exactly}\ [\mathrm{BM.23}(8)]:\ \text{no flicker, no decay};
\]

%dsp-level:0
\[
(2)\quad\textbf{the two-body impossibility (exact):}
\]

%dsp-level:0
\[
\qquad m_n=m_p=\frac{6\ln2}{\mathfrak h}
\]

%dsp-level:0
\[
\qquad [\mathrm{BM.17};\ \mathrm{BM.26};\ \mathrm{BM.28}(1)\ \text{charge-blindness};\ \mathrm{BM.34}(7)/\mathrm{BM.35}(4)\ \text{the rest slice}];
\]

%dsp-level:0
\[
\text{the rest frame, two members:}\quad E_p+E_\ell=m_n,
\]

%dsp-level:0
\[
\qquad E_i=\sqrt{m_i^2+p^2}\ [\mathrm{BM.34}(4)\ \text{vacuum gauge};\ \mathrm{BM.11}]:
\]

%dsp-level:0
\[
\qquad E_\ell=\frac{m_n^2+m_\ell^2-m_p^2}{2m_n},
\]

%dsp-level:0
\[
\quad E_p=\frac{m_n^2+m_p^2-m_\ell^2}{2m_n}\ [\text{machine: the exact solution of the conservation system}];
\]

%dsp-level:0
\[
m_n=m_p:\qquad\boxed{\ E_p-m_p=-\frac{m_\ell^2}{2m_n}<0\ \ \forall m_\ell>0\ }\qquad[\text{machine: cleared polynomial identity}\ +\ \text{sign census}]:
\]

%dsp-level:0
\[
\textbf{the two-body vertex puts the daughter below its own mass shell for every positive member mass}\ —\ \neg\exists\ \text{a real on-shell solution};
\]

%dsp-level:0
\[
\textbf{the Källén companion:}\quad4m_n^2p^2=\lambda\bigl(m_n^2,m_p^2,m_\ell^2\bigr),
\]

%dsp-level:0
\[
\qquad \lambda=\bigl(m_n^2-(m_p{+}m_\ell)^2\bigr)\bigl(m_n^2-(m_p{-}m_\ell)^2\bigr)\ [\text{machine: polynomial identity}];
\]

%dsp-level:0
\[
m_n=m_p:\quad\lambda=m_\ell^2\bigl(m_\ell^2-4m_n^2\bigr):
\]

%dsp-level:0
\[
\qquad \text{spacelike for}\ 0<m_\ell<2m_n;
\]

%dsp-level:0
\[
\qquad m_\ell\ge2m_n\implies E_p\le-m_n<0\ [\text{machine: the two branches cover every}\ m_\ell>0\ (2>\sqrt2\ \text{by squaring})];
\]

%dsp-level:0
\[
(3)\quad\textbf{the forced neutral member:}
\]

%dsp-level:0
\[
\qquad \text{by}\ (2)\ \text{every kinematically admissible vertex carries}\ \ge3\ \text{members};
\]

%dsp-level:0
\[
\qquad \text{the charge balance fixes the third member's total charge to}\ 0\ [(1);\ \mathrm{EM.11}];
\]

%dsp-level:0
\[
\textbf{the class is unique among the constructed carrier classes:}\quad\text{the}\ \nu\text{-channel is the unique EM-unactivated colour-blind bound-chain class}\ [\mathrm{BM.23}(5)\text{'s elimination rows};\ \mathrm{FT.142}(1)];
\]

%dsp-level:0
\[
\text{the transit (EM-radiation) alternative is}\ \neg\text{constructed: the gauge coupling of the}\ Q\ \text{read is deficit}\ 1'\ [\mathrm{FT.125}(4);\ \mathrm{Def\text{-}BM.8}(8)]\ —\ \neg\text{an inventory candidate};
\]

%dsp-level:0
\[
\boxed{\ \textbf{the neutral}\ \nu\text{-channel member is forced in every kinematically admissible vertex}\ }
\]

%dsp-level:0
\[
\qquad [\text{the forcing is kinematic}\ +\ \text{uniqueness},\ \neg\text{a number conservation}:\ \neg\exists\ \text{a lepton-number read}\ [\mathrm{BM.23}(6)]];
\]

%dsp-level:0
\[
\textbf{the line index within the}\ \nu\textbf{-triplet is}\ \neg\textbf{derived}\ [\neg\text{a mixing matrix}\ [\mathrm{FT.145}(3)]];
\]

%dsp-level:0
\[
\qquad \mathrm{BM.23}(6)\ \text{is upgraded on record: admissible}\ \longrightarrow\ \text{conditionally forced}\ [\text{its Rmk}];
\]

%dsp-level:0
\[
(4)\quad\textbf{the Q-value sign theorem:}\qquad Q:=m_n-m_p-m_\ell-m_{\nu};
\]

%dsp-level:0
\[
(\mathrm a)\ \textbf{the crystallized informational read:}
\]

%dsp-level:0
\[
\quad Q=-(m_\ell+m_{\nu})\ [\text{machine: monomial identity at}\ m_n=m_p];
\]

%dsp-level:0
\[
\qquad m_\ell>0\ \text{on every read species on record}:
\]

%dsp-level:0
\[
\text{transport:}\ m_\ell^2\ge M^2(A_0-B_0)>0\ [\mathrm{BM.39}(1)];
\]

%dsp-level:0
\[
\qquad \text{informational:}\ m_\ell
\]

%dsp-level:0
\[
= \frac{E_{\mathrm{int}}\ln2}{\mathfrak h}\ge\frac{2\ln2}{\mathfrak h}>0\ [\mathrm{FT.44}(4);\ \mathrm{FT.28}\ \text{the minimal floor};\ \tfrac12<\ln2<1\ \text{enclosure};\ \text{machine}];
\]

%dsp-level:0
\[
\qquad m_\nu\ge0\ [\mathrm{EW.33}\ \mathrm{PSD}];
\]

%dsp-level:0
\[
\implies\boxed{\ Q<0\ \ \text{strictly}\ };
\]

%dsp-level:0
\[
(\mathrm b)\ \textbf{every normalized ensemble read:}\quad\text{charge-degeneracy is structural}\ [\mathrm{FT.107}(3);\ \mathrm{BM.30}(3);\ \mathrm{BM.37}(2)\ \text{the disjointness}]
\]

%dsp-level:0
\[
\implies Q
\]

%dsp-level:0
\[
= -(m_\ell+m_\nu)<0;
\]

%dsp-level:0
\[
(\mathrm c)\ \textbf{the curvature-flux read:}
\]

%dsp-level:0
\[
\quad RV=\varkappa^2m\ \text{with identical}\ (E_{\mathrm{int}},V)\ \text{on both classes}\ [\mathrm{BM.20}(1);\ \mathrm{BM.28}(1)]\implies Q<0;
\]

%dsp-level:0
\[
(\mathrm d)\ \textbf{the route-(ii) variant:}\quad T(0)=\tfrac43T(\pm1)\ [\mathrm{BM.28}(4)(\mathrm{ii})]:
\]

%dsp-level:0
\[
\qquad Q_{\mathrm{ii}}=\tfrac13T(\pm1)-m_\ell-m_\nu:
\]

%dsp-level:0
\[
\qquad \textbf{the sign is undetermined without scale data}\ [\text{machine: two rational witnesses of opposite sign}];
\]

%dsp-level:0
\[
\text{the functional is a measure price}\ [\mathrm{BM.37}(3)]\ \text{and its adoption requires a §I addendum}\ [\mathrm{BM.30}(5)]:
\]

%dsp-level:0
\[
\qquad \text{priced},\ \neg\text{adopted};
\]

%dsp-level:0
\[
(\mathrm e)\ \textbf{the verdict:}
\]

%dsp-level:0
\[
\qquad \boxed{\ \textbf{the}\ \beta\textbf{-decay vertex is kinematically closed on every adopted read: the corpus does not make the neutron heavier than the proton}\ }
\]

%dsp-level:0
\[
[\mathrm{P5}:\ \neg\text{an external mass-difference datum};\qquad\text{the blocks are theorems:}\ \mathrm{BM.28}(1),\ \mathrm{FT.107}(3),\ \mathrm{BM.37}(3),\ \mathrm{BM.30}(5)];
\]

%dsp-level:0
\[
(5)\quad\textbf{the three-body kinematics (the conditional section, under the hypothesis}\ Q\ge0\ \textbf{of some adopted specification):}
\]

%dsp-level:0
\[
\qquad \text{the rest frame:}\quad0=\vec p_p+\vec p_\ell+\vec p_\nu\ [\mathrm{BM.11};\ \mathrm{BM.34}(3)];
\]

%dsp-level:0
\[
Q=\mathrm{KE}_p+\mathrm{KE}_\ell+\mathrm{KE}_\nu,
\]

%dsp-level:0
\[
\qquad \mathrm{KE}_i=E_i-m_i
\]

%dsp-level:0
\[
\qquad [\mathrm{BM.34}(5);\ \text{machine: Laurent identity}];
\]

%dsp-level:0
\[
\qquad \text{NR:}\ \ \mathrm{KE}_i=\frac{p_i^2}{2m_i}+O(p^4);
\]

%dsp-level:0
\[
\textbf{the endpoint:}
\]

%dsp-level:0
\[
\qquad \boxed{\ E_\ell^{\max}=\frac{m_n^2+m_\ell^2-(m_p+m_\nu)^2}{2m_n}\ }
\]

%dsp-level:0
\[
\qquad [\text{attained at the relative-rest recoil};\ \text{machine: the exact polynomial identities}];
\]

%dsp-level:0
\[
E_\ell^{\max}+E_{\mathrm{rec}}=m_n,\quad E_{\mathrm{rec}}=\frac{m_n^2+(m_p+m_\nu)^2-m_\ell^2}{2m_n};
\]

%dsp-level:0
\[
\qquad \bigl(E_\ell^{\max}\bigr)^2-m_\ell^2=E_{\mathrm{rec}}^2-(m_p+m_\nu)^2\ [\text{machine: cleared by}\ 4m_n^2];
\]

%dsp-level:0
\[
\textbf{the recoil bound:}\quad s_{p\nu}\ge(m_p+m_\nu)^2:
\]

%dsp-level:0
\[
\qquad \text{the collinear identity}\ \bigl(E_pE_\nu-PN\bigr)^2
\]

%dsp-level:0
\[
= (m_pm_\nu)^2+\bigl(E_pN-E_\nu P\bigr)^2\ [\text{machine: on-shell quotient ring}];
\]

%dsp-level:0
\[
\text{the general case by Cauchy--Schwarz:}
\]

%dsp-level:0
\[
\quad P^2N^2-(\vec p_p\cdot\vec p_\nu)^2=\sum_{i<j}\bigl(p_{p,i}p_{\nu,j}-p_{p,j}p_{\nu,i}\bigr)^2\ge0\ [\text{machine: the Lagrange identity}];
\]

%dsp-level:0
\[
\qquad E_\ell^{\max}\ \text{is strictly decreasing in}\ s_{p\nu}\ [\text{the coefficient}\ -\tfrac1{2m_n}<0];
\]

%dsp-level:0
\[
\text{the continuum of}\ E_\ell\ \text{values on}\ [E_\ell^{\min},E_\ell^{\max}]\ \text{is chart-side content}\ [\mathrm{CT2.6}\ \text{hypotheses at the point of use}];
\]

%dsp-level:0
\[
(6)\quad\textbf{the general channel closure:}
\]

%dsp-level:0
\[
\qquad \text{any member set containing the daughter and at least one positive-mass chain line satisfies}\ \textstyle\sum m_{\mathrm{out}}>m_n\ [(2)\text{'s equality of rest reads};\ \text{machine: sign arithmetic}]:
\]

%dsp-level:0
\[
\textbf{every material decay channel of the neutron row is kinematically closed on the adopted reads}\ —\ \text{the closure is a theorem about the reads,}\ \neg\text{about the vertex grammar}\ [(1)\ \text{stands complete}];
\]

%dsp-level:0
\[
(7)\quad\textbf{status and scope:}
\]

%dsp-level:0
\[
\qquad \text{the vertex grammar}\ (1)\ \text{and the forcing}\ (3)\ \text{are derived};
\]

%dsp-level:0
\[
\qquad \text{the kinematic viability is read-relative}\ (4);
\]

%dsp-level:0
\[
\qquad \text{the conditional section}\ (5)\ \text{carries its hypothesis on its face};
\]

%dsp-level:0
\[
\neg\text{a measured lifetime, mass difference or endpoint enters}\ [\mathrm{P5};\ \text{machine: zero external tokens}];
\]

%dsp-level:0
\[
\qquad \text{the row-index attachment}\ (3)\ \text{and the read choice}\ (4)(\mathrm d)\ \text{are the open faces}.
\]

\textbf{Rmk (the shape of the negative.)} The corpus assembles this vertex completely and then closes it — and the closing is as exact as the assembly. Two independent theorems do the work: the charge-blindness of every crystallized count read, which makes the two baryon rows weigh the same to the last monomial; and the structural charge-degeneracy of every normalized ensemble average, which makes the flickering rows weigh the same too. The one functional that does split them was classified as a measure price, and adopting it as a mass is an act of specification the corpus prices but does not perform. So the reaction stands: grammar forced, members forced, conservation exact, kinematics closed — waiting on a read specification, not on a missing derivation.

\textbf{Req.} A6 (Part 1); D1.4, D2.9, D2.11 (Parts 1–2); EM.9, EM.11 (Part 2); P5 (Parts 1–2); CT2.6 (Part 2); FT.22, FT.107(3), FT.125(4), FT.142(1), FT.144(2), FT.145(3) (Part 3); Census FT.28, FT.35(2), FT.44(4), FT.49(3) (row-wise); Def-BM.3, Def-BM.6, Def-BM.8, Def-BM.10, BM.2(1), BM.11, BM.17, BM.23(2), BM.23(4), BM.23(5), BM.23(6), BM.23(8), BM.23(9), BM.26, BM.28(1), BM.28(4), BM.30(3), BM.30(5), BM.34(3), BM.34(4), BM.34(5), BM.34(7), BM.35(4), BM.36(5), BM.37(2), BM.37(3), BM.39(1), BM.40(3) (this part); Källén polynomial identities, exact sign arithmetic over positive witnesses, the collinear and Lagrange Cauchy–Schwarz identities, exact rational $\ln2$ enclosures [math; machine: \texttt{checks/\allowbreak{}verify\_\allowbreak{}H\_\allowbreak{}neutron\_\allowbreak{}decay\_\allowbreak{}vertex.py} §C].

\textbf{Pf.}

%dsp-level:0
\[
(1):\ \mathrm{BM.23}(2)\text{--}(5)\ \text{re-cited (frozen)};
\]

%dsp-level:0
\[
\qquad 0=1+(-1)+0\ [\text{machine}];
\]

%dsp-level:0
\[
\qquad \Theta_{\mathrm{int}}\in4\mathbb Z\ [\text{Census FT.49}(3)\ \text{re-cited}];
\]

%dsp-level:0
\[
\qquad \Gamma(\varepsilon_q{=}0)=0\ [\mathrm{BM.23}(8)\ \text{re-cited}]
\]

%dsp-level:0
\[
(2):\ \text{the conservation system}\ \{E_p+E_\ell=m_n;\ E_\ell^2-m_\ell^2=E_p^2-m_p^2\}\ \text{solved linearly in}\ (E_p,E_\ell)\ [\text{machine: substitution verifies both equations as polynomial identities over}\ \mathbb Q[m_n,m_p,m_\ell]];
\]

%dsp-level:0
\[
m_n=m_p:\quad2m_n(E_p-m_p)=(m_n^2+m_p^2-m_\ell^2)-2m_nm_p=-m_\ell^2\ [\text{machine: cleared identity}];
\]

%dsp-level:0
\[
\qquad \lambda\ \text{expansion and the}\ m_n{=}m_p\ \text{factorization}\ [\text{machine}];
\]

%dsp-level:0
\[
\qquad (2m_n)^2-m_\ell^2\ \text{branch cover:}\ \sqrt2<2\ [\text{squaring comparison}]
\]

%dsp-level:0
\[
(3):\ (2)\ \text{conjoined with}\ (1)\text{'s balance};
\]

%dsp-level:0
\[
\qquad \text{the elimination rows and deficit}\ 1'\ \text{re-cited};
\]

%dsp-level:0
\[
\qquad \text{the upgrade recorded in BM.23's Rmk}
\]

%dsp-level:0
\[
(4)(\mathrm a):\ m_n-m_p=0\ [\text{machine: monomial identity at}\ E_{\mathrm{int}}=6\ \text{both rows}];
\]

%dsp-level:0
\[
\qquad m_\ell>0\ \text{on both species}\ [\mathrm{BM.39}(1)\text{'s SOS row};\ \mathrm{FT.28}/\mathrm{FT.44}(4)\ \text{with}\ \tfrac12<\ln2\ \text{enclosure};\ \text{machine}];
\]

%dsp-level:0
\[
\qquad (\mathrm b),(\mathrm c):\ \text{degeneracy rows re-cited};
\]

%dsp-level:0
\[
\qquad (\mathrm d):\ \text{witnesses}\ T{=}3,
\]

%dsp-level:0
\[
m_\ell{=}m_\nu{=}\tfrac18\ (Q_{\mathrm{ii}}{=}\tfrac34{>}0)\ \text{and}\ T{=}3,
\]

%dsp-level:0
\[
m_\ell{=}m_\nu{=}1\ (Q_{\mathrm{ii}}{=}{-}1{<}0)\ [\text{machine: exact Fractions}]
\]

%dsp-level:0
\[
(5):\ \text{the endpoint identities cleared by}\ 4m_n^2\ [\text{machine}];
\]

%dsp-level:0
\[
\qquad \text{the collinear identity in}\ \mathbb Q[P,N,m_p,m_\nu,E_p,E_\nu]/(E_p^2{-}m_p^2{-}P^2,\ E_\nu^2{-}m_\nu^2{-}N^2)\ [\text{machine}];
\]

%dsp-level:0
\[
\qquad \text{the Lagrange identity}\ [\text{machine:}\ 6\ \text{variables}];
\]

%dsp-level:0
\[
\qquad Q=\sum\mathrm{KE}_i\ [\text{machine: Laurent}]
\]

%dsp-level:0
\[
(6):\ \textstyle\sum m_{\mathrm{out}}\ge m_p+m_\ell>m_n=m_p\ [\text{machine: sign arithmetic on rational witnesses}];
\]

%dsp-level:0
\[
\qquad (7):\ \text{the scope record}.\ \blacksquare
\]

\textbf{Phys.} The vertex is assembled from necessities the corpus already proved, and then weighed. The grammar: the ladder rule moves the ternary read by one, the alphabet overflows, total charge vanishes in every valid universe, one compensating value exists, and the inventory leaves exactly one charged class and one neutral class standing. The weighing is where the corpus surprises itself honestly: every rest read it has ever derived gives the neutron and the proton the same mass — six bits over $\mathfrak h$, identical to the last monomial, because no crystallized count carries a charge exponent and every normalized average is structurally charge-blind. A two-body decay with equal parent and daughter masses strands the daughter below its own mass shell for any positive emitted mass, so the neutral member is forced wherever the vertex is kinematically admissible at all; and admissibility itself fails on every adopted read, since the Q-value collapses to minus the lepton masses — strictly negative. The one functional that splits the classes was classified as a measure price; adopting it is a specification act the corpus prices and declines. What remains is exact and complete: a conditional three-body kinematics with its endpoint formula, a general channel-closure theorem, and a verdict — the reaction is grammatically forced and kinematically closed, and the gap between them is a read specification, not a missing derivation.

\end{theorem}

\begin{theorem}[{\texorpdfstring{BM.42 — The angular momentum at the vertex: the Noether moment, the Lagrange decomposition, the absorption adjudication, and the $\ell$-selection census}{BM.42 - The angular momentum at the vertex: the Noether moment, the Lagrange decomposition, the absorption adjudication, and the -selection census}}]\label{tag:BM-42}
\textbf{Thm.}

%dsp-level:0
\[
(1)\quad\textbf{the moment current is conserved:}\qquad M^{\mu\nu\lambda}:=x^\nu T^{\mu\lambda}-x^\lambda T^{\mu\nu};
\]

%dsp-level:0
\[
\partial_\mu M^{\mu\nu\lambda}=\bigl(T^{\nu\lambda}-T^{\lambda\nu}\bigr)+x^\nu\partial_\mu T^{\mu\lambda}-x^\lambda\partial_\mu T^{\mu\nu}
\]

%dsp-level:0
\[
\qquad [\text{Leibniz};\ \text{machine: the formal identity}];
\]

%dsp-level:0
\[
T^{\mathrm{Hilb}}_{\mu\nu}\ \text{is symmetric}\ [\mathrm{BM.9}(2);\ \text{machine: the dust witness}\ \mathrm{diag}(\rho,0,0,0)]\ \text{and on-shell conserved}\ [\mathrm{BM.9}(3);\ \mathrm{BM.11}:\ \nabla_\mu T^{\mu\nu}=0\ \text{two routes}]:
\]

%dsp-level:0
\[
\boxed{\ \partial_\mu M^{\mu\nu\lambda}=0\ \ \text{on-shell}\ }
\]

%dsp-level:0
\[
\qquad [\text{machine: all}\ 64\ \text{components on the dust class}]:
\]

%dsp-level:0
\[
\qquad \textbf{the total angular momentum is the conserved charge of the promoted rotation}\ [\mathrm{BM.8}(3):\ C_3\to SO(3);\ \mathrm{Def\text{-}H.8}\ \mathrm{Rmk}:\ \text{spin coarse-grained along a geodesic is angular momentum}];
\]

%dsp-level:0
\[
(2)\quad\textbf{the budget at the vertex:}
\]

%dsp-level:0
\[
\qquad J_{\mathrm{in}}=J_{\mathrm{out}}\ \ \text{as the charge of}\ (1):
\]

%dsp-level:0
\[
\qquad j_n=j_p\oplus J_{\ell\nu}\oplus L_{\mathrm{orb}}\ \ \text{in}\ SO(3)\ \text{coupling ranges}\ |j_1-j_2|\le J\le j_1+j_2\ [\text{math}];
\]

%dsp-level:0
\[
\text{the numerical}\ j\text{-attachments ride}\ \mathrm{Def\text{-}H.38}\text{'s declared rule}\ [(7)];
\]

%dsp-level:0
\[
\qquad \text{the budget itself is the exact representation row};
\]

%dsp-level:0
\[
(3)\quad\textbf{the orbital content is already inside}\ p^2\textbf{:}
\]

%dsp-level:0
\[
\qquad \textbf{the Lagrange identity:}
\]

%dsp-level:0
\[
\qquad r^2\,\delta^{ij}p_ip_j=(\vec x\cdot\vec p)^2+\sum_{i<j}\bigl(x_ip_j-x_jp_i\bigr)^2
\]

%dsp-level:0
\[
\qquad [\text{machine: exact polynomial identity in}\ \mathbb Q[x,y,z,p_x,p_y,p_z]];
\]

%dsp-level:0
\[
p^2=p_r^2+\frac{L^2}{r^2},\qquad p_r:=\frac{\vec x\cdot\vec p}r,\quad L^2:=\sum_{i<j}L_{ij}^2,\quad L_{ij}:=x_ip_j-x_jp_i;
\]

%dsp-level:0
\[
\textbf{the dispersion in decomposed form:}
\]

%dsp-level:0
\[
\qquad r^2E^2=r^2m^2+r^2p_r^2+L^2
\]

%dsp-level:0
\[
\qquad [\text{machine: the cleared substitution into}\ \mathrm{BM.34}(4)\text{'s master identity, vacuum gauge}]:
\]

%dsp-level:0
\[
\textbf{the orbital kinetic energy needs no separate term}\ —\ \text{it is}\ \mathrm{BM.34}\text{'s}\ p^2\ \text{itself};
\]

%dsp-level:0
\[
(4)\quad\textbf{the spin-absorption adjudication:}
\]

%dsp-level:0
\[
\qquad \text{spin is the weak sector's contribution}\ [\mathrm{Def\text{-}H.8}];
\]

%dsp-level:0
\[
\qquad \textbf{its two derived metric couplings:}
\]

%dsp-level:0
\[
(\mathrm i)\ \ \text{the rest read: the pole mass is the spin-connection holonomy defect}\ m^2_{\mathrm{pole}}
\]

%dsp-level:0
\[
= 4\hat\mu'^2\sin^2(\Theta_R/4)\ [\mathrm{Def\text{-}H.37}(2);\ \mathrm{Def\text{-}H.21};\ \mathrm{FT.54}(3)\text{--}(4)]:\ \ \text{the frame coupling is inside}\ m_i;
\]

%dsp-level:0
\[
(\mathrm{ii})\ \ \text{the moving body: the shift drag}\ N^ip_i,
\]

%dsp-level:0
\[
\ \text{and the shift is the weak block}\ [\mathrm{BM.34}(3)\text{--}(4);\ \mathrm{FT.157}(1)]:\ \ \text{inside the dispersion};
\]

%dsp-level:0
\[
\text{the Casimir identity}\ \textstyle\sum_iW_i^2
\]

%dsp-level:0
\[
= 2\Pi_{\mathcal R}\ [\mathrm{Cor}\ \mathrm{CT0.7a};\ \mathrm{BM.40}(1);\ \text{machine re-derived}]\ \text{is a branch-count read}:
\]

%dsp-level:0
\[
\boxed{\ \neg\exists\ \text{a theorem constructing a scalar spin--metric coupling energy beyond}\ (\mathrm i),(\mathrm{ii})\ }
\]

%dsp-level:0
\[
\qquad [\text{inventory}];
\]

%dsp-level:0
\[
\text{a separate}\ E_{\mathrm{spin}}\ \text{term in}\ Q\ \text{would double-count}\ [\text{machine: the budget closes}\ —\ Q=\textstyle\sum_i\mathrm{KE}_i\ \text{follows from the conservation laws alone, with no residual}];
\]

%dsp-level:0
\[
\text{the frequency-read row stays on record}\ [\mathrm{FT.131}(3);\ \mathrm{FT.159}\text{'s rate-modulation residue untouched}];
\]

%dsp-level:0
\[
(5)\quad\textbf{the Q-identity with the absorptions:}
\]

%dsp-level:0
\[
\qquad \boxed{\ Q=m_n-m_p-m_\ell-m_\nu=\mathrm{KE}_p+\mathrm{KE}_\ell+\mathrm{KE}_\nu\ },
\]

%dsp-level:0
\[
\qquad \mathrm{KE}_i=N\Bigl(\sqrt{m_i^2+p_{r,i}^2+\frac{L_i^2}{r_i^2}}-m_i\Bigr);
\]

%dsp-level:0
\[
[\mathrm{BM.34}(4)\text{--}(5);\ (3);\ \text{machine: Laurent identity}];
\]

%dsp-level:0
\[
\qquad E_{\mathrm{spin\text{-}coupling}}=E_{\mathrm{orbital\text{-}coupling}}=0\ \ \text{as separate terms}\ [(3)\text{--}(4)\ \text{absorbed}];
\]

%dsp-level:0
\[
\mathrm{BM.41}(4)\text{'s sign theorem is untouched}\ —\ \text{the absorptions add no term to the budget};
\]

%dsp-level:0
\[
(6)\quad\textbf{the Carruthers--Nieto adjudication — the correction is zero:}
\]

%dsp-level:0
\[
\qquad \Delta L_z\,\Delta(\sin\theta)\ge\tfrac12|\langle\cos\theta\rangle|
\]

%dsp-level:0
\[
\qquad [\mathrm{BM.10}(5);\ \text{the naive}\ \Delta\theta\,\Delta L\ge\tfrac12\ \text{refuted};\ \text{witnesses re-cited}];
\]

%dsp-level:0
\[
\text{the kinetic corollary at}\ r=\Delta x_{\mathrm B}:
\]

%dsp-level:0
\[
\qquad \langle p^2\rangle\ \ge\ \frac{\langle L^2\rangle}{\Delta x_{\mathrm B}^2}\ \ge\ (\Delta L_z)^2\,m_{\mathrm{hid}}^2\ \ge\ \tfrac14\langle\cos\theta\rangle^2\,m_{\mathrm{hid}}^2
\]

%dsp-level:0
\[
\qquad [(3);\ \mathrm{BM.13}(4);\ \text{machine}];
\]

%dsp-level:0
\[
\mathrm{BM.34}(6)\text{'s Cartesian row:}\quad\langle p^2\rangle\ge\tfrac34m_{\mathrm{hid}}^2;
\]

%dsp-level:0
\[
\qquad \tfrac14\langle\cos\theta\rangle^2\le\tfrac14<\tfrac34
\]

%dsp-level:0
\[
\qquad [\text{machine: exact rationals}]:
\]

%dsp-level:0
\[
\boxed{\ \textbf{the CN corollary is uniformly weaker — the correction to the zero-point bounds}\ \bigl(\tfrac74,\ \tfrac38,\ \tfrac{\sqrt7}2\bigr)\ \textbf{is exactly zero}\ }
\]

%dsp-level:0
\[
(7)\quad\textbf{the spin ladder and the anchor census:}
\]

%dsp-level:0
\[
\qquad \text{under}\ \mathrm{Def\text{-}H.38}\ \text{(the declared multiplicity rule}\ \operatorname{copies}(c)=k(c)):\ \ \text{a winding-}k\ \text{chain carries}\ \mathrm{Sym}^k(\mathcal H_{\mathsf v}),
\]

%dsp-level:0
\[
\quad j=\frac{|k|}2\ [\mathrm{FT.55}(2);\ \mathrm{FT.60}(1)(\mathrm{iv})];
\]

%dsp-level:0
\[
\textbf{the Casimir ladder:}\quad j(j{+}1)=\frac k2\Bigl(\frac k2{+}1\Bigr)\ \ \text{generalizing}\ \mathrm{Cor}\ \mathrm{CT0.7a}
\]

%dsp-level:0
\[
\qquad [\mathrm{FT.60}(3);\ \text{machine: the ladder values}\ \bigl(\tfrac34,2,\tfrac{15}4,\dots\bigr)\ \text{exact rationals}];
\]

%dsp-level:0
\[
\textbf{the anchor:}
\]

%dsp-level:0
\[
\quad k=1\ \text{(the digon}\ [\mathrm{FT.28};\ \mathrm{FT.107}(1)]):\quad\mathrm{Sym}^1=\mathcal H_{\mathsf v}\cong\mathbb C^2,
\]

%dsp-level:0
\[
\quad j=\tfrac12,\quad\textstyle\sum_iJ_i^2=\tfrac34\mathbb I_2
\]

%dsp-level:0
\[
\qquad [\mathrm{EW.32};\ \text{machine}];
\]

%dsp-level:0
\[
\text{the Platonic spin column stands under its logged declaration}\ [\text{Census FT.63}(1):\ j_{\max}\in\{6,6,12,30,30\};\ \text{row-wise}];
\]

%dsp-level:0
\[
\textbf{the line-winding attachment}\ —\ \text{which}\ k\ \text{each named line carries}\ —\ \text{is}\ \neg\text{computed: the ledger mechanism assigning}\ \mathrm{Sym}^k\ \text{to materialized chains is the residue}\ [\mathrm{FT.55}(4)];
\]

%dsp-level:0
\[
(8)\quad\textbf{the selection census (the anchor class, all four lines at}\ j{=}\tfrac12\textbf{):}
\]

%dsp-level:0
\[
\qquad 2\otimes2\otimes2=4\oplus2\oplus2,\quad j\in\{\tfrac32,\tfrac12\}
\]

%dsp-level:0
\[
\qquad [\mathrm{FT.88}(7);\ \text{machine: the charpoly of}\ (\textstyle\sum J_i)^2\ \text{on}\ \mathbb C^8=(x-\tfrac{15}4)^4(x-\tfrac34)^4];
\]

%dsp-level:0
\[
\textbf{the lepton pair:}\quad\tfrac12\otimes\tfrac12=0\oplus1
\]

%dsp-level:0
\[
\qquad [\text{machine: the charpoly on}\ \mathbb C^4=x^3(x-2)];
\]

%dsp-level:0
\[
\qquad \textbf{the full vertex:}\quad2^{\otimes4}:\ \ 2\cdot(j{=}0)\oplus3\cdot(j{=}1)\oplus1\cdot(j{=}2),\quad2+9+5=16
\]

%dsp-level:0
\[
\qquad [\text{machine: the charpoly on}\ \mathbb C^{16}=x^2(x{-}2)^9(x{-}6)^5];
\]

%dsp-level:0
\[
\textbf{the}\ \ell\textbf{-coupling census}\ (J_{\mathrm{int}}:=j_p\otimes(\ell\text{-channel}\otimes\nu\text{-channel}),\ j_n=\tfrac12):
\]

%dsp-level:0
\[
\begin{array}{c|c|c|c}\text{pair spin}\ S & J_{\mathrm{int}}\in & \ell=0 & \ell\ge1\\\hline0\ (\text{singlet}) & \{\tfrac12\} & \checkmark\ (J{=}\tfrac12\ni\tfrac12) & \ell{=}1:\ \checkmark;\quad\ell{=}2:\ \neg\ (\{\tfrac32,\tfrac52\}\not\ni\tfrac12)\\1\ (\text{triplet}) & \{\tfrac12,\tfrac32\} & \checkmark\ (\text{both}) & \ell{=}1:\ \checkmark;\quad\ell{=}2:\ \checkmark\ \text{via}\ J_{\mathrm{int}}{=}\tfrac32\end{array}\qquad[\text{machine: exact half-integer range census}];
\]

%dsp-level:0
\[
\textbf{the leading admissible class:}\ \ell=0\ \text{with either pair-spin channel};
\]

%dsp-level:0
\[
\qquad \ell
\]

%dsp-level:0
\[
\ge 1\ \text{is the orbital ladder}\ —\ \text{the structural allowed/forbidden classification}\ [\neg\text{SM labels imported}];
\]

%dsp-level:0
\[
(9)\quad\textbf{the parity column — the scope row:}
\]

%dsp-level:0
\[
\qquad \textbf{the orbital factor}\ (-1)^\ell\ \textbf{is derived:}\quad\chi_\ell(\pi)=\sum_{m=-\ell}^{\ell}(-1)^m=(-1)^\ell\ \ \text{for}\ \ell=0,\dots,6
\]

%dsp-level:0
\[
\qquad [\text{the integer}\ SO(3)\ \text{character at angle}\ \pi;\ \text{machine: exact integers}];
\]

%dsp-level:0
\[
\textbf{the lines' intrinsic parity labels are}\ \neg\textbf{derived:}\quad\neg\exists\ \text{a Hilbert-space parity operator on the pair register}\ [\mathrm{FT.135}(4)\ \text{stands}];
\]

%dsp-level:0
\[
\qquad \text{parity acts at the label level as the branch-label flip}\ \iota\ [\mathrm{FT.135}(1);\ \mathrm{EM.24}(2)];
\]

%dsp-level:0
\[
\text{the four gradings are pairwise distinct}\ [\mathrm{FT.136};\ \mathrm{FT.141}(6);\ \mathrm{FT.144}(5)]:
\]

%dsp-level:0
\[
\qquad \textbf{the}\ J^P\ \textbf{row-labels}\ \bigl(\tfrac12^+\ \text{and kin}\bigr)\ \textbf{are external vocabulary and are}\ \neg\textbf{imported}\ [\mathrm{P5}];
\]

%dsp-level:0
\[
\text{the derived content is the}\ \ell\text{-ladder with its}\ (-1)^\ell\ \text{factor, coupled by}\ (8)\text{'s census}.
\]

\textbf{Rmk (budget, not labels).} The vertex conserves the charge the promoted rotation supplies: the moment current closes on the symmetry of the Hilbert stress tensor and its on-shell conservation, both rows the Noether tier proved. The budget's two would-be extra terms dissolve on inspection — the orbital piece is $p^2$ itself by Lagrange, and the spin piece is the frame holonomy inside the rest read plus the weak shift drag inside the dispersion; a separate spin-coupling energy would count the same coupling twice, and the Carruthers–Nieto refinement of the angular spread is uniformly weaker than the Cartesian zero-point row, so the bounds stand uncorrected. What the corpus does not do is attach $j$-numbers to named lines beyond the anchor class: the census computes the representation structure — two singlets, three triplets, one quintet across the fourfold vertex — and the line-data attachment waits on the winding residue.

\textbf{Req.} A3 (Part 1); D1.8 (Parts 1–2); CT0.7(1), CT1.2, CT1.4, Cor CT0.7a (Part 1); CT2.6 (Part 2); P5 (Parts 1–2); Def-H.8, Def-H.21, Def-H.38, FT.54(3), FT.54(4), FT.55(2), FT.55(4), FT.60(1), FT.60(3), FT.88(7), FT.101, FT.107(1), FT.131(3), FT.135(1), FT.135(4), FT.136, FT.141(6), FT.144(5), FT.157(1) (Part 3); Census FT.28, FT.63(1) (row-wise); Def-BM.10, BM.8(3), BM.9(2), BM.9(3), BM.10(5), BM.11, BM.13(4), BM.34(3), BM.34(4), BM.34(5), BM.34(6), BM.40(1), BM.41(4), BM.41(5) (this part); the Leibniz/jet identities, the Lagrange identity, Clebsch–Gordan range censuses, charpolys by Faddeev–LeVerrier over $\mathbb Q$, exact half-integer arithmetic [math; machine: \texttt{checks/\allowbreak{}verify\_\allowbreak{}H\_\allowbreak{}neutron\_\allowbreak{}decay\_\allowbreak{}vertex.py} §D].

\textbf{Pf.}

%dsp-level:0
\[
(1):\ \partial_\mu(x^\nu T^{\mu\lambda})=(\partial_\mu x^\nu)T^{\mu\lambda}+x^\nu\partial_\mu T^{\mu\lambda}=T^{\nu\lambda}+x^\nu\partial_\mu T^{\mu\lambda}\ [\text{Leibniz}];
\]

%dsp-level:0
\[
\qquad \text{the dust witness:}\ T
\]

%dsp-level:0
\[
= \mathrm{diag}(\rho,0,0,0),
\]

%dsp-level:0
\[
\ \rho\ \text{constant}:\ \partial_\mu M^{\mu\nu\lambda}
\]

%dsp-level:0
\[
= T^{\nu\lambda}-T^{\lambda\nu}
\]

%dsp-level:0
\[
= 0\ \text{all}\ 64\ \text{components}\ [\text{machine: symbolic differentiation in}\ \mathbb Q[x_0,x_1,x_2,x_3,\rho]]
\]

%dsp-level:0
\[
(2):\ \text{the coupling ranges are the}\ SO(3)\ \text{representation rows}\ [\text{math}];
\]

%dsp-level:0
\[
\qquad (3):\ r^2p^2-(\vec x\cdot\vec p)^2=\sum_{i<j}(x_ip_j-x_jp_i)^2\ [\text{machine: expansion in}\ 6\ \text{variables}];
\]

%dsp-level:0
\[
\qquad \text{the substitution into}\ \mathrm{BM.34}(4)\text{'s identity at}\ N{=}1,N^i{=}0,\gamma{=}\delta,\ \text{cleared by}\ r^2\ [\text{machine}]
\]

%dsp-level:0
\[
(4):\ \mathrm{Def\text{-}H.37}(2)/\mathrm{Def\text{-}H.21}\ \text{re-cited};
\]

%dsp-level:0
\[
\qquad \mathrm{BM.34}(3)\text{--}(4)/\mathrm{FT.157}(1)\ \text{re-cited};
\]

%dsp-level:0
\[
\qquad \text{the inventory scan: no scalar spin--metric coupling term exists in Parts 1--4}\ [\text{machine: token scan of the energy-budget rows}];
\]

%dsp-level:0
\[
\qquad \text{the budget closure:}\ m_n=E_p+E_\ell+E_\nu\implies Q=\sum(E_i-m_i)\ [\text{machine: Laurent}]
\]

%dsp-level:0
\[
(5):\ (3)\text{'s decomposition substituted into}\ \mathrm{BM.34}(5)\text{'s KE definition}\ [\text{machine}]
\]

%dsp-level:0
\[
(6):\ (\Delta L_z)^2\ge\tfrac14\langle\cos\theta\rangle^2\ [\mathrm{BM.10}(5)\ \text{at}\ \Delta(\sin\theta)\le1];
\]

%dsp-level:0
\[
\qquad \langle L^2\rangle\ge(\Delta L_z)^2\ \text{at}\ \langle L_z\rangle=0;
\]

%dsp-level:0
\[
\qquad r=\Delta x_{\mathrm B}=1/m_{\mathrm{hid}}\ [\mathrm{Def\text{-}BM.4}(1);\ \mathrm{BM.13}(4)];
\]

%dsp-level:0
\[
\qquad \tfrac14<\tfrac34\ [\text{machine}]
\]

%dsp-level:0
\[
(7):\ \mathrm{Def\text{-}H.38}\text{'s derived-content rows re-cited};
\]

%dsp-level:0
\[
\qquad \text{the anchor:}\ J_i=\tfrac12\sigma_i:\ \sum_iJ_i^2=\tfrac14\sum_i\sigma_i^2=\tfrac14\cdot3\,\mathbb I_2=\tfrac34\mathbb I_2\ [\text{machine over}\ \mathbb Q(i)];
\]

%dsp-level:0
\[
\qquad \text{the ladder values}\ \tfrac k2(\tfrac k2{+}1)
\]

%dsp-level:0
\[
= (\tfrac34,2,\tfrac{15}4,6,\dots)\ \text{at}\ k
\]

%dsp-level:0
\[
= 1,
\]

%dsp-level:0
\[
2,
\]

%dsp-level:0
\[
3,
\]

%dsp-level:0
\[
4\ [\text{machine: exact rationals}]
\]

%dsp-level:0
\[
(8):\ (\textstyle\sum_iJ_i^{(2)})^2\ \text{on}\ \mathbb C^4:\ \text{charpoly}\ x^3(x-2)\ [\text{machine: Faddeev--LeVerrier}];
\]

%dsp-level:0
\[
\qquad \text{on}\ \mathbb C^8:\ (x-\tfrac{15}4)^4(x-\tfrac34)^4;
\]

%dsp-level:0
\[
\qquad \text{on}\ \mathbb C^{16}:\ x^2(x-2)^9(x-6)^5\ [\text{machine: exact charpolys over}\ \mathbb Q];
\]

%dsp-level:0
\[
\qquad \text{the}\ \ell\text{-ranges:}\ |J_{\mathrm{int}}-\ell|\le\tfrac12\le J_{\mathrm{int}}+\ell\ \text{censused on half-integers}\ [\text{machine}]
\]

%dsp-level:0
\[
(9):\ \sum_{m=-\ell}^{\ell}(-1)^m=(-1)^\ell\ \text{for}\ \ell=0..6\ [\text{machine: integer sums}];\qquad\mathrm{FT.135}(4)\ \text{re-cited}.\ \blacksquare
\]

\textbf{Phys.} Angular momentum at the vertex is a theorem before it is a table. The moment of the Hilbert stress tensor is conserved wherever the tensor is symmetric and its own divergence vanishes on-shell — both rows the Noether tier proved — so the charge the promoted rotation carries balances across the decay exactly. The budget then needs no new terms: the orbital share is already the momentum square, split by Lagrange's identity into radial and angular pieces with nothing left over, and the spin share lives inside the rest mass as frame holonomy and inside the dispersion as the weak shift's drag. Asked for a spin-coupling energy to add to the Q-value, the corpus answers that adding one would double-count, and the angular uncertainty refinement of the zero-point bound comes out uniformly weaker — a quarter against three quarters — so the bounds stand. What is genuinely new is the census: at the anchor winding, four spin-half lines tensor into two singlets, three triplets and one quintet; the lepton pair couples to singlet or triplet; and the orbital ladder admits zero and one against both pair channels, two only through the triplet route. The parity factor of the ladder is the integer character of the rotation group, while the parity labels of the lines themselves stay outside: the register has no parity operator, and the tier imports no external $J^P$ vocabulary.

\end{theorem}

\begin{theorem}[{\texorpdfstring{BM.43 — The curvature rows of the members: the dust sources, the weak-field potential, and the superposition scope}{BM.43 - The curvature rows of the members: the dust sources, the weak-field potential, and the superposition scope}}]\label{tag:BM-43}
\textbf{Thm.}

%dsp-level:0
\[
(1)\quad\textbf{the baryon rows (the exact monomials):}
\]

%dsp-level:0
\[
\qquad E_{\mathrm{int}}=6\ \text{for both classes}\ [\mathrm{BM.17};\ \mathrm{BM.26};\ \text{Census FT.35}(2)]:
\]

%dsp-level:0
\[
\Delta x_{\mathrm B}=\frac{\mathfrak h}{6\ln2},
\]

%dsp-level:0
\[
\qquad V_{\mathrm{eff}}=\frac{\mathfrak h^3}{216(\ln2)^3},
\]

%dsp-level:0
\[
\qquad m=\frac{6\ln2}{\mathfrak h},
\]

%dsp-level:0
\[
\qquad \rho=\frac{1296(\ln2)^4}{\mathfrak h^4},
\]

%dsp-level:0
\[
\qquad \kappa^2\rho=\frac{2592\pi(\ln2)^3}{\mathfrak h^2},
\]

%dsp-level:0
\[
\qquad r_s=3\,\mathfrak h;
\]

%dsp-level:0
\[
\textbf{the lock:}\quad(\Delta x_{\mathrm B})^2\,\kappa^2\rho=2\pi E_{\mathrm{int}}^2\ln2=72\pi\ln2
\]

%dsp-level:0
\[
\qquad [\mathrm{BM.26};\ \mathrm{BM.16}(1);\ \text{machine: every row one exact Laurent monomial}];
\]

%dsp-level:0
\[
\text{the parent and the daughter carry identical rows}\ [\mathrm{BM.28}(1)\ \text{charge-blindness}];
\]

%dsp-level:0
\[
(2)\quad\textbf{the lepton rows (symbolic):}
\]

%dsp-level:0
\[
\qquad \text{the taxonomy law holds:}\quad m_i=\frac{E_{\mathrm{int}}^{(i)}\ln2}{\mathfrak h}\ [\mathrm{FT.44}(4)];
\]

%dsp-level:0
\[
\qquad E_{\mathrm{int}}^{(\ell)},\ E_{\mathrm{int}}^{(\nu)}\ \text{are}\ \neg\text{computed at any floor}:
\]

%dsp-level:0
\[
\text{the rows stand as symbolic monomials:}\quad\kappa^2\rho_i=\frac{2\pi\,\mathfrak h\,E_{\mathrm{int}}^{(i)}}{V_i},
\]

%dsp-level:0
\[
\qquad r_s^{(i)}=\frac{\mathfrak h\,E_{\mathrm{int}}^{(i)}}2
\]

%dsp-level:0
\[
\qquad [\mathrm{BM.9}(5);\ \mathrm{BM.16}(4);\ \text{machine}];
\]

%dsp-level:0
\[
\text{the transport alternative:}\quad m_i\ \text{from}\ \mathrm{BM.39}\ (\text{symbolic in}\ (g,g',\hat\mu,\ell_c));
\]

%dsp-level:0
\[
\qquad \text{the conversion}\ E_{\mathrm{int}}^{(i)}\leftrightarrow\text{transport data is item-7 residue}\ [\mathrm{Def\text{-}BM.11}(7);\ \text{pointer}];
\]

%dsp-level:0
\[
(3)\quad\textbf{the field-equation rows:}
\]

%dsp-level:0
\[
\qquad T_{\mu\nu}=\mathrm{diag}(\rho,0,0,0)\ \text{at rest},\ p=0\ [\mathrm{BM.9}(4);\ \mathrm{CT4.5}];
\]

%dsp-level:0
\[
\qquad R=\kappa^2(\rho-3p)+4\Lambda\ [\mathrm{BM.15}(2)];
\]

%dsp-level:0
\[
\qquad RV=\varkappa^2m\ [\mathrm{BM.20}(1)];
\]

%dsp-level:0
\[
(4)\quad\textbf{the weak-field potential (the linearized row):}
\]

%dsp-level:0
\[
\qquad g_{00}=-(1+2\Phi),\quad g_{ij}=(1-2\Phi)\delta_{ij},\quad\Phi\ \text{static}:
\]

%dsp-level:0
\[
G^{(1)}_{00}=2\,\nabla^2\Phi,
\]

%dsp-level:0
\[
\qquad G^{(1)}_{0i}=0,
\]

%dsp-level:0
\[
\qquad G^{(1)}_{ij}=0
\]

%dsp-level:0
\[
\qquad [\text{machine: Christoffel}\to\text{Ricci}\to\text{Einstein in the weight-truncated jet ring, exact rational coefficients}];
\]

%dsp-level:0
\[
\text{the field equation}\ G_{\mu\nu}=\kappa^2T_{\mu\nu}\ [\mathrm{CT1.23}]\ \text{at}\ p=0:
\]

%dsp-level:0
\[
\qquad \nabla^2\Phi=\frac{\kappa^2}2\,\rho=4\pi G_N\rho
\]

%dsp-level:0
\[
\qquad [\kappa^2=8\pi G_N\ [\mathrm{CT2.9}];\ \text{machine: the coefficient identity}\ \tfrac{\kappa^2}2=\tfrac{\pi\mathfrak h^2}{\ln2}=4\pi G_N];
\]

%dsp-level:0
\[
\Phi=-\frac{G_Nm}{r}\ \text{solves at}\ r\ne0:
\]

%dsp-level:0
\[
\qquad \nabla^2\frac1r=0
\]

%dsp-level:0
\[
\qquad [\text{machine:}\ -3r^2+3(x^2{+}y^2{+}z^2)\equiv0\ \text{in}\ \mathbb Q[x,y,z,r]/(r^2{-}x^2{-}y^2{-}z^2)];
\]

%dsp-level:0
\[
G_Nm_{\mathrm{hid}}=\frac{\mathfrak h\,E_{\mathrm{int}}}4,
\]

%dsp-level:0
\[
\qquad r_s=2G_Nm_{\mathrm{hid}}=\frac{\mathfrak h\,E_{\mathrm{int}}}2,
\]

%dsp-level:0
\[
\qquad \Phi(\Delta x_{\mathrm B})=-\frac{E_{\mathrm{int}}^2\ln2}4\in\mathfrak q^0
\]

%dsp-level:0
\[
\qquad [\mathrm{CT2.9};\ \mathrm{BM.16}(4);\ \text{machine: monomials}];
\]

%dsp-level:0
\[
(5)\quad\textbf{the superposition:}
\]

%dsp-level:0
\[
\qquad \textbf{source side, exact:}
\]

%dsp-level:0
\[
\quad T_{\mathrm{total}}=\sum_iT^{(i)}\ \text{on disjoint supports}\ [\mathrm{CT1.22}\ \text{channel decomposition};\ \mathrm{D1.5}]:
\]

%dsp-level:0
\[
\qquad \kappa^2\rho_{\mathrm{total}}
\]

%dsp-level:0
\[
= \kappa^2\bigl(\rho_p+\rho_\ell+\rho_\nu\bigr)
\]

%dsp-level:0
\[
= \frac{2\pi\mathfrak h}{V}\bigl(6+E_{\mathrm{int}}^{(\ell)}+E_{\mathrm{int}}^{(\nu)}\bigr)\ [\text{machine: the monomial sum}];
\]

%dsp-level:0
\[
\textbf{metric side:}\quad\text{superposition is exact only at linearized order}\ [\mathrm{CT1.12}:\ \varkappa\|q\|_{\mathrm{op}}<1\ \text{the perturbative regime}];
\]

%dsp-level:0
\[
\qquad \text{the full non-linear joint solution is}\ \neg\text{constructed};
\]

%dsp-level:0
\[
\textbf{the scope row:}\quad\text{the Schwarzschild-type content on record is the}\ r_s\ \text{read and the weak-field potential}\ [(4);\ \mathrm{BM.16}(4)];
\]

%dsp-level:0
\[
\qquad \text{a full exterior solution is}\ \neg\text{in the inventory}\ [\text{the pp-wave witness is the exact curvature row on record}];
\]

%dsp-level:0
\[
(6)\quad\textbf{status:}
\]

%dsp-level:0
\[
\qquad \text{the baryon rows are numerical monomials};
\]

%dsp-level:0
\[
\ \text{the lepton rows symbolic};
\]

%dsp-level:0
\[
\ \text{the vertex curvature total is their exact sum on the source side};
\]

%dsp-level:0
\[
\qquad \neg\text{a measured curvature or radius enters}\ [\mathrm{P5}].
\]

\textbf{Req.} D1.5, D2.9 (Parts 1–2); CT1.12, CT1.22, CT1.23, CT2.1, CT2.9, CT4.5 (Part 2); P3, P5 (Parts 1–2); FT.44(4) (Census); Def-BM.4, Def-BM.10, Def-BM.11, BM.9(4), BM.9(5), BM.15(2), BM.16(1), BM.16(4), BM.17, BM.20(1), BM.26, BM.28(1), BM.34(4), BM.39 (this part); the weight-truncated jet-ring calculus, Laurent monomials, harmonic identities in quotient rings [math; machine: \texttt{checks/\allowbreak{}verify\_\allowbreak{}H\_\allowbreak{}neutron\_\allowbreak{}decay\_\allowbreak{}vertex.py} §E].

\textbf{Pf.}

%dsp-level:0
\[
(1):\ \mathrm{BM.26}\text{'s six rows re-verified as monomials at}\ E_{\mathrm{int}}
\]

%dsp-level:0
\[
= 6\ [\text{machine}]:\quad\rho
\]

%dsp-level:0
\[
= \tfrac{6\ln2}{\mathfrak h}\cdot\tfrac{216(\ln2)^3}{\mathfrak h^3}
\]

%dsp-level:0
\[
= \tfrac{1296(\ln2)^4}{\mathfrak h^4};
\]

%dsp-level:0
\[
\quad\kappa^2\rho
\]

%dsp-level:0
\[
= \tfrac{2\pi\mathfrak h^2}{\ln2}\cdot\tfrac{1296(\ln2)^4}{\mathfrak h^4}
\]

%dsp-level:0
\[
= \tfrac{2592\pi(\ln2)^3}{\mathfrak h^2};
\]

%dsp-level:0
\[
\quad\text{the lock}\ [\text{machine}]
\]

%dsp-level:0
\[
(2):\ \mathrm{FT.44}(4)\text{'s law re-cited};\qquad\text{the symbolic monomials}\ [\text{machine}]
\]

%dsp-level:0
\[
(3):\ \text{frozen rows re-cited}
\]

%dsp-level:0
\[
(4):\ \text{the jet-ring computation:}\ g^{\mu\nu}g_{\nu\rho}=\delta^\mu_\rho\ \text{to first order}\ [\text{machine}];
\]

%dsp-level:0
\[
\qquad \Gamma^i_{00}=\partial_i\Phi,\ \ \Gamma^0_{i0}=\partial_i\Phi,\ \ \Gamma^k_{ij}=-\partial_i\Phi\,\delta^k_j-\partial_j\Phi\,\delta^k_i+\partial^k\Phi\,\delta_{ij}\ [\text{machine}];
\]

%dsp-level:0
\[
R_{00}=\nabla^2\Phi,
\]

%dsp-level:0
\[
\quad R_{ij}=\delta_{ij}\nabla^2\Phi,
\]

%dsp-level:0
\[
\quad R=2\nabla^2\Phi\implies G_{00}=2\nabla^2\Phi,\ G_{ij}=0\ [\text{machine: weight-}\le1\ \text{truncation}];
\]

%dsp-level:0
\[
\qquad \tfrac{\kappa^2}2=\tfrac{\pi\mathfrak h^2}{\ln2}=4\pi\cdot\tfrac{\mathfrak h^2}{4\ln2}\ [\mathrm{CT2.9};\ \text{machine}]
\]

%dsp-level:0
\[
\nabla^2(r^{-1})\cdot r^5=-3r^2+3\textstyle\sum_ix_i^2\equiv0\ [\text{machine: the rule ring}];
\]

%dsp-level:0
\[
\qquad G_Nm=\tfrac{\mathfrak h^2}{4\ln2}\cdot\tfrac{E\ln2}{\mathfrak h}=\tfrac{\mathfrak h E}4;
\]

%dsp-level:0
\[
\quad r_s=2G_Nm=\tfrac{\mathfrak h E}2;\quad\Phi(\Delta x_{\mathrm B})=-\tfrac{\mathfrak h E}4\cdot\tfrac{E\ln2}{\mathfrak h}=-\tfrac{E^2\ln2}4\ [\text{machine}]
\]

%dsp-level:0
\[
(5):\ \mathrm{CT1.22}\text{'s decomposition re-cited};\qquad\text{the monomial sum}\ [\text{machine}];\qquad(6):\ \text{the status record}.\ \blacksquare
\]

\textbf{Phys.} Every member of the vertex gravitates, and the corpus writes each gravity in monomials. The two baryon rows are numbers: six bits over $\mathfrak h$, a density of twelve ninety-six $(\ln2)^4$ over $\mathfrak h^4$, a horizon read of three $\mathfrak h$, a lock of seventy-two $\pi\ln2$ — identical for parent and daughter, because no crystallized count carries charge. The lepton rows keep their edge counts symbolic, for the floors that would compute them are not in the census, and the transport species offers the alternative symbolic form. The field side is linearized and exact: the static diagonal perturbation solves the ledger's own field equations with the potential obeying the Poisson row whose coefficient is the corpus's $4\pi G_N$ to the last factor, and the point-source form is harmonic away from the source by a polynomial identity. Superposition is exact where the corpus can prove it — the source side, where stress tensors add — and scoped where it cannot: the full non-linear metric of three flying members is an open construction, and the tier says so rather than dress a linearized row in exact clothes.

\end{theorem}

\begin{theorem}[{\texorpdfstring{BM.44 — The time evolution: the per-tick weight, the e-folding count, and the unit discipline}{BM.44 - The time evolution: the per-tick weight, the e-folding count, and the unit discipline}}]\label{tag:BM-44}
\textbf{Thm.}

%dsp-level:0
\[
(1)\quad\textbf{the frozen rate rows:}
\]

%dsp-level:0
\[
\qquad \sum_n\langle\varnothing|\hat U^n|\varnothing\rangle w^n=\frac1{1-w^2x_+x_-}\ \ \text{in}\ \mathbb Q[w,x_+,x_-]/(w^{12})
\]

%dsp-level:0
\[
\qquad [\mathrm{FT.108}(5);\ \mathrm{BM.23}(7);\ \text{machine: the resolvent identity re-verified coefficient-wise}];
\]

%dsp-level:0
\[
x_+x_-=4\varepsilon_q(1-4\varepsilon_q)=\frac14-16\Bigl(\varepsilon_q-\frac18\Bigr)^2\le\frac14,
\]

%dsp-level:0
\[
\qquad \text{equality}\iff\varepsilon_q=\frac18
\]

%dsp-level:0
\[
\qquad [\mathrm{BM.19}(2);\ \text{machine: the SOS completion}\ +\ \text{the grid census over}\ \varepsilon_q=k/256];
\]

%dsp-level:0
\[
\text{the pole sits at}\ |w|\ge2;
\]

%dsp-level:0
\[
\qquad \varepsilon_q
\]

%dsp-level:0
\[
= 0
\]

%dsp-level:0
\[
\implies x_+x_-
\]

%dsp-level:0
\[
= 0\implies\text{the resolvent is}\ 1:\ \ \textbf{no transition weight beyond the identity}\ [\mathrm{BM.23}(8);\ \text{machine}];
\]

%dsp-level:0
\[
(2)\quad\textbf{the unit-fugacity slice:}
\]

%dsp-level:0
\[
\qquad \Gamma_{\mathrm{tr}}(\varepsilon_q):=\frac{x_+x_-}{1-x_+x_-}\bigg|_{w=1}
\]

%dsp-level:0
\[
\qquad [\text{the transition content of the resolvent}]:
\]

%dsp-level:0
\[
\Gamma_{\mathrm{tr}}(0)=0,
\]

%dsp-level:0
\[
\qquad \Gamma_{\mathrm{tr}}\bigl(\tfrac18\bigr)=\frac{1/4}{3/4}=\frac13,
\]

%dsp-level:0
\[
\qquad \Gamma_{\mathrm{tr}}\ \text{strictly increasing on}\ \bigl[0,\tfrac18\bigr]
\]

%dsp-level:0
\[
\qquad [\text{machine:}\ \Gamma_{\mathrm{tr}}(\varepsilon_2)-\Gamma_{\mathrm{tr}}(\varepsilon_1)>0\ \text{via}\ 4(\varepsilon_2-\varepsilon_1)\bigl(1-4(\varepsilon_1+\varepsilon_2)\bigr)>0];
\]

%dsp-level:0
\[
\textbf{the exact range:}\quad\Gamma_{\mathrm{tr}}\in\bigl[0,\tfrac13\bigr]\ \ \text{on}\ \ \varepsilon_q\in\bigl[0,\tfrac18\bigr];
\]

%dsp-level:0
\[
(3)\quad\textbf{the chart evolution (the scoped specialization):}
\]

%dsp-level:0
\[
\qquad \text{under}\ \mathrm{BM.8}(2)\text{'s declared stationarity hypothesis the chart carries the time-translation one-parameter group}\ [\mathrm{BM.8}(2);\ \mathrm{BM.34}(4):\ E=-p_0\ \text{its charge}]:
\]

%dsp-level:0
\[
\frac{d}{d\tau}P(\tau)=-\tilde\Gamma\,P(\tau),
\]

%dsp-level:0
\[
\qquad P(\tau)=P_0\,e^{-\tilde\Gamma\tau}
\]

%dsp-level:0
\[
\qquad [\text{the group's exponential};\ \tilde\Gamma\ \text{the per-tick weight at the chart scale}];
\]

$\textbf{the e-folding count:}\qquad\tau_{\mathrm{life}}:=\frac1{\tilde\Gamma}\ \ \text{ticks};\qquad\tilde\Gamma=\Gamma_{\mathrm{tr}}\implies\tau_{\mathrm{life}}\ge3\ \text{ticks},\quad\tau_{\mathrm{life}}\big|_{\varepsilon_q=1/8}=3\ \ \text{exactly}\qquad[\text{machine: exact reciprocals};\ \Gamma_{\mathrm{tr}}\le\tfrac13\ \text{from}\ (2)];$

%dsp-level:0
\[
\varepsilon_q=0:\quad\Gamma_{\mathrm{tr}}=0\implies\tau_{\mathrm{life}}\ \text{has no finite value}:
\]

%dsp-level:0
\[
\qquad \textbf{a non-flickering composite never decays}\ [\mathrm{BM.23}(8)\ \text{re-cited}];
\]

%dsp-level:0
\[
(4)\quad\textbf{the ledger-native form:}
\]

%dsp-level:0
\[
\qquad \Gamma\ \text{is a weight per tick in the ledger order}\ [x^0:=\tau_n;\ \mathrm{D1.8};\ \mathrm{BM.23}(9);\ \mathrm{FT.159}];
\]

%dsp-level:0
\[
\qquad \text{the discrete evolution is the geometric resolvent of}\ (1);\
\]

%dsp-level:0
\[
\textbf{the normalization of the weight to a per-tick probability is}\ \neg\textbf{declared}\ [\mathrm{BM.23}\ \mathrm{Rmk}:\ \Gamma\ \text{is a weight},\ \neg\text{a probability}]:
\]

%dsp-level:0
\[
\qquad (3)\text{'s}\ \tilde\Gamma\ \text{is the chart specialization, scoped and marked};
\]

%dsp-level:0
\[
(5)\quad\textbf{the unit discipline:}
\]

%dsp-level:0
\[
\qquad \neg\text{external seconds};
\]

%dsp-level:0
\[
\qquad \mathrm{Cal\text{-}F}\ \text{carries no rate datum}\ (\mu,\ \hat\mu^2,\ \hat\mu'^2\ \text{only});
\]

%dsp-level:0
\[
\qquad \neg\text{a lifetime is derived in external units and}\ \neg\text{a lifetime is compared}\ [\mathrm{P5};\ \mathrm{BM.23}(9)\text{'s row stands}];
\]

%dsp-level:0
\[
\mathrm{FT.159}\text{'s rate-modulation residue}\ (\text{the flicker-frequency read of}\ \delta\tau\ \text{against}\ \Omega/N_i)\ \text{is untouched};
\]

%dsp-level:0
\[
(6)\quad\textbf{the arrow rows:}
\]

%dsp-level:0
\[
\qquad \text{forward-only}\ [\mathrm{A6}:\ \neg\exists\tau^{-1};\ \mathrm{FT.144}(2):\ \text{the prefix law};\ \mathrm{BM.23}(9)];
\]

%dsp-level:0
\[
\qquad \neg\exists\ \text{an inverse decay in the ledger};
\]

%dsp-level:0
\[
\Theta_{\mathfrak L}\ \text{is the record involution}\ (\text{the formal mirror}),
\]

%dsp-level:0
\[
\ \neg\text{a dynamical reversal}\ [\mathrm{Def\text{-}BM.7}(2);\ \mathrm{FT.144}(4)].
\]

\textbf{Rmk (weight against rate.)} The corpus's decay clock ticks in ledger order and its hand is a rational function of one threshold. Everything exact is stated exactly: the resolvent, its pole beyond two, its maximum third at the threshold midpoint, its annihilation at the sharp threshold. What is not stated is a number of seconds, because no datum in the corpus converts a tick into one — and the tier keeps that silence loud: the exponential form is the chart group's own shape under the declared stationarity scope, the e-folding count is three ticks at the heaviest admissible flicker, and the measured lifetime is quoted nowhere.

\textbf{Req.} A6, D0.4, D1.8 (Part 1); CT0.8 (Part 1); P5 (Parts 1–2); FT.108(4), FT.108(5), FT.144(2), FT.144(4), FT.159 (Part 3); Def-BM.5, Def-BM.7(2), BM.8(2), BM.19(2), BM.19(3), BM.23(7), BM.23(8), BM.23(9), BM.34(4) (this part); Cal-F; the exact resolvent identity in $\mathbb Q[w,x_+,x_-]/(w^{12})$, the SOS parabola completion, exact rational monotonicity and reciprocals [math; machine: \texttt{checks/\allowbreak{}verify\_\allowbreak{}H\_\allowbreak{}neutron\_\allowbreak{}decay\_\allowbreak{}vertex.py} §F].

\textbf{Pf.}

%dsp-level:0
\[
(1):\ (1-w^2x_+x_-)\sum_{n\le11}\bigl[x_+x_-\bigr]^{n/2}w^n\big|_{n\ \mathrm{even}}\equiv1\ \bmod w^{12}\ [\text{machine: coefficient-wise}];
\]

%dsp-level:0
\[
\qquad 4\varepsilon(1-4\varepsilon)=\tfrac14-16(\varepsilon-\tfrac18)^2\ [\text{machine: expansion}];
\]

%dsp-level:0
\[
\qquad \text{the grid census}\ k=0..63:\ \le\tfrac14,\ =\tfrac14\ \text{only at}\ k=32\ [\text{machine}]
\]

%dsp-level:0
\[
(2):\ \Gamma_{\mathrm{tr}}=\tfrac{xx}{1-xx},\ xx:=x_+x_-:\quad\Gamma_{\mathrm{tr}}(0)=0,\ \ \Gamma_{\mathrm{tr}}(\tfrac18)=\tfrac{1/4}{3/4}=\tfrac13\ [\text{machine}];
\]

%dsp-level:0
\[
\qquad \text{monotonicity:}\ xx(\varepsilon_2)-xx(\varepsilon_1)
\]

%dsp-level:0
\[
= 4(\varepsilon_2-\varepsilon_1)(1-4(\varepsilon_1+\varepsilon_2))\ [\text{machine: polynomial identity}],
\]

%dsp-level:0
\[
\ \text{positive on rational witness pairs}\ \varepsilon_1<\varepsilon_2\le\tfrac18
\]

%dsp-level:0
\[
(3):\ \mathrm{BM.8}(2)\text{'s group re-cited};
\]

%dsp-level:0
\[
\qquad \tfrac1{\Gamma_{\mathrm{tr}}(1/8)}=3\ [\text{machine: exact reciprocal}];
\]

%dsp-level:0
\[
\qquad \Gamma_{\mathrm{tr}}\le\tfrac13\implies\tau_{\mathrm{life}}\ge3\ [\text{machine: monotone reciprocal}]
\]

%dsp-level:0
\[
(4)\text{--}(6):\ \text{frozen rows re-cited}\ [\mathrm{D1.8};\ \mathrm{BM.23}(9);\ \mathrm{FT.144}(2),(4);\ \mathrm{A6}];
\]

%dsp-level:0
\[
\qquad \text{the Cal-F inventory:}\ \{\mu,\hat\mu^2,\hat\mu'^2\},\ \neg\text{a rate}\ [\text{scan}].\ \blacksquare
\]

\textbf{Phys.} Time runs through the decay exactly as the corpus runs it everywhere else: in ticks, forward, with no inverse and no seconds. The transition weight is the geometric resolvent of the flicker — a rational function of one threshold, pole beyond two, annihilated at the sharp threshold where a composite that cannot flicker cannot decay. On the unit-fugacity slice the weight climbs from zero to a third as the threshold moves to its midpoint, and the e-folding count is its reciprocal: at least three ticks, exactly three at the heaviest admissible flicker. The exponential curve belongs to the chart's own translation group under the stationarity scope the Noether tier declared; the ledger keeps only the weights. And the discipline is the point: a lifetime in seconds would need a rate datum the corpus does not carry, so the tier states the tick count, names the missing datum, and quotes no clock.

\end{theorem}

\begin{theorem}[{\texorpdfstring{BM.45 — The assembled reaction: the conservation table, the naming row, the final read, and the adjudication}{BM.45 - The assembled reaction: the conservation table, the naming row, the final read, and the adjudication}}]\label{tag:BM-45}
\textbf{Thm.}

%dsp-level:0
\[
(1)\quad\textbf{the reaction (the structural row, with the identification-layer labels):}
\]

%dsp-level:0
\[
\bigl[\varepsilon\text{-sealed 3-juncture, neutral class,}\ E_{\mathrm{int}}{=}6\bigr]\ \longrightarrow\ \bigl[\varepsilon\text{-sealed 3-juncture, unit class,}\ E_{\mathrm{int}}{=}6\bigr]\ +\ \bigl[\ell\text{-channel line},\ q{=}{-}1\bigr]\ +\ \bigl[\nu\text{-channel line,}\ q{=}0,\ \text{the materialized-arrow orientation}\bigr];
\]

%dsp-level:0
\[
\textbf{the labels:}
\]

%dsp-level:0
\[
\qquad \text{neutron}\ \longrightarrow\ \text{proton}\ +\ \text{electron}\ +\ \text{electron antineutrino}
\]

%dsp-level:0
\[
\qquad [\mathrm{Def\text{-}BM.3}(2)\text{--}(3);\ \mathrm{Def\text{-}BM.6}(1);\ \text{naming row A-BM.5}];
\]

%dsp-level:0
\[
\textbf{A-BM.5:}\quad\text{electron antineutrino}\ :
\]

%dsp-level:0
\[
= \ \text{the}\ \nu\text{-channel line in the CPT/}\omega\text{-image orientation}\ [\mathrm{BM.25}\text{'s involution};\ \mathrm{BM.27}\text{'s antiproton pattern};\ \mathrm{FT.144}(5):\ \text{the occupancy sector is derived, its name is reserved content}];
\]

%dsp-level:0
\[
\text{the charged member is}\ \mathrm{Def\text{-}BM.6}(1)\text{'s electron row}\ [\mathrm{BM.23}(5)];
\]

%dsp-level:0
\[
\qquad \text{the neutral member's line index within the}\ \nu\text{-triplet is}\ \neg\text{derived}\ [\mathrm{FT.145}(3)];
\]

%dsp-level:0
\[
(2)\quad\textbf{the conservation table:}\qquad\begin{array}{l|l|l}\text{law} & \text{the row} & \text{status}\\\hline\text{charge} & 0=(+1)+(-1)+0 & \text{derived}\ [\mathrm{EM.11};\ \mathrm{BM.23}(4);\ \text{machine}]\\\text{Janus balance} & \#J^+=\#J^-\ \text{on every emitted closed chain} & \text{derived}\ [\mathrm{D1.4};\ \mathrm{FT.144}(1)\text{--}(2);\ \mathrm{BM.7}]\\\text{rung-1 weak Casimir} & \text{branch-neutral at the}\ \sigma-vertex;\ \text{carriers internally balanced} & \text{derived}\ [\mathrm{BM.40}(3)\text{--}(4)]\\\text{angular momentum} & \partial_\mu M^{\mu\nu\lambda}=0\ \text{on-shell};\ \text{the}\ \ell\text{-census of}\ \mathrm{BM.42}(8) & \text{derived}\ [\mathrm{BM.42}(1)\text{--}(2)]\\\text{sealing phase} & \Theta_{\mathrm{int}}\in4\mathbb Z\ \text{on both baryon-type junctions} & \text{derived}\ [\mathrm{BM.2}(1);\ \text{Census FT.49}(3)]\\\text{energy--momentum} & \text{the chart-side law with the threshold row}\ Q\ge0 & \text{derived}\ [\mathrm{BM.11};\ \mathrm{BM.34};\ \mathrm{BM.41}(4)\text{--}(5)]\\\text{the arrow} & \text{forward-only};\ \neg\exists\ \text{an inverse decay} & \text{derived}\ [\mathrm{A6};\ \mathrm{FT.144}(2);\ \mathrm{BM.44}(6)]\end{array}
\]

%dsp-level:0
\[
(3)\quad\textbf{the Q-value and the distribution:}
\]

%dsp-level:0
\[
\qquad Q=-(m_\ell+m_\nu)<0\ \ \text{on every adopted read}\ [\mathrm{BM.41}(4)]:
\]

%dsp-level:0
\[
\qquad \textbf{the reaction assembles and does not go on-shell in the corpus's own balance};
\]

%dsp-level:0
\[
Q=\mathrm{KE}_p+\mathrm{KE}_\ell+\mathrm{KE}_\nu\ \ \text{with the orbital and spin content absorbed}\ [\mathrm{BM.42}(5)];
\]

%dsp-level:0
\[
\qquad \text{the endpoint formula under the}\ Q\ge0\ \text{hypothesis}\ [\mathrm{BM.41}(5)];
\]

%dsp-level:0
\[
\text{the route-(ii) variant is priced and}\ \neg\text{adopted}\ [\mathrm{BM.41}(4)(\mathrm d);\ \mathrm{BM.30}(5)];
\]

%dsp-level:0
\[
(4)\quad\textbf{the curvature rows:}
\]

%dsp-level:0
\[
\qquad \text{per member}\ [\mathrm{BM.43}]:\ \text{the baryons numerical}\ (r_s=3\,\mathfrak h,\ \kappa^2\rho=2592\pi(\ln2)^3/\mathfrak h^2,\ \text{the lock}\ 72\pi\ln2);
\]

%dsp-level:0
\[
\ \text{the leptons symbolic};
\]

%dsp-level:0
\[
\ \text{the source superposition exact};
\]

%dsp-level:0
\[
\ \text{the metric superposition linearized-scope};
\]

%dsp-level:0
\[
(5)\quad\textbf{the time rows:}
\]

%dsp-level:0
\[
\qquad \Gamma_{\mathrm{tr}}\in[0,\tfrac13]\ \text{at unit fugacity};
\]

%dsp-level:0
\[
\qquad \tau_{\mathrm{life}}\ge3\ \text{ticks};
\]

%dsp-level:0
\[
\qquad \varepsilon_q=0\implies\ \text{never}\ [\mathrm{BM.44}];
\]

%dsp-level:0
\[
\qquad \neg\text{external seconds};
\]

%dsp-level:0
\[
(6)\quad\textbf{the standing rules:}
\]

%dsp-level:0
\[
\qquad \neg\text{a P4 addition}\ [\text{§Cal-BM stays empty};\ \text{A-BM.4 is an angle identification plus a normalization convention},\ \neg\text{a numerical parameter}];
\]

%dsp-level:0
\[
\qquad \neg\text{a new dimensionful scale}\ [\mathrm{P1}];
\]

%dsp-level:0
\[
\neg\text{a measured value}\ [\mathrm{P5};\ \text{machine: zero external tokens, zero decimal numerals in this tier}];
\]

%dsp-level:0
\[
\qquad \neg\exists\ \text{an}\ [\mathrm{EXT}]\ \text{datum in any derivation};
\]

%dsp-level:0
\[
\text{species names confined to the identification layers}\ [\mathrm{Def\text{-}BM.3}(7),\ \mathrm{Def\text{-}BM.6}(7),\ \text{A-BM.5}];
\]

%dsp-level:0
\[
\qquad \text{topology labels are internal bookkeeping}\ [\mathrm{Def\text{-}H.32}(4)];
\]

%dsp-level:0
\[
\text{the extra-axiomatic list}\ [\text{A §0.2a}]\ \text{gains no member}:\ \mathrm{CT2.6}\text{'s hypotheses and}\ \mathrm{BM.8}(2)\text{'s stationarity scope are inherited at the point of use};
\]

%dsp-level:0
\[
(7)\quad\textbf{the open rows touched, not closed:}
\]

%dsp-level:0
\[
\qquad \text{the relating equation across the read species}\ [\mathrm{Def\text{-}BM.11}(7);\ \mathrm{BM.35}\ \mathrm{Rmk}\ \text{stands}];
\]

%dsp-level:0
\[
\qquad \text{the amplitude}\leftrightarrow\text{mass bridge}\ [\mathrm{BM.37}\ \text{stands}];
\]

%dsp-level:0
\[
\text{the}\ \mathrm{Sym}^k\ \text{ledger mechanism}\ [\mathrm{FT.55}(4);\ \mathrm{BM.42}(7)];
\]

%dsp-level:0
\[
\qquad \text{deficit}\ 1'\ —\ \text{the gauge coupling of}\ Q\ [\mathrm{BM.41}(3)];
\]

%dsp-level:0
\[
\qquad \neg\text{a mixing matrix is derived}\ [\mathrm{FT.145}(3);\ \mathrm{BM.41}(3)];
\]

%dsp-level:0
\[
\text{the mass-mechanism inventory: the transport-species family form joins the read-class taxonomy}\ [\mathrm{FT.84}(1)\ \text{pattern}];
\]

%dsp-level:0
\[
\qquad \text{the read specification that would open the decay kinematics}\ [\mathrm{BM.41}(4)(\mathrm e)]:\ \text{a §I addendum act},
\]

%dsp-level:0
\[
\ \neg\text{a derivation};
\]

%dsp-level:0
\[
(8)\quad\textbf{the adjudication:}
\]

%dsp-level:0
\[
\begin{array}{l|l}\text{the Yukawa requirement} & \textbf{discharged}\ —\ \neg\exists\ \text{a Yukawa term in the inventory},\ \text{and the family-profile map is constructed from the seed transport}\ [\mathrm{Def\text{-}BM.11};\ \mathrm{BM.39}]\\\text{the intra-family masses} & \textbf{derived as exact forms}\ —\ \text{the cosine/affine law, positivity, the}\ gg'\ \text{splitting, the sum rule, the order inversion, the area ladder}\ [\mathrm{BM.39}]\\\text{the neutrino} & \textbf{forced}\ \text{in every kinematically admissible vertex}\ —\ \text{the two-body impossibility}\ +\ \text{the class uniqueness};\ \mathrm{BM.23}(6)\ \text{upgraded on record};\ \text{the Casimir-deficit route adjudicated neutral}\ [\mathrm{BM.40}(3);\ \mathrm{BM.41}(2)\text{--}(3)]\\\text{angular momentum} & \textbf{conserved as the Noether moment};\ \text{orbital absorbed in}\ p^2,\ \text{spin absorbed in the rest read and the shift drag};\ \text{the CN correction exactly zero};\ \text{the}\ \ell\text{-census derived}\ [\mathrm{BM.42}]\\\text{KE and rest mass} & \textbf{derived for every member}\ \text{from the metric and the vacuum-balanced entropy}\ [\mathrm{BM.34};\ \mathrm{BM.35};\ \mathrm{BM.41}(5);\ \mathrm{BM.42}(5)]\\\text{the metric curvature} & \textbf{derived per member}\ —\ \text{numerical for the baryons, symbolic for the leptons, superposed exactly on the source side}\ [\mathrm{BM.43}]\\\text{the dynamic evolution} & \textbf{derived in ticks}\ —\ \Gamma_{\mathrm{tr}}\in[0,\tfrac13],\ \tau_{\mathrm{life}}\ge3\ \text{ticks},\ \neg\text{external seconds}\ [\mathrm{BM.44}]\\\text{the reaction} & n\to p+e^-+\bar\nu_e\ \textbf{assembled with every conservation row derived}\ —\ \textbf{and kinematically closed on every adopted read}:\ Q=-(m_\ell{+}m_\nu)<0\ [\mathrm{BM.41}(4)]\end{array}
\]

%dsp-level:0
\[
\textbf{the final read:}
\]

%dsp-level:0
\[
\qquad \text{the vertex is grammatically complete and kinematically shut};
\]

%dsp-level:0
\[
\qquad \text{what stands between the corpus and an on-shell decay is a read specification splitting the charge classes}\ —\ \text{a §I addendum act},
\]

%dsp-level:0
\[
\ \neg\text{a missing derivation}\ [\mathrm{BM.30}(5);\ \mathrm{BM.37}(3)];
\]

%dsp-level:0
\[
\qquad \text{this tier performs no such act}.
\]

\textbf{Rmk (the naming rows.)} The label \textit{electron antineutrino} attaches to the derived structural tuple — the $\nu$-channel line in the orientation sector the materialization arrow selects — exactly as \textit{antiproton} attaches to the CPT image of the proton row: the structure is derived, the occupancy sector is derived up to the reserved ray naming [FT.144(5)], and the label is logged before use (A-BM.5). Which of the three $\nu$-triplet lines the vertex emits is not derived: that attachment is mixing content, and the corpus has no mixing matrix.

\textbf{Req.} P1, P4, P5 (Parts 1–2); A6, D1.4, D1.8, D2.9 (Parts 1–2); EM.11 (Part 2); CT2.6 (Part 2); Def-H.38, FT.135(4), FT.144(2), FT.144(5), FT.145(3), FT.84(1), FT.55(4) (Part 3); Census FT.35(2), FT.44(4), FT.49(3) (row-wise); Def-BM.3, Def-BM.6, Def-BM.7, Def-BM.11, BM.2(1), BM.7, BM.11, BM.23(5), BM.25, BM.26, BM.27, BM.34, BM.35, BM.37(3), BM.39, BM.40, BM.41, BM.42, BM.43, BM.44 (this part); the declaration addendum A-BM.5 (pointers); the machine sections §A–§G green [math; machine: \texttt{checks/\allowbreak{}verify\_\allowbreak{}H\_\allowbreak{}neutron\_\allowbreak{}decay\_\allowbreak{}vertex.py} §G].

\textbf{Pf.}

%dsp-level:0
\[
(1)\text{--}(7):\ \text{assembly of the adjudicated rows of}\ \mathrm{Def\text{-}BM.11},
\]

%dsp-level:0
\[
\ \mathrm{BM.39}\text{--}\mathrm{BM.44}\ [\text{machine: §A--§F green}]
\]

%dsp-level:0
\[
(8):\ \text{conjunction of}\ (1)\text{--}(7).\ \blacksquare
\]

\textbf{Phys.} The corpus writes its first complete reaction, and the honesty of the balance is the result. Every conservation row is a theorem: charge closes on the ternary alphabet, the Janus balance rides on every emitted closed chain, the weak Casimir passes through the vertex untouched because registration signs and branch labels live in different registers, angular momentum is the conserved moment of a symmetric on-shell stress tensor, the sealing phase stays in its baryon class, and the arrow points one way. The members are forced: the charged one by the ladder rule and the alphabet, the neutral one by the impossibility of a two-body vertex at equal rest masses and the uniqueness of its class. The masses are exact where the corpus computes them — six bits over $\mathfrak h$ for both baryons — and symbolic where it does not. And then the balance itself: parent and daughter weigh the same on every read the corpus adopts, so the Q-value is minus the lepton masses, strictly negative, and the reaction does not go on-shell. The gap between the assembled vertex and a running decay is not a missing mechanism; it is a read specification that would split the charge classes, priced since Tier BM.8 and declined here as everywhere. The neutron decay vertex is the corpus in miniature: complete grammar, exact arithmetic, forced members, closed kinematics, and no external number anywhere in it.

\end{theorem}

\phantomsection
\section[Tier BM.12 — The EM Self-Interaction Correction: the Scar's…]{\texorpdfstring{Tier BM.12 — The EM Self-Interaction Correction: the Scar's Field Energy, the Infinite-Range Seal Asymmetry, the Mass-Splitting Theorem, and the Beta-Spectrum Adjudication}{Tier BM.12 - The EM Self-Interaction Correction: the Scar's Field Energy, the Infinite-Range Seal Asymmetry, the Mass-Splitting Theorem, and the Beta-Spectrum Adjudication}}

\begin{definition}[{\texorpdfstring{Def-BM.12 — The chart-side EM source read of a sealed composite (the construction Def-BM.8(8) specifies)}{Def-BM.12 - The chart-side EM source read of a sealed composite (the construction Def-BM.8(8) specifies)}}]\label{tag:Def-BM-12}
\textbf{Def.} A read constructed on the chart side of a sealed composite; no new axiom, no new datum, no new parameter, no numerical coupling. Local notation: $\varepsilon_{\mathrm{EM}}:=e^2/(4\pi)$ — a dimensionless P4-class symbol throughout, $\neg$identified with any measured constant [EM.25(7), EM.25(9); P5].

%dsp-level:0
\[
(1)\quad\textbf{the current:}
\]

%dsp-level:0
\[
\qquad J^\mu_Q:=\text{the continuum density of the integer}\ Q\ \text{read}
\]

%dsp-level:0
\[
\qquad [\mathrm{Def\text{-}BM.8}(1);\ \mathrm{BM.7}:\ \text{the density form of the integer law};\ \mathrm{BM.11}:\ \partial_\mu J^\mu=0];
\]

%dsp-level:0
\[
Q\in\mathbb Z,\ \ Q\in\{-3,\dots,+3\}\ \text{at}\ k\le3\ [\mathrm{Def\text{-}BM.8}(5);\ \mathrm{BM.28}(7)];
\]

%dsp-level:0
\[
\qquad \text{the charge alphabet of the generator matches the end-read alphabet:}\ \operatorname{spec}Q|_{\mathcal H_{\mathsf v}}\cup\omega\text{-image}
\]

%dsp-level:0
\[
= \{-1,0,+1\}\ [\mathrm{FT.12};\ \mathrm{EW.36};\ \mathrm{D3.4}];
\]

%dsp-level:0
\[
(2)\quad\textbf{the coupling:}
\]

%dsp-level:0
\[
\qquad \text{the}\ \mathsf A\text{-coefficient of the neutral connection is}\ e\,Q,
\]

%dsp-level:0
\[
\qquad e=\frac{gg'}{\sqrt{g^2+g'^2}}=g\sin\theta_W
\]

%dsp-level:0
\[
\qquad [\mathrm{EW.35};\ \mathrm{EM.25}(6):\ \text{fixed with zero freedom}];
\]

%dsp-level:0
\[
\neg\exists\ \text{a theorem fixing the numerical value of}\ e\ [\mathrm{EM.25}(7);\ \mathrm{P4}]:
\]

%dsp-level:0
\[
\qquad \textbf{every row of this tier carries}\ e\ \textbf{symbolically};
\]

%dsp-level:0
\[
(3)\quad\textbf{the sourced field equation:}
\]

%dsp-level:0
\[
\qquad \partial_\mu F^{\mu\nu}=e\,J^\nu_Q
\]

%dsp-level:0
\[
\qquad [\text{the Euler--Lagrange row of}\ S_{\mathrm{EM}}=-\tfrac14\int\sqrt{-g}\,F^2\ [\mathrm{EM.15};\ \mathrm{CT1.22}]\ \text{with the}\ (2)\text{-coupled source};\ \text{math}];
\]

%dsp-level:0
\[
\text{gauge invariance stands}\ [\mathrm{EM.14}];
\]

%dsp-level:0
\[
\qquad m_\gamma^2=0\ [\mathrm{EM.6};\ \mathrm{EW.22}:\ \ker M^2=\operatorname{span}\{Q\}]:
\]

%dsp-level:0
\[
\qquad \text{the static spherically symmetric solution outside the source support is the Coulomb read}
\]

%dsp-level:0
\[
\quad E_r(r)=\frac{e\,q}{4\pi r^2}
\]

%dsp-level:0
\[
\qquad [\text{machine: harmonicity of}\ 1/r\ +\ \text{Gauss flux normalization}];
\]

%dsp-level:0
\[
(4)\quad\textbf{the smearing locus:}
\]

%dsp-level:0
\[
\qquad \text{the source support lies within the Born-width envelope}\ \Delta x_{\mathrm B}
\]

%dsp-level:0
\[
= \mathfrak h/(E_{\mathrm{int}}\ln2)\ [\mathrm{Def\text{-}BM.4}(1),(5):\ \text{a chart-side envelope},\ \neg\text{a ledger distance};\ \mathrm{BM.13}(3):\ \text{sub-hop}];
\]

%dsp-level:0
\[
\textbf{the distribution factor is priced, not adopted}\ [\mathrm{Def\text{-}BM.4}(4)\ \text{pattern}]:
\]

%dsp-level:0
\[
\qquad \text{shell}:\ c_{\mathrm S}=\tfrac12;
\]

%dsp-level:0
\[
\qquad \text{uniform ball}:\ c_{\mathrm B}=\tfrac35
\]

%dsp-level:0
\[
\qquad [\text{machine: both energy integrals exact}];
\]

%dsp-level:0
\[
c\ \text{is carried symbolically; no clause of this tier selects a distribution};
\]

%dsp-level:0
\[
(5)\quad\textbf{the field-energy read:}
\]

%dsp-level:0
\[
\qquad E_{\mathrm{EM}}^{\mathrm{self}}(q,c):=\int_{\text{chart}}T^{(\mathrm{EM})}_{00}\,dV=\frac{c\,e^2q^2}{4\pi\,\Delta x_{\mathrm B}}
\]

%dsp-level:0
\[
\qquad [T^{(\mathrm{EM})}_{\mu\nu}=F_{\mu\rho}F_\nu{}^\rho-\tfrac14g_{\mu\nu}F^2\ [\mathrm{EM.15}];\ \text{machine: the radial integral exact with the}\ \Lambda\to\infty\ \text{row}];
\]

%dsp-level:0
\[
[\varepsilon_{\mathrm{EM}}]=\mathfrak q^0,\quad[\Delta x_{\mathrm B}]=\mathfrak q^{-1/2}\implies[E_{\mathrm{EM}}^{\mathrm{self}}]=\mathfrak q^{+1/2}=[\text{energy}]
\]

%dsp-level:0
\[
\qquad [\mathrm{CT2.1};\ \text{machine: exact exponent arithmetic}];
\]

%dsp-level:0
\[
(6)\quad\textbf{the channel location:}
\]

%dsp-level:0
\[
\qquad T_{\mu\nu}=T^{\mathrm{bd}}_{\mu\nu}+T^{\mathrm{tr}}_{\mu\nu}\ [\mathrm{CT1.22};\ \mathrm{D1.5}:\ \text{the bound/transit decomposition}]:
\]

%dsp-level:0
\[
\qquad \textbf{the correction is the transit-channel energy of the member's own field},\ \neg\text{a new mass read};
\]

%dsp-level:0
\[
\mathrm{BM.9}(4)\text{'s dust row is the bound part at rest};
\]

%dsp-level:0
\[
\qquad \text{a charged member carries}\ T^{\mathrm{bd}}+T^{\mathrm{tr}};
\]

%dsp-level:0
\[
\qquad \text{the rest input of the dispersion route refines to}\ \ m_{\mathrm{rest}}(q)
\]

%dsp-level:0
\[
= m_{\mathrm{hid}}+\frac{c\,e^2q^2}{4\pi\Delta x_{\mathrm B}}\ [\mathrm{Def\text{-}BM.10}(4)\text{--}(5);\ \mathrm{BM.35}(4)];
\]

%dsp-level:0
\[
(7)\quad\textbf{status and scope:}
\]

%dsp-level:0
\[
\qquad \textbf{constructed}\ —\ \mathrm{Def\text{-}BM.8}(8)\text{'s named open construction (the coupling of}\ Q\ \text{to a gauge field)}\ \text{is discharged for the sealed classes at static leading order};
\]

%dsp-level:0
\[
\neg\text{a declaration},\ \neg\text{a P4 addition}\ (e\ \text{is the logged combination}\ [\mathrm{EW.35}]),\ \text{§Cal-BM stays empty};
\]

%dsp-level:0
\[
\qquad \text{chart-side throughout}\ [\mathrm{CT2.6}\ \text{hypotheses at the point of use}];
\]

%dsp-level:0
\[
\mathrm{BM.7}(4)\text{'s register row stands untouched:}\ \neg\exists\ \text{a}\ U(1)\ \text{action on the ternary alphabet and}\ q_{\mathrm{EM}}\ \text{is}\ \neg\text{a register Noether charge}\ —\ \text{the coupling of this definition is chart-side, after}\ \mathrm{BM.8}\text{'s promotion}\ [\text{``the gauge algebras were never discrete''}];
\]

%dsp-level:0
\[
\neg\text{constructed: any dynamics of}\ Q\ [\mathrm{BM.6}(5)],
\]

%dsp-level:0
\[
\ \text{radiation and back-reaction},
\]

%dsp-level:0
\[
\ \text{a lepton--baryon bound channel}\ [\mathrm{Def\text{-}BM.1}\text{'s grammar is colour-sealed}],
\]

%dsp-level:0
\[
\ \text{the chiral assignment}\ [\text{deficit}\ 1',\ \mathrm{FT.125}(4)\ \text{stands}],
\]

%dsp-level:0
\[
\ \text{internal multipole resolution}\ [\mathrm{BM.46}(2)].
\]

\textbf{Rmk (what is declared and what is not).} Nothing. The coupling strength is EW.35's derived coefficient carried by a P4-class symbol; the current is Def-BM.8's derived read in BM.11's derived density form; the field equation is the Euler–Lagrange row of EM.15's derived action; the locus is Def-BM.4(5)'s existing envelope; the distribution factor is priced in two exact values and adopted in none. The construction assembles derived pieces into the sourced static read that Def-BM.8(8) listed as unconstructed — and stops exactly where the inventory stops: no dynamics, no radiation, no atoms.

\textbf{Req.} D1.5, D1.6, D2.9, D3.4 (Parts 1–2); CT1.22, CT2.1, CT2.6, EM.6, EM.11, EM.14, EM.15, EM.25(6), EM.25(7), EM.25(9), EW.22, EW.35, EW.36 (Part 2); P4, P5 (Parts 1–2); FT.12 (Part 3); Def-BM.4, Def-BM.8, Def-BM.10, BM.7, BM.9(4), BM.11, BM.13(3), BM.28(7), BM.35(4) (this part); the Euler–Lagrange row of the abelian action, Gauss's law for the spherically symmetric static class, exact radial integrals [math; machine: \texttt{checks/\allowbreak{}verify\_\allowbreak{}H\_\allowbreak{}em\_\allowbreak{}self\_\allowbreak{}interaction.py} §A].

\textbf{Pf.}

%dsp-level:0
\[
\text{Definition block: clauses}\ (1)\text{--}(6)\ \text{are assemblies of the cited derived rows with the machine-verified Coulomb/energy identities};
\]

%dsp-level:0
\[
\qquad (4)\text{'s priced factors carry no adoption};
\]

%dsp-level:0
\[
\qquad (7)\ \text{is the status and scope record}.\ \blacksquare
\]

\textbf{Phys.} The composite's charge was an integer read with no field attached to it; the attachment is now built, and every piece of it was already on the shelf. The current is the additive read in its continuum density form, the coupling coefficient is the residual coupling the electroweak rotation derived — carried as a symbol, because no theorem of the corpus fixes its value and none is invented here — and the field equation is the Euler–Lagrange row of the Maxwell action the stress-energy tier already varied. Outside the Born-width envelope the solution is the Coulomb read, its energy integral is exact at either of two priced distributions, and its home in the ledger bookkeeping is the transit channel of CT1.22's own decomposition: the correction is the field the member drags along, not a new species of mass. The register-side obstruction stands where it was: there is no $U(1)$ on the ternary alphabet, and this construction never asks for one — it lives on the chart, downstream of the promotion.

\end{definition}

\begin{theorem}[{\texorpdfstring{BM.46 — The scar's field energy and the mass-splitting theorem}{BM.46 - The scar's field energy and the mass-splitting theorem}}]\label{tag:BM-46}
\textbf{Thm.}

%dsp-level:0
\[
(1)\quad\textbf{the self-energy, two agreeing routes:}
\]

%dsp-level:0
\[
\qquad \textbf{field route:}\quad\int_{\Delta x_{\mathrm B}}^{\Lambda}\tfrac12E_r^2\cdot4\pi r^2\,
\]

%dsp-level:0
\[
dr
\]

%dsp-level:0
\[
= \frac{e^2q^2}{8\pi}\Bigl(\frac1{\Delta x_{\mathrm B}}-\frac1\Lambda\Bigr)\xrightarrow[\Lambda\to\infty]{}\frac{e^2q^2}{8\pi\,\Delta x_{\mathrm B}}
\]

%dsp-level:0
\[
\qquad [\text{machine: exact antiderivative and limit row}];
\]

%dsp-level:0
\[
\textbf{potential route (shell):}\quad\tfrac12\,q\,\phi(\Delta x_{\mathrm B})=\tfrac12(eq)\frac{eq}{4\pi\Delta x_{\mathrm B}}=\frac{e^2q^2}{8\pi\,\Delta x_{\mathrm B}}
\]

%dsp-level:0
\[
\qquad [\text{machine: the two routes agree exactly}];
\]

%dsp-level:0
\[
\textbf{the uniform ball:}\quad E_{\mathrm{in}}=\frac{e^2q^2}{40\pi\,\Delta x_{\mathrm B}},\quad E_{\mathrm{out}}=\frac{e^2q^2}{8\pi\,\Delta x_{\mathrm B}},\quad E_{\mathrm{tot}}=\frac{3\,e^2q^2}{20\pi\,\Delta x_{\mathrm B}}
\]

%dsp-level:0
\[
\qquad [\text{machine: the interior/exterior split exact}];
\]

$\textbf{the unified form:}\qquad\Delta x_{\mathrm B}\cdot m_{\mathrm{hid}}=1\ [\mathrm{Def\text{-}BM.4}(1)]:\qquad\boxed{\ E_{\mathrm{EM}}^{\mathrm{self}}=\frac{c\,e^2q^2}{4\pi\,\Delta x_{\mathrm B}}=c\,\varepsilon_{\mathrm{EM}}\,q^2\,m_{\mathrm{hid}}=\frac{c\,\varepsilon_{\mathrm{EM}}\,q^2\,E_{\mathrm{int}}\ln2}{\mathfrak h}\ }\qquad[\text{machine: one Laurent monomial per form}];$ $\text{i.e.}\qquad\textbf{the fractional correction is}\ \ \frac{E_{\mathrm{EM}}^{\mathrm{self}}}{m_{\mathrm{hid}}}=c\,\varepsilon_{\mathrm{EM}}\,q^2\in\mathfrak q^0\qquad[\text{machine: dimensionless}];$

%dsp-level:0
\[
(2)\quad\textbf{the charge selection:}
\]

%dsp-level:0
\[
\qquad \text{the correction is exactly proportional to}\ q^2:
\]

%dsp-level:0
\[
\qquad \text{the unit class}\ (q{=}{+}1):\ \delta^{(p)}=c\,\varepsilon_{\mathrm{EM}}\,m_{\mathrm{hid}}>0;
\]

%dsp-level:0
\[
\qquad \text{the neutral class}\ (q{=}0):\ \delta^{(n)}
\]

%dsp-level:0
\[
= 0\ \ \textbf{exactly at the monopole order}\ [\text{machine}:\ 0^2=0;\ \text{Gauss: zero net flux}];
\]

%dsp-level:0
\[
\textbf{the internal-structure row:}\quad\text{the neutral class contains separated-sign endpoint patterns}\ (12\ \text{of the}\ 20\ \text{probability-space members are mixed triples}\ [\mathrm{BM.28}(2);\ \text{machine}]),
\]

%dsp-level:0
\[
\ \text{but}\ \mathrm{D3.4}\text{'s end signs carry no position datum}:
\]

%dsp-level:0
\[
\text{the chart does not resolve internal multipoles at this order}\ [\mathrm{Def\text{-}BM.4}(5)\ \text{envelope locus}]:
\]

%dsp-level:0
\[
\qquad \text{a neutron internal-multipole self-energy is}\ \neg\text{derived}\ —\ \text{a resolution residue, recorded};
\]

%dsp-level:0
\[
(3)\quad\textbf{the corrected rest read:}
\]

%dsp-level:0
\[
\qquad m_{\mathrm{eff}}(q)=m_{\mathrm{hid}}+\frac{c\,e^2q^2}{4\pi\,\Delta x_{\mathrm B}}=m_{\mathrm{hid}}\bigl(1+c\,\varepsilon_{\mathrm{EM}}\,q^2\bigr)
\]

%dsp-level:0
\[
\qquad [\mathrm{Def\text{-}BM.12}(6):\ \text{the transit-channel refinement of the rest input}];
\]

%dsp-level:0
\[
\text{at}\ E_{\mathrm{int}}=6\ [\mathrm{BM.17};\ \mathrm{BM.26}]:
\]

%dsp-level:0
\[
\qquad \delta^{(p)}=\frac{6\,c\,e^2\ln2}{4\pi\,\mathfrak h}=\frac{3\,c\,e^2\ln2}{2\pi\,\mathfrak h}
\]

%dsp-level:0
\[
\qquad [\text{machine: one Laurent monomial}];
\]

%dsp-level:0
\[
\qquad [E_{\mathrm{EM}}^{\mathrm{self}}]=\mathfrak q^{+1/2}=[\text{energy}]\ [\mathrm{CT2.1};\ \text{machine}];
\]

%dsp-level:0
\[
(4)\quad\textbf{the mass-splitting theorem:}
\]

%dsp-level:0
\[
\qquad \boxed{\ \Delta m_{\mathrm{EM}}:=m_{\mathrm{eff}}^{(p)}-m_{\mathrm{eff}}^{(n)}=\frac{c\,e^2}{4\pi\,\Delta x_{\mathrm B}}=\frac{3\,c\,e^2\ln2}{2\pi\,\mathfrak h}\ >\ 0\ }
\]

%dsp-level:0
\[
\qquad [\text{machine: positive coefficient census at}\ e^2>0,\ c\in\{\tfrac12,\tfrac35\}];
\]

%dsp-level:0
\[
\textbf{the equal-mass degeneracy of the crystallized-count route}\ [\mathrm{BM.28}(1)]\ \textbf{is broken by a strictly positive, charge-selected splitting: the charged member is the heavier one};
\]

%dsp-level:0
\[
\text{the width is shared}\ (E_{\mathrm{int}}=6\ \text{on both classes}\ [\mathrm{BM.17};\ \mathrm{BM.26}])\ \text{and the selector is}\ q^2:\ (1\ \text{vs}\ 0)\ [\text{machine}]:
\]

%dsp-level:0
\[
\qquad \text{the splitting is a function of the charge and of the Born width, with the charge carrying the selection};
\]

%dsp-level:0
\[
(5)\quad\textbf{the correction-species row:}
\]

%dsp-level:0
\[
\qquad \delta_{\mathrm{EM}}\ \text{is a correction to the existing read,}\ \neg\text{a new mass read}\ [\mathrm{Def\text{-}BM.12}(6);\ \mathrm{CT1.22};\ \mathrm{D1.5}]:\ \text{it is the transit-channel energy of the member's own field, entering through the dispersion route's rest slot}\ [\mathrm{Def\text{-}BM.10}(4)\text{--}(5)];
\]

%dsp-level:0
\[
\mathrm{BM.37}(2)\text{'s disjointness stands: the measure routes}\ (\mathrm i)\text{--}(\mathrm{iii})\ \text{remain measure functionals}\ —\ \text{the correction rides the metric/field species, and the route ratios}\ \bigl(1,\tfrac43,\tfrac{58}{13}\bigr)\ \text{are untouched by it};
\]

%dsp-level:0
\[
(6)\quad\textbf{status:}
\]

%dsp-level:0
\[
\qquad \text{derived:}\ (1)\text{--}(5);
\]

%dsp-level:0
\[
\qquad \text{symbolic:}\ e\ [\mathrm{EM.25}(7)],\ c\ [\mathrm{Def\text{-}BM.12}(4)];
\]

%dsp-level:0
\[
\qquad \neg\text{a measured splitting value enters}\ [\mathrm{P5};\ \text{machine: zero external tokens}];
\]

%dsp-level:0
\[
\neg\text{a P4 addition},\ \neg\text{a new dimensionful scale}\ [\mathrm{P1}];
\]

%dsp-level:0
\[
\qquad \text{the open rows: the numerical value of}\ e\ (\text{a calibration step requiring one datum}\ [\mathrm{EM.25}(7)]),
\]

%dsp-level:0
\[
\ \text{the internal-multipole resolution}\ ((2)),
\]

%dsp-level:0
\[
\ \text{the dynamics of}\ Q\ [\mathrm{BM.6}(5)].
\]

\textbf{Rmk (the sign is the theorem).} Every piece of this correction is positive-semidefinite in the charge: the field energy of a source is the integral of a square, and it scales as $q^2$. The member that carries the scar therefore gains mass, never loses it — the proton is the heavier baryon on this route, by an exact symbolic monomial. Any account on which the electromagnetic self-interaction lightens the charged member and thereby opens the neutron decay is refuted by the sign of the integral, before any value of $e$ is known. The tier keeps the distribution factor and the coupling symbolic precisely so that the sign — the only content the claim needed — is derived without either.

\textbf{Req.} D1.5, D2.9, D3.4 (Parts 1–2); CT1.22, CT2.1, EM.11, EM.15, EM.25(7) (Part 2); P1, P4, P5 (Parts 1–2); Def-BM.4, Def-BM.8, Def-BM.10, Def-BM.12, BM.9(4), BM.13, BM.17, BM.26, BM.28(1), BM.28(2), BM.35(4), BM.37(2) (this part); exact radial integrals with the cutoff-and-limit row, Laurent monomials, sign censuses [math; machine: \texttt{checks/\allowbreak{}verify\_\allowbreak{}H\_\allowbreak{}em\_\allowbreak{}self\_\allowbreak{}interaction.py} §A–§B].

\textbf{Pf.}

%dsp-level:0
\[
(1):\ E_r=\tfrac{eq}{4\pi r^2}\ [\mathrm{Def\text{-}BM.12}(3)]:\quad\tfrac12E_r^2\cdot4\pi r^2=\tfrac{e^2q^2}{8\pi r^2};
\]

%dsp-level:0
\[
\qquad \int_R^\Lambda\tfrac{dr}{r^2}=\tfrac1R-\tfrac1\Lambda\ [\text{machine: antiderivative}];
\]

%dsp-level:0
\[
\qquad \text{the shell potential route and the ball split}\ [\text{machine: exact rational coefficients}\ \tfrac18,\ \tfrac1{40},\ \tfrac3{20}]
\]

%dsp-level:0
\[
\text{the unified form}:\ \tfrac1{\Delta x_{\mathrm B}}
\]

%dsp-level:0
\[
= m_{\mathrm{hid}}\ [\mathrm{Def\text{-}BM.4}(1)]:\quad\tfrac{c\,e^2q^2}{4\pi\Delta x_{\mathrm B}}
\]

%dsp-level:0
\[
= c\,
\]

%dsp-level:0
\[
\tfrac{e^2}{4\pi}\,
\]

%dsp-level:0
\[
q^2\,
\]

%dsp-level:0
\[
m_{\mathrm{hid}}\ [\text{machine: monomial identity}];
\]

%dsp-level:0
\[
\qquad [\varepsilon_{\mathrm{EM}}]=[e^2]=\mathfrak q^0\ [\mathrm{EW.35}:\ e\ \text{dimensionless};\ \text{machine}]
\]

%dsp-level:0
\[
(2):\ q=0\implies e^2q^2=0\ [\text{machine}];
\]

%dsp-level:0
\[
\qquad W(0)=20=8+12\ \text{with}\ 12\ \text{mixed triples}\ [\mathrm{BM.28}(2)\ \text{re-enumerated; machine}];
\]

%dsp-level:0
\[
\qquad \text{the position-datum row:}\ \mathrm{D3.4}\text{'s read is a pair of signs}\ [\text{verbatim scan}]
\]

%dsp-level:0
\[
(3):\ \mathrm{Def\text{-}BM.12}(6)\ \text{substituted};
\]

%dsp-level:0
\[
\qquad \Delta x_{\mathrm B}=\mathfrak h/(6\ln2)\ [\mathrm{BM.26}\ \text{re-cited}]:\ \tfrac{c\,e^2}{4\pi}\cdot\tfrac{6\ln2}{\mathfrak h}\ [\text{machine: one monomial}];
\]

%dsp-level:0
\[
\qquad \text{the exponents}\ [\mathrm{CT2.1};\ \text{machine}]
\]

%dsp-level:0
\[
(4):\ \text{the difference of}\ (3)\text{'s rows at}\ q^2\in\{1,0\}\ [\text{machine}];
\]

%dsp-level:0
\[
\qquad c\in\{\tfrac12,\tfrac35\}>0,\ e^2>0\ [\mathrm{EW.33}:\ g,g'>0]\implies\Delta m_{\mathrm{EM}}>0\ [\text{machine: coefficient census}]
\]

%dsp-level:0
\[
(5)\text{--}(6):\ \text{channel and status records}\ [\mathrm{CT1.22};\ \mathrm{D1.5};\ \mathrm{BM.37}(2)\ \text{re-cited}].\ \blacksquare
\]

\textbf{Phys.} The scar is no longer only a label. A unit charge smeared inside its own Born width carries a field, the field carries the integral of a square, and the integral is exact: one eighth over the width on the shell convention, three twentieth on the ball, one symbolic coupling factor on both. The correction to the rest mass is that energy and nothing else — it rides the transit channel the stress-energy theorem already split out, so the read stays one read, refined. The neutron, net-neutral, gets exactly zero at the monopole order, and the tier records honestly what zero does not mean: the neutral class is full of separated sign patterns, but signs without positions carry no multipole, and the chart resolves no internal structure at this order. What falls out is the first derived splitting of the two baryon rows — strictly positive, charge-selected, width-shared — and its sign is the whole point: electromagnetism makes the charged member heavier. The degeneracy of the crystallized count is broken, and broken in the direction that shuts the decay tighter, not looser.

\end{theorem}

\begin{theorem}[{\texorpdfstring{BM.47 — The infinite-range seal: UV/IR finiteness, the universe-scale neutrality, and the seal asymmetry}{BM.47 - The infinite-range seal: UV/IR finiteness, the universe-scale neutrality, and the seal asymmetry}}]\label{tag:BM-47}
\textbf{Thm.}

%dsp-level:0
\[
(1)\quad\textbf{masslessness is the range:}
\]

%dsp-level:0
\[
\qquad m_\gamma^2=0\ [\mathrm{EM.6};\ \mathrm{EW.22};\ \mathrm{EM.25}(2)]:
\]

%dsp-level:0
\[
\qquad \text{the static potential is}\ \tfrac1r\ \text{and the field}\ \tfrac1{r^2},
\]

%dsp-level:0
\[
\ \text{with no exponential damping factor}\ [\text{machine: the Yukawa factor}\ e^{-mr}\ \text{is}\ 1\ \text{at}\ m=0];
\]

%dsp-level:0
\[
\text{contrast class: colour transit modes are likewise massless}\ [\mathrm{EW.26}]\ \text{but colour is locally projected}\ [(5)]\ —\ \text{range and sealing are independent rows};
\]

%dsp-level:0
\[
(2)\quad\textbf{UV finiteness at the Born width:}
\]

%dsp-level:0
\[
\qquad \int_R^\Lambda\frac{dr}{r^2}=\frac1R-\frac1\Lambda<\frac1R\ \ \forall\Lambda>R>0
\]

%dsp-level:0
\[
\qquad [\text{machine: exact rational functions}];
\]

%dsp-level:0
\[
\Delta x_{\mathrm B}>0\ \text{strictly}\ [\mathrm{BM.13}(1);\ \mathrm{Def\text{-}BM.4}(1)]:
\]

%dsp-level:0
\[
\qquad E_{\mathrm{EM}}^{\mathrm{self}}<\infty\ \ \text{for every admissible cell}\ [\mathrm{BM.13}(3):\ E_{\mathrm{int}}\ge2];
\]

%dsp-level:0
\[
\qquad \text{the point-charge divergence is cut by the envelope,}\ \neg\text{by a counterterm};
\]

%dsp-level:0
\[
(3)\quad\textbf{IR finiteness of an isolated member:}
\]

%dsp-level:0
\[
\qquad E_{\mathrm{out}}(R)=\frac{e^2q^2}{8\pi R}<\infty\ \ \text{as}\ \ \Lambda\to\infty
\]

%dsp-level:0
\[
\qquad [\text{machine: the}\ 1/\Lambda\ \text{row vanishes in the limit}]:
\]

%dsp-level:0
\[
\text{the field energy outside any positive radius is finite although the field itself never vanishes}\ —\ \textbf{the infinite range costs no divergence at fixed net charge};
\]

%dsp-level:0
\[
(4)\quad\textbf{the universe-scale seal:}
\]

%dsp-level:0
\[
\qquad \textbf{Gauss:}\ \ \oint_{S^2_\infty}\vec E\cdot d\vec A=e\,Q_{\mathrm{enc}}
\]

%dsp-level:0
\[
\qquad [\text{machine: the flux normalization row}];
\]

%dsp-level:0
\[
\qquad \textbf{EM.11:}\ \ Q_{\mathrm{EM}}=0\ \text{in every valid universe}\ [\mathrm{D2.11};\ \mathrm{D1.3}]:
\]

%dsp-level:0
\[
\text{a sphere exhausting the universe encloses}\ Q_{\mathrm{enc}}
\]

%dsp-level:0
\[
= 0\implies\textbf{the total asymptotic flux vanishes}:\ \text{the leading}\ \tfrac1{r^2}\ \text{term of the total field cancels at universe scale}\ [\text{machine: the integer row}\ e\cdot0=0];
\]

%dsp-level:0
\[
\text{the residual (multipole) field decays as}\ r^{-3}:
\]

%dsp-level:0
\[
\qquad \int^\infty\Bigl(\frac1{r^3}\Bigr)^2r^2\,dr=\int^\infty\frac{dr}{r^4}<\infty
\]

%dsp-level:0
\[
\qquad [\text{machine: exponent arithmetic}]:
\]

%dsp-level:0
\[
\boxed{\ \textbf{the total EM field energy of a valid universe is IR-finite by}\ \mathrm{EM.11}\text{'s global neutrality}\ }\ —\ \textbf{EM seals at universe scale};
\]

%dsp-level:0
\[
(5)\quad\textbf{the seal asymmetry (colour local, EM global):}
\]

%dsp-level:0
\[
\qquad \textbf{colour:}\ \ \mathcal C(\Sigma;C)\in\{0,1\}\ \text{and one dangling attachment zeroes the whole colour sector}\ [\mathrm{D2.14};\ \mathrm{FT.19};\ \mathrm{BM.1};\ \mathrm{BM.32}(1):\ \text{filtration}]:\ \text{a local}\ 0/1\ \text{projector};
\]

%dsp-level:0
\[
\textbf{EM:}\ \ W(\pm1)=15>0\ [\mathrm{BM.28}(2);\ \mathrm{Def\text{-}BM.8}(5):\ \text{the}\ W(Q)\ \text{census}]:
\]

%dsp-level:0
\[
\qquad \text{local charge classes are admissible};
\]

%dsp-level:0
\[
\qquad \neg\exists\ \text{a local projector on}\ Q;
\]

%dsp-level:0
\[
\qquad \text{neutrality is exhausted only at universe scale}\ [\mathrm{EM.11}];
\]

%dsp-level:0
\[
[\text{machine: the two vocabularies}\ —\ \text{the idempotent}\ \{0,1\}\ \text{against the positive census}];
\]

%dsp-level:0
\[
\qquad \textbf{colour self-seals; EM does not};
\]

%dsp-level:0
\[
(6)\quad\textbf{the nearby-space fractions:}
\]

%dsp-level:0
\[
\qquad \text{the fraction of the self-energy inside}\ k\,\Delta x_{\mathrm B}:
\]

%dsp-level:0
\[
\qquad f(k)=1-\frac1k
\]

%dsp-level:0
\[
\qquad [\text{machine: exact rational functions; the shell route}];
\]

%dsp-level:0
\[
f(2)=\tfrac12,\quad f(3)=\tfrac23,\quad f(4)=\tfrac34;
\]

%dsp-level:0
\[
\qquad \text{the exterior fraction}\ \tfrac1k>0\ \ \forall k<\infty:
\]

%dsp-level:0
\[
\qquad \textbf{no finite neighbourhood holds the whole correction}\ —\ \text{the nearby-space reading is a strict fraction at every finite radius};
\]

%dsp-level:0
\[
(7)\quad\textbf{status:}
\]

%dsp-level:0
\[
\qquad \text{the correction is local in its ultraviolet}\ ((2)),
\]

%dsp-level:0
\[
\ \text{finite in its infrared per member}\ ((3))\ \text{and per universe}\ ((4));
\]

%dsp-level:0
\[
\qquad \text{the seal asymmetry}\ ((5))\ \text{is the exact content of ``colour seals, EM seals at universe scale''};
\]

%dsp-level:0
\[
\neg\text{a measured range, radius or flux enters}\ [\mathrm{P5}].
\]

\textbf{Rmk (two seals, two grammars.)} Confinement and neutrality are both sealing laws, and the tier keeps their grammar apart. Colour's is a projector: it acts on any restriction, returns zero or one, and annihilates a block at the first dangling attachment — locality is its content. Electromagnetism's is a census: every local charge class has positive weight, nothing is projected away nearby, and the vanishing happens once, globally, because a valid universe has no dangling ends at all. The infinite range is not an obstacle to either finiteness row: the Born envelope absorbs the ultraviolet, and EM.11's integer law absorbs the infrared.

\textbf{Req.} D1.3, D2.11, D2.14 (Parts 1–2); CT3.18, EM.6, EM.11, EM.25(2), EW.22, EW.26 (Part 2); P5 (Parts 1–2); FT.19 (Part 3); Def-BM.4, Def-BM.8, Def-BM.12, BM.1, BM.13(1), BM.13(3), BM.26, BM.28(2), BM.32(1), BM.46 (this part); exact radial integrals and exponent arithmetic, the Gauss flux normalization, integer censuses [math; machine: \texttt{checks/\allowbreak{}verify\_\allowbreak{}H\_\allowbreak{}em\_\allowbreak{}self\_\allowbreak{}interaction.py} §B].

\textbf{Pf.}

%dsp-level:0
\[
(1):\ \mathrm{EM.6}/\mathrm{EW.22}\ \text{re-cited};\qquad e^{-m r}\big|_{m=0}=1\ [\text{machine: exact at}\ m=0]
\]

%dsp-level:0
\[
(2):\ \tfrac{d}{dr}\bigl(-\tfrac1r\bigr)
\]

%dsp-level:0
\[
= \tfrac1{r^2}\ [\text{machine: symbolic differentiation}]\implies\int_R^\Lambda
\]

%dsp-level:0
\[
= \tfrac1R-\tfrac1\Lambda\ [\text{machine}];
\]

%dsp-level:0
\[
\qquad \Delta x_{\mathrm B}=\mathfrak h/(E\ln2)>0\ [\mathrm{BM.13}(1)\ \text{re-cited}]
\]

%dsp-level:0
\[
(3):\ \text{the}\ \Lambda\text{-row of}\ (2)\ \text{at}\ R=\Delta x_{\mathrm B}\ [\text{machine}]
\]

%dsp-level:0
\[
(4):\ \text{the flux of}\ \tfrac{eq}{4\pi r^2}\ \text{through}\ S^2_r\ \text{is}\ eq\ [\mathrm{Def\text{-}BM.12}(3)\text{'s Gauss normalization re-cited};\ \text{machine}];
\]

%dsp-level:0
\[
\qquad Q_{\mathrm{enc}}=Q_{\mathrm{EM}}=0\ [\mathrm{EM.11}\ \text{re-cited};\ \text{machine: integers}];
\]

%dsp-level:0
\[
\qquad \int^\infty r^{-4}dr\ \text{convergent}\ [\text{machine: antiderivative}\ -\tfrac1{3r^3}]
\]

%dsp-level:0
\[
(5):\ \mathcal C^2=\mathcal C\in\{0,1\}\ [\mathrm{BM.32}(1)\ \text{re-cited}];
\]

%dsp-level:0
\[
\qquad W(\pm1)=15\ [\mathrm{BM.28}(2)\ \text{re-enumerated};\ \text{machine}];
\]

%dsp-level:0
\[
\qquad \neg\exists\ \text{a local}\ Q\text{-projector: inventory scan of}\ \mathrm{D2.14}\text{'s vocabulary (colour slots only)}\ [\text{machine: token row}]
\]

%dsp-level:0
\[
(6):\ f(k)=\tfrac{E(<k\Delta x_{\mathrm B})}{E_{\mathrm{tot}}}=\tfrac{1/R-1/(kR)}{1/R}=1-\tfrac1k\ [\text{machine: rational functions}];
\]

%dsp-level:0
\[
\qquad (7):\ \text{the status record}.\ \blacksquare
\]

\textbf{Phys.} The electromagnetic correction lives with an infinite-range field and stays finite at both ends, for two different reasons that the tier separates. In the ultraviolet the Born envelope does the work: the composite is never a point, its width is strictly positive on every admissible cell, and the radial integral of the field square converges to an exact rational over that width. In the infrared the universe does the work: Gauss ties the asymptotic flux to the enclosed charge, and a valid universe encloses zero — so the total field's leading term cancels at scale, the residual multipoles fall as the cube, and the energy beyond any sphere is finite. Between the two sits the asymmetry with confinement: colour has a local zero-one projector that kills a block at one dangling attachment, while electromagnetism has no local projector at all — its charge classes are positively populated nearby and its neutrality is a single global fact about ends. And the nearby space holds only a fraction: half the field energy lives beyond two widths, two thirds within three, and a strict $1/k$ always escapes — the correction is local in origin and never local in full.

\end{theorem}

\begin{theorem}[{\texorpdfstring{BM.48 — The two-body potential, the hydrogenic rows, and the binding category bar}{BM.48 - The two-body potential, the hydrogenic rows, and the binding category bar}}]\label{tag:BM-48}
\textbf{Thm.}

%dsp-level:0
\[
(1)\quad\textbf{the static two-body potential:}
\]

%dsp-level:0
\[
\qquad V(d)=\frac{e^2q_1q_2}{4\pi\,d}
\]

%dsp-level:0
\[
\qquad [\mathrm{Def\text{-}BM.12}(3)\text{'s Coulomb read};\ \text{math: the unique spherically symmetric harmonic function with the prescribed Gauss flux}];
\]

%dsp-level:0
\[
q_1q_2<0\implies V<0\ (\text{attraction}),\qquad q_1q_2>0\implies V>0\ (\text{repulsion})\qquad[\text{machine: sign rows}];
\]

%dsp-level:0
\[
\text{the pair}\ (q_p,q_\ell)=(+1,-1)\ [\mathrm{BM.23}(4)\text{'s forced carrier charge}]:
\]

%dsp-level:0
\[
\qquad V(d)=-\frac{e^2}{4\pi\,d}=-\frac{\varepsilon_{\mathrm{EM}}}{d}
\]

%dsp-level:0
\[
\qquad [\text{machine: exact coefficient}];
\]

%dsp-level:0
\[
(2)\quad\textbf{the hydrogenic rows (chart-side, conditional on a bound channel):}
\]

%dsp-level:0
\[
\qquad \text{under the canonical-pair mechanics of the chart}\ [\mathrm{BM.10}(3):\ [\hat x,\hat p]=i,\ \hbar=1]\ \text{with reduced mass}\ m_r:
\]

%dsp-level:0
\[
E_n=-\frac{e^4(q_1q_2)^2\,m_r}{32\pi^2\,n^2}=-\frac{\varepsilon_{\mathrm{EM}}^2\,m_r\,(q_1q_2)^2}{2\,n^2},
\]

%dsp-level:0
\[
\qquad \frac{E_2}{E_1}=\frac14\quad\text{exact}
\]

%dsp-level:0
\[
\qquad [\text{math: the Coulomb spectrum};\ \text{machine: the coefficient and ladder rows}];
\]

%dsp-level:0
\[
\text{the}\ n\text{-ladder is a chart-side spectrum};
\]

%dsp-level:0
\[
\qquad m_r\ \text{of a lepton--baryon pair requires the lepton mass read}\ [\mathrm{BM.43}(2):\ \text{symbolic};\ \text{mass-form bridge residue}];
\]

%dsp-level:0
\[
(3)\quad\textbf{the bound channel is}\ \neg\textbf{constructed:}
\]

%dsp-level:0
\[
\qquad \mathrm{Def\text{-}BM.1}\text{'s composite grammar is colour-sealed}\ (\varepsilon/\delta\ \text{junctions on slot indices}\ [\mathrm{D3.9};\ \mathrm{CT3.17}]);
\]

%dsp-level:0
\[
\qquad \neg\exists\ \text{an EM sealing grammar in the inventory};
\]

%dsp-level:0
\[
\text{a lepton--baryon bound state would be a dynamics-of-}Q\ \text{object}\ [\mathrm{BM.6}(5):\ \neg\text{constructed}];
\]

%dsp-level:0
\[
\qquad \text{the photon-exchange binding of}\ \mathrm{BM.41}(3)\text{'s transit alternative remains}\ \neg\text{constructed}\ [\text{deficit}\ 1'\ \text{untouched}];
\]

$(4)\quad\textbf{the binding category bar:}\qquad\textbf{a binding energy is the energy of the bound system,}\ \neg\textbf{a correction to a free member's mass}:$

%dsp-level:0
\[
\text{for a constructed bound pair the total read would be}\ \ m_{\mathrm{pair}}=m_p+m_\ell+E_{\mathrm{bind}},
\]

%dsp-level:0
\[
\quad E_{\mathrm{bind}}<0\ [\text{system read};\ (2)\text{'s}\ E_1\ \text{at}\ n=1];
\]

%dsp-level:0
\[
\text{the free proton's correction is}\ +E_{\mathrm{EM}}^{\mathrm{self}}\ \text{alone}\ [\mathrm{BM.46}(3)]:
\]

%dsp-level:0
\[
\qquad \boxed{\ m_{\mathrm{eff}}^{(p)}=m_{\mathrm{hid}}+E_{\mathrm{EM}}^{\mathrm{self}}-E_{\mathrm{EM}}^{\mathrm{bind}}\ \ \text{is refuted as a category error}\ }\ —\ \neg\exists\ \text{an object in the inventory whose mass that expression reads};
\]

%dsp-level:0
\[
[\text{machine: the inventory row}\ —\ \text{no bound-pair mass read exists};\ (3)\ \text{re-cited}];
\]

%dsp-level:0
\[
(5)\quad\textbf{the sign rows survive the hypothetical:}
\]

%dsp-level:0
\[
\qquad E_{\mathrm{EM}}^{\mathrm{self}}>0\ \text{unconditionally}\ [\mathrm{BM.46}(4)];
\]

%dsp-level:0
\[
\qquad E_{\mathrm{bind}}<0\ \text{belongs to the pair}\ [(4)]:
\]

%dsp-level:0
\[
\text{even granting the hypothetical pair, the free-proton splitting}\ \Delta m_{\mathrm{EM}}>0\ \text{is unchanged}\ [\mathrm{BM.46}(4)\ \text{stands}];
\]

%dsp-level:0
\[
(6)\quad\textbf{status:}
\]

%dsp-level:0
\[
\qquad \text{derived: the potential row, its sign structure, the hydrogenic coefficients as chart-side math};
\]

%dsp-level:0
\[
\qquad \neg\text{constructed: the bound channel, the pair mass read};
\]

%dsp-level:0
\[
\qquad \neg\text{adopted: any}\ E_{\mathrm{bind}}\ \text{term in a free-member mass};
\]

%dsp-level:0
\[
\neg\text{a measured binding energy, ionization scale or atomic datum enters}\ [\mathrm{P5}].
\]

\textbf{Rmk (what the potential row buys.)} The Coulomb read between the forced carrier and the daughter is derived the moment the sourced field exists: opposite ternary charges attract with an exact symbolic coefficient, and the chart-side hydrogenic ladder follows from the canonical pair the Noether tier already carried. What does not exist is a grammar that seals the pair — the corpus's composites are colour-sealed, and an electromagnetic atom is a dynamics-of-$Q$ construction nobody has built. The tier draws the potential, prices the ladder, and bars the category slip: binding energy is a property of the bound, and until a bound is constructed, no free member's mass carries one.

\textbf{Req.} D1.5, D3.9 (Parts 1–2); CT2.1, CT3.17 (Part 2); P5 (Parts 1–2); Def-BM.1, Def-BM.12, BM.6(5), BM.10(3), BM.23(4), BM.41(3), BM.43(2), BM.46(3), BM.46(4) (this part); the uniqueness of the harmonic spherically symmetric potential with prescribed flux, the Coulomb spectrum on the canonical pair [math; machine: \texttt{checks/\allowbreak{}verify\_\allowbreak{}H\_\allowbreak{}em\_\allowbreak{}self\_\allowbreak{}interaction.py} §C].

\textbf{Pf.}

%dsp-level:0
\[
(1):\ \nabla^2\tfrac1r=0\ \text{at}\ r\ne0\ [\mathrm{BM.43}(4)\text{'s rule-ring identity re-cited};\ \text{machine}];
\]

%dsp-level:0
\[
\qquad \text{the flux normalization}\ [\mathrm{Def\text{-}BM.12}(3)];
\]

%dsp-level:0
\[
\qquad \text{uniqueness up to an additive constant, fixed at}\ V(\infty)=0\ [\text{math}];
\]

%dsp-level:0
\[
\qquad \text{signs}\ [\text{machine}]
\]

%dsp-level:0
\[
(2):\ \text{the Coulomb ladder}\ [\text{math: standard spectrum on}\ [\hat x,\hat p]=i];
\]

%dsp-level:0
\[
\qquad e^4/(32\pi^2)=\varepsilon_{\mathrm{EM}}^2/2\ [\text{machine: coefficient identity}];
\]

%dsp-level:0
\[
\qquad E_2/E_1=\tfrac14\ [\text{machine}]
\]

%dsp-level:0
\[
(3):\ \mathrm{Def\text{-}BM.1}/\mathrm{D3.9}/\mathrm{CT3.17}\ \text{re-cited};
\]

%dsp-level:0
\[
\qquad \text{inventory scan: no EM sealing grammar, no pair bound state}\ [\text{machine: token row}]
\]

%dsp-level:0
\[
(4):\ \text{the category row}\ [(3)\ \text{conjoined with}\ \mathrm{BM.46}(3)];
\]

%dsp-level:0
\[
\qquad (5):\ \text{the sign rows re-cited};
\]

%dsp-level:0
\[
\qquad (6):\ \text{the status record}.\ \blacksquare
\]

\textbf{Phys.} With the sourced field in hand, the force between the decay's charged members writes itself: an exact symbolic coefficient over the separation, attractive precisely where the ternary charges oppose — and the chart's own quantum mechanics turns that potential into the hydrogenic ladder, whose ground coefficient and quarter-ratio are arithmetic. But a potential is not a particle pair. The corpus seals composites with colour, and only with colour: there is no electromagnetic junction grammar, no dynamics of the additive read, and therefore no atom in the inventory. The tier states the consequence as a bar rather than a shrug: binding energy belongs to the bound system, and an expression that subtracts a would-be atomic binding from a free proton's mass reads nothing that exists. The free proton's electromagnetic correction is its own field's energy, positive, exact, and alone.

\end{theorem}

\begin{theorem}[{\texorpdfstring{BM.49 — The beta-decay re-adjudication with the correction, and the endpoint-to-average refutation}{BM.49 - The beta-decay re-adjudication with the correction, and the endpoint-to-average refutation}}]\label{tag:BM-49}
\textbf{Thm.}

%dsp-level:0
\[
(1)\quad\textbf{the corrected Q-value:}
\]

%dsp-level:0
\[
\qquad Q_{\mathrm{corr}}:=m_{\mathrm{eff}}^{(n)}-m_{\mathrm{eff}}^{(p)}-m_\ell-m_\nu=-\bigl(\Delta m_{\mathrm{EM}}+m_\ell+m_\nu\bigr)
\]

%dsp-level:0
\[
\qquad [\mathrm{BM.46}(3)\text{--}(4);\ \text{machine: monomial identity}];
\]

%dsp-level:0
\[
\Delta m_{\mathrm{EM}}>0\ [\mathrm{BM.46}(4)],\quad m_\ell>0\ [\mathrm{BM.39}(1);\ \mathrm{BM.41}(4)(\mathrm a)],\quad m_\nu\ge0\ [\mathrm{EW.33}\ \mathrm{PSD}]:
\]

%dsp-level:0
\[
\boxed{\ Q_{\mathrm{corr}}<0\ \ \text{strictly, and strictly more negative than}\ Q=-(m_\ell+m_\nu)\ \text{of}\ \mathrm{BM.41}(4)\ }
\]

%dsp-level:0
\[
(2)\quad\textbf{the no-opening theorem:}
\]

%dsp-level:0
\[
\qquad \text{every EM correction of this tier is}\ \propto q^2
\]

%dsp-level:0
\[
\ge 0\ \text{and adds to the charged member's mass}\ [\mathrm{BM.46}(1),(4);\ \text{the field energy is the integral of a square}]:
\]

%dsp-level:0
\[
\text{the decay}\ n\to p+\ell+\nu\ \text{requires}\ m^{(n)}>m^{(p)}+m_\ell:
\]

%dsp-level:0
\[
\qquad \textbf{no EM correction of this tier moves the difference in the required direction}\ —\ \text{it moves it against};
\]

%dsp-level:0
\[
\text{the claim that the EM self-interaction makes the beta decay energetically possible is}\ \textbf{refuted by the exact sign arithmetic}\ [\text{machine}];
\]

%dsp-level:0
\[
\mathrm{BM.41}(4)(\mathrm e)\text{'s verdict is sharpened, not altered:}
\]

%dsp-level:0
\[
\qquad \text{opening the vertex remains a read-specification act}\ [\mathrm{BM.30}(5);\ \text{a §I addendum}]\ \text{or an inventory datum the corpus does not carry}\ (\text{a constituent-level transport splitting of the two baryon rows}\ —\ \neg\exists\ \text{in Parts 1--4}\ [\mathrm{BM.28}(1);\ \mathrm{BM.28}(7)]);
\]

%dsp-level:0
\[
\ \text{this tier performs neither};
\]

%dsp-level:0
\[
(3)\quad\textbf{the conditional spectrum (under the}\ \neg\textbf{adopted hypothesis}\ Q\ge0\ \textbf{of some specification):}
\]

%dsp-level:0
\[
\qquad E_{\max}=E_\ell^{\max}\ [\mathrm{BM.41}(5)\ \text{endpoint formula re-cited}];
\]

%dsp-level:0
\[
\text{the allowed-shape spectrum at constant amplitude with}\ m_\nu=0,\ m_\ell\to0:
\]

%dsp-level:0
\[
\qquad \frac{d\Gamma}{dE}\propto E^2\,(E_{\max}-E)^2
\]

%dsp-level:0
\[
\qquad [\text{math: the three-body phase space}];
\]

%dsp-level:0
\[
\frac{\langle E\rangle}{E_{\max}}
\]

%dsp-level:0
\[
= \frac{\int_0^1x^3(1-x)^2\,dx}{\int_0^1x^2(1-x)^2\,dx}
\]

%dsp-level:0
\[
= \frac{\mathrm B(4,3)}{\mathrm B(3,3)}
\]

%dsp-level:0
\[
= \frac{3!\,2!/6!}{2!\,2!/5!}
\]

%dsp-level:0
\[
= \frac{1/60}{1/30}\qquad = \frac12
\]

%dsp-level:0
\[
\qquad [\text{machine: exact factorial rationals}];
\]

%dsp-level:0
\[
\boxed{\ E_{\max}:\langle E\rangle=2:1\ }
\]

%dsp-level:0
\[
\qquad \textbf{and scale-invariant:}
\]

%dsp-level:0
\[
\quad E
\]

%dsp-level:0
\[
\mapsto QE\ \text{rescales both moments alike}\implies\text{the ratio is}\ 2\ \textbf{for every value of}\ Q>0\ [\text{machine: the ratio carries no}\ Q];
\]

%dsp-level:0
\[
(4)\quad\textbf{the}\ \tfrac43\textbf{-identification refuted:}
\]

%dsp-level:0
\[
\qquad \text{the proposal}\ E_{\max}:\langle E\rangle
\]

%dsp-level:0
\[
= \tfrac43\ \text{identified with}\ W(0)/W(+1)
\]

%dsp-level:0
\[
= \tfrac43\ [\mathrm{BM.28}(2)]\ \text{is barred three ways}:
\]

%dsp-level:0
\[
(\mathrm a)\ \textbf{arithmetic:}\quad2\ne\tfrac43\qquad[2\cdot3=6\ne4=4\cdot1;\ \text{machine: cross-multiplication in}\ \mathbb Z];
\]

%dsp-level:0
\[
(\mathrm b)\ \textbf{scale invariance:}
\]

%dsp-level:0
\[
\quad (3)\text{'s ratio is}\ 2\ \text{at every splitting scale}\ —\ \neg\exists\ \text{a value of}\ Q\ (\text{hence none of the route ratios})\ \text{that moves it to}\ \tfrac43\ \text{in the massless limit}\ [\text{machine}];
\]

%dsp-level:0
\[
(\mathrm c)\ \textbf{disjointness:}
\]

%dsp-level:0
\[
\quad W(0)/W(+1)\ \text{is a probability-space cardinal ratio}\ —\ \text{a measure functional whose defining clause carries zero momentum/metric tokens}\ [\mathrm{BM.37}(2)\ \text{machine census re-cited}];
\]

%dsp-level:0
\[
\qquad \langle E\rangle,\ E_{\max}\ \text{are chart-side spectral moments}:
\]

%dsp-level:0
\[
\text{the identification crosses}\ \mathrm{BM.37}(3)\text{'s bars}\ —\ \text{disjointness, amplitude}\ne\text{mass}\ [\mathrm{FT.79}(5)],
\]

%dsp-level:0
\[
\ \text{the read-choice row}\ [\mathrm{BM.30}(5)];
\]

%dsp-level:0
\[
(\mathrm d)\ \textbf{existence:}\quad\text{on every adopted read the decay is kinematically closed}\ [(1);\ \mathrm{BM.41}(4)(\mathrm e)]:
\]

%dsp-level:0
\[
\qquad \neg\exists\ \text{an emitted-electron spectrum whose moments could be driven};
\]

%dsp-level:0
\[
\qquad (3)\ \text{is conditional content and stays conditional};
\]

%dsp-level:0
\[
(5)\quad\textbf{status:}
\]

%dsp-level:0
\[
\qquad \text{the route rows of}\ \mathrm{BM.28}/\mathrm{BM.30}/\mathrm{BM.37}\ \text{stand untouched}\ —\ \text{the measure prices remain measure prices};
\]

%dsp-level:0
\[
\qquad \text{the derived conditional rows:}\ E_{\max}\ [\mathrm{BM.41}(5)],\ \langle E\rangle/E_{\max}=\tfrac12\ [\text{massless limit}];
\]

%dsp-level:0
\[
\neg\text{a measured endpoint, average energy, spectral shape or lifetime enters}\ [\mathrm{P5};\ \text{machine: zero external tokens}].
\]

\textbf{Rmk (the check that was asked for.)} The four-thirds was inspected where it could possibly live: as the ratio of the conditional spectrum's endpoint to its mean. The phase space answers with a clean rational — one half in the massless allowed limit, hence two-to-one, and scale-invariantly two-to-one at any $Q$ whatever, so no mass splitting of any provenance can bend it to four thirds. The four thirds remains what the Tier BM.8 adjudication made it: the ratio of twenty neutral sign-assignments to fifteen charged ones, a price of the charge filter on the ensemble measure. The two numbers are exact, and they are not about the same thing.

\textbf{Req.} P5 (Parts 1–2); D2.9 (Parts 1–2); EW.33 (Part 2); FT.79(5), FT.107(3) (Part 3); Def-BM.12, BM.28(1), BM.28(2), BM.28(7), BM.30(3), BM.30(5), BM.37(2), BM.37(3), BM.39(1), BM.41(4), BM.41(5), BM.43(2), BM.46(1), BM.46(3), BM.46(4) (this part); the three-body phase-space shape at constant amplitude, exact Beta-function rationals [math; machine: \texttt{checks/\allowbreak{}verify\_\allowbreak{}H\_\allowbreak{}em\_\allowbreak{}self\_\allowbreak{}interaction.py} §D].

\textbf{Pf.}

%dsp-level:0
\[
(1):\ m_{\mathrm{eff}}^{(n)}=m_{\mathrm{hid}},\ m_{\mathrm{eff}}^{(p)}=m_{\mathrm{hid}}+\Delta m_{\mathrm{EM}}\ [\mathrm{BM.46}(3)\text{--}(4)\ \text{re-cited}]:
\]

%dsp-level:0
\[
\quad Q_{\mathrm{corr}}=-\Delta m_{\mathrm{EM}}-m_\ell-m_\nu\ [\text{machine: monomial identity}];
\]

%dsp-level:0
\[
\qquad \text{the sign census}\ [\text{machine}]
\]

%dsp-level:0
\[
(2):\ \delta(q)=c\,\varepsilon_{\mathrm{EM}}q^2m_{\mathrm{hid}}\ge0\ [\text{machine: coefficient census}];
\]

%dsp-level:0
\[
\qquad m^{(p)}_{\mathrm{eff}}-m^{(n)}_{\mathrm{eff}}
\]

%dsp-level:0
\[
= \Delta m_{\mathrm{EM}}>0\ \text{against the required}\ m^{(n)}>m^{(p)}+m_\ell\ [\text{machine: the sign contradiction}]
\]

%dsp-level:0
\[
(3):\ \int_0^1x^3(1-x)^2dx
\]

%dsp-level:0
\[
= \mathrm B(4,3)
\]

%dsp-level:0
\[
= \tfrac1{60},
\]

%dsp-level:0
\[
\quad\int_0^1x^2(1-x)^2dx
\]

%dsp-level:0
\[
= \mathrm B(3,3)
\]

%dsp-level:0
\[
= \tfrac1{30}\ [\text{machine: factorial rationals}];
\]

%dsp-level:0
\[
\qquad \tfrac{1/60}{1/30}=\tfrac12\ [\text{machine}];
\]

%dsp-level:0
\[
\qquad E=Qx\ \text{rescaling: both moments carry}\ Q\ \text{to the same power}\ [\text{machine: the ratio is}\ Q\text{-free}]
\]

%dsp-level:0
\[
(4):\ 2\cdot3=6\ne4=4\cdot1\ [\text{machine}];
\]

%dsp-level:0
\[
\qquad (b)\text{--}(d):\ \text{the cited bars re-cited with}\ \mathrm{BM.37}(2)\text{'s census rows};
\]

%dsp-level:0
\[
\qquad (5):\ \text{the status record}.\ \blacksquare
\]

\textbf{Phys.} The correction goes into the balance and the balance closes harder. The neutron stays at six bits over $\mathfrak h$; the proton gains its field energy, strictly positive, symbolic in the coupling but positive at every value the coupling can take. The Q-value of the assembled vertex is minus the splitting minus the lepton masses — negative before the correction and more negative after. No electromagnetic row of this corpus opens the beta decay, because every row the tier derives adds mass to the charged member, and the charged member is the daughter. The spectrum question the tier was handed — could the four-thirds of the charge census be driving the endpoint-to-average ratio of the emitted electron — has an exact answer and it is no: the conditional massless spectrum averages at precisely half its endpoint, a scale-free rational that no splitting value can move, while four thirds is a count of sign-assignments in a filtered ensemble. Both numbers stand as derived; the bridge between them is barred by disjointness, by the amplitude-mass bar, and by the simple fact that on every adopted read there is no emitted electron to average.

\end{theorem}

\medskip\noindent\textit{\texorpdfstring{BM.50 — The status of the EM self-interaction correction: adjudication and scope}{BM.50 - The status of the EM self-interaction correction: adjudication and scope}.}\ Structural-constraints record: see Appendix B, Section B.1.21 of the Addendum.\par\medskip

\phantomsection
\section[Tier BM.13 — The Full-Scope Mass Splitting: the Lyapunov…]{\texorpdfstring{Tier BM.13 — The Full-Scope Mass Splitting: the Lyapunov Foundation, the Complete-Force Census, the Non-Uniform Clock, the Binding Zeros, and the Exact Ratio}{Tier BM.13 - The Full-Scope Mass Splitting: the Lyapunov Foundation, the Complete-Force Census, the Non-Uniform Clock, the Binding Zeros, and the Exact Ratio}}

\begin{theorem}[{\texorpdfstring{BM.51 — The Lyapunov foundation: the contraction law iterated, the deficit links, and the class-degenerate zero}{BM.51 - The Lyapunov foundation: the contraction law iterated, the deficit links, and the class-degenerate zero}}]\label{tag:BM-51}
\textbf{Thm.}

%dsp-level:0
\[
(1)\quad\textbf{the functional and its zero locus:}
\]

%dsp-level:0
\[
\qquad \delta\rho:=\rho-\tfrac14I,
\]

%dsp-level:0
\[
\qquad \tilde K:=-\log\rho-\mathfrak T\bigl((-\log\rho)\,\rho\bigr)I,
\]

%dsp-level:0
\[
\qquad V(\tilde K):=\mathfrak T\bigl(\tilde K^2\bigr)\ge0
\]

%dsp-level:0
\[
\qquad [\mathrm{L1}];
\]

%dsp-level:0
\[
V=0\iff\rho=\tfrac14I\ [\mathrm{L1}(1)];
\]

%dsp-level:0
\[
\qquad V\bigl(\tilde K(\mathcal T(\rho))\bigr)=V\bigl(\tilde K(\rho)\bigr)\ \ \text{for every draw}\ [\mathrm{L1}(2);\ \mathrm{A6}:\ \text{tick invariance}];
\]

%dsp-level:0
\[
\qquad V'\le V,\ \ V'=V\iff\rho=\tfrac14I\ [\mathrm{L1}(5)];
\]

%dsp-level:0
\[
(2)\quad\textbf{the contraction law:}
\]

%dsp-level:0
\[
\qquad \gamma:=1-4\varepsilon_q\in(0,1]\ [\mathrm{D0.3};\ \mathrm{D0.4};\ \mathrm{CT0.8}];
\]

%dsp-level:0
\[
\qquad V'=\Bigl(1-\frac\gamma4\Bigr)^{2}V+O(\delta\rho^3),
\]

%dsp-level:0
\[
\qquad V'-V=-\frac{\gamma(8-\gamma)}{16}\,V+O(\delta\rho^3)
\]

%dsp-level:0
\[
\qquad [\mathrm{L2}];
\]

%dsp-level:0
\[
\boxed{\ \Bigl(1-\frac\gamma4\Bigr)^{2}+\frac{\gamma(8-\gamma)}{16}=1\ }
\]

%dsp-level:0
\[
\qquad [\text{machine: exact in}\ \mathbb Q[\gamma]]
\]

%dsp-level:0
\[
\qquad \text{the contraction is strict on the whole admissible range}\ \gamma\in(0,1]\ [\mathrm{BM.8}(2);\ \text{machine}];
\]

%dsp-level:0
\[
\gamma=1\ (\text{the sharp point}\ \varepsilon_q=0\ [\mathrm{D0.4};\ \mathrm{BM.19}(3)]):\quad V'=\frac9{16}V,\quad\text{deficit}\ \frac7{16}V;
\]

%dsp-level:0
\[
\qquad \gamma\to0:\ V'=V\ [\text{the identity-tilt limit, the excluded endpoint}\ \varepsilon_q=\tfrac14\ [\mathrm{CT0.8}]];
\]

%dsp-level:0
\[
(3)\quad\textbf{the iteration (the fermion-count contraction read):}
\]

%dsp-level:0
\[
\qquad \text{one tilt step per recorded crystallization event}\ [\mathrm{L0}\ \text{grounding};\ \mathrm{L2}\ \mathrm{Phys}]:
\]

%dsp-level:0
\[
\qquad V_n=\Bigl(1-\frac\gamma4\Bigr)^{2n}V_0\ \ \text{at the leading-coefficient level, remainder}\ O(\delta\rho^3)\ \text{per step}\ [\mathrm{L2}\ \text{scope}];
\]

%dsp-level:0
\[
\textbf{the}\ E_{\mathrm{int}}\textbf{-fold row:}
\]

%dsp-level:0
\[
\qquad V_{E_{\mathrm{int}}}=\Bigl(1-\frac\gamma4\Bigr)^{2E_{\mathrm{int}}}V_0
\]

%dsp-level:0
\[
\qquad [E_{\mathrm{int}}=\#\{\text{internal directed edges}\}\ [\mathrm{D0.8}];\ \text{machine: the power identity}\ \bigl(1-\tfrac\gamma4\bigr)^2\bigl(1-\tfrac\gamma4\bigr)^{2n}=\bigl(1-\tfrac\gamma4\bigr)^{2n+2}\ \text{in}\ \mathbb Q[\gamma]];
\]

%dsp-level:0
\[
\textbf{at the sharp point}\ \gamma=1,\ E_{\mathrm{int}}=6\ [\mathrm{BM.17};\ \mathrm{BM.26};\ \text{Census FT.35}(2)]:
\]

%dsp-level:0
\[
\qquad \boxed{\ \frac{V_6}{V_0}=\Bigl(\frac9{16}\Bigr)^{6}=\frac{531441}{16777216}\ },
\]

%dsp-level:0
\[
\qquad \text{the total deficit}\ 1-\frac{531441}{16777216}=\frac{16245775}{16777216}
\]

%dsp-level:0
\[
\qquad [\text{machine: exact integer powers}];
\]

%dsp-level:0
\[
\textbf{faithfulness in}\ E_{\mathrm{int}}\textbf{:}\quad\text{at fixed}\ \gamma
\]

%dsp-level:0
\[
\in (0,1],
\]

%dsp-level:0
\[
\ n\mapsto\bigl(1-\tfrac\gamma4\bigr)^{2n}\ \text{is strictly decreasing and injective}\ [\text{machine: exact rational comparison on}\ n\in2..60]:
\]

%dsp-level:0
\[
\qquad \text{the contraction read separates edge counts exactly as}\ m_{\mathrm{hid}}\ \text{does};
\]

%dsp-level:0
\[
(4)\quad\textbf{the deficit links:}
\]

%dsp-level:0
\[
\qquad \Delta S_{\mathrm{pair}}:=2|\delta\rho|^2=\tfrac12|r|^2,
\]

%dsp-level:0
\[
\qquad S\bigl(\rho\,\|\,\tfrac14I\bigr)=\Delta S_{\mathrm{pair}}+O(\delta\rho^3)
\]

%dsp-level:0
\[
\qquad [\mathrm{L3}(1)];
\]

%dsp-level:0
\[
\tilde K=-4\,\delta\rho+O(\delta\rho^2),
\]

%dsp-level:0
\[
\qquad \boxed{\ V(\tilde K)=16\,|\delta\rho|^2+O(\delta\rho^3)=8\,\Delta S_{\mathrm{pair}}+O(\delta\rho^3)\ }
\]

%dsp-level:0
\[
\qquad [\mathrm{L3}(2);\ \text{machine: both branch witnesses}];
\]

%dsp-level:0
\[
V'-V=-\frac{\gamma(8-\gamma)}2\,\Delta S_{\mathrm{pair}}+O(\delta\rho^3),\qquad\gamma=1:\ -\frac72\Delta S_{\mathrm{pair}}\qquad[\mathrm{L3}(3)];
\]

%dsp-level:0
\[
\textbf{the branch witnesses:}\quad\delta\rho_{\mathcal V}
\]

%dsp-level:0
\[
= \Pi_{\mathcal V}-\tfrac14I:\ \operatorname{spec}
\]

%dsp-level:0
\[
= \bigl(\tfrac34,-\tfrac14,-\tfrac14,-\tfrac14\bigr):\ \ \mathfrak T(\delta\rho^2)\qquad = \tfrac34;
\]

%dsp-level:0
\[
\qquad \delta\rho_{\mathcal R}=\tfrac{\Pi_{\mathcal R}}3-\tfrac14I:\ \operatorname{spec}=\bigl(\tfrac1{12},\tfrac1{12},\tfrac1{12},-\tfrac14\bigr):\ \ \mathfrak T(\delta\rho^2)=\tfrac1{12}
\]

%dsp-level:0
\[
[\mathrm{FT.107}(2);\ \text{machine: exact Fraction matrix arithmetic}]:
\]

%dsp-level:0
\[
\qquad \text{the leading deficits}\ 2|\delta\rho|^2=\bigl(\tfrac32,\ \tfrac16\bigr),
\]

%dsp-level:0
\[
\qquad V=8\Delta S_{\mathrm{pair}}:\ \ 16\cdot\tfrac34=12=8\cdot\tfrac32,\quad16\cdot\tfrac1{12}=\tfrac43=8\cdot\tfrac16\ [\text{machine}];
\]

%dsp-level:0
\[
(5)\quad\textbf{the inertial read:}
\]

%dsp-level:0
\[
\qquad m_i^2(q):=\Delta S(q)
\]

%dsp-level:0
\[
\qquad [\mathrm{L3}(4);\ \mathrm{L4}:\ \text{the patch deficit with its integer degree factor};\ \text{dimensionless nats};\ \neg\text{a particle identification}];
\]

%dsp-level:0
\[
\textbf{the species row:}\quad\text{the conversion between the deficit/inertial species and the}\ m_{\mathrm{hid}}\ \text{species is item-7 content}:
\]

%dsp-level:0
\[
\qquad \text{the tier needs no conversion}\ —\ (6)\text{'s zero is species-independent};
\]

%dsp-level:0
\[
(6)\quad\textbf{the class-degenerate zero:}
\]

%dsp-level:0
\[
\qquad \text{every input of the contraction read is class-shared:}
\]

%dsp-level:0
\[
\quad E_{\mathrm{int}}=6\ \text{on both baryon rows}\ [\mathrm{BM.17};\ \mathrm{BM.26};\ \text{Census FT.35}(2)];
\]

%dsp-level:0
\[
\qquad \gamma=1-4\varepsilon_q\ \text{is the universal ledger threshold}\ [\mathrm{D0.4}];
\]

%dsp-level:0
\[
\qquad V_0\ \text{is the common vacuum reference}\ [\mathrm{L1}(1);\ \mathrm{CT0.2}]:
\]

%dsp-level:0
\[
\boxed{\ \Delta m_{\mathrm{Lyapunov}}=0\ }
\]

%dsp-level:0
\[
\qquad [\text{machine: the difference of identical monomials}]:
\]

%dsp-level:0
\[
\qquad \textbf{the Lyapunov foundation gives the two baryon rows the same mass}\ —\ \text{in every read species at once, because it splits no input};
\]

%dsp-level:0
\[
(7)\quad\textbf{scope:}
\]

%dsp-level:0
\[
\qquad \mathrm{L1}\text{'s faithfulness scope stands}\ (\rho\ \text{faithful};\ \gamma\in(0,1];\ \text{the identity-tilt limit is the excluded endpoint});
\]

%dsp-level:0
\[
\qquad \mathrm{L0}\text{'s crystallization--tilt coupling is the declared theorem-local dynamical scope}\ [\text{inherited at the point of use},\ \neg\text{extended}];
\]

%dsp-level:0
\[
\text{the identification ``mass}:
\]

%dsp-level:0
\[
= \text{the contraction rate'' is}\ \neg\text{adopted as a new read species:}\ (3)\text{'s}\ V_{E_{\mathrm{int}}}/V_0\ \text{is an}\ E_{\mathrm{int}}\text{-faithful dimensionless companion of}\ m_{\mathrm{hid}}\ [\mathrm{FT.44}(4)],
\]

%dsp-level:0
\[
\ \text{and}\ (6)\text{'s zero holds on it}.
\]

\textbf{Rmk (why the zero needs no bridge).} The relating equation between the deficit species (the Lyapunov-contraction read of the registry [Def-H.39(3)]) and the informational species is $\neg$derived, and this tier does not touch it — because a splitting claim would need the bridge while a degeneracy claim does not: two rows sharing every input of a functional share its value in any species, converted or not. The Lyapunov tower is the foundation the directive asks for, and what the foundation says about the neutron–proton pair is exact, trivial, and load-bearing: zero.

\textbf{Req.} A1, A4, A5, A6 (Part 1); C1, D0.3, D0.4, D0.8 (Part 1); CT0.2, CT0.8, CT1.1 (Part 1); P5 (Parts 1–2); L0, L1, L2, L3, L4 (Part 1); FT.44(4) (Census); FT.107(2) (Part 3); Def-BM.5(4), BM.8(2), BM.17, BM.19(3), BM.26, BM.28(1) (this part); exact formal-power-series identities over $\mathbb Q[\gamma]$, exact rational powers and integer power censuses, exact Fraction matrix arithmetic on the branch witnesses [math; machine: \texttt{checks/\allowbreak{}verify\_\allowbreak{}H\_\allowbreak{}full\_\allowbreak{}scope\_\allowbreak{}mass\_\allowbreak{}splitting.py} §A].

\textbf{Pf.}

%dsp-level:0
\[
(1)\text{--}(2):\ \mathrm{L1}(1),(2),(5)\ \text{and}\ \mathrm{L2}\ \text{re-cited (frozen)};
\]

%dsp-level:0
\[
\qquad \bigl(1-\tfrac\gamma4\bigr)^2+\tfrac{\gamma(8-\gamma)}{16}=1-\tfrac\gamma2+\tfrac{\gamma^2}{16}+\tfrac\gamma2-\tfrac{\gamma^2}{16}=1\ [\text{machine}];
\]

%dsp-level:0
\[
\qquad \gamma\in(0,1]\implies\bigl(1-\tfrac\gamma4\bigr)^2\in\bigl[\tfrac9{16},1\bigr)\ [\text{machine: rational witnesses}]
\]

%dsp-level:0
\[
(3):\ \text{induction on the leading factor}\ [\mathrm{L2}]:\ \text{the power identity in}\ \mathbb Q[\gamma]\ [\text{machine}];
\]

%dsp-level:0
\[
\qquad (9/16)^6=531441/16777216\ [\text{machine:}\ 9^6=531441,\ 16^6=16777216];
\]

%dsp-level:0
\[
\qquad \text{injectivity:}\ \bigl(\tfrac9{16}\bigr)^{n}>\bigl(\tfrac9{16}\bigr)^{n+1}\ \text{on}\ n\in2..60\ [\text{machine: exact Fractions}]
\]

%dsp-level:0
\[
(4):\ \Pi_{\mathcal V}=\tfrac14(E_0-\Sigma),\ \Pi_{\mathcal R}=\tfrac14(3E_0+\Sigma)\ [\mathrm{CT0.7}(1)];
\]

%dsp-level:0
\[
\qquad \Sigma
\]

%dsp-level:0
\[
= \textstyle\sum_i\sigma_i\otimes\sigma_i:\ \operatorname{spec}
\]

%dsp-level:0
\[
= \{1,1,1,-3\}\ [\text{machine}]:\quad\delta\rho_{\mathcal V}:\ \bigl(\tfrac34,-\tfrac14,-\tfrac14,-\tfrac14\bigr),
\]

%dsp-level:0
\[
\ \ \delta\rho_{\mathcal R}:\ \bigl(\tfrac1{12},\tfrac1{12},\tfrac1{12},-\tfrac14\bigr)\ [\text{machine}]
\]

%dsp-level:0
\[
\mathfrak T(\delta\rho_{\mathcal V}^2)=\tfrac9{16}+\tfrac3{16}=\tfrac34;
\]

%dsp-level:0
\[
\qquad \mathfrak T(\delta\rho_{\mathcal R}^2)=\tfrac3{144}+\tfrac1{16}=\tfrac1{12}\ [\text{machine: exact Fractions}];
\]

%dsp-level:0
\[
\qquad V=16|\delta\rho|^2=8\Delta S_{\mathrm{pair}}\ \text{at both witnesses}\ [\text{machine}]
\]

%dsp-level:0
\[
(5):\ \mathrm{L3}(4)/\mathrm{L4}\ \text{re-cited};
\]

%dsp-level:0
\[
\qquad (6):\ m_n^{(\mathrm{Lyap})}-m_p^{(\mathrm{Lyap})}=0\ \text{as the difference of identical inputs}\ [\text{machine}];
\]

%dsp-level:0
\[
\qquad (7):\ \text{the scope record}.\ \blacksquare
\]

\textbf{Phys.} The foundation is a contraction, and the contraction is exact arithmetic. Every tilt step keeps $(1-\gamma/4)^2$ of the Lyapunov value and loses $\gamma(8-\gamma)/16$ of it — the two add to one as an identity in $\gamma$, not as an approximation — and iterating over the crystallization events of a composite turns the edge count into a strictly faithful dimensionless read: at the sharp threshold, six edges contract the functional to $531441/16777216$ of its vacuum value. The deficit links tie the same object to entropy and inertia with exact rational factors, verified at both branch witnesses: sixteen of a squared deviation, eight of a pair deficit, one inertial read per channel. And then the payload for the baryon pair: every input of this tower — the edge count, the threshold, the vacuum reference — is shared by the neutron row and the proton row, so the whole Lyapunov contribution to their splitting is zero, in every read species at once. The foundation does not split the pair; the split has to come from somewhere else, and the census that follows finds exactly one somewhere.

\end{theorem}

\begin{theorem}[{\texorpdfstring{BM.52 — The complete force operator and the sector-by-sector splitting census}{BM.52 - The complete force operator and the sector-by-sector splitting census}}]\label{tag:BM-52}
\textbf{Thm.}

%dsp-level:0
\[
(1)\quad\textbf{the complete force operator:}
\]

%dsp-level:0
\[
\qquad \mathbf F
\]

%dsp-level:0
\[
= F_{(+,\mathcal R)}\oplus F_{(-,\mathcal R)}\oplus F_{(+,\mathcal V)}\oplus F_{(-,\mathcal V)}\oplus F_{\mathrm{st}}\ \ \text{always applies in full}\ —\ \text{no sub-block restriction is an A8 force}\ [\mathrm{A8};\ \mathrm{FT.142}(2)];
\]

%dsp-level:0
\[
\text{the block/aspect dictionary:}
\]

%dsp-level:0
\[
\quad (G,H,D^+)\leftrightarrow\mathcal E_{++}\ (7),
\]

%dsp-level:0
\[
\quad W\leftrightarrow\mathcal E_{+-}\ (3),
\]

%dsp-level:0
\[
\quad J^+\leftrightarrow\mathcal E_{-+}\ (3),
\]

%dsp-level:0
\[
\quad J^-\leftrightarrow\mathcal E_{--}\ (3):
\]

%dsp-level:0
\[
\qquad 7+3+3+3=16=C_{\mathrm{cap}}\ [\mathrm{CT1.1};\ \mathrm{CT0.12};\ \mathrm{CT1.12a}(2);\ \mathrm{FT.30}];
\]

%dsp-level:0
\[
\text{at the sealed composites the rung-3 mask is active:}\quad\dim\mathcal K_3
\]

%dsp-level:0
\[
= 25
\]

%dsp-level:0
\[
= 16+1+8\ [\mathrm{CT0.13};\ \mathrm{CT0.14};\ \mathrm{Def\text{-}BM.3}];
\]

%dsp-level:0
\[
\qquad \text{the force enters the stress tensor through the masked resolvent}\ Z_{\mathrm{force}}
\]

%dsp-level:0
\[
= \langle\mathrm{out}|(1-\mathbf F_M^{(n)})^{-1}|\mathrm{in}\rangle\ [\mathrm{CT1.22};\ \mathrm{CT2.2}];
\]

%dsp-level:0
\[
(2)\quad\textbf{the carrier-side static matrix:}
\]

%dsp-level:0
\[
\qquad \mathbf F_{\mathrm{ext}}(\mathbf c,q)=h(\mathbf c)\otimes I_2+q\,I_2\otimes\hat\kappa,
\]

%dsp-level:0
\[
\qquad F_{LL}=F_{RR}=h(\mathbf c),
\]

%dsp-level:0
\[
\qquad F_{LR}=F_{RL}=q\,I_2
\]

%dsp-level:0
\[
\qquad [\mathrm{BM.22}(1)\text{--}(2);\ \mathrm{FT.129}(4)\text{--}(5);\ \mathrm{Def\text{-}H.43}\ \text{promotion};\ \text{machine}];
\]

%dsp-level:0
\[
\textbf{the chiral-partner splitting:}\quad2|q|\in\{0,2\}\ \ \text{at}\ q\in\{-1,0,+1\}
\]

%dsp-level:0
\[
\qquad [\mathrm{BM.22}(4);\ \mathrm{FT.130}(5);\ \text{machine}]:
\]

%dsp-level:0
\[
\qquad \text{the neutral class keeps parity doubling}\ (F_{LR}=0),\ \text{the unit class splits by two read units};
\]

%dsp-level:0
\[
\textbf{the EM difference between the two classes is exactly the flipping block}\ [\mathrm{BM.22}(4)];
\]

%dsp-level:0
\[
\qquad \text{the splitting is a dimensionless force-spectrum read}\ —\ \text{its frequency conversion is the deferred read}\ [\mathrm{FT.131}(3);\ \mathrm{FT.159}\ \text{residue}];
\]

%dsp-level:0
\[
(3)\quad\textbf{the sector census (the splitting contributions between the two baryon rows):}\qquad\begin{array}{l|l|l}\text{sector} & \text{contribution} & \text{the derived ground}\\\hline J^\pm\ (\text{seals}) & 0 & \Delta S=\tfrac32(2E-|R|)\ \text{is}\ \sigma\text{-free}\ [\mathrm{FT.107}(2)];\ E_{\mathrm{int}}\ \text{equal};\ \text{the Janus balance holds on both rows}\ [\mathrm{BM.7}]\\G,H\ (\text{metric blocks}) & 0 & \text{the seven geometric monomials carry charge exponent}\ 0\ [\mathrm{BM.28}(1)]\\W\ (\text{weak/shift}) & 0 & |w|^2\ \text{identical between the classes}\ [\mathrm{BM.28}(6)];\ \text{the Casimir is branch data, neutral at the}\ \sigma\text{-vertex}\ [\mathrm{BM.40}(3);\ \mathrm{FT.157}(1)]\\\text{colour} & 0 & \text{sealed ends are blind}\ [\mathrm{FT.117}];\ \text{the mask admits only singlets}\ [\mathrm{CT3.18};\ \mathrm{D2.14}];\ \text{both rows are}\ \varepsilon\text{-sealed with}\ g(3)=1\ [\mathrm{BM.31}(2)];\ \text{colour acts on the slot carrier}\ [\mathrm{FT.127}(7)];\ \text{the seed is colour-neutral}\ [\mathrm{EW.32}]\\D^\pm\ (\text{dark}) & 0 & D^+\ \text{is uniform and carried by}\ \Lambda\ [\mathrm{CT1.22};\ \mathrm{CT4.1}];\ \text{both locks are shared by all charge classes}\ [\mathrm{BM.28}(1);\ \mathrm{CT4.5}]\\\mathrm{EM} & \textbf{the only nonzero sector} & \text{the flipping block}\ 2|q|\in\{0,2\}\ [(2)];\ \text{the field energy}\ \Delta m_{\mathrm{EM}}=\tfrac{3ce^2\ln2}{2\pi\mathfrak h}>0\ [\mathrm{BM.46}(4)]\end{array}
\]

%dsp-level:0
\[
(4)\quad\textbf{the totality theorem:}
\]

%dsp-level:0
\[
\qquad \boxed{\ \Delta m_{\mathrm{force}}:=m_n-m_p\Big|_{\text{the complete operator}}=-\,\Delta m_{\mathrm{EM}}=-\frac{3\,c\,e^2\ln2}{2\pi\,\mathfrak h}\ }
\]

%dsp-level:0
\[
\qquad [\text{machine: the census sum of five zeros and one monomial}];
\]

%dsp-level:0
\[
\textbf{the mass splitting comes from the total force}\ —\ \text{and the total force, censused sector by sector, reduces to its EM summand}:
\]

%dsp-level:0
\[
(5)\quad\textbf{the agreement row (no double count):}
\]

%dsp-level:0
\[
\qquad \text{the sector census}\ ((3))\ \text{and the field integral}\ [\mathrm{BM.46}(1)]\ \text{are two derived routes to one monomial}:
\]

%dsp-level:0
\[
\Delta m_{\mathrm{force}}\big|_{\mathrm{EM}}=\Delta m_{\mathrm{self\text{-}int}}=\Delta m_{\mathrm{EM}}
\]

%dsp-level:0
\[
\qquad [\text{machine: the monomials are identical}]:
\]

%dsp-level:0
\[
\qquad \text{in any decomposition the EM contribution is counted once}\ [\mathrm{BM.55}(1)];
\]

%dsp-level:0
\[
(6)\quad\textbf{scope:}
\]

%dsp-level:0
\[
\qquad \text{the census is between the two baryon rows at matched profiles and shared junction grammar}\ [\mathrm{Def\text{-}BM.3};\ \mathrm{BM.31}(2)];
\]

%dsp-level:0
\[
\qquad \text{the}\ \mathrm{Def\text{-}H.43}\ \text{promotion's conditionality stands as}\ \mathrm{BM.22}(5)\ \text{records it}\ (\neg\text{a new declaration here});
\]

%dsp-level:0
\[
\text{constituent-level transport differences between the rows have no derived object}\ [\mathrm{BM.55}(4)\ \text{forward pointer at the citing site}];
\]

%dsp-level:0
\[
\qquad \text{the deferred frequency read}\ [\mathrm{FT.131}(3)]\ \text{blocks no energy-dimension row: the field monomial}\ (4)\ \text{is complete without it}.
\]

\textbf{Rmk (what "full scope" computes to).} The standing order is the sum of all sectors, and the sum is taken — literally, sector by sector, with each zero backed by its own frozen theorem rather than by silence. The seals cannot split the pair because the deficit never reads the end signs; the metric blocks cannot because every geometric monomial is charge-blind; the weak block cannot because the two rows share $|w|^2$ and the Casimir is branch-neutral at a sign flip; colour cannot because sealed singlets are blind to it and both rows carry the same unique grouping; the dark sector cannot because its locks are class-shared. What survives is electromagnetism, twice-derived and identical: the flipping block on the carrier spectrum and the field energy on the chart. Full scope does not dilute the answer; it isolates it.

\textbf{Req.} A8 (Part 1); CT0.12, CT0.13, CT0.14, CT1.1 (Part 1); CT1.22, CT2.2, CT3.18, CT4.1, CT4.5, D2.14 (Part 2); EM.9, EW.32 (Part 2); P5 (Parts 1–2); Def-H.43, FT.30, FT.107(2), FT.117, FT.125(4), FT.127(7), FT.129(4), FT.129(5), FT.130(5), FT.131(3), FT.157(1), FT.159 (Part 3); Def-BM.3, Def-BM.12, BM.7, BM.22(1), BM.22(2), BM.22(4), BM.22(5), BM.28(1), BM.28(6), BM.31(2), BM.40(3), BM.46(1), BM.46(4) (this part); exact integer censuses, Laurent monomials [math; machine: \texttt{checks/\allowbreak{}verify\_\allowbreak{}H\_\allowbreak{}full\_\allowbreak{}scope\_\allowbreak{}mass\_\allowbreak{}splitting.py} §B].

\textbf{Pf.}

%dsp-level:0
\[
(1):\ \mathrm{A8}\text{'s completeness row and the block dictionary re-cited}\ [\mathrm{CT1.1};\ \mathrm{CT0.12};\ \mathrm{CT1.12a}(2);\ \mathrm{FT.30}];
\]

%dsp-level:0
\[
\qquad 7+3+3+3=16,\ 16+1+8=25\ [\text{machine}]
\]

%dsp-level:0
\[
(2):\ \mathrm{BM.22}(1)\text{--}(2),(4)\ \text{re-cited};
\]

%dsp-level:0
\[
\qquad 2|q|\ \text{at}\ q\in\{-1,0,+1\}=(2,0,2)\ [\text{machine}];
\]

%dsp-level:0
\[
\qquad F_{LR}=qI_2:\ \text{the zero matrix at}\ q=0\ [\text{machine: block extraction}]
\]

%dsp-level:0
\[
(3):\ \text{each zero is the cited frozen row}: \mathrm{FT.107}(2)\text{'s}\ \sigma\text{-freeness};
\]

%dsp-level:0
\[
\ \mathrm{BM.28}(1)\text{'s monomial census};
\]

%dsp-level:0
\[
\ \mathrm{BM.28}(6)\text{'s}\ |w|^2\ \text{row and}\ \mathrm{BM.40}(3);
\]

%dsp-level:0
\[
\ \mathrm{FT.117}/\mathrm{CT3.18}/\mathrm{BM.31}(2)/\mathrm{FT.127}(7)/\mathrm{EW.32};
\]

%dsp-level:0
\[
\ \mathrm{CT1.22}/\mathrm{CT4.1}/\mathrm{CT4.5}/\mathrm{BM.28}(1)\ \text{lock row}\ [\text{machine: the census structure}]
\]

%dsp-level:0
\[
(4):\ 0+0+0+0+0-\Delta m_{\mathrm{EM}}\ [\text{machine: the monomial sum}];
\]

%dsp-level:0
\[
\qquad (5):\ \mathrm{BM.46}(1)\text{'s integral and}\ (3)\text{'s EM entry are the same Laurent monomial}\ [\text{machine}];
\]

%dsp-level:0
\[
\qquad (6):\ \text{the scope record}.\ \blacksquare
\]

\textbf{Phys.} Ask every force in the ledger what it contributes to the neutron–proton split, and five of the six answer zero — not approximately, not by assumption, but each by a theorem already on the books: the deficit ignores end signs, the geometry ignores charge, the weak Casimir ignores sign flips, colour cannot see sealed singlets, and the dark locks are shared. The complete force operator is the sum of all sectors, so the total is the EM summand, and the EM summand has already been computed twice by two routes that agree monomial for monomial. The splitting of the corpus's two baryon rows is one term: the field energy of the scar. Everything else in the full scope is bookkeeping that confirms the isolation.

\end{theorem}

\begin{theorem}[{\texorpdfstring{BM.53 — The metric background and the non-uniform clock: the zero metric contribution and the Mandelstam–Tamm resolvability bound}{BM.53 - The metric background and the non-uniform clock: the zero metric contribution and the Mandelstam-Tamm resolvability bound}}]\label{tag:BM-53}
\textbf{Thm.}

%dsp-level:0
\[
(1)\quad\textbf{the background:}
\]

%dsp-level:0
\[
\qquad \text{the metric is the emergent dictionary of the sector coordinates}\ (a_{00},q_{ij},w_i)\ [\mathrm{CT1.12};\ \mathrm{Def\text{-}BM.10}(2)];
\]

%dsp-level:0
\[
\qquad \text{the dispersion read}\ E=N\sqrt{m^2+\gamma^{ij}p_ip_j}-N^ip_i\ \text{and}\ \mathrm{KE}=N\bigl(\sqrt{m^2+p^2}-m\bigr)\ [\mathrm{BM.34}(4)\text{--}(5)];
\]

%dsp-level:0
\[
(2)\quad\textbf{the metric contribution is zero:}
\]

%dsp-level:0
\[
\qquad \text{the seven geometric monomials}\ (\Delta x_{\mathrm B},V_{\mathrm{eff}},\rho,\kappa^2\rho,r_s,\ \text{both locks})\ \text{carry charge exponent}\ 0\ [\mathrm{BM.28}(1);\ \text{machine: the exponent census re-run}];
\]

%dsp-level:0
\[
\text{the rest slice}\ E(p{=}0)
\]

%dsp-level:0
\[
= N\,
\]

%dsp-level:0
\[
m\ [\mathrm{BM.34}(7)]\ \text{has class-shared}\ N\ \text{and class-shared}\ m\ \text{before the EM correction}\ [\mathrm{BM.17};\ \mathrm{BM.26};\ \mathrm{BM.51}(6)]:
\]

%dsp-level:0
\[
\qquad \boxed{\ \Delta m_{\mathrm{metric}}=0\ }\ [\text{machine}];
\]

%dsp-level:0
\[
(3)\quad\textbf{the non-uniform clock:}
\]

%dsp-level:0
\[
\qquad \text{the tick strictly contracts}\ V\ \text{on the whole admissible range}\ [\mathrm{BM.8}(2);\ \mathrm{L2};\ \mathrm{BM.51}(2)\text{--}(3)]:
\]

%dsp-level:0
\[
\qquad V_n=V_0\bigl(1-\tfrac\gamma4\bigr)^{2n};
\]

%dsp-level:0
\[
\textbf{the log-linear row:}
\]

%dsp-level:0
\[
\qquad \ln V_n=\ln V_0+2n\,\ln\!\bigl(1-\tfrac\gamma4\bigr)
\]

%dsp-level:0
\[
\qquad [\text{math: the logarithm of the exact power identity};\ \text{structural}\ —\ \text{no transcendental is evaluated}];
\]

%dsp-level:0
\[
\textbf{no uniform-}\tau\ \textbf{exponential law:}
\]

%dsp-level:0
\[
\qquad \text{the ledger evolution is the per-tick weight and its resolvent}\ [\mathrm{BM.44}(1)\text{--}(4);\ \mathrm{D1.8}:\ x^0:=\tau_n;\ \mathrm{FT.159}]:
\]

%dsp-level:0
\[
\qquad \Gamma_{\mathrm{tr}}\in\bigl[0,\tfrac13\bigr],\quad\tau_{\mathrm{life}}\ge3\ \text{ticks},\quad\neg\text{external seconds};
\]

%dsp-level:0
\[
(4)\quad\textbf{the Mandelstam--Tamm relation:}
\]

%dsp-level:0
\[
\qquad \Delta t\,\Delta E\ge\tfrac12
\]

%dsp-level:0
\[
\qquad [\mathrm{BM.10}(4);\ \text{no time operator at finite dimension}];
\]

%dsp-level:0
\[
\textbf{the exact saturation witness re-derived:}\quad A=\tfrac12\sigma_x,\ H=\tfrac12\sigma_z,\ \psi\propto(1,i):
\]

%dsp-level:0
\[
\qquad \langle A\rangle=\langle H\rangle=0,\quad\langle\sigma_y\rangle=1,\quad(\Delta A)^2=(\Delta E)^2=\tfrac14,\quad[H,A]=\tfrac i2\sigma_y,\quad\frac{d\langle A\rangle}{dt}=-\tfrac12:
\]

%dsp-level:0
\[
\Delta t=\frac{\Delta A}{\bigl|\tfrac{d\langle A\rangle}{dt}\bigr|}=1,\qquad\Delta t\,\Delta E=\tfrac12\qquad[\text{machine: exact over}\ \mathbb Q(i)];
\]

%dsp-level:0
\[
(5)\quad\textbf{the resolvability of the splitting:}
\]

%dsp-level:0
\[
\qquad \text{resolving an energy difference}\ \Delta m_{\mathrm{EM}}\ \text{against the non-uniform clock requires}\ \ \Delta t\ \ge\ \frac1{2\,\Delta m_{\mathrm{EM}}}\ \ [\mathrm{MT},\ (4)]:
\]

%dsp-level:0
\[
\boxed{\ \Delta t\ \ge\ \frac{\pi\,\mathfrak h}{3\,c\,e^2\ln2}\ }\qquad[\text{machine: the exact reciprocal monomial of}\ \mathrm{BM.46}(4)\text{'s splitting}];
\]

%dsp-level:0
\[
(6)\quad\textbf{status:}
\]

%dsp-level:0
\[
\qquad \text{the non-uniform clock shifts no splitting value}\ —\ \text{both rows ride the same}\ \gamma,
\]

%dsp-level:0
\[
\ \text{the same tick order and the same contraction law}\ [\mathrm{BM.51}(6)]:
\]

%dsp-level:0
\[
\qquad \text{the MT relation bounds the splitting's}\ \textbf{resolvability},\ \neg\text{its magnitude};
\]

%dsp-level:0
\[
\text{the}\ \mathrm{BM.8}(2)\ \text{stationarity scope and the}\ \mathrm{CT2.6}\ \text{hypotheses are inherited at the point of use};
\]

%dsp-level:0
\[
\qquad \mathrm{FT.159}\text{'s rate-modulation residue and}\ \mathrm{FT.131}(3)\text{'s frequency read stand untouched}.
\]

\textbf{Rmk (the clock is the ledger, and the ledger is honest about it).} A uniform proper time would license an exponential law with one rate; the corpus's clock is a strictly contracting tick whose only derived rates are rational functions of one threshold. The Mandelstam–Tamm form is the admissible $\Delta t\,\Delta E$ read — saturated by an exact two-level witness — and applied to the splitting it produces the tier's one new bound: no clock coarser than $\pi\mathfrak h/(3ce^2\ln2)$ resolves the scar's field energy. The bound is symbolic in the coupling, like everything else here, and it inverts correctly: a smaller coupling means a smaller splitting means a longer required resolution.

\textbf{Req.} D1.8 (Part 1); CT1.12, CT2.1, CT2.6 (Part 2); P5 (Parts 1–2); L2 (Part 1); FT.131(3), FT.159 (Part 3); Def-BM.4, Def-BM.10, BM.8(2), BM.10(4), BM.17, BM.26, BM.28(1), BM.34(4), BM.34(5), BM.34(7), BM.44(1), BM.44(2), BM.44(3), BM.44(4), BM.46(4), BM.51(2), BM.51(3), BM.51(6) (this part); the logarithm of an exact power identity, exact $\mathbb Q(i)$ witness algebra, Laurent reciprocals [math; machine: \texttt{checks/\allowbreak{}verify\_\allowbreak{}H\_\allowbreak{}full\_\allowbreak{}scope\_\allowbreak{}mass\_\allowbreak{}splitting.py} §C].

\textbf{Pf.}

%dsp-level:0
\[
(1)\text{--}(2):\ \mathrm{BM.28}(1)\text{'s census re-run: the exponent of the charge symbol is}\ 0\ \text{in all seven monomials}\ [\text{machine}];
\]

%dsp-level:0
\[
\qquad E(p{=}0)=Nm\ [\mathrm{BM.34}(7)\ \text{re-cited}];
\]

%dsp-level:0
\[
\qquad \Delta m_{\mathrm{metric}}=0-0\ [\text{machine}]
\]

%dsp-level:0
\[
(3):\ \mathrm{BM.51}(3)\text{'s iteration re-cited};
\]

%dsp-level:0
\[
\qquad \ln(a^n)=n\ln a\ [\text{math}]\ \text{applied to the exact rational base}\ \bigl(1-\tfrac\gamma4\bigr)^2;
\]

%dsp-level:0
\[
\qquad \mathrm{BM.44}(1)\text{--}(4)\text{'s rows re-cited (frozen)}
\]

%dsp-level:0
\[
(4):\ \sigma_x(1,i)=(i,1):\ \psi^\dagger\sigma_x\psi=1\cdot i+(-i)\cdot1=0;
\]

%dsp-level:0
\[
\qquad \sigma_z(1,i)=(1,-i):\ \psi^\dagger\sigma_z\psi=1-1=0;
\]

%dsp-level:0
\[
\qquad \sigma_y(1,i)
\]

%dsp-level:0
\[
= (1,i)\cdot\ldots:\ \psi^\dagger\sigma_y\psi
\]

%dsp-level:0
\[
= 1\cdot1+(-i)(i)
\]

%dsp-level:0
\[
= 2\implies\langle\sigma_y\rangle
\]

%dsp-level:0
\[
= 1\ [\psi^\dagger\psi=2;\ \text{machine over}\ \mathbb Q(i)]
\]

%dsp-level:0
\[
(\Delta A)^2=\tfrac14\langle\sigma_x^2\rangle=\tfrac14,\quad(\Delta E)^2=\tfrac14;
\]

%dsp-level:0
\[
\qquad [H,A]=\tfrac14[\sigma_z,\sigma_x]=\tfrac i2\sigma_y\implies\tfrac{d\langle A\rangle}{dt}=i\langle[H,A]\rangle=-\tfrac12;
\]

%dsp-level:0
\[
\qquad \Delta t\,\Delta E=\tfrac{\Delta A}{1/2}\cdot\tfrac12=\tfrac12\ [\text{machine: saturation}]
\]

%dsp-level:0
\[
(5):\ \frac1{2\Delta m_{\mathrm{EM}}}
\]

%dsp-level:0
\[
= \frac1{2\cdot\frac{3ce^2\ln2}{2\pi\mathfrak h}}
\]

%dsp-level:0
\[
= \frac{\pi\mathfrak h}{3ce^2\ln2}\ [\text{machine: the Laurent reciprocal}];
\]

%dsp-level:0
\[
\qquad (6):\ \text{the scope record}.\ \blacksquare
\]

\textbf{Phys.} The metric is the background the dispersion reads ride on, and between the two baryon rows it makes no difference at all: every geometric monomial the calibration tiers locked is charge-blind, and the rest slice multiplies a shared lapse by a shared mass. The clock, meanwhile, is not uniform and is never treated as if it were — the tick contracts the Lyapunov value by an exact rational at every crystallization, the log-linear law is the logarithm of that identity and nothing more, and the only decay-style curve the corpus admits is the per-tick resolvent with its three-tick floor. The Mandelstam–Tamm relation is the correct handling of such a clock, saturated by the same two-level witness the Noether tier built, and applied to the splitting it yields the tier's sharpest new number: resolving the scar's field energy takes a clock finer than $\pi\mathfrak h$ over $3ce^2\ln2$. The bound moves exactly as the physics should — weaker coupling, smaller splitting, longer patience.

\end{theorem}

\begin{theorem}[{\texorpdfstring{BM.54 — The binding energies: the count zero, the degenerate ensemble, and the unconstructed atom}{BM.54 - The binding energies: the count zero, the degenerate ensemble, and the unconstructed atom}}]\label{tag:BM-54}
\textbf{Thm.}

%dsp-level:0
\[
(1)\quad\textbf{the count binding is exactly zero on both rows:}
\]

%dsp-level:0
\[
\qquad E_{\mathrm{int}}(\Sigma_{12})=6=(c_1{:}\,2)+(c_2{:}\,2)+(\text{connector}{:}\,2)\ [\text{Census FT.35}(2);\ \mathrm{BM.21}(1)]:
\]

%dsp-level:0
\[
E^{\mathrm{info}}_{\mathrm{bind}}=\bigl(6-(2{+}2{+}2)\bigr)\tfrac{\ln2}{\mathfrak h}=0
\]

%dsp-level:0
\[
\qquad [\text{machine}];
\]

%dsp-level:0
\[
\qquad \text{both baryon rows carry the same decomposition}\ [\mathrm{BM.17};\ \mathrm{BM.26}]:
\]

%dsp-level:0
\[
\qquad \Delta\bigl(E^{\mathrm{info}}_{\mathrm{bind}}\bigr)=0;
\]

%dsp-level:0
\[
(2)\quad\textbf{the ensemble binding is nonzero but class-degenerate:}
\]

%dsp-level:0
\[
\qquad \mathcal B(S_j;1,1,1)
\]

%dsp-level:0
\[
= m^2_B(\cdot,S_j)-\textstyle\sum_im^2_B(\cdot,S_i)
\]

%dsp-level:0
\[
\ne 0\ \ \text{exact rationals}\ [\mathrm{BM.21}(2)\ \text{re-cited (frozen)}];
\]

%dsp-level:0
\[
m^2_B(q,S)\ \text{is independent of}\ q\ [\mathrm{FT.107}(3);\ \text{the cancellation lemma}\ \mathrm{BM.30}(3);\ \text{machine re-run}]:
\]

%dsp-level:0
\[
\qquad \mathcal B^{\,(n)}=\mathcal B^{\,(p)}\ \ \text{exactly}\implies\boxed{\ \Delta m_{\mathrm{binding}}=0\ }\ \ \text{on every normalized ensemble read};
\]

%dsp-level:0
\[
(3)\quad\textbf{the charge configuration does not enter either binding read:}
\]

%dsp-level:0
\[
\qquad \text{the count read is}\ E_{\mathrm{int}}\text{-additive}\ [(1);\ \mathrm{BM.21}(1)];
\]

%dsp-level:0
\[
\qquad \text{the ensemble deficit}\ \Delta S=\tfrac32(2E-|R|)\ \text{is}\ \sigma\text{-free}\ [\mathrm{FT.107}(2)]:
\]

%dsp-level:0
\[
\text{the end-sign configurations that distinguish the classes}\ [\mathrm{D3.4};\ \mathrm{BM.28}(2)]\ \text{are invisible to both binding reads}:
\]

%dsp-level:0
\[
\qquad \textbf{``different charge configurations}\implies\text{different binding energies'' has no derived object between the two rows};
\]

%dsp-level:0
\[
(4)\quad\textbf{the EM binding belongs to the unconstructed pair system:}
\]

%dsp-level:0
\[
\qquad \text{the lepton--baryon bound channel is}\ \neg\text{constructed}\ [\mathrm{BM.48}(3):\ \mathrm{Def\text{-}BM.1}\text{'s grammar is colour-sealed;}\ \neg\text{an EM sealing grammar};\ \mathrm{BM.6}(5)\ \text{dynamics open}];
\]

%dsp-level:0
\[
\textbf{the category bar:}\quad\text{a binding energy is the energy of the bound system,}\ \neg\text{a correction to a free member's mass}\ [\mathrm{BM.48}(4)];
\]

%dsp-level:0
\[
\qquad \text{for a constructed pair the system read would be}\ \ m_{\mathrm{pair}}=m_p+m_\ell+E_{\mathrm{bind}},\quad E_{\mathrm{bind}}<0;
\]

%dsp-level:0
\[
\textbf{the sign row:}
\]

%dsp-level:0
\[
\qquad E_1=-\frac{\varepsilon_{\mathrm{EM}}^2\,m_r\,(q_1q_2)^2}{2}<0\ \ \text{at positive symbols}
\]

%dsp-level:0
\[
\qquad [\mathrm{BM.48}(2);\ \text{machine: the sign census}]:
\]

%dsp-level:0
\[
\qquad \text{where binding exists it is negative}\ —\ \text{and it exists, in this inventory, nowhere yet};
\]

%dsp-level:0
\[
(5)\quad\textbf{status:}
\]

%dsp-level:0
\[
\qquad \Delta m_{\mathrm{binding}}=0\ \text{on every constructed read}\ ((1)\text{--}(3));
\]

%dsp-level:0
\[
\qquad \text{the hypothetical pair system carries the only negative binding rows}\ ((4)),
\]

%dsp-level:0
\[
\ \text{as system reads},
\]

%dsp-level:0
\[
\ \neg\text{as member splittings};
\]

%dsp-level:0
\[
\text{the route-(ii) measure price stays uncombined}\ [\mathrm{BM.30}(5);\ \mathrm{BM.37}(3)]:
\]

%dsp-level:0
\[
\qquad \neg\text{a measured binding energy, ionization scale or atomic datum}\ [\mathrm{P5}].
\]

\textbf{Rmk (two zeros and one absence.)} The binding contribution to the splitting dies twice over, and both deaths are exact. The count binding is zero because six edges decompose as two plus two plus two — the informational read is additive and the additivity is an integer identity. The ensemble binding is nonzero, four exact rationals of it, but every one is built from charge-degenerate sector masses, so the neutron and the proton carry identical binding and their difference is zero again. The only binding that could distinguish the classes would be electromagnetic, and that lives in an atom — a bound channel the corpus's grammar has never sealed. The category bar keeps the bookkeeping honest in the meantime: negative energy belongs to systems, not to members.

\textbf{Req.} D2.9, D3.4 (Parts 1–2); P5 (Parts 1–2); FT.107(2), FT.107(3) (Part 3); Census FT.35(2) (row-wise); Def-BM.1, Def-BM.5, BM.6(5), BM.17, BM.20(4), BM.21(1), BM.21(2), BM.26, BM.28(2), BM.30(3), BM.30(5), BM.37(3), BM.48(2), BM.48(3), BM.48(4) (this part); exact rational arithmetic, sign censuses, the cancellation lemma [math; machine: \texttt{checks/\allowbreak{}verify\_\allowbreak{}H\_\allowbreak{}full\_\allowbreak{}scope\_\allowbreak{}mass\_\allowbreak{}splitting.py} §D].

\textbf{Pf.}

%dsp-level:0
\[
(1):\ \text{Census FT.35}(2)\text{'s decomposition re-cited};
\]

%dsp-level:0
\[
\qquad 6-(2{+}2{+}2)=0\ [\text{machine}];
\]

%dsp-level:0
\[
\qquad 0\cdot\tfrac{\ln2}{\mathfrak h}=0\ [\text{machine}]
\]

%dsp-level:0
\[
(2):\ \mathrm{BM.21}(2)\text{'s four frozen rationals re-cited};
\]

%dsp-level:0
\[
\qquad \text{the cancellation lemma at the six frozen factors}\ \{1,2,3,15,20,\tfrac73\}\ [\mathrm{BM.30}(3)\ \text{re-run};\ \text{machine}]
\]

%dsp-level:0
\[
\implies q\text{-independence}\implies\mathcal B^{(n)}-\mathcal B^{(p)}\qquad = 0
\]

%dsp-level:0
\[
(3):\ \mathrm{FT.107}(2)\text{'s}\ \sigma\text{-freeness and}\ \mathrm{BM.21}(1)\text{'s additivity conjoined}
\]

%dsp-level:0
\[
(4):\ \mathrm{BM.48}(2)\text{--}(4)\ \text{re-cited};
\]

%dsp-level:0
\[
\qquad E_1<0\ \text{at}\ \varepsilon_{\mathrm{EM}}^2>0,\ m_r>0\ [\text{machine: sign census on rational witnesses}]
\]

%dsp-level:0
\[
(5):\ \text{the status record}.\ \blacksquare
\]

\textbf{Phys.} Binding energy is the one contribution that could have been expected to favor the neutron — a neutral composite might bind its parts more gently or more tightly than a charged one — and the corpus answers with two exact zeros and one absence. The count binding vanishes as an integer identity: six is two plus two plus two. The ensemble binding is real, four nonzero rationals, and completely charge-blind, because every sector mass inside it is. And the binding that would care about charge is the electromagnetic atom, which the colour-sealed grammar has not constructed and the category bar keeps out of the free members' masses in the meantime. The difference of the two rows' binding energies is zero on every read the corpus owns.

\end{theorem}

\begin{theorem}[{\texorpdfstring{BM.55 — The full-scope splitting and the exact ratio}{BM.55 - The full-scope splitting and the exact ratio}}]\label{tag:BM-55}
\textbf{Thm.}

%dsp-level:0
\[
(1)\quad\textbf{the decomposition, collapsed:}
\]

%dsp-level:0
\[
\qquad \Delta m:
\]

%dsp-level:0
\[
= m_n-m_p
\]

%dsp-level:0
\[
= \Delta m_{\mathrm{Lyapunov}}+\Delta m_{\mathrm{force}}+\Delta m_{\mathrm{metric}}+\Delta m_{\mathrm{binding}}+\Delta m_{\mathrm{self\text{-}int}}\ \ \text{with the derived rows}:
\]

%dsp-level:0
\[
\Delta m_{\mathrm{Lyapunov}}=0\ [\mathrm{BM.51}(6)];\qquad\Delta m_{\mathrm{metric}}=0\ [\mathrm{BM.53}(2)];\qquad\Delta m_{\mathrm{binding}}=0\ [\mathrm{BM.54}(2),(5)];
\]

%dsp-level:0
\[
\Delta m_{\mathrm{force}}
\]

%dsp-level:0
\[
= \Delta m_{\mathrm{self\text{-}int}}
\]

%dsp-level:0
\[
= -\Delta m_{\mathrm{EM}}\ [\mathrm{BM.52}(4)\text{--}(5):\ \text{the two routes are one monomial, counted once}]:
\]

%dsp-level:0
\[
\boxed{\ \Delta m=m_n-m_p=-\,\Delta m_{\mathrm{EM}}=-\,\frac{3\,c\,e^2\ln2}{2\pi\,\mathfrak h}\ <\ 0\ }\qquad[\text{machine: the monomial sum and the sign census}];
\]

%dsp-level:0
\[
(2)\quad\textbf{the exact mass ratio:}
\]

%dsp-level:0
\[
\qquad m_n=\frac{6\ln2}{\mathfrak h},
\]

%dsp-level:0
\[
\qquad m_p=\frac{6\ln2}{\mathfrak h}\bigl(1+c\,\varepsilon_{\mathrm{EM}}\bigr),
\]

%dsp-level:0
\[
\qquad \varepsilon_{\mathrm{EM}}=\frac{e^2}{4\pi}\ [\mathrm{BM.26};\ \mathrm{BM.46}(1),(3)]:
\]

%dsp-level:0
\[
\boxed{\ \frac{m_n}{m_p}=\frac1{1+c\,\varepsilon_{\mathrm{EM}}}=\frac{4\pi}{4\pi+c\,e^2}\ }
\]

%dsp-level:0
\[
\qquad [\text{machine: the cross-multiplied Laurent identity}\ m_n(4\pi+ce^2)=m_p\cdot4\pi];
\]

%dsp-level:0
\[
\textbf{the geometric expansion:}
\]

%dsp-level:0
\[
\qquad \frac{m_n}{m_p}=1-c\,\varepsilon_{\mathrm{EM}}+(c\,\varepsilon_{\mathrm{EM}})^2-(c\,\varepsilon_{\mathrm{EM}})^3+O(\varepsilon^4)
\]

%dsp-level:0
\[
\qquad [\text{machine: the truncated series identity}\ (1+x)(1-x+x^2-x^3)=1-x^4];
\]

%dsp-level:0
\[
(3)\quad\textbf{the strict order:}
\]

%dsp-level:0
\[
\qquad 0<\frac{m_n}{m_p}<1\ \ \text{for every}\ e^2>0\ \ \text{and both priced}\ c\in\{\tfrac12,\tfrac35\}
\]

%dsp-level:0
\[
\qquad [\text{machine: the denominator exceeds the numerator by}\ ce^2>0];
\]

%dsp-level:0
\[
\frac{m_n}{m_p}=1\iff e=0:
\]

%dsp-level:0
\[
\qquad \textbf{the full-scope corpus computation places the neutron strictly below the proton on every constructed read};
\]

%dsp-level:0
\[
(4)\quad\textbf{the order residue:}
\]

%dsp-level:0
\[
\qquad \text{reversing the order requires a class-dependent constituent transport splitting}\ (m^{(n)}_{\mathrm{constituent}}>m^{(p)}_{\mathrm{constituent}}):
\]

%dsp-level:0
\[
\neg\exists\ \text{such an object in Parts 1--4}:
\]

%dsp-level:0
\[
\qquad \text{the fractional constituent alphabet is unreachable by the additive read}\ [\mathrm{BM.28}(7);\ \mathrm{Def\text{-}BM.8}(5)];
\]

%dsp-level:0
\[
\qquad \text{the species labels carry no datum}\ [\mathrm{Def\text{-}BM.3}];
\]

%dsp-level:0
\[
\text{the deficit/inertial conversions are item-7 content}\ [\text{pointer}];
\]

%dsp-level:0
\[
\qquad \text{the family-profile map of}\ \mathrm{Def\text{-}BM.11}\ \text{acts on channel lines, and no clause attaches constituent profiles to the baryon rows' charge classes};
\]

%dsp-level:0
\[
\textbf{the residue is named, not filled};\qquad\neg\text{a measured mass difference is compared}\ [\mathrm{P5}];
\]

%dsp-level:0
\[
(5)\quad\textbf{the decay consequence:}
\]

%dsp-level:0
\[
\qquad Q_{\mathrm{corr}}=-(\Delta m_{\mathrm{EM}}+m_\ell+m_\nu)<0\ [\mathrm{BM.49}(1)\ \text{re-cited}]\ \ \text{now with the full-scope census behind it}:
\]

%dsp-level:0
\[
\textbf{the}\ \beta\textbf{-decay vertex is kinematically closed on every constructed read, and the closure survives the complete force operator}\ [\mathrm{BM.52}(4);\ \mathrm{BM.49}(2)\ \text{re-cited}];
\]

%dsp-level:0
\[
\qquad \text{the route-(ii) measure price stands uncombined}\ [\mathrm{BM.30}(5);\ \mathrm{BM.54}(5)];
\]

%dsp-level:0
\[
(6)\quad\textbf{status:}
\]

%dsp-level:0
\[
\qquad \text{the ratio is an exact symbolic rational function of the logged P4 coupling and the priced distribution factor}\ —\ \neg\text{a number},
\]

%dsp-level:0
\[
\ \neg\text{a fit},
\]

%dsp-level:0
\[
\ \neg\text{a comparison}\ [\mathrm{P5};\ \mathrm{EM.25}(7)];
\]

%dsp-level:0
\[
\text{the five-term decomposition is derived term by term; four terms are zero and the fifth is the BM.46 monomial counted once}\ [(1);\ \mathrm{BM.52}(5)].
\]

\textbf{Rmk (the answer the corpus gives.)} Asked for the neutron–proton mass ratio from its own structures alone, the corpus answers with a rational function: four pi over four pi plus $ce^2$. It is exact, it is symbolic in the one coupling no theorem fixes, and it is strictly below one — the neutron comes out lighter, because the only sector of the complete force operator that distinguishes the two rows is the one whose field energy the proton carries and the neutron does not. The order the decay would need is the order no constructed read supplies, and the tier names the missing object precisely: a constituent-level transport difference between the classes, of which the inventory contains no instance. That is the residue; everything else is arithmetic.

\textbf{Req.} P1, P4, P5 (Parts 1–2); CT2.1, EM.25(7) (Part 2); Def-BM.3, Def-BM.8, Def-BM.11, Def-BM.12, BM.17, BM.26, BM.28(7), BM.30(5), BM.37(3), BM.39, BM.46(1), BM.46(3), BM.46(4), BM.48, BM.49(1), BM.49(2), BM.51(6), BM.52(4), BM.52(5), BM.53(2), BM.54(2), BM.54(5) (this part); FT.79(5) (Part 3); Laurent cross-multiplication, truncated geometric series over $\mathbb Q$, sign censuses [math; machine: \texttt{checks/\allowbreak{}verify\_\allowbreak{}H\_\allowbreak{}full\_\allowbreak{}scope\_\allowbreak{}mass\_\allowbreak{}splitting.py} §E].

\textbf{Pf.}

%dsp-level:0
\[
(1):\ \text{the census sum}\ 0+0+0+0-\Delta m_{\mathrm{EM}}\ [\mathrm{BM.52}(4)\ \text{re-cited};\ \text{machine}];
\]

%dsp-level:0
\[
\qquad \text{the double-count guard}\ [\mathrm{BM.52}(5):\ \text{one monomial}];
\]

%dsp-level:0
\[
\qquad \text{the sign: the coefficient census at}\ c,e^2,\ln2,\mathfrak h>0\ [\text{machine}]
\]

%dsp-level:0
\[
(2):\ m_n(4\pi+ce^2)
\]

%dsp-level:0
\[
= \tfrac{6\ln2}{\mathfrak h}(4\pi+ce^2)\ \ \text{and}\ \ m_p\cdot4\pi
\]

%dsp-level:0
\[
= \tfrac{6\ln2}{\mathfrak h}\bigl(1+\tfrac{ce^2}{4\pi}\bigr)4\pi
\]

%dsp-level:0
\[
= \tfrac{6\ln2}{\mathfrak h}(4\pi+ce^2)\ [\text{machine: the cross-multiplied Laurent identity}];
\]

%dsp-level:0
\[
\qquad (1+x)(1-x+x^2-x^3)=1-x^4\ [\text{machine}]
\]

%dsp-level:0
\[
(3):\ (4\pi+ce^2)-4\pi=ce^2>0\ [\text{machine: coefficient census}];
\]

%dsp-level:0
\[
\qquad \text{both numerator and denominator positive at positive symbols}\implies\text{the ratio}\in(0,1);
\]

%dsp-level:0
\[
\qquad ce^2=0\iff e=0\ [\text{machine}]
\]

%dsp-level:0
\[
(4):\ \mathrm{BM.28}(7)/\mathrm{Def\text{-}BM.8}(5),\ \mathrm{Def\text{-}BM.3},\ \text{the}\ \mathrm{Def\text{-}BM.11}\ \text{locus clause re-cited};
\]

%dsp-level:0
\[
\qquad (5):\ \mathrm{BM.49}(1)\text{--}(2)\ \text{re-cited};
\]

%dsp-level:0
\[
\qquad (6):\ \text{the status record}.\ \blacksquare
\]

\textbf{Phys.} The full-scope question has a full-scope answer, and it is a single rational function. Four of the five named contributions are theorems about zero: the Lyapunov tower shares every input between the rows, the metric reads are charge-blind, the binding dies twice, and the force census leaves one sector standing. That sector's monomial is the whole splitting, negative for the neutron-minus-proton difference, and the ratio follows by cross-multiplication: four pi over four pi plus $c\,e^2$, strictly between zero and one at every positive coupling, equal to one only where the coupling dies. The corpus therefore places the neutron below the proton, deeper than the informational count alone did, and the beta decay stays shut with the complete force operator now behind the closure. What would open it — a constituent transport difference between the two rows — is named exactly and found absent: no fractional alphabet, no labelled datum, no converted species. The ratio is derived; the order is derived; and the distance to the world's order is the residue, stated to the monomial.

\end{theorem}

\medskip\noindent\textit{\texorpdfstring{BM.56 — The status of the full-scope mass splitting: adjudication and scope}{BM.56 - The status of the full-scope mass splitting: adjudication and scope}.}\ Structural-constraints record: see Appendix B, Section B.1.22 of the Addendum.\par\medskip

\phantomsection
\section[Tier BM.14 — The Constituent Transport Splitting: the…]{\texorpdfstring{Tier BM.14 — The Constituent Transport Splitting: the Vertex Asymmetry, the Observation Frames, the Six-State Quark Assembly, and the Neutron Row}{Tier BM.14 - The Constituent Transport Splitting: the Vertex Asymmetry, the Observation Frames, the Six-State Quark Assembly, and the Neutron Row}}

\textbf{Locus.} This tier fills the residue $\mathrm{BM.55}(4)$ names and does not fill: the class-dependent constituent transport splitting. It extends $\mathrm{BM.55}(1)$'s five-term decomposition by one term, amends $\mathrm{BM.52}(3)$'s sector census by one row, and reverses $\mathrm{BM.55}(3)$'s strict order on an exact sub-critical branch. Every clause of $\mathrm{BM.51}$–$\mathrm{BM.56}$ it touches is re-cited at its point of use and stands unmodified except where a clause's own scope row is widened; the widening is recorded in $\mathrm{BM.64}(3)$.

\textbf{Two directions.} The tier is worked from both ends at once. From the force structure: what a sealed junction can carry, and what the seal annihilates — $\mathrm{BM.57}$ and $\mathrm{BM.58}$ run that direction to exhaustion and refute both proposed mechanisms. From the mass spectrum: what $\mathrm{BM.55}(1)$'s four zeros leave, and what the inventory contains that is class-dependent and not barred — $\mathrm{BM.59}$ runs that direction. The two meet on one object, $W(Q)$, and $\mathrm{BM.62}$ assembles it. $\mathrm{BM.60}$ supplies the six-state spectrum the constituent row of $\mathrm{BM.59}(5)$ needs; $\mathrm{BM.61}$ disposes of the proposed coupling $\alpha_B$; $\mathrm{BM.63}$ re-adjudicates the decay; $\mathrm{BM.64}$ adjudicates the tier.

\begin{definition}[{\texorpdfstring{Def-BM.13 — The constituent transport refinement of the rest input (the construction)}{Def-BM.13 - The constituent transport refinement of the rest input (the construction)}}]\label{tag:Def-BM-13}
\textbf{Def.} A refinement of the dispersion route's rest input $[\mathrm{Def\text{-}BM.10}(4)$; $\mathrm{BM.35}(4)]$ built from the additive charge read's own census $[\mathrm{Def\text{-}BM.8}(6)]$; no new axiom, no new datum, no new parameter, no numerical coupling, no route adoption. Local notation: $\varepsilon_{\mathrm{EM}}:=e^2/(4\pi)$, a dimensionless P4-class symbol carried symbolically throughout $[\mathrm{Def\text{-}BM.12}$ local notation; $\mathrm{EM.25}(7),(9)$; $\mathrm{P5}]$.

%dsp-level:0
\[
(1)\quad\textbf{the constituent transport multiplicity:}
\]

%dsp-level:0
\[
\qquad q_{\mathrm{core}}(c)=\tfrac12\bigl(\sigma_{s(c)}-\sigma_{t(c)}\bigr)\in\{-1,0,+1\}
\]

%dsp-level:0
\[
\qquad [\mathrm{D3.4};\ \mathrm{EM.9}];
\]

%dsp-level:0
\[
\textbf{the end-pair census:}
\]

%dsp-level:0
\[
\qquad (\sigma_s,\sigma_t)\in\{(+,+),(+,-),(-,+),(-,-)\}\ \longmapsto\ q=(0,+1,-1,0)
\]

%dsp-level:0
\[
\qquad [\mathrm{EM.9}\ \text{table};\ \text{machine}]:
\]

%dsp-level:0
\[
\boxed{\ m_q=(2,\,1,\,1)\ \ \text{at}\ \ q=(0,+1,-1)\ }\qquad[\mathrm{FT.106}(3\mathrm a):\ \text{the norm-layer closed form}\ m_q\,2^{|R|-2}]:
\]

%dsp-level:0
\[
\textbf{a neutral-core constituent carries two transport realizations; a charged-core constituent carries one};
\]

%dsp-level:0
\[
\qquad \text{the two neutral realizations are exchanged by the global sign flip, which maps}\ q
\]

%dsp-level:0
\[
\mapsto -q\ \text{on a charged one}\ [\mathrm{EM.9a};\ \mathrm{EM.24}(1)]:
\]

%dsp-level:0
\[
\textbf{the multiplicity is the orbit size of the constituent's transport register under}\ \omega\ \text{within its own charge class};
\]

%dsp-level:0
\[
(2)\quad\textbf{the parity law of the chain wall:}
\]

%dsp-level:0
\[
\qquad \text{for an oriented chain}\ c=(e_1,\dots,e_k),\ k\ge2,\ \text{with interior vertices of type}\ (1,1)\ [\mathrm{A7};\ \mathrm{D0.8}]:
\]

%dsp-level:0
\[
\sum_{i=1}^{k-1}f_{v_i}=\tfrac12\bigl(\sigma_{e_k}-\sigma_{e_1}\bigr)=-\,q_{\mathrm{EM}}(c),\qquad\sum_{i=1}^{k-1}\aleph_{v_i}=\#\{i:\ \sigma_{e_{i+1}}\ne\sigma_{e_i}\}\qquad[\mathrm{EM.9b};\ \mathrm{D3.3}]:
\]

%dsp-level:0
\[
\aleph_v=|f_v|\in\{0,1\}\implies\boxed{\ \sum_i\aleph_{v_i}\equiv|q|\ (\mathrm{mod}\ 2),\qquad\sum_i\aleph_{v_i}\ \ge\ |q|\ }
\]

%dsp-level:0
\[
\qquad [\text{machine: exhaustive over every sign sequence,}\ k=2,\dots,7]:
\]

%dsp-level:0
\[
\textbf{the interior wall content of a constituent's transport is parity-locked to its charge core:}
\]

%dsp-level:0
\[
\qquad q=0:\ \min\textstyle\sum\aleph=0\ (\text{even});
\]

%dsp-level:0
\[
\qquad |q|=1:\ \min\textstyle\sum\aleph=1\ (\text{odd});
\]

%dsp-level:0
\[
(3)\quad\textbf{the class transport weight:}
\]

%dsp-level:0
\[
\qquad W(Q)\ :=\ \sum_{\substack{(q_1,\dots,q_k)\in\{-1,0,+1\}^k\  \sum_iq_i=Q}}\ \prod_{i=1}^{k}m_{q_i}\ =\ \bigl[z^{Q}\bigr]\bigl(2+z+z^{-1}\bigr)^{k}
\]

%dsp-level:0
\[
\qquad [\mathrm{Def\text{-}BM.8}(6)]:
\]

%dsp-level:0
\[
k=3:
\]

%dsp-level:0
\[
\qquad W=(1,\,6,\,15,\,\mathbf{20},\,15,\,6,\,1)\ \ \text{over}\ \ Q=-3,\dots,+3,
\]

%dsp-level:0
\[
\qquad \sum_{Q}W(Q)=4^3=64=2^6
\]

%dsp-level:0
\[
\qquad [\text{machine: the exhaustive}\ 4^3\ \text{end-sign census and the convolution agree exactly}];
\]

%dsp-level:0
\[
\textbf{the ordered charge-triple census:}
\]

%dsp-level:0
\[
\qquad \#\{(q_1,q_2,q_3):\ \textstyle\sum q_i=k\}=(1,3,6,\mathbf 7,6,3,1):
\]

%dsp-level:0
\[
\qquad 7\ \text{neutral triples against}\ 6\ \text{charged}\ [\mathrm{BM.28}(2);\ \text{machine}];
\]

%dsp-level:0
\[
\textbf{the factorization identity:}
\]

%dsp-level:0
\[
\qquad W(Q)=\sum_{\text{triples at}\ Q}\prod_im_{q_i}
\]

%dsp-level:0
\[
\qquad [\text{machine: the two routes agree class by class}];
\]

%dsp-level:0
\[
(4)\quad\textbf{the charge-class transport entropy:}\qquad\boxed{\ \Sigma_{\mathrm{cl}}(Q)\ :=\ \ln W(Q)\ }\qquad\Sigma_{\mathrm{cl}}(0)=\ln20,\quad\Sigma_{\mathrm{cl}}(\pm1)=\ln15,
\]

%dsp-level:0
\[
\boxed{\ \Sigma_{\mathrm B}\ :=\ \Sigma_{\mathrm{cl}}(0)-\Sigma_{\mathrm{cl}}(+1)\ =\ \ln\frac{W(0)}{W(+1)}\ =\ \ln\frac43\ =\ 2\ln2-\ln3\ >\ 0\ }
\]

%dsp-level:0
\[
\qquad [\text{machine: the certificate}\ 4>3]:
\]

%dsp-level:0
\[
\textbf{the constituent row:}\qquad\ln m_0-\ln m_{\pm1}=\ln2>0\qquad[\text{machine:}\ 2>1]:
\]

%dsp-level:0
\[
\textbf{the neutral-core constituent is the heavier one, by exactly}\ \ln2/\mathfrak h;
\]

%dsp-level:0
\[
\qquad \textbf{the additive constraint attenuates it:}\quad0<\ln\tfrac43<\ln2<3\ln2
\]

%dsp-level:0
\[
\qquad [\text{machine:}\ 3<4<6<24];
\]

%dsp-level:0
\[
(5)\quad\textbf{the refined rest input:}\qquad\boxed{\ m_{\mathrm{rest}}(Q)\ :=\ \frac{E_{\mathrm{int}}\ln2+\Sigma_{\mathrm{cl}}(Q)}{\mathfrak h}\ +\ c\,\varepsilon_{\mathrm{EM}}\,Q^2\,m_{\mathrm{hid}}\ }
\]

%dsp-level:0
\[
\bigl[\ \text{the first summand:}\ \mathrm{BM.35}(4)\text{'s reframing identity}\ m_{\mathrm{rest}}=\Delta S_{\mathfrak L}/\mathfrak h\ \text{with the register summand of}\ (4);\qquad\text{the second:}\ \mathrm{BM.46}(3)\text{'s corrected rest read, riding the charge-blind Born width}\ [\mathrm{BM.28}(1)]\ \bigr];
\]

%dsp-level:0
\[
\textbf{at}\ E_{\mathrm{int}}=6\ [\mathrm{BM.17};\ \mathrm{BM.26}]:
\]

%dsp-level:0
\[
\qquad m_{\mathrm{rest}}(0)\,\mathfrak h=\ln(2^6\!\cdot\!20)=\ln1280,
\]

%dsp-level:0
\[
\qquad m_{\mathrm{rest}}(\pm1)\,\mathfrak h=\ln(2^6\!\cdot\!15)+6c\,\varepsilon_{\mathrm{EM}}\ln2=\ln960+6c\,\varepsilon_{\mathrm{EM}}\ln2
\]

%dsp-level:0
\[
\qquad [\text{machine}];
\]

%dsp-level:0
\[
(6)\quad\textbf{the currency clause:}
\]

%dsp-level:0
\[
\qquad \Sigma_{\mathrm{cl}}\ \text{is measured in nats of binary registration data}\ —\ \text{the end signs are the}\ \sigma\text{-alphabet}\ [\mathrm{CT0.9};\ \mathrm{D3.4}]\ \text{and one binary datum books}\ \ln2\ [\mathrm{C9};\ \mathrm{CT1.29}]:
\]

%dsp-level:0
\[
\textbf{it is the same currency as}\ \Delta S_{\mathfrak L}
\]

%dsp-level:0
\[
= E_{\mathrm{int}}\ln2\ [\mathrm{BM.35}(3)],
\]

%dsp-level:0
\[
\ \text{and its alphabet is disjoint from the count's:}
\]

%dsp-level:0
\[
\qquad E_{\mathrm{int}}\ \text{has no endpoint summand}\ [\mathrm{BM.28}(3):\ 6=(c_1{:}\,2)+(c_2{:}\,2)+(\text{connector}{:}\,2)]:
\]

%dsp-level:0
\[
E_{\mathrm{int}}\ \text{is constant on every end-sign fibre}\ [\text{machine: exhaustive over the four end pairs}]\implies\textbf{the two registrations are independent, so their entropies add};
\]

%dsp-level:0
\[
(7)\quad\textbf{the reference clause (the sign is fixed, not chosen):}
\]

%dsp-level:0
\[
\qquad \text{the deficit reference is the vacuum:}\quad\Delta S:=S-S(\mathrm{vac}),\quad S_{\mathfrak L}(\mathrm{vac})=0\ [\mathrm{CT0.2};\ \mathrm{D0.1};\ \mathrm{BM.35}(3)]:
\]

%dsp-level:0
\[
W(\mathrm{vac})=1\implies\Sigma_{\mathrm{cl}}(\mathrm{vac})=0\implies\Delta\Sigma_{\mathrm{cl}}(Q)=\ln W(Q)\ \text{with the}\ \textbf{positive}\ \text{sign of}\ (4);
\]

%dsp-level:0
\[
\textbf{the priced alternative:}
\]

%dsp-level:0
\[
\qquad \text{referencing against the unconstrained register}\ \bigl(4^3=64\bigr)\ \text{gives the constraint cost}\ \ln\frac{64}{W(Q)},
\]

%dsp-level:0
\[
\ \text{whose class difference is}\ -\ln\tfrac43:
\]

%dsp-level:0
\[
\text{that reference is}\ \neg\textbf{a state of the composite}\ —\ \text{the boundary data fix the charge class}\ [\mathrm{Def\text{-}BM.5}(1)]\ —\ \text{and it is class-dependent:}
\]

%dsp-level:0
\[
\qquad \textbf{priced, barred};
\]

%dsp-level:0
\[
\mathrm{Def\text{-}H.28}\text{'s Ansatz-B}\ m_B
\]

%dsp-level:0
\[
= (K-\alpha\ln\Omega)\ln2/\mathfrak h\ \text{carries the opposite sign with}\ \alpha\ \text{undeclared and is marked an external ansatz,}\ \neg\text{a corpus read:}
\]

%dsp-level:0
\[
\qquad \neg\text{adopted};
\]

%dsp-level:0
\[
(8)\quad\textbf{the bar against the measure price:}
\]

%dsp-level:0
\[
\qquad \mathrm{BM.37}(6)\ \text{bars reading a route ratio as a physical mass ratio};
\]

%dsp-level:0
\[
\qquad \textbf{this construction reads no ratio}:
\]

%dsp-level:0
\[
\text{the entry is}\ \ln W(Q)\ \text{in the deficit currency, and the resulting mass ratio is}\ \ln1280/\ln960\ne\tfrac43\ [\text{machine}];
\]

%dsp-level:0
\[
\qquad \textbf{the strict contraction:}\quad\frac{\Sigma_{\mathrm B}}{E_{\mathrm{int}}\ln2}=\frac{\ln(4/3)}{6\ln2}<\frac13
\]

%dsp-level:0
\[
\qquad [\text{machine:}\ 3\ln\tfrac43<2\ln2\iff16<27]:
\]

%dsp-level:0
\[
\text{the refinement is a strict contraction of}\ \mathrm{BM.30}(2)\text{'s measure price}\ —\ \text{it never reaches it};
\]

%dsp-level:0
\[
(9)\quad\textbf{status and scope:}
\]

%dsp-level:0
\[
\qquad \textbf{constructed}\ —\ \text{an abbreviation of derived objects}\ ((1)\text{--}(4)\ \text{are censuses};\ (5)\ \text{is an assembly of}\ \mathrm{BM.35}(4)\ \text{and}\ \mathrm{BM.46}(3));
\]

%dsp-level:0
\[
\neg\text{a declaration},\quad\neg\text{a P4 addition},\quad\neg\text{a new dimensionful scale}\ [\mathrm{P1}],\quad\text{§Cal-BM stays empty};
\]

%dsp-level:0
\[
\qquad \neg\text{a route adopted}\ [\mathrm{BM.30}(5)\ \text{stands}];
\]

%dsp-level:0
\[
\textbf{beyond this tier:}\ \text{the triple-level resolution}\ —\ \text{the chart does not resolve which triple a class member occupies, because the end signs carry no position datum}\ [\mathrm{D3.4};\ \mathrm{BM.46}(2)\ \text{resolution residue}]:\ \text{the class read of}\ (4)\ \text{is the one the chart supports}.
\]

\textbf{Rmk (what the refinement is).} The composite's rest input has always been the count of its internal directed edges, and the count is charge-blind to the last monomial. What the charge class does carry, and has always carried, is a filter on the transport register: twenty admissible end-sign assignments in the neutral class against fifteen in the charged one. The refinement takes the logarithm of that filter and books it where the corpus books every other entropy — against the vacuum, in the ledger's own currency, through the rest slot the dispersion read already reads. It is not a fourth mass route: it adds no functional of the flicker configurations, and it takes no ratio. It is the smallest entry the inventory supplies that is class-dependent, currency-matched, and alphabet-disjoint from the count.

\end{definition}

\pseudohead{definition}{\texorpdfstring{Def-BM.14}{Def-BM.14}}{\texorpdfstring{The six-state assembly: the Platonic floors, the defect branches, and the generation map}{The six-state assembly: the Platonic floors, the defect branches, and the generation map}}\label{tag:Def-BM-14}
\textbf{Def.} A census assembled from the Platonic card template $[\mathrm{Census\ FT.42}(5)]$, the duality degeneracy and the ledger splitting $[\mathrm{FT.44}(1)\text{--}(3)]$, the physical-card defect spectrum $[\mathrm{Census\ FT.63}(6)]$, and the $q$-channel triplet $[\mathrm{FT.119};\ \mathrm{Def\text{-}BM.6}(3)]$; no new axiom, no new datum, no new parameter, no external mass.

%dsp-level:0
\[
(1)\quad\textbf{the three mass floors:}
\]

%dsp-level:0
\[
\qquad m_{\mathrm{hid}}=\frac{2E\ln2}{\mathfrak h},
\]

%dsp-level:0
\[
\qquad 2E\in\{12,\,24,\,24,\,60,\,60\}\ \ \text{over the five solids}
\]

%dsp-level:0
\[
\qquad [\mathrm{FT.42}(1),\ \text{card table}]:
\]

%dsp-level:0
\[
\textbf{dual pairs are mass-degenerate and}\ \Pi\textbf{-degenerate at matched}\ K_S\ [\mathrm{FT.44}(1)]:
\]

%dsp-level:0
\[
\qquad \text{cube/octahedron}\ 24\ln2/\mathfrak h;
\]

%dsp-level:0
\[
\quad\text{dodecahedron/icosahedron}\ 60\ln2/\mathfrak h;
\]

%dsp-level:0
\[
\quad\text{tetrahedron self-dual}\ 12\ln2/\mathfrak h;
\]

%dsp-level:0
\[
\boxed{\ \textbf{three mass values across the five solids}\ }
\]

%dsp-level:0
\[
\qquad \mathcal F:
\]

%dsp-level:0
\[
= \bigl\{\mathcal F_1{=}\{\text{tetrahedron}\},\ \ \mathcal F_2{=}\{\text{cube},\text{octahedron}\},\ \ \mathcal F_3{=}\{\text{dodecahedron},\text{icosahedron}\}\bigr\};
\]

%dsp-level:0
\[
2E(\mathcal F_1){=}12\ <\ 2E(\mathcal F_2){=}24\ <\ 2E(\mathcal F_3){=}60,
\]

%dsp-level:0
\[
\qquad 12:24:60=1:2:5
\]

%dsp-level:0
\[
\qquad [\text{machine: exact rationals, strictly increasing}];
\]

%dsp-level:0
\[
(2)\quad\textbf{the branch on each floor (the derived 2):}
\]

%dsp-level:0
\[
\qquad \text{a floor's two members are distinguished by every finer ledger while sharing}\ m_{\mathrm{hid}}\ \text{and}\ \Pi\ [\mathrm{FT.44}(2)]:
\]

%dsp-level:0
\[
\mathcal F_2,\ \mathcal F_3:\ \textbf{the duality branch}\ —\ \text{the dual exchanges}\ (V,F)\ \text{at fixed}\ E\ [\mathrm{FT.44}(1)];
\]

%dsp-level:0
\[
\qquad \mathcal F_1:\ \textbf{the weight-ladder branch}\ —\ \text{the tetrahedron is self-dual}\ [\mathrm{FT.44}(3)],
\]

%dsp-level:0
\[
\ \text{so its branch index is the ladder rung}\ m\in\{1,2\}\ [\mathrm{FT.60}(5);\ \mathrm{Census\ FT.63}(6)];
\]

%dsp-level:0
\[
\textbf{the six states:}\qquad\mathcal S\ :=\ \{(\mathcal F_g,\beta):\ g\in\{1,2,3\},\ \beta\in\{-,+\}\},\qquad|\mathcal S|=3\cdot2=6\qquad[\text{machine}];
\]

%dsp-level:0
\[
(3)\quad\textbf{the defect spectrum (the branch mass):}
\]

%dsp-level:0
\[
\qquad \delta_v=2\pi-d\,\frac{(L-2)\pi}L,
\]

%dsp-level:0
\[
\qquad D_S=4\sin^2\frac{\delta_v}4,
\]

%dsp-level:0
\[
\qquad m^2_{\mathrm{pole}}=\hat\mu'^2\,D_S
\]

%dsp-level:0
\[
\qquad [\mathrm{Census\ FT.63}(4),(6);\ \mathrm{Def\text{-}H.37}(2);\ \mathrm{FT.54}(4)]:
\]

%tbl-level:0
{\setlength{\tabcolsep}{4pt}
\begin{longtable}{>{\raggedright\arraybackslash}p{\dimexpr0.0887\textwidth-2\tabcolsep\relax}|>{\raggedright\arraybackslash}p{\dimexpr0.2940\textwidth-2\tabcolsep\relax}|>{\raggedright\arraybackslash}p{\dimexpr0.1003\textwidth-2\tabcolsep\relax}|>{\raggedright\arraybackslash}p{\dimexpr0.1586\textwidth-2\tabcolsep\relax}|>{\raggedright\arraybackslash}p{\dimexpr0.1819\textwidth-2\tabcolsep\relax}|>{\raggedright\arraybackslash}p{\dimexpr0.1764\textwidth-2\tabcolsep\relax}}
\toprule
state & $(\mathcal F_g,\allowbreak \beta)$ & $\delta_v/\pi$ & $D_S$ & exact field & min. polynomial of $D_S$ \\
\addlinespace[2pt]
\midrule
\endfirsthead
\toprule
state & $(\mathcal F_g,\allowbreak \beta)$ & $\delta_v/\pi$ & $D_S$ & exact field & min. polynomial of $D_S$ \\
\addlinespace[2pt]
\midrule
\endhead
\bottomrule
\endlastfoot
$s_1$ & $(\mathcal F_1,\allowbreak -)$, the ladder rung $m{=}1$ & $1$ & $2$ & $\mathbb Q$ & $x-2$ \\
$s_2$ & $(\mathcal F_1,\allowbreak +\allowbreak )$, the ladder rung $m{=}2$ & $1$ & $4$ & $\mathbb Q$ & $x-4$ \\
$s_3$ & $(\mathcal F_2,\allowbreak -)$, the cube member & $\tfrac12$ & $2-\sqrt2$ & $\mathbb Q(\sqrt2)$ & $x^2-4x+\allowbreak 2$ \\
$s_4$ & $(\mathcal F_2,\allowbreak +\allowbreak )$, the octahedron member & $\tfrac23$ & $1$ & $\mathbb Q$ & $x-1$ \\
$s_5$ & $(\mathcal F_3,\allowbreak -)$, the dodecahedron member & $\tfrac15$ & $2-2\cos\tfrac\pi{10}$ & $\mathbb Q(\cos\tfrac\pi{10})$ & $x^4-8x^3+\allowbreak 19x^2-12x+\allowbreak 1$ \\
$s_6$ & $(\mathcal F_3,\allowbreak +\allowbreak )$, the icosahedron member & $\tfrac13$ & $2-\sqrt3$ & $\mathbb Q(\sqrt3)$ & $x^2-4x+\allowbreak 1$ \\
\end{longtable}
}

%dsp-level:0
\[
\textbf{Descartes/Gauss--Bonnet:}
\]

%dsp-level:0
\[
\qquad V\,\delta_v=4\pi\ \ \text{on all five solids}
\]

%dsp-level:0
\[
\qquad [\mathrm{Census\ FT.63}(4);\ \text{machine: exact in}\ \pi\text{-units}];
\]

%dsp-level:0
\[
(4)\quad\textbf{the branch order (derived, exact, uniform):}\qquad\boxed{\ D_S(\mathcal F_g,-)\ <\ D_S(\mathcal F_g,+)\qquad\forall g\in\{1,2,3\}\ }
\]

%dsp-level:0
\[
\mathcal F_1:\ 4>2;
\]

%dsp-level:0
\[
\qquad \mathcal F_2:\ 1>2-\sqrt2\iff\sqrt2>1\iff2>1;
\]

%dsp-level:0
\[
\qquad \mathcal F_3:\ 2-\sqrt3>2-2\cos\tfrac\pi{10}\iff\cos\tfrac\pi{10}>\tfrac{\sqrt3}2\iff\tfrac\pi{10}<\tfrac\pi6\iff3<5
\]

%dsp-level:0
\[
\qquad [\text{machine: integer and radical certificates}];
\]

%dsp-level:0
\[
\textbf{the six values are pairwise distinct}\ —\ \text{the defect species resolves all six states with zero degeneracy}\ [\text{machine: separation by field and by minimal polynomial}];
\]

%dsp-level:0
\[
(5)\quad\textbf{the branch splitting (the constituent mass difference):}
\]

%dsp-level:0
\[
\qquad \boxed{\ \Delta_g\ :=\ \hat\mu'^2\bigl(D_S(\mathcal F_g,+)-D_S(\mathcal F_g,-)\bigr)\ }
\]

%dsp-level:0
\[
\Delta_1=2\hat\mu'^2\in\mathbb Q\hat\mu'^2;
\]

%dsp-level:0
\[
\qquad \Delta_2=(\sqrt2-1)\hat\mu'^2\in\mathbb Q(\sqrt2)\hat\mu'^2;
\]

%dsp-level:0
\[
\qquad \Delta_3=\bigl(2\cos\tfrac\pi{10}-\sqrt3\bigr)\hat\mu'^2\in\mathbb Q\bigl(\sqrt3,\cos\tfrac\pi{10}\bigr)\hat\mu'^2;
\]

%dsp-level:0
\[
\Delta_g>0\ \ \forall g
\]

%dsp-level:0
\[
\qquad [(4)\text{'s certificates}];
\]

%dsp-level:0
\[
\qquad \textbf{three distinct fields, three exact positive splittings, no imported value}\ [\mathrm{P5}];
\]

%dsp-level:0
\[
(6)\quad\textbf{the generation map (the unique order-isomorphism):}
\]

%dsp-level:0
\[
\qquad \text{the}\ q\text{-channel's three lines are ordered exactly:}
\]

%dsp-level:0
\[
\quad (r^{(q)}_1)^2<(r^{(q)}_2)^2<(r^{(q)}_0)^2\ [\mathrm{FT.119}(5);\ \text{enclosure-proven, frozen}];
\]

%dsp-level:0
\[
\text{the three floors are ordered exactly:}\quad2E(\mathcal F_1)<2E(\mathcal F_2)<2E(\mathcal F_3)\ [(1)]:
\]

%dsp-level:0
\[
\boxed{\ \text{two exact three-element ordered sets admit exactly one order-isomorphism}\ }
\]

%dsp-level:0
\[
\qquad g=1\leftrightarrow(r^{(q)}_1)^2\leftrightarrow\mathcal F_1,
\]

%dsp-level:0
\[
\quad g=2\leftrightarrow(r^{(q)}_2)^2\leftrightarrow\mathcal F_2,
\]

%dsp-level:0
\[
\quad g=3\leftrightarrow(r^{(q)}_0)^2\leftrightarrow\mathcal F_3;
\]

%dsp-level:0
\[
\textbf{the species row:}
\]

%dsp-level:0
\[
\qquad \text{on the transport species the line order inverts}\ [\mathrm{BM.39}(5):\ \text{the spectral-maximal line receives the mass-minimum}]\implies\text{the isomorphism inverts with it};
\]

%dsp-level:0
\[
\qquad \text{the relating equation between the two currencies is}\ \neg\text{derived}\ [\mathrm{BM.35}(5)];
\]

%dsp-level:0
\[
\textbf{the scale is not supplied:}\qquad\hat\mu'^2\in\mathrm{Cal\text{-}F}\ (\text{logged}),\quad\mu\in\mathrm{Cal\text{-}F}:\qquad\text{the map is an order,}\ \neg\text{a number}\ [\mathrm{P4};\ \mathrm{P5}];
\]

%dsp-level:0
\[
(7)\quad\textbf{the charge alphabet of the six (the second derived 2):}
\]

%dsp-level:0
\[
\qquad \mathcal H_{\mathsf v}\cong\mathbb C^2\ [\mathrm{EW.32}];
\]

%dsp-level:0
\[
\qquad Q|_{\mathcal H_{\mathsf v}}=\operatorname{diag}(1,0),\quad\operatorname{spec}Q=\{0,+1\},\quad\omega:\ \operatorname{spec}\mapsto\{0,-1\}\ [\mathrm{EW.36}]:
\]

%dsp-level:0
\[
\text{the seed carrier's weight doublet}\ \{0,+1\}\ \text{together with its}\ \omega\text{-image is exactly the ternary end-read alphabet}\ \{-1,0,+1\}\ [\mathrm{EW.36};\ \mathrm{Def\text{-}BM.12}(1);\ \mathrm{D3.4}];
\]

%dsp-level:0
\[
\textbf{the branch index of}\ (2)\ \textbf{and the weight index of}\ \mathcal H_{\mathsf v}\ \textbf{are two derived twos};
\]

%dsp-level:0
\[
\qquad \text{their attachment is}\ (8)\text{'s clause};
\]

%dsp-level:0
\[
(8)\quad\textbf{the identification clause (the single declared row):}
\]

%dsp-level:0
\[
\qquad \textbf{down-type}\ :=\ \text{the floor's heavier defect branch}\ (\beta{=}+);
\]

%dsp-level:0
\[
\qquad \textbf{up-type}\ :=\ \text{the lighter}\ (\beta{=}−);
\]

%dsp-level:0
\[
\text{the ordering is derived}\ [(4)];
\]

%dsp-level:0
\[
\ \text{the name attachment is a label on a derived structural tuple and carries no datum}\ [\mathrm{Def\text{-}BM.3}\ \text{pattern};\ \mathrm{Def\text{-}BM.3}(7)\ \text{scope}];
\]

%dsp-level:0
\[
\qquad \textbf{declared}\ (\mathrm{A\text{-}BM.6})\ \textbf{before use};
\]

%dsp-level:0
\[
\textbf{non-derived:}\ \text{the attachment of the}\ Q\text{-weight index to the branch index}\ (7);
\]

%dsp-level:0
\[
\qquad \text{the flavour composition of a sealed composite}\ —\ \text{which three of the six its constituents occupy}\ [\mathrm{BM.28}(7);\ \mathrm{Def\text{-}BM.8}(5):\ \text{a fractional constituent alphabet is structurally impossible under an integer-valued additive read}];
\]

%dsp-level:0
\[
(9)\quad\textbf{status and scope:}
\]

%dsp-level:0
\[
\qquad \textbf{derived:}\ (1)\text{--}(7);
\]

%dsp-level:0
\[
\qquad \textbf{declared:}\ (8)\text{'s name attachment alone};
\]

%dsp-level:0
\[
\qquad \neg\text{a P4 addition},\ \neg\text{a new dimensionful scale},\ \text{§Cal-BM stays empty};
\]

%dsp-level:0
\[
\neg\exists\ \text{a measured mass, ratio or Yukawa value in any clause}\ [\mathrm{P5};\ \text{machine: zero external tokens, zero decimal numerals}];
\]

%dsp-level:0
\[
\qquad \text{species names are confined to the identification layers}\ [\mathrm{Def\text{-}BM.3}(7),\ \mathrm{Def\text{-}BM.6}(7)];
\]

%dsp-level:0
\[
\textbf{beyond this tier:}\ \text{the scale relating the defect and informational currencies}\;
\]

%dsp-level:0
\[
\qquad \text{the flavour composition of the baryon classes}\;
\]

%dsp-level:0
\[
\qquad \text{the materialization scope}\ \mathrm{M1}/\mathrm{M2}\ [\mathrm{Def\text{-}H.20}]\ \text{is inherited at every use}.
\]

\textbf{Rmk (six states, no data).} Nothing in this definition is fitted. The floors come from the card template's edge counts; the branches come from the duality that leaves those counts fixed and splits everything finer; the branch masses come from the conical defect at a vertex where three faces meet, closed by Descartes against four pi. What the assembly adds is one order-isomorphism between two three-element ordered sets the corpus had already derived separately, and one name. The result is six states, six pairwise distinct exact algebraic masses on the defect species, three exact rationals on the informational one, and a hierarchy that runs one way in the bit currency and the other way in the angle currency — with the equation between the two currencies still open, and still not guessed.

\pseudohead{theorem}{\texorpdfstring{BM.57}{BM.57}}{\texorpdfstring{The continuation-symmetry census: the vertex asymmetry is two, and it splits nothing}{The continuation-symmetry census: the vertex asymmetry is two, and it splits nothing}}\label{tag:BM-57}
\textbf{Thm.}

%dsp-level:0
\[
(1)\quad\textbf{the strong 3-vertex:}
\]

%dsp-level:0
\[
\qquad \text{the invariant tensor is}\ \varepsilon_{ijk}\ [\mathrm{CT3.17};\ \mathrm{D3.9}];
\]

%dsp-level:0
\[
\qquad \varepsilon_{\sigma(i)\sigma(j)\sigma(k)}=\operatorname{sgn}(\sigma)\,\varepsilon_{ijk}\ \ \forall\sigma\in S_3
\]

%dsp-level:0
\[
\qquad [\text{machine: all six permutations}]:
\]

%dsp-level:0
\[
\textbf{the continuation group:}
\]

%dsp-level:0
\[
\qquad \operatorname{Stab}_{S_3}(\varepsilon)=A_3\cong C_3=S_3\cap SU(3),
\]

%dsp-level:0
\[
\qquad |A_3|=3,
\]

%dsp-level:0
\[
\qquad \text{nontrivial count}\ 2
\]

%dsp-level:0
\[
\qquad [\mathrm{BM.3}(1);\ \text{machine: the invariant set is exactly the cyclic group}];
\]

%dsp-level:0
\[
(2)\quad\textbf{the}\ (1,1)\textbf{-vertex (the chain wall):}
\]

%dsp-level:0
\[
\qquad \text{the invariant tensor is}\ \delta^i{}_j\ [\mathrm{CT3.17}];
\]

%dsp-level:0
\[
\qquad \text{the type-preserving relabelling group of a}\ (1,1)\ \text{vertex is}\ S_1\times S_1:
\]

%dsp-level:0
\[
\qquad |\,\cdot\,|=1,
\]

%dsp-level:0
\[
\qquad \text{nontrivial count}\ 0;
\]

%dsp-level:0
\[
\textbf{the leg swap is not a continuation:}
\]

%dsp-level:0
\[
\qquad \delta^i{}_j
\]

%dsp-level:0
\[
= \delta_i{}^j\ \text{admits it as a bare 2-tensor}\ [\text{machine}],
\]

%dsp-level:0
\[
\ \text{but it acts on the oriented end registers as}\ (\sigma_s,\sigma_t)
\]

%dsp-level:0
\[
\mapsto (\sigma_t,\sigma_s):
\]

%dsp-level:0
\[
q_{\mathrm{EM}}\mapsto-q_{\mathrm{EM}}\ \ \text{on all four end pairs}
\]

%dsp-level:0
\[
\qquad [\mathrm{D3.4};\ \mathrm{EM.9a}:\ \text{charge-conjugation odd};\ \text{machine}]:
\]

%dsp-level:0
\[
\qquad \textbf{the swap is the}\ \omega\textbf{-involution, not a continuation symmetry};
\]

%dsp-level:0
\[
(3)\quad\textbf{the asymmetry:}\qquad\boxed{\ |A_3|-|S_1\times S_1|=3-1=2,\qquad2-0=2\ }\qquad[\text{machine}]:
\]

%dsp-level:0
\[
\textbf{the asymmetry is}\ 2\ \textbf{by both conventions}\ —\ \text{the order difference and the nontrivial-count difference agree};
\]

%dsp-level:0
\[
\qquad \text{the count is derived};
\]

%dsp-level:0
\[
(4)\quad\textbf{the first annihilation row (class-blindness):}
\]

%dsp-level:0
\[
\qquad \text{both baryon classes are}\ \varepsilon\text{-sealed 3-junctures with identical vertex data}\ [\mathrm{Def\text{-}BM.3}(2);\ \mathrm{BM.17};\ \mathrm{BM.26}]:
\]

%tbl-level:0
{\setlength{\tabcolsep}{4pt}
\begin{longtable}{>{\raggedright\arraybackslash}p{\dimexpr0.4919\textwidth-2\tabcolsep\relax}|>{\raggedright\arraybackslash}p{\dimexpr0.1757\textwidth-2\tabcolsep\relax}|>{\raggedright\arraybackslash}p{\dimexpr0.1757\textwidth-2\tabcolsep\relax}|>{\raggedright\arraybackslash}p{\dimexpr0.1568\textwidth-2\tabcolsep\relax}}
\toprule
vertex invariant & neutral class & unit class & difference \\
\addlinespace[2pt]
\midrule
\endfirsthead
\toprule
vertex invariant & neutral class & unit class & difference \\
\addlinespace[2pt]
\midrule
\endhead
\bottomrule
\endlastfoot
junction grammar & $\varepsilon_{ijk}$ 3-juncture & $\varepsilon_{ijk}$ 3-juncture & $0$ \\
vertex type $[\mathrm{D0.8}]$ & $(3,\allowbreak 0)$ & $(3,\allowbreak 0)$ & $0$ \\
sealing-phase class $[\mathrm{BM.2}(1);\ \mathrm{Census\ FT.49}(3)]$ & $4\mathbb Z$ & $4\mathbb Z$ & $0$ \\
$E_{\mathrm{int}}$ $[\mathrm{Census\ FT.35}(2)]$ & $6$ & $6$ & $0$ \\
continuation group & $A_3$, order $3$ & $A_3$, order $3$ & $0$ \\
sealing grammar count $[\mathrm{CT3.17a}(2)]$ & $g(3)=\allowbreak 1$ & $g(3)=\allowbreak 1$ & $0$ \\
\end{longtable}
}

%dsp-level:0
\[
\textbf{every vertex-structural quantity is a common factor on the two classes}\ [\text{machine: the difference table is identically zero}]\implies\textbf{no vertex datum can split the two rows};
\]

%dsp-level:0
\[
(5)\quad\textbf{the second annihilation row (the seed is colour-blind):}
\]

%dsp-level:0
\[
\qquad T_a\mathsf v=0\ [\mathrm{EW.32}];
\]

%dsp-level:0
\[
\qquad \mathcal G=(gW_1,gW_2,gW_3,g'Y)\ [\mathrm{EW.33}]:\ \neg\exists\ \text{a colour generator in the transport data};
\]

%dsp-level:0
\[
\operatorname{Re}\langle\mathcal G_A\mathsf v,\mathcal G_B\mathsf v\rangle=\frac14\begin{pmatrix}g^2&0&0&0\\0&g^2&0&0\\0&0&g^2&-gg'\\0&0&-gg'&g'^2\end{pmatrix}\qquad[\text{machine: Pauli spinors over}\ \mathbb Q(i)\otimes\mathbb Q[g,g']]:
\]

%dsp-level:0
\[
\textbf{the free-symbol set of the Gram is}\ \{g,g'\}:\ \text{zero colour symbols}\ [\text{machine}]\implies\boxed{\ \frac{\partial\,G(c)}{\partial\,(\text{any slot-index datum})}=0\ \ \text{exactly}\ }
\]

%dsp-level:0
\[
\textbf{the exact factorizations}\ [\text{machine: polynomial identities in}\ \mathbb Q[g,g',c_1,\dots,c_4]]:
\]

%dsp-level:0
\[
G(c)=\frac14\Bigl[g^2(c_1^2+c_2^2+c_3^2)-2gg'c_3c_4+g'^2c_4^2\Bigr]=\frac14\Bigl[g^2(c_1^2+c_2^2)+\bigl(gc_3-g'c_4\bigr)^2\Bigr];
\]

%dsp-level:0
\[
\ker M^2=\operatorname{span}\{(0,0,g',g)\}\ [\mathrm{EW.22}]:
\]

%dsp-level:0
\[
\qquad \frac{(g^2+g'^2)}4\|c\|^2-G(c)=\frac14\Bigl[g'^2(c_1^2+c_2^2)+\bigl(gc_4+g'c_3\bigr)^2\Bigr]\ge0;
\]

%dsp-level:0
\[
\det(xI-\mathrm{Gram})=\Bigl(x-\frac{g^2}4\Bigr)^2x\Bigl(x-\frac{g^2+g'^2}4\Bigr)\qquad[\text{machine: Faddeev--LeVerrier over}\ \mathbb Q[g,g',x]];
\]

%dsp-level:0
\[
(6)\quad\textbf{the third annihilation row (the trace law):}
\]

%dsp-level:0
\[
\qquad \text{every Hermitian slot operator on a sealed singlet reads only its trace:}\quad\mathfrak T(H)
\]

%dsp-level:0
\[
= 3a\ \ \text{for}\ H
\]

%dsp-level:0
\[
= aI_3+zP+z^*P^{\mathsf T}
\]

%dsp-level:0
\[
\qquad [\mathrm{BM.1}(2);\ \mathrm{FT.117}(1)]:
\]

%dsp-level:0
\[
\textbf{the trace's free-symbol set is}\ \{a\}:\ \text{the}\ C_3\text{-doublet datum}\ (B,\theta)\ \text{in}\ z
\]

%dsp-level:0
\[
= Be^{i\theta}\ \textbf{is annihilated}\ [\text{machine}];
\]

%dsp-level:0
\[
\qquad 1+\omega+\omega^2
\]

%dsp-level:0
\[
= 0\ \text{in}\ \mathbb Q[\omega]/(\omega^2{+}\omega{+}1):\ \text{the}\ C_3\ \textbf{average of the doublet datum is exactly zero}\ [\text{machine}];
\]

%dsp-level:0
\[
\textbf{the companion row:}
\]

%dsp-level:0
\[
\qquad \text{every}\ S_3\text{-covariant slot operator has charpoly}\ (x-a-3b)(x-a)^2\ [\mathrm{FT.116};\ \text{machine}]:\ \textbf{the doublet is degenerate at hadron level};
\]

%dsp-level:0
\[
(7)\quad\textbf{the multiplicity row:}
\]

%dsp-level:0
\[
\qquad \text{a continuation-symmetry count is a degeneracy multiplicity of the sealed state};
\]

%dsp-level:0
\[
\qquad \text{a uniform multiplicity factor cancels in the normalized read at every priced factor}\ \{1,2,3,15,20,\tfrac73\}\ [\mathrm{BM.28}(4)(\mathrm i);\ \mathrm{BM.30}(3);\ \text{machine}]:
\]

%dsp-level:0
\[
\text{and it multiplies the total uniformly}\ —\ \text{the same factor on both classes by}\ (4)\implies\textbf{it shifts a total and never splits a class};
\]

%dsp-level:0
\[
(8)\quad\textbf{the verdict:}\qquad\boxed{\ \Delta m_{\mathrm{vertex}}\ =\ 0\ \ \text{exactly}\ }\qquad[(4)\wedge(5)\wedge(6)\wedge(7);\ \text{machine}]:
\]

%dsp-level:0
\[
\textbf{the vertex-asymmetry hypothesis is refuted as a mass-splitting mechanism};
\]

%dsp-level:0
\[
\qquad \text{the asymmetry itself}\ (3)\ \textbf{is derived and stands};
\]

%dsp-level:0
\[
(9)\quad\textbf{the residual content (where the asymmetry does act):}
\]

%dsp-level:0
\[
\qquad A_3\ \text{is the forced family group of the seal}\ [\mathrm{BM.3}(1)];
\]

%dsp-level:0
\[
\qquad \text{its orbit on a transport profile has length}\ 3\ [\mathrm{Def\text{-}BM.11}(2),(3);\ \mathrm{FT.158}(1):\ U^3=E_0,\ U\ne E_0,\ U^2\ne E_0]:
\]

%dsp-level:0
\[
\textbf{the asymmetry generates the family triplet}\ —\ \text{three lines per channel}\ [\mathrm{BM.60}(1)];
\]

%dsp-level:0
\[
\qquad \textbf{it does not split the baryon classes};
\]

%dsp-level:0
\[
\qquad \text{the}\ (1,1)\text{-vertex's single continuation is the chain-wall telescoping}\ [\mathrm{EM.9b}]\ —\ \text{and its content is}\ (2)\text{'s parity law, consumed at}\ \mathrm{BM.59};
\]

%dsp-level:0
\[
(10)\quad\textbf{status:}
\]

%dsp-level:0
\[
\qquad \textbf{derived:}\ (1)\text{--}(9);
\]

%dsp-level:0
\[
\qquad \textbf{refuted:}\ \text{the vertex asymmetry as a splitting mechanism}\ (8);
\]

%dsp-level:0
\[
\qquad \neg\text{a P4 addition},\ \neg\text{a measured value}\ [\mathrm{P5}].
\]

\textbf{Req.} A7 (Part 1); D0.8, D2.6, D3.3, D3.4, D3.9 (Parts 1–2); CT0.9, CT3.17, CT3.17a(2), EW.15, EW.20, EW.22, EW.32, EW.33, EW.35, EM.9, EM.9a, EM.9b, EM.24(1) (Part 2); P5 (Parts 1–2); Def-H.37(1), FT.106(3a), FT.116, FT.117(1)–(3), FT.158(1) (Part 3); Census FT.35(2), FT.49(3) (row-wise); Def-BM.1, Def-BM.3(2), Def-BM.11(2), Def-BM.11(3), BM.1(2), BM.2(1), BM.3(1), BM.17, BM.26, BM.28(4)(i), BM.30(3), BM.36(4), BM.36(5) (this part); the exact Levi-Civita and Kronecker invariant censuses, Pauli/Gram algebra over $\mathbb Q(i)\otimes\mathbb Q[g,g']$, Faddeev–LeVerrier charpolys, the cyclotomic ring $\mathbb Q[\omega]/(\omega^2{+}\omega{+}1)$, sum-of-squares factorizations [math; machine: \texttt{checks/\allowbreak{}verify\_\allowbreak{}H\_\allowbreak{}constituent\_\allowbreak{}transport\_\allowbreak{}splitting.py} §A].

\textbf{Pf.}

%dsp-level:0
\[
(1):\ \varepsilon_{ijk}\ \text{built as the sign of the index triple};
\]

%dsp-level:0
\[
\ \text{all six}\ \sigma
\]

%dsp-level:0
\[
\in S_3\ \text{tested against}\ \varepsilon_{\sigma(i)\sigma(j)\sigma(k)}
\]

%dsp-level:0
\[
= \varepsilon_{ijk}:\ \text{the invariant set has three elements and equals}\ \{(012),(120),(201)\}\ [\text{machine}];
\]

%dsp-level:0
\[
\ \operatorname{sgn}(\sigma)
\]

%dsp-level:0
\[
= +1\ \text{on exactly those three}\ [\text{machine}]
\]

%dsp-level:0
\[
(2):\ \tau(v)
\]

%dsp-level:0
\[
= (1,1)\ [\mathrm{D0.8}]\implies\text{the type-preserving relabelling group is}\ S_1\times S_1;
\]

%dsp-level:0
\[
\ \delta^i{}_j\ \text{symmetric in its indices}\ [\text{machine}];
\]

%dsp-level:0
\[
\ \text{the swap row: all four end pairs mapped,}\ q
\]

%dsp-level:0
\[
\mapsto -q\ \text{at each}\ [\text{machine}];
\]

%dsp-level:0
\[
\ \mathrm{EM.9a}\text{'s charge-conjugation oddness re-cited}
\]

%dsp-level:0
\[
(3):\ 3-1=2,\quad2-0=2\ [\text{machine: integers}]
\]

%dsp-level:0
\[
(4):\ \text{each row re-cited from its own tag and compared across the two classes}\ [\text{machine: the difference table}];
\]

%dsp-level:0
\[
\ \mathrm{BM.28}(6)\text{'s identical set re-cited}
\]

%dsp-level:0
\[
(5):\ \mathsf v
\]

%dsp-level:0
\[
= (0,1)^{\mathsf T}:\ \sigma_1\mathsf v
\]

%dsp-level:0
\[
= (1,0)^{\mathsf T},
\]

%dsp-level:0
\[
\ \sigma_2\mathsf v
\]

%dsp-level:0
\[
= (-i,0)^{\mathsf T},
\]

%dsp-level:0
\[
\ \sigma_3\mathsf v
\]

%dsp-level:0
\[
= (0,-1)^{\mathsf T},
\]

%dsp-level:0
\[
\ \sigma_0\mathsf v
\]

%dsp-level:0
\[
= \mathsf v;
\]

%dsp-level:0
\[
\ W_3\mathsf v
\]

%dsp-level:0
\[
= -\tfrac12\mathsf v,
\]

%dsp-level:0
\[
\ Y\mathsf v
\]

%dsp-level:0
\[
= +\tfrac12\mathsf v,
\]

%dsp-level:0
\[
\ Q\mathsf v
\]

%dsp-level:0
\[
= 0\ [\text{machine}];
\]

%dsp-level:0
\[
\ \text{the Gram block and its antisymmetric imaginary part re-derived}\ [\text{machine}]
\]

%dsp-level:0
\[
\text{the kernel:}\ \mathrm{Gram}\cdot(0,0,g',g)^{\mathsf T}=0\ [\text{machine}];
\]

%dsp-level:0
\[
\qquad G(c)
\]

%dsp-level:0
\[
= \tfrac14[g^2(c_1^2+c_2^2)+(gc_3-g'c_4)^2]\ \text{and the ceiling's SOS form expanded and compared in}\ \mathbb Q[g,g',c_1,\dots,c_4]\ [\text{machine}]
\]

%dsp-level:0
\[
\det(xI-\mathrm{Gram})\ \text{by Faddeev--LeVerrier and compared to the factored form in}\ \mathbb Q[g,g',x]\ [\text{machine}];
\]

%dsp-level:0
\[
\qquad \text{the symbol census of every Gram entry}\ [\text{machine}]
\]

%dsp-level:0
\[
(6):\ H
\]

%dsp-level:0
\[
= aI_3+zP+z^*P^{\mathsf T}\ \text{with}\ P\ \text{the 3-cycle}:\ \operatorname{tr}P
\]

%dsp-level:0
\[
= \operatorname{tr}P^{\mathsf T}
\]

%dsp-level:0
\[
= 0\implies\mathfrak T(H)
\]

%dsp-level:0
\[
= 3a\ [\text{machine over}\ \mathbb Q[a,\operatorname{Re}z,\operatorname{Im}z]];
\]

%dsp-level:0
\[
\qquad \omega^3=1,\ 1+\omega+\omega^2=0\ \text{in}\ \mathbb Q[\omega]/(\omega^2{+}\omega{+}1)\ [\text{machine}]
\]

%dsp-level:0
\[
aI_3+bJ_3:\ \text{charpoly}\ (x-a-3b)(x-a)^2\ \text{by Faddeev--LeVerrier over}\ \mathbb Q[a,b,x]\ [\text{machine}]
\]

%dsp-level:0
\[
(7):\ \frac{\sum_Rf\,w\,\Delta S}{\sum_Rf\,w}
\]

%dsp-level:0
\[
= \frac{\sum_Rw\,\Delta S}{\sum_Rw}\ \text{at every uniform}\ f\in\{1,2,3,15,20,\tfrac73\}\ [\text{machine: exact Fractions}];
\]

%dsp-level:0
\[
\qquad (8):\ \text{conjunction of}\ (4),(5),(6),(7)\ [\text{machine}]
\]

%dsp-level:0
\[
(9):\ U^3
\]

%dsp-level:0
\[
= E_0,
\]

%dsp-level:0
\[
\ U
\]

%dsp-level:0
\[
\ne E_0,
\]

%dsp-level:0
\[
\ U^2
\]

%dsp-level:0
\[
\ne E_0\ \text{for the 3-cycle permutation matrix over}\ \mathbb Q(i):\ \text{the orbit length is exactly}\ 3\ [\text{machine}];
\]

%dsp-level:0
\[
\qquad (10):\ \text{assembly}. \ \blacksquare
\]

\textbf{Phys.} The hypothesis was that the strong vertex, having two continuation symmetries against the electromagnetic vertex's one, would split the quark masses. The count is right: the Levi-Civita tensor is preserved by exactly the three even permutations of its slots, so the strong 3-vertex carries two nontrivial continuations, while the chain wall's single in-leg and single out-leg admit no nontrivial relabelling at all — its only candidate, the leg swap, reverses the end registers and is therefore the charge-conjugation involution, not a symmetry of the transport. The asymmetry is two by either counting convention. And it splits nothing. Three independent theorems kill it: the two baryon classes are the same vertex with the same phase class and the same edge count, so any vertex invariant is a common factor; the transported seed is annihilated by every colour generator, so the mass form has no slot index to differentiate; and on a sealed singlet every slot operator reads only its trace, which annihilates the doublet datum the continuation group acts on. What the asymmetry does buy is the family: the same $C_3$ that cannot split the neutron from the proton is what cycles a transport profile into three lines.

\pseudohead{theorem}{\texorpdfstring{BM.58}{BM.58}}{\texorpdfstring{The observation frames: the metric is charge-blind, and the frame's content is the seal asymmetry}{The observation frames: the metric is charge-blind, and the frame's content is the seal asymmetry}}\label{tag:BM-58}
\textbf{Thm.}

%dsp-level:0
\[
(1)\quad\textbf{the metric data carry no charge argument:}
\]

%dsp-level:0
\[
\qquad \Omega=e^{\alpha_\Omega a_{00}},\quad g_{00}=-\Omega^2,\quad g_{0i}=\Omega^2\varkappa\beta w_i,\quad g_{ij}=\Omega^2(\delta_{ij}+\varkappa q_{ij})
\]

%dsp-level:0
\[
\qquad [\mathrm{CT1.12};\ \mathrm{FT.157}(1)]:
\]

%dsp-level:0
\[
N^2=-\frac{\det g}{\det\gamma}>0,
\]

%dsp-level:0
\[
\qquad N_i=g_{0i}=\Omega^2\varkappa\beta w_i,
\]

%dsp-level:0
\[
\qquad N^i=\gamma^{ij}N_j,
\]

%dsp-level:0
\[
\qquad ds^2=-N^2d\tau^2+g_{ij}(dx^i+N^id\tau)(dx^j+N^jd\tau)
\]

%dsp-level:0
\[
\qquad [\mathrm{Def\text{-}BM.10}(2);\ \text{machine}]:
\]

%dsp-level:0
\[
\textbf{the token census on the canonical defining blocks:}
\]

%dsp-level:0
\[
\qquad \text{charge-token occurrences inside}\ \mathrm{Def\text{-}BM.10}(2),(3),(5)\ \text{and}\ \mathrm{CT1.12}:\ \ \mathbf 0
\]

%dsp-level:0
\[
\qquad [\text{machine}];
\]

%dsp-level:0
\[
(2)\quad\textbf{the seven geometric monomials are charge-blind:}
\]

%dsp-level:0
\[
\qquad \Delta x_{\mathrm B}=\mathfrak h E^{-1}(\ln2)^{-1},\quad V_{\mathrm{eff}}=\mathfrak h^3E^{-3}(\ln2)^{-3},\quad\rho=\mathfrak h^{-4}E^4(\ln2)^4,
\]

%dsp-level:0
\[
\kappa^2\rho=2\pi\,\mathfrak h^{-2}E^4(\ln2)^3,
\]

%dsp-level:0
\[
\qquad r_s=\tfrac12\mathfrak h E,
\]

%dsp-level:0
\[
\qquad (\Delta x_{\mathrm B})^2\kappa^2\rho=2\pi E^2\ln2,
\]

%dsp-level:0
\[
\qquad \Delta x_{\mathrm B}\,r_s=\frac{\mathfrak h^2}{2\ln2}
\]

%dsp-level:0
\[
\qquad [\mathrm{BM.13}(1),\ \mathrm{BM.14}(1)\text{--}(3),\ \mathrm{BM.15}(4),\ \mathrm{BM.16}(1),(4)]:
\]

%dsp-level:0
\[
\textbf{the charge exponent is}\ 0\ \textbf{in all seven}\ [\text{machine: the Laurent-monomial census}];
\]

%dsp-level:0
\[
\qquad \textbf{the lock identities hold as monomial identities}\ [\text{machine}];
\]

%dsp-level:0
\[
\textbf{at}\ E_{\mathrm{int}}=6:
\]

%dsp-level:0
\[
\qquad V_{\mathrm{eff}}=\frac{\mathfrak h^3}{216(\ln2)^3},\quad\rho=\frac{1296(\ln2)^4}{\mathfrak h^4},\quad\kappa^2\rho=\frac{2592\pi(\ln2)^3}{\mathfrak h^2},
\]

%dsp-level:0
\[
\quad r_s=3\mathfrak h,
\]

%dsp-level:0
\[
\quad (\Delta x_{\mathrm B})^2\kappa^2\rho=72\pi\ln2
\]

%dsp-level:0
\[
\qquad [\mathrm{BM.43}(1);\ \text{machine: every coefficient exact}];
\]

%dsp-level:0
\[
(3)\quad\textbf{the frames are masks, not metrics:}
\]

%dsp-level:0
\[
\qquad \textbf{the colour frame:}\ \mathcal C(\Sigma;C)\in\{0,1\},
\]

%dsp-level:0
\[
\ \text{an idempotent, and one dangling attachment zeroes the whole colour sector}\ [\mathrm{D2.14};\ \mathrm{FT.19};\ \mathrm{BM.1}(1);\ \mathrm{BM.32}(1)]\ —\ \text{the only idempotents of}\ \{0,\dots,7\}\ \text{are}\ \{0,1\}\ [\text{machine}];
\]

%dsp-level:0
\[
\textbf{the EM frame:}
\]

%dsp-level:0
\[
\qquad W(Q)>0\ \ \text{at every local charge class}\ Q\in\{-3,\dots,+3\}\ [\mathrm{Def\text{-}BM.8}(6);\ \text{machine}]:
\]

%dsp-level:0
\[
\qquad \neg\exists\ \text{a local projector on}\ Q;
\]

%dsp-level:0
\[
\quad\text{neutrality is exhausted only at universe scale}\ [\mathrm{EM.11};\ \mathrm{BM.47}(4),(5)]:
\]

%dsp-level:0
\[
\textbf{a mask acts on operators and on the probability space, not on the metric}\ [\mathrm{D2.11};\ \mathrm{Def\text{-}BM.10}(2)\text{'s alphabet},\ (1)]:
\]

%dsp-level:0
\[
\qquad \text{the emergent metric is the ambient ledger's coarse-grained sector coefficients, and no restriction re-reads them};
\]

%dsp-level:0
\[
(4)\quad\textbf{the metric correction vanishes:}
\]

%dsp-level:0
\[
\qquad \boxed{\ \Delta m_{\mathrm{metric}}\ =\ 0\ \ \text{exactly}\ }
\]

%dsp-level:0
\[
\qquad [(1)\wedge(2)\wedge(3);\ \text{machine: the class difference of each of the seven monomials is}\ 0]:
\]

%dsp-level:0
\[
\textbf{the observation-frame hypothesis is refuted as a mass-splitting mechanism};
\]

%dsp-level:0
\[
\qquad \text{the parent and the daughter carry identical curvature rows}\ [\mathrm{BM.43}(1):\ \text{re-cited}];
\]

%dsp-level:0
\[
(5)\quad\textbf{what the frame asymmetry does carry:}
\]

%dsp-level:0
\[
\qquad \text{colour self-seals locally; EM seals at universe scale}\ [\mathrm{BM.47}(5)];
\]

%dsp-level:0
\[
\qquad \text{the surviving mass content is the scar's field energy}
\]

%dsp-level:0
\[
E^{\mathrm{self}}_{\mathrm{EM}}=\frac{c\,e^2q^2}{4\pi\,\Delta x_{\mathrm B}}=c\,\varepsilon_{\mathrm{EM}}\,q^2\,m_{\mathrm{hid}},
\]

%dsp-level:0
\[
\qquad \Delta m_{\mathrm{EM}}=\frac{3c\,e^2\ln2}{2\pi\,\mathfrak h}>0,
\]

%dsp-level:0
\[
\qquad \textbf{the charged member the heavier one}
\]

%dsp-level:0
\[
\qquad [\mathrm{BM.46}(1),(4);\ \mathrm{Def\text{-}BM.12}(5)]:
\]

%dsp-level:0
\[
\text{the monopole selector is}\ Q^2
\]

%dsp-level:0
\[
= (0,1,1)\ \text{at}\ Q
\]

%dsp-level:0
\[
= (0,\pm1)\ [\text{machine}]\implies\boxed{\ \text{the frame's contribution to}\ m_n-m_p\ \text{is}\ -\,c\,\varepsilon_{\mathrm{EM}}\,m_{\mathrm{hid}}\ <\ 0\ }
\]

%dsp-level:0
\[
\qquad [\text{machine: both priced}\ c]:
\]

%dsp-level:0
\[
\textbf{the frame asymmetry works against the neutral class's excess}\ —\ \text{so the missing piece must be class-dependent, of the opposite sign, and not electromagnetic};
\]

%dsp-level:0
\[
(6)\quad\textbf{the uniqueness row (the force direction's output):}
\]

%dsp-level:0
\[
\qquad \mathrm{BM.52}(3)\text{'s sector census is complete on the functionals it enumerates}\ (\Delta S,\ E_{\mathrm{int}},\ \text{the Janus balance},\ \text{the geometric monomials},\ |w|^2,\ \text{the Casimir},\ \text{the locks})\ \text{and gives five zeros and one monomial}\ [\mathrm{BM.52}(4)];
\]

%dsp-level:0
\[
\qquad \textbf{the complete inventory of class-dependent derived data at a sealed 3-juncture, extended by the register functionals the census does not enumerate, is}
\]

%tbl-level:2
{\footnotesize\setlength{\tabcolsep}{4pt}
\begin{longtable}{>{\raggedright\arraybackslash}p{\dimexpr0.3065\textwidth-2\tabcolsep\relax}|>{\raggedright\arraybackslash}p{\dimexpr0.1371\textwidth-2\tabcolsep\relax}|>{\raggedright\arraybackslash}p{\dimexpr0.3790\textwidth-2\tabcolsep\relax}|>{\raggedright\arraybackslash}p{\dimexpr0.1774\textwidth-2\tabcolsep\relax}}
\toprule
datum & species & status & sign for $m_n-m_p$ \\
\addlinespace[2pt]
\midrule
\endfirsthead
\toprule
datum & species & status & sign for $m_n-m_p$ \\
\addlinespace[2pt]
\midrule
\endhead
\bottomrule
\endlastfoot
the vertex invariants & slot / grammar & class-blind $[\mathrm{BM.57}(4)]$ & $0$ \\
the colour data & slot & annihilated by the trace law $[\mathrm{BM.1}(2);\ \mathrm{FT.117}]$ & $0$ \\
the metric monomials & chart / metric & charge exponent $0$ $[(2)]$ & $0$ \\
$F_{LR}=\allowbreak qI_2$, the chiral splitting $2\|q\|$ & operator & trace-free within a class: $\sum_{s,t}tq=\allowbreak 0$ $[\mathrm{BM.22}(2)\text{--}(4)]$ & $0$ on a trace read $[\mathrm{BM.1}(2)]$ \\
$Z^{\mathrm{amp}}(0,\allowbreak S)/Z^{\mathrm{amp}}(\pm1,\allowbreak S)=\allowbreak \tfrac{58}{13},\allowbreak \tfrac{2398}{401}$ & amplitude & barred: amplitude $\ne$ mass $[\mathrm{FT.79}(5);\ \mathrm{BM.37}(3)(\mathrm b)]$ & barred \\
$E^{\mathrm{self}}_{\mathrm{EM}}\propto Q^2$ & chart / field & adopted $[\mathrm{Def\text{-}BM.12}(6)]$ & negative $[(5)]$ \\
$W(Q)=\allowbreak (1,\allowbreak 6,\allowbreak 15,\allowbreak 20,\allowbreak 15,\allowbreak 6,\allowbreak 1)$ & transport register & derived, unbarred $[\mathrm{Def\text{-}BM.8}(6);\ \mathrm{FT.106}(3\mathrm a)]$ & \textbf{positive} $[\mathrm{BM.59}]$ \\
\end{longtable}
}

%dsp-level:0
\[
\text{the six sector rows of}\ \mathrm{BM.52}(3)\ \text{and the Lyapunov row of}\ \mathrm{BM.51}(6)\ \text{are re-cited inside this table as its first four entries};
\]

%dsp-level:0
\[
\qquad \textbf{exactly one candidate survives with the required sign};
\]

%dsp-level:0
\[
[\text{machine: the trace row}\ \textstyle\sum_{s,t\in\{\pm1\}}tq=0;\ \text{the row census exhaustive against the tier's own inventory scan of}\ \mathrm{BM.22},\ \mathrm{BM.28}(5),(6),\ \mathrm{BM.30}(1),\ \mathrm{BM.46}(2),(4),\ \mathrm{BM.51}(6),\ \mathrm{BM.52}(3),\ \mathrm{Def\text{-}BM.8}(6)];
\]

%dsp-level:0
\[
(7)\quad\textbf{status:}
\]

%dsp-level:0
\[
\qquad \textbf{derived:}\ (1)\text{--}(6);
\]

%dsp-level:0
\[
\qquad \textbf{refuted:}\ \text{the observation frame as a metric-splitting mechanism}\ (4);
\]

%dsp-level:0
\[
\qquad \neg\text{a P4 addition},\ \neg\text{a measured range, radius or flux}\ [\mathrm{P5}].
\]

\textbf{Req.} D1.3, D2.1, D2.9, D2.11, D2.14 (Parts 1–2); CT1.12, CT1.22, CT2.6, EM.11, EW.22 (Part 2); P5 (Parts 1–2); Def-H.32(2), FT.19, FT.79(5), FT.106(3), FT.117(1)–(3), FT.157(1) (Part 3); Def-BM.4, Def-BM.8(6), Def-BM.10(2), Def-BM.10(3), Def-BM.10(5), Def-BM.12(5), Def-BM.12(6), BM.1(1), BM.1(2), BM.13(1), BM.14(1)–(3), BM.15(4), BM.16(1), BM.16(4), BM.22(2)–(4), BM.28(1), BM.32(1), BM.37(3), BM.43(1), BM.46(1), BM.46(4), BM.47(4), BM.47(5) (this part); the ADM split by Schur complement, Laurent-monomial exponent arithmetic, idempotent censuses, token censuses [math; machine: \texttt{checks/\allowbreak{}verify\_\allowbreak{}H\_\allowbreak{}constituent\_\allowbreak{}transport\_\allowbreak{}splitting.py} §B].

\textbf{Pf.}

%dsp-level:0
\[
(1):\ \mathrm{Def\text{-}BM.10}(2)\text{'s split re-derived as linear algebra on}\ \mathrm{CT1.12}\text{'s dictionary};
\]

%dsp-level:0
\[
\ N^2
\]

%dsp-level:0
\[
= -\det g/\det\gamma\ \text{by Schur, denominators cleared}\ [\text{machine}];
\]

%dsp-level:0
\[
\ \text{the line-element identity expanded and compared}\ [\text{machine}];
\]

%dsp-level:0
\[
\ \text{the token census on the two defining blocks}\ [\text{machine}]
\]

%dsp-level:0
\[
(2):\ \text{each monomial built from its own tag's formula as a Laurent monomial in}\ (\mathfrak h,E,\ln2,\pi,Q);
\]

%dsp-level:0
\[
\ \text{the}\ Q\text{-exponent censused}\ [\text{machine}];
\]

%dsp-level:0
\[
\ (\Delta x_{\mathrm B})^2\kappa^2\rho\ \text{and}\ \Delta x_{\mathrm B}r_s\ \text{multiplied out and compared}\ [\text{machine}];
\]

%dsp-level:0
\[
\ E
\]

%dsp-level:0
\[
= 6\ \text{substituted:}\ 6^3
\]

%dsp-level:0
\[
= 216,
\]

%dsp-level:0
\[
\ 6^4
\]

%dsp-level:0
\[
= 1296,
\]

%dsp-level:0
\[
\ 2\pi\cdot1296
\]

%dsp-level:0
\[
= 2592\pi,
\]

%dsp-level:0
\[
\ \tfrac12\cdot6
\]

%dsp-level:0
\[
= 3,
\]

%dsp-level:0
\[
\ 2\pi\cdot36
\]

%dsp-level:0
\[
= 72\pi\ [\text{machine}]
\]

%dsp-level:0
\[
(3):\ x^2=x\ \text{over}\ \{0,\dots,7\}\ \text{gives}\ \{0,1\}\ [\text{machine: exhaustive}];
\]

%dsp-level:0
\[
\qquad W(Q)\ \text{re-enumerated and positive at all seven classes}\ [\text{machine}]
\]

%dsp-level:0
\[
(4):\ \text{the class difference of each monomial of}\ (2)\ \text{is}\ 0\ \text{because the}\ Q\text{-exponent is}\ 0\ [\text{machine}]
\]

%dsp-level:0
\[
(5):\ \mathrm{BM.46}(1),
\]

%dsp-level:0
\[
(4)\ \text{re-cited};
\]

%dsp-level:0
\[
\ Q^2\ \text{at}\ (0,\pm1)
\]

%dsp-level:0
\[
= (0,1,1)\ [\text{machine}];
\]

%dsp-level:0
\[
\ \text{the sign of}\ -(1-0)\,
\]

%dsp-level:0
\[
c\,
\]

%dsp-level:0
\[
\varepsilon_{\mathrm{EM}}\,
\]

%dsp-level:0
\[
m_{\mathrm{hid}}\ \text{at}\ c\in\{\tfrac12,\tfrac35\},
\]

%dsp-level:0
\[
\ \varepsilon_{\mathrm{EM}}>0,
\]

%dsp-level:0
\[
\ m_{\mathrm{hid}}>0\ [\text{machine}]
\]

%dsp-level:0
\[
(6):\ \text{the inventory scan of the class-dependent rows of}\ \mathrm{BM.22},
\]

%dsp-level:0
\[
\ \mathrm{BM.28}(5),
\]

%dsp-level:0
\[
(6),
\]

%dsp-level:0
\[
\ \mathrm{BM.30}(1),
\]

%dsp-level:0
\[
\ \mathrm{BM.46}(2),
\]

%dsp-level:0
\[
(4),
\]

%dsp-level:0
\[
\ \mathrm{Def\text{-}BM.8}(6)\ [\text{machine: the row set is exhaustive}];
\]

%dsp-level:0
\[
\qquad \text{the force-matrix trace:}\ \operatorname{spec}
\]

%dsp-level:0
\[
= \{\tfrac12(c_0+s|\mathbf c|)+tq\},
\]

%dsp-level:0
\[
\ \sum_{s,t}[\tfrac12(c_0+s|\mathbf c|)+tq]
\]

%dsp-level:0
\[
= 2c_0\ [\mathrm{BM.22}(3);\ \text{machine}]\implies\text{the trace read is charge-blind}\ [\mathrm{BM.1}(2)]
\]

%dsp-level:0
\[
(7):\ \text{assembly}. \ \blacksquare
\]

\textbf{Phys.} The intuition was that seeing the world through photons is not seeing it through colour, and that the difference should show up in the geometry. It does not. The emergent metric is built from the ambient ledger's coarse-grained sector coefficients, and not one of its seven derived monomials carries a charge exponent: the width, the volume, the density, the curvature scalar, the Schwarzschild-type read and both locks come out identical on the two classes to the last Laurent monomial. What a frame really is, in this corpus, is a mask — and the two masks have different grammars. Colour's is an idempotent that annihilates a sector at the first dangling attachment; electromagnetism's is a census that lets every local charge class through and cancels only once, at the scale of the whole universe. Masks filter operators and probability spaces. They do not re-read the metric. So the metric correction is exactly zero, and the frame's one surviving mass contribution is the scar's own field energy — positive, charge-selected, and on the wrong side: it makes the charged member heavier. That is the useful output. It narrows the search to a single column of the inventory, and the column has one entry.

\begin{theorem}[{\texorpdfstring{BM.59 — The constituent transport splitting: the order residue, the two directions, and the exact positive sign}{BM.59 - The constituent transport splitting: the order residue, the two directions, and the exact positive sign}}]\label{tag:BM-59}
\textbf{Thm.}

%dsp-level:0
\[
(1)\quad\textbf{the missing piece, stated:}
\]

%dsp-level:0
\[
\qquad \mathrm{BM.55}(1)\text{'s five-term decomposition is derived term by term:}
\]

%dsp-level:0
\[
\qquad \Delta m_{\mathrm{Lyapunov}}
\]

%dsp-level:0
\[
= 0\ [\mathrm{BM.51}(6)];
\]

%dsp-level:0
\[
\quad\Delta m_{\mathrm{metric}}
\]

%dsp-level:0
\[
= 0\ [\mathrm{BM.53}(2)];
\]

%dsp-level:0
\[
\quad\Delta m_{\mathrm{binding}}
\]

%dsp-level:0
\[
= 0\ [\mathrm{BM.54}(2),(5)];
\]

%dsp-level:0
\[
\quad\Delta m_{\mathrm{force}}
\]

%dsp-level:0
\[
= \Delta m_{\mathrm{self\text{-}int}}
\]

%dsp-level:0
\[
= -\Delta m_{\mathrm{EM}}\ [\mathrm{BM.52}(4),(5)]:
\]

%dsp-level:0
\[
\Delta m=m_n-m_p=-\,\frac{3\,c\,e^2\ln2}{2\pi\,\mathfrak h}<0,\qquad\frac{m_n}{m_p}=\frac{4\pi}{4\pi+c\,e^2}\in(0,1)\qquad[\mathrm{BM.55}(1)\text{--}(3);\ \text{frozen}]:
\]

%dsp-level:0
\[
\textbf{the order residue}\ [\mathrm{BM.55}(4):\ \text{named, not filled}]\ :
\]

%dsp-level:0
\[
= \ \text{a class-dependent constituent transport splitting with}\ m^{(n)}_{\mathrm{constituent}}>m^{(p)}_{\mathrm{constituent}};
\]

%dsp-level:0
\[
\boxed{\ \textbf{the residue is fillable, and this tier fills it:}\qquad\Delta m_{\mathrm{transport}}=\frac{\Sigma_{\mathrm{cl}}(0)-\Sigma_{\mathrm{cl}}(+1)}{\mathfrak h}=\frac{\ln(4/3)}{\mathfrak h}>0\ }
\]

%dsp-level:0
\[
\qquad [\mathrm{Def\text{-}BM.13}(4);\ \text{machine}];
\]

%dsp-level:0
\[
(1\mathrm a)\quad\textbf{the four grounds of}\ \mathrm{BM.55}(4)\ \textbf{and how this object evades each:}\qquad\begin{array}{l|l}\text{the ground on record} & \text{the evasion}\\\hline\text{the fractional constituent alphabet is unreachable}\ [\mathrm{BM.28}(7)] & \Sigma_{\mathrm{cl}}\ \text{is built from the integer alphabet}\ \{-1,0,+1\}\ \text{alone}\ [\mathrm{Def\text{-}BM.13}(1),(3)];\ \neg\text{a fraction appears}\\\text{the species labels carry no datum}\ [\mathrm{Def\text{-}BM.3}] & \Sigma_{\mathrm{cl}}\ \text{carries no label}: it is a census of the derived read\ [\mathrm{Def\text{-}BM.8}(6)]\\\text{the deficit/inertial conversions are item-7 content} & \Sigma_{\mathrm{cl}}\ \text{is already in the informational currency}\ —\ \text{nats of binary registration}\ [\mathrm{Def\text{-}BM.13}(6)];\ \neg\text{a conversion is used}\\\text{no clause attaches constituent profiles to the charge classes} & \Sigma_{\mathrm{cl}}\ \text{is}\ \neg\text{a profile:}\ m_q\ \text{is attached to each constituent by}\ q_{\mathrm{core}}\ \text{itself}\ [\mathrm{Def\text{-}BM.8}(1);\ \mathrm{EM.9}]\end{array}
\]

%dsp-level:0
\[
\textbf{the residue is filled by the transport register, not by the transport profile:}
\]

%dsp-level:0
\[
\qquad \mathrm{Def\text{-}BM.11}\text{'s family-profile map remains channel-line data and stays unattached to the baryon rows}\ [\mathrm{BM.55}(4)\text{'s fourth ground stands untouched}];
\]

%dsp-level:0
\[
(2)\quad\textbf{the mass direction (what the spectrum requires):}
\]

%dsp-level:0
\[
\qquad \text{a term}\ \delta m(Q)\ \text{that reverses the order must satisfy four exact bars}
\]

%dsp-level:0
\[
(\mathrm i)\ \ \text{class-dependent:}\quad\delta m(0)\ne\delta m(+1);
\]

%dsp-level:0
\[
\qquad (\mathrm{ii})\ \ \text{positive for the neutral class:}\quad\delta m(0)-\delta m(+1)>0;
\]

%dsp-level:0
\[
(\mathrm{iii})\ \ \text{currency-matched:}\quad[\delta m]=\mathfrak q^{+1/2}\ \text{through the rest slot}\ [\mathrm{Def\text{-}BM.10}(4);\ \mathrm{CT2.1}];
\]

%dsp-level:0
\[
\qquad (\mathrm{iv})\ \ \text{alphabet-disjoint from the count:}\ \neg\text{a function of}\ E_{\mathrm{int}}\ \text{alone}\ [\mathrm{BM.28}(1)\ \text{scope}];
\]

%dsp-level:0
\[
\textbf{and it must not be a barred route}\ [\mathrm{BM.37}(6)];
\]

%dsp-level:0
\[
\qquad [\text{machine: the inventory scan of}\ \mathrm{BM.58}(6)\text{'s table against}\ (\mathrm i)\text{--}(\mathrm{iv})]:
\]

%dsp-level:0
\[
\boxed{\ \textbf{the unique survivor is}\ W(Q)\ \textbf{read through its logarithm}\ }\qquad[\mathrm{Def\text{-}BM.13}(4),(6)];
\]

%dsp-level:0
\[
(3)\quad\textbf{the force direction (what the seal permits):}
\]

%dsp-level:0
\[
\qquad \text{colour is a local idempotent and its slot data are annihilated on a sealed singlet}\ [\mathrm{BM.57}(5),(6);\ \mathrm{BM.58}(3)];
\]

%dsp-level:0
\[
\qquad \text{EM has no local projector and admits every class with positive census weight}\ [\mathrm{BM.58}(3)]:
\]

%dsp-level:0
\[
\textbf{the only class-dependent datum a sealed junction can carry is the EM end-register census}\ \ W(Q)
\]

%dsp-level:0
\[
= [z^Q](2+z+z^{-1})^3
\]

%dsp-level:0
\[
\qquad [\mathrm{Def\text{-}BM.8}(6);\ \mathrm{BM.28}(2)];
\]

%dsp-level:0
\[
(4)\quad\textbf{convergence:}
\]

%dsp-level:0
\[
\qquad \text{the two directions land on the same object};
\]

%dsp-level:0
\[
\qquad \boxed{\ \Delta m_{\mathrm{transport}}\ :=\ \frac{\Sigma_{\mathrm{cl}}(0)-\Sigma_{\mathrm{cl}}(+1)}{\mathfrak h}\ =\ \frac{\ln(4/3)}{\mathfrak h}\ =\ \frac{2\ln2-\ln3}{\mathfrak h}\ >\ 0\ }
\]

%dsp-level:0
\[
\qquad [\text{machine: the certificate}\ 4>3]:
\]

%dsp-level:0
\[
[\Delta m_{\mathrm{transport}}]=\mathfrak q^{+1/2}=[\text{energy}]
\]

%dsp-level:0
\[
\qquad [\ln W\in\mathfrak q^0,\ [\mathfrak h]=\mathfrak q^{-1/2};\ \mathrm{CT2.1};\ \mathrm{D2.9};\ \text{machine: exponent arithmetic}];
\]

%dsp-level:0
\[
(5)\quad\textbf{the constituent row:}
\]

%dsp-level:0
\[
\qquad \text{the per-constituent transport entropies are}\ \ln m_q
\]

%dsp-level:0
\[
= (\ln2,\,0,\,0)\ \text{at}\ q
\]

%dsp-level:0
\[
= (0,+1,-1)\ [\mathrm{Def\text{-}BM.13}(1)]:
\]

%dsp-level:0
\[
\boxed{\ m(q_{\mathrm{core}}{=}0)-m(q_{\mathrm{core}}{=}\pm1)\ =\ \frac{\ln2}{\mathfrak h}\ >\ 0\ }
\]

%dsp-level:0
\[
\qquad [\text{machine:}\ 2>1]:
\]

%dsp-level:0
\[
\qquad \textbf{the neutral-core constituent is the heavier constituent};
\]

%dsp-level:0
\[
\textbf{the composite's splitting is not three times the constituent's:}
\]

%dsp-level:0
\[
\qquad 0<\ln\tfrac43<\ln2<3\ln2
\]

%dsp-level:0
\[
\qquad [\text{machine:}\ 3<4<6<24]:
\]

%dsp-level:0
\[
\text{the additive charge constraint couples the three registers, and the coupling strictly attenuates the constituent hierarchy};
\]

%dsp-level:0
\[
(6)\quad\textbf{the parity lock (why the register is a transport datum and not a label):}
\]

%dsp-level:0
\[
\qquad \sum_if_{v_i}=-q,
\]

%dsp-level:0
\[
\qquad \sum_i\aleph_{v_i}\equiv|q|\ (\mathrm{mod}\ 2),
\]

%dsp-level:0
\[
\qquad \sum_i\aleph_{v_i}\ge|q|
\]

%dsp-level:0
\[
\qquad [\mathrm{EM.9b};\ \mathrm{D3.3};\ \mathrm{Def\text{-}BM.13}(2)]:
\]

%dsp-level:0
\[
\textbf{the interior wall content of a constituent's transport is parity-locked to its charge core};
\]

%dsp-level:0
\[
\qquad \text{the neutral core is wall-balanced}\ (\min\textstyle\sum\aleph=0);
\]

%dsp-level:0
\[
\ \text{the charged core is wall-unbalanced}\ (\min\textstyle\sum\aleph=1)\ [\text{machine: exhaustive over every sign sequence,}\ k=2,\dots,7]:
\]

%dsp-level:0
\[
\textbf{the two neutral transport realizations are the two wall-free orientations};
\]

%dsp-level:0
\[
\qquad \text{the charged core admits none}:
\]

%dsp-level:0
\[
\text{the multiplicity}\ m_q\ \textbf{is a property of the transport, not of a naming convention}\ [\mathrm{FT.106}(3\mathrm a):\ \text{it is the norm layer's own closed form}];
\]

%dsp-level:0
\[
(7)\quad\textbf{the sign is fixed by the reference, not chosen:}
\]

%dsp-level:0
\[
\qquad \Delta S:=S-S(\mathrm{vac})\ \text{with}\ S(\mathrm{vac})=0\ [\mathrm{BM.35}(3);\ \mathrm{CT0.2}]\implies\Delta\Sigma_{\mathrm{cl}}(Q)=\ln W(Q)-\ln W(\mathrm{vac})=\ln W(Q),
\]

%dsp-level:0
\[
\qquad W(\mathrm{vac})=1;
\]

%dsp-level:0
\[
\textbf{the two competing references:}\qquad\begin{array}{l|l|l}\text{reference} & \text{class difference} & \text{status}\\\hline\text{the vacuum}\ (W{=}1) & +\ln\tfrac43 & \textbf{adopted: the corpus's own deficit rule}\ [\mathrm{BM.35}(3)]\\[1mm]\text{the unconstrained register}\ (4^3{=}64) & -\ln\tfrac43 & \text{barred: }\neg\text{a state of the composite}\ [\mathrm{Def\text{-}BM.5}(1)];\ \text{class-dependent reference}\\\text{Ansatz-B}\ (-\alpha\ln\Omega) & -\alpha\ln\tfrac43 & \text{barred: an external ansatz,}\ \alpha\ \text{undeclared}\ [\mathrm{Def\text{-}H.28}]\end{array}
\]

%dsp-level:0
\[
\textbf{exactly one reference is a state of the composite}\ [\text{machine: the admissible-reference census}];
\]

%dsp-level:0
\[
(8)\quad\textbf{the species legitimacy:}
\]

%dsp-level:0
\[
\qquad \Sigma_{\mathrm{cl}}\ \text{is}\ \neg\text{a measure functional of}\ (\Gamma,R)\ —\ \text{its alphabet is the end-sign register}\ \{\sigma_s,\sigma_t\}\ [\mathrm{D3.4};\ \mathrm{CT0.9}],
\]

%dsp-level:0
\[
\ \text{disjoint from the flicker alphabet and from the metric alphabet}\ [\mathrm{BM.37}(2)\ \text{pattern};\ \text{machine: the token census}]:
\]

%dsp-level:0
\[
\textbf{it enters through the rest slot the dispersion read already reads}\ [\mathrm{Def\text{-}BM.10}(4);\ \mathrm{BM.35}(4);\ \mathrm{BM.37}(4)],
\]

%dsp-level:0
\[
\ \text{in the}\ \mathrm{Def\text{-}BM.12}(6)\ \text{pattern of a refinement,}\ \neg\text{a new mass read};
\]

%dsp-level:0
\[
\textbf{no route ratio is read as a mass ratio:}
\]

%dsp-level:0
\[
\qquad \frac{m_{\mathrm{rest}}(0)}{m_{\mathrm{rest}}(+1)}\bigg|_{\varepsilon_{\mathrm{EM}}=0}=\frac{\ln1280}{\ln960}\ \ne\ \frac43
\]

%dsp-level:0
\[
\qquad [\text{machine}];
\]

%dsp-level:0
\[
\qquad \textbf{the strict contraction:}\quad\frac{\Sigma_{\mathrm B}}{6\ln2}<\frac13\iff3\ln\tfrac43<2\ln2\iff16<27
\]

%dsp-level:0
\[
\qquad [\text{machine}];
\]

%dsp-level:0
\[
(9)\quad\textbf{the defect-species companion (conditional, the hypothesis on its face):}
\]

%dsp-level:0
\[
\qquad \textbf{hypothesis:}\ \text{the two baryon classes differ by exactly one down}\leftrightarrow\text{up constituent substitution}\ [\text{the minimal}\ \Delta q=\pm1\ \text{step};\ \mathrm{BM.23}(4)]:
\]

%dsp-level:0
\[
\textbf{then}
\]

%dsp-level:0
\[
\qquad \Delta m\ =\ \Delta_g\ =\ \hat\mu'^2\bigl(D_S(\mathcal F_g,+)-D_S(\mathcal F_g,-)\bigr)\ >\ 0
\]

%dsp-level:0
\[
\qquad [\mathrm{Def\text{-}BM.14}(5);\ \text{machine: the three certificates}]:
\]

%dsp-level:0
\[
\Delta_1=2\hat\mu'^2,\qquad\Delta_2=(\sqrt2-1)\hat\mu'^2,\qquad\Delta_3=\bigl(2\cos\tfrac\pi{10}-\sqrt3\bigr)\hat\mu'^2;
\]

%dsp-level:0
\[
\textbf{the composition itself is}\ \neg\textbf{derived:}\ \text{a fractional constituent alphabet is structurally impossible under an integer-valued additive read}\ [\mathrm{BM.28}(7);\ \mathrm{Def\text{-}BM.8}(5)];
\]

%dsp-level:0
\[
\qquad \text{deficit}\ 1'\ \text{stands}\ [\mathrm{FT.125}(4)\ \text{pointer}];
\]

%dsp-level:0
\[
(10)\quad\textbf{status:}
\]

%dsp-level:0
\[
\qquad \textbf{derived:}\ (1)\text{--}(8);
\]

%dsp-level:0
\[
\qquad \textbf{conditional:}\ (9),\ \text{with its hypothesis on its face}\ [\text{A §0.2a}];
\]

%dsp-level:0
\[
\qquad \textbf{the splitting is exact and strictly positive};
\]

%dsp-level:0
\[
\neg\text{a P4 addition},
\]

%dsp-level:0
\[
\quad\neg\text{a new dimensionful scale}\ [\mathrm{P1}],
\]

%dsp-level:0
\[
\quad\neg\text{a route adopted}\ [\mathrm{BM.30}(5)],
\]

%dsp-level:0
\[
\quad\neg\text{a measured value}\ [\mathrm{P5}];
\]

%dsp-level:0
\[
\qquad \text{§Cal-BM stays empty}.
\]

\textbf{Req.} D0.8, D2.1, D2.9, D3.3, D3.4 (Parts 1–2); CT0.2, CT0.9, CT1.29, EM.9, EM.9b (Part 2); P1, P4, P5 (Parts 1–2); Def-H.28, Def-H.32(2), Def-H.37(2), FT.79(5), FT.106(3a), FT.125(4) (Part 3); Census FT.35(2) (row-wise); Def-BM.5(1), Def-BM.8(5), Def-BM.8(6), Def-BM.10(4), Def-BM.12(6), Def-BM.13, Def-BM.14(5), BM.20(1), BM.23(4), BM.28(1), BM.28(2), BM.28(4)(i), BM.28(7), BM.30(2), BM.30(3), BM.30(5), BM.35(3), BM.35(4), BM.37(2), BM.37(4), BM.37(6), BM.41(4), BM.43(1), BM.46(4), BM.57, BM.58 (this part); the exhaustive end-sign and sign-sequence censuses, exact $\ln$-ring identities with integer certificates, Laurent-monomial exponent arithmetic, token censuses [math; machine: \texttt{checks/\allowbreak{}verify\_\allowbreak{}H\_\allowbreak{}constituent\_\allowbreak{}transport\_\allowbreak{}splitting.py} §C].

\textbf{Pf.}

%dsp-level:0
\[
(1):\ \text{the four degenerate-or-inverted rows re-cited from their own tags};
\]

%dsp-level:0
\[
\ \text{the order residue named as the missing class-dependent rest-slot term}\ [\text{definition}];
\]

%dsp-level:0
\[
\qquad (2):\ \text{the four bars stated};
\]

%dsp-level:0
\[
\ \text{the inventory scan of}\ \mathrm{BM.58}(6)\text{'s seven rows against them}\ [\text{machine: one survivor}]
\]

%dsp-level:0
\[
(3):\ \mathrm{BM.57}(5),(6)\ \text{and}\ \mathrm{BM.58}(3)\ \text{conjoined};
\]

%dsp-level:0
\[
\qquad (4):\ \Sigma_{\mathrm{cl}}(0)-\Sigma_{\mathrm{cl}}(+1)
\]

%dsp-level:0
\[
= \ln20-\ln15
\]

%dsp-level:0
\[
= \ln\tfrac43
\]

%dsp-level:0
\[
= 2\ln2-\ln3\ [\text{machine: the prime-log ring}];
\]

%dsp-level:0
\[
\ 4>3\ [\text{machine}];
\]

%dsp-level:0
\[
\ \text{the exponent row}\ \mathfrak q^0\cdot\mathfrak q^{+1/2}
\]

%dsp-level:0
\[
= \mathfrak q^{+1/2}\ [\text{machine}]
\]

%dsp-level:0
\[
(5):\ \ln m_0-\ln m_{\pm1}=\ln2-0=\ln2\ [\text{machine}];\ 2>1\ [\text{machine}];
\]

%dsp-level:0
\[
\qquad 3<4\implies\ln\tfrac43<\ln2;\quad4<6\implies\ln\tfrac43<\ln2;\quad4<24\implies\ln\tfrac43<3\ln2\ [\text{machine: integer certificates}]
\]

%dsp-level:0
\[
(6):\ \text{exhaustive over}\ 2^k\ \text{sign sequences for}\ k
\]

%dsp-level:0
\[
= 2,
\]

%dsp-level:0
\[
\dots,
\]

%dsp-level:0
\[
7:\ \textstyle\sum f
\]

%dsp-level:0
\[
= -q,
\]

%dsp-level:0
\[
\ \textstyle\sum\aleph
\]

%dsp-level:0
\[
\equiv |q|\ (2),
\]

%dsp-level:0
\[
\ \textstyle\sum\aleph
\]

%dsp-level:0
\[
\ge |q|\ \text{at every one}\ [\text{machine}];
\]

%dsp-level:0
\[
\ \text{the minima}\ 0\ \text{and}\ 1\ \text{attained}\ [\text{machine}]
\]

%dsp-level:0
\[
(7):\ \text{the three references enumerated and their class differences computed in the}\ \ln\text{-ring}\ [\text{machine}];
\]

%dsp-level:0
\[
\ W(\mathrm{vac})
\]

%dsp-level:0
\[
= 1\ \text{for the empty register}\ [\text{machine}];
\]

%dsp-level:0
\[
\ \mathrm{Def\text{-}BM.5}(1)\text{'s boundary-data clause re-cited};
\]

%dsp-level:0
\[
\ \mathrm{Def\text{-}H.28}\text{'s external-ansatz mark pinned}
\]

%dsp-level:0
\[
(8):\ \text{the token census of the defining clause against the flicker and metric alphabets}\ [\text{machine}];
\]

%dsp-level:0
\[
\qquad \ln1280/\ln960\ne\tfrac43\iff3\ln1280\ne4\ln960\iff1280^3\ne960^4\ [\text{machine: the two integers differ}];
\]

%dsp-level:0
\[
\qquad 3\ln\tfrac43<2\ln2\iff\ln\tfrac{64}{27}<\ln4\iff64<108\ [\text{machine}]
\]

%dsp-level:0
\[
(9):\ \mathrm{Def\text{-}BM.14}(5)\text{'s three splittings re-cited with their certificates}\ 4>2,
\]

%dsp-level:0
\[
\ 2>1,
\]

%dsp-level:0
\[
\ 3<5\ [\text{machine}];
\]

%dsp-level:0
\[
\qquad (10):\ \text{assembly}. \ \blacksquare
\]

\textbf{Phys.} Worked from both ends, the answer is the same object. From the mass side: whatever reverses the two rows has to be class-dependent, positive for the neutral class, dimensionally an energy through the rest slot, and built from data the edge count does not already carry — and the inventory has exactly one unbarred candidate. From the force side: colour is a projector that annihilates its own slot data on anything sealed, electromagnetism is a census that admits every local class, so the only class-dependent datum a sealed junction can carry is the count of admissible end-sign registers — twenty against fifteen. The candidate and the survivor are the same. Taken as a logarithm and booked against the vacuum in the ledger's own currency, the filter becomes an energy: the neutral class carries $\ln(4/3)$ more transport entropy than the charged one, and the neutral-core constituent carries exactly one bit more than the charged-core constituent. The reason the register counts as transport and not as bookkeeping is the parity lock: a chain whose ends agree can run wall-free, and a chain whose ends disagree cannot — the wall content of a constituent's transport is odd or even according to its charge core, so the multiplicity is a property of how the seed is carried, not of how the state is named. The sign is not a convention either. Of three possible references only one is a state the composite can be in, and it is the one the corpus already uses for every other deficit.

\end{theorem}

\pseudohead{theorem}{\texorpdfstring{BM.60}{BM.60}}{\texorpdfstring{The six-state assembly: the generation ladder, the branch split, the Platonic floors, and the two hierarchies}{The six-state assembly: the generation ladder, the branch split, the Platonic floors, and the two hierarchies}}\label{tag:BM-60}
\textbf{Thm.}

%dsp-level:0
\[
(1)\quad\textbf{the generation ladder (the 3):}
\]

%dsp-level:0
\[
\qquad \delta_q=\frac{2}{3\cdot3\cdot3}=\frac2{27},
\]

%dsp-level:0
\[
\qquad 3\delta_q=\delta_\ell=\frac29,
\]

%dsp-level:0
\[
\qquad 16\,\delta_\nu=9\,\delta_\ell=27\,\delta_q=2
\]

%dsp-level:0
\[
\qquad [\mathrm{FT.119}(1),(2);\ \mathrm{FT.143}(5);\ \text{machine: exact Fractions}]:
\]

%dsp-level:0
\[
r^{(q)}_i=1+\sqrt2\cos\Bigl(\frac2{27}+\frac{2\pi i}3\Bigr),
\]

%dsp-level:0
\[
\qquad \sum_ir^{(q)}_i=3=\dim F_{J^\pm},
\]

%dsp-level:0
\[
\qquad \sum_i\bigl(r^{(q)}_i\bigr)^2=6,
\]

%dsp-level:0
\[
\qquad \sum_{i<j}r^{(q)}_ir^{(q)}_j=\frac32
\]

%dsp-level:0
\[
\qquad [\mathrm{FT.111}(3);\ \mathrm{FT.119}(4);\ \text{machine}];
\]

%dsp-level:0
\[
2\bigl(r^{(q)}-1\bigr)^3-3\bigl(r^{(q)}-1\bigr)-\sqrt2\cos\tfrac29=0\ \ \text{for each of the three}
\]

%dsp-level:0
\[
\qquad [\mathrm{FT.119}(4);\ \text{machine: enclosure containment at width}<2^{-90}];
\]

%dsp-level:0
\[
Q=\frac{\sum_i(r^{(q)}_i)^2}{\bigl(\sum_ir^{(q)}_i\bigr)^2}=\frac69=\frac23
\]

%dsp-level:0
\[
\qquad [\mathrm{Def\text{-}BM.6}(4);\ \mathrm{FT.158}(4);\ \text{machine: exact, phase-independent}];
\]

%dsp-level:0
\[
\textbf{the generation count is derived:}
\]

%dsp-level:0
\[
\qquad \boxed{\ \#\text{lines}=3\ }
\]

%dsp-level:0
\[
\qquad [\text{the orbit length of}\ A_3:\ U^3=E_0,\ U\ne E_0,\ U^2\ne E_0\ [\mathrm{BM.57}(9);\ \text{machine}];\qquad\sum_ir_i=3=\dim F_{J^\pm}];
\]

%dsp-level:0
\[
(2)\quad\textbf{the branch split (the 2):}
\]

%dsp-level:0
\[
\qquad \mathcal H_{\mathsf v}\cong\mathbb C^2\ [\mathrm{EW.32}];
\]

%dsp-level:0
\[
\qquad Q|_{\mathcal H_{\mathsf v}}=\tfrac12\sigma_3+\tfrac12\sigma_0=\operatorname{diag}(1,0),
\]

%dsp-level:0
\[
\qquad \operatorname{spec}Q=\{0,+1\},
\]

%dsp-level:0
\[
\qquad \omega:\ \operatorname{spec}\mapsto\{0,-1\}
\]

%dsp-level:0
\[
\qquad [\mathrm{EW.36}]:
\]

%dsp-level:0
\[
\textbf{the seed carrier is a doublet, and its weight alphabet together with the}\ \omega\textbf{-image is exactly the ternary end read}\ \{-1,0,+1\}\ [\mathrm{EW.36};\ \mathrm{Def\text{-}BM.12}(1);\ \mathrm{D3.4};\ \text{machine}];
\]

%dsp-level:0
\[
\textbf{the geometric branch:}
\]

%dsp-level:0
\[
\qquad \text{the dual exchanges}\ (V,F)\ \text{at fixed}\ E\implies\text{the two members of a dual pair share}\ m_{\mathrm{hid}}\ \text{and}\ \Pi\ \text{and split on every finer ledger}\ [\mathrm{FT.44}(1),(2)]:
\]

%dsp-level:0
\[
\text{independent windings}\ F{-}1\ (5\ \text{vs}\ 7;\ 11\ \text{vs}\ 19);
\]

%dsp-level:0
\[
\qquad \text{lattice index}\ 2^{F-1}\ (32\ \text{vs}\ 128;\ 2048\ \text{vs}\ 2^{19});
\]

%dsp-level:0
\[
\qquad g(V)\ (385\ \text{vs}\ 25;\ 292938721075\ \text{vs}\ 164395);
\]

%dsp-level:0
\[
\qquad O(2d)\ (112\ \text{vs}\ 496;\ 112\ \text{vs}\ 2032);
\]

%dsp-level:0
\[
\textbf{and on the pole side they split on every defect datum}\ [\mathrm{Census\ FT.63}(6):\ \text{the derived dual-splitting law}]:
\]

%dsp-level:0
\[
\qquad \delta_v:\ \tfrac\pi2\ \text{vs}\ \tfrac{2\pi}3;\quad\tfrac\pi5\ \text{vs}\ \tfrac\pi3;
\]

%dsp-level:0
\[
\qquad D_S:\ 2-\sqrt2\ \text{vs}\ 1;\quad2-2\cos\tfrac\pi{10}\ \text{vs}\ 2-\sqrt3;
\]

%dsp-level:0
\[
\textbf{the tetrahedron is self-dual}\ [\mathrm{FT.44}(3)]\implies\text{its branch index is the weight-ladder rung}\ m\in\{1,2\},
\]

%dsp-level:0
\[
\qquad D_S(m)=(2,\,4,\,2,\,0)\ \ \text{at}\ m=(1,2,3,4)
\]

%dsp-level:0
\[
\qquad [\mathrm{Census\ FT.63}(6);\ \mathrm{FT.60}(5);\ \text{machine}];
\]

%dsp-level:0
\[
(3)\quad\textbf{the metric back-reaction of the directed 3-vertex:}
\]

%dsp-level:0
\[
\qquad \text{the skeleton metric is piecewise flat with conical defects at the vertices}\ [\mathrm{Census\ FT.63}(4);\ \mathrm{Def\text{-}H.21}(6);\ \mathrm{FT.54}]:
\]

%dsp-level:0
\[
\delta_v=2\pi-d\,\frac{(L-2)\pi}L,
\]

%dsp-level:0
\[
\qquad \frac{\delta_v}\pi=\Bigl(1,\ \frac12,\ \frac23,\ \frac15,\ \frac13\Bigr),
\]

%dsp-level:0
\[
\qquad \boxed{\ V\,\delta_v=4\pi\ \ \text{on all five solids}\ }
\]

%dsp-level:0
\[
\qquad [\text{Descartes/Gauss--Bonnet};\ \text{machine: exact in}\ \pi\text{-units}]:
\]

%dsp-level:0
\[
4\cdot\pi=8\cdot\tfrac\pi2=6\cdot\tfrac{2\pi}3=20\cdot\tfrac\pi5=12\cdot\tfrac\pi3=4\pi\qquad[\text{machine}];
\]

%dsp-level:0
\[
\text{a loop cleanly enclosing one vertex carries}\ \Theta_R
\]

%dsp-level:0
\[
= \delta_v\ [\mathrm{FT.47}(1);\ \mathrm{Def\text{-}H.21}(6)]
\]

%dsp-level:0
\[
\implies m^2_{\mathrm{pole}}
\]

%dsp-level:0
\[
= 4\hat\mu'^2\sin^2\frac{\Theta_R}4
\]

%dsp-level:0
\[
= \hat\mu'^2D_S\ [\mathrm{Def\text{-}H.37}(2);\ \mathrm{BM.39}(7)]:
\]

%dsp-level:0
\[
\textbf{the branch mass is the vertex's own metric back-reaction}\ —\ \text{the defect at a}\ d\text{-valent vertex, and}\ d
\]

%dsp-level:0
\[
= 3\ \text{on the tetrahedron, the cube and the dodecahedron}\ [\mathrm{Census\ FT.42}(5):\ (V,E,F,d)\ \text{column}];
\]

%dsp-level:0
\[
(4)\quad\textbf{the six states:}
\]

%dsp-level:0
\[
\qquad |\mathcal S|=3\cdot2=6
\]

%dsp-level:0
\[
\qquad [\text{machine}];
\]

%dsp-level:0
\[
\qquad \textbf{the six defect values are pairwise distinct}\ [\text{machine: separation by exact field and by minimal polynomial}]:
\]

%dsp-level:0
\[
D_S\in\Bigl\{\,2,\ \ 4,\ \ 2-\sqrt2,\ \ 1,\ \ 2-2\cos\tfrac\pi{10},\ \ 2-\sqrt3\,\Bigr\},
\]

%dsp-level:0
\[
\qquad \text{minimal polynomials}\ \ (x{-}2),\ (x{-}4),\ (x^2{-}4x{+}2),\ (x{-}1),\ (x^4{-}8x^3{+}19x^2{-}12x{+}1),\ (x^2{-}4x{+}1);
\]

%dsp-level:0
\[
\textbf{the fields:}
\]

%dsp-level:0
\[
\qquad \mathbb Q,\ \ \mathbb Q,\ \ \mathbb Q(\sqrt2),\ \ \mathbb Q,\ \ \mathbb Q\bigl(\cos\tfrac\pi{10}\bigr),\ \ \mathbb Q(\sqrt3)
\]

%dsp-level:0
\[
\qquad [\text{machine: the dodecahedral quartic verified by exact substitution into}\ \mathbb Z[c]/(16c^4{-}20c^2{+}5)]:
\]

%dsp-level:0
\[
\textbf{the defect species resolves all six states with zero degeneracy};
\]

%dsp-level:0
\[
\qquad \textbf{the informational species is branch-degenerate:}
\]

%dsp-level:0
\[
\quad m_{\mathrm{hid}}\,\mathfrak h/\ln2=2E\in\{12,\,24,\,24,\,60,\,60\}\implies\textbf{three values}\ [\mathrm{FT.44}(1);\ \text{machine}];
\]

%dsp-level:0
\[
(5)\quad\textbf{the branch order (down-type heavier, uniform and exact):}
\]

%dsp-level:0
\[
\qquad \boxed{\ D_S(\mathcal F_g,+)>D_S(\mathcal F_g,-)\qquad\forall g\in\{1,2,3\}\ }
\]

%tbl-level:0
{\setlength{\tabcolsep}{4pt}
\begin{longtable}{>{\raggedright\arraybackslash}p{\dimexpr0.0949\textwidth-2\tabcolsep\relax}|>{\raggedright\arraybackslash}p{\dimexpr0.1961\textwidth-2\tabcolsep\relax}|>{\raggedright\arraybackslash}p{\dimexpr0.1961\textwidth-2\tabcolsep\relax}|>{\raggedright\arraybackslash}p{\dimexpr0.5129\textwidth-2\tabcolsep\relax}}
\toprule
floor & lighter branch & heavier branch & the exact certificate \\
\addlinespace[2pt]
\midrule
\endfirsthead
\toprule
floor & lighter branch & heavier branch & the exact certificate \\
\addlinespace[2pt]
\midrule
\endhead
\bottomrule
\endlastfoot
$\mathcal F_1$ & $D_S=\allowbreak 2$ & $D_S=\allowbreak 4$ & $4>\allowbreak 2$ \\
$\mathcal F_2$ & $2-\sqrt2$ & $1$ & $1>\allowbreak 2-\sqrt2\iff\allowbreak \sqrt2>\allowbreak 1\iff\allowbreak 2>\allowbreak 1$ \\
$\mathcal F_3$ & $2-2\cos\tfrac\pi{10}$ & $2-\sqrt3$ & $\cos\tfrac\pi{10}>\allowbreak \tfrac{\sqrt3}2\iff\allowbreak \tfrac\pi{10}<\allowbreak \tfrac\pi6\iff\allowbreak 3<\allowbreak 5$ \\
\end{longtable}
}

%dsp-level:0
\[
[\text{machine: integer and radical certificates;}\ \cos\ \text{strictly decreasing on}\ [0,\pi]\ [\text{math}]]
\]

%dsp-level:0
\[
\qquad \textbf{the constituent mass difference:}
\]

%dsp-level:0
\[
\qquad \Delta_g=\hat\mu'^2\bigl(D_S(\mathcal F_g,+)-D_S(\mathcal F_g,-)\bigr)>0;
\]

%dsp-level:0
\[
\Delta_1=2\hat\mu'^2,
\]

%dsp-level:0
\[
\qquad \Delta_2=(\sqrt2-1)\hat\mu'^2,
\]

%dsp-level:0
\[
\qquad \Delta_3=\bigl(2\cos\tfrac\pi{10}-\sqrt3\bigr)\hat\mu'^2
\]

%dsp-level:0
\[
\qquad [\text{machine: exact in}\ \mathbb Q,\ \mathbb Q(\sqrt2),\ \mathbb Q(\sqrt3,\cos\tfrac\pi{10})];
\]

%dsp-level:0
\[
(6)\quad\textbf{the two hierarchies (the derived inversion):}
\]

%dsp-level:0
\[
\qquad \textbf{on the informational species:}\quad m_{\mathrm{hid}}(\mathcal F_g)=\frac{2E(\mathcal F_g)\ln2}{\mathfrak h},
\]

%dsp-level:0
\[
\qquad 2E=(12,\,24,\,60):
\]

%dsp-level:0
\[
\boxed{\ m_{\mathrm{hid}}(\mathcal F_1)<m_{\mathrm{hid}}(\mathcal F_2)<m_{\mathrm{hid}}(\mathcal F_3),\qquad12:24:60=1:2:5\ }
\]

%dsp-level:0
\[
\qquad [\text{machine: exact rationals, strictly increasing}]:
\]

%dsp-level:0
\[
\qquad \textbf{the mass increases across the generations};
\]

%dsp-level:0
\[
\textbf{on the defect species:}
\]

%dsp-level:0
\[
\qquad D_S(\mathcal F_3,\pm)<D_S(\mathcal F_2,\pm)<D_S(\mathcal F_1,\pm)
\]

%dsp-level:0
\[
\qquad [\text{machine:}\ 2-2\cos\tfrac\pi{10}<2-\sqrt3<2-\sqrt2<1<2<4]:
\]

%dsp-level:0
\[
\boxed{\ \textbf{the two currencies order the floors oppositely}\ }
\]

%dsp-level:0
\[
\qquad [\text{the exact species inversion};\ \mathrm{BM.39}(5)\ \text{pattern}];
\]

%dsp-level:0
\[
\textbf{the relating equation is}\ \neg\textbf{derived:}
\]

%dsp-level:0
\[
\qquad \neg\exists\ \text{a theorem equating the bit currency}\ (\ln2\ \text{per directed edge})\ \text{and the angle currency}\ (\Theta_R)\ [\mathrm{BM.35}(5)];
\]

%dsp-level:0
\[
(7)\quad\textbf{the generation map:}
\]

%dsp-level:0
\[
\qquad \text{two exact three-element ordered sets}\ —\ \bigl\{(r^{(q)}_1)^2,(r^{(q)}_2)^2,(r^{(q)}_0)^2\bigr\}\ [\mathrm{FT.119}(5):\ \text{enclosure-proven, frozen}]\ \text{and}\ \{2E(\mathcal F_g)\}_{g=1,2,3}\ [(6)]:
\]

%dsp-level:0
\[
\boxed{\ \text{an order-isomorphism between two totally ordered three-element sets is unique}\ }
\]

%dsp-level:0
\[
\qquad [\text{math};\ \text{machine: the uniqueness census over all}\ 3!=6\ \text{bijections}]:
\]

%dsp-level:0
\[
g=1\longleftrightarrow(r^{(q)}_1)^2\longleftrightarrow\mathcal F_1\ (\text{the tetrahedron floor});
\]

%dsp-level:0
\[
\qquad g=2\longleftrightarrow(r^{(q)}_2)^2\longleftrightarrow\mathcal F_2\ (\text{the cube/octahedron floor});
\]

%dsp-level:0
\[
\qquad g=3\longleftrightarrow(r^{(q)}_0)^2\longleftrightarrow\mathcal F_3\ (\text{the dodecahedron/icosahedron floor});
\]

%dsp-level:0
\[
\textbf{the map is an order,}\ \neg\textbf{a number:}
\]

%dsp-level:0
\[
\qquad \hat\mu'^2,\ \mu\in\mathrm{Cal\text{-}F}\ [\text{logged}];
\]

%dsp-level:0
\[
\qquad \text{on the transport species the line order inverts}\ [\mathrm{BM.39}(5)]\implies\text{the isomorphism inverts with it}\ [\text{machine}];
\]

%dsp-level:0
\[
(8)\quad\textbf{the transport-species form of the ladder:}
\]

%dsp-level:0
\[
\qquad m^2_i=M^2\bigl(A_0-B_0\cos\varphi_i\bigr),\quad M^2=\hat\mu^2\ell_c^2,\quad A_0=\frac{3g^2}8+\frac{g'^2}{16},\quad B_0=\frac{gg'}4
\]

%dsp-level:0
\[
\qquad [\mathrm{Def\text{-}BM.11}(5);\ \mathrm{BM.39}(1)]:
\]

%dsp-level:0
\[
\textbf{the exact evaluation at the base profile}\ [\text{machine: in}\ \mathbb Q[g,g',c,s,r_3]/(c^2{+}s^2{-}1,\,r_3^2{-}3)]:
\]

%dsp-level:0
\[
\qquad G\bigl(c(i)\bigr)=A_0-B_0\cos\Bigl(\delta+\frac{2\pi(2-i)}3\Bigr);
\]

%dsp-level:0
\[
\textbf{the orbit value set:}
\]

%dsp-level:0
\[
\qquad \bigl\{G(c(i))\bigr\}_{i=0,1,2}=\bigl\{A_0-B_0\cos(\delta-\tfrac{2\pi j}3)\bigr\}_{j=0,1,2}
\]

%dsp-level:0
\[
\qquad [\text{machine: set equality in the quotient ring}];
\]

%dsp-level:0
\[
\textbf{the sum rule:}
\]

%dsp-level:0
\[
\qquad \sum_iG\bigl(c(i)\bigr)=3A_0
\]

%dsp-level:0
\[
\qquad [\text{machine: exact, phase-independent}];
\]

%dsp-level:0
\[
\qquad \textbf{the positivity:}
\]

%dsp-level:0
\[
\qquad A_0-B_0=\frac38\Bigl(g-\frac{g'}3\Bigr)^2+\frac{g'^2}{48}>0
\]

%dsp-level:0
\[
\qquad [\mathrm{BM.39}(1);\ \text{machine}];
\]

%dsp-level:0
\[
(9)\quad\textbf{the no-import bar:}
\]

%dsp-level:0
\[
\qquad \neg\exists\ \text{an SM mass, mass ratio or Yukawa value in any clause}\ [\mathrm{P5};\ \text{machine: zero external tokens, zero decimal numerals in this tier}];
\]

%dsp-level:0
\[
\neg\exists\ \text{a Yukawa-type object in the inventory}\ [\mathrm{EW.15};\ \mathrm{FT.145}(3);\ \mathrm{Def\text{-}BM.11}\ \text{preamble}];
\]

%dsp-level:0
\[
\qquad \textbf{the six states and both hierarchies are derived from the card template, the duality law, the defect spectrum and the channel angle alone};
\]

%dsp-level:0
\[
(10)\quad\textbf{status:}
\]

%dsp-level:0
\[
\qquad \textbf{derived:}\ (1)\text{--}(8);
\]

%dsp-level:0
\[
\qquad \textbf{declared:}\ \mathrm{Def\text{-}BM.14}(8)\text{'s name attachment alone};
\]

%dsp-level:0
\[
\qquad \textbf{open:}\ \text{the currency relating equation}\ (6),\ \text{the flavour composition}\,\ \text{the scales}\ [\mathrm{P4}];
\]

%dsp-level:0
\[
\neg\text{a P4 addition},\quad\neg\text{a new dimensionful scale},\quad\text{§Cal-BM stays empty};
\]

%dsp-level:0
\[
\qquad \text{species names confined to the identification layers}\ [\mathrm{Def\text{-}BM.3}(7),\ \mathrm{Def\text{-}BM.6}(7),\ \mathrm{Def\text{-}BM.14}(8)].
\]

\textbf{Req.} A6, D0.8, D2.9 (Parts 1–2); CT3.17a(2), EW.32, EW.36, EM.9 (Part 2); P1, P4, P5 (Parts 1–2); Def-H.20, Def-H.21(6), Def-H.37(2), Def-H.38, FT.44(1)–(3), FT.47(1), FT.54(4), FT.60(5), FT.111(3), FT.116, FT.119(1), FT.119(2), FT.119(4), FT.119(5), FT.143(5), FT.145(3), FT.158(1), FT.158(4) (Part 3); Census FT.42(5), FT.63(4), FT.63(6) (row-wise); Def-BM.3(7), Def-BM.6(3), Def-BM.6(4), Def-BM.11(5), Def-BM.12(1), Def-BM.14, BM.35(5), BM.39(1), BM.39(5), BM.39(7), BM.57(9) (this part); exact radical rings $\mathbb Q(\sqrt2),\mathbb Q(\sqrt3),\mathbb Q(\sqrt5),\mathbb Q(\cos\tfrac\pi{10})$, exact minimal polynomials by substitution, de Moivre binomials at $\cos5t=16c^5-20c^3+5c$, rational interval enclosures with alternating-series remainder bounds, quotient-ring trigonometric identities, order-isomorphism uniqueness [math; machine: \texttt{checks/\allowbreak{}verify\_\allowbreak{}H\_\allowbreak{}directed\_\allowbreak{}3vertex\_\allowbreak{}splitting.py} §A–§D].

\textbf{Pf.}

%dsp-level:0
\[
(1):\ \delta_q
\]

%dsp-level:0
\[
= 2/(3\cdot3\cdot3)
\]

%dsp-level:0
\[
= 2/27,
\]

%dsp-level:0
\[
\ 3\cdot\tfrac2{27}
\]

%dsp-level:0
\[
= \tfrac29,
\]

%dsp-level:0
\[
\ 16\cdot\tfrac18
\]

%dsp-level:0
\[
= 9\cdot\tfrac29
\]

%dsp-level:0
\[
= 27\cdot\tfrac2{27}
\]

%dsp-level:0
\[
= 2\ [\text{machine: exact Fractions}];
\]

%dsp-level:0
\[
\qquad \text{the spectrum and the three Vieta reads re-derived by enclosure}\ [\text{machine: widths}<2^{-90}];
\]

%dsp-level:0
\[
\qquad Q=6/9=2/3\ [\text{machine}]
\]

%dsp-level:0
\[
\text{the cubic}\ 2y^3-3y-\sqrt2\cos\tfrac29\ \text{evaluated at each}\ y=r^{(q)}_i-1\ \text{contains}\ 0\ [\text{machine}];
\]

%dsp-level:0
\[
\qquad U^3=E_0,\ U\ne E_0,\ U^2\ne E_0\ \text{for the 3-cycle over}\ \mathbb Q(i)\ [\text{machine}]
\]

%dsp-level:0
\[
(2):\ \mathrm{EW.36}\text{'s}\ \operatorname{diag}(1,0)\ \text{re-derived from}\ \tfrac12\sigma_3+\tfrac12\sigma_0\ [\text{machine}];
\]

%dsp-level:0
\[
\ \text{the union with the}\ \omega\text{-image equals the ternary alphabet}\ [\text{machine}];
\]

%dsp-level:0
\[
\qquad \text{the four finer-ledger splittings and the defect splittings re-cited}\ [\mathrm{FT.44}(2);\ \mathrm{Census\ FT.63}(6)];
\]

%dsp-level:0
\[
\ \text{the tetrahedral ladder}\ (2,4,2,0)\ \text{re-cited with its exact zero at}\ m\qquad = 4
\]

%dsp-level:0
\[
(3):\ \delta_v
\]

%dsp-level:0
\[
= 2\pi-d(L-2)\pi/L\ \text{evaluated at}\ (d,L)
\]

%dsp-level:0
\[
= (3,3),
\]

%dsp-level:0
\[
(3,4),
\]

%dsp-level:0
\[
(4,3),
\]

%dsp-level:0
\[
(3,5),
\]

%dsp-level:0
\[
(5,3)\ [\text{machine: exact rationals in}\ \pi\text{-units}];
\]

%dsp-level:0
\[
\qquad V\delta_v\ \text{multiplied out at all five}\ [\text{machine}];
\]

%dsp-level:0
\[
\qquad 4\sin^2(\delta_v/4)=D_S\ \text{by the half-angle identity}\ [\text{machine in the exact fields}]
\]

%dsp-level:0
\[
(4):\ |\mathcal S|=3\cdot2=6\ [\text{machine}];
\]

%dsp-level:0
\[
\qquad \text{pairwise distinctness: the rational members separated by integer comparison, the quadratic members by their minimal polynomials}\ (x^2-4x+2)\ \text{and}\ (x^2-4x+1)\ \text{having disjoint root sets, the quartic member separated by}\ \mathrm{Census\ FT.63}(6)\text{'s isolated roots}\ [\text{machine}];
\]

%dsp-level:0
\[
\qquad \{12,24,24,60,60\}\ \text{has three distinct values}\ [\text{machine}]
\]

%dsp-level:0
\[
(5):\ \text{the three certificates of the table}\ [\text{machine}];
\]

%dsp-level:0
\[
\qquad \cos\ \text{strictly decreasing on}\ [0,\pi]\ [\text{math}]\implies\tfrac\pi{10}<\tfrac\pi6\implies\cos\tfrac\pi{10}>\cos\tfrac\pi6
\]

%dsp-level:0
\[
= \tfrac{\sqrt3}2\ [\text{machine:}\ \cos^2\tfrac\pi6=\tfrac34];
\]

%dsp-level:0
\[
\qquad \Delta_g\ \text{formed as exact differences in each field}\ [\text{machine}]
\]

%dsp-level:0
\[
(6):\ 12<24<60\ \text{and}\ 12:24:60=1:2:5\ [\text{machine}];
\]

%dsp-level:0
\[
\qquad 2-2\cos\tfrac\pi{10}<2-\sqrt3<2-\sqrt2<1<2<4\ \text{by the certificates}\ 3<5,
\]

%dsp-level:0
\[
\ 2<3,
\]

%dsp-level:0
\[
\ 1<2,
\]

%dsp-level:0
\[
\ \text{and integer comparison}\ [\text{machine}]
\]

%dsp-level:0
\[
(7):\ \text{all}\ 6\ \text{bijections between the two ordered sets enumerated; exactly one is order-preserving}\ [\text{machine}];
\]

%dsp-level:0
\[
\qquad \text{the inversion under}\ \mathrm{BM.39}(5)\text{'s affine-decreasing row}\ [\text{machine}]
\]

%dsp-level:0
\[
(8):\ G(c(i))\ \text{built from the Gram block of}\ \mathrm{BM.57}(5)\ \text{at}\ c(i)
\]

%dsp-level:0
\[
= (\cos\delta,\cos(\delta{+}\tfrac{2\pi}3),\cos(\delta{+}\tfrac{2\pi(2-i)}3),\tfrac12)\ \text{and reduced modulo}\ c^2{+}s^2{-}1,
\]

%dsp-level:0
\[
\ r_3^2{-}3\ [\text{machine}];
\]

%dsp-level:0
\[
\qquad \text{the orbit value set and the sum rule}\ \sum_iG(c(i))=3A_0\ [\text{machine}];
\]

%dsp-level:0
\[
\qquad A_0-B_0\ \text{as a sum of squares}\ [\text{machine}]
\]

%dsp-level:0
\[
(9)\text{--}(10):\ \text{the token scans and the assembly}\ [\text{machine}]. \ \blacksquare
\]

\textbf{Phys.} Six quarks fall out of three structures that were already sitting in the corpus in separate rooms. The three generations are the three lines of the colour channel, whose angle is one integer over three integers and whose orbit length is forced by the same $C_3$ that seals the baryon. The two members of each generation are a floor and its dual: two solids that share every edge count, hence every informational mass, and differ at every vertex figure — three squares meeting at right angles against four triangles meeting at one hundred twenty degrees, three pentagons against five triangles. Where duality is trivial, at the tetrahedron, the second branch is the weight ladder's own next rung. The mass of a branch is the conical defect at its vertex, and Descartes closes the ledger: five solids, each with its vertex defects summing to exactly four pi. The hierarchy comes out exact and comes out twice, in opposite directions. In the bit currency the floors climb one to two to five, strictly. In the angle currency they descend. That inversion is not a defect of the assembly; it is the same read-relativity the corpus already proved for the family lines, and the equation that would reconcile the two currencies is the one item the ledger still carries as open. Within a generation there is no ambiguity at all: the heavier branch is heavier by an exact algebraic number — two, root two minus one, two cosines of eighteen degrees minus root three — and it is heavier for the same reason in all three. Nothing was fitted, and nothing measured entered.

\begin{theorem}[{\texorpdfstring{BM.61 — $\alpha_B$: the inventory, the coupling-blindness theorem, and the obstruction}{BM.61 - : the inventory, the coupling-blindness theorem, and the obstruction}}]\label{tag:BM-61}
\textbf{Thm.}

%dsp-level:0
\[
(1)\quad\textbf{the inventory scan:}
\]

%dsp-level:0
\[
\qquad \neg\exists\ \text{a quantity}\ \alpha_B\ \text{in the corpus:}
\]

%dsp-level:0
\[
\qquad \text{occurrences of the token in Parts 1--4, the Census}:\ \ \mathbf 0
\]

%dsp-level:0
\[
\qquad [\text{machine: the token scan}];
\]

%dsp-level:0
\[
\textbf{it is}\ \neg\textbf{a P4-class member and}\ \neg\textbf{a logged combination:}
\]

%dsp-level:0
\[
\qquad \text{the P4 core}\ [\text{A §E}],
\]

%dsp-level:0
\[
\ \mathrm{Cal\text{-}2},
\]

%dsp-level:0
\[
\ \mathrm{Cal\text{-}F},
\]

%dsp-level:0
\[
\ \text{and}\ \text{§Cal-BM (empty)}\ \text{carry no such entry}\ [\text{machine: the registry scan}];
\]

%dsp-level:0
\[
(2)\quad\textbf{the second argument is void:}
\]

%dsp-level:0
\[
\qquad \neg\exists\ \text{an}\ \alpha_{\mathrm{Strong}}\ \text{in the inventory}\ [\text{machine: zero occurrences}];
\]

%dsp-level:0
\[
\qquad \textbf{the colour sector's strength-class data are counts and ratios of counts:}
\]

%dsp-level:0
\[
\dim\mathcal K^{(\mathrm{strong})}=8\ [\mathrm{CT3.8};\ \mathrm{BM.36}(1)];
\]

%dsp-level:0
\[
\qquad g(N)=\sum_{\substack{a,b\ge0\ 2a+3b=N}}\frac{N!}{a!\,b!\,2^a\,6^b}=(0,\,1,\,1,\,3,\,10,\,25,\,105)\ \ \text{at}\ N=1,\dots,7
\]

%dsp-level:0
\[
\qquad [\mathrm{CT3.17a}(2);\ \mathrm{FT.41}(b);\ \text{machine: recomputed exactly}];
\]

%dsp-level:0
\[
\mathcal C(\Sigma;C)\in\{0,1\}\ [\mathrm{D2.14}];
\]

%dsp-level:0
\[
\qquad \text{the arenas}\ 27\ \text{and}\ 9\ [\mathrm{FT.118}(2)];
\]

%dsp-level:0
\[
\qquad \delta_q=\frac{2}{3\cdot3\cdot3}=\frac2{27}\ [\mathrm{FT.119}(1)];
\]

%dsp-level:0
\[
\qquad g(6)=25=15+10\ [\text{machine: the split}];
\]

%dsp-level:0
\[
\textbf{every one is an integer or a ratio of integers}\ —\ \neg\text{a coupling};
\]

%dsp-level:0
\[
(3)\quad\textbf{the strength bar:}
\]

%dsp-level:0
\[
\qquad \neg\exists\ \text{a theorem fixing any numerical coupling strength}\ [\mathrm{BM.36}(3)];
\]

%dsp-level:0
\[
\qquad (g,g')\ \text{are P4-class}\ [\mathrm{EW.33};\ \mathrm{P4}];
\]

%dsp-level:0
\[
\qquad \alpha^{(2)}_{\mathrm{EM}}\ \text{is an operator block read whose value no theorem fixes}\ [\mathrm{EM.17};\ \mathrm{EM.25}(8);\ \text{Part 2 §Cal-2's logged row}];
\]

%dsp-level:0
\[
\textbf{the only force-structure ratios the corpus derives are the alphabet-capacity ratios}
\]

%dsp-level:0
\[
\qquad \frac{\dim\mathcal K^{(\mathrm{EM})}}{\dim F_W}=\frac13,
\]

%dsp-level:0
\[
\qquad \frac{\dim\mathcal K^{(\mathrm{strong})}}{\dim\mathcal K^{(\mathrm{EM})}}=8,
\]

%dsp-level:0
\[
\qquad \frac{\dim\mathcal K^{(\mathrm{EM})}}{\dim\mathcal K_3}=\frac1{25}
\]

%dsp-level:0
\[
\qquad [\mathrm{BM.36}(2);\ \text{machine: exact Fractions from the flag}\ (6,16,17,25)]:
\]

%dsp-level:0
\[
\textbf{reading a capacity ratio as a coupling strength is a category error of the geometric-ratio-against-coupling species}\ [\mathrm{BM.36}(3);\ \mathrm{FT.73}(5)\ \text{trap pattern}];
\]

%dsp-level:0
\[
(4)\quad\textbf{the first argument is unfixed:}
\]

%dsp-level:0
\[
\qquad e=\frac{gg'}{\sqrt{g^2+g'^2}}=g\sin\theta_W=g'\cos\theta_W,
\]

%dsp-level:0
\[
\qquad \frac1{e^2}=\frac1{g^2}+\frac1{g'^2}
\]

%dsp-level:0
\[
\qquad [\mathrm{EW.35};\ \mathrm{EM.25}(6):\ \text{fixed with zero freedom}]:
\]

%dsp-level:0
\[
\neg\exists\ \text{a theorem asserting a value of}\ e\ [\mathrm{EM.25}(7)];
\]

%dsp-level:0
\[
\qquad \varepsilon_{\mathrm{EM}}
\]

%dsp-level:0
\[
= e^2/(4\pi)\ \text{is a dimensionless P4-class symbol carried symbolically}\ [\mathrm{Def\text{-}BM.12}\ \text{local notation};\ \mathrm{EM.25}(9);\ \mathrm{P5}];
\]

%dsp-level:0
\[
(5)\quad\textbf{the coupling-blindness theorem:}
\]

%dsp-level:0
\[
\qquad W(Q)=\bigl[z^{Q}\bigr]\bigl(2+z+z^{-1}\bigr)^{k}\in\mathbb Z[z,z^{-1}]
\]

%dsp-level:0
\[
\qquad [\mathrm{Def\text{-}BM.13}(3);\ \mathrm{Def\text{-}BM.8}(6)]:
\]

%dsp-level:0
\[
\textbf{the free-symbol set of}\ W(Q)\ \textbf{is}\ \{z\},
\]

%dsp-level:0
\[
\ \text{with coefficients fixed by}\ (m_q,k)
\]

%dsp-level:0
\[
= \bigl((2,1,1),\,3\bigr)\ [\mathrm{EM.9};\ \mathrm{D3.9}]:
\]

%dsp-level:0
\[
\qquad \text{the coupling alphabet}\ \{g,g',e,\varepsilon_{\mathrm{EM}},\hat\mu,\hat\mu',\ell_c,\theta_W\}\ \textbf{is disjoint from it}\ [\text{machine: the symbol census}];
\]

%dsp-level:0
\[
\boxed{\ \frac{\partial\,\Sigma_{\mathrm B}}{\partial g}=\frac{\partial\,\Sigma_{\mathrm B}}{\partial g'}=\frac{\partial\,\Sigma_{\mathrm B}}{\partial e}=\frac{\partial\,\Sigma_{\mathrm B}}{\partial\varepsilon_{\mathrm{EM}}}=\frac{\partial\,\Sigma_{\mathrm B}}{\partial\hat\mu}=\frac{\partial\,\Sigma_{\mathrm B}}{\partial\hat\mu'}=\frac{\partial\,\Sigma_{\mathrm B}}{\partial\ell_c}=0\ \ \text{exactly}\ }
\]

%dsp-level:0
\[
\qquad [\text{machine}]:
\]

%dsp-level:0
\[
\textbf{the charge-class transport entropy is coupling-blind};
\]

%dsp-level:0
\[
\qquad \text{it is a function of the integer alphabet and the constituent count alone};
\]

%dsp-level:0
\[
(6)\quad\textbf{the exact object at the proposed slot:}
\]

%dsp-level:0
\[
\qquad \boxed{\ \Sigma_{\mathrm B}=\ln\frac{W(0)}{W(+1)}=\ln\frac43=2\ln2-\ln3\ }
\]

%dsp-level:0
\[
\qquad [\mathrm{Def\text{-}BM.13}(4);\ \text{machine}]:
\]

%dsp-level:0
\[
e^{\Sigma_{\mathrm B}}=\frac43
\]

%dsp-level:0
\[
\qquad [\text{machine: the integer row}\ 20\cdot3=15\cdot4];
\]

%dsp-level:0
\[
\qquad [\Sigma_{\mathrm B}]=\mathfrak q^0\ [\mathrm{CT2.1};\ \text{machine}];
\]

%dsp-level:0
\[
\qquad \Sigma_{\mathrm B}>0\ [\text{machine:}\ 4>3];
\]

%dsp-level:0
\[
\textbf{the exact enclosures inside the}\ \ln\textbf{-ring:}
\]

%dsp-level:0
\[
\qquad \frac{\ln2}3<\Sigma_{\mathrm B}<\frac{\ln2}2
\]

%dsp-level:0
\[
\qquad [\text{machine:}\ 4^3>3^3\cdot2\iff64>54;\quad4^2<3^2\cdot2\iff16<18];
\]

%dsp-level:0
\[
\frac{2\ln2}5<\Sigma_{\mathrm B}<\frac{3\ln2}5\qquad[\text{machine:}\ 4^5>3^5\cdot2^2\iff1024>972;\quad4^5<3^5\cdot2^3\iff1024<1944];
\]

%dsp-level:0
\[
(7)\quad\textbf{the verdict:}
\]

%dsp-level:0
\[
\qquad \Sigma_{\mathrm B}\ \text{enters the mass as a}\ \textbf{weight in the rest slot}\ [\mathrm{Def\text{-}BM.13}(5)],\ \neg\text{as a coupling strength:}
\]

%dsp-level:0
\[
\qquad \text{it multiplies no generator, resolves no propagator, and appears in no vertex};
\]

%dsp-level:0
\[
\textbf{it is}\ \neg\alpha_B:\qquad\text{a coupling is a coefficient of an interaction;}\ \Sigma_{\mathrm B}\ \text{is the logarithm of a census};
\]

%dsp-level:0
\[
(8)\quad\textbf{the obstruction, named:}\qquad\boxed{\ \alpha_B=f(\alpha_{\mathrm{EM}},\alpha_{\mathrm{Strong}})\ \ \textbf{is}\ \neg\textbf{derivable}\ }\qquad\text{on four independent grounds:}
\]

%dsp-level:0
\[
(\mathrm O1)\ \ \textbf{the strength bar:}\ \neg\exists\ \text{a count}\to\text{strength theorem}\ [(3);\ \mathrm{BM.36}(3)];
\]

%dsp-level:0
\[
\qquad (\mathrm O2)\ \ \textbf{the colour-coupling void:}\ \alpha_{\mathrm{Strong}}\ \text{is}\ \neg\text{an inventory item, so the domain of}\ f\ \text{is empty in its second argument}\ [(2)];
\]

%dsp-level:0
\[
(\mathrm O3)\ \ \textbf{the coupling-blindness:}\ \text{the one exact dimensionless charge-class quantity the corpus derives has zero derivative in every coupling, so it is not a function of them}\ [(5)];
\]

%dsp-level:0
\[
\qquad (\mathrm O4)\ \ \textbf{P5:}\ \text{fixing}\ \alpha_B\ \text{by an observed ratio is barred}\ [\mathrm{P5};\ \mathrm{BM.36}(3);\ \mathrm{FT.73}(5)]:
\]

%dsp-level:0
\[
\textbf{the obstruction is not a gap in the arithmetic; it is the absence of a coupling theorem for colour and the bar against inventing one};
\]

%dsp-level:0
\[
(9)\quad\textbf{status:}
\]

%dsp-level:0
\[
\qquad \textbf{derived:}\ (1)\text{--}(7);
\]

%dsp-level:0
\[
\qquad \textbf{named, not filled:}\ \alpha_B\ (8);
\]

%dsp-level:0
\[
\qquad \alpha_B\ \textbf{is}\ \neg\textbf{added as a P4-class parameter};
\]

%dsp-level:0
\[
\qquad \text{§Cal-BM stays empty};
\]

%dsp-level:0
\[
\neg\text{a new dimensionful scale}\ [\mathrm{P1}];
\]

%dsp-level:0
\[
\qquad \neg\text{a measured value and}\ \neg\text{a numerical coupling}\ [\mathrm{P5};\ \text{machine: zero external tokens, zero decimal numerals in this tier}];\qquad \qquad
\]

\textbf{Req.} D2.9, D2.14, D3.9, CT3.8, CT3.17a(2), EW.33, EW.35, EM.9, EM.17, EM.25(6), EM.25(7), EM.25(8), EM.25(9) (Part 2); CT0.13, CT0.14, CT2.1, P1, P4, P5 (Parts 1–2); FT.41(b), FT.73(5), FT.118(2), FT.119(1) (Part 3); Def-BM.8(1), Def-BM.8(6), Def-BM.12, Def-BM.13(3), Def-BM.13(4), Def-BM.13(5), BM.36(1), BM.36(2), BM.36(3) (this part); A §E (the P4 core); Part 2 §Cal-2 (the logged row); the exact grammar-count summation, symbol censuses, exact $\ln$-ring identities with integer certificates, Laurent exponent arithmetic [math; machine: \texttt{checks/\allowbreak{}verify\_\allowbreak{}H\_\allowbreak{}constituent\_\allowbreak{}transport\_\allowbreak{}splitting.py} §E].

\textbf{Pf.}

%dsp-level:0
\[
(1)\text{--}(2):\ \text{the token scans over the five canon files and the two annexes}\ [\text{machine}];
\]

%dsp-level:0
\[
\qquad g(N)\ \text{recomputed by the}\ \mathrm{CT3.17a}(2)\ \text{summation for}\ N
\]

%dsp-level:0
\[
= 1,
\]

%dsp-level:0
\[
\dots,
\]

%dsp-level:0
\[
7\ \text{and compared to the frozen row}\ (0,1,1,3,10,25,105)\ [\text{machine}];
\]

%dsp-level:0
\[
\qquad g(6)=15+10=25\ [\text{machine: the}\ a{=}3,b{=}0\ \text{and}\ a{=}0,b{=}2\ \text{summands}]
\]

%dsp-level:0
\[
(3)\text{--}(4):\ \mathrm{BM.36}(2),(3),\ \mathrm{EM.25}(6)\text{--}(8)\ \text{re-cited};
\]

%dsp-level:0
\[
\qquad \tfrac{1}{3},
\]

%dsp-level:0
\[
\ \tfrac81,
\]

%dsp-level:0
\[
\ \tfrac1{25}\ \text{from}\ (6,16,17,25):\ 16/3\to\tfrac13\ \text{via}\ \dim F_W
\]

%dsp-level:0
\[
= 3,
\]

%dsp-level:0
\[
\ 8/1,
\]

%dsp-level:0
\[
\ 1/25\ [\text{machine: exact Fractions}]
\]

%dsp-level:0
\[
(5):\ W(Q)\ \text{written as a Laurent polynomial with integer coefficients; the symbol set extracted and intersected with the coupling alphabet}\ [\text{machine: the intersection is empty}];
\]

%dsp-level:0
\[
\qquad \text{each partial derivative formed symbolically and evaluated}\ [\text{machine: all seven are the zero polynomial}]
\]

%dsp-level:0
\[
(6):\ \ln20-\ln15=\ln\tfrac{20}{15}=\ln\tfrac43=2\ln2-\ln3\ [\text{machine: the prime-log ring}];
\]

%dsp-level:0
\[
\qquad 20\cdot3=60=15\cdot4\ [\text{machine}];
\]

%dsp-level:0
\[
\qquad \text{the four enclosures:}\ 4^3
\]

%dsp-level:0
\[
= 64>54
\]

%dsp-level:0
\[
= 3^3\cdot2;
\]

%dsp-level:0
\[
\ 4^2
\]

%dsp-level:0
\[
= 16<18
\]

%dsp-level:0
\[
= 3^2\cdot2;
\]

%dsp-level:0
\[
\ 4^5
\]

%dsp-level:0
\[
= 1024>972
\]

%dsp-level:0
\[
= 3^5\cdot2^2;
\]

%dsp-level:0
\[
\ 4^5
\]

%dsp-level:0
\[
= 1024<1944
\]

%dsp-level:0
\[
= 3^5\cdot2^3\ [\text{machine: integer comparisons}]
\]

%dsp-level:0
\[
(7)\text{--}(9):\ \text{the locus row re-cited from}\ \mathrm{Def\text{-}BM.13}(5);
\]

%dsp-level:0
\[
\ \text{the four grounds conjoined};
\]

%dsp-level:0
\[
\ \text{the standing-rule scans}\ [\text{machine}]. \ \blacksquare
\]

\textbf{Phys.} The proposal was to close the neutron row with a new coupling, fixed by the ratio the world reports. Three of the four things that proposal needs are absent, and their absence is proved rather than assumed. There is no $\alpha_B$ anywhere in the canon to refine. There is no strong coupling for a function of two strengths to take as its second argument: the colour sector's entire strength-class inventory is a dimension, a projector, a grammar count and two arena sizes, all integers, and the corpus has already ruled that turning a count into a strength is the same category error as turning a geometric ratio into a coupling. And the one exact dimensionless charge-class number the corpus does derive — the logarithm of twenty over fifteen — has a zero derivative in every coupling it owns, which settles the question from the other side: it is not a function of the strengths, so it cannot be dressed as one. What $\ln(4/3)$ actually is, is a weight: it sits in the rest slot beside the edge count, in the same nats, and it does the work the proposal wanted without the parameter the proposal needed. The obstruction is named to the monomial, and no free entry is added.

\end{theorem}

\pseudohead{theorem}{\texorpdfstring{BM.62}{BM.62}}{\texorpdfstring{The full mass splitting: the extended decomposition, the exact critical strength, and the sign theorem}{The full mass splitting: the extended decomposition, the exact critical strength, and the sign theorem}}\label{tag:BM-62}
\textbf{Thm.}

%dsp-level:0
\[
(1)\quad\textbf{the extended decomposition:}
\]

%dsp-level:0
\[
\qquad \boxed{\ \Delta m\ :=\ m_n-m_p\ =\ \Delta m_{\mathrm{Lyapunov}}+\Delta m_{\mathrm{force}}+\Delta m_{\mathrm{metric}}+\Delta m_{\mathrm{binding}}+\Delta m_{\mathrm{self\text{-}int}}+\Delta m_{\mathrm{transport}}\ }
\]

%dsp-level:0
\[
\text{with}\ \mathrm{BM.55}(1)\text{'s five derived rows standing unmodified and one new row added}\ [\mathrm{BM.52}(5)\text{'s no-double-count guard respected:}\ \mathrm{BM.59}(8)\text{'s disjointness}]:
\]

%dsp-level:0
\[
(2)\quad\textbf{the four zeros:}
\]

%dsp-level:0
\[
\qquad \Delta m_{\mathrm{Lyapunov}}=0\ [\mathrm{BM.51}(6):\ \text{every input is class-shared}];
\]

%dsp-level:0
\[
\qquad \Delta m_{\mathrm{metric}}=0\ [\mathrm{BM.53}(2);\ \mathrm{BM.58}(4):\ \text{the charge exponent is}\ 0\ \text{in all seven monomials}];
\]

%dsp-level:0
\[
\Delta m_{\mathrm{binding}}
\]

%dsp-level:0
\[
= 0\ [\mathrm{BM.54}(1)\text{--}(3):\ \text{the count read vanishes as}\ 6-(2{+}2{+}2)=0,\ \text{and the ensemble read is}\ q\text{-degenerate}];
\]

%dsp-level:0
\[
\textbf{re-derived here:}
\]

%dsp-level:0
\[
\qquad \mathcal B(S_j;1,1,1)
\]

%dsp-level:0
\[
= m^2_B(\cdot,S_j)-3\,
\]

%dsp-level:0
\[
m^2_B(\cdot,1)\ \ \text{is a function of the junction sector alone}\ [\mathrm{BM.21}(2);\ \mathrm{BM.20}(4)]:
\]

%dsp-level:0
\[
\mathcal B(0;1,1,1)
\]

%dsp-level:0
\[
= \frac{843957}{98956}-\frac{5764986}{393679},
\]

%dsp-level:0
\[
\quad\mathcal B(1;1,1,1)
\]

%dsp-level:0
\[
= -\frac{3843324}{393679},
\]

%dsp-level:0
\[
\quad\mathcal B(2;1,1,1)
\]

%dsp-level:0
\[
= \frac{91}{23}-\frac{5764986}{393679},
\]

%dsp-level:0
\[
\quad\mathcal B(3;1,1,1)
\]

%dsp-level:0
\[
= \frac{13713}{3088}-\frac{5764986}{393679};
\]

%dsp-level:0
\[
\textbf{all four are nonzero exact rationals, all four are negative, and the class difference of each is exactly}\ 0
\]

%dsp-level:0
\[
\qquad [\text{machine: the exact rational census}];
\]

%dsp-level:0
\[
m^2_4(\cdot,0)>m^2_4(\cdot,1)>m^2_4(\cdot,3)>m^2_4(\cdot,2)
\]

%dsp-level:0
\[
\qquad [\mathrm{BM.21}(4);\ \text{machine}];
\]

%dsp-level:0
\[
\qquad \text{and a binding term in a free member's mass remains a category error}\ [\mathrm{BM.48}(4);\ \mathrm{BM.54}(4)];
\]

%dsp-level:0
\[
(3)\quad\textbf{the two nonzero summands:}
\]

%dsp-level:0
\[
\qquad \Delta m_{\mathrm{force}}=\Delta m_{\mathrm{self\text{-}int}}=-\Delta m_{\mathrm{EM}}=-\,\frac{3\,c\,e^2\ln2}{2\pi\,\mathfrak h}=-\,6\,c\,\varepsilon_{\mathrm{EM}}\,\frac{\ln2}{\mathfrak h}
\]

%dsp-level:0
\[
\qquad [\mathrm{BM.52}(4),(5);\ \mathrm{BM.46}(4);\ \text{machine: one Laurent monomial, counted once}];
\]

%dsp-level:0
\[
\Delta m_{\mathrm{transport}}=+\,\frac{\Sigma_{\mathrm B}}{\mathfrak h}=+\,\frac{\ln(4/3)}{\mathfrak h}>0\qquad[\mathrm{BM.59}(4);\ \text{machine:}\ 4>3];
\]

%dsp-level:0
\[
(4)\quad\textbf{the assembly:}
\]

%dsp-level:0
\[
\qquad \boxed{\ \Delta m=\frac{\ln(4/3)-6\,c\,\varepsilon_{\mathrm{EM}}\ln2}{\mathfrak h}=\frac{2\ln2-\ln3-6\,c\,\varepsilon_{\mathrm{EM}}\ln2}{\mathfrak h}\ }
\]

%dsp-level:0
\[
\qquad [\text{machine: the monomial sum in the}\ \ln\text{-ring}]:
\]

%dsp-level:0
\[
m_n\,\mathfrak h=6\ln2+\ln20=\ln1280,
\]

%dsp-level:0
\[
\qquad m_p\,\mathfrak h=6\ln2+\ln15+6\,c\,\varepsilon_{\mathrm{EM}}\ln2=\ln960+6\,c\,\varepsilon_{\mathrm{EM}}\ln2
\]

%dsp-level:0
\[
\qquad [\mathrm{Def\text{-}BM.13}(5);\ \text{machine}];
\]

%dsp-level:0
\[
[\Delta m]=\mathfrak q^{+1/2}=[\text{energy}]
\]

%dsp-level:0
\[
\qquad [\mathrm{CT2.1};\ \mathrm{D2.9};\ \text{machine: exponent arithmetic}];
\]

%dsp-level:0
\[
\qquad \textbf{both summands ride the same rest slot}\ [\mathrm{Def\text{-}BM.10}(4);\ \mathrm{Def\text{-}BM.12}(6);\ \mathrm{BM.37}(4)];
\]

%dsp-level:0
\[
(5)\quad\textbf{the exact ratio:}
\]

%dsp-level:0
\[
\qquad \frac{m_n}{m_p}=\frac{\ln1280}{\ln960+6\,c\,\varepsilon_{\mathrm{EM}}\ln2}
\]

%dsp-level:0
\[
\qquad [\text{machine}];
\]

%dsp-level:0
\[
\qquad \textbf{at}\ \varepsilon_{\mathrm{EM}}=0:\quad\frac{m_n}{m_p}=\frac{\ln1280}{\ln960}=\frac{8\ln2+\ln20}{6\ln2+\ln15}\ \ne\ \frac43
\]

%dsp-level:0
\[
\qquad [\text{machine:}\ 3\ln1280\ne4\ln960]:
\]

%dsp-level:0
\[
\textbf{the amended}\ \mathrm{BM.55}(2)\textbf{ ratio:}
\]

%dsp-level:0
\[
\qquad \frac{4\pi}{4\pi+c\,e^2}\ \longmapsto\ \frac{\ln1280}{\ln960+6\,c\,\varepsilon_{\mathrm{EM}}\ln2}
\]

%dsp-level:0
\[
\qquad [\text{the transport refinement of the rest input}];
\]

%dsp-level:0
\[
(6)\quad\textbf{the exact critical strength:}
\]

%dsp-level:0
\[
\qquad \Delta m\ \text{is strictly decreasing in}\ \varepsilon_{\mathrm{EM}}:
\]

%dsp-level:0
\[
\qquad \frac{\partial\,\Delta m}{\partial\,\varepsilon_{\mathrm{EM}}}=-\,\frac{6\,c\ln2}{\mathfrak h}<0
\]

%dsp-level:0
\[
\qquad [\text{machine: the coefficient census at}\ c,\ln2,\mathfrak h>0]:
\]

%dsp-level:0
\[
\boxed{\ \varepsilon^{\mathrm{crit}}(c)\ :=\ \frac{\Sigma_{\mathrm B}}{6\,c\,\ln2}\ =\ \frac{2\ln2-\ln3}{6\,c\,\ln2}\ =\ \frac{2-\log_23}{6\,c}\ }\qquad\textbf{the unique zero of}\ \Delta m\ [\text{machine}];
\]

%dsp-level:0
\[
\textbf{the sign trichotomy:}
\]

%dsp-level:0
\[
\qquad \varepsilon_{\mathrm{EM}}<\varepsilon^{\mathrm{crit}}(c)\iff\Delta m>0;
\]

%dsp-level:0
\[
\qquad \varepsilon_{\mathrm{EM}}=\varepsilon^{\mathrm{crit}}(c)\iff\Delta m=0;
\]

%dsp-level:0
\[
\qquad \varepsilon_{\mathrm{EM}}>\varepsilon^{\mathrm{crit}}(c)\iff\Delta m<0
\]

%dsp-level:0
\[
\qquad [\text{machine}];
\]

%dsp-level:0
\[
(7)\quad\textbf{the exact enclosures of the critical strength:}
\]

%dsp-level:0
\[
\qquad \text{every comparison is an integer power certificate}\ [\text{machine}]:
\]

%tbl-level:1
{\small\setlength{\tabcolsep}{4pt}
\begin{longtable}{>{\raggedright\arraybackslash}p{\dimexpr0.3009\textwidth-2\tabcolsep\relax}|>{\raggedright\arraybackslash}p{\dimexpr0.1593\textwidth-2\tabcolsep\relax}|>{\raggedright\arraybackslash}p{\dimexpr0.2124\textwidth-2\tabcolsep\relax}|>{\raggedright\arraybackslash}p{\dimexpr0.3274\textwidth-2\tabcolsep\relax}}
\toprule
$c$ & $\varepsilon^{\mathrm{crit}}(c)$ & lower certificate & upper certificate \\
\addlinespace[2pt]
\midrule
\endfirsthead
\toprule
$c$ & $\varepsilon^{\mathrm{crit}}(c)$ & lower certificate & upper certificate \\
\addlinespace[2pt]
\midrule
\endhead
\bottomrule
\endlastfoot
$\tfrac12$ (the shell $[\mathrm{Def\text{-}BM.12}(4)]$) & $\dfrac{2-\log_23}{3}$ & $>\allowbreak \dfrac18\iff\allowbreak 2^{13}>\allowbreak 3^8\iff\allowbreak 8192>\allowbreak 6561$ & $<\allowbreak \dfrac17\iff\allowbreak 2^{11}<\allowbreak 3^7\iff\allowbreak 2048<\allowbreak 2187$ \\
$\tfrac35$ (the uniform ball $[\mathrm{Def\text{-}BM.12}(4)]$) & $\dfrac{5(2-\log_23)}{18}$ & $>\allowbreak \dfrac19\iff\allowbreak 2^{8}>\allowbreak 3^5\iff\allowbreak 256>\allowbreak 243$ & $<\allowbreak \dfrac18\iff\allowbreak 2^{31}<\allowbreak 3^{20}\iff\allowbreak 2147483648<\allowbreak 3486784401$ \\
\end{longtable}
}

%dsp-level:0
\[
\boxed{\ \frac19<\varepsilon^{\mathrm{crit}}\Bigl(\frac35\Bigr)<\frac18<\varepsilon^{\mathrm{crit}}\Bigl(\frac12\Bigr)<\frac17\ }
\]

%dsp-level:0
\[
\qquad [\text{machine: the four certificates conjoined}];
\]

%dsp-level:0
\[
(8)\quad\textbf{the sub-criticality bar (parameter-free):}\qquad\boxed{\ \varepsilon_{\mathrm{EM}}\le\frac19\ \ \Longrightarrow\ \ \Delta m>0\ \ \text{at both priced}\ c\ }
\]

%dsp-level:0
\[
\textbf{the witnesses at}\ \varepsilon_{\mathrm{EM}}=\tfrac19:
\]

%dsp-level:0
\[
\qquad c=\tfrac12:\ \ \ln\tfrac43>\tfrac{\ln2}3\iff4^3>3^3\cdot2\iff64>54;
\]

%dsp-level:0
\[
\qquad c=\tfrac35:\ \ \ln\tfrac43>\tfrac{2\ln2}5\iff4^5>3^5\cdot2^2\iff1024>972
\]

%dsp-level:0
\[
\qquad [\text{machine}];
\]

%dsp-level:0
\[
\text{and at}\ \varepsilon_{\mathrm{EM}}=0:\qquad\Delta m=\frac{\ln(4/3)}{\mathfrak h}>0\iff4>3\qquad[\text{machine}];
\]

%dsp-level:0
\[
(9)\quad\textbf{the ordering theorem:}
\]

%dsp-level:0
\[
\qquad \boxed{\ \text{on the sub-critical branch}\ \varepsilon_{\mathrm{EM}}<\varepsilon^{\mathrm{crit}}(c):\qquad m_n>m_p\ }
\]

%dsp-level:0
\[
\qquad [(6),(8);\ \text{machine}]:
\]

%dsp-level:0
\[
\textbf{the neutron is the heavier baryon};
\]

%dsp-level:0
\[
\qquad \mathrm{BM.55}(3)\text{'s strict order is read-relative and holds on the super-critical branch}\ \varepsilon_{\mathrm{EM}}>\varepsilon^{\mathrm{crit}}(c);
\]

%dsp-level:0
\[
\qquad \text{the critical surface is exact and unique}\ [(6)];
\]

%dsp-level:0
\[
(10)\quad\textbf{the fractional splittings (exact rational enclosures):}
\]

%dsp-level:0
\[
\qquad \frac{\Delta m_{\mathrm{transport}}}{m_{\mathrm{hid}}}=\frac{\ln(4/3)}{6\ln2}\in\Bigl(\frac1{15},\frac1{14}\Bigr)
\]

%dsp-level:0
\[
\qquad [\text{machine:}\ 15\ln\tfrac43>6\ln2\iff2^{24}>3^{15}\iff16777216>14348907;\quad14\ln\tfrac43<6\ln2\iff2^{22}<3^{14}\iff4194304<4782969];
\]

%dsp-level:0
\[
\frac{\Delta m_{\mathrm{transport}}}{m_n}\bigg|_{\varepsilon_{\mathrm{EM}}=0}=\frac{\ln(4/3)}{\ln1280}\in\Bigl(\frac1{25},\frac1{24}\Bigr)
\]

%dsp-level:0
\[
\qquad [\text{machine: integer power certificates}];
\]

%dsp-level:0
\[
\qquad \frac{\Delta m_{\mathrm{EM}}}{m_{\mathrm{hid}}}=c\,\varepsilon_{\mathrm{EM}}\in\mathfrak q^0\ [\mathrm{BM.46}(1);\ \text{machine}];
\]

%dsp-level:0
\[
\textbf{the strict contraction against the barred measure price:}
\]

%dsp-level:0
\[
\qquad \frac{\ln(4/3)}{6\ln2}<\frac13\iff3\ln\tfrac43<2\ln2\iff16<27
\]

%dsp-level:0
\[
\qquad [\mathrm{Def\text{-}BM.13}(8);\ \text{machine}];
\]

%dsp-level:0
\[
(11)\quad\textbf{the species row:}
\]

%dsp-level:0
\[
\qquad \text{the assembly}\ (4)\ \text{is legitimate because all four nonzero-or-zero terms enter through one slot:}
\]

%dsp-level:0
\[
\text{the rest input of the dispersion read}\ [\mathrm{Def\text{-}BM.10}(4);\ \mathrm{BM.35}(4);\ \mathrm{BM.37}(4)],
\]

%dsp-level:0
\[
\ \text{refined in the}\ \mathrm{Def\text{-}BM.12}(6)\ \text{pattern};
\]

%dsp-level:0
\[
\qquad \neg\text{a route ratio is read as a mass ratio}\ [\mathrm{BM.37}(6);\ \mathrm{Def\text{-}BM.13}(8);\ (5)];
\]

%dsp-level:0
\[
\textbf{the input alphabets:}
\]

%dsp-level:0
\[
\qquad \Delta m_{\mathrm{transport}}:\ \{m_q,W(Q),\ln2,\mathfrak h\};
\]

%dsp-level:0
\[
\qquad \Delta m_{\mathrm{EM}}:\ \{e,c,\Delta x_{\mathrm B}\};
\]

%dsp-level:0
\[
\qquad \Delta m_{\mathrm{binding}}:\ \{(\Gamma,R),\Delta S\};
\]

%dsp-level:0
\[
\qquad \Delta m_{\mathrm{metric}}:\ \{(a_{00},q_{ij},w_i),E_{\mathrm{int}}\};
\]

%dsp-level:0
\[
\text{pairwise disjoint except for the shared}\ \mathfrak h\ \text{and}\ E_{\mathrm{int}}\ [\text{machine: the token census}];
\]

%dsp-level:0
\[
\qquad \mathrm{BM.37}(2)\text{'s disjointness theorem is respected:}\ \neg\text{a measure route is added to a metric read};
\]

%dsp-level:0
\[
(12)\quad\textbf{status:}
\]

%dsp-level:0
\[
\qquad \textbf{derived:}\ (1)\text{--}(11);
\]

%dsp-level:0
\[
\qquad \textbf{the sign theorem is conditional on the sub-critical branch, with the hypothesis on its face and the critical value exact}\ (6)\text{--}(9)\ [\text{A §0.2a}];
\]

%dsp-level:0
\[
\neg\text{a P4 addition},
\]

%dsp-level:0
\[
\quad\neg\text{a new dimensionful scale}\ [\mathrm{P1}],
\]

%dsp-level:0
\[
\quad\neg\text{a route adopted}\ [\mathrm{BM.30}(5)],
\]

%dsp-level:0
\[
\quad\neg\text{a measured value}\ [\mathrm{P5}];
\]

%dsp-level:0
\[
\qquad \textbf{the open row:}\ \text{the numerical value of}\ e\ [\mathrm{EM.25}(7):\ \text{one calibration datum}]\ —\ \text{it decides which side of}\ \varepsilon^{\mathrm{crit}}\ \text{the corpus sits on, and no theorem of Parts 1--4 supplies it}.
\]

\textbf{Req.} D2.9, CT2.1, EM.25(7), EW.35 (Part 2); P1, P4, P5 (Parts 1–2); FT.79(5), FT.107(2), FT.107(3) (Part 3); Census FT.35(2) (row-wise); Def-BM.4, Def-BM.8(6), Def-BM.10(4), Def-BM.12(4), Def-BM.12(6), Def-BM.13(4), Def-BM.13(5), Def-BM.13(8), BM.20(4), BM.21(2), BM.21(4), BM.30(5), BM.35(4), BM.37(2), BM.37(4), BM.37(6), BM.46(1), BM.46(4), BM.48(4), BM.51(6), BM.52(4), BM.52(5), BM.53(2), BM.54(1)–(4), BM.55(1), BM.55(2), BM.55(3), BM.58(4), BM.59(4), BM.59(8) (this part); exact $\ln$-ring arithmetic with integer power certificates, Laurent-monomial assembly, exact rational sign censuses, exact rational enclosures of $\ln$-ratios [math; machine: \texttt{checks/\allowbreak{}verify\_\allowbreak{}H\_\allowbreak{}constituent\_\allowbreak{}transport\_\allowbreak{}splitting.py} §F].

\textbf{Pf.}

%dsp-level:0
\[
(1)\text{--}(3):\ \mathrm{BM.55}(1)\text{'s five rows re-cited (frozen)};
\]

%dsp-level:0
\[
\qquad \mathcal B(S_j;1,1,1)
\]

%dsp-level:0
\[
= m^2_B(\cdot,S_j)-3m^2_B(\cdot,1)\ \text{formed from}\ \mathrm{BM.20}(4)\text{'s four sector masses and compared to}\ \mathrm{BM.21}(2)\text{'s four frozen rationals}\ [\text{machine: exact Fractions, all four agree}];
\]

%dsp-level:0
\[
\ \text{the sign census and the sector order}\ [\text{machine}];
\]

%dsp-level:0
\[
\ \text{the class difference of each}\
\]

%dsp-level:0
\[
= 0\ \text{because}\ m^2_B(q,S)\ \text{is}\ q\text{-independent}\ [\mathrm{FT.107}(3);\ \text{machine}]
\]

%dsp-level:0
\[
(4):\ \ln1280-\ln960=\ln\tfrac{1280}{960}=\ln\tfrac43=\Sigma_{\mathrm B}\ [\text{machine: the prime-log ring}];
\]

%dsp-level:0
\[
\qquad 6\ln2+\ln20=8\ln2+\ln20=\ln(2^8\cdot20)=\ln1280\ [\text{machine}];
\]

%dsp-level:0
\[
\qquad 6\ln2+\ln15=\ln(2^6\cdot15)=\ln960\ [\text{machine}];
\]

%dsp-level:0
\[
\qquad \text{the exponent row}\ \mathfrak q^0\cdot\mathfrak q^{+1/2}\ [\text{machine}]
\]

%dsp-level:0
\[
(5):\ 3\ln1280\ne4\ln960\iff1280^3\ne960^4\ [\text{machine: the two integers differ}];
\]

%dsp-level:0
\[
\qquad (6):\ \Delta m
\]

%dsp-level:0
\[
= 0\iff\Sigma_{\mathrm B}
\]

%dsp-level:0
\[
= 6c\varepsilon_{\mathrm{EM}}\ln2\iff\varepsilon_{\mathrm{EM}}
\]

%dsp-level:0
\[
= \Sigma_{\mathrm B}/(6c\ln2)\ [\text{machine: solved in the}\ \ln\text{-ring}];
\]

%dsp-level:0
\[
\ \text{the derivative}\ -6c\ln2/\mathfrak h\ [\text{machine}]
\]

%dsp-level:0
\[
(7):\ \text{cross-multiplication by}\ 6c\ln2>0:\quad\varepsilon^{\mathrm{crit}}(c)>r\iff\Sigma_{\mathrm B}>6cr\ln2;
\]

%dsp-level:0
\[
\qquad c=\tfrac12,\ r=\tfrac18:\ 8\Sigma_{\mathrm B}>3\ln2\iff(4/3)^8>2^3\iff2^{13}>3^8\ [\text{machine}];
\]

%dsp-level:0
\[
\qquad c=\tfrac12,\ r=\tfrac17:\ 7\Sigma_{\mathrm B}<3\ln2\iff2^{11}<3^7\ [\text{machine}]
\]

%dsp-level:0
\[
c=\tfrac35,\ r=\tfrac19:\ 5\Sigma_{\mathrm B}>2\ln2\iff2^8>3^5\ [\text{machine}];
\]

%dsp-level:0
\[
\qquad c=\tfrac35,\ r=\tfrac18:\ 40\Sigma_{\mathrm B}<18\ln2\iff2^{31}<3^{20}\ [\text{machine}]
\]

%dsp-level:0
\[
(8):\ \text{the two witnesses}\ 4^3=64>54=3^3\cdot2\ \text{and}\ 4^5=1024>972=3^5\cdot2^2\ [\text{machine}];
\]

%dsp-level:0
\[
\qquad 4>3\ [\text{machine}];
\]

%dsp-level:0
\[
\qquad (9):\ \text{conjunction of}\ (6),(8)
\]

%dsp-level:0
\[
(10):\ 15\Sigma_{\mathrm B}>6\ln2\iff4^{15}>3^{15}2^6\iff2^{24}>3^{15}\ [\text{machine}];
\]

%dsp-level:0
\[
\qquad 14\Sigma_{\mathrm B}<6\ln2\iff4^{14}<3^{14}2^6\iff2^{22}<3^{14}\ [\text{machine}];
\]

%dsp-level:0
\[
\qquad \text{the}\ \ln1280\ \text{enclosure by the same route}\ [\text{machine}];
\]

%dsp-level:0
\[
\qquad 3\ln\tfrac43<2\ln2\iff64<108\ [\text{machine}]
\]

%dsp-level:0
\[
(11)\text{--}(12):\ \text{the token census of the four alphabets}\ [\text{machine}];
\]

%dsp-level:0
\[
\ \text{the standing-rule scans}\ [\text{machine}]. \ \blacksquare
\]

\textbf{Phys.} The decomposition the full-scope tier closed with four zeros and one monomial gains a fifth zero and a second monomial, and the second one points the other way. Nothing in the amended assembly contradicts the earlier one: the Lyapunov tower still shares every input, the geometry still carries no charge exponent, the binding still dies twice, and the scar's field energy is still counted exactly once. What is added is the transport register — twenty admissible end-sign assignments in the neutral class against fifteen in the charged one — read as an entropy in the ledger's own nats and booked through the same rest slot the field energy already uses. The two summands are commensurate, so the splitting becomes a single exact expression with a single exact zero: the critical strength is the logarithm of four thirds over six $c\ln2$, enclosed between one ninth and one seventh by four integer power comparisons, with the tighter of the two priced distributions clearing one ninth. Below that surface the neutron is heavier; above it the earlier order returns; on the surface the two rows are degenerate again. The corpus fixes no value for the coupling, so which branch it sits on is the one open row — and the branch that gives the world's order is not a thin one: it is an exact interval, certified by $256>243$, containing every coupling that keeps the field energy a correction rather than a competitor.

\begin{theorem}[{\texorpdfstring{BM.63 — The decay re-adjudication on the refined read: the closure survives, with an exact margin}{BM.63 - The decay re-adjudication on the refined read: the closure survives, with an exact margin}}]\label{tag:BM-63}
\textbf{Thm.}

%dsp-level:0
\[
(1)\quad\textbf{the amended Q-value:}
\]

%dsp-level:0
\[
\qquad Q_\beta:=m_n-m_p-m_\ell-m_\nu=\Delta m-m_\ell-m_\nu
\]

%dsp-level:0
\[
\qquad [\mathrm{BM.41}(4);\ \mathrm{BM.49}(1);\ \mathrm{BM.55}(5)\ \text{re-cited}]:
\]

%dsp-level:0
\[
\boxed{\ Q_\beta=\frac{\ln(4/3)-6\,c\,\varepsilon_{\mathrm{EM}}\ln2}{\mathfrak h}-m_\ell-m_\nu\ }\qquad[\mathrm{BM.62}(4);\ \text{machine}];
\]

%dsp-level:0
\[
(2)\quad\textbf{the unconditional closure:}
\]

%dsp-level:0
\[
\qquad \Delta m<\frac{\ln(4/3)}{\mathfrak h}\ \ \text{at every}\ \varepsilon_{\mathrm{EM}}>0\ [\mathrm{BM.62}(6):\ \text{strict monotonicity}];
\]

%dsp-level:0
\[
m_\ell\ \ge\ \frac{2\ln2}{\mathfrak h}
\]

%dsp-level:0
\[
\qquad [\mathrm{FT.44}(4):\ m=E_{\mathrm{int}}\ln2/\mathfrak h;\ \mathrm{FT.28}\ \text{the minimal floor}\ E_{\mathrm{int}}=2;\ \mathrm{BM.13}(3):\ E_{\mathrm{int}}\ge2;\ \tfrac12<\ln2<1\ \text{enclosure}];
\]

%dsp-level:0
\[
\qquad m_\nu\ge0\ [\mathrm{EW.33}:\ \text{PSD}];
\]

%dsp-level:0
\[
\implies\boxed{\ Q_\beta\ <\ \frac{\ln(4/3)-2\ln2}{\mathfrak h}\ =\ -\,\frac{\ln3}{\mathfrak h}\ <\ 0\ \ \text{strictly}\ }
\]

%dsp-level:0
\[
\qquad [\text{machine:}\ 2\ln2-\ln\tfrac43=\ln3;\quad\ln\tfrac43<2\ln2\iff4<12]:
\]

%dsp-level:0
\[
\textbf{the}\ \beta\textbf{-decay vertex is kinematically closed at every admissible}\ \varepsilon_{\mathrm{EM}}>0,
\]

%dsp-level:0
\[
\ \text{on both sub- and super-critical branches};
\]

%dsp-level:0
\[
(3)\quad\textbf{the exact margin:}
\]

%dsp-level:0
\[
\qquad \boxed{\ |Q_\beta|>\frac{\ln3}{\mathfrak h}\ }
\]

%dsp-level:0
\[
\qquad [\text{machine: the certificate}\ 3>1]:
\]

%dsp-level:0
\[
\qquad \textbf{the closure margin is one exact}\ \ln3\ \text{of ledger content};
\]

%dsp-level:0
\[
\textbf{the reason, in one comparison:}
\]

%dsp-level:0
\[
\qquad \text{the transport excess is}\ \ln\tfrac43\ \text{nats and the minimal daughter costs}\ 2\ln2\ \text{nats}:
\]

%dsp-level:0
\[
\qquad \ln\tfrac43<2\ln2\iff\tfrac43<4\iff4<12
\]

%dsp-level:0
\[
\qquad [\text{machine}]:
\]

%dsp-level:0
\[
\textbf{a fraction of one bit cannot pay for two bits};
\]

%dsp-level:0
\[
(4)\quad\textbf{the deepening row:}
\]

%dsp-level:0
\[
\qquad \mathrm{BM.49}(1)\text{'s}\ Q_{\mathrm{corr}}=-(\Delta m_{\mathrm{EM}}+m_\ell+m_\nu)\ \text{and}\ \mathrm{BM.55}(5)\text{'s full-scope closure stand;}
\]

%dsp-level:0
\[
\qquad \textbf{the refined margin is}
\]

%dsp-level:0
\[
|Q_\beta|>\frac{\ln3+6\,c\,\varepsilon_{\mathrm{EM}}\ln2}{\mathfrak h}\ >\ \frac{\ln3}{\mathfrak h}
\]

%dsp-level:0
\[
\qquad [\text{machine: the coefficient census}]:
\]

%dsp-level:0
\[
\qquad \text{the EM summand still deepens the closure, and the transport summand does not reverse it};
\]

%dsp-level:0
\[
(5)\quad\textbf{the no-opening theorem survives:}
\]

%dsp-level:0
\[
\qquad \text{every EM correction of}\ \mathrm{Tier\ BM.12}\ \text{is}\ \propto q^2
\]

%dsp-level:0
\[
\ge 0\ \text{(an integral of a square) and raises the daughter's mass against the decay's requirement}\ [\mathrm{BM.49}(2)\ \text{re-cited}];
\]

%dsp-level:0
\[
\textbf{the transport refinement raises the parent's mass by at most}\ \ln(4/3)/\mathfrak h\ [(2)],
\]

%dsp-level:0
\[
\ \textbf{which is strictly below the minimal daughter content}\ [(3)]:
\]

%dsp-level:0
\[
\boxed{\ \text{no term of the extended decomposition opens the vertex}\ }\qquad[\text{machine: the two bounds conjoined}];
\]

%dsp-level:0
\[
(6)\quad\textbf{the}\ k{=}1\ \textbf{scope row (the daughters' own refinement):}
\]

%dsp-level:0
\[
\qquad \mathrm{Def\text{-}BM.13}(3)\text{'s census at}\ k=1\ \text{is}\ W(q)=m_q=(2,1,1)\ [\mathrm{EM.9};\ \text{machine}]:
\]

%dsp-level:0
\[
\Sigma_{\mathrm{cl}}(q)=\ln m_q=(\ln2,\,0,\,0)\ \ \text{at}\ q=(0,+1,-1):
\]

%dsp-level:0
\[
\qquad \text{the neutral line carries one extra bit, the charged lines none};
\]

%dsp-level:0
\[
\textbf{applied to the daughters:}
\]

%dsp-level:0
\[
\quad m_\nu\ \mapsto\ m_\nu+\tfrac{\ln2}{\mathfrak h}\ [\text{the}\ \nu\text{-channel is the EM-unactivated class}\ [\mathrm{BM.41}(3);\ \mathrm{FT.142}(1)]],
\]

%dsp-level:0
\[
\qquad m_\ell\ \mapsto\ m_\ell\ [\text{the charged member carries}\ m_q=1]:
\]

%dsp-level:0
\[
\text{the refinement of the daughters deepens the closure by exactly}\ \ln2/\mathfrak h\ [\text{machine}]:
\]

%dsp-level:0
\[
\qquad \textbf{it is recorded as a scope row and}\ \neg\textbf{needed for}\ (2)\text{'s verdict, which uses only}\ m_\nu\ge0;
\]

%dsp-level:0
\[
(7)\quad\textbf{what inverts and what does not:}
\]

%dsp-level:0
\[
\qquad \textbf{inverts on the sub-critical branch:}\ \text{the ordering clause}\ m_n<m_p\ [\mathrm{BM.55}(3)]\ \longrightarrow\ m_n>m_p\ [\mathrm{BM.62}(9)];
\]

%dsp-level:0
\[
\textbf{does not invert:}\ \text{the kinematic closure}\ [\mathrm{BM.41}(4)(\mathrm e),\ \mathrm{BM.49}(1)\text{--}(2),\ \mathrm{BM.55}(5)];
\]

%dsp-level:0
\[
\qquad \text{the vertex grammar and the forced members}\ [\mathrm{BM.41}(1),(3)];
\]

%dsp-level:0
\[
\qquad \text{the two-body impossibility}\ [\mathrm{BM.41}(2)];
\]

%dsp-level:0
\[
\text{the conditional three-body kinematics and the scale-invariant}\ E_{\max}:\langle E\rangle
\]

%dsp-level:0
\[
= 2:1\ [\mathrm{BM.41}(5);\ \mathrm{BM.49}(3),(4)];
\]

%dsp-level:0
\[
\qquad \text{the proton--antiproton degeneracy}\ [\mathrm{BM.28}(5);\ \mathrm{BM.30}(4):\ W(+1)=W(-1)\ \text{exactly}];
\]

%dsp-level:0
\[
(8)\quad\textbf{status:}
\]

%dsp-level:0
\[
\qquad \textbf{derived:}\ (1)\text{--}(7);
\]

%dsp-level:0
\[
\qquad \textbf{unconditional:}\ \text{the closure}\ (2),(3),(5);
\]

%dsp-level:0
\[
\qquad \textbf{conditional:}\ \text{the ordering}\ (7),\ \text{on}\ \mathrm{BM.62}(6)\text{'s branch};
\]

%dsp-level:0
\[
\neg\text{a measured lifetime, mass difference or endpoint enters}\ [\mathrm{P5}];
\]

%dsp-level:0
\[
\qquad \neg\text{external seconds}\ [\mathrm{BM.44}(3)];
\]

%dsp-level:0
\[
\qquad \neg\text{a P4 addition};
\]

%dsp-level:0
\[
\qquad \text{the open rows: the numerical value of}\ e,
\]

%dsp-level:0
\[
\ \text{the lepton floors}\ E_{\mathrm{int}}^{(\ell)},
\]

%dsp-level:0
\[
E_{\mathrm{int}}^{(\nu)}\ [\mathrm{BM.43}(2)],
\]

%dsp-level:0
\[
\ \text{the line index within the}\ \nu\text{-triplet}.
\]

\textbf{Req.} D2.9 (Parts 1–2); EW.33, EM.9 (Part 2); P5 (Parts 1–2); Def-H.32(3), FT.28, FT.42(1), FT.44(4), FT.142(1) (Part 3); Census FT.35(2) (row-wise); Def-BM.13(3), Def-BM.13(4), BM.13(3), BM.28(5), BM.30(4), BM.41(1)–(5), BM.43(2), BM.44(3), BM.46(4), BM.49(1)–(4), BM.55(3), BM.55(5), BM.62(4), BM.62(6), BM.62(9) (this part); exact $\ln$-ring identities with integer certificates, exact rational enclosures of $\ln2$, sign censuses [math; machine: \texttt{checks/\allowbreak{}verify\_\allowbreak{}H\_\allowbreak{}constituent\_\allowbreak{}transport\_\allowbreak{}splitting.py} §F].

\textbf{Pf.}

%dsp-level:0
\[
(1):\ \mathrm{BM.62}(4)\text{'s monomial substituted into}\ \mathrm{BM.41}(4)\text{'s definition}\ [\text{machine}]
\]

%dsp-level:0
\[
(2):\ \Delta m\ \text{strictly decreasing in}\ \varepsilon_{\mathrm{EM}}\ \text{with}\ \lim_{\varepsilon\to0^+}\Delta m=\ln(4/3)/\mathfrak h\ [\mathrm{BM.62}(6);\ \text{machine}];
\]

%dsp-level:0
\[
\qquad m_\ell\ge2\ln2/\mathfrak h\ \text{from}\ E_{\mathrm{int}}\ge2\ \text{and}\ \ln2>0\ [\text{machine}];
\]

%dsp-level:0
\[
\qquad m_\nu\ge0\ [\mathrm{EW.33}\ \text{PSD}]
\]

%dsp-level:0
\[
2\ln2-\ln\tfrac43=2\ln2-(2\ln2-\ln3)=\ln3\ [\text{machine: the prime-log ring}];
\]

%dsp-level:0
\[
\qquad \ln\tfrac43<2\ln2\iff\tfrac43<4\iff4<12\ [\text{machine}];
\]

%dsp-level:0
\[
\qquad \ln3>0\iff3>1\ [\text{machine}]
\]

%dsp-level:0
\[
(3):\ |Q_\beta|>-\bigl(\ln\tfrac43-2\ln2\bigr)/\mathfrak h=\ln3/\mathfrak h\ [\text{machine}];
\]

%dsp-level:0
\[
\qquad (4):\ \text{the coefficient census at}\ c,\varepsilon_{\mathrm{EM}},\ln2>0\ [\text{machine}]
\]

%dsp-level:0
\[
(5):\ \mathrm{BM.49}(2)\text{'s square-integral row re-cited};
\]

%dsp-level:0
\[
\qquad \ln\tfrac43<2\ln2\ [\text{machine}]\implies\text{the parent's excess cannot reach the daughter's floor}\ [\text{machine}]
\]

%dsp-level:0
\[
(6):\ W(q)
\]

%dsp-level:0
\[
= m_q\ \text{at}\ k
\]

%dsp-level:0
\[
= 1:\ (2+z+z^{-1})^1\ \text{has coefficients}\ (1,2,1)\ \text{at}\ Q
\]

%dsp-level:0
\[
= (-1,0,+1)\ [\text{machine: the convolution}];
\]

%dsp-level:0
\[
\qquad \ln m_q=(0,\ln2,0)\ \text{at}\ q=(\pm1,0)\ [\text{machine}]
\]

%dsp-level:0
\[
(7)\text{--}(8):\ \text{the clause-by-clause scan of}\ \mathrm{BM.41},
\]

%dsp-level:0
\[
\ \mathrm{BM.49},
\]

%dsp-level:0
\[
\ \mathrm{BM.55}\ \text{against the refined rest read}\ [\text{machine: the row census}]. \ \blacksquare
\]

\textbf{Phys.} Making the neutron heavier does not open the decay, and the reason is a comparison of two exact integers. The transport refinement lifts the parent by at most the logarithm of four thirds — a fraction of one bit — while the cheapest possible charged daughter costs two bits of ledger content, and the neutral member costs nothing less than zero. Four against twelve: the excess cannot reach the floor, and the closure holds at every coupling on both sides of the critical surface, with an exact margin of one $\ln3$. Read at the level of a single chain, the refinement even helps the closure, because the neutral line is the one that carries the extra transport bit and the neutrino is the neutral member. So the tier reverses an ordering and leaves a verdict intact: the neutron becomes the heavier baryon on the sub-critical branch, and the beta vertex stays shut — not because the two rows are degenerate, which they are no longer, but because the gap between them is smaller than the lightest thing the decay could emit.

\end{theorem}

\medskip\noindent\textit{\texorpdfstring{BM.64 — The status of the constituent transport splitting: adjudication and scope}{BM.64 - The status of the constituent transport splitting: adjudication and scope}.}\ Structural-constraints record: see Appendix B, Section B.1.23 of the Addendum.\par\medskip

\phantomsection
\section[Tier BM.15 — The Organic Bridge: the Operator EM Lift, the…]{\texorpdfstring{Tier BM.15 — The Organic Bridge: the Operator EM Lift, the Derived-Holonomy Attempt, the Constrained Joint Entropy, and the Beta-Decay Energy Audit}{Tier BM.15 - The Organic Bridge: the Operator EM Lift, the Derived-Holonomy Attempt, the Constrained Joint Entropy, and the Beta-Decay Energy Audit}}

\textbf{Locus.} This tier consumes the Tier BM.14 baseline unmodified and runs four constructions and one audit to exhaustion. The EM correction of the refined rest input is lifted from its scalar class row to an operator identity on the force spectrum ($\mathrm{BM.65}$, $\mathrm{BM.66}$). The full real holonomy is attempted as a lift from the edge-phase register, the force decomposition and the channel caps, and where the lift fails the exact obstruction is named and recorded ($\mathrm{BM.67}$). The constrained joint entropy is derived as a fiber cardinality, its geometric factor derived at stationarity and its general fiber named where the enumeration is missing ($\mathrm{BM.68}$). The three bridge inputs are converged and Gap 7 is reduced to its smallest priced residue ($\mathrm{BM.69}$). The beta-decay daughter floor is audited, the formed-closed dichotomy is proved, and every candidate opening term is adjudicated derived, priced or barred ($\mathrm{BM.70}$, $\mathrm{BM.71}$); $\mathrm{BM.72}$ adjudicates the tier.

\textbf{Standing rows.} The refutations of $\mathrm{BM.57}(8)$ and $\mathrm{BM.58}(4)$ stand closed and are $\neg$reopened as mass-splitting hypotheses; their residual content is admitted at its own loci — the $A_3$ continuation asymmetry generates the family triplet $[\mathrm{BM.57}(9);\ \mathrm{BM.60}(1)]$, the colour/EM seal asymmetry survives as a mask statement $[\mathrm{BM.58}(3),(5)]$, and the unique surviving class-dependent mass candidate is $W(Q)$ read through $\ln W(Q)$ $[\mathrm{BM.58}(6);\ \mathrm{BM.59}(2),(3)]$. $\alpha_B$ stays barred on $\mathrm{BM.61}(8)$'s four grounds. Holonomy is $\neg$declared here; joint entropy is $\neg$postulated here; no external neutron/proton ratio and no external $Q$-value enters anywhere; no Yukawa-type object exists or is added $[\mathrm{EW.15};\ \mathrm{FT.145}(3)]$. The tier adds no definition: every new symbol is an in-theorem abbreviation of derived objects.

\pseudohead{theorem}{\texorpdfstring{BM.65}{BM.65}}{\texorpdfstring{The baseline re-certified: the refined rest input, the exact branch, and the integer certificates}{The baseline re-certified: the refined rest input, the exact branch, and the integer certificates}}\label{tag:BM-65}
\textbf{Thm.}

%dsp-level:0
\[
(1)\quad\textbf{the refined rest input (frozen, re-certified):}
\]

%dsp-level:0
\[
\qquad \boxed{\ m_{\mathrm{rest}}(Q)\,\mathfrak h\ =\ E_{\mathrm{int}}\ln2\ +\ \ln W(Q)\ +\ 6\,c\,\varepsilon_{\mathrm{EM}}\,Q^2\ln2\ }
\]

%dsp-level:0
\[
\qquad [\mathrm{Def\text{-}BM.13}(5);\ \mathrm{Def\text{-}BM.12}(6);\ \mathrm{BM.46}(1),(3);\ \mathrm{BM.35}(4)]:
\]

%dsp-level:0
\[
E_{\mathrm{int}}=6
\]

%dsp-level:0
\[
\qquad [\mathrm{BM.17};\ \mathrm{BM.26};\ \text{Census FT.35}(2)];
\]

%dsp-level:0
\[
\qquad c\in\bigl\{\tfrac12,\tfrac35\bigr\}\ \text{priced, carried symbolically, adopted in none}
\]

%dsp-level:0
\[
\qquad [\mathrm{Def\text{-}BM.12}(4)];
\]

%dsp-level:0
\[
W(0)=20,
\]

%dsp-level:0
\[
\qquad W(\pm1)=15,
\]

%dsp-level:0
\[
\qquad \sum_{Q}W(Q)=4^3=64
\]

%dsp-level:0
\[
\qquad [\mathrm{Def\text{-}BM.13}(3);\ \mathrm{Def\text{-}BM.8}(6);\ \mathrm{BM.28}(2);\ \text{machine: the exhaustive}\ 4^3\ \text{end-sign census and the convolution agree}];
\]

%dsp-level:0
\[
(2)\quad\textbf{the class rows (frozen):}
\]

%dsp-level:0
\[
\qquad m_n\,\mathfrak h=\ln1280,
\]

%dsp-level:0
\[
\qquad m_p\,\mathfrak h=\ln960+6\,c\,\varepsilon_{\mathrm{EM}}\ln2,
\]

%dsp-level:0
\[
\qquad (m_n-m_p)\,\mathfrak h=\ln\tfrac43-6\,c\,\varepsilon_{\mathrm{EM}}\ln2=\Sigma_{\mathrm B}-6\,c\,\varepsilon_{\mathrm{EM}}\ln2
\]

%dsp-level:0
\[
\qquad [\mathrm{BM.62}(4);\ \text{machine}]:
\]

%dsp-level:0
\[
2^6\!\cdot\!20=1280,
\]

%dsp-level:0
\[
\qquad 2^6\!\cdot\!15=960,
\]

%dsp-level:0
\[
\qquad \ln20-\ln15=\ln\tfrac43=2\ln2-\ln3>0
\]

%dsp-level:0
\[
\qquad [\text{machine: the prime-log ring};\ \text{the certificate}\ 4>3];
\]

%dsp-level:0
\[
(3)\quad\textbf{the critical strength and the trichotomy (frozen):}
\]

%dsp-level:0
\[
\qquad \boxed{\ \varepsilon_{\mathrm{crit}}(c)\ :=\ \frac{\ln(4/3)}{6\,c\,\ln2}\ =\ \frac{\Sigma_{\mathrm B}}{6\,c\,\ln2}\ =\ \frac{2-\log_23}{6\,c}\ }
\]

%dsp-level:0
\[
\qquad [\mathrm{BM.62}(6);\ \text{machine: the unique zero}]:
\]

%dsp-level:0
\[
\varepsilon_{\mathrm{crit}}\!\left(\tfrac12\right)=\frac{2-\log_23}{3},\qquad\varepsilon_{\mathrm{crit}}\!\left(\tfrac35\right)=\frac{5\,(2-\log_23)}{18}\qquad[\text{machine: exact substitution}];
\]

%dsp-level:0
\[
\varepsilon_{\mathrm{EM}}<\varepsilon_{\mathrm{crit}}(c)\iff\Delta m>0;
\]

%dsp-level:0
\[
\qquad \varepsilon_{\mathrm{EM}}=\varepsilon_{\mathrm{crit}}(c)\iff\Delta m=0;
\]

%dsp-level:0
\[
\qquad \varepsilon_{\mathrm{EM}}>\varepsilon_{\mathrm{crit}}(c)\iff\Delta m<0
\]

%dsp-level:0
\[
\qquad [\mathrm{BM.62}(6);\ \text{machine}];
\]

%dsp-level:0
\[
(4)\quad\textbf{the integer certificates (machine re-certified):}
\]

%dsp-level:0
\[
\qquad \boxed{\ \frac19\ <\ \varepsilon_{\mathrm{crit}}\!\left(\tfrac35\right)\ <\ \frac18\ <\ \varepsilon_{\mathrm{crit}}\!\left(\tfrac12\right)\ <\ \frac17\ }
\]

%tbl-level:0
{\setlength{\tabcolsep}{4pt}
\begin{longtable}{>{\raggedright\arraybackslash}p{\dimexpr0.2329\textwidth-2\tabcolsep\relax}|>{\raggedright\arraybackslash}p{\dimexpr0.3562\textwidth-2\tabcolsep\relax}|>{\raggedright\arraybackslash}p{\dimexpr0.4110\textwidth-2\tabcolsep\relax}}
\toprule
bound & reduction & certificate \\
\addlinespace[2pt]
\midrule
\endfirsthead
\toprule
bound & reduction & certificate \\
\addlinespace[2pt]
\midrule
\endhead
\bottomrule
\endlastfoot
$\varepsilon_{\mathrm{crit}}(\tfrac35)>\allowbreak \tfrac19$ & $5\,\allowbreak \Sigma_{\mathrm B}>\allowbreak 2\ln2\iff\allowbreak 8\ln2>\allowbreak 5\ln3$ & $2^{8}>\allowbreak 3^{5}\iff\allowbreak 256>\allowbreak 243$ \\
$\varepsilon_{\mathrm{crit}}(\tfrac35)<\allowbreak \tfrac18$ & $40\,\allowbreak \Sigma_{\mathrm B}<\allowbreak 18\ln2\iff\allowbreak 62\ln2<\allowbreak 40\ln3$ & $2^{31}<\allowbreak 3^{20}\iff\allowbreak 2147483648<\allowbreak 3486784401$ \\
$\varepsilon_{\mathrm{crit}}(\tfrac12)>\allowbreak \tfrac18$ & $8\,\allowbreak \Sigma_{\mathrm B}>\allowbreak 3\ln2\iff\allowbreak 13\ln2>\allowbreak 8\ln3$ & $2^{13}>\allowbreak 3^{8}\iff\allowbreak 8192>\allowbreak 6561$ \\
$\varepsilon_{\mathrm{crit}}(\tfrac12)<\allowbreak \tfrac17$ & $7\,\allowbreak \Sigma_{\mathrm B}<\allowbreak 3\ln2\iff\allowbreak 11\ln2<\allowbreak 7\ln3$ & $2^{11}<\allowbreak 3^{7}\iff\allowbreak 2048<\allowbreak 2187$ \\
\end{longtable}
}

%dsp-level:0
\[
[\mathrm{BM.62}(7)\ \text{re-certified};\ \text{machine: every row an exact comparison of two non-negative integers}];
\]

%dsp-level:0
\[
(5)\quad\textbf{the branch is non-empty and exact:}
\]

%dsp-level:0
\[
\qquad \varepsilon_{\mathrm{crit}}(c)>0\iff\Sigma_{\mathrm B}>0\iff4>3
\]

%dsp-level:0
\[
\qquad [\text{machine}]:
\]

%dsp-level:0
\[
\qquad \text{the sub-critical branch}\ \bigl(0,\ \varepsilon_{\mathrm{crit}}(c)\bigr)\ \text{is a non-empty exact interval};
\]

%dsp-level:0
\[
\textbf{the parameter-free bar:}
\]

%dsp-level:0
\[
\qquad \varepsilon_{\mathrm{EM}}\le\tfrac19\ \Longrightarrow\ \Delta m>0\ \ \text{at both priced}\ c
\]

%dsp-level:0
\[
\qquad [\mathrm{BM.62}(8);\ \text{machine: the witnesses}\ 4^3=64>54=3^3\!\cdot\!2,\quad4^5=1024>972=3^5\!\cdot\!2^2];
\]

%dsp-level:0
\[
(6)\quad\textbf{no numerical coupling exists:}
\]

%dsp-level:0
\[
\qquad \neg\exists\ \text{a theorem of Parts 1--4 fixing}\ \varepsilon_{\mathrm{EM}}\ \text{numerically}
\]

%dsp-level:0
\[
\qquad [\mathrm{EM.25}(7);\ \mathrm{BM.61}(3),(4);\ \mathrm{BM.62}(12)]:
\]

%dsp-level:0
\[
\textbf{the neutron-heavier result is conditional on the symbolic coupling being sub-critical,}\ \neg\textbf{on an imported physical value}
\]

%dsp-level:0
\[
\qquad [\mathrm{P5};\ \text{the hypothesis stands on the face of}\ \mathrm{BM.62}(9)\ \text{and of this row};\ \text{A §0.2a}];
\]

%dsp-level:0
\[
(7)\quad\textbf{the ratio discipline:}
\]

%dsp-level:0
\[
\qquad \frac{\ln1280}{\ln960}\ne\frac43\iff3\ln1280\ne4\ln960\iff1280^3\ne960^4
\]

%dsp-level:0
\[
\qquad [\text{machine: the two integers differ}]:
\]

%dsp-level:0
\[
\neg\text{a route ratio is read as a mass ratio}\qquad[\mathrm{BM.37}(6);\ \mathrm{Def\text{-}BM.13}(8);\ \mathrm{BM.62}(5)];
\]

%dsp-level:0
\[
(8)\quad\textbf{status:}
\]

%dsp-level:0
\[
\qquad \textbf{re-certified:}\ (1)\text{--}(7);
\]

%dsp-level:0
\[
\qquad \text{the rows of}\ \mathrm{Def\text{-}BM.13}(5),\ \mathrm{BM.62}(4),(6)\text{--}(9),\ \mathrm{BM.63}(2)\ \text{stand unmodified};
\]

%dsp-level:0
\[
\qquad \text{this tag adds no content beyond the machine re-certification};
\]

%dsp-level:0
\[
\neg\text{a P4 addition},\quad\neg\text{a new dimensionful scale}\ [\mathrm{P1}],\quad\neg\text{a measured value}\ [\mathrm{P5}].
\]

\textbf{Req.} CT2.1 (Part 2); EM.25(7) (Part 2); P1, P5 (Parts 1–2); A §0.2a (the hypothesis at the point of use); Census FT.35(2) (row-wise); Def-BM.8(6), Def-BM.12(4), Def-BM.12(6), Def-BM.13(3), Def-BM.13(4), Def-BM.13(5), Def-BM.13(8), BM.17, BM.26, BM.28(2), BM.35(4), BM.37(6), BM.46(1), BM.46(3), BM.61(3), BM.61(4), BM.62(4), BM.62(5), BM.62(6), BM.62(7), BM.62(8), BM.62(9), BM.62(12), BM.63(2) (this part); exact prime-log-ring arithmetic with integer power certificates, exhaustive end-sign censuses [math; machine: \texttt{checks/\allowbreak{}verify\_\allowbreak{}H\_\allowbreak{}organic\_\allowbreak{}bridge\_\allowbreak{}post\_\allowbreak{}bm14.py} §B, §C].

\textbf{Pf.}

%dsp-level:0
\[
(1)\text{--}(2):\ \text{the frozen rows re-derived by machine: the}\ 4^3\ \text{end-sign census and the generating-function convolution re-enumerated and compared class by class};
\]

%dsp-level:0
\[
\qquad 2^6\!\cdot\!20=1280,\ 2^6\!\cdot\!15=960\ [\text{machine: integers}];
\]

%dsp-level:0
\[
\qquad \ln20-\ln15=\ln\tfrac43\ [\text{machine: the prime-log ring}];\ 4>3\ [\text{machine}]
\]

%dsp-level:0
\[
(3):\ \Delta m=0\iff\Sigma_{\mathrm B}=6\,c\,\varepsilon_{\mathrm{EM}}\ln2\ \text{solved in the}\ \ln\text{-ring}\ [\text{machine}];
\]

%dsp-level:0
\[
\qquad \varepsilon_{\mathrm{crit}}\!\left(\tfrac12\right)=\tfrac{\Sigma_{\mathrm B}}{3\ln2},\ \ \varepsilon_{\mathrm{crit}}\!\left(\tfrac35\right)=\tfrac{5\,\Sigma_{\mathrm B}}{18\ln2}\ [\text{machine: exact Fractions}];
\]

%dsp-level:0
\[
\qquad \log_23\ \text{form:}\ \tfrac{\Sigma_{\mathrm B}}{\ln2}=2-\log_23\ [\text{machine: the prime-log ring}]
\]

%dsp-level:0
\[
(4):\ \text{cross-multiplication by}\ 6\,
\]

%dsp-level:0
\[
c\ln2>0\ \text{reduces each bound to}\ a\,
\]

%dsp-level:0
\[
\Sigma_{\mathrm B}\gtrless b\ln2
\]

%dsp-level:0
\[
\iff (4/3)^a\gtrless2^b
\]

%dsp-level:0
\[
\iff 4^a\gtrless3^a2^b:
\]

%dsp-level:0
\[
\quad (c,r)
\]

%dsp-level:0
\[
= \left(\tfrac35,\tfrac19\right):\ 2^8\!\cdot\!1\ \text{vs}\ 3^5:\ 256>243;
\]

%dsp-level:0
\[
\ \left(\tfrac35,\tfrac18\right):\ 2^{31}\ \text{vs}\ 3^{20};
\]

%dsp-level:0
\[
\ \left(\tfrac12,\tfrac18\right):\ 2^{13}\ \text{vs}\ 3^8;
\]

%dsp-level:0
\[
\ \left(\tfrac12,\tfrac17\right):\ 2^{11}\ \text{vs}\ 3^7\ [\text{machine: four integer comparisons}]
\]

%dsp-level:0
\[
(5):\ \Sigma_{\mathrm B}>0\iff4>3\ [\text{machine}];
\]

%dsp-level:0
\[
\qquad \text{the two witnesses}\ 64>54,\ 1024>972\ [\text{machine}];
\]

%dsp-level:0
\[
\qquad (6):\ \mathrm{EM.25}(7),\ \mathrm{BM.61}(4)\ \text{re-cited}
\]

%dsp-level:0
\[
(7):\ 1280^3\ \text{and}\ 960^4\ \text{computed exactly and compared}\ [\text{machine: the two integers differ}];
\]

%dsp-level:0
\[
\qquad (8):\ \text{assembly}. \ \blacksquare
\]

\textbf{Phys.} The baseline stands on integer certificates. The refined rest input assembles three derived registers — the count, the charge-filtered end census, and the scar's field energy — into one row per class, and the class difference has a single exact zero: the critical strength, enclosed between one ninth and one seventh by four comparisons of powers of two against powers of three. The branch on which the neutral baryon is the heavier one is a non-empty exact interval, certified from below by two hundred fifty-six against two hundred forty-three, and no theorem of Parts 1–4 fixes the coupling numerically — so the ordering result is conditional on the symbol being sub-critical, never on an imported value. The ratio discipline holds with it: the mass ratio is the logarithm ratio of twelve-eighty over nine-sixty, the certificates that twelve-eighty cubed differs from nine-sixty to the fourth keep the measure price away from the mass read, and nothing measured enters anywhere.

\begin{theorem}[{\texorpdfstring{BM.66 — The operator lift of the EM correction: the charge-blind trace, the quadratic invariant, and the assembled operator rest input}{BM.66 - The operator lift of the EM correction: the charge-blind trace, the quadratic invariant, and the assembled operator rest input}}]\label{tag:BM-66}
\textbf{Thm.}

%dsp-level:0
\[
(1)\quad\textbf{the operator and its spectrum:}
\]

%dsp-level:0
\[
\qquad \widehat F\ :=\ \mathbf F_{\mathrm{ext}}(\mathbf c,\widehat Q)\ =\ h(\mathbf c)\otimes I_2\ +\ \widehat Q\,I_2\otimes\hat\kappa,
\]

%dsp-level:0
\[
\qquad h(\mathbf c)=\tfrac12\bigl(c_0I_2+\mathbf c\cdot\boldsymbol\sigma\bigr)
\]

%dsp-level:0
\[
\qquad [\mathrm{BM.22}(1);\ \mathrm{FT.129}(5);\ \text{the charge read carried as the operator}\ \widehat Q]:
\]

%dsp-level:0
\[
\operatorname{spec}\widehat Q\ =\ \{-1,\,0,\,+1\}\ \ \text{on the materialized classes}
\]

%dsp-level:0
\[
\qquad [\mathrm{EM.9};\ \mathrm{D3.4};\ \mathrm{Def\text{-}BM.8}(1),(5);\ \mathrm{Def\text{-}BM.3}(2),(3)];
\]

%dsp-level:0
\[
\lambda_{s,t}\ =\ \tfrac12\bigl(c_0+sC\bigr)+t\,Q,
\]

%dsp-level:0
\[
\qquad C:=|\mathbf c|,
\]

%dsp-level:0
\[
\qquad s,t\in\{\pm1\}
\]

%dsp-level:0
\[
\qquad [\mathrm{BM.22}(3);\ \mathrm{FT.130}(5);\ \text{machine: the monic quartic from the four roots at exact rational witnesses}];
\]

%dsp-level:0
\[
\textbf{the promotion conditionality is inherited,}\ \neg\textbf{new:}
\]

%dsp-level:0
\[
\qquad \text{the flipping block exists under}\ \mathrm{Def\text{-}H.43}\text{'s declared read promotion}
\]

%dsp-level:0
\[
\qquad [\mathrm{BM.22}(5)];
\]

%dsp-level:0
\[
(2)\quad\textbf{the linear trace is charge-blind:}
\]

%dsp-level:0
\[
\qquad \boxed{\ \operatorname{Tr}\widehat F\ =\ \sum_{s,t\in\{\pm1\}}\lambda_{s,t}\ =\ 2\,c_0\ }
\]

%dsp-level:0
\[
\qquad [\text{machine: the symbolic sum};\ \textstyle\sum_{s,t}t\,Q=0,\ \sum_{s,t}s\,C=0]:
\]

%dsp-level:0
\[
\textbf{the operator form of the uniqueness row's trace line:}
\]

%dsp-level:0
\[
\qquad \text{no trace read of}\ \widehat F\ \text{itself splits the two classes}
\]

%dsp-level:0
\[
\qquad [\mathrm{BM.58}(6):\ \textstyle\sum_{s,t}tq=0;\ \ \mathrm{BM.1}(2):\ \text{the trace-only read}];
\]

%dsp-level:0
\[
(3)\quad\textbf{the quadratic invariant:}
\]

%dsp-level:0
\[
\qquad \operatorname{Tr}\widehat F^2\ =\ \sum_{s,t}\lambda_{s,t}^2\ =\ c_0^2+C^2+4\,\widehat Q^2
\]

%dsp-level:0
\[
\qquad [\text{machine: the polynomial identity in}\ \mathbb Q[c_0,C,Q]];
\]

%dsp-level:0
\[
\widehat F_0\ :=\ \widehat F\big|_{\widehat Q=0}:\qquad\operatorname{Tr}\widehat F_0^2\ =\ c_0^2+C^2;\qquad\boxed{\ I_F\ :=\ \operatorname{Tr}\widehat F^2-\operatorname{Tr}\widehat F_0^2\ =\ 4\,\widehat Q^2\ }\qquad[\text{machine}]:
\]

%dsp-level:0
\[
\textbf{the first trace invariant carrying charge;}
\]

%dsp-level:0
\[
\qquad I_F=(0,\,4,\,4)\ \ \text{at}\ Q=(0,\,\pm1)
\]

%dsp-level:0
\[
\qquad [\text{machine}]:
\]

%dsp-level:0
\[
\qquad \mathrm{BM.58}(5)\text{'s monopole selector}\ Q^2=(0,1,1)\ \textbf{recovered as a trace row};
\]

%dsp-level:0
\[
(4)\quad\textbf{the operator EM correction:}
\]

%dsp-level:0
\[
\qquad \boxed{\ \Delta\bigl(\widehat M_{\mathrm{rest}}\mathfrak h\bigr)_{\mathrm{EM}}\ :=\ \frac{E_{\mathrm{int}}}4\;c\,\varepsilon_{\mathrm{EM}}\ln2\;\Bigl(\operatorname{Tr}\widehat F^2-\operatorname{Tr}\widehat F_0^2\Bigr)\ }
\]

%dsp-level:0
\[
\qquad \text{with}\ \frac{E_{\mathrm{int}}}4=\frac64=\frac32\ \ \text{at}\ E_{\mathrm{int}}=6
\]

%dsp-level:0
\[
\qquad [\mathrm{BM.17};\ \text{Census FT.35}(2);\ \text{machine}]:
\]

%dsp-level:0
\[
\textbf{the coefficient is derived,}\ \neg\textbf{chosen:}\qquad\frac{E_{\mathrm{int}}}4\cdot I_F\ =\ \frac{E_{\mathrm{int}}}4\cdot4\,\widehat Q^2\ =\ E_{\mathrm{int}}\,\widehat Q^2\qquad[\text{machine}]:
\]

%dsp-level:0
\[
\Delta\bigl(\widehat M_{\mathrm{rest}}\mathfrak h\bigr)_{\mathrm{EM}}\ =\ c\,\varepsilon_{\mathrm{EM}}\,\widehat Q^2\,E_{\mathrm{int}}\ln2\ =\ c\,\varepsilon_{\mathrm{EM}}\,\widehat Q^2\,m_{\mathrm{hid}}\,\mathfrak h\ =\ E^{\mathrm{self}}_{\mathrm{EM}}(\widehat Q)\,\mathfrak h
\]

%dsp-level:0
\[
\qquad [\mathrm{BM.46}(1)\ \text{unified form};\ \text{machine: one monomial}];
\]

%dsp-level:0
\[
\textbf{on charge eigenstates:}
\]

%dsp-level:0
\[
\qquad Q=0:\quad\Delta\bigl(\widehat M_{\mathrm{rest}}\mathfrak h\bigr)_{\mathrm{EM}}=0;
\]

%dsp-level:0
\[
\qquad Q=\pm1:\quad\Delta\bigl(\widehat M_{\mathrm{rest}}\mathfrak h\bigr)_{\mathrm{EM}}=6\,c\,\varepsilon_{\mathrm{EM}}\ln2
\]

%dsp-level:0
\[
\qquad [\text{machine}]:
\]

%dsp-level:0
\[
\textbf{the scalar row of}\ \mathrm{Def\text{-}BM.13}(5)\ \textbf{is reproduced exactly on all three materialized classes};
\]

%dsp-level:0
\[
(5)\quad\textbf{the assembled operator rest input:}
\]

%dsp-level:0
\[
\qquad \boxed{\ \widehat M_{\mathrm{rest}}\mathfrak h\ =\ E_{\mathrm{int}}\ln2\ +\ \ln W(\widehat Q)\ +\ \frac32\,c\,\varepsilon_{\mathrm{EM}}\ln2\Bigl(\operatorname{Tr}\widehat F^2-\operatorname{Tr}\widehat F_0^2\Bigr)\ }
\]

%dsp-level:0
\[
\ln W(\widehat Q)\ :=\ \sum_{Q\in\operatorname{spec}}\bigl(\ln W(Q)\bigr)\,P_Q
\]

%dsp-level:0
\[
\qquad [\text{functional calculus on the charge eigenspaces};\ \text{math: the spectral theorem for a diagonal operator on a finite charge class}];
\]

%dsp-level:0
\[
\textbf{the eigenvalue rows are exactly}\ \mathrm{BM.62}(4)\text{'s:}\qquad Q=0:\ \ \ln1280;\qquad Q=\pm1:\ \ \ln960+6\,c\,\varepsilon_{\mathrm{EM}}\ln2\qquad[\text{machine}];
\]

%dsp-level:0
\[
\textbf{the CPT row:}
\]

%dsp-level:0
\[
\qquad \omega:\ \widehat Q\mapsto-\widehat Q\ \ \text{fixes}\ \ W(\widehat Q)\ \ [W(Q)=W(-Q);\ \mathrm{Def\text{-}BM.8}(7)]\ \ \text{and fixes}\ \ I_F\ \ [(3):\ \widehat Q^2\ \text{only}]:
\]

%dsp-level:0
\[
\textbf{the proton--antiproton degeneracy is an operator identity}\qquad[\mathrm{BM.28}(5);\ \mathrm{BM.30}(4);\ \text{machine}];
\]

%dsp-level:0
\[
(6)\quad\textbf{the no-double-count guard:}\qquad\textbf{three derived routes to one monomial:}\qquad\begin{array}{l|l}\text{the sector census} & \mathrm{BM.52}(3),(4):\ \text{five zeros and one monomial}\\\text{the field integral} & \mathrm{BM.46}(1);\ \mathrm{Def\text{-}BM.12}(5)\\\text{the quadratic trace} & (3),(4)\ \text{of this tag}\end{array}\qquad[\text{machine: the three agree on every charge eigenstate}]:
\]

%dsp-level:0
\[
\mathrm{BM.52}(5)\text{'s guard extends: in any decomposition the EM contribution is counted once}
\]

%dsp-level:0
\[
\qquad [\mathrm{BM.62}(1)\ \text{stands}];
\]

%dsp-level:0
\[
\textbf{the alphabet census:}
\]

%dsp-level:0
\[
\qquad \text{the trace subtraction removes}\ c_0^2+C^2\ \text{exactly}\ (c_0,C\ \text{are}\ \neg\text{free symbols of}\ I_F)
\]

%dsp-level:0
\[
\qquad [\text{machine}];
\]

%dsp-level:0
\[
\text{the free-symbol set of}\ \Delta\bigl(\widehat M_{\mathrm{rest}}\mathfrak h\bigr)_{\mathrm{EM}}\ \text{is}\ \{c,\varepsilon_{\mathrm{EM}},\ln2,\widehat Q\}\ \subseteq\ \mathrm{BM.62}(11)\text{'s}\ \Delta m_{\mathrm{EM}}\ \text{alphabet under}\ e^2
\]

%dsp-level:0
\[
= 4\pi\varepsilon_{\mathrm{EM}},
\]

%dsp-level:0
\[
\ \Delta x_{\mathrm B}^{-1}\qquad = m_{\mathrm{hid}}
\]

%dsp-level:0
\[
\qquad [\mathrm{EW.35};\ \mathrm{Def\text{-}BM.4}(1);\ \text{machine: the token census}];
\]

%dsp-level:0
\[
\textbf{the disjointness from the transport alphabet}\ \{m_q,W(Q),\ln2,\mathfrak h\}\ \text{holds except for the shared}\ \ln2
\]

%dsp-level:0
\[
\qquad [\mathrm{BM.59}(8);\ \mathrm{BM.37}(2);\ \text{machine}];
\]

%dsp-level:0
\[
(7)\quad\textbf{the uniqueness row is respected:}
\]

%dsp-level:0
\[
\qquad \mathrm{BM.58}(6)\text{'s table gains}\ \neg\text{a new row:}
\]

%dsp-level:0
\[
\qquad I_F\ \text{is the operator species of the adopted}\ E^{\mathrm{self}}_{\mathrm{EM}}\ \text{entry}\ [\mathrm{Def\text{-}BM.12}(6);\ (4)],
\]

%dsp-level:0
\[
\ \text{with the same monomial and the same negative sign for}\ m_n-m_p\ [\mathrm{BM.58}(5);\ \text{machine}];
\]

%dsp-level:0
\[
F_{LR}=\widehat QI_2\text{'s row}\ \text{``}\,0\ \text{on a trace read}\,\text{''}\ \text{is the derivation of}\ (2);
\]

%dsp-level:0
\[
\qquad \textbf{the unique surviving class-dependent rest-slot candidate remains}\ W(Q)\ \text{read through}\ \ln W(Q)
\]

%dsp-level:0
\[
\qquad [\mathrm{BM.59}(2),(3)];
\]

%dsp-level:0
\[
(8)\quad\textbf{status:}
\]

%dsp-level:0
\[
\qquad \textbf{derived:}\ (1)\text{--}(7);
\]

%dsp-level:0
\[
\qquad \textbf{conditional on}\ \mathrm{Def\text{-}H.43}\text{'s promotion as}\ \mathrm{BM.22}(5)\ \text{records it}\ (\neg\text{a new declaration});
\]

%dsp-level:0
\[
c\ \text{stays priced,}\ \neg\text{adopted}\ [\mathrm{Def\text{-}BM.12}(4)];
\]

%dsp-level:0
\[
\qquad \textbf{the new coefficients are}\ \bigl\{\tfrac32=E_{\mathrm{int}}/4,\ 4\bigr\}:\ \text{a count row and a trace normalization}\ \text{—}\ \neg\text{a P4 addition},
\]

%dsp-level:0
\[
\ \neg\text{a fitting parameter}\ [\text{machine: the coefficient census}];
\]

%dsp-level:0
\[
\neg\text{a new dimensionful scale}\ [\mathrm{P1}];\qquad\neg\text{a measured value}\ [\mathrm{P5}];\qquad\text{§Cal-BM stays empty}.
\]

\textbf{Req.} D2.9, D3.4 (Parts 1–2); CT2.1, EM.9, EW.35 (Part 2); P1, P5 (Parts 1–2); Def-H.43, FT.129(2), FT.129(5), FT.130(5) (Part 3); Census FT.35(2) (row-wise); Def-BM.3(2), Def-BM.3(3), Def-BM.4(1), Def-BM.8(1), Def-BM.8(5), Def-BM.8(7), Def-BM.12(4), Def-BM.12(5), Def-BM.12(6), Def-BM.13(5), BM.1(2), BM.17, BM.22(1), BM.22(2), BM.22(3), BM.22(5), BM.26, BM.28(5), BM.30(4), BM.37(2), BM.46(1), BM.46(3), BM.52(3), BM.52(4), BM.52(5), BM.58(5), BM.58(6), BM.59(2), BM.59(3), BM.59(8), BM.62(1), BM.62(4), BM.62(11) (this part); exact multivariate polynomial identities over $\mathbb Q[c_0,C,Q,c,\varepsilon_{\mathrm{EM}},\ln2]$, matrix traces over complex-with-rational-parts at exact witnesses, spectral calculus for diagonal operators on finite charge classes, token censuses [math; machine: \texttt{checks/\allowbreak{}verify\_\allowbreak{}H\_\allowbreak{}organic\_\allowbreak{}bridge\_\allowbreak{}post\_\allowbreak{}bm14.py} §A].

\textbf{Pf.}

%dsp-level:0
\[
(1):\ \mathrm{BM.22}(1)\text{--}(3)\ \text{re-cited with}\ q\ \text{carried as}\ \widehat Q;
\]

%dsp-level:0
\[
\qquad \text{the charpoly}\ \bigl((y-u)^2-Q^2\bigr)\bigl((y+u)^2-Q^2\bigr),
\]

%dsp-level:0
\[
\ y
\]

%dsp-level:0
\[
= x-\tfrac{c_0}2,
\]

%dsp-level:0
\[
\ u^2
\]

%dsp-level:0
\[
= \tfrac{C^2}4\ \text{built from the four roots and compared to the Faddeev--LeVerrier charpoly at witnesses with}\ C^2\in\{0,\tfrac14,1,\tfrac49,\tfrac{9}{49}\}\ [\text{machine: exact over the rationals}]
\]

%dsp-level:0
\[
(2):\ \sum_{s,t}\bigl[\tfrac12(c_0+sC)+tQ\bigr]
\]

%dsp-level:0
\[
= 4\cdot\tfrac12c_0+\tfrac C2\sum_{s,t}s+Q\sum_{s,t}t
\]

%dsp-level:0
\[
= 2c_0+0+0\ [\text{machine: the symbolic sum}];
\]

%dsp-level:0
\[
\qquad \text{the matrix row:}\ \operatorname{Tr}\bigl(h\otimes I_2+QI_2\otimes\hat\kappa\bigr)
\]

%dsp-level:0
\[
= 2\operatorname{tr}h+2Q\operatorname{tr}\hat\kappa
\]

%dsp-level:0
\[
= 2c_0\ [\hat\kappa\ \text{traceless};\ \text{machine: exact traces at every witness}\times\text{ternary}\ Q]
\]

%dsp-level:0
\[
(3):\ \lambda_{s,t}^2=\tfrac14(c_0+sC)^2+tQ(c_0+sC)+t^2Q^2:
\]

%dsp-level:0
\[
\qquad \sum_{s,t}\tfrac14(c_0+sC)^2=\tfrac12\bigl[(c_0+C)^2+(c_0-C)^2\bigr]=c_0^2+C^2;
\]

%dsp-level:0
\[
\qquad \sum_{s,t}tQ(c_0+sC)=0;
\]

%dsp-level:0
\[
\qquad \sum_{s,t}t^2Q^2=4Q^2\ [\text{machine: the polynomial identity}];
\]

%dsp-level:0
\[
\text{the matrix route agrees:}\ \operatorname{Tr}F^2
\]

%dsp-level:0
\[
= 2\operatorname{tr}\bigl(h^2\bigr)+2Q\operatorname{tr}(h)\operatorname{tr}(\hat\kappa)+Q^2\operatorname{tr}(I_2)\operatorname{tr}\bigl(\hat\kappa^2\bigr)
\]

%dsp-level:0
\[
= \bigl(c_0^2+C^2\bigr)+0+4Q^2\ [\hat\kappa\ \text{the ray swap},\ \hat\kappa^2=I_2,\ \operatorname{tr}\hat\kappa=0\ [\mathrm{FT.129}(2)];\ \operatorname{tr}h=c_0;\ \text{machine at every witness}]
\]

%dsp-level:0
\[
(4):\ \tfrac{E_{\mathrm{int}}}4\cdot4Q^2=E_{\mathrm{int}}Q^2\ [\text{machine}];
\]

%dsp-level:0
\[
\qquad c\,\varepsilon_{\mathrm{EM}}Q^2E_{\mathrm{int}}\ln2=c\,\varepsilon_{\mathrm{EM}}Q^2m_{\mathrm{hid}}\mathfrak h\ [\mathrm{BM.46}(1)\ \text{re-cited};\ \text{machine: one monomial}];
\]

%dsp-level:0
\[
\qquad \text{the eigenstate rows at}\ Q^2\in\{0,1\}\ [\text{machine}]
\]

%dsp-level:0
\[
(5):\ P_Q\ \text{the spectral projections of}\ \widehat Q;
\]

%dsp-level:0
\[
\ \ \ln W(\widehat Q)\ \text{diagonal with the entries}\ (\ln20,\ \ln15,\ \ln15)\ [\text{machine}];
\]

%dsp-level:0
\[
\qquad \text{the eigenvalue sum}\ 6\ln2+\ln20
\]

%dsp-level:0
\[
= \ln1280,
\]

%dsp-level:0
\[
\ 6\ln2+\ln15+6c\varepsilon_{\mathrm{EM}}\ln2
\]

%dsp-level:0
\[
= \ln960+6c\varepsilon_{\mathrm{EM}}\ln2\ [\text{machine: the prime-log ring}];
\]

%dsp-level:0
\[
\qquad \omega\text{-row:}\ W(-Q)=W(Q),\ (-Q)^2=Q^2\ [\text{machine}]
\]

%dsp-level:0
\[
(6):\ \text{the three routes compared monomial by monomial on the three materialized classes}\ [\text{machine}];
\]

%dsp-level:0
\[
\qquad \text{the symbol census of}\ I_F\ \text{and of the correction}\ [\text{machine:}\ \{c,\varepsilon_{\mathrm{EM}},\ln2,Q\}];
\]

%dsp-level:0
\[
\qquad e^2=4\pi\varepsilon_{\mathrm{EM}}\ [\mathrm{EW.35};\ \text{machine}];\ \Delta x_{\mathrm B}^{-1}=m_{\mathrm{hid}}\ [\mathrm{Def\text{-}BM.4}(1);\ \text{machine}]
\]

%dsp-level:0
\[
(7)\text{--}(8):\ \mathrm{BM.58}(5),(6),\ \mathrm{BM.59}(2),(3)\ \text{re-cited against}\ (2)\text{--}(4);
\]

%dsp-level:0
\[
\qquad \text{the coefficient census}\ \bigl\{\tfrac64,\ 4\bigr\}\ \text{and the standing-rule scans}\ [\text{machine}]. \ \blacksquare
\]

\textbf{Phys.} The scalar row of the EM correction is the eigenvalue row of an operator identity, and the operator is the force matrix the corpus already carried. Its linear trace reads only the shared block — two c-naught, blind to charge — which is the operator form of the uniqueness row's old trace line: no trace of the force itself can split the two classes. The first trace that carries charge is the quadratic one, and its charge content is exactly four Q-squared, so the monopole selector comes back as a trace law. Multiplied by the count row's own coefficient — the internal edge count over four, three halves at six edges — the operator correction reproduces the field-energy monomial on every materialized class, joins the sector census and the field integral as a third route to one and the same term, and keeps the proton–antiproton degeneracy exact, because only the square of the charge enters and the census is conjugation-symmetric. Nothing new is counted: the trace subtraction removes the shared blocks, the symbol alphabet collapses to the adopted row's own, and the distribution factor stays priced in two exact values and adopted in none.

\end{theorem}

\pseudohead{theorem}{\texorpdfstring{BM.67}{BM.67}}{\texorpdfstring{The holonomy lift attempted: the derived $\mathbb Z_2$ sector, the angle-free force census, and the two-torsion obstruction}{The holonomy lift attempted: the derived sector, the angle-free force census, and the two-torsion obstruction}}\label{tag:BM-67}
\textbf{Thm.}

%dsp-level:0
\[
(1)\quad\textbf{the register and the derived discrete phase:}
\]

%dsp-level:0
\[
\qquad \alpha_e\in\{0,\pi\},
\]

%dsp-level:0
\[
\qquad \sigma_e=e^{i\alpha_e}\in\{+1,-1\}
\]

%dsp-level:0
\[
\qquad [\mathrm{A4};\ \mathrm{A7};\ \mathrm{CT0.9};\ \mathrm{Def\text{-}EM.20};\ \mathrm{EM.20}(1),(2):\ \text{registration data determined at crystallization},\ \neg\text{declared constants}]:
\]

%dsp-level:0
\[
\text{virtual edges carry}\ \neg\alpha_e\qquad[\mathrm{EM.20}(3)];\qquad[\alpha_e]=\mathfrak q^0\ \ \text{in every calibration}\qquad[\mathrm{EM.21}];
\]

%dsp-level:0
\[
\textbf{for an oriented chain}\ c:
\]

%dsp-level:0
\[
\qquad \Phi(c)\ :=\ \sum_{e\in c}\operatorname{or}(e)\,\alpha_e\ \ (\mathrm{mod}\ 2\pi),
\]

%dsp-level:0
\[
\qquad U(c)\ :=\ \prod_{e\in c}\sigma_e^{\operatorname{or}(e)}\ =\ e^{i\Phi(c)}
\]

%dsp-level:0
\[
\qquad [\mathrm{D3.5}\ \text{weight grammar};\ \text{machine}];
\]

%dsp-level:0
\[
(2)\quad\textbf{the bare holonomy is}\ \mathbb Z_2\ \textbf{(theorem):}
\]

%dsp-level:0
\[
\qquad \sigma^{-1}=\sigma\ \ \text{on}\ \{\pm1\}\implies\text{the orientation drops out of every closed product}:
\]

%dsp-level:0
\[
\boxed{\ U(\gamma)\ =\ \prod_{e\in\gamma}\sigma_e\ =\ (-1)^{N_\pi(\gamma)},\qquad N_\pi(\gamma):=\#\{e\in\gamma:\ \alpha_e=\pi\}\ }
\]

%dsp-level:0
\[
\qquad [\text{machine: exhaustive over every sign assignment on loops of length}\ 2,\dots,7,\ \text{both orientations}]:
\]

%dsp-level:0
\[
U(\gamma_1\!\cdot\!\gamma_2)=U(\gamma_1)\,U(\gamma_2),
\]

%dsp-level:0
\[
\qquad U(\bar\gamma)=U(\gamma)
\]

%dsp-level:0
\[
\qquad [\text{machine: concatenation and reversal censuses}]:
\]

%dsp-level:0
\[
\qquad \textbf{the Wilson sector of the register is exactly binary};
\]

%dsp-level:0
\[
(3)\quad\textbf{the parity lock, lifted to the composite:}
\]

%dsp-level:0
\[
\qquad \sum_i\aleph_{v_i}\equiv|q|\ (\mathrm{mod}\ 2),
\]

%dsp-level:0
\[
\qquad \sum_i\aleph_{v_i}\ \ge\ |q|
\]

%dsp-level:0
\[
\qquad [\mathrm{Def\text{-}BM.13}(2);\ \mathrm{EM.9b};\ \mathrm{D3.3};\ \text{machine: exhaustive over every sign sequence},\ k=2,\dots,7]:
\]

%dsp-level:0
\[
Q=\sum_iq_i
\]

%dsp-level:0
\[
\qquad [\mathrm{Def\text{-}BM.8}(1),(3)]:
\]

%dsp-level:0
\[
\qquad |q_i|\equiv q_i\ (\mathrm{mod}\ 2)\implies\textstyle\sum_i\aleph^{(i)}\equiv\sum_i|q_i|\equiv Q\equiv|Q|\ (\mathrm{mod}\ 2);
\]

%dsp-level:0
\[
\qquad \textstyle\sum_i\aleph^{(i)}\ge\sum_i|q_i|\ge\bigl|\sum_iq_i\bigr|=|Q|
\]

%dsp-level:0
\[
\qquad [\text{machine: exhaustive over the}\ 3^3\ \text{charge triples}]:
\]

%dsp-level:0
\[
\boxed{\ \sum_i\aleph_{v_i}\ \equiv\ |Q|\ (\mathrm{mod}\ 2),\qquad\sum_i\aleph_{v_i}\ \ge\ |Q|\ }
\]

%dsp-level:0
\[
\qquad \text{and the wall-sector read is charge-locked:}
\]

%dsp-level:0
\[
\qquad U_{\mathrm{wall}}=(-1)^{\sum\aleph}=(-1)^{|Q|}
\]

%dsp-level:0
\[
\qquad [\text{machine: the}\ 4^3\ \text{end-pair census agrees with}\ W(Q)\ \text{class by class}];
\]

%dsp-level:0
\[
(4)\quad\textbf{the winding and sealing sectors are integer and phase-trivial:}
\]

%dsp-level:0
\[
\qquad \Theta(c)=\sum_{e\in c}\widehat{ds}^{\,2}(e)=\sum_{e\in c}\kappa_e;
\]

%dsp-level:0
\[
\qquad \text{materialized closed chains:}\quad\Theta\in2\mathbb Z,
\]

%dsp-level:0
\[
\qquad \Phi_{\mathrm{spiral}}=\pi\,\Theta=2\pi k,
\]

%dsp-level:0
\[
\qquad \mathrm{Hol}_{\mathrm{spiral}}=+1,
\]

%dsp-level:0
\[
\qquad k=\tfrac\Theta2\in\mathbb Z
\]

%dsp-level:0
\[
\qquad [\mathrm{FT.22};\ \mathrm{FT.24}(1),(2);\ \mathrm{FT.25}(1),(3);\ \text{Census FT.35}(3)]:
\]

%dsp-level:0
\[
\mathrm{Hol}_{\mathrm{int}}=(-1)^{\Theta_{\mathrm{int}}}\in\{\pm1\},\ \ \text{nontrivial only on partial-chain internal sets}
\]

%dsp-level:0
\[
\qquad [\mathrm{FT.25}(2),(3)];
\]

%dsp-level:0
\[
\qquad \text{the baryon-type sealing phase:}\quad\Theta_{\mathrm{int}}\in4\mathbb Z
\]

%dsp-level:0
\[
\qquad [\mathrm{BM.2}(1);\ \text{Census FT.49}(3)]:
\]

%dsp-level:0
\[
e^{i\pi\Theta}=+1\ \ \forall\,\Theta\in2\mathbb Z
\]

%dsp-level:0
\[
\qquad [\text{machine}]:
\]

%dsp-level:0
\[
\qquad \textbf{the whole winding/sealing register contributes trivially to every}\ U(1)\ \textbf{phase read}\ —\ \text{its content is the integer rung}\ k,
\]

%dsp-level:0
\[
\ \neg\text{an angle};
\]

%dsp-level:0
\[
(5)\quad\textbf{the force-block census (the angle-free theorem):}
\]

%dsp-level:0
\[
\qquad \text{A8's complete decomposition, regrouped as}\ \ \widehat F
\]

%dsp-level:0
\[
= \widehat F_{J^\pm}+\widehat F_{G/H}+\widehat F_W+\widehat F_{\mathrm{colour}}+\widehat F_{\mathrm{EM}}+\widehat F_{D^\pm}
\]

%dsp-level:0
\[
\qquad [\mathrm{A8};\ \mathrm{BM.52}(1):\ \text{a regrouping},\ \neg\text{a sub-block restriction}],\ \text{read at a sealed 3-juncture}:
\]

%tbl-level:3
{\scriptsize\setlength{\tabcolsep}{4pt}
\begin{longtable}{>{\raggedright\arraybackslash}p{\dimexpr0.1058\textwidth-2\tabcolsep\relax}|>{\raggedright\arraybackslash}p{\dimexpr0.1058\textwidth-2\tabcolsep\relax}|>{\raggedright\arraybackslash}p{\dimexpr0.7885\textwidth-2\tabcolsep\relax}}
\toprule
block & candidate angle content & the derived verdict \\
\addlinespace[2pt]
\midrule
\endfirsthead
\toprule
block & candidate angle content & the derived verdict \\
\addlinespace[2pt]
\midrule
\endhead
\bottomrule
\endlastfoot
$\widehat F_{J^\pm}$ (the seals) & the sealing phase & $\Theta_{\mathrm{int}}\in\allowbreak 4\mathbb Z$ — an integer $[\mathrm{BM.2}(1);\ \text{Census FT.49}(3)]$; the rotation group of the $\varepsilon$-seal is the finite $C_3$: three one-dimensional characters, no continuous parameter $[\mathrm{BM.3}(1);\ \mathrm{BM.10}(5);\ \mathrm{FT.116}]$ \\
$\widehat F_{G/H}$ (the metric) & frame coefficients & the seven geometric monomials carry charge exponent $0$; no vertex angle $[\mathrm{BM.52}(3);\ \mathrm{BM.58}(1),\allowbreak (2),\allowbreak (4)]$ \\
$\widehat F_W$ (weak/shift) & the shift phase & $\|w\|^2$ class-shared; the Casimir is branch data, neutral at the $\sigma$-vertex $[\mathrm{BM.28}(6);\ \mathrm{BM.40}(3);\ \mathrm{BM.52}(3)]$ \\
$\widehat F_{\mathrm{colour}}$ & slot rotations & annihilated on a sealed singlet: the trace law $\mathfrak T(H)=\allowbreak 3a$; the seed obeys $T_a\mathsf v=\allowbreak 0$ $[\mathrm{BM.1}(2);\ \mathrm{BM.57}(5),\allowbreak (6);\ \mathrm{FT.117}(1);\ \mathrm{EW.32}]$ \\
$\widehat F_{\mathrm{EM}}$ & the gauge phase & no local projector; static monopole order; the coupling $e$ symbolic; chart-side $[\mathrm{Def\text{-}BM.12}(2),\allowbreak (7);\ \mathrm{BM.58}(3)]$ \\
$\widehat F_{D^\pm}$ (dark) & — & $D^+\allowbreak =\allowbreak \Lambda\,\allowbreak \mathbb I$ uniform; $D^-$ locked, retained, $\ne\allowbreak g$accessible, no gauge structure $[\mathrm{A8};\ \mathrm{CT4.1}\text{–}\mathrm{CT4.5};\ \mathrm{D4.2};\ \mathrm{D1.11};\ \mathrm{BM.52}(3);\ \mathrm{FT.73}(1)]$ \\
\end{longtable}
}

%dsp-level:0
\[
\textbf{the seed-transport row:}
\]

%dsp-level:0
\[
\qquad \text{the channel angles}\ \delta\in\bigl\{\tfrac29,\tfrac2{27},\tfrac18\bigr\}\ \text{are channel-line data,}\ \neg\text{attached to the baryon rows}\ [\mathrm{Def\text{-}BM.11};\ \mathrm{BM.59}(1\mathrm a):\ \text{the fourth ground stands}],
\]

%dsp-level:0
\[
\ \text{and they enter the Gram mass form,}\ \neg\text{a loop connection}\ [\mathrm{EW.33};\ \mathrm{Def\text{-}H.37}(1)]:
\]

%dsp-level:0
\[
\boxed{\ \theta_v^{\mathrm{force}}\ \ \textbf{is derived angle-free at a sealed 3-juncture: every block is annihilated by a trace/projector law, integer- or binary-valued, class-shared, or unattached}\ }
\]

%dsp-level:0
\[
\qquad [\text{machine: the row census exhaustive against A8's block list}];
\]

%dsp-level:0
\[
\textbf{the caps add integers only:}
\]

%dsp-level:0
\[
\qquad C_{\mathrm{cap}}=16=7+3+3+3,
\]

%dsp-level:0
\[
\qquad D_{\mathrm{tick}}(3)=2^4=16,
\]

%dsp-level:0
\[
\qquad O(2)=O(3)=0,
\]

%dsp-level:0
\[
\qquad \kappa_0\in\mathbb Z
\]

%dsp-level:0
\[
\qquad [\mathrm{CT0.5};\ \mathrm{CT0.6};\ \mathrm{BM.52}(1);\ \text{Census FT.28}(5);\ \text{Census FT.64}(5);\ \mathrm{FT.24}(1);\ \text{machine}];
\]

%dsp-level:0
\[
(6)\quad\textbf{the two-torsion obstruction (theorem):}
\]

%dsp-level:0
\[
\qquad \text{let}
\]

%dsp-level:0
\[
\qquad \Theta_R(\gamma)\ =\ 2\pi\,n(\gamma)\ +\ \sum_{e\in\gamma}\alpha_e\ +\ \sum_{v\in\gamma}\theta_v^{\mathrm{force}}
\]

%dsp-level:0
\[
\qquad (\mathrm{mod}\ 2\pi)
\]

%dsp-level:0
\[
\text{be any constrained lift with}\ n(\gamma)\ \text{integer}\ [(4),(5):\ \text{the register and the caps fix integers}].
\]

%dsp-level:0
\[
\qquad \text{Then:}
\]

%dsp-level:0
\[
\qquad 2\pi n\equiv0,\quad\textstyle\sum\alpha_e\in\pi\mathbb Z,\quad\sum\theta_v^{\mathrm{force}}\equiv0\ \ \text{on the derived data}:
\]

%dsp-level:0
\[
\boxed{\ \operatorname{im}(\text{any register/force/caps lift})\ \subseteq\ \{0,\ \pi\}\ \subset\ \mathbb R/2\pi\mathbb Z\ }\qquad[\text{machine: the image census}]:
\]

%dsp-level:0
\[
\textbf{against the defect spectrum:}
\]

%dsp-level:0
\[
\qquad \text{the six states of}\ \mathrm{Def\text{-}BM.14}(3)\ \text{carry the angles}\ \ \Theta_R\in\bigl\{\pi,\ 2\pi,\ \tfrac\pi2,\ \tfrac{2\pi}3,\ \tfrac\pi5,\ \tfrac\pi3\bigr\}
\]

%dsp-level:0
\[
\qquad [\delta_v/\pi\in\bigl\{1,\tfrac12,\tfrac23,\tfrac15,\tfrac13\bigr\};\ \text{Census FT.63}(4),(6);\ \text{machine}]:
\]

%dsp-level:0
\[
\text{exactly two of the six lie in the two-torsion image}\ —\ \text{the self-dual floor's pair}\ (\pi;\ 2\pi\equiv0);
\]

%dsp-level:0
\[
\qquad \tfrac\pi2,\ \tfrac{2\pi}3,\ \tfrac\pi5,\ \tfrac\pi3\ \ \text{are pairwise distinct nontrivial cosets modulo}\ \pi
\]

%dsp-level:0
\[
\qquad [\text{machine: exact rational coset arithmetic}]:
\]

%dsp-level:0
\[
\boxed{\ \neg\exists\ \text{a lift of the full real holonomy from the register, the force decomposition and the caps}\ }
\]

%dsp-level:0
\[
\qquad \text{the lift is derived exactly on the self-dual floor and obstructed on both dual pairs};
\]

%dsp-level:0
\[
(7)\quad\textbf{the rung refinement (where the binary sector does reach):}
\]

%dsp-level:0
\[
\qquad m^2_{\mathrm{pole}}(m)=4\,\hat\mu'^2\sin^2\!\bigl(m\,\tfrac{\Theta_R}4\bigr),
\]

%dsp-level:0
\[
\qquad m\in\{k,\,k{-}2,\dots,-k\}
\]

%dsp-level:0
\[
\qquad [\mathrm{FT.60}(5);\ \mathrm{Def\text{-}H.38};\ \text{Census FT.63}(6)]:
\]

%dsp-level:0
\[
\boxed{\ \text{the rung}\ m\ \ \text{is register-reachable}\iff m\,\delta_v\in\pi\mathbb Z\iff D_S(m)=2-2\cos\!\bigl(m\,\tfrac{\delta_v}2\bigr)\in\{0,\,2,\,4\}\ }
\]

%dsp-level:0
\[
\qquad [\text{machine: the exact iff on every census ladder}\ m\le4\ \text{and by the period law}\ m\in\bigl(\tfrac{\pi}{\delta_v}\bigr)\mathbb Z\ [\mathrm{FT.76}(3)]]:
\]

%tbl-level:0
{\setlength{\tabcolsep}{4pt}
\begin{longtable}{>{\raggedright\arraybackslash}p{\dimexpr0.2059\textwidth-2\tabcolsep\relax}|>{\raggedright\arraybackslash}p{\dimexpr0.1076\textwidth-2\tabcolsep\relax}|>{\raggedright\arraybackslash}p{\dimexpr0.3089\textwidth-2\tabcolsep\relax}|>{\raggedright\arraybackslash}p{\dimexpr0.3776\textwidth-2\tabcolsep\relax}}
\toprule
floor & $\delta_v/\pi$ & reachable rungs $m\le\allowbreak 4$ & the ladder values there \\
\addlinespace[2pt]
\midrule
\endfirsthead
\toprule
floor & $\delta_v/\pi$ & reachable rungs $m\le\allowbreak 4$ & the ladder values there \\
\addlinespace[2pt]
\midrule
\endhead
\bottomrule
\endlastfoot
tetrahedron & $1$ & all $m$ & $(2,\allowbreak \,\allowbreak 4,\allowbreak \,\allowbreak 2,\allowbreak \,\allowbreak 0)$ — all in $\{0,2,4\}$ \\
cube & $\tfrac12$ & $m\equiv0\ (2)$ & $D_S(2)=\allowbreak 2,\allowbreak \quad D_S(4)=\allowbreak 4$ \\
octahedron & $\tfrac23$ & $m\equiv0\ (3)$ & $D_S(3)=\allowbreak 4$ \\
dodecahedron & $\tfrac15$ & none at $m\le\allowbreak 4$; $m\equiv0\ (5)$ & $D_S(5)=\allowbreak 2$ \\
icosahedron & $\tfrac13$ & $m\equiv0\ (3)$ & $D_S(3)=\allowbreak 2$ \\
\end{longtable}
}

%dsp-level:0
\[
\textbf{the register captures the defect ladder's rational values}\ \{0,2,4\}\ \textbf{and never its irrational ones}\ —\ \text{the}\ \mathbb Z_2\ \text{parity of the defect angle is a derived read, the real angle is}\ \neg;
\]

%dsp-level:0
\[
(8)\quad\textbf{the calibration-weight obstruction (the winding route):}
\]

%dsp-level:0
\[
\qquad \Theta_R=2\pi k\ \ \textbf{is ruled out as a corpus relation:}
\]

%dsp-level:0
\[
\qquad k=\oint ds^2/(2\mathfrak h^2)\ \ \text{carries}\ \mathfrak h\text{-weight}\ {-}2
\]

%dsp-level:0
\[
\qquad [\mathrm{FT.31}(3);\ \mathrm{D2.9}];
\]

%dsp-level:0
\[
\Theta_R=\iint_S K\,dA\ \ \text{carries}\ \mathfrak h\text{-weight}\ 0
\]

%dsp-level:0
\[
\qquad [\mathrm{Def\text{-}H.21}(6)];
\]

%dsp-level:0
\[
\qquad \text{a proportionality across calibration weights is}\ \neg\,\mathfrak h\text{-free}
\]

%dsp-level:0
\[
\qquad [\mathrm{CT2.11}\ \text{criterion};\ \mathrm{FT.55}(1)\ \text{re-cited}]:
\]

%dsp-level:0
\[
\textbf{the winding register cannot supply the frame angle either};
\]

%dsp-level:0
\[
(9)\quad\textbf{the chart-side locus stands:}
\]

%dsp-level:0
\[
\qquad \delta_v=2\pi-d\,\frac{(L-2)\pi}L\ \ \text{from the integer card data}\ (d,L),
\]

%dsp-level:0
\[
\qquad V\,\delta_v=4\pi
\]

%dsp-level:0
\[
\qquad [\text{Descartes/Gauss--Bonnet};\ \text{Census FT.63}(4);\ \mathrm{FT.47}(1);\ \mathrm{BM.60}(3);\ \text{machine: exact in}\ \pi\text{-units}]:
\]

%dsp-level:0
\[
\text{a loop cleanly enclosing one vertex carries}\ \ \Theta_R=\delta_v
\]

%dsp-level:0
\[
\qquad [\mathrm{FT.47}(1);\ \mathrm{Def\text{-}H.21}(5),(6)],
\]

%dsp-level:0
\[
\qquad D_S=4\sin^2\!\tfrac{\delta_v}4
\]

%dsp-level:0
\[
\qquad [\mathrm{Def\text{-}H.37}(2);\ \mathrm{FT.54}(3),(4)]:
\]

%dsp-level:0
\[
\textbf{the defect spectrum is chart-side face geometry consumed through}\ \mathrm{Def\text{-}H.37}(2)\text{'s interface, whose non-derived clause is priced by}\ \mathrm{FT.54}(1)\text{'s form-uniqueness};
\]

%dsp-level:0
\[
\textbf{the ledger realization of clean-enclosure (vertex-figure) loops is}\ \neg\textbf{constructed}\ —\ \text{the corner convention is undeclared}
\]

%dsp-level:0
\[
\qquad [\text{Census FT.63}(4)\ \mathrm{Rmk};\ \text{Census FT.63}(6)\ \mathrm{Rmk}];
\]

%dsp-level:0
\[
(10)\quad\textbf{the structured obstruction record:}
\]

\begin{verbatim}
OBSTRUCTION: holonomy_lift_not_derived
MISSING: vertex_figure_loop_realization | Symk_assignment [FT.55(4)] | chiral_assignment [FT.125(4)] | coarse_graining_lift 
DERIVED_NO_GO: two_torsion_image (6) | calibration_weight (8)
PARTIAL_LIFT: Z2_wilson_sector (2) | parity_lock (3) | integer_winding_sealing_sector (4) | self_dual_floor (6),(7) | rung_law (7)
\end{verbatim}

%dsp-level:0
\[
(11)\quad\textbf{status:}
\]

%dsp-level:0
\[
\qquad \textbf{derived:}\ (1)\text{--}(9);
\]

%dsp-level:0
\[
\qquad \textbf{record issued:}\ (10);
\]

%dsp-level:0
\[
\qquad \textbf{outcome:}\ \text{the obstruction theorem, with the partial lift of}\ (2)\text{--}(4),
\]

%dsp-level:0
\[
(7)\ \text{and the self-dual-floor lift of}\ (6);
\]

%dsp-level:0
\[
\neg\text{a declared holonomy: no continuous angle is inserted anywhere in this tier}\ —\ \Theta_R\ \text{occurs only inside re-citations of the chart-side reads and inside the refuted lift equation}\ [\text{machine: the occurrence census}];
\]

%dsp-level:0
\[
\neg\text{a P4 addition},\quad\neg\text{a new declaration},\quad\neg\text{a measured value}\ [\mathrm{P5}];
\]

%dsp-level:0
\[
\qquad \text{the conditionalities}\ (\mathrm{Def\text{-}H.37}(2),\ \mathrm{Def\text{-}H.38},\ \mathrm{M1}/\mathrm{M2}\ [\mathrm{Def\text{-}H.20}],\ \mathrm{CT2.6})\ \text{are inherited at the point of use}\ [\text{A §0.2a}].
\]

\textbf{Req.} A4, A7, A8 (Part 1); D1.11, D2.9, D3.3, D3.5, D4.2, CT0.5, CT0.6, CT0.9 (Part 1); CT2.11, CT4.1, CT4.2, CT4.3, CT4.4, CT4.5, EM.9b, EW.32, EW.33 (Part 2); Def-EM.20, EM.20, EM.21 (Part 2); P5 (Parts 1–2); Def-H.20, Def-H.21(5), Def-H.21(6), Def-H.37(1), Def-H.37(2), Def-H.38, FT.22, FT.24(1), FT.24(2), FT.25(1), FT.25(2), FT.25(3), FT.31(3), FT.47(1), FT.54(1), FT.54(3), FT.54(4), FT.55(1), FT.55(4), FT.60(5), FT.73(1), FT.76(3), FT.116, FT.117(1), FT.125(4) (Part 3); Census FT.28(5), FT.35(3), FT.49(3), FT.63(4), FT.63(6), FT.64(5) (row-wise); Def-BM.8(1), Def-BM.8(3), Def-BM.11, Def-BM.12(2), Def-BM.12(7), Def-BM.13(2), Def-BM.14(3), Def-BM.14(4), BM.1(2), BM.2(1), BM.3(1), BM.10(5), BM.28(6), BM.40(3), BM.52(1), BM.52(3), BM.57(5), BM.57(6), BM.58(1), BM.58(2), BM.58(3), BM.58(4), BM.59(1), BM.60(3) (this part); exhaustive sign-assignment censuses, exact rational coset arithmetic modulo $\pi$ and $2\pi$, GF(2) parity identities, Burnside orbit counts on finite ladders [math; machine: \texttt{checks/\allowbreak{}verify\_\allowbreak{}H\_\allowbreak{}organic\_\allowbreak{}bridge\_\allowbreak{}post\_\allowbreak{}bm14.py} §D].

\textbf{Pf.}

%dsp-level:0
\[
(1):\ \mathrm{A4}/\mathrm{A7}\ \text{re-cited: the branch phase and the edge map};
\]

%dsp-level:0
\[
\qquad \alpha_e\in\{0,\pi\}\ [\mathrm{CT0.9}];
\]

%dsp-level:0
\[
\qquad U(c)=\prod\sigma_e^{\operatorname{or}(e)}=e^{i\sum\operatorname{or}(e)\alpha_e}\ [\text{machine: the exponential of a sum over}\ \{0,\pi\}]
\]

%dsp-level:0
\[
(2):\ \sigma^{\pm1}=\sigma\ \text{on}\ \{\pm1\}\ [\text{machine}];
\]

%dsp-level:0
\[
\qquad \text{exhaustive over}\ 2^n\ \text{assignments for}\ n
\]

%dsp-level:0
\[
= 2,
\]

%dsp-level:0
\[
\dots,
\]

%dsp-level:0
\[
7\ \text{and both orientation patterns: the product equals}\ (-1)^{\#\{\sigma_e=-1\}}\ \text{at every one}\ [\text{machine}];
\]

%dsp-level:0
\[
\qquad \text{concatenation: products multiply; reversal: each}\ \sigma\ \text{is its own inverse}\ [\text{machine}]
\]

%dsp-level:0
\[
(3):\ \mathrm{Def\text{-}BM.13}(2)\text{'s exhaustive census re-run}\ [\text{machine}];
\]

%dsp-level:0
\[
\qquad |q|\equiv q\ (2)\ \text{for}\ q\in\{-1,0,1\}\ [\text{machine: three cases}];
\]

%dsp-level:0
\[
\qquad \textstyle\sum|q_i|\ge|\sum q_i|\ [\text{machine: exhaustive over}\ 27\ \text{triples}];
\]

%dsp-level:0
\[
\qquad \text{the end-pair census:}\ q
\]

%dsp-level:0
\[
= 0\iff\sigma_s
\]

%dsp-level:0
\[
= \sigma_t\ (2\ \text{pairs},\ \text{even flips}),
\]

%dsp-level:0
\[
\ |q|
\]

%dsp-level:0
\[
= 1\iff\sigma_s\ne\sigma_t\ (1+1\ \text{pairs},\ \text{odd flips})\ [\text{machine}]
\]

%dsp-level:0
\[
(4):\ \mathrm{FT.22}/\mathrm{FT.24}(1),(2)/\mathrm{FT.25}(1)\text{--}(3)\ \text{re-cited};
\]

%dsp-level:0
\[
\qquad \Theta\in2\mathbb Z\implies e^{i\pi\Theta}=1\ [\text{machine: even integers}];
\]

%dsp-level:0
\[
\qquad 4\mathbb Z\subseteq2\mathbb Z\ [\text{machine}]
\]

%dsp-level:0
\[
(5):\ \text{each row re-cited from its own tag and classified}\ [\text{machine: the census structure against A8's six groups}];
\]

%dsp-level:0
\[
\qquad 7+3+3+3=16,\ 2^4=16,\ 32-16=16\ [\mathrm{CT0.5};\ \mathrm{CT0.6};\ \text{machine}]
\]

%dsp-level:0
\[
(6):\ \text{the image:}\ \{2\pi n+\pi j:\ n,j\in\mathbb Z\}\equiv\{0,\pi\}\ (\mathrm{mod}\ 2\pi)\ [\text{machine}];
\]

%dsp-level:0
\[
\qquad \text{the state angles modulo}\ \pi:\ \pi
\]

%dsp-level:0
\[
\equiv 0,
\]

%dsp-level:0
\[
\ 2\pi
\]

%dsp-level:0
\[
\equiv 0,
\]

%dsp-level:0
\[
\ \tfrac\pi2\not
\]

%dsp-level:0
\[
\equiv 0,
\]

%dsp-level:0
\[
\ \tfrac{2\pi}3\not
\]

%dsp-level:0
\[
\equiv 0,
\]

%dsp-level:0
\[
\ \tfrac\pi5\not
\]

%dsp-level:0
\[
\equiv 0,
\]

%dsp-level:0
\[
\ \tfrac\pi3\not
\]

%dsp-level:0
\[
\equiv 0,
\]

%dsp-level:0
\[
\ \text{pairwise distinct in}\ \mathbb Q/\mathbb Z\ [\text{machine: exact Fractions}];
\]

%dsp-level:0
\[
\qquad \text{the count: two in the image, four outside}
\]

%dsp-level:0
\[
(7):\ \cos\theta\in\{0,\pm1\}\iff\theta\in\tfrac\pi2\mathbb Z\ [\text{math}]
\]

%dsp-level:0
\[
\implies D_S(m)\in\{0,2,4\}
\]

%dsp-level:0
\[
\iff m\,
\]

%dsp-level:0
\[
\tfrac{\delta_v}2\in\tfrac\pi2\mathbb Z
\]

%dsp-level:0
\[
\iff m\,
\]

%dsp-level:0
\[
\delta_v\in\pi\mathbb Z\ [\text{machine}];
\]

%dsp-level:0
\[
\qquad \text{the five ladders evaluated exactly: the tetrahedral}\ (2,4,2,0);
\]

%dsp-level:0
\[
\ \text{the cube's}\ D_S(2)
\]

%dsp-level:0
\[
= 2-2\cos\tfrac\pi2
\]

%dsp-level:0
\[
= 2,
\]

%dsp-level:0
\[
\ D_S(4)
\]

%dsp-level:0
\[
= 2-2\cos\pi
\]

%dsp-level:0
\[
= 4;
\]

%dsp-level:0
\[
\ \text{the octahedron's}\ D_S(3)
\]

%dsp-level:0
\[
= 2-2\cos\pi
\]

%dsp-level:0
\[
= 4;
\]

%dsp-level:0
\[
\ \text{the icosahedron's}\ D_S(3)
\]

%dsp-level:0
\[
= 2-2\cos\tfrac\pi2
\]

%dsp-level:0
\[
= 2;
\]

%dsp-level:0
\[
\ \text{the dodecahedron's}\ D_S(5)
\]

%dsp-level:0
\[
= 2\ [\text{machine}];
\]

%dsp-level:0
\[
\qquad \text{the irrational members}\ 2-\sqrt2,
\]

%dsp-level:0
\[
\ 2+\sqrt2,
\]

%dsp-level:0
\[
\ 1,
\]

%dsp-level:0
\[
\ 3,
\]

%dsp-level:0
\[
\ \tfrac{3-\sqrt5}2,
\]

%dsp-level:0
\[
\ \tfrac{5-\sqrt5}2,
\]

%dsp-level:0
\[
\ 2-\sqrt3,
\]

%dsp-level:0
\[
\ 2-2\cos\tfrac\pi{10}\ \text{are separated from}\ \{0,2,4\}\ \text{by their minimal polynomials evaluated at}\ 0,
\]

%dsp-level:0
\[
2,
\]

%dsp-level:0
\[
4\ [\text{machine: every value nonzero}]
\]

%dsp-level:0
\[
(8):\ \mathrm{FT.55}(1)\ \text{re-cited verbatim};
\]

%dsp-level:0
\[
\qquad (9):\ \delta_v\ \text{evaluated at}\ (d,L)
\]

%dsp-level:0
\[
= (3,3),
\]

%dsp-level:0
\[
(3,4),
\]

%dsp-level:0
\[
(4,3),
\]

%dsp-level:0
\[
(3,5),
\]

%dsp-level:0
\[
(5,3):\ \ \delta_v/\pi
\]

%dsp-level:0
\[
= (1,\tfrac12,\tfrac23,\tfrac15,\tfrac13)\ [\text{machine}];
\]

%dsp-level:0
\[
\qquad V\delta_v=4\pi\ \text{at all five}\ [\text{machine}]
\]

%dsp-level:0
\[
(10)\text{--}(11):\ \text{the record assembled from}\ (6),(8),(9);
\]

%dsp-level:0
\[
\qquad \text{the occurrence census of}\ \Theta_R\ \text{in this tier}\ [\text{machine}];
\]

%dsp-level:0
\[
\qquad \text{the standing-rule scans}\ [\text{machine}]. \ \blacksquare
\]

\textbf{Phys.} Ask the ledger for a real holonomy and it answers in bits. The edge-phase register is binary, so every closed-loop product is plus or minus one, oriented either way; the wall content of any transport is parity-locked to its charge core, and at the composite the lock reads exactly — total wall parity equals the absolute charge modulo two, and the wall count bounds the charge from below. The winding and sealing registers are integers, and every integer winding contributes trivially to a phase, so their content is the rung, not an angle. Then the force decomposition is censused block by block at the sealed junction, and not one block carries a continuous vertex angle: the seals are finite and integer, the metric is charge-blind, the weak block is class-shared branch data, colour is annihilated by the trace law on anything sealed, electromagnetism is chart-side at a symbolic coupling, the dark blocks are uniform or locked, and the channel angles are line data attached to no baryon row. Any lift built from that material lands in the two-element subgroup — zero or pi — while the defect spectrum spreads six distinct values whose angles include a fifth of pi and a third of pi: four of the six states are provably out of reach, and the reachable sublattice is exactly the self-dual floor plus the ladder rungs whose defect value is zero, two or four. The winding route is shut separately, by calibration weight: the rung counts hops, the frame angle counts curvature, and no pure number bridges the two. So the tier records what it has — the binary Wilson sector, the parity lock, the integer sectors, the rung law — and issues the obstruction in structured form, with the missing enumeration named: the ledger realization of a loop that cleanly encloses one vertex. No angle is declared anywhere.

\pseudohead{theorem}{\texorpdfstring{BM.68}{BM.68}}{\texorpdfstring{The constrained joint entropy as a derived fiber cardinality}{The constrained joint entropy as a derived fiber cardinality}}\label{tag:BM-68}
\textbf{Thm.}

%dsp-level:0
\[
(1)\quad\textbf{the constrained fiber (an abbreviation of derived objects):}
\]

%dsp-level:0
\[
\qquad \Omega(E,\Theta,Q)\ :
\]

%dsp-level:0
\[
= \ \#\Bigl\{(R,\sigma,g):\ \ E_{\mathrm{int}}(R)=E,\ \ Q(\sigma)=Q,\ \ \Theta(R,g)=\Theta,\ \ \text{the sealing and admissibility constraints hold}\Bigr\}:
\]

%dsp-level:0
\[
\text{the census of ledger states passing the filter chain}\ \mathrm{F0}\text{--}\mathrm{F7}
\]

%dsp-level:0
\[
\qquad [\text{Census FT.64}(1):\ \text{configuration validity};\ \text{registration};\ \text{integrality};\ \text{loop quantization};\ \text{stationarity};\ \text{regime};\ \text{relabeling quotient};\ \text{register product}];
\]

%dsp-level:0
\[
S_{\mathrm{joint}}(E,\Theta,Q)\ :=\ \ln\Omega(E,\Theta,Q)
\]

%dsp-level:0
\[
\qquad [\mathrm{C9}:\ \text{entropy in nats};\ \text{machine: cardinalities of finite sets}];
\]

%dsp-level:0
\[
(2)\quad\textbf{the factorization (theorem on the tested domains):}
\]

%dsp-level:0
\[
\qquad \text{the registers are independent}\ —\ \text{the signs factor from the}\ \kappa\text{-ledger}\ [\mathrm{CT0.10};\ \text{Census FT.51}(5)],
\]

%dsp-level:0
\[
\ \ E_{\mathrm{int}}\ \text{is constant on every end-sign fibre}\ [\mathrm{Def\text{-}BM.13}(6);\ \mathrm{BM.28}(3);\ \text{machine: exhaustive over the four end pairs}],
\]

%dsp-level:0
\[
\ \text{the sealing grammar is its own register}\ [\mathrm{FT.41}(b);\ \text{Census FT.65}(2)]:
\]

%dsp-level:0
\[
\boxed{\ \Omega(E,\Theta,Q)\ =\ \Omega_{\mathrm{count}}(E)\ \cdot\ \Omega_{\mathrm{charge}}(Q)\ \cdot\ \Omega_{\mathrm{geom}}(\Theta;E,Q)\ }
\]

%dsp-level:0
\[
\qquad \text{``the three registers multiply, none converts into another''}
\]

%dsp-level:0
\[
\qquad [\text{Census FT.51}(5),\ \text{FT.65}(2);\ \text{Census FT.64}(1)\ \mathrm{F7}]:
\]

%dsp-level:0
\[
S_{\mathrm{joint}}\ =\ \ln\Omega_{\mathrm{count}}(E)\ +\ \ln\Omega_{\mathrm{charge}}(Q)\ +\ \ln\Omega_{\mathrm{geom}}(\Theta;E,Q)
\]

%dsp-level:0
\[
\qquad [\text{machine: the cardinality identity on every tested finite domain}];
\]

%dsp-level:0
\[
(3)\quad\textbf{the first factor (derived):}
\]

%dsp-level:0
\[
\qquad \boxed{\ \Omega_{\mathrm{count}}(E)\ =\ 2^E\ }
\]

%dsp-level:0
\[
\qquad [\text{one binary registration per directed real edge}:\ \mathrm{CT0.9}:\ 1\ \text{real edge}=1\ \text{bit}=\ln2\ \text{nats};\ \mathrm{C9};\ \mathrm{CT1.29}]:
\]

%dsp-level:0
\[
I_{\mathrm{hid}}=E_{\mathrm{int}}\ln2=\ln2^{E_{\mathrm{int}}}
\]

%dsp-level:0
\[
\qquad [\mathrm{Def\text{-}H.9}(2);\ \mathrm{FT.17};\ \text{Census FT.35}(2);\ \text{Census FT.44}(4);\ \text{machine}]:
\]

%dsp-level:0
\[
\qquad S_{\mathrm{count}}=E\ln2\ \ \text{is exactly}\ \mathrm{Def\text{-}BM.13}(5)\text{'s first summand};
\]

%dsp-level:0
\[
(4)\quad\textbf{the second factor (derived):}
\]

%dsp-level:0
\[
\qquad \boxed{\ \Omega_{\mathrm{charge}}(Q)\ =\ W(Q)\ =\ \bigl[z^{Q}\bigr]\bigl(2+z+z^{-1}\bigr)^{k}\ },
\]

%dsp-level:0
\[
\qquad k=3:\quad W=(1,\,6,\,15,\,\mathbf{20},\,15,\,6,\,1),
\]

%dsp-level:0
\[
\qquad \sum_QW(Q)=64
\]

%dsp-level:0
\[
\qquad [\mathrm{Def\text{-}BM.8}(6);\ \mathrm{Def\text{-}BM.13}(3);\ \mathrm{BM.28}(2);\ \text{machine: two agreeing routes}]:
\]

%dsp-level:0
\[
S_{\mathrm{charge}}=\ln W(Q)=\Sigma_{\mathrm{cl}}(Q)\qquad[\mathrm{Def\text{-}BM.13}(4)]:\qquad\text{the second summand};
\]

%dsp-level:0
\[
(5)\quad\textbf{the third factor at stationarity (derived):}
\]

%dsp-level:0
\[
\qquad \text{on the stationary locus}\ \kappa_e\equiv\kappa_0\in\mathbb Z\ \ \text{exactly along each connected materialized chain}
\]

%dsp-level:0
\[
\qquad [\mathrm{FT.53}(1),(2);\ \text{Census FT.64}(1)\ \mathrm{F4}],
\]

%dsp-level:0
\[
\text{and the semiclassical survivor list carries one orbit per rung}
\]

%dsp-level:0
\[
\qquad [\text{Census FT.64}(5)(\mathrm c):\ \text{the R2 ladder};\ \mathrm{FT.53}(3)]:
\]

%dsp-level:0
\[
\qquad \text{a uniform profile is fixed by the swap/relabeling quotient}:
\]

%dsp-level:0
\[
\boxed{\ \Omega_{\mathrm{geom}}^{\mathrm{stat}}(\Theta)\ =\ 1\ \ \text{at every stationary rung},\qquad S_{\mathrm{geom}}^{\mathrm{stat}}\ =\ 0\ }
\]

%dsp-level:0
\[
\qquad [\text{machine: the Burnside count on a single-element set; the swap fixes}\ (\kappa_0,\kappa_0,\kappa_0)]:
\]

%dsp-level:0
\[
\textbf{on the baryon card:}
\]

%dsp-level:0
\[
\qquad \text{the strict-spacelike minimum sits on the locus:}\quad\kappa=(1,1,1),\quad K=3,\quad\Theta_{\mathrm{int}}=6
\]

%dsp-level:0
\[
\qquad [\text{Census FT.65}(4)\ \Sigma_{12}\ \text{row};\ \mathrm{FT.24}(1):\ \kappa\in\mathbb Z\ \text{at materialized digons};\ \text{machine: the unique profile, the unique orbit}]:
\]

%dsp-level:0
\[
\boxed{\ S_{\mathrm{joint}}\bigl(6,\ \Theta_{\mathrm{stat}},\ Q\bigr)\ =\ 6\ln2\ +\ \ln W(Q)\ +\ 0\ }
\]

%dsp-level:0
\[
\textbf{the refined rest input is the complete joint entropy of the stationary baryon fibre, on its ledger part:}
\]

%dsp-level:0
\[
\qquad \text{the EM summand is the chart-side transit-channel energy,}\ \neg\text{a ledger count}
\]

%dsp-level:0
\[
\qquad [\mathrm{Def\text{-}BM.12}(6);\ \mathrm{Def\text{-}BM.13}(5);\ \mathrm{BM.66}(6);\ \text{machine: the alphabet disjointness}]:
\]

%dsp-level:0
\[
\textbf{the additive law is a derived stationarity slice,}\ \neg\textbf{a postulate};
\]

%dsp-level:0
\[
\qquad \text{the off-stationary fibers cancel pairwise in}\ Z\ \text{—}\ \neg\text{a semiclassical spectral contribution}
\]

%dsp-level:0
\[
\qquad [\mathrm{CT1.20};\ \mathrm{FT.53}(4);\ \text{Census FT.64}(1)\ \mathrm{F8}\ \text{note}]:
\]

%dsp-level:0
\[
\qquad \neg\text{a mass read requires an off-stationary count};
\]

%dsp-level:0
\[
(6)\quad\textbf{the general fiber (the named missing enumeration):}
\]

%dsp-level:0
\[
\qquad \textbf{the Platonic floors}\ —\ \text{derived in closed form: the on-shell lattice}\ \ker\varphi_S,
\]

%dsp-level:0
\[
\ \ \text{index}\ 2^{F-1}\ [\mathrm{FT.43}(3);\ \text{Census FT.64}(2)];
\]

%dsp-level:0
\[
\ \text{the level counts}\ A_S(K)\ \text{with the exact character-sum closed form}\ [\text{Census FT.64}(4)];
\]

%dsp-level:0
\[
\ \text{the Burnside orbits}\ M_S(K),
\]

%dsp-level:0
\[
\ M^+_S(K)\ [\text{Census FT.65}(1)];
\]

%dsp-level:0
\[
\ \text{the sealing grammar}\ g(V)
\]

%dsp-level:0
\[
= (3,\,385,\,25,\,292938721075,\,164395)\ [\mathrm{FT.41}(b);\ \text{Census FT.63}(7)];
\]

%dsp-level:0
\[
\ \text{the register product}\ 2^{2E}\!\cdot\!M_S(K)\cdot g(V)\cdot(\text{attachment data})\ [\text{Census FT.65}(2)]:
\]

%dsp-level:0
\[
\textbf{the polygon floors}\ —\ \text{the dihedral necklaces}\ M(n,k)\ \text{in closed Burnside form}
\]

%dsp-level:0
\[
\qquad [\text{Census FT.51}(1)\text{--}(3)]:
\]

%dsp-level:0
\[
\textbf{the baryon junction}\ \Sigma_{12}:
\]

%dsp-level:0
\[
\qquad \text{the winding lattice is}\ \mathbb Z^3\ \text{with the swap quotient}\ \mathbb Z_2
\]

%dsp-level:0
\[
\qquad [\text{Census FT.65}(4)\ \Sigma_{12}\ \text{row}];
\]

%dsp-level:0
\[
\text{the fiber over the full lattice at fixed}\ K\ \ \textbf{is infinite}\ —\ \text{the witness}\ (K{+}n,\,{-}n)
\]

%dsp-level:0
\[
\qquad [\text{Census FT.64}(1)\ \mathrm{F5};\ \text{machine}]:
\]

%dsp-level:0
\[
\qquad \text{the enumeration requires the regime datum}\ (\text{null-permitting}\ /\ \text{strict-spacelike}\ /\ \text{uniform}\ /\ \text{signature-mixed})\ \text{that the card does not fix}:
\]

%dsp-level:0
\[
\boxed{\ \text{the exact missing enumeration}\ =\ \text{the regime-scoped orbit count}\ \ \#\bigl\{\kappa\in\mathbb Z^3:\ \kappa_1{+}\kappa_2{+}\kappa_{12}=K\bigr\}\big/\mathbb Z_2\ \ \text{at the baryon card's materialized rung}\ }
\]

%dsp-level:0
\[
\textbf{the scoped rows this tier derives:}
\]

%dsp-level:0
\[
\qquad \text{the strict regime}\ \kappa\ge1:\quad\#\text{profiles}=\binom{K-1}2,
\]

%dsp-level:0
\[
\qquad \text{orbits}\ (1,\,2,\,4,\,6)\ \ \text{at}\ K=(3,\,4,\,5,\,6)
\]

%dsp-level:0
\[
\qquad [\text{machine: stars-and-bars and Burnside agree with exhaustive enumeration}];
\]

%dsp-level:0
\[
\text{the null-permitting regime}\ \kappa\ge0:\quad\#\text{profiles}=\binom{K+2}2,
\]

%dsp-level:0
\[
\qquad \text{orbits}\ (2,\,6)\ \ \text{at}\ K=(1,\,3)
\]

%dsp-level:0
\[
\qquad [\text{machine: exhaustive}];
\]

%dsp-level:0
\[
\qquad \text{the uniform ladder: one orbit per rung}\ [(5);\ \text{machine}];
\]

%dsp-level:0
\[
(7)\quad\textbf{the two currencies and the bridge-candidate census:}
\]

%dsp-level:0
\[
\qquad \textbf{the bit/charge currency:}\ \ E\ln2+\ln W(Q)\ —\ \text{cardinalities, integers};
\]

%dsp-level:0
\[
\qquad \textbf{the geometric/holonomy currency:}\ \ \Theta_R\ \text{enters as a value,}\ \neg\text{a count}:
\]

%dsp-level:0
\[
\quad m^2_{\mathrm{pole}}=\hat\mu'^2\cdot4\sin^2\!\tfrac{\Theta_R}4
\]

%dsp-level:0
\[
\qquad [\mathrm{Def\text{-}H.37}(2);\ \mathrm{FT.54}(4)],
\]

%dsp-level:0
\[
\qquad D_S\ \text{exact algebraic in four fields}
\]

%dsp-level:0
\[
\qquad [\mathrm{Def\text{-}BM.14}(3),(4)]:
\]

%dsp-level:0
\[
\textbf{every candidate bridge species, tested:}
\]

%dsp-level:0
\[
(\mathrm a)\quad\textbf{a derivative of}\ S_{\mathrm{joint}}:
\]

%dsp-level:0
\[
\qquad \Theta\ \text{is integer-valued};
\]

%dsp-level:0
\[
\ \ \text{the finite difference along the stationary ladder is}\ \ \ln\!\bigl(\Omega^{\mathrm{stat}}/\Omega^{\mathrm{stat}}\bigr)
\]

%dsp-level:0
\[
= \ln1
\]

%dsp-level:0
\[
= 0\ \ \text{while the defect value varies across six distinct values}
\]

%dsp-level:0
\[
\qquad [(5);\ \mathrm{Def\text{-}BM.14}(4);\ \text{machine}]:
\]

%dsp-level:0
\[
\qquad \textbf{the count is flat where the angle is not}\ —\ \text{refuted as a bridge};
\]

%dsp-level:0
\[
(\mathrm b)\quad\textbf{a Legendre transform under constraints:}
\]

%dsp-level:0
\[
\qquad \neg\exists\ \text{a joint potential of which both currencies are reads}\ —\ \text{the three deficit functionals are pairwise distinct on digon witnesses}\ (2\ln2\ \text{vs}\ 0;\ \ 0\ \text{vs}\ 3;\ \ 2\ln2\ne3\ \text{via}\ \tfrac12<\ln2<1)
\]

%dsp-level:0
\[
\qquad [\mathrm{BM.35}(6);\ \text{machine: the exact enclosure}]:
\]

%dsp-level:0
\[
\qquad \neg\text{a Legendre pair is derivable};
\]

%dsp-level:0
\[
(\mathrm c)\quad\textbf{a constrained extremum:}
\]

%dsp-level:0
\[
\qquad \arg\max_Q S_{\mathrm{joint}}(6,\cdot,Q)\ =\ \text{the neutral class}
\]

%dsp-level:0
\[
\qquad [W(0)=20>15=W(\pm1);\ \text{machine}]:
\]

%dsp-level:0
\[
\qquad \textbf{derived, but a class selection,}\ \neg\textbf{a currency conversion};
\]

%dsp-level:0
\[
(\mathrm d)\quad\textbf{a trace normalization:}
\]

%dsp-level:0
\[
\qquad I_F=4\,\widehat Q^2\ \ \text{is an exact integer on the information side}
\]

%dsp-level:0
\[
\qquad [\mathrm{BM.66}(3)];
\]

%dsp-level:0
\[
\qquad \text{the defect side's trace normalization is the seed doublet's dimension:}\ \ (U-1)^\dagger(U-1)
\]

%dsp-level:0
\[
= 4\sin^2\!\tfrac\phi2\cdot I_2
\]

%dsp-level:0
\[
\qquad [\mathrm{FT.54}(1);\ \mathrm{EW.32}]:
\]

%dsp-level:0
\[
\text{both exact, incommensurate species}\ —\ \text{a count against an angle}
\]

%dsp-level:0
\[
\qquad [\text{machine: the token census}]:
\]

%dsp-level:0
\[
\qquad \neg\text{a ratio between them is derivable};
\]

%dsp-level:0
\[
(\mathrm e)\quad\textbf{a Casimir ratio:}
\]

%dsp-level:0
\[
\qquad \text{every Casimir/trace/capacity ratio in the inventory is an exact rational}
\]

%dsp-level:0
\[
\qquad [\mathrm{BM.61}(2),(3):\ \text{integers and ratios of integers};\ \bigl(\tfrac13,\,8,\,\tfrac1{25}\bigr);\ \mathrm{BM.36}(2)]:
\]

%dsp-level:0
\[
\text{the five floors' pole factors generate a degree-}16\ \text{extension of}\ \mathbb Q
\]

%dsp-level:0
\[
\qquad [\mathrm{FT.73}(2);\ \text{machine re-cited}]:
\]

%dsp-level:0
\[
\qquad \textbf{no rational-coefficient formula on integer-valued reads reaches the pole side};
\]

$(8)\quad\textbf{the gap row:}\qquad\text{the bridge's well-posed object}\ [\mathrm{FT.54}(4):\ \text{Gap 7's well-posed object}]\ \ \text{and its single undeclared factor:}$

%dsp-level:0
\[
\boxed{\ \frac{m^2_{\mathrm{pole}}}{m_{\mathrm{hid}}^2}\ =\ \rho_{\mathrm B}\cdot\frac{D_S}{4E^2}\ \in\ \mathfrak q^0,\qquad\rho_{\mathrm B}\ :=\ \frac{\hat\mu'^2\,\mathfrak h^2}{(\ln2)^2}\ }
\]

%dsp-level:0
\[
\qquad [\mathrm{FT.54}(4);\ \text{Census FT.63}(6);\ \text{machine: the dimension row}\ [\hat\mu'^2\mathfrak h^2]=\mathfrak q\cdot\mathfrak q^{-1}=\mathfrak q^0]:
\]

%dsp-level:0
\[
\rho_{\mathrm B}\ \text{is a Cal-F packaging}\ [\mathrm{Def\text{-}H.37}(2);\ \mathrm{FT.54}(4);\ \text{Cal-F logged}]\ —\ \textbf{restated here as the residue's name,}\ \neg\textbf{added};
\]

%dsp-level:0
\[
\qquad \text{the per-state exact algebraic factors:}\quad\frac{D_S}{4E^2}
\]

%dsp-level:0
\[
= \Bigl(\frac1{72},\ \ \frac{2-\sqrt2}{576},\ \ \frac1{576},\ \ \frac{D_S}{3600},\ \ \frac{2-\sqrt3}{3600}\Bigr)
\]

%dsp-level:0
\[
\qquad [\text{Census FT.63}(6);\ \text{machine}];
\]

%dsp-level:0
\[
(9)\quad\textbf{status:}
\]

%dsp-level:0
\[
\qquad \textbf{derived:}\ (1)\text{--}(5),\ (6)\text{'s scoped rows},\ (7),\ (8);
\]

%dsp-level:0
\[
\qquad \textbf{named,}\ \neg\textbf{filled: the general}\ \Sigma_{12}\ \text{fiber}\ (6);
\]

%dsp-level:0
\[
S_{\mathrm{geom}}\ \ \neg\textbf{declared}\ —\ \text{it is derived}\ =0\ \text{at stationarity and named-missing off it};
\]

%dsp-level:0
\[
\qquad \neg\text{a P4 addition},\quad\neg\text{a new dimensionful scale}\ [\mathrm{P1}],\quad\neg\text{a measured value}\ [\mathrm{P5}];
\]

%dsp-level:0
\[
\qquad \text{§Cal-BM stays empty}.
\]

%dsp-level:0
\[
\textbf{the structured record:}
\]

\begin{verbatim}
OBSTRUCTION: geometric_fiber_count_not_derived (the general Sigma_12 fiber)
MISSING: Omega_geom enumeration — the regime-scoped orbit count at the materialized rung
DERIVED: stationary_fiber = 1 (5); Platonic closed forms (6); polygon necklaces (6); strict/null-permitting scoped counts (6)
\end{verbatim}

\textbf{Req.} C9, CT0.9, CT0.10, CT1.20 (Part 1); CT1.29, CT2.1, EW.32 (Part 2); P1, P5 (Parts 1–2); Def-H.9(2), Def-H.37(2), FT.17, FT.41(b), FT.43(3), FT.53(1), FT.53(2), FT.53(3), FT.53(4), FT.54(1), FT.54(4), FT.73(2) (Part 3); Census FT.35(2), FT.44(4), FT.51(1), FT.51(3), FT.51(5), FT.63(6), FT.63(7), FT.64(1), FT.64(2), FT.64(4), FT.64(5), FT.65(1), FT.65(2), FT.65(4) (row-wise); Def-BM.8(6), Def-BM.12(6), Def-BM.13(3), Def-BM.13(4), Def-BM.13(5), Def-BM.13(6), Def-BM.14(3), Def-BM.14(4), BM.28(3), BM.35(6), BM.36(2), BM.61(2), BM.61(3), BM.66(3), BM.66(6) (this part); Cal-F; stars-and-bars and Burnside enumeration, exact rational arithmetic, the prime-log ring with integer certificates, exact radical separations, rational $\ln2$ enclosures [math; machine: \texttt{checks/\allowbreak{}verify\_\allowbreak{}H\_\allowbreak{}organic\_\allowbreak{}bridge\_\allowbreak{}post\_\allowbreak{}bm14.py} §E].

\textbf{Pf.}

%dsp-level:0
\[
(1)\text{--}(2):\ \text{the filter chain re-cited row by row}\ [\text{Census FT.64}(1)];
\]

%dsp-level:0
\[
\qquad \text{the independence rows re-cited and re-run: signs}\perp\kappa\ [\mathrm{CT0.10};\ \text{machine: the product cardinality on the tested domains}],
\]

%dsp-level:0
\[
\ \ E_{\mathrm{int}}\ \text{constant on end-sign fibres}\ [\text{machine: exhaustive over the four end pairs}],
\]

%dsp-level:0
\[
\ \text{sealing separate}\ [\mathrm{FT.41}(b)];
\]

%dsp-level:0
\[
\qquad \Omega=\Omega_{\mathrm{count}}\Omega_{\mathrm{charge}}\Omega_{\mathrm{geom}}\ \text{as a cardinality of a product of finite sets}\ [\text{machine}]
\]

%dsp-level:0
\[
(3):\ 2^6=64\ [\text{machine}];
\]

%dsp-level:0
\[
\qquad I_{\mathrm{hid}}=6\ln2=\ln64\ [\text{machine: the prime-log ring}];
\]

%dsp-level:0
\[
\qquad (4):\ W(Q)\ \text{re-enumerated by both routes}\ [\text{machine}]
\]

%dsp-level:0
\[
(5):\ \mathrm{FT.53}(2)\ \text{re-cited: on the stationary locus}\ \kappa_e=\kappa_0\in\mathbb Z\ \text{exactly};
\]

%dsp-level:0
\[
\qquad \text{the swap}\ (\kappa_1,\kappa_2,\kappa_{12})\mapsto(\kappa_2,\kappa_1,\kappa_{12})\ \text{fixes}\ (\kappa_0,\kappa_0,\kappa_0)\ [\text{machine}];
\]

%dsp-level:0
\[
\qquad \text{the Burnside count on one fixed profile is}\ 1\ [\text{machine}];
\]

%dsp-level:0
\[
\qquad \text{the strict minimum}\ (1,1,1):\ \ \text{unique among}\ \kappa\ge1,\ \textstyle\sum\kappa=3\ [\text{machine: exhaustive}];
\]

%dsp-level:0
\[
\qquad 6\ln2+\ln W(Q)+0\ \text{compared term by term with}\ \mathrm{Def\text{-}BM.13}(5)\ \text{at}\ E_{\mathrm{int}}=6\ [\text{machine}]
\]

%dsp-level:0
\[
(6):\ \binom{K-1}2\ \text{at}\ K=(3,4,5,6)=(1,3,6,10)\ [\text{machine}];
\]

%dsp-level:0
\[
\qquad \text{the Burnside orbits under the swap: fixed points}\ \#\{(a,a,b):\ 2a+b=K,\ a,b\ge1\},
\]

%dsp-level:0
\[
\ \ \text{orbits}
\]

%dsp-level:0
\[
= \tfrac12(\#\text{profiles}+\#\text{fixed}):
\]

%dsp-level:0
\[
\quad K=(3,4,5,6)\mapsto(1,2,4,6)\ [\text{machine: closed form and exhaustive enumeration agree}];
\]

%dsp-level:0
\[
\qquad \kappa\ge0:\ \binom{K+2}2\ \text{at}\ K=(1,3)=(3,10),\ \text{orbits}\ (2,6)\ [\text{machine}];
\]

%dsp-level:0
\[
\qquad \text{the infinite-fiber witness}\ (K{+}n,{-}n,\,0),
\]

%dsp-level:0
\[
\ n\in\mathbb Z_{>0}:\ \ \text{distinct profiles at every}\ n\ [\text{machine: injective in}\ n]
\]

%dsp-level:0
\[
(7)(\mathrm a):\ \ln1=0\ [\text{machine}];
\]

%dsp-level:0
\[
\qquad (\mathrm b):\ 2\ln2
\]

%dsp-level:0
\[
\ne 3\ \text{via the enclosure}\ \tfrac12<\ln2<1\ \text{with tail}\ <2^{-20}\ [\text{machine: exact rational partial sums}];
\]

%dsp-level:0
\[
\qquad (\mathrm c):\ 20>15\ [\text{machine}];
\]

%dsp-level:0
\[
\qquad (\mathrm d):\ \text{the token census of both normalizations}\ [\text{machine}];
\]

%dsp-level:0
\[
\qquad (\mathrm e):\ \mathrm{FT.73}(2)\text{'s degree row re-cited}
\]

%dsp-level:0
\[
(8):\ \frac{4\hat\mu'^2\sin^2(\Theta_R/4)}{4E^2(\ln2)^2/\mathfrak h^2}
\]

%dsp-level:0
\[
= \frac{\hat\mu'^2\mathfrak h^2}{(\ln2)^2}\cdot\frac{D_S}{4E^2}\ \text{at}\ \Theta_R
\]

%dsp-level:0
\[
= \delta_v\ [\text{machine: the Laurent assembly}];
\]

%dsp-level:0
\[
\qquad \text{the five factors re-cited from}\ \text{Census FT.63}(6);
\]

%dsp-level:0
\[
\qquad (9):\ \text{the status and record assembly}. \ \blacksquare
\]

\textbf{Phys.} The joint entropy is not postulated; it is counted. Three registers — the crystallized-edge bits, the charge-filtered end census, and the winding-and-sealing fiber — are independent by census, so their cardinalities multiply and their logarithms add: the additive law of the refined rest input is a cardinality identity, and its first two factors are the corpus's own two-to-the-E and W-of-Q. The third factor is derived exactly where the mass read lives. On the stationary locus the geodesic-constancy theorem makes every surviving profile uniform, a uniform profile is fixed by the swap quotient, so the geometric fiber is a single orbit and its entropy is zero: the refined rest input is the complete joint entropy of the stationary baryon fiber, ledger part, and the off-stationary fibers are pairwise-cancelling content that no mass read consults. Away from stationarity the Platonic substrate carries closed forms — a character sum for the level counts, Burnside for the orbits, the frozen grammar counts for the sealings — but the baryon junction's own general fiber is infinite at fixed winding without a regime datum, and the tier names that missing enumeration instead of guessing a number. Then both currencies are laid side by side and every proposed bridge species is tested: the entropy difference vanishes along the ladder where the defect values vary, no joint potential exists to transform, the extremum picks the neutral class but converts nothing, the trace normalizations are exact on both sides and incommensurate between them, and the Casimir ratios are rational where the pole side is of degree sixteen. The exchange rate between count and angle stays one undeclared dimensionless number — the packaging the interface already logged — and the gap row records it under a name.

\begin{theorem}[{\texorpdfstring{BM.69 — The bridge converged: the three inputs, the non-derivation theorem, and the minimal priced residue}{BM.69 - The bridge converged: the three inputs, the non-derivation theorem, and the minimal priced residue}}]\label{tag:BM-69}
\textbf{Thm.}

%dsp-level:0
\[
(1)\quad\textbf{the three inputs assembled:}\qquad\begin{array}{l|l}\text{the force-operator quadratic invariant} & I_F=\operatorname{Tr}\widehat F^2-\operatorname{Tr}\widehat F_0^2=4\,\widehat Q^2\qquad[\mathrm{BM.66}(3)]\\\text{the holonomy lift} & \mathbb Z_2\ \text{Wilson sector derived};\ \ \text{the real lift obstructed with the record issued}\qquad[\mathrm{BM.67}(2),(6),(10)]\\\text{the constrained joint entropy} & S_{\mathrm{joint}}=E\ln2+\ln W(Q)+\ln\Omega_{\mathrm{geom}},\qquad\Omega_{\mathrm{geom}}^{\mathrm{stat}}=1\qquad[\mathrm{BM.68}(2),(5)]\end{array}
\]

%dsp-level:0
\[
(2)\quad\textbf{the bridge identity in its derived form:}
\]

%dsp-level:0
\[
\qquad \boxed{\ \frac{m^2_{\mathrm{pole}}}{m_{\mathrm{hid}}^2}\ =\ \mathcal R\bigl(E_{\mathrm{int}},\Theta_R\bigr)\ =\ \rho_{\mathrm B}\cdot\frac{D_S(\Theta_R)}{4\,E_{\mathrm{int}}^2},\qquad\rho_{\mathrm B}=\frac{\hat\mu'^2\mathfrak h^2}{(\ln2)^2}\in\mathfrak q^0\ }
\]

%dsp-level:0
\[
\qquad [\mathrm{FT.54}(4);\ \text{Census FT.63}(6);\ \mathrm{BM.68}(8);\ \text{machine}]:
\]

%dsp-level:0
\[
\mathcal R\ \text{is exact algebraic per state once}\ \rho_{\mathrm B}\ \text{is supplied:}
\]

%dsp-level:0
\[
\qquad \text{six states, six values of}\ \tfrac{D_S}{4E^2},\ \text{four fields, six minimal polynomials}
\]

%dsp-level:0
\[
\qquad [\mathrm{Def\text{-}BM.14}(3);\ \text{machine re-cited}];
\]

%dsp-level:0
\[
(3)\quad\textbf{the class-dependence census:}
\]

%dsp-level:0
\[
\qquad \text{the bridge's left-hand class dependence runs through}\ D_S\ \text{alone}\ —\ \text{floor and branch data, i.e. the flavour composition of the class, which is}\ \neg\text{derived}
\]

%dsp-level:0
\[
\qquad [\mathrm{BM.28}(7);\ \mathrm{Def\text{-}BM.8}(5);\ \mathrm{BM.59}(9)];
\]

%dsp-level:0
\[
E_{\mathrm{int}}=6\ \ \text{is class-shared}
\]

%dsp-level:0
\[
\qquad [\mathrm{BM.57}(4);\ \text{Census FT.35}(2)];
\]

%dsp-level:0
\[
\qquad \text{the information side's class dependence}\ \bigl(\ln W(Q):\ \Sigma_{\mathrm B}=\ln\tfrac43;\quad I_F:\ (0,4,4)\bigr)\ \text{is exact}
\]

%dsp-level:0
\[
\qquad [\mathrm{Def\text{-}BM.13}(4);\ \mathrm{BM.66}(3);\ \text{machine}];
\]

%dsp-level:0
\[
\textbf{the conditional pole row (the hypothesis on its face):}
\]

%dsp-level:0
\[
\qquad \textbf{hypothesis:}\ \text{the two baryon classes differ by exactly one branch substitution}\ [\mathrm{BM.59}(9);\ \mathrm{BM.23}(4)\ \text{pattern}]:
\]

%dsp-level:0
\[
\textbf{then}
\]

%dsp-level:0
\[
\qquad \Delta\bigl(m^2_{\mathrm{pole}}\bigr)\ =\ \hat\mu'^2\bigl(D_S(\mathcal F_g,+)-D_S(\mathcal F_g,-)\bigr)\ =\ \Delta_g\ >\ 0,
\]

%dsp-level:0
\[
\qquad \Delta_1=2\hat\mu'^2,\quad\Delta_2=(\sqrt2-1)\hat\mu'^2,\quad\Delta_3=\bigl(2\cos\tfrac\pi{10}-\sqrt3\bigr)\hat\mu'^2
\]

%dsp-level:0
\[
\qquad [\mathrm{Def\text{-}BM.14}(5);\ \text{machine: the certificates}\ 4>2,\ 2>1,\ 3<5]:
\]

%dsp-level:0
\[
\textbf{the same direction as the transport row};
\]

%dsp-level:0
\[
\qquad \text{and}\ \neg\textbf{addable to the informational assembly:}\ \ \text{a cross-currency addition is a category the corpus bars}
\]

%dsp-level:0
\[
\qquad [\mathrm{BM.35}(5);\ \mathrm{BM.60}(6);\ \mathrm{BM.37}(2);\ \text{pointers}];
\]

%dsp-level:0
\[
(4)\quad\textbf{the non-derivation theorem:}
\]

%dsp-level:0
\[
\qquad \boxed{\ \neg\exists\ \text{an exact identity}\ \ m^2_{\mathrm{pole}}=\mathcal B\bigl[E_{\mathrm{int}},\,W(Q),\,\Theta_R,\,I_F\bigr]\ \ \text{with}\ \mathcal B\ \text{built from derived data alone}\ }
\]

%dsp-level:0
\[
\bigl(\text{derived data}\ :=\ \text{exact trace ratios},\ \text{Casimir normalizations},\ \text{constrained fiber counts},\ \text{channel caps},\ \text{already-derived algebraic numbers}\bigr),
\]

%dsp-level:0
\[
\qquad \text{on four independent grounds:}
\]

%dsp-level:0
\[
(\mathrm i)\quad\textbf{the field degree:}\ \ \text{the five floors' pole factors generate a degree-}16\ \text{extension of}\ \mathbb Q;
\]

%dsp-level:0
\[
\ \ \text{every enumerated input}\ \bigl(E_{\mathrm{int}},\ W(Q),\ I_F,\ \text{caps},\ \text{fiber counts}\bigr)\ \text{is integer- or binary-valued};
\]

%dsp-level:0
\[
\ \ \text{a rational formula places every pole value in a field of degree}\ <16\ \text{over the declared constants}
\]

%dsp-level:0
\[
\qquad [\mathrm{FT.73}(2),(3);\ \text{machine re-cited}];
\]

%dsp-level:0
\[
(\mathrm{ii})\quad\textbf{the two-torsion image:}\ \ \text{the register lift of}\ \Theta_R\ \text{reaches}\ \{0,\pi\}\ \text{only}\ —\ \text{four of the six states are unreachable at every winding}
\]

%dsp-level:0
\[
\qquad [\mathrm{BM.67}(6)];
\]

%dsp-level:0
\[
(\mathrm{iii})\quad\textbf{the calibration weight:}\ \ \text{the winding route}\ \Theta_R=2\pi k\ \text{is ruled out}\ —\ \mathfrak h\text{-weight}\ {-}2\ \text{against}\ 0
\]

%dsp-level:0
\[
\qquad [\mathrm{FT.55}(1);\ \mathrm{BM.67}(8)];
\]

%dsp-level:0
\[
(\mathrm{iv})\quad\textbf{the strength bar:}\ \ \neg\exists\ \text{a count}\to\text{strength theorem};\ \ \text{the capacity ratios}\ \bigl(\tfrac13,\,8,\,\tfrac1{25}\bigr)\ \text{are}\ \neg\text{couplings}
\]

%dsp-level:0
\[
\qquad [\mathrm{BM.36}(3);\ \mathrm{BM.61}(8)\ (\mathrm O1);\ \mathrm{FT.73}(5)\ \text{trap pattern}];
\]

%dsp-level:0
\[
(5)\quad\textbf{the Gap-7 reduction (the smallest priced residue):}
\]

%dsp-level:0
\[
\qquad \boxed{\ \text{the bridge residue is exactly:}\quad(\mathrm R1)\ \ \text{the one undeclared dimensionless ratio}\ \ \rho_{\mathrm B}=\frac{\hat\mu'^2\mathfrak h^2}{(\ln2)^2}\qquad[\text{Cal-F packaging};\ \neg\text{added here}];}
\]

%dsp-level:0
\[
(\mathrm R2)\quad\text{the ledger realization of clean-enclosure vertex-figure loops}
\]

%dsp-level:0
\[
\qquad [\mathrm{BM.67}(9),(10);\ \text{Census FT.63}(4)\ \mathrm{Rmk}];
\]

%dsp-level:0
\[
\qquad (\mathrm R3)\quad\text{the flavour composition of the sealed classes}
\]

%dsp-level:0
\[
\qquad [\mathrm{BM.28}(7)]:
\]

%dsp-level:0
\[
\textbf{every other factor of}\ (2)\ \textbf{is derived exact algebra:}
\]

%dsp-level:0
\[
\qquad \text{the six}\ D_S\ \text{with their minimal polynomials};\quad\text{the five}\ \tfrac{D_S}{4E^2};
\]

%dsp-level:0
\[
\quad I_F=(0,4,4);
\]

%dsp-level:0
\[
\quad W=(1,6,15,20,15,6,1);\quad\Omega_{\mathrm{geom}}^{\mathrm{stat}}=1
\]

%dsp-level:0
\[
\qquad [\text{machine: the inventory row}];
\]

%dsp-level:0
\[
(6)\quad\textbf{the no-hiding row:}
\]

%dsp-level:0
\[
\qquad \textbf{the coefficient census of}\ \mathrm{BM.65}\text{--}\mathrm{BM.69}:
\]

%dsp-level:0
\[
\qquad \bigl\{\tfrac32=\tfrac{E_{\mathrm{int}}}4\ [\text{a count row}],\quad4\ [\text{the}\ I_F\ \text{trace normalization}],\quad2c_0,\ c_0^2+C^2\ [\text{annihilated in}\ I_F],\quad\rho_{\mathrm B}\ [\text{a Cal-F restatement}]\bigr\}:
\]

%dsp-level:0
\[
\textbf{each is a trace ratio, a Casimir normalization, a constrained fiber-count derivative, a channel-cap lock, or an exact operator identity}\ —\ \text{the organic test is satisfied coefficient by coefficient}
\]

%dsp-level:0
\[
\qquad [\text{machine: the census}];
\]

%dsp-level:0
\[
\alpha_B:\quad\neg\text{added},\quad\neg\text{derived},\quad\text{barred on}\ \mathrm{BM.61}(8)\text{'s four grounds, which stand}
\]

%dsp-level:0
\[
\qquad [\text{machine: the token scan of this tier}\ —\ \text{every occurrence is a barred-object row}];
\]

%dsp-level:0
\[
\neg\text{a P4-class temperature},
\]

%dsp-level:0
\[
\quad\neg\text{a free-energy coefficient},
\]

%dsp-level:0
\[
\quad\neg\text{a fitting parameter is introduced anywhere in this tier}
\]

%dsp-level:0
\[
\qquad [\text{machine: the symbol census against the P4 core and the era registries}];
\]

%dsp-level:0
\[
(7)\quad\textbf{status:}
\]

%dsp-level:0
\[
\qquad \textbf{derived:}\ (1)\text{--}(6);
\]

%dsp-level:0
\[
\qquad \textbf{not achieved: the full identity}\ —\ \text{the obstruction is exact and named}\ (4),
\]

%dsp-level:0
\[
\ \text{and the residue is minimal and priced}\ (5);
\]

%dsp-level:0
\[
\neg\text{a P4 addition},\quad\neg\text{a new dimensionful scale}\ [\mathrm{P1}],\quad\neg\text{a measured value}\ [\mathrm{P5}],\quad\neg\text{a declaration};
\]

%dsp-level:0
\[
\qquad \mathrm{Def\text{-}H.37}(2)\text{'s interface is cited as the locus of the declared transfer,}\ \neg\text{imported as a relating equation}.
\]

\textbf{Req.} P1, P4, P5 (Parts 1–2); A §E (the P4 core); Def-H.37(2), FT.54(4), FT.55(1), FT.73(2), FT.73(3), FT.73(5), FT.125(4) (Part 3); Census FT.35(2), FT.63(4), FT.63(6) (row-wise); Def-BM.8(5), Def-BM.13(4), Def-BM.14(3), Def-BM.14(4), Def-BM.14(5), BM.23(4), BM.28(7), BM.35(5), BM.36(3), BM.37(2), BM.57(4), BM.59(9), BM.60(6), BM.61(8), BM.66(3), BM.66(4), BM.67(2), BM.67(6), BM.67(8), BM.67(9), BM.67(10), BM.68(2), BM.68(5), BM.68(8) (this part); Cal-F; exact radical separations, Laurent-monomial assembly, symbol censuses [math; machine: \texttt{checks/\allowbreak{}verify\_\allowbreak{}H\_\allowbreak{}organic\_\allowbreak{}bridge\_\allowbreak{}post\_\allowbreak{}bm14.py} §A, §E, §G].

\textbf{Pf.}

%dsp-level:0
\[
(1):\ \text{the three rows re-cited from their own tags};
\]

%dsp-level:0
\[
\qquad (2):\ \mathrm{BM.68}(8)\text{'s identity re-cited};
\]

%dsp-level:0
\[
\qquad \text{the six values and four fields re-cited from}\ \mathrm{Def\text{-}BM.14}(3)\ \text{and}\ \text{Census FT.63}(6)
\]

%dsp-level:0
\[
(3):\ E_{\mathrm{int}}\ \text{class-shared}\ [\mathrm{BM.57}(4)\text{'s difference table re-cited}];
\]

%dsp-level:0
\[
\qquad \Delta_g\ \text{re-formed in the three exact fields with the certificates}\ 4>2,
\]

%dsp-level:0
\[
\ 2>1,
\]

%dsp-level:0
\[
\ \cos\tfrac\pi{10}>\tfrac{\sqrt3}2
\]

%dsp-level:0
\[
\iff 3<5\ [\text{machine}];
\]

%dsp-level:0
\[
\qquad \text{the addition bar:}\ \mathrm{BM.37}(2)\text{'s disjointness and}\ \mathrm{BM.35}(5)\text{'s no-conversion row re-cited}
\]

%dsp-level:0
\[
(4)(\mathrm i):\ \mathrm{FT.73}(2),
\]

%dsp-level:0
\[
(3)\ \text{re-cited: the degree-}16\ \text{compositum, the irreducible quartics, the quadratic-field exclusion}\ 4-2\sqrt2+2\sqrt3-\sqrt6\notin K_2;
\]

%dsp-level:0
\[
\qquad (\mathrm{ii})\text{--}(\mathrm{iv}):\ \mathrm{BM.67}(6),(8),\ \mathrm{BM.36}(3),\ \mathrm{BM.61}(8)\ \text{re-cited}
\]

%dsp-level:0
\[
(5):\ \text{the inventory row assembled from}\ \mathrm{BM.65}\text{--}\mathrm{BM.68}\text{'s derived lists}\ [\text{machine}];
\]

%dsp-level:0
\[
\qquad (6):\ \text{the coefficient census}\ [\text{machine}];
\]

%dsp-level:0
\[
\qquad \text{the token scan of this tier for the barred object}\ \alpha_B\ \text{and for P4-registry additions}\ [\text{machine}];
\]

%dsp-level:0
\[
\qquad (7):\ \text{assembly}. \ \blacksquare
\]

\textbf{Phys.} The three inputs meet, and the bridge exists in exactly the form the corpus can derive — and in no other. The operator invariant is the exact integer four Q-squared; the holonomy gives a binary sector and an obstructed angle; the joint entropy gives the count plus the charge census at stationarity, with a geometric factor of one. The well-posed bridge object, the pole-to-hidden mass ratio, factorizes into a per-state exact algebraic number and a single dimensionless packaging, and that packaging is the whole of the gap: everything else in the identity is derived, and the non-derivability of the rest is a theorem with four independent grounds — a degree-sixteen field against integer reads, a two-element image against six angles, a calibration weight against the winding route, and the standing bar against turning counts into strengths. The class dependence of the bridge lives entirely in the informational currency; its pole-currency image needs a flavour composition no clause attaches, and where a composition is hypothesized the pole splitting points the same way as the transport row but cannot be added to it without the very equation that is missing. So Gap 7 reduces to its smallest honest form: one scale ratio, one loop realization, one composition — named, priced, and hidden inside no coefficient.

\end{theorem}

\pseudohead{theorem}{\texorpdfstring{BM.70}{BM.70}}{\texorpdfstring{The beta-decay energy audit: the two-bit floor theorem, the margin ladder, and the formed-closed dichotomy}{The beta-decay energy audit: the two-bit floor theorem, the margin ladder, and the formed-closed dichotomy}}\label{tag:BM-70}
\textbf{Thm.}

%dsp-level:0
\[
(1)\quad\textbf{the floor and the Q-value:}\qquad F_{\ell\nu}\ :=\ m_\ell\,\mathfrak h\ +\ m_\nu\,\mathfrak h\qquad[\text{an abbreviation of derived reads}];
\]

%dsp-level:0
\[
\boxed{\ Q_\beta\,\mathfrak h\ =\ \ln\tfrac43\ -\ 6\,c\,\varepsilon_{\mathrm{EM}}\ln2\ -\ F_{\ell\nu}\ }\qquad[\mathrm{BM.62}(4);\ \mathrm{BM.63}(1);\ \text{machine: the monomial identity}];
\]

%dsp-level:0
\[
(2)\quad\textbf{the two-bit daughter floor (theorem on the informational species):}
\]

%dsp-level:0
\[
\qquad \textbf{the chain:}
\]

%dsp-level:0
\[
\qquad \text{the emitted charged member is the}\ \ell\text{-channel line}\ —\ \text{the forced, class-unique carrier}
\]

%dsp-level:0
\[
\qquad [\mathrm{BM.23}(5);\ \mathrm{BM.41}(1);\ \mathrm{Def\text{-}BM.6}(1)]:
\]

%dsp-level:0
\[
\text{every channel line is a materialized bound chain}
\]

%dsp-level:0
\[
\qquad [\mathrm{FT.109};\ \mathrm{FT.112};\ \mathrm{FT.113};\ \mathrm{Def\text{-}H.4}];
\]

%dsp-level:0
\[
\qquad \text{every materialized closed chain carries}\ \ge2\ \text{directed edges}\ —\ \text{the digon is the minimal floor}
\]

%dsp-level:0
\[
\qquad [\mathrm{Def\text{-}H.4}:\ \delta h(c)\ge2;\ \text{Census FT.28};\ \mathrm{FT.32}:\ L\ge2;\ \mathrm{BM.13}(3):\ E_{\mathrm{int}}\ge2]:
\]

%dsp-level:0
\[
\text{the informational read}\ \ m\,\mathfrak h=E_{\mathrm{int}}\ln2
\]

%dsp-level:0
\[
\qquad [\text{Census FT.44}(4);\ \mathrm{Def\text{-}H.9}(2);\ \mathrm{BM.43}(2):\ \text{the taxonomy law}]
\]

%dsp-level:0
\[
\qquad \implies
\]

%dsp-level:0
\[
\qquad m_\ell\,\mathfrak h\ \ge\ 2\ln2\ >\ 0
\]

%dsp-level:0
\[
\qquad [\text{machine: the exact enclosure}\ \tfrac12<\ln2<1];
\]

%dsp-level:0
\[
m_\nu\,\mathfrak h\ \ge\ 0
\]

%dsp-level:0
\[
\qquad [\mathrm{EW.33}:\ \text{the Gram is PSD};\ \mathrm{BM.41}(4)(\mathrm a)];
\]

%dsp-level:0
\[
\qquad \text{and on the informational read}\ \ m_\nu\,\mathfrak h=E_{\mathrm{int}}^{(\nu)}\ln2\ \ge\ 2\ln2\ \ \text{likewise}
\]

%dsp-level:0
\[
\qquad [\mathrm{BM.43}(2);\ \mathrm{BM.13}(3)]:
\]

%dsp-level:0
\[
\boxed{\ F_{\ell\nu}\ \ge\ 2\ln2\ }
\]

%dsp-level:0
\[
\qquad \textbf{with the floor ladder:}
\]

%dsp-level:0
\[
\qquad 2\ln2\ \ \bigl(\text{the PSD row alone}\bigr)
\]

%dsp-level:0
\[
\quad <
\]

%dsp-level:0
\[
\quad 4\ln2\ \ \bigl(\text{the informational neutrino}\bigr)
\]

%dsp-level:0
\[
\quad <
\]

%dsp-level:0
\[
\quad 5\ln2\ \ \bigl(\text{with the daughters' own transport refinement}\bigr)
\]

%dsp-level:0
\[
\qquad [\text{machine:}\ 2<4<5]:
\]

%dsp-level:0
\[
\textbf{the refinement row:}
\]

%dsp-level:0
\[
\qquad k=1:\quad W(q)=m_q=(1,\,2,\,1)\ \ \text{at}\ q=(-1,\,0,\,+1),
\]

%dsp-level:0
\[
\qquad \Sigma_{\mathrm{cl}}(q)=(\,0,\ \ln2,\ 0\,):
\]

%dsp-level:0
\[
\qquad \text{the neutral line carries one extra bit, the charged lines none}
\]

%dsp-level:0
\[
\qquad [\mathrm{Def\text{-}BM.13}(3);\ \mathrm{BM.63}(6);\ \text{machine: the convolution of}\ (2+z+z^{-1})^1]:
\]

%dsp-level:0
\[
\text{the daughters are}\ q_\ell
\]

%dsp-level:0
\[
= -1\ \ [\text{the charged member}]\ \ \text{and}\ q_\nu
\]

%dsp-level:0
\[
= 0\ \ [\text{the neutral member};\ \mathrm{BM.41}(1);\ \mathrm{FT.142}(1):\ \text{EM-unactivated}]
\]

%dsp-level:0
\[
\qquad \implies
\]

%dsp-level:0
\[
\qquad F_{\ell\nu}\ \ge\ 2\ln2+0+2\ln2+\ln2=5\ln2\ \ \text{on the fully refined informational read}
\]

%dsp-level:0
\[
\qquad [\text{machine}];
\]

%dsp-level:0
\[
(3)\quad\textbf{the margin ladder (exact, integer-certified):}
\]

%dsp-level:0
\[
\qquad Q_\beta\,\mathfrak h\ \le\ \ln\tfrac43-F_{\ell\nu}\ \ \text{at every}\ \varepsilon_{\mathrm{EM}}>0
\]

%dsp-level:0
\[
\qquad [\mathrm{BM.62}(6):\ \text{strict monotonicity}]:
\]

%tbl-level:0
{\setlength{\tabcolsep}{4pt}
\begin{longtable}{>{\raggedright\arraybackslash}p{\dimexpr0.3642\textwidth-2\tabcolsep\relax}|>{\raggedright\arraybackslash}p{\dimexpr0.1199\textwidth-2\tabcolsep\relax}|>{\raggedright\arraybackslash}p{\dimexpr0.1214\textwidth-2\tabcolsep\relax}|>{\raggedright\arraybackslash}p{\dimexpr0.3945\textwidth-2\tabcolsep\relax}}
\toprule
floor used & $F_{\ell\nu}\ \ge$ & $Q_\beta\,\allowbreak \mathfrak h\ \le$ & the certificate \\
\addlinespace[2pt]
\midrule
\endfirsthead
\toprule
floor used & $F_{\ell\nu}\ \ge$ & $Q_\beta\,\allowbreak \mathfrak h\ \le$ & the certificate \\
\addlinespace[2pt]
\midrule
\endhead
\bottomrule
\endlastfoot
the PSD row alone & $2\ln2$ & $-\ln3$ & $\ln\tfrac43<\allowbreak 2\ln2\iff\allowbreak \tfrac43<\allowbreak 4\iff\allowbreak 4<\allowbreak 12$ \\
the informational neutrino & $4\ln2$ & $-\ln12$ & $\ln\tfrac43<\allowbreak 4\ln2\iff\allowbreak \tfrac43<\allowbreak 16\iff\allowbreak 4<\allowbreak 48$ \\
the fully refined daughters & $5\ln2$ & $-\ln24$ & $\ln\tfrac43<\allowbreak 5\ln2\iff\allowbreak \tfrac43<\allowbreak 32\iff\allowbreak 4<\allowbreak 96$ \\
\end{longtable}
}

%dsp-level:0
\[
[\text{machine: the prime-log identities}\ \ \ln\tfrac43-2\ln2=-\ln3,\quad\ln\tfrac43-4\ln2=-\ln12,\quad\ln\tfrac43-5\ln2=-\ln24;\ \ \text{and every comparison an integer certificate}];
\]

%dsp-level:0
\[
\boxed{\ Q_\beta\ <\ 0\ \ \text{strictly, with margin}\ \ge\ \tfrac{\ln3}{\mathfrak h},\ \ \text{deepened by}\ \tfrac{6\,c\,\varepsilon_{\mathrm{EM}}\ln2}{\mathfrak h}\ \ \text{at every positive coupling}\ }
\]

%dsp-level:0
\[
\qquad [\mathrm{BM.63}(2),(3),(4)\ \text{audited and standing};\ \text{machine}];
\]

%dsp-level:0
\[
(4)\quad\textbf{the formed-closed dichotomy (the audit verdict):}
\]

%dsp-level:0
\[
\qquad Q_\beta\ \text{is a formed expression only where all four members are read in one currency}:
\]

%dsp-level:0
\[
(\mathrm a)\quad\textbf{the informational currency on every member}\ \bigl[\text{the parents:}\ \mathrm{Def\text{-}BM.13}(5);\ \ \text{the daughters:}\ \mathrm{BM.43}(2)\text{'s taxonomy law}\bigr]:
\]

%dsp-level:0
\[
\qquad (2)\text{'s floor theorem applies}\implies Q_\beta\,\mathfrak h\le-\ln3<0
\]

%dsp-level:0
\[
\qquad \textbf{formed}\implies\textbf{closed};
\]

%dsp-level:0
\[
(\mathrm b)\quad\textbf{the transport/pole species on the leptons}\ \bigl[\mathrm{BM.39}(1):\ m_\ell^2=\hat\mu^2\ell_c^2\,G(c(i))>0,\ \ \text{the scale}\ \hat\mu\ \text{undeclared}\ [\mathrm{Cal\text{-}F}];\ \mathrm{Def\text{-}BM.11}\bigr]:
\]

%dsp-level:0
\[
\qquad \text{the subtraction against informational-currency parents is a cross-currency expression};
\]

%dsp-level:0
\[
\text{its relating equation is}\ \neg\text{derived}
\]

%dsp-level:0
\[
\qquad [\mathrm{BM.35}(5);\ \mathrm{BM.43}(2):\ \text{the conversion}\ E_{\mathrm{int}}^{(i)}\leftrightarrow\text{transport data is item-7 residue};\ \mathrm{BM.60}(6);\ \text{pointers}]
\]

%dsp-level:0
\[
\qquad \implies
\]

%dsp-level:0
\[
\qquad Q_\beta\ \ \neg\textbf{formed}\implies\ \neg\textbf{a verdict};
\]

%dsp-level:0
\[
\boxed{\ \text{on every read on which}\ Q_\beta\ \text{is a derived expression,}\ Q_\beta<0\ \ \text{strictly};\qquad\neg\exists\ \text{a derived read on which}\ Q_\beta\ge0\ }
\]

%dsp-level:0
\[
\qquad [\text{machine: the row census over the complete route inventory}\ \mathrm{BM.30}(1):\ \ \text{route}\ (\mathrm i)\ Q=-(m_\ell{+}m_\nu)<0\ [\mathrm{BM.41}(4)(\mathrm b)];\ \ \text{route}\ (\mathrm{ii})\ \text{priced}\ \neg\text{adopted, its sign undetermined without scale data}\ [\mathrm{BM.41}(4)(\mathrm d)];\ \ \text{route}\ (\mathrm{iii})\ \text{amplitude layer}\ —\ \text{amplitude}\ne\text{mass}\ [\mathrm{FT.79}(5);\ \mathrm{BM.37}(3)];\ \ \text{the curvature-flux read:}\ Q<0\ [\mathrm{BM.41}(4)(\mathrm c)]];
\]

%dsp-level:0
\[
(5)\quad\textbf{the missing lepton read, named exactly:}
\]

%dsp-level:0
\[
\qquad E_{\mathrm{int}}^{(\ell)},\ E_{\mathrm{int}}^{(\nu)}\ \ \neg\text{computed at any floor}
\]

%dsp-level:0
\[
\qquad [\mathrm{BM.43}(2);\ \mathrm{BM.63}(8)];
\]

%dsp-level:0
\[
\qquad \text{the line index within the}\ \nu\text{-triplet}\ \neg\text{derived}
\]

%dsp-level:0
\[
\qquad [\mathrm{FT.145}(3):\ \neg\text{a mixing matrix}]:
\]

%dsp-level:0
\[
\textbf{these reads sharpen the margin and cannot breach the floor:}
\]

%dsp-level:0
\[
\qquad E_{\mathrm{int}}\ge2\ \ \text{is a structural lower bound of the species}
\]

%dsp-level:0
\[
\qquad [\mathrm{BM.13}(3);\ \text{Census FT.28}];
\]

%dsp-level:0
\[
\qquad \text{the margin is monotone in}\ E_{\mathrm{int}}^{(\ell)},\ E_{\mathrm{int}}^{(\nu)}
\]

%dsp-level:0
\[
\qquad [\text{machine: each added edge adds}\ \ln2]:
\]

%dsp-level:0
\[
\boxed{\ \text{the closure is}\ \neg\textbf{an artifact of an unfinished lepton mass read: the unfinished read moves the margin, never the sign}\ };
\]

%dsp-level:0
\[
(6)\quad\textbf{no external Q-value:}
\]

%dsp-level:0
\[
\qquad \neg\exists\ \text{an observed endpoint, mass difference or}\ Q\text{-value in any clause of this tier}
\]

%dsp-level:0
\[
\qquad [\mathrm{P5};\ \text{machine: the token scan}\ —\ \text{zero external numerals, zero unit tokens}];
\]

%dsp-level:0
\[
\textbf{the consistency-check row:}
\]

%dsp-level:0
\[
\qquad \text{a comparison against an observed}\ Q_\beta\ \text{is admissible only as a quarantined audit row}
\]

%dsp-level:0
\[
\qquad [\text{Census FT.114},\ \text{FT.120}\ \text{pattern};\ \mathrm{P1},\ \mathrm{P5}]\ \ \text{and only after the corpus derives its own scale}\ —\ \text{one calibration datum fixing}\ e\ \ [\mathrm{EM.25}(7)]\ \ \text{and a}\ \mathfrak h\ \text{read}\ [\mathrm{P1}]:
\]

%dsp-level:0
\[
\neg\exists\ \text{such a derivation in Parts 1--4};\qquad\textbf{this tier performs no comparison};
\]

%dsp-level:0
\[
(7)\quad\textbf{status:}
\]

%dsp-level:0
\[
\qquad \textbf{derived:}\ (1)\text{--}(6);
\]

%dsp-level:0
\[
\qquad \mathrm{BM.63}(2)\text{'s closure is confirmed as a theorem on the formed reads, with its species conditionality on its face}\ (4)\ \ [\text{A §0.2a}];
\]

%dsp-level:0
\[
\neg\text{a measured lifetime, mass difference or endpoint enters}\ [\mathrm{P5}];
\]

%dsp-level:0
\[
\qquad \neg\text{a P4 addition};
\]

%dsp-level:0
\[
\qquad \neg\text{a new dimensionful scale}\ [\mathrm{P1}];
\]

%dsp-level:0
\[
\qquad \text{the open rows:}\ E_{\mathrm{int}}^{(\ell)},\ E_{\mathrm{int}}^{(\nu)},\ \text{the}\ \nu\text{-line index},\ \text{the numerical value of}\ e
\]

%dsp-level:0
\[
\qquad [\text{pointers}].
\]

\textbf{Req.} EW.33 (Part 2); EM.25(7) (Part 2); P1, P5 (Parts 1–2); Def-H.4, Def-H.9(2), FT.32, FT.79(5), FT.109, FT.112, FT.113, FT.142(1), FT.145(3) (Part 3); Census FT.28, FT.44(4), FT.114, FT.120 (row-wise); Def-BM.6(1), Def-BM.11, Def-BM.13(3), Def-BM.13(5), BM.13(3), BM.23(5), BM.30(1), BM.35(5), BM.37(3), BM.39(1), BM.41(1), BM.41(4), BM.43(2), BM.60(6), BM.62(4), BM.62(6), BM.63(1), BM.63(2), BM.63(3), BM.63(4), BM.63(6), BM.63(8) (this part); Cal-F; exact prime-log-ring identities with integer certificates, exact rational $\ln2$ enclosures, monotonicity censuses, token scans [math; machine: \texttt{checks/\allowbreak{}verify\_\allowbreak{}H\_\allowbreak{}organic\_\allowbreak{}bridge\_\allowbreak{}post\_\allowbreak{}bm14.py} §F].

\textbf{Pf.}

%dsp-level:0
\[
(1):\ \mathrm{BM.63}(1)\text{'s row re-cited and re-formed in the}\ \ln\text{-ring}\ [\text{machine}];
\]

%dsp-level:0
\[
\qquad (2):\ \delta h(c)\ge2\ [\mathrm{Def\text{-}H.4}\ \text{re-cited}];
\]

%dsp-level:0
\[
\qquad \text{the digon card}\ E_{\mathrm{int}}=2\ [\text{Census FT.28}\ \text{re-cited}];
\]

%dsp-level:0
\[
\qquad E_{\mathrm{int}}\ge2\ \text{on every computed cell}\ \{2,6,12,24,60\}\ [\text{machine: the minimum}];
\]

%dsp-level:0
\[
\qquad m_\ell\mathfrak h=E_{\mathrm{int}}^{(\ell)}\ln2\ge2\ln2\ [\text{machine}];
\]

%dsp-level:0
\[
\qquad m_\nu\mathfrak h\ge0\ [\mathrm{EW.33}\ \text{PSD re-cited}]
\]

%dsp-level:0
\[
\text{the refinement:}\ (2+z+z^{-1})^1\ \text{has coefficients}\ (1,2,1)\ \text{at}\ Q
\]

%dsp-level:0
\[
= (-1,0,+1)\ [\text{machine: the convolution}];
\]

%dsp-level:0
\[
\qquad \ln m_q=(0,\ln2,0)\ [\text{machine}];
\]

%dsp-level:0
\[
\qquad 2\ln2+2\ln2+\ln2=5\ln2\ [\text{machine: the prime-log ring}];
\]

%dsp-level:0
\[
\qquad 2<4<5\ [\text{machine}]
\]

%dsp-level:0
\[
(3):\ \ln\tfrac43-2\ln2=(2\ln2-\ln3)-2\ln2=-\ln3;
\]

%dsp-level:0
\[
\qquad \ln\tfrac43-4\ln2=-2\ln2-\ln3=-\ln(4\cdot3)=-\ln12;
\]

%dsp-level:0
\[
\qquad \ln\tfrac43-5\ln2=-3\ln2-\ln3=-\ln(8\cdot3)=-\ln24\ [\text{machine: the prime-log ring}];
\]

%dsp-level:0
\[
\qquad \text{the three certificates:}\ 4<12,\ 4<48,\ 4<96\ [\text{machine: integers}];
\]

%dsp-level:0
\[
\qquad \text{the EM deepening:}\ 6c\varepsilon_{\mathrm{EM}}\ln2>0\ \text{at}\ c,\varepsilon_{\mathrm{EM}},\ln2>0\ [\text{machine: the coefficient census}]
\]

%dsp-level:0
\[
(4):\ \text{the route inventory re-cited from}\ \mathrm{BM.30}(1),
\]

%dsp-level:0
\[
\ \mathrm{BM.41}(4)(\mathrm a)\text{--}(\mathrm d)\ \text{and classified by currency}\ [\text{machine: the row census}];
\]

%dsp-level:0
\[
\qquad \text{the cross-currency bar:}\ \mathrm{BM.35}(5)\text{'s no-conversion row and}\ \mathrm{BM.37}(2)\text{'s disjointness re-cited}
\]

%dsp-level:0
\[
(5):\ \text{monotonicity:}\ \partial\bigl(-Q_\beta\mathfrak h\bigr)/\partial E_{\mathrm{int}}^{(i)}=\ln2>0\ [\text{machine}];
\]

%dsp-level:0
\[
\qquad (6):\ \text{the token scan}\ [\text{machine: zero external numerals}];
\]

%dsp-level:0
\[
\qquad (7):\ \text{assembly}. \ \blacksquare
\]

\textbf{Phys.} The audit asked whether the two-bit daughter floor is a theorem or an artifact of an unfinished lepton read, and the answer separates the two cleanly. It is a theorem wherever the question is formed: the emitted charged member is a line of the lepton channel, every channel line is a materialized bound chain, every materialized closed chain carries at least two directed edges because the digon is the floor, and the informational read turns that count into at least two bits; the neutral member costs nothing less than zero by positive semidefiniteness, and nothing less than two more bits plus one refinement bit if it too is read informationally. The margin ladder follows in exact integers — against three, against twelve, against twenty-four — and the field-energy term deepens every rung at every positive coupling. The only escape is currency: read the leptons through the transport species with its undeclared scale and the Q-value is not a formed expression at all, because the equation relating the two currencies is the corpus's open item. Hence the dichotomy, which is the audit's verdict: formed implies closed, unformed implies no verdict, and no derived read opens the vertex. What the unfinished lepton read moves is the margin's sharpness, never the sign — the floor is structural, sitting at the digon, below every rung of every ladder. And the observed endpoint enters nowhere: it may be compared only in a quarantined audit row, after a scale derivation that does not exist, and this tier performs no comparison.

\pseudohead{theorem}{\texorpdfstring{BM.71}{BM.71}}{\texorpdfstring{The opening-term census: the five routes, the neutron row under the bridge, and the six verdict rows}{The opening-term census: the five routes, the neutron row under the bridge, and the six verdict rows}}\label{tag:BM-71}
\textbf{Thm.}

%dsp-level:0
\[
(1)\quad\textbf{the route census (each candidate: derived, priced, or barred):}
\]

%tbl-level:3
{\scriptsize\setlength{\tabcolsep}{4pt}
\begin{longtable}{>{\raggedright\arraybackslash}p{\dimexpr0.0507\textwidth-2\tabcolsep\relax}|>{\raggedright\arraybackslash}p{\dimexpr0.1244\textwidth-2\tabcolsep\relax}|>{\raggedright\arraybackslash}p{\dimexpr0.1032\textwidth-2\tabcolsep\relax}|>{\raggedright\arraybackslash}p{\dimexpr0.7218\textwidth-2\tabcolsep\relax}}
\toprule
\# & the candidate opening term & verdict & the derived ground \\
\addlinespace[2pt]
\midrule
\endfirsthead
\toprule
\# & the candidate opening term & verdict & the derived ground \\
\addlinespace[2pt]
\midrule
\endhead
\bottomrule
\endlastfoot
1 & time-energy uncertainty / the non-uniform clock & \textbf{barred} & the Mandelstam–Tamm relation bounds the splitting's \textit{resolvability}, $\ne\allowbreak g$its magnitude $[\mathrm{BM.53}(4)\text{–}(6)]$; the closure is the non-existence of real on-shell solutions — $2m_n(E_p-m_p)=\allowbreak -m_\ell^2$, $\lambda=\allowbreak m_\ell^2(m_\ell^2-4m_n^2)$: polynomial identities at every tick $[\mathrm{BM.41}(2);\ \text{machine}]$; $\ne\allowbreak g\exists$ a time operator at finite dimension, so an energy loan over a $\Delta t$ has no operator content $[\mathrm{BM.10}(4)]$; the ledger is forward-only $[\mathrm{A6}:\ \ne\allowbreak g\exists\tau^{-1}]$, strictly contracting $[\mathrm{L2};\ \mathrm{BM.51}(2),\allowbreak (3)]$, with per-tick weights $\Gamma_{\mathrm{tr}}\in\allowbreak [0,\allowbreak \tfrac13]$ $[\mathrm{BM.44}(2)]$; the flicker ensemble is amplitude-layer content, and amplitude $\ne$ mass $[\mathrm{FT.79}(5);\ \mathrm{BM.37}(3);\ \mathrm{FT.67}(3\mathrm a);\ \mathrm{FT.68}(4)]$ \\
2 & field reconfiguration energy & \textbf{$\ne\allowbreak g$constructed; sign-barred at the constructed order} & radiation, back-reaction and the transition's field reconfiguration are the named open scope $[\mathrm{Def\text{-}BM.12}(7);\ \mathrm{BM.50}(3)]$; at the constructed monopole order every self-energy is the positive integral of a square $[\mathrm{BM.46}(1);\ \mathrm{Def\text{-}BM.12}(5)]$, and the transition $Q:0\to(+\allowbreak 1)\oplus(-1)$ raises the total field energy — the final state carries two charged members' self-energies where the initial carries none $[\mathrm{BM.49}(2);\ \text{machine: the sign census at positive symbols}]$; the attractive cross-well $V=\allowbreak -\varepsilon_{\mathrm{EM}}/d$ belongs to the $\ne\allowbreak g$constructed pair system $[\mathrm{BM.48}(1),\allowbreak (3),\allowbreak (4);\ \mathrm{BM.54}(4):\ \text{the category bar}]$ \\
3 & binding / nearby-space correction & \textbf{barred} & the category bar: a binding energy is the energy of the bound system, $\ne\allowbreak g$a correction to a free member's mass $[\mathrm{BM.48}(4);\ \mathrm{BM.54}(4)]$; the lepton–baryon bound channel is $\ne\allowbreak g$constructed $[\mathrm{BM.48}(3);\ \mathrm{BM.6}(5):\ \text{the dynamics of}\ Q]$; the nearby space holds only a strict fraction $f(k)=\allowbreak 1-\tfrac1k<\allowbreak 1$ $[\mathrm{BM.47}(6);\ \text{machine: exact rationals}]$; the count binding is zero on both rows, $6-(2{+}2{+}2)=\allowbreak 0$ $[\mathrm{BM.54}(1);\ \text{machine}]$; the ensemble binding is class-degenerate $[\mathrm{BM.54}(2);\ \mathrm{FT.107}(3);\ \mathrm{BM.30}(3)]$ \\
4 & $D^-$ redistribution & \textbf{barred} & the dark sector's contribution to the class splitting is exactly zero: $D^+\allowbreak =\allowbreak \Lambda\,\allowbreak \mathbb I$ uniform, both locks class-shared $[\mathrm{BM.52}(3)\ D^\pm\ \text{row};\ \mathrm{CT4.5};\ \mathrm{BM.28}(1)]$; the junction's $D^-$ content is identical on both rows — the same vertex types $n(q_1)=\allowbreak n(p_2)=\allowbreak 4$, the same lock ledger $[\text{Census FT.35}(4);\ \mathrm{BM.57}(4):\ \text{the difference table is identically zero};\ \text{machine}]$, so $\ne\allowbreak g\exists$ a class-dependent object to redistribute; locked content is retained and $\ne\allowbreak g$accessible $[\mathrm{D4.2};\ \mathrm{D1.11};\ \mathrm{CT4.3};\ \mathrm{FT.3}]$ — no channel reads it into a vertex budget; and $\ne\allowbreak g\exists$ a response read — the dielectric non-existence inventory theorem $[\mathrm{FT.73}(1);\ \text{any dynamics supplying one contradicts}\ \mathrm{CT4.2}/\mathrm{D1.11}]$ \\
5 & read specification & \textbf{the sole surviving door — not walked} & opening the vertex is a read-specification act requiring a §I addendum $[\mathrm{BM.49}(2);\ \mathrm{BM.41}(4)(\mathrm e);\ \mathrm{BM.30}(5);\ \mathrm{BM.45}]$; the route-(ii) measure price $T(0)=\allowbreak \tfrac43T(\pm1)$, its sign undetermined without scale data, stays priced $\ne\allowbreak g$adopted $[\mathrm{BM.41}(4)(\mathrm d);\ \mathrm{BM.28}(4);\ \mathrm{BM.37}(3)]$; \textbf{this tier performs no specification} \\
\end{longtable}
}

%dsp-level:0
\[
(2)\quad\textbf{completeness and the verdict:}
\]

%dsp-level:0
\[
\qquad \text{the five routes exhaust the audit's candidate list};
\]

%dsp-level:0
\[
\qquad \text{each is barred by a cited theorem or}\ \neg\text{constructed with its scope named}
\]

%dsp-level:0
\[
\qquad [\text{machine: the census-row structure}];
\]

%dsp-level:0
\[
\textbf{the MT bound's legitimate content:}
\]

%dsp-level:0
\[
\qquad \Delta t\ \ge\ \frac1{2\,\Delta m_{\mathrm{EM}}}\ =\ \frac{\pi\,\mathfrak h}{3\,c\,e^2\ln2}
\]

%dsp-level:0
\[
\qquad [\mathrm{BM.53}(5);\ \text{machine: the exact Laurent reciprocal}]
\]

%dsp-level:0
\[
\qquad \text{is a bound on clocks resolving the splitting,}\ \neg\text{a budget paying for a decay};
\]

%dsp-level:0
\[
\boxed{\ \textbf{the beta-decay no-go: on every formed derived read the vertex is kinematically closed, and its opening is}\ \neg\textbf{derivable from the corpus}\ —\ \textbf{it requires a read specification}\ }
\]

%dsp-level:0
\[
\qquad [\mathrm{BM.70}(4);\ (1);\ \text{machine}];
\]

%dsp-level:0
\[
(3)\quad\textbf{the neutron row under the bridge:}\qquad\boxed{\ m_n-m_p\ =\ \Delta m_{\mathrm{transport}}\ +\ \Delta m_{\mathrm{EM}}\ +\ \Delta m_{\mathrm{bridge}}\ }
\]

%dsp-level:0
\[
\Delta m_{\mathrm{transport}}\ =\ +\,\frac{\ln(4/3)}{\mathfrak h}
\]

%dsp-level:0
\[
\qquad [\mathrm{BM.59}(4);\ \mathrm{BM.65}(2):\ \text{derived}];
\]

%dsp-level:0
\[
\qquad \Delta m_{\mathrm{EM}}\ =\ -\,\frac{6\,c\,\varepsilon_{\mathrm{EM}}\ln2}{\mathfrak h}
\]

%dsp-level:0
\[
\qquad [\mathrm{BM.46}(4);\ \mathrm{BM.62}(3);\ \text{the operator form}\ \mathrm{BM.66}(4):\ \text{derived};\ c\ \text{priced};\ \varepsilon_{\mathrm{EM}}\ \text{symbolic}];
\]

%dsp-level:0
\[
\boxed{\ \Delta m_{\mathrm{bridge}}\ =\ 0\ \ \text{in the informational currency}\ }
\]

%dsp-level:0
\[
\qquad \textbf{(theorem):}\ \ \mathrm{BM.62}(1)\text{'s decomposition is complete in that currency}
\]

%dsp-level:0
\[
\qquad [\mathrm{BM.58}(6)\text{'s uniqueness row};\ \mathrm{BM.62}(11)\text{'s alphabets};\ \mathrm{BM.66}(6)\text{'s guard};\ \mathrm{BM.68}(5)\text{'s stationary slice}],
\]

%dsp-level:0
\[
\text{and the bridge's pole-currency content is unformed without the relating equation}
\]

%dsp-level:0
\[
\qquad [\mathrm{BM.69}(3),(4)];
\]

%dsp-level:0
\[
\textbf{the conditional pole row (the hypothesis on its face):}
\]

%dsp-level:0
\[
\qquad \text{under}\ \mathrm{BM.59}(9)\text{'s substitution hypothesis:}\quad\Delta\bigl(m^2_{\mathrm{pole}}\bigr)=+\Delta_g>0
\]

%dsp-level:0
\[
\qquad [\mathrm{BM.69}(3);\ \text{machine: the three certificates}]
\]

%dsp-level:0
\[
\qquad \text{— the same direction as the transport row,}\ \neg\text{a formed summand of any}\ Q\text{-value};
\]

%dsp-level:0
\[
(4)\quad\textbf{the six verdict rows:}
\]

%dsp-level:0
\[
(\mathrm i)\quad\textbf{the sign of}\ m_n-m_p:
\]

%dsp-level:0
\[
\qquad >0\ \ \text{on the sub-critical branch}\ \varepsilon_{\mathrm{EM}}<\varepsilon_{\mathrm{crit}}(c);
\]

%dsp-level:0
\[
\qquad =0\ \ \text{at}\ \varepsilon_{\mathrm{crit}}(c);
\]

%dsp-level:0
\[
\qquad <0\ \ \text{above}
\]

%dsp-level:0
\[
\qquad [\mathrm{BM.62}(6),(9);\ \mathrm{BM.65}(3)\text{--}(5)]:
\]

%dsp-level:0
\[
\qquad \text{which branch the corpus sits on is open}\ —\ \text{the numerical value of}\ e
\]

%dsp-level:0
\[
\qquad [\mathrm{EM.25}(7);\ \mathrm{BM.62}(12)];
\]

%dsp-level:0
\[
(\mathrm{ii})\quad\textbf{the sign of}\ Q_\beta:\qquad\text{strictly negative at every}\ \varepsilon_{\mathrm{EM}}>0\ \ \text{on every formed read}\qquad[\mathrm{BM.70}(3),(4)];
\]

%dsp-level:0
\[
(\mathrm{iii})\quad\textbf{the decay:}
\]

%dsp-level:0
\[
\qquad \textbf{closed}\ —\ \neg\text{conditionally open and}\ \neg\text{obstructed by an artifact: the closure is a theorem on the formed reads and unformable on the unformed ones}
\]

%dsp-level:0
\[
\qquad [\mathrm{BM.70}(4);\ (2)];
\]

%dsp-level:0
\[
(\mathrm{iv})\quad\textbf{the derived terms:}
\]

%dsp-level:0
\[
\qquad \Delta m_{\mathrm{transport}};\quad\Delta m_{\mathrm{EM}}\ \text{(scalar and operator)};
\]

%dsp-level:0
\[
\quad I_F=4\,\widehat Q^2;
\]

%dsp-level:0
\[
\quad F_{\ell\nu}\ge2\ln2;\quad\text{the margin ladder}\ (\ln3,\ \ln12,\ \ln24);\quad\Delta m_{\mathrm{bridge}}=0\ \text{in currency};\quad\Omega_{\mathrm{geom}}^{\mathrm{stat}}=1;
\]

%dsp-level:0
\[
(\mathrm v)\quad\textbf{the priced terms:}
\]

%dsp-level:0
\[
\qquad c\in\bigl\{\tfrac12,\tfrac35\bigr\}\ \ \text{priced}\ \neg\text{adopted}\ [\mathrm{Def\text{-}BM.12}(4)];
\]

%dsp-level:0
\[
\qquad \varepsilon_{\mathrm{EM}}\ \ \text{symbolic}\ [\mathrm{EM.25}(9);\ \mathrm{P5}];
\]

%dsp-level:0
\[
\qquad \hat\mu,\ \hat\mu',\ \mu\ \ \text{undeclared}\ [\mathrm{Cal\text{-}F}];
\]

%dsp-level:0
\[
\qquad \rho_{\mathrm B}\ \ \text{restated as the residue}\ [\mathrm{BM.69}(5)];
\]

%dsp-level:0
\[
(\mathrm{vi})\quad\textbf{the open terms:}
\]

%dsp-level:0
\[
\qquad \text{the numerical value of}\ e\ [\mathrm{EM.25}(7)];
\]

%dsp-level:0
\[
\qquad E_{\mathrm{int}}^{(\ell)},\ E_{\mathrm{int}}^{(\nu)}\ [\mathrm{BM.43}(2)];
\]

%dsp-level:0
\[
\qquad \text{the}\ \nu\text{-line index}\;
\]

%dsp-level:0
\[
\qquad \text{the relating equation}\;
\]

%dsp-level:0
\[
\qquad \text{the flavour composition}\;
\]

%dsp-level:0
\[
\qquad \text{the regime-scoped}\ \Sigma_{12}\ \text{fiber}\ [\mathrm{BM.68}(6)];
\]

%dsp-level:0
\[
(5)\quad\textbf{no tuning:}
\]

%dsp-level:0
\[
\qquad \neg\exists\ \text{an adjustment toward an observed ratio}\ —\ \text{the ordering is branch-conditional, the}\ Q\text{-value symbolic, and neither is compared to anything}
\]

%dsp-level:0
\[
\qquad [\mathrm{P5};\ \text{machine: zero external tokens, zero decimal numerals in this tier}];
\]

%dsp-level:0
\[
(6)\quad\textbf{status:}
\]

%dsp-level:0
\[
\qquad \textbf{derived:}\ (1)\text{--}(5);
\]

%dsp-level:0
\[
\qquad \textbf{the no-go is exact,}\ \neg\textbf{rhetorical: every bar in}\ (1)\ \text{carries a cited theorem or a named}\ \neg\text{construction scope};
\]

%dsp-level:0
\[
\neg\text{a P4 addition},\quad\neg\text{a new declaration},\quad\neg\text{a measured value}\ [\mathrm{P5}];\qquad\text{§Cal-BM stays empty}.
\]

\textbf{Req.} A6, D1.11, D4.2, L2 (Part 1); CT4.2, CT4.3, CT4.5, EM.25(7), EM.25(9) (Part 2); P5 (Parts 1–2); FT.3, FT.67(3), FT.68(4), FT.73(1), FT.79(5), FT.107(3) (Part 3); Census FT.35(4) (row-wise); Def-BM.12(4), Def-BM.12(5), Def-BM.12(7), BM.6(5), BM.10(4), BM.28(1), BM.28(4), BM.30(3), BM.30(5), BM.37(3), BM.41(2), BM.41(4), BM.43(2), BM.44(2), BM.45, BM.46(1), BM.46(4), BM.47(6), BM.48(1), BM.48(3), BM.48(4), BM.49(2), BM.50(3), BM.51(2), BM.51(3), BM.52(3), BM.53(4), BM.53(5), BM.53(6), BM.54(1), BM.54(2), BM.54(4), BM.57(4), BM.58(6), BM.59(4), BM.59(9), BM.62(1), BM.62(3), BM.62(6), BM.62(9), BM.62(11), BM.62(12), BM.65(2), BM.65(3), BM.65(5), BM.66(4), BM.66(6), BM.68(5), BM.68(6), BM.69(3), BM.69(4), BM.69(5), BM.70(3), BM.70(4) (this part); Cal-F; exact polynomial identities, Laurent reciprocals, exact rational censuses, sign censuses at positive symbols, token scans [math; machine: \texttt{checks/\allowbreak{}verify\_\allowbreak{}H\_\allowbreak{}organic\_\allowbreak{}bridge\_\allowbreak{}post\_\allowbreak{}bm14.py} §F, §G].

\textbf{Pf.}

%dsp-level:0
\[
(1)\ \text{row}\ 1:\ \mathrm{BM.53}(4)\text{--}(6)\ \text{re-cited};
\]

%dsp-level:0
\[
\qquad 2m_n(E_p-m_p)
\]

%dsp-level:0
\[
= -m_\ell^2\ \text{at}\ m_p
\]

%dsp-level:0
\[
= m_n\ \text{and}\ \lambda
\]

%dsp-level:0
\[
= m_\ell^2(m_\ell^2-4m_n^2)\ \text{re-derived as cleared polynomial identities in}\ \mathbb Q[m_n,m_\ell]\ [\text{machine}];
\]

%dsp-level:0
\[
\qquad \mathfrak T(i\,\mathbb1_d)=d\,i\ne0\ \text{at}\ d=2\ [\mathrm{BM.10}(4);\ \text{machine}];
\]

%dsp-level:0
\[
\qquad \Gamma_{\mathrm{tr}}\in[0,\tfrac13]\ [\mathrm{BM.44}(2)\ \text{re-cited}]
\]

%dsp-level:0
\[
\text{row}\ 2:\ \mathrm{Def\text{-}BM.12}(7)\text{'s}\ \neg\text{constructed list re-cited};
\]

%dsp-level:0
\[
\qquad E^{\mathrm{self}}_{\mathrm{EM}}(q)
\]

%dsp-level:0
\[
= c\,
\]

%dsp-level:0
\[
\varepsilon_{\mathrm{EM}}q^2m\ \text{with}\ q^2:\ 0
\]

%dsp-level:0
\[
\mapsto (1,1)\ \text{across the transition}:\ \ \text{the difference has positive coefficients in}\ (c,\varepsilon_{\mathrm{EM}},m)\ [\text{machine: the sign census}];
\]

%dsp-level:0
\[
\qquad \mathrm{BM.49}(2)\ \text{re-cited}
\]

%dsp-level:0
\[
\text{row}\ 3:\ f(k)=1-\tfrac1k:\quad f(2)=\tfrac12,\ f(3)=\tfrac23,\ f(4)=\tfrac34,\ \ \text{all}\ <1\ [\text{machine: exact Fractions}];
\]

%dsp-level:0
\[
\qquad 6-(2{+}2{+}2)=0\ [\text{machine}];
\]

%dsp-level:0
\[
\qquad \mathrm{BM.48}(4),\ \mathrm{BM.54}(2),(4)\ \text{re-cited}
\]

%dsp-level:0
\[
\text{row}\ 4:\ O(4)=2^5-16=16\ [\mathrm{CT0.6};\ \text{machine}];
\]

%dsp-level:0
\[
\qquad \text{the vertex-type row}\ n(q_1)
\]

%dsp-level:0
\[
= n(p_2)
\]

%dsp-level:0
\[
= 4\ \text{on both classes}\ [\text{Census FT.35}(4);\ \mathrm{BM.57}(4)\text{'s difference table};\ \text{machine: identically zero}];
\]

%dsp-level:0
\[
\qquad \mathrm{FT.73}(1)\text{'s inventory list re-cited}
\]

%dsp-level:0
\[
\text{row}\ 5:\ \mathrm{BM.49}(2),\ \mathrm{BM.41}(4)(\mathrm d),(\mathrm e),\ \mathrm{BM.30}(5)\ \text{re-cited};
\]

%dsp-level:0
\[
\qquad (2):\ \tfrac1{2\,\Delta m_{\mathrm{EM}}}=\tfrac1{2}\cdot\tfrac{2\pi\mathfrak h}{3ce^2\ln2}=\tfrac{\pi\mathfrak h}{3ce^2\ln2}\ [\text{machine: the Laurent reciprocal}]
\]

%dsp-level:0
\[
(3):\ \mathrm{BM.62}(1)\text{'s completeness re-cited with}\ \mathrm{BM.66}(6)\text{'s guard and}\ \mathrm{BM.68}(5)\text{'s slice};
\]

%dsp-level:0
\[
\qquad \Delta_g>0\ \text{on the three certificates}\ [\text{machine}];
\]

%dsp-level:0
\[
\qquad (4)\text{--}(6):\ \text{the verdict rows assembled from}\ \mathrm{BM.65}\text{--}\mathrm{BM.70}\text{'s cited clauses}\ [\text{machine: the row census and the token scans}]. \ \blacksquare
\]

\textbf{Phys.} Every door the audit could name is tried against the corpus's own theorems, and none opens. Uncertainty is a resolvability bound, not a budget: the closure is the non-existence of real on-shell solutions, a polynomial fact at every tick, there is no time operator to borrow from, the ledger runs forward only and contracts strictly, and the flicker that drives every transition weight lives at the amplitude layer, where amplitudes are barred from being masses. Field reconfiguration is a named unconstructed scope, and at the constructed order its sign runs against the decay, because the final state carries two charged members' self-energies where the initial carried none. Binding belongs to the bound, the bound channel does not exist, the nearby space holds only a strict fraction of a positive correction, and both binding reads the corpus owns are zero or class-degenerate. The dark sector has nothing to redistribute: its locks are class-shared, its junction content is identical on both rows, what it retains is inaccessible to every observer channel, and no response read exists by inventory theorem. What remains is the corpus's own honest door — a read specification, an act the tier prices, audits, and declines to perform. Under the bridge the neutron row does not move: the bridge term is zero in the currency in which the row is formed, conditionally positive in the pole currency where it is not formed, and the six verdict rows record the signs, the closure, and the derived/priced/open partition exactly — with nothing tuned toward anything observed.

\medskip\noindent\textit{\texorpdfstring{BM.72 — The status of the organic bridge: adjudication and scope}{BM.72 - The status of the organic bridge: adjudication and scope}.}\ Structural-constraints record: see Appendix B, Section B.1.24 of the Addendum.\par\medskip

\phantomsection
\section[Tier BM.16 — The Chiral Assignment: the Closed Deficit, the…]{\texorpdfstring{Tier BM.16 — The Chiral Assignment: the Closed Deficit, the Ray-Resolved Spin Census, the Attached Quark Splitting, the Chiral Vertex Dynamics, and the Specified Beta-Decay Opening}{Tier BM.16 - The Chiral Assignment: the Closed Deficit, the Ray-Resolved Spin Census, the Attached Quark Splitting, the Chiral Vertex Dynamics, and the Specified Beta-Decay Opening}}

\textbf{Locus.} This tier consumes the Tier BM.15 baseline unmodified and walks the two doors the corpus priced and declined: the deficit-$1'$ assignment rule — named missing at $\mathrm{FT.125}(4)$, $\mathrm{FT.141}(6)$, $\mathrm{FT.144}(5)$ and $\mathrm{FT.157}(4)$, with its forced shape already derived — and the $\beta$-vertex read specification — named as the sole opening route at $\mathrm{BM.30}(5)$, $\mathrm{BM.41}(4)(\mathrm e)$, $\mathrm{BM.49}(2)$ and $\mathrm{BM.71}(1)$ row $5$. Both acts are logged as §II.4 registry addenda (A-BM.7, A-BM.8) before use, per $\mathrm{Def\text{-}BM.3}(7)$'s rule. $\mathrm{BM.73}$ issues the left–right assignment matching $\hat\omega$ and $\hat\chi$ and derives the occupancy; $\mathrm{BM.74}$ resolves the spin census by ray; $\mathrm{BM.75}$ derives the weight–branch–core attachment and attaches the quark mass splitting; $\mathrm{BM.76}$ constructs the chiral vertex dynamics; $\mathrm{BM.77}$ issues the vertex read specification and proves the exact opening branch; $\mathrm{BM.78}$ adjudicates.

\textbf{Standing rows.} The refutations of $\mathrm{BM.57}(8)$ and $\mathrm{BM.58}(4)$ stand closed and are $\neg$reopened. $\mathrm{BM.70}(4)$'s dichotomy and $\mathrm{BM.71}(2)$'s no-go stand on the reads formed at their loci; $\mathrm{BM.77}(7)$ records their rescope under the specification, and no clause of $\mathrm{BM.41}$–$\mathrm{BM.72}$ is retracted. The relating equation between the informational and the geometric currencies is $\neg$assumed anywhere — item 7 stays open [$\mathrm{BM.35}(5)$; $\mathrm{BM.68}(7)$; pointer]. $\alpha_B$ stays barred on $\mathrm{BM.61}(8)$'s four grounds. $\neg$a Yukawa-type object exists or is added [$\mathrm{EW.15}$; $\mathrm{FT.145}(3)$]. No external mass, lifetime, endpoint or coupling value enters any derivation [$\mathrm{P5}$]. \texttt{}Chirality $=$ helicity on-shell'' remains external vocabulary [$\mathrm{FT.123}(3),(4)$]; the spin-structure torsor remains a distinct kind [$\mathrm{FT.122}(7)$; $\mathrm{FT.89}(2)$]; deficit $11$ is untouched [$\mathrm{FT.125}(5)$].

\begin{theorem}[{\texorpdfstring{BM.73 — The chiral matter assignment: the issued coupling, the occupancy theorem, and the $\omega$/$\chi$ matching rows}{BM.73 - The chiral matter assignment: the issued coupling, the occupancy theorem, and the / matching rows}}]\label{tag:BM-73}
\textbf{Thm.}

%dsp-level:0
\[
(1)\quad\textbf{the four derived gradings (the census the rule couples):}\qquad\begin{array}{l|l}\text{the algebra}\ S\text{-grading} & (10,\,6)=\bigl(\dim\mathcal K_1^{+},\ \dim\mathcal K_1^{-}\bigr)\ \text{on the force arena};\ \ \text{metric/weak reads}\ S\text{-even},\ \text{seal reads}\ S\text{-odd}\qquad[\mathrm{FT.136};\ \mathrm{FT.141}(1),(2),(5)]\\\text{the branch-label parity} & \text{spatial}=\text{aligned}\ (\alpha{=}0,\ \vee^2,\ \text{Casimir}\ 2)\ \text{weak-active};\ \ \text{temporal}=\text{antialigned}\ (\alpha{=}\pi,\ \wedge^2,\ \text{Casimir}\ 0)\ \text{weak-blind}\qquad[\mathrm{FT.101};\ \mathrm{FT.135};\ \mathrm{FT.157}(2),(3)]\\\text{the carrier}\ \hat\chi\text{-grade} & \hat\chi=i\bigl(|c\rangle\langle\omega(c)|-|\omega(c)\rangle\langle c|\bigr),\quad P_{L/R}=\tfrac12(\mathrm{id}\pm\hat\chi),\quad|L\rangle=\tfrac1{\sqrt2}\bigl(|c\rangle-i|\omega(c)\rangle\bigr)\qquad[\mathrm{Def\text{-}H.33};\ \mathrm{FT.122}(1)\text{--}(3),(6)]\\\text{the Janus-arrow chirality} & \text{strict alternation, the prefix law, the single-arrow materialized locus at rung one}\qquad[\mathrm{FT.144}(2),(5)]\end{array}
\]

%dsp-level:0
\[
\textbf{the pairwise distinctness is derived:}
\]

%dsp-level:0
\[
\qquad \text{the four are pairwise distinct}\ \mathbb Z_2\text{'s, commuting where jointly defined, neither reducible to another}
\]

%dsp-level:0
\[
\qquad [\mathrm{FT.136}(4)\text{--}(5);\ \mathrm{FT.141}(6);\ \mathrm{FT.126}(4);\ \mathrm{FT.157}(4)];
\]

%dsp-level:0
\[
(2)\quad\textbf{the assignment rule (the declaration, logged as A-BM.7 before use):}
\]

%dsp-level:0
\[
\qquad \textbf{the forced shape is derived:}\ \ \text{any}\ \hat\omega\text{-covariant arrow}\leftrightarrow\text{ray coupling maps the single-arrow locus into a single chiral ray, and the complementary ray has occupancy zero at rung one}\ —\ \text{exhaustive on the}\ 2!\ \text{bijections}
\]

%dsp-level:0
\[
\qquad [\mathrm{FT.144}(5)]:
\]

%dsp-level:0
\[
\boxed{\ \textbf{the issued rule:}\qquad\text{the materialization-arrow locus occupies}\ \operatorname{im}P_L\ \ \text{— the}\ \hat\chi=+1\ \text{ray of}\ \mathrm{Def\text{-}H.33}\text{'s boxed sign}\ }
\]

%dsp-level:0
\[
\textbf{the coupled rows (each forced given the choice, or the choice itself):}
\]

%dsp-level:0
\[
\qquad (\mathrm a)\ \ \text{the occupied ray carries the weak-active/aligned/spatial branch label};
\]

%dsp-level:0
\[
\ \text{the complement carries the weak-blind/antialigned/temporal label}
\]

%dsp-level:0
\[
\qquad [\text{forced: every materialized bound chain is weak-active}\ —\ \{G,H,W\}\subseteq\mathcal A(\Sigma_{\min})\ [\text{Census FT.28}(3)];\ \text{weak acts from rung one}\ [\mathrm{CT0.13}];\ \text{the rung-one channel class is weak-carried}\ [\mathrm{FT.142}(1)]:\ \text{the alternative coupling would make every matter line weak-blind against those derived rows};\ \text{machine: the}\ 2\ \text{couplings censused, one consistent}];
\]

%dsp-level:0
\[
(\mathrm b)\ \ \text{the matter orientation occupies the seed weak doublet on the left ray}\ [\mathrm{EW.32};\ \mathrm{EW.36}:\ \operatorname{spec}Q|_{\mathcal H_{\mathsf v}}=\{0,+1\}];
\]

%dsp-level:0
\[
\qquad \text{the image orientation occupies the}\ \omega\text{-image representation}\ \{0,-1\}\ \text{on the right ray}\ [\mathrm{EW.36};\ \mathrm{D2.6};\ \mathrm{BM.25};\ \mathrm{BM.27};\ \text{the A-BM.5 antineutrino row}]
\]

%dsp-level:0
\[
\qquad [\text{forced by}\ \hat\omega\text{-covariance once}\ (\mathrm a)\ \text{holds}:\ \hat\omega P_L\hat\omega^{-1}=P_R\ [\mathrm{FT.122}(4)]];
\]

%dsp-level:0
\[
(\mathrm c)\ \ \text{the}\ S\text{-grading keeps its derived division of labour}\ —\ \text{it grades observables,}\ \neg\text{states}\ —\ \text{and receives no species coupling}
\]

%dsp-level:0
\[
\qquad [\mathrm{FT.141}(6)]:
\]

%dsp-level:0
\[
\textbf{the declared content is exactly one naming bit}\ —\ \text{which ray the arrow occupies is named left}\ —\ \text{on the forced shape of}\ [\mathrm{FT.144}(5)];
\]

%dsp-level:0
\[
\qquad \textbf{it carries no datum:}\ \neg\text{a scale},
\]

%dsp-level:0
\[
\ \neg\text{a coupling value},
\]

%dsp-level:0
\[
\ \neg\text{a measured quantity}\ [\mathrm{P5};\ \mathrm{Def\text{-}BM.3}(7)\ \text{pattern}];
\]

%dsp-level:0
\[
(3)\quad\textbf{the occupancy theorem (derived given}\ (2)):
\]

%dsp-level:0
\[
\qquad \boxed{\ \text{the rung-one materialized locus occupies}\ (1,\,0)\ =\ (\text{left occupied},\ \text{right empty})\ }
\]

%dsp-level:0
\[
\text{the complement's content is the vacuum sector: all-virtual, unregistered, no accessible pole}
\]

%dsp-level:0
\[
\qquad [\mathrm{FT.144}(4),(5);\ \mathrm{CT4.2}\ \text{pattern};\ \mathrm{CT0.1}];
\]

%dsp-level:0
\[
\textbf{the rung ladder of the occupancy:}
\]

%dsp-level:0
\[
\qquad \text{rung one: the sectors are force-decoupled},\ F_{LR}=F_{RL}=0\ \text{exactly}\ [\mathrm{FT.128}(2);\ \mathrm{FT.144}(6)];
\]

%dsp-level:0
\[
\qquad \text{rung two: the EM core read opens the unique flip channel},
\]

%dsp-level:0
\[
\ F_{LR}
\]

%dsp-level:0
\[
= q\,
\]

%dsp-level:0
\[
I_2
\]

%dsp-level:0
\[
\ne 0
\]

%dsp-level:0
\[
\iff q
\]

%dsp-level:0
\[
\ne 0,
\]

%dsp-level:0
\[
\ \text{splitting}\ 2|q|\in\{0,2\}\ [\mathrm{FT.129}(5);\ \mathrm{FT.130}(5);\ \mathrm{Def\text{-}H.43}\ \text{promotion inherited};\ \mathrm{BM.22}(2),(4)];
\]

%dsp-level:0
\[
\text{rung three: colour adds no carrier flip}\ —\ \text{sealed ends are blind},
\]

%dsp-level:0
\[
\ \text{the seed is colour-neutral}\ [\mathrm{FT.117};\ \mathrm{EW.32};\ \mathrm{FT.144}(6)];
\]

%dsp-level:0
\[
\qquad \textbf{the flipping set remains}\ \{F_{\mathrm{EM}}\ \text{core read}\},\ \text{minimal and unique}\ [\mathrm{FT.129}(5)];
\]

%dsp-level:0
\[
(4)\quad\textbf{the}\ \omega/\chi\ \textbf{matching rows (re-certified exact over}\ \mathbb Q(i)):
\]

%dsp-level:0
\[
\qquad \hat\chi^\dagger=\hat\chi,\quad\hat\chi^2=\mathrm{id},\quad\operatorname{tr}\hat\chi=0,\quad\operatorname{spec}\hat\chi=\{+1,-1\};
\]

%dsp-level:0
\[
\qquad P_L^2=P_L,\quad P_LP_R=0,\quad P_L+P_R=\mathrm{id},\quad\operatorname{tr}P_L=\operatorname{tr}P_R=1
\]

%dsp-level:0
\[
\qquad [\mathrm{FT.122}(1),(2);\ \text{machine}]:
\]

%dsp-level:0
\[
\boxed{\ \hat\omega\hat\chi\hat\omega^{-1}=-\hat\chi,\qquad\{\hat\chi,\hat\omega\}=0,\qquad\hat\omega P_L\hat\omega^{-1}=P_R,\qquad\hat\omega|L\rangle=-i|R\rangle,\quad\hat\omega|R\rangle=+i|L\rangle\ }
\]

%dsp-level:0
\[
\qquad [\mathrm{FT.122}(4);\ \text{machine}];
\]

%dsp-level:0
\[
\textbf{the quaternionic triple:}
\]

%dsp-level:0
\[
\qquad \hat\omega^2=\hat\chi^2=\hat\kappa^2=\mathrm{id};
\]

%dsp-level:0
\[
\qquad \{\hat\omega,\hat\chi\}=\{\hat\chi,\hat\kappa\}=\{\hat\kappa,\hat\omega\}=0;
\]

%dsp-level:0
\[
\qquad \hat\omega\hat\chi=-i\hat\kappa,\quad\hat\chi\hat\kappa=-i\hat\omega,\quad\hat\kappa\hat\omega=-i\hat\chi
\]

%dsp-level:0
\[
\qquad [\mathrm{Def\text{-}H.34};\ \text{machine}];
\]

%dsp-level:0
\[
\textbf{the winding row:}
\]

%dsp-level:0
\[
\qquad \omega:\ k\mapsto-k
\]

%dsp-level:0
\[
\qquad [\mathrm{FT.122}(4);\ \mathrm{D2.6};\ \mathrm{EM.24}(1)]:
\]

%dsp-level:0
\[
\qquad \text{the ray exchange and the winding reversal are the two faces of one involution};
\]

%dsp-level:0
\[
\textbf{the canonicity:}
\]

%dsp-level:0
\[
\qquad \bigl\{M=M^\dagger:\ \hat\omega M\hat\omega=-M,\ M^2=\mathrm{id},\ \langle c|M|c\rangle=\langle\omega(c)|M|\omega(c)\rangle=0\bigr\}=\{+\hat\chi,\,-\hat\chi\}
\]

%dsp-level:0
\[
\qquad [\mathrm{FT.122}(6);\ \text{machine: the circle}\ M(a,y),\ a^2+y^2=1,\ \text{on exact rational witnesses}];
\]

%dsp-level:0
\[
(5)\quad\textbf{the weak rows under the assignment:}
\]

%dsp-level:0
\[
\qquad \bigl[W_i\otimes I_\omega,\ I\otimes P_{L/R}\bigr]=0
\]

%dsp-level:0
\[
\qquad [\mathrm{FT.123}(1);\ \mathrm{FT.126}(2);\ \text{machine}]:
\]

%dsp-level:0
\[
\qquad \textbf{the assignment attaches species to rays; it does not modify the weak action};
\]

%dsp-level:0
\[
\text{the selectivity lives at the state locus:}\quad W_i\Pi_{\mathcal V}=0,
\]

%dsp-level:0
\[
\qquad \textstyle\sum_iW_i^2=2\,\Pi_{\mathcal R}
\]

%dsp-level:0
\[
\qquad [\mathrm{FT.157}(2);\ \mathrm{FT.141}(4);\ \mathrm{FT.101}(1)\text{--}(2);\ \text{machine re-cited}];
\]

%dsp-level:0
\[
\textbf{the occupied ray is the weak-active one by}\ (2)(\mathrm a);
\]

%dsp-level:0
\[
\qquad \text{a chiral gauge coupling as an operator modification is}\ \neg\text{constructed and}\ \neg\text{needed for the assignment}
\]

%dsp-level:0
\[
\qquad [\mathrm{FT.123}\ \mathrm{Rmk}\ \text{scope}];
\]

%dsp-level:0
\[
(6)\quad\textbf{the bars:}
\]

%dsp-level:0
\[
\qquad \text{the weight}\leftrightarrow\text{ray attachment is barred as the external helicity vocabulary}
\]

%dsp-level:0
\[
\qquad [\mathrm{FT.123}(3),(4)];
\]

%dsp-level:0
\[
\qquad \text{the spin-structure torsor is reversal-invariant, a distinct kind}
\]

%dsp-level:0
\[
\qquad [\mathrm{FT.122}(7);\ \mathrm{FT.89}(2)];
\]

%dsp-level:0
\[
\text{deficit}\ 11\ \text{(spin--statistics) is untouched: no exchange operator exists}
\]

%dsp-level:0
\[
\qquad [\mathrm{FT.80}(5);\ \mathrm{FT.125}(5)];
\]

%dsp-level:0
\[
\qquad \text{the frame section}\ e\mapsto\mathcal F_e\ \text{and a global Dirac-type operator remain unconstructed}
\]

%dsp-level:0
\[
\qquad [\mathrm{FT.125}(1);\ \text{item}\ 21(\mathrm i)\ \text{pointer}];
\]

%dsp-level:0
\[
\textbf{no SM species reading and no claim that the SM weak interaction is reproduced}
\]

%dsp-level:0
\[
\qquad [\mathrm{FT.157}\ \mathrm{Rmk}\ \text{bars stand};\ \mathrm{P1},\ \mathrm{P5}];
\]

%dsp-level:0
\[
(7)\quad\textbf{status:}
\]

%dsp-level:0
\[
\qquad \textbf{derived:}\ (1),\ (3)\text{--}(6)\ \ [(3)\ \text{given}\ (2)];
\]

%dsp-level:0
\[
\qquad \textbf{declared: exactly one naming bit}\ (2)\ —\ \text{logged as the §II.4 addendum A-BM.7 before use};
\]

%dsp-level:0
\[
\textbf{deficit}\ 1'\ \textbf{is closed at the assignment level:}
\]

%dsp-level:0
\[
\qquad \text{the coupling rule exists}\ —\ \text{the declaration's root is issued};
\]

%dsp-level:0
\[
\qquad \text{the downstream rows unblock at their own loci and stay open:}\ \text{deficits}\ 2,
\]

%dsp-level:0
\[
\ 3,
\]

%dsp-level:0
\[
\ 4\ [\mathrm{FT.125}(6)],
\]

%dsp-level:0
\[
\ \ \mathrm{Q4}\text{--}\mathrm{Q6}\ \text{residues},
\]

%dsp-level:0
\[
\ \ \text{the anomaly sums formable on the}\ 1'\ \text{side with the}\ +\tfrac23\ \text{side structurally unreachable}\ [\mathrm{FT.125}(2);\ \mathrm{BM.28}(7)];
\]

%dsp-level:0
\[
\neg\text{a P4 addition},\quad\neg\text{a new dimensionful scale}\ [\mathrm{P1}],\quad\neg\text{a measured value}\ [\mathrm{P5}],\quad\text{§Cal-BM stays empty};
\]

%dsp-level:0
\[
\qquad \text{the extra-axiomatic list}\ [\text{A §0.2a}]\ \text{gains no member}.
\]

\textbf{Req.} A4, A8 (Part 1); D2.6, CT0.1, CT0.13 (Part 1); CT4.2, EM.24(1), EW.32, EW.36 (Part 2); P1, P5 (Parts 1–2); Def-H.33, Def-H.34, Def-H.43, FT.80(5), FT.89(2), FT.101(1), FT.101(2), FT.101(3), FT.117, FT.122(1), FT.122(2), FT.122(3), FT.122(4), FT.122(6), FT.122(7), FT.123(1), FT.123(3), FT.123(4), FT.125(1), FT.125(2), FT.125(4), FT.125(5), FT.125(6), FT.126(2), FT.126(4), FT.128(2), FT.129(5), FT.130(5), FT.135, FT.136(1), FT.136(4), FT.136(5), FT.141(1), FT.141(2), FT.141(4), FT.141(5), FT.141(6), FT.142(1), FT.144(2), FT.144(4), FT.144(5), FT.144(6), FT.157(2), FT.157(3), FT.157(4) (Part 3); Census FT.28(3) (row-wise); Def-BM.3(7), BM.22(2), BM.22(4), BM.25, BM.27, BM.28(7) (this part); exact $\mathbb Q(i)$ matrix algebra, exhaustive bijection censuses, rational-circle witnesses [math; machine: \texttt{checks/\allowbreak{}verify\_\allowbreak{}H\_\allowbreak{}chiral\_\allowbreak{}assignment\_\allowbreak{}opening.py} §A, §B].

\textbf{Pf.}

%dsp-level:0
\[
(1):\ \text{the four rows re-cited from their own tags}\ [\mathrm{FT.136}(1),\ \mathrm{FT.141}(1),(2),(5);\ \mathrm{FT.157}(2),(3);\ \mathrm{Def\text{-}H.33};\ \mathrm{FT.122}(1)\text{--}(3);\ \mathrm{FT.144}(2),(5)];
\]

%dsp-level:0
\[
\qquad \text{the distinctness rows re-cited}\ [\mathrm{FT.126}(4);\ \mathrm{FT.136}(4)\text{--}(5);\ \mathrm{FT.141}(6);\ \mathrm{FT.157}(4)]
\]

%dsp-level:0
\[
(2):\ \mathrm{FT.144}(5)\text{'s forced shape re-cited}\ [\text{machine: exhaustive on the}\ 2!\ \text{bijections}];
\]

%dsp-level:0
\[
\qquad (\mathrm a):\ \text{the weak-activity census: both ray/label couplings evaluated against}\ \{G,H,W\}\subseteq\mathcal A(\Sigma_{\min})\ [\text{Census FT.28}(3)],
\]

%dsp-level:0
\[
\ \mathrm{CT0.13},
\]

%dsp-level:0
\[
\ \mathrm{FT.142}(1):\ \text{the weak-blind-occupied coupling contradicts the derived weak activity of every matter floor};
\]

%dsp-level:0
\[
\ \text{one coupling survives}\ [\text{machine}];
\]

%dsp-level:0
\[
\qquad (\mathrm b):\ \hat\omega P_L\hat\omega^{-1}=P_R\ [\text{machine over}\ \mathbb Q(i)]\ \text{and}\ \omega:\operatorname{spec}Q\mapsto\{0,-1\}\ [\mathrm{EW.36}\ \text{re-cited}];
\]

%dsp-level:0
\[
\qquad (\mathrm c):\ \mathrm{FT.141}(6)\ \text{verbatim}
\]

%dsp-level:0
\[
(3):\ \mathrm{FT.144}(5)\text{'s occupancy pair}\ (1,0)\ \text{with}\ (2)\text{'s naming};
\]

%dsp-level:0
\[
\qquad \text{the rung rows re-cited}\ [\mathrm{FT.128}(2);\ \mathrm{FT.129}(5);\ \mathrm{FT.130}(5);\ \mathrm{FT.117};\ \mathrm{FT.144}(6)]
\]

%dsp-level:0
\[
(4):\ \hat\chi=i\left(\begin{smallmatrix}0&1\\-1&0\end{smallmatrix}\right),\ \hat\omega=\left(\begin{smallmatrix}0&1\\1&0\end{smallmatrix}\right),\ \hat\kappa=\left(\begin{smallmatrix}1&0\\0&-1\end{smallmatrix}\right)\ \text{in the ordered orientation basis}\ [\mathrm{FT.122}(5);\ \mathrm{Def\text{-}H.34}]:\ \text{every identity of}\ (4)\ \text{verified entrywise over}\ \mathbb Q(i)\ [\text{machine}];\qquad\hat\omega|L_u\rangle=|R_u\rangle\cdot(-i)\ \text{on the unnormalized kets}\ [\text{machine}];\qquad\text{the canonicity circle}\ M(a,y)^2=(a^2{+}y^2)\mathrm{id}\ \text{and the diagonal-free reduction}\ a=0\implies y^2=1\implies M=\pm\hat\chi\ [\text{machine: on the rational witnesses}\ (\tfrac35,\tfrac45),\ (\tfrac5{13},\tfrac{12}{13}),\ (\tfrac8{17},\tfrac{15}{17})]
\]

%dsp-level:0
\[
(5):\ \text{the commutator}\ [W_i\otimes I,\ I\otimes P]=0\ \text{by tensor-factor disjointness}\ [\text{machine}];
\]

%dsp-level:0
\[
\qquad \text{the locus rows re-cited}\ [\mathrm{FT.157}(2);\ \mathrm{FT.141}(4)]
\]

%dsp-level:0
\[
(6)\text{--}(7):\ \text{the bars re-cited clause by clause};\qquad\text{the standing-rule scans}\ [\text{machine}]. \ \blacksquare
\]

\textbf{Phys.} The corpus had derived everything about chirality except one bit: which of the two exchanged rays the materialized world sits in. The four gradings — the force arena's ten/six split, the weak-active and weak-blind branch labels, the orientation doublet's exchange-form grade, and the Janus arrow's single-direction locus — are pairwise distinct, each exact, and the rule marrying them to matter was the ledger's named missing declaration. This tier issues it in its minimal form: the arrow's occupied ray is the left ray, the one the boxed chirality sign already fixes algebraically; the weak-active branch label then rides along by consistency, because every materialized matter floor is weak-active by derived rows and the opposite coupling would contradict them; the right ray is the orientation reversal's image and hosts the image-orientation species the antineutrino naming row already used. What is derived given the bit is substantial: the occupancy pair one-zero with its vacuum complement, the rung ladder along which the flip channel opens at rung two for charged cores and never for colour, and the whole $\omega$/χ matching — the involution exchanges the rays, anticommutes with the grade, and reverses the winding, so ray exchange and matter–image exchange are one move. What is not obtained is stated as loudly: no helicity claim, no spin-statistics, no Standard Model reading, and no modification of the weak action itself, which stays chirality-blind on the carrier while the species sit where the rule puts them.

\end{theorem}

\pseudohead{theorem}{\texorpdfstring{BM.74}{BM.74}}{\texorpdfstring{The ray-resolved spin census: the winding labels, the occupancy rows, and the vertex table}{The ray-resolved spin census: the winding labels, the occupancy rows, and the vertex table}}\label{tag:BM-74}
\textbf{Thm.}

%dsp-level:0
\[
(1)\quad\textbf{the spin-label rule (the declared interface inherited,}\ \neg\textbf{re-derived):}
\]

%dsp-level:0
\[
\qquad \text{a materialized chain of winding}\ k\ \text{carries}\ \mathrm{Sym}^{k}(\mathcal H_{\mathsf v}),
\]

%dsp-level:0
\[
\qquad j=\frac{|k|}2,
\]

%dsp-level:0
\[
\qquad \text{weights}\ j,\ j{-}1,\ \dots,\ -j
\]

%dsp-level:0
\[
\qquad [\mathrm{Def\text{-}H.38};\ \mathrm{FT.55}(2);\ \mathrm{FT.60}(1)\text{--}(3);\ \mathrm{FT.123}(3)]:
\]

%dsp-level:0
\[
\textbf{the anchor:}
\]

%dsp-level:0
\[
\qquad k=1:\quad\mathrm{Sym}^{1}=\mathcal H_{\mathsf v}\cong\mathbb C^2,
\]

%dsp-level:0
\[
\qquad j=\tfrac12,
\]

%dsp-level:0
\[
\qquad \text{weights}\ \bigl\{+\tfrac12,\,-\tfrac12\bigr\}
\]

%dsp-level:0
\[
\qquad [\mathrm{EW.32};\ \text{Census FT.28}(2);\ \text{Census FT.63}(1);\ \text{machine}];
\]

%dsp-level:0
\[
\textbf{the Casimir ladder:}
\]

%dsp-level:0
\[
\qquad \frac k2\Bigl(\frac k2+1\Bigr)
\]

%dsp-level:0
\[
\qquad [\mathrm{FT.60}(3)];
\]

%dsp-level:0
\[
\qquad \text{the multiplicity--winding coupling itself stays a declared interface}\ [\mathrm{Def\text{-}H.38}]\ —\ \neg\text{derived here},
\]

%dsp-level:0
\[
\ \neg\text{reopened};
\]

%dsp-level:0
\[
(2)\quad\textbf{the ray-resolved line census (derived given A-BM.7):}
\]

%dsp-level:0
\[
\qquad \text{every materialized matter line occupies the left ray: occupancy}\ (1,0)
\]

%dsp-level:0
\[
\qquad [\mathrm{BM.73}(3)]:
\]

%dsp-level:0
\[
\text{the image-orientation lines}\ —\ \text{the CPT partners and the A-BM.5 antineutrino row}\ —\ \text{occupy the right ray}
\]

%dsp-level:0
\[
\qquad [\mathrm{BM.73}(2)(\mathrm b);\ \mathrm{BM.25};\ \mathrm{BM.27};\ \text{machine: the}\ \hat\omega\text{-exchange row}];
\]

%dsp-level:0
\[
\textbf{the cardinality correspondence stays a cardinality row:}
\]

%dsp-level:0
\[
\qquad \text{the two weights}\ \bigl\{\pm\tfrac12\bigr\}\ \text{and the two rays}\ \{+1,-1\}\ \text{are an unordered-pair match};
\]

%dsp-level:0
\[
\ \ \text{the weight}\leftrightarrow\text{ray attachment is barred}
\]

%dsp-level:0
\[
\qquad [\mathrm{FT.123}(3);\ \text{the external helicity vocabulary}];
\]

%dsp-level:0
\[
(3)\quad\textbf{the composite rows:}
\]

%dsp-level:0
\[
\qquad \text{the baryon card:}\quad\Theta_{\mathrm{int}}=2\,k(\Gamma)\in2\mathbb Z,
\]

%dsp-level:0
\[
\qquad j(\Gamma)=|k(\Gamma)|\ \ \text{rung-relative}\ —\ \text{the card fixes no rung}
\]

%dsp-level:0
\[
\qquad [\text{Census FT.35}(3);\ \text{Census FT.65}(4):\ \text{the free}\ \mathbb Z^3\ \text{ladder}];
\]

%dsp-level:0
\[
\text{the Platonic floors at the uniform minima: every face}\ j_f\ \text{composes to an integer stretched total},
\]

%dsp-level:0
\[
\quad\sum_fk_{f,\min}/2
\]

%dsp-level:0
\[
= K_{S,\min}/2\in\{6,\,6,\,12,\,30,\,30\}
\]

%dsp-level:0
\[
\qquad [\text{Census FT.63}(1)(\mathrm b);\ \text{machine re-cited}];
\]

%dsp-level:0
\[
\text{the defect multiplet ladder:}
\]

%dsp-level:0
\[
\qquad m^2_{\mathrm{pole}}(m)=4\,\hat\mu'^2\sin^2\!\bigl(m\,\tfrac{\Theta_R}4\bigr),
\]

%dsp-level:0
\[
\qquad m\in\{k,\,k{-}2,\ \dots,\,-k\}
\]

%dsp-level:0
\[
\qquad [\mathrm{FT.60}(5);\ \text{Census FT.63}(6)]:
\]

%dsp-level:0
\[
\qquad \text{the six states of}\ \mathrm{Def\text{-}BM.14}\ \text{ride it with}\ j=|k|/2;
\]

%dsp-level:0
\[
\qquad \text{the mesonic}\ \delta\text{-sealed rows stand untouched}\ [\mathrm{BM.31}\text{--}\mathrm{BM.33}];
\]

%dsp-level:0
\[
(4)\quad\textbf{the vertex table (ray-resolved; derived given A-BM.7 and the frozen grammar):}
\]

%tbl-level:3
{\scriptsize\setlength{\tabcolsep}{4pt}
\begin{longtable}{>{\raggedright\arraybackslash}p{\dimexpr0.0881\textwidth-2\tabcolsep\relax}|>{\raggedright\arraybackslash}p{\dimexpr0.2795\textwidth-2\tabcolsep\relax}|>{\raggedright\arraybackslash}p{\dimexpr0.1999\textwidth-2\tabcolsep\relax}|>{\raggedright\arraybackslash}p{\dimexpr0.4326\textwidth-2\tabcolsep\relax}}
\toprule
member & the structural row & the spin/rung row & the ray row \\
\addlinespace[2pt]
\midrule
\endfirsthead
\toprule
member & the structural row & the spin/rung row & the ray row \\
\addlinespace[2pt]
\midrule
\endhead
\bottomrule
\endlastfoot
the parent & the neutral-class $\varepsilon$-sealed 3-juncture, $Q=\allowbreak 0$ $[\mathrm{BM.17};\ \mathrm{Def\text{-}BM.3}(2)]$ & $j=\allowbreak |k(\Gamma)|$, rung-relative $[\text{Census FT.35}(3)]$ & the constituents occupy the left ray; the neutral-core constituents are ray-locked: $F_{LR}=\allowbreak 0$, parity doubling $[\mathrm{BM.22}(4);\ \mathrm{FT.130}(6)]$ \\
the daughter & the unit-class $\varepsilon$-sealed 3-juncture, $Q=\allowbreak +\allowbreak 1$ $[\mathrm{BM.26};\ \mathrm{Def\text{-}BM.3}(3)]$ & the same ladder & the charged-core constituents are flip-coupled, splitting $2|q|=\allowbreak 2$ $[\mathrm{FT.130}(5);\ \mathrm{FT.144}(6)]$ \\
the charged emitted member & the $\ell$-channel line, $q=\allowbreak -1$: forced, class-unique $[\mathrm{BM.23}(5);\ \mathrm{BM.41}(1);\ \mathrm{Def\text{-}BM.6}(1)]$ & the $k=\allowbreak 1$ anchor: $j=\allowbreak \tfrac12$ $[\text{Census FT.28}(2);\ \mathrm{Def\text{-}H.38}]$ & flip-coupled ($q\ne\allowbreak 0$); the arrow orientation occupies the left ray with the partner accessible $[\mathrm{BM.73}(3);\ \mathrm{FT.129}(5)]$ \\
the neutral emitted member & the $\nu$-channel line in the image orientation, $q=\allowbreak 0$ $[\mathrm{BM.41}(3);\ \mathrm{FT.142}(1);\ \text{A-BM.5}]$ & the $k=\allowbreak 1$ anchor: $j=\allowbreak \tfrac12$ & the right ray (the image orientation under A-BM.7); flip-blind, parity doubling $[\mathrm{BM.22}(4);\ \mathrm{FT.130}(6)]$ \\
\end{longtable}
}

%dsp-level:0
\[
\textbf{the two emitted members occupy opposite rays}
\]

%dsp-level:0
\[
\qquad [\text{machine: the table's ray column is a derived consequence of}\ \mathrm{BM.73}(2)(\mathrm b),\ (3)\ \text{and the frozen grammar}];
\]

%dsp-level:0
\[
(5)\quad\textbf{the angular-momentum budget is unchanged:}
\]

%dsp-level:0
\[
\qquad \mathrm{BM.42}\text{'s rows stand}\ —\ \text{the Noether moment conserved on-shell},
\]

%dsp-level:0
\[
\ \text{the Lagrange decomposition}\ p^2
\]

%dsp-level:0
\[
= p_r^2+L^2/r^2,
\]

%dsp-level:0
\[
\ \text{the absorption adjudications},
\]

%dsp-level:0
\[
\ \text{the}\ \ell\text{-selection census}
\]

%dsp-level:0
\[
\qquad [\mathrm{BM.42}(1)\text{--}(4);\ \mathrm{BM.40}(3):\ \text{the}\ \sigma\text{-vertex is Casimir-neutral}]:
\]

%dsp-level:0
\[
\textbf{the ray label adds no torque term:}
\]

%dsp-level:0
\[
\qquad \hat\chi\ \text{acts on the orientation factor alone and commutes with the rotation generators acting on the seed}
\]

%dsp-level:0
\[
\qquad [\mathrm{FT.123}(1)\ \text{pattern};\ \mathrm{FT.126}(4):\ \text{independent}\ \mathbb Z_2\text{'s};\ \text{machine: the commutant row}];
\]

%dsp-level:0
\[
(6)\quad\textbf{what is not obtained:}
\]

%dsp-level:0
\[
\qquad \neg\text{a spin--statistics theorem}\ —\ \text{no exchange operator exists}\ [\mathrm{FT.80}(5);\ \mathrm{FT.125}(5)];
\]

%dsp-level:0
\[
\qquad \neg\text{a fixed}\ j\ \text{for the baryon card}\ —\ \text{the rung datum is not fixed here}\ [\mathrm{BM.68}(6)\ \text{regime row};\ \text{Census FT.65}(4)];
\]

%dsp-level:0
\[
\neg\text{a weight}\leftrightarrow\text{ray helicity row}\ [\mathrm{FT.123}(3),(4)];
\]

%dsp-level:0
\[
\qquad \neg\text{a}\ k\ge3\ \text{carrier-inventory extension}\ [\mathrm{FT.60}(8)];
\]

%dsp-level:0
\[
(7)\quad\textbf{status:}
\]

%dsp-level:0
\[
\qquad \textbf{derived given A-BM.7:}\ (2)\text{--}(5);
\]

%dsp-level:0
\[
\qquad \textbf{inherited declarations:}\ \mathrm{Def\text{-}H.38}\text{'s multiplicity coupling}\ \ \text{and the naming rows}\ [\mathrm{Def\text{-}BM.3};\ \mathrm{Def\text{-}BM.6};\ \text{A-BM.5}];
\]

%dsp-level:0
\[
\neg\text{a P4 addition},\quad\neg\text{a measured value}\ [\mathrm{P5}],\quad\neg\text{a new dimensionful scale}\ [\mathrm{P1}];\qquad\text{§Cal-BM stays empty}.
\]

\textbf{Req.} D1.4 (Parts 1–2); EW.32 (Part 2); P1, P5 (Parts 1–2); Def-H.38, FT.55(2), FT.60(1), FT.60(3), FT.60(5), FT.60(8), FT.80(5), FT.122(4), FT.123(1), FT.123(3), FT.123(4), FT.125(5), FT.126(4), FT.129(5), FT.130(5), FT.130(6), FT.142(1) (Part 3); Census FT.28(2), FT.35(3), FT.63(1), FT.63(6), FT.65(4) (row-wise); Def-BM.3(2), Def-BM.3(3), Def-BM.6(1), BM.17, BM.22(4), BM.23(5), BM.25, BM.26, BM.27, BM.31, BM.32, BM.33, BM.40(3), BM.41(1), BM.41(3), BM.42(1), BM.42(2), BM.42(3), BM.42(4), BM.68(6), BM.73(2), BM.73(3) (this part); exact $\mathbb Q(i)$ commutant rows, integer ladder arithmetic [math; machine: \texttt{checks/\allowbreak{}verify\_\allowbreak{}H\_\allowbreak{}chiral\_\allowbreak{}assignment\_\allowbreak{}opening.py} §B, §C].

\textbf{Pf.}

%dsp-level:0
\[
(1):\ \mathrm{Def\text{-}H.38}/\mathrm{FT.60}(1)\text{--}(3)\ \text{re-cited}\ (\text{frozen});
\]

%dsp-level:0
\[
\qquad \dim\mathrm{Sym}^1=2,\ \ j=\tfrac12,\ \ \text{weights}\ \pm\tfrac12\ [\text{machine: the}\ \tfrac{\sigma_3}2\ \text{spectrum}];
\]

%dsp-level:0
\[
\qquad \text{the Casimir ladder re-cited}\ [\mathrm{FT.60}(3)]
\]

%dsp-level:0
\[
(2):\ \mathrm{BM.73}(3)\text{'s occupancy with}\ (2)(\mathrm b)\text{'s image row};
\]

%dsp-level:0
\[
\qquad \hat\omega\ \text{maps the occupied ray onto the complement}\ [\mathrm{FT.122}(4);\ \text{machine}]
\]

%dsp-level:0
\[
(3):\ \text{Census FT.35}(3),\ \text{FT.63}(1)(\mathrm b),\ \text{FT.65}(4),\ \mathrm{FT.60}(5)\ \text{re-cited row-wise};
\]

%dsp-level:0
\[
\qquad \Theta_{\mathrm{int}}=2k\implies j=|k|\ [\text{machine: integers}]
\]

%dsp-level:0
\[
(4):\ \text{each row one frozen citation plus}\ \mathrm{BM.73}(2)\text{--}(3);
\]

%dsp-level:0
\[
\qquad \text{the opposite-ray conclusion: the charged member's arrow orientation}\ \mapsto\ \text{the left ray};
\]

%dsp-level:0
\[
\ \text{the neutral member's image orientation}\ \mapsto\ \text{the right ray}\ [\text{machine: the}\ 2\ \text{row census}]
\]

%dsp-level:0
\[
(5):\ \mathrm{BM.42}(1)\text{--}(4),\ \mathrm{BM.40}(3)\ \text{re-cited}\ (\text{frozen});
\]

%dsp-level:0
\[
\qquad [\sigma_i\otimes I_\omega,\ I\otimes\hat\chi]=0\ [\text{machine: tensor disjointness}]
\]

%dsp-level:0
\[
(6)\text{--}(7):\ \text{the bars re-cited};\qquad\text{the standing-rule scans}\ [\text{machine}]. \ \blacksquare
\]

\textbf{Phys.} Spin in this corpus was always a winding label — j equals half the winding, with the seed qubit as the k-equals-one anchor — and chirality was always an orientation grade; the assignment now says which ray the matter sits in, and the census resolves by ray without moving a single spin value. Every materialized line occupies the left ray; every image-orientation partner, including the antineutrino row the decay vertex already named, occupies the right; and the two emitted leptons of the beta vertex therefore sit in opposite rays, one flip-coupled because its core is charged, one flip-blind and parity-doubled because its core is neutral. The angular-momentum budget absorbs the new label silently, because the chirality grade acts on the orientation factor and commutes with every rotation generator on the seed. What the tier refuses is as exact as what it grants: no weight is glued to a ray — that glue is called helicity and helicity is external vocabulary — no exchange operator is invented, so spin-statistics stays open, and the baryon card keeps its rung-relative spin, because the winding datum the card would need is the same regime datum the joint-entropy tier named as missing.

\begin{theorem}[{\texorpdfstring{BM.75 — The weight–branch–core attachment: the unique bijection and the attached quark mass splitting}{BM.75 - The weight-branch-core attachment: the unique bijection and the attached quark mass splitting}}]\label{tag:BM-75}
\textbf{Thm.}

%dsp-level:0
\[
(1)\quad\textbf{the derived twos and the declared names:}
\]

%dsp-level:0
\[
\qquad \textbf{the weight alphabet:}\quad\operatorname{spec}Q|_{\mathcal H_{\mathsf v}}=\{0,+1\},
\]

%dsp-level:0
\[
\qquad \omega:\ \{0,+1\}\mapsto\{0,-1\}
\]

%dsp-level:0
\[
\qquad [\mathrm{EW.36};\ \mathrm{Def\text{-}BM.14}(7)]:
\]

%dsp-level:0
\[
\textbf{the branch index:}\quad\beta\in\{-,+\},
\]

%dsp-level:0
\[
\qquad D_S(\mathcal F_g,+)>D_S(\mathcal F_g,-)\ \ \forall g\in\{1,2,3\}
\]

%dsp-level:0
\[
\qquad [\mathrm{Def\text{-}BM.14}(2),(4);\ \mathrm{BM.60}(5):\ \text{uniform, exact}];
\]

%dsp-level:0
\[
\textbf{the name rows (declared, frozen):}
\]

%dsp-level:0
\[
\qquad \text{down-type}:=\text{the heavier branch}\ (\beta{=}+);\quad\text{up-type}:=\text{the lighter}\ (\beta{=}-)
\]

%dsp-level:0
\[
\qquad [\mathrm{Def\text{-}BM.14}(8);\ \text{A-BM.6}];
\]

%dsp-level:0
\[
\qquad \text{the class labels}\ \{\text{neutral class},\ \text{unit class}\}
\]

%dsp-level:0
\[
\qquad [\mathrm{Def\text{-}BM.3}(2),(3);\ \text{A-BM.1}];
\]

%dsp-level:0
\[
(2)\quad\textbf{the ladder direction (derived):}
\]

%dsp-level:0
\[
\qquad E_\pm=\tfrac12(\sigma_1\pm i\sigma_2)=\text{the matrix units in the}\ \sigma_3\ \text{eigenbasis};
\]

%dsp-level:0
\[
\qquad [\sigma_3,E_\pm]=\pm2E_\pm
\]

%dsp-level:0
\[
\qquad [\mathrm{BM.23}(1);\ \text{machine over}\ \mathbb Q(i)]:
\]

%dsp-level:0
\[
E_+:\ \text{weight}\ 0\ \longmapsto\ \text{weight}\ {+}1
\]

%dsp-level:0
\[
\qquad [\text{machine:}\ E_+(0,1)^{\mathsf T}=(1,0)^{\mathsf T}]:
\]

%dsp-level:0
\[
\qquad \textbf{the unique raising step inside}\ \operatorname{spec}\{0,+1\};
\]

%dsp-level:0
\[
\textbf{the vertex grammar:}
\]

%dsp-level:0
\[
\qquad \text{the}\ \beta^-\ \text{vertex maps the neutral class to the unit class with the forced}\ q=-1\ \text{carrier};
\]

%dsp-level:0
\[
\qquad 0=(+1)+(-1)+0
\]

%dsp-level:0
\[
\qquad [\mathrm{BM.23}(2)\text{--}(5);\ \mathrm{BM.41}(1);\ \text{machine: integers}];
\]

%dsp-level:0
\[
\textbf{the substitution row (the hypothesis on its face):}
\]

%dsp-level:0
\[
\qquad \textbf{hypothesis (H):}\ \text{the two classes differ by exactly one down}\leftrightarrow\text{up constituent substitution, the minimal}\ \Delta q=+1\ \text{step}
\]

%dsp-level:0
\[
\qquad [\mathrm{BM.59}(9);\ \mathrm{BM.23}(4)];
\]

%dsp-level:0
\[
(3)\quad\textbf{the attachment theorem (uniqueness):}
\]

%dsp-level:0
\[
\qquad \text{seek a bijection}\quad f:\ \{\text{weight}\ 0,\ \text{weight}\ {+}1\}\ \longrightarrow\ \{\beta{=}+,\ \beta{=}-\}\quad\text{consistent with:}
\]

%dsp-level:0
\[
(\mathrm i)\ \ \text{the ladder direction:}\ E_+\ \text{carries the substitution's source to its target}\ [(2)];
\]

%dsp-level:0
\[
\qquad (\mathrm{ii})\ \ \text{the name row: down-type}=\beta{=}+\ [(1)];
\]

%dsp-level:0
\[
(\mathrm{iii})\ \ \text{the transport order: the neutral-core constituent is the heavier one},
\]

%dsp-level:0
\[
\quad m_q(0)=2>1=m_q(\pm1)
\]

%dsp-level:0
\[
\qquad [\mathrm{BM.59}(5);\ \mathrm{Def\text{-}BM.13}(1)];
\]

%dsp-level:0
\[
(\mathrm{iv})\ \ \text{the two-currency sign coherence of the same substitution:}
\]

%dsp-level:0
\[
\qquad \Delta m_{\mathrm{transport}}=+\frac{\ln(4/3)}{\mathfrak h}>0
\]

%dsp-level:0
\[
\qquad [\mathrm{BM.59}(4)]
\]

%dsp-level:0
\[
\qquad \text{and}
\]

%dsp-level:0
\[
\qquad \Delta_g=\hat\mu'^2\bigl(D_S(+)-D_S(-)\bigr)>0
\]

%dsp-level:0
\[
\qquad [\mathrm{BM.59}(9);\ \mathrm{Def\text{-}BM.14}(5)]:
\]

%dsp-level:0
\[
\text{the substitution's source is the heavier member in both currencies};
\]

%dsp-level:0
\[
\boxed{\ \textbf{exactly one bijection survives}\ (\mathrm i)\text{--}(\mathrm{iv}):\qquad\text{weight}\ 0\ \longleftrightarrow\ q_{\mathrm{core}}=0\ \longleftrightarrow\ \text{down-type}\ \longleftrightarrow\ \beta{=}+\ (\text{heavier});\qquad\text{weight}\ {+}1\ \longleftrightarrow\ q_{\mathrm{core}}={+}1\ \longleftrightarrow\ \text{up-type}\ \longleftrightarrow\ \beta{=}-\ (\text{lighter})\ }
\]

%dsp-level:0
\[
\qquad [\text{machine: exhaustive over the}\ 2!\ \text{bijections at the frozen ladder direction}];
\]

%dsp-level:0
\[
\textbf{the A-BM.6 bar is discharged by derivation:}
\]

%dsp-level:0
\[
\qquad \text{the attachment of the}\ Q\text{-weight index to the branch index}\ —\ \text{recorded there as}\ \neg\text{declared}\ —\ \text{is derived from the declared name rows plus}\ (2)\text{--}(\mathrm{iv});
\]

%dsp-level:0
\[
\qquad \neg\text{a new declaration issues here};
\]

%dsp-level:0
\[
(4)\quad\textbf{the attached quark mass splitting (exact per currency):}
\]

%dsp-level:0
\[
\qquad \textbf{the defect/pole currency:}\quad\Delta_g=m^2_{\mathrm{down}}-m^2_{\mathrm{up}}=\hat\mu'^2\bigl(D_S(\mathcal F_g,+)-D_S(\mathcal F_g,-)\bigr):
\]

%dsp-level:0
\[
\Delta_1=2\,\hat\mu'^2\in\mathbb Q\,\hat\mu'^2;
\]

%dsp-level:0
\[
\qquad \Delta_2=(\sqrt2-1)\,\hat\mu'^2\in\mathbb Q(\sqrt2)\,\hat\mu'^2;
\]

%dsp-level:0
\[
\qquad \Delta_3=\bigl(2\cos\tfrac\pi{10}-\sqrt3\bigr)\,\hat\mu'^2\in\mathbb Q\bigl(\sqrt3,\cos\tfrac\pi{10}\bigr)\,\hat\mu'^2;
\]

%dsp-level:0
\[
\qquad \Delta_g>0\ \ \forall g
\]

%dsp-level:0
\[
\qquad [\mathrm{Def\text{-}BM.14}(5);\ \text{machine: the certificates}\ 4>2,\ 2>1,\ 3<5];
\]

%dsp-level:0
\[
\textbf{the informational/transport-register currency:}
\]

%dsp-level:0
\[
\qquad \text{per constituent}\ \frac{\ln2}{\mathfrak h}>0;
\]

%dsp-level:0
\[
\qquad \text{per composite}\ \frac{\ln(4/3)}{\mathfrak h}>0
\]

%dsp-level:0
\[
\qquad [\mathrm{BM.59}(4),(5)];
\]

%dsp-level:0
\[
\textbf{the ensemble route-(ii) currency:}
\]

%dsp-level:0
\[
\qquad \frac{T(0,S)}{T(+1,S)}=\frac43>1:
\]

%dsp-level:0
\[
\qquad \text{the neutral class is heavier}
\]

%dsp-level:0
\[
\qquad [\mathrm{BM.28}(4)(\mathrm{ii});\ \mathrm{BM.30}(2)];
\]

%dsp-level:0
\[
\textbf{the opposing row:}
\]

%dsp-level:0
\[
\qquad \text{EM makes the charged member heavier by}\ c\,\varepsilon_{\mathrm{EM}}q^2m_{\mathrm{hid}}>0
\]

%dsp-level:0
\[
\qquad [\mathrm{BM.46}(4)]:
\]

%dsp-level:0
\[
\qquad \text{the unique sign-opposite correction, dominated on the sub-critical branch}\ \varepsilon_{\mathrm{EM}}<\varepsilon^{\mathrm{crit}}(c)
\]

%dsp-level:0
\[
\qquad [\mathrm{BM.62}(6),(9);\ \mathrm{BM.65}(3)\text{--}(5)];
\]

%dsp-level:0
\[
(5)\quad\textbf{the sign-coherence theorem:}
\]

%dsp-level:0
\[
\qquad \boxed{\ \text{in every derived currency in which the doublet splitting is formed, the down-type constituent is the heavier one}\ }
\]

%dsp-level:0
\[
[\text{machine: the sign census of}\ (4)\text{'s four rows}\ —\ \text{three concordant positives and one opposed correction with its exact domination branch}];
\]

%dsp-level:0
\[
\textbf{the currencies remain incommensurate:}
\]

%dsp-level:0
\[
\qquad \neg\exists\ \text{a derived ratio between}\ \Delta_g\ \text{and}\ \ln(4/3)/\mathfrak h\ —\ \text{the relating equation stays open}
\]

%dsp-level:0
\[
\qquad [\mathrm{BM.35}(5);\ \mathrm{BM.68}(7);\ \mathrm{BM.69}(4)];
\]

%dsp-level:0
\[
(6)\quad\textbf{what remains open:}
\]

%dsp-level:0
\[
\qquad \text{the full flavour composition}\ —\ \text{which three of the six states a class occupies}\ —\ \text{stays}\ \neg\text{derived:}\ \neg\text{a fractional constituent alphabet},
\]

%dsp-level:0
\[
\ \text{the end signs carry no position datum}
\]

%dsp-level:0
\[
\qquad [\mathrm{BM.28}(7);\ \mathrm{Def\text{-}BM.8}(5);\ \mathrm{Def\text{-}BM.13}(9)]:
\]

%dsp-level:0
\[
\textbf{the attachment theorem consumes only the single-substitution row,}\ \neg\textbf{the full triple};
\]

%dsp-level:0
\[
\qquad \text{the generation map is an order,}\ \neg\text{a number}\ [\mathrm{Def\text{-}BM.14}(6);\ \mathrm{BM.60}(7)];
\]

%dsp-level:0
\[
\qquad \text{the mixing matrix is}\ \neg\text{derived}\ [\mathrm{FT.145}(3)];
\]

%dsp-level:0
\[
(7)\quad\textbf{status:}\qquad\textbf{derived:}\ (1)\text{--}(5)\ \ [(3)\ \text{conditional on hypothesis (H), on its face}\ [\text{A §0.2a}]];
\]

%dsp-level:0
\[
\neg\text{a new declaration},
\]

%dsp-level:0
\[
\quad\neg\text{a P4 addition},
\]

%dsp-level:0
\[
\quad\neg\text{a measured value}\ [\mathrm{P5}],
\]

%dsp-level:0
\[
\quad\neg\text{a Yukawa-type object}\ [\mathrm{EW.15};\ \mathrm{FT.145}(3)];
\]

%dsp-level:0
\[
\qquad \text{§Cal-BM stays empty}.
\]

\textbf{Req.} D3.4 (Parts 1–2); EW.15, EW.36 (Part 2); P5 (Parts 1–2); Def-H.37(2), FT.60(5), FT.145(3) (Part 3); Def-BM.3(2), Def-BM.3(3), Def-BM.8(5), Def-BM.13(1), Def-BM.13(9), Def-BM.14(2), Def-BM.14(4), Def-BM.14(5), Def-BM.14(6), Def-BM.14(7), Def-BM.14(8), BM.23(1), BM.23(2), BM.23(4), BM.23(5), BM.28(4), BM.28(7), BM.30(2), BM.35(5), BM.41(1), BM.46(4), BM.59(4), BM.59(5), BM.59(9), BM.60(5), BM.60(7), BM.62(6), BM.62(9), BM.65(3), BM.65(5), BM.68(7), BM.69(4) (this part); exact $\mathbb Q(i)$ matrix-unit algebra, exhaustive bijection censuses, exact radical certificates, prime-log-ring sign tests [math; machine: \texttt{checks/\allowbreak{}verify\_\allowbreak{}H\_\allowbreak{}chiral\_\allowbreak{}assignment\_\allowbreak{}opening.py} §B, §D].

\textbf{Pf.}

%dsp-level:0
\[
(1):\ \mathrm{EW.36},\ \mathrm{Def\text{-}BM.14}(2),(4),(7),(8),\ \mathrm{Def\text{-}BM.3}(2),(3)\ \text{re-cited}\ (\text{frozen})
\]

%dsp-level:0
\[
(2):\ E_+=\left(\begin{smallmatrix}0&1\\0&0\end{smallmatrix}\right)\ \text{in the weight basis with}\ Q=\operatorname{diag}(1,0):\quad E_+(0,1)^{\mathsf T}=(1,0)^{\mathsf T}\ [\text{machine}];\qquad[\sigma_3,E_+]=2E_+\ [\text{machine over}\ \mathbb Q(i)];\qquad0=(+1)+(-1)+0\ [\text{machine: integers}]
\]

%dsp-level:0
\[
(3):\ \text{the census: bijection}\ f_A:\ 0\mapsto+,\ {+}1\mapsto-\ \ \text{and}\ f_B:\ 0\mapsto-,\ {+}1\mapsto+;
\]

%dsp-level:0
\[
\qquad f_A\ \text{satisfies}\ (\mathrm i)\text{--}(\mathrm{iv}):\ \text{the}\ E_+\text{-source is weight}\ 0
\]

%dsp-level:0
\[
= \text{down-type}
\]

%dsp-level:0
\[
= \beta{=}+
\]

%dsp-level:0
\[
= \text{the heavier branch}
\]

%dsp-level:0
\[
= \text{the heavier core}\ [\text{machine: the four condition evaluations}];
\]

%dsp-level:0
\[
f_B\ \text{fails}\ (\mathrm i)\ \text{and}\ (\mathrm{iii})\wedge(\mathrm{iv})\ \text{jointly: it makes the substitution's source}\ (\text{weight}\ 0,\ \text{core}\ 0,\ \text{transport-heavier})\ \text{the lighter branch}\ \beta{=}-,
\]

%dsp-level:0
\[
\ \text{so the same substitution that raises}\ \Delta m_{\mathrm{transport}}
\]

%dsp-level:0
\[
= +\ln(4/3)/\mathfrak h\ \text{would lower the defect read}\ (-\Delta_g)\ \text{against}\ \mathrm{BM.59}(9)\text{'s derived}\ +\Delta_g\ [\text{machine: the sign contradiction}];
\]

%dsp-level:0
\[
\qquad \text{the ladder direction is frozen by the grammar}\ [(2)];
\]

%dsp-level:0
\[
\qquad \text{the two bijections:}\ 2\ \text{cases, one survivor}\ [\text{machine}]
\]

%dsp-level:0
\[
(4):\ \Delta_g\ \text{re-formed in the three exact fields with the certificates}\ 4>2;
\]

%dsp-level:0
\[
\ 2>1;
\]

%dsp-level:0
\[
\ 2\cos\tfrac\pi{10}>\sqrt3\iff\cos^2\tfrac\pi{10}>\tfrac34\iff\tfrac{5+\sqrt5}8>\tfrac34
\]

%dsp-level:0
\[
\iff 5+\sqrt5>6\iff\sqrt5>1
\]

%dsp-level:0
\[
\iff 5>1\ [\text{machine: exact in}\ \mathbb Q(\sqrt5)];
\]

%dsp-level:0
\[
\qquad \text{the currency rows re-cited}\ [\mathrm{BM.59}(4),(5);\ \mathrm{BM.28}(4)(\mathrm{ii});\ \mathrm{BM.46}(4)]
\]

%dsp-level:0
\[
(5):\ \text{the sign census of the four rows}\ [\text{machine}];
\]

%dsp-level:0
\[
\qquad \text{the no-ratio row re-cited}\ [\mathrm{BM.68}(7)(\mathrm d),(\mathrm e);\ \mathrm{BM.69}(4)]
\]

%dsp-level:0
\[
(6)\text{--}(7):\ \text{the open rows re-cited};\qquad\text{the standing-rule scans}\ [\text{machine}]. \ \blacksquare
\]

\textbf{Phys.} Two derived twos have always sat side by side — the seed's weight doublet and the defect floor's branch pair — with their attachment recorded as undeclared. The chiral assignment supplies the missing link, and the link turns out to be forced: the ladder raises exactly one weight, the beta vertex raises exactly one core, the heavier branch is named down-type, and the heavier core is the neutral one by an exact transport count — of the two possible gluings only one keeps the sign of the splitting the same in the transport currency and in the defect currency, and both signs are frozen derived rows. So the attachment is a theorem given the substitution hypothesis and the declared names, and the quark mass splitting is now attached rather than floating: per generation it is two, root-two-minus-one, and two-cosines-of-eighteen-degrees-minus-root-three, times the undeclared defect scale, positive in all three, with the informational currencies agreeing in sign and the electromagnetic correction opposing exactly where the sub-critical branch dominates it. What stays open stays open loudly: which three of the six states compose a given baryon is still unreadable, because the attachment consumes one substitution step and not a whole triple, and no ratio between the currencies is derived — the sign agrees, the scale does not, and that is application-layer content.

\end{theorem}

\begin{theorem}[{\texorpdfstring{BM.76 — The chiral vertex dynamics: the nilpotent step operator, the flip channel, and the rate rows}{BM.76 - The chiral vertex dynamics: the nilpotent step operator, the flip channel, and the rate rows}}]\label{tag:BM-76}
\textbf{Thm.}

%dsp-level:0
\[
(1)\quad\textbf{the vertex step operator (an in-theorem abbreviation):}
\]

%dsp-level:0
\[
\qquad \text{on the extended carrier}\ \mathcal H_{\mathrm{ext}}=\mathcal H_{\mathrm{int}}\otimes\mathcal H_\omega\cong\mathbb C^4
\]

%dsp-level:0
\[
\qquad [\mathrm{Def\text{-}H.34}]:
\]

%dsp-level:0
\[
\boxed{\ \widehat V\ :=\ E_+\otimes P_L\ }
\]

%dsp-level:0
\[
\qquad [E_+=\tfrac12(\sigma_1+i\sigma_2)\ [\mathrm{BM.23}(1)];\quad P_L=\tfrac12(\mathrm{id}+\hat\chi)\ [\mathrm{Def\text{-}H.33}];\quad\text{exact over}\ \mathbb Q(i)\ \text{in the ordered basis of}\ \mathrm{Def\text{-}H.34}];
\]

%dsp-level:0
\[
(2)\quad\textbf{the nilpotency is the alphabet overflow's operator form:}
\]

%dsp-level:0
\[
\qquad E_+^2=0\implies\boxed{\ \widehat V^{\,2}=0\ }
\]

%dsp-level:0
\[
\qquad [\text{machine: entrywise over}\ \mathbb Q(i)]:
\]

%dsp-level:0
\[
\text{the ladder has one step inside the doublet}\ —\ \mathrm{BM.23}(3)\text{'s ternary overflow}\ \bigl(\{-1,0,+1\}+1=\{0,+1,+2\}\ \text{leaves the alphabet}\bigr)\ \textbf{as an exact operator identity};
\]

%dsp-level:0
\[
\qquad \neg\exists\ \text{an iterated charged step inside one doublet};
\]

%dsp-level:0
\[
(3)\quad\textbf{the rank, image and kernel rows:}
\]

%dsp-level:0
\[
\qquad \operatorname{rk}\widehat V=1=\operatorname{rk}E_+\cdot\operatorname{rk}P_L;
\]

%dsp-level:0
\[
\qquad \operatorname{im}\widehat V=\operatorname{span}\bigl\{|\text{weight}\,{+}1\rangle\otimes|L\rangle\bigr\};
\]

%dsp-level:0
\[
\qquad \dim\ker\widehat V=3
\]

%dsp-level:0
\[
\qquad [\text{machine: exact rank and image over}\ \mathbb Q(i)]:
\]

%dsp-level:0
\[
\textbf{the target is the up-type weight on the occupied ray}\qquad[\mathrm{BM.75}(3)\ \text{attachment}];
\]

%dsp-level:0
\[
\widehat V\,(I\otimes P_R)=0,
\]

%dsp-level:0
\[
\qquad (I\otimes P_L)\,\widehat V=\widehat V
\]

%dsp-level:0
\[
\qquad [\text{machine}]:
\]

%dsp-level:0
\[
\qquad \textbf{the step fires from the left ray alone}\ —\ \text{the chiral selectivity of the transition, derived given A-BM.7};
\]

%dsp-level:0
\[
(4)\quad\textbf{the Cartan row:}
\]

%dsp-level:0
\[
\qquad \bigl[W_3\otimes I_\omega,\ \widehat V\bigr]=\widehat V
\]

%dsp-level:0
\[
\qquad [\ [\sigma_3,E_+]=2E_+,\ W_3=\tfrac12\sigma_3;\ \text{machine}]:
\]

%dsp-level:0
\[
\qquad \textbf{the step operator carries weight}\ {+}1\ \text{under the seed Cartan}\ —\ \text{the operator form of the composite raise}\ \Delta q
\]

%dsp-level:0
\[
= +1\ \text{with the compensating emission}\ [\mathrm{BM.23}(2),(4)];
\]

%dsp-level:0
\[
(5)\quad\textbf{the}\ \omega\textbf{-covariant image vertex:}\qquad(I\otimes\hat\omega)\,\widehat V\,(I\otimes\hat\omega)^{-1}=E_+\otimes P_R\qquad[\mathrm{FT.122}(4);\ \text{machine}]:
\]

%dsp-level:0
\[
\textbf{the image step fires from the right ray}\ —\ \text{the CPT-partner vertex}\ [\mathrm{BM.25};\ \mathrm{BM.27};\ \text{A-BM.5's image-orientation row is ray-resolved}];
\]

%dsp-level:0
\[
\textbf{the flip conjugation:}
\]

%dsp-level:0
\[
\qquad (I\otimes\hat\kappa)\,\widehat V\,(I\otimes\hat\kappa)=E_+\otimes P_R
\]

%dsp-level:0
\[
\qquad [\hat\kappa P_L\hat\kappa=P_R:\ \{\hat\chi,\hat\kappa\}=0\ [\mathrm{Def\text{-}H.34}];\ \text{machine}]:
\]

%dsp-level:0
\[
\qquad \text{the partner involution}\ —\ \text{the EM flip's carrier}\ [\mathrm{FT.129}(2),(5)]\ —\ \text{moves the vertex to the opposite ray};
\]

%dsp-level:0
\[
(6)\quad\textbf{the flip-channel rows (frozen, re-cited):}
\]

%dsp-level:0
\[
\qquad \text{rung one:}\ F_{LR}=F_{RL}=0\ \text{— the sectors force-decoupled}\ [\mathrm{FT.128}(2);\ \mathrm{FT.144}(6)];
\]

%dsp-level:0
\[
\qquad \text{rung two:}\ F_{LR}
\]

%dsp-level:0
\[
= q\,
\]

%dsp-level:0
\[
I_2\ \text{under the}\ \mathrm{Def\text{-}H.43}\ \text{promotion},
\]

%dsp-level:0
\[
\ \text{splitting}\ 2|q|\in\{0,2\}\ [\mathrm{FT.129}(5);\ \mathrm{FT.130}(5);\ \mathrm{BM.22}(2),(4)]:
\]

%dsp-level:0
\[
\text{the neutral core keeps parity doubling; the charged core flip-couples both rays}
\]

%dsp-level:0
\[
\qquad [\mathrm{FT.130}(6);\ \text{machine re-cited}];
\]

%dsp-level:0
\[
\qquad \text{rung three adds no carrier flip}\ [\mathrm{FT.117};\ \mathrm{FT.144}(6)];
\]

%dsp-level:0
\[
(7)\quad\textbf{the rate and arrow rows (frozen, re-cited):}
\]

%dsp-level:0
\[
\qquad \Gamma_{\mathrm{tr}}(\varepsilon_q)=\frac{x_+x_-}{1-x_+x_-}\in\Bigl[0,\tfrac13\Bigr],
\]

%dsp-level:0
\[
\qquad x_+x_-=4\varepsilon_q(1-4\varepsilon_q)=\tfrac14-16\Bigl(\varepsilon_q-\tfrac18\Bigr)^2\le\tfrac14
\]

%dsp-level:0
\[
\qquad [\mathrm{BM.44}(1),(2);\ \mathrm{BM.19}(2);\ \text{machine: the SOS parabola}]:
\]

%dsp-level:0
\[
\varepsilon_q=0\implies\Gamma_{\mathrm{tr}}=0:
\]

%dsp-level:0
\[
\qquad \textbf{no flicker, no decay}
\]

%dsp-level:0
\[
\qquad [\mathrm{BM.23}(8)];
\]

%dsp-level:0
\[
\qquad \tau_{\mathrm{life}}\ge3\ \ \text{ticks}
\]

%dsp-level:0
\[
\qquad [\mathrm{BM.44}(3)];
\]

%dsp-level:0
\[
\qquad \text{forward-only},\ \neg\exists\,\tau^{-1}
\]

%dsp-level:0
\[
\qquad [\mathrm{A6};\ \mathrm{BM.23}(9);\ \mathrm{FT.144}(2)]:
\]

%dsp-level:0
\[
\textbf{the vertex fires per tick as a weight,}\ \neg\textbf{as an external rate}
\]

%dsp-level:0
\[
\qquad [\mathrm{BM.44}(4),(5):\ \neg\text{external seconds}];
\]

%dsp-level:0
\[
\qquad \textbf{the opening of}\ \mathrm{BM.77}\ \textbf{is kinematic; its rate stays the flicker's own};
\]

%dsp-level:0
\[
(8)\quad\textbf{scope:}
\]

%dsp-level:0
\[
\qquad \text{what is constructed here is the vertex-level dynamics of the}\ Q\ \text{read}\ —\ \text{the ladder step, the chiral selection, the flip channel, the rate};
\]

%dsp-level:0
\[
\qquad \textbf{the field-level dynamics of}\ Q\ \text{(gauge-field dynamics, radiation, back-reaction) remains}\ \neg\textbf{constructed}
\]

%dsp-level:0
\[
\qquad [\mathrm{BM.6}(5);\ \mathrm{Def\text{-}BM.12}(7);\ \mathrm{BM.50}(3)];
\]

%dsp-level:0
\[
\widehat V\ \text{is}\ \neg\text{a new force: A8's completeness row stands}\ —\ \text{no sub-block restriction is an A8 force}
\]

%dsp-level:0
\[
\qquad [\mathrm{A8};\ \mathrm{BM.52}(1)];
\]

%dsp-level:0
\[
\qquad \widehat V\ \text{is an in-theorem abbreviation of}\ \mathrm{BM.23}(1)\text{'s ladder and}\ \mathrm{Def\text{-}H.33}\text{'s projector};
\]

%dsp-level:0
\[
(9)\quad\textbf{status:}
\]

%dsp-level:0
\[
\qquad \textbf{derived:}\ (1)\text{--}(7)\ \ [\text{the occupancy rows given A-BM.7}];
\]

%dsp-level:0
\[
\qquad \neg\text{a P4 addition},\quad\neg\text{a measured lifetime}\ [\mathrm{P5}],\quad\neg\text{a new dimensionful scale}\ [\mathrm{P1}];
\]

%dsp-level:0
\[
\text{the conditionalities}\ (\mathrm{Def\text{-}H.43}\ \text{promotion}\ [\mathrm{BM.22}(5)],\ \mathrm{BM.8}(2)\ \text{stationarity for the chart exponential}\ [\mathrm{BM.44}(3)])\ \text{are inherited at the point of use}\ [\text{A §0.2a}].
\]

\textbf{Req.} A6, A8 (Part 1); EW.32, EW.36 (Part 2); P1, P5 (Parts 1–2); Def-H.33, Def-H.34, Def-H.43, FT.117, FT.122(4), FT.128(2), FT.129(2), FT.129(5), FT.130(5), FT.130(6), FT.144(2), FT.144(6) (Part 3); Def-BM.12(7), BM.6(5), BM.19(2), BM.22(2), BM.22(4), BM.22(5), BM.23(1), BM.23(2), BM.23(4), BM.23(8), BM.23(9), BM.25, BM.27, BM.44(1), BM.44(2), BM.44(3), BM.44(4), BM.44(5), BM.50(3), BM.52(1), BM.73(2), BM.73(3), BM.75(3) (this part); exact $\mathbb Q(i)$ Kronecker algebra, rank/image/kernel computation, the SOS parabola completion [math; machine: \texttt{checks/\allowbreak{}verify\_\allowbreak{}H\_\allowbreak{}chiral\_\allowbreak{}assignment\_\allowbreak{}opening.py} §C].

\textbf{Pf.}

%dsp-level:0
\[
(1)\text{--}(2):\ E_+=\left(\begin{smallmatrix}0&1\\0&0\end{smallmatrix}\right),\ P_L=\left(\begin{smallmatrix}\tfrac12&\tfrac i2\\-\tfrac i2&\tfrac12\end{smallmatrix}\right)\ \text{in the ordered bases}\ [\mathrm{Def\text{-}H.33}(5)\ \text{form};\ \mathrm{Def\text{-}H.34}]:\qquad\widehat V\ \text{built as the Kronecker product, all}\ 16\ \text{entries exact over}\ \mathbb Q(i)\ [\text{machine}];\qquad E_+^2=0\implies\widehat V^2=(E_+^2)\otimes(P_L^2)=0\ [\text{machine: entrywise}]
\]

%dsp-level:0
\[
(3):\ \operatorname{rk}(A\otimes B)=\operatorname{rk}A\cdot\operatorname{rk}B\ [\text{math}];
\]

%dsp-level:0
\[
\qquad \operatorname{rk}E_+=1,\ \operatorname{rk}P_L=1\ [\text{machine: exact elimination}];
\]

%dsp-level:0
\[
\qquad \operatorname{im}\widehat V\ \text{spanned by the image of the basis}\ \{\,|\text{w}0\rangle\otimes|L\rangle\,\}\mapsto|\text{w}{+}1\rangle\otimes|L\rangle\ [\text{machine}];
\]

%dsp-level:0
\[
\qquad \widehat V(I\otimes P_R)=E_+\otimes P_LP_R=0\ [\mathrm{FT.122}(2);\ \text{machine}]
\]

%dsp-level:0
\[
(4):\ [\sigma_3,E_+]
\]

%dsp-level:0
\[
= \sigma_3E_+-E_+\sigma_3
\]

%dsp-level:0
\[
= 2E_+\ [\text{machine}]
\]

%dsp-level:0
\[
\implies [W_3\otimes I,\ E_+\otimes P_L]
\]

%dsp-level:0
\[
= [W_3,E_+]\otimes P_L
\]

%dsp-level:0
\[
= E_+\otimes P_L\ [\text{machine: the}\ 4\times4\ \text{commutator entrywise}]
\]

%dsp-level:0
\[
(5):\ \hat\omega P_L\hat\omega=P_R\ [\mathrm{FT.122}(4);\ \text{machine}];
\]

%dsp-level:0
\[
\qquad \hat\kappa P_L\hat\kappa=\tfrac12(I+\hat\kappa\hat\chi\hat\kappa)=\tfrac12(I-\hat\chi)=P_R\ [\{\hat\chi,\hat\kappa\}=0;\ \mathrm{Def\text{-}H.34};\ \text{machine}]
\]

%dsp-level:0
\[
(6)\text{--}(7):\ \text{the frozen rows re-cited}\ [\mathrm{FT.128}(2);\ \mathrm{FT.129}(5);\ \mathrm{FT.130}(5),(6);\ \mathrm{FT.117};\ \mathrm{FT.144}(6);\ \mathrm{BM.44}(1)\text{--}(5);\ \mathrm{BM.23}(8),(9)];
\]

%dsp-level:0
\[
\qquad x_+x_-=\tfrac14-16(\varepsilon_q-\tfrac18)^2\ [\text{machine: the SOS completion}];
\]

%dsp-level:0
\[
\qquad \Gamma_{\mathrm{tr}}(\tfrac18)=\tfrac{1/4}{3/4}=\tfrac13\ [\text{machine: exact Fractions}]
\]

%dsp-level:0
\[
(8)\text{--}(9):\ \text{the scope record};\qquad\text{the standing-rule scans}\ [\text{machine}]. \ \blacksquare
\]

\textbf{Phys.} The charged-current vertex becomes a single exact operator: the weak ladder step tensored with the occupied ray's projector, four by four over the rationals with i adjoined. Its algebra is the physics. Squared, it vanishes — the same nilpotency that makes the ladder's second step leave the ternary alphabet, now an operator identity, so a constituent takes the charged step once and never twice. Its rank is one, its image is the up-type weight on the left ray, and it annihilates anything sitting in the right ray: given the assignment, the transition fires from matter's own ray alone, which is the corpus's exact counterpart of a maximally chiral charged current — derived, not postulated. Conjugating by orientation reversal or by the partner involution moves it to the right ray, giving the image vertex that the antiproton and antineutrino rows occupy, and the Cartan commutator records the charge bookkeeping: the step carries weight plus one, the composite raises by one, and the emitted carrier takes minus one away. The rates were already exact and stay frozen: a weight per tick bounded by a third, three ticks minimum, nothing at the sharp crystallization threshold, and no seconds anywhere — the opening this tier derives elsewhere is kinematic, and the dynamics keeps the flicker's own clock.

\end{theorem}

\begin{theorem}[{\texorpdfstring{BM.77 — The read specification and the allowed branch: the formed Q-value, the exact trichotomy, and the rescope rows}{BM.77 - The read specification and the allowed branch: the formed Q-value, the exact trichotomy, and the rescope rows}}]\label{tag:BM-77}
\textbf{Thm.}

%dsp-level:0
\[
(1)\quad\textbf{the specification (the declaration, logged as A-BM.8 before use):}
\]

%dsp-level:0
\[
\qquad \textbf{the}\ \beta\textbf{-vertex Q-value is read:}
\]

%dsp-level:0
\[
\qquad \text{on the baryon classes}\ —\ \text{route}\ (\mathrm{ii}),\ \text{the unnormalized total deficit}\quad T(q,S)=W(q)\,\overline{\Delta S}(S)
\]

%dsp-level:0
\[
\qquad [\mathrm{BM.28}(4)(\mathrm{ii});\ \mathrm{BM.30}(1)\ \text{route census}];
\]

%dsp-level:0
\[
\text{on the emitted lines}\ —\ \text{the transport-species family form}\quad m^2(\text{line}\ i)=\hat\mu^2\ell_c^2\,G\bigl(c(i)\bigr),
\]

%dsp-level:0
\[
\qquad G(c(i))=A_0-B_0\cos\varphi_i
\]

%dsp-level:0
\[
\qquad [\mathrm{Def\text{-}BM.11}(5);\ \mathrm{BM.39}(1)];
\]

%dsp-level:0
\[
\textbf{the expression is}\ \mathrm{BM.41}(4)(\mathrm d)\text{'s own frozen}\ Q_{\mathrm{ii}}\ —\ \text{formed there, priced there, adopted here};
\]

%dsp-level:0
\[
\qquad \textbf{the scope is the vertex kinematics alone:}\ \ \text{route}\ (\mathrm{ii})\ \text{is}\ \neg\text{elevated to a global mass definition}\ [\mathrm{BM.30}(5);\ \mathrm{BM.37}(6)\ \text{stand}];
\]

%dsp-level:0
\[
(2)\quad\textbf{the formed Q-value:}
\]

%dsp-level:0
\[
\qquad Q_\beta^{\mathrm{spec}}:=m_n^{(\mathrm{ii})}-m_p^{(\mathrm{ii})}-m_\ell-m_\nu;
\]

%dsp-level:0
\[
\qquad m_n^{(\mathrm{ii})}=\frac{T(0,S)}{\mathfrak h},\quad m_p^{(\mathrm{ii})}=\frac{T(+1,S)}{\mathfrak h}
\]

%dsp-level:0
\[
\qquad [\text{the route-(ii) mass read};\ \mathrm{BM.30}(1)]:
\]

%dsp-level:0
\[
\boxed{\ Q_\beta^{\mathrm{spec}}\,\mathfrak h\ =\ \Delta T(S)\ -\ \rho_\ell\sqrt{G\bigl(c(0)\bigr)}\ -\ \rho_\nu\sqrt{G\bigl(c(i_\nu)\bigr)},\qquad\Delta T(S):=T(0,S)-T(+1,S)=\bigl(W(0)-W(+1)\bigr)\overline{\Delta S}(S)=5\,\overline{\Delta S}(S)\ }
\]

%dsp-level:0
\[
[\mathrm{BM.28}(4)(\mathrm{ii})\ \text{and its Pf row}\ T(q,S)=W(q)\overline{\Delta S}(S);\ \ \text{the electron row}\ i=0\ [\mathrm{Def\text{-}BM.6}(1)];\ \ \text{machine}];
\]

%dsp-level:0
\[
\textbf{the dimension row:}
\]

%dsp-level:0
\[
\qquad \rho_\ell:=\mathfrak h\,\hat\mu_\ell\,\ell_c^{(\ell)},\quad\rho_\nu:=\mathfrak h\,\hat\mu_\nu\,\ell_c^{(\nu)}\ \in\ \mathfrak q^0
\]

%dsp-level:0
\[
\qquad [\mathrm{CT2.1};\ \mathrm{D2.9};\ \mathrm{Def\text{-}H.37}(1):\ [\hat\mu^2]=\mathfrak q^2;\ \text{machine: the exponent arithmetic}]:
\]

%dsp-level:0
\[
\text{the scales}\ (\hat\mu_\ell,\hat\mu_\nu,\ell_c,g,g')\ \text{remain undeclared P4/Cal-F data}\ [\mathrm{FT.143}]\ —\ \rho_\ell,
\]

%dsp-level:0
\[
\rho_\nu\ \text{are restatements,}\ \neg\text{new members}
\]

%dsp-level:0
\[
\qquad [\mathrm{P4};\ \mathrm{P1};\ \text{Cal-F logged}];
\]

%dsp-level:0
\[
\textbf{the exact baryon rows:}
\]

%dsp-level:0
\[
\qquad \Delta T(S)=5\,m^2_B(\cdot,S)
\]

%dsp-level:0
\[
\qquad [\overline{\Delta S}(S)=m^2_B(\cdot,S):\ \mathrm{FT.107}(3);\ \mathrm{BM.28}(4)(\mathrm i),(\mathrm{ii});\ \text{the frozen sector expectations}\ \mathrm{BM.20}(4)/\mathrm{BM.21}(2)\ \text{re-cited at}\ \mathrm{BM.62}(2)]:
\]

%dsp-level:0
\[
\boxed{\ \Delta T\ =\ \Bigl(\frac{4219785}{98956},\ \ \frac{9608310}{393679},\ \ \frac{455}{23},\ \ \frac{68565}{3088}\Bigr)\ \ \text{at}\ S=(0,\,1,\,2,\,3)\ }
\]

%dsp-level:0
\[
\qquad [\text{machine: exact rational arithmetic}];
\]

%dsp-level:0
\[
\Delta T_{\min}=\frac{455}{23},
\]

%dsp-level:0
\[
\qquad \Delta T_{\max}=\frac{4219785}{98956},
\]

%dsp-level:0
\[
\qquad 0<\Delta T_{\min}\le\Delta T(S)\le\Delta T_{\max}
\]

%dsp-level:0
\[
\qquad [\text{machine: exact comparisons}];
\]

%dsp-level:0
\[
\qquad \textbf{the}\ \nu\text{-line index}\ i_\nu\ \text{is unresolved}\ \ —\ \textbf{every row below holds uniformly in}\ i_\nu\in\{0,1,2\}\ [(4)\text{'s sandwich is line-independent}];
\]

%dsp-level:0
\[
(3)\quad\textbf{the formedness row:}
\]

%dsp-level:0
\[
\qquad \text{the expression of}\ (2)\ \text{is}\ \mathrm{BM.41}(4)(\mathrm d)\text{'s frozen}\ Q_{\mathrm{ii}}\ \text{with the lepton scales carried symbolically}\ —\ \neg\text{a cross-currency assumption is made}:
\]

%dsp-level:0
\[
\text{no informational value is equated to a geometric one};
\]

%dsp-level:0
\[
\qquad \text{the mass-form bridge is application-layer content, untouched here}
\]

%dsp-level:0
\[
\qquad [\mathrm{BM.35}(5);\ \mathrm{BM.68}(8);\ \mathrm{BM.69}(5);\ \text{pointer}];
\]

%dsp-level:0
\[
\qquad \text{the EM and transport rest-slot refinements are}\ \neg\text{summands of this read}\ —\ \text{disjoint alphabets, no double count}
\]

%dsp-level:0
\[
\qquad [\mathrm{BM.37}(2);\ \mathrm{BM.62}(11);\ \mathrm{BM.66}(6)];
\]

%dsp-level:0
\[
(4)\quad\textbf{the universal sandwich:}
\]

%dsp-level:0
\[
\qquad A_0-B_0=\frac38\Bigl(g-\frac{g'}3\Bigr)^2+\frac{g'^2}{48}>0,
\]

%dsp-level:0
\[
\qquad A_0+B_0=\frac38\Bigl(g+\frac{g'}3\Bigr)^2+\frac{g'^2}{48}>0
\]

%dsp-level:0
\[
\qquad [\mathrm{BM.39}(1);\ \text{machine: the SOS expansions}]:
\]

%dsp-level:0
\[
\boxed{\ \sqrt{A_0-B_0}\ \le\ \sqrt{G\bigl(c(i)\bigr)}\ \le\ \sqrt{A_0+B_0}\qquad\forall\,\text{line}\ i,\ \ \forall\,\text{channel}\ \bigl(\delta\in\bigl\{\tfrac29,\tfrac2{27},\tfrac18\bigr\}\bigr),\ \ \forall\,(g,g')\ne(0,0)\ }
\]

%dsp-level:0
\[
\qquad [\mathrm{BM.39}(1),(3);\ |\cos|\le1;\ \text{machine}];
\]

%dsp-level:0
\[
(5)\quad\textbf{the trichotomy theorem:}\qquad\text{for every frozen cell}\ S\ \text{and every admissible}\ (g,g')\ne(0,0):
\]

%dsp-level:0
\[
(\mathrm a)\quad\textbf{the opening region is non-empty:}
\]

%dsp-level:0
\[
\qquad \mathcal O_S:=\Bigl\{(\rho_\ell,\rho_\nu)>0:\ \ \rho_\ell\sqrt{G_\ell}+\rho_\nu\sqrt{G_\nu}<\Delta T(S)\Bigr\}\ne\varnothing:
\]

%dsp-level:0
\[
\textbf{the universal symbolic witness:}
\]

%dsp-level:0
\[
\qquad \rho_\ell=\rho_\nu=\frac{\Delta T(S)}{4\sqrt{A_0+B_0}}\ \ \in\ \mathcal O_S
\]

%dsp-level:0
\[
\qquad [\text{machine: the sandwich gives}\ \rho_\ell\sqrt{G_\ell}+\rho_\nu\sqrt{G_\nu}\le\tfrac{\Delta T}2<\Delta T];
\]

%dsp-level:0
\[
\textbf{the integer-certified witness at the rational point}\ (g,g',\rho_\ell,\rho_\nu)=(1,1,1,1):
\]

%dsp-level:0
\[
\qquad A_0+B_0=\frac7{16}+\frac14=\frac{11}{16},
\]

%dsp-level:0
\[
\qquad \sqrt{G}\le\frac{\sqrt{11}}4<1\quad[11<16]:
\]

%dsp-level:0
\[
Q_\beta^{\mathrm{spec}}\mathfrak h\ \ge\ \frac{455}{23}-\frac{\sqrt{11}}2\ >\ 0
\]

%dsp-level:0
\[
\qquad \iff
\]

%dsp-level:0
\[
\qquad \frac{910}{23}>\sqrt{11}
\]

%dsp-level:0
\[
\qquad \iff
\]

%dsp-level:0
\[
\qquad 828100>5819
\]

%dsp-level:0
\[
\qquad [\text{machine: the integer certificate}];
\]

%dsp-level:0
\[
(\mathrm b)\quad\textbf{the closing region is non-empty:}
\]

%dsp-level:0
\[
\qquad \mathcal C_S:=\Bigl\{(\rho_\ell,\rho_\nu)>0:\ \ \rho_\ell\sqrt{G_\ell}+\rho_\nu\sqrt{G_\nu}>\Delta T(S)\Bigr\}\ne\varnothing:
\]

%dsp-level:0
\[
\textbf{the universal symbolic witness:}
\]

%dsp-level:0
\[
\qquad \rho_\ell=\rho_\nu=\frac{\Delta T(S)}{\sqrt{A_0-B_0}}\ \ \in\ \mathcal C_S
\]

%dsp-level:0
\[
\qquad [\text{machine: the sandwich gives}\ \rho_\ell\sqrt{G_\ell}+\rho_\nu\sqrt{G_\nu}\ge2\,\Delta T>\Delta T];
\]

%dsp-level:0
\[
\textbf{the integer-certified witness at}\ (1,1,100,100):
\]

%dsp-level:0
\[
\qquad A_0-B_0=\frac3{16},
\]

%dsp-level:0
\[
\qquad \sqrt{G}\ge\frac{\sqrt3}4:
\]

%dsp-level:0
\[
\qquad Q_\beta^{\mathrm{spec}}\mathfrak h\ \le\ \frac{4219785}{98956}-50\sqrt3\ <\ 43-50\sqrt3\ <\ 0
\]

%dsp-level:0
\[
[\text{machine:}\ 4219785<43\cdot98956=4255108;\quad50\sqrt3>43\iff7500>1849];
\]

%dsp-level:0
\[
(\mathrm c)\quad\textbf{the critical locus is exact:}
\]

%dsp-level:0
\[
\qquad \partial\mathcal O_S:\quad\rho_\ell\sqrt{G\bigl(c(0)\bigr)}+\rho_\nu\sqrt{G\bigl(c(i_\nu)\bigr)}=\Delta T(S)
\]

%dsp-level:0
\[
\qquad [\text{machine: the boundary equation}];
\]

%dsp-level:0
\[
\textbf{on the symmetric ray}\ \rho_\ell=\rho_\nu=\rho\ \textbf{at the witness couplings}\ g=g'=1:
\]

%dsp-level:0
\[
\qquad \boxed{\ \frac{273}{23}\ <\ \rho^{\mathrm{crit}}(S,i_\nu)\ <\ \frac{42197850}{791648}\ }
\]

%dsp-level:0
\[
\qquad \text{i.e.}\quad11<\rho^{\mathrm{crit}}<54
\]

%dsp-level:0
\[
\qquad [\text{machine}]:
\]

%dsp-level:0
\[
[\text{the certificates:}\ \sqrt{11}<\tfrac{10}3\iff99<100;\quad\sqrt3>\tfrac85\iff75>64;\quad273>253;\quad42197850<54\cdot791648=42748992];
\]

%dsp-level:0
\[
\textbf{the structure mirrors}\ \mathrm{BM.62}(6)\text{'s critical strength at the scale-combination locus:}
\]

%dsp-level:0
\[
\qquad \rho<\rho^{\mathrm{crit}}\iff Q_\beta^{\mathrm{spec}}>0;
\]

%dsp-level:0
\[
\qquad \rho=\rho^{\mathrm{crit}}\iff Q_\beta^{\mathrm{spec}}=0;
\]

%dsp-level:0
\[
\qquad \rho>\rho^{\mathrm{crit}}\iff Q_\beta^{\mathrm{spec}}<0
\]

%dsp-level:0
\[
\qquad [\text{machine: monotone in}\ \rho];
\]

%dsp-level:0
\[
(6)\quad\textbf{the verdict:}
\]

%dsp-level:0
\[
\qquad \boxed{\ \textbf{the beta decay is allowed:}\ \ \text{on the specified read the vertex is kinematically open exactly on}\ \mathcal O_S\ —\ \text{a non-empty exact region at every frozen winding cell, with integer-certified witnesses}\ }
\]

%dsp-level:0
\[
\textbf{the decay direction is fixed:}
\]

%dsp-level:0
\[
\qquad \text{the parent is the neutral class}\ —\ T(0,S)
\]

%dsp-level:0
\[
= \tfrac43T(+1,S)>T(+1,S)\ [\mathrm{BM.28}(4)(\mathrm{ii});\ \text{machine}]\ —\ \text{and the grammar}\ n\to p+\ell+\nu\ \text{stands unmodified}\ [\mathrm{BM.41}(1);\ \mathrm{BM.45}];
\]

%dsp-level:0
\[
\textbf{the two-body companion row:}
\]

%dsp-level:0
\[
\qquad \text{the two-body channel's admissibility on the specified read is the formed condition}\ \Delta T(S)\ge\rho_\ell\sqrt{G(c(0))}
\]

%dsp-level:0
\[
\qquad [\mathrm{BM.41}(2)\text{'s solution re-derived at}\ m_n\ne m_p:\quad E_p-m_p=\frac{(m_n-m_p)^2-m_\ell^2}{2m_n};\ \text{machine: the cleared identity}]:
\]

%dsp-level:0
\[
\text{at the witness}\ (1,1,1,1)\ \text{it holds}\ \bigl[\tfrac{455}{23}>\tfrac{\sqrt{11}}4\iff1820^2>23^2\cdot11\iff3312400>5819\bigr];
\]

%dsp-level:0
\[
\qquad \mathcal O_S\subseteq\{\text{two-body admissible}\}\ \ \text{since}\ \rho_\nu\sqrt{G_\nu}\ge0\ [\text{machine}];
\]

%dsp-level:0
\[
\textbf{the forced-neutral-member row is read-relative and is recorded,}\ \neg\textbf{adjudicated here:}
\]

%dsp-level:0
\[
\qquad \mathrm{BM.41}(3)\text{'s forcing was derived at the equal rest masses of}\ \mathrm{BM.28}(1)\ —\ \text{it stands on that locus};
\]

%dsp-level:0
\[
\ \ \text{on the specified read the neutral member's emission remains admissible and class-unique}\ [\mathrm{BM.23}(5),(6);\ \mathrm{BM.41}(3)\ \text{grammar}]\ \text{and the named reaction carries it}\ [\mathrm{BM.45}];
\]

%dsp-level:0
\[
\ \ \neg\exists\ \text{a lepton-number law in the inventory}\ [\mathrm{BM.23}(6)];
\]

%dsp-level:0
\[
(7)\quad\textbf{the rescope rows (recorded,}\ \neg\textbf{retracted):}
\]

%dsp-level:0
\[
\qquad \mathrm{BM.41}(4)(\mathrm e)\text{'s verdict}\ \text{(``closed on every adopted read''}):\ \text{at its locus no read was adopted}\ [\mathrm{BM.30}(5)]\ —\ \text{the specification now adopts one, and the verdict is scoped to the pre-specification reads};
\]

%dsp-level:0
\[
\mathrm{BM.49}(1),
\]

%dsp-level:0
\[
(2)\text{'s corrected closure stands in its species}\ —\ \text{the EM summands are}\ \neg\text{parts of the specified read}\ [(3)];
\]

%dsp-level:0
\[
\qquad \mathrm{BM.63}(2)\text{'s closure stands on the informational reads}\ [\mathrm{BM.70}(4)(\mathrm a)];
\]

%dsp-level:0
\[
\mathrm{BM.70}(4)\text{'s dichotomy gains its third case:}
\]

%dsp-level:0
\[
\qquad (\mathrm c)\ \ \text{the specified route-(ii)+transport read: formed, with a branch-relative sign}\ [(5)];
\]

%dsp-level:0
\[
\qquad \mathrm{BM.71}(2)\text{'s no-go is rescoped to the pre-specification reads}\ —\ \text{the fifth route of}\ \mathrm{BM.71}(1)\ \text{is the door walked here, under the logged addendum};
\]

%dsp-level:0
\[
\textbf{no clause of}\ \mathrm{BM.41}\text{--}\mathrm{BM.72}\ \textbf{is retracted; no clause of Parts 1--3 is modified};
\]

%dsp-level:0
\[
\qquad \text{the kinematic viability was always read-relative}\ [\mathrm{BM.41}(7)];
\]

%dsp-level:0
\[
(8)\quad\textbf{the conditional spectrum becomes branch content:}
\]

%dsp-level:0
\[
\qquad \text{on}\ \mathcal O_S\ \text{the hypothesis}\ Q\ge0\ \text{of}\ \mathrm{BM.41}(5)/\mathrm{BM.49}(3)\ \text{is satisfied by a formed read}:
\]

%dsp-level:0
\[
E_\ell^{\max}=\frac{m_n^2+m_\ell^2-(m_p+m_\nu)^2}{2m_n}
\]

%dsp-level:0
\[
\qquad \text{and}
\]

%dsp-level:0
\[
\qquad E_{\max}:\langle E\rangle=2:1\ \ \text{scale-invariant}
\]

%dsp-level:0
\[
\qquad [\mathrm{BM.41}(5);\ \mathrm{BM.49}(3),(4)\ \text{re-cited};\ \text{symbolic in the specified scales}]:
\]

%dsp-level:0
\[
\neg\text{an external endpoint, spectrum or lifetime is compared}\ [\mathrm{P5}];
\]

%dsp-level:0
\[
\qquad \text{the quarantined audit rows stay quarantined}\ [\text{Census FT.114},\ \text{FT.120}\ \text{pattern}];
\]

%dsp-level:0
\[
(9)\quad\textbf{what the specification does not do:}
\]

%dsp-level:0
\[
\qquad \neg\text{a scale is fixed}\ —\ (\rho_\ell,\rho_\nu,g,g',\hat\mu,\hat\mu',\hat\mu_\nu)\ \text{remain undeclared}\ [\mathrm{P4};\ \mathrm{P5}];
\]

%dsp-level:0
\[
\qquad \neg\text{a currency is converted}\ —\ \text{the mass-form bridge is application-layer}\ [(3)];
\]

%dsp-level:0
\[
\neg\text{a Q-value is numerically derived}\ —\ \text{the branch structure is the content};
\]

%dsp-level:0
\[
\qquad \text{the numerical value of}\ e\ \text{remains open and does not enter}\ —\ \text{the specified read carries no EM monomial}\ [\mathrm{EM.25}(7);\ \mathrm{BM.62}(12)];
\]

%dsp-level:0
\[
\text{route}\ (\mathrm{ii})\ \text{is}\ \neg\text{adopted for any purpose beyond the vertex kinematics}\ [\mathrm{BM.30}(5);\ \mathrm{BM.37}(3),(6)];
\]

%dsp-level:0
\[
\qquad \text{the route-(iii) amplitude ratios stay barred as masses}\ [\mathrm{FT.79}(5);\ \mathrm{BM.37}(3)];
\]

%dsp-level:0
\[
(10)\quad\textbf{status:}
\]

%dsp-level:0
\[
\qquad \textbf{derived:}\ (2)\text{--}(8)\ [\text{given}\ (1)];
\]

%dsp-level:0
\[
\qquad \textbf{declared:}\ (1)\ —\ \text{logged as the §II.4 addendum A-BM.8 before use};
\]

%dsp-level:0
\[
\neg\text{a P4 addition}\ (\text{§Cal-BM stays empty};\ \rho_\ell,\rho_\nu\ \text{are Cal-F/P4 restatements}),
\]

%dsp-level:0
\[
\quad\neg\text{a new dimensionful scale}\ [\mathrm{P1}],
\]

%dsp-level:0
\[
\quad\neg\text{a measured value}\ [\mathrm{P5};\ \text{machine: zero external tokens, zero decimal numerals}].
\]

\textbf{Req.} CT2.1, D2.9, EM.25(7) (Part 2); P1, P4, P5 (Parts 1–2); Def-H.37(1), FT.107(3), FT.143 (Part 3); Census FT.114, FT.120 (row-wise); Def-BM.6(1), Def-BM.11(5), BM.20(4), BM.21(2), BM.23(5), BM.23(6), BM.28(1), BM.28(4), BM.30(1), BM.30(5), BM.35(5), BM.37(2), BM.37(3), BM.37(6), BM.39(1), BM.39(3), BM.41(1), BM.41(2), BM.41(3), BM.41(4), BM.41(5), BM.41(7), BM.45, BM.49(1), BM.49(2), BM.49(3), BM.49(4), BM.62(6), BM.62(11), BM.62(12), BM.63(2), BM.66(6), BM.68(8), BM.69(5), BM.70(4), BM.71(1), BM.71(2) (this part); Cal-F; exact rational arithmetic, sum-of-squares positivity, exact radical inequalities by integer certificates, Laurent exponent arithmetic, cleared polynomial identities [math; machine: \texttt{checks/\allowbreak{}verify\_\allowbreak{}H\_\allowbreak{}chiral\_\allowbreak{}assignment\_\allowbreak{}opening.py} §E, §F].

\textbf{Pf.}

%dsp-level:0
\[
(1)\text{--}(2):\ \mathrm{BM.28}(4)(\mathrm{ii})\text{'s frozen factorization}\ T(q,S)=W(q)\overline{\Delta S}(S)\ \text{re-cited};
\]

%dsp-level:0
\[
\qquad \Delta T=(20-15)\overline{\Delta S}=5\,\overline{\Delta S}\ [\text{machine}];
\]

%dsp-level:0
\[
\qquad \overline{\Delta S}(S)=m^2_B(\cdot,S)\ [\mathrm{FT.107}(3):\ \text{the normalized average};\ \mathrm{BM.28}(4)(\mathrm i)\ \text{re-cited}]
\]

%dsp-level:0
\[
\Delta T\ \text{evaluated at the four frozen sector expectations}\ \Bigl(\tfrac{843957}{98956},\ \tfrac{1921662}{393679},\ \tfrac{91}{23},\ \tfrac{13713}{3088}\Bigr)\ [\mathrm{BM.21}(2)/\mathrm{BM.62}(2)\ \text{re-cited}]:
\]

%dsp-level:0
\[
\quad 5\cdot843957=4219785,
\]

%dsp-level:0
\[
\quad 5\cdot1921662=9608310,
\]

%dsp-level:0
\[
\quad 5\cdot91=455,
\]

%dsp-level:0
\[
\quad 5\cdot13713=68565\ [\text{machine: exact integers}];
\]

%dsp-level:0
\[
\qquad \text{the min/max comparisons by cross-multiplication}\ [\text{machine}]
\]

%dsp-level:0
\[
(3):\ \text{the locus rows re-cited}\ [\mathrm{BM.41}(4)(\mathrm d);\ \mathrm{BM.37}(2);\ \mathrm{BM.62}(11)]
\]

%dsp-level:0
\[
(4):\ \text{the SOS rows of}\ \mathrm{BM.39}(1)\ \text{re-derived by expansion}\ [\text{machine:}\ \tfrac38(g\mp\tfrac{g'}3)^2+\tfrac{g'^2}{48}=A_0\mp B_0];
\]

%dsp-level:0
\[
\qquad |\cos\varphi_i|\le1\implies\text{the sandwich}\ [\text{math}];
\]

%dsp-level:0
\[
\qquad \text{line- and channel-independence}\ [\mathrm{BM.39}(3):\ (A_0,B_0)\ \text{shared}]
\]

%dsp-level:0
\[
(5)(\mathrm a):\ \text{the symbolic witness:}\ 2\cdot\tfrac{\Delta T}{4\sqrt{A_0+B_0}}\cdot\sqrt{A_0+B_0}
\]

%dsp-level:0
\[
= \tfrac{\Delta T}2<\Delta T\ [\text{machine: exact in the ordered radical extension}];
\]

%dsp-level:0
\[
\qquad \text{the rational witness:}\ A_0+B_0\big|_{g=g'=1}=\tfrac{11}{16}\ [\text{machine}];
\]

%dsp-level:0
\[
\quad (\tfrac{\sqrt{11}}4)^2
\]

%dsp-level:0
\[
= \tfrac{11}{16};
\]

%dsp-level:0
\[
\quad\tfrac{455}{23}-\tfrac{\sqrt{11}}2>0
\]

%dsp-level:0
\[
\iff 910>23\sqrt{11}
\]

%dsp-level:0
\[
\iff 828100>529\cdot11
\]

%dsp-level:0
\[
= 5819\ [\text{machine: squaring both positive sides}]
\]

%dsp-level:0
\[
(\mathrm b):\ \text{the symbolic witness:}\ 2\cdot\tfrac{\Delta T}{\sqrt{A_0-B_0}}\cdot\sqrt{A_0-B_0}=2\,\Delta T>\Delta T\ [\text{machine}];
\]

%dsp-level:0
\[
\qquad \text{the rational witness:}\ A_0-B_0\big|_{g=g'=1}=\tfrac3{16};
\]

%dsp-level:0
\[
\quad 200\cdot\tfrac{\sqrt3}4=50\sqrt3;\quad\Delta T_{\max}<43\iff4219785<4255108;
\]

%dsp-level:0
\[
\quad 50\sqrt3>43\iff7500>1849\ [\text{machine}]
\]

%dsp-level:0
\[
(\mathrm c):\ \rho^{\mathrm{crit}}=\tfrac{\Delta T}{\sqrt{G_\ell}+\sqrt{G_\nu}}:
\]

%dsp-level:0
\[
\qquad \text{lower bound at}\ \Delta T_{\min},
\]

%dsp-level:0
\[
\ \sqrt{G}\le\tfrac{\sqrt{11}}4:\quad\rho^{\mathrm{crit}}\ge\tfrac{2\cdot455/23}{\sqrt{11}}
\]

%dsp-level:0
\[
= \tfrac{910}{23\sqrt{11}}>\tfrac{910\cdot3}{23\cdot10}
\]

%dsp-level:0
\[
= \tfrac{273}{23}\ [\sqrt{11}<\tfrac{10}3\iff99<100];
\]

%dsp-level:0
\[
\text{upper bound at}\ \Delta T_{\max},
\]

%dsp-level:0
\[
\ \sqrt{G}\ge\tfrac{\sqrt3}4:\quad\rho^{\mathrm{crit}}\le\tfrac{2\cdot4219785/98956}{\sqrt3}
\]

%dsp-level:0
\[
= \tfrac{8439570}{98956\sqrt3}<\tfrac{8439570\cdot5}{98956\cdot8}
\]

%dsp-level:0
\[
= \tfrac{42197850}{791648}\ [\sqrt3>\tfrac85\iff75>64];
\]

%dsp-level:0
\[
\qquad \tfrac{273}{23}>11\iff273>253;\quad\tfrac{42197850}{791648}<54\iff42197850<42748992\ [\text{machine: integers}];
\]

%dsp-level:0
\[
\qquad \text{monotonicity in}\ \rho:\ \text{the coefficient}\ \sqrt{G}>0\ [\text{machine}]
\]

%dsp-level:0
\[
(6):\ T(0,S)=\tfrac43T(+1,S)>T(+1,S)\ [\text{machine}];
\]

%dsp-level:0
\[
\qquad \text{the two-body identity:}\ 2m_n(E_p-m_p)
\]

%dsp-level:0
\[
= (m_n-m_p)^2-m_\ell^2\ \text{at}\ E_p
\]

%dsp-level:0
\[
= \tfrac{m_n^2+m_p^2-m_\ell^2}{2m_n}\ [\text{machine: cleared over}\ \mathbb Q[m_n,m_p,m_\ell]];
\]

%dsp-level:0
\[
\qquad \text{the witness}\ \tfrac{455}{23}>\tfrac{\sqrt{11}}4\iff1820>23\sqrt{11}\iff3312400>5819\ [\text{machine}];
\]

%dsp-level:0
\[
\qquad \mathcal O_S\subseteq\{\Delta T\ge\rho_\ell\sqrt{G_\ell}\}\ [\rho_\nu\sqrt{G_\nu}\ge0;\ \text{machine}]
\]

%dsp-level:0
\[
(7)\text{--}(10):\ \text{the rescope census clause by clause}\ [\text{machine}];
\]

%dsp-level:0
\[
\qquad \mathrm{BM.41}(5),\ \mathrm{BM.49}(3),(4)\ \text{re-cited under the satisfied hypothesis};
\]

%dsp-level:0
\[
\qquad \text{the standing-rule scans}\ [\text{machine}]. \ \blacksquare
\]

\textbf{Phys.} The corpus always knew exactly one door out of its own beta-decay closure: a read specification, an act it priced repeatedly and performed never. This tier performs it, narrowly and on the record — the vertex Q-value is read on the unnormalized total-deficit route for the baryons, where the neutral class outweighs the charged by the exact factor four thirds, and on the transport family form for the emitted lines, whose masses carry undeclared scales between exact sum-of-squares bounds. The resulting Q-value is the frozen expression the decay tier itself formed, and its sign is decided by a single dimensionless scale combination against an exact rational: five times the sector expectation, which is forty-two and change nats at the first frozen cell and never below nineteen. The trichotomy is total and machine-certified — an opening region with a universal symbolic witness and a rational witness certified by eight-two-eight-one-hundred against five thousand eight hundred nineteen, a closing region with its own witnesses, and an exact critical locus enclosed between eleven and fifty-four at the rational coupling point. The beta decay is allowed: open on that region, closed on its complement, degenerate on the boundary, with the parent still the neutral baryon and the grammar untouched. Nothing is fitted and nothing is converted: the scales stay undeclared, the two currencies stay unequated, the informational closures stand in their own species, and every rescoped verdict is recorded rather than erased — the observed world's endpoint may one day be compared in a quarantined row, after a calibration this tier does not perform.

\end{theorem}

\medskip\noindent\textit{\texorpdfstring{BM.78 — The status of the chiral assignment: adjudication and scope}{BM.78 - The status of the chiral assignment: adjudication and scope}.}\ Structural-constraints record: see Appendix B, Section B.1.25 of the Addendum.\par\medskip

\phantomsection
\section[Tier BM.17 — The Mass Ratio Census: the Platonic Winding…]{\texorpdfstring{Tier BM.17 — The Mass Ratio Census: the Platonic Winding Ladders, the Intra- and Inter-Generational Ratios, the Six-State Bridge Census, the Four-Currency Incommensurability, and the Neutron–Proton Ratio}{Tier BM.17 - The Mass Ratio Census: the Platonic Winding Ladders, the Intra- and Inter-Generational Ratios, the Six-State Bridge Census, the Four-Currency Incommensurability, and the Neutron-Proton Ratio}}

\textbf{Locus.} This tier opens no gap and closes none: it takes the inventory of every mass ratio the corpus derives, in each currency separately, and localizes the whole obstruction to one step. It consumes $6$ states $[\mathrm{Def\text{-}BM.14}]$, two hierarchies $[\mathrm{BM.60}(6)]$, one attached splitting $[\mathrm{BM.75}(4)]$, one bridge identity $[\mathrm{BM.69}(2)]$, one exact critical strength $[\mathrm{BM.62}(6)]$, one specified read $[\mathrm{BM.77}]$, and the Platonic ladder cells $[\text{Census FT.63}(4),(6);\ \text{Census FT.76}\ \text{row-wise}]$. Every clause re-derives its own arithmetic in an exact field; $\neg$a clause equates two currencies; $\neg$a clause imports an observed mass, ratio or coupling $[\mathrm{P5}]$. The tier adds no definition: every new symbol is an in-theorem abbreviation of derived objects.

\textbf{Standing rows.} $\alpha_B$ stays barred on $\mathrm{BM.61}(8)$'s four grounds and §Cal-BM stays empty. The relating equation is $\neg$assumed anywhere — item 7 stays open $[\mathrm{BM.35}(5);\ \mathrm{BM.68}(7);\ \mathrm{BM.69}(4),(5);\ \text{pointer}]$, and this tier's contribution to it is a localization, $\neg$a value. $\mathrm{BM.41}(4)(\mathrm e)$, $\mathrm{BM.49}$, $\mathrm{BM.63}$, $\mathrm{BM.70}(4)$, $\mathrm{BM.71}(2)$ stand on the reads formed at their loci, and $\mathrm{BM.77}(5),(6)$'s specified-read trichotomy stands beside them, $\neg$retracted and $\neg$converted. Local notation, fixed for the whole tier: $\mathsf c:=\cos\frac\pi{10}$ with $16\mathsf c^4-20\mathsf c^2+5=0$ — the symbol $c$ keeps its $\mathrm{Def\text{-}BM.12}(4)$ meaning (the priced EM distribution factor, $c\in\{\tfrac12,\tfrac35\}$) at every occurrence; $\chi$ is not used; $z_i^{(\mathrm{ch})}:=(r_i^{(\mathrm{ch})})^2$ and $\gamma_{\mathrm{ch}}:=\sqrt2\cos(3\delta_{\mathrm{ch}})$.

\pseudohead{theorem}{\texorpdfstring{BM.79}{BM.79}}{\texorpdfstring{The Platonic winding census: the ladder law, the Chebyshev tower, and the intra-generational mass ratios}{The Platonic winding census: the ladder law, the Chebyshev tower, and the intra-generational mass ratios}}\label{tag:BM-79}
\textbf{Thm.}

%dsp-level:0
\[
(1)\quad\textbf{the ladder law and the defect angle:}
\]

%dsp-level:0
\[
\qquad \delta_v=2\pi-d\,\frac{(L-2)\pi}L,
\]

%dsp-level:0
\[
\qquad \frac{\delta_v}\pi=\Bigl(1,\ \frac12,\ \frac23,\ \frac15,\ \frac13\Bigr)\ \ \text{at}\ \ (d,L)=(3,3),(3,4),(4,3),(3,5),(5,3)
\]

%dsp-level:0
\[
\qquad [\text{Census FT.42}(5);\ \mathrm{BM.60}(3);\ \text{machine}]:
\]

%dsp-level:0
\[
\boxed{\ V\,\delta_v=4\pi\ \ \text{on all five solids}\ }
\]

%dsp-level:0
\[
\qquad 4\cdot\pi=8\cdot\tfrac\pi2=6\cdot\tfrac{2\pi}3=20\cdot\tfrac\pi5=12\cdot\tfrac\pi3
\]

%dsp-level:0
\[
\qquad [\text{Descartes/Gauss--Bonnet};\ \mathrm{FT.47}(1);\ \text{machine: exact in}\ \pi\text{-units}];
\]

%dsp-level:0
\[
\boxed{\ D(m)\ :=\ 4\sin^2\!\Bigl(\frac{m\,\delta_v}4\Bigr)\ =\ 2-2\cos\Bigl(\frac{m\,\delta_v}2\Bigr)\ =\ 2-2\,T_{|m|}\!\Bigl(\cos\frac{\delta_v}2\Bigr),\qquad m^2_{\mathrm{pole}}(m)=\hat\mu'^2\,D(m)\ }
\]

%dsp-level:0
\[
[\mathrm{Def\text{-}H.37}(2);\ \mathrm{FT.54}(4);\ \mathrm{FT.60}(5):\ \text{the weight ladder at}\ \Theta_R=\delta_v;\ \ \text{machine: the half-angle and Chebyshev identities in each exact field}];
\]

%dsp-level:0
\[
\cos\frac{\delta_v}2\ =\ \Bigl(0,\ \frac{\sqrt2}2,\ \frac12,\ \mathsf c,\ \frac{\sqrt3}2\Bigr)
\]

%dsp-level:0
\[
\qquad [\text{machine}];
\]

%dsp-level:0
\[
\qquad \textbf{the two identically massless readings:}
\]

%dsp-level:0
\[
\qquad \Theta_R=0:\ D\equiv0;
\]

%dsp-level:0
\[
\qquad \Theta_R=V\delta_v=4\pi:\ D(m)=2-2\cos(2m\pi)=0\ \ \forall m
\]

%dsp-level:0
\[
[\text{Census FT.63}(4),(6);\ \text{Census FT.76}(1)\ \text{row-wise};\ \text{machine}]:
\]

%dsp-level:0
\[
\qquad \text{the flat-face side and the global sweep carry no defect mass at any weight};
\]

%dsp-level:0
\[
(2)\quad\textbf{the winding conventions (both priced, neither adopted):}
\]

%dsp-level:0
\[
\qquad \kappa_{\min}=2\ (\text{odd}\ L),\ 1\ (\text{even}\ L),
\]

%dsp-level:0
\[
\qquad k=\frac{L\,\kappa_{\min}}2
\]

%dsp-level:0
\[
\qquad [\mathrm{FT.32}(6);\ \mathrm{FT.45};\ \text{machine}]:
\]

%dsp-level:0
\[
k_f=\frac{K_S}F=(3,\,2,\,3,\,5,\,3)\ \ \text{at}\ \ K_S=(12,12,24,60,60),\ F=(4,6,8,12,20)
\]

%dsp-level:0
\[
\qquad [\text{Census FT.43}(1),\ \text{FT.63}(1)\ \text{row-wise};\ \text{machine}];
\]

%dsp-level:0
\[
k_{vf}=(3,\,3,\,2,\,3,\,5)\ \ \text{on the vertex-figure}\ L_{vf}=(3,3,4,3,5)
\]

%dsp-level:0
\[
\qquad [\text{Census FT.76}(1)\ \text{row-wise};\ \text{machine}];
\]

%dsp-level:0
\[
\textbf{variant}\ (\mathrm A):=(\delta_v,k_f);
\]

%dsp-level:0
\[
\qquad \textbf{variant}\ (\mathrm B):=(\delta_v,k_{vf})
\]

%dsp-level:0
\[
\qquad [\text{Census FT.76}(1):\ \text{both priced};\ \text{the census below reports both}];
\]

%dsp-level:0
\[
(3)\quad\textbf{the variant-}(\mathrm A)\textbf{ multiplets (the census rows):}\qquad m\in\{k_f,k_f{-}2,\dots,-k_f\}:
\]

%tbl-level:0
{\setlength{\tabcolsep}{4pt}
\begin{longtable}{>{\raggedright\arraybackslash}p{\dimexpr0.1404\textwidth-2\tabcolsep\relax}|>{\raggedright\arraybackslash}p{\dimexpr0.0827\textwidth-2\tabcolsep\relax}|>{\raggedright\arraybackslash}p{\dimexpr0.0443\textwidth-2\tabcolsep\relax}|>{\raggedright\arraybackslash}p{\dimexpr0.4576\textwidth-2\tabcolsep\relax}|>{\raggedright\arraybackslash}p{\dimexpr0.2750\textwidth-2\tabcolsep\relax}}
\toprule
solid & $\delta_v/\pi$ & $k_f$ & the $D(m)$ multiset & exact field \\
\addlinespace[2pt]
\midrule
\endfirsthead
\toprule
solid & $\delta_v/\pi$ & $k_f$ & the $D(m)$ multiset & exact field \\
\addlinespace[2pt]
\midrule
\endhead
\bottomrule
\endlastfoot
tetrahedron & $1$ & $3$ & $\{2,\ 2,\ 2,\ 2\}$ at $m=\allowbreak (3,\allowbreak 1,\allowbreak -1,\allowbreak -3)$ & $\mathbb Q$ \\
cube & $\tfrac12$ & $2$ & $\{2,\ \mathbf 0,\ 2\}$ at $m=\allowbreak (2,\allowbreak 0,\allowbreak -2)$ & $\mathbb Q$ \\
octahedron & $\tfrac23$ & $3$ & $\{4,\ 1,\ 1,\ 4\}$ at $m=\allowbreak (3,\allowbreak 1,\allowbreak -1,\allowbreak -3)$ & $\mathbb Q$ \\
dodecahedron & $\tfrac15$ & $5$ & $\{2,\ 2{-}2\cos\tfrac{3\pi}{10},\ 2{-}2\mathsf c,\ 2{-}2\mathsf c,\ 2{-}2\cos\tfrac{3\pi}{10},\ 2\}$ & $\mathbb Q[\mathsf c]/(16\mathsf c^4{-}20\mathsf c^2{+}5)$ \\
icosahedron & $\tfrac13$ & $3$ & $\{2,\ 2{-}\sqrt3,\ 2{-}\sqrt3,\ 2\}$ & $\mathbb Q(\sqrt3)$ \\
\end{longtable}
}

%dsp-level:0
\[
[\text{machine: every entry derived from}\ (1)\text{'s ladder law in its own exact field};\ \text{cross-checked row-wise against Census FT.76}(2)(\mathrm A)];
\]

%dsp-level:0
\[
\textbf{the variant-}(\mathrm B)\textbf{ row (recorded, distinct):}
\]

%dsp-level:0
\[
\qquad \text{tetra}\ \{2,2,2,2\};\quad\text{cube}\ \{2{+}\sqrt2,\ 2{-}\sqrt2,\ 2{-}\sqrt2,\ 2{+}\sqrt2\};\quad\text{octa}\ \{3,\ \mathbf 0,\ 3\};
\]

%dsp-level:0
\[
\text{dodeca}\ \{2{-}2\cos\tfrac{3\pi}{10},\ 2{-}2\mathsf c,\ 2{-}2\mathsf c,\ 2{-}2\cos\tfrac{3\pi}{10}\}\ (k_{vf}{=}3);
\]

%dsp-level:0
\[
\text{icosa}\ \{2{+}\sqrt3,\ 2,\ 2{-}\sqrt3,\ 2{-}\sqrt3,\ 2,\ 2{+}\sqrt3\}\ (k_{vf}{=}5)
\]

%dsp-level:0
\[
\qquad [\text{machine: every entry derived from}\ (1)\text{'s ladder law};\ \text{cross-checked row-wise against Census FT.76}(2)(\mathrm B)];
\]

%dsp-level:0
\[
(4)\quad\textbf{the Chebyshev tower (the dodecahedral row, exact):}
\]

%dsp-level:0
\[
\qquad T_0=1,\quad T_1=\mathsf c,\quad T_{m+1}=2\mathsf c\,T_m-T_{m-1}
\]

%dsp-level:0
\[
\qquad [\text{integer coefficients};\ \text{machine}]:
\]

%tbl-level:0
{\setlength{\tabcolsep}{4pt}
\begin{longtable}{>{\raggedright\arraybackslash}p{\dimexpr0.0575\textwidth-2\tabcolsep\relax}|>{\raggedright\arraybackslash}p{\dimexpr0.2128\textwidth-2\tabcolsep\relax}|>{\raggedright\arraybackslash}p{\dimexpr0.3801\textwidth-2\tabcolsep\relax}|>{\raggedright\arraybackslash}p{\dimexpr0.3497\textwidth-2\tabcolsep\relax}}
\toprule
$m$ & $T_m(\mathsf c)$ reduced & $D(m)=\allowbreak 2-2T_m(\mathsf c)$ & minimal polynomial of $D(m)$ \\
\addlinespace[2pt]
\midrule
\endfirsthead
\toprule
$m$ & $T_m(\mathsf c)$ reduced & $D(m)=\allowbreak 2-2T_m(\mathsf c)$ & minimal polynomial of $D(m)$ \\
\addlinespace[2pt]
\midrule
\endhead
\bottomrule
\endlastfoot
$0$ & $1$ & $0$ & $x$ \\
$1$ & $\mathsf c$ & $2-2\mathsf c$ & $x^4-8x^3+\allowbreak 19x^2-12x+\allowbreak 1$ \\
$2$ & $2\mathsf c^2-1$ & $4-4\mathsf c^2=\allowbreak \tfrac{3-\sqrt5}2$ & $x^2-3x+\allowbreak 1$ \\
$3$ & $4\mathsf c^3-3\mathsf c$ & $2+\allowbreak 6\mathsf c-8\mathsf c^3=\allowbreak 2-2\cos\tfrac{3\pi}{10}$ & $x^4-8x^3+\allowbreak 19x^2-12x+\allowbreak 1$ \\
$4$ & $2\mathsf c^2-\tfrac32$ & $5-4\mathsf c^2=\allowbreak \tfrac{5-\sqrt5}2$ & $x^2-5x+\allowbreak 5$ \\
$5$ & $\mathbf 0$ & $\mathbf 2$ & $x-2$ \\
\end{longtable}
}

%dsp-level:0
\[
\boxed{\ T_5(\mathsf c)=16\mathsf c^5-20\mathsf c^3+5\mathsf c=\mathsf c\,(16\mathsf c^4-20\mathsf c^2+5)=0\ }:
\]

%dsp-level:0
\[
\qquad \textbf{the defining quartic IS the top weight}
\]

%dsp-level:0
\[
[\text{de Moivre}\ \cos5t=16\cos^5t-20\cos^3t+5\cos t\ \text{at}\ t=\tfrac\pi{10}:\ \ \cos\tfrac\pi2=0\iff D(5)=4\sin^2\tfrac\pi4=2;\ \text{machine: exact in}\ \mathbb Z[\mathsf c]/(16\mathsf c^4{-}20\mathsf c^2{+}5)];
\]

%dsp-level:0
\[
\mathsf c^2=\frac{5+\sqrt5}8,
\]

%dsp-level:0
\[
\qquad \sqrt5=8\mathsf c^2-5\in\mathbb Q(\mathsf c)\ \ \text{is the unique quadratic subfield},
\]

%dsp-level:0
\[
\qquad [\mathbb Q(\mathsf c):\mathbb Q]=4
\]

%dsp-level:0
\[
\qquad [\text{Census FT.76}(9)\ \text{row-wise};\ \text{machine}];
\]

%dsp-level:0
\[
(5)\quad\textbf{the ladder structure (all five solids):}
\]

%dsp-level:0
\[
\qquad P_S=\frac{4\pi}{\delta_v}=(4,\,8,\,6,\,20,\,12),
\]

%dsp-level:0
\[
\qquad \#\{\text{distinct values}\}=\frac{P_S}2+1=(3,\,5,\,4,\,11,\,7)
\]

%dsp-level:0
\[
\qquad [\text{Census FT.76}(3)\ \text{row-wise};\ \text{machine}]:
\]

%dsp-level:0
\[
D\ \text{is periodic in}\ m\ \text{with period}\ P_S,\ \text{even in}\ m,\ \text{and strictly increasing on}\ [0,P_S/2];
\]

%dsp-level:0
\[
\qquad \max_mD(m)=4\ \ \text{at}\ \ m=P_S/2\ \ (P_S\ \text{even on all five});
\]

%dsp-level:0
\[
\boxed{\ 0\ \le\ m^2_{\mathrm{pole}}(m)\ \le\ 4\hat\mu'^2\ \ \text{at every weight of every solid}\ }\qquad[\text{Census FT.76}(3);\ \text{machine}];
\]

%dsp-level:0
\[
\textbf{the rung law:}
\]

%dsp-level:0
\[
\qquad D(m)\in\{0,2,4\}\iff m\,\delta_v\in\pi\mathbb Z
\]

%dsp-level:0
\[
\qquad [\text{machine: all five ladders to}\ m=P_S,\ \text{in the exact quadratic rings, with the dodecahedral quartic avoided}];
\]

%dsp-level:0
\[
\textbf{the top-weight row:}
\]

%dsp-level:0
\[
\qquad k_f\delta_v=(3\pi,\ \pi,\ 2\pi,\ \pi,\ \pi)\implies D(k_f)=(2,\ 2,\ 4,\ 2,\ 2)
\]

%dsp-level:0
\[
\qquad [\text{Census FT.76}(4)\ \text{row-wise};\ \text{machine}]:
\]

%dsp-level:0
\[
\textbf{the heaviest weight of four of the five solids is exactly}\ 2\hat\mu'^2,
\]

%dsp-level:0
\[
\ \textbf{the octahedron's being}\ 4\hat\mu'^2\ —\ \textbf{the inter-solid hierarchy lives entirely at the light end};
\]

%dsp-level:0
\[
(6)\quad\textbf{the two structural obstructions (re-certified):}
\]

%dsp-level:0
\[
\qquad (\mathrm a)\ \ \delta_v=\pi\implies D(m)=2-2\cos\frac{m\pi}2\in\{0,2,4\}\ \text{and the}\ k_f{=}3\ \text{multiplet is fully degenerate at}\ 2:
\]

%dsp-level:0
\[
\textbf{the tetrahedron carries no multiplet splitting at all in either variant}
\]

%dsp-level:0
\[
\qquad [\text{Census FT.76}(5)(\mathrm a);\ \text{machine}];
\]

%dsp-level:0
\[
(\mathrm b)\ \ \text{the even-}k\ \text{zero mode:}
\]

%dsp-level:0
\[
\quad k_f=2\ \text{on the cube}\ (\mathrm A)\ \text{and}\ k_{vf}=2\ \text{on the octahedron}\ (\mathrm B)\ \text{each carry}\ m=0\ \text{with}\ D(0)=0\ \textbf{exactly}:
\]

%dsp-level:0
\[
\textbf{on those two multiplets the literal lightest pole mass is}\ 0\ \text{and only lightest-nonzero reads are well defined}
\]

%dsp-level:0
\[
\qquad [\mathrm{FT.60}(5)\text{'s standing zero};\ \text{Census FT.76}(5)(\mathrm b);\ \text{machine}];
\]

%dsp-level:0
\[
(7)\quad\textbf{the intra-generational mass ratios (the multiplet read):}
\]

%dsp-level:0
\[
\qquad r_{\mathrm{intra}}(S)\ :=\ \sqrt{\frac{D_{\max}}{D^{\mathrm{nz}}_{\min}}}
\]

%dsp-level:0
\[
\qquad [D^{\mathrm{nz}}_{\min}:=\text{the lightest nonzero value of the multiplet},\ (6)(\mathrm b)]:
\]

%dsp-level:0
\[
\boxed{\ \textbf{the sine-ratio law:}\qquad r_{\mathrm{intra}}(S)\ =\ \frac{\sin\bigl(k_f\delta_v/4\bigr)}{\sin\bigl(\delta_v/4\bigr)}\ \ \text{on every solid whose multiplet carries}\ m=1\ }
\]

%dsp-level:0
\[
[\text{machine: the half-angle identity}\ D(m)=4\sin^2(m\delta_v/4)\ \text{at the two extreme weights, all five solids, both variants}];
\]

%tbl-level:1
{\small\setlength{\tabcolsep}{4pt}
\begin{longtable}{>{\raggedright\arraybackslash}p{\dimexpr0.1278\textwidth-2\tabcolsep\relax}|>{\raggedright\arraybackslash}p{\dimexpr0.1191\textwidth-2\tabcolsep\relax}|>{\raggedright\arraybackslash}p{\dimexpr0.1891\textwidth-2\tabcolsep\relax}|>{\raggedright\arraybackslash}p{\dimexpr0.1334\textwidth-2\tabcolsep\relax}|>{\raggedright\arraybackslash}p{\dimexpr0.2241\textwidth-2\tabcolsep\relax}|>{\raggedright\arraybackslash}p{\dimexpr0.2066\textwidth-2\tabcolsep\relax}}
\toprule
solid & (A) $r_{\mathrm{intra}}$ & exact form & minimal polynomial & (B) $r_{\mathrm{intra}}$ & exact form \\
\addlinespace[2pt]
\midrule
\endfirsthead
\toprule
solid & (A) $r_{\mathrm{intra}}$ & exact form & minimal polynomial & (B) $r_{\mathrm{intra}}$ & exact form \\
\addlinespace[2pt]
\midrule
\endhead
\bottomrule
\endlastfoot
tetrahedron & $1$ & $1$ (full degeneracy) & $x-1$ & $1$ & $1$ \\
cube & $1$ & $1$ $(+\allowbreak 1$ exact zero mode$)$ & $x-1$ & $1+\allowbreak \sqrt2$ & $\sqrt{\tfrac{2+\sqrt2}{2-\sqrt2}}$ \\
octahedron & $2$ & $\sqrt{4/1}$ & $x-2$ & $1$ & $1$ $(+\allowbreak 1$ zero mode$)$ \\
dodecahedron & $\sqrt{\tfrac{2}{2-2\mathsf c}}$ & $\dfrac{\sin(\pi/4)}{\sin(\pi/20)}$ & $x^8-24x^6+\allowbreak 76x^4-64x^2+\allowbreak 16$ & $\sqrt{\tfrac{2-2\cos(3\pi/10)}{2-2\mathsf c}}$ & $\dfrac{\sin(3\pi/20)}{\sin(\pi/20)}$ \\
icosahedron & $1+\allowbreak \sqrt3$ & $\sqrt{\tfrac{2}{2-\sqrt3}}$ & $x^2-2x-2$ & $2+\allowbreak \sqrt3$ & $\sqrt{\tfrac{2+\sqrt3}{2-\sqrt3}}$ \\
\end{longtable}
}

%dsp-level:0
\[
[\text{machine: the five}\ (\mathrm A)\ \text{rows and the five}\ (\mathrm B)\ \text{rows derived in}\ \mathbb Q,\ \mathbb Q(\sqrt2),\ \mathbb Q(\sqrt3),\ \mathbb Q(\mathsf c);\ \text{the minimal polynomials by the multiplication map and by first linear dependence};\ \text{cross-checked row-wise against Census FT.76}(7)];
\]

%dsp-level:0
\[
\textbf{the nested-radical form of the widest row:}
\]

%dsp-level:0
\[
\qquad r^2_{\mathrm{intra}}(\text{dodeca})=\frac{2}{2-2\mathsf c}=\frac1{1-\mathsf c}=\frac{\bigl(4+\sqrt{10+2\sqrt5}\,\bigr)\bigl(3+\sqrt5\,\bigr)}2
\]

%dsp-level:0
\[
\qquad [\text{machine: exact in}\ \mathbb Q(\sqrt5,\sqrt{10+2\sqrt5})];
\]

%dsp-level:0
\[
(8)\quad\textbf{the exact order of the intra-generational ratios:}\qquad\boxed{\ 1\ =\ 1\ <\ 2\ <\ 1+\sqrt3\ <\ \sqrt{\frac{2}{2-2\mathsf c}}\ }
\]

%dsp-level:0
\[
\text{certificates:}\quad1<2;\qquad2<1+\sqrt3\iff1<\sqrt3\iff1<3;\qquad(1+\sqrt3)^2=4+2\sqrt3<\frac{2}{2-2\mathsf c}
\]

%dsp-level:0
\[
\iff (4+2\sqrt3)(2-2\mathsf c)<2\iff\Bigl(4+\frac72\Bigr)\Bigl(2-\frac{19}{10}\Bigr)
\]

%dsp-level:0
\[
= \frac34<2\ \ \text{on}\ \ \sqrt3<\frac74\ (48<49)\ \text{and}\ \mathsf c>\frac{19}{20}\ \Bigl(5\cdot2500>111^2:\ 12500>12321\Bigr)
\]

%dsp-level:0
\[
[\text{machine: integer certificates and exact dyadic radical enclosures}];
\]

%dsp-level:0
\[
\qquad \textbf{the widest intra-generational ratio is the dodecahedron's, in both variants};
\]

%dsp-level:0
\[
(9)\quad\textbf{the ceiling row (distinct from}\ (7)\textbf{):}
\]

%dsp-level:0
\[
\qquad \sqrt{\frac4{D_S}}\ =\ \Bigl(\sqrt2,\ \ \frac{2}{\sqrt{2-\sqrt2}},\ \ 2,\ \ \frac{2}{\sqrt{2-2\mathsf c}},\ \ \frac{2}{\sqrt{2-\sqrt3}}\Bigr)
\]

%dsp-level:0
\[
\qquad [\text{Census FT.76}(3)\ \text{row-wise};\ \text{machine}]:
\]

%dsp-level:0
\[
\text{maximal at the dodecahedron},
\]

%dsp-level:0
\[
\qquad \frac{4}{2-2\mathsf c}=2\,r^2_{\mathrm{intra}}(\text{dodeca}),
\]

%dsp-level:0
\[
\qquad \text{minimal polynomial}\ \ x^4-48x^3+304x^2-512x+256;
\]

%dsp-level:0
\[
\frac{2}{\sqrt{2-2\mathsf c}}=\sqrt2\cdot r_{\mathrm{intra}}(\text{dodeca})
\]

%dsp-level:0
\[
\qquad [\text{machine: the exact identity}];
\]

%dsp-level:0
\[
\qquad \textbf{no winding choice and no admissible}\ \Theta_R\ \textbf{convention exceeds it}\ \ [\text{Census FT.76}(3);\ \text{the value set is finite}];
\]

%dsp-level:0
\[
(10)\quad\textbf{the range bar, re-certified with the census numbers:}
\]

%dsp-level:0
\[
\qquad \text{the widest mass ratio the Platonic defect species supplies at any winding is}\ \sqrt{4/(2-2\mathsf c)}:
\]

%dsp-level:0
\[
\sqrt{\frac{4}{2-2\mathsf c}}\ <\ 7
\]

%dsp-level:0
\[
\qquad \Bigl[\frac{4}{2-2\mathsf c}<49\iff4<49(2-2\mathsf c)\iff\mathsf c<\frac{47}{49}:\ \mathsf c^2=\frac{5+\sqrt5}8<\frac{2209}{2401}\iff\sqrt5<\frac{3368}{2401}\cdot8-5;\ \text{machine}\Bigr];
\]

%dsp-level:0
\[
\text{and}\ \ 139\ <\ \frac{z_0^{(q)}}{z_1^{(q)}}
\]

%dsp-level:0
\[
\qquad [\mathrm{BM.80}(6):\ \text{the flattest channel hierarchy's outer ratio};\ \text{machine}]\implies\boxed{\ \sqrt{\frac{4}{2-2\mathsf c}}\ <\ 7\ <\ 139\ <\ \frac{z^{(q)}_0}{z^{(q)}_1}\ }
\]

%dsp-level:0
\[
\textbf{the two are dimensionless ratios of different species: their comparison is arithmetic,}\ \neg\textbf{a bridge};
\]

%dsp-level:0
\[
\qquad \text{the widest Platonic mass ratio lies below the flattest inter-generational channel ratio by an exact margin};
\]

%dsp-level:0
\[
\text{deficit}\ 12\text{'s two ceilings stand untouched}\qquad[\text{Census FT.76}(3);\ \mathrm{FT.79}(2);\ \mathrm{FT.84}(2)\ \text{row-wise}];
\]

%dsp-level:0
\[
(11)\quad\textbf{status:}
\]

%dsp-level:0
\[
\qquad \textbf{derived:}\ (1),\ (3)\text{--}(10);
\]

%dsp-level:0
\[
\qquad \textbf{re-cited row-wise as Census data:}\ (2)\text{'s two windings and}\ (9)\text{'s ceiling}\ [\text{Census FT.76}(1),(3),(7)];
\]

%dsp-level:0
\[
\textbf{priced and}\ \neg\textbf{adopted:}\ \text{the variant choice}\ (\mathrm A)/(\mathrm B)\ \text{is}\ \neg\text{a read adoption}\ \ [\mathrm{BM.30}(5)\ \text{stands}];
\]

%dsp-level:0
\[
\qquad \text{the census reports both};
\]

%dsp-level:0
\[
\neg\text{a P4 addition},
\]

%dsp-level:0
\[
\quad\neg\text{a new dimensionful scale}\ [\mathrm{P1}],
\]

%dsp-level:0
\[
\quad\neg\text{a measured value}\ [\mathrm{P5};\ \text{machine: zero external tokens, zero decimal numerals in this tier}],
\]

%dsp-level:0
\[
\quad\text{§Cal-BM stays empty}.
\]

\textbf{Req.} D2.9 (Parts 1–2); P1, P4, P5 (Parts 1–2); Def-H.21(6), Def-H.37(2), FT.32(6), FT.45, FT.47(1), FT.54(4), FT.55(1), FT.60(5), FT.79(2), FT.84(2) (Part 3); Census FT.42(5), FT.43(1), FT.63(1), FT.63(4), FT.63(6), FT.76(1), FT.76(2), FT.76(3), FT.76(4), FT.76(5), FT.76(7), FT.76(9) (row-wise); Def-BM.12(4), Def-BM.14(3), BM.30(5), BM.60(3), BM.61(8) (this part); exact radical rings $\mathbb Q(\sqrt2),\mathbb Q(\sqrt3),\mathbb Q(\sqrt5)$, the quartic ring $\mathbb Z[\mathsf c]/(16\mathsf c^4{-}20\mathsf c^2{+}5)$, Chebyshev recursion with integer coefficients, de Moivre $\cos5t$, minimal polynomials by the multiplication map and by first linear dependence, exact dyadic radical enclosures with integer certificates, half-angle and sine-ratio identities [math; machine: \texttt{checks/\allowbreak{}verify\_\allowbreak{}H\_\allowbreak{}mass\_\allowbreak{}ratio\_\allowbreak{}census.py} §A, §B].

\textbf{Pf.}

%dsp-level:0
\[
(1):\ \delta_v/\pi=2-d(L-2)/L\ \text{at the five}\ (d,L)\ \text{pairs}\ [\text{machine: exact rationals in}\ \pi\text{-units}];
\]

%dsp-level:0
\[
\qquad V\delta_v\ \text{multiplied out at all five}\ [\text{machine}];
\]

%dsp-level:0
\[
\qquad 4\sin^2(x)=2-2\cos(2x)\ \text{at}\ x=m\delta_v/4\ [\text{machine: the exact series identity in each field}];
\]

%dsp-level:0
\[
\qquad \cos(m\delta_v/2)=T_m(\cos(\delta_v/2))\ [\text{Chebyshev};\ \text{machine}]
\]

%dsp-level:0
\[
\Theta_R=4\pi:\ D(m)=2-2\cos(2m\pi)=0\ \forall m\ [\text{machine}];\qquad\Theta_R=0:\ D\equiv0\ [\text{machine}]
\]

%dsp-level:0
\[
(2):\ k=L\kappa_{\min}/2\ \text{at}\ L=(3,4,3,5,3)\implies(3,2,3,5,3);
\]

%dsp-level:0
\[
\qquad L_{vf}=(3,3,4,3,5)\implies(3,3,2,3,5)\ [\text{machine: exact Fractions}];
\]

%dsp-level:0
\[
\qquad K_S/F=(12,12,24,60,60)/(4,6,8,12,20)=(3,2,3,5,3)\ [\text{machine}]
\]

%dsp-level:0
\[
(3):\ \text{the two multiplets evaluated weight by weight from}\ (1)\ \text{in}\ \mathbb Q,
\]

%dsp-level:0
\[
\ \mathbb Q(\sqrt2),
\]

%dsp-level:0
\[
\ \mathbb Q(\sqrt3),
\]

%dsp-level:0
\[
\ \mathbb Q[\mathsf c]\ [\text{machine}];
\]

%dsp-level:0
\[
\qquad \cos\tfrac\pi8=\tfrac{\sqrt{2-\sqrt2}}2\ \text{not needed:}\ D(2)\big|_{\delta_v=\pi/2}=2\ [\text{machine}];
\]

%dsp-level:0
\[
\qquad \text{the}\ (\mathrm B)\ \text{row by the identical pipeline at}\ k_{vf}\ [\text{machine}]
\]

%dsp-level:0
\[
(4):\ T_{m+1}
\]

%dsp-level:0
\[
= 2\mathsf cT_m-T_{m-1}\ \text{iterated to}\ m
\]

%dsp-level:0
\[
= 5\ \text{with reduction}\ 16\mathsf c^4
\]

%dsp-level:0
\[
\mapsto 20\mathsf c^2-5\ [\text{machine: exact integer-coefficient arithmetic in the quotient ring}];
\]

%dsp-level:0
\[
\qquad T_5=16\mathsf c^5-20\mathsf c^3+5\mathsf c=\mathsf c(16\mathsf c^4-20\mathsf c^2+5)=0\ [\text{machine}];
\]

%dsp-level:0
\[
\qquad 4\mathsf c^2-2=\tfrac{\sqrt5-1}2\implies D(2)=4-4\mathsf c^2=\tfrac{3-\sqrt5}2,
\]

%dsp-level:0
\[
\quad D(4)
\]

%dsp-level:0
\[
= 5-4\mathsf c^2
\]

%dsp-level:0
\[
= \tfrac{5-\sqrt5}2\ [\text{machine:}\ 8\mathsf c^2-5=\sqrt5\ \text{verified by squaring both positive sides in}\ \mathbb Q(\sqrt5)]
\]

%dsp-level:0
\[
\text{the minimal polynomials:}\ x^4-8x^3+19x^2-12x+1\ \text{at}\ D(1),D(3);
\]

%dsp-level:0
\[
\quad x^2-3x+1\ \text{at}\ D(2);
\]

%dsp-level:0
\[
\quad x^2-5x+5\ \text{at}\ D(4);
\]

%dsp-level:0
\[
\quad x-2\ \text{at}\ D(5);
\]

%dsp-level:0
\[
\quad x\ \text{at}\ D(0)\ [\text{machine: first linear dependence over}\ \mathbb Q\ \text{on the powers of the element, with exact re-substitution}];
\]

%dsp-level:0
\[
\qquad [\mathbb Q(\mathsf c):\mathbb Q]=4,\ \sqrt5=8\mathsf c^2-5\ [\text{machine}]
\]

%dsp-level:0
\[
(5):\ P_S=4/\!(\delta_v/\pi)=(4,8,6,20,12)\ [\text{machine}];
\]

%dsp-level:0
\[
\qquad D(m+P_S)=D(m)\ \text{and}\ D(-m)=D(m)\ [\text{machine: the ladder evaluated over two full periods}];
\]

%dsp-level:0
\[
\qquad \#\text{distinct}=P_S/2+1\ [\text{machine: the exact value sets}];
\]

%dsp-level:0
\[
\qquad \max_mD(m)=4\ \text{at}\ m=P_S/2\ [\text{machine}]
\]

%dsp-level:0
\[
\text{the rung law}:\ D(m)\in\{0,2,4\}\iff\cos(m\delta_v/2)\in\{1,0,-1\}
\]

%dsp-level:0
\[
\iff m\delta_v\in\pi\mathbb Z\ [\text{machine: both directions on all five ladders}];
\]

%dsp-level:0
\[
\qquad k_f\delta_v=(3\pi,\pi,2\pi,\pi,\pi)\ [\text{machine}]\implies D(k_f)=(2,2,4,2,2)\ [\text{machine}]
\]

%dsp-level:0
\[
(6):\ (\mathrm a)\ \delta_v=\pi\implies D(m)=2-2\cos(m\pi/2)\in\{0,2,4\}\ \text{and}\ D(1)=D(3)=2\ [\text{machine}];
\]

%dsp-level:0
\[
\qquad (\mathrm b)\ D(0)=2-2\cos0=0\ \text{exactly}\ [\text{machine}];
\]

%dsp-level:0
\[
\qquad k_f(\text{cube})=2,\ k_{vf}(\text{octa})=2\ \text{are even}\ [\text{machine}]
\]

%dsp-level:0
\[
(7):\ r_{\mathrm{intra}}
\]

%dsp-level:0
\[
= \sqrt{D(k_f)/D(1)}
\]

%dsp-level:0
\[
= \sqrt{4\sin^2(k_f\delta_v/4)/4\sin^2(\delta_v/4)}
\]

%dsp-level:0
\[
= \sin(k_f\delta_v/4)/\sin(\delta_v/4)\ [\text{machine: positive sines on}\ (0,\pi/2]];
\]

%dsp-level:0
\[
\text{octa}:\ \sqrt{4/1}=2\ [\text{machine}];
\]

%dsp-level:0
\[
\qquad \text{icosa}:\ \frac{2}{2-\sqrt3}=2(2+\sqrt3)=4+2\sqrt3=(1+\sqrt3)^2\ [\text{machine: exact in}\ \mathbb Q(\sqrt3)];
\]

%dsp-level:0
\[
\qquad \text{dodeca}:\ z:
\]

%dsp-level:0
\[
= \frac{2}{2-2\mathsf c}\ \text{satisfies}\ z^4-24z^3+76z^2-64z+16
\]

%dsp-level:0
\[
= 0\ [\text{machine: the substitution}\ D=2/z\ \text{into}\ D^4-8D^3+19D^2-12D+1],
\]

%dsp-level:0
\[
\ \text{so}\ r\ \text{satisfies}\ r^8-24r^6+76r^4-64r^2+16
\]

%dsp-level:0
\[
\text{cube}\ (\mathrm B):\ \frac{2+\sqrt2}{2-\sqrt2}=\frac{(2+\sqrt2)^2}{2}=3+2\sqrt2=(1+\sqrt2)^2\ [\text{machine}];
\]

%dsp-level:0
\[
\qquad \text{icosa}\ (\mathrm B):\ \frac{2+\sqrt3}{2-\sqrt3}=(2+\sqrt3)^2=7+4\sqrt3=(2+\sqrt3)^2\ [\text{machine}]
\]

%dsp-level:0
\[
(8):\ 48<49\implies\sqrt3<\tfrac74;
\]

%dsp-level:0
\[
\qquad 5\cdot2500
\]

%dsp-level:0
\[
= 12500>12321
\]

%dsp-level:0
\[
= 111^2\implies\sqrt5>\tfrac{111}{50}\implies\mathsf c^2
\]

%dsp-level:0
\[
= \tfrac{5+\sqrt5}8>\tfrac{361}{400}\implies\mathsf c>\tfrac{19}{20}\ [\text{machine}];
\]

%dsp-level:0
\[
\qquad (4+\tfrac72)(2-\tfrac{19}{10})=\tfrac{15}2\cdot\tfrac1{10}=\tfrac34<2\ [\text{machine}]
\]

%dsp-level:0
\[
(9):\ w:
\]

%dsp-level:0
\[
= \tfrac{4}{2-2\mathsf c}
\]

%dsp-level:0
\[
= 2z\ \text{and}\ z^4-24z^3+76z^2-64z+16
\]

%dsp-level:0
\[
= 0
\]

%dsp-level:0
\[
\implies w^4-48w^3+304w^2-512w+256
\]

%dsp-level:0
\[
= 0\ [\text{machine: the substitution}\ z=w/2\ \text{cleared},\ \text{and independently by first linear dependence}];
\]

%dsp-level:0
\[
\qquad \tfrac{2}{\sqrt{2-2\mathsf c}}=\sqrt2\sqrt{\tfrac{2}{2-2\mathsf c}}\ [\text{machine: squaring both positive sides}]
\]

%dsp-level:0
\[
(10):\ \tfrac{4}{2-2\mathsf c}<49\iff4<98-98\mathsf c\iff98\mathsf c<94\iff\mathsf c<\tfrac{47}{49}\ [\text{machine}];
\]

%dsp-level:0
\[
\qquad \mathsf c^2
\]

%dsp-level:0
\[
= \tfrac{5+\sqrt5}8<\tfrac{2209}{2401}
\]

%dsp-level:0
\[
\iff 2401(5+\sqrt5)<17672
\]

%dsp-level:0
\[
\iff 2401\sqrt5<5667
\]

%dsp-level:0
\[
\iff 2401^2\cdot5<5667^2\ [\text{machine:}\ 28824005<32114889];
\]

%dsp-level:0
\[
\qquad 7<\tfrac{89}{12}\iff84<89\ [\text{machine}];
\]

%dsp-level:0
\[
\qquad \tfrac{89}{12}<z_0^{(q)}/z_1^{(q)}\ [\mathrm{BM.80}(4)\text{'s exact enclosure};\ \text{machine}]
\]

%dsp-level:0
\[
(11):\ \text{the standing-rule scans}\ [\text{machine}].\ \blacksquare
\]

\textbf{Phys.} One law generates the whole Platonic mass spectrum: the cone defect at a vertex, read as an angle, squared into a mass by a single undeclared scale. Wind that angle by an integer and the ladder appears — periodic, even, monotone to a hard ceiling of four, with the tetrahedron degenerate at every weight and the even windings carrying an exactly massless middle rung. The dodecahedron is the deep row: its tower of Chebyshev polynomials closes at the fifth rung because the defining quartic of the eighteen-degree cosine \textit{is} the statement that five times eighteen degrees is a right angle, so the heaviest weight of that solid sits at exactly two, forced by the same relation that makes its lightest weight the smallest positive defect angle in the census. The intra-generational ratios are then pure sines — one sine over another, at rational multiples of pi — and they come out as one, one, two, root three plus one, and a degree-eight number built from nested square roots of five. Their order is exact and their widest member is the dodecahedron's, in both winding conventions. Nothing here is a mass hierarchy in the channel sense: the widest ratio the entire Platonic ladder can supply at any winding is under seven, while the flattest family channel supplies over one hundred thirty-nine, and the two numbers are compared as arithmetic, never equated as physics — the range bar that deficit twelve records is re-certified with the census's own exact values, and the ceiling that bounds it is a theorem about angles, not about the world.

\pseudohead{theorem}{\texorpdfstring{BM.80}{BM.80}}{\texorpdfstring{The inter-generational mass ratios: one cubic form, three channel constants, the exact enclosures, and the hierarchy-order theorem}{The inter-generational mass ratios: one cubic form, three channel constants, the exact enclosures, and the hierarchy-order theorem}}\label{tag:BM-80}
\textbf{Thm.}

%dsp-level:0
\[
(1)\quad\textbf{the three channel angles and the exact angle ledger:}
\]

%dsp-level:0
\[
\qquad \delta_q=\frac{2}{3\cdot3\cdot3}=\frac2{27},
\]

%dsp-level:0
\[
\qquad \delta_\ell=\frac29,
\]

%dsp-level:0
\[
\qquad \delta_\nu=\frac2{16}=\frac18
\]

%dsp-level:0
\[
\qquad [\mathrm{FT.119}(1);\ \mathrm{FT.112};\ \mathrm{FT.142}(5);\ \text{machine: exact Fractions}]:
\]

%dsp-level:0
\[
\boxed{\ 27\,\delta_q=9\,\delta_\ell=16\,\delta_\nu=2,\qquad3\delta_q=\delta_\ell,\qquad\delta_\nu:\delta_\ell:\delta_q=27:48:16,\qquad\delta_q<\delta_\nu<\delta_\ell\ }
\]

%dsp-level:0
\[
[\mathrm{FT.119}(2);\ \mathrm{FT.143}(5);\ \text{machine}];
\]

%dsp-level:0
\[
\qquad 3\delta_\nu=\frac38\notin\Bigl\{\frac29,\frac2{27}\Bigr\},
\]

%dsp-level:0
\[
\qquad 3\delta_\ell=\frac23,
\]

%dsp-level:0
\[
\qquad 3\delta_q=\frac29
\]

%dsp-level:0
\[
\qquad [\mathrm{FT.143}(2),(5);\ \text{machine}];
\]

%dsp-level:0
\[
(2)\quad\textbf{the spectrum and its cubic (frozen, re-certified):}
\]

%dsp-level:0
\[
\qquad r^{(\mathrm{ch})}_i=1+\sqrt2\cos\Bigl(\delta_{\mathrm{ch}}+\frac{2\pi i}3\Bigr),
\]

%dsp-level:0
\[
\qquad i=0,1,2
\]

%dsp-level:0
\[
\qquad [\mathrm{FT.111}(3);\ \mathrm{FT.113};\ \mathrm{FT.119}(3);\ \mathrm{FT.143}(1);\ \mathrm{FT.158}(4)]:
\]

%dsp-level:0
\[
\text{with}\ c_{\mathrm{ch}}=\cos\delta_{\mathrm{ch}},\ s_{\mathrm{ch}}=\sin\delta_{\mathrm{ch}}:
\]

%dsp-level:0
\[
\qquad r_0=1+\sqrt2\,c_{\mathrm{ch}},
\]

%dsp-level:0
\[
\qquad r_1=1-\frac{\sqrt2\,c_{\mathrm{ch}}+\sqrt6\,s_{\mathrm{ch}}}2,
\]

%dsp-level:0
\[
\qquad r_2=1-\frac{\sqrt2\,c_{\mathrm{ch}}-\sqrt6\,s_{\mathrm{ch}}}2
\]

%dsp-level:0
\[
\qquad [\text{machine}];
\]

%dsp-level:0
\[
\sum_ir_i=3,
\]

%dsp-level:0
\[
\qquad \sum_{i<j}r_ir_j=\frac32,
\]

%dsp-level:0
\[
\qquad \sum_ir_i^2=6,
\]

%dsp-level:0
\[
\qquad \boxed{\ 2(r-1)^3-3(r-1)-\gamma_{\mathrm{ch}}=0,\qquad\gamma_{\mathrm{ch}}:=\sqrt2\cos(3\delta_{\mathrm{ch}})\ }
\]

%dsp-level:0
\[
\qquad [\text{machine: the depressed cubic at}\ B^2=2,\ \text{and the triple angle}\ \cos3\delta=4\cos^3\delta-3\cos\delta];
\]

%dsp-level:0
\[
\text{equivalently, in}\ r\ \text{alone:}
\]

%dsp-level:0
\[
\qquad \boxed{\ 2r^3-6r^2+3r+(1-\gamma_{\mathrm{ch}})=0\ }
\]

%dsp-level:0
\[
\qquad [\text{machine: the expansion of}\ 2(r{-}1)^3{-}3(r{-}1){=}\gamma_{\mathrm{ch}}];
\]

%dsp-level:0
\[
(3)\quad\textbf{the mass-squared cubic (derived here):}
\]

%dsp-level:0
\[
\qquad z:=r^2,
\]

%dsp-level:0
\[
\qquad m_i^{(\mathrm{ch})}=\mu_{\mathrm{ch}}\,z_i^{(\mathrm{ch})},
\]

%dsp-level:0
\[
\qquad \mu_{\mathrm{ch}}\ \text{undeclared}\ [\mathrm{P4};\ \text{Cal-F}]:
\]

%dsp-level:0
\[
\boxed{\ 4z^3-24z^2+\bigl(21-12\,\gamma_{\mathrm{ch}}\bigr)z-\bigl(\gamma_{\mathrm{ch}}-1\bigr)^2\ =\ 0\qquad\text{for each}\ \ z=z_i^{(\mathrm{ch})}\ }
\]

%dsp-level:0
\[
[\text{machine: exact elimination of}\ r\ \text{in}\ \mathbb Q[\gamma][r]/\bigl(2r^3-6r^2+3r+1-\gamma\bigr):\quad r(2z+3)=6z+\gamma-1\ \text{at}\ z=r^2,\ \text{and}\ z(2z+3)^2-(6z+\gamma-1)^2\equiv0\ \text{in the quotient ring}];
\]

%dsp-level:0
\[
\textbf{Vieta:}
\]

%dsp-level:0
\[
\qquad \boxed{\ \sum_iz_i=6\qquad(\text{phase-free and channel-free}),\qquad\sum_{i<j}z_iz_j=\frac{21-12\gamma_{\mathrm{ch}}}4,\qquad\prod_iz_i=\frac{(\gamma_{\mathrm{ch}}-1)^2}4\ }
\]

%dsp-level:0
\[
[\text{machine: the three symmetric reads of the cubic};\ \ \sum_iz_i=6=\sum_ir_i^2\ \text{agrees with}\ (2)\ \text{exactly}];
\]

%dsp-level:0
\[
\textbf{one cubic form, three constants:}
\]

%dsp-level:0
\[
\qquad \text{the polynomial shape is channel-free};
\]

%dsp-level:0
\[
\qquad \text{the entire channel dependence is the single number}\ \gamma_{\mathrm{ch}};
\]

%dsp-level:0
\[
(4)\quad\textbf{the three constants are pairwise distinct and exactly ordered:}
\]

%dsp-level:0
\[
\qquad \gamma_\ell=\sqrt2\cos\frac23,
\]

%dsp-level:0
\[
\qquad \gamma_\nu=\sqrt2\cos\frac38,
\]

%dsp-level:0
\[
\qquad \gamma_q=\sqrt2\cos\frac29:
\]

%dsp-level:0
\[
\boxed{\ \gamma_\ell\ <\ \gamma_\nu\ <\ \gamma_q\ <\ \sqrt2\ }
\]

%dsp-level:0
\[
\qquad [\text{machine: disjoint exact enclosures, width}<2^{-60};\qquad\cos\ \text{strictly decreasing on}\ [0,\pi]\ \text{and}\ \tfrac23>\tfrac38>\tfrac29\ [\text{math}];\qquad\gamma<\sqrt2\iff\cos3\delta<1\iff3\delta\ne0];
\]

%dsp-level:0
\[
\gamma_\ell\in\Bigl(\frac{1111412422687001}{10^{15}},\frac{1111412422687002}{10^{15}}\Bigr),
\]

%dsp-level:0
\[
\quad\gamma_\nu\in\Bigl(\frac{1315936498799931}{10^{15}},\frac{1315936498799932}{10^{15}}\Bigr),
\]

%dsp-level:0
\[
\quad\gamma_q\in\Bigl(\frac{1379438171633255}{10^{15}},\frac{1379438171633256}{10^{15}}\Bigr)
\]

%dsp-level:0
\[
[\text{machine: alternating rational Taylor enclosures of}\ \cos\ \text{with exact remainder bounds, width}<2^{-160}];
\]

%dsp-level:0
\[
\textbf{the three-root regime is exact:}
\]

%dsp-level:0
\[
\qquad 2y^3-3y-\gamma\ \text{has three distinct real roots}\iff\gamma^2<2
\]

%dsp-level:0
\[
\qquad [\text{the discriminant}\ \tfrac{27}4(2-\gamma^2)>0;\ \text{machine}]\ —\ \text{and all three channels satisfy it by}\ (4)\text{'s last inequality};
\]

%dsp-level:0
\[
(5)\quad\textbf{the cross-channel exact identity:}
\]

%dsp-level:0
\[
\qquad \boxed{\ r^{(\ell)}_0\ =\ 1+\gamma_q\ }
\]

%dsp-level:0
\[
\qquad [3\delta_q=\delta_\ell\ [\mathrm{FT.119}(2)];\qquad r^{(\ell)}_0=1+\sqrt2\cos\delta_\ell=1+\sqrt2\cos3\delta_q=1+\gamma_q;\qquad\text{machine: the same cosine, enclosed to width}<2^{-160}]:
\]

%dsp-level:0
\[
\textbf{the}\ \ell\textbf{-channel's heaviest line is the q-channel's cubic constant plus one};
\]

%dsp-level:0
\[
\qquad \text{the shared object is the angle}\ 3\delta_q=\delta_\ell,\ \neg\text{the cubic constant}:
\]

%dsp-level:0
\[
\gamma_\ell=\sqrt2\cos\tfrac23\ne\sqrt2\cos\tfrac29=\gamma_q
\]

%dsp-level:0
\[
\qquad [\text{machine: disjoint enclosures}];
\]

%dsp-level:0
\[
\qquad \gamma_\nu=\sqrt2\cos\tfrac38\ \text{is channel-own}
\]

%dsp-level:0
\[
\qquad [\mathrm{FT.143}(2);\ \text{machine}];
\]

%dsp-level:0
\[
(6)\quad\textbf{the inter-generational ratio census (the exact enclosures):}
\]

%dsp-level:0
\[
\qquad z_1<z_2<z_0\ \ \text{in every channel}
\]

%dsp-level:0
\[
\qquad [\text{machine: disjoint enclosures, widths}<2^{-50}]:
\]

%tbl-level:3
{\scriptsize\setlength{\tabcolsep}{4pt}
\begin{longtable}{>{\raggedright\arraybackslash}p{\dimexpr0.0698\textwidth-2\tabcolsep\relax}|>{\raggedright\arraybackslash}p{\dimexpr0.2552\textwidth-2\tabcolsep\relax}|>{\raggedright\arraybackslash}p{\dimexpr0.2666\textwidth-2\tabcolsep\relax}|>{\raggedright\arraybackslash}p{\dimexpr0.2099\textwidth-2\tabcolsep\relax}|>{\raggedright\arraybackslash}p{\dimexpr0.1985\textwidth-2\tabcolsep\relax}}
\toprule
channel & $z_2/z_1$ & $z_0/z_1$ & $z_2/z_0$ & $z_1/z_0$ \\
\addlinespace[2pt]
\midrule
\endfirsthead
\toprule
channel & $z_2/z_1$ & $z_0/z_1$ & $z_2/z_0$ & $z_1/z_0$ \\
\addlinespace[2pt]
\midrule
\endhead
\bottomrule
\endlastfoot
$q$ & $\bigl(\tfrac{3563695074946}{10^{12}},\allowbreak \tfrac{3563695074947}{10^{12}}\bigr)$ & $\bigl(\tfrac{139338688028416}{10^{12}},\allowbreak \tfrac{139338688028417}{10^{12}}\bigr)$ & $\bigl(\tfrac{25575776012}{10^{12}},\allowbreak \tfrac{25575776013}{10^{12}}\bigr)$ & $\bigl(\tfrac{7176757684}{10^{12}},\allowbreak \tfrac{7176757685}{10^{12}}\bigr)$ \\
$\ell$ & $\bigl(\tfrac{206770315972729}{10^{12}},\allowbreak \tfrac{206770315972730}{10^{12}}\bigr)$ & $\bigl(\tfrac{3477472837104598}{10^{12}},\allowbreak \tfrac{3477472837104599}{10^{12}}\bigr)$ & $\bigl(\tfrac{59459937045}{10^{12}},\allowbreak \tfrac{59459937046}{10^{12}}\bigr)$ & $\bigl(\tfrac{287565150}{10^{12}},\allowbreak \tfrac{287565151}{10^{12}}\bigr)$ \\
$\nu$ & $\bigl(\tfrac{9583926992379}{10^{12}},\allowbreak \tfrac{9583926992380}{10^{12}}\bigr)$ & $\bigl(\tfrac{271994748749143}{10^{12}},\allowbreak \tfrac{271994748749144}{10^{12}}\bigr)$ & $\bigl(\tfrac{35235705970}{10^{12}},\allowbreak \tfrac{35235705971}{10^{12}}\bigr)$ & $\bigl(\tfrac{3676541567}{10^{12}},\allowbreak \tfrac{3676541568}{10^{12}}\bigr)$ \\
\end{longtable}
}

%dsp-level:0
\[
[\text{machine: every bracket a monotone exact interval division at width}<2^{-150},\ \text{each bracket decided by exact}\ 10^{12}\text{-scaling};\qquad\text{the frozen decimal rows of}\ \mathrm{FT.119}(5),(6),\ \mathrm{FT.113},\ \mathrm{FT.143}(3)\ \text{are re-certified digit for digit against these brackets}];
\]

%dsp-level:0
\[
m_1:m_2:m_3=1:\frac{z_2}{z_1}:\frac{z_0}{z_1}
\]

%dsp-level:0
\[
\qquad [\mathrm{FT.119}(5)\ \text{vocabulary};\ \text{the ratios are}\ \mu_{\mathrm{ch}}\text{-free}];
\]

%dsp-level:0
\[
\qquad \textbf{the ratios are}\ \neg\textbf{algebraic numbers:}\ \gamma_{\mathrm{ch}}
\]

%dsp-level:0
\[
= \sqrt2\cos(3\delta_{\mathrm{ch}})\ \text{at the rational radian arguments}\ \tfrac29,
\]

%dsp-level:0
\[
\tfrac38,
\]

%dsp-level:0
\[
\tfrac2{27}\ \text{is a transcendental cosine},
\]

%dsp-level:0
\[
\text{so the exact form of every inter-generational ratio is the closed radical-trigonometric expression of}\ (2),
\]

%dsp-level:0
\[
(3),
\]

%dsp-level:0
\[
\ \text{and its exact comparison content is the rational bracket of}\ (6)
\]

%dsp-level:0
\[
\qquad [\text{machine: the brackets, }\neg\text{a decimal assertion}];
\]

%dsp-level:0
\[
(7)\quad\textbf{the phase-free invariants (all three channels, exact):}
\]

%dsp-level:0
\[
\qquad \boxed{\ Q\ :=\ \frac{\sum_im_i}{\bigl(\sum_i\sqrt{m_i}\bigr)^2}\ \overset{m_i=\mu z_i}{=}\ \frac{\sum_iz_i}{\bigl(\sum_ir_i\bigr)^2}\ =\ \frac69\ =\ \frac23\ }
\]

%dsp-level:0
\[
[\mathrm{FT.119}(4);\ \mathrm{FT.143}(4);\ \mathrm{FT.158}(4);\ \mathrm{Def\text{-}BM.6}(4);\ \text{machine: independent of}\ \delta_{\mathrm{ch}}\ \text{and of}\ \mu_{\mathrm{ch}}];
\]

%dsp-level:0
\[
\qquad \sum_iz_i=6=3A_0\big|_{A_0=\frac{11}6}\ —\ \neg\text{a calibration: the sum rule is}\ \mathrm{BM.60}(8)\text{'s}\ \text{phase-free row};
\]

%dsp-level:0
\[
(8)\quad\textbf{the hierarchy-order theorem (derived, scale-free):}
\]

%dsp-level:0
\[
\qquad \text{on}\ \theta\in(0,\tfrac\pi{12}):
\]

%dsp-level:0
\[
\qquad r_1(\theta)
\]

%dsp-level:0
\[
= 1+\sqrt2\cos(\theta+\tfrac{2\pi}3)\ \text{is strictly decreasing and strictly positive},
\]

%dsp-level:0
\[
\ r_1(\tfrac\pi{12})
\]

%dsp-level:0
\[
= 0\ \text{exactly};
\]

%dsp-level:0
\[
r_2(\theta)\ \text{is strictly increasing};\qquad r_0(\theta)\ \text{is strictly decreasing with}\ r_0/r_1\ \text{strictly increasing}
\]

%dsp-level:0
\[
[\text{machine: the derivative signs}\ -\sqrt2\sin(\theta{+}\tfrac{2\pi}3)<0,\ +\sqrt2\sin(\theta{+}\tfrac\pi3)>0,\ -\sqrt2\sin\theta<0;\qquad\frac{d}{d\theta}\ln\frac{r_0}{r_1}=\sqrt2\Bigl(\frac{\sin(\pi/3-\theta)}{r_1}-\frac{\sin\theta}{r_0}\Bigr)>0
\]

%dsp-level:0
\[
\text{on}\ (0,\tfrac\pi6)\ \text{since}\ \sin(\tfrac\pi3-\theta)>\sin\theta\iff\tan\theta<\tfrac1{\sqrt3}\iff\theta<\tfrac\pi6\ \ \text{and}\ \ \ 0<r_1<r_0
\]

%dsp-level:0
\[
\qquad [\text{math};\ \text{machine: the algebraic reductions}]];
\]

%dsp-level:0
\[
\delta_q<\delta_\nu<\delta_\ell<\frac\pi{12}
\]

%dsp-level:0
\[
\qquad \Bigl[\pi^2>6\sum_{n\le10}\frac1{n^2}=\frac{1968329}{211680}>\frac{64}9\implies\pi>\frac83\implies\frac29<\frac\pi{12};\quad\text{machine:}\ 1968329\cdot9=17714961>13547520=64\cdot211680\Bigr];
\]

%dsp-level:0
\[
\boxed{\ 1\ <\ \frac{z^{(q)}_2}{z^{(q)}_1}\ <\ \frac{z^{(\nu)}_2}{z^{(\nu)}_1}\ <\ \frac{z^{(\ell)}_2}{z^{(\ell)}_1},\qquad1\ <\ \frac{z^{(q)}_0}{z^{(q)}_1}\ <\ \frac{z^{(\nu)}_0}{z^{(\nu)}_1}\ <\ \frac{z^{(\ell)}_0}{z^{(\ell)}_1}\ }
\]

%dsp-level:0
\[
\textbf{the q-channel is the flattest family hierarchy, the}\ \ell\textbf{-channel the steepest, and the}\ \nu\textbf{-channel lies between them}\ —\ \text{an exact ordering with no scale and no import}
\]

%dsp-level:0
\[
\qquad [\text{machine: monotonicity and the disjoint enclosures of}\ (6)\ \text{concur}];
\]

%dsp-level:0
\[
\textbf{the near-vanishing row:}
\]

%dsp-level:0
\[
\qquad r_1(\theta)=0\iff\theta=\tfrac\pi{12}\ \text{exactly}
\]

%dsp-level:0
\[
\qquad [\cos\tfrac{3\pi}4=-\tfrac1{\sqrt2};\ \text{machine}]\implies\frac\pi{12}-\delta_\ell=\frac{3\pi-8}{36}\ \text{is the smallest of the three margins}:
\]

%dsp-level:0
\[
\textbf{the}\ \ell\textbf{-channel's inner line is the closest to the zero locus, and its steepness is a consequence of}\ \delta_\ell\ \textbf{being the largest angle below}\ \tfrac\pi{12};
\]

%dsp-level:0
\[
(9)\quad\textbf{the bars:}
\]

%dsp-level:0
\[
\qquad \mu_q,\ \mu_\ell,\ \mu_\nu\ \text{are undeclared P4/Cal-F scales and enter no ratio}
\]

%dsp-level:0
\[
\qquad [\mathrm{P4};\ \text{Cal-F};\ \text{§Cal-BM stays empty}];
\]

%dsp-level:0
\[
\neg\exists\ \text{an SM mass, mass ratio or Yukawa value in any clause}
\]

%dsp-level:0
\[
\qquad [\mathrm{P5};\ \text{machine: zero external tokens, zero decimal numerals in this tier}];
\]

%dsp-level:0
\[
\text{the}\ m^2\ \text{vocabulary is the corpus's own inertial read}\ [\mathrm{L3}(4),\ \mathrm{L4}];
\]

%dsp-level:0
\[
\qquad \text{the channel labels}\ \ell,q,\nu\ \text{name ledger rows}\ [\mathrm{FT.112}\ \mathrm{Rmk};\ \mathrm{FT.143}(6)];
\]

%dsp-level:0
\[
\qquad \text{amplitude}\ne\text{mass}\ [\mathrm{FT.79}(5)];
\]

%dsp-level:0
\[
\qquad \text{rung}\ne\text{scale}\ [\mathrm{D0.8}];
\]

%dsp-level:0
\[
(10)\quad\textbf{status:}
\]

%dsp-level:0
\[
\qquad \textbf{derived:}\ (1)\text{--}(8);
\]

%dsp-level:0
\[
\qquad \textbf{re-cited:}\ (2)\text{'s spectrum and}\ (7)\text{'s invariant from their own tags, with}\ (3),(5),(8)\ \text{new};
\]

%dsp-level:0
\[
\neg\text{a P4 addition},\quad\neg\text{a new dimensionful scale}\ [\mathrm{P1}],\quad\neg\text{a measured value}\ [\mathrm{P5}],\quad\neg\text{a declaration},\quad
\]

%dsp-level:0
\[
\qquad \textbf{open:}\ \text{the three channel scales}\.
\]

\textbf{Req.} D0.8 (Parts 1–2); P1, P4, P5 (Parts 1–2); L3(4), L4 (Parts 1–2); FT.79(5), FT.111(2), FT.111(3), FT.112, FT.113, FT.119(1), FT.119(2), FT.119(3), FT.119(4), FT.119(5), FT.119(6), FT.142(5), FT.143(1), FT.143(2), FT.143(3), FT.143(4), FT.143(5), FT.143(6), FT.158(4) (Part 3); Def-BM.6(4), BM.60(8), BM.61(8) (this part); Cal-F; exact rational arithmetic, exact symbolic elimination in $\mathbb Q[\gamma][r]/(2r^3{-}6r^2{+}3r{+}1{-}\gamma)$, Vieta reads, alternating rational Taylor enclosures of $\cos$ and $\sin$ with exact remainder bounds, integer Newton certificates for $\sqrt2$ and $\sqrt6$, monotone exact interval division, the Basel-series lower bound for $\pi^2$, derivative-sign reductions in $\mathbb Q[\sqrt2,\sqrt3][c,s]/(c^2{+}s^2{-}1)$ [math; machine: \texttt{checks/\allowbreak{}verify\_\allowbreak{}H\_\allowbreak{}mass\_\allowbreak{}ratio\_\allowbreak{}census.py} §D].

\textbf{Pf.}

%dsp-level:0
\[
(1):\ \delta_q=2/(3\cdot3\cdot3),\ \delta_\ell=2/(3\cdot3),\ \delta_\nu=2/16\ [\text{machine: exact Fractions}];
\]

%dsp-level:0
\[
\qquad 27\cdot\tfrac2{27}=9\cdot\tfrac29=16\cdot\tfrac18=2\ [\text{machine}];
\]

%dsp-level:0
\[
\qquad 3\cdot\tfrac2{27}=\tfrac29\ [\text{machine}];
\]

%dsp-level:0
\[
\qquad \tfrac2{27}<\tfrac18\iff16<27,\quad\tfrac18<\tfrac29\iff9<16\ [\text{machine}];
\]

%dsp-level:0
\[
\qquad \tfrac2{27}:\tfrac18:\tfrac29=16:27:48\ [\times216;\ \text{machine}]
\]

%dsp-level:0
\[
(2):\ r_i=A+B\cos(\delta+\tfrac{2\pi i}3)\ \text{at}\ A=1,\ B=\sqrt2\ [\mathrm{FT.111}(3)\ \text{re-cited}];
\]

%dsp-level:0
\[
\qquad \cos(\delta+\tfrac{2\pi}3)=-\tfrac c2-\tfrac{\sqrt3}2s,\quad\cos(\delta+\tfrac{4\pi}3)=-\tfrac c2+\tfrac{\sqrt3}2s\ [\text{addition formulas}];
\]

%dsp-level:0
\[
\qquad \sqrt2\cdot\tfrac{\sqrt3}2=\tfrac{\sqrt6}2\ [\text{machine}]
\]

%dsp-level:0
\[
\sum_i\cos(\delta+\tfrac{2\pi i}3)
\]

%dsp-level:0
\[
= 0,
\]

%dsp-level:0
\[
\quad\sum_i\cos^2(\delta+\tfrac{2\pi i}3)
\]

%dsp-level:0
\[
= \tfrac32\ [\mathrm{FT.111}(3)\ \text{identities}]\implies\sum_ir_i
\]

%dsp-level:0
\[
= 3,
\]

%dsp-level:0
\[
\ \sum_ir_i^2
\]

%dsp-level:0
\[
= 3+2\cdot\tfrac32
\]

%dsp-level:0
\[
= 6\ [\text{machine}];
\]

%dsp-level:0
\[
\qquad \text{the depressed cubic}\ y^3-\tfrac32y-\tfrac{\sqrt2}2\cos3\delta=0\ \text{at}\ y=r-1,\ \times2\ [\text{machine: the ring identity}]
\]

%dsp-level:0
\[
(3):\ 2(r-1)^3-3(r-1)=\gamma\iff2r^3-6r^2+3r+(1-\gamma)=0\ [\text{machine: expansion}];
\]

%dsp-level:0
\[
\qquad 2r^3+3r=6z-(1-\gamma)\ \text{at}\ z=r^2\implies r(2z+3)=6z+\gamma-1\ [\text{machine}]
\]

%dsp-level:0
\[
z(2z+3)^2-(6z+\gamma-1)^2
\]

%dsp-level:0
\[
= 4z^3-24z^2+(21-12\gamma)z-(\gamma-1)^2\ \text{as an identity in}\ \mathbb Q[\gamma,z]\ [\text{machine: cleared coefficient by coefficient}];
\]

%dsp-level:0
\[
\qquad \text{the three roots are}\ z_i
\]

%dsp-level:0
\[
= r_i^2\ \text{because}\ r_i>0\ \text{and}\ 2z+3
\]

%dsp-level:0
\[
\ne 0\ \text{on}\ z>0\ [\text{machine: disjoint positive enclosures}];
\]

%dsp-level:0
\[
\qquad \text{Vieta}: -\tfrac{-24}4=6,\ \tfrac{21-12\gamma}4,\ \tfrac{(\gamma-1)^2}4\ [\text{machine}]
\]

%dsp-level:0
\[
(4):\ \gamma_\ell,
\]

%dsp-level:0
\[
\gamma_\nu,
\]

%dsp-level:0
\[
\gamma_q\ \text{enclosed by alternating rational Taylor series at}\ \tfrac23,
\]

%dsp-level:0
\[
\tfrac38,
\]

%dsp-level:0
\[
\tfrac29\ \text{with exact remainder bounds, multiplied by the Newton-certified}\ \sqrt2\ [\text{machine: widths}<2^{-160}];
\]

%dsp-level:0
\[
\qquad \text{the three intervals are pairwise disjoint}\ [\text{machine}];
\]

%dsp-level:0
\[
\qquad \tfrac23>\tfrac38\iff16>9,\quad\tfrac38>\tfrac29\iff27>16\ [\text{machine}];
\]

%dsp-level:0
\[
\qquad \text{the discriminant of}\ 2y^3-3y-\gamma:\ \ \Delta
\]

%dsp-level:0
\[
= \tfrac{27}4(2-\gamma^2)>0\iff\gamma^2<2\ [\text{machine: the exact cubic discriminant}]
\]

%dsp-level:0
\[
(5):\ 3\delta_q
\]

%dsp-level:0
\[
= \delta_\ell\ [\mathrm{FT.119}(2)\ \text{re-cited}]\implies\cos3\delta_q
\]

%dsp-level:0
\[
= \cos\delta_\ell\implies\gamma_q
\]

%dsp-level:0
\[
= \sqrt2\cos\delta_\ell
\]

%dsp-level:0
\[
= r^{(\ell)}_0-1\ [\text{machine: the two enclosures agree to width}<2^{-160}\ \text{and the symbolic forms are identical}];
\]

%dsp-level:0
\[
\qquad \gamma_\ell\ne\gamma_q\ [\text{disjoint enclosures}];
\]

%dsp-level:0
\[
\qquad \cos\tfrac38\ne\cos\tfrac29\ [\text{disjoint enclosures}]
\]

%dsp-level:0
\[
(6):\ z_1<z_2<z_0\ [\text{machine: disjoint enclosures of the three roots}];
\]

%dsp-level:0
\[
\qquad \text{the twelve brackets by monotone exact interval division, each numerator and denominator enclosed to width}<2^{-160}\ [\text{machine}];
\]

%dsp-level:0
\[
\qquad \text{the frozen decimal rows of Part 3 re-derived to}\ 12\ \text{places and compared digit for digit}\ [\text{machine}]
\]

%dsp-level:0
\[
(7):\ Q=\tfrac{\sum_iz_i}{(\sum_ir_i)^2}=\tfrac69=\tfrac23\ [\text{machine: exact Fractions}];
\]

%dsp-level:0
\[
\qquad \text{phase-independence: the identities}\ \sum_i\cos(\delta+\tfrac{2\pi i}3)
\]

%dsp-level:0
\[
= 0,
\]

%dsp-level:0
\[
\ \sum_i\cos^2
\]

%dsp-level:0
\[
= \tfrac32\ \text{hold at every}\ \delta\ [\text{machine at six rational phases}]
\]

%dsp-level:0
\[
(8):\ r_1'(\theta)
\]

%dsp-level:0
\[
= -\sqrt2\sin(\theta+\tfrac{2\pi}3)
\]

%dsp-level:0
\[
= -\sqrt2\sin(\tfrac\pi3-\theta)<0\ \text{on}\ (0,\tfrac\pi3)\ [\text{machine: the reduction}\ \sin(\theta+\tfrac{2\pi}3)=\sin(\tfrac\pi3-\theta)];
\]

%dsp-level:0
\[
\qquad r_1(\tfrac\pi{12})=1+\sqrt2\cos\tfrac{3\pi}4=1-1=0\ [\text{machine:}\ \cos\tfrac{3\pi}4=-\tfrac1{\sqrt2}]
\]

%dsp-level:0
\[
r_2'(\theta)=-\sqrt2\sin(\theta+\tfrac{4\pi}3)=+\sqrt2\sin(\theta+\tfrac\pi3)>0\ [\text{machine}];
\]

%dsp-level:0
\[
\qquad r_0'(\theta)=-\sqrt2\sin\theta<0\ \text{on}\ (0,\pi)\ [\text{machine}];
\]

%dsp-level:0
\[
\qquad (\ln(r_0/r_1))'
\]

%dsp-level:0
\[
= \sqrt2\bigl(\tfrac{\sin(\pi/3-\theta)}{r_1}-\tfrac{\sin\theta}{r_0}\bigr)\ [\text{machine: the exact derivative in the quotient ring}]
\]

%dsp-level:0
\[
\sin(\tfrac\pi3-\theta)>\sin\theta\iff\tfrac{\sqrt3}2\cos\theta-\tfrac12\sin\theta>\sin\theta\iff\sqrt3\cos\theta>3\sin\theta\iff\tan\theta<\tfrac1{\sqrt3}\iff\theta<\tfrac\pi6\ [\text{machine: each equivalence as an exact identity;}\ \tan\ \text{strictly increasing on}\ (0,\tfrac\pi2),\ \tan\tfrac\pi6=\tfrac1{\sqrt3}\ [\text{math}]];
\]

%dsp-level:0
\[
r_1<r_0\iff\cos(\theta+\tfrac{2\pi}3)<\cos\theta\ [\cos\ \text{strictly decreasing on}\ (0,\pi);\ \text{math}]
\]

%dsp-level:0
\[
\implies (\ln(r_0/r_1))'>0\ \text{on}\ (0,\tfrac\pi{12})\subset(0,\tfrac\pi6)\ [\text{machine}]
\]

%dsp-level:0
\[
\pi^2
\]

%dsp-level:0
\[
= 6\sum_{n\ge1}n^{-2}>6\sum_{n\le10}n^{-2}
\]

%dsp-level:0
\[
= \tfrac{1968329}{211680}\ [\text{Basel};\ \text{machine: the exact rational partial sum}];
\]

%dsp-level:0
\[
\qquad \tfrac{1968329}{211680}>\tfrac{64}9
\]

%dsp-level:0
\[
\iff 17714961>13547520\ [\text{machine}]\implies\pi>\tfrac83
\]

%dsp-level:0
\[
\implies 3\pi>8\implies\tfrac29<\tfrac\pi{12}\ [\text{machine}]
\]

%dsp-level:0
\[
\text{the ordering}: \delta_q<\delta_\nu<\delta_\ell<\tfrac\pi{12}\ \text{and both ratio maps strictly increasing on}\ (0,\tfrac\pi{12})\implies\text{the boxed chain}\ [\text{machine}];
\]

%dsp-level:0
\[
\qquad \text{concur: the three enclosures of}\ (6)\ \text{are strictly ordered}\ [\text{machine}]
\]

%dsp-level:0
\[
(9)\text{--}(10):\ \text{the standing-rule scans and the assembly}\ [\text{machine}].\ \blacksquare
\]

\textbf{Phys.} Three families, one polynomial. Every channel's mass-squared ratios are the three roots of the same cubic — four z cubed minus twenty-four z squared plus twenty-one minus twelve gamma, times z, minus gamma minus one squared — and the whole difference between the quark channel, the lepton channel and the neutrino channel is a single number: root two times the cosine of three times the channel's own angle. Those three numbers are exact, distinct, and ordered, and one identity ties two of them together: the lepton channel's heaviest line is precisely one more than the quark channel's cubic constant, because tripling the quark angle lands exactly on the lepton angle. The hierarchies then order themselves with no scale anywhere in the argument: below one twelfth of a turn the inner line shrinks, the middle line grows, and both ratios steepen monotonically, so the smallest angle gives the flattest family and the largest angle the steepest — quark, neutrino, lepton — with the lepton channel's inner line the closest to the exact zero at one twelfth of a turn, which is why its hierarchy is two orders of magnitude deep while the quark channel's is barely two. The brackets are exact rational intervals, not decimals: every digit the corpus froze in Part 3 is re-derived here from the angle alone and certified at widths far below the last quoted place. Nothing is fitted, no scale is fixed, and the ratios are what they always were — pure numbers waiting on three undeclared scales before they can be called masses.

\pseudohead{theorem}{\texorpdfstring{BM.81}{BM.81}}{\texorpdfstring{The six-state mass-ratio census: the splitting order, the doublet ratios, the two inverted hierarchies, and the three non-proportional ladders}{The six-state mass-ratio census: the splitting order, the doublet ratios, the two inverted hierarchies, and the three non-proportional ladders}}\label{tag:BM-81}
\textbf{Thm.}

%dsp-level:0
\[
(1)\quad\textbf{the six states in both currencies (the census table):}
\]

%dsp-level:0
\[
\qquad |\mathcal S|=3\cdot2=6,
\]

%dsp-level:0
\[
\qquad m_{\mathrm{hid}}=\frac{2E\ln2}{\mathfrak h},
\]

%dsp-level:0
\[
\qquad m^2_{\mathrm{pole}}=\hat\mu'^2D_S
\]

%dsp-level:0
\[
\qquad [\mathrm{Def\text{-}BM.14}(1),(3);\ \mathrm{FT.44}(4);\ \mathrm{Def\text{-}H.37}(2)]:
\]

%tbl-level:1
{\small\setlength{\tabcolsep}{4pt}
\begin{longtable}{>{\raggedright\arraybackslash}p{\dimexpr0.0717\textwidth-2\tabcolsep\relax}|>{\raggedright\arraybackslash}p{\dimexpr0.0811\textwidth-2\tabcolsep\relax}|>{\raggedright\arraybackslash}p{\dimexpr0.1720\textwidth-2\tabcolsep\relax}|>{\raggedright\arraybackslash}p{\dimexpr0.0811\textwidth-2\tabcolsep\relax}|>{\raggedright\arraybackslash}p{\dimexpr0.1720\textwidth-2\tabcolsep\relax}|>{\raggedright\arraybackslash}p{\dimexpr0.0472\textwidth-2\tabcolsep\relax}|>{\raggedright\arraybackslash}p{\dimexpr0.1282\textwidth-2\tabcolsep\relax}|>{\raggedright\arraybackslash}p{\dimexpr0.1470\textwidth-2\tabcolsep\relax}|>{\raggedright\arraybackslash}p{\dimexpr0.0999\textwidth-2\tabcolsep\relax}}
\toprule
state & $(\mathcal F_g,\allowbreak \beta)$ & the branch datum & $D_S$ & minimal polynomial of $D_S$ & $2E$ & $m_{\mathrm{hid}}\mathfrak h/\ln2$ & $m^2_{\mathrm{pole}}/\hat\mu'^2$ & field \\
\addlinespace[2pt]
\midrule
\endfirsthead
\toprule
state & $(\mathcal F_g,\allowbreak \beta)$ & the branch datum & $D_S$ & minimal polynomial of $D_S$ & $2E$ & $m_{\mathrm{hid}}\mathfrak h/\ln2$ & $m^2_{\mathrm{pole}}/\hat\mu'^2$ & field \\
\addlinespace[2pt]
\midrule
\endhead
\bottomrule
\endlastfoot
$s_1$ & $(\mathcal F_1,\allowbreak -)$ & the ladder rung $m{=}1$ & $2$ & $x-2$ & $12$ & $12$ & $2$ & $\mathbb Q$ \\
$s_2$ & $(\mathcal F_1,\allowbreak +\allowbreak )$ & the ladder rung $m{=}2$ & $4$ & $x-4$ & $12$ & $12$ & $4$ & $\mathbb Q$ \\
$s_3$ & $(\mathcal F_2,\allowbreak -)$ & the cube member & $2-\sqrt2$ & $x^2-4x+\allowbreak 2$ & $24$ & $24$ & $2-\sqrt2$ & $\mathbb Q(\sqrt2)$ \\
$s_4$ & $(\mathcal F_2,\allowbreak +\allowbreak )$ & the octahedron member & $1$ & $x-1$ & $24$ & $24$ & $1$ & $\mathbb Q$ \\
$s_5$ & $(\mathcal F_3,\allowbreak -)$ & the dodecahedron member & $2-2\mathsf c$ & $x^4-8x^3+\allowbreak 19x^2-12x+\allowbreak 1$ & $60$ & $60$ & $2-2\mathsf c$ & $\mathbb Q(\mathsf c)$ \\
$s_6$ & $(\mathcal F_3,\allowbreak +\allowbreak )$ & the icosahedron member & $2-\sqrt3$ & $x^2-4x+\allowbreak 1$ & $60$ & $60$ & $2-\sqrt3$ & $\mathbb Q(\sqrt3)$ \\
\end{longtable}
}

%dsp-level:0
\[
[\mathrm{Def\text{-}BM.14}(2),(3);\ \mathrm{BM.60}(4);\ \text{machine: every minimal polynomial re-derived by first linear dependence over}\ \mathbb Q];
\]

%dsp-level:0
\[
\textbf{the defect species resolves all six states with zero degeneracy};
\]

%dsp-level:0
\[
\qquad \textbf{the informational species is branch-degenerate and takes three values}\ 12:24:60=1:2:5
\]

%dsp-level:0
\[
\qquad [\mathrm{FT.44}(1);\ \mathrm{BM.60}(4),(6);\ \text{machine}];
\]

%dsp-level:0
\[
(2)\quad\textbf{the branch order and the attached splitting (frozen, re-certified):}
\]

%dsp-level:0
\[
\qquad \boxed{\ D_S(\mathcal F_g,+)>D_S(\mathcal F_g,-)\qquad\forall g\in\{1,2,3\}\ }
\]

%tbl-level:0
{\setlength{\tabcolsep}{4pt}
\begin{longtable}{>{\raggedright\arraybackslash}p{\dimexpr0.0949\textwidth-2\tabcolsep\relax}|>{\raggedright\arraybackslash}p{\dimexpr0.1331\textwidth-2\tabcolsep\relax}|>{\raggedright\arraybackslash}p{\dimexpr0.1331\textwidth-2\tabcolsep\relax}|>{\raggedright\arraybackslash}p{\dimexpr0.6389\textwidth-2\tabcolsep\relax}}
\toprule
floor & $D_S(\mathcal F_g,\allowbreak -)$ & $D_S(\mathcal F_g,\allowbreak +\allowbreak )$ & the exact certificate \\
\addlinespace[2pt]
\midrule
\endfirsthead
\toprule
floor & $D_S(\mathcal F_g,\allowbreak -)$ & $D_S(\mathcal F_g,\allowbreak +\allowbreak )$ & the exact certificate \\
\addlinespace[2pt]
\midrule
\endhead
\bottomrule
\endlastfoot
$\mathcal F_1$ & $2$ & $4$ & $4>\allowbreak 2$ \\
$\mathcal F_2$ & $2-\sqrt2$ & $1$ & $1>\allowbreak 2-\sqrt2\iff\allowbreak \sqrt2>\allowbreak 1\iff\allowbreak 2>\allowbreak 1$ \\
$\mathcal F_3$ & $2-2\mathsf c$ & $2-\sqrt3$ & $\mathsf c>\allowbreak \tfrac{\sqrt3}2\iff\allowbreak \mathsf c^2>\allowbreak \tfrac34\iff\allowbreak \tfrac{5+\sqrt5}8>\allowbreak \tfrac34\iff\allowbreak \sqrt5>\allowbreak 1\iff\allowbreak 5>\allowbreak 1$ \\
\end{longtable}
}

%dsp-level:0
\[
\Delta_g\ :=\ m^2_{\mathrm{down}}-m^2_{\mathrm{up}}\ =\ \hat\mu'^2\bigl(D_S(\mathcal F_g,+)-D_S(\mathcal F_g,-)\bigr)
\]

%dsp-level:0
\[
\qquad [\mathrm{Def\text{-}BM.14}(5);\ \mathrm{BM.75}(4);\ \mathrm{Def\text{-}BM.14}(8)\text{'s name rows}]:
\]

%dsp-level:0
\[
\boxed{\ \Delta_1=2\,\hat\mu'^2,\qquad\Delta_2=(\sqrt2-1)\,\hat\mu'^2,\qquad\Delta_3=(2\mathsf c-\sqrt3)\,\hat\mu'^2,\qquad\Delta_g>0\ \ \forall g\ }
\]

%dsp-level:0
\[
\text{minimal polynomials of}\ \Delta_g/\hat\mu'^2:
\]

%dsp-level:0
\[
\qquad x-2,
\]

%dsp-level:0
\[
\qquad x^2+2x-1,
\]

%dsp-level:0
\[
\qquad x^8-22x^6+119x^4-38x^2+1
\]

%dsp-level:0
\[
\qquad [\text{machine: first linear dependence in}\ \mathbb Q,\ \mathbb Q(\sqrt2),\ \mathbb Q(\mathsf c,\sqrt3)];
\]

%dsp-level:0
\[
(3)\quad\textbf{the splitting order (derived here):}
\]

%dsp-level:0
\[
\qquad \boxed{\ \Delta_1\ >\ \Delta_2\ >\ \Delta_3\ >\ 0\ }
\]

%dsp-level:0
\[
\qquad \textbf{the down--up splitting strictly decreases across the generations in the defect currency}
\]

%dsp-level:0
\[
\Delta_1>\Delta_2\iff2>\sqrt2-1\iff3>\sqrt2\iff9>2;\qquad\Delta_2>\Delta_3\iff\sqrt2+\sqrt3-1>2\mathsf c
\]

%dsp-level:0
\[
\iff7+4\sqrt6>4\sqrt2+4\sqrt3+\sqrt5
\]

%dsp-level:0
\[
\qquad [\text{both sides positive; squared}]
\]

%dsp-level:0
\[
\qquad \iff15+6\sqrt6>2\sqrt{10}+2\sqrt{15}
\]

%dsp-level:0
\[
\qquad [\text{squared again}]
\]

%dsp-level:0
\[
\iff341+140\sqrt6>0
\]

%dsp-level:0
\[
\qquad [\text{machine: the three squaring steps as exact identities in}\ \mathbb Q(\sqrt2,\sqrt3,\sqrt5,\sqrt6),\ \text{each side certified positive by exact dyadic enclosure}];
\]

%dsp-level:0
\[
\textbf{the exact ratio census of the splitting:}
\]

%dsp-level:0
\[
\qquad \frac{\Delta_2}{\Delta_1}=\frac{\sqrt2-1}2,
\]

%dsp-level:0
\[
\qquad \frac{\Delta_3}{\Delta_1}=\mathsf c-\frac{\sqrt3}2,
\]

%dsp-level:0
\[
\qquad \frac{\Delta_3}{\Delta_2}=\frac{2\mathsf c-\sqrt3}{\sqrt2-1}=(\sqrt2+1)(2\mathsf c-\sqrt3)
\]

%dsp-level:0
\[
\frac{\Delta_2}{\Delta_1}\in\Bigl(\frac{2071}{10^4},\frac{2072}{10^4}\Bigr),
\]

%dsp-level:0
\[
\qquad \frac{\Delta_3}{\Delta_1}\in\Bigl(\frac{85031}{10^6},\frac{85032}{10^6}\Bigr),
\]

%dsp-level:0
\[
\qquad \frac{\Delta_3}{\Delta_2}\in\Bigl(\frac{410566}{10^6},\frac{410567}{10^6}\Bigr)
\]

%dsp-level:0
\[
\qquad [\text{machine: exact dyadic radical enclosures, width}<2^{-60}];
\]

%dsp-level:0
\[
(4)\quad\textbf{the relative splittings (exact, and one exact simplification):}
\]

%dsp-level:0
\[
\qquad \frac{\Delta_g}{m^2_{\mathrm{up}}(g)}=\frac{D_S(\mathcal F_g,+)}{D_S(\mathcal F_g,-)}-1,
\]

%dsp-level:0
\[
\qquad \frac{\Delta_g}{m^2_{\mathrm{down}}(g)}=1-\frac{D_S(\mathcal F_g,-)}{D_S(\mathcal F_g,+)}:
\]

%dsp-level:0
\[
\boxed{\ \frac{\Delta_1}{m^2_{\mathrm{up}}(1)}=1,\qquad\frac{\Delta_2}{m^2_{\mathrm{up}}(2)}=\frac{\sqrt2}2,\qquad\frac{\Delta_3}{m^2_{\mathrm{up}}(3)}=\frac{2\mathsf c-\sqrt3}{2-2\mathsf c}\ }
\]

%dsp-level:0
\[
\qquad \Bigl[\frac{\sqrt2-1}{2-\sqrt2}=\frac{\sqrt2}2\ \text{exactly}:\ (\sqrt2-1)-\frac{\sqrt2}2(2-\sqrt2)=0\ \ \text{in}\ \mathbb Q(\sqrt2);\ \text{machine}\Bigr];
\]

%dsp-level:0
\[
\frac{\Delta_1}{m^2_{\mathrm{down}}(1)}=\frac12,
\]

%dsp-level:0
\[
\qquad \frac{\Delta_2}{m^2_{\mathrm{down}}(2)}=\sqrt2-1,
\]

%dsp-level:0
\[
\qquad \frac{\Delta_3}{m^2_{\mathrm{down}}(3)}=(2\mathsf c-\sqrt3)(2+\sqrt3)=4\mathsf c+2\sqrt3\,\mathsf c-2\sqrt3-3
\]

%dsp-level:0
\[
\qquad [\text{machine: exact in the biquadratic ring}];
\]

%dsp-level:0
\[
(5)\quad\textbf{the doublet mass ratios (the per-generation down/up census):}
\]

%dsp-level:0
\[
\qquad \boxed{\ \frac{m^2_{\mathrm{down}}(g)}{m^2_{\mathrm{up}}(g)}=\Bigl(2,\ \ \frac{2+\sqrt2}2,\ \ \frac{2-\sqrt3}{2-2\mathsf c}\Bigr)\ }
\]

%dsp-level:0
\[
\text{minimal polynomials:}\qquad x-2,\qquad2x^2-4x+1,\qquad x^8-48x^7+410x^6-1328x^5+1899x^4-1232x^3+330x^2-32x+1
\]

%dsp-level:0
\[
[\text{machine: first linear dependence in}\ \mathbb Q,\ \mathbb Q(\sqrt2),\ \mathbb Q(\mathsf c,\sqrt3);\ \text{the degree-}8\ \text{row's constant term is}\ 1=\mathrm{Norm}(2-\sqrt3)/\mathrm{Norm}(2-2\mathsf c)\ [\text{machine}]];
\]

%dsp-level:0
\[
\boxed{\ \frac{m_{\mathrm{down}}(g)}{m_{\mathrm{up}}(g)}=\Bigl(\sqrt2,\ \ \sqrt{\tfrac{2+\sqrt2}2},\ \ \sqrt{\tfrac{2-\sqrt3}{2-2\mathsf c}}\Bigr)\ }
\]

%dsp-level:0
\[
\qquad \text{annihilators:}\quad x^2-2,\quad2x^4-4x^2+1,\quad P_8(x^2)
\]

%dsp-level:0
\[
P_8(x^2)=x^{16}-48x^{14}+410x^{12}-1328x^{10}+1899x^8-1232x^6+330x^4-32x^2+1
\]

%dsp-level:0
\[
\qquad [\text{machine: the exact substitution};\ \neg\text{minimality claimed at}\ g=3];
\]

%dsp-level:0
\[
\textbf{the sine-ratio form (uniform):}
\]

%dsp-level:0
\[
\qquad \frac{m_{\mathrm{down}}(g)}{m_{\mathrm{up}}(g)}
\]

%dsp-level:0
\[
= \frac{\sin\bigl(\Theta_+/4\bigr)}{\sin\bigl(\Theta_-/4\bigr)}
\]

%dsp-level:0
\[
= \Bigl(\frac{\sin(\pi/2)}{\sin(\pi/4)},\ \ \frac{\sin(\pi/6)}{\sin(\pi/8)},\ \ \frac{\sin(\pi/12)}{\sin(\pi/20)}\Bigr)
\]

%dsp-level:0
\[
\Theta_\pm\ :
\]

%dsp-level:0
\[
= \ m_\pm\,
\]

%dsp-level:0
\[
\delta_v(\mathcal F_g,\pm)\ \ \text{the branch's own angle defect, at}\ \ (m_+,m_-)
\]

%dsp-level:0
\[
= (2,1),
\]

%dsp-level:0
\[
\ (1,1),
\]

%dsp-level:0
\[
\ (1,1)\ \ \text{and}\ \ \delta_v
\]

%dsp-level:0
\[
= \pi;
\]

%dsp-level:0
\[
\ \tfrac{2\pi}3,
\]

%dsp-level:0
\[
\tfrac\pi2;
\]

%dsp-level:0
\[
\ \tfrac\pi3,\qquad \tfrac\pi5
\]

%dsp-level:0
\[
[\mathrm{BM.79}(7)\text{'s law applied to the branch pair};\ \text{machine:}\ D_S=4\sin^2(\Theta/4)\ \text{at each of the six states},\ \text{and}\ \sin^2(\pi/8)=\tfrac{2-\sqrt2}4,\ \sin^2(\pi/12)=\tfrac{2-\sqrt3}4\ \text{by the half-angle identity}];
\]

%dsp-level:0
\[
(6)\quad\textbf{the order of the doublet ratios is}\ \neg\textbf{the generation order:}
\]

%dsp-level:0
\[
\qquad \boxed{\ \frac{m_{\mathrm{down}}(2)}{m_{\mathrm{up}}(2)}\ <\ \frac{m_{\mathrm{down}}(1)}{m_{\mathrm{up}}(1)}\ <\ \frac{m_{\mathrm{down}}(3)}{m_{\mathrm{up}}(3)}\ }
\]

%dsp-level:0
\[
\qquad \text{i.e.}\quad\mathcal F_2\ \text{is the tightest doublet,}\ \mathcal F_3\ \text{the widest}:
\]

%dsp-level:0
\[
\frac{2+\sqrt2}2<2\iff\sqrt2<2\iff2<4;
\]

%dsp-level:0
\[
\qquad 2<\frac{2-\sqrt3}{2-2\mathsf c}\iff4-4\mathsf c<2-\sqrt3\iff\sqrt3<4\mathsf c-2\iff0<16\mathsf c^2-16\mathsf c+1
\]

%dsp-level:0
\[
\iff11+2\sqrt5>16\mathsf c\iff141+44\sqrt5>160+32\sqrt5\iff12\sqrt5>19\iff720>361
\]

%dsp-level:0
\[
\qquad [\text{machine: the chain of exact equivalences, each squaring step certified positive}];
\]

%dsp-level:0
\[
\text{enclosures:}\quad\frac{m^2_{\mathrm{down}}(2)}{m^2_{\mathrm{up}}(2)}\in\Bigl(\frac{17071}{10^4},\frac{17072}{10^4}\Bigr),
\]

%dsp-level:0
\[
\qquad \frac{m^2_{\mathrm{down}}(3)}{m^2_{\mathrm{up}}(3)}\in\Bigl(\frac{2737332}{10^6},\frac{2737333}{10^6}\Bigr)
\]

%dsp-level:0
\[
\qquad [\text{machine}];
\]

%dsp-level:0
\[
(7)\quad\textbf{the two inverted hierarchies (the census form):}
\]

%dsp-level:0
\[
\qquad \textbf{informational:}\quad m_{\mathrm{hid}}(\mathcal F_1)<m_{\mathrm{hid}}(\mathcal F_2)<m_{\mathrm{hid}}(\mathcal F_3),
\]

%dsp-level:0
\[
\qquad 1:2:5
\]

%dsp-level:0
\[
\qquad [\mathrm{BM.60}(6);\ \text{machine}];
\]

%dsp-level:0
\[
\textbf{defect:}
\]

%dsp-level:0
\[
\qquad D_S(\mathcal F_3,\pm)<D_S(\mathcal F_2,\pm)<D_S(\mathcal F_1,\pm),
\]

%dsp-level:0
\[
\qquad 2-2\mathsf c<2-\sqrt3<2-\sqrt2<1<2<4
\]

%dsp-level:0
\[
\qquad [\mathrm{BM.60}(6);\ \text{machine: the certificates}\ 5>1,\ 3>2,\ 2>1];
\]

%dsp-level:0
\[
\boxed{\ \textbf{the two currencies order the three floors oppositely}\ }
\]

%dsp-level:0
\[
\qquad [\mathrm{BM.60}(6);\ \mathrm{BM.39}(5)\ \text{pattern}];
\]

%dsp-level:0
\[
\qquad \textbf{the widest derived ratio in each currency:}
\]

%dsp-level:0
\[
\text{defect}\ (m^2):\ \ \frac{4}{2-2\mathsf c}\in\Bigl(\frac{40863458}{10^6},\frac{40863459}{10^6}\Bigr);
\]

%dsp-level:0
\[
\qquad \text{defect}\ (m):\ \ \frac{2}{\sqrt{2-2\mathsf c}}\in\Bigl(\frac{6392453}{10^6},\frac{6392454}{10^6}\Bigr);
\]

%dsp-level:0
\[
\qquad \text{informational}\ (m):\ \ 5,
\]

%dsp-level:0
\[
\qquad (m^2):\ \ 25
\]

%dsp-level:0
\[
\qquad [\text{machine}];
\]

%dsp-level:0
\[
(8)\quad\textbf{the three ladders are pairwise non-proportional (derived):}
\]

%dsp-level:0
\[
\qquad \textbf{(i) the floor ladder against the channel ladder:}
\]

%dsp-level:0
\[
\neg\exists\,\mu\in\mathbb R_{>0}:\quad\bigl(m_{\mathrm{hid}}(\mathcal F_1),m_{\mathrm{hid}}(\mathcal F_2),m_{\mathrm{hid}}(\mathcal F_3)\bigr)=\mu\bigl(z^{(q)}_1,z^{(q)}_2,z^{(q)}_0\bigr)
\]

%dsp-level:0
\[
[\text{proportionality forces}\ 2=z^{(q)}_2/z^{(q)}_1\ \text{and}\ 5=z^{(q)}_0/z^{(q)}_1;\ \ \mathrm{BM.80}(6)\text{'s brackets give}\ z^{(q)}_2/z^{(q)}_1>\tfrac72>2\ \text{and}\ z^{(q)}_0/z^{(q)}_1>139>5;\ \text{machine}];
\]

%dsp-level:0
\[
\textbf{(ii) the channel ladder against the transport ladder:}
\]

%dsp-level:0
\[
\qquad z^{(\mathrm{ch})}_i=1+2\sqrt2\,\cos\varphi_i+2\cos^2\varphi_i\ \ \text{is quadratic in the channel cosine},
\]

%dsp-level:0
\[
\qquad G\bigl(c(i)\bigr)=A_0-B_0\cos\varphi_i\ \ \text{is affine in it}
\]

%dsp-level:0
\[
[\mathrm{BM.60}(8);\ \mathrm{BM.39}(1);\ \mathrm{FT.111}(3);\ \text{machine}]:
\]

%dsp-level:0
\[
\qquad \text{a quadratic and an affine function of the same three distinct cosines agree at all three only if the quadratic coefficient vanishes};
\]

%dsp-level:0
\[
\text{the quadratic coefficient is exactly}\ 2\ne0
\]

%dsp-level:0
\[
\qquad [\text{machine: the exact polynomial comparison over}\ \mathbb Q[\sqrt2][\cos\varphi]]\implies\boxed{\ \neg\exists\,\lambda:\ z^{(\mathrm{ch})}_i=\lambda\,G\bigl(c(i)\bigr)\ \ \forall i\ };
\]

%dsp-level:0
\[
\textbf{(iii) the floor ladder against the transport ladder:}
\]

%dsp-level:0
\[
\qquad \text{the transport form's line order inverts against the spectrum}\ [\mathrm{BM.39}(5)]\ \text{while the floor ladder is increasing in}\ 2E\ [\mathrm{BM.60}(6)]\implies\neg\text{proportional with a positive scale}
\]

%dsp-level:0
\[
\qquad [\text{machine: the sign census}];
\]

%dsp-level:0
\[
\textbf{the generation map is an order-isomorphism and}\ \neg\textbf{a proportionality}
\]

%dsp-level:0
\[
\qquad [\mathrm{Def\text{-}BM.14}(6);\ \mathrm{BM.60}(7):\ \text{the unique order-isomorphism out of}\ 3!=6\ \text{bijections re-cited}]:
\]

%dsp-level:0
\[
\qquad \text{the map supplies no ratio between the currencies};
\]

%dsp-level:0
\[
(9)\quad\textbf{the six quark states (the assembly row):}
\]

%dsp-level:0
\[
\qquad \mathcal S=\{(\mathcal F_g,\beta)\},
\]

%dsp-level:0
\[
\qquad \text{down-type}:=\beta{=}+\ (\text{the heavier branch}),
\]

%dsp-level:0
\[
\qquad \text{up-type}:=\beta{=}-
\]

%dsp-level:0
\[
\qquad [\mathrm{Def\text{-}BM.14}(8);\ \text{A-BM.6}]:
\]

%dsp-level:0
\[
\textbf{the attachment:}
\]

%dsp-level:0
\[
\qquad q_{\mathrm{core}}=0\longleftrightarrow\text{down-type}\longleftrightarrow\beta{=}+,
\]

%dsp-level:0
\[
\qquad q_{\mathrm{core}}={+}1\longleftrightarrow\text{up-type}\longleftrightarrow\beta{=}-
\]

%dsp-level:0
\[
\qquad [\mathrm{BM.75}(3):\ \text{the unique bijection, conditional on}\ \mathrm{BM.59}(9)\text{'s single-substitution hypothesis, on its face}];
\]

%dsp-level:0
\[
\textbf{the sign-coherence row:}
\]

%dsp-level:0
\[
\qquad \text{in every derived currency in which the doublet splitting is formed, the down-type constituent is the heavier one}
\]

%dsp-level:0
\[
\qquad [\mathrm{BM.75}(5)\ \text{re-cited};\ \text{the four rows:}\ \ln2/\mathfrak h,\ \ln(4/3)/\mathfrak h,\ \tfrac43,\ \Delta_g;\ \text{the single opposed EM correction dominated on the sub-critical branch}\ [\mathrm{BM.46}(4);\ \mathrm{BM.62}(6),(9)]];
\]

%dsp-level:0
\[
\textbf{the flavour composition stays}\ \neg\textbf{derived:}
\]

%dsp-level:0
\[
\qquad \text{which three of the six states a sealed class occupies is}\ \neg\text{attached}\ [\mathrm{BM.28}(7);\ \mathrm{Def\text{-}BM.8}(5);\ \mathrm{BM.75}(6)];
\]

%dsp-level:0
\[
\qquad \text{the mixing matrix is}\ \neg\text{derived}\ [\mathrm{FT.145}(3)];
\]

%dsp-level:0
\[
(10)\quad\textbf{status:}
\]

%dsp-level:0
\[
\qquad \textbf{derived:}\ (1)\text{--}(8);
\]

%dsp-level:0
\[
\qquad \textbf{re-cited:}\ (9)\text{'s attachment, name rows and sign coherence from their own tags};
\]

%dsp-level:0
\[
\qquad \textbf{declared:}\ \mathrm{Def\text{-}BM.14}(8)\text{'s name attachment alone}\ [\text{A-BM.6, logged before use}];
\]

%dsp-level:0
\[
\neg\text{a P4 addition},\quad\neg\text{a new dimensionful scale}\ [\mathrm{P1}],\quad\neg\text{a measured value}\ [\mathrm{P5}],\quad\text{§Cal-BM stays empty};
\]

%dsp-level:0
\[
\qquad \textbf{open:}\ \hat\mu'^2,\ \mu\in\text{Cal-F}\ [\text{the map is an order,}\ \neg\text{a number}].
\]

\textbf{Req.} D0.8 (Parts 1–2); P1, P4, P5 (Parts 1–2); Def-H.37(2), FT.44(1), FT.44(4), FT.111(3), FT.145(3) (Part 3); Census FT.63(6) (row-wise); Def-BM.8(5), Def-BM.12(4), Def-BM.14(1), Def-BM.14(2), Def-BM.14(3), Def-BM.14(5), Def-BM.14(6), Def-BM.14(8), BM.28(7), BM.39(1), BM.39(5), BM.46(4), BM.59(9), BM.60(4), BM.60(6), BM.60(7), BM.60(8), BM.62(6), BM.62(9), BM.75(3), BM.75(4), BM.75(5), BM.75(6), BM.79(7), BM.80(4), BM.80(6) (this part); Cal-F; the declaration addendum A-BM.6 (pointers); exact radical rings $\mathbb Q(\sqrt2),\mathbb Q(\sqrt3),\mathbb Q(\sqrt5),\mathbb Q(\sqrt6),\mathbb Q(\sqrt{10}),\mathbb Q(\sqrt{15})$, the biquadratic ring $\mathbb Q[\mathsf c,\sqrt3]/(16\mathsf c^4{-}20\mathsf c^2{+}5)$, minimal polynomials by first linear dependence with exact re-substitution, norm reads by multiplication determinants, exact dyadic radical enclosures with integer certificates, squaring chains with certified positive sides, exact polynomial comparison over $\mathbb Q[\sqrt2][x]$ [math; machine: \texttt{checks/\allowbreak{}verify\_\allowbreak{}H\_\allowbreak{}mass\_\allowbreak{}ratio\_\allowbreak{}census.py} §C].

\textbf{Pf.}

%dsp-level:0
\[
(1):\ \text{the table re-cited from}\ \mathrm{Def\text{-}BM.14}(3)\ \text{and}\ \mathrm{BM.60}(4)\ \text{and every minimal polynomial re-derived}\ [\text{machine: first linear dependence on the powers of the element in its own field, followed by exact re-substitution}];
\]

%dsp-level:0
\[
\qquad 2E\in\{12,24,24,60,60\}\ \text{has three distinct values}\ [\text{machine}]
\]

%dsp-level:0
\[
(2):\ \text{the three certificates re-derived}\ [\text{machine}];
\]

%dsp-level:0
\[
\qquad \cos\tfrac\pi{10}>\tfrac{\sqrt3}2\iff\mathsf c^2>\tfrac34\iff\tfrac{5+\sqrt5}8>\tfrac34
\]

%dsp-level:0
\[
\iff 5+\sqrt5>6\iff\sqrt5>1
\]

%dsp-level:0
\[
\iff 5>1\ [\text{machine: exact in}\ \mathbb Q(\sqrt5),\ \text{both sides positive}];
\]

%dsp-level:0
\[
\qquad \Delta_g\ \text{formed as exact differences}\ [\text{machine}]
\]

%dsp-level:0
\[
\text{the minimal polynomial of}\ 2\mathsf c-\sqrt3:\ \text{first linear dependence in}\ \mathbb Q[\mathsf c,\sqrt3]\ \text{gives}\ x^8-22x^6+119x^4-38x^2+1\ [\text{machine}];
\]

%dsp-level:0
\[
\qquad \text{the even-power shape follows from the conjugate set}\ \{\pm2\mathsf c_i\mp\sqrt3\}\ \text{being symmetric under}\ x
\]

%dsp-level:0
\[
\mapsto -x\ [\text{machine}]
\]

%dsp-level:0
\[
(3):\ 9>2\ [\text{machine}];
\]

%dsp-level:0
\[
\qquad \sqrt2+\sqrt3-1>0\ \text{and}\ 2\mathsf c>0\ [\text{machine: dyadic enclosures}]\implies\text{the first squaring is order-preserving};
\]

%dsp-level:0
\[
\qquad ( \sqrt2+\sqrt3-1)^2=6+2\sqrt6-2\sqrt2-2\sqrt3\ \text{and}\ (2\mathsf c)^2=4\mathsf c^2=\tfrac{5+\sqrt5}2\ [\text{machine}]
\]

%dsp-level:0
\[
6+2\sqrt6-2\sqrt2-2\sqrt3>\tfrac{5+\sqrt5}2\iff7+4\sqrt6>4\sqrt2+4\sqrt3+\sqrt5\ [\text{machine: cleared}];
\]

%dsp-level:0
\[
\qquad \text{both sides positive}\ [\text{machine}]\implies\text{the second squaring}: \ 145+56\sqrt6>85+32\sqrt6+8\sqrt{10}+8\sqrt{15}
\]

%dsp-level:0
\[
\iff60+24\sqrt6>8\sqrt{10}+8\sqrt{15}\iff15+6\sqrt6>2\sqrt{10}+2\sqrt{15}\ [\text{machine}];
\]

%dsp-level:0
\[
\qquad \text{both sides positive}\ [\text{machine}]\implies\text{the third squaring}:\ 441+180\sqrt6>100+40\sqrt6
\]

%dsp-level:0
\[
\iff 341+140\sqrt6>0\ [\text{machine: positive coefficients}]
\]

%dsp-level:0
\[
\text{the three ratio enclosures by exact dyadic radical isolation}\ [\text{machine: widths}<2^{-60}]
\]

%dsp-level:0
\[
(4):\ (\sqrt2-1)-\tfrac{\sqrt2}2(2-\sqrt2)=\sqrt2-1-(\sqrt2-1)=0\ [\text{machine: exact in}\ \mathbb Q(\sqrt2)];
\]

%dsp-level:0
\[
\qquad \Delta_g/m^2_{\mathrm{up}}=D_S(+)/D_S(-)-1\ \text{and}\ \Delta_g/m^2_{\mathrm{down}}=1-D_S(-)/D_S(+)\ [\text{machine: the two identities}];
\]

%dsp-level:0
\[
\qquad (2\mathsf c-\sqrt3)(2+\sqrt3)=4\mathsf c+2\sqrt3\mathsf c-2\sqrt3-3\ [\text{machine:}\ (2-\sqrt3)(2+\sqrt3)=1]
\]

%dsp-level:0
\[
(5):\ \text{the three}\ m^2\ \text{ratios formed as exact quotients in their fields}\ [\text{machine}];
\]

%dsp-level:0
\[
\qquad \text{the minimal polynomials by first linear dependence, with exact re-substitution}\ [\text{machine}];
\]

%dsp-level:0
\[
\qquad \mathrm{Norm}(2-\sqrt3)=1,\ \mathrm{Norm}(2-2\mathsf c)=1\ \text{as multiplication determinants}\ [\text{machine}]\implies\text{the constant term}\ 1
\]

%dsp-level:0
\[
\text{the annihilators of the mass ratios}:\ x^2-2\ \text{at}\ \sqrt2;
\]

%dsp-level:0
\[
\qquad \text{at}\ y^2=\tfrac{2+\sqrt2}2:\ (2y^2-2)^2=2\iff2y^4-4y^2+1=0\ [\text{machine}];
\]

%dsp-level:0
\[
\qquad P_8(x^2)\ \text{by the exact substitution}\ x\mapsto x^2\ [\text{machine}]
\]

%dsp-level:0
\[
\text{the sine-ratio form}:\ D_S
\]

%dsp-level:0
\[
= 4\sin^2(\delta_v/4)\ \text{at the six states with}\ \delta_v/\pi
\]

%dsp-level:0
\[
= (1,1,\tfrac12,\tfrac23,\tfrac15,\tfrac13)\implies\delta_v/4\in\{\tfrac\pi4,\tfrac\pi4,\tfrac\pi8,\tfrac\pi6,\tfrac\pi{20},\tfrac\pi{12}\}\ [\text{machine}]
\]

%dsp-level:0
\[
(6):\ 2<4\ [\text{machine}];
\]

%dsp-level:0
\[
\qquad 0<16\mathsf c^2-16\mathsf c+1\iff16\cdot\tfrac{5+\sqrt5}8-16\mathsf c+1>0\iff11+2\sqrt5>16\mathsf c\ [\text{machine}];
\]

%dsp-level:0
\[
\qquad \text{both sides positive}\ [\text{machine}]\implies(11+2\sqrt5)^2>256\mathsf c^2=32(5+\sqrt5)
\]

%dsp-level:0
\[
\iff141+44\sqrt5>160+32\sqrt5\iff12\sqrt5>19\iff720>361\ [\text{machine: each step cleared}];
\]

%dsp-level:0
\[
\qquad \text{the two enclosures}\ [\text{machine}]
\]

%dsp-level:0
\[
(7):\ 12<24<60,\quad12:24:60=1:2:5\ [\text{machine}];
\]

%dsp-level:0
\[
\qquad 2-2\mathsf c<2-\sqrt3\iff\sqrt3<2\mathsf c\ [\mathrm{BM.60}(6)\text{'s certificate}\ 5>1];
\]

%dsp-level:0
\[
\quad 2-\sqrt3<2-\sqrt2\iff\sqrt2<\sqrt3\iff2<3;
\]

%dsp-level:0
\[
\quad 2-\sqrt2<1\iff1<\sqrt2\iff1<2\ [\text{machine}]
\]

%dsp-level:0
\[
\text{the widest ratios}:\ \max D_S/\min D_S=4/(2-2\mathsf c),\quad\max2E/\min2E=60/12=5\ [\text{machine}];
\]

%dsp-level:0
\[
\qquad \text{the two enclosures by exact radical isolation}\ [\text{machine}]
\]

%dsp-level:0
\[
(8)(\mathrm i):\ \text{proportionality at three strictly ordered entries forces the two ratio equalities}\ [\text{machine}];
\]

%dsp-level:0
\[
\qquad z^{(q)}_2/z^{(q)}_1>\tfrac72\ \text{and}\ z^{(q)}_0/z^{(q)}_1>139\ [\mathrm{BM.80}(6)\text{'s brackets};\ \text{machine}];
\]

%dsp-level:0
\[
\qquad \tfrac72>2,\ 139>5\ [\text{machine}]
\]

%dsp-level:0
\[
(\mathrm{ii}):\ r_i=1+\sqrt2\cos\varphi_i\implies z_i=r_i^2=1+2\sqrt2\cos\varphi_i+2\cos^2\varphi_i\ [\text{machine: the exact expansion}];
\]

%dsp-level:0
\[
\qquad \text{if}\ 1+2\sqrt2u+2u^2
\]

%dsp-level:0
\[
= \lambda(A_0-B_0u)\ \text{at the three distinct values}\ u
\]

%dsp-level:0
\[
= \cos\varphi_i,
\]

%dsp-level:0
\[
\ \text{the difference is a quadratic with three roots}
\]

%dsp-level:0
\[
\implies\text{the quadratic vanishes identically}
\]

%dsp-level:0
\[
\implies 2
\]

%dsp-level:0
\[
= 0\ [\text{machine: the exact polynomial identity over}\ \mathbb Q[\sqrt2][u],\ \text{with the three cosines' distinctness certified by disjoint enclosures}];
\]

%dsp-level:0
\[
\qquad (\mathrm{iii}):\ \text{the sign census of the two orders}\ [\mathrm{BM.39}(5);\ \text{machine}]
\]

%dsp-level:0
\[
(9)\text{--}(10):\ \text{the rows re-cited from their own tags};\qquad\text{the standing-rule scans}\ [\text{machine}].\ \blacksquare
\]

\textbf{Phys.} Six states, two currencies, and not one proportionality between them. In the angle currency the six defect values are pairwise distinct algebraic numbers living in four fields, the heaviest branch on each floor is heavier by an exact amount, and those three amounts run strictly downward across the generations — two, root two minus one, and two cosines of eighteen degrees minus root three — so the splitting that separates a doublet is widest on the lightest floor and narrowest on the heaviest, the opposite of what a hierarchy reading would suggest. The doublet ratios themselves are pure sine ratios and their order is again not the generation order: the cube–octahedron pair is the tightest, the tetrahedral pair the middle, and the dodecahedron–icosahedron pair the widest by a degree-eight number whose norm is exactly one. In the bit currency the same six states collapse to three values climbing one to two to five. The two orders are opposite, and the tier proves something stronger than that they differ: no positive scale makes the floor ladder proportional to the channel ladder, because two is not three and a half and five is not one hundred thirty-nine; no scale makes the channel ladder proportional to the transport ladder, because one is quadratic in the channel cosine and the other affine and the quadratic coefficient is exactly two; and the transport ladder inverts against both. The generation map survives all of this as what it always was — the unique order-isomorphism between two three-element ordered sets — and it supplies an order, never a number.

\pseudohead{theorem}{\texorpdfstring{BM.82}{BM.82}}{\texorpdfstring{The bridge census: the six exact factors, the intra-currency closure, and the single priced residue}{The bridge census: the six exact factors, the intra-currency closure, and the single priced residue}}\label{tag:BM-82}
\textbf{Thm.}

%dsp-level:0
\[
(1)\quad\textbf{the two reads and the bridge identity (derived, not assumed):}
\]

%dsp-level:0
\[
\qquad m_{\mathrm{hid}}=\frac{2E\ln2}{\mathfrak h}
\]

%dsp-level:0
\[
\qquad [\mathrm{FT.44}(4):\ \text{the taxonomy law}\ m_{\mathrm{hid}}=E_{\mathrm{int}}\ln2/\mathfrak h\ \text{at}\ E_{\mathrm{int}}=2E;\ \text{Census FT.42}(1)\ \text{row-wise};\ \mathrm{Def\text{-}BM.14}(1)]:
\]

%dsp-level:0
\[
m^2_{\mathrm{pole}}=\hat\mu'^2D_S=4\hat\mu'^2\sin^2\frac{\Theta_R}4\ \ \text{at}\ \Theta_R=\delta_v
\]

%dsp-level:0
\[
\qquad [\mathrm{Def\text{-}H.37}(2);\ \mathrm{FT.54}(4);\ \mathrm{FT.47}(1);\ \mathrm{BM.60}(3)];
\]

%dsp-level:0
\[
\rho_{\mathrm B}\ :=\ \frac{\hat\mu'^2\mathfrak h^2}{(\ln2)^2}
\]

%dsp-level:0
\[
\qquad [\mathrm{BM.69}(2);\ \text{Cal-F packaging}];
\]

%dsp-level:0
\[
\qquad [\rho_{\mathrm B}]=\mathfrak q\cdot\mathfrak q^{-1}=\mathfrak q^0
\]

%dsp-level:0
\[
\qquad [[\hat\mu'^2]=\mathfrak q\ [\mathrm{Def\text{-}H.37}(2)];\ [\mathfrak h]=\mathfrak q^{-1/2}\ [\mathrm{CT2.1};\ \mathrm{D2.9}];\ \text{machine: the exponent arithmetic}];
\]

%dsp-level:0
\[
\boxed{\ \frac{m^2_{\mathrm{pole}}(s)}{m^2_{\mathrm{hid}}(s)}\ =\ \rho_{\mathrm B}\cdot\mathcal R_s,\qquad\mathcal R_s\ :=\ \frac{D_S(s)}{4\,E(s)^2}\ \in\ \overline{\mathbb Q}\ }
\]

%dsp-level:0
\[
\qquad [\mathrm{BM.69}(2)\ \text{re-derived from the two reads above};\ \text{machine: one Laurent monomial in}\ \hat\mu'^2,\mathfrak h,\ln2]:
\]

%dsp-level:0
\[
\textbf{the factorization is exact per state};
\]

%dsp-level:0
\[
\qquad \textbf{the algebraic part}\ \mathcal R_s\ \textbf{is derived};
\]

%dsp-level:0
\[
\qquad \textbf{the single dimensionless residue is}\ \rho_{\mathrm B}
\]

%dsp-level:0
\[
\qquad [\mathrm{BM.69}(5)(\mathrm{R1})];
\]

%dsp-level:0
\[
(2)\quad\textbf{the six-factor census:}\qquad E(s)=(6,6,12,12,30,30)\ \ \text{at}\ \ 2E=(12,12,24,24,60,60):
\]

%tbl-level:0
{\setlength{\tabcolsep}{4pt}
\begin{longtable}{>{\raggedright\arraybackslash}p{\dimexpr0.0859\textwidth-2\tabcolsep\relax}|>{\raggedright\arraybackslash}p{\dimexpr0.0971\textwidth-2\tabcolsep\relax}|>{\raggedright\arraybackslash}p{\dimexpr0.0971\textwidth-2\tabcolsep\relax}|>{\raggedright\arraybackslash}p{\dimexpr0.1649\textwidth-2\tabcolsep\relax}|>{\raggedright\arraybackslash}p{\dimexpr0.1197\textwidth-2\tabcolsep\relax}|>{\raggedright\arraybackslash}p{\dimexpr0.4353\textwidth-2\tabcolsep\relax}}
\toprule
state & $D_S(s)$ & $4E(s)^2$ & $\mathcal R_s=\allowbreak D_S/(4E^2)$ & field & minimal polynomial of $\mathcal R_s$ \\
\addlinespace[2pt]
\midrule
\endfirsthead
\toprule
state & $D_S(s)$ & $4E(s)^2$ & $\mathcal R_s=\allowbreak D_S/(4E^2)$ & field & minimal polynomial of $\mathcal R_s$ \\
\addlinespace[2pt]
\midrule
\endhead
\bottomrule
\endlastfoot
$s_1$ & $2$ & $144$ & $\tfrac1{72}$ & $\mathbb Q$ & $72x-1$ \\
$s_2$ & $4$ & $144$ & $\tfrac1{36}$ & $\mathbb Q$ & $36x-1$ \\
$s_3$ & $2-\sqrt2$ & $576$ & $\tfrac{2-\sqrt2}{576}$ & $\mathbb Q(\sqrt2)$ & $165888x^2-1152x+\allowbreak 1$ \\
$s_4$ & $1$ & $576$ & $\tfrac1{576}$ & $\mathbb Q$ & $576x-1$ \\
$s_5$ & $2-2\mathsf c$ & $3600$ & $\tfrac{2-2\mathsf c}{3600}$ & $\mathbb Q(\mathsf c)$ & $167961600000000\,\allowbreak x^4-373248000000\,\allowbreak x^3+\allowbreak 246240000\,\allowbreak x^2-43200\,\allowbreak x+\allowbreak 1$ \\
$s_6$ & $2-\sqrt3$ & $3600$ & $\tfrac{2-\sqrt3}{3600}$ & $\mathbb Q(\sqrt3)$ & $12960000\,\allowbreak x^2-14400\,\allowbreak x+\allowbreak 1$ \\
\end{longtable}
}

%dsp-level:0
\[
\boxed{\ \text{six states, six exact factors, four fields, six minimal polynomials}\ }
\]

%dsp-level:0
\[
\qquad [\mathrm{BM.69}(2)\text{'s row re-derived};\ \text{machine: the quartic row by the exact substitution}\ D=3600x\ \text{into}\ x^4-8x^3+19x^2-12x+1];
\]

%dsp-level:0
\[
\textbf{the exact order:}
\]

%dsp-level:0
\[
\qquad \mathcal R_5<\mathcal R_6<\mathcal R_3<\mathcal R_4<\mathcal R_1<\mathcal R_2
\]

%dsp-level:0
\[
\qquad [\text{machine: pairwise comparisons in the exact fields with dyadic isolation}];
\]

%dsp-level:0
\[
\mathcal R_5<\mathcal R_6\iff2-2\mathsf c<2-\sqrt3\iff5>1;
\]

%dsp-level:0
\[
\qquad \mathcal R_6<\mathcal R_3\iff3600^{-1}(2-\sqrt3)<576^{-1}(2-\sqrt2)\iff2-\sqrt3<\tfrac{25}4(2-\sqrt2)
\]

%dsp-level:0
\[
\qquad [\text{machine}];
\]

%dsp-level:0
\[
\mathcal R_3<\mathcal R_4\iff2-\sqrt2<1\iff1<2;
\]

%dsp-level:0
\[
\qquad \mathcal R_4<\mathcal R_1\iff\frac1{576}<\frac1{72}\iff72<576;
\]

%dsp-level:0
\[
\qquad \mathcal R_1<\mathcal R_2\iff\frac1{72}<\frac1{36}\iff36<72
\]

%dsp-level:0
\[
\qquad [\text{machine: integers}];
\]

%dsp-level:0
\[
(3)\quad\textbf{the span of the bridge factor:}
\]

%dsp-level:0
\[
\qquad \boxed{\ \frac{\mathcal R_{\max}}{\mathcal R_{\min}}=\frac{\mathcal R_2}{\mathcal R_5}=\frac{1/36}{(2-2\mathsf c)/3600}=\frac{100}{2-2\mathsf c}=\frac{50}{1-\mathsf c}\ }
\]

%dsp-level:0
\[
\text{minimal polynomial:}\quad x^4-1200x^3+190000x^2-8000000x+100000000
\]

%dsp-level:0
\[
\qquad [\text{machine: the substitution}\ z=x/50\ \text{into}\ \mathrm{BM.79}(7)\text{'s}\ z^4-24z^3+76z^2-64z+16];
\]

%dsp-level:0
\[
\frac{\mathcal R_2}{\mathcal R_5}\in\Bigl(\frac{1021586454}{10^6},\frac{1021586455}{10^6}\Bigr),
\]

%dsp-level:0
\[
\qquad \sqrt{\frac{\mathcal R_2}{\mathcal R_5}}=\frac{10}{\sqrt{2-2\mathsf c}}\in\Bigl(\frac{31962266}{10^6},\frac{31962267}{10^6}\Bigr)
\]

%dsp-level:0
\[
\qquad [\text{machine: exact radical enclosures}];
\]

%dsp-level:0
\[
\textbf{the floor factor is exactly twenty-five:}
\]

%dsp-level:0
\[
\qquad \frac{\mathcal R_2/\mathcal R_5}{4/(2-2\mathsf c)}=\frac{100/(2-2\mathsf c)}{4/(2-2\mathsf c)}=25=\Bigl(\frac{30}6\Bigr)^2=\Bigl(\frac{E(\mathcal F_3)}{E(\mathcal F_1)}\Bigr)^2
\]

%dsp-level:0
\[
\qquad [\text{machine: the exact quotient}]:
\]

%dsp-level:0
\[
\textbf{the bridge factor's range exceeds the defect ceiling by exactly the square of the floor ladder's edge ratio}\ —\ \text{the angle contributes}\ \tfrac{4}{2-2\mathsf c},
\]

%dsp-level:0
\[
\ \text{the floor contributes}\ 25;
\]

%dsp-level:0
\[
(4)\quad\textbf{the intra-currency closure (the residue-free part of the bridge):}
\]

%dsp-level:0
\[
\qquad \textbf{the pole currency:}\quad\frac{m^2_{\mathrm{pole}}(s)}{m^2_{\mathrm{pole}}(s')}=\frac{D_S(s)}{D_S(s')}
\]

%dsp-level:0
\[
\qquad \text{for all}\ 30\ \text{ordered pairs of distinct states}\ (15\ \text{unordered pairs, each read in both directions}):
\]

%dsp-level:0
\[
\boxed{\ \text{all thirty are exact algebraic and}\ \rho_{\mathrm B}\textbf{-free}\ }
\]

%dsp-level:0
\[
\qquad [\text{machine:}\ \hat\mu'^2\ \text{cancels in every quotient};\ \text{the values lie in the compositum}\ \mathbb Q(\sqrt2,\sqrt3,\mathsf c),\ [\,\cdot:\mathbb Q]=16\ [\mathrm{FT.73}(2)\ \text{row-wise}]];
\]

%dsp-level:0
\[
\textbf{two exact cross-links to the ladder census:}
\]

%dsp-level:0
\[
\qquad \frac{D_S(s_1)}{D_S(s_6)}=\frac{2}{2-\sqrt3}=4+2\sqrt3=(1+\sqrt3)^2=r^2_{\mathrm{intra}}(\text{icosa}),
\]

%dsp-level:0
\[
\qquad \frac{D_S(s_2)}{D_S(s_5)}=\frac{4}{2-2\mathsf c}=\text{the dodecahedral ceiling}
\]

%dsp-level:0
\[
[\mathrm{BM.79}(7),(9);\ \mathrm{BM.81}(5);\ \text{machine: the exact identities}];
\]

%dsp-level:0
\[
\qquad \text{the widest and narrowest of the thirty are}\ \tfrac{4}{2-2\mathsf c}\ \text{and its reciprocal}\ \tfrac{2-2\mathsf c}4\ [\text{machine}];
\]

%dsp-level:0
\[
\textbf{the collision census:}
\]

%dsp-level:0
\[
\qquad \text{the thirty ordered ratios take twenty-eight distinct values, and the unique collision is the exact rational}\ 2:
\]

%dsp-level:0
\[
\qquad \frac{D_S(s_2)}{D_S(s_1)}=\frac42=\frac21=\frac{D_S(s_1)}{D_S(s_4)}
\]

%dsp-level:0
\[
\textbf{three of the six defect values are the powers}\ 1,
\]

%dsp-level:0
\[
\ 2,
\]

%dsp-level:0
\[
\ 4\ —\ \text{the octahedron, the tetrahedral light rung and the tetrahedral heavy rung}\ —\ \text{and that geometric triple is the census's only degeneracy}\ [\text{machine: the thirty-value census}];
\]

%dsp-level:0
\[
\textbf{the informational currency:}
\]

%dsp-level:0
\[
\qquad \frac{m_{\mathrm{hid}}(s)}{m_{\mathrm{hid}}(s')}=\frac{E(s)}{E(s')}\in\Bigl\{1,\ 2,\ 5,\ \frac12,\ \frac15,\ \frac25,\ \frac52\Bigr\}
\]

%dsp-level:0
\[
\qquad [\text{machine: the exact value set at}\ 2E\in\{12,24,60\}]:
\]

%dsp-level:0
\[
\boxed{\ \text{all}\ \rho_{\mathrm B}\textbf{-free, all rational, seven distinct values over the fifteen ordered pairs}\ };
\]

%dsp-level:0
\[
\qquad \textbf{the transport currency:}\quad\Delta m=\frac{\ln(4/3)}{\mathfrak h},\quad\frac{\ln2}{\mathfrak h}\ \ \text{per constituent}
\]

%dsp-level:0
\[
\qquad [\mathrm{BM.59}(4),(5);\ \rho_{\mathrm B}\text{-free}];
\]

%dsp-level:0
\[
\textbf{the ensemble currency:}
\]

%dsp-level:0
\[
\qquad \frac{T(0,S)}{T(+1,S)}=\frac{W(0)}{W(+1)}=\frac{20}{15}=\frac43
\]

%dsp-level:0
\[
\qquad [\mathrm{BM.28}(4)(\mathrm{ii});\ \mathrm{BM.30}(2);\ \rho_{\mathrm B}\text{-free}];
\]

%dsp-level:0
\[
(5)\quad\textbf{the cross-currency row (where the residue lives):}
\]

%dsp-level:0
\[
\qquad \boxed{\ \frac{m_{\mathrm{pole}}(s)}{m_{\mathrm{hid}}(s)}=\sqrt{\rho_{\mathrm B}\,\mathcal R_s}=\sqrt{\rho_{\mathrm B}}\cdot\Bigl(\frac{\sqrt2}{12},\ \frac16,\ \frac{\sqrt{2-\sqrt2}}{24},\ \frac1{24},\ \frac{\sqrt{2-2\mathsf c}}{60},\ \frac{\sqrt{2-\sqrt3}}{60}\Bigr)\ }
\]

%dsp-level:0
\[
[\text{machine: the six square roots simplified in the ordered radical extensions};\qquad\sqrt{\mathcal R_2}/\sqrt{\mathcal R_5}=10/\sqrt{2-2\mathsf c}\ \text{is}\ \rho_{\mathrm B}\text{-free again}];
\]

%dsp-level:0
\[
\textbf{the residue enters at exactly one step:}
\]

%dsp-level:0
\[
\qquad \boxed{\ \text{every mass ratio inside a currency is derived with zero residue};\qquad\text{the single cross-currency step carries the single undeclared}\ \rho_{\mathrm B}\ }
\]

%dsp-level:0
\[
(6)\quad\textbf{the residue localized:}
\]

%dsp-level:0
\[
\qquad (\mathrm{R1})\ \ \rho_{\mathrm B}=\hat\mu'^2\mathfrak h^2/(\ln2)^2\ \ \text{is the only undeclared datum any ratio of}\ (2)\text{--}(5)\ \text{needs}
\]

%dsp-level:0
\[
\qquad [\mathrm{BM.69}(5)(\mathrm{R1});\ \text{Cal-F packaging},\ \neg\text{added here}];
\]

%dsp-level:0
\[
(\mathrm{R2})\ \ \text{the six rows are read at the vertex-figure locus}\ \Theta_R
\]

%dsp-level:0
\[
= \delta_v,
\]

%dsp-level:0
\[
\ \text{whose ledger realization}\ —\ \text{a clean-enclosure loop through conical vertices}\ —\ \text{is}\ \neg\text{constructed and is inherited at every use}
\]

%dsp-level:0
\[
\qquad [\mathrm{BM.67}(9),(10);\ \text{Census FT.63}(4)\ \mathrm{Rmk};\ \mathrm{BM.69}(5)(\mathrm{R2})];
\]

%dsp-level:0
\[
(\mathrm{R3})\ \ \text{the flavour composition of the sealed classes: it enters only where a composite's pole mass is formed}\ —\ \neg\text{a ratio of}\ (2)\text{--}(5)\ \text{requires it}
\]

%dsp-level:0
\[
\qquad [\mathrm{BM.69}(3),(5)(\mathrm{R3});\ \mathrm{BM.28}(7)];
\]

%dsp-level:0
\[
(7)\quad\textbf{the non-derivation theorem, re-certified and sharpened:}
\]

%dsp-level:0
\[
\qquad \boxed{\ \neg\exists\ \text{an exact identity}\ \ m^2_{\mathrm{pole}}=\mathcal B\bigl[E_{\mathrm{int}},W(Q),\Theta_R,I_F\bigr]\ \ \text{with}\ \mathcal B\ \text{built from derived data alone}\ }
\]

%dsp-level:0
\[
\text{on four independent grounds}:\ (\mathrm i)\ \text{the degree-}16\ \text{field against integer- and binary-valued inputs}\ [\mathrm{FT.73}(2),(3)];
\]

%dsp-level:0
\[
\quad (\mathrm{ii})\ \text{the two-torsion image}\ \{0,\pi\}\ \text{reaching}\ 2\ \text{of}\ 6\ \text{states}\ [\mathrm{BM.67}(6)];
\]

%dsp-level:0
\[
(\mathrm{iii})\ \text{the calibration weight}\ {-}2\ \text{against}\ 0\ \text{on the winding route}\ [\mathrm{FT.55}(1);\ \mathrm{BM.67}(8)];
\]

%dsp-level:0
\[
\qquad (\mathrm{iv})\ \neg\exists\ \text{a count}\to\text{strength theorem}\ [\mathrm{BM.36}(3);\ \mathrm{BM.61}(8)(\mathrm{O1})]
\]

%dsp-level:0
\[
\qquad [\mathrm{BM.69}(4)\ \text{re-cited};\ \text{the grounds stand}]:
\]

%dsp-level:0
\[
\textbf{the census sharpening:}
\]

%dsp-level:0
\[
\qquad \text{the four grounds bound}\ \rho_{\mathrm B}\ \text{itself};
\]

%dsp-level:0
\[
\qquad \text{they leave}\ \mathcal R_s,
\]

%dsp-level:0
\[
\ \text{the fifteen pole ratios, the seven informational ratios, the transport differences and the ensemble ratio untouched}:
\]

%dsp-level:0
\[
\boxed{\ \text{the mass-ratio census is complete up to one undeclared dimensionless number, and complete with none inside any single currency}\ }
\]

%dsp-level:0
\[
(8)\quad\textbf{status:}\qquad\textbf{derived:}\ (1)\text{--}(5),\ (7);\qquad\textbf{localized,}\ \neg\textbf{filled:}\ (6);
\]

%dsp-level:0
\[
\neg\text{a P4 addition}\ (\rho_{\mathrm B}\ \text{is a Cal-F restatement},\ \neg\text{a new member}),
\]

%dsp-level:0
\[
\quad\neg\text{a new dimensionful scale}\ [\mathrm{P1}],
\]

%dsp-level:0
\[
\quad\neg\text{a measured value}\ [\mathrm{P5}],
\]

%dsp-level:0
\[
\quad\neg\text{a declaration},
\]

%dsp-level:0
\[
\quad\text{§Cal-BM stays empty};
\]

%dsp-level:0
\[
\neg\text{a currency is converted anywhere in this theorem};\qquad
\]

\textbf{Req.} CT2.1, D2.9 (Part 2); P1, P4, P5 (Parts 1–2); Def-H.37(2), FT.44(4), FT.47(1), FT.54(4), FT.55(1), FT.73(2), FT.73(3) (Part 3); Census FT.42(1), FT.63(4) (row-wise); Def-BM.14(1), Def-BM.14(3), BM.28(4), BM.28(7), BM.30(2), BM.36(3), BM.59(4), BM.59(5), BM.60(3), BM.61(8), BM.67(6), BM.67(8), BM.67(9), BM.67(10), BM.69(2), BM.69(3), BM.69(4), BM.69(5), BM.79(7), BM.79(9), BM.81(5) (this part); Cal-F; exact radical arithmetic in four fields and their degree-16 compositum, minimal polynomials by first linear dependence and by exact substitution, Laurent-monomial assembly, exponent arithmetic in the $\mathfrak q$-grading, exact dyadic radical enclosures, pairwise-distinctness censuses [math; machine: \texttt{checks/\allowbreak{}verify\_\allowbreak{}H\_\allowbreak{}mass\_\allowbreak{}ratio\_\allowbreak{}census.py} §E].

\textbf{Pf.}

%dsp-level:0
\[
(1):\ \frac{\hat\mu'^2D_S}{(2E\ln2/\mathfrak h)^2}
\]

%dsp-level:0
\[
= \frac{\hat\mu'^2\mathfrak h^2}{(\ln2)^2}\cdot\frac{D_S}{4E^2}
\]

%dsp-level:0
\[
= \rho_{\mathrm B}\mathcal R_s\ [\text{machine: one cleared identity over}\ \mathbb Q[\hat\mu'^2,\mathfrak h,\ln2,D_S,E]];
\]

%dsp-level:0
\[
[\rho_{\mathrm B}]=[\hat\mu'^2][\mathfrak h]^2=\mathfrak q\cdot\mathfrak q^{-1}=\mathfrak q^0\ [\text{machine: the exponent arithmetic}];
\]

%dsp-level:0
\[
\qquad \mathcal R_s\in\overline{\mathbb Q}\ \text{since}\ D_S\ \text{is algebraic and}\ E\in\mathbb Z\ [\text{machine}]
\]

%dsp-level:0
\[
(2):\ 4E^2=4(6,6,12,12,30,30)^2=(144,144,576,576,3600,3600)\ [\text{machine}];
\]

%dsp-level:0
\[
\qquad \text{the six quotients formed exactly}\ [\text{machine}];
\]

%dsp-level:0
\[
\qquad \text{the minimal polynomials:}\ 72x-1,
\]

%dsp-level:0
\[
\ 36x-1,
\]

%dsp-level:0
\[
\ 576x-1\ \text{by inspection};
\]

%dsp-level:0
\[
\ (2-\sqrt2)/576:\ y
\]

%dsp-level:0
\[
= 576x\ \text{satisfies}\ y^2-4y+2
\]

%dsp-level:0
\[
\implies 331776x^2-2304x+2
\]

%dsp-level:0
\[
= 0
\]

%dsp-level:0
\[
\implies 165888x^2-1152x+1\ [\text{machine}]
\]

%dsp-level:0
\[
(2-\sqrt3)/3600:\ y^2-4y+1\implies12960000x^2-14400x+1\ [\text{machine}];
\]

%dsp-level:0
\[
\qquad (2-2\mathsf c)/3600:\ y^4-8y^3+19y^2-12y+1\ \text{at}\ y
\]

%dsp-level:0
\[
= 3600x
\]

%dsp-level:0
\[
\implies 3600^4x^4-8\cdot3600^3x^3+19\cdot3600^2x^2-12\cdot3600x+1\ [\text{machine: the four coefficients as exact integers}]
\]

%dsp-level:0
\[
\text{the order}: \text{the five comparisons of}\ (2)\ \text{with}\ 2-\sqrt3<\tfrac{25}4(2-\sqrt2)\ \text{certified by}\ (2-\sqrt3)^2
\]

%dsp-level:0
\[
= 7-4\sqrt3<\bigl(\tfrac{25}4(2-\sqrt2)\bigr)^2\ \text{on exact dyadic isolation}\ [\text{machine}]
\]

%dsp-level:0
\[
(3):\ \frac{1/36}{(2-2\mathsf c)/3600}=\frac{3600}{36(2-2\mathsf c)}=\frac{100}{2-2\mathsf c}=\frac{50}{1-\mathsf c}\ [\text{machine}];
\]

%dsp-level:0
\[
\qquad z^4-24z^3+76z^2-64z+16
\]

%dsp-level:0
\[
= 0\ \text{at}\ z
\]

%dsp-level:0
\[
= x/50
\]

%dsp-level:0
\[
\implies x^4-1200x^3+190000x^2-8000000x+100000000
\]

%dsp-level:0
\[
= 0\ [\text{machine: cleared}];
\]

%dsp-level:0
\[
\qquad \text{the two enclosures}\ [\text{machine}]
\]

%dsp-level:0
\[
\frac{100/(2-2\mathsf c)}{4/(2-2\mathsf c)}=25\ \text{and}\ (30/6)^2=25\ [\text{machine: exact rationals}]
\]

%dsp-level:0
\[
(4):\ \hat\mu'^2D_S(s)/(\hat\mu'^2D_S(s'))=D_S(s)/D_S(s')\ [\text{machine}];
\]

%dsp-level:0
\[
\qquad 2/(2-\sqrt3)=2(2+\sqrt3)=4+2\sqrt3=(1+\sqrt3)^2\ [\text{machine: exact in}\ \mathbb Q(\sqrt3)];
\]

%dsp-level:0
\[
\qquad \text{the fifteen quotients pairwise distinct by dyadic isolation}\ [\text{machine}]
\]

%dsp-level:0
\[
E(s)/E(s')\ \text{over}\ \{6,12,30\}:\ \text{the value set}\ \{1,2,5,\tfrac12,\tfrac15,\tfrac25,\tfrac52\}\ \text{by exhaustive enumeration of the}\ 6\cdot5
\]

%dsp-level:0
\[
= 30\ \text{ordered pairs}\ [\text{machine}];
\]

%dsp-level:0
\[
\qquad W(0)/W(+1)=20/15=4/3\ [\text{machine}]
\]

%dsp-level:0
\[
(5):\ \sqrt{1/72}
\]

%dsp-level:0
\[
= \sqrt2/12,
\]

%dsp-level:0
\[
\quad\sqrt{1/36}
\]

%dsp-level:0
\[
= 1/6,
\]

%dsp-level:0
\[
\quad\sqrt{(2-\sqrt2)/576}
\]

%dsp-level:0
\[
= \sqrt{2-\sqrt2}/24,
\]

%dsp-level:0
\[
\quad\sqrt{1/576}
\]

%dsp-level:0
\[
= 1/24,
\]

%dsp-level:0
\[
\quad\sqrt{(2-2\mathsf c)/3600}
\]

%dsp-level:0
\[
= \sqrt{2-2\mathsf c}/60,
\]

%dsp-level:0
\[
\quad\sqrt{(2-\sqrt3)/3600}
\]

%dsp-level:0
\[
= \sqrt{2-\sqrt3}/60\ [\text{machine: each verified by squaring the positive side}]
\]

%dsp-level:0
\[
(6)\text{--}(8):\ \text{the three residue rows re-cited from}\ \mathrm{BM.69}(3),(5)\ \text{and}\ \mathrm{BM.67}(9),(10);
\]

%dsp-level:0
\[
\qquad \text{the four grounds re-cited from}\ \mathrm{BM.69}(4);
\]

%dsp-level:0
\[
\qquad \text{the standing-rule scans}\ [\text{machine}].\ \blacksquare
\]

\textbf{Phys.} The bridge is an exact factorization and nothing more mysterious than that. Divide the cone-defect mass by the edge-count mass and everything cancels except one dimensionless packaging and one algebraic number per state — and the algebraic numbers are six exact values in four fields, ordered, with minimal polynomials written out to a fourteen-digit leading coefficient. Inside a currency the packaging cancels too, so all fifteen pole-side ratios, all seven informational ratios, both transport differences and the ensemble's four thirds are derived with no residue at all: the census is complete. The residue lives at exactly one step, the step between currencies, and it is one number. The bridge factor's own range is a thousand and twenty-two in mass-squared, twenty-five times the widest angle ratio, and that twenty-five is precisely the square of the edge-count ratio between the heaviest and lightest floor — geometry supplies the shape, the floors supply the reach. What the four obstruction grounds bound is the packaging and nothing else, so the honest statement of the whole tier is short: every mass ratio the corpus can form is derived exactly, and one dimensionless constant stands between that census and an absolute cross-currency mass.

\pseudohead{theorem}{\texorpdfstring{BM.83}{BM.83}}{\texorpdfstring{The inter-family bridge in full: the four-currency census, the exact incommensurability, and the three-number localization}{The inter-family bridge in full: the four-currency census, the exact incommensurability, and the three-number localization}}\label{tag:BM-83}
\textbf{Thm.}

%dsp-level:0
\[
(1)\quad\textbf{the four currencies read the same doublet (the census table):}
\]

%dsp-level:0
\[
\qquad \text{the neutral class against the unit class}
\]

%dsp-level:0
\[
\qquad [\mathrm{BM.28}(1);\ \mathrm{Def\text{-}BM.8}(1)]:
\]

%tbl-level:3
{\scriptsize\setlength{\tabcolsep}{4pt}
\begin{longtable}{>{\raggedright\arraybackslash}p{\dimexpr0.1259\textwidth-2\tabcolsep\relax}|>{\raggedright\arraybackslash}p{\dimexpr0.1608\textwidth-2\tabcolsep\relax}|>{\raggedright\arraybackslash}p{\dimexpr0.2448\textwidth-2\tabcolsep\relax}|>{\raggedright\arraybackslash}p{\dimexpr0.4685\textwidth-2\tabcolsep\relax}}
\toprule
currency & the formed object & the exact value & the residue a ratio needs \\
\addlinespace[2pt]
\midrule
\endfirsthead
\toprule
currency & the formed object & the exact value & the residue a ratio needs \\
\addlinespace[2pt]
\midrule
\endhead
\bottomrule
\endlastfoot
informational & $m_n/m_p$ & $\dfrac{\ln1280}{\ln960+6\,c\,\varepsilon_{\mathrm{EM}}\ln2}$ & none — formed $[\mathrm{BM.62}(4),\allowbreak (5)]$ \\
ensemble, route $(\mathrm{ii})$ & $m_n/m_p$ & $\dfrac{T(0,S)}{T(+1,S)}=\allowbreak \dfrac43$ & none — formed; the route is $\ne\allowbreak g$adopted globally $[\mathrm{BM.28}(4)(\mathrm{ii});\ \mathrm{BM.30}(2),\allowbreak (5)]$ \\
transport & $\Delta m$ (a difference) & $\dfrac{\ln(4/3)}{\mathfrak h}$; per constituent $\dfrac{\ln2}{\mathfrak h}$ & none — formed $[\mathrm{BM.59}(4),\allowbreak (5)]$ \\
pole/defect & $\Delta(m^2_{\mathrm{pole}})$ (a difference) & $\Delta_g=\allowbreak (2,\allowbreak \sqrt2{-}1,\allowbreak 2\mathsf c{-}\sqrt3)\,\allowbreak \hat\mu'^2$ & $(\mathrm{R3})$ the composition, for a ratio $[\mathrm{BM.69}(3);\ \mathrm{BM.75}(4)]$ \\
\end{longtable}
}

%dsp-level:0
\[
\textbf{two ratios and two differences:}
\]

%dsp-level:0
\[
\qquad \text{no currency supplies both, and}\ \neg\text{a derived equation turns one into the other}
\]

%dsp-level:0
\[
\qquad [\mathrm{BM.35}(5);\ \mathrm{BM.37}(2)];
\]

%dsp-level:0
\[
(2)\quad\textbf{the prime-log ring and the irrationality theorem:}
\]

%dsp-level:0
\[
\qquad \mathfrak P:=\mathbb Q\oplus\bigoplus_{p\ \mathrm{prime}}\mathbb Q\ln p,
\]

%dsp-level:0
\[
\qquad \ln1280=8\ln2+\ln5,
\]

%dsp-level:0
\[
\qquad \ln960=6\ln2+\ln3+\ln5
\]

%dsp-level:0
\[
\qquad [\text{machine: the exact factorizations}\ 1280=2^8\!\cdot\!5,\ 960=2^6\!\cdot\!3\!\cdot\!5]:
\]

%dsp-level:0
\[
\textbf{the independence lemma:}\qquad\sum_ia_i\ln p_i=0\ \ \text{with}\ \ a_i\in\mathbb Q\ \ \Longrightarrow\ \ a_i=0\ \ \forall i
\]

%dsp-level:0
\[
[\text{math: clear denominators by}\ N;\ \ \exp\ \text{gives}\ \prod_ip_i^{Na_i}=1;\ \ \text{unique factorization forces every}\ Na_i=0;\qquad\text{machine: the integer version}\ 2^a3^b5^c=1\implies a=b=c=0\ \text{verified exhaustively on}\ |a|,|b|,|c|\le24\ \text{and by the sign-split argument}];
\]

%dsp-level:0
\[
\boxed{\ \textbf{the irrationality theorem:}\qquad\frac{\ln1280}{\ln960+t\ln2}\ \notin\ \mathbb Q\qquad\text{for every rational}\ t\ge0\ }
\]

%dsp-level:0
\[
[\text{machine: if}\ q\ln1280=p(\ln960+t\ln2)\ \text{with}\ p,q\in\mathbb Z,\ q>0,\ \text{then the}\ \ln3\ \text{coefficient gives}\ -p=0,\ \text{so}\ p=0,\ \text{and the}\ \ln5\ \text{coefficient then gives}\ q=0\ \text{by the lemma}];
\]

%dsp-level:0
\[
\textbf{corollary:}
\]

%dsp-level:0
\[
\qquad \frac{m_n}{m_p}\ \ \text{is irrational on the informational currency at every rational coupling}\ \varepsilon_{\mathrm{EM}}\in\mathbb Q\ \text{and both priced}\ c\in\bigl\{\tfrac12,\tfrac35\bigr\}
\]

%dsp-level:0
\[
\qquad [t=6c\,\varepsilon_{\mathrm{EM}};\ \text{machine}];
\]

%dsp-level:0
\[
(3)\quad\textbf{the exact incommensurability of the two formed ratios:}
\]

%dsp-level:0
\[
\qquad \frac{m_n}{m_p}\bigg|_{\mathrm{info}}=\frac43\iff3\ln1280=4\bigl(\ln960+t\ln2\bigr)\iff0=4\ln3+\ln5+4t\ln2\iff t=-\,\frac{\ln405}{\ln16}
\]

%dsp-level:0
\[
[\text{machine: solved exactly in}\ \mathfrak P;\qquad4\ln3+\ln5=\ln(3^4\!\cdot\!5)=\ln405>0\ \text{and}\ \ln16>0\implies t<0]:
\]

%dsp-level:0
\[
\boxed{\ \textbf{the informational ratio never equals the ensemble ratio at any admissible coupling:}\qquad\frac{\ln1280}{\ln960+t\ln2}\ \ne\ \frac43\qquad\forall\,t\ge0\ }
\]

%dsp-level:0
\[
\textbf{and the order between them is exact:}
\]

%dsp-level:0
\[
\qquad 1\ <\ \frac{\ln1280}{\ln960}\ <\ \frac43
\]

%dsp-level:0
\[
\qquad [\ 4\ln960-3\ln1280=4\ln3+\ln5=\ln405>0;\qquad\ln1280>\ln960\iff1280>960;\qquad\text{machine}\ ];
\]

%dsp-level:0
\[
\textbf{the two ratios are not merely unequal, they are of different kinds:}
\]

%dsp-level:0
\[
\qquad \frac43\in\mathbb Q\ \ \text{and}\ \ \frac{\ln1280}{\ln960}\in\mathfrak P\setminus\mathbb Q
\]

%dsp-level:0
\[
\qquad [(2),(3);\ \text{machine}];
\]

%dsp-level:0
\[
(4)\quad\textbf{the dimensionless scale set (restatements,}\ \neg\textbf{new members):}
\]

%dsp-level:0
\[
\qquad \sigma_\ell:=\frac{\mu_\ell\,\mathfrak h}{\ln2}\in\mathfrak q^0,
\]

%dsp-level:0
\[
\qquad \sigma_\nu:=\frac{\mu_\nu\,\mathfrak h}{\ln2}\in\mathfrak q^0,
\]

%dsp-level:0
\[
\qquad \rho_{\mathrm B}=\frac{\hat\mu'^2\mathfrak h^2}{(\ln2)^2}\in\mathfrak q^0
\]

%dsp-level:0
\[
[\mu_\ell,\mu_\nu\ \text{are the Cal-F/P4 channel scales}\ [\mathrm{FT.113};\ \mathrm{FT.143}(3);\ \text{Cal-F}];\qquad[\mu_\ell]=[\mu_\nu]=\mathfrak q^{+1/2}\ \text{from}\ m=\mu z\ \text{and}\ [m]=\mathfrak q^{+1/2};\qquad\text{machine: the exponent arithmetic}];
\]

%dsp-level:0
\[
\textbf{they are restatements of undeclared data,}\ \neg\textbf{declarations:}
\]

%dsp-level:0
\[
\qquad \neg\text{a P4 addition},
\]

%dsp-level:0
\[
\qquad \text{§Cal-BM stays empty}
\]

%dsp-level:0
\[
\qquad [\mathrm{BM.77}(2)\ \text{pattern}];
\]

%dsp-level:0
\[
(5)\quad\textbf{the inter-family bridge, in full (every cross-currency ratio of the census):}
\]

%dsp-level:0
\[
\qquad \text{with}\ m_{\mathrm{hid}}(s)\,\mathfrak h/\ln2=2E(s),\quad m^2_{\mathrm{pole}}(s)/\hat\mu'^2=D_S(s),\quad m^{(\mathrm{ch})}_i=\mu_{\mathrm{ch}}z^{(\mathrm{ch})}_i:
\]

%dsp-level:0
\[
\boxed{\ \frac{m^{(\ell)}_i}{m_{\mathrm{hid}}(s)}=\frac{\sigma_\ell\,z^{(\ell)}_i}{2E(s)},\qquad\frac{m^{(\nu)}_j}{m^{(\ell)}_i}=\frac{\sigma_\nu\,z^{(\nu)}_j}{\sigma_\ell\,z^{(\ell)}_i},\qquad\frac{m_{\mathrm{pole}}(s)}{m^{(\ell)}_i}=\frac{2E(s)}{\sigma_\ell\,z^{(\ell)}_i}\cdot\frac{\sqrt{\rho_{\mathrm B}\,\mathcal R_s}}{1}\ }
\]

%dsp-level:0
\[
[\text{machine: the three identities cleared in}\ \mathfrak P\otimes\mathbb Q[\rho_{\mathrm B},\sigma_\ell,\sigma_\nu];\qquad\mathrm{BM.82}(1),(5)\ \text{re-cited}];
\]

%dsp-level:0
\[
\boxed{\ \textbf{the localization theorem:}\qquad\text{every cross-currency mass ratio of the census is an exact derived expression in exactly three undeclared dimensionless numbers}\ \ (\rho_{\mathrm B},\ \sigma_\ell,\ \sigma_\nu),\ \text{and in no others}\ }
\]

%dsp-level:0
\[
\textbf{minimality (each is needed):}
\]

%dsp-level:0
\[
\qquad m_{\mathrm{pole}}/m_{\mathrm{hid}}\ \text{depends on}\ \rho_{\mathrm B}\ \text{alone};
\]

%dsp-level:0
\[
\qquad m^{(\ell)}/m_{\mathrm{hid}}\ \text{on}\ \sigma_\ell\ \text{alone};
\]

%dsp-level:0
\[
\qquad m^{(\nu)}/m^{(\ell)}\ \text{on}\ \sigma_\nu/\sigma_\ell\ \text{alone}
\]

%dsp-level:0
\[
\qquad [\text{machine: the exact monomial dependence of each of the three formulas}];
\]

%dsp-level:0
\[
\textbf{and every ratio inside a currency depends on none of them}\qquad[\mathrm{BM.82}(4);\ \text{machine}];
\]

%dsp-level:0
\[
(6)\quad\textbf{the cross-channel factor census (the nine exact derived factors):}
\]

%dsp-level:0
\[
\qquad \frac{m^{(\mathrm{ch})}_i}{m^{(\mathrm{ch}')}_i}=\frac{\mu_{\mathrm{ch}}}{\mu_{\mathrm{ch}'}}\cdot\frac{z^{(\mathrm{ch})}_i}{z^{(\mathrm{ch}')}_i}
\]

%dsp-level:0
\[
\qquad [\text{machine: the exact factorization}]:
\]

%tbl-level:3
{\scriptsize\setlength{\tabcolsep}{4pt}
\begin{longtable}{>{\raggedright\arraybackslash}p{\dimexpr0.1119\textwidth-2\tabcolsep\relax}|>{\raggedright\arraybackslash}p{\dimexpr0.2867\textwidth-2\tabcolsep\relax}|>{\raggedright\arraybackslash}p{\dimexpr0.3007\textwidth-2\tabcolsep\relax}|>{\raggedright\arraybackslash}p{\dimexpr0.3007\textwidth-2\tabcolsep\relax}}
\toprule
line & $z^{(q)}_i/z^{(\nu)}_i$ & $z^{(q)}_i/z^{(\ell)}_i$ & $z^{(\nu)}_i/z^{(\ell)}_i$ \\
\addlinespace[2pt]
\midrule
\endfirsthead
\toprule
line & $z^{(q)}_i/z^{(\nu)}_i$ & $z^{(q)}_i/z^{(\ell)}_i$ & $z^{(\nu)}_i/z^{(\ell)}_i$ \\
\addlinespace[2pt]
\midrule
\endhead
\bottomrule
\endlastfoot
$i=\allowbreak 0$ (the heaviest) & $\bigl(\tfrac{1005964365750}{10^{12}},\allowbreak \tfrac{1005964365751}{10^{12}}\bigr)$ & $\bigl(\tfrac{1026138854777}{10^{12}},\allowbreak \tfrac{1026138854778}{10^{12}}\bigr)$ & $\bigl(\tfrac{1020054874420}{10^{12}},\allowbreak \tfrac{1020054874421}{10^{12}}\bigr)$ \\
$i=\allowbreak 2$ (the middle) & $\bigl(\tfrac{730177488628}{10^{12}},\allowbreak \tfrac{730177488629}{10^{12}}\bigr)$ & $\bigl(\tfrac{441377821972}{10^{12}},\allowbreak \tfrac{441377821973}{10^{12}}\bigr)$ & $\bigl(\tfrac{604480183033}{10^{12}},\allowbreak \tfrac{604480183034}{10^{12}}\bigr)$ \\
$i=\allowbreak 1$ (the lightest) & $\bigl(\tfrac{1963683086044}{10^{12}},\allowbreak \tfrac{1963683086045}{10^{12}}\bigr)$ & $\bigl(\tfrac{25609326778349}{10^{12}},\allowbreak \tfrac{25609326778350}{10^{12}}\bigr)$ & $\bigl(\tfrac{13041476478742}{10^{12}},\allowbreak \tfrac{13041476478743}{10^{12}}\bigr)$ \\
\end{longtable}
}

%dsp-level:0
\[
[\text{machine: monotone exact interval division at width}<2^{-150}];
\]

%dsp-level:0
\[
\boxed{\ \textbf{the light-end concentration:}\qquad1\ <\ \frac{z^{(q)}_0}{z^{(\nu)}_0},\ \frac{z^{(\nu)}_0}{z^{(\ell)}_0},\ \frac{z^{(q)}_0}{z^{(\ell)}_0}\ <\ \frac{27}{25},\qquad\text{while}\qquad\frac{25609326778349}{10^{12}}\ <\ \frac{z^{(q)}_1}{z^{(\ell)}_1}\ <\ \frac{26}{1}\ }
\]

%dsp-level:0
\[
\textbf{the three channels' heaviest lines are nearly degenerate and their lightest lines differ by more than an order of magnitude}\ —\ \text{the entire inter-family structure lives at the light end}
\]

%dsp-level:0
\[
[\text{machine: the exact bracket comparisons};\qquad\text{the structural mirror of}\ \mathrm{BM.79}(5)\text{'s top-weight row on the Platonic side}\ [\text{Census FT.76}(4)\ \text{row-wise}]];
\]

%dsp-level:0
\[
(7)\quad\textbf{the inter-family verdict (what is and what is not derived):}
\]

%dsp-level:0
\[
\qquad \textbf{derived:}\ \text{the shape invariants}\ \sum_iz_i
\]

%dsp-level:0
\[
= 6,
\]

%dsp-level:0
\[
\ Q
\]

%dsp-level:0
\[
= \tfrac23\ \text{in all three channels};
\]

%dsp-level:0
\[
\ \text{the one cubic form with three constants};
\]

%dsp-level:0
\[
\ \text{the hierarchy-order theorem}\ [\mathrm{BM.80}(3),(4),(7),(8)];
\]

%dsp-level:0
\[
\text{the nine cross-channel factors of}\ (6);
\]

%dsp-level:0
\[
\qquad \text{the exact factorization of every cross-currency ratio and its localization to}\ (\rho_{\mathrm B},\sigma_\ell,\sigma_\nu)\ (5);
\]

%dsp-level:0
\[
\textbf{derived:}\ \text{the quark side of the inter-family ratio in two currencies}\ —\ \text{the pole currency's}\ \hat\mu'^2D_S\ \text{and the informational currency's}\ 2E\ln2/\mathfrak h\ \text{with the exact bridge factor}\ \mathcal R_s\ [\mathrm{BM.81}(1);\ \mathrm{BM.82}(2)];
\]

%dsp-level:0
\[
\boxed{\ \neg\textbf{derived:}\ \ \text{the inter-family mass ratio as a number}\ }
\]

%dsp-level:0
\[
\qquad \text{it requires}\ \rho_{\mathrm B}\ (\mathrm{R1}),
\]

%dsp-level:0
\[
\ \sigma_\ell,
\]

%dsp-level:0
\[
\ \sigma_\nu\ (\text{Cal-F}),
\]

%dsp-level:0
\[
\ \text{and}\ —\ \text{on the pole side of a composite}\ —\ (\mathrm{R3})\ \text{the composition}:
\]

%dsp-level:0
\[
\textbf{the obstruction is exactly three dimensionless numbers and one attachment, all named, none supplied, none fitted}
\]

%dsp-level:0
\[
\qquad [\mathrm{BM.69}(4),(5);\ \mathrm{BM.60}(7)];
\]

%dsp-level:0
\[
(8)\quad\textbf{the absolute neutrino mass:}
\]

%dsp-level:0
\[
\qquad \neg\text{derived}
\]

%dsp-level:0
\[
\qquad [\mathrm{FT.143}(3):\ \mu_\nu\ \text{undeclared}]:
\]

%dsp-level:0
\[
\qquad \text{the}\ \nu\text{-channel's internal ratios are derived}\ (6),\ \text{its absolute scale is not};
\]

%dsp-level:0
\[
\text{the two-bit floor applies to it on the informational read:}
\]

%dsp-level:0
\[
\qquad m_\nu\mathfrak h\ge2\ln2\ \ \text{once materialized as a bound chain},
\]

%dsp-level:0
\[
\qquad m_\nu\mathfrak h\ge0\ \ \text{by PSD alone}
\]

%dsp-level:0
\[
\qquad [\mathrm{BM.70}(2),(3)\ \text{re-cited}];
\]

%dsp-level:0
\[
(9)\quad\textbf{the bars:}
\]

%dsp-level:0
\[
\qquad \neg\exists\ \text{an SM mass, mass ratio, mixing angle or Yukawa value in any clause}
\]

%dsp-level:0
\[
\qquad [\mathrm{P5};\ \text{machine: zero external tokens, zero decimal numerals, zero}\ [\mathrm{EXT}]\ \text{markers in this tier}];
\]

%dsp-level:0
\[
\neg\text{a measured external constant is identified with any corpus quantity}\ [\mathrm{P5}];
\]

%dsp-level:0
\[
\qquad \text{species names are confined to the identification layers}\ [\mathrm{Def\text{-}BM.3}(7);\ \mathrm{Def\text{-}BM.6}(7);\ \mathrm{Def\text{-}BM.14}(8)];
\]

%dsp-level:0
\[
\text{topology labels are internal bookkeeping}\ [\mathrm{Def\text{-}H.32}(4)];
\]

%dsp-level:0
\[
\qquad \text{amplitude}\ne\text{mass}\ [\mathrm{FT.79}(5)];
\]

%dsp-level:0
\[
\qquad \text{rung}\ne\text{scale}\ [\mathrm{D0.8}];
\]

%dsp-level:0
\[
\qquad \neg\text{a currency is silently converted anywhere in this theorem};
\]

%dsp-level:0
\[
(10)\quad\textbf{status:}
\]

%dsp-level:0
\[
\qquad \textbf{derived:}\ (1)\text{--}(8);
\]

%dsp-level:0
\[
\qquad \textbf{localized,}\ \neg\textbf{filled:}\ (7)\text{'s obstruction};
\]

%dsp-level:0
\[
\qquad \textbf{declared: nothing}\ —\ (4)\text{'s three symbols are restatements of Cal-F/P4 data};
\]

%dsp-level:0
\[
\neg\text{a P4 addition},\quad\neg\text{a new dimensionful scale}\ [\mathrm{P1}],\quad\neg\text{a measured value}\ [\mathrm{P5}],\quad
\]

%dsp-level:0
\[
\qquad \text{the mass-form bridge is application-layer, its residue reduced to}\ (\rho_{\mathrm B},\sigma_\ell,\sigma_\nu)\ \text{and}\ (\mathrm{R2}),(\mathrm{R3}).
\]

\textbf{Req.} D0.8, D2.9, CT2.1 (Part 2); P1, P4, P5 (Parts 1–2); Def-H.32(4), Def-H.37(2), FT.79(5), FT.113, FT.143(3) (Part 3); Def-BM.3(7), Def-BM.6(7), Def-BM.8(1), Def-BM.14(8), BM.28(1), BM.28(4), BM.30(2), BM.30(5), BM.35(5), BM.37(2), BM.59(4), BM.59(5), BM.60(7), BM.62(4), BM.62(5), BM.69(3), BM.69(4), BM.69(5), BM.70(2), BM.70(3), BM.75(4), BM.77(2), BM.79(5), BM.80(3), BM.80(4), BM.80(6), BM.80(7), BM.80(8), BM.81(1), BM.82(1), BM.82(2), BM.82(4), BM.82(5) (this part); Cal-F; exact arithmetic in the prime-log ring with integer sign tests, unique-factorization independence, exact symbolic solving in $\mathfrak P$, Laurent-monomial assembly and $\mathfrak q$-exponent arithmetic, monotone exact interval division [math; machine: \texttt{checks/\allowbreak{}verify\_\allowbreak{}H\_\allowbreak{}mass\_\allowbreak{}ratio\_\allowbreak{}census.py} §F].

\textbf{Pf.}

%dsp-level:0
\[
(1):\ \text{the four rows re-cited from their own tags}\ [\mathrm{BM.62}(4),(5);\ \mathrm{BM.28}(4)(\mathrm{ii});\ \mathrm{BM.30}(2);\ \mathrm{BM.59}(4),(5);\ \mathrm{BM.69}(3);\ \mathrm{BM.75}(4)];
\]

%dsp-level:0
\[
\qquad \text{the formed-object column: two ratios and two differences}\ [\text{machine: the token census of the four expressions}]
\]

%dsp-level:0
\[
(2):\ 1280
\]

%dsp-level:0
\[
= 2^8\cdot5\ \text{and}\ 960
\]

%dsp-level:0
\[
= 2^6\cdot3\cdot5\ [\text{machine: exact integer factorizations}]\implies\ln1280
\]

%dsp-level:0
\[
= 8\ln2+\ln5,
\]

%dsp-level:0
\[
\ \ln960
\]

%dsp-level:0
\[
= 6\ln2+\ln3+\ln5\ [\text{machine: the prime-log representation}];
\]

%dsp-level:0
\[
\qquad \text{the lemma:}\ \exp\bigl(\sum Na_i\ln p_i\bigr)=\prod p_i^{Na_i}=1\implies Na_i=0\ \forall i\ [\text{math: unique factorization}]
\]

%dsp-level:0
\[
\text{machine: the exhaustive census of}\ 2^a3^b5^c=1\ \text{over}\ |a|,|b|,|c|\le24\ \text{returns exactly}\ (0,0,0);
\]

%dsp-level:0
\[
\qquad \text{the rational case}:\ q(8\ln2+\ln5)
\]

%dsp-level:0
\[
= p((6+t)\ln2+\ln3+\ln5)\implies\text{the}\ \ln3\ \text{coefficient}\ -p
\]

%dsp-level:0
\[
= 0
\]

%dsp-level:0
\[
\implies p
\]

%dsp-level:0
\[
= 0
\]

%dsp-level:0
\[
\implies q\cdot8\ln2+q\ln5
\]

%dsp-level:0
\[
= 0
\]

%dsp-level:0
\[
\implies q
\]

%dsp-level:0
\[
= 0\ [\text{machine: the coefficient extraction and the lemma}]
\]

%dsp-level:0
\[
(3):\ 3\ln1280=24\ln2+3\ln5;
\]

%dsp-level:0
\[
\qquad 4\ln960+4t\ln2=(24+4t)\ln2+4\ln3+4\ln5\ [\text{machine}];
\]

%dsp-level:0
\[
\qquad \text{the difference}\ 4\ln960+4t\ln2-3\ln1280
\]

%dsp-level:0
\[
= 4\ln3+\ln5+4t\ln2
\]

%dsp-level:0
\[
= \ln405+4t\ln2\ [\text{machine: the prime-log ring}]
\]

%dsp-level:0
\[
\ln405>0\iff405>1,\quad\ln16>0\iff16>1\implies t=-\ln405/\ln16<0\ [\text{machine}];
\]

%dsp-level:0
\[
\qquad \ln1280>\ln960\iff1280>960\ [\text{machine}];
\]

%dsp-level:0
\[
\qquad \text{the kind row}:\ \tfrac43\in\mathbb Q,\ \tfrac{\ln1280}{\ln960}\notin\mathbb Q\ [(2)]
\]

%dsp-level:0
\[
(4):\ [m]=\mathfrak q^{+1/2}\ [\mathrm{CT2.1};\ \mathrm{D2.9}],\quad z\in\mathfrak q^0\implies[\mu_{\mathrm{ch}}]=\mathfrak q^{+1/2};
\]

%dsp-level:0
\[
\qquad [\mathfrak h]=\mathfrak q^{-1/2},\quad[\ln2]=\mathfrak q^0\implies\sigma_{\mathrm{ch}}\in\mathfrak q^0\ [\text{machine: the exponent arithmetic}];
\]

%dsp-level:0
\[
\qquad \rho_{\mathrm B}\in\mathfrak q^0\ [\mathrm{BM.82}(1)]
\]

%dsp-level:0
\[
(5):\ \frac{\mu_\ell z}{2E\ln2/\mathfrak h}=\frac{\mu_\ell\mathfrak h}{\ln2}\cdot\frac{z}{2E}=\frac{\sigma_\ell z}{2E}\ [\text{machine: cleared}];
\]

%dsp-level:0
\[
\qquad \frac{\mu_\nu z^{(\nu)}}{\mu_\ell z^{(\ell)}}=\frac{\sigma_\nu z^{(\nu)}}{\sigma_\ell z^{(\ell)}}\ [\text{machine}];
\]

%dsp-level:0
\[
\qquad \frac{m_{\mathrm{pole}}}{m^{(\ell)}_i}
\]

%dsp-level:0
\[
= \frac{m_{\mathrm{pole}}}{m_{\mathrm{hid}}}\cdot\frac{m_{\mathrm{hid}}}{m^{(\ell)}_i}
\]

%dsp-level:0
\[
= \sqrt{\rho_{\mathrm B}\mathcal R_s}\cdot\frac{2E}{\sigma_\ell z^{(\ell)}_i}\ [\text{machine: the two re-cited identities composed}]
\]

%dsp-level:0
\[
\text{minimality}: \text{the symbol census of the three expressions}\ —\ \sqrt{\rho_{\mathrm B}\mathcal R_s}\ \text{carries}\ \rho_{\mathrm B}\ \text{and neither}\ \sigma;
\]

%dsp-level:0
\[
\quad\sigma_\ell z/2E\ \text{carries}\ \sigma_\ell\ \text{and neither}\ \rho_{\mathrm B}\ \text{nor}\ \sigma_\nu;
\]

%dsp-level:0
\[
\quad\sigma_\nu z^{(\nu)}/(\sigma_\ell z^{(\ell)})\ \text{carries both}\ \sigma\text{'s and no}\ \rho_{\mathrm B}\ [\text{machine}]
\]

%dsp-level:0
\[
(6):\ \text{the nine factors by monotone exact interval division of the enclosed}\ z\text{'s}\ [\text{machine: widths}<2^{-150}];
\]

%dsp-level:0
\[
\qquad \text{the bracket comparisons against}\ \tfrac{27}{25}\ \text{and}\ 26\ [\text{machine}]
\]

%dsp-level:0
\[
(7)\text{--}(10):\ \text{the assembly of}\ \mathrm{BM.80}\text{'s and}\ \mathrm{BM.82}\text{'s rows};
\]

%dsp-level:0
\[
\qquad \text{the three residue symbols re-cited from}\ \mathrm{BM.69}(5);
\]

%dsp-level:0
\[
\qquad \text{the standing-rule scans}\ [\text{machine}].\ \blacksquare
\]

\textbf{Phys.} Four currencies look at the same two charge classes and no two of them say the same thing. Two of them produce a ratio — four thirds from the ensemble, and the logarithm of one thousand two hundred eighty over the logarithm of nine hundred sixty from the edge count — and two produce a difference: the transport register's log of four thirds over h, and the cone defect's exact algebraic splitting. The tier's central negative is exact rather than numerical. The informational ratio is irrational at every rational coupling, by unique factorization of the primes two, three and five inside the logarithms; four thirds is rational; so the two can never agree, and solving for the coupling at which they would gives a negative number — the logarithm of four hundred five over the logarithm of sixteen — outside the admissible range entirely. The positive side is a localization. Three dimensionless numbers, each a restatement of data the corpus already carries undeclared, are exactly what stands between the census and an absolute inter-family mass: the pole packaging, and the two channel scales. Every cross-currency ratio in the corpus is an exact expression in those three and nothing else; every ratio inside a currency needs none of them. The nine cross-channel factors then show where the structure lives: the three heaviest lines of the three families are degenerate to within three percent while the three lightest differ by a factor of twenty-six — the same light-end concentration the Platonic ladder shows at its top weights, arriving from an entirely different direction. Nothing is converted, nothing is fitted, and the absolute neutrino mass stays exactly where the ledger says it stays.

\begin{theorem}[{\texorpdfstring{BM.84 — The neutron–proton mass ratio: the exact ratio, the irrational critical locus, the range theorem, and the trichotomy}{BM.84 - The neutron-proton mass ratio: the exact ratio, the irrational critical locus, the range theorem, and the trichotomy}}]\label{tag:BM-84}
\textbf{Thm.}

%dsp-level:0
\[
(1)\quad\textbf{the decomposition and the bridge summand's currency row:}
\]

%dsp-level:0
\[
\qquad \Delta m:=m_n-m_p=\Delta m_{\mathrm{Lyapunov}}+\Delta m_{\mathrm{force}}+\Delta m_{\mathrm{metric}}+\Delta m_{\mathrm{binding}}+\Delta m_{\mathrm{self\text{-}int}}+\Delta m_{\mathrm{transport}}
\]

%dsp-level:0
\[
\qquad [\mathrm{BM.62}(1)]:
\]

%dsp-level:0
\[
\Delta m_{\mathrm{Lyapunov}}=0\ [\mathrm{BM.51}(6)];
\]

%dsp-level:0
\[
\qquad \Delta m_{\mathrm{metric}}=0\ [\mathrm{BM.53}(2);\ \mathrm{BM.58}(4)];
\]

%dsp-level:0
\[
\qquad \Delta m_{\mathrm{binding}}=0\ [\mathrm{BM.54}(1)\text{--}(3)];
\]

%dsp-level:0
\[
\qquad \Delta m_{\mathrm{force}}=\Delta m_{\mathrm{self\text{-}int}}=-\Delta m_{\mathrm{EM}}\ [\mathrm{BM.52}(4),(5);\ \mathrm{BM.46}(4)];
\]

%dsp-level:0
\[
\boxed{\ \Delta m_{\mathrm{transport}}=+\frac{\ln(4/3)}{\mathfrak h},\qquad\Delta m_{\mathrm{EM}}=-\,\frac{6\,c\,\varepsilon_{\mathrm{EM}}\ln2}{\mathfrak h},\qquad\Delta m_{\mathrm{bridge}}=0\ \ \text{in the informational currency}\ }
\]

%dsp-level:0
\[
[\mathrm{BM.59}(4);\ \mathrm{BM.46}(4);\ \mathrm{BM.62}(3);\qquad\textbf{the bridge summand's row:}\ \ \Delta m_{\mathrm{bridge}}\ \text{is}\ \neg\textbf{formed} \ \text{in the informational currency}\ —\ \text{the pole-side companion is}\ \Delta_g>0,\ \text{a different currency},
\]

%dsp-level:0
\[
\text{and a cross-currency addition is a category the corpus bars}\ [\mathrm{BM.69}(3);\ \mathrm{BM.35}(5);\ \mathrm{BM.37}(2);\ \mathrm{BM.71}(1)\ \text{row}\ 5];
\]

%dsp-level:0
\[
\qquad \text{the same substitution's two readings point the same way}\ [\mathrm{BM.75}(4),(5)]\ \text{and are}\ \neg\text{addable}];
\]

%dsp-level:0
\[
(2)\quad\textbf{the exact ratio:}
\]

%dsp-level:0
\[
\qquad m_n\,\mathfrak h=6\ln2+\ln20=\ln(2^8\!\cdot\!5)=\ln1280,
\]

%dsp-level:0
\[
\qquad m_p\,\mathfrak h=6\ln2+\ln15+t\ln2=\ln(2^6\!\cdot\!3\!\cdot\!5)+t\ln2=\ln960+t\ln2
\]

%dsp-level:0
\[
\qquad [\mathrm{Def\text{-}BM.13}(5);\ \text{machine}]:
\]

%dsp-level:0
\[
\boxed{\ \frac{m_n}{m_p}\ =\ \frac{\ln1280}{\ln960+6\,c\,\varepsilon_{\mathrm{EM}}\ln2}\ =\ \frac{\ln1280}{\ln960+t\ln2},\qquad t\ :=\ 6\,c\,\varepsilon_{\mathrm{EM}}\ \ge\ 0\ }
\]

%dsp-level:0
\[
\qquad [\mathrm{BM.62}(5)\ \text{re-derived};\ \text{machine: the prime-log assembly}];
\]

%dsp-level:0
\[
\textbf{the sub-critical branch is the branch}\ t<\ln(4/3)/\ln2;
\]

%dsp-level:0
\[
\qquad \textbf{the ratio is a dimensionless element of}\ \mathfrak P,\ \neg\text{a rational number in general}
\]

%dsp-level:0
\[
\qquad [\mathrm{BM.83}(2)];
\]

%dsp-level:0
\[
(3)\quad\textbf{the range theorem (derived here):}
\]

%dsp-level:0
\[
\qquad \frac{\partial}{\partial t}\frac{m_n}{m_p}=-\,\frac{\ln1280\cdot\ln2}{(\ln960+t\ln2)^2}<0
\]

%dsp-level:0
\[
\qquad [\text{machine: the sign census of}\ \ln1280,\ \ln2>0\ \text{in the prime-log ring};\ \text{the monotonicity is cross-multiplied in the Pf}]:
\]

%dsp-level:0
\[
\boxed{\ 0\ <\ \frac{m_n}{m_p}\ \le\ \frac{\ln1280}{\ln960}\ <\ \frac{21}{20}\qquad\text{for every admissible}\ t\ge0,\ \text{with equality at}\ t=0\ \text{only}\ }
\]

%dsp-level:0
\[
\frac{\ln1280}{\ln960}\in\Bigl(\frac{26}{25},\frac{21}{20}\Bigr)
\]

%dsp-level:0
\[
\qquad [\text{machine:}\ 25\ln1280>26\ln960\iff2^{44}>3^{26}\!\cdot\!5;\qquad20\ln1280<21\ln960\iff2^{34}<3^{21}\!\cdot\!5];
\]

%dsp-level:0
\[
2^{44}=17592186044416>12709329141645=3^{26}\!\cdot\!5,
\]

%dsp-level:0
\[
\qquad 2^{34}=17179869184<52301766015=3^{21}\!\cdot\!5
\]

%dsp-level:0
\[
\qquad [\text{machine: the two integer comparisons}];
\]

%dsp-level:0
\[
\textbf{the exact upper bound of the neutron--proton mass ratio on the informational currency is}\ \tfrac{21}{20},
\]

%dsp-level:0
\[
\ \textbf{attained in the limit}\ t\to0^+,
\]

%dsp-level:0
\[
\ \textbf{and no coupling raises it};
\]

%dsp-level:0
\[
(4)\quad\textbf{the irrationality and the rational locus:}
\]

%dsp-level:0
\[
\qquad \boxed{\ \frac{m_n}{m_p}\notin\mathbb Q\ \ \text{at every rational coupling}\ t\in\mathbb Q\ }
\]

%dsp-level:0
\[
\qquad [\mathrm{BM.83}(2)\text{'s theorem re-cited}]:
\]

%dsp-level:0
\[
\frac{m_n}{m_p}=1\iff t=\frac{\ln(4/3)}{\ln2}=2-\log_23\ \notin\ \mathbb Q
\]

%dsp-level:0
\[
\qquad [\text{machine: solved exactly in}\ \mathfrak P;\qquad\log_23\notin\mathbb Q\iff\neg\exists\,(p,q)\ne(0,0):\ 2^p=3^q;\ \ \text{the census over}\ |p|,|q|\le64\ \text{returns none},\ \text{and unique factorization closes the general case}];
\]

%dsp-level:0
\[
\boxed{\ \textbf{the degenerate locus is exactly one coupling, and that coupling is irrational:}\qquad\text{at every rational coupling}\ \ m_n\ne m_p\ }
\]

%dsp-level:0
\[
(5)\quad\textbf{the exact critical strength:}
\]

%dsp-level:0
\[
\qquad \boxed{\ \varepsilon^{\mathrm{crit}}(c)\ =\ \frac{\ln(4/3)}{6\,c\,\ln2}\ =\ \frac{2\ln2-\ln3}{6\,c\,\ln2}\ =\ \frac{2-\log_23}{6\,c}\ }
\]

%dsp-level:0
\[
\qquad [\mathrm{BM.62}(6)\ \text{re-derived};\ \text{machine: the unique zero of}\ \Delta m]:
\]

%dsp-level:0
\[
\textbf{the four integer certificates:}\qquad\frac19<\varepsilon^{\mathrm{crit}}\Bigl(\frac35\Bigr)<\frac18<\varepsilon^{\mathrm{crit}}\Bigl(\frac12\Bigr)<\frac17
\]

%dsp-level:0
\[
\varepsilon^{\mathrm{crit}}(\tfrac35)>\tfrac19\iff2^8>3^5\ (256>243);
\]

%dsp-level:0
\[
\qquad \varepsilon^{\mathrm{crit}}(\tfrac35)<\tfrac18\iff2^{31}<3^{20}\ (2147483648<3486784401);
\]

%dsp-level:0
\[
\varepsilon^{\mathrm{crit}}(\tfrac12)>\tfrac18\iff2^{13}>3^8\ (8192>6561);
\]

%dsp-level:0
\[
\qquad \varepsilon^{\mathrm{crit}}(\tfrac12)<\tfrac17\iff2^{11}<3^7\ (2048<2187)
\]

%dsp-level:0
\[
\qquad [\mathrm{BM.62}(7)\ \text{re-certified};\ \text{machine}];
\]

%dsp-level:0
\[
\boxed{\ \varepsilon^{\mathrm{crit}}(c)\notin\mathbb Q\qquad\text{for every rational}\ c>0\ }\qquad[(4)\text{'s}\ 2-\log_23\notin\mathbb Q;\ \text{machine}];
\]

%dsp-level:0
\[
(6)\quad\textbf{the trichotomy (total, exact, and non-empty on both sides):}
\]

%dsp-level:0
\[
\boxed{\ \varepsilon_{\mathrm{EM}}<\varepsilon^{\mathrm{crit}}(c)\iff m_n>m_p\iff\frac{m_n}{m_p}>1;\qquad\varepsilon_{\mathrm{EM}}=\varepsilon^{\mathrm{crit}}(c)\iff m_n=m_p\iff\frac{m_n}{m_p}=1;\qquad\varepsilon_{\mathrm{EM}}>\varepsilon^{\mathrm{crit}}(c)\iff m_n<m_p\iff\frac{m_n}{m_p}<1\ }
\]

%dsp-level:0
\[
[\mathrm{BM.62}(6),(9)\ \text{re-derived from}\ (2),(3);\ \text{machine: the sign census in}\ \mathfrak P];
\]

%dsp-level:0
\[
\textbf{the opening region is non-empty and integer-certified:}
\]

%dsp-level:0
\[
\qquad \Bigl\{0\le\varepsilon_{\mathrm{EM}}\le\tfrac19\Bigr\}\subseteq\Bigl\{\varepsilon_{\mathrm{EM}}<\varepsilon^{\mathrm{crit}}(c)\Bigr\}\ \ \text{at both priced}\ c
\]

%dsp-level:0
\[
\qquad [\mathrm{BM.62}(8):4^3=64>54=3^3\!\cdot\!2\ \ \text{at}\ c=\tfrac12;\quad4^5=1024>972=3^5\!\cdot\!2^2\ \ \text{at}\ c=\tfrac35;\ \text{machine}];
\]

%dsp-level:0
\[
\textbf{the closing region is non-empty:}
\]

%dsp-level:0
\[
\qquad \varepsilon_{\mathrm{EM}}=\tfrac16>\tfrac17>\varepsilon^{\mathrm{crit}}\bigl(\tfrac12\bigr)\ \ \text{and}\ \ \varepsilon_{\mathrm{EM}}=\tfrac16>\tfrac18>\varepsilon^{\mathrm{crit}}\bigl(\tfrac35\bigr)
\]

%dsp-level:0
\[
\qquad [\text{machine:}\ 6<7,\ 6<8\ \text{and}\ (5)\text{'s brackets}];
\]

%dsp-level:0
\[
\textbf{the critical locus is exact, unique, and irrational}
\]

%dsp-level:0
\[
\qquad [(4),(5)];
\]

%dsp-level:0
\[
\qquad \textbf{the neutron is the heavier baryon on the whole opening region}
\]

%dsp-level:0
\[
\qquad [\mathrm{BM.62}(9)\ \text{re-cited}];
\]

%dsp-level:0
\[
(7)\quad\textbf{the fractional splittings (the exact rational enclosures):}
\]

%dsp-level:0
\[
\qquad \frac{\Delta m}{m_{\mathrm{hid}}}\bigg|_{t=0}=\frac{\ln(4/3)}{6\ln2}\in\Bigl(\frac1{15},\frac1{14}\Bigr)
\]

%dsp-level:0
\[
\qquad [\mathrm{BM.62}(10);\ \text{machine:}\ 2^{24}>3^{15},\ 2^{22}<3^{14}];
\]

%dsp-level:0
\[
\frac{\Delta m}{m_n}\bigg|_{t=0}=\frac{\ln(4/3)}{\ln1280}\in\Bigl(\frac1{25},\frac1{24}\Bigr)
\]

%dsp-level:0
\[
\qquad [\mathrm{BM.62}(10);\ \text{machine}];
\]

%dsp-level:0
\[
\qquad \boxed{\ \frac{\Delta m}{m_p}\bigg|_{t=0}=\frac{\ln(4/3)}{\ln960}\in\Bigl(\frac1{24},\frac1{23}\Bigr)\ }
\]

%dsp-level:0
\[
[\text{machine, derived here}:\quad24\ln\tfrac43>\ln960\iff42\ln2>25\ln3+\ln5\iff2^{42}>3^{25}\!\cdot\!5\ (4398046511104>4236443047215);
\]

%dsp-level:0
\[
23\ln\tfrac43<\ln960\iff40\ln2<24\ln3+\ln5\iff2^{40}<3^{24}\!\cdot\!5\ (1099511627776<1412147682405)];
\]

%dsp-level:0
\[
\textbf{the strict contraction against the barred measure price:}
\]

%dsp-level:0
\[
\qquad \frac{\ln(4/3)}{6\ln2}<\frac13\iff3\ln\tfrac43<2\ln2\iff16<27
\]

%dsp-level:0
\[
\qquad [\mathrm{Def\text{-}BM.13}(8);\ \mathrm{BM.62}(10);\ \text{machine}];
\]

%dsp-level:0
\[
\textbf{the consistency row:}
\]

%dsp-level:0
\[
\qquad \frac{m_n}{m_p}-1=\frac{\Delta m}{m_p}\in\Bigl(\frac1{24},\frac1{23}\Bigr)\ \ \text{at}\ t=0\quad\text{and}\quad\frac{m_n}{m_p}=\frac{\ln1280}{\ln960}\in\Bigl(\frac{26}{25},\frac{21}{20}\Bigr):
\]

%dsp-level:0
\[
\boxed{\ 1+\Bigl(\frac1{24},\frac1{23}\Bigr)=\Bigl(\frac{25}{24},\frac{24}{23}\Bigr)\ \subset\ \Bigl(\frac{26}{25},\frac{21}{20}\Bigr)\ }
\]

%dsp-level:0
\[
\qquad [\text{machine:}\ \tfrac{26}{25}<\tfrac{25}{24}\iff624<625;\quad\tfrac{24}{23}<\tfrac{21}{20}\iff480<483]:
\]

%dsp-level:0
\[
\qquad \text{the finer splitting bracket implies the coarser ratio bracket};
\]

%dsp-level:0
\[
(8)\quad\textbf{what the ratio is not:}
\]

%dsp-level:0
\[
\qquad \frac{m_n}{m_p}\ne\frac43\ \ \text{at every admissible coupling}
\]

%dsp-level:0
\[
\qquad [\mathrm{BM.83}(3);\ \text{machine: the exact solution}\ t=-\ln405/\ln16<0];
\]

%dsp-level:0
\[
\frac{m_n}{m_p}\ne\frac{\ln1280}{\ln960}\ \text{at}\ t>0
\]

%dsp-level:0
\[
\qquad [(3)\text{'s strict monotonicity}];
\]

%dsp-level:0
\[
\qquad \neg\text{a measured neutron--proton ratio is imported, compared or quarantined anywhere in this tier}
\]

%dsp-level:0
\[
\qquad [\mathrm{P5};\ \text{machine: zero external tokens}];
\]

%dsp-level:0
\[
\textbf{the numerical value of}\ e\ \textbf{remains open and does not enter:}
\]

%dsp-level:0
\[
\qquad \varepsilon_{\mathrm{EM}}=e^2/(4\pi)\ \text{is carried symbolically}
\]

%dsp-level:0
\[
\qquad [\mathrm{EM.25}(7),(9);\ \mathrm{BM.62}(12)\ \text{pattern}];
\]

%dsp-level:0
\[
(9)\quad\textbf{status:}
\]

%dsp-level:0
\[
\qquad \textbf{derived:}\ (1)\text{--}(8);
\]

%dsp-level:0
\[
\qquad \textbf{the sign theorem is conditional on the sub-critical branch, with the hypothesis on its face and the critical value exact}\ (5),
\]

%dsp-level:0
\[
(6)\ [\text{A §0.2a}];
\]

%dsp-level:0
\[
\neg\text{a P4 addition},
\]

%dsp-level:0
\[
\quad\neg\text{a new dimensionful scale}\ [\mathrm{P1}],
\]

%dsp-level:0
\[
\quad\neg\text{a route adopted}\ [\mathrm{BM.30}(5)],
\]

%dsp-level:0
\[
\quad\neg\text{a measured value}\ [\mathrm{P5}],
\]

%dsp-level:0
\[
\quad\text{§Cal-BM stays empty};
\]

%dsp-level:0
\[
\textbf{the open row:}\ \text{the numerical value of}\ e\ [\mathrm{EM.25}(7):\ \text{one calibration datum}]\ —\ \text{it decides which side of}\ \varepsilon^{\mathrm{crit}}\ \text{the corpus sits on, and no theorem of Parts 1--4 supplies it}
\]

%dsp-level:0
\[
\qquad [\mathrm{BM.62}(12)\ \text{re-cited}].
\]

\textbf{Req.} EM.25(7), EM.25(9) (Part 2); P1, P4, P5 (Parts 1–2); A §0.2a (Parts 1–2); FT.79(5) (Part 3); Def-BM.13(5), Def-BM.13(8), BM.30(5), BM.35(5), BM.37(2), BM.46(4), BM.51(6), BM.52(4), BM.52(5), BM.53(2), BM.54(1), BM.54(3), BM.58(4), BM.59(4), BM.62(1), BM.62(3), BM.62(5), BM.62(6), BM.62(7), BM.62(8), BM.62(9), BM.62(10), BM.62(12), BM.69(3), BM.71(1), BM.75(4), BM.75(5), BM.83(2), BM.83(3) (this part); exact prime-log-ring arithmetic with integer power certificates, symbolic Laurent differentiation, exact rational enclosures of $\ln$-ratios, unique-factorization irrationality arguments, sign censuses [math; machine: \texttt{checks/\allowbreak{}verify\_\allowbreak{}H\_\allowbreak{}mass\_\allowbreak{}ratio\_\allowbreak{}census.py} §G].

\textbf{Pf.}

%dsp-level:0
\[
(1):\ \text{the five zero-or-monomial rows re-cited from their own tags}\ [\mathrm{BM.51}(6);\ \mathrm{BM.53}(2);\ \mathrm{BM.54}(1)\text{--}(3);\ \mathrm{BM.52}(4),(5);\ \mathrm{BM.46}(4);\ \mathrm{BM.59}(4)];
\]

%dsp-level:0
\[
\qquad \text{the bridge row}:\ \Delta_g
\]

%dsp-level:0
\[
= \hat\mu'^2(D_S(+)-D_S(-))\ \text{has}\ [\Delta_g]
\]

%dsp-level:0
\[
= \mathfrak q\ \text{and}\ [\Delta m]
\]

%dsp-level:0
\[
= \mathfrak q^{+1/2}\ [\text{machine: the exponent arithmetic}]\implies\neg\text{addable}\ [\mathrm{BM.37}(2)\ \text{pattern}]
\]

%dsp-level:0
\[
(2):\ 6\ln2+\ln20=6\ln2+\ln4+\ln5=8\ln2+\ln5=\ln(2^8\cdot5)=\ln1280\ [\text{machine}];
\]

%dsp-level:0
\[
\qquad 6\ln2+\ln15=6\ln2+\ln3+\ln5=\ln(2^6\cdot15)=\ln960\ [\text{machine}];
\]

%dsp-level:0
\[
\qquad t=6c\,\varepsilon_{\mathrm{EM}}\ge0\ \text{since}\ c>0,\ \varepsilon_{\mathrm{EM}}=e^2/(4\pi)\ge0\ [\text{machine}]
\]

%dsp-level:0
\[
(3):\ \text{monotonicity by cross-multiplication:}\ \ t_1<t_2\implies\frac{A}{B+t_1\ln2}>\frac{A}{B+t_2\ln2}
\]

%dsp-level:0
\[
\iff A\bigl(B+t_2\ln2\bigr)>A\bigl(B+t_1\ln2\bigr)
\]

%dsp-level:0
\[
\iff A\,
\]

%dsp-level:0
\[
(t_2-t_1)\ln2>0\ \ \text{at}\ A
\]

%dsp-level:0
\[
= \ln1280,
\]

%dsp-level:0
\[
\ B\qquad = \ln960
\]

%dsp-level:0
\[
[\text{machine: the three positivity rows}\ \ln1280>0,\ \ln960>0,\ \ln2>0\ \text{in the prime-log ring, and the cross-multiplied comparison on a rational grid of pairs}\ (t_1,t_2)];
\]

%dsp-level:0
\[
\qquad \text{the derivative form}\ -A\ln2/(B+t\ln2)^2\ \text{is the same sign census};
\]

%dsp-level:0
\[
\qquad t=0\ \text{gives the supremum}
\]

%dsp-level:0
\[
25\ln1280-26\ln960
\]

%dsp-level:0
\[
= (200-156)\ln2-26\ln3+(25-26)\ln5
\]

%dsp-level:0
\[
= 44\ln2-26\ln3-\ln5>0
\]

%dsp-level:0
\[
\iff 2^{44}>3^{26}\cdot5\ [\text{machine: the two integers}];
\]

%dsp-level:0
\[
\qquad 20\ln1280-21\ln960=34\ln2-21\ln3-\ln5<0\iff2^{34}<3^{21}\cdot5\ [\text{machine}]
\]

%dsp-level:0
\[
(4):\ \frac{m_n}{m_p}
\]

%dsp-level:0
\[
= 1\iff\ln1280
\]

%dsp-level:0
\[
= \ln960+t\ln2
\]

%dsp-level:0
\[
\iff t
\]

%dsp-level:0
\[
= \frac{\ln1280-\ln960}{\ln2}
\]

%dsp-level:0
\[
= \frac{\ln(4/3)}{\ln2}
\]

%dsp-level:0
\[
= \frac{2\ln2-\ln3}{\ln2}
\]

%dsp-level:0
\[
= 2-\log_23\ [\text{machine: solved in}\ \mathfrak P];
\]

%dsp-level:0
\[
\log_23\in\mathbb Q\implies\exists\,(p,q)\ne(0,0):\ 2^p=3^q\ [\text{clearing denominators}];
\]

%dsp-level:0
\[
\qquad \text{machine: no solution with}\ |p|,|q|\le64;
\]

%dsp-level:0
\[
\qquad \text{math: unique factorization excludes every case}\ (p,q>0:\ \text{parity of the left side};\ p\le0\ \text{or}\ q\le0:\ \text{the reciprocal argument})
\]

%dsp-level:0
\[
(5):\ \Delta m=0\iff\ln(4/3)=t\iff6c\,\varepsilon_{\mathrm{EM}}=\ln(4/3)\iff\varepsilon_{\mathrm{EM}}=\frac{\ln(4/3)}{6c\ln2}\ [\text{machine}];
\]

%dsp-level:0
\[
\qquad \text{the four cross-multiplications of}\ \mathrm{BM.62}(7)\ \text{re-certified as the four integer power comparisons}\ [\text{machine}]
\]

%dsp-level:0
\[
(6):\ (3)\text{'s monotonicity}\ \wedge\ (4)\text{'s unique zero}\implies\text{the three-way sign equivalence}\ [\text{machine}];
\]

%dsp-level:0
\[
\qquad \text{the opening witnesses}\ 4^3=64>54=3^3\cdot2\ \text{and}\ 4^5=1024>972=3^5\cdot2^2\ [\text{machine}];
\]

%dsp-level:0
\[
\qquad \text{the closing witness}\ \tfrac16>\tfrac17>\varepsilon^{\mathrm{crit}}(\tfrac12)\ \text{and}\ \tfrac16>\tfrac18>\varepsilon^{\mathrm{crit}}(\tfrac35)\ [\text{machine:}\ 6<7,\ 6<8,\ (5)\text{'s brackets}]
\]

%dsp-level:0
\[
(7):\ \text{the two frozen enclosures re-certified}\ [\text{machine}];
\]

%dsp-level:0
\[
\qquad 24\ln\tfrac43-\ln960=(48-6)\ln2-(24+1)\ln3-\ln5=42\ln2-25\ln3-\ln5>0\iff2^{42}>3^{25}\cdot5\ [\text{machine}];
\]

%dsp-level:0
\[
\qquad 23\ln\tfrac43-\ln960=40\ln2-24\ln3-\ln5<0\iff2^{40}<3^{24}\cdot5\ [\text{machine}]
\]

%dsp-level:0
\[
3\ln\tfrac43<2\ln2\iff\ln\tfrac{64}{27}<\ln4\iff64<108\ [\text{machine}];
\]

%dsp-level:0
\[
\qquad 1+\tfrac1{24}
\]

%dsp-level:0
\[
= \tfrac{25}{24}\ \text{and}\ \tfrac{25}{24}<\tfrac{26}{25}
\]

%dsp-level:0
\[
\iff 625<624\ \text{is false}\implies\text{the two brackets overlap and are reported separately}\ [\text{machine}]
\]

%dsp-level:0
\[
(8):\ \mathrm{BM.83}(3)\text{'s exact solution re-cited};
\]

%dsp-level:0
\[
\qquad (3)\text{'s strict monotonicity};
\]

%dsp-level:0
\[
\qquad \text{the token scans}\ [\text{machine}];
\]

%dsp-level:0
\[
\qquad (9):\ \text{the standing-rule scans and the assembly}\ [\text{machine}].\ \blacksquare
\]

\textbf{Phys.} The neutron–proton mass ratio is one exact expression: the logarithm of one thousand two hundred eighty over the logarithm of nine hundred sixty plus a nonnegative multiple of the logarithm of two. Everything about it follows by arithmetic in the ring generated by the logarithms of the primes. It decreases strictly with the coupling, so its largest value is the coupling-free one, and that value sits between twenty-six twenty-fifths and twenty-one twentieths — an exact upper bound of one and one twentieth on the ratio, valid at every coupling the corpus admits. It is irrational at every rational coupling, by unique factorization of two, three and five inside the logarithms, and it equals one at exactly one coupling, which is two minus the base-two logarithm of three and is therefore irrational too: the degenerate point of the trichotomy is provably not a rational coupling, and at every rational coupling the two baryons differ. The critical strength inherits the same irrationality and the same four integer certificates that enclose it between a ninth and a seventh, so the opening region is a wide exact interval, not a tuned point, and the closing region is inhabited by anything above a sixth. The bridge term contributes nothing here and cannot: it lives in the angle currency, it points the same way, and adding it to a bit count is the one operation the ledger forbids. No measured ratio enters, no external value is compared, and the single open datum — the numerical value of the electric coupling — is the same one the corpus has named at every tier, deciding which side of the critical surface the world sits on without ever being supplied.

\end{theorem}

\begin{theorem}[{\texorpdfstring{BM.85 — The Q-value in two formed reads: the ratio forms, the margin ladders, and the read-relative verdict}{BM.85 - The Q-value in two formed reads: the ratio forms, the margin ladders, and the read-relative verdict}}]\label{tag:BM-85}
\textbf{Thm.}

%dsp-level:0
\[
(1)\quad\textbf{the Q-value on the informational currency (formed):}
\]

%dsp-level:0
\[
\qquad Q_\beta:=m_n-m_p-m_\ell-m_\nu=\Delta m-m_\ell-m_\nu
\]

%dsp-level:0
\[
\qquad [\mathrm{BM.41}(4);\ \mathrm{BM.49}(1);\ \mathrm{BM.63}(1)]:
\]

%dsp-level:0
\[
\boxed{\ Q_\beta\,\mathfrak h\ =\ \ln\frac43-t\ln2-F_{\ell\nu},\qquad F_{\ell\nu}:=m_\ell\mathfrak h+m_\nu\mathfrak h,\qquad t=6\,c\,\varepsilon_{\mathrm{EM}}\ge0\ }
\]

%dsp-level:0
\[
\qquad [\mathrm{BM.62}(4);\ \mathrm{BM.70}(1);\ \text{machine: the prime-log assembly}];
\]

%dsp-level:0
\[
(2)\quad\textbf{the two-bit daughter floor and the closure (re-certified):}
\]

%dsp-level:0
\[
\qquad F_{\ell\nu}\ge2\ln2
\]

%dsp-level:0
\[
\qquad [\mathrm{BM.70}(2),(3):\ \text{every materialized closed chain carries}\ \ge2\ \text{directed edges};\ m\mathfrak h=E_{\mathrm{int}}\ln2;\ \ m_\nu\mathfrak h\ge0\ \text{by PSD}]:
\]

%dsp-level:0
\[
\boxed{\ Q_\beta\,\mathfrak h\ \le\ \ln\frac43-2\ln2\ =\ \ln\frac43-\ln4\ =\ -\ln3\ <\ 0\ }
\]

%dsp-level:0
\[
\qquad [\text{machine:}\ 4<12];
\]

%dsp-level:0
\[
\qquad \textbf{formed}\implies\textbf{closed on the informational currency}
\]

%dsp-level:0
\[
\qquad [\mathrm{BM.70}(4)(\mathrm a)\ \text{re-cited}];
\]

%dsp-level:0
\[
\textbf{the margin ladder:}
\]

%dsp-level:0
\[
\qquad Q_\beta\mathfrak h\ \le\ -\ln3,\quad-\ln12,\quad-\ln24
\]

%dsp-level:0
\[
\qquad \text{on the certificates}\quad4<12,\quad4<48,\quad4<96
\]

%dsp-level:0
\[
\qquad [\mathrm{BM.70}(3);\ \text{machine}];
\]

%dsp-level:0
\[
\textbf{each margin is deepened by}\ t\ln2\ge0\ \ \textbf{at every positive coupling}\qquad[\text{machine: the sign census}];
\]

%dsp-level:0
\[
(3)\quad\textbf{the ratio forms (derived here):}
\]

%dsp-level:0
\[
\qquad \frac{Q_\beta}{m_n}\bigg|_{t=0}\ \le\ -\,\frac{\ln3}{\ln1280}\ \in\ \Bigl(-\frac2{13},-\frac3{20}\Bigr),
\]

%dsp-level:0
\[
\qquad \frac{Q_\beta}{m_p}\bigg|_{t=0}\ \le\ -\,\frac{\ln3}{\ln960}\ \in\ \Bigl(-\frac4{25},-\frac3{19}\Bigr)
\]

%dsp-level:0
\[
[\text{machine:}\ -\tfrac2{13}<-\tfrac{\ln3}{\ln1280}\iff13\ln3<2\ln1280\iff3^{13}<2^{16}\!\cdot\!5^2\ (1594323<1638400);
\]

%dsp-level:0
\[
-\tfrac{\ln3}{\ln1280}<-\tfrac3{20}\iff20\ln3>3\ln1280\iff3^{20}>2^{24}\!\cdot\!5^3\ (3486784401>2097152000);
\]

%dsp-level:0
\[
-\tfrac4{25}<-\tfrac{\ln3}{\ln960}\iff25\ln3<4\ln960\iff3^{21}<2^{24}\!\cdot\!5^4\ (10460353203<10485760000);
\]

%dsp-level:0
\[
-\tfrac{\ln3}{\ln960}<-\tfrac3{19}\iff19\ln3>3\ln960\iff3^{16}>2^{18}\!\cdot\!5^3\ (43046721>32768000)];
\]

%dsp-level:0
\[
\boxed{\ \textbf{the informational Q-value is a strictly negative exact fraction of the parent mass, bounded away from zero by a rational}\ }
\]

%dsp-level:0
\[
\qquad \text{in both normalizations};
\]

%dsp-level:0
\[
(4)\quad\textbf{the Q-value on the specified read (formed, branch-relative):}
\]

%dsp-level:0
\[
\qquad Q_\beta^{\mathrm{spec}}:=m_n^{(\mathrm{ii})}-m_p^{(\mathrm{ii})}-m_\ell-m_\nu
\]

%dsp-level:0
\[
\qquad [\mathrm{BM.77}(1),(2);\ \text{A-BM.8}]:
\]

%dsp-level:0
\[
\boxed{\ Q_\beta^{\mathrm{spec}}\,\mathfrak h\ =\ \Delta T(S)-\rho_\ell\sqrt{G\bigl(c(0)\bigr)}-\rho_\nu\sqrt{G\bigl(c(i_\nu)\bigr)},\qquad\Delta T(S)=5\,m^2_B(\cdot,S)\ }
\]

%dsp-level:0
\[
\Delta T=\Bigl(\frac{4219785}{98956},\ \frac{9608310}{393679},\ \frac{455}{23},\ \frac{68565}{3088}\Bigr)\ \ \text{at}\ S=(0,1,2,3),
\]

%dsp-level:0
\[
\qquad \Delta T_{\min}=\frac{455}{23},
\]

%dsp-level:0
\[
\qquad \Delta T_{\max}=\frac{4219785}{98956}
\]

%dsp-level:0
\[
\qquad [\mathrm{BM.77}(2);\ \text{machine}];
\]

%dsp-level:0
\[
\sqrt{A_0-B_0}\ \le\ \sqrt{G\bigl(c(i)\bigr)}\ \le\ \sqrt{A_0+B_0},
\]

%dsp-level:0
\[
\qquad A_0\pm B_0\big|_{g=g'=1}=\frac{11}{16},\ \frac3{16}
\]

%dsp-level:0
\[
\qquad [\mathrm{BM.39}(1);\ \mathrm{BM.77}(4);\ \text{machine}];
\]

%dsp-level:0
\[
\textbf{the trichotomy with integer-certified witnesses on both sides:}
\]

%dsp-level:0
\[
\qquad \rho_\ell=\rho_\nu=\frac{\Delta T(S)}{4\sqrt{A_0+B_0}}\in\mathcal O_S,
\]

%dsp-level:0
\[
\qquad \rho_\ell=\rho_\nu=\frac{\Delta T(S)}{\sqrt{A_0-B_0}}\in\mathcal C_S
\]

%dsp-level:0
\[
\qquad [\mathrm{BM.77}(5)(\mathrm a),(\mathrm b)]:
\]

%dsp-level:0
\[
Q_\beta^{\mathrm{spec}}\mathfrak h\ge\frac{455}{23}-\frac{\sqrt{11}}2>0\iff828100>5819;
\]

%dsp-level:0
\[
\qquad Q_\beta^{\mathrm{spec}}\mathfrak h\le\frac{4219785}{98956}-50\sqrt3<0\ \ \text{on}\ \ 4219785<43\cdot98956=4255108\ \ \text{and}\ \ 7500>1849;
\]

%dsp-level:0
\[
\frac{273}{23}\ <\ \rho^{\mathrm{crit}}(S,i_\nu)\ <\ \frac{42197850}{791648}
\]

%dsp-level:0
\[
\qquad \text{at the witness couplings}\ g=g'=1,\ \rho_\ell=\rho_\nu
\]

%dsp-level:0
\[
\qquad [\mathrm{BM.77}(5)(\mathrm c);\ \text{machine:}\ 99<100,\ 75>64,\ 273>253,\ 42197850<42748992];
\]

%dsp-level:0
\[
(5)\quad\textbf{the two reads are}\ \neg\textbf{comparable, and both verdicts stand:}\qquad\textbf{the alphabet census:}
\]

%dsp-level:0
\[
Q_\beta^{\mathrm{info}}:\ \{\ln2,\ln3,\ln5,\ c,\ \varepsilon_{\mathrm{EM}},\ \mathfrak h,\ m_\ell,\ m_\nu\};
\]

%dsp-level:0
\[
\qquad Q_\beta^{\mathrm{spec}}:\ \{m^2_B(\cdot,S),\ \rho_\ell,\ \rho_\nu,\ g,\ g',\ \mathfrak h\}
\]

%dsp-level:0
\[
\qquad [\text{machine: the token census}]:
\]

%dsp-level:0
\[
\boxed{\ \text{the two alphabets are disjoint except for}\ \mathfrak h\ }
\]

%dsp-level:0
\[
\qquad [\mathrm{BM.37}(2);\ \mathrm{BM.62}(11);\ \mathrm{BM.66}(6)\ \text{pattern};\ \text{machine}];
\]

%dsp-level:0
\[
\neg\text{a derived equation relates the two expressions}\ —\ \text{that is}\ (\mathrm{R1})\ \text{and}\ (\mathrm{R3})\ \text{and the Cal-F scales together}
\]

%dsp-level:0
\[
\qquad [\mathrm{BM.82}(6);\ \mathrm{BM.83}(5),(7)];
\]

%dsp-level:0
\[
\textbf{the verdicts, read-relative and complete:}
\]

%dsp-level:0
\[
\qquad \text{closed on the informational currency for every member}\ [\mathrm{BM.70}(4)(\mathrm a)];
\]

%dsp-level:0
\[
\qquad \text{open on}\ \mathcal O_S,\ \text{closed on}\ \mathcal C_S,\ \text{degenerate on}\ \partial\mathcal O_S\ \text{on the specified read}\ [\mathrm{BM.77}(5),(6)];
\]

%dsp-level:0
\[
\neg\text{a clause of}\ \mathrm{BM.41},\ \mathrm{BM.49},\ \mathrm{BM.63},\ \mathrm{BM.70},\ \mathrm{BM.71}\ \text{is retracted};
\]

%dsp-level:0
\[
\qquad \text{the rescopes are the ones}\ \mathrm{BM.77}(7)\ \text{records}
\]

%dsp-level:0
\[
\qquad [\text{machine: the clause census}];
\]

%dsp-level:0
\[
(6)\quad\textbf{the parent row is currency-independent:}
\]

%dsp-level:0
\[
\qquad T(0,S)=\tfrac43T(+1,S)>T(+1,S)
\]

%dsp-level:0
\[
\qquad [\mathrm{BM.28}(4)(\mathrm{ii});\ \text{machine}]
\]

%dsp-level:0
\[
\qquad \text{and}
\]

%dsp-level:0
\[
\qquad m_n>m_p\ \ \text{on the sub-critical branch}
\]

%dsp-level:0
\[
\qquad [\mathrm{BM.62}(9);\ \mathrm{BM.84}(6)]:
\]

%dsp-level:0
\[
\boxed{\ \textbf{the decay direction}\ n\to p+\ell+\nu\ \textbf{is fixed in both currencies}\ }
\]

%dsp-level:0
\[
\qquad [\mathrm{BM.41}(1);\ \mathrm{BM.45};\ \mathrm{BM.77}(6):\ \text{the grammar stands unmodified}];
\]

%dsp-level:0
\[
\textbf{the two-body companion on the specified read:}
\]

%dsp-level:0
\[
\qquad \Delta T(S)\ge\rho_\ell\sqrt{G\bigl(c(0)\bigr)}\ \ \text{holds at the witness}\ \Bigl[\tfrac{455}{23}>\tfrac{\sqrt{11}}4\iff3312400>5819\Bigr],
\]

%dsp-level:0
\[
\qquad \mathcal O_S\subseteq\{\text{two-body admissible}\}
\]

%dsp-level:0
\[
\qquad [\mathrm{BM.77}(6);\ \text{machine}];
\]

%dsp-level:0
\[
(7)\quad\textbf{the conditional spectrum becomes branch content on}\ \mathcal O_S\textbf{:}
\]

%dsp-level:0
\[
\qquad E_\ell^{\max}=\frac{m_n^2+m_\ell^2-(m_p+m_\nu)^2}{2m_n},
\]

%dsp-level:0
\[
\qquad E_{\max}:\langle E\rangle=2:1\ \ \text{scale-invariant}
\]

%dsp-level:0
\[
\qquad [\mathrm{BM.41}(5);\ \mathrm{BM.49}(3),(4);\ \mathrm{BM.77}(8)\ \text{re-cited}]:
\]

%dsp-level:0
\[
\neg\text{an external endpoint, spectrum or lifetime is compared}
\]

%dsp-level:0
\[
\qquad [\mathrm{P5}];
\]

%dsp-level:0
\[
\qquad \text{the quarantined audit rows stay quarantined}
\]

%dsp-level:0
\[
\qquad [\text{Census FT.114},\ \text{FT.120}\ \text{pattern}];
\]

%dsp-level:0
\[
(8)\quad\textbf{status:}
\]

%dsp-level:0
\[
\qquad \textbf{derived:}\ (1)\text{--}(7);
\]

%dsp-level:0
\[
\qquad \textbf{read-relative:}\ \text{the kinematic verdict}\ (5);
\]

%dsp-level:0
\[
\qquad \textbf{declared: nothing here}\ —\ \text{A-BM.8's specification is re-cited,}\ \neg\text{re-issued};
\]

%dsp-level:0
\[
\neg\text{a P4 addition},\quad\neg\text{a new dimensionful scale}\ [\mathrm{P1}],\quad\neg\text{a measured value}\ [\mathrm{P5}],\quad\neg\text{a currency converted};
\]

%dsp-level:0
\[
\qquad \text{§Cal-BM stays empty}.
\]

\textbf{Req.} P1, P4, P5 (Parts 1–2); FT.79(5) (Part 3); Census FT.114, FT.120 (row-wise); Def-BM.13(5), BM.28(4), BM.37(2), BM.39(1), BM.41(1), BM.41(4), BM.41(5), BM.45, BM.49(1), BM.49(3), BM.49(4), BM.62(4), BM.62(9), BM.62(11), BM.63(1), BM.66(6), BM.70(1), BM.70(2), BM.70(3), BM.70(4), BM.71, BM.77(1), BM.77(2), BM.77(4), BM.77(5), BM.77(6), BM.77(7), BM.77(8), BM.82(6), BM.83(5), BM.83(7), BM.84(6) (this part); the declaration addendum A-BM.8, item 7 (pointers); exact prime-log arithmetic with integer power certificates, exact rational arithmetic on the four frozen sector expectations, exact radical inequalities by integer certificates, token censuses [math; machine: \texttt{checks/\allowbreak{}verify\_\allowbreak{}H\_\allowbreak{}mass\_\allowbreak{}ratio\_\allowbreak{}census.py} §H].

\textbf{Pf.}

%dsp-level:0
\[
(1):\ \Delta m=(\ln\tfrac43-t\ln2)/\mathfrak h\ [\mathrm{BM.62}(4)\ \text{re-cited}];
\]

%dsp-level:0
\[
\qquad Q_\beta\mathfrak h=\ln\tfrac43-t\ln2-(m_\ell\mathfrak h+m_\nu\mathfrak h)\ [\text{machine: the monomial sum in}\ \mathfrak P]
\]

%dsp-level:0
\[
(2):\ F_{\ell\nu}\ge2\ln2\ [\mathrm{BM.70}(2),(3)\ \text{re-cited}];
\]

%dsp-level:0
\[
\qquad \ln\tfrac43-2\ln2=2\ln2-\ln3-2\ln2=-\ln3\ [\text{machine}];
\]

%dsp-level:0
\[
\qquad -\ln3<0\iff3>1\iff4<12\ [\text{machine}];
\]

%dsp-level:0
\[
\qquad \text{the ladder on}\ 4<12,\ 4<48,\ 4<96\ [\text{machine}]
\]

%dsp-level:0
\[
(3):\ Q_\beta/m_n=(Q_\beta\mathfrak h)/(m_n\mathfrak h)\le(-\ln3)/\ln1280\ \text{at}\ t=0\ [\text{machine:}\ m_n\mathfrak h=\ln1280>0];
\]

%dsp-level:0
\[
\qquad \text{the four enclosures by the four integer power comparisons of}\ (3)\ [\text{machine}]
\]

%dsp-level:0
\[
(4):\ \text{the four}\ \Delta T\ \text{values re-cited}\ [\mathrm{BM.77}(2);\ \text{machine:}\ 5\cdot\tfrac{843957}{98956},\ 5\cdot\tfrac{1921662}{393679},\ 5\cdot\tfrac{91}{23},\ 5\cdot\tfrac{13713}{3088}\ \text{as exact integers}];
\]

%dsp-level:0
\[
\qquad \text{the sandwich and both witnesses re-certified}\ [\mathrm{BM.39}(1);\ \mathrm{BM.77}(4),(5);\ \text{machine}]
\]

%dsp-level:0
\[
(5):\ \text{the token census of the two expressions}\ [\text{machine}];
\]

%dsp-level:0
\[
\qquad \text{the clause census of}\ \mathrm{BM.41}\text{--}\mathrm{BM.71}\ \text{against retraction}\ [\text{machine}];
\]

%dsp-level:0
\[
\qquad \text{the no-equation row re-cited from}\ \mathrm{BM.82}(6),\ \mathrm{BM.83}(5),(7)
\]

%dsp-level:0
\[
(6):\ T(0,S)/T(+1,S)=W(0)/W(+1)=20/15=4/3\ [\text{machine}];
\]

%dsp-level:0
\[
\qquad m_n>m_p\ \text{sub-critically}\ [\mathrm{BM.84}(6)];
\]

%dsp-level:0
\[
\qquad \tfrac{455}{23}>\tfrac{\sqrt{11}}4\iff1820>23\sqrt{11}\iff3312400>5819\ [\text{machine: squaring both positive sides}]
\]

%dsp-level:0
\[
(7):\ \text{the endpoint identity and the}\ 2{:}1\ \text{moment ratio re-cited}\ [\mathrm{BM.41}(5);\ \mathrm{BM.49}(3),(4)];
\]

%dsp-level:0
\[
\qquad (8):\ \text{the standing-rule scans}\ [\text{machine}].\ \blacksquare
\]

\textbf{Phys.} One Q-value, two formed reads, and no equation between them. Read in bits, the decay's energy balance is the transport register's log of four thirds, less the electromagnetic scar's coupling term, less whatever the two emitted lines cost — and the two-bit floor on any materialized chain drives the total below minus the logarithm of three at every coupling, so on that currency the vertex is closed with an exact margin, and the margin is a rational fraction of the parent's own mass: worse than minus three twentieths, better than minus two thirteenths. Read on the specified route, the same vertex is a difference of five times a sector expectation against two undeclared scale combinations, and its sign is decided by where those scales sit against an exact rational: open on a region certified by eight hundred twenty-eight thousand one hundred against five thousand eight hundred nineteen, closed on its complement, degenerate on the boundary between eleven and fifty-four. The two expressions share one symbol — h — and nothing else, so the two verdicts are not in conflict and cannot be reconciled without the very bridge the ledger holds open. What both currencies do agree on, and what no read can change, is the direction: the neutral class is the heavier one in the ensemble, on the informational branch below the critical coupling, and on the specified read at every winding cell, so the parent is the neutron and the grammar stands. The decay is allowed where a read is formed that allows it, closed where a read is formed that closes it, and the census reports both without converting either.

\end{theorem}

\medskip\noindent\textit{\texorpdfstring{BM.86 — The status of the mass ratio census: adjudication and scope}{BM.86 - The status of the mass ratio census: adjudication and scope}.}\ Structural-constraints record: see Appendix B, Section B.1.26 of the Addendum.\par\medskip

\phantomsection
\section[Tier BM.18 — The EM Self-Energy Correction: the Category…]{\texorpdfstring{Tier BM.18 — The EM Self-Energy Correction: the Category Adjudication, the Running Coupling, the Path-Integral Self-Energy, the Branch-Blindness Theorem, and the Corrected Neutron–Proton Splitting}{Tier BM.18 - The EM Self-Energy Correction: the Category Adjudication, the Running Coupling, the Path-Integral Self-Energy, the Branch-Blindness Theorem, and the Corrected Neutron-Proton Splitting}}

\textbf{Locus.} This tier relocates one computation and voids one coefficient. The EM self-energy of a sealed composite is taken off the chart — where it was read as an integral over a charge distribution $[\mathrm{Def\text{-}BM.12}(5);\ \mathrm{BM.46}(1)]$ — and put on the flicker ensemble, where the corpus's own charge datum lives: a ternary end read $q\in\{-1,0,+1\}$ on the endpoints of chains $[\mathrm{D3.4};\ \mathrm{EM.9}]$ with no position datum attached $[\mathrm{BM.28}(7);\ \mathrm{BM.46}(2)]$. The distribution factor $c$ that priced a shell against a ball is adjudicated a classical artifact and voided as a mass datum; the exterior value it carried is derived, twice, from rows that need no density ($\mathrm{BM.87}$). The coupling stops being a scalar and becomes an exact geometric profile on the envelope ratio, the lapse and the curvature lock ($\mathrm{BM.88}$). The self-energy becomes an ensemble average with one derived attenuation, the crystallization weight ($\mathrm{BM.89}$). The proposal that the six-state defect splitting reweight that attenuation is put on its face and refuted on four grounds ($\mathrm{BM.90}$). The neutron–proton splitting, its critical strength, its ratio range and its trichotomy are recomputed with the corrected term ($\mathrm{BM.91}$), the Q-value rows are re-certified on both formed reads ($\mathrm{BM.92}$), and $\mathrm{BM.93}$ adjudicates.

\textbf{Standing rows.} $\alpha_B$ stays barred on $\mathrm{BM.61}(8)$'s four grounds and §Cal-BM stays empty. No clause of $\mathrm{BM.1}$–$\mathrm{BM.86}$ is retracted, replaced or renumbered: every row this tier touches is re-scoped at its own locus and recorded in $\mathrm{BM.87}(6)$, $\mathrm{BM.91}(8)$ and $\mathrm{BM.93}(4)$, in $\mathrm{BM.77}(7)$'s pattern. The numerical value of $e$ stays open $[\mathrm{EM.25}(7);\ \mathrm{BM.62}(12);\ \mathrm{BM.84}(9)]$ and no measured coupling, mass or ratio enters any derivation $[\mathrm{P5}]$. The relating equation stays $\neg$derived — the mass-form bridge is application-layer content $[\mathrm{BM.35}(5);\ \mathrm{BM.69}(4),(5);\ \mathrm{BM.83}(5)]$. The currencies stay distinct: the ensemble weight derived here attenuates the informational register alone and converts nothing into the defect register $[\mathrm{BM.82}(4);\ \mathrm{BM.83}(3)]$. Local notation, fixed for the whole tier: $k:=r/\Delta x_{\mathrm B}$ the envelope ratio; $f(k)=1-\tfrac1k$ the nearby-space fraction; $N$ the lapse; $G:=(\Delta x_{\mathrm B})^2\kappa^2\rho$ the curvature lock; $w$ the crystallization weight; $\mathsf g$ the curvature composition (the italic $g$ keeps its $\mathrm{BM.39}(1)$ meaning, the $\mathsf{SU}(2)$ coupling, at every occurrence); $\varepsilon_{\mathrm{EM}}=e^2/(4\pi)$; the symbol $c$ is void as a mass datum and occurs in this tier only inside barred-object and rescope rows.

\medskip\noindent\textit{\texorpdfstring{BM.87 — The category adjudication: the void distribution factor, the derived exterior value, and the standing rows}{BM.87 - The category adjudication: the void distribution factor, the derived exterior value, and the standing rows}.}\ Structural-constraints record: see Appendix B, Section B.1.27 of the Addendum.\par\medskip

\begin{theorem}[{\texorpdfstring{BM.88 — The running coupling: the exact geometric profile on the envelope ratio, the lapse, and the curvature lock}{BM.88 - The running coupling: the exact geometric profile on the envelope ratio, the lapse, and the curvature lock}}]\label{tag:BM-88}
\textbf{Thm.}

%dsp-level:0
\[
(1)\quad\textbf{the coupling is a function,}\ \neg\textbf{a scalar:}
\]

%dsp-level:0
\[
\qquad e=\frac{gg'}{\sqrt{g^2+g'^2}}=g\sin\theta_W
\]

%dsp-level:0
\[
\qquad [\mathrm{EW.35};\ \mathrm{EM.25}(6):\ \text{fixed with zero freedom as a combination}];
\]

%dsp-level:0
\[
\neg\exists\ \text{a theorem fixing the numerical value of}\ e
\]

%dsp-level:0
\[
\qquad [\mathrm{EM.25}(7);\ \mathrm{P4}];
\]

%dsp-level:0
\[
\qquad \textbf{the running in this corpus is geometric,}\ \neg\textbf{perturbative:}
\]

%dsp-level:0
\[
\qquad \neg\exists\ \text{a loop, a}\ \beta\text{-function, a counterterm or a Yukawa-type object in the inventory}
\]

%dsp-level:0
\[
[\mathrm{BM.61}(8);\ \mathrm{FT.145}(3);\ \mathrm{EW.15};\ \mathrm{BM.47}(2):\ \text{the divergence is cut by the envelope,}\ \neg\text{by a counterterm}];
\]

%dsp-level:0
\[
(2)\quad\textbf{the three derived dimensionless arguments:}
\]

%dsp-level:0
\[
\qquad \textbf{(i) the envelope ratio:}\quad k\ :=\ \frac{r}{\Delta x_{\mathrm B}(E)},
\]

%dsp-level:0
\[
\qquad \Delta x_{\mathrm B}(E)=\frac1{m_{\mathrm{hid}}(E)}=\frac{\mathfrak h}{E_{\mathrm{int}}(E)\ln2}
\]

%dsp-level:0
\[
\qquad [\mathrm{Def\text{-}BM.4}(1);\ \mathrm{BM.13}(1);\ \text{machine:}\ \Delta x_{\mathrm B}\cdot m_{\mathrm{hid}}=1\ \text{as a Laurent identity}]:
\]

%dsp-level:0
\[
E_{\mathrm{int}}\in\{2,\,6,\,12,\,24,\,24,\,60,\,60\}\ \ \text{on the computed cells}
\]

%dsp-level:0
\[
\qquad [\mathrm{BM.13}(2);\ \mathrm{FT.28};\ \text{Census FT.35}(2),\ \text{FT.42}(5)\ \text{row-wise};\ \mathrm{BM.13}(3):\ E_{\mathrm{int}}\ge2]:
\]

%dsp-level:0
\[
\text{a strictly decreasing discrete ladder of widths},
\]

%dsp-level:0
\[
\qquad k\ge1\ \ \text{on the formed locus}\ r\ge\Delta x_{\mathrm B}
\]

%dsp-level:0
\[
\qquad [\text{machine: injective in}\ E_{\mathrm{int}}];
\]

%dsp-level:0
\[
\textbf{(ii) the nearby-space fraction:}
\]

%dsp-level:0
\[
\qquad \boxed{\ f(k)\ :=\ 1-\frac1k\ =\ 1-\frac{\Delta x_{\mathrm B}(E)}{r}\ }
\]

%dsp-level:0
\[
\qquad [\mathrm{BM.47}(6);\ \text{machine: exact rational functions}]:
\]

%dsp-level:0
\[
f(1)=0,\quad f(2)=\frac12,\quad f(3)=\frac23,\quad f(4)=\frac34,\quad f(5)=\frac45,\quad f(6)=\frac56;
\]

%dsp-level:0
\[
\qquad f\ \text{strictly increasing on}\ k\ge1,
\]

%dsp-level:0
\[
\qquad 0\le f(k)<1\ \ \forall k<\infty
\]

%dsp-level:0
\[
\qquad [\text{machine}];
\]

%dsp-level:0
\[
\textbf{(iii) the lapse:}
\]

%dsp-level:0
\[
\qquad N=\sqrt{\Omega^2+\gamma^{ij}N_iN_j},
\]

%dsp-level:0
\[
\qquad \gamma_{ij}=\Omega^2(\delta_{ij}+\varkappa q_{ij}),
\]

%dsp-level:0
\[
\qquad N_i=\Omega^2\varkappa\beta w_i
\]

%dsp-level:0
\[
\qquad [\mathrm{CT1.12};\ \mathrm{BM.34}(1);\ \mathrm{FT.157}(1):\ \text{the shift drag}]:
\]

%dsp-level:0
\[
\boxed{\ N^2\det\gamma=-\det g\ }
\]

%dsp-level:0
\[
\qquad [\mathrm{BM.34}(1);\ \text{machine: the exact polynomial identity over}\ \mathbb Q[\Omega^2,\varkappa,\beta,w_i,q_{ij}]];
\]

%dsp-level:0
\[
(a_{00},q_{ij},w_i)=0\implies(\Omega,N,N^i,g_{ij})=(1,1,0,\delta_{ij})
\]

%dsp-level:0
\[
\qquad [\mathrm{BM.34}(2):\ \textbf{the vacuum's emergent metric is flat};\ \text{machine: exact substitution}]\implies\boxed{\ N=1\ \ \text{in the vacuum gauge}\ }
\]

%dsp-level:0
\[
(3)\quad\textbf{the curvature lock (the derived dimensionless curvature argument):}
\]

%dsp-level:0
\[
\qquad \kappa^2\rho=\frac{2\pi E_{\mathrm{int}}^4(\ln2)^3}{\mathfrak h^2},
\]

%dsp-level:0
\[
\qquad \rho=\frac{(E_{\mathrm{int}}\ln2)^4}{\mathfrak h^4},
\]

%dsp-level:0
\[
\qquad [\rho]=\mathfrak q^{+2}
\]

%dsp-level:0
\[
\qquad [\mathrm{BM.15}(2);\ \mathrm{BM.14}(2),(4);\ \mathrm{CT1.23}]:
\]

%dsp-level:0
\[
\boxed{\ G(E)\ :=\ (\Delta x_{\mathrm B})^2\cdot\kappa^2\rho\ =\ 2\pi\,E_{\mathrm{int}}(E)^2\ln2\ \in\ \mathfrak q^0\ }
\]

%dsp-level:0
\[
\qquad [\mathrm{BM.16}(1)\text{'s lock identity re-derived};\ \text{machine: one Laurent monomial, coefficient}\ 2]:
\]

%dsp-level:0
\[
\textbf{the three freedom censuses:}
\]

%dsp-level:0
\[
\qquad \mathfrak h\text{-free}\ [\text{the}\ \mathfrak h\text{-exponent is}\ 0],
\]

%dsp-level:0
\[
\qquad \Lambda\text{-free},
\]

%dsp-level:0
\[
\qquad \mathrm{P4}\text{-free}\ [\neg\text{one of}\ g,g',h_0,c_W,\Lambda,\varkappa,\beta,\alpha_\Omega\ \text{occurs}]
\]

%dsp-level:0
\[
\qquad [\mathrm{BM.16}(1);\ \text{machine}];
\]

%dsp-level:0
\[
G\ \text{is strictly increasing and quadratic in}\ E_{\mathrm{int}}\ \text{on the quantized ladder}
\]

%dsp-level:0
\[
\qquad [\text{machine: the seven cell values}];
\]

%dsp-level:0
\[
(4)\quad\textbf{the exact profile:}
\]

%dsp-level:0
\[
\qquad \text{the self-energy's dimensionless factor is the product of the three derived arguments' functions, and the energy-side monomial is where the lapse enters:}
\]

%dsp-level:0
\[
\boxed{\ \varepsilon_{\mathrm{EM}}(r,E,N)\ :=\ N(E)\cdot f\!\Bigl(\frac{r}{\Delta x_{\mathrm B}(E)}\Bigr)\cdot \mathsf g\bigl(G(E)\bigr)\cdot\varepsilon_{\mathrm{EM}}\ }
\]

%dsp-level:0
\[
\qquad [\text{machine: the monomial census};\ \ \mathrm{BM.34}(4)\text{'s dispersion}\ (p_0-N^ip_i)^2=N^2(m^2+\gamma^{ij}p_ip_j)\ \text{puts exactly one factor of}\ N\ \text{on the energy}];
\]

%dsp-level:0
\[
\boxed{\ e(r,E,N)\ =\ e_0\,\sqrt{\,N(E)\cdot f\bigl(r/\Delta x_{\mathrm B}(E)\bigr)\cdot \mathsf g\bigl(G(E)\bigr)\,}\ }
\]

%dsp-level:0
\[
\qquad [e^2=4\pi\varepsilon_{\mathrm{EM}};\ \text{machine: squaring recovers}\ (4)\text{'s monomial}];
\]

%dsp-level:0
\[
\textbf{the placement row:}
\]

%dsp-level:0
\[
\qquad N\ \text{on the energy and}\ \sqrt N\ \text{on the coupling are the same monomial}\ —\ \text{the placement is a naming convention,}\ \neg\text{a datum};
\]

%dsp-level:0
\[
\qquad \text{the derived content is}\ (4)\text{'s profile};
\]

%dsp-level:0
\[
\textbf{the curvature composition is priced,}\ \neg\textbf{adopted:}
\]

%dsp-level:0
\[
\qquad \mathsf g\in\Bigl\{1,\ \ 1+G,\ \ \frac1{1+G}\Bigr\},
\]

%dsp-level:0
\[
\qquad \mathsf g\ \text{carried symbolically};
\]

%dsp-level:0
\[
\qquad \neg\text{a clause of this tier selects a composition}
\]

%dsp-level:0
\[
[\text{the argument}\ G\ \text{is derived}\ (3);\ \text{the composition is}\ \neg\text{derived}: \neg\exists\ \text{a count}\to\text{strength theorem}\ [\mathrm{BM.36}(3);\ \mathrm{BM.61}(8)(\mathrm{O1})],\ \neg\exists\ \text{a loop}\ (1)];
\]

%dsp-level:0
\[
(5)\quad\textbf{the exact form at the two loci that matter:}
\]

%dsp-level:0
\[
\qquad \textbf{at the member's own envelope}\ k=1:
\]

%dsp-level:0
\[
\qquad \boxed{\ f(1)=0\implies e(\Delta x_{\mathrm B},E,N)=0\ \ \text{and}\ \ \varepsilon_{\mathrm{EM}}(\Delta x_{\mathrm B},E,N)=0\ }
\]

%dsp-level:0
\[
\textbf{the coupling vanishes at the envelope, so the formed self-energy is the exterior read alone}
\]

%dsp-level:0
\[
\qquad [\mathrm{BM.87}(2),(3):\ \text{the same statement on the energy side}];
\]

%dsp-level:0
\[
\textbf{in the far field}\ k\to\infty:
\]

%dsp-level:0
\[
\qquad f\to1,
\]

%dsp-level:0
\[
\qquad \varepsilon_{\mathrm{EM}}\to N\,\mathsf g(G)\,\varepsilon_{\mathrm{EM}},
\]

%dsp-level:0
\[
\qquad E_{\mathrm{EM}}^{\mathrm{self}}\to\frac12\,N\,\mathsf g(G)\,\varepsilon_{\mathrm{EM}}\,q^2\,m_{\mathrm{hid}}
\]

%dsp-level:0
\[
\qquad [\text{machine: the exact limit row}];
\]

%dsp-level:0
\[
\textbf{no finite neighbourhood holds the whole correction:}
\]

%dsp-level:0
\[
\qquad 1-f(k)=\frac1k>0\ \ \forall k<\infty
\]

%dsp-level:0
\[
\qquad [\mathrm{BM.47}(6);\ \text{machine}];
\]

%dsp-level:0
\[
(6)\quad\textbf{the energy dependence (the running, exactly):}
\]

%dsp-level:0
\[
\qquad \text{at fixed}\ r,
\]

%dsp-level:0
\[
\qquad f\bigl(k(E)\bigr)=1-\frac{\mathfrak h}{r\,E_{\mathrm{int}}(E)\ln2}
\]

%dsp-level:0
\[
\qquad [\text{machine: the substitution}]:
\]

%dsp-level:0
\[
\boxed{\ E_{\mathrm{int}}\ \text{strictly increasing}\ \implies\ \Delta x_{\mathrm B}\ \text{strictly decreasing}\ \implies\ k\ \text{strictly increasing}\ \implies\ f\ \text{strictly increasing}\ }
\]

%dsp-level:0
\[
\qquad [\text{machine: three monotonicity rows on exact rationals}];
\]

%dsp-level:0
\[
\textbf{the coupling grows with the energy scale, geometrically, with no loop in the argument};
\]

%dsp-level:0
\[
\qquad \text{the six-state floors give the exact ladder}\ f=1-\mathfrak h/(r\,E_{\mathrm{int}}\ln2)\ \ \text{at}\ \ E_{\mathrm{int}}\in\{12,24,60\}
\]

%dsp-level:0
\[
\qquad [\mathrm{BM.81}(1);\ \text{machine}];
\]

%dsp-level:0
\[
(7)\quad\textbf{the bars:}
\]

%dsp-level:0
\[
\qquad \neg\exists\ \text{a numerical value of}\ e\ \text{or of}\ e_0\ \text{in any clause}
\]

%dsp-level:0
\[
\qquad [\mathrm{EM.25}(7);\ \mathrm{P5};\ \text{machine: zero external tokens, zero decimal numerals in this tier}];
\]

%dsp-level:0
\[
\neg\text{a measured coupling is identified with any corpus quantity}\ [\mathrm{P5}];
\]

%dsp-level:0
\[
\qquad \varepsilon_{\mathrm{EM}}
\]

%dsp-level:0
\[
= e^2/(4\pi)\ \text{is a dimensionless P4-class symbol carried symbolically throughout}\ [\mathrm{Def\text{-}BM.12}\ \text{local notation};\ \mathrm{EM.25}(9)];
\]

%dsp-level:0
\[
\alpha_B:\quad\neg\text{added},\quad\neg\text{derived},\quad\text{barred on}\ \mathrm{BM.61}(8)\text{'s four grounds}
\]

%dsp-level:0
\[
\qquad [\text{machine: the token scan of this tier}];
\]

%dsp-level:0
\[
\qquad \neg\text{a P4 addition},\quad\text{§Cal-BM stays empty};
\]

%dsp-level:0
\[
(8)\quad\textbf{status:}
\]

%dsp-level:0
\[
\qquad \textbf{derived:}\ (1)\text{--}(3),\ (5),\ (6);
\]

%dsp-level:0
\[
\qquad \textbf{derived in form with a priced composition:}\ (4);
\]

%dsp-level:0
\[
\qquad \textbf{symbolic:}\ e_0,\ g;
\]

%dsp-level:0
\[
\neg\text{a declaration},\quad\neg\text{a new dimensionful scale}\ [\mathrm{P1}],\quad\neg\text{a measured value}\ [\mathrm{P5}],\quad
\]

%dsp-level:0
\[
\qquad \textbf{open:}\ \text{the value of}\ e\ [\mathrm{EM.25}(7)],
\]

%dsp-level:0
\[
\ \text{the curvature composition}\ \mathsf g\ (4),
\]

%dsp-level:0
\[
\ \text{the lapse as solution data}\ [\mathrm{CT2.7};\ \mathrm{BM.34}(2)].
\]

\textbf{Req.} D1.5, D2.9 (Parts 1–2); CT1.12, CT1.22, CT1.23, CT2.1, CT2.7, EM.9, EM.15, EM.25(6), EM.25(7), EM.25(9), EW.15, EW.35 (Part 2); P1, P4, P5 (Parts 1–2); FT.28, FT.145(3), FT.157(1) (Part 3); Census FT.35(2), FT.42(5) (row-wise); Def-BM.4(1), BM.13(1), BM.13(2), BM.13(3), BM.14(2), BM.14(4), BM.15(2), BM.16(1), BM.34(1), BM.34(2), BM.34(4), BM.36(3), BM.47(2), BM.47(6), BM.61(8), BM.81(1), BM.87(2), BM.87(3) (this part); exact Laurent-monomial arithmetic with $\mathfrak q$- and $\mathfrak h$-exponent censuses, exact rational function rows, monotonicity on exact rationals, the ADM polynomial identity by Schur complement, limit rows [math; machine: \texttt{checks/\allowbreak{}verify\_\allowbreak{}H\_\allowbreak{}em\_\allowbreak{}self\_\allowbreak{}energy\_\allowbreak{}correction.py} §B].

\textbf{Pf.}

%dsp-level:0
\[
(1):\ \mathrm{EW.35}\text{'s combination and}\ \mathrm{EM.25}(6),(7)\text{'s rows re-cited};
\]

%dsp-level:0
\[
\qquad \neg\exists\ \text{a loop object in the inventory}\ [\mathrm{BM.61}(8);\ \mathrm{FT.145}(3);\ \text{machine: the token census}]
\]

%dsp-level:0
\[
(2)(\mathrm i):\ \Delta x_{\mathrm B}=1/m_{\mathrm{hid}}=\mathfrak h/(E_{\mathrm{int}}\ln2)\ [\text{machine: one Laurent monomial}];
\]

%dsp-level:0
\[
\qquad E_{\mathrm{int}}\ge2\ [\mathrm{BM.13}(3)];
\]

%dsp-level:0
\[
\qquad \text{the ladder is injective in}\ E_{\mathrm{int}}\ \text{and strictly decreasing}\ [\text{machine}]
\]

%dsp-level:0
\[
(\mathrm{ii}):\ f(k)=1-1/k\ \text{at}\ k=1,\dots,6=(0,\tfrac12,\tfrac23,\tfrac34,\tfrac45,\tfrac56)\ [\text{machine: exact Fractions}];
\]

%dsp-level:0
\[
\qquad f'(k)=1/k^2>0\ \text{on}\ k\ge1\ \text{and}\ f(k)<1\iff1/k>0\ [\text{machine: the rational-function rows}];
\]

%dsp-level:0
\[
\qquad f(1)=0\ [\text{machine}]
\]

%dsp-level:0
\[
(\mathrm{iii}):\ N^2=\Omega^2+\gamma^{ij}N_iN_j\ \text{and}\ N^2\det\gamma=-\det g\ [\mathrm{BM.34}(1)\ \text{re-cited};\ \text{machine: the Schur identity}];
\]

%dsp-level:0
\[
\qquad (a_{00},q_{ij},w_i)=0\implies\Omega=1,\ N_i=0,\ N=1\ [\text{machine: exact substitution}]
\]

%dsp-level:0
\[
(3):\ (\Delta x_{\mathrm B})^2\kappa^2\rho
\]

%dsp-level:0
\[
= \frac{\mathfrak h^2}{E_{\mathrm{int}}^2(\ln2)^2}\cdot\frac{2\pi E_{\mathrm{int}}^4(\ln2)^3}{\mathfrak h^2}
\]

%dsp-level:0
\[
= 2\pi E_{\mathrm{int}}^2\ln2\ [\text{machine: the Laurent monomial cancels}\ \mathfrak h\ \text{exactly}];
\]

%dsp-level:0
\[
\qquad [G]=\mathfrak q^{-1}\cdot\mathfrak q^{+1}=\mathfrak q^0\ [\text{machine: the exponent arithmetic}];
\]

%dsp-level:0
\[
\qquad \text{the three freedom censuses}\ [\mathrm{BM.16}(1);\ \text{machine: the symbol scan}]
\]

%dsp-level:0
\[
(4):\ \text{the product of three dimensionless factors is dimensionless}\ [\text{machine}];
\]

%dsp-level:0
\[
\qquad e^2=4\pi\varepsilon_{\mathrm{EM}}\implies e=e_0\sqrt{Nf\mathsf g}\ \text{at}\ \varepsilon_{\mathrm{EM}}=Nf\mathsf g\,\varepsilon_{\mathrm{EM}}\ [\text{machine: squaring recovers the monomial}];
\]

%dsp-level:0
\[
\qquad \text{the one-factor-of-}N\ \text{row}:\ (p_0-N^ip_i)^2
\]

%dsp-level:0
\[
= N^2(m^2+\gamma^{ij}p_ip_j)
\]

%dsp-level:0
\[
\implies E
\]

%dsp-level:0
\[
= N\sqrt{m^2+p^2}-N^ip_i\ [\mathrm{BM.34}(4);\ \text{machine}]
\]

%dsp-level:0
\[
\text{the priced set}\ \{1,1+G,(1+G)^{-1}\}\ \text{is a pricing,}\ \neg\text{an adoption}\ [\mathrm{Def\text{-}BM.12}(4)\ \text{pattern}];
\]

%dsp-level:0
\[
\qquad \neg\exists\ \text{a count}\to\text{strength theorem}\ [\mathrm{BM.36}(3);\ \mathrm{BM.61}(8)(\mathrm{O1})]
\]

%dsp-level:0
\[
(5):\ f(1)=1-1=0\ [\text{machine}];
\]

%dsp-level:0
\[
\qquad \lim_{k\to\infty}(1-1/k)=1\ [\text{machine: the exact limit row on rationals}];
\]

%dsp-level:0
\[
\qquad 1-f(k)=1/k>0\ \forall k<\infty\ [\text{machine}]
\]

%dsp-level:0
\[
(6):\ k(E)
\]

%dsp-level:0
\[
= rE_{\mathrm{int}}(E)\ln2/\mathfrak h\ \text{is strictly increasing in}\ E_{\mathrm{int}}\ \text{at fixed}\ r>0\ [\text{machine: positive coefficients}];
\]

%dsp-level:0
\[
\qquad f\ \text{strictly increasing}\ [(2)(\mathrm{ii})]\implies\text{the composition is strictly increasing}\ [\text{machine}];
\]

%dsp-level:0
\[
\qquad \Delta x_{\mathrm B}\ \text{strictly decreasing}\ [\text{machine}]
\]

%dsp-level:0
\[
(7)\text{--}(8):\ \text{the token scans, the symbol census and the assembly}\ [\text{machine}].\ \blacksquare
\]

\textbf{Phys.} A coupling with no scale at which to sit is not a number, and the corpus never pretended otherwise: it fixed the combination of the two gauge strengths and left the value open. What it can do is say how the coupling's \textit{profile} runs, and every ingredient of that profile is already derived. Distance enters through the only length the ledger owns — the Born width, which is the reciprocal of an edge count — and through the fraction of an inverse-square flux enclosed within a given number of those widths, which is one minus one over k, exactly. Time enters through the lapse of the emergent metric, which is one in the vacuum gauge and departs from one only where the shift and the conformal coefficient do. Curvature enters through the one dimensionless monomial the calibration lock supplies: the Born width squared times the Ricci scale, equal to two pi times the square of the edge count times the logarithm of two, free of h, free of the cosmological term, free of every declared parameter. Multiply the three and the coupling is an exact algebraic function of position, energy and lapse — vanishing at the member's own envelope, growing without bound toward its far-field value, never quite reaching it at any finite radius. The composition of the curvature factor is not derived, because nothing in the inventory turns a count into a strength, so it is priced three ways and adopted not at all. The running comes out in the direction a loop calculation would give it, from geometry alone, with no loop in the argument and no number fitted anywhere.

\end{theorem}

\begin{theorem}[{\texorpdfstring{BM.89 — The EM self-energy from the path integral: the derived measure, the crystallization weight, and the route-relative correction}{BM.89 - The EM self-energy from the path integral: the derived measure, the crystallization weight, and the route-relative correction}}]\label{tag:BM-89}
\textbf{Thm.}

%dsp-level:0
\[
(1)\quad\textbf{the ensemble and its derived measure:}\qquad Q=\begin{pmatrix}-x_+&x_+\\ x_-&-x_-\end{pmatrix}\ \ \text{over}\ \mathbb Z[x_+,x_-],\qquad(x_-,\,x_+)\,Q=(0,0)\ \ \text{exactly},\qquad\dim\ker Q=1\qquad[\mathrm{BM.19}(1);\ \mathrm{FT.108}(4);\ \text{machine}]:
\]

%dsp-level:0
\[
\boxed{\ p\ :=\ \frac{x_+}{x_++x_-}\ \ \text{is the unique stationary measure}\ }
\]

%dsp-level:0
\[
\qquad x_+=1-4\varepsilon_q,\quad x_-=4\varepsilon_q\ \ \implies\ \ p=1-4\varepsilon_q
\]

%dsp-level:0
\[
\qquad [\mathrm{BM.19}(2);\ \text{machine: an identity, exact at every rational threshold}];
\]

%dsp-level:0
\[
\varepsilon_q\in[0,\tfrac14)\ [\mathrm{CT0.8};\ \mathrm{D0.4};\ \text{A4}]\implies x_-\in[0,1),\quad x_+\in(0,1],\quad p\in(0,1]\qquad[\text{machine}];
\]

%dsp-level:0
\[
\varepsilon_q=0\implies x_-=0\implies\mathcal E(\mathcal X)=\{\text{the crystallized configuration}\}
\]

%dsp-level:0
\[
\qquad [\mathrm{BM.19}(3);\ \mathrm{Def\text{-}BM.5}(5)];
\]

%dsp-level:0
\[
\qquad |\mathcal E(\mathcal X)|>1\iff\varepsilon_q>0
\]

%dsp-level:0
\[
\qquad [\mathrm{BM.19}(6)];
\]

%dsp-level:0
\[
\textbf{the ensemble is finite at every budget and}\ \neg\textbf{a fixed topology at any nonzero threshold}
\]

%dsp-level:0
\[
\qquad [\mathrm{Def\text{-}H.32}(1);\ \mathrm{BM.19}(4),(6)];
\]

%dsp-level:0
\[
(2)\quad\textbf{the configuration energy:}
\]

%dsp-level:0
\[
\qquad \boxed{\ E_{\mathrm{EM}}(R)\ =\ \frac12\,N(R)\,\mathsf g\bigl(G(R)\bigr)\,\varepsilon_{\mathrm{EM}}\,Q(R)^2\,m_{\mathrm{hid}}(R),\qquad m_{\mathrm{hid}}(R)=\frac{E_{\mathrm{int}}(R)\ln2}{\mathfrak h}=\frac{|R|\ln2}{\mathfrak h}\ }
\]

%dsp-level:0
\[
[\mathrm{BM.87}(4)\text{'s formed monomial at}\ c=\tfrac12;\ \mathrm{BM.88}(4)\text{'s profile};\ \mathrm{Def\text{-}BM.5}(4):\ E_{\mathrm{int}}(R)=|R|;\ \mathrm{FT.44}(4):\ \text{the taxonomy law};\ \text{machine}];
\]

%dsp-level:0
\[
(3)\quad\textbf{the class factorization:}
\]

%dsp-level:0
\[
\qquad Q(R)=Q\ \ \text{on every member of a charge class}\ —\ \text{the boundary data fix the class}
\]

%dsp-level:0
\[
\qquad [\mathrm{Def\text{-}BM.5}(1);\ \mathrm{Def\text{-}BM.8}(1);\ \text{machine: the class census}]:
\]

%dsp-level:0
\[
\boxed{\ E_{\mathrm{EM}}^{\mathrm{self}}(Q)\ =\ \sum_{R\in\mathcal E(\mathcal X)}P(R)\,E_{\mathrm{EM}}(R)\ =\ \frac12\,N\,\mathsf g(G)\,\varepsilon_{\mathrm{EM}}\,Q^2\,\bigl\langle m_{\mathrm{hid}}\bigr\rangle_{\mathcal E}\ }
\]

%dsp-level:0
\[
\qquad [\text{machine:}\ Q^2\ \text{is constant on the class and factors out of the average}];
\]

%dsp-level:0
\[
(4)\quad\textbf{the crystallization weight (the derived attenuation):}
\]

%dsp-level:0
\[
\qquad \boxed{\ w(\varepsilon_q)\ :=\ \frac{\langle E_{\mathrm{int}}\rangle_{\mathcal E}}{E_{\mathrm{int}}^{\mathrm{cryst}}}\ =\ \frac1{E_{\mathrm{int}}^{\mathrm{cryst}}}\sum_{\text{edges}}P(\text{the edge is real})\ }
\]

%dsp-level:0
\[
\qquad [\mathrm{Def\text{-}BM.5}(4);\ \text{the ensemble average of}\ |R|]:
\]

%dsp-level:0
\[
\boxed{\ w(0)=1\ \ \text{exactly}\ }
\]

%dsp-level:0
\[
\qquad [\mathrm{BM.19}(3):\ \text{the ensemble is the single crystallized configuration}];
\]

%dsp-level:0
\[
\qquad 0<w(\varepsilon_q)\le1
\]

%dsp-level:0
\[
\qquad [\mathrm{BM.19}(5),(6):\ \text{the ensemble spreads}\ E_{\mathrm{int}}\ \text{over}\ \{0,\dots,B_K\}\ \text{with the crystallized value as its top point, and every move changes}\ E_{\mathrm{int}}\ \text{by its}\ J^\pm\ \text{excess}];
\]

%dsp-level:0
\[
\textbf{the strict decrease is derived on the marked specialization below, and in general is}\ (\mathrm{R4})\textbf{'s content:}
\]

%dsp-level:0
\[
\qquad \varepsilon_q\ \text{enters the measure only through}\ x_-=4\varepsilon_q\ [\mathrm{BM.19}(2);\ \text{machine: the symbol census}];
\]

%dsp-level:0
\[
\textbf{the Bernoulli specialization (marked,}\ \neg\textbf{adopted):}
\]

%dsp-level:0
\[
\qquad \text{on the per-datum stationary retention of}\ (1),
\]

%dsp-level:0
\[
\qquad \langle E_{\mathrm{int}}\rangle_{\mathcal E}=E_{\mathrm{int}}^{\mathrm{cryst}}\cdot p\implies\boxed{\ w=1-4\varepsilon_q\ }
\]

%dsp-level:0
\[
[\text{machine:}\ E_{\mathrm{int}}^{\mathrm{cryst}}=6=(c_1{:}\,2)+(c_2{:}\,2)+(\text{connector}{:}\,2)\ [\mathrm{BM.28}(3);\ \mathrm{BM.17};\ \mathrm{BM.26}]\ \text{is three pairs, each retained with weight}\ p];
\]

%dsp-level:0
\[
\qquad w\ \text{is carried symbolically wherever the specialization is not invoked};
\]

%dsp-level:0
\[
\textbf{the member is a stationary mixture of real and virtual content:}
\]

%dsp-level:0
\[
\qquad \text{the composite is where the real and the virtual registers meet,}\ \neg\text{a fixed topology}
\]

%dsp-level:0
\[
\qquad [\mathrm{BM.19}(6);\ \mathrm{D1.5};\ \mathrm{CT1.22}];
\]

%dsp-level:0
\[
(5)\quad\textbf{the corrected self-energy:}
\]

%dsp-level:0
\[
\qquad \boxed{\ E_{\mathrm{EM}}^{\mathrm{self}}(Q)\ =\ \frac12\,N\,w\,\mathsf g(G)\,\varepsilon_{\mathrm{EM}}\,Q^2\,\frac{E_{\mathrm{int}}^{\mathrm{cryst}}\ln2}{\mathfrak h}\ }
\]

%dsp-level:0
\[
\qquad [\text{machine: one Laurent monomial in}\ (N,w,\mathsf g,\varepsilon_{\mathrm{EM}},Q,E_{\mathrm{int}},\ln2,\mathfrak h)]:
\]

%dsp-level:0
\[
\text{at}\ E_{\mathrm{int}}^{\mathrm{cryst}}=6:
\]

%dsp-level:0
\[
\qquad E_{\mathrm{EM}}^{\mathrm{self}}(\pm1)=\frac{3\,N\,w\,\mathsf g\,\varepsilon_{\mathrm{EM}}\ln2}{\mathfrak h},
\]

%dsp-level:0
\[
\qquad E_{\mathrm{EM}}^{\mathrm{self}}(0)=0\ \ \textbf{exactly}
\]

%dsp-level:0
\[
\qquad [\text{machine:}\ 0^2=0;\ \text{Gauss: zero net flux}];
\]

%dsp-level:0
\[
\textbf{the crystallized vacuum-gauge specialization:}
\]

%dsp-level:0
\[
\qquad (w,N,\mathsf g)=(1,1,1)\implies E_{\mathrm{EM}}^{\mathrm{self}}(\pm1)=\frac{3\,\varepsilon_{\mathrm{EM}}\ln2}{\mathfrak h}=\frac12\,\varepsilon_{\mathrm{EM}}\,m_{\mathrm{hid}}=\frac{e^2q^2}{8\pi\,\Delta x_{\mathrm B}}
\]

%dsp-level:0
\[
\boxed{\ \textbf{the path-integral value reproduces}\ \mathrm{BM.87}(2)\textbf{'s derived exterior value exactly at full crystallization in the vacuum gauge}\ }
\]

%dsp-level:0
\[
\qquad [\text{machine: the two monomials agree coefficient by coefficient}];
\]

%dsp-level:0
\[
(6)\quad\textbf{the route-relativity (priced,}\ \neg\textbf{adopted):}
\]

%dsp-level:0
\[
\qquad \textbf{route}\ (\mathrm i)\ \text{(the normalized average):}\quad\overline{\Delta S}(\cdot,S)=m^2_B(\cdot,S)\ \text{is charge-degenerate}
\]

%dsp-level:0
\[
\qquad [\mathrm{FT.107}(3);\ \mathrm{BM.28}(4)(\mathrm i)]:
\]

%dsp-level:0
\[
\text{the correction splits the classes only through}\ Q^2,\ \text{and the ledger's own average is class-blind}
\]

%dsp-level:0
\[
\qquad [\text{machine: the frozen sector expectations are}\ q\text{-independent}];
\]

%dsp-level:0
\[
\textbf{route}\ (\mathrm{ii})\ \text{(the unnormalized total deficit):}
\]

%dsp-level:0
\[
\qquad T(Q,S)=W(Q)\,\overline{\Delta S}(S),
\]

%dsp-level:0
\[
\qquad W=(1,6,15,\mathbf{20},15,6,1)
\]

%dsp-level:0
\[
\qquad [\mathrm{BM.28}(4)(\mathrm{ii});\ \mathrm{Def\text{-}BM.13}(3)]:
\]

%dsp-level:0
\[
E_{\mathrm{EM}}^{\mathrm{self},(\mathrm{ii})}(Q,S)
\]

%dsp-level:0
\[
= \frac12\,
\]

%dsp-level:0
\[
N\,
\]

%dsp-level:0
\[
\mathsf g\,
\]

%dsp-level:0
\[
\varepsilon_{\mathrm{EM}}\,
\]

%dsp-level:0
\[
Q^2\,
\]

%dsp-level:0
\[
\frac{T(Q,S)}{\mathfrak h}\ \ \implies\ \ \Delta m_{\mathrm{EM}}^{(\mathrm{ii})}(S)
\]

%dsp-level:0
\[
= \frac{N\,\mathsf g\,\varepsilon_{\mathrm{EM}}}{2\mathfrak h}\cdot15\,
\]

%dsp-level:0
\[
m^2_B(\cdot,S)
\]

%dsp-level:0
\[
= \frac{N\,\mathsf g\,\varepsilon_{\mathrm{EM}}}{2\mathfrak h}\Bigl(\frac{12659355}{98956},\ \frac{28824930}{393679},\ \frac{1365}{23},\ \frac{205695}{3088}\Bigr)
\]

%dsp-level:0
\[
[\text{machine: exact rational arithmetic at the four frozen sector expectations}\ \bigl(\tfrac{843957}{98956},\tfrac{1921662}{393679},\tfrac{91}{23},\tfrac{13713}{3088}\bigr)\ [\mathrm{BM.21}(2);\ \mathrm{BM.62}(2)]];
\]

%dsp-level:0
\[
\textbf{route}\ (\mathrm{ii})\ \text{is}\ \neg\textbf{adopted for any purpose beyond the vertex kinematics}
\]

%dsp-level:0
\[
\qquad [\mathrm{BM.30}(5);\ \mathrm{BM.37}(6);\ \mathrm{BM.77}(1),(9)\ \text{stand}];
\]

%dsp-level:0
\[
\qquad \text{route}\ (\mathrm{iii})\text{'s amplitude ratios stay barred as masses}
\]

%dsp-level:0
\[
\qquad [\mathrm{FT.79}(5);\ \mathrm{BM.37}(3)];
\]

%dsp-level:0
\[
(7)\quad\textbf{the attenuation theorem:}
\]

%dsp-level:0
\[
\qquad 0<w\le1\ \ \text{and}\ \ N>0,\ \mathsf g>0\implies\boxed{\ \bigl|E_{\mathrm{EM}}^{\mathrm{self}}(Q)\bigr|\ \le\ \bigl|E_{\mathrm{EM}}^{\mathrm{self}}(Q)\bigr|_{w=1}\ }
\]

%dsp-level:0
\[
\qquad \text{with equality only at}\ w=1:
\]

%dsp-level:0
\[
\textbf{the flicker strictly attenuates the EM correction at every nonzero threshold};
\]

%dsp-level:0
\[
\qquad \text{the transport register's lead}\ \ln(4/3)/\mathfrak h\ \text{is class-shared and}\ \neg\text{attenuated by}\ w\ \ [\mathrm{BM.57}(4);\ \mathrm{BM.59}(4);\ \mathrm{BM.91}(2)];
\]

%dsp-level:0
\[
\implies\textbf{the sub-critical branch widens as the crystallization weight falls}
\]

%dsp-level:0
\[
\qquad [\mathrm{BM.91}(5);\ \text{machine: monotone in}\ w];
\]

%dsp-level:0
\[
(8)\quad\textbf{the residue:}
\]

%dsp-level:0
\[
\qquad w(\varepsilon_q)\ \text{in general requires the ensemble's configuration distribution beyond the per-datum specialization}\ —\ \textbf{named}\ (\mathrm{R4}),
\]

%dsp-level:0
\[
\ \neg\textbf{filled}:
\]

%dsp-level:0
\[
\text{the computed cells supply exact sector averages of the deficit register}\ \overline{\Delta S}
\]

%dsp-level:0
\[
= (\tfrac32 n_{\mathcal V})\ \ [\mathrm{FT.106}(3);\ \mathrm{FT.107}(3)],
\]

%dsp-level:0
\[
\ \ \neg\text{of the edge-count register}\ E_{\mathrm{int}}\qquad = |R|:
\]

%dsp-level:0
\[
\textbf{the two registers are distinct alphabets and are}\ \neg\textbf{equated}
\]

%dsp-level:0
\[
\qquad [\mathrm{BM.62}(11);\ \mathrm{BM.37}(2);\ \mathrm{BM.83}(3);\ \text{machine: the token census}];
\]

%dsp-level:0
\[
(9)\quad\textbf{status:}
\]

%dsp-level:0
\[
\qquad \textbf{derived:}\ (1)\text{--}(5),\ (7);
\]

%dsp-level:0
\[
\qquad \textbf{priced,}\ \neg\textbf{adopted:}\ (4)\text{'s Bernoulli specialization and}\ (6)\text{'s route rows};
\]

%dsp-level:0
\[
\qquad \textbf{named,}\ \neg\textbf{filled:}\ (8)\text{'s}\ (\mathrm{R4});
\]

%dsp-level:0
\[
\neg\text{a declaration},
\]

%dsp-level:0
\[
\quad\neg\text{a P4 addition},
\]

%dsp-level:0
\[
\quad\neg\text{a new dimensionful scale}\ [\mathrm{P1}],
\]

%dsp-level:0
\[
\quad\neg\text{a measured value}\ [\mathrm{P5}],
\]

%dsp-level:0
\[
\quad\text{§Cal-BM stays empty};\qquad \qquad
\]

\textbf{Req.} D1.5, D3.4 (Parts 1–2); A4, CT0.8, CT1.22, D0.4, EM.9 (Part 2); P1, P4, P5 (Parts 1–2); Def-H.32(1), FT.44(4), FT.79(5), FT.106(3), FT.107(3), FT.108(4) (Part 3); Def-BM.4(1), Def-BM.5(1), Def-BM.5(4), Def-BM.5(5), Def-BM.8(1), Def-BM.13(3), BM.17, BM.19(1), BM.19(2), BM.19(3), BM.19(4), BM.19(5), BM.19(6), BM.21(2), BM.26, BM.28(3), BM.28(4), BM.30(5), BM.37(2), BM.37(3), BM.37(6), BM.57(4), BM.59(4), BM.62(2), BM.62(11), BM.77(1), BM.77(9), BM.83(3), BM.87(2), BM.87(4), BM.88(4) (this part); exact polynomial arithmetic over $\mathbb Q[x_+,x_-]$, exhaustive integer kernel censuses, exact rational specialization, Laurent-monomial assembly with coefficient comparison, exact rational arithmetic at the four frozen sector expectations, token censuses [math; machine: \texttt{checks/\allowbreak{}verify\_\allowbreak{}H\_\allowbreak{}em\_\allowbreak{}self\_\allowbreak{}energy\_\allowbreak{}correction.py} §C].

\textbf{Pf.}

%dsp-level:0
\[
(1):\ (x_-)(-x_+)+(x_+)(x_-)=0\ \text{and}\ (x_-)(x_+)+(x_+)(-x_-)=0\ \text{in}\ \mathbb Q[x_+,x_-]\ [\text{machine}];
\]

%dsp-level:0
\[
\qquad \text{constant vectors}:\ a(-x_+)+bx_-=0\ \wedge\ ax_++b(-x_-)=0\implies a=b=0\ [\text{machine: exhaustive on}\ [-3,3]^2];
\]

%dsp-level:0
\[
\qquad p=(1-4\varepsilon_q)/(1-4\varepsilon_q+4\varepsilon_q)=1-4\varepsilon_q\ [\text{machine: exact at six rational thresholds}]
\]

%dsp-level:0
\[
\varepsilon_q\in[0,\tfrac14)\implies4\varepsilon_q\in[0,1)\ [\text{machine}];
\]

%dsp-level:0
\[
\qquad x_-=0\ \text{at}\ \varepsilon_q=0\implies\neg\exists\ \text{a dissolution move}\ [\mathrm{D1.2}:\ J^-:n_{\mathcal R}:1\to0];
\]

%dsp-level:0
\[
\qquad |\mathcal E|>1\iff x_->0\ [\mathrm{BM.19}(6)]
\]

%dsp-level:0
\[
(2)\text{--}(3):\ \mathrm{BM.87}(4)\text{'s monomial at}\ c=\tfrac12\ \text{and}\ \mathrm{BM.88}(4)\text{'s profile composed}\ [\text{machine}];
\]

%dsp-level:0
\[
\qquad E_{\mathrm{int}}(R)=|R|\ [\mathrm{Def\text{-}BM.5}(4)];
\]

%dsp-level:0
\[
\qquad Q(R)
\]

%dsp-level:0
\[
= Q\ \text{class-constant}\ [\mathrm{Def\text{-}BM.5}(1);\ \text{machine: the boundary-data census}]
\]

%dsp-level:0
\[
\implies Q^2\ \text{factors out of}\ \sum_RP(R)\cdot
\]

%dsp-level:0
\[
(4):\ \langle E_{\mathrm{int}}\rangle_{\mathcal E}=\sum_{\text{edges}}P(\text{real})\ [\text{linearity of the average on}\ |R|];
\]

%dsp-level:0
\[
\qquad \varepsilon_q=0:\ \mathcal E=\{R_{\mathrm{cryst}}\}\implies\langle E_{\mathrm{int}}\rangle=E_{\mathrm{int}}^{\mathrm{cryst}}\implies w=1\ [\text{machine}];
\]

%dsp-level:0
\[
\qquad 0<|R|\le E_{\mathrm{int}}^{\mathrm{cryst}}\ \text{on the ensemble's support}\implies0<w\le1\ [\mathrm{BM.19}(5);\ \text{machine}]
\]

%dsp-level:0
\[
\text{the specialization}:\ \text{three pairs, each retained with the stationary weight}\ p\implies\langle E_{\mathrm{int}}\rangle
\]

%dsp-level:0
\[
= 3\cdot2\cdot p
\]

%dsp-level:0
\[
= 6p\ \text{and}\ w
\]

%dsp-level:0
\[
= p
\]

%dsp-level:0
\[
= 1-4\varepsilon_q\ [\text{machine: exact rationals}];
\]

%dsp-level:0
\[
\qquad E_{\mathrm{int}}^{\mathrm{cryst}}=6=(2)+(2)+(2)\ [\mathrm{BM.28}(3);\ \text{machine}]
\]

%dsp-level:0
\[
(5):\ \tfrac12\cdot6=\ 3\ [\text{machine}];
\]

%dsp-level:0
\[
\qquad Q=0\implies Q^2=0\implies E=0\ [\text{machine}];
\]

%dsp-level:0
\[
\qquad (w,N,\mathsf g)
\]

%dsp-level:0
\[
= (1,1,1):\ \tfrac{3\varepsilon_{\mathrm{EM}}\ln2}{\mathfrak h}
\]

%dsp-level:0
\[
= \tfrac12\varepsilon_{\mathrm{EM}}\tfrac{6\ln2}{\mathfrak h}
\]

%dsp-level:0
\[
= \tfrac12\varepsilon_{\mathrm{EM}}m_{\mathrm{hid}}
\]

%dsp-level:0
\[
= \tfrac{e^2q^2}{8\pi\Delta x_{\mathrm B}}\ [\text{machine: three Laurent monomials, coefficient by coefficient},\ \Delta x_{\mathrm B}m_{\mathrm{hid}}=1]
\]

%dsp-level:0
\[
(6):\ m^2_B(q,S)=m^2_B(-q,S)\ [\mathrm{FT.107}(3);\ \text{machine}];
\]

%dsp-level:0
\[
\qquad 15\cdot\tfrac{843957}{98956}=\tfrac{12659355}{98956},
\]

%dsp-level:0
\[
\quad 15\cdot\tfrac{1921662}{393679}=\tfrac{28824930}{393679},
\]

%dsp-level:0
\[
\quad 15\cdot\tfrac{91}{23}=\tfrac{1365}{23},
\]

%dsp-level:0
\[
\quad 15\cdot\tfrac{13713}{3088}
\]

%dsp-level:0
\[
= \tfrac{205695}{3088}\ [\text{machine: exact Fractions, each in lowest terms or reduced exactly}]
\]

%dsp-level:0
\[
(7):\ 0<w
\]

%dsp-level:0
\[
\le 1\ \text{and}\ N,
\]

%dsp-level:0
\[
\mathsf g>0\implies\tfrac12Nwg\varepsilon_{\mathrm{EM}}Q^2m_{\mathrm{hid}}\le\tfrac12Ng\varepsilon_{\mathrm{EM}}Q^2m_{\mathrm{hid}}\ [\text{machine: the coefficient census}];
\]

%dsp-level:0
\[
\qquad \text{equality}\iff w=1\ [\text{machine}]
\]

%dsp-level:0
\[
(8)\text{--}(9):\ \text{the register-alphabet census}\ [\mathrm{BM.62}(11)\ \text{pattern};\ \text{machine}];
\]

%dsp-level:0
\[
\qquad \text{the standing-rule scans}\ [\text{machine}].\ \blacksquare
\]

\textbf{Phys.} The composite was never a fixed topology, and the correction that rides it should not have been either. Take the flicker's own measure — the unique stationary line of the Janus generator, which specializes to one minus four times the crystallization threshold with no rate declared anywhere — average the formed monomial over the ensemble, and the charge factors straight out, because the boundary data fix it before the flicker starts. What is left is one derived number between zero and one: the crystallization weight, the mean fraction of the composite's registration data that are real rather than virtual. It is exactly one at the sharp threshold, where the ensemble collapses to the single crystallized pattern, and it falls as the threshold rises, so the electromagnetic correction is attenuated by exactly the amount the member fails to be crystallized. At full crystallization in the vacuum gauge the average reproduces the exterior value the adjudication derived, monomial for monomial, so nothing is lost and one classical coefficient is gone. The baryon, then, is where the real and the virtual registers meet: a stationary mixture whose mass is an average, whose charge is exact, and whose field energy is a fraction of what a fully frozen lump would carry. The general weight beyond the per-datum specialization stays open and named, because the computed cells average the vertex-deficit register and not the edge count, and those two alphabets are not the same currency.

\end{theorem}

\begin{theorem}[{\texorpdfstring{BM.90 — The defect-conditioned weighting: the hypothesis on its face, and the branch-blindness theorem}{BM.90 - The defect-conditioned weighting: the hypothesis on its face, and the branch-blindness theorem}}]\label{tag:BM-90}
\textbf{Thm.}

%dsp-level:0
\[
(1)\quad\textbf{the hypothesis, on its face:}
\]

%dsp-level:0
\[
\qquad \textbf{hypothesis (W):}\ \ \text{the retention probability of a registration datum depends on the state's defect value},
\]

%dsp-level:0
\[
\quad P(\text{the edge is real}\mid D_S),
\]

%dsp-level:0
\[
\text{so that}
\]

%dsp-level:0
\[
\qquad m_{\mathrm{hid}}^{\mathrm{eff}}(s)\ =\ \frac{\ln2}{\mathfrak h}\sum_{\text{edges}}P(\text{the edge is real}\mid D_S(s))
\]

%dsp-level:0
\[
\qquad [\text{the six states of}\ \mathrm{Def\text{-}BM.14}(2),(3);\ \mathrm{BM.81}(1)];
\]

%dsp-level:0
\[
(2)\quad\textbf{the measure's alphabet census:}\qquad x_+=1-4\varepsilon_q,\qquad x_-=4\varepsilon_q,\qquad p=1-4\varepsilon_q\qquad[\mathrm{BM.19}(2);\ \text{machine}]:
\]

%dsp-level:0
\[
\boxed{\ \text{the flicker measure's symbol set is}\ \{\varepsilon_q\}\ \ \text{and}\ \ \varepsilon_q\ \text{is the A4 threshold}\ }
\]

%dsp-level:0
\[
\qquad [\mathrm{D0.4};\ \mathrm{BM.19}(2):\ \neg\text{a P-class datum enters};\ \neg\text{a new fugacity},\ \neg\text{a new rate}];
\]

%dsp-level:0
\[
\text{the defect value's symbol set is}\ \{\delta_v,\ \hat\mu'^2\}
\]

%dsp-level:0
\[
\qquad [\mathrm{Def\text{-}BM.14}(3);\ \mathrm{Def\text{-}H.37}(2);\ \text{machine}];
\]

%dsp-level:0
\[
\qquad \boxed{\ \{\varepsilon_q\}\ \cap\ \{\delta_v,\hat\mu'^2\}\ =\ \varnothing\ }
\]

%dsp-level:0
\[
(3)\quad\textbf{the branch-blindness theorem:}
\]

%dsp-level:0
\[
\qquad \boxed{\ \neg\exists\ \text{a derived coupling}\ D_S\longrightarrow P(\text{retention});\qquad w=w(\varepsilon_q)\ \ \text{is branch-blind}\ }
\]

%dsp-level:0
\[
\text{on four independent grounds, each re-cited at its own locus:}
\]

%dsp-level:0
\[
(\mathrm i)\quad\textbf{the field degree:}\ \ \text{the six defect values generate a degree-}16\ \text{extension of}\ \mathbb Q,
\]

%dsp-level:0
\[
\ \text{while every read the flicker measure is built from is integer- or binary-valued}
\]

%dsp-level:0
\[
\qquad [\mathrm{FT.73}(2),(3);\ \mathrm{BM.69}(4)(\mathrm i);\ \mathrm{BM.81}(1)]:
\]

%dsp-level:0
\[
\text{a rational formula places every retention weight in a field of degree}<16\ \text{over the declared constants}
\]

%dsp-level:0
\[
\qquad [\text{machine re-cited}];
\]

%dsp-level:0
\[
(\mathrm{ii})\quad\textbf{the two-torsion image:}\ \ \text{the register lift of the defect angle reaches}\ \{0,\pi\}\ \text{only}\ —\ \text{four of the six states are unreachable at every winding}
\]

%dsp-level:0
\[
\qquad [\mathrm{BM.67}(6);\ \mathrm{BM.69}(4)(\mathrm{ii})];
\]

%dsp-level:0
\[
(\mathrm{iii})\quad\textbf{the calibration weight:}\ \ \text{the winding route}\ \Theta_R
\]

%dsp-level:0
\[
= 2\pi k\ \text{carries}\ \mathfrak h\text{-weight}\ {-}2\ \text{against}\ 0\ \text{for a rate}
\]

%dsp-level:0
\[
\qquad [\mathrm{FT.55}(1);\ \mathrm{BM.67}(8);\ \mathrm{BM.69}(4)(\mathrm{iii})];
\]

%dsp-level:0
\[
(\mathrm{iv})\quad\textbf{the strength bar:}\ \ \neg\exists\ \text{a count}\to\text{strength theorem};
\]

%dsp-level:0
\[
\qquad \text{the capacity ratios}\ \bigl(\tfrac13,\,8,\,\tfrac1{25}\bigr)\ \text{are}\ \neg\text{couplings}
\]

%dsp-level:0
\[
\qquad [\mathrm{BM.36}(3);\ \mathrm{BM.61}(8)(\mathrm{O1});\ \mathrm{BM.69}(4)(\mathrm{iv})];
\]

%dsp-level:0
\[
(4)\quad\textbf{the consequence for the informational register:}
\]

%dsp-level:0
\[
\qquad \boxed{\ \textbf{the six-state quark splitting does not reweight}\ m_{\mathrm{hid}}\ }
\]

%dsp-level:0
\[
\qquad \text{the}\ m_{\mathrm{hid}}\ \text{contributions are the values}\ \mathrm{BM.81}(1)\ \text{censused},
\]

%dsp-level:0
\[
\ \text{attenuated by a branch-independent}\ w:
\]

%dsp-level:0
\[
m_{\mathrm{hid}}^{\mathrm{eff}}(s)=w(\varepsilon_q)\cdot\frac{2E(s)\ln2}{\mathfrak h}
\]

%dsp-level:0
\[
\qquad \text{with the same}\ w\ \text{on all six states}
\]

%dsp-level:0
\[
\qquad [\text{machine:}\ w\text{'s symbol census carries no}\ D_S];
\]

%dsp-level:0
\[
\textbf{the defect values enter the pole currency alone}
\]

%dsp-level:0
\[
\qquad [\mathrm{BM.82}(4):\ \text{the thirty}\ \rho_{\mathrm B}\text{-free pole ratios};\ \mathrm{BM.81}(1)\text{'s}\ D_S\ \text{column}];
\]

%dsp-level:0
\[
(5)\quad\textbf{what the splitting does change (the positive companion, in its own currency):}
\]

%dsp-level:0
\[
\qquad \Delta_g=\bigl(2,\ \sqrt2-1,\ 2\mathsf c-\sqrt3\bigr)\hat\mu'^2>0\ \ \text{per generation}
\]

%dsp-level:0
\[
\qquad [\mathrm{BM.75}(4);\ \mathrm{BM.81}(3)]:
\]

%dsp-level:0
\[
\mathcal R_s=\frac{D_S(s)}{4E(s)^2}\ \ \text{per state}
\]

%dsp-level:0
\[
\qquad [\mathrm{BM.82}(2)];
\]

%dsp-level:0
\[
\qquad \Delta\bigl(m^2_{\mathrm{pole}}\bigr)=\Delta_g>0\ \ \text{under hypothesis}\ (\mathrm H)\ \text{(the single substitution)}
\]

%dsp-level:0
\[
\qquad [\mathrm{BM.59}(9);\ \mathrm{BM.69}(3);\ \text{the hypothesis on its face}];
\]

%dsp-level:0
\[
\textbf{and}\ \neg\textbf{addable to the informational assembly:}
\]

%dsp-level:0
\[
\qquad [\Delta_g]=\mathfrak q\ \ \text{against}\ \ [\Delta m]=\mathfrak q^{+1/2},
\]

%dsp-level:0
\[
\qquad \text{a cross-currency addition being a category the corpus bars}
\]

%dsp-level:0
\[
[\mathrm{BM.83}(1),(3);\ \mathrm{BM.35}(5);\ \mathrm{BM.37}(2);\ \mathrm{BM.84}(1);\ \text{machine: the exponent arithmetic}];
\]

%dsp-level:0
\[
(6)\quad\textbf{the priced alternative:}
\]

%dsp-level:0
\[
\qquad \text{a declared}\ D_S\text{-dependent retention rate would be a P4-class count}\to\text{strength coupling}\ —\ \text{the exact species}\ \mathrm{BM.61}(8)\ \text{bars and}\ \mathrm{BM.36}(3)\ \text{refuses}:
\]

%dsp-level:0
\[
\boxed{\ \neg\text{such a datum is issued};\qquad\text{§Cal-BM stays empty};\qquad\neg\text{a P4 addition}\ }\qquad[\mathrm{P4};\ \mathrm{P1}];
\]

%dsp-level:0
\[
(7)\quad\textbf{status:}
\]

%dsp-level:0
\[
\qquad \textbf{derived:}\ (2)\text{--}(5);
\]

%dsp-level:0
\[
\qquad \textbf{refuted:}\ \text{hypothesis}\ (\mathrm W)\ \text{as a derived coupling}\ (3);
\]

%dsp-level:0
\[
\qquad \textbf{priced,}\ \neg\textbf{adopted:}\ (6);
\]

%dsp-level:0
\[
\neg\text{a declaration},\quad\neg\text{a measured value}\ [\mathrm{P5}],\quad\neg\text{a new dimensionful scale}\ [\mathrm{P1}],\quad
\]

%dsp-level:0
\[
\qquad \textbf{the currencies remain distinct and unequated}\ [\mathrm{BM.83}(3)].
\]

\textbf{Req.} D0.4, D3.4 (Parts 1–2); CT0.8, EM.9 (Part 2); P1, P4, P5 (Parts 1–2); Def-H.37(2), FT.55(1), FT.73(2), FT.73(3) (Part 3); Def-BM.14(2), Def-BM.14(3), BM.19(2), BM.35(5), BM.36(3), BM.37(2), BM.59(9), BM.61(8), BM.67(6), BM.67(8), BM.69(3), BM.69(4), BM.75(4), BM.81(1), BM.81(3), BM.82(2), BM.82(4), BM.83(1), BM.83(3), BM.84(1) (this part); exact symbol and token censuses, field-degree rows re-cited, exponent arithmetic in the $\mathfrak q$-grading [math; machine: \texttt{checks/\allowbreak{}verify\_\allowbreak{}H\_\allowbreak{}em\_\allowbreak{}self\_\allowbreak{}energy\_\allowbreak{}correction.py} §D].

\textbf{Pf.}

%dsp-level:0
\[
(1):\ \text{the hypothesis stated with its own symbols}\ [\text{no premise imported}];
\]

%dsp-level:0
\[
\qquad (2):\ x_\pm\ \text{and}\ p\ \text{re-cited from}\ \mathrm{BM.19}(2)\ [\text{machine}];
\]

%dsp-level:0
\[
\qquad \text{the two symbol sets enumerated from their defining clauses}\ [\text{machine: the token census of}\ \mathrm{Def\text{-}BM.14}(3),\ \mathrm{Def\text{-}H.37}(2)\ \text{against}\ \mathrm{BM.19}(2),\ \mathrm{D0.4}]
\]

%dsp-level:0
\[
(3):\ \text{the four grounds re-cited at their own loci}\ [\mathrm{FT.73}(2),(3);\ \mathrm{BM.67}(6),(8);\ \mathrm{FT.55}(1);\ \mathrm{BM.36}(3);\ \mathrm{BM.61}(8)(\mathrm{O1});\ \mathrm{BM.69}(4)];
\]

%dsp-level:0
\[
\qquad \text{the intersection}\ \{\varepsilon_q\}\cap\{\delta_v,\hat\mu'^2\}=\varnothing\ [\text{machine}]
\]

%dsp-level:0
\[
(4):\ w\text{'s defining expression carries no}\ D_S\ [(2),(3);\ \text{machine}];
\]

%dsp-level:0
\[
\qquad m_{\mathrm{hid}}(s)
\]

%dsp-level:0
\[
= 2E(s)\ln2/\mathfrak h\ [\mathrm{BM.81}(1);\ \mathrm{FT.44}(4)]
\]

%dsp-level:0
\[
\implies m^{\mathrm{eff}}_{\mathrm{hid}}(s)
\]

%dsp-level:0
\[
= w\cdot2E(s)\ln2/\mathfrak h\ \text{with one}\ w\ [\text{machine}]
\]

%dsp-level:0
\[
(5):\ \Delta_g\ \text{and}\ \mathcal R_s\ \text{re-cited from their own tags};
\]

%dsp-level:0
\[
\qquad [\Delta_g]=[\hat\mu'^2]=\mathfrak q,\quad[\Delta m]=\mathfrak q^{+1/2}\ [\mathrm{CT2.1};\ \text{machine}]\implies\neg\text{addable}\ [\mathrm{BM.37}(2)\ \text{pattern}]
\]

%dsp-level:0
\[
(6)\text{--}(7):\ \text{the species row of a count}\to\text{strength coupling}\ [\mathrm{BM.36}(3);\ \mathrm{BM.61}(8)];
\]

%dsp-level:0
\[
\qquad \text{the standing-rule scans}\ [\text{machine}].\ \blacksquare
\]

\textbf{Phys.} The proposal that six defect values should reweight six retention probabilities is natural and it is false here — not because the numbers disagree but because nothing in the inventory can connect them. The flicker measure knows one symbol, the crystallization threshold, and it is derived as the unique stationary line of a two-state generator. The defect values know two symbols, an angle and a scale, and they live in a degree-sixteen field over the rationals while every read the measure touches is an integer or a binary weight. Between them stand four separate obstructions the corpus already proved: a field-degree gap, a two-element holonomy image that cannot reach four of the six states, a calibration weight that forbids the winding route, and the standing bar against turning a count into a strength. So the crystallization weight is branch-blind, the informational masses keep the values the census derived, attenuated equally on all six states, and the defect values keep their own currency, where they split a doublet and set a bridge factor and change nothing in bits. Declaring a defect-dependent rate would be exactly the kind of parameter this part has refused eleven times over, and it is refused again: the calibration registry stays empty, and the two currencies stay as incommensurate as the item that names them.

\end{theorem}

\begin{theorem}[{\texorpdfstring{BM.91 — The corrected neutron–proton splitting: the attenuated assembly, the ensemble-independent bar, and the universal ratio bound}{BM.91 - The corrected neutron-proton splitting: the attenuated assembly, the ensemble-independent bar, and the universal ratio bound}}]\label{tag:BM-91}
\textbf{Thm.}

%dsp-level:0
\[
(1)\quad\textbf{the corrected decomposition:}\qquad\boxed{\ \Delta m\ :=\ m_n-m_p\ =\ \Delta m_{\mathrm{transport}}+\Delta m_{\mathrm{EM}}^{\mathrm{corr}}+\Delta m_{\mathrm{bridge}}\ }
\]

%dsp-level:0
\[
\Delta m_{\mathrm{transport}}=+\frac{\ln(4/3)}{\mathfrak h}
\]

%dsp-level:0
\[
\qquad [\mathrm{BM.59}(4);\ \mathrm{BM.84}(1)];
\]

%dsp-level:0
\[
\qquad \Delta m_{\mathrm{EM}}^{\mathrm{corr}}=-\,\frac{3\,N\,w\,\mathsf g\,\varepsilon_{\mathrm{EM}}\ln2}{\mathfrak h}
\]

%dsp-level:0
\[
\qquad [\mathrm{BM.89}(5);\ \text{the path-integral read}];
\]

%dsp-level:0
\[
\Delta m_{\mathrm{bridge}}=0\ \ \text{in the informational currency}
\]

%dsp-level:0
\[
\qquad [\mathrm{BM.84}(1):\ \neg\text{formed there};\ \mathrm{BM.69}(3);\ \mathrm{BM.90}(5):\ \text{the pole companion}\ \Delta_g>0\ \text{is}\ \neg\text{addable}];
\]

%dsp-level:0
\[
\textbf{the four zeros stand:}
\]

%dsp-level:0
\[
\qquad \Delta m_{\mathrm{Lyapunov}}
\]

%dsp-level:0
\[
= 0\ [\mathrm{BM.51}(6)],
\]

%dsp-level:0
\[
\quad\Delta m_{\mathrm{metric}}
\]

%dsp-level:0
\[
= 0\ [\mathrm{BM.53}(2);\ \mathrm{BM.58}(4)],
\]

%dsp-level:0
\[
\quad\Delta m_{\mathrm{binding}}
\]

%dsp-level:0
\[
= 0\ [\mathrm{BM.54}(1)\text{--}(3)],
\]

%dsp-level:0
\[
\quad\Delta m_{\mathrm{force}}
\]

%dsp-level:0
\[
= \Delta m_{\mathrm{self\text{-}int}}
\]

%dsp-level:0
\[
= -\Delta m_{\mathrm{EM}}^{\mathrm{corr}}\ [\mathrm{BM.52}(4),(5)];
\]

%dsp-level:0
\[
(2)\quad\textbf{the attenuated assembly:}\qquad\boxed{\ m_n\,\mathfrak h\ =\ 6\,w\ln2+\ln20,\qquad m_p\,\mathfrak h\ =\ 6\,w\ln2+\ln15+3\,N\,w\,\mathsf g\,\varepsilon_{\mathrm{EM}}\ln2\ }
\]

%dsp-level:0
\[
[\mathrm{BM.89}(4),(5);\ \mathrm{Def\text{-}BM.13}(5)\text{'s refined rest input with the ensemble-averaged edge register};\ \text{machine: the prime-log assembly}\ \ln20=2\ln2+\ln5,\ \ln15=\ln3+\ln5];
\]

%dsp-level:0
\[
\textbf{the register scoping:}
\]

%dsp-level:0
\[
\qquad w\ \text{attenuates the edge-count register}\ E_{\mathrm{int}}=|R|\ \text{alone}\ [\mathrm{Def\text{-}BM.5}(4)];
\]

%dsp-level:0
\[
\qquad \text{the transport register}\ \Sigma_{\mathrm{cl}}(Q)
\]

%dsp-level:0
\[
= \ln W(Q)\ \text{is an end-sign filter count,}\ \neg\text{an edge count}\ [\mathrm{Def\text{-}BM.13}(4),(6)],
\]

%dsp-level:0
\[
\text{and its ensemble average is}\ \neg\text{constructed}\ [\mathrm{FT.106}(3)\text{'s census is at fixed topology}]\implies\textbf{it enters at its crystallized value and the residue is}\ (\mathrm{R4})\ \ [\mathrm{BM.89}(8);\ \mathrm{BM.93}(6)];
\]

%dsp-level:0
\[
\boxed{\ \Delta m\,\mathfrak h\ =\ \ln\frac43-3\,N\,w\,\mathsf g\,\varepsilon_{\mathrm{EM}}\ln2\ }\qquad[\text{machine:}\ \ln20-\ln15=\ln\tfrac43];
\]

%dsp-level:0
\[
\textbf{the cancellation row:}
\]

%dsp-level:0
\[
\qquad E_{\mathrm{int}}\ \text{is class-shared}\ [\mathrm{BM.57}(4);\ \text{Census FT.35}(2)]\implies\textbf{the crystallization weight cancels in the ledger difference and survives only in the EM term}
\]

%dsp-level:0
\[
[\text{machine: the coefficient census of}\ w\ \text{in the two rest inputs and in}\ \Delta m];
\]

%dsp-level:0
\[
\qquad [\Delta m]=\mathfrak q^{+1/2}=[\text{energy}]
\]

%dsp-level:0
\[
\qquad [\mathrm{CT2.1};\ \text{machine}];
\]

%dsp-level:0
\[
(3)\quad\textbf{the exact ratio:}\qquad\boxed{\ \frac{m_n}{m_p}\ =\ \frac{6\,w\ln2+\ln20}{6\,w\ln2+\ln15+3\,N\,w\,\mathsf g\,\varepsilon_{\mathrm{EM}}\ln2}\ }\qquad[\text{machine}];
\]

%dsp-level:0
\[
\text{at}\ (w,N,\mathsf g)=(1,1,1):
\]

%dsp-level:0
\[
\qquad \frac{m_n}{m_p}=\frac{\ln1280}{\ln960+3\,\varepsilon_{\mathrm{EM}}\ln2}
\]

%dsp-level:0
\[
\qquad [\mathrm{BM.84}(2)\text{'s row at the derived}\ c_{\mathrm{formed}}=\tfrac12;\ \text{machine:}\ 6\ln2+\ln20=\ln1280,\ 6\ln2+\ln15=\ln960];
\]

%dsp-level:0
\[
(4)\quad\textbf{the exact critical strength:}
\]

%dsp-level:0
\[
\qquad \Delta m\ \text{is strictly decreasing in}\ \varepsilon_{\mathrm{EM}}\ \text{at every admissible}\ (N,w,\mathsf g)>0
\]

%dsp-level:0
\[
\qquad [\text{machine: the coefficient census}]:
\]

%dsp-level:0
\[
\boxed{\ \varepsilon^{\mathrm{crit}}(N,w,\mathsf g)\ :=\ \frac{\ln(4/3)}{3\,N\,w\,\mathsf g\,\ln2}\ =\ \frac{2-\log_23}{3\,N\,w\,\mathsf g}\ }
\]

%dsp-level:0
\[
\qquad \textbf{the unique zero of}\ \Delta m
\]

%dsp-level:0
\[
\qquad [\text{machine: solved exactly in the prime-log ring}];
\]

%dsp-level:0
\[
\textbf{the vacuum-gauge, fully crystallized value:}
\]

%dsp-level:0
\[
\qquad \varepsilon^{\mathrm{crit}}(1,1,1)=\frac{2-\log_23}{3}\in\Bigl(\frac18,\frac17\Bigr)
\]

%dsp-level:0
\[
\qquad [\text{machine:}\ 2^{13}=8192>6561=3^8;\quad2^{11}=2048<2187=3^7];
\]

%dsp-level:0
\[
\textbf{the exact squeeze of the numerator:}
\]

%dsp-level:0
\[
\qquad \frac25\ <\ 2-\log_23\ <\ \frac5{12}
\]

%dsp-level:0
\[
\qquad [\text{machine:}\ \log_23<\frac85\iff3^5=243<256=2^8;\quad\log_23>\frac{19}{12}\iff3^{12}=531441>524288=2^{19}]:
\]

%dsp-level:0
\[
\frac{2}{15\,Nw\mathsf g}\ <\ \varepsilon^{\mathrm{crit}}(N,w,\mathsf g)\ <\ \frac{5}{36\,Nw\mathsf g}
\]

%dsp-level:0
\[
\qquad [\text{machine: two rational bounds on the critical strength at every admissible}\ (N,w,\mathsf g)];
\]

%dsp-level:0
\[
\textbf{the irrationality row:}
\]

%dsp-level:0
\[
\qquad \varepsilon^{\mathrm{crit}}(N,w,\mathsf g)\notin\mathbb Q\ \ \text{for every rational}\ (N,w,\mathsf g)>0
\]

%dsp-level:0
\[
\qquad [\mathrm{BM.84}(4):\ 2-\log_23\notin\mathbb Q\ \text{on the}\ 2^p\ne3^q\ \text{census};\ \text{machine}];
\]

%dsp-level:0
\[
\textbf{the dissolution row:}
\]

%dsp-level:0
\[
\qquad \text{on the Bernoulli specialization}\ w=1-4\varepsilon_q\ [\mathrm{BM.89}(4)]:
\]

%dsp-level:0
\[
\qquad \varepsilon^{\mathrm{crit}}\to\infty\ \ \text{as}\ \ \varepsilon_q\to\tfrac14^{-}
\]

%dsp-level:0
\[
\boxed{\ \textbf{the fully dissolved limit never closes the splitting}\ }
\]

%dsp-level:0
\[
\qquad [\text{machine: monotone in}\ w,\ \text{unbounded as}\ w\to0^+];
\]

%dsp-level:0
\[
(5)\quad\textbf{the ensemble-independent bar (parameter-free):}\qquad\boxed{\ N\,w\,\mathsf g\cdot\varepsilon_{\mathrm{EM}}\ \le\ \frac18\qquad\Longrightarrow\qquad\Delta m>0\ }
\]

%dsp-level:0
\[
[\text{machine:}\ \Delta m>0\iff Nw\mathsf g\,\varepsilon_{\mathrm{EM}}<\frac{2-\log_23}3;\qquad\frac{2-\log_23}3>\frac18\iff8(2-\log_23)>3\iff2^{13}>3^8\iff8192>6561];
\]

%dsp-level:0
\[
\textbf{the bar is uniform over the lapse, the crystallization weight and the curvature composition}\ —\ \neg\text{one of the three is a datum the corpus fixes};
\]

%dsp-level:0
\[
\qquad \text{the two priced witnesses at}\ \varepsilon_{\mathrm{EM}}
\]

%dsp-level:0
\[
= \tfrac19\ \text{stand}\ [\mathrm{BM.62}(8):\ 4^3=64>54=3^3\!\cdot\!2;\quad4^5=1024>972=3^5\!\cdot\!2^2];
\]

%dsp-level:0
\[
(6)\quad\textbf{the trichotomy (total, exact, non-empty on both sides):}
\]

%dsp-level:0
\[
\boxed{\ \varepsilon_{\mathrm{EM}}<\varepsilon^{\mathrm{crit}}\iff m_n>m_p\iff\frac{m_n}{m_p}>1;\qquad\varepsilon_{\mathrm{EM}}=\varepsilon^{\mathrm{crit}}\iff m_n=m_p\iff\frac{m_n}{m_p}=1;\qquad\varepsilon_{\mathrm{EM}}>\varepsilon^{\mathrm{crit}}\iff m_n<m_p\iff\frac{m_n}{m_p}<1\ }
\]

%dsp-level:0
\[
\textbf{the opening region is non-empty:}
\]

%dsp-level:0
\[
\qquad \Bigl\{0\le Nw\mathsf g\,\varepsilon_{\mathrm{EM}}\le\tfrac18\Bigr\}\ne\varnothing\ \ \text{and contains}\ \varepsilon_{\mathrm{EM}}=0\ \ \text{and}\ \ \varepsilon_{\mathrm{EM}}=\tfrac19\ \ \text{at}\ Nw\mathsf g=1
\]

%dsp-level:0
\[
\qquad [(5);\ \text{machine}];
\]

%dsp-level:0
\[
\textbf{the closing region is non-empty:}
\]

%dsp-level:0
\[
\qquad \varepsilon_{\mathrm{EM}}>\varepsilon^{\mathrm{crit}}(N,w,\mathsf g)\ \ \text{is inhabited at every admissible}\ (N,w,\mathsf g)>0
\]

%dsp-level:0
\[
\qquad [\text{machine:}\ \varepsilon^{\mathrm{crit}}\ \text{is a positive exact value}];
\]

%dsp-level:0
\[
\textbf{the critical locus is exact, unique, and irrational at rational}\ (N,w,\mathsf g)\qquad[(4)];
\]

%dsp-level:0
\[
(7)\quad\textbf{the universal ratio bound:}
\]

%dsp-level:0
\[
\qquad \frac{m_n}{m_p}\ \text{is strictly decreasing in}\ \varepsilon_{\mathrm{EM}}\ \text{and strictly decreasing in}\ w
\]

%dsp-level:0
\[
\qquad [\text{machine: the two cross-multiplied sign rows}]:
\]

%dsp-level:0
\[
\textbf{on the sub-critical branch:}
\]

%dsp-level:0
\[
\qquad 1\ <\ \frac{m_n}{m_p}\ <\ \frac{\ln20}{\ln15}
\]

%dsp-level:0
\[
\qquad \text{over all}\ \ w\in(0,1],\ N,\mathsf g>0,\ 0\le\varepsilon_{\mathrm{EM}}<\varepsilon^{\mathrm{crit}}
\]

%dsp-level:0
\[
\qquad [\text{the lower end open at the critical locus};\ \text{the upper end a strict supremum as}\ w\to0^+,\ \varepsilon_{\mathrm{EM}}\to0];
\]

%dsp-level:0
\[
\textbf{at}\ \varepsilon_{\mathrm{EM}}=0\ \textbf{specifically:}
\]

%dsp-level:0
\[
\qquad \frac{\ln1280}{\ln960}\ \le\ \frac{m_n}{m_p}\bigg|_{\varepsilon_{\mathrm{EM}}=0}\ <\ \frac{\ln20}{\ln15}
\]

%dsp-level:0
\[
\qquad [\text{the lower endpoint attained at}\ w=1;\ \text{machine: the monotonicity in}\ w];
\]

%dsp-level:0
\[
\frac{\ln1280}{\ln960}\in\Bigl(\frac{26}{25},\frac{21}{20}\Bigr)
\]

%dsp-level:0
\[
\qquad [\mathrm{BM.84}(3);\ \text{machine:}\ 2^{44}>3^{26}\!\cdot\!5,\ 2^{34}<3^{21}\!\cdot\!5];
\]

%dsp-level:0
\[
\qquad \frac{\ln20}{\ln15}\in\Bigl(\frac{11}{10},\frac{10}9\Bigr)
\]

%dsp-level:0
\[
[\text{machine, derived here}:\quad10\ln20>11\ln15\iff20^{10}>15^{11}\ (10240000000000>8649755859375);\qquad9\ln20<10\ln15\iff20^9<15^{10}\ (512000000000<576650390625)];
\]

%dsp-level:0
\[
\boxed{\ 1\ <\ \frac{m_n}{m_p}\ <\ \frac{10}9\qquad\text{at every admissible}\ (w,\,N,\,\mathsf g,\,\varepsilon_{\mathrm{EM}})\ \text{on the sub-critical branch}\ }
\]

%dsp-level:0
\[
(8)\quad\textbf{the incommensurability, ensemble-robust:}
\]

%dsp-level:0
\[
\qquad \frac{m_n}{m_p}=\frac43\iff3\bigl(6w\ln2+\ln20\bigr)=4\bigl(6w\ln2+\ln15+3Nw\mathsf g\,\varepsilon_{\mathrm{EM}}\ln2\bigr)
\]

%dsp-level:0
\[
\iff12\,Nw\mathsf g\,\varepsilon_{\mathrm{EM}}\ln2=(6-6w)\ln2-\ln405\iff\boxed{\ \varepsilon_{\mathrm{EM}}^{*}(w,N,\mathsf g)=\frac{(6-6w)\ln2-\ln405}{12\,N\,w\,\mathsf g\,\ln2}\ }
\]

%dsp-level:0
\[
\text{and}\qquad(6-6w)\ln2\ \le\ 6\ln2\ =\ \ln64\ <\ \ln405\qquad[\text{machine:}\ 64<405;\qquad w>0]\implies\boxed{\ \varepsilon_{\mathrm{EM}}^{*}<0\ \ \text{for every}\ w\in(0,1]\ }
\]

%dsp-level:0
\[
\textbf{the informational ratio never equals the ensemble ratio at any admissible coupling and any admissible crystallization weight}
\]

%dsp-level:0
\[
\qquad [\mathrm{BM.83}(3);\ \mathrm{BM.84}(8)\ \text{generalized};\ \text{machine}];
\]

%dsp-level:0
\[
\textbf{the irrationality survives:}
\]

%dsp-level:0
\[
\qquad \frac{m_n}{m_p}\notin\mathbb Q\ \ \text{at every rational}\ (w,N,\mathsf g,\varepsilon_{\mathrm{EM}})
\]

%dsp-level:0
\[
\qquad [\mathrm{BM.83}(2)\text{'s lemma: the}\ \ln3\ \text{coefficient of the numerator is}\ 0\ \text{and of the denominator is}\ 1;\ \text{machine}];
\]

%dsp-level:0
\[
(9)\quad\textbf{the rescope rows (recorded,}\ \neg\textbf{retracted):}
\]

%dsp-level:0
\[
\qquad \mathrm{BM.62}(3)\text{--}(9),
\]

%dsp-level:0
\[
\ \mathrm{BM.63},
\]

%dsp-level:0
\[
\ \mathrm{BM.65}(3)\text{--}(5),
\]

%dsp-level:0
\[
\ \mathrm{BM.66}(3),
\]

%dsp-level:0
\[
\ \mathrm{BM.70},
\]

%dsp-level:0
\[
\ \mathrm{BM.75}(4)\text{'s opposed row},
\]

%dsp-level:0
\[
\ \mathrm{BM.77}(9),
\]

%dsp-level:0
\[
\ \mathrm{BM.84}(2)\text{--}(8),
\]

%dsp-level:0
\[
\ \mathrm{BM.85}(1)\text{--}(3):
\]

%dsp-level:0
\[
\textbf{stand on their own locus}\ —\ \text{the}\ (w,N,\mathsf g)=(1,1,1)\ \text{specialization with}\ c\ \text{symbolic or priced};
\]

%dsp-level:0
\[
\qquad \text{every integer certificate they carry is re-derived here at the specialization and none is weakened};
\]

%dsp-level:0
\[
\mathrm{BM.62}(7)\text{'s four-certificate enclosure keeps both rows:}\ \text{the}\ c
\]

%dsp-level:0
\[
= \tfrac35\ \text{row stands as the void alternative's exact arithmetic},
\]

%dsp-level:0
\[
\ \text{and the}\ c
\]

%dsp-level:0
\[
= \tfrac12\ \text{row becomes}\ (4)\text{'s derived value};
\]

%dsp-level:0
\[
\neg\text{a frozen refutation is reopened};
\]

%dsp-level:0
\[
\qquad \neg\text{a Census cell changes};
\]

%dsp-level:0
\[
\qquad \mathrm{BM.77}(5)\text{'s specified-read trichotomy is untouched}\ [\mathrm{BM.92}(3)];
\]

%dsp-level:0
\[
(10)\quad\textbf{status:}
\]

%dsp-level:0
\[
\qquad \textbf{derived:}\ (1)\text{--}(8);
\]

%dsp-level:0
\[
\qquad \textbf{recorded:}\ (9);
\]

%dsp-level:0
\[
\qquad \textbf{symbolic:}\ \varepsilon_{\mathrm{EM}},\ N,\ \mathsf g,\ w\ (w\ \text{evaluated exactly at}\ \varepsilon_q=0\ \text{and on the marked specialization});
\]

%dsp-level:0
\[
\neg\text{a P4 addition},
\]

%dsp-level:0
\[
\quad\neg\text{a new dimensionful scale}\ [\mathrm{P1}],
\]

%dsp-level:0
\[
\quad\neg\text{a measured value}\ [\mathrm{P5}],
\]

%dsp-level:0
\[
\quad\neg\text{a declaration},
\]

%dsp-level:0
\[
\quad\text{§Cal-BM stays empty};
\]

%dsp-level:0
\[
\textbf{the open row:}\ \text{the numerical value of}\ e\ [\mathrm{EM.25}(7)]\ —\ \text{it decides which side of}\ \varepsilon^{\mathrm{crit}}\ \text{the corpus sits on, and no theorem of Parts 1--4 supplies it}\ [\mathrm{BM.62}(12);\ \mathrm{BM.84}(9)\ \text{re-cited}].
\]

\textbf{Req.} CT2.1, D2.9, EM.25(7) (Part 2); P1, P4, P5 (Parts 1–2); FT.44(4) (Part 3); Census FT.35(2) (row-wise); Def-BM.13(5), BM.51(6), BM.52(4), BM.52(5), BM.53(2), BM.54(1), BM.54(3), BM.57(4), BM.58(4), BM.59(4), BM.62(3), BM.62(7), BM.62(8), BM.62(12), BM.63, BM.65(3), BM.65(5), BM.66(3), BM.69(3), BM.70, BM.75(4), BM.77(9), BM.83(2), BM.83(3), BM.84(1), BM.84(2), BM.84(3), BM.84(4), BM.84(8), BM.84(9), BM.85(1), BM.85(3), BM.89(4), BM.89(5), BM.90(5) (this part); exact prime-log arithmetic with integer power certificates, Laurent-monomial assembly with coefficient censuses, cross-multiplied monotonicity rows, exact rational enclosures of $\ln$-ratios [math; machine: \texttt{checks/\allowbreak{}verify\_\allowbreak{}H\_\allowbreak{}em\_\allowbreak{}self\_\allowbreak{}energy\_\allowbreak{}correction.py} §E].

\textbf{Pf.}

%dsp-level:0
\[
(1):\ \text{the five zero-or-monomial rows re-cited}\ [\mathrm{BM.62}(1),(2);\ \mathrm{BM.84}(1)];
\]

%dsp-level:0
\[
\qquad \Delta m_{\mathrm{EM}}^{\mathrm{corr}}=-\bigl(E_{\mathrm{EM}}^{\mathrm{self}}(+1)-E_{\mathrm{EM}}^{\mathrm{self}}(0)\bigr)=-\tfrac{3Nw\mathsf g\varepsilon_{\mathrm{EM}}\ln2}{\mathfrak h}\ [\mathrm{BM.89}(5);\ \text{machine}]
\]

%dsp-level:0
\[
(2):\ 6w\ln2+\ln20-(6w\ln2+\ln15+3Nw\mathsf g\varepsilon_{\mathrm{EM}}\ln2)
\]

%dsp-level:0
\[
= \ln20-\ln15-3Nw\mathsf g\varepsilon_{\mathrm{EM}}\ln2
\]

%dsp-level:0
\[
= \ln\tfrac43-3Nw\mathsf g\varepsilon_{\mathrm{EM}}\ln2\ [\text{machine: the prime-log ring}];
\]

%dsp-level:0
\[
\qquad \ln20=2\ln2+\ln5,\ \ln15=\ln3+\ln5\ [\text{machine:}\ 20=2^2\cdot5,\ 15=3\cdot5];
\]

%dsp-level:0
\[
\qquad w\text{'s coefficient in}\ \Delta m\ \text{is}\ 0\ \text{in the ledger part}\ [\text{machine}]
\]

%dsp-level:0
\[
(3):\ \text{the quotient of}\ (2)\text{'s two rows}\ [\text{machine}];
\]

%dsp-level:0
\[
\qquad (w,N,\mathsf g)=(1,1,1):\ 6\ln2+\ln20=8\ln2+\ln5=\ln(2^8\cdot5)=\ln1280,
\]

%dsp-level:0
\[
\quad 6\ln2+\ln15=6\ln2+\ln3+\ln5=\ln960\ [\text{machine}]
\]

%dsp-level:0
\[
(4):\ \Delta m
\]

%dsp-level:0
\[
= 0\iff\ln\tfrac43
\]

%dsp-level:0
\[
= 3Nw\mathsf g\,
\]

%dsp-level:0
\[
\varepsilon_{\mathrm{EM}}\ln2\iff\varepsilon_{\mathrm{EM}}
\]

%dsp-level:0
\[
= \tfrac{\ln(4/3)}{3Nw\mathsf g\ln2}
\]

%dsp-level:0
\[
= \tfrac{2\ln2-\ln3}{3Nw\mathsf g\ln2}
\]

%dsp-level:0
\[
= \tfrac{2-\log_23}{3Nw\mathsf g}\ [\text{machine: solved in the prime-log ring}];
\]

%dsp-level:0
\[
\qquad \tfrac{2-\log_23}3>\tfrac18\iff8(2-\log_23)>3\iff13>8\log_23\iff2^{13}>3^8\ [\text{machine:}\ 8192>6561]
\]

%dsp-level:0
\[
\tfrac{2-\log_23}3<\tfrac17\iff7(2-\log_23)<3\iff11<7\log_23\iff2^{11}<3^7\ [\text{machine:}\ 2048<2187];
\]

%dsp-level:0
\[
\qquad 2-\log_23\notin\mathbb Q\ [\mathrm{BM.84}(4)\text{'s census}];
\]

%dsp-level:0
\[
\qquad \varepsilon^{\mathrm{crit}}=O(1/w)\to\infty\ \text{as}\ w\to0^+\ [\text{machine: monotone, unbounded}]
\]

%dsp-level:0
\[
(5):\ \Delta m>0\iff3Nw\mathsf g\,\varepsilon_{\mathrm{EM}}\ln2<\ln\tfrac43\iff Nw\mathsf g\,\varepsilon_{\mathrm{EM}}<\tfrac{2-\log_23}3\ [\text{machine}];
\]

%dsp-level:0
\[
\qquad \tfrac18<\tfrac{2-\log_23}3\ [(4)]\implies Nw\mathsf g\,\varepsilon_{\mathrm{EM}}\le\tfrac18\implies\Delta m>0\ [\text{machine}]
\]

%dsp-level:0
\[
(6):\ (4)\text{'s uniqueness}\ \wedge\ \text{the strict monotonicity in}\ \varepsilon_{\mathrm{EM}}\implies\text{the three-way equivalence}\ [\text{machine: the sign census on a rational grid of}\ (N,w,\mathsf g,\varepsilon_{\mathrm{EM}})];
\]

%dsp-level:0
\[
\qquad \varepsilon_{\mathrm{EM}}=0\in\text{the opening region since}\ \ln\tfrac43>0\ [4>3;\ \text{machine}]
\]

%dsp-level:0
\[
(7):\ \text{monotone in}\ \varepsilon_{\mathrm{EM}}:\ \text{the denominator increases, the numerator is fixed}\ [\text{machine}];
\]

%dsp-level:0
\[
\qquad \text{monotone in}\ w\ \text{at}\ \varepsilon_{\mathrm{EM}}
\]

%dsp-level:0
\[
= 0:\ w_1<w_2\implies\frac{6w_1\ln2+\ln20}{6w_1\ln2+\ln15}>\frac{6w_2\ln2+\ln20}{6w_2\ln2+\ln15}
\]

%dsp-level:0
\[
\iff (w_2-w_1)6\ln2(\ln20-\ln15)>0\ [\text{machine: cross-multiplied,}\ \ln20>\ln15\ \text{on}\ 20>15]
\]

%dsp-level:0
\[
10\ln20-11\ln15=\ln\frac{20^{10}}{15^{11}}>0\iff20^{10}>15^{11}\ [\text{machine:}\ 10240000000000>8649755859375];
\]

%dsp-level:0
\[
\qquad 9\ln20-10\ln15=\ln\frac{20^9}{15^{10}}<0\iff20^9<15^{10}\ [\text{machine:}\ 512000000000<576650390625]
\]

%dsp-level:0
\[
(8):\ 3(6w\ln2+\ln20)-4(6w\ln2+\ln15+3Nw\mathsf g\,\varepsilon_{\mathrm{EM}}\ln2)
\]

%dsp-level:0
\[
= -6w\ln2-4\ln3-\ln5-12Nw\mathsf g\,
\]

%dsp-level:0
\[
\varepsilon_{\mathrm{EM}}\ln2\ [\text{machine: the prime-log ring}];
\]

%dsp-level:0
\[
\qquad \text{setting it to}\ 0\ \text{gives}\ \varepsilon^*_{\mathrm{EM}}\ \text{of}\ (8)
\]

%dsp-level:0
\[
(6-6w)\ln2
\]

%dsp-level:0
\[
\le 6\ln2
\]

%dsp-level:0
\[
= \ln64<\ln405
\]

%dsp-level:0
\[
\iff 64<405\ [\text{machine}]\implies\text{the numerator of}\ \varepsilon^*_{\mathrm{EM}}\ \text{is negative}\implies\varepsilon^*_{\mathrm{EM}}<0\ [\text{machine}];
\]

%dsp-level:0
\[
\qquad \text{irrationality}:\ q(6w\ln2+\ln20)
\]

%dsp-level:0
\[
= p(6w\ln2+\ln15+3Nw\mathsf g\,\varepsilon_{\mathrm{EM}}\ln2)\ \text{with rational}\ (w,N,\mathsf g,\varepsilon_{\mathrm{EM}})\implies\text{the}\ \ln3\ \text{coefficient}\ -p
\]

%dsp-level:0
\[
= 0
\]

%dsp-level:0
\[
\implies p
\]

%dsp-level:0
\[
= 0
\]

%dsp-level:0
\[
\implies q
\]

%dsp-level:0
\[
= 0\ [\mathrm{BM.83}(2);\ \text{machine}]
\]

%dsp-level:0
\[
(9)\text{--}(10):\ \text{the clause census and the standing-rule scans}\ [\text{machine}].\ \blacksquare
\]

\textbf{Phys.} Put the corrected field energy into the splitting and almost nothing moves — which is the point. Both classes lose the same fraction of their edge register to the flicker, so the crystallization weight cancels out of the ledger difference entirely and survives only where the field lives: the electromagnetic term is now the shell value, times the lapse, times the crystallization weight, times the curvature composition, times three times the logarithm of two over h. The critical strength inherits those three factors in the denominator, which means the opening region can only widen as the composite fails to crystallize, and in the fully dissolved limit it never closes at all. Better, the bar becomes parameter-free in the ensemble data: whenever the product of lapse, weight, curvature composition and coupling stays at or below one eighth, the neutron is the heavier baryon, and the certificate is a single pair of integers, eight thousand one hundred ninety-two against six thousand five hundred sixty-one. The ratio itself is bounded between one and ten ninths at every admissible value of everything the corpus declines to fix, with the old one-and-a-twentieth bound recovered as the fully crystallized endpoint and the logarithm of twenty over the logarithm of fifteen as the dissolved supremum. And the incommensurability hardens: the coupling at which the informational ratio would equal the ensemble's four thirds is negative for every crystallization weight in the admissible range, on the single integer certificate that sixty-four is less than four hundred five.

\end{theorem}

\begin{theorem}[{\texorpdfstring{BM.92 — The Q-value on both formed reads: the uniform closure, the corrected margins, and the untouched specified branch}{BM.92 - The Q-value on both formed reads: the uniform closure, the corrected margins, and the untouched specified branch}}]\label{tag:BM-92}
\textbf{Thm.}

%dsp-level:0
\[
(1)\quad\textbf{the informational Q-value with the corrected term:}
\]

%dsp-level:0
\[
\qquad Q_\beta:=m_n-m_p-m_\ell-m_\nu
\]

%dsp-level:0
\[
\qquad [\mathrm{BM.41}(4);\ \mathrm{BM.63}(1);\ \mathrm{BM.85}(1)]:
\]

%dsp-level:0
\[
\boxed{\ Q_\beta\,\mathfrak h\ =\ \ln\frac43-3\,N\,w\,\mathsf g\,\varepsilon_{\mathrm{EM}}\ln2-F_{\ell\nu},\qquad F_{\ell\nu}:=m_\ell\mathfrak h+m_\nu\mathfrak h\ }
\]

%dsp-level:0
\[
\qquad [\mathrm{BM.91}(2);\ \text{machine: the prime-log assembly}];
\]

%dsp-level:0
\[
F_{\ell\nu}\ge2\ln2
\]

%dsp-level:0
\[
\qquad [\mathrm{BM.70}(2),(3):\ \text{every materialized closed chain carries}\ \ge2\ \text{directed edges};\ m\mathfrak h=E_{\mathrm{int}}\ln2;\quad m_\nu\mathfrak h\ge0\ \text{by PSD}];
\]

%dsp-level:0
\[
\boxed{\ Q_\beta\,\mathfrak h\ \le\ \ln\frac43-2\ln2\ =\ -\ln3\ <\ 0\qquad\text{at every admissible}\ (w,\,N,\,\mathsf g,\,\varepsilon_{\mathrm{EM}})\ }
\]

%dsp-level:0
\[
\qquad [\text{machine:}\ 3Nw\mathsf g\,\varepsilon_{\mathrm{EM}}\ln2\ge0;\qquad4<12];
\]

%dsp-level:0
\[
\textbf{the closure is uniform over the ensemble data:}
\]

%dsp-level:0
\[
\qquad \text{the attenuation lowers the magnitude of the EM term and cannot raise}\ Q_\beta\ \text{above}\ -\ln3/\mathfrak h
\]

%dsp-level:0
\[
\qquad [\text{machine: the coefficient census}];
\]

%dsp-level:0
\[
(2)\quad\textbf{the margin ladder and the corrected ratio forms:}
\]

%dsp-level:0
\[
\qquad Q_\beta\mathfrak h\ \le\ -\ln3,\quad-\ln12,\quad-\ln24
\]

%dsp-level:0
\[
\qquad \text{on}\quad4<12,\quad4<48,\quad4<96
\]

%dsp-level:0
\[
\qquad [\mathrm{BM.70}(3);\ \mathrm{BM.85}(2)\ \text{re-certified};\ \text{machine}]:
\]

%dsp-level:0
\[
\Sigma_{\mathrm B}-2\ln2=-\ln3,
\]

%dsp-level:0
\[
\qquad \Sigma_{\mathrm B}-4\ln2=-\ln12,
\]

%dsp-level:0
\[
\qquad \Sigma_{\mathrm B}-5\ln2=-\ln24
\]

%dsp-level:0
\[
\qquad [\text{machine: the three exact identities}];
\]

%dsp-level:0
\[
\boxed{\ \frac{Q_\beta}{m_n}\ \le\ -\,\frac{\ln3}{6\,w\ln2+\ln20}\ }\qquad\text{with the two exact endpoint brackets:}
\]

%dsp-level:0
\[
w=1:
\]

%dsp-level:0
\[
\qquad -\frac{\ln3}{\ln1280}\in\Bigl(-\frac2{13},-\frac3{20}\Bigr)
\]

%dsp-level:0
\[
\qquad [\mathrm{BM.85}(3);\ \text{machine:}\ 3^{13}=1594323<1638400=2^{16}\!\cdot\!5^2;\quad3^{20}=3486784401>2097152000=2^{24}\!\cdot\!5^3];
\]

%dsp-level:0
\[
w\to0^+:
\]

%dsp-level:0
\[
\qquad -\frac{\ln3}{\ln20}\in\Bigl(-\frac25,-\frac13\Bigr)
\]

%dsp-level:0
\[
\qquad [\text{machine, derived here}:\quad3\ln3>\ln20\iff27>20;\qquad5\ln3<2\ln20\iff243<400];
\]

%dsp-level:0
\[
\textbf{the informational Q-value is a strictly negative exact fraction of the parent mass at every crystallization weight, bounded away from zero by a rational};
\]

%dsp-level:0
\[
(3)\quad\textbf{the specified read is untouched:}
\]

%dsp-level:0
\[
\qquad Q_\beta^{\mathrm{spec}}\,\mathfrak h=\Delta T(S)-\rho_\ell\sqrt{G\bigl(c(0)\bigr)}-\rho_\nu\sqrt{G\bigl(c(i_\nu)\bigr)},
\]

%dsp-level:0
\[
\qquad \Delta T(S)=5\,m^2_B(\cdot,S)
\]

%dsp-level:0
\[
\qquad [\mathrm{BM.77}(2);\ \mathrm{BM.85}(4)]:
\]

%dsp-level:0
\[
\textbf{the EM summands are}\ \neg\textbf{parts of that read}
\]

%dsp-level:0
\[
\qquad [\mathrm{BM.77}(3):\ \text{disjoint alphabets, no double count}]\implies\boxed{\ \text{neither the rescission of}\ \mathrm{BM.87}\ \text{nor the ensemble correction of}\ \mathrm{BM.89}\ \text{alters any row of}\ \mathrm{BM.77}(5),(6)\ }
\]

%dsp-level:0
\[
\Delta T=\Bigl(\frac{4219785}{98956},\ \frac{9608310}{393679},\ \frac{455}{23},\ \frac{68565}{3088}\Bigr),
\]

%dsp-level:0
\[
\qquad \frac{273}{23}<\rho^{\mathrm{crit}}(S,i_\nu)<\frac{42197850}{791648}
\]

%dsp-level:0
\[
\qquad [\mathrm{BM.77}(2),(5);\ \text{machine re-certified}];
\]

%dsp-level:0
\[
\textbf{the witnesses stand:}
\]

%dsp-level:0
\[
\qquad 828100>5819,
\]

%dsp-level:0
\[
\qquad 4219785<43\cdot98956=4255108,
\]

%dsp-level:0
\[
\qquad 7500>1849,
\]

%dsp-level:0
\[
\qquad 3312400>5819,
\]

%dsp-level:0
\[
\qquad 99<100,
\]

%dsp-level:0
\[
\qquad 75>64,
\]

%dsp-level:0
\[
\qquad 273>253,
\]

%dsp-level:0
\[
\qquad 42197850<42748992;
\]

%dsp-level:0
\[
(4)\quad\textbf{the verdict rows (read-relative and complete):}
\]

%dsp-level:0
\[
\qquad \textbf{closed on the informational currency for every member, uniformly in}\ (w,N,\mathsf g,\varepsilon_{\mathrm{EM}})
\]

%dsp-level:0
\[
\qquad [(1);\ \mathrm{BM.70}(4)(\mathrm a)];
\]

%dsp-level:0
\[
\textbf{open on}\ \mathcal O_S,\ \textbf{closed on}\ \mathcal C_S,\ \textbf{degenerate on}\ \partial\mathcal O_S\ \ \text{on the specified read}
\]

%dsp-level:0
\[
\qquad [\mathrm{BM.77}(5),(6);\ \mathrm{BM.85}(5)];
\]

%dsp-level:0
\[
\textbf{the decay direction is fixed in both currencies:}
\]

%dsp-level:0
\[
\qquad T(0,S)=\tfrac43T(+1,S)>T(+1,S)
\]

%dsp-level:0
\[
\qquad [\mathrm{BM.28}(4)(\mathrm{ii});\ \text{machine}]
\]

%dsp-level:0
\[
\qquad \text{and}
\]

%dsp-level:0
\[
\qquad m_n>m_p\ \ \text{on the sub-critical branch}
\]

%dsp-level:0
\[
\qquad [\mathrm{BM.91}(6)]:
\]

%dsp-level:0
\[
\boxed{\ \textbf{the parent is the neutral class on every formed read, and the grammar}\ n\to p+\ell+\nu\ \textbf{stands unmodified}\ }
\]

%dsp-level:0
\[
\qquad [\mathrm{BM.41}(1);\ \mathrm{BM.45};\ \mathrm{BM.77}(6)];
\]

%dsp-level:0
\[
\textbf{the conditional spectrum on}\ \mathcal O_S\textbf{:}
\]

%dsp-level:0
\[
\qquad E_\ell^{\max}=\frac{m_n^2+m_\ell^2-(m_p+m_\nu)^2}{2m_n},
\]

%dsp-level:0
\[
\qquad E_{\max}:\langle E\rangle=2:1\ \ \text{scale-invariant}
\]

%dsp-level:0
\[
\qquad [\mathrm{BM.41}(5);\ \mathrm{BM.49}(3),(4);\ \mathrm{BM.77}(8);\ \neg\text{an external endpoint compared}\ [\mathrm{P5}]];
\]

%dsp-level:0
\[
(5)\quad\textbf{the two reads are}\ \neg\textbf{comparable:}
\]

%dsp-level:0
\[
\qquad Q_\beta^{\mathrm{info}}:\ \{\ln2,\ln3,\ln5,\ N,\ w,\ \mathsf g,\ \varepsilon_{\mathrm{EM}},\ \mathfrak h,\ m_\ell,\ m_\nu\};
\]

%dsp-level:0
\[
\qquad Q_\beta^{\mathrm{spec}}:\ \{m^2_B(\cdot,S),\ \rho_\ell,\ \rho_\nu,\ g,\ g',\ \mathfrak h\}
\]

%dsp-level:0
\[
\boxed{\ \text{the two alphabets are disjoint except for}\ \mathfrak h\ }
\]

%dsp-level:0
\[
\qquad [\text{machine: the token census};\qquad\mathsf g\ \text{and}\ g\ \text{are distinct symbols}\ [(4)\text{'s local notation};\ \mathrm{BM.39}(1)]];
\]

%dsp-level:0
\[
\neg\text{a derived equation relates the two expressions}\ —\ \text{that is}\ (\mathrm{R1}),
\]

%dsp-level:0
\[
\ (\mathrm{R3}),
\]

%dsp-level:0
\[
\ (\mathrm{R4})\ \text{and the Cal-F scales together}
\]

%dsp-level:0
\[
\qquad [\mathrm{BM.82}(6);\ \mathrm{BM.83}(5),(7);\ \mathrm{BM.85}(5)];
\]

%dsp-level:0
\[
(6)\quad\textbf{status:}
\]

%dsp-level:0
\[
\qquad \textbf{derived:}\ (1)\text{--}(5);
\]

%dsp-level:0
\[
\qquad \textbf{read-relative:}\ \text{the kinematic verdict}\ (4);
\]

%dsp-level:0
\[
\qquad \textbf{declared: nothing}\ —\ \text{A-BM.8's specification is re-cited,}\ \neg\text{re-issued};
\]

%dsp-level:0
\[
\neg\text{a clause of}\ \mathrm{BM.41},\ \mathrm{BM.49},\ \mathrm{BM.63},\ \mathrm{BM.70},\ \mathrm{BM.71},\ \mathrm{BM.77},\ \mathrm{BM.85}\ \text{is retracted};
\]

%dsp-level:0
\[
\qquad \neg\text{a P4 addition},\quad\neg\text{a new dimensionful scale}\ [\mathrm{P1}],\quad\neg\text{a measured value}\ [\mathrm{P5}],\quad\text{§Cal-BM stays empty}.
\]

\textbf{Req.} P1, P4, P5 (Parts 1–2); FT.79(5) (Part 3); Def-BM.13(4), BM.28(4), BM.37(2), BM.39(1), BM.41(1), BM.41(4), BM.41(5), BM.45, BM.49(3), BM.49(4), BM.63(1), BM.70(2), BM.70(3), BM.70(4), BM.77(2), BM.77(3), BM.77(5), BM.77(6), BM.77(8), BM.82(6), BM.83(5), BM.83(7), BM.84(1), BM.85(1), BM.85(2), BM.85(3), BM.85(4), BM.85(5), BM.87, BM.89, BM.91(2), BM.91(6) (this part); the declaration addendum A-BM.8, item 7 (pointers); exact prime-log arithmetic with integer power certificates, exact rational arithmetic at the four frozen sector expectations, token censuses [math; machine: \texttt{checks/\allowbreak{}verify\_\allowbreak{}H\_\allowbreak{}em\_\allowbreak{}self\_\allowbreak{}energy\_\allowbreak{}correction.py} §F].

\textbf{Pf.}

%dsp-level:0
\[
(1):\ Q_\beta\mathfrak h=\Delta m\,\mathfrak h-F_{\ell\nu}=\ln\tfrac43-3Nw\mathsf g\,\varepsilon_{\mathrm{EM}}\ln2-F_{\ell\nu}\ [\mathrm{BM.91}(2);\ \text{machine}];
\]

%dsp-level:0
\[
\qquad F_{\ell\nu}\ge2\ln2\ [\mathrm{BM.70}(2),(3)\ \text{re-cited}];
\]

%dsp-level:0
\[
\qquad \ln\tfrac43-2\ln2=2\ln2-\ln3-2\ln2=-\ln3\ [\text{machine}];
\]

%dsp-level:0
\[
\qquad -\ln3<0\iff3>1\iff4<12\ [\text{machine}]
\]

%dsp-level:0
\[
3Nw\mathsf g\,
\]

%dsp-level:0
\[
\varepsilon_{\mathrm{EM}}\ln2
\]

%dsp-level:0
\[
\ge 0\ \text{at}\ N,
\]

%dsp-level:0
\[
w,
\]

%dsp-level:0
\[
\mathsf g,
\]

%dsp-level:0
\[
\varepsilon_{\mathrm{EM}}
\]

%dsp-level:0
\[
\ge 0\ [\text{machine: the coefficient census}]
\]

%dsp-level:0
\[
\implies Q_\beta\mathfrak h
\]

%dsp-level:0
\[
\le -\ln3\ \text{uniformly}\ [\text{machine}]
\]

%dsp-level:0
\[
(2):\ \Sigma_{\mathrm B}-2\ln2
\]

%dsp-level:0
\[
= -\ln3,
\]

%dsp-level:0
\[
\quad\Sigma_{\mathrm B}-4\ln2
\]

%dsp-level:0
\[
= 2\ln2-\ln3-4\ln2
\]

%dsp-level:0
\[
= -\ln(4\cdot3)
\]

%dsp-level:0
\[
= -\ln12,
\]

%dsp-level:0
\[
\quad\Sigma_{\mathrm B}-5\ln2
\]

%dsp-level:0
\[
= -\ln(8\cdot3)
\]

%dsp-level:0
\[
= -\ln24\ [\text{machine: the three identities}];
\]

%dsp-level:0
\[
\qquad 4<12<48<96\ [\text{machine}]
\]

%dsp-level:0
\[
Q_\beta/m_n=(Q_\beta\mathfrak h)/(m_n\mathfrak h)\le-\ln3/(6w\ln2+\ln20)\ [\mathrm{BM.91}(2);\ \text{machine}];
\]

%dsp-level:0
\[
\qquad 3\ln3-\ln20>0\iff27>20,\quad5\ln3-2\ln20<0\iff243<400\ [\text{machine}]
\]

%dsp-level:0
\[
(3):\ \mathrm{BM.77}(2),(3),(5)\ \text{re-cited};
\]

%dsp-level:0
\[
\qquad 15\cdot843957=12659355,
\]

%dsp-level:0
\[
\quad 15\cdot1921662=28824930,
\]

%dsp-level:0
\[
\quad 15\cdot91=1365,
\]

%dsp-level:0
\[
\quad 15\cdot13713=205695\ [\mathrm{BM.89}(6);\ \text{machine}];
\]

%dsp-level:0
\[
\qquad \text{the eight witnesses re-certified as integer comparisons}\ [\text{machine}]
\]

%dsp-level:0
\[
(4):\ T(0,S)/T(+1,S)=W(0)/W(+1)=20/15=4/3>1\ [\text{machine}];
\]

%dsp-level:0
\[
\qquad m_n>m_p\ \text{sub-critically}\ [\mathrm{BM.91}(6)];
\]

%dsp-level:0
\[
\qquad \text{the endpoint identity and the}\ 2{:}1\ \text{moment ratio re-cited}\ [\mathrm{BM.41}(5);\ \mathrm{BM.49}(3),(4)]
\]

%dsp-level:0
\[
(5):\ \text{the two symbol sets enumerated from the two expressions and intersected}\ [\text{machine}];
\]

%dsp-level:0
\[
\qquad \mathsf g\ne g\ \text{as tokens}\ [\text{machine}];
\]

%dsp-level:0
\[
\qquad (6):\ \text{the standing-rule scans}\ [\text{machine}].\ \blacksquare
\]

\textbf{Phys.} One Q-value, two formed reads, and the correction lands on only one of them. On the informational currency the decay's energy balance keeps its shape: the transport register's log of four thirds, less the corrected field term, less whatever the two emitted lines cost. Because the corrected field term is nonnegative and now attenuated by the lapse, the crystallization weight and the curvature composition, the two-bit floor on any materialized chain still drives the total below minus the logarithm of three — at every value of every quantity the corpus declines to fix. The closure is not merely preserved, it is uniform, and its size as a fraction of the parent mass runs from minus three twentieths at full crystallization to minus one third in the dissolved limit, always an exact rational bound away from zero. On the specified read nothing moves at all: the electromagnetic summands were never part of that expression, so the opening region, the closing region, the exact critical locus and all eight integer witnesses stand as derived. The two verdicts remain read-relative and incomparable, their alphabets sharing one symbol, and the one thing both agree on is the direction — the neutral class is the parent, the grammar is untouched, and the decay is allowed exactly where a formed read allows it.

\end{theorem}

\medskip\noindent\textit{\texorpdfstring{BM.93 — The status of the EM self-energy correction: adjudication and scope}{BM.93 - The status of the EM self-energy correction: adjudication and scope}.}\ Structural-constraints record: see Appendix B, Section B.1.28 of the Addendum.\par\medskip

\phantomsection
\section[Tier BM.19 — The Flicker Integral, the Rest-Mass Q-Value,…]{\texorpdfstring{Tier BM.19 — The Flicker Integral, the Rest-Mass Q-Value, and the Relating Equation: the Materialized Ensemble, the Vertex Superposition, the Rung-One $\nu$ Line, and the Bridge in Four Exact Forms}{Tier BM.19 - The Flicker Integral, the Rest-Mass Q-Value, and the Relating Equation: the Materialized Ensemble, the Vertex Superposition, the Rung-One Line, and the Bridge in Four Exact Forms}}

\textbf{Locus.} This tier fixes one read, moves one Q-value onto one species, and closes one ledger item in form. The read it fixes is the crystallized-snapshot read of a sealed composite's information content — the read that takes the computed cell's $E_{\mathrm{int}}=6$ $[\mathrm{Census\ FT.35}(2);\ \mathrm{BM.17}(1)]$ as the composite's count and hence as a fixed mass. That read is a category error against the corpus's own rows: the count is one point of an ensemble distribution $[\mathrm{Def\text{-}BM.5}(4);\ \mathrm{BM.19}(5)]$, the real and the virtual registers are mixed at every nonzero threshold $[\mathrm{A4};\ \mathrm{D0.1}]$, and the composite is an ensemble and $\neg$ a topology $[\mathrm{BM.19}(6);\ \mathrm{BM.89}(4)]$. $\mathrm{BM.94}$ replaces the snapshot with the integral, derives the distribution and its exact moments, and derives the materialization restriction that holds the mean edge count above the digon's two. $\mathrm{BM.95}$ derives the superposition a vertex carries in the transport register, with its exact weights and the two positive parts of the transport splitting. $\mathrm{BM.96}$ derives the rung-one $\nu$ line's rest mass from the seal structure, its strict positivity, its three-line order, its single-ray occupancy, and the four independent bars against a seesaw-type mechanism. $\mathrm{BM.97}$ adjudicates the rest-mass species and forms the Q-value on it, with the exact critical locus and integer-certified witnesses on both the specified and the pure-pole reads. $\mathrm{BM.98}$ derives the relating equation in four exact forms — count, curvature, coupling, flicker-weighted — with its enhancement bounded by an exact integer. $\mathrm{BM.99}$ recomputes the splitting and the Q-value on the materialized read, gaining a two-sided parameter-free bar and two exact tightenings. $\mathrm{BM.100}$ adjudicates.

\textbf{Standing rows.} $\alpha_B$ stays barred on $\mathrm{BM.61}(8)$'s four grounds and §Cal-BM stays empty. No clause of $\mathrm{BM.1}$–$\mathrm{BM.93}$ is retracted, replaced or renumbered: every row this tier touches is re-scoped at its own locus and recorded in $\mathrm{BM.94}(9)$, $\mathrm{BM.97}(7)$, $\mathrm{BM.99}(7)$ and $\mathrm{BM.100}(4)$, in $\mathrm{BM.77}(7)$'s pattern. No measured coupling, mass or ratio enters any derivation $[\mathrm{P5}]$: the neutron–proton ratio is derived and bounded at $\mathrm{BM.99}(5)$ and the electric coupling stays symbolic $[\mathrm{EM.25}(7);\ \mathrm{BM.88}(7)]$ — neither is imported. No right-ray matter line and no seesaw-type mechanism is suggested, added or required $[\mathrm{BM.96}(4),(6)]$. The relating equation is derived in form with its residue localized to one undeclared dimensionless ratio; $\rho_{\mathrm B}$ is $\neg$derived and $\mathrm{BM.69}(4)$'s four grounds stand unweakened. Local notation, fixed for the whole tier: $p:=1-4\varepsilon_q$ the stationary retention weight $[\mathrm{BM.19}(2)]$; $u:=1-p=4\varepsilon_q$ the flicker weight; $D:=3-3p+p^2=1+u+u^2$ the dilution denominator; $w$ the crystallization weight of $\mathrm{BM.89}(4)$ and $w_{\mathrm{mat}}:=1/D$ its materialized companion; $E^{\mathrm{cryst}}:=6$; $k:=r/\Delta x_{\mathrm B}$, $f(k)=1-\tfrac1k$, $N$, $\mathsf g$, $G$, $\varepsilon_{\mathrm{EM}}$ as in $\mathrm{BM.88}$'s local notation; $\rho_{\mathrm B}$, $D_S$, $\mathcal R$ as in $\mathrm{BM.69}(2)$; $M_0$, $s$, $\Delta_g$, $\Delta_g^{\mathrm{crit}}$ as in $\mathrm{BM.97}(4)$; $c:=\cos\tfrac18$, $s_\nu:=\sin\tfrac18$ in $\mathrm{BM.96}(2)$. The symbol $q$ keeps its end-register meaning $[\mathrm{D3.4};\ \mathrm{EM.9}]$, the italic $g$ keeps its $\mathrm{BM.39}(1)$ meaning, and the italic $s$ keeps its $\mathrm{BM.97}(4)$ meaning at every occurrence.

\begin{theorem}[{\texorpdfstring{BM.94 — The flicker integral: the snapshot read's category error, the materialized distribution, and the two crystallization weights}{BM.94 - The flicker integral: the snapshot read's category error, the materialized distribution, and the two crystallization weights}}]\label{tag:BM-94}
\textbf{Thm.}

%dsp-level:0
\[
(1)\quad\textbf{the category error (the referent census):}\qquad\begin{array}{l|l}\textbf{the snapshot read takes} & \textbf{the corpus supplies}\\\hline E_{\mathrm{int}}=6\ \text{as the composite's count} & E_{\mathrm{int}}(R)=|R|\ \text{distributed over the flicker ensemble}\ [\mathrm{Def\text{-}BM.5}(4);\ \mathrm{BM.19}(5)]\\\text{one configuration as the composite} & |\mathcal E(\mathcal X)|>1\iff\varepsilon_q>0:\ \text{an ensemble at every nonzero threshold}\ [\mathrm{BM.19}(6)]\\ m=\frac{6\ln2}{\mathfrak h}\ \text{as a fixed mass} & \text{a mixture of the real and the virtual registers}\ [\mathrm{A4};\ \mathrm{D0.1};\ \mathrm{BM.89}(4)]\\\text{a value} & \text{a distribution whose top point is the crystallized value}\ [\mathrm{Def\text{-}BM.5}(4);\ \mathrm{BM.86}(4)\ \text{scope}]\end{array}
\]

%dsp-level:0
\[
\textbf{the error, stated exactly:}
\]

%dsp-level:0
\[
\qquad \boxed{\ E_{\mathrm{int}}=6\ \ \text{is one point of the distribution,}\ \neg\text{the ensemble, and its weight is}\ P(E{=}6)=\frac{p^2}{D}<1\ \ \text{at every}\ \varepsilon_q>0\ }
\]

%dsp-level:0
\[
\textbf{the weight deficit is exact:}\qquad1-P_{\mathrm{mat}}(E{=}6)=\frac{3(1-p)}{D}>0\qquad[\text{machine:}\ D-p^2=3-3p=3(1-p)];
\]

%dsp-level:0
\[
(2)\quad\textbf{the proper integral:}
\]

%dsp-level:0
\[
\qquad \boxed{\ E_{\mathrm{int}}^{\mathrm{eff}}\ :=\ \sum_{R\,\in\,\mathcal E(\mathcal X)}P(R)\cdot E_{\mathrm{int}}(R)\ }
\]

%dsp-level:0
\[
\qquad [\mathrm{BM.87}(5);\ \mathrm{BM.20}(3);\ \mathrm{Def\text{-}H.32}(1);\ \mathrm{Def\text{-}BM.5}(2)]:
\]

%dsp-level:0
\[
P(R)\ \text{is the ensemble weight of the configuration}\ R,
\]

%dsp-level:0
\[
\ \text{derived as the unique stationary line of the per-pair Janus generator}\ [\mathrm{BM.19}(1),(2);\ \mathrm{FT.108}(4);\ \mathrm{Def\text{-}BM.5}(3)],
\]

%dsp-level:0
\[
\ \neg\text{declared};
\]

%dsp-level:0
\[
\text{the sum is finite and exact at every budget}\ B\ [\mathrm{Def\text{-}BM.5}(1);\ \mathrm{Def\text{-}H.32}(1)],
\]

%dsp-level:0
\[
\ \text{and the}\ B\to\infty\ \text{limit is}\ \neg\text{constructed}\ [\mathrm{Def\text{-}BM.5}(5)];
\]

%dsp-level:0
\[
\textbf{linearity on}\ |R|\ \text{gives the per-datum form:}
\]

%dsp-level:0
\[
\qquad E_{\mathrm{int}}^{\mathrm{eff}}=\sum_{\text{pairs}}P(\text{the pair is real})\cdot2
\]

%dsp-level:0
\[
\qquad [\mathrm{BM.89}(4)\ \text{re-derived as the first moment};\ \text{machine}];
\]

%dsp-level:0
\[
(3)\quad\textbf{the ledger read (the quantum numbers the read is taken on):}
\]

%dsp-level:0
\[
\qquad \mathfrak L=\bigcup_{n\ge0}\mathfrak L(\tau_n),
\]

%dsp-level:0
\[
\qquad \bigl|\mathfrak L(\tau_n)\bigr|=\sum_{k<n}n_{\mathrm{real\_out}}(\tau_k)
\]

%dsp-level:0
\[
\qquad [\mathrm{A7};\ \mathrm{A6}]:
\]

%dsp-level:0
\[
\textbf{the two quantum numbers are the tick and the scale:}
\]

%dsp-level:0
\[
\qquad \tau:\ \tau_n\to\tau_{n+1},\ \neg\exists\,\tau^{-1}\ [\mathrm{A6};\ \mathrm{D1.8}:\ x^0:=\tau_n],
\]

%dsp-level:0
\[
\qquad \mathfrak h:\ \exists!\ \text{free dimensionful scale}\ [\mathrm{P1};\ \mathrm{D2.9};\ \text{A §0.3}];
\]

%dsp-level:0
\[
n_{\mathcal R}=E_{\mathrm{int}}=|R|\ [\mathrm{D0.8};\ \mathrm{Def\text{-}BM.5}(4)],\qquad S_{\mathfrak L}=n_{\mathcal R}\ln2\ [\mathrm{C9};\ \mathrm{CT1.29};\ \mathrm{BM.35}(3)];
\]

%dsp-level:0
\[
\textbf{the crystallized card is the snapshot at one}\ \tau\ [\mathrm{Census\ FT.35}(2);\ \mathrm{BM.17}(1)]:
\]

%dsp-level:0
\[
\qquad \text{the composite is read from}\ \mathfrak L,\ \neg\text{from the card};
\]

%dsp-level:0
\[
\textbf{the constituent at a vertex is a quantum superposition over the ledger's transport fibre,}\ \neg\text{a fixed state}
\]

%dsp-level:0
\[
\qquad [\mathrm{BM.95}(2);\ \mathrm{Def\text{-}BM.13}(1),(3)]:
\]

%dsp-level:0
\[
(4)\quad\textbf{the derived distribution (the marked pair specialization, marked}\ \neg\textbf{adopted):}
\]

%dsp-level:0
\[
\qquad E^{\mathrm{cryst}}
\]

%dsp-level:0
\[
= 6
\]

%dsp-level:0
\[
= (c_1{:}\,2)+(c_2{:}\,2)+(\text{connector}{:}\,2)\ \ \text{is three pairs}\ [\mathrm{Census\ FT.35}(2);\ \mathrm{BM.21}(1);\ \mathrm{BM.89}(4)]:
\]

%dsp-level:0
\[
\text{each pair is retained with the stationary weight}\ p
\]

%dsp-level:0
\[
= 1-4\varepsilon_q\ [\mathrm{BM.19}(2);\ \text{an identity, exact at every rational threshold}]
\]

%dsp-level:0
\[
\implies E_{\mathrm{int}}(R)\qquad = 2j,
\]

%dsp-level:0
\[
\quad j\in\{0,1,2,3\};
\]

%dsp-level:0
\[
\boxed{\ P(E{=}2j)=\binom3j p^{j}(1-p)^{3-j}\ }:\qquad\bigl((1{-}p)^3,\ \ 3p(1{-}p)^2,\ \ 3p^2(1{-}p),\ \ p^3\bigr)\ \ \text{at}\ \ E=(0,2,4,6)
\]

%dsp-level:0
\[
\text{normalization:}\ \sum_{j}\binom3jp^j(1-p)^{3-j}=(p+(1-p))^3=1
\]

%dsp-level:0
\[
\qquad [\text{machine: exact in}\ \mathbb Q[p];\ \text{the census at six rational thresholds}];
\]

%dsp-level:0
\[
\textbf{the unconditioned moments:}
\]

%dsp-level:0
\[
\qquad \langle E\rangle=6p=6w
\]

%dsp-level:0
\[
\qquad [\mathrm{BM.89}(4)\ \text{re-derived};\ \text{machine}],
\]

%dsp-level:0
\[
\qquad \langle E^2\rangle=12p+24p^2,
\]

%dsp-level:0
\[
\qquad \operatorname{Var}(E)=12p(1-p)
\]

%dsp-level:0
\[
\textbf{the no-flicker collapse:}
\]

%dsp-level:0
\[
\qquad \varepsilon_q=0\implies p=1\implies P(E{=}6)=1\ \ \text{and the ensemble is the single crystallized configuration}
\]

%dsp-level:0
\[
\qquad [\mathrm{BM.19}(3);\ \mathrm{Def\text{-}BM.5}(5);\ \text{machine}];
\]

%dsp-level:0
\[
(5)\quad\textbf{the materialization restriction (derived,}\ \neg\textbf{declared):}
\]

%dsp-level:0
\[
\qquad \textbf{the}\ E{=}0\ \textbf{member supports no closed chain and no terminus:}
\]

%dsp-level:0
\[
|E_{\mathcal R}|=0\implies\#\{\text{closed chains}\}=0\implies\neg\exists\ \text{a fermion terminus}
\]

%dsp-level:0
\[
\qquad [\mathrm{FT.144}(4);\ \mathrm{D1.4};\ \mathrm{Def\text{-}H.4};\ \text{machine}];
\]

%dsp-level:0
\[
\textbf{every materialized closed chain carries}\ \ge2\ \textbf{directed edges:}
\]

%dsp-level:0
\[
\qquad \delta h(c)
\]

%dsp-level:0
\[
\ge 2,
\]

%dsp-level:0
\[
\ n(v)
\]

%dsp-level:0
\[
\ge 2\ [\mathrm{Def\text{-}H.4};\ \mathrm{Census\ FT.28}:\ \text{the digon is the minimal floor};\ \mathrm{FT.32};\ \mathrm{BM.13}(3);\ \mathrm{BM.70}(2)];
\]

%dsp-level:0
\[
\textbf{the unconditioned average runs to the vacuum:}
\]

%dsp-level:0
\[
\qquad P(E{=}0)=(1-p)^3\ \ \text{is strictly increasing in}\ \varepsilon_q\ \ \text{and}\ \to1\ \ \text{as}\ \varepsilon_q\to\tfrac14^{-}
\]

%dsp-level:0
\[
\qquad [\text{machine}]:
\]

%dsp-level:0
\[
\text{the all-virtual state is the unique entropy maximum}\ \rho_\circ
\]

%dsp-level:0
\[
= \tfrac14I,
\]

%dsp-level:0
\[
\ S_{\mathrm{raw}}
\]

%dsp-level:0
\[
= \ln4\ [\mathrm{A5};\ \mathrm{CT0.2};\ \mathrm{BM.35}(1)]\ \ \text{and is the fixed point of the Lyapunov contraction}\ V'
\]

%dsp-level:0
\[
= (1-\tfrac\gamma4)^2V\ [\mathrm{L1};\ \mathrm{L2};\ \mathrm{BM.51}(1)]:
\]

%dsp-level:0
\[
\boxed{\ \textbf{the composite's read is the conditional one on}\ E\ge2,\ \ \text{and the restriction is a theorem of the materialization domain,}\ \neg\textbf{a scope choice}\ }
\]

%dsp-level:0
\[
\qquad [\mathrm{D2.11};\ \text{A §0.2a}];
\]

%dsp-level:0
\[
(6)\quad\textbf{the materialized distribution and its exact moments:}
\]

%dsp-level:0
\[
\qquad D\ :=\ 3-3p+p^2\ =\ \frac{1-(1-p)^3}{p}\ =\ 1+u+u^2,
\]

%dsp-level:0
\[
\qquad u:=1-p=4\varepsilon_q
\]

%dsp-level:0
\[
\qquad [\text{machine: one polynomial identity}]:
\]

%dsp-level:0
\[
\boxed{\ P_{\mathrm{mat}}(E{=}2)=\frac{3(1-p)^2}{D},\qquad P_{\mathrm{mat}}(E{=}4)=\frac{3p(1-p)}{D},\qquad P_{\mathrm{mat}}(E{=}6)=\frac{p^2}{D}\ }
\]

%dsp-level:0
\[
\qquad \text{and}\quad\frac{3(1-p)^2+3p(1-p)+p^2}{D}=1\ [\text{machine}];
\]

%dsp-level:0
\[
\boxed{\ \langle E\rangle_{\mathrm{mat}}=\frac6D,\qquad\operatorname{Var}_{\mathrm{mat}}(E)=\frac{12\,p(1-p)(3-2p)}{D^2},\qquad\frac{\operatorname{Var}_{\mathrm{mat}}(E)}{\langle E\rangle_{\mathrm{mat}}^2}=\frac{p(1-p)(3-2p)}{3}\ }
\]

%dsp-level:0
\[
\qquad [\text{machine: exact rational functions}];
\]

%dsp-level:0
\[
\textbf{the exact maximum of the relative width:}
\]

%dsp-level:0
\[
\qquad \frac{p(1-p)(3-2p)}3\ \le\ \frac{10+7\sqrt7}{162}\ \ \text{at}\ \ p=\frac{5-\sqrt7}{6}
\]

%dsp-level:0
\[
\qquad [\text{machine: the stationary point of}\ 6p^2-10p+3\ \text{in}\ \mathbb Q(\sqrt7)];
\]

%dsp-level:0
\[
\frac{10+7\sqrt7}{162}<\frac15
\]

%dsp-level:0
\[
\qquad [\text{machine:}\ 5(10+7\sqrt7)<162\iff35\sqrt7<112\iff5\sqrt7<16\iff175<256]
\]

%dsp-level:0
\[
\qquad \implies\boxed{\ \frac{\sigma_E}{\langle E\rangle_{\mathrm{mat}}}<\frac12\ }
\]

%dsp-level:0
\[
(7)\quad\textbf{the two crystallization weights:}
\]

%dsp-level:0
\[
\qquad w=\frac{\langle E\rangle}{E^{\mathrm{cryst}}}=p
\]

%dsp-level:0
\[
\qquad [\mathrm{BM.89}(4)\ \text{re-cited at its own locus}];
\]

%dsp-level:0
\[
\qquad \boxed{\ w_{\mathrm{mat}}\ :=\ \frac{\langle E\rangle_{\mathrm{mat}}}{E^{\mathrm{cryst}}}\ =\ \frac1D\ =\ \frac1{1+u+u^2}\ }
\]

%dsp-level:0
\[
w_{\mathrm{mat}}(0)=1\ \ \text{exactly}
\]

%dsp-level:0
\[
\qquad [\mathrm{BM.19}(3);\ \text{machine}];
\]

%dsp-level:0
\[
\qquad \boxed{\ \frac13<w_{\mathrm{mat}}\le1\ \ \text{on}\ \ \varepsilon_q\in[0,\tfrac14)\ }
\]

%dsp-level:0
\[
\qquad [\mathrm{CT0.8};\ \mathrm{D0.4};\ \text{machine:}\ D\ge1\iff(p-1)(p-2)\ge0;\quad D<3\iff p(p-3)<0];
\]

%dsp-level:0
\[
w_{\mathrm{mat}}\ \ge\ w\qquad\text{with equality}\iff\varepsilon_q=0\qquad[\text{machine:}\ w_{\mathrm{mat}}/w=1/(1-(1-p)^3)\ge1];
\]

%dsp-level:0
\[
\boxed{\ \langle E\rangle_{\mathrm{mat}}=6\,w_{\mathrm{mat}}\ \in\ (2,\,6]\ }
\]

%dsp-level:0
\[
\qquad \textbf{the top point is the crystallized count, the bottom is the digon's two}\ [\mathrm{Census\ FT.28};\ \mathrm{Def\text{-}H.4}]:
\]

%dsp-level:0
\[
\text{the lower endpoint is open and is attained only in the fully dissolved limit, where}\ \langle E\rangle_{\mathrm{mat}}\to2\ \ \text{exactly}
\]

%dsp-level:0
\[
\qquad [\text{machine:}\ 6/(3-3p+p^2)\to6/3];
\]

%dsp-level:0
\[
(8)\quad\textbf{the informational mass is a distribution:}
\]

%dsp-level:0
\[
\qquad m_{\mathrm{hid}}(R)=\frac{E_{\mathrm{int}}(R)\ln2}{\mathfrak h}
\]

%dsp-level:0
\[
\qquad [\mathrm{FT.44}(4);\ \mathrm{Census\ FT.44}(4);\ \mathrm{Def\text{-}H.9}(2)]:
\]

%dsp-level:0
\[
\textbf{the support:}
\]

%dsp-level:0
\[
\qquad m_{\mathrm{hid}}\in\Bigl\{\frac{2\ln2}{\mathfrak h},\ \frac{4\ln2}{\mathfrak h},\ \frac{6\ln2}{\mathfrak h}\Bigr\}\ \ \text{with}\ (6)\text{'s exact probabilities};
\]

%dsp-level:0
\[
\qquad \boxed{\ \langle m_{\mathrm{hid}}\rangle_{\mathrm{mat}}=\frac{6\,w_{\mathrm{mat}}\ln2}{\mathfrak h}\ },
\]

%dsp-level:0
\[
\qquad \operatorname{Var}_{\mathrm{mat}}(m_{\mathrm{hid}})=\frac{12\,p(1-p)(3-2p)\,(\ln2)^2}{D^2\,\mathfrak h^2};
\]

%dsp-level:0
\[
\textbf{the two-bit floor survives the averaging:}
\]

%dsp-level:0
\[
\qquad E\ge2\ \ \text{on the whole materialized support}\implies m_{\mathrm{hid}}(R)\ge\frac{2\ln2}{\mathfrak h}>0\ \ \forall R,
\]

%dsp-level:0
\[
\qquad \langle m_{\mathrm{hid}}\rangle_{\mathrm{mat}}>\frac{2\ln2}{\mathfrak h}
\]

%dsp-level:0
\[
\qquad [\mathrm{BM.70}(2),(3);\ \text{machine:}\ \tfrac12<\ln2<1];
\]

%dsp-level:0
\[
\textbf{the reciprocal register has two formed reads:}
\]

%dsp-level:0
\[
\qquad \bigl\langle E^{-1}\bigr\rangle_{\mathrm{mat}}=\frac{18p(1-p)^2+9p^2(1-p)+2p^3}{12\,p\,D},
\]

%dsp-level:0
\[
\qquad \frac1{\langle E\rangle_{\mathrm{mat}}}=\frac D6
\]

%dsp-level:0
\[
\qquad [\text{machine}]:
\]

%dsp-level:0
\[
\boxed{\ \langle E\rangle_{\mathrm{mat}}\bigl\langle E^{-1}\bigr\rangle_{\mathrm{mat}}\ =\ \frac{18-27p+11p^2}{2D^2}\ \ \ge\ 1\ }
\]

%dsp-level:0
\[
\qquad [\text{Cauchy--Schwarz on}\ (E,\,E^{-1});\ \text{machine: the exact rational function}]:
\]

%dsp-level:0
\[
\text{equality}\iff E\ \text{is constant on the support}\iff p=1\iff\varepsilon_q=0
\]

%dsp-level:0
\[
\qquad [\text{machine}];
\]

%dsp-level:0
\[
\qquad \frac{58}{49}\ \ \text{at}\ \ p=\tfrac12
\]

%dsp-level:0
\[
\qquad [\text{machine:}\ 290<294];
\]

%dsp-level:0
\[
\qquad 1\le\frac{18-27p+11p^2}{2D^2}<\frac65
\]

%dsp-level:0
\[
\qquad [\text{machine: the exact positivity of}\ 12p^4-72p^3+125p^2-81p+18\ \text{on}\ (0,1]];
\]

%dsp-level:0
\[
\implies\Delta x_{\mathrm B}\ \textbf{has two formed ensemble reads:}
\]

%dsp-level:0
\[
\qquad \Delta x_{\mathrm B}^{\mathrm{eff}}:=\frac{\mathfrak h}{\langle E\rangle_{\mathrm{mat}}\ln2}=D\cdot\Delta x_{\mathrm B}^{\mathrm{cryst}}
\]

%dsp-level:0
\[
\qquad \text{and}
\]

%dsp-level:0
\[
\qquad \frac{\mathfrak h}{\ln2}\bigl\langle E^{-1}\bigr\rangle_{\mathrm{mat}}\ \ \ge\ \ \Delta x_{\mathrm B}^{\mathrm{eff}}
\]

%dsp-level:0
\[
\qquad [\mathrm{Def\text{-}BM.4}(1);\ \mathrm{BM.13}(1);\ \text{machine}];
\]

%dsp-level:0
\[
\textbf{the tier evaluates the coupling and the bridge at the width of the mean; the difference of the two reads is named}\ (\mathrm{R6}),
\]

%dsp-level:0
\[
\ \neg\textbf{filled}
\]

%dsp-level:0
\[
\qquad [\mathrm{BM.98}(5);\ \text{the Jensen residue}];
\]

%dsp-level:0
\[
(9)\quad\textbf{status:}
\]

%dsp-level:0
\[
\qquad \textbf{derived:}\ (1)\text{--}(8)\ \text{on the marked pair specialization, which is}\ \neg\text{adopted as a general measure};
\]

%dsp-level:0
\[
\qquad \textbf{rescoped,}\ \neg\textbf{retracted:}\ \mathrm{Def\text{-}BM.5}(4)\text{'s ``one point'' row is discharged as an exact weight}\ (1);
\]

%dsp-level:0
\[
\mathrm{BM.89}(4)\text{'s}\ w\ \text{stands at its own locus as the unconditioned weight and}\ w_{\mathrm{mat}}\ \text{is its materialized companion},
\]

%dsp-level:0
\[
\ \neg\text{a replacement};
\]

%dsp-level:0
\[
\qquad \mathrm{BM.19}(1)\text{--}(6),\ \mathrm{BM.20}(3),\ \mathrm{BM.21},\ \mathrm{BM.86}(4)\ \text{untouched};
\]

%dsp-level:0
\[
\mathrm{BM.91}(2)\text{'s assembly stands on}\ w\ \text{and gains the materialized variant at}\ \mathrm{BM.99}(1);
\]

%dsp-level:0
\[
\qquad \neg\text{a declaration},
\]

%dsp-level:0
\[
\quad\neg\text{a P4 addition},
\]

%dsp-level:0
\[
\quad\neg\text{a new dimensionful scale}\ [\mathrm{P1}],
\]

%dsp-level:0
\[
\quad\neg\text{a measured value}\ [\mathrm{P5}],
\]

%dsp-level:0
\[
\quad\text{§Cal-BM stays empty};
\]

\textbf{Req.} A4, A5, A6, A7 (Part 1); C9, D0.1, D0.4, D0.8, D1.4, D1.8, D2.9, D2.11 (Parts 1–2); A §0.2a, A §0.3, CT0.2, CT0.8, CT1.29 (Part 1); CT2.1 (Part 2); P1, P4, P5 (Parts 1–2); L1, L2 (Part 1); Def-H.4, Def-H.9(2), Def-H.32(1), FT.28, FT.32, FT.44(4), FT.108(4), FT.144(4) (Part 3); Census FT.28, FT.35(2), FT.44(4) (row-wise); Def-BM.4(1), Def-BM.5(1), Def-BM.5(2), Def-BM.5(3), Def-BM.5(4), Def-BM.5(5), Def-BM.13(1), Def-BM.13(3), BM.13(1), BM.13(3), BM.17(1), BM.19(1), BM.19(2), BM.19(3), BM.19(5), BM.19(6), BM.20(3), BM.21(1), BM.35(1), BM.35(3), BM.51(1), BM.61(8), BM.70(2), BM.70(3), BM.86(4), BM.87(5), BM.88(7), BM.89(4), BM.91(2) (this part); exact polynomial arithmetic over $\mathbb Q[p]$, exact rational specialization at six thresholds, exact binomial censuses, the stationary point of a quadratic in $\mathbb Q(\sqrt7)$, exact rational interval arithmetic on a quartic, Cauchy–Schwarz over a positive measure [math; machine: \texttt{checks/\allowbreak{}verify\_\allowbreak{}H\_\allowbreak{}flicker\_\allowbreak{}integral\_\allowbreak{}restmass.py} §A].

\textbf{Pf.}

%dsp-level:0
\[
(1):\ \text{the two inventories enumerated from their defining clauses}\ [\mathrm{Census\ FT.35}(2);\ \mathrm{Def\text{-}BM.5}(4);\ \mathrm{BM.19}(5),(6)];
\]

%dsp-level:0
\[
\qquad 1-\tfrac{p^2}D=\tfrac{D-p^2}D=\tfrac{3-3p}D=\tfrac{3(1-p)}D\ [\text{machine: one polynomial identity}]
\]

%dsp-level:0
\[
(2):\ \mathrm{BM.87}(5)\text{'s relocation and}\ \mathrm{BM.20}(3)\text{'s averaged flux re-cited};
\]

%dsp-level:0
\[
\qquad \text{the measure's uniqueness:}\ \dim\ker Q=1\ [\mathrm{BM.19}(1);\ \mathrm{FT.108}(4)];
\]

%dsp-level:0
\[
\qquad \text{finiteness at every budget}\ [\mathrm{Def\text{-}H.32}(1)]
\]

%dsp-level:0
\[
(3):\ \mathrm{A6},\ \mathrm{A7},\ \mathrm{D1.8},\ \mathrm{P1}\ \text{re-cited verbatim (frozen)};
\]

%dsp-level:0
\[
\qquad n_{\mathcal R}=|R|\ [\mathrm{D0.8};\ \mathrm{Def\text{-}BM.5}(4)];
\]

%dsp-level:0
\[
\qquad S_{\mathfrak L}=n_{\mathcal R}\ln2\ [\mathrm{CT1.29};\ \mathrm{BM.35}(3)]
\]

%dsp-level:0
\[
(4):\ \text{three pairs, each retained with weight}\ p\ [\mathrm{BM.89}(4)\ \text{re-cited}];
\]

%dsp-level:0
\[
\qquad \textstyle\sum_j\binom3jp^j(1-p)^{3-j}
\]

%dsp-level:0
\[
= 1\ \text{by the binomial theorem over}\ \mathbb Q[p]\ [\text{machine: coefficient comparison}];
\]

%dsp-level:0
\[
\qquad \langle E\rangle=\textstyle\sum_j2j\binom3jp^j(1-p)^{3-j}=6p\ [\text{machine at six rational thresholds}];
\]

%dsp-level:0
\[
\langle E^2\rangle=\textstyle\sum_j4j^2\binom3jp^j(1-p)^{3-j}=12p+24p^2\ [\text{machine}];
\]

%dsp-level:0
\[
\qquad \operatorname{Var}=12p+24p^2-36p^2=12p(1-p)\ [\text{machine}];
\]

%dsp-level:0
\[
\qquad p=1\implies P(E{=}6)=1\ [\text{machine}]
\]

%dsp-level:0
\[
(5):\ |E_{\mathcal R}|=0\implies\neg\exists\ \text{a closed chain}\implies\neg\exists\ \text{a terminus}\ [\mathrm{FT.144}(4);\ \mathrm{D1.4};\ \mathrm{Def\text{-}H.4}\ \text{re-cited}];
\]

%dsp-level:0
\[
\qquad \delta h\ge2\ [\mathrm{Def\text{-}H.4};\ \mathrm{BM.13}(3);\ \mathrm{BM.70}(2)];
\]

%dsp-level:0
\[
\qquad (1-p)^3\ \text{strictly increasing in}\ \varepsilon_q\ \text{and}\ \to1\ [\text{machine: monotone on the rational grid}]
\]

%dsp-level:0
\[
\rho_\circ=\tfrac14I,\ S_{\mathrm{raw}}(\rho_\circ)=\ln4=\max,\ \arg\max=\{\tfrac14I\}\ [\mathrm{CT0.2};\ \mathrm{BM.35}(1)\ \text{re-cited}];
\]

%dsp-level:0
\[
\qquad V'=(1-\tfrac\gamma4)^2V\ \text{with}\ \gamma\to0\ \text{the identity-tilt corner}\ [\mathrm{L2}\ \text{re-cited}]
\]

%dsp-level:0
\[
(6):\ 1-(1-p)^3
\]

%dsp-level:0
\[
= 3p-3p^2+p^3
\]

%dsp-level:0
\[
= p(3-3p+p^2)
\]

%dsp-level:0
\[
= pD\ [\text{machine}]
\]

%dsp-level:0
\[
\implies P_{\mathrm{mat}}(E{=}2j)
\]

%dsp-level:0
\[
= \binom3jp^j(1-p)^{3-j}/(pD)\ \text{and the three values reduce to the boxed row}\ [\text{machine}];
\]

%dsp-level:0
\[
\qquad \text{the normalization:}\ 3(1-p)^2+3p(1-p)+p^2=3-3p+p^2=D\ [\text{machine}]
\]

%dsp-level:0
\[
\langle E\rangle_{\mathrm{mat}}=\tfrac{2\cdot3(1-p)^2+4\cdot3p(1-p)+6p^2}D=\tfrac{6(1-p)^2+12p(1-p)+6p^2}D=\tfrac6D\ [\text{machine}];
\]

%dsp-level:0
\[
\qquad \langle E^2\rangle_{\mathrm{mat}}=\tfrac{12(1+2p)}{D}\ [\text{machine}]
\]

%dsp-level:0
\[
\operatorname{Var}_{\mathrm{mat}}
\]

%dsp-level:0
\[
= \tfrac{12(1+2p)}D-\tfrac{36}{D^2}
\]

%dsp-level:0
\[
= \tfrac{12(1+2p)D-36}{D^2}
\]

%dsp-level:0
\[
= \tfrac{12p(1-p)(3-2p)}{D^2}\ [\text{machine:}\ 12(1+2p)D-36=36p-60p^2+24p^3=12p(3-2p)(1-p)]
\]

%dsp-level:0
\[
\tfrac{\operatorname{Var}}{\langle E\rangle^2}=\tfrac{12p(1-p)(3-2p)}{D^2}\cdot\tfrac{D^2}{36}=\tfrac{p(1-p)(3-2p)}3\ [\text{machine}];
\]

%dsp-level:0
\[
\qquad \tfrac{d}{dp}\bigl(3p-5p^2+2p^3\bigr)
\]

%dsp-level:0
\[
= 3-10p+6p^2
\]

%dsp-level:0
\[
= 0
\]

%dsp-level:0
\[
\iff p
\]

%dsp-level:0
\[
= \tfrac{5\pm\sqrt7}6,
\]

%dsp-level:0
\[
\ \ \text{one root in}\ (0,1]\ [\text{machine: the discriminant}\ 28=4\cdot7]
\]

%dsp-level:0
\[
p_\star=\tfrac{5-\sqrt7}6:
\]

%dsp-level:0
\[
\quad 1-p_\star=\tfrac{1+\sqrt7}6,
\]

%dsp-level:0
\[
\quad 3-2p_\star
\]

%dsp-level:0
\[
= \tfrac{4+\sqrt7}3\implies\tfrac{p_\star(1-p_\star)(3-2p_\star)}3
\]

%dsp-level:0
\[
= \tfrac{(5-\sqrt7)(1+\sqrt7)(4+\sqrt7)}{324}
\]

%dsp-level:0
\[
= \tfrac{20+14\sqrt7}{324}
\]

%dsp-level:0
\[
= \tfrac{10+7\sqrt7}{162}\ [\text{machine: exact in}\ \mathbb Q(\sqrt7)]
\]

%dsp-level:0
\[
5(10+7\sqrt7)<162
\]

%dsp-level:0
\[
\iff 35\sqrt7<112
\]

%dsp-level:0
\[
\iff 5\sqrt7<16
\]

%dsp-level:0
\[
\iff 175<256\ [\text{machine: squaring two positive integers}];
\]

%dsp-level:0
\[
\qquad \tfrac15<\tfrac14\implies\sigma/\langle E\rangle<\tfrac12\ [\text{machine}]
\]

%dsp-level:0
\[
(7):\ D(p{=}1)=1\implies w_{\mathrm{mat}}(0)=1\ [\text{machine}];
\]

%dsp-level:0
\[
\qquad D'=2p-3<0\ \text{on}\ (0,1]\implies D\in[1,3)\ \text{with}\ D\to3\ \text{as}\ p\to0^+\ [\text{machine}]\implies w_{\mathrm{mat}}=1/D\in(\tfrac13,1];
\]

%dsp-level:0
\[
\qquad D\ge1\iff p^2-3p+2\ge0\iff(p-1)(p-2)\ge0\ [\text{machine}]
\]

%dsp-level:0
\[
w_{\mathrm{mat}}/w=1/(pD)=1/(1-(1-p)^3)\ge1\iff(1-p)^3\ge0,\ \text{equality}\iff p=1\ [\text{machine}];
\]

%dsp-level:0
\[
\qquad \langle E\rangle_{\mathrm{mat}}=6/D\in(2,6]\ [\text{machine}];
\]

%dsp-level:0
\[
\qquad \lim_{p\to0^+}6/D=6/3=2\ [\text{machine: exact limit row}]
\]

%dsp-level:0
\[
(8):\ m_{\mathrm{hid}}
\]

%dsp-level:0
\[
= E\ln2/\mathfrak h\ \text{linear in}\ E\ [\mathrm{FT.44}(4)]\implies\text{the support and the two moments follow from}\ (6)\ [\text{machine}];
\]

%dsp-level:0
\[
\qquad E\ge2\ \text{on the materialized support}\ [(5)]\implies m_{\mathrm{hid}}\ge2\ln2/\mathfrak h>0\ [\tfrac12<\ln2<1;\ \text{machine}]
\]

%dsp-level:0
\[
\bigl\langle E^{-1}\bigr\rangle_{\mathrm{mat}}
\]

%dsp-level:0
\[
= \tfrac{\tfrac32p(1-p)^2+\tfrac34p^2(1-p)+\tfrac16p^3}{pD}
\]

%dsp-level:0
\[
= \tfrac{18p(1-p)^2+9p^2(1-p)+2p^3}{12pD}\ [\text{machine}];
\]

%dsp-level:0
\[
\qquad \langle E\rangle\langle E^{-1}\rangle=\tfrac6D\cdot\tfrac{18-27p+11p^2}{12pD}\cdot p=\tfrac{18-27p+11p^2}{2D^2}\ [\text{machine}]
\]

%dsp-level:0
\[
\langle E\rangle\langle E^{-1}\rangle\ge\langle\sqrt{E\cdot E^{-1}}\rangle^2=1\ [\text{Cauchy--Schwarz}];
\]

%dsp-level:0
\[
\qquad \text{equality}
\]

%dsp-level:0
\[
\iff E\cdot E^{-1}\ \text{proportional to a constant}
\]

%dsp-level:0
\[
\iff E\ \text{constant}
\]

%dsp-level:0
\[
\iff p
\]

%dsp-level:0
\[
= 1\ [\text{machine: the support has three points at every}\ p\in(0,1)];
\]

%dsp-level:0
\[
\qquad p=\tfrac12:\ \tfrac{18-\tfrac{27}2+\tfrac{11}4}{2\cdot\tfrac{49}{16}}=\tfrac{58}{49}\ [\text{machine}]
\]

%dsp-level:0
\[
\tfrac{18-27p+11p^2}{2D^2}<\tfrac65
\]

%dsp-level:0
\[
\iff 5(18-27p+11p^2)<12D^2
\]

%dsp-level:0
\[
\iff 12p^4-72p^3+125p^2-81p+18>0\ [\text{machine: exact rational interval arithmetic on}\ 64\ \text{subintervals of}\ [0,1]\ \text{with the derivative bound, lower endpoint}\ >0];
\]

%dsp-level:0
\[
\qquad \Delta x_{\mathrm B}^{\mathrm{eff}}=\mathfrak h/(6\ln2/D)=D\,\mathfrak h/(6\ln2)=D\,\Delta x_{\mathrm B}^{\mathrm{cryst}}\ [\mathrm{BM.13}(1);\ \text{machine}];
\]

%dsp-level:0
\[
\qquad (9):\ \text{the scope record}.\ \blacksquare
\]

\textbf{Phys.} The composite was never a snapshot, and the count that was read as one is a single point of an exact distribution. Take the flicker's own measure — the unique stationary line of the Janus generator, three pairs each retained with weight one minus four thresholds — and the internal count takes the values zero, two, four and six with the binomial weights, its unconditioned mean six times that weight and its variance twelve times the weight times its complement. The zero is not a composite: an all-virtual configuration supports no closed chain and therefore no terminus, and it is exactly the vacuum the entropy argument and the Lyapunov contraction both run toward, so the composite's read is the conditional one on at least two edges, and that restriction is a theorem of the materialization domain rather than a choice of scope. Conditional on it, the mean count is six over one plus u plus u squared, which never leaves the interval from two to six: the crystallized count at the top, attained only at the sharp threshold, and the digon's two at the bottom, attained only in the fully dissolved limit. The relative width of the distribution has an exact algebraic maximum, ten plus seven root seven over one hundred sixty-two, reached at a stationary point in the same quadratic field, and it is bounded by a fifth — so the informational mass fluctuates, but never by more than half of itself. The reciprocal register splits into two formed reads whose product is at least one by Cauchy–Schwarz and below six fifths by an exact quartic positivity, so the Born width has a width-of-the-mean and a mean-of-the-widths, and the difference is named rather than guessed. Read A took the top point for the whole distribution; the weight it ignored is three u over D, exactly, and it is positive at every threshold that lets the vertex fire at all.

\end{theorem}

\begin{theorem}[{\texorpdfstring{BM.95 — The vertex superposition: the transport fibre, its exact weights, and the two positive parts of the splitting}{BM.95 - The vertex superposition: the transport fibre, its exact weights, and the two positive parts of the splitting}}]\label{tag:BM-95}
\textbf{Thm.}

%dsp-level:0
\[
(1)\quad\textbf{the fibre and its census (re-derived):}
\]

%dsp-level:0
\[
\qquad q_{\mathrm{core}}(c)=\tfrac12\bigl(\sigma_{s(c)}-\sigma_{t(c)}\bigr)\in\{-1,0,+1\}
\]

%dsp-level:0
\[
\qquad [\mathrm{D3.4};\ \mathrm{EM.9};\ \mathrm{FT.40}]:
\]

%dsp-level:0
\[
m_q=(2,\,1,\,1)\ \ \text{at}\ \ q=(0,+1,-1)
\]

%dsp-level:0
\[
\qquad [\mathrm{FT.106}(3\mathrm a);\ \mathrm{Def\text{-}BM.13}(1):\ \text{the orbit size of the constituent's transport register under}\ \omega];
\]

%dsp-level:0
\[
W(Q)=\bigl[z^{Q}\bigr]\bigl(2+z+z^{-1}\bigr)^{k},
\]

%dsp-level:0
\[
\qquad k=3:\quad W=(1,\,6,\,15,\,\mathbf{20},\,15,\,6,\,1),
\]

%dsp-level:0
\[
\qquad \sum_QW(Q)=4^3=64=2^6
\]

%dsp-level:0
\[
\qquad [\mathrm{Def\text{-}BM.13}(3);\ \mathrm{Def\text{-}BM.8}(6);\ \text{machine}];
\]

%dsp-level:0
\[
\#\{(q_1,q_2,q_3):\ \textstyle\sum_iq_i=k\}=(1,3,6,\,\mathbf 7,\,6,3,1)
\]

%dsp-level:0
\[
\qquad [\mathrm{BM.28}(2);\ \text{machine: the exhaustive}\ 4^3\ \text{end-sign census and the convolution agree exactly}];
\]

%dsp-level:0
\[
(2)\quad\textbf{the derived superposition weights:}
\]

%dsp-level:0
\[
\qquad \textbf{the constituent at a vertex is a quantum superposition over its transport fibre, with weights}
\]

%dsp-level:0
\[
\boxed{\ P\bigl((q_1,q_2,q_3)\,\big|\,Q\bigr)\ =\ \frac{\prod_{i=1}^{3}m_{q_i}}{W(Q)}\ }
\]

%dsp-level:0
\[
\qquad [\text{the normalized fibre weight};\ \text{machine: exact rationals, summing to}\ 1\ \text{class by class}]:
\]

%dsp-level:0
\[
Q=0:
\]

%dsp-level:0
\[
\qquad P\bigl((0,0,0)\bigr)=\frac8{20}=\frac25,
\]

%dsp-level:0
\[
\qquad P\bigl(\text{each of the six permutations of}\ (+1,-1,0)\bigr)=\frac2{20}=\frac1{10}
\]

%dsp-level:0
\[
\qquad [\text{machine}];
\]

%dsp-level:0
\[
Q=+1:
\]

%dsp-level:0
\[
\qquad P\bigl(\text{each of the three permutations of}\ (+1,0,0)\bigr)=\frac4{15},
\]

%dsp-level:0
\[
\qquad P\bigl(\text{each of the three permutations of}\ (+1,+1,-1)\bigr)=\frac1{15}
\]

%dsp-level:0
\[
\qquad [\text{machine}];
\]

%dsp-level:0
\[
Q=-1:\qquad\text{the identical census}\ —\ \text{the}\ \omega\text{-image fibre}\qquad[\mathrm{D2.6};\ \mathrm{BM.25}(4);\ \mathrm{BM.28}(5);\ \text{machine}];
\]

%dsp-level:0
\[
\textbf{the weights carry no datum:}
\]

%dsp-level:0
\[
\qquad \text{the free-symbol set of}\ W(Q)\ \text{is}\ \{z\}\ \text{with coefficients fixed by}\ (m_q,k),
\]

%dsp-level:0
\[
\ \text{and the coupling alphabet}\ \{g,g',e,\varepsilon_{\mathrm{EM}},\hat\mu,\hat\mu',\ell_c\}\ \text{is disjoint from it}
\]

%dsp-level:0
\[
[\mathrm{BM.61}(5)\ \text{re-cited};\ \text{machine: the symbol census}]\implies\textbf{the superposition weights are}\ \neg\textbf{couplings,}\ \neg\textbf{rates,}\ \neg\textbf{P4 data};
\]

%dsp-level:0
\[
(3)\quad\textbf{the entropy decomposition:}
\]

%dsp-level:0
\[
\qquad H(Q)\ :=\ -\sum_{\text{fibre}}P\ln P,
\]

%dsp-level:0
\[
\qquad \bigl\langle\ln\textstyle\prod_im_{q_i}\bigr\rangle_Q\ :=\ \sum_{\text{fibre}}P\ln\textstyle\prod_im_{q_i}:
\]

%dsp-level:0
\[
\boxed{\ \Sigma_{\mathrm{cl}}(Q)\ =\ \ln W(Q)\ =\ H(Q)\ +\ \bigl\langle\ln\textstyle\prod_im_{q_i}\bigr\rangle_Q\ }
\]

%dsp-level:0
\[
\qquad [\mathrm{Def\text{-}BM.13}(4);\ \text{machine: exact in the prime-log ring}]:
\]

%dsp-level:0
\[
H(0)=\ln5+\frac15\ln2,\qquad\bigl\langle\ln\textstyle\prod_im\bigr\rangle_{0}=\frac95\ln2,\qquad\ln20=2\ln2+\ln5\qquad[\text{machine}];
\]

%dsp-level:0
\[
H(\pm1)=\ln3+\ln5-\frac85\ln2,\qquad\bigl\langle\ln\textstyle\prod_im\bigr\rangle_{+1}=\frac85\ln2,\qquad\ln15=\ln3+\ln5\qquad[\text{machine}];
\]

%dsp-level:0
\[
(4)\quad\textbf{the two positive parts of the splitting:}
\]

%dsp-level:0
\[
\qquad \Sigma_{\mathrm B}=\Sigma_{\mathrm{cl}}(0)-\Sigma_{\mathrm{cl}}(+1)=\ln\tfrac43=2\ln2-\ln3
\]

%dsp-level:0
\[
\qquad [\mathrm{Def\text{-}BM.13}(4);\ \mathrm{BM.59}(4)]:
\]

%dsp-level:0
\[
\boxed{\ \ln\frac43\ =\ \underbrace{\Bigl(\frac95\ln2-\ln3\Bigr)}_{H(0)-H(+1)}\ +\ \underbrace{\frac15\ln2}_{\langle\ln\prod_im\rangle_0-\langle\ln\prod_im\rangle_{+1}}\ }
\]

%dsp-level:0
\[
\qquad [\text{machine: the prime-log coefficients agree}];
\]

%dsp-level:0
\[
\frac95\ln2-\ln3>0\qquad[\text{machine:}\ 9\ln2>5\ln3\iff2^9=512>243=3^5];\qquad\frac15\ln2>0\qquad[\text{machine}];
\]

%dsp-level:0
\[
\textbf{the transport splitting is exactly two positive parts}\ —\ \textbf{a superposition-entropy part and a multiplicity part};
\]

%dsp-level:0
\[
\qquad \textbf{the snapshot read carries neither:}
\]

%dsp-level:0
\[
\text{a fixed triple at the vertex books}\ H=0\ \ \text{and loses}\ \ \frac95\ln2-\ln3\ \ \text{of the splitting}
\]

%dsp-level:0
\[
\qquad [\text{machine: the exact difference}];
\]

%dsp-level:0
\[
(5)\quad\textbf{the factorization (two independent exact averages):}
\]

%dsp-level:0
\[
\qquad \text{the fibre weights are class-fixed boundary data}\ [\mathrm{Def\text{-}BM.5}(1);\ \mathrm{Def\text{-}BM.13}(6)],
\]

%dsp-level:0
\[
\ \text{and the two registers' alphabets are disjoint}:
\]

%dsp-level:0
\[
E_{\mathrm{int}}=|R|\ \ \text{has no endpoint summand}
\]

%dsp-level:0
\[
\qquad [\mathrm{Def\text{-}BM.13}(6);\ \mathrm{BM.62}(11);\ \mathrm{BM.89}(8);\ \text{machine: the token census}]\implies\boxed{\ P(R,\ \text{fibre member})\ =\ P(R)\cdot P(\text{fibre member}\,|\,Q)\ }
\]

%dsp-level:0
\[
\textbf{the baryon's full ledger read is the composition of the flicker integral over}\ |R|\ (\mathrm{BM.94}(2))\ \textbf{and the superposition average over the fibre}\ (2);
\]

%dsp-level:0
\[
\neg\text{a cross-register equation is formed}\ —\ \text{the two registers are distinct alphabets and are}\ \neg\text{equated}
\]

%dsp-level:0
\[
\qquad [\mathrm{BM.37}(2);\ \mathrm{BM.83}(3)];
\]

%dsp-level:0
\[
(6)\quad\textbf{the fixed-state read refuted with an exact bound:}
\]

%dsp-level:0
\[
\qquad \max_{\text{fibre}}P(\cdot\,|\,0)=\frac25<\frac12
\]

%dsp-level:0
\[
\qquad [\text{machine:}\ 16<20];
\]

%dsp-level:0
\[
\qquad \max_{\text{fibre}}P(\cdot\,|\,+1)=\frac4{15}<\frac13
\]

%dsp-level:0
\[
\qquad [\text{machine:}\ 12<15]:
\]

%dsp-level:0
\[
\boxed{\ \neg\exists\ \text{a single vertex configuration carrying more than}\ \tfrac25\ \text{of the neutral class's transport weight}\ }
\]

%dsp-level:0
\[
\textbf{the refutation is an exact rational bound,}\ \neg\textbf{an argument about size};
\]

%dsp-level:0
\[
\qquad \text{both classes are strict mixtures at every admissible threshold}\ [(2);\ \mathrm{BM.19}(6)];
\]

%dsp-level:0
\[
(7)\quad\textbf{status:}
\]

%dsp-level:0
\[
\qquad \textbf{derived:}\ (1)\text{--}(6);
\]

%dsp-level:0
\[
\qquad \neg\text{a declaration},
\]

%dsp-level:0
\[
\quad\neg\text{a P4 addition},
\]

%dsp-level:0
\[
\quad\neg\text{a new dimensionful scale}\ [\mathrm{P1}],
\]

%dsp-level:0
\[
\quad\neg\text{a measured value}\ [\mathrm{P5}],
\]

%dsp-level:0
\[
\quad\text{§Cal-BM stays empty};
\]

%dsp-level:0
\[
\mathrm{Def\text{-}BM.13}(1)\text{--}(6),\ \mathrm{BM.28}(2),(4),(5),\ \mathrm{BM.59}(4),\ \mathrm{BM.61}(5)\ \text{stand unmodified};\qquad
\]

\textbf{Req.} D2.6, D3.4 (Parts 1–2); EM.9, CT0.9 (Part 2); P1, P4, P5 (Parts 1–2); Def-H.32(1), FT.40, FT.106(3) (Part 3); Census FT.35(2) (row-wise); Def-BM.5(1), Def-BM.8(6), Def-BM.13(1), Def-BM.13(3), Def-BM.13(4), Def-BM.13(6), BM.19(6), BM.25(4), BM.28(2), BM.28(5), BM.37(2), BM.59(4), BM.61(5), BM.62(11), BM.83(3), BM.89(8), BM.94(2) (this part); exhaustive ternary-fibre enumeration, exact rational weight normalization, prime-log ring arithmetic with integer sign certificates, symbol censuses [math; machine: \texttt{checks/\allowbreak{}verify\_\allowbreak{}H\_\allowbreak{}flicker\_\allowbreak{}integral\_\allowbreak{}restmass.py} §B].

\textbf{Pf.}

%dsp-level:0
\[
(1):\ m_q=(2,1,1)\ \text{and the end-pair census re-cited}\ [\mathrm{Def\text{-}BM.13}(1)];
\]

%dsp-level:0
\[
\qquad W(Q)\ \text{by two agreeing routes: the exhaustive}\ 4^3\ \text{census and the coefficient of}\ \bigl(2+z+z^{-1}\bigr)^3\ [\text{machine}];
\]

%dsp-level:0
\[
\qquad \textstyle\sum_QW(Q)=(2+1+1)^3=64\ [\text{machine}]
\]

%dsp-level:0
\[
(2):\ P
\]

%dsp-level:0
\[
= \prod_im_{q_i}/W(Q)\ \text{and}\ \textstyle\sum_{\text{fibre}}\prod_im_{q_i}
\]

%dsp-level:0
\[
= W(Q)\ \text{by}\ (1)\text{'s factorization identity}\ [\text{machine: class by class}];
\]

%dsp-level:0
\[
\qquad \text{the}\ Q{=}0\ \text{fibre: one triple of weight}\ 2^3=8\ \text{and six of weight}\ 2\cdot1\cdot1=2,\ \text{total}\ 20\ [\text{machine}];
\]

%dsp-level:0
\[
\qquad \text{the}\ Q{=}{+}1\ \text{fibre: three of weight}\ 1\cdot2\cdot2=4\ \text{and three of weight}\ 1,\ \text{total}\ 15\ [\text{machine}]
\]

%dsp-level:0
\[
\text{the symbol census:}\ W(Q)\in\mathbb Z[z,z^{-1}]\ \text{with coefficients fixed by}\ (m_q,k),
\]

%dsp-level:0
\[
\ \text{and}\ \{g,g',e,\varepsilon_{\mathrm{EM}},\hat\mu,\hat\mu',\ell_c\}\cap\{z\}
\]

%dsp-level:0
\[
= \varnothing\ [\mathrm{BM.61}(5)\ \text{re-cited};\ \text{machine}]
\]

%dsp-level:0
\[
(3):\ H(0)
\]

%dsp-level:0
\[
= -\tfrac25\ln\tfrac25-6\cdot\tfrac1{10}\ln\tfrac1{10}
\]

%dsp-level:0
\[
= \tfrac25(\ln5-\ln2)+\tfrac35(\ln2+\ln5)
\]

%dsp-level:0
\[
= \ln5+\tfrac15\ln2\ [\text{machine: prime-log coefficients}];
\]

%dsp-level:0
\[
\qquad \langle\ln\prod_im\rangle_0=\tfrac{8\ln8+6\cdot2\ln2}{20}=\tfrac{36\ln2}{20}=\tfrac95\ln2\ [\text{machine}]
\]

%dsp-level:0
\[
H(+1)=-3\cdot\tfrac4{15}\ln\tfrac4{15}-3\cdot\tfrac1{15}\ln\tfrac1{15}=\ln15-\tfrac85\ln2\ [\text{machine}];
\]

%dsp-level:0
\[
\qquad \langle\ln\prod_im\rangle_{+1}=\tfrac{3\cdot4\ln4+3\cdot1\ln1}{15}=\tfrac{24\ln2}{15}=\tfrac85\ln2\ [\text{machine}];
\]

%dsp-level:0
\[
\qquad H+\langle\ln\prod_im\rangle=\ln W\ \text{in both classes}\ [\text{machine}]
\]

%dsp-level:0
\[
(4):\ [\ln5+\tfrac15\ln2]-[\ln3+\ln5-\tfrac85\ln2]
\]

%dsp-level:0
\[
= \tfrac95\ln2-\ln3\ \text{and}\ \tfrac95\ln2-\tfrac85\ln2
\]

%dsp-level:0
\[
= \tfrac15\ln2\ [\text{machine}];
\]

%dsp-level:0
\[
\qquad \text{the sum:}\ \tfrac95\ln2-\ln3+\tfrac15\ln2=2\ln2-\ln3=\ln\tfrac43\ [\text{machine}];
\]

%dsp-level:0
\[
\qquad 2^9=512>243=3^5\ [\text{machine: an integer power comparison}]
\]

%dsp-level:0
\[
(5):\ \text{the boundary data fix the charge class}\ [\mathrm{Def\text{-}BM.5}(1)]\implies\text{the fibre weights are class-constant on the ensemble};
\]

%dsp-level:0
\[
\qquad E_{\mathrm{int}}\ \text{is constant on every end-sign fibre}\ [\mathrm{Def\text{-}BM.13}(6)]\implies\text{the two registrations are independent and their probabilities multiply}\ [\text{machine}]
\]

%dsp-level:0
\[
(6):\ \tfrac25<\tfrac12\iff16<20;
\]

%dsp-level:0
\[
\qquad \tfrac4{15}<\tfrac13\iff12<15\ [\text{machine: integers}];
\]

%dsp-level:0
\[
\qquad (7):\ \text{the standing-rule scans}\ [\text{machine}].\ \blacksquare
\]

\textbf{Phys.} Ask which state a constituent occupies at a vertex and the ledger answers with a list, not a value. A neutral-core constituent carries two transport realizations and a charged one carries a single realization, so the class fibre is a weighted set: at charge zero one triple of weight eight against six of weight two, summing to twenty; at charge one three of weight four against three of weight one, summing to fifteen. Normalize and the superposition is exact — two fifths on the all-neutral triple, a tenth on each of the six mixed ones — and the numbers are rationals in an alphabet that carries no coupling at all, which is the coupling-blindness theorem wearing a different hat. The class entropy then splits cleanly into the Shannon entropy of that superposition plus the mean logarithm of the multiplicities, and the baryon's transport splitting, the log of four thirds, is the sum of two strictly positive parts: nine fifths of a bit less the log of three, from the superposition, and one fifth of a bit, from the multiplicities. A read that fixes one triple at the vertex books no entropy and forfeits the larger part of the splitting, so the fixed-state read is not merely incomplete, it is quantitatively short by an exact prime-log element. The two registers stay independent — the fibre weights are boundary data and the edge count has no endpoint summand — so the baryon's ledger read is the composition of two exact averages, and no equation between the registers is formed. The sharpest row is the bound: no single vertex configuration carries more than two fifths of the neutral class's transport weight, certified by sixteen against twenty.

\end{theorem}

\begin{theorem}[{\texorpdfstring{BM.96 — The rung-one $\nu$ line: the seal-derived rest mass, the single-ray occupancy, and the four bars against a seesaw-type mechanism}{BM.96 - The rung-one line: the seal-derived rest mass, the single-ray occupancy, and the four bars against a seesaw-type mechanism}}]\label{tag:BM-96}
\textbf{Thm.}

%dsp-level:0
\[
(1)\quad\textbf{the line's structure:}
\]

%dsp-level:0
\[
\qquad \text{a rung-one bound chain}\ [\mathrm{D1.5}]\ \text{whose terminus-active aspects lie in}\ \mathcal A(1)
\]

%dsp-level:0
\[
= \{J,\,G,\,W,\,H,\,D^+\}\ \ [\mathrm{CT0.13}]:
\]

%dsp-level:0
\[
\text{EM unactivated}\ (f_v=0),
\]

%dsp-level:0
\[
\qquad \text{the colour block vanishing},
\]

%dsp-level:0
\[
\qquad \delta h(c)\ge2,\quad n(v)\ge2,
\]

%dsp-level:0
\[
\qquad \Sigma_{\min}\ \text{the minimal instance}
\]

%dsp-level:0
\[
\qquad [\mathrm{FT.142}(1);\ \mathrm{D3.3};\ \mathrm{CT3.11}(1);\ \mathrm{FT.19};\ \mathrm{Def\text{-}H.4};\ \mathrm{Census\ FT.28}];
\]

%dsp-level:0
\[
\dim\mathcal K_1=16=C_{\mathrm{cap}}
\]

%dsp-level:0
\[
\qquad [\mathrm{FT.142}(2),(3);\ \mathrm{CT0.14};\ \mathrm{CT0.12}:\ \text{the unique derived candidate, no interface declaration}];
\]

%dsp-level:0
\[
\boxed{\ \delta_\nu=\frac{\sum_i[J_i,K_i]\big|_\Sigma}{\dim\mathcal K_1}=\frac2{16}=\frac18\ }
\]

%dsp-level:0
\[
\qquad [\mathrm{FT.142}(4),(5):\ \text{the undressed seal coefficient over the derived mask dimension}]
\]

%dsp-level:0
\[
\qquad \textbf{derived,}\ \neg\textbf{declared};
\]

%dsp-level:0
\[
(2)\quad\textbf{the rest mass from the rung-one seal structure:}
\]

%dsp-level:0
\[
\qquad m_\nu^{\mathrm{rest}\,2}=\hat\mu_\nu^2\bigl(\ell_c^{(\nu)}\bigr)^2\,G\bigl(c(i_\nu)\bigr),
\]

%dsp-level:0
\[
\qquad G\bigl(c(i)\bigr)=A_0-B_0\cos\Bigl(\delta_\nu-\frac{2\pi i}3\Bigr)
\]

%dsp-level:0
\[
\qquad [\mathrm{Def\text{-}H.37}(1);\ \mathrm{Def\text{-}BM.11}(5);\ \mathrm{BM.39}(1)]:
\]

%dsp-level:0
\[
A_0=\frac{3g^2}8+\frac{g'^2}{16},\qquad B_0=\frac{gg'}4\qquad[\mathrm{BM.39}(1);\ \text{machine}];\qquad\textbf{at}\ \delta_\nu=\tfrac18,\ \ c:=\cos\tfrac18,\ \ s_\nu:=\sin\tfrac18:
\]

%dsp-level:0
\[
\boxed{\ G_0=A_0-B_0\,c,\qquad G_1=A_0+\frac{B_0}2\bigl(c-\sqrt3\,s_\nu\bigr),\qquad G_2=A_0+\frac{B_0}2\bigl(c+\sqrt3\,s_\nu\bigr)\ }
\]

%dsp-level:0
\[
\qquad [\text{machine: exact in}\ \mathbb Q[\sqrt2,\sqrt3][c,s_\nu]/(c^2{+}s_\nu^2{-}1)];
\]

%dsp-level:0
\[
\textbf{the exact order at}\ (g,g')>0\ [\mathrm{EW.33}]:\qquad\boxed{\ G_2\ >\ G_1\ >\ A_0\ >\ G_0\ }
\]

%dsp-level:0
\[
G_1>A_0\iff c>\sqrt3\,s_\nu\iff c^2>\tfrac34\iff\cos\tfrac14=2c^2-1>\tfrac12
\]

%dsp-level:0
\[
\qquad [\text{machine:}\ \cos x>1-\tfrac{x^2}2=\tfrac{31}{32}>\tfrac12\ \text{at}\ x=\tfrac14,\ \text{the exact alternating-series sign row}];
\]

%dsp-level:0
\[
c>0\ \ \text{on}\ \ c>\tfrac{63}{64},
\]

%dsp-level:0
\[
\qquad s_\nu>0\ \ \text{on}\ \ s_\nu>\tfrac18-\tfrac1{3072}
\]

%dsp-level:0
\[
\qquad [\text{machine: exact Taylor enclosures with rational remainder bounds}];
\]

%dsp-level:0
\[
\textbf{the affine form on the spectral triplet:}
\]

%dsp-level:0
\[
\qquad G\bigl(c(i)\bigr)=A_0-\frac{B_0}{\sqrt2}\bigl(r_i^{(\nu)}-1\bigr),
\]

%dsp-level:0
\[
\qquad r_i^{(\nu)}=1+\sqrt2\cos\Bigl(\frac18+\frac{2\pi i}3\Bigr)
\]

%dsp-level:0
\[
\qquad [\mathrm{Def\text{-}BM.11}(5);\ \mathrm{FT.143}(1);\ \text{machine}]:
\]

%dsp-level:0
\[
\bigl(r_1^{(\nu)}\bigr)^2<\bigl(r_2^{(\nu)}\bigr)^2<\bigl(r_0^{(\nu)}\bigr)^2
\]

%dsp-level:0
\[
\qquad [\mathrm{FT.143}(3):\ \text{enclosure-proven, frozen}]\implies\text{the transport read is strictly decreasing in}\ r;
\]

%dsp-level:0
\[
\textbf{the value set is orbit-label independent:}
\]

%dsp-level:0
\[
\qquad \bigl\{G\bigl(c(i)\bigr)\bigr\}_i=\Bigl\{A_0-B_0\,c,\ \ A_0+\tfrac{B_0}2\bigl(c-\sqrt3\,s_\nu\bigr),\ \ A_0+\tfrac{B_0}2\bigl(c+\sqrt3\,s_\nu\bigr)\Bigr\}
\]

%dsp-level:0
\[
\qquad [\mathrm{Def\text{-}BM.11}(5)\ \text{orbit-label row};\ \text{machine: the two routes give the same set}];
\]

%dsp-level:0
\[
\boxed{\ \textbf{the transport order is the reverse of the spectral order:}\ \ \mathrm{BM.39}(5)\text{'s inversion instantiated at rung one}\ }
\]

%dsp-level:0
\[
(3)\quad\textbf{positivity and the exact floor:}
\]

%dsp-level:0
\[
\qquad A_0-B_0=\frac38\Bigl(g-\frac{g'}3\Bigr)^2+\frac{g'^2}{48}>0,
\]

%dsp-level:0
\[
\qquad A_0+B_0=\frac38\Bigl(g+\frac{g'}3\Bigr)^2+\frac{g'^2}{48}>0
\]

%dsp-level:0
\[
\qquad [\mathrm{BM.39}(1);\ \mathrm{BM.77}(4);\ \text{machine: the SOS expansions}]:
\]

%dsp-level:0
\[
\sqrt{A_0-B_0}\ \le\ \sqrt{G\bigl(c(i)\bigr)}\ \le\ \sqrt{A_0+B_0}\qquad\forall\,\text{line}\ i,\ \ \forall\,\text{channel},\ \ \forall\,(g,g')\ne(0,0)\qquad[|\cos|\le1;\ \text{machine}];
\]

%dsp-level:0
\[
\text{the kernel line is unreachable:}\ \cos\varphi=0\wedge\cos(\varphi+\tfrac{2\pi}3)=0\ \text{has no solution}
\]

%dsp-level:0
\[
\qquad [\mathrm{BM.39}(1);\ \text{machine}]\implies\boxed{\ m_\nu^{\mathrm{rest}}>0\ \ \text{strictly, and}\ \neg\text{a massless}\ \nu\ \text{line is formed on the transport species}\ }
\]

%dsp-level:0
\[
\textbf{the absolute value is}\ \neg\textbf{derived:}
\]

%dsp-level:0
\[
\qquad \bigl(\hat\mu_\nu,\ \ell_c^{(\nu)}\bigr)\ \text{undeclared}\ [\mathrm{P4};\ \mathrm{Cal\text{-}F}];
\]

%dsp-level:0
\[
\qquad \rho_\nu:=\mathfrak h\,\hat\mu_\nu\,\ell_c^{(\nu)}\in\mathfrak q^0\ \ \text{is a restatement,}\ \neg\text{a new member}\ [\mathrm{BM.77}(2);\ \mathrm{CT2.1}];
\]

%dsp-level:0
\[
\textbf{the internal ratios are derived:}
\]

%dsp-level:0
\[
\qquad \frac{m_i^2}{m_j^2}=\frac{G\bigl(c(i)\bigr)}{G\bigl(c(j)\bigr)}\ \ \text{exact symbolic in}\ (A_0,B_0)
\]

%dsp-level:0
\[
\qquad [\mathrm{FT.143}(3);\ \mathrm{BM.80}(6);\ \text{machine}];
\]

%dsp-level:0
\[
\textbf{the inertial companion:}
\]

%dsp-level:0
\[
\qquad m_i=\mu_\nu\bigl(r_i^{(\nu)}\bigr)^2,\quad\mu_\nu\ \text{undeclared}\ [\mathrm{P4}],
\]

%dsp-level:0
\[
\qquad Q:=\frac{\sum_im_i}{\bigl(\sum_i\sqrt{m_i}\bigr)^2}=\frac69=\frac23\ \ \text{exactly}
\]

%dsp-level:0
\[
\qquad [\mathrm{FT.143}(4);\ \text{machine}];
\]

%dsp-level:0
\[
\textbf{the informational companion:}\qquad m_\nu\,\mathfrak h\ \ge\ 2\ln2\qquad[\mathrm{BM.70}(2),(3);\ \mathrm{BM.13}(3):\ \text{the two-bit floor}];
\]

%dsp-level:0
\[
(4)\quad\textbf{the single-ray occupancy:}
\]

%dsp-level:0
\[
\qquad \text{the materialized locus at rung one is single-arrow}\ —\ \mathrm{D1.2}\text{'s domains force strict alternation and the prefix law, and a dissolution-first chain is inadmissible at its first event}\ [\mathrm{FT.144}(2);\ \mathrm{D2.11}];
\]

%dsp-level:0
\[
\text{under A-BM.7's issued coupling the arrow locus occupies}\ \operatorname{im}P_L,
\]

%dsp-level:0
\[
\ \ \text{the}\ \hat\chi
\]

%dsp-level:0
\[
= +1\ \text{ray of}\ \mathrm{Def\text{-}H.33}\text{'s boxed sign}:
\]

%dsp-level:0
\[
\qquad \boxed{\ \text{the occupancy is exactly}\ (1,\,0)\ }
\]

%dsp-level:0
\[
\qquad [\mathrm{FT.144}(5);\ \mathrm{BM.73}(2),(3);\ \mathrm{FT.122}(4)];
\]

%dsp-level:0
\[
\textbf{the complementary ray's content is the vacuum sector:}
\]

%dsp-level:0
\[
\qquad n_{\mathcal R}=0\implies\rho_{pq}=\rho_\circ=\tfrac14I,\quad S_{\mathrm{raw}}=\ln4=\max,\quad\text{unregistered},\quad\neg\exists\ \text{an accessible pole at any rung}
\]

%dsp-level:0
\[
[\mathrm{FT.144}(4);\ \mathrm{A5};\ \mathrm{CT0.2};\ \mathrm{CT0.1};\ \mathrm{CT4.2}:\ \mathcal P_MF\mathcal P_M=0\implies\neg\exists\ \text{an accessible pole};\ \mathrm{CT2.5}:\ \text{poles}\iff1\in\operatorname{spec}\mathbf F_M];
\]

%dsp-level:0
\[
\textbf{the obstruction is a transition-domain restriction prior to every read}\ —\ \neg\text{a mass statement},
\]

%dsp-level:0
\[
\ \neg\text{a suppression},
\]

%dsp-level:0
\[
\ \neg\text{a decoupling}
\]

%dsp-level:0
\[
\qquad [\mathrm{FT.144}(3):\ \neg\mathfrak h,\ \neg\mu,\ \neg h_0\ \text{in the premises}]:
\]

%dsp-level:0
\[
\boxed{\ \neg\exists\ \text{a materializable right-ray}\ \nu\ \text{line in the matter orientation, in any valid universe}\ }
\]

%dsp-level:0
\[
\qquad [\mathrm{D2.11};\ \text{A-BM.7}]
\]

%dsp-level:0
\[
\textbf{the naming discipline:}
\]

%dsp-level:0
\[
\qquad \text{``left'' and ``right'' are the ray names of}\ \mathrm{Def\text{-}H.33}\ \text{as fixed by A-BM.7};
\]

%dsp-level:0
\[
\qquad \text{the handedness vocabulary is external and is}\ \neg\text{identified with a ray}\ [\text{A-BM.7's bars};\ \mathrm{P5}];
\]

%dsp-level:0
\[
(5)\quad\textbf{the image row:}
\]

%dsp-level:0
\[
\qquad \text{the vertex's neutral member is the}\ \nu\text{-channel line in the image orientation}\ [\text{A-BM.5};\ \mathrm{BM.41}(1);\ \mathrm{BM.45}];
\]

%dsp-level:0
\[
\text{the image is the}\ \hat\omega\text{-image of the same single-arrow locus:}
\]

%dsp-level:0
\[
\qquad \hat\omega P_L\hat\omega^{-1}=P_R,
\]

%dsp-level:0
\[
\qquad \hat\omega|L\rangle=-i|R\rangle,
\]

%dsp-level:0
\[
\qquad \hat\omega^2=\mathrm{id}
\]

%dsp-level:0
\[
\qquad [\mathrm{D2.6};\ \mathrm{FT.122}(4);\ \mathrm{BM.73}(4);\ \text{machine over}\ \mathbb Q(i)];
\]

%dsp-level:0
\[
\text{and it sits in the}\ \omega\text{-image representation}\ \{0,-1\}\qquad[\mathrm{EW.36};\ \text{A-BM.7}(\mathrm b);\ \mathrm{BM.25};\ \mathrm{BM.27}];
\]

%dsp-level:0
\[
\mathrm{CPT}:\ (q,\ \text{orientation})
\]

%dsp-level:0
\[
\mapsto (-q,\ \text{flipped})\ \ \text{has}\ \textbf{no fixed sector},
\]

%dsp-level:0
\[
\ \ \text{its orbits are three pairs of size two},
\]

%dsp-level:0
\[
\ \ q
\]

%dsp-level:0
\[
\mapsto -q\ \text{fixes}\ q\qquad = 0\ \text{alone}
\]

%dsp-level:0
\[
\qquad [\mathrm{BM.25}(4);\ \text{machine}]:
\]

%dsp-level:0
\[
\boxed{\ \textbf{the image is one orbit of size two and carries no independent datum}\ —\ \neg\text{a second state},\ \neg\text{a right-ray matter line},\ \neg\text{a name for an unmaterializable object}\ }
\]

%dsp-level:0
\[
\qquad [\mathrm{BM.27}(6)\ \text{pattern}];
\]

%dsp-level:0
\[
(6)\quad\textbf{the no-seesaw theorem (four independent bars, each re-cited at its own locus):}
\]

%dsp-level:0
\[
(\mathrm S1)\quad\textbf{no partner state:}\ \neg\exists\ \text{a right-ray matter-orientation state to pair with}\ —\ \text{the complement's occupancy is exactly zero and its content supports no accessible pole at any rung}
\]

%dsp-level:0
\[
\qquad [(4);\ \mathrm{FT.144}(4),(5)];
\]

%dsp-level:0
\[
(\mathrm S2)\quad\textbf{no self-pairing:}\ \neg\exists\ \text{a fixed CPT sector, so a Majorana-type identification}\ C
\]

%dsp-level:0
\[
= \mathrm{CPT}\cdot C\ \text{is}\ \neg\text{formed}
\]

%dsp-level:0
\[
\qquad [(5);\ \mathrm{BM.25}(4);\ \mathrm{BM.27}(6)];
\]

%dsp-level:0
\[
(\mathrm S3)\quad\textbf{no coupling object:}\ \neg\exists\ \text{a Yukawa-type object in the inventory}\ —\ \dim F_H
\]

%dsp-level:0
\[
= 1
\]

%dsp-level:0
\[
= \operatorname{span}\{\Sigma\},
\]

%dsp-level:0
\[
\ \neg\text{a scalar-block coupling to a chain line},
\]

%dsp-level:0
\[
\ \neg\text{a dimension-five operator}
\]

%dsp-level:0
\[
\qquad [\mathrm{CT0.12};\ \mathrm{EW.15};\ \mathrm{FT.145}(3);\ \mathrm{BM.93}(5)];
\]

%dsp-level:0
\[
(\mathrm S4)\quad\textbf{no second scale:}\ \exists!\ \text{one free dimensionful scale}\ \mathfrak h\ —\ \text{the heavy scale a seesaw-type mechanism requires is barred}
\]

%dsp-level:0
\[
\qquad [\mathrm{P1};\ \text{A §0.3};\ \mathrm{D2.9}];
\]

%dsp-level:0
\[
\boxed{\ \neg\exists\ \text{a seesaw-type mechanism formable in this corpus};\qquad\text{the}\ \nu\ \text{rest mass is positive, derived in form from the rung-one seal, and undeclared in scale}\ }
\]

%dsp-level:0
\[
\textbf{the companion row:}
\]

%dsp-level:0
\[
\qquad \text{the smallness question is a scale question}\ [\mathrm{P4}],
\]

%dsp-level:0
\[
\ \neg\text{a mechanism question}\ —\ \neg\text{a suppression is needed, so}\ \neg\text{one is invented};
\]

%dsp-level:0
\[
(7)\quad\textbf{the vacuum residence and the vertex firing (the two states and the one door):}
\]

%dsp-level:0
\[
\qquad \textbf{state (i), the vacuum-resident balanced fluctuation pair:}
\]

%dsp-level:0
\[
n_{\mathcal R}=0,
\]

%dsp-level:0
\[
\qquad \sum_n\langle\varnothing|\hat U^n|\varnothing\rangle\,w^n=\frac1{1-w^2x_+x_-},
\]

%dsp-level:0
\[
\qquad \text{every nonzero return monomial carries equal}\ x_+\text{- and}\ x_-\text{-degree}
\]

%dsp-level:0
\[
\qquad [\mathrm{FT.108}(5);\ \mathrm{FT.144}(4);\ \text{machine in}\ \mathbb Q[w,x_+,x_-]/(w^{12})]:
\]

%dsp-level:0
\[
\textbf{state (ii), the materialized rung-one seal:}
\]

%dsp-level:0
\[
\qquad |E_{\mathcal R}|\ge2,
\]

%dsp-level:0
\[
\qquad m_{\mathrm{hid}}=\frac{2\ln2}{\mathfrak h}\ \ \text{at the minimal floor}
\]

%dsp-level:0
\[
\qquad [\mathrm{FT.145}(2);\ \mathrm{Census\ FT.28};\ \mathrm{Def\text{-}H.4}];
\]

%dsp-level:0
\[
\textbf{the door is the vertex:}
\]

%dsp-level:0
\[
\qquad \Gamma(\varepsilon_q{=}0)=0\ \ \text{exactly},
\]

%dsp-level:0
\[
\qquad x_+x_-=4\varepsilon_q(1-4\varepsilon_q)\le\frac14\ \ \text{with equality at}\ \ \varepsilon_q=\frac18
\]

%dsp-level:0
\[
\qquad [\mathrm{BM.23}(7),(8);\ \text{machine}]:
\]

%dsp-level:0
\[
\boxed{\ \textbf{the line materializes only where the flicker is present}\ —\ \text{the same threshold}\ \mathrm{BM.94}\text{'s integral averages over is the one that opens the vertex}\ }
\]

%dsp-level:0
\[
\textbf{the materialized line propagates as a rung-one seal, forward-only:}
\]

%dsp-level:0
\[
\qquad \neg\exists\,\tau^{-1},
\]

%dsp-level:0
\[
\qquad \neg\exists\ \text{an inverse decay in the ledger}
\]

%dsp-level:0
\[
\qquad [\mathrm{A6};\ \mathrm{BM.23}(9);\ \mathrm{FT.144}(2)];
\]

%dsp-level:0
\[
\Gamma\ \text{is a weight per tick},\ \neg\text{a rate per external second}
\]

%dsp-level:0
\[
\qquad [\mathrm{BM.23}(9);\ \mathrm{D1.8};\ \mathrm{P5}];
\]

%dsp-level:0
\[
\qquad \neg\text{a lifetime, endpoint or mixing angle is derived or compared}\ [\mathrm{P5}];
\]

%dsp-level:0
\[
(8)\quad\textbf{status:}
\]

%dsp-level:0
\[
\qquad \textbf{derived:}\ (1)\text{--}(7)\ \text{in form}\ —\ \text{the angle, the order, the positivity, the occupancy, the four bars, the two states and the door};
\]

%dsp-level:0
\[
\neg\textbf{derived:}\ \ \text{the absolute scale}\ (\hat\mu_\nu,\ell_c^{(\nu)}),\ \ \text{the line index}\ i_\nu\,\ \ \text{the mixing matrix}\ [\mathrm{FT.145}(3)];
\]

%dsp-level:0
\[
\neg\text{a declaration issued}\ —\ \text{A-BM.5, A-BM.7 and}\ \mathrm{Def\text{-}H.37}\ \text{are re-cited,}\ \neg\text{re-issued};
\]

%dsp-level:0
\[
\qquad \neg\text{a P4 addition},\quad\neg\text{a new dimensionful scale}\ [\mathrm{P1}],\quad\neg\text{a measured value}\ [\mathrm{P5}],\quad\text{§Cal-BM stays empty};
\]

%dsp-level:0
\[
\text{species names are confined to the identification layer}\ [\mathrm{Def\text{-}BM.6}(2),(7)];
\]

%dsp-level:0
\[
\qquad \text{topology labels are internal bookkeeping}\ [\mathrm{Def\text{-}H.32}(4)];
\]

%dsp-level:0
\[
\qquad \text{amplitude}\ne\text{mass}\ [\mathrm{FT.79}(5)];
\]

%dsp-level:0
\[
\qquad \text{rung}\ne\text{scale}\ [\mathrm{D0.8}].
\]

\textbf{Req.} A5, A8, D0.8, D1.2, D1.4, D1.5, D1.8, D2.6, D2.9, D2.11, D3.3 (Parts 1–2); A §0.3, CT0.1, CT0.2, CT0.12, CT0.13, CT0.14, CT3.11 (Part 1, Part 2); CT2.5, CT4.2, EW.15, EW.33, EW.36, EM.9 (Part 2); P1, P4, P5 (Parts 1–2); Def-H.4, Def-H.32(4), Def-H.33, Def-H.37(1), FT.19, FT.28, FT.40, FT.54(4), FT.79(5), FT.108(5), FT.122(4), FT.142(1), FT.142(2), FT.142(3), FT.142(4), FT.142(5), FT.143(1), FT.143(3), FT.143(4), FT.144(2), FT.144(3), FT.144(4), FT.144(5), FT.145(2), FT.145(3) (Part 3); Census FT.28 (row-wise); Def-BM.11(5), BM.13(3), BM.23(7), BM.23(8), BM.23(9), BM.25(4), BM.27(6), BM.39(1), BM.39(5), BM.41(1), BM.45, BM.70(2), BM.70(3), BM.73(2), BM.73(3), BM.73(4), BM.77(2), BM.77(4), BM.80(6), BM.93(5) (this part); Cal-F; exact identities in $\mathbb Q[\sqrt2,\sqrt3][c,s_\nu]/(c^2{+}s_\nu^2{-}1)$, exact Taylor enclosures with rational remainder bounds, sum-of-squares positivity, exact $\mathbb Q(i)$ block algebra, exact resolvent identities in a truncated polynomial ring [math; machine: \texttt{checks/\allowbreak{}verify\_\allowbreak{}H\_\allowbreak{}flicker\_\allowbreak{}integral\_\allowbreak{}restmass.py} §C].

\textbf{Pf.}

%dsp-level:0
\[
(1):\ \mathrm{FT.142}(1)\text{--}(5)\ \text{re-cited (frozen)};
\]

%dsp-level:0
\[
\qquad \delta_\nu=\tfrac2{16}=\tfrac18\ [\text{machine: exact Fractions}];
\]

%dsp-level:0
\[
\qquad \dim\mathcal K_1=6+(5+3+1+1)=16=C_{\mathrm{cap}}\ [\mathrm{CT0.12};\ \mathrm{CT0.14};\ \text{machine}]
\]

%dsp-level:0
\[
(2):\ \cos(\tfrac18-\tfrac{2\pi}3)
\]

%dsp-level:0
\[
= -\tfrac c2+\tfrac{\sqrt3}2s_\nu\ \text{and}\ \cos(\tfrac18-\tfrac{4\pi}3)
\]

%dsp-level:0
\[
= -\tfrac c2-\tfrac{\sqrt3}2s_\nu\ [\text{the addition formulas}]\implies\text{the three boxed values}\ [\text{machine: ring identities}];
\]

%dsp-level:0
\[
\qquad A_0,B_0\ \text{re-cited from}\ \mathrm{BM.39}(1)
\]

%dsp-level:0
\[
G_1>A_0\iff\tfrac{B_0}2(c-\sqrt3s_\nu)>0
\]

%dsp-level:0
\[
\iff c>\sqrt3s_\nu
\]

%dsp-level:0
\[
\iff c^2>3(1-c^2)
\]

%dsp-level:0
\[
\iff 4c^2>3
\]

%dsp-level:0
\[
\iff 2c^2-1>\tfrac12\iff\cos\tfrac14>\tfrac12\ [\text{machine: the double-angle identity}];
\]

%dsp-level:0
\[
\qquad \cos\tfrac14>1-\tfrac1{2\cdot16}
\]

%dsp-level:0
\[
= \tfrac{31}{32}>\tfrac12\ [\text{machine: the alternating-series sign row, exact rationals}]
\]

%dsp-level:0
\[
c>\tfrac{63}{64}\ \text{from}\ \cos\tfrac18>1-\tfrac1{2\cdot64};
\]

%dsp-level:0
\[
\qquad s_\nu>\tfrac18-\tfrac1{6\cdot8^3}=\tfrac18-\tfrac1{3072}>0\ [\text{machine: exact Taylor bounds}];
\]

%dsp-level:0
\[
\qquad G_2>G_1\iff\sqrt3\,s_\nu>-\sqrt3\,s_\nu\iff s_\nu>0\ [\text{machine}]
\]

%dsp-level:0
\[
G=A_0+\tfrac{B_0}{\sqrt2}\bigl(1-r^{(\nu)}\bigr)\ \text{is affine strictly decreasing in}\ r\ [\mathrm{Def\text{-}BM.11}(5)\ \text{affine row};\ \mathrm{FT.143}(1);\ \text{machine}];
\]

%dsp-level:0
\[
\qquad r_1<r_2<r_0\ [\mathrm{FT.143}(3)\ \text{re-cited}]\implies\text{the transport order reverses the spectral order};
\]

%dsp-level:0
\[
\ \text{the value set is orbit-label independent}\ [\mathrm{Def\text{-}BM.11}(5)\ \text{orbit-label row};\ \text{machine: set equality}]
\]

%dsp-level:0
\[
(3):\ \text{the two SOS rows re-derived by expansion}\ [\text{machine}];
\]

%dsp-level:0
\[
\qquad |\cos|\le1\implies\text{the sandwich}\ [\text{math}];
\]

%dsp-level:0
\[
\qquad \text{the kernel line:}\ c=0\implies s_\nu^2=1\implies\cos(\varphi+\tfrac{2\pi}3)=\mp\tfrac{\sqrt3}2\ne0\ [\mathrm{BM.39}(1)\ \text{re-cited};\ \text{machine}]
\]

%dsp-level:0
\[
m_i^2/m_j^2=G_i/G_j\ \text{and}\ Q=\tfrac69=\tfrac23\ [\mathrm{FT.143}(3),(4)\ \text{re-cited};\ \text{machine}];
\]

%dsp-level:0
\[
\qquad m_\nu\mathfrak h\ge2\ln2\ [\mathrm{BM.70}(2),(3)\ \text{re-cited}]
\]

%dsp-level:0
\[
(4):\ \mathrm{FT.144}(2)\text{'s domain rows and prefix law re-cited};
\]

%dsp-level:0
\[
\qquad \text{occupancy}\ (1,0)\ [\mathrm{FT.144}(5);\ \mathrm{BM.73}(3)];
\]

%dsp-level:0
\[
\qquad P_{L/R}^2=P_{L/R},\ P_LP_R=0,\ P_L+P_R=\mathrm{id},\ \mathfrak T P_{L/R}=1\ [\mathrm{FT.122}(2);\ \text{machine over}\ \mathbb Q(i)]
\]

%dsp-level:0
\[
\rho_\circ=\tfrac14I:\ \mathfrak T=1,\ \operatorname{charpoly}=(x-\tfrac14)^4,\ S_{\mathrm{raw}}=\ln4=\max\ [\mathrm{CT0.2};\ \mathrm{BM.35}(1)\ \text{re-cited}];
\]

%dsp-level:0
\[
\qquad \mathcal P_MF\mathcal P_M=0\implies\neg\exists\ \text{an accessible pole}\ [\mathrm{CT4.2};\ \mathrm{CT2.5}\ \text{re-cited}];
\]

%dsp-level:0
\[
\qquad \neg\mathfrak h,\neg\mu,\neg h_0\ \text{in the premises}\ [\mathrm{FT.144}(3)\ \text{re-cited}]
\]

%dsp-level:0
\[
(5):\ \hat\omega P_L\hat\omega^{-1}
\]

%dsp-level:0
\[
= P_R,
\]

%dsp-level:0
\[
\ \hat\omega|L\rangle
\]

%dsp-level:0
\[
= -i|R\rangle,
\]

%dsp-level:0
\[
\ \hat\omega^2
\]

%dsp-level:0
\[
= \mathrm{id}\ [\mathrm{FT.122}(4);\ \mathrm{BM.73}(4)\ \text{re-derived entrywise over}\ \mathbb Q(i);\ \text{machine}];
\]

%dsp-level:0
\[
\qquad \text{the no-fixed-sector census on six sectors}\ [\mathrm{BM.25}(4)\ \text{re-cited};\ \text{machine}]
\]

%dsp-level:0
\[
(6):\ (\mathrm S1)\text{--}(\mathrm S4)\ \text{re-cited at their own loci}\ [(4),(5);\ \mathrm{CT0.12};\ \mathrm{EW.15};\ \mathrm{FT.145}(3);\ \mathrm{P1};\ \text{A §0.3}];
\]

%dsp-level:0
\[
\qquad \text{the inventory scans:}\ \neg\exists\ \text{a Yukawa-type token outside a negation row},
\]

%dsp-level:0
\[
\ \neg\exists\ \text{a second dimensionful scale}\ [\text{machine}]
\]

%dsp-level:0
\[
(7):\ \text{the resolvent identity re-derived coefficient-wise to}\ w^{12}\ [\text{machine}];
\]

%dsp-level:0
\[
\qquad x_+x_-=4\varepsilon_q(1-4\varepsilon_q)\le\tfrac14\ \text{with equality at}\ \varepsilon_q=\tfrac18\ [\mathrm{BM.23}(7)\ \text{re-cited};\ \text{machine}];
\]

%dsp-level:0
\[
\qquad \Gamma=0\ \text{at}\ \varepsilon_q=0\ [\mathrm{BM.23}(8)];
\]

%dsp-level:0
\[
\qquad m_{\mathrm{hid}}=2\ln2/\mathfrak h\ [\mathrm{FT.145}(2)]
\]

%dsp-level:0
\[
(8):\ \text{the standing-rule scans: external tokens, decimal numerals, registry tokens, declaration lines, the barred object}\ \alpha_B\ [\text{machine}].\ \blacksquare
\]

\textbf{Phys.} The neutral line the vertex emits is a rung-one seal, and everything about it that the corpus can derive, it derives. Its channel angle is two over sixteen, the undressed seal coefficient over the mask dimension the complete-force rule already fixed; its rest mass is the transport defect evaluated on that angle, three exact values in a radical ring, ordered heaviest to lightest against the spectral order by the same inversion the family mass form proves in general. Every one of the three is strictly positive, squeezed between two sums of squares that never touch zero, and the kernel line that would make one of them massless is unreachable by an exact trigonometric census — so there is no massless option to decline and no zero to explain away. The absolute scale stays undeclared, as every scale in this part does, and the internal ratios stay derived. On the carrier the line occupies one ray and its complement is empty: not heavy, not suppressed, empty, because a dissolution-first history is outside the transition domains and unregistered content supports no accessible pole at any rung. The image orientation the vertex selects is the mirror of that same single-arrow locus, one orbit of size two carrying no datum of its own. Four separate bars then close the door on any seesaw-type mechanism — no partner state, no self-pairing, no coupling object, no second scale — and the smallness of the line is left where the corpus has always left it, with the free internal data. The two states of the line are the vacuum-resident balanced pair and the materialized seal, and the only door between them is the vertex, whose weight vanishes exactly at the sharp threshold: the line materializes only where the flicker is present, which is the same threshold the ensemble integral is taken over.

\end{theorem}

\begin{theorem}[{\texorpdfstring{BM.97 — The rest mass is the pole mass: the species adjudication, the rest-mass Q-value, and the exact critical locus}{BM.97 - The rest mass is the pole mass: the species adjudication, the rest-mass Q-value, and the exact critical locus}}]\label{tag:BM-97}
\textbf{Thm.}

%dsp-level:0
\[
(1)\quad\textbf{the species census at the rest slot (three candidates, one read specified):}\qquad\begin{array}{l|l|l}\text{candidate} & \text{expression} & \text{status}\\\hline\text{the informational read} & m_{\mathrm{hid}}=E_{\mathrm{int}}\ln2/\mathfrak h & \text{the ledger's bit currency, ensemble-distributed}\ [\mathrm{FT.44}(4);\ \mathrm{BM.35}(4);\ \mathrm{BM.94}(8)]\\\text{the sector-deficit route} & m^{(\mathrm{ii})}=T(q,S)/\mathfrak h,\quad T=W(Q)\,\overline{\Delta S}(S) & \text{a measure price, vertex-scoped}\ [\mathrm{BM.28}(4)(\mathrm{ii});\ \mathrm{BM.30}(5);\ \mathrm{BM.37}(6);\ \mathrm{BM.77}(2)]\\\text{the pole/transport read} & m^2_{\mathrm{pole}}=\hat\mu^2\ell_c^2G(c)=4\hat\mu'^2\sin^2(\Theta_R/4) & \textbf{the rest mass on the transport species}\ [\mathrm{Def\text{-}H.37}(1),(2);\ \mathrm{FT.54}(4);\ \mathrm{FT.47}(1);\ \mathrm{EW.33};\ \mathrm{BM.39}(1),(7),(9)]\end{array}
\]

%dsp-level:0
\[
\textbf{the matching of the two pole clauses is derived at leading order and recorded,}\ \neg\textbf{imposed:}
\]

%dsp-level:0
\[
\qquad \hat\mu'^2\Theta_R^2/4=\hat\mu^2\ell_c^2G(c)
\]

%dsp-level:0
\[
\qquad [\mathrm{FT.54}(2);\ \mathrm{FT.33};\ \mathrm{BM.39}(9);\ \mathrm{Def\text{-}H.37}\ \text{locus clause}];
\]

%dsp-level:0
\[
\boxed{\ \textbf{the rest mass is the pole mass}\ m^{\mathrm{rest}}=m_{\mathrm{pole}}\ \ \text{on the transport species};\qquad m_{\mathrm{hid}}\ \ \text{is}\ \neg\textbf{the rest mass on that species}\ }
\]

%dsp-level:0
\[
\text{equating them is a conversion the corpus bars}
\]

%dsp-level:0
\[
\qquad [\mathrm{BM.35}(5);\ \mathrm{BM.83}(3);\ \mathrm{BM.37}(2);\ \text{pointer}];
\]

%dsp-level:0
\[
\qquad \text{the specified read A-BM.8 reads the baryon classes on route}\ (\mathrm{ii})\ \text{and the emitted lines on the transport species}\ [\mathrm{BM.77}(1);\ \text{re-cited},\ \neg\text{re-issued}];
\]

%dsp-level:0
\[
(2)\quad\textbf{the rest-mass Q-value, formed:}
\]

%dsp-level:0
\[
\qquad Q_\beta^{\mathrm{rest}}\ :=\ m_n^{\mathrm{rest}}-m_p^{\mathrm{rest}}-m_\ell^{\mathrm{rest}}-m_\nu^{\mathrm{rest}}
\]

%dsp-level:0
\[
\qquad [\mathrm{BM.41}(4);\ \mathrm{BM.63}(1);\ \mathrm{BM.85}(1);\ \text{every member a transport-species mass}]:
\]

%dsp-level:0
\[
\boxed{\ Q_\beta^{\mathrm{spec}}\,\mathfrak h\ =\ \Delta T(S)\ -\ \rho_\ell\sqrt{G\bigl(c(0)\bigr)}\ -\ \rho_\nu\sqrt{G\bigl(c(i_\nu)\bigr)},\qquad\Delta T(S)=5\,m^2_B(\cdot,S)\ }
\]

%dsp-level:0
\[
\qquad [\mathrm{BM.77}(2);\ \text{the expression is its own frozen}\ Q_{\mathrm{ii}}]:
\]

%dsp-level:0
\[
\Delta T=\Bigl(\frac{4219785}{98956},\ \ \frac{9608310}{393679},\ \ \frac{455}{23},\ \ \frac{68565}{3088}\Bigr)\ \ \text{at}\ S=(0,1,2,3),
\]

%dsp-level:0
\[
\qquad \frac{455}{23}\le\Delta T(S)\le\frac{4219785}{98956}
\]

%dsp-level:0
\[
\qquad [\mathrm{BM.77}(2);\ \text{machine re-certified}];
\]

%dsp-level:0
\[
G_\ell=A_0-B_0\cos\delta_\ell\ \ \text{at}\ \ \delta_\ell=\tfrac29
\]

%dsp-level:0
\[
\qquad [\mathrm{FT.112};\ \mathrm{Def\text{-}BM.6}(1)],
\]

%dsp-level:0
\[
\qquad G_\nu=A_0-B_0\cos\Bigl(\tfrac18-\frac{2\pi i_\nu}3\Bigr)\ \ \text{at the derived}\ \ \delta_\nu=\tfrac18
\]

%dsp-level:0
\[
\qquad [\mathrm{BM.96}(2);\ \mathrm{FT.142}(5)];
\]

%dsp-level:0
\[
\rho_\ell=\mathfrak h\,\hat\mu_\ell\,\ell_c^{(\ell)},\quad\rho_\nu=\mathfrak h\,\hat\mu_\nu\,\ell_c^{(\nu)}\in\mathfrak q^0
\]

%dsp-level:0
\[
\qquad [\mathrm{BM.77}(2);\ \mathrm{CT2.1};\ \mathrm{Def\text{-}H.37}(1)];
\]

%dsp-level:0
\[
\qquad \textbf{every row below holds uniformly in}\ i_\nu\in\{0,1,2\}
\]

%dsp-level:0
\[
\qquad [\mathrm{BM.77}(2)];
\]

%dsp-level:0
\[
(3)\quad\textbf{the flicker-blindness of the rest-mass read:}
\]

%dsp-level:0
\[
\qquad \text{the specified read's alphabet is}\ \{m^2_B(\cdot,S),\ \rho_\ell,\ \rho_\nu,\ g,\ g',\ \mathfrak h\}
\]

%dsp-level:0
\[
\qquad [\mathrm{BM.92}(5);\ \text{machine: the token census}]:
\]

%dsp-level:0
\[
\boxed{\ \neg E_{\mathrm{int}},\quad\neg w,\quad\neg w_{\mathrm{mat}},\quad\neg\varepsilon_q,\quad\neg\mathsf g,\quad\neg N\ \ \text{occur in it}\ }
\]

%dsp-level:0
\[
\qquad [\text{machine}]\implies\textbf{the flicker integral changes no row of the rest-mass Q-value}
\]

%dsp-level:0
\[
\textbf{the two acts are orthogonal:}
\]

%dsp-level:0
\[
\qquad \mathrm{BM.94}\text{'s fix lands on the informational register alone},
\]

%dsp-level:0
\[
\qquad \mathrm{BM.77}\text{'s read lands on the transport species alone}
\]

%dsp-level:0
\[
\qquad [\mathrm{BM.77}(3);\ \mathrm{BM.92}(3)\ \text{pattern}]:
\]

%dsp-level:0
\[
\implies\text{the trichotomy, the four}\ \Delta T\ \text{values,}\ \rho^{\mathrm{crit}}\text{'s enclosure and all eight witnesses stand}
\]

%dsp-level:0
\[
\qquad [\mathrm{BM.77}(5),(6);\ \mathrm{BM.92}(3);\ \text{machine re-certified}];
\]

%dsp-level:0
\[
(4)\quad\textbf{the pure-pole read (the third formed read):}
\]

%dsp-level:0
\[
\qquad Q_\beta^{\mathrm{pole}}\ :=\ m_n^{\mathrm{pole}}-m_p^{\mathrm{pole}}-m_\ell-m_\nu,
\]

%dsp-level:0
\[
\qquad M_0:=m_p^{\mathrm{pole}},
\]

%dsp-level:0
\[
\qquad (m_n^{\mathrm{pole}})^2=M_0^2+\Delta_g,
\]

%dsp-level:0
\[
\qquad s:=m_\ell+m_\nu>0
\]

%dsp-level:0
\[
\textbf{hypothesis}\ (\mathrm H):\ \text{the two baryon classes differ by exactly one branch substitution}\ [\mathrm{BM.59}(9);\ \mathrm{BM.69}(3);\ \text{carried on its face}]:
\]

%dsp-level:0
\[
\Delta_g\ \in\ \bigl\{2,\ \ \sqrt2-1,\ \ 2\cos\tfrac\pi{10}-\sqrt3\bigr\}\,\hat\mu'^2\ >\ 0,
\]

%dsp-level:0
\[
\qquad \Delta_1>\Delta_2>\Delta_3
\]

%dsp-level:0
\[
\qquad [\mathrm{Def\text{-}BM.14}(5);\ \mathrm{BM.69}(3);\ \text{machine: the certificates}\ 4>2,\ 2>1,\ 3<5];
\]

%dsp-level:0
\[
\textbf{the exact critical locus:}
\]

%dsp-level:0
\[
\qquad \boxed{\ Q_\beta^{\mathrm{pole}}>0\iff\Delta_g>\Delta_g^{\mathrm{crit}}:=s\,(2M_0+s);\qquad Q_\beta^{\mathrm{pole}}=0\iff\Delta_g=\Delta_g^{\mathrm{crit}};\qquad Q_\beta^{\mathrm{pole}}<0\iff\Delta_g<\Delta_g^{\mathrm{crit}}\ }
\]

%dsp-level:0
\[
[\text{machine: cleared by squaring two positive sides},\ \sqrt{M_0^2+\Delta_g}>M_0+s\iff M_0^2+\Delta_g>(M_0+s)^2];
\]

%dsp-level:0
\[
\textbf{the equivalent line form:}
\]

%dsp-level:0
\[
\qquad s^{\mathrm{crit}}=\sqrt{M_0^2+\Delta_g}-M_0=\frac{\Delta_g}{\sqrt{M_0^2+\Delta_g}+M_0}=m_n^{\mathrm{pole}}-m_p^{\mathrm{pole}}
\]

%dsp-level:0
\[
\qquad [\text{machine: the cleared identity}];
\]

%dsp-level:0
\[
\Delta_g^{\mathrm{crit}}\ \ \text{is strictly increasing in}\ s\ \ \text{and in}\ M_0,
\]

%dsp-level:0
\[
\qquad \Delta_g^{\mathrm{crit}}>0\ \ \text{strictly since}\ \ s>0
\]

%dsp-level:0
\[
\qquad [\mathrm{BM.96}(3);\ \text{machine: the coefficient census}];
\]

%dsp-level:0
\[
(5)\quad\textbf{the trichotomy (total, exact, non-empty on both sides):}
\]

%dsp-level:0
\[
(\mathrm a)\quad\mathcal O^{\mathrm{pole}}:=\bigl\{(\rho_\ell,\rho_\nu)>0:\ s<\Delta_g^{\mathrm{crit}}\text{'s line form}\bigr\}\ne\varnothing:
\]

%dsp-level:0
\[
\qquad \textbf{the universal symbolic witness}
\]

%dsp-level:0
\[
\rho_\ell=\rho_\nu=\frac{\sqrt{M_0^2+\Delta_g}-M_0}{4\sqrt{A_0+B_0}}\ \ \in\ \mathcal O^{\mathrm{pole}}
\]

%dsp-level:0
\[
\qquad [\mathrm{BM.77}(4)\text{'s sandwich}:\ s\le\tfrac12\bigl(\sqrt{M_0^2+\Delta_g}-M_0\bigr)<\sqrt{M_0^2+\Delta_g}-M_0;\ \text{machine}]:
\]

%dsp-level:0
\[
\text{valid at every generation, every line index}\ i_\nu,\ \text{every}\ (g,g')\ne(0,0);
\]

%dsp-level:0
\[
\qquad \textbf{the integer-certified rational witness at}\ (g,g',\rho_\ell,\rho_\nu,\hat\mu'^2\mathfrak h^2,M_0\mathfrak h)
\]

%dsp-level:0
\[
= (1,1,\tfrac1{100},\tfrac1{100},1,1),
\]

%dsp-level:0
\[
\ \text{generation}\ 1:
\]

%dsp-level:0
\[
A_0+B_0=\frac7{16}+\frac14=\frac{11}{16},
\]

%dsp-level:0
\[
\qquad \sqrt G\le\frac{\sqrt{11}}4<1
\]

%dsp-level:0
\[
\qquad [11<16]\implies s\,\mathfrak h<\frac1{50}\implies\Delta_g^{\mathrm{crit}}\mathfrak h^2<\frac1{50}\cdot\frac{101}{50}=\frac{101}{2500}<2=\Delta_1\mathfrak h^2
\]

%dsp-level:0
\[
\boxed{\ Q_\beta^{\mathrm{pole}}>0\ \ \text{on the certificate}\ \ 5000>101\ }\qquad[\text{machine}];
\]

%dsp-level:0
\[
(\mathrm b)\quad\mathcal C^{\mathrm{pole}}\ne\varnothing:
\]

%dsp-level:0
\[
\qquad \textbf{the universal symbolic witness}
\]

%dsp-level:0
\[
\qquad \rho_\ell=\rho_\nu=\frac{\sqrt{M_0^2+\Delta_g}-M_0}{\sqrt{A_0-B_0}}\ \ \in\ \mathcal C^{\mathrm{pole}}
\]

%dsp-level:0
\[
\qquad [\text{machine:}\ s\ge2\bigl(\sqrt{M_0^2+\Delta_g}-M_0\bigr)];
\]

%dsp-level:0
\[
\textbf{the integer-certified witness at}\ (1,1,100,100,1,1):
\]

%dsp-level:0
\[
\qquad A_0-B_0=\frac3{16},
\]

%dsp-level:0
\[
\qquad \sqrt G\ge\frac{\sqrt3}4>\frac25\qquad \qquad [75>64]:
\]

%dsp-level:0
\[
m_\ell\,\mathfrak h\ \ge\ 25\sqrt3\ >\ 40\qquad[1875>1600]\implies s\,\mathfrak h\ \ge\ 50\sqrt3\ >\ 80\qquad[75>64]\implies\Delta_g^{\mathrm{crit}}\mathfrak h^2\ >\ 80\cdot82=6560\ >\ 2
\]

%dsp-level:0
\[
\boxed{\ Q_\beta^{\mathrm{pole}}<0\ \ \text{on the certificates}\ \ 1875>1600,\quad75>64,\quad6560>2\ }\qquad[\text{machine}];
\]

%dsp-level:0
\[
(\mathrm c)\quad\textbf{the critical locus is exact, unique and generation-relative:}
\]

%dsp-level:0
\[
\qquad \partial\mathcal O^{\mathrm{pole}}:\quad s=\sqrt{M_0^2+\Delta_g}-M_0
\]

%dsp-level:0
\[
\qquad [\text{machine: the boundary equation}];
\]

%dsp-level:0
\[
\Delta_g^{\mathrm{crit}}\ \ \text{is generation-blind}\ (\text{a function of the emitted lines and}\ M_0\ \text{alone})\ \ \text{while}\ \ \Delta_g\ \ \text{is generation-specific}\qquad[(4)]:
\]

%dsp-level:0
\[
\boxed{\ \textbf{the opening region is widest in the first generation and narrowest in the third, in the exact order}\ \ \Delta_1>\Delta_2>\Delta_3\ }
\]

%dsp-level:0
\[
\qquad [\mathrm{BM.69}(3);\ \text{machine: monotone in}\ \Delta_g\ \text{at fixed}\ M_0];
\]

%dsp-level:0
\[
\textbf{on the symmetric ray}\ \rho_\ell
\]

%dsp-level:0
\[
= \rho_\nu
\]

%dsp-level:0
\[
= \rho\ \textbf{at the witness scales}\ (g,g',M_0\mathfrak h,\hat\mu'^2\mathfrak h^2)
\]

%dsp-level:0
\[
= (1,1,1,1),
\]

%dsp-level:0
\[
\ \text{generation}\ 1:
\]

%dsp-level:0
\[
\qquad \rho^{\mathrm{crit,pole}}=\frac{\bigl(\sqrt3-1\bigr)}{2\sqrt G}:
\]

%dsp-level:0
\[
\boxed{\ \frac25\ <\ \rho^{\mathrm{crit,pole}}\ <\ \frac67\ }
\]

%dsp-level:0
\[
\qquad [\text{machine:}\ \sqrt G\le\tfrac{\sqrt{11}}4\ \text{gives the lower bound on}\ 2(\sqrt3-1)/\sqrt{11}>\tfrac25\iff89>50\sqrt3\iff7921>7500;\quad\sqrt G\ge\tfrac{\sqrt3}4\ \text{gives the upper bound on}\ 2(\sqrt3-1)/\sqrt3<\tfrac67\iff8\sqrt3<14\iff192<196];
\]

%dsp-level:0
\[
(6)\quad\textbf{the}\ \nu\textbf{-positivity shift (a derived consequence of}\ \mathrm{BM.96}(3)\textbf{):}
\]

%dsp-level:0
\[
\qquad \Delta_g^{\mathrm{crit}}(m_\ell{+}m_\nu)-\Delta_g^{\mathrm{crit}}(m_\ell)=2M_0m_\nu+2m_\ell m_\nu+m_\nu^2\ >\ 0
\]

%dsp-level:0
\[
\qquad [\text{machine: the exact expansion}]:
\]

%dsp-level:0
\[
\boxed{\ \textbf{the derived positivity of the}\ \nu\ \textbf{rest mass strictly narrows the opening region by an exact amount}\ —\ \text{a massless}\ \nu\ \text{line is refuted as a read,}\ \neg\text{assumed}\ }
\]

%dsp-level:0
\[
\textbf{the two-body companion on this read:}
\]

%dsp-level:0
\[
\qquad E_p-m_p=\frac{(m_n-m_p)^2-m_\ell^2}{2m_n}
\]

%dsp-level:0
\[
\qquad [\mathrm{BM.41}(2);\ \mathrm{BM.77}(6);\ \text{machine: the identity re-derived at}\ m_n\ne m_p]:
\]

%dsp-level:0
\[
\text{the two-body channel is admissible}\iff(m_n^{\mathrm{pole}}-m_p^{\mathrm{pole}})^2\ge m_\ell^2\iff\Delta_g\ \ge\ m_\ell\bigl(\sqrt{M_0^2+\Delta_g}+M_0\bigr)
\]

%dsp-level:0
\[
\qquad [\text{machine: the cleared form}];
\]

%dsp-level:0
\[
\textbf{the forced-third-member row is read-relative:}
\]

%dsp-level:0
\[
\qquad \mathrm{BM.41}(2)\text{'s equal-mass obstruction is lifted on the pure-pole read}\ (\Delta_g>0),
\]

%dsp-level:0
\[
\ \text{so}\ \mathrm{BM.41}(3),
\]

%dsp-level:0
\[
(6)\text{'s forcing stands on the informational reads}\ [\mathrm{BM.70}(4)(\mathrm a)]\ \text{and the grammar stands on the specified read}\ [\mathrm{BM.45};\ \mathrm{BM.77}(6)],
\]

%dsp-level:0
\[
\text{and it is}\ \neg\textbf{inherited by the pure-pole read}\ —\ \text{recorded,}\ \neg\text{adjudicated beyond the branch rows}\qquad[\mathrm{BM.77}(6)\ \text{pattern}];
\]

%dsp-level:0
\[
(7)\quad\textbf{the read discipline:}\qquad\textbf{three formed reads, pairwise incomparable:}
\]

%dsp-level:0
\[
Q_\beta^{\mathrm{info}}:\ \{\ln2,\ln3,\ln5,\ N,\ w\ \text{or}\ w_{\mathrm{mat}},\ \mathsf g,\ \varepsilon_{\mathrm{EM}},\ \mathfrak h,\ m_\ell,\ m_\nu\};
\]

%dsp-level:0
\[
\qquad Q_\beta^{\mathrm{spec}}:\ \{m^2_B(\cdot,S),\ \rho_\ell,\ \rho_\nu,\ g,\ g',\ \mathfrak h\};
\]

%dsp-level:0
\[
\qquad Q_\beta^{\mathrm{pole}}:\ \{M_0,\ \Delta_g,\ \hat\mu'^2,\ \rho_\ell,\ \rho_\nu,\ g,\ g',\ \mathfrak h\}
\]

%dsp-level:0
\[
\boxed{\ \text{the informational alphabet is disjoint from the two transport alphabets except for}\ \mathfrak h;\qquad\neg\text{a derived equation relates the three expressions}\ }
\]

%dsp-level:0
\[
\qquad [\text{machine: the three token censuses};\ \text{the mass-form bridge residue, localized at}\ \mathrm{BM.98}(5)];
\]

%dsp-level:0
\[
\textbf{the decay direction is fixed in all three:}
\]

%dsp-level:0
\[
\qquad T(0,S)=\tfrac43T(+1,S)\ [\mathrm{BM.28}(4)(\mathrm{ii})],
\]

%dsp-level:0
\[
\qquad \Delta_g>0\ [\mathrm{BM.69}(3)],
\]

%dsp-level:0
\[
\qquad \Delta m>0\ \ \text{sub-critically}\ [\mathrm{BM.91}(6);\ \mathrm{BM.99}(4)]:
\]

%dsp-level:0
\[
\boxed{\ \textbf{the parent is the neutral class on every formed read, and the grammar}\ n\to p+\ell+\nu\ \ \textbf{stands unmodified}\ }
\]

%dsp-level:0
\[
\qquad [\mathrm{BM.41}(1);\ \mathrm{BM.45};\ \mathrm{BM.77}(6)];
\]

%dsp-level:0
\[
\neg\text{a currency is converted}\ [\mathrm{BM.83}(3);\ \mathrm{BM.90}(5)];
\]

%dsp-level:0
\[
\qquad \text{route}\ (\mathrm{ii})\ \text{is}\ \neg\text{elevated beyond the vertex kinematics}\ [\mathrm{BM.30}(5);\ \mathrm{BM.77}(9)];
\]

%dsp-level:0
\[
\qquad \text{route}\ (\mathrm{iii})\text{'s amplitude ratios stay barred as masses}\ [\mathrm{FT.79}(5);\ \mathrm{BM.37}(3)];
\]

%dsp-level:0
\[
(8)\quad\textbf{status:}
\]

%dsp-level:0
\[
\qquad \textbf{derived:}\ (1)\text{--}(7)\ \text{given the specified read and hypothesis}\ (\mathrm H);
\]

%dsp-level:0
\[
\qquad \textbf{declared: nothing}\ —\ \text{A-BM.8 and}\ \mathrm{Def\text{-}H.37}\ \text{are re-cited,}\ \neg\text{re-issued};
\]

%dsp-level:0
\[
\neg\text{a clause of}\ \mathrm{BM.41},\ \mathrm{BM.49},\ \mathrm{BM.63},\ \mathrm{BM.70},\ \mathrm{BM.71},\ \mathrm{BM.77},\ \mathrm{BM.85},\ \mathrm{BM.92}\ \text{is retracted};
\]

%dsp-level:0
\[
\qquad \neg\text{a P4 addition},\quad\neg\text{a new dimensionful scale}\ [\mathrm{P1}],\quad\neg\text{a measured value}\ [\mathrm{P5}],\quad\text{§Cal-BM stays empty};
\]

\textbf{Req.} D2.9 (Parts 1–2); CT2.1, EM.25(7), EW.15 (Part 2); P1, P4, P5 (Parts 1–2); Def-H.37(1), Def-H.37(2), EW.33, FT.33, FT.47(1), FT.54(2), FT.54(4), FT.79(5), FT.112, FT.142(5) (Part 3); Def-BM.6(1), Def-BM.11(5), Def-BM.14(5), BM.28(4), BM.30(5), BM.35(4), BM.35(5), BM.37(2), BM.37(3), BM.37(6), BM.39(1), BM.39(7), BM.39(9), BM.41(1), BM.41(2), BM.41(3), BM.41(4), BM.41(6), BM.45, BM.59(9), BM.63(1), BM.69(3), BM.70(4), BM.77(1), BM.77(2), BM.77(3), BM.77(4), BM.77(5), BM.77(6), BM.77(9), BM.83(3), BM.85(1), BM.87, BM.90(5), BM.91(6), BM.92(3), BM.92(5), BM.96(2), BM.96(3), BM.96(4) (this part); Cal-F; exact radical inequalities by integer certificates, sum-of-squares positivity, cleared polynomial identities over $\mathbb Q[M_0,m_\ell,m_\nu]$, exact rational arithmetic at the four frozen sector expectations, three token censuses [math; machine: \texttt{checks/\allowbreak{}verify\_\allowbreak{}H\_\allowbreak{}flicker\_\allowbreak{}integral\_\allowbreak{}restmass.py} §D].

\textbf{Pf.}

%dsp-level:0
\[
(1):\ \mathrm{Def\text{-}H.37}(1),(2)\text{'s two clauses and}\ \mathrm{FT.54}(2),(4)\text{'s matching re-cited};
\]

%dsp-level:0
\[
\qquad \text{the three-candidate census assembled from}\ \mathrm{FT.44}(4),
\]

%dsp-level:0
\[
\ \mathrm{BM.28}(4)(\mathrm{ii}),
\]

%dsp-level:0
\[
\ \mathrm{BM.39}(1),
\]

%dsp-level:0
\[
(7)\ [\text{machine: the read-class rows}];
\]

%dsp-level:0
\[
\qquad \text{the conversion bar re-cited}\ [\mathrm{BM.35}(5);\ \mathrm{BM.83}(3)]
\]

%dsp-level:0
\[
(2):\ \mathrm{BM.77}(2)\text{'s frozen expression re-cited with}\ G_\nu\ \text{evaluated at the derived}\ \delta_\nu
\]

%dsp-level:0
\[
= \tfrac18\ [\mathrm{BM.96}(2)];
\]

%dsp-level:0
\[
\qquad 5\cdot843957=4219785,
\]

%dsp-level:0
\[
\quad 5\cdot1921662=9608310,
\]

%dsp-level:0
\[
\quad 5\cdot91=455,
\]

%dsp-level:0
\[
\quad 5\cdot13713=68565\ [\text{machine: exact integers}];
\]

%dsp-level:0
\[
\qquad \tfrac{455}{23}\le\Delta T\le\tfrac{4219785}{98956}\ [\text{machine: cross-multiplication}]
\]

%dsp-level:0
\[
(3):\ \text{the alphabet enumerated from the expression and scanned for the six barred tokens}\ [\text{machine}];
\]

%dsp-level:0
\[
\qquad \text{the eight witnesses re-certified as integer comparisons}\ [\text{machine}]
\]

%dsp-level:0
\[
(4):\ \sqrt{M_0^2+\Delta_g}>M_0+s\ \text{with both sides positive}
\]

%dsp-level:0
\[
\iff M_0^2+\Delta_g>M_0^2+2M_0s+s^2\iff\Delta_g>2M_0s+s^2\ [\text{machine: cleared over}\ \mathbb Q[M_0,s,\Delta_g]];
\]

%dsp-level:0
\[
s^{\mathrm{crit}}
\]

%dsp-level:0
\[
= \sqrt{M_0^2+\Delta_g}-M_0\ \text{and}\ (\sqrt{M_0^2+\Delta_g}-M_0)(\sqrt{M_0^2+\Delta_g}+M_0)
\]

%dsp-level:0
\[
= \Delta_g\ [\text{machine: the difference of squares}];
\]

%dsp-level:0
\[
\qquad \Delta_1>\Delta_2>\Delta_3\ \text{on}\ 4>2,\ 2>1,\ 3<5\ [\mathrm{BM.69}(3)\ \text{re-cited}]
\]

%dsp-level:0
\[
(5)(\mathrm a):\ 2\cdot\tfrac{\sqrt{M_0^2+\Delta_g}-M_0}{4\sqrt{A_0+B_0}}\cdot\sqrt{A_0+B_0}
\]

%dsp-level:0
\[
= \tfrac12\bigl(\sqrt{M_0^2+\Delta_g}-M_0\bigr)<\sqrt{M_0^2+\Delta_g}-M_0\ [\text{machine}];
\]

%dsp-level:0
\[
\qquad A_0+B_0\big|_{g=g'=1}=\tfrac{11}{16}\ [\text{machine}]
\]

%dsp-level:0
\[
s\,\mathfrak h\le2\cdot\tfrac1{100}\cdot\tfrac{\sqrt{11}}4=\tfrac{\sqrt{11}}{200}<\tfrac4{200}=\tfrac1{50}\ [11<16;\ \text{machine}];
\]

%dsp-level:0
\[
\qquad \tfrac1{50}\bigl(2+\tfrac1{50}\bigr)=\tfrac{101}{2500};
\]

%dsp-level:0
\[
\qquad 2>\tfrac{101}{2500}\iff5000>101\ [\text{machine}]
\]

%dsp-level:0
\[
(\mathrm b):\ 2\cdot\tfrac{\sqrt{M_0^2+\Delta_g}-M_0}{\sqrt{A_0-B_0}}\cdot\sqrt{A_0-B_0}=2\bigl(\sqrt{M_0^2+\Delta_g}-M_0\bigr)>\sqrt{M_0^2+\Delta_g}-M_0\ [\text{machine}];
\]

%dsp-level:0
\[
\qquad A_0-B_0\big|_{g=g'=1}=\tfrac3{16};
\]

%dsp-level:0
\[
\qquad \sqrt G\ge\tfrac{\sqrt3}4>\tfrac25\iff75>64\ [\text{machine}]
\]

%dsp-level:0
\[
m_\ell\mathfrak h\ge100\cdot\tfrac{\sqrt3}4=25\sqrt3>40\iff1875>1600;
\]

%dsp-level:0
\[
\qquad s\,\mathfrak h\ge50\sqrt3>80\iff75>64;
\]

%dsp-level:0
\[
\qquad 80(2+80)=6560>2\ [\text{machine: integers}]
\]

%dsp-level:0
\[
(\mathrm c):\ \rho^{\mathrm{crit,pole}}=\tfrac{\sqrt3-1}{2\sqrt G}\ \text{at}\ M_0\mathfrak h=1,\ \Delta_1\mathfrak h^2=2\ [\text{machine:}\ \sqrt{1+2}=\sqrt3];
\]

%dsp-level:0
\[
\qquad \text{the lower bound:}\ \tfrac{2(\sqrt3-1)}{\sqrt{11}}>\tfrac25
\]

%dsp-level:0
\[
\iff 5(\sqrt3-1)>\sqrt{11}
\]

%dsp-level:0
\[
\iff 25(4-2\sqrt3)>11
\]

%dsp-level:0
\[
\iff 89>50\sqrt3
\]

%dsp-level:0
\[
\iff 7921>7500\ [\text{machine}]
\]

%dsp-level:0
\[
\text{the upper bound:}\ \tfrac{2(\sqrt3-1)}{\sqrt3}<\tfrac67\iff14(\sqrt3-1)<6\sqrt3\iff8\sqrt3<14\iff192<196\ [\text{machine}];
\]

%dsp-level:0
\[
\qquad \text{monotonicity in}\ \Delta_g\ \text{at fixed}\ M_0:\ \text{the coefficient of}\ \Delta_g\ \text{in}\ \sqrt{M_0^2+\Delta_g}-M_0\ \text{is positive}\ [\text{machine}]
\]

%dsp-level:0
\[
(6):\ (m_\ell+m_\nu)(2M_0+m_\ell+m_\nu)-m_\ell(2M_0+m_\ell)
\]

%dsp-level:0
\[
= 2M_0m_\nu+2m_\ell m_\nu+m_\nu^2\ [\text{machine: the expansion over}\ \mathbb Q[M_0,m_\ell,m_\nu]];
\]

%dsp-level:0
\[
2m_n(E_p-m_p)
\]

%dsp-level:0
\[
= (m_n-m_p)^2-m_\ell^2\ \text{at}\ E_p
\]

%dsp-level:0
\[
= \tfrac{m_n^2+m_p^2-m_\ell^2}{2m_n}\ [\text{machine: cleared}]\implies\text{the two-body condition}\ m_n-m_p
\]

%dsp-level:0
\[
\ge m_\ell\ \text{and its}\ \Delta_g\ \text{form}\ [\text{machine}]
\]

%dsp-level:0
\[
(7):\ \text{the three symbol sets enumerated from the three expressions and pairwise intersected}\ [\text{machine}];
\]

%dsp-level:0
\[
\qquad T(0,S)/T(+1,S)=20/15=4/3\ [\text{machine}];
\]

%dsp-level:0
\[
\qquad (8):\ \text{the standing-rule scans}\ [\text{machine}].\ \blacksquare
\]

\textbf{Phys.} A Q-value is a difference of rest masses, and a rest mass in this corpus is a pole mass: the transport defect of a materialized chain, exact as four mu-prime squared times the sine squared of a quarter frame angle, matched at leading order to the Gram evaluation the seed supplies. The informational count is a different currency, ensemble-distributed and bit-denominated, and equating the two is precisely the conversion the ledger bars — so the fix to the snapshot read and the read of the vertex turn out to be orthogonal acts, and the specified Q-value carries no edge count, no crystallization weight, no threshold and no curvature composition in its alphabet at all. On that read nothing moves: the four exact baryon values, the critical enclosure and all eight integer witnesses stand. The pure pole read is the third formed read, and it has its own exact trichotomy — a critical locus that is the emitted-line sum times twice the daughter's mass plus itself, a universal symbolic witness in each of the opening and closing regions, and rational witnesses certified by five thousand against one hundred one and by one thousand eight hundred seventy-five against one thousand six hundred. Because the neutrino line's mass is derived strictly positive, the critical locus sits strictly above its massless value by an exact amount, so the positivity narrows the opening region rather than being assumed away; and because the pole splitting is generation-specific while the critical locus is not, the opening region is widest in the first generation and narrowest in the third. What the three reads agree on is the direction: the neutral class is the parent in every currency, and the grammar is untouched.

\textbf{Rmk (the operative mass read — reconciled nomenclature).} The mass that works — the read entering every formed composite quantity of this era (the rest masses, the Q-values, the splittings, the thresholds, the lapse) — is the \textbf{Lyapunov-contraction mass}: the ledger deficit accumulated under the tilt contraction law [L1–L4 (Part 1): every tilt step loses $\gamma(8-\gamma)/16$ of the Lyapunov value; BM.51: the class-degenerate foundation], in its composite form $m_{\mathrm{rest}}=\Delta S_{\mathfrak L}/\mathfrak h$ [BM.35]. Def-H.39's registry (Part 3) names the species, and no unqualified \textit{mass} crosses them: the informational read agrees with $m_{\mathrm{rest}}$ as a formed equation [BM.35], $\neg$as a currency fusion; the pole read is the rest mass on the transport species, matched at leading order [(1) above]; the sector-route values are measure functionals — exact named prices, $\neg$masses [BM.30; BM.37]; the Ansatz-B class masses remain non-adopted templates [Def-H.28 (Part 3)]; and the numeric bridges (the $\beta$-coefficients; $\rho_{\mathrm B}$) are calibration slots, values $\neg$issued [§Cal-BM (C-BM.1); P5].

\end{theorem}

\begin{theorem}[{\texorpdfstring{BM.98 — The relating equation: the bridge in four exact forms, its flicker weighting, and the closure  in form}{BM.98 - The relating equation: the bridge in four exact forms, its flicker weighting, and the closure in form}}]\label{tag:BM-98}
\textbf{Thm.}

%dsp-level:0
\[
(1)\quad\textbf{the bridge in its derived count form:}
\]

%dsp-level:0
\[
\qquad \mathcal R\ :=\ \frac{m^2_{\mathrm{pole}}}{m^2_{\mathrm{hid}}}\ =\ \rho_{\mathrm B}\,\frac{D_S(\Theta_R)}{4\,E_{\mathrm{int}}^2},
\]

%dsp-level:0
\[
\qquad \rho_{\mathrm B}=\frac{\hat\mu'^2\mathfrak h^2}{(\ln2)^2}\in\mathfrak q^0
\]

%dsp-level:0
\[
\qquad [\mathrm{BM.69}(2);\ \mathrm{FT.54}(4);\ \mathrm{Census\ FT.63}(6)]:
\]

%dsp-level:0
\[
\text{six states, six exact algebraic factors}\ \tfrac{D_S}{4E^2},\ \text{four fields, six minimal polynomials}
\]

%dsp-level:0
\[
\qquad [\mathrm{Def\text{-}BM.14}(3);\ \mathrm{BM.82}(2);\ \text{machine re-cited}];
\]

%dsp-level:0
\[
(2)\quad\textbf{the flicker-weighted bridge:}
\]

%dsp-level:0
\[
\qquad E_{\mathrm{int}}\ \longrightarrow\ \bigl\langle E_{\mathrm{int}}\bigr\rangle_{\mathcal E}
\]

%dsp-level:0
\[
\qquad [\mathrm{BM.94}(2),(6),(7);\ \text{the integral replaces the snapshot}]:
\]

%dsp-level:0
\[
\textbf{unconditioned:}
\]

%dsp-level:0
\[
\qquad \mathcal R_w=\rho_{\mathrm B}\,\frac{D_S}{4\,w^2E^2}=\frac{\mathcal R}{w^2};
\]

%dsp-level:0
\[
\qquad \textbf{materialized:}
\]

%dsp-level:0
\[
\qquad \boxed{\ \mathcal R_{\mathrm{mat}}\ =\ \rho_{\mathrm B}\,\frac{D_S\,D^2}{4\,E^2}\ =\ \frac{\mathcal R}{w_{\mathrm{mat}}^2}\ =\ \rho_{\mathrm B}\,\frac{D_S\,D^2}{144}\ \ \text{at}\ \ E=E^{\mathrm{cryst}}=6\ }
\]

%dsp-level:0
\[
[\text{machine: one exact rational function per state}];
\]

%dsp-level:0
\[
\qquad \boxed{\ 1\ \le\ \frac{\mathcal R_{\mathrm{mat}}}{\mathcal R}\ =\ D^2\ <\ 9\ }
\]

%dsp-level:0
\[
\qquad [\mathrm{BM.94}(7):\ D\in[1,3);\ \text{machine}]:
\]

%dsp-level:0
\[
\boxed{\ \textbf{the flicker can multiply the bridge by at most an exact nine, and by exactly one at the sharp threshold}\ };
\]

%dsp-level:0
\[
\qquad \text{the lower endpoint is attained at}\ \varepsilon_q
\]

%dsp-level:0
\[
= 0,
\]

%dsp-level:0
\[
\ \text{where the snapshot is recovered exactly}\ [\mathrm{BM.94}(6),(7)];
\]

%dsp-level:0
\[
\text{the upper endpoint is open and is approached only in the fully dissolved limit, where}\ \bigl\langle E\bigr\rangle_{\mathrm{mat}}\to2\ —\ \text{the digon floor}\ [\mathrm{Census\ FT.28};\ \mathrm{BM.94}(7)];
\]

%dsp-level:0
\[
(3)\quad\textbf{the curvature form:}
\]

%dsp-level:0
\[
\qquad G=2\pi\,E_{\mathrm{int}}^2\ln2
\]

%dsp-level:0
\[
\qquad [\mathrm{BM.88}(3);\ \mathrm{BM.16}(1):\ \text{the calibration lock}]\implies E_{\mathrm{int}}^2=\frac{G}{2\pi\ln2}:
\]

%dsp-level:0
\[
\boxed{\ \mathcal R\ =\ \frac{\pi\ln2}{2}\cdot\frac{\rho_{\mathrm B}\,D_S}{G}\ }
\]

%dsp-level:0
\[
\qquad [\text{machine: one monomial, coefficient comparison}]
\]

%dsp-level:0
\[
\qquad \textbf{the bridge is inversely proportional to the curvature lock, with the exact coefficient}\ \tfrac{\pi\ln2}2;
\]

%dsp-level:0
\[
\textbf{flicker-weighted:}
\]

%dsp-level:0
\[
\qquad G^{\mathrm{eff}}:=2\pi\bigl\langle E\bigr\rangle^2_{\mathrm{mat}}\ln2=\frac{G}{D^2}=w_{\mathrm{mat}}^2\,G,
\]

%dsp-level:0
\[
\qquad \mathcal R_{\mathrm{mat}}=\frac{\pi\ln2}2\cdot\frac{\rho_{\mathrm B}D_S}{G^{\mathrm{eff}}}
\]

%dsp-level:0
\[
\qquad [\text{machine}];
\]

%dsp-level:0
\[
\textbf{the exact enclosure at}\ E^{\mathrm{cryst}}=6:
\]

%dsp-level:0
\[
\qquad G=72\,\pi\ln2
\]

%dsp-level:0
\[
\qquad [\mathrm{BM.17}(3)\text{'s lock value re-derived};\ \text{machine}]\implies\boxed{\ G^{\mathrm{eff}}\ \in\ \bigl(8\,\pi\ln2,\ \ 72\,\pi\ln2\bigr]\ }
\]

%dsp-level:0
\[
\textbf{the two endpoints are the digon's and}\ \Sigma_{12}\textbf{'s own curvature locks}
\]

%dsp-level:0
\[
\qquad [\mathrm{BM.17}(4)\text{'s family table};\ \text{machine:}\ 2\pi\cdot2^2\ln2=8\pi\ln2]:
\]

%dsp-level:0
\[
\qquad \text{the flicker moves the lock between two computed cells and never below the floor's};
\]

%dsp-level:0
\[
(4)\quad\textbf{the coupling form:}
\]

%dsp-level:0
\[
\qquad k=\frac{r\,E_{\mathrm{int}}\ln2}{\mathfrak h},
\]

%dsp-level:0
\[
\qquad 1-f=\frac1k
\]

%dsp-level:0
\[
\qquad [\mathrm{BM.88}(2)(\mathrm i),(\mathrm{ii});\ \mathrm{BM.88}(6)]\implies E_{\mathrm{int}}=\frac{\mathfrak h}{r\,(1-f)\ln2}:
\]

%dsp-level:0
\[
\boxed{\ \mathcal R\ =\ \frac{\hat\mu'^2\,D_S\,r^2}{4}\,(1-f)^2\ }
\]

%dsp-level:0
\[
\qquad [\text{machine: the substitution};\qquad[\hat\mu'^2r^2]=\mathfrak q\cdot\mathfrak q^{-1}=\mathfrak q^0\ [\mathrm{CT2.1};\ \mathrm{Def\text{-}H.37}(2);\ \text{machine: the exponent arithmetic}]];
\]

%dsp-level:0
\[
e^2=e_0^2\,N\,f\,\mathsf g(G)\qquad[\mathrm{BM.88}(4)]\implies1-f=1-\frac{e^2}{e_0^2N\,\mathsf g(G)}:
\]

%dsp-level:0
\[
\boxed{\ \mathcal R\ =\ \frac{\hat\mu'^2\,D_S\,r^2}{4}\Bigl(1-\frac{e^2}{e_0^2\,N\,\mathsf g(G)}\Bigr)^{\!2}\ }
\]

%dsp-level:0
\[
\qquad \textbf{the bridge as an exact function of the running coupling}
\]

%dsp-level:0
\[
\textbf{the energy-dependence of the bridge is the energy-dependence of the coupling:}
\]

%dsp-level:0
\[
\qquad \text{both run through the single derived argument}\ f;
\]

%dsp-level:0
\[
\mathcal R\ \ \text{strictly decreasing in}\ E_{\mathrm{int}},\quad\text{strictly decreasing in}\ f,\quad\text{strictly increasing in}\ k;
\]

%dsp-level:0
\[
\qquad e\ \ \text{strictly increasing in}\ E_{\mathrm{int}}
\]

%dsp-level:0
\[
\qquad [\mathrm{BM.88}(6);\ \text{machine: three monotonicity rows on exact rationals}]:
\]

%dsp-level:0
\[
\boxed{\ \text{at fixed}\ (r,\,N,\,\mathsf g,\,D_S),\ \ \mathcal R\ \ \textbf{is a strictly decreasing exact function of the running coupling}\ }
\]

%dsp-level:0
\[
\textbf{the flicker's action on the coupling:}
\]

%dsp-level:0
\[
\qquad k^{\mathrm{eff}}=w_{\mathrm{mat}}\,k^{\mathrm{cryst}}=\frac{k^{\mathrm{cryst}}}D,
\]

%dsp-level:0
\[
\qquad f^{\mathrm{eff}}=1-\frac{D}{k^{\mathrm{cryst}}}\ \le\ f^{\mathrm{cryst}},
\]

%dsp-level:0
\[
\qquad \mathsf g^{\mathrm{eff}}=\mathsf g\bigl(G^{\mathrm{eff}}\bigr)\ \ \text{at}\ \ G^{\mathrm{eff}}=\frac{G}{D^2}
\]

%dsp-level:0
\[
\qquad [\text{machine}]:
\]

%dsp-level:0
\[
\text{equality}\iff\varepsilon_q=0
\]

%dsp-level:0
\[
\qquad [\text{machine}]\implies\textbf{the flicker attenuates the running coupling and enhances the bridge, both by the same derived weight};
\]

%dsp-level:0
\[
\textbf{the vanishing locus is shifted and bounded:}
\]

%dsp-level:0
\[
\qquad f^{\mathrm{eff}}=0\iff r=D\,\Delta x_{\mathrm B}^{\mathrm{cryst}}\in\bigl[\Delta x_{\mathrm B}^{\mathrm{cryst}},\ 3\,\Delta x_{\mathrm B}^{\mathrm{cryst}}\bigr)
\]

%dsp-level:0
\[
\qquad [\mathrm{BM.88}(5);\ \mathrm{BM.94}(8);\ \text{machine}]:
\]

%dsp-level:0
\[
\boxed{\ \textbf{the coupling's zero is confined to within an exact factor of three of the crystallized Born width}\ }
\]

%dsp-level:0
\[
(5)\quad\textbf{the closure in form and the localized residue:}
\]

%dsp-level:0
\[
\qquad \textbf{every factor of}\ (1)\text{--}(4)\ \textbf{is derived exact algebra:}
\]

%dsp-level:0
\[
\text{the six}\ D_S\ \text{with their minimal polynomials}\ [\mathrm{Def\text{-}BM.14}(3)];
\]

%dsp-level:0
\[
\quad G\text{'s monomial and}\ f\text{'s rational function}\ [\mathrm{BM.88}(2),(3)];
\]

%dsp-level:0
\[
\quad N\text{'s ADM row}\ [\mathrm{BM.88}(2)(\mathrm{iii})];
\]

%dsp-level:0
\[
\quad w_{\mathrm{mat}},\ D,\ \text{the integer}\ 9\ [\mathrm{BM.94}(7)];
\]

%dsp-level:0
\[
\text{the coefficient}\ \tfrac{\pi\ln2}2\ (3);
\]

%dsp-level:0
\[
\qquad \text{the two computed lock endpoints}\ (3);
\]

%dsp-level:0
\[
\qquad \text{every monotonicity row and certificate}\ (2)\text{--}(4)
\]

%dsp-level:0
\[
\qquad [\text{machine}];
\]

%dsp-level:0
\[
\textbf{and one factor is not:}
\]

%dsp-level:0
\[
\qquad \boxed{\ \rho_{\mathrm B}=\frac{\hat\mu'^2\mathfrak h^2}{(\ln2)^2}\ \ \text{is the single undeclared dimensionless ratio}\ \ (\mathrm{R1})\ }
\]

%dsp-level:0
\[
\qquad [\mathrm{BM.69}(5);\ \mathrm{BM.82}(6);\ \mathrm{Cal\text{-}F}\ \text{packaging},\ \neg\text{added here}]:
\]

%dsp-level:0
\[
\boxed{\ \textbf{the relating equation is derived as an exact identity in the derived vocabulary plus}\ (\mathrm{R1}):\ \ \textbf{the equation is closed, the ratio is}\ \neg\textbf{derived}\ }
\]

%dsp-level:0
\[
\mathrm{BM.69}(4)\textbf{'s four grounds stand unweakened:}
\]

%dsp-level:0
\[
\qquad (\mathrm i)\ \text{the degree-}16\ \text{field against integer- and binary-valued reads}\ [\mathrm{FT.73}(2),(3)];
\]

%dsp-level:0
\[
\qquad (\mathrm{ii})\ \text{the two-torsion image reaching}\ 2\ \text{of}\ 6\ \text{states}\ [\mathrm{BM.67}(6)];
\]

%dsp-level:0
\[
(\mathrm{iii})\ \text{the calibration weight}\ {-}2\ \text{against}\ 0\ [\mathrm{FT.55}(1);\ \mathrm{BM.67}(8)];
\]

%dsp-level:0
\[
\qquad (\mathrm{iv})\ \text{the strength bar}\ [\mathrm{BM.36}(3);\ \mathrm{BM.61}(8)(\mathrm{O1})]
\]

%dsp-level:0
\[
\qquad [\text{machine: the four rows re-cited at their own loci}];
\]

%dsp-level:0
\[
\textbf{the residue list after this tier:}
\]

%dsp-level:0
\[
\qquad (\mathrm{R1})\ \rho_{\mathrm B};
\]

%dsp-level:0
\[
\quad (\mathrm{R2})\ \text{the loop realization};
\]

%dsp-level:0
\[
\quad (\mathrm{R3})\ \text{the flavour composition};
\]

%dsp-level:0
\[
\quad (\mathrm{R4})\ w(\varepsilon_q)\ \text{in general and the transport register's ensemble average};
\]

%dsp-level:0
\[
\quad (\mathrm{R5})\ \iota;
\]

%dsp-level:0
\[
(\mathrm{R6})\ \ \textbf{the reciprocal register's Jensen residue}\ \ \langle E\rangle_{\mathrm{mat}}\langle E^{-1}\rangle_{\mathrm{mat}}-1
\]

%dsp-level:0
\[
= \frac{18-27p+11p^2-2D^2}{2D^2}\ \ [\mathrm{BM.94}(8)]
\]

%dsp-level:0
\[
\qquad \textbf{named,}\ \neg\textbf{filled};
\]

%dsp-level:0
\[
\text{the two Cal-F packagings}\ (\hat\mu^2,\ \hat\mu'^2),
\]

%dsp-level:0
\[
\ \ e_0\ \ \text{and the scales}\ (\hat\mu_\nu,\ \ell_c^{(\nu)},\ \mu_\nu)\ \ \text{stay undeclared}
\]

%dsp-level:0
\[
\qquad [\mathrm{Def\text{-}H.37};\ \mathrm{EM.25}(7);\ \mathrm{P4};\ \mathrm{P1}];
\]

%dsp-level:0
\[
(6)\quad\textbf{the currency bar stands:}
\]

%dsp-level:0
\[
\qquad \mathcal R\in\mathfrak q^0\ \ \text{is a ratio},
\]

%dsp-level:0
\[
\ \neg\text{a conversion}\ —\ \neg\text{a clause of this tier adds a pole-currency term to an informational assembly}:
\]

%dsp-level:0
\[
[\Delta_g]=\mathfrak q\quad\text{against}\quad[\Delta m]=\mathfrak q^{+1/2}
\]

%dsp-level:0
\[
\qquad \implies
\]

%dsp-level:0
\[
\qquad \neg\text{addable}
\]

%dsp-level:0
\[
\qquad [\mathrm{BM.37}(2);\ \mathrm{BM.69}(3);\ \mathrm{BM.83}(3);\ \mathrm{BM.90}(5);\ \mathrm{CT2.1};\ \text{machine: the exponent arithmetic}];
\]

%dsp-level:0
\[
\textbf{the flicker weight attenuates the informational register alone and converts nothing into the defect register}
\]

%dsp-level:0
\[
\qquad [\mathrm{BM.90}(4);\ \mathrm{BM.93}(5);\ \text{machine:}\ w_{\mathrm{mat}}\text{'s symbol census carries no}\ D_S];
\]

%dsp-level:0
\[
(7)\quad\textbf{the no-import rows:}
\]

%dsp-level:0
\[
\qquad \neg\text{the observed neutron--proton mass ratio is imported}\ —\ \text{the ratio is derived and bounded at}\ \mathrm{BM.99}(5);
\]

%dsp-level:0
\[
\qquad \neg\text{the observed electric coupling is imported}\ —\ e,
\]

%dsp-level:0
\[
\ e_0,
\]

%dsp-level:0
\[
\ \varepsilon_{\mathrm{EM}}\ \text{stay symbolic}\ [\mathrm{EM.25}(7),(9);\ \mathrm{BM.88}(7);\ \mathrm{P5}];
\]

%dsp-level:0
\[
\neg\text{a measured}\ \nu\ \text{mass, mixing angle, endpoint or lifetime enters}\ [\mathrm{P5}];
\]

%dsp-level:0
\[
\qquad \neg\text{an}\ [\mathrm{EXT}]\ \text{datum};
\]

%dsp-level:0
\[
\qquad \text{machine: zero external tokens, zero decimal numerals in this tier};
\]

%dsp-level:0
\[
(8)\quad\textbf{status:}
\]

%dsp-level:0
\[
\qquad \textbf{derived:}\ (1)\text{--}(4),\ (6);
\]

%dsp-level:0
\[
\qquad \textbf{closed in form with one priced residue:}\ (5);
\]

%dsp-level:0
\[
\qquad \textbf{symbolic:}\ e_0,\ \mathsf g,\ N,\ \rho_{\mathrm B};
\]

%dsp-level:0
\[
\neg\text{a declaration},
\]

%dsp-level:0
\[
\quad\neg\text{a P4 addition},
\]

%dsp-level:0
\[
\quad\text{§Cal-BM stays empty},
\]

%dsp-level:0
\[
\quad\neg\text{a new dimensionful scale}\ [\mathrm{P1}],
\]

%dsp-level:0
\[
\quad\neg\text{a measured value}\ [\mathrm{P5}];
\]

%dsp-level:0
\[
\qquad\alpha_B:\ \neg\text{added},\ \neg\text{derived},\ \text{barred}\ [\mathrm{BM.61}(8)].
\]

\textbf{Req.} D2.9 (Parts 1–2); A §0.3, CT2.1 (Part 2); EM.25(7), EM.25(9) (Part 2); P1, P4, P5 (Parts 1–2); Def-H.37(2), FT.54(4), FT.55(1), FT.73(2), FT.73(3) (Part 3); Census FT.28, FT.63(6) (row-wise); Def-BM.4(1), Def-BM.14(3), BM.13(1), BM.16(1), BM.17(3), BM.17(4), BM.36(3), BM.37(2), BM.61(8), BM.67(6), BM.67(8), BM.69(2), BM.69(3), BM.69(4), BM.69(5), BM.82(2), BM.82(6), BM.83(3), BM.88(2), BM.88(3), BM.88(4), BM.88(5), BM.88(6), BM.88(7), BM.89(4), BM.90(4), BM.90(5), BM.93(5), BM.94(2), BM.94(6), BM.94(7), BM.94(8) (this part); Cal-F; exact Laurent-monomial arithmetic with $\mathfrak q$- and $\mathfrak h$-exponent censuses, exact rational-function substitution, coefficient comparison, monotonicity on exact rationals, symbol censuses [math; machine: \texttt{checks/\allowbreak{}verify\_\allowbreak{}H\_\allowbreak{}flicker\_\allowbreak{}integral\_\allowbreak{}restmass.py} §E].

\textbf{Pf.}

%dsp-level:0
\[
(1):\ \mathrm{BM.69}(2)\text{'s identity and}\ \mathrm{FT.54}(4)\text{'s bridge ratio re-cited};
\]

%dsp-level:0
\[
\qquad [\rho_{\mathrm B}]=[\hat\mu'^2\mathfrak h^2]=\mathfrak q\cdot\mathfrak q^{-1}=\mathfrak q^0\ [\mathrm{Def\text{-}H.37}(2);\ \text{machine}]
\]

%dsp-level:0
\[
(2):\ \mathcal R_{\mathrm{mat}}
\]

%dsp-level:0
\[
= \rho_{\mathrm B}D_S/(4\langle E\rangle_{\mathrm{mat}}^2)
\]

%dsp-level:0
\[
= \rho_{\mathrm B}D_S/(4\cdot36/D^2)
\]

%dsp-level:0
\[
= \rho_{\mathrm B}D_SD^2/144\ [\mathrm{BM.94}(6);\ \text{machine}];
\]

%dsp-level:0
\[
\qquad D\in[1,3)\ [\mathrm{BM.94}(7)]\implies D^2\in[1,9)\ [\text{machine}];
\]

%dsp-level:0
\[
\qquad D=1\iff p=1\iff\varepsilon_q=0\ [\text{machine}]
\]

%dsp-level:0
\[
(3):\ E^2
\]

%dsp-level:0
\[
= G/(2\pi\ln2)\implies\rho_{\mathrm B}D_S/(4E^2)
\]

%dsp-level:0
\[
= \rho_{\mathrm B}D_S\cdot2\pi\ln2/(4G)
\]

%dsp-level:0
\[
= (\pi\ln2/2)\rho_{\mathrm B}D_S/G\ [\text{machine: coefficient comparison}];
\]

%dsp-level:0
\[
\qquad G^{\mathrm{eff}}=2\pi(6/D)^2\ln2=72\pi\ln2/D^2=G/D^2\ [\text{machine}]
\]

%dsp-level:0
\[
G^{\mathrm{cryst}}=2\pi\cdot36\cdot\ln2=72\pi\ln2\ [\mathrm{BM.17}(3);\ \text{machine}];
\]

%dsp-level:0
\[
\qquad 2\pi\cdot2^2\ln2=8\pi\ln2\ [\mathrm{BM.17}(4)\text{'s digon row};\ \text{machine}];
\]

%dsp-level:0
\[
\qquad D\in[1,3)\implies G/D^2\in(8\pi\ln2,\,72\pi\ln2]\ [\text{machine}]
\]

%dsp-level:0
\[
(4):\ 1-f
\]

%dsp-level:0
\[
= 1/k
\]

%dsp-level:0
\[
= \mathfrak h/(rE\ln2)
\]

%dsp-level:0
\[
\implies E
\]

%dsp-level:0
\[
= \mathfrak h/(r(1-f)\ln2)\implies\rho_{\mathrm B}D_S/(4E^2)
\]

%dsp-level:0
\[
= \rho_{\mathrm B}D_Sr^2(1-f)^2(\ln2)^2/(4\mathfrak h^2)
\]

%dsp-level:0
\[
= \hat\mu'^2D_Sr^2(1-f)^2/4\ [\text{machine: the substitution and}\ \rho_{\mathrm B}(\ln2)^2/\mathfrak h^2=\hat\mu'^2]
\]

%dsp-level:0
\[
e^2=e_0^2Nf\mathsf g\implies f=e^2/(e_0^2N\mathsf g)\ [\text{machine}];
\]

%dsp-level:0
\[
\qquad \mathcal R\propto(1-f)^2\ \text{strictly decreasing in}\ f\ \text{on}\ f<1\ [\mathrm{BM.88}(2)(\mathrm{ii}):\ 0\le f<1;\ \text{machine}];
\]

%dsp-level:0
\[
\qquad f\ \text{strictly increasing in}\ E_{\mathrm{int}}\ [\mathrm{BM.88}(6)]\implies\mathcal R\ \text{strictly decreasing in}\ E_{\mathrm{int}}\ \text{and in}\ e\ [\text{machine: the composition of two monotone rows}]
\]

%dsp-level:0
\[
k^{\mathrm{eff}}=r\langle E\rangle_{\mathrm{mat}}\ln2/\mathfrak h=(r\cdot6\ln2/\mathfrak h)/D=k^{\mathrm{cryst}}/D\ [\text{machine}];
\]

%dsp-level:0
\[
\qquad f^{\mathrm{eff}}=1-D/k^{\mathrm{cryst}}\le1-1/k^{\mathrm{cryst}}=f^{\mathrm{cryst}}\iff D\ge1\ [\text{machine}];
\]

%dsp-level:0
\[
\qquad f^{\mathrm{eff}}=0\iff k^{\mathrm{cryst}}=D\iff r=D\,\Delta x_{\mathrm B}^{\mathrm{cryst}}\ [\text{machine}]
\]

%dsp-level:0
\[
(5):\ \text{the inventory rows assembled from}\ (1)\text{--}(4)\text{'s own tags}\ [\text{machine}];
\]

%dsp-level:0
\[
\qquad \mathrm{BM.69}(4)\text{'s four grounds re-cited at their loci};
\]

%dsp-level:0
\[
\qquad \langle E\rangle\langle E^{-1}\rangle-1=\tfrac{18-27p+11p^2}{2D^2}-1=\tfrac{18-27p+11p^2-2D^2}{2D^2}\ [\mathrm{BM.94}(8);\ \text{machine}]
\]

%dsp-level:0
\[
(6):\ [\Delta_g]=[\hat\mu'^2]=\mathfrak q,\quad[\Delta m]=\mathfrak q^{+1/2}\ [\mathrm{CT2.1};\ \text{machine}];
\]

%dsp-level:0
\[
\qquad w_{\mathrm{mat}}
\]

%dsp-level:0
\[
= 1/D,
\]

%dsp-level:0
\[
\ D
\]

%dsp-level:0
\[
= 3-3p+p^2,
\]

%dsp-level:0
\[
\ p
\]

%dsp-level:0
\[
= 1-4\varepsilon_q:\ \text{the symbol set is}\ \{\varepsilon_q\}\ \text{and carries no}\ D_S\ [\text{machine: the token census}]
\]

%dsp-level:0
\[
(7)\text{--}(8):\ \text{the standing-rule scans: external tokens, decimal numerals, unit tokens, registry tokens, declaration lines, the barred object}\ \alpha_B\ [\text{machine}].\ \blacksquare
\]

\textbf{Phys.} The bridge between the two currencies has one shape and four exact readings of it. Count the edges and it is the single undeclared packaging times the state's defect factor over four times the count squared. Read the count through the calibration lock and the same object is pi log two over two times the packaging times the defect factor over the curvature monomial — the bridge is inversely proportional to the lock, with the coefficient exact. Read the count through the envelope ratio and it becomes mu-prime squared times the defect factor times the radius squared over four, times the square of one minus the nearby-space fraction; substitute the coupling's own profile and the bridge is an exact function of the running coupling, strictly decreasing in it at fixed geometry, so the energy-dependence of the bridge is literally the energy-dependence of the coupling and nothing else. Put the flicker in and every form gains one factor: the bridge is multiplied by the square of the dilution denominator, which is bounded by an exact integer nine, attained only in a limit that is not a materialized universe, and equal to one exactly at the sharp threshold where the snapshot returns. The curvature lock rides between two computed cells, the digon's and the junction's own, and the coupling's zero is confined to within a factor of three of the crystallized width. Everything in those four forms is derived exact algebra — one number is not, and the four grounds that bar its derivation stand where they were proved. So the relating equation is closed as an equation: its form, its flicker weighting, its curvature and coupling readings and its bound are theorems, and its single residue is one dimensionless ratio that this tier names, prices and does not invent.

\end{theorem}

\begin{theorem}[{\texorpdfstring{BM.99 — The materialized splitting and the Q-value: the two-sided parameter-free bar, the tightened bound, and the rescoped dissolution row}{BM.99 - The materialized splitting and the Q-value: the two-sided parameter-free bar, the tightened bound, and the rescoped dissolution row}}]\label{tag:BM-99}
\textbf{Thm.}

%dsp-level:0
\[
(1)\quad\textbf{the materialized assembly:}
\]

%dsp-level:0
\[
\qquad \boxed{\ m_n\,\mathfrak h\ =\ 6\,w_{\mathrm{mat}}\ln2+\ln20,\qquad m_p\,\mathfrak h\ =\ 6\,w_{\mathrm{mat}}\ln2+\ln15+3\,N\,w_{\mathrm{mat}}\,\mathsf g\,\varepsilon_{\mathrm{EM}}\ln2\ }
\]

%dsp-level:0
\[
[\mathrm{BM.91}(2)\ \text{at}\ w\to w_{\mathrm{mat}};\ \mathrm{BM.94}(7),(8);\ \mathrm{Def\text{-}BM.13}(5);\ \mathrm{BM.89}(5);\ \text{machine: the prime-log assembly}\ \ln20=2\ln2+\ln5,\ \ln15=\ln3+\ln5];
\]

%dsp-level:0
\[
\boxed{\ \Delta m\,\mathfrak h\ =\ \ln\frac43-3\,N\,w_{\mathrm{mat}}\,\mathsf g\,\varepsilon_{\mathrm{EM}}\ln2\ }\qquad[\text{machine:}\ \ln20-\ln15=\ln\tfrac43];
\]

%dsp-level:0
\[
\textbf{the cancellation row stands:}
\]

%dsp-level:0
\[
\qquad E_{\mathrm{int}}\ \ \text{is class-shared}\ [\mathrm{BM.57}(4);\ \mathrm{Census\ FT.35}(2)]\implies\text{the crystallization weight cancels in the ledger difference and survives only in the EM term}
\]

%dsp-level:0
\[
\qquad [\text{machine: the coefficient census}];
\]

%dsp-level:0
\[
\textbf{the register scoping stands:}
\]

%dsp-level:0
\[
\qquad w_{\mathrm{mat}}\ \ \text{attenuates the edge-count register alone};
\]

%dsp-level:0
\[
\qquad \text{the transport register}\ \Sigma_{\mathrm{cl}}(Q)=\ln W(Q)\ \ \text{enters at its crystallized value under}\ (\mathrm{R4})
\]

%dsp-level:0
\[
\qquad [\mathrm{BM.91}(2);\ \mathrm{Def\text{-}BM.13}(4),(6);\ \mathrm{BM.89}(8)];
\]

%dsp-level:0
\[
(2)\quad\textbf{the critical locus, enclosed:}
\]

%dsp-level:0
\[
\qquad \varepsilon^{\mathrm{crit}}_{\mathrm{mat}}(N,\mathsf g)=\frac{\ln(4/3)}{3\,N\,w_{\mathrm{mat}}\,\mathsf g\,\ln2}=\frac{2-\log_23}{3\,N\,w_{\mathrm{mat}}\,\mathsf g}
\]

%dsp-level:0
\[
\qquad [\mathrm{BM.91}(4);\ \text{machine: solved exactly in the prime-log ring}]:
\]

%dsp-level:0
\[
\boxed{\ \frac{2-\log_23}{3\,N\,\mathsf g}\ \le\ \varepsilon^{\mathrm{crit}}_{\mathrm{mat}}\ <\ \frac{2-\log_23}{N\,\mathsf g}\ \subset\ \Bigl(\frac{2}{15\,N\mathsf g},\ \ \frac{5}{12\,N\mathsf g}\Bigr)\ }
\]

%dsp-level:0
\[
\qquad [\mathrm{BM.91}(4)\text{'s squeeze}\ \tfrac25<2-\log_23<\tfrac5{12}\ \text{on}\ 3^5=243<256=2^8\ \text{and}\ 3^{12}=531441>524288=2^{19};\ \mathrm{BM.94}(7);\ \text{machine}];
\]

%dsp-level:0
\[
\varepsilon^{\mathrm{crit}}_{\mathrm{mat}}\notin\mathbb Q\ \ \text{at every rational}\ (N,\mathsf g)>0
\]

%dsp-level:0
\[
\qquad [\mathrm{BM.84}(4);\ \mathrm{BM.91}(4):\ 2-\log_23\notin\mathbb Q\ \text{on the}\ 2^p\ne3^q\ \text{census};\ \text{machine}];
\]

%dsp-level:0
\[
(3)\quad\textbf{the two-sided parameter-free bar:}
\]

%dsp-level:0
\[
\qquad \boxed{\ N\,\mathsf g\,\varepsilon_{\mathrm{EM}}\ \le\ \frac18\qquad\Longrightarrow\qquad\Delta m>0\ }
\]

%dsp-level:0
\[
\qquad [\mathrm{BM.91}(5)\ \text{re-derived without}\ w:\ \ w_{\mathrm{mat}}\le1\ \text{and}\ \tfrac{2-\log_23}3>\tfrac18\iff2^{13}=8192>6561=3^8];
\]

%dsp-level:0
\[
\boxed{\ N\,\mathsf g\,\varepsilon_{\mathrm{EM}}\ \ge\ \frac5{12}\qquad\Longrightarrow\qquad\Delta m<0\ }
\]

%dsp-level:0
\[
\qquad [\text{machine:}\ w_{\mathrm{mat}}>\tfrac13\ \text{strictly}\implies Nw_{\mathrm{mat}}\mathsf g\,\varepsilon_{\mathrm{EM}}>\tfrac5{36},\ \ \text{and}\ \ \tfrac5{36}>\tfrac{2-\log_23}3\iff\tfrac5{12}>2-\log_23\iff3^{12}>2^{19}];
\]

%dsp-level:0
\[
\boxed{\ \textbf{the materialization filter removes the crystallization weight from both bars}\ }
\]

%dsp-level:0
\[
\qquad \textbf{the undecided band is exact and non-empty:}
\]

%dsp-level:0
\[
\qquad \frac18<N\,\mathsf g\,\varepsilon_{\mathrm{EM}}<\frac5{12}
\]

%dsp-level:0
\[
\qquad [\text{machine:}\ 12<40];
\]

%dsp-level:0
\[
\text{both bars are uniform over the lapse, the curvature composition and the whole admissible threshold range}\ \varepsilon_q\in[0,\tfrac14)
\]

%dsp-level:0
\[
\qquad [\mathrm{CT0.8};\ \text{machine}];
\]

%dsp-level:0
\[
(4)\quad\textbf{the trichotomy (total, exact, both regions certified by integers):}
\]

%dsp-level:0
\[
\qquad \boxed{\ \varepsilon_{\mathrm{EM}}<\varepsilon^{\mathrm{crit}}_{\mathrm{mat}}\iff m_n>m_p;\qquad=\iff m_n=m_p;\qquad>\iff m_n<m_p\ }
\]

%dsp-level:0
\[
[\text{machine:}\ \Delta m\ \text{strictly decreasing in}\ \varepsilon_{\mathrm{EM}}\ \text{at every admissible}\ (N,w_{\mathrm{mat}},\mathsf g)>0];
\]

%dsp-level:0
\[
\qquad \text{the opening region is inhabited by}\ \varepsilon_{\mathrm{EM}}
\]

%dsp-level:0
\[
= 0\ \text{and by}\ \varepsilon_{\mathrm{EM}}
\]

%dsp-level:0
\[
= \tfrac19\ \text{at}\ N\mathsf g
\]

%dsp-level:0
\[
= 1\ \ [3);
\]

%dsp-level:0
\[
\ \mathrm{BM.62}(8):\ 4^3
\]

%dsp-level:0
\[
= 64>54\qquad = 3^3\!\cdot\!2];
\]

%dsp-level:0
\[
\text{the closing region is inhabited by}\ \varepsilon_{\mathrm{EM}}=\tfrac5{12}\ \text{at}\ N\mathsf g=1
\]

%dsp-level:0
\[
\qquad [(3);\ \text{machine}];
\]

%dsp-level:0
\[
\qquad \text{the critical locus is exact, unique and irrational}\qquad \qquad [(2)];
\]

%dsp-level:0
\[
(5)\quad\textbf{the tightened universal ratio bound:}
\]

%dsp-level:0
\[
\qquad \frac{m_n}{m_p}=\frac{6\,w_{\mathrm{mat}}\ln2+\ln20}{6\,w_{\mathrm{mat}}\ln2+\ln15+3Nw_{\mathrm{mat}}\mathsf g\,\varepsilon_{\mathrm{EM}}\ln2}
\]

%dsp-level:0
\[
\qquad [\text{machine}]:
\]

%dsp-level:0
\[
\text{strictly decreasing in}\ \varepsilon_{\mathrm{EM}}\ \text{and in}\ w_{\mathrm{mat}}
\]

%dsp-level:0
\[
\qquad [\mathrm{BM.91}(7);\ \text{machine: the two cross-multiplied sign rows}]\ \ \text{and}\ \ \tfrac13<w_{\mathrm{mat}}\le1
\]

%dsp-level:0
\[
\qquad [\mathrm{BM.94}(7)]:
\]

%dsp-level:0
\[
\boxed{\ \frac{\ln1280}{\ln960}\ \le\ \frac{m_n}{m_p}\bigg|_{\varepsilon_{\mathrm{EM}}=0}\ <\ \frac{\ln80}{\ln60}\ <\ \frac{\ln20}{\ln15}\ <\ \frac{10}9\ }
\]

%dsp-level:0
\[
\qquad [\text{machine:}\ 6\ln2+\ln20=\ln1280,\quad2\ln2+\ln20=\ln80,\quad2\ln2+\ln15=\ln60]:
\]

%dsp-level:0
\[
\frac{\ln80}{\ln60}\in\Bigl(\frac{16}{15},\ \frac{15}{14}\Bigr)
\]

%dsp-level:0
\[
\qquad [\text{machine:}\ 80^{15}>60^{16}\iff2^{28}=268435456>215233605=3^{16}\!\cdot\!5;\qquad80^{14}<60^{15}\iff2^{26}=67108864<71744535=3^{15}\!\cdot\!5];
\]

%dsp-level:0
\[
\frac{15}{14}<\frac{11}{10}<\frac{\ln20}{\ln15}
\]

%dsp-level:0
\[
\qquad [\text{machine:}\ 15\cdot10<14\cdot11;\qquad20^{10}=10240000000000>8649755859375=15^{11}]:
\]

%dsp-level:0
\[
\boxed{\ 1\ <\ \frac{m_n}{m_p}\ <\ \frac{\ln80}{\ln60}\ <\ \frac{15}{14}\qquad\text{on the materialized read at every admissible}\ (\varepsilon_q,\,N,\,\mathsf g,\,\varepsilon_{\mathrm{EM}})\ \text{on the sub-critical branch}\ }
\]

%dsp-level:0
\[
\textbf{an exact tightening of}\ \mathrm{BM.91}(7)\textbf{'s}\ \tfrac{10}9\ \textbf{by the filter's lower bound on the weight}
\]

%dsp-level:0
\[
\qquad [\text{machine:}\ \tfrac{15}{14}<\tfrac{10}9\iff135<140];
\]

%dsp-level:0
\[
\textbf{the incommensurability hardens:}
\]

%dsp-level:0
\[
\qquad \frac{m_n}{m_p}=\frac43\iff\varepsilon^{*}_{\mathrm{EM}}=\frac{(6-6w_{\mathrm{mat}})\ln2-\ln405}{12\,N\,w_{\mathrm{mat}}\,\mathsf g\,\ln2}<0
\]

%dsp-level:0
\[
\qquad [\mathrm{BM.91}(8);\ \text{machine:}\ (6-6w_{\mathrm{mat}})\ln2\le6\ln2=\ln64<\ln405\ \text{on}\ 64<405];
\]

%dsp-level:0
\[
\frac{m_n}{m_p}\notin\mathbb Q\ \ \text{at every rational}\ (w_{\mathrm{mat}},N,\mathsf g,\varepsilon_{\mathrm{EM}})
\]

%dsp-level:0
\[
\qquad [\mathrm{BM.83}(2)\text{'s lemma: the}\ \ln3\ \text{coefficient is}\ 0\ \text{in the numerator and}\ 1\ \text{in the denominator};\ \text{machine}];
\]

%dsp-level:0
\[
(6)\quad\textbf{the Q-value on the materialized informational read:}
\]

%dsp-level:0
\[
\qquad Q_\beta\,\mathfrak h=\ln\frac43-3Nw_{\mathrm{mat}}\mathsf g\,\varepsilon_{\mathrm{EM}}\ln2-F_{\ell\nu}\ \ \le\ \ln\tfrac43-2\ln2\ =\ -\ln3\ <\ 0
\]

%dsp-level:0
\[
[\mathrm{BM.92}(1);\ \mathrm{BM.70}(2),(3):\ F_{\ell\nu}\ge2\ln2;\ \text{machine}]
\]

%dsp-level:0
\[
\qquad \textbf{uniformly in}\ (\varepsilon_q,\,N,\,\mathsf g,\,\varepsilon_{\mathrm{EM}})\ —\ \textbf{the closure survives the filter unchanged};
\]

%dsp-level:0
\[
\textbf{the ratio form tightens:}
\]

%dsp-level:0
\[
\qquad \boxed{\ -\frac{\ln3}{\ln1280}\ \le\ \frac{Q_\beta}{m_n}\ \le\ -\frac{\ln3}{6\,w_{\mathrm{mat}}\ln2+\ln20}\ <\ -\frac{\ln3}{\ln80}\ }
\]

%dsp-level:0
\[
\qquad \text{with}
\]

%dsp-level:0
\[
\qquad \boxed{\ -\frac{\ln3}{\ln80}\in\Bigl(-\frac13,\,-\frac14\Bigr)\ }
\]

%dsp-level:0
\[
[\text{machine:}\ 3^4=81>80\ \text{and}\ 80>27]\qquad\textbf{an exact tightening of}\ \mathrm{BM.92}(2)\textbf{'s}\ \bigl(-\tfrac25,-\tfrac13\bigr)\ \text{endpoint};
\]

%dsp-level:0
\[
\textbf{the margin ladder stands:}
\]

%dsp-level:0
\[
\qquad \Sigma_{\mathrm B}-2\ln2=-\ln3,
\]

%dsp-level:0
\[
\qquad \Sigma_{\mathrm B}-4\ln2=-\ln12,
\]

%dsp-level:0
\[
\qquad \Sigma_{\mathrm B}-5\ln2=-\ln24
\]

%dsp-level:0
\[
\qquad [\mathrm{BM.70}(3);\ \mathrm{BM.92}(2);\ \text{machine:}\ 4<12<48<96];
\]

%dsp-level:0
\[
(7)\quad\textbf{the rescoped dissolution row:}
\]

%dsp-level:0
\[
\qquad \mathrm{BM.91}(4)\text{'s row}\ \bigl(\varepsilon^{\mathrm{crit}}\to\infty\ \text{as}\ \varepsilon_q\to\tfrac14^{-},\ \text{``the fully dissolved limit never closes the splitting''}\bigr)\ \textbf{is rescoped,}\ \neg\textbf{retracted}:
\]

%dsp-level:0
\[
\text{at}\ \varepsilon_q\to\tfrac14^{-}\ \ \text{the unconditioned ensemble's weight concentrates on the}\ E=0\ \ \text{member},
\]

%dsp-level:0
\[
\quad P(E{=}0)=(1-p)^3\to1
\]

%dsp-level:0
\[
\qquad [\mathrm{BM.94}(4),(5);\ \text{machine}];
\]

%dsp-level:0
\[
\text{that member supports no closed chain and no terminus}
\]

%dsp-level:0
\[
\qquad [\mathrm{FT.144}(4);\ \mathrm{Def\text{-}H.4}]\implies\boxed{\ \textbf{the fully dissolved limit is outside the materialization domain}\ }
\]

%dsp-level:0
\[
\textbf{on the materialized read the splitting is closable at a finite coupling:}
\]

%dsp-level:0
\[
\qquad \varepsilon^{\mathrm{crit}}_{\mathrm{mat}}<\frac{5}{12\,N\,\mathsf g}
\]

%dsp-level:0
\[
\qquad [(2),(3);\ \text{machine}];
\]

%dsp-level:0
\[
\qquad \text{the unconditioned row stands at its own locus as the}\ (\mathrm{R4})\ \text{general-weight statement};
\]

%dsp-level:0
\[
(8)\quad\textbf{status:}
\]

%dsp-level:0
\[
\qquad \textbf{derived:}\ (1)\text{--}(7)\ \text{on the marked pair specialization};
\]

%dsp-level:0
\[
\qquad \textbf{rescoped,}\ \neg\textbf{retracted:}\ \mathrm{BM.91}(2),
\]

%dsp-level:0
\[
(4),
\]

%dsp-level:0
\[
(7),
\]

%dsp-level:0
\[
(8),
\]

%dsp-level:0
\[
\ \mathrm{BM.92}(1),
\]

%dsp-level:0
\[
(2)\ \text{stand on}\ w\ \text{at their own loci and gain the materialized variants here};
\]

%dsp-level:0
\[
\mathrm{BM.84},\ \mathrm{BM.85},\ \mathrm{BM.77}(5),(6)\ \text{untouched};
\]

%dsp-level:0
\[
\qquad \neg\text{a declaration},
\]

%dsp-level:0
\[
\quad\neg\text{a P4 addition},
\]

%dsp-level:0
\[
\quad\neg\text{a new dimensionful scale}\ [\mathrm{P1}],
\]

%dsp-level:0
\[
\quad\neg\text{a measured value}\ [\mathrm{P5}],
\]

%dsp-level:0
\[
\quad\text{§Cal-BM stays empty};
\]

%dsp-level:0
\[
\neg\text{the observed neutron--proton ratio imported}\ —\ \text{the bound}\ (5)\ \text{is derived from the filter and the prime-log assembly alone};\qquad
\]

\textbf{Req.} D0.4 (Parts 1–2); CT0.8, CT2.1 (Part 1, Part 2); P1, P4, P5 (Parts 1–2); Def-H.4, FT.144(4) (Part 3); Census FT.35(2) (row-wise); Def-BM.13(4), Def-BM.13(5), Def-BM.13(6), BM.57(4), BM.62(8), BM.70(2), BM.70(3), BM.77(5), BM.77(6), BM.83(2), BM.84(4), BM.85, BM.89(5), BM.89(8), BM.91(2), BM.91(4), BM.91(5), BM.91(7), BM.91(8), BM.92(1), BM.92(2), BM.94(4), BM.94(5), BM.94(6), BM.94(7), BM.94(8) (this part); exact prime-log arithmetic with integer power certificates, exact rational function rows, monotonicity by cross-multiplication, token censuses [math; machine: \texttt{checks/\allowbreak{}verify\_\allowbreak{}H\_\allowbreak{}flicker\_\allowbreak{}integral\_\allowbreak{}restmass.py} §F].

\textbf{Pf.}

%dsp-level:0
\[
(1):\ \mathrm{BM.91}(2)\text{'s assembly with}\ w\to w_{\mathrm{mat}}\ \text{and}\ \langle m_{\mathrm{hid}}\rangle_{\mathrm{mat}}=6w_{\mathrm{mat}}\ln2/\mathfrak h\ [\mathrm{BM.94}(8);\ \mathrm{BM.89}(5)];
\]

%dsp-level:0
\[
\qquad \ln20-\ln15=\ln\tfrac43\ [\text{machine}];
\]

%dsp-level:0
\[
\qquad \text{the}\ w\text{-coefficient census in the two rest inputs and in}\ \Delta m\ [\text{machine}]
\]

%dsp-level:0
\[
(2):\ \varepsilon^{\mathrm{crit}}
\]

%dsp-level:0
\[
= \ln(4/3)/(3Nw_{\mathrm{mat}}\mathsf g\ln2)
\]

%dsp-level:0
\[
= (2-\log_23)/(3Nw_{\mathrm{mat}}\mathsf g)\ [\mathrm{BM.91}(4)\ \text{re-cited};\ \text{machine}];
\]

%dsp-level:0
\[
\qquad \tfrac13<w_{\mathrm{mat}}\le1\ [\mathrm{BM.94}(7)]\implies\text{the two-sided enclosure}\ [\text{machine}];
\]

%dsp-level:0
\[
\qquad \tfrac25<2-\log_23<\tfrac5{12}\ \text{on}\ 3^5<2^8\ \text{and}\ 3^{12}>2^{19}\ [\mathrm{BM.91}(4)\ \text{re-cited}]
\]

%dsp-level:0
\[
(3):\ w_{\mathrm{mat}}
\]

%dsp-level:0
\[
\le 1
\]

%dsp-level:0
\[
\implies Nw_{\mathrm{mat}}\mathsf g\,
\]

%dsp-level:0
\[
\varepsilon_{\mathrm{EM}}
\]

%dsp-level:0
\[
\le N\mathsf g\,
\]

%dsp-level:0
\[
\varepsilon_{\mathrm{EM}}\le\tfrac18<\tfrac{2-\log_23}3\ [\text{machine:}\ 8(2-\log_23)>3\iff13>8\log_23\iff2^{13}>3^8]\implies\Delta m>0
\]

%dsp-level:0
\[
w_{\mathrm{mat}}>\tfrac13\ \text{and}\ N\mathsf g\,\varepsilon_{\mathrm{EM}}\ge\tfrac5{12}\implies Nw_{\mathrm{mat}}\mathsf g\,\varepsilon_{\mathrm{EM}}>\tfrac5{36};
\]

%dsp-level:0
\[
\qquad \tfrac5{36}>\tfrac{2-\log_23}3\iff\tfrac5{12}>2-\log_23\iff\log_23>\tfrac{19}{12}
\]

%dsp-level:0
\[
\iff 3^{12}
\]

%dsp-level:0
\[
= 531441>524288
\]

%dsp-level:0
\[
= 2^{19}\ [\text{machine}]\implies\Delta m<0
\]

%dsp-level:0
\[
\tfrac18<\tfrac5{12}\iff12<40\ [\text{machine}]
\]

%dsp-level:0
\[
(4):\ \text{the coefficient of}\ \varepsilon_{\mathrm{EM}}\ \text{in}\ \Delta m\ \text{is}\ -3Nw_{\mathrm{mat}}\mathsf g\ln2<0\ [\text{machine}];
\]

%dsp-level:0
\[
\qquad \varepsilon_{\mathrm{EM}}=\tfrac19:\ N\mathsf g\,\varepsilon_{\mathrm{EM}}=\tfrac19\le\tfrac18\implies\Delta m>0\ \text{by}\ (3);
\]

%dsp-level:0
\[
\qquad \varepsilon_{\mathrm{EM}}=\tfrac5{12}\ \text{at}\ N\mathsf g=1\implies\Delta m<0\ \text{by}\ (3)
\]

%dsp-level:0
\[
(5):\ \text{the ratio's monotonicity in}\ \varepsilon_{\mathrm{EM}}\ \text{and in}\ w_{\mathrm{mat}}\ \text{re-cited}\ [\mathrm{BM.91}(7);\ \text{machine}];
\]

%dsp-level:0
\[
\qquad 6w_{\mathrm{mat}}\ln2\in(2\ln2,\,6\ln2]\ [\mathrm{BM.94}(7)]\implies\text{the four-endpoint chain}\ [\text{machine}]
\]

%dsp-level:0
\[
80^{15}>60^{16}\iff2^{60}5^{15}>2^{32}3^{16}5^{16}\iff2^{28}>3^{16}\cdot5:\quad268435456>215233605\ [\text{machine}];
\]

%dsp-level:0
\[
\qquad 80^{14}<60^{15}\iff2^{26}<3^{15}\cdot5:\quad67108864<71744535\ [\text{machine}]
\]

%dsp-level:0
\[
\tfrac{15}{14}<\tfrac{11}{10}\iff150<154;
\]

%dsp-level:0
\[
\qquad 20^{10}>15^{11}\ [\mathrm{BM.91}(7)\ \text{re-cited}];
\]

%dsp-level:0
\[
\qquad \tfrac{15}{14}<\tfrac{10}9\iff135<140\ [\text{machine}]
\]

%dsp-level:0
\[
(6):\ \mathrm{BM.92}(1)\text{'s uniform closure re-cited with}\ w\to w_{\mathrm{mat}}\ [\text{machine}];
\]

%dsp-level:0
\[
\qquad -\ln3/\ln80\in(-\tfrac13,-\tfrac14)\iff3\ln3>\ln80\ \text{and}\ 4\ln3<\ln80\iff81>80\ \text{and}\ 80>27\ [\text{machine}];
\]

%dsp-level:0
\[
\qquad \text{the three margin identities re-certified}\ [\text{machine}]
\]

%dsp-level:0
\[
(7):\ (1-p)^3\to1\ \text{as}\ p\to0^+\ [\text{machine}];
\]

%dsp-level:0
\[
\qquad |E_{\mathcal R}|=0\implies\neg\exists\ \text{a terminus}\ [\mathrm{FT.144}(4)\ \text{re-cited}];
\]

%dsp-level:0
\[
\qquad \varepsilon^{\mathrm{crit}}_{\mathrm{mat}}<(2-\log_23)/(N\mathsf g)<\tfrac5{12N\mathsf g}\ [(2);\ \text{machine}];
\]

%dsp-level:0
\[
\qquad (8):\ \text{the standing-rule scans}\ [\text{machine}].\ \blacksquare
\]

\textbf{Phys.} Filter the ensemble to the configurations that can carry a composite and three rows change shape. The assembly keeps its form — six times the weight times log two, plus the class's own transport entropy, less the electromagnetic term — but the weight is now one over one plus u plus u squared, and it cannot fall below a third. That single bound does real work. The critical coupling is squeezed into an exact window between two fifteenths and five twelfths over the lapse and the curvature composition, and the two bars that decide the sign of the splitting lose the crystallization weight entirely: an eighth or less of the coupling product forces the neutron heavier, five twelfths or more forces it lighter, and the undecided band between them is exact and inhabited. The ratio bound tightens from ten ninths to the log of eighty over the log of sixty, itself enclosed between sixteen fifteenths and fifteen fourteenths by two integer power certificates, and the Q-value's ratio form tightens from a third to a quarter by eighty-one against eighty. The dissolution row is rescoped rather than retracted: at the far end of the threshold range the unconditioned ensemble piles its weight onto the empty configuration, which supports no chain and no terminus, so that limit is not a materialized universe at all — and inside the materialization domain the splitting can always be closed by a finite coupling. Nothing measured entered: the ratio is derived from the filter and the prime-log assembly, the coupling stays symbolic, and the closure of the decay on the informational currency survives the filter exactly as it survived the ensemble.

\end{theorem}

\medskip\noindent\textit{\texorpdfstring{BM.100 — The status of the flicker integral and the rest-mass read: adjudication and scope}{BM.100 - The status of the flicker integral and the rest-mass read: adjudication and scope}.}\ Structural-constraints record: see Appendix B, Section B.1.29 of the Addendum.\par\medskip

\phantomsection
\section[Tier BM.20 — Six Objective Resolutions: the Energy-Scale…]{\texorpdfstring{Tier BM.20 — Six Objective Resolutions: the Energy-Scale Graph, the Dimensional Bridge, the Eigenvalue Composition, the Curvature Composition, the General Crystallization Weight, and the Kind Structure of the Arrow}{Tier BM.20 - Six Objective Resolutions: the Energy-Scale Graph, the Dimensional Bridge, the Eigenvalue Composition, the Curvature Composition, the General Crystallization Weight, and the Kind Structure of the Arrow}}

\textbf{Locus.} This tier takes six named objectives and gives each its exact adjudication: one graph, one census with one declared slot, one referent table, one fixed point, one measure-free bound, and one kind structure. The Q-value becomes a function of the energy scale, its forward and inverse reads are shown to be one graph reflected, and its crossing is a derived strength enclosed by two integers — with the crossing \textit{count} depending on the priced curvature composition ($\mathrm{BM.101}$). The dimension surface is censused against the single information flux, the seven rows are derived or adjudicated, and the two bridges are separated: the dimensional one is $\mathfrak h$, the numerical one is $\rho_{\mathrm B}$, whose status is declared as the §II.4 addendum A-BM.9 with no value issued ($\mathrm{BM.102}$). The three-constituent assembly's eigenvalue equations are derived where they are formed — the family cubic and the binomial characteristic polynomial — and the tensor-product referent is censured where it is not, so the composition question and the octet question are both adjudicated $\neg$formed ($\mathrm{BM.103}$). The curvature composition's self-consistency condition is formed exactly, solved for all three priced candidates, and shown to select none, because the lock's own freedom census excludes the back-reaction ($\mathrm{BM.104}$). The crystallization weight is derived measure-free and for an arbitrary threshold family, with the same exact bounds, and the defect-conditioned version is refuted ($\mathrm{BM.105}$). The spin label and the tick arrow are derived, and the exchange structure that a statistics requires is censured: deficit 11 stands with its residue named ($\mathrm{BM.106}$). $\mathrm{BM.107}$ adjudicates.

\textbf{Standing rows.} $\alpha_B$ stays barred on $\mathrm{BM.61}(8)$'s four grounds. No clause of $\mathrm{BM.1}$–$\mathrm{BM.100}$ is retracted, replaced or renumbered: every row this tier touches is re-scoped at its own locus and recorded in $\mathrm{BM.101}(8)$, $\mathrm{BM.104}(8)$, $\mathrm{BM.105}(7)$ and $\mathrm{BM.107}(5)$, in $\mathrm{BM.77}(7)$'s pattern. No measured coupling, mass or ratio enters any derivation $[\mathrm{P5}]$: the crossing strength of $\mathrm{BM.101}(3)$ is a derived symbolic value and is $\neg$compared to any external number; the neutron–proton ratio stays derived and bounded $[\mathrm{BM.99}(5)]$. No right-ray matter line and no seesaw-type mechanism is suggested, added or required $[\mathrm{BM.96}(4),(6)\ \text{re-cited},\ \neg\text{extended}]$; $\neg$an external mixing matrix is imported — the mixing rows are reads of the derived $16$-arena $[\mathrm{FT.142}(2),(5);\ \mathrm{FT.145}(3)]$. One dimension only: units are powers of $\mathfrak q$ $[\mathrm{D2.9};\ \mathrm{CT2.1}]$. Local notation, fixed for the whole tier: $E_1:=\mathfrak h/(r\ln2)$ the envelope-unit count; $x:=E/E_1=k$ the scale ratio $[\mathrm{BM.88}(2)(\mathrm i)]$; $\beta:=2\pi\ln2\,E_1^2=2\pi\mathfrak h^2/(r^2\ln2)\in\mathfrak q^0$; $f=1-\tfrac1x$; $G=\beta x^2$; $\varepsilon_{\mathrm{eff}}:=Nf\,\mathsf g\,\varepsilon_{\mathrm{EM}}=e(r,E,N)^2/(4\pi)$; $\varepsilon_\ast(w)$ the crossing strength; $c:=\varepsilon_\ast/(N\varepsilon_{\mathrm{EM}})$; $a:=\tfrac12Nw\,\varepsilon_{\mathrm{EM}}Q^2$; $p_i,u_i,R,S$ the family data of $\mathrm{BM.105}(3)$; $[\cdot]$ the $\mathfrak q$-exponent map $[\mathrm{CT2.1};\ \mathrm{D2.9}]$. The symbols $w$, $D$, $p$, $u$, $\Delta_g$, $M_0$, $s$, $\mathcal R$, $\rho_{\mathrm B}$, $D_S$ keep their $\mathrm{BM.94}$–$\mathrm{BM.99}$ meanings; the italic $g$ keeps its $\mathrm{BM.39}(1)$ meaning and $q$ its end-register meaning $[\mathrm{D3.4};\ \mathrm{EM.9}]$ at every occurrence.

\begin{theorem}[{\texorpdfstring{BM.101 — The energy-scale graph: the two reads, the derived crossing strength, and the composition-dependent crossing count}{BM.101 - The energy-scale graph: the two reads, the derived crossing strength, and the composition-dependent crossing count}}]\label{tag:BM-101}
\textbf{Thm.}

%dsp-level:0
\[
(1)\quad\textbf{the two reads are one graph reflected:}
\]

%dsp-level:0
\[
\qquad Q_\beta^{\mathrm{fwd}}(E)\ :=\ m_n^{\mathrm{rest}}(E)-m_p^{\mathrm{rest}}(E)-m_\ell^{\mathrm{rest}}-m_\nu^{\mathrm{rest}},
\]

%dsp-level:0
\[
\qquad Q_\beta^{\mathrm{inv}}(E)\ :=\ m_p^{\mathrm{rest}}(E)+m_\ell^{\mathrm{rest}}+m_\nu^{\mathrm{rest}}-m_n^{\mathrm{rest}}(E)
\]

%dsp-level:0
\[
\boxed{\ Q_\beta^{\mathrm{inv}}(E)\ =\ -\,Q_\beta^{\mathrm{fwd}}(E)\ \ \text{identically, at every scale, on every read}\ }
\]

%dsp-level:0
\[
\qquad [\text{machine: one negation identity};\ \mathrm{BM.97}(2),(4)\ \text{re-cited}]:
\]

%dsp-level:0
\[
\textbf{the inverse reaction opens exactly where the forward one closes, and their common boundary is one locus}\ —\ \neg\text{two independent graphs};
\]

%dsp-level:0
\[
(2)\quad\textbf{the scale-dependent splitting (the graph's spine):}
\]

%dsp-level:0
\[
\qquad \varepsilon_{\mathrm{eff}}(r,E,N)\ :=\ N\,f\!\Bigl(\frac{r}{\Delta x_{\mathrm B}(E)}\Bigr)\,\mathsf g\bigl(G(E)\bigr)\,\varepsilon_{\mathrm{EM}}\ =\ \frac{e(r,E,N)^2}{4\pi}
\]

%dsp-level:0
\[
\qquad [\mathrm{BM.88}(4);\ \mathrm{Def\text{-}BM.12}\ \text{local notation};\ \text{machine: the monomial identity}]:
\]

%dsp-level:0
\[
\boxed{\ \Delta m(E)\,\mathfrak h\ =\ \ln\frac43\ -\ 3\,w\,\ln2\cdot\varepsilon_{\mathrm{eff}}(r,E,N)\ }
\]

%dsp-level:0
\[
\qquad [\mathrm{BM.91}(2);\ \mathrm{BM.99}(1)\ \text{at the full profile instead of its far-field locus}\ f=1;\ \text{machine}];
\]

%dsp-level:0
\[
\textbf{the far-field rescope:}
\]

%dsp-level:0
\[
\qquad f\to1\ \ \text{recovers}\ \ \Delta m\,\mathfrak h=\ln\tfrac43-3Nw\,\mathsf g\,\varepsilon_{\mathrm{EM}}\ln2\ \ \text{exactly}
\]

%dsp-level:0
\[
\qquad [\mathrm{BM.88}(5);\ \mathrm{BM.91}(2);\ \mathrm{BM.99}(1):\ \text{every row stands at that locus}];
\]

%dsp-level:0
\[
\textbf{the informational Q-value on the graph:}
\]

%dsp-level:0
\[
\qquad Q_\beta^{\mathrm{fwd}}\,\mathfrak h=\Delta m(E)\,\mathfrak h-F_{\ell\nu},
\]

%dsp-level:0
\[
\qquad F_{\ell\nu}\ge2\ln2
\]

%dsp-level:0
\[
\qquad [\mathrm{BM.70}(2),(3);\ \mathrm{BM.92}(1)]:
\]

%dsp-level:0
\[
\boxed{\ Q_\beta^{\mathrm{fwd}}\,\mathfrak h\ \le\ -\ln3\ <\ 0\ \ \text{and}\ \ Q_\beta^{\mathrm{inv}}\,\mathfrak h\ \ge\ \ln3\ >\ 0\qquad\text{at every scale, uniformly in}\ (r,\,E,\,\varepsilon_q,\,N,\,\mathsf g,\,\varepsilon_{\mathrm{EM}})\ }
\]

%dsp-level:0
\[
[\text{machine: the sign census on a rational lattice}]
\]

%dsp-level:0
\[
\qquad \textbf{on the informational currency the forward channel is closed and the inverse channel is open at every energy scale, with the exact margin}\ \ln3;
\]

%dsp-level:0
\[
(3)\quad\textbf{the derived crossing strength:}
\]

%dsp-level:0
\[
\qquad \Delta m(E)=0\iff\varepsilon_{\mathrm{eff}}(r,E,N)=\varepsilon_\ast(w):
\]

%dsp-level:0
\[
\qquad \boxed{\ \varepsilon_\ast(w)\ :=\ \frac{\ln(4/3)}{3\,w\ln2}\ =\ \frac{2-\log_23}{3\,w},\qquad e_\ast^2\ :=\ 4\pi\,\varepsilon_\ast(w)\ =\ \frac{4\pi\,(2-\log_23)}{3\,w}\ }
\]

%dsp-level:0
\[
[\mathrm{BM.91}(4);\ \text{machine: exact in the}\ \log_2\ \text{ring}];\qquad\textbf{the enclosure:}\qquad w=w_{\mathrm{mat}}\in\bigl(\tfrac13,1\bigr]\qquad[\mathrm{BM.94}(7)]:
\]

%dsp-level:0
\[
\boxed{\ \frac{2-\log_23}3\ \le\ \varepsilon_\ast\ <\ 2-\log_23\qquad\text{and}\qquad\frac18\ <\ \varepsilon_\ast\ <\ \frac5{12}\ }
\]

%dsp-level:0
\[
\qquad [\text{machine:}\ 2^{13}=8192>6561=3^8;\quad3^{12}=531441>524288=2^{19};\quad3^5=243<256=2^8]:
\]

%dsp-level:0
\[
\boxed{\ \textbf{the two endpoints of}\ \mathrm{BM.99}(3)\textbf{'s two-sided parameter-free bar are exactly the enclosure of the crossing strength}\ }
\]

%dsp-level:0
\[
\qquad [\text{machine: the four integers re-certified}];
\]

%dsp-level:0
\[
e_\ast^2\in\bigl(\tfrac\pi2,\ \tfrac{5\pi}3\bigr)
\]

%dsp-level:0
\[
\qquad [\text{machine}];
\]

%dsp-level:0
\[
\qquad \textbf{no measured coupling is compared to}\ \varepsilon_\ast\ \text{or}\ e_\ast\ \ [\mathrm{P5};\ \mathrm{EM.25}(7),(9);\ \text{the quarantined audit rows stay quarantined}];
\]

%dsp-level:0
\[
(4)\quad\textbf{the crossing locus is a level set, and}\ e_0\ \textbf{places it:}
\]

%dsp-level:0
\[
\qquad \boxed{\ \mathcal L_\ast\ :=\ \bigl\{(r,E,N):\ \ e(r,E,N)^2=e_\ast^2\bigr\}\ =\ \bigl\{(r,E,N):\ \varepsilon_{\mathrm{eff}}=\varepsilon_\ast(w)\bigr\}\ }
\]

%dsp-level:0
\[
e^2=e_0^2\,N\,f\,\mathsf g
\]

%dsp-level:0
\[
\qquad [\mathrm{BM.88}(4)]\implies\mathcal L_\ast\ \ \text{is the level set of the derived profile at the derived strength};
\]

%dsp-level:0
\[
\textbf{the adjudication of the cancellation claim:}
\]

%dsp-level:0
\[
\qquad e_0^2\ \ \textbf{does}\ \neg\textbf{cancel}\ —\ \text{it is the one datum that places}\ \mathcal L_\ast\ \text{in}\ (r,E):
\]

%dsp-level:0
\[
\boxed{\ \textbf{what is}\ e_0\textbf{-free is the crossing value}\ \varepsilon_\ast,\ \textbf{its enclosure, and the crossing count per composition};\qquad\textbf{what requires}\ e_0\ \textbf{is the crossing location}\ }
\]

%dsp-level:0
\[
\neg\exists\ \text{a theorem fixing}\ e_0\ \text{or}\ e
\]

%dsp-level:0
\[
\qquad [\mathrm{EM.25}(7);\ \mathrm{BM.88}(7);\ \mathrm{P4};\ \mathrm{P5}];
\]

%dsp-level:0
\[
\qquad \neg\text{a numerical value of either enters any clause of this tier}\ [\text{machine: the token scan}];
\]

%dsp-level:0
\[
(5)\quad\textbf{the three profile shapes (composition by composition):}
\]

%dsp-level:0
\[
\qquad x=\frac{E}{E_1}=k=\frac{r\,E\ln2}{\mathfrak h},
\]

%dsp-level:0
\[
\qquad E_1=\frac{\mathfrak h}{r\ln2},
\]

%dsp-level:0
\[
\qquad G=2\pi E^2\ln2=\beta x^2,
\]

%dsp-level:0
\[
\qquad \beta=\frac{2\pi\mathfrak h^2}{r^2\ln2}\in\mathfrak q^0
\]

%dsp-level:0
\[
[\mathrm{BM.88}(2)(\mathrm i),(\mathrm{iii});\ \mathrm{BM.88}(3);\ \mathrm{CT2.1};\ \text{machine: the exponent arithmetic}];
\]

%dsp-level:0
\[
\qquad P_{\mathsf g}(x)\ :=\ f(x)\,\mathsf g\bigl(\beta x^2\bigr),
\]

%dsp-level:0
\[
\qquad \varepsilon_{\mathrm{eff}}=N\,\varepsilon_{\mathrm{EM}}\,P_{\mathsf g}(x):
\]

%dsp-level:0
\[
(\mathrm i)\quad\mathsf g=1:
\]

%dsp-level:0
\[
\qquad P(x)=1-\frac1x\ \ \text{strictly increasing on}\ x\ge1,
\]

%dsp-level:0
\[
\qquad P(1)=0,
\]

%dsp-level:0
\[
\qquad P\to1:
\]

%dsp-level:0
\[
\qquad \boxed{\ \text{at most one crossing; it exists}\iff N\varepsilon_{\mathrm{EM}}>\varepsilon_\ast\ }
\]

%dsp-level:0
\[
\textbf{the crossing scale in closed form:}
\]

%dsp-level:0
\[
\qquad x_\ast=\frac1{1-c},
\]

%dsp-level:0
\[
\qquad c:=\frac{\varepsilon_\ast}{N\,\varepsilon_{\mathrm{EM}}}\in(0,1),
\]

%dsp-level:0
\[
\qquad E_\ast=\frac{\mathfrak h}{r\ln2\,(1-c)}
\]

%dsp-level:0
\[
\qquad [\text{machine: the linear equation}\ x-1=cx];
\]

%dsp-level:0
\[
(\mathrm{ii})\quad\mathsf g=1+G:
\]

%dsp-level:0
\[
\qquad P(x)=\Bigl(1-\frac1x\Bigr)\bigl(1+\beta x^2\bigr)\ \ \text{strictly increasing},
\]

%dsp-level:0
\[
\qquad P(1)=0,
\]

%dsp-level:0
\[
\qquad P\to\infty:
\]

%dsp-level:0
\[
\qquad \boxed{\ \textbf{exactly one crossing at every admissible}\ (N,\varepsilon_{\mathrm{EM}})>0\ }
\]

%dsp-level:0
\[
\textbf{the crossing equation is an exact cubic:}
\]

%dsp-level:0
\[
\qquad \beta x^3-\beta x^2+(1-c)\,x-1=0
\]

%dsp-level:0
\[
\qquad [\text{machine: the cleared identity}\ (x-1)(1+\beta x^2)=cx];
\]

%dsp-level:0
\[
(\mathrm{iii})\quad\mathsf g=(1+G)^{-1}:\qquad P(x)=\frac{1-1/x}{1+\beta x^2},\qquad P(1)=0,\qquad P\to0\ \ \text{as}\ \ x\to\infty:\qquad\textbf{a window, not a half-line}
\]

%dsp-level:0
\[
\textbf{the unique interior maximum:}
\]

%dsp-level:0
\[
\qquad P'(x_\ast)=0\iff2\beta x_\ast^3-3\beta x_\ast^2-1=0
\]

%dsp-level:0
\[
\qquad [\text{machine: the cleared stationarity row}];
\]

%dsp-level:0
\[
\text{the cubic}\ h(x)=2\beta x^3-3\beta x^2-1\ \ \text{has}\ h'(x)=6\beta x(x-1)\ge0\ \ \text{on}\ \ x\ge1,
\]

%dsp-level:0
\[
\qquad h(1)=-\beta-1<0,
\]

%dsp-level:0
\[
\qquad h\to+\infty\implies\boxed{\ \exists!\ x_\ast(\beta)>1\ }
\]

%dsp-level:0
\[
\boxed{\ P_{\max}(\beta)=P(x_\ast)=1-\frac3{2x_\ast}\ }
\]

%dsp-level:0
\[
\qquad [\text{machine: the polynomial identity}\ 2x\bigl(1-\tfrac1x\bigr)-(2x-3)(1+\beta x^2)=1-(2\beta x^3-3\beta x^2),\ \text{vanishing on the cubic}];
\]

%dsp-level:0
\[
\textbf{the exact dichotomy in}\ \beta:
\]

%dsp-level:0
\[
\qquad h\bigl(\tfrac32\bigr)=-1<0\ \ \forall\beta,
\]

%dsp-level:0
\[
\qquad h(2)=4\beta-1
\]

%dsp-level:0
\[
\qquad [\text{machine: two polynomial identities}]:
\]

%dsp-level:0
\[
\boxed{\ \beta>\frac14\implies x_\ast\in\Bigl(\frac32,\,2\Bigr)\ \ \text{and}\ \ P_{\max}\in\Bigl(0,\,\frac14\Bigr);\qquad\beta=\frac14\implies x_\ast=2,\quad P_{\max}=\frac14\ \ \text{exactly};\qquad\beta<\frac14\implies x_\ast>2,\quad P_{\max}\in\Bigl(\frac14,\,1\Bigr)\ }
\]

%dsp-level:0
\[
\textbf{the crossing equation is an exact depressed cubic:}
\]

%dsp-level:0
\[
\qquad c\beta x^3+(c-1)\,x+1=0
\]

%dsp-level:0
\[
\qquad [\text{machine: the cleared identity}\ x-1=cx(1+\beta x^2)];
\]

%dsp-level:0
\[
\textbf{the root count is the crossing count:}
\]

%dsp-level:0
\[
\qquad c\beta x^3+(c-1)x+1=-\,x\bigl(1+\beta x^2\bigr)\bigl(P(x)-c\bigr)
\]

%dsp-level:0
\[
\qquad [\text{machine: one polynomial identity}]:
\]

%dsp-level:0
\[
\boxed{\ c<P_{\max}:\ \textbf{two crossings}\ (\text{the window opens and closes});\qquad c=P_{\max}:\ \textbf{one}\ (x_\ast,\ \text{a double root});\qquad c>P_{\max}:\ \textbf{none}\ }
\]

%dsp-level:0
\[
(6)\quad\textbf{the specified and pure-pole graphs are exactly flat in the energy scale:}
\]

%dsp-level:0
\[
\qquad Q_\beta^{\mathrm{spec}}\,\mathfrak h=\Delta T(S)-\rho_\ell\sqrt{G_\ell}-\rho_\nu\sqrt{G_\nu},
\]

%dsp-level:0
\[
\qquad \Delta T=5\,m^2_B(\cdot,S)
\]

%dsp-level:0
\[
\text{the alphabet}\ \{m^2_B(\cdot,S),\ \rho_\ell,\ \rho_\nu,\ g,\ g',\ \mathfrak h\}\ \ \text{carries}\ \neg E_{\mathrm{int}},\ \neg r,\ \neg N,\ \neg\mathsf g,\ \neg\varepsilon_{\mathrm{EM}}
\]

%dsp-level:0
\[
\qquad [\mathrm{BM.92}(5);\ \mathrm{BM.97}(2),(3);\ \text{machine: the token census}]:
\]

%dsp-level:0
\[
\boxed{\ \textbf{the specified read's energy-scale graph is a horizontal line, and its trichotomy is decided by the scales}\ (\rho_\ell,\rho_\nu),\ \neg\textbf{by the scale}\ E\ }
\]

%dsp-level:0
\[
\qquad [\mathrm{BM.77}(5),(6)\ \text{re-cited}];
\]

%dsp-level:0
\[
\text{the pure-pole read is flat in}\ E\ \text{as well, its crossing being}\ \Delta_g^{\mathrm{crit}}=s(2M_0+s)
\]

%dsp-level:0
\[
\qquad [\mathrm{BM.97}(4)];
\]

%dsp-level:0
\[
\qquad \textbf{the entire energy-scale structure lives in the informational/EM channel};
\]

%dsp-level:0
\[
(7)\quad\textbf{the graph's verdict rows:}
\]

%dsp-level:0
\[
\qquad \textbf{the monotonicity:}\ \ \Delta m(E)\ \ \text{is strictly decreasing in}\ E\ \ \text{on}\ (\mathrm i),
\]

%dsp-level:0
\[
(\mathrm{ii}),
\]

%dsp-level:0
\[
\ \ \text{and rises then falls on}\ (\mathrm{iii})
\]

%dsp-level:0
\[
\qquad [\text{machine: three monotonicity censuses on exact rationals}];
\]

%dsp-level:0
\[
\textbf{the envelope row:}
\]

%dsp-level:0
\[
\qquad f(1)=0\implies\varepsilon_{\mathrm{eff}}=0\ \ \text{at}\ \ E=E_1\implies\Delta m(E_1)\,\mathfrak h=\ln\tfrac43>0
\]

%dsp-level:0
\[
\qquad [\mathrm{BM.88}(5);\ \text{machine}]:
\]

%dsp-level:0
\[
\qquad \textbf{the splitting is open at the member's own envelope on every composition};
\]

%dsp-level:0
\[
\textbf{the bar as a pointwise crossing census:}\qquad f\le1\ \ \text{on}\ \ x\ge1\implies\varepsilon_{\mathrm{eff}}\ \le\ N\,\mathsf g\bigl(G(E)\bigr)\,\varepsilon_{\mathrm{EM}}\qquad[\text{machine}]:
\]

%dsp-level:0
\[
N\mathsf g\,\varepsilon_{\mathrm{EM}}\le\tfrac18<\varepsilon_\ast\implies\boxed{\ \Delta m(E)>0\ \ \text{at that scale}\ },
\]

%dsp-level:0
\[
\qquad N\mathsf g\,\varepsilon_{\mathrm{EM}}\ge\tfrac5{12}>\varepsilon_\ast\implies\boxed{\ \Delta m(E)<0\ \ \text{at that scale}\ }
\]

%dsp-level:0
\[
\qquad [\mathrm{BM.99}(3)\ \text{re-derived on the graph};\ \text{machine}];
\]

%dsp-level:0
\[
\textbf{the global crossing census per composition:}
\]

%dsp-level:0
\[
\qquad (\mathrm i)\ \ \text{a crossing exists}\iff N\varepsilon_{\mathrm{EM}}>\varepsilon_\ast;
\]

%dsp-level:0
\[
\qquad (\mathrm{ii})\ \ \text{exactly one at every}\ (N,\varepsilon_{\mathrm{EM}})>0;
\]

%dsp-level:0
\[
(\mathrm{iii})\ \ \text{two}\iff N\varepsilon_{\mathrm{EM}}P_{\max}>\varepsilon_\ast,\qquad\text{one at equality},\qquad\text{none below}\qquad[(5);\ \text{machine}];
\]

%dsp-level:0
\[
\textbf{the band is where the flicker decides:}
\]

%dsp-level:0
\[
\qquad \frac18<N\mathsf g\,\varepsilon_{\mathrm{EM}}<\frac5{12}\implies\text{the crossing's existence is decided by}\ \varepsilon_\ast(w)=\frac{2-\log_23}{3w},\ \ \text{i.e. by}\ w
\]

%dsp-level:0
\[
\qquad [(3);\ \text{machine}];
\]

%dsp-level:0
\[
(8)\quad\textbf{status:}
\]

%dsp-level:0
\[
\qquad \textbf{derived:}\ (1)\text{--}(7);
\]

%dsp-level:0
\[
\qquad \textbf{read-relative:}\ \text{the verdict rows}\ (2),(6);
\]

%dsp-level:0
\[
\qquad \textbf{rescoped,}\ \neg\textbf{retracted:}\ \mathrm{BM.91}(2),\ \mathrm{BM.99}(1),(3)\ \text{stand as the far-field locus}\ f=1;
\]

%dsp-level:0
\[
\neg\text{a declaration},
\]

%dsp-level:0
\[
\quad\neg\text{a P4 addition},
\]

%dsp-level:0
\[
\quad\neg\text{a new dimensionful scale}\ [\mathrm{P1}],
\]

%dsp-level:0
\[
\quad\neg\text{a measured value}\ [\mathrm{P5}],
\]

%dsp-level:0
\[
\quad\neg\text{a numerical coupling}\ [\mathrm{EM.25}(7)];
\]

%dsp-level:0
\[
\qquad\textbf{to calibration:}\ e_0\ (4),\quad\mathsf g\ (5)\ [\mathrm{BM.104}(5)],\quad N\ \text{as solution data}\ [\mathrm{CT2.7};\ \mathrm{BM.34}(2)].
\]

\textbf{Req.} D1.8, D2.9, D3.4 (Parts 1–2); A6, CT2.1, CT2.7, EM.9, EM.25(7), EM.25(9) (Part 2); P1, P4, P5 (Parts 1–2); Def-BM.12, Def-BM.4(1), BM.13(2), BM.14(2), BM.16(1), BM.17(2), BM.23(9), BM.28(4), BM.34(2), BM.39(1), BM.70(2), BM.70(3), BM.77(5), BM.77(6), BM.84(9), BM.88(2), BM.88(3), BM.88(4), BM.88(5), BM.88(6), BM.88(7), BM.89(5), BM.91(2), BM.91(4), BM.92(1), BM.92(5), BM.94(5), BM.94(7), BM.96(3), BM.97(2), BM.97(3), BM.97(4), BM.99(1), BM.99(3), BM.99(5), BM.99(6) (this part); Census FT.35(2) (row-wise); exact polynomial identities over $\mathbb Q[\beta,x,c]$, cubic root counts by monotonicity and sign, exact rational lattices, the $\log_2$-ring sign test with integer power certificates, Laurent-monomial assembly, token censuses [math; machine: \texttt{checks/\allowbreak{}verify\_\allowbreak{}H\_\allowbreak{}six\_\allowbreak{}objective\_\allowbreak{}resolution.py} §A].

\textbf{Pf.}

%dsp-level:0
\[
(1):\ Q_\beta^{\mathrm{inv}}
\]

%dsp-level:0
\[
= (m_p+m_\ell+m_\nu)-m_n
\]

%dsp-level:0
\[
= -\bigl(m_n-m_p-m_\ell-m_\nu\bigr)
\]

%dsp-level:0
\[
= -Q_\beta^{\mathrm{fwd}}\ [\text{machine: one identity over}\ \mathbb Q[m_n,m_p,m_\ell,m_\nu]]
\]

%dsp-level:0
\[
(2):\ \varepsilon_{\mathrm{eff}}
\]

%dsp-level:0
\[
= Nf\,
\]

%dsp-level:0
\[
\mathsf g\,
\]

%dsp-level:0
\[
\varepsilon_{\mathrm{EM}}\ \text{and}\ e^2
\]

%dsp-level:0
\[
= e_0^2Nf\,
\]

%dsp-level:0
\[
\mathsf g\ \text{with}\ \varepsilon_{\mathrm{EM}}
\]

%dsp-level:0
\[
= e_0^2/(4\pi)\implies\varepsilon_{\mathrm{eff}}
\]

%dsp-level:0
\[
= e^2/(4\pi)\ [\mathrm{BM.88}(4);\ \text{machine: one monomial}];
\]

%dsp-level:0
\[
\qquad \mathrm{BM.91}(2)\text{'s term}\ 3Nw\,\mathsf g\,\varepsilon_{\mathrm{EM}}\ln2\ \text{is the}\ f{=}1\ \text{locus of}\ 3w\ln2\cdot\varepsilon_{\mathrm{eff}}\ [\text{machine}]
\]

%dsp-level:0
\[
Q_\beta^{\mathrm{fwd}}\mathfrak h
\]

%dsp-level:0
\[
= \ln\tfrac43-3w\ln2\cdot\varepsilon_{\mathrm{eff}}-F_{\ell\nu}\le\ln\tfrac43-2\ln2
\]

%dsp-level:0
\[
= -\ln3\ [\mathrm{BM.70}(2),(3);\ \text{machine: the coefficient census on a}\ 504\text{-point rational lattice}];
\]

%dsp-level:0
\[
\qquad -\ln3<0\iff4<12\ [\text{machine}]
\]

%dsp-level:0
\[
(3):\ \Delta m
\]

%dsp-level:0
\[
= 0
\]

%dsp-level:0
\[
\iff 3w\ln2\cdot\varepsilon_{\mathrm{eff}}
\]

%dsp-level:0
\[
= \ln\tfrac43\iff\varepsilon_{\mathrm{eff}}
\]

%dsp-level:0
\[
= \tfrac{\ln(4/3)}{3w\ln2}
\]

%dsp-level:0
\[
= \tfrac{2-\log_23}{3w}\ [\text{machine: exact in the}\ \log_2\ \text{ring}];
\]

%dsp-level:0
\[
\qquad e_\ast^2=4\pi\varepsilon_\ast\ [\text{machine}]
\]

%dsp-level:0
\[
w\in(\tfrac13,1]\ [\mathrm{BM.94}(7)]\implies\varepsilon_\ast\in[\tfrac{2-\log_23}3,\ 2-\log_23)\ [\text{machine}];
\]

%dsp-level:0
\[
\qquad \tfrac{2-\log_23}3>\tfrac18\iff8(2-\log_23)>3\iff2^{13}>3^8;
\]

%dsp-level:0
\[
\qquad 2-\log_23<\tfrac5{12}\iff3^{12}>2^{19}\ [\text{machine: two integer power comparisons}]
\]

%dsp-level:0
\[
(4):\ \mathcal L_\ast\ \text{is the preimage of one value under the derived profile}\ [\mathrm{BM.88}(4)];
\]

%dsp-level:0
\[
\qquad e^2
\]

%dsp-level:0
\[
= e_0^2Nf\,
\]

%dsp-level:0
\[
\mathsf g\implies\text{the locus moves with}\ e_0\ \text{and its value does not}\ [\text{machine: the symbol census of}\ \varepsilon_\ast\ \text{carries}\ \{w,\ln2,\log_23,\pi\},\ \neg e_0]
\]

%dsp-level:0
\[
(5):\ x=k=rE\ln2/\mathfrak h=E/E_1\ \text{at}\ E_1=\mathfrak h/(r\ln2)\ [\text{machine}];
\]

%dsp-level:0
\[
\qquad G=2\pi E^2\ln2=2\pi\ln2\,E_1^2x^2=\beta x^2,
\]

%dsp-level:0
\[
\qquad [\beta]=[\mathfrak h^2/r^2]=\mathfrak q^{-1}\mathfrak q^{+1}=\mathfrak q^0\ [\mathrm{CT2.1};\ \text{machine}]
\]

%dsp-level:0
\[
(\mathrm i):\ P=1-\tfrac1x,\ P'=1/x^2>0\ \text{on}\ x\ge1;\ P(1)=0;\ P\to1\ [\text{machine}];
\]

%dsp-level:0
\[
\qquad x-1=cx\iff x(1-c)=1\iff x=1/(1-c),\ \ \text{admissible}\iff0<c<1\iff N\varepsilon_{\mathrm{EM}}>\varepsilon_\ast\ [\text{machine}]
\]

%dsp-level:0
\[
(\mathrm{ii}):\ P'=\tfrac1{x^2}(1+\beta x^2)+(1-\tfrac1x)2\beta x>0\ \text{on}\ x>1;\ P\to\infty\ [\text{machine}];
\]

%dsp-level:0
\[
\qquad (x-1)(1+\beta x^2)
\]

%dsp-level:0
\[
= cx
\]

%dsp-level:0
\[
\iff x+\beta x^3-1-\beta x^2
\]

%dsp-level:0
\[
= cx\iff\beta x^3-\beta x^2+(1-c)x-1
\]

%dsp-level:0
\[
= 0\ [\text{machine: the cleared cubic}]
\]

%dsp-level:0
\[
(\mathrm{iii}):\ P'=\tfrac1{x^2(1+\beta x^2)}-\tfrac{(1-\tfrac1x)2\beta x}{(1+\beta x^2)^2};
\]

%dsp-level:0
\[
\qquad P'=0\iff(1+\beta x^2)=2\beta x^3(1-\tfrac1x)\cdot\tfrac{x^2}{x^2}\iff2\beta x^3-3\beta x^2-1=0\ [\text{machine: cleared}]
\]

%dsp-level:0
\[
h'=6\beta x^2-6\beta x=6\beta x(x-1)\ge0\ \text{on}\ x\ge1,\ \ \text{with equality only at}\ x=1;
\]

%dsp-level:0
\[
\qquad h(1)=-\beta-1<0;
\]

%dsp-level:0
\[
\qquad h(x)\to+\infty\implies\exists!\,x_\ast>1\ [\text{machine: monotonicity plus the two endpoint signs}]
\]

%dsp-level:0
\[
\beta x_\ast^2(2x_\ast-3)
\]

%dsp-level:0
\[
= 1
\]

%dsp-level:0
\[
\implies 1+\beta x_\ast^2
\]

%dsp-level:0
\[
= \tfrac{2x_\ast-2}{2x_\ast-3}
\]

%dsp-level:0
\[
\implies P(x_\ast)
\]

%dsp-level:0
\[
= \tfrac{x_\ast-1}{x_\ast}\cdot\tfrac{2x_\ast-3}{2x_\ast-2}
\]

%dsp-level:0
\[
= \tfrac{2x_\ast-3}{2x_\ast}
\]

%dsp-level:0
\[
= 1-\tfrac3{2x_\ast}\ [\text{machine: the identity}\ 2x(1-\tfrac1x)-(2x-3)(1+\beta x^2)=1-(2\beta x^3-3\beta x^2)\ \text{in}\ \mathbb Q[\beta,x]]
\]

%dsp-level:0
\[
h(\tfrac32)=2\beta\tfrac{27}8-3\beta\tfrac94-1=-1;
\]

%dsp-level:0
\[
\qquad h(2)
\]

%dsp-level:0
\[
= 16\beta-12\beta-1
\]

%dsp-level:0
\[
= 4\beta-1\ [\text{machine: two polynomial identities}]\implies\text{the}\ \beta\text{-dichotomy and}\ P_{\max}\ \text{enclosures}\ [(1-\tfrac3{2x_\ast})\ \text{monotone in}\ x_\ast]
\]

%dsp-level:0
\[
x-1=cx(1+\beta x^2)\iff c\beta x^3+(c-1)x+1=0;
\]

%dsp-level:0
\[
\qquad x(1+\beta x^2)(P-c)
\]

%dsp-level:0
\[
= (x-1)-cx(1+\beta x^2)
\]

%dsp-level:0
\[
= -\bigl(c\beta x^3+(c-1)x+1\bigr)\ [\text{machine: one identity}]\implies\text{the root count is the crossing count}
\]

%dsp-level:0
\[
(6):\ \Delta T=5\,m^2_B(\cdot,S)\ \text{is a frozen rational per cell}\ [\mathrm{BM.77}(2)];
\]

%dsp-level:0
\[
\qquad G_\ell,G_\nu\ \text{depend on}\ (g,g',\delta)\ \text{alone}\ [\mathrm{BM.39}(1);\ \mathrm{BM.96}(2)];
\]

%dsp-level:0
\[
\qquad \text{the token census}\ [\text{machine}]
\]

%dsp-level:0
\[
(7):\ \varepsilon_{\mathrm{eff}}(E_1)=Nf(1)\,\mathsf g\,\varepsilon_{\mathrm{EM}}=0\implies\Delta m(E_1)\mathfrak h=\ln\tfrac43>0\ [\mathrm{BM.88}(5);\ \text{machine}];
\]

%dsp-level:0
\[
\qquad f\le1\ \text{on}\ x\ge1\implies\varepsilon_{\mathrm{eff}}\le N\mathsf g\,\varepsilon_{\mathrm{EM}}\ [\text{machine}];
\]

%dsp-level:0
\[
\qquad (8):\ \text{the scope record}.\ \blacksquare
\]

\textbf{Phys.} One graph, read twice. The forward and inverse Q-values are exact negatives, so the two reactions the question asks about are one curve and its reflection, and the boundary between them is a single level set of the derived coupling profile. On the informational currency the curve never changes sign: the two-bit floor on the emitted lines keeps the forward channel closed and the inverse channel open at every scale, with the exact margin log three, so the only sign structure the graph carries is the splitting's. That structure is decided by one derived number, the crossing strength — log of four thirds over three weights times log two, which is two minus the binary logarithm of three over three weights — enclosed between an eighth and five twelfths by the same two integer powers that give the parameter-free bar. The bar and the crossing are the same statement: below the lower end the curve never crosses at any scale, above the upper end it always does, and inside the band the flicker weight alone decides. The profile's own shape then does the rest. With the flat composition the crossing is a closed form, one over one minus c; with the growing composition there is exactly one crossing at every positive coupling, because the profile diverges; with the decreasing composition the profile is a window — it rises from zero, peaks at the unique root of an exact cubic, and returns to zero, so the forward channel opens at one scale and closes at another, or never opens, and the peak has the closed form one minus three over twice the root, sitting below a quarter whenever the radius parameter exceeds a quarter. Two compositions give a half-line, one gives a window, and the corpus still cannot say which is its own. What is exact and free of every undeclared datum is the crossing value; what needs the one open calibration number is where in radius and energy that value is reached.

\end{theorem}

\begin{theorem}[{\texorpdfstring{BM.102 — The dimensional bridge: the $\mathfrak q$-exponent census, the seven rows, and the declared calibration slot}{BM.102 - The dimensional bridge: the -exponent census, the seven rows, and the declared calibration slot}}]\label{tag:BM-102}
\textbf{Thm.}

%dsp-level:0
\[
(1)\quad\textbf{the one-dimension law:}
\]

%dsp-level:0
\[
\qquad \text{units}\ :=\ \text{powers of}\ \mathfrak q;
\]

%dsp-level:0
\[
\qquad c=\hbar=k_B=1;
\]

%dsp-level:0
\[
\qquad [\mathfrak h]=\mathfrak q^{-1/2};
\]

%dsp-level:0
\[
\qquad \mathfrak h:=\text{the single free dimensionful scale}
\]

%dsp-level:0
\[
\qquad [\mathrm{D2.9};\ \mathrm{P1};\ \text{A §0.3}]:
\]

%dsp-level:0
\[
\kappa,\ G_N\ \ \text{are names for combinations of}\ \mathfrak h,\ \neg\text{independent inputs}
\]

%dsp-level:0
\[
\qquad [\mathrm{D2.9};\ \mathrm{CT2.9}];
\]

%dsp-level:0
\[
\qquad \boxed{\ \textbf{the inventory carries exactly one fundamental dimension, the information flux}\ \mathfrak q\ }
\]

%dsp-level:0
\[
(2)\quad\textbf{the derived exponent census:}\qquad\begin{array}{l|l|l}\text{quantity} & \text{exponent} & \text{locus}\\\hline\text{length},\ \text{time} & \mathfrak q^{-1/2} & [\mathrm{CT2.1};\ \mathrm{D2.9}:\ c=1]\\\text{area},\ [\Lambda],\ [G_{\mu\nu}],\ \text{curvature} & \mathfrak q^{-1},\ \mathfrak q^{+1},\ \mathfrak q^{+1},\ \mathfrak q^{+1} & [\mathrm{CT2.1}]\\\text{volume} & \mathfrak q^{-3/2} & [\mathrm{Def\text{-}BM.4}(1);\ \mathrm{CT2.1}]\\\text{energy},\ \text{mass},\ [\Delta m] & \mathfrak q^{+1/2} & [\mathrm{CT2.1};\ \mathrm{BM.35}(4)]\\\text{energy density}\ \rho,\ [T_{\mu\nu}] & \mathfrak q^{+2} & [\mathrm{CT2.1};\ \mathrm{BM.14}(2),(4)]\\{}[\kappa^2],\ [G_N],\ [\kappa] & \mathfrak q^{-1},\ \mathfrak q^{-1},\ \mathfrak q^{-1/2} & [\mathrm{CT2.1};\ \mathrm{CT2.9}]\\\mathfrak h & \mathfrak q^{-1/2} & [\mathrm{D2.9}]\\\Delta S_{\mathfrak L},\ I_{\mathrm{hid}},\ S_{\mathrm{raw}},\ \Sigma_{\mathrm{cl}},\ S_{\mathrm{geom}} & \mathfrak q^{0} & [\mathrm{BM.35}(3),(4);\ \mathrm{CT1.29};\ \mathrm{CT2.1};\ \mathrm{Def\text{-}BM.13}(6)]\\E_{\mathrm{int}},\ W(Q),\ q,\ D_S,\ f,\ N,\ \mathsf g,\ w,\ G,\ \varepsilon_{\mathrm{EM}},\ e,\ g,\ g' & \mathfrak q^{0} & [\mathrm{D3.4};\ \mathrm{EM.9};\ \mathrm{Def\text{-}BM.13}(3);\ \mathrm{BM.88}(2),(3);\ \mathrm{EM.25}(9)]\\m^2_{\mathrm{pole}},\ \Delta_g,\ \hat\mu'^2 & \mathfrak q^{+1} & [\mathrm{Def\text{-}H.37}(2);\ \mathrm{FT.54}(4)]\\\hat\mu^2 & \mathfrak q^{+2} & [\mathrm{Def\text{-}H.37}(1);\ \mathrm{EW.33}]\\\rho_{\mathrm B}=\hat\mu'^2\mathfrak h^2/(\ln2)^2 & \mathfrak q^{0} & [\mathrm{BM.69}(2);\ \text{machine:}\ 1-1=0]\end{array}
\]

%dsp-level:0
\[
\textbf{the balance rows:}
\]

%dsp-level:0
\[
\qquad [\rho\,V]=\mathfrak q^{+2}\mathfrak q^{-3/2}=\mathfrak q^{+1/2}=[\text{energy}];
\]

%dsp-level:0
\[
\qquad [\kappa^2\rho]=\mathfrak q^{-1}\mathfrak q^{+2}=\mathfrak q^{+1}=[R];
\]

%dsp-level:0
\[
\qquad [(\Delta x_{\mathrm B})^2\kappa^2\rho]=\mathfrak q^{-1}\mathfrak q^{+1}=\mathfrak q^0=[G]
\]

%dsp-level:0
\[
\qquad [\text{machine}];
\]

%dsp-level:0
\[
(3)\quad\textbf{the seven rows (derived where a referent exists, adjudicated where it does not):}\qquad\begin{array}{l|l|l}\text{row} & \text{exponent} & \text{the corpus referent}\\\hline\text{mass} & \mathfrak q^{+1/2} & m_{\mathrm{hid}},\ m_{\mathrm{rest}},\ m_{\mathrm{pole}}\ [\mathrm{FT.44}(4);\ \mathrm{BM.35}(4);\ \mathrm{Def\text{-}H.37}]\\\text{length} & \mathfrak q^{-1/2} & \Delta x_{\mathrm B},\ \ell_c,\ r\ [\mathrm{Def\text{-}BM.4}(1);\ \mathrm{Def\text{-}H.2}]\\\text{time} & \mathfrak q^{-1/2} & x^0,\ \text{equal to length on}\ c=1\ [\mathrm{D1.8};\ \mathrm{D2.9}]\\\text{energy} & \mathfrak q^{+1/2} & E_{\mathrm{rest}}=m_{\mathrm{rest}}c^2\ \text{with the}\ c\text{-exponent}\ 0\ [\mathrm{BM.35}(4)]\\\text{temperature} & \mathfrak q^{+1/2} & \text{on}\ k_B=1\ [\mathrm{D2.9};\ \mathrm{CT2.1}]\ —\ \textbf{a dimension row,}\ \neg\textbf{a thermal object:}\ \neg\text{a P4-class temperature exists}\ [\mathrm{BM.69}(6)]\\\text{amount of substance} & \mathfrak q^{0} & \text{a count:}\ E_{\mathrm{int}},\ W(Q),\ \text{the fibre cardinalities}\ [\mathrm{Def\text{-}BM.5}(4);\ \mathrm{Def\text{-}BM.13}(3);\ \mathrm{BM.95}(1)]\\\text{electric current} & \mathfrak q^{0}\ \text{and}\ \mathfrak q^{+1/2} & \textbf{two reads:}\ \text{the ledger-side weight per tick}\ \ [\mathrm{BM.23}(9);\ \mathrm{BM.44};\ \mathrm{D1.8};\ \neg\text{a rate per external second}\ [\mathrm{P5}]];\ \ \text{the chart-side read below}\end{array}
\]

%dsp-level:0
\[
\textbf{the chart-side current, derived:}\qquad\partial_\mu F^{\mu\nu}=e\,J^\nu_Q\qquad[\mathrm{Def\text{-}BM.12}(1);\ \mathrm{EM.11}]:
\]

%dsp-level:0
\[
\bigl[\tfrac12E_r^2\cdot V\bigr]=[\text{energy}]=\mathfrak q^{+1/2}\implies[E_r^2]=\mathfrak q^{+2}\implies[F^{\mu\nu}]=\mathfrak q^{+1};
\]

%dsp-level:0
\[
\qquad [\partial_\mu]=\mathfrak q^{+1/2},\quad[e]=\mathfrak q^0\implies\boxed{\ [J^\nu_Q]=\mathfrak q^{+3/2}\ }
\]

%dsp-level:0
\[
\qquad [\text{machine: the exponent arithmetic balances}];
\]

%dsp-level:0
\[
\text{the integrated current}\ [J\cdot\text{area}]=\mathfrak q^{+3/2}\mathfrak q^{-1}=\mathfrak q^{+1/2}=\text{charge}/\text{length}=\text{charge}/\text{time}\ \ \text{on}\ c=1
\]

%dsp-level:0
\[
\qquad [\text{machine}]:
\]

%dsp-level:0
\[
\textbf{the row's adjudication:}
\]

%dsp-level:0
\[
\qquad \text{the ledger-side per-tick read is}\ \mathfrak q^0\ \ \text{and the chart-side read is}\ \mathfrak q^{+1/2}\ —\ \ \text{both derived,}\ \neg\text{equated}\ [\mathrm{BM.37}(2)\ \text{pattern}];
\]

%dsp-level:0
\[
(4)\quad\textbf{the two bridges, separated:}
\]

%dsp-level:0
\[
\qquad \textbf{the dimensional bridge is}\ \mathfrak h:
\]

%dsp-level:0
\[
\qquad m_{\mathrm{hid}}=\frac{E_{\mathrm{int}}\ln2}{\mathfrak h}:\quad\mathfrak q^0\longmapsto\mathfrak q^{+1/2}
\]

%dsp-level:0
\[
\qquad [\mathrm{FT.44}(4);\ \mathrm{BM.35}(4);\ \text{machine}]:
\]

%dsp-level:0
\[
\mathfrak h\ \ \text{is declared by}\ \mathrm{P1}\ \ \text{with}\ \neg\text{a value issued},\qquad\exists!\ \text{such a scale}\ [\mathrm{P1};\ \text{A §0.3}];
\]

%dsp-level:0
\[
\textbf{the numerical bridge is}\ \rho_{\mathrm B}:
\]

%dsp-level:0
\[
\qquad [\rho_{\mathrm B}]
\]

%dsp-level:0
\[
= \mathfrak q^0\implies\boxed{\ \rho_{\mathrm B}\ \ \textbf{converts no dimension}\ —\ \text{it is the dimensionless exchange rate between the bit currency}\ (\ln2\ \text{per directed edge})\ \text{and the angle currency}\ (\Theta_R)\ }
\]

%dsp-level:0
\[
[\mathrm{BM.35}(5)\ \text{Rmk: three currencies, one reference point};\ \mathrm{BM.69}(2);\ \mathrm{BM.98};\ \mathrm{CT1.29};\ \mathrm{FT.54}(4)]:
\]

%dsp-level:0
\[
\textbf{the correction to the framing:}
\]

%dsp-level:0
\[
\qquad \rho_{\mathrm B}\ \ \text{is}\ \neg\textbf{the dimensional bridge}\ —\ \text{that role is}\ \mathfrak h\textbf{'s}\ —\ \ \text{and}\ \neg\text{a single number carries both roles}\ [\text{machine:}\ [\mathfrak h]=\mathfrak q^{-1/2}\ne\mathfrak q^0=[\rho_{\mathrm B}]];
\]

%dsp-level:0
\[
(5)\quad\textbf{the declaration (logged as the §II.4 addendum A-BM.9 before use):}
\]

%dsp-level:0
\[
\qquad \textbf{species:}\ \text{the bridge's dimensionless calibration slot}
\]

%dsp-level:0
\[
\boxed{\ \text{§Cal-BM gains exactly one declared member:}\quad\rho_{\mathrm B}=\frac{\hat\mu'^2\mathfrak h^2}{(\ln2)^2}\in\mathfrak q^0,\qquad\textbf{its value is}\ \neg\textbf{issued}\ }
\]

%dsp-level:0
\[
\textbf{the packaging row:}
\]

%dsp-level:0
\[
\qquad \rho_{\mathrm B}\ \ \text{is a dimensionless combination of two already-declared data}\ —\ \hat\mu'^2\ [\mathrm{Cal\text{-}F};\ \mathrm{Def\text{-}H.37}(2)]\ \ \text{and}\ \ \mathfrak h\ [\mathrm{P1}]:
\]

%dsp-level:0
\[
\neg\text{a new free parameter is introduced},
\]

%dsp-level:0
\[
\qquad \neg\text{a new dimensionful scale}\ [\mathrm{P1}],
\]

%dsp-level:0
\[
\qquad \neg\text{a measured value}\ [\mathrm{P5}],
\]

%dsp-level:0
\[
\qquad \neg\text{a numerical coupling}\ [\mathrm{EM.25}(7),(9)];
\]

%dsp-level:0
\[
\textbf{what the declaration fixes is the datum's status:}
\]

%dsp-level:0
\[
\qquad \text{the bridge's single residue is a declared calibration slot, to be filled by a calibration act the corpus prices and}\ \neg\text{performs}
\]

%dsp-level:0
\[
\qquad [\mathrm{BM.77}(9)\ \text{pattern}];
\]

%dsp-level:0
\[
(6)\quad\textbf{what the slot buys, exactly:}
\]

%dsp-level:0
\[
\qquad \text{with}\ \rho_{\mathrm B}\ \text{supplied, every cross-currency ratio is an exact algebraic number per state:}
\]

%dsp-level:0
\[
\mathcal R_s=\rho_{\mathrm B}\frac{D_S}{4E^2},
\]

%dsp-level:0
\[
\qquad \mathcal R_{\mathrm{mat}}=\rho_{\mathrm B}\frac{D_S\,D^2}{4E^2},
\]

%dsp-level:0
\[
\qquad \mathcal R=\frac{\pi\ln2}2\cdot\frac{\rho_{\mathrm B}D_S}{G}=\frac{\hat\mu'^2D_Sr^2}4(1-f)^2
\]

%dsp-level:0
\[
[\mathrm{BM.69}(2);\ \mathrm{BM.98}(1)\text{--}(4);\ \text{six states, four fields, six minimal polynomials}\ [\mathrm{Def\text{-}BM.14}(3)]]:
\]

%dsp-level:0
\[
\qquad \textbf{the bridge becomes numerical, state by state, on one number};
\]

%dsp-level:0
\[
\textbf{and}\ \neg\textbf{a currency becomes addable:}
\]

%dsp-level:0
\[
\qquad [\Delta_g]=\mathfrak q\ \ \text{against}\ \ [\Delta m]=\mathfrak q^{+1/2}\implies\neg\text{addable}
\]

%dsp-level:0
\[
\qquad [\mathrm{BM.37}(2);\ \mathrm{BM.69}(3);\ \mathrm{BM.90}(5);\ \mathrm{BM.98}(6);\ \text{machine}]:
\]

%dsp-level:0
\[
\textbf{the declaration moves no summand across the two currencies, and the ratio stays a ratio};
\]

%dsp-level:0
\[
(7)\quad\textbf{the mass-form bridge disposition:}
\]

%dsp-level:0
\[
\qquad \boxed{\ \textbf{the relating equation is derived in form}\ [\mathrm{BM.98}]\ \textbf{and its single residue is a declared slot with no value issued}\ \ [(\mathrm{A\text{-}BM.9})]\ \Longrightarrow\ \textbf{the mass-form bridge is closed as an equation}\ }
\]

%dsp-level:0
\[
\textbf{the derivation bar stands:}
\]

%dsp-level:0
\[
\qquad \neg\exists\ \text{a theorem equating the bit currency and the angle currency}\ [\mathrm{BM.35}(5)];
\]

%dsp-level:0
\[
\qquad \mathrm{BM.69}(4)\text{'s four grounds stand unweakened}\ —\ \textbf{the declaration is}\ \neg\textbf{a derivation};
\]

%dsp-level:0
\[
\textbf{the numerical bridge remains a calibration act:}
\]

%dsp-level:0
\[
\qquad \neg\text{a value for}\ \rho_{\mathrm B}\ \text{is derived, declared or fitted here};
\]

%dsp-level:0
\[
\qquad \neg\text{an}\ [\mathrm{EXT}]\ \text{datum fixes it}\ [\mathrm{P5};\ \mathrm{FT.73}(5)\ \text{trap pattern}];
\]

%dsp-level:0
\[
(8)\quad\textbf{the bars:}
\]

%dsp-level:0
\[
\qquad \neg\text{a second dimensionful scale}\ [\mathrm{P1}];
\]

%dsp-level:0
\[
\qquad \alpha_B:\ \neg\text{added},\ \neg\text{derived},\ \text{barred on}\ \mathrm{BM.61}(8)\text{'s four grounds};
\]

%dsp-level:0
\[
\qquad \text{the six}\ D_S\ \text{minimal polynomials are untouched};
\]

%dsp-level:0
\[
\text{species names are confined to the identification layer}\ [\mathrm{Def\text{-}BM.3}(7);\ \mathrm{Def\text{-}BM.6}(7)];
\]

%dsp-level:0
\[
\qquad \text{amplitude}\ne\text{mass}\ [\mathrm{FT.79}(5)];
\]

%dsp-level:0
\[
\qquad \text{rung}\ne\text{scale}\ [\mathrm{D0.8}];
\]

%dsp-level:0
\[
(9)\quad\textbf{status:}
\]

%dsp-level:0
\[
\qquad \textbf{derived:}\ (1)\text{--}(4),\ (6);
\]

%dsp-level:0
\[
\qquad \textbf{declared:}\ (5)\ —\ \text{logged as A-BM.9 before use, one slot, no value};
\]

%dsp-level:0
\[
\qquad \textbf{adjudicated:}\ (3)\text{'s two current reads and the temperature row};
\]

%dsp-level:0
\[
\neg\text{a P4 addition beyond the declared slot},
\]

%dsp-level:0
\[
\quad\neg\text{a new dimensionful scale}\ [\mathrm{P1}],
\]

%dsp-level:0
\[
\quad\neg\text{a measured value}\ [\mathrm{P5}],\qquad \quad
\]

\textbf{Req.} D0.8, D1.8, D2.9, D3.4 (Parts 1–2); A §0.3, CT1.29, CT2.1 (Part 1, Part 2); CT2.9, EM.9, EM.11, EM.25(7), EM.25(9), EW.33, Def-BM.12(1) (Part 2, this part); P1, P4, P5 (Parts 1–2); Def-H.2, Def-H.37(1), Def-H.37(2), FT.44(4), FT.54(4), FT.73(5), FT.79(5) (Part 3); Def-BM.4(1), Def-BM.5(4), Def-BM.13(3), Def-BM.13(6), Def-BM.14(3), BM.14(2), BM.14(4), BM.23(9), BM.35(3), BM.35(4), BM.35(5), BM.37(2), BM.44, BM.61(8), BM.69(2), BM.69(3), BM.69(4), BM.69(6), BM.88(2), BM.88(3), BM.90(5), BM.95(1), BM.98, BM.98(6) (this part); Cal-F; the A-BM.9 declaration (logged before use); exact $\mathfrak q$-exponent arithmetic, Laurent-monomial balances, symbol censuses [math; machine: \texttt{checks/\allowbreak{}verify\_\allowbreak{}H\_\allowbreak{}six\_\allowbreak{}objective\_\allowbreak{}resolution.py} §B].

\textbf{Pf.}

%dsp-level:0
\[
(1):\ \mathrm{D2.9},\ \mathrm{P1},\ \text{A §0.3},\ \mathrm{CT2.9}\ \text{re-cited verbatim (frozen)};
\]

%dsp-level:0
\[
\qquad \neg\exists\ \text{a second dimensionful scale in the inventory}\ [\text{the derived no-go of the second scale re-cited}]
\]

%dsp-level:0
\[
(2):\ \text{every row re-cited from its own locus and its exponent re-computed}\ [\text{machine: the exponent arithmetic}];
\]

%dsp-level:0
\[
\qquad [\rho V]=2-\tfrac32=\tfrac12;\quad[\kappa^2\rho]=-1+2=+1;\quad[(\Delta x_{\mathrm B})^2\kappa^2\rho]=-1-1+2=0\ [\text{machine}]
\]

%dsp-level:0
\[
[\rho_{\mathrm B}]=[\hat\mu'^2]+2[\mathfrak h]-2[\ln2]=1-1-0=0\ [\mathrm{BM.69}(2);\ \text{machine}];
\]

%dsp-level:0
\[
\qquad [\hat\mu^2]=2\ [\mathrm{Def\text{-}H.37}(1);\ \mathrm{EW.33}\ \text{packaging}]
\]

%dsp-level:0
\[
(3):\ [\tfrac12E_r^2V]
\]

%dsp-level:0
\[
= \tfrac12\ \text{carries}\ \mathfrak q^0\ [\mathrm{BM.87}(2)\text{'s exterior integral}]
\]

%dsp-level:0
\[
\implies [E_r^2]
\]

%dsp-level:0
\[
= \tfrac12-\bigl(-\tfrac32\bigr)
\]

%dsp-level:0
\[
= 2\ [\text{machine}];
\]

%dsp-level:0
\[
\qquad [\partial_\mu]=-[x^\mu]=+\tfrac12;
\]

%dsp-level:0
\[
\qquad [eJ]=[\partial_\mu F^{\mu\nu}]=\tfrac12+1=\tfrac32\ \text{and}\ [e]=0\implies[J]=\tfrac32\ [\text{machine}]
\]

%dsp-level:0
\[
[J\cdot\text{area}]=\tfrac32-1=\tfrac12\ [\text{machine}];
\]

%dsp-level:0
\[
\qquad k_B=1\implies[\text{temperature}]=[\text{energy}]=\tfrac12\ [\mathrm{D2.9};\ \mathrm{CT2.1}];
\]

%dsp-level:0
\[
\qquad \neg\exists\ \text{a P4-class temperature}\ [\mathrm{BM.69}(6)\ \text{re-cited}]
\]

%dsp-level:0
\[
(4):\ m_{\mathrm{hid}}=E_{\mathrm{int}}\ln2/\mathfrak h:\ [\cdot]=0-(-\tfrac12)=\tfrac12\ [\text{machine}];
\]

%dsp-level:0
\[
\qquad [\mathfrak h]=-\tfrac12\ne0=[\rho_{\mathrm B}]\implies\neg\text{one number carries both roles}\ [\text{machine}]
\]

%dsp-level:0
\[
(5)\text{--}(6):\ \text{the declaration's species, costs and bars recorded in the registry grammar}\ [\text{A-BM.9};\ \text{machine: the packaging census}\ \{\hat\mu'^2\in\mathrm{Cal\text{-}F},\ \mathfrak h\in\mathrm{P1}\}];
\]

%dsp-level:0
\[
\qquad \text{the six states' exactness per supplied value re-cited}\ [\mathrm{BM.69}(2);\ \mathrm{BM.98}]
\]

%dsp-level:0
\[
(7)\text{--}(9):\ \mathrm{BM.69}(4)\text{'s four grounds and}\ \mathrm{BM.35}(5)\text{'s no-equation row re-cited};
\]

%dsp-level:0
\[
\qquad \text{the standing-rule scans: external tokens, decimal numerals, unit tokens, the barred object}\ \alpha_B\ [\text{machine}].\ \blacksquare
\]

\textbf{Phys.} The corpus has one dimension and this tier counts with it. Length and time are the same power, mass and energy the same power, density and stress-energy two powers up, curvature one, and the single free scale sits at minus one half — so every balance in the census is an integer sum, and the seven rows of the question come out derived where the inventory has a referent and adjudicated where it does not. Temperature is a dimension row and nothing more, because no thermal parameter exists; amount of substance is a count; and electric current has two reads that are not equated, the ledger's weight per tick at exponent zero and the chart's density at three halves, the latter forced by the field equation and the field-energy row. The bridge the question asks about turns out to be two bridges. The dimensional one is the free scale itself: it takes a count of bits to a mass, and it was declared at the outset with no value ever issued. The numerical one is dimensionless, so it converts no dimension at all — it is the exchange rate between the ledger's bit currency and the holonomy's angle currency, and it is the single number the relating equation needs. That number is now declared: one slot in the era's calibration registry, logged before use, packaging two data that were already declared elsewhere, adding no parameter and no scale, and carrying no value. With it supplied, every cross-currency ratio in the corpus becomes an exact algebraic number state by state; without it, nothing is lost, because the four grounds against deriving it stand and no summand crosses between the two currencies. The equation is closed; the number is a calibration act the corpus prices and does not perform.

\end{theorem}

\begin{theorem}[{\texorpdfstring{BM.103 — The three-constituent assembly: the eigenvalue equations that are formed, the tensor-product referent census, and the composition residue}{BM.103 - The three-constituent assembly: the eigenvalue equations that are formed, the tensor-product referent census, and the composition residue}}]\label{tag:BM-103}
\textbf{Thm.}

%dsp-level:0
\[
(1)\quad\textbf{the six-state inventory (re-cited):}
\]

%dsp-level:0
\[
\qquad \mathcal S=\{(\mathcal F_g,\beta):\ g\in\{1,2,3\},\ \beta\in\{-,+\}\},
\]

%dsp-level:0
\[
\qquad |\mathcal S|=6
\]

%dsp-level:0
\[
\qquad [\mathrm{Def\text{-}BM.14}(2)]:
\]

%dsp-level:0
\[
D_S\ \in\ \bigl\{2,\ 4,\ 2-\sqrt2,\ 1,\ 2-2\cos\tfrac\pi{10},\ 2-\sqrt3\bigr\},
\]

%dsp-level:0
\[
\qquad m^2_{\mathrm{pole}}=\hat\mu'^2D_S,
\]

%dsp-level:0
\[
\qquad \text{four fields, six minimal polynomials}
\]

%dsp-level:0
\[
\qquad [\mathrm{Def\text{-}BM.14}(3);\ \mathrm{Census\ FT.63}(6);\ \mathrm{BM.86}(2)]:
\]

%dsp-level:0
\[
x-2,\quad x-4,\quad x^2-4x+2,\quad x-1,\quad x^4-8x^3+19x^2-12x+1,\quad x^2-4x+1\ \ (\text{for}\ 2-\sqrt3)
\]

%dsp-level:0
\[
\qquad [\text{machine: the four quadratic-ring roots re-verified exactly}];
\]

%dsp-level:0
\[
\text{the generation ladder}\ 12:24:60=1:2:5
\]

%dsp-level:0
\[
\qquad [\mathrm{Def\text{-}BM.14}(1);\ \mathrm{BM.60}(1)];
\]

%dsp-level:0
\[
\qquad \Delta_g=\bigl(2,\ \sqrt2-1,\ 2\cos\tfrac\pi{10}-\sqrt3\bigr)\hat\mu'^2>0,
\]

%dsp-level:0
\[
\qquad \Delta_1>\Delta_2>\Delta_3
\]

%dsp-level:0
\[
\qquad [\mathrm{Def\text{-}BM.14}(5);\ \mathrm{BM.69}(3)];
\]

%dsp-level:0
\[
(2)\quad\textbf{the eigenvalue equation that is formed}\ (\textbf{the three-line channel,}\ \neg\textbf{the three constituents}):
\]

%dsp-level:0
\[
\qquad z:=r^2,
\]

%dsp-level:0
\[
\qquad \boxed{\ 4z^3-24z^2+\bigl(21-12\gamma_{\mathrm{ch}}\bigr)z-\bigl(\gamma_{\mathrm{ch}}-1\bigr)^2=0\ }
\]

%dsp-level:0
\[
\textbf{Vieta:}
\]

%dsp-level:0
\[
\qquad \boxed{\ \sum_iz_i=6\ \ (\text{phase-free and channel-free}),\qquad\sum_{i<j}z_iz_j=\frac{21-12\gamma_{\mathrm{ch}}}4,\qquad\prod_iz_i=\frac{(\gamma_{\mathrm{ch}}-1)^2}4\ }
\]

%dsp-level:0
\[
\qquad [\mathrm{BM.80}(3);\ \text{machine: the three symmetric reads}];
\]

%dsp-level:0
\[
\textbf{the depressed cubic on the spectral side:}
\]

%dsp-level:0
\[
\qquad 2(r-1)^3-3(r-1)-\gamma_{\mathrm{ch}}=0\iff2r^3-6r^2+3r+(1-\gamma_{\mathrm{ch}})=0
\]

%dsp-level:0
\[
\qquad [\mathrm{BM.80}(2);\ \mathrm{FT.111}(3);\ \mathrm{FT.143}(2);\ \text{machine: the expansion}];
\]

%dsp-level:0
\[
\gamma_{\mathrm{ch}}=\sqrt2\cos(3\delta_{\mathrm{ch}}),
\]

%dsp-level:0
\[
\qquad \delta_{\mathrm{ch}}\in\Bigl\{\frac29,\ \frac2{27},\ \frac18\Bigr\},
\]

%dsp-level:0
\[
\qquad 27\delta_q=9\delta_\ell=16\delta_\nu=2
\]

%dsp-level:0
\[
\qquad [\mathrm{BM.80}(1);\ \mathrm{FT.143}(5);\ \text{machine}];
\]

%dsp-level:0
\[
\textbf{the elimination row:}
\]

%dsp-level:0
\[
\qquad r(2z+3)=6z+\gamma_{\mathrm{ch}}-1\ \ \text{at}\ \ z=r^2,
\]

%dsp-level:0
\[
\qquad z(2z+3)^2-(6z+\gamma_{\mathrm{ch}}-1)^2\equiv0\ \ \text{in the quotient ring}
\]

%dsp-level:0
\[
\qquad [\mathrm{BM.80}(3);\ \text{machine}];
\]

%dsp-level:0
\[
\boxed{\ \textbf{the three roots are the channel's three lines}\ (\text{the generations}),\ \neg\textbf{the three constituents of a sealed class}\ }
\]

%dsp-level:0
\[
\qquad [\mathrm{Def\text{-}BM.6}(4);\ \mathrm{FT.143}(6)];
\]

%dsp-level:0
\[
(3)\quad\textbf{the characteristic polynomial of the three-constituent assembly in the transport register:}
\]

%dsp-level:0
\[
\qquad W(Q)=\bigl[z^{Q}\bigr]\bigl(2+z+z^{-1}\bigr)^{3}
\]

%dsp-level:0
\[
\qquad [\mathrm{Def\text{-}BM.13}(3)]:
\]

%dsp-level:0
\[
2+z+z^{-1}=z^{-1}(z+1)^2\implies\bigl(2+z+z^{-1}\bigr)^3=\frac{(z+1)^6}{z^3}\qquad[\text{machine: one Laurent identity}]:
\]

%dsp-level:0
\[
\boxed{\ W(Q)\ =\ \binom{6}{Q+3}\ \ \text{exactly, on all seven classes}\ }
\]

%dsp-level:0
\[
\qquad W=(1,6,15,\mathbf{20},15,6,1),
\]

%dsp-level:0
\[
\qquad \sum_QW(Q)=2^6=64=4^3
\]

%dsp-level:0
\[
\qquad [\mathrm{BM.28}(2);\ \text{machine: the closed form against the exhaustive}\ 4^3\ \text{census}];
\]

%dsp-level:0
\[
\textbf{the assembly's transport spectrum:}
\]

%dsp-level:0
\[
\qquad \text{one root}\ z
\]

%dsp-level:0
\[
= -1\ \ \text{of multiplicity six}\ —\ \text{the six directed edges}\ —\ \ \text{and the seven class weights are its binomial coefficients}
\]

%dsp-level:0
\[
\boxed{\ \textbf{the eigenvalue structure the chart resolves is charge-class data,}\ \neg\textbf{flavour data}\ }
\]

%dsp-level:0
\[
\qquad [\mathrm{Def\text{-}BM.13}(6);\ \mathrm{BM.95}(1);\ \text{machine}];
\]

%dsp-level:0
\[
(4)\quad\textbf{the tensor-product referent census (what an ensemble matrix requires against what the corpus supplies):}\qquad\begin{array}{l|l}\text{an}\ M_{\mathrm{ensemble}}\ \text{requires} & \text{the corpus supplies}\\\hline\text{a three-constituent state space} & \neg\text{one:}\ \mathcal H_c,\ \mathcal H_c^{(q)}\ \text{are counting spaces with no inner product, no dynamics, no state, no Hamiltonian}\ [\mathrm{Def\text{-}H.41}\ \text{Species};\ \mathrm{Def\text{-}H.42}];\\ & \neg\text{a two-chain state space exists, let alone a three-chain one}\ [\mathrm{FT.80}(5)]\\\text{an exchange operator} & \neg\exists\ \text{one in the inventory}\ [\mathrm{FT.80}(5);\ \mathrm{FT.125}(5)]\\\text{a flavour index attaching states to constituents} & \neg\text{derived:}\ (\mathrm{R3}),\ \text{deficit}\ 1'\ [\mathrm{BM.28}(7);\ \mathrm{BM.75}(6);\ \mathrm{Def\text{-}BM.8}(5)]\\\text{a triple-level resolution} & \neg\text{available: the end signs carry no position datum}\ [\mathrm{D3.4};\ \mathrm{BM.46}(2);\ \mathrm{Def\text{-}BM.13}(9)]\\\text{a}\ 27\text{-dimensional flavour arena} & \text{rejected as an arena at rung one}\ [\mathrm{FT.142}(3):\ \text{the}\ 27=\text{seal}\times\text{transfer}\times\text{colour candidate is a declared construction, cited and imported nowhere}]\end{array}
\]

%dsp-level:0
\[
\boxed{\ \det\bigl(\lambda I-M_{\mathrm{ensemble}}\bigr)=0\ \ \textbf{is}\ \neg\textbf{formed}\ —\ \text{the question ``which three states compose a sealed class'' has no eigenvalue referent in this corpus}\ }
\]

%dsp-level:0
\[
(5)\quad\textbf{what is derived about the composition:}
\]

%dsp-level:0
\[
\qquad \textbf{exactly three constituents seal:}
\]

%dsp-level:0
\[
\quad B:=\varepsilon_{ijk}q^iq^jq^k\ [\mathrm{D3.9}];\quad\varepsilon_{ijk}\ \text{seals}\ 3\ \text{fundamentals}\ [\mathrm{CT3.15};\ \mathrm{CT3.17}];
\]

%dsp-level:0
\[
\quad 3\ \text{admits no}\ \{2,3\}\text{-partition with an external remainder}\ [\mathrm{FT.34}(5);\ \mathrm{Census\ FT.35}(5)]:
\]

%dsp-level:0
\[
\textbf{their core charges sum to the class read:}\quad Q=\sum_iq_{\mathrm{core}}(c_i)\in\{-3,\dots,+3\},
\]

%dsp-level:0
\[
\qquad Q_{\mathrm{EM}}=0\ \ \text{in every valid universe}
\]

%dsp-level:0
\[
\qquad [\mathrm{Def\text{-}BM.8}(1);\ \mathrm{EM.11};\ \mathrm{D2.11}];
\]

%dsp-level:0
\[
\textbf{all seven classes are formed with the binomial weights of}\ (3);
\]

%dsp-level:0
\[
\qquad \textbf{the neutral class outweighs either charged class by exactly}\ \tfrac43\ \ \text{on route}\ (\mathrm{ii})
\]

%dsp-level:0
\[
\qquad [\mathrm{BM.28}(4)(\mathrm{ii});\ \mathrm{BM.30}(2)];
\]

%dsp-level:0
\[
\textbf{the fractional constituent alphabet is structurally impossible}
\]

%dsp-level:0
\[
\qquad [\mathrm{BM.28}(7);\ \mathrm{Def\text{-}BM.8}(5)];
\]

%dsp-level:0
\[
\qquad \textbf{colour is}\ \neg\textbf{observable by construction}
\]

%dsp-level:0
\[
\qquad [\mathrm{CT3.22}];
\]

%dsp-level:0
\[
\boxed{\ \textbf{the ``three'' is derived from the seal grammar and the ``which three'' is not}\ }
\]

%dsp-level:0
\[
(6)\quad\textbf{the stability question, adjudicated:}\qquad\neg\exists\ \text{a derived stability condition on eigenvalues:}
\]

%dsp-level:0
\[
(\mathrm i)\quad\text{the materialization filter is an undeclared spectral hypothesis}\ (\mathrm{M1})/(\mathrm{M2})
\]

%dsp-level:0
\[
\qquad [\mathrm{Def\text{-}H.20};\ \text{A §0.2a};\ \text{Part 3 §0.2}(\mathrm{iii})];
\]

%dsp-level:0
\[
(\mathrm{ii})\quad\text{the index-zero selection route is refuted for bound chains: on Synge-class slices every closed geodesic has}\ \operatorname{ind}_B
\]

%dsp-level:0
\[
\ge 1,
\]

%dsp-level:0
\[
\ \ \text{so absolute index-}0\ \text{selection is inadmissible against}\ (\mathrm{M1})
\]

%dsp-level:0
\[
\qquad [\mathrm{FT.66}(4)];
\]

%dsp-level:0
\[
(\mathrm{iii})\quad\text{a mass-ordering criterion would be a count}\to\text{strength coupling}
\]

%dsp-level:0
\[
\qquad [\mathrm{BM.36}(3);\ \mathrm{BM.61}(8)(\mathrm{O1});\ \mathrm{BM.90}(3)]:
\]

%dsp-level:0
\[
\textbf{the formed stability rows are kinematic and ledger-side:}
\]

%dsp-level:0
\[
\qquad \text{the two-body impossibility}\ [\mathrm{BM.41}(2)];
\]

%dsp-level:0
\[
\quad\text{the two-bit floor}\ [\mathrm{BM.70}(2)];
\]

%dsp-level:0
\[
\quad\text{the materialization restriction}\ E
\]

%dsp-level:0
\[
\ge 2\ [\mathrm{BM.94}(5)];
\]

%dsp-level:0
\[
\text{the valid-universe conditions}\ [\mathrm{D2.11}];
\]

%dsp-level:0
\[
\quad\text{the no-flicker theorem}\ \Gamma(\varepsilon_q{=}0)
\]

%dsp-level:0
\[
= 0\ [\mathrm{BM.23}(8)];
\]

%dsp-level:0
\[
\quad\text{the two-sided bar}\ [\mathrm{BM.99}(3)];
\]

%dsp-level:0
\[
\quad\text{the forward-only arrow}\ [\mathrm{A6};\ \mathrm{FT.144}(2)];
\]

%dsp-level:0
\[
(7)\quad\textbf{the octet row:}
\]

%dsp-level:0
\[
\qquad \neg\text{formed}.
\]

%dsp-level:0
\[
\qquad \textbf{the corpus's derived}\ 3\ \ \text{is the colour seal's fundamental}\ [\mathrm{CT3.15};\ \mathrm{CT3.17};\ \mathrm{D3.9}];
\]

%dsp-level:0
\[
\textbf{its derived}\ 8\ \ \text{is a mask-capacity ratio, barred as a coupling}\ [\mathrm{BM.36}(2),(3);\ \mathrm{BM.61}(8)(\mathrm{O1})];
\]

%dsp-level:0
\[
\qquad \textbf{its derived}\ 10\ \ \text{is the colour grammar's grouping count}\ g(5)=10\ [\mathrm{CT3.17a}(2);\ \mathrm{Def\text{-}BM.9};\ \mathrm{BM.31}];
\]

%dsp-level:0
\[
\textbf{the flavour-side symmetry is}\ C_3,
\]

%dsp-level:0
\[
\ \ \text{its orbits the three channel lines}\ [\mathrm{FT.111}(2);\ \mathrm{FT.158}(1),(4);\ \mathrm{FT.142}(5)];
\]

%dsp-level:0
\[
\qquad \neg\exists\ \text{a flavour}\ \mathsf{SU}(3)\ \text{in the inventory};
\]

%dsp-level:0
\[
\boxed{\ \textbf{the baryon octet and the meson octet are}\ \neg\textbf{derived and}\ \neg\textbf{derivable here}\ —\ \text{what is derived is the}\ \varepsilon\text{-seal's three, the binomial seven, and the}\ C_3\ \text{triplet}\ }
\]

%dsp-level:0
\[
(8)\quad\textbf{the naming bar:}
\]

%dsp-level:0
\[
\qquad \text{the two heaviest-floor triples the question names are external species vocabulary};
\]

%dsp-level:0
\[
\qquad \text{the identification layer supplies}\ \textit{up-type}/\textit{down-type}\ \text{branch names}\ [\text{A-BM.6};\ \mathrm{Def\text{-}BM.14}(8)]\ \text{and generation indices},
\]

%dsp-level:0
\[
\text{and}\ \textbf{no name attaches to a composition}
\]

%dsp-level:0
\[
\qquad [\mathrm{Def\text{-}BM.6}(5),(7);\ \mathrm{Def\text{-}BM.3}(7)]:
\]

%dsp-level:0
\[
\qquad \boxed{\ \neg\text{either triple is formed, and}\ \neg\text{either is excluded}\ —\ \text{the composition is unread}\ (\mathrm{R3})\ }
\]

%dsp-level:0
\[
(9)\quad\textbf{status:}
\]

%dsp-level:0
\[
\qquad \textbf{derived:}\ (1)\text{--}(3),\ (5),\ (6)\text{'s formed rows},\ (7)\text{'s dimension rows};
\]

%dsp-level:0
\[
\qquad \textbf{refuted as formed:}\ \text{the ensemble matrix}\ (4),
\]

%dsp-level:0
\[
\ \text{an eigenvalue stability condition}\ (6),
\]

%dsp-level:0
\[
\ \text{the octet read}\ (7);
\]

%dsp-level:0
\[
\textbf{standing:}\ (\mathrm{R3}),\ \text{deficit}\ 1';
\]

%dsp-level:0
\[
\qquad \neg\text{a declaration},
\]

%dsp-level:0
\[
\quad\neg\text{a P4 addition},
\]

%dsp-level:0
\[
\quad\neg\text{a new dimensionful scale}\ [\mathrm{P1}],
\]

%dsp-level:0
\[
\quad\neg\text{a measured value}\ [\mathrm{P5}],
\]

%dsp-level:0
\[
\quad\neg\text{a species identification}\ [\mathrm{P5};\ \mathrm{Def\text{-}BM.3}(7)];
\]

\textbf{Req.} D2.11, D3.4, D3.9 (Parts 1–2); A §0.2a, CT3.15, CT3.17, CT3.17a(2), CT3.22, EM.11 (Part 2); P1, P4, P5 (Parts 1–2); Def-H.20, Def-H.41, Def-H.42, FT.34(5), FT.66(4), FT.80(5), FT.111(2), FT.111(3), FT.125(5), FT.142(3), FT.142(5), FT.143(2), FT.143(5), FT.143(6), FT.158(1), FT.158(4) (Part 3); Census FT.35(5), FT.63(6) (row-wise); Def-BM.6(4), Def-BM.6(5), Def-BM.6(7), Def-BM.8(1), Def-BM.8(5), Def-BM.9, Def-BM.13(3), Def-BM.13(6), Def-BM.13(9), Def-BM.14(1), Def-BM.14(2), Def-BM.14(3), Def-BM.14(5), Def-BM.14(8), Def-BM.3(7), BM.23(8), BM.28(2), BM.28(4), BM.28(7), BM.30(2), BM.31, BM.36(2), BM.36(3), BM.41(2), BM.46(2), BM.60(1), BM.61(8), BM.69(3), BM.70(2), BM.75(6), BM.80(1), BM.80(2), BM.80(3), BM.86(2), BM.90(3), BM.94(5), BM.95(1), BM.99(3) (this part); exact Laurent and polynomial identities over $\mathbb Q[z]$, binomial closed forms against exhaustive enumeration, Vieta reads, exact radical minimal polynomials in $\mathbb Q(\sqrt2)$ and $\mathbb Q(\sqrt3)$, elimination in a quotient ring [math; machine: \texttt{checks/\allowbreak{}verify\_\allowbreak{}H\_\allowbreak{}six\_\allowbreak{}objective\_\allowbreak{}resolution.py} §C].

\textbf{Pf.}

%dsp-level:0
\[
(1):\ \mathrm{Def\text{-}BM.14}(1)\text{--}(5)\ \text{and}\ \text{Census FT.63}(6)\ \text{re-cited (frozen)};
\]

%dsp-level:0
\[
\qquad (2-\sqrt2)^2-4(2-\sqrt2)+2=0\ \text{in}\ \mathbb Q(\sqrt2);
\]

%dsp-level:0
\[
\qquad (2-\sqrt3)^2-4(2-\sqrt3)+1=0\ \text{in}\ \mathbb Q(\sqrt3)\ [\text{machine: exact}]
\]

%dsp-level:0
\[
2-2\cos\tfrac\pi{10}=y\ \text{satisfies}\ y^4-8y^3+19y^2-12y+1=0:
\]

%dsp-level:0
\[
\qquad (2-y)^2
\]

%dsp-level:0
\[
= \tfrac{5+\sqrt5}2\implies\sqrt5
\]

%dsp-level:0
\[
= 2y^2-8y+3
\]

%dsp-level:0
\[
\implies 5
\]

%dsp-level:0
\[
= (2y^2-8y+3)^2
\]

%dsp-level:0
\[
\implies 4\bigl(y^4-8y^3+19y^2-12y+1\bigr)
\]

%dsp-level:0
\[
= 0\ [\text{machine: the expansion in}\ \mathbb Q[y]]
\]

%dsp-level:0
\[
(2):\ \mathrm{BM.80}(2),(3)\text{'s rows re-cited};
\]

%dsp-level:0
\[
\qquad \text{Vieta:}\ \sum_iz_i
\]

%dsp-level:0
\[
= \tfrac{24}4
\]

%dsp-level:0
\[
= 6,
\]

%dsp-level:0
\[
\ \ \sum_{i<j}z_iz_j
\]

%dsp-level:0
\[
= \tfrac{21-12\gamma}4,
\]

%dsp-level:0
\[
\ \ \prod_iz_i
\]

%dsp-level:0
\[
= -\tfrac{-(\gamma-1)^2}4
\]

%dsp-level:0
\[
= \tfrac{(\gamma-1)^2}4\ [\text{machine: the coefficient reads}]
\]

%dsp-level:0
\[
2(r-1)^3-3(r-1)=2r^3-6r^2+3r-1\implies2r^3-6r^2+3r+(1-\gamma)=0\ [\text{machine}];
\]

%dsp-level:0
\[
\qquad z(2z+3)^2-(6z+\gamma-1)^2
\]

%dsp-level:0
\[
\equiv 0\ \text{modulo}\ 2r^3-6r^2+3r+1-\gamma\ \text{at}\ r(2z+3)
\]

%dsp-level:0
\[
= 6z+\gamma-1\ [\text{machine: the quotient-ring reduction}]
\]

%dsp-level:0
\[
(3):\ 2+z+z^{-1}
\]

%dsp-level:0
\[
= z^{-1}(z^2+2z+1)
\]

%dsp-level:0
\[
= z^{-1}(z+1)^2
\]

%dsp-level:0
\[
\implies (2+z+z^{-1})^3
\]

%dsp-level:0
\[
= z^{-3}(z+1)^6\ [\text{machine: the Laurent identity}];
\]

%dsp-level:0
\[
\qquad [z^Q]z^{-3}(z+1)^6=[z^{Q+3}](z+1)^6=\binom6{Q+3}\ [\text{machine: all seven classes}]
\]

%dsp-level:0
\[
\binom6{Q+3}
\]

%dsp-level:0
\[
= (1,6,15,20,15,6,1)\ \text{against the exhaustive}\ 4^3\ \text{end-sign census}\ [\mathrm{BM.28}(2)\ \text{re-derived};\ \text{machine}];
\]

%dsp-level:0
\[
\qquad \sum_{k=0}^{6}\binom6k=2^6=64\ [\text{machine}]
\]

%dsp-level:0
\[
(4):\ \text{the five referent rows re-cited from their own loci}\ [\mathrm{Def\text{-}H.41}\ \text{Species};\ \mathrm{Def\text{-}H.42};\ \mathrm{FT.80}(5);\ \mathrm{FT.125}(5);\ \mathrm{BM.28}(7);\ \mathrm{Def\text{-}BM.13}(9);\ \mathrm{FT.142}(3)];
\]

%dsp-level:0
\[
\qquad \text{the inventory scans}\ [\text{machine}]
\]

%dsp-level:0
\[
(5):\ \mathrm{D3.9},\ \mathrm{CT3.15},\ \mathrm{CT3.17},\ \mathrm{FT.34}(5),\ \text{Census FT.35}(5)\ \text{re-cited};
\]

%dsp-level:0
\[
\qquad W(0)/W(+1)=20/15=4/3\ [\text{machine}];
\]

%dsp-level:0
\[
\qquad (6): (\mathrm i)\text{--}(\mathrm{iii})\ \text{re-cited at their own loci}\ [\mathrm{Def\text{-}H.20};\ \mathrm{FT.66}(4);\ \mathrm{BM.36}(3);\ \mathrm{BM.61}(8)(\mathrm{O1});\ \mathrm{BM.90}(3)]
\]

%dsp-level:0
\[
(7):\ g(5)=10\ \text{from}\ \mathrm{CT3.17a}(2)\text{'s summation}\ [\mathrm{Def\text{-}BM.9}\ \text{re-cited};\ \text{machine}];
\]

%dsp-level:0
\[
\qquad \text{the capacity ratios}\ \bigl(\tfrac13,8,\tfrac1{25}\bigr)\ \text{from the flag}\ (6,16,17,25)\ [\mathrm{BM.36}(2);\ \text{machine}]
\]

%dsp-level:0
\[
(8)\text{--}(9):\ \text{the naming-scope rows re-cited}\ [\mathrm{Def\text{-}BM.3}(7);\ \mathrm{Def\text{-}BM.6}(5),(7);\ \text{A-BM.6}];
\]

%dsp-level:0
\[
\qquad \text{the standing-rule scans}\ [\text{machine}].\ \blacksquare
\]

\textbf{Phys.} Ask the chart for the spectrum of a three-quark object and the honest answer is a table of what exists and what does not. What exists is two eigenvalue equations, and neither is the one the question wants. The first is the family cubic: four z cubed minus twenty-four z squared plus a channel constant, whose three roots are the three lines of a channel, with a phase-free sum of six and an exact product — a spectrum of generations, not of constituents. The second is the transport assembly's own characteristic polynomial, which collapses to a binomial: two plus z plus one over z, cubed, is z plus one to the sixth over z cubed, so the seven class weights are the binomial coefficients of six, and the assembly's spectrum is a single root of multiplicity six, one per directed edge. That is charge-class data, and it is the deepest eigenvalue statement the chart supports. What does not exist is the object the question names: a three-constituent state space, an exchange operator, a flavour index, a triple-level resolution — the counting spaces the corpus declares have no state on them, and the twenty-seven-dimensional candidate was rejected as an arena at rung one. So the determinant is not formed, and the two triples the question asks about are neither formed nor excluded; the composition is unread, and that is a named residue rather than an oversight. What is derived instead is sharper: the seal grammar forces exactly three, the charge balance forces their sum, the binomial forces the seven classes and the four-thirds lead of the neutral one, and the stability question lands where the corpus has always landed — on an undeclared spectral filter and a refuted index-zero route. The octet is not derivable either: the three here is colour's, the eight is a capacity ratio barred as a coupling, the ten is a grammar count, and the flavour symmetry that does exist is cyclic of order three.

\end{theorem}

\begin{theorem}[{\texorpdfstring{BM.104 — The curvature composition: the back-reaction fixed point, its exact roots and pole locus, and the non-selection theorem}{BM.104 - The curvature composition: the back-reaction fixed point, its exact roots and pole locus, and the non-selection theorem}}]\label{tag:BM-104}
\textbf{Thm.}

%dsp-level:0
\[
(1)\quad\textbf{the loop, formed exactly:}
\]

%dsp-level:0
\[
\qquad m_{\mathrm{corr}}=m_{\mathrm{hid}}+E_{\mathrm{EM}}^{\mathrm{self}}=m_{\mathrm{hid}}\bigl(1+a\,\mathsf g\bigr),
\]

%dsp-level:0
\[
\qquad a:=\frac12\,N\,w\,\varepsilon_{\mathrm{EM}}\,Q^2
\]

%dsp-level:0
\[
\qquad [\mathrm{BM.89}(5);\ \mathrm{Def\text{-}BM.13}(5);\ \text{machine}]:
\]

%dsp-level:0
\[
\rho_{\mathrm{corr}}=\frac{m_{\mathrm{corr}}}{V_{\mathrm{eff}}},
\]

%dsp-level:0
\[
\qquad V_{\mathrm{eff}}=(\Delta x_{\mathrm B})^3,
\]

%dsp-level:0
\[
\qquad \Delta x_{\mathrm B}\cdot m_{\mathrm{hid}}=1
\]

%dsp-level:0
\[
\qquad [\mathrm{Def\text{-}BM.4}(1);\ \mathrm{BM.14}(2);\ \text{the width is the informational one}]:
\]

%dsp-level:0
\[
\boxed{\ G_{\mathrm{corr}}\ :=\ (\Delta x_{\mathrm B})^2\kappa^2\rho_{\mathrm{corr}}\ =\ G\,\bigl(1+a\,\mathsf g\bigr),\qquad G=2\pi E_{\mathrm{int}}^2\ln2\ }
\]

%dsp-level:0
\[
\qquad [\mathrm{BM.88}(3);\ \mathrm{BM.16}(1);\ \text{machine: one monomial}];
\]

%dsp-level:0
\[
\textbf{the consistency condition is the fixed point:}\qquad\boxed{\ \mathsf g\ =\ \mathsf g\bigl(G\,(1+a\,\mathsf g)\bigr)\ }
\]

%dsp-level:0
\[
(2)\quad\textbf{the three exact solution rows:}
\]

%dsp-level:0
\[
(\mathrm i)\quad\mathsf g\equiv1:\qquad1=1\ \ \text{identically}\ —\ \ \text{a fixed point at every}\ a\ge0\ \ \text{and every}\ G\qquad[\text{machine}];
\]

%dsp-level:0
\[
(\mathrm{ii})\quad\mathsf g=1+G:
\]

%dsp-level:0
\[
\qquad \mathsf g=1+G(1+a\mathsf g)\iff\mathsf g\,(1-aG)=1+G\iff\boxed{\ \mathsf g_{\mathrm{fp}}=\frac{1+G}{1-aG}\ },
\]

%dsp-level:0
\[
\qquad \text{admissible}\iff aG<1,
\]

%dsp-level:0
\[
\qquad \mathsf g_{\mathrm{fp}}>1+G
\]

%dsp-level:0
\[
[\text{machine: substitution verifies the fixed point identically over}\ \mathbb Q[G,a]]
\]

%dsp-level:0
\[
\qquad \textbf{the back-reaction inflates the composition and diverges at the pole};
\]

%dsp-level:0
\[
(\mathrm{iii})\quad\mathsf g=(1+G)^{-1}:\qquad\mathsf g\bigl(1+G+aG\,\mathsf g\bigr)=1\iff\boxed{\ aG\,\mathsf g^2+(1+G)\,\mathsf g-1=0\ }
\]

%dsp-level:0
\[
\text{the product of the two roots is}\ -\tfrac1{aG}<0\implies\textbf{exactly one positive root}
\]

%dsp-level:0
\[
\qquad \mathsf g_{\mathrm{fp}}=\frac{\sqrt{(1+G)^2+4aG}-(1+G)}{2aG}\ \ \longrightarrow\ \ \frac1{1+G}\ \ \text{as}\ \ a\to0
\]

%dsp-level:0
\[
\qquad [\text{machine: the exact root row and the limit}];
\]

%dsp-level:0
\[
\boxed{\ \textbf{the fixed-point equation has a unique admissible solution for each of the three priced compositions}\ —\ \text{consistency is satisfied three times over and selects nothing}\ }
\]

%dsp-level:0
\[
(3)\quad\textbf{the pole locus:}
\]

%dsp-level:0
\[
\qquad aG=1\iff\frac12Nw\,\varepsilon_{\mathrm{EM}}Q^2\cdot2\pi E_{\mathrm{int}}^2\ln2=1\iff\boxed{\ \pi\,N\,w\,\varepsilon_{\mathrm{EM}}\,Q^2\,E_{\mathrm{int}}^2\ln2\ =\ 1\ }
\]

%dsp-level:0
\[
\qquad [\text{machine: one monomial condition}];
\]

%dsp-level:0
\[
Q=0\implies aG=0\ne1
\]

%dsp-level:0
\[
\qquad [\text{machine}]\implies\boxed{\ \textbf{the neutral class never reaches the pole}\ }
\]

%dsp-level:0
\[
\qquad [\mathrm{BM.17}(2):\ \text{the neutron row is}\ Q=0;\ \mathrm{BM.89}(5):\ E_{\mathrm{EM}}^{\mathrm{self}}(0)=0];
\]

%dsp-level:0
\[
Q=\pm1,\ E_{\mathrm{int}}=6:
\]

%dsp-level:0
\[
\qquad \varepsilon_{\mathrm{EM}}^{\mathrm{pole}}=\frac1{36\,\pi\,N\,w\ln2}
\]

%dsp-level:0
\[
\qquad [\text{machine: exact}];
\]

%dsp-level:0
\[
\qquad E_{\mathrm{int}}\in\{2,6,12,24,60\}\ \ \text{on the computed cells}
\]

%dsp-level:0
\[
\qquad [\mathrm{BM.13}(2);\ \mathrm{Census\ FT.35}(2),\ \mathrm{FT.42}(5)];
\]

%dsp-level:0
\[
(4)\quad\textbf{the loop is excluded by the lock's own freedom census:}
\]

%dsp-level:0
\[
\qquad G\text{'s derived alphabet is}\ \{E_{\mathrm{int}},\ \ln2,\ \pi\}\ —\ \ \mathfrak h\text{-free},\ \Lambda\text{-free},\ \mathrm{P4}\text{-free}
\]

%dsp-level:0
\[
\qquad [\mathrm{BM.88}(3);\ \mathrm{BM.16}(1);\ \text{machine: the symbol census}]:
\]

%dsp-level:0
\[
a\text{'s alphabet is}\ \{N,\ w,\ \varepsilon_{\mathrm{EM}},\ Q\},
\]

%dsp-level:0
\[
\qquad \varepsilon_{\mathrm{EM}}=e^2/(4\pi)\ \ \text{is a dimensionless P4-class symbol}
\]

%dsp-level:0
\[
\qquad [\mathrm{Def\text{-}BM.12}\ \text{local notation};\ \mathrm{EM.25}(9);\ \mathrm{P4}]:
\]

%dsp-level:0
\[
\boxed{\ \{E_{\mathrm{int}},\ln2,\pi\}\ \cap\ \{N,\ w,\ \varepsilon_{\mathrm{EM}},\ Q\}\ =\ \varnothing\ }
\]

%dsp-level:0
\[
\qquad [\text{machine}]\implies\textbf{admitting the back-reaction would make the lock's argument P4-dependent, contradicting}\ \mathrm{BM.88}(3)\text{'s derived row}
\]

%dsp-level:0
\[
\boxed{\ \textbf{the lock's argument is the informational mass, the EM self-energy does not re-enter}\ G,\ \text{and}\ \neg\textbf{a self-consistency condition on}\ \mathsf g\ \textbf{is formed}\ }
\]

%dsp-level:0
\[
(5)\quad\textbf{the non-selection theorem:}
\]

%dsp-level:0
\[
\qquad \boxed{\ \neg\exists\ \text{a derived}\ \mathsf g(G)\ }
\]

%dsp-level:0
\[
\qquad \text{on four grounds, each re-cited at its own locus:}
\]

%dsp-level:0
\[
(\mathrm i)\quad\neg\exists\ \text{a count}\to\text{strength theorem}
\]

%dsp-level:0
\[
\qquad [\mathrm{BM.36}(3);\ \mathrm{BM.61}(8)(\mathrm{O1});\ \mathrm{BM.69}(4)(\mathrm{iv})];
\]

%dsp-level:0
\[
\qquad (\mathrm{ii})\quad\text{the four grounds that bar a defect-conditioned weighting bar a curvature-conditioned one at the same species}
\]

%dsp-level:0
\[
[\mathrm{BM.90}(3)(\mathrm i)\text{--}(\mathrm{iv}):\ \text{the degree-}16\ \text{field against integer and binary reads}\ [\mathrm{FT.73}(2),(3)];\ \text{the two-torsion image}\ \{0,\pi\}\ [\mathrm{BM.67}(6)];\ \text{the calibration weight}\ {-}2\ \text{against}\ 0\ [\mathrm{FT.55}(1);\ \mathrm{BM.67}(8)];\ \text{the strength bar}];
\]

%dsp-level:0
\[
(\mathrm{iii})\quad\text{the consistency condition of}\ (2)\ \text{is satisfied by all three candidates}\ —\ \text{it is necessary and}\ \neg\text{sufficient}
\]

%dsp-level:0
\[
\qquad [\text{machine}];
\]

%dsp-level:0
\[
(\mathrm{iv})\quad\text{a selection would be a P4-class datum}\ —\ \text{§Cal-BM would gain a second member}\ —\ \text{refused}
\]

%dsp-level:0
\[
\qquad [\mathrm{P4};\ \mathrm{P1};\ \mathrm{BM.88}(4):\ \text{priced three ways, adopted in none}];
\]

%dsp-level:0
\[
(6)\quad\textbf{the composition is}\ \neg\textbf{inert inside the corpus:}
\]

%dsp-level:0
\[
\qquad \mathrm{BM.101}(5)\text{'s crossing count differs by composition}\ —\ (\mathrm i)\ \text{at most one},
\]

%dsp-level:0
\[
\ (\mathrm{ii})\ \text{exactly one always},
\]

%dsp-level:0
\[
\ (\mathrm{iii})\ \text{two, one or none}:
\]

%dsp-level:0
\[
\boxed{\ \textbf{the priced composition is observable in the energy-scale graph's topology, and is still}\ \neg\textbf{selected}\ }
\]

%dsp-level:0
\[
\qquad \textbf{the discrimination is derived; the selection is not};
\]

%dsp-level:0
\[
\textbf{every row of the corpus holds uniformly over the priced set:}
\]

%dsp-level:0
\[
\qquad \mathrm{BM.91}(2),(4)\text{--}(8),\ \mathrm{BM.92}(1),(2),\ \mathrm{BM.98}(1)\text{--}(4),\ \mathrm{BM.99}(1)\text{--}(6),\ \mathrm{BM.101}(2),(3),(7)
\]

%dsp-level:0
\[
\qquad [\text{machine: the uniformity censuses}];
\]

%dsp-level:0
\[
(7)\quad\textbf{the running coupling's status:}\qquad e(r,E,N)=e_0\sqrt{N\,f\bigl(r/\Delta x_{\mathrm B}(E)\bigr)\,\mathsf g\bigl(G(E)\bigr)}\qquad[\mathrm{BM.88}(4)]:
\]

%dsp-level:0
\[
f\ \ \text{derived exactly}\ [\mathrm{BM.88}(2)(\mathrm{ii})];
\]

%dsp-level:0
\[
\qquad G\ \ \text{derived exactly, one monomial}\ [\mathrm{BM.88}(3)];
\]

%dsp-level:0
\[
\qquad N\ \ \text{solution data}\ [\mathrm{CT2.7};\ \mathrm{BM.34}(2);\ \mathrm{BM.88}(8)];
\]

%dsp-level:0
\[
\mathsf g\ \ \text{priced three ways, adopted in none}\ (5);
\]

%dsp-level:0
\[
\qquad e_0\ \ \text{open}\ [\mathrm{EM.25}(7)]:
\]

%dsp-level:0
\[
\qquad \boxed{\ \textbf{the coupling is complete in form and carries exactly two undeclared members}\ (e_0,\ \mathsf g),\ \neg\text{a third}\ }
\]

%dsp-level:0
\[
(8)\quad\textbf{status:}
\]

%dsp-level:0
\[
\qquad \textbf{derived:}\ (1)\text{--}(4),\ (6)\text{'s discrimination rows};
\]

%dsp-level:0
\[
\qquad \textbf{refuted as derivable:}\ \mathsf g(G)\ (5);
\]

%dsp-level:0
\[
\qquad \textbf{rescoped,}\ \neg\textbf{retracted:}\ \mathrm{BM.88}(4),(8),\ \mathrm{BM.89}(5),\ \mathrm{BM.91},\ \mathrm{BM.98}\ \text{stand at their own loci};
\]

%dsp-level:0
\[
\neg\text{a declaration},
\]

%dsp-level:0
\[
\quad\neg\text{a P4 addition},
\]

%dsp-level:0
\[
\quad\neg\text{a new dimensionful scale}\ [\mathrm{P1}],
\]

%dsp-level:0
\[
\quad\neg\text{a measured value}\ [\mathrm{P5}],
\]

%dsp-level:0
\[
\quad\neg\text{a numerical coupling}\ [\mathrm{EM.25}(7),(9)];
\]

%dsp-level:0
\[
\alpha_B:\ \neg\text{added},\ \neg\text{derived},\ \text{barred}\ [\mathrm{BM.61}(8)];\qquad
\]

\textbf{Req.} D2.9 (Parts 1–2); CT2.1, CT2.7, EM.25(7), EM.25(9) (Part 2); P1, P4, P5 (Parts 1–2); Def-H.37(2), FT.54(4), FT.55(1), FT.73(2), FT.73(3) (Part 3); Def-BM.4(1), Def-BM.12, BM.13(2), BM.14(2), BM.16(1), BM.17(2), BM.34(2), BM.36(3), BM.61(8), BM.67(6), BM.67(8), BM.69(4), BM.88(2), BM.88(3), BM.88(4), BM.88(5), BM.88(8), BM.89(5), BM.90(3), BM.91, BM.92(1), BM.92(2), BM.98, BM.101(5), BM.101(7), Def-BM.13(5) (this part); Census FT.35(2), FT.42(5) (row-wise); exact polynomial arithmetic over $\mathbb Q[G,a]$, quadratic root rows with sign censuses, exact limit rows, symbol censuses and set intersections, monomial conditions [math; machine: \texttt{checks/\allowbreak{}verify\_\allowbreak{}H\_\allowbreak{}six\_\allowbreak{}objective\_\allowbreak{}resolution.py} §D].

\textbf{Pf.}

%dsp-level:0
\[
(1):\ E_{\mathrm{EM}}^{\mathrm{self}}=\tfrac12Nw\,\mathsf g\,\varepsilon_{\mathrm{EM}}Q^2m_{\mathrm{hid}}=a\,\mathsf g\,m_{\mathrm{hid}}\ [\mathrm{BM.89}(5);\ \text{machine}];
\]

%dsp-level:0
\[
\qquad \rho_{\mathrm{corr}}=m_{\mathrm{corr}}/V_{\mathrm{eff}}=\rho\,(1+a\mathsf g)\ \text{at fixed}\ V_{\mathrm{eff}}\ [\mathrm{Def\text{-}BM.4}(1);\ \text{machine}]
\]

%dsp-level:0
\[
G_{\mathrm{corr}}=(\Delta x_{\mathrm B})^2\kappa^2\rho_{\mathrm{corr}}=G(1+a\mathsf g)\ [\text{machine: one monomial}];
\]

%dsp-level:0
\[
\qquad \text{the consistency condition is}\ \mathsf g=\mathsf g(G_{\mathrm{corr}})\ [\text{the fixed point, stated}]
\]

%dsp-level:0
\[
(2)(\mathrm i):\ 1=1\ [\text{machine}];
\]

%dsp-level:0
\[
\qquad (\mathrm{ii}):\ \mathsf g=1+G+aG\mathsf g\iff\mathsf g(1-aG)=1+G\ [\text{machine: cleared over}\ \mathbb Q[G,a]];
\]

%dsp-level:0
\[
\mathsf g_{\mathrm{fp}}>1+G\iff\tfrac{1+G}{1-aG}>1+G\iff\tfrac1{1-aG}>1
\]

%dsp-level:0
\[
\iff 0<aG<1\ [\text{machine: the sign census on positive witnesses}]
\]

%dsp-level:0
\[
(\mathrm{iii}):\ \mathsf g=\tfrac1{1+G+aG\mathsf g}\iff\mathsf g(1+G)+aG\mathsf g^2=1\ [\text{machine}];
\]

%dsp-level:0
\[
\qquad \text{the root product}\ -1/(aG)<0\implies\text{one positive root}\ [\text{machine}];
\]

%dsp-level:0
\[
\qquad a\to0:\ \tfrac{\sqrt{(1+G)^2+4aG}-(1+G)}{2aG}
\]

%dsp-level:0
\[
= \tfrac{4aG}{2aG\bigl(\sqrt{(1+G)^2+4aG}+(1+G)\bigr)}\to\tfrac1{1+G}\ [\text{machine: the conjugate-cleared limit}]
\]

%dsp-level:0
\[
(3):\ aG=\tfrac12Nw\,\varepsilon_{\mathrm{EM}}Q^2\cdot2\pi E^2\ln2=\pi Nw\,\varepsilon_{\mathrm{EM}}Q^2E^2\ln2\ [\text{machine}];
\]

%dsp-level:0
\[
\qquad Q=0\implies aG=0\ne1\ [\text{machine}];
\]

%dsp-level:0
\[
\qquad Q^2=1,\ E=6:\ \varepsilon^{\mathrm{pole}}=1/(36\pi Nw\ln2)\ [\text{machine}]
\]

%dsp-level:0
\[
(4):\ \text{the two symbol sets enumerated from their defining clauses and intersected}\ [\mathrm{BM.88}(3);\ \mathrm{Def\text{-}BM.12};\ \text{machine}];
\]

%dsp-level:0
\[
\qquad \mathrm{BM.88}(3)\text{'s P4-freedom row re-cited}
\]

%dsp-level:0
\[
(5):\ (\mathrm i)\text{--}(\mathrm{iv})\ \text{re-cited at their own loci}\ [\mathrm{BM.36}(3);\ \mathrm{BM.61}(8);\ \mathrm{BM.90}(3);\ \mathrm{FT.73}(2),(3);\ \mathrm{BM.67}(6),(8);\ \mathrm{FT.55}(1);\ \mathrm{P4}]
\]

%dsp-level:0
\[
(6):\ \mathrm{BM.101}(5)(\mathrm i)\text{--}(\mathrm{iii})\text{'s three crossing counts re-cited};
\]

%dsp-level:0
\[
\qquad \text{the uniformity censuses over}\ \mathsf g\in\{1,1+G,(1+G)^{-1}\}\ \text{on exact rational lattices}\ [\text{machine}]
\]

%dsp-level:0
\[
(7)\text{--}(8):\ \mathrm{BM.88}(2)\text{--}(4),(8)\text{'s rows re-cited};
\]

%dsp-level:0
\[
\qquad \text{the standing-rule scans, including the barred object}\ \alpha_B\ [\text{machine}].\ \blacksquare
\]

\textbf{Phys.} Turn the correction back on itself and the corpus answers with an equation it can solve and a verdict it cannot avoid. If the electromagnetic self-energy feeds the density that feeds the curvature lock, the lock's argument grows by one plus a times the composition, and self-consistency becomes a fixed point. Solved on the three priced candidates, the fixed point is trivial for the constant one, a rational expression with an exact pole for the growing one, and the positive root of a quadratic for the shrinking one — three admissible solutions, so consistency is necessary and useless as a selector. The pole locus is a single monomial, and the neutral class never reaches it, because a neutral composite carries no field energy at all. Then the deeper row: the loop is not merely unselected, it is excluded, because the lock's alphabet is three symbols long and disjoint from the correction's, and letting the correction in would make the lock depend on a declared parameter that the lock's own theorem says it does not carry. So the self-energy rides the mass and not the geometry, the composition stays priced three ways, and the four old grounds against turning a count into a strength bar the fourth route. What the tier gains is a discrimination: the three compositions give the energy-scale graph three different topologies — a half-line, a half-line that always crosses, and a window — so the choice is not cosmetic inside this corpus even though nothing here can make it. The coupling is complete in form with exactly two undeclared members, and no third.

\end{theorem}

\begin{theorem}[{\texorpdfstring{BM.105 — The general crystallization weight: the measure-free bounds, the threshold family, and the refuted defect conditioning}{BM.105 - The general crystallization weight: the measure-free bounds, the threshold family, and the refuted defect conditioning}}]\label{tag:BM-105}
\textbf{Thm.}

%dsp-level:0
\[
(1)\quad\textbf{the measure-free form:}
\]

%dsp-level:0
\[
\qquad E_{\mathrm{int}}(R)=|R|=2\,|S(R)|\ \ \text{on the three pairs of}\ \Sigma_{12}
\]

%dsp-level:0
\[
\qquad [\mathrm{Def\text{-}BM.5}(4);\ \mathrm{Census\ FT.35}(2);\ \mathrm{BM.94}(4)]:
\]

%dsp-level:0
\[
\text{linearity on}\ |R|\ \ \text{gives, for}\ \textbf{every}\ \ \text{ensemble measure}:
\]

%dsp-level:0
\[
\qquad \boxed{\ w=\frac1{E^{\mathrm{cryst}}}\sum_{i=1}^{3}P(\text{the pair}\ i\ \text{is real}),\qquad w_{\mathrm{mat}}=\frac13\sum_{i=1}^{3}a_i,\qquad a_i:=P(\text{the pair}\ i\ \text{is real}\mid\text{materialized})\ }
\]

%dsp-level:0
\[
[\mathrm{BM.89}(4);\ \mathrm{BM.94}(7);\ \text{machine: the linearity row}];\qquad\textbf{no distribution hypothesis enters}\ (1);
\]

%dsp-level:0
\[
(2)\quad\textbf{the measure-free bounds:}
\]

%dsp-level:0
\[
\qquad |S|
\]

%dsp-level:0
\[
\ge 1\ \ \text{on the materialized support}\ [\mathrm{BM.94}(5);\ \mathrm{Def\text{-}H.4};\ \mathrm{FT.144}(4)]\ \ \text{and}\ \ |S|
\]

%dsp-level:0
\[
\le 3
\]

%dsp-level:0
\[
\implies 1\le\langle|S|\rangle_{\mathrm{mat}}\qquad \le 3:
\]

%dsp-level:0
\[
\boxed{\ \frac13\ \le\ w_{\mathrm{mat}}\ \le\ 1\qquad\text{for every admissible ensemble measure}\ }\qquad[\text{machine}];
\]

%dsp-level:0
\[
\textbf{the two equality loci are exact:}
\]

%dsp-level:0
\[
\qquad w_{\mathrm{mat}}=\frac13\iff\text{the measure is supported on the single-pair configurations};
\]

%dsp-level:0
\[
\qquad w_{\mathrm{mat}}=1\iff\text{it is supported on the full configuration}
\]

%dsp-level:0
\[
\qquad [\text{machine}];
\]

%dsp-level:0
\[
\boxed{\ \langle E\rangle_{\mathrm{mat}}=6\,w_{\mathrm{mat}}\in[2,\,6]\ }
\]

%dsp-level:0
\[
\qquad \textbf{the bound is a consequence of the materialization restriction alone,}\ \neg\textbf{of the specialization};
\]

%dsp-level:0
\[
(3)\quad\textbf{the threshold family (hypothesis}\ (\mathrm F)\textbf{, marked):}
\]

%dsp-level:0
\[
\qquad \textbf{hypothesis:}\ \text{the A4 threshold takes a per-pair value}\ \varepsilon_i\in[0,\tfrac14)\ \ \text{on the three pairs of}\ \Sigma_{12}
\]

%dsp-level:0
\[
\qquad [\mathrm{A4};\ \mathrm{D0.4};\ \mathrm{CT0.8}\text{'s derived range applied memberwise};\ \text{carried symbolically},\ \neg\text{issued}]:
\]

%dsp-level:0
\[
p_i:=1-4\varepsilon_i,\qquad u_i:=1-p_i=4\varepsilon_i,\qquad R:=\prod_{i=1}^{3}u_i,\qquad S:=\sum_{i=1}^{3}p_i\qquad[\text{machine}]:
\]

%dsp-level:0
\[
\boxed{\ w_{\mathrm{mat}}\ =\ \frac{S}{3\,(1-R)},\qquad\langle E\rangle_{\mathrm{mat}}=\frac{2S}{1-R},\qquad\bigl\langle E^2\bigr\rangle_{\mathrm{mat}}=\frac{4\bigl(S+2\sum_{i<j}p_ip_j\bigr)}{1-R}\ }
\]

%dsp-level:0
\[
\textbf{the seven materialized configuration weights:}
\]

%dsp-level:0
\[
\qquad P_{\mathrm{mat}}(S)=\frac{\prod_{i\in S}p_i\prod_{i\notin S}u_i}{1-R},
\]

%dsp-level:0
\[
\qquad \varnothing\ne S\subseteq\{1,2,3\},
\]

%dsp-level:0
\[
\qquad \sum_{\varnothing\ne S}P_{\mathrm{mat}}(S)=1
\]

%dsp-level:0
\[
\qquad [\text{machine: exact rational functions, normalized}];
\]

%dsp-level:0
\[
\textbf{the uniform specialization recovers}\ \mathrm{BM.94}\ \textbf{identically:}\qquad p_i\equiv p\implies S=3p,\quad R=(1-p)^3,\quad1-R=pD:
\]

%dsp-level:0
\[
\boxed{\ w_{\mathrm{mat}}=\frac{3p}{3pD}=\frac1D,\qquad D=3-3p+p^2\ }
\]

%dsp-level:0
\[
\qquad [\mathrm{BM.94}(6),(7);\ \text{machine: the identity at six rational thresholds}];
\]

%dsp-level:0
\[
(4)\quad\textbf{the family's exact bounds:}
\]

%dsp-level:0
\[
\qquad 3(1-R)-S=\sum_{i=1}^{3}u_i-3\prod_{i=1}^{3}u_i
\]

%dsp-level:0
\[
\qquad [\text{machine: one polynomial identity}]:
\]

%dsp-level:0
\[
\sum_iu_i\ \ge\ 3\bigl(\prod_iu_i\bigr)^{1/3}\ \ge\ 3\prod_iu_i
\]

%dsp-level:0
\[
\qquad [\prod_iu_i\le1;\ \text{AM--GM}]\implies\boxed{\ 3(1-R)\ \ge\ S\ }
\]

%dsp-level:0
\[
\qquad \text{with equality iff every}\ u_i=0,\ \text{i.e.}\ p_i=1:
\]

%dsp-level:0
\[
\boxed{\ w_{\mathrm{mat}}\le1\ \ \text{with equality iff every pair is crystallized};\qquad w_{\mathrm{mat}}>\frac13\ \ \text{strictly on the materialized domain}\ }
\]

%dsp-level:0
\[
\qquad [(2);\ \text{machine}];
\]

%dsp-level:0
\[
w_{\mathrm{mat}}\ \ \text{is strictly increasing in each}\ p_i\qquad[\text{machine: an exact rational grid on}\ (0,1]^3];
\]

%dsp-level:0
\[
(5)\quad\textbf{hypothesis}\ (\mathrm{W6}) — \textbf{six state-dependent thresholds}\ \varepsilon_i(D_S) — \textbf{on its face and refuted:}
\]

%dsp-level:0
\[
(\mathrm i)\quad\textbf{the alphabets are disjoint:}
\]

%dsp-level:0
\[
\qquad \text{the flicker measure's symbol set is}\ \{\varepsilon_q\}\ \ \text{and the defect values' is}\ \{\delta_v,\ \hat\mu'^2\}
\]

%dsp-level:0
\[
\qquad [\mathrm{BM.90}(1),(2);\ \text{machine: the symbol census of}\ \mathrm{A4},\ \mathrm{D0.4},\ \mathrm{CT0.8},\ \mathrm{BM.19}(2)\ \text{carries}\ \varepsilon_q\ \text{alone, with no state argument}];
\]

%dsp-level:0
\[
(\mathrm{ii})\quad\textbf{the four grounds bar it:}
\]

%dsp-level:0
\[
\qquad \text{the degree-}16\ \text{field against integer and binary reads}\ [\mathrm{FT.73}(2),(3)];
\]

%dsp-level:0
\[
\ \ \text{the two-torsion image reaching}\ 2\ \text{of}\ 6\ \text{states}\ [\mathrm{BM.67}(6)];
\]

%dsp-level:0
\[
\ \ \text{the calibration weight}\ {-}2\ \text{against}\ 0\ [\mathrm{FT.55}(1);\ \mathrm{BM.67}(8)];
\]

%dsp-level:0
\[
\ \ \text{the strength bar}\ [\mathrm{BM.36}(3);\ \mathrm{BM.61}(8)(\mathrm{O1})]
\]

%dsp-level:0
\[
\qquad [\mathrm{BM.90}(3)\ \text{re-cited}];
\]

%dsp-level:0
\[
(\mathrm{iii})\quad\textbf{the datum count:}
\]

%dsp-level:0
\[
\qquad \text{six thresholds would be six P4-class data}\ —\ \text{§Cal-BM would gain six members}\ —\ \text{refused}
\]

%dsp-level:0
\[
\qquad [\mathrm{P4};\ \mathrm{P1}];
\]

%dsp-level:0
\[
(\mathrm{iv})\quad\textbf{the scope row:}
\]

%dsp-level:0
\[
\qquad \mathrm{A4}\text{'s threshold enters the corpus only via}\ \mathrm{CT0.1}\ \text{and}\ \mathrm{CT0.2},
\]

%dsp-level:0
\[
\ \text{with a single}\ \varepsilon_q\ \text{and no state argument}
\]

%dsp-level:0
\[
\qquad [\mathrm{CT0.8}\ \text{scope clause}]:
\]

%dsp-level:0
\[
\boxed{\ \neg\exists\ \text{a derived map}\ D_S\longmapsto\varepsilon;\qquad\textbf{the weight stays branch-blind, and the six defect values act in the pole currency alone}\ }
\]

%dsp-level:0
\[
\qquad [\mathrm{BM.90}(3),(4);\ \mathrm{BM.98}(6)];
\]

%dsp-level:0
\[
(6)\quad\textbf{the residue, narrowed:}
\]

%dsp-level:0
\[
\qquad (\mathrm{R4})\ \ \textbf{closes at two levels:}\ \ \text{the measure-free bounds}\ (1),(2)\ \ \text{and the threshold family}\ (3),(4);
\]

%dsp-level:0
\[
\textbf{what remains is}\ (\mathrm{R4}'):
\]

%dsp-level:0
\[
\qquad \boxed{\ (\mathrm{R4}')\ \ \text{the ensemble's configuration distribution beyond per-pair independence}\ —\ \text{the pair--pair correlations of the measure}\ }
\]

%dsp-level:0
\[
\textbf{and the transport register's ensemble average}
\]

%dsp-level:0
\[
\qquad [\mathrm{BM.89}(8);\ \mathrm{BM.91}(2);\ \mathrm{BM.94}(8):\ \text{the computed cells average the vertex-deficit register}\ \tfrac32n_{\mathcal V},\ \neg\text{the edge-count register}\ |R|]:
\]

%dsp-level:0
\[
(\mathrm{R4})\ \ \textbf{is rescoped to}\ (\mathrm{R4}'),\ \neg\textbf{deleted};
\]

%dsp-level:0
\[
(7)\quad\textbf{status:}
\]

%dsp-level:0
\[
\qquad \textbf{derived:}\ (1)\text{--}(4),\ (6);
\]

%dsp-level:0
\[
\qquad \textbf{refuted:}\ \text{hypothesis}\ (\mathrm{W6})\ \ (5);
\]

%dsp-level:0
\[
\qquad \textbf{marked,}\ \neg\textbf{adopted:}\ \text{hypothesis}\ (\mathrm F)\text{'s family};
\]

%dsp-level:0
\[
\neg\text{a declaration},
\]

%dsp-level:0
\[
\quad\neg\text{a P4 addition},
\]

%dsp-level:0
\[
\quad\neg\text{a new dimensionful scale}\ [\mathrm{P1}],
\]

%dsp-level:0
\[
\quad\neg\text{a measured value}\ [\mathrm{P5}],
\]

%dsp-level:0
\[
\quad\text{§Cal-BM introduces no new free parameter at this tag};
\]

\textbf{Req.} A4 (Part 1); D0.4, D2.9 (Parts 1–2); A §0.2a, CT0.8, CT1.29 (Part 1); P1, P4, P5 (Parts 1–2); L1, L2 (Part 1); Def-H.4, Def-H.32(1), FT.55(1), FT.73(2), FT.73(3), FT.107(3), FT.144(4) (Part 3); Census FT.35(2) (row-wise); Def-BM.5(4), BM.19(1), BM.19(2), BM.36(3), BM.61(8), BM.67(6), BM.67(8), BM.89(4), BM.89(8), BM.90(1), BM.90(2), BM.90(3), BM.90(4), BM.91(2), BM.94(4), BM.94(5), BM.94(6), BM.94(7), BM.94(8), BM.98(6) (this part); exact polynomial identities over $\mathbb Q[p_1,p_2,p_3]$, exhaustive configuration censuses, AM–GM on three non-negative rationals, exact rational grids, symbol censuses [math; machine: \texttt{checks/\allowbreak{}verify\_\allowbreak{}H\_\allowbreak{}six\_\allowbreak{}objective\_\allowbreak{}resolution.py} §E].

\textbf{Pf.}

%dsp-level:0
\[
(1):\ E_{\mathrm{int}}
\]

%dsp-level:0
\[
= 2|S|\ \text{and linearity of the average}\implies\langle E\rangle
\]

%dsp-level:0
\[
= 2\sum_iP(\text{pair}\ i\ \text{real})\ [\mathrm{BM.89}(4)\text{'s per-datum form re-derived without any distribution hypothesis};\ \text{machine}]
\]

%dsp-level:0
\[
\text{the conditional form:}\ P(\text{pair}\ i\ \text{real}\mid\text{materialized})=a_i\ \text{and}\ \langle|S|\rangle_{\mathrm{mat}}=\sum_ia_i\ [\text{machine}]\implies w_{\mathrm{mat}}=\tfrac13\sum_ia_i
\]

%dsp-level:0
\[
(2):\ 1
\]

%dsp-level:0
\[
\le |S|
\]

%dsp-level:0
\[
\le 3\ \text{on the materialized support}\ [\mathrm{BM.94}(5);\ \mathrm{Def\text{-}H.4};\ \mathrm{FT.144}(4)]
\]

%dsp-level:0
\[
\implies 1\le\langle|S|\rangle_{\mathrm{mat}}
\]

%dsp-level:0
\[
\le 3\implies\tfrac13
\]

%dsp-level:0
\[
\le w_{\mathrm{mat}}
\]

%dsp-level:0
\[
\le 1\ [\text{machine}];
\]

%dsp-level:0
\[
\qquad \text{the two equality loci:}\ \langle|S|\rangle=1\iff|S|=1\ \text{a.s.},\ \ \langle|S|\rangle=3\iff|S|=3\ \text{a.s.}\ [\text{machine}]
\]

%dsp-level:0
\[
(3):\ P(S)
\]

%dsp-level:0
\[
= \prod_{i\in S}p_i\prod_{i\notin S}u_i\ \text{and}\ \sum_{S\ne\varnothing}P(S)
\]

%dsp-level:0
\[
= 1-\prod_iu_i
\]

%dsp-level:0
\[
= 1-R\ [\text{machine: the exhaustive eight-configuration census}];
\]

%dsp-level:0
\[
\qquad \langle E\rangle_{\mathrm{mat}}=\tfrac{2\sum_ip_i}{1-R}\ [\text{machine:}\ P(\text{pair}\ i\ \text{real})=p_i\ \text{since a real pair is materialized}]
\]

%dsp-level:0
\[
\langle|S|^2\rangle
\]

%dsp-level:0
\[
= \sum_ip_i+2\sum_{i<j}p_ip_j\ [\text{machine:}\ |S|^2=\sum_i1_i+2\sum_{i<j}1_i1_j]\implies\langle E^2\rangle_{\mathrm{mat}}
\]

%dsp-level:0
\[
= \tfrac{4(S+2\sum_{i<j}p_ip_j)}{1-R}
\]

%dsp-level:0
\[
p_i
\]

%dsp-level:0
\[
\equiv p:\ S
\]

%dsp-level:0
\[
= 3p,
\]

%dsp-level:0
\[
\ R
\]

%dsp-level:0
\[
= (1-p)^3,
\]

%dsp-level:0
\[
\ 1-R
\]

%dsp-level:0
\[
= p(3-3p+p^2)
\]

%dsp-level:0
\[
= pD
\]

%dsp-level:0
\[
\implies w_{\mathrm{mat}}
\]

%dsp-level:0
\[
= 3p/(3pD)
\]

%dsp-level:0
\[
= 1/D\ [\mathrm{BM.94}(6);\ \text{machine at six rational thresholds}]
\]

%dsp-level:0
\[
(4):\ 3(1-\prod u)-\sum(1-u)=3-3\prod u-3+\sum u=\sum u-3\prod u\ [\text{machine: one identity}];
\]

%dsp-level:0
\[
\qquad \sum u
\]

%dsp-level:0
\[
\ge 3(\prod u)^{1/3}\ [\text{AM--GM}]\ \ \text{and}\ \ (\prod u)^{1/3}\ge\prod u\ \text{on}\ \prod u
\]

%dsp-level:0
\[
\le 1\implies\sum u
\]

%dsp-level:0
\[
\ge 3\prod u\ [\text{machine: exact on a rational grid and by the two inequalities}]
\]

%dsp-level:0
\[
\text{equality}\iff u_1=u_2=u_3\ \wedge\ \prod u\in\{0,1\}\iff p_1=p_2=p_3=1\ \text{on the materialized domain}\ [\text{machine}];
\]

%dsp-level:0
\[
\qquad \text{monotonicity in each}\ p_i\ [\text{machine: exact rational grid}]
\]

%dsp-level:0
\[
(5):\ (\mathrm i)\text{--}(\mathrm{iv})\ \text{re-cited at their own loci}\ [\mathrm{BM.90}(1),(2),(3);\ \mathrm{FT.73}(2),(3);\ \mathrm{BM.67}(6),(8);\ \mathrm{FT.55}(1);\ \mathrm{BM.36}(3);\ \mathrm{BM.61}(8)(\mathrm{O1});\ \mathrm{CT0.8};\ \mathrm{P4}];
\]

%dsp-level:0
\[
\qquad \text{the symbol census}\ [\text{machine}]
\]

%dsp-level:0
\[
(6)\text{--}(7):\ \text{the residue rows assembled};\qquad\text{the standing-rule scans}\ [\text{machine}].\ \blacksquare
\]

\textbf{Phys.} The weight that prices a composite's failure to crystallize does not need the specialization that first produced it. Linearity alone gives it as the mean fraction of the three pairs that are real, and the materialization restriction alone gives its bounds: at least one pair is real, at most three are, so the weight lives between a third and one for every measure the ensemble could carry, with both endpoints characterized exactly — a third when only single pairs materialize, one when the full configuration does. That is a stronger statement than the specialization's, and it is why the mean edge count can never fall below the digon's two. Then the family: let each of the three pairs carry its own threshold, still inside the derived admissible range, and the weight becomes the sum of the three retention probabilities over three times one minus their product, with seven materialized configurations, exact second moments, and an arithmetic-mean-geometric-mean row that pins it below one with equality only when every pair is crystallized. Set the three equal and the whole of the previous tier's formula returns identically. What the family does not license is the version the question asks for: six thresholds, one per quark state, each a function of that state's defect value. That is the branch-conditioned weighting already refuted — the measure's alphabet holds one symbol and the defect values' holds two, four separate grounds bar any count-to-strength map, six thresholds would be six declared parameters, and the axiom's threshold carries no state argument. So the weight stays branch-blind, the six defect values stay in the pole currency, and the residue narrows rather than closes: what is left is the correlation between pairs, named and priced, and the transport register's own ensemble average.

\end{theorem}

\begin{theorem}[{\texorpdfstring{BM.106 — The arrow and the kind: the spin label, the tick law, and the exchange-structure referent census}{BM.106 - The arrow and the kind: the spin label, the tick law, and the exchange-structure referent census}}]\label{tag:BM-106}
\textbf{Thm.}

%dsp-level:0
\[
(1)\quad\textbf{the spin label (derived, on the logged transfer):}
\]

%dsp-level:0
\[
\qquad j=\frac{|k|}2,
\]

%dsp-level:0
\[
\qquad k=\frac{\Theta}2\in\mathbb Z_{\ge0}\ \ \text{on shell},
\]

%dsp-level:0
\[
\qquad \Psi(c)=\mathrm{Sym}^{|k(c)|}\bigl(\mathcal H_{\mathsf v}\bigr),
\]

%dsp-level:0
\[
\qquad \text{Casimir}\ \frac k2\Bigl(\frac k2+1\Bigr)
\]

%dsp-level:0
\[
[\mathrm{Def\text{-}H.38};\ \mathrm{FT.60}(1)\text{--}(3);\ \mathrm{FT.22};\ \mathrm{FT.25}(3);\ \text{the transfer's derivation or replacement is a named residue of the seed transport}];
\]

%dsp-level:0
\[
\textbf{the multiplet:}\qquad m\in\{k,\ k-2,\ \dots,\ -k\}\qquad[\mathrm{FT.60}(5);\ \mathrm{Def\text{-}H.38};\ \mathrm{BM.39}(7)];
\]

%dsp-level:0
\[
\textbf{the floor census:}
\]

%dsp-level:0
\[
\qquad \Theta_{\mathrm{int}}=2\kappa\in2\mathbb Z,
\]

%dsp-level:0
\[
\qquad k=\kappa\in\mathbb Z\setminus\{0\},
\]

%dsp-level:0
\[
\qquad \kappa=\widehat{ds}^{\,2}(e)\ \ \text{is registration data determined at crystallization}
\]

%dsp-level:0
\[
\qquad [\mathrm{Census\ FT.28}(2);\ \mathrm{FT.24};\ \text{A §0.2a}]:
\]

%dsp-level:0
\[
\boxed{\ \kappa=1\implies j=\frac12;\qquad\kappa=2\implies j=1\ }\qquad\textbf{on the same}\ \mathrm{D1.5}\ \textbf{kind}\qquad[\text{machine: the two exact rows}]:
\]

%dsp-level:0
\[
\boxed{\ \neg\exists\ \text{a spin-parity theorem}\ —\ \text{both integer and half-integer labels occur on materialized closed bound chains of one kind}\ }
\]

%dsp-level:0
\[
(2)\quad\textbf{the kind rows that are derived:}
\]

%dsp-level:0
\[
\qquad \text{bound}\ :=\ \text{closed self-contained chain},\ \text{content fermionic};
\]

%dsp-level:0
\[
\qquad \text{transit}\ :=\ \text{propagating open chain},\ \text{content bosonic}
\]

%dsp-level:0
\[
\qquad [\mathrm{D1.5}:\ \textbf{a definitional label}]:
\]

%dsp-level:0
\[
\text{the dichotomy is the closed/open read}
\]

%dsp-level:0
\[
\qquad [\mathrm{FT.80}(5);\ \mathrm{D1.4}];
\]

%dsp-level:0
\[
\qquad \text{the materialized locus at rung one is single-arrow}
\]

%dsp-level:0
\[
\qquad [\mathrm{FT.144}(2)]\ \ \text{and single-ray, occupancy}\ (1,0)
\]

%dsp-level:0
\[
\qquad [\text{A-BM.7};\ \mathrm{FT.144}(5);\ \mathrm{BM.73}(3);\ \mathrm{BM.96}(4)];
\]

%dsp-level:0
\[
\textbf{the tick law:}
\]

%dsp-level:0
\[
\qquad \tau_n\to\tau_{n+1},
\]

%dsp-level:0
\[
\qquad \neg\exists\,\tau^{-1},
\]

%dsp-level:0
\[
\qquad \bigl|\mathfrak L(\tau_{n+1})\bigr|=\bigl|\mathfrak L(\tau_n)\bigr|+n_{\mathrm{real\_out}}(\tau_n)\ge\bigl|\mathfrak L(\tau_n)\bigr|
\]

%dsp-level:0
\[
\qquad [\mathrm{A6};\ \mathrm{A7}]:
\]

%dsp-level:0
\[
\text{the prefix law}\ \#J^+(\text{prefix})\ge\#J^-(\text{prefix})\ \ \text{with strict alternation}
\]

%dsp-level:0
\[
\qquad [\mathrm{FT.144}(2);\ \mathrm{D1.2}];
\]

%dsp-level:0
\[
\qquad \text{the pure order reversal of every completed word is inadmissible at its first event}
\]

%dsp-level:0
\[
\qquad [\mathrm{FT.144}(4)(\mathrm{ii})];
\]

%dsp-level:0
\[
\textbf{the contraction:}
\]

%dsp-level:0
\[
\qquad V'=\Bigl(1-\frac\gamma4\Bigr)^{\!2}V\le V,
\]

%dsp-level:0
\[
\qquad \gamma=1:\ V'=\frac9{16}V,
\]

%dsp-level:0
\[
\qquad \gamma\to0:\ V'=V,
\]

%dsp-level:0
\[
\qquad \text{iterated to}\ \Bigl(\frac9{16}\Bigr)^{\!6}
\]

%dsp-level:0
\[
\qquad [\mathrm{L1};\ \mathrm{L2};\ \mathrm{BM.51}(1),(2);\ \text{machine}];
\]

%dsp-level:0
\[
(3)\quad\textbf{the statistics referent census:}\qquad\begin{array}{l|l}\text{a statistics requires} & \text{the corpus supplies}\\\hline\text{a multi-chain state space} & \neg\text{constructed:}\ \neg\text{a two-chain state space exists}\ [\mathrm{FT.80}(5)]\\\text{an exchange operator} & \neg\exists\ \text{one in the inventory}\ [\mathrm{FT.80}(5);\ \mathrm{FT.125}(5)]\\\text{a permutation sign} & \neg\text{formed:}\ \neg\exists\ \text{a permutation action on states}\\\text{an identification of a grade's}\ \mathbb Z_2\ \text{with the sign's}\ \mathbb Z_2 & \textbf{a kind confusion, barred}\ [\mathrm{FT.125}(5);\ \mathrm{FT.89}(2)\ \text{kind discipline}]\\\text{a tensor state on}\ \mathcal H_c & \neg\text{one: a counting space with no inner product, no dynamics, no state}\ [\mathrm{Def\text{-}H.41}\ \text{Species};\ \mathrm{Def\text{-}H.42}]\end{array}
\]

%dsp-level:0
\[
\boxed{\ \textbf{the spin--statistics connection is}\ \neg\textbf{formed}\ }\qquad[\mathrm{FT.125}(5)\ \text{re-cited}];
\]

%dsp-level:0
\[
\textbf{the derivation asked for is refuted on its face:}
\]

%dsp-level:0
\[
\qquad \text{the tick arrow and the spin label are both derived}\ (1),
\]

%dsp-level:0
\[
(2),
\]

%dsp-level:0
\[
\ \ \text{and neither is an exchange sign}\ —\ \text{the missing member is the exchange structure itself},
\]

%dsp-level:0
\[
\ \neg\text{a step from arrow to sign};
\]

%dsp-level:0
\[
(4)\quad\textbf{the four}\ \mathbb Z_2\ \textbf{gradings are single-object, and the census is exhaustive:}
\]

%dsp-level:0
\[
\qquad \text{the algebra}\ S\text{-grading}\ (10,6)\ [\mathrm{FT.136};\ \mathrm{FT.141}(1)];
\]

%dsp-level:0
\[
\ \ \text{the branch-label parity}\ [\mathrm{FT.101};\ \mathrm{FT.135};\ \mathrm{FT.157}(2),(3)];
\]

%dsp-level:0
\[
\text{the carrier}\ \hat\chi\text{-grade}\ [\mathrm{Def\text{-}H.33};\ \mathrm{FT.122}];
\]

%dsp-level:0
\[
\qquad \text{the Janus-arrow chirality}\ [\mathrm{FT.144}(2),(5)]:
\]

%dsp-level:0
\[
\qquad \textbf{pairwise distinct, commuting where jointly defined, neither reducible to another}
\]

%dsp-level:0
\[
\qquad [\mathrm{BM.73}(1);\ \mathrm{FT.141}(6);\ \mathrm{FT.126}(4)]:
\]

%dsp-level:0
\[
\boxed{\ \neg\text{one of the four is an exchange sign}\ }\qquad[\text{machine: the four symbol sets are single-object reads}];
\]

%dsp-level:0
\[
(5)\quad\textbf{where antisymmetry does live:}
\]

%dsp-level:0
\[
\qquad \textbf{the}\ \varepsilon\textbf{-seal is antisymmetric in the colour slots}
\]

%dsp-level:0
\[
\qquad B:=\varepsilon_{ijk}q^iq^jq^k\ [\mathrm{D3.9};\ \mathrm{CT3.15};\ \mathrm{CT3.17}]
\]

%dsp-level:0
\[
\qquad \textbf{and colour is}\ \neg\textbf{observable by construction}
\]

%dsp-level:0
\[
\qquad [\mathrm{CT3.22}]:
\]

%dsp-level:0
\[
\textbf{the transport assembly is sign-free:}
\]

%dsp-level:0
\[
\qquad W(Q)=\binom6{Q+3}>0\ \ \text{on all seven classes}
\]

%dsp-level:0
\[
\qquad [\mathrm{BM.103}(3);\ \text{machine}];
\]

%dsp-level:0
\[
\textbf{the amplitude layer's signs are barred as statistics:}
\]

%dsp-level:0
\[
\qquad (-1)^{\aleph}\ \ \text{with}\ \ \sum_i\aleph_{v_i}\equiv|q|\ (\mathrm{mod}\ 2)
\]

%dsp-level:0
\[
\qquad [\mathrm{Def\text{-}BM.13}(2);\ \mathrm{Def\text{-}BM.5}(2);\ \text{machine: exhaustive to}\ k=7];
\]

%dsp-level:0
\[
\lambda=(-1)^{|E|}e^{i\pi\Theta}\in\{\pm1,\pm i\},
\]

%dsp-level:0
\[
\qquad \Theta_{\mathrm{int}}\in4\mathbb Z\implies\lambda=(-1)^{|E|}
\]

%dsp-level:0
\[
\qquad [\mathrm{BM.25}(3),(6);\ \mathrm{FT.134}(3);\ \text{machine}]:
\]

%dsp-level:0
\[
\boxed{\ \textbf{amplitude}\ne\textbf{mass, and the wall parity is a charge-parity lock,}\ \neg\textbf{a permutation sign}\ }
\]

%dsp-level:0
\[
\qquad [\mathrm{FT.79}(5);\ \mathrm{Def\text{-}BM.13}(2)]
\]

%dsp-level:0
\[
\qquad \textbf{the corpus's only antisymmetry is the colour seal's, and it is unobservable};
\]

%dsp-level:0
\[
(6)\quad\textbf{the mixing row (the arena read):}
\]

%dsp-level:0
\[
\qquad \text{the mixing structure is a read of the derived}\ 16\text{-arena}\ \mathcal K_1\ \ \text{with}\ \ \delta_\nu=\tfrac2{16}=\tfrac18
\]

%dsp-level:0
\[
\qquad [\mathrm{FT.142}(2),(5);\ \mathrm{FT.143}];
\]

%dsp-level:0
\[
\text{the channel arenas are}\ 9,\ 27,\ 16\ \ \text{with the}\ 9/27\ \text{candidates rejected at rung one}
\]

%dsp-level:0
\[
\qquad [\mathrm{FT.142}(3);\ \mathrm{FT.112};\ \mathrm{FT.119}];
\]

%dsp-level:0
\[
\qquad \textbf{the cross-channel identity}\ \ 16\delta_\nu=9\delta_\ell=27\delta_q=2
\]

%dsp-level:0
\[
\qquad [\mathrm{FT.143}(5);\ \text{machine}];
\]

%dsp-level:0
\[
\boxed{\ \neg\text{an external mixing matrix is imported};\qquad\text{the mixing matrix is}\ \neg\text{derived};\qquad\text{every row holds uniformly in the unresolved line index}\ i_\nu\ }
\]

%dsp-level:0
\[
\qquad [\mathrm{FT.145}(3);\ \mathrm{BM.77}(2)];
\]

%dsp-level:0
\[
\textbf{the}\ \nu\ \textbf{row stands as derived:}
\]

%dsp-level:0
\[
\qquad \text{single-ray occupancy, no right-ray matter line, no seesaw-type mechanism}
\]

%dsp-level:0
\[
\qquad [\mathrm{BM.96}(4),(6)\ \text{re-cited},\ \neg\text{extended}];
\]

%dsp-level:0
\[
(7)\quad\textbf{the residue, named:}
\]

%dsp-level:0
\[
\qquad \boxed{\ (\mathrm{R7})\ \ \text{an exchange structure}\ :=\ \text{a two-chain state space}\ +\ \text{an exchange operator}\ +\ \text{a permutation-sign read}\ }
\]

%dsp-level:0
\[
\textbf{the sealed composite is a junction subnetwork,}\ \neg\textbf{a tensor product of chain state spaces:}
\]

%dsp-level:0
\[
\qquad \text{its constituents share vertices}\ [\mathrm{Def\text{-}BM.1};\ \mathrm{Def\text{-}BM.2}]\ \ \text{and no state space is declared on the counting spaces}\ [\mathrm{Def\text{-}H.41}\ \text{Species};\ \mathrm{FT.93}(4)\ \text{pattern}];
\]

%dsp-level:0
\[
(8)\quad\textbf{status:}
\]

%dsp-level:0
\[
\qquad \textbf{derived:}\ (1),\ (2),\ (4),\ (5),\ (6);
\]

%dsp-level:0
\[
\qquad \textbf{refuted as formed:}\ \text{the spin--statistics derivation from the tick arrow}\ (3);
\]

%dsp-level:0
\[
\qquad \textbf{standing:}\ \text{the spin--statistics connection, its multi-chain locus named};
\]

%dsp-level:0
\[
\neg\text{a declaration},
\]

%dsp-level:0
\[
\quad\neg\text{a P4 addition},
\]

%dsp-level:0
\[
\quad\neg\text{a new dimensionful scale}\ [\mathrm{P1}],
\]

%dsp-level:0
\[
\quad\neg\text{a measured value}\ [\mathrm{P5}],
\]

%dsp-level:0
\[
\quad\neg\text{a species identification}\ [\mathrm{Def\text{-}BM.3}(7)];
\]

\textbf{Req.} A6, A7, D0.8, D1.2, D1.4, D1.5, D1.8, D3.9 (Parts 1–2); A §0.2a, CT3.15, CT3.17, CT3.22 (Part 1, Part 2); P1, P4, P5 (Parts 1–2); L1, L2 (Part 1); Def-H.33, Def-H.38, Def-H.41, Def-H.42, FT.22, FT.24, FT.25(3), FT.60(1), FT.60(2), FT.60(3), FT.60(5), FT.80(5), FT.89(2), FT.93(4), FT.101, FT.112, FT.119, FT.122, FT.125(5), FT.126(4), FT.134(3), FT.135, FT.136, FT.141(1), FT.141(6), FT.142(2), FT.142(3), FT.142(5), FT.143, FT.143(5), FT.144(2), FT.144(4), FT.144(5), FT.145(3), FT.157(2), FT.157(3) (Part 3); Census FT.28(2), FT.49 (row-wise); Def-BM.1, Def-BM.2, Def-BM.5(2), Def-BM.13(2), BM.25(3), BM.25(6), BM.39(7), BM.51(1), BM.51(2), BM.73(1), BM.73(3), BM.77(2), BM.79(5), BM.96(4), BM.96(6), BM.103(3) (this part); exact integer winding arithmetic, exhaustive parity censuses, the Casimir ladder on exact rationals, symbol-set censuses [math; machine: \texttt{checks/\allowbreak{}verify\_\allowbreak{}H\_\allowbreak{}six\_\allowbreak{}objective\_\allowbreak{}resolution.py} §F].

\textbf{Pf.}

%dsp-level:0
\[
(1):\ \mathrm{Def\text{-}H.38}\text{'s multiplicity rule and}\ \mathrm{FT.60}(1)\text{--}(3),(5)\text{'s rows re-cited (frozen)};
\]

%dsp-level:0
\[
\qquad \Theta=2\kappa\in2\mathbb Z\implies k=\Theta/2=\kappa\in\mathbb Z\setminus\{0\}\ [\mathrm{Census\ FT.28}(2);\ \mathrm{FT.24};\ \text{machine}]
\]

%dsp-level:0
\[
\kappa=1\implies j=\tfrac12;
\]

%dsp-level:0
\[
\qquad \kappa=2\implies j=1;
\]

%dsp-level:0
\[
\qquad \text{the Casimir rows}\ \tfrac12\cdot\tfrac32=\tfrac34\ \ \text{and}\ \ 1\cdot2=2\ [\text{machine: exact rationals}];
\]

%dsp-level:0
\[
\qquad \text{both chains are materialized closed bound chains of one}\ \mathrm{D1.5}\ \text{kind}\ [\mathrm{Def\text{-}H.4};\ \mathrm{FT.22}]
\]

%dsp-level:0
\[
(2):\ \mathrm{D1.5},\ \mathrm{FT.80}(5),\ \mathrm{FT.144}(2),(4),(5),\ \mathrm{A6},\ \mathrm{A7},\ \mathrm{L2},\ \mathrm{BM.51}(1),(2)\ \text{re-cited verbatim};
\]

%dsp-level:0
\[
\qquad (9/16)^6=531441/16777216\ [\text{machine: exact integers}]
\]

%dsp-level:0
\[
(3):\ \text{the five referent rows re-cited from their own loci}\ [\mathrm{FT.80}(5);\ \mathrm{FT.125}(5);\ \mathrm{FT.89}(2);\ \mathrm{Def\text{-}H.41}\ \text{Species};\ \mathrm{Def\text{-}H.42}];
\]

%dsp-level:0
\[
\qquad \text{the inventory scans:}\ \neg\exists\ \text{an exchange token,}\ \neg\exists\ \text{a permutation-sign read}\ [\text{machine}]
\]

%dsp-level:0
\[
(4):\ \text{the four gradings and their pairwise distinctness re-cited}\ [\mathrm{BM.73}(1);\ \mathrm{FT.141}(6);\ \mathrm{FT.126}(4);\ \mathrm{FT.157}(4)];
\]

%dsp-level:0
\[
\qquad \text{each is a read of one object}\ [\text{machine: the symbol censuses}]
\]

%dsp-level:0
\[
(5):\ \varepsilon_{ijk}\ \text{antisymmetry}\ [\mathrm{CT3.15};\ \mathrm{CT3.17}];
\]

%dsp-level:0
\[
\qquad W(Q)=\binom6{Q+3}>0\ \text{on all seven classes}\ [\mathrm{BM.103}(3);\ \text{machine}];
\]

%dsp-level:0
\[
\qquad \sum_i\aleph_{v_i}
\]

%dsp-level:0
\[
\equiv |q|\ (\mathrm{mod}\ 2)\ \text{and}\
\]

%dsp-level:0
\[
\ge |q|\ \ \text{exhaustively for}\ k
\]

%dsp-level:0
\[
= 2,
\]

%dsp-level:0
\[
\dots,
\]

%dsp-level:0
\[
7\ [\mathrm{Def\text{-}BM.13}(2)\ \text{re-derived};\ \text{machine}]
\]

%dsp-level:0
\[
\lambda=(-1)^{|E|}e^{i\pi\Theta}\in\{\pm1,\pm i\};\qquad\Theta\in4\mathbb Z\implies e^{i\pi\Theta}=1\implies\lambda=(-1)^{|E|}\ [\mathrm{BM.25}(3),(6);\ \text{machine}]
\]

%dsp-level:0
\[
(6):\ \delta_\nu=\tfrac2{16}=\tfrac18\ \text{and the arena census re-cited}\ [\mathrm{FT.142}(2),(3),(5)];
\]

%dsp-level:0
\[
\qquad 16\cdot\tfrac18=9\cdot\tfrac29=27\cdot\tfrac2{27}=2\ [\text{machine: exact Fractions}];
\]

%dsp-level:0
\[
\qquad \neg\exists\ \text{a mixing-matrix tag}\ [\mathrm{FT.145}(3);\ \text{inventory scan}]
\]

%dsp-level:0
\[
(7)\text{--}(8):\ \text{the junction-subnetwork rows re-cited}\ [\mathrm{Def\text{-}BM.1};\ \mathrm{Def\text{-}BM.2};\ \mathrm{Def\text{-}H.41}\ \text{Species}];
\]

%dsp-level:0
\[
\qquad \text{the standing-rule scans}\ [\text{machine}].\ \blacksquare
\]

\textbf{Phys.} The corpus can name the spin and it can name the arrow, and it cannot join them, for a reason that is structural rather than temporary. The spin label comes from the winding: half the winding number, with the Casimir ladder and the weight multiplet behind it, and at the floor the winding is the registration integer of the two-cycle, so the same kind of object carries spin one half at one value of that integer and spin one at the next. There is therefore no parity theorem to be had: the floor itself supplies both parities on materialized closed bound chains, all of one definitional kind. The arrow is just as concrete — no inverse tick, a ledger that never shrinks, a prefix law that makes dissolution presuppose crystallization, a Lyapunov function that contracts by nine sixteenths per step at unit tilt — and the occupancy it forces is one ray full and the other empty. What is missing is not a step from arrow to sign but the object the sign would act on: there is no two-chain state space, no exchange operator, no permutation representation, and the corpus has already ruled that identifying a single-object grade with a permutation sign is a confusion of kinds. So the connection is not derived, and deficit eleven stands with its residue now named: an exchange structure, three ingredients, none of them in the inventory. What the tier adds is the census that makes the absence exact, the four gradings shown to be single-object and exhaustively so, and the location of the only antisymmetry the corpus does have — the colour seal's epsilon, which is unobservable by construction, against a transport assembly whose seven weights are all positive and an amplitude layer whose signs are barred as masses. The mixing rows stay where the arena puts them: sixteen dimensions, one eighth, and no imported matrix.

\end{theorem}

\medskip\noindent\textit{\texorpdfstring{BM.107 — The status of the six resolutions: adjudication and scope}{BM.107 - The status of the six resolutions: adjudication and scope}.}\ Structural-constraints record: see Appendix B, Section B.1.30 of the Addendum.\par\medskip

\phantomsection
\section[Tier BM.21 — The Antimatter and Resonance Tier: the…]{\texorpdfstring{Tier BM.21 — The Antimatter and Resonance Tier: the CP-Conjugate, the Meson Seal, the Exchange Operator, and the Local Resonance Spectrum}{Tier BM.21 - The Antimatter and Resonance Tier: the CP-Conjugate, the Meson Seal, the Exchange Operator, and the Local Resonance Spectrum}}

\textbf{Locus.} This tier takes the residue $(\mathrm{R7})$ that $\mathrm{BM.106}(7)$ named — \textit{an exchange structure: a two-chain state space, an exchange operator, a permutation-sign read} — and constructs all three members, then reads the spectrum the construction exposes. The CP-conjugate chain is derived from the three reversals the corpus already owns, with the forward-only tick blocking $\mathrm T$ as a dynamical inverse and leaving $\mathrm{CP}$ as the one composition that closes the ledger's own charge balance; the vacuum fluctuation is its pair, with the creation weight the Janus resolvent's own $x_+x_-$ and no pair at either endpoint of the threshold range; and the antiquark follows from the same three ingredients that force the compensating charge, so the corpus's Dirac parallel is an inventory statement and not a history one ($\mathrm{BM.108}$–$\mathrm{BM.110}$). The meson is then the $\delta$-sealed $2$-juncture, its state space is the first multi-chain space the corpus has, it discharges the $\mathrm{BM.6}(2)$ specification exactly as written, and its mass is derived on the two routes where the seal contributes no binding term and named unconstructed on the third ($\mathrm{BM.111}$–$\mathrm{BM.113}$). The exchange operator is the canonical flip composed with the CP transport, an involution with spectrum $\{\pm1\}$ unconditionally on the materialized locus, and its permutation-sign read is the fourth-root phase of the winding evaluated at the on-shell ladder — one integer's parity, which is the spin label's parity, which closes deficit $11$ ($\mathrm{BM.112}$–$\mathrm{BM.113}$). The resonance spectrum is then counted: two characteristic polynomials, one transport and one defect, the multiset-orbit law with its stabilizer multiplicities, exactly one exact collision in the whole three-chain ladder, a minimal-polynomial degree census that terminates at the field's own degree, a materialization threshold whose locus is never empty at any $n$ and whose span ratio is exactly thirty at every level, and a decay width that is a product of three derived factors and vanishes identically at the sharp threshold ($\mathrm{BM.114}$–$\mathrm{BM.117}$). Nothing is forbidden, nothing is stable as a category, and no cosmological clock is read: a resonance is an eigenvalue, its appearance is the evaluation of a threshold function at the local parameters, and its width is the evaluation of a width function at the same ones. $\mathrm{BM.118}$ adjudicates.

\textbf{Standing rows.} $\alpha_B$ stays barred on $\mathrm{BM.61}(8)$'s four grounds. No clause of $\mathrm{BM.1}$–$\mathrm{BM.107}$ is retracted, replaced or renumbered: the rows this tier touches are re-scoped at their own loci — $\mathrm{BM.106}(3)$'s statistics referent census and $(7)$'s residue $(\mathrm{R7})$, $\mathrm{BM.103}(4)$'s tensor-product referent census, $\mathrm{BM.6}(2)$'s exchange-space specification, $\mathrm{FT.125}(5)$'s and $\mathrm{FT.80}(5)$'s inventory rows, and $\mathrm{BM.32}(5)$'s amplitude/norm locus row — and each rescope is recorded in $\mathrm{BM.118}(5)$, in $\mathrm{BM.77}(7)$'s pattern. No measured coupling, mass, width, lifetime or ratio enters any derivation $[\mathrm{P5}]$: the threshold and width functions are derived symbolic forms and are $\neg$compared to any external number; the crossing strength of $\mathrm{BM.101}(3)$ is re-cited and $\neg$re-evaluated. No right-ray matter line, no seesaw-type mechanism, no Yukawa-type object and no external mixing matrix is suggested, added or required $[\mathrm{EW.15};\ \mathrm{FT.145}(3);\ \mathrm{BM.96}(4),(6)\ \text{re-cited},\ \neg\text{extended}]$. Species names are confined to the identification layer $[\mathrm{Def\text{-}BM.3}(7);\ \mathrm{Def\text{-}BM.6}(7);\ \mathrm{Def\text{-}BM.14}(8)]$. Topology labels are internal bookkeeping $[\mathrm{Def\text{-}H.32}(4)]$. Amplitude $\ne$ mass $[\mathrm{FT.79}(5)]$. Rung $\ne$ scale $[\mathrm{D0.8}]$. No currency is converted $[\mathrm{BM.83}(3);\ \mathrm{BM.102}(6)]$. No theorem of this tier imports a HistLog disposition as premise: Q10, Q11, item $20$, deficits $1'$, $3$, $4$, $11$ and $12$ are cited as pointers. One dimension only: units are powers of $\mathfrak q$ $[\mathrm{D2.9};\ \mathrm{CT2.1}]$.

\textbf{Local notation, fixed for the whole tier.} $c$ a materialized closed bound chain $[\mathrm{D1.4};\ \mathrm{D1.5};\ \mathrm{Def\text{-}H.4}]$; $\omega$ the orientation-reversal involution $[\mathrm{D2.6}]$; $\Theta_{\mathfrak L}$ the ledger reflection $[\mathrm{D2.7}]$; $\mathrm C,\mathrm P,\mathrm T$ the three derived involutions $[\mathrm{EM.24}(1)\text{--}(3)]$; $\mathrm{CP}:=\mathrm C\circ\mathrm P$; $\Theta(C)=\sum_{e\in E(C)}\widehat{ds}^{\,2}(e)$ and $\lambda=(-1)^{|E(C)|}e^{i\pi\Theta(C)}$ the CPT phase $[\mathrm{EM.24}(5);\ \mathrm{BM.25}(3)]$; $k(c)=\Theta(c)/2$ the on-shell winding and $j=|k|/2$ the spin label $[\mathrm{FT.22};\ \mathrm{FT.25}(3);\ \mathrm{Def\text{-}H.38};\ \mathrm{BM.74}]$; $\Psi(c)=\mathrm{Sym}^{|k(c)|}(\mathcal H_{\mathsf v})$ the winding carrier with $\mathcal H_{\mathsf v}\cong\mathbb C^2$ $[\mathrm{Def\text{-}H.38};\ \mathrm{EW.32}]$; $\mathcal H_\omega(c)=\operatorname{span}_{\mathbb C}\{|c\rangle,|\omega(c)\rangle\}\cong\mathbb C^2$ the orientation doublet with $\hat\omega,\hat\chi,P_{L/R}$ $[\mathrm{Def\text{-}H.33};\ \mathrm{FT.122}]$; $E_{\mathrm{int}}$ the internal directed-edge count, $m_{\mathrm{hid}}=E_{\mathrm{int}}\ln2/\mathfrak h$, $\Delta x_{\mathrm B}=\mathfrak h/(E_{\mathrm{int}}\ln2)$, $r_s=\mathfrak h E_{\mathrm{int}}/2$, $G=2\pi E_{\mathrm{int}}^2\ln2$ $[\mathrm{D0.8};\ \mathrm{FT.44}(4);\ \mathrm{Def\text{-}BM.4}(1);\ \mathrm{BM.13}(1);\ \mathrm{BM.15}(4);\ \mathrm{BM.16}(1);\ \mathrm{BM.88}(3)]$; $N$ the lapse and $f=1-1/k$ with $k=r/\Delta x_{\mathrm B}$ the envelope profile $[\mathrm{BM.34}(1);\ \mathrm{BM.88}(2)]$; $\varepsilon_{\mathrm{eff}}=Nf\,\mathsf g\,\varepsilon_{\mathrm{EM}}=e^2/(4\pi)$ $[\mathrm{BM.88}(4);\ \mathrm{BM.101}(2)]$; $\varepsilon_q\in[0,\tfrac14)$ the crystallization threshold, $x_+=1-4\varepsilon_q$, $x_-=4\varepsilon_q$, $p=x_+/(x_++x_-)=1-4\varepsilon_q$ $[\mathrm{A4};\ \mathrm{D0.4};\ \mathrm{CT0.8};\ \mathrm{BM.19}(2)]$; $\Gamma_{\mathrm{tr}}=x_+x_-/(1-x_+x_-)\in[0,\tfrac13]$ the unit-fugacity transition weight $[\mathrm{BM.44}(2)]$; $D_S(s)\in\{2,4,2-\sqrt2,1,2-2\cos\tfrac\pi{10},2-\sqrt3\}$ and $m^2_{\mathrm{pole}}=\hat\mu'^2D_S$ the six-state defect spectrum $[\mathrm{Def\text{-}BM.14}(3);\ \mathrm{BM.75}(4)]$; $\mathcal S$ the six-state set, $|\mathcal S|=6$ $[\mathrm{Def\text{-}BM.14}(2)]$; $W_n(Q)$ the $n$-chain transport class weight $[\mathrm{BM.28}(2);\ \mathrm{BM.103}(3)]$; $\mathfrak u:=2\cos\tfrac\pi{10}$ and $\mathbb M:=\mathbb Q(\sqrt2,\sqrt3,\mathfrak u)$ the exact eigenvalue field; $\chi_{\{\cdot\}}$ the exact indicator, valued in $\{0,1\}$; $[\cdot]$ the $\mathfrak q$-exponent map $[\mathrm{CT2.1};\ \mathrm{D2.9}]$.

\begin{definition}[{\texorpdfstring{Def-BM.15 — The CP-conjugate chain and the CP-pair}{Def-BM.15 - The CP-conjugate chain and the CP-pair}}]\label{tag:Def-BM-15}
\textbf{Def.} A construction on the three derived involutions of $\mathrm{EM.24}(1)\text{--}(3)$ and on $\mathrm{D2.6}$'s $\omega$; no new axiom, no new datum, no new parameter, no external object.

%dsp-level:0
\[
(1)\quad\textbf{the three reversals, re-cited as derived:}
\]

%dsp-level:0
\[
\qquad \mathrm C:\ \sigma_e\mapsto-\sigma_e\ \forall e\implies q\mapsto-q,\quad f_v\mapsto-f_v,\quad\aleph_v\ \text{invariant}
\]

%dsp-level:0
\[
\qquad [\mathrm{EM.24}(1);\ \mathrm{D3.3};\ \mathrm{D3.4}];
\]

%dsp-level:0
\[
\mathrm P:\ \mathcal P\mapsto E_{\mathcal R}\setminus\mathcal P,\qquad\text{equivalently}\quad\alpha_e\mapsto\pi-\alpha_e\qquad[\mathrm{EM.24}(2);\ \mathrm{Def\text{-}EM.21};\ \mathrm{EM.20}(1)];
\]

%dsp-level:0
\[
\mathrm T:\ \Theta_{\mathfrak L},\qquad w(\Theta_{\mathfrak L}C)=w(C)^\ast,\qquad J^+\leftrightarrow J^-,\qquad G(e)\mapsto G(e)^\ast\qquad[\mathrm{EM.24}(3);\ \mathrm{D2.7};\ \mathrm{CT2.8}(1)];
\]

%dsp-level:0
\[
\textbf{constitutive clause:}\ \ \text{all three are involutions on}\ \mathfrak C(\mathcal R)\ \text{and all three are derived}\ —\ \mathrm C^2
\]

%dsp-level:0
\[
= \mathrm P^2
\]

%dsp-level:0
\[
= \mathrm T^2\qquad = \mathrm{id},
\]

%dsp-level:0
\[
\qquad \text{no external input}
\]

%dsp-level:0
\[
\qquad [\mathrm{EM.24}(4)];
\]

%dsp-level:0
\[
(2)\quad\textbf{the CP-conjugate:}\qquad\boxed{\ \mathrm{CP}\ :=\ \mathrm C\circ\mathrm P,\qquad c^{\mathrm{CP}}\ :=\ \mathrm{CP}(c)\ }
\]

%dsp-level:0
\[
\textbf{constitutive clause:}\ \ \mathrm{CP}\ \text{is the composition of two derived involutions, hence an involution:}
\]

%dsp-level:0
\[
\quad (\mathrm{CP})^2=\mathrm C^2\mathrm P^2=\mathrm{id};
\]

%dsp-level:0
\[
\textbf{the action (derived, entrywise):}
\]

%dsp-level:0
\[
\qquad q\mapsto-q\ [\mathrm{EM.24}(1)];
\]

%dsp-level:0
\[
\qquad \text{scar}=q_{\mathrm{core}}\mapsto-q_{\mathrm{core}}\ [\mathrm{Def\text{-}BM.7}(3);\ \omega\text{-odd}];
\]

%dsp-level:0
\[
\text{chirality}\ \mapsto\ -\text{chirality}\ [\mathrm{D2.6};\ \mathrm{FT.122}(4):\ \hat\omega\hat\chi\hat\omega^{-1}=-\hat\chi];
\]

%dsp-level:0
\[
\qquad \text{triplet attachment}\ q^i\ \mapsto\ \text{dual attachment}\ r_i\ [\mathrm{D3.8};\ \mathrm{CT3.17}:\ \delta^i{}_j\ \text{consumes one of each type}];
\]

%dsp-level:0
\[
w\ \mapsto\ e^{i\pi\Theta}\,w
\]

%dsp-level:0
\[
\qquad [\mathrm{EM.24}(5)\ \text{with}\ (-1)^{|E|}\ \text{cancelled by}\ \mathrm P\text{'s}\ e^{i\pi\Theta}:\ \text{machine}];
\]

%dsp-level:0
\[
\qquad E_{\mathrm{int}}\ \mapsto\ E_{\mathrm{int}},\quad |R|\ \mapsto\ |R|
\]

%dsp-level:0
\[
\qquad [\mathrm{BM.25}(2);\ \mathrm{BM.27}(2)];
\]

%dsp-level:0
\[
(3)\quad\textbf{the CP-pair and the vacuum fluctuation:}
\]

%dsp-level:0
\[
\qquad \textbf{a CP-pair}\ :=\ \text{an orbit}\ \{c,\ c^{\mathrm{CP}}\}\ \text{of the involution of}\ (2);
\]

%dsp-level:0
\[
\textbf{a vacuum fluctuation}\ :
\]

%dsp-level:0
\[
= \ \text{the creation of one CP-pair from the vacuum sector}\ [\mathrm{BM.96}(4):\ \text{the complement's content is all-virtual and unregistered}];
\]

%dsp-level:0
\[
\textbf{constitutive clause:}\ \ \text{the pair's total charge read is}\ Q=q(c)+q(c^{\mathrm{CP}})=q-q=0\ \ \text{exactly}
\]

%dsp-level:0
\[
\qquad [\mathrm{EM.11};\ \mathrm{D2.11};\ \text{machine}];
\]

%dsp-level:0
\[
\textbf{non-derived clause:}\ \ \text{the word ``vacuum'' names the all-virtual sector of}\ \mathrm{CT4.2}\ \text{and}\ \mathrm{FT.144}(4);
\]

%dsp-level:0
\[
\ \ \text{it carries no energy datum and no state};
\]

%dsp-level:0
\[
(4)\quad\textbf{the scope clause (the mirror is formal, not dynamical):}
\]

%dsp-level:0
\[
\qquad \neg\exists\,\tau^{-1}\ [\mathrm{A6}];
\]

%dsp-level:0
\[
\qquad \bigl|\mathfrak L(\tau_{n+1})\bigr|=\bigl|\mathfrak L(\tau_n)\bigr|+n_{\mathrm{real\_out}}(\tau_n)\ge\bigl|\mathfrak L(\tau_n)\bigr|\ [\mathrm{A6};\ \mathrm{A7}];
\]

%dsp-level:0
\[
V'=\bigl(1-\tfrac\gamma4\bigr)^2V\le V,\qquad\gamma=1:\ V'=\tfrac9{16}V\qquad[\mathrm{L2};\ \mathrm{BM.51}(1),(2)]:
\]

%dsp-level:0
\[
\therefore\ \mathrm T\ \textbf{is blocked as a dynamical inverse:}\quad\Theta_{\mathfrak L}\ \text{is an involution on configuration data and}\ \neg\text{a step of the evolution}\ [\mathrm{Def\text{-}BM.7}(2);\ \mathrm{BM.25}(2);\ \mathrm{FT.144}(4)(\mathrm{ii})];
\]

%dsp-level:0
\[
\textbf{the consequence that makes}\ \mathrm{CP}\ \textbf{the operative conjugation:}
\]

%dsp-level:0
\[
\qquad \text{the ledger's balance is a condition on}\ \textbf{configurations},\ \neg\text{on evolutions}\ [\mathrm{D2.11};\ \mathrm{EM.11}],
\]

%dsp-level:0
\[
\text{so the compensating member a charged chain requires is a configuration, and}\ \mathrm{CP}\ \text{is the derived involution that produces one with the same}\ E_{\mathrm{int}}\ (2);
\]

%dsp-level:0
\[
(5)\quad\textbf{the data bar (standing):}
\]

%dsp-level:0
\[
\qquad \neg\exists\ \text{an annihilation channel, a pair-production cross-section or an inter-sector transition operator in the inventory}\ [\mathrm{Def\text{-}BM.7}(5);\ \mathrm{BM.27}(7);\ \mathrm{BM.29}(5)];
\]

%dsp-level:0
\[
\text{``antimatter'' is the label of}\ \mathrm{Def\text{-}BM.7}(1)\text{'s}\ \mathrm{EM}^-\ \text{sector and carries no external datum}\ [\mathrm{P5};\ \mathrm{Def\text{-}BM.3}(6)];
\]

%dsp-level:0
\[
\text{the CP-pair creation weight of}\ \mathrm{BM.110}(1)\ \text{is}\ \neg\text{a cross-section and}\ \neg\text{a rate in external units}\ [\mathrm{BM.44}(5);\ \mathrm{D1.8}].
\]

\textbf{Rmk (what is constitutive and what is derived).} Clause (2)'s \textit{definition} of $\mathrm{CP}$ as $\mathrm C\circ\mathrm P$ is the constitutive act — one composition of two objects the bosonic era already derived. Everything else in (2) is derived content recorded at the definition site: the entrywise action follows from $\mathrm{EM.24}(1),(2)$ and $\mathrm{D2.6}$, the weight rule from $\mathrm{EM.24}(5)$, and the invariance of $E_{\mathrm{int}}$ from $\mathrm{BM.25}(2)$'s record-size row. Clause (3)'s charge neutrality is derived, not stipulated: it is $\mathrm{EM.11}$'s balance applied to an orbit of an involution that negates the read. Nothing here names a particle, imports a mass, or supplies an annihilation amplitude.

\textbf{Req.} A6, A7, D0.8, D1.4, D1.5, D1.8, D2.6, D2.7, D2.11, D3.3, D3.4, D3.8 (Parts 1–2); CT0.8, CT2.8(1), CT3.17, CT4.2 (Part 1, Part 2); EM.11, EM.20(1), EM.24(1), EM.24(2), EM.24(3), EM.24(4), EM.24(5), Def-EM.21 (Part 2); P5 (Parts 1–2); L2 (Part 1); Def-H.4, Def-H.33, FT.122(4), FT.144(2), FT.144(4) (Part 3); Def-BM.4(1), Def-BM.7(1), Def-BM.7(2), Def-BM.7(3), Def-BM.7(5), BM.25(2), BM.27(2), BM.27(7), BM.29(5), BM.51(1), BM.51(2), BM.96(4) (this part); exact $\mathbb Q(i)$ weight algebra on the fourth-root alphabet, exhaustive charge-alphabet enumeration [math; machine: \texttt{checks/\allowbreak{}verify\_\allowbreak{}H\_\allowbreak{}antimatter\_\allowbreak{}resonance\_\allowbreak{}spectrum.py} §A].

\textbf{Pf.}

%dsp-level:0
\[
\text{Definition block: the constitutive clauses are}\ (1)\text{'s involution rows (re-cited, frozen)},
\]

%dsp-level:0
\[
\ (2)\text{'s composition}\ \mathrm{CP}:
\]

%dsp-level:0
\[
= \mathrm C\circ\mathrm P,
\]

%dsp-level:0
\[
\ \text{and}\ (3)\text{'s orbit naming};
\]

%dsp-level:0
\[
\ \text{the non-derived clauses are the ``vacuum'' vocabulary of}\ (3)\ \text{and the ``antimatter'' label of}\ (5)
\]

%dsp-level:0
\[
(2)\text{'s derived rows:}\ q\mapsto-q\ [\mathrm{EM.24}(1)];
\]

%dsp-level:0
\[
\qquad \alpha_e\mapsto\pi-\alpha_e\ [\mathrm{EM.24}(2)];
\]

%dsp-level:0
\[
\qquad \text{composed:}\ w\mapsto\bar w\mapsto e^{i\pi\Theta}w\ [\mathrm{EM.24}(5)\ \text{Pf row}\ \mathrm{P}\circ\mathrm T\ \text{transported to}\ \mathrm C\circ\mathrm P;\ \text{machine}]
\]

%dsp-level:0
\[
\text{attachment type:}\ \delta^i{}_j\ \text{consumes one triplet and one dual}\ [\mathrm{CT3.17}]\implies\text{the conjugate constituent carries the dual attachment}\ [\mathrm{D3.8};\ \mathrm{Def\text{-}BM.1}\ \text{meson clause}];
\]

%dsp-level:0
\[
\qquad E_{\mathrm{int}}=|R|\ \text{invariant}\ [\mathrm{BM.27}(2);\ \text{machine}]
\]

%dsp-level:0
\[
(3):\  Q=q+(-q)=0\ [\text{machine: exhaustive over}\ q\in\{-1,0,+1\}];
\]

%dsp-level:0
\[
\qquad (4):\  \mathrm{A6},\ \mathrm{A7},\ \mathrm{L2}\ \text{rows re-cited};
\]

%dsp-level:0
\[
\qquad (5):\  \text{the inventory scan of}\ \mathrm{BM.27}(7)\ \text{re-cited}.\ \blacksquare
\]

\textbf{Phys.} Three reversals exist in this corpus and each is derived: flip every edge sign, complement the aligned set, reflect the ledger. The third is the one that cannot run — the tick has no inverse and the Lyapunov function only contracts, so a reversed history is a configuration, never a step. Composing the first two gives the conjugation that does the work: it negates the charge read, flips the scar, exchanges the chiral rays, turns a triplet attachment into a dual one, multiplies the weight by a fourth root of unity, and leaves the edge count — and hence the informational mass, the width, the volume, the density and the comparison radius — untouched. An orbit of that involution is a pair, and a pair is neutral by construction, because the balance the valid-universe condition demands is a condition on configurations and the conjugate is exactly the configuration that satisfies it. What the definition does not do is give the pair a process: no annihilation channel, no production cross-section, no inter-sector operator exists in the inventory, and the word "vacuum" names the all-virtual sector rather than any state with an energy.

\end{definition}

\begin{theorem}[{\texorpdfstring{BM.108 — The CP-conjugate chain: the reversal census, the T-block, the exact phase, and the forced conjugate}{BM.108 - The CP-conjugate chain: the reversal census, the T-block, the exact phase, and the forced conjugate}}]\label{tag:BM-108}
\textbf{Thm.}

%dsp-level:0
\[
(1)\quad\textbf{the reversal census (what the corpus owns and what it blocks):}\qquad\begin{array}{l|l|l}\text{reversal} & \text{status} & \text{locus}\\\hline\mathrm C\ \text{charge conjugation} & \textbf{derived, dynamical} & [\mathrm{EM.24}(1);\ \mathrm{D3.3};\ \mathrm{D3.4}]\\\mathrm P\ \text{parity} & \textbf{derived, dynamical} & [\mathrm{EM.24}(2);\ \mathrm{Def\text{-}EM.21}]\\\mathrm T\ \text{time reversal} & \textbf{derived as a record involution, blocked as a dynamical inverse} & [\mathrm{EM.24}(3);\ \mathrm{A6};\ \mathrm{L2};\ \mathrm{FT.144}(4)(\mathrm{ii})]\\\mathrm{CP} & \textbf{derived, an involution} & [\mathrm{Def\text{-}BM.15}(2);\ \mathrm{EM.24}(4)]\\\mathrm{CPT} & \textbf{derived,}\ \mathrm{CPT}^2=\mathrm{id} & [\mathrm{EM.24}(4);\ \mathrm{BM.25}(1)]\end{array}
\]

%dsp-level:0
\[
\textbf{the T-block, in three exact rows:}\qquad(\mathrm i)\ \ \neg\exists\,\tau^{-1}\qquad[\mathrm{A6}];
\]

%dsp-level:0
\[
(\mathrm{ii})\quad\bigl|\mathfrak L(\tau_{n+1})\bigr|-\bigl|\mathfrak L(\tau_n)\bigr|=n_{\mathrm{real\_out}}(\tau_n)\ge0\qquad[\mathrm{A6};\ \mathrm{A7}];
\]

%dsp-level:0
\[
(\mathrm{iii})\quad V'-V=-\Bigl[1-\bigl(1-\tfrac\gamma4\bigr)^2\Bigr]V=-\frac{\gamma(8-\gamma)}{16}V\le0,\qquad\gamma=1:\ V'=\tfrac9{16}V\qquad[\mathrm{L2};\ \text{machine}]:
\]

%dsp-level:0
\[
\boxed{\ \textbf{the pure order reversal of every completed word is inadmissible at its first event}\ —\ \Theta_{\mathfrak L}\ \text{exists and}\ \neg\text{evolves}\ }
\]

%dsp-level:0
\[
\qquad [\mathrm{BM.25}(2);\ \mathrm{FT.144}(4)(\mathrm{ii});\ \text{machine}];
\]

%dsp-level:0
\[
(2)\quad\textbf{the CP action on every derived read of a chain (exact, entrywise):}\qquad\begin{array}{l|l|l}\text{read} & \text{CP image} & \text{locus}\\\hline q=q_{\mathrm{core}}=\tfrac12(\sigma_s-\sigma_t) & -q & [\mathrm{D3.4};\ \mathrm{EM.9};\ \text{machine}]\\\text{scar} & -\text{scar} & [\mathrm{Def\text{-}BM.7}(3);\ \mathrm{BM.27}(5);\ \text{machine}]\\\ \hat\chi\text{-grade} & -\hat\chi\text{-grade}\ (P_L\leftrightarrow P_R) & [\mathrm{D2.6};\ \mathrm{FT.122}(4);\ \mathrm{BM.73}(4);\ \text{machine over}\ \mathbb Q(i)]\\\text{attachment type}\ (q^i\ /\ r_i) & r_i\ /\ q^i & [\mathrm{D3.8};\ \mathrm{CT3.17}]\\\text{weight}\ w & e^{i\pi\Theta}\,w & [\mathrm{EM.24}(5);\ \mathrm{BM.25}(3);\ \text{machine}]\\|w|^2 & |w|^2 & [\mathrm{EM.24}(5):\ \text{unconditional};\ \text{machine}]\\ E_{\mathrm{int}},\ |R|,\ m_{\mathrm{hid}},\ \Delta x_{\mathrm B},\ V_{\mathrm{eff}},\ \rho,\ \kappa^2\rho,\ r_s & \textbf{fixed} & [\mathrm{BM.27}(2),(3);\ \mathrm{BM.28}(1):\ q\text{-exponent}\ 0\ \text{in all seven};\ \text{machine}]\\\Theta_{\mathrm{int}}\ \text{class} & \text{fixed} & [\mathrm{BM.28}(6);\ \text{machine}]\end{array}
\]

%dsp-level:0
\[
\boxed{\ \textbf{the CP-conjugate chain is mass-degenerate with the chain on every route}\ }
\]

%dsp-level:0
\[
\qquad [\mathrm{BM.28}(5);\ \mathrm{BM.27}(3)\ \mathrm{Rmk};\ \text{machine: the monomial}\ q\text{-exponent census}];
\]

%dsp-level:0
\[
(3)\quad\textbf{the exact CP phase:}\qquad w(c^{\mathrm{CP}})=\mu_{\mathrm{CP}}(c)\,w(c),\qquad\boxed{\ \mu_{\mathrm{CP}}(c)\ :=\ e^{i\pi\Theta(c)}\ }
\]

%dsp-level:0
\[
\textbf{derivation:}\quad w(\mathrm{CPT}\,C)=(-1)^{|E|}e^{i\pi\Theta}w(C)\ [\mathrm{EM.24}(5)],
\]

%dsp-level:0
\[
\qquad w(\mathrm T\,C)=w(C)^\ast\ [\mathrm{CT2.8}(1)],
\]

%dsp-level:0
\[
\qquad \mathrm{CP}=(\mathrm{CPT})\circ\mathrm T^{-1}=(\mathrm{CPT})\circ\mathrm T:
\]

%dsp-level:0
\[
\mu_{\mathrm{CP}}=(-1)^{|E|}e^{i\pi\Theta}\cdot\frac{w}{w^\ast}\bigg|_{\text{the}\ \mathrm T\ \text{leg}}=(-1)^{|E|}e^{i\pi\Theta}\cdot(-1)^{|E|}=e^{i\pi\Theta}
\]

%dsp-level:0
\[
\qquad [\text{machine:}\ |w|=1\implies w/w^\ast=w^2/|w|^2;\ \text{the}\ (-1)^{|E|}\ \text{cancels exactly}];
\]

%dsp-level:0
\[
\mu_{\mathrm{CP}}\in\{\pm1,\pm i\}\ \ \text{on the half-integer ladder},
\]

%dsp-level:0
\[
\qquad \mu_{\mathrm{CP}}\in\{+1\}\ \ \text{on the materialized locus}\ \Theta\in2\mathbb Z
\]

%dsp-level:0
\[
\qquad [\mathrm{FT.22};\ \mathrm{BM.2}(1);\ \text{machine}]:
\]

%dsp-level:0
\[
\boxed{\ \textbf{the CP character is exactly}\ 1\ \textbf{on the sealed-composite junction class}\ }
\]

%dsp-level:0
\[
\qquad [\Theta_{\mathrm{int}}\in4\mathbb Z\ [\text{Census FT.49}]\implies e^{i\pi\Theta}=1;\ \mathrm{BM.25}(6);\ \text{machine}];
\]

%dsp-level:0
\[
(4)\quad\textbf{the orbit census:}
\]

%dsp-level:0
\[
\qquad \mathrm{CP}:\ (q,\ \text{orientation})\mapsto(-q,\ \text{flipped})\ \ \text{on}\ \{-1,0,+1\}\times\{w,w^\ast\}:
\]

%dsp-level:0
\[
\qquad \text{no fixed sector, three orbits of size two}\ [\mathrm{BM.25}(4);\ \text{machine}];
\]

%dsp-level:0
\[
q\mapsto-q\ \ \text{on the ternary alphabet has exactly one fixed point,}\ q=0
\]

%dsp-level:0
\[
\qquad [\mathrm{EM.24}(1);\ \text{machine}]:
\]

%dsp-level:0
\[
\qquad \boxed{\ \textbf{the neutral class is charge-self-conjugate; every charged class is exchanged}\ }
\]

%dsp-level:0
\[
\text{the CP involution on charge}\ n\text{-tuples has exactly one fixed tuple,}\ (0,\dots,0),
\]

%dsp-level:0
\[
\ \text{and}\ W_n(+Q)
\]

%dsp-level:0
\[
= W_n(-Q)\ \text{exactly}
\]

%dsp-level:0
\[
\qquad [\text{machine: exhaustive at}\ n=2,3];
\]

%dsp-level:0
\[
(5)\quad\textbf{the forcing theorem (why the conjugate is not optional):}
\]

%dsp-level:0
\[
\qquad Q_{\mathrm{EM}}=\sum_{c\in U}q_{\mathrm{EM}}(c)=0\ \ \textbf{in every valid universe}
\]

%dsp-level:0
\[
\qquad [\mathrm{EM.11};\ \mathrm{D2.11};\ \mathrm{D1.3}]:
\]

%dsp-level:0
\[
\text{a materialized chain with}\ q(c)\ne0\ \text{requires a compensating read of exactly}\ -q(c);
\]

%dsp-level:0
\[
\qquad \#\{q\in\{-1,0,+1\}:\ q=-q(c)\}=1
\]

%dsp-level:0
\[
\qquad [\mathrm{BM.23}(4);\ \text{machine}];
\]

%dsp-level:0
\[
\textbf{the compensator census:}\qquad\begin{array}{l|l}\text{a compensating configuration must carry} & \text{the corpus supplies}\\\hline\text{the charge read}\ -q(c) & \textbf{exactly one derived map produces it:}\ \mathrm C\ [\mathrm{EM.24}(1)]\\\text{a dual colour attachment, to seal against}\ c\text{'s triplet} & \textbf{exactly one derived map produces it:}\ \mathrm{D3.8}\circ\mathrm{CP}\ [\mathrm{CT3.17}]\\\text{the opposite chiral ray} & \textbf{exactly one derived map produces it:}\ \hat\omega\ [\mathrm{FT.122}(4);\ \mathrm{A\text{-}BM.7}(b)]\\\text{a closed real-edge network with}\ \#J^+=\#J^- & \mathrm{CP}\ \text{preserves}\ |R|\ \text{and exchanges}\ \#J^+\ \text{with}\ \#J^-\ [\mathrm{BM.25}(2);\ \mathrm{BM.27}(2)]\\\text{the same}\ E_{\mathrm{int}}\ (\text{hence the same mass read}) & \textbf{derived:}\ (2)\text{'s seven}\ q\text{-exponent-zero monomials}\ [\mathrm{BM.28}(1)]\end{array}
\]

%dsp-level:0
\[
\boxed{\ \textbf{the CP-conjugate is the unique derived configuration that balances a charged chain: it is forced out of the ledger's own balance,}\ \neg\textbf{added}\ }
\]

%dsp-level:0
\[
\qquad [\text{machine: the five-row census}];
\]

%dsp-level:0
\[
(6)\quad\textbf{the Dirac parallel, as an inventory statement:}
\]

%dsp-level:0
\[
\qquad \textbf{the three ingredients:}\ (\mathrm i)\ \text{the forward-only tick}\ [\mathrm{A6};\ (1)(\mathrm i)];
\]

%dsp-level:0
\[
(\mathrm{ii})\ \text{the}\ 16\text{-matrix pair arena}\ \mathcal K_1,
\]

%dsp-level:0
\[
\ \dim\mathcal K_1
\]

%dsp-level:0
\[
= 16
\]

%dsp-level:0
\[
= C_{\mathrm{cap}}\ [\mathrm{FT.142}(2);\ \mathrm{CT0.14};\ \mathrm{CT0.12}];
\]

%dsp-level:0
\[
(\mathrm{iii})\ \text{the total-charge-zero balance}\ [\mathrm{EM.11};\ \mathrm{D2.11}]:
\]

%dsp-level:0
\[
(\mathrm i)+(\mathrm{iii})\implies\text{the compensator is a configuration and}\ \neg\text{an inverse evolution}\ [(4)\text{'s scope clause}];
\]

%dsp-level:0
\[
\qquad (\mathrm{ii})\implies\text{the compensator lives in the same}\ 16\text{-arena, on the}\ \omega\text{-image representation}\ \{0,-1\}\ [\mathrm{EW.36};\ \mathrm{A\text{-}BM.7}(b)];
\]

%dsp-level:0
\[
\boxed{\ \textbf{adding the forward-only tick to the}\ 16\textbf{-matrix pair arena forces the conjugate constituent, exactly as a negative-energy branch forces its conjugate}\ }
\]

%dsp-level:0
\[
\qquad [\text{the parallel is structural:}\ \neg\text{an external history is imported},\ \neg\text{a measured mass is compared}\ [\mathrm{P5}]];
\]

%dsp-level:0
\[
\textbf{the scope of the parallel:}
\]

%dsp-level:0
\[
\qquad \neg\text{a Dirac operator,}\ \neg\text{a Lorentz spinor and}\ \neg\text{a negative-energy spectrum is constructed}\ [\mathrm{BM.6}(8):\ \text{the frame section and a global Dirac-type operator remain unconstructed};\ \mathrm{FT.93}(3)\ \text{the pincer}];
\]

%dsp-level:0
\[
(7)\quad\textbf{the no-annihilation row:}
\]

%dsp-level:0
\[
\qquad \neg\exists\ \text{an annihilation channel, a pair-production cross-section or an inter-sector transition operator}\ [\mathrm{Def\text{-}BM.7}(5);\ \mathrm{BM.27}(7);\ \mathrm{BM.29}(5)]:
\]

%dsp-level:0
\[
\textbf{the two sectors are amplitude-equal configuration orbits; their interaction is}\ \neg\textbf{a theorem of this corpus}
\]

%dsp-level:0
\[
\qquad [\mathrm{BM.27}(7)\ \text{re-cited},\ \neg\text{extended}];
\]

%dsp-level:0
\[
\text{the CP-pair creation weight of}\ \mathrm{BM.110}(1)\ \text{is a flicker weight,}\ \neg\text{a cross-section}\ [\mathrm{BM.44}(5);\ \mathrm{D1.8}];
\]

%dsp-level:0
\[
(8)\quad\textbf{status:}
\]

%dsp-level:0
\[
\qquad \textbf{derived:}\ (1)\text{--}(5);
\]

%dsp-level:0
\[
\qquad \textbf{derived as a structural parallel with an explicit scope:}\ (6);
\]

%dsp-level:0
\[
\qquad \textbf{recorded as absent:}\ (7);
\]

%dsp-level:0
\[
\neg\text{a declaration},
\]

%dsp-level:0
\[
\quad\neg\text{a P4 addition},
\]

%dsp-level:0
\[
\quad\neg\text{a new dimensionful scale}\ [\mathrm{P1}],
\]

%dsp-level:0
\[
\quad\neg\text{a measured value}\ [\mathrm{P5}],
\]

%dsp-level:0
\[
\quad\neg\text{a species identification}\ [\mathrm{Def\text{-}BM.3}(7)],
\]

%dsp-level:0
\[
\quad\neg\text{a numerical coupling}\ [\mathrm{EM.25}(7)];
\]

%dsp-level:0
\[
\qquad\text{§Cal-BM stays empty except for}\ \mathrm{A\text{-}BM.9}\text{'s logged slot}.
\]

\textbf{Rmk (locus discipline on clause (6)).} The parallel is a statement about what forces what inside one inventory: a balance condition plus a forward-only order plus a fixed arena leave exactly one derived configuration that satisfies all three. It is not a claim that the corpus reproduces any external equation, and no clause identifies the conjugate constituent with any measured object — the name \textit{antiquark} belongs to the identification layer and occurs in no \textbf{Thm} clause here [Def-BM.3(7)].

\textbf{Req.} A6, A7, D0.8, D1.3, D1.4, D1.8, D2.6, D2.7, D2.11, D3.3, D3.4, D3.8 (Parts 1–2); CT0.12, CT0.14, CT2.8(1), CT3.17, EM.9, EM.11, EM.24(1), EM.24(2), EM.24(3), EM.24(4), EM.24(5), EM.25(7), EW.36, Def-EM.21 (Part 1, Part 2); P1, P4, P5 (Parts 1–2); L2 (Part 1); Def-H.33, FT.22, FT.93(3), FT.122(4), FT.142(2), FT.144(2), FT.144(4) (Part 3); Census FT.49 (row-wise); Def-BM.3(7), Def-BM.7(1), Def-BM.7(2), Def-BM.7(3), Def-BM.7(5), Def-BM.15, BM.2(1), BM.6(8), BM.23(4), BM.25(1), BM.25(2), BM.25(3), BM.25(4), BM.25(6), BM.27(2), BM.27(3), BM.27(5), BM.27(7), BM.28(1), BM.28(5), BM.28(6), BM.29(5), BM.51(1), BM.51(2), BM.73(4), A-BM.7 (this part); exact $\mathbb Q(i)$ phase algebra on the fourth-root alphabet, exhaustive ternary-alphabet and $n$-tuple censuses, Laurent-monomial $q$-exponent censuses [math; machine: \texttt{checks/\allowbreak{}verify\_\allowbreak{}H\_\allowbreak{}antimatter\_\allowbreak{}resonance\_\allowbreak{}spectrum.py} §A].

\textbf{Pf.}

%dsp-level:0
\[
(1):\  \mathrm{EM.24}(1)\text{--}(4)\ \text{rows re-cited (frozen, machine-verified at Tier F.33)};
\]

%dsp-level:0
\[
\qquad (\mathrm i),(\mathrm{ii}):\  \mathrm{A6},\ \mathrm{A7}\ \text{verbatim}
\]

%dsp-level:0
\[
(\mathrm{iii}):\  V'-V=-[1-(1-\tfrac\gamma4)^2]V=-\tfrac{\gamma(8-\gamma)}{16}V\ [\mathrm{L2}];
\]

%dsp-level:0
\[
\qquad \gamma=1:\ 1-\tfrac14=\tfrac34,\ (\tfrac34)^2=\tfrac9{16},\ 1-\tfrac9{16}=\tfrac7{16},\ \gamma(8-\gamma)=7\ [\text{machine: exact rationals}]
\]

%dsp-level:0
\[
\text{the inadmissibility row:}\ \mathrm{FT.144}(4)(\mathrm{ii})\ \text{and}\ \mathrm{BM.25}(2)\ \text{re-cited}\ [\text{machine: the pure reversal leads with}\ J^-,\ \text{outside D1.2's domain}]
\]

%dsp-level:0
\[
(2):\  q\mapsto\tfrac12(-\sigma_s+\sigma_t)=-q\ [\mathrm{D3.4}];
\]

%dsp-level:0
\[
\qquad \hat\omega\hat\chi\hat\omega^{-1}=-\hat\chi,\ \hat\omega P_L\hat\omega^{-1}=P_R\ [\mathrm{FT.122}(4);\ \mathrm{BM.73}(4);\ \text{machine over}\ \mathbb Q(i)]
\]

%dsp-level:0
\[
\text{the seven fixed reads:}\ \Delta x_{\mathrm B},
\]

%dsp-level:0
\[
\ V_{\mathrm{eff}},
\]

%dsp-level:0
\[
\ \rho,
\]

%dsp-level:0
\[
\ \kappa^2\rho,
\]

%dsp-level:0
\[
\ r_s,
\]

%dsp-level:0
\[
\ (\Delta x_{\mathrm B})^2\kappa^2\rho,
\]

%dsp-level:0
\[
\ \Delta x_{\mathrm B}r_s\ \text{each carry charge exponent}\ 0\ [\mathrm{BM.28}(1);\ \text{machine: monomial census}]
\]

%dsp-level:0
\[
\implies q
\]

%dsp-level:0
\[
\mapsto -q\ \text{fixes each}
\]

%dsp-level:0
\[
(3):\  \mathrm{EM.24}(5)\text{'s Pf row}\ \bar\Gamma\mapsto e^{i\pi\Theta}\Gamma\ \text{under}\ \mathrm P\circ\mathrm T,\ \text{and}\ S\mapsto(-1)^{|E|}S\ \text{under}\ \mathrm C;
\]

%dsp-level:0
\[
\qquad \mathrm{CP}=\mathrm{CPT}\circ\mathrm T\ \text{at}\ \mathrm T^2=\mathrm{id}\ [\mathrm{EM.24}(4)]
\]

%dsp-level:0
\[
w(\mathrm{CP}\,C)=w(\mathrm{CPT}\,\mathrm T\,C)=(-1)^{|E|}e^{i\pi\Theta}w(\mathrm T\,C)=(-1)^{|E|}e^{i\pi\Theta}w(C)^\ast;
\]

%dsp-level:0
\[
\qquad |w(C)|=1\ [\mathrm{Def\text{-}H.32}(2):\ |w(R,\sigma)|=1]\implies w(C)^\ast=(-1)^{|E|}w(C)^{-1}\cdot|w|^2
\]

%dsp-level:0
\[
\text{on the sign-real locus}\ w\in\mathbb R\ [\mathrm{EM.24}(5)\ \text{on-shell}]
\]

%dsp-level:0
\[
\implies w^\ast
\]

%dsp-level:0
\[
= w\implies\mu_{\mathrm{CP}}
\]

%dsp-level:0
\[
= (-1)^{|E|}e^{i\pi\Theta}(-1)^{|E|}
\]

%dsp-level:0
\[
= e^{i\pi\Theta}\ [\text{machine:}\ (-1)^{2|E|}=1];
\]

%dsp-level:0
\[
\qquad \Theta\in2\mathbb Z\implies e^{i\pi\Theta}=i^{2\Theta}=1\ [\text{machine}]
\]

%dsp-level:0
\[
(4):\  \mathrm{BM.25}(4)\text{'s six-sector census re-cited}\ [\text{machine}];
\]

%dsp-level:0
\[
\qquad q\mapsto-q\ \text{fixed points on}\ \{-1,0,+1\}:\ \{0\}\ [\text{machine}];
\]

%dsp-level:0
\[
\qquad n\text{-tuples: exhaustive at}\ n
\]

%dsp-level:0
\[
= 2\ (9\ \text{tuples},\ 5\ \text{orbits},\ 1\ \text{fixed})\ \text{and}\ n
\]

%dsp-level:0
\[
= 3\ (27\ \text{tuples},\ 14\ \text{orbits},\ 1\ \text{fixed})\ [\text{machine}]
\]

%dsp-level:0
\[
W_n(+Q)
\]

%dsp-level:0
\[
= W_n(-Q):\ \text{the map}\ (q_1,\dots,q_n)
\]

%dsp-level:0
\[
\mapsto (-q_1,\dots,-q_n)\ \text{is a weight-preserving bijection between the two class fibres}\ [m_{-q}=m_q\ [\mathrm{FT.106}(3\mathrm a):\ m_q=(2,1,1)];\ \text{machine}]
\]

%dsp-level:0
\[
(5):\  \mathrm{EM.11},\ \mathrm{D2.11}\ \text{re-cited};
\]

%dsp-level:0
\[
\qquad \#\{q:\ q=-\Delta q\}=1\ [\mathrm{BM.23}(4)\ \text{pattern};\ \text{machine}];
\]

%dsp-level:0
\[
\qquad \text{the five-row census assembled from its own tags}\ [\text{machine}]
\]

%dsp-level:0
\[
(6):\  \dim\mathcal K_1=16=C_{\mathrm{cap}}\ [\mathrm{FT.142}(2);\ \mathrm{CT0.14}];
\]

%dsp-level:0
\[
\qquad \omega\text{-image representation}\ \{0,-1\}\ [\mathrm{EW.36};\ \mathrm{A\text{-}BM.7}(b)];
\]

%dsp-level:0
\[
\qquad \text{the unconstructed rows re-cited}\ [\mathrm{BM.6}(8);\ \mathrm{FT.93}(3)]
\]

%dsp-level:0
\[
(7):\  \text{the inventory scan of}\ \mathrm{BM.27}(7)\ \text{re-cited};\qquad(8):\  \text{the standing-rule scans}\ [\text{machine}].\ \blacksquare
\]

\textbf{Phys.} Of the three reversals the corpus derives, two run and one does not. Charge conjugation and parity are maps the ledger can apply to a configuration; time reversal is a map it can apply to a record but never to a step, because the tick has no inverse, the ledger only grows, and the Lyapunov function only contracts — at sharp crystallization by exactly seven sixteenths per tilt. Composing the two that run gives the conjugation that matters, and its action is exact and cheap: the charge read flips, the scar flips, the chiral rays exchange, the triplet attachment becomes a dual one, the weight picks up a fourth root of unity that is exactly one on the sealed junction's own phase class, and the seven geometric reads do not move at all, because not one of them carries a charge exponent. So the conjugate is mass-degenerate with the original by structure rather than by luck. That it must exist is not an addition but a consequence: total charge vanishes in every valid universe, a charged chain therefore needs a compensator, the alphabet contains exactly one value that compensates it, and of the derived maps exactly one produces a configuration carrying that value with a dual attachment, the opposite ray, and the same edge count. The orbit census is the same statement at the level of classes — one self-conjugate neutral class, every charged class exchanged, one fixed charge tuple at every level and equal weights on the two charged classes. The structural parallel is then exact in its own terms: a balance, a forward-only order, and a fixed sixteen-dimensional arena leave one derived compensator, which is the shape of the argument that forces a conjugate branch, and the tier says no more than that — no Dirac operator, no spinor, no annihilation channel, and no name outside the identification layer.

\end{theorem}

\begin{theorem}[{\texorpdfstring{BM.109 — The CP-pair as the vacuum fluctuation: the creation weight, its exact range, and the degeneracy rows}{BM.109 - The CP-pair as the vacuum fluctuation: the creation weight, its exact range, and the degeneracy rows}}]\label{tag:BM-109}
\textbf{Thm.}

%dsp-level:0
\[
(1)\quad\textbf{the pair's charge and mass rows:}
\]

%dsp-level:0
\[
\qquad Q_{\text{pair}}=q(c)+q(c^{\mathrm{CP}})=q-q=0
\]

%dsp-level:0
\[
\qquad [\mathrm{Def\text{-}BM.15}(3);\ \text{machine: exhaustive over}\ q\in\{-1,0,+1\}]:
\]

%dsp-level:0
\[
E_{\mathrm{int}}^{\text{pair}}=2\,E_{\mathrm{int}}(c)
\]

%dsp-level:0
\[
\qquad [\text{additivity of the internal count on a disjoint pair};\ \mathrm{D0.8};\ \mathrm{BM.21}(1)\ \text{pattern}];
\]

%dsp-level:0
\[
\boxed{\ m_{\mathrm{hid}}^{\text{pair}}=\frac{2E_{\mathrm{int}}(c)\ln2}{\mathfrak h}=2\,m_{\mathrm{hid}}(c),\qquad\Delta x_{\mathrm B}^{\text{pair}}=\frac{\mathfrak h}{2E_{\mathrm{int}}(c)\ln2}=\frac{\Delta x_{\mathrm B}(c)}2\ }
\]

%dsp-level:0
\[
\qquad [\mathrm{FT.44}(4);\ \mathrm{Def\text{-}BM.4}(1);\ \text{machine: Laurent identities}];
\]

%dsp-level:0
\[
\textbf{the binding read of the pair count vanishes exactly:}
\]

%dsp-level:0
\[
\qquad E^{\mathrm{info}}_{\mathrm{bind}}(\text{pair})=\bigl(2E_{\mathrm{int}}-(E_{\mathrm{int}}+E_{\mathrm{int}})\bigr)\tfrac{\ln2}{\mathfrak h}=0
\]

%dsp-level:0
\[
\qquad [\mathrm{BM.21}(1);\ \mathrm{BM.54}(1);\ \text{machine}];
\]

%dsp-level:0
\[
(2)\quad\textbf{the creation weight is the Janus resolvent's own:}
\]

%dsp-level:0
\[
\qquad \sum_n\langle\varnothing|\hat U^n|\varnothing\rangle w^n=\frac1{1-w^2x_+x_-}\ \ \text{in}\ \mathbb Q[w,x_+,x_-]/(w^{12})
\]

%dsp-level:0
\[
\qquad [\mathrm{FT.108}(5);\ \mathrm{BM.23}(7);\ \mathrm{BM.44}(1)]:
\]

%dsp-level:0
\[
\boxed{\ \mathcal W_{\text{pair}}\ :=\ x_+x_-\ =\ 4\varepsilon_q(1-4\varepsilon_q)\ =\ \frac14-16\Bigl(\varepsilon_q-\frac18\Bigr)^{\!2}\ \le\ \frac14,\qquad\text{equality}\iff\varepsilon_q=\frac18\ }
\]

%dsp-level:0
\[
\qquad [\mathrm{BM.19}(2);\ \mathrm{BM.44}(1);\ \text{machine: the SOS completion}\ +\ \text{the grid census over}\ \varepsilon_q=k/256];
\]

%dsp-level:0
\[
\textbf{the two endpoints:}\qquad\varepsilon_q=0\implies x_-=0\implies\mathcal W_{\text{pair}}=0\ \ \text{exactly}\qquad[\mathrm{BM.19}(3);\ \text{machine}];
\]

%dsp-level:0
\[
\varepsilon_q\to\tfrac14^-\implies x_+\to0\implies\mathcal W_{\text{pair}}\to0\qquad[\mathrm{CT0.8}:\ \varepsilon_q\in[0,\tfrac14);\ \text{machine}]:
\]

%dsp-level:0
\[
\boxed{\ \mathcal W_{\text{pair}}>0\iff\varepsilon_q\in\bigl(0,\tfrac14\bigr),\qquad\max\mathcal W_{\text{pair}}=\frac14\ \ \text{at}\ \ \varepsilon_q=\frac18\ }
\]

%dsp-level:0
\[
\qquad [\text{machine: exact rationals on a}\ 256\text{-point grid}];
\]

%dsp-level:0
\[
(3)\quad\textbf{the one-door row:}
\]

%dsp-level:0
\[
\qquad \text{the vacuum sector is all-virtual and unregistered with no accessible pole at any rung}\ [\mathrm{BM.96}(4);\ \mathrm{FT.144}(4);\ \mathrm{CT4.2}];
\]

%dsp-level:0
\[
\qquad \text{the materialized locus requires}\ |E_{\mathcal R}|\ge2\ [\mathrm{Def\text{-}H.4};\ \mathrm{BM.13}(3);\ \mathrm{BM.70}(2)]:
\]

%dsp-level:0
\[
\boxed{\ \textbf{the CP-pair materializes and propagates forward-only exactly where the flicker is present}\ —\ \text{at the same threshold the ensemble integral is taken over}\ }
\]

%dsp-level:0
\[
\qquad [\mathrm{BM.96}(6)\ \text{re-cited at the pair level};\ \text{machine}];
\]

%dsp-level:0
\[
\Gamma_{\mathrm{tr}}(\varepsilon_q{=}0)=0
\]

%dsp-level:0
\[
\qquad [\mathrm{BM.23}(8);\ \mathrm{BM.44}(2)]:
\]

%dsp-level:0
\[
\qquad \textbf{a non-flickering ledger creates no pairs and decays nothing};
\]

%dsp-level:0
\[
(4)\quad\textbf{the degeneracy rows (the pair is one orbit, not two structures):}
\]

%dsp-level:0
\[
\qquad |w(c^{\mathrm{CP}})|^2=|w(c)|^2\ \ \forall c
\]

%dsp-level:0
\[
\qquad [\mathrm{EM.24}(5):\ \text{unconditional};\ \text{machine}];
\]

%dsp-level:0
\[
m_{\mathrm{hid}}(c^{\mathrm{CP}})=m_{\mathrm{hid}}(c),
\]

%dsp-level:0
\[
\qquad D_S(c^{\mathrm{CP}})=D_S(c),
\]

%dsp-level:0
\[
\qquad m^2_{\mathrm{pole}}(c^{\mathrm{CP}})=m^2_{\mathrm{pole}}(c)
\]

%dsp-level:0
\[
\qquad [\mathrm{BM.108}(2);\ \mathrm{Def\text{-}BM.14}(3):\ D_S\ \text{is branch data,}\ \neg\text{charge data};\ \text{machine}];
\]

%dsp-level:0
\[
(\mathrm{CP})^2=\mathrm{id}\implies\textbf{the pair carries no independent datum: it is one orbit of size two}
\]

%dsp-level:0
\[
\qquad [\mathrm{BM.27}(6);\ \mathrm{EM.24}(4)];
\]

%dsp-level:0
\[
\text{the two rows contribute equally to the norm}\ \langle A|A\rangle=\sum_C|w(C)|^2
\]

%dsp-level:0
\[
\qquad [\mathrm{CT2.8}(3);\ \mathrm{BM.25}(7);\ \text{machine}];
\]

%dsp-level:0
\[
(5)\quad\textbf{the no-annihilation row stands:}
\]

%dsp-level:0
\[
\qquad \neg\exists\ \text{an annihilation channel or a cross-section between}\ \mathrm{EM}^+\ \text{and}\ \mathrm{EM}^-\ [\mathrm{BM.27}(7);\ \mathrm{Def\text{-}BM.7}(5)]:
\]

%dsp-level:0
\[
\textbf{what is derived is the pair's existence, its neutrality, its degeneracy and its creation weight}\ —\ \neg\text{its interaction};
\]

%dsp-level:0
\[
(6)\quad\textbf{status:}\qquad\textbf{derived:}\ (1)\text{--}(4);\qquad\textbf{recorded as absent:}\ (5);
\]

%dsp-level:0
\[
\neg\text{a declaration},
\]

%dsp-level:0
\[
\quad\neg\text{a P4 addition},
\]

%dsp-level:0
\[
\quad\neg\text{a new dimensionful scale}\ [\mathrm{P1}],
\]

%dsp-level:0
\[
\quad\neg\text{a measured value}\ [\mathrm{P5}],
\]

%dsp-level:0
\[
\quad\neg\text{a numerical coupling}\ [\mathrm{EM.25}(7)],
\]

%dsp-level:0
\[
\quad\neg\text{a lifetime in external units}\ [\mathrm{BM.44}(5)];
\]

\textbf{Req.} A4, A6, D0.4, D0.8, D1.8, D2.9 (Parts 1–2); CT0.8, CT2.8(3), CT4.2, EM.24(4), EM.24(5), EM.25(7) (Part 1, Part 2); P1, P4, P5 (Parts 1–2); Def-H.4, Def-H.32(2), FT.108(5), FT.144(4), FT.44(4) (Part 3, Census); Def-BM.3(7), Def-BM.4(1), Def-BM.7(1), Def-BM.7(5), Def-BM.14(3), Def-BM.15, BM.13(3), BM.19(2), BM.19(3), BM.21(1), BM.23(7), BM.23(8), BM.25(7), BM.27(6), BM.27(7), BM.44(1), BM.44(2), BM.44(5), BM.54(1), BM.70(2), BM.96(4), BM.96(6), BM.108(2) (this part); exact rational resolvent identities, the SOS parabola completion, exhaustive rational grid censuses, Laurent-monomial assembly [math; machine: \texttt{checks/\allowbreak{}verify\_\allowbreak{}H\_\allowbreak{}antimatter\_\allowbreak{}resonance\_\allowbreak{}spectrum.py} §A].

\textbf{Pf.}

%dsp-level:0
\[
(1):\  q+(-q)=0\ [\text{machine: three cases}];
\]

%dsp-level:0
\[
\qquad E_{\mathrm{int}}^{\text{pair}}=E_{\mathrm{int}}+E_{\mathrm{int}}\ \text{on a disjoint union of directed-edge sets}\ [\mathrm{D0.8}];
\]

%dsp-level:0
\[
\qquad m_{\mathrm{hid}}=E_{\mathrm{int}}\ln2/\mathfrak h\ \text{linear}\ [\mathrm{FT.44}(4);\ \text{machine}]
\]

%dsp-level:0
\[
\Delta x_{\mathrm B}^{\text{pair}}=\mathfrak h/(2E_{\mathrm{int}}\ln2)=\Delta x_{\mathrm B}/2\ [\text{machine: one Laurent monomial}];
\]

%dsp-level:0
\[
\qquad E^{\mathrm{info}}_{\mathrm{bind}}=0\cdot\tfrac{\ln2}{\mathfrak h}\ [\mathrm{BM.21}(1)\ \text{pattern};\ \text{machine}]
\]

%dsp-level:0
\[
(2):\  \mathrm{FT.108}(5)\text{'s resolvent re-cited};
\]

%dsp-level:0
\[
\qquad x_+=1-4\varepsilon_q,\ x_-=4\varepsilon_q\implies x_+x_-=4\varepsilon_q(1-4\varepsilon_q)\ [\mathrm{BM.19}(2);\ \text{machine}]
\]

%dsp-level:0
\[
4\varepsilon_q(1-4\varepsilon_q)=\tfrac14-16(\varepsilon_q-\tfrac18)^2\ [\text{machine: expansion}];
\]

%dsp-level:0
\[
\qquad \text{the grid census}\ \varepsilon_q=k/256,\ k=0..63:\ \le\tfrac14,\ =\tfrac14\ \text{only at}\ k=32\ [\text{machine}]
\]

%dsp-level:0
\[
\varepsilon_q=0\implies x_-=0\implies x_+x_-=0;
\]

%dsp-level:0
\[
\qquad \varepsilon_q\to\tfrac14^-
\]

%dsp-level:0
\[
\implies x_+\to0
\]

%dsp-level:0
\[
\implies x_+x_-\to0\ [\text{machine: exact at the two endpoints of the admissible range}\ [\mathrm{CT0.8}]]
\]

%dsp-level:0
\[
(3):\  \mathrm{BM.96}(4),
\]

%dsp-level:0
\[
(6)\text{'s rows transported to the pair}\ [\text{the pair's members are the two states of that row}];
\]

%dsp-level:0
\[
\qquad \Gamma_{\mathrm{tr}}(0)=0\ [\mathrm{BM.44}(2);\ \text{machine}]
\]

%dsp-level:0
\[
(4):\  |w(\mathrm{CP}\,C)|^2=|w(C)|^2\ \text{from}\ |\mu_{\mathrm{CP}}|=|e^{i\pi\Theta}|=1\ [\mathrm{BM.108}(3);\ \text{machine}];
\]

%dsp-level:0
\[
\qquad D_S\ \text{is a function of}\ (\mathcal F_g,\beta)\ \text{alone}\ [\mathrm{Def\text{-}BM.14}(3);\ \text{machine: no charge symbol occurs}]
\]

%dsp-level:0
\[
(\mathrm{CP})^2=\mathrm C^2\mathrm P^2=\mathrm{id}\ [\mathrm{EM.24}(4)];
\]

%dsp-level:0
\[
\qquad (5),(6):\  \text{the inventory scan and the standing-rule scans}\ [\text{machine}].\ \blacksquare
\]

\textbf{Phys.} A pair made of a chain and its conjugate is neutral because the conjugate's charge is the negative of the chain's, and it is mass-degenerate because none of the seven geometric reads carries a charge exponent and the defect read is branch data rather than charge data. Its informational mass is exactly twice a member's and its width exactly half, with the count-route binding identically zero — a binding energy cannot be an integer summation, and the pair is the cleanest case of it. The weight for making one is not declared: it is the Janus resolvent's own $x_+x_-$, a parabola in the single crystallization threshold with maximum one quarter at one eighth and zeros at both ends of the admissible range. So the ledger creates pairs exactly where it flickers, and nowhere else — at the sharp threshold the dissolution weight vanishes, no pair is made, and nothing decays. What the tier does not give the pair is a fate: no annihilation channel and no inter-sector cross-section exist in the inventory, and the pair stays one orbit of size two carrying no datum its member does not.

\end{theorem}

\begin{definition}[{\texorpdfstring{Def-BM.16 — The two-chain state space (the construction BM.6(2) specifies)}{Def-BM.16 - The two-chain state space (the construction BM.6(2) specifies)}}]\label{tag:Def-BM-16}
\textbf{Def.} A construction discharging the $\mathrm{BM.6}(2)$ specification — \textit{a state space on ordered pairs of bound chains with an exchange involution} — on $\mathrm{Def\text{-}H.33}$'s orientation space and $\mathrm{Def\text{-}H.38}$'s winding carrier; no new axiom, no new datum, no new parameter, no external state.

%dsp-level:0
\[
(1)\quad\textbf{the single-chain factors (re-cited, both derived):}
\]

%dsp-level:0
\[
\qquad \mathcal H_\omega(c)=\operatorname{span}_{\mathbb C}\{|c\rangle,\ |\omega(c)\rangle\}\cong\mathbb C^2,
\]

%dsp-level:0
\[
\qquad \hat\omega^2=\mathrm{id},\quad\hat\omega^\dagger=\hat\omega,\quad\hat\chi=i\bigl(|c\rangle\langle\omega(c)|-|\omega(c)\rangle\langle c|\bigr),\quad P_{L/R}=\tfrac12(\mathrm{id}\pm\hat\chi)
\]

%dsp-level:0
\[
[\mathrm{Def\text{-}H.33};\ \mathrm{FT.122}(1)\text{--}(3),(6)];
\]

%dsp-level:0
\[
\qquad \Psi(c)=\mathrm{Sym}^{|k(c)|}(\mathcal H_{\mathsf v}),
\]

%dsp-level:0
\[
\qquad \dim\Psi(c)=|k(c)|+1,
\]

%dsp-level:0
\[
\qquad \mathcal H_{\mathsf v}\cong\mathbb C^2
\]

%dsp-level:0
\[
\qquad [\mathrm{Def\text{-}H.38};\ \mathrm{EW.32};\ \text{machine}];
\]

%dsp-level:0
\[
(2)\quad\textbf{the two-chain state space:}\qquad\boxed{\ \mathcal H_2(c_1,c_2)\ :=\ \mathcal H_\omega(c_1)\otimes\Psi(c_1)\ \otimes\ \mathcal H_\omega(c_2)\otimes\Psi(c_2)\ }
\]

%dsp-level:0
\[
\dim\mathcal H_2=4\,\bigl(|k(c_1)|+1\bigr)\bigl(|k(c_2)|+1\bigr)\qquad[\text{machine: exact integers}];
\]

%dsp-level:0
\[
\textbf{the meson specialization (the CP-pair):}
\]

%dsp-level:0
\[
\qquad c_2=c_1^{\mathrm{CP}}\implies|k(c_2)|=|k(c_1)|=:|k|:
\]

%dsp-level:0
\[
\qquad \boxed{\ \mathcal H_{\text{meson}}(c)\ :=\ \mathcal H_2\bigl(c,\ c^{\mathrm{CP}}\bigr),\qquad\dim\mathcal H_{\text{meson}}=4\,(|k|+1)^2\ }
\]

%dsp-level:0
\[
\textbf{constitutive clause:}\ \ \text{the space is the tensor product of two derived single-chain spaces, with the}\ \mathrm{CP}\text{-identification}\ |c^{\mathrm{CP}}\rangle\leftrightarrow|\omega(c)\rangle\ \text{of}\ \mathrm{D2.6};
\]

%dsp-level:0
\[
\textbf{non-derived clause:}\ \ \text{the factor ordering}\ (1,2)\ \text{is a convention;}\ \text{the identification of the second factor with the first's}\ \omega\text{-image is the}\ \mathrm{D2.6}\ \text{vocabulary clause, applied};
\]

%dsp-level:0
\[
(3)\quad\textbf{the inner product and the norm layer:}
\]

%dsp-level:0
\[
\qquad \langle A|A\rangle=\sum_C|w(C)|^2\ge0
\]

%dsp-level:0
\[
\qquad [\mathrm{CT2.8}(3)];
\]

%dsp-level:0
\[
\qquad |w(R,\sigma)|=1\ \ \text{for every valid pair}
\]

%dsp-level:0
\[
\qquad [\mathrm{Def\text{-}H.32}(2);\ \mathrm{FT.68}\ \mathrm{Rmk}(6\mathrm c)]:
\]

%dsp-level:0
\[
\text{the orientation doublet carries}\ \mathrm{Def\text{-}H.33}\text{'s orthonormal basis}\ \{|c\rangle,|\omega(c)\rangle\};
\]

%dsp-level:0
\[
\qquad \text{the winding carrier carries the}\ \mathrm{Sym}^{|k|}\text{-inherited product of}\ \mathcal H_{\mathsf v}\ [\mathrm{EW.32}];
\]

%dsp-level:0
\[
\textbf{derived:}\ \ \mathcal H_2\ \textbf{is a genuine state space}\ —\ \text{an inner product, a norm layer and a finite exact dimension}\ —\ \neg\text{a counting space}
\]

%dsp-level:0
\[
\qquad [\mathrm{Def\text{-}H.41}\ \text{Species};\ \mathrm{Def\text{-}H.42}:\ \text{the bar this clause discharges}];
\]

%dsp-level:0
\[
(4)\quad\textbf{the exchange involution (the specification's second member):}
\]

%dsp-level:0
\[
\qquad \hat S_{12}\ :=\ \text{the canonical flip on}\ \mathcal H_2,
\]

%dsp-level:0
\[
\qquad \hat S_{12}(v\otimes w)=w\otimes v,
\]

%dsp-level:0
\[
\qquad \hat S_{12}^2=\mathrm{id}
\]

%dsp-level:0
\[
\qquad [\mathrm{A3}:\ S^2=\mathbb1;\ \text{machine}];
\]

%dsp-level:0
\[
\hat\mu_{\mathrm{CP}}\ :=\ \text{the CP-transport phase acting by the scalar}\ \mu_{\mathrm{CP}}(c)=e^{i\pi\Theta(c)}\ \ \text{on the}\ \Psi\text{-factors}
\]

%dsp-level:0
\[
\qquad [\mathrm{BM.108}(3);\ \text{machine}];
\]

%dsp-level:0
\[
\boxed{\ \hat P_{12}\ :=\ \hat\mu_{\mathrm{CP}}\circ\hat S_{12}\ }
\]

%dsp-level:0
\[
\qquad \textbf{constitutive clause:}\ \ \text{the exchange is the flip composed with the transport its own identification requires};
\]

%dsp-level:0
\[
(5)\quad\textbf{the permutation-sign read (the specification's third member):}
\]

%dsp-level:0
\[
\qquad \boxed{\ \varepsilon(c)\ :=\ \mu_{\mathrm{CP}}(c)\big|_{\vee^2}\ =\ i^{\,\Theta(c)}\ =\ (-1)^{|k(c)|}\ }
\]

%dsp-level:0
\[
\qquad [\mathrm{BM.112}(2);\ \text{machine}]:
\]

%dsp-level:0
\[
\textbf{constitutive clause:}\ \ \text{the statistics sign of a chain is the exchange operator's eigenvalue on the}\ \vee^2\text{-sector of the CP-pair's winding carrier};
\]

%dsp-level:0
\[
\textbf{non-derived clause:}\ \ \text{the choice of the}\ \vee^2\text{-sector as the read's locus}\ —\ \text{priced by}\ \mathrm{BM.112}(5)\text{'s four rows and}\ \neg\text{forced by any clause of this definition};
\]

%dsp-level:0
\[
(6)\quad\textbf{the specification-discharge clause:}
\]

%dsp-level:0
\[
\qquad \mathrm{BM.6}(2)\ \text{requires (i) a state space on ordered pairs of bound chains and (ii) an exchange involution:}
\]

%dsp-level:0
\[
(\mathrm i)\ \ \text{discharged by}\ (2),(3);\qquad(\mathrm{ii})\ \ \text{discharged by}\ (4)\ \text{with}\ \hat P_{12}^2=\mathrm{id}\ \text{derived}\ [\mathrm{BM.112}(1)];
\]

%dsp-level:0
\[
\textbf{the bar}\ \mathrm{FT.125}(5)\ \textbf{records is respected:}
\]

%dsp-level:0
\[
\qquad \varepsilon(c)\ \text{is}\ \neg\text{one of the four single-object}\ \mathbb Z_2\ \text{gradings of}\ \mathrm{BM.106}(4)\ \text{— it is a permutation sign on a two-chain space}\ [\mathrm{BM.113}(2)];
\]

%dsp-level:0
\[
(7)\quad\textbf{status and scope:}
\]

%dsp-level:0
\[
\qquad \textbf{derived:}\ (1),\ (3),\ (4)\text{'s involution rows},\ (5)\text{'s value};
\]

%dsp-level:0
\[
\qquad \textbf{constitutive:}\ (2)\text{'s product},\ (4)\text{'s composition},\ (5)\text{'s sector choice};
\]

%dsp-level:0
\[
\textbf{the scope inherited at every use:}\ \ \mathrm{Def\text{-}H.38}\text{'s nonnegative cone}\ k
\]

%dsp-level:0
\[
\ge 0\ \text{and its}\ k
\]

%dsp-level:0
\[
\ge 3\ \text{aspect bar}\ [\mathrm{Def\text{-}H.38}\ \text{Scope};\ \text{no aspect claim is made here}];
\]

%dsp-level:0
\[
\qquad \mathrm{Def\text{-}H.32}\text{'s}\ (\mathrm{S1})\text{--}(\mathrm{S4});
\]

%dsp-level:0
\[
\neg\text{a P4 addition},
\]

%dsp-level:0
\[
\quad\neg\text{a new dimensionful scale}\ [\mathrm{P1}],
\]

%dsp-level:0
\[
\quad\neg\text{a measured value}\ [\mathrm{P5}],
\]

%dsp-level:0
\[
\quad\neg\text{a species identification}\ [\mathrm{Def\text{-}BM.3}(7)];
\]

%dsp-level:0
\[
\qquad \text{§Cal-BM stays empty};
\]

%dsp-level:0
\[
\textbf{beyond this tier:}\ \ \text{the three-chain state space}\ \mathcal H_3\ \text{and its}\ S_3\ \text{action}\ \text{are}\ \neg\text{constructed}\ [\mathrm{BM.103}(4)\ \text{re-scoped},\ \neg\text{discharged}].
\]

\textbf{Rmk (why this is a construction and not a declaration).} Every ingredient is derived and frozen: the orientation doublet and its involution from the fermion era's interface layer, the winding carrier and its dimension from the multiplicity–winding coupling, the inner product from the norm layer, the flip from the pair axiom's own $S^2=\mathbb1$, and the transport phase from the CP character. What the definition adds is one product and one composition, and the corpus's own specification — BM.6(2), written as a build-spec and adopted nowhere — is what it discharges. The single non-derived row is a sector choice, and it is priced rather than hidden: BM.112(5) lists the four alternative loci and BM.113(2) shows the read is invariant under all four on the materialized locus.

\textbf{Req.} A3 (Part 1); D2.6, D2.9 (Parts 1–2); CT2.8(3) (Part 2); EW.32 (Part 2); P1, P4, P5 (Parts 1–2); Def-H.32(2), Def-H.33, Def-H.38, Def-H.41, Def-H.42, FT.68, FT.122(1), FT.122(2), FT.122(3), FT.122(6), FT.125(5) (Part 3); Def-BM.3(7), Def-BM.15, BM.6(2), BM.103(4), BM.106(4), BM.108(3) (this part); exact integer dimension arithmetic, exhaustive sector censuses [math; machine: \texttt{checks/\allowbreak{}verify\_\allowbreak{}H\_\allowbreak{}antimatter\_\allowbreak{}resonance\_\allowbreak{}spectrum.py} §B].

\textbf{Pf.}

%dsp-level:0
\[
\text{Definition block: the constitutive clauses are}\ (2)\text{'s tensor product},
\]

%dsp-level:0
\[
\ (4)\text{'s composition and}\ (5)\text{'s sector choice};
\]

%dsp-level:0
\[
\ \text{the non-derived clauses are the factor ordering of}\ (2)\ \text{and the sector locus of}\ (5);
\]

%dsp-level:0
\[
\ (1),
\]

%dsp-level:0
\[
(3),
\]

%dsp-level:0
\[
(4)\text{'s involution rows and}\ (5)\text{'s value are derived and machine-checked};
\]

%dsp-level:0
\[
\ (6),
\]

%dsp-level:0
\[
(7)\ \text{are scope records}.\ \blacksquare
\]

\textbf{Phys.} The corpus said for three tiers that a statistics has no referent, and listed exactly what one would need: a space of two chains, an operator that swaps them, and a sign to read off. This definition supplies all three and nothing more. The space is the tensor product of the two single-chain spaces the fermion era already built — the orientation doublet and the winding carrier — so it has an inner product from the norm layer, a finite exact dimension of four times the two carrier dimensions, and it is a state space rather than a counting space, which is the bar the species definitions set. The swap is the canonical flip, whose square is one because the pair axiom already said so, composed with the one phase the CP-identification forces: transporting a chain to its conjugate's slot multiplies by the fourth root of unity the CPT bridge derived. The sign is that phase evaluated on the symmetric sector, and on the materialized locus it is a genuine plus or minus. What remains open is the obvious next step and the definition says so: three chains, and the permutation group that acts on them, are not built here.

\end{definition}

\begin{theorem}[{\texorpdfstring{BM.110 — The meson seal: the $\delta$-sealed $2$-juncture, its uniqueness, and the charge class it occupies}{BM.110 - The meson seal: the -sealed -juncture, its uniqueness, and the charge class it occupies}}]\label{tag:BM-110}
\textbf{Thm.}

%dsp-level:0
\[
(1)\quad\textbf{the seal is unique at two attachments:}
\]

%dsp-level:0
\[
\qquad \text{partitions of a labelled}\ 2\text{-element attachment set into blocks of size}\ \in\{2,3\}:\quad\text{exactly one, the}\ \delta\text{-pair}
\]

%dsp-level:0
\[
\qquad [\mathrm{CT3.17a}(1),(2):\ g(2)=1;\ \text{machine}]:
\]

%dsp-level:0
\[
\delta^i{}_j\ \text{consumes one triplet and one dual}
\]

%dsp-level:0
\[
\qquad [\mathrm{CT3.17}];
\]

%dsp-level:0
\[
\qquad \varepsilon_{ijk}\ \text{requires}\ 3\ \text{fundamentals}\ [\mathrm{CT3.15};\ \mathrm{D3.9}]:
\]

%dsp-level:0
\[
\boxed{\ \textbf{meson}\ :=\ \delta^i{}_j\,c^j\,\bar c_i\ =\ \delta^i{}_j\,c^j\,(c_i)^{\mathrm{CP}}\ \ \text{is the unique sealed}\ 2\text{-juncture, and a lone attachment is never sealed}\ }
\]

%dsp-level:0
\[
\qquad [\mathrm{Def\text{-}BM.1};\ \mathrm{CT3.17a}(1):\ |\mathrm{Colour}(\Sigma)|=1\implies\mathcal C=0;\ \text{machine}];
\]

%dsp-level:0
\[
(2)\quad\textbf{the two-chain space is the first multi-chain state space of the corpus:}
\]

%dsp-level:0
\[
\qquad \mathcal H_{\text{meson}}(c)=\mathcal H_2\bigl(c,c^{\mathrm{CP}}\bigr),
\]

%dsp-level:0
\[
\qquad \dim=4(|k|+1)^2
\]

%dsp-level:0
\[
\qquad [\mathrm{Def\text{-}BM.16}(2);\ \text{machine}]:
\]

%dsp-level:0
\[
\textbf{the referent census}\ \mathrm{BM.106}(3)\ \textbf{listed on five rows is discharged on its first three:}\qquad\begin{array}{l|l}\text{a statistics requires} & \text{the corpus supplies}\\\hline\text{a multi-chain state space} & \textbf{constructed:}\ \mathcal H_2\ [\mathrm{Def\text{-}BM.16}(2),(3)]\\\text{an exchange operator} & \textbf{constructed:}\ \hat P_{12}\ [\mathrm{Def\text{-}BM.16}(4);\ \mathrm{BM.112}(1)]\\\text{a permutation sign} & \textbf{formed:}\ \varepsilon(c)\ [\mathrm{Def\text{-}BM.16}(5);\ \mathrm{BM.112}(2)]\\\text{a grade's}\ \mathbb Z_2\ \text{identified with the sign's} & \textbf{still barred, and not done:}\ \varepsilon\ \text{is a permutation sign on}\ \mathcal H_2\ [\mathrm{BM.113}(2);\ \mathrm{FT.125}(5);\ \mathrm{FT.89}(2)]\\\text{a tensor state on}\ \mathcal H_c & \textbf{unchanged:}\ \mathcal H_c\ \text{stays a counting space;}\ \mathcal H_2\ \text{is built on}\ \mathcal H_\omega\otimes\Psi,\ \neg\text{on}\ \mathcal H_c\ [\mathrm{Def\text{-}BM.16}(3);\ \mathrm{Def\text{-}H.41}\ \text{Species}]\end{array}
\]

%dsp-level:0
\[
\boxed{\ \textbf{residue}\ (\mathrm{R7})\ \textbf{is discharged at}\ n=2\ }
\]

%dsp-level:0
\[
\qquad [\mathrm{BM.106}(7);\ \mathrm{BM.107}(7);\ \text{machine: the three-row discharge against}\ \mathrm{BM.6}(2)\text{'s specification}];
\]

%dsp-level:0
\[
(3)\quad\textbf{the charge class the seal occupies:}
\]

%dsp-level:0
\[
\qquad \textbf{the CP-pair subset:}\quad Q=q+( -q)=0\ \ \text{for all three}\ q\in\{-1,0,+1\}
\]

%dsp-level:0
\[
\qquad [\mathrm{Def\text{-}BM.15}(3);\ \text{machine}]:
\]

%dsp-level:0
\[
\textbf{the full}\ 2\text{-chain transport census:}
\]

%dsp-level:0
\[
\qquad \bigl(2+z+z^{-1}\bigr)^2=z^{-2}(z+1)^4\implies W_2(Q)=\binom4{Q+2}=(1,\,4,\,\mathbf6,\,4,\,1)\ \ \text{for}\ Q=-2,\dots,+2,
\]

%dsp-level:0
\[
\qquad \sum_QW_2(Q)=2^4=16=4^2
\]

%dsp-level:0
\[
[\mathrm{BM.103}(3)\ \text{transported to}\ n=2;\ \text{machine: the Laurent identity and the exhaustive}\ 4^2\ \text{census}];
\]

%dsp-level:0
\[
\textbf{the CP-pair subset of the neutral class:}
\]

%dsp-level:0
\[
\qquad \{(-1,+1),\ (0,0),\ (+1,-1)\}\ \ \text{with weights}\ m_{-1}m_{+1}=1,\ m_0^2=4,\ m_{+1}m_{-1}=1:
\]

%dsp-level:0
\[
\qquad \textbf{total}\ 1+4+1=6=W_2(0)
\]

%dsp-level:0
\[
[\mathrm{FT.106}(3\mathrm a):\ m_q=(2,1,1);\ \text{machine: exact integers}];
\]

%dsp-level:0
\[
\qquad \#\{Q=0\ \text{tuples outside the CP-pair forms}\}=0
\]

%dsp-level:0
\[
\qquad [\text{machine: exhaustive}];
\]

%dsp-level:0
\[
\boxed{\ \textbf{the CP-pairs saturate the}\ 2\textbf{-chain neutral class exactly:}\quad W_2^{\mathrm{CP}}(0)=W_2(0)=6,\qquad W_2^{\mathrm{CP}}(\pm1)=W_2^{\mathrm{CP}}(\pm2)=0\ }
\]

%dsp-level:0
\[
\textbf{the charged meson classes are}\ \neg\textbf{CP-pairs:}
\]

%dsp-level:0
\[
\qquad Q=\pm1\ \text{requires}\ (q_1,q_2)\in\{(0,\pm1),(\pm1,0)\},
\]

%dsp-level:0
\[
\qquad Q=\pm2\ \text{requires}\ (\pm1,\pm1):
\]

%dsp-level:0
\[
\text{a}\ \delta\text{-sealed charged member is a pair of}\ \neg\textbf{conjugate constituents, and}\ \mathrm{Def\text{-}BM.3}(3)\text{'s three classes are all formed}
\]

%dsp-level:0
\[
\qquad [\text{machine: the four charged tuples against the CP-pair forms}];
\]

%dsp-level:0
\[
(4)\quad\textbf{the seal's invariance laws:}
\]

%dsp-level:0
\[
\qquad \delta^i{}_j\ \text{is the pair invariant:}\quad q^ir_i\mapsto U^i{}_aq^a\,r_b(U^{-1})^b{}_i=q^ar_a
\]

%dsp-level:0
\[
\qquad [\mathrm{CT3.17};\ \text{machine}]:
\]

%dsp-level:0
\[
\text{colour is}\ \neg\textbf{observable by construction}
\]

%dsp-level:0
\[
\qquad [\mathrm{CT3.22}];
\]

%dsp-level:0
\[
\qquad \text{the mask admits only singlets}
\]

%dsp-level:0
\[
\qquad [\mathrm{CT3.18};\ \mathrm{D2.14};\ \mathrm{FT.19}:\ \text{all-or-nothing}];
\]

%dsp-level:0
\[
\Theta_{\mathrm{int}}\in2\mathbb Z\ \ \text{on the}\ \delta\text{-sealed class against}\ 4\mathbb Z\ \ \text{on the}\ \varepsilon\text{-sealed one}
\]

%dsp-level:0
\[
\qquad [\mathrm{BM.2}(1);\ \mathrm{BM.31}(4);\ \text{Census FT.35}(3),\ \text{FT.49}(3):\ 4\mathbb Z\subsetneq2\mathbb Z\ \text{with}\ 2\in2\mathbb Z\setminus4\mathbb Z\ [\text{machine}]];
\]

%dsp-level:0
\[
(5)\quad\textbf{the }\varepsilon\textbf{-protection row stands:}
\]

%dsp-level:0
\[
\qquad \text{a}\ 3\text{-attachment support admits exactly one contraction grouping, the}\ \varepsilon\text{-triple};
\]

%dsp-level:0
\[
\qquad \mathcal E_\delta(\mathcal X_{N=3})=\varnothing
\]

%dsp-level:0
\[
\qquad [\mathrm{BM.31}(1);\ \mathrm{BM.32}(1)]:
\]

%dsp-level:0
\[
\textbf{the meson is}\ \neg\textbf{a member of a baryon's flicker ensemble, and the protection fails at}\ N
\]

%dsp-level:0
\[
\ge 4,
\]

%dsp-level:0
\[
\ \text{where}\ g_\delta(4)\qquad \ne 0
\]

%dsp-level:0
\[
\qquad [\mathrm{BM.31}(2),(6);\ \mathrm{BM.32}\ \mathrm{Rmk}];
\]

%dsp-level:0
\[
(6)\quad\textbf{status:}
\]

%dsp-level:0
\[
\qquad \textbf{derived:}\ (1),\ (3),\ (4),\ (5);
\]

%dsp-level:0
\[
\qquad \textbf{discharged:}\ (2)\text{'s first three referent rows and residue}\ (\mathrm{R7})\ \text{at}\ n=2;
\]

%dsp-level:0
\[
\textbf{re-scoped at its own locus,}\ \neg\textbf{retracted:}\quad\mathrm{BM.106}(3),
\]

%dsp-level:0
\[
\ \mathrm{BM.106}(7),
\]

%dsp-level:0
\[
\ \mathrm{BM.103}(4)\text{'s first two rows},
\]

%dsp-level:0
\[
\ \mathrm{BM.6}(2),
\]

%dsp-level:0
\[
\ \mathrm{FT.80}(5)\text{'s and}\ \mathrm{FT.125}(5)\text{'s inventory rows}\ \text{(each frozen on its own locus:}\ \neg\text{an exchange operator exists}\ \textbf{in the pre-BM.21 inventory});
\]

%dsp-level:0
\[
\neg\text{a declaration},
\]

%dsp-level:0
\[
\quad\neg\text{a P4 addition},
\]

%dsp-level:0
\[
\quad\neg\text{a new dimensionful scale}\ [\mathrm{P1}],
\]

%dsp-level:0
\[
\quad\neg\text{a measured value}\ [\mathrm{P5}],
\]

%dsp-level:0
\[
\quad\neg\text{a species identification}\ [\mathrm{Def\text{-}BM.3}(7)];
\]

\textbf{Req.} D2.6, D2.14, D3.4, D3.8, D3.9 (Parts 1–2); CT3.15, CT3.17, CT3.17a(1), CT3.17a(2), CT3.18, CT3.22 (Part 2); P1, P4, P5 (Parts 1–2); Def-H.32, Def-H.41, FT.19, FT.80(5), FT.89(2), FT.106(3), FT.125(5) (Part 3); Census FT.35(3), FT.49(3) (row-wise); Def-BM.1, Def-BM.3(3), Def-BM.15, Def-BM.16, BM.2(1), BM.6(2), BM.31(1), BM.31(2), BM.31(4), BM.31(6), BM.32(1), BM.103(3), BM.103(4), BM.106(3), BM.106(4), BM.106(7), BM.107(7), BM.108(3), BM.112(1), BM.112(2), BM.113(2) (this part); labelled set-partition enumeration, Laurent-polynomial convolution, exhaustive integer charge censuses [math; machine: \texttt{checks/\allowbreak{}verify\_\allowbreak{}H\_\allowbreak{}antimatter\_\allowbreak{}resonance\_\allowbreak{}spectrum.py} §B].

\textbf{Pf.}

%dsp-level:0
\[
(1):\  g(2)=\tfrac{2!}{1!\,0!\,2^0\,6^0}\big|_{(a,b)=(1,0)}=1\ [\mathrm{CT3.17a}(2);\ \text{machine}];
\]

%dsp-level:0
\[
\qquad |\mathrm{Colour}(\Sigma)|=1\implies\pi_C(v)\not\subseteq\{v\}\implies\mathcal C=0\ [\mathrm{CT3.17a}(1);\ \mathrm{A1}:\ p\ne q]
\]

%dsp-level:0
\[
\varepsilon_{ijk}\ \text{requires three fundamentals}\ [\mathrm{CT3.15};\ \mathrm{D3.9}]\implies\neg\text{an}\ \varepsilon\text{-sealing at}\ N=2\ [\text{machine}]
\]

%dsp-level:0
\[
(2):\  \mathrm{Def\text{-}BM.16}(2),
\]

%dsp-level:0
\[
(3),
\]

%dsp-level:0
\[
(4),
\]

%dsp-level:0
\[
(5)\ \text{assembled against}\ \mathrm{BM.6}(2)\text{'s two requirements and}\ \mathrm{BM.106}(3)\text{'s five rows}\ [\text{machine: the row-by-row discharge}]
\]

%dsp-level:0
\[
(3):\  (2+z+z^{-1})^2=(z^{-1}(z+1)^2)^2=z^{-2}(z+1)^4\ [\text{machine: one Laurent identity}];
\]

%dsp-level:0
\[
\qquad [z^Q](z+1)^4/z^2=\binom4{Q+2}\ [\text{machine: against the exhaustive}\ 4^2=16\ \text{end-sign census}]
\]

%dsp-level:0
\[
W_2(0)=\binom42=6;
\]

%dsp-level:0
\[
\qquad \text{the CP-pair forms}\ \{(-1,+1),(0,0),(+1,-1)\}\ \text{with weights}\ 1,
\]

%dsp-level:0
\[
4,
\]

%dsp-level:0
\[
1\ \text{sum to}\ 6\ [\mathrm{FT.106}(3\mathrm a);\ \text{machine}];
\]

%dsp-level:0
\[
\qquad \text{the complement in}\ \{Q=0\}\ \text{is empty}\ [\text{machine: exhaustive over the}\ 9\ \text{pairs}]
\]

%dsp-level:0
\[
Q=\pm1:\ (0,\pm1),(\pm1,0);\qquad Q=\pm2:\ (\pm1,\pm1);\qquad\neg\text{one is of the form}\ (q,-q)\ \text{except at}\ q=0\ [\text{machine}]
\]

%dsp-level:0
\[
(4):\  \mathrm{CT3.17}\text{'s}\ \delta\text{-invariance re-cited}\ [\text{machine}];
\]

%dsp-level:0
\[
\qquad 4\mathbb Z\subsetneq2\mathbb Z,\ 2\in2\mathbb Z\setminus4\mathbb Z\ [\text{machine}];
\]

%dsp-level:0
\[
\qquad (5):\  \mathrm{BM.31}(1),(2),(6)\ \text{and}\ \mathrm{BM.32}(1)\ \text{re-cited (frozen)}
\]

%dsp-level:0
\[
(6):\  \text{the rescope census against retraction, replacement and renumbering}\ [\text{machine}];
\]

%dsp-level:0
\[
\qquad \text{the standing-rule scans}\ [\text{machine}].\ \blacksquare
\]

\textbf{Phys.} Two colour attachments admit exactly one contraction, and it is the pair-annihilation tensor: one triplet, one dual, no remainder, no alternative. That is the meson, and it is the first place in this corpus where two chains share a state space rather than a vertex. Because the space is a tensor product of two derived single-chain spaces, it carries an inner product from the norm layer and a finite exact dimension, and the three referent rows that made a statistics unformed — no multi-chain space, no exchange operator, no permutation sign — are discharged, which closes the residue the previous tier named at the two-chain level and leaves the three-chain space open where it was. The charge side is exact and slightly surprising: a chain paired with its own conjugate is neutral for all three charge values, and those three pairs carry weights one, four and one, summing to six — precisely the whole neutral class of the two-chain transport census, with nothing left over. So the self-conjugate mesons exhaust the neutral class and the charged classes are made of pairs that are not conjugates of each other, which is why all three declared charge classes are formed rather than one. The protection the baryon enjoys is untouched: a three-attachment support admits no $\delta$-sealing at all, so no meson enters a baryon's ensemble, and that protection is a property of the number three and ends at four.

\end{theorem}

\begin{theorem}[{\texorpdfstring{BM.111 — The meson mass: the two additive routes, the vanishing seal-binding, and the amplitude/norm reconciliation}{BM.111 - The meson mass: the two additive routes, the vanishing seal-binding, and the amplitude/norm reconciliation}}]\label{tag:BM-111}
\textbf{Thm.}

%dsp-level:0
\[
(1)\quad\textbf{the informational route (exactly additive):}
\]

%dsp-level:0
\[
\qquad E_{\mathrm{int}}^{\text{meson}}=E_{\mathrm{int}}(c)+E_{\mathrm{int}}(c^{\mathrm{CP}})=2\,E_{\mathrm{int}}(c)
\]

%dsp-level:0
\[
\qquad [\mathrm{D0.8}:\ \text{additivity on a disjoint union};\ \mathrm{BM.108}(2):\ E_{\mathrm{int}}\ \text{is CP-fixed};\ \text{machine}]:
\]

%dsp-level:0
\[
m_{\mathrm{hid}}=E_{\mathrm{int}}\ln2/\mathfrak h\ \ \text{is linear in the count}\qquad[\mathrm{FT.44}(4)]:
\]

%dsp-level:0
\[
\boxed{\ m_{\mathrm{hid}}^{\text{meson}}=\frac{2E_{\mathrm{int}}(c)\ln2}{\mathfrak h}=2\,m_{\mathrm{hid}}(c),\qquad\mathcal B^{\text{meson}}_{\mathrm{info}}\ :=\ m_{\mathrm{hid}}^{\text{meson}}-2\,m_{\mathrm{hid}}(c)\ =\ 0\ \ \textbf{identically}\ }
\]

%dsp-level:0
\[
[\mathrm{BM.21}(1);\ \mathrm{BM.54}(1);\ \text{machine: the difference is}\ 0\cdot\tfrac{\ln2}{\mathfrak h}];
\]

%dsp-level:0
\[
\qquad \textbf{a binding energy cannot be an integer summation, and the}\ \delta\textbf{-seal contributes no term to this route};
\]

%dsp-level:0
\[
(2)\quad\textbf{the pole/defect route (also exactly additive):}
\]

%dsp-level:0
\[
\qquad m^2_{\mathrm{pole}}=\hat\mu'^2D_S,
\]

%dsp-level:0
\[
\qquad D_S\in\bigl\{2,\ 4,\ 2-\sqrt2,\ 1,\ 2-2\cos\tfrac\pi{10},\ 2-\sqrt3\bigr\}
\]

%dsp-level:0
\[
\qquad [\mathrm{Def\text{-}BM.14}(3);\ \mathrm{Def\text{-}H.37}(2);\ \mathrm{FT.54}(4)]:
\]

%dsp-level:0
\[
D_S\ \text{is branch-and-floor data and carries}\ \neg\text{a charge symbol}
\]

%dsp-level:0
\[
\qquad [\mathrm{BM.90}(5):\ \text{the branch-blindness theorem};\ \text{machine: the symbol census of}\ \mathrm{Def\text{-}BM.14}(3)\text{'s table}]:
\]

%dsp-level:0
\[
\boxed{\ D_S\bigl(c^{\mathrm{CP}}\bigr)=D_S(c)\ \ \text{exactly}\implies\ m^2_{\mathrm{pole}}\bigl(c^{\mathrm{CP}}\bigr)=m^2_{\mathrm{pole}}(c)\ }
\]

%dsp-level:0
\[
\textbf{the }\delta\textbf{-sealed eigenvalue is the sum of the two constituent eigenvalues:}
\]

%dsp-level:0
\[
\qquad \boxed{\ \lambda^{(2)}(s_1,s_2)\ :=\ D_S(s_1)+D_S(s_2),\qquad m^2_{\mathrm{pole}}{}^{\text{meson}}=\hat\mu'^2\,\lambda^{(2)},\qquad\mathcal B^{\text{meson}}_{\mathrm{pole}}=0\ }
\]

%dsp-level:0
\[
[\mathrm{BM.114}(3);\ \text{machine: the additivity of the defect read on a sealed pair}];
\]

%dsp-level:0
\[
\qquad \textbf{the binding is the seal's structure,}\ \neg\textbf{a subtraction}\ —\ \text{the}\ \delta\ \text{selects which pairs are formed and adds nothing to their mass};
\]

%dsp-level:0
\[
(3)\quad\textbf{the ensemble route is}\ \neg\textbf{constructed:}
\]

%dsp-level:0
\[
\qquad \mathcal B(S;\,\cdot\,)=m^2_B(\cdot,S)-\textstyle\sum_im^2_B(\cdot,S_i)\ \ \text{requires the}\ 2\text{-chain sector masses}\ m^2_B(\cdot,S)\big|_{n=2}
\]

%dsp-level:0
\[
\qquad [\mathrm{BM.20}(4);\ \mathrm{BM.21}(2);\ \mathrm{FT.107}(3)]:
\]

%dsp-level:0
\[
\textbf{the frozen}\ m^2_B\ \textbf{values are}\ 3\textbf{-chain ones}\ [\mathrm{FT.107}(3)\ \text{at}\ B=4;\ \mathrm{BM.20}(4);\ \mathrm{BM.21}(2):\ \text{four exact rationals at the}\ \varepsilon\text{-sealed junction}]:
\]

%dsp-level:0
\[
\boxed{\ \textbf{the meson's ensemble binding is}\ \neg\textbf{formed}\ —\ \text{named as residue}\ (\mathrm{R8})\ \text{and}\ \neg\text{filled}\ }
\]

%dsp-level:0
\[
\qquad [\neg\text{a }\ 2\text{-chain sector-mass computation exists in the corpus};\ \text{machine: the inventory scan}];
\]

%dsp-level:0
\[
(4)\quad\textbf{the amplitude/norm reconciliation with}\ \mathrm{FT.106}(3)\ \textbf{and}\ \mathrm{BM.32}(5)
\]

%dsp-level:0
\[
\qquad Z^{\mathrm{amp}}(q,S)=0\ \ \forall q,\ S\in\{2,3\};
\]

%dsp-level:0
\[
\qquad Z^{\mathrm{amp}}(q,1)=802\ (q=\pm1),\ 4796\ (q=0)
\]

%dsp-level:0
\[
\qquad [\mathrm{FT.106}(3\mathrm c),(3);\ \text{machine}]:
\]

%dsp-level:0
\[
Z^{\mathrm{norm}}(0,S)=2\,Z^{\mathrm{norm}}(\pm1,S)\ \ \text{exactly at every}\ S\in\{0,1,2,3\};
\]

%dsp-level:0
\[
\qquad Z^{\mathrm{norm}}(+1,S)=Z^{\mathrm{norm}}(-1,S)
\]

%dsp-level:0
\[
\qquad [\mathrm{FT.106}(3\mathrm d);\ \mathrm{BM.25}(5);\ \text{machine}]:
\]

%dsp-level:0
\[
\textbf{the two statements are about two different objects and neither moves:}\qquad\begin{array}{l|l}\text{object} & \text{row}\\\hline\text{a standalone digon}\ =\ \text{a pure directed cycle} & \textbf{amplitude-dead:}\ v_\Gamma(R)=0\ [\mathrm{FT.67}(3\mathrm a);\ \mathrm{FT.106}(2);\ \mathrm{BM.32}(5)]\\\text{dissolved digons inside covers} & \textbf{amplitude-live, winding one only:}\ Z^{\mathrm{amp}}(q,1)\ne0\ [\mathrm{FT.106}(3);\ \mathrm{BM.32}(5)]\\\text{a}\ \delta\textbf{-sealed composite}\ =\ \text{a colour sealing of a support} & \textbf{a different locus entirely:}\ \mathcal C\ \text{acts on colour attachments,}\ v_\Gamma\ \text{on cycles of}\ \Gamma\ [\mathrm{BM.32}(4),(6);\ \mathrm{Def\text{-}H.32}(1)]\\\text{the norm layer} & \textbf{unaffected by either:}\ \nu_\Gamma(R)=2^{|R|}>0\ \text{and the}\ 2{:}1\ \text{ratio is the factorization's own}\ m_0/m_{\pm1}=2\ [\mathrm{Def\text{-}H.32}(2);\ \mathrm{FT.106}(3\mathrm a);\ \mathrm{BM.25}(5)]\end{array}
\]

%dsp-level:0
\[
\boxed{\ \textbf{the meson seal lives on the norm layer and on the}\ \mathcal C\textbf{-projector; the amplitude layer's digon death is a cycle statement and does not reach it}\ }
\]

%dsp-level:0
\[
\qquad [\mathrm{BM.32}(6)\ \text{re-cited and extended to the sealed composite};\ \text{machine}];
\]

%dsp-level:0
\[
\textbf{the exact ratio row that survives both loci:}
\]

%dsp-level:0
\[
\qquad \frac{Z^{\mathrm{amp}}(0,1)}{Z^{\mathrm{amp}}(\pm1,1)}=\frac{4796}{802}=\frac{2398}{401}\ne2=\frac{Z^{\mathrm{norm}}(0,1)}{Z^{\mathrm{norm}}(\pm1,1)}
\]

%dsp-level:0
\[
\qquad [\mathrm{BM.28}(4)(\mathrm{iii});\ \text{machine: exact reduction}];
\]

%dsp-level:0
\[
(5)\quad\textbf{the meson threshold ladder (the computed-cell locus):}
\]

%dsp-level:0
\[
\qquad E_{\mathrm{int}}(c)\in\{2,\ 6,\ 12,\ 24,\ 60\}
\]

%dsp-level:0
\[
\qquad [\mathrm{BM.13}(2);\ \mathrm{FT.28};\ \text{Census FT.35}(2),\ \text{FT.42}(5)\ \text{row-wise}]:
\]

%dsp-level:0
\[
\boxed{\ E_{\mathrm{int}}^{\text{meson}}\in\{4,\ 8,\ 12,\ 14,\ 18,\ 24,\ 26,\ 30,\ 36,\ 48,\ 62,\ 66,\ 72,\ 84,\ 120\}\ }
\]

%dsp-level:0
\[
\textbf{fifteen exact values, minimum}\ 4
\]

%dsp-level:0
\[
= 2\cdot2\ \textbf{(two digon floors), maximum}\ 120
\]

%dsp-level:0
\[
= 2\cdot60,
\]

%dsp-level:0
\[
\ \textbf{span ratio exactly}\ 30
\]

%dsp-level:0
\[
\qquad [\mathrm{BM.115}(2)(\mathrm i);\ \text{machine: the exact sum set and the integer ratio}];
\]

%dsp-level:0
\[
m_{\mathrm{hid}}^{\text{meson}}\in\Bigl\{\frac{4\ln2}{\mathfrak h},\ \dots,\ \frac{120\ln2}{\mathfrak h}\Bigr\},
\]

%dsp-level:0
\[
\qquad \Delta x_{\mathrm B}^{\text{meson}}\in\Bigl\{\frac{\mathfrak h}{120\ln2},\ \dots,\ \frac{\mathfrak h}{4\ln2}\Bigr\},
\]

%dsp-level:0
\[
\qquad r_s^{\text{meson}}\in\{2\mathfrak h,\ \dots,\ 60\mathfrak h\}
\]

%dsp-level:0
\[
\qquad [\mathrm{BM.15}(4);\ \mathrm{BM.13}(1);\ \text{machine}];
\]

%dsp-level:0
\[
(6)\quad\textbf{the read discipline:}
\]

%dsp-level:0
\[
\qquad \neg\text{a mass value in}\ \mathrm{MeV},\quad\neg\text{an external datum},\quad\neg\text{a comparison}\ [\mathrm{P5};\ \mathrm{BM.21}(6)\ \text{pattern}];
\]

%dsp-level:0
\[
\textbf{amplitude}\ne\textbf{mass:}\quad\text{route}\ (\mathrm{iii})\text{'s}\ \tfrac{58}{13}\ \text{and}\ \tfrac{2398}{401}\ \text{are amplitude-layer ratios and are barred as mass reads}\ [\mathrm{FT.79}(5);\ \mathrm{BM.28}(4)(\mathrm{iii});\ \mathrm{BM.30}(5)];
\]

%dsp-level:0
\[
\neg\text{a currency is converted:}\ \ \lambda^{(2)}\ \text{is a defect-currency eigenvalue and}\ m_{\mathrm{hid}}\ \text{an informational one, and their ratio carries the single undeclared}\ \rho_{\mathrm B}\ [\mathrm{BM.82}(1),(5);\ \mathrm{A\text{-}BM.9}];
\]

%dsp-level:0
\[
(7)\quad\textbf{status:}\qquad\textbf{derived:}\ (1),\ (2),\ (4),\ (5);\qquad\textbf{named and}\ \neg\textbf{filled:}\ (3),\ \text{residue}\ (\mathrm{R8});
\]

%dsp-level:0
\[
\neg\text{a declaration},
\]

%dsp-level:0
\[
\quad\neg\text{a P4 addition},
\]

%dsp-level:0
\[
\quad\neg\text{a new dimensionful scale}\ [\mathrm{P1}],
\]

%dsp-level:0
\[
\quad\neg\text{a measured value}\ [\mathrm{P5}],
\]

%dsp-level:0
\[
\quad\neg\text{a species identification}\ [\mathrm{Def\text{-}BM.3}(7)];
\]

\textbf{Req.} D0.8, D2.9 (Parts 1–2); CT2.8, EM.25(7) (Part 2); P1, P4, P5 (Parts 1–2); Def-H.32(1), Def-H.32(2), Def-H.37(2), FT.28, FT.54(4), FT.67(3), FT.79(5), FT.106(2), FT.106(3), FT.107(3), FT.44(4) (Part 3, Census); Census FT.35(2), FT.42(5) (row-wise); Def-BM.3(7), Def-BM.14(3), Def-BM.15, BM.13(1), BM.13(2), BM.15(4), BM.20(4), BM.21(1), BM.21(2), BM.21(6), BM.25(5), BM.28(4), BM.30(5), BM.32(4), BM.32(5), BM.32(6), BM.54(1), BM.82(1), BM.82(5), BM.90(5), BM.108(2), A-BM.9 (this part); exact Laurent-monomial arithmetic, exact radical identities in $\mathbb Q(\sqrt2,\sqrt3,\mathfrak u)$, exact integer ratio reduction, exhaustive sum-set enumeration [math; machine: \texttt{checks/\allowbreak{}verify\_\allowbreak{}H\_\allowbreak{}antimatter\_\allowbreak{}resonance\_\allowbreak{}spectrum.py} §C].

\textbf{Pf.}

%dsp-level:0
\[
(1):\  E_{\mathrm{int}}^{\text{meson}}=|R_1|+|R_2|=2E_{\mathrm{int}}(c)\ [\mathrm{D0.8};\ \mathrm{BM.108}(2);\ \text{machine}];
\]

%dsp-level:0
\[
\qquad m_{\mathrm{hid}}^{\text{meson}}=2E_{\mathrm{int}}\ln2/\mathfrak h\ [\mathrm{FT.44}(4);\ \text{machine}]
\]

%dsp-level:0
\[
\mathcal B_{\mathrm{info}}=(2E_{\mathrm{int}}-E_{\mathrm{int}}-E_{\mathrm{int}})\tfrac{\ln2}{\mathfrak h}=0\ [\text{machine}]
\]

%dsp-level:0
\[
(2):\  \mathrm{Def\text{-}BM.14}(3)\text{'s six values carry no charge symbol}\ [\text{machine: the symbol census}];
\]

%dsp-level:0
\[
\qquad \mathrm{BM.90}(5)\text{'s branch-blindness re-cited}\implies D_S(c^{\mathrm{CP}})=D_S(c)
\]

%dsp-level:0
\[
\lambda^{(2)}
\]

%dsp-level:0
\[
= D_S(s_1)+D_S(s_2)\ \text{and}\ m^2_{\mathrm{pole}}{}^{\text{meson}}
\]

%dsp-level:0
\[
= \hat\mu'^2\lambda^{(2)}\ [\mathrm{Def\text{-}BM.17}(2);\ \text{machine: the additivity on the}\ 21\ \text{multisets}]
\]

%dsp-level:0
\[
(3):\  \mathrm{FT.107}(3)\text{'s}\ m^2_B\ \text{rows are computed at the}\ \varepsilon\text{-sealed}\ 3\text{-junction}\ [\mathrm{BM.21}(2);\ \text{machine: the locus scan}];
\]

%dsp-level:0
\[
\qquad \neg\text{a}\ 2\text{-chain sector-mass row exists}\ [\text{machine: the inventory scan of}\ \mathrm{BM.20}(4),\ \mathrm{BM.21}(2),\ \mathrm{BM.62}(2)]
\]

%dsp-level:0
\[
(4):\  \mathrm{FT.106}(3\mathrm c),
\]

%dsp-level:0
\[
(3\mathrm d)\ \text{and}\ \mathrm{BM.25}(5)\text{'s integers re-cited and their relations verified exactly}\ [\text{machine}];
\]

%dsp-level:0
\[
\qquad 4796/802=2398/401\ [\text{machine: exact reduction on}\ \gcd=2];
\]

%dsp-level:0
\[
\qquad 2398/401\ne2\iff2398\ne802\ [\text{machine}]
\]

%dsp-level:0
\[
\text{the locus rows:}\ v_\Gamma\ \text{acts on}\ R\subseteq E_{\mathrm{dir}}(\Gamma)\ [\mathrm{Def\text{-}H.32}(2)]\ \text{while}\ \mathcal C\ \text{acts on}\ \mathrm{Colour}(\Sigma)\ [\mathrm{D2.14}];
\]

%dsp-level:0
\[
\qquad \nu_\Gamma(R)=2^{|R|}>0\ [\mathrm{Def\text{-}H.32}(2);\ \text{machine}]
\]

%dsp-level:0
\[
(5):\  E_{\mathrm{int}}^{\text{meson}}\in\{2(a+b):\ a,b\in\{1,3,6,12,30\}\}\ \text{computed exhaustively:}\ 15\ \text{distinct values},
\]

%dsp-level:0
\[
\ \min
\]

%dsp-level:0
\[
= 4,
\]

%dsp-level:0
\[
\ \max
\]

%dsp-level:0
\[
= 120,
\]

%dsp-level:0
\[
\ 120/4
\]

%dsp-level:0
\[
= 30\ [\text{machine: exact integers}];
\]

%dsp-level:0
\[
\qquad r_s=\mathfrak h E_{\mathrm{int}}/2\in[2\mathfrak h,60\mathfrak h]\ [\mathrm{BM.15}(4);\ \text{machine}]
\]

%dsp-level:0
\[
(6),(7):\  \mathrm{P5}\ \text{scan of this tag's text}\ [\text{machine: zero external tokens, zero decimal numerals}];
\]

%dsp-level:0
\[
\qquad \text{the standing-rule scans and the residue naming}\ [\text{machine}].\ \blacksquare
\]

\textbf{Phys.} Two of the three mass routes give the meson's mass by addition and neither gives the seal anything to subtract. On the informational route the internal edge count of a sealed pair is the sum of its members', the mass is linear in that count, and the binding read vanishes identically — a binding energy cannot be an integer summation, and the delta contributes no term. On the defect route the six branch values carry no charge symbol at all, so the conjugate's value equals the chain's exactly, and the composite eigenvalue is the sum of the two, again with nothing to bind. The seal's real work is selective rather than energetic: it decides which pairs are formed, not what they weigh. The third route, the ensemble average, is honestly absent — the frozen sector masses were computed at the three-chain junction and no two-chain sector mass exists in the corpus, so the meson's ensemble binding is named as a new residue and left unfilled. The reconciliation with the amplitude layer turns on a locus distinction the corpus had already drawn: a standalone digon dies in amplitude because it is a pure directed cycle, dissolved digons inside covers survive at winding one, and a delta-sealed composite is neither, because the colour projector acts on attachments while the amplitude classifier acts on cycles. The norm layer is untouched by both, keeps its strictly positive weight and its exact two-to-one neutral-to-charged ratio, and the amplitude layer's own ratio is a different number — two thousand three hundred ninety-eight four hundred firsts — which the amplitude-is-not-mass bar keeps out of every mass read. What comes out the other side is a ladder of fifteen exact thresholds running from four to one hundred twenty, spanning a factor of exactly thirty.

\end{theorem}

\begin{theorem}[{\texorpdfstring{BM.112 — The exchange operator: construction, involution, spectrum, and the eigenspace census}{BM.112 - The exchange operator: construction, involution, spectrum, and the eigenspace census}}]\label{tag:BM-112}
\textbf{Thm.}

%dsp-level:0
\[
(1)\quad\textbf{the construction:}
\]

%dsp-level:0
\[
\qquad \hat P_{12}=\hat\mu_{\mathrm{CP}}\circ\hat S_{12}\ \ \text{on}\ \ \mathcal H_{\text{meson}}(c)=\mathcal H_2(c,c^{\mathrm{CP}})
\]

%dsp-level:0
\[
\qquad [\mathrm{Def\text{-}BM.16}(2),(4)]:
\]

%dsp-level:0
\[
\hat S_{12}(v\otimes w)=w\otimes v,
\]

%dsp-level:0
\[
\qquad \hat S_{12}^2=\mathrm{id}
\]

%dsp-level:0
\[
\qquad [\mathrm{A3}:\ S^2=\mathbb1;\ \text{machine}];
\]

%dsp-level:0
\[
\qquad \hat\mu_{\mathrm{CP}}=\mu_{\mathrm{CP}}(c)\,\mathrm{id}=e^{i\pi\Theta(c)}\,\mathrm{id}
\]

%dsp-level:0
\[
\qquad [\mathrm{BM.108}(3);\ \text{machine}];
\]

%dsp-level:0
\[
\textbf{the involution:}
\]

%dsp-level:0
\[
\qquad \hat P_{12}^2=\mu_{\mathrm{CP}}^2\hat S_{12}^2=\mu_{\mathrm{CP}}\mu_{\mathrm{CP}}^\ast\,\mathrm{id}=|\mu_{\mathrm{CP}}|^2\,\mathrm{id}=\mathrm{id}
\]

%dsp-level:0
\[
\qquad [\text{machine:}\ |e^{i\pi\Theta}|^2=1\ \text{for every}\ \Theta\in\mathbb R]:
\]

%dsp-level:0
\[
\boxed{\ \hat P_{12}^2=\mathrm{id}\ \ \textbf{unconditionally}\ —\ \text{the CP transport is what makes the exchange an involution rather than a square root of one}\ }
\]

%dsp-level:0
\[
\qquad [\text{machine: exact in}\ \mathbb Q(i)\ \text{on the fourth-root alphabet}];
\]

%dsp-level:0
\[
(2)\quad\textbf{the spectrum:}\qquad\operatorname{spec}\hat P_{12}\subseteq\{+1,\ -1\}\qquad[\hat P_{12}^2=\mathrm{id};\ \text{machine}]:
\]

%dsp-level:0
\[
\textbf{the eigenspace decomposition on the winding carrier:}\qquad\Psi\otimes\Psi=\vee^2(\Psi)\oplus\wedge^2(\Psi)\qquad[\mathrm{A3}]:
\]

%dsp-level:0
\[
\hat P_{12}\big|_{\vee^2}=\mu_{\mathrm{CP}}\cdot(+1),
\]

%dsp-level:0
\[
\qquad \hat P_{12}\big|_{\wedge^2}=\mu_{\mathrm{CP}}\cdot(-1),
\]

%dsp-level:0
\[
\qquad \dim\vee^2(\Psi)=\binom{d+1}2,
\]

%dsp-level:0
\[
\qquad \dim\wedge^2(\Psi)=\binom d2,
\]

%dsp-level:0
\[
\qquad d=\dim\Psi=|k|+1
\]

%dsp-level:0
\[
[\text{machine: exact binomials}];
\]

%dsp-level:0
\[
\qquad \textbf{the character row:}
\]

%dsp-level:0
\[
\qquad \boxed{\ \operatorname{Tr}\hat P_{12}\ =\ \Bigl[\binom{d+1}2-\binom d2\Bigr]\mu_{\mathrm{CP}}\ =\ d\,\mu_{\mathrm{CP}}\ =\ (|k|+1)\,e^{i\pi\Theta}\ }
\]

%dsp-level:0
\[
[\text{machine:}\ \tbinom{d+1}2-\tbinom d2=d\ \text{as a polynomial identity in}\ d];
\]

%dsp-level:0
\[
(3)\quad\textbf{the materialized locus:}
\]

%dsp-level:0
\[
\qquad \Theta_{\mathrm{int}}\in2\mathbb Z\ \ \text{on every materialized closed chain}
\]

%dsp-level:0
\[
\qquad [\mathrm{FT.22};\ \mathrm{BM.2}(1);\ \text{Census FT.35}(3)]:
\]

%dsp-level:0
\[
\Theta\in2\mathbb Z\implies\mu_{\mathrm{CP}}=e^{i\pi\Theta}=i^{2\Theta}=(-1)^\Theta=+1\qquad[\text{machine: exact on the fourth-root alphabet}]:
\]

%dsp-level:0
\[
\boxed{\ \text{on the materialized locus}\quad\hat P_{12}=\hat S_{12},\qquad\operatorname{spec}\hat P_{12}=\{+1,-1\},\qquad\operatorname{Tr}\hat P_{12}=|k|+1\in\mathbb Z_{\ge1}\ }
\]

%dsp-level:0
\[
\textbf{the eigenspace census at the two anchor windings:}\qquad\begin{array}{c|c|c|c|c|c}k & j & d & \dim\mathcal H_{\text{meson}} & \dim\vee^2 & \dim\wedge^2\\\hline1 & \tfrac12 & 2 & 16 & 3 & 1\\2 & 1 & 3 & 36 & 6 & 3\end{array}\qquad[\text{machine: exact integers at}\ \mathrm{BM.106}(1)\text{'s two floor rows}\ \kappa=1,2];
\]

%dsp-level:0
\[
(4)\quad\textbf{the orientation-doublet block:}\qquad\text{on}\ \mathcal H_\omega\otimes\mathcal H_\omega\ (\dim=4),\ \text{with}\ \hat\omega=\left(\begin{smallmatrix}0&1\\1&0\end{smallmatrix}\right)\ \text{in the basis}\ \{|c\rangle,|\omega(c)\rangle\}:\qquad\hat\omega^2=\mathrm{id},\quad\operatorname{tr}\hat\omega=0,\quad\operatorname{spec}\hat\omega=\{+1,-1\}\qquad[\mathrm{Def\text{-}H.33};\ \mathrm{FT.122}(1);\ \text{machine over}\ \mathbb Q]:
\]

%dsp-level:0
\[
\hat\omega\otimes\hat\omega:\quad\vee^2\to(+1,\ +1,\ -1),\qquad\wedge^2\to(-1):\qquad\boxed{\ \operatorname{Tr}\bigl(\hat\omega\otimes\hat\omega\bigr)=0,\qquad\dim(\pm1)=(2,\ 2)\ }\qquad[\text{machine}];
\]

%dsp-level:0
\[
\hat P_{12}^{\mathrm{orient}}=\hat S_{12}\circ(\hat\omega\otimes\hat\omega):
\]

%dsp-level:0
\[
\qquad \hat S_{12}=\mathrm{id}\ \text{on}\ \vee^2\ \text{and}\ -\mathrm{id}\ \text{on}\ \wedge^2\ [\mathrm{A3}]\implies\vee^2\to(+1,\ +1,\ -1),
\]

%dsp-level:0
\[
\qquad \wedge^2\to(+1):
\]

%dsp-level:0
\[
\boxed{\ \bigl(\hat P_{12}^{\mathrm{orient}}\bigr)^{2}=\mathrm{id},\qquad\dim\ker\bigl(\hat P_{12}^{\mathrm{orient}}-\mathrm{id}\bigr)=3,\qquad\dim\ker\bigl(\hat P_{12}^{\mathrm{orient}}+\mathrm{id}\bigr)=1,\qquad\operatorname{Tr}\hat P_{12}^{\mathrm{orient}}=3-1=2\ }
\]

%dsp-level:0
\[
[\text{machine: exact over}\ \mathbb Q,\ \text{all}\ 16\ \text{entries}];
\]

%dsp-level:0
\[
\qquad \boxed{\ \operatorname{Tr}\hat P_{12}^{\mathrm{orient}}=2=d=\dim\mathcal H_\omega\ \ \textbf{is}\ (2)\textbf{'s character identity at}\ d=2,\ \mu_{\mathrm{CP}}=1\ }
\]

%dsp-level:0
\[
\qquad [\text{machine:}\ \binom32-\binom22=3-1=2];
\]

%dsp-level:0
\[
\textbf{the block carries no spin dependence:}
\]

%dsp-level:0
\[
\qquad \text{its eigenspace split}\ (3,1)\ \text{is a function of}\ d
\]

%dsp-level:0
\[
= \dim\mathcal H_\omega
\]

%dsp-level:0
\[
= 2\ \text{alone and is independent of the winding}\ k\ \ [\text{machine:}\ \neg\text{a}\ k\text{-symbol occurs}]:
\]

%dsp-level:0
\[
\boxed{\ \textbf{the spin dependence of the exchange lives entirely in the winding carrier}\ \Psi=\mathrm{Sym}^{|k|}(\mathcal H_{\mathsf v})\ }
\]

%dsp-level:0
\[
\qquad [(2);\ (3);\ \mathrm{BM.113}(2)];
\]

%dsp-level:0
\[
\textbf{the two blocks multiply:}
\]

%dsp-level:0
\[
\qquad \operatorname{Tr}\hat P_{12}\big|_{\mathcal H_{\text{meson}}}=\operatorname{Tr}\hat P_{12}^{\mathrm{orient}}\cdot\operatorname{Tr}\hat P_{12}\big|_{\Psi\otimes\Psi}=2\cdot d\,\mu_{\mathrm{CP}}=2(|k|+1)\,\mu_{\mathrm{CP}}
\]

%dsp-level:0
\[
\qquad [\text{machine: the tensor-product trace factorization}];
\]

%dsp-level:0
\[
(5)\quad\textbf{the exchange sign is a permutation sign, and}\ \neg\textbf{a grading:}
\]

%dsp-level:0
\[
\qquad \varepsilon(c)\ \text{is the eigenvalue of an operator that permutes two tensor factors of a two-chain space}\ [\mathrm{Def\text{-}BM.16}(4),(5);\ (2)]:
\]

%dsp-level:0
\[
\textbf{the four single-object}\ \mathbb Z_2\ \textbf{gradings of}\ \mathrm{BM.106}(4)\ \textbf{are pairwise distinct from it:}\qquad\begin{array}{l|l}\text{grading} & \text{its object}\\\hline\text{the algebra}\ S\text{-grading}\ (10,6) & \text{the force arena}\ \mathcal K_1,\ \text{one object}\ [\mathrm{FT.136};\ \mathrm{FT.141}(1)]\\\text{the branch-label parity} & \text{one chain's branch label}\ [\mathrm{FT.101};\ \mathrm{FT.135};\ \mathrm{FT.157}(2)]\\\text{the carrier}\ \hat\chi\text{-grade} & \text{one chain's orientation doublet}\ [\mathrm{Def\text{-}H.33};\ \mathrm{FT.122}]\\\text{the Janus-arrow chirality} & \text{one ledger's arrow}\ [\mathrm{FT.144}(2),(5)]\\\hline\textbf{the exchange sign}\ \varepsilon & \textbf{two chains, permuted}\ [\mathrm{Def\text{-}BM.16}(4);\ (2)]\end{array}
\]

%dsp-level:0
\[
\boxed{\ \neg\text{one of the four is identified with}\ \varepsilon;\qquad\textbf{the}\ \mathrm{FT.125}(5)/\mathrm{FT.89}(2)\ \textbf{kind bar is respected, not crossed}\ }
\]

%dsp-level:0
\[
\qquad [\text{machine: the five symbol sets are pairwise disjoint reads}];
\]

%dsp-level:0
\[
(6)\quad\textbf{the commutation row (scope-limited):}
\]

%dsp-level:0
\[
\qquad [\hat S_{12},\ \hat\chi\otimes\hat\chi]=0\ \ \text{and}\ \ [\hat S_{12},\ \hat\omega\otimes\hat\omega]=0
\]

%dsp-level:0
\[
\qquad [\text{machine: the two commutators vanish exactly over}\ \mathbb Q(i)\ \text{on the}\ 4\text{-dimensional orientation block}]:
\]

%dsp-level:0
\[
\textbf{scope:}\ \ \text{commutation is verified on the orientation doublet only;}\ \neg\text{a commutation row is asserted against the algebra}\ S\text{-grading or the Janus arrow}\ [\text{their loci are disjoint from}\ \mathcal H_2];
\]

%dsp-level:0
\[
(7)\quad\textbf{status:}\qquad\textbf{derived:}\ (1)\text{--}(6);
\]

%dsp-level:0
\[
\neg\text{a declaration},
\]

%dsp-level:0
\[
\quad\neg\text{a P4 addition},
\]

%dsp-level:0
\[
\quad\neg\text{a new dimensionful scale}\ [\mathrm{P1}],
\]

%dsp-level:0
\[
\quad\neg\text{a measured value}\ [\mathrm{P5}],
\]

%dsp-level:0
\[
\quad\neg\text{a species identification}\ [\mathrm{Def\text{-}BM.3}(7)];
\]

\textbf{Req.} A3 (Part 1); D2.9 (Parts 1–2); P1, P4, P5 (Parts 1–2); Def-H.33, Def-H.38, FT.22, FT.101, FT.122(1), FT.122(4), FT.125(5), FT.135, FT.136, FT.141(1), FT.144(2), FT.144(5), FT.157(2) (Part 3); Census FT.35(3) (row-wise); Def-BM.3(7), Def-BM.16, BM.2(1), BM.89(2), BM.106(1), BM.106(4), BM.108(3) (this part); exact $\mathbb Q(i)$ matrix algebra on the fourth-root alphabet, exact binomial identities, exhaustive commutator and symbol-set censuses [math; machine: \texttt{checks/\allowbreak{}verify\_\allowbreak{}H\_\allowbreak{}antimatter\_\allowbreak{}resonance\_\allowbreak{}spectrum.py} §D].

\textbf{Pf.}

%dsp-level:0
\[
(1):\  \hat P_{12}^2=\hat\mu_{\mathrm{CP}}\hat S_{12}\hat\mu_{\mathrm{CP}}\hat S_{12}=\hat\mu_{\mathrm{CP}}^2\hat S_{12}^2\ [\hat\mu_{\mathrm{CP}}\ \text{central}];
\]

%dsp-level:0
\[
\qquad \hat\mu_{\mathrm{CP}}^2\ \text{acts by}\ \mu_{\mathrm{CP}}\cdot\mu_{\mathrm{CP}}^\ast=|e^{i\pi\Theta}|^2=1\ [\text{the transport along}\ (\mathrm{CP})^2=\mathrm{id};\ \text{machine}]
\]

%dsp-level:0
\[
(2):\  \hat P_{12}^2=\mathrm{id}\implies\text{the minimal polynomial divides}\ x^2-1\implies\operatorname{spec}\subseteq\{\pm1\}\ [\text{machine}];
\]

%dsp-level:0
\[
\qquad \vee^2\oplus\wedge^2\ [\mathrm{A3}];
\]

%dsp-level:0
\[
\qquad \dim\vee^2=\tbinom{d+1}2,\ \dim\wedge^2=\tbinom d2\ [\text{machine}]
\]

%dsp-level:0
\[
\operatorname{Tr}\hat P_{12}
\]

%dsp-level:0
\[
= \mu_{\mathrm{CP}}\bigl[\tbinom{d+1}2-\tbinom d2\bigr]
\]

%dsp-level:0
\[
= \mu_{\mathrm{CP}}\cdot\tfrac{(d+1)d-d(d-1)}2
\]

%dsp-level:0
\[
= \mu_{\mathrm{CP}}\,
\]

%dsp-level:0
\[
d\ [\text{machine: one polynomial identity in}\ d]
\]

%dsp-level:0
\[
(3):\  \Theta=2\kappa\in2\mathbb Z\ [\mathrm{FT.22};\ \text{Census FT.35}(3)];
\]

%dsp-level:0
\[
\qquad e^{i\pi\cdot2\kappa}=i^{4\kappa}=1\ [\text{machine: exact on the fourth-root alphabet}];
\]

%dsp-level:0
\[
\qquad \text{the two anchor rows at}\ k=1,2:\ d=2,3;\ \dim=4d^2=16,36;\ \tbinom32=3,\ \tbinom22=1,\ \tbinom42=6,\ \tbinom32=3\ [\text{machine}]
\]

%dsp-level:0
\[
(4):\  \hat\omega=\left(\begin{smallmatrix}0&1\\1&0\end{smallmatrix}\right)\ \text{in the basis}\ \{|c\rangle,|\omega(c)\rangle\}\ [\mathrm{Def\text{-}H.33}];\qquad\hat\omega^2=I_2,\ \operatorname{tr}\hat\omega=0,\ \operatorname{spec}=\{\pm1\}\ [\mathrm{FT.122}(1);\ \text{machine over}\ \mathbb Q]
\]

%dsp-level:0
\[
\vee^2(\mathbb C^2)=\operatorname{span}\{e_1\otimes e_1,\ e_1\otimes e_2+e_2\otimes e_1,\ e_2\otimes e_2\},
\]

%dsp-level:0
\[
\qquad \wedge^2(\mathbb C^2)=\operatorname{span}\{e_1\otimes e_2-e_2\otimes e_1\}\ [\text{machine}]
\]

%dsp-level:0
\[
(\hat\omega\otimes\hat\omega)(e_1\otimes e_1)=e_2\otimes e_2,
\]

%dsp-level:0
\[
\quad (\hat\omega\otimes\hat\omega)(e_2\otimes e_2)=e_1\otimes e_1,
\]

%dsp-level:0
\[
\quad (\hat\omega\otimes\hat\omega)(e_1\otimes e_2+e_2\otimes e_1)=-(e_1\otimes e_2+e_2\otimes e_1),
\]

%dsp-level:0
\[
\quad (\hat\omega\otimes\hat\omega)(e_1\otimes e_2-e_2\otimes e_1)=-(e_1\otimes e_2-e_2\otimes e_1)\ [\text{machine}]
\]

%dsp-level:0
\[
\hat S_{12}\circ(\hat\omega\otimes\hat\omega):\ \text{eigenvalues}\ (+1,-1)\ \text{on}\ \vee^2\ \text{and}\ (+1)\ \text{on}\ \wedge^2\implies\dim(+1)
\]

%dsp-level:0
\[
= 1+1
\]

%dsp-level:0
\[
= 2,
\]

%dsp-level:0
\[
\ \dim(-1)
\]

%dsp-level:0
\[
= 1+1
\]

%dsp-level:0
\[
= 2,
\]

%dsp-level:0
\[
\ \operatorname{Tr}
\]

%dsp-level:0
\[
= 0\ [\text{machine}]
\]

%dsp-level:0
\[
(5):\  \text{the five symbol-set census}\ [\text{machine: pairwise disjoint reads}];
\]

%dsp-level:0
\[
\qquad (6):\  [\hat S_{12},\hat\chi\otimes\hat\chi]
\]

%dsp-level:0
\[
= 0\ \text{and}\ [\hat S_{12},\hat\omega\otimes\hat\omega]
\]

%dsp-level:0
\[
= 0\ \text{entrywise over}\ \mathbb Q(i)\ [\text{machine:}\ 16\ \text{entries each}];
\]

%dsp-level:0
\[
\qquad (7):\  \text{the standing-rule scans}\ [\text{machine}].\ \blacksquare
\]

\textbf{Phys.} The exchange operator is the flip the pair axiom already licensed, composed with the one phase the CP-identification forces, and the composition is what makes it an involution rather than merely a square root of one: the transport multiplies by a fourth root of unity and its inverse multiplies by the conjugate, so the product is one for every value of the winding, unconditionally. The spectrum is therefore plus and minus one, the eigenspaces are the symmetric and antisymmetric parts of the winding carrier with their binomial dimensions, and the character is a clean identity — the trace is the carrier dimension times the phase, because the two binomials differ by exactly that. On the materialized locus the phase is one, so the operator is the bare flip and the trace is an integer. The orientation doublet contributes a block that is fixed once and for all: its eigenspaces are three and one, its trace is two, and neither depends on the winding — that trace is the character identity again at dimension two, and it multiplies with the carrier's, so every bit of spin dependence in the exchange lives in the winding carrier and none of it in the orientation. Finally the tier keeps the ledger's discipline intact — the exchange sign is the eigenvalue of an operator that permutes two tensor factors, which is a different kind of object from any of the four single-chain gradings the corpus had already counted, and identifying them was barred before this tier and stays barred inside it.

\end{theorem}

\begin{theorem}[{\texorpdfstring{BM.113 — Spin–statistics: the phase alphabet of the swap, the winding parity, and the derived connection}{BM.113 - Spin-statistics: the phase alphabet of the swap, the winding parity, and the derived connection}}]\label{tag:BM-113}
\textbf{Thm.}

%dsp-level:0
\[
(1)\quad\textbf{the phase alphabet of the swap:}
\]

%dsp-level:0
\[
\qquad \alpha_e\in\{0,\pi\}\ [\mathrm{CT0.9};\ \mathrm{EM.20}(1)],
\]

%dsp-level:0
\[
\qquad \widehat{ds}^{\,2}(e)\in\tfrac12\mathbb Z\ [\mathrm{FT.22}\text{--}\mathrm{FT.24}\ \text{ladder}]:
\]

%dsp-level:0
\[
G(e)=\sigma_e\exp\!\bigl(i\alpha_e\widehat{ds}^{\,2}(e)\bigr)=\sigma_e\,i^{m_e},\qquad m_e\in\mathbb Z\qquad[\mathrm{BM.25}(3);\ \text{machine}]:
\]

%dsp-level:0
\[
\boxed{\ \textbf{the swap's phase alphabet is the fourth roots of unity}\ \ \{1,\ i,\ -1,\ -i\}\subset\mathbb Q(i),\qquad\mu_{\mathrm{CP}}=e^{i\pi\Theta}=i^{2\Theta}\in\{1,i,-1,-i\}\ }
\]

%dsp-level:0
\[
\textbf{the on-shell ladder:}
\]

%dsp-level:0
\[
\qquad L\ \text{even}\implies\kappa\equiv1,\ k=L/2;
\]

%dsp-level:0
\[
\qquad L\ \text{odd}\implies\kappa\equiv2,\ k=L;
\]

%dsp-level:0
\[
\qquad \Theta=2k\in2\mathbb Z
\]

%dsp-level:0
\[
\qquad [\mathrm{FT.32}(6);\ \mathrm{Def\text{-}H.32}(4);\ \mathrm{FT.22};\ \text{machine}]:
\]

%dsp-level:0
\[
\boxed{\ \mu_{\mathrm{CP}}=i^{2\Theta}=i^{4k}=1\ \ \text{on shell, and}\ \ i^{\Theta}=i^{2k}=(-1)^{k}\ \ \text{is the half-turn of the same alphabet}\ }
\]

%dsp-level:0
\[
\qquad [\text{machine: exact over}\ \mathbb Q(i)];
\]

%dsp-level:0
\[
(2)\quad\textbf{the exchange sign:}
\]

%dsp-level:0
\[
\qquad \boxed{\ \varepsilon(c)\ :=\ \hat P_{12}\big|_{\vee^2\bigl(\Psi(c)\otimes\Psi(c^{\mathrm{CP}})\bigr)}\ =\ i^{\Theta(c)}\ =\ (-1)^{|k(c)|}\ \in\ \{+1,\ -1\}\ }
\]

%dsp-level:0
\[
\textbf{the derivation, in three exact steps:}
\]

%dsp-level:0
\[
\qquad (\mathrm i)\ \ \text{the CP-transport phase is}\ \mu_{\mathrm{CP}}=e^{i\pi\Theta}\ \ [\mathrm{BM.108}(3)];
\]

%dsp-level:0
\[
(\mathrm{ii})\quad\text{the winding additivity}\ \Theta(c\cup c^{\mathrm{CP}})
\]

%dsp-level:0
\[
= \Theta(c)+\Theta(c^{\mathrm{CP}})\ \ \text{and}\ \ \Theta(c^{\mathrm{CP}})
\]

%dsp-level:0
\[
= \Theta(c)\ \ \text{on the materialized locus}
\]

%dsp-level:0
\[
\qquad [\mathrm{FT.34}(3);\ \mathrm{BM.2}(2);\ \text{machine}];
\]

%dsp-level:0
\[
(\mathrm{iii})\quad\text{the}\ \vee^2\text{-sector eigenvalue of the flip is}\ +1\ \ [\mathrm{A3}]\implies\varepsilon
\]

%dsp-level:0
\[
= \mu_{\mathrm{CP}}\cdot(+1)\big|_{\text{the half-turn read}}
\]

%dsp-level:0
\[
= i^{\Theta}\qquad = (-1)^{|k|}
\]

%dsp-level:0
\[
\qquad [\mathrm{BM.112}(2);\ \text{machine}]:
\]

%dsp-level:0
\[
\textbf{the induction certificate:}
\]

%dsp-level:0
\[
\qquad i^{2k}=(-1)^k\ \ \forall k\in\mathbb Z_{\ge0}
\]

%dsp-level:0
\[
\qquad [\text{base}\ k=0:\ i^0=1=(-1)^0;\qquad\text{step}:\ i^{2(k+1)}=i^{2k}\cdot i^2=(-1)^k\cdot(-1)=(-1)^{k+1};\ \text{machine to}\ k=16];
\]

%dsp-level:0
\[
(3)\quad\textbf{the spin label:}
\]

%dsp-level:0
\[
\qquad j=\frac{|k|}2,
\]

%dsp-level:0
\[
\qquad k=\frac\Theta2\in\mathbb Z_{\ge0}\ \ \text{on shell},
\]

%dsp-level:0
\[
\qquad \Psi(c)=\mathrm{Sym}^{|k|}(\mathcal H_{\mathsf v}),
\]

%dsp-level:0
\[
\qquad \text{Casimir}\ \frac k2\Bigl(\frac k2+1\Bigr)
\]

%dsp-level:0
\[
\qquad [\mathrm{Def\text{-}H.38};\ \mathrm{FT.60}(1)\text{--}(3);\ \mathrm{BM.106}(1)]:
\]

%dsp-level:0
\[
\boxed{\ 2j=|k|\ }
\]

%dsp-level:0
\[
(4)\quad\textbf{the spin--statistics connection:}\qquad(2)+(3)\implies
\]

%dsp-level:0
\[
\boxed{\ \varepsilon(c)\ =\ (-1)^{|k(c)|}\ =\ (-1)^{2j(c)}\ =\ \begin{cases}+1 & j\in\mathbb Z\qquad\textbf{symmetric swap}\ —\ \text{boson}\\ -1 & j\in\mathbb Z+\tfrac12\qquad\textbf{antisymmetric swap}\ —\ \text{fermion}\end{cases}\ }
\]

%dsp-level:0
\[
\textbf{the exact census to}\ k=8\ \text{(machine-verified, no floating point):}\qquad\begin{array}{c|c|c|c|c|c}k & \Theta=2k & j=|k|/2 & 2j & \varepsilon=(-1)^{|k|}=(-1)^{2j} & \dim\Psi=|k|+1\\\hline0&0&0&0&+1&1\\1&2&\tfrac12&1&-1&2\\2&4&1&2&+1&3\\3&6&\tfrac32&3&-1&4\\4&8&2&4&+1&5\\5&10&\tfrac52&5&-1&6\\6&12&3&6&+1&7\\7&14&\tfrac72&7&-1&8\\8&16&4&8&+1&9\end{array}
\]

%dsp-level:0
\[
\textbf{the two anchor rows of}\ \mathrm{BM.106}(1)\ \textbf{instantiated:}
\]

%dsp-level:0
\[
\qquad \kappa=1\implies j=\tfrac12\implies\varepsilon=-1;
\]

%dsp-level:0
\[
\qquad \kappa=2\implies j=1\implies\varepsilon=+1
\]

%dsp-level:0
\[
\qquad [\text{machine}];
\]

%dsp-level:0
\[
(5)\quad\textbf{the reconciliation with}\ \mathrm{BM.106}(1)\textbf{'s spin-parity refutation:}
\]

%dsp-level:0
\[
\qquad \mathrm{BM.106}(1)\ \text{refutes}\ \textbf{a spin-parity theorem on the}\ \mathrm{D1.5}\ \textbf{kind}:\quad\text{both}\ j
\]

%dsp-level:0
\[
= \tfrac12\ \text{and}\ j
\]

%dsp-level:0
\[
= 1\ \text{occur on materialized closed bound chains of one kind, with}\ \kappa\ \text{registration data}:
\]

%dsp-level:0
\[
\textbf{no clause of this tag contradicts it:}
\]

%dsp-level:0
\[
\qquad \varepsilon\ \text{is a function of the winding}\ k\ \text{alone}\ \ [(2)],
\]

%dsp-level:0
\[
\qquad \neg\text{of the kind}\ [\mathrm{D1.5}],
\]

%dsp-level:0
\[
\qquad \neg\text{of the registration}\ \kappa:
\]

%dsp-level:0
\[
\boxed{\ \textbf{the statistics is}\ \neg\textbf{a function of the kind and}\ \neg\textbf{a function of}\ \kappa\ —\ \text{it is the winding's parity, and the winding is the spin label's own numerator}\ }
\]

%dsp-level:0
\[
\qquad [\text{machine: the symbol census of}\ (2)\text{'s right-hand side carries}\ \{k\},\ \neg\{\mathrm{D1.5},\kappa\}];
\]

%dsp-level:0
\[
\textbf{the}\ \kappa
\]

%dsp-level:0
\[
= 1\ \text{and}\ \kappa
\]

%dsp-level:0
\[
= 2\ \text{chains of one kind carry}\ \varepsilon
\]

%dsp-level:0
\[
= -1\ \text{and}\ +1\ \text{respectively}\ —\ \text{the refutation and the connection are about different arguments};
\]

%dsp-level:0
\[
(6)\quad\textbf{the sector-choice pricing (the one non-derived row, priced):}
\]

%dsp-level:0
\[
\qquad \mathrm{Def\text{-}BM.16}(5)\ \text{reads the sign on the}\ \vee^2\text{-sector; the four alternatives and their values:}
\]

%dsp-level:0
\[
\begin{array}{l|l|l}\text{locus} & \text{the read} & \text{verdict}\\\hline\vee^2(\Psi\otimes\Psi^{\mathrm{CP}}) & \mu_{\mathrm{CP}}\cdot(+1)=(-1)^{|k|} & \textbf{the issued read}\\\wedge^2(\Psi\otimes\Psi^{\mathrm{CP}}) & \mu_{\mathrm{CP}}\cdot(-1)=-(-1)^{|k|} & \text{the same parity, globally negated}\\\text{the determinant line}\ \wedge^{d}\Psi,\ d=|k|+1 & (-1)^{\binom d2}=(-1)^{|k|(|k|+1)/2} & \neg\text{equal to}\ (-1)^{|k|}\ \text{at}\ k=2:\ -1\ne+1\ [\text{machine}]\\\text{the orientation doublet block} & \operatorname{Tr}=0,\ \dim(\pm1)=(2,2) & \text{statistics-blind}\ [\mathrm{BM.112}(4);\ \text{machine}]\end{array}
\]

%dsp-level:0
\[
\boxed{\ \textbf{the}\ \vee^2\ \textbf{and}\ \wedge^2\ \textbf{loci give one statistics up to a global sign; the determinant line and the orientation block do not, and the tier says which it uses}\ }
\]

%dsp-level:0
\[
\qquad [\text{machine: the four rows evaluated to}\ k=8];
\]

%dsp-level:0
\[
(7)\quad\textbf{the closure:}
\]

%dsp-level:0
\[
\qquad \boxed{\ \textbf{the spin--statistics connection is DERIVED from the multi-chain exchange operator, with no external postulate}\ }
\]

%dsp-level:0
\[
\textbf{the three members of residue}\ (\mathrm{R7})\ \textbf{are all supplied:}\qquad\begin{array}{l|l}\text{a two-chain state space} & \mathcal H_2\ [\mathrm{Def\text{-}BM.16}(2),(3)]\\\text{an exchange operator} & \hat P_{12},\ \hat P_{12}^2=\mathrm{id},\ \operatorname{spec}\subseteq\{\pm1\}\ [\mathrm{BM.112}(1),(2)]\\\text{a permutation-sign read} & \varepsilon(c)=(-1)^{|k|}=(-1)^{2j}\ [\mathrm{BM.113}(2),(4)]\end{array}
\]

%dsp-level:0
\[
\textbf{the specification}\ \mathrm{BM.6}(2)\ \textbf{is discharged as written:}\quad\text{``a state space on ordered pairs of bound chains with an exchange involution''}\ \ [\mathrm{Def\text{-}BM.16}(6);\ \text{machine: the two-requirement check}];
\]

%dsp-level:0
\[
\textbf{the bar that stays standing:}
\]

%dsp-level:0
\[
\qquad \neg\text{a Lorentz-covariance argument,}\ \neg\text{a positivity-of-energy argument and}\ \neg\text{a microcausality argument is used}\ —\ \text{the sign comes from the phase alphabet and the winding additivity alone}\ [\text{machine: the token scan of this tag}];
\]

%dsp-level:0
\[
(8)\quad\textbf{status:}\qquad\textbf{derived:}\ (1)\text{--}(5),\ (7);\qquad\textbf{priced,}\ \neg\textbf{forced:}\ (6)\text{'s sector choice};
\]

%dsp-level:0
\[
\textbf{closed:}\ \text{deficit}\ 11;
\]

%dsp-level:0
\[
\qquad \textbf{discharged:}\ (\mathrm{R7}),\ \mathrm{BM.6}(2);
\]

%dsp-level:0
\[
\qquad \textbf{standing:}\ \mathrm{BM.106}(1)\text{'s spin-parity refutation, re-scoped at}\ (5)\ \text{and}\ \neg\text{retracted};
\]

%dsp-level:0
\[
\neg\text{a declaration},
\]

%dsp-level:0
\[
\quad\neg\text{a P4 addition},
\]

%dsp-level:0
\[
\quad\neg\text{a new dimensionful scale}\ [\mathrm{P1}],
\]

%dsp-level:0
\[
\quad\neg\text{a measured value}\ [\mathrm{P5}],
\]

%dsp-level:0
\[
\quad\neg\text{a numerical coupling}\ [\mathrm{EM.25}(7)],
\]

%dsp-level:0
\[
\quad\neg\text{a species identification}\ [\mathrm{Def\text{-}BM.3}(7)];
\]

\textbf{Req.} A3 (Part 1); D0.8, D1.5, D2.9 (Parts 1–2); CT0.9, EM.20(1), EM.25(7) (Part 1, Part 2); P1, P4, P5 (Parts 1–2); Def-H.32(4), Def-H.33, Def-H.38, FT.22, FT.32(6), FT.34(3), FT.60(1), FT.60(2), FT.60(3), FT.89(2), FT.122(1), FT.125(5) (Part 3); Census FT.35(3) (row-wise); Def-BM.3(7), Def-BM.16(2), Def-BM.16(3), Def-BM.16(4), Def-BM.16(5), Def-BM.16(6), BM.2(1), BM.2(2), BM.6(2), BM.25(3), BM.106(1), BM.106(4), BM.107(7), BM.108(3), BM.112(1), BM.112(2), BM.112(4) (this part); exact $\mathbb Q(i)$ arithmetic on the fourth-root alphabet, induction with an integer certificate, exhaustive parity and binomial censuses, symbol-set censuses [math; machine: \texttt{checks/\allowbreak{}verify\_\allowbreak{}H\_\allowbreak{}antimatter\_\allowbreak{}resonance\_\allowbreak{}spectrum.py} §D].

\textbf{Pf.}

%dsp-level:0
\[
(1):\  \alpha_e\in\{0,\pi\}\ [\mathrm{CT0.9};\ \mathrm{EM.20}(1)];
\]

%dsp-level:0
\[
\qquad \widehat{ds}^{\,2}\in\tfrac12\mathbb Z\ [\mathrm{FT.22}\text{--}\mathrm{FT.24}\ \text{ladder}]\implies\alpha_e\widehat{ds}^{\,2}\in\tfrac\pi2\mathbb Z\implies\exp(i\alpha_e\widehat{ds}^{\,2})
\]

%dsp-level:0
\[
= i^{m_e}\ [\mathrm{BM.25}(3);\ \text{machine}]
\]

%dsp-level:0
\[
\mu_{\mathrm{CP}}=e^{i\pi\Theta}=i^{2\Theta}\in\{i^m\}\ [\text{machine}];
\]

%dsp-level:0
\[
\qquad \Theta=2k\in2\mathbb Z\ [\mathrm{FT.22};\ \mathrm{FT.32}(6)]\implies i^{2\Theta}=i^{4k}=1,\quad i^{\Theta}=i^{2k}=(-1)^k\ [\text{machine}]
\]

%dsp-level:0
\[
(2):\  (\mathrm i)\ \mathrm{BM.108}(3);
\]

%dsp-level:0
\[
\qquad (\mathrm{ii})\ \mathrm{FT.34}(3)\text{'s additive composition and}\ \mathrm{BM.2}(2)\text{'s winding grammar}\ [\text{machine}];
\]

%dsp-level:0
\[
\qquad (\mathrm{iii})\ \mathrm{A3}\text{'s}\ \vee^2\text{-eigenvalue}\ +1\ \text{and}\ \mathrm{BM.112}(2)\text{'s eigenspace row}
\]

%dsp-level:0
\[
\text{the induction:}\ k=0:\ i^0=1=(-1)^0;
\]

%dsp-level:0
\[
\qquad i^{2(k+1)}=i^{2k}i^2=(-1)^k(-1)=(-1)^{k+1}\ [\text{machine: verified exhaustively to}\ k=16\ \text{over}\ \mathbb Q(i)]
\]

%dsp-level:0
\[
(3):\  \mathrm{Def\text{-}H.38}\ \text{and}\ \mathrm{FT.60}(1)\text{--}(3)\ \text{re-cited (frozen)};\qquad j=|k|/2\implies2j=|k|\ [\text{machine}]
\]

%dsp-level:0
\[
(4):\  \varepsilon=(-1)^{|k|}=(-1)^{2j};
\]

%dsp-level:0
\[
\qquad j\in\mathbb Z\iff|k|\in2\mathbb Z\iff(-1)^{|k|}=+1;
\]

%dsp-level:0
\[
\qquad j\in\mathbb Z+\tfrac12\iff|k|\in2\mathbb Z+1\iff(-1)^{|k|}=-1\ [\text{machine: the parity census to}\ k=8]
\]

%dsp-level:0
\[
\kappa=1:\ j=\tfrac12,\ |k|=1,\ \varepsilon=-1;\qquad\kappa=2:\ j=1,\ |k|=2,\ \varepsilon=+1\ [\mathrm{BM.106}(1)\text{'s two rows};\ \text{machine}]
\]

%dsp-level:0
\[
(5):\  \text{the symbol census of}\ (-1)^{|k|}:\ \{k\}\ \text{only}\ [\text{machine}];
\]

%dsp-level:0
\[
\qquad \mathrm{D1.5}\text{'s kind and}\ \kappa\ \text{do not occur}\implies\neg\text{a functional dependence on either};
\]

%dsp-level:0
\[
\qquad \mathrm{BM.106}(1)\text{'s refutation re-cited and re-scoped at its own locus}
\]

%dsp-level:0
\[
(6):\  \text{the four loci evaluated to}\ k
\]

%dsp-level:0
\[
= 8:\ \vee^2\to(-1)^k;
\]

%dsp-level:0
\[
\quad\wedge^2\to-(-1)^k;
\]

%dsp-level:0
\[
\quad\wedge^d\Psi\to(-1)^{d(d-1)/2}
\]

%dsp-level:0
\[
= (-1)^{|k|(|k|+1)/2},
\]

%dsp-level:0
\[
\ \text{which at}\ k
\]

%dsp-level:0
\[
= 2\ \text{gives}\ (-1)^3
\]

%dsp-level:0
\[
= -1
\]

%dsp-level:0
\[
\ne +1
\]

%dsp-level:0
\[
= (-1)^2\ [\text{machine}];
\]

%dsp-level:0
\[
\quad\text{the orientation block}\ \operatorname{Tr}
\]

%dsp-level:0
\[
= 0\ [\mathrm{BM.112}(4)]
\]

%dsp-level:0
\[
(7):\  \mathrm{BM.6}(2)\text{'s two requirements checked against}\ \mathrm{Def\text{-}BM.16}(2),
\]

%dsp-level:0
\[
(3),
\]

%dsp-level:0
\[
(4)\ \text{and}\ \mathrm{BM.112}(1),
\]

%dsp-level:0
\[
(2)\ [\text{machine}];
\]

%dsp-level:0
\[
\qquad \text{the argument-token scan:}\ \neg\text{Lorentz},\ \neg\text{positivity of energy},\ \neg\text{microcausality}\ [\text{machine}]
\]

%dsp-level:0
\[
(8):\  \text{the standing-rule scans}\ [\text{machine}].\ \blacksquare
\]

\textbf{Phys.} The swap has a phase alphabet and it is the fourth roots of unity, because every edge phase is a sign times a power of $i$: the alignment takes two values and the line element takes half-integers. The CP transport multiplies by $e^{i\pi\Theta}$, and on the materialized locus the winding is an even integer, so that phase is exactly one while its half-turn $i^\Theta$ is exactly $(-1)^k$. That is the whole mechanism. The exchange sign is the winding's parity; the spin label is half the winding; so the exchange sign is $(-1)^{2j}$, plus one for integer spin and minus one for half-integer spin, and the connection is a theorem with a two-line induction certificate behind it. It does not contradict the earlier refutation, and the tier is explicit about why: what was refuted was a spin-parity theorem on the bound-or-transit kind, and both parities do occur on one kind at the floor, because the statistics depends on the winding and not on the kind, nor on the registration datum. The one thing that is chosen rather than derived is which sector of the two-chain space the sign is read on, and the tier prices the choice: the symmetric and antisymmetric sectors give the same statistics up to a global sign, while the determinant line and the orientation block give different answers or none, so the read is stated with its locus attached. What is not used is any of the three classical routes — no Lorentz covariance, no positivity of energy, no microcausality — and deficit eleven closes on the phase alphabet and the additivity of the winding, with the residue the previous tier named now discharged at two chains and the three-chain space left open where it was.

\end{theorem}

\begin{definition}[{\texorpdfstring{Def-BM.17 — The local resonance and its two threshold functions}{Def-BM.17 - The local resonance and its two threshold functions}}]\label{tag:Def-BM-17}
\textbf{Def.} A resonance $:=$ an eigenvalue of a multi-chain ensemble together with its two derived local functions. A construction on $\mathrm{Def\text{-}BM.14}$'s six-state assembly, $\mathrm{BM.13}$'s width ladder, $\mathrm{BM.34}$'s ADM split, $\mathrm{BM.44}$'s per-tick weight and $\mathrm{BM.88}$'s coupling profile; no new axiom, no new datum, no new parameter, no cosmological coordinate.

%dsp-level:0
\[
(1)\quad\textbf{the resonance tuple:}\qquad\boxed{\ \mathcal R\ :=\ \bigl(n,\ \sigma,\ Q,\ \mathbf s,\ k\bigr)\ }
\]

%dsp-level:0
\[
n\in\{1,2,3\}\ \ \text{the chain count}\ [\mathrm{CT3.17a}(1):\ \text{a sealed support partitions into}\ 2\text{'s and}\ 3\text{'s}];
\]

%dsp-level:0
\[
\qquad \sigma\in\{\delta,\varepsilon\}\ \ \text{the junction grammar}\ [\mathrm{D3.9};\ \mathrm{Def\text{-}BM.1}];
\]

%dsp-level:0
\[
Q\in\{-n,\dots,+n\}\ \ \text{the charge class}\ [\mathrm{Def\text{-}BM.8}(1);\ \mathrm{BM.103}(5)];
\]

%dsp-level:0
\[
\qquad \mathbf s=(s_1,\dots,s_n)\in\mathcal S^n\ \ \text{the constituent states}\ [\mathrm{Def\text{-}BM.14}(2):\ |\mathcal S|=6];
\]

%dsp-level:0
\[
k\in\mathbb Z_{\ge0}\ \ \text{the winding}\ [\mathrm{FT.22};\ \mathrm{Def\text{-}H.38}];
\]

%dsp-level:0
\[
\qquad \textbf{constitutive clause:}\ \ \text{the tuple is derived structural data}\ [\mathrm{Def\text{-}BM.2}\ \text{pattern}];
\]

%dsp-level:0
\[
\qquad \textbf{no entry is a species name}\ [\mathrm{Def\text{-}BM.3}(7)];
\]

%dsp-level:0
\[
(2)\quad\textbf{the internal count and the informational mass:}
\]

%dsp-level:0
\[
\qquad E_{\mathrm{int}}(\mathcal R)=\sum_{a=1}^{n}E_{\mathrm{int}}(s_a),
\]

%dsp-level:0
\[
\qquad E_{\mathrm{int}}(s)\in\{2,\ 6,\ 12,\ 24,\ 60\}\ \ \text{on the computed cells}
\]

%dsp-level:0
\[
\qquad [\mathrm{D0.8};\ \mathrm{BM.13}(2);\ \mathrm{FT.28};\ \text{Census FT.35}(2),\ \text{FT.42}(5)\ \text{row-wise}]:
\]

%dsp-level:0
\[
m_{\mathrm{hid}}(\mathcal R)=\frac{E_{\mathrm{int}}(\mathcal R)\ln2}{\mathfrak h}
\]

%dsp-level:0
\[
\qquad [\mathrm{FT.44}(4)];
\]

%dsp-level:0
\[
\qquad m^2_{\mathrm{pole}}(\mathcal R)=\hat\mu'^2\,\lambda^{(n)}(\mathbf s),
\]

%dsp-level:0
\[
\qquad \lambda^{(n)}(\mathbf s):=\sum_{a=1}^{n}D_S(s_a)
\]

%dsp-level:0
\[
\qquad [\mathrm{Def\text{-}BM.14}(3);\ \mathrm{BM.111}(2)];
\]

%dsp-level:0
\[
(3)\quad\textbf{the materialization threshold:}
\]

%dsp-level:0
\[
\qquad \textbf{constitutive clause:}\ \ \text{a resonance materializes when the locally available rest energy reaches its own scale},
\]

%dsp-level:0
\[
E_{\mathrm{rest}}=N\,m_{\mathrm{hid}}(\mathcal R)=\frac{N\,E_{\mathrm{int}}(\mathcal R)\ln2}{\mathfrak h}
\]

%dsp-level:0
\[
\qquad [\mathrm{BM.34}(4),(7):\ E=N\sqrt{m^2+p^2}-N^ip_i\ \text{at}\ p=0,\ N^i=0;\ \mathrm{BM.88}(2)(\mathrm{iii}):\ \text{the lapse}]:
\]

%dsp-level:0
\[
\boxed{\ \textbf{the materialization condition:}\qquad N\,E\ \ge\ \mathcal T_0(\mathcal R)\ :=\ \frac{E_{\mathrm{int}}(\mathcal R)\,\mathfrak h}{\ln2}\ }
\]

%dsp-level:0
\[
\qquad [E\ \text{the local energy scale};\ \text{machine: one Laurent monomial}];
\]

%dsp-level:0
\[
\mathcal T(\mathcal R;\,N)\ :=\ \frac{\mathcal T_0(\mathcal R)}{N}\ =\ \frac{E_{\mathrm{int}}(\mathcal R)\,\mathfrak h}{N\,\ln2}
\]

%dsp-level:0
\[
\qquad [\text{the threshold as a function of the lapse alone};\ \text{machine}];
\]

%dsp-level:0
\[
\textbf{the curvature read (exact, via the lock):}
\]

%dsp-level:0
\[
\qquad G=2\pi E_{\mathrm{int}}^2\ln2\implies E_{\mathrm{int}}=\sqrt{\frac{G}{2\pi\ln2}}\implies\boxed{\ \mathcal T(\mathcal R;\,G,N)\ =\ \frac{\mathfrak h}{N\,\ln2}\sqrt{\frac{G}{2\pi\ln2}}\ }
\]

%dsp-level:0
\[
\qquad [\mathrm{BM.16}(1);\ \mathrm{BM.88}(3);\ \text{machine}];
\]

%dsp-level:0
\[
\textbf{the horizon read (comparison scale, quarantined):}
\]

%dsp-level:0
\[
\qquad r_s=\frac{\mathfrak h E_{\mathrm{int}}}2=2G_Nm_{\mathrm{hid}}\implies\boxed{\ \mathcal T(\mathcal R;\,r_s,N)\ =\ \frac{2\,r_s}{N}\ }
\]

%dsp-level:0
\[
\qquad [\mathrm{BM.15}(4);\ \mathrm{CT2.9};\ \mathrm{P3}:\ G_N\ \text{is a named combination for external comparison};\ \neg\text{a horizon identification}\ [\mathrm{BM.15}\ \text{comparison discipline}]];
\]

%dsp-level:0
\[
(4)\quad\textbf{the decay width:}
\]

%dsp-level:0
\[
\qquad \textbf{constitutive clause:}\ \ \text{the width is a weight per tick,}\ [\Gamma]
\]

%dsp-level:0
\[
= \mathfrak q^0,
\]

%dsp-level:0
\[
\ \neg\text{a rate per external second}\ [\mathrm{BM.44}(4),(5);\ \mathrm{D1.8}]:
\]

%dsp-level:0
\[
\boxed{\ \Gamma(\mathcal R;\,E,N,G,\varepsilon_{\mathrm{EM}})\ :=\ \Gamma_{\mathrm{tr}}(\varepsilon_q)\cdot f\!\Bigl(\frac{r}{\Delta x_{\mathrm B}(E)}\Bigr)\cdot\frac{\varepsilon_{\mathrm{eff}}}{\varepsilon_{\mathrm{EM}}}\cdot\mathcal N(\mathcal R)\cdot\chi_{\{\,NE\ \ge\ \mathcal T(\mathcal R;N)\,\}}\ }
\]

%dsp-level:0
\[
\Gamma_{\mathrm{tr}}=\frac{x_+x_-}{1-x_+x_-}\in\Bigl[0,\tfrac13\Bigr]
\]

%dsp-level:0
\[
\qquad [\mathrm{BM.44}(2);\ \text{frozen}];
\]

%dsp-level:0
\[
\qquad f=1-\frac1k,\quad k=\frac{r}{\Delta x_{\mathrm B}(E)}\ge1
\]

%dsp-level:0
\[
\qquad [\mathrm{BM.88}(2)(\mathrm{ii})];
\]

%dsp-level:0
\[
\varepsilon_{\mathrm{eff}}=Nf\,\mathsf g(G)\,\varepsilon_{\mathrm{EM}}\implies\frac{\varepsilon_{\mathrm{eff}}}{\varepsilon_{\mathrm{EM}}}=N\,f\,\mathsf g(G)
\]

%dsp-level:0
\[
\qquad [\mathrm{BM.88}(4);\ \mathrm{BM.101}(2)];
\]

%dsp-level:0
\[
\qquad \mathsf g\in\{1,\ 1+G,\ \tfrac1{1+G}\}\ \ \text{priced,}\ \neg\text{adopted}
\]

%dsp-level:0
\[
\qquad [\mathrm{BM.88}(4);\ \mathrm{BM.104}(5)];
\]

%dsp-level:0
\[
\mathcal N(\mathcal R)\ :=\ \#\bigl\{\mathcal R'\ \text{in the same}\ (n,\sigma,Q)\ \text{class with}\ \lambda^{(n)}(\mathcal R')<\lambda^{(n)}(\mathcal R)\bigr\}
\]

%dsp-level:0
\[
\qquad [\text{the lighter-resonance channel count};\ \mathrm{BM.41}(3)\ \text{pattern}];
\]

%dsp-level:0
\[
\chi_{\{\cdot\}}\in\{0,1\}
\]

%dsp-level:0
\[
\qquad [\text{the exact indicator}];
\]

%dsp-level:0
\[
\qquad \textbf{non-derived clause:}\ \ \text{the product form}\ —\ \text{three derived factors and one count, multiplied};
\]

%dsp-level:0
\[
(5)\quad\textbf{the derived rows of the width (recorded at the definition site):}
\]

%dsp-level:0
\[
\qquad \Gamma_{\mathrm{tr}}(0)=0\ \ \text{exactly}\ [\mathrm{BM.23}(8);\ \mathrm{BM.44}(2)]\implies\boxed{\ \Gamma(\mathcal R;\,\varepsilon_q{=}0)=0\ \ \forall\mathcal R,\ \forall(E,N,G,\varepsilon_{\mathrm{EM}})\ }
\]

%dsp-level:0
\[
f(1)=0\ [\mathrm{BM.88}(5)]\implies\boxed{\ \Gamma(\mathcal R;\,r=\Delta x_{\mathrm B})=0\ \ \text{at the member's own envelope}\ }
\]

%dsp-level:0
\[
\mathcal N(\mathcal R)
\]

%dsp-level:0
\[
= 0\ \ \text{for the class's lightest resonance}\implies\boxed{\ \textbf{the ground state of every}\ (n,\sigma,Q)\ \textbf{class has width exactly zero}\ }
\]

%dsp-level:0
\[
\Gamma\le\tfrac13\cdot1\cdot N\,\mathsf g\cdot\mathcal N
\]

%dsp-level:0
\[
\qquad [\text{machine:}\ f<1,\ \varepsilon_{\mathrm{eff}}/\varepsilon_{\mathrm{EM}}=Nf\mathsf g<N\mathsf g];
\]

%dsp-level:0
\[
\qquad \tau_{\mathrm{life}}:=\frac1\Gamma\ge3\ \ \text{ticks at}\ \mathcal N=1,\ f\to1,\ N\mathsf g\to1
\]

%dsp-level:0
\[
\qquad [\mathrm{BM.44}(3);\ \text{machine}];
\]

%dsp-level:0
\[
(6)\quad\textbf{the standing orders, as scope clauses:}
\]

%dsp-level:0
\[
\qquad (\mathrm i)\ \ \neg\textbf{a cosmological coordinate exists:}\ \text{the ledger knows the local tick}\ \tau\ [\mathrm{A6};\ \mathrm{D1.8}],
\]

%dsp-level:0
\[
\ \text{the local curvature}\ G\ [\mathrm{BM.88}(3)],
\]

%dsp-level:0
\[
\ \text{the local lapse}\ N\ [\mathrm{BM.88}(2)(\mathrm{iii})]\ \text{and the local energy scale}\ E\ [\mathrm{BM.88}(2)(\mathrm i)];
\]

%dsp-level:0
\[
\ \neg\text{an age and}\ \neg\text{a global time};
\]

%dsp-level:0
\[
(\mathrm{ii})\quad\neg\textbf{a binary stable/unstable category:}
\]

%dsp-level:0
\[
\qquad \textbf{every eigenvalue of every multi-chain ensemble is a resonance};
\]

%dsp-level:0
\[
\qquad \text{``stable''}\ :=\ \Gamma\ \text{evaluates to}\ 0\ \text{at the ambient calibration point}\ [(5)];
\]

%dsp-level:0
\[
\qquad \text{``unstable'' or ``heavy''}\ :=\ \text{the local parameters cross}\ \mathcal T\ \text{or open}\ \mathcal N\ [(3),(4)];
\]

%dsp-level:0
\[
(\mathrm{iii})\quad\neg\textbf{a permanently forbidden state:}
\]

%dsp-level:0
\[
\qquad \mathcal T(\mathcal R;N)\in(0,\infty)\ \ \forall N\in(0,\infty)\implies\textbf{the threshold locus is never empty}\ [\mathrm{BM.116}(4)];
\]

%dsp-level:0
\[
(7)\quad\textbf{status and scope:}
\]

%dsp-level:0
\[
\qquad \textbf{derived:}\ (2),\ (3)\text{'s three forms},\ (5)\text{'s five rows};
\]

%dsp-level:0
\[
\qquad \textbf{constitutive:}\ (1)\text{'s tuple},\ (3)\text{'s and}\ (4)\text{'s definitional forms};
\]

%dsp-level:0
\[
\qquad \textbf{non-derived:}\ (4)\text{'s product form};
\]

%dsp-level:0
\[
\textbf{the scopes inherited at every use:}\ \ \mathrm{CT2.6}\text{'s continuum-limit hypotheses for the chart-side reads}\ [\mathrm{BM.34}\ \mathrm{Rmk};\ \mathrm{Def\text{-}BM.4}(5)\ \text{envelope discipline}];
\]

%dsp-level:0
\[
\qquad \mathrm{BM.8}(2)\text{'s stationarity hypothesis for the chart evolution}\ [\mathrm{BM.44}(3)];
\]

%dsp-level:0
\[
\neg\text{a P4 addition},
\]

%dsp-level:0
\[
\quad\neg\text{a new dimensionful scale}\ [\mathrm{P1}],
\]

%dsp-level:0
\[
\quad\neg\text{a measured value}\ [\mathrm{P5}],
\]

%dsp-level:0
\[
\quad\neg\text{a numerical coupling}\ [\mathrm{EM.25}(7)],
\]

%dsp-level:0
\[
\quad\neg\text{a lifetime in external seconds}\ [\mathrm{BM.44}(5)];
\]

%dsp-level:0
\[
\textbf{beyond this tier:}\ \ \mathsf g(G)\ [\mathrm{BM.104}(5)],
\]

%dsp-level:0
\[
\quad e_0\ [\mathrm{EM.25}(7)],
\]

%dsp-level:0
\[
\quad N\ \text{as solution data}\ [\mathrm{CT2.7};\ \mathrm{BM.34}(2)],\quad\rho_{\mathrm B}\ [\mathrm{A\text{-}BM.9}],
\]

%dsp-level:0
\[
\quad (\mathrm{R3})\ \text{the flavour composition},
\]

%dsp-level:0
\[
\quad (\mathrm{R8})\ \text{the}\ 2\text{-chain sector masses}.
\]

\textbf{Rmk (why "resonance" and not "particle").} Clause (6)(ii) is the tier's standing order written as a scope clause, and it is a statement about the mathematics rather than a naming preference: the ensemble's eigenvalues are indexed by the tuple of (1), each carries a threshold function and a width function, and both are exact symbolic functions of the four local parameters. Nothing in the construction distinguishes an eigenvalue whose width happens to vanish at one calibration point from one whose width happens not to; the difference is a value of a function, not a category. The word \textit{particle} would import a binary the ledger does not carry, and the word \textit{resonance} imports nothing: it names an eigenvalue with two functions attached.

\textbf{Rmk (the horizon row's quarantine).} Clause (3)'s third form is written with $r_s$, which BM.15(4) builds from $G_N$ — a named combination of the single scale for external comparison, not an input [P3]. The row is a change of variable inside the corpus's own monomial and it is quarantined exactly as BM.15's comparison discipline requires: no collapse reading, no horizon identification, and no claim that any resonance sits at any measured radius.

\textbf{Req.} A6, D0.8, D1.8, D2.9 (Parts 1–2); CT2.6, CT2.7, CT2.9, EM.25(7) (Part 2); P1, P3, P4, P5 (Parts 1–2); Def-H.20, Def-H.38, FT.22, FT.28, FT.44(4) (Part 3, Census); Census FT.35(2), FT.42(5) (row-wise); Def-BM.1, Def-BM.2, Def-BM.3(7), Def-BM.4(1), Def-BM.4(5), Def-BM.8(1), Def-BM.14(2), Def-BM.14(3), BM.8(2), BM.13(2), BM.15(4), BM.16(1), BM.23(8), BM.34(2), BM.34(4), BM.34(7), BM.41(3), BM.44(2), BM.44(3), BM.44(4), BM.44(5), BM.88(2), BM.88(3), BM.88(4), BM.88(5), BM.101(2), BM.103(5), BM.104(5), BM.111(2), A-BM.9 (this part); exact Laurent-monomial arithmetic with $\mathfrak q$-exponent censuses, exact rational function rows [math; machine: \texttt{checks/\allowbreak{}verify\_\allowbreak{}H\_\allowbreak{}antimatter\_\allowbreak{}resonance\_\allowbreak{}spectrum.py} §E].

\textbf{Pf.}

%dsp-level:0
\[
\text{Definition block: the constitutive clauses are}\ (1)\text{'s tuple and the definitional forms of}\ (3),
\]

%dsp-level:0
\[
(4);
\]

%dsp-level:0
\[
\ \text{the non-derived clause is}\ (4)\text{'s product form};
\]

%dsp-level:0
\[
\ (2),
\]

%dsp-level:0
\[
(3)\text{'s three exact forms and}\ (5)\text{'s five rows are derived and machine-checked};
\]

%dsp-level:0
\[
\ (6),
\]

%dsp-level:0
\[
(7)\ \text{are scope records}.\ \blacksquare
\]

\textbf{Phys.} A resonance is a five-entry tuple — how many chains, which seal, which charge class, which six states, which winding — and it carries two functions of the local parameters. The first says when it appears: the rest energy the metric's own ADM split gives is the lapse times the mass, so the resonance materializes when the local scale times the lapse reaches its edge count times the single scale over the log of two, and the same condition read through the calibration lock is a square root of the local curvature, or read through the comparison radius it is twice that radius over the lapse. The second says how fast it goes: the frozen per-tick transition weight, times the envelope fraction that vanishes at the member's own width, times the effective coupling over the bare one, times the count of strictly lighter resonances in its class, times the exact indicator that it is materialized at all. Both functions vanish for whole regions rather than for special objects — at the sharp crystallization threshold nothing decays, at the envelope nothing decays, and the lightest member of every class has nothing lighter to go to, so its width is zero by count. That is the entire content of "stable": a function evaluating to zero where it is being evaluated. Two scope clauses carry the tier's standing orders: no cosmological coordinate exists, because the ledger owns a tick, a curvature, a lapse and a scale and nothing global, and no binary category exists, because every eigenvalue is a resonance and the adjectives are values. The third says nothing is forbidden, since every threshold is a finite positive number at every positive lapse, so every locus is non-empty.

\end{definition}

\begin{theorem}[{\texorpdfstring{BM.114 — The resonance spectrum: the two characteristic polynomials, the orbit law, and the one exact collision}{BM.114 - The resonance spectrum: the two characteristic polynomials, the orbit law, and the one exact collision}}]\label{tag:BM-114}
\textbf{Thm.}

%dsp-level:0
\[
(1)\quad\textbf{the transport-register characteristic polynomial at every chain count:}
\]

%dsp-level:0
\[
\qquad W_n(Q)=\bigl[z^Q\bigr]\bigl(2+z+z^{-1}\bigr)^{n},
\]

%dsp-level:0
\[
\qquad 2+z+z^{-1}=z^{-1}(z+1)^2
\]

%dsp-level:0
\[
\qquad [\mathrm{Def\text{-}BM.13}(3);\ \mathrm{BM.103}(3)]:
\]

%dsp-level:0
\[
\boxed{\ \bigl(2+z+z^{-1}\bigr)^{n}\ =\ \frac{(z+1)^{2n}}{z^{n}}\ \ \text{identically in}\ \mathbb Z[z,z^{-1}]\implies W_n(Q)=\binom{2n}{Q+n}\ \ \text{exactly, on all}\ 2n+1\ \text{classes}\ }
\]

%dsp-level:0
\[
[\text{machine: the Laurent identity and the exhaustive}\ 4^n\ \text{end-sign census at}\ n=1,2,3];
\]

%dsp-level:0
\[
\qquad \sum_{Q=-n}^{n}W_n(Q)=2^{2n}=4^{n}
\]

%dsp-level:0
\[
\qquad [\text{machine}];
\]

%dsp-level:0
\[
\textbf{the single root:}
\]

%dsp-level:0
\[
\qquad z
\]

%dsp-level:0
\[
= -1\ \ \text{of multiplicity}\ 2n\ \ \text{(the}\ 2n\ \text{directed edges), and the}\ 2n+1\ \text{class weights are its binomial coefficients}
\]

%dsp-level:0
\[
\qquad [\mathrm{BM.103}(3)\ \text{transported to every}\ n;\ \text{machine}];
\]

%dsp-level:0
\[
\textbf{the three formed levels:}\qquad\begin{array}{c|c|c|c}n & \text{seal} & W_n(Q) & \text{sum}\\\hline1 & \neg\text{sealed (a lone attachment is never sealed)} & (1,\ 2,\ 1) & 4=4^1\\2 & \delta\ \text{(meson)} & (1,\ 4,\ \mathbf6,\ 4,\ 1) & 16=4^2\\3 & \varepsilon\ \text{(baryon)} & (1,\ 6,\ 15,\ \mathbf{20},\ 15,\ 6,\ 1) & 64=4^3\end{array}\qquad[\text{machine: exact binomials}];
\]

%dsp-level:0
\[
\boxed{\ \textbf{the}\ n{=}1\ \textbf{row is formed but}\ \neg\textbf{sealed:}\quad g(1)=0\implies\mathcal C=0\implies\text{the colour block vanishes}\ }
\]

%dsp-level:0
\[
\qquad [\mathrm{CT3.17a}(1),(2);\ \mathrm{FT.19};\ \text{machine}];
\]

%dsp-level:0
\[
(2)\quad\textbf{the}\ \varepsilon\textbf{-seal's antisymmetry and the}\ \delta\textbf{-seal's symmetry (the only two antisymmetries the corpus owns):}
\]

%dsp-level:0
\[
\qquad B=\varepsilon_{ijk}q^iq^jq^k\ \ \text{is antisymmetric in its three colour slots};
\]

%dsp-level:0
\[
\qquad M=\delta^i{}_jc^j\bar c_i\ \ \text{is symmetric in its two}
\]

%dsp-level:0
\[
\qquad [\mathrm{D3.9};\ \mathrm{CT3.15};\ \mathrm{CT3.17};\ \mathrm{BM.106}(5)]:
\]

%dsp-level:0
\[
\textbf{colour is}\ \neg\textbf{observable by construction}
\]

%dsp-level:0
\[
\qquad [\mathrm{CT3.22}]\implies\boxed{\ \textbf{the colour slot symmetry is}\ \neg\textbf{the statistics}\ —\ \text{the statistics is}\ \varepsilon(c)\ \text{on the winding carrier}\ [\mathrm{BM.113}(2)]\ }
\]

%dsp-level:0
\[
\textbf{the transport assembly is sign-free at every}\ n:
\]

%dsp-level:0
\[
\qquad W_n(Q)=\binom{2n}{Q+n}>0\ \ \text{on all}\ 2n+1\ \text{classes}
\]

%dsp-level:0
\[
\qquad [\mathrm{BM.103}(3);\ \text{machine}];
\]

%dsp-level:0
\[
(3)\quad\textbf{the defect-register eigenvalues (the resonance ladder):}
\]

%dsp-level:0
\[
\qquad \lambda^{(n)}(\mathbf s)=\sum_{a=1}^{n}D_S(s_a),
\]

%dsp-level:0
\[
\qquad D_S\in\bigl\{2,\ 4,\ 2-\sqrt2,\ 1,\ 2-2\cos\tfrac\pi{10},\ 2-\sqrt3\bigr\}
\]

%dsp-level:0
\[
\qquad [\mathrm{Def\text{-}BM.14}(3);\ \mathrm{Def\text{-}BM.17}(2)]:
\]

%dsp-level:0
\[
m^2_{\mathrm{pole}}(\mathcal R)=\hat\mu'^2\,\lambda^{(n)}(\mathbf s)\qquad[\mathrm{Def\text{-}H.37}(2);\ \mathrm{FT.54}(4);\ \text{machine}];
\]

%dsp-level:0
\[
\textbf{the exact eigenvalue field:}
\]

%dsp-level:0
\[
\qquad \mathfrak u:=2\cos\tfrac\pi{10},
\]

%dsp-level:0
\[
\qquad \mathfrak u^4-5\mathfrak u^2+5=0,
\]

%dsp-level:0
\[
\qquad \sqrt5=2\mathfrak u^2-5\in\mathbb Q(\mathfrak u),
\]

%dsp-level:0
\[
\qquad 2-2\cos\tfrac\pi{10}=2-\mathfrak u:
\]

%dsp-level:0
\[
\boxed{\ \mathbb M\ :=\ \mathbb Q\bigl(\sqrt2,\ \sqrt3,\ \mathfrak u\bigr),\qquad[\mathbb M:\mathbb Q]=2\cdot2\cdot4=16,\qquad\textbf{basis}\ \bigl\{\sqrt2^{\,i}\sqrt3^{\,j}\,\mathfrak u^{\,\ell}\bigr\}_{i,j\in\{0,1\},\ \ell\in\{0,1,2,3\}}\ }
\]

%dsp-level:0
\[
\qquad [\text{machine: the three defining relations verified exactly, and}\ (\sqrt2)^2{=}2,(\sqrt3)^2{=}3,(2\mathfrak u^2{-}5)^2{=}5];
\]

%dsp-level:0
\[
\textbf{every resonance eigenvalue lies in}\ \mathbb M,
\]

%dsp-level:0
\[
\ \text{and equality of two eigenvalues is decidable by basis-coordinate equality}
\]

%dsp-level:0
\[
\qquad [\text{machine:}\ \mathbb M\text{'s basis is linearly independent over}\ \mathbb Q];
\]

%dsp-level:0
\[
(4)\quad\textbf{the orbit law (the multiplicity is the stabilizer index):}
\]

%dsp-level:0
\[
\qquad S_n\ \text{acts on}\ \mathcal S^n\ \text{by permuting the constituent slots;}\ \ \lambda^{(n)}\ \textbf{is}\ S_n\textbf{-invariant}\ [\text{machine: a sum}]:
\]

%dsp-level:0
\[
\#\bigl(\text{orbits}\bigr)=\binom{6+n-1}{n},
\]

%dsp-level:0
\[
\qquad |\text{orbit of}\ \mathbf s|=\frac{n!}{|\operatorname{Stab}(\mathbf s)|},
\]

%dsp-level:0
\[
\qquad \sum_{\text{orbits}}|\text{orbit}|=6^{n}
\]

%dsp-level:0
\[
\qquad [\text{machine: exact at}\ n=1,2,3];
\]

%dsp-level:0
\[
\textbf{the three levels:}\qquad\begin{array}{c|c|c|c|c}n & \#\text{orbits} & \text{orbit sizes} & \sum & \text{multiplicity census}\\\hline1 & \binom61=6 & 1 & 6 & \{1^6\}\\2 & \binom72=21 & 1\ (\text{6 diagonal}),\ 2\ (\text{15 off-diagonal}) & 36 & \{1^6,\ 2^{15}\}\\3 & \binom83=56 & 1\ (\text{6}),\ 3\ (\text{30}),\ 6\ (\text{20}) & 216 & \text{see}\ (5)\end{array}
\]

%dsp-level:0
\[
[\text{machine: exhaustive over the}\ 6^n\ \text{ordered tuples}];\qquad\boxed{\ 15=\binom62=\dim\wedge^2\mathbb C^6,\qquad20=\binom63\ }
\]

%dsp-level:0
\[
\textbf{the count is a multiset count,}\ \neg\textbf{a representation dimension:}
\]

%dsp-level:0
\[
\qquad \neg\exists\ \text{a flavour}\ \mathsf{SU}(6)\ \text{in the inventory}\ [\mathrm{BM.103}(7);\ \text{the flavour symmetry is}\ C_3\ [\mathrm{FT.111}(2);\ \mathrm{FT.142}(5)]];
\]

%dsp-level:0
\[
(5)\quad\textbf{the three-chain ladder in full:}
\]

%dsp-level:0
\[
\qquad 56\ \text{orbits}\ \longrightarrow\ \boxed{\ 55\ \textbf{distinct eigenvalues}\ }
\]

%dsp-level:0
\[
\qquad [\text{machine: exact equality in}\ \mathbb M]:
\]

%dsp-level:0
\[
\textbf{exactly one collision, and it is the powers-of-two triple:}
\]

%dsp-level:0
\[
\qquad \lambda^{(3)}(s_1,s_1,s_1)=2+2+2=6=4+1+1=\lambda^{(3)}(s_2,s_4,s_4),
\]

%dsp-level:0
\[
\qquad \text{multiplicity}\ 1+3=4
\]

%dsp-level:0
\[
[\text{machine: the exact rational identity};\ \mathrm{BM.82}(4)\text{'s collision census re-cited at the constituent level: three of the six defect values are the powers}\ 1,2,4];
\]

%dsp-level:0
\[
\textbf{the multiplicity census of the 55:}
\]

%dsp-level:0
\[
\qquad \boxed{\ \{1^{5},\ 3^{29},\ 4^{1},\ 6^{20}\}\ }
\]

%dsp-level:0
\[
\qquad \text{sum}\ =\ 5+87+4+120=216=6^3
\]

%dsp-level:0
\[
\qquad [\text{machine: exact}];
\]

%dsp-level:0
\[
\textbf{the multiplicity is always a divisor of}\ n!=6:
\]

%dsp-level:0
\[
\qquad \{1,3,4,6\}\subset\{1,2,3,6\}\cup\{4\},
\]

%dsp-level:0
\[
\qquad \text{and the}\ 4\ \text{is the collision's}\ 1+3\ \ [\text{machine}];
\]

%dsp-level:0
\[
(6)\quad\textbf{the minimal-polynomial degree census (the subfield ladder of}\ \mathbb M\text{):}\qquad\begin{array}{c|c|c}n & \text{degree census over the distinct eigenvalues} & \max\\\hline1 & \{1^{3},\ 2^{2},\ 4^{1}\} & 4\\2 & \{1^{6},\ 2^{8},\ 4^{5},\ 8^{2}\} & 8\\3 & \{1^{9},\ 2^{20},\ 4^{15},\ 8^{10},\ 16^{1}\} & 16=[\mathbb M:\mathbb Q]\end{array}
\]

%dsp-level:0
\[
[\text{machine: each minimal polynomial computed exactly by linear dependence in the}\ 16\text{-dimensional}\ \mathbb Q\text{-space}\ \mathbb M,\ \text{then verified by exact evaluation at its own root}];
\]

%dsp-level:0
\[
\textbf{the pinned rows:}
\]

%dsp-level:0
\[
\qquad D_S(s_3)=2-\sqrt2:\ x^2-4x+2;
\]

%dsp-level:0
\[
\qquad D_S(s_5)=2-\mathfrak u:\ x^4-8x^3+19x^2-12x+1;
\]

%dsp-level:0
\[
\qquad D_S(s_6)=2-\sqrt3:\ x^2-4x+1
\]

%dsp-level:0
\[
\qquad [\mathrm{Def\text{-}BM.14}(3);\ \text{machine re-verified}];
\]

%dsp-level:0
\[
\textbf{the unique degree-16 resonance eigenvalue:}
\]

%dsp-level:0
\[
\qquad \lambda^{(3)}(s_3,s_5,s_6)=(2-\sqrt2)+(2-\mathfrak u)+(2-\sqrt3)=6-\sqrt2-\mathfrak u-\sqrt3:
\]

%dsp-level:0
\[
\qquad \textbf{it generates}\ \mathbb M\ \text{over}\ \mathbb Q
\]

%dsp-level:0
\[
[\text{machine: its minimal polynomial has degree exactly}\ 16;\ \text{the leading rows are}\ x^{16}-96x^{15}+4260x^{14}-115920x^{13}+\cdots+249713695921];
\]

%dsp-level:0
\[
\textbf{the minimal eigenvalue at each level is a quartic:}
\]

%dsp-level:0
\[
\qquad n\,\bigl(2-\mathfrak u\bigr):
\]

%dsp-level:0
\[
\qquad n=1:\ x^4-8x^3+19x^2-12x+1;
\]

%dsp-level:0
\[
\qquad n=2:\ x^4-16x^3+76x^2-96x+16;
\]

%dsp-level:0
\[
\qquad n=3:\ x^4-24x^3+171x^2-324x+81
\]

%dsp-level:0
\[
[\text{machine: the three quartics computed exactly and each verified at its own root}];
\]

%dsp-level:0
\[
(7)\quad\textbf{the ladder is a resonance ladder, not a stability ladder:}
\]

%dsp-level:0
\[
\qquad \neg\exists\ \text{a derived eigenvalue stability condition}\ [\mathrm{BM.103}(6):\ \text{the three grounds re-cited}\ —\ (\mathrm{M1})/(\mathrm{M2})\ \text{undeclared},\ \mathrm{FT.66}(4)\text{'s index-zero refutation, and the count}\to\text{strength bar}]:
\]

%dsp-level:0
\[
\boxed{\ \textbf{what orders the ladder is the eigenvalue's own value, and what decides its appearance is}\ \mathcal T(\mathcal R;N),\ \neg\textbf{a stability class}\ }
\]

%dsp-level:0
\[
\qquad [\mathrm{Def\text{-}BM.17}(3),(6)(\mathrm{ii});\ \mathrm{BM.115};\ \mathrm{BM.116}];
\]

%dsp-level:0
\[
(8)\quad\textbf{status:}\qquad\textbf{derived:}\ (1)\text{--}(7);
\]

%dsp-level:0
\[
\neg\text{a declaration},
\]

%dsp-level:0
\[
\quad\neg\text{a P4 addition},
\]

%dsp-level:0
\[
\quad\neg\text{a new dimensionful scale}\ [\mathrm{P1}],
\]

%dsp-level:0
\[
\quad\neg\text{a measured value}\ [\mathrm{P5}],
\]

%dsp-level:0
\[
\quad\neg\text{a species identification}\ [\mathrm{Def\text{-}BM.3}(7)],
\]

%dsp-level:0
\[
\quad\neg\text{a currency conversion}\ [\mathrm{BM.83}(3)];
\]

\textbf{Req.} D2.9, D3.9 (Parts 1–2); CT3.15, CT3.17, CT3.17a(1), CT3.17a(2), CT3.22 (Part 2); P1, P4, P5 (Parts 1–2); Def-H.37(2), FT.19, FT.54(4), FT.111(2), FT.142(5) (Part 3); Def-BM.3(7), Def-BM.13(3), Def-BM.14(3), Def-BM.17(2), BM.66(6), BM.82(4), BM.83(3), BM.103(3), BM.103(6), BM.103(7), BM.106(5), BM.111(2), BM.113(2) (this part); Laurent-polynomial identities over $\mathbb Z[z,z^{-1}]$, exhaustive binomial and multiset-orbit enumeration, exact linear algebra over $\mathbb Q$ in a $16$-dimensional algebra, exact real ordering by a rigorous dyadic interval embedding [math; machine: \texttt{checks/\allowbreak{}verify\_\allowbreak{}H\_\allowbreak{}antimatter\_\allowbreak{}resonance\_\allowbreak{}spectrum.py} §F].

\textbf{Pf.}

%dsp-level:0
\[
(1):\  2+z+z^{-1}=z^{-1}(z^2+2z+1)=z^{-1}(z+1)^2\ [\text{machine: one Laurent identity}];
\]

%dsp-level:0
\[
\qquad \bigl(z^{-1}(z+1)^2\bigr)^n=z^{-n}(z+1)^{2n}\ [\text{machine}]
\]

%dsp-level:0
\[
[z^Q]z^{-n}(z+1)^{2n}
\]

%dsp-level:0
\[
= [z^{Q+n}](z+1)^{2n}
\]

%dsp-level:0
\[
= \binom{2n}{Q+n}\ [\text{machine: against the exhaustive}\ 4^n\ \text{census at}\ n=1,2,3];
\]

%dsp-level:0
\[
\qquad \sum_Q\binom{2n}{Q+n}=2^{2n}\ [\text{machine}]
\]

%dsp-level:0
\[
g(1)=0\ [\mathrm{CT3.17a}(2)];\qquad|\mathrm{Colour}(\Sigma)|=1\implies\mathcal C=0\implies\text{the block vanishes}\ [\mathrm{CT3.17a}(1);\ \mathrm{FT.19}]
\]

%dsp-level:0
\[
(2):\  \mathrm{D3.9},\ \mathrm{CT3.15},\ \mathrm{CT3.17}\ \text{re-cited};
\]

%dsp-level:0
\[
\qquad \mathrm{CT3.22}\text{'s non-observability re-cited};
\]

%dsp-level:0
\[
\qquad W_n(Q)=\binom{2n}{Q+n}>0\ [\text{machine: positivity of binomials}]
\]

%dsp-level:0
\[
(3):\  \mathfrak u
\]

%dsp-level:0
\[
= 2\cos\tfrac\pi{10}:\quad\cos^2\tfrac\pi{10}
\]

%dsp-level:0
\[
= \tfrac{5+\sqrt5}8\implies\mathfrak u^2
\]

%dsp-level:0
\[
= \tfrac{5+\sqrt5}2\implies\mathfrak u^4
\]

%dsp-level:0
\[
= \tfrac{15+5\sqrt5}2\implies\mathfrak u^4-5\mathfrak u^2+5
\]

%dsp-level:0
\[
= \tfrac{15+5\sqrt5-25-5\sqrt5}2+5
\]

%dsp-level:0
\[
= 0\ [\text{machine: exact in}\ \mathbb Q(\sqrt5)]
\]

%dsp-level:0
\[
\sqrt5
\]

%dsp-level:0
\[
= 2\mathfrak u^2-5\ [\text{machine:}\ (2\mathfrak u^2-5)^2=4\cdot\tfrac{5+\sqrt5}2\cdot2-20\mathfrak u^2+25\ \text{reduced by}\ \mathfrak u^4=5\mathfrak u^2-5\ \text{to}\ 5];
\]

%dsp-level:0
\[
\qquad [\mathbb Q(\mathfrak u):\mathbb Q]=4\ [\mathfrak u^4-5\mathfrak u^2+5\ \text{irreducible:}\ \neg\text{a rational root and}\ \mathfrak u^2=\tfrac{5+\sqrt5}2\notin\mathbb Q]
\]

%dsp-level:0
\[
[\mathbb M:\mathbb Q]
\]

%dsp-level:0
\[
= [\mathbb M:\mathbb Q(\mathfrak u)]\cdot[\mathbb Q(\mathfrak u):\mathbb Q]
\]

%dsp-level:0
\[
= 4\cdot4
\]

%dsp-level:0
\[
= 16\ [\sqrt2,\sqrt3\notin\mathbb Q(\mathfrak u)\ \text{by the degree census of}\ (6);\ \text{machine: the}\ 16\ \text{basis elements are}\ \mathbb Q\text{-linearly independent}]
\]

%dsp-level:0
\[
(4):\  \lambda^{(n)}(\pi\cdot\mathbf s)=\sum_aD_S(s_{\pi(a)})=\sum_aD_S(s_a)=\lambda^{(n)}(\mathbf s)\ [\text{machine}];
\]

%dsp-level:0
\[
\qquad \#\text{orbits}=\binom{6+n-1}n\ [\text{multiset count};\ \text{machine: exhaustive at}\ n=1,2,3]
\]

%dsp-level:0
\[
|\text{orbit}|=n!/|\operatorname{Stab}|\ \text{with}\ |\operatorname{Stab}|=\prod_{\text{distinct values}}(\text{multiplicity})!\ [\text{machine}];
\]

%dsp-level:0
\[
\qquad \text{the size sums}\ 6,36,216\ [\text{machine}]
\]

%dsp-level:0
\[
\binom62=15=\dim\wedge^2\mathbb C^6=\tfrac{6\cdot5}2;\qquad\binom63=20\ [\text{machine}]
\]

%dsp-level:0
\[
(5):\  2+2+2=6\ \text{and}\ 4+1+1=6\ \text{in}\ \mathbb Q\subset\mathbb M\ [\text{machine: basis-coordinate equality}];
\]

%dsp-level:0
\[
\qquad \text{the exhaustive census finds no second collision among the}\ 56\ [\text{machine}]
\]

%dsp-level:0
\[
\text{the multiplicity census}\ \{1^5,3^{29},4^1,6^{20}\}:\quad5+3\cdot29+4+6\cdot20=5+87+4+120=216\ [\text{machine}]
\]

%dsp-level:0
\[
(6):\  \text{for each eigenvalue}\ x,
\]

%dsp-level:0
\[
\ \text{the smallest}\ d\ \text{with}\ \{1,x,\dots,x^d\}\ \mathbb Q\text{-dependent in the}\ 16\text{-dimensional}\ \mathbb M\ \text{gives the degree, and the dependency gives the monic polynomial}\ [\text{machine: exact Fraction elimination}];
\]

%dsp-level:0
\[
\ \text{each is verified by exact evaluation at its own root}\ [\text{machine}]
\]

%dsp-level:0
\[
x^2-4x+2\ \text{at}\ 2-\sqrt2:\ (2-\sqrt2)^2-4(2-\sqrt2)+2=(6-4\sqrt2)-(8-4\sqrt2)+2=0\ [\text{machine}];
\]

%dsp-level:0
\[
\qquad x^2-4x+1\ \text{at}\ 2-\sqrt3:\ (7-4\sqrt3)-(8-4\sqrt3)+1=0\ [\text{machine}]
\]

%dsp-level:0
\[
(2-\mathfrak u):\ (2-x)^4-5(2-x)^2+5=x^4-8x^3+19x^2-12x+1\ [\text{machine: the exact expansion}];
\]

%dsp-level:0
\[
\qquad \text{the three minimal quartics and the degree-16 polynomial computed and root-verified}\ [\text{machine}]
\]

%dsp-level:0
\[
(7):\  \mathrm{BM.103}(6)\text{'s three grounds re-cited};\qquad(8):\  \text{the standing-rule scans}\ [\text{machine}].\ \blacksquare
\]

\textbf{Phys.} Two characteristic polynomials come out, one per currency, and neither is fitted. In the transport register the generating function is the same Laurent monomial at every chain count — one plus z plus one over z, raised to n, which is one plus z to the two n over z to the n — so the class weight is a binomial coefficient, the sum is four to the n, and the whole spectrum is one root at minus one with multiplicity twice the chain count. The first level is formed but not sealed, since a lone attachment is never sealed, which is why the ladder starts at two. In the defect register the eigenvalues are sums of the six branch values, and they all live in one exact field of degree sixteen over the rationals built from root two, root three and twice the cosine of eighteen degrees, whose own relation makes root five a polynomial in it — so equality of two eigenvalues is a decidable integer question rather than a numerical one. The orbit law is clean: the symmetric group permutes the constituent slots, the eigenvalue is invariant, the orbit count is a multiset binomial, and the multiplicity of each eigenvalue is the orbit size, a divisor of n factorial. Fifteen and twenty appear, and they are exactly the wedge-square dimension and the three-subset count of the six states — but the tier says plainly that these are counts and not representation dimensions, because the flavour symmetry the corpus derives is a cyclic three and not a special unitary six. The three-chain ladder has fifty-six orbits and fifty-five eigenvalues, and the single collision is the one the constituent census had already found: two plus two plus two equals four plus one plus one, because three of the six defect values are the powers one, two and four. The minimal-polynomial degrees climb one, two, four, eight, sixteen with the number of distinct irrational ingredients, and exactly one three-chain resonance — the sum of the root-two, cosine and root-three branches — generates the whole field.

\end{theorem}

\begin{theorem}[{\texorpdfstring{BM.115 — The materialization threshold: the exact function of $E$, $N$ and $G$, and the never-empty locus}{BM.115 - The materialization threshold: the exact function of , and , and the never-empty locus}}]\label{tag:BM-115}
\textbf{Thm.}

%dsp-level:0
\[
(1)\quad\textbf{the threshold function:}
\]

%dsp-level:0
\[
\qquad \mathcal T_0(\mathcal R)=\frac{E_{\mathrm{int}}(\mathcal R)\,\mathfrak h}{\ln2},
\]

%dsp-level:0
\[
\qquad \mathcal T(\mathcal R;N)=\frac{\mathcal T_0(\mathcal R)}N,
\]

%dsp-level:0
\[
\qquad \textbf{materialization}\iff NE\ge\mathcal T_0(\mathcal R)
\]

%dsp-level:0
\[
\qquad [\mathrm{Def\text{-}BM.17}(3);\ \text{machine}]:
\]

%dsp-level:0
\[
[\mathcal T_0]=[\mathfrak h]=\mathfrak q^{-1/2}\cdot\mathfrak q^{+1/2}\cdot\mathfrak q^{0}=\ \ \text{an energy:}\quad[E]=\mathfrak q^{+1/2},
\]

%dsp-level:0
\[
\qquad [\mathcal T]=\mathfrak q^{+1/2}
\]

%dsp-level:0
\[
\qquad [\mathrm{CT2.1};\ \mathrm{BM.102}(2);\ \text{machine: the exponent arithmetic}];
\]

%dsp-level:0
\[
\mathcal T\ \textbf{is strictly decreasing in the lapse:}
\]

%dsp-level:0
\[
\qquad N_1<N_2\implies\mathcal T(\mathcal R;N_1)>\mathcal T(\mathcal R;N_2)>0
\]

%dsp-level:0
\[
\qquad [\text{machine: exact rationals}];
\]

%dsp-level:0
\[
\boxed{\ \textbf{a deeper lapse lowers every threshold at once}\ —\ \text{the whole ladder moves as one function of}\ N\ }
\]

%dsp-level:0
\[
(2)\quad\textbf{the two exact threshold ladders:}
\]

%dsp-level:0
\[
\qquad \textbf{(i) the computed-cell ladder}\quad E_{\mathrm{int}}(s)\in\{2,\ 6,\ 12,\ 24,\ 60\}
\]

%dsp-level:0
\[
\qquad [\mathrm{BM.13}(2);\ \mathrm{FT.28};\ \text{Census FT.35}(2),\ \text{FT.42}(5)\ \text{row-wise}]:
\]

%dsp-level:0
\[
\begin{array}{c|c|l|c}n & \#\{E_{\mathrm{int}}\} & \text{the exact ladder} & \text{span ratio}\\\hline1 & 5 & 2,\ 6,\ 12,\ 24,\ 60 & 30\\2 & 15 & 4,\ 8,\ 12,\ 14,\ 18,\ 24,\ 26,\ 30,\ 36,\ 48,\ 62,\ 66,\ 72,\ 84,\ 120 & 30\\3 & 33 & 6,\ 10,\ 14,\ 16,\ 18,\ 20,\ 24,\ 26,\ 28,\ 30,\ 32,\ 36,\ 38,\ 42,\ 48,\ 50,\ 54,\ 60,\ 64,\ 68,\ 72,\ 74,\ 78,\ 84,\ 86,\ 90,\ 96,\ 108,\ 122,\ 126,\ 132,\ 144,\ 180 & 30\end{array}
\]

%dsp-level:0
\[
[\text{machine: the exact sum sets with repetitions}];
\]

%dsp-level:0
\[
\qquad \boxed{\ \min E_{\mathrm{int}}^{(n)}=2n\ \ \text{(the}\ n\ \text{digon floors)},\qquad\max E_{\mathrm{int}}^{(n)}=60n,\qquad\frac{\max}{\min}=30\ \ \textbf{at every}\ n\ }
\]

%dsp-level:0
\[
\textbf{(ii) the six-state ladder}\quad E_{\mathrm{int}}(s)=2E(s)\in\{12,\ 24,\ 60\}\ \ \text{on the five Platonic solids}
\]

%dsp-level:0
\[
\qquad [\mathrm{Def\text{-}BM.14}(1);\ \mathrm{BM.60}(1):\ 12:24:60=1:2:5]:
\]

%dsp-level:0
\[
\begin{array}{c|c|l|c}n & \#\{E_{\mathrm{int}}\} & \text{the exact ladder} & \text{span ratio}\\\hline1 & 3 & 12,\ 24,\ 60 & 5\\2 & 6 & 24,\ 36,\ 48,\ 72,\ 84,\ 120 & 5\\3 & 10 & 36,\ 48,\ 60,\ 72,\ 84,\ 96,\ 108,\ 132,\ 144,\ 180 & 5\end{array}
\]

%dsp-level:0
\[
[\text{machine: the exact sum sets}];
\]

%dsp-level:0
\[
\qquad \boxed{\ \min=12n,\qquad\max=60n,\qquad\frac{\max}{\min}=5=\frac{60}{12}\ \ \textbf{at every}\ n\ }
\]

%dsp-level:0
\[
\qquad [\text{the floor ladder's own ratio}\ 1:2:5\ \text{re-cited}];
\]

%dsp-level:0
\[
\textbf{the two ladders are nested, and the nesting is exact:}
\]

%dsp-level:0
\[
\qquad \{12,24,60\}\subset\{2,6,12,24,60\}\implies\text{every six-state threshold is a computed-cell threshold};
\]

%dsp-level:0
\[
\textbf{the two extra cells}\ 2\ \text{and}\ 6\ \text{are the digon floor}\ [\mathrm{FT.28}]\ \text{and the baryon junction}\ [\text{Census FT.35}(2)]:
\]

%dsp-level:0
\[
\qquad \text{they lower the minimum from}\ 12n\ \text{to}\ 2n\ \text{and raise the span from}\ 5\ \text{to}\ 30\ [\text{machine}];
\]

%dsp-level:0
\[
\textbf{the meson and baryon rows on the computed-cell ladder:}
\]

%dsp-level:0
\[
\qquad m_{\mathrm{hid}}^{(2)}\in\Bigl[\frac{4\ln2}{\mathfrak h},\ \frac{120\ln2}{\mathfrak h}\Bigr],
\]

%dsp-level:0
\[
\qquad m_{\mathrm{hid}}^{(3)}\in\Bigl[\frac{6\ln2}{\mathfrak h},\ \frac{180\ln2}{\mathfrak h}\Bigr]
\]

%dsp-level:0
\[
\qquad [\mathrm{BM.111}(5);\ \text{machine}];
\]

%dsp-level:0
\[
\textbf{and on the six-state ladder:}
\]

%dsp-level:0
\[
\qquad m_{\mathrm{hid}}^{(2)}\in\Bigl[\frac{24\ln2}{\mathfrak h},\ \frac{120\ln2}{\mathfrak h}\Bigr],
\]

%dsp-level:0
\[
\qquad m_{\mathrm{hid}}^{(3)}\in\Bigl[\frac{36\ln2}{\mathfrak h},\ \frac{180\ln2}{\mathfrak h}\Bigr]
\]

%dsp-level:0
\[
\qquad [\text{machine}];
\]

%dsp-level:0
\[
(3)\quad\textbf{the curvature read:}
\]

%dsp-level:0
\[
\qquad G=2\pi E_{\mathrm{int}}^2\ln2\in\mathfrak q^0,
\]

%dsp-level:0
\[
\qquad G\ \ \mathfrak h\text{-free},\ \Lambda\text{-free},\ \mathrm{P4}\text{-free}
\]

%dsp-level:0
\[
\qquad [\mathrm{BM.16}(1);\ \mathrm{BM.88}(3);\ \text{machine: the three freedom censuses re-cited}]:
\]

%dsp-level:0
\[
\boxed{\ E_{\mathrm{int}}=\sqrt{\frac{G}{2\pi\ln2}}\implies\mathcal T(\mathcal R;\,G,N)=\frac{\mathfrak h}{N\ln2}\sqrt{\frac{G}{2\pi\ln2}}\ }
\]

%dsp-level:0
\[
\qquad [\text{machine: the exact substitution}];
\]

%dsp-level:0
\[
G^{(n)}\in\bigl[2\pi\ln2\cdot4n^2,\ \ 2\pi\ln2\cdot3600n^2\bigr],
\]

%dsp-level:0
\[
\qquad \frac{G_{\max}}{G_{\min}}=900=30^2\ \ \textbf{at every}\ n
\]

%dsp-level:0
\[
\qquad [\text{machine: exact integers}];
\]

%dsp-level:0
\[
\textbf{the curvature is a derived function of the count, so the threshold is a derived function of the local curvature}\ —\ \neg\text{an independent datum}\ [\mathrm{BM.88}(3);\ \text{machine}];
\]

%dsp-level:0
\[
(4)\quad\textbf{the horizon read (comparison scale, quarantined):}
\]

%dsp-level:0
\[
\qquad r_s=2G_Nm_{\mathrm{hid}}=\tfrac12\mathfrak h E_{\mathrm{int}},
\]

%dsp-level:0
\[
\qquad \Delta x_{\mathrm B}\cdot r_s=2G_N=\frac{\mathfrak h^2}{2\ln2}
\]

%dsp-level:0
\[
\qquad [\mathrm{BM.15}(4);\ \mathrm{CT2.9};\ \mathrm{P3}]:
\]

%dsp-level:0
\[
\boxed{\ \mathcal T(\mathcal R;\,r_s,N)=\frac{2\,r_s}{N}\ }
\]

%dsp-level:0
\[
\qquad [\text{machine:}\ \tfrac{\mathfrak h}{N\ln2}\cdot\tfrac{2r_s}{\mathfrak h}=\tfrac{2r_s}N];
\]

%dsp-level:0
\[
\qquad r_s^{(n)}\in[n\mathfrak h,\ 30n\mathfrak h],
\]

%dsp-level:0
\[
\qquad \frac{r_{s,\max}}{r_{s,\min}}=30
\]

%dsp-level:0
\[
\qquad [\text{machine}];
\]

%dsp-level:0
\[
\textbf{the comparison discipline:}\quad\neg\text{a collapse reading},
\]

%dsp-level:0
\[
\quad\neg\text{a horizon identification},
\]

%dsp-level:0
\[
\quad\neg\text{a measured radius}\ [\mathrm{BM.15}(4)\ \text{comparison discipline};\ \mathrm{P5};\ \mathrm{Def\text{-}BM.3}(6)];
\]

%dsp-level:0
\[
\textbf{what the row is:}\quad\text{a change of variable inside the corpus's own monomial, which exhibits the threshold as}\ \textbf{twice the comparison radius over the lapse};
\]

%dsp-level:0
\[
(5)\quad\textbf{the never-empty locus (no forbidden state):}
\]

%dsp-level:0
\[
\qquad \forall\mathcal R,\ \forall N\in(0,\infty):
\]

%dsp-level:0
\[
\qquad \mathcal T(\mathcal R;N)\in(0,\infty)
\]

%dsp-level:0
\[
\qquad [E_{\mathrm{int}}\in\mathbb Z_{\ge2},\ \mathfrak h>0,\ \ln2>0\ [\mathrm{BM.13}(3):\ \tfrac12<\ln2<1]]:
\]

%dsp-level:0
\[
\{E>0:\ NE\ge\mathcal T_0(\mathcal R)\}=\Bigl[\frac{\mathcal T_0(\mathcal R)}N,\ \infty\Bigr)\ne\varnothing
\]

%dsp-level:0
\[
\qquad [\text{machine: the half-line is nonempty at every admissible}\ N];
\]

%dsp-level:0
\[
\boxed{\ \neg\exists\ \textbf{a permanently forbidden resonance: every eigenvalue's threshold locus is a nonempty half-line in the local energy scale}\ }
\]

%dsp-level:0
\[
\textbf{the two named triples are neither formed nor excluded:}
\]

%dsp-level:0
\[
\qquad \neg\exists\ \text{a triple-selection rule in the corpus}\ [\mathrm{BM.103}(8):\ (\mathrm{R3})]:
\]

%dsp-level:0
\[
\textbf{but both are admissible threshold-wise:}\quad\text{any}\ \mathbf s\in\mathcal S^3\ \text{has}\ E_{\mathrm{int}}(\mathbf s)
\]

%dsp-level:0
\[
\in [6,180]\subset\mathbb Z_{>0}\ \text{and hence a nonempty locus}\ [(5)\text{'s first row};\ \text{machine: all}\ 216\ \text{ordered triples}];
\]

%dsp-level:0
\[
\boxed{\ \textbf{there are no forbidden compositions; there are compositions whose thresholds require a higher local scale}\ }
\]

%dsp-level:0
\[
(6)\quad\textbf{the accelerator and horizon regimes are evaluations, not mechanisms:}
\]

%dsp-level:0
\[
\qquad \textbf{a high local scale}\ E\gg E_1:\quad NE\ge\mathcal T_0(\mathcal R)\ \text{holds for every}\ \mathcal R\ \text{with}\ E_{\mathrm{int}}(\mathcal R)\le NE\ln2/\mathfrak h
\]

%dsp-level:0
\[
\qquad [\text{machine}]:
\]

%dsp-level:0
\[
\textbf{a small lapse}\ N\ll1\ \text{(the deep-gravity read):}\quad\mathcal T(\mathcal R;N)
\]

%dsp-level:0
\[
= \mathcal T_0/N\to\infty\implies\textbf{thresholds rise and fewer resonances materialize}
\]

%dsp-level:0
\[
\qquad [\text{machine: monotonicity from}\ (1)];
\]

%dsp-level:0
\[
\textbf{a large comparison radius}\ r_s:\quad\mathcal T=2r_s/N\ \text{grows linearly}\qquad[(4);\ \text{machine}];
\]

%dsp-level:0
\[
\boxed{\ \textbf{both regimes are the same evaluation of the same function at different local parameters}\ —\ \neg\text{two mechanisms,}\ \neg\text{a new object}\ }
\]

%dsp-level:0
\[
\neg\text{a cosmological coordinate is used:}\ \text{the rows are stated in}\ (E,\ N,\ G,\ r_s)\ \text{and}\ \neg\text{in any global time}\ [\mathrm{Def\text{-}BM.17}(6)(\mathrm i);\ \mathrm{A6};\ \mathrm{D1.8}];
\]

%dsp-level:0
\[
(7)\quad\textbf{the crossing strength re-cited and}\ \neg\textbf{re-evaluated:}
\]

%dsp-level:0
\[
\qquad \varepsilon_\ast(w)=\frac{\ln(4/3)}{3w\ln2}=\frac{2-\log_23}{3w},
\]

%dsp-level:0
\[
\qquad \tfrac18<\varepsilon_\ast<\tfrac5{12}
\]

%dsp-level:0
\[
\qquad [\mathrm{BM.101}(3);\ \text{machine: the enclosure re-certified on}\ 2^{13}>3^8,\ 3^{12}>2^{19}]:
\]

%dsp-level:0
\[
\textbf{its role here:}\quad\text{the crossing locus}\ \mathcal L_\ast\ \text{is where the mass splitting changes sign, and it is a level set of the same profile}\ \varepsilon_{\mathrm{eff}}\ \text{that the width function carries}\ [\mathrm{BM.101}(4);\ \mathrm{Def\text{-}BM.17}(4)]:
\]

%dsp-level:0
\[
\neg\text{a numerical coupling enters},
\]

%dsp-level:0
\[
\qquad \neg\text{a measured value is compared}
\]

%dsp-level:0
\[
\qquad [\mathrm{P5};\ \mathrm{EM.25}(7);\ \text{machine: the token scan}];
\]

%dsp-level:0
\[
(8)\quad\textbf{status:}\qquad\textbf{derived:}\ (1)\text{--}(6);\qquad\textbf{re-cited,}\ \neg\textbf{re-evaluated:}\ (7);
\]

%dsp-level:0
\[
\neg\text{a declaration},
\]

%dsp-level:0
\[
\quad\neg\text{a P4 addition},
\]

%dsp-level:0
\[
\quad\neg\text{a new dimensionful scale}\ [\mathrm{P1}],
\]

%dsp-level:0
\[
\quad\neg\text{a measured value}\ [\mathrm{P5}],
\]

%dsp-level:0
\[
\quad\neg\text{a species identification}\ [\mathrm{Def\text{-}BM.3}(7)];
\]

\textbf{Req.} D0.8, D1.8, D2.9 (Parts 1–2); A6, CT2.1, CT2.9, EM.25(7) (Part 1, Part 2); P1, P3, P4, P5 (Parts 1–2); FT.28, FT.44(4) (Part 3, Census); Census FT.35(2), FT.42(5) (row-wise); Def-BM.3(6), Def-BM.3(7), Def-BM.4(1), Def-BM.17(3), Def-BM.17(4), BM.13(2), BM.13(3), BM.15(4), BM.16(1), BM.88(3), BM.101(3), BM.101(4), BM.102(2), BM.103(8), BM.111(5) (this part); exact integer sum-set enumeration, exact rational monotonicity, Laurent-monomial exponent censuses, integer power certificates [math; machine: \texttt{checks/\allowbreak{}verify\_\allowbreak{}H\_\allowbreak{}antimatter\_\allowbreak{}resonance\_\allowbreak{}spectrum.py} §E].

\textbf{Pf.}

%dsp-level:0
\[
(1):\  \mathcal T_0=E_{\mathrm{int}}\mathfrak h/\ln2\ [\mathrm{Def\text{-}BM.17}(3);\ \text{machine}];
\]

%dsp-level:0
\[
\qquad [\mathcal T_0]=[\mathfrak h]=\mathfrak q^{-1/2}\ \text{and}\ [E_{\mathrm{int}}]=[\ln2]=\mathfrak q^0\implies[\mathcal T_0]=\mathfrak q^{-1/2}\cdot\mathfrak q^{0}
\]

%dsp-level:0
\[
\text{with}\ E_{\mathrm{int}}=\mathfrak h^{-1}E\ \text{dimensionless:}
\]

%dsp-level:0
\[
\quad [\mathcal T_0]=[\mathfrak h]\cdot[E]/[\mathfrak h]=\mathfrak q^{+1/2}\ [\mathrm{CT2.1};\ \mathrm{BM.102}(2);\ \text{machine: the exponent arithmetic}];
\]

%dsp-level:0
\[
\qquad N_1<N_2\implies\mathcal T_0/N_1>\mathcal T_0/N_2\ [\text{machine: exact rationals}]
\]

%dsp-level:0
\[
(2):\  \text{the three sum sets computed exhaustively from}\ \{2,6,12,24,60\}\ \text{with repetitions,}\ n
\]

%dsp-level:0
\[
= 1,
\]

%dsp-level:0
\[
2,\qquad 3:
\]

%dsp-level:0
\[
\quad 5,\ 15,\ 33\ \text{distinct values};
\]

%dsp-level:0
\[
\qquad \min=2n,\ \max=60n,\ 60n/2n=30\ [\text{machine: exact integers}]
\]

%dsp-level:0
\[
(3):\  G=2\pi E_{\mathrm{int}}^2\ln2\implies E_{\mathrm{int}}=\sqrt{G/(2\pi\ln2)}\ [\text{machine: the positive branch at}\ E_{\mathrm{int}}>0];
\]

%dsp-level:0
\[
\qquad \text{substitution into}\ (1)\ [\text{machine}]
\]

%dsp-level:0
\[
G_{\min}=2\pi\ln2\cdot(2n)^2=4n^2\cdot2\pi\ln2,
\]

%dsp-level:0
\[
\quad G_{\max}=2\pi\ln2\cdot(60n)^2=3600n^2\cdot2\pi\ln2,
\]

%dsp-level:0
\[
\quad 3600/4=900=30^2\ [\text{machine}]
\]

%dsp-level:0
\[
(4):\  r_s
\]

%dsp-level:0
\[
= \tfrac12\mathfrak h E_{\mathrm{int}}
\]

%dsp-level:0
\[
\implies E_{\mathrm{int}}
\]

%dsp-level:0
\[
= 2r_s/\mathfrak h\implies\mathcal T_0
\]

%dsp-level:0
\[
= \tfrac{\mathfrak h}{\ln2}\cdot\tfrac{2r_s}{\mathfrak h}
\]

%dsp-level:0
\[
= \tfrac{2r_s}{\ln2}\cdot\ln2\ [\text{machine}];
\]

%dsp-level:0
\[
\qquad \mathcal T=\mathcal T_0/N\ \text{with}\ \mathcal T_0=2r_s\ \text{in energy units}\ [\text{machine}]
\]

%dsp-level:0
\[
r_s\in[\tfrac12\mathfrak h\cdot2n,\ \tfrac12\mathfrak h\cdot60n]=[n\mathfrak h,\ 30n\mathfrak h]\ [\text{machine}]
\]

%dsp-level:0
\[
(5):\  E_{\mathrm{int}}\ge2\ [\mathrm{FT.28}]\implies\mathcal T_0\ge2\mathfrak h/\ln2>0\ [\tfrac12<\ln2<1\ [\mathrm{BM.13}(3)];\ \text{machine}];
\]

%dsp-level:0
\[
\qquad \text{the half-line}\ [\mathcal T_0/N,\infty)\ \text{is nonempty}\ \forall N>0\ [\text{machine}]
\]

%dsp-level:0
\[
\text{all}\ 216\ \text{ordered triples}: E_{\mathrm{int}}\in[6,180]\subset\mathbb Z_{>0}\ [\text{machine: exhaustive}];
\]

%dsp-level:0
\[
\qquad \neg\text{a selection rule exists}\ [\mathrm{BM.103}(8)]
\]

%dsp-level:0
\[
(6):\  NE\ge\mathcal T_0\iff E_{\mathrm{int}}\le NE\ln2/\mathfrak h\ [\text{machine}];
\]

%dsp-level:0
\[
\qquad N\downarrow\implies\mathcal T\uparrow\ [(1)];
\]

%dsp-level:0
\[
\qquad (7):\  \mathrm{BM.101}(3),
\]

%dsp-level:0
\[
(4)\ \text{re-cited, the enclosure re-certified on the two integer powers}\ [\text{machine}];
\]

%dsp-level:0
\[
\qquad (8):\  \text{the standing-rule scans}\ [\text{machine}].\ \blacksquare
\]

\textbf{Phys.} The threshold is one Laurent monomial and it has three exact readings. Read in the local scale it is the resonance's edge count times the single free scale over the log of two, divided by the lapse, so a deeper lapse lowers every threshold in the ladder at once and the whole spectrum moves as one function. Read in the local curvature it is a square root, because the calibration lock makes the curvature a derived quadratic function of the same count — and the lock is free of the scale, free of the cosmological term and free of every declared parameter, so the curvature reading adds no datum. Read in the comparison radius it is twice the radius over the lapse, and that row is quarantined exactly as the corpus quarantines every $G_N$-built quantity: no collapse reading, no horizon identification, no measured radius. The ladders themselves are exact integer sets — five values for one chain, fifteen for a meson, thirty-three for a baryon — running from twice the chain count to sixty times it, so the span ratio is exactly thirty at every level and the curvature span is exactly nine hundred. Then the row that answers the question about forbidden states: every threshold is a finite positive number at every positive lapse, so every locus is a nonempty half-line, and all two hundred sixteen ordered constituent triples sit inside it. Nothing is excluded; what differs is where in the local scale the half-line starts. Both regimes the question names — a collider and a deep-gravity region — are the same evaluation of the same function at different local parameters, one raising the scale and one lowering the lapse, and neither needs a mechanism, a cosmological clock, or a new object.

\end{theorem}

\begin{theorem}[{\texorpdfstring{BM.116 — The decay width: the exact function of the local parameters, its three exact zeros, and the lifetime floor}{BM.116 - The decay width: the exact function of the local parameters, its three exact zeros, and the lifetime floor}}]\label{tag:BM-116}
\textbf{Thm.}

%dsp-level:0
\[
(1)\quad\textbf{the width function:}
\]

%dsp-level:0
\[
\qquad \boxed{\ \Gamma(\mathcal R;\,E,N,G,\varepsilon_{\mathrm{EM}})\ =\ \underbrace{\frac{x_+x_-}{1-x_+x_-}}_{\Gamma_{\mathrm{tr}}(\varepsilon_q)}\ \cdot\ \underbrace{\Bigl(1-\frac{\Delta x_{\mathrm B}(E)}{r}\Bigr)}_{f}\ \cdot\ \underbrace{N\,\mathsf g(G)}_{\varepsilon_{\mathrm{eff}}/\varepsilon_{\mathrm{EM}}\,\cdot\,f^{-1}\cdot f}\ \cdot\ \underbrace{\mathcal N(\mathcal R)}_{\text{the channel count}}\ \cdot\ \underbrace{\chi_{\{NE\ge\mathcal T_0(\mathcal R)\}}}_{\text{materialization}}\ }
\]

%dsp-level:0
\[
\textbf{assembled from four frozen rows and one count:}
\]

%dsp-level:0
\[
\qquad \Gamma_{\mathrm{tr}}=\tfrac{x_+x_-}{1-x_+x_-}\in[0,\tfrac13]\ [\mathrm{BM.44}(2)];
\]

%dsp-level:0
\[
\qquad f=1-\tfrac1k,\ k=\tfrac{r}{\Delta x_{\mathrm B}(E)}\ge1\ [\mathrm{BM.88}(2)(\mathrm{ii})];
\]

%dsp-level:0
\[
\frac{\varepsilon_{\mathrm{eff}}}{\varepsilon_{\mathrm{EM}}}=Nf\,\mathsf g(G)\ [\mathrm{BM.88}(4);\ \mathrm{BM.101}(2)];
\]

%dsp-level:0
\[
\qquad \mathcal N(\mathcal R)=\#\{\mathcal R'\ \text{in the same}\ (n,\sigma,Q)\ \text{class}:\ \lambda^{(n)}(\mathcal R')<\lambda^{(n)}(\mathcal R)\}\ [\mathrm{Def\text{-}BM.17}(4)];
\]

%dsp-level:0
\[
[\Gamma]=\mathfrak q^0,\qquad[\tau_{\mathrm{life}}]=\mathfrak q^0\ \ \text{ticks}\qquad[\mathrm{BM.44}(4);\ \mathrm{D1.8};\ \mathrm{BM.102}(2);\ \text{machine: the exponent census}];
\]

%dsp-level:0
\[
(2)\quad\textbf{the flicker factor and its exact range:}
\]

%dsp-level:0
\[
\qquad x_+x_-=4\varepsilon_q(1-4\varepsilon_q)=\frac14-16\Bigl(\varepsilon_q-\frac18\Bigr)^{\!2}\le\frac14,
\]

%dsp-level:0
\[
\qquad \text{equality}\iff\varepsilon_q=\frac18
\]

%dsp-level:0
\[
\qquad [\mathrm{BM.19}(2);\ \mathrm{BM.44}(1);\ \text{machine: the SOS completion and the}\ 256\text{-point grid}]:
\]

%dsp-level:0
\[
\boxed{\ \Gamma_{\mathrm{tr}}(0)=0,\qquad\Gamma_{\mathrm{tr}}\bigl(\tfrac18\bigr)=\frac{1/4}{3/4}=\frac13,\qquad\Gamma_{\mathrm{tr}}\in\Bigl[0,\tfrac13\Bigr]\ \ \text{on}\ \ \varepsilon_q\in\Bigl[0,\tfrac18\Bigr],\qquad\Gamma_{\mathrm{tr}}\ \text{strictly increasing there}\ }
\]

%dsp-level:0
\[
[\mathrm{BM.44}(2);\ \text{machine:}\ \Gamma_{\mathrm{tr}}(\varepsilon_2)-\Gamma_{\mathrm{tr}}(\varepsilon_1)>0\ \text{via}\ 4(\varepsilon_2-\varepsilon_1)(1-4(\varepsilon_1+\varepsilon_2))>0];
\]

%dsp-level:0
\[
(3)\quad\textbf{the three exact zeros (the derived ``stable'' loci):}
\]

%dsp-level:0
\[
\qquad \boxed{\ (\mathrm Z1)\quad\varepsilon_q=0\implies\Gamma\equiv0\ \ \text{for every resonance at every local parameter}\ }
\]

%dsp-level:0
\[
\qquad [\mathrm{BM.23}(8);\ \mathrm{BM.19}(3);\ \text{machine}]:
\]

%dsp-level:0
\[
\textbf{the mechanism:}
\]

%dsp-level:0
\[
\quad x_-
\]

%dsp-level:0
\[
= 4\varepsilon_q
\]

%dsp-level:0
\[
= 0\implies\neg\exists\ \text{a dissolution move}\implies\text{the ensemble collapses to its crystallized member}\implies\neg\text{a transition weight beyond the identity};
\]

%dsp-level:0
\[
(\mathrm Z2)\quad r=\Delta x_{\mathrm B}(E)\implies k=1\implies f=0\implies\Gamma=0
\]

%dsp-level:0
\[
\qquad [\mathrm{BM.88}(5);\ \text{machine}]:
\]

%dsp-level:0
\[
\qquad \textbf{at the member's own Born envelope the width vanishes};
\]

%dsp-level:0
\[
(\mathrm Z3)\quad\mathcal N(\mathcal R)=0\implies\Gamma=0
\]

%dsp-level:0
\[
\qquad [\mathrm{Def\text{-}BM.17}(5);\ \text{machine}]:
\]

%dsp-level:0
\[
\qquad \boxed{\ \textbf{the lightest resonance of every}\ (n,\sigma,Q)\ \textbf{class has width exactly zero, at every local parameter}\ }
\]

%dsp-level:0
\[
\textbf{the count of zero-width resonances at the three formed levels:}\qquad\begin{array}{c|c|c}\text{level} & \#\{\mathcal N=0\} & \text{locus}\\\hline n=2,\ \sigma=\delta & 3 & \text{one per charge class}\ Q\in\{0,\pm1\}\ [\mathrm{BM.110}(3)]\\n=3,\ \sigma=\varepsilon & 7 & \text{one per charge class}\ Q\in\{-3,\dots,+3\}\ [\mathrm{BM.114}(1)]\end{array}\qquad[\text{machine: exact}];
\]

%dsp-level:0
\[
(4)\quad\textbf{the fourth zero is the threshold's, and it is not a width:}
\]

%dsp-level:0
\[
\qquad NE<\mathcal T_0(\mathcal R)\implies\chi=0\implies\Gamma=0
\]

%dsp-level:0
\[
\qquad [\text{machine}]:
\]

%dsp-level:0
\[
\boxed{\ \textbf{an unmaterialized resonance has no width because it has no ensemble}\ —\ \text{the zero is the indicator's,}\ \neg\text{a decay}\ }
\]

%dsp-level:0
\[
\qquad [\mathrm{Def\text{-}BM.17}(3);\ \mathrm{BM.94}(5):\ \text{the materialization restriction}];
\]

%dsp-level:0
\[
\textbf{the distinction the tier keeps:}\quad\Gamma=0\ \text{from}\ (\mathrm Z1),(Z2),(Z3)\ \textbf{is a stable resonance};
\]

%dsp-level:0
\[
\qquad \chi=0\ \textbf{is an absent one};
\]

%dsp-level:0
\[
\qquad \neg\text{a binary category covers both}\ [\mathrm{Def\text{-}BM.17}(6)(\mathrm{ii})];
\]

%dsp-level:0
\[
(5)\quad\textbf{the bound and the lifetime floor:}\qquad f<1\ \forall k<\infty\ [\mathrm{BM.88}(5);\ \text{machine}]\implies
\]

%dsp-level:0
\[
\boxed{\ 0\ \le\ \Gamma(\mathcal R)\ \le\ \frac13\cdot N\,\mathsf g(G)\cdot\mathcal N(\mathcal R)\ }\qquad[\text{machine: the three factor bounds composed}];
\]

%dsp-level:0
\[
\mathcal N=1,\ f\to1,\ N\mathsf g\to1:
\]

%dsp-level:0
\[
\qquad \Gamma\to\Gamma_{\mathrm{tr}}\le\tfrac13\implies\boxed{\ \tau_{\mathrm{life}}:=\frac1\Gamma\ \ge\ 3\ \ \textbf{ticks}\ }
\]

%dsp-level:0
\[
\qquad [\mathrm{BM.44}(3);\ \text{machine: exact reciprocal}];
\]

%dsp-level:0
\[
\varepsilon_q=0:\qquad\Gamma=0\implies\tau_{\mathrm{life}}\ \textbf{has no finite value}\qquad[\mathrm{BM.44}(3);\ \text{machine}];
\]

%dsp-level:0
\[
\neg\text{a lifetime in external seconds is derived and}\ \neg\text{a lifetime is compared}
\]

%dsp-level:0
\[
\qquad [\mathrm{BM.44}(5);\ \mathrm{P5};\ \text{machine: zero external unit tokens}];
\]

%dsp-level:0
\[
(6)\quad\textbf{the energy dependence (the width grows with the local scale):}
\]

%dsp-level:0
\[
\qquad \text{at fixed}\ r,
\]

%dsp-level:0
\[
\qquad k(E)=\frac{r\,E_{\mathrm{int}}(E)\ln2}{\mathfrak h},
\]

%dsp-level:0
\[
\qquad f=1-\frac{\mathfrak h}{r\,E_{\mathrm{int}}(E)\ln2}
\]

%dsp-level:0
\[
\qquad [\mathrm{BM.88}(6);\ \text{machine}]:
\]

%dsp-level:0
\[
\boxed{\ E\uparrow\implies k\uparrow\implies f\uparrow\implies\Gamma\uparrow\ }
\]

%dsp-level:0
\[
\qquad [\mathrm{BM.88}(6)\text{'s three monotonicity rows re-cited};\ \text{machine: positive coefficients}];
\]

%dsp-level:0
\[
\textbf{the calibration-scale reading:}
\]

%dsp-level:0
\[
\qquad \text{``low energy''}\ :=\ E\ \text{near}\ E_1=\mathfrak h/(r\ln2),\ \text{where}\ f(1)=0;
\]

%dsp-level:0
\[
\qquad \text{``high energy''}\ :=\ E\gg E_1,\ \text{where}\ f\to1
\]

%dsp-level:0
\[
\qquad [\mathrm{BM.101}(5);\ \mathrm{BM.88}(2)(\mathrm i);\ \text{machine}]:
\]

%dsp-level:0
\[
\boxed{\ \textbf{at the low calibration point every width vanishes through}\ f\ \textbf{and through}\ \chi\ \textbf{for the unmaterialized ones; at the high one the widths open in the order of}\ \mathcal N\ \textbf{and}\ \mathcal T_0\ }
\]

%dsp-level:0
\[
\textbf{and both are evaluations of one function,}\ \neg\textbf{two regimes of two theories}
\]

%dsp-level:0
\[
\qquad [\mathrm{Def\text{-}BM.17}(6)(\mathrm{ii});\ \text{machine}];
\]

%dsp-level:0
\[
(7)\quad\textbf{the two-sided bar and the width:}
\]

%dsp-level:0
\[
\qquad N\mathsf g\,\varepsilon_{\mathrm{EM}}\le\tfrac18\implies\Delta m>0;
\]

%dsp-level:0
\[
\qquad \ge\tfrac5{12}\implies\Delta m<0
\]

%dsp-level:0
\[
\qquad [\mathrm{BM.99}(3);\ \mathrm{BM.101}(3):\ \text{the same enclosure}]:
\]

%dsp-level:0
\[
\textbf{the consequence for the width:}\quad\frac{\varepsilon_{\mathrm{eff}}}{\varepsilon_{\mathrm{EM}}}
\]

%dsp-level:0
\[
= Nf\,
\]

%dsp-level:0
\[
\mathsf g
\]

%dsp-level:0
\[
\le N\mathsf g\implies\boxed{\ \text{the bar's band}\ \bigl[\tfrac18,\tfrac5{12}\bigr]\ \text{bounds the width's coupling factor exactly}\ }
\]

%dsp-level:0
\[
\textbf{and the sign of the mass splitting flips inside it}\ —\ \text{one function, two readings}\ [\mathrm{BM.101}(1),(2);\ \text{machine}];
\]

%dsp-level:0
\[
(8)\quad\textbf{status:}\qquad\textbf{derived:}\ (1)\text{--}(7);
\]

%dsp-level:0
\[
\textbf{non-derived and named:}\ \ (1)\text{'s product form}\ [\mathrm{Def\text{-}BM.17}(4)\ \text{non-derived clause}];
\]

%dsp-level:0
\[
\qquad \textbf{priced,}\ \neg\textbf{adopted:}\ \ \mathsf g\in\{1,1+G,\tfrac1{1+G}\}\ [\mathrm{BM.88}(4);\ \mathrm{BM.104}(5)];
\]

%dsp-level:0
\[
\neg\text{a declaration},
\]

%dsp-level:0
\[
\quad\neg\text{a P4 addition},
\]

%dsp-level:0
\[
\quad\neg\text{a new dimensionful scale}\ [\mathrm{P1}],
\]

%dsp-level:0
\[
\quad\neg\text{a measured value}\ [\mathrm{P5}],
\]

%dsp-level:0
\[
\quad\neg\text{a numerical coupling}\ [\mathrm{EM.25}(7)],
\]

%dsp-level:0
\[
\quad\neg\text{a species identification}\ [\mathrm{Def\text{-}BM.3}(7)];
\]

\textbf{Req.} D0.8, D1.8, D2.9 (Parts 1–2); CT0.8, EM.25(7) (Part 1, Part 2); P1, P4, P5 (Parts 1–2); Def-H.20, FT.44(4) (Part 3); Def-BM.3(7), Def-BM.4(1), Def-BM.17(3), Def-BM.17(4), Def-BM.17(5), Def-BM.17(6), BM.19(2), BM.19(3), BM.23(8), BM.44(1), BM.44(2), BM.44(3), BM.44(4), BM.44(5), BM.88(2), BM.88(4), BM.88(5), BM.88(6), BM.94(5), BM.99(3), BM.101(1), BM.101(2), BM.101(3), BM.101(5), BM.102(2), BM.104(5), BM.110(3), BM.114(1), BM.115(1) (this part); exact rational resolvent identities, the SOS parabola completion, exhaustive rational grid censuses, exact monotonicity by polynomial sign, Laurent-monomial exponent censuses [math; machine: \texttt{checks/\allowbreak{}verify\_\allowbreak{}H\_\allowbreak{}antimatter\_\allowbreak{}resonance\_\allowbreak{}spectrum.py} §E].

\textbf{Pf.}

%dsp-level:0
\[
(1):\  \text{the four frozen rows substituted into}\ \mathrm{Def\text{-}BM.17}(4)\text{'s product}\ [\text{machine: the monomial assembly}];
\]

%dsp-level:0
\[
\qquad \frac{\varepsilon_{\mathrm{eff}}}{\varepsilon_{\mathrm{EM}}}=\frac{Nf\,\mathsf g\,\varepsilon_{\mathrm{EM}}}{\varepsilon_{\mathrm{EM}}}=Nf\,\mathsf g\ [\text{machine}];
\]

%dsp-level:0
\[
\qquad [\Gamma]=\mathfrak q^0\ [\text{every factor dimensionless};\ \text{machine}]
\]

%dsp-level:0
\[
(2):\  \mathrm{BM.44}(1),
\]

%dsp-level:0
\[
(2)\text{'s rows re-cited}\ [\text{machine:}\ 4\varepsilon_q(1-4\varepsilon_q)=\tfrac14-16(\varepsilon_q-\tfrac18)^2\ \text{by expansion};\ \text{the}\ 256\text{-point grid census with equality only at}\ k=32]
\]

%dsp-level:0
\[
\Gamma_{\mathrm{tr}}(\tfrac18)=\tfrac{1/4}{1-1/4}=\tfrac{1/4}{3/4}=\tfrac13\ [\text{machine: exact rationals}];
\]

%dsp-level:0
\[
\qquad \text{monotonicity via}\ 4(\varepsilon_2-\varepsilon_1)(1-4(\varepsilon_1+\varepsilon_2))>0\ \text{on}\ \varepsilon_1<\varepsilon_2\le\tfrac18\ [\text{machine: the polynomial identity and a rational witness grid}]
\]

%dsp-level:0
\[
(3):\  (\mathrm Z1)\ x_-=4\cdot0=0\implies x_+x_-=0\implies\Gamma_{\mathrm{tr}}=0/(1-0)=0\ [\text{machine}];
\]

%dsp-level:0
\[
\qquad (\mathrm Z2)\ k=1\implies f=1-1=0\ [\mathrm{BM.88}(5);\ \text{machine}];
\]

%dsp-level:0
\[
\qquad (\mathrm Z3)\ \mathcal N=0\implies\text{the product is}\ 0\ [\text{machine}]
\]

%dsp-level:0
\[
\#\{\mathcal N=0\}:\ \text{at}\ n
\]

%dsp-level:0
\[
= 2\ \text{the three classes}\ Q\in\{0,\pm1\}\ \text{each have a lightest member}\ (Q=\pm2\ \text{has one member and is its own lightest, counted once each});
\]

%dsp-level:0
\[
\ \text{the exact count over the}\ \lambda^{(2)}\ \text{ladder of}\ \mathrm{BM.114}(3)\ \text{is}\ 3\ \text{on the sealed classes}\ [\text{machine}]
\]

%dsp-level:0
\[
\text{at}\ n=3\ \text{the seven classes}\ Q=-3,\dots,+3\ \text{each contribute one}\ [\text{machine}]
\]

%dsp-level:0
\[
(4):\  \chi_{\{NE<\mathcal T_0\}}=0\implies\Gamma=0\ [\text{machine}];
\]

%dsp-level:0
\[
\qquad \mathrm{BM.94}(5)\text{'s materialization restriction re-cited}:\ E\ge2\ \text{supports no closed chain below it}
\]

%dsp-level:0
\[
(5):\  f=1-\tfrac1k<1\iff\tfrac1k>0\iff k<\infty\ [\mathrm{BM.88}(5);\ \text{machine}];
\]

%dsp-level:0
\[
\qquad \Gamma_{\mathrm{tr}}\le\tfrac13\ [(2)];
\]

%dsp-level:0
\[
\qquad \text{the product bound}\ [\text{machine}]
\]

%dsp-level:0
\[
\mathcal N=1,\ f\to1,\ N\mathsf g\to1:\ \Gamma\to\tfrac13\implies\tau_{\mathrm{life}}\to3\ [\mathrm{BM.44}(3);\ \text{machine: the exact reciprocal}];
\]

%dsp-level:0
\[
\qquad \Gamma=0\implies1/\Gamma\ \text{has no finite value}\ [\text{machine}]
\]

%dsp-level:0
\[
(6):\  \mathrm{BM.88}(6)\text{'s three monotonicity rows re-cited};
\]

%dsp-level:0
\[
\qquad f\ \text{is a factor of}\ \Gamma\ \text{with all other factors}\
\]

%dsp-level:0
\[
\ge 0\implies\Gamma\ \text{increases with}\ f\ [\text{machine: positive-coefficient composition}]
\]

%dsp-level:0
\[
(7):\  \mathrm{BM.99}(3),
\]

%dsp-level:0
\[
\ \mathrm{BM.101}(3)\ \text{re-cited, the enclosure re-certified on}\ 2^{13}>3^8\ \text{and}\ 3^{12}>2^{19}\ [\text{machine}];
\]

%dsp-level:0
\[
\qquad f\le1\implies Nf\,\mathsf g\le N\mathsf g\ [\text{machine}];
\]

%dsp-level:0
\[
\qquad (8):\  \text{the standing-rule scans}\ [\text{machine}].\ \blacksquare
\]

\textbf{Phys.} The width is a product of four things the corpus had already frozen and one thing it can count. The flicker factor is the transition content of the Janus resolvent at unit fugacity, a rational function of the single crystallization threshold that climbs from zero to exactly a third at the threshold's midpoint and is strictly increasing all the way. The envelope factor is one minus one over k, which vanishes at the member's own Born width. The coupling factor is the effective coupling over the bare one, the lapse times the envelope fraction times the priced curvature composition. The count is how many strictly lighter resonances sit in the same chain-count, seal and charge class. And the indicator says whether the resonance is materialized at all. Three of those give exact zeros with distinct mechanisms: no flicker means no dissolution move and hence no transition weight; sitting at your own envelope means the enclosed flux fraction is zero; and being the lightest in your class means there is nothing to go to, so the count is zero and the width is zero by arithmetic. The fourth zero is different in kind and the tier keeps it apart — an unmaterialized resonance has no width because it has no ensemble, which is absence rather than stability. The bound follows by composing the factor bounds, and the lifetime floor is three ticks at the extreme, exactly as the vertex tier derived, with no second anywhere in it. Energy enters monotonically: raising the local scale raises the envelope ratio, raises the enclosed fraction, and raises every width, so "low energy" and "high energy" are two points on one function where the widths are respectively all zero and open in the order of their thresholds and their counts. The two-sided bar on the mass splitting is the same product seen from the other side, and its band is exactly the bound on the width's coupling factor.

\end{theorem}

\begin{theorem}[{\texorpdfstring{BM.117 — The groupings, the ceilings, and the octet adjudication}{BM.117 - The groupings, the ceilings, and the octet adjudication}}]\label{tag:BM-117}
\textbf{Thm.}

%dsp-level:0
\[
(1)\quad\textbf{the groupings that are derived:}\qquad\begin{array}{l|l|l}\text{grouping} & \text{the exact value} & \text{locus}\\\hline\text{the charge classes at}\ n & 2n+1\ \text{classes with the binomial weights}\ \binom{2n}{Q+n} & [\mathrm{BM.114}(1)]\\\text{the sealing grammar} & g(N)=\sum_{2a+3b=N}\frac{N!}{a!\,b!\,2^a6^b}=(0,1,1,3,10,25,105)\ \text{for}\ N=1..7 & [\mathrm{CT3.17a}(2);\ \mathrm{Def\text{-}BM.9}]\\\text{the multiset orbits at}\ n & \binom{6+n-1}n=(6,\ 21,\ 56) & [\mathrm{BM.114}(4)]\\\text{the all-distinct orbits} & \binom6n=(6,\ 15,\ 20)\ \ \text{with multiplicity}\ n! & [\mathrm{BM.114}(4);\ \text{machine}]\\\text{the generation floors} & 3,\ \text{ordered exactly by}\ 2E=12<24<60 & [\mathrm{Def\text{-}BM.14}(1),(6);\ \mathrm{BM.60}(1)]\\\text{the branches} & 2\ \text{per floor, ordered exactly by}\ D_S(\mathcal F_g,-)<D_S(\mathcal F_g,+) & [\mathrm{Def\text{-}BM.14}(4)]\\\text{the flavour symmetry} & C_3,\ \text{its orbits the three channel lines} & [\mathrm{FT.111}(2);\ \mathrm{FT.158}(1);\ \mathrm{FT.142}(5)]\\\text{the defect field} & \mathbb M,\ [\mathbb M:\mathbb Q]=16,\ \text{with the subfield ladder}\ 1,2,4,8,16 & [\mathrm{BM.114}(3),(6)]\end{array}
\]

%dsp-level:0
\[
(2)\quad\textbf{the grouping the spectrum actually carries:}
\]

%dsp-level:0
\[
\qquad \textbf{the three-level grouping of the}\ n\textbf{-chain ladder is}\ (\text{floor},\ \text{branch},\ \text{charge class})\ \text{with the orbit multiplicities of}\ \mathrm{BM.114}(4):
\]

%dsp-level:0
\[
\boxed{\ \textbf{at}\ n=3:\quad56\ \text{orbits}\ =\ 6\ \text{of size}\ 1\ +\ 30\ \text{of size}\ 3\ +\ 20\ \text{of size}\ 6,\qquad55\ \text{eigenvalues after the one collision}\ }
\]

%dsp-level:0
\[
\textbf{the}\ 20\ \textbf{all-distinct orbits are the}\ \varepsilon\textbf{-seal's fully antisymmetric slot pattern realized on distinct states:}
\]

%dsp-level:0
\[
\qquad \varepsilon_{ijk}\ \text{seals three fundamentals antisymmetrically}\ [\mathrm{D3.9};\ \mathrm{CT3.15}],
\]

%dsp-level:0
\[
\ \text{and}\ S_3\ \text{acts freely exactly on the all-distinct tuples}\ [\text{machine}]:
\]

%dsp-level:0
\[
\boxed{\ 20=\binom63\ \ \textbf{is derived as an orbit count and is}\ \neg\textbf{an}\ \mathsf{SU}(3)\ \textbf{decuplet dimension}\ }
\]

%dsp-level:0
\[
\qquad [\neg\exists\ \text{a flavour}\ \mathsf{SU}(3)\ \text{in the inventory}\ [\mathrm{BM.103}(7)];\ \text{machine: the symbol census}];
\]

%dsp-level:0
\[
(3)\quad\textbf{the octet adjudication (re-certified, and what replaces it):}
\]

%dsp-level:0
\[
\qquad \textbf{the derived}\ 3\ \textbf{is the colour seal's fundamental}\ [\mathrm{CT3.15};\ \mathrm{CT3.17};\ \mathrm{D3.9}];
\]

%dsp-level:0
\[
\qquad \textbf{the derived}\ 8\ \textbf{is a mask-capacity ratio, barred as a coupling}\ [\mathrm{BM.36}(2),(3);\ \mathrm{BM.61}(8)(\mathrm{O1})];
\]

%dsp-level:0
\[
\textbf{the derived}\ 10\ \textbf{is the grammar count}\ g(5)=10\ [\mathrm{CT3.17a}(2);\ \mathrm{Def\text{-}BM.9};\ \mathrm{BM.31}];
\]

%dsp-level:0
\[
\qquad \textbf{the derived}\ 15\ \textbf{and}\ 20\ \textbf{are the multiset orbit counts of}\ (1);
\]

%dsp-level:0
\[
\boxed{\ \textbf{the baryon octet and the meson octet are}\ \neg\textbf{derived, and the groupings that are derived are the binomial charge classes, the multiset orbits, the}\ C_3\ \textbf{triplet and the grammar counts}\ }
\]

%dsp-level:0
\[
[\mathrm{BM.103}(7)\ \text{re-cited and extended with the resonance census};\ \text{machine:}\ 8\notin\{2n+1,\ \binom{6+n-1}n,\ \binom6n,\ g(N)\}\ \text{at}\ n\le3,\ N\le7];
\]

%dsp-level:0
\[
(4)\quad\textbf{the ceilings:}
\]

%dsp-level:0
\[
\qquad \textbf{the per-constituent pole ceiling}\quad D_S(s)\le4\implies m^2_{\mathrm{pole}}(s)\le4\hat\mu'^2
\]

%dsp-level:0
\[
\qquad [\mathrm{FT.76}(3);\ \mathrm{Def\text{-}BM.14}(3):\ \max D_S=4\ \text{at}\ s_2;\ \text{machine}]:
\]

%dsp-level:0
\[
\textbf{the}\ n\textbf{-chain resonance ceiling is its}\ n\textbf{-fold sum:}
\]

%dsp-level:0
\[
\qquad \boxed{\ \lambda^{(n)}_{\max}=4n,\qquad m^2_{\mathrm{pole}}(\mathcal R)\le4n\,\hat\mu'^2,\qquad\text{attained exactly by}\ \mathbf s=(s_2,\dots,s_2)\ }
\]

%dsp-level:0
\[
[\text{machine:}\ 4n\ \text{is in the ladder at every}\ n=1,2,3,\ \text{with multiplicity}\ 1];
\]

%dsp-level:0
\[
\textbf{the locus discipline:}
\]

%dsp-level:0
\[
\qquad \mathrm{FT.76}(3)\text{'s ceiling is a property of a single chain's pole-side read and}\ \mathrm{FT.79}(2)\text{'s of a single chain's colour-side read}\ [\mathrm{BM.103}(6);\ \mathrm{FT.84}(2):\ \text{the two-ceiling synthesis}]:
\]

%dsp-level:0
\[
\boxed{\ \textbf{an}\ n\textbf{-chain composite eigenvalue is a sum of}\ n\ \textbf{single-chain reads and is}\ \neg\textbf{a single-chain read}\ —\ \neg\text{a ceiling is violated and}\ \neg\text{a ceiling is extended to the composite as such}\ }
\]

%dsp-level:0
\[
\textbf{the informational ceiling:}
\]

%dsp-level:0
\[
\qquad E_{\mathrm{int}}^{(n)}\in[12n,\ 60n]\ \text{on the six-state ladder}\implies m_{\mathrm{hid}}^{(n)}\in\Bigl[\frac{12n\ln2}{\mathfrak h},\ \frac{60n\ln2}{\mathfrak h}\Bigr]
\]

%dsp-level:0
\[
\qquad [\mathrm{BM.115}(2)(\mathrm{ii});\ \text{machine}];
\]

%dsp-level:0
\[
\textbf{the two currencies stay incommensurate:}
\]

%dsp-level:0
\[
\qquad \neg\exists\ \text{a derived ratio between}\ \lambda^{(n)}\ \text{and}\ E_{\mathrm{int}}^{(n)}\ —\ \text{the single residue is}\ \rho_{\mathrm B}
\]

%dsp-level:0
\[
\qquad [\mathrm{BM.82}(1),(5),(7);\ \mathrm{BM.69}(4);\ \mathrm{A\text{-}BM.9}];
\]

%dsp-level:0
\[
(5)\quad\textbf{the deficit-12 row is untouched:}
\]

%dsp-level:0
\[
\qquad \text{the two mass ceilings stand as properties of the reads they were proved for}\ [\mathrm{FT.76}(3);\ \mathrm{FT.79}(2);\ \mathrm{FT.97}(5):\ \text{the three-ceiling synthesis}]:
\]

%dsp-level:0
\[
\textbf{the resonance ladder adds no fourth ceiling and lifts none:}
\]

%dsp-level:0
\[
\qquad \lambda^{(n)}\ \text{is a sum of ceiling-bounded terms and its own bound}\ 4n\ \text{is the sum's}\ [(4)];
\]

%dsp-level:0
\[
\neg\text{a measured mass, hierarchy or range is compared}
\]

%dsp-level:0
\[
\qquad [\mathrm{P5};\ \mathrm{FT.84}(5\mathrm b);\ \text{machine: zero external tokens}];
\]

%dsp-level:0
\[
(6)\quad\textbf{the no-forbidden-state row, at the grouping level:}
\]

%dsp-level:0
\[
\qquad \textbf{every one of the}\ 6^n\ \textbf{ordered constituent tuples is a formed resonance with a nonempty threshold locus:}
\]

%dsp-level:0
\[
\#\{\text{ordered tuples at}\ n\}=6^n=(6,\ 36,\ 216),
\]

%dsp-level:0
\[
\qquad \#\{\text{distinct eigenvalues}\}=(6,\ 21,\ 55),
\]

%dsp-level:0
\[
\qquad \#\{\text{tuples with an empty locus}\}=0
\]

%dsp-level:0
\[
\qquad [\mathrm{BM.115}(5);\ \text{machine: exhaustive}]:
\]

%dsp-level:0
\[
\boxed{\ \textbf{there is no excluded composition at any level}\ —\ \text{the two named heaviest-floor triples are unread,}\ \neg\text{unreachable}\ }
\]

%dsp-level:0
\[
\qquad [\mathrm{BM.103}(8)];
\]

%dsp-level:0
\[
(7)\quad\textbf{status:}\qquad\textbf{derived:}\ (1),\ (2),\ (4),\ (6);\qquad\textbf{re-certified and extended:}\ (3);\qquad\textbf{untouched:}\ (5);
\]

%dsp-level:0
\[
\neg\text{a declaration},
\]

%dsp-level:0
\[
\quad\neg\text{a P4 addition},
\]

%dsp-level:0
\[
\quad\neg\text{a new dimensionful scale}\ [\mathrm{P1}],
\]

%dsp-level:0
\[
\quad\neg\text{a measured value}\ [\mathrm{P5}],
\]

%dsp-level:0
\[
\quad\neg\text{a currency conversion}\ [\mathrm{BM.83}(3)],
\]

%dsp-level:0
\[
\quad\neg\text{a species identification}\ [\mathrm{Def\text{-}BM.3}(7)];
\]

%dsp-level:0
\[
\alpha_B:\quad\neg\text{added},\quad\neg\text{derived},\quad\text{barred on}\ \mathrm{BM.61}(8)\text{'s four grounds}\ [\text{machine: the token scan of this tier}];
\]

\textbf{Req.} D2.9, D3.9 (Parts 1–2); CT3.15, CT3.17, CT3.17a(2) (Part 2); P1, P4, P5 (Parts 1–2); FT.76(3), FT.79(2), FT.84(2), FT.84(5), FT.97(5), FT.111(2), FT.142(5), FT.158(1) (Part 3); Def-BM.3(7), Def-BM.9, Def-BM.14(1), Def-BM.14(3), Def-BM.14(4), Def-BM.14(6), BM.31, BM.36(2), BM.36(3), BM.60(1), BM.61(8), BM.69(4), BM.82(1), BM.82(5), BM.82(7), BM.83(3), BM.103(7), BM.103(8), BM.114(1), BM.114(3), BM.114(4), BM.114(6), BM.115(2), BM.115(5), A-BM.9 (this part); exact integer orbit and grammar enumeration, exhaustive tuple censuses, exact radical comparison in $\mathbb M$ [math; machine: \texttt{checks/\allowbreak{}verify\_\allowbreak{}H\_\allowbreak{}antimatter\_\allowbreak{}resonance\_\allowbreak{}spectrum.py} §F].

\textbf{Pf.}

%dsp-level:0
\[
(1):\  \text{the eight rows re-cited from their own tags and their exact values re-computed}\ [\text{machine}];
\]

%dsp-level:0
\[
\qquad g(N)\ \text{for}\ N=1..7\ \text{by the labelled set-partition formula}\ [\mathrm{CT3.17a}(2);\ \text{machine}]
\]

%dsp-level:0
\[
(2):\  6+30+20=56,
\]

%dsp-level:0
\[
\qquad 6\cdot1+30\cdot3+20\cdot6=6+90+120=216=6^3\ [\text{machine}];
\]

%dsp-level:0
\[
\qquad S_3\ \text{acts freely on the all-distinct tuples}: |\operatorname{Stab}|
\]

%dsp-level:0
\[
= 1
\]

%dsp-level:0
\[
\implies |\text{orbit}|
\]

%dsp-level:0
\[
= 3!
\]

%dsp-level:0
\[
= 6\ \text{and}\ \#\text{orbits}
\]

%dsp-level:0
\[
= \tfrac{6\cdot5\cdot4}{6}
\]

%dsp-level:0
\[
= 20\ [\text{machine}]
\]

%dsp-level:0
\[
\varepsilon_{ijk}\ \text{antisymmetric in three slots}\ [\mathrm{D3.9};\ \mathrm{CT3.15}];
\]

%dsp-level:0
\[
\qquad \neg\exists\ \text{a flavour}\ \mathsf{SU}(3)\ [\mathrm{BM.103}(7);\ \text{machine: the symbol census}]\implies20\ \text{is an orbit count}
\]

%dsp-level:0
\[
(3):\  \mathrm{BM.103}(7)\text{'s four rows re-cited};
\]

%dsp-level:0
\[
\qquad 8\notin\{2n+1\}_{n\le3}\cup\{\binom{6+n-1}n\}_{n\le3}\cup\{\binom6n\}_{n\le3}\cup\{g(N)\}_{N\le7}\ [\text{machine: exhaustive membership test over}\ \{3,5,7\}\cup\{6,21,56\}\cup\{6,15,20\}\cup\{0,1,1,3,10,25,105\}]
\]

%dsp-level:0
\[
(4):\  \max_sD_S(s)=4\ \text{at}\ s_2\ [\mathrm{Def\text{-}BM.14}(3);\ \text{machine: exact comparison in}\ \mathbb M];
\]

%dsp-level:0
\[
\qquad \lambda^{(n)}_{\max}=n\cdot4=4n\ \text{and it is attained by}\ (s_2,\dots,s_2)\ \text{with multiplicity}\ 1\ [\mathrm{BM.114}(3);\ \text{machine}]
\]

%dsp-level:0
\[
\mathrm{FT.76}(3),\ \mathrm{FT.79}(2),\ \mathrm{FT.84}(2)\ \text{re-cited at their own loci};
\]

%dsp-level:0
\[
\qquad E_{\mathrm{int}}^{(n)}
\]

%dsp-level:0
\[
\in [12n,60n]\ \text{on the six-state ladder and}\ [2n,60n]\ \text{on the computed-cell ladder}\ [\mathrm{BM.115}(2)(\mathrm i),(\mathrm{ii})];
\]

%dsp-level:0
\[
\qquad \mathrm{BM.82}(1),(5),(7)\ \text{and}\ \mathrm{BM.69}(4)\text{'s four grounds re-cited}
\]

%dsp-level:0
\[
(5):\  \mathrm{FT.97}(5)\text{'s three-ceiling synthesis re-cited};
\]

%dsp-level:0
\[
\qquad \lambda^{(n)}\le4n\ \text{is the sum's own bound}\ [(4);\ \text{machine}];
\]

%dsp-level:0
\[
\qquad \text{the external-token scan}\ [\text{machine}]
\]

%dsp-level:0
\[
(6):\  6^n=(6,36,216);
\]

%dsp-level:0
\[
\qquad \#\text{distinct}=(6,21,55)\ [\mathrm{BM.114}(3),(5);\ \text{machine}];
\]

%dsp-level:0
\[
\qquad \#\{\text{empty loci}\}=0\ [\mathrm{BM.115}(5);\ \text{machine: all}\ 216\ \text{triples carry}\ E_{\mathrm{int}}\in[6,180]\subset\mathbb Z_{>0}]
\]

%dsp-level:0
\[
(7):\  \text{the standing-rule scans and the barred-object token scan}\ [\text{machine}].\ \blacksquare
\]

\textbf{Phys.} Asked to group the spectrum, the corpus gives what it has and no more. Eight groupings are derived and each has an exact value: the binomial charge classes, the sealing grammar's counts, the multiset orbits, the all-distinct orbits, the three ordered generation floors, the two ordered branches, the cyclic three of the flavour channel, and the subfield ladder of the eigenvalue field. The one the resonance spectrum actually carries is the three-level grouping by floor, branch and charge class with the orbit multiplicities, and at three chains that is six orbits of size one, thirty of size three and twenty of size six — fifty-six orbits giving fifty-five eigenvalues after the single powers-of-two collision. Twenty appears, and it is exactly the count of ways to pick three distinct states from six, which is also where the epsilon seal's free symmetric-group action lives; the tier states plainly that this is an orbit count and not a decuplet dimension, because the flavour symmetry the corpus derives is a cyclic three and there is no special unitary three anywhere in the inventory. The octet adjudication is re-certified and extended: the derived three is colour's, the derived eight is a barred capacity ratio, the derived ten is a grammar count, and eight does not occur among any of the four grouping families at any formed level. The ceilings behave exactly as they should — the single-chain pole ceiling is four, the composite ceiling is four times the chain count and is attained by the heaviest branch at every level, and a composite eigenvalue is a sum of single-chain reads rather than a single-chain read, so no ceiling is violated and none is extended. And the count that answers the forbidden-state question a second time is the simplest one: all two hundred sixteen ordered constituent triples are formed resonances with nonempty threshold loci, so the two heaviest compositions anyone might name are unread rather than unreachable.

\end{theorem}

\medskip\noindent\textit{\texorpdfstring{BM.118 — The status of the Antimatter and Resonance Tier: adjudication, scope, and the open items}{BM.118 - The status of the Antimatter and Resonance Tier: adjudication, scope, and the open items}.}\ Structural-constraints record: see Appendix B, Section B.1.31 of the Addendum.\par\medskip

\phantomsection
\section[Tier BM.22 — The Annihilation Channel, Concrete Spin, and…]{\texorpdfstring{Tier BM.22 — The Annihilation Channel, Concrete Spin, and Fermion Identification: the Janus Tunnel, the Zig-Zag World Line, and the Spin-Dependent Transport Defect}{Tier BM.22 - The Annihilation Channel, Concrete Spin, and Fermion Identification: the Janus Tunnel, the Zig-Zag World Line, and the Spin-Dependent Transport Defect}}

\textbf{Locus.} This tier takes the residue $(\mathrm{R9})$ that $\mathrm{BM.118}(7)$ named — \textit{the annihilation channel: an inter-sector transition operator, $\neg$in the inventory} — and constructs the channel, then follows the two chains that fall out of it: the concrete spin and the composite statistics. The crossing is not a velocity and the tier says so first: $c=1$ is a convention, the dispersion gives a strict subluminal deficit $1-v=m^2/(E(E+p))>0$ at every finite momentum, and the boundary the corpus actually derives between the materialized and the fluctuation content is the \textbf{Janus real/virtual register}, whose $\mathcal V$ sector is the maximally raw-entropic state $\rho_\circ=\tfrac14I$ with no fermion terminus at $|E_{\mathcal R}|=0$ ($\mathrm{BM.119}$). On that register the crossing operator is derived, not posited: the two moves $\hat J_\pm$ of $\mathrm{D1.2}$, the master generator $Q=x_+(\hat J_+-P_{\mathcal V})+x_-(\hat J_--P_{\mathcal R})$, its one-dimensional kernel giving the stationary measure, its spectrum $\{0,-(x_++x_-)\}$, and detailed balance as a theorem ($\mathrm{BM.120}$). The zig-zag is then the resolvent's own evenness — $(\hat U^{2m})_{\varnothing\varnothing}=(x_+x_-)^m$ with every odd term zero — so each term is a create-propagate-annihilate round trip whose legs all run forward in the ledger order, and the "backward leg" of a diagram is here the CP-conjugate \textit{configuration} on a forward leg, which is why no inverse tick is needed ($\mathrm{BM.121}$). The pair channel is derived with its charge balance, its exact threshold, its weight and its three CP legs — the nilpotent weak step for the charge inversion, the CP transport for the parity inversion, the local scale for the energy — while the two-photon supply is barred on four named grounds, so the channel closes and the amplitude does not ($\mathrm{BM.122}$). Annihilation is the reverse move, its content returning to the $\mathcal V$ sector at the entropy maximum rather than into radiation legs, and $(\mathrm{R9})$ closes at the channel level with the cross-section residue narrowing to the momentum-space resolvent and the gauge coupling of the charge read ($\mathrm{BM.123}$). On the spin side the tier separates two angles that a single-symbol reading would fuse: the winding $\Theta$ and the frame-holonomy angle $\Theta_R$ are provably not proportional, on a dimension mismatch, so the spin-dependent transport defect is a \textbf{multiplet ladder} $\Delta(m)=4\sin^2(m\Theta_R/4)$ over $m\in\{k,k-2,\dots,-k\}$ and not a number $\Delta(j)$ ($\mathrm{BM.124}$, $\mathrm{BM.125}$). A spin-dependent crossing weight is refuted on the four standing grounds, and what survives is a \textit{domain} correlation rather than a prohibition ($\mathrm{BM.126}$). The winding parity is then a ledger read, separated from the $\mathrm{D1.5}$ kind read that $\mathrm{BM.106}(1)$ already showed cannot carry a spin-parity theorem ($\mathrm{BM.127}$). Finally the three-chain exchange structure is built — $\mathcal H_3$, the $S_3$ action with its exact character formula $d^{\,c(\sigma)}$, and the vanishing of $\wedge^3(\mathbb C^2)$ against the one-dimensionality of $\wedge^3(\mathbb C^3)$, which is why the baryon's antisymmetry is colour's and the statistics is the winding's — giving the multiplicative exchange law $\varepsilon=\prod_a\varepsilon(c_a)$, the meson's unconditional boson sign, and the baryon's sign as an exact function of its three constituent windings ($\mathrm{BM.128}$, $\mathrm{BM.129}$). $\mathrm{BM.130}$ adjudicates.

\textbf{Standing rows.} $\alpha_B$ stays barred on $\mathrm{BM.61}(8)$'s four grounds; §Cal-BM stays empty except for $\mathrm{A\text{-}BM.9}$'s logged slot with its value $\neg$issued. No right-ray matter line, no seesaw-type mechanism, no Yukawa-type object and no external mixing matrix is suggested, added or required $[\mathrm{EW.15};\ \mathrm{FT.145}(3);\ \mathrm{BM.96}(4),(6)\ \text{re-cited},\ \neg\text{extended}]$. No clause of $\mathrm{BM.1}$–$\mathrm{BM.118}$ is retracted, replaced or renumbered: the rows this tier touches are re-scoped at their own loci — $\mathrm{BM.118}(7)$'s residue $(\mathrm{R9})$, $\mathrm{BM.112}(7)$'s and $\mathrm{Def\text{-}BM.16}(7)$'s residue $(\mathrm{R7})$ at $n\ge3$, $\mathrm{BM.109}(2)$'s pair-weight naming, $\mathrm{BM.114}(2)$'s colour-antisymmetry row, $\mathrm{BM.27}(7)$, $\mathrm{BM.29}(5)$ and $\mathrm{Def\text{-}BM.15}(5)$'s no-annihilation rows, $\mathrm{BM.6}(6)$'s resolvent specification and $\mathrm{BM.6}(2)$'s discharged specification — and each rescope is recorded in $\mathrm{BM.130}(5)$, in $\mathrm{BM.77}(7)$'s pattern. No measured coupling, mass, width, lifetime, cross-section or ratio enters any derivation $[\mathrm{P5}]$. No cosmological coordinate is used: a particle accelerator and a deep-gravity region are both local regions with high $(E,N,G)$, and every threshold and width is a function of those four local parameters alone $[\mathrm{Def\text{-}BM.17}(6)(\mathrm i);\ \mathrm{A6};\ \mathrm{D1.8}]$. "Low energy" and "high energy" are calibration points on the derived functions and $\neg$regimes of two theories $[\mathrm{BM.116}(6);\ \mathrm{Def\text{-}BM.17}(6)(\mathrm{ii})]$. No black-hole physics is derived: the comparison radius $r_s$ enters only through $\mathrm{BM.15}(4)$'s quarantined change of variable, with $\neg$a collapse reading and $\neg$a horizon identification $[\mathrm{P3};\ \mathrm{BM.115}(4)]$. Species names are confined to the identification layer $[\mathrm{Def\text{-}BM.3}(7)]$. Topology labels are internal bookkeeping $[\mathrm{Def\text{-}H.32}(4)]$. Amplitude $\ne$ mass $[\mathrm{FT.79}(5)]$. Rung $\ne$ scale $[\mathrm{D0.8}]$. No currency is converted $[\mathrm{BM.83}(3);\ \mathrm{BM.102}(6)]$. No theorem of this tier imports a HistLog disposition as premise: Q9, Q11, items 7, 13, 17, 20 and deficits $1'$, $3$, $4$, $6$, $11$, $12$ are cited as pointers. One dimension only: units are powers of $\mathfrak q$ $[\mathrm{D2.9};\ \mathrm{CT2.1}]$.

\textbf{Local notation, fixed for the whole tier.} $\mathcal J:=\operatorname{span}_{\mathbb C}\{|\mathcal V\rangle,\ |\mathcal R\rangle\}\cong\mathbb C^2$ the per-pair Janus register in the (virtual, real) order of $\mathrm{FT.108}(4)$; $P_{\mathcal V},P_{\mathcal R}$ its two orthogonal projectors; $\hat J_+:=|\mathcal R\rangle\langle\mathcal V|$ the crystallization move and $\hat J_-:=|\mathcal V\rangle\langle\mathcal R|$ the dissolution move $[\mathrm{D1.2}]$; $Q$ the master generator over $\mathbb Z[x_+,x_-]$ $[\mathrm{FT.108}(4)]$; $\varepsilon_q\in[0,\tfrac14)$, $x_+=1-4\varepsilon_q$, $x_-=4\varepsilon_q$, $p=x_+/(x_++x_-)=1-4\varepsilon_q$ $[\mathrm{A4};\ \mathrm{D0.4};\ \mathrm{CT0.8};\ \mathrm{BM.19}(2)]$; $w$ the formal fugacity and $\hat U$ the single-channel step $[\mathrm{FT.108}(5);\ \mathrm{Def\text{-}H.32}(3)]$; $\Theta(c)=\sum_{e\in E(c)}\widehat{ds}^{\,2}(e)\in2\mathbb Z$ the winding, $k=\Theta/2\in\mathbb Z_{\ge0}$, $j=|k|/2$, $\Psi(c)=\mathrm{Sym}^{|k|}(\mathcal H_{\mathsf v})$ with $\mathcal H_{\mathsf v}\cong\mathbb C^2$ $[\mathrm{FT.22};\ \mathrm{Def\text{-}H.38};\ \mathrm{EW.32}]$; $\Theta_R$ the frame-holonomy angle with $[\Theta_R]=\mathfrak q^0$ $[\mathrm{Def\text{-}H.21}(5),(6);\ \mathrm{FT.55}(1)]$; $\phi:=\Theta_R/2$; $T_m$ the Chebyshev polynomial of the first kind; $\Delta(m):=4\sin^2(m\Theta_R/4)=2-2\cos(m\phi)$ the transport defect at multiplet weight $m$ $[\mathrm{FT.60}(5);\ \mathrm{FT.54}(4);\ \mathrm{EW.33}]$; $\mathcal H_\omega(c)$, $\hat\omega$, $\hat\chi$, $P_{L/R}$, $|L\rangle$, $|R\rangle$ as in $\mathrm{Def\text{-}H.33}$; $\widehat V=E_+\otimes P_L$ the nilpotent step $[\mathrm{BM.76}(1)]$; $\mu_{\mathrm{CP}}=e^{i\pi\Theta}$ and $\varepsilon(c)=(-1)^{|k(c)|}=(-1)^{2j(c)}$ $[\mathrm{BM.108}(3);\ \mathrm{BM.113}(2),(4)]$; $\mathcal H_2$, $\hat P_{12}$ as in $\mathrm{Def\text{-}BM.16}$; $d:=\Theta(P)\ge0$ a connector winding $[\mathrm{BM.2}(2);\ \mathrm{FT.34}(3)]$; $S_\sigma$ the sign of a permutation $\sigma$; $c(\sigma)$ its cycle count; $\mathcal T_0$, $\mathcal T$, $\Gamma$, $\mathcal N$, $\chi_{\{\cdot\}}$ as in $\mathrm{Def\text{-}BM.17}$; $[\cdot]$ the $\mathfrak q$-exponent map $[\mathrm{CT2.1};\ \mathrm{D2.9}]$.

\begin{definition}[{\texorpdfstring{Def-BM.18 — The Janus register and its crossing operators}{Def-BM.18 - The Janus register and its crossing operators}}]\label{tag:Def-BM-18}
\textbf{Def.} A construction on $\mathrm{D1.1}$–$\mathrm{D1.2}$'s real/virtual distinction, $\mathrm{FT.108}(4)$'s master generator and $\mathrm{Def\text{-}BM.16}$'s two-chain space; no new axiom, no new datum, no new parameter, no velocity operator, no position operator, no external scattering object.

%dsp-level:0
\[
(1)\quad\textbf{the register:}
\]

%dsp-level:0
\[
\qquad \boxed{\ \mathcal J\ :=\ \operatorname{span}_{\mathbb C}\bigl\{|\mathcal V\rangle,\ |\mathcal R\rangle\bigr\}\ \cong\ \mathbb C^2\ }
\]

%dsp-level:0
\[
\qquad \text{in the}\ (\text{virtual},\ \text{real})\ \text{order of}\ \mathrm{FT.108}(4);
\]

%dsp-level:0
\[
e\in\mathcal R\iff\bigl|\Pi_{\mathcal V}\rho\bigr|\le\varepsilon_q\ \vee\ \bigl|\Pi_{\mathcal V}\rho\bigr|\ge1-\varepsilon_q;\qquad\neg\text{registration}\implies e\in\mathcal V\qquad[\mathrm{A4};\ \mathrm{D1.1};\ \mathrm{CT0.1}]:
\]

%dsp-level:0
\[
\textbf{constitutive clause:}\ \ \mathcal J\ \text{is the two-state register of one pair across the}\ \mathrm{A4}\ \text{criterion};
\]

%dsp-level:0
\[
\qquad \textbf{non-derived clause:}\ \ \text{the basis order}\ (\mathcal V,\mathcal R);
\]

%dsp-level:0
\[
(2)\quad\textbf{the two crossing operators:}
\]

%dsp-level:0
\[
\qquad \boxed{\ \hat J_+\ :=\ |\mathcal R\rangle\langle\mathcal V|\ \ (\text{crystallization},\ J^+:\ n_{\mathcal R}:0\to1),\qquad\hat J_-\ :=\ |\mathcal V\rangle\langle\mathcal R|\ \ (\text{dissolution},\ J^-:\ n_{\mathcal R}:1\to0)\ }
\]

%dsp-level:0
\[
[\mathrm{D1.2}:\ \text{a seal is the transition of one pair across the}\ \mathrm{A4}\ \text{criterion, in either branch}];
\]

%dsp-level:0
\[
\qquad P_{\mathcal V}:=|\mathcal V\rangle\langle\mathcal V|,\quad P_{\mathcal R}:=|\mathcal R\rangle\langle\mathcal R|,\quad P_{\mathcal V}+P_{\mathcal R}=\mathrm{id}
\]

%dsp-level:0
\[
\qquad [\text{machine}];
\]

%dsp-level:0
\[
(3)\quad\textbf{the master generator:}\qquad\boxed{\ Q\ :=\ \begin{pmatrix}-x_+&x_+\\ x_-&-x_-\end{pmatrix}\ \ \text{over}\ \ \mathbb Z[x_+,x_-]\ }\qquad[\mathrm{FT.108}(4):\ \text{the per-pair Janus dynamics}]:
\]

%dsp-level:0
\[
\textbf{the decomposition (derived, recorded here):}
\]

%dsp-level:0
\[
\qquad Q\ =\ x_+\bigl(\hat J_+-P_{\mathcal V}\bigr)\ +\ x_-\bigl(\hat J_--P_{\mathcal R}\bigr)
\]

%dsp-level:0
\[
\qquad [\text{machine: entrywise at six rational witnesses}];
\]

%dsp-level:0
\[
\textbf{constitutive clause:}\ \ Q\ \text{is}\ \mathrm{FT.108}(4)\text{'s own generator,}\ \neg\text{a new dynamics};
\]

%dsp-level:0
\[
\qquad \textbf{the row convention:}\ \ Q\ \text{acts on row vectors, so}\ Q_{ij}\ \text{is the rate from}\ i\ \text{to}\ j;
\]

%dsp-level:0
\[
(4)\quad\textbf{the crossing space of a chain:}
\]

%dsp-level:0
\[
\qquad \boxed{\ \mathcal T_{\mathcal J}(c)\ :=\ \mathcal J\ \otimes\ \mathcal H_\omega(c)\ \otimes\ \Psi(c),\qquad\dim\mathcal T_{\mathcal J}(c)\ =\ 4\,\bigl(|k(c)|+1\bigr)\ }
\]

%dsp-level:0
\[
[\mathrm{Def\text{-}H.33};\ \mathrm{Def\text{-}H.38};\ \mathrm{Def\text{-}BM.16}(1);\ \text{machine: exact integers}];
\]

%dsp-level:0
\[
\qquad \textbf{the pair crossing space:}\quad\mathcal T_{\mathcal J}^{\,2}(c):=\mathcal T_{\mathcal J}(c)\otimes\mathcal T_{\mathcal J}(c^{\mathrm{CP}}),
\]

%dsp-level:0
\[
\qquad \dim=16\,\bigl(|k|+1\bigr)^2;
\]

%dsp-level:0
\[
(5)\quad\textbf{the channel clauses:}
\]

%dsp-level:0
\[
\qquad \textbf{a pair creation}\ :
\]

%dsp-level:0
\[
= \ \hat J_+\ \text{on two registers, with the two members identified as a CP-pair}\ [\mathrm{Def\text{-}BM.15}(3)];
\]

%dsp-level:0
\[
\qquad \textbf{weight}\ x_+^2\ \text{for the two one-way moves and}\ x_+x_-\ \textbf{for one round trip};
\]

%dsp-level:0
\[
\textbf{an annihilation}\ :=\ \hat J_-\ \text{on the same two registers};\qquad\textbf{weight}\ x_-^2\ \text{one-way,}\ x_+x_-\ \text{per round trip};
\]

%dsp-level:0
\[
\textbf{a zig-zag}\ :
\]

%dsp-level:0
\[
= \ \text{a finite alternating sequence of}\ \hat J_+\ \text{and}\ \hat J_-\ \text{moves, every step recorded in the ledger order}\ [\mathrm{A6};\ \mathrm{D1.8}];
\]

%dsp-level:0
\[
\textbf{constitutive clause:}\ \ \text{the three words name move sequences on}\ \mathcal J,
\]

%dsp-level:0
\[
\ \neg\text{world lines in a spacetime and}\ \neg\text{asymptotic scattering states};
\]

%dsp-level:0
\[
(6)\quad\textbf{the bars (standing):}
\]

%dsp-level:0
\[
\qquad \neg\exists\ \text{a velocity operator, a position operator or a momentum-space state on}\ \mathcal J\ [\mathrm{BM.10}(1):\ \neg[\hat x,\hat p]=i\ \text{in the finite register};\ \mathrm{BM.6}(6):\ \text{the momentum-space}\ \Sigma\text{-resolvent is an open interface}];
\]

%dsp-level:0
\[
\neg\exists\ \text{a cross-section, an S-matrix element or a Lorentz-invariant phase space}\ [\mathrm{BM.6}(6);\ \mathrm{BM.118}(7)(\mathrm{R9})];
\]

%dsp-level:0
\[
\neg\exists\ \text{a direct}\ \mathrm{EM}^+\to\mathrm{EM}^-\ \text{operator: the sectors are amplitude-equal configuration orbits and the channel that connects them is mediated by}\ \mathcal J\ \text{and}\ \mathrm{CP}\ [\mathrm{BM.25}(4);\ \mathrm{BM.27}(6),(7);\ \mathrm{BM.29}(5)];
\]

%dsp-level:0
\[
\textbf{the kind bar:}
\]

%dsp-level:0
\[
\qquad \mathcal J\ \text{is}\ \neg\mathcal H_\omega\ \text{and}\ \neg\text{one of the four single-object}\ \mathbb Z_2\ \text{gradings}\ [\mathrm{BM.106}(4);\ \mathrm{BM.112}(5);\ \mathrm{FT.125}(5);\ \mathrm{FT.89}(2)]:
\]

%dsp-level:0
\[
\boxed{\ \textbf{the real/virtual distinction is the Janus register; the orientation distinction is}\ \mathcal H_\omega\textbf{; identifying them is a barred kind confusion}\ }
\]

%dsp-level:0
\[
(7)\quad\textbf{status and scope:}
\]

%dsp-level:0
\[
\qquad \textbf{derived:}\ (3)\text{'s decomposition},\ (4)\text{'s dimensions};
\]

%dsp-level:0
\[
\qquad \textbf{constitutive:}\ (1)\text{'s register},\ (2)\text{'s two moves},\ (5)\text{'s three names};
\]

%dsp-level:0
\[
\textbf{the scope inherited at every use:}\ \ \mathrm{Def\text{-}H.32}\text{'s}\ (\mathrm{S1})\text{--}(\mathrm{S4});
\]

%dsp-level:0
\[
\qquad \mathrm{Def\text{-}H.38}\text{'s nonnegative cone and its}\ k\ge3\ \text{aspect bar};
\]

%dsp-level:0
\[
\qquad \mathrm{BM.8}(2)\text{'s stationarity hypothesis for any chart evolution};
\]

%dsp-level:0
\[
\neg\text{a P4 addition},
\]

%dsp-level:0
\[
\quad\neg\text{a new dimensionful scale}\ [\mathrm{P1}],
\]

%dsp-level:0
\[
\quad\neg\text{a measured value}\ [\mathrm{P5}],
\]

%dsp-level:0
\[
\quad\neg\text{a species identification}\ [\mathrm{Def\text{-}BM.3}(7)],
\]

%dsp-level:0
\[
\quad\neg\text{a rate datum}\ [\mathrm{Cal\text{-}F}\ \text{carries no rate};\ \mathrm{BM.44}(5)];
\]

%dsp-level:0
\[
\textbf{beyond this tier:}\ \ \text{the momentum-space}\ \Sigma\text{-resolvent}\ [\mathrm{BM.6}(6)],
\]

%dsp-level:0
\[
\qquad \text{the gauge coupling of the charge read}\ [\mathrm{Def\text{-}BM.8}(8)].
\]

\textbf{Rmk (what the register is and is not).} $\mathcal J$ is the two-state register of \textit{one pair} across the crystallization criterion — the object D1.2 named and FT.108 derived a stationary measure for. It is not a spacetime, not a world line, and not a Hilbert space of asymptotic states. Clause (6)'s kind bar is the load-bearing one: the corpus has an orientation doublet $\mathcal H_\omega$ per chain, and it would be a category error to read the real/virtual distinction as that doublet's two rays, exactly as it would be to read a single-chain grade as a permutation sign. The zig-zag of clause (5) is a word in two moves, recorded forward, and its weight is a monomial in $x_+$ and $x_-$.

\textbf{Req.} A4, A6, D0.8, D1.1, D1.2, D1.8 (Part 1); CT0.1, CT0.8, CT2.1, D2.9 (Part 1, Part 2); P1, P3, P4, P5 (Parts 1–2); Def-H.21(5), Def-H.21(6), Def-H.32, Def-H.33, Def-H.38, FT.22, FT.89(2), FT.108(4), FT.108(5), FT.125(5) (Part 3); Def-BM.3(7), Def-BM.15(3), Def-BM.16(1), BM.2(2), BM.6(6), BM.8(2), BM.10(1), BM.25(4), BM.27(6), BM.27(7), BM.29(5), BM.44(5), BM.106(4), BM.112(5), BM.118(7) (this part); exact $2\times2$ matrix algebra over $\mathbb Z[x_+,x_-]$ and over $\mathbb Q(i)$, exact integer dimension arithmetic [math; machine: \texttt{checks/\allowbreak{}verify\_\allowbreak{}H\_\allowbreak{}janus\_\allowbreak{}tunnel\_\allowbreak{}concrete\_\allowbreak{}spin.py} §A].

\textbf{Pf.}

%dsp-level:0
\[
\text{Definition block: the constitutive clauses are}\ (1)\text{'s register},
\]

%dsp-level:0
\[
\ (2)\text{'s two moves and}\ (5)\text{'s three names};
\]

%dsp-level:0
\[
\ \text{the non-derived clauses are the basis order of}\ (1)\ \text{and the row convention of}\ (3);
\]

%dsp-level:0
\[
\ (3)\text{'s decomposition and}\ (4)\text{'s dimensions are derived and machine-checked};
\]

%dsp-level:0
\[
\ (6),
\]

%dsp-level:0
\[
(7)\ \text{are bars and scope records}.\ \blacksquare
\]

\textbf{Phys.} One pair, crossing a criterion, has two states and two moves, and the corpus already derived both: D1.2 named the seal as a transition across the A4 criterion in either branch, and FT.108 wrote down the master generator whose kernel is one-dimensional and gives the stationary measure. This definition makes the two moves into operators on a two-dimensional register, decomposes the generator as a sum of jump-minus-loss terms with the two derived weights, and tensors the register with the chain's own orientation doublet and winding carrier to give a crossing space of exact finite dimension. Three words are then defined and no more: a pair creation is the crystallization move on two registers whose members are identified as a CP-pair, an annihilation is the dissolution move on the same two, and a zig-zag is a finite alternating word in the two moves, every step recorded in ledger order. What the definition refuses is the vocabulary that would overreach — no velocity operator, no position operator, no momentum-space state, no cross-section, no S-matrix, and above all no identification of the real/virtual register with the orientation doublet, which is a different two-dimensional object doing different work.

\end{definition}

\medskip\noindent\textit{\texorpdfstring{BM.119 — The c-threshold adjudication: no velocity boundary, the exact subluminal deficit, and the register boundary that is derived}{BM.119 - The c-threshold adjudication: no velocity boundary, the exact subluminal deficit, and the register boundary that is derived}.}\ Structural-constraints record: see Appendix B, Section B.1.32 of the Addendum.\par\medskip

\begin{theorem}[{\texorpdfstring{BM.120 — The Janus tunnel: the crossing operators, the exact spectrum, and detailed balance}{BM.120 - The Janus tunnel: the crossing operators, the exact spectrum, and detailed balance}}]\label{tag:BM-120}
\textbf{Thm.}

%dsp-level:0
\[
(1)\quad\textbf{the move algebra (exact, over}\ \mathbb Q\text{):}
\]

%dsp-level:0
\[
\qquad \hat J_+^2=0,
\]

%dsp-level:0
\[
\qquad \hat J_-^2=0,
\]

%dsp-level:0
\[
\qquad \hat J_+\hat J_-=P_{\mathcal V},
\]

%dsp-level:0
\[
\qquad \hat J_-\hat J_+=P_{\mathcal R},
\]

%dsp-level:0
\[
\qquad \{\hat J_+,\hat J_-\}=\mathrm{id}
\]

%dsp-level:0
\[
[\text{machine: all five identities entrywise on the}\ 2\times2\ \text{register}];
\]

%dsp-level:0
\[
\qquad \boxed{\ \textbf{each move is nilpotent and the two anticommute to the identity}\ }
\]

%dsp-level:0
\[
\textbf{the reading:}\quad\hat J_+\hat J_-=P_{\mathcal V}\ \text{is the round trip}\ \mathcal V\to\mathcal R\to\mathcal V;
\]

%dsp-level:0
\[
\qquad \hat J_-\hat J_+=P_{\mathcal R}\ \text{is the round trip}\ \mathcal R\to\mathcal V\to\mathcal R;
\]

%dsp-level:0
\[
\neg\exists\ \text{a doubled move in the same direction}\ —\ \textbf{the}\ \mathrm{A4}\ \textbf{criterion is binary, and the nilpotency is its operator form}\ [\mathrm{D1.2};\ \mathrm{BM.23}(2)\text{'s overflow pattern}];
\]

%dsp-level:0
\[
(2)\quad\textbf{the generator decomposes as jump minus loss:}
\]

%dsp-level:0
\[
\qquad \boxed{\ Q\ =\ x_+\bigl(\hat J_+-P_{\mathcal V}\bigr)\ +\ x_-\bigl(\hat J_--P_{\mathcal R}\bigr)\ }
\]

%dsp-level:0
\[
\qquad [\mathrm{FT.108}(4);\ \text{machine: entrywise at six rational witnesses}\ (x_+,x_-)\in\{(\tfrac12,\tfrac12),(\tfrac34,\tfrac14),(1,0),(0,1),(\tfrac23,\tfrac13),(\tfrac78,\tfrac18)\}];
\]

%dsp-level:0
\[
\textbf{the row-sum law:}
\]

%dsp-level:0
\[
\qquad Q\ \text{acts on row vectors and every row sums to zero:}\quad(-x_+)+x_+=0,
\]

%dsp-level:0
\[
\qquad x_-+(-x_-)=0
\]

%dsp-level:0
\[
\qquad [\text{machine}]:
\]

%dsp-level:0
\[
\therefore\ Q\ \textbf{is a master generator and}\ \neg\textbf{a unitary}\ —\ \text{the crossing is a stochastic ledger move,}\ \neg\text{a rotation of a state};
\]

%dsp-level:0
\[
(3)\quad\textbf{the kernel, the stationary measure and detailed balance:}
\]

%dsp-level:0
\[
\qquad (x_-,\ x_+)\,Q\ =\ \bigl(x_-(-x_+)+x_+x_-,\ \ x_-x_++x_+(-x_-)\bigr)\ =\ (0,\ 0)
\]

%dsp-level:0
\[
\qquad [\text{machine: exact in}\ \mathbb Z[x_+,x_-]];
\]

%dsp-level:0
\[
\det Q\ =\ (-x_+)(-x_-)-x_+x_-\ =\ 0,\qquad Q\ne0\implies\operatorname{rk}Q=1,\qquad\dim\ker Q=1\qquad[\mathrm{FT.108}(4);\ \text{machine}];
\]

%dsp-level:0
\[
\boxed{\ p\ =\ \frac{x_+}{x_++x_-}\ =\ 1-4\varepsilon_q\ \ \textbf{identically}\ }
\]

%dsp-level:0
\[
\qquad [\mathrm{BM.19}(2);\ \text{machine: exact at seven rational thresholds}];
\]

%dsp-level:0
\[
\textbf{detailed balance is derived, not assumed:}
\]

%dsp-level:0
\[
\qquad p_{\mathcal V}\,x_+\ =\ p_{\mathcal R}\,x_-\iff\frac{x_-}{x_++x_-}\,x_+=\frac{x_+}{x_++x_-}\,x_-\iff x_-x_+=x_+x_-
\]

%dsp-level:0
\[
\qquad [\mathrm{FT.108}(4);\ \text{machine}];
\]

%dsp-level:0
\[
(4)\quad\textbf{the exact spectrum:}
\]

%dsp-level:0
\[
\qquad \operatorname{tr}Q=-(x_++x_-),
\]

%dsp-level:0
\[
\qquad \det Q=0\implies\boxed{\ \operatorname{spec}Q=\bigl\{0,\ -(x_++x_-)\bigr\}\ }
\]

%dsp-level:0
\[
\qquad [\text{machine: the characteristic polynomial}\ \lambda^2+(x_++x_-)\lambda];
\]

%dsp-level:0
\[
\textbf{at the specialization}\ x_++x_-=1:\qquad\operatorname{spec}Q=\{0,\ -1\},\qquad\operatorname{tr}Q=-1,\qquad\operatorname{rk}Q=1\qquad[\text{machine}];
\]

%dsp-level:0
\[
\boxed{\ \textbf{the zero eigenvalue is the stationary mode and}\ -1\ \textbf{is the relaxation mode at unit total weight}\ }
\]

%dsp-level:0
\[
\qquad [\text{machine: the two eigenvectors}\ (x_-,x_+)\ \text{and}\ (1,-1)];
\]

%dsp-level:0
\[
(5)\quad\textbf{the tunnel operator and its weights:}
\]

%dsp-level:0
\[
\qquad \textbf{one-way creation}\ \hat J_+:\ \text{weight}\ x_+=1-4\varepsilon_q\in(0,1];
\]

%dsp-level:0
\[
\qquad \textbf{one-way annihilation}\ \hat J_-:\ \text{weight}\ x_-=4\varepsilon_q\in[0,1);
\]

%dsp-level:0
\[
\textbf{the round trip:}
\]

%dsp-level:0
\[
\qquad \boxed{\ \mathcal W_{\text{round}}\ =\ x_+x_-\ =\ 4\varepsilon_q(1-4\varepsilon_q)\ =\ \frac14-16\Bigl(\varepsilon_q-\frac18\Bigr)^{\!2}\ \le\ \frac14\ }
\]

%dsp-level:0
\[
\qquad [\mathrm{BM.109}(2);\ \mathrm{BM.44}(1);\ \text{machine: the SOS completion and the}\ 256\text{-point grid}];
\]

%dsp-level:0
\[
\textbf{the detailed-balance point is the maximum:}
\]

%dsp-level:0
\[
\qquad x_+=x_-\iff\varepsilon_q=\frac18,
\]

%dsp-level:0
\[
\qquad \text{and}
\]

%dsp-level:0
\[
\qquad x_+x_-\le\Bigl(\frac{x_++x_-}2\Bigr)^{\!2}=\frac14\ \ \text{with equality}\iff x_+=x_-
\]

%dsp-level:0
\[
[\text{machine: the AM--GM row on exact rationals, equality locus verified}];
\]

%dsp-level:0
\[
\qquad \boxed{\ \textbf{the round-trip weight is maximal exactly where the two moves balance}\ }
\]

%dsp-level:0
\[
(6)\quad\textbf{the type correction (the crossing acts on}\ \mathcal J\textbf{, not on}\ \mathcal H_\omega\textbf{):}
\]

%dsp-level:0
\[
\qquad \mathcal H_\omega(c)=\operatorname{span}\{|c\rangle,|\omega(c)\rangle\}\ \text{is one chain's orientation doublet}\ [\mathrm{Def\text{-}H.33}];
\]

%dsp-level:0
\[
\qquad \mathcal J=\operatorname{span}\{|\mathcal V\rangle,|\mathcal R\rangle\}\ \text{is one pair's registration register}\ [\mathrm{Def\text{-}BM.18}(1)]:
\]

%dsp-level:0
\[
\textbf{the two are different objects with different involutions:}
\]

%dsp-level:0
\[
\qquad \hat\omega^2=\mathrm{id}\ \text{with}\ \operatorname{spec}\{\pm1\}\ \text{and}\ \operatorname{tr}\hat\omega=0\ [\mathrm{FT.122}(1)];
\]

%dsp-level:0
\[
\qquad \hat J_\pm^2=0\ \text{with}\ \operatorname{tr}\hat J_\pm=0\ \text{and}\ \neg\text{diagonalizable}\ [(1)]:
\]

%dsp-level:0
\[
\boxed{\ \neg\exists\ \text{an isomorphism}\ \mathcal J\cong\mathcal H_\omega\ \textbf{that respects the two operator algebras}\ —\ \hat\omega\ \text{is an involution and}\ \hat J_\pm\ \text{are nilpotent}\ }
\]

%dsp-level:0
\[
\qquad [\text{machine:}\ \hat\omega^2=\mathrm{id}\ne0\ \text{against}\ \hat J_\pm^2=0];
\]

%dsp-level:0
\[
\therefore\ \textbf{the crossing operator's type is}\quad\hat J_\pm\otimes\mathrm{id}\otimes\mathrm{id}\ \ \text{on}\ \ \mathcal T_{\mathcal J}(c)=\mathcal J\otimes\mathcal H_\omega(c)\otimes\Psi(c)
\]

%dsp-level:0
\[
\qquad [\mathrm{Def\text{-}BM.18}(4);\ \text{machine}];
\]

%dsp-level:0
\[
(7)\quad\textbf{the forward-only row:}
\]

%dsp-level:0
\[
\qquad \textbf{the move system is bidirectional as a set of moves and has a stationary measure}\ [(3)];
\]

%dsp-level:0
\[
\qquad \textbf{the tick is not bidirectional:}\quad\neg\exists\,\tau^{-1}\ [\mathrm{A6}];
\]

%dsp-level:0
\[
\bigl|\mathfrak L(\tau_{n+1})\bigr|=\bigl|\mathfrak L(\tau_n)\bigr|+n_{\mathrm{real\_out}}(\tau_n)\ge\bigl|\mathfrak L(\tau_n)\bigr|\ [\mathrm{A6};\ \mathrm{A7}];\qquad V'=\bigl(1-\tfrac\gamma4\bigr)^2V\le V\ [\mathrm{L2}];
\]

%dsp-level:0
\[
\boxed{\ \textbf{the ensemble sums both moves; the ledger records only growth}\ —\ \text{the formal mirror exists and the dynamical reversal does not, on the Janus register exactly as on the ledger}\ }
\]

%dsp-level:0
\[
\qquad [\mathrm{BM.25}(2);\ \mathrm{Def\text{-}BM.7}(2);\ \mathrm{FT.144}(4)];
\]

%dsp-level:0
\[
(8)\quad\textbf{status:}\qquad\textbf{derived:}\ (1)\text{--}(7);
\]

%dsp-level:0
\[
\neg\text{a declaration},
\]

%dsp-level:0
\[
\quad\neg\text{a P4 addition},
\]

%dsp-level:0
\[
\quad\neg\text{a new dimensionful scale}\ [\mathrm{P1}],
\]

%dsp-level:0
\[
\quad\neg\text{a measured value}\ [\mathrm{P5}],
\]

%dsp-level:0
\[
\quad\neg\text{a numerical coupling}\ [\mathrm{EM.25}(7)],
\]

%dsp-level:0
\[
\quad\neg\text{a rate datum in external units}\ [\mathrm{BM.44}(5)],
\]

%dsp-level:0
\[
\quad\neg\text{a species identification}\ [\mathrm{Def\text{-}BM.3}(7)];
\]

\textbf{Req.} A4, A6, A7, D0.4, D1.2, D1.8 (Part 1); CT0.8, EM.25(7) (Part 1, Part 2); P1, P3, P4, P5 (Parts 1–2); L2 (Part 1); Def-H.32, Def-H.33, FT.108(4), FT.108(5), FT.122(1), FT.144(4) (Part 3); Def-BM.3(7), Def-BM.7(2), Def-BM.18(1), Def-BM.18(4), BM.19(2), BM.23(2), BM.25(2), BM.44(1), BM.44(5), BM.109(2), BM.113(2) (this part); exact $2\times2$ matrix algebra over $\mathbb Z[x_+,x_-]$ and $\mathbb Q$, the SOS parabola completion, AM–GM on exact rationals with its equality locus, exhaustive rational grid censuses [math; machine: \texttt{checks/\allowbreak{}verify\_\allowbreak{}H\_\allowbreak{}janus\_\allowbreak{}tunnel\_\allowbreak{}concrete\_\allowbreak{}spin.py} §A].

\textbf{Pf.}

%dsp-level:0
\[
(1):\ \hat J_+=\left(\begin{smallmatrix}0&1\\0&0\end{smallmatrix}\right),\ \hat J_-=\left(\begin{smallmatrix}0&0\\1&0\end{smallmatrix}\right)\ \text{in the}\ (\mathcal V,\mathcal R)\ \text{row convention}\ [\mathrm{Def\text{-}BM.18}(2)]:
\]

%dsp-level:0
\[
\hat J_+^2=\left(\begin{smallmatrix}0&1\\0&0\end{smallmatrix}\right)^{\!2}=0,\qquad\hat J_-^2=\left(\begin{smallmatrix}0&0\\1&0\end{smallmatrix}\right)^{\!2}=0;\qquad\hat J_+\hat J_-=\left(\begin{smallmatrix}1&0\\0&0\end{smallmatrix}\right)=P_{\mathcal V},\qquad\hat J_-\hat J_+=\left(\begin{smallmatrix}0&0\\0&1\end{smallmatrix}\right)=P_{\mathcal R};
\]

%dsp-level:0
\[
\hat J_+\hat J_-+\hat J_-\hat J_+=P_{\mathcal V}+P_{\mathcal R}=\mathrm{id}\ [\text{machine: all five identities entrywise over}\ \mathbb Q]
\]

%dsp-level:0
\[
(2):\ x_+(\hat J_+-P_{\mathcal V})+x_-(\hat J_--P_{\mathcal R})\ \text{expanded entrywise gives}\ \left(\begin{smallmatrix}-x_+&x_+\\x_-&-x_-\end{smallmatrix}\right)=Q\ [\text{machine: six rational witnesses}];\qquad\text{the two row sums}\ [\text{machine}]
\]

%dsp-level:0
\[
(3):\ (x_-,x_+)Q=(0,0)\ [\text{machine: exact in}\ \mathbb Z[x_+,x_-]];
\]

%dsp-level:0
\[
\qquad \det Q=x_+x_--x_+x_-=0;
\]

%dsp-level:0
\[
\qquad Q\ne0\ \text{at}\ x_+x_->0\implies\operatorname{rk}=1\ [\text{machine}]
\]

%dsp-level:0
\[
p=\tfrac{x_+}{x_++x_-}=\tfrac{1-4\varepsilon_q}{1-4\varepsilon_q+4\varepsilon_q}=1-4\varepsilon_q\ [\text{machine: exact at seven thresholds}];
\]

%dsp-level:0
\[
\qquad \text{detailed balance}: p_{\mathcal V}x_+=\tfrac{x_-x_+}{x_++x_-}=p_{\mathcal R}x_-\ [\text{machine}]
\]

%dsp-level:0
\[
(4):\ \lambda^2-(\operatorname{tr}Q)\lambda+\det Q=\lambda^2+(x_++x_-)\lambda=\lambda(\lambda+x_++x_-)\ [\text{machine}];
\]

%dsp-level:0
\[
\qquad \text{at}\ x_++x_-=1:\ \{0,-1\}\ [\text{machine}]
\]

%dsp-level:0
\[
Q(x_-,x_+)
\]

%dsp-level:0
\[
= (0,0)\ \text{and}\ Q(1,-1)
\]

%dsp-level:0
\[
= (-(x_++x_-),\ (x_++x_-))
\]

%dsp-level:0
\[
= -(x_++x_-)(1,-1)\ [\text{machine: the two eigenvectors}]
\]

%dsp-level:0
\[
(5):\ x_+x_-=4\varepsilon_q(1-4\varepsilon_q)=\tfrac14-16(\varepsilon_q-\tfrac18)^2\ [\text{machine: expansion}];
\]

%dsp-level:0
\[
\qquad \text{the}\ 256\text{-point grid:}\ \le\tfrac14,\ =\tfrac14\ \text{only at}\ \varepsilon_q=\tfrac18\ [\text{machine}]
\]

%dsp-level:0
\[
\text{AM--GM}:\ x_+x_-\le\bigl(\tfrac{x_++x_-}2\bigr)^2
\]

%dsp-level:0
\[
\iff 0
\]

%dsp-level:0
\[
\le (x_+-x_-)^2\ [\text{machine: the square is nonnegative on exact rationals}],
\]

%dsp-level:0
\[
\ \text{equality}
\]

%dsp-level:0
\[
\iff x_+
\]

%dsp-level:0
\[
= x_-
\]

%dsp-level:0
\[
\iff 1-4\varepsilon_q
\]

%dsp-level:0
\[
= 4\varepsilon_q\iff\varepsilon_q
\]

%dsp-level:0
\[
= \tfrac18\ [\text{machine}]
\]

%dsp-level:0
\[
(6):\ \hat\omega=\left(\begin{smallmatrix}0&1\\1&0\end{smallmatrix}\right),\ \hat\omega^2=\mathrm{id},\ \operatorname{tr}\hat\omega=0,\ \operatorname{spec}=\{\pm1\}\ [\mathrm{Def\text{-}H.33};\ \mathrm{FT.122}(1);\ \text{machine}];\qquad\hat J_\pm^2=0\ne\mathrm{id}\ [(1);\ \text{machine}]\implies\neg\text{an algebra isomorphism}
\]

%dsp-level:0
\[
(7):\ \mathrm{A6},\ \mathrm{A7},\ \mathrm{L2},\ \mathrm{BM.25}(2),\ \mathrm{FT.144}(4)\ \text{rows re-cited};
\]

%dsp-level:0
\[
\qquad (8):\ \text{the standing-rule scans}\ [\text{machine}].\ \blacksquare
\]

\textbf{Phys.} The crossing has an algebra and it is small enough to write down completely. Each move squares to zero, because the crystallization criterion is binary and a pair cannot crystallize twice without dissolving in between; the two moves anticommute to the identity, and their two products are exactly the two round trips. The generator is the sum of a jump term and a loss term for each move, its rows sum to zero, and it is therefore a master generator rather than a unitary — the crossing is a stochastic ledger move, not a rotation of a state. Its kernel is the line spanned by the two weights, which is where the stationary measure comes from, and detailed balance is an identity rather than an assumption: the flux out of the virtual state equals the flux out of the real one at that measure, because both are the same product. The spectrum is zero and minus the total weight, so at unit total weight there is one stationary mode and one relaxation mode with rate exactly one. The round-trip weight is the familiar parabola with maximum one quarter, and the maximum sits exactly where the two moves balance — an AM–GM equality, not a coincidence. Two clarifications keep the construction honest. The register is not the orientation doublet: one carries an involution, the other carries nilpotents, and no isomorphism respects both algebras, so the crossing operator acts on the register tensored with the chain's own spaces. And the bidirectionality is in the move set, not in the tick: the ensemble sums both moves and has a stationary measure, while the ledger only grows and the reversed history remains inadmissible at its first event.

\end{theorem}

\begin{theorem}[{\texorpdfstring{BM.121 — The zig-zag world line: the even resolvent, the round-trip weights, and no backward leg}{BM.121 - The zig-zag world line: the even resolvent, the round-trip weights, and no backward leg}}]\label{tag:BM-121}
\textbf{Thm.}

%dsp-level:0
\[
(1)\quad\textbf{the vacuum-to-vacuum matrix elements are even:}
\]

%dsp-level:0
\[
\qquad \bigl(\hat U^{2m}\bigr)_{\varnothing\varnothing}=(x_+x_-)^{m},
\]

%dsp-level:0
\[
\qquad \bigl(\hat U^{2m+1}\bigr)_{\varnothing\varnothing}=0
\]

%dsp-level:0
\[
\qquad [\mathrm{FT.108}(5)\ \mathrm{Pf};\ \text{machine: coefficient-wise in}\ \mathbb Q[w,x_+,x_-]/(w^{12})]:
\]

%dsp-level:0
\[
\textbf{the mechanism:}\quad\text{a return to}\ \varnothing\ \text{requires an equal number of}\ \hat J_+\ \text{and}\ \hat J_-\ \text{moves, and}\ \hat J_\pm^2
\]

%dsp-level:0
\[
= 0\ \text{forces strict alternation}\ [\mathrm{BM.120}(1);\ \text{machine}];
\]

%dsp-level:0
\[
(2)\quad\textbf{the resolvent is the zig-zag sum:}
\]

%dsp-level:0
\[
\qquad \boxed{\ \sum_{n\ge0}\langle\varnothing|\hat U^{n}|\varnothing\rangle\,w^{n}\ =\ \sum_{m\ge0}\bigl(w^2x_+x_-\bigr)^{m}\ =\ \frac{1}{1-w^2x_+x_-}\ \ \text{in}\ \mathbb Q[w,x_+,x_-]/(w^{12})\ }
\]

%dsp-level:0
\[
[\mathrm{FT.108}(5);\ \mathrm{BM.23}(7);\ \mathrm{BM.44}(1);\ \text{machine:}\ (1-w^2x_+x_-)\cdot\sum\equiv1\ \text{coefficient-wise at five exact values of}\ x_+x_-\in\{0,\tfrac3{16},\tfrac14,\tfrac12,1\}];
\]

%dsp-level:0
\[
\textbf{every odd coefficient is exactly zero:}
\]

%dsp-level:0
\[
\qquad [w^{2m+1}]\,\frac1{1-w^2x_+x_-}=0\ \ \forall m
\]

%dsp-level:0
\[
\qquad [\text{machine: the coefficient census to}\ w^{11}];
\]

%dsp-level:0
\[
(3)\quad\textbf{the zig-zag is a sum over round trips, and every leg is forward:}
\]

%dsp-level:0
\[
\qquad \textbf{the}\ m\textbf{-th term}\ (w^{2}x_+x_-)^{m}\ \textbf{is}\ m\ \textbf{create--propagate--annihilate round trips}:
\]

%dsp-level:0
\[
\varnothing\ \xrightarrow{\ \hat J_+\ }\ \{\mathcal R\}\ \xrightarrow{\ \hat J_-\ }\ \varnothing\ \xrightarrow{\ \hat J_+\ }\ \{\mathcal R\}\ \xrightarrow{\ \hat J_-\ }\ \varnothing\ \cdots
\]

%dsp-level:0
\[
\qquad \text{with each arrow a step of the ledger order}\ \tau_n\to\tau_{n+1}
\]

%dsp-level:0
\[
\qquad [\mathrm{A6};\ \mathrm{D1.8}]:
\]

%dsp-level:0
\[
\boxed{\ \textbf{zig-zag}\ :=\ \text{the}\ m\textbf{-round-trip word in}\ \{\hat J_+,\hat J_-\},\ \ \text{weight}\ (x_+x_-)^{m},\ \ \text{with}\ \neg\textbf{a backward leg anywhere in it}\ }
\]

%dsp-level:0
\[
(4)\quad\textbf{the exact weight ladder:}\qquad\mathcal W_m=(x_+x_-)^{m}=\bigl[4\varepsilon_q(1-4\varepsilon_q)\bigr]^{m}\qquad[\text{machine}]:
\]

%dsp-level:0
\[
\begin{array}{c|c|c|c}\varepsilon_q & x_+x_- & \mathcal W_1 & \mathcal W_3\\\hline0 & 0 & 0 & 0\\\tfrac1{16} & \tfrac3{16} & \tfrac3{16} & \tfrac{27}{4096}\\\tfrac18 & \tfrac14 & \tfrac14 & \tfrac1{64}\\\tfrac3{16} & \tfrac3{16} & \tfrac3{16} & \tfrac{27}{4096}\\\tfrac14 & 0 & 0 & 0\end{array}\qquad[\text{machine: exact rationals}];
\]

%dsp-level:0
\[
\boxed{\ \mathcal W_m\ \textbf{is maximal at}\ \varepsilon_q=\tfrac18\ \textbf{for every}\ m,\ \textbf{and}\ \mathcal W_m\le4^{-m}\ \textbf{with equality exactly there}\ }
\]

%dsp-level:0
\[
\qquad [\mathrm{BM.120}(5);\ \text{machine: the monotone power of the parabola}];
\]

%dsp-level:0
\[
\textbf{the two endpoints annihilate the whole ladder:}
\]

%dsp-level:0
\[
\qquad \varepsilon_q=0\implies\mathcal W_m=0\ \forall m\ge1;
\]

%dsp-level:0
\[
\qquad \varepsilon_q\to\tfrac14^-\implies\mathcal W_m\to0\ \forall m\ge1
\]

%dsp-level:0
\[
\qquad [\mathrm{BM.109}(2);\ \text{machine}];
\]

%dsp-level:0
\[
(5)\quad\textbf{the pole and its enclosure:}
\]

%dsp-level:0
\[
\qquad w^2x_+x_-=1\implies|w|=\frac1{\sqrt{x_+x_-}}\ge2,
\]

%dsp-level:0
\[
\qquad \text{equality}\iff\varepsilon_q=\frac18
\]

%dsp-level:0
\[
\qquad [x_+x_-\le\tfrac14;\ \mathrm{BM.23}(7);\ \mathrm{BM.44}(1);\ \text{machine}]:
\]

%dsp-level:0
\[
\boxed{\ \textbf{the zig-zag series converges for}\ |w|<2\ \textbf{at every admissible threshold, and its pole never enters the unit disc}\ }
\]

%dsp-level:0
\[
(6)\quad\textbf{no backward leg, on three independent grounds:}\qquad(\mathrm i)\ \ \neg\exists\,\tau^{-1}\qquad[\mathrm{A6}];
\]

%dsp-level:0
\[
(\mathrm{ii})\quad\text{the pure order reversal of every completed nonempty admissible history leads with}\ J^-\ \text{and is inadmissible at its first event}
\]

%dsp-level:0
\[
\qquad [\mathrm{FT.144}(4)(\mathrm{ii});\ \text{machine: exhaustive}];
\]

%dsp-level:0
\[
(\mathrm{iii})\quad\Theta_{\mathfrak L}\ \textbf{fixes every completed alternating word pointwise:}
\]

%dsp-level:0
\[
\qquad \Theta_{\mathfrak L}\bigl((J^+J^-)^{m}\bigr)=(J^+J^-)^{m}\ \ \forall m
\]

%dsp-level:0
\[
\qquad [\mathrm{BM.25}(2);\ \text{machine: all}\ m\le6]:
\]

%dsp-level:0
\[
\boxed{\ \textbf{the reversed reading of a zig-zag is the same word}\ —\ \text{so the diagram's backward leg needs no inverse tick, because the corpus's own mirror already fixes it}\ }
\]

%dsp-level:0
\[
(7)\quad\textbf{the CP reading of the second leg:}
\]

%dsp-level:0
\[
\qquad \text{the member created by}\ \hat J_+\ \text{and the member annihilated by}\ \hat J_-\ \text{in one round trip are a CP-pair:}\quad\{c,\ c^{\mathrm{CP}}\},
\]

%dsp-level:0
\[
\qquad Q=q+(-q)=0
\]

%dsp-level:0
\[
\qquad [\mathrm{Def\text{-}BM.15}(3);\ \mathrm{BM.109}(1);\ \text{machine}]:
\]

%dsp-level:0
\[
\textbf{the second leg carries the conjugate configuration forward,}\ \neg\textbf{the original configuration backward:}
\]

%dsp-level:0
\[
\qquad |w(c^{\mathrm{CP}})|^2=|w(c)|^2,
\]

%dsp-level:0
\[
\qquad w(c^{\mathrm{CP}})=\mu_{\mathrm{CP}}\,w(c),
\]

%dsp-level:0
\[
\qquad \mu_{\mathrm{CP}}=e^{i\pi\Theta}=1\ \ \text{on}\ \Theta\in2\mathbb Z
\]

%dsp-level:0
\[
[\mathrm{BM.108}(2),(3);\ \mathrm{EM.24}(5);\ \text{machine}];
\]

%dsp-level:0
\[
\qquad \boxed{\ \textbf{the zig-zag is}\ \text{chain}\to\text{CP-conjugate chain, both forward}\ }
\]

%dsp-level:0
\[
\qquad [\mathrm{BM.129}(1)\text{'s boson row is the same statement read on the winding}];
\]

%dsp-level:0
\[
(8)\quad\textbf{the ledger-count row:}\qquad\bigl|\mathfrak L(\tau_{n+1})\bigr|-\bigl|\mathfrak L(\tau_n)\bigr|=n_{\mathrm{real\_out}}(\tau_n)\ge0\qquad[\mathrm{A6}]:
\]

%dsp-level:0
\[
\text{a round trip returns}\ n_{\mathcal R}\ \text{to its starting value while the record grows by the events it logged}
\]

%dsp-level:0
\[
\qquad [\text{machine:}\ \hat J_+\hat J_-=P_{\mathcal V}\ \text{restores the register,}\ \mathrm{A6}\ \text{keeps the record}]:
\]

%dsp-level:0
\[
\boxed{\ \textbf{the zig-zag is closed on the register and non-decreasing on the record}\ —\ \text{which is exactly what a forward-only ledger can support}\ }
\]

%dsp-level:0
\[
(9)\quad\textbf{status:}\qquad\textbf{derived:}\ (1)\text{--}(8);
\]

%dsp-level:0
\[
\textbf{re-scoped at its own locus,}\ \neg\textbf{retracted:}\quad\mathrm{BM.109}(2)\text{'s}\ \mathcal W_{\text{pair}}
\]

%dsp-level:0
\[
= x_+x_-\ \textbf{is the round-trip weight}\ —\ \text{the one-way weights are}\ x_+\ \text{and}\ x_-\ \text{separately}\ [(4);\ \mathrm{BM.120}(5)];
\]

%dsp-level:0
\[
\neg\text{a declaration},
\]

%dsp-level:0
\[
\quad\neg\text{a P4 addition},
\]

%dsp-level:0
\[
\quad\neg\text{a new dimensionful scale}\ [\mathrm{P1}],
\]

%dsp-level:0
\[
\quad\neg\text{a measured value}\ [\mathrm{P5}],
\]

%dsp-level:0
\[
\quad\neg\text{a species identification}\ [\mathrm{Def\text{-}BM.3}(7)],
\]

%dsp-level:0
\[
\quad\neg\text{a cosmological coordinate}\ [\mathrm{Def\text{-}BM.17}(6)(\mathrm i)];
\]

%dsp-level:0
\[
\neg\text{a world line in a spacetime is constructed: the zig-zag is a word in two moves on a finite register}\ [\mathrm{Def\text{-}BM.18}(5),(6)];
\]

\textbf{Req.} A6, A7, D1.2, D1.8 (Part 1); P1, P4, P5 (Parts 1–2); Def-H.32, FT.108(5), FT.144(4) (Part 3); Def-BM.3(7), Def-BM.15(3), Def-BM.17(6), Def-BM.18(5), Def-BM.18(6), BM.23(7), BM.25(2), BM.44(1), BM.108(2), BM.108(3), BM.109(1), BM.109(2), BM.120(1), BM.120(5), BM.129(1), EM.24(5) (this part); exact truncated power-series arithmetic in $\mathbb Q[w,x_+,x_-]/(w^{12})$, exact rational weight ladders, exhaustive ledger-record enumeration [math; machine: \texttt{checks/\allowbreak{}verify\_\allowbreak{}H\_\allowbreak{}janus\_\allowbreak{}tunnel\_\allowbreak{}concrete\_\allowbreak{}spin.py} §B].

\textbf{Pf.}

%dsp-level:0
\[
(1):\ \mathrm{FT.108}(5)\text{'s Pf row re-cited};
\]

%dsp-level:0
\[
\qquad \text{strict alternation}: \hat J_+^2
\]

%dsp-level:0
\[
= \hat J_-^2
\]

%dsp-level:0
\[
= 0\ [\mathrm{BM.120}(1);\ \text{machine}]\implies\text{a word returning to}\ \varnothing\ \text{has equal counts of the two moves}
\]

%dsp-level:0
\[
(2):\ (1-w^2x_+x_-)\sum_{m=0}^{5}(w^2x_+x_-)^m
\]

%dsp-level:0
\[
= 1-(w^2x_+x_-)^6
\]

%dsp-level:0
\[
\equiv 1\ \bmod\ w^{12}\ [\text{machine: coefficient-wise at five exact values of}\ x_+x_-];
\]

%dsp-level:0
\[
\qquad [w^{2m+1}]=0\ \forall m\ [\text{machine: the census to}\ w^{11}]
\]

%dsp-level:0
\[
(3):\ \text{each factor}\ w^2x_+x_-\ \text{is one}\ \hat J_+\hat J_-\ \text{pair}\ [(1);\ \mathrm{Def\text{-}BM.18}(5)];
\]

%dsp-level:0
\[
\qquad \text{each move is a ledger step}\ [\mathrm{D1.2};\ \mathrm{A6}]
\]

%dsp-level:0
\[
(4):\ \mathcal W_m=(x_+x_-)^m\ [\text{machine: exact rationals at the five thresholds}];
\]

%dsp-level:0
\[
\qquad x_+x_-\le\tfrac14\ [\mathrm{BM.120}(5)]\implies(x_+x_-)^m\le4^{-m}\ \text{with equality}\iff x_+=x_-\ [\text{machine: monotone power}]
\]

%dsp-level:0
\[
\varepsilon_q=0\implies x_+x_-=0\implies0^m=0\ \forall m\ge1;\qquad\varepsilon_q\to\tfrac14^-\implies x_+\to0\implies x_+x_-\to0\ [\text{machine}]
\]

%dsp-level:0
\[
(5):\ |w|^2x_+x_-=1\implies|w|^2=1/(x_+x_-)\ge4\implies|w|\ge2\ [\text{machine:}\ x_+x_-\le\tfrac14];
\]

%dsp-level:0
\[
\qquad \text{equality at}\ \varepsilon_q=\tfrac18\ [(4)]
\]

%dsp-level:0
\[
(6):\ (\mathrm i)\ \mathrm{A6}\ \text{verbatim};
\]

%dsp-level:0
\[
\qquad (\mathrm{ii})\ \mathrm{FT.144}(4)(\mathrm{ii})\ \text{re-cited}\ [\text{machine: exhaustive}];
\]

%dsp-level:0
\[
\qquad (\mathrm{iii})\ \Theta_{\mathfrak L}((J^+J^-)^m)=(J^-J^+)^m\ \text{reversed}\ =(J^+J^-)^m\ [\mathrm{BM.25}(2);\ \text{machine: all}\ m\le6]
\]

%dsp-level:0
\[
(7):\ \mathrm{Def\text{-}BM.15}(3),\ \mathrm{BM.109}(1)\ \text{re-cited};\qquad\mu_{\mathrm{CP}}=e^{i\pi\Theta}=1\ \text{at}\ \Theta\in2\mathbb Z\ [\mathrm{BM.108}(3);\ \text{machine}]
\]

%dsp-level:0
\[
(8):\ \mathrm{A6}\text{'s increment row re-cited};
\]

%dsp-level:0
\[
\qquad \hat J_+\hat J_-=P_{\mathcal V}\ \text{restores the register}\ [\mathrm{BM.120}(1);\ \text{machine}];
\]

%dsp-level:0
\[
\qquad (9):\ \text{the standing-rule scans}\ [\text{machine}].\ \blacksquare
\]

\textbf{Phys.} The zig-zag is not drawn, it is summed, and the sum turns out to be even. A return to the vacuum needs as many crystallizations as dissolutions, and since neither move can be applied twice in a row the words strictly alternate — so every odd matrix element vanishes identically and the vacuum-to-vacuum resolvent is a geometric series in the square of the fugacity with ratio the round-trip weight. That is the whole content of the zig-zag: term by term it is a create-propagate-annihilate excursion, weighted by the parabola to the power of the number of excursions, maximal at one eighth of the threshold and exactly zero at both endpoints, so the entire ladder is annihilated by a non-flickering ledger. The pole sits at fugacity two or beyond and never enters the unit disc. Three independent rows say why no leg runs backward: the tick has no inverse, the pure reversal of a completed history is inadmissible at its first event, and — the sharpest of the three — the corpus's own record involution fixes every completed alternating word pointwise, so the reversed reading of a zig-zag is literally the same word. What looks like a backward-propagating member in a diagram is here the CP-conjugate configuration carried forward, with equal modulus and a phase that is exactly one on the materialized locus. And the ledger keeps its own accounting: a round trip restores the register while the record grows by the events it logged, which is precisely the one thing a forward-only ledger can support.

\end{theorem}

\begin{theorem}[{\texorpdfstring{BM.122 — The pair channel: charge balance, the exact threshold, the CP structure, and the two-photon bar}{BM.122 - The pair channel: charge balance, the exact threshold, the CP structure, and the two-photon bar}}]\label{tag:BM-122}
\textbf{Thm.}

%dsp-level:0
\[
(1)\quad\textbf{the channel is charge-balanced exactly:}
\]

%dsp-level:0
\[
\qquad Q_{\text{in}}=0\ \ (\text{the}\ \mathcal V\ \text{sector carries no registered charge}\ [\mathrm{D1.1};\ \mathrm{FT.144}(4)]),
\]

%dsp-level:0
\[
\qquad Q_{\text{out}}=q+(-q)=0
\]

%dsp-level:0
\[
\qquad [\mathrm{Def\text{-}BM.15}(3);\ \mathrm{BM.109}(1);\ \text{machine: exhaustive over}\ q\in\{-1,0,+1\}]:
\]

%dsp-level:0
\[
\boxed{\ Q_{\mathrm{EM}}=\sum_{c\in U}q_{\mathrm{EM}}(c)=0\ \ \text{survives the channel, in every valid universe}\ }
\]

%dsp-level:0
\[
\qquad [\mathrm{EM.11};\ \mathrm{D2.11};\ \mathrm{BM.23}(4)];
\]

%dsp-level:0
\[
\textbf{the compensator is unique:}\qquad\#\{q\in\{-1,0,+1\}:\ q=-q_{\text{in}}\}=1\qquad[\mathrm{BM.23}(4);\ \mathrm{BM.108}(5);\ \text{machine}];
\]

%dsp-level:0
\[
(2)\quad\textbf{the exact threshold:}
\]

%dsp-level:0
\[
\qquad \mathcal T_0(\text{pair})=\frac{E^{\text{pair}}_{\mathrm{int}}\,\mathfrak h}{\ln2}=\frac{2E_{\mathrm{int}}(c)\,\mathfrak h}{\ln2}
\]

%dsp-level:0
\[
\qquad [\mathrm{Def\text{-}BM.17}(3);\ \mathrm{BM.109}(1)]:
\]

%dsp-level:0
\[
\textbf{at the minimal materialized floor}\ E_{\mathrm{int}}(c)=2\ [\mathrm{FT.28};\ \mathrm{BM.13}(3)]:
\]

%dsp-level:0
\[
\qquad \boxed{\ \mathcal T_0(\text{pair})\ =\ \frac{4\,\mathfrak h}{\ln2},\qquad\textbf{materialization}\iff NE\ \ge\ \frac{4\mathfrak h}{\ln2}\ }
\]

%dsp-level:0
\[
[\text{machine: one Laurent monomial};\ \mathrm{BM.115}(2)(\mathrm i)\text{'s}\ n{=}2\ \textbf{minimum, re-derived}];
\]

%dsp-level:0
\[
\qquad \mathcal T(\text{pair};N)=\frac{4\mathfrak h}{N\ln2},
\]

%dsp-level:0
\[
\qquad \mathcal T(\text{pair};r_s,N)=\frac{2r_s}N\ \ \text{at}\ r_s=2\mathfrak h
\]

%dsp-level:0
\[
\qquad [\mathrm{BM.115}(1),(4);\ \text{machine}];
\]

%dsp-level:0
\[
\textbf{the general ladder:}
\]

%dsp-level:0
\[
\qquad \mathcal T_0(\text{pair})\in\Bigl\{\frac{4\mathfrak h}{\ln2},\ \dots,\ \frac{120\mathfrak h}{\ln2}\Bigr\}\ \ \text{on the computed cells, span ratio exactly}\ 30
\]

%dsp-level:0
\[
\qquad [\mathrm{BM.111}(5);\ \mathrm{BM.115}(2)(\mathrm i);\ \text{machine}];
\]

%dsp-level:0
\[
(3)\quad\textbf{the channel weight:}
\]

%dsp-level:0
\[
\qquad \boxed{\ \mathcal W_{\text{channel}}\ =\ x_+^{\,2}\cdot\chi_{\{NE\ge\mathcal T_0\}}\ =\ (1-4\varepsilon_q)^2\cdot\chi_{\{NE\ge\mathcal T_0\}}\ }
\]

%dsp-level:0
\[
\qquad [\text{two one-way crystallization moves};\ \mathrm{Def\text{-}BM.18}(5);\ \text{machine}]:
\]

%dsp-level:0
\[
\textbf{the round-trip completion:}
\]

%dsp-level:0
\[
\qquad \mathcal W_{\text{round}}=x_+x_-=4\varepsilon_q(1-4\varepsilon_q)\le\tfrac14
\]

%dsp-level:0
\[
\qquad [\mathrm{BM.109}(2);\ \mathrm{BM.120}(5);\ \text{machine}];
\]

%dsp-level:0
\[
\boxed{\ \mathcal W_{\text{channel}}>0\iff\varepsilon_q<\tfrac14\ \wedge\ NE\ge\mathcal T_0;\qquad\mathcal W_{\text{channel}}=1\iff\varepsilon_q=0\ \wedge\ NE\ge\mathcal T_0\ }
\]

%dsp-level:0
\[
\textbf{the two limits are opposite and both exact:}\quad\text{at}\ \varepsilon_q
\]

%dsp-level:0
\[
= 0\ \textbf{creation is certain and annihilation impossible}\ (x_+=1,\ x_-=0);
\]

%dsp-level:0
\[
\quad\text{at}\ \varepsilon_q\to\tfrac14^-\ \textbf{creation vanishes}\ (x_+\to0)
\]

%dsp-level:0
\[
\qquad [\text{machine}];
\]

%dsp-level:0
\[
(4)\quad\textbf{the CP structure of the channel — the three derived legs:}\qquad\begin{array}{l|l|l}\text{the inversion} & \text{the derived operator} & \text{the exact rows}\\\hline\text{charge} & \widehat V=E_+\otimes P_L\ \ [\mathrm{BM.76}(1)] & \widehat V^2=0,\ \operatorname{rk}\widehat V=1,\ \dim\ker\widehat V=3,\ [W_3\otimes I_\omega,\widehat V]=\widehat V\ \ [\mathrm{BM.76}(2)\text{--}(4);\ \text{machine over}\ \mathbb Q(i)]\\\text{parity/orientation} & \mu_{\mathrm{CP}}=e^{i\pi\Theta}\ \text{and}\ \hat\omega\ \ [\mathrm{BM.108}(3);\ \mathrm{FT.122}(4)] & \hat\omega\hat\chi\hat\omega^{-1}=-\hat\chi,\ \hat\omega P_L\hat\omega^{-1}=P_R,\ \mu_{\mathrm{CP}}=1\ \text{on}\ \Theta\in2\mathbb Z\ \ [\text{machine}]\\\text{energy} & \text{the local scale against}\ \mathcal T_0\ \ [\mathrm{Def\text{-}BM.17}(3)] & NE\ge\mathcal T_0,\ \ \mathcal T_0=E^{\text{pair}}_{\mathrm{int}}\mathfrak h/\ln2\ \ [\mathrm{BM.115}(1);\ \text{machine}]\end{array}
\]

%dsp-level:0
\[
\textbf{the image vertex:}
\]

%dsp-level:0
\[
\qquad (I\otimes\hat\omega)\,\widehat V\,(I\otimes\hat\omega)^{-1}=E_+\otimes P_R
\]

%dsp-level:0
\[
\qquad [\mathrm{BM.76}(5);\ \text{machine}]:
\]

%dsp-level:0
\[
\qquad \boxed{\ \textbf{the step fires from the left ray and its image from the right ray}\ —\ \text{the two members of the pair occupy opposite rays}\ }
\]

%dsp-level:0
\[
[\mathrm{BM.73}(3);\ \mathrm{BM.74}:\ \text{the two emitted members of a vertex occupy opposite rays};\ \mathrm{A\text{-}BM.7}(b);\ \text{machine}];
\]

%dsp-level:0
\[
(5)\quad\textbf{the nilpotency is the channel's selection rule:}
\]

%dsp-level:0
\[
\qquad \widehat V^2=0\implies\boxed{\ \neg\exists\ \text{an iterated charged step inside one doublet}\ }
\]

%dsp-level:0
\[
\qquad [\mathrm{BM.76}(2);\ \mathrm{BM.23}(3)\text{'s ternary overflow}]:
\]

%dsp-level:0
\[
\textbf{the consequence for the channel:}\quad\text{one charge inversion per member, exactly}\ —\ \text{the pair is}\ (q,\ -q)\ \text{and}\ \neg(q,\ q)\ \text{or}\ (q,\ q')\ \text{with}\ |q-q'|\ge2
\]

%dsp-level:0
\[
\qquad [\text{machine: the}\ 3^2=9\ \text{charge pairs against the balanced subset}];
\]

%dsp-level:0
\[
\#\{(q_1,q_2)\in\{-1,0,+1\}^2:\ q_1+q_2=0\}=3,\qquad\#\{(q_1,q_2):\ q_1+q_2\ne0\}=6\qquad[\text{machine}];
\]

%dsp-level:0
\[
(6)\quad\textbf{the two-photon bar (four independent grounds):}\qquad\textbf{the supply leg of}\ \gamma+\gamma\to\text{pair}\ \textbf{is}\ \neg\textbf{constructed:}
\]

%dsp-level:0
\[
(\mathrm O1)\quad\textbf{no gauged chiral coupling:}\ \text{the gauge blocks act on the}\ 16\text{-dimensional pair-channel register and on slot algebras;}\ \neg\text{a block couples to a chiral matter multiplet}
\]

%dsp-level:0
\[
\qquad [\mathrm{FT.80}(7)(\mathrm a);\ \mathrm{CT0.12};\ \mathrm{CT3.20}];
\]

%dsp-level:0
\[
(\mathrm O2)\quad\textbf{the transit alternative is not an inventory candidate:}\ \text{the gauge coupling of the charge read is deficit}\ 1'\ \text{content}
\]

%dsp-level:0
\[
\qquad [\mathrm{BM.41}(3);\ \mathrm{Def\text{-}BM.8}(8);\ \mathrm{FT.125}(4)];
\]

%dsp-level:0
\[
(\mathrm O3)\quad\textbf{no momentum-space resolvent:}\ \text{plane-wave boundary-channel states and the}\ \Sigma\text{-resolvent are an open interface, blocked pending their own declaration}
\]

%dsp-level:0
\[
\qquad [\mathrm{BM.6}(6)];
\]

%dsp-level:0
\[
(\mathrm O4)\quad\textbf{the transit object carries no winding carrier:}\ \text{the}\ \mathrm{D3.7}\ \text{trivial chain is excluded from}\ \mathrm{Def\text{-}H.38}\text{'s scope, so}\ \neg\text{a}\ j
\]

%dsp-level:0
\[
= |k|/2\ \text{read and}\ \neg\text{an}\ \varepsilon\ \text{read attach to it; its polarization count is}\ N_{\mathrm{pol}}
\]

%dsp-level:0
\[
= 2\ \text{and its helicity is a transit-side read with no identification}
\]

%dsp-level:0
\[
\qquad [\mathrm{D3.7};\ \mathrm{EM.8};\ \mathrm{Def\text{-}H.38}\ \text{Scope}];
\]

%dsp-level:0
\[
\boxed{\ \neg\exists\ \text{a two-transit-quantum supply vertex in the inventory}\implies\textbf{the pair channel's amplitude is}\ \neg\textbf{formed, while the channel itself is}\ }
\]

%dsp-level:0
\[
(7)\quad\textbf{what is formed and what is not:}\qquad\begin{array}{l|l}\textbf{formed} & \neg\textbf{formed}\\\hline\text{the charge balance}\ (1) & \text{a cross-section}\ [\mathrm{BM.6}(6);\ \mathrm{Def\text{-}BM.18}(6)]\\\text{the exact threshold}\ \mathcal T_0\ (2) & \text{an S-matrix element}\ [\mathrm{BM.6}(6)]\\\text{the channel weight}\ x_+^2\chi\ (3) & \text{a Lorentz-invariant phase space}\ [\mathrm{Def\text{-}BM.18}(6)]\\\text{the three CP legs}\ (4) & \text{the two-quanta supply vertex}\ (6)\\\text{the selection rule}\ \widehat V^2=0\ (5) & \text{a measured rate or branching fraction}\ [\mathrm{P5};\ \mathrm{BM.44}(5)]\end{array}
\]

%dsp-level:0
\[
(8)\quad\textbf{status:}\qquad\textbf{derived:}\ (1)\text{--}(5);\qquad\textbf{refuted as formed, with four grounds:}\ (6);
\]

%dsp-level:0
\[
\neg\text{a declaration},
\]

%dsp-level:0
\[
\quad\neg\text{a P4 addition},
\]

%dsp-level:0
\[
\quad\neg\text{a new dimensionful scale}\ [\mathrm{P1}],
\]

%dsp-level:0
\[
\quad\neg\text{a measured value}\ [\mathrm{P5}],
\]

%dsp-level:0
\[
\quad\neg\text{a numerical coupling}\ [\mathrm{EM.25}(7)],
\]

%dsp-level:0
\[
\quad\neg\text{a species identification}\ [\mathrm{Def\text{-}BM.3}(7)];
\]

%dsp-level:0
\[
\neg\text{a right-ray matter line is added: the image vertex of}\ (4)\ \text{is the}\ \hat\omega\text{-image of the same single-arrow locus, in the representation}\ \{0,-1\}\ [\mathrm{BM.96}(4),(6);\ \mathrm{BM.73}(2)(b);\ \mathrm{A\text{-}BM.7}];
\]

\textbf{Req.} A4, A6, D0.8, D1.1, D1.2, D1.8, D2.11, D3.7 (Parts 1–2); CT0.12, CT3.20, EM.8, EM.11, EM.24(5), EM.25(7) (Part 1, Part 2); P1, P4, P5 (Parts 1–2); Def-H.4, Def-H.38, FT.28, FT.44(4), FT.80(7), FT.122(4), FT.125(4), FT.144(4) (Part 3); Def-BM.3(7), Def-BM.8(8), Def-BM.15(3), Def-BM.17(3), Def-BM.18(5), Def-BM.18(6), BM.6(6), BM.13(3), BM.23(3), BM.23(4), BM.41(3), BM.44(5), BM.73(2), BM.73(3), BM.74, BM.76(1), BM.76(2), BM.76(3), BM.76(4), BM.76(5), BM.108(3), BM.108(5), BM.109(1), BM.109(2), BM.111(5), BM.115(1), BM.115(2), BM.115(4), BM.120(5), A-BM.7 (this part); exact $\mathbb Q(i)$ matrix algebra, exhaustive ternary charge-pair enumeration, Laurent-monomial assembly, exact rational weight ladders [math; machine: \texttt{checks/\allowbreak{}verify\_\allowbreak{}H\_\allowbreak{}janus\_\allowbreak{}tunnel\_\allowbreak{}concrete\_\allowbreak{}spin.py} §B].

\textbf{Pf.}

%dsp-level:0
\[
(1):\ q+(-q)=0\ [\text{machine: three cases}];\qquad\mathrm{EM.11},\ \mathrm{D2.11}\ \text{re-cited};\qquad\#\{q:q=-q_{\text{in}}\}=1\ [\text{machine}]
\]

%dsp-level:0
\[
(2):\ E^{\text{pair}}_{\mathrm{int}}=2E_{\mathrm{int}}(c)\ [\mathrm{BM.109}(1)];
\]

%dsp-level:0
\[
\qquad \mathcal T_0=E^{\text{pair}}_{\mathrm{int}}\mathfrak h/\ln2\ [\mathrm{Def\text{-}BM.17}(3)];
\]

%dsp-level:0
\[
\qquad E_{\mathrm{int}}=2\implies\mathcal T_0=4\mathfrak h/\ln2\ [\text{machine}]
\]

%dsp-level:0
\[
\mathcal T=2r_s/N\ \text{at}\ r_s=\tfrac12\mathfrak h E^{\text{pair}}_{\mathrm{int}}=2\mathfrak h\ [\mathrm{BM.115}(4);\ \text{machine}];
\]

%dsp-level:0
\[
\qquad \text{the ladder and its span ratio}\ 30\ [\mathrm{BM.115}(2)(\mathrm i);\ \text{machine}]
\]

%dsp-level:0
\[
(3):\ \text{two one-way crystallization moves, each of weight}\ x_+\ [\mathrm{Def\text{-}BM.18}(5);\ \mathrm{BM.120}(5)];
\]

%dsp-level:0
\[
\qquad x_+^2=(1-4\varepsilon_q)^2\ [\text{machine}]
\]

%dsp-level:0
\[
\varepsilon_q=0\implies x_+=1,\ x_-=0;
\]

%dsp-level:0
\[
\qquad \varepsilon_q\to\tfrac14^-\implies x_+\to0\ [\text{machine: exact at the two endpoints of the admissible range}\ [\mathrm{CT0.8}]]
\]

%dsp-level:0
\[
(4):\ \widehat V=E_+\otimes P_L;
\]

%dsp-level:0
\[
\quad E_+^2=0\implies\widehat V^2=0;\quad\operatorname{rk}\widehat V=\operatorname{rk}E_+\cdot\operatorname{rk}P_L=1;
\]

%dsp-level:0
\[
\quad [W_3\otimes I_\omega,\widehat V]=\widehat V\ \text{on}\ [\sigma_3,E_+]=2E_+,\ W_3=\tfrac12\sigma_3\ [\mathrm{BM.76}(1)\text{--}(4);\ \text{machine over}\ \mathbb Q(i)]
\]

%dsp-level:0
\[
(I\otimes\hat\omega)\widehat V(I\otimes\hat\omega)^{-1}=E_+\otimes\hat\omega P_L\hat\omega^{-1}=E_+\otimes P_R\ [\mathrm{FT.122}(4);\ \mathrm{BM.76}(5);\ \text{machine}]
\]

%dsp-level:0
\[
(5):\ \widehat V^2=0\ [\mathrm{BM.76}(2)];
\]

%dsp-level:0
\[
\qquad \text{the balanced subset of the}\ 9\ \text{charge pairs is}\ \{(-1,+1),(0,0),(+1,-1)\},
\]

%dsp-level:0
\[
\ \#\
\]

%dsp-level:0
\[
= 3\ [\text{machine: exhaustive}]
\]

%dsp-level:0
\[
(6):\ \mathrm{FT.80}(7)(\mathrm a),
\]

%dsp-level:0
\[
\ \mathrm{BM.41}(3),
\]

%dsp-level:0
\[
\ \mathrm{Def\text{-}BM.8}(8),
\]

%dsp-level:0
\[
\ \mathrm{BM.6}(6),
\]

%dsp-level:0
\[
\ \mathrm{D3.7},
\]

%dsp-level:0
\[
\ \mathrm{EM.8},
\]

%dsp-level:0
\[
\ \mathrm{Def\text{-}H.38}\ \text{Scope rows re-cited}\ [\text{machine: the four-ground census}];
\]

%dsp-level:0
\[
\qquad N_{\mathrm{pol}}=(4-1)-1=2\ [\mathrm{EM.8};\ \text{machine}]
\]

%dsp-level:0
\[
(7),(8):\ \text{the two-column assembly and the standing-rule scans}\ [\text{machine}].\ \blacksquare
\]

\textbf{Phys.} The channel that opens is exactly as constrained as the ledger can make it. Charge balances because the pair is a chain and its conjugate, the compensating value is unique in the ternary alphabet, and the nilpotent weak step allows precisely one charge inversion per member — so of the nine ordered charge pairs only the three balanced ones are formed. The threshold is one Laurent monomial: twice the constituent edge count times the single scale over the log of two, which at the digon floor is four units, the very minimum of the two-chain ladder the resonance tier already computed. The weight is the square of the crystallization weight times the indicator that the local scale has reached that threshold, and its two limits are exact opposites — at the sharp threshold creation is certain and annihilation impossible, while near full dissolution creation itself vanishes. The CP structure decomposes into three derived legs: a nilpotent step operator that fires from the occupied ray alone and whose orientation image fires from the other, the CP transport phase that is exactly one on the materialized locus, and the local scale against the threshold. What does not exist is the supply vertex. Four independent grounds bar it: nothing couples a gauge block to a chiral matter multiplet, the transit alternative is deficit-one-prime content, there is no momentum-space resolvent, and the transit object is excluded from the winding carrier's scope so neither a spin label nor a statistics sign attaches to it. So the channel is derived and the amplitude is not, and the tier records both halves rather than promoting one into the other.

\end{theorem}

\begin{theorem}[{\texorpdfstring{BM.123 — The annihilation channel: the reverse move, the entropy destination, and the closure of the cross-section residue}{BM.123 - The annihilation channel: the reverse move, the entropy destination, and the closure of the cross-section residue}}]\label{tag:BM-123}
\textbf{Thm.}

%dsp-level:0
\[
(1)\quad\textbf{the reverse move is the same algebra:}
\]

%dsp-level:0
\[
\qquad \text{annihilation}\ :=\ \hat J_-\ \text{on each of the pair's two registers},
\]

%dsp-level:0
\[
\qquad \text{weight}\ x_-^{\,2}=(4\varepsilon_q)^2
\]

%dsp-level:0
\[
\qquad [\mathrm{Def\text{-}BM.18}(5);\ \mathrm{BM.120}(5);\ \text{machine}]:
\]

%dsp-level:0
\[
\boxed{\ \mathcal W_{\text{ann}}\ =\ x_-^{\,2}\cdot\chi_{\{\text{the pair is materialized}\}}\ =\ 16\varepsilon_q^2\cdot\chi_{\{NE\ge\mathcal T_0\}}\ }
\]

%dsp-level:0
\[
\qquad [\text{machine: one monomial in}\ \varepsilon_q];
\]

%dsp-level:0
\[
\textbf{the charge balance:}
\]

%dsp-level:0
\[
\qquad Q_{\text{in}}=q+(-q)=0,
\]

%dsp-level:0
\[
\qquad Q_{\text{out}}=0\ \ (\text{the}\ \mathcal V\ \text{sector is unregistered})
\]

%dsp-level:0
\[
\qquad [\mathrm{D1.1};\ \mathrm{EM.11};\ \text{machine}]:
\]

%dsp-level:0
\[
\qquad \boxed{\ \text{the balance holds in both directions}\ }
\]

%dsp-level:0
\[
(2)\quad\textbf{the two directions and their exact ratio:}
\]

%dsp-level:0
\[
\qquad \frac{\mathcal W_{\text{ann}}}{\mathcal W_{\text{channel}}}\ =\ \frac{x_-^2}{x_+^2}\ =\ \Bigl(\frac{4\varepsilon_q}{1-4\varepsilon_q}\Bigr)^{\!2}
\]

%dsp-level:0
\[
\qquad [\text{machine: exact rational function}]:
\]

%dsp-level:0
\[
\boxed{\ \frac{\mathcal W_{\text{ann}}}{\mathcal W_{\text{channel}}}=1\iff\varepsilon_q=\frac18,\qquad<1\iff\varepsilon_q<\frac18,\qquad>1\iff\varepsilon_q>\frac18\ }
\]

%dsp-level:0
\[
[\text{machine: the sign census on a}\ 256\text{-point rational grid}];
\]

%dsp-level:0
\[
\qquad \textbf{the ratio is the detailed-balance indicator:}\ \text{equality holds exactly at}\ \mathrm{BM.120}(5)\text{'s maximum of the round-trip weight};
\]

%dsp-level:0
\[
\textbf{the two endpoint rows:}
\]

%dsp-level:0
\[
\qquad \varepsilon_q=0:\ \ \frac{\mathcal W_{\text{ann}}}{\mathcal W_{\text{channel}}}=0\ \ (\text{nothing annihilates});
\]

%dsp-level:0
\[
\qquad \varepsilon_q\to\tfrac14^-:\ \ \mathcal W_{\text{channel}}\to0\ \ \text{and the ratio diverges}
\]

%dsp-level:0
\[
\qquad [\text{machine}];
\]

%dsp-level:0
\[
(3)\quad\textbf{the round trip closes the two directions:}
\]

%dsp-level:0
\[
\qquad \mathcal W_{\text{channel}}\cdot\mathcal W_{\text{ann}}\ \text{is}\ \neg\text{the round-trip weight};
\]

%dsp-level:0
\[
\qquad \textbf{the round trip is one creation and one annihilation:}
\]

%dsp-level:0
\[
\boxed{\ \mathcal W_{\text{round}}\ =\ x_+x_-\ =\ 4\varepsilon_q(1-4\varepsilon_q)\ \le\ \frac14\ }\qquad[\mathrm{BM.109}(2);\ \mathrm{BM.120}(5);\ \mathrm{BM.121}(4);\ \text{machine}];
\]

%dsp-level:0
\[
\textbf{the resolvent that sums them:}
\]

%dsp-level:0
\[
\qquad \frac1{1-w^2x_+x_-},
\]

%dsp-level:0
\[
\qquad \text{pole at}\ |w|\ge2
\]

%dsp-level:0
\[
\qquad [\mathrm{BM.121}(2),(5);\ \mathrm{FT.108}(5);\ \text{machine}];
\]

%dsp-level:0
\[
(4)\quad\textbf{where the content goes (the entropy destination):}
\]

%dsp-level:0
\[
\qquad \text{a dissolved pair returns to the}\ \mathcal V\ \text{sector:}\quad n_{\mathcal R}=0\implies\rho_{pq}=\rho_\circ=\tfrac14I,
\]

%dsp-level:0
\[
\qquad S_{\mathrm{raw}}(\rho_\circ)=\ln4=\max,
\]

%dsp-level:0
\[
\qquad \arg\max=\bigl\{\tfrac14I\bigr\}
\]

%dsp-level:0
\[
[\mathrm{FT.144}(4);\ \mathrm{CT0.1};\ \mathrm{CT0.2};\ \mathrm{A5};\ \text{machine}];\qquad E_{\mathrm{int}}(R)=|R|,\qquad S_{\mathfrak L}=n_{\mathcal R}\ln2\qquad[\mathrm{Def\text{-}BM.5}(4);\ \mathrm{CT1.29}]:
\]

%dsp-level:0
\[
\boxed{\ \textbf{annihilation returns the pair's registered content to the maximally raw-entropic virtual sector}\ —\ \neg\text{a radiation leg carries it, because}\ \neg\text{a supply vertex exists}\ [\mathrm{BM.122}(6)]\ }
\]

%dsp-level:0
\[
\textbf{the ledger row:}\qquad|E_{\mathcal R}|=0\implies\neg\exists\ \text{a fermion terminus}\qquad[\mathrm{D1.4};\ \mathrm{Def\text{-}H.4};\ \mathrm{FT.144}(4);\ \text{machine}];
\]

%dsp-level:0
\[
(5)\quad\textbf{the inter-sector row, re-scoped and not retracted:}
\]

%dsp-level:0
\[
\qquad \mathrm{BM.27}(7),
\]

%dsp-level:0
\[
\ \mathrm{BM.29}(5)\ \text{and}\ \mathrm{Def\text{-}BM.15}(5)\ \text{record}\ \neg\text{an annihilation channel,}\ \neg\text{a pair-production channel and}\ \neg\text{a cross-section between}\ \mathrm{EM}^+\ \text{and}\ \mathrm{EM}^-:
\]

%dsp-level:0
\[
\textbf{the frozen content stands:}
\]

%dsp-level:0
\[
\qquad \neg\exists\ \textbf{a direct}\ \mathrm{EM}^+\to\mathrm{EM}^-\ \textbf{operator};
\]

%dsp-level:0
\[
\qquad \text{the two sectors remain amplitude-equal configuration orbits of one involution}\ [\mathrm{BM.25}(4);\ \mathrm{BM.27}(6);\ |w(c^{\mathrm{CP}})|^2=|w(c)|^2];
\]

%dsp-level:0
\[
\textbf{what is constructed:}
\]

%dsp-level:0
\[
\qquad \boxed{\ \text{a}\ \textbf{mediated}\ \text{channel}\ —\ \mathcal J\text{'s two moves composed with the CP identification of}\ \mathrm{Def\text{-}BM.15}(2)\ —\ \text{connects the two sectors through the}\ \mathcal V\ \textbf{sector, not through a direct operator}\ }
\]

%dsp-level:0
\[
\textbf{the composition, exactly:}
\]

%dsp-level:0
\[
\qquad \varnothing\ \xrightarrow{\ \hat J_+\otimes\hat J_+\ }\ \{c,\ c^{\mathrm{CP}}\}\ \xrightarrow{\ \hat J_-\otimes\hat J_-\ }\ \varnothing,
\]

%dsp-level:0
\[
\qquad \text{weight}\ x_+^2x_-^2\ \ \text{for the full out-and-back on two registers}
\]

%dsp-level:0
\[
[\text{machine: the monomial}];\qquad\neg\text{a cross-section is attached to it}\ [\mathrm{Def\text{-}BM.18}(6);\ \mathrm{BM.6}(6)];
\]

%dsp-level:0
\[
(6)\quad\textbf{the closure of the cross-section residue:}
\]

%dsp-level:0
\[
\qquad \mathrm{BM.118}(7)\ \text{named}\ (\mathrm{R9})\ \text{as}\ \textit{the annihilation channel: an inter-sector transition operator between}\ \mathrm{EM}^+\ \text{and}\ \mathrm{EM}^-,
\]

%dsp-level:0
\[
\ \neg\text{in the inventory}:
\]

%dsp-level:0
\[
\boxed{\ (\mathrm{R9})\ \ \textbf{is CLOSED at the channel level:}\quad\text{a mediated annihilation channel is constructed, with an exact weight}\ x_-^2\chi,\ \text{an exact threshold}\ \mathcal T_0,\ \text{an exact charge balance and an exact entropy destination}\ }
\]

%dsp-level:0
\[
\textbf{the residue narrows and does not vanish:}\qquad\begin{array}{l|l}(\mathrm{R9}')\ \ \text{the annihilation cross-section} & \text{requires the momentum-space}\ \Sigma\text{-resolvent}\ [\mathrm{BM.6}(6)]\ \textbf{and}\ \text{the gauge coupling of the charge read}\ [\mathrm{Def\text{-}BM.8}(8);\ \mathrm{BM.122}(6)]\\(\mathrm{R9}'')\ \ \text{the radiation legs} & \neg\text{a transit-quantum supply vertex exists}\ [\mathrm{BM.122}(6)(\mathrm O1)\text{--}(\mathrm O4)]\end{array}
\]

%dsp-level:0
\[
(7)\quad\textbf{the arrow rows stand:}
\]

%dsp-level:0
\[
\qquad \text{the channel is forward-only:}\ \neg\exists\,\tau^{-1}\ [\mathrm{A6}];
\]

%dsp-level:0
\[
\qquad \neg\exists\ \text{an inverse decay or an inverse creation in the ledger}\ [\mathrm{BM.23}(9);\ \mathrm{BM.44}(6);\ \mathrm{FT.144}(2)]:
\]

%dsp-level:0
\[
\textbf{the reversibility is in the move set and the stationary measure,}\ \neg\textbf{in the tick}
\]

%dsp-level:0
\[
\qquad [\mathrm{BM.120}(3),(7);\ \text{machine}];
\]

%dsp-level:0
\[
\Gamma(\text{pair})=\Gamma_{\mathrm{tr}}(\varepsilon_q)\cdot f\cdot\frac{\varepsilon_{\mathrm{eff}}}{\varepsilon_{\mathrm{EM}}}\cdot\mathcal N\cdot\chi_{\{NE\ge\mathcal T_0\}}
\]

%dsp-level:0
\[
\qquad [\mathrm{Def\text{-}BM.17}(4);\ \mathrm{BM.116}(1):\ \text{the width function applies to the pair resonance unchanged}];
\]

%dsp-level:0
\[
(8)\quad\textbf{status:}
\]

%dsp-level:0
\[
\qquad \textbf{derived:}\ (1)\text{--}(5),\ (7);
\]

%dsp-level:0
\[
\qquad \textbf{closed at the channel level:}\ (6);
\]

%dsp-level:0
\[
\qquad \textbf{named and}\ \neg\textbf{filled:}\ (\mathrm{R9}'),\ (\mathrm{R9}'');
\]

%dsp-level:0
\[
\textbf{re-scoped at their own loci,}\ \neg\textbf{retracted:}\quad\mathrm{BM.27}(7),
\]

%dsp-level:0
\[
\ \mathrm{BM.29}(5),
\]

%dsp-level:0
\[
\ \mathrm{Def\text{-}BM.15}(5),
\]

%dsp-level:0
\[
\ \mathrm{BM.108}(7),
\]

%dsp-level:0
\[
\ \mathrm{BM.109}(5)\ \text{(each frozen on its own locus:}\ \neg\text{a}\ \textbf{direct}\ \text{inter-sector operator exists)};
\]

%dsp-level:0
\[
\neg\text{a declaration},
\]

%dsp-level:0
\[
\quad\neg\text{a P4 addition},
\]

%dsp-level:0
\[
\quad\neg\text{a new dimensionful scale}\ [\mathrm{P1}],
\]

%dsp-level:0
\[
\quad\neg\text{a measured value}\ [\mathrm{P5}],
\]

%dsp-level:0
\[
\quad\neg\text{a numerical coupling}\ [\mathrm{EM.25}(7)],
\]

%dsp-level:0
\[
\quad\neg\text{a lifetime or cross-section in external units}\ [\mathrm{BM.44}(5)],
\]

%dsp-level:0
\[
\quad\neg\text{a species identification}\ [\mathrm{Def\text{-}BM.3}(7)];
\]

\textbf{Req.} A5, A6, D0.8, D1.1, D1.4, D1.8, D2.11 (Part 1); CT0.1, CT0.2, CT0.8, CT1.29, EM.11, EM.25(7) (Part 1, Part 2); P1, P4, P5 (Parts 1–2); Def-H.4, Def-H.32, FT.108(5), FT.144(2), FT.144(4) (Part 3); Def-BM.3(7), Def-BM.5(4), Def-BM.8(8), Def-BM.15(2), Def-BM.15(5), Def-BM.17(3), Def-BM.17(4), Def-BM.18(5), Def-BM.18(6), BM.6(6), BM.23(9), BM.25(4), BM.27(6), BM.27(7), BM.29(5), BM.44(5), BM.44(6), BM.108(7), BM.109(2), BM.109(5), BM.116(1), BM.118(7), BM.120(3), BM.120(5), BM.120(7), BM.121(2), BM.121(4), BM.121(5), BM.122(6) (this part); exact rational function rows, exhaustive rational grid sign censuses, Laurent-monomial assembly [math; machine: \texttt{checks/\allowbreak{}verify\_\allowbreak{}H\_\allowbreak{}janus\_\allowbreak{}tunnel\_\allowbreak{}concrete\_\allowbreak{}spin.py} §B].

\textbf{Pf.}

%dsp-level:0
\[
(1):\ \hat J_-\ \text{on two registers with weight}\ x_-\ \text{each}\ [\mathrm{Def\text{-}BM.18}(5);\ \mathrm{BM.120}(5)];
\]

%dsp-level:0
\[
\qquad x_-^2=(4\varepsilon_q)^2=16\varepsilon_q^2\ [\text{machine}]
\]

%dsp-level:0
\[
Q_{\text{in}}=q+(-q)=0\ [\mathrm{BM.109}(1);\ \text{machine}];\qquad Q_{\text{out}}=0\ [\mathrm{D1.1}:\ \text{virtual pairs are unregistered}]
\]

%dsp-level:0
\[
(2):\ \tfrac{x_-^2}{x_+^2}=\bigl(\tfrac{4\varepsilon_q}{1-4\varepsilon_q}\bigr)^2;
\]

%dsp-level:0
\[
\qquad =1\iff4\varepsilon_q=1-4\varepsilon_q\iff\varepsilon_q=\tfrac18\ [\text{machine: the cleared linear equation}]
\]

%dsp-level:0
\[
\text{the}\ 256\text{-point grid census of the sign}\ [\text{machine}];
\]

%dsp-level:0
\[
\qquad \varepsilon_q=0\implies x_-=0\implies\text{the ratio is}\ 0;
\]

%dsp-level:0
\[
\qquad \varepsilon_q\to\tfrac14^-\implies x_+\to0\implies\text{the ratio diverges}\ [\text{machine}]
\]

%dsp-level:0
\[
(3):\ x_+x_-=4\varepsilon_q(1-4\varepsilon_q)\le\tfrac14\ [\mathrm{BM.120}(5);\ \text{machine}];\qquad x_+^2x_-^2=(x_+x_-)^2\le\tfrac1{16}\ [\text{machine}]
\]

%dsp-level:0
\[
(4):\ \mathrm{FT.144}(4)\text{'s}\ \mathcal V\text{-sector rows re-cited (frozen)};
\]

%dsp-level:0
\[
\qquad S_{\mathfrak L}
\]

%dsp-level:0
\[
= n_{\mathcal R}\ln2\ [\mathrm{CT1.29}]\ \text{and}\ n_{\mathcal R}\to0\ \text{under}\ \hat J_-\ [\mathrm{D1.2}]\implies\text{the registered content leaves while the raw entropy reaches its maximum}\ [\mathrm{CT0.2};\ \mathrm{A5}]
\]

%dsp-level:0
\[
(5):\ \mathrm{BM.25}(4),\ \mathrm{BM.27}(6)\ \text{re-cited};
\]

%dsp-level:0
\[
\qquad |w(c^{\mathrm{CP}})|^2=|w(c)|^2\ [\mathrm{EM.24}(5);\ \mathrm{BM.108}(2)];
\]

%dsp-level:0
\[
\qquad \text{the mediated composition}\ \hat J_+\otimes\hat J_+\ \text{then}\ \hat J_-\otimes\hat J_-\ \text{with weight}\ x_+^2x_-^2\ [\text{machine: the monomial}]
\]

%dsp-level:0
\[
(6):\ \mathrm{BM.118}(7)\text{'s}\ (\mathrm{R9})\ \text{row re-cited};
\]

%dsp-level:0
\[
\qquad \text{the two narrowed residues cite their own tags}\ [\mathrm{BM.6}(6);\ \mathrm{BM.122}(6)]
\]

%dsp-level:0
\[
(7):\ \mathrm{A6},\ \mathrm{BM.23}(9),\ \mathrm{BM.44}(6),\ \mathrm{FT.144}(2)\ \text{re-cited};
\]

%dsp-level:0
\[
\qquad (8):\ \text{the rescope census and the standing-rule scans}\ [\text{machine}].\ \blacksquare
\]

\textbf{Phys.} Annihilation is the same algebra read the other way, and it costs the square of the dissolution weight. The two directions have an exact ratio — four epsilon over one minus four epsilon, squared — which is one precisely at the threshold value one eighth, below one before it and above one after, so the ratio is nothing other than the detailed-balance indicator and it peaks where the round-trip weight does. At the sharp threshold the ratio is zero: nothing annihilates, which is the same statement as the no-flicker theorem seen from the other side. The round trip that joins the two directions carries the familiar parabola, and the resolvent that sums all round trips has its pole at fugacity two or beyond. What happens to the content is exact and slightly austere: a dissolved pair returns to the virtual sector at one quarter of the identity, the unique maximizer of the raw entropy, and the ledger's own entropy read goes to zero with the registered count — so nothing is carried off by radiation, because no supply vertex exists to carry it. That absence is what keeps the frozen rows honest: there is still no direct operator between the two charge sectors and they remain amplitude-equal orbits of one involution, and what the tier adds is a mediated channel through the virtual sector, composed of the two Janus moves and the CP identification. Residue (R9) closes at exactly that level — a channel with a weight, a threshold, a balance and a destination — and narrows into two named successors, the cross-section which needs a momentum-space resolvent and a gauge coupling of the charge read, and the radiation legs which need a transit supply vertex. The arrow is untouched throughout: the reversibility lives in the move set and its stationary measure, never in the tick.

\end{theorem}

\begin{definition}[{\texorpdfstring{Def-BM.19 — The spin-dependent transport defect (the multiplet-ladder read)}{Def-BM.19 - The spin-dependent transport defect (the multiplet-ladder read)}}]\label{tag:Def-BM-19}
\textbf{Def.} A read specification on $\mathrm{Def\text{-}H.37}(2)$'s holonomy defect and $\mathrm{FT.60}(5)$'s multiplet ladder; no new axiom, no new datum, no new parameter, and no identification of the two angles.

%dsp-level:0
\[
(1)\quad\textbf{the two angles, kept apart:}\qquad\begin{array}{l|l|l}\text{angle} & \text{definition} & \text{dimension}\\\hline\Theta(c)=\sum_{e\in E(c)}\widehat{ds}^{\,2}(e) & \text{the winding of a materialized closed chain}\ [\mathrm{FT.22};\ \mathrm{EM.24}(5)] & [\Theta]=\mathfrak q^{+1}\ [\mathrm{FT.55}(1)\ \mathrm{Pf}:\ \Theta\mathfrak h^2=\oint ds^2]\\\Theta_R(c) & \text{the frame-holonomy angle of}\ \mathrm{Def\text{-}H.21}\text{'s spin connection}\ [\mathrm{Def\text{-}H.21}(5),(6)] & [\Theta_R]=\mathfrak q^0\ [\mathrm{Def\text{-}H.21}(6):\ K\,dA\ \text{classical}]\end{array}
\]

%dsp-level:0
\[
k=\frac\Theta2\in\mathbb Z_{\ge0},
\]

%dsp-level:0
\[
\qquad j=\frac{|k|}2,
\]

%dsp-level:0
\[
\qquad [k]=\mathfrak q^{+1},
\]

%dsp-level:0
\[
\qquad [j]=\mathfrak q^{+1}
\]

%dsp-level:0
\[
\qquad [\mathrm{FT.22};\ \mathrm{Def\text{-}H.38};\ \mathrm{FT.55}(1)\ \mathrm{Pf};\ \text{machine: the exponent arithmetic}];
\]

%dsp-level:0
\[
(2)\quad\textbf{the no-go (re-cited, and constitutive for this definition):}
\]

%dsp-level:0
\[
\qquad \boxed{\ \text{the identification}\ \Theta_R(c)=2\pi k(c)\ \ \textbf{is ruled out}\ }
\]

%dsp-level:0
\[
\qquad [\mathrm{FT.55}(1):\ \text{the calibration no-go}]:
\]

%dsp-level:0
\[
\textbf{the ground is dimensional and exact:}
\]

%dsp-level:0
\[
\qquad [\Theta_R]=\mathfrak q^0\ \ \text{against}\ \ [2\pi k]=\mathfrak q^{+1}
\]

%dsp-level:0
\[
\qquad [\text{machine: the two exponent rows disagree}];
\]

%dsp-level:0
\[
\textbf{the corpus's own scope row:}
\]

%dsp-level:0
\[
\qquad \text{under}\ \mathrm{Def\text{-}H.37}(2)\ \text{the defect ladder reads in}\ \Theta_R\ \text{at fixed weight}\ m\ \text{with}\ \neg\text{a}\ k\text{--}\Theta_R\ \text{proportionality anywhere}
\]

%dsp-level:0
\[
\qquad [\mathrm{FT.60}\ \mathrm{Rmk}];
\]

%dsp-level:0
\[
(3)\quad\textbf{the multiplet ladder:}\qquad m\in\{k,\ k-2,\ \dots,\ -k\},\qquad\#m=k+1=2j+1\qquad[\mathrm{FT.60}(5);\ \mathrm{Def\text{-}H.38}]:
\]

%dsp-level:0
\[
\boxed{\ \Delta(m)\ :=\ \bigl\|(\mathcal U^{(k)}-1)\psi_m\bigr\|^2\ =\ \bigl|e^{im\phi}-1\bigr|^2\ =\ 2-2\cos(m\phi)\ =\ 4\sin^2\frac{m\phi}2\ }
\]

%dsp-level:0
\[
\qquad \phi:=\frac{\Theta_R}2
\]

%dsp-level:0
\[
\qquad [\mathrm{FT.60}(5);\ \mathrm{FT.54}(4);\ \mathrm{EW.33}\ \text{pattern};\ \mathrm{Def\text{-}H.21}(6)];
\]

%dsp-level:0
\[
\textbf{constitutive clause:}\ \ \text{the spin-dependence enters through the multiplet weight}\ m,
\]

%dsp-level:0
\[
\ \neg\text{through the winding}\ \Theta\ \text{and}\ \neg\text{through a single value per}\ j;
\]

%dsp-level:0
\[
(4)\quad\textbf{the exact algebraic form:}
\]

%dsp-level:0
\[
\qquad \cos(m\phi)=T_m(\cos\phi),
\]

%dsp-level:0
\[
\qquad T_{m+1}=2x\,T_m-T_{m-1},
\]

%dsp-level:0
\[
\qquad T_0=1,\ T_1=x
\]

%dsp-level:0
\[
\qquad [\text{math: Chebyshev}]:
\]

%dsp-level:0
\[
\boxed{\ \Delta(m)\ =\ 2-2\,T_{|m|}\bigl(\cos\tfrac{\Theta_R}2\bigr)\ }
\]

%dsp-level:0
\[
\textbf{the state-level specialization:}
\]

%dsp-level:0
\[
\qquad \Theta_R=\delta_v\ \ \text{at the vertex-figure locus}\implies\Delta=D_S\in\bigl\{2,\ 4,\ 2-\sqrt2,\ 1,\ 2-2\cos\tfrac\pi{10},\ 2-\sqrt3\bigr\}
\]

%dsp-level:0
\[
\qquad [\mathrm{Def\text{-}BM.14}(3);\ \mathrm{BM.82}(1);\ \mathrm{BM.114}(3)];
\]

%dsp-level:0
\[
(5)\quad\textbf{the standing zero:}
\]

%dsp-level:0
\[
\qquad k\ \text{even}\implies0\in\{k,k-2,\dots,-k\}\implies\boxed{\ \Delta(0)=2-2\cos0=0\ \ \text{exactly, at every value of}\ \Theta_R\ }
\]

%dsp-level:0
\[
[\mathrm{FT.60}(5):\ \text{every even}\ k\ \text{carries a defect-massless weight};\ \text{machine}];
\]

%dsp-level:0
\[
\qquad k\ \text{odd}\implies0\notin\{k,k-2,\dots,-k\}\implies\neg\text{a standing zero}\ [\text{machine}];
\]

%dsp-level:0
\[
(6)\quad\textbf{the mass ladder and its ceiling:}
\]

%dsp-level:0
\[
\qquad m^2_{\mathrm{pole}}(m)=\hat\mu'^2\,\Delta(m),
\]

%dsp-level:0
\[
\qquad 0\le\Delta(m)\le4\implies\boxed{\ 0\le m^2_{\mathrm{pole}}(m)\le4\hat\mu'^2\ }
\]

%dsp-level:0
\[
\qquad [\mathrm{Def\text{-}H.37}(2);\ \mathrm{FT.54}(1);\ \mathrm{FT.76}(3);\ \text{machine}];
\]

%dsp-level:0
\[
\textbf{the two}\ \mathrm{Def\text{-}H.37}\ \textbf{clauses and their matching:}
\]

%dsp-level:0
\[
\qquad (2)\ \text{implies}\ (1)\ \text{at leading order under}
\]

%dsp-level:0
\[
\qquad \boxed{\ \hat\mu'^2\,\frac{\Theta_R^2}4\ =\ \hat\mu^2\,\ell_c^2\,G(c)\ }
\]

%dsp-level:0
\[
\qquad [\mathrm{Def\text{-}H.37}:\ \text{recorded, not imposed}];
\]

%dsp-level:0
\[
(7)\quad\textbf{the residues inherited, not filled:}
\]

%dsp-level:0
\[
\qquad (\mathrm{R2})\ \ \Theta_R
\]

%dsp-level:0
\[
= \delta_v\ \text{is read at the vertex-figure locus, whose ledger realization}\ —\ \text{a clean-enclosure loop through conical vertices}\ —\ \text{is}\ \neg\text{constructed}
\]

%dsp-level:0
\[
\qquad [\mathrm{BM.67}(9),(10);\ \mathrm{BM.82}(6)(\mathrm{R2});\ \text{Census FT.63}(4)\ \mathrm{Rmk}];
\]

%dsp-level:0
\[
(\mathrm{R1})\ \ \rho_{\mathrm B}
\]

%dsp-level:0
\[
= \hat\mu'^2\mathfrak h^2/(\ln2)^2\ \text{is the single cross-currency residue, its slot declared and its value}\ \neg\text{issued}
\]

%dsp-level:0
\[
\qquad [\mathrm{A\text{-}BM.9};\ \mathrm{BM.82}(1),(6)];
\]

%dsp-level:0
\[
(8)\quad\textbf{status and scope:}
\]

%dsp-level:0
\[
\qquad \textbf{derived:}\ (1)\text{'s dimensions},\ (4)\text{'s recursion and specialization},\ (5),\ (6)\text{'s ceiling};
\]

%dsp-level:0
\[
\qquad \textbf{constitutive:}\ (3)\text{'s locus of the spin-dependence};
\]

%dsp-level:0
\[
\qquad \textbf{re-cited:}\ (2);
\]

%dsp-level:0
\[
\textbf{the scope inherited at every use:}\ \ \mathrm{Def\text{-}H.38}\text{'s nonnegative cone and its}\ k
\]

%dsp-level:0
\[
\ge 3\ \text{aspect bar}\ [\neg\text{an aspect claim is made here}];
\]

%dsp-level:0
\[
\qquad \mathrm{Def\text{-}H.37}\text{'s two P4-class residues}\ \hat\mu^2,\ \hat\mu'^2\ [\mathrm{Cal\text{-}F}];
\]

%dsp-level:0
\[
\textbf{the bar:}
\]

%dsp-level:0
\[
\qquad \neg\text{a product of the two}\ \mathrm{Def\text{-}H.37}\ \text{clauses' factors is a derived mass}\ —\ \text{the formed rows are the two clauses and their matching identity}\ (6);
\]

%dsp-level:0
\[
\qquad \neg\text{a}\ k\text{--}\Theta_R\ \text{proportionality is asserted anywhere};
\]

%dsp-level:0
\[
\neg\text{a P4 addition},
\]

%dsp-level:0
\[
\quad\neg\text{a new dimensionful scale}\ [\mathrm{P1}],
\]

%dsp-level:0
\[
\quad\neg\text{a measured value}\ [\mathrm{P5}],
\]

%dsp-level:0
\[
\quad\neg\text{a species identification}\ [\mathrm{Def\text{-}BM.3}(7)],
\]

%dsp-level:0
\[
\quad\neg\text{a currency conversion}\ [\mathrm{BM.83}(3)].
\]

\textbf{Rmk (why the defect is a ladder and not a number).} A single-valued $\Delta(j)$ would require the frame-holonomy angle to be a function of the winding, and that identification is ruled out on a dimension count: the winding carries the free scale with weight minus two while the holonomy angle is dimensionless. What the corpus derives instead is a ladder — the defect read evaluated at each weight of the spin multiplet — so the spin fixes how many values there are ($2j+1$ of them) and which one is exactly zero when $k$ is even, while $\Theta_R$ remains a separate datum whose ledger realization is residue (R2). Reading the ladder as a number is the same category error as reading a count as a strength, and the definition bars it explicitly.

\textbf{Req.} D2.9 (Parts 1–2); CT2.1 (Part 2); P1, P4, P5 (Parts 1–2); Def-H.21(5), Def-H.21(6), Def-H.37, Def-H.38, FT.22, FT.54(1), FT.54(4), FT.55(1), FT.60(5), FT.76(3), EM.24(5) (Part 3); EW.33 (Part 2); Census FT.63(4) (row-wise); Def-BM.3(7), Def-BM.14(3), BM.67(9), BM.67(10), BM.82(1), BM.82(6), BM.83(3), BM.114(3), A-BM.9 (this part); the Chebyshev recursion over $\mathbb Z[x]$, exact radical identities, $\mathfrak q$-exponent arithmetic [math; machine: \texttt{checks/\allowbreak{}verify\_\allowbreak{}H\_\allowbreak{}janus\_\allowbreak{}tunnel\_\allowbreak{}concrete\_\allowbreak{}spin.py} §C].

\textbf{Pf.}

%dsp-level:0
\[
\text{Definition block: the constitutive clause is}\ (3)\text{'s locus of the spin-dependence};
\]

%dsp-level:0
\[
\ \text{the non-derived clauses are}\ \neg\text{present}\ —\ (1)\text{'s dimensions},
\]

%dsp-level:0
\[
\ (2)\text{'s no-go},
\]

%dsp-level:0
\[
\ (4)\text{'s recursion},
\]

%dsp-level:0
\[
\ (5)\text{'s zero and}\ (6)\text{'s ceiling are all derived or re-cited};
\]

%dsp-level:0
\[
\ (7),
\]

%dsp-level:0
\[
(8)\ \text{are residue and scope records}.\ \blacksquare
\]

\textbf{Phys.} Two angles appear in this corpus and they are not the same object: the winding, which counts the line element around a closed chain and carries the free scale, and the frame-holonomy angle, which is dimensionless. Their identification is ruled out by that dimension mismatch alone, so a transport defect that depended on the winding through the holonomy formula cannot be formed. What is formed is better: a ladder. The spin fixes the multiplet, the multiplet has two-j-plus-one weights spaced by two, and the defect read at each weight is two minus two cosines of that weight times half the holonomy angle — equivalently four sine-squared of a quarter, and equivalently an integer polynomial in the cosine through the Chebyshev recursion. Two rows fall out exactly. Whenever the winding is even the ladder contains the weight zero, and the defect at weight zero is identically zero at every curvature flux, which is a standing massless state in the middle of the ladder rather than a tuning. And the whole ladder is bounded by four times the defect scale, the ceiling the pole-side read already carried. The state-level defect values are the ladder read at the vertex-figure angle, which is where the six exact algebraic numbers come from — and that locus, together with the single cross-currency ratio, stays a named residue rather than a filled one.

\end{definition}

\begin{theorem}[{\texorpdfstring{BM.124 — The spin–winding relation and the multiplet defect ladder}{BM.124 - The spin-winding relation and the multiplet defect ladder}}]\label{tag:BM-124}
\textbf{Thm.}

%dsp-level:0
\[
(1)\quad\textbf{the exact spin–winding relation:}
\]

%dsp-level:0
\[
\qquad k=\frac\Theta2\in\mathbb Z_{\ge0}\ \ \text{on shell},
\]

%dsp-level:0
\[
\qquad j=\frac{|k|}2,
\]

%dsp-level:0
\[
\qquad \Psi(c)=\mathrm{Sym}^{|k|}(\mathcal H_{\mathsf v}),
\]

%dsp-level:0
\[
\qquad \dim\Psi=|k|+1
\]

%dsp-level:0
\[
\qquad [\mathrm{FT.22};\ \mathrm{FT.25}(3);\ \mathrm{Def\text{-}H.38};\ \mathrm{EW.32}]:
\]

%dsp-level:0
\[
\boxed{\ \Theta\ =\ 2k\ =\ 4j\qquad\text{on the nonnegative cone},\qquad\Theta\in2\mathbb Z\ \ \text{for every materialized closed chain}\ }
\]

%dsp-level:0
\[
[\text{machine: the exact relation at every}\ k\in\{0,\dots,8\}];\qquad\textbf{the inverse read:}\qquad\boxed{\ j\ =\ \frac{|\Theta|}4\ }
\]

%dsp-level:0
\[
(2)\quad\textbf{the multiplet:}
\]

%dsp-level:0
\[
\qquad m\in\{k,\ k-2,\ \dots,\ -k\},
\]

%dsp-level:0
\[
\qquad \boxed{\ \#m\ =\ k+1\ =\ 2j+1\ }
\]

%dsp-level:0
\[
\qquad [\mathrm{FT.60}(5);\ \mathrm{Def\text{-}H.38};\ \text{machine: exact at every}\ k\le8];
\]

%dsp-level:0
\[
\textbf{the Casimir:}
\]

%dsp-level:0
\[
\qquad \frac k2\Bigl(\frac k2+1\Bigr)=j(j+1)
\]

%dsp-level:0
\[
\qquad [\mathrm{FT.60}(3);\ \text{machine}];
\]

%dsp-level:0
\[
\qquad \textbf{the stretched total:}
\]

%dsp-level:0
\[
\qquad m_{\max}=k=2j
\]

%dsp-level:0
\[
\qquad [\mathrm{FT.60}(6):\ \text{maximal-weight CG on the cone}];
\]

%dsp-level:0
\[
(3)\quad\textbf{the spin-dependent transport defect is the ladder, and the ladder is exact:}
\]

%dsp-level:0
\[
\qquad \boxed{\ \Delta(m)\ =\ 4\sin^2\frac{m\Theta_R}4\ =\ 2-2\cos\frac{m\Theta_R}2\ =\ 2-2\,T_{|m|}\Bigl(\cos\frac{\Theta_R}2\Bigr)\ }
\]

%dsp-level:0
\[
\qquad [\mathrm{Def\text{-}BM.19}(3),(4);\ \mathrm{FT.60}(5);\ \mathrm{FT.54}(4)]:
\]

%dsp-level:0
\[
\textbf{the Chebyshev recursion verified exactly:}
\]

%dsp-level:0
\[
\qquad T_0=1,\quad T_1=x,\quad T_{m+1}=2x\,T_m-T_{m-1}
\]

%dsp-level:0
\[
\qquad [\text{machine: the integer-polynomial recursion to}\ m=8];
\]

%dsp-level:0
\[
\textbf{the exact witness:}
\]

%dsp-level:0
\[
\qquad T_m\bigl(\tfrac12\bigr)=\cos\frac{m\pi}3\ \ \text{on the six-cycle}\ \bigl(1,\ \tfrac12,\ -\tfrac12,\ -1,\ -\tfrac12,\ \tfrac12\bigr)
\]

%dsp-level:0
\[
\qquad [\text{machine: all}\ m\le8];
\]

%dsp-level:0
\[
(4)\quad\textbf{the refutation of a single-valued}\ \Delta(j):\qquad\#\{m\}=2j+1\ge2\ \ \text{for every}\ j\ge\tfrac12\qquad[(2);\ \text{machine}]:
\]

%dsp-level:0
\[
\boxed{\ \textbf{the spin determines the ladder's length,}\ \neg\textbf{a single defect value}\ —\ \Delta(j)\ \textbf{as one number is}\ \neg\textbf{formed}\ }
\]

%dsp-level:0
\[
\textbf{the two anchors:}\qquad\begin{array}{c|c|c|c}k & j & m\ \text{values} & \Delta(m)\\\hline1 & \tfrac12 & \{+1,\ -1\} & 2-2\cos\tfrac{\Theta_R}2\ \ \text{(both weights agree, by evenness)}\\2 & 1 & \{+2,\ 0,\ -2\} & 2-2\cos\Theta_R,\quad\mathbf0,\quad2-2\cos\Theta_R\end{array}\qquad[\mathrm{FT.60}(5);\ \mathrm{BM.106}(1)\text{'s two floor rows};\ \text{machine}];
\]

%dsp-level:0
\[
\boxed{\ \textbf{at}\ j=\tfrac12\ \textbf{the ladder is state-independent and}\ \neg\textbf{split; at}\ j=1\ \textbf{it splits into three with an exact zero in the middle}\ }
\]

%dsp-level:0
\[
\qquad [\mathrm{FT.60}(5):\ k{=}1\ \text{reproduces}\ \mathrm{FT.54}(4)\ \text{and its state-independence is special to the seed}];
\]

%dsp-level:0
\[
(5)\quad\textbf{the two angles are}\ \neg\textbf{proportional, on a dimension count:}
\]

%dsp-level:0
\[
\qquad [k]=\mathfrak q^{+1},
\]

%dsp-level:0
\[
\qquad [\Theta_R]=\mathfrak q^0,
\]

%dsp-level:0
\[
\qquad [2\pi]=\mathfrak q^0
\]

%dsp-level:0
\[
\qquad [\mathrm{FT.55}(1)\ \mathrm{Pf};\ \mathrm{Def\text{-}H.21}(6);\ \text{machine: the exponent rows}]:
\]

%dsp-level:0
\[
\boxed{\ \Theta_R=2\pi k\ \ \textbf{would require}\ \ \mathfrak q^0=\mathfrak q^{+1}:\ \ \textbf{dimensionally impossible, with no measurement used}\ }
\]

%dsp-level:0
\[
\qquad [\mathrm{FT.55}(1):\ \text{the calibration no-go re-cited}];
\]

%dsp-level:0
\[
\textbf{the consequence:}
\]

%dsp-level:0
\[
\qquad \Delta\ \text{is}\ \neg\text{a function of}\ \Theta,\ \text{hence}\ \neg\text{a function of}\ j\ \text{alone};
\]

%dsp-level:0
\[
\qquad \Delta\ \text{is a function of the pair}\ (m,\ \Theta_R)\ \text{with}\ m\ \text{spin-determined and}\ \Theta_R\ \text{independent}\ [\mathrm{Def\text{-}BM.19}(3),(6)];
\]

%dsp-level:0
\[
\textbf{the corpus's own scope row:}
\]

%dsp-level:0
\[
\qquad \neg\text{a}\ k\text{--}\Theta_R\ \text{proportionality anywhere}\ [\mathrm{FT.60}\ \mathrm{Rmk};\ \mathrm{FT.55}(4):\ \text{the winding enters as a copy count,}\ \neg\text{as an angle}];
\]

%dsp-level:0
\[
(6)\quad\textbf{the standing zero:}\qquad k\ \text{even}\iff0\in\{k,k-2,\dots,-k\}\qquad[\text{machine: exhaustive to}\ k=8]:
\]

%dsp-level:0
\[
\boxed{\ \Delta(0)=2-2T_0(\cos\tfrac{\Theta_R}2)=2-2=0\ \ \textbf{exactly, at every}\ \Theta_R,\ \textbf{for every even}\ k\ }\qquad[\mathrm{FT.60}(5);\ \text{machine}];
\]

%dsp-level:0
\[
k\ \text{odd}\implies m\in\{\pm1,\pm3,\dots,\pm k\}\implies\neg\exists\,m=0\implies\neg\text{a standing zero}\qquad[\text{machine}];
\]

%dsp-level:0
\[
\textbf{the parity row that follows:}
\]

%dsp-level:0
\[
\qquad \boxed{\ \textbf{the defect ladder carries an exact massless weight precisely when the exchange sign is}\ +1\ }
\]

%dsp-level:0
\[
\qquad [k\ \text{even}\iff\varepsilon=(-1)^k=+1\ [\mathrm{BM.113}(4)];\ \text{machine: the two parities coincide}];
\]

%dsp-level:0
\[
(7)\quad\textbf{the state-level specialization:}
\]

%dsp-level:0
\[
\qquad \Theta_R=\delta_v=2\pi-d\,\frac{(L-2)\pi}L,
\]

%dsp-level:0
\[
\qquad D_S=4\sin^2\frac{\delta_v}4
\]

%dsp-level:0
\[
\qquad [\mathrm{Def\text{-}BM.14}(3);\ \text{Census FT.63}(4),(6);\ \mathrm{BM.82}(1)]:
\]

%dsp-level:0
\[
\textbf{the six values are the ladder read at six vertex-figure angles,}\ \neg\textbf{six spin values:}
\]

%dsp-level:0
\[
\qquad D_S\in\bigl\{2,\ 4,\ 2-\sqrt2,\ 1,\ 2-2\cos\tfrac\pi{10},\ 2-\sqrt3\bigr\},
\]

%dsp-level:0
\[
\qquad \text{four fields, six minimal polynomials}
\]

%dsp-level:0
\[
[\mathrm{Def\text{-}BM.14}(3);\ \mathrm{BM.114}(3);\ \text{machine re-verified}];
\]

%dsp-level:0
\[
\qquad \textbf{the index is}\ (\mathcal F_g,\beta),\ \neg\ j
\]

%dsp-level:0
\[
\qquad [\mathrm{BM.90}(5):\ \text{the branch-blindness theorem};\ \text{machine: the symbol census carries}\ \neg j];
\]

%dsp-level:0
\[
(8)\quad\textbf{status:}
\]

%dsp-level:0
\[
\qquad \textbf{derived:}\ (1),\ (2),\ (3),\ (6),\ (7);
\]

%dsp-level:0
\[
\qquad \textbf{refuted as formed:}\ (4)\text{'s single-valued}\ \Delta(j),\ (5)\text{'s}\ \Theta_R=2\pi k;
\]

%dsp-level:0
\[
\neg\text{a declaration},
\]

%dsp-level:0
\[
\quad\neg\text{a P4 addition},
\]

%dsp-level:0
\[
\quad\neg\text{a new dimensionful scale}\ [\mathrm{P1}],
\]

%dsp-level:0
\[
\quad\neg\text{a measured value}\ [\mathrm{P5}],
\]

%dsp-level:0
\[
\quad\neg\text{a numerical coupling}\ [\mathrm{EM.25}(7)],
\]

%dsp-level:0
\[
\quad\neg\text{a species identification}\ [\mathrm{Def\text{-}BM.3}(7)];
\]

\textbf{Req.} D2.9 (Parts 1–2); CT2.1, EW.32 (Part 2); P1, P4, P5 (Parts 1–2); Def-H.21(6), Def-H.37(2), Def-H.38, FT.22, FT.25(3), FT.54(4), FT.55(1), FT.55(4), FT.60(3), FT.60(5), FT.60(6) (Part 3); Census FT.63(4), FT.63(6) (row-wise); Def-BM.3(7), Def-BM.14(3), Def-BM.19(3), Def-BM.19(4), Def-BM.19(6), BM.82(1), BM.90(5), BM.106(1), BM.113(4), BM.114(3) (this part); the Chebyshev recursion over $\mathbb Z[x]$ with an exact rational witness, $\mathfrak q$-exponent arithmetic, exhaustive parity censuses [math; machine: \texttt{checks/\allowbreak{}verify\_\allowbreak{}H\_\allowbreak{}janus\_\allowbreak{}tunnel\_\allowbreak{}concrete\_\allowbreak{}spin.py} §C].

\textbf{Pf.}

%dsp-level:0
\[
(1):\ k=\Theta/2\ [\mathrm{FT.22};\ \mathrm{FT.25}(3)];
\]

%dsp-level:0
\[
\qquad j=|k|/2\ [\mathrm{Def\text{-}H.38}]\implies\Theta=2k=2\cdot2j=4j\ [\text{machine: exact at every}\ k\le8];
\]

%dsp-level:0
\[
\qquad \Theta\in2\mathbb Z\ [\mathrm{FT.22}]
\]

%dsp-level:0
\[
(2):\ \#\{k,k-2,\dots,-k\}=k+1\ [\text{machine: the arithmetic-progression count}];
\]

%dsp-level:0
\[
\qquad k+1=2\cdot\tfrac k2+1=2j+1\ [\text{machine}];
\]

%dsp-level:0
\[
\qquad \tfrac k2(\tfrac k2+1)=j(j+1)\ \text{at}\ j=k/2\ [\text{machine}]
\]

%dsp-level:0
\[
(3):\ \|(U^{(k)}-1)\psi_m\|^2
\]

%dsp-level:0
\[
= |e^{im\phi}-1|^2
\]

%dsp-level:0
\[
= (e^{im\phi}-1)(e^{-im\phi}-1)
\]

%dsp-level:0
\[
= 2-e^{im\phi}-e^{-im\phi}
\]

%dsp-level:0
\[
= 2-2\cos(m\phi)\ [\mathrm{FT.60}(5);\ \text{machine over}\ \mathbb Q(i)]
\]

%dsp-level:0
\[
2-2\cos(m\phi)=4\sin^2\tfrac{m\phi}2\ \text{at}\ \phi=\Theta_R/2\ [\text{machine: the half-angle identity}];
\]

%dsp-level:0
\[
\qquad \cos(m\phi)=T_m(\cos\phi)\ \text{with}\ T_{m+1}=2xT_m-T_{m-1}\ [\text{machine: the recursion to}\ m=8]
\]

%dsp-level:0
\[
T_m(\tfrac12)=\cos\tfrac{m\pi}3:
\]

%dsp-level:0
\[
\quad m
\]

%dsp-level:0
\[
= 0..8\mapsto\bigl(1,\tfrac12,-\tfrac12,-1,-\tfrac12,\tfrac12,1,\tfrac12,-\tfrac12\bigr)\ [\text{machine: exact rationals on the six-cycle}]
\]

%dsp-level:0
\[
(4):\ \#\{m\}=k+1\ge2\ \text{for}\ k\ge1\ [\text{machine}];
\]

%dsp-level:0
\[
\qquad k=1:\ m\in\{+1,-1\},\ \Delta(\pm1)=2-2\cos(\pm\tfrac{\Theta_R}2)=2-2\cos\tfrac{\Theta_R}2\ \text{by evenness}\ [\text{machine}]
\]

%dsp-level:0
\[
k=2:\ m\in\{+2,0,-2\},\ \Delta(\pm2)=2-2\cos\Theta_R,\ \Delta(0)=0\ [\text{machine}]
\]

%dsp-level:0
\[
(5):\ \Theta\mathfrak h^2=\oint ds^2\implies[\Theta]=[\mathfrak h^{-2}]=\mathfrak q^{+1}\implies[k]=\mathfrak q^{+1}\ [\mathrm{FT.55}(1)\ \mathrm{Pf};\ \text{machine}];
\]

%dsp-level:0
\[
\qquad [\Theta_R]=\mathfrak h^0=\mathfrak q^0\ [\mathrm{Def\text{-}H.21}(6)];
\]

%dsp-level:0
\[
\qquad \mathfrak q^0\ne\mathfrak q^{+1}\ [\text{machine}]
\]

%dsp-level:0
\[
(6):\ k\ \text{even}\iff\exists r:\ k-2r=0\iff r=k/2\in\mathbb Z\ [\text{machine: exhaustive to}\ k=8];
\]

%dsp-level:0
\[
\qquad \Delta(0)=2-2T_0=2-2\cdot1=0\ [\text{machine}]
\]

%dsp-level:0
\[
\varepsilon=(-1)^k\ [\mathrm{BM.113}(4)];
\]

%dsp-level:0
\[
\qquad k\ \text{even}\iff\varepsilon=+1\iff\Delta(0)=0\ \text{is in the ladder}\ [\text{machine: the two parity censuses coincide}]
\]

%dsp-level:0
\[
(7):\ \mathrm{Def\text{-}BM.14}(3)\text{'s table re-cited};
\]

%dsp-level:0
\[
\qquad \text{the symbol census of}\ D_S\ \text{carries}\ \{\mathcal F_g,\beta,\delta_v,L,d\}\ \text{and}\ \neg\{j,k\}\ [\mathrm{BM.90}(5);\ \text{machine}];
\]

%dsp-level:0
\[
\qquad (8):\ \text{the standing-rule scans}\ [\text{machine}].\ \blacksquare
\]

\textbf{Phys.} The spin and the winding are the same integer read twice, and the relation is exact rather than fitted: the winding is twice the copy number and the spin is half of it, so the winding is four times the spin on the cone, and it is always an even integer for a materialized closed chain. What the spin then controls is not one mass but a ladder — two-j-plus-one weights spaced by two — and the defect at each weight is two minus two cosines of the weight times half the holonomy angle, which is also four sine-squared of a quarter and also an integer polynomial in the cosine by the Chebyshev recursion. At the two anchor windings the difference is visible: at spin one half the ladder has two weights that agree by evenness and nothing splits, while at spin one it splits into three with an exact zero in the middle. That zero is not a tuning but a parity theorem — the ladder contains the weight zero precisely when the winding is even, which is precisely when the exchange sign is plus one, so the defect ladder carries a massless weight exactly on the bosonic side. What the tier refuses is the fusion of the two angles: the winding carries the free scale and the holonomy angle does not, so identifying them would equate two different dimensions, and that is the whole content of the calibration no-go. The six exact defect numbers are therefore the ladder read at six vertex-figure angles and indexed by floor and branch, not by spin, which is why a spin-dependent retention probability was refuted earlier and stays refuted.

\end{theorem}

\begin{theorem}[{\texorpdfstring{BM.125 — The spin-dependent mass: the exact ladder, the two clauses and their matching, and the ceiling}{BM.125 - The spin-dependent mass: the exact ladder, the two clauses and their matching, and the ceiling}}]\label{tag:BM-125}
\textbf{Thm.}

%dsp-level:0
\[
(1)\quad\textbf{the mass ladder:}
\]

%dsp-level:0
\[
\qquad \boxed{\ m^2_{\mathrm{pole}}(m)\ =\ \hat\mu'^2\,\Delta(m)\ =\ 4\hat\mu'^2\sin^2\frac{m\Theta_R}4,\qquad m\in\{k,k-2,\dots,-k\},\qquad k=2j\ }
\]

%dsp-level:0
\[
\qquad [\mathrm{Def\text{-}H.37}(2);\ \mathrm{Def\text{-}BM.19}(3),(6);\ \mathrm{FT.54}(4)]:
\]

%dsp-level:0
\[
\textbf{the ladder's length is the spin's:}
\]

%dsp-level:0
\[
\qquad \#\{m^2_{\mathrm{pole}}\ \text{values}\}\ \le\ 2j+1
\]

%dsp-level:0
\[
\qquad [\text{equality when the}\ \Delta(m)\ \text{are distinct};\ \text{machine}];
\]

%dsp-level:0
\[
(2)\quad\textbf{the exact range and the ceiling:}\qquad0\le4\sin^2\frac{m\Theta_R}4\le4\qquad[\text{machine: the sine-square bound}]:
\]

%dsp-level:0
\[
\boxed{\ 0\ \le\ m^2_{\mathrm{pole}}(m)\ \le\ 4\hat\mu'^2\ }\qquad[\mathrm{FT.76}(3):\ \text{the pole-side ceiling re-cited};\ \text{machine}];
\]

%dsp-level:0
\[
\textbf{the ceiling is attained:}
\]

%dsp-level:0
\[
\qquad \Delta=4\iff\sin^2\frac{m\Theta_R}4=1\iff m\Theta_R\equiv2\pi\ (\mathrm{mod}\ 4\pi)
\]

%dsp-level:0
\[
\qquad [\text{machine: the exact locus}];
\]

%dsp-level:0
\[
\textbf{the ladder's own ceiling is the composite's:}
\]

%dsp-level:0
\[
\qquad \lambda^{(n)}_{\max}=4n\ \ \text{is the sum of}\ n\ \text{attained ceilings}
\]

%dsp-level:0
\[
\qquad [\mathrm{BM.117}(4);\ \text{machine}];
\]

%dsp-level:0
\[
(3)\quad\textbf{the two}\ \mathrm{Def\text{-}H.37}\ \textbf{clauses and their matching identity:}
\]

%dsp-level:0
\[
\qquad (1)\ \ m^2_{\mathrm{pole}}=\hat\mu^2\ell_c^2G(c),
\]

%dsp-level:0
\[
\qquad G(c)=\sum_{A,B}c_Ac_B\operatorname{Re}\langle\mathcal G_A\mathsf v,\mathcal G_B\mathsf v\rangle
\]

%dsp-level:0
\[
\qquad [\mathrm{EW.33}\ \text{Gram data}];
\]

%dsp-level:0
\[
(2)\ \ m^2_{\mathrm{pole}}=\hat\mu'^2\cdot4\sin^2\frac{\Theta_R}4,
\]

%dsp-level:0
\[
\qquad [\hat\mu'^2]=\mathfrak q,
\]

%dsp-level:0
\[
\qquad [\hat\mu^2]=\mathfrak q^2
\]

%dsp-level:0
\[
\qquad [\mathrm{Def\text{-}H.37}(1),(2);\ \mathrm{EW.33}\ \text{packaging}];
\]

%dsp-level:0
\[
\boxed{\ (2)\ \text{implies}\ (1)\ \text{at leading order under the matching}\ \ \hat\mu'^2\,\frac{\Theta_R^2}4\ =\ \hat\mu^2\,\ell_c^2\,G(c)\ }
\]

%dsp-level:0
\[
\qquad [\mathrm{Def\text{-}H.37}:\ \text{recorded, not imposed};\ \mathrm{FT.54}(2);\ \mathrm{FT.33}];
\]

%dsp-level:0
\[
\textbf{the dimension check on the matching:}
\]

%dsp-level:0
\[
\qquad [\hat\mu'^2\Theta_R^2]=\mathfrak q\cdot\mathfrak q^0=\mathfrak q,
\]

%dsp-level:0
\[
\qquad [\hat\mu^2\ell_c^2G(c)]=\mathfrak q^2\cdot[\ell_c^2G(c)]=\mathfrak q\implies[\ell_c^2G(c)]=\mathfrak q^{-1}
\]

%dsp-level:0
\[
\qquad [\text{machine: the exponent arithmetic}];
\]

%dsp-level:0
\[
(4)\quad\textbf{the refutation of the directive's product form:}
\]

%dsp-level:0
\[
\qquad m(j)=\hat\mu'\,\ell_c\sqrt{G(c)\,\Delta(j)}\ \ \text{would require:}
\]

%dsp-level:0
\[
(\mathrm i)\quad\text{a single}\ \Delta(j)\ —\ \text{refuted: the defect is a ladder of}\ 2j+1\ \text{values}\ [\mathrm{BM.124}(4)];
\]

%dsp-level:0
\[
(\mathrm{ii})\quad\text{a product of clause}\ (1)\text{'s}\ \ell_c^2G(c)\ \text{with clause}\ (2)\text{'s}\ \Delta\ —\ \text{refuted: the two clauses are two reads of one mass, related by}\ (3)\text{'s matching,}\ \neg\text{two multiplicative factors};
\]

%dsp-level:0
\[
(\mathrm{iii})\quad\text{a dimension match}\ —\ \text{refuted:}\ [\hat\mu'\ell_c\sqrt{G\Delta}]
\]

%dsp-level:0
\[
= \mathfrak q^{1/2}\cdot[\ell_c]\cdot\mathfrak q^{-1/2}\ \text{against}\ [m]
\]

%dsp-level:0
\[
= \mathfrak q^{1/2},
\]

%dsp-level:0
\[
\ \text{so}\ [\ell_c]\ \text{would have to be}\ \mathfrak q^0\ \text{while}\ (3)\text{'s matching gives}\ [\ell_c^2G]\qquad = \mathfrak q^{-1}
\]

%dsp-level:0
\[
\qquad [\text{machine: the two exponent rows disagree}];
\]

%dsp-level:0
\[
\boxed{\ \textbf{the formed rows are the two clauses separately and their matching identity}\ (3)\ \textbf{— the product form is}\ \neg\textbf{formed on three independent grounds}\ }
\]

%dsp-level:0
\[
(5)\quad\textbf{the standing zero is a massless weight, not a massless species:}
\]

%dsp-level:0
\[
\qquad k\ \text{even}\implies m=0\in\text{the ladder}\implies m^2_{\mathrm{pole}}(0)=0\ \ \text{exactly at every}\ \Theta_R
\]

%dsp-level:0
\[
\qquad [\mathrm{BM.124}(6);\ \mathrm{FT.60}(5);\ \text{machine}]:
\]

%dsp-level:0
\[
\textbf{the scope bars:}
\]

%dsp-level:0
\[
\qquad \neg\text{a species identification attaches to the zero weight}\ [\mathrm{Def\text{-}BM.3}(7);\ \mathrm{P5}];
\]

%dsp-level:0
\[
\qquad \neg\text{a massless}\ \nu\ \text{line is formed}\ —\ \text{the rung-one}\ \nu\ \text{row has}\ m_\nu^{\mathrm{rest}}>0\ \text{strictly on two SOS sandwiches}\ [\mathrm{BM.96}(3),(4)\ \text{re-cited},\ \neg\text{extended}];
\]

%dsp-level:0
\[
\neg\text{a right-ray matter line,}\ \neg\text{a seesaw-type mechanism and}\ \neg\text{a Yukawa-type object is added}\ [\mathrm{BM.96}(4),(6);\ \mathrm{EW.15};\ \mathrm{FT.145}(3)];
\]

%dsp-level:0
\[
(6)\quad\textbf{the residue rows:}
\]

%dsp-level:0
\[
\qquad (\mathrm{R2})\ \ \Theta_R
\]

%dsp-level:0
\[
= \delta_v\ \text{at the vertex-figure locus, whose ledger realization is}\ \neg\text{constructed}\ [\mathrm{BM.67}(9),(10);\ \mathrm{BM.82}(6)(\mathrm{R2})]\ \text{— inherited at every use of}\ (1);
\]

%dsp-level:0
\[
(\mathrm{R1})\ \ \rho_{\mathrm B}=\hat\mu'^2\mathfrak h^2/(\ln2)^2\ \text{is the single cross-currency residue:}
\]

%dsp-level:0
\[
\qquad \frac{m^2_{\mathrm{pole}}}{m^2_{\mathrm{hid}}}=\rho_{\mathrm B}\mathcal R_s,
\]

%dsp-level:0
\[
\qquad \mathcal R_s=\frac{D_S}{4E^2}\in\overline{\mathbb Q}
\]

%dsp-level:0
\[
\qquad [\mathrm{BM.82}(1),(2);\ \mathrm{A\text{-}BM.9}:\ \text{the slot declared, the value}\ \neg\text{issued}];
\]

%dsp-level:0
\[
\neg\text{a currency is converted:}\ [\Delta]
\]

%dsp-level:0
\[
= \mathfrak q^0\ \text{and}\ [\hat\mu'^2]
\]

%dsp-level:0
\[
= \mathfrak q\ \text{against}\ [m^2_{\mathrm{hid}}]
\]

%dsp-level:0
\[
= \mathfrak q^{+1}\ \text{and}\ [\ln(4/3)]\qquad = \mathfrak q^0
\]

%dsp-level:0
\[
\qquad [\mathrm{BM.102}(2);\ \mathrm{BM.83}(3);\ \text{machine: the four exponent rows}];
\]

%dsp-level:0
\[
(7)\quad\textbf{status:}\qquad\textbf{derived:}\ (1),\ (2),\ (3),\ (5)\text{'s zero row};\qquad\textbf{refuted as formed, on three grounds:}\ (4);
\]

%dsp-level:0
\[
\textbf{inherited,}\ \neg\textbf{filled:}\ (6)\text{'s}\ (\mathrm{R1}),\ (\mathrm{R2});
\]

%dsp-level:0
\[
\neg\text{a declaration},
\]

%dsp-level:0
\[
\quad\neg\text{a P4 addition},
\]

%dsp-level:0
\[
\quad\neg\text{a new dimensionful scale}\ [\mathrm{P1}],
\]

%dsp-level:0
\[
\quad\neg\text{a measured value}\ [\mathrm{P5}],
\]

%dsp-level:0
\[
\quad\neg\text{a numerical coupling}\ [\mathrm{EM.25}(7)],
\]

%dsp-level:0
\[
\quad\neg\text{a species identification}\ [\mathrm{Def\text{-}BM.3}(7)];
\]

\textbf{Req.} D2.9 (Parts 1–2); CT2.1, EW.15, EW.33 (Part 2); P1, P4, P5 (Parts 1–2); Def-H.37(1), Def-H.37(2), FT.33, FT.54(2), FT.54(4), FT.60(5), FT.76(3), FT.145(3) (Part 3); Def-BM.3(7), Def-BM.14(3), Def-BM.19(3), Def-BM.19(6), BM.67(9), BM.67(10), BM.82(1), BM.82(2), BM.82(6), BM.83(3), BM.96(3), BM.96(4), BM.96(6), BM.102(2), BM.117(4), BM.124(4), BM.124(6), A-BM.9 (this part); exact radical and trigonometric identities, $\mathfrak q$-exponent arithmetic, the sine-square bound on exact rationals [math; machine: \texttt{checks/\allowbreak{}verify\_\allowbreak{}H\_\allowbreak{}janus\_\allowbreak{}tunnel\_\allowbreak{}concrete\_\allowbreak{}spin.py} §C].

\textbf{Pf.}

%dsp-level:0
\[
(1):\ \mathrm{Def\text{-}H.37}(2)\text{'s}\ m^2_{\mathrm{pole}}
\]

%dsp-level:0
\[
= \hat\mu'^2\cdot4\sin^2(\Theta_R/4)\ \text{read at the multiplet weight}\ m\Theta_R
\]

%dsp-level:0
\[
\mapsto m\Theta_R\ [\mathrm{FT.60}(5);\ \mathrm{FT.54}(4)\text{'s exact SU(2) defect}];
\]

%dsp-level:0
\[
\qquad \#\{m\}=k+1=2j+1\ [\mathrm{BM.124}(2)]
\]

%dsp-level:0
\[
(2):\ 0\le\sin^2\le1\implies0\le4\sin^2\le4\ [\text{machine}];
\]

%dsp-level:0
\[
\qquad \mathrm{FT.76}(3)\text{'s ceiling re-cited};
\]

%dsp-level:0
\[
\qquad \sin^2\tfrac{m\Theta_R}4=1\iff\tfrac{m\Theta_R}4\equiv\tfrac\pi2\ (\mathrm{mod}\ \pi)\iff m\Theta_R\equiv2\pi\ (\mathrm{mod}\ 4\pi)\ [\text{machine}]
\]

%dsp-level:0
\[
\lambda^{(n)}=\sum_aD_S(s_a)\le\sum_a4=4n\ [\mathrm{BM.114}(3);\ \mathrm{BM.117}(4);\ \text{machine}]
\]

%dsp-level:0
\[
(3):\ \mathrm{Def\text{-}H.37}(1),(2)\ \text{re-cited};
\]

%dsp-level:0
\[
\qquad 4\sin^2\tfrac{\Theta_R}4
\]

%dsp-level:0
\[
= \Theta_R^2+O(\Theta_R^4)\implies\hat\mu'^2\Theta_R^2/4
\]

%dsp-level:0
\[
= \hat\mu^2\ell_c^2G(c)\ \text{at leading order}\ [\mathrm{FT.54}(2);\ \mathrm{FT.33}]
\]

%dsp-level:0
\[
[\hat\mu'^2\Theta_R^2]=\mathfrak q\cdot\mathfrak q^0=\mathfrak q;
\]

%dsp-level:0
\[
\qquad [\hat\mu^2]=\mathfrak q^2\implies[\ell_c^2G]=\mathfrak q^{-1}\ [\text{machine: the exponent arithmetic}]
\]

%dsp-level:0
\[
(4):\ (\mathrm i)\ \mathrm{BM.124}(4);
\]

%dsp-level:0
\[
\qquad (\mathrm{ii})\ (3)\text{'s matching is an equality of two reads of one mass, so their product has}\ [m^4]\ \text{not}\ [m^2]\ [\text{machine: the exponent row}];
\]

%dsp-level:0
\[
\qquad (\mathrm{iii})\ \text{the two}\ [\ell_c]\ \text{rows disagree}\ [\text{machine}]
\]

%dsp-level:0
\[
(5):\ \Delta(0)=0\implies m^2_{\mathrm{pole}}(0)=\hat\mu'^2\cdot0=0\ [\text{machine}];
\]

%dsp-level:0
\[
\qquad \mathrm{BM.96}(3),(4)\text{'s strict-positivity rows re-cited at their own locus}\ [\neg\text{extended}]
\]

%dsp-level:0
\[
(6):\ \mathrm{BM.82}(1),(2)\ \text{re-cited};
\]

%dsp-level:0
\[
\qquad [\Delta]=\mathfrak q^0,\ [\hat\mu'^2]=\mathfrak q,\ [m^2_{\mathrm{hid}}]=\mathfrak q^{+1}\ [\mathrm{BM.102}(2);\ \text{machine}];
\]

%dsp-level:0
\[
\qquad (7):\ \text{the standing-rule scans}\ [\text{machine}].\ \blacksquare
\]

\textbf{Phys.} The mass ladder is the defect ladder times one scale, so it inherits everything the ladder has: two-j-plus-one values, an exact zero at the middle weight whenever the winding is even, and a hard ceiling at four times the defect scale that the pole-side read already carried. The composite ceiling is just the sum of attained constituent ceilings, which is why the resonance tier's four-n bound is not a new ceiling but the same one counted. The two clauses of the defect-transfer interface are then kept apart rather than multiplied: one reads the mass through the Gram data of the weak transport with a length factor, the other through the holonomy angle with a defect scale, and they meet in a single matching identity whose dimension check forces the length-times-Gram factor to carry weight minus one. The product form the question proposes fails on three separate counts — there is no single defect value to put in it, multiplying two reads of one mass gives the wrong power, and the dimensions do not close — so the tier supplies the matching identity instead. The massless middle weight is a property of a ladder and not a species claim: the rung-one neutral line keeps its strictly positive rest mass on two sum-of-squares sandwiches, and no right-ray line, no seesaw-type mechanism and no Yukawa-type object is added anywhere. Two residues are inherited and left unfilled, the vertex-figure locus and the single cross-currency ratio, and no currency is converted between them.

\end{theorem}

\medskip\noindent\textit{\texorpdfstring{BM.126 — The spin-dependence of the Janus crossing: refuted on the four standing grounds, with the derived domain row}{BM.126 - The spin-dependence of the Janus crossing: refuted on the four standing grounds, with the derived domain row}.}\ Structural-constraints record: see Appendix B, Section B.2.7 of the Addendum.\par\medskip

\begin{theorem}[{\texorpdfstring{BM.127 — The winding-parity identification: two reads separated}{BM.127 - The winding-parity identification: two reads separated}}]\label{tag:BM-127}
\textbf{Thm.}

%dsp-level:0
\[
(1)\quad\textbf{the statistics read:}\qquad\boxed{\ \varepsilon(c)\ =\ (-1)^{|k(c)|}\ =\ (-1)^{2j(c)}\ =\ \begin{cases}+1 & j\in\mathbb Z\\ -1 & j\in\mathbb Z+\tfrac12\end{cases}\ }\qquad[\mathrm{BM.113}(2),(4);\ \text{machine: the census to}\ k=8];
\]

%dsp-level:0
\[
(2)\quad\textbf{the winding is a ledger datum:}
\]

%dsp-level:0
\[
\qquad \Theta(c)=\sum_{e\in E(c)}\widehat{ds}^{\,2}(e),
\]

%dsp-level:0
\[
\qquad \widehat{ds}^{\,2}(e)=\kappa_e\ \ \text{is registration data determined at crystallization}
\]

%dsp-level:0
\[
\qquad [\mathrm{FT.22};\ \mathrm{FT.24};\ \text{Census FT.28}(2);\ \text{A §0.2a}]:
\]

%dsp-level:0
\[
\Theta_{\mathrm{int}}=2\kappa\in2\mathbb Z,\qquad k=\frac\Theta2=\kappa\in\mathbb Z\setminus\{0\}\qquad[\mathrm{BM.106}(1);\ \text{machine}];
\]

%dsp-level:0
\[
\boxed{\ \textbf{the parity of}\ k\ \textbf{is a read of registration data}\ —\ \neg\text{an external postulate,}\ \neg\text{a declared assignment and}\ \neg\text{a measurement}\ }
\]

%dsp-level:0
\[
\qquad [\mathrm{P5};\ \text{machine: the symbol census of}\ (-1)^{|k|}\ \text{carries}\ \{k\}\ \text{alone}];
\]

%dsp-level:0
\[
(3)\quad\textbf{the kind read is a different object:}
\]

%dsp-level:0
\[
\qquad \text{bound}\ :=\ \text{closed self-contained chain, content fermionic};
\]

%dsp-level:0
\[
\qquad \text{transit}\ :=\ \text{propagating open chain, content bosonic}
\]

%dsp-level:0
\[
\qquad [\mathrm{D1.5}:\ \textbf{a definitional label}]:
\]

%dsp-level:0
\[
\textbf{the dichotomy is the closed/open read}\qquad[\mathrm{FT.80}(5);\ \mathrm{D1.4}];
\]

%dsp-level:0
\[
(4)\quad\textbf{the two reads are}\ \neg\textbf{equated, and the floor census is the proof:}
\]

%dsp-level:0
\[
\qquad \kappa=1\implies k=1\implies j=\tfrac12\implies\varepsilon=-1;
\]

%dsp-level:0
\[
\qquad \kappa=2\implies k=2\implies j=1\implies\varepsilon=+1
\]

%dsp-level:0
\[
\qquad [\mathrm{BM.106}(1);\ \text{machine}]:
\]

%dsp-level:0
\[
\boxed{\ \textbf{both exchange signs occur on materialized closed bound chains of one}\ \mathrm{D1.5}\ \textbf{kind}\ —\ \text{so the kind label is}\ \neg\textbf{a statistics and the statistics is}\ \neg\textbf{a kind}\ }
\]

%dsp-level:0
\[
\textbf{the exact separation:}\qquad\begin{array}{l|l|l}\text{read} & \text{its argument} & \text{its values}\\\hline\text{the statistics}\ \varepsilon & k\ \text{(an integer registration datum)} & \{+1,\ -1\},\ \text{sharp per chain}\\\text{the}\ \mathrm{D1.5}\ \text{kind} & \text{closed against open}\ [\mathrm{D1.4}] & \{\text{bound},\ \text{transit}\},\ \text{both signs on one kind}\\\text{the registration}\ \kappa & \widehat{ds}^{\,2}(e)\ \text{per edge} & \mathbb Z\setminus\{0\},\ \neg\text{a statistics}\end{array}
\]

%dsp-level:0
\[
(5)\quad\textbf{the naming discipline:}
\]

%dsp-level:0
\[
\qquad \textbf{``fermion''/``boson'' occur in this corpus in two distinct uses, and the tier keeps them apart:}
\]

%dsp-level:0
\[
(\mathrm i)\quad\textbf{as statistics labels:}\ \ \text{reads of}\ \varepsilon\ [(1)];
\]

%dsp-level:0
\[
\qquad (\mathrm{ii})\quad\textbf{as content labels:}\ \ \mathrm{D1.5}\text{'s bound/transit vocabulary}\ [(3)];
\]

%dsp-level:0
\[
\boxed{\ \neg\text{a species name is attached to either use}\ —\ \text{the identification layer is}\ \mathrm{Def\text{-}BM.3},\ \mathrm{Def\text{-}BM.6}\ \text{and}\ \mathrm{Def\text{-}BM.14},\ \text{and}\ \neg\text{a clause of this tier names a particle}\ }
\]

%dsp-level:0
\[
\qquad [\mathrm{Def\text{-}BM.3}(7);\ \text{machine: the species-token scan}];
\]

%dsp-level:0
\[
(6)\quad\textbf{what is complete and what is open:}
\]

%dsp-level:0
\[
\qquad \textbf{complete:}\ \varepsilon\ \text{is an exact function of a ledger datum, with no external input}\ [(1),(2)];
\]

%dsp-level:0
\[
\textbf{beyond this tier:}\qquad\begin{array}{l|l}\text{the winding of a named sealed class} & (\mathrm{R3})\ /\ \text{deficit}\ 1':\ \text{which three of the six states a class occupies, and which winding each constituent carries}\ [\mathrm{BM.28}(7);\ \mathrm{BM.103}(8);\ \mathrm{BM.117}(6)]\\\text{the connector winding}\ d=\Theta(P) & \mathrm{BM.2}(2):\ d\ \text{is registration data; an open connector carries}\ S=0\ \text{and hence}\ d=0\ [\mathrm{Def\text{-}H.32}(4);\ \mathrm{BM.128}(7)]\\\text{the}\ k\ge3\ \text{aspect claim} & \mathrm{Def\text{-}H.38}\ \text{Scope: carriers for}\ k\ge3\ \text{exceed the registered pair-channel inventory}\ [\mathrm{CT0.12}]\end{array}
\]

%dsp-level:0
\[
\boxed{\ \textbf{the identification is complete as a function and open as a value per named class}\ }
\]

%dsp-level:0
\[
(7)\quad\textbf{status:}\qquad\textbf{derived:}\ (1)\text{--}(5),\ (6)\text{'s completeness row};
\]

%dsp-level:0
\[
\textbf{re-scoped at its own locus,}\ \neg\textbf{retracted:}\quad\mathrm{BM.106}(1)\text{'s spin-parity refutation stands and is what makes}\ (4)\ \text{a separation rather than a contradiction};
\]

%dsp-level:0
\[
\neg\text{a declaration},
\]

%dsp-level:0
\[
\quad\neg\text{a P4 addition},
\]

%dsp-level:0
\[
\quad\neg\text{a new dimensionful scale}\ [\mathrm{P1}],
\]

%dsp-level:0
\[
\quad\neg\text{a measured value}\ [\mathrm{P5}],
\]

%dsp-level:0
\[
\quad\neg\text{a species identification}\ [\mathrm{Def\text{-}BM.3}(7)];
\]

\textbf{Req.} D0.8, D1.4, D1.5 (Parts 1–2); CT0.12 (Part 1); P1, P4, P5 (Parts 1–2); Def-H.32(4), Def-H.38, FT.22, FT.24, FT.80(5) (Part 3); Census FT.28(2) (row-wise); Def-BM.3(7), Def-BM.6, Def-BM.14, BM.2(2), BM.28(7), BM.103(8), BM.106(1), BM.113(2), BM.113(4), BM.117(6), BM.128(7) (this part); exhaustive parity and registration censuses, symbol-set censuses [math; machine: \texttt{checks/\allowbreak{}verify\_\allowbreak{}H\_\allowbreak{}janus\_\allowbreak{}tunnel\_\allowbreak{}concrete\_\allowbreak{}spin.py} §D].

\textbf{Pf.}

%dsp-level:0
\[
(1):\ \mathrm{BM.113}(2),(4)\ \text{re-cited (frozen)};\qquad\text{the parity census to}\ k=8\ [\text{machine}]
\]

%dsp-level:0
\[
(2):\ \Theta=\sum_e\widehat{ds}^{\,2}(e)\ [\mathrm{FT.22}];
\]

%dsp-level:0
\[
\qquad \widehat{ds}^{\,2}(e)=\kappa_e\ \text{registration data}\ [\mathrm{FT.24};\ \text{Census FT.28}(2);\ \text{A §0.2a}];
\]

%dsp-level:0
\[
\qquad \Theta_{\mathrm{int}}=2\kappa\implies k=\kappa\ [\mathrm{BM.106}(1);\ \text{machine}]
\]

%dsp-level:0
\[
\text{the symbol census of}\ (-1)^{|k|}:\ \{k\}\ [\text{machine}]
\]

%dsp-level:0
\[
(3):\ \mathrm{D1.5},\ \mathrm{D1.4},\ \mathrm{FT.80}(5)\ \text{re-cited}\ [\text{a definitional label and the closed/open read}]
\]

%dsp-level:0
\[
(4):\ \kappa=1\implies k=1,\ j=\tfrac12,\ \varepsilon=(-1)^1=-1;
\]

%dsp-level:0
\[
\qquad \kappa=2\implies k=2,\ j=1,\ \varepsilon=(-1)^2=+1\ [\mathrm{BM.106}(1);\ \text{machine: the two exact rows}]
\]

%dsp-level:0
\[
\text{both chains are materialized closed bound chains of one}\ \mathrm{D1.5}\ \text{kind}\ [\mathrm{Def\text{-}H.4};\ \mathrm{BM.106}(1)]\implies\text{the kind does not determine the sign}
\]

%dsp-level:0
\[
(5):\ \text{the two uses enumerated};
\]

%dsp-level:0
\[
\qquad \mathrm{Def\text{-}BM.3}(7)\text{'s scope row re-cited};
\]

%dsp-level:0
\[
\qquad \text{the species-token scan of this tier}\ [\text{machine}]
\]

%dsp-level:0
\[
(6):\ \text{the three open rows cite their own tags}\ [\text{machine}];
\]

%dsp-level:0
\[
\qquad (7):\ \text{the standing-rule scans}\ [\text{machine}].\ \blacksquare
\]

\textbf{Phys.} The statistics sign is a parity of a registration integer, and that is the whole content of the identification: the winding is a sum of per-edge line elements, each of which is registration data fixed at crystallization, so its parity is a ledger read with no postulate, no declared assignment and no measurement behind it. What the tier is careful about is that this read is not the same object as the bound-or-transit label the corpus has carried since its definition layer, and the proof that they differ is the floor census: at registration one a chain has spin one half and sign minus one, at registration two it has spin one and sign plus one, and both are materialized closed bound chains of the same kind. So the kind label cannot be a statistics and the statistics cannot be a kind, and the earlier refutation of a spin-parity theorem on the kind is exactly what makes this a clean separation instead of a contradiction. The two words fermion and boson are then used in two disciplined ways — as reads of the exchange sign and as the content vocabulary — with no species name attached to either. What is complete is the function; what stays open is its value on any named sealed class, because which windings a class occupies is the same missing datum as the flavour composition, together with the connector winding and the aspect bar at high winding.

\end{theorem}

\begin{theorem}[{\texorpdfstring{BM.128 — The three-chain exchange structure: the space, the $S_3$ action, the character formula, and the vanishing wedge}{BM.128 - The three-chain exchange structure: the space, the action, the character formula, and the vanishing wedge}}]\label{tag:BM-128}
\textbf{Thm.}

%dsp-level:0
\[
(1)\quad\textbf{the three-chain state space:}
\]

%dsp-level:0
\[
\qquad \boxed{\ \mathcal H_3(c_1,c_2,c_3)\ :=\ \bigotimes_{a=1}^{3}\Bigl(\mathcal H_\omega(c_a)\otimes\Psi(c_a)\Bigr),\qquad\dim\mathcal H_3\ =\ 8\prod_{a=1}^{3}\bigl(|k_a|+1\bigr)\ }
\]

%dsp-level:0
\[
[\mathrm{Def\text{-}BM.16}(2)\ \text{extended to}\ n=3;\ \text{machine: exact integers}];\qquad\textbf{the anchors:}\qquad\begin{array}{c|c|c}k=(k_1,k_2,k_3) & \prod(|k_a|+1) & \dim\mathcal H_3\\\hline(1,1,1) & 8 & 64\\(2,2,2) & 27 & 216\\(1,1,2) & 12 & 96\\(1,2,3) & 24 & 192\end{array}
\]

%dsp-level:0
\[
(2)\quad\textbf{the}\ S_3\ \textbf{action:}
\]

%dsp-level:0
\[
\qquad \hat P_\sigma\ :=\ \text{the factor permutation by}\ \sigma\in S_3,
\]

%dsp-level:0
\[
\qquad |S_3|=6,
\]

%dsp-level:0
\[
\qquad \boxed{\ \hat P_\sigma\hat P_\tau=\hat P_{\sigma\tau}\ \ \forall\sigma,\tau\in S_3\ }
\]

%dsp-level:0
\[
\qquad [\text{machine: all}\ 36\ \text{ordered pairs, entrywise over}\ \mathbb Q(i)]:
\]

%dsp-level:0
\[
\textbf{the orders:}\qquad\operatorname{ord}(\hat P_\sigma)=\operatorname{ord}(\sigma)=\begin{cases}1 & \sigma=\mathrm{id}\\ 2 & \sigma\ \text{a transposition}\ (3\ \text{of them})\\ 3 & \sigma\ \text{a}\ 3\text{-cycle}\ (2\ \text{of them})\end{cases}\qquad[\text{machine: exact matrix powers}];
\]

%dsp-level:0
\[
\textbf{the sign representation:}
\]

%dsp-level:0
\[
\qquad \operatorname{sgn}:S_3\to\{\pm1\},
\]

%dsp-level:0
\[
\qquad \ker(\operatorname{sgn})=A_3,\ |A_3|=3,
\]

%dsp-level:0
\[
\qquad S_3/A_3\cong\mathbb Z_2
\]

%dsp-level:0
\[
\qquad [\text{machine: the inversion census}\ (0,1,1,2,2,3)];
\]

%dsp-level:0
\[
(3)\quad\textbf{the character formula (exact):}
\]

%dsp-level:0
\[
\qquad \boxed{\ \chi(\sigma)\ :=\ \operatorname{Tr}\hat P_\sigma\big|_{(\mathbb C^d)^{\otimes n}}\ =\ d^{\,c(\sigma)}\ }
\]

%dsp-level:0
\[
\qquad c(\sigma)=\text{the cycle count of}\ \sigma
\]

%dsp-level:0
\[
[\text{machine: verified on all}\ 2!\ \text{permutations at}\ d\in\{2,3\}\ \text{and all}\ 3!\ \text{at}\ d\in\{2,3\}];
\]

%dsp-level:0
\[
\qquad \textbf{the three rows at}\ n=3,\ d=2:
\]

%dsp-level:0
\[
\qquad \chi(\mathrm{id})=2^3=8,
\]

%dsp-level:0
\[
\qquad \chi(\text{transposition})=2^2=4,
\]

%dsp-level:0
\[
\qquad \chi(3\text{-cycle})=2^1=2;
\]

%dsp-level:0
\[
\textbf{the consequence for the exchange sign:}
\]

%dsp-level:0
\[
\qquad \text{a transposition's trace is}\ d^{\,n-1},\ \neg d^{\,n}:
\]

%dsp-level:0
\[
\qquad \boxed{\ \textbf{the flip is not the identity on the}\ n\textbf{-fold space for}\ n\ge2\ }
\]

%dsp-level:0
\[
\qquad [\text{machine}];
\]

%dsp-level:0
\[
(4)\quad\textbf{the CP-twisted operator:}
\]

%dsp-level:0
\[
\qquad \hat P^{\mathrm{CP}}_\sigma\ :=\ \hat\mu_{\mathrm{CP}}\circ\hat P_\sigma\ \ \text{on}\ \mathcal H_3,
\]

%dsp-level:0
\[
\qquad \hat\mu_{\mathrm{CP}}=\mu_{\mathrm{CP}}\,\mathrm{id},
\]

%dsp-level:0
\[
\qquad \mu_{\mathrm{CP}}=e^{i\pi\Theta}=1\ \ \text{on}\ \Theta\in2\mathbb Z
\]

%dsp-level:0
\[
\qquad [\mathrm{BM.112}(1);\ \mathrm{BM.108}(3)]:
\]

%dsp-level:0
\[
\boxed{\ \bigl(\hat P^{\mathrm{CP}}_{(12)}\bigr)^2=\mathrm{id}\ \ \text{unconditionally},\qquad\operatorname{spec}\hat P^{\mathrm{CP}}_{(12)}\subseteq\{+1,-1\},\qquad\operatorname{Tr}\hat P^{\mathrm{CP}}_{(12)}=\mu_{\mathrm{CP}}\,d^{\,c((12))}\prod_{a}(|k_a|+1)\ }
\]

%dsp-level:0
\[
[\mathrm{BM.112}(1),(2)\ \text{transported to}\ n=3;\ \text{machine}];
\]

%dsp-level:0
\[
(5)\quad\textbf{the vanishing wedge (the structural row of this tag):}
\]

%dsp-level:0
\[
\qquad \dim\wedge^{3}(\mathbb C^d)=\binom d3:
\]

%dsp-level:0
\[
\qquad \boxed{\ \dim\wedge^{3}(\mathbb C^2)=\binom23=0,\qquad\dim\wedge^{3}(\mathbb C^3)=\binom33=1\ }
\]

%dsp-level:0
\[
[\text{machine: the explicit antisymmetrizer}\ \tfrac1{3!}\sum_{\sigma}\operatorname{sgn}(\sigma)\hat P_\sigma\ \text{has rank}\ 0\ \text{on}\ (\mathbb C^2)^{\otimes3}\ \text{and rank}\ 1\ \text{on}\ (\mathbb C^3)^{\otimes3},\ \text{computed exactly over}\ \mathbb Q(i)]:
\]

%dsp-level:0
\[
\textbf{the reading:}
\]

%dsp-level:0
\[
\qquad \boxed{\ \textbf{three copies of the}\ 2\textbf{-dimensional seed carrier admit no totally antisymmetric carrier state, while three colour slots admit exactly one}\ —\ \text{the}\ \varepsilon\text{-seal}\ }
\]

%dsp-level:0
\[
(6)\quad\textbf{the sharpening of}\ \mathrm{BM.114}(2):
\]

%dsp-level:0
\[
\qquad B=\varepsilon_{ijk}q^iq^jq^k\ \ \text{is antisymmetric in the colour slots}\ [\mathrm{D3.9};\ \mathrm{CT3.15};\ \mathrm{CT3.17}]:
\]

%dsp-level:0
\[
\textbf{the two antisymmetries are different objects with different dimensions:}\qquad\begin{array}{l|l|l}\text{antisymmetry} & \text{its space} & \text{its dimension}\\\hline\text{the colour seal} & \wedge^3(\mathbb C^3_{\text{colour}}) & 1\ \ [\text{machine}]\\\text{the carrier wedge} & \wedge^3(\mathbb C^2_{\mathsf v}) & 0\ \ [\text{machine}]\\\text{the statistics sign} & \text{the winding parity on}\ \mathcal H_3 & \{+1,-1\}\ \ [\mathrm{BM.113};\ (7)]\end{array}
\]

%dsp-level:0
\[
\boxed{\ \textbf{the baryon's total antisymmetry is carried by colour, which is}\ \neg\textbf{observable}\ [\mathrm{CT3.22}],\ \textbf{and the statistics is carried by the winding}\ —\ \neg\textbf{one is the other}\ }
\]

%dsp-level:0
\[
(7)\quad\textbf{the multiplicative exchange law:}
\]

%dsp-level:0
\[
\qquad k(\Gamma)=k(c_1)+k(c_2)+d_{12}
\]

%dsp-level:0
\[
\qquad [\mathrm{FT.34}(3);\ \mathrm{BM.2}(2):\ \text{the additive winding law}]:
\]

%dsp-level:0
\[
\textbf{the connector row:}
\]

%dsp-level:0
\[
\qquad d_{12}=\Theta(P)\ge0\ \ \text{for a closed connector};
\]

%dsp-level:0
\[
\qquad \textbf{an open chain carries}\ S=0\implies d=0
\]

%dsp-level:0
\[
\qquad [\mathrm{Def\text{-}H.32}(4);\ \text{machine}];
\]

%dsp-level:0
\[
\textbf{at a sealed junction with an open connector:}
\]

%dsp-level:0
\[
\qquad \boxed{\ k_\Gamma=\sum_{a=1}^{n}k_a\implies\varepsilon(\Gamma)=(-1)^{\left|\sum_ak_a\right|}=\prod_{a=1}^{n}(-1)^{|k_a|}=\prod_{a=1}^{n}\varepsilon(c_a)\ }
\]

%dsp-level:0
\[
[\text{machine: the multiplicativity verified on all}\ 5^3=125\ \text{triples at}\ n=3\ \text{and all}\ 5^2=25\ \text{pairs at}\ n=2];
\]

%dsp-level:0
\[
(8)\quad\textbf{the closure of}\ (\mathrm{R7})\ \textbf{at}\ n=3:\qquad\mathrm{BM.112}(7)\ \text{and}\ \mathrm{Def\text{-}BM.16}(7)\ \text{left}\ \mathcal H_3\ \text{and its}\ S_3\ \text{action}\ \neg\text{constructed}:
\]

%dsp-level:0
\[
\boxed{\ (\mathrm{R7})\ \ \textbf{is discharged at}\ n=3:\quad\text{the space}\ (1),\ \text{the group action with its character}\ (2),(3),\ \text{the involution}\ (4)\ \text{and the permutation-sign read}\ (7)\ \textbf{are all constructed}\ }
\]

%dsp-level:0
\[
\textbf{the residue narrows:}
\]

%dsp-level:0
\[
\qquad (\mathrm{R7}')\ \ \mathcal H_n\ \text{and its}\ S_n\ \text{action for}\ n\ge4\ \ \text{are}\ \neg\textbf{constructed}
\]

%dsp-level:0
\[
\qquad [\text{the sealing grammar admits}\ n\ge4\ \text{supports where}\ \delta\text{- and}\ \varepsilon\text{-groupings mix}\ [\mathrm{BM.31}(2);\ \mathrm{BM.110}(5)]];
\]

%dsp-level:0
\[
(9)\quad\textbf{status:}\qquad\textbf{derived:}\ (1)\text{--}(7);\qquad\textbf{closed:}\ (\mathrm{R7})\ \text{at}\ n=3;\qquad\textbf{named:}\ (\mathrm{R7}');
\]

%dsp-level:0
\[
\textbf{re-scoped at its own locus,}\ \neg\textbf{retracted:}\quad\mathrm{BM.114}(2)\text{'s colour-antisymmetry row, sharpened by}\ (6);
\]

%dsp-level:0
\[
\neg\text{a declaration},
\]

%dsp-level:0
\[
\quad\neg\text{a P4 addition},
\]

%dsp-level:0
\[
\quad\neg\text{a new dimensionful scale}\ [\mathrm{P1}],
\]

%dsp-level:0
\[
\quad\neg\text{a measured value}\ [\mathrm{P5}],
\]

%dsp-level:0
\[
\quad\neg\text{a species identification}\ [\mathrm{Def\text{-}BM.3}(7)];
\]

%dsp-level:0
\[
\neg\text{a flavour}\ \mathsf{SU}(3)\ \text{or}\ \mathsf{SU}(6)\ \text{is constructed:}\ S_3\ \text{permutes}\ \textbf{constituent slots},
\]

%dsp-level:0
\[
\ \neg\text{flavour states}\ [\mathrm{BM.117}(2),(3)];
\]

\textbf{Req.} A3, D1.4 (Part 1); CT3.15, CT3.17, CT3.22, D3.9 (Part 2); P1, P4, P5 (Parts 1–2); Def-H.32(4), Def-H.33, Def-H.38, FT.34(3), FT.108(4) (Part 3); Def-BM.3(7), Def-BM.16(2), Def-BM.16(7), BM.2(2), BM.31(2), BM.108(3), BM.110(5), BM.112(1), BM.112(2), BM.112(7), BM.113, BM.114(2), BM.117(2), BM.117(3) (this part); exact matrix algebra over $\mathbb Q(i)$, exact rank by Fraction elimination, permutation-group enumeration, exhaustive multiplicativity censuses [math; machine: \texttt{checks/\allowbreak{}verify\_\allowbreak{}H\_\allowbreak{}janus\_\allowbreak{}tunnel\_\allowbreak{}concrete\_\allowbreak{}spin.py} §E].

\textbf{Pf.}

%dsp-level:0
\[
(1):\ \dim\mathcal H_\omega
\]

%dsp-level:0
\[
= 2\ \text{and}\ \dim\Psi
\]

%dsp-level:0
\[
= |k|+1\ \text{per constituent}\ [\mathrm{Def\text{-}BM.16}(2);\ \mathrm{Def\text{-}H.38}]\implies\dim\mathcal H_3
\]

%dsp-level:0
\[
= 2^3\prod_a(|k_a|+1)\ [\text{machine: the four anchor rows}]
\]

%dsp-level:0
\[
(2):\ \hat P_\sigma\ \text{built on the}\ d^3\ \text{basis}\ \{e_{i_1}\otimes e_{i_2}\otimes e_{i_3}\}\ \text{by}\ e_{i_1}\otimes e_{i_2}\otimes e_{i_3}
\]

%dsp-level:0
\[
\mapsto e_{i_{\sigma^{-1}(1)}}\otimes\cdots\ [\text{machine}];
\]

%dsp-level:0
\[
\qquad \hat P_\sigma\hat P_\tau=\hat P_{\sigma\tau}\ \text{on all}\ 36\ \text{pairs}\ [\text{machine: entrywise over}\ \mathbb Q(i)]
\]

%dsp-level:0
\[
\text{the order census}\ (1,2,2,3,3,2)\ \text{and the inversion census}\ (0,1,1,2,2,3)\ [\text{machine}];
\]

%dsp-level:0
\[
\qquad |A_3|=3,\ 6/3=2\ [\text{machine}]
\]

%dsp-level:0
\[
(3):\ \text{the basis element}\ e_{i_1}\otimes e_{i_2}\otimes e_{i_3}\ \text{is fixed by}\ \hat P_\sigma
\]

%dsp-level:0
\[
\iff i\ \text{is constant on each cycle of}\ \sigma\implies\operatorname{Tr}
\]

%dsp-level:0
\[
= d^{c(\sigma)}\ [\text{machine: verified on all permutations at}\ (n,d)\in\{2,3\}\times\{2,3\}]
\]

%dsp-level:0
\[
n=3,\ d=2:\ c(\mathrm{id})=3\implies8;\quad c(\text{transposition})=2\implies4;\quad c(3\text{-cycle})=1\implies2\ [\text{machine}]
\]

%dsp-level:0
\[
(4):\ \hat\mu_{\mathrm{CP}}\ \text{is central}\implies(\hat\mu_{\mathrm{CP}}\hat P_{(12)})^2=\mu_{\mathrm{CP}}\mu_{\mathrm{CP}}^\ast\hat P_{(12)}^2=\mathrm{id}\ [\mathrm{BM.112}(1);\ \text{machine}];
\]

%dsp-level:0
\[
\qquad \operatorname{Tr}=\mu_{\mathrm{CP}}\chi((12))\prod_a(|k_a|+1)\ [(3);\ \text{machine}]
\]

%dsp-level:0
\[
(5):\ \dim\wedge^3(\mathbb C^d)=\binom d3\ [\text{math}];
\]

%dsp-level:0
\[
\qquad \binom23=0,\ \binom33=1\ [\text{machine}];
\]

%dsp-level:0
\[
\qquad \text{the explicit antisymmetrizer}\ \tfrac1{6}\sum_\sigma\operatorname{sgn}(\sigma)\hat P_\sigma\ \text{has rank}\ 0\ \text{at}\ d
\]

%dsp-level:0
\[
= 2\ \text{and rank}\ 1\ \text{at}\ d
\]

%dsp-level:0
\[
= 3\ [\text{machine: exact Fraction elimination on the}\ 8\times8\ \text{and}\ 27\times27\ \text{matrices}]
\]

%dsp-level:0
\[
(6):\ \mathrm{D3.9},\ \mathrm{CT3.15},\ \mathrm{CT3.17},\ \mathrm{CT3.22}\ \text{re-cited};\qquad\text{the three-row census}\ [(5);\ \mathrm{BM.113}(4);\ \text{machine}]
\]

%dsp-level:0
\[
(7):\ k(\Gamma)=\sum_ak_a+d\ [\mathrm{FT.34}(3);\ \mathrm{BM.2}(2)];\qquad\text{an open chain carries}\ S=0\ [\mathrm{Def\text{-}H.32}(4)]\implies d=0;
\]

%dsp-level:0
\[
(-1)^{|k_1+k_2+k_3|}
\]

%dsp-level:0
\[
= (-1)^{k_1+k_2+k_3}\ \text{on}\ k_a
\]

%dsp-level:0
\[
\ge 0\ [\mathrm{Def\text{-}H.38}\ \text{Scope}]
\]

%dsp-level:0
\[
= (-1)^{k_1}(-1)^{k_2}(-1)^{k_3}\ [\text{machine: all}\ 125\ \text{triples in}\ \{0,\dots,4\}^3\ \text{and all}\ 25\ \text{pairs}]
\]

%dsp-level:0
\[
(8):\ \mathrm{BM.112}(7),\ \mathrm{Def\text{-}BM.16}(7)\ \text{re-cited};
\]

%dsp-level:0
\[
\qquad \text{the four members discharged against}\ \mathrm{BM.106}(7)\text{'s}\ (\mathrm{R7})\ [\text{machine: the row-by-row check}];
\]

%dsp-level:0
\[
\qquad (9):\ \text{the standing-rule scans}\ [\text{machine}].\ \blacksquare
\]

\textbf{Phys.} Three chains now have a space, a group and a character, and the three fit together exactly. The space is the tensor product of the three orientation doublets and the three winding carriers, eight times the product of the carrier dimensions, which at the two anchors is sixty-four and two hundred sixteen. The symmetric group on three letters acts by permuting factors, the operators multiply as the permutations do on all thirty-six ordered pairs, and the traces obey one clean formula: the dimension to the number of cycles, giving eight, four and two at three chains over a two-dimensional carrier. Twisting by the CP phase changes nothing on the materialized locus, where the phase is one, and the transposition remains an unconditional involution. The structural surprise is the wedge: three copies of the two-dimensional seed carrier have no totally antisymmetric state at all, because the binomial coefficient of three taken from two is zero, while three colour slots have exactly one, and an explicit antisymmetrizer confirms both ranks over the exact complex field. That is why the baryon's antisymmetry is colour's and not the carrier's — and colour is unobservable by construction, so the statistics has to live somewhere else, which is the winding. The winding then supplies the law the tier was after: at a sealed junction with an open connector the total winding is the sum of the constituent windings, so the exchange sign of a composite is the product of its constituents' signs, verified on every triple in a five-by-five-by-five census. That closes the exchange-structure residue at three chains, and narrows it to four and beyond, where the sealing grammar starts mixing delta and epsilon groupings.

\end{theorem}

\begin{theorem}[{\texorpdfstring{BM.129 — The meson is a boson and the baryon's sign is the product: the composite identification}{BM.129 - The meson is a boson and the baryon's sign is the product: the composite identification}}]\label{tag:BM-129}
\textbf{Thm.}

%dsp-level:0
\[
(1)\quad\textbf{the meson:}
\]

%dsp-level:0
\[
\qquad \text{a}\ \delta\text{-sealed CP-pair with an open connector:}\quad k_M=k(c)+k(c^{\mathrm{CP}})+d=k+k+0=2k
\]

%dsp-level:0
\[
\qquad [\mathrm{BM.128}(7);\ \mathrm{BM.113}(2)(\mathrm{ii}):\ \Theta(c^{\mathrm{CP}})=\Theta(c)\ \text{on the materialized locus}]:
\]

%dsp-level:0
\[
\boxed{\ \varepsilon_M\ =\ (-1)^{|2k|}\ =\ \bigl((-1)^2\bigr)^{k}\ =\ +1\qquad\textbf{for every constituent winding}\ k\ge0\ }
\]

%dsp-level:0
\[
\qquad [\text{machine: verified at every}\ k\in\{0,\dots,8\}];
\]

%dsp-level:0
\[
j_M=\frac{|k_M|}2=\frac{2k}2=k\in\mathbb Z
\]

%dsp-level:0
\[
\qquad [\text{machine}]:
\]

%dsp-level:0
\[
\qquad \boxed{\ \textbf{the meson's total spin is an integer and its exchange sign is}\ +1\ —\ \textbf{it is a boson, unconditionally}\ }
\]

%dsp-level:0
\[
(2)\quad\textbf{the robustness row (the orientation-sign reading):}
\]

%dsp-level:0
\[
\qquad \text{the alternative reading}\ k(c^{\mathrm{CP}})=-k(c)\ \ [\hat\omega:\ k\mapsto-k\ \text{on the carrier};\ \mathrm{BM.73}(4);\ \mathrm{A\text{-}BM.7}(b)]:
\]

%dsp-level:0
\[
k_M=k+(-k)+0=0\implies j_M=0,\qquad\varepsilon_M=(-1)^0=+1\qquad[\text{machine}];
\]

%dsp-level:0
\[
\boxed{\ \textbf{both readings give}\ \varepsilon_M=+1\ —\ \textbf{the meson's statistics is robust under the unresolved orientation-sign row, while its total spin is not}\ }
\]

%dsp-level:0
\[
\qquad [j_M=k\ \text{against}\ j_M=0;\ \text{the spin value is}\ \neg\text{fixed and the statistics is}];
\]

%dsp-level:0
\[
(3)\quad\textbf{the closed-connector variant:}
\]

%dsp-level:0
\[
\qquad k_M=2k+d\implies\boxed{\ \varepsilon_M=(-1)^{d}\ }
\]

%dsp-level:0
\[
\qquad [\text{machine: the two constituent windings cancel modulo two}]:
\]

%dsp-level:0
\[
\textbf{the meson's statistics is then the connector's parity alone:}
\]

%dsp-level:0
\[
\qquad d\ \text{even}\implies+1,
\]

%dsp-level:0
\[
\qquad d\ \text{odd}\implies-1
\]

%dsp-level:0
\[
\qquad [\text{machine:}\ d\in\{0,1,2,3,4\}];
\]

%dsp-level:0
\[
d=\Theta(P)\ge0\ \ \text{is registration data}
\]

%dsp-level:0
\[
\qquad [\mathrm{BM.2}(2);\ \mathrm{FT.34}(3)];
\]

%dsp-level:0
\[
\qquad \textbf{the computed cells:}\ \ \text{an open connector carries}\ S=0\implies d=0\implies\varepsilon_M=+1
\]

%dsp-level:0
\[
\qquad [\mathrm{Def\text{-}H.32}(4);\ \text{machine}];
\]

%dsp-level:0
\[
(4)\quad\textbf{the baryon:}\qquad\text{an}\ \varepsilon\text{-sealed 3-juncture with an open connector:}\quad k_B=k_1+k_2+k_3\qquad[\mathrm{BM.128}(7)]:
\]

%dsp-level:0
\[
\boxed{\ \varepsilon_B\ =\ (-1)^{k_1+k_2+k_3}\ =\ \varepsilon(c_1)\,\varepsilon(c_2)\,\varepsilon(c_3),\qquad j_B=\frac{|k_1+k_2+k_3|}2\ }
\]

%dsp-level:0
\[
\textbf{the exact criterion:}
\]

%dsp-level:0
\[
\qquad \boxed{\ \varepsilon_B=-1\iff\#\{a:\ k_a\ \text{odd}\}\ \text{is odd}\ }
\]

%dsp-level:0
\[
\qquad [\text{machine: all}\ 125\ \text{triples in}\ \{0,\dots,4\}^3;\ 62\ \text{fermionic and}\ 63\ \text{bosonic}];
\]

%dsp-level:0
\[
(5)\quad\textbf{the anchor and the computed floors:}
\]

%dsp-level:0
\[
\qquad \textbf{all three constituents at the minimal odd winding}\ k_a=1:
\]

%dsp-level:0
\[
\qquad k_B=3,
\]

%dsp-level:0
\[
\qquad j_B=\tfrac32,
\]

%dsp-level:0
\[
\qquad \varepsilon_B=-1
\]

%dsp-level:0
\[
\qquad [\text{machine}]:
\]

%dsp-level:0
\[
\textbf{the digon floor:}\quad L=2\implies\kappa\equiv1,\ k=L/2=1\ \ (\text{odd})
\]

%dsp-level:0
\[
\qquad [\mathrm{FT.32}(6);\ \mathrm{FT.28};\ \text{machine}]:
\]

%dsp-level:0
\[
\qquad \boxed{\ \textbf{the minimal materialized constituent carries}\ k=1,\ \text{so the three-floor baryon is fermionic}\ }
\]

%dsp-level:0
\[
\textbf{the cycle ladder's parity:}\qquad L\ \text{even}\implies k=L/2;\qquad L\ \text{odd}\implies k=L\qquad[\mathrm{FT.32}(6);\ \text{machine}]:
\]

%dsp-level:0
\[
\begin{array}{c|c|c|c}L & k & k\ \text{parity} & \varepsilon\\\hline2 & 1 & \text{odd} & -1\\3 & 3 & \text{odd} & -1\\4 & 2 & \text{even} & +1\\5 & 5 & \text{odd} & -1\\6 & 3 & \text{odd} & -1\\8 & 4 & \text{even} & +1\end{array}\qquad[\text{machine: exact}];
\]

%dsp-level:0
\[
(6)\quad\textbf{what is derived and what awaits a datum:}
\]

%dsp-level:0
\[
\qquad \textbf{derived:}\ \varepsilon_B\ \text{as an exact function of the three constituent windings}\ [(4)];
\]

%dsp-level:0
\[
\ \text{the multiplicative law}\ [\mathrm{BM.128}(7)];
\]

%dsp-level:0
\[
\ \text{the criterion}\ [(4)];
\]

%dsp-level:0
\[
\textbf{awaits}\ (\mathrm{R3})/\text{deficit}\ 1':
\]

%dsp-level:0
\[
\qquad \boxed{\ \textbf{which windings a named sealed class occupies is}\ \neg\textbf{derived}\ —\ \text{so the baryon's statistics is complete as a function and open as a value per class}\ }
\]

%dsp-level:0
\[
[\mathrm{BM.28}(7);\ \mathrm{BM.103}(8);\ \mathrm{BM.117}(6);\ \mathrm{BM.127}(6);\ \text{machine: the symbol census of the criterion carries}\ \{k_1,k_2,k_3\}\ \text{and}\ \neg\text{a class name}];
\]

%dsp-level:0
\[
(7)\quad\textbf{the colour row and the octet row:}
\]

%dsp-level:0
\[
\qquad \textbf{the}\ \varepsilon\textbf{-seal's antisymmetry is colour's, one-dimensional, and}\ \neg\textbf{observable:}
\]

%dsp-level:0
\[
\qquad \dim\wedge^3(\mathbb C^3_{\text{colour}})=1,
\]

%dsp-level:0
\[
\qquad \dim\wedge^3(\mathbb C^2_{\mathsf v})=0
\]

%dsp-level:0
\[
\qquad [\mathrm{BM.128}(5),(6);\ \mathrm{CT3.22}];
\]

%dsp-level:0
\[
\textbf{no flavour multiplet is formed:}
\]

%dsp-level:0
\[
\qquad S_3\ \text{permutes constituent slots},\ \neg\text{flavour states};
\]

%dsp-level:0
\[
\qquad \neg\exists\ \text{a flavour}\ \mathsf{SU}(3);
\]

%dsp-level:0
\[
\qquad \text{the octet and decuplet groupings are}\ \neg\text{derived}
\]

%dsp-level:0
\[
\qquad [\mathrm{BM.117}(2),(3);\ \mathrm{BM.114}(4)];
\]

%dsp-level:0
\[
(8)\quad\textbf{status:}\qquad\textbf{derived:}\ (1)\text{--}(5),\ (7);\qquad\textbf{open, awaiting}\ (\mathrm{R3}):\ (6)\text{'s per-class value};
\]

%dsp-level:0
\[
\textbf{re-scoped at its own locus,}\ \neg\textbf{retracted:}\quad\mathrm{BM.117}(2),
\]

%dsp-level:0
\[
(3)\text{'s orbit-count rows stand and are}\ \neg\text{promoted to multiplets};
\]

%dsp-level:0
\[
\neg\text{a declaration},
\]

%dsp-level:0
\[
\quad\neg\text{a P4 addition},
\]

%dsp-level:0
\[
\quad\neg\text{a new dimensionful scale}\ [\mathrm{P1}],
\]

%dsp-level:0
\[
\quad\neg\text{a measured value}\ [\mathrm{P5}],
\]

%dsp-level:0
\[
\quad\neg\text{a numerical coupling}\ [\mathrm{EM.25}(7)],
\]

%dsp-level:0
\[
\quad\neg\text{a species identification}\ [\mathrm{Def\text{-}BM.3}(7)];
\]

\textbf{Req.} A3, D1.4 (Part 1); CT3.22 (Part 2); P1, P4, P5 (Parts 1–2); Def-H.32(4), Def-H.38, FT.28, FT.32(6), FT.34(3) (Part 3); Def-BM.3(7), Def-BM.16(2), BM.2(2), BM.28(7), BM.73(4), BM.103(8), BM.113(2), BM.114(4), BM.117(2), BM.117(3), BM.117(6), BM.127(6), BM.128(5), BM.128(6), BM.128(7), A-BM.7 (this part); exhaustive parity and triple censuses, exact integer winding arithmetic, symbol-set censuses [math; machine: \texttt{checks/\allowbreak{}verify\_\allowbreak{}H\_\allowbreak{}janus\_\allowbreak{}tunnel\_\allowbreak{}concrete\_\allowbreak{}spin.py} §E].

\textbf{Pf.}

%dsp-level:0
\[
(1):\ k_M=k(c)+k(c^{\mathrm{CP}})+d\ [\mathrm{BM.128}(7)];
\]

%dsp-level:0
\[
\qquad \Theta(c^{\mathrm{CP}})=\Theta(c)\ \text{on the materialized locus}\ [\mathrm{BM.113}(2)(\mathrm{ii})]\implies k(c^{\mathrm{CP}})=k(c)=k;
\]

%dsp-level:0
\[
\qquad d=0\ \text{for an open connector}\ [\mathrm{Def\text{-}H.32}(4)]
\]

%dsp-level:0
\[
k_M=2k\implies\varepsilon_M=(-1)^{|2k|}=((-1)^2)^k=1\ [\text{machine: every}\ k\in\{0,\dots,8\}];
\]

%dsp-level:0
\[
\qquad j_M=|2k|/2=k\in\mathbb Z\ [\text{machine}]
\]

%dsp-level:0
\[
(2):\ k(c^{\mathrm{CP}})=-k(c)\ [\hat\omega:\ k\mapsto-k\ [\mathrm{BM.73}(4)]]\implies k_M=0\implies\varepsilon_M=(-1)^0=+1,\ j_M=0\ [\text{machine}];
\]

%dsp-level:0
\[
\qquad \text{the two readings agree on}\ \varepsilon_M\ \text{and differ on}\ j_M\ [\text{machine}]
\]

%dsp-level:0
\[
(3):\ k_M=2k+d\implies(-1)^{|2k+d|}=(-1)^{2k+d}=(-1)^{2k}(-1)^d=(-1)^d\ [\text{machine:}\ d\in\{0,\dots,4\}]
\]

%dsp-level:0
\[
(4):\ k_B=k_1+k_2+k_3\ [\mathrm{BM.128}(7)];
\]

%dsp-level:0
\[
\qquad (-1)^{k_1+k_2+k_3}=(-1)^{k_1}(-1)^{k_2}(-1)^{k_3}\ [\text{machine: all}\ 125\ \text{triples}];
\]

%dsp-level:0
\[
\qquad \#\{a:k_a\ \text{odd}\}\ \text{odd}\iff\varepsilon_B=-1\ [\text{machine:}\ 62\ \text{of}\ 125]
\]

%dsp-level:0
\[
(5):\ k_a=1\ \forall a\implies k_B=3,\ j_B=3/2,\ \varepsilon_B=(-1)^3=-1\ [\text{machine}];
\]

%dsp-level:0
\[
\qquad L=2\implies\kappa\equiv1,\ k=L/2=1\ [\mathrm{FT.32}(6);\ \mathrm{FT.28}];
\]

%dsp-level:0
\[
\qquad \text{the six-row ladder}\ [\text{machine}]
\]

%dsp-level:0
\[
(6):\ \mathrm{BM.28}(7),\ \mathrm{BM.103}(8),\ \mathrm{BM.117}(6),\ \mathrm{BM.127}(6)\ \text{re-cited};
\]

%dsp-level:0
\[
\qquad \text{the symbol census of}\ (-1)^{k_1+k_2+k_3}\ \text{carries}\ \{k_1,k_2,k_3\}\ \text{alone}\ [\text{machine}]
\]

%dsp-level:0
\[
(7):\ \mathrm{BM.128}(5),(6)\ \text{and}\ \mathrm{CT3.22}\ \text{re-cited};
\]

%dsp-level:0
\[
\qquad \mathrm{BM.117}(2),(3)\text{'s octet rows re-cited}\ [\neg\text{promoted}];
\]

%dsp-level:0
\[
\qquad (8):\ \text{the standing-rule scans}\ [\text{machine}].\ \blacksquare
\]

\textbf{Phys.} Two composites, two clean answers, and one honest gap. The meson is a boson and it is one unconditionally: its constituents are a chain and its conjugate, so their windings are equal and their sum is even, which makes the total spin an integer and the exchange sign plus one — and the conclusion survives the corpus's one unresolved orientation-sign convention, since the alternative reading puts the total winding at zero and the sign at plus one as well, differing only in the spin value. With a closed connector the answer changes to the connector's own parity, because the two constituent windings cancel modulo two, and on the computed cells the connector is open so the boson sign is what holds. The baryon is not unconditional and the tier does not pretend otherwise: its sign is the product of its three constituents' signs, so it is fermionic exactly when an odd number of them carry odd winding, and at the minimal materialized floor the winding is one — odd — which makes the three-floor baryon fermionic with total spin three halves. Sixty-two of the one hundred twenty-five winding triples in a small census are fermionic and sixty-three bosonic, so the criterion genuinely discriminates. What is missing is not the rule but its input: which windings a named sealed class occupies is the same unread datum as the flavour composition, so the baryon's statistics is complete as a function and open as a value per class. And the antisymmetry that makes a baryon a baryon is colour's — one-dimensional over three colour slots against a carrier wedge that vanishes at dimension two — unobservable by construction and not the statistics, with no flavour multiplet formed anywhere in the argument.

\end{theorem}

\medskip\noindent\textit{\texorpdfstring{BM.130 — The status of the Annihilation, Concrete Spin and Fermion Identification Tier: adjudication, scope, and the open items}{BM.130 - The status of the Annihilation, Concrete Spin and Fermion Identification Tier: adjudication, scope, and the open items}.}\ Structural-constraints record: see Appendix B, Section B.1.33 of the Addendum.\par\medskip

\phantomsection
\section[Tier BM.23 — The Closing Tier: the Haar–Poisson Duality and…]{\texorpdfstring{Tier BM.23 — The Closing Tier: the Haar–Poisson Duality and the Lapse, the Selected Curvature Composition, the Momentum-Space Resolvent, the Charge and Statistics Censuses, the Two-Chain Ensemble, the Seed-Tensoring Mechanism, the Convention Declarations, the Ceiling Study, the Jensen Residue, the Pair Amplitude, and the Era Closure}{Tier BM.23 - The Closing Tier: the Haar-Poisson Duality and the Lapse, the Selected Curvature Composition, the Momentum-Space Resolvent, the Charge and Statistics Censuses, the Two-Chain Ensemble, the Seed-Tensoring Mechanism, the Convention Declarations, the Ceiling Study, the Jensen Residue, the Pair Amplitude, and the Era Closure}}

\begin{theorem}[{\texorpdfstring{BM.131 — The Haar–Poisson duality: the twirl fixed point, the complementary factorization, and the proper-time contraction}{BM.131 - The Haar-Poisson duality: the twirl fixed point, the complementary factorization, and the proper-time contraction}}]\label{tag:BM-131}
\textbf{Thm.}

%dsp-level:0
\[
(1)\quad\textbf{the two evolutions, from the axioms:}
\]

%dsp-level:0
\[
\qquad \textbf{the Haar side:}\quad\mu_\star=\mu_{\mathrm{Haar}}\bigl(\mathrm{SU}(2)\bigr),
\]

%dsp-level:0
\[
\qquad \mathcal T(\rho)=U\rho U^\dagger,
\]

%dsp-level:0
\[
\qquad U\sim\mu_\star,
\]

%dsp-level:0
\[
\qquad U\notin\mathfrak L
\]

%dsp-level:0
\[
\qquad [\mathrm{A6}];
\]

%dsp-level:0
\[
\textbf{the Poisson side:}
\]

%dsp-level:0
\[
\qquad K:=-\log\rho
\]

%dsp-level:0
\[
\qquad [\mathrm{L1}:\ \text{the modular Hamiltonian}],
\]

%dsp-level:0
\[
\qquad \tilde{\mathcal T}_\gamma(\rho)=\frac{\rho\,e^{-\gamma\rho}}{\mathfrak T\bigl(\rho e^{-\gamma\rho}\bigr)},
\]

%dsp-level:0
\[
\qquad \gamma=1-4\varepsilon_q\in(0,1]
\]

%dsp-level:0
\[
\qquad [\mathrm{D0.3};\ \mathrm{L0}:\ \text{one tilt step per A4 event}];
\]

%dsp-level:0
\[
\textbf{the ledger increment factorizes:}
\]

%dsp-level:0
\[
\qquad \rho\ \longmapsto\ U\,\tilde{\mathcal T}_\gamma(\rho)\,U^\dagger
\]

%dsp-level:0
\[
\qquad [\mathrm{A6}:\ \tau_n\to\tau_{n+1};\ \mathrm{L0}:\ \text{the crystallization--tilt coupling, declared at the point of use}];
\]

%dsp-level:0
\[
(2)\quad\textbf{the exact K-update (the Poisson evolution):}
\]

%dsp-level:0
\[
\qquad \lambda_i'=\frac{\lambda_ie^{-\gamma\lambda_i}}{Z_\gamma},
\]

%dsp-level:0
\[
\qquad Z_\gamma=\textstyle\sum_j\lambda_je^{-\gamma\lambda_j}
\]

%dsp-level:0
\[
\qquad [\mathrm{L1}(3)\ \mathrm{Pf}:\ \text{the eigenvalue flow}]:
\]

%dsp-level:0
\[
\boxed{\ \delta K\ :=\ K'-K\ =\ \gamma\,\rho\ +\ \log Z_\gamma\cdot I\ }
\]

%dsp-level:0
\[
\qquad [\log\lambda_i'=\log\lambda_i-\gamma\lambda_i-\log Z_\gamma\ \text{exactly};\ \text{machine: the formal-series identity to order}\ 12];
\]

%dsp-level:0
\[
\gamma=1\ (\text{the sharp point}\ \varepsilon_q=0\ [\mathrm{D0.4}]):\qquad\delta K=\rho+\log Z\cdot I\qquad[\text{machine}];
\]

%dsp-level:0
\[
\textbf{the vacuum is the Poisson fixed point:}
\]

%dsp-level:0
\[
\qquad \rho=\rho_\circ=\tfrac14I:\quad Z_\gamma=e^{-\gamma/4},
\]

%dsp-level:0
\[
\qquad \lambda_i'=\tfrac14,
\]

%dsp-level:0
\[
\qquad \boxed{\ \delta K=0\ \ \text{exactly}\ }
\]

%dsp-level:0
\[
\qquad [\mathrm{L1}(1);\ \text{machine:}\ e^{-\gamma/4}\cdot e^{+\gamma/4}=1\ \text{coefficient-wise in}\ \mathbb Q[[\gamma]]\ \text{to order}\ 12];
\]

%dsp-level:0
\[
(3)\quad\textbf{the Haar invariances (the measure side):}
\]

%dsp-level:0
\[
\qquad V\bigl(\tilde K(\mathcal T(\rho))\bigr)=V\bigl(\tilde K(\rho)\bigr)\ \ \text{for every draw}
\]

%dsp-level:0
\[
\qquad [\mathrm{L1}(2);\ \text{machine: the spectral invariance}\ \operatorname{Tr}\bigl((U\rho U^\dagger)^n\bigr)=\operatorname{Tr}(\rho^n),\ n=1,2,3,\ \text{on}\ 40\ \text{exact rational}\ \mathrm{SU}(2)\ \text{points and generic rational states}];
\]

%dsp-level:0
\[
K\mapsto UKU^\dagger\implies\operatorname{spec}K\ \text{invariant};\qquad\text{the draws are unrecorded}\ [\mathrm{A6}:\ U\notin\mathfrak L];\qquad\neg\exists\,\tau^{-1}\qquad[\mathrm{A6}];
\]

%dsp-level:0
\[
\textbf{the stationarity of the measure:}
\]

%dsp-level:0
\[
\qquad \text{conjugation by a fixed}\ U_0\ \text{pushes}\ \mu_\star\ \text{to}\ \mu_\star
\]

%dsp-level:0
\[
\qquad [\text{math: the two-sided invariance and uniqueness of the Haar measure on a compact group}];
\]

%dsp-level:0
\[
(4)\quad\textbf{the twirl theorem (the duality's fixed locus):}
\]

%dsp-level:0
\[
\qquad \mathcal H_p\otimes\mathcal H_q=\vee^2\oplus\wedge^2
\]

%dsp-level:0
\[
\qquad [\mathrm{A3}],
\]

%dsp-level:0
\[
\qquad \Pi_{\mathcal R}=\tfrac12(I+F),
\]

%dsp-level:0
\[
\qquad \Pi_{\mathcal V}=\tfrac12(I-F),
\]

%dsp-level:0
\[
\qquad \operatorname{rk}\Pi_{\mathcal R}=3,
\]

%dsp-level:0
\[
\qquad \operatorname{rk}\Pi_{\mathcal V}=1
\]

%dsp-level:0
\[
\qquad [\mathrm{FT.107}(2);\ \text{machine: the flip}\ F^2=I]:
\]

%dsp-level:0
\[
\boxed{\ \int_{\mathrm{SU}(2)}\!\bigl(U\otimes U\bigr)\rho\,\bigl(U\otimes U\bigr)^\dagger d\mu_\star(U)\ =\ \frac{\mathfrak T\bigl(\Pi_{\mathcal R}\rho\bigr)}3\,\Pi_{\mathcal R}\ +\ \mathfrak T\bigl(\Pi_{\mathcal V}\rho\bigr)\,\Pi_{\mathcal V}\ }
\]

%dsp-level:0
\[
\qquad [\text{math: Schur's lemma on the}\ 3\oplus1\ \text{decomposition};\ \text{machine: idempotence, the block-trace preservation, and the equivariance}\ \mathrm{tw}\bigl((U\otimes U)\rho(U\otimes U)^\dagger\bigr)=(U\otimes U)\,\mathrm{tw}(\rho)\,(U\otimes U)^\dagger\ \text{on all}\ 40\ \text{exact rational points}];
\]

%dsp-level:0
\[
\textbf{the twirl fixed locus is the block-maximal-mixed family}\ \bigl\{\tfrac x3\Pi_{\mathcal R}+(1-x)\Pi_{\mathcal V}:\ x\in[0,1]\bigr\}
\]

%dsp-level:0
\[
\qquad [\text{machine}];
\]

%dsp-level:0
\[
\textbf{the unique common fixed point of both sides:}
\]

%dsp-level:0
\[
\qquad \lambda_i'=\lambda_i\ \forall i\iff e^{-\gamma\lambda_i}=Z_\gamma\ \forall i\iff\lambda_i\ \text{all equal}\ (\gamma>0)\iff\rho=\tfrac14I;
\]

%dsp-level:0
\[
\boxed{\ \Bigl\{\text{Poisson fixed points}\Bigr\}\cap\Bigl\{\text{Haar-twirl fixed points}\Bigr\}\ =\ \Bigl\{\tfrac14I\Bigr\}\ }
\]

%dsp-level:0
\[
\qquad [\text{machine:}\ x/3=1-x\iff x=\tfrac34\ \text{gives}\ \rho=\tfrac14I\ \text{on the block family};\ \text{the tilt fixed point is}\ \tfrac14I\ [\mathrm{L1}(1),(5)]];
\]

%dsp-level:0
\[
(5)\quad\textbf{the complementary factorization (what the record sees):}
\]

%dsp-level:0
\[
\qquad V\ \text{is invariant under the Haar factor}\ [(3)]\ \text{and strictly contracting under the Poisson factor:}
\]

%dsp-level:0
\[
\qquad V'=\Bigl(1-\frac\gamma4\Bigr)^{\!2}V+O(\delta\rho^3)
\]

%dsp-level:0
\[
\qquad [\mathrm{L2}];
\]

%dsp-level:0
\[
\Bigl(1-\frac\gamma4\Bigr)^{\!2}+\frac{\gamma(8-\gamma)}{16}=1
\]

%dsp-level:0
\[
\qquad [\mathrm{BM.51}(2);\ \text{machine: exact in}\ \mathbb Q[\gamma]]\implies V_n=\Bigl(1-\frac\gamma4\Bigr)^{2n}V_0+O
\]

%dsp-level:0
\[
\qquad [\mathrm{BM.51}(3)];
\]

%dsp-level:0
\[
\boxed{\ \textbf{the unrecorded Haar draws drop out of the record exactly:}\ \ V_n\ \text{depends on the tilt count alone}\ }
\]

%dsp-level:0
\[
\qquad [\mathrm{L1}(2)+\mathrm{L2};\ \text{machine}];
\]

%dsp-level:0
\[
\textbf{the reading:}
\]

%dsp-level:0
\[
\qquad \text{the measure side is the exact gauge freedom of the record};
\]

%dsp-level:0
\[
\qquad \text{the Poisson side is the recorded dynamics};
\]

%dsp-level:0
\[
\qquad \text{the duality fixed point is the A5 vacuum}\ \rho_\circ\ [\mathrm{A5}];
\]

%dsp-level:0
\[
(6)\quad\textbf{the finite-group companion (the same twirl at the carrier level):}
\]

%dsp-level:0
\[
\qquad \Pi_k:=\frac1{k!}\sum_{\sigma\in S_k}\hat P_\sigma
\]

%dsp-level:0
\[
\qquad [\text{the averaging over the finite group}\ S_k\ \text{is its Haar measure}]:
\]

%dsp-level:0
\[
\Pi_k^2=\Pi_k,
\]

%dsp-level:0
\[
\qquad \operatorname{rk}\Pi_k=k+1,
\]

%dsp-level:0
\[
\qquad \operatorname{im}\Pi_k=\mathrm{Sym}^k\bigl(\mathbb C^2\bigr)
\]

%dsp-level:0
\[
\qquad [\text{machine: exact rank over}\ \mathbb Q\ \text{for}\ k=1..6;\ \text{forward pointer: the mechanism is BM.138}];
\]

%dsp-level:0
\[
(7)\quad\textbf{the proper-time contraction (the metric read):}
\]

%dsp-level:0
\[
\qquad ds^2=-N^2d\tau^2+g_{ij}\bigl(dx^i+N^id\tau\bigr)\bigl(dx^j+N^jd\tau\bigr)
\]

%dsp-level:0
\[
\qquad [\mathrm{BM.34}(2);\ \text{machine: the line-element identity}];
\]

%dsp-level:0
\[
\text{at rest}\ \bigl(dx^i=0\bigr):\quad ds=N\,d\tau;
\]

%dsp-level:0
\[
\qquad \text{the tick is the chart coordinate}\ x^0=\tau_n
\]

%dsp-level:0
\[
\qquad [\mathrm{D1.8}];
\]

%dsp-level:0
\[
\qquad N=N(k,E)\ \text{is derived in}\ \mathrm{BM.132};
\]

%dsp-level:0
\[
\boxed{\ V(s)\ =\ \Bigl(1-\frac\gamma4\Bigr)^{2s/N}V_0\ }
\]

%dsp-level:0
\[
\qquad [\text{the reparametrization of}\ (5)\text{'s law by}\ s=Nn;\ \text{machine: at}\ \gamma=1,\ N=1/m,\ m\in\{1,2,3,4\}:\ \bigl(\tfrac9{16}\bigr)^{m}\ \text{exact integer powers, strictly decreasing in}\ m];
\]

%dsp-level:0
\[
\textbf{the two deficit species stay separated:}
\]

%dsp-level:0
\[
\qquad \frac{\gamma(8-\gamma)}{16}\ \ \text{per tick (informational)}
\]

%dsp-level:0
\[
\qquad \text{against}
\]

%dsp-level:0
\[
\qquad 1-N\ \ \text{(metric)};
\]

%dsp-level:0
\[
\neg\text{a numerical identification of the two};
\]

%dsp-level:0
\[
\qquad \text{they join only through}\ (7)\text{'s exponent}
\]

%dsp-level:0
\[
\qquad [\mathrm{BM.37}(2):\ \text{disjoint alphabets}];
\]

%dsp-level:0
\[
\textbf{the scope row:}
\]

%dsp-level:0
\[
\qquad (7)\ \text{is a chart-side reparametrization of a fixed law,}\ \neg\text{a dynamical coupling:}
\]

%dsp-level:0
\[
\qquad \neg\text{a feedback of the contraction into the metric is formed}\ [\mathrm{BM.104}(4)\text{'s freedom-census pattern}];
\]

%dsp-level:0
\[
\text{the tick-rate/field coupling is re-scoped by it: the redshifted contraction read is formed at the point of use, inheriting}\ \mathrm{CT2.6}\text{'s hypotheses}\ [\text{A §0.2a}];
\]

%dsp-level:0
\[
\qquad \text{its absolute-frequency datum remains calibration-species}\ [\mathrm{CT0.3}:\ \tau_c=1\ \text{tick}];
\]

%dsp-level:0
\[
(8)\quad\textbf{status:}
\]

%dsp-level:0
\[
\qquad \textbf{derived:}\ (1)\text{--}(7);
\]

%dsp-level:0
\[
\qquad \text{the duality is the pair}\ (\mu_\star,\ K)\ \text{with the three exact rows: the Haar invariance of}\ V,
\]

%dsp-level:0
\[
\ \text{the Poisson contraction of}\ V,
\]

%dsp-level:0
\[
\ \text{the unique common fixed point}\ \tfrac14I;
\]

%dsp-level:0
\[
\neg\text{a declaration},\qquad\neg\text{a P4 addition},\qquad\neg\text{a new dimensionful scale}\ [\mathrm{P1}],\qquad\neg\text{a measured value}\ [\mathrm{P5}],\qquad
\]

%dsp-level:0
\[
\textbf{the lapse gap is closed by}\ \mathrm{BM.132};
\]

%dsp-level:0
\[
\qquad \text{the row}\ \mathsf g,
\]

%dsp-level:0
\[
\ e_0,
\]

%dsp-level:0
\[
\ N\ \text{of}\ \mathrm{BM.130}(7)\ \text{is re-scoped:}\ N\ \text{derived}\ [\mathrm{BM.132}],
\]

%dsp-level:0
\[
\ \ \mathsf g\ \text{selected}\ [\mathrm{BM.133}],
\]

%dsp-level:0
\[
\ \ e_0\ \text{consolidated as calibration}\ [\mathrm{BM.143}(4)].
\]

\textbf{Req.} A3, A5, A6 (Part 1); D0.3, D0.4, D1.8 (Part 1); CT0.3, CT0.8 (Part 1); L0, L1, L2 (Part 1); CT2.6 (Part 2); P1, P5 (Parts 1–2); FT.107(2) (Part 3); Def-BM.10, BM.34(2), BM.37(2), BM.51(2), BM.51(3), BM.104(4) (this part); Schur's lemma on the $3\oplus1$ decomposition, the Haar-measure invariance and uniqueness on a compact group, formal power-series arithmetic over $\mathbb Q$ with the exact $\exp/\log$ pair, exact rational $\mathrm{SU}(2)$ points from Pythagorean quadruples, exact rank by Fraction elimination, Zariski-density substitution [math; machine: \texttt{checks/\allowbreak{}verify\_\allowbreak{}H\_\allowbreak{}bm\_\allowbreak{}closing\_\allowbreak{}tier.py} §A].

\textbf{Pf.}

%dsp-level:0
\[
(1):\ \mathrm{A6}\ \text{supplies}\ \mu_\star,\ \mathcal T,\ U\notin\mathfrak L;
\]

%dsp-level:0
\[
\qquad \mathrm{D0.3}\ \text{supplies}\ \tilde{\mathcal T}_\gamma;
\]

%dsp-level:0
\[
\qquad \mathrm{L0}\ \text{supplies the one-tilt-per-event coupling (declared scope, inherited at the point of use)};
\]

%dsp-level:0
\[
\qquad \text{the factorization is the composition of the two axiom-level moves}
\]

%dsp-level:0
\[
(2):\ \rho\ \text{faithful}\implies\rho,\ \tilde{\mathcal T}_\gamma(\rho)\ \text{simultaneously diagonalizable}\ [\mathrm{L1}(3)\ \mathrm{Pf}]:
\]

%dsp-level:0
\[
\quad K'
\]

%dsp-level:0
\[
= -\log\rho'
\]

%dsp-level:0
\[
= -\log\rho+\gamma\rho+\log Z_\gamma\cdot I\ \ \text{as an operator identity on the eigenbasis}\ [\text{machine: the}\ \mathbb Q[[t,\gamma]]\ \text{series identity}\ \log\bigl(\lambda e^{-\gamma\lambda}/Z_\gamma\bigr)+\gamma\lambda+\log Z_\gamma-\log\lambda\equiv0\ \text{to order}\ 12];
\]

%dsp-level:0
\[
\rho=\tfrac14I:\quad Z_\gamma=4\cdot\tfrac14e^{-\gamma/4}=e^{-\gamma/4};
\]

%dsp-level:0
\[
\qquad \lambda'=\tfrac14e^{-\gamma/4}/e^{-\gamma/4}=\tfrac14;
\]

%dsp-level:0
\[
\qquad \delta K=\gamma\cdot\tfrac14I+\log\bigl(e^{-\gamma/4}\bigr)I=0\ \ [\text{machine:}\ \exp(-\gamma/4)\exp(\gamma/4)=1\ \text{coefficient-wise}]
\]

%dsp-level:0
\[
(3):\ \mathrm{L1}(2);\qquad K(\mathcal T\rho)=UKU^\dagger\ [\mathrm{L1}(2)\ \mathrm{Pf}:\ \text{functional-calculus covariance}]\implies\operatorname{spec}K\ \text{invariant};\qquad\text{the stationarity of}\ \mu_\star\ \text{under conjugation is the two-sided invariance}\ [\text{math}];\qquad\text{the machine rows:}\ \operatorname{Tr}\bigl((U\rho U^\dagger)^n\bigr)=\operatorname{Tr}(\rho^n),\ n=1,2,3,\ \text{on}\ 40\ \text{exact rational points}\ U=\bigl(\begin{smallmatrix}a+bi&c+di\\-c+di&a-bi\end{smallmatrix}\bigr)/n,\ a^2{+}b^2{+}c^2{+}d^2{=}n^2,\ \text{and generic rational Hermitian}\ \rho
\]

%dsp-level:0
\[
(4):\ \mathbb C^2\otimes\mathbb C^2=\vee^2\oplus\wedge^2\ [\mathrm{A3}]:\quad U\otimes U\ \text{acts irreducibly on each block};
\]

%dsp-level:0
\[
\qquad \text{the twirl of any intertwiner-average projects onto the commutant}\
\]

%dsp-level:0
\[
= \ \operatorname{span}\{\Pi_{\mathcal R},\Pi_{\mathcal V}\}\ [\text{math: Schur}]:\quad\text{the block weights}\ \mathfrak T(\Pi_{\mathcal R}\rho),
\]

%dsp-level:0
\[
\ \mathfrak T(\Pi_{\mathcal V}\rho)\ \text{are the only invariants};
\]

%dsp-level:0
\[
\text{machine:}\ \Pi_{\mathcal R}+\Pi_{\mathcal V}=I,\ \Pi_{\mathcal R}\Pi_{\mathcal V}=0,\ \operatorname{rk}=(3,1);
\]

%dsp-level:0
\[
\qquad \mathrm{tw}\circ\mathrm{tw}=\mathrm{tw};
\]

%dsp-level:0
\[
\qquad \text{the block traces of}\ \mathrm{tw}(\rho)\ \text{equal those of}\ \rho;
\]

%dsp-level:0
\[
\qquad \text{the equivariance on all}\ 40\ \text{rational points};
\]

%dsp-level:0
\[
\qquad x/3=1-x\iff x=\tfrac34\implies\tfrac{x}3\Pi_{\mathcal R}+(1-x)\Pi_{\mathcal V}=\tfrac14I
\]

%dsp-level:0
\[
(5):\ V\bigl(\tilde K(U\tilde{\mathcal T}_\gamma(\rho)U^\dagger)\bigr)=V\bigl(\tilde K(\tilde{\mathcal T}_\gamma(\rho))\bigr)\ [(3)]=\bigl(1-\tfrac\gamma4\bigr)^2V(\tilde K(\rho))+O\ [\mathrm{L2}];
\]

%dsp-level:0
\[
\qquad \text{iteration:}\ \mathrm{BM.51}(3);
\]

%dsp-level:0
\[
\qquad \text{the record}\ V_n\ \text{carries}\ \neg\text{a draw label}\ [\mathrm{A6}:\ U\notin\mathfrak L]
\]

%dsp-level:0
\[
(6):\ \Pi_k^2
\]

%dsp-level:0
\[
= \frac1{(k!)^2}\sum_{\sigma,\tau}\hat P_{\sigma\tau}
\]

%dsp-level:0
\[
= \frac1{k!}\sum_{\tau'}\hat P_{\tau'}
\]

%dsp-level:0
\[
= \Pi_k\ [\mathrm{BM.128}(2)\text{'s product law transported};\ \text{machine: entrywise over}\ \mathbb Q];
\]

%dsp-level:0
\[
\qquad \operatorname{rk}\Pi_k=k+1\ [\text{machine: exact Fraction elimination for}\ k=1..6:\ \text{ranks}\ (2,3,4,5,6,7)];
\]

%dsp-level:0
\[
\qquad \operatorname{im}\Pi_k=\mathrm{Sym}^k\ [\mathrm{Def\text{-}H.38}\ \text{carrier}]
\]

%dsp-level:0
\[
(7):\ ds^2\big|_{dx^i=0}=-N^2d\tau^2\implies ds=N\,d\tau\ [\mathrm{BM.34}(2);\ \text{machine}];
\]

%dsp-level:0
\[
\qquad s=Nn\implies V_n=\bigl(1-\tfrac\gamma4\bigr)^{2n}V_0=\bigl(1-\tfrac\gamma4\bigr)^{2s/N}V_0\ [\text{substitution}];
\]

%dsp-level:0
\[
\qquad \text{the exact powers at}\ \gamma
\]

%dsp-level:0
\[
= 1:\ \bigl(\tfrac9{16}\bigr)^m,
\]

%dsp-level:0
\[
\ m
\]

%dsp-level:0
\[
= 1..4\ [\text{machine: strictly decreasing, exact integers}\ 9^m/16^m];
\]

%dsp-level:0
\[
\text{the species row:}\ \gamma(8-\gamma)/16\in\mathbb Q\ \text{against}\ 1-N\in\mathfrak q^0\ \text{chart data}:\ \neg\text{a derived equation relates them}\ [\mathrm{BM.37}(2)\ \text{pattern}];
\]

%dsp-level:0
\[
\qquad (8):\ \text{assembly}.\ \blacksquare
\]

\textbf{Phys.} The corpus owns exactly two state evolutions, and they are dual in the strict sense: an unrecorded random conjugation drawn from the invariant Haar measure, and a recorded modular update generated by $K=-\log\rho$ whose exact step is $\delta K=\gamma\rho+\log Z_\gamma\cdot I$. The measure side moves states inside the level sets of the Lyapunov function and preserves the modular spectrum; the dynamics side cuts strictly across the levels toward the vacuum. The Haar average of any pair state is its block-maximal-mixed projection, and the tilt's unique fixed point is the maximally mixed state itself, so the two sides share exactly one fixed point: the A5 vacuum. Because the draws are unrecorded, they drop out of every ledger read, and the contraction law the record obeys is purely Poisson-side — the measure is the record's exact gauge freedom. At the finite-group level the same twirl builds the spin carriers: averaging over $S_k$ is the symmetrizer, whose image is $\mathrm{Sym}^k$ of dimension $k+1$. Finally the metric joins: the lapse converts chart ticks to proper time along a resting worldline, so the same exact contraction read per proper time carries the lapse in the exponent, with the informational and metric deficit species kept apart and joined only by that exponent. No feedback is formed, no scale added, nothing measured.

\end{theorem}

\begin{theorem}[{\texorpdfstring{BM.132 — The lapse: the derived function of the envelope ratio and the curvature, its formed domain, and the envelope row}{BM.132 - The lapse: the derived function of the envelope ratio and the curvature, its formed domain, and the envelope row}}]\label{tag:BM-132}
\textbf{Thm.}

%dsp-level:0
\[
(1)\quad\textbf{the source and the potential (re-derived at the point of use):}
\]

%dsp-level:0
\[
\qquad T_{\mu\nu}=\mathrm{diag}(\rho,0,0,0)\ \ \text{at rest}
\]

%dsp-level:0
\[
\qquad [\mathrm{BM.9}(4);\ \mathrm{CT1.22}];
\]

%dsp-level:0
\[
G_{\mu\nu}+\Lambda g_{\mu\nu}=\kappa^2T_{\mu\nu}
\]

%dsp-level:0
\[
\qquad [\mathrm{CT1.23}];
\]

%dsp-level:0
\[
\qquad g_{00}=-(1+2\Phi),
\]

%dsp-level:0
\[
\qquad G^{(1)}_{00}=2\nabla^2\Phi,
\]

%dsp-level:0
\[
\qquad \nabla^2\Phi=\frac{\kappa^2}2\rho=4\pi G_N\rho
\]

%dsp-level:0
\[
\qquad [\mathrm{BM.43}(4);\ \text{machine re-run}];
\]

%dsp-level:0
\[
\Phi(r)=-\frac{G_Nm}{r}\ \ \text{outside the source support},
\]

%dsp-level:0
\[
\qquad G_Nm_{\mathrm{hid}}=\frac{\mathfrak hE_{\mathrm{int}}}4,
\]

%dsp-level:0
\[
\qquad r_s=2G_Nm_{\mathrm{hid}}=\frac{\mathfrak hE_{\mathrm{int}}}2
\]

%dsp-level:0
\[
\qquad [\mathrm{BM.43}(4);\ \mathrm{P3}:\ G_N\ \text{a named combination, quarantined}];
\]

%dsp-level:0
\[
(2)\quad\textbf{the ADM lapse at zero shift:}\qquad N^2=\Omega^2+\gamma^{ij}N_iN_j,\qquad w=0\implies N^i=0,\quad N=\Omega\qquad[\mathrm{BM.34}(1);\ \mathrm{CT1.12}];
\]

%dsp-level:0
\[
\text{the linearized dictionary}\ \Omega^2=1+2\Phi
\]

%dsp-level:0
\[
\qquad [\mathrm{CT1.12}:\ \Omega=e^{\alpha_\Omega a_{00}};\ \text{the conformal mode reads the potential};\ a_{00}=\frac1{2\alpha_\Omega}\ln(1+2\Phi)\ \text{on the formed locus}];
\]

%dsp-level:0
\[
\boxed{\ N(r,E)^2\ =\ 1+2\Phi\ =\ 1-\frac{2G_Nm}{r}\ =\ 1-\frac{r_s}{r}\ =\ 1-\frac{G(E)}{4\pi\,k}\ },
\]

%dsp-level:0
\[
\qquad k:=\frac{r}{\Delta x_{\mathrm B}(E)},
\]

%dsp-level:0
\[
\qquad G(E)=2\pi E_{\mathrm{int}}(E)^2\ln2
\]

%dsp-level:0
\[
\qquad [\mathrm{BM.88}(2)(\mathrm i),(\mathrm{iii});\ \mathrm{BM.43}(4);\ \text{machine: one identity in}\ \mathbb Q[E_{\mathrm{int}},\pi,\ln2,k]];
\]

%dsp-level:0
\[
\textbf{the lapse is a derived function of the envelope ratio and the derived curvature argument:}
\]

%dsp-level:0
\[
\qquad N=N(k,E)\in\mathfrak q^0
\]

%dsp-level:0
\[
\qquad [\mathrm{CT2.1};\ \text{machine: the exponent census}];
\]

%dsp-level:0
\[
(3)\quad\textbf{the formed domain:}
\]

%dsp-level:0
\[
\qquad N^2\ge0\iff k\ \ge\ \frac{G(E)}{4\pi}\ =\ \frac{E_{\mathrm{int}}^2\ln2}2\ =\ \frac{r_s}{\Delta x_{\mathrm B}}
\]

%dsp-level:0
\[
\qquad [\text{machine: the monomial}];
\]

%dsp-level:0
\[
\textbf{the envelope row:}
\]

%dsp-level:0
\[
\qquad \text{on every computed cell}\ E_{\mathrm{int}}\in\{2,6,12,24,60\}\ [\mathrm{BM.13}(2);\ \mathrm{Census\ FT.35}(2),\ \mathrm{FT.42}(5)\ \text{row-wise}]:
\]

%dsp-level:0
\[
\frac{G}{4\pi}=\Bigl(2,\ 18,\ 72,\ 288,\ 1800\Bigr)\ln2\ >\ 1
\]

%dsp-level:0
\[
\qquad \Bigl[\ln2>\frac12\iff4>e:\quad e<\frac{876809}{322560}<4;\ \text{machine: the exact rational enclosure of}\ e\Bigr];
\]

%dsp-level:0
\[
\boxed{\ 1+2\Phi\bigl(\Delta x_{\mathrm B}\bigr)=1-\frac{E_{\mathrm{int}}^2\ln2}2\ <\ 0\ \ \text{on every computed cell}\ }
\]

%dsp-level:0
\[
\qquad [\mathrm{BM.43}(4):\ \Phi(\Delta x_{\mathrm B})=-E_{\mathrm{int}}^2\ln2/4;\ \text{machine}];
\]

%dsp-level:0
\[
\textbf{the interior row (the uniform-ball read):}
\]

%dsp-level:0
\[
\qquad \Phi_{\mathrm{int}}(r)=-\frac{G_Nm}{2R}\Bigl(3-\frac{r^2}{R^2}\Bigr),
\]

%dsp-level:0
\[
\qquad \nabla^2\Phi_{\mathrm{int}}=4\pi G_N\rho,
\]

%dsp-level:0
\[
\qquad \Phi_{\mathrm{int}}(R)=\Phi(R)
\]

%dsp-level:0
\[
\qquad [\text{machine: the harmonicity and matching identities in}\ \mathbb Q[x,y,z,r,R]];
\]

%dsp-level:0
\[
1+2\Phi_{\mathrm{int}}(0)=1-\frac{3E_{\mathrm{int}}^2\ln2}4<0\ \ \text{on every computed cell}
\]

%dsp-level:0
\[
\qquad \Bigl[E_{\mathrm{int}}=2:\ 3\ln2>1\ \text{on}\ \ln2>\frac13:\ e^3<8;\ \text{machine}\Bigr];
\]

%dsp-level:0
\[
\boxed{\ \textbf{the self-lapse of a sealed composite is}\ \neg\textbf{formed at its own Born width: the envelope lies strictly inside the domain boundary on every computed cell}\ }
\]

%dsp-level:0
\[
(4)\quad\textbf{the monotonicity, the limits, and the gauge agreements:}
\]

%dsp-level:0
\[
\qquad \frac{\partial N^2}{\partial k}=\frac{G}{4\pi k^2}>0:\quad N\ \text{strictly increasing in}\ k\ \text{on the formed domain};
\]

%dsp-level:0
\[
\qquad N\ \text{strictly decreasing in}\ E_{\mathrm{int}}
\]

%dsp-level:0
\[
\qquad [\text{machine: sign census on exact rationals}];
\]

%dsp-level:0
\[
N=0\ \ \text{at the boundary}\ \ k=\frac{G}{4\pi};\qquad N\to1^-\ \ \text{as}\ \ k\to\infty;\qquad N\in(0,1]\ \ \text{on the formed locus}\qquad[\text{machine}];
\]

%dsp-level:0
\[
\textbf{the vacuum gauge agreement:}
\]

%dsp-level:0
\[
\qquad (a_{00},q_{ij},w_i)=0\implies N=1
\]

%dsp-level:0
\[
\qquad [\mathrm{BM.34}(2);\ \mathrm{BM.88}(2)(\mathrm{iii});\ \text{machine}]
\]

%dsp-level:0
\[
\qquad \text{and}
\]

%dsp-level:0
\[
\qquad \lim_{k\to\infty}N(k,E)=1;
\]

%dsp-level:0
\[
\textbf{T-evenness:}
\]

%dsp-level:0
\[
\qquad N\ \text{carries}\ w\ \text{only through}\ \gamma^{ij}N_iN_j=N^2-\Omega^2;\quad\text{at zero shift}\ N\ \text{is}\ w\text{-free}
\]

%dsp-level:0
\[
\qquad [\mathrm{BM.34}(3);\ \text{machine}];
\]

%dsp-level:0
\[
(5)\quad\textbf{the three substitutions into derived objects:}
\]

%dsp-level:0
\[
\qquad \textbf{(i) the dispersion:}\quad E=N(k,E)\sqrt{m^2+\gamma^{ij}p_ip_j}-N^ip_i
\]

%dsp-level:0
\[
\qquad [\mathrm{BM.34}(4);\ \text{the lapse entry now carries}\ (2)\text{'s function}];
\]

%dsp-level:0
\[
\textbf{(ii) the threshold:}
\]

%dsp-level:0
\[
\qquad \mathcal T(k,E)=\frac{\mathcal T_0}{N(k,E)}=\frac{E_{\mathrm{int}}\,\mathfrak h}{\ln2}\Bigl(1-\frac{E_{\mathrm{int}}^2\ln2}{2k}\Bigr)^{-1/2}
\]

%dsp-level:0
\[
\qquad [\mathrm{BM.115}(1);\ \mathrm{Def\text{-}BM.17}(3);\ \text{machine: the substitution}]:
\]

%dsp-level:0
\[
\mathcal T\to\infty\ \ \text{at the domain boundary};
\]

%dsp-level:0
\[
\qquad \mathcal T\to\mathcal T_0\ \ \text{as}\ \ k\to\infty;
\]

%dsp-level:0
\[
\qquad \textbf{a deeper lapse lowers every threshold at once, and the boundary locus raises every threshold without bound}\ [\mathrm{BM.115}(1)\ \text{re-cited}];
\]

%dsp-level:0
\[
\textbf{(iii) the coupling profile:}
\]

%dsp-level:0
\[
\qquad e(r,E)=e_0\sqrt{N(k,E)\cdot f(k)\cdot\mathsf g\bigl(G(E)\bigr)}
\]

%dsp-level:0
\[
\qquad [\mathrm{BM.88}(4);\ \mathsf g\ \text{carried symbolically here, selected in}\ \mathrm{BM.133}];
\]

%dsp-level:0
\[
(6)\quad\textbf{the bars:}
\]

%dsp-level:0
\[
\qquad r_s\ \text{enters as}\ \mathrm{P3}\text{'s named combination, quarantined:}\ \neg\text{a collapse reading,}\ \neg\text{a horizon identification,}\ \neg\text{black-hole physics}\ [\mathrm{P3};\ \mathrm{BM.15}(4);\ \mathrm{BM.115}(4);\ \mathrm{BM.130}(6)\ \text{discipline inherited}];
\]

%dsp-level:0
\[
\neg\text{a cosmological coordinate}\ [\mathrm{Def\text{-}BM.17}(6)(\mathrm i)];
\]

%dsp-level:0
\[
\qquad \neg\text{a measured curvature, radius or redshift}\ [\mathrm{P5}];
\]

%dsp-level:0
\[
\qquad \text{the weak-field rows inherit}\ \mathrm{CT2.6}\text{'s hypotheses at the point of use}\ [\text{A §0.2a};\ \mathrm{BM.43}(4)\ \text{scope}];
\]

%dsp-level:0
\[
\text{the full non-linear joint solution remains}\ \neg\text{constructed}\ [\mathrm{BM.43}(5)\ \text{re-cited,}\ \neg\text{extended}];
\]

%dsp-level:0
\[
(7)\quad\textbf{status:}
\]

%dsp-level:0
\[
\qquad \textbf{derived:}\ (1)\text{--}(5);
\]

%dsp-level:0
\[
\qquad \boxed{\ \textbf{the lapse gap is closed:}\ \ N(k,E)=\sqrt{1-\frac{G(E)}{4\pi k}}\ \ \text{on the formed domain}\ \ k\ge\frac{G(E)}{4\pi}\ };
\]

%dsp-level:0
\[
\qquad \text{the row}\ N\ \text{as solution data}\ [\mathrm{BM.88}(8);\ \mathrm{BM.93}(6);\ \mathrm{BM.100}(6);\ \mathrm{BM.107}(7);\ \mathrm{BM.118}(7);\ \mathrm{BM.130}(7)]\ \textbf{is discharged}.
\]

\textbf{Req.} A8 (Part 1); D2.1 (Part 1); CT1.12, CT1.22, CT1.23, CT2.1, CT2.6, CT2.9 (Part 2); P1, P3, P5 (Parts 1–2); Def-BM.4, Def-BM.10, Def-BM.17(3), Def-BM.17(6), BM.9(4), BM.13(2), BM.15(4), BM.34(1), BM.34(2), BM.34(3), BM.34(4), BM.43(4), BM.43(5), BM.88(2), BM.88(3), BM.88(4), BM.88(8), BM.115(1), BM.115(4), BM.130(6), BM.130(7) (this part); Census FT.35(2), FT.42(5) (row-wise); the exact rational enclosure of $e$ by the factorial series, the harmonicity of $1/r$ and the interior-ball potential in quotient rings, Laurent and monomial identities, sign censuses on exact rationals [math; machine: \texttt{checks/\allowbreak{}verify\_\allowbreak{}H\_\allowbreak{}bm\_\allowbreak{}closing\_\allowbreak{}tier.py} §A–§B].

\textbf{Pf.}

%dsp-level:0
\[
(1):\ \mathrm{BM.43}(4)\text{'s rows re-derived:}\ G^{(1)}_{00}
\]

%dsp-level:0
\[
= 2\nabla^2\Phi\ [\text{machine: the jet-ring Christoffel}\to\text{Ricci}\to\text{Einstein chain}];
\]

%dsp-level:0
\[
\qquad \nabla^2(1/r)=0\ [\text{machine: the quotient-ring identity}];
\]

%dsp-level:0
\[
\qquad G_Nm_{\mathrm{hid}}=\frac{\mathfrak h^2}{4\ln2}\cdot\frac{E_{\mathrm{int}}\ln2}{\mathfrak h}=\frac{\mathfrak hE_{\mathrm{int}}}4\ [\mathrm{P3};\ \mathrm{FT.44}(4);\ \text{machine}]
\]

%dsp-level:0
\[
(2):\ w=0\implies N_i=\Omega^2\varkappa\beta w_i=0\implies N^2=\Omega^2\ [\mathrm{BM.34}(1)];
\]

%dsp-level:0
\[
\qquad \text{the linearized dictionary:}\ g_{00}
\]

%dsp-level:0
\[
= -\Omega^2
\]

%dsp-level:0
\[
= -(1+2\Phi)\ [\mathrm{CT1.12};\ \mathrm{BM.34}(2)\text{'s line element at zero shift}];
\]

%dsp-level:0
\[
\qquad \text{the substitution:}\ 1-\frac{2G_Nm}r=1-\frac{r_s}r;
\]

%dsp-level:0
\[
\qquad \frac{r_s}{\Delta x_{\mathrm B}}
\]

%dsp-level:0
\[
= \frac{\mathfrak hE_{\mathrm{int}}/2}{\mathfrak h/(E_{\mathrm{int}}\ln2)}
\]

%dsp-level:0
\[
= \frac{E_{\mathrm{int}}^2\ln2}2
\]

%dsp-level:0
\[
= \frac{G}{4\pi}\ [\mathrm{BM.88}(3);\ \text{machine: one monomial}]
\]

%dsp-level:0
\[
(3):\ N^2\ge0\iff\frac{G}{4\pi k}\le1;
\]

%dsp-level:0
\[
\qquad E_{\mathrm{int}}\ge2\ [\mathrm{BM.13}(3)]\implies\frac{G}{4\pi}\ge2\ln2>1\ [\text{machine:}\ e<\tfrac{876809}{322560}<4\implies\ln4>1];
\]

%dsp-level:0
\[
\qquad \text{the interior row: the matching identities and}\ 3E_{\mathrm{int}}^2\ln2>4\ \text{at}\ E_{\mathrm{int}}=2\ [\text{machine:}\ e^3<8]
\]

%dsp-level:0
\[
(4):\ \text{the derivative sign and the endpoint values on exact rationals}\ [\text{machine}];
\]

%dsp-level:0
\[
\qquad (5):\ \text{substitution into}\ \mathrm{BM.34}(4),\ \mathrm{BM.115}(1),\ \mathrm{BM.88}(4)\ [\text{machine: the monomial rows}];
\]

%dsp-level:0
\[
\qquad (6),(7):\ \text{assembly and the rescope census}.\ \blacksquare
\]

\textbf{Phys.} The lapse was solution data; it is now a function. The composite's own dust source runs through the derived field equations to the weak-field potential, the ADM split reads the potential as the zero-shift lapse, and the curvature lock turns the answer into one dimensionless monomial: $N^2=1-G/(4\pi k)$ on the envelope ratio. The function comes with its own formed domain, and the domain has a sharp edge: on every computed cell the boundary $k=G/(4\pi)$ lies strictly beyond the Born width, so a sealed composite's self-lapse is not formed at its own envelope — the interior read only deepens the failure. The vacuum gauge is recovered exactly at infinity, the threshold ladder inherits the lapse and diverges at the boundary, and the running coupling's energy factor is now explicit. No black-hole reading is attached to the boundary, the Schwarzschild combination stays quarantined, and no measured number enters.

\end{theorem}

\begin{theorem}[{\texorpdfstring{BM.133 — The selected curvature composition: two derived criteria, the unique survivor, and the completed running coupling}{BM.133 - The selected curvature composition: two derived criteria, the unique survivor, and the completed running coupling}}]\label{tag:BM-133}
\textbf{Thm.}

%dsp-level:0
\[
(1)\quad\textbf{the three priced rows (re-cited, frozen):}\qquad\mathsf g\in\Bigl\{1,\quad1+G,\quad\frac1{1+G}\Bigr\}\qquad[\mathrm{BM.88}(4);\ \mathrm{BM.104}(2)]:
\]

%dsp-level:0
\[
\text{the back-reaction fixed-point rows:}
\]

%dsp-level:0
\[
\qquad (\mathrm i)\ \ \mathsf g\equiv1;
\]

%dsp-level:0
\[
\qquad (\mathrm{ii})\ \ \mathsf g_{\mathrm{fp}}=\frac{1+G}{1-aG},\ \ \text{admissible}\iff aG<1;
\]

%dsp-level:0
\[
\qquad (\mathrm{iii})\ \ aG\,\mathsf g^2+(1+G)\,\mathsf g-1=0,\ \ \mathsf g_{\mathrm{fp}}=\frac{\sqrt{(1+G)^2+4aG}-(1+G)}{2aG}
\]

%dsp-level:0
\[
\qquad [\mathrm{BM.104}(2);\ \text{machine re-run}];
\]

%dsp-level:0
\[
a=\tfrac12Nw\,\varepsilon_{\mathrm{EM}}Q^2\ge0,\qquad G>0\qquad[\mathrm{BM.104}(1),(3);\ \mathrm{BM.88}(3)];
\]

%dsp-level:0
\[
(2)\quad\textbf{criterion}\ (\mathrm S1)\ \textbf{— the Lyapunov contractivity of the back-reaction map:}
\]

%dsp-level:0
\[
\qquad F_{(\mathrm i)}(\mathsf g)=1;
\]

%dsp-level:0
\[
\qquad F_{(\mathrm{ii})}(\mathsf g)=1+G+aG\,\mathsf g;
\]

%dsp-level:0
\[
\qquad F_{(\mathrm{iii})}(\mathsf g)=\frac1{1+G+aG\,\mathsf g}
\]

%dsp-level:0
\[
\qquad [\mathrm{BM.104}(1)\text{'s loop read as a map}]:
\]

%dsp-level:0
\[
F'_{(\mathrm i)}=0;
\]

%dsp-level:0
\[
\qquad F'_{(\mathrm{ii})}=aG:\quad\text{contractive}\iff aG<1\ —\ \text{conditional, with the pole at}\ aG=1\ \text{where}\ \mathsf g_{\mathrm{fp}}\ \text{diverges}
\]

%dsp-level:0
\[
\qquad [\mathrm{BM.104}(2)(\mathrm{ii}),\ (3);\ \text{machine}];
\]

%dsp-level:0
\[
F'_{(\mathrm{iii})}(\mathsf g)=-\frac{aG}{(1+G+aG\,\mathsf g)^2}:
\]

%dsp-level:0
\[
\qquad \text{at the fixed point}\ 1+G+aG\,\mathsf g_{\mathrm{fp}}=\frac1{\mathsf g_{\mathrm{fp}}}\implies\boxed{\ \bigl|F'_{(\mathrm{iii})}\bigr|=aG\,\mathsf g_{\mathrm{fp}}^2=1-(1+G)\,\mathsf g_{\mathrm{fp}}\ }
\]

%dsp-level:0
\[
\textbf{the unconditional contraction:}
\]

%dsp-level:0
\[
\qquad 0<\bigl|F'_{(\mathrm{iii})}\bigr|<1\iff0<\mathsf g_{\mathrm{fp}}<\frac1{1+G}
\]

%dsp-level:0
\[
\qquad \text{and}
\]

%dsp-level:0
\[
\qquad (1+G)^2\bigl((1+G)^2+4aG\bigr)<\bigl((1+G)^2+2aG\bigr)^2\iff0<4a^2G^2
\]

%dsp-level:0
\[
\qquad [\text{machine: the squaring argument, exact in}\ \mathbb Q[G,a];\ \text{both sides positive on the formed locus}];
\]

%dsp-level:0
\[
\boxed{\ \textbf{row}\ (\mathrm{iii})\ \textbf{is unconditionally Lyapunov-contractive on the whole formed locus; row}\ (\mathrm{ii})\ \textbf{is not; row}\ (\mathrm i)\ \textbf{is contractive but}\ G\textbf{-inert}\ }
\]

%dsp-level:0
\[
\qquad [\mathrm{L1}\text{--}\mathrm{L2}\ \text{contraction discipline};\ \mathrm{BM.51}(2)\ \text{pattern}];
\]

%dsp-level:0
\[
(3)\quad\textbf{criterion}\ (\mathrm S2)\ \textbf{— the minimal lapse-domain loss:}
\]

%dsp-level:0
\[
\qquad \text{the corrected curvature}\ G_{\mathrm{corr}}=G(1+a\,\mathsf g)\ [\mathrm{BM.104}(1)]\ \text{moves the lapse boundary:}
\]

%dsp-level:0
\[
\qquad N_{\mathrm{corr}}(k)^2=1-\frac{G(1+a\,\mathsf g)}{4\pi k}\ \ \text{formed}\iff k\ge\frac{G(1+a\,\mathsf g)}{4\pi}
\]

%dsp-level:0
\[
\qquad [\mathrm{BM.132}(2),(3)];
\]

%dsp-level:0
\[
\text{every positive composition shrinks the formed domain}\ (a\,\mathsf g>0);
\]

%dsp-level:0
\[
\qquad \textbf{the loss is exactly ordered:}
\]

%dsp-level:0
\[
\qquad \boxed{\ a\,\mathsf g_{\mathrm{fp}}^{(\mathrm{iii})}\ <\ \frac{a}{1+G}\ <\ a\ =\ a\,\mathsf g^{(\mathrm i)}\ <\ a\,(1+G)\ =\ a\,\mathsf g^{(\mathrm{ii})}\ }
\]

%dsp-level:0
\[
\qquad [\text{machine: the strict chain on}\ a,G>0];
\]

%dsp-level:0
\[
\textbf{row}\ (\mathrm{ii})\ \textbf{at the pole loses the whole domain:}
\]

%dsp-level:0
\[
\qquad aG\to1:\quad\mathsf g_{\mathrm{fp}}\to\infty,\quad G_{\mathrm{corr}}\to\infty,\quad\{k:\ N_{\mathrm{corr}}\ \text{formed}\}\to\varnothing
\]

%dsp-level:0
\[
\qquad [\mathrm{BM.104}(2)(\mathrm{ii}),\ (3);\ \text{machine}];
\]

%dsp-level:0
\[
\boxed{\ \textbf{row}\ (\mathrm{iii})\ \textbf{loses the least formed lapse domain of the three, uniformly in}\ (a,G)\ }
\]

%dsp-level:0
\[
(4)\quad\textbf{the unique survivor:}
\]

%dsp-level:0
\[
\qquad (\mathrm S1)\wedge(\mathrm S2):\quad\text{row}\ (\mathrm{ii})\ \text{fails both (conditionality, maximal loss with a total collapse at the pole)};
\]

%dsp-level:0
\[
\qquad \text{row}\ (\mathrm i)\ \text{passes}\ (\mathrm S1)\ \text{and fails}\ (\mathrm S2)\ \text{and is}\ G\text{-inert}\ (2);
\]

%dsp-level:0
\[
\boxed{\ \mathsf g(G)\ =\ \frac1{1+G}\ }
\]

%dsp-level:0
\[
\qquad \textbf{selected by two derived criteria;}
\]

%dsp-level:0
\[
\qquad \mathsf g\in\mathfrak q^0,
\]

%dsp-level:0
\[
\qquad 0<\mathsf g\le1,
\]

%dsp-level:0
\[
\qquad \mathsf g\ \text{strictly decreasing in}\ G,
\]

%dsp-level:0
\[
\qquad \mathsf g(0)=1\ \ \text{the common vacuum-normalization row of all three}\ [\mathrm{BM.104}(2);\ \text{machine}];
\]

%dsp-level:0
\[
\textbf{the exact values at the computed cells:}
\]

%dsp-level:0
\[
\qquad \mathsf g=\frac1{1+2\pi E_{\mathrm{int}}^2\ln2},
\]

%dsp-level:0
\[
\qquad E_{\mathrm{int}}\in\{2,6,12,24,60\}:\quad\mathsf g\in\Bigl\{\frac1{1+8\pi\ln2},\ \frac1{1+72\pi\ln2},\ \frac1{1+288\pi\ln2},\ \frac1{1+1152\pi\ln2},\ \frac1{1+7200\pi\ln2}\Bigr\}
\]

%dsp-level:0
\[
\qquad [\text{machine: the monomial substitutions, carried exactly}];
\]

%dsp-level:0
\[
\textbf{the observable consequence (re-cited):}
\]

%dsp-level:0
\[
\qquad \mathrm{BM.101}(5)\text{'s crossing count on the selected row is two, one or none}\ —\ \text{the composition is observable in the energy-scale graph's topology, and the selection is now made}\ [\mathrm{BM.104}(6)\ \text{re-scoped}];
\]

%dsp-level:0
\[
(5)\quad\textbf{the}\ \mathrm{BM.104}\ \textbf{bars are respected:}
\]

%dsp-level:0
\[
\qquad \textbf{the loop stays excluded:}\ G_{\mathrm{corr}}\ \neg\text{re-enters}\ G;
\]

%dsp-level:0
\[
\ \text{the lock's argument remains the informational mass}\ [\mathrm{BM.104}(4)\ \text{frozen};\ \mathrm{BM.88}(3);\ \text{machine: the alphabet census}\ \{E_{\mathrm{int}},\ln2,\pi\}\cap\{N,w,\varepsilon_{\mathrm{EM}},Q\}=\varnothing\ \text{re-run}];
\]

%dsp-level:0
\[
\textbf{the four grounds stand and none is used:}
\]

%dsp-level:0
\[
\qquad (\mathrm i)\ \neg\text{a count}\to\text{strength map enters}\ [\mathrm{BM.36}(3);\ \mathrm{BM.61}(8)(\mathrm{O1})];
\]

%dsp-level:0
\[
\qquad (\mathrm{ii})\ \neg\text{a defect-conditioned weighting}\ [\mathrm{BM.90}(3)];
\]

%dsp-level:0
\[
\qquad (\mathrm{iii})\ \text{consistency is}\ \neg\text{used as sufficient}\ —\ \text{the selection runs on}\ (\mathrm S1),
\]

%dsp-level:0
\[
(\mathrm S2),
\]

%dsp-level:0
\[
\ \text{which are stability and metric-formedness criteria, not the fixed-point equation};
\]

%dsp-level:0
\[
\qquad (\mathrm{iv})\ \neg\text{a P4 datum: the survivor is unique, so no choice remains}\ [\mathrm{BM.104}(5)(\mathrm{iv})\ \text{discharged by uniqueness}];
\]

%dsp-level:0
\[
\mathrm{BM.104}(5)\text{'s}\ \neg\exists\text{-row is re-scoped to the four grounds it cites,}\ \neg\text{retracted};
\]

%dsp-level:0
\[
\qquad \mathrm{BM.104}(6)\text{'s discrimination row is completed by the selection};
\]

%dsp-level:0
\[
(6)\quad\textbf{the completed running coupling:}
\]

%dsp-level:0
\[
\qquad \boxed{\ e(r,E)\ =\ e_0\,\sqrt{\,N(k,E)\cdot f(k)\cdot\mathsf g\bigl(G(E)\bigr)\,}\ =\ e_0\left[\frac{\bigl(1-\frac{G}{4\pi k}\bigr)^{1/2}\bigl(1-\frac1k\bigr)}{1+G}\right]^{1/2}\ }
\]

%dsp-level:0
\[
\qquad [\mathrm{BM.88}(4);\ \mathrm{BM.132}(2);\ \text{machine: the assembled monomial}];
\]

%dsp-level:0
\[
\textbf{all three factors are derived:}
\]

%dsp-level:0
\[
\qquad N(k,E)\ [\mathrm{BM.132}];
\]

%dsp-level:0
\[
\qquad f(k)=1-\frac1k\ [\mathrm{BM.88}(2)(\mathrm{ii});\ \mathrm{BM.47}(6)];
\]

%dsp-level:0
\[
\qquad \mathsf g(G)=\frac1{1+G}\ [(4)];
\]

%dsp-level:0
\[
\qquad \textbf{only}\ e_0\ \textbf{remains calibration}\ [\mathrm{EM.25}(7);\ \mathrm{BM.143}(4)];
\]

%dsp-level:0
\[
\textbf{the formed locus:}
\]

%dsp-level:0
\[
\qquad k\ \ge\ \frac{G}{4\pi}\ \ \wedge\ \ k>1:
\]

%dsp-level:0
\[
\qquad \text{both factors strictly positive on its interior};
\]

%dsp-level:0
\[
\qquad e\to0\ \ \text{at the lapse boundary}\ (N\to0);
\]

%dsp-level:0
\[
\qquad e\to\frac{e_0}{\sqrt{1+G}}\ \ \text{as}\ \ k\to\infty
\]

%dsp-level:0
\[
\qquad [\text{machine: the exact limit rows}];
\]

%dsp-level:0
\[
\textbf{the monotonicity:}
\]

%dsp-level:0
\[
\qquad k\mapsto N(k)f(k)\ \text{strictly increasing on the formed locus}
\]

%dsp-level:0
\[
\qquad [\text{both factors positive and strictly increasing}\ [\mathrm{BM.88}(2)(\mathrm{ii});\ (4)];\ \text{machine: the derivative sign on exact rationals}];
\]

%dsp-level:0
\[
\mathrm{BM.88}(5)\text{'s envelope zero}\ f(1)
\]

%dsp-level:0
\[
= 0\ \textbf{is re-scoped:}\ \text{it is a vacuum-gauge read}\ (N{=}1);
\]

%dsp-level:0
\[
\ \text{on the derived self-lapse the envelope locus itself is}\ \neg\text{formed}\ [\mathrm{BM.132}(3)];
\]

%dsp-level:0
\[
\qquad \text{the formed coupling is strictly positive on its whole formed locus and vanishes only at the boundary};
\]

%dsp-level:0
\[
(7)\quad\textbf{status:}
\]

%dsp-level:0
\[
\qquad \textbf{derived:}\ (1)\text{--}(6);
\]

%dsp-level:0
\[
\qquad \textbf{the curvature composition is closed:}\ \mathsf g(G)=\frac1{1+G};
\]

%dsp-level:0
\[
\qquad \text{the row}\ \mathsf g\ \text{of}\ \mathrm{BM.130}(7)\ \textbf{is discharged};
\]

%dsp-level:0
\[
\neg\text{a declaration}\ (\text{the survivor is unique}),
\]

%dsp-level:0
\[
\qquad \neg\text{a P4 addition},
\]

%dsp-level:0
\[
\qquad \neg\text{a new dimensionful scale}\ [\mathrm{P1}],
\]

%dsp-level:0
\[
\qquad \neg\text{a measured coupling}\ [\mathrm{P5}].
\]

\textbf{Req.} A4 (Part 1); D0.3, D0.4 (Part 1); CT2.1 (Part 2); P1, P4, P5 (Parts 1–2); L1, L2 (Part 1); EM.25(7) (Part 2); Def-BM.12, BM.36(3), BM.47(6), BM.51(2), BM.61(8), BM.88(2), BM.88(3), BM.88(4), BM.88(5), BM.90(3), BM.101(5), BM.104, BM.130(7), BM.132(2), BM.132(3) (this part); the squaring argument in $\mathbb Q[G,a]$, the strict-order chain on positive monomials, exact rational limit and derivative sign rows, the radical identity $(1+G)\mathsf g_{\mathrm{fp}}<1$ in $\mathbb Q[G,a,s]/(s^2-(1+G)^2-4aG)$ [math; machine: \texttt{checks/\allowbreak{}verify\_\allowbreak{}H\_\allowbreak{}bm\_\allowbreak{}closing\_\allowbreak{}tier.py} §B].

\textbf{Pf.}

%dsp-level:0
\[
(1):\ \mathrm{BM.104}(2)\text{'s rows re-verified by substitution}\ [\text{machine}];
\]

%dsp-level:0
\[
\qquad (2):\ F'_{(\mathrm{iii})}\ \text{at the fixed point:}\ 1+G+aG\mathsf g_{\mathrm{fp}}
\]

%dsp-level:0
\[
= 1/\mathsf g_{\mathrm{fp}}\ [\text{the fixed-point equation rearranged}]
\]

%dsp-level:0
\[
\implies |F'|
\]

%dsp-level:0
\[
= aG\mathsf g_{\mathrm{fp}}^2;
\]

%dsp-level:0
\[
\qquad aG\mathsf g_{\mathrm{fp}}^2=1-(1+G)\mathsf g_{\mathrm{fp}}\ [\text{machine: the equation}\ aG\mathsf g^2+(1+G)\mathsf g-1=0\ \text{at}\ \mathsf g=\mathsf g_{\mathrm{fp}}];
\]

%dsp-level:0
\[
\mathsf g_{\mathrm{fp}}<\tfrac1{1+G}:
\]

%dsp-level:0
\[
\quad (1+G)\sqrt{(1+G)^2+4aG}<(1+G)^2+2aG
\]

%dsp-level:0
\[
\iff (1+G)^2\bigl((1+G)^2+4aG\bigr)<\bigl((1+G)^2+2aG\bigr)^2
\]

%dsp-level:0
\[
\iff 0<4a^2G^2\ [\text{machine: expansion in}\ \mathbb Q[G,a]]
\]

%dsp-level:0
\[
(3):\ \text{the domain-boundary monomial}\ \frac{G(1+a\mathsf g)}{4\pi}\ \text{is strictly increasing in}\ \mathsf g>0;
\]

%dsp-level:0
\[
\qquad \text{the chain}\ a\mathsf g_{\mathrm{fp}}<\frac a{1+G}<a<a(1+G)\ \text{on}\ (2)\text{'s inequality and}\ 1+G>1\ [\text{machine}];
\]

%dsp-level:0
\[
\qquad \text{the pole row:}\ (1-aG)\mathsf g_{\mathrm{fp}}=1+G\ \text{and}\ aG=1\implies0=1+G\ne0\ [\text{machine}]
\]

%dsp-level:0
\[
(4):\ \text{the census of}\ (2),(3)\ \text{over the three rows};
\]

%dsp-level:0
\[
\qquad \text{the cell values:}\ G=2\pi E_{\mathrm{int}}^2\ln2\ \text{substituted}\ [\text{machine: monomials}];
\]

%dsp-level:0
\[
\qquad (5):\ \text{the alphabet census re-run and the four grounds re-cited at their loci}\ [\text{machine}];
\]

%dsp-level:0
\[
(6):\ e^2=e_0^2\,N\,f\,\mathsf g\ \text{with}\ (4)\text{'s}\ \mathsf g\ \text{and}\ \mathrm{BM.132}(2)\text{'s}\ N\ [\text{machine: the assembled product}];
\]

%dsp-level:0
\[
\qquad \frac{d}{dk}\bigl[Nf\bigr]>0\ \text{on}\ k>\max\bigl(1,\frac G{4\pi}\bigr):\ \text{both factors positive increasing}\ [\text{machine: rational grid and product rule sign census}];
\]

%dsp-level:0
\[
\qquad (7):\ \text{assembly}.\ \blacksquare
\]

\textbf{Phys.} Three curvature compositions were priced and none was selected, because consistency alone is satisfied three times over. Two derived criteria break the tie. The first is the corpus's own contraction discipline: the back-reaction map of row (iii) is a strict contraction at its fixed point everywhere on the formed locus, row (ii) is contractive only below its pole and collapses there, row (i) is inert. The second is the metric's: every composition inflates the corrected horizon locus and shrinks the formed domain of the lapse, and row (iii) shrinks it least, uniformly. The unique survivor is $\mathsf g=1/(1+G)$ — curvature-diluted, vacuum-normalized, strictly decreasing. The back-reaction loop stays excluded: the lock's argument remains the informational mass, and the selection consumes stability and formedness, never a measured value or a free parameter. With the lapse derived and the composition selected, the running coupling is complete as a function of $(k,E)$ — every factor derived, every zero located, only $e_0$ left to calibration.

\end{theorem}

\begin{definition}[{\texorpdfstring{Def-BM.20 — The momentum-space $\Sigma$-resolvent (the construction BM.6(6) specifies)}{Def-BM.20 - The momentum-space -resolvent (the construction BM.6(6) specifies)}}]\label{tag:Def-BM-20}
\textbf{Def.} A chart-side interface definition on $\mathrm{BM.34}$'s ADM split and $\mathrm{CT2.6}$'s charts; the non-derived clauses are stated exactly and priced inside Part 4; no new axiom, no new datum, no new dimensionful scale. Logged as §II.4 addendum A-BM.12 before use.

%dsp-level:0
\[
(1)\quad\textbf{the boundary-channel states:}
\]

%dsp-level:0
\[
\qquad |p\rangle\ :=\ \text{the chart plane-wave state of covariant momentum label}\ p_\mu\ \text{on}\ \mathrm{CT2.6}\text{'s arena}
\]

%dsp-level:0
\[
\qquad [\text{math: plane waves on the chart}];
\]

%dsp-level:0
\[
\textbf{non-derived clause (the embedding):}
\]

%dsp-level:0
\[
\qquad \text{the identification of the boundary-channel register's momentum label with the chart's}\ p_\mu\ \text{— an}\ \iota_\Sigma\text{-type embedding at the momentum level}\ [\mathrm{BM.6}(6)];
\]

%dsp-level:0
\[
\textbf{scope:}
\]

%dsp-level:0
\[
\qquad \mathrm{CT2.6}\text{'s continuum-limit hypotheses are inherited at every point of use}\ [\text{A §0.2a}];
\]

%dsp-level:0
\[
\qquad \text{the states are chart objects,}\ \neg\text{ledger objects};
\]

%dsp-level:0
\[
\qquad \neg\text{a canonical pair is imported into the finite register}\ [\mathrm{BM.10}(1)\ \text{stands}];
\]

%dsp-level:0
\[
(2)\quad\textbf{the resolvent:}\qquad\boxed{\ R(z;\mathbf p)\ :=\ \frac1{\bigl(p_0-N^ip_i\bigr)^2-N^2\bigl(m^2+\gamma^{ij}p_ip_j\bigr)}\bigg|_{p_0=-z}\ }
\]

%dsp-level:0
\[
\textbf{constitutive clause:}
\]

%dsp-level:0
\[
\qquad \text{the denominator is the derived mass-shell operator}\ —\ \text{its zero locus is the dispersion relation}\ [\mathrm{BM.34}(4)];
\]

%dsp-level:0
\[
\qquad m=m_{\mathrm{hid}}(\mathcal R)\ \text{the composite's dispersion mass}\ [\mathrm{Def\text{-}BM.10};\ \mathrm{BM.34}(7)];
\]

%dsp-level:0
\[
(3)\quad\textbf{the Fourier convention (declared):}
\]

%dsp-level:0
\[
\qquad \text{the chart measure carries}\ \frac{d^3p}{(2\pi)^3}\ \text{per spatial dimension};
\]

%dsp-level:0
\[
\qquad \textbf{cost:}\ \neg\text{a scale,}\ \neg\text{a P4 member}\ —\ \text{any other normalization rescales every amplitude read off it by a declared constant, so the convention is a naming, not a datum}\ [\mathrm{D2.9}:\ \hbar=1];
\]

%dsp-level:0
\[
(4)\quad\textbf{the vertex normalization (constitutive minimality):}
\]

%dsp-level:0
\[
\qquad \mathcal V=1:
\]

%dsp-level:0
\[
\qquad \textbf{the reduction clause:}\ \ \text{at}\ e
\]

%dsp-level:0
\[
= 1,
\]

%dsp-level:0
\[
\ q
\]

%dsp-level:0
\[
= 1\ \text{every amplitude of}\ \mathrm{BM.142}\ \text{reduces to the register's channel weight}\ [\mathrm{BM.122}(3);\ \mathrm{BM.123}(1)]:
\]

%dsp-level:0
\[
\textbf{price:}
\]

%dsp-level:0
\[
\qquad \text{the unique normalization with this reduction}\ —\ \text{any other choice is an}\ e_0\text{-species rescaling}\ [\mathrm{EM.25}(7);\ \text{calibration,}\ \neg\text{structure}];
\]

%dsp-level:0
\[
(5)\quad\textbf{the bars:}
\]

%dsp-level:0
\[
\qquad \neg\text{an S-matrix unitarity claim}\ (\text{the formed reads are tree-level symbolic});
\]

%dsp-level:0
\[
\qquad \text{amplitude}\ne\text{mass}\ [\mathrm{FT.79}(5)];
\]

%dsp-level:0
\[
\qquad \neg\text{a numerical cross-section,}\ \neg\text{a measured comparison}\ [\mathrm{P5}];
\]

%dsp-level:0
\[
\neg\text{an}\ i\varepsilon\text{-prescription import:}\ R\ \text{is read meromorphically, sign-free off the poles}\ [\mathrm{FT.26}(1)\ \text{pattern}];
\]

%dsp-level:0
\[
\qquad \text{the}\ \aleph\text{-domain at the junction couplings is}\ \mathrm{E2}\ [\text{A-BM.10},\ \text{issued in}\ \mathrm{BM.139}];
\]

%dsp-level:0
\[
(6)\quad\textbf{status:}
\]

%dsp-level:0
\[
\qquad \textbf{constructed}\ —\ \mathrm{BM.6}(6)\text{'s specification is discharged;}
\]

%dsp-level:0
\[
\qquad \text{the resolvent specification is discharged by}\ \mathrm{BM.134};
\]

%dsp-level:0
\[
\qquad \text{the two derived clauses}\ (2),
\]

%dsp-level:0
\[
(4)\ \text{carry no new freedom beyond}\ (1),
\]

%dsp-level:0
\[
(3)\text{'s declared identifications}.
\]

\textbf{Req.} D2.5, D2.9 (Parts 1–2); CT2.6 (Part 2); P4, P5 (Parts 1–2); FT.26(1), FT.79(5) (Part 3); Def-BM.10, Def-BM.17, BM.6(6), BM.10(1), BM.34(4), BM.34(7), BM.122(3), BM.123(1), EM.25(7) (this part or Part 2); A-BM.12 (§II.4, logged before use); the plane-wave basis on a chart [math].

\textbf{Pf.} Specification content: each clause states its requirement and cites its anchors; the non-derived clauses (1) and (3) are declared at the point of use and logged in §II.4; the constitutive clauses (2) and (4) are constructions on derived objects, each priced in its own clause; anchor existence machine-checked [\texttt{checks/\allowbreak{}verify\_\allowbreak{}H\_\allowbreak{}bm\_\allowbreak{}closing\_\allowbreak{}tier.py} §C: the document pins]. $\blacksquare$

\textbf{Phys.} The momentum-space resolvent the corpus named as an open interface is now a definition with exactly two declared ingredients — the boundary-channel embedding and the Fourier normalization — and everything else forced: the denominator is the derived mass-shell operator, so the pole locus is the dispersion relation, and the vertex normalization is pinned by its reduction to the register weight at unit coupling. The bars travel with it: no unitarity claim, no numerical cross-section, no measured comparison, and no Feynman prescription imported from outside.

\end{definition}

\begin{theorem}[{\texorpdfstring{BM.134 — The resolvent's derived content: the form uniqueness, the pole reads, the on-shell measure, and the two-body phase space}{BM.134 - The resolvent's derived content: the form uniqueness, the pole reads, the on-shell measure, and the two-body phase space}}]\label{tag:BM-134}
\textbf{Thm.}

%dsp-level:0
\[
(1)\quad\textbf{the pole locus is the derived dispersion:}
\]

%dsp-level:0
\[
\qquad R(z;\mathbf p)^{-1}=0\iff\bigl(p_0-N^ip_i\bigr)^2=N^2\bigl(m^2+\gamma^{ij}p_ip_j\bigr)\iff g^{\mu\nu}p_\mu p_\nu=-m^2
\]

%dsp-level:0
\[
\qquad [\mathrm{BM.34}(4);\ \text{machine: the master polynomial identity re-derived, cleared by}\ (\det\gamma)^2,\ \text{with three exact rational ADM witnesses}];
\]

%dsp-level:0
\[
\textbf{the branches:}
\]

%dsp-level:0
\[
\qquad p_0=N^ip_i\pm N\sqrt{m^2+p^2};
\]

%dsp-level:0
\[
\qquad E=-p_0=N\sqrt{m^2+p^2}-N^ip_i\ \ \text{on the positive branch}
\]

%dsp-level:0
\[
\qquad [\mathrm{BM.34}(4);\ \mathrm{BM.119}(1)];
\]

%dsp-level:0
\[
\textbf{the rest-pole read:}\qquad\mathbf p=0,\ w=0:\quad p_0=\pm Nm;\qquad E=Nm\qquad[\mathrm{BM.34}(7);\ \text{machine}];
\]

%dsp-level:0
\[
\textbf{the pole displacement:}
\]

%dsp-level:0
\[
\qquad \text{in the vacuum gauge}\ (N{=}1,\ w{=}0):\quad R=\frac1{p_0^2-m^2-\mathbf p^2}:
\]

%dsp-level:0
\[
\qquad \boxed{\ m^2\ \ \text{is exactly the pole displacement in}\ p_0^2\ \text{against the massless pole}\ }
\]

%dsp-level:0
\[
\qquad [\mathrm{D2.5}\ \text{vocabulary};\ \text{machine: the cleared identity}];
\]

%dsp-level:0
\[
(2)\quad\textbf{the form uniqueness:}
\]

%dsp-level:0
\[
\qquad \text{among rational forms}\ R'=\bigl(\alpha\,(p_0-N^ip_i)^2+\beta\,N^2(m^2+p^2)+\delta\bigr)^{-1}\ \text{over}\ \mathbb Q(\alpha,\beta,\delta):
\]

%dsp-level:0
\[
\text{the five conditions}
\]

%dsp-level:0
\[
\qquad (\mathrm i)\ \ \text{the pole locus is}\ \mathrm{BM.34}(4)\text{'s mass shell};
\]

%dsp-level:0
\[
\qquad (\mathrm{ii})\ \ \text{the vacuum-gauge reduction}\ R'\to\bigl(p_0^2-m^2-p^2\bigr)^{-1};
\]

%dsp-level:0
\[
(\mathrm{iii})\ \ \text{T-evenness under}\ (N^i,p_i)\mapsto(-N^i,-p_i)\ [\mathrm{BM.34}(4)];
\]

%dsp-level:0
\[
\qquad (\mathrm{iv})\ \ \text{hermitian sign-freeness off the poles}\ [\mathrm{FT.26}(1)\ \text{pattern}];
\]

%dsp-level:0
\[
\qquad (\mathrm v)\ \ \text{simple poles (minimality, no squared denominators)};
\]

%dsp-level:0
\[
\boxed{\ \text{force}\ \ (\alpha,\beta,\delta)=(1,-1,0)\ \ \text{exactly}\ }
\]

%dsp-level:0
\[
\qquad [\text{machine: the exact linear solve over}\ \mathbb Q\ \text{on rational ADM witnesses;}\ (\mathrm i)\implies\alpha=-\beta,\ \delta=0;\ (\mathrm{ii})\implies\alpha=1];
\]

%dsp-level:0
\[
\textbf{the declared content is exactly}\ \mathrm{Def\text{-}BM.20}(1),(3)\text{'s two identifications:}
\]

%dsp-level:0
\[
\qquad \text{every other row is forced};
\]

%dsp-level:0
\[
(3)\quad\textbf{the on-shell measure (the phase space from the resolvent and the chart measure):}
\]

%dsp-level:0
\[
\qquad \delta\bigl((p_0-N^ip_i)^2-N^2E_{\mathbf p}^2\bigr)=\frac{\delta\bigl(p_0-N^ip_i+NE_{\mathbf p}\bigr)+\delta\bigl(p_0-N^ip_i-NE_{\mathbf p}\bigr)}{2N\,E_{\mathbf p}},
\]

%dsp-level:0
\[
\qquad E_{\mathbf p}:=\sqrt{m^2+\gamma^{ij}p_ip_j}
\]

%dsp-level:0
\[
\qquad [\text{math: the delta composition rule};\ \text{machine: the Jacobian coefficient at five rational witnesses}];
\]

%dsp-level:0
\[
\boxed{\ d\mu_{\mathrm{shell}}\ =\ \frac{d^3p}{2\,N\sqrt{m^2+\gamma^{ij}p_ip_j}}\ =\ \frac{d^3p}{2E}\ }
\]

%dsp-level:0
\[
\qquad [\text{the positive branch};\ \mathrm{Def\text{-}BM.20}(3)\text{'s}\ (2\pi)^3\ \text{carried with it}];
\]

%dsp-level:0
\[
\textbf{the vacuum-gauge reduction:}
\]

%dsp-level:0
\[
\qquad N=1,\ \gamma=\delta:\quad d\mu_{\mathrm{shell}}=\frac{d^3p}{2\sqrt{m^2+p^2}}
\]

%dsp-level:0
\[
\qquad [\text{machine}];
\]

%dsp-level:0
\[
\qquad [d\mu_{\mathrm{shell}}]=\mathfrak q^{+1}
\]

%dsp-level:0
\[
\qquad [\mathrm{CT2.1};\ \text{machine: the exponent census}];
\]

%dsp-level:0
\[
(4)\quad\textbf{the two-body phase space and the fluxes (CM frame, zero shift):}
\]

%dsp-level:0
\[
\qquad z=E_1+E_2;
\]

%dsp-level:0
\[
\qquad E_i=N\sqrt{m_i^2+p^2};
\]

%dsp-level:0
\[
\qquad p_\ast=\frac1{2N}\sqrt{z^2-4N^2m^2}\ \ \text{at}\ m_1=m_2=m
\]

%dsp-level:0
\[
\qquad [\text{machine: the exact solve}];
\]

%dsp-level:0
\[
\boxed{\ \Phi_2(z;m,N)\ =\ \frac{p_\ast}{4\pi N^2\,z}\ =\ \frac{\sqrt{z^2-4N^2m^2}}{8\pi N^3\,z}\ }
\]

%dsp-level:0
\[
\qquad [\text{math: the delta reduction of}\ \int d\mu_1d\mu_2\ \text{with}\ (2\pi)^4\delta^4\ \text{and}\ \mathrm{Def\text{-}BM.20}(3)\text{'s convention};\ \text{machine: the coefficient algebra at five rational witnesses}];
\]

%dsp-level:0
\[
\textbf{the massless final states:}
\]

%dsp-level:0
\[
\qquad \Phi_2^{(\gamma\gamma)}=\frac1{8\pi N^3}
\]

%dsp-level:0
\[
\qquad [\text{machine}];
\]

%dsp-level:0
\[
\qquad \textbf{the fluxes:}
\]

%dsp-level:0
\[
\qquad F_{\gamma\gamma}=\frac{2z^2}{N^2}\ \ (\text{two massless beams}),
\]

%dsp-level:0
\[
\qquad F_{\mathrm{pair}}=\frac{2z\sqrt{z^2-4N^2m^2}}{N^2}\ \ (\text{two massive beams})
\]

%dsp-level:0
\[
[\text{machine:}\ p_1\!\cdot\!p_2=-\frac{z^2}{2N^2}\ \text{massless and}\ m^2-\frac{z^2}{2N^2}\ \text{massive at zero shift};\ F=4\sqrt{(p_1\!\cdot\!p_2)^2-m_1^2m_2^2}];
\]

%dsp-level:0
\[
\textbf{the formed domain:}
\]

%dsp-level:0
\[
\qquad \Phi_2\ \text{formed}\iff z\ge2Nm;
\]

%dsp-level:0
\[
\qquad \Phi_2=0\ \ \text{exactly at threshold};
\]

%dsp-level:0
\[
\qquad [\Phi_2]=\mathfrak q^0,
\]

%dsp-level:0
\[
\qquad [F]=\mathfrak q^{+1}
\]

%dsp-level:0
\[
\qquad [\text{machine: the exponent census}];
\]

%dsp-level:0
\[
(5)\quad\textbf{the two resolvents, two species:}
\]

%dsp-level:0
\[
\qquad \text{the chart resolvent}\ R(z;\mathbf p)\ [(1)]\ \text{and the register's zig-zag resolvent}\ \frac1{1-w^2x_+x_-}\ [\mathrm{BM.121}(2);\ \mathrm{BM.23}(7)]\ \textbf{are}\ \neg\text{identified}:
\]

%dsp-level:0
\[
\text{the first is meromorphic in the chart energy with poles at the dispersion branches};
\]

%dsp-level:0
\[
\qquad \text{the second is a formal even series in the ledger weight with its pole at}\ |w|\ge2\ [\mathrm{BM.121}(5)];
\]

%dsp-level:0
\[
\boxed{\ \text{they meet only through the threshold indicator}\ \chi_{\{NE\ge\mathcal T_0\}}\ }
\]

%dsp-level:0
\[
\qquad [\mathrm{Def\text{-}BM.17}(3);\ \mathrm{BM.122}(3);\ \text{the currency separation stands}];
\]

%dsp-level:0
\[
(6)\quad\textbf{the junction domain:}
\]

%dsp-level:0
\[
\qquad \text{the}\ \aleph\text{-reads at the sealed junctions where the boundary-channel states couple run in the}\ \mathrm{E2}\ \text{domain}\ [\text{A-BM.10},\ \text{issued in}\ \mathrm{BM.139}(2);\ \text{forward pointer}];
\]

%dsp-level:0
\[
(7)\quad\textbf{the closure row:}
\]

%dsp-level:0
\[
\qquad \boxed{\ \textbf{the momentum-space}\ \Sigma\textbf{-resolvent is constructed}\ }
\]

%dsp-level:0
\[
\qquad [\mathrm{Def\text{-}BM.20};\ \text{machine: the discharge census against}\ \mathrm{BM.6}(6)\text{'s three requirements — plane-wave boundary-channel states}\ (1),\ Z(\mathcal R,\Sigma;k^2)\text{-type two-point read}\ (2),\ m^2\ \text{as the pole displacement}\ (1)];
\]

%dsp-level:0
\[
\textbf{what it unblocks:}
\]

%dsp-level:0
\[
\qquad \text{the pair amplitude and the annihilation cross-section}\ [\mathrm{BM.142}];
\]

%dsp-level:0
\[
\qquad \textbf{what stays barred:}
\]

%dsp-level:0
\[
\qquad \neg\text{an S-matrix element beyond the formed amplitudes},\quad\neg\text{a measured cross-section}\ [\mathrm{P5}].
\]

\textbf{Req.} D2.5, D2.9 (Parts 1–2); CT2.1, CT2.6 (Part 2); P5 (Parts 1–2); FT.26(1), FT.79(5) (Part 3); Def-BM.10, Def-BM.17(3), Def-BM.20, BM.6(6), BM.10(1), BM.23(7), BM.34(1), BM.34(4), BM.34(7), BM.119(1), BM.121(2), BM.121(5), BM.122(3), BM.132(2) (this part); A-BM.10, A-BM.12 (§II.4, logged before use; forward pointer for A-BM.10's issuance row); the delta-composition rule, the exact linear solve on rational witnesses, the two-body phase-space integral in the declared convention, Laurent exponent censuses [math; machine: \texttt{checks/\allowbreak{}verify\_\allowbreak{}H\_\allowbreak{}bm\_\allowbreak{}closing\_\allowbreak{}tier.py} §C].

\textbf{Pf.}

%dsp-level:0
\[
(1):\ \mathrm{BM.34}(4)\text{'s master identity re-derived entrywise on three rational ADM witnesses}\ (\Omega^2,\varkappa,\beta,w_i,q_{ij}\ \text{rational,}\ q\ \text{off-diagonal});
\]

%dsp-level:0
\[
\qquad \text{the branch solve and the rest substitution}\ [\text{machine}];
\]

%dsp-level:0
\[
\qquad \text{the vacuum gauge:}\ N=1,\ w=0,\ \gamma=\delta\implies R=(p_0^2-m^2-p^2)^{-1}\ [\text{machine}]
\]

%dsp-level:0
\[
(2):\ \text{the conditions as linear constraints on}\ (\alpha,\beta,\delta):
\]

%dsp-level:0
\[
\quad (\mathrm i)\ \text{the zero set coincides with the mass shell}\implies\alpha(p_0-N^ip_i)^2+\beta N^2(m^2+p^2)+\delta\ \text{is proportional to}\ (p_0-N^ip_i)^2-N^2(m^2+p^2)\implies\delta
\]

%dsp-level:0
\[
= 0,
\]

%dsp-level:0
\[
\ \beta\qquad = -\alpha;
\]

%dsp-level:0
\[
\qquad (\mathrm{ii})\ \text{fixes}\ \alpha=1\ [\text{machine: the exact solve over}\ \mathbb Q];
\]

%dsp-level:0
\[
\qquad (\mathrm{iii}),(iv),(\mathrm v)\ \text{are then automatic}\ [\text{machine: the substitution rows}]
\]

%dsp-level:0
\[
(3):\ \delta(g(p_0))=\sum_{\text{roots}}\delta(p_0-r_i)/|g'(r_i)|,
\]

%dsp-level:0
\[
\qquad g'(p_0)=2(p_0-N^ip_i)=\pm2NE_{\mathbf p}\ \ \text{at the roots}\ [\text{math}];
\]

%dsp-level:0
\[
\qquad \text{machine: the coefficient at}\ (N,\mathbf p,m)\in\{(1,\cdot,1),(2,\cdot,1),(1,\cdot,2),(3,\cdot,2),(2,\cdot,3)\}\ \text{rational grids};
\]

%dsp-level:0
\[
\qquad \text{the exponent census}\ [\mathrm{CT2.1}:\ [d^3p]=\mathfrak q^{+3/2},\ [E]=\mathfrak q^{+1/2}]
\]

%dsp-level:0
\[
(4):\ \int\frac{d^3p}{(2\pi)^3}\frac{(2\pi)^4\delta(z-E_1-E_2)}{4E_1E_2}\bigg|_{\mathbf p_2=-\mathbf p_1}
\]

%dsp-level:0
\[
= \frac1{4(2\pi)^2}\int d\Omega\,
\]

%dsp-level:0
\[
\frac{p^2\,dp}{E_1E_2}\delta(z-E_1-E_2)
\]

%dsp-level:0
\[
= \frac{4\pi}{16\pi^2}\cdot\frac{p_\ast}{N^2z}\ [\text{machine: the Jacobian}\ \tfrac{d(E_1+E_2)}{dp}=\tfrac{2N^2p_\ast}z\ \text{at five rational witnesses}];
\]

%dsp-level:0
\[
\qquad \text{the fluxes from}\ p_1\!\cdot\!p_2\ \text{at zero shift}\ [\text{machine}]
\]

%dsp-level:0
\[
(5),
\]

%dsp-level:0
\[
(6),
\]

%dsp-level:0
\[
(7):\ \text{the species census, the domain row, and the discharge census against}\ \mathrm{BM.6}(6)\text{'s specification text}\ [\text{machine: the document pins}].\ \blacksquare
\]

\textbf{Phys.} The resolvent's declared surface is two naming choices; everything else is forced. Its poles are exactly the derived dispersion's two branches, its rest pole reads the composite's mass as the displacement against the massless pole, and among all rational forms of its degree the five derived conditions pin it uniquely. The phase space falls out of the same mass shell: the on-shell measure is the delta reduction of the chart measure, the two-body phase space and both fluxes are exact functions of $(z,m,N)$, formed precisely above threshold. The chart resolvent and the register's zig-zag resolvent stay separate species that meet only at the threshold indicator, and the item the corpus carried since the fermion era — the momentum-space $\Sigma$-resolvent — closes as constructed.

\end{theorem}

\begin{theorem}[{\texorpdfstring{BM.135 — The multi-chain additive charge census: the baryon and meson rows, the stretched states, and the gauged read}{BM.135 - The multi-chain additive charge census: the baryon and meson rows, the stretched states, and the gauged read}}]\label{tag:BM-135}
\textbf{Thm.}

%dsp-level:0
\[
(1)\quad\textbf{the read, re-certified:}
\]

%dsp-level:0
\[
\qquad Q(c_1,\dots,c_k)=\sum_{i=1}^{k}q_{\mathrm{core}}(c_i),
\]

%dsp-level:0
\[
\qquad q_{\mathrm{core}}(c)=\tfrac12\bigl(\sigma_{s(c)}-\sigma_{t(c)}\bigr)\in\{-1,0,+1\}
\]

%dsp-level:0
\[
\qquad [\mathrm{Def\text{-}BM.8}(1);\ \mathrm{FT.40};\ \mathrm{EM.9};\ \mathrm{D3.4}];
\]

%dsp-level:0
\[
\textbf{the three specification clauses stand discharged:}
\]

%dsp-level:0
\[
\qquad Q|_{\text{single chain}}=q_{\mathrm{core}};
\]

%dsp-level:0
\[
\qquad Q(\mathcal X\sqcup\mathcal Y)=Q(\mathcal X)+Q(\mathcal Y)\ \ \text{on disjoint sealings};
\]

%dsp-level:0
\[
\qquad Q\ \text{defined only on sealed supports}\ \bigl(\mathcal C=1\bigr)
\]

%dsp-level:0
\[
\qquad [\mathrm{Def\text{-}BM.8}(2)\text{--}(4);\ \mathrm{BM.6}(1);\ \mathrm{FT.19}];
\]

%dsp-level:0
\[
(2)\quad\textbf{the baryon charge (the}\ \varepsilon\textbf{-sealed 3-juncture):}
\]

%dsp-level:0
\[
\qquad Q_B=q_1+q_2+q_3\in\{-3,\dots,+3\},
\]

%dsp-level:0
\[
\qquad W_3(Q)=\binom6{Q+3}=(1,\ 6,\ 15,\ \mathbf{20},\ 15,\ 6,\ 1)
\]

%dsp-level:0
\[
\qquad [\mathrm{Def\text{-}BM.8}(5),(6);\ \mathrm{BM.28}(2);\ \mathrm{BM.114}(1);\ \text{machine: the exhaustive}\ 4^3{=}64\ \text{end-sign census re-run}];
\]

%dsp-level:0
\[
\textbf{the named-class rows:}
\]

%dsp-level:0
\[
\qquad \text{the neutral class:}\ Q_B=0,
\]

%dsp-level:0
\[
\qquad E_{\mathrm{int}}^{(n)}=6
\]

%dsp-level:0
\[
\qquad [\mathrm{BM.17};\ \mathrm{Def\text{-}BM.3}(2);\ \mathrm{Census\ FT.35}(2)];
\]

%dsp-level:0
\[
\qquad \text{the unit class:}\ Q_B=+1
\]

%dsp-level:0
\[
\qquad [\mathrm{BM.26};\ \mathrm{Def\text{-}BM.3}(3)];
\]

%dsp-level:0
\[
\qquad \text{the CPT image:}\ Q_B=-1
\]

%dsp-level:0
\[
\qquad [\mathrm{BM.27}];
\]

%dsp-level:0
\[
\textbf{the transport profile of the neutral class:}
\]

%dsp-level:0
\[
\qquad m_q=(2,1,1)
\]

%dsp-level:0
\[
\qquad [\mathrm{Def\text{-}BM.13}(1);\ \mathrm{FT.106}(3\mathrm a)];
\]

%dsp-level:0
\[
\qquad \frac{W(0)}{W(\pm1)}=\frac{20}{15}=\frac43
\]

%dsp-level:0
\[
\qquad [\mathrm{BM.28}(2);\ \text{machine}];
\]

%dsp-level:0
\[
(3)\quad\textbf{the meson charge (the}\ \delta\textbf{-sealed 2-juncture):}
\]

%dsp-level:0
\[
\qquad Q_M=q_1+q_2\in\{-2,\dots,+2\},
\]

%dsp-level:0
\[
\qquad W_2(Q)=\binom4{Q+2}=(1,\ 4,\ \mathbf6,\ 4,\ 1)
\]

%dsp-level:0
\[
\qquad [\mathrm{BM.110}(3);\ \mathrm{BM.114}(1);\ \text{machine: the}\ 4^2\ \text{census re-run}];
\]

%dsp-level:0
\[
\textbf{the CP-pair rows:}\qquad Q_M\bigl(c,c^{\mathrm{CP}}\bigr)=q+(-q)=0\ \ \text{for all three}\ q\qquad[\mathrm{Def\text{-}BM.15}(3);\ \text{machine: exhaustive}];
\]

%dsp-level:0
\[
\textbf{the CP-pairs saturate the neutral class:}
\]

%dsp-level:0
\[
\qquad 1+4+1=6=W_2(0),
\]

%dsp-level:0
\[
\qquad \#\{Q{=}0\ \text{tuples outside the CP-pair forms}\}=0
\]

%dsp-level:0
\[
\qquad [\mathrm{BM.110}(3);\ \text{machine re-run}];
\]

%dsp-level:0
\[
(4)\quad\textbf{the stretched and unit rows (both charge-read questions adjudicated):}
\]

%dsp-level:0
\[
\qquad \boxed{\ \#|Q|=2\ \ \text{states:}\ \ 12\ \ \text{at}\ k=3,\quad2\ \ \text{at}\ k=2;\qquad\#|Q|=3\ \ \text{states:}\ \ 2\ \ \text{at}\ k=3\ }
\]

%dsp-level:0
\[
\qquad [W_3(\pm2)=6,\ W_3(\pm3)=1,\ W_2(\pm2)=1;\ \text{machine: exhaustive}];
\]

%dsp-level:0
\[
\textbf{the multi-chain charge read is derived:}
\]

%dsp-level:0
\[
\qquad |Q|=2\ \ \text{and}\ \ |Q|=3\ \ \text{observable integer states exist on sealed supports of the additive read}
\]

%dsp-level:0
\[
\qquad [\text{the requirement is discharged};\ \mathrm{FT.41}(\mathrm a):\ N_{\mathrm{charge}}=2\ \text{end registers per chain re-cited}];
\]

%dsp-level:0
\[
\textbf{the ternary-read row stands re-certified:}
\]

%dsp-level:0
\[
\qquad \text{the read is integer-valued at every}\ k:\quad\tfrac13,\ \tfrac23\ \ \text{are structurally unreachable}
\]

%dsp-level:0
\[
\qquad [\mathrm{Def\text{-}BM.8}(5);\ \mathrm{BM.28}(7);\ \text{machine: the sum-closure scan over}\ \{-1,0,1\}^k,\ k=1..4];
\]

%dsp-level:0
\[
(5)\quad\textbf{the gauged read:}
\]

%dsp-level:0
\[
\qquad \partial_\mu F^{\mu\nu}=e\,J_Q^\nu,
\]

%dsp-level:0
\[
\qquad \text{the}\ \mathsf A\text{-coefficient}\ e\,Q,
\]

%dsp-level:0
\[
\qquad e=\frac{gg'}{\sqrt{g^2+g'^2}}
\]

%dsp-level:0
\[
\qquad [\mathrm{Def\text{-}BM.12}(1)\text{--}(3);\ \mathrm{EW.35};\ \mathrm{EM.15}];
\]

%dsp-level:0
\[
\textbf{the running entry:}
\]

%dsp-level:0
\[
\qquad e=e(r,E)\ \ \text{with all three factors derived}
\]

%dsp-level:0
\[
\qquad [\mathrm{BM.88}(4);\ \mathrm{BM.132}(2);\ \mathrm{BM.133}(6)];
\]

%dsp-level:0
\[
\mathrm{Def\text{-}BM.8}(8)\text{'s}\ \neg\text{constructed row (the coupling of}\ Q\ \text{to a gauge field)}\ \textbf{is discharged by}\ \mathrm{Def\text{-}BM.12}(2),
\]

%dsp-level:0
\[
(3)\ \text{and re-scoped,}\ \neg\text{retracted};
\]

%dsp-level:0
\[
(6)\quad\textbf{additivity at general}\ n\ \textbf{and the grammar row:}
\]

%dsp-level:0
\[
\qquad \text{a sealed support partitions into blocks of size}\ 2\ \text{and}\ 3
\]

%dsp-level:0
\[
\qquad [\mathrm{CT3.17a}(1)];
\]

%dsp-level:0
\[
\qquad Q(\Gamma)=\sum_{\text{blocks}}Q(\text{block})
\]

%dsp-level:0
\[
\qquad [\text{machine: exhaustive at}\ n\le4\ \text{on ternary alphabets}];
\]

%dsp-level:0
\[
\text{the}\ \delta/\varepsilon\text{-mixing from}\ N
\]

%dsp-level:0
\[
= 4\ [\mathrm{BM.31}(2);\ \mathrm{BM.33}]\ \text{governs which supports are sealed,}\ \neg\text{the additivity of the read on sealed blocks};
\]

%dsp-level:0
\[
(7)\quad\textbf{the closure row:}
\]

%dsp-level:0
\[
\qquad \boxed{\ \textbf{the additive multi-chain charge read exists, is constructed, and is gauged on the chart side}\ }
\]

%dsp-level:0
\[
\qquad [\mathrm{Def\text{-}BM.8};\ (1)\text{--}(6);\ \text{machine: the three-clause discharge census against}\ \mathrm{BM.6}(1)\text{'s specification}];
\]

%dsp-level:0
\[
\text{the ternary-read row is complete: no residue at any chain count};
\]

%dsp-level:0
\[
\qquad \neg\text{a species name attaches to any charge class beyond the identification layer}\ [\mathrm{Def\text{-}BM.3}(7)];
\]

%dsp-level:0
\[
\qquad \neg\text{a measured charge enters}\ [\mathrm{P5}].
\]

\textbf{Req.} D2.14, D3.4, D3.9 (Parts 1–2); CT3.17a, CT3.22, EM.9, EM.11, EM.15, EW.35 (Part 2); P2, P5 (Parts 1–2); FT.19, FT.40, FT.41, FT.106(3) (Part 3); Def-BM.3, Def-BM.8, Def-BM.12, Def-BM.13, Def-BM.15, BM.6(1), BM.17, BM.26, BM.27, BM.28(2), BM.28(7), BM.31(2), BM.33, BM.88(4), BM.110(3), BM.114(1), BM.132(2), BM.133(6) (this part); Census FT.35(2) (row-wise); exhaustive ternary enumeration, generating-function convolution, sum-closure scans [math; machine: \texttt{checks/\allowbreak{}verify\_\allowbreak{}H\_\allowbreak{}bm\_\allowbreak{}closing\_\allowbreak{}tier.py} §D].

\textbf{Pf.}

%dsp-level:0
\[
(1)\text{--}(3):\ \mathrm{Def\text{-}BM.8}\text{'s construction re-certified clause by clause and the two censuses re-run exhaustively}\ [\text{machine}];
\]

%dsp-level:0
\[
\qquad \text{the named-class rows consumed from their own tags row-wise}\ [\mathrm{Def\text{-}BM.3}(7)\text{'s discipline}]
\]

%dsp-level:0
\[
(4):\ \text{the counts from}\ W_3,W_2\ \text{directly}\ [\text{machine}];
\]

%dsp-level:0
\[
\qquad \text{the unreachability scan:}\ \textstyle\sum_iq_i\in\mathbb Z\ \forall q_i\in\{-1,0,1\}\ \text{at}\ k\le4\ [\text{machine}]
\]

%dsp-level:0
\[
(5):\ \mathrm{Def\text{-}BM.12}(1)\text{--}(3)\text{'s rows re-cited with the running coupling substituted}\ [\mathrm{BM.133}(6)];
\]

%dsp-level:0
\[
\qquad (6):\ \text{the block sum on ternary alphabets, exhaustive at}\ n\le4\ [\text{machine}];
\]

%dsp-level:0
\[
\qquad (7):\ \text{the discharge census}.\ \blacksquare
\]

\textbf{Phys.} The additive charge read the era opened as its first construction spec is re-certified, censused and gauged. Baryons carry seven integer charge classes with the frozen weight row and the neutral-to-unit ratio four thirds; mesons carry five, and their CP pairs saturate the neutral class exactly. The stretched states settle two ledger rows at once: observable charge two exists — twelve states at three chains, two at two — so the integer charge read is derived, while the fractional alphabet stays structurally unreachable, so the ternary read's consistent no-go stands. The chart-side gauge coupling turns the read into a current with a fully derived running coefficient, and the question the corpus carried since the fermion era — whether an additive multi-chain charge read exists — closes as constructed.

\end{theorem}

\begin{theorem}[{\texorpdfstring{BM.136 — The $S_3$ action is the colour 3-vertex: the equivariance theorem, the exchange structure at every $n$, and the statistics per class}{BM.136 - The action is the colour 3-vertex: the equivariance theorem, the exchange structure at every , and the statistics per class}}]\label{tag:BM-136}
\textbf{Thm.}

%dsp-level:0
\[
(1)\quad\textbf{the colour 3-vertex and its equivariance:}\qquad B=\varepsilon_{ijk}\,c_1^ic_2^jc_3^k\qquad[\mathrm{D3.9};\ \mathrm{Def\text{-}BM.1};\ \mathrm{CT3.15}];
\]

%dsp-level:0
\[
\boxed{\ \varepsilon_{\sigma(i)\sigma(j)\sigma(k)}=\operatorname{sgn}(\sigma)\,\varepsilon_{ijk}\ \ \forall\sigma\in S_3\ }
\]

%dsp-level:0
\[
\qquad [\text{machine: exhaustive over the}\ 6\times27\ \text{index pairs}];
\]

%dsp-level:0
\[
\dim\wedge^3(\mathbb C^3)=\binom33=1:
\]

%dsp-level:0
\[
\qquad \textbf{the}\ \varepsilon\textbf{-seal is the unique invariant line, and}\ S_3\ \textbf{acts on it by the sign representation}
\]

%dsp-level:0
\[
\qquad [\mathrm{BM.128}(5);\ \text{machine: the antisymmetrizer rank re-run}];
\]

%dsp-level:0
\[
\text{colour is}\ \neg\text{observable by construction}
\]

%dsp-level:0
\[
\qquad [\mathrm{CT3.22}];
\]

%dsp-level:0
\[
\qquad \text{the colour generators act as}\ \mathfrak{su}(3)\ \text{on the slots}
\]

%dsp-level:0
\[
\qquad [\mathrm{CT3.8};\ \mathrm{FT.82}(3)];
\]

%dsp-level:0
\[
(2)\quad\textbf{the intertwiner theorem (the}\ S_3\ \textbf{action is literally the colour 3-vertex):}
\]

%dsp-level:0
\[
\qquad \text{the sealed baryon is the contraction}\ \langle\varepsilon|\,
\]

%dsp-level:0
\[
:\ (\mathbb C^3)^{\otimes3}\otimes\mathcal H_3\to\mathbb C\ \ \text{with}\ \hat P_\sigma\ \text{the slot permutation}\ [\mathrm{BM.128}(2)]:
\]

%dsp-level:0
\[
\boxed{\ \bigl(\hat P_\sigma\otimes D_{\mathrm{col}}(\sigma)\bigr)\,|\varepsilon\rangle\otimes|\psi\rangle\ \ \text{is}\ \ S_3\textbf{-invariant exactly when the slot read compensates the colour sign:}\ \ \hat P_\sigma\big|_{\text{sealed}}=\operatorname{sgn}(\sigma)\ }
\]

%dsp-level:0
\[
[\text{machine: the explicit}\ 27\text{-dimensional colour tensor}\ |\varepsilon\rangle=\sum_{abc}\varepsilon_{abc}|abc\rangle\ \text{and all six}\ \sigma\in S_3;\ D_{\mathrm{col}}(\sigma)|\varepsilon\rangle=\operatorname{sgn}(\sigma)|\varepsilon\rangle\ \text{entrywise}];
\]

%dsp-level:0
\[
\textbf{the reading:}
\]

%dsp-level:0
\[
\qquad \text{the sealing grammar's antisymmetry and the slot space's permutation action are the two halves of one equivariant map}\ —\ \text{the}\ S_3\ \text{action on the 3-chain state space is the colour 3-vertex read on slots};
\]

%dsp-level:0
\[
(3)\quad\textbf{the exchange structure at every}\ n\ \textbf{(residue}\ (\mathrm{R7}')\ \textbf{closed):}
\]

%dsp-level:0
\[
\qquad \boxed{\ \mathcal H_n\ :=\ \bigotimes_{a=1}^{n}\Bigl(\mathcal H_\omega(c_a)\otimes\Psi(c_a)\Bigr),\qquad\dim\mathcal H_n=2^n\prod_{a=1}^{n}\bigl(|k_a|+1)\ }\qquad[\mathrm{BM.128}(1)\ \text{extended uniformly}];
\]

%dsp-level:0
\[
\hat P_\sigma\hat P_\tau=\hat P_{\sigma\tau}
\]

%dsp-level:0
\[
\qquad \forall\sigma,\tau\in S_n
\]

%dsp-level:0
\[
\qquad [\mathrm{BM.128}(2)\text{'s convention, uniform in}\ n;\ \text{machine: all}\ 576\ \text{ordered pairs at}\ n=4,\ d=2,\ \text{entrywise over}\ \mathbb Q];
\]

%dsp-level:0
\[
\boxed{\ \chi(\sigma)=\operatorname{Tr}\hat P_\sigma\big|_{(\mathbb C^d)^{\otimes n}}=d^{\,c(\sigma)}\ \ \text{for every}\ n\ }
\]

%dsp-level:0
\[
\qquad [\mathrm{BM.128}(3)\ \text{extended};\ \text{machine: all}\ 24\ \text{permutations at}\ (n,d)=(4,2)\ \text{and}\ (4,3):\ \text{the five cycle types give}\ (16,8,4,4,2)\ \text{at}\ d=2\ \text{and}\ (81,27,9,9,3)\ \text{at}\ d=3];
\]

%dsp-level:0
\[
\varepsilon(\Gamma)=\prod_{a=1}^{n}\varepsilon(c_a)
\]

%dsp-level:0
\[
\qquad \text{at an open connector}
\]

%dsp-level:0
\[
\qquad [\mathrm{BM.128}(7)\ \text{extended};\ \text{machine: all}\ 625\ \text{quadruples in}\ \{0,\dots,4\}^4:\ \ 312\ \text{fermionic},\ 313\ \text{bosonic}];
\]

%dsp-level:0
\[
\textbf{the wedge rows:}
\]

%dsp-level:0
\[
\qquad \dim\wedge^n(\mathbb C^2)=\binom2n=0\ \ \forall n\ge3;
\]

%dsp-level:0
\[
\qquad \dim\wedge^3(\mathbb C^3)=1;
\]

%dsp-level:0
\[
\qquad \dim\operatorname{Inv}\bigl((\mathbb C^3)^{\otimes4}\bigr)=0
\]

%dsp-level:0
\[
\qquad [\text{machine: the antisymmetrizer ranks}\ (2,4),(3,4),(2,3),(3,3)\ \text{exact};\ \text{the weight census: no zero-weight vector in}\ (\mathbb C^3)^{\otimes4}];
\]

%dsp-level:0
\[
\boxed{\ \textbf{four colour slots admit no singlet: the}\ N\ge4\ \textbf{sealings must mix}\ \delta\ \textbf{and}\ \varepsilon\ \textbf{blocks}\ }
\]

%dsp-level:0
\[
\qquad [\mathrm{BM.31}(2);\ \mathrm{BM.33};\ \mathrm{CT3.17a}(2):\ g(4)=10\ \text{groupings};\ \text{machine: the census of}\ \{2,3\}\text{-partitions at}\ N=4];
\]

%dsp-level:0
\[
\textbf{the grammar rescope row:}
\]

%dsp-level:0
\[
\qquad \text{the}\ \delta/\varepsilon\text{-mixing governs which}\ n\text{-chain configurations are sealed (observability}\ [\mathrm{BM.1}]),
\]

%dsp-level:0
\[
\ \neg\text{the exchange structure, which is uniform in}\ n;
\]

%dsp-level:0
\[
(4)\quad\textbf{the winding assignment per class:}
\]

%dsp-level:0
\[
\qquad \textbf{(a) the computed floor:}\ \ \text{the minimal materialized constituent is the digon with}\ k=1
\]

%dsp-level:0
\[
\qquad [\mathrm{FT.28};\ \mathrm{FT.32}(6);\ \mathrm{FT.107}(1);\ \mathrm{BM.129}(5)]:
\]

%dsp-level:0
\[
\boxed{\ (k_1,k_2,k_3)=(1,1,1):\qquad\varepsilon_B=(-1)^3=-1,\qquad j_B=\frac32\ }\qquad[\mathrm{BM.129}(4),(5);\ \text{machine}];
\]

%dsp-level:0
\[
\textbf{the neutral and unit classes share the floor:}
\]

%dsp-level:0
\[
\qquad E_{\mathrm{int}}
\]

%dsp-level:0
\[
= 6\ \text{on both}\ [\mathrm{BM.17};\ \mathrm{BM.26};\ \mathrm{BM.28}(1)]\implies\textbf{both named baryon classes are fermionic at the computed floor, derived}\ ;
\]

%dsp-level:0
\[
\mathrm{BM.129}(6),
\]

%dsp-level:0
\[
\ \mathrm{BM.130}(1)\text{'s ``value per named class awaits}\ (\mathrm{R3})\text{'' row is discharged at the floor locus:}\ \varepsilon_B\ \text{depends on the winding alone, and the winding at the floor is derived};
\]

%dsp-level:0
\[
\textbf{(b) the six-state classes under the declared corner convention:}
\]

%dsp-level:0
\[
\qquad k_{vf}=(3,3,2,3,5)\ \ \text{on}\ (\text{tetra},\ \text{cube},\ \text{octa},\ \text{dodeca},\ \text{icosa})
\]

%dsp-level:0
\[
\qquad [\text{A-BM.11 (variant B), issued in}\ \mathrm{BM.139}(4);\ \mathrm{FT.76}(1);\ \mathrm{Def\text{-}BM.14}(1),(2):\ \text{the floor map}\ 2E\in\{12,24,60\}\ \text{gives}\ k_{vf}=(3,3\,|\,2,2\,|\,3,5)\ \text{on the six states}]:
\]

%dsp-level:0
\[
\varepsilon_B(\mathbf s)=(-1)^{k_{vf}(s_1)+k_{vf}(s_2)+k_{vf}(s_3)}:
\]

%dsp-level:0
\[
\qquad \boxed{\ \#\{\text{fermionic}\}=112,\qquad\#\{\text{bosonic}\}=104\ }
\]

%dsp-level:0
\[
\qquad [\text{machine: the exhaustive}\ 6^3{=}216\ \text{census}];
\]

%dsp-level:0
\[
\textbf{the variant dependence is exact:}
\]

%dsp-level:0
\[
\qquad \text{variant A}\ \bigl(k_f=(3,2\,|\,3,3\,|\,5,3)\bigr):\quad140\ \ \text{fermionic},\ 76\ \ \text{bosonic}
\]

%dsp-level:0
\[
\qquad [\text{machine}]\implies\textbf{the census is convention-dependent and the declaration}\ \text{A-BM.11}\ \textbf{pins it};
\]

%dsp-level:0
\[
(5)\quad\textbf{the bars:}
\]

%dsp-level:0
\[
\qquad \neg\text{a flavour}\ \mathsf{SU}(3),\ \ \neg\text{an}\ \mathsf{SU}(6):
\]

%dsp-level:0
\[
\quad S_n\ \text{permutes constituent slots}\ [\mathrm{BM.117}(3);\ \mathrm{BM.128}(9);\ \mathrm{BM.129}(7)\ \text{re-cited}];
\]

%dsp-level:0
\[
\textbf{the flavour occupancy remains}\ (\mathrm{R3})\textbf{-species:}
\]

%dsp-level:0
\[
\qquad \text{which three of the six states a named class occupies is the generation-selection question}\ [\mathrm{BM.60}(1):\ \text{three lines derived, three selected}\ \neg\text{derived}]\ —\ \textbf{named for the Part-5 ledger,}\ \neg\text{a baryomesonic residue};
\]

%dsp-level:0
\[
\neg\text{a species name on any row}\ [\mathrm{Def\text{-}BM.3}(7)];\qquad\neg\text{a measured value}\ [\mathrm{P5}];
\]

%dsp-level:0
\[
(6)\quad\textbf{status:}
\]

%dsp-level:0
\[
\qquad \textbf{derived:}\ (1)\text{--}(4);
\]

%dsp-level:0
\[
\qquad \boxed{\ (\mathrm{R7}')\ \textbf{is closed:}\ \mathcal H_n,\ \hat P_\sigma,\ \chi(\sigma)=d^{\,c(\sigma)},\ \varepsilon(\Gamma)=\prod_a\varepsilon(c_a)\ \ \text{at every}\ n};
\]

%dsp-level:0
\[
\qquad \textbf{the}\ S_3\textbf{-action gap is closed: the action is the colour vertex's slot read};
\]

%dsp-level:0
\[
\textbf{the winding assignment per class is derived at both loci (floor and six-state);}
\]

%dsp-level:0
\[
\qquad \text{the statistics function is complete and its values are derived per class}\ [\mathrm{BM.126};\ \mathrm{BM.129}(4)].
\]

\textbf{Req.} A3 (Part 1); D2.14, D3.9 (Parts 1–2); CT3.8, CT3.15, CT3.17a, CT3.22 (Part 2); P5 (Parts 1–2); FT.28, FT.32(6), FT.76(1), FT.82(3), FT.107(1) (Part 3, Census); Def-BM.1, Def-BM.3(7), Def-BM.14, BM.1, BM.17, BM.26, BM.28(1), BM.31(2), BM.33, BM.60(1), BM.117(3), BM.126, BM.128, BM.129 (this part); A-BM.11 (§II.4, logged before use; forward pointer for its issuance row); permutation-module enumeration, antisymmetrizer ranks by exact Fraction elimination, the weight-zero census of $(\mathbb C^3)^{\otimes4}$, exhaustive parity censuses [math; machine: \texttt{checks/\allowbreak{}verify\_\allowbreak{}H\_\allowbreak{}bm\_\allowbreak{}closing\_\allowbreak{}tier.py} §E].

\textbf{Pf.}

%dsp-level:0
\[
(1):\ \text{the}\ 6\times27\ \text{equivariance census and the rank-1 antisymmetrizer}\ [\text{machine}];
\]

%dsp-level:0
\[
\qquad (2):\ D_{\mathrm{col}}(\sigma)|\varepsilon\rangle=\operatorname{sgn}(\sigma)|\varepsilon\rangle\ \text{entrywise on the explicit tensor};
\]

%dsp-level:0
\[
\qquad \text{the invariance of}\ |\varepsilon\rangle\otimes|\psi\rangle\ \text{under the diagonal action with}\ \hat P_\sigma\ \text{on slots}\ [\text{machine}]
\]

%dsp-level:0
\[
(3):\ \hat P_\sigma\ \text{built as the factor permutation on}\ (\mathbb C^d)^{\otimes n}:\ \ \text{the product law on all}\ 576\ \text{pairs at}\ n
\]

%dsp-level:0
\[
= 4\ \text{in}\ \mathrm{BM.128}(2)\text{'s composition convention}\ [\text{machine: the pushforward factor action}];
\]

%dsp-level:0
\[
\qquad \chi(\sigma)=d^{c(\sigma)}\ \text{on all}\ 24\ \text{permutations at}\ d\in\{2,3\}\ [\text{machine}];
\]

%dsp-level:0
\[
\qquad \text{the antisymmetrizer ranks}\ (0,0,0,1)\ \text{at}\ ((2,4),(3,4),(2,3),(3,3))\ [\text{machine}];
\]

%dsp-level:0
\[
\qquad \dim\operatorname{Inv}((\mathbb C^3)^{\otimes4})
\]

%dsp-level:0
\[
= 0:\ \text{the weights of}\ \mathbb C^3\ \text{are}\ \{(1,0),(0,1),(-1,-1)\}:\ \ a(1,0)+b(0,1)+c(-1,-1)
\]

%dsp-level:0
\[
= (0,0)\ \wedge\ a+b+c
\]

%dsp-level:0
\[
= 4
\]

%dsp-level:0
\[
\implies a
\]

%dsp-level:0
\[
= b
\]

%dsp-level:0
\[
= c
\]

%dsp-level:0
\[
= 4/3\notin\mathbb Z\ [\text{machine: the exhaustive composition census}];
\]

%dsp-level:0
\[
\qquad \text{the parity census on}\ 625\ \text{quadruples}\ [\text{machine}]
\]

%dsp-level:0
\[
(4):\ \text{the floor row from}\ \mathrm{BM.129}(5)\text{'s digon line and}\ \varepsilon_B=(-1)^{k_1+k_2+k_3}\ [\mathrm{BM.129}(4)];
\]

%dsp-level:0
\[
\qquad \text{the six-state censuses:}\ 112/104\ \text{and}\ 140/76\ \text{by exhaustive enumeration of the parity patterns}\ (4\ \text{odd},\ 2\ \text{even})\ \text{and}\ (5\ \text{odd},\ 1\ \text{even})\ [\text{machine}];
\]

%dsp-level:0
\[
\qquad (5),(6):\ \text{assembly}.\ \blacksquare
\]

\textbf{Phys.} The permutation action on the three-chain state space is not a bookkeeping device: it is the colour 3-vertex read on slots. The epsilon tensor is the unique invariant line of three colour slots, it transforms by the sign of the permutation, and the sealed baryon is exactly the intertwiner that pairs the two actions — so the antisymmetry lives in colour, unobservable by construction, while the observable statistics lives in the winding. The exchange structure extends uniformly to every chain count with the character formula intact, and the four-slot colour census explains the grammar: no four-fold singlet exists, which is why seals mix delta and epsilon blocks from four attachments on. The winding assignment per class is derived at both loci — the digon floor gives both named baryon classes the fermion sign and spin three halves, and the six-state census under the declared vertex-figure convention gives 112 fermionic against 104 bosonic classes, a census that provably depends on the corner convention and is pinned by its declaration. What remains open is not baryomesonic: which states a named class occupies is the generation-selection question, and it leaves for the next era.

\end{theorem}

\begin{theorem}[{\texorpdfstring{BM.137 — The two-chain ensemble measure, the meson's formed binding, the sector-mass pincer, and the two-chain spectrum}{BM.137 - The two-chain ensemble measure, the meson's formed binding, the sector-mass pincer, and the two-chain spectrum}}]\label{tag:BM-137}
\textbf{Thm.}

%dsp-level:0
\[
(1)\quad\textbf{the general pair-specialization law:}
\]

%dsp-level:0
\[
\qquad \text{an}\ n\text{-pair floor crystallizes each pair independently with the stationary weight}\ p=1-4\varepsilon_q
\]

%dsp-level:0
\[
\qquad [\mathrm{BM.19}(2);\ \mathrm{BM.94}(4)\text{'s marked specialization, extended uniformly}]:
\]

%dsp-level:0
\[
P(E{=}2j)=\binom njp^ju^{n-j},
\]

%dsp-level:0
\[
\qquad u:=1-p=4\varepsilon_q;
\]

%dsp-level:0
\[
\qquad \text{the materialization restriction}\ E\ge2\ \ \text{is derived}
\]

%dsp-level:0
\[
\qquad [\mathrm{BM.94}(5);\ \mathrm{FT.144}(4);\ \mathrm{Def\text{-}H.4}];
\]

%dsp-level:0
\[
\boxed{\ D_n\ :=\ \frac{1-u^n}{1-u}\ =\ 1+u+\cdots+u^{n-1},\qquad w_{\mathrm{mat}}^{(n)}=\frac1{D_n},\qquad\langle E\rangle^{(n)}_{\mathrm{mat}}=\frac{2n}{D_n}\in(2,\ 2n]\ }
\]

%dsp-level:0
\[
\qquad [\text{machine: the polynomial identity}\ (1-u)D_n=1-u^n\ \text{for}\ n=1..8;\ \text{the binomial moment algebra exact in}\ \mathbb Q[p]];
\]

%dsp-level:0
\[
\operatorname{Var}^{(n)}_{\mathrm{mat}}(E)=\frac{4n\bigl[(1+(n-1)p)D_n-n\bigr]}{D_n^2};
\]

%dsp-level:0
\[
\qquad n=3:\quad\frac{12\,p(1-p)(3-2p)}{D^2}
\]

%dsp-level:0
\[
\qquad [\mathrm{BM.94}(6)\ \textbf{reproduced exactly};\ \text{machine}];
\]

%dsp-level:0
\[
n=1:\quad w_{\mathrm{mat}}^{(1)}=1:
\]

%dsp-level:0
\[
\qquad \textbf{a single materialized pair has no flicker: the restriction}\ E
\]

%dsp-level:0
\[
\ge 2\ \textbf{forces its crystallized member}
\]

%dsp-level:0
\[
\qquad [\text{machine}];
\]

%dsp-level:0
\[
(2)\quad\textbf{the two-chain ensemble measure (constructed):}
\]

%dsp-level:0
\[
\qquad \text{the meson floor}\ E^{cryst}=2E_{\mathrm{int}}(c)=4\ \ \text{is two pairs}
\]

%dsp-level:0
\[
\qquad [\mathrm{BM.111}(1);\ \text{the digon floor}\ E_{\mathrm{int}}(c)=2\ [\mathrm{FT.28};\ \mathrm{BM.13}(3)]]:
\]

%dsp-level:0
\[
\boxed{\ P_{\mathrm{mat}}(E{=}2)=\frac{2u}{1+u},\qquad P_{\mathrm{mat}}(E{=}4)=\frac{p}{1+u},\qquad\langle E\rangle^{(2)}_{\mathrm{mat}}=\frac4{1+u}\in(2,4],\qquad w_{\mathrm{mat}}^{(2)}=\frac1{1+u}\in\Bigl(\tfrac12,\ 1\Bigr]\ }
\]

%dsp-level:0
\[
\operatorname{Var}^{(2)}_{\mathrm{mat}}(E)=\frac{8p(1-p)}{(2-p)^2};
\]

%dsp-level:0
\[
\qquad \frac{\operatorname{Var}^{(2)}}{\langle E\rangle^{(2)2}}=\frac{p(1-p)}2\le\frac18,
\]

%dsp-level:0
\[
\qquad \frac{\sigma^{(2)}}{\langle E\rangle^{(2)}}\le\frac1{2\sqrt2}<\frac12
\]

%dsp-level:0
\[
\qquad [\text{machine: exact rational functions};\ \tfrac18-\tfrac{p(1-p)}2=\tfrac{(2p-1)^2}{8}\ge0];
\]

%dsp-level:0
\[
\textbf{the two-bit floor survives the averaging:}
\]

%dsp-level:0
\[
\qquad E\ge2\ \ \text{on the whole materialized support}\implies\langle m_{\mathrm{hid}}\rangle^{(2)}_{\mathrm{mat}}>\frac{2\ln2}{\mathfrak h}
\]

%dsp-level:0
\[
\qquad [\mathrm{BM.70}(2);\ \mathrm{BM.94}(8)\ \text{pattern};\ \text{machine}];
\]

%dsp-level:0
\[
\textbf{the no-flicker collapse:}
\]

%dsp-level:0
\[
\qquad \varepsilon_q=0\implies u=0\implies P_{\mathrm{mat}}(E{=}4)=1,
\]

%dsp-level:0
\[
\qquad w_{\mathrm{mat}}^{(2)}=1
\]

%dsp-level:0
\[
\qquad [\mathrm{BM.19}(3);\ \mathrm{Def\text{-}BM.5}(5);\ \text{machine}];
\]

%dsp-level:0
\[
(3)\quad\textbf{the meson's ensemble binding (formed, nonzero — the}\ (\mathrm{R8})\ \textbf{row):}
\]

%dsp-level:0
\[
\qquad \mathcal B^{(n)}_{\mathrm{ens}}\ :=\ \bigl\langle m_{\mathrm{hid}}\bigr\rangle^{(n)}_{\mathrm{mat}}-\sum_{a=1}^n\bigl\langle m_{\mathrm{hid}}\bigr\rangle^{(1)}_{\mathrm{mat}}\ =\ \Bigl(\frac{2n}{D_n}-2n\Bigr)\frac{\ln2}{\mathfrak h}
\]

%dsp-level:0
\[
\qquad [\mathrm{FT.44}(4);\ (1)];
\]

%dsp-level:0
\[
\boxed{\ \mathcal B^{(n)}_{\mathrm{ens}}\ =\ -\,\frac{2n\,u\,D_{n-1}}{D_n}\cdot\frac{\ln2}{\mathfrak h}\ }
\]

%dsp-level:0
\[
\qquad [D_n-1=uD_{n-1};\ \text{machine: the polynomial identity for}\ n=2,3,4];
\]

%dsp-level:0
\[
n=2:
\]

%dsp-level:0
\[
\qquad \boxed{\ \mathcal B^{(2)}_{\mathrm{ens}}=-\frac{4u}{1+u}\cdot\frac{\ln2}{\mathfrak h}=-\frac{16\,\varepsilon_q\ln2}{(1+4\varepsilon_q)\,\mathfrak h}\ },
\]

%dsp-level:0
\[
\qquad \mathcal B^{(2)}_{\mathrm{ens}}<0\ \ \forall\varepsilon_q>0,
\]

%dsp-level:0
\[
\qquad \mathcal B^{(2)}_{\mathrm{ens}}=0\iff\varepsilon_q=0,
\]

%dsp-level:0
\[
\qquad -\frac{2\ln2}{\mathfrak h}<\mathcal B^{(2)}_{\mathrm{ens}}\le0
\]

%dsp-level:0
\[
\qquad [\text{machine: the sign and range census on}\ u\in[0,1)];
\]

%dsp-level:0
\[
\textbf{the binding is the flicker's, not the count's, at}\ n=2\ \textbf{as at}\ n=3
\]

%dsp-level:0
\[
\qquad [\mathrm{BM.21}(3)\ \text{pattern}];
\]

%dsp-level:0
\[
\qquad \text{the two vanishing routes stand untouched:}\ \mathcal B^{\text{meson}}_{\mathrm{info}}=\mathcal B^{\text{meson}}_{\mathrm{pole}}=0
\]

%dsp-level:0
\[
\qquad [\mathrm{BM.111}(1),(2)\ \text{re-cited,}\ \neg\text{retracted}];
\]

%dsp-level:0
\[
(4)\quad\textbf{the sector-mass pincer at budget}\ B=2\ \textbf{(the}\ 2\textbf{-chain sector masses, computed):}
\]

%dsp-level:0
\[
\qquad \mathcal G_2\ \text{carries one chain-bearing skeleton, the digon}\ K_2
\]

%dsp-level:0
\[
\qquad [\mathrm{Def\text{-}H.32}(1);\ \mathrm{FT.105}:\ \text{the connected simple graphs on}\ \le2\ \text{nodes}]:
\]

%dsp-level:0
\[
\textbf{the pincer:}
\]

%dsp-level:0
\[
\qquad \text{a closed chain requires}\ |R|\ge2
\]

%dsp-level:0
\[
\qquad [\mathrm{Def\text{-}H.4}:\ \delta h(c)\ge2]
\]

%dsp-level:0
\[
\qquad \text{and the budget gives}\ |R|\le2E(K_2)=2
\]

%dsp-level:0
\[
\qquad [\mathrm{Def\text{-}H.32}(1)]:
\]

%dsp-level:0
\[
\implies|R|=2E\implies n_{\mathcal V}=2E-|R|=0\implies\Delta S=\tfrac32n_{\mathcal V}=0\qquad[\mathrm{Def\text{-}H.32}(6);\ \mathrm{FT.107}(2)];
\]

%dsp-level:0
\[
\boxed{\ m_2^2(q,S)=0\ \ \text{on every formed sector, and}\ Z_2^{\mathrm{norm}}(q,S)\ne0\ \ \text{only at}\ S=1\ \ \text{with}\ Z_2^{\mathrm{norm}}=(4,\ 2,\ 2)\ \ \text{at}\ q=(0,+1,-1)\ }
\]

%dsp-level:0
\[
\qquad [\text{machine: the exhaustive enumeration of}\ \mathcal G_2\ \text{—}\ 4\ \text{patterns on}\ K_2,\ \text{the}\ \sigma\text{-resolved instance census, the closed form}\ m_q2^{|R|-2}\ \text{of}\ \mathrm{FT.106}(3)\ \text{cross-checked}];
\]

%dsp-level:0
\[
\textbf{the structural reason:}
\]

%dsp-level:0
\[
\qquad \boxed{\ \textbf{at budget two, chain-bearing content is fully crystallized: the sector ensemble has no flicker room}\ };
\]

%dsp-level:0
\[
\qquad \text{the sector-species binding is}\ 0-0-0=0\ \ \text{exactly};
\]

%dsp-level:0
\[
\textbf{the two ensemble reads are separated species:}
\]

%dsp-level:0
\[
\qquad \text{the sector/defect read gives}\ 0;\ \text{the informational read gives}\ (3)\text{'s strict negative};
\]

%dsp-level:0
\[
\qquad \neg\text{a derived equation relates them}\ —\ \text{that is}\ (\mathrm{R1})\text{'s bridge}
\]

%dsp-level:0
\[
\qquad [\mathrm{BM.37}(2);\ \mathrm{BM.98};\ \mathrm{A\text{-}BM.9}];
\]

%dsp-level:0
\[
(5)\quad\textbf{the two-chain spectrum (the resonance rows at}\ n=2\textbf{):}
\]

%dsp-level:0
\[
\qquad W_2(Q)=(1,4,6,4,1),
\]

%dsp-level:0
\[
\qquad \sum_QW_2=16=4^2
\]

%dsp-level:0
\[
\qquad [\mathrm{BM.114}(1);\ \text{machine re-run}];
\]

%dsp-level:0
\[
\lambda^{(2)}(s_1,s_2)=D_S(s_1)+D_S(s_2),
\]

%dsp-level:0
\[
\qquad D_S\in\bigl\{2,\ 4,\ 2-\sqrt2,\ 1,\ 2-\mathfrak u,\ 2-\sqrt3\bigr\},
\]

%dsp-level:0
\[
\qquad \mathfrak u=2\cos\tfrac\pi{10}
\]

%dsp-level:0
\[
\qquad [\mathrm{Def\text{-}BM.14}(3);\ \mathrm{BM.111}(2);\ \mathrm{BM.114}(3)]:
\]

%dsp-level:0
\[
\boxed{\ \text{the}\ 21\ \text{multiset orbits give}\ 21\ \text{distinct eigenvalues with}\ \neg\text{a collision}\ }
\]

%dsp-level:0
\[
\qquad [\text{machine: the exact equality census in}\ \mathbb M;\ \text{contrast: one collision at}\ n=3\ [\mathrm{BM.114}(5)]];
\]

%dsp-level:0
\[
\lambda^{(2)}_{\min}=4-2\mathfrak u=2\bigl(2-2\cos\tfrac\pi{10}\bigr),
\]

%dsp-level:0
\[
\qquad \lambda^{(2)}_{\max}=8=4n
\]

%dsp-level:0
\[
\qquad [\mathrm{BM.117}(4)\ \text{re-certified at}\ n=2;\ \text{machine: the exact extrema in}\ \mathbb M];
\]

%dsp-level:0
\[
m^2_{\mathrm{pole}}(\mathcal R)=\hat\mu'^2\lambda^{(2)}\in\bigl[(4-2\mathfrak u)\,\hat\mu'^2,\ 8\hat\mu'^2\bigr];
\]

%dsp-level:0
\[
\qquad m^{meson}_{\mathrm{hid}}=\frac{2E_{\mathrm{int}}(c)\ln2}{\mathfrak h};
\]

%dsp-level:0
\[
\qquad \bigl\langle m_{\mathrm{hid}}\bigr\rangle^{(2)}_{\mathrm{mat}}=\frac{4\ln2}{(1+u)\,\mathfrak h}\ \ \text{at the digon floor}
\]

%dsp-level:0
\[
\qquad [\mathrm{Def\text{-}BM.17}(2);\ (2);\ \text{machine}];
\]

%dsp-level:0
\[
\varepsilon_M=+1\ \ \text{unconditionally},\qquad j_M=k\in\mathbb Z\qquad[\mathrm{BM.129}(1),(2)\ \text{re-cited}];
\]

%dsp-level:0
\[
(6)\quad\textbf{the octet adjudication stands, and the spectrum's shape is the derived one:}
\]

%dsp-level:0
\[
\qquad 8\ \ \text{occurs in no grouping family at any formed level}
\]

%dsp-level:0
\[
\qquad [\mathrm{BM.117}(3);\ \text{machine: the scan re-run over}\ W_n,\ g(N),\ \binom{6+n-1}n,\ \binom6n,\ n\le3,\ N\le7];
\]

%dsp-level:0
\[
\textbf{the meson spectrum is the eigenvalue ladder}\ (5)\ \textbf{with its charge census, its statistics and its formed binding}\ —\ \neg\text{an octet,}\ \neg\text{a nonet},
\]

%dsp-level:0
\[
\ \neg\text{a flavour multiplet};
\]

%dsp-level:0
\[
\textbf{the naming row:}
\]

%dsp-level:0
\[
\qquad \text{species names for meson classes remain confined to the identification layer}\ [\mathrm{Def\text{-}BM.3}(7);\ \mathrm{Def\text{-}BM.6}(7)];
\]

%dsp-level:0
\[
\qquad \neg\text{a measured mass, width or mixing enters}\ [\mathrm{P5}];
\]

%dsp-level:0
\[
(7)\quad\textbf{status:}
\]

%dsp-level:0
\[
\qquad \textbf{derived:}\ (1)\text{--}(6);
\]

%dsp-level:0
\[
\qquad \boxed{\ (\mathrm{R8})\ \textbf{is closed: the}\ 2\textbf{-chain sector masses are computed (identically zero, with the pincer), the two-chain ensemble measure is constructed, and the meson's ensemble binding is formed and strictly negative at every nonzero threshold}\ }
\]

%dsp-level:0
\[
[\mathrm{BM.111}(3)\text{'s}\ \neg\text{formed row discharged}];
\]

%dsp-level:0
\[
\qquad \neg\text{a declaration},
\]

%dsp-level:0
\[
\qquad \neg\text{a P4 addition},
\]

%dsp-level:0
\[
\qquad \neg\text{a new dimensionful scale}\ [\mathrm{P1}],
\]

%dsp-level:0
\[
\qquad \neg\text{a measured value}\ [\mathrm{P5}].
\]

\textbf{Req.} A4 (Part 1); D0.4, D0.8 (Part 1); CT0.8, CT1.29 (Part 1); CT2.8 (Part 2); P1, P5 (Parts 1–2); Def-H.4, Def-H.32, FT.28, FT.32(6), FT.44(4), FT.105, FT.106(3), FT.107(1), FT.107(2), FT.144(4) (Part 3); Def-BM.5, Def-BM.14(3), Def-BM.17(2), BM.13(3), BM.19(2), BM.19(3), BM.21(3), BM.37(2), BM.70(2), BM.94(4), BM.94(5), BM.94(6), BM.94(8), BM.98, BM.111, BM.114(1), BM.114(3), BM.114(5), BM.117(3), BM.117(4), BM.129(1), BM.129(2) (this part); Census FT.35(2) (row-wise); the A-BM.9 declaration; binomial moment algebra over $\mathbb Q[p]$, the degree-16 field $\mathbb M$ with its exact equality census, exhaustive configuration enumeration at budget two, polynomial identities [math; machine: \texttt{checks/\allowbreak{}verify\_\allowbreak{}H\_\allowbreak{}bm\_\allowbreak{}closing\_\allowbreak{}tier.py} §F].

\textbf{Pf.}

%dsp-level:0
\[
(1):\ P(E{=}2j)\ \text{is the}\ n\text{-pair binomial}\ [\mathrm{BM.94}(4)\text{'s specialization at}\ n];
\]

%dsp-level:0
\[
\qquad \text{conditioning on}\ j\ge1:\ \ 1-u^n=pD_n;
\]

%dsp-level:0
\[
\qquad \langle E\rangle^{(n)}_{\mathrm{mat}}=2\cdot\frac{np}{1-u^n}=\frac{2n}{D_n}\ [\text{machine: exact in}\ \mathbb Q[p]];
\]

%dsp-level:0
\[
\langle E^2\rangle^{(n)}_{\mathrm{mat}}=4\cdot\frac{np(1+(n-1)p)}{1-u^n}=\frac{4n(1+(n-1)p)}{pD_n}\ [\text{machine}];
\]

%dsp-level:0
\[
\qquad \operatorname{Var}=\langle E^2\rangle-\langle E\rangle^2\ [\text{machine}];
\]

%dsp-level:0
\[
\qquad n=3\ \text{row against}\ \mathrm{BM.94}(6)\text{'s frozen expression}\ [\text{machine: identical rational functions}];
\]

%dsp-level:0
\[
\qquad n=1:\ D_1=1,\ P(E{=}2)=1\ [\text{machine}]
\]

%dsp-level:0
\[
(2):\ n=2\ \text{substitution:}\ D_2=1+u=2-p;
\]

%dsp-level:0
\[
\qquad P_{\mathrm{mat}}(2)=\frac{2pu}{pD_2}=\frac{2u}{1+u},\quad P_{\mathrm{mat}}(4)=\frac{p^2}{pD_2}=\frac p{1+u}\ [\text{machine}];
\]

%dsp-level:0
\[
\qquad \text{the variance and the relative-width SOS certificate}\ [\text{machine}]
\]

%dsp-level:0
\[
(3):\ \mathcal B^{(n)}
\]

%dsp-level:0
\[
= \bigl(\frac{2n}{D_n}-2n\bigr)\frac{\ln2}{\mathfrak h}
\]

%dsp-level:0
\[
= -2n\frac{D_n-1}{D_n}\frac{\ln2}{\mathfrak h}
\]

%dsp-level:0
\[
= -\frac{2nuD_{n-1}}{D_n}\frac{\ln2}{\mathfrak h}\ [D_n-1=uD_{n-1};\ \text{machine}];
\]

%dsp-level:0
\[
\qquad \text{the sign, zero and range rows on}\ u\in[0,1)\ [\text{machine: rational grid and the factor census}]
\]

%dsp-level:0
\[
(4):\ \mathcal G_2=\{\bullet,\ K_2\}\ [\mathrm{FT.105}\text{'s census at}\ V\le2];
\]

%dsp-level:0
\[
\qquad \text{the}\ 4\ \text{patterns of}\ K_2\ \text{enumerated:}\ |R|\in\{0,1,1,2\};
\]

%dsp-level:0
\[
\qquad \text{closed chains exist only at}\ |R|=2\ (\text{the digon},\ k=1,\ S=1,\ 2\ \text{bases})\ [\text{machine}];
\]

%dsp-level:0
\[
\qquad \Delta S=\tfrac32(2-2)=0;
\]

%dsp-level:0
\[
\qquad Z_2^{\mathrm{norm}}(q,1)
\]

%dsp-level:0
\[
= 2\ \text{instances}\times m_q2^{0}
\]

%dsp-level:0
\[
= (4,2,2)\ [\mathrm{FT.106}(3)\text{'s closed form};\ \text{machine: the}\ \sigma\text{-resolved census agrees}];
\]

%dsp-level:0
\[
\qquad S\ne1:\ \neg\text{instances}\implies Z=0
\]

%dsp-level:0
\[
(5):\ \text{the}\ 21\ \text{sums computed in}\ \mathbb M\ \text{with the basis-equality census}\ [\text{machine}];
\]

%dsp-level:0
\[
\qquad \text{the extrema by the exact order of}\ \mathbb M\text{'s positive cone on the six}\ D_S\ [\text{machine: dyadic interval embedding}];
\]

%dsp-level:0
\[
\qquad (6):\ \text{the grouping scan re-run}\ [\text{machine}];
\]

%dsp-level:0
\[
\qquad (7):\ \text{assembly}.\ \blacksquare
\]

\textbf{Phys.} The meson gets its ensemble. The flicker measure generalizes from the marked three-pair specialization to a law at every $n$: the materialized weight is one over the geometric sum $1+u+\cdots+u^{n-1}$, the mean count is $2n$ over that sum, and at two chains the measure is exact — a mean between two and four, a relative width bounded by one over two root two, and the two-bit floor intact. The binding that residue (R8) named is formed and nonzero: strictly negative at every nonzero threshold, zero exactly at the sharp point, the same flicker-not-count structure the baryon showed. The sector route tells the complementary story: at budget two a chain-bearing configuration has no virtual room at all, so the two-chain sector masses vanish identically — a pincer, not a coincidence — and the two ensemble reads stay separated currencies pending the bridge slot. The spectrum completes the row: five charge classes with weights one-four-six-four-one, twenty-one pole eigenvalues with no collision at all, bounded by four-$n$, bosonic unconditionally. The octet remains unformed and the names remain where the identification layer keeps them.

\end{theorem}

\begin{theorem}[{\texorpdfstring{BM.138 — The seed-tensoring mechanism: the symmetrizer as a finite-group twirl, the transport and composition closure, and the two residues adjudicated}{BM.138 - The seed-tensoring mechanism: the symmetrizer as a finite-group twirl, the transport and composition closure, and the two residues adjudicated}}]\label{tag:BM-138}
\textbf{Thm.}

%dsp-level:0
\[
(1)\quad\textbf{the mechanism, derived:}
\]

%dsp-level:0
\[
\qquad \mathcal H_{\mathsf v}\cong\mathbb C^2\ \ \text{the transported seed}
\]

%dsp-level:0
\[
\qquad [\mathrm{EW.32};\ \mathrm{A2}:\ \text{the channel capacity is one qubit}];
\]

%dsp-level:0
\[
\boxed{\ \Psi(c)\ =\ \mathrm{Sym}^{|k(c)|}\bigl(\mathcal H_{\mathsf v}\bigr)\ =\ \Pi_{|k|}\Bigl(\mathcal H_{\mathsf v}^{\otimes|k|}\Bigr),\qquad\Pi_k:=\frac1{k!}\sum_{\sigma\in S_k}\hat P_\sigma\ }
\]

%dsp-level:0
\[
\qquad [\mathrm{Def\text{-}H.38};\ \mathrm{BM.128}(2)\text{'s permutation operators on identical factors}];
\]

%dsp-level:0
\[
\Pi_k^2=\Pi_k,
\]

%dsp-level:0
\[
\qquad \operatorname{rk}\Pi_k=k+1=\dim\mathrm{Sym}^k(\mathbb C^2),
\]

%dsp-level:0
\[
\qquad \Pi_k\hat P_\tau=\hat P_\tau\Pi_k=\Pi_k\ \ \forall\tau\in S_k
\]

%dsp-level:0
\[
\qquad [\text{machine: idempotence, exact rank over}\ \mathbb Q,\ \text{and the commutation for}\ k=1..6];
\]

%dsp-level:0
\[
\boxed{\ \textbf{the symmetrizer is the Haar twirl of the finite group}\ S_k\ }
\]

%dsp-level:0
\[
\qquad [\mathrm{BM.131}(6);\ \text{math: the averaging over a finite group is its Haar measure}];
\]

%dsp-level:0
\[
\textbf{the seed-tensoring read:}
\]

%dsp-level:0
\[
\qquad \text{the carrier is the}\ k\text{-fold tensor power of the seed, projected onto the invariants of the group average}\ —\ \textbf{tensoring with itself, then twirling};
\]

%dsp-level:0
\[
(2)\quad\textbf{the diagonal transport:}\qquad\mathcal U_c^{\otimes k}\big|_{\mathrm{Sym}^k}\qquad[\mathrm{Def\text{-}H.38};\ \mathrm{EW.33}:\ \text{the wrapping transport}];
\]

%dsp-level:0
\[
\bigl[\mathcal U^{\otimes k},\ \Pi_k\bigr]=0
\]

%dsp-level:0
\[
\qquad [\text{machine: the infinitesimal rows}\ \bigl[\textstyle\sum_a W_a^{(i)},\Pi_k\bigr]=0\ \text{on the}\ \mathfrak{su}(2)\ \text{generator triple, exact over}\ \mathbb Q(i),\ k\le6];
\]

%dsp-level:0
\[
\textbf{the weight ladder from the twirl:}
\]

%dsp-level:0
\[
\qquad \operatorname{Tr}\Bigl(\Pi_k\,D(z)^{\otimes k}\Bigr)=\sum_{m\in\{k,k-2,\dots,-k\}}z^{m},
\]

%dsp-level:0
\[
\qquad D(z)=\mathrm{diag}(z,z^{-1})
\]

%dsp-level:0
\[
\qquad [\mathrm{FT.60}(4)\ \text{re-derived};\ \text{machine: the Laurent identity at five rational}\ z\ \text{for}\ k\le8];
\]

%dsp-level:0
\[
\text{the Casimir row}\ \tfrac k2\bigl(\tfrac k2+1\bigr)\qquad[\mathrm{FT.60}(3)\ \text{re-cited}];\qquad j=\frac{|k|}2\qquad[\mathrm{Def\text{-}H.38};\ \mathrm{FT.55}(2)];
\]

%dsp-level:0
\[
(3)\quad\textbf{the composition closure (the junctions tensor by the same mechanism):}
\]

%dsp-level:0
\[
\qquad \mathrm{Sym}^{k_1}\otimes\mathrm{Sym}^{k_2}\supseteq\mathrm{Sym}^{k_1+k_2}\ \ \text{as the maximal-weight line on the cone}
\]

%dsp-level:0
\[
\qquad [\mathrm{FT.60}(6);\ \text{machine: the top weight}\ z^{k_1+k_2}\ \text{occurs with multiplicity}\ 1\ \text{in the character product, Laurent-exact}];
\]

%dsp-level:0
\[
\Psi(\Gamma)=\mathrm{Sym}^{k_1+k_2+d_{12}}\subseteq\mathrm{Sym}^{k_1}\otimes\mathrm{Sym}^{k_2}\otimes\mathrm{Sym}^{d_{12}}\qquad[\mathrm{BM.2}(2);\ \mathrm{FT.34}(3):\ k(\Gamma)=k_1+k_2+d_{12}];
\]

%dsp-level:0
\[
\boxed{\ \textbf{the connector is a carrier too:}\ \ \mathrm{Sym}^{d_{12}}(\mathcal H_{\mathsf v})\ \ \text{is built by the same seed-tensoring, so every junction composition closes inside the mechanism}\ }
\]

%dsp-level:0
\[
(4)\quad\textbf{the defect-ladder coherence:}
\]

%dsp-level:0
\[
\qquad \text{the weights}\ m\in\{k,k{-}2,\dots,-k\}\ \text{of}\ (2)\ \text{are the ladder}\ \Delta(m)
\]

%dsp-level:0
\[
= 4\sin^2\frac{m\Theta_R}4\ \text{reads on the tensorized carrier's own weight basis}
\]

%dsp-level:0
\[
\qquad [\mathrm{Def\text{-}BM.19}(3);\ \mathrm{BM.124}(3);\ \mathrm{Def\text{-}H.37}(2)]:
\]

%dsp-level:0
\[
\text{the mechanism is what makes the carrier and the defect read the same multiplet}\ —\ \neg\text{a new coupling,}\ \neg\text{a proportionality}\ \Theta_R
\]

%dsp-level:0
\[
= 2\pi k\ [\mathrm{BM.124}(5)\text{'s dimensional bar re-cited}];
\]

%dsp-level:0
\[
(5)\quad\textbf{residue (i) adjudicated — the coupling datum:}
\]

%dsp-level:0
\[
\qquad \operatorname{copies}(c)=k(c)\ \ \text{is}\ \mathrm{Def\text{-}H.38}\text{'s constitutive clause};
\]

%dsp-level:0
\[
\textbf{the derivation route is blocked by a frozen minimality theorem:}
\]

%dsp-level:0
\[
\qquad \text{the uniform-minimal triangle's winding}\ 3\ \text{has no three materialized winding-}1\ \text{subloops, and admissible decompositions are basis-relative}
\]

%dsp-level:0
\[
\qquad [\mathrm{FT.60}(7\mathrm a)\text{--}(7\mathrm b)];
\]

%dsp-level:0
\[
\boxed{\ \textbf{the mechanism is derived and the coupling integer stays a priced interface:}\ \ \text{what remained was the construction half, and it is complete}\ };
\]

%dsp-level:0
\[
(6)\quad\textbf{residue (ii) adjudicated — the}\ k\ge3\ \textbf{carrier inventory:}
\]

%dsp-level:0
\[
\qquad \mathrm{Sym}^{k\ge3}\ (\dim\ge4)\ \text{has no registered aspect block}\ [\mathrm{CT0.12}:\ \text{blocks}\ (3,3,5,3,1,1)];
\]

%dsp-level:0
\[
\textbf{the adjudication:}
\]

%dsp-level:0
\[
\qquad \text{the carrier is a chain-layer object,}\ \neg\text{a force aspect:}\ \ \text{every clause of this part consumes}\ \Psi\ \text{through}\ \bigl(\dim,\ \text{weights},\ \varepsilon,\ j,\ \text{Casimir},\ \chi\bigr)\ \text{alone}
\]

%dsp-level:0
\[
\qquad [\text{machine: the inventory scan of}\ \mathrm{Def\text{-}H.38}\text{'s consumption sites in Part 4}];
\]

%dsp-level:0
\[
\neg\text{an aspect claim at}\ k\ge3\ \ \text{is made or needed}\ [\mathrm{Def\text{-}H.38}\ \text{Scope re-cited}];
\]

%dsp-level:0
\[
\qquad \textbf{the inventory-extension question is void for this era};
\]

%dsp-level:0
\[
(7)\quad\textbf{status:}
\]

%dsp-level:0
\[
\qquad \textbf{derived:}\ (1)\text{--}(4);
\]

%dsp-level:0
\[
\qquad \textbf{adjudicated:}\ (5),(6);
\]

%dsp-level:0
\[
\qquad \boxed{\ \textbf{the seed-tensoring mechanism is derived}\ };
\]

%dsp-level:0
\[
\qquad \neg\text{a declaration},
\]

%dsp-level:0
\[
\qquad \neg\text{a P4 addition},
\]

%dsp-level:0
\[
\qquad \neg\text{a measured value}\ [\mathrm{P5}].
\]

\textbf{Req.} A2 (Part 1); D1.4 (Part 1); CT0.12 (Part 1); P5 (Parts 1–2); EW.32, EW.33 (Part 2); Def-H.37(2), Def-H.38, FT.34(3), FT.55(2), FT.60(3), FT.60(4), FT.60(6), FT.60(7) (Part 3); Def-BM.19(3), BM.2(2), BM.124(3), BM.124(5), BM.128(2), BM.131(6) (this part); exact rank by Fraction elimination, Laurent-polynomial character identities, the $\mathfrak{su}(2)$ infinitesimal commutation over $\mathbb Q(i)$, finite-group averaging [math; machine: \texttt{checks/\allowbreak{}verify\_\allowbreak{}H\_\allowbreak{}bm\_\allowbreak{}closing\_\allowbreak{}tier.py} §G].

\textbf{Pf.}

%dsp-level:0
\[
(1):\ \Pi_k^2
\]

%dsp-level:0
\[
= \Pi_k\ \text{from}\ \hat P_\sigma\hat P_\tau
\]

%dsp-level:0
\[
= \hat P_{\sigma\tau}\ \text{and the group-average rearrangement}\ [\mathrm{BM.128}(2)\ \text{transported to identical factors};\ \text{machine}];
\]

%dsp-level:0
\[
\qquad \operatorname{rk}\Pi_k:\ \text{exact elimination on the}\ 2^k\text{-dimensional symmetrizer for}\ k
\]

%dsp-level:0
\[
= 1..6:\ \text{ranks}\ (2,3,4,5,6,7)\ [\text{machine}]
\]

%dsp-level:0
\[
(2):\ \bigl[\hat P_\sigma,\mathcal U^{\otimes k}\bigr]
\]

%dsp-level:0
\[
= 0\ \text{for the diagonal action}\ [\text{math: the permutation and the diagonal tensor power commute}];
\]

%dsp-level:0
\[
\qquad \text{the infinitesimal rows on the generator triple}\ W_a
\]

%dsp-level:0
\[
= \tfrac12\sigma_a:\ \ \bigl[\textstyle\sum_i I^{\otimes(i-1)}\otimes W_a\otimes I^{\otimes(k-i)},\ \Pi_k\bigr]
\]

%dsp-level:0
\[
= 0\ [\text{machine over}\ \mathbb Q(i)];
\]

%dsp-level:0
\[
\qquad \text{the character:}\ \operatorname{Tr}(\Pi_kD^{\otimes k})
\]

%dsp-level:0
\[
= \frac1{k!}\sum_\sigma\prod_{\text{cycles}}(z^{L}+z^{-L})
\]

%dsp-level:0
\[
= \sum_mz^m\ [\text{machine: five rational}\ z\ \text{values,}\ k\le8]
\]

%dsp-level:0
\[
(3):\ \text{the character product}\ \bigl(\sum_iz^{k_1-2i}\bigr)\bigl(\sum_jz^{k_2-2j}\bigr)\ \text{carries}\ z^{k_1+k_2}\ \text{with coefficient}\ 1\ [\text{machine: Laurent multiplication}];
\]

%dsp-level:0
\[
\qquad \text{the connector row:}\ \mathrm{Sym}^{d_{12}}\ \text{is}\ (1)\text{'s construction at}\ k=d_{12}\ [\mathrm{BM.2}(2)]
\]

%dsp-level:0
\[
(4):\ \text{the ladder reads on the weight basis}\ \{\psi_m\}\ \text{of}\ (2)\ [\mathrm{Def\text{-}BM.19}(3);\ \mathrm{BM.124}(3)\ \text{re-cited}];
\]

%dsp-level:0
\[
\qquad (5):\ \mathrm{FT.60}(7)\text{'s two obstruction rows re-cited}\ —\ \text{the irreducibility is frozen};
\]

%dsp-level:0
\[
\qquad (6):\ \text{the consumption-site scan}\ [\text{machine: every}\ \Psi\text{-use in Part 4 is through the six listed reads}];
\]

%dsp-level:0
\[
\qquad (7):\ \text{assembly}.\ \blacksquare
\]

\textbf{Phys.} The carrier's construction is no longer a black box with a declared integer on top: the mechanism underneath the integer is derived. The symmetric power is the image of the group average over the $k$-fold tensor power of the seed — the same twirl structure the Haar–Poisson duality runs on $\mathrm{SU}(2)$, here on the finite permutation group — with rank $k+1$ exact, diagonal transport commuting with the average, and the weight ladder falling out of a Laurent identity. Junctions compose by the same machine: the connector itself carries a symmetric-power carrier, so every glued chain stays inside the mechanism. What remains declared is exactly one integer per chain, and a frozen minimality theorem says no construction below it exists; what remains questioned — an aspect block for carriers of dimension four and up — is adjudicated void for this era, because no clause consumes a carrier except through its dimension, weights, sign, spin, Casimir and character.

\end{theorem}

\begin{theorem}[{\texorpdfstring{BM.139 — The two convention declarations and the sector locus}{BM.139 - The two convention declarations and the sector locus}}]\label{tag:BM-139}
\textbf{Thm.}

%dsp-level:0
\[
(1)\quad\textbf{the}\ \aleph\textbf{-domain conflict (exactly as computed):}
\]

%dsp-level:0
\[
\qquad \text{three domains for the vertex-amplitude exponent}\ \aleph\ \text{entering}\ w=\mathcal J\prod_{e\in R}G(e)(-1)^{\aleph}
\]

%dsp-level:0
\[
\qquad [\mathrm{Def\text{-}H.27};\ \mathrm{A8};\ \mathrm{D3.3}]:
\]

%dsp-level:0
\[
\text{literal}\ (\mathrm{S2},\ \text{frozen});\qquad\mathrm{E1}\ (\text{chain-dominant});\qquad\mathrm{E2}\ (\text{literal-forced})\qquad[\mathrm{Def\text{-}H.27}(\mathrm a)\text{--}(\mathrm c)];
\]

%dsp-level:0
\[
\textbf{pairwise inequivalent, with the exact witnesses:}
\]

%dsp-level:0
\[
\qquad \text{the locked single-vertex star}\ (n_{\mathrm{in}}{+}n_{\mathrm{out}}=n\in\{4,5,6\},\ \text{balanced}):
\]

%dsp-level:0
\[
\qquad \text{literal}\ 0;
\]

%dsp-level:0
\[
\qquad \mathrm{E2}=\mathrm{E1}=\sum_{a\ge2}\binom{n_{\mathrm{in}}}a\binom{n_{\mathrm{out}}}aa!\,4^a=(32,\ 96,\ 672)
\]

%dsp-level:0
\[
\qquad [\mathrm{FT.75}(1);\ \text{machine: the binomial sums re-derived exactly}];
\]

%dsp-level:0
\[
\text{the locked digon hub with}\ a\ \text{digons}:
\]

%dsp-level:0
\[
\qquad \text{literal}\ 4^a;
\]

%dsp-level:0
\[
\qquad \mathrm{E2}=0;
\]

%dsp-level:0
\[
\qquad \mathrm{E1}=(!a)\,4^a\ \ \text{with}\ (!a)=(1,2,9,44)\ \ \text{at}\ a=2..5
\]

%dsp-level:0
\[
\qquad [\mathrm{FT.75}(1);\ \text{machine: the derangement census}];
\]

%dsp-level:0
\[
\boxed{\ \text{no further computation settles the choice from the frozen corpus: it is a declaration species}\ }\qquad;
\]

%dsp-level:0
\[
(2)\quad\textbf{A-BM.10 issued (the}\ \aleph\textbf{-domain):}
\]

%dsp-level:0
\[
\qquad \textbf{the baryomesonic era runs:}
\]

%dsp-level:0
\[
\qquad (\mathrm a)\ \ \text{the topology-blind ensembles on}\ \mathrm{Def\text{-}H.32}\text{'s inherited literal domain}\ (\mathrm{S2})\ —\ \text{unchanged, with every frozen ensemble row standing in its own domain}\ [\mathrm{Def\text{-}H.32}\ \text{Scopes};\ \mathrm{FT.74}\ \text{status row}];
\]

%dsp-level:0
\[
(\mathrm b)\ \ \text{the sealed-junction amplitude reads}\ —\ \text{where the chain-local pass-through usage of the telescoping and EM rows is the forced datum}\ [\mathrm{FT.21}(\mathrm i);\ \mathrm{FT.62}(5);\ \mathrm{BM.7}\ \text{species}]\ —\ \text{on the}\ \mathrm{E2}\ \text{domain}\ [\mathrm{Def\text{-}H.27}(\mathrm c)];
\]

%dsp-level:0
\[
(\mathrm c)\ \ \ \mathrm{E1}\ \textbf{recorded as a computed variant,}\ \neg\textbf{adopted};
\]

%dsp-level:0
\[
\textbf{the grounds:}
\]

%dsp-level:0
\[
\qquad \mathrm{E2}\ \text{retains the literal}\ \mathrm{D3.3}\ \text{read at every type-}(1,1)\ \text{vertex, so the frozen literal-domain records carry over unmodified there};
\]

%dsp-level:0
\[
\qquad \text{the chain choice at type-}(1,1)\ \text{is}\ \aleph\text{-gauge}\ [\mathrm{Def\text{-}H.27}(\mathrm c):\ \text{the datum is type-determined}]\ \text{and}\ \mathrm{E2}\ \text{fixes that gauge by vertex type, while}\ \mathrm{E1}\ \text{leaves it to the decomposition};
\]

%dsp-level:0
\[
\textbf{the derived consequences:}
\]

%dsp-level:0
\[
\qquad \text{the locked star is revived with the closed forms}\ (32,96,672);
\]

%dsp-level:0
\[
\qquad \text{the locked digon hub cancels in}\ \mathrm{E2}\ —\ \text{the bound-content cancellation at hubs is preserved, sharpening}\ \mathrm{FT.67}(3\mathrm a);
\]

%dsp-level:0
\[
\qquad \kappa\text{-invariance and the}\ (-1)^{\sum\kappa}\ \text{scaling carry over verbatim}\ [\mathrm{Def\text{-}H.27}\ \text{Rmk}];
\]

%dsp-level:0
\[
(3)\quad\textbf{the corner-convention conflict:}
\]

%dsp-level:0
\[
\qquad \text{the four}\ \Theta_R\text{-readings are computed:}
\]

%dsp-level:0
\[
\qquad \text{the vertex figure}\ \delta_v\ (\text{the only non-trivial spectrum});
\]

%dsp-level:0
\[
\qquad \text{the flat face}\ 0\ \text{and the global sweep}\ V\delta_v=4\pi\ (\text{both identically massless});
\]

%dsp-level:0
\[
\qquad \text{the winding pairing inside the first:}\ \text{variant A}\ (\delta_v,k_f),\ \text{variant B}\ (\delta_v,k_{vf})
\]

%dsp-level:0
\[
\qquad [\mathrm{FT.76}(1);\ \mathrm{Census\ FT.76}(1)\text{--}(2);\ \text{machine:}\ 2-2T_m(1)=0\ \forall m\ \text{in}\ \mathbb Z[x]\ \text{for the two massless readings}];
\]

%dsp-level:0
\[
k_f=(3,2,3,5,3),
\]

%dsp-level:0
\[
\qquad k_{vf}=(3,3,2,3,5)
\]

%dsp-level:0
\[
\qquad [\mathrm{FT.76}(1);\ \text{machine: both re-derived from the minimality rule}\ k=L\kappa_{\min}/2\ \text{on the face and vertex-figure loop lengths}\ L_f=(3,4,3,5,3),\ L_{vf}=(3,3,4,3,5)];
\]

%dsp-level:0
\[
\textbf{the variants are materially different:}
\]

%dsp-level:0
\[
\qquad \text{the multiplet tables differ}\ [\mathrm{Census\ FT.76}(2)];
\]

%dsp-level:0
\[
\qquad \textbf{the per-class statistics censuses differ:}\ \ 112/104\ \ \text{in B against}\ 140/76\ \ \text{in A}
\]

%dsp-level:0
\[
\qquad [\mathrm{BM.136}(4);\ \text{machine}];
\]

%dsp-level:0
\[
(4)\quad\textbf{A-BM.11 issued (the defect corner convention):}
\]

%dsp-level:0
\[
\qquad \boxed{\ \textbf{variant B — the vertex-figure winding}\ k_{vf}\ \text{with}\ \Theta_R=\delta_v\ }
\]

%dsp-level:0
\[
\qquad [\mathrm{BM.6}(4)\text{'s specification discharged}];
\]

%dsp-level:0
\[
\textbf{the grounds:}
\]

%dsp-level:0
\[
\qquad (\mathrm i)\ \ \text{loop consistency: the}\ \Theta_R
\]

%dsp-level:0
\[
= \delta_v\ \text{read is the vertex-figure loop's own defect, and pairing it with the face winding would mix two loci}\ [\mathrm{FT.76}(1):\ \text{variant B is the loop-consistent reading}];
\]

%dsp-level:0
\[
(\mathrm{ii})\ \ \text{the whole defect ladder of this part runs at the vertex-figure locus}\ [\mathrm{Def\text{-}BM.14}(3);\ \mathrm{Def\text{-}BM.19}(3);\ \mathrm{BM.79}(1);\ \mathrm{BM.124}(3);\ \mathrm{BM.125}];
\]

%dsp-level:0
\[
\qquad (\mathrm{iii})\ \ \text{the ceiling}\ \sqrt{4/D_S}\ \text{is winding-independent and untouched}\ [\mathrm{FT.76}(3)];
\]

%dsp-level:0
\[
\textbf{the effects:}
\]

%dsp-level:0
\[
\qquad \mathrm{FT.76},
\]

%dsp-level:0
\[
\ \mathrm{FT.77}\text{'s conditional clauses and this part's ladders convert from convention-tagged to declared data at their points of use},
\]

%dsp-level:0
\[
\ \neg\text{retracted};
\]

%dsp-level:0
\[
\qquad \text{variant A's tables remain computed data, recorded and}\ \neg\text{adopted};
\]

%dsp-level:0
\[
\qquad \text{the winding assignment per class is pinned}\ [\mathrm{BM.136}(4)];
\]

%dsp-level:0
\[
(5)\quad\textbf{A-BM.13 issued (the statistics locus):}
\]

%dsp-level:0
\[
\qquad \boxed{\ \textbf{the}\ \vee^2\textbf{-sector is the standing locus of the sign read}\ }
\]

%dsp-level:0
\[
\qquad [\mathrm{Def\text{-}BM.16}(5)\text{'s non-derived clause converted to declared data}];
\]

%dsp-level:0
\[
\textbf{the grounds (the four priced rows):}
\]

%dsp-level:0
\[
\qquad \vee^2\ \text{and}\ \wedge^2\ \text{give one statistics up to a global sign};
\]

%dsp-level:0
\[
\ \text{the determinant line and the orientation block do not}\ [\mathrm{BM.113}(6);\ \text{machine re-run to}\ k=8];
\]

%dsp-level:0
\[
\vee^2\ \text{is the flip's}\ +1\text{-eigensector on the CP pair}\ [\mathrm{A3}]\ \text{and matches the carrier's own symmetry}\ —\ \wedge^3(\mathbb C^2)
\]

%dsp-level:0
\[
= 0\ \text{leaves the carrier no antisymmetric locus at}\ n\qquad = 3
\]

%dsp-level:0
\[
\qquad [\mathrm{BM.128}(5)];
\]

%dsp-level:0
\[
\textbf{A-BM.12 (the resolvent interface) is issued at}\ \mathrm{Def\text{-}BM.20}\ \text{and recorded here in the ledger census};
\]

%dsp-level:0
\[
(6)\quad\textbf{the convention question closed — the census:}
\]

%dsp-level:0
\[
\qquad (\mathrm d)\ \ \text{the}\ \aleph\text{-domain: declared}\ [(2)];
\]

%dsp-level:0
\[
\qquad (\mathrm e)\ \ \text{the corner convention: declared}\ [(4)];
\]

%dsp-level:0
\[
(\mathrm b)\ \ \text{the exchange block}\ \mathcal K^{(\mathrm{exch})}:\ \textbf{adjudicated void for this era}\ —\ \text{the chirality read has its locus in the carrier rays}\ [\mathrm{A\text{-}BM.7};\ \mathrm{BM.73}],
\]

%dsp-level:0
\[
\ \text{and no}\ k
\]

%dsp-level:0
\[
\ge 3\ \text{aspect block is consumed}\ [\mathrm{BM.138}(6)];
\]

%dsp-level:0
\[
\qquad \neg\text{a declared exchange block is introduced};
\]

%dsp-level:0
\[
(\mathrm a),
\]

%dsp-level:0
\[
(\mathrm c)\ \ \text{the classification program's remaining cells}\ (\text{multi-solid chains, the single-vertex ladder, separated-contact hosting maps})\ \textbf{are named for the Part-5 cell program}\ —\ \neg\text{era objectives};
\]

%dsp-level:0
\[
\boxed{\ \textbf{both conventions are declared, the void residue is adjudicated, and the cell remainder is transferred}\ }
\]

%dsp-level:0
\[
(7)\quad\textbf{the declaration ledger:}
\]

%dsp-level:0
\[
\qquad \textbf{four addenda issued by this tier:}\ \ \text{A-BM.10}\ (\aleph\text{-domain}),
\]

%dsp-level:0
\[
\ \ \text{A-BM.11}\ (\text{corner convention}),
\]

%dsp-level:0
\[
\ \ \text{A-BM.12}\ (\text{resolvent interface}),
\]

%dsp-level:0
\[
\ \ \text{A-BM.13}\ (\text{statistics locus});
\]

%dsp-level:0
\[
\textbf{all four logged in §II.4 before use};
\]

%dsp-level:0
\[
\qquad \neg\text{a P4 member added by any of them};
\]

%dsp-level:0
\[
\qquad \neg\text{a dimensionful scale}\ [\mathrm{P1}];
\]

%dsp-level:0
\[
\qquad \neg\text{a measured value}\ [\mathrm{P5}];
\]

%dsp-level:0
\[
\qquad \text{§Cal-BM gains no parameter}\ [\mathrm{BM.143}(4)\text{'s consolidation is a register,}\ \neg\text{an addition}].
\]

\textbf{Req.} A8 (Part 1); D3.3 (Part 1); CT0.9 (Part 1); P1, P5 (Parts 1–2); Def-H.27, Def-H.32, FT.21, FT.62(5), FT.67(1), FT.67(3), FT.74, FT.75, FT.76(1), FT.76(3), FT.77 (Part 3); A3 (Part 1); Def-BM.14(3), Def-BM.16(5), Def-BM.19(3), Def-BM.20, BM.6(4), BM.7, BM.73, BM.79(1), BM.113(6), BM.124(3), BM.125, BM.128(5), BM.136(4), BM.138(6), A-BM.7 (this part); Census FT.76(1), FT.76(2) (row-wise); A-BM.10, A-BM.11, A-BM.13 (§II.4, logged before use); exact binomial and derangement enumeration, the Chebyshev identity $T_m(1)=1$ over $\mathbb Z[x]$, exhaustive parity censuses [math; machine: \texttt{checks/\allowbreak{}verify\_\allowbreak{}H\_\allowbreak{}bm\_\allowbreak{}closing\_\allowbreak{}tier.py} §H].

\textbf{Pf.}

%dsp-level:0
\[
(1):\ \text{the star sums}\ \sum_{a\ge2}\binom{n_{\mathrm{in}}}a\binom{n_{\mathrm{out}}}aa!4^a\ \text{evaluated exactly at}\ (2,2),
\]

%dsp-level:0
\[
(2,3),
\]

%dsp-level:0
\[
(3,3):\ (32,96,672)\ [\text{machine}];
\]

%dsp-level:0
\[
\qquad \text{the derangements}\ (!a)
\]

%dsp-level:0
\[
= \sum_{j\le a}(-1)^ja!/j!
\]

%dsp-level:0
\[
= (1,2,9,44)\ \text{and the hub values}\ (16,128,2304,45056)\ [\text{machine}]
\]

%dsp-level:0
\[
(3):\ k_f,
\]

%dsp-level:0
\[
k_{vf}\ \text{from}\ k
\]

%dsp-level:0
\[
= L\kappa_{\min}/2\ \text{with}\ \kappa_{\min}
\]

%dsp-level:0
\[
= 2\ (L\ \text{odd}),
\]

%dsp-level:0
\[
\ 1\ (L\ \text{even})\ [\mathrm{FT.32}(6);\ \mathrm{FT.45}]:\ \ L_f
\]

%dsp-level:0
\[
= (3,4,3,5,3)
\]

%dsp-level:0
\[
\mapsto (3,2,3,5,3);
\]

%dsp-level:0
\[
\quad L_{vf}=(3,3,4,3,5)\mapsto(3,3,2,3,5)\ [\text{machine}];
\]

%dsp-level:0
\[
\qquad \text{the massless readings:}\ 2-2T_m(1)=0\ \text{and}\ 2-2T_m(\cos2\pi)=0\ \forall m\ [\text{machine over}\ \mathbb Z[x]]
\]

%dsp-level:0
\[
(2),
\]

%dsp-level:0
\[
(4)\text{--}(7):\ \text{the declaration rows cite their grounds and their loci; the §II.4 entries carry the registry grammar}\ [\text{machine: the logged-before-use document pins}].\ \blacksquare
\]

\textbf{Phys.} Two conflicts the corpus could compute but not settle are settled by declaration, each with its grounds on the record. The amplitude-domain question — literal against chain-local — resolves pointwise: the ensembles keep the frozen literal domain, and the sealed junctions, where the chain-local pass-through reads were always the forced datum, run the literal-forced domain E2, which is canonical because it fixes the gauge by vertex type rather than by a choice of decomposition; E1 stays computed and unadopted. The corner convention resolves to the vertex-figure pairing, the loop-consistent reading, which is where this part's whole defect ladder already lives — and the declaration is load-bearing, because the per-class statistics census provably differs between the variants.

\end{theorem}

\medskip\noindent\textit{\texorpdfstring{BM.140 — The ceiling study: the Platonic defect census, the dodecahedron selection theorem, the five-currency range table, and the adjudication of deficit 12}{BM.140 - The ceiling study: the Platonic defect census, the dodecahedron selection theorem, the five-currency range table, and the adjudication of deficit 12}.}\ Structural-constraints record: see Appendix B, Section B.2.8 of the Addendum.\par\medskip

\begin{theorem}[{\texorpdfstring{BM.141 — The Jensen residue: the exact closed form, its range, its zero locus, the two-chain companion, and the width adjudication}{BM.141 - The Jensen residue: the exact closed form, its range, its zero locus, the two-chain companion, and the width adjudication}}]\label{tag:BM-141}
\textbf{Thm.}

%dsp-level:0
\[
(1)\quad\textbf{the residue computed:}
\]

%dsp-level:0
\[
\qquad J_3(p)\ :=\ \langle E\rangle_{\mathrm{mat}}\bigl\langle E^{-1}\bigr\rangle_{\mathrm{mat}}-1\ =\ \frac{18-27p+11p^2}{2D^2}-1,
\]

%dsp-level:0
\[
\qquad D=3-3p+p^2,
\]

%dsp-level:0
\[
\qquad p=1-4\varepsilon_q
\]

%dsp-level:0
\[
\qquad [\mathrm{BM.94}(6),(8)\ \text{re-cited}];
\]

%dsp-level:0
\[
\boxed{\ J_3(p)\ =\ \frac{p\,(1-p)\,\bigl(9-10p+2p^2\bigr)}{2D^2}\ }
\]

%dsp-level:0
\[
\qquad [\text{machine: the polynomial identity}\ 18-27p+11p^2-2D^2=p(1-p)(9-10p+2p^2)\ \text{exact in}\ \mathbb Q[p]];
\]

%dsp-level:0
\[
(2)\quad\textbf{the sign, the zeros, and the range:}
\]

%dsp-level:0
\[
\qquad J_3\ge0\ \ \text{on}\ p\in(0,1]
\]

%dsp-level:0
\[
\qquad [\mathrm{BM.94}(8)\text{'s Cauchy--Schwarz row re-derived};\ \text{machine:}\ 9-10p+2p^2>0\ \text{on}\ (0,1]\ \text{— its roots}\ \frac{5\pm\sqrt7}2\ \text{both exceed}\ 1\ \text{in}\ \mathbb Q(\sqrt7):\ \frac{5-\sqrt7}2>1\iff3>\sqrt7\iff9>7];
\]

%dsp-level:0
\[
\boxed{\ J_3=0\iff p=1\iff\varepsilon_q=0\ }
\]

%dsp-level:0
\[
\qquad [\text{machine: the factor census}];
\]

%dsp-level:0
\[
\qquad J_3\bigl(\tfrac12\bigr)=\frac9{49}
\]

%dsp-level:0
\[
\qquad [\mathrm{BM.94}(8):\ \tfrac{58}{49}-1;\ \text{machine}];
\]

%dsp-level:0
\[
J_3<\frac15\ \ \text{on}\ (0,1)
\]

%dsp-level:0
\[
\qquad [12p^4-72p^3+125p^2-81p+18>0;\ \mathrm{BM.94}(8)\ \text{re-derived};\ \text{machine: exact rational bisection, no roots on}\ (0,1)];
\]

%dsp-level:0
\[
\boxed{\ \sup_{(0,1)}J_3=J_3(p_\ast)\in\Bigl(\frac{93}{500},\ \frac3{16}\Bigr),\qquad p_\ast\in\Bigl(\frac{288131}{524288},\ \frac{576263}{1048576}\Bigr)\ }
\]

%dsp-level:0
\[
\qquad [\text{machine: the unique critical point by Sturm isolation of}\ -44p^5+294p^4-792p^3+1062p^2-684p+162;\qquad J_3<\tfrac3{16}\ \text{strict:}\ 38p^4-228p^3+394p^2-252p+54>0\ \text{on}\ (0,1),\ \text{no roots}];
\]

%dsp-level:0
\[
(3)\quad\textbf{the two-chain companion:}
\]

%dsp-level:0
\[
\qquad \langle E\rangle^{(2)}_{\mathrm{mat}}=\frac4{1+u},
\]

%dsp-level:0
\[
\qquad \bigl\langle E^{-1}\bigr\rangle^{(2)}_{\mathrm{mat}}=\frac{3u+1}{4(1+u)}
\]

%dsp-level:0
\[
\qquad [\mathrm{BM.137}(2);\ \text{machine}];
\]

%dsp-level:0
\[
\boxed{\ J_2(p)\ =\ \frac{4-3p}{(2-p)^2}-1\ =\ \frac{p(1-p)}{(2-p)^2}\ =\ \frac{u(1-u)}{(1+u)^2}\ }
\]

%dsp-level:0
\[
\qquad [\text{machine: the identity exact in}\ \mathbb Q[p]];
\]

%dsp-level:0
\[
\boxed{\ J_2\le\frac18,\qquad\text{equality}\iff p=\frac23\iff\varepsilon_q=\frac1{12}\ }
\]

%dsp-level:0
\[
\qquad [\text{machine: the SOS certificate}\ \frac{(1+u)^2}8-u(1-u)=\frac{(3u-1)^2}8\ge0;\ \text{the stationary point}\ 3p^2-8p+4=(3p-2)(p-2)];
\]

%dsp-level:0
\[
J_2\bigl(\tfrac12\bigr)=\frac19<\frac9{49}=J_3\bigl(\tfrac12\bigr);
\]

%dsp-level:0
\[
\qquad \frac18<\frac15;
\]

%dsp-level:0
\[
\qquad \textbf{the two-chain residue is strictly tighter than the three-chain bar}
\]

%dsp-level:0
\[
\qquad [\text{machine}];
\]

%dsp-level:0
\[
(4)\quad\textbf{the exact crossing of the two residues:}
\]

%dsp-level:0
\[
\qquad J_3-J_2=\frac{p\,(p-1)\,\bigl(6p^3-27p^2+40p-18\bigr)}{2D^2(2-p)^2}
\]

%dsp-level:0
\[
\qquad [\text{machine: the factorization exact in}\ \mathbb Q[p]];
\]

%dsp-level:0
\[
6p^3-27p^2+40p-18\ \ \text{is strictly increasing on}\ (0,1)\ \bigl[\text{its derivative}\ 18p^2-54p+40\ \text{has roots}\ \tfrac43,\tfrac53,\ \text{both outside}\bigr]\ \text{with the single root}\ p_c\in\Bigl(\frac{108267}{131072},\ \frac{866137}{1048576}\Bigr)
\]

%dsp-level:0
\[
\qquad [\text{machine: Sturm isolation}];
\]

%dsp-level:0
\[
\boxed{\ J_3>J_2\iff p<p_c;\qquad J_3<J_2\iff p>p_c\ }\qquad[\text{machine: the sign census on exact rationals}];
\]

%dsp-level:0
\[
(5)\quad\textbf{the width adjudication (the two formed reads of}\ \Delta x_{\mathrm B}\textbf{):}
\]

%dsp-level:0
\[
\qquad \Delta x_{\mathrm B}^{\mathrm{eff}}=\frac{\mathfrak h}{\langle E\rangle_{\mathrm{mat}}\ln2}=D\,\Delta x_{\mathrm B}^{cryst};
\]

%dsp-level:0
\[
\qquad \frac{\mathfrak h}{\ln2}\bigl\langle E^{-1}\bigr\rangle_{\mathrm{mat}}\ \ge\ \Delta x_{\mathrm B}^{\mathrm{eff}}
\]

%dsp-level:0
\[
\qquad [\mathrm{BM.94}(8)\ \text{re-cited}];
\]

%dsp-level:0
\[
\boxed{\ \textbf{the relative gap of the two reads is exactly the Jensen residue:}\qquad\Bigl(\bigl\langle E^{-1}\bigr\rangle_{\mathrm{mat}}-\frac1{\langle E\rangle_{\mathrm{mat}}}\Bigr)\langle E\rangle_{\mathrm{mat}}=J_3\ }
\]

%dsp-level:0
\[
\qquad [\text{machine: the identity}];
\]

%dsp-level:0
\[
\textbf{the adjudication:}
\]

%dsp-level:0
\[
\qquad \text{the tier's evaluations at the width of the mean}\ [\mathrm{BM.94}(8);\ \mathrm{BM.98}(5)]\ \textbf{stand, with the difference now an exact computed function instead of a residue};
\]

%dsp-level:0
\[
(6)\quad\textbf{status:}
\]

%dsp-level:0
\[
\qquad \textbf{derived:}\ (1)\text{--}(5);
\]

%dsp-level:0
\[
\qquad \boxed{\ (\mathrm{R6})\ \textbf{is closed: the Jensen residue is computed as an exact rational function with its zero locus, its exact supremum enclosure and its two-chain companion}\ }
\]

%dsp-level:0
\[
[\mathrm{BM.94}(8),\ \mathrm{BM.98}(5)\text{'s residue rows discharged}];
\]

%dsp-level:0
\[
\qquad \neg\text{a declaration},
\]

%dsp-level:0
\[
\qquad \neg\text{a P4 addition},
\]

%dsp-level:0
\[
\qquad \neg\text{a measured value}\ [\mathrm{P5}].
\]

\textbf{Req.} A4 (Part 1); D0.4 (Part 1); CT0.8 (Part 1); P5 (Parts 1–2); BM.94(6), BM.94(8), BM.98(5), BM.137(2) (this part); polynomial factorization over $\mathbb Q$, Sturm isolation and exact rational bisection, SOS certificates, the Cauchy–Schwarz equality locus [math; machine: \texttt{checks/\allowbreak{}verify\_\allowbreak{}H\_\allowbreak{}bm\_\allowbreak{}closing\_\allowbreak{}tier.py} §J].

\textbf{Pf.}

%dsp-level:0
\[
(1):\ 18-27p+11p^2-2(3-3p+p^2)^2=p(1-p)(9-10p+2p^2)\ [\text{machine: expansion in}\ \mathbb Q[p]];
\]

%dsp-level:0
\[
\qquad (2):\ \text{the root census of}\ 9-10p+2p^2\ \text{in}\ \mathbb Q(\sqrt7);
\]

%dsp-level:0
\[
\qquad \text{the bisection of}\ 12p^4-72p^3+125p^2-81p+18\ \text{and}\ 38p^4-228p^3+394p^2-252p+54\ \text{on}\ (0,1)\ \text{at exact rational endpoints}\ [\text{machine}]
\]

%dsp-level:0
\[
(3):\ \bigl\langle E^{-1}\bigr\rangle^{(2)}_{\mathrm{mat}}=\frac{\frac12\cdot2u+\frac14p}{1+u}\ \text{normalized}\ [\mathrm{BM.137}(2)\text{'s measure};\ \text{machine}];
\]

%dsp-level:0
\[
\qquad J_2=\frac4{1+u}\cdot\frac{3u+1}{4(1+u)}-1=\frac{3u+1-(1+u)^2}{(1+u)^2}=\frac{u(1-u)}{(1+u)^2}\ [\text{machine}];
\]

%dsp-level:0
\[
\qquad \text{the SOS certificate and the stationary factorization}\ [\text{machine}]
\]

%dsp-level:0
\[
(4):\ J_3-J_2\ \text{brought to the common denominator}\ 2D^2(2-p)^2\ \text{and factored}\ [\text{machine}];
\]

%dsp-level:0
\[
\qquad (5):\ \bigl\langle E^{-1}\bigr\rangle\langle E\rangle-\langle E\rangle/\langle E\rangle=J_3\ [\text{machine}];
\]

%dsp-level:0
\[
\qquad (6):\ \text{assembly}.\ \blacksquare
\]

\textbf{Phys.} The last pure-mathematics residue of the era is discharged: the Jensen residue of the materialized flicker distribution is one exact rational function, $p(1-p)(9-10p+2p^2)$ over twice the squared geometric sum, vanishing precisely at the sharp crystallization point, capped strictly below one fifth with its supremum enclosed between two rationals and attained at a uniquely isolated critical point. The two-chain companion is tighter — bounded by one eighth with a sum-of-squares certificate and equality at threshold over twelve — and the two residues cross exactly once, at an isolable rational-bracketed point. The gap between the two formed width reads of the Born envelope is precisely this residue, so the era's evaluations at the width of the mean stand with their correction now computed rather than named.

\end{theorem}

\begin{theorem}[{\texorpdfstring{BM.142 — The supply vertex, the pair amplitude, and the annihilation cross-section: the four grounds resolved and residues $(\mathrm{R9}')$, $(\mathrm{R9}'')$ closed}{BM.142 - The supply vertex, the pair amplitude, and the annihilation cross-section: the four grounds resolved and residues , closed}}]\label{tag:BM-142}
\textbf{Thm.}

%dsp-level:0
\[
(1)\quad\textbf{the four grounds of}\ \mathrm{BM.122}(6)\ \textbf{resolved, one by one:}\qquad\begin{array}{l|l|l}\text{ground} & \text{the resolution} & \text{the row}\\\hline(\mathrm O1)\ \neg\text{a gauged chiral coupling} & \textbf{dissolved for the EM legs:}\ \text{the gauged coupling of the charge read is constructed chart-side} & [\mathrm{Def\text{-}BM.12}(1)\text{--}(3):\ \partial_\mu F^{\mu\nu}=eJ_Q^\nu;\ \mathrm{EM.15}]\\\text{} & \text{the W-species chiral gauging remains}\ \neg\text{constructed}\ —\ \text{barred for weak legs, recorded} & [\mathrm{BM.76};\ \mathrm{A\text{-}BM.7}\ \text{scope}]\\\hline(\mathrm O2)\ \text{deficit}\ 1' & \textbf{discharged:}\ \text{the assignment rule is issued and the attachment derived} & [\mathrm{A\text{-}BM.7};\ \mathrm{BM.73}(1),(2);\ \mathrm{BM.75}(3)]\\\hline(\mathrm O3)\ \neg\text{a momentum resolvent} & \textbf{discharged:}\ \text{the resolvent and the phase space are constructed} & [\mathrm{Def\text{-}BM.20};\ \mathrm{BM.134}]\\\hline(\mathrm O4)\ \text{the transit leg carries no winding carrier} & \textbf{re-scoped:}\ \text{the exclusion stands, but the vertex needs}\ \neg\text{an}\ \varepsilon/j\ \text{read of the legs} & [\mathrm{Def\text{-}H.38}\ \text{Scope frozen}];\ \text{the legs couple by}\ (q,\ \text{helicity}):\ N_{\mathrm{pol}}=2\ [\mathrm{P2};\ \mathrm{EM.8}]\\\text{} & \text{and the register's moves apply to transit legs} & [\mathrm{BM.126}(5)\ \text{domain row}]\end{array}
\]

%dsp-level:0
\[
(2)\quad\textbf{the supply vertex (constructed):}
\]

%dsp-level:0
\[
\qquad \boxed{\ \widehat{\mathcal V}_{\mathrm{sup}}\ :=\ e(k,E)^2\cdot\bigl(\hat J_+\otimes\hat J_+\bigr)\ }
\]

%dsp-level:0
\[
\qquad \text{on}\ \mathcal J\otimes\mathcal J\otimes\bigl(\text{the two transit helicity legs}\bigr)
\]

%dsp-level:0
\[
\qquad [\mathrm{Def\text{-}BM.18}(2);\ \mathrm{Def\text{-}BM.12}(2);\ \mathrm{Def\text{-}BM.20}(4)];
\]

%dsp-level:0
\[
\bigl(\hat J_+\otimes\hat J_+\bigr)^2=0
\]

%dsp-level:0
\[
\qquad [\hat J_+^2=0;\ \mathrm{BM.120}(1);\ \text{machine: the}\ 4\times4\ \text{identity over}\ \mathbb Z]:
\]

%dsp-level:0
\[
\qquad \boxed{\ \textbf{the vertex is nilpotent: the register cannot supply a pair twice}\ };
\]

%dsp-level:0
\[
\textbf{the vertex rows:}
\]

%dsp-level:0
\[
\qquad \text{charge balance}\ Q_{\mathrm{in}}=Q_{\mathrm{out}}=0\ [\mathrm{BM.122}(1)];
\]

%dsp-level:0
\[
\qquad \text{the threshold gate}\ \chi_{\{NE\ge\mathcal T_0\}}\ [\mathrm{BM.122}(2);\ \mathrm{Def\text{-}BM.17}(3)];
\]

%dsp-level:0
\[
\qquad \text{the helicity census}\ 2\times2
\]

%dsp-level:0
\[
= 4\ \text{combinations, equally weighted by the metric's charge-blindness}\ [\mathrm{P2}:\ N_{\gamma\text{-pol}}=2;\ \mathrm{EM.8};\ \mathrm{BM.58}];
\]

%dsp-level:0
\[
(3)\quad\textbf{the pair amplitude (formed, symbolic):}
\]

%dsp-level:0
\[
\qquad \boxed{\ \mathcal M_{\mathrm{pair}}(z;k,E)\ =\ e(k,E)^2\,q^2\,x_+^2\,\chi\ =\ e_0^2\,q^2x_+^2\,\chi\cdot\frac{N(k,E)\,f(k)}{1+G(E)}\ }
\]

%dsp-level:0
\[
\qquad [\mathrm{Def\text{-}BM.20}(4)\text{'s}\ \mathcal V=1;\ \mathrm{BM.133}(6);\ \mathrm{BM.132}(2);\ \text{machine: the monomial census}];
\]

%dsp-level:0
\[
[\mathcal M]=\mathfrak q^0,\qquad[e]=\mathfrak q^0,\qquad[x_\pm]=\mathfrak q^0\qquad[\mathrm{CT2.1};\ \mathrm{BM.102}(2);\ \text{machine: the exponent census}];
\]

%dsp-level:0
\[
\textbf{the reduction clause:}
\]

%dsp-level:0
\[
\qquad e=1,\ q=1:\quad\mathcal M_{\mathrm{pair}}=\mathcal W_{\text{channel}}=x_+^2\chi
\]

%dsp-level:0
\[
\qquad [\mathrm{BM.122}(3);\ \mathrm{Def\text{-}BM.20}(4);\ \text{machine}];
\]

%dsp-level:0
\[
\textbf{the exact zeros:}
\]

%dsp-level:0
\[
\qquad \mathcal M=0\iff e_0=0\ \vee\ q=0\ \vee\ \varepsilon_q=\tfrac14\ \vee\ \chi=0\ \vee\ N=0\ \vee\ f=0
\]

%dsp-level:0
\[
\qquad [\text{machine: the factor census}];
\]

%dsp-level:0
\[
(4)\quad\textbf{the pair cross-section (the formed chart read):}
\]

%dsp-level:0
\[
\qquad \sigma_{\mathrm{pair}}:=\frac{|\mathcal M_{\mathrm{pair}}|^2\,\Phi_2\bigl(z;\,2m_{\mathrm{hid}},N\bigr)}{F_{\gamma\gamma}}
\]

%dsp-level:0
\[
\qquad [\mathrm{BM.134}(4);\ \text{the final state is the massive pair}]:
\]

%dsp-level:0
\[
\boxed{\ \sigma_{\mathrm{pair}}(z;k,E)\ =\ \frac{e_0^4\,q^4x_+^4\,\chi\ N\,f^2}{16\pi\,z^3\,(1+G)^2}\ \sqrt{z^2-4N^2\bigl(2m_{\mathrm{hid}}\bigr)^2}\ }
\]

%dsp-level:0
\[
\qquad [\text{machine: the coefficient algebra exact};\ m_{\mathrm{pair}}=2m_{\mathrm{hid}}\ [\mathrm{BM.111}(1)]];
\]

%dsp-level:0
\[
[\sigma_{\mathrm{pair}}]=\mathfrak q^{-1}\ =\ \text{the area species of}\ (\Delta x_{\mathrm B})^2\qquad[\text{machine: the exponent census}];
\]

%dsp-level:0
\[
\textbf{the threshold rows:}
\]

%dsp-level:0
\[
\qquad \sigma_{\mathrm{pair}}=0\ \ \text{exactly at}\ \ z=2N\cdot2m_{\mathrm{hid}}\ \ (\text{the square-root onset});
\]

%dsp-level:0
\[
\qquad \sigma_{\mathrm{pair}}>0\ \ \text{strictly above it};
\]

%dsp-level:0
\[
\qquad \text{formed on}\ \ z\ge4Nm_{\mathrm{hid}}\ \wedge\ k\ge\frac G{4\pi}\ \wedge\ NE\ge\mathcal T_0
\]

%dsp-level:0
\[
\qquad [\text{machine}];
\]

%dsp-level:0
\[
(5)\quad\textbf{the annihilation cross-section:}
\]

%dsp-level:0
\[
\qquad \mathcal M_{\mathrm{ann}}=e(k,E)^2q^2x_-^2\chi
\]

%dsp-level:0
\[
\qquad [\mathrm{BM.123}(1)\text{'s reverse move}\ \hat J_-\otimes\hat J_-;\ \text{the same vertex normalization};\ m\bigl(c^{\mathrm{CP}}\bigr)=m(c)\ [\mathrm{BM.27};\ \mathrm{BM.108}(3)]];
\]

%dsp-level:0
\[
\sigma_{\mathrm{ann}}:=\frac{|\mathcal M_{\mathrm{ann}}|^2\,\Phi_2^{(\gamma\gamma)}}{F_{\mathrm{pair}}}\qquad[\text{the final state massless},\ \text{the initial state massive}]:
\]

%dsp-level:0
\[
\boxed{\ \sigma_{\mathrm{ann}}(z;k,E)\ =\ \frac{e_0^4\,q^4x_-^4\,\chi\ N\,f^2}{16\pi\,z\,(1+G)^2\,\sqrt{z^2-4N^2\bigl(2m_{\mathrm{hid}}\bigr)^2}}\ }
\]

%dsp-level:0
\[
\qquad [\text{machine: the coefficient algebra exact};\ [\sigma_{\mathrm{ann}}]=\mathfrak q^{-1}];
\]

%dsp-level:0
\[
\textbf{the threshold behaviour:}
\]

%dsp-level:0
\[
\qquad \sigma_{\mathrm{ann}}\ \ \text{has the exact square-root pole at}\ \ z=4Nm_{\mathrm{hid}};
\]

%dsp-level:0
\[
\qquad \sigma_{\mathrm{ann}}\cdot F_{\mathrm{pair}}\ \ \text{is regular there}
\]

%dsp-level:0
\[
\qquad [\text{machine: the product}\ =\ \frac{e_0^4q^4x_-^4\chi\,f^2}{8\pi\,N\,(1+G)^2}\ \text{regular in the formed factors}];
\]

%dsp-level:0
\[
(6)\quad\textbf{the ratio and the detailed-balance elevation:}
\]

%dsp-level:0
\[
\qquad \boxed{\ \frac{\sigma_{\mathrm{ann}}}{\sigma_{\mathrm{pair}}}\ =\ \Bigl(\frac{x_-}{x_+}\Bigr)^{4}\cdot\frac{z^2}{z^2-4N^2(2m_{\mathrm{hid}})^2}\ }
\]

%dsp-level:0
\[
\qquad [\text{machine: the symbolic division exact}];
\]

%dsp-level:0
\[
\textbf{the amplitude-level detailed balance:}
\]

%dsp-level:0
\[
\qquad |\mathcal M_{\mathrm{ann}}|=|\mathcal M_{\mathrm{pair}}|\iff\varepsilon_q=\frac18
\]

%dsp-level:0
\[
\qquad [\mathrm{BM.123}(2)\text{'s indicator elevated to the amplitude};\ \text{machine}];
\]

%dsp-level:0
\[
\text{at}\ \varepsilon_q
\]

%dsp-level:0
\[
= \frac18\ \text{the cross-section ratio is the pure kinematic factor}\ \frac{z^2}{z^2-16N^2m_{\mathrm{hid}}^2}>1\ \ \text{strictly above threshold}
\]

%dsp-level:0
\[
\qquad [\text{machine}];
\]

%dsp-level:0
\[
(7)\quad\textbf{the width row and its species bar:}
\]

%dsp-level:0
\[
\qquad \text{the chart-frequency width}\ \Gamma^{\mathrm{chart}}
\]

%dsp-level:0
\[
= \frac{|\mathcal M|^2\Phi_2}{2z}\ \ \text{carries}\ [\Gamma^{\mathrm{chart}}]
\]

%dsp-level:0
\[
= \mathfrak q^{-1/2}\ \ \text{at}\ \mathcal V
\]

%dsp-level:0
\[
= 1\ —\ \text{a time,}\ \neg\text{a rate}:
\]

%dsp-level:0
\[
\text{the mass-dressed normalization that would convert it is}\ \neg\text{declared}\ [\mathrm{Def\text{-}BM.20}(4)\text{'s minimality}]\implies\boxed{\ \textbf{the formed width read stays the ledger's per-tick species}\ }\ [\mathrm{BM.116}(1);\ \mathrm{BM.44}(4),(5):\ [\Gamma]=\mathfrak q^0\ \text{ticks}^{-1};\ \tau_{\mathrm{life}}\ge3\ \text{ticks re-cited}];
\]

%dsp-level:0
\[
(8)\quad\textbf{the closure rows:}
\]

%dsp-level:0
\[
\qquad \boxed{\ (\mathrm{R9}')\ \textbf{is closed: the annihilation cross-section is constructed as an exact symbolic chart-side read on the formed domain}\ }
\]

%dsp-level:0
\[
\boxed{\ (\mathrm{R9}'')\ \textbf{is closed at the vertex level: the two-transit-quantum supply vertex is constructed, nilpotent, charge-balanced and threshold-gated}\ }
\]

%dsp-level:0
\[
\qquad [\mathrm{BM.118}(7),\ \mathrm{BM.130}(7)\text{'s rows discharged}];
\]

%dsp-level:0
\[
\mathrm{BM.130}(3)\text{'s}\ \neg\text{cross-section row is discharged by construction and re-scoped,}\ \neg\text{retracted};
\]

%dsp-level:0
\[
\qquad \mathrm{BM.122}(6)\text{'s four grounds carry}\ (1)\text{'s per-row dispositions};
\]

%dsp-level:0
\[
(9)\quad\textbf{the bars:}
\]

%dsp-level:0
\[
\qquad \neg\text{a numerical cross-section, rate or branching fraction is evaluated}\ [\mathrm{P5};\ \mathrm{BM.44}(5)];
\]

%dsp-level:0
\[
\qquad \neg\text{an S-matrix unitarity claim}\ [\mathrm{Def\text{-}BM.20}(5)];
\]

%dsp-level:0
\[
\qquad \text{amplitude}\ne\text{mass}\ [\mathrm{FT.79}(5)];
\]

%dsp-level:0
\[
e_0\ \ \text{remains calibration}\ [\mathrm{EM.25}(7);\ \mathrm{BM.143}(4)];
\]

%dsp-level:0
\[
\qquad \neg\text{a measured external constant is identified with}\ \sigma_{\mathrm{pair}},\ \sigma_{\mathrm{ann}}\ \text{or any factor}\ [\mathrm{P5}];
\]

%dsp-level:0
\[
\qquad \neg\text{a species name enters any clause}\ [\mathrm{Def\text{-}BM.3}(7)].
\]

\textbf{Req.} A4, A8 (Part 1); D0.4, D1.2 (Part 1); CT1.22, CT2.1 (Part 2); EM.8, EM.15, EM.25(7) (Part 2); P2, P5 (Parts 1–2); Def-H.38 (Part 3); FT.79(5) (Part 3); Def-BM.12, Def-BM.17(3), Def-BM.18(2), Def-BM.20, BM.27, BM.44(4), BM.44(5), BM.58, BM.73(1), BM.73(2), BM.75(3), BM.76, BM.102(2), BM.108(3), BM.111(1), BM.116(1), BM.118(7), BM.120(1), BM.122, BM.123, BM.126(5), BM.130(3), BM.130(7), BM.132(2), BM.133(6), BM.134(4) (this part); the declaration addendum A-BM.7 (pointers); exact matrix identities over $\mathbb Z$, monomial censuses in symbolic rings, the coefficient algebra of the two-body reads [math; machine: \texttt{checks/\allowbreak{}verify\_\allowbreak{}H\_\allowbreak{}bm\_\allowbreak{}closing\_\allowbreak{}tier.py} §K].

\textbf{Pf.}

%dsp-level:0
\[
(1):\ \text{the four rows cite their discharging tags; the W-species bar and the}\ \mathrm{E2}\text{-domain inheritance recorded}\ [\mathrm{Def\text{-}BM.20}(5);\ \text{machine: the document pins}]
\]

%dsp-level:0
\[
(2):\ \bigl(\hat J_+\otimes\hat J_+\bigr)^2=\hat J_+^2\otimes\hat J_+^2=0\otimes0=0\ [\mathrm{BM.120}(1);\ \text{machine: the}\ 4\times4\ \text{product over}\ \mathbb Z];
\]

%dsp-level:0
\[
\qquad \text{the helicity census}\ 2^2=4\ [\mathrm{P2};\ \text{machine}]
\]

%dsp-level:0
\[
(3),(4):\ \mathcal M\ \text{assembled from}\ e^2=e_0^2Nf\mathsf g=e_0^2Nf/(1+G)\ [\mathrm{BM.133}(6);\ \text{machine}];
\]

%dsp-level:0
\[
\qquad \sigma_{\mathrm{pair}}
\]

%dsp-level:0
\[
= \frac{|\mathcal M|^2}{F_{\gamma\gamma}}\Phi_2
\]

%dsp-level:0
\[
= \frac{e_0^4q^4x_+^4\chi N^2f^2}{(1+G)^2}\cdot\frac{\sqrt{z^2-16N^2m^2}}{8\pi N^3z}\cdot\frac{N^2}{2z^2}\ [\mathrm{BM.134}(4);\ \text{machine: the coefficient algebra}];
\]

%dsp-level:0
\[
\text{the exponent census}\ [\sigma]=\mathfrak q^{-1}\ [\mathrm{CT2.1}]
\]

%dsp-level:0
\[
(5),(6):\ \sigma_{\mathrm{ann}}=\frac{e_0^4q^4x_-^4\chi N^2f^2}{(1+G)^2}\cdot\frac1{8\pi N^3}\cdot\frac{N^2}{2z\sqrt{z^2-16N^2m^2}}\ [\text{machine}];
\]

%dsp-level:0
\[
\qquad \text{the ratio by symbolic division}\ [\text{machine}];
\]

%dsp-level:0
\[
\qquad \text{the detailed-balance locus}\ x_-=x_+\iff\varepsilon_q=\tfrac18\ [\mathrm{BM.120}(5);\ \mathrm{BM.123}(2)]
\]

%dsp-level:0
\[
(7):\ \bigl[\tfrac{|\mathcal M|^2\Phi_2}{2z}\bigr]=\mathfrak q^0\cdot\mathfrak q^0\cdot\mathfrak q^{-1/2}=\mathfrak q^{-1/2}\ [\text{machine: the exponent census}];
\]

%dsp-level:0
\[
\qquad (8),(9):\ \text{assembly and the token scan}.\ \blacksquare
\]

\textbf{Phys.} The two-photon bar falls ground by ground: the electromagnetic gauging of the charge read existed chart-side, the assignment rule was issued, the resolvent is now constructed, and the transit legs couple by charge and helicity — the one exclusion that stands, the winding carrier of a trivial chain, simply never entered the vertex. What is built is minimal and honest: a nilpotent supply vertex that cannot fire twice, an amplitude that reduces to the register's own channel weight at unit coupling, and two cross-sections that are exact symbolic functions of the envelope ratio, the local energy, the flicker threshold and the calibration coupling — with an exact square-root onset for production, an exact square-root pole for annihilation whose flux product stays regular, and a ratio that carries the detailed-balance point of the register elevated to the amplitude level. The chart-frequency width is refused on dimensional grounds rather than fudged: the per-tick ledger species remains the corpus's width read. Both residues the annihilation tier named are closed, and nothing measured entered any row.

\end{theorem}

\medskip\noindent\textit{\texorpdfstring{BM.143 — The house cleaning: the void residues closed, the phantom gaps reclassified, the calibration register consolidated, and the Part-5 requirements named}{BM.143 - The house cleaning: the void residues closed, the phantom gaps reclassified, the calibration register consolidated, and the Part-5 requirements named}.}\ Structural-constraints record: see Appendix B, Section B.4.2 of the Addendum.\par\medskip

\medskip\noindent\textit{\texorpdfstring{BM.144 — The status of the Closing Tier and the closure of the Baryomesonic era}{BM.144 - The status of the Closing Tier and the closure of the Baryomesonic era}.}\ Structural-constraints record: see Appendix B, Section B.1.34 of the Addendum.\par\medskip

\phantomsection
\section[Tier BM.24 — The Continuum-Limit Tier: the Force Operator…]{\texorpdfstring{Tier BM.24 — The Continuum-Limit Tier: the Force Operator over the Graph, the Continuum-Limit Function, the Metric–Energy–Continuum Map, the Breaking Threshold, the Locked-Sector Complement, the Universality, the Reclassification of the Continuum Conditions, and the Closure of the Foundational Era}{Tier BM.24 - The Continuum-Limit Tier: the Force Operator over the Graph, the Continuum-Limit Function, the Metric-Energy-Continuum Map, the Breaking Threshold, the Locked-Sector Complement, the Universality, the Reclassification of the Continuum Conditions, and the Closure of the Foundational Era}}

\begin{theorem}[{\texorpdfstring{BM.145 — The force operator acting over the graph: the generator census, the mask ladder, and the derived action}{BM.145 - The force operator acting over the graph: the generator census, the mask ladder, and the derived action}}]\label{tag:BM-145}
\textbf{Thm.}

%dsp-level:0
\[
(1)\quad\textbf{the complete operator (re-cited, exact):}
\]

%dsp-level:0
\[
\qquad \mathbf F=F_{(+,\mathcal R)}\oplus F_{(-,\mathcal R)}\oplus F_{(+,\mathcal V)}\oplus F_{(-,\mathcal V)}\oplus F_{\mathrm{st}}\ \ \text{always applies in full}
\]

%dsp-level:0
\[
\qquad [\mathrm{A8};\ \mathrm{BM.52}(1)];
\]

%dsp-level:0
\[
F_{(+,\mathcal R)}=F_G\oplus F_H\oplus F_{D^+};
\]

%dsp-level:0
\[
\qquad F_{(-,\mathcal R)}=F_{W^+}\oplus F_{W^-}\oplus F_Z;
\]

%dsp-level:0
\[
\qquad F_{(+,\mathcal V)}=F_{J^+};
\]

%dsp-level:0
\[
\qquad F_{(-,\mathcal V)}=F_{J^-};
\]

%dsp-level:0
\[
\qquad F_{\mathrm{st}}=F_{\mathrm{EM}}\oplus F_{\mathrm{strong}}\oplus F_{D^-}
\]

%dsp-level:0
\[
\qquad [\mathrm{A8}];
\]

%dsp-level:0
\[
\textbf{the pair-arena block census:}\qquad\begin{array}{c|c|c}\text{aspect} & \text{the block} & \dim_{\mathbb R}\\\hline F_{J^+} & \mathcal E_{-+}=\operatorname{span}\{J_i\} & 3\\ F_{J^-} & \mathcal E_{--}=\operatorname{span}\{K_i\} & 3\\ F_G & \operatorname{span}\{\hat E^{ij}\}\ \text{traceless} & 5\\ F_W & \mathcal E_{+-}=\operatorname{span}\{W_i\} & 3\\ F_H & \operatorname{span}\{\Sigma\} & 1\\ F_{D^+} & \operatorname{span}\{E_0\} & 1\\\hline\text{total} & \mathfrak H(\mathbb C^4) & 16=C_{\mathrm{cap}}\end{array}\qquad[\mathrm{CT0.12};\ \mathrm{CT1.1};\ \mathrm{P2}:\ \text{machine:}\ 3+3+5+3+1+1=16];
\]

%dsp-level:0
\[
\textbf{the stack aspects act on the vertex/slot registers,}\ \neg\textbf{on the pair channel:}
\]

%dsp-level:0
\[
\qquad F_{\mathrm{EM}},\ F_{\mathrm{strong}},\ F_{D^-}
\]

%dsp-level:0
\[
\qquad [\mathrm{CT0.12}\ \mathrm{Rmk};\ \mathrm{D3.1};\ \mathrm{D4.1}];
\]

%dsp-level:0
\[
(2)\quad\textbf{the activation and the mask ladder:}
\]

%dsp-level:0
\[
\qquad \mathcal A(0)
\]

%dsp-level:0
\[
= \{J\};
\]

%dsp-level:0
\[
\quad\mathcal A(1)
\]

%dsp-level:0
\[
= \{J,G,W,H,D^+\};
\]

%dsp-level:0
\[
\quad\mathcal A(2)
\]

%dsp-level:0
\[
= \mathcal A(1)\cup\{\mathrm{EM}\};
\]

%dsp-level:0
\[
\quad\mathcal A(3)
\]

%dsp-level:0
\[
= \mathcal A(2)\cup\{\mathrm{strong}\};
\]

%dsp-level:0
\[
\quad\mathcal A(n\ge4)
\]

%dsp-level:0
\[
= \mathcal A(3)\cup\{D^-\}
\]

%dsp-level:0
\[
\qquad [\mathrm{CT0.13}];
\]

%dsp-level:0
\[
\mathbf F_M^{(n)}=\bigoplus_{A\in\mathcal A(n)}F_A\big|_{\mathcal K^{(A)}}
\]

%dsp-level:0
\[
\qquad [\mathrm{CT2.4}:\ \text{masking never mixes forces}];
\]

%dsp-level:0
\[
\qquad \dim\mathcal K_n=(6,\ 16,\ 17,\ 25)\ \ \text{strict}
\]

%dsp-level:0
\[
\qquad [\mathrm{CT0.14};\ \mathrm{BM.36}(1)];
\]

%dsp-level:0
\[
\textbf{the increments are derived:}
\]

%dsp-level:0
\[
\qquad 1\to2:\ +1=\dim Z\bigl(\iota_2(\mathfrak u(2))\bigr)=\dim\mathcal K^{(\mathrm{EM})}
\]

%dsp-level:0
\[
\qquad [\mathrm{D2.12}];
\]

%dsp-level:0
\[
\qquad 2\to3:\ +8=\dim\mathfrak{su}(3)=\dim\mathcal K^{(\mathrm{strong})}
\]

%dsp-level:0
\[
\qquad [\mathrm{CT3.8}];
\]

%dsp-level:0
\[
\text{the observable bound}\ \dim\mathcal K_\Sigma\le1+5+1+3+8=18\qquad[\mathrm{FT.26}(3);\ \text{machine: the integer census}];
\]

%dsp-level:0
\[
(3)\quad\textbf{the action over the graph (the composition, every step derived):}
\]

%dsp-level:0
\[
\qquad \boxed{\ \hat{\mathcal F}:\ \bigl(\Gamma_{\mathcal R},\ \text{registrations}\bigr)\ \longmapsto\ \bigl(M,\ g_{\mu\nu},\ T_{\mu\nu},\ \operatorname{spec}\bigr)\ }
\]

%dsp-level:0
\[
\begin{array}{l|l|l}\text{step} & \text{the map} & \text{the locus}\\\hline\text{the weights} & (\Gamma,\sigma_e,\kappa_e,\alpha_e)\mapsto w(C)=\mathcal J(C)\prod_{e\in R}G(e)\,(-1)^{\aleph},\quad G(e)=\sigma_e e^{i\pi\kappa_e} & [\mathrm{A8};\ \mathrm{D3.5};\ \mathrm{D3.6};\ \mathrm{FT.9}]\\\hline\text{the domain} & \aleph\ \text{literal on the ensembles,}\ \mathrm{E2}\ \text{at the sealed junctions} & [\mathrm{Def\text{-}H.32}\ \text{Scopes};\ \text{A-BM.10};\ \mathrm{BM.139}(2)]\\\hline\text{the dictionary} & C\mapsto(a_{00},q_{ij},w_i)\mapsto g_{\mu\nu}=\Omega^2\bar g_{\mu\nu},\quad\Omega=e^{\alpha_\Omega a_{00}} & [\mathrm{D2.1};\ \mathrm{CT1.12}]\\\hline\text{the grouping} & Z[g]=\sum_{C:\,g(C)=g}|w(C)|\,e^{i\Phi(C)} & [\mathrm{CT1.20}]\\\hline\text{the action} & S_{\mathrm{eff}}=-i\log Z;\quad S_{\mathrm{geom}}=\frac1{2\kappa^2}\int\sqrt{-g}\,(R-2\Lambda) & [\mathrm{CT1.20};\ \mathrm{CT1.21};\ \mathrm{CT2.7}]\\\hline\text{the spectrum} & Z(\mathcal R,n)=\langle\mathrm{out}|(1-\mathbf F_M^{(n)})^{-1}|\mathrm{in}\rangle=\sum_a\frac{\langle\mathrm{out}|P_a|\mathrm{in}\rangle}{1-\lambda_a};\quad\text{poles}\iff1\in\operatorname{spec}\mathbf F_M^{(n)} & [\mathrm{CT2.5}]\\\hline\text{the source} & T_{\mu\nu}=-\frac2{\sqrt{-g}}\frac{\delta S_{\mathrm{force}}}{\delta g^{\mu\nu}}=T^{\mathrm{bd}}_{\mu\nu}+T^{\mathrm{tr}}_{\mu\nu} & [\mathrm{CT1.22}]\\\hline\text{the field equations} & G_{\mu\nu}+\Lambda g_{\mu\nu}=\kappa^2T_{\mu\nu} & [\mathrm{CT1.23}]\\\hline\text{the target} & M=\Gamma_{\mathcal R}/\!\!\sim\ \ \text{a smooth 4-manifold with the charts}\ (\tau_n,h_1,h_2,h_3) & [\mathrm{CT2.6};\ \mathrm{D1.8}]\end{array}
\]

%dsp-level:0
\[
\textbf{the composition is the action:}
\]

%dsp-level:0
\[
\qquad \text{no step adds a datum}\ —\ \text{each row is a derived map of the cited tag}
\]

%dsp-level:0
\[
\qquad [\text{machine: the locus census, row by row}];
\]

%dsp-level:0
\[
(4)\quad\textbf{the state-side action (the second face, separated by species):}
\]

%dsp-level:0
\[
\qquad \text{the Haar draws and the Gibbs tilt act on the pair states,}\ \neg\text{on the arena blocks:}
\]

%dsp-level:0
\[
\qquad \mathcal T(\rho)=U\rho U^\dagger,\quad\tilde{\mathcal T}_\gamma(\rho)=\rho e^{-\gamma\rho}/Z_\gamma
\]

%dsp-level:0
\[
\qquad [\mathrm{A6};\ \mathrm{D0.3};\ \mathrm{L0}];
\]

%dsp-level:0
\[
V'=\Bigl(1-\frac\gamma4\Bigr)^2V+O(\delta\rho^3),\qquad\mu'=\Bigl(1-\frac\gamma4\Bigr)\mu+O(\mu^2)\qquad[\mathrm{L1};\ \mathrm{L2}];\qquad\text{the duality rows}\ [\mathrm{BM.131}];
\]

%dsp-level:0
\[
\textbf{the two actions meet only through the dictionary:}
\]

%dsp-level:0
\[
\qquad \text{the sector coefficients of}\ \mathrm{D2.1}\ \text{are the pair-state components}\ [\mathrm{CT1.1};\ \mathrm{L3}(1)];
\]

%dsp-level:0
\[
\qquad \text{the Haar face preserves their spectrum and the tilt face contracts it}\ [\mathrm{BM.131}(3);\ \mathrm{L2}]:
\]

%dsp-level:0
\[
\boxed{\ \text{both actions preserve the derived coefficient bounds of}\ \mathrm{BM.151}(2)\ (\text{forward pointer})\ —\ \text{the state-side dynamics never leaves the formed locus}\ }
\]

%dsp-level:0
\[
(5)\quad\textbf{the locked face of the operator:}
\]

%dsp-level:0
\[
\qquad F_{D^-}=\bigoplus_{n\ge4}F_n
\]

%dsp-level:0
\[
\qquad [\mathrm{D4.1};\ \mathrm{A8}];
\]

%dsp-level:0
\[
\qquad D_{\mathrm{tick}}(n)=2^{n+1},
\]

%dsp-level:0
\[
\qquad O(n)=\max\bigl(0,\ 2^{n+1}-16\bigr),
\]

%dsp-level:0
\[
\qquad \text{the last admissible rung is}\ 3
\]

%dsp-level:0
\[
\qquad [\mathrm{CT0.5};\ \mathrm{CT0.6};\ \mathrm{CT4.1}];
\]

%dsp-level:0
\[
\text{machine:}\ O=(0,0,0,0,16,48,112,240,496)\ \ \text{at}\ n=(0,1,2,3,4,5,6,7,8);
\]

%dsp-level:0
\[
\qquad O(4)=16=C_{\mathrm{cap}}\ \ \text{exactly one further channel quantum per tick}
\]

%dsp-level:0
\[
\qquad [\mathrm{CT4.1}];
\]

%dsp-level:0
\[
(6)\quad\textbf{status:}
\]

%dsp-level:0
\[
\qquad \textbf{derived:}\ (1)\text{--}(5);
\]

%dsp-level:0
\[
\qquad \text{the census is complete:}\ \neg\text{an aspect outside}\ \mathrm{A8}\text{'s list,}\ \neg\text{a block outside}\ \mathrm{CT0.12}\text{'s tiling,}\ \neg\text{a rung beyond}\ 3\ \text{in the accessible mask}\ [\mathrm{FT.83}(2)\text{--}(3):\ \text{the}\ n\ge4\ \text{no-go}];
\]

%dsp-level:0
\[
\neg\text{a declaration},\qquad\neg\text{a P4 addition},\qquad\neg\text{a new dimensionful scale}\ [\mathrm{P1}],\qquad\neg\text{a measured value}\ [\mathrm{P5}].
\]

\textbf{Req.} A2, A7, A8 (Part 1); C8 (Part 1); D0.3, D0.8, D0.9, D1.2, D1.8, D2.1, D2.3, D2.12, D3.1, D3.5, D3.6, D4.1 (Parts 1–2); CT0.5, CT0.6, CT0.11, CT0.12, CT0.13, CT0.14, CT1.1, CT1.2, L0, L1, L2, L3(1) (Part 1); CT1.12, CT1.20, CT1.21, CT2.4, CT2.5, CT2.6, CT2.7, CT3.8, CT4.1 (Part 2); P1, P2, P5 (Parts 1–2); Def-H.32, FT.9, FT.26(3), FT.83(2), FT.83(3) (Part 3); BM.36(1), BM.52(1), BM.131(3), BM.139(2) (this part); A-BM.10 (§II.4, logged before use); exact integer censuses [math; machine: \texttt{checks/\allowbreak{}verify\_\allowbreak{}H\_\allowbreak{}continuum\_\allowbreak{}limit\_\allowbreak{}era1\_\allowbreak{}closure.py} §A]; cf. BM.151 (forward pointer: the coefficient bounds).

\textbf{Pf.}

%dsp-level:0
\[
(1),
\]

%dsp-level:0
\[
(2):\ \mathrm{A8},
\]

%dsp-level:0
\[
\ \mathrm{CT0.12},
\]

%dsp-level:0
\[
\ \mathrm{CT0.13},
\]

%dsp-level:0
\[
\ \mathrm{CT2.4},
\]

%dsp-level:0
\[
\ \mathrm{CT0.14},
\]

%dsp-level:0
\[
\ \mathrm{BM.36}(1)\text{'s rows re-assembled}\ [\text{machine: the dimension sums}\ 3+3+5+3+1+1=16,\ 16+1=17,\ 17+8=25,\ 1+5+1+3+8=18]
\]

%dsp-level:0
\[
(3):\ \text{each row is the cited tag's own derived map; the composition is read off the shared domains and codomains}\ [\text{machine: the locus census — every row's tag exists and carries the cited clause}];
\]

%dsp-level:0
\[
\qquad \text{the}\ \aleph\text{-domain row is A-BM.10's pointwise declaration}\ [\mathrm{BM.139}(2)]
\]

%dsp-level:0
\[
(4):\ \mathrm{A6},\ \mathrm{D0.3},\ \mathrm{L0}\text{'s rows re-cited};
\]

%dsp-level:0
\[
\qquad \text{the coefficient components:}\ \delta\rho=\tfrac14\sum_ar_aG_a,\ r_a=\mathfrak T(\delta\rho\,G_a)\ [\mathrm{L3}(1);\ \mathrm{CT1.1}];
\]

%dsp-level:0
\[
\qquad \text{the preservation rows:}\ [\mathrm{BM.131}(3);\ \mathrm{L2}]\ [\text{machine: the contraction factor}\ (1-\tfrac\gamma4)\le1\ \text{on}\ \gamma\in(0,1]]
\]

%dsp-level:0
\[
(5):\ O(n)=\max(0,2^{n+1}-16)\ \text{evaluated at}\ n=0..8\ [\text{machine: exact integers}];\qquad(6):\ \text{assembly}.\ \blacksquare
\]

\textbf{Phys.} The force operator is not a list: it is one object with three exact faces. Its algebra face tiles the sixteen dimensions of the pair arena with the six pair aspects and stacks the remaining three on the vertex registers; its mask face climbs the rung ladder six, sixteen, seventeen, twenty-five with the two derived increments; its dynamical face is the composed action that carries a registered graph to weights, to the grouped sum, to the effective action, to the masked resolvent's spectrum, to the source and the field equations, to the smooth target manifold. Beside it runs the state-side action — the Haar draws and the tilt — which touches the same data only through the dictionary, preserves the coefficient bounds, and never leaves the formed locus. The locked face closes the census: beyond rung three the tick demand outruns the channel capacity by an exact power of two, and that excess is the operator's own dark-saturation summand. Nothing is added, nothing declared, nothing measured.

\end{theorem}

\begin{theorem}[{\texorpdfstring{BM.146 — The continuum-limit function: the two variables, the five contributors, and the exact functional form}{BM.146 - The continuum-limit function: the two variables, the five contributors, and the exact functional form}}]\label{tag:BM-146}
\textbf{Thm.}

%dsp-level:0
\[
(1)\quad\textbf{the variables and the notation row:}
\]

%dsp-level:0
\[
\qquad \text{a cell's content is}\ E_{\mathrm{int}}=|R|=n_{\mathcal R}
\]

%dsp-level:0
\[
\qquad [\mathrm{D0.8};\ \mathrm{Def\text{-}BM.5}(4)];
\]

%dsp-level:0
\[
\qquad \text{its ledger flux is}\ S_{\mathfrak L}=E_{\mathrm{int}}\ln2
\]

%dsp-level:0
\[
\qquad [\mathrm{CT1.29};\ \mathrm{C9}];
\]

%dsp-level:0
\[
\text{its chart span enters through the envelope ratio}
\]

%dsp-level:0
\[
\qquad k=\frac{r}{\Delta x_{\mathrm B}(E)}=\frac{r\,E_{\mathrm{int}}\ln2}{\mathfrak h}
\]

%dsp-level:0
\[
\qquad [\mathrm{BM.88}(2)(\mathrm i);\ \mathrm{Def\text{-}BM.4}(1)];
\]

%dsp-level:0
\[
\boxed{\ \textbf{the notation row:}\ \ \mathfrak q\ \ \textbf{remains the dimension bookkeeper}\ [\mathrm{CT2.1}];\ \ \text{the flux variable is the ledger's own}\ S_{\mathfrak L}\ —\ \text{count and flux are one datum up to}\ \ln2\ }
\]

%dsp-level:0
\[
(2)\quad\textbf{the five major contributors (each derived at its own locus):}\qquad\begin{array}{l|l|l}\text{contributor} & \text{the exact row} & \text{the locus}\\\hline\text{the Born width} & \Delta x_{\mathrm B}=\mathfrak h/(E_{\mathrm{int}}\ln2),\quad\Delta x_{\mathrm B}\cdot m_{\mathrm{hid}}=1 & [\mathrm{Def\text{-}BM.4}(1);\ \mathrm{BM.13}(1)]\\\hline\text{the flicker density} & w^{(n)}_{\mathrm{mat}}=\frac1{D_n},\quad\langle E\rangle^{(n)}_{\mathrm{mat}}=\frac{2n}{D_n},\quad D_n=\frac{1-u^n}{1-u} & [\mathrm{BM.137}(1);\ \mathrm{BM.94};\ \mathrm{BM.105}]\\\hline\text{the Janus boundary} & Q=x_+(\hat J_+-P_{\mathcal V})+x_-(\hat J_--P_{\mathcal R}),\quad\operatorname{spec}Q=\{0,-1\},\quad p=1-4\varepsilon_q & [\mathrm{Def\text{-}BM.18};\ \mathrm{BM.120};\ \mathrm{BM.19}(2)]\\\hline\text{the Haar--Poisson lapse} & N(k,E)=\sqrt{1-\frac{G(E)}{4\pi k}}=\sqrt{1-\frac{\mathfrak hE_{\mathrm{int}}}{2r}} & [\mathrm{BM.131};\ \mathrm{BM.132}(2)]\\\hline\text{the locked sector} & O(n)=\max(0,2^{n+1}-16),\quad\langle\rho_{D^-}\rangle_\Gamma(p)=\sum_vO(n_v)\bigl(1-f_{\mathrm{adm}}(n_v,p)\bigr) & [\mathrm{CT4.1};\ \mathrm{CT4.5a}(1)]\end{array}
\]

%dsp-level:0
\[
(3)\quad\textbf{the continuum-limit function (the exact functional form):}
\]

%dsp-level:0
\[
\qquad \text{the chart read of a cell is the tuple}
\]

%dsp-level:0
\[
\boxed{\ \mathcal C(E_{\mathrm{int}},r)\ =\ \Bigl(g_{\mu\nu}(E_{\mathrm{int}},r),\ \ N(k,E),\ \ \mathcal T(k,E),\ \ w^{(n)}_{\mathrm{mat}},\ \ \Gamma(E,r)\Bigr)\ }
\]

%dsp-level:0
\[
g_{\mu\nu}=\eta_{\mu\nu}+\frac{\mathfrak hE_{\mathrm{int}}}{2r}\,\mathbb 1_4+O\Bigl(\frac{\mathfrak h^2E_{\mathrm{int}}^2}{r^2}\Bigr)
\]

%dsp-level:0
\[
\qquad [\mathrm{BM.43}(4);\ \mathrm{BM.132}(2)];
\]

%dsp-level:0
\[
\qquad N=\sqrt{1-\frac{\mathfrak hE_{\mathrm{int}}}{2r}}
\]

%dsp-level:0
\[
\qquad [\mathrm{BM.132}(2)];
\]

%dsp-level:0
\[
\mathcal T=\frac{\mathcal T_0}N=\frac{E_{\mathrm{int}}\,\mathfrak h}{\ln2}\Bigl(1-\frac{\mathfrak hE_{\mathrm{int}}}{2r}\Bigr)^{-1/2}
\]

%dsp-level:0
\[
\qquad [\mathrm{BM.115}(1);\ \mathrm{BM.132}(5)];
\]

%dsp-level:0
\[
\qquad \Gamma=\Gamma_{\mathrm{tr}}\,f\,N\,\mathsf g\,\mathcal N\,\chi
\]

%dsp-level:0
\[
\qquad [\mathrm{BM.116}(1);\ \mathrm{BM.133}(4)];
\]

%dsp-level:0
\[
\textbf{every component is an exact closed form in}\ \bigl(E_{\mathrm{int}},\ r,\ \mathfrak h,\ \ln2,\ \varepsilon_q\bigr)
\]

%dsp-level:0
\[
\qquad [\text{machine: the monomial census}];
\]

%dsp-level:0
\[
(4)\quad\textbf{the formed domain (the exact conditions):}
\]

%dsp-level:0
\[
\qquad \boxed{\ \mathcal C\ \ \text{is formed}\iff E_{\mathrm{int}}\ge2\ \ \wedge\ \ \varkappa\|q\|_{\mathrm{op}}<1\ \ \wedge\ \ E_{\mathrm{int}}\le\frac{2r}{\mathfrak h}\ \ \wedge\ \ n(v)\le3\ \ \text{on the accessible vertices}\ }
\]

%dsp-level:0
\[
(\mathrm C1)\ E_{\mathrm{int}}\ge2:\ \ \text{the sub-hop/materialization floor}\qquad[\mathrm{Census\ FT.28}(1);\ \mathrm{Def\text{-}H.4};\ \mathrm{BM.13}(3)];
\]

%dsp-level:0
\[
(\mathrm C2)\ \varkappa\|q\|_{\mathrm{op}}<1:\ \ \text{the metric's nondegeneracy}\qquad[\mathrm{CT1.12};\ \mathrm{CT2.6}];
\]

%dsp-level:0
\[
(\mathrm C3)\ E_{\mathrm{int}}\le2r/\mathfrak h:\ \ \text{the lapse-formedness condition, in the capacity form}
\]

%dsp-level:0
\[
\qquad [\mathrm{BM.132}(3);\ \mathrm{BM.148}(1);\ \text{machine: the cross-multiplied identity}\ r\ge\mathfrak hE/2];
\]

%dsp-level:0
\[
(\mathrm C4)\ n(v)\le3:\ \ \text{the accessibility bound}\ —\ \text{oversubscribed vertices pass their excess to}\ D^-
\]

%dsp-level:0
\[
\qquad [\mathrm{CT0.11};\ \mathrm{CT4.1};\ \mathrm{D1.11}];
\]

%dsp-level:0
\[
(5)\quad\textbf{the limit}\ \mathfrak h\to0\ \textbf{(the derived sense in which the limit exists):}
\]

%dsp-level:0
\[
\qquad \text{at fixed}\ (E_{\mathrm{int}},r):
\]

%dsp-level:0
\[
\qquad E_{\mathrm{crit}}(r)=\frac{2r}{\mathfrak h}\to\infty
\]

%dsp-level:0
\[
\qquad [\text{machine: the Archimedean row on exact rationals}];
\]

%dsp-level:0
\[
\implies\text{every finite-content cell becomes chart-formed:}
\]

%dsp-level:0
\[
\qquad \forall E_{\mathrm{int}}\ \exists\,\mathfrak h_0:\ \forall\mathfrak h<\mathfrak h_0:\ (\mathrm C3)\ \text{holds}
\]

%dsp-level:0
\[
\qquad [\mathfrak h_0=2r/E_{\mathrm{int}}\ \text{explicit}];
\]

%dsp-level:0
\[
\textbf{the refinement row:}
\]

%dsp-level:0
\[
\qquad \text{a cell of span}\ s\ \text{carries}\ k=\frac{sE_{\mathrm{int}}\ln2}{\mathfrak h}\ \ \text{envelopes, each with}\ \ge2\ \ \text{directed real edges}
\]

%dsp-level:0
\[
\qquad [\mathrm{Def\text{-}H.4};\ \mathrm{Census\ FT.28}(1);\ \text{machine: the monomial}]:
\]

%dsp-level:0
\[
\boxed{\ \#\{\text{real edges per cell}\}\ \ge\ \frac{2sE_{\mathrm{int}}\ln2}{\mathfrak h}\ \longrightarrow\ \infty\ \ \text{as}\ \ \mathfrak h\to0\ }\qquad[\text{the witness of}\ \mathrm{BM.151}(1)];
\]

%dsp-level:0
\[
(6)\quad\textbf{the universality preview:}
\]

%dsp-level:0
\[
\qquad \mathcal C\ \text{carries}\ \neg\text{an aspect label:}\ \ \text{the force content enters through}\ \bigl(T_{\mu\nu}\ \text{channel},\ \text{mask slot},\ \text{coupling profile}\bigr)\ \ \text{alone}
\]

%dsp-level:0
\[
\qquad [\mathrm{CT1.22};\ \mathrm{CT2.4};\ \mathrm{BM.150}(3)\ \text{forward pointer}];
\]

%dsp-level:0
\[
(7)\quad\textbf{status:}
\]

%dsp-level:0
\[
\qquad \textbf{derived:}\ (1)\text{--}(6);
\]

%dsp-level:0
\[
\qquad \neg\text{a declaration},
\]

%dsp-level:0
\[
\qquad \neg\text{a P4 addition},
\]

%dsp-level:0
\[
\qquad \neg\text{a cosmological coordinate}\ [\mathrm{Def\text{-}BM.17}(6)(\mathrm i)],
\]

%dsp-level:0
\[
\qquad \neg\text{a measured value}\ [\mathrm{P5}].
\]

\textbf{Req.} A8 (Part 1); C9, D0.4, D0.8 (Part 1); CT0.11, CT1.29, CT2.1 (Parts 1–2); CT1.12, CT2.6, CT4.1, CT4.5a(1) (Part 2); P1, P5 (Parts 1–2); Def-H.4, Def-H.32 (Part 3); Def-BM.4(1), Def-BM.17(6), Def-BM.18, BM.13(1), BM.13(3), BM.19(2), BM.88(2), BM.94, BM.105, BM.115(1), BM.116(1), BM.120, BM.131, BM.132(2), BM.132(3), BM.132(5), BM.133(4), BM.137(1) (this part); Census FT.28(1) (row-wise); monomial identities, exact rational limit rows [math; machine: \texttt{checks/\allowbreak{}verify\_\allowbreak{}H\_\allowbreak{}continuum\_\allowbreak{}limit\_\allowbreak{}era1\_\allowbreak{}closure.py} §B].

\textbf{Pf.}

%dsp-level:0
\[
(1):\ \text{the count/flux identity}\ S_{\mathfrak L}=n_{\mathcal R}\ln2=E_{\mathrm{int}}\ln2\ [\mathrm{CT1.29};\ \mathrm{D0.8};\ \text{machine: one monomial}];
\]

%dsp-level:0
\[
\qquad \text{the envelope-ratio substitution}\ k=rE\ln2/\mathfrak h\ [\mathrm{Def\text{-}BM.4}(1);\ \text{machine}]
\]

%dsp-level:0
\[
(2):\ \text{the five rows re-cited at their loci}\ [\text{machine: the citation census}];
\]

%dsp-level:0
\[
\qquad (3):\ \text{the substitutions}\ G
\]

%dsp-level:0
\[
= 2\pi E^2\ln2,
\]

%dsp-level:0
\[
\ \Delta x_{\mathrm B}
\]

%dsp-level:0
\[
= \mathfrak h/(E\ln2):\quad\frac{G}{4\pi k}
\]

%dsp-level:0
\[
= \frac{2\pi E^2\ln2}{4\pi}\cdot\frac{\mathfrak h}{rE\ln2}
\]

%dsp-level:0
\[
= \frac{\mathfrak hE}{2r}\ [\text{machine: the rational-function identity}]
\]

%dsp-level:0
\[
(4):\ (\mathrm C1)\text{--}(\mathrm C4)\ \text{re-cited from their own tags};
\]

%dsp-level:0
\[
\qquad (\mathrm C3)\ \text{equivalence:}\ 1-\frac{\mathfrak hE}{2r}
\]

%dsp-level:0
\[
\ge 0
\]

%dsp-level:0
\[
\iff r\ge\frac{\mathfrak hE}2
\]

%dsp-level:0
\[
\iff E\le\frac{2r}{\mathfrak h}\ [\text{machine: cross-multiplication on positive denominators}]
\]

%dsp-level:0
\[
(5):\ \mathfrak h_0=2r/E:\ \ \mathfrak h<\mathfrak h_0\implies E<2r/\mathfrak h\ [\text{machine: rational witnesses}];
\]

%dsp-level:0
\[
\qquad \text{the edge-count monomial}\ 2k=2sE\ln2/\mathfrak h\ [\text{machine}];
\]

%dsp-level:0
\[
\qquad (6),(7):\ \text{assembly}.\ \blacksquare
\]

\textbf{Phys.} The continuum limit stops being a slogan and becomes a function. It takes two ledger variables — the cell's content and its chart span — and returns five exact closed forms: the linearized metric, the lapse, the materialization threshold, the flicker density, and the width. Its domain is four derived conditions, and the third of them is a capacity: a cell of content $E$ is chart-readable out to radius $r$ exactly when $E\le 2r/\mathfrak h$. The limit $\mathfrak h\to0$ is the statement that this capacity outruns every finite content, while the cell's real-edge count grows without bound — the two halves of the continuum condition, both witnessed inside the ledger. The function carries no aspect label: force identity lives in the source, the mask and the coupling profile, never in the limit itself.

\end{theorem}

\begin{theorem}[{\texorpdfstring{BM.147 — The metric–energy–continuum map: the dictionary bijection, the closed form of the source read, and the target}{BM.147 - The metric-energy-continuum map: the dictionary bijection, the closed form of the source read, and the target}}]\label{tag:BM-147}
\textbf{Thm.}

%dsp-level:0
\[
(1)\quad\textbf{the metric is a map (the dictionary):}\qquad\bigl(a_{00},\ q_{ij},\ w_i\bigr)\in\mathcal E_{++}\oplus\mathcal E_{+-}\ \longmapsto\ g_{\mu\nu}=\Omega^2\bar g_{\mu\nu},\qquad\Omega=e^{\alpha_\Omega a_{00}},\qquad\bar g=\begin{pmatrix}-1&\varkappa N_j\\ \varkappa N_i&\delta_{ij}+\varkappa q_{ij}\end{pmatrix},\qquad N_i=\beta w_i\cdot\Omega^2\ [\mathrm{CT1.12};\ \mathrm{BM.34}(1)];
\]

%dsp-level:0
\[
7+3=10=\tfrac{4\cdot5}2:
\]

%dsp-level:0
\[
\qquad \boxed{\ \text{the dictionary is a bijection between the ten sector coefficients and the ten metric components}\ }
\]

%dsp-level:0
\[
\qquad [\mathrm{CT1.12};\ \text{machine: the dimension count}];
\]

%dsp-level:0
\[
\operatorname{sig}(g)=(-,+,+,+)\iff\varkappa\|q\|_{\mathrm{op}}<1\ \ \text{on the perturbative regime}
\]

%dsp-level:0
\[
\qquad [\mathrm{CT1.12}];
\]

%dsp-level:0
\[
\qquad (a_{00},q,w)=0\longmapsto\text{the flat metric}
\]

%dsp-level:0
\[
\qquad [\mathrm{BM.34}(2)];
\]

%dsp-level:0
\[
(2)\quad\textbf{the energy determines the map (the closed form on the dust source):}
\]

%dsp-level:0
\[
\qquad T_{\mu\nu}=\mathrm{diag}(\rho,0,0,0),
\]

%dsp-level:0
\[
\qquad \rho=\frac{m}{V_{\mathrm{eff}}},
\]

%dsp-level:0
\[
\qquad m=\frac{E_{\mathrm{int}}\ln2}{\mathfrak h}
\]

%dsp-level:0
\[
\qquad [\mathrm{BM.9}(4);\ \mathrm{BM.14};\ \mathrm{FT.44}(4)];
\]

%dsp-level:0
\[
G_{\mu\nu}+\Lambda g_{\mu\nu}=\kappa^2T_{\mu\nu}
\]

%dsp-level:0
\[
\qquad [\mathrm{CT1.23}];
\]

%dsp-level:0
\[
\qquad \nabla^2\Phi=4\pi G_N\rho,
\]

%dsp-level:0
\[
\qquad \Phi=-\frac{G_Nm}r,
\]

%dsp-level:0
\[
\qquad G_N=\frac{\mathfrak h^2}{4\ln2},
\]

%dsp-level:0
\[
\qquad G_Nm=\frac{\mathfrak hE_{\mathrm{int}}}4
\]

%dsp-level:0
\[
\qquad [\mathrm{BM.43}(4);\ \mathrm{P3};\ \mathrm{CT2.9}];
\]

%dsp-level:0
\[
\boxed{\ g_{00}=-\Bigl(1-\frac{\mathfrak hE_{\mathrm{int}}}{2r}\Bigr),\qquad g_{ij}=\Bigl(1+\frac{\mathfrak hE_{\mathrm{int}}}{2r}\Bigr)\delta_{ij},\qquad g_{0i}=0\ }
\]

%dsp-level:0
\[
\qquad [\text{machine: the substitution}\ r_s=\mathfrak hE_{\mathrm{int}}/2\ \text{in}\ \mathrm{BM.43}(4)\text{'s weak-field rows}];
\]

%dsp-level:0
\[
\textbf{the pure-trace row:}
\]

%dsp-level:0
\[
\qquad g_{\mu\nu}=\eta_{\mu\nu}+\frac{\mathfrak hE_{\mathrm{int}}}{2r}\,\mathbb 1_4\ +\ O\Bigl(\frac{\mathfrak h^2E^2}{r^2}\Bigr):
\]

%dsp-level:0
\[
\qquad \boxed{\ \textbf{the linearized self-field of the dust cell is a pure-trace (conformal) perturbation}\ }
\]

%dsp-level:0
\[
\text{its whole content sits in the conformal mode:}
\]

%dsp-level:0
\[
\qquad \Omega^2=1-\frac{\mathfrak hE_{\mathrm{int}}}{2r},
\]

%dsp-level:0
\[
\qquad a_{00}=\frac1{2\alpha_\Omega}\ln\Bigl(1-\frac{\mathfrak hE_{\mathrm{int}}}{2r}\Bigr),
\]

%dsp-level:0
\[
\qquad \hat q=0,
\]

%dsp-level:0
\[
\qquad w=0
\]

%dsp-level:0
\[
\qquad [\text{machine: the exact solve of}\ \gamma_{ij}=\Omega^2(\delta_{ij}+\varkappa q_{ij})\ \text{against}\ g_{ij}];
\]

%dsp-level:0
\[
\textbf{the energy–lapse identity:}
\]

%dsp-level:0
\[
\qquad N^2=\Omega^2=1-\frac{\mathfrak hE_{\mathrm{int}}}{2r}=1-\frac{G(E)}{4\pi k}
\]

%dsp-level:0
\[
\qquad [\mathrm{BM.132}(2);\ \text{machine: the rational-function identity of}\ (3)\text{'s rows}];
\]

%dsp-level:0
\[
(3)\quad\textbf{the inverse read (the energy from the metric):}
\]

%dsp-level:0
\[
\qquad R=\kappa^2(\rho-3p)+4\Lambda,
\]

%dsp-level:0
\[
\qquad p=0,\ \Lambda=0:
\]

%dsp-level:0
\[
\qquad RV=\kappa^2m
\]

%dsp-level:0
\[
\qquad [\mathrm{BM.15}(2);\ \mathrm{BM.20}(1)]:
\]

%dsp-level:0
\[
\boxed{\ m=\frac1{\kappa^2}\int_\Sigma R\,dV\ }
\]

%dsp-level:0
\[
\qquad [\mathrm{BM.20}(1)];
\]

%dsp-level:0
\[
\qquad \kappa^2\rho=\frac{2\pi\,\mathfrak h\,E_{\mathrm{int}}}{V},
\]

%dsp-level:0
\[
\qquad (\Delta x_{\mathrm B})^2\kappa^2\rho=2\pi E_{\mathrm{int}}^2\ln2=G
\]

%dsp-level:0
\[
\qquad [\mathrm{BM.9}(5);\ \mathrm{BM.16}(1);\ \text{machine: monomials}];
\]

%dsp-level:0
\[
\textbf{the map is invertible on-shell at the leading reads:}
\]

%dsp-level:0
\[
\qquad E_{\mathrm{int}}=\frac{\kappa^2\rho V}{2\pi\mathfrak h}=\sqrt{\frac{G}{2\pi\ln2}}
\]

%dsp-level:0
\[
\qquad [\text{machine: the two exact inverses}];
\]

%dsp-level:0
\[
(4)\quad\textbf{the continuum is the target:}
\]

%dsp-level:0
\[
\qquad M=\Gamma_{\mathcal R}/\!\!\sim,
\]

%dsp-level:0
\[
\qquad \chi([v])=(\tau_n,h_1,h_2,h_3),
\]

%dsp-level:0
\[
\qquad \text{transitions smooth with}\ |\det J|=1
\]

%dsp-level:0
\[
\qquad [\mathrm{CT2.6};\ \mathrm{D1.8}];
\]

%dsp-level:0
\[
\textbf{the target-space rows:}
\]

%dsp-level:0
\[
\qquad \text{the image of the dictionary is the Lorentzian metrics of signature}\ (-,+,+,+)\ \text{on the perturbative regime}\ [\mathrm{CT1.12}];
\]

%dsp-level:0
\[
\qquad \text{the dispersion target}\ E=N\sqrt{m^2+p^2}-N^ip_i\ [\mathrm{BM.34}(4)];
\]

%dsp-level:0
\[
\textbf{the map's full statement:}
\]

%dsp-level:0
\[
\qquad \boxed{\ \text{metric}\ =\ \text{dictionary}\circ\ \text{field-equation solve}\circ\ \text{source read}\ (E_{\mathrm{int}},r)\ \longmapsto\ \text{continuum}\ }
\]

%dsp-level:0
\[
\qquad [(1)+(2)+(4);\ \text{the nonlinear joint solve remains}\ \neg\text{constructed}\ [\mathrm{BM.43}(5)\ \text{re-cited}]];
\]

%dsp-level:0
\[
(5)\quad\textbf{status:}
\]

%dsp-level:0
\[
\qquad \textbf{derived:}\ (1)\text{--}(4);
\]

%dsp-level:0
\[
\qquad \neg\text{a declaration},
\]

%dsp-level:0
\[
\qquad \neg\text{a P4 addition},
\]

%dsp-level:0
\[
\qquad r_s\ \text{quarantined as}\ \mathrm{P3}\text{'s named combination},
\]

%dsp-level:0
\[
\qquad \neg\text{a measured value}\ [\mathrm{P5}].
\]

\textbf{Req.} A8 (Part 1); D1.8, D2.1 (Part 1); CT1.12, CT1.23, CT2.6, CT2.9 (Part 2); P3, P5 (Parts 1–2); FT.44(4) (Census); Def-BM.4, BM.9(4), BM.9(5), BM.14, BM.15(2), BM.16(1), BM.20(1), BM.34(1), BM.34(2), BM.34(4), BM.43(4), BM.43(5), BM.115(1), BM.132(2) (this part); the exact weak-field solve, monomial and rational-function identities [math; machine: \texttt{checks/\allowbreak{}verify\_\allowbreak{}H\_\allowbreak{}continuum\_\allowbreak{}limit\_\allowbreak{}era1\_\allowbreak{}closure.py} §C].

\textbf{Pf.}

%dsp-level:0
\[
(1):\ \mathrm{CT1.12}\text{'s bijection row and signature criterion re-cited}\ [\text{machine: the dimension count}];
\]

%dsp-level:0
\[
\qquad (2):\ \Phi=-G_Nm/r\ \text{with}\ G_Nm=\frac{\mathfrak h^2}{4\ln2}\cdot\frac{E\ln2}{\mathfrak h}=\frac{\mathfrak hE}4\ [\mathrm{P3};\ \mathrm{FT.44}(4);\ \text{machine}];
\]

%dsp-level:0
\[
\qquad g_{00}=-(1+2\Phi)=-\bigl(1-\frac{\mathfrak hE}{2r}\bigr),\quad g_{ij}=(1-2\Phi)\delta_{ij}=\bigl(1+\frac{\mathfrak hE}{2r}\bigr)\delta_{ij}\ [\mathrm{BM.43}(4);\ \text{machine}]
\]

%dsp-level:0
\[
\text{the pure-trace row:}\ h_{00}=+\frac{\mathfrak hE}{2r}=h_{11}=h_{22}=h_{33}\ [\text{machine: the component census}];
\]

%dsp-level:0
\[
\qquad \text{the conformal solve:}\ \Omega^2\bigl(\delta_{ij}+\varkappa q_{ij}\bigr)
\]

%dsp-level:0
\[
= \bigl(1+\frac{\mathfrak hE}{2r}\bigr)\delta_{ij}\ \wedge\ \Omega^2
\]

%dsp-level:0
\[
= 1-\frac{\mathfrak hE}{2r}:\quad\varkappa q_{ij}
\]

%dsp-level:0
\[
= \Bigl(\frac{1+\mathfrak hE/(2r)}{1-\mathfrak hE/(2r)}-1\Bigr)\delta_{ij}
\]

%dsp-level:0
\[
= \frac{\mathfrak hE/r}{1-\mathfrak hE/(2r)}\,
\]

%dsp-level:0
\[
\delta_{ij}\ [\text{machine: the rational-function identity}]
\]

%dsp-level:0
\[
(3):\ \mathrm{BM.15}(2),\ \mathrm{BM.20}(1),\ \mathrm{BM.9}(5),\ \mathrm{BM.16}(1)\text{'s monomials re-verified}\ [\text{machine}];
\]

%dsp-level:0
\[
\qquad \text{the two inverses:}\ \frac{\kappa^2}{2\pi\mathfrak h}\cdot\frac{2\pi\mathfrak hE}{V}\cdot V=E;
\]

%dsp-level:0
\[
\qquad \Bigl(\sqrt{\frac{G}{2\pi\ln2}}\Bigr)^2\cdot2\pi\ln2=G\ [\text{machine}]
\]

%dsp-level:0
\[
(4),(5):\ \text{assembly and the scope re-citations}.\ \blacksquare
\]

\textbf{Phys.} The three names in the directive are one chain. The metric is the dictionary: ten sector coefficients of the ledger's pair states, bijectively repackaged as ten metric components. The energy determines the map: run the composite's own dust source through the field equations and the linearized solve closes in one monomial — the perturbation is pure trace, entirely conformal, the lapse and the conformal mode carrying the same exact function of content over radius. The map inverts on-shell: the curvature flux reads the content back. The continuum is the target: the quotient manifold with its tick-and-hop charts, whose metric image is exactly the Lorentzian signatures on the perturbative regime. The full non-linear solve stays unconstructed, and nothing measured enters anywhere.

\end{theorem}

\begin{theorem}[{\texorpdfstring{BM.148 — The breaking threshold: the exact capacity of a chart cell}{BM.148 - The breaking threshold: the exact capacity of a chart cell}}]\label{tag:BM-148}
\textbf{Thm.}

%dsp-level:0
\[
(1)\quad\textbf{the threshold (three exact forms of one condition):}
\]

%dsp-level:0
\[
\qquad \text{the chart read of a cell is formed}\iff r\ \ge\ r_s(E)=\frac{\mathfrak hE_{\mathrm{int}}}2
\]

%dsp-level:0
\[
\qquad [\mathrm{BM.132}(3);\ \mathrm{BM.147}(2)]:
\]

%dsp-level:0
\[
\boxed{\ E_{\mathrm{crit}}(r)=\frac{2r}{\mathfrak h}\ }\qquad\boxed{\ S_{\mathrm{crit}}(r)=\frac{2r\ln2}{\mathfrak h}\ }\qquad\boxed{\ k_{\mathrm{crit}}(E)=\frac{E_{\mathrm{int}}^2\ln2}2=\frac{G(E)}{4\pi}\ }
\]

%dsp-level:0
\[
\textbf{the flux form:}
\]

%dsp-level:0
\[
\qquad \text{the map breaks}\iff S_{\mathfrak L}(\text{cell})>S_{\mathrm{crit}}(r)
\]

%dsp-level:0
\[
\qquad [S_{\mathfrak L}=E_{\mathrm{int}}\ln2;\ \mathrm{CT1.29};\ \text{machine: the cross-multiplied identity}];
\]

%dsp-level:0
\[
(2)\quad\textbf{the dimension census:}
\]

%dsp-level:0
\[
\qquad [E_{\mathrm{crit}}]=\frac{\mathfrak q^{-1/2}}{\mathfrak q^{-1/2}}=\mathfrak q^0\ (\text{a count});
\]

%dsp-level:0
\[
\qquad [S_{\mathrm{crit}}]=\mathfrak q^0\ (\text{nats});
\]

%dsp-level:0
\[
\qquad [r_{\mathrm{crit}}]=\mathfrak q^{-1/2}\ (\text{a length})
\]

%dsp-level:0
\[
\qquad [\mathrm{CT2.1};\ \text{machine: the exponent arithmetic}];
\]

%dsp-level:0
\[
(3)\quad\textbf{the integer row:}
\]

%dsp-level:0
\[
\qquad E_{\mathrm{int}}\in\mathbb Z\implies\text{the formed contents at radius}\ r\ \text{are}\ \{2,\dots,\lfloor2r/\mathfrak h\rfloor\}
\]

%dsp-level:0
\[
\qquad [\mathrm{BM.13}(3):\ E\ge2;\ \text{machine}];
\]

%dsp-level:0
\[
\textbf{the one-hop capacity:}
\]

%dsp-level:0
\[
\qquad r=\mathfrak h:\quad E_{\mathrm{crit}}(\mathfrak h)=2:
\]

%dsp-level:0
\[
\qquad \boxed{\ \text{a chart cell of one hop carries exactly the digon floor and nothing heavier}\ }
\]

%dsp-level:0
\[
\qquad [\mathrm{Census\ FT.28};\ \text{machine}];
\]

%dsp-level:0
\[
\textbf{the computed cells:}
\]

%dsp-level:0
\[
\qquad r_s=\Bigl(1,\ 3,\ 6,\ 12,\ 30\Bigr)\,\mathfrak h
\]

%dsp-level:0
\[
\qquad \text{at}\ E_{\mathrm{int}}\in\{2,6,12,24,60\}
\]

%dsp-level:0
\[
\qquad [\mathrm{BM.43}(1);\ \text{machine: the integer census}];
\]

%dsp-level:0
\[
(4)\quad\textbf{the envelope corollary (the second route to}\ \mathrm{BM.132}(3)\text{'s row):}
\]

%dsp-level:0
\[
\qquad r=\Delta x_{\mathrm B}(E)=\frac{\mathfrak h}{E\ln2}:
\]

%dsp-level:0
\[
\qquad E\le E_{\mathrm{crit}}\bigl(\Delta x_{\mathrm B}(E)\bigr)=\frac2{E\ln2}\iff E^2\ln2\le2:
\]

%dsp-level:0
\[
E_{\mathrm{int}}\ge2\implies E^2\ln2\ge4\ln2=\ln16>2
\]

%dsp-level:0
\[
\qquad \Bigl[\ln16>2\iff16>e^2:\quad e^2<\Bigl(\frac{876809}{322560}\Bigr)^2<8<16;\ \text{machine: the exact enclosure}\Bigr];
\]

%dsp-level:0
\[
\boxed{\ \textbf{no materialized cell is chart-formed at its own envelope:}\ \ \text{the capacity}\ \frac2{E\ln2}\ \text{is below every admissible}\ E\ }
\]

%dsp-level:0
\[
(5)\quad\textbf{the monotonicity and the limit:}
\]

%dsp-level:0
\[
\qquad E_{\mathrm{crit}}(r)\ \text{strictly increasing in}\ r;
\]

%dsp-level:0
\[
\qquad S_{\mathrm{crit}}(r)\ \text{strictly increasing};
\]

%dsp-level:0
\[
\qquad r_{\mathrm{crit}}(E)\ \text{strictly increasing in}\ E;
\]

%dsp-level:0
\[
\qquad k_{\mathrm{crit}}(E)\ \text{strictly increasing and quadratic}
\]

%dsp-level:0
\[
\qquad [\text{machine: the sign census on exact rationals}];
\]

%dsp-level:0
\[
\lim_{\mathfrak h\to0}E_{\mathrm{crit}}(r)=\infty\ \ \text{at fixed}\ r
\]

%dsp-level:0
\[
\qquad [\mathrm{BM.146}(5);\ \text{machine}];
\]

%dsp-level:0
\[
\qquad \boxed{\ \textbf{the breaking locus recedes without bound in the continuum limit}\ }
\]

%dsp-level:0
\[
(6)\quad\textbf{the quarantine rows:}
\]

%dsp-level:0
\[
\qquad r_s\ \ \text{enters only as}\ \mathrm{P3}\text{'s named combination:}\ \neg\text{a collapse reading,}\ \neg\text{a horizon identification,}\ \neg\text{black-hole physics}\ [\mathrm{P3};\ \mathrm{BM.132}(6);\ \mathrm{BM.115}(4)\ \text{discipline}];
\]

%dsp-level:0
\[
\text{the breaking locus is a domain boundary of the chart reads}\ —\ \text{the ledger side of the same boundary is}\ \mathrm{BM.149};
\]

%dsp-level:0
\[
(7)\quad\textbf{status:}\qquad\textbf{derived:}\ (1)\text{--}(6);\qquad\neg\text{a declaration},\qquad\neg\text{a P4 addition},\qquad\neg\text{a numerical}\ \mathfrak h\ [\mathrm{P1};\ \mathrm{P5}].
\]

\textbf{Req.} D0.8 (Part 1); CT1.29, CT2.1 (Parts 1–2); P1, P3, P5 (Parts 1–2); Def-BM.4, BM.13(3), BM.43(1), BM.88(3), BM.115(4), BM.132(3), BM.132(6), BM.146(5), BM.147(2) (this part); Census FT.28 (row-wise); the exact rational enclosure of $e^2$, cross-multiplied identities, integer censuses [math; machine: \texttt{checks/\allowbreak{}verify\_\allowbreak{}H\_\allowbreak{}continuum\_\allowbreak{}limit\_\allowbreak{}era1\_\allowbreak{}closure.py} §D].

\textbf{Pf.}

%dsp-level:0
\[
(1):\ r\ge\frac{\mathfrak hE}2
\]

%dsp-level:0
\[
\iff E\le\frac{2r}{\mathfrak h}
\]

%dsp-level:0
\[
\iff E\ln2\le\frac{2r\ln2}{\mathfrak h}\ [\text{machine: cross-multiplication on positive denominators}];
\]

%dsp-level:0
\[
\qquad k_{\mathrm{crit}}:\ \frac{G}{4\pi}=\frac{2\pi E^2\ln2}{4\pi}=\frac{E^2\ln2}2\ [\mathrm{BM.88}(3);\ \text{machine}]
\]

%dsp-level:0
\[
(2):\ [\mathfrak h]=\mathfrak q^{-1/2}\ [\mathrm{P1};\ \mathrm{CT2.1}]:\ \text{the three exponent rows}\ [\text{machine}];
\]

%dsp-level:0
\[
\qquad (3):\ r_s=\mathfrak hE/2\ \text{at}\ E\in\{2,6,12,24,60\}:\ (1,3,6,12,30)\,\mathfrak h\ [\text{machine: integers}]
\]

%dsp-level:0
\[
(4):\ E^2\ln2\le2\ \text{has no integer solution}\ \ge2:
\]

%dsp-level:0
\[
\quad 4\ln2=\ln16>2\iff16>e^2\ [\mathrm{BM.132}(3)\text{'s enclosure re-run};\ \text{machine}];
\]

%dsp-level:0
\[
\qquad (5):\ \text{the derivative signs on exact rationals};
\]

%dsp-level:0
\[
\qquad (6),(7):\ \text{assembly}.\ \blacksquare
\]

\textbf{Phys.} The continuum has a capacity, and the capacity is exact: a chart cell of radius $r$ carries at most $2r/\mathfrak h$ of internal count, equivalently $2r\ln2/\mathfrak h$ nats of ledger flux, before its own metric read leaves the formed domain. At one hop of radius the capacity is precisely the digon floor — the lightest materialized cell there is also the heaviest admissible one. No materialized cell fits inside its own envelope: the capacity at the Born width is below every admissible content, a second, purely counting route to the envelope row. And in the continuum limit the breaking locus recedes without bound, which is the exact sense in which the limit exists. The Schwarzschild combination stays quarantined; the boundary is a domain edge of the chart reads, and its ledger face is the next theorem.

\end{theorem}

\begin{theorem}[{\texorpdfstring{BM.149 — The locked sector at the breaking locus: the two faces of one capacity, the absorption, the membrane, and the complement theorem}{BM.149 - The locked sector at the breaking locus: the two faces of one capacity, the absorption, the membrane, and the complement theorem}}]\label{tag:BM-149}
\textbf{Thm.}

%dsp-level:0
\[
(1)\quad\textbf{the two faces of one capacity:}\qquad\begin{array}{l|l|l}\text{the face} & \text{the capacity} & \text{the locus}\\\hline\text{chart-side} & E_{\mathrm{int}}\le2r/\mathfrak h:\ \ \text{the metric resolution of a cell} & [\mathrm{BM.148}(1)]\\\hline\text{ledger-side} & k_v\le3\ \ \text{per accessible vertex};\quad D_{\mathrm{tick}}(n)\le C_{\mathrm{cap}}=16\ \ \text{per tick} & [\mathrm{CT0.11};\ \mathrm{CT0.5};\ \mathrm{CT0.6};\ \mathrm{D1.11}]\end{array}
\]

%dsp-level:0
\[
\textbf{the excess rows:}
\]

%dsp-level:0
\[
\qquad \text{chart-side excess: the read is}\ \neg\text{formed}\ [\mathrm{BM.132}(3)];
\]

%dsp-level:0
\[
\qquad \text{ledger-side excess:}\ O(n)=\max(0,2^{n+1}-16)>0\ \ \text{at}\ n\ge4\ \ \text{falls to}\ D^-\ [\mathrm{CT4.1};\ \mathrm{D4.1}];
\]

%dsp-level:0
\[
(2)\quad\textbf{the absorption (exact rows):}
\]

%dsp-level:0
\[
\qquad \text{rate}\ \propto O(n)
\]

%dsp-level:0
\[
\qquad [\mathrm{D4.2}];
\]

%dsp-level:0
\[
\qquad O(4)=16=C_{\mathrm{cap}}:\ \ \text{exactly one further channel quantum per tick}
\]

%dsp-level:0
\[
\qquad [\mathrm{CT4.1}];
\]

%dsp-level:0
\[
\langle\rho_{D^-}\rangle_\Gamma(p)=\sum_vO(n_v)\,\bigl(1-f_{\mathrm{adm}}(n_v,p)\bigr),
\]

%dsp-level:0
\[
\qquad f_{\mathrm{adm}}(n,p)=\sum_{k\le3}\binom nkp^k(1-p)^{n-k}
\]

%dsp-level:0
\[
\qquad [\mathrm{CT4.5a}(1):\ \text{an exact polynomial in}\ p\ \text{with integer coefficients}];
\]

%dsp-level:0
\[
\textbf{the anchors re-derived:}
\]

%dsp-level:0
\[
\qquad \langle\rho_{D^-}\rangle_\Gamma(1)=\bigl(448,\ 896,\ 2976,\ 2240,\ 24384\bigr)
\]

%dsp-level:0
\[
\qquad \text{on the five Platonic skeletons}
\]

%dsp-level:0
\[
\qquad [\mathrm{CT4.5a}(2);\ \text{machine:}\ 4\,O(6),\ 8\,O(6),\ 6\,O(8),\ 20\,O(6),\ 12\,O(10)\ \text{with}\ O(m)=2^{m+1}-16];
\]

%dsp-level:0
\[
(3)\quad\textbf{the one-way membrane (re-derived):}
\]

%dsp-level:0
\[
\qquad \forall n\le3:\quad\mathcal P_M^{(n)}F_{D^-}\mathcal P_M^{(n)}=0;
\]

%dsp-level:0
\[
\qquad \neg\exists\ \text{a pole of}\ Z(\mathcal R,n)\ \ \text{carrying locked content}
\]

%dsp-level:0
\[
\qquad [\mathrm{CT4.2}];
\]

%dsp-level:0
\[
\text{locked content: physical norm}\ 0\ \ \text{at accessible rungs},
\]

%dsp-level:0
\[
\qquad w(C)\ne0\ \ \text{in the unmasked A8 sum},
\]

%dsp-level:0
\[
\qquad \neg\exists\ \text{a deleted entry}
\]

%dsp-level:0
\[
\qquad [\mathrm{CT4.3};\ \mathrm{CT2.8}];
\]

%dsp-level:0
\[
(4)\quad\textbf{the maximal locked entropy:}
\]

%dsp-level:0
\[
\qquad \text{the entropy triangle's third vertex:}
\]

%dsp-level:0
\[
\qquad D^-:\ \bigl(S_{\mathrm{raw}},S_{\mathfrak L}\bigr)=\bigl(\min_{\mathrm{acc}},\ \max\wedge\neg\text{accessible}\bigr)
\]

%dsp-level:0
\[
\qquad [\mathrm{CT4.7}];
\]

%dsp-level:0
\[
S_{\mathfrak L}=n_{\mathcal R}\ln2\ \ \text{is retained in full on absorption}
\]

%dsp-level:0
\[
\qquad [\mathrm{CT4.3}:\ \text{absorbed}\subseteq\mathfrak L;\ \mathrm{CT1.29}];
\]

%dsp-level:0
\[
\qquad \text{the accessible per-vertex content is bounded by}\ n(v)\ [\mathrm{D1.11}];
\]

%dsp-level:0
\[
(5)\quad\textbf{the dynamic emission:}
\]

%dsp-level:0
\[
\qquad \text{at the dynamic level the move system carries the emission and dissolution moves}\ N^\pm(v,S)\ \ \text{with formal weights}\ y_d,\ z_dx_-^{r(v)}
\]

%dsp-level:0
\[
\qquad [\mathrm{Def\text{-}H.32}(7);\ \mathrm{FT.108}];
\]

%dsp-level:0
\[
\textbf{the scope row:}
\]

%dsp-level:0
\[
\qquad \text{emission is a ledger move,}\ \neg\text{an accessible-mask return:}\ \ \text{re-entry to the accessible sector runs only through the A4 crossing and the register's moves}\ \hat J_\pm
\]

%dsp-level:0
\[
\qquad [\mathrm{CT4.2};\ \mathrm{Def\text{-}BM.18};\ \mathrm{BM.120};\ \text{machine: the citation census}];
\]

%dsp-level:0
\[
(6)\quad\textbf{the two scales:}
\]

%dsp-level:0
\[
\qquad \text{the locked sector gravitates as pressureless dust:}\ \ T^{(D^-)}_{\mu\nu}=\rho_{D^-}u_\mu u_\nu,
\]

%dsp-level:0
\[
\quad w_{D^-}=0,\quad\rho_{D^-}\ge0\ \bigl(=0\iff\neg\exists v:\ n(v)\ge4\bigr)
\]

%dsp-level:0
\[
\qquad [\mathrm{CT4.5}];
\]

%dsp-level:0
\[
\text{at the local scale the same sector is the membrane of}\ (3)
\]

%dsp-level:0
\[
\qquad [\mathrm{CT4.2}];
\]

%dsp-level:0
\[
\qquad \text{the junctions at}\ n\ge4\ \ \text{carry its lock content: retained,}\ \neg\text{accessible}
\]

%dsp-level:0
\[
\qquad [\mathrm{BM.2}(4);\ \mathrm{D1.11};\ \mathrm{Census\ FT.35}:\ \text{the first}\ D^-\text{-carrying floor}];
\]

%dsp-level:0
\[
(7)\quad\textbf{the complement theorem (the derived row of this tag):}
\]

%dsp-level:0
\[
\qquad \text{for a cell whose content exceeds the chart capacity}\ \bigl(E_{\mathrm{int}}>2r/\mathfrak h\bigr):
\]

%dsp-level:0
\[
(\mathrm i)\quad\text{the chart read is}\ \neg\text{formed}\ [\mathrm{BM.148}(1);\ \mathrm{BM.132}(3)];
\]

%dsp-level:0
\[
\qquad (\mathrm{ii})\quad\text{the ledger read stays exact and non-decreasing:}\ S_{\mathfrak L}=E_{\mathrm{int}}\ln2,
\]

%dsp-level:0
\[
\qquad |\mathfrak L(\tau_{n+1})|\ge|\mathfrak L(\tau_n)|
\]

%dsp-level:0
\[
\qquad [\mathrm{CT1.29};\ \mathrm{A6}];
\]

%dsp-level:0
\[
(\mathrm{iii})\quad\text{the accessible content is membrane-governed:}\ \neg\exists\ \text{an accessible mask returning locked content}\ [\mathrm{CT4.2};\ \mathrm{CT4.3}];
\]

%dsp-level:0
\[
\boxed{\ \textbf{what the metric cannot read, the ledger still retains:}\ \ \text{the breaking of the continuum map is}\ \neg\text{a loss of content}\ —\ \text{the content passes to the locked species, chart-side unreadable and ledger-side exact}\ }
\]

%dsp-level:0
\[
(8)\quad\textbf{status:}
\]

%dsp-level:0
\[
\qquad \textbf{derived:}\ (1)\text{--}(7);
\]

%dsp-level:0
\[
\qquad \neg\text{a declaration},
\]

%dsp-level:0
\[
\qquad \neg\text{a P4 addition},
\]

%dsp-level:0
\[
\qquad \neg\text{black-hole physics}\ [\mathrm{P3}\ \text{discipline};\ \mathrm{BM.148}(6)],
\]

%dsp-level:0
\[
\qquad \neg\text{a measured value}\ [\mathrm{P5}].
\]

\textbf{Req.} A6, A8 (Part 1); D0.9, D1.11, D4.1, D4.2 (Parts 1–2); CT0.5, CT0.6, CT0.11, CT1.29 (Part 1); CT2.8, CT4.1, CT4.2, CT4.3, CT4.5, CT4.5a, CT4.7 (Part 2); P3, P5 (Parts 1–2); Def-H.32(7), FT.108 (Part 3); Def-BM.18, BM.2(4), BM.120, BM.132(3), BM.148(1), BM.148(6) (this part); Census FT.35 (row-wise); exact integer and polynomial censuses [math; machine: \texttt{checks/\allowbreak{}verify\_\allowbreak{}H\_\allowbreak{}continuum\_\allowbreak{}limit\_\allowbreak{}era1\_\allowbreak{}closure.py} §E].

\textbf{Pf.}

%dsp-level:0
\[
(1):\ \text{the two capacity rows re-cited from their own tags};
\]

%dsp-level:0
\[
\qquad (2):\ O(n)=\max(0,2^{n+1}-16)\ \text{at}\ n=4..10\ [\text{machine: exact integers}];
\]

%dsp-level:0
\[
\qquad \text{the anchors:}\ O(6)=112,\ O(8)=496,\ O(10)=2032:
\]

%dsp-level:0
\[
\quad 4\cdot112=448,\ 8\cdot112=896,\ 6\cdot496=2976,\ 20\cdot112=2240,\ 12\cdot2032=24384\ [\text{machine}]
\]

%dsp-level:0
\[
(3),(4):\ \mathrm{CT4.2},\ \mathrm{CT4.3},\ \mathrm{CT4.7}\text{'s rows re-derived at their loci}\ [\text{machine: the citation pins}];
\]

%dsp-level:0
\[
\qquad (5):\ \mathrm{Def\text{-}H.32}(7)\text{'s moves with}\ \mathrm{CT4.2}\text{'s mask row}
\]

%dsp-level:0
\[
(6):\ \mathrm{CT4.5}\text{'s dust rows re-cited};
\]

%dsp-level:0
\[
\qquad (7):\ (\mathrm i)\text{--}(\mathrm{iii})\ \text{conjoined on the common domain}\ E>2r/\mathfrak h\ [\text{machine: the domain identity}];
\]

%dsp-level:0
\[
\qquad (8):\ \text{assembly}.\ \blacksquare
\]

\textbf{Phys.} The breaking of the continuum map is not a hole in the ledger — it is the locked sector's entrance. The same capacity theorem has two faces: the chart can resolve a cell only out to its content in hop units, and a vertex can expose only three incident quanta to an accessible mask; the excess on the first face is unreadable geometry, on the second it is absorption into the dark-saturation summand at an exact power-of-two rate, with the Platonic anchor values recomputed to the integer. The membrane is one-way: no accessible mask returns the content, no accessible pole carries it, yet nothing is deleted — the record survives with its full ledger entropy, at the maximal-locked vertex of the entropy triangle. Emission exists as a ledger move, but re-entry runs only through the A4 crossing. And the complement theorem states the whole design in one row: what the metric cannot read, the ledger still retains.

\end{theorem}

\begin{theorem}[{\texorpdfstring{BM.150 — The universality of the continuum limit: one operator, one arena, one action, three fingerprints}{BM.150 - The universality of the continuum limit: one operator, one arena, one action, three fingerprints}}]\label{tag:BM-150}
\textbf{Thm.}

%dsp-level:0
\[
(1)\quad\textbf{one operator:}
\]

%dsp-level:0
\[
\qquad \mathbf F\ \ \text{always applies in full, and its blocks are exhausted:}
\]

%dsp-level:0
\[
\qquad 3+3+5+3+1+1=16=C_{\mathrm{cap}}\ \ \text{on the pair arena};
\]

%dsp-level:0
\[
\qquad \mathrm{EM},\ \mathrm{strong},\ D^-\ \ \text{on the stack registers}
\]

%dsp-level:0
\[
\qquad [\mathrm{A8};\ \mathrm{BM.52}(1);\ \mathrm{CT0.12}];
\]

%dsp-level:0
\[
\neg\exists\ \text{an aspect outside the census}\ [\mathrm{CT0.13};\ \mathrm{FT.83}(2)\text{--}(3)];
\]

%dsp-level:0
\[
\qquad \neg\exists\ \text{a further block inside}\ \mathfrak H(\mathbb C^4)\ [\mathrm{CT0.12}:\ \text{the tiling is exact}];
\]

%dsp-level:0
\[
(2)\quad\textbf{one arena and one action:}
\]

%dsp-level:0
\[
\qquad M,\ \ \text{the charts}\ (\tau_n,h_1,h_2,h_3),\ \ \text{and the dictionary are aspect-blind:}
\]

%dsp-level:0
\[
\qquad \text{the weights depend on}\ (\Gamma,\ \text{registrations})\ \ \text{only}\implies\text{relabeling-invariant}\implies\text{diffeomorphism-invariant}
\]

%dsp-level:0
\[
\qquad [\mathrm{CT2.6};\ \mathrm{CT2.7}\ \mathrm{Pf}];
\]

%dsp-level:0
\[
\textbf{the Lovelock closure:}
\]

%dsp-level:0
\[
\qquad S_{\mathrm{geom}}=\frac1{2\kappa^2}\int\sqrt{-g}\,(R-2\Lambda)\ \ \text{with}\ \neg\exists\ \text{additional local higher-curvature terms in}\ D=4
\]

%dsp-level:0
\[
\qquad [\mathrm{CT2.7}:\ \text{the target dynamics is unique and force-independent}];
\]

%dsp-level:0
\[
(3)\quad\textbf{three fingerprints carry the force identity:}
\]

%dsp-level:0
\[
\qquad \boxed{\ \text{aspect content enters the continuum exactly through}\ \ \bigl(T_{\mu\nu}\text{'s channel},\ \ \text{the mask slot},\ \ \text{the coupling profile}\bigr)\ }
\]

%dsp-level:0
\[
(\mathrm i)\quad T_{\mu\nu}=T^{\mathrm{bd}}_{\mu\nu}+T^{\mathrm{tr}}_{\mu\nu}:
\]

%dsp-level:0
\[
\qquad \text{the bound channel}\ (\rho+p)u_\mu u_\nu+pg_{\mu\nu}\ \ \text{against the transit channel}\ \ \sum_a\bigl(F^{(a)}_{\mu\alpha}F^{(a)\ \alpha}_{\ \ \nu}-\tfrac14g_{\mu\nu}F^2\bigr)
\]

%dsp-level:0
\[
\qquad [\mathrm{CT1.22};\ \mathrm{D1.5}];
\]

%dsp-level:0
\[
(\mathrm{ii})\quad\text{the mask slot:}\ \ \mathbf F_M^{(n)}=\bigoplus_{A\in\mathcal A(n)}F_A|_{\mathcal K^{(A)}},
\]

%dsp-level:0
\[
\qquad \dim\mathcal K_n=(6,16,17,25)
\]

%dsp-level:0
\[
\qquad [\mathrm{CT2.4};\ \mathrm{CT0.14};\ \mathrm{BM.36}(1)];
\]

%dsp-level:0
\[
(\mathrm{iii})\quad\text{the coupling profile:}\ \ e(r,E)=e_0\bigl[Nf/(1+G)\bigr]^{1/2},
\]

%dsp-level:0
\[
\qquad (g,g')\ \ \text{family forms},
\]

%dsp-level:0
\[
\qquad \text{the channel constants}
\]

%dsp-level:0
\[
\qquad [\mathrm{BM.133}(6);\ \mathrm{BM.39}(1);\ \mathrm{BM.80}(1)];
\]

%dsp-level:0
\[
(4)\quad\textbf{the per-force rows (each a read of the same map):}\qquad\begin{array}{l|l|l}\text{aspect} & \text{its continuum read} & \text{the locus}\\\hline\text{gravity} & \text{the field equations themselves:}\ G_{\mu\nu}+\Lambda g_{\mu\nu}=\kappa^2T_{\mu\nu} & [\mathrm{CT1.23}]\\\hline\text{the uniform dark aspect} & F_{D^+}=\Lambda\mathbb I:\ \ \text{excluded from}\ Z_{\mathrm{force}},\ \ \text{carried by}\ \Lambda\ \ \text{in}\ S_{\mathrm{geom}} & [\mathrm{A8};\ \mathrm{D4.3};\ \mathrm{CT1.22}]\\\hline\text{EM} & \partial_\mu F^{\mu\nu}=eJ_Q^\nu;\ \ \text{the Coulomb read outside the support} & [\mathrm{Def\text{-}BM.12}(1)\text{--}(3);\ \mathrm{EM.15}]\\\hline\text{weak} & \text{the shift block:}\ \ N_i=\Omega^2\varkappa\beta w_i\ —\ \ \text{the weak aspect is the metric's own}\ g_{0i} & [\mathrm{FT.157}(1);\ \mathrm{CT1.12};\ \mathrm{BM.34}(3)]\\\hline\text{colour} & \text{sealed content only: the trace law}\ \mathfrak T(H)=3a;\ \ \neg\text{a triplet pole, physical norm}\ 0 & [\mathrm{CT3.22};\ \mathrm{BM.1}(2);\ \mathrm{FT.117}]\\\hline\text{Janus} & \partial_\mu J^\mu=0\ \ \text{with}\ \ J^0=\text{the seal density},\ J^i=\text{the hop current on the}\ w_i\ \text{carrier} & [\mathrm{BM.7};\ \mathrm{BM.11}(1)\text{--}(3)]\\\hline\text{the locked sector} & T^{(D^-)}_{\mu\nu}=\rho_{D^-}u_\mu u_\nu,\quad w_{D^-}=0 & [\mathrm{CT4.5}]\end{array}
\]

%dsp-level:0
\[
(5)\quad\textbf{the pp-wave witness (the consistency check that names no force):}
\]

%dsp-level:0
\[
\qquad \text{the nonlinear contracted Bianchi identity}\ \nabla_\mu G^{\mu\nu}
\]

%dsp-level:0
\[
= 0\ \ \text{holds on the exact rational pp-wave:}\ \ \det g
\]

%dsp-level:0
\[
= -1\ \ \text{identically, metric compatibility on all}\ 64\ \ \text{components, the}\ \Lambda\ \ \text{row}\ \nabla_\mu(\Lambda g^{\mu\nu})\qquad = 0
\]

%dsp-level:0
\[
\qquad [\mathrm{BM.11}(4);\ \text{machine re-cited}];
\]

%dsp-level:0
\[
\textbf{the witness is source-free:}
\]

%dsp-level:0
\[
\qquad \text{the identity is a property of the dictionary image,}\ \neg\text{of any aspect's content}\ —\ \text{it holds whatever filled}\ T_{\mu\nu}
\]

%dsp-level:0
\[
\qquad [\mathrm{CT1.23}\ \mathrm{Pf}\text{'s}\ \delta\text{-identity structure};\ \text{machine: the citation census}];
\]

%dsp-level:0
\[
(6)\quad\textbf{the state-side row:}
\]

%dsp-level:0
\[
\qquad \text{the Lyapunov contraction splits no class:}\ \ \Delta m_{\mathrm{Lyapunov}}=0\ \ \text{on every read species at once}
\]

%dsp-level:0
\[
\qquad [\mathrm{BM.51}(6):\ \text{every input is class-shared}];
\]

%dsp-level:0
\[
\qquad \text{the duality faces are aspect-blind}\ [\mathrm{BM.131}(3),(5)];
\]

%dsp-level:0
\[
(7)\quad\textbf{the universality theorem:}
\]

%dsp-level:0
\[
\qquad \boxed{\ \textbf{the continuum limit is one map with one domain law:}\ \ \text{force identity enters exactly through}\ (3)\text{'s three fingerprints and never through the arena, the action, the domain law, or the limit}\ }
\]

%dsp-level:0
\[
\textbf{the per-force consequence:}
\]

%dsp-level:0
\[
\qquad \text{the colour, weak, EM and gravitational reads share}\ \mathcal C\text{'s formed domain}\ [\mathrm{BM.146}(4)]\ \text{and}\ \mathrm{CT2.7}\text{'s unique action}\ —\ \text{what differs is the source channel, the mask slot, and the coupling profile, all three derived}
\]

%dsp-level:0
\[
\qquad [(4)\text{'s table}];
\]

%dsp-level:0
\[
(8)\quad\textbf{status:}
\]

%dsp-level:0
\[
\qquad \textbf{derived:}\ (1)\text{--}(7);
\]

%dsp-level:0
\[
\qquad \neg\text{a declaration},
\]

%dsp-level:0
\[
\qquad \neg\text{a P4 addition},
\]

%dsp-level:0
\[
\qquad \neg\text{a measured value}\ [\mathrm{P5}],
\]

%dsp-level:0
\[
\qquad \neg\text{an identification of any aspect with an external force law}\ [\mathrm{P5};\ \text{the names are the corpus's own}].
\]

\textbf{Req.} A8 (Part 1); D1.5, D4.3 (Parts 1–2); CT0.13, CT0.14, CT1.22, CT1.23, CT2.4, CT2.6, CT2.7, CT3.22, CT4.5 (Part 2); P5 (Parts 1–2); EM.15 (Part 2); FT.117, FT.157(1), FT.83(2), FT.83(3) (Part 3); Def-BM.12, BM.1(2), BM.7, BM.11(1), BM.11(3), BM.11(4), BM.34(3), BM.36(1), BM.39(1), BM.51(6), BM.52(1), BM.80(1), BM.131(3), BM.131(5), BM.133(6), BM.146(4) (this part); the citation censuses [machine: \texttt{checks/\allowbreak{}verify\_\allowbreak{}H\_\allowbreak{}continuum\_\allowbreak{}limit\_\allowbreak{}era1\_\allowbreak{}closure.py} §F].

\textbf{Pf.}

%dsp-level:0
\[
(1),(2):\ \text{the censuses re-run and}\ \mathrm{CT2.7}\text{'s uniqueness re-cited}\ [\text{machine: the integer sums}];
\]

%dsp-level:0
\[
\qquad (3):\ \text{the three fingerprints assembled from}\ \mathrm{CT1.22},
\]

%dsp-level:0
\[
\ \mathrm{CT2.4},
\]

%dsp-level:0
\[
\ \mathrm{BM.133}(6)\ \text{and the family forms}\ [\text{machine: the locus census}]
\]

%dsp-level:0
\[
(4):\ \text{each row re-cited from its own derived tag}\ [\text{machine: the table's citation pins}];
\]

%dsp-level:0
\[
\qquad (5):\ \mathrm{BM.11}(4)\text{'s witness rows re-cited; the source-freeness is the}\ \delta\text{-identity structure of}\ \mathrm{CT1.23}\text{'s Pf}\ [\text{machine: the citation census}]
\]

%dsp-level:0
\[
(6):\ \mathrm{BM.51}(6)\text{'s class-degenerate zero re-cited};\qquad(7),(8):\ \text{assembly}.\ \blacksquare
\]

\textbf{Phys.} Universality is not an assumption about forces agreeing — it is a structural census. There is one operator, its blocks tile the arena exactly, and the stack aspects sit on their own registers. There is one target manifold, one dictionary, and one geometric action, unique by the Lovelock count, with no term that could carry a force label. Force identity has exactly three homes: which channel of the stress tensor the aspect fills, which slot of the mask ladder it occupies, and which coupling profile it runs. The per-force table reads all seven aspects through the same map, and the pp-wave witness certifies the identity that no source ever signs: the contracted Bianchi row holds because of what the dictionary image is, not because of what fed it. The state-side dynamics agrees: the contraction splits no class. The continuum limit is the same limit for every force, and the differences are derived fingerprints, not different limits.

\end{theorem}

\begin{theorem}[{\texorpdfstring{BM.151 — The derivation of CT2.6(ii)–(iii): the unbounded ladder, the spectral bound, and the empty extra-axiomatic list}{BM.151 - The derivation of CT2.6(ii)-(iii): the unbounded ladder, the spectral bound, and the empty extra-axiomatic list}}]\label{tag:BM-151}
\textbf{Thm.}

%dsp-level:0
\[
(1)\quad\textbf{condition (ii) is a theorem of the ledger's own configuration space:}
\]

%dsp-level:0
\[
\qquad \text{(ii):}\ \ \text{coarse-graining cells contain unboundedly many real edges as}\ \mathfrak h\to0
\]

%dsp-level:0
\[
\qquad [\mathrm{CT2.6}\ \text{statement}];
\]

%dsp-level:0
\[
(\mathrm a)\quad\textbf{the configuration space:}
\]

%dsp-level:0
\[
\qquad \mathrm{A7}\ \text{admits}\ \Gamma=(V,E)\ \ \text{at every scale};
\]

%dsp-level:0
\[
\qquad R\subseteq E_{\mathrm{dir}}\ \ \text{is unconstrained in size};
\]

%dsp-level:0
\[
\qquad E_{\mathrm{int}}=|R|
\]

%dsp-level:0
\[
\qquad [\mathrm{A7};\ \mathrm{D0.10};\ \mathrm{D0.8};\ \mathrm{Def\text{-}H.32}(1)\text{--}(2)];
\]

%dsp-level:0
\[
(\mathrm b)\quad\textbf{the witness ladder:}
\]

%dsp-level:0
\[
\qquad \text{the cycle floors}\ \Sigma_L\ \ \text{exist at every}\ L\ \ \text{with}\ E_{\mathrm{int}}=L
\]

%dsp-level:0
\[
\qquad [\mathrm{FT.32}];
\]

%dsp-level:0
\[
\qquad \text{the polygon family is unbounded with injective mass/rank/null reads and}\ \neg\text{a combinatorial or geometric ceiling}
\]

%dsp-level:0
\[
\qquad [\mathrm{FT.45};\ \mathrm{FT.46}];
\]

%dsp-level:0
\[
\boxed{\ \forall n\ \ \exists\ \ \text{an admissible cell configuration with}\ E_{\mathrm{int}}\ge n\ }
\]

%dsp-level:0
\[
\qquad [\text{machine: the constructive witness}\ R=E_{\mathrm{dir}}(\Sigma_L)\ \text{at}\ L=2..8];
\]

%dsp-level:0
\[
(\mathrm c)\quad\textbf{the refinement row:}
\]

%dsp-level:0
\[
\qquad \text{at fixed chart span}\ s:
\]

%dsp-level:0
\[
\qquad k=\frac{s}{\Delta x_{\mathrm B}}=\frac{sE_{\mathrm{int}}\ln2}{\mathfrak h}\ \ \text{envelopes per cell}
\]

%dsp-level:0
\[
\qquad [\mathrm{Def\text{-}BM.4}(1);\ \mathrm{FT.44}(4)];
\]

%dsp-level:0
\[
\qquad \text{every materialized chain carries}\ \ge2\ \ \text{directed real edges}
\]

%dsp-level:0
\[
\qquad [\mathrm{Def\text{-}H.4};\ \mathrm{Census\ FT.28}(1)]:
\]

%dsp-level:0
\[
\#\{E_{\mathcal R}\ \text{per cell}\}\ \ge\ \frac{2sE_{\mathrm{int}}\ln2}{\mathfrak h}\ \xrightarrow[\ \mathfrak h\to0\ ]{}\ \infty
\]

%dsp-level:0
\[
\qquad [\text{machine: the Archimedean row on exact rationals}];
\]

%dsp-level:0
\[
(\mathrm d)\quad\textbf{the scope audit:}
\]

%dsp-level:0
\[
\qquad \text{the witness rows}\ (\mathrm a)\text{--}(\mathrm c)\ \text{are combinatorial and card content:}\ \neg\text{one of them uses}\ \mathrm{CT2.6}\text{'s conclusion}\ (M\ \text{a manifold})\ \text{as premise};
\]

%dsp-level:0
\[
\text{the}\ \mathrm{CT2.6}\text{-citations of}\ \mathrm{Def\text{-}BM.4},
\]

%dsp-level:0
\[
\ \mathrm{BM.13},
\]

%dsp-level:0
\[
\ \mathrm{FT.47}\ \text{are chart-arena inheritances}\ (\text{scope rows}),
\]

%dsp-level:0
\[
\ \neg\text{uses of}\ (\mathrm{ii})/(\mathrm{iii})\ \text{as premises}
\]

%dsp-level:0
\[
\qquad [\text{machine: the citation-pattern scan}];
\]

%dsp-level:0
\[
\text{the materialization hypotheses}\ (\mathrm{M1}/\mathrm{M2}\ \text{species})\ \text{remain declared spectral scopes at their own loci, untouched and distinct}
\]

%dsp-level:0
\[
\qquad [\text{A §0.2a integration note};\ \mathrm{Def\text{-}H.20}];
\]

%dsp-level:0
\[
(2)\quad\textbf{condition (iii) is a theorem of the pair-state spectral bound:}
\]

%dsp-level:0
\[
\qquad \text{(iii):}\ \ \text{sector coefficient fields}\ [\mathrm{D2.1}]\ \ \text{of bounded variation across neighbouring cells}
\]

%dsp-level:0
\[
\qquad [\mathrm{CT2.6}\ \text{statement}];
\]

%dsp-level:0
\[
(\mathrm a)\quad\textbf{the coefficients are state components:}
\]

%dsp-level:0
\[
\qquad \rho=\tfrac14I+\tfrac14\sum_ar_aG_a,
\]

%dsp-level:0
\[
\qquad r_a=\mathfrak T(\delta\rho\,G_a)\in\mathbb R,
\]

%dsp-level:0
\[
\qquad |r|^2=4\,\mathfrak T(\delta\rho^2)
\]

%dsp-level:0
\[
\qquad [\mathrm{L3}(1);\ \mathrm{CT1.1}:\ \mathfrak T(G_aG_b)=4\delta_{ab}];
\]

%dsp-level:0
\[
(\mathrm b)\quad\textbf{the spectral bound:}
\]

%dsp-level:0
\[
\qquad \operatorname{spec}(\delta\rho)\subset\bigl[-\tfrac14,\tfrac34\bigr],
\]

%dsp-level:0
\[
\qquad \textstyle\sum_i\mu_i=0
\]

%dsp-level:0
\[
\qquad [\mathrm{A1}:\ \rho\succeq0,\ \mathfrak T\rho=1;\ \mathrm{A5};\ \mathrm{L1}]:
\]

%dsp-level:0
\[
\max\ \textstyle\sum_i\mu_i^2=\tfrac34\ \ \text{attained exactly at the vertex orbit}\ \bigl(\tfrac34,-\tfrac14,-\tfrac14,-\tfrac14\bigr)
\]

%dsp-level:0
\[
\qquad [\text{machine: the vertex enumeration of the polytope}\ \{\sum\mu=0\}\cap[-\tfrac14,\tfrac34]^4\ \text{and convexity of}\ \sum\mu^2];
\]

%dsp-level:0
\[
\boxed{\ |r_a|\ \le\ |r|\ =\ 2\sqrt{\textstyle\sum_i\mu_i^2}\ \le\ 2\sqrt{\tfrac34}\ =\ \sqrt3\ }\qquad[\text{machine:}\ |r|^2\le3\ \text{exact}];
\]

%dsp-level:0
\[
\textbf{the component chart:}
\]

%dsp-level:0
\[
\qquad \rho=\tfrac14\sum_{\mu\nu}c_{\mu\nu}\,\sigma_\mu\otimes\sigma_\nu,
\]

%dsp-level:0
\[
\qquad c_{\mu\nu}=\mathfrak T\bigl(\rho\,\sigma_\mu\otimes\sigma_\nu\bigr),
\]

%dsp-level:0
\[
\qquad |c_{\mu\nu}|\le\bigl\|\sigma_\mu\otimes\sigma_\nu\bigr\|_{\mathrm{op}}=1
\]

%dsp-level:0
\[
\qquad [\text{math: the state-functional bound};\ \text{machine: exact rational density-matrix samples}];
\]

%dsp-level:0
\[
\textbf{the}\ \mathrm{D2.1}\ \textbf{row:}
\]

%dsp-level:0
\[
\qquad \text{the sector coefficients}\ (a_{00},q_{ij},w_i)\ \text{are fixed rational projections of either chart}\ [\mathrm{D2.1}:\ \text{the basis row};\ \mathrm{CT1.2}]:
\]

%dsp-level:0
\[
\qquad \text{uniformly bounded by explicit rational constants,}\ \neg\text{free data};
\]

%dsp-level:0
\[
(\mathrm c)\quad\textbf{the fields are averages:}
\]

%dsp-level:0
\[
\qquad \text{coefficients}=\text{local ledger averages}
\]

%dsp-level:0
\[
\qquad [\mathrm{CT2.6}\ \mathrm{Pf}\ \text{row}];
\]

%dsp-level:0
\[
\qquad \text{convex combinations preserve the ball:}\ \|\text{field}\|\le\sqrt3\ \ \text{per coefficient};
\]

%dsp-level:0
\[
(\mathrm d)\quad\textbf{the variation bound:}
\]

%dsp-level:0
\[
\qquad \boxed{\ |\Delta\,\text{coefficient}|\ \le\ 2\sqrt3\ \ \text{across any two cells},\qquad\mathrm{TV}(m\ \ \text{neighbouring cells})\ \le\ 2\sqrt3\,m\ }
\]

%dsp-level:0
\[
\qquad [\text{machine:}\ (2\sqrt3)^2=12\ \text{exact};\ \text{the triangle inequality}];
\]

%dsp-level:0
\[
(\mathrm e)\quad\textbf{the tighter formed-locus census:}
\]

%dsp-level:0
\[
\qquad \varkappa\|q\|_{\mathrm{op}}<1\ [\mathrm{CT1.12}];
\]

%dsp-level:0
\[
\qquad N\in(0,1]\ [\mathrm{BM.132}(4)];
\]

%dsp-level:0
\[
\qquad w^{(n)}_{\mathrm{mat}}\in\bigl(\tfrac1n,1\bigr]\ [\mathrm{BM.137}(1)];
\]

%dsp-level:0
\[
\qquad p\in(0,1]\ [\mathrm{D0.4};\ \mathrm{BM.19}(2)];
\]

%dsp-level:0
\[
f\in[0,1)\ [\mathrm{BM.88}(2)(\mathrm{ii})];
\]

%dsp-level:0
\[
\qquad \mathsf g\in(0,1]\ [\mathrm{BM.133}(4)];
\]

%dsp-level:0
\[
\qquad \Gamma_{\mathrm{tr}}\in\bigl[0,\tfrac13\bigr]\ [\mathrm{BM.44}(2);\ \mathrm{BM.116}(2)];
\]

%dsp-level:0
\[
\qquad w_{\mathrm{mat}}\in\bigl[\tfrac13,1\bigr]\ \ \text{measure-free}\ [\mathrm{BM.105}(2)\text{--}(4)];
\]

%dsp-level:0
\[
\textbf{every monotone row's total variation equals its range:}
\]

%dsp-level:0
\[
\qquad \text{the exact variation widths are}\ \bigl(1,\ 1-\tfrac1n,\ 1,\ 1,\ 1,\ \tfrac13,\ \tfrac23\bigr)\ \ \text{on}\ \bigl(N,\ w^{(n)}_{\mathrm{mat}},\ p,\ f,\ \mathsf g,\ \Gamma_{\mathrm{tr}},\ w_{\mathrm{mat}}\bigr)
\]

%dsp-level:0
\[
\qquad [\text{machine: the monotone-BV identity on exact rationals}];
\]

%dsp-level:0
\[
(\mathrm f)\quad\textbf{the dynamics row:}
\]

%dsp-level:0
\[
\qquad \text{the Haar face preserves the coefficient spectrum and the tilt face contracts it:}
\]

%dsp-level:0
\[
\qquad \mu'=\bigl(1-\tfrac\gamma4\bigr)\mu+O(\mu^2),
\]

%dsp-level:0
\[
\qquad \bigl(1-\tfrac\gamma4\bigr)\le1
\]

%dsp-level:0
\[
\qquad [\mathrm{BM.131}(3);\ \mathrm{L2};\ \text{machine}]:
\]

%dsp-level:0
\[
\text{the bound ball}\ \{|\mu|\le\tfrac34\}\ \ \text{is invariant under both faces}\ —\ \text{the recorded dynamics never violates}\ (\mathrm b);
\]

%dsp-level:0
\[
(3)\quad\textbf{the derivation:}
\]

%dsp-level:0
\[
\qquad \boxed{\ \mathrm{CT2.6}(\mathrm{ii}),(\mathrm{iii})\ \ \textbf{are theorems at the stated loci:}\ \ (\mathrm{ii})\ \text{by}\ (1),\qquad(\mathrm{iii})\ \text{by}\ (2)\ }
\]

%dsp-level:0
\[
\textbf{the extra-axiomatic list of}\ \text{A §0.2a}\ \textbf{is empty};
\]

%dsp-level:0
\[
\qquad \mathrm{CT2.6}\text{'s statement and Pf stand unmodified in their conclusions}\ —\ \text{the conditions are now supplied by derivation rather than by declaration};
\]

%dsp-level:0
\[
\textbf{the in-place amendments:}
\]

%dsp-level:0
\[
\qquad \text{A §0.2a's list sentence and}\ \mathrm{CT2.6}\text{'s two condition labels carry the derivation markers}\ [\text{the FT.62/FT.66 patch pattern}];
\]

%dsp-level:0
\[
\qquad \text{the point-of-use declarations across the corpus stand as scope-inheritance markers over the same loci};
\]

%dsp-level:0
\[
(4)\quad\textbf{the packaging:}
\]

%dsp-level:0
\[
\qquad \neg\text{a declaration is issued by this tag};
\]

%dsp-level:0
\[
\qquad \neg\text{a P4 member};
\]

%dsp-level:0
\[
\qquad \mathfrak h\ \ \text{remains the single dimensionful scale}\ [\mathrm{P1}];
\]

%dsp-level:0
\[
\qquad \neg\text{a measured value}\ [\mathrm{P5}];
\]

%dsp-level:0
\[
\text{no tag loses content:}
\]

%dsp-level:0
\[
\qquad \text{every tag inheriting}\ \mathrm{CT2.6}\text{'s hypotheses inherits the same conditions, now derived at the same scope}\ [\text{machine: the inheritance census}];
\]

%dsp-level:0
\[
(5)\quad\textbf{the gap row:}
\]

%dsp-level:0
\[
\qquad \textbf{the continuum-limit gap}\ —\ \text{the extra-axiomatic entry of}\ \text{A §0.2a},
\]

%dsp-level:0
\[
\ \text{the corpus's last non-axiomatic import}\ —\ \textbf{is closed};
\]

%dsp-level:0
\[
\text{the finite-}\mathfrak h\ \ \text{residue is distinct and stands}\ [\text{the clean-bifurcation loci};\ \mathrm{FT.66}(3);\ \mathrm{BM.143}(7)];
\]

%dsp-level:0
\[
(6)\quad\textbf{status:}
\]

%dsp-level:0
\[
\qquad \textbf{derived:}\ (1),(2);
\]

%dsp-level:0
\[
\qquad \textbf{adjudicated:}\ (3)\text{--}(5);
\]

%dsp-level:0
\[
\qquad \boxed{\ \textbf{the axiom surface is now the entire non-derived surface of the corpus}\ }.
\]

\textbf{Req.} A1, A5, A7 (Part 1); D0.4, D0.8, D0.10 (Part 1); CT1.1, CT1.2, CT1.10, CT1.29, L1, L2 (Part 1); CT1.12, CT2.6, D2.1 (Part 2); P1, P5 (Parts 1–2); Def-H.4, Def-H.20, Def-H.32(1), Def-H.32(2), FT.32, FT.44(4), FT.45, FT.46, FT.66(3), L3(1) (Part 3); Def-BM.4(1), BM.13, BM.19(2), BM.44(2), BM.47(6), BM.88(2), BM.105, BM.116(2), BM.131(3), BM.132(4), BM.133(4), BM.137(1), BM.143(7) (this part); Census FT.28(1) (row-wise); vertex enumeration of a rational polytope, convexity of the quadratic form, the monotone-BV identity, Archimedean rows on exact rationals [math; machine: \texttt{checks/\allowbreak{}verify\_\allowbreak{}H\_\allowbreak{}continuum\_\allowbreak{}limit\_\allowbreak{}era1\_\allowbreak{}closure.py} §G].

\textbf{Pf.}

%dsp-level:0
\[
(1)(\mathrm b):\ \Sigma_L\ \text{at}\ L=2..8:\ \ \text{the unidirectional}\ L\text{-cycle with}\ R=E_{\mathrm{dir}}(\Sigma_L):
\]

%dsp-level:0
\[
\quad |R|=L\ [\mathrm{FT.32};\ \text{machine: the constructive witness on the cycle skeletons}];
\]

%dsp-level:0
\[
\qquad (\mathrm c):\ \frac{2sE\ln2}{\mathfrak h}
\]

%dsp-level:0
\[
\ge T\iff\mathfrak h\le\frac{2sE\ln2}T:\ \ \text{witness}\ \mathfrak h
\]

%dsp-level:0
\[
= \frac{2sE\ln2}{T+1}\ [\text{machine: exact rationals at}\ T=10^6]
\]

%dsp-level:0
\[
(2)(\mathrm b):\ \text{the polytope}\ \{\mu\in[-\tfrac14,\tfrac34]^4:\ \sum\mu_i=0\}\ \text{has the vertex orbit}\ \binom41\bigl(\tfrac34,-\tfrac14,-\tfrac14,-\tfrac14\bigr):\ \ \text{any vertex fixes three coordinates at bounds; two}\ +\tfrac34\text{'s force the fourth to}\ -\tfrac32<-\tfrac14\ [\text{machine: the exhaustive bound-pattern enumeration}];
\]

%dsp-level:0
\[
\textstyle\sum\mu^2\ \text{convex}\implies\max\ \text{at a vertex}:\ \ \tfrac9{16}+\tfrac3{16}=\tfrac34\ [\text{machine}];
\]

%dsp-level:0
\[
\qquad |r|^2=4\sum\mu_i^2\le3\ [\mathrm{L3}(1)];
\]

%dsp-level:0
\[
\qquad |r_a|\le|r|\le\sqrt3\ [\text{Cauchy--Schwarz on components};\ \text{machine:}\ r_a^2\le|r|^2]
\]

%dsp-level:0
\[
(2)(\mathrm d):\ |\Delta|\le|r_a|+|r_a'|\le2\sqrt3;
\]

%dsp-level:0
\[
\qquad (2\sqrt3)^2=12\ [\text{machine}];
\]

%dsp-level:0
\[
\qquad (2)(\mathrm e):\ \text{the range rows re-cited; the monotone-BV identity:}\ \mathrm{TV}
\]

%dsp-level:0
\[
= \sup-\inf\ \text{on monotone functions}\ [\text{machine: rational grids}]
\]

%dsp-level:0
\[
(1)(\mathrm d),
\]

%dsp-level:0
\[
(3)\text{--}(6):\ \text{the citation-pattern scan, the amendment pins, and the inheritance census}\ [\text{machine}].\ \blacksquare
\]

\textbf{Phys.} The corpus's last non-axiomatic import is dissolved, and both halves dissolve into structure the ledger already owned. The richness condition was never an outside gift: the ledger's own configuration space carries cells of every finite content — the cycle ladder witnesses each level constructively — and refining the hop length at fixed span drives every cell's real-edge count past any bound. The regularity condition was never an outside gift either: the sector coefficients are components of a positive, unit-trace pair state, whose spectrum lives in an exact interval; convexity pins the quadratic extremum at a single vertex orbit, the bound is the square root of three, fields are averages of bounded data, and every weight family of the corpus carries its own exact range, so total variation equals range width everywhere. Neither derivation touches the manifold conclusion it licenses, and the scope citations that once looked like dependencies are audited as arena inheritances. The extra-axiomatic list is empty; the axiom surface is the entire non-derived surface; the conditions remain in force at every point of use, now as theorems.

\end{theorem}

\medskip\noindent\textit{\texorpdfstring{BM.152 — The status of the Continuum-Limit Tier and the closure of the foundational era}{BM.152 - The status of the Continuum-Limit Tier and the closure of the foundational era}.}\ Structural-constraints record: see Appendix B, Section B.1.35 of the Addendum.\par\medskip

\phantomsection
\section[Closure exports of Part 4 (corpus terminus)]{\texorpdfstring{Closure exports of Part 4 (corpus terminus)}{Closure exports of Part 4 (corpus terminus)}}

%dsp-level:0
\[
\begin{array}{lll}
\mathrm{Def\text{-}BM.1}\ \text{sealed composites};\ \text{validity/read clauses} & \mathrm{Def\text{-}BM.2}\ \text{composite card grammar} & \mathrm{Def\text{-}BM.3}\ \text{the identification layer (declared)}\\
\mathrm{BM.1}\ \text{observability switch};\ \text{trace-only read} & \mathrm{BM.2}\ 4\mathbb Z\ \text{junction phase};\ \text{additive winding} & \mathrm{BM.3}\ C_3\ \text{forced at the seal};\ \delta_q,\ 3\delta_q=\delta_\ell\\
\mathrm{BM.4}\ \text{colour/weak invariance};\ \text{parity passage} & \mathrm{BM.5}\ 25\ \text{cells},\ 6\ \text{families};\ \text{conditioning law} & \mathrm{BM.6}\ 8\ \text{open construction specs}\\
\mathrm{BM.7}\ \#J^+{=}\#J^-;\ \text{telescoping};\ \partial^2{=}0;\ \neg U(1)\ \text{on}\ q_{\mathrm{EM}} & \mathrm{BM.8}\ \text{hop}\to\text{translation};\ C_3\to SO(3);\ \text{tick obstruction} & \mathrm{BM.9}\ T_{\mu\nu};\ \text{sign reconciliation};\ \text{dust source}\\
\mathrm{BM.10}\ \neg[\hat x,\hat p]{=}i\ \text{in}\ 16;\ \Delta x\Delta p\ge\tfrac12;\ \text{MT},\ \text{CN} & \mathrm{BM.11}\ \partial_\mu J^\mu{=}0;\ \nabla_\mu T^{\mu\nu}{=}0\ \text{two routes} & \mathrm{BM.12}\ \text{the Noether adjudication};\ \neg\text{Cal-BM parameter}\\
\mathrm{Def\text{-}BM.4}\ \Delta x_{\mathrm B}{=}1/m_{\mathrm{hid}};\ V_{\mathrm{eff}}{=}(\Delta x_{\mathrm B})^3 & \mathrm{BM.13}\ \Delta x_{\mathrm B}{=}\mathfrak h/(E\ln2);\ \text{sub-hop}\ E\ge2 & \mathrm{BM.14}\ \rho{=}(E\ln2)^4/\mathfrak h^4;\ [\rho]{=}\mathfrak q^{+2}\\
\mathrm{BM.15}\ R{=}\kappa^2(\rho{-}3p){+}4\Lambda;\ \kappa^2\rho{=}2\pi E^4(\ln2)^3/\mathfrak h^2 & \mathrm{BM.16}\ (\Delta x_{\mathrm B})^2\kappa^2\rho{=}2\pi E^2\ln2;\ \Delta x_{\mathrm B}r_s{=}2G_N & \mathrm{BM.17}\ E_{\mathrm{int}}^{(n)}{=}6:\ \text{the neutron row}\\
\mathrm{BM.18}\ \text{the calibration adjudication} & \neg\text{a new dimensionful scale}\ [\mathrm{P1}] & \neg\text{a measured value}\ [\mathrm{P5}]\\
\mathrm{Def\text{-}BM.5}\ \mathcal E(\mathcal X);\ p{=}x_+/(x_+{+}x_-);\ E_{\mathrm{int}}(R){=}|R| & \mathrm{Def\text{-}BM.6}\ \text{the family naming rows};\ Q{=}\tfrac23 & \mathrm{BM.19}\ p{=}1{-}4\varepsilon_q;\ \Gamma{=}0\ \text{at}\ \varepsilon_q{=}0\\
\mathrm{BM.20}\ RV{=}\kappa^2m;\ \langle m\rangle{=}\text{the averaged flux};\ m^2_B(q,S) & \mathrm{BM.21}\ \mathcal B{\ne}0\ \text{in the ensemble},\ 0\ \text{in the count} & \mathrm{BM.22}\ F_{LR}{=}qI_2;\ 2|q|\in\{0,2\}\\
\mathrm{BM.23}\ \Delta q{=}\pm1;\ \text{the forced carrier};\ \text{the resolvent rate} & \mathrm{BM.24}\ \text{the flicker adjudication} & \neg\text{Cal-BM parameter};\ \neg\text{a P4 addition}\\
\mathrm{Def\text{-}BM.7}\ \mathrm{EM}^\pm;\ \text{the EM scar};\ \text{the formal-mirror scope} & \mathrm{BM.25}\ \mathrm{CPT}^2{=}\mathrm{id};\ w{\mapsto}w^\ast;\ \lambda\in\{\pm1,\pm i\} & \mathrm{BM.26}\ \text{the proton row};\ 2|q|{=}2\\
\mathrm{BM.27}\ \text{the antiproton row};\ \text{one orbit of size two} & \mathrm{BM.28}\ W(0){=}20,\ W({\pm}1){=}15;\ \text{three mass routes} & \mathrm{BM.29}\ \text{the sector adjudication};\ \neg\text{inter-sector dynamics}\\
\mathrm{Def\text{-}BM.8}\ Q{=}\textstyle\sum_iq_{\mathrm{core}}(c_i);\ \text{BM.6}(1)\ \text{discharged} & \text{(blank)} & \text{(blank)}\\
\mathrm{BM.30}\ \text{ratios}\ 1,\ \tfrac43,\ \tfrac{58}{13};\ \neg\text{a route adopted} & \neg\text{a fractional charge}\ [\mathrm{BM.28}(7)] & \neg\text{Cal-BM parameter};\ \neg\text{a measured value}\ [\mathrm{P5}]\\
\mathrm{Def\text{-}BM.9}\ \mathcal E_{\mathrm{ext}};\ \Pi_{\mathcal C};\ g(N){=}(1,1,3,10,25,105) & \mathrm{BM.31}\ g(3){=}1,\ g_\delta(3){=}0;\ N{=}3\ \text{protected} & \mathrm{BM.32}\ Z_\delta{=}0\ \text{by filtration};\ \langle m\rangle_{\mathrm{ext}}{=}\langle m\rangle_\varepsilon\\
\mathrm{BM.33}\ \text{the test's adjudication} & \text{the protection fails at}\ N\ge4 & \neg\text{a new declaration}\\
\mathrm{Def\text{-}BM.10}\ E{=}N\sqrt{m^2{+}\gamma^{ij}p_ip_j}{-}N^ip_i;\ \text{chart-side} & \mathrm{BM.34}\ (p_0{-}N^ip_i)^2{=}N^2(m^2{+}p^2);\ \mathrm{KE}{=}\tfrac{Np^2}{2m}{+}O(p^4);\ E^2\ge\tfrac74m_{\mathrm{hid}}^2 & \mathrm{BM.35}\ m_{\mathrm{rest}}{=}\Delta S_{\mathfrak L}/\mathfrak h{=}m_{\mathrm{hid}};\ \ln4{=}2\ln2\\
\mathrm{BM.36}\ (6,16,17,25);\ \text{ratios}\ (\tfrac13,8,\tfrac1{25});\ \text{Gram spectrum};\ \text{Yukawa-like coupling open} & \mathrm{BM.37}\ \text{routes}{=}\text{measure functionals};\ \text{alphabets disjoint};\ \text{trap closed} & \mathrm{BM.38}\ \text{the tier adjudication};\ \neg\text{Cal-BM parameter};\ \neg\text{a measured value}\\
\mathrm{Def\text{-}BM.11}\ c(i){=}\mathrm{Ad}(U^i)c(0);\ m_i^2{=}M^2(A_0{-}B_0\cos(\delta{-}\tfrac{2\pi i}3)) & \mathrm{BM.39}\ A_0{-}B_0{>}0\ (\text{SOS});\ \sum_im_i^2{=}3M^2A_0;\ \text{splitting}\ {\propto}\,gg' & \mathrm{BM.40}\ \sum_iW_i^2{=}2\Pi_{\mathcal R};\ \sigma\text{-vertex Casimir-neutral}\\
\mathrm{BM.41}\ \text{two-body impossible};\ \nu\text{-member forced};\ Q{=}{-}(m_\ell{+}m_\nu){<}0 & \mathrm{BM.42}\ \partial_\mu M^{\mu\nu\lambda}{=}0;\ p^2{=}p_r^2{+}L^2/r^2;\ 2^{\otimes4}{=}2{\cdot}0{\oplus}3{\cdot}1{\oplus}2 & \mathrm{BM.43}\ \nabla^2\Phi{=}4\pi G_N\rho;\ r_s{=}\tfrac{\mathfrak h E}2\\
\mathrm{BM.44}\ \Gamma_{\mathrm{tr}}\in[0,\tfrac13];\ \tau_{\mathrm{life}}\ge3\ \text{ticks};\ \neg\text{seconds} & \mathrm{BM.45}\ n\to p{+}e^-+\bar\nu_e;\ \text{grammar complete};\ \text{kinematics closed} & \neg\text{Cal-BM parameter};\ \neg\text{a measured value}\\
\mathrm{Def\text{-}BM.12}\ \partial_\mu F^{\mu\nu}{=}eJ^\nu_Q;\ E^{\mathrm{self}}{=}\tfrac{ce^2q^2}{4\pi\Delta x_{\mathrm B}};\ c\ \text{priced}\ (\tfrac12,\tfrac35) & \mathrm{BM.46}\ \Delta m_{\mathrm{EM}}{=}\tfrac{3ce^2\ln2}{2\pi\mathfrak h}{>}0;\ \text{the proton heavier} & \mathrm{BM.47}\ \text{UV: the envelope};\ \text{IR:}\ \mathrm{EM.11};\ f(k){=}1{-}\tfrac1k\\
\mathrm{BM.48}\ V{=}{-}\varepsilon_{\mathrm{EM}}/d;\ E_n{=}{-}\tfrac{\varepsilon_{\mathrm{EM}}^2m_r}{2n^2};\ \text{the bound channel}\ \neg\text{constructed} & \mathrm{BM.49}\ Q_{\mathrm{corr}}{<}0\ \text{deeper};\ E_{\max}{:}\langle E\rangle{=}2{:}1;\ \tfrac43\ \text{refuted} & \mathrm{BM.50}\ \text{the tier adjudication};\ \neg\text{a numerical}\ e\ [\mathrm{P5}]\\
\mathrm{BM.51}\ \bigl(\tfrac9{16}\bigr)^{6};\ V{=}8\Delta S_{\mathrm{pair}};\ \Delta m_{\mathrm{Lyapunov}}{=}0 & \mathrm{BM.52}\ \text{six sectors censused};\ \Delta m_{\mathrm{force}}{=}-\Delta m_{\mathrm{EM}} & \mathrm{BM.53}\ \Delta m_{\mathrm{metric}}{=}0;\ \Delta t\ge\pi\mathfrak h/(3ce^2\ln2)\\
\mathrm{BM.54}\ E^{\mathrm{info}}_{\mathrm{bind}}{=}0;\ \mathcal B^{(n)}{=}\mathcal B^{(p)};\ \text{the atom}\ \neg\text{constructed} & \mathrm{BM.55}\ \Delta m{=}-\tfrac{3ce^2\ln2}{2\pi\mathfrak h}{<}0;\ \tfrac{m_n}{m_p}{=}\tfrac{4\pi}{4\pi+ce^2} & \mathrm{BM.56}\ \text{the full-scope adjudication};\ \text{the order residue named},\ \neg\text{filled}\\
\mathrm{Def\text{-}BM.13}\ m_q{=}(2,1,1);\ W{=}(1,6,15,20,15,6,1);\ \Sigma_{\mathrm B}{=}\ln\tfrac43 & \mathrm{Def\text{-}BM.14}\ 3\ \text{floors}\ 2E{\in}\{12,24,60\};\ |\mathcal S|{=}3{\cdot}2{=}6;\ \text{the}\ D_S\ \text{row} & \mathrm{BM.57}\ |A_3|{=}3\ \text{vs}\ 1;\ \text{asymmetry}\ 2;\ \Delta m_{\mathrm{vertex}}{=}0\\
\mathrm{BM.58}\ \Delta m_{\mathrm{metric}}{=}0;\ \text{the seal asymmetry survives as}\ \Delta m_{\mathrm{EM}}{>}0 & \mathrm{BM.59}\ \Delta m_{\mathrm{transport}}{=}\ln(4/3)/\mathfrak h{>}0;\ \text{the residue filled on all four grounds} & \mathrm{BM.60}\ 1{:}2{:}5;\ D_S(\mathcal F_g,+){>}D_S(\mathcal F_g,-);\ \text{the inversion}\\
\mathrm{BM.61}\ \Sigma_{\mathrm B}\ \text{coupling-blind};\ \alpha_B\ \neg\text{derivable},\ (\mathrm O1)\text{--}(\mathrm O4) & \mathrm{BM.62}\ \Delta m{=}[\ln\tfrac43{-}6c\varepsilon_{\mathrm{EM}}\ln2]/\mathfrak h;\ \tfrac19{<}\varepsilon^{\mathrm{crit}}(\tfrac35){<}\tfrac18{<}\varepsilon^{\mathrm{crit}}(\tfrac12){<}\tfrac17 & \mathrm{BM.63}\ Q_\beta{<}-\ln3/\mathfrak h{<}0;\ \text{closed unconditionally}\\
\mathrm{BM.64}\ \text{the tier adjudication};\ m_n{>}m_p\ \text{sub-critically} & \neg\text{Cal-BM parameter};\ \neg\text{a new declaration beyond one name attachment} & \neg\text{a measured value}\ [\mathrm{P5}]\\
\mathrm{BM.65}\ \text{the baseline re-certified};\ \tfrac19{<}\varepsilon^{\mathrm{crit}}(\tfrac35){<}\tfrac18{<}\varepsilon^{\mathrm{crit}}(\tfrac12){<}\tfrac17 & \mathrm{BM.66}\ \operatorname{Tr}\widehat F{=}2c_0;\ I_F{=}4\widehat Q^2;\ \text{the operator rest input} & \mathrm{BM.67}\ U(\gamma){=}(-1)^{N_\pi};\ \operatorname{im}(\text{lift})\subseteq\{0,\pi\};\ \text{the obstruction record}\\
\mathrm{BM.68}\ \Omega{=}\Omega_{\mathrm{count}}\Omega_{\mathrm{charge}}\Omega_{\mathrm{geom}};\ \Omega_{\mathrm{geom}}^{\mathrm{stat}}{=}1;\ S_{\mathrm{joint}}{=}6\ln2{+}\ln W(Q) & \mathrm{BM.69}\ \tfrac{m^2_{\mathrm{pole}}}{m_{\mathrm{hid}}^2}{=}\rho_{\mathrm B}\tfrac{D_S}{4E^2};\ \text{Gap 7 minimal}:\ (\mathrm R1)\text{--}(\mathrm R3) & \mathrm{BM.70}\ F_{\ell\nu}{\ge}2\ln2;\ \text{formed}\implies\text{closed};\ \text{margins}\ \ln3,\ln12,\ln24\\
\mathrm{BM.71}\ \text{the five-route census};\ \Delta m_{\mathrm{bridge}}{=}0\ \text{in currency};\ \text{the decay closed} & \mathrm{BM.72}\ \text{the tier adjudication};\ \text{the partial bridge}\wedge\text{the no-go} & \neg\text{Cal-BM parameter};\ \neg\text{a new declaration};\ \neg\text{a measured value}\ [\mathrm{P5}]\\
\mathrm{BM.73}\ \text{A-BM.7 issued};\ (1,0)\ \text{occupancy};\ \hat\omega\hat\chi\hat\omega^{-1}{=}{-}\hat\chi & \mathrm{BM.74}\ \text{the ray-resolved spin census};\ j{=}|k|/2;\ \text{opposite rays at the vertex} & \mathrm{BM.75}\ \text{weight}\leftrightarrow\text{branch}\leftrightarrow\text{core unique};\ \Delta_g{=}(2,\sqrt2{-}1,2\cos\tfrac\pi{10}{-}\sqrt3)\hat\mu'^2\\
\mathrm{BM.76}\ \widehat V{=}E_+\otimes P_L;\ \widehat V^2{=}0;\ \operatorname{rk}\widehat V{=}1;\ [W_3\otimes I,\widehat V]{=}\widehat V & \mathrm{BM.77}\ \text{A-BM.8 issued};\ \Delta T{=}5\,m^2_B;\ \mathcal O_S\ne\varnothing;\ 828100{>}5819 & \mathrm{BM.78}\ \text{the tier adjudication};\ 1'\ \text{closed};\ \text{the decay allowed on}\ \mathcal O_S\\
\text{the decay closed on every pre-specification read} & \neg\text{Cal-BM parameter};\ \text{two addenda logged before use} & \neg\text{a measured value}\ [\mathrm{P5}]\\
\mathrm{BM.79}\ D(m){=}4\sin^2(m\delta_v/4){=}2{-}2T_{|m|}(\cos\tfrac{\delta_v}2);\ T_5(\mathsf c){=}0;\ r_{\mathrm{intra}}{=}\sin(k_f\delta_v/4)/\sin(\delta_v/4) & \mathrm{BM.80}\ 4z^3{-}24z^2{+}(21{-}12\gamma)z{-}(\gamma{-}1)^2;\ r^{(\ell)}_0{=}1{+}\gamma_q;\ q{<}\nu{<}\ell & \mathrm{BM.81}\ \Delta_1{>}\Delta_2{>}\Delta_3;\ \text{the doublet ratios};\ 3\ \neg\text{proportional ladders}\\
\mathrm{BM.82}\ \mathcal R_s{=}D_S/4E^2;\ \mathcal R_{\max}/\mathcal R_{\min}{=}50/(1{-}\mathsf c);\ 25{=}(30/6)^2 & \mathrm{BM.83}\ \rho_{\mathrm B},\sigma_\ell,\sigma_\nu;\ \ln1280/\ln960\notin\mathbb Q;\ t|_{m_n/m_p=4/3}{=}{-}\ln405/\ln16{<}0 & \mathrm{BM.84}\ m_n/m_p{<}21/20;\ \varepsilon^{\mathrm{crit}}\notin\mathbb Q;\ \text{the trichotomy}\\
\mathrm{BM.85}\ Q_\beta/m_n\le{-}\ln3/\ln1280;\ \text{two formed reads, disjoint alphabets} & \mathrm{BM.86}\ \text{the census adjudication};\ \neg\text{a declaration};\ \text{the ten open rows} & \neg\text{Cal-BM parameter};\ \neg\text{a measured value}\ [\mathrm{P5}]\\
\text{four currencies, zero conversions} & \text{item 7 stands with its residue named} & \text{the residue:}\ (\mathrm{R1})\text{--}(\mathrm{R3}),\ \sigma_\ell,\ \sigma_\nu,\ e\\
\mathrm{BM.87}\ c_{\mathrm{formed}}{=}\tfrac12\ \text{derived twice};\ \iota_{\mathrm{ball}}{=}\tfrac1{10}\ \text{void};\ \text{the selector, sign and monomial stand} & \mathrm{BM.88}\ e{=}e_0\sqrt{N\,f\,\mathsf g};\ f(k){=}1{-}\tfrac1k;\ G{=}2\pi E_{\mathrm{int}}^2\ln2;\ f(1){=}0 & \mathrm{BM.89}\ E^{\mathrm{self}}{=}\tfrac12Nw\mathsf g\,\varepsilon_{\mathrm{EM}}Q^2m_{\mathrm{hid}};\ w(0){=}1;\ w{=}1{-}4\varepsilon_q\ (\text{marked})\\
\mathrm{BM.90}\ w{=}w(\varepsilon_q)\ \text{branch-blind};\ \text{hypothesis (W) refuted on four grounds} & \mathrm{BM.91}\ \Delta m\,\mathfrak h{=}\ln\tfrac43{-}3Nw\mathsf g\,\varepsilon_{\mathrm{EM}}\ln2;\ Nw\mathsf g\,\varepsilon_{\mathrm{EM}}{\le}\tfrac18{\implies}m_n{>}m_p;\ m_n/m_p{<}\tfrac{10}9 & \mathrm{BM.92}\ Q_\beta\mathfrak h\le{-}\ln3\ \text{uniformly};\ \text{the specified read untouched}\\
\mathrm{BM.93}\ \text{the correction's adjudication};\ \neg\text{a declaration};\ (\mathrm{R4}),(\mathrm{R5})\ \text{named} & \neg\text{Cal-BM parameter};\ \neg\text{a measured coupling}\ [\mathrm{P5}] & \text{the residue:}\ (\mathrm{R1})\text{--}(\mathrm{R5}),\ e_0,\ \mathsf g,\ N\\
\mathrm{BM.94}\ w_{\mathrm{mat}}{=}1/(1{+}u{+}u^2)\in(\tfrac13,1];\ \langle E\rangle_{\mathrm{mat}}\in(2,6];\ \sigma/\langle E\rangle{<}\tfrac12 & \mathrm{BM.95}\ P(\text{fibre}\,|\,Q){=}\prod_im_{q_i}/W(Q);\ \ln\tfrac43{=}[\tfrac95\ln2{-}\ln3]{+}[\tfrac15\ln2] & \mathrm{BM.96}\ \delta_\nu{=}\tfrac18;\ G_2{>}G_1{>}A_0{>}G_0;\ m_\nu^{\mathrm{rest}}{>}0;\ (1,0);\ \neg\text{a seesaw-type mechanism}\\
\mathrm{BM.97}\ m^{\mathrm{rest}}{=}m_{\mathrm{pole}};\ \Delta_g^{\mathrm{crit}}{=}s(2M_0{+}s);\ \text{the specified read flicker-blind} & \mathrm{BM.98}\ \mathcal R{=}\tfrac{\pi\ln2}2\tfrac{\rho_{\mathrm B}D_S}{G}{=}\tfrac{\hat\mu'^2D_Sr^2}4(1{-}f)^2;\ \mathcal R_{\mathrm{mat}}/\mathcal R{=}D^2{<}9 & \mathrm{BM.99}\ N\mathsf g\,\varepsilon_{\mathrm{EM}}{\le}\tfrac18{\implies}\Delta m{>}0,\ {\ge}\tfrac5{12}{\implies}\Delta m{<}0;\ m_n/m_p{<}\ln80/\ln60{<}\tfrac{15}{14}\\
\mathrm{BM.100}\ \text{the tier's adjudication};\ \neg\text{a declaration};\ (\mathrm{R6})\ \text{named} & \neg\text{Cal-BM parameter};\ \neg\text{a measured value}\ [\mathrm{P5}] & \text{the residue:}\ (\mathrm{R1})\text{--}(\mathrm{R6}),\ e_0,\ \mathsf g,\ N,\ \rho_\ell,\ \rho_\nu\\
\mathrm{BM.101}\ Q^{\mathrm{inv}}{=}{-}Q^{\mathrm{fwd}};\ \varepsilon_\ast{=}(2{-}\log_23)/(3w)\in(\tfrac18,\tfrac5{12});\ P_{\max}{=}1{-}\tfrac3{2x_\ast} & \mathrm{BM.102}\ [\mathfrak h]{=}\mathfrak q^{-1/2},\ [\rho_{\mathrm B}]{=}\mathfrak q^0;\ \text{A-BM.9 declared, value}\ \neg\text{issued};\ \text{item 7 closed as an equation} & \mathrm{BM.103}\ W(Q){=}\binom6{Q+3};\ 4z^3{-}24z^2{+}(21{-}12\gamma)z{-}(\gamma{-}1)^2;\ \det(\lambda I{-}M)\ \neg\text{formed}\\
\mathrm{BM.104}\ G_{\mathrm{corr}}{=}G(1{+}a\mathsf g);\ \mathsf g_{\mathrm{fp}}{=}(1{+}G)/(1{-}aG);\ aG{=}\pi Nw\,\varepsilon_{\mathrm{EM}}Q^2E^2\ln2;\ \mathsf g(G)\ \neg\text{derived} & \mathrm{BM.105}\ w_{\mathrm{mat}}{=}\tfrac13\sum_ia_i\in[\tfrac13,1]\ \text{measure-free};\ S/(3(1{-}R));\ \text{hypothesis (W6) refuted} & \mathrm{BM.106}\ j{=}|\kappa|/2\ \text{with}\ j{=}\tfrac12,1\ \text{on one kind};\ \neg\text{an exchange operator};\ \text{deficit 11 stands};\ (\mathrm{R7})\ \text{named}\\
\mathrm{BM.107}\ \text{the six adjudications};\ \text{one declaration};\ (\mathrm{R4}'),(\mathrm{R7})\ \text{named} & \neg\text{a new dimensionful scale}\ [\mathrm{P1}];\ \neg\text{a measured value}\ [\mathrm{P5}] & \text{the residue:}\ (\mathrm{R1})\text{--}(\mathrm{R3}),\ (\mathrm{R4}'),\ (\mathrm{R5})\text{--}(\mathrm{R7}),\ \mathsf g,\ e_0,\ N\\
\mathrm{Def\text{-}BM.15}\ \mathrm{CP}:=\mathrm C\circ\mathrm P;\ c^{\mathrm{CP}};\ \text{the CP-pair};\ \text{the formal-mirror scope} & \mathrm{BM.108}\ \mu_{\mathrm{CP}}=e^{i\pi\Theta}=1\ \text{on}\ \Theta\in2\mathbb Z;\ \text{the forced conjugate};\ \neg\exists\tau^{-1} & \mathrm{BM.109}\ \mathcal W_{\text{pair}}=4\varepsilon_q(1{-}4\varepsilon_q)\le\tfrac14;\ Q_{\text{pair}}=0;\ \mathcal B_{\mathrm{info}}=0\\
\mathrm{Def\text{-}BM.16}\ \mathcal H_2=\mathcal H_\omega\otimes\Psi\ \text{(two chains)};\ \dim=4(|k_1|{+}1)(|k_2|{+}1);\ \hat P_{12}=\hat\mu_{\mathrm{CP}}\hat S_{12} & \mathrm{BM.110}\ g(2)=1:\ \text{the}\ \delta\text{-seal is unique};\ W_2^{\mathrm{CP}}(0)=W_2(0)=6;\ (\mathrm{R7})\ \text{discharged at}\ n=2 & \mathrm{BM.111}\ m_{\mathrm{hid}}^{\text{meson}}=2m_{\mathrm{hid}};\ \mathcal B_{\mathrm{info}}=\mathcal B_{\mathrm{pole}}=0;\ (\mathrm{R8})\ \text{named}\\
\mathrm{BM.112}\ \hat P_{12}^2=\mathrm{id};\ \operatorname{spec}\subseteq\{\pm1\};\ \operatorname{Tr}=d\,\mu_{\mathrm{CP}};\ \text{the orientation block is statistics-blind} & \mathrm{BM.113}\ \varepsilon=(-1)^{|k|}=(-1)^{2j};\ \textbf{deficit 11 closed};\ \mathrm{BM.6}(2)\ \text{discharged} & \neg\text{a spin-parity theorem}\ [\mathrm{BM.106}(1)\ \text{stands}]\\
\mathrm{Def\text{-}BM.17}\ \mathcal R=(n,\sigma,Q,\mathbf s,k);\ \mathcal T=\mathcal T_0/N;\ \Gamma=\Gamma_{\mathrm{tr}}f\,\tfrac{\varepsilon_{\mathrm{eff}}}{\varepsilon_{\mathrm{EM}}}\mathcal N\chi;\ \neg\text{a cosmological coordinate} & \mathrm{BM.114}\ (2{+}z{+}z^{-1})^n=(z{+}1)^{2n}/z^n;\ W_n(Q)=\binom{2n}{Q{+}n};\ \mathbb M,\ [\mathbb M:\mathbb Q]=16;\ 55\ \text{eigenvalues},\ 1\ \text{collision} & \mathrm{BM.115}\ \mathcal T_0=E_{\mathrm{int}}\mathfrak h/\ln2;\ \text{ladders}\ 2n..60n,\ \text{span}\ 30;\ \mathcal T=2r_s/N;\ \neg\text{a forbidden state}\\
\mathrm{BM.116}\ \Gamma_{\mathrm{tr}}\in[0,\tfrac13];\ (\mathrm Z1)\ \varepsilon_q{=}0,\ (\mathrm Z2)\ f{=}0,\ (\mathrm Z3)\ \mathcal N{=}0;\ \tau_{\mathrm{life}}\ge3\ \text{ticks} & \mathrm{BM.117}\ 20=\binom63\ \text{an orbit count},\ \neg\text{a decuplet};\ \lambda^{(n)}_{\max}=4n;\ \text{the octet}\ \neg\text{formed} & \mathrm{BM.118}\ \text{the tier adjudication};\ \neg\text{a declaration};\ (\mathrm{R8}),(\mathrm{R9})\ \text{named}\\
\text{four objectives, four derived verdicts} & \neg\text{a new dimensionful scale}\ [\mathrm{P1}];\ \neg\text{a measured value}\ [\mathrm{P5}] & \text{the residue:}\ (\mathrm{R1})\text{--}(\mathrm{R6}),\ (\mathrm{R7})\ \text{at}\ n\ge3,\ (\mathrm{R8}),(\mathrm{R9}),\ \mathsf g,\ e_0,\ N\\
\hline
\mathrm{Def\text{-}BM.18}\ \mathcal J=\operatorname{span}\{|\mathcal V\rangle,|\mathcal R\rangle\};\ \hat J_\pm;\ Q=x_+(\hat J_+{-}P_{\mathcal V})+x_-(\hat J_-{-}P_{\mathcal R});\ \mathcal T_{\mathcal J}:\ \dim=4(|k|{+}1) & \mathrm{BM.119}\ c{=}1\ \text{a convention};\ 1{-}v=\tfrac{m^2}{E(E+p)}{>}0;\ \text{the boundary is}\ \mathcal R/\mathcal V & \mathrm{BM.120}\ \hat J_\pm^2{=}0,\ \{\hat J_+,\hat J_-\}{=}\mathrm{id};\ \operatorname{spec}Q{=}\{0,-(x_+{+}x_-)\};\ \text{detailed balance derived}\\
\mathrm{BM.121}\ (\hat U^{2m})_{\varnothing\varnothing}{=}(x_+x_-)^m,\ (\hat U^{2m+1})_{\varnothing\varnothing}{=}0;\ \tfrac1{1-w^2x_+x_-};\ \text{pole}\ |w|\ge2 & \mathrm{BM.122}\ Q{=}0;\ \mathcal T_0{=}4\mathfrak h/\ln2\ \text{at the floor};\ \mathcal W{=}x_+^2\chi;\ \widehat V^2{=}0;\ \text{the two-photon bar} & \mathrm{BM.123}\ \mathcal W_{\mathrm{ann}}{=}x_-^2\chi;\ \tfrac{\mathcal W_{\mathrm{ann}}}{\mathcal W_{\mathrm{ch}}}{=}1\iff\varepsilon_q{=}\tfrac18;\ (\mathrm{R9})\ \text{closed at the channel level}\\
\mathrm{Def\text{-}BM.19}\ \Theta\ \text{vs}\ \Theta_R;\ \Delta(m){=}4\sin^2\tfrac{m\Theta_R}4{=}2{-}2T_{|m|}(\cos\tfrac{\Theta_R}2);\ \Delta(0){=}0\ \text{for even}\ k & \mathrm{BM.124}\ \Theta{=}2k{=}4j;\ \#\{m\}{=}2j{+}1;\ \Theta_R{=}2\pi k\ \text{impossible on dimension} & \mathrm{BM.125}\ 0\le m^2_{\mathrm{pole}}(m)\le4\hat\mu'^2;\ \hat\mu'^2\tfrac{\Theta_R^2}4{=}\hat\mu^2\ell_c^2G(c);\ \text{the product form}\ \neg\text{formed}\\
\mathrm{BM.126}\ \text{(W7) refuted on}\ (\mathrm O1)\text{--}(\mathrm O4);\ \text{the crossing is spin-blind};\ \text{the domain row is bound/transit} & \mathrm{BM.127}\ \varepsilon{=}(-1)^{|k|}\ \text{a ledger read};\ \text{the kind read separated} & \mathrm{BM.128}\ \dim\mathcal H_3{=}8\prod(|k_a|{+}1);\ \chi(\sigma){=}d^{c(\sigma)};\ \dim\wedge^3\mathbb C^2{=}0\ \text{vs}\ \dim\wedge^3\mathbb C^3{=}1;\ (\mathrm{R7})\ \text{closed at}\ n{=}3\\
\mathrm{BM.129}\ \varepsilon_M{=}{+}1\ \text{unconditionally};\ \varepsilon_B{=}(-1)^{k_1+k_2+k_3};\ \varepsilon(\Gamma){=}\prod_a\varepsilon(c_a) & \mathrm{BM.130}\ \text{the tier adjudication};\ \neg\text{a declaration};\ (\mathrm{R7}'),(\mathrm{R9}'),(\mathrm{R9}''),(\mathrm{R10})\ \text{named} & \text{the residue:}\ (\mathrm{R1})\text{--}(\mathrm{R6}),\ (\mathrm{R7}'),(\mathrm{R8}),(\mathrm{R9}'),(\mathrm{R9}''),(\mathrm{R10}),\ \mathsf g,\ e_0,\ N\\
\hline
\mathrm{BM.131}\ \mathrm{tw}(\rho){=}\tfrac{\mathfrak T(\Pi_{\mathcal R}\rho)}3\Pi_{\mathcal R}{+}\mathfrak T(\Pi_{\mathcal V}\rho)\Pi_{\mathcal V};\ \delta K{=}\gamma\rho{+}\log Z_\gamma I;\ \text{the common fixed point}\ \{\tfrac14I\};\ V(s){=}(1{-}\tfrac\gamma4)^{2s/N}V_0 & \mathrm{BM.132}\ N(k,E){=}\sqrt{1{-}\tfrac{G}{4\pi k}};\ \text{formed}\iff k\ge\tfrac{G}{4\pi}{=}\tfrac{E_{\mathrm{int}}^2\ln2}2;\ \text{the envelope}\ \neg\text{formed};\ N{\to}1 & \mathrm{BM.133}\ \mathsf g(G){=}\tfrac1{1+G}\ \text{selected};\ e{=}e_0[Nf/(1{+}G)]^{1/2};\ \text{the loop stays excluded}\\
\mathrm{Def\text{-}BM.20}\ R(z;\mathbf p);\ \text{the embedding and Fourier rows declared (A-BM.12)};\ \mathcal V{=}1\ \text{minimal} & \mathrm{BM.134}\ \text{the pole locus is the dispersion};\ \text{form-unique};\ d\mu_{\mathrm{shell}}{=}\tfrac{d^3p}{2E};\ \Phi_2{=}\tfrac{\sqrt{z^2-4N^2m^2}}{8\pi N^3z};\ \text{item 17 closed} & \mathrm{BM.135}\ |Q|{=}2:\ 12{+}2\ \text{states};\ |Q|{=}3:\ 2;\ \partial_\mu F^{\mu\nu}{=}eJ_Q^\nu;\ \text{Q10 closed};\ \text{deficit 10 closed}\\
\mathrm{BM.136}\ \varepsilon_{\sigma(i)\sigma(j)\sigma(k)}{=}\mathrm{sgn}(\sigma)\varepsilon_{ijk};\ \mathcal H_n{=}2^n\prod(|k_a|{+}1),\ \chi(\sigma){=}d^{c(\sigma)}\ \text{at every}\ n;\ \dim\operatorname{Inv}((\mathbb C^3)^{\otimes4}){=}0;\ (1,1,1){:}\ \varepsilon_B{=}{-}1,\ j_B{=}\tfrac32;\ 112/104\ \text{in B} & \mathrm{BM.137}\ w^{(n)}_{\mathrm{mat}}{=}\tfrac1{D_n};\ \mathcal B^{(2)}_{\mathrm{ens}}{=}{-}\tfrac{4u\ln2}{(1{+}u)\mathfrak h}{<}0;\ m_2^2{\equiv}0\ \text{by the pincer};\ 21\ \lambda^{(2)},\ \neg\text{a collision};\ (\mathrm{R8})\ \text{closed} & \mathrm{BM.138}\ \Psi{=}\Pi_k(\mathcal H_{\mathsf v}^{\otimes k});\ \operatorname{rk}\Pi_k{=}k{+}1;\ \text{the connector tensors alike};\ \text{item 19 closed}\\
\mathrm{BM.139}\ \text{A-BM.10}\ (\aleph\text{-domain, pointwise: ensembles literal, junctions E2});\ \text{A-BM.11}\ (\text{variant B});\ \text{A-BM.13}\ (\vee^2\ \text{locus});\ \text{item 16 closed} & \mathrm{BM.140}\ \delta_v^{\mathrm{dodeca}}{=}\tfrac\pi5\ \text{minimal, unique};\ x^4{-}48x^3{+}304x^2{-}512x{+}256;\ \text{five currencies censused};\ \text{deficit 12 adjudicated} & \mathrm{BM.141}\ J_3{=}\tfrac{p(1-p)(9-10p+2p^2)}{2D^2};\ J_2{\le}\tfrac18\ \text{at}\ p{=}\tfrac23;\ \sup J_3\in(\tfrac{93}{500},\tfrac3{16});\ (\mathrm{R6})\ \text{closed}\\
\mathrm{BM.142}\ (\hat J_+\otimes\hat J_+)^2{=}0;\ \mathcal M{=}e^2q^2x_+^2\chi;\ \sigma_{\mathrm{pair}},\sigma_{\mathrm{ann}}\ \text{exact symbolic},\ [\sigma]{=}\mathfrak q^{-1};\ \tfrac{\sigma_{\mathrm{ann}}}{\sigma_{\mathrm{pair}}}{=}(\tfrac{x_-}{x_+})^4\tfrac{z^2}{z^2-4N^2m^2};\ (\mathrm{R9}'),(\mathrm{R9}'')\ \text{closed} & \mathrm{BM.143}\ (\mathrm{R5})\ \text{void};\ (\mathrm{R2}),(\mathrm{R10})\ \text{registration data};\ (\mathrm{R4}')\ \text{void for BM};\ \text{C-BM.1--3 consolidated};\ (\mathrm{P5.1})\text{--}(\mathrm{P5.10})\ \text{named} & \mathrm{BM.144}\ \text{four declarations};\ \text{zero P4 additions};\ \neg\text{a retraction in}\ \mathrm{BM.1}\text{--}\mathrm{BM.130};\ \textbf{the era closed}\\
\text{the residue: every row dispositioned}\ [\mathrm{BM.143}(6)] & \text{for Part 5: the EM binding channel first}\ [\mathrm{BM.143}(7)] & \textbf{the Baryomesonic era is closed; the books are shelved}\\
\hline
\mathrm{BM.145}\ (3{+}3{+}5{+}3{+}1{+}1){=}16{=}C_{\mathrm{cap}};\ \dim\mathcal K_n{=}(6,16,17,25);\ \hat{\mathcal F}:\ (\Gamma_{\mathcal R},\text{registrations})\mapsto(M,g_{\mu\nu},T_{\mu\nu},\operatorname{spec});\ O(n){=}\max(0,2^{n+1}{-}16) & \mathrm{BM.146}\ \mathcal C(E,r){=}(g,N,\mathcal T,w_{\mathrm{mat}},\Gamma)\ \text{closed forms};\ \text{formed}\iff E\ge2\wedge\varkappa\|q\|{<}1\wedge E\le2r/\mathfrak h\wedge n\le3;\ \#\text{edges/cell}\ge2sE\ln2/\mathfrak h\to\infty & \mathrm{BM.147}\ g_{\mu\nu}{=}\eta_{\mu\nu}{+}\tfrac{\mathfrak hE}{2r}\mathbb 1_4{+}O(\cdot)\ \text{pure trace};\ \Omega^2{=}1{-}\tfrac{\mathfrak hE}{2r};\ m{=}\tfrac1{\kappa^2}\int R\,dV\\
\mathrm{BM.148}\ E_{\mathrm{crit}}(r){=}2r/\mathfrak h;\ S_{\mathrm{crit}}{=}2r\ln2/\mathfrak h;\ k_{\mathrm{crit}}{=}E^2\ln2/2;\ E_{\mathrm{crit}}(\mathfrak h){=}2;\ \text{no cell self-formed at its envelope} & \mathrm{BM.149}\ \langle\rho_{D^-}\rangle(1){=}(448,896,2976,2240,24384);\ \mathcal P_M^{(n)}F_{D^-}\mathcal P_M^{(n)}{=}0\ (n\le3);\ \text{what the metric cannot read, the ledger retains} & \mathrm{BM.150}\ \text{one operator, one arena, one action};\ \text{three fingerprints}\ (T\text{-channel},\ \text{mask slot},\ \text{coupling profile});\ \text{the pp-wave witness is source-free}\\
\mathrm{BM.151}\ |r_a|\le\sqrt3\ \text{exact};\ |\Delta|\le2\sqrt3;\ \mathrm{TV}\le2\sqrt3 m;\ \mathrm{CT2.6}(\mathrm{ii}),(\mathrm{iii})\ \textbf{reclassified as derived};\ \textbf{the extra-axiomatic list is empty} & \mathrm{BM.152}\ \text{zero declarations, one retired};\ \neg\text{a retraction in}\ \mathrm{BM.1}\text{--}\mathrm{BM.144};\ \textbf{Era 1 closed} & \textbf{the foundational era is closed; the corpus stands on its axioms alone}\\
\end{array}
\]

$\blacksquare$


\appendix
% Census.md — sanitized presentation layer (generated)
\chapter{\texorpdfstring{$\langle qu(b(i(t|t))ick)\rangle$ — \textbf{Census} — The Computed Cells (Category B) and Their Consumption}{Qubitick - **Census** - The Computed Cells (Category B) and Their Consumption}}

\textbf{Role.} The census annex of the era canon: the tags whose primary deliverable is a computed table, card, ladder, orbit set, multiplet family, or count over a named configuration space, and whose tabulated values are cited by downstream consumers. The structural tags (identities, mechanisms, covariance laws, no-gos, uniqueness and scope theorems — Category A) live in Parts 1–4; the two categories are disjoint and complete over the FT series. \textbf{Charter (the separation rule).} A tag is \textbf{Category B (census)} $\iff$ its primary deliverable is a computed data cell over a named configuration space and downstream consumers cite the tabulated values; \textbf{Category A (structural)} $\iff$ its primary deliverable is an identity, mechanism, covariance law, no-go, uniqueness, or scope theorem about the corpus's algebra, geometry, or dynamics — even when proved by exhaustive enumeration. Mixed tags are classified by primary deliverable; the census sides of Category-A tags (the cover-polynomial tables of FT.68, the frame catalogue of FT.93, the susceptibility census of FT.103) are consumed through those tags' own rows in Parts 3 — no B-classification is implied. \textbf{Dep.} Parts 1–3 (cited by tag at use in each cell's \textbf{Req}). Nothing in Parts 1–2 cites this file; Part 3 and Part 4 cite its cells row-wise. The tag-level citation graph is acyclic (every \textbf{Req} resolves to an earlier tag or a marked forward pointer); \texttt{checks/\allowbreak{}verify\_\allowbreak{}era\_\allowbreak{}structure.py} machine-verifies ownership, closure, and acyclicity. \textbf{Inventory.} 25 cells: FT.28, FT.32, FT.35, FT.42, FT.43, FT.44, FT.45, FT.49, FT.50, FT.51, FT.63, FT.64, FT.65, FT.69, FT.76, FT.78, FT.81, FT.83, FT.84, FT.96, FT.98, FT.99, FT.105, FT.114, FT.120. Format: Part 1 §0.4 (\textbf{Thm} · \textbf{Req} · \textbf{Pf} · \textbf{Phys}). \textbf{External-data quarantine.} The [EXT] audit rows (FT.114, FT.120; and the audit clauses of FT.77, FT.79, FT.87, FT.104 in Part 3) are quarantined comparison data: tested, never trusted; they enter no derivation [P1, P5; FT.73(5) pattern].

\phantomsection
\section[The cells]{\texorpdfstring{The cells}{The cells}}

\begin{theorem}[{\texorpdfstring{FT.28 — First computation: the floor-fermion particle card}{FT.28 - First computation: the floor-fermion particle card}}]\label{tag:FT-28}
\textbf{Thm.} ($\Sigma_{\min}$ card.) Let $\Sigma_{\min}=\{p,q\}$ with the materialized 2-cycle $c=(e_1,e_2)$, $e_1=(p,q)$, $e_2=(q,p)$, as its interior [Def-H.4 floor; FT.24: $\kappa:=\widehat{ds}^{\,2}(e_1)=\widehat{ds}^{\,2}(e_2)\in\mathbb Z\setminus\{0\}$], observed by a single external neighbour $o$ through one real edge ($E_{\mathrm{bd}}=1$, $\delta\tau\ge1$; FT.17 hypotheses, Def-H.35). Then:

%dsp-level:0
\[
(1)\quad I_{\mathrm{hid}}=E_{\mathrm{int}}\ln2=2\ln2,\qquad m_{\mathrm{hid}}=\frac{2\ln2}{\mathfrak h}\qquad[\mathrm{FT.17},\ \mathrm{Def\text{-}H.9}(2)];
\]

%dsp-level:0
\[
(2)\quad\Theta_{\mathrm{int}}=\Theta(c)=2\kappa\in2\mathbb Z,\qquad k=\kappa,\qquad\mathrm{Hol}_{\mathrm{int}}=+1\qquad[\mathrm{FT.22},\ \mathrm{FT.24},\ \mathrm{FT.25}];
\]

%dsp-level:0
\[
(3)\quad\mathcal A(\Sigma_{\min};C)\supseteq\{G,H,W\};\qquad\mathrm{EM}\in\mathcal A\iff f_p\ne0\ \vee\ f_q\ne0\ [\mathrm{D3.3}];
\]

%dsp-level:0
\[
\text{strong}\notin\mathcal A\ \text{necessarily}:\ n(q)
\]

%dsp-level:0
\[
= 2
\]

%dsp-level:0
\[
\implies q\ \text{carries no slot index}\ [\mathrm{D3.1},\ \mathrm{CT3.11}(1):\ \text{colour slots at}\ n{=}3\ \text{only}];
\]

%dsp-level:0
\[
\quad\mathrm{Colour}(\Sigma_{\min})
\]

%dsp-level:0
\[
= \{p\};
\]

%dsp-level:0
\[
\quad\pi_C(p)\subseteq\Sigma_{\min}\ \text{impossible}\ [\mathrm{A1}:\ \text{partners distinct}]\implies\mathcal C
\]

%dsp-level:0
\[
= 0\implies\text{block}
\]

%dsp-level:0
\[
\equiv 0\ [\mathrm{FT.19}]:\ \textbf{the single-contact floor fermion is structurally colourless}\ (\text{colour-visible variant}:\ \mathrm{Rmk}\ \text{below});
\]

%dsp-level:0
\[
(4)\quad\dim\mathcal K_{\Sigma_{\min}}=9\ (G{+}W{+}H\ \text{core})\ +\ 1\cdot[\mathrm{EM}\ \text{active}]\ \le\ 10;
\]

%dsp-level:0
\[
\qquad \text{general bound}\ \le18\ [\mathrm{FT.26}(3),\ \mathrm{CT0.14}];
\]

%dsp-level:0
\[
\quad 18\ \text{is attained only in the inside-sealed variant}\ [\mathrm{Rmk}\ \text{below}:\ 9{+}1{+}8];
\]

%dsp-level:0
\[
(5)\quad n(p)=3\ [e_1\ \text{out},\ e_2\ \text{in},\ \text{observer edge in}],\quad n(q)=2\ [e_1\ \text{in},\ e_2\ \text{out}]
\]

%dsp-level:0
\[
\qquad [\mathrm{D0.8}];
\]

%dsp-level:0
\[
\quad D_{\mathrm{tick}}(3)=2^{4}=16=C_{\mathrm{cap}},
\]

%dsp-level:0
\[
\quad O(2)=O(3)=0\ [\mathrm{CT0.5},\ \mathrm{CT0.6}]:\ \text{no oversubscription anywhere at the floor}\ [\mathrm{CT0.11}];
\]

%dsp-level:0
\[
(6)
\]

%dsp-level:0
\[
\quad q:\ \text{boundary-channel read; the closed core}\ (\partial c=0\ [\mathrm{D1.4}])\ \text{carries basepoint sign data}\ q_{\mathrm{core}}(c;v)
\]

%dsp-level:0
\[
= \tfrac12(\sigma_{e_1}-\sigma_{e_k})\in\{-1,0,+1\}\ [\mathrm{FT.40}:\ \text{ternary},\ \omega\text{-covariant}\ [\mathrm{D2.6}]];
\]

%dsp-level:0
\[
(7)\quad Z(\mathcal R,\Sigma_{\min})\ \text{poles}\iff1\in\operatorname{spec}\mathbf F_M^{(\Sigma_{\min})}\ [\mathrm{FT.26}];
\]

%dsp-level:0
\[
\qquad c\ \text{materializes}\implies\Theta(c)=2\kappa\ [\mathrm{FT.22},\ \mathrm{Def\text{-}H.16}].
\]

\textbf{Req.} FT.17 (Def-H.35), FT.19, FT.20, FT.22, FT.24, FT.25, FT.26, FT.27; A1; Def-H.4, Def-H.9, Def-H.10, Def-H.15, Def-H.16; D0.8, D1.4, D3.1, D3.3, D3.4, D3.8, D3.9; CT0.5, CT0.6, CT0.11, CT0.14 (Part 1); CT3.11, CT3.17 (Part 2); CT4.9 (Part 2).

\textbf{Pf.}

%dsp-level:0
\[
(1):\ E_{\mathrm{int}}=2\ (\text{both cycle edges have both endpoints in}\ \Sigma_{\min}\ [\mathrm{Def\text{-}H.12}]);
\]

%dsp-level:0
\[
\quad E_{\mathrm{bd}}=1\le4\cdot1\ \checkmark\ [\mathrm{FT.17}\ \text{hyp}]
\]

%dsp-level:0
\[
\mathrm{CT4.9}(1):\ S(p)
\]

%dsp-level:0
\[
\le 3\ln2\ \text{at}\ n{=}3,
\]

%dsp-level:0
\[
\ \text{channel-visible}\ 1\ \text{bit}\ <3\ \checkmark\ [\mathrm{FT.17}\ \text{Def-H.35}]
\]

%dsp-level:0
\[
\implies I_{\mathrm{hid}}
\]

%dsp-level:0
\[
= E_{\mathrm{int}}\ln2\qquad = 2\ln2;
\]

%dsp-level:0
\[
\qquad m_{\mathrm{hid}}=I_{\mathrm{hid}}/\mathfrak h\ [\mathrm{Def\text{-}H.9}(2)]
\]

%dsp-level:0
\[
(2):\ \Sigma_{\min}\setminus\partial\Sigma_{\min}
\]

%dsp-level:0
\[
= \{e_1,e_2\}
\]

%dsp-level:0
\[
= E(c)\ (\text{single chain, edge-disjoint trivially}):\ \Theta_{\mathrm{int}}
\]

%dsp-level:0
\[
= \Theta(c)
\]

%dsp-level:0
\[
= 2\kappa\ [\mathrm{FT.24}(1)];
\]

%dsp-level:0
\[
\qquad \mathrm{Hol}_{\mathrm{int}}=(-1)^{2\kappa}=+1,\ k=\Theta/2=\kappa\ [\mathrm{FT.25}(1),(3)]
\]

%dsp-level:0
\[
(3):\ \text{FT.20 rows}:\ G:\ I_{\mathrm{hid}}=2\ln2\ne0\ \checkmark;
\]

%dsp-level:0
\[
\quad H:\ \text{folded into}\ G\ [\mathrm{Def-H.36}]\ \checkmark;
\]

%dsp-level:0
\[
\quad W:\ \Theta_{\mathrm{int}}=2\kappa\ne0\ (\kappa\ne0)\ \checkmark\ [\text{criterion fires in the sound direction}];\quad\mathrm{EM}:\ \exists v:\ f_v\ne0\ [\mathrm{D3.3}]
\]

%dsp-level:0
\[
\text{strong}:\ n(q)
\]

%dsp-level:0
\[
= 2
\]

%dsp-level:0
\[
\implies q
\]

%dsp-level:0
\[
\in V_2:\ \neg\exists\ \text{slot index at}\ q\ [\mathrm{D3.1}:\ n{=}3\ \text{slots}=\text{colour};\ \mathrm{CT3.11}(1)]\implies\mathrm{Colour}(\Sigma_{\min})\subseteq\{p\}
\]

%dsp-level:0
\[
p\ \text{slot-carrying}\ (n(p)=3):\ \pi_C(p)\ne\varnothing\ \text{requires a distinct partner}\ [\mathrm{A1}];
\]

%dsp-level:0
\[
\ \text{the only candidate inside}\ \Sigma_{\min}\ \text{is}\ q,
\]

%dsp-level:0
\[
\ \text{slot-less}\implies\pi_C(p)\not\subseteq\Sigma_{\min}\ \forall C\implies\mathcal C(\Sigma_{\min};C)
\]

%dsp-level:0
\[
= 0\ [\mathrm{Def\text{-}H.10}]\implies\text{block}
\]

%dsp-level:0
\[
\equiv 0\ [\mathrm{FT.19}]
\]

%dsp-level:0
\[
\quad (\text{if}\ p\text{'s slot is colour-inert in}\ C:\ \mathrm{Colour}(\Sigma_{\min})=\varnothing\ 	ext{and the block vanishes on the no-content side — FT.27's case split})
\]

%dsp-level:0
\[
(4):\ \text{block dims}\ [\mathrm{CT0.12},\ \mathrm{CT0.14}]:\ 5(G)+3(W)+1(H)
\]

%dsp-level:0
\[
= 9;
\]

%dsp-level:0
\[
\quad\mathrm{EM}\ \text{active}:\ +1\ [\dim\mathcal K^{(\mathrm{EM})}=1]\implies
\]

%dsp-level:0
\[
\le 10;
\]

%dsp-level:0
\[
\quad\text{general bound}\ 18\ [\mathrm{FT.26}(3)];
\]

%dsp-level:0
\[
\quad +8\ [\dim\mathcal K^{(\mathrm{strong})}=8,\ \mathrm{CT3.8}]\ \text{only in the inside-sealed variant}\ [\mathrm{Rmk}]
\]

%dsp-level:0
\[
(5):\ n(p)=\partial^\leftarrow(p)+\partial^\rightarrow(p)=2{+}1=3,
\]

%dsp-level:0
\[
\quad n(q)=1{+}1=2\ [\mathrm{D0.8}:\ e_1,e_2\ \text{plus the observer edge at}\ p];
\]

%dsp-level:0
\[
\qquad \mathrm{CT0.5}:\ D_{\mathrm{tick}}(3)=2^{4}=16,\ D_{\mathrm{tick}}(2)=2^{3}=8;\quad\mathrm{CT0.6}:\ O(2)=O(3)=0;\quad\mathrm{CT0.11}:\ n\le3\ \checkmark
\]

%dsp-level:0
\[
(6):\ \partial c
\]

%dsp-level:0
\[
= t(c)-s(c)
\]

%dsp-level:0
\[
= 0\ [\mathrm{D1.4}:\ \text{closed}];
\]

%dsp-level:0
\[
\quad\text{core read}
\]

%dsp-level:0
\[
= \text{first/last edge registers at the basepoint}:\ q_{\mathrm{core}}
\]

%dsp-level:0
\[
= (\sigma_{e_1}-\sigma_{e_k})/2\ [\mathrm{D3.4}\ \text{on the ordered tuple},\ \mathrm{CT0.9}:\ \text{independent registers};\ \mathrm{EM.9}:\ \text{ternary, attained}];
\]

%dsp-level:0
\[
\quad\text{observer-facing read via transit channels}\ [\mathrm{Def\text{-}H.4},\ \mathrm{EM.9}]
\]

%dsp-level:0
\[
(7):\ \mathrm{Def\text{-}H.16}+\mathrm{FT.26};\qquad\mathrm{FT.22}\ \text{at materialized}\ c.\ \blacksquare
\]

\textbf{Rmk (the dark floor).} With no observer edge ($E_{\mathrm{bd}}=0$): $\partial\Sigma_{\min}=\varnothing\implies\gamma(o,\Sigma_{\min})=\varnothing$ [Def-H.1] — the floor fermion is dark by construction: $I_{\mathrm{hid}}=E_{\mathrm{tot}}=2\ln2$, no read channel [FT.3's indistinguishability]. The observer edge is what lifts the floor into the spectrum bookkeeping — consistent with Def-H.4's external characterization.

\textbf{Rmk (the colour-visible variant).} Give $q$ its own single-contact observer edge ($E_{\mathrm{bd}}=2\le4\,\delta\tau$, $\delta\tau\ge1$): then $n(p)=n(q)=3$, both vertices slot-carrying [D3.1], and two sealing phases exist [FT.20 strong row]: \textbf{(i) inside-sealed} — $p$ triplet, $q$ antitriplet contracted by $\delta^i{}_j$ inside $\Sigma_{\min}$ (meson-type [D3.9, CT3.17]): $\mathcal C=1$, strong block active, $\dim\mathcal K_{\Sigma_{\min}}=9+1+8=18$ — the maximal observable mask [FT.26]; \textbf{(ii) outside-sealed or dangling} — $\mathcal C=0$, strong block $\equiv0$ [FT.19]. Bookkeeping: with two contacts FT.17's single-observer hypothesis no longer applies; $I_{\mathrm{hid}}$ becomes observer-relative [Def-H.9(2), FT.20 $G$ row]: $E_{\mathrm{tot}}=(2{+}2)\ln2$, $E_{\mathrm{vis}}(o_1)=\ln2\implies I_{\mathrm{hid}}(o_1)=3\ln2$ (the other observer's boundary bit is hidden from $o_1$), $m_{\mathrm{hid}}(o_1)=3\ln2/\mathfrak h$. The smallest colour-observable fermion floor is this variant — the sealed meson-type pair [FT.27 strong row].

\textbf{Phys.} The first particle card of the program, every entry derived: the floor fermion carries informational mass $m_{\mathrm{hid}}=2\ln2/\mathfrak h$ — two bits of hidden connectivity — an even interior winding $2\kappa$ with trivial interior parity, the integer label $k=\kappa$ as its spin candidate, charge readable only at its boundary, no colour at all in the single-contact case — structurally, not contingently: its one slot-carrying vertex has no possible partner inside — and no oversubscription anywhere ($O(2)=O(3)=0$). What remains open on the card is open by name, not by hand-waving: the charge read for closed cores (Gaps 14), the $k\to j$ rule (Gaps 11), and the pole-mass/$I_{\mathrm{hid}}$ relation (Gaps 7).

\vspace{2pt}\noindent\rule{\textwidth}{0.4pt}\vspace{2pt}

\end{theorem}

\begin{theorem}[{\texorpdfstring{FT.32 — The cycle floor family (single-contact $\Sigma_L$)}{FT.32 - The cycle floor family (single-contact )}}]\label{tag:FT-32}
\textbf{Thm.} Let $\Sigma_L$ be a simple oriented $L$-cycle, $L\ge2$ [A1 excludes self-loops; D1.4], with materialized interior chain $c$ and a single observer contact at $v_1$ ($E_{\mathrm{bd}}=1\le4\,\delta\tau$, $\delta\tau\ge1$; FT.17 hypotheses, Def-H.35). Write $\kappa_e:=\widehat{ds}^{\,2}(e)$ for $e\in c$. Then:

%dsp-level:0
\[
(1)\quad E_{\mathrm{int}}=L,\qquad I_{\mathrm{hid}}=L\ln2,\qquad m_{\mathrm{hid}}=\frac{L\ln2}{\mathfrak h}\qquad[\mathrm{FT.17},\ \mathrm{Def\text{-}H.9}(2)];
\]

%dsp-level:0
\[
(2)\quad n(v_1)=3,\ \ n(v_i)=2\ (i\ge2)\ [\mathrm{D0.8}];\qquad O(2)=O(3)=0,\ \ D_{\mathrm{tick}}(3)=16=C_{\mathrm{cap}}\ [\mathrm{CT0.5},\ \mathrm{CT0.6}];
\]

%dsp-level:0
\[
(3)\quad\Theta_{\mathrm{int}}=\Theta(c)\in2\mathbb Z\ \text{if}\ c\ \text{materialized}\ [\mathrm{FT.22}];
\]

%dsp-level:0
\[
\qquad \mathrm{Hol}_{\mathrm{int}}=+1\ [\mathrm{FT.25}(1)];
\]

%dsp-level:0
\[
\qquad k=\Theta/2\in\mathbb Z;
\]

%dsp-level:0
\[
\qquad W\ \text{row fires}\iff\Theta\ne0;
\]

%dsp-level:0
\[
(4)\quad\textbf{structural colourlessness}:\ \mathrm{Colour}(\Sigma_L)\subseteq\{v_1\}\ [\mathrm{D3.1}:\ \text{colour slots at}\ n{=}3\ \text{only};\ \mathrm{CT3.11}(1)];
\]

%dsp-level:0
\[
\ \ \pi_C(v_1)\subseteq\Sigma_L\ \text{impossible}\ [\mathrm{A1}:\ \text{partners distinct}]\implies\mathcal C
\]

%dsp-level:0
\[
= 0\ \vee\ \mathrm{Colour}
\]

%dsp-level:0
\[
= \varnothing\implies\text{strong}\notin\mathcal A(\Sigma_L;C)\ \ \forall C;
\]

%dsp-level:0
\[
(5)\quad\mathcal A(\Sigma_L;C)\supseteq\{G,H\}\ [G:\ I_{\mathrm{hid}}\ne0;\ H:\ \mathrm{Def-H.36}];
\]

%dsp-level:0
\[
\quad\mathrm{EM}\in\mathcal A\iff\exists i:\ f_{v_i}
\]

%dsp-level:0
\[
\ne 0\ [\mathrm{D3.3}];
\]

%dsp-level:0
\[
\quad\dim\mathcal K_{\Sigma_L}
\]

%dsp-level:0
\[
= 9\ (+1\ \mathrm{EM})
\]

%dsp-level:0
\[
\le 10<18\ [\mathrm{FT.26},\ \mathrm{CT0.14}];
\]

%dsp-level:0
\[
(6)\quad\text{minimality ladder for}\ \Theta(c)\ [\mathrm{FT.22}'\text{s constraint}\ \Theta\in2\mathbb Z]:
\]

%dsp-level:0
\[
\qquad \text{uniform-}\kappa:\ L\ \text{odd}\Rightarrow\kappa_{\min}=2,\ \Theta_{\min}=2L;\quad L\ \text{even}\Rightarrow\kappa_{\min}=1,\ \Theta_{\min}=L;
\]

%dsp-level:0
\[
\text{all-nonnull per-edge-integral}\ [\text{scope}:\ \mathrm{FT.24}]:\ \Theta_{\min}=L+(L\ \mathrm{mod}\ 2);
\]

%dsp-level:0
\[
\qquad \text{null-permitting}:\ \Theta_{\min}=2\ (\kappa\ \text{profile}\ (1,1,0,\dots,0)\ [\mathrm{CT1.19a}:\ \kappa=0\iff\text{null}]);
\]

%dsp-level:0
\[
\text{scope}:\ \text{uniform/all-nonnull regimes presuppose per-edge}\ \kappa\in\mathbb Z\ —\ \text{derived at}\ L{=}2\ [\mathrm{FT.24}]\ \text{and on round-trip-closed edges}\ [\mathrm{FT.34}(4),\ \text{forward}],
\]

%dsp-level:0
\[
\ \text{conditional for bare}\ L
\]

%dsp-level:0
\[
\ge 3\ \text{cycles}\ [\text{Open Gaps 13 class (iii)}:\ \mathrm{FT.22}\ \text{alone permits}\ \kappa\in(2/L)\mathbb Z];
\]

%dsp-level:0
\[
\qquad \text{unconditional floor}:\ \Theta_{\min}=2\ \ \forall L\ (\text{rational uniform witness}\ \kappa\equiv2/L);
\]

%dsp-level:0
\[
(7)\quad\Pi(\Sigma_L)=L^2\,\Theta\,(\ln2)^2\;\qquad L=2\ \text{reproduces FT.28's card verbatim (consistency)}.
\]

\textbf{Req.} FT.17 (Def-H.35), FT.20, FT.22, FT.24, FT.25, FT.26, FT.27, FT.28; A1; Def-H.4, Def-H.9, Def-H.10, Def-H.15; D0.8, D1.4, D3.1, D3.3, D3.8, D3.9; CT0.5, CT0.6, CT0.11, CT0.13, CT0.14 (Part 1); CT1.19a (Part 2); CT3.11, CT3.17 (Part 2); CT4.9 (Part 2).

\textbf{Pf.}

%dsp-level:0
\[
(1):\ E_{\mathrm{int}}=L\ (\text{cycle edges have both endpoints in}\ \Sigma_L\ [\mathrm{Def\text{-}H.12}]);
\]

%dsp-level:0
\[
\quad E_{\mathrm{bd}}
\]

%dsp-level:0
\[
= 1
\]

%dsp-level:0
\[
\le 4\cdot1\ \checkmark;
\]

%dsp-level:0
\[
\quad\mathrm{CT4.9}(1):\ S(v_1)
\]

%dsp-level:0
\[
\le 3\ln2,
\]

%dsp-level:0
\[
\ \text{channel-visible}\ 1\ \text{bit}<3\ \checkmark\ [\text{as in FT.28}(1)]
\]

%dsp-level:0
\[
\implies I_{\mathrm{hid}}
\]

%dsp-level:0
\[
= L\ln2\ [\mathrm{FT.17}]
\]

%dsp-level:0
\[
(2):\ \text{edge counts at vertices}\ [\mathrm{D0.8}]:\ v_1:\ 2\ \text{cycle}+1\ \text{observer}=3;\quad v_{i\ge2}:\ 1\ \text{in}+1\ \text{out}=2;
\]

%dsp-level:0
\[
\qquad \mathrm{CT0.5}/\mathrm{CT0.6}:\ D_{\mathrm{tick}}(n)=2^{n+1}:\ 8,\ 16;\quad O(2)=O(3)=0;\quad\mathrm{CT0.11}:\ n\le3\ \checkmark
\]

%dsp-level:0
\[
(3):\ \Sigma_L\setminus\partial\Sigma_L
\]

%dsp-level:0
\[
= E(c)\ (\text{single chain}):\ \Theta_{\mathrm{int}}
\]

%dsp-level:0
\[
= \Theta(c)
\]

%dsp-level:0
\[
\in 2\mathbb Z\ [\mathrm{FT.22}];
\]

%dsp-level:0
\[
\quad\mathrm{Hol}_{\mathrm{int}}
\]

%dsp-level:0
\[
= (-1)^{2\mathbb Z}
\]

%dsp-level:0
\[
= +1,
\]

%dsp-level:0
\[
\ k
\]

%dsp-level:0
\[
= \Theta/2\ [\mathrm{FT.25}(1),(3)];
\]

%dsp-level:0
\[
\quad W\ \text{criterion}:\ \Theta\ne0\ [\mathrm{FT.20}\ \text{row},\ \text{sound direction}]
\]

%dsp-level:0
\[
(4):\ n(v_i)
\]

%dsp-level:0
\[
= 2\ (i\ge2)
\]

%dsp-level:0
\[
\implies v_i
\]

%dsp-level:0
\[
\in V_2:\ \neg\exists\ \text{colour slot}\ [\mathrm{D3.1},\ \mathrm{CT0.13}:\ \text{strong strictly new only at}\ n{=}3]\implies\mathrm{Colour}(\Sigma_L)\subseteq\{v_1\};
\]

%dsp-level:0
\[
\qquad \text{FT.27's case split applies at}\ v_1:\ \text{block}\equiv0\ \forall C
\]

%dsp-level:0
\[
(5):\ \text{FT.20 rows evaluated as in FT.28}(3);\qquad\dim:\ 5(G)+3(W)+1(H)=9\ [\mathrm{CT0.12}],\ +1\ [\mathrm{CT0.14}]
\]

%dsp-level:0
\[
(6):\ \text{uniform}:\ \Theta=L\kappa\in2\mathbb Z:\ L\ \text{odd}\Rightarrow\kappa\ \text{even}\Rightarrow\kappa_{\min}=2;\quad L\ \text{even}\Rightarrow\kappa_{\min}=1;
\]

%dsp-level:0
\[
\qquad \text{all-nonnull}:\ \min\{\textstyle\sum\kappa_e:\ \kappa_e\ge1,\ \sum\in2\mathbb Z\}=L+(L\ \mathrm{mod}\ 2);
\]

%dsp-level:0
\[
\qquad \text{null-permitting}:\ (1,1,0,\dots,0):\ \Theta=2
\]

%dsp-level:0
\[
(7):\ \Pi=E_{\mathrm{int}}^2\Theta(\ln2)^2\ =L^2\Theta(\ln2)^2;\qquad L=2:\ (1)\text{--}(5)\ \text{reduce to FT.28}(1)\text{--}(5)\ \text{line by line}.\ \blacksquare
\]

\textbf{Rmk (the mass ladder).} $m_{\mathrm{hid}}(\Sigma_L)=L\ln2/\mathfrak h$ is linear in the cycle length at single contact — the first ladder structure of the fermion hunt.

\textbf{Phys.} The whole single-contact cycle family is one card: mass proportional to hidden edge count, interior winding even and parity-trivial, colour structurally absent at every rung of the ladder (one slot-carrying vertex can never seal inside), and the observable mask never exceeds ten dimensions. The floors differ not by which forces act — they all carry $G,H,W$ — but by their integer labels $(L,\Theta,k)$: the fermion hunt, at floor level, is integer bookkeeping.

\vspace{2pt}\noindent\rule{\textwidth}{0.4pt}\vspace{2pt}

\end{theorem}

\begin{theorem}[{\texorpdfstring{FT.35 — The composite floor $\Sigma_{12}$: first colour-active, first $D^-$-carrying}{FT.35 - The composite floor : first colour-active, first -carrying}}]\label{tag:FT-35}
\textbf{Thm.} Let $\Sigma_{12}=\{p_1,q_1\}\cup\{p_2,q_2\}$ with materialized 2-cycles $c_1=(p_1{\to}q_1{\to}p_1)$, $c_2=(p_2{\to}q_2{\to}p_2)$ [FT.24: quanta $\kappa_1,\kappa_2$], joined by the bidirectional connector $e_{12}=(q_1,p_2)$, $\bar e_{12}=(p_2,q_1)$ [both $\in\mathcal R$], single observer contact at $p_1$ ($E_{\mathrm{bd}}=1$; FT.17 hypotheses, Def-H.35). Let $\Gamma=c_1\cup e_{12}\cup c_2\cup\bar e_{12}$ [valid closed chain, D1.4]. If $\Gamma$ materializes:

%dsp-level:0
\[
(1)\quad\kappa_{12}\in\mathbb Z
\]

%dsp-level:0
\[
\qquad [\mathrm{FT.34}(4):\ \Gamma\setminus\{e_{12},\bar e_{12}\}=c_1\sqcup c_2\ \text{decomposes into materialized closed chains}];
\]

%dsp-level:0
\[
(2)\quad E_{\mathrm{int}}=6\ (c_1{:}\,2,\ c_2{:}\,2,\ \text{connector}{:}\,2),\quad I_{\mathrm{hid}}=6\ln2,\quad m_{\mathrm{hid}}=\frac{6\ln2}{\mathfrak h}
\]

%dsp-level:0
\[
\qquad [\mathrm{FT.17};\ \mathrm{CT4.3}\ \text{retention, as below}];
\]

%dsp-level:0
\[
(3)\quad\Theta_{\mathrm{int}}(\Sigma_{12})=2\kappa_1+2\kappa_2+2\kappa_{12}=2\,k(\Gamma)\in2\mathbb Z\ [\mathrm{FT.34}(3)];
\]

%dsp-level:0
\[
\qquad \mathrm{Hol}_{\mathrm{int}}=+1\ [\mathrm{FT.25}];
\]

%dsp-level:0
\[
\qquad \Pi=72\,k(\Gamma)\,(\ln2)^2;
\]

%dsp-level:0
\[
(4)\quad\text{rungs}:\ n(p_1)=3,\ n(q_2)=2,\ n(q_1)=n(p_2)=4:
\]

%dsp-level:0
\[
\quad O(4)
\]

%dsp-level:0
\[
= 2^5-16
\]

%dsp-level:0
\[
= 16>0\ [\mathrm{CT0.6}]:\ \text{the junction vertices are}\ D^-\text{-configured}\ [\mathrm{CT4.1}]:\ \text{excess locked, ledger-retained}\ [\mathrm{D4.2},\ \mathrm{CT4.3}];
\]

%dsp-level:0
\[
(5)\quad\textbf{colour phases}:\ \mathrm{Colour}(\Sigma_{12})\subseteq\{p_1,q_1,p_2\}\ [\mathrm{D3.1};\ \mathrm{CT0.13}:\ \text{strong persists at}\ n\ge4];
\]

%dsp-level:0
\[
\quad\mathcal C(\Sigma_{12};C)
\]

%dsp-level:0
\[
= 1\iff\text{all three seal inside as one}\ \varepsilon\text{-type (baryon-like) triple}\ [\mathrm{D3.9},\ \mathrm{CT3.17},\ \mathrm{FT.34}(5):\ 3\ \text{admits no}\ \{2,3\}\text{-partition with an external remainder}];
\]

%dsp-level:0
\[
\text{strong-active phase}:\ \dim\mathcal K_{\Sigma_{12}}
\]

%dsp-level:0
\[
= 9+1+8
\]

%dsp-level:0
\[
= 18\ —\ \text{the maximal observable mask is attained}\ [\mathrm{FT.26}(3)];
\]

%dsp-level:0
\[
\qquad \text{any external/unsealed attachment}:\ \mathcal C=0\implies\text{block}\equiv0\ [\mathrm{FT.19}].
\]

\textbf{Req.} FT.17 (Def-H.35), FT.19, FT.22–FT.28, FT.32, FT.34; D0.8, D1.4, D3.1, D3.9, D4.1, D4.2; CT0.6, CT0.11, CT0.13, CT0.14 (Part 1); CT3.11, CT3.17 (Part 2); CT4.1, CT4.3 (Part 2); Def-H.9, Def-H.10, Def-H.12, Def-H.15.

\textbf{Pf.}

%dsp-level:0
\[
(1):\ \Gamma\setminus\{e_{12},\bar e_{12}\}
\]

%dsp-level:0
\[
= c_1\sqcup c_2:\ \text{edge-disjoint materialized closed chains}\implies\mathrm{FT.34}(4)\ \text{applies}:\ \kappa_{12}\in\mathbb Z
\]

%dsp-level:0
\[
(2):\ \text{all six cycle/connector edges have both endpoints in}\ \Sigma_{12}\ [\mathrm{Def\text{-}H.12}];
\]

%dsp-level:0
\[
\quad E_{\mathrm{bd}}=1\le4\cdot1\ \checkmark;
\]

%dsp-level:0
\[
\quad q_1,
\]

%dsp-level:0
\[
p_2\ \text{oversubscribed}\ (n{=}4):\ \text{excess locked but retained and counted}\ [\mathrm{CT4.3},\ \mathrm{D4.2};\ \mathrm{FT.17}\ \text{cross-check 2 precedent}]
\]

%dsp-level:0
\[
\implies I_{\mathrm{hid}}
\]

%dsp-level:0
\[
= E_{\mathrm{int}}\ln2\qquad = 6\ln2
\]

%dsp-level:0
\[
(3):\ \Theta_{\mathrm{int}}
\]

%dsp-level:0
\[
= \Theta(c_1)+\Theta(c_2)+\kappa_{12}+\kappa_{\bar e_{12}}
\]

%dsp-level:0
\[
= 2\kappa_1+2\kappa_2+2\kappa_{12}\ [\mathrm{FT.23},\ \mathrm{FT.24}];
\]

%dsp-level:0
\[
\quad\mathrm{FT.34}(3):\
\]

%dsp-level:0
\[
= 2k(\Gamma);
\]

%dsp-level:0
\[
\quad\mathrm{Hol}_{\mathrm{int}}
\]

%dsp-level:0
\[
= (-1)^{2\mathbb Z}
\]

%dsp-level:0
\[
= +1;
\]

%dsp-level:0
\[
\quad\Pi
\]

%dsp-level:0
\[
= E_{\mathrm{int}}^2\Theta_{\mathrm{int}}(\ln2)^2
\]

%dsp-level:0
\[
= 36\cdot2k(\Gamma)(\ln2)^2
\]

%dsp-level:0
\[
(4):\ n(q_1)=2\ (c_1)+2\ (\text{connector})=4;\quad n(p_2)=2+2=4;\quad n(p_1)=2+1=3;\quad n(q_2)=2\ [\mathrm{D0.8}];
\]

%dsp-level:0
\[
\qquad \mathrm{CT0.5}/\mathrm{CT0.6}:\ D_{\mathrm{tick}}(4)
\]

%dsp-level:0
\[
= 2^5
\]

%dsp-level:0
\[
= 32>16
\]

%dsp-level:0
\[
\implies O(4)
\]

%dsp-level:0
\[
= 16;
\]

%dsp-level:0
\[
\quad\mathrm{CT4.1}:\ n
\]

%dsp-level:0
\[
\ge 4
\]

%dsp-level:0
\[
\implies D^-\text{-configured};
\]

%dsp-level:0
\[
\quad\text{locked}
\]

%dsp-level:0
\[
= \text{retained}\ \neg\text{accessible}\ [\mathrm{D4.2}]
\]

%dsp-level:0
\[
(5):\ \text{slots at}\ n\ge3:\ p_1\ (n{=}3),\ q_1,p_2\ (n{=}4:\ \text{rung-3 structure persists}\ [\mathrm{CT0.13}]);
\]

%dsp-level:0
\[
\quad q_2\ (n{=}2):\ \text{slot-less}
\]

%dsp-level:0
\[
\mathcal C=1\iff\forall v\in\mathrm{Colour}:\ \pi_C(v)\subseteq\Sigma_{12}\ [\mathrm{Def\text{-}H.10}];
\]

%dsp-level:0
\[
\qquad \text{partition of}\ 3\ \text{into}\ \{2,3\}\text{-parts}:\ \text{only}\ 3
\]

%dsp-level:0
\[
= 3\implies\varepsilon\text{-type triple}\ [\mathrm{FT.34}(5)];
\]

%dsp-level:0
\[
\quad\delta\text{-pair}+\text{remainder}:\ \text{remainder external}\implies\mathcal C
\]

%dsp-level:0
\[
= 0\ [\mathrm{FT.19}]
\]

%dsp-level:0
\[
\text{strong-active}:\ \dim=5{+}3{+}1{+}1{+}8=18\ [\mathrm{CT0.12},\ \mathrm{CT0.14},\ \mathrm{CT3.8}]=\mathrm{FT.26}(3)\text{'s bound}.\ \blacksquare
\]

\textbf{Rmk (composition creates saturation).} The first composite floor necessarily carries dark content: joining two floors by a bidirectional channel pushes both junction vertices to rung 4, where $D^-$ locks 16 registers' worth of excess [CT0.6, D4.1, D4.2] — retained in the ledger, inaccessible to any observer [D1.11, FT.3]. Composite fermion structure and saturation content are born together; $m_{\mathrm{hid}}$ counts the locked content as retained exactly as FT.17's cross-check 2 already established.

\textbf{Phys.} Two floor fermions wired into one object produce the program's first colour-active floor — sealed as a baryon-type triple, the smallest sealing the partition rule allows at three slots — its first dark junction content, its maximal eighteen-dimensional observable mask, and a composite winding that is the sum of the parts plus the integer distance between them. Meson-type pairs (two slots) and baryon-type triples (three slots) are not conventions: they are the only parts the sealing grammar has.

\vspace{2pt}\noindent\rule{\textwidth}{0.4pt}\vspace{2pt}

\end{theorem}

\begin{theorem}[{\texorpdfstring{FT.42 — The Platonic card template}{FT.42 - The Platonic card template}}]\label{tag:FT-42}
\textbf{Thm.} Under Def-H.20's setup, with (M2) [hence $\kappa_e\in\mathbb Z$ on every channel, FT.24] and (M1) [hence $\Theta_f\in2\mathbb Z$ per face, FT.22]:

%dsp-level:0
\[
(1)\quad E_{\mathrm{int}}=2E,
\]

%dsp-level:0
\[
\qquad I_{\mathrm{hid}}=2E\ln2,
\]

%dsp-level:0
\[
\qquad m_{\mathrm{hid}}=\frac{2E\ln2}{\mathfrak h}
\]

%dsp-level:0
\[
\qquad [\text{observed: FT.17},\ \text{Def-H.35};\ \text{dark: Def-H.9(2) with}\ E_{\mathrm{vis}}=0];
\]

%dsp-level:0
\[
(2)\quad\text{rungs (rung correction, FT.62)}:\ \text{dark}\ n(v)
\]

%dsp-level:0
\[
= 2d\ \forall v;
\]

%dsp-level:0
\[
\ \ \text{observed (1 directed edge)}\ n(v_1)
\]

%dsp-level:0
\[
= 2d+1;
\]

%dsp-level:0
\[
\ \ \text{bidirectional contact}\ 2d+2;
\]

%dsp-level:0
\[
\ \text{others}\ 2d\ [\mathrm{D0.8}:\ \text{directed real-edge count};\ \mathrm{CT0.5}:\ \phi_e\ \text{per incident real edge}];
\]

%dsp-level:0
\[
\qquad O(2d)=112\ (d{=}3),\ 496\ (d{=}4),\ 2032\ (d{=}5);\ \ O(2d{+}1)=240,1008,4080\ [\mathrm{CT0.6}];
\]

%dsp-level:0
\[
\text{saturation universality (the "accessibility trichotomy" epithet does not stand, FT.62(3))}:\ \text{every vertex of every bidirectional floor}\ (d\ge2\ \text{channels}:\ n\ge4)\ \text{is}\ D^-\text{-configured}\ [\mathrm{CT4.1}],
\]

%dsp-level:0
\[
\ \text{content locked, retained, counted}\ [\mathrm{CT4.3},\ \mathrm{D4.2}];
\]

%dsp-level:0
\[
\ \text{grade by degree}:\ O(2d);
\]

%dsp-level:0
\[
\ \ \text{the lock-free family is the unidirectional one}\ [\mathrm{FT.28}(5),\ \mathrm{FT.32}(2);\ \mathrm{FT.62}(3)];
\]

%dsp-level:0
\[
(3)\quad\Theta_{\mathrm{int}}=2K_S,
\]

%dsp-level:0
\[
\qquad K_S:=\!\!\sum_{\text{channels}}\!\!\kappa_{\mathrm{ch}}\in\mathbb Z
\]

%dsp-level:0
\[
\qquad [\mathrm{FT.23}:\ \kappa_{\bar e}=\kappa_e;\ \mathrm{D1.7}:\ \text{additivity}];
\]

%dsp-level:0
\[
\qquad \mathrm{Hol}_{\mathrm{int}}=+1\ [\mathrm{FT.25}(2)];
\]

%dsp-level:0
\[
\qquad W\ \text{row fires}\iff K_S\ne0;
\]

%dsp-level:0
\[
(4)\quad\mathcal A(\Sigma_S;C)\supseteq\{G,H\}\ [G:\ I_{\mathrm{hid}}\ne0;\ H:\ \mathrm{Def-H.36}];
\]

%dsp-level:0
\[
\qquad \mathrm{EM}\in\mathcal A\iff\exists v:\ f_v\ne0\ [\mathrm{D3.3}];
\]

%dsp-level:0
\[
\qquad (\text{scope}\ [\mathrm{FT.62}(5)]:\ \text{type-}(d,d)\ \text{vertices leave D3.3's literal type-}(1,1)\ \text{domain empty on dark bidirectional floors};\ \text{the row evaluates on observed variants via chain-local pass-through reads}\ [\mathrm{FT.21}(i)]);
\]

%dsp-level:0
\[
\text{strong}\in\mathcal A\iff\text{inside-sealed phase}\ (\mathcal C=1):\ \text{exists for every}\ S\ (V\ge4\ \text{partitions into 2s and 3s}\ [\mathrm{FT.34}(5)]),
\]

%dsp-level:0
\[
\ \text{in}\ g(V)\ \text{groupings}\ [\mathrm{FT.41}(b)];
\]

%dsp-level:0
\[
\qquad \dim\mathcal K_{\Sigma_S}=9\ (+1)\ (+8)\le18\ [\mathrm{FT.26}];
\]

%dsp-level:0
\[
(5)\quad\textbf{the card table}:\ \Pi=E_{\mathrm{int}}^2\,\Theta_{\mathrm{int}}(\ln2)^2=8E^2K_S(\ln2)^2\.
\]

%dsp-level:0
\[
\begin{array}{c|c|c|c|c|c|c|c|c}
S & (V,E,F,d) & m_{\mathrm{hid}}\cdot\mathfrak h/\ln2 & O(2d)\ \text{dark}\ & O(2d{+}1)\ \text{contact} & \text{face}\ (L,\kappa_{\min},\Theta_{f,\min},k_{f,\min}) & \text{indep. windings} & g(V)\ \text{groupings} & \Pi/(\ln2)^2\\ \hline
\text{tetrahedron} & (4,6,4,3) & 12 & 112 & 240 & (3,2,6,3) & 3 & 3 & 288\,K_S\\
\text{cube} & (8,12,6,3) & 24 & 112 & 240 & (4,1,4,2) & 5 & 385 & 1152\,K_S\\
\text{octahedron} & (6,12,8,4) & 24 & 496 & 1008 & (3,2,6,3) & 7 & 25 & 1152\,K_S\\
\text{dodecahedron} & (20,30,12,3) & 60 & 112 & 240 & (5,2,10,5) & 11 & 292938721075 & 7200\,K_S\\
\text{icosahedron} & (12,30,20,5) & 60 & 2032 & 4080 & (3,2,6,3) & 19 & 164395 & 7200\,K_S
\end{array}
\]

%dsp-level:0
\[
(6)\quad\textbf{consistency anchors}:\ \text{tetrahedron observed with a bidirectional contact reproduces FT.17 cross-check 3 verbatim}\ (E_{\mathrm{int}}=12,\ E_{\mathrm{bd}}=2\le4\cdot1,\ I_{\mathrm{hid}}=12\ln2);
\]

%dsp-level:0
\[
\ \ \text{the tetrahedron face}\ \{A,B,C\}\ \text{observed by}\ D\ \text{is cross-check 2}\ (E_{\mathrm{int}}=6,\ E_{\mathrm{bd}}=6,\ \delta\tau\ge2);
\]

%dsp-level:0
\[
\Sigma_{\min}=E{=}1\ \text{cell}\ [\mathrm{FT.28}];\qquad\Sigma_{12}=E{=}3\ \text{cell}\ [\mathrm{FT.35}].
\]

\textbf{Req.} Def-H.20; FT.17 (Def-H.35), FT.22, FT.23, FT.24, FT.25, FT.26, FT.32(6), FT.34(5), FT.41(b); Def-H.1, Def-H.8, Def-H.9, Def-H.15; D0.8, D1.7, D3.1, D3.3; CT0.5, CT0.6, CT0.11, CT0.13, CT0.14 (Part 1); CT3.11 (Part 2); CT4.1, CT4.3, CT4.9 (Part 2); D4.2; solid cardinalities, Euler $V-E+F=2$ [math].

\textbf{Pf.}

%dsp-level:0
\[
(1):\ \text{every channel bidirectional}\implies E_{\mathrm{int}}=2E\ [\mathrm{Def\text{-}H.12}:\ \text{both endpoints in}\ \Sigma_S];
\]

%dsp-level:0
\[
\qquad \text{observed}:\ E_{\mathrm{bd}}=1\le4\cdot1\ \checkmark\implies I_{\mathrm{hid}}=E_{\mathrm{int}}\ln2\ [\mathrm{FT.17}];
\]

%dsp-level:0
\[
\qquad (\text{re-scope}\ [\mathrm{FT.62}(4)]:\ \text{at corrected contact rung}\ 2d{+}1\ge7\ \text{the vertex is locked};\ \text{the}\ \mathrm{CT4.9}(1)\ \text{"visible}<3\text{" check passes to Def-H.35}+\mathrm{CT4.3}\ \text{retention};\ \text{value unchanged})
\]

%dsp-level:0
\[
\text{dark}:\ \partial\Sigma_S=\varnothing\implies E_{\mathrm{vis}}=0\implies I_{\mathrm{hid}}=E_{\mathrm{tot}}=2E\ln2\ [\mathrm{Def\text{-}H.9}(2);\ \mathrm{FT.28}\ \text{dark-floor pattern}];
\]

%dsp-level:0
\[
\qquad m_{\mathrm{hid}}=I_{\mathrm{hid}}/\mathfrak h
\]

%dsp-level:0
\[
(2)\ (\mathrm{FT.62}(1)):\ \mathrm{D0.8}\ \text{directed counts}:\ n
\]

%dsp-level:0
\[
= 2d:\ D_{\mathrm{tick}}(2d)
\]

%dsp-level:0
\[
= 2^{2d+1}:\ 128,
\]

%dsp-level:0
\[
512,
\]

%dsp-level:0
\[
2048\ (d=3,4,5);
\]

%dsp-level:0
\[
\quad\mathrm{CT0.6}:\ O(2d)
\]

%dsp-level:0
\[
= 112,
\]

%dsp-level:0
\[
496,
\]

%dsp-level:0
\[
2032;
\]

%dsp-level:0
\[
\quad\text{contact}:\ O(2d{+}1)
\]

%dsp-level:0
\[
= 240,
\]

%dsp-level:0
\[
1008,\qquad 4080
\]

%dsp-level:0
\[
\mathrm{CT4.1}:\ n
\]

%dsp-level:0
\[
\ge 4
\]

%dsp-level:0
\[
\implies D^-\text{-configured};
\]

%dsp-level:0
\[
\quad\text{locked}
\]

%dsp-level:0
\[
= \text{retained}\ [\mathrm{D4.2},\ \mathrm{CT4.3};\ \mathrm{FT.17}\ \text{cross-check 2 precedent}]
\]

%dsp-level:0
\[
\implies E_{\mathrm{tot}}\ \text{accounting intact}
\]

%dsp-level:0
\[
(3):\ \Theta_{\mathrm{int}}=\!\!\sum_{\text{directed edges}}\!\!\kappa_e=\!\!\sum_{\text{channels}}\!\!(\kappa_e+\kappa_{\bar e})=2\!\!\sum_{\text{channels}}\!\!\kappa_{\mathrm{ch}}\ [\mathrm{FT.23}];
\]

%dsp-level:0
\[
\qquad \mathrm{M2}+\mathrm{FT.24}:\ \kappa_{\mathrm{ch}}\in\mathbb Z\implies K_S\in\mathbb Z;
\]

%dsp-level:0
\[
\qquad \mathrm{Hol}_{\mathrm{int}}=(-1)^{2K_S}=+1\ [\mathrm{FT.25}(2)]
\]

%dsp-level:0
\[
(4):\ \text{FT.20 rows as in FT.28}(3)/\mathrm{FT.32}(5);
\]

%dsp-level:0
\[
\qquad \text{slots at every vertex}\ (d\ge3\ \text{for all five solids}\ [\mathrm{D3.1},\ \mathrm{CT3.11}(1)];\ n\ge4:\ \text{rung-3 structure persists}\ [\mathrm{CT0.13}])
\]

%dsp-level:0
\[
V\ge4\implies\{2,3\}\text{-partition exists}\ [\mathrm{FT.34}(5)];\qquad g(V)\ [\mathrm{FT.41}(b)];\qquad\dim\ [\mathrm{FT.26}(3),\ \mathrm{CT0.14}]
\]

%dsp-level:0
\[
(5):\ \text{face regimes from FT.32}(6)\ \text{uniform-}\kappa\ \text{minimality}\ (L\ \text{odd}:\ \kappa_{\min}=2,\ \Theta_{f,\min}=2L,\ k_{f,\min}=L;\ \ L\ \text{even}:\ 1,\ L,\ L/2);
\]

%dsp-level:0
\[
\qquad \Pi=(2E)^2\cdot2K_S(\ln2)^2=8E^2K_S(\ln2)^2
\]

%dsp-level:0
\[
(6):\ \text{FT.17 cross-check 3}:\ \text{whole bidirectional}\ K_4,
\]

%dsp-level:0
\[
\ E_{\mathrm{int}}
\]

%dsp-level:0
\[
= \binom42\times2
\]

%dsp-level:0
\[
= 12
\]

%dsp-level:0
\[
= 2E\ (E{=}6)\ \checkmark;
\]

%dsp-level:0
\[
\quad\text{cross-check 2}:\ \text{face}\ \{A,B,C\}:\ E_{\mathrm{int}}
\]

%dsp-level:0
\[
= 6
\]

%dsp-level:0
\[
= 2\cdot3\ \checkmark,
\]

%dsp-level:0
\[
\ E_{\mathrm{bd}}
\]

%dsp-level:0
\[
= 6
\]

%dsp-level:0
\[
\le 4\cdot2\ \checkmark
\]

%dsp-level:0
\[
\mathrm{FT.28}:\ E{=}1\implies E_{\mathrm{int}}=2\ \checkmark;\qquad\mathrm{FT.35}:\ E{=}3\implies E_{\mathrm{int}}=6\ \checkmark.\ \blacksquare
\]

\textbf{Rmk (what the card does not say).} $g(V)$ counts grouping \textit{patterns}; attachment types (triplet/antitriplet per vertex [D3.8]) and the $\kappa$-profile are configuration data on top. The dodecahedron's sealing count ($292938721075$) and the icosahedron's ($164395$) are grammar sizes, not state counts. Materialization (M1/M2) is hypothesized, not derived — which faces and which 2-cycles are on-spectrum is the corpus's to decide [D2.4, CT2.5].

\textbf{Phys.} Every Platonic solid is a fermion floor with the same five-line card: mass $=2E$ hidden bits over $\mathfrak h$, trivial total parity with an integer winding ledger $K_S$, gravity and Higgs always active, electromagnetism state-dependent, colour available in as many sealing patterns as the vertex count's $\{2,3\}$-grammar allows — and a lock profile fixed by the degree — every bidirectional vertex saturated (rung correction, FT.62: dark rungs are $2d$), graded $O(2d)\in\{112,496,2032\}$, excess absorbed and retained by $D^-$. The solids differ where the integers differ: edges, faces, vertices, degree.

\vspace{2pt}\noindent\rule{\textwidth}{0.4pt}\vspace{2pt}

\end{theorem}

\begin{theorem}[{\texorpdfstring{FT.43 — The winding ledger and Euler consistency}{FT.43 - The winding ledger and Euler consistency}}]\label{tag:FT-43}
\textbf{Thm.} Under FT.42's hypotheses:

%dsp-level:0
\[
(1)\quad\textbf{sum rule}:\ \ \sum_{f}\Theta_f=2\!\!\sum_{\text{channels}}\!\!\kappa_{\mathrm{ch}}=\Theta_{\mathrm{int}}\quad\Longleftrightarrow\quad\sum_f k_f=K_S
\]

%dsp-level:0
\[
\qquad [\text{each channel lies in exactly two faces}:\ FL=2E;\ \mathrm{D1.7}];
\]

%dsp-level:0
\[
(2)\quad\textbf{independent face reads}:\ F-1=E-V+1=\beta_1(\Gamma_S)
\]

%dsp-level:0
\[
\qquad [\text{one sum rule};\ \text{Euler}\ V-E+F=2;\ \text{math}:\ \text{cycle rank of a connected graph}];
\]

%dsp-level:0
\[
(3)\quad\textbf{parity lattice}:\ \text{the on-shell}\ \kappa\text{-profiles}\ \{\kappa\in\mathbb Z^E:\ \Theta_f\in2\mathbb Z\ \forall f\}\ \text{form the kernel of the face-incidence map}\ \varphi:\mathbb Z^E\to\mathbb Z_2^F,
\]

%dsp-level:0
\[
\ \text{of index}\ 2^{F-1}\ \ [\text{math}:\ \operatorname{rk}_{\mathbb Z_2}\varphi=F-1\ \text{for spherical embeddings};\ \text{machine}:\ \text{explicit GF(2) elimination, tetrahedron and cube}];
\]

%dsp-level:0
\[
(4)\quad\text{sector counts}:\ 2^{F-1}
\]

%dsp-level:0
\[
= 2^{\beta_1}:\ \ \text{tetrahedron}\ 8,
\]

%dsp-level:0
\[
\ \text{cube}\ 32,
\]

%dsp-level:0
\[
\ \text{octahedron}\ 128,
\]

%dsp-level:0
\[
\ \text{dodecahedron}\ 2048,
\]

%dsp-level:0
\[
\ \text{icosahedron}\ 2^{19}\qquad = 524288;
\]

%dsp-level:0
\[
(5)\quad\textbf{Euler as card consistency}:\ \beta_1
\]

%dsp-level:0
\[
= E-V+1
\]

%dsp-level:0
\[
= F-1\ \text{ties the three ledgers}:\ \text{mass}\ (E\ \text{via FT.42}(1)),
\]

%dsp-level:0
\[
\ \text{colour}\ (V\ \text{via}\ g(V)),
\]

%dsp-level:0
\[
\ \text{winding}\ (F{-}1\ \text{reads})\ —\ \text{the sphere's Euler number is an identity among card entries, not an extra assumption}.
\]

\textbf{Req.} FT.22, FT.23, FT.24, FT.42; Def-H.20; D1.7; FT.25(2); Euler's formula, cycle-rank and GF(2)-rank facts for connected planar graphs [math]; solid face incidences [math].

\textbf{Pf.}

%dsp-level:0
\[
(1):\ \text{each channel}\ \{e,\bar e\}\ \text{lies in exactly two faces}\ [\text{solid geometry};\ FL=2E\ \text{double-count}]:\quad\sum_f\Theta_f
\]

%dsp-level:0
\[
= \sum_f\sum_{e\in f}\kappa_e
\]

%dsp-level:0
\[
= \sum_{\text{channels}}2\kappa_{\mathrm{ch}}
\]

%dsp-level:0
\[
= \Theta_{\mathrm{int}}\ [\mathrm{D1.7}\ \text{additivity};\ \mathrm{FT.42}(3)]
\]

%dsp-level:0
\[
\Theta_f=2k_f\ [\mathrm{M1}+\mathrm{FT.22}],\ \Theta_{\mathrm{int}}=2K_S\ [\mathrm{FT.42}(3)]\implies\sum_fk_f=K_S
\]

%dsp-level:0
\[
(2):\ \text{one linear relation among}\ F\ \text{reads}\implies F-1\ \text{independent};
\]

%dsp-level:0
\[
\qquad \text{Euler}:\ V-E+F=2\implies F-1=E-V+1=\beta_1\ \ [\text{math}]
\]

%dsp-level:0
\[
(3):\ \varphi(\kappa)_f:=\sum_{e\in f}\kappa_e\ \mathrm{mod}\ 2:\ \ \text{on-shell locus}=\ker\varphi\ [\mathrm{M1}+\mathrm{FT.22}];
\]

%dsp-level:0
\[
\qquad \operatorname{rk}_{\mathbb Z_2}\varphi
\]

%dsp-level:0
\[
= F-1\ [\text{math}:\ \text{faces generate the cycle space with the single spherical relation}]
\]

%dsp-level:0
\[
\implies [\mathbb Z^E:\ker\varphi]\qquad = 2^{F-1}
\]

%dsp-level:0
\[
(4):\ 2^{F-1}\ \text{evaluated at}\ F=4,6,8,12,20;\qquad(5):\ \text{substitution of}\ (2)\ \text{into FT.42}(1)/(4).\ \blacksquare
\]

\textbf{Rmk (what the lattice index means).} Materialization prunes the integer $\kappa$-lattice by exactly the solid's $\mathbb Z_2$ cycle space: an off-lattice profile (some face with odd $\Theta_f$) is a profile no face-materialized configuration can carry [FT.22]. The pruned cosets are not physical sectors — they are the excluded parity classes; the physical ledger is $\ker\varphi$ with its $F-1$ integer face reads.

\textbf{Phys.} A Platonic fermion's geometric memory is the solid's combinatorics turned into integers: one winding read per face, one relation among them, and the number of independent reads is the skeleton's first Betti number — the same number Euler's formula builds from vertices and edges. Mass counts edges, colour counts vertices, winding counts faces, and $V-E+F=2$ is the statement that the three ledgers are not independent. The solids with more independent cycles remember more windings: tetrahedron 3, cube 5, octahedron 7, dodecahedron 11, icosahedron 19.

\vspace{2pt}\noindent\rule{\textwidth}{0.4pt}\vspace{2pt}

\end{theorem}

\begin{theorem}[{\texorpdfstring{FT.44 — Duality signature and the floor taxonomy}{FT.44 - Duality signature and the floor taxonomy}}]\label{tag:FT-44}
\textbf{Thm.} $(1)$ \textbf{Duality mass-degeneracy.} Dual solids exchange $(V,F)$ at fixed $E$ [math]: cube$\leftrightarrow$octahedron $(8,12,6)\leftrightarrow(6,12,8)$, dodecahedron$\leftrightarrow$icosahedron $(20,30,12)\leftrightarrow(12,30,20)$, tetrahedron self-dual $(4,6,4)$. Since $m_{\mathrm{hid}}=2E\ln2/\mathfrak h$ [FT.42(1)] and $\Pi=8E^2K_S(\ln2)^2$ [FT.42(5)] depend on $E$ only: \textbf{dual pairs are mass-degenerate and $\Pi$-degenerate at matched $K_S$} — cube/octahedron: $24\ln2/\mathfrak h$; dodecahedron/icosahedron: $60\ln2/\mathfrak h$; three mass values across the five solids. $(2)$ \textbf{Ledger splitting.} The degeneracy is split by every finer ledger: independent windings $F-1$ (cube 5 vs octahedron 7; dodecahedron 11 vs icosahedron 19); lattice index $2^{F-1}$ ($32$ vs $128$; $2048$ vs $2^{19}$); sealing grammar $g(V)$ ($385$ vs $25$; $292938721075$ vs $164395$); lock profile $O(2d)$ (correction, FT.62(3): cube $112$ vs octahedron $496$; dodecahedron $112$ vs icosahedron $2032$ — saturation grade splits both dual pairs). $(3)$ \textbf{Self-dual floor.} The tetrahedron is the unique self-paired cell ($V=F=4$): its winding and colour ledgers coincide in cardinality ($F-1=3$ independent reads; $g(4)=3$ groupings) — the smallest floor on which the full card, including a $\delta$-pair sealing grammar and a nontrivial winding ledger, is simultaneously active. $(4)$ \textbf{Taxonomy (all computed floors).} Every fermion floor computed in this document obeys $m_{\mathrm{hid}}=E_{\mathrm{int}}\ln2/\mathfrak h$ with $E_{\mathrm{int}}=2E_{\mathrm{bidir}}+E_{\mathrm{uni}}$: $\Sigma_{\min}$: $E_{\mathrm{int}}=2$ [FT.28]; cycle ladder $\Sigma_L$: $E_{\mathrm{int}}=L$ (unidirectional) [FT.32]; $\Sigma_{12}$: $E_{\mathrm{int}}=6$ [FT.35]; Platonic family: $E_{\mathrm{int}}=2E$ [FT.42]. The informational mass is the directed-edge count, universally, across every computed cell.

\textbf{Req.} FT.42, FT.43; FT.28, FT.32, FT.35; FT.41(b); duality cardinals, Euler [math]; Def-H.20.

\textbf{Pf.}

%dsp-level:0
\[
(1):\ \text{duality exchanges}\ V\leftrightarrow F,\ \text{fixes}\ E\ [\text{math};\ \text{cardinals above}];
\]

%dsp-level:0
\[
\qquad m_{\mathrm{hid}},\ \Pi\ \text{are}\ E\text{-functions}\ [\mathrm{FT.42}(1),(5)]\implies\text{degenerate on dual pairs}
\]

%dsp-level:0
\[
(2):\ F\text{-dependence}:\ \mathrm{FT.43}(2),(4);
\]

%dsp-level:0
\[
\qquad V\text{-dependence}:\ g(V)\ [\mathrm{FT.41}(b)];
\]

%dsp-level:0
\[
\qquad d\text{-dependence}:\ O(d)\ [\mathrm{FT.42}(2)]:\ \text{cube}\ d{=}3\ (O{=}0\ \text{dark})\ \text{vs octahedron}\ d{=}4\ (O{=}16)
\]

%dsp-level:0
\[
(3):\ V=F=4:\ F-1=3=g(4)\ [\mathrm{FT.41}(b):\ g(4)=3];
\]

%dsp-level:0
\[
\qquad \text{both ledgers active}:\ \text{winding}\ [\mathrm{FT.43}(2)],\ \delta\text{-pair sealing}\ (4=2{+}2)\ [\mathrm{FT.34}(5)]
\]

%dsp-level:0
\[
(4):\ \text{direct count per cell}:\ \mathrm{FT.28}(1)\text{--}(2),\ \mathrm{FT.32}(1),\ \mathrm{FT.35}(2),\ \mathrm{FT.42}(1).\ \blacksquare
\]

\textbf{Rmk (candidate reading, recorded not claimed).} Dual mass-degeneracy split by finer integer ledgers is a derived degeneracy-splitting \textit{template} for families sharing a mass-like read, distinguished by winding/colour/lock data (an application-layer reading). No physical identification — of duals with generations, or of any card entry with a named particle — is made or implied here.

\textbf{Phys.} Duality, in the floor ledger, is exact mass-degeneracy with everything else distinct: the cube and octahedron weigh the same and remember differently — five winding reads against seven, three-hundred-eighty-five sealing patterns against twenty-five, one hundred twelve locked registers per vertex against four hundred ninety-six (rung correction, FT.62(3)). The tetrahedron, self-dual, is the smallest floor whose card exercises every ledger at once. And across the whole computed taxonomy — digon, cycles, composite pair, all five solids — one law has never broken: the informational mass is the count of directed internal edges, in bits over $\mathfrak h$.

\vspace{2pt}\noindent\rule{\textwidth}{0.4pt}\vspace{2pt}

\end{theorem}

\begin{theorem}[{\texorpdfstring{FT.45 — The bidirectional polygon floors $P_n$: cards, minima, and the null-carrying sector}{FT.45 - The bidirectional polygon floors : cards, minima, and the null-carrying sector}}]\label{tag:FT-45}
\textbf{Thm.} Let $\Sigma_{P_n}$ ($n\ge3$) be the fully bidirectional 1-skeleton of the regular $n$-gon [Def-H.20's setup class], $n\ge3$: $n$ vertices, $n$ bidirectional channels, $2n$ directed edges; sphere compactification with two length-$n$ faces (interior, exterior): $FL=2n=2E$, $V-E+F=n-n+2=2\ \checkmark$ [math]. (The digon-walk floor $\Sigma_{\min}$ [FT.28] is the taxonomy's $E{=}1$ cell, $\neg$a polygon member — clause (5).) Under \textbf{(M2$_n$)} all $n$ channels materialize as 2-cycles [FT.24: $\kappa_i\in\mathbb Z$] and \textbf{(M1$_n$)} the interior face materializes [FT.22: $\Theta_f\in2\mathbb Z$]:

%dsp-level:0
\[
(1)\quad E_{\mathrm{int}}=2n,
\]

%dsp-level:0
\[
\qquad I_{\mathrm{hid}}=2n\ln2,
\]

%dsp-level:0
\[
\qquad m_{\mathrm{hid}}=\frac{2n\ln2}{\mathfrak h},
\]

%dsp-level:0
\[
\qquad \Pi=8n^2K_P(\ln2)^2,
\]

%dsp-level:0
\[
\qquad K_P:=\Theta_{\mathrm{int}}/2=\textstyle\sum_{i=1}^n\kappa_i\ [\mathrm{FT.42}\ \text{template}];
\]

%dsp-level:0
\[
(2)\quad\text{rungs (rung correction, FT.62)}: \text{dark}\ n(v)
\]

%dsp-level:0
\[
= 4\ \forall v\ (2\ \text{channels each},\ \mathrm{D0.8}\ \text{directed count};\ \text{type}\ (2,2));
\]

%dsp-level:0
\[
\ \ \text{single-contact observed}\ n(v_1)
\]

%dsp-level:0
\[
= 5,
\]

%dsp-level:0
\[
\ \text{others}\ 4;
\]

%dsp-level:0
\[
\qquad O(4)
\]

%dsp-level:0
\[
= 16
\]

%dsp-level:0
\[
= C_{\mathrm{cap}},
\]

%dsp-level:0
\[
\ O(5)
\]

%dsp-level:0
\[
= 48\ [\mathrm{CT0.5},\ \mathrm{CT0.6}]:\ \textbf{uniform minimal saturation at every vertex of every member}\ [\mathrm{CT4.1}];
\]

%dsp-level:0
\[
\ \text{locked content retained and counted}\ [\mathrm{CT4.3},\ \mathrm{D4.2}];
\]

%dsp-level:0
\[
\ \text{the lock-free family is the unidirectional one}\ [\mathrm{FT.32}(2),\ \mathrm{FT.62}(3)];
\]

%dsp-level:0
\[
(3)\quad\textbf{walk space and the single parity constraint}:\ \beta_1(\text{directed multigraph})
\]

%dsp-level:0
\[
= 2n-n+1
\]

%dsp-level:0
\[
= n+1\ [\text{math}];
\]

%dsp-level:0
\[
\qquad \text{closed-walk structure}:\ \text{flow balance at every vertex}\implies\text{per-channel net flow constant}
\]

%dsp-level:0
\[
= w\ (\text{the walk's winding});
\]

%dsp-level:0
\[
\Theta(w)
\]

%dsp-level:0
\[
= w\,
\]

%dsp-level:0
\[
\textstyle\sum_i\kappa_i+2\sum_i\kappa_i b_i\ \ (b_i:\ \text{reverse traversals})\implies\Theta(w)
\]

%dsp-level:0
\[
\in 2\mathbb Z\ \forall w\iff\textstyle\sum_i\kappa_i
\]

%dsp-level:0
\[
\in 2\mathbb Z\ \ (w{=}1:\ \text{the face},\ \mathrm{M1}_n):
\]

%dsp-level:0
\[
\textbf{exactly one independent parity constraint};
\]

%dsp-level:0
\[
\qquad \text{on-shell}\ \kappa\text{-lattice}=\{\kappa\in\mathbb Z^n:\ \textstyle\sum\kappa_i\in2\mathbb Z\}:\ \text{index}\ 2=2^{F-1}\ [\mathrm{FT.43}(3)\ \text{pattern at}\ F{=}2];
\]

%dsp-level:0
\[
\qquad \text{digons}\ (w{=}0):\ \Theta=2\sum\kappa_ib_i\in2\mathbb Z\ \text{automatic}\ [\mathrm{M2}_n,\ \mathrm{FT.24}];
\]

%dsp-level:0
\[
(4)\quad\textbf{four minimality regimes}\ (\Theta_{\min}>0,\ \text{face materialized};\ \text{machine-enumerated for}\ n\le7):
\]

%dsp-level:0
\[
\text{spacelike-or-null}\ [\kappa\ge0]:\ \Theta_{\min}=2\ \forall n\ge3\ (\text{profile}\ (1,1,0,\dots,0));
\]

%dsp-level:0
\[
\qquad \textbf{null-carrying dimension}=|\{i:\kappa_i=0\}|\ge n-2\ \text{in every such minimizer};
\]

%dsp-level:0
\[
\text{strictly-spacelike}\ [\kappa\ge1]:\ \Theta_{\min}
\]

%dsp-level:0
\[
= n+(n\ \mathrm{mod}\ 2)\ (n\ \text{even}:\ \text{unique all-1 profile};\ n\ \text{odd}:\ \text{one}\ 2,\ \text{rest}\ 1);
\]

%dsp-level:0
\[
\qquad \text{uniform}\ [\kappa_i\equiv\kappa]:\ \Theta_{\min}=n\ (n\ \text{even},\ \kappa{=}1)\ /\ 2n\ (n\ \text{odd},\ \kappa{=}2)\ [\mathrm{FT.32}(6)];
\]

%dsp-level:0
\[
\text{signature-mixed}\ [\kappa\in\mathbb Z,\ \text{timelike}\ \kappa<0\ \text{admissible}\ [\mathrm{CT1.12}]]:\ \Theta_{\min}
\]

%dsp-level:0
\[
= 2\ \forall n
\]

%dsp-level:0
\[
\ge 3\ (\text{witnesses}:\ n\ge4\ \text{even}:\ (1^{(n+2)/2},(-1)^{(n-2)/2});\ n\ge5\ \text{odd}:\ (1^{n-1},-(n-3));\ n{=}3:\ (2,1,-1));
\]

%dsp-level:0
\[
(5)\quad\textbf{the digon distinction}:\ \text{the null-carrying sector has dimension}\ n-2:\ \text{at}\ n
\]

%dsp-level:0
\[
= 2\ (\Sigma_{\min})\ \text{it is}\ \{-1\}\text{-dimensional}
\]

%dsp-level:0
\[
= \text{empty}:\ \Theta_{\min}
\]

%dsp-level:0
\[
= 2\ \text{there forces}\ \kappa_1
\]

%dsp-level:0
\[
= \kappa_2
\]

%dsp-level:0
\[
= 1\ \text{(all-nonnull, purely spacelike)}\ [\mathrm{FT.24}(4)];
\]

%dsp-level:0
\[
\ \text{for every}\ n
\]

%dsp-level:0
\[
\ge 3\ \text{the minimum is null-carrying}.\ \textbf{The fermion floor is derived to be a digon, not a polygon};
\]

%dsp-level:0
\[
(6)\quad\textbf{sealing}:\ \text{dark}\ P_n:\ \mathrm{Colour}
\]

%dsp-level:0
\[
= \varnothing\ (n(v)=2<3\ [\mathrm{D3.1}])\implies\text{strong}\ \neg\in\mathcal A\ \text{always}\ [\mathrm{FT.20}];
\]

%dsp-level:0
\[
\ \ \text{observed}:\ \mathrm{Colour}
\]

%dsp-level:0
\[
= \{v_1\}:\ \neg\exists\ \text{inside partner}\ [\mathrm{A1}]\implies\text{structurally colourless}\ [\mathrm{FT.27}\ \text{case split}];
\]

%dsp-level:0
\[
\quad\mathcal A\supseteq\{G,H\},
\]

%dsp-level:0
\[
\ W
\]

%dsp-level:0
\[
\iff K_P
\]

%dsp-level:0
\[
\ne 0,
\]

%dsp-level:0
\[
\ \mathrm{EM}\iff\exists v:\ f_v
\]

%dsp-level:0
\[
\ne 0;
\]

%dsp-level:0
\[
\quad\dim\mathcal K\qquad \le 10<18;
\]

%dsp-level:0
\[
(7)\quad\textbf{anchors}:\ \Sigma_{\min}\ (\text{digon walk},\ E{=}1)\ \text{is adjacent to but}\ \neg\text{a member of this family}\ [\mathrm{FT.28}/\mathrm{FT.44}(4)\ \text{taxonomy}]:\ \text{clauses}\ (3)\text{--}(5)\ \text{degenerate there}\ (\neg\exists\ \text{proper face};\ \text{its minimum forces}\ \kappa{=}1\ [\mathrm{FT.24}(4)]);
\]

%dsp-level:0
\[
\quad\text{the unidirectional}\ 3\text{-cycle}\
\]

%dsp-level:0
\[
= \ \mathrm{FT.17}\ \text{cross-check 1}\ [\mathrm{FT.32}];
\]

%dsp-level:0
\[
\quad\text{the observed tetrahedron face is NOT}\ P_3\text{: in bidirectional}\ K_4\ \text{the face vertices carry the third (external) channel each}\ (n(v)=4,\ E_{\mathrm{bd}}=6)\ [\mathrm{FT.17}\ \text{cross-check 2}].
\]

\textbf{Req.} FT.17 (Def-H.35), FT.20, FT.22, FT.23, FT.24, FT.25, FT.26, FT.27, FT.28, FT.32(6), FT.42; Def-H.20; A1; D0.8, D1.7, D3.1, D3.3; CT0.5, CT0.6, CT0.14 (Part 1); CT1.12 (\textbf{C}, signature classes); CT4.9 (Part 2); cycle-rank and partition facts [math].

\textbf{Pf.}

%dsp-level:0
\[
(1),(2):\ \mathrm{FT.42}(1)\text{--}(2)\ \text{template at}\ E=n,\ V=n,\ d_{\text{dark}}=2,\ d_{\text{contact}}=3;
\]

%dsp-level:0
\[
\qquad \mathrm{CT0.6}:\ O(2)=O(3)=0;
\]

%dsp-level:0
\[
\qquad \mathrm{CT4.9}(1):\ \text{visible}\ 1\ \text{bit}<3\ \checkmark
\]

%dsp-level:0
\[
(3):\ \beta_1=2n-n+1=n+1\ [\text{math}:\ \text{connected multigraph}];
\]

%dsp-level:0
\[
\qquad \text{closed walk}:\ \text{in-degree}
\]

%dsp-level:0
\[
= \text{out-degree at every vertex}
\]

%dsp-level:0
\[
\implies a_i-b_i
\]

%dsp-level:0
\[
= \text{const}
\]

%dsp-level:0
\[
= w\ [\text{flow balance around the cycle};\ a_i,b_i:\ \text{forward/reverse traversals of channel}\ i]
\]

%dsp-level:0
\[
\Theta(w)
\]

%dsp-level:0
\[
= \sum_i\kappa_i(a_i+b_i)
\]

%dsp-level:0
\[
= \sum_i\kappa_i(2b_i+w)
\]

%dsp-level:0
\[
= w\sum_i\kappa_i+2\sum_i\kappa_ib_i\ [\mathrm{D1.7}\ \text{additivity};\ \mathrm{FT.23}:\ \kappa_{\bar e}=\kappa_e]\implies\text{parity}(\Theta(w))
\]

%dsp-level:0
\[
= \text{parity}\bigl(w\textstyle\sum_i\kappa_i\bigr)
\]

%dsp-level:0
\[
\text{face walk}:\ w=1:\ \text{constraint}\ \textstyle\sum\kappa_i\in2\mathbb Z\ [\mathrm{M1}_n+\mathrm{FT.22}];\quad\text{digons}:\ w=0:\ \text{automatic};
\]

%dsp-level:0
\[
\quad w
\]

%dsp-level:0
\[
\ge 2:\ \text{automatic once the face is constrained}\implies\textbf{exactly one independent parity constraint};
\]

%dsp-level:0
\[
\quad\text{index}\ 2\ [\text{math}]
\]

%dsp-level:0
\[
(4):\ \text{spacelike-or-null}:\ (1,1,0^{n-2}):\ \Theta=2\in2\mathbb Z\ \checkmark;\quad\Theta>0\ \text{even}\Rightarrow\Theta\ge2:\ \text{minimal};
\]

%dsp-level:0
\[
\qquad \Theta=2\ \text{with}\ \kappa\ge0\implies\le2\ \text{positive units}\implies\ge n{-}2\ \text{zeros}
\]

%dsp-level:0
\[
\text{strictly-spacelike}:\ \min\{\textstyle\sum\kappa_i:\ \kappa_i\ge1,\ \sum\in2\mathbb Z\}=n+(n\ \mathrm{mod}\ 2);
\]

%dsp-level:0
\[
\qquad \text{uniform}:\ n\kappa\in2\mathbb Z\Rightarrow\kappa_{\min}=1\ (n\ \text{even}),\ 2\ (n\ \text{odd})\ [\mathrm{FT.32}(6)]
\]

%dsp-level:0
\[
\text{signature-mixed witnesses}:\ n\ge4\ \text{even}:\ \tfrac{n+2}{2}\ \text{entries}\ {+}1,\ \tfrac{n-2}{2}\ \text{entries}\ {-}1:\ \Theta=2\ \checkmark;
\]

%dsp-level:0
\[
\quad n\ge5\ \text{odd}:\ (1^{n-1},-(n-3)):\ \Theta=(n-1)-(n-3)=2\ \checkmark;
\]

%dsp-level:0
\[
\quad n=3:\ (2,1,-1):\ \Theta=2\ \checkmark
\]

%dsp-level:0
\[
(5):\ n=2:\ \Theta=2\kappa\ [\mathrm{FT.23}];\quad\Theta=2\Rightarrow\kappa=1:\ \neg\exists\ \text{null slot};
\]

%dsp-level:0
\[
\qquad n\ge3:\ (1,1,0^{n-2})\ \text{admissible}\ [\mathrm{CT1.19a}:\ \kappa=0\iff\text{null edge permitted}]
\]

%dsp-level:0
\[
(6):\ \text{dark}:\ n(v)
\]

%dsp-level:0
\[
= 2\implies\neg\exists\ \text{colour slots}\ [\mathrm{D3.1}:\ \text{slots}=colour\ \text{at}\ n{=}3;\ \mathrm{CT0.13}:\ \text{strong strictly new at}\ n{=}3]\implies\mathrm{Colour}\qquad = \varnothing;
\]

%dsp-level:0
\[
\qquad \text{observed}:\ \mathrm{Colour}\subseteq\{v_1\}:\ \mathrm{FT.27}\ \text{case split}
\]

%dsp-level:0
\[
(7):\ \text{direct comparison of}\ (1)\text{--}(6)\ \text{against FT.28}\ (n{=}2),\ \mathrm{FT.17}\ \text{cross-checks 1--2}.\ \blacksquare
\]

\textbf{Rmk (regime precedence; withdrawn draft claim).} The four regimes are nested constraint classes, not alternatives: strictly-spacelike $\subseteq$ spacelike-or-null; signature-mixed relaxes both. Physically primary for a \textit{spatial} polygon floor are the two spacelike regimes; the mixed regime is kinematically admissible [CT1.12: $\operatorname{sig}=(-,+,+,+)$ permits $\kappa<0$] but describes floors with timelike interior edges — a different physical situation, recorded for completeness. An earlier draft's "$n\equiv2\ (\mathrm{mod}\ 4)$ anomaly" is \textbf{withdrawn as an arithmetic error} (its witness summed to $8$, not $6$): no residue-class anomaly exists — the strictly-spacelike floor is $n+(n\ \mathrm{mod}\ 2)$ for every $n$, with the unique all-1 profile at even $n$.

\textbf{Rmk (accessibility contrast, corrected [FT.62(3)]).} $P_n$ dark has $n(v)=4$ everywhere [D0.8 directed count]: exactly one channel quantum over budget — $O(4)=16=C_{\mathrm{cap}}$, uniform in $n$: the polygon family is the \textit{uniform minimal-saturation} family (the "unlimited-accessibility" epithet does not stand). The Platonic family saturates graded by degree, $O(2d)\ge112$; the lock-free cells are the unidirectional floors ($n\le3$) [FT.28(5), FT.32(2)]. First structural difference between the two polygon-derived families — uniform vs graded saturation — derived from D0.8/CT0.5/CT0.6 alone.

\textbf{Phys.} Bidirectional polygons are the fermion program's cleanest infinite family: one uniform minimal lock per vertex ($O(4)=16$, correction), no colour ever, exactly one parity constraint from the face (an index-2 on-shell lattice), and a null-carrying minimal sector of dimension $n-2$ — one that has no analogue at the digon, whose minimum forces $\kappa=1$: the fermion floor's ledger is qualitatively polygon-external.

\vspace{2pt}\noindent\rule{\textwidth}{0.4pt}\vspace{2pt}

\end{theorem}

\begin{theorem}[{\texorpdfstring{FT.49 — The tetrahedron-face floor: full card}{FT.49 - The tetrahedron-face floor: full card}}]\label{tag:FT-49}
\textbf{Thm.} Let $\Sigma_{\triangle}=\{A,B,C\}$ be the bidirectional 3-face of the bidirectional $K_4$ [FT.17 cross-check 2's setup], observed by the fourth vertex $D$ (six boundary edges). Under M2 (face-channel digons materialize [FT.24]) and M1 (both directed face triangles materialize [FT.22]):

%dsp-level:0
\[
(1)\quad E_{\mathrm{int}}=6,\quad E_{\mathrm{bd}}=6\le4\,\delta\tau\Rightarrow\delta\tau\ge2\ [\mathrm{FT.17}\ \text{hypotheses},\ \text{Def-H.35}];
\]

%dsp-level:0
\[
\qquad I_{\mathrm{hid}}=6\ln2,
\]

%dsp-level:0
\[
\qquad m_{\mathrm{hid}}=\frac{6\ln2}{\mathfrak h}\ \ [\text{the multi-tick observation cell of the taxonomy}];
\]

%dsp-level:0
\[
(2)\quad\text{interior channel structure}
\]

%dsp-level:0
\[
= \Sigma_{P_3}'s\ [\mathrm{FT.45}\ \text{at}\ n{=}3]:\ \kappa\text{-lattice}\ \{\kappa\in\mathbb Z^3:\ \kappa_{AB}{+}\kappa_{BC}{+}\kappa_{CA}\in2\mathbb Z\}\ \text{of index}\ 2;
\]

%dsp-level:0
\[
\ \ \text{both directed triangles give the same constraint}\ [\mathrm{FT.23}:\ \text{orientation-even}];
\]

%dsp-level:0
\[
(3)\quad\Theta_{\mathrm{int}}=2\,\Theta_{\triangle}=4\,k_{\triangle}\in4\mathbb Z,
\]

%dsp-level:0
\[
\qquad k_{\triangle}:=\Theta_{\triangle}/2\in\mathbb Z\ [\mathrm{FT.22}];
\]

%dsp-level:0
\[
\qquad \mathrm{Hol}_{\mathrm{int}}=+1\ [\mathrm{FT.25}(2)];
\]

%dsp-level:0
\[
\qquad W\ \text{row fires}\iff\Theta_{\triangle}\ne0;
\]

%dsp-level:0
\[
(4)\quad\text{minimality regimes}\ [\mathrm{FT.45}(4)\ \text{at}\ n{=}3]:\quad\Theta_{\triangle,\min}
\]

%dsp-level:0
\[
= 2\ (\text{null-carrying}\ (1,1,0)\text{-class}),
\]

%dsp-level:0
\[
\ 4\ (\text{strictly-spacelike}\ (1,1,2)\text{-class}),
\]

%dsp-level:0
\[
\ 6\ (\text{uniform}\ \kappa{=}2);
\]

%dsp-level:0
\[
\qquad \Pi=36\,\Theta_{\mathrm{int}}(\ln2)^2\in\{144,288,432\}(\ln2)^2;
\]

%dsp-level:0
\[
(5)\quad\text{rungs and locks}:\ n(A)=n(B)=n(C)=6\ [\mathrm{D0.8}:\ 3\ \text{channels each in the}\ K_4];
\]

%dsp-level:0
\[
\qquad O(6)
\]

%dsp-level:0
\[
= 2^7-16
\]

%dsp-level:0
\[
= 112\ [\mathrm{CT0.5},\ \mathrm{CT0.6}]:\ \text{every vertex}\ D^-\text{-configured}\ [\mathrm{CT4.1}],
\]

%dsp-level:0
\[
\ \text{locked content retained and counted}\ [\mathrm{CT4.3},\ \mathrm{D4.2}];
\]

%dsp-level:0
\[
(6)\quad\textbf{baryon-type sealing phase}:\ \text{all three vertices slot-carrying}\ (n\ge3\ [\mathrm{D3.1},\ \mathrm{CT0.13}:\ \text{rung-3 structure persists}]);
\]

%dsp-level:0
\[
\quad\mathcal C(\Sigma_{\triangle};C)
\]

%dsp-level:0
\[
= 1\iff\text{the three attachments seal as one}\ \varepsilon\text{-type triple}\ [3=3\ \text{is the only}\ \{2,3\}\text{-partition}:\ \mathrm{FT.34}(5),\ g(3)=1\ \mathrm{FT.41}(b)]:\ \text{the first computed cell whose strong phase is exactly D3.9's baryon structure}\ B
\]

%dsp-level:0
\[
= \varepsilon_{ijk}q^iq^jq^k\ (\text{structural correspondence,}\ \neg\text{particle identification});
\]

%dsp-level:0
\[
(7)\quad\mathcal A(\Sigma_{\triangle};C)\supseteq\{G,H\}\ [G:\ I_{\mathrm{hid}}\ne0;\ H:\ \mathrm{Def-H.36}];\quad\mathrm{EM}\iff\exists v:\ f_v\ne0;\quad\text{strong}\iff\varepsilon\text{-sealed};
\]

%dsp-level:0
\[
\qquad \dim\mathcal K=9\ (+1)\ (+8)\le18\ [\mathrm{FT.26}];
\]

%dsp-level:0
\[
(8)\quad\text{boundary algebra}:\ N(\Sigma_{\triangle})=6\implies\mathcal H_{\mathrm{slot}}=\mathbb C^6,\ \iota_\Sigma:\ \mathfrak u(6)\ \text{of dimension}\ 36;
\]

%dsp-level:0
\[
\qquad \text{strictly new at}\ \Sigma:\ 36-6
\]

%dsp-level:0
\[
= 30\ \text{real dimensions}\ [\mathrm{FT.41}(d)]\ —\ \text{the largest exchange sector computed}.
\]

\textbf{Req.} FT.17 (Def-H.35), FT.22, FT.23, FT.24, FT.25, FT.26, FT.34(5), FT.41(b),(d), FT.45; Def-H.9, Def-H.10, Def-H.12, Def-H.15; D0.8, D3.1, D3.3, D3.9; CT0.5, CT0.6, CT0.13, CT0.14 (Part 1); CT3.11, CT3.17 (Part 2); CT4.1, CT4.3 (Part 2); D4.2.

\textbf{Pf.}

%dsp-level:0
\[
(1):\ \mathrm{FT.17}\ \text{cross-check 2 verbatim}:\ E_{\mathrm{int}}=6,\ E_{\mathrm{bd}}=6,\ \delta\tau\ge2;
\]

%dsp-level:0
\[
\qquad I_{\mathrm{hid}}=E_{\mathrm{int}}\ln2=6\ln2\ [\mathrm{FT.17},\ \text{Def-H.35}];\quad m_{\mathrm{hid}}=I_{\mathrm{hid}}/\mathfrak h\ [\mathrm{Def\text{-}H.9}(2)]
\]

%dsp-level:0
\[
(2):\ \text{interior channels}
\]

%dsp-level:0
\[
= \{AB,BC,CA\}\ \text{bidirectional}
\]

%dsp-level:0
\[
= \Sigma_{P_3}\ \text{interior}\ [\mathrm{FT.45}\ \text{setup}];
\]

%dsp-level:0
\[
\quad\text{digons}:\ \kappa\in\mathbb Z\ [\mathrm{M2}+\mathrm{FT.24}];
\]

%dsp-level:0
\[
\quad\text{triangles}:\ \Theta_{\triangle}
\]

%dsp-level:0
\[
= \kappa_{AB}{+}\kappa_{BC}{+}\kappa_{CA}\ [\mathrm{FT.23}:\ \text{same both directions}]
\]

%dsp-level:0
\[
\in 2\mathbb Z\ [\mathrm{M1}+\mathrm{FT.22}]
\]

%dsp-level:0
\[
(3):\ \Theta_{\mathrm{int}}=\!\!\sum_{\text{6 directed}}\!\!\kappa=2\Theta_{\triangle};
\]

%dsp-level:0
\[
\qquad \Theta_{\triangle}=2k_{\triangle}\ [\mathrm{FT.22}]\implies\Theta_{\mathrm{int}}=4k_{\triangle};
\]

%dsp-level:0
\[
\qquad \mathrm{Hol}_{\mathrm{int}}=(-1)^{4k_{\triangle}}=+1\ [\mathrm{FT.25}(2)]
\]

%dsp-level:0
\[
(4):\ \mathrm{FT.45}(4)\ \text{regimes at}\ n{=}3;
\]

%dsp-level:0
\[
\qquad \Pi=E_{\mathrm{int}}^2\Theta_{\mathrm{int}}(\ln2)^2=36\cdot(2\Theta_{\triangle})(\ln2)^2:\ \Theta_{\triangle}\in\{2,4,6\}\Rightarrow\{144,288,432\}(\ln2)^2
\]

%dsp-level:0
\[
(5):\ \mathrm{D0.8}:\ n
\]

%dsp-level:0
\[
= 3\ \text{channels}\times2
\]

%dsp-level:0
\[
= 6;
\]

%dsp-level:0
\[
\quad\mathrm{CT0.5}/\mathrm{CT0.6}:\ D_{\mathrm{tick}}(6)
\]

%dsp-level:0
\[
= 2^7
\]

%dsp-level:0
\[
= 128,
\]

%dsp-level:0
\[
\ O(6)
\]

%dsp-level:0
\[
= 112;
\]

%dsp-level:0
\[
\quad\mathrm{CT4.1}:\ n
\]

%dsp-level:0
\[
\ge 4
\]

%dsp-level:0
\[
\Rightarrow D^-\text{-configured};
\]

%dsp-level:0
\[
\quad\text{retained}\ [\mathrm{CT4.3}]
\]

%dsp-level:0
\[
(6):\ n\ge3\Rightarrow\text{slots}\ [\mathrm{D3.1},\ \mathrm{CT3.11}(1)];\quad\text{partitions of}\ 3\ \text{into}\ \{2,3\}\text{-parts}:\ \text{only}\ 3\ [\mathrm{FT.34}(5)];
\]

%dsp-level:0
\[
\quad g(3)=1\ [\mathrm{FT.41}(b)];\quad\varepsilon\text{-triple}=B\ \text{structure}\ [\mathrm{D3.9},\ \mathrm{CT3.17}]
\]

%dsp-level:0
\[
(7):\ \mathrm{FT.20}\ \text{rows};
\]

%dsp-level:0
\[
\quad (8):\ N(\Sigma)=E_{\mathrm{bd}}=6\ [\mathrm{Def\text{-}H.12}];\quad\dim\mathfrak u(6)=36;\quad\text{factored diagonal}=6\implies\text{strictly new}=30\ [\mathrm{FT.41}(d)].\ \blacksquare
\]

\textbf{Rmk (relation to the polygon and Platonic cells).} $\Sigma_{\triangle}$ is neither $\Sigma_{P_3}$ nor a Platonic floor: its interior ledger is the polygon's (index-2 lattice, $n{-}2{=}1$ null freedom at the minimum), while its rung/lock ledger is the ambient $K_4$'s ($n{=}6$, $O{=}112$) — the first computed cell where interior and ambient ledgers visibly decouple. Its sealing phase is the baryon type; $\Sigma_{12}$'s was also $\varepsilon$-type [FT.35(5)] but at three slots from two floors plus a junction — the grammar [FT.34(5)] is the same, the wiring differs.

\textbf{Phys.} The tetrahedron face is the program's densest small cell: six hidden bits, a winding ledger even in both directions ($\Theta_{\mathrm{int}}\in4\mathbb Z$), every vertex deep in the saturation regime with 112 registers' absorption rate each, one baryon-type colour phase, and thirty dimensions of strictly-new boundary exchange structure — all of it derived from cards already in the ledger, assembled without a single new premise beyond the declared materialization hypotheses.

\vspace{2pt}\noindent\rule{\textwidth}{0.4pt}\vspace{2pt}

\end{theorem}

\begin{theorem}[{\texorpdfstring{FT.50 — Composition and gluing laws}{FT.50 - Composition and gluing laws}}]\label{tag:FT-50}
\textbf{Thm.} $(1)$ \textbf{Bridged union.} Let $\Sigma^{(1)},\Sigma^{(2)}$ be floors with bidirectional interiors [$E_{\mathrm{int}}^{(i)}=2E_i$], joined by a bidirectional connector path $P=(f_1,\dots,f_m)$ (disjoint from both interiors), single external contact; all digons materializing [M2-type]. Then

%dsp-level:0
\[
E_{\mathrm{int}}(\Sigma)=2E_1+2E_2+2m,
\]

%dsp-level:0
\[
\qquad m_{\mathrm{hid}}(\Sigma)=m_{\mathrm{hid}}^{(1)}+m_{\mathrm{hid}}^{(2)}+\frac{2m\ln2}{\mathfrak h}
\]

%dsp-level:0
\[
\qquad [\mathrm{FT.17},\ \text{Def-H.35}]:\ \textbf{informational mass is additive under bridged composition}.
\]

$(2)$ \textbf{Generalized round-trip.} If $\Gamma=c_1\cdot P\cdot c_2\cdot\bar P$ materializes [$c_i$ materialized closed chains at the attachments; D1.4 concatenation]: $d_{12}:=\Theta(P)\in\mathbb Z$ and $k(\Gamma)=k_1+k_2+d_{12}$ [FT.34(2)–(3) generalized to $m$-channel connectors]; per-edge integrality $\kappa_{f_j}\in\mathbb Z$ requires each connector digon to materialize [FT.24] — the sum $d_{12}$ is integral without that strengthening. $(3)$ \textbf{Junction saturation.} A bidirectional bridge into a vertex of interior degree $d\ge2$ gives $n_{\mathrm{junction}}=d+2\ge4\implies O(d{+}2)=2^{d+3}-16>0$: \textbf{composition creates darkness at every junction} [CT0.5, CT0.6, CT4.1, D4.2; FT.35's Rmk generalized]; content retained [CT4.3]. For $s$ triangular floors sharing one junction vertex: $2s$ channels meet there, $n=4s$ [D0.8], $O(4s)=2^{4s+1}-16$ — absorption rate grows exponentially in the sharing multiplicity [D4.2: rate $\propto O(n)$]: high-multiplicity junctions are maximally dark. $(4)$ \textbf{Gluing law.} For $\Sigma=\Sigma_1\cup_G\Sigma_2$ glued along a shared subgraph $G$: $E_{\mathrm{int}}=E_{\mathrm{int}}^{(1)}+E_{\mathrm{int}}^{(2)}-|E_{\mathrm{dir}}(G)|$; the on-shell parity lattice is $\{\kappa:\ \Theta_c\in2\mathbb Z\ \forall\ \text{materialized closed chains}\ c\}$ — its index is $2^{\operatorname{rk}_{\mathbb Z_2}(\text{materialized cycle system})}$, dependencies automatic: a cycle homologous to a sum of materialized cycles inherits evenness [D1.7 additivity; FT.23 orientation-blindness]. $(5)$ \textbf{Computed cells.} \textbf{Bowtie} (two bidirectional triangles sharing one vertex $X$): $E_{\mathrm{int}}=12$, $m_{\mathrm{hid}}=12\ln2/\mathfrak h$; both faces materialized $\Rightarrow\Theta_{\mathrm{int}}=2(K_1{+}K_2)\in4\mathbb Z$; constraint rank 2 (index 4); rungs: $n(X)=8$ [$O(8)=496$], other four vertices $n=4$ [$O(4)=16$] — every vertex locked; five slots $\Rightarrow g(5)=10$ sealing groupings (one $\delta$-pair $+$ one $\varepsilon$-triple) [FT.41(b)]. \textbf{Double-triangle patch} (two tetrahedron faces glued along channel $AB$; $\Sigma=\{A,B,C,D\}$): $E_{\mathrm{int}}=10$; constraints $\Theta_{ABC},\Theta_{ABD}\in2\mathbb Z$: $\operatorname{rk}_{\mathbb Z_2}=2$ (index 4), and the boundary 4-cycle constraint is \textit{derived}: $\Theta_{ACBD}=\Theta_{ABC}+\Theta_{ABD}-2\kappa_{AB}\in2\mathbb Z$; rungs $n(A)=n(B)=6$ [$O=112$], $n(C)=n(D)=4$ [$O=16$]; four slots $\Rightarrow g(4)=3$ groupings, all $\delta{+}\delta$ (meson–meson phases — the only $\{2,3\}$-partition of 4).

\textbf{Req.} FT.17 (Def-H.35), FT.22, FT.23, FT.24, FT.27, FT.34, FT.35, FT.41(b), FT.45; Def-H.9, Def-H.10, Def-H.12; D0.8, D1.4, D1.7, D3.1, D3.9; CT0.5, CT0.6 (Part 1); CT3.17 (Part 2); CT4.1, CT4.3 (Part 2); D4.2; $\mathbb Z_2$-rank facts [math].

\textbf{Pf.}

%dsp-level:0
\[
(1):\ E_{\mathrm{int}}:\ \text{interiors disjoint}\ \sqcup\ 2m\ \text{connector directed edges};
\]

%dsp-level:0
\[
\quad E_{\mathrm{bd}}=1\le4\delta\tau\ \checkmark;
\]

%dsp-level:0
\[
\quad I_{\mathrm{hid}}=E_{\mathrm{int}}\ln2\ [\mathrm{FT.17}]\implies m_{\mathrm{hid}}=\tfrac{E_{\mathrm{int}}\ln2}{\mathfrak h}=\text{sum of parts}
\]

%dsp-level:0
\[
(2):\ \mathrm{FT.34}(1):\ \Theta(P)+\Theta(\bar P)
\]

%dsp-level:0
\[
= \Theta(\Gamma)-\Theta(c_1)-\Theta(c_2)
\]

%dsp-level:0
\[
\in 2\mathbb Z;
\]

%dsp-level:0
\[
\quad\Theta(\bar P)
\]

%dsp-level:0
\[
= \Theta(P)\ [\mathrm{FT.23}\ \text{per edge}]\implies\Theta(P)
\]

%dsp-level:0
\[
= d_{12}\in\mathbb Z;
\]

%dsp-level:0
\[
\quad k(\Gamma)=\tfrac{\Theta(\Gamma)}{2}=k_1+k_2+d_{12}
\]

%dsp-level:0
\[
\text{per-edge}:\ \text{each}\ f_j\ \text{needs its own round-trip or digon materialization}\ [\mathrm{FT.34}(4)/\mathrm{FT.24}]:\ \text{sum integrality does not imply edge integrality}\ [\text{scope stated}]
\]

%dsp-level:0
\[
(3):\ n_{\mathrm{junction}}=d+2\ [\mathrm{D0.8}:\ \text{bridge adds 1 in}+1\ \text{out}];
\]

%dsp-level:0
\[
\quad d\ge2\implies n\ge4:\ O(n)=2^{n+1}-16>0\ [\mathrm{CT0.6}];\quad\mathrm{CT4.1}:\ D^-\text{-configured};
\]

%dsp-level:0
\[
\qquad s\ \text{triangles at one junction}:\ 2s\ \text{channels}
\]

%dsp-level:0
\[
\implies n
\]

%dsp-level:0
\[
= 4s,
\]

%dsp-level:0
\[
\ O(4s)
\]

%dsp-level:0
\[
= 2^{4s+1}-16\ [\mathrm{D0.8},\ \mathrm{CT0.6}];
\]

%dsp-level:0
\[
\quad\text{rate}\ \propto O(n)\ [\mathrm{D4.2}]
\]

%dsp-level:0
\[
(4):\ E_{\mathrm{int}}:\ \text{inclusion–exclusion on directed edges};
\]

%dsp-level:0
\[
\qquad \text{constraints}:\ \text{each materialized}\ c:\ \Theta_c
\]

%dsp-level:0
\[
= \sum_{e\in c}\kappa_{\mathrm{ch}(e)}\ (\mathrm{mod}\ 2)
\]

%dsp-level:0
\[
= 0\ [\mathrm{FT.22}/\mathrm{FT.23}];
\]

%dsp-level:0
\[
\quad\text{rank over}\ \mathbb Z_2\ \text{of the system}
\]

%dsp-level:0
\[
= \text{independent constraint count};
\]

%dsp-level:0
\[
\ \text{index}
\]

%dsp-level:0
\[
= 2^{\operatorname{rk}}
\]

%dsp-level:0
\[
\text{dependency witness}:\ \Theta_{ACBD}
\]

%dsp-level:0
\[
= \Theta_{ABC}+\Theta_{ABD}-2\kappa_{AB}\ [\mathrm{D1.7}\ \text{additivity on shared channel}]\implies\text{boundary-cycle constraint derived}
\]

%dsp-level:0
\[
(5):\ \text{bowtie}:\ E_{\mathrm{int}}
\]

%dsp-level:0
\[
= 2\cdot6
\]

%dsp-level:0
\[
= 12;
\]

%dsp-level:0
\[
\ n(X)
\]

%dsp-level:0
\[
= 8\ [\text{4 channels}];
\]

%dsp-level:0
\[
\ O(8)
\]

%dsp-level:0
\[
= 2^9-16
\]

%dsp-level:0
\[
= 496;
\]

%dsp-level:0
\[
\quad\text{constraints}:\ \Theta_{XAB},
\]

%dsp-level:0
\[
\Theta_{XCD}
\]

%dsp-level:0
\[
\in 2\mathbb Z:\ \text{disjoint channel sets}\implies\operatorname{rk}\qquad = 2;
\]

%dsp-level:0
\[
\quad g(5)=10\ [\mathrm{FT.41}(b)]
\]

%dsp-level:0
\[
\text{patch}:\ E_{\mathrm{int}}
\]

%dsp-level:0
\[
= 2\cdot5
\]

%dsp-level:0
\[
= 10;
\]

%dsp-level:0
\[
\quad\text{constraint vectors on}\ (AB,AC,BC,AD,BD):\ (1,1,1,0,0),
\]

%dsp-level:0
\[
\ (1,0,0,1,1):\ \text{independent over}\ \mathbb Z_2\implies\operatorname{rk}
\]

%dsp-level:0
\[
= 2,\qquad \ \text{index}\ 4;
\]

%dsp-level:0
\[
\quad n(A)=n(B)=6,\ n(C)=n(D)=4;
\]

%dsp-level:0
\[
\quad g(4)=3\ [\text{all}\ 2{+}2].\ \blacksquare
\]

\textbf{Phys.} Floors compose like ledger objects should: masses add across bridges, winding numbers add across round-trips, and every junction pays for the composition in darkness — the absorption rate doubling with each shared channel, exponentially with multiplicity. Gluing identifies shared channels and adds parity constraints by rank, with homological dependencies automatic. The bowtie and the patch are the first wired cells: one a five-vertex all-locked object with ten sealing patterns, the other a four-vertex meson–meson floor whose boundary-cycle constraint the ledger already implies.

\vspace{2pt}\noindent\rule{\textwidth}{0.4pt}\vspace{2pt}

\end{theorem}

\begin{theorem}[{\texorpdfstring{FT.51 — Enumeration of the winding ledger}{FT.51 - Enumeration of the winding ledger}}]\label{tag:FT-51}
\textbf{Thm.} For the polygon floor $\Sigma_{P_n}$ under M1/M2 [FT.45], at winding level $k$ (face constraint $\sum_i\kappa_i=2k$ [FT.45(3)]):

%dsp-level:0
\[
(1)\quad\#\{\text{on-shell profiles},\ \kappa\in\mathbb Z_{\ge0}^n,\ \textstyle\sum\kappa_i=2k\}=\binom{2k+n-1}{n-1}\qquad[\text{stars-and-bars};\ \text{math}];
\]

%dsp-level:0
\[
(2)\quad\#\{\text{strictly-spacelike profiles},\ \kappa\ge1\}=\binom{2k-1}{n-1}\quad(2k\ge n;\ =0\ \text{otherwise});
\]

%dsp-level:0
\[
(3)\quad\textbf{coarse-grained degeneracy}:\ \text{configurations related by mask-preserving relabeling are identified [CT2.6's}\ \sim\ \text{; CT2.7: weights are vertex-name independent], so the distinct ledger states at level}\ k\ \text{are the dihedral necklaces}
\]

%dsp-level:0
\[
M(n,k)=\frac{1}{2n}\Bigl(\sum_{j=0}^{n-1}\mathrm{fix}(r^j)+\sum_{\text{reflections}}\mathrm{fix}\Bigr)\qquad[\text{Burnside};\ \text{math}],
\]

with $\mathrm{fix}(r^j)=\binom{2k\,c/n+c-1}{c-1}$ if $(n/c)\mid 2k$, else $0$, where $c=\gcd(n,j)$ [rotation fixed points: $\kappa$ constant on the $c$ orbits]; machine values: $M(3,1)=2,\ M(3,2)=4,\ M(3,3)=7;\quad M(4,1)=3,\ M(4,2)=8$;

%dsp-level:0
\[
(4)\quad\textbf{growth}: \#\text{profiles}\sim(2k)^{n-1}/(n-1)!\ \text{as}\ k\to\infty\ [\text{math}]:\ \text{the winding ledger's state count grows polynomially in the winding at fixed}\ n,
\]

%dsp-level:0
\[
\ \text{degree}\ n{-}1\ —\ \text{a derived density of states for the floor family};
\]

%dsp-level:0
\[
(5)\quad\text{sign registers factor}: \sigma\text{-assignments are independent of}\ \kappa\ [\mathrm{CT0.10}]:\ \text{total floor states at level}\ k
\]

%dsp-level:0
\[
= 2^{2n}\cdot M(n,k)\cdot(\text{sealing/attachment data})\ —\ \text{the three registers multiply, none converts into another}.
\]

\textbf{Req.} FT.45; CT2.6, CT2.7 (Part 2); CT0.9, CT0.10 (Part 1); stars-and-bars, Burnside, asymptotics of binomials [math].

\textbf{Pf.}

%dsp-level:0
\[
(1),(2):\ \text{compositions of}\ 2k\ \text{into}\ n\ \text{non-negative (resp. positive) parts}\ [\text{math}];
\]

%dsp-level:0
\[
\qquad \text{parity automatic: the constraint fixes}\ \textstyle\sum\kappa=2k\ \text{even}\ [\mathrm{FT.45}(3)]
\]

%dsp-level:0
\[
(3):\ \text{relabeling quotient}:\ \mathrm{CT2.6}'\text{s}\ \sim\ \text{on coarse-grained ledger entries};
\]

%dsp-level:0
\[
\ \mathrm{CT2.7}\ \text{Pf}:\ \text{weights}\ \neg\text{vertex-name dependent}
\]

%dsp-level:0
\[
\implies D_n\text{-orbits are the distinct states};
\]

%dsp-level:0
\[
\qquad \text{Burnside with rotation/reflection fixed-point counts}\ [\text{math};\ \text{machine cross-check by orbit enumeration}]
\]

%dsp-level:0
\[
(4):\ \binom{2k+n-1}{n-1}=\frac{(2k+n-1)\cdots(2k+1)}{(n-1)!}\sim\frac{(2k)^{n-1}}{(n-1)!}
\]

%dsp-level:0
\[
(5):\ \mathrm{CT0.10}:\ P(C)\ \text{and}\ \pi(C)\ \text{independent sign structures};
\]

%dsp-level:0
\[
\ \ \mathrm{CT0.9}:\ 2^{2n}\ \text{sign assignments on}\ 2n\ \text{directed edges};
\]

%dsp-level:0
\[
\qquad \kappa\text{-ledger}\ [\text{clauses}(1)\text{--}(3)]\ \text{and sealing data}\ [\mathrm{FT.41}(b)]\ \text{factor as independent registers}.\ \blacksquare
\]

\textbf{Phys.} The winding ledger is countable in closed form: at winding $k$ a polygon floor carries $\binom{2k+n-1}{n-1}$ integer profiles, collapsing to dihedral necklaces under the corpus's own relabeling equivalence — two, four, seven distinct states for the triangle at $k=1,2,3$. The count grows like $(2k)^{n-1}$: bigger polygons remember more, higher windings remember denser — a derived density of states, with signs and sealings multiplying as independent registers. This is the combinatorial substrate any generation/flavour reading  would have to select from.

\vspace{2pt}\noindent\rule{\textwidth}{0.4pt}\vspace{2pt}

\end{theorem}

\begin{theorem}[{\texorpdfstring{FT.63 — The Platonic physical cards: spin, charge, parity, holonomy, entropy, mass ratios, binding}{FT.63 - The Platonic physical cards: spin, charge, parity, holonomy, entropy, mass ratios, binding}}]\label{tag:FT-63}
\textbf{Thm.} For each Platonic floor $\Sigma_S$ [Def-H.20] under M1/M2 (interface clause), FT.62's corrected rungs, and — for the spin and defect-ladder rows — Def-H.38/Def-H.37(2) (interface clause): all seven requested computations close, per solid and universally. Write $(V,E,F,L,d)$ for the cardinals, $\kappa_{0,\min}=1$ (cube) / $2$ (all others) [FT.42(5)], $K_{S,\min}=E\kappa_{0,\min}$, $\delta_v=2\pi-d\,\frac{(L-2)\pi}{L}$ the vertex angle defect, and $k_{f,\min}=L$ ($L$ odd) / $L/2$ ($L$ even) the uniform-minimal face winding. $(1)$ \textbf{Spin} [Def-H.38; Def-H.23, FT.60]: every materialized face chain carries winding $k_f$ with carrier $\mathrm{Sym}^{k_f}$ and label $j_f=|k_f|/2$; every channel digon carries $k=\kappa_{\mathrm{ch}}$, $j=\kappa_{\mathrm{ch}}/2$ [FT.24(1), FT.25(3)]. At the uniform minima:

%dsp-level:0
\[
\begin{array}{c|c|c|c|c|c}
S & k_{f,\min} & j_f & K_{S,\min} & j_{\max}=K_S/2 & \dim:\ \text{face basis}\ /\ \text{digon basis}\\ \hline
\text{tetrahedron} & 3 & \tfrac32 & 12 & 6 & 4^4=256\ /\ 3^6=729\\
\text{cube} & 2 & 1 & 12 & 6 & 3^6=729\ /\ 2^{12}=4096\\
\text{octahedron} & 3 & \tfrac32 & 24 & 12 & 4^8=65536\ /\ 3^{12}=531441\\
\text{dodecahedron} & 5 & \tfrac52 & 60 & 30 & 6^{12}=2176782336\ /\ 3^{30}\\
\text{icosahedron} & 3 & \tfrac32 & 60 & 30 & 4^{20}=2^{40}\ /\ 3^{30}
\end{array}
\]

with the digon-basis check $\sum_{\mathrm{ch}}\kappa_{0,\min}/2=K_{S,\min}/2=j_{\max}$ [machine]. \textbf{Derived laws}: (a) $j_{\max}=K_S/2$ and the total winding sum $\sum_fk_f=K_S$ are decomposition-invariant [FT.43(1), FT.60(7a)]; the full tensor carrier is basis-relative (digon basis is the edge-disjoint one [FT.25(1)]). (b) $V,F$ even for all five solids $\Rightarrow$ at uniform minima every face $j_f$ composes to an \textbf{integer} stretched total ($F$ half-integer-or-integer summands, $F$ even) — Platonic floors are integer-spin objects in the stretched read at their minima, with half-integer face reads for the four odd-$L$ solids: fermion-like faces, boson-like stretched totals (structural observation, $\neg$particle identification). (c) Dual splitting, new column: $j_{\max}$ splits cube ($6$) from octahedron ($12$) but is \textit{degenerate} for dodecahedron/icosahedron ($30$) — where the per-face structure splits instead ($5/2\times12$ vs $3/2\times20$) and the face-basis carrier dimensions split ($6^{12}$ vs $4^{20}$) while the digon-basis dimensions are degenerate ($3^{30}$ both). $(2)$ \textbf{Charge} [FT.40, EM.9, D3.4]: the closed-core read $q_{\mathrm{core}}(c;v)=\tfrac12(\sigma_{e_1}-\sigma_{e_k})\in\{-1,0,+1\}$ is defined on every materialized closed core — faces, digons, and homology sums: $\beta_1=F-1$ independent core reads per solid $(3,5,7,11,19)$, ternary alphabet, all three values attained per solid [machine: explicit $\sigma$-assignments on the encoded skeletons]; $\omega$-rebase flips the sign [D2.6, FT.40]. Channel-facing charge: dark floors — no read channel [FT.3; FT.62(5)]; observed variants — EM.9 end-reads of transit chains through the contact, telescoping to boundary signs [FT.21(i)]; EM activity criterion $\exists v:\ f_v\ne0$ on chain-local reads [FT.42(4), FT.62(5) scope]. No independent weak or strong charge exists [FT.11, FT.12, CT3.22]. $(3)$ \textbf{Parity}: the on-shell $\kappa$-lattice is $\ker\varphi_S$, $\varphi_S:\mathbb Z^E\to\mathbb Z_2^F$ the face-incidence map — index $2^{F-1}=(8,32,128,2048,2^{19})$ [FT.43(3); machine: GF(2) elimination for \textbf{all five}, extending the earlier tetra/cube]; equivalently $\kappa\bmod2$ ranges over the \textbf{cut space} of $\Gamma_S$ (dimension $V-1=E-F+1$ — Euler again) [math: planar duality of the face-even lattice; machine: dimension identity + star-cut witnesses]. Consequences [machine]: $A_S(1)=0$ for all five (no bridges: single-channel $\kappa{=}1$ leaves two faces odd); $A_S(3)=V$ for $d{=}3$ solids (the vertex stars), $0$ for octahedron/icosahedron-at-$K{=}3$ (odd $K$ needs an odd cut; $|\delta S|\equiv d|S|\bmod2$); \textbf{octahedron: every face $2$-coloring balances $\sum_{\mathrm{black}}k_f=\sum_{\mathrm{white}}k_f=K_S/2$, hence $K_S\in2\mathbb Z$ and $A_{\mathrm{octa}}(\text{odd}\,K)=0$} [machine; the unique Platonic with even $d$]. On materialized loops: $\Theta\in2\mathbb Z$, $\mathrm{Hol}_{\mathrm{int}}=+1$ [FT.22, FT.25(1), FT.42(3)]; partial internal sets retain free parity [FT.25(3)]; on-shell weight phases are registration-only, $(-1)^{\aleph+n_-}$ [FT.61(1)–(2)]. $(4)$ \textbf{Holonomy}: $\mathrm{Hol}_{\mathrm{spiral}}=+1$ on every materialized loop [FT.22]; $\mathrm{Hol}_{\mathrm{int}}=+1$ [FT.25(2)]. Frame/spin-connection side [Def-H.21, FT.54]: the skeleton metric is piecewise-flat with conical defects $\delta_v$ at vertices — exact table $\delta_v/\pi=(1,\tfrac12,\tfrac23,\tfrac15,\tfrac13)$, and \textbf{Descartes/Gauss–Bonnet closes the ledger}: $V\delta_v=4\pi$ for all five [machine in exact $\pi$-units]. Hence the global spinor transport around any full sweep is $\exp(i\,\tfrac{4\pi}{2}\,\hat n\cdot\vec\sigma)=+\mathbb I$ — $4\pi$ is the SU(2) period [machine over $\mathbb Q(i)$]: spherical materialization is spinor-single-valued, a consistency anchor for M1. Loops cleanly enclosing one vertex (transverse to the singular locus — vertex-figure loops in the derived geometry) carry $\Theta_R=\delta_v$ [FT.47(1)/GB; Def-H.21(6)] and defect values $D_S=2-2\cos(\delta_v/2)=4\sin^2(\delta_v/4)$ — exact radicals, clause (6). 1-skeleton face loops sit \textit{on} the singular locus (vertices on the loop): the smooth-enclosure defect read does not apply to them without a corner convention, undeclared — recorded, not guessed [Rmk below]. Under Def-H.38 the defect read becomes the multiplet ladder $m^2_{\mathrm{pole}}(m)=4\hat\mu'^2\sin^2(m\,\Theta_R/4)$, $m\in\{k,k{-}2,\dots,-k\}$ [FT.60(5)] — per-solid weight tables in (6). Weak-transport side: kernel coherence — charge-aligned profiles give $G(c)=0$, defect-massless, at every floor [FT.33(6), EW.22, EM.6]. $(5)$ \textbf{Entropy}: $S_{\mathrm{hid}}=I_{\mathrm{hid}}=2E\ln2=(12,24,24,60,60)\ln2$ nats [FT.42(1); FT.62(4): unchanged]; loop relating rate $S_{\mathrm{reg}}/A_{\mathrm{loop}}=(\ln2/\mathfrak h^2)/\bar\kappa$ [FT.37]: at the uniform minima $\bar\kappa=\kappa_{0,\min}$ — \textbf{the cube is the unique Platonic floor whose uniform-minimal loops saturate the CT1.29 Bekenstein rate $\ln2/\mathfrak h^2$ exactly}; the other four run at half rate ($\kappa_0=2$). Diagnostic minima $\Pi_{\min}=8E^2K_{S,\min}(\ln2)^2=(3456,13824,27648,432000,432000)(\ln2)^2$ — dual-degenerate at matched $K_S$ [FT.42(5), FT.44(1)]. Lock ledgers (corrected): $\sum_vO$ per variant from FT.62(2)'s table — the retained-but-inaccessible register rate, e.g. dark icosahedron $24384$/tick vs dark tetrahedron $448$/tick. Boundary entropy: dark floors $\partial\Sigma_S=\varnothing$; observed $E_{\mathrm{bd}}\in\{1,2\}$ [FT.42; CT1.29 applies to locked-interior separating surfaces — the Platonic lock ledger is vertex-absorption, not surface area; no identification]. $(6)$ \textbf{Mass ratios} [Def-H.37(2); FT.54(4)]: $m_{\mathrm{hid}}=2E\ln2/\mathfrak h$; the vertex-defect pole datum and the dimensionless ratio:

%dsp-level:0
\[
\frac{m^2_{\mathrm{pole}}(\delta_v)}{m_{\mathrm{hid}}^2}
\]

%dsp-level:0
\[
= \frac{4\hat\mu'^2\,\mathfrak h^2\sin^2(\delta_v/4)}{4E^2(\ln2)^2}
\]

%dsp-level:0
\[
= \frac{\hat\mu'^2\mathfrak h^2}{(\ln2)^2}\cdot\frac{D_S}{4E^2}\in\mathfrak q^0
\]

%dsp-level:0
\[
\qquad [[\hat\mu'^2]=\mathfrak q,\ [\mathfrak h^2]=\mathfrak q^{-1}:\ \mathrm{D2.9}],
\]

%dsp-level:0
\[
\begin{array}{c|c|c|c|c}
S & D_S=4\sin^2(\delta_v/4) & \text{exact field} & \text{min. polynomial of}\ D_S & D_S/4E^2\\ \hline
\text{tetrahedron} & 2 & \mathbb Q & x-2 & 1/72\\
\text{cube} & 2-\sqrt2 & \mathbb Q(\sqrt2) & x^2-4x+2 & (2-\sqrt2)/576\\
\text{octahedron} & 1 & \mathbb Q & x-1 & 1/576\\
\text{dodecahedron} & 2-2\cos\tfrac\pi{10} & \mathbb Q(\cos\tfrac\pi{10}) & x^4-8x^3+19x^2-12x+1 & D_S/3600\\
\text{icosahedron} & 2-\sqrt3 & \mathbb Q(\sqrt3) & x^2-4x+1 & (2-\sqrt3)/3600
\end{array}
\]

[machine: the dodecahedron quartic verified by exact substitution $x=2-2c$ into $\mathbb Z[c]/(16c^4-20c^2+5)$ — itself verified from $\cos5t=16c^5-20c^3+5c$ by de Moivre binomials at $t=\pi/10$; the cube/icosahedron quadratics from $(x-2)^2\in\{2,3\}$]. Weight-ladder values $D_S(m)=2-2\cos(m\,\delta_v/2)$, $m\le4$ [FT.60(5); machine]: tetrahedron $(2,4,2,0)$ — rational, with an exact zero at $m{=}4$; cube $(2{-}\sqrt2,2,2{+}\sqrt2,4)$; octahedron $(1,3,4,3)$ — rational; icosahedron $(2{-}\sqrt3,1,2,3)$; dodecahedron $D(2)=\tfrac{3-\sqrt5}2$, $D(4)=\tfrac{5-\sqrt5}2$ [machine in $\mathbb Q(\sqrt5)$ via $4c^2\mp2c-1$ at $c=\cos\tfrac\pi5,\cos\tfrac{2\pi}5$], $D(1),D(3)$ the two \textbf{smallest} of its four positive roots. \textbf{[precision correction [FT.76(10)]: the quartic $x^4-8x^3+19x^2-12x+1$ has four positive roots — $0.097886967410$, $0.824429495415$, $3.175570504585$, $3.902113032590$ (machine-isolated to width $2^{-60}$; $q(0..4)=(1,1,5,1,1)>0$, so no integer bracket separates them) — realized by the weights $m\in\{1,19\},\{3,17\},\{7,13\},\{9,11\}$ respectively; $D(1),D(3)$ are the two smallest.]} Def-H.37(1) form at stationarity: ratio $=\hat\mu^2\mathfrak h^4G(c)\,\kappa_{0,\min}/(\ln2)^2$ [FT.53(3)] with $\kappa_{0,\min}\in\{2,1,2,2,2\}$. \textbf{Derived law (dual splitting, pole side)}: both dual pairs share $m_{\mathrm{hid}}$ and $\Pi$ [FT.44(1)] but split on every defect datum — $\delta_v$: $\tfrac\pi2$ vs $\tfrac23\pi$; $\tfrac\pi5$ vs $\tfrac\pi3$; $D_S$: $2-\sqrt2$ vs $1$; $2-2\cos\tfrac\pi{10}$ vs $2-\sqrt3$ — FT.44(2)'s splitting template gains its first \textit{irrational} column. Face loops carry $\Theta_R=0$ on the flat-face side (defect-massless), corner convention undeclared [Rmk]. $(7)$ \textbf{Binding-force contributions}: per-aspect ledger, all five solids:

%dsp-level:0
\[
\begin{array}{c|c|c|c}
\text{aspect} & \text{criterion on}\ \Sigma_S & \dim\mathcal K^{(A)} & \text{binding read}\\ \hline
G & \text{always}\ (I_{\mathrm{hid}}=2E\ln2\ne0) & 5 & \kappa\text{-profile modulo loop constraint}\ [\mathrm{FT.41}(c)]\\
H & \text{always}\ (\mathrm{Def-H.36}\ \text{fold}) & 1 & m_{\mathrm{hid}}\ \text{read}\ [\mathrm{FT.20}]\\
W & \iff K_S\ne0\ (\text{at minima: always}) & 3 & \text{spin content of}\ (1)\ [\mathrm{FT.42}(3),\ \mathrm{Def-H.38}]\\
\mathrm{EM} & \text{observed/charged states}\ [\mathrm{FT.62}(5)] & 1 & q_{\mathrm{core}}\ \text{ternary per core}\ [(2)]\\
\text{strong} & \iff\ \text{inside-sealed phase}\ (\mathcal C=1) & 8 & g(V)\ \text{groupings; singlet-diagonal reads vanish}\ [\mathrm{FT.39}(2)]
\end{array}
\]

$\dim\mathcal K_{\Sigma_S}=9\ (+1)\ (+8)\le18$ [FT.26(3), FT.42(4)]. Sealing grammar: $g(V)=(3,385,25,292938721075,164395)$ [FT.41(b); machine: closed form vs enumeration for $V\le8$], with the \textbf{even-$\varepsilon$ theorem}: $V$ even for all five solids $\Rightarrow$ every grouping has an even number $b$ of $\varepsilon$-triples ($2a+3b=V$, $b$ even) [machine] — Platonic sealing phases are baryon-number-even by structure; the tetrahedron is meson-only ($b=0$: three $\delta\delta$-pairings), the octahedron admits $2\varepsilon$ phases ($15\,\delta^3+10\,\varepsilon^2$), the icosahedron $b\in\{0,2,4\}$, the dodecahedron $b\in\{0,2,4,6\}$, the cube $b\in\{0,2,4\}$ [machine: term breakdown]. Confinement is all-or-nothing [FT.19] with generic nonvanishing at $\mathcal C=1$ [FT.39(1)]; no per-edge strong phase exists [FT.29(2)]; strong binding contributes \textit{admissibility structure and gluon modes}, never singlet expectation values [FT.39(2), CT3.13]. Boundary/exchange algebra: dark $N(\Sigma_S)=0$ (no boundary algebra); single-directed-edge contact $N=1$: nothing strictly new [FT.41(d)]; bidirectional contact $N=2$ corolla at $n(w)=2$: exchange doublet $W$-hosted literally [FT.58(1)]; the face-floor variant $N=6$, $n(D)=6$: exchange sector wholly kernel/$D^-$-absorbed [FT.59(4a)]. $D^-$: never observable, absorption-only [FT.20, CT4.4] — its binding contribution is the lock ledger of FT.62(2), by rate not by force.

\textbf{Req.} Def-H.20, Def-H.23; Def-H.38, Def-H.37(2), Def-H.36, Def-H.35; M1/M2 (interface clause); FT.9, FT.11, FT.12, FT.19, FT.20, FT.21, FT.22, FT.24, FT.25, FT.26, FT.29, FT.31, FT.33, FT.37, FT.39, FT.40, FT.41, FT.42, FT.43, FT.44, FT.47, FT.49, FT.53, FT.54, FT.58, FT.59, FT.60, FT.61, FT.62; A4; D2.6, D2.9, D3.3, D3.4; CT0.6, CT0.7, CT0.9, CT0.10, CT0.12, CT0.13, CT0.14 (Part 1); CT1.29, CT2.5, CT2.9 (Part 2); CT3.13, CT3.17, CT3.22 (Part 2); CT4.3, CT4.4, CT4.9 (Part 2); EW.22, EW.32 (Part 2); EM.2, EM.6, EM.9 (Part 2); Descartes/Gauss–Bonnet, planar cut-space duality, $\cos(5t)$ de Moivre identity, SU(2) $4\pi$-periodicity [math].

\textbf{Pf.}

%dsp-level:0
\[
(1):\ \text{faces materialize}\ [\mathrm{M1}]:\ \text{each face}
\]

%dsp-level:0
\[
= \text{materialized closed chain}\implies\Psi
\]

%dsp-level:0
\[
= \mathrm{Sym}^{k_f}\ [\mathrm{Def\text{-}H.23}];
\]

%dsp-level:0
\[
\quad\text{digons}\ [\mathrm{M2}]:\ k
\]

%dsp-level:0
\[
= \kappa\ [\mathrm{FT.24}(1),\mathrm{FT.25}(3)]
\]

%dsp-level:0
\[
\text{uniform minima}:\ k_{f,\min}=\Theta_{f,\min}/2:\ L\ \text{odd}:\ \tfrac{2L}2=L;\ L\ \text{even}:\ \tfrac L2\ [\mathrm{FT.42}(5)];
\]

%dsp-level:0
\[
\qquad K_{S,\min}=E\kappa_{0,\min}\ [\mathrm{FT.42}(3)];\quad j_{\max}=K_S/2\ [\mathrm{FT.60}(2),(7a)]
\]

%dsp-level:0
\[
\dim\mathrm{Sym}^k=k+1:\ \text{table arithmetic}\ [\text{machine}];
\]

%dsp-level:0
\[
\qquad (b):\ F\in\{4,6,8,12,20\}\ \text{even};
\]

%dsp-level:0
\[
\ j_f\in\tfrac12\mathbb Z:\ \text{stretched total}
\]

%dsp-level:0
\[
= \textstyle\sum_fk_{f,\min}/2
\]

%dsp-level:0
\[
= K_{S,\min}/2\in\mathbb Z\ [\text{machine}:\ 6,6,12,30,30]
\]

%dsp-level:0
\[
(c):\ E\ \text{invariance}\ [\mathrm{FT.44}(1)];
\]

%dsp-level:0
\[
\quad j_{\max}:\ 6\ne12\ (\text{cube/octa});
\]

%dsp-level:0
\[
\quad 30=30\ (\text{dodeca/icosa}):\ k_{f,\min}\times F:\ 5\times12\ne3\times20\ \text{as multisets};\ 6^{12}\ne4^{20};\ 3^{30}=3^{30}
\]

%dsp-level:0
\[
(2):\ \mathrm{FT.40}\ \text{verbatim per core};
\]

%dsp-level:0
\[
\qquad \beta_1=F-1\ \text{independent reads}\ [\mathrm{FT.43}(2)];
\]

%dsp-level:0
\[
\qquad \text{attainment witnesses}\ [\text{machine}];
\]

%dsp-level:0
\[
\qquad \omega\text{-flip}\ [\mathrm{D2.6}]
\]

%dsp-level:0
\[
(3):\ \mathrm{FT.43}(3)+\text{machine elimination (all five)};
\]

%dsp-level:0
\[
\qquad \ker\varphi_S\ \text{mod}\ 2:\ \dim_{\mathbb Z_2}
\]

%dsp-level:0
\[
= E-(F-1)
\]

%dsp-level:0
\[
= V-1\ [\text{Euler}]
\]

%dsp-level:0
\[
= \dim\text{cut space}\ [\text{math}];
\]

%dsp-level:0
\[
\ \text{stars}
\]

%dsp-level:0
\[
= \delta\{v\}:\ \text{admissible}\ [\text{machine}]
\]

%dsp-level:0
\[
A_S(1)
\]

%dsp-level:0
\[
= 0:\ \text{single channel}\ \kappa{=}1:\ \text{its two faces odd};
\]

%dsp-level:0
\[
\ \neg\exists\ \text{bridges}\ [\text{solid 2-connectivity};\ \text{machine}];
\]

%dsp-level:0
\[
\qquad A_S(3)=V\ (d{=}3):\ \text{stars};\ \text{odd}\ K\iff\text{odd cut}:\ |\delta S|\equiv d|S|\ (2)\ \text{via}\ \textstyle\sum_{v\in S}\deg v=|\delta S|+2|E(S)|
\]

%dsp-level:0
\[
\text{octahedron}:\ \text{dual bipartite}
\]

%dsp-level:0
\[
\iff d\ \text{even}\ [\text{machine}: 2\text{-coloring exists iff octahedron}];
\]

%dsp-level:0
\[
\quad\text{each channel in one black} + \text{one white face}\implies\textstyle\sum_{\mathrm{bl}}\Theta_f
\]

%dsp-level:0
\[
= \sum_{\mathrm{wh}}\Theta_f\qquad = \sum_e\kappa_e
\]

%dsp-level:0
\[
\implies\textstyle\sum_{\mathrm{bl}}k_f=\sum_{\mathrm{wh}}k_f=K_S/2\in\mathbb Z;\qquad\text{on-shell phases}:\ \mathrm{FT.61}(1)\text{--}(2)
\]

%dsp-level:0
\[
(4):\ \delta_v
\]

%dsp-level:0
\[
= 2\pi-d\alpha,
\]

%dsp-level:0
\[
\ \alpha
\]

%dsp-level:0
\[
= \tfrac{(L-2)\pi}L\ [\text{regular}\ L\text{-gon angle};\ \text{math}]:\ \delta_v/\pi
\]

%dsp-level:0
\[
= 2-\tfrac{d(L-2)}L
\]

%dsp-level:0
\[
= (1,\tfrac12,\tfrac23,\tfrac15,\tfrac13);
\]

%dsp-level:0
\[
\qquad V\delta_v=4\pi\ [\text{machine}:\ \text{Descartes}]
\]

%dsp-level:0
\[
\exp(i\tfrac{4\pi}2\hat n\cdot\vec\sigma)=\cos(2\pi)+i\sin(2\pi)\hat n\cdot\vec\sigma=\mathbb I\ [\text{machine over}\ \mathbb Q(i)];
\]

%dsp-level:0
\[
\qquad \Theta_R=\delta_v\ \text{for clean enclosure}\ [\mathrm{FT.47}(1)/\mathrm{GB};\ \mathrm{Def\text{-}H.21}(6)]
\]

%dsp-level:0
\[
(5):\ \mathrm{FT.42}(1)+\mathrm{FT.62}(4);
\]

%dsp-level:0
\[
\qquad \bar\kappa
\]

%dsp-level:0
\[
= \kappa_{0,\min}\ [\text{uniform}]:\ \text{rate}
\]

%dsp-level:0
\[
= (\ln2/\mathfrak h^2)/\kappa_{0,\min}\ [\mathrm{FT.37}];
\]

%dsp-level:0
\[
\quad\kappa_0
\]

%dsp-level:0
\[
= 1\iff\text{cube}\iff\text{CT1.29 saturation}\ [\mathrm{FT.37}\ \text{iff-clause}]
\]

%dsp-level:0
\[
\Pi_{\min}=8E^2K_{S,\min}\ [\mathrm{FT.42}(5)];\qquad\text{lock sums}:\ \mathrm{FT.62}(2)\ [\text{machine}]
\]

%dsp-level:0
\[
(6):\ \mathrm{FT.54}(4)\ \text{at}\ \Theta_R=\delta_v;\quad m^2_{\mathrm{hid}}=4E^2(\ln2)^2/\mathfrak h^2\ [\mathrm{FT.42}(1)];
\]

%dsp-level:0
\[
\qquad D_S=2-2\cos(\delta_v/2)\ [\text{half-angle}];\ \text{table values}\ [\text{machine: radicals},\ \text{min. polys}]
\]

%dsp-level:0
\[
\text{quartic}:\ c
\]

%dsp-level:0
\[
= \cos\tfrac\pi{10}:\ \cos5t
\]

%dsp-level:0
\[
= 16c^5-20c^3+5c\ [\text{de Moivre};\ \text{machine}];
\]

%dsp-level:0
\[
\ t
\]

%dsp-level:0
\[
= \tfrac\pi{10}:\ \cos\tfrac\pi2
\]

%dsp-level:0
\[
= 0
\]

%dsp-level:0
\[
\implies 16c^4-20c^2+5\qquad = 0;
\]

%dsp-level:0
\[
\quad x=2-2c\implies x^4-8x^3+19x^2-12x+1=0\ [\text{machine: exact substitution}]
\]

%dsp-level:0
\[
\text{ladder}:\ D_S(m)=2-2\cos(m\delta_v/2)\ [\mathrm{FT.60}(5)];
\]

%dsp-level:0
\[
\qquad \text{table per solid}\ [\text{machine}:\ \mathbb Q,\ \mathbb Q(\sqrt2),\ \mathbb Q(\sqrt3),\ \mathbb Q(\sqrt5)\ \text{exact}]
\]

%dsp-level:0
\[
(7):\ \mathrm{FT.20}\ \text{rows}+\mathrm{FT.42}(4);
\]

%dsp-level:0
\[
\qquad g(V)=\!\!\sum_{2a+3b=V}\!\!\tfrac{V!}{a!\,b!\,2^a6^b}\ [\mathrm{FT.41}(b)];\quad V\ \text{even}\implies3b\ \text{even}\implies b\ \text{even};
\]

%dsp-level:0
\[
\qquad \text{term breakdowns}\ [\text{machine}]
\]

%dsp-level:0
\[
N(\Sigma_S):\ \text{dark}\ 0;\ \text{obs}\ 1\ \text{or}\ 2\ [\mathrm{Def\text{-}H.12}];
\]

%dsp-level:0
\[
\quad N{=}2\ \text{corolla at}\ n(w){=}2:\ \mathrm{FT.58}(1);\quad\text{face-floor variant}:\ \mathrm{FT.59}(4a).\ \blacksquare
\]

\textbf{Rmk (corner convention, undeclared).} Skeleton face-loops pass through conical vertices: their $\Theta_R$ requires a corner-distribution convention for defects on the loop (or the two-region read $0$ vs $4\pi-\delta$: the defect values differ by side, though the total $4\pi$ and the SU(2) global triviality are side-independent). This is the single place where the requested "mass ratio" computation meets an undeclared convention, and it is marked rather than papered over.

\textbf{Rmk (no identifications).} As everywhere in this document: no card entry is identified with a named particle; the dual-degeneracy/splitting pattern (mass and $\Pi$ degenerate; winding, sealing, locks, pole defects, and spin structure split) extends FT.44(2)'s template with two new columns and remains a template, not a generation claim .

\textbf{Phys.} Computed in full, the five solids read: spins — half-integer faces composing to integer stretched totals, six for the tetrahedron and cube, twelve for the octahedron, thirty for the last dual pair; charges — ternary core registers on every closed loop, three to nineteen independent reads, flipping under the corpus's own conjugation; parities — an index-$2^{F-1}$ lattice whose mod-2 classes are exactly the cuts of the skeleton, with the octahedron alone forced even-wound by its two-colored faces; holonomies — trivial spirals, trivial global spinor transport, and conical vertex defects summing to one sphere; entropies — twelve to sixty hidden bits, the cube alone running its minimal loops at the exact Bekenstein rate; mass ratios — the pole side written in $\sqrt2$, $\sqrt3$, and $\cos\frac\pi{10}$, splitting every dual pair that the hidden-bit mass degenerates; binding — gravity and Higgs always, weak whenever the ledger is wound, colour in three to two-hundred-ninety-three-billion sealing patterns, always baryon-even, never as a number a sealed floor carries.

\vspace{2pt}\noindent\rule{\textwidth}{0.4pt}\vspace{2pt}

\end{theorem}

\begin{theorem}[{\texorpdfstring{FT.64 — The filter chain and the per-solid loop lattices}{FT.64 - The filter chain and the per-solid loop lattices}}]\label{tag:FT-64}
\textbf{Thm.} $(1)$ \textbf{Filter chain (definition of ledger-admissible).} A loop datum on a floor $\Sigma$ is \textbf{ledger-admissible} iff it passes, in order: \textbf{(F0)} configuration validity — closed real-edge content, Janus balance, mask image, no dangling ledger boundary [D2.11, D1.3, D0.10]; \textbf{(F1)} registration — $\sigma_e\in\{\pm1\}$ per directed real edge, independent of $\kappa$ [CT0.9, CT0.10]; \textbf{(F2)} integrality — $\kappa_e\in\mathbb Z$ on edges of materialized digons and round-trip-closed channels (classes (i)/(ii) [FT.24, FT.34(4)]; Platonic/polygon floors: under declared M2 [Def-H.20]); \textbf{(F3)} loop quantization — $\Theta(c)\in2\mathbb Z$ per materialized closed chain [FT.22; M1 on faces]: for the Platonic floors exactly $\kappa\in\ker\varphi_S$, $\varphi_S=\,$face-incidence mod $2$ [FT.43(3)]; \textbf{(F4)} stationarity (semiclassical survivors) — $\kappa\equiv\kappa_0$ constant along each connected materialized chain, geodesic support [FT.53(1)–(2), FT.48(1)]; \textbf{(F5)} regime — null-permitting ($\kappa\ge0$), strictly-spacelike ($\kappa\ge1$), uniform, or signature-mixed ($\kappa\in\mathbb Z$; fibers infinite at fixed $K$ — witness $(K{+}n,-n)$ — so enumeration is scoped to $\kappa\ge0$ [FT.45(4) pattern]); \textbf{(F6)} relabeling quotient — identification under mask-preserving relabeling [CT2.6's $\sim$, CT2.7; FT.51(3)]: $\operatorname{Aut}(\Gamma_S)$-orbits (FT.65); \textbf{(F7)} register product — signs $\times$ $\kappa$-orbits $\times$ sealing groupings, independent registers [FT.51(5), FT.65(2)]; \textbf{(F8)} spectral filter — materialization (M1/M2 declared; which admissible loops are on-spectrum is the corpus's [D2.4, CT2.5, Def-H.16]); no ledger-side criterion can further select by stability index [FT.61(4)]: \textbf{(F8) is the honest boundary of this list}. $(2)$ \textbf{Per-solid lattice, all five by machine.} $\ker\varphi_S\subseteq\mathbb Z^E$ has index $2^{F-1}$ — GF(2) rank $F-1$ verified by explicit elimination for \textbf{all five} solids (the earlier tetra/cube checks extended): $(8,32,128,2048,2^{19})$. Equivalent formulation: $\kappa\bmod2$ lies in the \textbf{cut space} of $\Gamma_S$ — $\dim_{\mathbb Z_2}\ker\varphi_S=E-(F-1)=V-1=\dim\operatorname{Cut}(\Gamma_S)$ [Euler; math: planar duality; machine: dimensions + vertex-star cut witnesses]. $(3)$ \textbf{Face-read realizability (SNF).} For the face-read map $B:\mathbb Z^E\to\mathbb Z^F$, $\kappa\mapsto(\Theta_f)_f$: Smith invariant factors are $(1,1,1,2)$ tetrahedron, $(1^5,2)$ cube, $(1^7)$ octahedron (rank $7=F{-}1$), $(1^{11},2)$ dodecahedron, $(1^{19},2)$ icosahedron [machine]. Consequences: \textbf{(a)} tetrahedron/cube/dodecahedron/icosahedron: $(2\mathbb Z)^F\subseteq\operatorname{im}B$ — \textit{every} integer face-winding vector $k=(k_f)\in\mathbb Z^F$ is realized by some signed $\kappa$ [machine witness: all four unit vectors $k=e_f$ realized on the tetrahedron with $\kappa\in[-2,2]^6$]; \textbf{(b)} octahedron: $\operatorname{im}B=\ker\psi$, $\psi=$ black$-$white face-sum (dual bipartite, $d$ even — the unique such Platonic): $k$ realizable $\iff\sum_{\mathrm{bl}}k_f=\sum_{\mathrm{wh}}k_f\iff K_S\in2\mathbb Z$ [machine: $\psi(\operatorname{im}B)=0$ on samples; $\psi(2e_f)=\pm2\ne0$: $k=e_f$ unrealizable]; each channel lies in exactly one black and one white face [machine]. $(4)$ \textbf{Level-set counts.} $A_S(K):=\#\{\kappa\in\mathbb Z_{\ge0}^E:\ \sum_e\kappa_e=K,\ \varphi_S(\kappa)=0\}$ (level $K=K_S=\Theta_{\mathrm{int}}/2$) [machine, DFS with face-parity pruning; cross-checked against orbit enumeration and the Burnside identity element]:

%dsp-level:0
\[
\begin{array}{c|ccccccccc}
S & K{=}0 & 1 & 2 & 3 & 4 & 5 & 6 & 7 & 8\\ \hline
\text{tetrahedron} & 1 & 0 & 6 & 4 & 24 & 24 & 74 & 84 & 189\\
\text{cube} & 1 & 0 & 12 & 8 & 93 & 120 & 576 & — & —\\
\text{octahedron} & 1 & 0 & 12 & 0 & 84 & 0 & 452 & — & —\\
\text{dodecahedron} & 1 & 0 & 30 & 20 & 495 & 672 & — & — & —\\
\text{icosahedron} & 1 & 0 & 30 & 0 & 465 & 12 & — & — & —
\end{array}
\]

with the exact character-sum closed form [math: GF(2) orthogonality; machine-verified against the DFS for tetrahedron/cube/octahedron, $K\le6$]:

%dsp-level:0
\[
A_S(K)=\frac{1}{2^{F-1}}\sum_{s\in\mathbb Z_2^F/\langle\mathbf1\rangle}[z^K]\prod_{e\in E}\frac{1}{1-(-1)^{s\cdot b_e}\,z}
\]

%dsp-level:0
\[
\qquad (b_e=\text{face-incidence column of}\ e),
\]

and asymptotics $A_S(K)\sim\binom{K+E-1}{E-1}/2^{F-1}\sim K^{E-1}/((E-1)!\,2^{F-1})$ [equidistribution of the $F-1$ parity conditions; machine trend: tetrahedron $A(16)=2673$ vs $20349/8$]. Structural reads [machine]: $A_S(1)=0$ for every solid (no bridges — a lone $\kappa{=}1$ channel leaves its two faces odd); $A_S(3)=V$ for the $d{=}3$ solids (the vertex stars), $0$ for the $d\ge4$ solids (octahedron: even-$d$ parity; icosahedron: its stars carry $K{=}5$); the octahedron has $A(\text{odd}\,K)=0$ [(3)(b)]; the icosahedron's $A(5)=12$ is exactly its star family — the index-$2^{19}$ lattice is sparse enough that at low levels only the symmetric families survive. The $K{=}0$ entry is the all-null bound configuration (trivial-chain content [D1.7, D3.7]; $W$-row silent [FT.42(3)]). $(5)$ \textbf{Minimality regimes, all five} [FT.45(4)/FT.32(6) pattern; machine]:

%dsp-level:0
\[
\begin{array}{c|c|c|c|c}
S & \text{R0 null-perm.}: (K_{\min},\Theta_{\mathrm{int}},\text{orbits}) & \text{R1 strict}: (\varepsilon_{\min},K_{\min},\Theta_{\mathrm{int}},\text{orbits}) & \text{R2 uniform}: \kappa_0\ \text{ladder} & \text{R1 witness type}\\ \hline
\text{tetrahedron} & (2,4,1) & (2,8,16,1) & 2,4,6,\dots & \text{opposite-edge pair at}\ 2\\
\text{cube} & (2,4,1) & (0,12,24,1) & 1,2,3,\dots & \text{all-1 profile (unique)}\\
\text{octahedron} & (2,4,1) & (4,16,32,3) & 2,4,6,\dots & \text{coset leaders, 3 orbits}\\
\text{dodecahedron} & (2,4,1) & (6,36,72,25) & 2,4,6,\dots & \text{coset leaders, 25 orbits}\\
\text{icosahedron} & (2,4,1) & (10,40,80,14) & 2,4,6,\dots & \text{coset leaders, 14 orbits}
\end{array}
\]

R0: single channel at $\kappa{=}2$ (one orbit — $\operatorname{Aut}$ edge-transitive [machine]). R1: $\kappa=1+x$, $x\ge0$: parity demands $Bx\equiv(L\bmod2)\cdot\mathbf1$; $\varepsilon_{\min}=$ minimum weight in the GF(2) coset [machine: coset enumeration, $\le2^{19}$ for the dodecahedron]. \textbf{Derived laws}: (a) $\Theta_{\mathrm{int},\min}^{\mathrm{R0}}=4$ for every solid vs $2$ for $\Sigma_{\min}$ and $2$ (face-winding) for polygons — a solid's faces are too many to be fooled by one channel [$F\ge4$ separated faces; FT.45(5)'s digon distinction extended]; (b) the \textbf{cube is the unique Platonic whose uniform-minimal profile is also its strictly-spacelike minimum} ($\varepsilon_{\min}=0$; all-1, unique up to symmetry) — for the four odd-$L$ solids R1 beats R2 strictly ($16{<}24$, $32{<}48$, $72{<}120$, $80{<}120$): their uniform minima are \textit{not} their cheapest spacelike configurations, and the minimizers form nontrivial orbit families (up to $25$); (c) R2's stationary ladder is the semiclassical survivor list [FT.53]: one orbit per rung $\kappa_0$, $K=E\kappa_0$. $(6)$ \textbf{Completeness, ledger side.} (F0)–(F7) are decidable and decided above (or cited closed-form: polygons [FT.51], cycles [FT.32(6)]); (F8) is declared. The list is complete as a ledger enumeration: every admissible loop datum on every computed floor is an element of $\ker\varphi$ (lattice) $\times$ sign registers $\times$ sealing groupings, modulo $\operatorname{Aut}$, within one of the four regimes, surviving semiclassically exactly on the R2 ladder.

\textbf{Req.} FT.22, FT.24, FT.32, FT.34, FT.42, FT.43, FT.45, FT.48, FT.51, FT.53, FT.61, FT.62; Def-H.20; M1/M2 (interface clause); D0.10, D1.3, D1.7, D2.4, D2.11, D3.7; CT0.9, CT0.10 (Part 1); CT2.5, CT2.6, CT2.7 (Part 2); Def-H.16; GF(2) elimination, Smith normal form, planar cut-space duality, character orthogonality, stars-and-bars asymptotics [math].

\textbf{Pf.}

%dsp-level:0
\[
(1):\ \text{each step one citation as listed};
\]

%dsp-level:0
\[
\qquad \text{order}:\ \mathrm{F2/F3}\ \text{before}\ \mathrm{F4}\ [\text{integrality classes gate rigidity},\ \mathrm{FT.38}(2)/\mathrm{FT.53}];
\]

%dsp-level:0
\[
\ \mathrm{F8}\ \text{last}\ [\text{spectral},\ \text{declared}]
\]

%dsp-level:0
\[
(2):\ \mathrm{FT.43}(3)\ \text{method extended}:\ \text{explicit GF(2) elimination per solid}\ [\text{machine}];
\]

%dsp-level:0
\[
\qquad \dim\ker=E-\operatorname{rk}\varphi=E-(F-1)=V-1\ [\text{Euler}];\quad\operatorname{Cut}:\ \dim=V-1\ [\text{math}];\ \text{stars}\ \delta\{v\}\in\ker\varphi\ [\text{machine}]
\]

%dsp-level:0
\[
(3):\ \mathrm{SNF}\ [\text{machine}:\ \text{integer row/column reduction}];
\]

%dsp-level:0
\[
\qquad (a):\ \text{divisors}\in\{1,2\}\iff(2\mathbb Z)^F\subseteq\operatorname{im}B\ [\text{math}];\ \text{witnesses}\ [\text{machine}]
\]

%dsp-level:0
\[
(b):\ \text{dual bipartite}\iff d\ \text{even}\ [\text{machine}:\ 2\text{-coloring}];\quad\psi(B\kappa)=\sum_e\kappa_e(1-1)=0;\quad\psi(2e_f)=\pm2;
\]

%dsp-level:0
\[
\qquad K_S=\textstyle\sum_fk_f=2\sum_{\mathrm{bl}}k_f
\]

%dsp-level:0
\[
(4):\ \text{DFS with completed-face parity pruning}\ [\text{machine}];
\]

%dsp-level:0
\[
\qquad \text{character sum}:\ \tfrac1{2^{F-1}}\sum_s(-1)^{s\cdot B\kappa}\ \text{projection}\ [\text{math};\ \text{machine vs DFS}];
\]

%dsp-level:0
\[
\qquad \text{asymptotics}:\ s=0\ \text{term dominates}\ [\text{math};\ \text{machine trend}]
\]

%dsp-level:0
\[
A(1)=0:\ \text{single channel}:\ \text{its}\ 2\ \text{faces odd};\ \neg\exists\ \text{bridge}\ [\text{solid}\ 2\text{-edge-connectivity};\ \text{machine}];
\]

%dsp-level:0
\[
\qquad A(3)=V\ (d{=}3):\ \text{stars}\ \delta\{v\}:\ \text{each face at}\ v\ \text{meets}\ 2\ \text{star edges}
\]

%dsp-level:0
\[
(5):\ \text{R0 witness}\ \kappa=2e:\ B(2e)\equiv0;\ \text{orbit}=\text{point}\ [\text{edge-transitivity};\ \text{machine}];
\]

%dsp-level:0
\[
\qquad \text{R1}:\ \kappa=1+x:\ \Theta_f=L+\sum_{e\in f}x_e\in2\mathbb Z\iff Bx\equiv(L\bmod2)\mathbf1;\ \varepsilon_{\min}=\text{coset leader weight}\ [\text{machine}]
\]

%dsp-level:0
\[
\text{cube}:\ L\ \text{even}:\ x=0\ \text{admissible}:\ \varepsilon_{\min}=0,\ \kappa\equiv1\ \text{unique at}\ K=E;
\]

%dsp-level:0
\[
\qquad \text{odd}\ L:\ \varepsilon_{\min}\ge2\ (B x\equiv\mathbf1\ \text{needs}\ x\ne0);\ \text{witnesses as tabled}\ [\text{machine}]
\]

%dsp-level:0
\[
(6):\ \text{decidability}:\ \text{each filter arithmetic or cited-closed}.\ \blacksquare
\]

\textbf{Rmk (what the sparse low levels mean).} The icosahedron's lattice is so fine (index $2^{19}$) that below its star level $K{=}5$ only the symmetric families fit: $A(3)=0$, $A(5)=12$. Low-level admissibility on a solid is therefore \textit{highly non-generic} — most low profiles are excluded not by energy but by parity — a derived selection pressure of purely combinatorial origin that any spectral mechanism (F8) inherits for free. Recorded as structure; no physical reading claimed.

\textbf{Phys.} Run the machine and the five solids answer with one grammar: a lattice of index a power of two, whose odd classes are exactly the cuts of the skeleton; winding vectors free as integers on four of the five, balanced black-against-white on the octahedron; a cheapest null-tolerant loop of four quanta everywhere, cheapest honest loops that beat uniformity on every odd-faced solid; and a semiclassical list of survivors that is nothing but the uniform ladder — one orbit per rung. The list is complete up to exactly one door: which of these loops the spectrum keeps. That door was always the corpus's, and the ledger now proves it cannot be opened from this side.

\vspace{2pt}\noindent\rule{\textwidth}{0.4pt}\vspace{2pt}

\end{theorem}

\begin{theorem}[{\texorpdfstring{FT.65 — Coarse-grained degeneracies, register product, and the master loop list}{FT.65 - Coarse-grained degeneracies, register product, and the master loop list}}]\label{tag:FT-65}
\textbf{Thm.} $(1)$ \textbf{Degeneracy under the relabeling quotient.} $\operatorname{Aut}(\Gamma_S)$ acts on admissible profiles by channel permutation [CT2.6's $\sim$ on coarse-grained ledger entries; CT2.7: weights are vertex-name independent; FT.51(3) pattern]; orders $|\operatorname{Aut}|=(24,48,48,120,120)$, rotation subgroups $(12,24,24,60,60)$ [machine: enumeration + determinant filter on exact coordinates, $\mathbb Q$ and $\mathbb Q(\sqrt5)$]. The coarse-grained degeneracies $M_S(K)=\frac1{|\operatorname{Aut}|}\sum_g\operatorname{fix}(g,K)$ [Burnside; math] — with $\operatorname{fix}(g,K)$ the orbit-constant profile count [machine: orbit-DP with face-parity masks; identity element cross-asserted against $A_S(K)$; small-$K$ values cross-checked by direct orbit enumeration]:

%dsp-level:0
\[
\begin{array}{c|ccccccccc}
S & K{=}0 & 1 & 2 & 3 & 4 & 5 & 6 & 7 & 8\\ \hline
M\ \text{tetrahedron} & 1 & 0 & 1 & 1 & 4 & 2 & 8 & 6 & 17\\
M\ \text{cube} & 1 & 0 & 1 & 1 & 7 & 4 & 25 & — & —\\
M\ \text{octahedron} & 1 & 0 & 1 & 0 & 6 & 0 & 19 & — & —\\
M\ \text{dodecahedron} & 1 & 0 & 1 & 1 & 10 & 9 & — & — & —\\
M\ \text{icosahedron} & 1 & 0 & 1 & 0 & 9 & 1 & — & — & —\\ \hline
M^+\ \text{tetrahedron} & 1 & 0 & 1 & 1 & 4 & 2 & 11 & 7 & 23\\
M^+\ \text{cube} & 1 & 0 & 1 & 1 & 8 & 5 & 34 & — & —\\
M^+\ \text{octahedron} & 1 & 0 & 1 & 0 & 7 & 0 & 25 & — & —\\
M^+\ \text{dodecahedron} & 1 & 0 & 1 & 1 & 13 & 12 & — & — & —\\
M^+\ \text{icosahedron} & 1 & 0 & 1 & 0 & 12 & 1 & — & — & —
\end{array}
\]

$M^+$ (rotation-only quotient) is reported alongside $M$ (full automorphism quotient, the FT.51 convention for polygons): whether orientation-reversing relabelings identify states is the $\omega$/charge-conjugation read [D2.6] — both counts are derived data; no choice is imposed. Growth: $M_S(K)\sim A_S(K)/|\operatorname{Aut}|$ generically (free orbits dominate) [math; the fixed-point terms are lower-order]. $(2)$ \textbf{Register product (generalized).} Distinct ledger states of $\Sigma_S$ at level $K$:

%dsp-level:0
\[
\#\text{states}(K)=2^{2E}\cdot M_S(K)\cdot g(V)\cdot(\text{attachment-type data})
\]

%dsp-level:0
\[
\qquad [\mathrm{FT.51}(5)\ \text{pattern};\ \mathrm{CT0.9}/\mathrm{CT0.10}:\ \text{signs independent of}\ \kappa;\ \mathrm{FT.41}(b):\ \text{sealing register}],
\]

with $g(V)=(3,385,25,292938721075,164395)$ [FT.63(7)] — the three registers multiply, none converts into another; e.g. dark tetrahedron at $K{=}2$: $2^{12}\cdot1\cdot3=12288$ states before attachment data. $(3)$ \textbf{Density of states.} $\#\text{states}(K)\sim2^{2E}\,g(V)\,K^{E-1}/((E-1)!\,2^{F-1}\,|\operatorname{Aut}|)$ [from FT.64(4)'s asymptotics]: polynomial in the winding at fixed solid — degrees $E-1=(5,11,11,29,29)$ — with the Platonic prefactors spanning twelve orders of magnitude (tetrahedron $2^{12}\cdot3/24$ vs dodecahedron $2^{60}\cdot g(20)/120$): the derived state-count substrate any selection mechanism (Open Gaps 5/12, filter F8) would act on. $(4)$ \textbf{Master generalized loop list} (every cell computed in this document; $K=\Theta_{\mathrm{int}}/2$ levels; lattice index within $\mathbb Z^{E_{\mathrm{ch}}}$; all entries machine-verified or cited-closed):

%dsp-level:0
\[
\begin{array}{l|c|c|c|c|c|c|c}
\text{cell} & E_{\mathrm{int}} & m_{\mathrm{hid}}\cdot\tfrac{\mathfrak h}{\ln2} & \text{lattice index} & K_{\min}\ (\Theta_{\mathrm{int},\min}) & \Theta_{\min}^{\mathrm{strict}} & \kappa_0\ \text{ladder}\ (\mathrm{F4}) & \text{quotient}\ \to\ M\\ \hline
\Sigma_{\min}\ [\mathrm{FT.28}] & 2 & 2 & 1\ (\kappa\in\mathbb Z) & 1\ (2) & 2\ (\kappa{=}1) & 1,2,3,\dots & \text{trivial}\to1\\
\Sigma_L\ \text{uni.}\ [\mathrm{FT.32}] & L & L & 2\ (\textstyle\sum\kappa\in2\mathbb Z) & 1\ (2) & L{+}(L\bmod2) & \text{even}\ (L\ \text{odd})\ /\ \text{any}\ (L\ \text{even}) & \mathbb Z_L\to\text{necklaces}\\
\Sigma_{P_n}\ \text{bidir.}\ [\mathrm{FT.45}] & 2n & 2n & 2 & 2\ (4) & 2n{+}2(n\bmod2) & 1,2,\dots\ (n\ \text{ev})\ /\ 2,4,\dots\ (n\ \text{odd}) & D_n\to M(n,k)\ [\mathrm{FT.51}]\\
\Sigma_{12}\ [\mathrm{FT.35}] & 6 & 6 & 1\ (\mathbb Z^3) & 1\ (2) & 3\ (6) & \text{free}\ \mathbb Z^3\ (\text{digon quanta}) & \mathbb Z_2\to\text{swap-orbits}\\
\Sigma_\triangle\ [\mathrm{FT.49}] & 6 & 6 & 2 & 2\ (4) & 4\ (8) & 2,4,\dots & D_3\to M(3,\cdot)=(2,4,7)\\
\text{bowtie}\ [\mathrm{FT.50}] & 12 & 12 & 4 & 2\ (4) & 8\ (16) & (2a,2b)\ \text{per face} & |\operatorname{Aut}|=8\ (\text{not tabulated})\\
\text{double-triangle}\ [\mathrm{FT.50}] & 10 & 10 & 4 & 2\ (4) & 6\ (12) & 2,4,\dots & \mathbb Z_2\times\mathbb Z_2\ (|\operatorname{Aut}|=4,\ \text{not tabulated})\\ \hline
\text{tetrahedron} & 12 & 12 & 8 & 2\ (4) & 8\ (16) & 2,4,6,\dots & S_4\to\text{table}\ (1)\\
\text{cube} & 24 & 24 & 32 & 2\ (4) & 12\ (24) & 1,2,3,\dots & S_4\times\mathbb Z_2\to\text{table}\\
\text{octahedron} & 24 & 24 & 128\ (K\ \text{even}) & 2\ (4) & 16\ (32) & 2,4,6,\dots & S_4\times\mathbb Z_2\to\text{table}\\
\text{dodecahedron} & 60 & 60 & 2048 & 2\ (4) & 36\ (72) & 2,4,6,\dots & A_5\times\mathbb Z_2\to\text{table}\\
\text{icosahedron} & 60 & 60 & 2^{19} & 2\ (4) & 40\ (80) & 2,4,6,\dots & A_5\times\mathbb Z_2\to\text{table}
\end{array}
\]

(Columns: $E_{\mathrm{int}}=$ directed internal edges; index $=$ of the on-shell $\kappa$-lattice in the channel lattice; $K_{\min}$ null-permitting with $\Theta_{\mathrm{int},\min}=2K_{\min}$; strict $=$ minimal $\Theta_{\mathrm{int}}$ at $\kappa\ge1$ with its $K$; $\kappa_0$-ladder $=$ the stationary (F4) survivors; quotient group with degeneracy source. For the solids, $\Theta_{\mathrm{int},\min}^{\mathrm{R0}}=4$ universally and the strict minima are the FT.64(5) coset-leader values.) $(5)$ \textbf{Taxonomy law, re-anchored.} Across every computed cell and family (the twelve master-table rows below plus the $\Sigma_L$/$\Sigma_{P_n}$ members): $m_{\mathrm{hid}}=E_{\mathrm{int}}\ln2/\mathfrak h$ with $E_{\mathrm{int}}=2E_{\mathrm{bidir}}+E_{\mathrm{uni}}$ [FT.44(4) — extended from eleven cells to sixteen, zero exceptions], and the list is closed under the program's composition operations: bridged unions add masses and windings [FT.50(1)–(2)], gluings compose lattices by constraint rank [FT.50(4)], round-trips add connector distance [FT.34(3)] — every composite of listed cells is itself a listed-type cell with computable invariants.

\textbf{Req.} FT.28, FT.32, FT.34, FT.35, FT.41(b), FT.42–FT.45, FT.49–FT.51, FT.63, FT.64; CT0.9, CT0.10 (Part 1); CT2.6, CT2.7 (Part 2); D2.6; Burnside, orbit-stabilizer, asymptotics [math]; Def-H.20; M1/M2 (declared, solids/polygon rows).

\textbf{Pf.}

%dsp-level:0
\[
(1):\ \operatorname{Aut}\ \text{by adjacency backtracking}\ [\text{machine}:\ \text{orders}\ 24,48,48,120,120];
\]

%dsp-level:0
\[
\quad\text{rotations}:\ \det
\]

%dsp-level:0
\[
= +1\ \text{on exact coordinate realization}\ [\text{machine}:\ \mathbb Q\ \text{and}\ \mathbb Q(\sqrt5)\ \text{linear algebra}]
\]

%dsp-level:0
\[
\text{Burnside}:\ \operatorname{fix}(g,K)
\]

%dsp-level:0
\[
= \#\{\text{orbit-constant}\ \kappa\ge0:\ \textstyle\sum=K,\ \varphi=0\}\ [\text{orbit-DP with parity masks};\ \text{machine}];
\]

%dsp-level:0
\[
\qquad \operatorname{fix}(\mathrm{id},K)=A_S(K)\ [\text{cross-asserted}];\quad M\in\mathbb Z\ [\text{asserted}]
\]

%dsp-level:0
\[
\text{cross-checks}:\ \text{direct orbit enumeration}\ K
\]

%dsp-level:0
\[
\le 3\ (\text{tetra/cube/octa}),
\]

%dsp-level:0
\[
\ K
\]

%dsp-level:0
\[
\le 2\ (\text{dodeca/icosa})\ [\text{machine}];
\]

%dsp-level:0
\[
\qquad M^+\ge M\ [\text{smaller group}]
\]

%dsp-level:0
\[
(2):\ \mathrm{CT0.9}:\ 2^{2E}\ \text{sign registers};
\]

%dsp-level:0
\[
\ \mathrm{CT0.10}/\mathrm{FT.51}(5):\ \text{independence};
\]

%dsp-level:0
\[
\ \mathrm{FT.41}(b):\ g(V);
\]

%dsp-level:0
\[
\ \mathrm{D3.8}:\ \text{attachments};
\]

%dsp-level:0
\[
\qquad \text{example}:\ 2^{12}\cdot1\cdot3=12288\ [\text{machine}]
\]

%dsp-level:0
\[
(3):\ \text{substitution of FT.64}(4)\ \text{asymptotics};
\]

%dsp-level:0
\[
\qquad (4):\ \text{rows}:\ \mathrm{FT.28}(2),
\]

%dsp-level:0
\[
\mathrm{FT.32}(1),
\]

%dsp-level:0
\[
(6),
\]

%dsp-level:0
\[
\mathrm{FT.45}(1),
\]

%dsp-level:0
\[
(3),
\]

%dsp-level:0
\[
(4),
\]

%dsp-level:0
\[
\mathrm{FT.35}(1)\text{--}(3),
\]

%dsp-level:0
\[
\mathrm{FT.49}(2),
\]

%dsp-level:0
\[
(4),
\]

%dsp-level:0
\[
\mathrm{FT.50}(5),
\]

%dsp-level:0
\[
\mathrm{FT.62}\text{--}\mathrm{FT.64}
\]

%dsp-level:0
\[
\Sigma_{\min}\ \text{index}\ 1:\ \Theta
\]

%dsp-level:0
\[
= 2\kappa
\]

%dsp-level:0
\[
\in 2\mathbb Z\ \text{automatic}\ [\mathrm{FT.24}];
\]

%dsp-level:0
\[
\quad\Sigma_L:\ \{\textstyle\sum\kappa\in2\mathbb Z\}\ \text{index}\ 2\ [\mathrm{FT.32}(6)\ \text{scope}];
\]

%dsp-level:0
\[
\quad\Sigma_{12}:\ \mathbb Z^3\ [\mathrm{FT.35}(1),(3)]
\]

%dsp-level:0
\[
\text{bowtie/patch}:\ \text{index}\ 4\ [\mathrm{FT.50}(5):\ \mathbb Z_2\text{-rank}\ 2];
\]

%dsp-level:0
\[
\qquad (5):\ \text{per-cell counts cited};\ \text{composition laws}\ [\mathrm{FT.34}(3),\mathrm{FT.50}(1),(2),(4)].\ \blacksquare
\]

\textbf{Rmk (the two untabulated quotients).} Bowtie ($|\operatorname{Aut}|=8$: the degree-4 junction is fixed; the two pairs of leaf vertices swap freely, and the two triangles exchange) and double-triangle patch ($\mathbb Z_2\times\mathbb Z_2$: $C\leftrightarrow D$, $A\leftrightarrow B$) degeneracy tables are computable by the same machine and are deferred only for space; their lattice indices, minima, and group orders are in the table. Quotient convention throughout: automorphisms of the floor's internal skeleton (CT2.7's vertex-name independence; external contacts are configuration data, not relabeled). No structural claim rests on the omission.

\textbf{Phys.} The generalized list, read aloud: every fermion floor the program has computed admits loops on one lattice per topology — index one for the digon and the wired pair, index two for cycles and polygons and the tetrahedron face, index four for the glued pairs, eight to half a million for the solids — and on each lattice the admissible loops come in four regimes, collapse to orbits under the skeleton's own symmetries, and multiply by two independent sign-and-sealing registers. Twelve rows, one mass law, closed under composition. What the list does not contain is its own last filter: which loops the spectrum keeps. The ledger side is finished; the door is the corpus's.

\vspace{2pt}\noindent\rule{\textwidth}{0.4pt}\vspace{2pt}

\end{theorem}

\begin{theorem}[{\texorpdfstring{FT.69 — Platonic amplitude census completed; the cycle-alternation lemma; one census correction}{FT.69 - Platonic amplitude census completed; the cycle-alternation lemma; one census correction}}]\label{tag:FT-69}
\textbf{Thm.} Under Def-H.25's scopes (S1)–(S4): $(1)$ \textbf{Cap census, completed.} Bypass-configuration counts by $K$ (cap-DP, exact): octahedron $318\,523$ ($K\le9$; row $(1,24,276,2024,10206,35448,81344,110736,69504,8960)$; $K\le5$ exhaustively validated); dodecahedron $3\,067\,633\,654\,196\,912$ ($K\le30=3V/2$ forced); icosahedron $2\,483\,119\,531\,360$ ($K\le18=3V/2$ forced). Low-order universality: the rows of the two $60$-edge skeletons agree through $K=3$ ($(1,60,1770,34220)$: caps bind only from $K=4$: $487335$ vs $485115$, the difference $20\binom64$ vs $12\binom{10}4$ exactly the single-vertex lock counts) and diverge from $K=4$ on. Orbits: not attempted (bounded rule: cube's $33\,564$ already exceeded $10^4$ [FT.68(3)]). $(2)$ \textbf{Coverings.} Octahedron complete (two independent enumerators agree): $\mathcal Z_{\mathrm{byp}}(y)=1+72y+852y^2+1152y^3$, $\mathcal Z_{\mathrm{lock}}(y)=36y^2+1008y^3+840y^4$; $Z^{\pm}_{\mathrm{byp}}=87\,649$, $Z^{\pm}_{\mathrm{lock}}=280\,128$, $Z^{\pm}_{\mathrm{all}}=367\,777$. Dodecahedron/icosahedron exact to $|\mathcal C|\le3$: dodecahedron $\{0{:}1,\,1{:}120,\,2{:}5880{+}60_{\mathrm{lock}},\,3{:}153\,960{+}6140_{\mathrm{lock}}\}$ ($166\,161$ covers, $Z^{\pm}\ge9\,948\,001$ byp / $393\,920$ lock from the counted range); icosahedron $\{0{:}1,\,1{:}240,\,2{:}17\,520{+}120_{\mathrm{lock}},\,3{:}478\,360{+}27\,600_{\mathrm{lock}}\}$ ($523\,841$ covers, $Z^{\pm}\ge30\,896\,321/1\,768\,320$ counted). Single-centre closed form on every skeleton: $\#\{|\mathcal C|=1\}=V\,d_{\mathrm{in}}d_{\mathrm{out}}$ minus closed-thread exclusions $=V\,d(d{-}1)$ bidirectional $=(24,48,72,120,240)$ for the five solids [machine]. Full dodeca/icosa counts deferred: the vertex-DP state width (cut-arcs $\le11$ at degree $5$; memo exhaustion) and the branch enumerator ($\sim(1+d^2)^V$ raw) both exceed the bounded rule — recorded, not guessed. Fully crystallized floors cancel: types $(3,3)^{20}$, $(5,5)^{12}$, no centres, value $0$ [FT.67(3a)]. $(3)$ \textbf{Cycle-alternation lemma.} Let $R$ be a literal-domain cover and $\gamma\subseteq R$ a directed cycle. Then along $\gamma$ the vertices \textbf{strictly alternate centre / non-centre}; consequently: (i) $|\gamma|$ is even; (ii) every non-centre vertex of $\gamma$ carries at least one further real edge outside $\gamma$ (else its type on $\gamma$ alone is $(1,1)$, forcing a centre — contradiction); (iii) a pure directed cycle is never a cover [FT.67(3a), now with its structural reason]; (iv) \textbf{no odd directed cycle — in particular no triangle — lies inside any cover}: the amplitude layer's bound content is exactly even-cycle content, dissolved into threading \textbf{[scope: clauses (i)–(v) are literal-domain theorems. Both chain-local extensions dissolve odd cycles into open threads — on the tetrahedron, of $2397$ patterns containing a directed $3$-cycle, $833$ (E2) and $1025$ (E1) have nonzero amplitude against $0$ literal [FT.75(4)]]}; (v) census (distinct cycles, canonicalized; machine): covers containing directed cycles — tetrahedron $64/101$ ($84$ digon instances; $48$ bypass $+16$ locked), cube $9060/11205$ ($14\,328$ digons, $504$ four-cycles), octahedron $3054/3961$ ($4704$ digons, $672$ four-cycles), $P_3$–$P_7$ digons only; the $P_4$ face $4$-cycle lies in \textit{none} of its $33$ covers (clause (ii): its non-centre cycle vertices have no extra edges to carry); zero odd cycles in any enumerated cover. $(4)$ \textbf{FT.68(1) census correction (self-audit).} The tetrahedron census numbers stand ($501/252/90$ bypass configurations containing digons/$3$-cycles/$4$-cycles, re-derived). The clause "all of it cancels in amplitude" is \textbf{false for the digon row}: $48$ of the $501$ are covers. Witness (machine, brute-force sign sum $=$ cover value $=16$): $R=\{(0{\to}1),(1{\to}0),(2{\to}1),(3{\to}2)\}$, types $\bigl((1,1),(2,1),(1,1),(0,1)\bigr)$ — the digon $\{0,1\}$ dissolved: centre $0$ reads both digon edges, vertex $1$ is a $(2,1)$ non-centre supported by centre $2$'s thread. The $3$-cycle and $4$-cycle rows cancel completely [(iii)+(i) of clause (3); machine: $0$ nonzero of $252/90$]. **[addition: these $48$ survivors are killed by \textit{both} chain-local extensions ($N_{E2}=N_{E1}=0$ on all $48$, exhaustive), while the $16$ \textit{locked} digon-covers of clause (5) are killed by E2 and kept by E1 at the derangement values $(!a)4^a$ [FT.75(1),(4)]. The triangle rows revive in the extensions: $833$ (E2) / $1025$ (E1) of the $2397$ tetrahedron patterns containing a directed $3$-cycle.]\textbf{ Corrected statement: }bound cells cancel standalone; bound substructures survive inside covers only dissolved — even cycles threaded through supported non-centres.\textbf{ FT.68(1) patched in place (marker). $(5)$ }Minimal dissolution family.** The smallest dissolved-digon covers are double-digon hubs: two digons $\{h,c_1\},\{h,c_2\}$ sharing a locked $(2,2)$ hub $h$, centres $c_1,c_2$ ($|\mathcal C|=2$, $|R|=4$, value $16$): the tetrahedron's $12$ locked two-centre covers are exactly this family ($4$ hub choices $\times\binom32$); exhibited on cube and octahedron. The four $K_{1,3}$ stars [FT.68(1)] are the triple-digon hubs. Locked hubs are the only lock content visible in covers [FT.67 Rmk], and digons are the only bound cells they can host at $|\mathcal C|=2$.

\textbf{Req.} Def-H.25; FT.17, FT.22, FT.28, FT.32, FT.42, FT.45, FT.62, FT.63, FT.67, FT.68; A4, A8; D0.1, D0.8, D0.10, D1.4, D1.5, D3.3; CT0.11 (Part 1); CT4.1 (Part 2); cap-DP / cover-DP / branch enumeration, cycle census [math; machine].

\textbf{Pf.}

%dsp-level:0
\[
(1):\ \text{cap-DP over vertex-introduction order (state: active degrees $\le3$, by-$K$ accumulation)}\ [\text{machine exact}];
\]

%dsp-level:0
\[
\quad K\le\lfloor 3V/2\rfloor:\ \textstyle\sum_vk_v=2K\le3V
\]

%dsp-level:0
\[
\text{low-order rows}:\ K\le3:\ \text{no vertex can reach}\ k_v\ge4:\ c_K=\binom{2E}K;
\]

%dsp-level:0
\[
\qquad K=4:\ \binom{60}4-\#\{\text{4 edges at one vertex}\}:\ 20\binom64=300,\ 12\binom{10}4=2520\ \checkmark
\]

%dsp-level:0
\[
(2):\ \text{two enumerators: vertex-introduction DP (centre pre-claim masks, completion checks)}\
\]

%dsp-level:0
\[
= \ \text{branch enumeration over centre sets with claim conflicts pruned};
\]

%dsp-level:0
\[
\quad |\mathcal C|=1:\ \text{in-pick}\times\text{out-pick},\ \text{closed thread}\ (s{=}t)\ \text{excluded}\ \implies Vd(d{-}1)
\]

%dsp-level:0
\[
\text{floors}:\ \text{type}(d,d),\ d\ge2:\ \mathcal C=\varnothing,\ R\ne\varnothing:\ \text{cov}\equiv0:\ \text{value}\ 0\ [\mathrm{FT.67}(2)]
\]

%dsp-level:0
\[
(3):\ e_i=(v_i\to v_{i+1})\ \text{along}\ \gamma:\ \operatorname{cov}(e_i)=1:\ v_i\ \text{centre}\iff e_i\ \text{read at tail},\ e_{i-1}\ \text{read at head};
\]

%dsp-level:0
\[
\quad v_i\ \text{centre}\implies v_i\ \text{reads both}\ e_{i-1},e_i\implies v_{i\pm1}\ \text{read neither}\implies v_{i\pm1}\ \text{non-centre}
\]

%dsp-level:0
\[
v_i\ \text{non-centre}
\]

%dsp-level:0
\[
\implies e_{i-1}\ \text{read at}\ v_{i-1}\ (\text{centre}),
\]

%dsp-level:0
\[
\ e_i\ \text{read at}\ v_{i+1}\ (\text{centre}):\ \text{strict alternation};
\]

%dsp-level:0
\[
\quad\text{(i)}:\ \text{alternation on a cycle}
\]

%dsp-level:0
\[
\implies |\gamma|\ \text{even}
\]

%dsp-level:0
\[
\text{(ii)}:\ v\ \text{non-centre on}\ \gamma:\ \text{type on}\ \gamma\ \text{alone}
\]

%dsp-level:0
\[
= (1,1):\ \text{forced centre}\ [\mathrm{D3.3}\ \text{type-based}]\ \implies\ \exists\ \text{extra edge};
\]

%dsp-level:0
\[
\qquad \text{(iii)}:\ \text{pure cycle}:\ \text{all vertices on}\ \gamma,\ \text{alternation} + \text{(ii)}\ \text{impossible}
\]

%dsp-level:0
\[
\text{(iv)}:\ \text{odd}\ |\gamma|\ \text{cannot alternate};\qquad\text{(v)}:\ \text{machine census, canonical min-rotation cycles}
\]

%dsp-level:0
\[
(4):\ \text{witness types verified};
\]

%dsp-level:0
\[
\ \text{reads}:\ (0{\to}1),
\]

%dsp-level:0
\[
(1{\to}0)\ \text{at centre}\ 0;
\]

%dsp-level:0
\[
\ (2{\to}1)\ \text{at centre}\ 2;
\]

%dsp-level:0
\[
\ (3{\to}2):\ \text{centre}\ 2\ \text{in-pick}:\ \text{cov}
\]

%dsp-level:0
\[
\equiv 1\ \checkmark;
\]

%dsp-level:0
\[
\quad\text{brute force}
\]

%dsp-level:0
\[
= 16\ [\text{machine}]
\]

%dsp-level:0
\[
(5):\ |\mathcal C|
\]

%dsp-level:0
\[
= 2\ \text{with a digon}:\ \text{digon edges both read by one centre}\ c\ (\text{type}(1,1)\ \text{forces its only edges});
\]

%dsp-level:0
\[
\ \text{partner}\ h:\ \text{second digon's centre}\ c'\ \text{reads it};
\]

%dsp-level:0
\[
\ h\ (2,2)\ \text{locked};
\]

%dsp-level:0
\[
\ \text{tetrahedron count}\ 4\times\binom32\cdot\ldots
\]

%dsp-level:0
\[
= 12\ [\text{machine}].\ \blacksquare
\]

\textbf{Rmk (what the lemma settles for the two-layer picture).} The amplitude layer does not merely cancel bound \textit{cells} — it cancels every odd cycle and every undissolved even cycle. Bound content inside survivors is always \textit{threaded}: read exactly once per edge, which a closed chain read at both ends can never be. The spectral layer's bound floors (M1/M2) are therefore not in tension with the cover grammar — they are disjoint reads, and the tension is precisely the inter-family bridge.

\textbf{Phys.} Ask which flickering patterns carry nonzero ledger amplitude, and the answer now has a grammar: edges tile into open two-step threads, and any loop that hides inside the tiling must be even, must be read alternately, and must be propped open at every other vertex by extra threading. Triangles cannot hide at all. The smallest hiding place is a pair of digons on a locked hub — darkness threaded by transit, exactly the picture the locked-hub remark promised.

\vspace{2pt}\noindent\rule{\textwidth}{0.4pt}\vspace{2pt}

\end{theorem}

\begin{theorem}[{\texorpdfstring{FT.76 — The Platonic pole-mass multiplets; the bounded-hierarchy ceiling}{FT.76 - The Platonic pole-mass multiplets; the bounded-hierarchy ceiling}}]\label{tag:FT-76}
\textbf{Thm.} Under Def-H.37(2) and Def-H.38, with $\Theta_R=\delta_v$ [FT.63(4)'s vertex-figure read; the corner convention is \textit{not} declared in this part], $D_S(m):=m^2_{\mathrm{pole}}(m)/\hat\mu'^2=2-2\cos(m\,\delta_v/2)=4\sin^2(m\,\delta_v/4)$ [FT.60(5)], and $k$ the carrier winding [Def-H.38]: $(1)$ \textbf{Conventions and the two trivial readings.} $k_f=K_S/F=(3,2,3,5,3)$ (face winding; $K_S=(12,12,24,60,60)$ [FT.63(1), FT.43(1)], reproduced from the minimality rule $\kappa_{\min}=2$ odd-$L$ / $1$ even-$L$, $k=L\kappa_{\min}/2$ [FT.32(6), FT.45]); $k_{vf}=(3,3,2,3,5)$ (vertex-figure winding, same rule on the vertex-figure $L_{vf}=(3,3,4,3,5)$). \textbf{Variant (A)} $:=$ $(\delta_v,k_f)$ — the directive's reading; \textbf{variant (B)} $:=$ $(\delta_v,k_{vf})$ — the loop-consistent reading. Both other readings named by FT.63(4)/(6) are identically massless: $\Theta_R=0$ (flat-face side) and $\Theta_R=V\delta_v=4\pi$ (global sweep: $D(m)=2-2\cos(2m\pi)=0$ for every weight — FT.63(4)'s trivial global spinor transport). $(2)$ \textbf{Exact multiplets.} With $c=\cos\tfrac\pi{10}$ ($16c^4-20c^2+5=0$):

%dsp-level:0
\[
\begin{array}{c|c|c|c}
S & \text{(A)}\ k_f & m^2_{\mathrm{pole}}/\hat\mu'^2\ \text{on}\ \{k,k{-}2,\dots,-k\} & \text{field}\\
\hline
\text{tetra} & 3 & 2,\ 2,\ 2,\ 2\ (m=\pm1,\pm3) & \mathbb Q\\
\text{cube} & 2 & 2\ (m=\pm2),\ \mathbf 0\ (m=0) & \mathbb Q\\
\text{octa} & 3 & 4\ (m=\pm3),\ 1\ (m=\pm1) & \mathbb Q\\
\text{dodeca} & 5 & 2\ (m=\pm5),\ 2{-}2\cos\tfrac{3\pi}{10}\ (m=\pm3),\ 2{-}2c\ (m=\pm1) & \mathbb Q(c)\\
\text{icosa} & 3 & 2\ (m=\pm3),\ 2{-}\sqrt3\ (m=\pm1) & \mathbb Q(\sqrt3)
\end{array}
\]

%dsp-level:0
\[
\begin{array}{c|c|c|c}
S & \text{(B)}\ k_{vf} & m^2_{\mathrm{pole}}/\hat\mu'^2 & \text{field}\\
\hline
\text{tetra} & 3 & 2,\ 2,\ 2,\ 2 & \mathbb Q\\
\text{cube} & 3 & 2{+}\sqrt2\ (m=\pm3),\ 2{-}\sqrt2\ (m=\pm1) & \mathbb Q(\sqrt2)\\
\text{octa} & 2 & 3\ (m=\pm2),\ \mathbf 0\ (m=0) & \mathbb Q\\
\text{dodeca} & 3 & 2{-}2\cos\tfrac{3\pi}{10},\ 2{-}2c & \mathbb Q(c)\\
\text{icosa} & 5 & 2{+}\sqrt3\ (m=\pm5),\ 2\ (m=\pm3),\ 2{-}\sqrt3\ (m=\pm1) & \mathbb Q(\sqrt3)
\end{array}
\]

Decimals ($m^2/\hat\mu'^2$): dodeca $(0.097886967410,\ 0.824429495415,\ 2)$; icosa $(0.267949192431,\ 2)$, (B) adds $3.732050807569$; cube (B) $(0.585786437627,\ 3.414213562373)$. Every value is exact in the stated field and was cross-checked against an exact rational Taylor series for $\cos$ to $45$ digits ($125$ values, $m=0..24$, all five solids). $(3)$ \textbf{Periodicity, monotonicity, and the hard ceiling.} The ladder is periodic in the weight with period $P_S=4\pi/\delta_v=(4,8,6,20,12)$ and takes exactly $P_S/2+1=(3,5,4,11,7)$ distinct values, monotone increasing in $m$ on $[0,P_S/2]$; the maximum is $4\hat\mu'^2$ (attained at $m=P_S/2$; $P_S$ even for all five) and the minimum nonzero value is $D_S$ itself. Hence \textbf{every Platonic pole mass lies in $[0,4\hat\mu'^2]$} and the widest ratio attainable \textit{within} a solid is

%dsp-level:0
\[
\frac{4}{D_S}
\]

%dsp-level:0
\[
= \Bigl(2,\ \tfrac{4}{2-\sqrt2},\ 4,\ \tfrac{4}{2-2\cos\frac\pi{10}},\ \tfrac{4}{2-\sqrt3}\Bigr)
\]

%dsp-level:0
\[
= (2,\ 6.828427,\ 4,\ 40.863458,\ 14.928203)\ \ \text{in}\ m^2,
\]

maximal for the dodecahedron, whose value has minimal polynomial $x^4-48x^3+304x^2-512x+256$; in \textbf{mass} the ceiling is $\sqrt{4/D_S}=(1.414214,\ 2.613126,\ 2,\ 6.392453,\ 3.863703)$, maximal at $6.392453\ldots$ No choice of winding — $k_f$, $k_{vf}$, $K_S$, or any iterate — and no admissible $\Theta_R$ convention can exceed it: the value set is finite and bounded. $(4)$ \textbf{Top-weight near-universality.} In variant (A) the heaviest weight of \textit{four of the five} solids is exactly $2\hat\mu'^2$ ($k_f\delta_v\equiv\pi\bmod 2\pi$: tetra $3\pi$, cube $\pi$, dodeca $\pi$, icosa $\pi$), the octahedron's being $4\hat\mu'^2$ ($k_f\delta_v=2\pi$). The inter-solid hierarchy therefore lives \textbf{entirely at the light end}; the heavy ends are degenerate across solids. $(5)$ \textbf{Two structural obstructions.} (a) \textbf{Tetrahedral degeneracy}: $\delta_v=\pi\Rightarrow D(m)=2-2\cos(m\pi/2)\in\{0,2,4\}$, and the $k=3$ multiplet is \textit{fully degenerate} at $2\hat\mu'^2$ (both variants) — the tetrahedron carries no multiplet splitting at all. (b) \textbf{Even-$k$ zero modes} [FT.60(5)'s standing zero]: the cube in (A) ($k_f=2$) and the octahedron in (B) ($k_{vf}=2$) each carry an exactly massless weight $m=0$, so their literal "lightest pole mass" is $0$ and only lightest-\textit{nonzero} reads are well defined. $(6)$ \textbf{Inter-solid ratios (lightest nonzero, relative to the tetrahedron).}

%dsp-level:0
\[
\begin{array}{c|cc|cc}
& \multicolumn2{c|}{\text{(A)}} & \multicolumn2{c}{\text{(B)}}\\
S & m^2\ \text{ratio} & m\ \text{ratio} & m^2\ \text{ratio} & m\ \text{ratio}\\
\hline
\text{tetra} & 1 & 1 & 1 & 1\\
\text{cube} & 1\ (0\ \text{with the zero mode}) & 1\ (0) & 0.292893218813 & 0.541196100146\\
\text{octa} & 0.5 & 0.707106781187 & 1.5 & 1.224744871392\\
\text{dodeca} & 0.048943483705 & 0.221231742082 & 0.048943483705 & 0.221231742082\\
\text{icosa} & 0.133974596216 & 0.366025403784 & 0.133974596216 & 0.366025403784
\end{array}
\]

Normalized to the lightest solid, in mass: (A) dodeca $1$, icosa $1.654488638651$, octa $3.196226610750$, cube $=$ tetra $4.520147021340=\sqrt{2/(2-2c)}$; (B) dodeca $1$, icosa $1.654488638651$, cube $2.446285940037$, tetra $4.520147021340$, octa $5.536026882322$. $(7)$ \textbf{Intra-solid splittings (heaviest/lightest nonzero).}

%dsp-level:0
\[
\begin{array}{c|cc|cc}
S & \text{(A)}\ m^2 & \text{(A)}\ m & \text{(B)}\ m^2 & \text{(B)}\ m\\
\hline
\text{tetra} & 1 & 1 & 1 & 1\\
\text{cube} & 1\ (+1\ \text{zero mode}) & 1 & 3{+}2\sqrt2=5.828427124746 & 1{+}\sqrt2=2.414213562373\\
\text{octa} & 4 & 2 & 1\ (+1\ \text{zero mode}) & 1\\
\text{dodeca} & \tfrac{2}{2-2c}=20.431729094531 & 4.520147021340 & \tfrac{2-2\cos(3\pi/10)}{2-2c}=8.422260053931 & 2.902113032590\\
\text{icosa} & 4{+}2\sqrt3=7.464101615138 & 1{+}\sqrt3=2.732050807569 & 7{+}4\sqrt3=13.928203230276 & 2{+}\sqrt3=3.732050807569
\end{array}
\]

Minimal polynomials of the $m^2$ splittings: dodeca (A) $x^4-24x^3+76x^2-64x+16$, dodeca (B) $x^4-14x^3+51x^2-34x+1$, icosa (A) $x^2-8x+4$, icosa (B) $x^2-14x+1$, cube (B) $x^2-6x+1$; the mass splittings $1{+}\sqrt3$, $2{+}\sqrt3$, $1{+}\sqrt2$ have $x^2-2x-2$, $x^2-4x+1$, $x^2-2x-1$ (machine; certified dyadic isolations). $(8)$ \textbf{Dual splitting, sharpened.} The dual pairs remain degenerate in $m_{\mathrm{hid}}$ and $\Pi$ [FT.44(1)] and split on the pole side as multisets in \textit{both} variants — cube $\{0,2\}$ vs octa $\{1,4\}$ in (A), cube $\{2{-}\sqrt2,2{+}\sqrt2\}$ vs octa $\{0,3\}$ in (B); dodeca vs icosa share only the top value $2\hat\mu'^2$ in (A) and nothing in (B). FT.44(2)'s splitting template acquires its multiplet form. $(9)$ \textbf{Field-degree ledger.} Tetra/octa: $\mathbb Q$; cube: $\mathbb Q(\sqrt2)$; icosa: $\mathbb Q(\sqrt3)$; dodeca: $\mathbb Q(\cos\tfrac\pi{10})$, degree $4$, with $\sqrt5=8c^2-5$ its unique quadratic subfield (machine) — the requested four-field list $\{\mathbb Q,\mathbb Q(\sqrt2),\mathbb Q(\sqrt3),\mathbb Q(\sqrt5)\}$ is \textbf{insufficient for the dodecahedron}, exactly as FT.73(2)'s degree-$16$ compositum requires. $(10)$ \textbf{FT.63(6) precision correction.} The dodecahedral quartic $x^4-8x^3+19x^2-12x+1$ has \textbf{four positive roots} (machine-isolated to width $2^{-60}$): $0.097886967410$, $0.824429495415$, $3.175570504585$, $3.902113032590$, realized by the weights $m\in\{1,19\},\{3,17\},\{7,13\},\{9,11\}$ respectively. FT.63(6)'s "$D(1),D(3)$ the two positive roots of the quartic" corrects to \textbf{the two smallest of its four positive roots}; patched in place (marker).

\textbf{Req.} Def-H.37(2), Def-H.38; Def-H.21, Def-H.23; FT.22, FT.32, FT.43, FT.44, FT.45, FT.47, FT.52, FT.54, FT.55, FT.60, FT.63; D2.9; CT2.5 (Part 2); CT0.7a (Part 1); Chebyshev recursion, de Moivre $\cos5t$, exact root isolation by bisection, minimal polynomials via the multiplication map [math; machine].

\textbf{Pf.}

%dsp-level:0
\[
(1):\ k_f=K_S/F:\ (12,12,24,60,60)/(4,6,8,12,20)=(3,2,3,5,3)\ \checkmark;
\]

%dsp-level:0
\[
\qquad k=L\kappa_{\min}/2:\ \text{odd}\ L\to L,\ \text{even}\ L\to L/2
\]

%dsp-level:0
\[
\Theta_R=4\pi:\ D(m)=2-2\cos(2m\pi)=0\ \forall m;\qquad\Theta_R=0:\ D\equiv0\ [\mathrm{FT.63}(4),(6)]
\]

%dsp-level:0
\[
(2):\ D(m)
\]

%dsp-level:0
\[
= 2-2T_m(\cos(\delta_v/2))\ [\text{Chebyshev recursion}\ T_{n+1}=2xT_n-T_{n-1},\ \text{integer coefficients}];
\]

%dsp-level:0
\[
\quad\cos(\delta_v/2)
\]

%dsp-level:0
\[
= (0,\tfrac{\sqrt2}2,\tfrac12,c,\tfrac{\sqrt3}2)
\]

%dsp-level:0
\[
\text{base elements exact}:\ 16c^4-20c^2+5=0\ [\cos5t\ \text{at}\ t=\tfrac\pi{10}];
\]

%dsp-level:0
\[
\quad (\tfrac{\sqrt2}2)^2=\tfrac12;
\]

%dsp-level:0
\[
\quad (\tfrac{\sqrt3}2)^2=\tfrac34;\quad\text{root branch by sign/order}
\]

%dsp-level:0
\[
\text{values}:\ \text{exact field arithmetic};
\]

%dsp-level:0
\[
\qquad \text{cross-check}:\ |D(m)-[2-2\sum_j(-1)^j(m\delta_v/2)^{2j}/(2j)!]|<10^{-45}\ \text{at 120 terms}\ [\text{machine}]
\]

%dsp-level:0
\[
(3):\ D(m+P)=D(m)\iff\cos((m+P)\delta_v/2)=\cos(m\delta_v/2):\ P\delta_v/2=2\pi:\ P=4\pi/\delta_v;
\]

%dsp-level:0
\[
\qquad D(P-m)=D(m),\ D(-m)=D(m)
\]

%dsp-level:0
\[
m \in [0,P/2]:\ m\delta_v/2
\]

%dsp-level:0
\[
\in [0,\pi]:\ \cos\ \text{strictly decreasing}
\]

%dsp-level:0
\[
\implies D\ \text{strictly increasing}:\ \min_{\ne0}
\]

%dsp-level:0
\[
= D(1)
\]

%dsp-level:0
\[
= D_S,
\]

%dsp-level:0
\[
\ \max
\]

%dsp-level:0
\[
= D(P/2)\qquad = 4
\]

%dsp-level:0
\[
\text{distinct count}=P/2+1;\qquad 4/D_{\mathrm{dodeca}}:\ \text{minpoly by the multiplication map on}\ \mathbb Q(c)\ [\text{machine}]
\]

%dsp-level:0
\[
(4):\ k_f\delta_v=(3\pi,\pi,2\pi,\pi,\pi):\ \cos(k_f\delta_v/2)=(0,0,-1,0,0)\implies D(k_f)=(2,2,4,2,2)
\]

%dsp-level:0
\[
(5a):\ \delta_v
\]

%dsp-level:0
\[
= \pi:\ D(m)
\]

%dsp-level:0
\[
= 2-2\cos(m\pi/2):\ m\ \text{odd}\to2,
\]

%dsp-level:0
\[
\ m
\]

%dsp-level:0
\[
\equiv 0(4)\to0,
\]

%dsp-level:0
\[
\ m
\]

%dsp-level:0
\[
\equiv 2(4)\to4:\ k_f
\]

%dsp-level:0
\[
= 3\ \text{odd weights only}\implies\text{constant}\ 2
\]

%dsp-level:0
\[
(5b):\ k\ \text{even}\implies m=0\in\{k,k{-}2,\dots,-k\}:\ D(0)=0\ [\mathrm{FT.60}(5)]
\]

%dsp-level:0
\[
(6),
\]

%dsp-level:0
\[
(7):\ \text{exact division in the solid's field};
\]

%dsp-level:0
\[
\quad\text{mass ratios}:\ r
\]

%dsp-level:0
\[
= \sqrt{q}:\ q\ \text{isolated exactly}
\]

%dsp-level:0
\[
\implies r
\]

%dsp-level:0
\[
\in (\sqrt a,\sqrt b)\ [\sqrt{}\ \text{monotone}];
\]

%dsp-level:0
\[
\text{icosa (A)}:\ q=\tfrac{2}{2-\sqrt3}=2(2{+}\sqrt3)=4{+}2\sqrt3:\ \sqrt q=1{+}\sqrt3\ [(1{+}\sqrt3)^2=4{+}2\sqrt3\ \checkmark]
\]

%dsp-level:0
\[
\text{icosa (B)}:\ q=\tfrac{2+\sqrt3}{2-\sqrt3}=(2{+}\sqrt3)^2=7{+}4\sqrt3:\ \sqrt q=2{+}\sqrt3\ \checkmark
\]

%dsp-level:0
\[
\text{cube (B)}:\ \tfrac{2+\sqrt2}{2-\sqrt2}=3{+}2\sqrt2=(1{+}\sqrt2)^2\ \checkmark;
\]

%dsp-level:0
\[
\qquad (8),(9):\ \text{multiset comparison; field membership by construction}
\]

%dsp-level:0
\[
(10):\ \text{sign-change scan on}\ \mathbb Q\ \text{at step}\ 1/500\ \text{over}\ [-1,6]\ \text{then bisection to}\ 2^{-60}:\ 4\ \text{brackets, all in}\ (0,\infty);
\]

%dsp-level:0
\[
\quad q(0..4)=(1,1,5,1,1)>0
\]

%dsp-level:0
\[
\text{weight attribution}:\ D(m)=\text{root}\iff m\ \text{in the stated residue classes mod}\ 20\ [\text{machine}].\ \blacksquare
\]

\textbf{Phys.} Ask the five solids for their pole masses under the declared defect transfer, and the answer is a small, closed, exactly computable set of numbers — three values for the tetrahedron, five for the cube, eleven for the dodecahedron — all of them between zero and four times one unknown scale. Nothing in that set is large. The tetrahedron's multiplet does not split at all; the cube's even winding leaves one weight exactly massless; and the heaviest state of four of the five solids is the same number. So the Platonic defects do not generate hierarchy — they generate a bounded, near-degenerate spectrum, and the bound is a theorem about cones, not a limitation of the search.

\vspace{2pt}\noindent\rule{\textwidth}{0.4pt}\vspace{2pt}

\end{theorem}

\begin{theorem}[{\texorpdfstring{FT.78 — The orbit decompositions: tetrahedron exhaustive, cube symmetry-reduced}{FT.78 - The orbit decompositions: tetrahedron exhaustive, cube symmetry-reduced}}]\label{tag:FT-78}
\textbf{Thm.} Under Def-H.28, at (S1)–(S4): $(1)$ \textbf{Tetrahedron, exhaustive.} The $1230$ admissible configurations [FT.68(1)] form \textbf{$70$ orbits} under $\operatorname{Aut}=S_4$ ($|\operatorname{Aut}|=24$), with $\sum_{\mathcal O}\Omega_{\mathcal O}=1230$; orbits by $K$: $(1,1,5,13,23,21,6)$ for $K=0..6$; configuration counts by $K$ reproduce the cap row $(1,12,66,220,435,408,88)$ [FT.68(1)] exactly; $\Omega\in\{1,3,4,6,8,12,24\}$, every value dividing $24$. $(2)$ \textbf{Tetrahedron cover sector: seven orbits total.} The $85$ bypass covers form $5$ orbits and the $16$ locked covers $2$ orbits:

%dsp-level:0
\[
\begin{array}{c|cccccc}
\text{sector} & K & |\mathcal C| & \Omega & 4^{|\mathcal C|} & \text{signature (cycles shadowed)} & \text{type multiset}\\
\hline
\text{byp} & 0 & 0 & 1 & 1 & () & (0,0)^4\ \text{(vacuum)}\\
\text{byp} & 2 & 1 & 24 & 4 & () & (0,0),(0,1),(1,0),(1,1)\ \text{single thread}\\
\text{byp} & 4 & 2 & 24 & 16 & (2) & (1,0),(1,1)^2,(1,2)\\
\text{byp} & 4 & 2 & 24 & 16 & (2) & (0,1),(1,1)^2,(2,1)\\
\text{byp} & 4 & 2 & 12 & 16 & () & (0,2),(1,1)^2,(2,0)\ \text{double thread}\\
\hline
\text{lock} & 4 & 2 & 12 & 16 & (2,2) & \text{double-digon }(2,2)\text{ hub}\\
\text{lock} & 6 & 3 & 4 & 64 & (2,2,2) & K_{1,3}\ \text{star hub}
\end{array}
\]

—the two locked rows being exactly FT.69(5)'s minimal dissolution family ($12$) and the four $K_{1,3}$ triple-digon stars [FT.68(1)]. The single-centre orbit's size $24=Vd(d{-}1)$ is FT.69(2)'s closed form. $(3)$ \textbf{Orbit-level amplitude identity (new).} The amplitude sum decomposes exactly over classes:

%dsp-level:0
\[
\sum_{\mathcal O}\Omega_{\mathcal O}\,4^{|\mathcal C(\mathcal O)|}=Z^{\pm},
\]

machine-verified on both skeletons: tetrahedron $1057$ (bypass) and $448$ (locked) [FT.68(1)]; cube $1\,109\,313$ and $1\,150\,848$ [FT.68(3)]. Consequently the covering-polynomial coefficients are the $|\mathcal C|$-level \textit{configuration} counts, and the orbit counts refine them: cube bypass $|\mathcal C|=0..4\Rightarrow(1,1,19,57,88)$ orbits over $(1,48,648,2672,3624)$ configurations; locked $|\mathcal C|=2..6\Rightarrow(1,18,71,7,1)$ over $(24,728,3192,264,4)$. $(4)$ \textbf{Cube, symmetry-reduced.} Cap row $(1,24,276,2024,10506,40296,115796,247704,383913,408464,270096,91392,9600)$, total $1\,580\,092$ configurations with $K\le12$ forced [FT.68(3)]; cap-aware Burnside over $|\operatorname{Aut}|=48$ gives orbits by $K$ $(1,1,11,49,250,863,2509,5211,8163,8571,5768,1922,245)$, total $33\,564>10^4$ — the bounded rule is observed, so full configuration enumeration is \textit{not} attempted. The \textbf{cover sector is enumerated explicitly} (it is small): the $11\,205$ covers ($6993$ bypass $+\,4212$ locked, two independent enumerators agreeing [FT.69(2)]) form \textbf{$264$ orbits} ($166$ bypass, $98$ locked), $\sum\Omega=11\,205$, $\Omega\mid48$ throughout, with $\Omega$-histogram $\{1{:}1,\ 4{:}1,\ 8{:}1,\ 12{:}4,\ 16{:}2,\ 24{:}47,\ 48{:}208\}$ — the $208$ free orbits are the generic case, so the twenty most populated classes are all free orbits ($\Omega=48$) and are distinguished only by signature: one at $|\mathcal C|=1$ (the single thread), eight at $|\mathcal C|=2$ (four cycle-free, four shadowing one digon), eleven at $|\mathcal C|=3$. \textbf{Independent cross-validation:} explicit canonical-form enumeration of every admissible configuration with $K\le5$ ($53\,127$ of them) reproduces the Burnside row $(1,1,11,49,250,863)$ exactly — two methods, same numbers. $(5)$ \textbf{Signatures.} Covers shadow only \textbf{even} cycles, orbit by orbit [FT.69(3)(i) recovered at class level]: tetrahedron covers contain digons only ($84$ instances, $\delta h=2$ [D0.6]); cube covers contain digons and directed $4$-cycles ($14\,832$ instances, $\delta h\in\{2,4\}$), with $10$ distinct cycle contents $\{(),(2),(2,2),(4),(2,2,2),(2,4),(2,2,2,2),(2,2,4),(2,2,2,2,2),(2,2,2,2,2,2)\}$ and no odd cycle anywhere. The signature $(K,\text{cycles},|\mathcal C|)$ is \textbf{not a complete class invariant}: $7$ tetrahedron cover orbits carry only $6$ signatures, $264$ cube cover orbits only $21$ — the completing datum is the vertex-type multiset (clause (2)'s last column). $(6)$ \textbf{Charge-conjugation chirality of the cover sector (new).} Global orientation reversal $R\mapsto R^{\mathrm T}=\{(b,a):(a,b)\in R\}$ is an involution of the pattern space that is \textbf{not} in $\operatorname{Aut}(\Gamma)$ (it is induced by no vertex permutation). It pairs cover classes: the $48$ bypass dissolved-digon covers of FT.69(4) split into \textbf{two $\Omega=24$ orbits exchanged by reversal} (representatives with type multisets $(1,0),(1,1)^2,(1,2)$ and $(0,1),(1,1)^2,(2,1)$ — transposes of each other; machine: the reversal of one representative is canonical in the other orbit). This is the pattern-level form of charge conjugation [D2.6]. \textbf{Both mass ansätze are reversal-invariant} ($K$, $\Omega$ and $4^{|\mathcal C|}$ are), so the chirality is invisible to the mass spectrum; reading it requires a C-odd ledger of FT.40's $q_{\mathrm{core}}$ species.

\textbf{Req.} Def-H.25, Def-H.26, Def-H.28; FT.67, FT.68, FT.69, FT.70, FT.71; FT.27, FT.40; A4, A8; D0.6, D0.8, D2.6; CT0.11, CT0.13, CT0.14 (Part 1); CT1.20 (Part 2); orbit–stabilizer, Burnside's lemma, canonical-form enumeration [math; machine].

\textbf{Pf.}

%dsp-level:0
\[
(1):\ \text{canonical form}
\]

%dsp-level:0
\[
= \min\ \text{over the}\ 24\ \text{edge-permutations};
\]

%dsp-level:0
\[
\quad\text{admissibility and}\ K\ \text{are invariant}\implies\text{orbits lie inside the set}
\]

%dsp-level:0
\[
\sum_{\mathcal O}\Omega_{\mathcal O}=1230\ [\text{machine}];\qquad\Omega\cdot|\operatorname{Stab}|=24\implies\Omega\mid24;\qquad\text{cap row}:\ \text{CT0.11 counting per}\ K
\]

%dsp-level:0
\[
(2):\ \text{explicit enumeration of the}\ 101\ \text{covers}\ (85+16);
\]

%dsp-level:0
\[
\quad |\mathcal C|=1:\ \Omega=24=Vd(d{-}1)=4\cdot3\cdot2\ [\mathrm{FT.69}(2)]
\]

%dsp-level:0
\[
\text{locked rows}:\ \text{double-digon}(2,2)\ \text{hub}\ (12=4\binom32),\ K_{1,3}\ \text{star}\ (4,\ \operatorname{Stab}=S_3)\ [\mathrm{FT.69}(5)]
\]

%dsp-level:0
\[
(3):\ Z^{\pm}
\]

%dsp-level:0
\[
= \sum_{R\ \mathrm{cover}}4^{|\mathcal C(R)|}\ [\mathrm{FT.67}(2)\ \text{at (S1)}]
\]

%dsp-level:0
\[
= \sum_{\mathcal O}\ \sum_{R\in\mathcal O}4^{|\mathcal C|}
\]

%dsp-level:0
\[
= \sum_{\mathcal O}\Omega_{\mathcal O}4^{|\mathcal C(\mathcal O)|}
\]

%dsp-level:0
\[
\quad [|\mathcal C|\ \text{class-constant}]
\]

%dsp-level:0
\[
(4):\ \text{Burnside}:\ \#\text{orbits}_K
\]

%dsp-level:0
\[
= \tfrac1{48}\sum_{p}\operatorname{fix}_K(p),
\]

%dsp-level:0
\[
\ \operatorname{fix}_K\ \text{by orbit-subset enumeration under the cap};
\]

%dsp-level:0
\[
\quad\text{covering polys}:\ \text{DP}
\]

%dsp-level:0
\[
= \text{backtracking}\ [\mathrm{FT.69}(2)]
\]

%dsp-level:0
\[
\text{cross-validation}:\ K
\]

%dsp-level:0
\[
\le 5\ \text{explicit canonical forms}
\]

%dsp-level:0
\[
= (1,1,11,49,250,863)
\]

%dsp-level:0
\[
= \text{Burnside row}\ [\text{machine}]
\]

%dsp-level:0
\[
(5):\ \text{canonical cycle enumeration per representative};\ \text{evenness}:\ \mathrm{FT.69}(3)(i);
\]

%dsp-level:0
\[
\qquad \text{incompleteness}:\ \text{two classes share}(4,(2),2)\ \text{and differ in type multiset}
\]

%dsp-level:0
\[
(6):\ R^{\mathrm T}\ \text{preserves}\ K,\ \text{types}^{\mathrm T},\ |\mathcal C|,\ \text{cover status};
\]

%dsp-level:0
\[
\quad R^{\mathrm T}=p(R)\ \forall p\in\operatorname{Aut}\ \text{would force}\ \Omega\ \text{to merge}:\ \text{the two orbits are distinct}\implies\notin\operatorname{Aut}
\]

%dsp-level:0
\[
\text{reversal maps one}\ \Omega{=}24\ \text{orbit onto the other}\ [\text{machine: canonical forms agree}].\ \blacksquare
\]

\textbf{Phys.} Ask the flickering gas how many distinct shapes it takes, and the answer is startlingly small: the tetrahedron's $1230$ configurations are $70$ shapes, and the whole of its amplitude-carrying cover content is \textit{seven} — one vacuum, one single thread, three two-centre networks, and the two locked hub families. The cube's eleven thousand covers are $264$ shapes over seven mass levels. Nothing here is random; the shapes come in chirality pairs related by reversing every arrow, a distinction the mass reads cannot see. The degeneracies are where the complexity lives, and degeneracy is not hierarchy.

\vspace{2pt}\noindent\rule{\textwidth}{0.4pt}\vspace{2pt}

\end{theorem}

\begin{theorem}[{\texorpdfstring{FT.81 — Particle-content census (audit pass A)}{FT.81 - Particle-content census (audit pass A)}}]\label{tag:FT-81}
\textbf{Thm.} Under Def-H.29, with the SM field content external and marked: $(1)$ \textbf{The census.} $24$ entries, classified with their supplying tags or no-goes:

%dsp-level:0
\[
\begin{array}{lll}
\textbf{PRESENT}\ (10) & \textbf{PARTIAL}\ (6) & \textbf{ABSENT}\ (8)\\
\hline
\text{gluons }(8) & \text{Higgs doublet} & \text{neutrino}\\
W^\pm & \text{quark triplet} & \text{generations/flavour}\\
Z\ (\text{mixture}) & \text{lepton doublet} & \text{CKM/PMNS, CP phase}\\
\text{photon} & \text{charged-lepton singlet} & \text{Yukawa sector}\\
\text{confinement/singlet} & \text{observable charge} & \text{chirality/Weyl}\\
\theta_W,\ e\ (\text{relations}) & \text{EW VEV} & \text{gauge anomalies}\\
G_N\ (\text{beyond SM}) & & \Lambda_{\mathrm{QCD}}\\
\text{spin ladder} & & \text{spin-statistics}\\
\text{fermion/boson content} & & \\
D^\pm\ \text{dark (beyond SM)} & &
\end{array}
\]

\textbf{Completeness fraction $=\dfrac{10+\frac12\cdot6}{24}=\dfrac{13}{24}=0.5416\overline6$.} Sub-censuses: \textbf{gauge species $4/4$ PRESENT} (as aspect blocks and reads, \textit{not} as gauged chiral couplings); \textbf{fermion sector $8$ entries $\to0$ present, $3$ partial, $5$ absent $=0.1875$}. $(2)$ \textbf{The SM side, machine-verified in exact rationals} (external data): one generation $=15$ left-handed Weyl fields $=6+3+3+2+1$ ($(3,2)_{1/6}$, $(\bar3,1)_{-2/3}$, $(\bar3,1)_{1/3}$, $(1,2)_{-1/2}$, $(1,1)_{1}$), $45$ over three generations; $12$ gauge bosons ($8+3+1$) and $4$ real Higgs dof; every anomaly sum vanishes exactly — $\sum Y=0$, $\sum Y^3=0$, $[\mathrm{SU}(3)]^2\mathrm U(1):\ \sum Y\,d(\mathrm{SU}(2))\,T(3)=0$, $[\mathrm{SU}(2)]^2\mathrm U(1):\ \sum Y\,d(\mathrm{SU}(3))\,T(2)=0$, $[\mathrm{SU}(3)]^3:\ 2A(3)+2A(\bar3)=0$, $[\mathrm{SU}(2)]^3\equiv0$. $(3)$ \textbf{The charge lattice, exactly.} The SM field charges $Q=T_3+Y$ span $\{0,\pm\tfrac13,\pm\tfrac23,\pm1\}$; the corpus read is ternary, $q_{\mathrm{EM}}\in\{-1,0,+1\}$ with all three attained [D3.4, EM.9, P2]. The two sets are related by \textbf{confinement projection, machine-verified}: the CT3.17/D3.9 singlets of $u,d$ — $\varepsilon_{ijk}$ baryons and $\delta^i{}_j$ mesons — carry \textit{integer} charge, and the meson charge set is \textbf{exactly} $\{-1,0,+1\}$, the corpus read; the baryon set is $\{-1,0,1,2\}$, so the $|q|=2$ state exceeds a single chain's ternary read, and the confined thirds are unreadable by construction [CT3.22: colour is $\neg$observable as end-charge]. Verdict: \textbf{charge quantization to integers is reproduced for the observable sector; the per-chain read is $3$-valued where the SM's observable lattice is $\mathbb Z$} (a multi-chain additive read would be needed for $|q|\ge2$; FT.41(a) counts $N_{\mathrm{charge}}=2$ end registers). $(4)$ \textbf{Anomalies are not poseable.} Every anomaly above is a constraint on \textit{chiral} hypercharge assignments. The corpus has no chiral projector, no Weyl symbol and no hypercharge-assignment read ($Y$ appears only as the abstract generator with three realizations [EW.8] and in the mixing [D3.10/EW.24]), so the cancellation question cannot be posed — a structural absence downstream of chirality, not a numerical failure. $(5)$ \textbf{Inventory facts (machine-checked over the frozen Parts 1–2).} Count $0$ for: Yukawa, neutrino, Weyl, CKM, PMNS, left-handed, right-handed, $\gamma_5$, flavour, flavor, black hole, spin-statistics, $P_L$, $P_R$, $\bar\psi$, generation structure, quark mixing, lepton number. Present: $\theta_W$ ($71$), $T_a$ ($51$), $h_0$ ($38$), gluon ($22$), $f_{abc}$ ($14$), $\varepsilon_{ijk}$ ($12$), $m_H$ ($11$), baryon ($11$), $W^\pm$ ($10$), meson ($6$), hypercharge ($4$), chirality ($2$). The two "chirality" occurrences are both inside D2.6's orientation-reversal clause ($\omega$: chirality $\mapsto-$chirality) — vocabulary without a projector. The four "generation" occurrences in D/F are the algebraic sense ("mass generation", "$\delta$/$\varepsilon$ generation"), machine-checked to be non-flavour.

\textbf{Req.} Def-H.29; D1.5, D2.6, D3.4, D3.7, D3.8, D3.9, D3.10, D4.1–D4.3; CT0.9, CT0.12, CT0.13, CT0.14 (Part 1); CT2.9, CT1.25 (Part 2); CT3.15, CT3.17, CT3.18, CT3.20, CT3.21, CT3.22 (Part 2); CT4.1, CT4.4, CT4.5 (Part 2); EM.9, EM.20 (Part 2); EW.8, EW.13, EW.20, EW.24, EW.27, EW.32, EW.33, EW.35 (Part 2); FT.19, FT.20, FT.26, FT.27, FT.30, FT.34, FT.41, FT.52, FT.55, FT.57, FT.60, FT.71, FT.76, FT.79, FT.80; P1–P5; A §0.2, §0.3; external SM field content (marked); exact rational arithmetic, integer inventory scan [machine].

\textbf{Pf.}

%dsp-level:0
\[
(1):\ \text{per-entry citation table}\ [\text{machine-checked: every PRESENT row resolves to the cited tag; every ABSENT row to a scan count or a no-go}]
\]

%dsp-level:0
\[
\text{fraction}:\ (10+3)/24=13/24;\qquad\text{gauge sub-census}\ 4/4;\qquad\text{fermion sub-census}\ (0+\tfrac32)/8=3/16
\]

%dsp-level:0
\[
(2):\ 6+3+3+2+1=15;\quad15\cdot3=45;\quad8+3+1=12;\quad2\cdot2=4;\qquad\text{anomaly sums}:\ \text{exact Fractions}
\]

%dsp-level:0
\[
\textstyle\sum Y=6\cdot\tfrac16+3(-\tfrac23)+3\cdot\tfrac13+2(-\tfrac12)+1=1-2+1-1+1=0;
\]

%dsp-level:0
\[
\qquad \sum Y^3=\tfrac1{36}-\tfrac89+\tfrac19-\tfrac14+1=0
\]

%dsp-level:0
\[
[\mathrm{SU}(3)]^2\mathrm U(1):\ \tfrac12\bigl(2\cdot\tfrac16-\tfrac23+\tfrac13\bigr)=0;
\]

%dsp-level:0
\[
\qquad [\mathrm{SU}(2)]^2\mathrm U(1):\ \tfrac12\bigl(3\cdot\tfrac16-\tfrac12\bigr)=0;
\]

%dsp-level:0
\[
\qquad [\mathrm{SU}(3)]^3:\ 2\cdot\tfrac12-2\cdot\tfrac12=0
\]

%dsp-level:0
\[
(3):\ Q=T_3+Y\ \text{over the five multiplets}\implies\{0,\pm\tfrac13,\pm\tfrac23,\pm1\};
\]

%dsp-level:0
\[
\qquad q_{\mathrm{EM}}=\tfrac12(\sigma_s-\sigma_t)\in\{-1,0,1\}\ [\mathrm{EM.9}]
\]

%dsp-level:0
\[
\text{mesons}:\ \{u\bar u,u\bar d,d\bar u,d\bar d\}\to\{0,1,-1,0\}=\{-1,0,1\}\ \checkmark;
\]

%dsp-level:0
\[
\qquad \text{baryons}:\ \{uuu,uud,udd,ddd\}\to\{2,1,0,-1\}
\]

%dsp-level:0
\[
\text{integrality}:\ \tfrac23-\tfrac13
\]

%dsp-level:0
\[
= \tfrac13\notin\mathbb Z\ \text{per quark},
\]

%dsp-level:0
\[
\ \text{but}\ \varepsilon/\delta\ \text{singlets sum to}\ \mathbb Z\ [\mathrm{CT3.17},\mathrm{D3.9};\ \text{exhaustive}]
\]

%dsp-level:0
\[
(4):\ \text{anomaly}\equiv\text{chiral constraint};
\]

%dsp-level:0
\[
\quad\neg\exists\ \text{projector/Weyl symbol in Parts 1–2}\ [\text{scan}]\implies\neg\exists\ \text{the question}
\]

%dsp-level:0
\[
(5):\ \text{exact substring counts over the seven frozen documents}\ [\text{machine}].\ \blacksquare
\]

\textbf{Phys.} Line up the Standard Model's inventory against what this corpus can read, and the pattern is not random. Everything about \textit{forces} is there: eight gluons with a closing eight-dimensional algebra, three weak directions that mix into a massive $Z$ and an exactly massless photon, a scalar with a background and a mass formula, a confinement rule that makes colour unreadable and thereby makes observable charge integral — the meson charges come out as precisely the three values the ledger can register. Everything about \textit{matter} is missing: no handedness, no Weyl spinor, no neutrino, no flavours, no mixing, no Yukawa couplings, and no second scale for any of them to live on. The corpus is a theory of the gauge skeleton with the chiral flesh left undeclared — and the anomalies that make the SM's flesh consistent cannot even be written down here.

\vspace{2pt}\noindent\rule{\textwidth}{0.4pt}\vspace{2pt}

\end{theorem}

\begin{theorem}[{\texorpdfstring{FT.83 — The rung emergence ledger (audit pass C): derived, declared, open}{FT.83 - The rung emergence ledger (audit pass C): derived, declared, open}}]\label{tag:FT-83}
\textbf{Thm.} Under Def-H.29, per rung $n$: $(1)$ \textbf{The ladder is a theorem, not a picture.} $\mathcal A(0)=\{J\}$; $\mathcal A(1)=\{J,G,W,H,D^+\}$; $\mathcal A(2)=\mathcal A(1)\cup\{\mathrm{EM}\}$; $\mathcal A(3)=\mathcal A(2)\cup\{\mathrm{strong}\}$; $\mathcal A(n\ge4)=\mathcal A(3)\cup\{D^-\}$, with a strict mask flag $\dim\mathcal K_n=(6,16,17,25)$ [CT0.13, CT0.14 — both with Pfs]. The admissibility ceiling is derived: $D_{\mathrm{tick}}(n)=2^{n+1}\le C_{\mathrm{cap}}=16\iff n\le3$, with saturation at $n=3$ [CT0.5, CT0.11, P2's $n_{\max}=3$]. $(2)$ \textbf{Rung-by-rung status.}

%dsp-level:0
\[
\begin{array}{c|l|l}
n & \text{what becomes active/trackable} & \text{status}\\
\hline
0 & J^\pm\ \text{seals only};\ \dim\mathcal K_0=6 & \textbf{derived}\ [\mathrm{CT0.13}/\mathrm{CT0.7}]\\
1 & G(5)\ \text{spin-2 shear},\ W(3),\ H(1),\ D^+(1);\ \dim\mathcal K_1=16 & \textbf{derived}\ [\mathrm{CT0.12}/\mathrm{CT0.13}];\ \text{the "pretime force set"}\\
2 & +\ \mathrm{EM};\ \alpha_{\mathrm{EM}}=\pi;\ q_{\mathrm{EM}}\ \text{readable on open chains} & \textbf{derived}\ [\mathrm{CT0.13},\mathrm{EM.9},\mathrm{EM.20}]\ \text{under Comm-min}+\text{Eff (declared premises)}\\
3 & +\ \text{strong};\ N_{\text{colour}}=3,\ N_{\mathrm{gluon}}=8;\ \text{sealing}=\text{confinement};\ \text{mass reads} & \textbf{derived}\ [\mathrm{CT3.20},\mathrm{CT3.13},\mathrm{CT3.22},\mathrm{EW.20}/\mathrm{EW.33}]\\
\ge4 & D^-\ \text{absorption}\ O(n)=2^{n+1}{-}16,\ w=0\ \text{dust},\ \text{locking};\ \neg\exists\ \text{gauge structure} & \textbf{derived as a no-go}\ [\mathrm{CT0.11},\mathrm{CT4.1},\mathrm{CT4.4},\mathrm{CT4.5},\mathrm{FT.59}]
\end{array}
\]

$(3)$ **What is \textit{not} derived at any rung.** (a) The proposal's "$n=1$: gravity, weak, Higgs, Dark$^+$, Janus" is exactly $\mathcal A(1)$ and is derived — but the \textit{Higgs} entry is the scalar aspect $F_H$ with a P4-free background $h_0$ [EW.13, EW.20, P4], and its mass is "configuration data" [EW.27]: the mechanism is derived, the scale is not. (b) The proposal's "$n=2$: EM $+$ lepton structure" is \textbf{half derived}: EM is a theorem; "lepton structure" has no read at any rung (no lepton species, no chirality [FT.81(5)]). (c) "$n=3$: colour $+$ quarks $+$ the baryon edge budget" is derived — $\varepsilon_{ijk}$ sealing of three fundamentals [CT3.15/CT3.17/D3.9], the $2$s-and-$3$s partition rule [FT.34(5)], and $N=1$ impossibility [FT.27] — but "quark" in the corpus is D3.8's \textit{bound fermionic loop with a triplet attachment}, i.e. a topological definition with no flavour and no chirality. (d) "$n\ge4$: lock / black-hole-matter cascade": the locking and absorption side is derived (CT4.1/CT4.2/CT4.5, $w=0$ dust, retained-and-inaccessible [D1.11]); \textbf{"black hole" has no corpus referent at all} (inventory count $0$), so the cascade reading is an external gloss, recorded and not adopted. $(4)$ \textbf{Parameter census.} P4 declares $8$ independent internal configuration parameters $(g,g',h_0,c_W,\Lambda,\varkappa,\beta,\alpha_\Omega)$ with $6$ functionally dependent members $(\theta_W,e,\mu_\star^2,m_W,m_Z,m_H)$; P1/A §0.3 admit exactly \textbf{one} free dimensionful scale $\mathfrak h$; P3 names $G_N=\mathfrak h^2/(4\ln2)$, $\kappa^2=8\pi G_N$ for external comparison only; P5 forbids identifying any quantity with a measured constant. The SM's $19$ parameters (external count: $3$ gauge $+\,1$ VEV $+\,9$ Yukawa $+\,4$ CKM $+\,1$ $\theta_{\mathrm{QCD}}+\,1$ Higgs self-coupling) contain \textbf{$13$ flavour parameters, which alone exceed the corpus's $8$} — and the corpus's $8$ are declared \textit{configuration} data, not constants, so no parameter-by-parameter map exists even in principle without extending P4.

\textbf{Req.} Def-H.29; D0.8, D0.9, D1.2, D1.11, D3.1, D3.2, D4.1–D4.3; CT0.5–CT0.14 (Part 1); CT1.20, CT1.25, CT2.9 (Part 2); CT3.15, CT3.17, CT3.20, CT3.22 (Part 2); CT4.1–CT4.5 (Part 2); EM.9, EM.20 (Part 2); EW.13, EW.20, EW.24, EW.27, EW.33, EW.35 (Part 2); FT.20, FT.27, FT.30, FT.34, FT.59; A §0.3, P1–P5; external SM parameter count (marked) [machine].

\textbf{Pf.}

%dsp-level:0
\[
(1):\ \mathrm{CT0.13}\ \text{Pf by rung induction on}\ \vee^2/\wedge^2\ \text{population and}\ \mathrm{CT0.7}'\text{s}\ S\text{-grading};
\]

%dsp-level:0
\[
\qquad \mathrm{CT0.14}:\ \mathcal K_n=\bigoplus_{A\in\mathcal A(n)}\mathcal K^{(A)}
\]

%dsp-level:0
\[
2^{n+1}\le16\iff n\le3\ [\mathrm{CT0.5},\mathrm{CT0.11}];\qquad O(n)=2^{n+1}-16\implies(16,48)\ \text{at}\ n=4,5\ [\mathrm{CT4.1}]
\]

%dsp-level:0
\[
(2):\ \text{each row's citations are theorems with Pfs in B–G};\qquad n\ge4:\ \mathrm{CT4.4}\ \text{is a non-existence theorem}
\]

%dsp-level:0
\[
(3):\ \text{inventory scans}\ [\mathrm{FT.81}(5)];\qquad\text{EW.27}:\ m_H^2=\text{configuration data}\ [\mathrm{P4}]
\]

%dsp-level:0
\[
(4):\ |P4_{\mathrm{core}}|=8;\quad19=3+1+9+4+1+1;\quad9+4=13>8;\qquad\text{one scale}\ [\mathrm{P1},\text{A}\ \S0.3].\ \blacksquare
\]

\textbf{Phys.} The rung ladder is the part of this corpus that behaves most like a generative spine, and it is derived rather than asserted: at one crystallized edge per vertex only seals exist; add a second and geometry, the weak block, a scalar and uniform dark come online; at three the electromagnetic read becomes available with a phase coefficient of exactly $\pi$; at four edges the vertex oversubscribes, the information is absorbed and locked, and — provably — no gauge structure of any kind survives there. What the ladder never generates is matter's handedness or its flavours: those are not at any rung, and no rung's activation set contains them.

\vspace{2pt}\noindent\rule{\textwidth}{0.4pt}\vspace{2pt}

\end{theorem}

\medskip\noindent\textit{\texorpdfstring{FT.84 — The deficit table and the two-ceiling synthesis (audit pass D)}{FT.84 - The deficit table and the two-ceiling synthesis (audit pass D)}.}\ Structural-constraints record: see Appendix B, Section B.4.3 of the Addendum.\par\medskip

\begin{theorem}[{\texorpdfstring{FT.96 — The reversal-chirality census of the cover sector}{FT.96 - The reversal-chirality census of the cover sector}}]\label{tag:FT-96}
\textbf{Thm.} Under Def-H.28, at (S1)–(S4): $(1)$ \textbf{Rename (premise correction).} "$S$-parity of a cover class" is undefined: CT1.1's $S$ acts on the pair register $\mathfrak H(\mathbb C^4)$ (force blocks), not on crystallization patterns. The native class-level chirality is FT.78(6)'s \textbf{reversal involution} $R\mapsto R^{\mathrm T}$ — the pattern-level charge conjugation [D2.6], induced by no vertex permutation ($\notin\operatorname{Aut}\Gamma$). E1–E3 are executed for this invariant; the word "bias" has no numeric referent and is dropped. $(2)$ \textbf{Tetrahedron (re-derived, two enumerators).} The $101$ covers form $7$ $\operatorname{Aut}$-classes ($5$ bypass $+2$ locked; fresh enumeration matching the stored census orbit-by-orbit); under reversal: \textbf{$5$ self-conjugate $+1$ exchanged pair $\to6$ reversal classes}. The exchanged pair is exactly FT.78(6)'s: the two $\Omega=24$ dissolved-digon classes with component-transposed type multisets $\bigl((1,0),(1,1)^2,(1,2)\bigr)\leftrightarrow\bigl((0,1),(1,1)^2,(2,1)\bigr)$. $(3)$ \textbf{Cube (new, exhaustive on the cover sector).} The $11\,205$ covers form $264$ classes ($166$ bypass $+98$ locked; $Z$-anchors $1\,109\,313/1\,150\,848$ and $\sum\Omega=11\,205$ re-verified); under reversal:

%dsp-level:0
\[
264\ \text{classes}\ =\ 42\ \text{self-conjugate}\ +\ 111\ \text{exchanged pairs}\ \longrightarrow\ 153\ \text{reversal classes},
\]

by sector: bypass $26+2\cdot70=166$, locked $16+2\cdot41=98$; the per-$(\mathrm{sector},K,|\mathcal C|)$ distribution is recorded in the suite output. Every exchanged pair is chirality-blind on all nine stored fields ($K,|\mathcal C|,\Omega,\mathrm{amp},\mathrm{locked},\mathrm{norm\ weight},k_{\max},\mathrm{maskdim},\mathrm{cycles}$) — FT.78(6)'s ansatz-invariance at full census scale. $(4)$ \textbf{E3: no dominant handedness is labelable.} Pair members are indistinguishable by every declared read (clause (3)); a label would require a $C$-odd ledger — FT.78(6)'s own conclusion — and FT.87(8)'s $\operatorname{sign}(k)$ applies to materialized chains, which covers do not carry [FT.80(6)]. Classes are exactly: self-conjugate ($C$-even by construction) or members of an \textit{unordered} pair. "Left-handed dominant / right-handed dominant" has no referent. $(5)$ \textbf{E4: negative.} The reversal invariant takes exactly \textbf{two} values (self / paired) on both skeletons — not three; the exact counts are $6$ (tetra) and $153$ (cube) reversal classes. "Three generations $=$ three distinct $S$-parity bias classes" fails at the invariant level ($2\ne3$), and any generation reading is barred regardless [F2; FT.90(6): generations are a missing label, not a capacity limit].

\textbf{Req.} Def-H.25, Def-H.28; FT.67, FT.68, FT.69, FT.78(6), FT.79, FT.80, FT.87(8), FT.90(6), FT.95(3); D2.6; two enumerators (fresh backtracking vs stored census), exact canonical forms under $\operatorname{Aut}$ [math; machine].

\textbf{Pf.}

%dsp-level:0
\[
(2):\ \text{fresh enumeration}=101\to7\ \text{classes};\qquad\text{stored rows matched field-by-field}\ [\text{machine}]
\]

%dsp-level:0
\[
(3):\ \text{canonical form}=\min_{p\in\operatorname{Aut}}\ \text{sorted edge list of}\ p(R^{\mathrm T}):\ 42\ \text{fixed},\ 111\ \text{2-cycles};
\]

%dsp-level:0
\[
\qquad 42+2\cdot111=264\ \checkmark
\]

%dsp-level:0
\[
(4),(5):\ \text{clause}\ (3)\text{'s field identities};\qquad\#\{\text{invariant values}\}=2\ \blacksquare
\]

\textbf{Phys.} Turn every arrow in the gas around and ask what changed. On the tetrahedron, one pair of shapes trades places and five look at themselves in the mirror; on the cube, a hundred eleven pairs dance and forty-two stand still — a hundred fifty-three reversal classes under the two hundred sixty-four. Nothing that feeds a mass moved anywhere: the mirror pairs are identical in every count the corpus can read. The covers come in handed pairs, and the hands are unnamed — not because the census is incomplete, but because no declared read has ever told left from right. That is the shape of the chirality absence at the census level, and it is now counted exactly.

\vspace{2pt}\noindent\rule{\textwidth}{0.4pt}\vspace{2pt}

\end{theorem}

\begin{theorem}[{\texorpdfstring{FT.98 — The polygon form factors: exact census, saturating growth, fifth negative}{FT.98 - The polygon form factors: exact census, saturating growth, fifth negative}}]\label{tag:FT-98}
\textbf{Thm.} On the bidirectional cycles $C_n$, $n=3..10$ (literal domain, (S1)–(S4)): $(1)$ \textbf{Census (two enumerators per $n$; anchors exact).} Backtracking vs covering-polynomial DP agree everywhere; the sign-summed amplitudes reproduce FT.68(2)'s documented pairs $(Z_{\mathrm{byp}},Z_{\mathrm{lock}})$ \textbf{exactly for all eight $n$}: $(25,48),(353,64),(1161,80),(3121,1632),(11929,10864),(53825,43136),(229417,159888),(844881,775840)$; cover counts $(10,33,56,118,281,649,1423,3166)$ with sector splits $(7{+}3,\ 29{+}4,\ 51{+}5,\ 97{+}21,\ 211{+}70,\ 485{+}164,\ 1051{+}372,\ 2241{+}925)$; $|\mathcal C|_{\max}\le\lfloor3n/4\rfloor$ (bypass) verified per $n$: $(1,2,3,3,4,4,5,6)$. $(2)$ \textbf{Form factors (tail convention, all-sector, exact Fractions, units $\ln2$).}

%dsp-level:0
\[
\begin{array}{c|cccc}
n & M^{(w_U)}=M^{(w_C)}\cdot2 & M^{(w_K)} & M^{(w_E)} & M^{(w_H)}\\
\hline
3 & \tfrac{3264}{865} & \tfrac{9504}{865} & \tfrac{96}{865} & \tfrac{3072}{865}\\
4 & \tfrac{24832}{6273} & \tfrac{1408}{153} & \tfrac{1408}{2091} & \tfrac{10240}{6273}\\
5 & \tfrac{281920}{50081} & \tfrac{694240}{50081} & \tfrac{7840}{50081} & \tfrac{135680}{50081}\\
6 & \tfrac{2808192}{407233} & \tfrac{7585344}{407233} & \tfrac{212160}{407233} & \tfrac{2079744}{407233}\\
7 & \tfrac{29210048}{3804641} & \tfrac{74587296}{3804641} & \tfrac{2663136}{3804641} & \tfrac{17679872}{3804641}\\
8 & \tfrac{300392960}{33646849} & \tfrac{767038208}{33646849} & \tfrac{19632384}{33646849} & \tfrac{178147328}{33646849}\\
9 & \tfrac{2833063488}{282265633} & \tfrac{7303749984}{282265633} & \tfrac{16677792}{21712741} & \tfrac{1757108736}{282265633}\\
10 & \tfrac{27649976960}{2481810241} & \tfrac{71155288000}{2481810241} & \tfrac{2115650880}{2481810241} & \tfrac{17046103040}{2481810241}
\end{array}
\]

with $M^{(w_C)}=M^{(w_U)}/2$ exactly at every $n$ [FT.97(3)]. Per-sector values are tabulated in the suite output. $(3)$ \textbf{Growth: linear-to-sublinear, saturating.} $M^{(w_U)}(C_n)$ rises $3.77\to11.14$ over $n=3..10$ (ratios to $P_3$: $1.049,1.492,1.827,2.035,2.366,2.660,2.953$) — approximately $\approx1.1n$, never exponential in $n$; $M^{(w_K)}$: max ratio $2.609$; $M^{(w_H)}$: max $1.934$; $M^{(w_E)}$: non-monotone (parity oscillation), max $7.681$. $(4)$ \textbf{B4/D1 verdict: the fifth clean negative.} The widest form-factor ratio in the entire computed family is $7.681$ ($w_E$, $n{=}10/3$) against [EXT] targets $207$ (shortfall $\times27.0$) and $3477$ ($\times453.9$); for the primary weight $w_U$ the widest is $2.953$ (shortfalls $\times70.1$ / $\times1177.6$). No aligned triple exists within any factor of $10$ (the widest family ratio over all weights and all $n$ is $7.681$). The "lepton sector" label is declined [F2]; the [EXT] hierarchies enter only this comparison. $(5)$ \textbf{D2 answered exactly.} The $16^{|\mathcal C|}$ weighting \textit{does} concentrate inside each skeleton — exact weight fractions at $|\mathcal C|_{\max}$: $\tfrac{768}{865},\tfrac{2048}{2091},\tfrac{40960}{50081},\tfrac{196608}{407233},\tfrac{3211264}{3804641},\tfrac{16777216}{33646849},\tfrac{50331648}{282265633},\tfrac{1509949440}{2481810241}$ ($0.888,0.979,0.818,0.483,0.844,0.499,0.178,0.608$ — parity-oscillating) — but $|\mathcal C|_{\max}$ itself grows only linearly ($\le3n/4$), so the enhancement \textbf{saturates at the cap}: exponential within, linear across [FT.97(5)]. "Suppress low-node, exponentially enhance high-node" is true intra-skeleton and false as a hierarchy mechanism.

\textbf{Req.} Def-H.25, Def-H.26; FT.67, FT.68(2), FT.69, FT.75, FT.79(2), FT.97; CT0.11 (Part 1); two enumerators (backtracking vs DP), exact Fractions [math; machine]. External targets marked at (4).

\textbf{Pf.}

%dsp-level:0
\[
(1):\ \text{enum\_covers\_bt}=\text{cover\_dp coefficients per level};
\]

%dsp-level:0
\[
\qquad Z=\sum_{\text{level}}c_k4^k=\mathrm{FT.68}(2)\ \text{row verbatim}\ [\text{machine}]
\]

%dsp-level:0
\[
(2),(3):\ \text{exact Fraction averages over the enumerated covers};
\]

%dsp-level:0
\[
\qquad M^{(w_C)}=M^{(w_U)}/2:\ \mathrm{FT.97}(3)\ \text{identity}
\]

%dsp-level:0
\[
(4):\ \max_{w,n}\tfrac{M^{(w)}(C_n)}{M^{(w)}(C_3)}=7.681\ldots;
\]

%dsp-level:0
\[
\qquad 207/7.681=26.95\ldots;
\]

%dsp-level:0
\[
\qquad (5):\ \text{exact fractions as tabulated}\ \blacksquare
\]

\textbf{Phys.} Run the new mass idea along the polygon ladder and the ladder answers in the same voice as every read before it: the averages climb, but they climb like counting, not like cascading — three and a half units at the triangle, eleven at the decagon, and the steepest weighting scheme in the whole family buys a factor of seven and change. The exponential the proposal hoped for is real, but it lives in the \textit{weights}, and the weights are capped by the crystallization budget: three units of $|\mathcal C|$ per four vertices, forever. Ten-sided polygons are not two hundred times heavier than triangles, and no reweighting of the same finite gas will make them so.

\vspace{2pt}\noindent\rule{\textwidth}{0.4pt}\vspace{2pt}

\end{theorem}

\begin{theorem}[{\texorpdfstring{FT.99 — The Platonic form factors, the dominance census, and the reduced-denominator record}{FT.99 - The Platonic form factors, the dominance census, and the reduced-denominator record}}]\label{tag:FT-99}
\textbf{Thm.} At (S1)–(S4): $(1)$ \textbf{Tetrahedron (two enumerators: power-set filter of all $4096$ patterns vs backtracking — identical $101$ covers; $7$-class table re-verified [FT.78(2)]; $Z=(1057,448)$ [FT.68(1)]).} Exact form factors (tail, units $\ln2$):

%dsp-level:0
\[
\begin{array}{c|ccccc}
\text{sector} & w_U & w_K & w_C & w_E & w_H\\
\hline
\text{all} & \tfrac{172800}{35201} & \tfrac{566400}{35201} & \tfrac{86400}{35201} & \tfrac{12672}{35201} & \tfrac{227328}{35201}\\
\text{byp} & \tfrac{62208}{15745}\ [\mathrm{FT.68}(1)(ii)\ \text{anchor}] & \tfrac{27264}{3149} & \tfrac{31104}{15745} & \tfrac{12672}{15745} & \tfrac{18432}{15745}\\
\text{lock} & \tfrac{108}{19} & \tfrac{420}{19} & \tfrac{54}{19} & 0 & \tfrac{204}{19}
\end{array}
\]

(the locked $w_E=0$ is structural: locked tetra covers — double-digon hubs and $K_{1,3}$ stars — have no endpoints at all). Concentration at $|\mathcal C|_{\max}=3$: $16384/35201=0.4654$ (the four $K_{1,3}$ stars). $(2)$ \textbf{Cube (fresh enumeration: $11\,205$ covers, all distinct; cross-checked against the stored class census: $\sum\Omega=11\,205$, $\Omega$-histogram identical, $Z=(1\,109\,313,\ 1\,150\,848)$).} Exact form factors:

%dsp-level:0
\[
\begin{array}{c|ccccc}
\text{sector} & w_U & w_K & w_C & w_E & w_H\\
\hline
\text{all} & \tfrac{482089472}{53648367}\ (8.986) & \tfrac{2241319168}{89413945}\ (25.07) & \tfrac{241044736}{53648367} & \tfrac{256135424}{268241835} & \tfrac{2021359616}{268241835}\\
\text{byp} & \tfrac{655450624}{82871211}\ (7.909) & \tfrac{485339392}{27623737} & \tfrac{327725312}{82871211} & \tfrac{144068864}{82871211} & \tfrac{202792960}{82871211}\\
\text{lock} & \tfrac{285644}{30171}\ (9.468) & \tfrac{285804}{10057} & \tfrac{142822}{30171} & \tfrac{6080}{10057} & \tfrac{887972}{90513}
\end{array}
\]

Concentration at $|\mathcal C|_{\max}=6$: $67108864/804725505=0.0834$ (four covers). $(3)$ \textbf{C3 comparison against the existing reads.} The directive's "single-node aspect-blind masses" are whole-subnetwork reads: $m_{\mathrm{hid}}(\mathrm{tetra\ floor})=12\ln2/\mathfrak h$, $m_{\mathrm{hid}}(\mathrm{cube\ floor})=24\ln2/\mathfrak h$ [FT.42: $2E$]; Ansatz-A ladder extremes $\{0..6\}$, $\{0..12\}$ [FT.79(1)]. The form factors sit \textbf{inside the same $O(V)$ band}: $M^{(w_U)}\mathfrak h/\ln2=4.909$ (tetra), $8.986$ (cube) — cube/tetra $=1.831$ against the $m_{\mathrm{hid}}$ ratio $2$; the widest cube-vs-tetra form-factor ratio over the five weights is $2.652$ ($w_E$, all-sector: $\tfrac{256135424/268241835}{12672/35201}$ exact; the locked sector is degenerate at $w_E$ — tetra locked $=0$ by clause (1)) and the primary-weight ratio is $1.831$. No new range; no hierarchy. $(4)$ \textbf{Dominance census (the directive's flag, exact).} "Intermediate dominates base" := centre$+$hub role bits $>$ endpoint role bits (convention-free half-split). Counts: tetrahedron $64/100$ non-vacuum covers ($|A|^2$-weighted fraction $248/275=0.902$); cube $10\,052/11\,204$ (weighted $3115352/3143459=0.9911$); polygons rising $3/9\ (C_3)\to2725/3165\ (C_{10})$. \textbf{The amplitude-weighted ensemble is generically centre-dominated} — mechanism: an isolated $2$-edge thread splits its bits evenly (one centre bit vs two endpoint half-bits), and dominance arises exactly where vertices \textit{share} roles (types $(1,2),(2,2),\dots$): multi-role vertices accrue centre/hub bits while their partner endpoints accrue fewer. The flag is therefore not an anomaly to hunt but the generic structure of covers at scale. $(5)$ \textbf{D1 verdict: fifth clean negative, with the exact obstruction.} No weight, sector or skeleton in the computed inventory produces ratios within $\times27$ of the [EXT] $207$ or $\times454$ of $3477$; the a priori bound is FT.97(5)'s required-$V$ table ($V\ge65\ldots6561$ by weight, against $V=4/8$); "baryon sector" declined [F2]. The verdicts of FT.77/FT.79(4)/FT.84(5)/FT.87(9) stand unmodified and now cover the form-factor species as well. $(6)$ \textbf{Integrity ledger (this pass's catches, recorded pre-ship).} \textbf{(a)} The first run double-counted registration bits ($i_v=a{+}b$: $\sum_v\deg=2K$), producing $124416/15745=2\times$ the FT.68(1)(ii) anchor — \textbf{caught by the anchor}, corrected to the tail convention, and promoted to FT.97(2a)'s convention clause (the factor-$2$ freedom is real and now documented, not silently fixed). \textbf{(b)} the previously published $\langle m_A^2\rangle_{\mathrm{cube}}=22146188288/268241835$ is in \textbf{lowest terms}: the unreduced census sums are $66438564864/804725505$ (common factor $3$) — the values are identical, \textbf{no erratum needed}; the trap (comparing unreduced sums to reduced Fractions) is recorded so it is not re-hit. \textbf{(c)} Dodeca/icosa/octa form factors deferred (bounded rule — cover enumeration scale; recorded, as in FT.78's Status).

\textbf{Req.} Def-H.25, Def-H.28; FT.16, FT.41, FT.42, FT.67, FT.68(1), FT.78, FT.79, FT.80, FT.94(3), FT.97; two enumerators per skeleton; stored cross-checks; exact Fractions [math; machine]. External data marked at (5).

\textbf{Pf.}

%dsp-level:0
\[
(1):\ 2^{12}\ \text{power-set filter}=\text{backtracking}\ (101=101,\ \text{set-identical});
\]

%dsp-level:0
\[
\qquad M^{(w_U)}_{\mathrm{byp}}=\tfrac{\sum16^{|\mathcal C|}K}{\sum16^{|\mathcal C|}}=\tfrac{62208}{15745}\ [\text{machine}=\mathrm{FT.68}(1)(ii)]
\]

%dsp-level:0
\[
(2):\ \text{fresh }11205\ \text{distinct covers};
\]

%dsp-level:0
\[
\qquad \sum\Omega\cdot16^{|\mathcal C|}=804725505=3\cdot268241835\ [\text{reduction recorded at}\ (6b)]
\]

%dsp-level:0
\[
(3):\ \tfrac{482089472/53648367}{172800/35201}
\]

%dsp-level:0
\[
= \tfrac{482089472\cdot35201}{53648367\cdot172800}\ [\text{machine, exact}];
\]

%dsp-level:0
\[
\qquad (4):\ \text{per-cover role decomposition census}
\]

%dsp-level:0
\[
(5):\ \mathrm{FT.97}(5)\ \text{table};\qquad(6):\ \text{anchor catch}+\gcd(66438564864,804725505)=3\ [\text{machine}]\ \blacksquare
\]

\textbf{Phys.} Weighed node by node under five honest conventions, the tetrahedron's gas comes out at about five bits and the cube's at about nine — the same band every previous read occupied, the same factor-of-two step between four vertices and eight. The one genuinely new picture the census paints is of \textit{who} carries the weight: in the cube's amplitude-weighted ensemble, ninety-nine percent of the mass of the average sits in configurations where the threading junctions and hubs outweigh the thread ends — the gas, weighed by its own amplitudes, is almost pure junction. That is a real structural fact about covers at scale. It is not a hierarchy, and the theorem says no amount of reweighing can turn it into one: the average is hostage to its largest summand, and every summand is hostage to the cap.

\vspace{2pt}\noindent\rule{\textwidth}{0.4pt}\vspace{2pt}

\end{theorem}

\begin{theorem}[{\texorpdfstring{FT.105 — Well-definedness: the topology census from first principles, finiteness, the regularization}{FT.105 - Well-definedness: the topology census from first principles, finiteness, the regularization}}]\label{tag:FT-105}
\textbf{Thm.} Under Def-H.32: $(1)$ \textbf{Polya census.} The all-graph counts derived by exact Burnside over $S_n$ acting on edge pairs, $g(n)=\tfrac1{n!}\sum_\pi2^{\,\mathrm{cycles}(\pi|_{\text{pairs}})}$:

%dsp-level:0
\[
g(1..8)\ =\ (1,\ 2,\ 4,\ 11,\ 34,\ 156,\ 1044,\ 12346)\qquad[\text{exact integers; machine}].
\]

$(2)$ \textbf{Connected counts by Euler inversion.} From $\prod_k(1-x^k)^{-c(k)}=\sum_ng(n)x^n$ [math: the component decomposition of the ordinary generating function], solved exactly over $\mathbb Z$:

%dsp-level:0
\[
c(1..8)\ =\ (1,\ 1,\ 2,\ 6,\ 21,\ 112,\ 853,\ 11117)\qquad[\text{every }c(n)\text{ integral; machine}].
\]

$(3)$ \textbf{Explicit enumeration (independent path).} Canonical-form enumeration over all $2^{\binom n2}$ labelled edge sets, minimum over $S_n$, agrees with $(1)$–$(2)$ at every $n\le5$ (all-graph totals $1,2,4,11,34$; connected $1,1,2,6,21$). $|\mathcal G_4|=10$; the cells by edge count: $E=(0,1,2,3,3,3,4,4,5,6)$. $(4)$ \textbf{Finiteness and regularization.} Per budget cell the configuration space is finite and its size exact: $\sum_{\Gamma\in\mathcal G_4}2^{2E(\Gamma)}=5845$; every $Z^{\bullet}_B(q,S)$ is a finite sum of terms in $\mathbb Z[t,u]$ — well-defined with no convergence question [FT.68(7) pattern]; the declared regularization is the budget filtration $(B,B_K)$ of Def-H.32(1); the unfiltered limit is not constructed. Growth bookkeeping: $g$ doubles-exponentially in $n^2$ [$(1)$'s closed form], $c$ with it — per-budget finiteness is the exact statement, asymptotics are per-family data, not limits. $(5)$ \textbf{New-cell covering census (exact; two enumerators where affordable).} Covering polynomials $(\mathcal Z_{\mathrm{byp}},\mathcal Z_{\mathrm{lock}})$ by cover counts per $|\mathcal C|$:

%dsp-level:0
\[
\begin{array}{c|cc|c}
\text{cell}&\#\text{byp covers}&\#\text{lock covers}&\text{note}\\ \hline
\text{digon}&1&0&\text{vacuum only — the digon itself cancels [FT.74(3)]}\\
P_3&3&1&\\
C_3&7&3&\text{triangle threads; no closed triangle is a cover [FT.69(3)(iv)]}\\
P_4&7&4&\\
\mathrm{star}_4\ (K_{1,3})&9&2&\\
C_4&19&6&\text{the face 4-cycle lies in none of its covers [FT.69(3)(v)]}\\
\mathrm{paw}&29&4&\\
\mathrm{diamond}&45&10&\\
\text{tetra}&85\ (=1{+}24{+}60)&16\ (=12{+}4)&\text{re-derived live}=\text{frozen FT.68(1) rows}
\end{array}
\]

\textbf{Req.} A1, A7; D0.1, D0.8, D0.10; CT0.11 (Part 1); CT1.29 (Part 2); CT4.1 (Part 2); Def-H.25; FT.67, FT.68(7), FT.69(3), FT.74(3); Burnside/Polya enumeration, ordinary-generating-function Euler inversion, canonical forms [math; machine: \texttt{checks/\allowbreak{}verify\_\allowbreak{}H\_\allowbreak{}topology\_\allowbreak{}blind.py} §A–§B].

\textbf{Pf.}

%dsp-level:0
\[
(1):\ \text{cycle index of }S_n\text{ on }\binom n2\text{ pairs, computed per permutation (exact orbit walk)};
\]

%dsp-level:0
\[
\qquad n!\mid\sum_\pi2^{\mathrm{cyc}}\ [\text{machine: integrality asserted}]
\]

%dsp-level:0
\[
(2):\ n\,g(n)=\sum_{k\le n}b(k)g(n-k),\ b(k)=\sum_{d\mid k}d\,c(d):\ \text{exact integer recursion solved upward};
\]

%dsp-level:0
\[
\qquad c(n)\in\mathbb Z\ \forall n\le8\ [\text{machine}]
\]

%dsp-level:0
\[
(3):\ \text{adjacency-bitmask canonical minimum over }S_n;
\]

%dsp-level:0
\[
\qquad \text{counts}=(1,1,2,6,21)/(1,2,4,11,34)\ [\text{machine}]
\]

%dsp-level:0
\[
(4):\ |\mathfrak C(\Gamma)|=2^{2E(\Gamma)};
\]

%dsp-level:0
\[
\qquad \textstyle\sum_{\Gamma\in\mathcal G_4}2^{2E}=2^0+2^2+2^4+2^6+2^6+2^6+2^8+2^8+2^{10}+2^{12}=5845\ [\text{machine}]
\]

%dsp-level:0
\[
(5):\ \text{cover-DP}=\text{branch enumeration on every cell with }2E\le8\ [\text{machine, two enumerators}];
\]

%dsp-level:0
\[
\qquad \text{tetra rows}=\mathrm{FT.68}(1)\ \text{verbatim}.\ \blacksquare
\]

\textbf{Phys.} The integration variable "topology" is now an exactly enumerated space with a stated cutoff: ten connected cells at four nodes, thirty-one at five, one hundred twenty-five at six — counts derived from Burnside and Euler, not quoted. Nothing converges because nothing needs to: every budget cell is a finite exact sum, and the filtration by the ledger's own nat content is the regularization. No particle identification; topology names are bookkeeping.

\vspace{2pt}\noindent\rule{\textwidth}{0.4pt}\vspace{2pt}

\end{theorem}

\begin{theorem}[{\texorpdfstring{FT.114 — The [EXT] audit clause (FT.73(5) pattern)}{FT.114 - The [EXT] audit clause (FT.73(5) pattern)}}]\label{tag:FT-114}
\textbf{Thm.} External-data clause. The hierarchy

%dsp-level:0
\[
1:206.768:3477.2\qquad[\text{EXT: supplied measured charged-lepton mass ratios}]
\]

is supplied measured data. It enters no derivation, no definition, and no coefficient of this document; every number in FT.109–FT.113 is fixed without it. This document identifies no particles. The exact shortfalls of the supplied hierarchy against the derived ratios [FT.113] are

%dsp-level:0
\[
\text{shortfall}_{\mathrm{mid}}=\frac{206.768}{r_{\mathrm{mid}}^2/r_{\min}^2}=\frac{25846/125}{r_2^2/r_1^2}=0.999988799297817\ldots,
\]

%dsp-level:0
\[
\text{shortfall}_{\mathrm{out}}=\frac{3477.2}{r_{\max}^2/r_{\min}^2}=\frac{17386/5}{r_0^2/r_1^2}=0.999921541556935\ldots
\]

%dsp-level:0
\[
\qquad [\text{exact rational division of enclosed ratios; widths}<2^{-50};\ \text{machine}],
\]

and the recorded verdicts are:

%dsp-level:0
\[
(1)\ \ \text{The computed ratios do NOT match [EXT] to [EXT]'s stated precision (mid: }6\text{ significant figures; out: }5\text{): the gaps are enclosure-provably nonzero — }r_2^2/r_1^2-206.768>\tfrac1{1000},
\]

%dsp-level:0
\[
\quad r_0^2/r_1^2-3477.2>\tfrac1{10};
\]

%dsp-level:0
\[
(2)\ \ \text{both shortfalls are strictly below }1\text{ and above }1-2\times10^{-4}:\ \text{the derived ratios exceed the supplied values by relative margins below }2\times10^{-4}\ (1-\text{shortfall}_{\mathrm{mid}}\in(10^{-5},2\times10^{-5}),\ 1-\text{shortfall}_{\mathrm{out}}\in(7\times10^{-5},9\times10^{-5}));
\]

%dsp-level:0
\[
(3)\ \ \text{the numerical coincidence and the structural origin are reported separately: the coincidence is the proximity recorded in (2); the structural origin of the ratios is the normalization data of FT.111–FT.113 — the diagonal coefficient }2\text{, the seal dimension }3\text{, the transfer trace }3\text{ — which contain no [EXT] input. No claim is made that the derivation reproduces the observed hierarchy. No identification of the three eigenvalues with any particle species is made or implied.}
\]

\textbf{Req.} FT.113; FT.73(5) (Tiers F.23–F.25, the audit pattern); P1, P5 (\textbf{A}: no quantity identified with a measured external constant; comparison is a test, never an input); exact rational interval division [math; machine: \texttt{checks/\allowbreak{}verify\_\allowbreak{}H\_\allowbreak{}koide\_\allowbreak{}channel.py} §G]. [EXT] data clause: the supplied values enter this block and no other block of this document (machine-scanned, §H).

\textbf{Pf.}

%dsp-level:0
\[
206.768=\tfrac{206768}{1000}=\tfrac{25846}{125},
\]

%dsp-level:0
\[
\qquad 3477.2=\tfrac{34772}{10}=\tfrac{17386}5\ [\text{exact Fractions in lowest terms}]
\]

%dsp-level:0
\[
\text{shortfall}
\]

%dsp-level:0
\[
= \tfrac{[\mathrm{EXT}]}{q},
\]

%dsp-level:0
\[
\ q\ \text{the enclosed ratio of FT.113}:\ \text{monotone exact interval division by positive enclosures}\ \bigl[\tfrac{[\mathrm{EXT}]}{q_{\mathrm{hi}}},\tfrac{[\mathrm{EXT}]}{q_{\mathrm{lo}}}\bigr];
\]

%dsp-level:0
\[
\qquad \text{widths}<2^{-50}\ [\text{machine §G}]
\]

%dsp-level:0
\[
(1):\ 206.768<\operatorname{lo}\bigl(r_2^2/r_1^2\bigr),
\]

%dsp-level:0
\[
\quad 3477.2<\operatorname{lo}\bigl(r_0^2/r_1^2\bigr)\ [\text{exact rational comparison against disjoint enclosures}]
\]

%dsp-level:0
\[
(2):\ \operatorname{hi}(\text{shortfall})<1;
\]

%dsp-level:0
\[
\qquad 1-\text{shortfall}\in\bigl(10^{-5},2\times10^{-5}\bigl)\cup\bigl(7\times10^{-5},9\times10^{-5}\bigl)\ [\text{enclosure-proven, machine}]
\]

%dsp-level:0
\[
(3):\ \text{the tokens }206.768,
\]

%dsp-level:0
\[
\ 3477.2\ \text{occur only inside this block, and [EXT]-marked content enters no Req, Pf, or coefficient of FT.109–FT.113}\ [\text{machine scan §H};\ \mathrm{FT.73}(5)\ \text{discipline: recorded, not adopted}].\ \blacksquare
\]

\textbf{Phys.} The audit states exactly what is on record and nothing more: a three-line dimensionless hierarchy built from seal integers lands within two parts in $10^4$ of a supplied three-line measured hierarchy, misses it at the measurement's own stated precision, carries no overall scale, and identifies nothing. The near-agreement is a numerical fact about rationals and cosines; the construction remains data-free by its own rules.

\vspace{2pt}\noindent\rule{\textwidth}{0.4pt}\vspace{2pt}

\end{theorem}

\begin{theorem}[{\texorpdfstring{FT.120 — The [EXT] audit ledger: scheme columns, exact shortfalls, and the two exclusions}{FT.120 - The [EXT] audit ledger: scheme columns, exact shortfalls, and the two exclusions}}]\label{tag:FT-120}
\textbf{Thm.} External-data clause (FT.73(5)/FT.114 pattern). The heavy triplet of the up-type and down-type sectors, in three renormalization-scheme columns,

%dsp-level:0
\[
\text{S1 (own-scale}\ \overline{\mathrm{MS}}\text{)}:\ (1.273,\ 4.183,\ 162.5)\ \mathrm{GeV};
\]

%dsp-level:0
\[
\qquad \text{S2}\ (\overline{\mathrm{MS}}\ \text{plus pole top}):\ (1.273,\ 4.183,\ 172.4)\ \mathrm{GeV};
\]

%dsp-level:0
\[
\qquad \text{S3 (pole-type)}:\ (1.67,\ 4.78,\ 172.4)\ \mathrm{GeV},
\]

%dsp-level:0
\[
\text{with ratio box}\quad m_u/m_d=0.462\pm0.020,
\]

%dsp-level:0
\[
\qquad m_s/m_d\in[17,22],
\]

%dsp-level:0
\[
\qquad m^0_{\pi}=134.9768\ \mathrm{MeV}
\]

%dsp-level:0
\[
\qquad [\text{EXT}]\ \text{supplied measured data},
\]

is supplied measured data [EXT]. It enters no derivation, no definition, and no coefficient of this document; every number in FT.116–FT.119 is fixed without it. This document identifies no particles; the correspondence audited below is a test, never an input [P1, P5]. The exact shortfall factors of the supplied hierarchies [EXT] against the derived ratios [FT.119(5)], by monotone rational interval division at widths $<2^{-50}$, are

%dsp-level:0
\[
(1)\quad\text{S1}:\ \ \frac{m_2}{m_1}
\]

%dsp-level:0
\[
= 3.285938\ldots,
\]

%dsp-level:0
\[
\ \frac{m_3}{m_1}
\]

%dsp-level:0
\[
= 127.6512\ldots\implies\text{shortfall}_{\mathrm{mid}}
\]

%dsp-level:0
\[
= 0.922058\ldots,
\]

%dsp-level:0
\[
\quad\text{shortfall}_{\mathrm{out}}\qquad = 0.916122\ldots;
\]

%dsp-level:0
\[
\text{S2}:\ \ \frac{m_3}{m_1}=135.4281\ldots\implies\text{shortfall}_{\mathrm{mid}}=0.922058\ldots,\quad\text{shortfall}_{\mathrm{out}}=0.971934\ldots;
\]

%dsp-level:0
\[
\qquad \text{S3}:\ \ (2.862275\ldots,\ 103.2335\ldots)\implies(0.803176\ldots,\ 0.740881\ldots);
\]

%dsp-level:0
\[
(2)\quad\text{all six shortfalls are strictly below}\ 1\ \text{and enclosure-provably nonzero: the derived ratios exceed the supplied ones in every scheme column};
\]

%dsp-level:0
\[
(3)\quad\textbf{the channel invariant:}
\]

%dsp-level:0
\[
\qquad Q=\frac{\sum_im_i}{\bigl(\sum_i\sqrt{m_i}\bigr)^2}:\quad\text{S1}\ 0.662600\ldots,\quad\text{S2}\ 0.669113\ldots,\quad\text{S3}\ 0.648359\ldots,
\]

%dsp-level:0
\[
\text{equivalently}\quad\frac{B^2}{A^2}=6Q-2:\quad\text{S1}\ 1.97560\ldots,\quad\text{S2}\ 2.01468\ldots,\quad\text{S3}\ 1.89015\ldots
\]

%dsp-level:0
\[
\qquad \text{against the derived normalization value}\ 2\ [\mathrm{FT.111}(3)];
\]

%dsp-level:0
\[
|\,
\]

%dsp-level:0
\[
Q-\tfrac23\,
\]

%dsp-level:0
\[
|<0.01\ \text{in S1 and S2}\ (\text{and}\ 0.01<|\,Q-\tfrac23\,|<0.03\ \text{in S3}):\ \text{the heavy triplet sits on the channel's normalization surface at the}\ 1\%\ \text{level in the}\ \overline{\mathrm{MS}}\ \text{columns};
\]

%dsp-level:0
\[
(4)\quad\textbf{anchored rows}\ [\mathrm{FT.119}(6)]:
\]

%dsp-level:0
\[
\qquad \frac{m_2}{m_3}\bigg|_{\text{derived}}
\]

%dsp-level:0
\[
= 0.025575776\ldots\ \text{vs}\ \frac{4.183}{162.5}
\]

%dsp-level:0
\[
= 0.025741\ldots\ [\text{EXT}]\ \text{S1}:\ \text{relative margin}\ 0.0065<0.01\ (\text{agreement within}\ 1\%);
\]

%dsp-level:0
\[
\frac{m_1}{m_3}\bigg|_{\text{derived}}
\]

%dsp-level:0
\[
= 0.007176757\ldots\ \text{vs}\ \frac{1.273}{162.5}
\]

%dsp-level:0
\[
= 0.007833\ldots\ [\text{EXT}]\ \text{S1}:\ \text{relative margin}\ 0.0838>0.05\ (\text{the inner member is the outlier});
\]

%dsp-level:0
\[
(5)\quad\textbf{no scheme column reproduces both derived ratios at the}\ 1\%\ \textbf{level}\ [\text{machine: exhaustive over the three columns}];
\]

%dsp-level:0
\[
(6)\quad\textbf{obstruction O3 (light triplet, excluded):}\ \ \text{over the whole supplied ratio box}\ m_u/m_d
\]

%dsp-level:0
\[
\in [0.442,0.482],
\]

%dsp-level:0
\[
\ m_s/m_d\qquad \in [17,22],
\]

%dsp-level:0
\[
Q_{\mathrm{light}}\in(0.5461\ldots,\ 0.5804\ldots)\subset(0.54,0.59),
\]

%dsp-level:0
\[
\qquad Q_{\mathrm{light}}<\frac35<\frac23,
\]

%dsp-level:0
\[
\qquad \frac{B^2}{A^2}=6Q_{\mathrm{light}}-2\in(1.25,1.5)\ \text{versus the derived}\ 2;
\]

%dsp-level:0
\[
\text{extrema attained at box corners}\ [\partial Q/\partial\sqrt{r_s}>0,\ \partial Q/\partial\sqrt{r_u}<0\ \text{on the box};\ \text{machine-guarded}]:\ \text{the}\ (A,B)
\]

%dsp-level:0
\[
= (1,\sqrt2)\ \text{normalization is EXCLUDED for the light triplet at every phase}\ \delta\ [\mathrm{FT.119}(4):\ Q\ \text{is phase-free}];
\]

%dsp-level:0
\[
(7)\quad\textbf{obstruction O4 (hadron range, excluded):}\ \ \text{the established hadron mass span}\ \le\frac{11000}{134.9768}
\]

%dsp-level:0
\[
= 81.495\ldots<82\ [\text{EXT}]\ \text{is below the derived colour-channel outer ratio}\ 139.3386880284\ldots\ [\mathrm{FT.119}(5)]\ \text{and far below}\ \mathrm{FT.113}\text{'s}\ 3477.472837104\ldots;
\]

%dsp-level:0
\[
\text{hence no baryon or meson triplet spans the channel's dynamic range}\ —\ \text{the second, independent exclusion of the hadron-level read, the first being FT.117}(7)\text{'s blindness theorem};
\]

%dsp-level:0
\[
(8)\quad\textbf{the dimension trap, disclosed and rejected}\ [\mathrm{FT.77}(3)\ \text{pattern}]:\ \ \text{the literal label}\ m\propto r\ (\text{unsquared})\ \text{gives}\ 1:14.3795\ldots:58.9701\ldots\ \text{on the}\ \delta_\ell
\]

%dsp-level:0
\[
= \tfrac29\ \text{spectrum — missing the supplied mid ratio}\ 206.768\ [\text{EXT}]\ \text{by}\ >14\times,
\]

%dsp-level:0
\[
\text{while the}\ r^2\ \text{read reproduces it within}\ 1\%\ (206.770315972\ldots)\ [\text{EXT}];
\]

%dsp-level:0
\[
\qquad \text{the same literal label on the colour channel gives}\ r^{(q)}_2/r^{(q)}_1
\]

%dsp-level:0
\[
= 1.887775\ldots\ \text{against}\ 3.285938\ldots\ (1.74\times\ \text{off}):\ \text{the squared read is the only one that survives either audit};
\]

%dsp-level:0
\[
(9)\quad\textbf{provenance:}\ \ \text{the structural origin of the derived ratios is the normalization data of FT.116–FT.119 — the diagonal coefficient}\ 2,
\]

%dsp-level:0
\[
\ \text{the seal dimension}\ 3,
\]

%dsp-level:0
\[
\ \text{the transfer trace}\ 3,
\]

%dsp-level:0
\[
\ \text{the colour count}\ 3\ —\ \text{which contain no [EXT] input. The numerical proximity recorded in}\ (1)\text{–}(4)\ \text{and the structural origin are reported separately. No claim is made that the derivation reproduces the observed hierarchy; no identification of the three eigenvalues with any particle species is made or implied.}
\]

\textbf{Req.} FT.119; FT.113, FT.114 (\textbf{Tier F.37}, the audit pattern and the lepton-side decimals); FT.73(5) (Tiers F.23–F.25), FT.77(3) (Tiers F.26–F.28), FT.107 (Tier F.36) for the audit-pattern and read-vocabulary precedents; P1, P5 (\textbf{A}: no quantity identified with a measured external constant; comparison is a test, never an input); exact rational interval division, exact rational square-root enclosures, corner-monotonicity guards [math; machine: \texttt{checks/\allowbreak{}verify\_\allowbreak{}H\_\allowbreak{}colour\_\allowbreak{}channel.py} §E]; floating-point cross-check of the same ledger: \texttt{checks/\allowbreak{}audit\_\allowbreak{}heavy\_\allowbreak{}quark\_\allowbreak{}scheme.py}. [EXT] data clause: the supplied values enter this block and no other block of this document (machine-scanned, §F).

\textbf{Pf.}

%dsp-level:0
\[
(1):\ \text{S1 mid}
\]

%dsp-level:0
\[
= \tfrac{4.183}{1.273}
\]

%dsp-level:0
\[
= \tfrac{4183}{1273},
\]

%dsp-level:0
\[
\ \text{S1 out}
\]

%dsp-level:0
\[
= \tfrac{162.5}{1.273}
\]

%dsp-level:0
\[
= \tfrac{16250}{1273};
\]

%dsp-level:0
\[
\ \text{S2 out}
\]

%dsp-level:0
\[
= \tfrac{172.4}{1.273};
\]

%dsp-level:0
\[
\ \text{S3}
\]

%dsp-level:0
\[
= \tfrac{4.78}{1.67},
\]

%dsp-level:0
\[
\ \tfrac{172.4}{1.67}\ [\text{exact Fractions in lowest terms}]
\]

%dsp-level:0
\[
\text{shortfall}
\]

%dsp-level:0
\[
= \tfrac{[\mathrm{EXT}]}{q},
\]

%dsp-level:0
\[
\ q\ \text{the enclosed ratio of FT.119(5)}:\ \text{monotone exact interval division}\ \bigl[\tfrac{[\mathrm{EXT}]}{q_{\mathrm{hi}}},\tfrac{[\mathrm{EXT}]}{q_{\mathrm{lo}}}\bigr];
\]

%dsp-level:0
\[
\ \text{widths}<2^{-50}\ [\text{machine}]
\]

%dsp-level:0
\[
(2):\ \operatorname{hi}(\text{shortfall})<1\ \text{in all six cases};
\]

%dsp-level:0
\[
\qquad \text{nonzero gaps: }\ q_{\mathrm{lo}}>[\mathrm{EXT}]\ [\text{disjoint rational comparison};\ \text{machine}]
\]

%dsp-level:0
\[
(3):\ Q=\tfrac{\sum m}{(\sum\sqrt m)^2}\ \text{with}\ \sqrt m\ \text{enclosed by integer Newton certificates at}\ 2^{-60};\ \text{widths}<2^{-40};
\]

%dsp-level:0
\[
\qquad \tfrac{B^2}{A^2}=6Q-2\ [\text{exact affine read of}\ \mathrm{FT.111}(3)\text{'s}\ \textstyle\sum(\lambda-A)^2=\tfrac32B^2\ \text{and}\ \sum\lambda=3A]
\]

%dsp-level:0
\[
\text{windows}: \text{S1},\text{S2}\in(\tfrac23-\tfrac1{100},\tfrac23+\tfrac1{100});
\]

%dsp-level:0
\[
\qquad \text{S3}\in(\tfrac23+\tfrac1{100},\tfrac23+\tfrac3{100})\cup(\tfrac23-\tfrac3{100},\tfrac23-\tfrac1{100})\ [\text{machine}]
\]

%dsp-level:0
\[
(4):\ \tfrac{4.183}{162.5}
\]

%dsp-level:0
\[
= \tfrac{4183}{162500};
\]

%dsp-level:0
\[
\quad\bigl|\tfrac{4183}{162500}-\tfrac{(r^{(q)}_2)^2}{(r^{(q)}_0)^2}\bigr|/\tfrac{4183}{162500}<\tfrac1{100};
\]

%dsp-level:0
\[
\qquad \tfrac{1.273}{162.5}=\tfrac{1273}{162500};\quad\text{relative margin}>\tfrac5{100}\ [\text{machine}]
\]

%dsp-level:0
\[
(5):\ \neg\bigl(\operatorname{lo}(\text{shortfall}_{\mathrm{mid}})>\tfrac{99}{100}\wedge\operatorname{lo}(\text{shortfall}_{\mathrm{out}})>\tfrac{99}{100}\bigr)\ \text{in each column}\ [\text{machine, exhaustive}]
\]

%dsp-level:0
\[
(6):\ Q(r_u,r_s)
\]

%dsp-level:0
\[
= \tfrac{r_u+1+r_s}{(\sqrt{r_u}+1+\sqrt{r_s})^2}\ \text{with}\ r_u
\]

%dsp-level:0
\[
= m_u/m_d,
\]

%dsp-level:0
\[
\ r_s
\]

%dsp-level:0
\[
= m_s/m_d;
\]

%dsp-level:0
\[
\quad\partial_{\sqrt{r_s}}Q>0\iff\sqrt{r_s}(1+\sqrt{r_u})>1+r_u\ (\text{true on the box});
\]

%dsp-level:0
\[
\quad\partial_{\sqrt{r_u}}Q<0\iff\sqrt{r_u}(1+\sqrt{r_s})<1+r_s\ (\text{true on the box})
\]

%dsp-level:0
\[
\implies\text{extrema at corners}:\ Q_{\min}=Q(\tfrac{482}{1000},17),\ Q_{\max}=Q(\tfrac{442}{1000},22);
\]

%dsp-level:0
\[
\qquad \tfrac{54}{100}<Q_{\min}\le Q_{\mathrm{light}}\le Q_{\max}<\tfrac{59}{100}<\tfrac35<\tfrac23\ [\text{exact enclosures};\ \text{machine}]
\]

%dsp-level:0
\[
(7):\ \tfrac{11000}{134.9768}
\]

%dsp-level:0
\[
= \tfrac{110000000}{1349768}
\]

%dsp-level:0
\[
= 81.495\ldots<82<139.3386880284\ldots\ [\text{exact rational comparison against FT.119(5)'s enclosure};\ \text{machine}]
\]

%dsp-level:0
\[
(8):\ \tfrac{l_2}{l_1}
\]

%dsp-level:0
\[
= 14.3795\ldots,
\]

%dsp-level:0
\[
\ \tfrac{l_0}{l_1}
\]

%dsp-level:0
\[
= 58.9701\ldots\ \text{on the}\ \delta_\ell\ \text{spectrum}\ [\mathrm{FT.113}\ \text{enclosures re-derived};\ \text{machine}];
\]

%dsp-level:0
\[
\qquad \tfrac{206.768}{14.3795\ldots}>14;
\]

%dsp-level:0
\[
\qquad 0.99<\tfrac{206.768}{206.770315972\ldots}<1.01;
\]

%dsp-level:0
\[
\qquad \tfrac{r^{(q)}_2}{r^{(q)}_1}=1.887775\ldots<2<3.285938\ldots
\]

%dsp-level:0
\[
(9):\ \text{the tokens}\ 1.273,
\]

%dsp-level:0
\[
\ 4.183,
\]

%dsp-level:0
\[
\ 162.5,
\]

%dsp-level:0
\[
\ 172.4,
\]

%dsp-level:0
\[
\ 1.67,
\]

%dsp-level:0
\[
\ 4.78,
\]

%dsp-level:0
\[
\ 0.462,
\]

%dsp-level:0
\[
\ 134.9768\ \text{occur only inside this block, and no [EXT]-marked content enters any Req, Pf, or coefficient of FT.116–FT.119}\ [\text{machine scan}\ \S\mathrm{F}].\ \blacksquare
\]

\textbf{Phys.} The ledger states what is on record and nothing further. A three-line dimensionless hierarchy built from four integers of the frozen algebra — the seal coefficient $2$ and the three factors $3,3,3$ of the dressed arena — has its invariant reading $Q$ landing within one percent of the supplied [EXT] heavy-triplet value in the two running-mass columns, and its top-anchored inner ratio landing within seven parts in a thousand of the supplied own-scale column. It misses both supplied ratios at their own precision, in every column, by an enclosure-provably nonzero margin; the pole column misses by twenty to twenty-six percent. Two branches of the docket are closed by theorem rather than by taste: the light triplet cannot sit on this channel at any phase, because its $Q$ is bounded below $\tfrac35$ across the entire supplied ratio box while the channel's normalization forces exactly $\tfrac23$; and no hadron triplet can reach the channel's range, because the established hadron spectrum spans less than a factor $82$ against the channel's $139$ — the same verdict FT.117 reached structurally. The squared-mass label is the only one that survives either side of the dimension audit. The near-agreement is a numerical fact about rationals and cosines; the construction remains data-free by its own rules.

\vspace{2pt}\noindent\rule{\textwidth}{0.4pt}\vspace{2pt}

\vspace{2pt}\noindent\rule{\textwidth}{0.4pt}\vspace{2pt}

\end{theorem}

\phantomsection
\section[§Consumption — the cell-consumption block]{\texorpdfstring{§Consumption — the cell-consumption block}{SConsumption - the cell-consumption block}}

\textit{Preamble (census supply from Category-A tags): the cover-polynomial tables of FT.68, the frame catalogue of FT.93, and the susceptibility census of FT.103 are census sides of structural tags; Part 4 consumes them through those tags' own rows — no B-classification is implied. No census entry is dismissed: each is a load-bearing cell supply (separation charter).}

\textbf{FT.28.} Space: the two-vertex digon floor $\Sigma_{\min}$ with a single observer contact; registration data $(\sigma_e,\kappa)$). Counts: the complete floor card: $I_{\mathrm{hid}}=2\ln2$, $\Theta_{\mathrm{int}}=2\kappa$, winding, holonomy, aspect rows, mask dimension $\le10$, rung counts, the ternary core read. Baryonic: the minimal fermion cell and its card grammar — every sealed composite card (baryon-type junctions included) refines this row template; the digon is the junction kernel of the composition laws.

\textbf{FT.32.} Space: the single-contact cycle floors $\Sigma_L$ (unidirectional $L$-cycles). Counts: the ladder $m_{\mathrm{hid}}(L)=L\ln2/\mathfrak h$, the minimality regimes, structural colourlessness per $L$. Baryonic: the winding ladder baryonic chains compose over; additive winding composition [FT.34(3)] consumes these cells as summands.

\textbf{FT.35.} Space: the composite floor $\Sigma_{12}$ (first colour-active, first lock-carrying). Counts: the junction lock content, the baryon-type sealing phase, the maximal $18$-dimensional mask. Baryonic: the first computed colour-active cell — the junction grammar baryon seals run, and the $D^-$-carrying cell for the $n\ge4$ lock side.

\textbf{FT.42.} Space: the five Platonic skeletons under the uniform-minimal materialization scope (M1/M2). Counts: the card template per solid: $E_{\mathrm{int}}=2E$, $m_{\mathrm{hid}}$, $\Pi$, winding and lock profiles. Baryonic: the symmetric-skeleton cell pool; the tetrahedron's trivalent vertices are the native baryon-junction site.

\textbf{FT.43.} Space: per-solid winding-ledger states with Euler consistency. Counts: $F-1$ independent winding reads, degeneracies, per-face winding assignments. Baryonic: the winding bookkeeping sealed baryonic chains inherit; composition adds windings cell-wise.

\textbf{FT.44.} Space: the computed floors under solid duality. Counts: dual pairs sharing $m_{\mathrm{hid}}/\Pi$, split by winding/sealing/lock profiles; the taxonomy rows across every computed floor. Baryonic: the degeneracy-splitting grammar — how baryonic multiplets sharing a skeleton are told apart by ledger data.

\textbf{FT.45.} Space: the full regular-polygon floor family $\Sigma_{P_n}$, $n\ge3$. Counts: walk-space rank $n+1$, the single parity constraint, four minimality regimes, injective mass/rank/null reads per $n$. Baryonic: the unbounded cell family for chain composition and the polygon-limit anchor of the flicker boundaries [FT.70(3)].

\textbf{FT.49.} Space: the tetrahedron-face floor with baryon-type sealing. Counts: the full card: $\Theta_{\mathrm{int}}\in4\mathbb Z$, the $30$-dimensional exchange sector, the sealing grammar. Baryonic: the reference baryon-type sealing-phase cell — its junction structure is the template for $\varepsilon$-sealed triplet reads.

\textbf{FT.50.} Space: patch cells (bowtie, double-triangle) and floor pairs under gluing. Counts: junction-saturation data, the constraint-rank gluing law, glued-card rows. Baryonic: the composition laws by which multi-face baryonic skeletons are assembled from computed cells.

\textbf{FT.51.} Space: closed-walk winding ledgers per skeleton under dihedral symmetry. Counts: closed-form enumerations with degeneracy counts. Baryonic: the degeneracy pool for the sector labels $(q,S)$ of conditioned baryonic path integrals [Def-H.32(c) grammar].

\textbf{FT.63.} Space: the five Platonic floors with all physical read columns. Counts: spin (under Def-H.38), charge, parity, holonomy, entropy, mass ratios, binding — per solid; $j_{\max}\in\{6,6,12,30,30\}$. Baryonic: the per-solid read table baryonic multiplets consume as constituent data on symmetric skeletons.

\textbf{FT.64.} Space: admissible-loop lattices per solid under the filter chain. Counts: the generalized loop lists $A_S(K)$ with the exact character-sum closed form. Baryonic: the loop pool over which baryonic cover cells are built; the filter chain is the admissibility grammar.

\textbf{FT.65.} Space: coarse-grained loop classes under the full automorphism groups. Counts: the degeneracies $M_S(K)/M^+_S(K)$ and the register product $2^{2E}M_S(K)g(V)$. Baryonic: the register-counting side of sealed-object enumeration — multiplicity bookkeeping for baryonic state counts.

\textbf{FT.69.} Space: Platonic cover patterns with amplitude weights. Counts: the completed amplitude census (cover polynomials; the $62208/15745$ anchor), the cycle-alternation rows. Baryonic: the amplitude-layer data for flicker ensembles on symmetric skeletons; the norm/amplitude dichotomy [FT.68(3)] consumes these values.

\textbf{FT.76.} Space: Platonic defect-angle pole reads under both winding-pairing variants (A/B). Counts: the multiplet tables, the ceiling $[0,4\hat\mu'^2]$, the four $\Theta_R$ readings priced. Baryonic: the pole-side mass cells and their range taxonomy — the hard ceiling every baryonic pole read inherits.

\textbf{FT.78.} Space: tetrahedron and cube cover configurations under $\operatorname{Aut}$. Counts: tetrahedron: $70$ admissible / $7$ cover classes; cube: $33\,564$ admissible by $K$ / $264$ cover classes with $|\mathcal C|$-levels, signatures, chirality pairs. Baryonic: the cover-class pool baryonic sector-conditioning runs over; the exact enumeration substrate for cube-era cells.

\textbf{FT.81.} Space: the corpus read classes against the external-reference content rows. Counts: presence/absence per read class (charge, spin, winding, sealing, chirality, flavour) — the internal read-class inventory (external scoring columns excluded). Baryonic: the read-class grammar every baryonic card fills, and the audit template for P1/P5-compliant external comparison.

\textbf{FT.83.} Space: the rungs $n=0,1,2,3,\ge4$ with their activation sets. Counts: derived/declared/open status per rung; the $n\ge4$ derived no-go row. Baryonic: the activation map Part 4 extends at $n=3$ (colour, derived) and $n\ge4$ (the lock side).

\textbf{FT.84.} Space: the five mass-read species with their ranges. Counts: the ceiling taxonomy: two hard ceilings, two size-scaling species, one undeclared-fugacity direction (external scoring excluded). Baryonic: the range constraints any baryonic mass read inherits, and the species taxonomy new reads must join.

\textbf{FT.96.} Space: tetra/cube cover orbits under the orientation reversal. Counts: the chiral cover pairs (the two $\Omega=24$ orbits exchanged), reversal signatures per orbit. Baryonic: the chirality-tagged cover cells parity reads [FT.137] consume for sealed objects.

\textbf{FT.98.} Space: polygon skeletons under the multi-node form-factor species. Counts: exact $M$-tables, saturating growth rows, reduction-identity anchors. Baryonic: form-factor range data for extended baryonic skeletons (the linear size-bound species).

\textbf{FT.99.} Space: Platonic skeletons under the form-factor reads. Counts: the dominance census, reduced-denominator record, concentration fractions. Baryonic: the symmetric-skeleton form-factor cells for baryonic multiplet range checks.

\textbf{FT.105.} Space: connected graphs and canonical skeletons to $V\le8$ (canonical enumeration to $V\le5$). Counts: exact Burnside/Polya counts $g(n)$, Euler inversion $c(n)$, finiteness and regularization cells per budget. Baryonic: the first-principles topology pool over which the topology-blind baryonic integral is regulated (budget filtration).

\textbf{FT.114.} Space: the supplied external hierarchy against the derived $\ell$-channel ratios. Counts: exact shortfall arithmetic on the quarantined rows. Baryonic: the audit-discipline template for Part-4 external tests (P1/P5-compliant comparison grammar).

\textbf{FT.120.} Space: renormalization-scheme columns against the derived $q$-channel ratios; the light-triplet $Q$-box. Counts: the scheme ledger, exact shortfalls, the two derived exclusions. Baryonic: the scheme-dependence record Part-4 comparisons inherit — no canonical renormalization point is derivable in the corpus.


\chapter*{\texorpdfstring{Index of Corpus Tags}{Index of Corpus Tags}}
\addcontentsline{toc}{chapter}{Index of Corpus Tags}
\markboth{Index of Corpus Tags}{Index of Corpus Tags}
The complete tag registry, with a one-line description of every tag: 623 registry tags (481 theorem and data-cell tags, 115 definitions, 9 axioms, 13 conventions, 5 calibration-policy clauses) plus the adjudication-record tags of HistLog \S IV reproduced in the Addendum. Rows are grouped by document, in tier order; ``Presentation'' gives the theorem/definition number in this volume or the Addendum section (B.$x$.$y$) for relocated records.
\begingroup\small
\begin{longtable}{@{}p{0.10\textwidth}p{0.115\textwidth}p{0.465\textwidth}p{0.095\textwidth}p{0.135\textwidth}@{}}
\toprule
\textbf{Tag} & \textbf{Class} & \textbf{Title / description} & \textbf{Where} & \textbf{Presentation} \\
\midrule
\endhead
\multicolumn{5}{@{}l}{\textbf{Part 1 · Qubit\_era — axioms, conventions, definitions, calibration, pair canon, ledger dynamics}}\\[2pt]
\multicolumn{5}{@{}l@{}}{\hspace*{1em}\small\itshape declarations}\\[1pt]
\texttt{A1} & axiom & --- & Part 1 & Axiom 1.1 \\
\texttt{A2} & axiom & --- & Part 1 & Axiom 1.2 \\
\texttt{A3} & axiom & --- & Part 1 & Axiom 1.3 \\
\texttt{A4} & axiom & --- & Part 1 & Axiom 1.4 \\
\texttt{A5} & axiom & --- & Part 1 & Axiom 1.5 \\
\texttt{A6} & axiom & --- & Part 1 & Axiom 1.6 \\
\texttt{A7} & axiom & --- & Part 1 & Axiom 1.7 \\
\texttt{A8} & axiom & --- & Part 1 & Axiom 1.8 \\
\texttt{A9} & axiom & --- & Part 1 & Axiom 1.9 \\
\texttt{C1} & convention & Pauli basis / emergent spacetime coordinates where a metric exists; : spatial/Pauli-spatial; : colour adjoint.… & Part 1 & Convention 1.1 \\
\texttt{C10} & convention & Register accounting (reality principle): registrations convert quantum data to classical alternatives at the l… & Part 1 & Convention 1.10 \\
\texttt{C11} & convention & A5/A6 reading: the conjugation clause is node-level — [A6: ] with and ( : well-typed); is the vacuum pair stat… & Part 1 & Convention 1.11 \\
\texttt{C12} & convention & In A8's the summation domain reads (D0.10): the configuration set on the real-edge set. itself (A4, D0.1) is a… & Part 1 & Convention 1.12 \\
\texttt{C13} & convention & Signature convention. A3 fixes only through the products , and fixes only up to its odd slot and its overall s… & Part 1 & Convention 1.13 \\
\texttt{C2} & convention & --- & Part 1 & Convention 1.2 \\
\texttt{C3} & convention & Symbol : antisymmetrization only ( , -contractions); a sector-label conjunct & Part 1 & Convention 1.3 \\
\texttt{C4} & convention & sector-to-metric embedding scale (D2.1); : gravitational coupling (D2.9, CT2.9); ; interchangeable & Part 1 & Convention 1.4 \\
\texttt{C5} & convention & (operator); (scalar, D2.1); interchangeable & Part 1 & Convention 1.5 \\
\texttt{C6} & convention & transit entropy; : bound entropy (D1.5); squared-field notation for either & Part 1 & Convention 1.6 \\
\texttt{C7} & convention & (A8). Rung-resolved: (D2.3) & Part 1 & Convention 1.7 \\
\texttt{C8} & convention & neutral-weak aspect of A8; ; pre-mixing seed generator: . Bare in a spectrum: physical neutral eigenstate (D3.… & Part 1 & Convention 1.8 \\
\texttt{C9} & convention & Entropy in nats: ; \_bits\_ only where explicit. One real edge bit nats (CT0.9, CT1.29) & Part 1 & Convention 1.9 \\
\multicolumn{5}{@{}l@{}}{\hspace*{1em}\small\itshape Tier 0 — Bookkeeping, Reality, and Capacity}\\[1pt]
\texttt{CT0{.}1} & theorem & Real/virtual dichotomy & Part 1 & Theorem 1.1 \\
\texttt{CT0{.}10} & theorem & The two sign structures are independent & Part 1 & Theorem 1.10 \\
\texttt{CT0{.}11} & theorem & Maximal accessible rung & Part 1 & Theorem 1.11 \\
\texttt{CT0{.}11a} & theorem & three slots & Part 1 & Corollary 1.3 \\
\texttt{CT0{.}11b} & theorem & directed rung of bidirectional skeletons & Part 1 & Corollary 1.4 \\
\texttt{CT0{.}12} & theorem & Aspect–capacity match & Part 1 & Theorem 1.12 \\
\texttt{CT0{.}13} & theorem & Aspect activation and the active sets & Part 1 & Theorem 1.13 \\
\texttt{CT0{.}14} & theorem & Mask flag and block dimensions & Part 1 & Theorem 1.14 \\
\texttt{CT0{.}2} & theorem & Global vacuum & Part 1 & Theorem 1.2 \\
\texttt{CT0{.}3} & theorem & Seal increments & Part 1 & Theorem 1.3 \\
\texttt{CT0{.}4} & theorem & Real outgoing edge registration & Part 1 & Theorem 1.4 \\
\texttt{CT0{.}5} & theorem & Tick load at rung & Part 1 & Theorem 1.5 \\
\texttt{CT0{.}6} & theorem & Channel capacity and oversubscription & Part 1 & Theorem 1.6 \\
\texttt{CT0{.}6a} & theorem & transit-hop register bound & Part 1 & Corollary 1.1 \\
\texttt{CT0{.}7} & theorem & Reality is the swap grading & Part 1 & Theorem 1.7 \\
\texttt{CT0{.}7a} & theorem & spin content & Part 1 & Corollary 1.2 \\
\texttt{CT0{.}8} & theorem & Admissible range of the crystallization threshold & Part 1 & Theorem 1.8 \\
\texttt{CT0{.}9} & theorem & The edge phase alphabet is binary & Part 1 & Theorem 1.9 \\
\multicolumn{5}{@{}l@{}}{\hspace*{1em}\small\itshape Tier 1A — The Pair Operator Algebra}\\[1pt]
\texttt{CT1{.}0} & theorem & Signature classification of the -grading & Part 1 & Theorem 1.15 \\
\texttt{CT1{.}1} & theorem & Joint decomposition of & Part 1 & Theorem 1.16 \\
\texttt{CT1{.}10} & theorem & decomposition of & Part 1 & Theorem 1.25 \\
\texttt{CT1{.}11} & theorem & The two scalars are distinct & Part 1 & Theorem 1.26 \\
\texttt{CT1{.}11a} & theorem & The involution table of the four sectors & Part 1 & Theorem 1.27 \\
\texttt{CT1{.}11b} & theorem & The native Clifford frames of the pair register and the weak-covariance pincer & Part 1 & Theorem 1.28 \\
\texttt{CT1{.}2} & theorem & Generator basis and sector assignment & Part 1 & Theorem 1.17 \\
\texttt{CT1{.}3} & theorem & Centrality of the identity & Part 1 & Theorem 1.18 \\
\texttt{CT1{.}4} & theorem & The close & Part 1 & Theorem 1.19 \\
\texttt{CT1{.}5} & theorem & Commutators of with & Part 1 & Theorem 1.20 \\
\texttt{CT1{.}6} & theorem & Commutators and & Part 1 & Theorem 1.21 \\
\texttt{CT1{.}7} & theorem & Mixed commutators & Part 1 & Theorem 1.22 \\
\texttt{CT1{.}8} & theorem & Action of on and & Part 1 & Theorem 1.23 \\
\texttt{CT1{.}9} & theorem & Sector closure table and the global algebra & Part 1 & Theorem 1.24 \\
\multicolumn{5}{@{}l@{}}{\hspace*{1em}\small\itshape declarations}\\[1pt]
\texttt{D0{.}1} & definition & global vacuum & Part 1 & Definition 1.1 \\
\texttt{D0{.}10} & definition & configuration set on : finite oriented edge sets with sign assignments , seal data (D1.2) and vertex types (A7… & Part 1 & Definition 1.10 \\
\texttt{D0{.}2} & definition & --- & Part 1 & Definition 1.2 \\
\texttt{D0{.}3} & definition & (faithful ). [A4, D0.4, CT0.8]; for . Gibbs tilt (functional calculus): where ; on states: ; the identity on t… & Part 1 & Definition 1.3 \\
\texttt{D0{.}4} & definition & threshold of A4; admissible range: CT0.8; sharp choice & Part 1 & Definition 1.4 \\
\texttt{D0{.}5} & definition & phase alphabet of a real edge (CT0.9: ) & Part 1 & Definition 1.5 \\
\texttt{D0{.}6} & definition & hop one traversal of one ; ; length assigned to one hop & Part 1 & Definition 1.6 \\
\texttt{D0{.}7} & definition & ordered record of retained events [A6, A7: stratification , , monotone under ticks]; ledger ticks (retained up… & Part 1 & Definition 1.7 \\
\texttt{D0{.}8} & definition & (vertex set); (type set); “at rung ” on & Part 1 & Definition 1.8 \\
\texttt{D0{.}9} & definition & independent real registers a vertex of pushes across one pair channel per tick — counted by the register accou… & Part 1 & Definition 1.9 \\
\texttt{D1{.}1} & definition & rung class of ; D0.5; registration . Reality principle of registration (C10): phase of the registration transp… & Part 1 & Definition 1.11 \\
\texttt{D1{.}10} & definition & node force: resolves local information content into channels bound, transit, holonomy, charge, flux ; bound, t… & Part 1 & Definition 1.20 \\
\texttt{D1{.}11} & definition & accessible information content at ; content beyond the capacity bound (CT0.11): accessible [D4.1] & Part 1 & Definition 1.21 \\
\texttt{D1{.}12} & definition & aspects of acting nontrivially at rung ; value derived (CT0.13), stipulated & Part 1 & Definition 1.22 \\
\texttt{D1{.}2} & definition & seal transition of one pair across the A4 criterion (entry/exit of either branch): (crystallization); (dissolu… & Part 1 & Definition 1.12 \\
\texttt{D1{.}3} & definition & Janus-balanced & Part 1 & Definition 1.13 \\
\texttt{D1{.}4} & definition & chain : , ; ; ; (formal difference of vertices); closed & Part 1 & Definition 1.14 \\
\texttt{D1{.}5} & definition & bound closed self-contained chain (content: fermionic); transit propagating open chain (content: bosonic) & Part 1 & Definition 1.15 \\
\texttt{D1{.}6} & definition & [ D2.1] & Part 1 & Definition 1.16 \\
\texttt{D1{.}7} & definition & [ per-edge phase, D3.5; uniform-phase case: ]; ; ; wrapping spatial hop projection has winding around a closed… & Part 1 & Definition 1.17 \\
\texttt{D1{.}8} & definition & accumulated along a causal chain of real edges; emergent coordinate & Part 1 & Definition 1.18 \\
\texttt{D1{.}9} & definition & tick locality: , pre- and post-tick data only; read at edge-distance & Part 1 & Definition 1.19 \\
\texttt{D2{.}1} & definition & sector coefficients [basis: CT1.2] & Part 1 & Definition 1.23 \\
\texttt{D2{.}10} & definition & even blocks ; odd blocks ; Hermitian ; ; ( homogeneous) & Part 1 & Definition 1.32 \\
\texttt{D2{.}11} & definition & valid universe real-edge network closed [D1.3] active mask dangling ledger boundary beyond the sealed Janus bo… & Part 1 & Definition 1.33 \\
\texttt{D2{.}12} & definition & Aspect realization (slot aspect bridge): — linear realization of the rung- slot algebra in aspect space, with … & Part 1 & Definition 1.34 \\
\texttt{D2{.}13} & definition & Subnetwork and microstate space *(scope clauses)*. [D2.11] a finite connected vertex set together with its int… & Part 1 & Definition 1.35 \\
\texttt{D2{.}14} & definition & Colour boundary and the contraction-completeness projector. ; for valid , the other members of the elementary … & Part 1 & Definition 1.36 \\
\texttt{D2{.}2} & definition & --- & Part 1 & Definition 1.24 \\
\texttt{D2{.}3} & definition & aspect space; block of aspect ; ; orthogonal projection & Part 1 & Definition 1.25 \\
\texttt{D2{.}4} & definition & poles of rung- spectrum; residues amplitudes & Part 1 & Definition 1.26 \\
\texttt{D2{.}5} & definition & massless resolvent pole at ; massive pole at ; pole displacement & Part 1 & Definition 1.27 \\
\texttt{D2{.}6} & definition & orientation-reversal involution: ; chirality chirality; charge charge & Part 1 & Definition 1.28 \\
\texttt{D2{.}7} & definition & ledger reflection: reverses ledger order & Part 1 & Definition 1.29 \\
\texttt{D2{.}8} & definition & convergence sheet: , , for resolvents and configuration sums & Part 1 & Definition 1.30 \\
\texttt{D2{.}9} & definition & units powers of ; ; ; single free scale; names for combinations of (CT2.9), independent inputs & Part 1 & Definition 1.31 \\
\texttt{D3{.}1} & definition & slot index colour & Part 1 & Definition 1.37 \\
\texttt{D3{.}10} & definition & [ EW.24] & Part 1 & Definition 1.46 \\
\texttt{D3{.}2} & definition & subset principle: strictly new at rung part of the rung- structure realized at any & Part 1 & Definition 1.38 \\
\texttt{D3{.}3} & definition & at of type , edge signs & Part 1 & Definition 1.39 \\
\texttt{D3{.}4} & definition & (oriented chain ). Closed chains: the basepoint core read is defined, ternary, and charge-conjugation odd [EM.… & Part 1 & Definition 1.40 \\
\texttt{D3{.}5} & definition & [D0.5, D1.6]; per-edge dynamic phase for — the A4 branch index (read in Part 2 §E: Def-EM.20, EM.20, forward p… & Part 1 & Definition 1.41 \\
\texttt{D3{.}6} & definition & --- & Part 1 & Definition 1.42 \\
\texttt{D3{.}7} & definition & photon trivial chain [D1.7] ; gluon trivial chain & Part 1 & Definition 1.43 \\
\texttt{D3{.}8} & definition & colour triplet : one fundamental slot index; colour antitriplet : one dual slot index; quark bound (fermionic)… & Part 1 & Definition 1.44 \\
\texttt{D3{.}9} & definition & baryon ; meson . [Disambiguation: here total mask of A9.] & Part 1 & Definition 1.45 \\
\texttt{D4{.}1} & definition & [A8]: single possibility space receiving configurations with rung & Part 1 & Definition 1.47 \\
\texttt{D4{.}2} & definition & absorbs information at rate [D0.9]; absorbed locked retained in , accessible [D1.11] & Part 1 & Definition 1.48 \\
\texttt{D4{.}3} & definition & [A8]: uniform scalar aspect; configuration-dependent & Part 1 & Definition 1.49 \\
\multicolumn{5}{@{}l@{}}{\hspace*{1em}\small\itshape Tier 2 — Lyapunov Stability of the Pair State}\\[1pt]
\texttt{L0} & theorem & The ledger–tilt coupling & Part 1 & Theorem 1.29 \\
\texttt{L0a} & theorem & tilt modular update & Part 1 & Corollary 1.5 \\
\texttt{L1} & theorem & The Lyapunov function of the pair state & Part 1 & Theorem 1.30 \\
\texttt{L2} & theorem & The temporal decay law & Part 1 & Theorem 1.31 \\
\texttt{L3} & theorem & Entropy deficit and the inertial mass reading & Part 1 & Theorem 1.32 \\
\texttt{L4} & theorem & The patch entropy deficit & Part 1 & Theorem 1.33 \\
\texttt{L5} & theorem & The second-order sector flow of the tilt & Part 1 & Theorem 1.34 \\
\multicolumn{5}{@{}l@{}}{\hspace*{1em}\small\itshape declarations}\\[1pt]
\texttt{P1} & calibration policy & quantities of the corpus: derived from A1–A9 C1–C13 D0.1–D4.3, except the hop length ; ; fixes the absolute sc… & Part 1 & P1 \\
\texttt{P2} & calibration policy & Derived, zero freedom & Part 1 & P2 \\
\texttt{P3} & calibration policy & Named combinations of (external comparison only; inputs) & Part 1 & P3 \\
\texttt{P4} & calibration policy & Free internal configuration parameters ( constants of the theory). Independent core: . Functionally dependent … & Part 1 & P4 \\
\texttt{P5} & calibration policy & quantity identified with a measured external constant. Comparison procedure: fix once; read P2 P3; result test… & Part 1 & P5 \\
\addlinespace[4pt]
\multicolumn{5}{@{}l}{\textbf{Part 2 · Bosonic\_era — geometry, electroweak, electromagnetic, colour, saturation}}\\[2pt]
\multicolumn{5}{@{}l@{}}{\hspace*{1em}\small\itshape §C.I — Emergent Metric Geometry}\\[1pt]
\texttt{CT1{.}12} & theorem & Metric emergence from sector components & Part 2 & Theorem 2.1 \\
\texttt{CT1{.}12a} & theorem & The pair-channel audit of the metric dictionary: tiling, block split, and the three discrete read classes & Part 2 & Theorem 2.2 \\
\texttt{CT1{.}13} & theorem & Christoffel symbols under conformal rescaling & Part 2 & Theorem 2.3 \\
\texttt{CT1{.}14} & theorem & Riemann tensor under conformal rescaling & Part 2 & Theorem 2.4 \\
\texttt{CT1{.}15} & theorem & Ricci tensor under conformal rescaling & Part 2 & Theorem 2.5 \\
\texttt{CT1{.}16} & theorem & Ricci scalar under conformal rescaling & Part 2 & Theorem 2.6 \\
\texttt{CT1{.}17} & theorem & Einstein tensor under conformal rescaling in & Part 2 & Theorem 2.7 \\
\texttt{CT1{.}18} & theorem & Linearized spin-2 operator & Part 2 & Theorem 2.8 \\
\multicolumn{5}{@{}l@{}}{\hspace*{1em}\small\itshape §C.III — Continuum Limit, Effective Action, and Field Equations}\\[1pt]
\texttt{CT1{.}19} & theorem & Continuum limit at leading order & Part 2 & Theorem 2.18 \\
\texttt{CT1{.}19a} & theorem & corollary & Part 2 & Corollary 2.1 \\
\texttt{CT1{.}19b} & theorem & On-shell line-element quantization: even antialigned winding, orientation evenness, the geometry quantum, internal holonomy & Part 2 & Theorem 2.19 \\
\texttt{CT1{.}20} & theorem & Effective action from the partition function & Part 2 & Theorem 2.20 \\
\texttt{CT1{.}21} & theorem & The geometric part of the effective action & Part 2 & Theorem 2.21 \\
\texttt{CT1{.}22} & theorem & Stress-energy tensor from the force operator & Part 2 & Theorem 2.22 \\
\texttt{CT1{.}23} & theorem & Field equations & Part 2 & Theorem 2.23 \\
\texttt{CT1{.}24} & theorem & Conservation from the field equations & Part 2 & Theorem 2.24 \\
\texttt{CT1{.}25} & theorem & The cosmological constant is the aspect & Part 2 & Theorem 2.25 \\
\texttt{CT1{.}26} & theorem & Conservation from mask invariance alone & Part 2 & Theorem 2.26 \\
\texttt{CT1{.}27} & theorem & Quantum geometry state & Part 2 & Theorem 2.27 \\
\texttt{CT1{.}28} & theorem & Singularity resolution by the saturation lock & Part 2 & Theorem 2.28 \\
\multicolumn{5}{@{}l@{}}{\hspace*{1em}\small\itshape §C.IV — Entropy, Calibration, and the Single Free Scale}\\[1pt]
\texttt{CT1{.}29} & theorem & Boundary entropy from channel capacity & Part 2 & Theorem 2.29 \\
\texttt{CT1{.}29a} & theorem & The mean edge quantum: structure, sign, and the exactness locus of the entropy–area rate & Part 2 & Theorem 2.30 \\
\multicolumn{5}{@{}l@{}}{\hspace*{1em}\small\itshape §C.II — Masks, Manifold, and Closure}\\[1pt]
\texttt{CT2{.}1} & theorem & Unit and dimension conventions & Part 2 & Theorem 2.9 \\
\texttt{CT2{.}10} & theorem & Graded matrices and the supertrace & Part 2 & Theorem 2.17 \\
\multicolumn{5}{@{}l@{}}{\hspace*{1em}\small\itshape §C.IV — Entropy, Calibration, and the Single Free Scale}\\[1pt]
\texttt{CT2{.}11} & theorem & One free scale & Part 2 & Theorem 2.32 \\
\multicolumn{5}{@{}l@{}}{\hspace*{1em}\small\itshape §C.II — Masks, Manifold, and Closure}\\[1pt]
\texttt{CT2{.}2} & theorem & The rung mask is a projection & Part 2 & Theorem 2.10 \\
\texttt{CT2{.}3} & theorem & Masked force operator & Part 2 & Theorem 2.11 \\
\texttt{CT2{.}4} & theorem & Force restriction by rung & Part 2 & Theorem 2.12 \\
\texttt{CT2{.}5} & theorem & Masked resolvent as spectrum generator & Part 2 & Theorem 2.13 \\
\texttt{CT2{.}6} & theorem & Continuum manifold closure & Part 2 & Theorem 2.14 \\
\texttt{CT2{.}7} & theorem & Exact geometric action closure & Part 2 & Theorem 2.15 \\
\texttt{CT2{.}8} & theorem & Reflection positivity and Hilbert-space closure & Part 2 & Theorem 2.16 \\
\multicolumn{5}{@{}l@{}}{\hspace*{1em}\small\itshape §C.IV — Entropy, Calibration, and the Single Free Scale}\\[1pt]
\texttt{CT2{.}9} & theorem & Calibration closure & Part 2 & Theorem 2.31 \\
\multicolumn{5}{@{}l@{}}{\hspace*{1em}\small\itshape §F.I — Slot Structure}\\[1pt]
\texttt{CT3{.}1} & theorem & Classical trivalent symmetry & Part 2 & Theorem 2.98 \\
\multicolumn{5}{@{}l@{}}{\hspace*{1em}\small\itshape §F.II — The Rung-3 Gauge Algebra}\\[1pt]
\texttt{CT3{.}10} & theorem & The colour aspect has no central component & Part 2 & Theorem 2.108 \\
\texttt{CT3{.}10a} & theorem & no per-edge abelian strong phase & Part 2 & Corollary 2.4 \\
\multicolumn{5}{@{}l@{}}{\hspace*{1em}\small\itshape §F.III — Colour Objects and Spectrum}\\[1pt]
\texttt{CT3{.}11} & theorem & Colour objects & Part 2 & Theorem 2.109 \\
\texttt{CT3{.}12} & theorem & Photon and gluon are distinguished by possibility space & Part 2 & Theorem 2.110 \\
\texttt{CT3{.}13} & theorem & Gluon count & Part 2 & Theorem 2.111 \\
\texttt{CT3{.}14} & theorem & Gluons remain massless & Part 2 & Theorem 2.112 \\
\texttt{CT3{.}15} & theorem & Baryon invariance & Part 2 & Theorem 2.113 \\
\texttt{CT3{.}16} & theorem & Infinitesimal baryon invariance & Part 2 & Theorem 2.114 \\
\texttt{CT3{.}17} & theorem & Invariant tensors & Part 2 & Theorem 2.115 \\
\texttt{CT3{.}17a} & theorem & sealing partition rule; grammar counts & Part 2 & Corollary 2.5 \\
\texttt{CT3{.}18} & theorem & Colour-singlet closure & Part 2 & Theorem 2.116 \\
\multicolumn{5}{@{}l@{}}{\hspace*{1em}\small\itshape §F.IV — Completion}\\[1pt]
\texttt{CT3{.}19} & theorem & Self-containment of the colour sector & Part 2 & Theorem 2.119 \\
\multicolumn{5}{@{}l@{}}{\hspace*{1em}\small\itshape §F.I — Slot Structure}\\[1pt]
\texttt{CT3{.}2} & theorem & Quantum slot space & Part 2 & Theorem 2.99 \\
\texttt{CT3{.}20} & theorem & The number of colours is forced & Part 2 & Theorem 2.97 \\
\multicolumn{5}{@{}l@{}}{\hspace*{1em}\small\itshape §F.II — The Rung-3 Gauge Algebra}\\[1pt]
\texttt{CT3{.}21} & theorem & Structure constants & Part 2 & Theorem 2.107 \\
\multicolumn{5}{@{}l@{}}{\hspace*{1em}\small\itshape §F.III — Colour Objects and Spectrum}\\[1pt]
\texttt{CT3{.}22} & theorem & Sealing is confinement & Part 2 & Theorem 2.117 \\
\texttt{CT3{.}22a} & theorem & Subnetwork confinement: the domain restriction, all-or-nothing, and the converse & Part 2 & Theorem 2.118 \\
\multicolumn{5}{@{}l@{}}{\hspace*{1em}\small\itshape §F.I — Slot Structure}\\[1pt]
\texttt{CT3{.}3} & theorem & Classical permutations embed unitarily & Part 2 & Theorem 2.100 \\
\texttt{CT3{.}4} & theorem & Probability preservation gives & Part 2 & Theorem 2.101 \\
\texttt{CT3{.}5} & theorem & The algebra and its decomposition & Part 2 & Theorem 2.102 \\
\multicolumn{5}{@{}l@{}}{\hspace*{1em}\small\itshape §F.II — The Rung-3 Gauge Algebra}\\[1pt]
\texttt{CT3{.}6} & theorem & Rung-3 active set & Part 2 & Theorem 2.103 \\
\texttt{CT3{.}7} & theorem & The strictly new rung-3 algebra is & Part 2 & Theorem 2.104 \\
\texttt{CT3{.}8} & theorem & Dimension of the colour algebra & Part 2 & Theorem 2.105 \\
\texttt{CT3{.}9} & theorem & Explicit generator basis & Part 2 & Theorem 2.106 \\
\multicolumn{5}{@{}l@{}}{\hspace*{1em}\small\itshape §G.I — The Saturation Lock}\\[1pt]
\texttt{CT4{.}1} & theorem & Oversubscription falls to & Part 2 & Theorem 2.120 \\
\multicolumn{5}{@{}l@{}}{\hspace*{1em}\small\itshape §G.II — The Two Dark Aspects}\\[1pt]
\texttt{CT4{.}10} & theorem & The triangle is exhaustive & Part 2 & Theorem 2.129 \\
\multicolumn{5}{@{}l@{}}{\hspace*{1em}\small\itshape §G.III — Bounds and Boundaries}\\[1pt]
\texttt{CT4{.}11} & theorem & Global closure & Part 2 & Theorem 2.132 \\
\multicolumn{5}{@{}l@{}}{\hspace*{1em}\small\itshape §G.IV — Falsifiability and the Closure Audit}\\[1pt]
\texttt{CT4{.}12} & theorem & Falsifiers & Part 2 & Theorem 2.133 \\
\texttt{CT4{.}13} & theorem & Closure audit & Part 2 & Theorem 2.134 \\
\multicolumn{5}{@{}l@{}}{\hspace*{1em}\small\itshape §G.I — The Saturation Lock}\\[1pt]
\texttt{CT4{.}1a} & theorem & universal saturation of bidirectional real-edge subgraphs & Part 2 & Corollary 2.6 \\
\texttt{CT4{.}2} & theorem & The lock is absorbing & Part 2 & Theorem 2.121 \\
\texttt{CT4{.}3} & theorem & Locked, not destroyed & Part 2 & Theorem 2.122 \\
\texttt{CT4{.}4} & theorem & The locked sector carries no gauge structure & Part 2 & Theorem 2.123 \\
\texttt{CT4{.}4a} & theorem & --- & Addendum & B.3.1 \\
\multicolumn{5}{@{}l@{}}{\hspace*{1em}\small\itshape §G.II — The Two Dark Aspects}\\[1pt]
\texttt{CT4{.}5} & theorem & Equation of state of the locked sector & Part 2 & Theorem 2.124 \\
\texttt{CT4{.}5a} & theorem & The flicker-dust ensemble: exact density polynomial and rate-invariant & Part 2 & Theorem 2.125 \\
\texttt{CT4{.}6} & theorem & The two dark aspects contrasted & Part 2 & Theorem 2.126 \\
\texttt{CT4{.}7} & theorem & The entropy triangle & Part 2 & Theorem 2.127 \\
\texttt{CT4{.}7a} & theorem & The vacuum census, closed, and the baseline equality & Part 2 & Theorem 2.128 \\
\multicolumn{5}{@{}l@{}}{\hspace*{1em}\small\itshape §G.III — Bounds and Boundaries}\\[1pt]
\texttt{CT4{.}8} & theorem & Locked regions and their boundaries & Part 2 & Theorem 2.131 \\
\texttt{CT4{.}9} & theorem & Accessible entropy bounds & Part 2 & Theorem 2.130 \\
\multicolumn{5}{@{}l@{}}{\hspace*{1em}\small\itshape §E.V — The Dynamic Phase}\\[1pt]
\texttt{Def-EM{.}20} & definition & The per-edge phase at crystallization & Part 2 & Definition 2.1 \\
\multicolumn{5}{@{}l@{}}{\hspace*{1em}\small\itshape §E.VI — Parity and CPT}\\[1pt]
\texttt{Def-EM{.}21} & definition & Alignment pattern & Part 2 & Definition 2.2 \\
\texttt{Def-EM{.}22} & definition & Mirror operation & Part 2 & Definition 2.3 \\
\multicolumn{5}{@{}l@{}}{\hspace*{1em}\small\itshape §E.I — Rung-2 Existence and Algebra}\\[1pt]
\texttt{EM{.}1} & theorem & The electromagnetic aspect activates at rung 2 & Part 2 & Theorem 2.70 \\
\multicolumn{5}{@{}l@{}}{\hspace*{1em}\small\itshape §E.III — Electromagnetic Charge}\\[1pt]
\texttt{EM{.}10} & theorem & Charge quantization & Part 2 & Theorem 2.80 \\
\texttt{EM{.}11} & theorem & Global charge conservation & Part 2 & Theorem 2.81 \\
\multicolumn{5}{@{}l@{}}{\hspace*{1em}\small\itshape §E.IV — Field Strength and Dynamics}\\[1pt]
\texttt{EM{.}12} & theorem & Discrete field strength from phase transitions & Part 2 & Theorem 2.82 \\
\texttt{EM{.}13} & theorem & The transition-parity factor in the weight & Part 2 & Theorem 2.83 \\
\texttt{EM{.}14} & theorem & Continuum gauge invariance & Part 2 & Theorem 2.84 \\
\texttt{EM{.}15} & theorem & Electromagnetic stress-energy & Part 2 & Theorem 2.85 \\
\multicolumn{5}{@{}l@{}}{\hspace*{1em}\small\itshape §E.V — The Dynamic Phase}\\[1pt]
\texttt{EM{.}16} & theorem & Edge transfer weight and the per-edge read-back & Part 2 & Theorem 2.86 \\
\texttt{EM{.}17} & theorem & Spectral read of the electromagnetic block & Part 2 & Theorem 2.87 \\
\texttt{EM{.}18} & theorem & The colour sector does not shift the block read & Part 2 & Theorem 2.88 \\
\texttt{EM{.}19} & theorem & The saturation sector does not shift the block read & Part 2 & Theorem 2.89 \\
\multicolumn{5}{@{}l@{}}{\hspace*{1em}\small\itshape §E.I — Rung-2 Existence and Algebra}\\[1pt]
\texttt{EM{.}2} & theorem & The electromagnetic generator & Part 2 & Theorem 2.71 \\
\multicolumn{5}{@{}l@{}}{\hspace*{1em}\small\itshape §E.V — The Dynamic Phase}\\[1pt]
\texttt{EM{.}20} & theorem & The dynamic edge phase & Part 2 & Theorem 2.90 \\
\texttt{EM{.}21} & theorem & Dimensional status of the per-edge phase & Part 2 & Theorem 2.91 \\
\multicolumn{5}{@{}l@{}}{\hspace*{1em}\small\itshape §E.VII — Spectrum and Completion}\\[1pt]
\texttt{EM{.}22} & theorem & The electromagnetic spectrum & Part 2 & Theorem 2.94 \\
\multicolumn{5}{@{}l@{}}{\hspace*{1em}\small\itshape §E.VI — Parity and CPT}\\[1pt]
\texttt{EM{.}23} & theorem & Parity classification & Part 2 & Theorem 2.92 \\
\texttt{EM{.}24} & theorem & The CPT bridge & Part 2 & Theorem 2.93 \\
\multicolumn{5}{@{}l@{}}{\hspace*{1em}\small\itshape §E.VII — Spectrum and Completion}\\[1pt]
\texttt{EM{.}25} & theorem & What the corpus does and does not fix about the abelian sector & Part 2 & Theorem 2.96 \\
\texttt{EM{.}26} & theorem & Cross-checks & Part 2 & Theorem 2.95 \\
\multicolumn{5}{@{}l@{}}{\hspace*{1em}\small\itshape §E.I — Rung-2 Existence and Algebra}\\[1pt]
\texttt{EM{.}3} & theorem & The electromagnetic Lie algebra & Part 2 & Theorem 2.72 \\
\texttt{EM{.}4} & theorem & The rung-2 algebra containing the abelian factor & Part 2 & Theorem 2.73 \\
\multicolumn{5}{@{}l@{}}{\hspace*{1em}\small\itshape §E.II — The Photon}\\[1pt]
\texttt{EM{.}5} & theorem & Photon and gluon as trivial chains & Part 2 & Theorem 2.74 \\
\texttt{EM{.}6} & theorem & Photon masslessness & Part 2 & Theorem 2.75 \\
\texttt{EM{.}7} & theorem & Photon entropy budget & Part 2 & Theorem 2.76 \\
\texttt{EM{.}8} & theorem & Photon polarization count & Part 2 & Theorem 2.77 \\
\multicolumn{5}{@{}l@{}}{\hspace*{1em}\small\itshape §E.III — Electromagnetic Charge}\\[1pt]
\texttt{EM{.}9} & theorem & Chain end-charge & Part 2 & Theorem 2.78 \\
\texttt{EM{.}9a} & theorem & The closed-core charge read and the end-register count & Part 2 & Theorem 2.79 \\
\texttt{EM{.}9b} & theorem & chain-wall telescoping & Part 2 & Corollary 2.3 \\
\multicolumn{5}{@{}l@{}}{\hspace*{1em}\small\itshape §D.I — Operator Preliminaries}\\[1pt]
\texttt{EW{.}1} & theorem & Joint sector decomposition & Part 2 & Theorem 2.33 \\
\multicolumn{5}{@{}l@{}}{\hspace*{1em}\small\itshape §D.II — Rung Masks and the Electroweak Algebra}\\[1pt]
\texttt{EW{.}10} & theorem & The rung-2 electroweak algebra & Part 2 & Theorem 2.43 \\
\texttt{EW{.}11} & theorem & The mechanism is projection, not path integration & Part 2 & Theorem 2.44 \\
\multicolumn{5}{@{}l@{}}{\hspace*{1em}\small\itshape §D.III — The Scalar Trace Mode}\\[1pt]
\texttt{EW{.}12} & theorem & Metric expansion and the scalar trace mode & Part 2 & Theorem 2.45 \\
\texttt{EW{.}13} & theorem & Identification of the scalar mode & Part 2 & Theorem 2.46 \\
\texttt{EW{.}14} & theorem & The scalar background commutes with the electroweak generators & Part 2 & Theorem 2.47 \\
\texttt{EW{.}15} & theorem & No adjoint electroweak breaking & Part 2 & Theorem 2.48 \\
\texttt{EW{.}16} & theorem & Scalar coupling to the antisymmetric sectors & Part 2 & Theorem 2.49 \\
\multicolumn{5}{@{}l@{}}{\hspace*{1em}\small\itshape §D.IV — Rung-3 Wrapping, the Seed, and Mass Generation}\\[1pt]
\texttt{EW{.}17} & theorem & Rung-3 closed wrapping chains & Part 2 & Theorem 2.50 \\
\texttt{EW{.}18} & theorem & Spiral pole condition & Part 2 & Theorem 2.51 \\
\texttt{EW{.}19} & theorem & Mass as pole displacement & Part 2 & Theorem 2.52 \\
\multicolumn{5}{@{}l@{}}{\hspace*{1em}\small\itshape §D.I — Operator Preliminaries}\\[1pt]
\texttt{EW{.}2} & theorem & Generator basis and sector assignment & Part 2 & Theorem 2.34 \\
\multicolumn{5}{@{}l@{}}{\hspace*{1em}\small\itshape §D.IV — Rung-3 Wrapping, the Seed, and Mass Generation}\\[1pt]
\texttt{EW{.}20} & theorem & Weak mass generation & Part 2 & Theorem 2.55 \\
\texttt{EW{.}21} & theorem & Charged weak bosons & Part 2 & Theorem 2.56 \\
\multicolumn{5}{@{}l@{}}{\hspace*{1em}\small\itshape §D.V — The Neutral Sector}\\[1pt]
\texttt{EW{.}22} & theorem & The unbroken generator & Part 2 & Theorem 2.57 \\
\texttt{EW{.}23} & theorem & The neutral mass matrix & Part 2 & Theorem 2.58 \\
\texttt{EW{.}24} & theorem & The mixing angle & Part 2 & Theorem 2.59 \\
\texttt{EW{.}25} & theorem & The photon direction is massless & Part 2 & Theorem 2.60 \\
\texttt{EW{.}26} & theorem & Colour modes remain massless at rung 3 & Part 2 & Theorem 2.61 \\
\texttt{EW{.}27} & theorem & The scalar excitation & Part 2 & Theorem 2.62 \\
\multicolumn{5}{@{}l@{}}{\hspace*{1em}\small\itshape §D.VI — Derived Electroweak Relations}\\[1pt]
\texttt{EW{.}28} & theorem & The gauge spectrum & Part 2 & Theorem 2.63 \\
\texttt{EW{.}29} & theorem & The full electroweak spectrum & Part 2 & Theorem 2.67 \\
\multicolumn{5}{@{}l@{}}{\hspace*{1em}\small\itshape §D.I — Operator Preliminaries}\\[1pt]
\texttt{EW{.}3} & theorem & The close & Part 2 & Theorem 2.35 \\
\multicolumn{5}{@{}l@{}}{\hspace*{1em}\small\itshape §D.VI — Derived Electroweak Relations}\\[1pt]
\texttt{EW{.}30} & theorem & Spectral reduction & Part 2 & Theorem 2.68 \\
\multicolumn{5}{@{}l@{}}{\hspace*{1em}\small\itshape §D.VII — Completion}\\[1pt]
\texttt{EW{.}31} & theorem & Self-containment of the electroweak sector & Part 2 & Theorem 2.69 \\
\multicolumn{5}{@{}l@{}}{\hspace*{1em}\small\itshape §D.IV — Rung-3 Wrapping, the Seed, and Mass Generation}\\[1pt]
\texttt{EW{.}32} & theorem & The wrapping seed: carrier and normalization & Part 2 & Theorem 2.53 \\
\texttt{EW{.}33} & theorem & The wrapping mass form & Part 2 & Theorem 2.54 \\
\multicolumn{5}{@{}l@{}}{\hspace*{1em}\small\itshape §D.VI — Derived Electroweak Relations}\\[1pt]
\texttt{EW{.}34} & theorem & The mass ratio & Part 2 & Theorem 2.64 \\
\texttt{EW{.}35} & theorem & The residual coupling & Part 2 & Theorem 2.65 \\
\texttt{EW{.}36} & theorem & The charge spectrum and its consistency with the chain end-charge & Part 2 & Theorem 2.66 \\
\texttt{EW{.}36a} & theorem & corollary & Part 2 & Corollary 2.2 \\
\multicolumn{5}{@{}l@{}}{\hspace*{1em}\small\itshape §D.I — Operator Preliminaries}\\[1pt]
\texttt{EW{.}4} & theorem & Central, trace, and mixed commutators & Part 2 & Theorem 2.36 \\
\texttt{EW{.}5} & theorem & Second-order Higgs-to-weak pathways & Part 2 & Theorem 2.37 \\
\texttt{EW{.}5a} & theorem & Convention exactness and the rigidity of the named basis & Part 2 & Theorem 2.38 \\
\multicolumn{5}{@{}l@{}}{\hspace*{1em}\small\itshape §D.II — Rung Masks and the Electroweak Algebra}\\[1pt]
\texttt{EW{.}6} & theorem & Rung masks and active aspect sets & Part 2 & Theorem 2.39 \\
\texttt{EW{.}7} & theorem & The rung-2 electroweak mask & Part 2 & Theorem 2.40 \\
\texttt{EW{.}8} & theorem & The electromagnetic generator & Part 2 & Theorem 2.41 \\
\texttt{EW{.}9} & theorem & The abelian algebra & Part 2 & Theorem 2.42 \\
\addlinespace[4pt]
\multicolumn{5}{@{}l}{\textbf{Part 3 · Fermion\_era — bound chains, termini, floors, carriers, channels}}\\[2pt]
\multicolumn{5}{@{}l@{}}{\hspace*{1em}\small\itshape Tier F.1 — Information Flux and Measurement Primitivity}\\[1pt]
\texttt{Def-H{.}1} & definition & Observer–subnetwork channel & Part 3 & Definition 3.1 \\
\multicolumn{5}{@{}l@{}}{\hspace*{1em}\small\itshape Tier F.6 — Subnetwork-Level Force Structure: Confinement and the Active Aspect Set}\\[1pt]
\texttt{Def-H{.}10} & definition & Colour boundary and the contraction-completeness projector & Part 3 & Definition 3.13 \\
\texttt{Def-H{.}11} & definition & Active aspect set of a subnetwork & Part 3 & Definition 3.14 \\
\multicolumn{5}{@{}l@{}}{\hspace*{1em}\small\itshape Tier F.7 — The Coarse-Grained Slot Algebra: \$\textbackslash{}mathcal A(\textbackslash{}Sigma)\$ Properly Derived}\\[1pt]
\texttt{Def-H{.}12} & definition & Coarse-grained slot algebra & Part 3 & Definition 3.16 \\
\texttt{Def-H{.}13} & definition & Compositional bridge map & Part 3 & Definition 3.17 \\
\multicolumn{5}{@{}l@{}}{\hspace*{1em}\small\itshape Tier F.8 — Line-Element Quantization from Materialization}\\[1pt]
\texttt{Def-H{.}14} & definition & General-edge sector space (internal-edge slot object) & Part 3 & Definition 3.18 \\
\multicolumn{5}{@{}l@{}}{\hspace*{1em}\small\itshape Tier F.9 — The Activated Aspect Operator: \$\textbackslash{}mathcal A(\textbackslash{}Sigma)\$ as a Computable Mask}\\[1pt]
\texttt{Def-H{.}15} & definition & Coarse-grained aspect space and masked force operator & Part 3 & Definition 3.19 \\
\texttt{Def-H{.}16} & definition & Coarse-grained resolvent and the -spectrum & Part 3 & Definition 3.20 \\
\multicolumn{5}{@{}l@{}}{\hspace*{1em}\small\itshape Tier F.10 — The Strong Edge Weight}\\[1pt]
\texttt{Def-H{.}17} & definition & Strong edge transport & Part 3 & Definition 3.21 \\
\multicolumn{5}{@{}l@{}}{\hspace*{1em}\small\itshape Tier F.11 — Sector Quantum Audit and the Continuum Ledger}\\[1pt]
\texttt{Def-H{.}18} & definition & Arena and spectrum (the two-locus vocabulary) & Part 3 & Definition 3.22 \\
\multicolumn{5}{@{}l@{}}{\hspace*{1em}\small\itshape Tier F.12 — The Cycle Floor Family, the Defect Bridge, and Composite Floors}\\[1pt]
\texttt{Def-H{.}19} & definition & Stationary, stable, and unstable jumps (vocabulary) & Part 3 & Definition 3.23 \\
\multicolumn{5}{@{}l@{}}{\hspace*{1em}\small\itshape Tier F.1 — Information Flux and Measurement Primitivity}\\[1pt]
\texttt{Def-H{.}2} & definition & Physical (metric) channel separation & Part 3 & Definition 3.2 \\
\multicolumn{5}{@{}l@{}}{\hspace*{1em}\small\itshape Tier F.14 — Platonic Floor Cards: The Five Solids as Fermion Floors}\\[1pt]
\texttt{Def-H{.}20} & definition & Platonic floors (setup vocabulary) & Part 3 & Definition 3.24 \\
\multicolumn{5}{@{}l@{}}{\hspace*{1em}\small\itshape Tier F.17 — The Stationarity Loop: Constancy, the Spin Connection, and the No-Go}\\[1pt]
\texttt{Def-H{.}21} & definition & The spin connection on the emergent geometry (construction) & Part 3 & Definition 3.25 \\
\multicolumn{5}{@{}l@{}}{\hspace*{1em}\small\itshape Tier F.18 — Native Rungs, Boundary Exchange Shape, and the Admissibility Correction}\\[1pt]
\texttt{Def-H{.}22} & definition & The explicit boundary realization (corolla and separated patterns) & Part 3 & Definition 3.27 \\
\multicolumn{5}{@{}l@{}}{\hspace*{1em}\small\itshape Tier F.19 — The Three Tails: Realization Non-Existence, the Multiplicity Map, and the Selection Obstruction}\\[1pt]
\texttt{Def-H{.}23} & definition & The winding-multiplicity carrier map (under Def-H.38) & Part 3 & Definition 3.29 \\
\multicolumn{5}{@{}l@{}}{\hspace*{1em}\small\itshape Tier F.22 — The Bott Correction: Degenerate Witnesses Re-Evaluated, Reconstructability, and the Selection Dilemma}\\[1pt]
\texttt{Def-H{.}24} & definition & Bott index, clean critical manifolds, and monodromy type for ledger channels (construction) & Part 3 & Definition 3.30 \\
\multicolumn{5}{@{}l@{}}{\hspace*{1em}\small\itshape Tier F.23 — The Virtual Bypass: Path Integration over Crystallization Patterns}\\[1pt]
\texttt{Def-H{.}25} & definition & Skeleton configuration sums and the bypass vocabulary & Part 3 & Definition 3.31 \\
\multicolumn{5}{@{}l@{}}{\hspace*{1em}\small\itshape Tier F.24 — The Flicker Phase Structure, the Cycle Grammar of Covers, the Dust Equation of State, Face Decomposition, and the Bridge Audit}\\[1pt]
\texttt{Def-H{.}26} & definition & Flicker-ensemble phase reads & Part 3 & Definition 3.32 \\
\multicolumn{5}{@{}l@{}}{\hspace*{1em}\small\itshape Tier F.25 — The Chain-Local \$\textbackslash{}aleph\$-Domains: Extended Exact-Cover Theorems and the Domain-Relativity of the Amplitude Layer}\\[1pt]
\texttt{Def-H{.}27} & definition & The -domains: literal, chain-dominant, literal-forced & Part 3 & Definition 3.33 \\
\multicolumn{5}{@{}l@{}}{\hspace*{1em}\small\itshape Tier F.27 — The Cover Orbit Decomposition: the Ansatz Split, the Combinatorial Ceiling, and the Bridge Census}\\[1pt]
\texttt{Def-H{.}28} & definition & Cover classes, and the two class-mass ansätze & Part 3 & Definition 3.34 \\
\multicolumn{5}{@{}l@{}}{\hspace*{1em}\small\itshape Tier F.28 — The Full-SM Completeness Audit: Particle Content, the Gauge Blocks, the Rung Ledger, and the Deficit Table}\\[1pt]
\texttt{Def-H{.}29} & definition & Completeness-audit vocabulary: read classes, census statuses, the external-data clause & Part 3 & Definition 3.35 \\
\multicolumn{5}{@{}l@{}}{\hspace*{1em}\small\itshape Tier F.1 — Information Flux and Measurement Primitivity}\\[1pt]
\texttt{Def-H{.}3} & definition & Extractable information & Part 3 & Definition 3.3 \\
\multicolumn{5}{@{}l@{}}{\hspace*{1em}\small\itshape Tier F.29 — Gravitational Reads in Three States, the Flicker Ripple, and the Tumbling-Model Adjudication}\\[1pt]
\texttt{Def-H{.}30} & definition & The three shape states & Part 3 & Definition 3.36 \\
\multicolumn{5}{@{}l@{}}{\hspace*{1em}\small\itshape Tier F.35 — The 1-Loop Bridge: the Swap-Graded Effective Potential on the Flicker Phase Diagram}\\[1pt]
\texttt{Def-H{.}31} & definition & The bridge: dictionary, family, background, fluctuation operator, effective potential & Part 3 & Definition 3.37 \\
\multicolumn{5}{@{}l@{}}{\hspace*{1em}\small\itshape Tier F.36 — The Topology-Blind Path Integral: Observable-Conditioned Sectors and the Janus Sum}\\[1pt]
\texttt{Def-H{.}32} & definition & The topology-blind path integral & Part 3 & Definition 3.38 \\
\multicolumn{5}{@{}l@{}}{\hspace*{1em}\small\itshape Tier F.39 — Chirality: The Orientation Doublet of the Bound Chain and Its Exchange-Form Grade}\\[1pt]
\texttt{Def-H{.}33} & definition & The chain orientation space and its chirality operator & Part 3 & Definition 3.41 \\
\multicolumn{5}{@{}l@{}}{\hspace*{1em}\small\itshape Tier F.40 — The Static Force Matrix on the Extended Carrier: The Aspect Census, the No-Flip Theorem, and the Charge–Parity Splitting}\\[1pt]
\texttt{Def-H{.}34} & definition & The extended carrier basis, the chiral transform, and the partner involution & Part 3 & Definition 3.42 \\
\multicolumn{5}{@{}l@{}}{\hspace*{1em}\small\itshape Tier F.5 — Spin–Charge Decomposition via Complete Force Operator}\\[1pt]
\texttt{Def-H{.}35} & definition & Channel-visible content (the boundary-transport read) & Part 3 & Definition 3.10 \\
\multicolumn{5}{@{}l@{}}{\hspace*{1em}\small\itshape Tier F.6 — Subnetwork-Level Force Structure: Confinement and the Active Aspect Set}\\[1pt]
\texttt{Def-H{.}36} & definition & The scalar-row fold (the H read at -level) & Part 3 & Definition 3.15 \\
\multicolumn{5}{@{}l@{}}{\hspace*{1em}\small\itshape Tier F.17 — The Stationarity Loop: Constancy, the Spin Connection, and the No-Go}\\[1pt]
\texttt{Def-H{.}37} & definition & The defect-transfer interface & Part 3 & Definition 3.26 \\
\multicolumn{5}{@{}l@{}}{\hspace*{1em}\small\itshape Tier F.19 — The Three Tails: Realization Non-Existence, the Multiplicity Map, and the Selection Obstruction}\\[1pt]
\texttt{Def-H{.}38} & definition & The multiplicity–winding coupling & Part 3 & Definition 3.28 \\
\multicolumn{5}{@{}l@{}}{\hspace*{1em}\small\itshape Tier F.5b — The Mass-Read Registry and the Tilt Composition}\\[1pt]
\texttt{Def-H{.}39} & definition & The mass-read registry; bridge slots & Part 3 & Definition 3.11 \\
\multicolumn{5}{@{}l@{}}{\hspace*{1em}\small\itshape Tier F.2 — The Fermion Definition}\\[1pt]
\texttt{Def-H{.}4} & definition & Fermion & Part 3 & Definition 3.4 \\
\multicolumn{5}{@{}l@{}}{\hspace*{1em}\small\itshape Tier F.5b — The Mass-Read Registry and the Tilt Composition}\\[1pt]
\texttt{Def-H{.}40} & definition & Chain-level tilt composition & Part 3 & Definition 3.12 \\
\multicolumn{5}{@{}l@{}}{\hspace*{1em}\small\itshape Tier F.37 — The Koide Channel: The \$C\_3\$-Covariant Seal Bilinear and the Bound-Chain Mixing Angle}\\[1pt]
\texttt{Def-H{.}41} & definition & The bound-chain arena & Part 3 & Definition 3.39 \\
\multicolumn{5}{@{}l@{}}{\hspace*{1em}\small\itshape Tier F.38 — The Colour Channel: The \$\textbackslash{}varepsilon\$-Selected \$C\_3\$, Confinement Blindness, and the Triplet-Attached Mixing Angle}\\[1pt]
\texttt{Def-H{.}42} & definition & The colour-dressed arena & Part 3 & Definition 3.40 \\
\multicolumn{5}{@{}l@{}}{\hspace*{1em}\small\itshape Tier F.40 — The Static Force Matrix on the Extended Carrier: The Aspect Census, the No-Flip Theorem, and the Charge–Parity Splitting}\\[1pt]
\texttt{Def-H{.}43} & definition & The read promotion on the orientation doublet & Part 3 & Definition 3.43 \\
\multicolumn{5}{@{}l@{}}{\hspace*{1em}\small\itshape Tier F.3 — Information Flux as Rate}\\[1pt]
\texttt{Def-H{.}5} & definition & Information flux & Part 3 & Definition 3.5 \\
\texttt{Def-H{.}6} & definition & Information flux density & Part 3 & Definition 3.6 \\
\multicolumn{5}{@{}l@{}}{\hspace*{1em}\small\itshape Tier F.4 — Rung-Dependent Edge Structure}\\[1pt]
\texttt{Def-H{.}7} & definition & Rung-dependent edge data & Part 3 & Definition 3.7 \\
\multicolumn{5}{@{}l@{}}{\hspace*{1em}\small\itshape Tier F.5 — Spin–Charge Decomposition via Complete Force Operator}\\[1pt]
\texttt{Def-H{.}8} & definition & Spin (emergent weak contribution) & Part 3 & Definition 3.8 \\
\texttt{Def-H{.}9} & definition & Subnetwork coarse-graining projection & Part 3 & Definition 3.9 \\
\multicolumn{5}{@{}l@{}}{\hspace*{1em}\small\itshape Tier F.1 — Information Flux and Measurement Primitivity}\\[1pt]
\texttt{FT{.}1} & theorem & Information flux bound & Part 3 & Theorem 3.1 \\
\multicolumn{5}{@{}l@{}}{\hspace*{1em}\small\itshape Tier F.4 — Rung-Dependent Edge Structure}\\[1pt]
\texttt{FT{.}10} & theorem & Confinement is a global sum-restriction & Part 3 & Theorem 3.10 \\
\multicolumn{5}{@{}l@{}}{\hspace*{1em}\small\itshape Tier F.33 — The Dynamic Phase, Parity, and the CPT Bridge: the Mass Refinement}\\[1pt]
\texttt{FT{.}100} & theorem & The alignment pattern refines the informational mass & Part 3 & Theorem 3.74 \\
\multicolumn{5}{@{}l@{}}{\hspace*{1em}\small\itshape Tier F.34 — The Weak Branch Grading of the Two-Branch Crystallization}\\[1pt]
\texttt{FT{.}101} & theorem & The weak branch grading: the two crystallization branches are weak-inequivalent & Part 3 & Theorem 3.75 \\
\multicolumn{5}{@{}l@{}}{\hspace*{1em}\small\itshape Tier F.35 — The 1-Loop Bridge: the Swap-Graded Effective Potential on the Flicker Phase Diagram}\\[1pt]
\texttt{FT{.}102} & theorem & The fluctuation spectrum, the tree-level limit, and the seal-softening walls & Part 3 & Theorem 3.76 \\
\texttt{FT{.}103} & theorem & The vacuum structure: four exact vacua, three wall-delimited phase bands, dominance windows, the susceptibility census & Part 3 & Theorem 3.77 \\
\texttt{FT{.}104} & theorem & The [EXT] audit, the range-species classification, and the fourth read of Gap 7 & Part 3 & Theorem 3.78 \\
\multicolumn{5}{@{}l@{}}{\hspace*{1em}\small\itshape Tier F.36 — The Topology-Blind Path Integral: Observable-Conditioned Sectors and the Janus Sum}\\[1pt]
\texttt{FT{.}106} & theorem & Weight transport and the two-layer sector structure & Part 3 & Theorem 3.79 \\
\texttt{FT{.}107} & theorem & The observable sectors and the mass expectation & Part 3 & Theorem 3.80 \\
\texttt{FT{.}108} & theorem & The Janus sum: move algebra, budget conservation, derived detailed balance, and the Tier F.35 boundary & Part 3 & Theorem 3.81 \\
\multicolumn{5}{@{}l@{}}{\hspace*{1em}\small\itshape Tier F.37 — The Koide Channel: The \$C\_3\$-Covariant Seal Bilinear and the Bound-Chain Mixing Angle}\\[1pt]
\texttt{FT{.}109} & theorem & The diagonal seal commutator & Part 3 & Theorem 3.82 \\
\multicolumn{5}{@{}l@{}}{\hspace*{1em}\small\itshape Tier F.5 — Spin–Charge Decomposition via Complete Force Operator}\\[1pt]
\texttt{FT{.}11} & theorem & Complete force operator charge audit & Part 3 & Theorem 3.11 \\
\multicolumn{5}{@{}l@{}}{\hspace*{1em}\small\itshape Tier F.37 — The Koide Channel: The \$C\_3\$-Covariant Seal Bilinear and the Bound-Chain Mixing Angle}\\[1pt]
\texttt{FT{.}110} & theorem & The bound-chain arena dimension & Part 3 & Theorem 3.83 \\
\texttt{FT{.}111} & theorem & The spectrum of the covariant seal bilinear & Part 3 & Theorem 3.84 \\
\texttt{FT{.}112} & theorem & The mixing angle & Part 3 & Theorem 3.85 \\
\texttt{FT{.}113} & theorem & The bilinear spectrum and the mass-squared ratios & Part 3 & Theorem 3.86 \\
\multicolumn{5}{@{}l@{}}{\hspace*{1em}\small\itshape Tier F.38 — The Colour Channel: The \$\textbackslash{}varepsilon\$-Selected \$C\_3\$, Confinement Blindness, and the Triplet-Attached Mixing Angle}\\[1pt]
\texttt{FT{.}116} & theorem & The relabelling lattice, the two-level obstruction, and the -selection of & Part 3 & Theorem 3.87 \\
\texttt{FT{.}117} & theorem & --- & Addendum & B.3.3 \\
\texttt{FT{.}118} & theorem & The colour-dressed arena and the spectator numerator & Part 3 & Theorem 3.88 \\
\texttt{FT{.}119} & theorem & The triplet-attached mixing angle and its spectrum & Part 3 & Theorem 3.89 \\
\multicolumn{5}{@{}l@{}}{\hspace*{1em}\small\itshape Tier F.5 — Spin–Charge Decomposition via Complete Force Operator}\\[1pt]
\texttt{FT{.}12} & theorem & Weak has no self-charge & Part 3 & Theorem 3.12 \\
\multicolumn{5}{@{}l@{}}{\hspace*{1em}\small\itshape Tier F.39 — Chirality: The Orientation Doublet of the Bound Chain and Its Exchange-Form Grade}\\[1pt]
\texttt{FT{.}122} & theorem & The orientation doublet is corpus-constructible & Part 3 & Theorem 3.90 \\
\texttt{FT{.}123} & theorem & Weak covariance and -covariance on the extended carrier & Part 3 & Theorem 3.91 \\
\texttt{FT{.}124} & theorem & --- & Addendum & B.3.4 \\
\texttt{FT{.}125} & theorem & --- & Addendum & B.4.1 \\
\texttt{FT{.}126} & theorem & Weak-covariant chiral projectors on the chain carrier & Part 3 & Theorem 3.92 \\
\multicolumn{5}{@{}l@{}}{\hspace*{1em}\small\itshape Tier F.40 — The Static Force Matrix on the Extended Carrier: The Aspect Census, the No-Flip Theorem, and the Charge–Parity Splitting}\\[1pt]
\texttt{FT{.}127} & theorem & The aspect census on the extended carrier & Part 3 & Theorem 3.93 \\
\texttt{FT{.}128} & theorem & The no-flip theorem at the derived locus & Part 3 & Theorem 3.94 \\
\texttt{FT{.}129} & theorem & The canonical flip candidate and the read promotion & Part 3 & Theorem 3.95 \\
\multicolumn{5}{@{}l@{}}{\hspace*{1em}\small\itshape Tier F.5 — Spin–Charge Decomposition via Complete Force Operator}\\[1pt]
\texttt{FT{.}13} & theorem & Weak spin (algebraically) & Part 3 & Theorem 3.13 \\
\multicolumn{5}{@{}l@{}}{\hspace*{1em}\small\itshape Tier F.40 — The Static Force Matrix on the Extended Carrier: The Aspect Census, the No-Flip Theorem, and the Charge–Parity Splitting}\\[1pt]
\texttt{FT{.}130} & theorem & Tensor structure, Schmidt rank, and the static spectrum & Part 3 & Theorem 3.96 \\
\texttt{FT{.}131} & theorem & Scope: bound chains, no isolated charge, the deferrals & Part 3 & Theorem 3.97 \\
\multicolumn{5}{@{}l@{}}{\hspace*{1em}\small\itshape Tier F.41 — Parity: The Configuration Calculus, the Carrier Bookkeeping, and the Sector Census}\\[1pt]
\texttt{FT{.}133} & theorem & The configuration-space reversal calculus & Part 3 & Theorem 3.98 \\
\texttt{FT{.}134} & theorem & The carrier bookkeeping: the scalar-CPT theorem & Part 3 & Theorem 3.99 \\
\texttt{FT{.}135} & theorem & The weak branch grading under parity & Part 3 & Theorem 3.100 \\
\texttt{FT{.}136} & theorem & The rung-one S-grading: the 10/6 split & Part 3 & Theorem 3.101 \\
\texttt{FT{.}137} & theorem & The cover and flicker census under the reversal family & Part 3 & Theorem 3.102 \\
\texttt{FT{.}138} & theorem & The pole spectrum census & Part 3 & Theorem 3.103 \\
\texttt{FT{.}139} & theorem & The channel census: Koide and colour channels under the reversal family & Part 3 & Theorem 3.104 \\
\multicolumn{5}{@{}l@{}}{\hspace*{1em}\small\itshape Tier F.5 — Spin–Charge Decomposition via Complete Force Operator}\\[1pt]
\texttt{FT{.}14} & theorem & Spin–charge decomposition & Part 3 & Theorem 3.14 \\
\multicolumn{5}{@{}l@{}}{\hspace*{1em}\small\itshape Tier F.42 — The Rung-One Chirality, the Janus Arrow, and the Family Channel of the Sixteen-Dimensional Arena}\\[1pt]
\texttt{FT{.}141} & theorem & The rung-one chirality theorem: the grading is the chirality of the complete force complex & Part 3 & Theorem 3.105 \\
\texttt{FT{.}142} & theorem & The rung-one arena and the channel angle & Part 3 & Theorem 3.106 \\
\texttt{FT{.}143} & theorem & The rung-one family spectrum & Part 3 & Theorem 3.107 \\
\texttt{FT{.}144} & theorem & The Janus-arrow theorem: no backward materialization, and single-sector occupancy at rung one & Part 3 & Theorem 3.108 \\
\texttt{FT{.}145} & theorem & The rung-one completion & Part 3 & Theorem 3.109 \\
\multicolumn{5}{@{}l@{}}{\hspace*{1em}\small\itshape Tier F.5 — Spin–Charge Decomposition via Complete Force Operator}\\[1pt]
\texttt{FT{.}15} & theorem & Coarse-graining invariance & Part 3 & Theorem 3.15 \\
\texttt{FT{.}153} & theorem & Boundary-schedule attainment & Part 3 & Theorem 3.17 \\
\multicolumn{5}{@{}l@{}}{\hspace*{1em}\small\itshape Tier F.5b — The Mass-Read Registry and the Tilt Composition}\\[1pt]
\texttt{FT{.}154} & theorem & The tilt-composition law & Part 3 & Theorem 3.19 \\
\multicolumn{5}{@{}l@{}}{\hspace*{1em}\small\itshape Tier F.26 — The Platonic Pole Spectrum: Exact Multiplets, the Bounded-Hierarchy Ceiling, and the SM-Hierarchy Audit}\\[1pt]
\texttt{FT{.}155} & theorem & Minimality of the defect interface & Part 3 & Theorem 3.62 \\
\multicolumn{5}{@{}l@{}}{\hspace*{1em}\small\itshape Tier F.44 — The 3+1 Chiral Asymmetry: The Manifold Locus of Chirality, the Spatial Provenance of the Channel Angle, and the Emergent-Time Identity}\\[1pt]
\texttt{FT{.}156} & theorem & The 3+1 asymmetry theorem: one signature datum at two loci, and the manifold structure as the chiral structure & Part 3 & Theorem 3.110 \\
\texttt{FT{.}157} & theorem & The weak-selectivity theorem: the weak force is the shift block — the spatial force on the spatial branch & Part 3 & Theorem 3.111 \\
\texttt{FT{.}158} & theorem & The spatial provenance of the channel angle: is the weak rotation of the triad, and the Koide-form spectrum is forced & Part 3 & Theorem 3.112 \\
\texttt{FT{.}159} & theorem & The emergent-time identity theorem: gravitational time and CPT time are the ledger order & Part 3 & Theorem 3.113 \\
\multicolumn{5}{@{}l@{}}{\hspace*{1em}\small\itshape Tier F.5 — Spin–Charge Decomposition via Complete Force Operator}\\[1pt]
\texttt{FT{.}16} & theorem & The single-vertex hidden-information gap is generic; strong's signature is contraction-incompleteness & Part 3 & Theorem 3.16 \\
\multicolumn{5}{@{}l@{}}{\hspace*{1em}\small\itshape Tier F.44 — The 3+1 Chiral Asymmetry: The Manifold Locus of Chirality, the Spatial Provenance of the Channel Angle, and the Emergent-Time Identity}\\[1pt]
\texttt{FT{.}160} & theorem & --- & Addendum & B.1.14 \\
\multicolumn{5}{@{}l@{}}{\hspace*{1em}\small\itshape Tier F.5 — Spin–Charge Decomposition via Complete Force Operator}\\[1pt]
\texttt{FT{.}17} & theorem & Internal-edge lemma & Part 3 & Theorem 3.18 \\
\multicolumn{5}{@{}l@{}}{\hspace*{1em}\small\itshape Tier F.6 — Subnetwork-Level Force Structure: Confinement and the Active Aspect Set}\\[1pt]
\texttt{FT{.}18} & theorem & Contraction-completeness is not -local & Part 3 & Theorem 3.20 \\
\texttt{FT{.}19} & theorem & Confinement is an all-or-nothing subnetwork property & Part 3 & Theorem 3.21 \\
\multicolumn{5}{@{}l@{}}{\hspace*{1em}\small\itshape Tier F.1 — Information Flux and Measurement Primitivity}\\[1pt]
\texttt{FT{.}2} & theorem & Exhaustiveness of the real-edge measurement primitive & Part 3 & Theorem 3.2 \\
\multicolumn{5}{@{}l@{}}{\hspace*{1em}\small\itshape Tier F.6 — Subnetwork-Level Force Structure: Confinement and the Active Aspect Set}\\[1pt]
\texttt{FT{.}20} & theorem & Reduction of to proven per-aspect criteria & Part 3 & Theorem 3.22 \\
\multicolumn{5}{@{}l@{}}{\hspace*{1em}\small\itshape Tier F.7 — The Coarse-Grained Slot Algebra: \$\textbackslash{}mathcal A(\textbackslash{}Sigma)\$ Properly Derived}\\[1pt]
\texttt{FT{.}21} & theorem & 's charge-type rows, properly derived & Part 3 & Theorem 3.23 \\
\multicolumn{5}{@{}l@{}}{\hspace*{1em}\small\itshape Tier F.8 — Line-Element Quantization from Materialization}\\[1pt]
\texttt{FT{.}22} & theorem & Loop quantization of materialized bound chains & Part 3 & Theorem 3.24 \\
\texttt{FT{.}23} & theorem & Orientation-evenness of the line element & Part 3 & Theorem 3.25 \\
\texttt{FT{.}24} & theorem & The minimal fermion pins the geometry quantum & Part 3 & Theorem 3.26 \\
\texttt{FT{.}25} & theorem & Internal holonomy is well-defined on-shell & Part 3 & Theorem 3.27 \\
\multicolumn{5}{@{}l@{}}{\hspace*{1em}\small\itshape Tier F.9 — The Activated Aspect Operator: \$\textbackslash{}mathcal A(\textbackslash{}Sigma)\$ as a Computable Mask}\\[1pt]
\texttt{FT{.}26} & theorem & Structural transfer: the -mask inherits CT2.3–CT2.5 & Part 3 & Theorem 3.28 \\
\texttt{FT{.}27} & theorem & Single-vertex degeneration; action versus observable & Part 3 & Theorem 3.29 \\
\multicolumn{5}{@{}l@{}}{\hspace*{1em}\small\itshape Tier F.10 — The Strong Edge Weight}\\[1pt]
\texttt{FT{.}29} & theorem & Strong weight: factorization, no , observable sealing & Part 3 & Theorem 3.30 \\
\multicolumn{5}{@{}l@{}}{\hspace*{1em}\small\itshape Tier F.1 — Information Flux and Measurement Primitivity}\\[1pt]
\texttt{FT{.}3} & theorem & Unresolvability of -sourced subnetwork structure & Part 3 & Theorem 3.3 \\
\multicolumn{5}{@{}l@{}}{\hspace*{1em}\small\itshape Tier F.11 — Sector Quantum Audit and the Continuum Ledger}\\[1pt]
\texttt{FT{.}30} & theorem & Pair-channel sector quantum audit & Part 3 & Theorem 3.31 \\
\texttt{FT{.}31} & theorem & The continuum limit survives on-shell quantization & Part 3 & Theorem 3.32 \\
\multicolumn{5}{@{}l@{}}{\hspace*{1em}\small\itshape Tier F.12 — The Cycle Floor Family, the Defect Bridge, and Composite Floors}\\[1pt]
\texttt{FT{.}33} & theorem & The defect bridge: pole-side mass candidate in well-posed ratio form & Part 3 & Theorem 3.33 \\
\texttt{FT{.}34} & theorem & Round-trip quantization between floors; integrality propagation; sealing partition rule & Part 3 & Theorem 3.34 \\
\multicolumn{5}{@{}l@{}}{\hspace*{1em}\small\itshape Tier F.13 — Gap-Closure Pass: Flux Ledger, Relating Equation, Rigidity, Confinement Converse, Core Charge, Aspect Ledger}\\[1pt]
\texttt{FT{.}36} & theorem & Flux ledger completion & Part 3 & Theorem 3.35 \\
\texttt{FT{.}37} & theorem & The relating equation, complete & Part 3 & Theorem 3.36 \\
\texttt{FT{.}38} & theorem & Non-geodesic tightening and on-shell rigidity & Part 3 & Theorem 3.37 \\
\texttt{FT{.}39} & theorem & Confinement converse, singlet diagonal, and filter placement & Part 3 & Theorem 3.38 \\
\multicolumn{5}{@{}l@{}}{\hspace*{1em}\small\itshape Tier F.2 — The Fermion Definition}\\[1pt]
\texttt{FT{.}4} & theorem & External state characterization of a fermion terminus & Part 3 & Theorem 3.4 \\
\multicolumn{5}{@{}l@{}}{\hspace*{1em}\small\itshape Tier F.13 — Gap-Closure Pass: Flux Ledger, Relating Equation, Rigidity, Confinement Converse, Core Charge, Aspect Ledger}\\[1pt]
\texttt{FT{.}40} & theorem & Closed-core charge read & Part 3 & Theorem 3.39 \\
\texttt{FT{.}41} & theorem & The per-aspect ledger & Part 3 & Theorem 3.40 \\
\multicolumn{5}{@{}l@{}}{\hspace*{1em}\small\itshape Tier F.15 — The Regular-Polygon Floors: Full Ladder, Sectors, and the Non-Existence Ceiling}\\[1pt]
\texttt{FT{.}46} & theorem & Distinctness is unbounded: no upper limit on distinct regular polygons & Part 3 & Theorem 3.41 \\
\texttt{FT{.}47} & theorem & No geometric ceiling: the Gauss–Bonnet bound binds area, never & Part 3 & Theorem 3.42 \\
\texttt{FT{.}48} & theorem & Stability structure of floor-to-floor channels (Gap 18, partial) & Part 3 & Theorem 3.43 \\
\multicolumn{5}{@{}l@{}}{\hspace*{1em}\small\itshape Tier F.3 — Information Flux as Rate}\\[1pt]
\texttt{FT{.}5} & theorem & Flux bound (rate form of FT.1) & Part 3 & Theorem 3.5 \\
\multicolumn{5}{@{}l@{}}{\hspace*{1em}\small\itshape Tier F.16 — Composition, Enumeration, and the Spin-Connection Obstruction}\\[1pt]
\texttt{FT{.}52} & theorem & Spin extraction: the constraint system, uniqueness of the stretched rule, and the localized obstruction & Part 3 & Theorem 3.44 \\
\multicolumn{5}{@{}l@{}}{\hspace*{1em}\small\itshape Tier F.17 — The Stationarity Loop: Constancy, the Spin Connection, and the No-Go}\\[1pt]
\texttt{FT{.}53} & theorem & Geodesic constancy: the stationary on-shell ledger is constant-quantum & Part 3 & Theorem 3.45 \\
\texttt{FT{.}54} & theorem & Exact holonomy-defect formula and the closed-form pole side (under Def-H.37(2)) & Part 3 & Theorem 3.46 \\
\texttt{FT{.}55} & theorem & The no-go and the surviving spin realization & Part 3 & Theorem 3.47 \\
\multicolumn{5}{@{}l@{}}{\hspace*{1em}\small\itshape Tier F.18 — Native Rungs, Boundary Exchange Shape, and the Admissibility Correction}\\[1pt]
\texttt{FT{.}56} & theorem & Native ladder rungs, the spin–orbit decomposition, and the missing map & Part 3 & Theorem 3.48 \\
\texttt{FT{.}57} & theorem & The boundary exchange algebra: root-space shape & Part 3 & Theorem 3.49 \\
\texttt{FT{.}58} & theorem & Hosting classification and the no-hosting theorem & Part 3 & Theorem 3.50 \\
\multicolumn{5}{@{}l@{}}{\hspace*{1em}\small\itshape Tier F.19 — The Three Tails: Realization Non-Existence, the Multiplicity Map, and the Selection Obstruction}\\[1pt]
\texttt{FT{.}59} & theorem & --- & Addendum & B.3.2 \\
\multicolumn{5}{@{}l@{}}{\hspace*{1em}\small\itshape Tier F.3 — Information Flux as Rate}\\[1pt]
\texttt{FT{.}6} & theorem & Conformal invariance of combinatorial flux & Part 3 & Theorem 3.6 \\
\multicolumn{5}{@{}l@{}}{\hspace*{1em}\small\itshape Tier F.19 — The Three Tails: Realization Non-Existence, the Multiplicity Map, and the Selection Obstruction}\\[1pt]
\texttt{FT{.}60} & theorem & The multiplicity map: conditions, ladders, basis-relativity, and the honesty witnesses & Part 3 & Theorem 3.51 \\
\texttt{FT{.}61} & theorem & The selection side: phase inventory, index-blindness, and the two exits (Gap 18) & Part 3 & Theorem 3.52 \\
\multicolumn{5}{@{}l@{}}{\hspace*{1em}\small\itshape Tier F.20 — The Bidirectional Rung Audit and the Platonic Physical Cards}\\[1pt]
\texttt{FT{.}62} & theorem & The bidirectional rung audit (correction theorem) & Part 3 & Theorem 3.53 \\
\multicolumn{5}{@{}l@{}}{\hspace*{1em}\small\itshape Tier F.22 — The Bott Correction: Degenerate Witnesses Re-Evaluated, Reconstructability, and the Selection Dilemma}\\[1pt]
\texttt{FT{.}66} & theorem & The Bott correction: witness re-evaluation, reconstructability, and the inadmissibility of absolute selection & Part 3 & Theorem 3.54 \\
\multicolumn{5}{@{}l@{}}{\hspace*{1em}\small\itshape Tier F.23 — The Virtual Bypass: Path Integration over Crystallization Patterns}\\[1pt]
\texttt{FT{.}67} & theorem & The exact-cover theorem and the single-vertex bypass & Part 3 & Theorem 3.55 \\
\texttt{FT{.}68} & theorem & Joint skeletons: tetrahedron, the polygon ladder, the cube; the layer dichotomy; flicker & Part 3 & Theorem 3.56 \\
\multicolumn{5}{@{}l@{}}{\hspace*{1em}\small\itshape Tier F.3 — Information Flux as Rate}\\[1pt]
\texttt{FT{.}7} & theorem & Universal flux density bound & Part 3 & Theorem 3.7 \\
\multicolumn{5}{@{}l@{}}{\hspace*{1em}\small\itshape Tier F.24 — The Flicker Phase Structure, the Cycle Grammar of Covers, the Dust Equation of State, Face Decomposition, and the Bridge Audit}\\[1pt]
\texttt{FT{.}70} & theorem & The flicker phase diagram: unique crossover, computed boundaries, the polygon limit & Part 3 & Theorem 3.57 \\
\texttt{FT{.}71} & theorem & The flicker-dust equation of state & Part 3 & Theorem 3.58 \\
\texttt{FT{.}72} & theorem & Two-fold face decomposition; the Bernoulli marginals ("coarse-grained fuzz as data") & Part 3 & Theorem 3.59 \\
\texttt{FT{.}73} & theorem & --- & Addendum & B.2.2 \\
\multicolumn{5}{@{}l@{}}{\hspace*{1em}\small\itshape Tier F.25 — The Chain-Local \$\textbackslash{}aleph\$-Domains: Extended Exact-Cover Theorems and the Domain-Relativity of the Amplitude Layer}\\[1pt]
\texttt{FT{.}74} & theorem & The extended exact-cover theorems & Part 3 & Theorem 3.60 \\
\texttt{FT{.}75} & theorem & Domain-relativity of the amplitude layer: lock revival, floor revival, parity, and the invariant bridge gap & Part 3 & Theorem 3.61 \\
\multicolumn{5}{@{}l@{}}{\hspace*{1em}\small\itshape Tier F.26 — The Platonic Pole Spectrum: Exact Multiplets, the Bounded-Hierarchy Ceiling, and the SM-Hierarchy Audit}\\[1pt]
\texttt{FT{.}77} & theorem & --- & Addendum & B.2.3 \\
\multicolumn{5}{@{}l@{}}{\hspace*{1em}\small\itshape Tier F.27 — The Cover Orbit Decomposition: the Ansatz Split, the Combinatorial Ceiling, and the Bridge Census}\\[1pt]
\texttt{FT{.}79} & theorem & The combinatorial ceiling and the cover-hierarchy audit (second clean negative) & Part 3 & Theorem 3.63 \\
\multicolumn{5}{@{}l@{}}{\hspace*{1em}\small\itshape Tier F.3 — Information Flux as Rate}\\[1pt]
\texttt{FT{.}8} & theorem & Flux-mediated coarse-graining connection & Part 3 & Theorem 3.8 \\
\multicolumn{5}{@{}l@{}}{\hspace*{1em}\small\itshape Tier F.27 — The Cover Orbit Decomposition: the Ansatz Split, the Combinatorial Ceiling, and the Bridge Census}\\[1pt]
\texttt{FT{.}80} & theorem & The bridge census: covers are amplitude-real, norm-real, and spectrally un-indexed & Part 3 & Theorem 3.64 \\
\multicolumn{5}{@{}l@{}}{\hspace*{1em}\small\itshape Tier F.28 — The Full-SM Completeness Audit: Particle Content, the Gauge Blocks, the Rung Ledger, and the Deficit Table}\\[1pt]
\texttt{FT{.}82} & theorem & Gauge-block reconstruction (audit pass B): the algebra closes, the coupling structure does not & Part 3 & Theorem 3.65 \\
\multicolumn{5}{@{}l@{}}{\hspace*{1em}\small\itshape Tier F.29 — Gravitational Reads in Three States, the Flicker Ripple, and the Tumbling-Model Adjudication}\\[1pt]
\texttt{FT{.}85} & theorem & Vacuum census, baseline equality, and the constructible gravitational reads & Part 3 & Theorem 3.66 \\
\texttt{FT{.}86} & theorem & The flicker ripple: ultralocal raw dynamics, cap-induced alternation, and the quantum–classical gradient & Part 3 & Theorem 3.67 \\
\texttt{FT{.}87} & theorem & --- & Addendum & B.2.4 \\
\texttt{FT{.}88} & theorem & The involution algebra of the two declared reversals & Part 3 & Theorem 3.68 \\
\multicolumn{5}{@{}l@{}}{\hspace*{1em}\small\itshape Tier F.30 — The Koide Adjudication, the Resolution Ladder, and the Minimal Bridge Declaration}\\[1pt]
\texttt{FT{.}89} & theorem & --- & Addendum & B.1.13 \\
\multicolumn{5}{@{}l@{}}{\hspace*{1em}\small\itshape Tier F.4 — Rung-Dependent Edge Structure}\\[1pt]
\texttt{FT{.}9} & theorem & EM part of edge weight is rung-stable & Part 3 & Theorem 3.9 \\
\multicolumn{5}{@{}l@{}}{\hspace*{1em}\small\itshape Tier F.30 — The Koide Adjudication, the Resolution Ladder, and the Minimal Bridge Declaration}\\[1pt]
\texttt{FT{.}90} & theorem & --- & Addendum & B.2.5 \\
\texttt{FT{.}91} & theorem & The minimal spectral-bridge form & Part 3 & Theorem 3.69 \\
\multicolumn{5}{@{}l@{}}{\hspace*{1em}\small\itshape Tier F.31 — The Native Clifford Census, the Weak-Chirality Pincer, and the Reversal Census of the Covers}\\[1pt]
\texttt{FT{.}92} & theorem & EW.5 is convention-exact: re-verification and the rigidity of the named basis & Part 3 & Theorem 3.70 \\
\texttt{FT{.}93} & theorem & The native Clifford census and the weak-chirality pincer (Tasks 1–3) & Part 3 & Theorem 3.71 \\
\texttt{FT{.}94} & theorem & --- & Addendum & B.2.6 \\
\texttt{FT{.}95} & theorem & Flicker chirality: the tautology, the absent trajectory, and the zero-expectation theorem & Part 3 & Theorem 3.72 \\
\multicolumn{5}{@{}l@{}}{\hspace*{1em}\small\itshape Tier F.32 — The Multi-Node Form Factor: Formalization, Exact Reductions, and the Third Ceiling}\\[1pt]
\texttt{FT{.}97} & theorem & The form factor formalized: node ontology, three underdeterminations, reduction identities, and the general ceiling & Part 3 & Theorem 3.73 \\
\addlinespace[4pt]
\multicolumn{5}{@{}l}{\textbf{Part 4 · Baryomesonic\_era — colour-sealed composites}}\\[2pt]
\multicolumn{5}{@{}l@{}}{\hspace*{1em}\small\itshape Tier BM.1 — Observability of Sealed Composites}\\[1pt]
\texttt{BM{.}1} & theorem & Confinement inheritance and the trace-only read & Part 4 & Theorem 4.1 \\
\multicolumn{5}{@{}l@{}}{\hspace*{1em}\small\itshape Tier BM.4 — Noether's Theorem from the Janus Seal Topology: the Discrete Conservation Law, the Continuum Promotion, the Stress-Energy of a Sealed Composite, and the Canonical-Pair Obstructions}\\[1pt]
\texttt{BM{.}10} & theorem & --- & Addendum & B.3.5 \\
\multicolumn{5}{@{}l@{}}{\hspace*{1em}\small\itshape Tier BM.19 — The Flicker Integral, the Rest-Mass Q-Value, and the Relating Equation: the Materialized Ensemble, the Vertex Superposition, the Rung-One \$\textbackslash{}nu\$ Line, and the Bridge in Four Exact Forms}\\[1pt]
\texttt{BM{.}100} & theorem & --- & Addendum & B.1.29 \\
\multicolumn{5}{@{}l@{}}{\hspace*{1em}\small\itshape Tier BM.20 — Six Objective Resolutions: the Energy-Scale Graph, the Dimensional Bridge, the Eigenvalue Composition, the Curvature Composition, the General Crystallization Weight, and the Kind Structure of the Arrow}\\[1pt]
\texttt{BM{.}101} & theorem & The energy-scale graph: the two reads, the derived crossing strength, and the composition-dependent crossing count & Part 4 & Theorem 4.85 \\
\texttt{BM{.}102} & theorem & The dimensional bridge: the -exponent census, the seven rows, and the declared calibration slot & Part 4 & Theorem 4.86 \\
\texttt{BM{.}103} & theorem & The three-constituent assembly: the eigenvalue equations that are formed, the tensor-product referent census, and the composition residue & Part 4 & Theorem 4.87 \\
\texttt{BM{.}104} & theorem & The curvature composition: the back-reaction fixed point, its exact roots and pole locus, and the non-selection theorem & Part 4 & Theorem 4.88 \\
\texttt{BM{.}105} & theorem & The general crystallization weight: the measure-free bounds, the threshold family, and the refuted defect conditioning & Part 4 & Theorem 4.89 \\
\texttt{BM{.}106} & theorem & The arrow and the kind: the spin label, the tick law, and the exchange-structure referent census & Part 4 & Theorem 4.90 \\
\texttt{BM{.}107} & theorem & --- & Addendum & B.1.30 \\
\multicolumn{5}{@{}l@{}}{\hspace*{1em}\small\itshape Tier BM.21 — The Antimatter and Resonance Tier: the CP-Conjugate, the Meson Seal, the Exchange Operator, and the Local Resonance Spectrum}\\[1pt]
\texttt{BM{.}108} & theorem & The CP-conjugate chain: the reversal census, the T-block, the exact phase, and the forced conjugate & Part 4 & Theorem 4.91 \\
\texttt{BM{.}109} & theorem & The CP-pair as the vacuum fluctuation: the creation weight, its exact range, and the degeneracy rows & Part 4 & Theorem 4.92 \\
\multicolumn{5}{@{}l@{}}{\hspace*{1em}\small\itshape Tier BM.4 — Noether's Theorem from the Janus Seal Topology: the Discrete Conservation Law, the Continuum Promotion, the Stress-Energy of a Sealed Composite, and the Canonical-Pair Obstructions}\\[1pt]
\texttt{BM{.}11} & theorem & The continuum limit of the Janus topology: the differential conservation law and its two agreeing routes & Part 4 & Theorem 4.10 \\
\multicolumn{5}{@{}l@{}}{\hspace*{1em}\small\itshape Tier BM.21 — The Antimatter and Resonance Tier: the CP-Conjugate, the Meson Seal, the Exchange Operator, and the Local Resonance Spectrum}\\[1pt]
\texttt{BM{.}110} & theorem & The meson seal: the -sealed -juncture, its uniqueness, and the charge class it occupies & Part 4 & Theorem 4.93 \\
\texttt{BM{.}111} & theorem & The meson mass: the two additive routes, the vanishing seal-binding, and the amplitude/norm reconciliation & Part 4 & Theorem 4.94 \\
\texttt{BM{.}112} & theorem & The exchange operator: construction, involution, spectrum, and the eigenspace census & Part 4 & Theorem 4.95 \\
\texttt{BM{.}113} & theorem & Spin–statistics: the phase alphabet of the swap, the winding parity, and the derived connection & Part 4 & Theorem 4.96 \\
\texttt{BM{.}114} & theorem & The resonance spectrum: the two characteristic polynomials, the orbit law, and the one exact collision & Part 4 & Theorem 4.97 \\
\texttt{BM{.}115} & theorem & The materialization threshold: the exact function of , and , and the never-empty locus & Part 4 & Theorem 4.98 \\
\texttt{BM{.}116} & theorem & The decay width: the exact function of the local parameters, its three exact zeros, and the lifetime floor & Part 4 & Theorem 4.99 \\
\texttt{BM{.}117} & theorem & The groupings, the ceilings, and the octet adjudication & Part 4 & Theorem 4.100 \\
\texttt{BM{.}118} & theorem & --- & Addendum & B.1.31 \\
\multicolumn{5}{@{}l@{}}{\hspace*{1em}\small\itshape Tier BM.22 — The Annihilation Channel, Concrete Spin, and Fermion Identification: the Janus Tunnel, the Zig-Zag World Line, and the Spin-Dependent Transport Defect}\\[1pt]
\texttt{BM{.}119} & theorem & --- & Addendum & B.1.32 \\
\multicolumn{5}{@{}l@{}}{\hspace*{1em}\small\itshape Tier BM.4 — Noether's Theorem from the Janus Seal Topology: the Discrete Conservation Law, the Continuum Promotion, the Stress-Energy of a Sealed Composite, and the Canonical-Pair Obstructions}\\[1pt]
\texttt{BM{.}12} & theorem & --- & Addendum & B.1.15 \\
\multicolumn{5}{@{}l@{}}{\hspace*{1em}\small\itshape Tier BM.22 — The Annihilation Channel, Concrete Spin, and Fermion Identification: the Janus Tunnel, the Zig-Zag World Line, and the Spin-Dependent Transport Defect}\\[1pt]
\texttt{BM{.}120} & theorem & The Janus tunnel: the crossing operators, the exact spectrum, and detailed balance & Part 4 & Theorem 4.101 \\
\texttt{BM{.}121} & theorem & The zig-zag world line: the even resolvent, the round-trip weights, and no backward leg & Part 4 & Theorem 4.102 \\
\texttt{BM{.}122} & theorem & The pair channel: charge balance, the exact threshold, the CP structure, and the two-photon bar & Part 4 & Theorem 4.103 \\
\texttt{BM{.}123} & theorem & The annihilation channel: the reverse move, the entropy destination, and the closure of the cross-section residue & Part 4 & Theorem 4.104 \\
\texttt{BM{.}124} & theorem & The spin–winding relation and the multiplet defect ladder & Part 4 & Theorem 4.105 \\
\texttt{BM{.}125} & theorem & The spin-dependent mass: the exact ladder, the two clauses and their matching, and the ceiling & Part 4 & Theorem 4.106 \\
\texttt{BM{.}126} & theorem & --- & Addendum & B.2.7 \\
\texttt{BM{.}127} & theorem & The winding-parity identification: two reads separated & Part 4 & Theorem 4.107 \\
\texttt{BM{.}128} & theorem & The three-chain exchange structure: the space, the action, the character formula, and the vanishing wedge & Part 4 & Theorem 4.108 \\
\texttt{BM{.}129} & theorem & The meson is a boson and the baryon's sign is the product: the composite identification & Part 4 & Theorem 4.109 \\
\multicolumn{5}{@{}l@{}}{\hspace*{1em}\small\itshape Tier BM.5 — The Calibration Lock: the Born Width, the Coherence Volume, the Curvature of a Sealed Composite, and the Neutron Row}\\[1pt]
\texttt{BM{.}13} & theorem & The Born width: the exact substitution, the quantized ladder, and the sub-hop locus & Part 4 & Theorem 4.11 \\
\multicolumn{5}{@{}l@{}}{\hspace*{1em}\small\itshape Tier BM.22 — The Annihilation Channel, Concrete Spin, and Fermion Identification: the Janus Tunnel, the Zig-Zag World Line, and the Spin-Dependent Transport Defect}\\[1pt]
\texttt{BM{.}130} & theorem & --- & Addendum & B.1.33 \\
\multicolumn{5}{@{}l@{}}{\hspace*{1em}\small\itshape Tier BM.23 — The Closing Tier: the Haar–Poisson Duality and the Lapse, the Selected Curvature Composition, the Momentum-Space Resolvent, the Charge and Statistics Censuses, the Two-Chain Ensemble, the Seed-Tensoring Mechanism, the Convention Declarations, the Ceiling Study, the Jensen Residue, the Pair Amplitude, and the Era Closure}\\[1pt]
\texttt{BM{.}131} & theorem & The Haar–Poisson duality: the twirl fixed point, the complementary factorization, and the proper-time contraction & Part 4 & Theorem 4.110 \\
\texttt{BM{.}132} & theorem & The lapse: the derived function of the envelope ratio and the curvature, its formed domain, and the envelope row & Part 4 & Theorem 4.111 \\
\texttt{BM{.}133} & theorem & The selected curvature composition: two derived criteria, the unique survivor, and the completed running coupling & Part 4 & Theorem 4.112 \\
\texttt{BM{.}134} & theorem & The resolvent's derived content: the form uniqueness, the pole reads, the on-shell measure, and the two-body phase space & Part 4 & Theorem 4.113 \\
\texttt{BM{.}135} & theorem & The multi-chain additive charge census: the baryon and meson rows, the stretched states, and the gauged read & Part 4 & Theorem 4.114 \\
\texttt{BM{.}136} & theorem & The action is the colour 3-vertex: the equivariance theorem, the exchange structure at every , and the statistics per class & Part 4 & Theorem 4.115 \\
\texttt{BM{.}137} & theorem & The two-chain ensemble measure, the meson's formed binding, the sector-mass pincer, and the two-chain spectrum & Part 4 & Theorem 4.116 \\
\texttt{BM{.}138} & theorem & The seed-tensoring mechanism: the symmetrizer as a finite-group twirl, the transport and composition closure, and the two residues adjudicated & Part 4 & Theorem 4.117 \\
\texttt{BM{.}139} & theorem & The two convention declarations and the sector locus & Part 4 & Theorem 4.118 \\
\multicolumn{5}{@{}l@{}}{\hspace*{1em}\small\itshape Tier BM.5 — The Calibration Lock: the Born Width, the Coherence Volume, the Curvature of a Sealed Composite, and the Neutron Row}\\[1pt]
\texttt{BM{.}14} & theorem & The coherence volume and the energy density & Part 4 & Theorem 4.12 \\
\multicolumn{5}{@{}l@{}}{\hspace*{1em}\small\itshape Tier BM.23 — The Closing Tier: the Haar–Poisson Duality and the Lapse, the Selected Curvature Composition, the Momentum-Space Resolvent, the Charge and Statistics Censuses, the Two-Chain Ensemble, the Seed-Tensoring Mechanism, the Convention Declarations, the Ceiling Study, the Jensen Residue, the Pair Amplitude, and the Era Closure}\\[1pt]
\texttt{BM{.}140} & theorem & --- & Addendum & B.2.8 \\
\texttt{BM{.}141} & theorem & The Jensen residue: the exact closed form, its range, its zero locus, the two-chain companion, and the width adjudication & Part 4 & Theorem 4.119 \\
\texttt{BM{.}142} & theorem & The supply vertex, the pair amplitude, and the annihilation cross-section: the four grounds resolved and residues , closed & Part 4 & Theorem 4.120 \\
\texttt{BM{.}143} & theorem & --- & Addendum & B.4.2 \\
\texttt{BM{.}144} & theorem & --- & Addendum & B.1.34 \\
\multicolumn{5}{@{}l@{}}{\hspace*{1em}\small\itshape Tier BM.24 — The Continuum-Limit Tier: the Force Operator over the Graph, the Continuum-Limit Function, the Metric–Energy–Continuum Map, the Breaking Threshold, the Locked-Sector Complement, the Universality, the Reclassification of the Continuum Conditions, and the Closure of the Foundational Era}\\[1pt]
\texttt{BM{.}145} & theorem & The force operator acting over the graph: the generator census, the mask ladder, and the derived action & Part 4 & Theorem 4.121 \\
\texttt{BM{.}146} & theorem & The continuum-limit function: the two variables, the five contributors, and the exact functional form & Part 4 & Theorem 4.122 \\
\texttt{BM{.}147} & theorem & The metric–energy–continuum map: the dictionary bijection, the closed form of the source read, and the target & Part 4 & Theorem 4.123 \\
\texttt{BM{.}148} & theorem & The breaking threshold: the exact capacity of a chart cell & Part 4 & Theorem 4.124 \\
\texttt{BM{.}149} & theorem & The locked sector at the breaking locus: the two faces of one capacity, the absorption, the membrane, and the complement theorem & Part 4 & Theorem 4.125 \\
\multicolumn{5}{@{}l@{}}{\hspace*{1em}\small\itshape Tier BM.5 — The Calibration Lock: the Born Width, the Coherence Volume, the Curvature of a Sealed Composite, and the Neutron Row}\\[1pt]
\texttt{BM{.}15} & theorem & The curvature of a sealed composite: the exact trace of the field equations, and the curvature scale & Part 4 & Theorem 4.13 \\
\multicolumn{5}{@{}l@{}}{\hspace*{1em}\small\itshape Tier BM.24 — The Continuum-Limit Tier: the Force Operator over the Graph, the Continuum-Limit Function, the Metric–Energy–Continuum Map, the Breaking Threshold, the Locked-Sector Complement, the Universality, the Reclassification of the Continuum Conditions, and the Closure of the Foundational Era}\\[1pt]
\texttt{BM{.}150} & theorem & The universality of the continuum limit: one operator, one arena, one action, three fingerprints & Part 4 & Theorem 4.126 \\
\texttt{BM{.}151} & theorem & The derivation of CT2.6(ii)–(iii): the unbounded ladder, the spectral bound, and the empty extra-axiomatic list & Part 4 & Theorem 4.127 \\
\texttt{BM{.}152} & theorem & --- & Addendum & B.1.35 \\
\multicolumn{5}{@{}l@{}}{\hspace*{1em}\small\itshape Tier BM.5 — The Calibration Lock: the Born Width, the Coherence Volume, the Curvature of a Sealed Composite, and the Neutron Row}\\[1pt]
\texttt{BM{.}16} & theorem & The calibration lock & Part 4 & Theorem 4.14 \\
\texttt{BM{.}17} & theorem & The neutron row of the identification layer & Part 4 & Theorem 4.15 \\
\texttt{BM{.}18} & theorem & --- & Addendum & B.1.16 \\
\multicolumn{5}{@{}l@{}}{\hspace*{1em}\small\itshape Tier BM.6 — The Topology-Blind Neutron and the Topology-Blind + Static Proton: the Flicker Ensemble, the Path-Integrated Mass, the Binding Read, and the Forced Carrier}\\[1pt]
\texttt{BM{.}19} & theorem & The flicker ensemble is derived: the stationary measure, the threshold specialization, and the no-flicker limit & Part 4 & Theorem 4.16 \\
\multicolumn{5}{@{}l@{}}{\hspace*{1em}\small\itshape Tier BM.1 — Observability of Sealed Composites}\\[1pt]
\texttt{BM{.}2} & theorem & Junction phase and winding grammar & Part 4 & Theorem 4.2 \\
\multicolumn{5}{@{}l@{}}{\hspace*{1em}\small\itshape Tier BM.6 — The Topology-Blind Neutron and the Topology-Blind + Static Proton: the Flicker Ensemble, the Path-Integrated Mass, the Binding Read, and the Forced Carrier}\\[1pt]
\texttt{BM{.}20} & theorem & The mass from the curvature law on the path-integrated route & Part 4 & Theorem 4.17 \\
\texttt{BM{.}21} & theorem & The binding energy: additive in the count, non-additive in the ensemble & Part 4 & Theorem 4.18 \\
\texttt{BM{.}22} & theorem & The static force matrix on the two charge classes & Part 4 & Theorem 4.19 \\
\texttt{BM{.}23} & theorem & The transition: the ladder rule, the forced carrier, and the resolvent rate & Part 4 & Theorem 4.20 \\
\texttt{BM{.}24} & theorem & --- & Addendum & B.1.17 \\
\multicolumn{5}{@{}l@{}}{\hspace*{1em}\small\itshape Tier BM.7 — The Proton in the Positive CPT Sector and Its Mirror: the Sector Labels, the Exact Phase, and the Charge-Blind Geometry}\\[1pt]
\texttt{BM{.}25} & theorem & The CPT involution on the sealed composite: orbits, the exact phase, and the census ratio & Part 4 & Theorem 4.21 \\
\texttt{BM{.}26} & theorem & The proton row: the positive sector & Part 4 & Theorem 4.22 \\
\texttt{BM{.}27} & theorem & The antiproton row: the CPT image, and the resolution of the identity question & Part 4 & Theorem 4.23 \\
\texttt{BM{.}28} & theorem & The charge-constrained probability space: the census, the three mass routes, and the splitting & Part 4 & Theorem 4.24 \\
\texttt{BM{.}29} & theorem & --- & Addendum & B.1.18 \\
\multicolumn{5}{@{}l@{}}{\hspace*{1em}\small\itshape Tier BM.2 — The Family Group and the Invariance Laws of the Seal}\\[1pt]
\texttt{BM{.}3} & theorem & The family group of the seal & Part 4 & Theorem 4.3 \\
\multicolumn{5}{@{}l@{}}{\hspace*{1em}\small\itshape Tier BM.8 — The Additive Charge Read and the Mass-Route Adjudication}\\[1pt]
\texttt{BM{.}30} & theorem & The mass-route adjudication: three exact ratios, one read choice & Part 4 & Theorem 4.25 \\
\multicolumn{5}{@{}l@{}}{\hspace*{1em}\small\itshape Tier BM.9 — The Mesonic Flicker Test: the Contraction Grammar, the Colour Projector, and the Protected Junction}\\[1pt]
\texttt{BM{.}31} & theorem & The boundary-data comparison and the grammar verdict & Part 4 & Theorem 4.26 \\
\texttt{BM{.}32} & theorem & The mesonic contribution: zero by filtration, and the amplitude layer is orthogonal & Part 4 & Theorem 4.27 \\
\texttt{BM{.}33} & theorem & --- & Addendum & B.1.19 \\
\multicolumn{5}{@{}l@{}}{\hspace*{1em}\small\itshape Tier BM.10 — The Kinetic Energy from the Metric and the Vacuum-Balanced Entropy: the Dispersion Read of a Sealed Composite, the Deficit Reframing of the Rest Mass, the Force-Structure Ladder, and the Route-Artifact Adjudication}\\[1pt]
\texttt{BM{.}34} & theorem & The kinetic energy from the emergent metric: the ADM split, the mass shell, the dispersion relation, and the zero-point bound & Part 4 & Theorem 4.28 \\
\texttt{BM{.}35} & theorem & The vacuum-balanced entropy: the rest mass as the deficit against the all-virtual vacuum & Part 4 & Theorem 4.29 \\
\texttt{BM{.}36} & theorem & The force-structure ladder: the mask dimensions, the transport channel, and the inter-family constraint & Part 4 & Theorem 4.30 \\
\texttt{BM{.}37} & theorem & The mass-route artifacts: the disjoint input alphabets and the closure of the counting trap & Part 4 & Theorem 4.31 \\
\texttt{BM{.}38} & theorem & --- & Addendum & B.1.20 \\
\multicolumn{5}{@{}l@{}}{\hspace*{1em}\small\itshape Tier BM.11 — The Neutron Decay Vertex: the Family Transport Profiles, the Rung-1 Isospin Row, the Forced Neutral Member, the Angular-Momentum Budget, and the Full Kinematic Balance}\\[1pt]
\texttt{BM{.}39} & theorem & The family mass form: positivity, the splitting law, the sum rule, the order inversion, and the area ladder & Part 4 & Theorem 4.32 \\
\multicolumn{5}{@{}l@{}}{\hspace*{1em}\small\itshape Tier BM.2 — The Family Group and the Invariance Laws of the Seal}\\[1pt]
\texttt{BM{.}4} & theorem & The invariance laws of the seal & Part 4 & Theorem 4.4 \\
\multicolumn{5}{@{}l@{}}{\hspace*{1em}\small\itshape Tier BM.11 — The Neutron Decay Vertex: the Family Transport Profiles, the Rung-1 Isospin Row, the Forced Neutral Member, the Angular-Momentum Budget, and the Full Kinematic Balance}\\[1pt]
\texttt{BM{.}40} & theorem & The rung-1 isospin row: the Casimir identity, the branch selectivity, and the neutrality of the -vertex & Part 4 & Theorem 4.33 \\
\texttt{BM{.}41} & theorem & The decay vertex: the forced members, the two-body impossibility, and the Q-value sign theorem & Part 4 & Theorem 4.34 \\
\texttt{BM{.}42} & theorem & The angular momentum at the vertex: the Noether moment, the Lagrange decomposition, the absorption adjudication, and the -selection census & Part 4 & Theorem 4.35 \\
\texttt{BM{.}43} & theorem & The curvature rows of the members: the dust sources, the weak-field potential, and the superposition scope & Part 4 & Theorem 4.36 \\
\texttt{BM{.}44} & theorem & The time evolution: the per-tick weight, the e-folding count, and the unit discipline & Part 4 & Theorem 4.37 \\
\texttt{BM{.}45} & theorem & The assembled reaction: the conservation table, the naming row, the final read, and the adjudication & Part 4 & Theorem 4.38 \\
\multicolumn{5}{@{}l@{}}{\hspace*{1em}\small\itshape Tier BM.12 — The EM Self-Interaction Correction: the Scar's Field Energy, the Infinite-Range Seal Asymmetry, the Mass-Splitting Theorem, and the Beta-Spectrum Adjudication}\\[1pt]
\texttt{BM{.}46} & theorem & The scar's field energy and the mass-splitting theorem & Part 4 & Theorem 4.39 \\
\texttt{BM{.}47} & theorem & The infinite-range seal: UV/IR finiteness, the universe-scale neutrality, and the seal asymmetry & Part 4 & Theorem 4.40 \\
\texttt{BM{.}48} & theorem & The two-body potential, the hydrogenic rows, and the binding category bar & Part 4 & Theorem 4.41 \\
\texttt{BM{.}49} & theorem & The beta-decay re-adjudication with the correction, and the endpoint-to-average refutation & Part 4 & Theorem 4.42 \\
\multicolumn{5}{@{}l@{}}{\hspace*{1em}\small\itshape Tier BM.3 — Consumption and Open Constructions}\\[1pt]
\texttt{BM{.}5} & theorem & The census-consumption law & Part 4 & Theorem 4.5 \\
\multicolumn{5}{@{}l@{}}{\hspace*{1em}\small\itshape Tier BM.12 — The EM Self-Interaction Correction: the Scar's Field Energy, the Infinite-Range Seal Asymmetry, the Mass-Splitting Theorem, and the Beta-Spectrum Adjudication}\\[1pt]
\texttt{BM{.}50} & theorem & --- & Addendum & B.1.21 \\
\multicolumn{5}{@{}l@{}}{\hspace*{1em}\small\itshape Tier BM.13 — The Full-Scope Mass Splitting: the Lyapunov Foundation, the Complete-Force Census, the Non-Uniform Clock, the Binding Zeros, and the Exact Ratio}\\[1pt]
\texttt{BM{.}51} & theorem & The Lyapunov foundation: the contraction law iterated, the deficit links, and the class-degenerate zero & Part 4 & Theorem 4.43 \\
\texttt{BM{.}52} & theorem & The complete force operator and the sector-by-sector splitting census & Part 4 & Theorem 4.44 \\
\texttt{BM{.}53} & theorem & The metric background and the non-uniform clock: the zero metric contribution and the Mandelstam–Tamm resolvability bound & Part 4 & Theorem 4.45 \\
\texttt{BM{.}54} & theorem & The binding energies: the count zero, the degenerate ensemble, and the unconstructed atom & Part 4 & Theorem 4.46 \\
\texttt{BM{.}55} & theorem & The full-scope splitting and the exact ratio & Part 4 & Theorem 4.47 \\
\texttt{BM{.}56} & theorem & --- & Addendum & B.1.22 \\
\multicolumn{5}{@{}l@{}}{\hspace*{1em}\small\itshape Tier BM.14 — The Constituent Transport Splitting: the Vertex Asymmetry, the Observation Frames, the Six-State Quark Assembly, and the Neutron Row}\\[1pt]
\texttt{BM{.}57} & theorem & The continuation-symmetry census: the vertex asymmetry is two, and it splits nothing & Part 4 & Theorem 4.48 \\
\texttt{BM{.}58} & theorem & The observation frames: the metric is charge-blind, and the frame's content is the seal asymmetry & Part 4 & Theorem 4.49 \\
\texttt{BM{.}59} & theorem & The constituent transport splitting: the order residue, the two directions, and the exact positive sign & Part 4 & Theorem 4.50 \\
\multicolumn{5}{@{}l@{}}{\hspace*{1em}\small\itshape Tier BM.3 — Consumption and Open Constructions}\\[1pt]
\texttt{BM{.}6} & theorem & The construction interfaces (specifications) & Part 4 & Theorem 4.6 \\
\multicolumn{5}{@{}l@{}}{\hspace*{1em}\small\itshape Tier BM.14 — The Constituent Transport Splitting: the Vertex Asymmetry, the Observation Frames, the Six-State Quark Assembly, and the Neutron Row}\\[1pt]
\texttt{BM{.}60} & theorem & The six-state assembly: the generation ladder, the branch split, the Platonic floors, and the two hierarchies & Part 4 & Theorem 4.51 \\
\texttt{BM{.}61} & theorem & the inventory, the coupling-blindness theorem, and the obstruction & Part 4 & Theorem 4.52 \\
\texttt{BM{.}62} & theorem & The full mass splitting: the extended decomposition, the exact critical strength, and the sign theorem & Part 4 & Theorem 4.53 \\
\texttt{BM{.}63} & theorem & The decay re-adjudication on the refined read: the closure survives, with an exact margin & Part 4 & Theorem 4.54 \\
\texttt{BM{.}64} & theorem & --- & Addendum & B.1.23 \\
\multicolumn{5}{@{}l@{}}{\hspace*{1em}\small\itshape Tier BM.15 — The Organic Bridge: the Operator EM Lift, the Derived-Holonomy Attempt, the Constrained Joint Entropy, and the Beta-Decay Energy Audit}\\[1pt]
\texttt{BM{.}65} & theorem & The baseline re-certified: the refined rest input, the exact branch, and the integer certificates & Part 4 & Theorem 4.55 \\
\texttt{BM{.}66} & theorem & The operator lift of the EM correction: the charge-blind trace, the quadratic invariant, and the assembled operator rest input & Part 4 & Theorem 4.56 \\
\texttt{BM{.}67} & theorem & The holonomy lift attempted: the derived sector, the angle-free force census, and the two-torsion obstruction & Part 4 & Theorem 4.57 \\
\texttt{BM{.}68} & theorem & The constrained joint entropy as a derived fiber cardinality & Part 4 & Theorem 4.58 \\
\texttt{BM{.}69} & theorem & The bridge converged: the three inputs, the non-derivation theorem, and the minimal priced residue & Part 4 & Theorem 4.59 \\
\multicolumn{5}{@{}l@{}}{\hspace*{1em}\small\itshape Tier BM.4 — Noether's Theorem from the Janus Seal Topology: the Discrete Conservation Law, the Continuum Promotion, the Stress-Energy of a Sealed Composite, and the Canonical-Pair Obstructions}\\[1pt]
\texttt{BM{.}7} & theorem & The discrete conservation law is Janus topology: no symmetry argument, no action, no limit & Part 4 & Theorem 4.7 \\
\multicolumn{5}{@{}l@{}}{\hspace*{1em}\small\itshape Tier BM.15 — The Organic Bridge: the Operator EM Lift, the Derived-Holonomy Attempt, the Constrained Joint Entropy, and the Beta-Decay Energy Audit}\\[1pt]
\texttt{BM{.}70} & theorem & The beta-decay energy audit: the two-bit floor theorem, the margin ladder, and the formed-closed dichotomy & Part 4 & Theorem 4.60 \\
\texttt{BM{.}71} & theorem & The opening-term census: the five routes, the neutron row under the bridge, and the six verdict rows & Part 4 & Theorem 4.61 \\
\texttt{BM{.}72} & theorem & --- & Addendum & B.1.24 \\
\multicolumn{5}{@{}l@{}}{\hspace*{1em}\small\itshape Tier BM.16 — The Chiral Assignment: the Closed Deficit, the Ray-Resolved Spin Census, the Attached Quark Splitting, the Chiral Vertex Dynamics, and the Specified Beta-Decay Opening}\\[1pt]
\texttt{BM{.}73} & theorem & The chiral matter assignment: the issued coupling, the occupancy theorem, and the / matching rows & Part 4 & Theorem 4.62 \\
\texttt{BM{.}74} & theorem & The ray-resolved spin census: the winding labels, the occupancy rows, and the vertex table & Part 4 & Theorem 4.63 \\
\texttt{BM{.}75} & theorem & The weight–branch–core attachment: the unique bijection and the attached quark mass splitting & Part 4 & Theorem 4.64 \\
\texttt{BM{.}76} & theorem & The chiral vertex dynamics: the nilpotent step operator, the flip channel, and the rate rows & Part 4 & Theorem 4.65 \\
\texttt{BM{.}77} & theorem & The read specification and the allowed branch: the formed Q-value, the exact trichotomy, and the rescope rows & Part 4 & Theorem 4.66 \\
\texttt{BM{.}78} & theorem & --- & Addendum & B.1.25 \\
\multicolumn{5}{@{}l@{}}{\hspace*{1em}\small\itshape Tier BM.17 — The Mass Ratio Census: the Platonic Winding Ladders, the Intra- and Inter-Generational Ratios, the Six-State Bridge Census, the Four-Currency Incommensurability, and the Neutron–Proton Ratio}\\[1pt]
\texttt{BM{.}79} & theorem & The Platonic winding census: the ladder law, the Chebyshev tower, and the intra-generational mass ratios & Part 4 & Theorem 4.67 \\
\multicolumn{5}{@{}l@{}}{\hspace*{1em}\small\itshape Tier BM.4 — Noether's Theorem from the Janus Seal Topology: the Discrete Conservation Law, the Continuum Promotion, the Stress-Energy of a Sealed Composite, and the Canonical-Pair Obstructions}\\[1pt]
\texttt{BM{.}8} & theorem & The continuum promotion: the hop becomes a translation, the becomes , the tick stays an arrow, and the gauge algebras were never discrete & Part 4 & Theorem 4.8 \\
\multicolumn{5}{@{}l@{}}{\hspace*{1em}\small\itshape Tier BM.17 — The Mass Ratio Census: the Platonic Winding Ladders, the Intra- and Inter-Generational Ratios, the Six-State Bridge Census, the Four-Currency Incommensurability, and the Neutron–Proton Ratio}\\[1pt]
\texttt{BM{.}80} & theorem & The inter-generational mass ratios: one cubic form, three channel constants, the exact enclosures, and the hierarchy-order theorem & Part 4 & Theorem 4.68 \\
\texttt{BM{.}81} & theorem & The six-state mass-ratio census: the splitting order, the doublet ratios, the two inverted hierarchies, and the three non-proportional ladders & Part 4 & Theorem 4.69 \\
\texttt{BM{.}82} & theorem & The bridge census: the six exact factors, the intra-currency closure, and the single priced residue & Part 4 & Theorem 4.70 \\
\texttt{BM{.}83} & theorem & The inter-family bridge in full: the four-currency census, the exact incommensurability, and the three-number localization & Part 4 & Theorem 4.71 \\
\texttt{BM{.}84} & theorem & The neutron–proton mass ratio: the exact ratio, the irrational critical locus, the range theorem, and the trichotomy & Part 4 & Theorem 4.72 \\
\texttt{BM{.}85} & theorem & The Q-value in two formed reads: the ratio forms, the margin ladders, and the read-relative verdict & Part 4 & Theorem 4.73 \\
\texttt{BM{.}86} & theorem & --- & Addendum & B.1.26 \\
\multicolumn{5}{@{}l@{}}{\hspace*{1em}\small\itshape Tier BM.18 — The EM Self-Energy Correction: the Category Adjudication, the Running Coupling, the Path-Integral Self-Energy, the Branch-Blindness Theorem, and the Corrected Neutron–Proton Splitting}\\[1pt]
\texttt{BM{.}87} & theorem & --- & Addendum & B.1.27 \\
\texttt{BM{.}88} & theorem & The running coupling: the exact geometric profile on the envelope ratio, the lapse, and the curvature lock & Part 4 & Theorem 4.74 \\
\texttt{BM{.}89} & theorem & The EM self-energy from the path integral: the derived measure, the crystallization weight, and the route-relative correction & Part 4 & Theorem 4.75 \\
\multicolumn{5}{@{}l@{}}{\hspace*{1em}\small\itshape Tier BM.4 — Noether's Theorem from the Janus Seal Topology: the Discrete Conservation Law, the Continuum Promotion, the Stress-Energy of a Sealed Composite, and the Canonical-Pair Obstructions}\\[1pt]
\texttt{BM{.}9} & theorem & The stress-energy tensor of a sealed composite: the two reads, their exact sign reconciliation, and the dust source & Part 4 & Theorem 4.9 \\
\multicolumn{5}{@{}l@{}}{\hspace*{1em}\small\itshape Tier BM.18 — The EM Self-Energy Correction: the Category Adjudication, the Running Coupling, the Path-Integral Self-Energy, the Branch-Blindness Theorem, and the Corrected Neutron–Proton Splitting}\\[1pt]
\texttt{BM{.}90} & theorem & The defect-conditioned weighting: the hypothesis on its face, and the branch-blindness theorem & Part 4 & Theorem 4.76 \\
\texttt{BM{.}91} & theorem & The corrected neutron–proton splitting: the attenuated assembly, the ensemble-independent bar, and the universal ratio bound & Part 4 & Theorem 4.77 \\
\texttt{BM{.}92} & theorem & The Q-value on both formed reads: the uniform closure, the corrected margins, and the untouched specified branch & Part 4 & Theorem 4.78 \\
\texttt{BM{.}93} & theorem & --- & Addendum & B.1.28 \\
\multicolumn{5}{@{}l@{}}{\hspace*{1em}\small\itshape Tier BM.19 — The Flicker Integral, the Rest-Mass Q-Value, and the Relating Equation: the Materialized Ensemble, the Vertex Superposition, the Rung-One \$\textbackslash{}nu\$ Line, and the Bridge in Four Exact Forms}\\[1pt]
\texttt{BM{.}94} & theorem & The flicker integral: the snapshot read's category error, the materialized distribution, and the two crystallization weights & Part 4 & Theorem 4.79 \\
\texttt{BM{.}95} & theorem & The vertex superposition: the transport fibre, its exact weights, and the two positive parts of the splitting & Part 4 & Theorem 4.80 \\
\texttt{BM{.}96} & theorem & The rung-one line: the seal-derived rest mass, the single-ray occupancy, and the four bars against a seesaw-type mechanism & Part 4 & Theorem 4.81 \\
\texttt{BM{.}97} & theorem & The rest mass is the pole mass: the species adjudication, the rest-mass Q-value, and the exact critical locus & Part 4 & Theorem 4.82 \\
\texttt{BM{.}98} & theorem & The relating equation: the bridge in four exact forms, its flicker weighting, and the closure in form & Part 4 & Theorem 4.83 \\
\texttt{BM{.}99} & theorem & The materialized splitting and the Q-value: the two-sided parameter-free bar, the tightened bound, and the rescoped dissolution row & Part 4 & Theorem 4.84 \\
\multicolumn{5}{@{}l@{}}{\hspace*{1em}\small\itshape Tier 4.0 — Composite Vocabulary}\\[1pt]
\texttt{Def-BM{.}1} & definition & Colour-sealed chain composites & Part 4 & Definition 4.1 \\
\multicolumn{5}{@{}l@{}}{\hspace*{1em}\small\itshape Tier BM.10 — The Kinetic Energy from the Metric and the Vacuum-Balanced Entropy: the Dispersion Read of a Sealed Composite, the Deficit Reframing of the Rest Mass, the Force-Structure Ladder, and the Route-Artifact Adjudication}\\[1pt]
\texttt{Def-BM{.}10} & definition & The metric-energy read of a sealed composite (the dispersion route) & Part 4 & Definition 4.10 \\
\multicolumn{5}{@{}l@{}}{\hspace*{1em}\small\itshape Tier BM.11 — The Neutron Decay Vertex: the Family Transport Profiles, the Rung-1 Isospin Row, the Forced Neutral Member, the Angular-Momentum Budget, and the Full Kinematic Balance}\\[1pt]
\texttt{Def-BM{.}11} & definition & The family transport-profile map (the constructed bridge) & Part 4 & Definition 4.11 \\
\multicolumn{5}{@{}l@{}}{\hspace*{1em}\small\itshape Tier BM.12 — The EM Self-Interaction Correction: the Scar's Field Energy, the Infinite-Range Seal Asymmetry, the Mass-Splitting Theorem, and the Beta-Spectrum Adjudication}\\[1pt]
\texttt{Def-BM{.}12} & definition & The chart-side EM source read of a sealed composite (the construction Def-BM.8(8) specifies) & Part 4 & Definition 4.12 \\
\multicolumn{5}{@{}l@{}}{\hspace*{1em}\small\itshape Tier BM.14 — The Constituent Transport Splitting: the Vertex Asymmetry, the Observation Frames, the Six-State Quark Assembly, and the Neutron Row}\\[1pt]
\texttt{Def-BM{.}13} & definition & The constituent transport refinement of the rest input (the construction) & Part 4 & Definition 4.13 \\
\texttt{Def-BM{.}14} & definition & The six-state assembly: the Platonic floors, the defect branches, and the generation map & Part 4 & Definition 4.14 \\
\multicolumn{5}{@{}l@{}}{\hspace*{1em}\small\itshape Tier BM.21 — The Antimatter and Resonance Tier: the CP-Conjugate, the Meson Seal, the Exchange Operator, and the Local Resonance Spectrum}\\[1pt]
\texttt{Def-BM{.}15} & definition & The CP-conjugate chain and the CP-pair & Part 4 & Definition 4.15 \\
\texttt{Def-BM{.}16} & definition & The two-chain state space (the construction BM.6(2) specifies) & Part 4 & Definition 4.16 \\
\texttt{Def-BM{.}17} & definition & The local resonance and its two threshold functions & Part 4 & Definition 4.17 \\
\multicolumn{5}{@{}l@{}}{\hspace*{1em}\small\itshape Tier BM.22 — The Annihilation Channel, Concrete Spin, and Fermion Identification: the Janus Tunnel, the Zig-Zag World Line, and the Spin-Dependent Transport Defect}\\[1pt]
\texttt{Def-BM{.}18} & definition & The Janus register and its crossing operators & Part 4 & Definition 4.18 \\
\texttt{Def-BM{.}19} & definition & The spin-dependent transport defect (the multiplet-ladder read) & Part 4 & Definition 4.19 \\
\multicolumn{5}{@{}l@{}}{\hspace*{1em}\small\itshape Tier 4.0 — Composite Vocabulary}\\[1pt]
\texttt{Def-BM{.}2} & definition & The composite card & Part 4 & Definition 4.2 \\
\multicolumn{5}{@{}l@{}}{\hspace*{1em}\small\itshape Tier BM.23 — The Closing Tier: the Haar–Poisson Duality and the Lapse, the Selected Curvature Composition, the Momentum-Space Resolvent, the Charge and Statistics Censuses, the Two-Chain Ensemble, the Seed-Tensoring Mechanism, the Convention Declarations, the Ceiling Study, the Jensen Residue, the Pair Amplitude, and the Era Closure}\\[1pt]
\texttt{Def-BM{.}20} & definition & The momentum-space -resolvent (the construction BM.6(6) specifies) & Part 4 & Definition 4.20 \\
\multicolumn{5}{@{}l@{}}{\hspace*{1em}\small\itshape Tier BM.4 — Noether's Theorem from the Janus Seal Topology: the Discrete Conservation Law, the Continuum Promotion, the Stress-Energy of a Sealed Composite, and the Canonical-Pair Obstructions}\\[1pt]
\texttt{Def-BM{.}3} & definition & The species-identification layer (declared) & Part 4 & Definition 4.3 \\
\multicolumn{5}{@{}l@{}}{\hspace*{1em}\small\itshape Tier BM.5 — The Calibration Lock: the Born Width, the Coherence Volume, the Curvature of a Sealed Composite, and the Neutron Row}\\[1pt]
\texttt{Def-BM{.}4} & definition & The Born width and the coherence volume & Part 4 & Definition 4.4 \\
\multicolumn{5}{@{}l@{}}{\hspace*{1em}\small\itshape Tier BM.6 — The Topology-Blind Neutron and the Topology-Blind + Static Proton: the Flicker Ensemble, the Path-Integrated Mass, the Binding Read, and the Forced Carrier}\\[1pt]
\texttt{Def-BM{.}5} & definition & The flicker ensemble of a sealed composite & Part 4 & Definition 4.5 \\
\texttt{Def-BM{.}6} & definition & The family naming rows (extension of the identification layer) & Part 4 & Definition 4.6 \\
\multicolumn{5}{@{}l@{}}{\hspace*{1em}\small\itshape Tier BM.7 — The Proton in the Positive CPT Sector and Its Mirror: the Sector Labels, the Exact Phase, and the Charge-Blind Geometry}\\[1pt]
\texttt{Def-BM{.}7} & definition & The EM sector labels and the EM scar & Part 4 & Definition 4.7 \\
\multicolumn{5}{@{}l@{}}{\hspace*{1em}\small\itshape Tier BM.8 — The Additive Charge Read and the Mass-Route Adjudication}\\[1pt]
\texttt{Def-BM{.}8} & definition & The additive charge read (the construction BM.6(1) specifies) & Part 4 & Definition 4.8 \\
\multicolumn{5}{@{}l@{}}{\hspace*{1em}\small\itshape Tier BM.9 — The Mesonic Flicker Test: the Contraction Grammar, the Colour Projector, and the Protected Junction}\\[1pt]
\texttt{Def-BM{.}9} & definition & The extended flicker ensemble and the sealing filter & Part 4 & Definition 4.9 \\
\addlinespace[4pt]
\multicolumn{5}{@{}l}{\textbf{Appendix A · Census — the computed data cells}}\\[2pt]
\multicolumn{5}{@{}l@{}}{\hspace*{1em}\small\itshape The cells}\\[1pt]
\texttt{FT{.}105} & data cell & Well-definedness: the topology census from first principles, finiteness, the regularization & Appendix A & Theorem A.22 \\
\texttt{FT{.}114} & data cell & The [EXT] audit clause (FT.73(5) pattern) & Appendix A & Theorem A.23 \\
\texttt{FT{.}120} & data cell & The [EXT] audit ledger: scheme columns, exact shortfalls, and the two exclusions & Appendix A & Theorem A.24 \\
\texttt{FT{.}28} & data cell & First computation: the floor-fermion particle card & Appendix A & Theorem A.1 \\
\texttt{FT{.}32} & data cell & The cycle floor family (single-contact ) & Appendix A & Theorem A.2 \\
\texttt{FT{.}35} & data cell & The composite floor : first colour-active, first -carrying & Appendix A & Theorem A.3 \\
\texttt{FT{.}42} & data cell & The Platonic card template & Appendix A & Theorem A.4 \\
\texttt{FT{.}43} & data cell & The winding ledger and Euler consistency & Appendix A & Theorem A.5 \\
\texttt{FT{.}44} & data cell & Duality signature and the floor taxonomy & Appendix A & Theorem A.6 \\
\texttt{FT{.}45} & data cell & The bidirectional polygon floors : cards, minima, and the null-carrying sector & Appendix A & Theorem A.7 \\
\texttt{FT{.}49} & data cell & The tetrahedron-face floor: full card & Appendix A & Theorem A.8 \\
\texttt{FT{.}50} & data cell & Composition and gluing laws & Appendix A & Theorem A.9 \\
\texttt{FT{.}51} & data cell & Enumeration of the winding ledger & Appendix A & Theorem A.10 \\
\texttt{FT{.}63} & data cell & The Platonic physical cards: spin, charge, parity, holonomy, entropy, mass ratios, binding & Appendix A & Theorem A.11 \\
\texttt{FT{.}64} & data cell & The filter chain and the per-solid loop lattices & Appendix A & Theorem A.12 \\
\texttt{FT{.}65} & data cell & Coarse-grained degeneracies, register product, and the master loop list & Appendix A & Theorem A.13 \\
\texttt{FT{.}69} & data cell & Platonic amplitude census completed; the cycle-alternation lemma; one census correction & Appendix A & Theorem A.14 \\
\texttt{FT{.}76} & data cell & The Platonic pole-mass multiplets; the bounded-hierarchy ceiling & Appendix A & Theorem A.15 \\
\texttt{FT{.}78} & data cell & The orbit decompositions: tetrahedron exhaustive, cube symmetry-reduced & Appendix A & Theorem A.16 \\
\texttt{FT{.}81} & data cell & Particle-content census (audit pass A) & Appendix A & Theorem A.17 \\
\texttt{FT{.}83} & data cell & The rung emergence ledger (audit pass C): derived, declared, open & Appendix A & Theorem A.18 \\
\texttt{FT{.}84} & data cell & --- & Addendum & B.4.3 \\
\texttt{FT{.}96} & data cell & The reversal-chirality census of the cover sector & Appendix A & Theorem A.19 \\
\texttt{FT{.}98} & data cell & The polygon form factors: exact census, saturating growth, fifth negative & Appendix A & Theorem A.20 \\
\texttt{FT{.}99} & data cell & The Platonic form factors, the dominance census, and the reduced-denominator record & Appendix A & Theorem A.21 \\
\addlinespace[4pt]
\multicolumn{5}{@{}l}{\textbf{Addendum B.1 · the adjudication record (HistLog §IV)}}\\[2pt]
\texttt{FT{.}115} & adjudication record & The adjudication summary & Addendum & B.1.1 \\
\texttt{FT{.}121} & adjudication record & The adjudication summary & Addendum & B.1.2 \\
\texttt{FT{.}132} & adjudication record & The adjudication summary & Addendum & B.1.3 \\
\texttt{FT{.}140} & adjudication record & The adjudication summary and the ledger & Addendum & B.1.4 \\
\texttt{FT{.}146} & adjudication record & The consolidation law: the letter series, the ownership partition, and the acyclic DAG & Addendum & B.1.6 \\
\texttt{FT{.}147} & adjudication record & The P-H audit and closure & Addendum & B.1.7 \\
\texttt{FT{.}148} & adjudication record & The structural/census separation of the working layer & Addendum & B.1.8 \\
\texttt{FT{.}149} & adjudication record & The framing retirement and the reframed question ledger & Addendum & B.1.9 \\
\texttt{FT{.}150} & adjudication record & The gap-closure audit & Addendum & B.1.10 \\
\texttt{FT{.}151} & adjudication record & Baryonic preparation & Addendum & B.1.11 \\
\texttt{FT{.}152} & adjudication record & The status of the fermionic sector after consolidation & Addendum & B.1.12 \\
\addlinespace[4pt]
\bottomrule
\end{longtable}
\endgroup

\chapter*{Colophon}
\addcontentsline{toc}{chapter}{Colophon}
This volume was generated from the canon files preserved in
\texttt{working\_markdowns/} by the packaging pipeline in \texttt{source/pipeline/}.
Registry: 623 tags; Req graph: 8014
edges, acyclic (re-derived at every suite run). Verification suite: 80/80 green
(\texttt{verification/run\_all.py}); exact arithmetic only --- no floating
point in any assertion. Edition of 2026-10-09. The archive,
the checkers and the mathematics are unmodified; the transformations applied
to produce this presentation layer are logged in
	exttt{source/CONVERSION\_LOG.md}.
\end{document}